\chapter{Masses, covolumes and an arithmetic volume estimate}\label{chapter-covolumes}
The comparison inequality of Chapter~\ref{chapter-inert-prime-balance} needs an upper
bound for the covolume of the lattice of $j$-th lifted forms on each indefinite unitary
group $U(W_j)$ of the chain. This chapter proves that bound, with main term
$\sum_\rho(p_j-q_j)e_j(\rho)\log\Omega$ and an error of size about
$e^{x^{\varepsilon}}$; masses enter it through the volumes of the level subgroups of the
$U(W_j)$. The chapter also proves the counting results on the definite side that use
Siegel's mass formula, and states the two forms of the Chowla--Selberg formula, one used
on each side.

\section{Normalized covolumes and slopes}\label{section-cv-slopes}

\begin{definition}\label{cv-hermitian-norms}
Let $K$ be a totally complex number field with $r=[K:\mathbf Q]/2$ infinite places, and
let $V$ be a finite-dimensional $K$-vector space. A \emph{Hermitian norm system} on $V$ is
a family $N=(N_\sigma)_{\sigma:K\to\mathbf C}$ of functions $N_\sigma:V\to\mathbf R$ for
which there are a basis $b_1,\dots,b_n$ of $V$ and invertible matrices
$X_v\in GL_n(\mathbf C)$, one for each infinite place $v$ with chosen embedding $\sigma_v$,
such that $N_{\sigma_v}(x)=|X_v^{-1}\sigma_v(x_b)|$ for all $x\in V$. Here $x_b\in K^n$ is
the coordinate vector of $x$ and $|\cdot|$ is the standard Hermitian norm on
$\mathbf C^n$. The \emph{realization} of $V$ is the map
$V\to\bigoplus_v\mathbf C^n\cong\mathbf R^{2nr}$,
$x\mapsto(X_v^{-1}\sigma_v(x_b))_v$, with the standard Euclidean structure.

Let $L\subset V$ be a nonzero finitely generated $\mathcal O_K$-submodule, and let
$\operatorname{rk}L=\dim_K KL$. Let $\operatorname{covol}(L)$ be the square root of the
Gram determinant of a $\mathbf Z$-basis of the image of $L$ in the realization, and $0$ if
that image is not discrete. Define
\[
 \mu(L)=-\frac{2\log\operatorname{covol}(L)}{[K:\mathbf Q]\operatorname{rk}L}-\log2
 +\frac{\log|d_K|}{[K:\mathbf Q]},\qquad
 \mu_{\max}(L)=\sup\{\mu(L'):0\ne L'\subset L\},
\]
where $L'$ runs over the $\mathcal O_K$-submodules of $L$. These do not depend on the
choice of $b$ and $X$. A \emph{metric} on $K^d$ is a family $X=(X_v)_v$ of matrices
$X_v\in GL_d(\mathbf C)$, one for each infinite place $v$; it gives the Hermitian norm
system $N_{\sigma_v}(c)=|X_v^{-1}\sigma_v(c)|$ on $K^d$ with its standard basis. For a
metric $X$, the \emph{normalized height} and the \emph{normalized log-covolume} of
$L\subset K^d$ are
\[
 \operatorname{ht}_X(L)=2^{r\operatorname{rk}L}\,|d_K|^{-\operatorname{rk}L/2}
 \operatorname{covol}(L),\qquad
 \operatorname{lcv}_X(L)=\frac{2}{[K:\mathbf Q]}\log\operatorname{ht}_X(L).
\]
If the image of $L$ is discrete, then $\operatorname{lcv}_X(L)=-\operatorname{rk}L\cdot\mu(L)$.
\end{definition}
\leanref{Development/AbelianVarietyTheory/Definitions/Foundations/IsProbabilityMeasureMuPart.lean}{GlobalTheory.AllWt.HasGram}
\leanref{Development/AbelianVarietyTheory/Definitions/Foundations/IsProbabilityMeasureMuPart.lean}{GlobalTheory.AllWt.slope_sq}
\leanref{Development/AbelianVarietyTheory/Definitions/Foundations/IsProbabilityMeasureMuPart.lean}{GlobalTheory.AllWt.slopeMax_sq}
\leanref{Formalism/HarderNarasimhan/StatementPl.lean}{Formalism.HN.muDef}
\leanref{Formalism/HarderNarasimhan/StatementPl.lean}{Formalism.HN.ht}
\leanref{Formalism/QuotientMetricDivisors/Statement.lean}{Formalism.DivisionValues.lcv}

\begin{lemma}\label{cv-rescaling}
Let $K$ be a totally complex number field, let $N$ be a Hermitian norm system on a
$K$-vector space $V$ (Definition~\ref{cv-hermitian-norms}), and let $L\subset V$ be an
$\mathcal O_K$-submodule. Let $b_1,\dots,b_d\in V$ be linearly independent over $K$ with
$L\subset\sum_kKb_k$, and let $\beta:K^d\to V$ be $c\mapsto\sum_kc_kb_k$. Let $X=(X_v)_v$
be a metric on $K^d$ (Definition~\ref{cv-hermitian-norms}) and let $0<s\le1$ be such that
$|X_v^{-1}\sigma_v(c)|=s\,N_{\sigma_v}(\beta(c))$ for every infinite place $v$ and every
$c\in K^d$. If $\beta^{-1}(L)$ is a lattice in $K^d$, then
\[
 \frac1d\operatorname{lcv}_X\big(\beta^{-1}(L)\big)\le-\mu(L).
\]
\end{lemma}
\begin{proof}
The map $\beta$ identifies $\beta^{-1}(L)$ with $L$, and the realization of $K^d$ for $X$
is $s$ times the realization of $V$ for $N$, restricted to $\sum_kKb_k$. Hence the
covolume of $\beta^{-1}(L)$ is $s^{[K:\mathbf Q]d}\operatorname{covol}(L)$. By the last
assertion of Definition~\ref{cv-hermitian-norms},
$\frac1d\operatorname{lcv}_X(\beta^{-1}(L))=-\mu(L)+2\log s$. Since $s\le1$, the claim
follows.
\end{proof}
\leanref{Development/AbelianVarietyTheory/Proofs/DivisorProofs/CovolumePerRankLeNegSlopeScaledStatement.lean}{GlobalTheory.T.covolumePerRank_le_neg_slope_scaled}

\begin{lemma}\label{cv-slope-base-change}
Let $K\subset K'$ be totally complex number fields. Let $V$ be a $K$-vector space and $V'$
a $K'$-vector space with Hermitian norm systems $N$ and $N'$
(Definition~\ref{cv-hermitian-norms}), and suppose that
$N_{\bar\sigma}=N_\sigma$ for every embedding $\sigma:K\to\mathbf C$, where
$\bar\sigma$ is its complex conjugate. Let $\iota:V\to V'$ be an additive map with
$\iota(av)=a\,\iota(v)$ for $a\in K$ and $v\in V$, which carries every $K$-linearly
independent family to a $K'$-linearly independent family. Let $L\subset V$ be a nonzero
finitely generated $\mathcal O_K$-submodule, let $L'\subset V'$ be the
$\mathcal O_{K'}$-submodule generated by $\iota(L)$, and let $B\ge0$.
\begin{enumerate}
\item If $N'_{\sigma'}(\iota(x))\le e^{B}N_{\sigma'|K}(x)$ for every embedding
  $\sigma':K'\to\mathbf C$ and every $x\in V$, then $\mu(L)-2B\le\mu(L')$.
\item If there are numbers $c_{\sigma'}>0$ with $|\log c_{\sigma'}|\le B$ and
  $N'_{\sigma'}(\iota(x))=c_{\sigma'}N_{\sigma'|K}(x)$ for every $\sigma'$ and every
  $x\in V$, then $\mu_{\max}(L')\le\mu_{\max}(L)+2B$.
\end{enumerate}
\end{lemma}
\leanref{Development/AbelianVarietyTheory/Proofs/DivisorProofs/SlopeSqSubLeSlopeSqSpan.lean}{GlobalTheory.SlopeExt.slope_sq_sub_le_slope_sq_span}
\leanref{Development/AbelianVarietyTheory/Proofs/DivisorProofs/MaxSlopeSqFieldExtensionLeComparableNormsStatement.lean}{GlobalTheory.T.maxSlope_sq_fieldExtension_le_comparable_norms}
\begin{proof}
First a formula for the slope. Let $M\subset V$ be a nonzero finitely generated
$\mathcal O_K$-submodule of rank $j$, let $e_1,\dots,e_j\in M$ be a basis of $KM$, and
let $M_0=\sum_i\mathcal O_Ke_i$, a subgroup of finite index in $M$. For an embedding
$\sigma$ let $G_\sigma(e)$ be the Gram matrix of $e$ for the Hermitian form of $N_\sigma$.
The image of $\mathcal O_K^j$ in $\bigoplus_v\mathbf C^j$ has covolume
$(2^{-r}|d_K|^{1/2})^j$, where $r=[K:\mathbf Q]/2$, and the map
$c\mapsto\sum_ic_ie_i$ multiplies volumes by $\det G_{\sigma_v}(e)$ at the place $v$.
Hence $\operatorname{covol}(M)=[M:M_0]^{-1}(2^{-r}|d_K|^{1/2})^j\prod_v\det G_{\sigma_v}(e)$.
Since $N_{\bar\sigma}=N_\sigma$, the determinants at $\sigma_v$ and $\bar\sigma_v$ agree, and
Definition~\ref{cv-hermitian-norms} gives
\[
 j\,\mu(M)=\frac{1}{[K:\mathbf Q]}\Big(2\log[M:M_0]-\sum_{\sigma:K\to\mathbf C}\log\det G_\sigma(e)\Big).
\]

Proof of (1). Let $j=\operatorname{rk}L$ and choose $e$ and $L_0$ for $L$ as above. The
vectors $\iota(e_i)$ are linearly independent over $K'$ and lie in $L'$, and
$L'_0=\sum_i\mathcal O_{K'}\iota(e_i)$. The map $\mathcal O_{K'}\otimes_{\mathcal O_K}L\to V'$
induced by $\iota$ is injective with image $L'$: the module $\mathcal O_{K'}\otimes L$ embeds
in $K'\otimes_KKL$, which $\iota$ maps injectively since it preserves linear independence. So
$\operatorname{rk}L'=j$, and $L'/L'_0\cong\mathcal O_{K'}\otimes_{\mathcal O_K}(L/L_0)$, which
has $[L:L_0]^{[K':K]}$ elements because $\mathcal O_{K'}$ is locally free of rank $[K':K]$
over $\mathcal O_K$. Fix $\sigma'$ and put $\sigma=\sigma'|K$. The image $\sigma(K)$ is dense
in $\mathbf C$, so the hypothesis extends by continuity: the $\mathbf C$-linear map
$V\otimes_{K,\sigma}\mathbf C\to V'\otimes_{K',\sigma'}\mathbf C$ induced by $\iota$ has
operator norm at most $e^B$. Hence $G_{\sigma'}(\iota e)\le e^{2B}G_\sigma(e)$ as Hermitian
matrices, and $\log\det G_{\sigma'}(\iota e)\le2Bj+\log\det G_\sigma(e)$. Each $\sigma$
has $[K':K]$ extensions $\sigma'$. The formula above, for $L'$ and for $L$, gives
\[
 j\,\mu(L')\ge\frac{2[K':K]\log[L:L_0]-2Bj[K':\mathbf Q]-[K':K]\sum_\sigma\log\det G_\sigma(e)}{[K':\mathbf Q]}
 =j\,\mu(L)-2Bj.
\]

Proof of (2). Let $F$ be the normal closure of $K'$ over $K$. Choose a sublattice
$M\subset L'$ with $\mu(M)=\mu_{\max}(L')$; it exists since the slopes of sublattices of
$L'$ are bounded above and take a discrete set of values. Write $V'$ in coordinates
$K'^n$ for a basis, let $\varphi:V'\to F^n$ be the extension of scalars, and give $F^n$ the
norm system $N''$ with $N''_\theta(\varphi(y))=c_{\theta|K'}^{-1}N'_{\theta|K'}(y)$ for
$\theta:F\to\mathbf C$. Then $N''_\theta(\varphi\iota(x))=N_{\theta|K}(x)$ for $x\in V$, and
$N''_\theta(\varphi(y))\le e^BN'_{\theta|K'}(y)$. By (1) for $K'\subset F$ and $\varphi$,
$\mu(M)-2B\le\mu(M'')$, where $M''$ is the $\mathcal O_F$-span of $\varphi(M)$. This span lies
in the $\mathcal O_F$-span $L''$ of $\varphi\iota(L)$, so $\mu(M'')\le\mu_{\max}(L'')$. The
lattice $L''$ is the extension of $L$ to the Galois extension $F$ of $K$ with the extended
norms, and $\mu_{\max}(L'')\le\mu_{\max}(L)$. Indeed, the sublattice $M_1$ of $L''$ of largest rank
among those of maximal slope is unique, so it is stable under $\operatorname{Gal}(F/K)$, acting on the
coordinates; hence its $F$-span is defined over $K$, and $M_1$ is the extension of the sublattice
$M_0$ of $L$ cut out by that span. Since the norms of $L''$ restrict to those of $L$, the Gram
determinants in the proof of (1), with $B=0$, are equal, so $\mu(M_1)=\mu(M_0)\le\mu_{\max}(L)$.
Together, $\mu_{\max}(L')=\mu(M)\le\mu_{\max}(L)+2B$.
\end{proof}
\leanref{Development/AbelianVarietyTheory/Proofs/DivisorProofs/MaxSlopeSqFieldExtensionLeComparableNorms.lean}{GlobalTheory.SlopeMaxExt.slopeMax_sq_span_le}

\begin{lemma}\label{cv-cut}
Let $K$ be a totally complex number field, let $N$ be a Hermitian norm system on a
$K$-vector space $V$ (Definition~\ref{cv-hermitian-norms}), and let $L\subset M\subset V$
be $\mathcal O_K$-submodules with $M$ finitely generated and $L\ne0$. Let $e:V\to V$ be a
$K$-linear idempotent that is orthogonal for every $N_\sigma$, that is,
$N_\sigma(e(x)+y-e(y))^2=N_\sigma(e(x))^2+N_\sigma(y-e(y))^2$ for all $\sigma$, $x$ and
$y$. Let $g\ge1$ be an integer with $g\,e(L)\subset L$, and let $\delta$ and $Q$ be real
numbers with $0\le\delta\le1/12$,
\[
 \operatorname{rk}(L\cap\ker e)\le3\delta\operatorname{rk}L,\qquad
 \mu_{\max}(M)\le Q,\qquad -Q\le\mu(L),\qquad 2\log g\le Q.
\]
Then $\mu(L)-12\delta Q\le\mu(L\cap\operatorname{im}e)$.
\end{lemma}
\begin{proof}
Write $L_g=L\cap\operatorname{im}e$ and $L_b=L\cap\ker e$, with ranks $d_g$ and $d_b$. An
orthogonal idempotent is self-adjoint for the Hermitian forms of $N$, so $1-e$ is the orthogonal
projection onto $\ker e$ and has operator norm at most $1$ at every place. For $x\in L$ both
$g\,e(x)$ and $g(x-e(x))$ lie in $L$, so $gL\subset L_g\oplus L_b\subset L$ and
$d_g+d_b=\operatorname{rk}L$. Let $\varphi=g(1-e)$ and $I=\varphi(L)$. Then $I\subset L_b$, and
$I$ has rank $d_b$, since the kernel of $\varphi$ on $KL$ is $KL\cap\operatorname{im}e=KL_g$. A
sublattice of the same rank has covolume at least that of the lattice, so
$\operatorname{covol}(L_b)\le\operatorname{covol}(I)$. The map $\varphi$ vanishes on $KL_g$, so $I$ is
the image under $\varphi$ of the lattice $L/L_g$, realized in the orthogonal complement of $KL_g$
with the quotient norms; as $L_g$ is saturated in $L$, its covolume is
$\operatorname{covol}(L)/\operatorname{covol}(L_g)$. The operator norm of $\varphi$ is at most $g$ at
each of the $[K:\mathbf Q]/2$ infinite places, and $\varphi$ is injective on the $K$-span of
$L/L_g$, of rank $d_b$; so
$\operatorname{covol}(I)\le g^{[K:\mathbf Q]d_b}\operatorname{covol}(L)/\operatorname{covol}(L_g)$.
Taking logarithms gives
\[
 \operatorname{rk}L\cdot\mu(L)\le d_g\,\mu(L_g)+d_b\big(\mu(L_b)+2\log g\big).
\]
Also $\mu(L_b)\le\mu_{\max}(M)\le Q$, because $L_b\subset M$. If $d_g=0$ then
$d_b=\operatorname{rk}L>3\delta\operatorname{rk}L$, against the hypothesis; so $d_g>0$.
The displayed inequality gives
$d_g(\mu(L_g)-\mu(L))\ge d_b(\mu(L)-\mu(L_b)-2\log g)\ge-3Qd_b$. Since
$d_b\le3\delta\operatorname{rk}L$ and $d_g\ge(1-3\delta)\operatorname{rk}L\ge\frac34\operatorname{rk}L$,
this gives $\mu(L_g)-\mu(L)\ge-12\delta Q$.
\end{proof}
\leanref{Development/AbelianVarietyTheory/Proofs/DivisorProofs/MajAtAddRight.lean}{GlobalTheory.CutNum.slope_good_ge}
\leanref{Formalism/MeanHeightDescent/MuMaxMinkowski.lean}{Formalism.ExtensionFiltration.GoodBadCut}
\leanref{Formalism/MeanHeightDescent/AverageDegreeBounds/GoodBadCut.lean}{Formalism.ExtensionFiltration.good_bad_split}

\begin{remark}\label{cv-cut-index}
The displayed inequality in the proof of Lemma~\ref{cv-cut} also follows from an index count. The
lattices $L_g$ and $L_b$ are orthogonal, so
$\operatorname{covol}(L_g\oplus L_b)=\operatorname{covol}(L_g)\operatorname{covol}(L_b)$. The quotient
$L/(L_g\oplus L_b)$ embeds, through $x\mapsto x-e(x)$, into $(1-e)(L)/L_b\subset g^{-1}L_b/L_b$, so its
order is at most $g^{[K:\mathbf Q]d_b}$, and
$\operatorname{covol}(L_g)\operatorname{covol}(L_b)\le g^{[K:\mathbf Q]d_b}\operatorname{covol}(L)$.
\end{remark}

\begin{lemma}\label{cv-slope-descent}
Let $K$ be a totally complex number field, let $V$ and $W$ be finite-dimensional
$K$-vector spaces with Hermitian norm systems $N^V$ and $N^W$
(Definition~\ref{cv-hermitian-norms}), and let $A\subset V$ and $B\subset W$ be finitely
generated $\mathcal O_K$-submodules with $A\ne0$. Let $\varphi_1,\dots,\varphi_n:V\to W$
be $K$-linear maps with $\varphi_h(A)\subset B$ for every $h$, such that $f=0$ whenever
$\varphi_h(f)=0$ for all $h$, and such that
\[
 \sum_{h=1}^nN^W_\sigma(\varphi_h(f))^2\le e^{2P}N^V_\sigma(f)^2
 \qquad(\sigma:K\to\mathbf C,\ f\in V)
\]
for a real number $P$. Suppose that $B$ is contained in the $K$-span of a set of vectors
$y\in B$ with $N^W_\sigma(y)\le e^{\Lambda}$ for all $\sigma$, where $\Lambda\ge0$. Let
$\varepsilon\ge0$ and $t$ be real numbers with
$n\operatorname{rk}B\le(1+\varepsilon)\operatorname{rk}A$ and $t\le\mu(A)-2P$. Then
\[
 t-\varepsilon\big(|t|+2\Lambda\big)\le\mu(B).
\]
\end{lemma}
\begin{proof}
Let $B^n\subset W^n$ be the product lattice, with the orthogonal sum of the norms, and let
$\Phi=(\varphi_h)_h:V\to W^n$. The map $\Phi$ is injective and sends $A$ into $B^n$. The
norm hypothesis says that $\Phi$ has operator norm at most $e^P$ at every place, so its
height, the mean over the infinite places of the logarithm of the squared operator norm,
is at most $2P$. The lattice $B^n$ is spanned over $K$ by the vectors of $B$ placed in one
coordinate, which have norm at most $e^\Lambda$, so its minimal slope is at least
$-2\Lambda$. The slope inequality for an injective map $\Phi$ with $\Phi(A)\subset B^n$
(Theorem~\ref{hs-bost-slope-inequality}, in the form used here) gives
\[
 a\big(\mu(A)-2P\big)+(a'-a)(-2\Lambda)\le a'\mu(B^n),
\]
where $a=\operatorname{rk}A$ and $a'=\operatorname{rk}B^n=n\operatorname{rk}B$. Finally
$\mu(B^n)=\mu(B)$, and $a\le a'\le(1+\varepsilon)a$. Using $t\le\mu(A)-2P$ and
$a'-a\le\varepsilon a'$, one gets $a't-\varepsilon a'(|t|+2\Lambda)\le a'\mu(B)$.
\end{proof}
\leanref{Development/AbelianVarietyTheory/Proofs/DivisorProofs/HtMapLeOfPoints.lean}{GlobalTheory.DescentFamily.descent}
\leanref{Formalism/MeanHeightDescent/MuMaxMinkowski.lean}{Formalism.ExtensionFiltration.DescentSlope}

\begin{lemma}\label{cv-minkowski-slope}
Let $K$ be a totally complex number field, let $N$ be a Hermitian norm system on a $K$-vector space
$V$ (Definition~\ref{cv-hermitian-norms}) with $N_{\bar\sigma}=N_\sigma$ for every embedding
$\sigma:K\to\mathbf C$, let $L\subset V$ be a nonzero finitely generated $\mathcal O_K$-submodule, and
let $\lambda$ be a real number such that
\[
 \frac1{[K:\mathbf Q]}\sum_{\sigma:K\to\mathbf C}\log N_\sigma(x)\ge-\lambda
\]
for every nonzero $x\in L$. Then
$\mu_{\max}(L)\le2\lambda+\log\operatorname{rk}L+\log|d_K|/[K:\mathbf Q]$.
\end{lemma}
\leanref{Formalism/MeanHeightDescent/MuMaxMinkowski.lean}{Formalism.ExtensionFiltration.MuMaxMinkowski}
\leanref{Formalism/MeanHeightDescent/DivisionInequalityWithCovolume.lean}{Formalism.ExtensionFiltration.muMaxMinkowski}
\leanref{Development/AbelianVarietyTheory/Proofs/DivisorProofs/MaxSlopeSqLeOfNoShortVector.lean}{GlobalTheory.SlopeMink.slopeMax_sq_le_of_no_short}
\begin{proof}
Let $r=[K:\mathbf Q]/2$, and let $x\in L$ be nonzero. Its image in the realization has squared
length $\sum_vN_{\sigma_v}(x)^2$, the sum over the $r$ infinite places. Since $N_{\bar\sigma}=N_\sigma$,
the hypothesis says $\frac1r\sum_v\log N_{\sigma_v}(x)^2\ge-2\lambda$, and by the inequality between
arithmetic and geometric means
\[
 \log\Big(\frac1r\sum_vN_{\sigma_v}(x)^2\Big)\ge\frac1r\sum_v\log N_{\sigma_v}(x)^2\ge-2\lambda .
\]
So every nonzero vector of the image of $L$ has length at least $\rho=\sqrt r\,e^{-\lambda}$. Now let
$L'\subset L$ be a nonzero $\mathcal O_K$-submodule of rank $j$, with covolume $c$. Its image spans a
real space of dimension $m=2jr$, in which it is a lattice of covolume $c$. The ball of radius
$\sqrt m\,c^{1/m}$ contains a cube of side $2c^{1/m}$, of volume $2^mc$, so by Minkowski's theorem $L'$
has a nonzero element of length at most $\sqrt m\,c^{1/m}$. Hence $\rho^2\le mc^{2/m}$, that is,
$-2\lambda\le\log(2j)+\frac{\log c}{jr}$. By Definition~\ref{cv-hermitian-norms},
$\frac{\log c}{jr}=-\mu(L')-\log2+\frac{\log|d_K|}{[K:\mathbf Q]}$. So
$\mu(L')\le2\lambda+\log j+\log|d_K|/[K:\mathbf Q]$. As $j\le\operatorname{rk}L$, taking the supremum
over $L'$ gives the claim.
\end{proof}

\section{Masses and class numbers}\label{section-cv-masses}

\begin{definition}\label{cv-mass}
Let $K$ be a CM field with complex conjugation $c$, let $V$ be a finite-dimensional
$K$-vector space with a Hermitian form $h$, and let $U(h)$ be the group of $K$-linear
automorphisms of $V$ preserving $h$. Let $L\subset V$ be an $\mathcal O_K$-lattice. A
lattice $M\subset V$ lies in the \emph{genus} of $L$ if for every maximal ideal
$\mathfrak p$ of $\mathcal O_K$ there is $g\in U(h)$ with $g(L_{\mathfrak p})=M_{\mathfrak p}$,
where the subscript denotes localization at the elements $s\notin\mathfrak p$ with
$c(s)\notin\mathfrak p$. The \emph{mass} of $L$ is
\[
 \operatorname{mass}(L)=\sum_{M}\frac1{|\operatorname{Aut}(M)|},
\]
where $M$ runs over a set of representatives of the $U(h)$-orbits in the genus of $L$ and
$\operatorname{Aut}(M)$ is the stabilizer of $M$ in $U(h)$. For $M\ge1$ let $I_M$ denote
$\mathcal O_K^M\subset K^M$ with the form $h(x,y)=\sum_ic(x_i)y_i$.
\end{definition}
\leanref{Formalism/HermitianLattices/StatementMass.lean}{Formalism.HermitianLatticesTheory.mass}
\leanref{Formalism/HermitianLattices/StatementLat.lean}{Formalism.HermitianLatticesTheory.genus}
\leanref{Formalism/SharpRestriction/GlobalSchurBounds/GeneralLinearActions/GapAtLeast.lean}{Formalism.MainNotions.StandardLattice.massI}

\begin{theorem}\label{cv-siegel-mass-ratio}
Let $K$ be a CM field with totally real subfield $F$ and complex conjugation $c$, and
suppose that for every maximal ideal $\mathfrak P$ of $\mathcal O_K$ there is
$\theta\in\mathcal O_K$ with $\theta-c(\theta)\notin\mathfrak P$. Put
$v_K=|d_K|^{1/(2[F:\mathbf Q])}/2$, half the root discriminant of $K$. Let $M\ge3$. Then
(Definition~\ref{cv-mass})
\[
 \frac{\operatorname{mass}(I_M)}{\operatorname{mass}(I_{M-1})}
 =\Big((M-1)!\,\Big(\frac{v_K}{\pi}\Big)^M\Big)^{[F:\mathbf Q]}
 (2v_K)^{-[F:\mathbf Q]/2}\,\Lambda_K(M),
\]
where $\Lambda_K(M)=\prod_{\mathfrak P}e_{\mathfrak P}$ over the maximal ideals of
$\mathcal O_K$, with $e_{\mathfrak P}=(1-(-1)^MN\mathfrak P^{-M/2})^{-1}$ if
$c(\mathfrak P)=\mathfrak P$ and $e_{\mathfrak P}=(1-N\mathfrak P^{-M})^{-1/2}$
otherwise. Thus $\Lambda_K(M)=L_F(M,\eta^M)$ for the quadratic character $\eta$ of $K/F$.
\end{theorem}
\leanref{Development/AbelianVarietyTheory/Proofs/GlobalInterpolation/MassRatioAllUnramifiedCmFields.lean}{GlobalTheory.Mass.massRatio_all_unramified_cmFields}
\leanref{Formalism/SharpRestriction/GlobalSchurBounds/StandardLattice.lean}{Formalism.MainNotions.StandardLattice.SiegelMassRatioCM}
\leanref{Formalism/HermitianLattices/StatementMass.lean}{Formalism.HermitianLatticesTheory.massRatioRightSide}

The hypothesis on $\theta$ says that $K/F$ is unramified at every finite prime, so that
$I_M$ and $I_{M-1}$ are unimodular at every place. The formula is the ratio of the mass
formulas for $I_M$ and $I_{M-1}$ of Gan, Hanke and Yu \cite{GHY01}; see also
\cite{Shi97}. In the form of Gan and Yu \cite[Theorem~10.20]{GY00}, with the Tamagawa number $2$ of a
unitary group and the local densities of unimodular lattices at the places unramified in $K$, the mass
formula reads
\[
 \operatorname{mass}(I_n)=2\,d_F^{\,n^2/2}\prod_{i=1}^n\Big(\frac{(i-1)!}{(2\pi)^i}\Big)^{[F:\mathbf Q]}L_F(i,\eta^i)
\]
when $K/F$ is unramified at every finite prime, with $\eta^i$ trivial for even $i$; as $|d_K|=d_F^2$, the
ratio for $n=M$ and $n=M-1$ is the right side above.

\begin{lemma}\label{cv-transference}
Let $\Gamma\subset\mathbf R^N$ be a lattice of rank $N\ge1$ and let
$\Gamma^\vee=\{y\in\mathbf R^N:\langle y,x\rangle\in\mathbf Z\text{ for all }x\in\Gamma\}$ be its
dual lattice. Let $\lambda>0$, and suppose that every nonzero vector of $\Gamma$ has length at
least $\lambda$. Then $\Gamma^\vee$ contains $N$ linearly independent vectors of length at most
$N^2/\lambda$.
\end{lemma}
\leanref{ReciprocityTheory/ReductiveGroups/ArithmeticGeometry/HeightTheory/Transference.lean}{ReciprocityTheory.ChartRestriction.G.T0.transference_dual}

This is a weak form of the transference theorems of the geometry of numbers: Banaszczyk proves
that the $N$-th successive minimum of $\Gamma^\vee$ is at most $N$ divided by the first successive
minimum of $\Gamma$ \cite{Ban93}.

\begin{lemma}\label{cv-short-basis}
Let $E$ be a CM field with complex conjugation $c$, let $n=[E:\mathbf Q]$, $f=n/2$ and
$\operatorname{rd}_E=|d_E|^{1/n}$, and let $M\ge1$. Let $h(x,y)=\sum_ic(x_i)y_i$ on $E^M$, and let
$\Lambda\subset E^M$ be a finitely generated $\mathcal O_E$-submodule that is \emph{self-dual}:
$y\in\Lambda$ if and only if $h(x,y)\in\mathcal O_E$ for every $x\in\Lambda$. Then there are
$v_1,\dots,v_M\in\Lambda$, linearly independent over $E$, with
\[
 |\sigma(h(v_i,v_k))|\le4f^3M^4\operatorname{rd}_E^2
\]
for every embedding $\sigma:E\to\mathbf C$ and all $i,k$.
\end{lemma}
\leanref{Development/AbelianVarietyTheory/Proofs/FoundationProofs/ExistsFinsetSelfDualClassesRootDiscr.lean}{GlobalTheory.RootClassCount.exists_bounded_basis}
\leanref{Development/AbelianVarietyTheory/Proofs/CuspFormEvaluations/TraceHermNondegenerate.lean}{GlobalTheory.exists_short_vectors}
\leanref{Development/AbelianVarietyTheory/Proofs/CuspFormEvaluations/EmbeddingHermSelf.lean}{GlobalTheory.trace_dual_min}
\begin{proof}
Let $T(x,y)=\operatorname{Tr}_{E/\mathbf Q}h(x,y)$, a symmetric $\mathbf Q$-bilinear form on
$V=E^M$, a $\mathbf Q$-vector space of dimension $N=nM$. Since $\sigma\circ c=\overline\sigma$ for
every embedding $\sigma$,
\[
 T(x,x)=\sum_{\sigma:E\to\mathbf C}\sigma(h(x,x))=\sum_\sigma\sum_k|\sigma(x_k)|^2,
\]
so $T$ is positive definite. Since $\Lambda$ is self-dual it spans $V$, and it is a lattice of
rank $N$ in $V\otimes\mathbf R$ with the Euclidean structure $T$. Its dual for $T$ is
$\Lambda^T=\{u\in V:T(u,y)\in\mathbf Z\text{ for all }y\in\Lambda\}$, and the dual of $\Lambda^T$ is
$\Lambda$.

Step 1: every nonzero $u\in\Lambda^T$ has $T(u,u)\ge n/\operatorname{rd}_E^2$. Let $\mathfrak d$ be
the different of $E$, so that $\mathfrak d^{-1}=\{z\in E:\operatorname{Tr}(az)\in\mathbf Z\text{ for
all }a\in\mathcal O_E\}$. For $y\in\Lambda$ and $a\in\mathcal O_E$ the vector $ay$ lies in
$\Lambda$, and $\operatorname{Tr}(a\,h(u,y))=T(u,ay)\in\mathbf Z$, so $h(u,y)\in\mathfrak d^{-1}$.
Let $\delta\in\mathfrak d$. Since $h(x,u)=c(h(u,x))$ and $\mathfrak d$ is stable under $c$, we get
$h(x,\delta u)=\delta\,c(h(u,x))\in\mathcal O_E$ for every $x\in\Lambda$, so $\delta u\in\Lambda$ by
self-duality. Then $\delta h(u,u)=h(u,\delta u)\in\mathfrak d^{-1}$ for every $\delta\in\mathfrak d$, so
$h(u,u)\in\mathfrak d^{-2}$. Now $h(u,u)$ is nonzero, since its
conjugates are the positive numbers $\sum_k|\sigma(u_k)|^2$, so
$|N_{E/\mathbf Q}h(u,u)|\ge N(\mathfrak d)^{-2}=|d_E|^{-2}$. By the inequality between arithmetic
and geometric means,
$T(u,u)=\sum_\sigma\sigma(h(u,u))\ge n|N_{E/\mathbf Q}h(u,u)|^{1/n}\ge n|d_E|^{-2/n}=n/\operatorname{rd}_E^2$.

Step 2: short vectors. By Step 1 and Lemma~\ref{cv-transference}, applied to $\Lambda^T$ with
$\lambda=\sqrt n/\operatorname{rd}_E$, the lattice $\Lambda$ contains $N$ vectors $w_1,\dots,w_N$,
linearly independent over $\mathbf Q$, with $T(w_i,w_i)\le N^4\operatorname{rd}_E^2/n$.

Step 3: the bound at one embedding. Let $\tau$ be an embedding. The embeddings $\tau$ and
$\overline\tau$ are distinct, since $E$ is totally complex, and contribute the same term
$\sum_k|\tau(w_k)|^2$ to $T(w,w)$. Hence
$2\sum_k|\tau(w_{i,k})|^2\le T(w_i,w_i)\le N^4\operatorname{rd}_E^2/n$, that is,
$\sum_k|\tau(w_{i,k})|^2\le n^3M^4\operatorname{rd}_E^2/2=4f^3M^4\operatorname{rd}_E^2$.

Step 4: a basis. The $w_i$ span $V$ over $\mathbf Q$, so they span $E^M$ over $E$, and $M$ of them,
$v_1,\dots,v_M$, are linearly independent over $E$. Finally, for every $\sigma$,
\[
 |\sigma(h(v_i,v_k))|\le\sum_j|\sigma(v_{i,j})|\,|\sigma(v_{k,j})|
 \le\frac12\Big(\sum_j|\sigma(v_{i,j})|^2+\sum_j|\sigma(v_{k,j})|^2\Big)\le4f^3M^4\operatorname{rd}_E^2.
\]
\end{proof}

\begin{lemma}\label{cv-class-count}
Let $E$ be a CM field with complex conjugation $c$ and $n=[E:\mathbf Q]$, let $M\ge1$, let
$h(x,y)=\sum_ic(x_i)y_i$ on $E^M$, and let $B\ge0$. Let $\mathcal L$ be a set of finitely generated
$\mathcal O_E$-submodules $\Lambda\subset E^M$ such that $h(\Lambda,\Lambda)\subset\mathcal O_E$ and
$\Lambda$ contains $v_1,\dots,v_M$, linearly independent over $E$, with
$|\sigma(h(v_i,v_k))|\le B$ for every embedding $\sigma$ and all $i,k$. Then there is a finite subset
$R\subset\mathcal L$ with
\[
 |R|\le\big((1+2B)\,MB\big)^{nM^2}
\]
such that every $\Lambda\in\mathcal L$ is $\gamma(\Lambda_0)$ for some $\Lambda_0\in R$ and some
$\gamma\in GL_M(E)$ preserving $h$.
\end{lemma}
\leanref{Development/AbelianVarietyTheory/Proofs/CuspFormEvaluations/ExistsFinsetSelfDualClassesHerm.lean}{GlobalTheory.ClassCount.exists_finset_classes}
\leanref{Development/AbelianVarietyTheory/Proofs/CuspFormEvaluations/ExistsFinsetSelfDualClassesHerm.lean}{GlobalTheory.ClassCount.exists_finset_types}
\begin{proof}
For $\Lambda\in\mathcal L$ choose $v$ as in the statement, let
$A=(h(v_i,v_k))_{i,k}\in M_M(\mathcal O_E)$ be its Gram matrix, and let
\[
 P(\Lambda)=\{(h(x,v_k))_k:x\in\Lambda\}\subset\mathcal O_E^M,
\]
an $\mathcal O_E$-submodule. The form $h$ is nondegenerate and the $v_i$ are a basis, so
$\det A\ne0$. Since $h(\sum_io_iv_i,v_k)=\sum_ic(o_i)A_{ik}$ and $c(\mathcal O_E)=\mathcal O_E$, the
module $P(\Lambda)$ contains the row space $\mathcal O_E^MA$.

Step 1: the pair $(A,P(\Lambda))$ determines $\Lambda$ up to isometry. Let $\Lambda'\in\mathcal L$,
with basis $v'$, give the same $A$ and $P$. Let $\gamma\in GL_M(E)$ with $\gamma v_i=v'_i$; it
preserves $h$, because the Gram matrices agree. For $x\in\Lambda$ the vector $(h(x,v_k))_k$ lies in
$P(\Lambda')$, so there is $x'\in\Lambda'$ with $h(x',v'_k)=h(x,v_k)=h(\gamma x,v'_k)$ for all $k$;
since the $v'_k$ are a basis and $h$ is nondegenerate, $x'=\gamma x$. So
$\gamma\Lambda\subset\Lambda'$, and in the same way $\gamma^{-1}\Lambda'\subset\Lambda$.

Step 2: the number of matrices. An element $a\in\mathcal O_E$ with $|\sigma(a)|\le B$ for all
$\sigma$ lies in a set of at most $(1+2B)^n$ elements. Indeed, a nonzero element of $\mathcal O_E$
has norm of absolute value at least $1$, so some conjugate has absolute value at least $1$; hence
the open sets $\{z\in E\otimes\mathbf R:|\sigma(z-a)|<\frac12\text{ for all }\sigma\}$ around these
$a$ are disjoint, and they lie in $\{z:|\sigma(z)|<B+\frac12\text{ for all }\sigma\}$, whose volume
is $(1+2B)^n$ times theirs. So there are at most $(1+2B)^{nM^2}$ matrices $A$.

Step 3: the number of modules for one matrix. Fix $A$ and put $a=\det A$. Then
$|\mathcal O_E^M/\mathcal O_E^MA|=|N_{E/\mathbf Q}(a)|=\prod_\sigma|\sigma(a)|\le(MB)^{nM}$, since
$|\sigma(a)|\le M!\,B^M$. A submodule $P$ with $\mathcal O_E^MA\subset P\subset\mathcal O_E^M$ is
determined by its image in this quotient, and that image is generated by $M$ elements: $P$ is a
projective $\mathcal O_E$-module of rank $M$ and $\mathcal O_E/a$ is a principal ideal ring, so there
are $x_1,\dots,x_M\in P$ with $P=\sum_i\mathcal O_Ex_i+aP$, and
$aP\subset a\mathcal O_E^M\subset\mathcal O_E^MA$ by the adjugate formula. So there are at most
$|\mathcal O_E^M/\mathcal O_E^MA|^M\le(MB)^{nM^2}$ such modules.

Let $R$ contain one $\Lambda\in\mathcal L$ for each pair $(A,P)$ that occurs. By Steps 2 and 3 it
has at most $((1+2B)MB)^{nM^2}$ elements, and by Step 1 it has the stated property.
\end{proof}

\begin{lemma}\label{cv-mass-bounds}
Let $E$ be a CM field with complex conjugation $c$, root discriminant
$\operatorname{rd}_E=|d_E|^{1/[E:\mathbf Q]}$ and $n=[E:\mathbf Q]$, and let $M\ge1$.
Then (Definition~\ref{cv-mass})
\[
 3^{-nM^2}\le\operatorname{mass}(I_M)\le\big(3\,n^2M^4\operatorname{rd}_E^2\big)^{3nM^2}.
\]
\end{lemma}
\leanref{Development/AbelianVarietyTheory/Proofs/DivisorProofs/EqZeroOfNormLtOne.lean}{GlobalTheory.StdMassLower.main}
\leanref{Development/AbelianVarietyTheory/Proofs/DivisorProofs/HermitianMassUpperBound.lean}{GlobalTheory.mass_stdLat_le}
\leanref{Development/AbelianVarietyTheory/Proofs/FoundationProofs/ExistsFinsetSelfDualClassesRootDiscr.lean}{GlobalTheory.exists_finset_selfDual_classes_rootDiscr}
\begin{proof}
Lower bound. The class of $I_M$ contributes $1/|\operatorname{Aut}(I_M)|$ to the mass. Let
$g\in\operatorname{Aut}(I_M)$. Its matrix $(g_{ab})$ has entries in $\mathcal O_E$, and
$h(ge_b,ge_b)=1$ for each standard basis vector $e_b$. Since $\varphi\circ c=\overline\varphi$
for every embedding $\varphi:E\to\mathbf C$, this says
$\sum_a|\varphi(g_{ab})|^2=1$, so $|\varphi(g_{ab})|\le1$ for all $\varphi$. Two elements
$x,y\in\mathcal O_E$ with all conjugates of absolute value at most $1$ and $x\equiv y$
modulo $3\mathcal O_E$ are equal: $(x-y)/3$ lies in $\mathcal O_E$ and all its conjugates have
absolute value at most $2/3$, so its norm is an integer of absolute value less than $1$. So
$g$ is determined by its reduction modulo $3$, and
$|\operatorname{Aut}(I_M)|\le|\mathcal O_E/3|^{M^2}=3^{nM^2}$.

Upper bound. Put $f=n/2$ and $B=4f^3M^4\operatorname{rd}_E^2$. A lattice in the genus of $I_M$ is
self-dual in the sense of Lemma~\ref{cv-short-basis}, since it is locally isometric to $I_M$, and
$h$ takes integral values on it. By Lemma~\ref{cv-short-basis}, the set $\mathcal L$ of self-dual
lattices satisfies the hypothesis of Lemma~\ref{cv-class-count} with this $B$. Let $R$ be the set
given by that lemma. Every lattice $M'$ in the genus is $\gamma(\Lambda_0)$ with $\Lambda_0\in R$
and $\gamma$ an isometry, so lies in the class of $\Lambda_0$; two lattices sent to the same
$\Lambda_0$ lie in the same class. Each class contributes at most $1$ to the mass, so
$\operatorname{mass}(I_M)\le|R|\le((1+2B)MB)^{nM^2}$. It remains to compare this with the stated
bound. Put $p=n^2M^4\operatorname{rd}_E^2=4f^2M^4\operatorname{rd}_E^2$, so that $B=fp$; as
$\operatorname{rd}_E\ge1$, we have $1\le f^2M\le p$ and $B\ge1$. Then
\[
 (1+2B)MB\le3B\cdot MB=3f^2M\,p^2\le3p^3\le(3p)^3,
\]
and $((1+2B)MB)^{nM^2}\le(3p)^{3nM^2}$, which is the stated bound.
\end{proof}

\begin{lemma}\label{cv-volume-mass-haar}
Let $P,Q\ge0$ with $N=P+Q\ge1$, let $J_{P,Q}$ be the diagonal matrix with $P$ entries $1$
followed by $Q$ entries $-1$, and let $U(P,Q)$ be the group of $g\in GL_N(\mathbf C)$ with
$\bar g^{\mathsf T}J_{P,Q}g=J_{P,Q}$, so that $U(N,0)$ is the unitary group $U(N)$. Let $\lambda$ be
the Lebesgue measure on the Lie algebra $\mathfrak u(P,Q)$ for which $\lambda\times\lambda$
corresponds to the Lebesgue measure on $M_N(\mathbf C)$ under $(X,Y)\mapsto X+iY$. The
\emph{density} of a Haar measure $m$ on $U(P,Q)$ is the limit of $m(\exp B_t)/\lambda(B_t)$ as
$t\to0^+$, where $B_t$ is the set of $X\in\mathfrak u(P,Q)$ with $\sum_{i,j}|X_{ij}|<t$. Let
$m^{\mathrm c}_{P,Q}$ be the Haar measure on $U(P,Q)$ whose density is that of the Haar probability
measure on $U(N)$, and put $\operatorname{vol}U(N)=\prod_{i=0}^{N-1}2\pi^{i+1}/i!$. Then:
\begin{enumerate}
\item the Haar probability measure on $U(N)$ has density $2^{-N(N-1)/2}/\operatorname{vol}U(N)$;
\item the Haar measure on $U(P,Q)$ of density $2^{-N(N-1)/2}$ is
  $\operatorname{vol}U(N)\cdot m^{\mathrm c}_{P,Q}$.
\end{enumerate}
\end{lemma}
\leanref{ReciprocityTheory/ReductiveGroups/DiscreteSeries/FormalDegree/HaarVolumes/BridgeA.lean}{ReciprocityTheory.ChartRestriction.DS.PolarCoordinateComparison.omegaDensity_main}
\leanref{Development/AbelianVarietyTheory/Proofs/DivisorProofs/StandardMeasureUnitaryPQEqVolumeSmulCompactDual.lean}{GlobalTheory.LowRow.CompactDualScale.isMinf_eq_volU_smul_cd}
\begin{proof}
Part (1) is the case of $U(N)$ of Macdonald's formula for the volume of a compact Lie group
\cite{Mac80}. For the Haar measure whose density is $1$ with respect to a Lebesgue measure on
$\mathfrak u(N)$, that formula gives $U(N)$ the volume $\operatorname{vol}U(N)$ times the covolume
of the lattice in $\mathfrak u(N)$ spanned by the matrices $iE_{kk}$, $E_{jk}-E_{kj}$ and
$i(E_{jk}+E_{kj})$ for $j<k$. Under $(X,Y)\mapsto X+iY$, the diagonal entries of $X$ and $Y$ are
carried to a diagonal entry with Jacobian $1$, and the entries of $X$ and $Y$ in positions $(j,k)$
and $(k,j)$ to the two entries of $X+iY$ there with Jacobian $4$. So this lattice has covolume
$2^{N(N-1)/2}$ for $\lambda$, and covolume $1$ for $2^{-N(N-1)/2}\lambda$, which gives (1). For
(2), two Haar measures on $U(P,Q)$ differ by a positive factor, which is the ratio of their
densities. By (1), $\operatorname{vol}U(N)\cdot m^{\mathrm c}_{P,Q}$ has density $2^{-N(N-1)/2}$.
\end{proof}

\begin{definition}\label{cv-volume-mass-normalized}
Let $E$ be a CM field with totally real subfield $F$, and let $W$ be a Hermitian space over $E$
with correct signature (Definition~\ref{th-spaces}), with form $h$, unimodular lattice $L$, and
signature $(P_\rho,Q_\rho)$ at the embeddings of $E$ above each real place $\rho$ of $F$. Let
$K_L\subset U(W)(\mathbf A_{F,f})$ be the stabilizer of
$\widehat L=L\otimes_{\mathcal O_F}\widehat{\mathcal O}_F$, and for each $\rho$ identify
$U(W)(F_\rho)$ with $U(P_\rho,Q_\rho)$ through a basis of $W\otimes_{F,\rho}\mathbf R$ in which
the form has matrix $J_{P_\rho,Q_\rho}$. The \emph{normalized measure} $\mu_W$ is the Haar measure on
$U(W)(\mathbf A_F)$ such that
\[
 \mu_W\{g:g_f\in K_L,\ g_\rho\in B_\rho\text{ for every }\rho\}
 =\prod_\rho m^{\mathrm c}_{P_\rho,Q_\rho}(B_\rho)
\]
for all measurable $B_\rho\subset U(P_\rho,Q_\rho)$, with $m^{\mathrm c}$ as in
Lemma~\ref{cv-volume-mass-haar}. The \emph{volume} of $W$ is $\operatorname{vol}(W)=\mu_W(\mathcal F)$
for a measurable fundamental domain $\mathcal F$ of $U(W)(F)$ in $U(W)(\mathbf A_F)$.
\end{definition}
\leanref{Formalism/SharpRestriction/GlobalSchurBounds/StandardLattice.lean}{Formalism.MainNotions.StandardLattice.IsVolMeas}
\leanref{Development/AbelianVarietyTheory/Definitions/Foundations/CdScaleTieIff.lean}{GlobalTheory.Mem.Vol.trueVol}

\begin{lemma}\label{cv-volume-mass-local-lattices}
Let $E$ be a CM field with totally real subfield $F$, and let $v$ be a finite place of $F$ that
does not ramify in $E$. Let $F_v$ be the completion, $\mathcal O_v$ its ring of integers,
$E_v=E\otimes_FF_v$ and $\mathcal O_{E,v}$ the integral closure of $\mathcal O_v$ in $E_v$. Let $h$
be a nondegenerate Hermitian form on $E_v^M$ with $M\ge1$, let $U(h)$ be its unitary group, and let
$L\subset E_v^M$ be a self-dual $\mathcal O_{E,v}$-lattice, that is,
$L=\{y:h(x,y)\in\mathcal O_{E,v}\text{ for all }x\in L\}$. Then:
\begin{enumerate}
\item $L$ has a basis $e_1,\dots,e_M$ with $h(e_i,e_j)=\delta_{ij}$;
\item every self-dual lattice $L'\subset E_v^M$ is $g(L)$ for some $g\in U(h)$;
\item let $\beta\in\mathcal O_v^\times$. There is $x\in L$ with $h(x,x)=\beta$. For each such $x$,
  $L=\mathcal O_{E,v}x\oplus(L\cap x^\perp)$ and $L\cap x^\perp$ is self-dual in $x^\perp$; the
  stabilizer $K_L$ of $L$ in $U(h)$ acts transitively on the set of such $x$; and if $U(h)_x$ is
  the stabilizer of $x$, then $K_L\cap U(h)_x$ is the stabilizer of $L\cap x^\perp$ in $U(h)_x$, and
  restriction to $x^\perp$ is an isomorphism from $K_L\cap U(h)_x$ onto the stabilizer of $L\cap x^\perp$
  in the unitary group of $x^\perp$;
\item let $\beta\in\mathcal O_v^\times$. The group $U(h)$ acts transitively on the set of
  $x\in E_v^M$ with $h(x,x)=\beta$, and for each such $x$ there are a self-dual lattice containing $x$
  and a self-dual lattice in $x^\perp$.
\end{enumerate}
\end{lemma}
\leanref{Development/AbelianVarietyTheory/Proofs/GlobalInterpolation/SelfDualLatticesUnramifiedPlaceNoHypothesis.lean}{GlobalTheory.UnramSelfDualClosed.selfDualLatticeLemmas}
\begin{proof}
Part (1) is Jacobowitz's description of unimodular Hermitian lattices over an unramified
quadratic extension \cite{Jac62}; when $v$ splits in $E$ it is elementary. Part (2) follows from
(1), with $g$ carrying an orthonormal basis of $L$ to one of $L'$.

For (3), the extension $E_v/F_v$ is unramified or split, so $\beta=a\bar a$ for some
$a\in\mathcal O_{E,v}^\times$, and $x=ae_1$ has $h(x,x)=\beta$. Let $x\in L$ with $h(x,x)=\beta$.
For $y\in L$ the vector $y-\beta^{-1}h(x,y)x$ is orthogonal to $x$ and lies in $L$, since
$h(x,y)\in\mathcal O_{E,v}$; this gives the decomposition, and $L\cap x^\perp$ is self-dual
because $L$ is. If $x'$ is a second such vector, $L\cap x^\perp$ and $L\cap x'^\perp$ have
orthonormal bases by (1), and the linear map that sends $x$ to $x'$ and the one basis to the other
is an isometry preserving $L$. An element of $U(h)_x$ preserves $L$ if and only if it preserves
$L\cap x^\perp$, by the decomposition. An element of $U(h)_x$ is determined by its restriction to
$x^\perp$, and a unitary transformation of $x^\perp$ that preserves $L\cap x^\perp$ extends, by fixing $x$,
to an element of $K_L\cap U(h)_x$; this gives the isomorphism.

(4) By Witt's theorem, $U(h)$ acts transitively on the vectors $x$ with $h(x,x)=\beta$. Let $x_0\in L$ be
a vector with $h(x_0,x_0)=\beta$, which exists by (3), and let $x=gx_0$ with $g\in U(h)$. Then $g(L)$ is a
self-dual lattice containing $x$, and $g(L\cap x_0^\perp)$, which is self-dual in $x_0^\perp$ by (3), is a
self-dual lattice in $g(x_0^\perp)=x^\perp$.
\end{proof}

\begin{lemma}\label{cv-volume-mass-invariance}
Let $E$ be a CM field with totally real subfield $F$ such that no finite prime of $F$ ramifies in
$E$, and let $W$ be a Hermitian space over $E$ with correct signature
(Definition~\ref{th-spaces}).
\begin{enumerate}
\item The normalized measure $\mu_W$ of Definition~\ref{cv-volume-mass-normalized} exists and does
  not depend on the identifications $U(W)(F_\rho)\cong U(P_\rho,Q_\rho)$. The volume
  $\operatorname{vol}(W)$ is positive and finite and does not depend on the fundamental domain.
\item If $L'$ is another unimodular lattice in $W$, then $W$ with the lattice $L'$ has the same
  normalized measure and the same volume.
\item If $W'$ is a Hermitian space over $E$ with correct signature and there is an isometry
  $W\to W'$, then $\operatorname{vol}(W')=\operatorname{vol}(W)$.
\end{enumerate}
\end{lemma}
\leanref{Development/AbelianVarietyTheory/Proofs/IwasawaDecomposition/UnitaryAdelicQuotientVolumeIsVolume.lean}{GlobalTheory.Mem.Vol.volIs_trueVol}
\leanref{Development/AbelianVarietyTheory/Proofs/GlobalInterpolation/IsFundamentalDomainMap.lean}{GlobalTheory.VolumeIndep.latticeIndependence}
\leanref{Development/AbelianVarietyTheory/Proofs/GlobalInterpolation/IsFundamentalDomainMap.lean}{GlobalTheory.VolumeIndep.volumeInvariance}
\begin{proof}
(1) Let $\mu$ be a Haar measure on $U(W)(\mathbf A_F)$. The group $K_L$ is compact and open, so
$B\mapsto\mu\{g:g_f\in K_L,\ (g_\rho)_\rho\in B\}$ is a Haar measure on $\prod_\rho U(P_\rho,Q_\rho)$,
equal to $c\prod_\rho m^{\mathrm c}_{P_\rho,Q_\rho}$ for some $c>0$, and $\mu_W=c^{-1}\mu$. It is
unique because Haar measures are proportional. Two identifications at $\rho$ differ by conjugation
in $U(P_\rho,Q_\rho)$, which preserves the Haar measure of this unimodular group. The volume is
finite by the theorem of Borel and Harish-Chandra \cite{BHC62}, and it is positive because the
translates of $\mathcal F$ by the countable group $U(W)(F)$ cover $U(W)(\mathbf A_F)$. Two
fundamental domains $\mathcal F,\mathcal F'$ have the same measure, since
$\mu_W(\mathcal F)=\sum_\gamma\mu_W(\mathcal F\cap\gamma\mathcal F')
=\sum_\gamma\mu_W(\gamma^{-1}\mathcal F\cap\mathcal F')=\mu_W(\mathcal F')$.

(2) By Lemma~\ref{cv-volume-mass-local-lattices}(2), at each finite place $v$ the completions of
$L$ and $L'$ are conjugate under $U(W)(F_v)$, and they are equal for almost all $v$. So there is
$k\in U(W)(\mathbf A_{F,f})$ with $k\widehat L=\widehat{L'}$, and $K_{L'}=kK_Lk^{-1}$. The group
$U(W)(\mathbf A_F)$ is unimodular, so conjugation by $k$ preserves $\mu_W$, and $\mu_W$ satisfies the
defining condition for $L'$.

(3) An isometry $W\to W'$ carries $U(W)(\mathbf A_F)$, $U(W)(F)$, the stabilizer of $\widehat L$
and the identifications at the real places to the corresponding objects for $W'$ with the image of
$L$, which is a unimodular lattice in $W'$. So $\operatorname{vol}(W)$ is the volume of $W'$ with
that lattice, which is $\operatorname{vol}(W')$ by (2).
\end{proof}

\begin{theorem}\label{cv-volume-mass-landherr}
Let $E$ be a CM field with totally real subfield $F$ such that no finite prime of $F$ ramifies in
$E$, and let $W$ and $W'$ be Hermitian spaces over $E$ with correct signature
(Definition~\ref{th-spaces}), of the same rank $M$, with the same signature at every complex
embedding of $E$. Then there is an isometry $W\to W'$. If the signature of $W$ is $(P,Q)$ at every
complex embedding, then $W$ has a basis $e_1,\dots,e_M$ with $h(e_i,e_j)=0$ for $i\ne j$,
$h(e_i,e_i)=1$ for $i\le P$ and $h(e_i,e_i)=-1$ for $i>P$.
\end{theorem}
\leanref{Development/AbelianVarietyTheory/Proofs/GlobalInterpolation/DetNeZeroOfHasSignature.lean}{GlobalTheory.LandherrSignature.exists_isometry_host}
\leanref{Development/AbelianVarietyTheory/Proofs/GlobalInterpolation/HasComplementOfBasis.lean}{GlobalTheory.UnitPerp.exists_isometry}
\begin{proof}
By Landherr's theorem \cite{Lan36}, two Hermitian spaces over $E$ of the same dimension are
isometric if they have the same signature at every complex embedding and the determinants of their
Gram matrices, which lie in $F^\times$, have the same class modulo $N_{E/F}E^\times$. Let $d$ and $d'$
be these determinants for $W$ and $W'$. At a finite place $v$ of $F$, a basis of the completion of
the unimodular lattice of $W$ has a Gram matrix whose determinant is a unit of $F_v$, and $d$
differs from that determinant by a norm from $E_v$. Units of $F_v$ are norms from $E_v$, since
$E_v/F_v$ is unramified or split, so $d$, and in the same way $d'$, is a norm at $v$. At a real
place $d/d'$ is positive, since the signatures agree. By the Hasse norm theorem for the quadratic
extension $E/F$, $d/d'$ is a norm from $E$. For the second statement, take for $W'$ the space $E^M$
with the form $\sum_{i\le P}\bar x_iy_i-\sum_{i>P}\bar x_iy_i$, whose lattice $\mathcal O_E^M$ is
unimodular.
\end{proof}

\begin{lemma}\label{cv-volume-mass-unit-vector}
Let $E$ be a CM field with totally real subfield $F$ such that no finite prime of $F$ ramifies in
$E$, and let $W$ be a Hermitian space over $E$ with correct signature
(Definition~\ref{th-spaces}), of rank $M\ge3$ and signature $(P,Q)$ at every complex embedding. Let
$\beta=1$ and $(P',Q')=(P-1,Q)$ if $P\ge1$, or $\beta=-1$ and $(P',Q')=(P,Q-1)$ if $Q\ge1$. Then there
are $x\in W$ with $h(x,x)=\beta$ and a unimodular lattice $L\subset W$ with $x\in L$ and
$L=\mathcal O_Ex\oplus(L\cap x^\perp)$. The orthogonal complement $x^\perp$, with the lattice
$L\cap x^\perp$, is a Hermitian space with correct signature $(P',Q')$ at every complex embedding.
\end{lemma}
\leanref{Development/AbelianVarietyTheory/Proofs/GlobalInterpolation/HasComplementOfBasis.lean}{GlobalTheory.UnitPerp.unitVectorPerp}
\begin{proof}
Let $e_1,\dots,e_M$ be the basis of Theorem~\ref{cv-volume-mass-landherr}, and let $L$ be the
$\mathcal O_E$-module it spans, which is unimodular because the Gram matrix of the $e_i$ is diagonal
with entries $\pm1$. Take $x=e_1$ if $\beta=1$ and $x=e_M$ if $\beta=-1$. Then $x^\perp$ is spanned by
the other $e_i$, $L\cap x^\perp$ is their $\mathcal O_E$-span, and the signature of $x^\perp$ is
$(P',Q')$.
\end{proof}

\begin{lemma}\label{cv-volume-mass-definite}
Let $E$ be a CM field with totally real subfield $F$ such that no finite prime of $F$ ramifies in
$E$, and let $W$ be a Hermitian space over $E$ with correct signature
(Definition~\ref{th-spaces}), of rank $M\ge1$, with form $h$ and unimodular lattice $L$.
\begin{enumerate}
\item If the signature of $W$ is $(M,0)$ at every complex embedding, then
  $\operatorname{vol}(W)=\operatorname{mass}(L)=\operatorname{mass}(I_M)$
  (Definitions~\ref{cv-volume-mass-normalized} and~\ref{cv-mass}).
\item If the signature of $W$ is $(0,M)$ at every complex embedding, then
  $\operatorname{vol}(W)=\operatorname{vol}(W^-)$, where $W^-$ is $W$ with the form $-h$ and the
  lattice $L$.
\end{enumerate}
\end{lemma}
\leanref{Formalism/LocalToGlobal/GlobalTheory/GL5/FamilyMassVolumeIdentityOfClassDecomp.lean}{Formalism.L2.GlobalTheoryTheory.GenusMassVolumes.mass_volume_identity_of_class_decomp}
\leanref{Development/AbelianVarietyTheory/Proofs/FoundationProofs/FiniteAndMass.lean}{GlobalTheory.DefiniteDecomp.defClassDecomp}
\leanref{Development/AbelianVarietyTheory/Proofs/GlobalInterpolation/UnramifiedOfIsCMUnram.lean}{GlobalTheory.DefiniteMass.mass_eq_massI}
\leanref{Development/AbelianVarietyTheory/Proofs/GlobalInterpolation/VolumeNegativeDefiniteEqOpposite.lean}{GlobalTheory.NegVol.vol_neg_eq}
\begin{proof}
(1) Put $G=U(W)(\mathbf A_F)$, $\Gamma=U(W)(F)$ and $K=K_L\times\prod_\rho U(W)(F_\rho)$, an open
subgroup of $G$. Since $m^{\mathrm c}_{M,0}$ is the Haar probability measure on $U(M)$,
$\mu_W(K)=1$. Since $W$ is totally definite, the double coset space $\Gamma\backslash G/K$ is finite,
say with representatives $g_1,\dots,g_n$, and the groups $\Gamma_i=\Gamma\cap g_iKg_i^{-1}$ are
finite. The map sending $g$ to the lattice $W\cap g_f\widehat L$ induces a bijection from
$\Gamma\backslash G/K$ to the classes in the genus of $L$, under which $\Gamma_i$ is the automorphism
group of $W\cap g_{i,f}\widehat L$; so $\sum_i1/|\Gamma_i|=\operatorname{mass}(L)$. A fundamental
domain $\mathcal F$ meets $\Gamma g_iK$ in a fundamental domain for the action of $\Gamma$ on
$\Gamma g_iK$, whose measure is $\mu_W(g_iK)/|\Gamma_i|=1/|\Gamma_i|$. So
$\operatorname{vol}(W)=\sum_i1/|\Gamma_i|=\operatorname{mass}(L)$.

By Theorem~\ref{cv-volume-mass-landherr} there is an isometry from $W$ to $E^M$ with the form
$\sum_i\bar x_iy_i$. It carries $L$ to a self-dual lattice $L_1$, and the mass does not change.
The lattice $L_1$ lies in the genus of $I_M$: for a maximal ideal $\mathfrak p$ of $\mathcal O_E$
over the prime $p$, the completions of $L_1$ at the places of $F$ above $p$ have orthonormal bases
by Lemma~\ref{cv-volume-mass-local-lattices}(1), and weak approximation in the unitary group gives
an element of $U(E^M)$ close to the matrices of these bases at those places, which carries $I_M$ to
$L_1$ after localization at $\mathfrak p$. The mass is constant on a genus, so
$\operatorname{mass}(L)=\operatorname{mass}(I_M)$.

(2) The groups $U(W)$ and $U(W^-)$ are equal, with the same stabilizer $K_L$ and the same subgroup
$K$, so they have the same double cosets and the same groups $\Gamma_i$. The measure
$m^{\mathrm c}_{0,M}$ on $U(0,M)=U(M)$ is the Haar probability measure, so $\mu_W(K)=\mu_{W^-}(K)=1$,
and the computation in (1) gives $\operatorname{vol}(W)=\sum_i1/|\Gamma_i|=\operatorname{vol}(W^-)$.
\end{proof}

\begin{lemma}\label{cv-volume-mass-unfolding}
Let $E$ be a CM field with totally real subfield $F$ such that no finite prime of $F$ ramifies in
$E$, and let $W$ be a Hermitian space over $E$ with correct signature
(Definition~\ref{th-spaces}), of rank $M\ge3$, with form $h$, unimodular lattice $L$ and signature
$(P,Q)$ at every complex embedding. Let $\beta=\pm1$, let $x\in L$ with $h(x,x)=\beta$, and suppose
that $W'=x^\perp$, with the lattice $L'=L\cap x^\perp$, is a Hermitian space with correct signature
$(P',Q')$ at every complex embedding, where $(P',Q')$ is $(P-1,Q)$ if $\beta=1$ and $(P,Q-1)$ if
$\beta=-1$. Let $G_\infty=\prod_\rho U(W)(F_\rho)$, let $H_\infty\subset G_\infty$ be the stabilizer
of $x$, which is $\prod_\rho U(W')(F_\rho)$, and let $\nu$ be the measure on
$H_\infty\backslash G_\infty$ with
$\int_{G_\infty}\phi=\int_{H_\infty\backslash G_\infty}\int_{H_\infty}\phi(hg)\,dh\,d\nu(H_\infty g)$
for measurable $\phi\ge0$, where $G_\infty$ carries $\prod_\rho m^{\mathrm c}_{P,Q}$ and $H_\infty$
carries $\prod_\rho m^{\mathrm c}_{P',Q'}$ (Lemma~\ref{cv-volume-mass-haar}). Let $\Phi_\infty\ge0$
be a measurable function on $W\otimes_F(F\otimes_{\mathbf Q}\mathbf R)$, and let $\Phi$ be the
function on $W\otimes_F\mathbf A_F$ that is $\Phi_\infty$ at the infinite places times the
indicator of $\widehat L$ at the finite places. Let $\mu_W$ and $\operatorname{vol}$ be as in
Definition~\ref{cv-volume-mass-normalized}. Then for every measurable fundamental domain
$\mathcal F$ of $U(W)(F)$ in $U(W)(\mathbf A_F)$,
\[
 \int_{\mathcal F}\sum_{y\in W,\ h(y,y)=\beta}\Phi(g^{-1}y)\,d\mu_W(g)
 =\operatorname{vol}(W')\int_{H_\infty\backslash G_\infty}\Phi_\infty(g^{-1}x)\,d\nu(g).
\]
\end{lemma}
\leanref{Development/AbelianVarietyTheory/Proofs/GlobalInterpolation/Unfolding.lean}{GlobalTheory.UnfoldingObj.unfolding}
\leanref{Development/AbelianVarietyTheory/Proofs/GlobalInterpolation/StabSelfDualLatticeTransitiveOnLength.lean}{GlobalTheory.StabOrbit.stabTransitive}
\begin{proof}
Let $G=U(W)(\mathbf A_F)$, $\Gamma=U(W)(F)$, and let $H\subset G$ be the stabilizer of $x$, which is
$U(W')(\mathbf A_F)$ acting on $x^\perp$. By Witt's theorem, $\Gamma$ acts transitively on the
$y\in W$ with $h(y,y)=\beta$, and the stabilizer of $x$ in $\Gamma$ is $\Gamma'=U(W')(F)$. So the sum
is $\sum_{\gamma\in\Gamma'\backslash\Gamma}\Phi(g^{-1}\gamma^{-1}x)$, and the left side is the
integral of $\Phi(g^{-1}x)$ over $\Gamma'\backslash G$. Both $G$ and $H$ are unimodular. Integrating
first over $\Gamma'\backslash H$ for $\mu_{W'}$, where the integrand is constant, and then over
$H\backslash G$ for the quotient measure $\bar\nu$ of $\mu_W$ by $\mu_{W'}$, the left side is
$\operatorname{vol}(W')\int_{H\backslash G}\Phi(g^{-1}x)\,d\bar\nu(g)$. The integrand is
$\Phi_\infty(g_\infty^{-1}x)$ times the indicator of the set of $g$ with
$g_f^{-1}x\in\widehat L$, and $\bar\nu$ is the product of $\nu$ with the quotient $\bar\nu_f$ of the
finite parts $\mu_{W,f}$ and $\mu_{W',f}$ of the normalized measures, which give $K_L$ and
$K_{L'}$ the measure $1$. By Lemma~\ref{cv-volume-mass-local-lattices}(3) at every finite place,
$K_L$ acts transitively on the $y\in\widehat L$ with $h(y,y)=\beta$; so the set of $g_f$ with
$g_f^{-1}x\in\widehat L$ is $H_fK_L$, where $H_f$ is the finite part of $H$, and $K_L\cap H_f$ is the
stabilizer $K_{L'}$ of $\widehat{L'}$ in $H_f$. So the finite factor is
$\bar\nu_f(H_f\backslash H_fK_L)=\mu_{W,f}(K_L)/\mu_{W',f}(K_{L'})=1$, which gives the formula.
\end{proof}

\begin{theorem}\label{cv-volume-mass-siegel-weil}
Let $E$ be a CM field with totally real subfield $F$, let $h$ be a nondegenerate Hermitian form on
$V=E^M$ with $M\ge3$, let $\beta\in F^\times$, and let $\psi$ be a nontrivial continuous character
of $\mathbf A_F/F$. For each finite place $v$ of $F$ let $dx_v$ be the Haar measure on
$V_v=V\otimes_FF_v$ that is self-dual for the pairing $\psi_v(\operatorname{Tr}_{E_v/F_v}h(x,y))$, let
$dt_v$ be the Haar measure on $F_v$ that is self-dual for $\psi_v(st)$, and for $t\in F_v^\times$ let
$dx_v^t$ be measures on $\{x\in V_v:h(x,x)=t\}$ with
$\int_{V_v}\varphi\,dx_v=\int_{F_v}\int\varphi\,dx_v^t\,dt_v$ for measurable $\varphi\ge0$, such that
$t\mapsto\int\varphi\,dx_v^t$ is continuous for continuous compactly supported $\varphi$. Define
$dx_\infty$, $dt_\infty$ and $dx_\infty^t$ in the same way for $F_\infty=F\otimes_{\mathbf Q}\mathbf R$
and $\psi_\infty$. Let $\Phi_\infty\ge0$ be measurable on $V_\infty$, let $\Phi_v\ge0$ be locally
constant and compactly supported on $V_v$, equal to the indicator of $\mathcal O_{E,v}^M$ for almost
all $v$, and let $\Phi=\Phi_\infty\otimes\bigotimes_v\Phi_v$. Let $\mu$ be a Haar measure on
$U(h)(\mathbf A_F)$ and $\mathcal F$ a measurable fundamental domain of $U(h)(F)$. Then the product
over $v$ below converges, and
\[
 \frac1{\mu(\mathcal F)}\int_{\mathcal F}\sum_{x\in V,\ h(x,x)=\beta}\Phi(g^{-1}x)\,d\mu(g)
 =\int\Phi_\infty\,dx_\infty^\beta\cdot\prod_v\int\Phi_v\,dx_v^\beta.
\]
\end{theorem}
\leanref{CompanionTheory/AutomorphicTheory/SiegelUnitTheory/UnitsAinf.lean}{CompanionTheory.HU.sw1b2}
\leanref{Formalism/SharpRestriction/GlobalSchurBounds/StandardLattice.lean}{Formalism.MainNotions.StandardLattice.SW1b2}

The left side is the mean over $U(h)(F)\backslash U(h)(\mathbf A_F)$ of a theta series, and the
right side is a product of local densities of representations of $\beta$ by $h$. The identity is
Siegel's formula for the representations of a number by a Hermitian form, in the adelic form of
Weil \cite{Wei65}.

\begin{lemma}\label{cv-volume-mass-local-density}
Let $E$ be a CM field with totally real subfield $F$, let $v$ be a finite place of $F$ that does
not ramify in $E$, with residue field of $q_v$ elements, and let $\eta_v=-1$ if $E\otimes_FF_v$ is a
field and $\eta_v=1$ otherwise. Let $\psi_v$ be a nontrivial character of $F_v$, let $c$ be the
integer with $\{t\in F_v:\psi_v(t\mathcal O_v)=1\}=\mathfrak p_v^{-c}$, and put $N_v=q_v^c$. Let $h$
be a nondegenerate Hermitian form on $E_v^M$ with $M\ge1$, let $L\subset E_v^M$ be a self-dual
$\mathcal O_{E,v}$-lattice (Lemma~\ref{cv-volume-mass-local-lattices}), and let $dx_v$, $dt_v$ and
$dx_v^t$ be as in Theorem~\ref{cv-volume-mass-siegel-weil}. Then $dx_v(L)=N_v^{-M}$,
$dt_v(\mathcal O_v)=N_v^{-1/2}$, and for every $\beta\in\mathcal O_v^\times$
\[
 \int\mathbf 1_L\,dx_v^\beta=N_v^{1/2-M}\big(1-\eta_v^Mq_v^{-M}\big).
\]
\end{lemma}
\leanref{Development/AbelianVarietyTheory/Proofs/GlobalInterpolation/SelfDualLatticeDensity.lean}{GlobalTheory.LatticeDensityClosed.selfDualLatticeDensity}
\begin{proof}
Since $E_v/F_v$ is unramified or split, the dual of $\mathcal O_{E,v}$ for $\psi_v\circ\operatorname{Tr}$
is $\mathfrak p_v^{-c}\mathcal O_{E,v}$, so the dual of $L$ for the pairing of $dx_v$ is
$\mathfrak p_v^{-c}L$, of index $q_v^{2cM}$ over $L$. A self-dual measure gives a lattice and its
dual volumes with product $1$, so $dx_v(L)=N_v^{-M}$; in the same way $dt_v(\mathcal O_v)=N_v^{-1/2}$.

For a unit $u$ of $\mathcal O_v$ there is $a\in\mathcal O_{E,v}^\times$ with $a\bar a=u$, and
$x\mapsto ax$ preserves $L$ and $dx_v$ and carries the set where $h(x,x)=t$ to the set where
$h(x,x)=ut$. So $\int\mathbf 1_L\,dx_v^t$ is the same for all units $t$, and
\[
 dx_v\{x\in L:h(x,x)\in\mathcal O_v^\times\}=\int\mathbf 1_L\,dx_v^\beta\cdot dt_v(\mathcal O_v^\times)
 =\int\mathbf 1_L\,dx_v^\beta\cdot N_v^{-1/2}(1-q_v^{-1}).
\]
In an orthonormal basis of $L$ (Lemma~\ref{cv-volume-mass-local-lattices}(1)), $h(x,x)$ is the sum
of the norms of the coordinates, so the left side is $N_v^{-M}q_v^{-2M}$ times the number of
$x\in(\mathcal O_{E,v}/\mathfrak p_v\mathcal O_{E,v})^M$ whose sum of norms is nonzero. The residue ring is
$\mathbf F_{q_v^2}$ with its norm form if $\eta_v=-1$, and $\mathbf F_{q_v}\times\mathbf F_{q_v}$ with
the form $(a,b)\mapsto ab$ if $\eta_v=1$, and in both cases that number is
$(q_v-1)(q_v^{2M-1}-\eta_v^Mq_v^{M-1})$. Dividing gives the formula.
\end{proof}

\begin{lemma}\label{cv-volume-mass-euler-product}
Let $E$ be a CM field with totally real subfield $F$, $f=[F:\mathbf Q]$, such that no finite prime
of $F$ ramifies in $E$, and put $v_E=|d_E|^{1/(2f)}/2$. Let $\psi$ be the character of
$\mathbf A_F/F$ obtained by composing the standard character of $\mathbf A_{\mathbf Q}/\mathbf Q$ with
the trace, so that $\psi_\rho(t)=e^{-2\pi it}$ at every real place $\rho$, and for each finite place
$v$ let $N_v$, $q_v$ and $\eta_v$ be as in Lemma~\ref{cv-volume-mass-local-density} for $\psi_v$.
Then $\prod_vN_v=|d_F|=(2v_E)^f$, and for every $M\ge2$ the product
$\prod_vN_v^{1/2-M}(1-\eta_v^Mq_v^{-M})$ converges to $(2v_E)^{f(1/2-M)}\Lambda_E(M)^{-1}$, with
$\Lambda_E(M)$ as in Theorem~\ref{cv-siegel-mass-ratio}.
\end{lemma}
\leanref{Development/AbelianVarietyTheory/Proofs/GlobalInterpolation/EulerProductOfDensities.lean}{GlobalTheory.DensityProductClosed.eulerProductOfDensities}
\begin{proof}
For this $\psi$, the set of $t\in F_v$ with $\psi_v(t\mathcal O_v)=1$ is the inverse of the local
different of $F$ at $v$, so $N_v$ is the norm of the local different, and the product of these
norms is $|d_F|$. No finite prime ramifies in $E/F$, so $|d_E|=|d_F|^2$, that is,
$|d_F|=|d_E|^{1/2}=(2v_E)^f$. A finite place $v$ of $F$ is either inert in $E$, with one prime
$\mathfrak P$ above it, $c(\mathfrak P)=\mathfrak P$ and $N\mathfrak P=q_v^2$, or split, with two
primes $\mathfrak P\ne c(\mathfrak P)$ of norm $q_v$. In both cases the factors of $\Lambda_E(M)$ at
the primes above $v$ multiply to $(1-\eta_v^Mq_v^{-M})^{-1}$, and the product over $v$ converges
absolutely for $M\ge2$.
\end{proof}

\begin{lemma}\label{cv-volume-mass-density-gauges}
Let $P',Q'\ge0$ and $\beta=\pm1$, let $(P,Q)$ be $(P'+1,Q')$ if $\beta=1$ and $(P',Q'+1)$ if $\beta=-1$,
and put $M=P+Q$. Let $h_{P,Q}$, $q$, $dx$ and the measures $dx^t$ be as in
Lemma~\ref{cv-volume-mass-archimedean}, and let $\lambda_{P,Q}$ and the sets $B_r$ be as in
Lemma~\ref{cv-volume-mass-haar}. Write a vector of $\mathbf C^M$ as $y=(y_0,z)$, where $y_0$ is a
coordinate at which $h_{P,Q}$ has the sign $\beta$ and $z\in\mathbf C^{M-1}$ are the others, ordered so
that $q(y)=\beta|y_0|^2+q'(z)$ with $q'(z)=h_{P',Q'}(z,z)$. Let $e_0=(1,0)$, so that $q(e_0)=\beta$, and
let $\iota_0\colon U(P',Q')\to U(P,Q)$ send $g'$ to the element that fixes $e_0$ and acts by $g'$ on the
coordinates $z$. For $g$ in a unitary group, $X$ in its Lie algebra, $y\in\mathbf C^M$ and
$(u,z)\in\mathbf R\times\mathbf C^{M-1}$ put
\[
 \|g\|=\sum_{a,b}|(g-1)_{ab}|,\qquad|X|_1=\sum_{a,b}|X_{ab}|,\qquad
 \|y\|_0=|y_0-1|+2\sum_a|z_a|,\qquad|(u,z)|=|u|+2\sum_a|z_a|,
\]
and let $\operatorname{vol}$ be the Lebesgue measure on $\mathbf R\times\mathbf C^{M-1}=\mathbf R^{2M-1}$.
Then:
\begin{enumerate}
\item for $r,s>0$, $0<\lambda_{P,Q}(B_r)<\infty$ and $\lambda_{P,Q}(B_{sr})=s^{M^2}\lambda_{P,Q}(B_r)$,
  and $0<\operatorname{vol}\{|(u,z)|<r\}<\infty$ with
  $\operatorname{vol}\{|(u,z)|<sr\}=s^{2M-1}\operatorname{vol}\{|(u,z)|<r\}$; and for every $t>1$ and every
  small $r>0$, $\exp(B_{r/t})\subset\{g:\|g\|<r\}\subset\exp(B_{tr})$ in $U(P,Q)$;
\item for every measurable set $S\subset\mathbf C^M$ on which $y_0\ne0$,
  $dx^\beta(S)=2^{M-1}\operatorname{vol}\{(u,z):-\pi<u<\pi,\ p(u,z)\in S\}$, where
  $p(u,z)=\big(\sqrt{1-\beta q'(z)}\,e^{iu},z\big)$;
\item for every $t>1$, every small $r>0$ and every $(u,z)$ with $-\pi<u<\pi$: if $|(u,z)|<r/t$ then
  $\|p(u,z)\|_0<r$, and if $\|p(u,z)\|_0<r$ then $|(u,z)|<tr$;
\item $\lambda_{P,Q}(B_r)=2^{M-1}(\operatorname{vol}\times\lambda_{P',Q'})
  \{((u,z),X'):|(u,z)|+|X'|_1<r\}$ for every $r$;
\item there is a measurable map $\sigma\colon\mathbf C^M\to U(P,Q)$ with $\sigma(y)e_0=y$ whenever
  $q(y)=\beta$, such that if moreover $y_0\ne-1$, then $|\sigma(y)_{0b}|=|z_b|$ and
  $|(\sigma(y)-1)_{ab}|=|z_a||z_b|/|1+y_0|$ for all coordinates $a,b$ of $z$;
\item for every $t>1$ there is $\delta>0$ such that for every $y$ with $\|y\|_0<\delta$, every
  $g\in U(P,Q)$ with $ge_0=y$ whose entries satisfy the equalities of (5) for $y$, and every
  $g'\in U(P',Q')$, $\|y\|_0+\|g'\|\le t\|g\iota_0(g')\|$ and $\|g\iota_0(g')\|\le t(\|y\|_0+\|g'\|)$;
  and $\|ge_0\|_0\le2\|g\|$ for every $g\in U(P,Q)$;
\item let $\mu_k$ and $\nu_k$, for $k=1,2$, be s-finite measures and $n_k$ and $N_k$ nonnegative
  measurable functions with $0<\nu_k\{N_k<r\}<\infty$ for $r>0$. If
  $\mu_k\{n_k<r\}/\nu_k\{N_k<r\}$ tends to $c_k<\infty$ as $r\to0^+$ for $k=1,2$, then
  $(\mu_1\times\mu_2)\{n_1+n_2<r\}/(\nu_1\times\nu_2)\{N_1+N_2<r\}$ tends to $c_1c_2$.
\end{enumerate}
\end{lemma}
\leanref{Development/AbelianVarietyTheory/Proofs/GlobalInterpolation/IsCompactMatrixEntryAbsSumSublevel.lean}{GlobalTheory.CompactDualQuot.ballU_props}
\leanref{Development/AbelianVarietyTheory/Proofs/GlobalInterpolation/IsCompactMatrixEntryAbsSumSublevel.lean}{GlobalTheory.CompactDualQuot.tanBall_props}
\leanref{Development/AbelianVarietyTheory/Proofs/GlobalInterpolation/UnitaryExpLieBallComparableEntryGaugeStatement.lean}{GlobalTheory.T.unitaryExpLieBall_comparable_entryGauge}
\leanref{Development/AbelianVarietyTheory/Proofs/GlobalInterpolation/HermitianNormFibreMeasureEqPolarCoordinatesStatement.lean}{GlobalTheory.T.hermitianNormFibreMeasure_eq_polar_coordinates}
\leanref{Development/AbelianVarietyTheory/Proofs/GlobalInterpolation/QuadricPointGaugePolarCoordinatesCompareTangentGStatement.lean}{GlobalTheory.T.quadricPointGauge_polarCoordinates_compare_tangentGauge}
\leanref{Development/AbelianVarietyTheory/Proofs/GlobalInterpolation/UnitaryLieVolumeBallSplitBasisVectorStatement.lean}{GlobalTheory.T.unitaryLie_volume_ball_split_basis_vector}
\leanref{Development/AbelianVarietyTheory/Proofs/GlobalInterpolation/ExistsTangentBorelSectionQuadricStatement.lean}{GlobalTheory.T.exists_tangent_borelSection_quadric}
\leanref{Development/AbelianVarietyTheory/Proofs/GlobalInterpolation/QuadricPointGaugeSectionMulBlockCompareStatement.lean}{GlobalTheory.T.quadricPointGauge_section_mul_block_compare}
\leanref{Development/AbelianVarietyTheory/Proofs/GlobalInterpolation/IsCompactMatrixEntryAbsSumSublevel.lean}{GlobalTheory.CompactDualQuot.ptGauge_act_le}
\leanref{Development/AbelianVarietyTheory/Proofs/GlobalInterpolation/TendstoSublevelDensityProdMeasureStatement.lean}{GlobalTheory.T.tendsto_sublevel_density_prod_measure}

Each part is a self-contained computation in coordinates. In (2), $p$ is the parametrization of the
quadric by the argument of $y_0$ and the coordinates $z$, and the factor $2^{M-1}$ is the ratio of the
normalizations of $dx$ and of the Lebesgue measure. In (4), $\mathfrak u(P,Q)$ is the sum of
$\mathfrak u(P',Q')$, the line $i\mathbf R$ at the entry $(0,0)$, and the vectors $z$ of the row and column
of $y_0$, which enter $|X|_1$ with the weight $2$. In (5), $\sigma(y)$ is built from the unitary
reflection that exchanges $e_0$ and $-y$.

\begin{lemma}\label{cv-volume-mass-archimedean}
Let $P,Q\ge0$ and $M=P+Q$, let $h_{P,Q}(x,y)=\sum_{j\le P}\bar x_jy_j-\sum_{j>P}\bar x_jy_j$ on
$\mathbf C^M$ and $q(x)=h_{P,Q}(x,x)$, let $dx$ be $2^M$ times the Lebesgue measure on
$\mathbf C^M=\mathbf R^{2M}$, and for $t\ne0$ let $dx^t$ be measures on $\{q=t\}$ with
$\int\varphi\,dx=\int_{\mathbf R}\int\varphi\,dx^t\,dt$ for measurable $\varphi\ge0$, continuous in
$t$ as in Theorem~\ref{cv-volume-mass-siegel-weil}. For $M\ge1$ put $c_M=(2\pi)^M/(M-1)!$, and let
$m^{\mathrm c}$ be as in Lemma~\ref{cv-volume-mass-haar}. Then:
\begin{enumerate}
\item for $e=\pm1$, $dx$ is the Haar measure on $\mathbf C^M$ that is self-dual for the pairing
  $e^{2\pi ie\operatorname{Tr}_{\mathbf C/\mathbf R}h_{P,Q}(x,y)}$, and the Lebesgue measure is the
  Haar measure on $\mathbf R$ that is self-dual for $e^{2\pi iest}$;
\item if $Q=0$ and $M\ge1$, the sphere $\{q=1\}$ has mass $c_M$ for $dx^1$;
\item let $\beta=\pm1$, let $x_0\in\mathbf C^M$ with $q(x_0)=\beta$, let $(P',Q')$ be $(P-1,Q)$ if
  $\beta=1$ and $(P,Q-1)$ if $\beta=-1$, let $j$ be an isometry from $(\mathbf C^{M-1},h_{P',Q'})$
  onto the orthogonal complement of $x_0$, and let $\iota\colon U(P',Q')\to U(P,Q)$ send $g'$ to the
  element fixing $x_0$ and acting as $jg'j^{-1}$ on its complement. Then $U(P,Q)$ acts transitively
  on $\{q=\beta\}$, the stabilizer of $x_0$ is $\iota(U(P',Q'))$, and for every measurable map $s$
  from $\{q=\beta\}$ to $U(P,Q)$ with $s(y)^{-1}x_0=y$ and every measurable $\phi\ge0$ on $U(P,Q)$,
  \[
   \int\phi\,dm^{\mathrm c}_{P,Q}
   =c_M^{-1}\int_{\{q=\beta\}}\int\phi(\iota(g')s(y))\,dm^{\mathrm c}_{P',Q'}(g')\,dx^\beta(y).
  \]
\end{enumerate}
\end{lemma}
\leanref{Development/AbelianVarietyTheory/Proofs/GlobalInterpolation/RealPlaceMeasures.lean}{GlobalTheory.RealPlaceMeasuresProof.realPlaceMeasures}
\leanref{Development/AbelianVarietyTheory/Proofs/GlobalInterpolation/IsCompactMatrixEntryAbsSumSublevel.lean}{GlobalTheory.CompactDualQuot.quotFormula}
\begin{proof}
(1) In real coordinates the pairing $\operatorname{Tr}h_{P,Q}(x,y)=2\operatorname{Re}h_{P,Q}(x,y)$ has
a diagonal matrix with entries $\pm2$, so the self-dual measure is $2^M$ times the Lebesgue measure;
the case of $\mathbf R$ is the same with the entry $1$. (2) For $Q=0$ write $x=r\omega$ with $r>0$
and $\omega$ in the unit sphere $S^{2M-1}$. Then $dx=2^Mr^{2M-1}\,dr\,d\omega$ with $d\omega$ the
surface measure, and with $t=r^2$ this is $2^{M-1}t^{M-1}\,dt\,d\omega$. So $dx^1$ is
$2^{M-1}d\omega$, of mass $2^{M-1}\cdot2\pi^M/(M-1)!=c_M$.

(3) Transitivity and the stabilizer follow from Witt's theorem. The groups $U(P,Q)$ and $U(P',Q')$ are
unimodular, so $g\mapsto g^{-1}$ preserves their Haar measures. Applying this to $\phi$ and replacing $s$
by $y\mapsto s(y)^{-1}$, it suffices to show that
\[
 (*)\qquad\int_{\{q=\beta\}}\int f(s(y)\iota(g'))\,dm^{\mathrm c}_{P',Q'}(g')\,dx^\beta(y)
 =c_M\int f\,dm^{\mathrm c}_{P,Q}
\]
for every measurable $f\ge0$ on $U(P,Q)$ and every measurable $s$ with $s(y)x_0=y$ on $\{q=\beta\}$.

Reduction to a basis vector. Let the coordinates $(y_0,z)$, the vector $e_0$ and the embedding
$\iota_0$ be as in Lemma~\ref{cv-volume-mass-density-gauges} for $(P',Q')$ and $\beta$. Let $g_0$ be the
linear map with $g_0e_0=x_0$ and $g_0(0,z)=j(z)$. It preserves $h_{P,Q}$, so $g_0\in U(P,Q)$, and
$g_0\iota_0(g')g_0^{-1}=\iota(g')$, since both fix $x_0$ and act as $jg'j^{-1}$ on its complement. If $s$ is
as in $(*)$, then $y\mapsto s(y)g_0$ carries $e_0$ to $y$, and
$f(s(y)\iota(g'))=f_0(s(y)g_0\iota_0(g'))$ with $f_0(g)=f(gg_0^{-1})$, whose integral against
$m^{\mathrm c}_{P,Q}$ is that of $f$. So it suffices to prove $(*)$ for $x_0=e_0$ and $\iota=\iota_0$.

The constant. The measure $\nu=dx^\beta$ is finite on compact sets and invariant under $U(P,Q)$, and
$\iota_0$ is a closed embedding onto the stabilizer of $e_0$. By the uniqueness of invariant measures on
the quotient of a unimodular group by a closed unimodular subgroup \cite[\S2.6]{Fol95}, applied to
$\{q=\beta\}$, there is $\kappa\ge0$ such that the left side of $(*)$ is
$\kappa\int f\,dm^{\mathrm c}_{P,Q}$ for every such $s$ and every $f$.

The densities: $\kappa a=a'$, where $a$ and $a'$ are the densities of $m=m^{\mathrm c}_{P,Q}$ and
$m'=m^{\mathrm c}_{P',Q'}$ (Lemma~\ref{cv-volume-mass-haar}). Write $\lambda=\lambda_{P,Q}$ and
$\lambda'=\lambda_{P',Q'}$, and use the gauges and the parts of
Lemma~\ref{cv-volume-mass-density-gauges}.
\begin{enumerate}
\item[(i)] $m\{\|g\|<r\}/\lambda(B_r)\to a$ as $r\to0^+$. Indeed, for $t>1$ and small $r$, part (1) puts this
  quotient between $t^{-M^2}m(\exp B_{r/t})/\lambda(B_{r/t})$ and $t^{M^2}m(\exp B_{tr})/\lambda(B_{tr})$,
  which tend to $t^{-M^2}a$ and $t^{M^2}a$. In the same way $m'\{\|g'\|<r\}/\lambda'(B_r)\to a'$.
\item[(ii)] $\nu\{\|y\|_0<r\}/\operatorname{vol}\{|(u,z)|<r\}\to2^{M-1}$. For $r<1$ the set $\{\|y\|_0<r\}$
  lies where $y_0\ne0$, so part (2) computes its measure, and for $t>1$ and small $r$ part (3) puts its
  preimage under $p$ between $\{|(u,z)|<r/t\}$ and $\{|(u,z)|<tr\}$, whose volumes are
  $t^{\mp(2M-1)}\operatorname{vol}\{|(u,z)|<r\}$ by part (1).
\item[(iii)] By (i), (ii) and part (7),
  $(\nu\times m')\{\|y\|_0+\|g'\|<r\}$ divided by
  $(\operatorname{vol}\times\lambda')\{|(u,z)|+|X'|_1<r\}$ tends to $2^{M-1}a'$. By part (4) the
  denominator is $2^{1-M}\lambda(B_r)$, so $(\nu\times m')\{\|y\|_0+\|g'\|<r\}/\lambda(B_r)\to a'$.
\item[(iv)] Let $\sigma$ be the section of part (5). For $f$ the indicator of $\{\|g\|<r\}$ the constant
  $\kappa$ gives $(\nu\times m')\{\|\sigma(y)\iota_0(g')\|<r\}=\kappa\,m\{\|g\|<r\}$. Let $t>1$, let $\delta$
  be as in part (6), and let $r<\min(\delta,1)/2$. For $\nu$-almost every $y$, $q(y)=\beta$. If
  $\|\sigma(y)\iota_0(g')\|<r$, then $\|y\|_0=\|\sigma(y)\iota_0(g')e_0\|_0<2r$ by part (6), so $y_0\ne-1$ and
  $\|y\|_0<\delta$, and part (6) with $g=\sigma(y)$ gives $\|y\|_0+\|g'\|<tr$. If $\|y\|_0+\|g'\|<r/t$, then
  $\|y\|_0<2r$ and part (6) gives $\|\sigma(y)\iota_0(g')\|<r$. So, up to null sets,
  \[
   \{\|y\|_0+\|g'\|<r/t\}\subset\{\|\sigma(y)\iota_0(g')\|<r\}\subset\{\|y\|_0+\|g'\|<tr\},
  \]
  and with (iii) and $\lambda(B_{sr})=s^{M^2}\lambda(B_r)$, $\kappa\,m\{\|g\|<r\}/\lambda(B_r)\to a'$.
\end{enumerate}
By (i) the last quotient also tends to $\kappa a$, so $\kappa a=a'$.

The value of $\kappa$. By definition $a$ and $a'$ are the densities of the Haar probability measures on
$U(M)$ and $U(M-1)$, and $a\ne0$ by Lemma~\ref{cv-volume-mass-haar}(1). The same argument for
$(P,Q)=(M,0)$ and $\beta=1$ gives a constant $\kappa_0$ with $\kappa_0a=a'$, so $\kappa=\kappa_0$. For
$(M,0)$, taking $f=1$ gives $\kappa_0\,m^{\mathrm c}_{M,0}(U(M))=dx^1\{q=1\}\cdot m^{\mathrm c}_{M-1,0}(U(M-1))$,
and both groups have measure $1$; so $\kappa_0=c_M$ by (2).
\end{proof}
% Lean: realPlaceMeasures = RealPlaceCompare.typings_agree with selfDualReal, SphereMassProof.sphereMass and
% CompactDualQuot.compactDualQuotient; MeasurePreservingGlue.quotient_iff is the inversion. quotFormula
% does the reduction (exists_unitaryBlockEmbedding_onto_stabiliser, orthogonalComplement_unitary_embedding_conj_block,
% lintegral_mul_right_eq_self); formula_base gives kappa from invariant_measure_quadric_mul_haar_stabilizer and
% hermitianNormFibreMeasure_finite_invariant, then const_mul_density twice (for (P,Q) and for (M,0) through
% CDHaar.trick) and positiveDefiniteUnitSphereMass_proof; const_mul_density is (i) tendsto_gauge, (ii)
% tendsto_fib, (iii) unitaryLie_volume_ball_split_basis_vector with tendsto_sublevel_density_prod_measure, and
% (iv) quadricPointGauge_section_mul_block_compare with ptGauge_act_le.

\begin{proposition}\label{cv-volume-mass-ratio}
Let $E$ be a CM field with totally real subfield $F$, $f=[F:\mathbf Q]$, such that no finite prime
of $F$ ramifies in $E$, and put $v_E=|d_E|^{1/(2f)}/2$. Let $W$ be a Hermitian space over $E$ with
correct signature (Definition~\ref{th-spaces}), of rank $M\ge3$ and signature $(P,Q)$ at every
complex embedding. Let $\beta=1$ and $(P',Q')=(P-1,Q)$ if $P\ge1$, or $\beta=-1$ and
$(P',Q')=(P,Q-1)$ if $Q\ge1$. Then there is $x\in W$ with $h(x,x)=\beta$ such that $x^\perp$ carries a
unimodular lattice, has correct signature $(P',Q')$ at every complex embedding, and
(Definition~\ref{cv-volume-mass-normalized})
\[
 \frac{\operatorname{vol}(W)}{\operatorname{vol}(x^\perp)}
 =\Big((M-1)!\Big(\frac{v_E}\pi\Big)^M\Big)^f(2v_E)^{-f/2}\,\Lambda_E(M),
\]
the right side of Theorem~\ref{cv-siegel-mass-ratio} for $K=E$.
\end{proposition}
\leanref{Development/AbelianVarietyTheory/Proofs/GlobalInterpolation/VolumeRatioOfUnfolding.lean}{GlobalTheory.VolRatioUnfold.volumeRatio_of_unfolding}
\leanref{Development/AbelianVarietyTheory/Proofs/GlobalInterpolation/VolumeRatioOfUnfolding.lean}{GlobalTheory.VolRatioUnfold.ratio_of_setting}
\leanref{Formalism/SharpRestriction/GlobalSchurBounds/ThickPlacesCorollaryOfAnyTranslate.lean}{Formalism.MainNotions.StandardLattice.volumeMass_of_unfolding}
\begin{proof}
Take $x$ and the unimodular lattice $L\ni x$ of Lemma~\ref{cv-volume-mass-unit-vector}, and put
$W'=x^\perp$ with the lattice $L\cap x^\perp$. By Lemma~\ref{cv-volume-mass-invariance}(2) the
volume of $W$ does not change when its lattice is replaced by $L$. Let $\psi$ be the character of
Lemma~\ref{cv-volume-mass-euler-product}, and $\mathcal F$ a fundamental domain of $U(W)(F)$.

Step 1: a test function. For each real place $\rho$ identify $W\otimes_{F,\rho}\mathbf R$ with
$\mathbf C^M$ so that $h$ becomes $h_{P,Q}$, and let $x_\rho$ be the image of $x$. Choose a continuous
$\varphi_\rho\ge0$ with compact support in $\mathbf C^M\smallsetminus\{0\}$ and
$\varphi_\rho(x_\rho)>0$, and let $\Phi_\infty=\prod_\rho\varphi_\rho$. Since $\psi_\rho(t)=e^{-2\pi it}$,
Lemma~\ref{cv-volume-mass-archimedean}(1) shows that the self-dual measures of
Theorem~\ref{cv-volume-mass-siegel-weil} at $\rho$ are $dx$ and the Lebesgue measure, so the fibre
measure $dx_\infty^\beta$ is the product of the measures $dx^\beta$ of
Lemma~\ref{cv-volume-mass-archimedean}. Then $a=\int\Phi_\infty\,dx_\infty^\beta$ is positive and
finite. Let $\Phi$ be $\Phi_\infty$ times the indicator of $\widehat L$.

Step 2: unfolding. At each $\rho$, Lemma~\ref{cv-volume-mass-archimedean}(3) says that in the
coordinates $(g',y)\mapsto\iota(g')s(y)$ on $U(P,Q)$ the measure $m^{\mathrm c}_{P,Q}$ is
$c_M^{-1}\,dm^{\mathrm c}_{P',Q'}\,dx^\beta$. So the measure $\nu$ of
Lemma~\ref{cv-volume-mass-unfolding} is carried by $H_\infty g\mapsto g^{-1}x$ to
$c_M^{-f}dx_\infty^\beta$, and $\int_{H_\infty\backslash G_\infty}\Phi_\infty(g^{-1}x)\,d\nu=c_M^{-f}a$.
By Lemma~\ref{cv-volume-mass-unfolding},
\[
 \int_{\mathcal F}\sum_{y\in W,\ h(y,y)=\beta}\Phi(g^{-1}y)\,d\mu_W(g)=\operatorname{vol}(W')\,c_M^{-f}a.
\]

Step 3: Siegel's formula. By Theorem~\ref{cv-volume-mass-siegel-weil}, in the coordinates of a basis
of $W$ and with $\mu=\mu_W$, the same integral is
$\operatorname{vol}(W)\cdot a\cdot\prod_v\int\mathbf 1_{L_v}\,dx_v^\beta$, where $L_v$ is the completion
of $L$. By Lemma~\ref{cv-volume-mass-local-density}, as $\beta=\pm1$ is a unit at every $v$, and by
Lemma~\ref{cv-volume-mass-euler-product}, the product is $(2v_E)^{f(1/2-M)}\Lambda_E(M)^{-1}$.

Step 4: the ratio. Comparing Steps 2 and 3 and dividing by $a$,
\[
 \frac{\operatorname{vol}(W')}{\operatorname{vol}(W)}
 =c_M^f\,(2v_E)^{f(1/2-M)}\,\Lambda_E(M)^{-1}.
\]
Since $c_M^{-1}(2v_E)^M=(M-1)!\,(v_E/\pi)^M$, the reciprocal is the stated formula.
\end{proof}

\begin{theorem}\label{cv-volume-mass}
Let $D$ be a definite datum (Definition~\ref{th-definite-datum}) with host fields
$F\subset E$, $f=[F:\mathbf Q]$, and let $V$ be a Hermitian space over $E$ of dimension
$r\ge1$, with signature $(P,Q)$ at every real place $\rho$ of $F$, containing a self-dual
$\mathcal O_E$-lattice. Let $K_V\subset U(V)(\mathbf A_{F,f})$ be the stabilizer of that
lattice, and fix isomorphisms $U(V)(F_\rho)\cong U(P,Q)$. On $U(P,Q)$ let $m$ be the Haar
measure whose density at the identity is $2^{-r(r-1)/2}$ with respect to the Lebesgue
measure $\lambda$ on its Lie algebra $\mathfrak u(P,Q)$ for which $\lambda\times\lambda$
corresponds to the Lebesgue measure on $M_r(\mathbf C)$ under $(X,Y)\mapsto X+iY$. Put
\[
 \operatorname{vol}U(r)=\prod_{i=0}^{r-1}\frac{2\pi^{i+1}}{i!}.
\]
Let $\mu$ be a Haar measure on $U(V)(\mathbf A_F)$ and $\mathcal F$ a measurable fundamental
domain for $U(V)(F)$. Then for every measurable $B\subset\prod_\rho U(P,Q)$,
\[
 \mu(\mathcal F)\prod_\rho m(B_\rho)
 =\mu\{h:h_f\in K_V,\ h_\infty\in B\}\cdot\operatorname{vol}U(r)^f\cdot\operatorname{mass}(I_r).
\]
\end{theorem}
\leanref{Development/WeilBound/RepresentationTheoryGL/TubeMeasures/DisjointTubes/ProductMapEquality/PiCdEqMap.lean}{MainStatements.Tam.vol}
\leanref{Development/WeilBound/RepresentationTheoryGL/TubeMeasures/DisjointTubes/TubeSt.lean}{MainStatements.Tam.VolSt}
\leanref{Development/AbelianVarietyTheory/Proofs/GlobalInterpolation/UnitaryVolumeEqRankMassAllSignatures.lean}{GlobalTheory.Glob.vol_eq_massI_all_sig}
\leanref{Formalism/SharpRestriction/GlobalSchurBounds/ThickPlacesCorollaryOfAnyTranslate.lean}{Formalism.MainNotions.StandardLattice.massRatio_of_definiteVolume}
\leanref{Formalism/SharpRestriction/GlobalSchurBounds/ThickPlacesCorollaryOfAnyTranslate.lean}{Formalism.MainNotions.StandardLattice.volume_eq_mass_of_definite}
\leanref{Development/AbelianVarietyTheory/Proofs/GlobalInterpolation/VolEqMassSelfDualRankGeTwo.lean}{GlobalTheory.Glob.vol_eq_massI_step}
\leanref{Development/AbelianVarietyTheory/Proofs/GlobalInterpolation/VolumeRatioUpPositiveLine.lean}{GlobalTheory.VolRatioUp.volRatio_up}

\begin{proof}
The host fields satisfy the hypotheses of the lemmas above: $E$ is a CM field with totally real
subfield $F$, and no finite prime ramifies in $E/F$ (Proposition~\ref{hf-host-field-properties}).
For a Hermitian space $W$ over $E$ with correct signature (Definition~\ref{th-spaces}) let
$\operatorname{vol}(W)$ be its volume (Definition~\ref{cv-volume-mass-normalized}), and for $M\ge3$
let $R(M)$ be the right side of Proposition~\ref{cv-volume-mass-ratio}. Steps 1 to 3 show that
$\operatorname{vol}(W)=\operatorname{mass}(I_{P+Q})$ whenever $W$ has signature $(P,Q)$ at every
complex embedding and $P+Q\ge1$.

Step 1: the definite signatures and the ratio of masses. If $Q=0$, then
$\operatorname{vol}(W)=\operatorname{mass}(I_P)$ by Lemma~\ref{cv-volume-mass-definite}(1). If
$(P,Q)=(0,1)$, then $\operatorname{vol}(W)=\operatorname{vol}(W^-)=\operatorname{mass}(I_1)$ by
Lemma~\ref{cv-volume-mass-definite}(2) and then (1) for $W^-$. For $M\ge3$, apply
Proposition~\ref{cv-volume-mass-ratio} to $E^M$ with the form $\sum_i\bar x_iy_i$ and the lattice
$\mathcal O_E^M$, with $\beta=1$. The space and the complement $x^\perp$ are positive definite, so
by Lemma~\ref{cv-volume-mass-definite}(1)
\[
 \frac{\operatorname{mass}(I_M)}{\operatorname{mass}(I_{M-1})}=R(M)\qquad(M\ge3).
\]
This is Theorem~\ref{cv-siegel-mass-ratio} for $K=E$, obtained here from the volumes.

Step 2: $P\ge2$, by induction on $Q$. The case $Q=0$ is Step 1. Let $W$ have signature $(P,Q+1)$.
Its rank is at least $3$, and Proposition~\ref{cv-volume-mass-ratio} with $\beta=-1$ gives $x$ with
$x^\perp$ of signature $(P,Q)$ and $\operatorname{vol}(W)=R(P+Q+1)\operatorname{vol}(x^\perp)$. By
induction $\operatorname{vol}(x^\perp)=\operatorname{mass}(I_{P+Q})$, and by Step 1
$\operatorname{vol}(W)=\operatorname{mass}(I_{P+Q+1})$.

Step 3: $P\le1$ and $P+Q\ge2$, by adding a positive line. Let $U$ be $E^{P+Q+1}$ with the form
$\sum_{i\le P+1}\bar x_iy_i-\sum_{i>P+1}\bar x_iy_i$ and the unimodular lattice
$\mathcal O_E^{P+Q+1}$; it has signature $(P+1,Q)$ and rank at least $3$.
Proposition~\ref{cv-volume-mass-ratio} with $\beta=1$ gives $x\in U$ with $x^\perp$ of signature
$(P,Q)$ and $\operatorname{vol}(U)=R(P+Q+1)\operatorname{vol}(x^\perp)$. By
Theorem~\ref{cv-volume-mass-landherr} the space $x^\perp$ is isometric to $W$, so
$\operatorname{vol}(x^\perp)=\operatorname{vol}(W)$ by Lemma~\ref{cv-volume-mass-invariance}(3). If
$P=1$, then $\operatorname{vol}(U)=\operatorname{mass}(I_{P+Q+1})$ by Step 2. If $P=0$, then $Q\ge2$
and $U$ has signature $(1,Q)$, so the case $P=1$, applied to $U$, gives the same. In both cases
Step 1 gives $\operatorname{vol}(W)=\operatorname{mass}(I_{P+Q+1})/R(P+Q+1)=\operatorname{mass}(I_{P+Q})$.

Steps 1 to 3 cover every $(P,Q)$ with $P+Q\ge1$.

Step 4: the measures. Let $V$, $K_V$, $m$, $\mu$, $\mathcal F$ and $B$ be as in the statement, and
write $m^f$ and $(m^{\mathrm c}_{P,Q})^f$ for the product measures on $\prod_\rho U(P,Q)$, so that
$m^f(B)$ is the product written $\prod_\rho m(B_\rho)$ in the statement. By
Lemma~\ref{cv-volume-mass-invariance}(1) the normalized measure of $V$ is $\mu_V=c^{-1}\mu$ for some
$c>0$; then $\mathcal F$ is a fundamental domain for $\mu_V$, and
$\mu_V(\mathcal F)=\operatorname{vol}(V)=\operatorname{mass}(I_r)$. By the definition of $\mu_V$,
$\mu\{h:h_f\in K_V,\ h_\infty\in B\}=c\,(m^{\mathrm c}_{P,Q})^f(B)$, and by
Lemma~\ref{cv-volume-mass-haar}(2), $m=\operatorname{vol}U(r)\cdot m^{\mathrm c}_{P,Q}$. So
\[
 \mu(\mathcal F)\,m^f(B)=c\operatorname{mass}(I_r)\operatorname{vol}U(r)^f(m^{\mathrm c}_{P,Q})^f(B)
 =\mu\{h:h_f\in K_V,\ h_\infty\in B\}\cdot\operatorname{vol}U(r)^f\cdot\operatorname{mass}(I_r).
\]
\end{proof}

\begin{remark}\label{cv-volume-mass-literature}
For $Q=0$, Theorem~\ref{cv-volume-mass} is related to the mass formula of Gan, Hanke and Yu
\cite{GHY01}, which computes the mass of a lattice in a definite Hermitian space as a product of local
factors and rests on Prasad's volume formula \cite{Pra89}.
\end{remark}

\begin{definition}\label{cv-level-sizes}
Let $D$ be a definite datum (Definition~\ref{th-definite-datum}) with ranks $k$, $m$, $t$, $f=2n$
and $r=k+m-1$ (Definition~\ref{th-ranks}), host fields $F\subset E$, level primes $g_0,g_1$,
auxiliary prime $d$, Harish-Chandra parameters $\lambda(\tau)=(u,v,b,e)$ for $\tau\in\mathcal T$
(Definition~\ref{th-hc-parameter}), and spaces $W_j$ of signature $(p_j,q_j)=(k+j-1,m-j)$ for
$1\le j\le m$ (Definition~\ref{th-spaces}), so that $p_j+q_j=r$. For $\tau\in\mathcal T$ put
$\ell_a=v_a$ for $1\le a\le k-1$ and $\ell_a=-b_{a-k+1}$ for $k\le a\le r$, the entries of
$\lambda(\tau)$, and let the \emph{weight} $\kappa^\lambda_j(\tau)\in\mathbf Z^r$ of $W_j$ at $\tau$ be
\[
 \kappa^\lambda_j(\tau)_a=\Big\lfloor\ell_a+\frac{q_j}2-\frac{p_j-1}2+a-1\Big\rfloor,\qquad
 \kappa^\lambda_j(\tau)_{p_j+c}=\Big\lfloor\ell_{p_j+c}-\frac{p_j}2-\frac{q_j-1}2+c-1\Big\rfloor
\]
for $1\le a\le p_j$ and $1\le c\le q_j$. Let $\chi_d$ be the quadratic Dirichlet character modulo $d$,
given by the Legendre symbol, $L(s,\chi_d)$ its $L$-function, and
$\operatorname{rd}_E=|d_E|^{1/[E:\mathbf Q]}$ the root discriminant of $E$. A real number $s$ is a
\emph{size bound} for $D$ if:
\begin{enumerate}
\item $fr(r+t)\le\exp(s)$;
\item $\sum_{\tau\in\mathcal T}\sum_{a=1}^r|\kappa^\lambda_j(\tau)_a+e_j(\tau)|\le\exp(2s)$ for
  $1\le j\le m$;
\item $\log(g_0g_1)\le30s$;
\item $\log\operatorname{rd}_E+\big|\operatorname{Re}\frac{L'}{L}(1,\chi_d)\big|\le2s$.
\end{enumerate}
\end{definition}
\leanref{Development/AbelianVarietyTheory/BlockTheory/DivisorBase/OneMeasureSizes.lean}{GlobalTheory.IndefiniteRowSideFullHecke.OneMeasureSizes.Sizes}
\leanref{Development/AbelianVarietyTheory/BlockTheory/DivisorBase/OneMeasureSizes.lean}{GlobalTheory.IndefiniteRowSideFullHecke.OneMeasureSizes.wtSize}
\leanref{Formalism/LocalToGlobal/HeckeThetaElements/DefsWt.lean}{Formalism.L2.NTheta.wtOf}
\leanref{Formalism/LocalToGlobal/HeckeThetaElements/DefsWt.lean}{Formalism.L2.NTheta.kapQ}
\leanref{Formalism/LocalRingCardinality/DataParams₂.lean}{Formalism.BasicNotions.dlogL}

\begin{proposition}\label{cv-level-volume}
Let $D$ be a definite datum (Definition~\ref{th-definite-datum}) with host fields
$F\subset E$, $f=[F:\mathbf Q]$, and level primes $g_0,g_1$. Let $V$ be a Hermitian space over
$E$ of dimension $r\ge2$, with signature $(P,Q)$ at every real place $\rho$ of $F$, containing a
self-dual $\mathcal O_E$-lattice with stabilizer $K_V\subset U(V)(\mathbf A_{F,f})$, and fix
isomorphisms $U(V)(F_\rho)\cong U(P,Q)$. Let $K\subset K_V\subset K_0$ be compact open subgroups of
$U(V)(\mathbf A_{F,f})$ with $[K_0:K]\le(g_0g_1)^{2fr^2}$. Let $m$ be the Haar measure on $U(P,Q)$
with density $2^{-r(r-1)/2}$ at the identity (Theorem~\ref{cv-volume-mass}), let $\mu$ be the Haar
measure on $U(V)(\mathbf A_F)$ with $\mu\{h:h_f\in K,\ h_\infty\in B\}=\prod_\rho m(B_\rho)$ for every
measurable $B\subset\prod_\rho U(P,Q)$, and let $\mathcal F$ be a measurable fundamental domain for
$U(V)(F)$. Then
\[
 \mu(\mathcal F)\ge\exp\Big(-3\big(fr^2(\log(fr+3)+1)+\log|d_E|\big)\Big).
\]
If moreover $r$ is the integer $r$ of the rank datum of $D$ (Definition~\ref{th-ranks}), and $s$ is
a size bound for $D$ (Definition~\ref{cv-level-sizes}), then
\[
 \log\mu(\mathcal F)\le fr^2\log(r+3)+119fr^2s.
\]
\end{proposition}
\leanref{Development/WeilBound/RepresentationTheoryGL/TubeMeasures/DisjointTubes/BoxExistence/ExistsBox.lean}{MainStatements.Tam.low}
\leanref{Development/WeilBound/RepresentationTheoryGL/TubeMeasures/DisjointTubes/ExponentialThreeBound/LeExpThree.lean}{MainStatements.Tam.A.up_core}
\begin{proof}
Choose a measurable box $B\subset\prod_\rho U(P,Q)$ with $0<\prod_\rho m(B_\rho)<\infty$.
Put $X=\operatorname{vol}U(r)^f\operatorname{mass}(I_r)$. By Theorem~\ref{cv-volume-mass},
$\mu(\mathcal F)\prod_\rho m(B_\rho)=\mu\{h:h_f\in K_V,\ h_\infty\in B\}\cdot X$.

Lower bound. Since $K\subset K_V$, the right side is at least $\prod_\rho m(B_\rho)\cdot X$, so
$\mu(\mathcal F)\ge X$. From the product defining it, $r^{-r^2}\le\operatorname{vol}U(r)\le13^{r^2}$:
each factor is at least $1/i!\ge r^{-r}$, and the product is at most
$2^r\pi^{r(r+1)/2}\le13^{r^2}$. With Lemma~\ref{cv-mass-bounds} and $[E:\mathbf Q]=2f$,
$X\ge r^{-fr^2}3^{-2fr^2}$, which is at least the stated bound.

Upper bound. The group $K_0$ is the disjoint union of $n=[K_0:K]$ cosets $kK$, and by
invariance of $\mu$ each set $\{h:h_f\in kK,\ h_\infty\in B\}$ has measure
$\prod_\rho m(B_\rho)$. As $K_V\subset K_0$, the right side is at most
$n\prod_\rho m(B_\rho)\cdot X$, so $\mu(\mathcal F)\le nX$. Three of the conditions on $s$ in
Definition~\ref{cv-level-sizes} are used, in these forms. By (1), as $r\le r+t$,
$fr^2\le fr(r+t)\le\exp(s)$. By (3), $\log(g_0g_1)\le30s$. By (4), as the second term there is
nonnegative, $\log\operatorname{rd}_E\le2s$. Condition (2) on the weights is not used, and neither is
the bound that (4) gives for $|\operatorname{Re}\frac{L'}{L}(1,\chi_d)|$. Now
$\log n\le2fr^2\log(g_0g_1)\le60fr^2s$, $\log\operatorname{vol}U(r)^f\le fr^2\log13$, and by
Lemma~\ref{cv-mass-bounds} $\log\operatorname{mass}(I_r)\le6fr^2\log(12f^2r^4\operatorname{rd}_E^2)$.
Since $f\ge2$ and $r\ge2$, $fr^2\le\exp(s)$ gives $s\ge1$, and with $\log\operatorname{rd}_E\le2s$ the
three terms add up to at most $117fr^2s$. This is at most $fr^2\log(r+3)+119fr^2s$.
% Lean: up_core takes OneMeasureSizes.Sizes D L and uses its first, third and fourth parts (h1, h3, h4);
% hfr is f r^2 <= e^L from r <= nV, hrdL drops |dlogL D 1|, and the second part, on wtSize, is discarded.
\end{proof}

\begin{proposition}\label{cv-class-number-formula}
Let $D$ be a definite datum (Definition~\ref{th-definite-datum}) with ranks $k$ and $f$,
host fields $F\subset E$ and level primes $g_0,g_1$. Suppose that
$g_i>2fk+1$ for some $i$, and that $g_0$ and $g_1$ split completely in $F$. Let $u$ be a
word of $D$ whose auxiliary prime $T_u$ splits completely in $F$, and let
$\operatorname{Cl}_u$ be the set of lattice classes at the bottom of the chain of $u$.
Then
\[
 |\operatorname{Cl}_u|=\operatorname{mass}(I_k)\prod_{i=0,1}
 \Big(\frac{|GL_k(\mathbf F_{g_i})|}{|N_k(\mathbf F_{g_i})|}\Big)^{f}
 \,(T_u^k-1)^{f},
\]
where $I_k$ is the standard lattice over $E$ (Definition~\ref{cv-mass}) and
$N_k\subset GL_k$ is the subgroup of upper unitriangular matrices.
\end{proposition}
\begin{proof}
There is a unimodular lattice $M_0$ in the bottom space $X_0$ of rank $k$ whose stabilizer
in the finite adeles contains the level group $U_u$ of $u$. Unimodular lattices of rank
$k$ form one genus with the standard lattice, so $\operatorname{mass}(M_0)=\operatorname{mass}(I_k)$.
The index of $U_u$ in the stabilizer of $\widehat M_0$ is a product of local indices. At
each of the $f$ places of $F$ above $g_i$ the group is $GL_k(\mathbf F_{g_i})$ modulo its
upper unitriangular subgroup, and at each of the $f$ places above $T_u$ it is the
stabilizer of a nonzero vector in $\mathbf F_{T_u}^k$, of index $T_u^k-1$. The mass is the
sum over the double cosets $U(X_0)(F)\backslash U(X_0)(\mathbf A_f)/\operatorname{Stab}(\widehat M_0)$
of the reciprocals of the orders of the global stabilizers $\Gamma_y$. Let $\gamma\in U(X_0)(F)$
lie in a conjugate $yU_uy^{-1}$. Since $X_0$ is definite, $\gamma$ has finite order. At a
place above $g_i$ the group $U_u$ is unipotent modulo $g_i$, so the order of $\gamma$ is a
power of $g_i$. If $\gamma\ne1$, a primitive $g_i$-th root of unity is an eigenvalue of
$\gamma$, whose characteristic polynomial over $\mathbf Q$ has degree $[E:\mathbf Q]k=2fk$;
this needs $g_i-1\le2fk$. So when $g_i>2fk+1$ every such $\gamma$ is trivial, the double
coset of $y$ for $\operatorname{Stab}(\widehat M_0)$ splits into
$[\operatorname{Stab}(\widehat M_0):U_u]/|\Gamma_y|$ double cosets for $U_u$, and summing
over $y$ gives $|\operatorname{Cl}_u|=[\operatorname{Stab}(\widehat M_0):U_u]\operatorname{mass}(M_0)$.
\end{proof}
\leanref{Formalism/SharpRestriction/GlobalSchurBounds/StandardLattice.lean}{Formalism.MainNotions.StandardLattice.ClassNumberFormula₂}
\leanref{Development/AbelianVarietyTheory/Proofs/FoundationProofs/ClassNumberAll.lean}{GlobalTheory.ColClassNumber.classNumber_all}
\leanref{Development/AbelianVarietyTheory/Proofs/FoundationProofs/MassSelfDualLatticeFirstSpaceEq.lean}{GlobalTheory.ColClassNumber.classNumber}

\begin{remark}\label{cv-where-mass-enters}
Theorem~\ref{cv-siegel-mass-ratio} and Proposition~\ref{cv-class-number-formula} are not
used in the proof of the covolume bound of Theorem~\ref{cv-covolume-bound}, and neither is
the Rallis inner product formula. Masses do enter that proof, through
Theorem~\ref{cv-volume-mass}. Proposition~\ref{cv-level-volume} bounds the volume of a
fundamental domain of each $U(W_j)$, normalized by its level subgroup, above and below by
way of $\operatorname{mass}(I_r)$ and Lemma~\ref{cv-mass-bounds}. The norm of a lifted form
is its Petersson norm times the square root of that volume times the dimensions of the
archimedean types, so these bounds give the comparison of norms in
Proposition~\ref{cv-point-conditions}(5) and the bound for $\mu_{\max}$ in (4). The proof of
Theorem~\ref{cv-volume-mass} passes from the definite signature to the others through the ratio of
the volumes of unitary groups of consecutive ranks, which comes from a rank-one case of Siegel's
formula, the representations of a number by a Hermitian form of rank at least $3$; this ratio is
the volume form of Theorem~\ref{cv-siegel-mass-ratio}. So a form of Siegel's mass formula enters
the covolume bound through Theorem~\ref{cv-volume-mass}. The upper bound of
Lemma~\ref{cv-mass-bounds} also appears in the bound for the rank of integral cusp forms in
Proposition~\ref{cv-point-conditions}(4).

The three results above are used on the definite side, and the mass formula enters there
in two places. First, Proposition~\ref{cv-class-number-formula}
gives the number of lattice classes at the bottom of each chain, hence the rank of the
space of algebraic modular forms there. Second, telescoping
Theorem~\ref{cv-siegel-mass-ratio} from rank $k$ to rank $r=k+m-1$ gives
\[
 \frac{\operatorname{mass}(I_r)}{\operatorname{mass}(I_k)}
 =\prod_{M=k+1}^{r}\frac{\operatorname{mass}(I_M)}{\operatorname{mass}(I_{M-1})},
\]
and this ratio bounds below the growth of the space at the top of the chain against the
space at the bottom. For every word $u$, let $d^*_u$ be the dimension of the space of forms
at the bottom of the chain of $u$, given by Proposition~\ref{cv-class-number-formula} and
the dimensions of the archimedean types, and let $D^*$ be a lower bound for the dimension
at the top, given by $\operatorname{mass}(I_r)$ in the same way; they play the role of the
numbers $d^*_u$ and $D^*$ of Definition~\ref{th-definite-datum}. Then
$m\,x^{\varepsilon}\le\log\big((D^*/d^*_u)^{1/(fk)}\Gamma^{-2/(fk)}\big)$ for an
explicit level factor $\Gamma$. This lower bound enters the choice of the families of
definite-side data (Chapter~\ref{chapter-assembly}).
\end{remark}
\leanref{Development/AbelianVarietyTheory/Proofs/GlobalInterpolation/WindowLower.lean}{GlobalTheory.EX.WinLow.window_lower}
\leanref{Development/AutomorphicArithmetic/CharacterBoxes/ClassSel.lean}{UpperTheory.Top.classSel}

\section{The Chowla--Selberg formula and the period $\Omega$}\label{section-cv-chowla-selberg}

\subsection{The Chowla--Selberg formula and constant sections}\label{subsection-cv-chowla-selberg-constant-sections}

\begin{theorem}\label{cv-chowla-selberg}
Let $K$ be the imaginary quadratic field of discriminant $-d$ with $d>4$, let
$\varphi:K\to\mathbf C$ be an embedding, and let $\chi$ be the Kronecker character of $K$,
the Dirichlet character modulo $d$ with $\chi(n)=\big(\frac{-d}{n}\big)$ for odd $n$. Let
$h_K$ be the class number of $K$, and for each ideal class $c$ let $\mathfrak a_c$ be a
fractional ideal in $c$.
\begin{enumerate}
\item Let $\tau_c$ be a point of the upper half-plane with
  $\varphi(\mathfrak a_c)$ homothetic to $\mathbf Z+\mathbf Z\tau_c$, and let $\eta$ be
  Dedekind's eta function. Then
  \[
   \frac1{h_K}\sum_c\log\big(4\pi\operatorname{Im}(\tau_c)\,|\eta(\tau_c)|^4\big)
   =\gamma+\log\frac{2\pi}{\sqrt d}-\frac{L'}{L}(1,\chi).
  \]
\item Suppose that $d$ is prime, so that $\chi$ is the Legendre symbol modulo $d$. Let
  $\Delta(\Lambda)=g_2(\Lambda)^3-27g_3(\Lambda)^2$ be the discriminant of a lattice
  $\Lambda\subset\mathbf C$. Then
  \[
   \frac1{h_K}\sum_c\log\big(N\mathfrak a_c^{6}\,|\Delta(\varphi(\mathfrak a_c))|\big)
   =12\log2\pi-6\log d+6\gamma-6\frac{L'}{L}(1,\chi).
  \]
\end{enumerate}
\end{theorem}
\leanref{Development/AbelianVarietyTheory/Proofs/HomogeneousIdeals/DedekindZetaEqSumRealEisenstein.lean}{GlobalTheory.chowla_selberg_eta}
\leanref{Development/AbelianVarietyTheory/Proofs/DivisorProofs/ClassMeanLogNormDiscChowlaSelbergStatement.lean}{GlobalTheory.T.classMean_log_norm_disc_chowlaSelberg}

Part~(1) is the Chowla--Selberg formula, a consequence of Kronecker's limit formula
\cite{CS67}. For prime $d$ the two parts are equivalent, since
$|\Delta(\mathbf Z+\mathbf Z\tau)|=(2\pi)^{12}|\eta(\tau)|^{24}$, $\Delta$ is homogeneous
of degree $-12$, and $N\mathfrak a$ is the area of $\varphi(\mathfrak a)$ divided by
$\sqrt d/2$.

\begin{definition}\label{cv-log-omega}
Let $E$ be a number field with root discriminant $\operatorname{rd}_E=|d_E|^{1/[E:\mathbf Q]}$,
and let $\chi$ be a Dirichlet character with $L(1,\chi)\ne0$. Put
\[
 \log\Omega_E(\chi)=\gamma+\log\frac{2\pi}{\operatorname{rd}_E}-\frac{L'}{L}(1,\chi),
\]
a complex number, which is real when $\chi$ is real.
\end{definition}
\leanref{Development/AbelianVarietyTheory/Definitions/CMLattices/ExistsIsThetaDegOfIsLocalModel.lean}{GlobalTheory.CM.logOm}
\leanref{Formalism/LocalToGlobal/Divisibility/Numbers.lean}{Formalism.L2.DIV.NumDatum.logOm}

\begin{definition}\label{cv-constant-sections}
Let $E$ be a CM field with a CM type $\Phi$, a set of embeddings $E\to\mathbf C$ containing one
from each pair of complex conjugate embeddings, and let a PEL datum of signature
$(p,q)$ over $E$, a level, its moduli scheme $X$ over its coefficient field $F_c$, and
uniformizations of $X$ at the complex embeddings of $F_c$ be given; $F_c$ contains the images
of $E$ under the embeddings in $\Phi$. Let $Z$ be a connected component of $X$, let
$\lambda$ be a weight, and let $\delta$ be a real function on the finite places of $F_c$. A
\emph{constant section} of weight $\lambda$ and defect $\delta$ on $Z$ is a rational section
$s$ of the automorphic line bundle $\omega^\lambda$ of weight $\lambda$, supported on $Z$,
together with a rational section $s'$ of $\omega^{-\lambda}$ with $ss'=1_Z$, for which there
are numbers $c_\sigma(s)>0$ and a finitely supported function $\rho_s$ on the finite places
of $F_c$ such that:
\begin{enumerate}
\item for every complex embedding $\sigma$ of $F_c$, the norm of $s$ in the uniformization at
  $\sigma$ equals $c_\sigma(s)$ at every complex point of $Z$;
\item for every finite place $w$ of $F_c$, the valuation of $s$ at every integral point of
  $Z$ over a complete discrete valuation ring above $w$ equals $\delta(w)+\rho_s(w)$.
\end{enumerate}
Its \emph{degree} is
\[
 \deg s=\frac{1}{[F_c:\mathbf Q]}\Big(\sum_w\rho_s(w)\log Nw-\sum_{\sigma:F_c\to\mathbf C}\log c_\sigma(s)\Big).
\]
For $\tau\in\Phi$ let $\lambda_\tau$ be the weight equal to $-1$ at $\tau$ and $0$ at the other
elements of $\Phi$, and let $\delta_\tau(w)=q\big(v_w(\mathfrak d_\tau)-v_w(d)\big)$ for a prime $d$, where
$\mathfrak d_\tau$ is the ideal of $\mathcal O_{F_c}$ generated by the image of the different of
$E$ under the embedding of $E$ in $F_c$ given by $\tau$. A \emph{flat section} of type $\tau$
on $Z$ is a constant section of weight $\lambda_\tau$ and defect $\delta_\tau$.
\end{definition}
\leanref{Development/AbelianVarietyTheory/Definitions/Foundations/ExistsFlatSecRankZero.lean}{GlobalTheory.FlatLine.ConstSec}
\leanref{Development/AbelianVarietyTheory/Definitions/Foundations/ExistsFlatSecRankZero.lean}{GlobalTheory.FlatLine.FlatSec}
\leanref{Development/AbelianVarietyTheory/Definitions/Foundations/ExistsFlatSecRankZero.lean}{GlobalTheory.FlatLine.lineDeg}

The defect sums to a multiple of the root discriminant: since the norm of the different of
$E$ is $|d_E|$,
\[
 \frac{1}{[F_c:\mathbf Q]}\sum_w\delta_\tau(w)\log Nw=q\big(\log\operatorname{rd}_E-\log d\big).
\]

\begin{lemma}\label{cv-constant-sections-proportional}
Let a PEL datum over a CM field $E$, its moduli scheme $X$ over its coefficient field $F_c$, and
uniformizations of $X(\mathbf C)$ at the complex embeddings of $F_c$ be given, with $F_c$
algebraically closed in the function field of every connected component of $X$, and let $Z$ be
a connected component (Definition~\ref{cv-constant-sections}). If $s$ and $t$ are constant sections of the same
weight and the same defect on $Z$, then $t=as$ for some $a\in F_c^\times$, and
$\deg t=\deg s$.
\end{lemma}
\leanref{Development/AbelianVarietyTheory/Proofs/FoundationProofs/IsAlgebraicOfNormConst.lean}{GlobalTheory.FlatLine.ConstSec.exists_eq_smul}
\leanref{Development/AbelianVarietyTheory/Proofs/FoundationProofs/IsAlgebraicOfNormConst.lean}{GlobalTheory.FlatLine.ConstSec.deg_eq}
\begin{proof}
The product $u=ts'$ is a constant section of weight $0$ and defect $0$ on $Z$: a rational
function supported on $Z$ whose absolute value at the complex points of $Z$, for a fixed
embedding $\sigma_0$ of $F_c$, is a constant $b$. Let $\Omega$ be the function field of $Z$ and
$g\in\Omega$ the value of $u$ at the generic point. Every $F_c$-embedding
$P:\Omega\to\mathbf C$ is evaluation at a complex point of $Z$, so $|P(g)|=b$ for all $P$.

The element $g$ is algebraic over $F_c$. Otherwise $g$ belongs to a transcendence basis $T$ of
$\Omega$ over $F_c$. Fix $P_0$. The numbers $2P_0(y)$, $y\in T$, are algebraically independent
over $F_c$, so there is an $F_c$-embedding $P$ with $P(y)=2P_0(y)$ on $T$. Then
$b=|P(g)|=2|P_0(g)|=2b$, so $b=0$ and $g=0$, which is algebraic.

Since $F_c$ is algebraically closed in $\Omega$, $g=a$ for some $a\in F_c$. A rational function
is determined by its values at the generic points of the components, so $u=a1_Z$. Then
$t=t\,s's=as$, and $a\ne0$ because $t\ne0$. Finally $c_\sigma(as)=|\sigma(a)|c_\sigma(s)$ and
$\rho_{as}(w)=\rho_s(w)+v_w(a)$, so
$\deg(as)-\deg s=\frac1{[F_c:\mathbf Q]}\big(\sum_wv_w(a)\log Nw-\sum_\sigma\log|\sigma(a)|\big)=0$
by the product formula.
\end{proof}

\begin{theorem}\label{cv-gaga}
Let $X$ be a scheme, proper and smooth over $\mathbf C$, and let $L$ be a line bundle on $X$, trivialized
on an open cover $(U_i)$ with transition functions $g_{ij}$ on $U_i\cap U_j$. For each $i$ let $h_i$
be a complex function on the complex points of $U_i$ that is holomorphic in every holomorphic
family of such points, and suppose that $h_i=g_{ij}h_j$ at every complex point of $U_i\cap U_j$.
Then there is a global section $s$ of $L$ whose value in the trivialization over $U_i$ is $h_i$ at
every complex point of $U_i$, for every $i$.
\end{theorem}
\leanref{Development/AutomorphicSpectrum/PublishedTheorems/Geometry₂/GagaSections.lean}{GlobalTheory.CompanionTheory.Ar.Flat.GagaSections}
\leanref{Development/AutomorphicSpectrum/Proofs/GlobalProofs/GagaSections.lean}{GlobalTheory.CompanionTheory.Gaga.gagaSections}

This is the statement on global sections of Serre's GAGA \cite{Ser56}, in the form for proper
schemes of \cite[Expos\'e~XII, Th\'eor\`eme~4.4]{Gro71}.

\begin{lemma}\label{cv-resolution-coordinates}
Let $K$ be a field of characteristic zero, let $M,D\ge0$, let $X\subset\mathbf P^M_K$ be a reduced
closed subscheme, and let $U\subset X$ be a dense open subscheme that is smooth of relative dimension
$D$ over $K$. Then there are $r\ge0$, a closed subscheme $X'\subset\mathbf P^M_K\times\mathbf P^r_K$
and a morphism $f:X'\to X$ that is the restriction of the first projection, such that:
\begin{enumerate}
\item $X'$ is smooth of relative dimension $D$ over $K$;
\item $f$ is an isomorphism over $U$, and $f^{-1}(U)$ is dense in $X'$;
\item $X'$ is covered by finitely many affine open subschemes $V$, on each of which there are
  functions $z_1,\dots,z_D$ whose differentials form a basis of $\Omega^1_{V/K}$ and an integer
  $k\le D$ such that $V\setminus f^{-1}(U)$ is the set of points of $V$ at which one of
  $z_1,\dots,z_k$ vanishes.
\end{enumerate}
\end{lemma}
\leanref{ReciprocityTheory/ReductiveGroups/ArithmeticGeometry/BirationalGeometry/BlowupProjections/BlowupProjStatement.lean}{ReciprocityTheory.ChartRestriction.G.Geo.ResK.ResKStatement}
\leanref{Development/WeilBound/CoreStatements/ResolutionCharZero/WeakPair.lean}{PositivityTheory.ResPlug.resK_of_resolution_of_principalisation}
\begin{proof}
By Theorem~\ref{cv-functorial-resolution}, the functorial form of Hironaka's theorem \cite{Hir64}
(compare Corollary~\ref{lit-resolution} and \cite[Tag~R02G]{Res26}), there is a finite sequence of
blow-ups $s:X_1\to X$ with $X_1$ smooth over $K$, $s$ an isomorphism over the smooth locus of $X$,
which contains $U$, and the inverse image of the smooth locus dense in $X_1$. A blow-up of a
projective scheme is projective, so $X_1$ is projective over $K$. Since $U$ is dense in the smooth
locus, $s^{-1}(U)$ is dense in $X_1$, and $X_1$ has relative dimension $D$ because
$s^{-1}(U)\cong U$ does. The closed subset $X_1\setminus s^{-1}(U)$ contains no irreducible
component of $X_1$. Hironaka's principalization, Theorem~\ref{lit-principalization} with
$E=\emptyset$ (compare \cite[Tag~R021]{Res26}), applied on $X_1$ to the ideal of this subset, gives
a sequence of blow-ups in smooth centres $t:X_2\to X_1$ with $X_2$ smooth and projective over $K$,
$t$ an isomorphism over $s^{-1}(U)$, and the reduced inverse image of $X_1\setminus s^{-1}(U)$ a
simple normal crossings divisor. The complement of a simple normal crossings divisor is dense,
which gives (2) for $f=s\circ t$, and (1) follows as for $X_1$. At a closed point $x$ of $X_2$, a
regular system of parameters of which some cut out the divisor has differentials that form a basis
of $\Omega^1_{X_2/K}$ near $x$, since $X_2$ is smooth over $K$ and the residue field of $x$ is finite
and separable over $K$. These neighbourhoods cover $X_2$, and quasi-compactness gives the finite
cover of (3). Finally $X_2$ is projective over $X$, so it is a
closed subscheme of $\mathbf P^r_X\subset\mathbf P^M_K\times\mathbf P^r_K$ for some $r$,
compatibly with $f$.
\end{proof}

\subsection{Periods of CM points}\label{subsection-cv-cm-points}

\begin{definition}\label{cv-cm-point-periods}
Let $E$ be a CM field with complex conjugation $b\mapsto\bar b$ and different $\mathfrak d_E$, let $d$
be a prime with $\sqrt{-d}\in E$, and let $\theta\in\mathbf C$ with $\theta^2=-d$. Let $L$ be a
finitely generated projective $\mathcal O_E$-module of rank $p+q$ with a perfect alternating form
$\langle\ ,\ \rangle:L\times L\to\mathbf Z$ such that $\langle bx,y\rangle=\langle x,\bar by\rangle$
for $b\in\mathcal O_E$.
\begin{enumerate}
\item The \emph{$E$-valued form} of $L$ is the form $\langle x,y\rangle_E\in E$ with
  $\operatorname{Tr}_{E/\mathbf Q}(b\langle x,y\rangle_E)=\langle bx,y\rangle$ for all $b\in\mathcal O_E$.
  The \emph{Gram matrix} of a family $u=(u_1,\dots,u_n)$ in $L$ is
  $G(u)=(\langle u_m,u_{m'}\rangle_E)_{m,m'}$. The family is a \emph{local basis at} a prime $\ell$ if
  it is linearly independent over $\mathcal O_E$ and $NL\subset\sum_m\mathcal O_Eu_m$ for an integer
  $N$ prime to $\ell$.
\item A \emph{type of signature $(p,q)$ for $\theta$} is a map $k\mapsto\sigma_k$ from
  $\{1,\dots,g\}$ to the embeddings $E\to\mathbf C$ that takes each embedding $\sigma$ with
  $\sigma(\sqrt{-d})=\theta$ as a value $p$ times and each other embedding $q$ times. For
  $z\in\mathbf C^g$ and an embedding $\sigma$, the \emph{$\sigma$-part} $z_\sigma\in\mathbf C^g$ has
  the coordinates $z_k$ with $\sigma_k=\sigma$ and $0$ at the others.
\item A \emph{period} of $L$ of this type is an injective additive map $P:L\to\mathbf C^g$ with
  $P(bx)_k=\sigma_k(b)P(x)_k$ for $b\in\mathcal O_E$, whose image is a lattice in $\mathbf C^g$, and
  for which some positive definite Hermitian form $H$ on $\mathbf C^g$ has
  $\operatorname{Im}H(P(x),P(y))=\langle x,y\rangle$. Its \emph{pure sublattices} are
  \[
   L_\pm(P)=\{x\in L:\ P(x)_k=0\text{ whenever }\sigma_k(\sqrt{-d})\ne\pm\theta\}.
  \]
\item Let $A=\mathbf C^g/\Lambda$ be a complex abelian variety with the polarization of a positive
  definite Hermitian form $H$ on $\mathbf C^g$. The \emph{Hodge Gram matrix} of tangent vectors
  $V_1,\dots,V_n\in\mathbf C^g$ of $A$ at the origin is $(\pi H(V_m,V_{m'}))_{m,m'}$, and $\pi H(V,V)$ is
  the \emph{squared Hodge norm} of $V$.
\item A \emph{complex structure of signature $(p,q)$ for $\theta$} on $L$ is an $\mathbf R$-linear map
  $J$ of $L\otimes\mathbf R$ with $J^2=-1$ that commutes with $\mathcal O_E$, such that
  $\langle Jx,Jy\rangle=\langle x,y\rangle$ for all $x,y$, $\langle Jx,x\rangle>0$ for $x\ne0$, and each
  $b\in\mathcal O_E$ has trace $\sum_\sigma\big(p\,\sigma(b)+q\,\overline{\sigma(b)}\big)$ on the complex
  vector space $L\otimes\mathbf R$ with $i$ acting as $J$, the sum over the embeddings $\sigma$ with
  $\sigma(\sqrt{-d})=\theta$.
\end{enumerate}
\end{definition}
\leanref{ReciprocityTheory/ReductiveGroups/ArithmeticGeometry/BirationalGeometry/ToroidalCompactification/TwoWeightEqOneAddRoot.lean}{ReciprocityTheory.ChartRestriction.G.Geo.Tor.Datum}
\leanref{Development/AutomorphicSpectrum/Definitions/GlobalGeometry/IsHodgeGram.lean}{GlobalTheory.CompanionTheory.Ar.Fl.sesq}
\leanref{Development/AutomorphicSpectrum/Definitions/GlobalGeometry/IsTransition.lean}{GlobalTheory.CompanionTheory.Ar.Fl.gramE}
\leanref{Development/AutomorphicSpectrum/Definitions/GlobalGeometry/IsTransition.lean}{GlobalTheory.CompanionTheory.Ar.Fl.IsLocBasis}
\leanref{ReciprocityTheory/ReductiveGroups/ArithmeticGeometry/BirationalGeometry/ToroidalCompactification/IsHolPeriods.lean}{ReciprocityTheory.ChartRestriction.G.Geo.Tor.IsTypeFn}
\leanref{Development/AutomorphicSpectrum/Definitions/GlobalGeometry/IsHodgeGram.lean}{GlobalTheory.CompanionTheory.Ar.Fl.tyPart}
\leanref{ReciprocityTheory/ReductiveGroups/ArithmeticGeometry/BirationalGeometry/ToroidalCompactification/IsHolPeriods.lean}{ReciprocityTheory.ChartRestriction.G.Geo.Tor.IsPeriod}
\leanref{Development/AbelianVarietyTheory/Statements/PrincipalLattices/FiniteDimensionalBase.lean}{GlobalTheory.PureType.lat}
\leanref{Development/AutomorphicSpectrum/Definitions/GlobalGeometry/IsHodgeGram.lean}{GlobalTheory.CompanionTheory.Ar.Fl.IsHodgeGram}
\leanref{Development/AbelianVarietyTheory/Statements/PrincipalLattices/LogHermFormOfUnitDiscriminant.lean}{GlobalTheory.IsLieNormSq}

\begin{lemma}\label{cv-cm-point-gram}
Let $E$, $d$, $\theta$ and $L$ be as in Definition~\ref{cv-cm-point-periods}. For a prime $\ell$ and
an embedding $\iota:E\to\mathbf C_\ell$, put
$\log|\mathfrak d_E|_\iota=\log\max_{b\in\mathfrak d_E}|\iota(b)|$, for the absolute value of
$\mathbf C_\ell$ with $|\ell|=\ell^{-1}$. Let $e=(e_1,\dots,e_{p+q})$ and $u=(u_1,\dots,u_{p+q})$ be
families in $L$, with $e$ linearly independent over $\mathcal O_E$.
\begin{enumerate}
\item The form $\langle x,y\rangle_E$ is $\mathcal O_E$-linear in $x$, conjugate-linear in $y$,
  satisfies $\langle y,x\rangle_E=-\overline{\langle x,y\rangle_E}$, and takes values in
  $\mathfrak d_E^{-1}$.
\item $\det G(e)\ne0$, and $e$ is a local basis at all but finitely many primes.
\item There is a unique matrix $T\in M_{p+q}(E)$ with
  $\langle u_m,y\rangle_E=\sum_nT_{mn}\langle e_n,y\rangle_E$ for all $m$ and all $y\in L$. It
  satisfies $G(u)=T\,G(e)\,\bar T^{\mathsf t}$, and $\det T\ne0$ if $u$ is linearly independent over
  $\mathcal O_E$.
\item If $u$ is a local basis at $\ell$, then
  $\log|\iota(\det G(u))|=-(p+q)\log|\mathfrak d_E|_\iota$ for every embedding
  $\iota:E\to\mathbf C_\ell$. If $u$ and $e$ are local bases at $\ell$, then
  $|\iota(\det T)|=1$ for every such $\iota$.
\item Let $P$ be a period of $L$ of a type of signature $(p,q)$ for $\theta$. Then
  $\langle x,y\rangle_E=\langle y,x\rangle_E=0$ for $x\in L_+(P)$ and $y\in L_-(P)$. If $u$ is a
  local basis at $\ell$ with $u_1,\dots,u_p\in L_+(P)$ and $u_{p+1},\dots,u_{p+q}\in L_-(P)$, then
  $G(u)$ is block diagonal with the blocks $G_+(u)=(\langle u_a,u_b\rangle_E)_{a,b\le p}$ and
  $G_-(u)=(\langle u_{p+i},u_{p+i'}\rangle_E)_{i,i'\le q}$, and
  $\log|\iota(\det G_-(u))|=-q\log|\mathfrak d_E|_\iota$ for every embedding
  $\iota:E\to\mathbf C_\ell$.
\end{enumerate}
\end{lemma}
\leanref{Development/AutomorphicSpectrum/Proofs/AlgebraProofs/SesqSesquilinearSkewHermitianStatement.lean}{GlobalTheory.T.sesq_sesquilinear_skew_hermitian}
\leanref{Development/AutomorphicSpectrum/Proofs/AlgebraProofs/GramDetNeZeroLocalBasisCofiniteStatement.lean}{GlobalTheory.T.gram_det_ne_zero_local_basis_cofinite}
\leanref{Development/AutomorphicSpectrum/Proofs/AlgebraProofs/ExistsTransitionMatrixGramStatement.lean}{GlobalTheory.T.exists_transition_matrix_gram}
\leanref{Development/AutomorphicSpectrum/Proofs/AlgebraProofs/GramDeterminantLogDifferentStatement.lean}{GlobalTheory.T.gram_determinant_log_different}
\leanref{Development/AutomorphicSpectrum/Proofs/AlgebraProofs/PadicNormTransitionDetOneStatement.lean}{GlobalTheory.T.padic_norm_transition_det_one}
\leanref{Development/AutomorphicSpectrum/Proofs/AlgebraProofs/PureTypeLatticesOrthogonalStatement.lean}{GlobalTheory.T.pure_type_lattices_orthogonal}
\leanref{Development/AutomorphicSpectrum/Proofs/AlgebraProofs/LocalBasisGramDetLogNormStatement.lean}{GlobalTheory.T.local_basis_gram_det_log_norm}
\begin{proof}
(1) For $b,c\in\mathcal O_E$, $\operatorname{Tr}(b\langle cx,y\rangle_E)=\langle bcx,y\rangle
=\operatorname{Tr}(bc\langle x,y\rangle_E)$ and
$\operatorname{Tr}(b\langle x,cy\rangle_E)=\langle\bar cbx,y\rangle=\operatorname{Tr}(b\bar c\langle x,y\rangle_E)$.
Since $\langle\ ,\ \rangle$ is alternating,
$\operatorname{Tr}(b\langle y,x\rangle_E)=-\langle x,by\rangle=-\langle\bar bx,y\rangle
=-\operatorname{Tr}(\bar b\langle x,y\rangle_E)=-\operatorname{Tr}(b\overline{\langle x,y\rangle_E})$.
The trace form of $E$ is nondegenerate, which gives the three identities. Finally
$\operatorname{Tr}(b\langle x,y\rangle_E)\in\mathbf Z$ for all $b\in\mathcal O_E$, so
$\langle x,y\rangle_E\in\mathfrak d_E^{-1}$.

(2) The family $e$ is a basis of $L\otimes\mathbf Q$ over $E$, since $L$ has rank $p+q$. If
$c\in E^{p+q}$ and $\sum_mc_m\langle e_m,e_n\rangle_E=0$ for all $n$, then $x=\sum_mc_me_m$
satisfies $\langle x,y\rangle_E=0$, hence $\langle x,y\rangle=0$, for all $y\in L\otimes\mathbf Q$,
so $x=0$ because $\langle\ ,\ \rangle$ is perfect, and $c=0$. So $\det G(e)\ne0$. Every $y\in L$
has a multiple $cy$ with $c\in\mathcal O_E\setminus0$ in $\sum_m\mathcal O_Ee_m$, hence also the
multiple $N_{E/\mathbf Q}(c)\,y$. Since $L$ is finitely generated, one integer $N\ge1$ has
$NL\subset\sum_m\mathcal O_Ee_m$, and $e$ is a local basis at every $\ell\nmid N$.

(3) Write $u_m=\sum_nT_{mn}e_n$ in $L\otimes\mathbf Q$; by (1) this $T$ has the stated property, and
it is the only one, since $\det G(e)\ne0$. By (1),
$\langle u_m,u_{m'}\rangle_E=\sum_{n,n'}T_{mn}\overline{T_{m'n'}}\langle e_n,e_{n'}\rangle_E$. If $u$ is
linearly independent, it is also a basis of $L\otimes\mathbf Q$ over $E$, and $T$ is invertible.

(4) Let $O'$ be the localization of $\mathcal O_E$ at the primes above $\ell$, a principal ideal
domain, and let $\delta$ generate $\mathfrak d_EO'$; then
$|\iota(\delta)|=\max_{b\in\mathfrak d_E}|\iota(b)|$ for every $\iota$. The module
$L'=L\otimes\mathbf Z_{(\ell)}$ is free over $O'$ with basis $u$, and $\delta\langle\ ,\ \rangle_E$
takes values in $O'$ on it by (1). This form is perfect: composing with the trace identifies
$\operatorname{Hom}_{O'}(L',\mathfrak d_E^{-1}O')$ with
$\operatorname{Hom}_{\mathbf Z_{(\ell)}}(L',\mathbf Z_{(\ell)})$, and $\langle\ ,\ \rangle$ is perfect.
Hence $\delta G(u)$ is invertible over $O'$, and
$(p+q)\log|\iota(\delta)|+\log|\iota(\det G(u))|=0$. If $e$ is also a local basis at $\ell$, then
$T$ and $T^{-1}$ have entries in $O'$, so $\iota(\det T)$ is a unit.

(5) Let $x\in L_+(P)$ and $y\in L_-(P)$. Then $P(\sqrt{-d}x)=\theta P(x)$ and
$P(\sqrt{-d}y)=-\theta P(y)$, so
$\langle\sqrt{-d}x,\sqrt{-d}y\rangle=\operatorname{Im}H(\theta P(x),-\theta P(y))=-d\langle x,y\rangle$,
while $\langle\sqrt{-d}x,\sqrt{-d}y\rangle=\langle x,dy\rangle=d\langle x,y\rangle$. Hence
$\langle x,y\rangle=0$. Since $L_+(P)$ is an $\mathcal O_E$-submodule, also $\langle bx,y\rangle=0$ for
all $b\in\mathcal O_E$, so $\langle x,y\rangle_E=0$, and $\langle y,x\rangle_E=0$ by (1). For $u$ as
stated, $G(u)$ is therefore block diagonal, and so is $\delta G(u)$ in the proof of (4). Its blocks
are invertible over $O'$, which gives $\log|\iota(\det G_-(u))|=-q\log|\iota(\delta)|$.
\end{proof}

\begin{definition}\label{cv-cm-point-curve}
Let $d\equiv3\pmod4$ be a prime, let $\mathcal O=\mathbf Z[w_0]$ with $w_0=(1+\sqrt{-d})/2$ be the
ring of integers of $\mathbf Q(\sqrt{-d})$, and let $\theta\in\mathbf C$ with $\theta^2=-d$. A
\emph{CM curve with root $\theta$} consists of:
\begin{enumerate}
\item a number field $M\subset\mathbf C$ containing $\theta$ and a principally polarized elliptic
  curve $C$ over $\mathcal O_M$ with an action of $\mathcal O$ whose Rosati involution is complex
  conjugation;
\item a basis $v$ of the Lie algebra of $C$ on which $\sqrt{-d}$ acts by $\theta$;
\item a uniformization $C(\mathbf C)=\mathbf C/c(\mathbf Z+\mathbf Z\frac{1+i\sqrt d}2)$ of the
  complex curve at the inclusion $M\subset\mathbf C$, with $c\in\mathbf C^\times$, and the number
  $z\in\mathbf C$ for which $v$ is the tangent vector $cz$ in this uniformization.
\end{enumerate}
Then $z\ne0$. For an embedding $s:M\to\mathbf C$, $n(s)>0$ is the squared Hodge norm of the image of
$v$ in the complex curve at $s$ (Definition~\ref{cv-cm-point-periods}).
\end{definition}
\leanref{Development/AutomorphicSpectrum/Definitions/GlobalGeometry/IsCMFrame.lean}{GlobalTheory.CompanionTheory.Ar.Fl.Curve0}
\leanref{Development/AbelianVarietyTheory/Definitions/CMLattices/IsDiffOf.lean}{GlobalTheory.CM.Ell}
\leanref{Development/AutomorphicSpectrum/Proofs/GeometryProofs/CurveZNeZeroStatement.lean}{GlobalTheory.T.curve_z_ne_zero}

\begin{lemma}\label{cv-cm-point-models}
Let $K$ be an imaginary quadratic field, let $\sigma_0:K\to\mathbf C$ be an embedding, and for each
ideal class $C$ of $K$ let $\mathfrak b_C$ be a nonzero ideal of $\mathcal O_K$ in $C$. For a lattice
$\Lambda\subset\mathbf C$ write $\Delta(\Lambda)=g_2(\Lambda)^3-27g_3(\Lambda)^2$. Then there are a
number field $M\subset\mathbf C$ containing $\sigma_0(K)$, with inclusion $\iota:M\to\mathbf C$ and
with $\kappa:K\to M$ the embedding such that $\iota\circ\kappa=\sigma_0$, and:
\begin{enumerate}
\item principally polarized elliptic curves $P$ and $Q_C$, one for each class $C$, over
  $\mathcal O_M$, with bases $v$ of the Lie algebra of $P$ and $t_C$ of that of $Q_C$;
\item for each embedding $s:M\to\mathbf C$, lattices $L_0(s)$ and $L_C(s)$ in $\mathbf C$ such that
  the complex curve of $P$ at $s$ is $\mathbf C/L_0(s)$, with the image of $v$ the tangent vector $1$,
  and the complex curve of $Q_C$ at $s$ is $\mathbf C/L_C(s)$, with the image of $t_C$ the tangent
  vector $1$;
\item units $\delta_C\in\mathcal O_M^\times$ with $\Delta(L_C(s))=s(\delta_C)$ for every $s$;
\item a number $c\in\mathbf C^\times$ with $L_0(\iota)=c\,\sigma_0(\mathcal O_K)$;
\item for each $a\in\mathcal O_K$, an endomorphism of the generic fibre $P_M$ that multiplies $v$ by
  $\kappa(a)$;
\item elements $\beta_C\in M^\times$ such that for every $x\in K$ with $x\mathfrak b_C\subset\mathcal O_K$
  there is a homomorphism $P_M\to Q_{C,M}$ carrying $v$ to $\beta_C\kappa(x)t_C$, and for every
  $y\in\mathfrak b_C$ there is a homomorphism $Q_{C,M}\to P_M$ carrying $t_C$ to
  $\beta_C^{-1}\kappa(y)v$.
\end{enumerate}
\end{lemma}
\leanref{Development/AbelianVarietyTheory/Proofs/CubicalLineBundles/ExistsNumberFieldDefiningCmCurveHomsStatement.lean}{GlobalTheory.T.exists_numberField_defining_cm_curve_homs}
\leanref{Development/AbelianVarietyTheory/Proofs/CubicalLineBundles/ExistsCmCurvesOfWeierstrassCurves.lean}{GlobalTheory.exists_cm_curves_of_weierstrass_curves}
\leanref{Development/AbelianVarietyTheory/Proofs/CubicalLineBundles/ExistsEllipticCurvePolarisedLatticeFamilyStatement.lean}{GlobalTheory.T.exists_ellipticCurve_polarised_lattice_family}
\leanref{Development/AbelianVarietyTheory/Proofs/CubicalLineBundles/ExistsFieldCmCurvesUnitDiscriminant.lean}{GlobalTheory.exists_field_cm_curves_unit_discriminant}
\begin{proof}
For a fractional ideal $\mathfrak a$ of $K$ the lattice $\sigma_0(\mathfrak a)$ has complex
multiplication by $\mathcal O_K$, so its $j$-invariant is an algebraic integer
\cite[Chapter~II, \S6]{Sil94}; it depends only on the class of $\mathfrak a$, so there are finitely
many of them. An elliptic curve over a number field with integral $j$-invariant has potentially good
reduction at every prime \cite[Proposition~VII.5.5]{Sil09}, so over a finite extension it has a
Weierstrass model with integral coefficients and discriminant $1$ \cite[\S VIII.8]{Sil09}. Let $M$ be
generated over $\sigma_0(K)$ by the coefficients of such models for the classes of $\mathcal O_K$ and
of the $\mathfrak b_C^{-1}$. A Weierstrass curve over $\mathcal O_M$ with unit discriminant is an
elliptic curve over $\mathcal O_M$ with its principal polarization, a basis of whose Lie algebra is
dual to the invariant differential; at an embedding $s$ its complex curve is $\mathbf C/L(s)$ with
$g_2(L(s))=s(c_4)/12$, $g_3(L(s))=s(c_6)/216$ and $\Delta(L(s))=s(\Delta)$, and the invariant
differential is $dz$ \cite[Chapter~VI]{Sil09}. This gives (1), (2) and (3). Lattices with the same
$j$-invariant are homothetic \cite[Chapter~VI]{Sil09}, which gives (4) and numbers $c_C\ne0$ with
$L_C(\iota)=c_C\,\sigma_0(\mathfrak b_C^{-1})$. Put $\beta_C=c_C/c$. Then
$\sigma_0(a)L_0(\iota)\subset L_0(\iota)$ for $a\in\mathcal O_K$,
$\beta_C\sigma_0(x)L_0(\iota)\subset L_C(\iota)$ for $x\mathfrak b_C\subset\mathcal O_K$, and
$\beta_C^{-1}\sigma_0(y)L_C(\iota)\subset L_0(\iota)$ for $y\in\mathfrak b_C$. A number $\alpha$ with
$\alpha L\subset L'$ is the derivative at the origin of a homomorphism $\mathbf C/L\to\mathbf C/L'$
\cite[Theorem~VI.4.1]{Sil09}. Such a homomorphism between the complex curves of curves over $M$ is
defined over a finite extension of $M$ \cite[Chapter~II, \S2]{Sil94}, and the numbers $\alpha$ with
$\alpha L\subset L'$ form a finitely generated group, so one finite extension of $M$ serves for all the
homomorphisms above. Replacing $M$ by it and the curves by their base changes gives (5) and (6); the
homomorphism for $x=1$ shows $\beta_C\in M$.
\end{proof}

\begin{lemma}\label{cv-cm-point-area}
Let $K$, $\sigma_0$, $\mathfrak b_C$, $M$, $\kappa$, $P$, $v$, $L_0$, $L_C$, $\delta_C$ and $\beta_C$ be as
in Lemma~\ref{cv-cm-point-models}, let $h$ be the class number of $K$, and let $A(\Lambda)$ be the area
of a fundamental domain of a lattice $\Lambda\subset\mathbf C$. For an embedding $s:M\to\mathbf C$ let
$n(s)$ be the squared Hodge norm of the image of $v$ in the complex curve of $P$ at $s$
(Definition~\ref{cv-cm-point-periods}).
\begin{enumerate}
\item $n(s)=\pi/A(L_0(s))$.
\item $L_0(s)=c_s\,s\kappa(\mathfrak a_s)$ for a nonzero ideal $\mathfrak a_s$ of $\mathcal O_K$ and a
  number $c_s\ne0$, and $L_C(s)=s(\beta_C)\,s\kappa(\mathfrak b_C^{-1})L_0(s)$; in particular
  $A(L_C(s))\,N\mathfrak b_C=|s(\beta_C)|^2A(L_0(s))$.
\item With $\Lambda_C=\sigma_0(\mathfrak b_C)$,
  \[
   \frac1{[M:\mathbf Q]}\sum_s\log A(L_0(s))=\frac1{6h}\sum_C\log\big(A(\Lambda_C)^6|\Delta(\Lambda_C)|\big).
  \]
\item If $\Lambda=a(\mathbf Z+\mathbf Z\tau)$ with $a\ne0$ and $\operatorname{Im}\tau>0$, then
  $\log\pi-\frac16\log\big(A(\Lambda)^6|\Delta(\Lambda)|\big)=-\log\big(4\pi\operatorname{Im}\tau\,|\eta(\tau)|^4\big)$,
  where $\eta$ is Dedekind's eta function.
\item If $\tau_C$ is a point of the upper half-plane with $\sigma_0(\mathfrak b_C)$ homothetic to
  $\mathbf Z+\mathbf Z\tau_C$ for each class $C$, then
  \[
   \frac1{[M:\mathbf Q]}\sum_s\log n(s)=-\frac1h\sum_C\log\big(4\pi\operatorname{Im}\tau_C\,|\eta(\tau_C)|^4\big).
  \]
\end{enumerate}
\end{lemma}
\leanref{Development/AbelianVarietyTheory/Proofs/CubicalLineBundles/IsLieNormSqIffEqPiDivCovolume.lean}{GlobalTheory.isLieNormSq_iff_eq_pi_div_covolume}
\leanref{Development/AbelianVarietyTheory/Proofs/CubicalLineBundles/LatticeCmEqIdealAndCovolume.lean}{GlobalTheory.lattice_cm_eq_ideal_and_covolume}
\leanref{Development/AbelianVarietyTheory/Proofs/CubicalLineBundles/MeanLogCovolumeEqMeanClassGroup.lean}{GlobalTheory.mean_log_covolume_eq_mean_classGroup}
\leanref{Development/AbelianVarietyTheory/Proofs/CubicalLineBundles/MeanLogCovolumeEqMeanClassGroup.lean}{GlobalTheory.log_pi_sub_log_latticeNormDisc}
\leanref{Development/AbelianVarietyTheory/Proofs/CubicalLineBundles/LogPiSubMeanLogLatticeNormDisc.lean}{GlobalTheory.log_pi_sub_mean_log_latticeNormDisc}
\leanref{Development/AbelianVarietyTheory/Proofs/CubicalLineBundles/MeanLogLieNormSqEqMeanLogEtaStatement.lean}{GlobalTheory.T.mean_log_lieNormSq_eq_mean_log_eta_alt}
\begin{proof}
(1) The principal polarization of $\mathbf C/L$ is the Hermitian form $H(z,z')=z\bar z'/A(L)$, whose
imaginary part is $\pm1$ on a basis of $L$. So $\pi H(1,1)=\pi/A(L)$.

(2) A homomorphism over $M$ that multiplies $v$ by $\kappa(a)$ gives at $s$ a homomorphism of complex
curves with derivative $s\kappa(a)$, so $s\kappa(a)L_0(s)\subset L_0(s)$ for $a\in\mathcal O_K$
(Lemma~\ref{cv-cm-point-models}(5) and \cite[Theorem~VI.4.1]{Sil09}). For $0\ne c'\in L_0(s)$,
$c'^{-1}L_0(s)$ is then the image under $s\kappa$ of a fractional ideal of $K$, and a multiple of it is
an ideal. In the same way Lemma~\ref{cv-cm-point-models}(6) gives
$s(\beta_C)s\kappa(x)L_0(s)\subset L_C(s)$ for $x\mathfrak b_C\subset\mathcal O_K$ and
$s\kappa(y)L_C(s)\subset s(\beta_C)L_0(s)$ for $y\in\mathfrak b_C$. Since
$\mathfrak b_C^{-1}\mathfrak b_C=\mathcal O_K$, the second inclusion gives
$L_C(s)\subset s(\beta_C)s\kappa(\mathfrak b_C^{-1})L_0(s)$, and the first gives the reverse inclusion.
The area of $c'\,s\kappa(\mathfrak a)$ is $|c'|^2N\mathfrak a$ times that of $s\kappa(\mathcal O_K)$, which
gives the area formula.

(3) Put $D(\Lambda)=A(\Lambda)^6|\Delta(\Lambda)|$; since $\Delta$ is homogeneous of degree $-12$,
$D(a\Lambda)=D(\Lambda)$ for $a\ne0$, and $D(\bar\Lambda)=D(\Lambda)$. By (2) and
Lemma~\ref{cv-cm-point-models}(3),
\[
 \log D(L_C(s))=6\big(2\log|s(\beta_C)|+\log A(L_0(s))-\log N\mathfrak b_C\big)+\log|s(\delta_C)|.
\]
The homomorphisms of Lemma~\ref{cv-cm-point-models}(6) extend over $\mathcal O_M$
\cite[Chapter~I, Proposition~2.7]{FC90}, so they carry $v$ and $t_C$ into the Lie algebras over
$\mathcal O_M$; hence $\beta_C\kappa(x)$ and $\beta_C^{-1}\kappa(y)$ are integral for $x\in\mathfrak b_C^{-1}$
and $y\in\mathfrak b_C$, and $\beta_C\mathcal O_M=\mathfrak b_C\mathcal O_M$. Taking norms,
$\sum_s\log|s(\beta_C)|=[M:K]\log N\mathfrak b_C$, and $\sum_s\log|s(\delta_C)|=0$ since $\delta_C$ is a
unit. As $[M:\mathbf Q]=2[M:K]$, summing the display over $s$ gives
$\sum_s\log D(L_C(s))=6\sum_s\log A(L_0(s))$ for each $C$. On the other hand, for fixed $s$ the lattice
$L_C(s)$ is homothetic by (2) to $s\kappa(\mathfrak b_C^{-1}\mathfrak a_s)$, and $s\kappa$ is $\sigma_0$ or
its complex conjugate. The class of $\mathfrak b_C^{-1}\mathfrak a_s$ is that of
$\mathfrak b_{C'}$ with $C'=[\mathfrak a_s]C^{-1}$, and $C'$ runs over all classes once as $C$ does;
so $\sum_C\log D(L_C(s))=\sum_C\log D(\Lambda_C)$. Summing over $C$ and $s$ in the two orders gives
$6h\sum_s\log A(L_0(s))=[M:\mathbf Q]\sum_C\log D(\Lambda_C)$.

(4) $A(\Lambda)=|a|^2\operatorname{Im}\tau$ and $\Delta(\Lambda)=a^{-12}(2\pi)^{12}\eta(\tau)^{24}$
\cite[Chapter~I]{Sil94}, so $\frac16\log D(\Lambda)=\log\operatorname{Im}\tau+2\log2\pi+4\log|\eta(\tau)|$,
and $\log\pi$ minus this is $-\log(4\pi\operatorname{Im}\tau\,|\eta(\tau)|^4)$.

(5) By (1), $\log n(s)=\log\pi-\log A(L_0(s))$. Average over $s$ with (3), and apply (4) to each
$\Lambda_C$, which is homothetic to $\mathbf Z+\mathbf Z\tau_C$.
\end{proof}

\begin{lemma}\label{cv-cm-point-action}
Let $d\equiv3\pmod4$ be a prime, let $\mathcal O=\mathbf Z[w_0]$ with $w_0=(1+\sqrt{-d})/2$, let
$M\subset\mathbf C$ be a number field containing $w=(1+i\sqrt d)/2$, and let $P$ be a principally
polarized elliptic curve over $\mathcal O_M$ with a basis $v$ of its Lie algebra. Suppose that the
generic fibre $P_M$ has an endomorphism $W$ that multiplies $v$ by $w$.
\begin{enumerate}
\item $W^2-W+\frac{1+d}4=0$, and there is an action of $\mathcal O$ on $P$ over $\mathcal O_M$ with
  $w_0$ acting on $P_M$ as $W$; under it $\sqrt{-d}$ acts on $v$ by $i\sqrt d$.
\item For every action of $\mathcal O$ on $P$, the Rosati involution of the principal polarization is
  complex conjugation on $\mathcal O$.
\item Composing the action of (1) with complex conjugation of $\mathcal O$ gives an action of
  $\mathcal O$ on $P$ under which $\sqrt{-d}$ acts on $v$ by $-i\sqrt d$.
\end{enumerate}
\end{lemma}
\leanref{Development/AbelianVarietyTheory/Proofs/CubicalLineBundles/QuadraticEquationOfDerivative.lean}{GlobalTheory.LieMultiple.quadratic_equation_of_derivative}
\leanref{Development/AbelianVarietyTheory/Proofs/CubicalLineBundles/ExistsCmAdjoinActionOfEquationStatement.lean}{GlobalTheory.T.exists_cmAdjoin_action_of_equation}
\leanref{Development/AbelianVarietyTheory/Proofs/CubicalLineBundles/ExistsOActExtendDedekind.lean}{GlobalTheory.CM.exists_oAct_extend_dedekind}
\leanref{Development/AbelianVarietyTheory/Proofs/CubicalLineBundles/CmActionPolarisedCurveAdjointConjStatement.lean}{GlobalTheory.T.cmAction_polarisedCurve_adjoint_conj}
\leanref{Development/AbelianVarietyTheory/Proofs/CubicalLineBundles/ExistsConjAction.lean}{GlobalTheory.CM.EllConj.exists_conj_action}
\begin{proof}
(1) The endomorphism $W^2-W+\frac{1+d}4$ multiplies $v$ by $w^2-w+\frac{1+d}4=0$, and an endomorphism
of an elliptic curve in characteristic zero is determined by its action on the Lie algebra
\cite[\S III.5]{Sil09}; so it is $0$, and $w_0\mapsto W$ defines an action of $\mathcal O$ on $P_M$.
It extends to $P$, since homomorphisms of abelian schemes over the normal base $\mathcal O_M$ extend
from the generic fibre \cite[Chapter~I, Proposition~2.7]{FC90}. Then $\sqrt{-d}=2w_0-1$ multiplies $v$
by $2w-1=i\sqrt d$.

(2) On a complex fibre $\mathbf C/\Lambda$ the polarization is $H(z,z')=z\bar z'/A$, with $A$ the area
of a fundamental domain of $\Lambda$, and $a\in\mathcal O$ acts by multiplication by $\varphi(a)$ for an
embedding $\varphi:\mathcal O\to\mathbf C$. So $H(\varphi(a)z,z')=H(z,\varphi(\bar a)z')$, and the Rosati
involution is $a\mapsto\bar a$ on this fibre. Homomorphisms of abelian schemes over $\mathcal O_M$ are
determined by their complex fibre at the inclusion $M\subset\mathbf C$
\cite[Chapter~I, Proposition~2.7]{FC90}, so the same holds over $\mathcal O_M$.

(3) Complex conjugation is a ring automorphism of $\mathcal O$, and after composing with it $w_0$
acts by $\bar w_0=1-w_0$, which multiplies $v$ by $1-w$; so $\sqrt{-d}$ multiplies $v$ by
$2(1-w)-1=-i\sqrt d$.
\end{proof}

\begin{lemma}\label{cv-cm-point-split}
Let $E$ be a CM field that is unramified over its totally real subfield, and let $d$, $\theta$ and
$L$ be as in Definition~\ref{cv-cm-point-periods}, with $L$ given a complex structure of signature
$(p,q)$ for $\theta$. Fix a type of signature $(p,q)$ for $\theta$, let $\Psi$ be the set of
embeddings $\tau:E\to\mathbf C$ with $\tau(\sqrt{-d})=\theta$, and let $\ell$ be a prime.
\begin{enumerate}
\item $L$ has a period of this type.
\item Let $P$ be a period of this type, $u=(u_1,\dots,u_{p+q})$ a family in $L$ that is linearly
  independent over $\mathcal O_E$, and $\tau\in\Psi$. Then the Hermitian matrix $i\tau(G(u))$ is
  invertible of signature $(p,q)$.
\item There is a local basis $u$ of $L$ at $\ell$ with $\det G(u_1,\dots,u_p)\ne0$ and
  $\log|\iota(\det G(u_1,\dots,u_p))|=-p\log|\mathfrak d_E|_\iota$ for every embedding
  $\iota:E\to\mathbf C_\ell$, in the notation of Lemma~\ref{cv-cm-point-gram}.
\item For $\tau\in\Psi$ let $M_\tau$ be a Hermitian matrix of size $p+q$ and $C_\tau$ a complex
  $p\times(p+q)$ matrix with $C_\tau M_\tau\bar C_\tau^{\mathsf t}$ positive definite. Then for every
  $X_0\in M_{p\times(p+q)}(\mathcal O_E)$ and every integer $n\ge1$ there is
  $Y\in M_{p\times(p+q)}(\mathcal O_E)$ such that, with $X=X_0+nY$, the matrix
  $\tau(X)M_\tau\overline{\tau(X)}^{\mathsf t}$ is positive definite for every $\tau\in\Psi$.
\item If $u_1,\dots,u_p\in L$ satisfy the conditions of (3) on $G(u_1,\dots,u_p)$, there are
  $u_{p+1},\dots,u_{p+q}\in L$ such that $u$ is a local basis at $\ell$ and
  $\langle u_a,u_{p+i}\rangle_E=0$ for $a\le p$ and $i\le q$.
\item Let $u$ be linearly independent over $\mathcal O_E$ with $\langle u_a,u_{p+i}\rangle_E=0$ for
  $a\le p$ and $i\le q$, and suppose that $i\tau(G_+(u))$ and $-i\tau(G_-(u))$ are positive definite
  for every $\tau\in\Psi$, with the blocks of Lemma~\ref{cv-cm-point-gram}(5). Then there is a period
  $P$ of this type with $u_1,\dots,u_p\in L_+(P)$ and $u_{p+1},\dots,u_{p+q}\in L_-(P)$.
\item There are a period $P$ of this type and a local basis $u$ of $L$ at $\ell$ with
  $u_1,\dots,u_p\in L_+(P)$ and $u_{p+1},\dots,u_{p+q}\in L_-(P)$, and then
  $NL\subset L_+(P)+L_-(P)$ for an integer $N\ge1$ prime to $\ell$.
\end{enumerate}
\end{lemma}
\leanref{Development/AbelianVarietyTheory/Proofs/ToricGeometry/ExistsEquivOfSumEq.lean}{GlobalTheory.PeriodNonempty.exists_isPeriod}
\leanref{Development/AutomorphicSpectrum/Proofs/AlgebraProofs/PeriodGramHermitianSignatureStatement.lean}{GlobalTheory.T.period_gram_hermitian_signature}
\leanref{Development/AutomorphicSpectrum/Proofs/AlgebraProofs/LocalBasisGramDeterminantBoundStatement.lean}{GlobalTheory.T.local_basis_gram_determinant_bound}
\leanref{Development/AutomorphicSpectrum/Proofs/AlgebraProofs/ExistsIntegralPosDefPerturbationStatement.lean}{GlobalTheory.T.exists_integral_posDef_perturbation}
\leanref{Development/AutomorphicSpectrum/Proofs/AlgebraProofs/ExistsOrthogonalLocalBasisStatement.lean}{GlobalTheory.T.exists_orthogonal_local_basis}
\leanref{Development/AutomorphicSpectrum/Proofs/AlgebraProofs/GramSignatureOrthogonalSplitStatement.lean}{GlobalTheory.T.gram_signature_orthogonal_split}
\leanref{Development/AutomorphicSpectrum/Proofs/AlgebraProofs/PurePeriodLocalBasisExistsStatement.lean}{GlobalTheory.T.pure_period_local_basis_exists}
\leanref{Development/AutomorphicSpectrum/Proofs/AlgebraProofs/PureFamilyDecompositionStatement.lean}{GlobalTheory.T.pure_family_decomposition}
\begin{proof}
(1) Let $J$ be the complex structure. Since $J$ commutes with $\mathcal O_E$, the ring $\mathcal O_E$
acts linearly on the complex vector space $W=(L\otimes\mathbf R,J)$, and by the trace condition its
characters on $W$ are the $\sigma_k$ of the type; as $E\otimes\mathbf R$ is a product of copies of
$\mathbf C$, this gives an isomorphism $W\cong\mathbf C^g$ under which $b\in\mathcal O_E$ acts on the
$k$-th coordinate by $\sigma_k(b)$. Let $P$ be the composite $L\to W\cong\mathbf C^g$; its image is a
lattice, and $H(x,y)=\langle Jx,y\rangle+i\langle x,y\rangle$ is a positive definite Hermitian form on
$W$ with imaginary part $\langle\ ,\ \rangle$.

(2) Let $H$ be a Hermitian form for $P$. For $b\in\mathcal O_E$,
$\operatorname{Im}H(bz,z')=\operatorname{Im}H(z,\bar bz')$ on $P(L)$, hence on $\mathbf C^g$, so
$H(bz,z')=H(z,\bar bz')$; as $\bar b$ acts on the $k$-th coordinate by $\overline{\sigma_k(b)}$, $H$ is
the orthogonal sum of positive definite forms $H_\sigma$ on the $\sigma$-parts. Put
$h_\sigma(x,y)=H_\sigma(P(x)_\sigma,P(y)_\sigma)$. Then
$\operatorname{Tr}_{E/\mathbf Q}(b\langle x,y\rangle_E)=\langle bx,y\rangle=\operatorname{Im}\sum_\sigma\sigma(b)h_\sigma(x,y)$
for all $b$, and comparing characters gives
$\sigma(\langle x,y\rangle_E)=\big(h_\sigma(x,y)-\overline{h_{\bar\sigma}(x,y)}\big)/2i$. For $\tau\in\Psi$
the matrix $i\tau(G(u))$ is therefore half the difference of a positive semidefinite matrix of rank at
most $p$, the Gram matrix of the $P(u_m)_\tau$ for $H_\tau$, and one of rank at most $q$, the conjugate
of the Gram matrix of the $P(u_m)_{\bar\tau}$ for $H_{\bar\tau}$. It is invertible by
Lemma~\ref{cv-cm-point-gram}(2), so its signature is $(p,q)$.

(3) Let $O'$ be the localization of $\mathcal O_E$ at the primes above $\ell$. Since $E$ is unramified
over its totally real subfield $E^+$, the different of $E$ is generated at these primes by an element
$\delta$ of $E^+$, and $f=\delta\langle\ ,\ \rangle_E$ is a perfect skew-Hermitian form on
$L'=L\otimes\mathbf Z_{(\ell)}$ with values in $O'$, as in the proof of
Lemma~\ref{cv-cm-point-gram}(4). For each maximal ideal $\mathfrak m$ of $O'$ there is $a\in\mathcal O_E$
with $a-\bar a\notin\mathfrak m$, again because $E/E^+$ is unramified. If $f(x,x)\in\mathfrak m$ for all
$x$, then $f(x,y)-\overline{f(x,y)}=f(x,y)+f(y,x)\in\mathfrak m$ for all $x,y$, and taking $f(x,y)=1$
and replacing $x$ by $ax$ gives $a-\bar a\in\mathfrak m$, which is impossible. With the Chinese
remainder theorem this gives $x\in L'$ with $f(x,x)$ a unit of $O'$; then $L'=O'x\oplus x^\perp$, and by
induction $L'$ has an orthogonal basis $u_1,\dots,u_{p+q}$ with every $f(u_m,u_m)$ a unit. After
multiplying by an integer prime to $\ell$, the $u_m$ lie in $L$ and form a local basis at $\ell$, and
$\det G(u_1,\dots,u_p)=\prod_{a\le p}\langle u_a,u_a\rangle_E$ has
$|\iota(\cdot)|=|\iota(\delta)|^{-p}$, with $\log|\iota(\delta)|=\log|\mathfrak d_E|_\iota$.

(4) The set $\Psi$ contains one embedding from each pair of complex conjugates, so
$E\otimes\mathbf R\cong\prod_{\tau\in\Psi}\mathbf C$, and $X_0+nM_{p\times(p+q)}(\mathcal O_E)$ is a
translate of a lattice in $M_{p\times(p+q)}(\mathbf C)^\Psi$; so it meets every ball of some fixed radius
$R$. The set of $(X_\tau)$ with every $X_\tau M_\tau\bar X_\tau^{\mathsf t}$ positive definite is open,
contains $(C_\tau)$ and is stable under multiplication by positive reals; so it contains the ball of
radius $R$ about $t(C_\tau)$ for large $t$, and that ball contains some $(\tau(X))_\tau$.

(5) With $O'$ and $\delta$ as in (3), the determinant of $\delta G(u_1,\dots,u_p)$ has absolute value $1$
at every $\iota$, so this matrix is invertible over $O'$. Hence $L'$ is the direct sum of the span of
$u_1,\dots,u_p$ and its orthogonal complement for $\langle\ ,\ \rangle_E$, which is free of rank $q$
over the principal ideal domain $O'$. A basis of it, multiplied by an integer prime to $\ell$, gives
$u_{p+1},\dots,u_{p+q}$.

(6) The family $u$ is a basis of $L\otimes\mathbf Q$ over $E$; write $x=\sum_mc_mu_m$ with $c_m\in E$.
For $\tau\in\Psi$ let the $\tau$-part of $P(x)$ be $(\tau(c_a))_{a\le p}$ and its $\bar\tau$-part be
$(\bar\tau(c_{p+i}))_{i\le q}$, in the order of the coordinates of the type. Then
$P(bx)_k=\sigma_k(b)P(x)_k$, $P$ extends to an isomorphism $L\otimes\mathbf R\to\mathbf C^g$, and the
$u_a$ and $u_{p+i}$ lie in $L_+(P)$ and $L_-(P)$. Let $H$ be the orthogonal sum of the forms with
matrices $2i\tau(G_+(u))$ on the $\tau$-parts and the conjugate of $-2i\tau(G_-(u))$ on the
$\bar\tau$-parts; it is positive definite, and since $G(u)$ is block diagonal the computation in (2),
read backwards, gives $\operatorname{Im}H(P(x),P(y))=\operatorname{Tr}_{E/\mathbf Q}\langle x,y\rangle_E
=\langle x,y\rangle$.

(7) By (3) there is a local basis $x$ at $\ell$ with the conditions of (3), and by (1) a period $P_0$. By
(2) for $P_0$, for each $\tau\in\Psi$ some $C_\tau$ makes $C_\tau\,i\tau(G(x))\bar C_\tau^{\mathsf t}$
positive definite, and (4) with $X_0=(1\mid0)$ and $n=\ell$ gives $X$ such that the vectors
$x'_a=\sum_mX_{am}x_m$, $a\le p$, have $i\tau(G(x'_1,\dots,x'_p))$ positive definite for all
$\tau\in\Psi$. Since $\delta G(x)$ has entries in $O'$ by Lemma~\ref{cv-cm-point-gram}(1) and
$x'_a\equiv x_a$ modulo $\ell L$, the determinant of $\delta G(x'_1,\dots,x'_p)$ is congruent to that of
$\delta G(x_1,\dots,x_p)$ modulo $\ell O'$, so it is a unit and $x'_1,\dots,x'_p$ satisfy the conditions of
(3). By (5) they extend to a local basis $u$ at $\ell$ with $\langle u_a,u_{p+i}\rangle_E=0$. Then
$i\tau(G(u))$ is block diagonal with positive definite first block, and by (2) it has signature
$(p,q)$, so $-i\tau(G_-(u))$ is positive definite. Now (6) gives $P$. Finally, if
$NL\subset\sum_m\mathcal O_Eu_m$ with $\ell\nmid N$, then $NL\subset L_+(P)+L_-(P)$ because
$L_\pm(P)$ are $\mathcal O_E$-submodules.
\end{proof}

\begin{lemma}\label{cv-cm-point-descent}
Let $E$ be a CM field and $d\equiv3\pmod4$ a prime with $\sqrt{-d}\in E$; let
$\mathcal O=\mathbf Z[w_0]\subset\mathcal O_E$ with $w_0=(1+\sqrt{-d})/2$ and $w=(1+i\sqrt d)/2$. Let a
PEL datum of signature $(p,q)$ over $E$, a level, its moduli scheme $X$ and its coefficient field
$F_c\subset\mathbf C$ be given, with $\sigma(E)\subset F_c$ for every embedding $\sigma$. Let
$\tau:E\to\mathbf C$ be an embedding, $\theta=\tau(\sqrt{-d})$, and let $L$ be as in
Definition~\ref{cv-cm-point-periods} with a complex structure of signature $(p,q)$ for $\theta$. Fix a
type of signature $(p,q)$ for $\theta$, and suppose that to each period $P$ of $L$ of this type are
attached a complex point $y_P$ of $X$ over $F_c\subset\mathbf C$ and a surjective holomorphic
homomorphism $\exp_P:\mathbf C^g\to A_{y_P}(\mathbf C)$ with kernel $P(L)$, compatible with the actions
of $\mathcal O_E$ and with the polarizations. Let $Z_+$ and $Z_-$ be CM curves with roots $\theta$ and
$-\theta$ over number fields $M_\pm$ (Definition~\ref{cv-cm-point-curve}), with complex curves
$\mathbf C/c_\pm(\mathbf Z+\mathbf Zw)$. Let $P$ be a period of this type with
$NL\subset L_+(P)+L_-(P)$ for an integer $N\ge1$, let $\varepsilon_\pm\in\{\pm1\}$ be the sign of
$\pm\operatorname{Im}\theta$, and for $b,b'\in L$ put
\[
 h_\pm(b,b')=\varepsilon_\pm\big(\langle w_0b,b'\rangle-\langle b,b'\rangle+\langle b,b'\rangle w_0\big)\in\mathcal O.
\]
\begin{enumerate}
\item For $b\in L_\pm(P)$ the map $\psi_\pm(b):c_\pm t\mapsto\exp_P(tP(b))$ from the complex curve of
  $Z_\pm$ to $A_{y_P}(\mathbf C)$ is a well-defined homomorphism. It is additive in $b$, satisfies
  $\psi_\pm(sb)=s\circ\psi_\pm(b)$ for $s\in\mathcal O_E$, commutes with the actions of $\mathcal O$, and
  the polarization of $A_{y_P}$ pulls back along the pair $(\psi_\pm(b),\psi_\pm(b'))$ to that of
  $Z_\pm$ composed with the action of $h_\pm(b,b')$.
\item For every number field $k_0\subset\mathbf C$ there are a number field $K\subset\mathbf C$
  containing $k_0$, $F_c$, $M_+$ and $M_-$, and a point $y\in X(K)$ over $F_c\subset K$ whose image
  under $K\subset\mathbf C$ is $y_P$, such that $A_y$, with its $\mathcal O_E$-action and polarization,
  extends to an abelian scheme $B$ over $\mathcal O_K$, and each $\psi_\pm(b)$ is the map on complex
  points of a homomorphism over $\mathcal O_K$ from the base change $C_\pm$ of $Z_\pm$ to $B$; these
  homomorphisms have the properties of (1).
\item Let $a_1,\dots,a_g\in L$ with $a_k\in L_{e_k}(P)$ for signs $e_k$, such that $N'L$ lies in the
  $\mathcal O$-span of the $a_k$ for an integer $N'\ge1$. Let $j:K\to\mathbf C_\ell$ be an embedding with
  $\ell\nmid N'$, and $V_j=\{b\in K:|j(b)|\le1\}$. Then the images of the given bases of the Lie
  algebras of the $C_{e_k}$ under the derivatives of the $\psi_{e_k}(a_k)$ form a basis of the Lie
  algebra of $B$ over $V_j$.
\end{enumerate}
\end{lemma}
\leanref{Development/AbelianVarietyTheory/Proofs/CubicalLineBundles/ExistsCmCurveHomFamilyToFibreAtPeriodStatement.lean}{GlobalTheory.T.exists_cmCurve_homFamily_to_fibre_at_period}
\leanref{Development/AbelianVarietyTheory/Proofs/CubicalLineBundles/ExistsUniqueCmHermitianLatticeFormStatement.lean}{GlobalTheory.T.exists_unique_cmHermitianLatticeForm}
\leanref{Development/AutomorphicSpectrum/Proofs/GeometryProofs/FamilyPurePeriodsSpanStatement.lean}{GlobalTheory.T.family_pure_periods_span}
\leanref{Development/AutomorphicSpectrum/Proofs/GeometryProofs/PeriodsSpanPureDecompositionStatement.lean}{GlobalTheory.T.periods_span_pure_decomposition}
\leanref{Development/AbelianVarietyTheory/Proofs/CubicalLineBundles/CmFamilyDescendsToNumberFieldStatement.lean}{GlobalTheory.T.cmFamily_descends_to_numberField}
\leanref{Development/AbelianVarietyTheory/Proofs/CubicalLineBundles/PrincipalTripleExtendsOfCurveHomomorphismsStatement.lean}{GlobalTheory.T.principal_triple_extends_of_curve_homomorphisms}
\leanref{Development/AbelianVarietyTheory/Proofs/CubicalLineBundles/CmFamilyExtendsOverDedekindDomainStatement.lean}{GlobalTheory.T.cmFamily_extends_over_dedekindDomain}
\leanref{Development/AutomorphicSpectrum/Proofs/GeometryProofs/CmPeriodAlgebraicStatement.lean}{GlobalTheory.T.exists_cmPack_field_ge_of_isPeriod}
\leanref{Development/AutomorphicSpectrum/Proofs/GeometryProofs/CmPackProductFrameStatement.lean}{GlobalTheory.T.cm_pack_product_frame}
\begin{proof}
(1) For $b\in L_+(P)$ only the coordinates $k$ with $\sigma_k(\sqrt{-d})=\theta$ of $P(b)$ can be
nonzero, so $P(w_0b)=\frac{1+\theta}2P(b)$, and $\frac{1+\theta}2$ is $w$ or $\bar w=1-w$; similarly
$P(w_0b)=\frac{1-\theta}2P(b)$ for $b\in L_-(P)$. Hence $(m+nw)P(b)\in P(L)$ for $m,n\in\mathbf Z$, so
$t\mapsto tP(b)$ carries $\mathbf Z+\mathbf Zw$ into $\ker\exp_P$, and $\psi_\pm(b)$ is a well-defined
holomorphic homomorphism. Additivity and compatibility with $\mathcal O_E$ come from those of $P$ and
$\exp_P$. On $Z_\pm$, $w_0$ acts on the tangent space by $\frac{1\pm\theta}2$, which is how it acts on
$P(b)$, so $\psi_\pm(b)$ commutes with $\mathcal O$. The Riemann form of $A_{y_P}$ is
$\langle\ ,\ \rangle$ on $P(L)$, so its pull-back along $(\psi_\pm(b),\psi_\pm(b'))$ to the lattice of
$Z_\pm$, identified with $\mathcal O$ through $c_\pm$, is $(x,x')\mapsto\langle xb,x'b'\rangle$. A
computation in the basis $1,w_0$ of $\mathcal O$ identifies it with the Riemann form of the
principal polarization of $Z_\pm$ composed with $h_\pm(b,b')$; the sign $\varepsilon_\pm$ is the one
for which that polarization is positive.

(2) Since $NL\subset L_+(P)+L_-(P)$ and $P(L)$ spans $\mathbf C^g$ over $\mathbf R$, there are
$a_1,\dots,a_g\in L_+(P)\cup L_-(P)$ with $P(a_1),\dots,P(a_g)$ a basis of $\mathbf C^g$; then the sum
of the $\psi(a_k)$ is a surjection with finite kernel from a product of the complex curves of $Z_\pm$
onto $A_{y_P}$. So $A_{y_P}$ is isogenous to a product of elliptic curves with complex multiplication
defined over number fields; it is then defined over a number field together with its
$\mathcal O_E$-action, polarization and level structure, and $y_P$ comes from a point $y$ over a number
field \cite{ShT61}. Homomorphisms between abelian varieties over number fields that are defined over
$\mathbf C$ are defined over a finite extension, and by additivity and $\mathcal O_E$-linearity finitely
many $\psi_\pm(b)$ determine all of them; let $K$ contain $k_0$, $F_c$, $M_\pm$, the field of $y$ and
fields of definition of these. Over $K$, $A_y$ is isogenous to a product of the curves $C_\pm$, which
have good reduction at every prime, so $A_y$ has good reduction at every prime
\cite[Corollary~2 of Theorem~1]{ST68} and extends to an abelian scheme $B$ over $\mathcal O_K$
\cite[\S1]{ST68}. Its $\mathcal O_E$-action, its polarization and the homomorphisms $\psi_\pm(b)$ extend
over $\mathcal O_K$ \cite[Chapter~I, Proposition~2.7]{FC90}, and the identities of (1) hold over
$\mathcal O_K$ because they hold on the generic fibre.

(3) Let $\Theta=\sum_k\psi_{e_k}(a_k):\prod_kC_{e_k}\to B$. Since $P(N'L)$ spans $\mathbf C^g$ over
$\mathbf R$ and lies in the span of the $P(a_k)$, these $g$ vectors are a basis of $\mathbf C^g$, and
$\Theta$ is surjective. If $\sum_k\psi_{e_k}(a_k)(c_{e_k}t_k)=0$ on complex points, then
$\sum_kt_kP(a_k)=P(y')$ for some $y'\in L$; writing $N'y'=\sum_ks_ka_k$ with $s_k\in\mathcal O$ and
comparing coefficients, $N't_k$ is the image of $s_k$ in $\mathbf Z+\mathbf Zw$, so $N'$ kills the kernel
of $\Theta$. Hence multiplication by $N'$ factors as $\chi\circ\Theta$ for a homomorphism
$\chi:B\to\prod_kC_{e_k}$, which extends over $\mathcal O_K$ \cite[Chapter~I, Proposition~2.7]{FC90}, and
then $\Theta\circ\chi=N'$ as well. Over $V_j$ the integer $N'$ is a unit, so the derivative of $\Theta$
is an isomorphism of Lie algebras over $V_j$, and the Lie algebra of $\prod_kC_{e_k}$ over $V_j$ has the
given bases of the factors as a basis.
\end{proof}

\begin{lemma}\label{cv-cm-point-trace}
Let $E$ be a CM field and $d\equiv3\pmod4$ a prime with $\sqrt{-d}\in E$, let
$F=\mathbf Q(\sqrt{-d})\subset E$ and $\mathcal O=\mathbf Z[w_0]$ with $w_0=(1+\sqrt{-d})/2$, and let
$(b_k)$ be a $\mathbf Z$-basis of $\mathcal O_E$ with trace-dual basis $(b_k^\vee)$, so that
$\operatorname{Tr}_{E/\mathbf Q}(b_kb_{k'}^\vee)$ is $1$ for $k=k'$ and $0$ otherwise.
\begin{enumerate}
\item For embeddings $\sigma,\sigma'$ of $E$ into a field, $\sum_k\sigma(b_k^\vee)\sigma'(b_k)$ is $1$
  if $\sigma=\sigma'$ and $0$ otherwise. Hence, for a period $P$ of a lattice $L$ as in
  Definition~\ref{cv-cm-point-periods} and $x\in L$, $\sum_k\sigma(b_k^\vee)P(b_kx)=P(x)_\sigma$; and if
  $\mathcal O_E$ acts on a vector space $U$ over a field $k$ through $\rho$, and $U$ is the sum of the
  eigenspaces $U_{\sigma'}$ of the embeddings $\sigma':E\to k$, then $\sum_k\sigma(b_k^\vee)\rho(b_k)$ is
  the projection $e_\sigma$ onto $U_\sigma$.
\item Let $L$ be as in Definition~\ref{cv-cm-point-periods}, let $\varepsilon\in\{\pm1\}$ and
  $h(x,y)=\varepsilon\big(\langle w_0x,y\rangle-\langle x,y\rangle+\langle x,y\rangle w_0\big)$. Then
  $h(x,y)=\varepsilon\sqrt{-d}\operatorname{Tr}_{E/F}\langle x,y\rangle_E$, and for every embedding
  $\sigma:E\to\mathbf C$,
  \[
   \sum_{k,k'}\sigma(b_k^\vee)\,\sigma\big(h(b_kx,b_{k'}y)\big)\,\overline{\sigma(b_{k'}^\vee)}
   =\varepsilon\,\sigma\big(\sqrt{-d}\langle x,y\rangle_E\big).
  \]
\item Let $C$ be a principally polarized complex elliptic curve with an action of $\mathcal O$, and $v$
  a tangent vector of $C$ of squared Hodge norm $r$ on which $w_0$ acts by $\kappa(w_0)$, for a ring map
  $\kappa:\mathcal O\to\mathbf C$. Let $A$ be a polarized complex abelian variety, $S$ a set, $h$ a map
  $S\times S\to\mathcal O$, and $\psi(a):C\to A$ homomorphisms for $a\in S$ such that the polarization
  of $A$ pulls back along $(\psi(a),\psi(a'))$ to that of $C$ composed with the action of $h(a,a')$.
  Then for $a_1,\dots,a_n\in S$ the Hodge Gram matrix of the vectors $d\psi(a_m)(v)$ is
  $\big(r\,\kappa(h(a_m,a_{m'}))\big)_{m,m'}$. If $V'_m=\sum_kT_{mk}V_k$ for tangent vectors $V_k$ of
  $A$ with Hodge Gram matrix $\Gamma$, the Hodge Gram matrix of the $V'_m$ is
  $T\Gamma\bar T^{\mathsf t}$.
\end{enumerate}
\end{lemma}
\leanref{Development/AutomorphicSpectrum/Proofs/GeometryProofs/CmPackEigenPeriod.lean}{GlobalTheory.CompanionTheory.Ar.Fl.CoordPart.sum_eq_tyPart}
\leanref{Development/AutomorphicSpectrum/Proofs/GeometryProofs/DualBasisHermitianSumStatement.lean}{GlobalTheory.T.dual_basis_hermitian_sum}
\leanref{Development/AutomorphicSpectrum/Proofs/GeometryProofs/LieNormSquareOfFamilyStatement.lean}{GlobalTheory.T.lie_norm_square_of_family}
\leanref{Development/AutomorphicSpectrum/Proofs/GeometryProofs/HodgeGramLinearCombinationStatement.lean}{GlobalTheory.T.hodge_gram_linear_combination}
\begin{proof}
(1) The map $x\mapsto\sum_k\operatorname{Tr}_{E/\mathbf Q}(xb_k^\vee)b_k$ is the identity of $E$, so
$\sum_{\rho}\rho(x)\sum_k\rho(b_k^\vee)\sigma'(b_k)=\sigma'(x)$ for all $x$, the outer sum over the
embeddings $\rho$; the characters $\rho$ are linearly independent, which gives the first claim. The
other two follow, since $b$ acts on the coordinate $P(\cdot)_j$ by $\sigma_j(b)$ and on $U_{\sigma'}$ by
$\sigma'(b)$.

(2) Let $f=\operatorname{Tr}_{E/F}\langle x,y\rangle_E$. For $s\in\mathcal O$,
$\langle sx,y\rangle=\operatorname{Tr}_{F/\mathbf Q}(sf)$, so $\langle x,y\rangle=f+\bar f$ and
$\langle w_0x,y\rangle=w_0f+\bar w_0\bar f$. As $w_0-1=-\bar w_0$, the bracket in $h$ is
$(w_0-\bar w_0)f=\sqrt{-d}f$. Hence
$\sigma(h(b_kx,b_{k'}y))=\varepsilon\sigma(\sqrt{-d})\sum_\rho\rho(b_k)\overline{\rho(b_{k'})}\rho(\langle x,y\rangle_E)$,
the sum over the embeddings $\rho$ that agree with $\sigma$ on $F$, and by (1) the double sum keeps
only $\rho=\sigma$.

(3) By the hypothesis on the polarizations, the Hermitian form $H_A$ of $A$ satisfies
$H_A(d\psi(a)(v),d\psi(a')(v))=\kappa(h(a,a'))H_C(v,v)$, and $\pi H_C(v,v)=r$. The last claim is the
sesquilinearity of $H_A$.
\end{proof}

\begin{lemma}\label{cv-cm-point-eigen}
Let $E$ be a CM field and $d\equiv3\pmod4$ a prime with $\sqrt{-d}\in E$, let $\mathcal O=\mathbf Z[w_0]$
with $w_0=(1+\sqrt{-d})/2$, and let $\ell$ be a prime. Use the notation $|\mathfrak d_E|_\iota$ of
Lemma~\ref{cv-cm-point-gram}.
\begin{enumerate}
\item There are $\gamma_1,\dots,\gamma_f\in\mathcal O_E$ with $2f=[E:\mathbf Q]$ that form a local basis of
  $\mathcal O_E$ over $\mathcal O$ at $\ell$, and $\gamma'_1,\dots,\gamma'_f\in E$ such that, for
  embeddings $\sigma,\sigma'$ of $E$ into a field with $\sigma(\sqrt{-d})=\sigma'(\sqrt{-d})$,
  $\sum_k\sigma(\gamma'_k)\sigma'(\gamma_k)$ is $1$ if $\sigma=\sigma'$ and $0$ otherwise, and
  $\log\max_k|\iota(\gamma'_k)|=\log|\iota(\sqrt{-d})|-\log|\mathfrak d_E|_\iota$ for every embedding
  $\iota:E\to\mathbf C_\ell$.
\item There is $\delta_0\in E^\times$ with
  $\log|\iota(\delta_0)|=\log|\iota(\bar\delta_0)|=\log|\mathfrak d_E|_\iota-\log|\iota(\sqrt{-d})|$ for
  every embedding $\iota:E\to\mathbf C_\ell$.
\item Let $K$ be a field of characteristic zero with $[E:\mathbf Q]$ embeddings $E\to K$, let
  $j:K\to\mathbf C_\ell$ be an embedding and $V=\{b\in K:|j(b)|\le1\}$. Let $U$ be a vector space over
  $K$ with an action $\rho$ of $\mathcal O_E$; it is the sum of the eigenspaces $U_{\sigma'}$ of the
  embeddings $\sigma':E\to K$. Let $\sigma:E\to K$ be an embedding and $e_\sigma$ the projection onto
  $U_\sigma$. Let $u_1,\dots,u_r\in U$ with $\rho(\sqrt{-d})u_a=\sigma(\sqrt{-d})u_a$ and
  $u'_1,\dots,u'_{r'}\in U$ with $\rho(\sqrt{-d})u'_i=-\sigma(\sqrt{-d})u'_i$, such that the vectors
  $\rho(\gamma_k)u_a$ and $\rho(\gamma_k)u'_i$ are linearly independent over $K$, and let
  $\delta\in K^\times$ with $\log|j(\delta)|=\log|\mathfrak d_E|_{j\sigma}-\log|j\sigma(\sqrt{-d})|$.
  Then the vectors $\delta e_\sigma(u_a)$ are linearly independent, and their span over $V$ is the
  intersection of $U_\sigma$ with the span over $V$ of the $\rho(\gamma_k)u_a$ and $\rho(\gamma_k)u'_i$.
\end{enumerate}
\end{lemma}
\leanref{Development/AutomorphicSpectrum/Proofs/AlgebraProofs/ExistsLocalBasisCmOrderStatement.lean}{GlobalTheory.T.exists_local_basis_cm_order}
\leanref{Development/AutomorphicSpectrum/Proofs/AlgebraProofs/DualBasisTraceLogNormStatement.lean}{GlobalTheory.T.dual_basis_trace_log_norm}
\leanref{Development/AutomorphicSpectrum/Proofs/AlgebraProofs/DifferentTwistedGeneratorExistsStatement.lean}{GlobalTheory.T.different_twisted_generator_exists}
\leanref{Development/AutomorphicSpectrum/Proofs/AlgebraProofs/EigenLocalBasisLogNormStatement.lean}{GlobalTheory.T.eigen_local_basis_log_norm}
\begin{proof}
Let $F=\mathbf Q(\sqrt{-d})$; as $d\equiv3\pmod4$, $\mathcal O=\mathcal O_F$ and its different is
$\sqrt{-d}\,\mathcal O$, so the different of $E$ is $\mathfrak d_{E/F}\sqrt{-d}$, with $\mathfrak d_{E/F}$
the relative different.

(1) The module $\mathcal O_E\otimes\mathbf Z_{(\ell)}$ is finitely generated and torsion free over the
semilocal Dedekind ring $\mathcal O\otimes\mathbf Z_{(\ell)}$, which is a principal ideal domain, so it is
free of rank $[E:F]=f$; a basis, multiplied by an integer prime to $\ell$, gives the $\gamma_k$. Let
$(\gamma'_k)$ be the dual basis for $\operatorname{Tr}_{E/F}$. The identity for $\sigma$ and $\sigma'$
follows as in Lemma~\ref{cv-cm-point-trace}(1), with $F$ in place of $\mathbf Q$. The span of the
$\gamma'_k$ over $\mathcal O\otimes\mathbf Z_{(\ell)}$ is the localization of $\mathfrak d_{E/F}^{-1}$,
so by the ultrametric inequality $\max_k|\iota(\gamma'_k)|$ is the largest absolute value on
$\iota(\mathfrak d_{E/F}^{-1})$, which is $|\iota(\sqrt{-d})|/|\mathfrak d_E|_\iota$.

(2) Complex conjugation preserves $F$, hence $\mathfrak d_{E/F}$, and permutes the primes above $\ell$.
Choose $\delta_0\in\mathfrak d_{E/F}$ generating it at every prime above $\ell$; then so does
$\bar\delta_0$, and both have absolute value $|\mathfrak d_E|_\iota/|\iota(\sqrt{-d})|$ at every $\iota$.

(3) Since $\sqrt{-d}$ acts on $u_a$ by $\sigma(\sqrt{-d})$, only eigenspaces of embeddings $\sigma'$ with
$\sigma'(\sqrt{-d})=\sigma(\sqrt{-d})$ meet the components of $u_a$, and (1) gives
$e_\sigma(u_a)=\sum_k\sigma(\gamma'_k)\rho(\gamma_k)u_a$; likewise $e_\sigma$ kills every
$\rho(\gamma_k)u'_i$. Let $k_0$ be an index with $|j\sigma(\gamma'_{k_0})|$ largest; by (1) and the
choice of $\delta$, $|j(\delta\sigma(\gamma'_k))|\le|j(\delta\sigma(\gamma'_{k_0}))|=1$ for all $k$. So
$\delta e_\sigma(u_a)$ lies in the span over $V$, and these vectors are linearly independent because
their coefficients at the independent vectors $\rho(\gamma_{k_0})u_a$ are nonzero. Conversely, let
$x=\sum_{k,a}c_{ka}\rho(\gamma_k)u_a+\sum_{k,i}c'_{ki}\rho(\gamma_k)u'_i$ with coefficients in $V$ lie in
$U_\sigma$. Since $e_\sigma\rho(\gamma)=\sigma(\gamma)e_\sigma$, $x=e_\sigma(x)=\sum_ax_ae_\sigma(u_a)$
with $x_a=\sum_kc_{ka}\sigma(\gamma_k)$. Expanding $e_\sigma(u_a)$ and comparing the coefficients of
$\rho(\gamma_{k_0})u_a$ gives $c_{k_0a}=x_a\sigma(\gamma'_{k_0})$, so
$|j(x_a\delta^{-1})|=|j(c_{k_0a})|\le1$, and $x=\sum_a(x_a\delta^{-1})\,\delta e_\sigma(u_a)$.
\end{proof}

\begin{lemma}\label{cv-cm-point}
Let $E$ be a CM field and let $d\equiv7\pmod{16}$ be a prime with $\sqrt{-d}\in E$, such that $E$ is
unramified over its totally real subfield. Let $\chi$ be the Kronecker character of
$\mathbf Q(\sqrt{-d})$, and put $x=\gamma+\log(2\pi/\sqrt d)-\frac{L'}{L}(1,\chi)$.
\begin{enumerate}
\item For every $\theta\in\mathbf C$ with $\theta^2=-d$ there is a CM curve with root $\theta$ over a
  number field $M$ (Definition~\ref{cv-cm-point-curve}) with
  $-\frac1{[M:\mathbf Q]}\sum_s\log n(s)=x$, the sum over the embeddings $s:M\to\mathbf C$.
\item Let a PEL datum of signature $(p,q)$ over $E$, a level, its moduli scheme $X$ and its
  coefficient field $F_c\subset\mathbf C$ be given; $F_c$ contains $\sigma(E)$ for every embedding
  $\sigma$. Let $\tau:E\to\mathbf C$ be an embedding, $\bar\tau$ its complex conjugate and
  $\theta=\tau(\sqrt{-d})$. Let $L$ be as in Definition~\ref{cv-cm-point-periods}, with a complex
  structure of signature $(p,q)$ for $\theta$, fix a type of
  signature $(p,q)$ for $\theta$, and suppose that to each period $P$ of $L$ of this type are
  attached a complex point $y_P$ of $X$ over $F_c\subset\mathbf C$ and a surjective holomorphic
  homomorphism $\exp_P:\mathbf C^g\to A_{y_P}(\mathbf C)$ onto the abelian variety at $y_P$, with
  kernel $P(L)$, compatible with the actions of $\mathcal O_E$ (on $\mathbf C^g$ through the type)
  and with the polarizations. Let $Z_+$ and $Z_-$ be CM curves with roots $\theta$ and $-\theta$,
  over number fields $M_\pm$, with numbers $z_\pm$ and squared Hodge norms $n_\pm$. Let $\ell$ be a
  prime and $k_0\subset\mathbf C$ a number field. Then there are a period $P$, a local basis
  $u=(u_1,\dots,u_{p+q})$ of $L$ at $\ell$ with $u_1,\dots,u_p\in L_+(P)$ and
  $u_{p+1},\dots,u_{p+q}\in L_-(P)$, a number field $K\subset\mathbf C$ containing $k_0$, $F_c$,
  $M_+$ and $M_-$, a point $y\in X(K)$ over $F_c\subset K$ whose image under $K\subset\mathbf C$ is
  $y_P$, and bases $v=(v_1,\dots,v_p)$ of $\operatorname{Lie}_\tau(y)$ and $w=(w_1,\dots,w_q)$ of
  $\operatorname{Lie}_{\bar\tau}(y)$, such that:
  \begin{enumerate}
  \item under $K\subset\mathbf C$, $v_a=d\exp_P\big(z_+P(u_a)_\tau\big)$ for $a\le p$ and
    $w_i=d\exp_P\big(z_-P(u_{p+i})_{\bar\tau}\big)$ for $i\le q$;
  \item for every embedding $s:K\to\mathbf C$, the Hodge Gram matrices of the images of $v$ and $w$
    in the abelian variety at $s(y)$ are
    \[
     \big(n_+(s)\,\varepsilon_+\,s(\tau(\sqrt{-d}\langle u_a,u_b\rangle_E))\big)_{a,b\le p},\qquad
     \big(n_-(s)\,\varepsilon_-\,s(\bar\tau(\sqrt{-d}\langle u_{p+i},u_{p+i'}\rangle_E))\big)_{i,i'\le q},
    \]
    where $\varepsilon_\pm\in\{\pm1\}$ is the sign of $\pm\operatorname{Im}\theta$ and $n_\pm(s)$ is
    $n_\pm$ at the restriction of $s$ to $M_\pm$;
  \item for every embedding $j:K\to\mathbf C_\ell$, the abelian variety at $y$, with its
    $\mathcal O_E$-action and polarization, extends to a semi-abelian scheme over the valuation ring
    $V_j=\{b\in K:|j(b)|\le1\}$, and there are $\delta,\delta'\in K^\times$ with
    $\log|j(\delta)|=\log|j(\delta')|=\log|\mathfrak d_E|_{j\tau}-\log|j\tau(\sqrt{-d})|$
    (Lemma~\ref{cv-cm-point-gram}) such that $\delta v$ and $\delta'w$ are bases of the $\tau$- and
    $\bar\tau$-eigenspaces of the Lie algebra of this scheme.
  \end{enumerate}
\end{enumerate}
\end{lemma}
\leanref{Development/AutomorphicSpectrum/Proofs/GeometryProofs/ChowlaSelbergPeriodStatement.lean}{GlobalTheory.T.chowla_selberg_period}
\leanref{Development/AutomorphicSpectrum/Proofs/GeometryProofs/ExistsCmPeriodEmbeddingStatement.lean}{GlobalTheory.T.exists_cm_period_embedding}
\leanref{Development/AutomorphicSpectrum/Definitions/GlobalGeometry/IsCMFrame.lean}{GlobalTheory.CompanionTheory.Ar.Fl.CMPack}
\leanref{Development/AutomorphicSpectrum/Proofs/GeometryProofs/CmPackPureFrameStatement.lean}{GlobalTheory.T.cm_pack_pure_frame}
\leanref{Development/AutomorphicSpectrum/Proofs/GeometryProofs/CmPackEigenPeriodStatement.lean}{GlobalTheory.T.cm_pack_eigen_period}
\leanref{Development/AutomorphicSpectrum/Proofs/GeometryProofs/CmPackPeriodProductStatement.lean}{GlobalTheory.T.cm_pack_period_product}
\leanref{Development/AutomorphicSpectrum/Proofs/GeometryProofs/CmPackPureNormStatement.lean}{GlobalTheory.T.cm_pack_pure_norm}
\begin{proof}
(1) Let $F=\mathbf Q(\sqrt{-d})$, let $\sigma_0:F\to\mathbf C$ be the embedding with
$\sigma_0(\sqrt{-d})=i\sqrt d$, and choose an ideal $\mathfrak b_C$ in each ideal class $C$ of $F$.
Lemma~\ref{cv-cm-point-models} gives a number field $M\supset\sigma_0(F)$, an elliptic curve $P$ over
$\mathcal O_M$ with a basis $v$ of its Lie algebra, and the other data there. As $d\equiv3\pmod4$,
$\mathcal O_F=\mathcal O=\mathbf Z[w_0]$ and $\sigma_0(\mathcal O_F)=\mathbf Z+\mathbf Zw$ with
$w=(1+i\sqrt d)/2\in M$. By Lemma~\ref{cv-cm-point-models}(4) the complex curve of $P$ at the inclusion
$M\subset\mathbf C$ is $\mathbf C/c(\mathbf Z+\mathbf Zw)$, in which $v$ is the tangent vector
$1=c\cdot c^{-1}$; so $z=c^{-1}$. By Lemma~\ref{cv-cm-point-models}(5) the generic fibre of $P$ has an
endomorphism that multiplies $v$ by $w$, and Lemma~\ref{cv-cm-point-action}(1) and (2) give an action
of $\mathcal O$ on $P$, with Rosati involution complex conjugation, under which $\sqrt{-d}$ acts on $v$
by $i\sqrt d$. This is a CM curve with root $i\sqrt d$. By Lemma~\ref{cv-cm-point-area}(5) and
Theorem~\ref{cv-chowla-selberg}(1), applied with $\varphi=\sigma_0$ and $\mathfrak a_C=\mathfrak b_C$,
\[
 -\frac1{[M:\mathbf Q]}\sum_s\log n(s)=\gamma+\log\frac{2\pi}{\sqrt d}-\frac{L'}{L}(1,\chi)=x.
\]
For $\theta=-i\sqrt d$, compose the action of $\mathcal O$ with complex conjugation:
Lemma~\ref{cv-cm-point-action}(2) and (3) show that this is a CM curve with root $-i\sqrt d$, with the
same $c$, $z$ and $n(s)$.

(2) \emph{Step 1: the period and the local basis.} Lemma~\ref{cv-cm-point-split}(7) gives a period $P$
and a local basis $u$ at $\ell$ with $u_1,\dots,u_p\in L_+(P)$ and $u_{p+1},\dots,u_{p+q}\in L_-(P)$,
and an integer $N$ prime to $\ell$ with $NL\subset\sum_m\mathcal O_Eu_m\subset L_+(P)+L_-(P)$.

\emph{Step 2: the CM point.} Lemma~\ref{cv-cm-point-descent} with this $P$ and with $k_0$ gives the
number field $K$, the point $y$, the abelian scheme $B$ over $\mathcal O_K$ extending $A_y$, the curves
$C_\pm$ over $\mathcal O_K$ with the bases $v_\pm$ of their Lie algebras given by $Z_\pm$, and the
homomorphisms $\psi_\pm(b)$ for $b\in L_\pm(P)$, with the forms $h_\pm$ and the signs
$\varepsilon_\pm$.

\emph{Step 3: the bases $v$ and $w$.} Let $(b_k)$ be a $\mathbf Z$-basis of $\mathcal O_E$ with
trace-dual basis $(b_k^\vee)$. Since $F_c\subset K$ contains the values of $\tau$ and $\bar\tau$, the
vectors
\[
 v_a=\sum_k\tau(b_k^\vee)\,d\psi_+(b_ku_a)(v_+),\qquad
 w_i=\sum_k\bar\tau(b_k^\vee)\,d\psi_-(b_ku_{p+i})(v_-)
\]
lie in the Lie algebra of $A_y$ over $K$.

(a) In the complex curve of $Z_\pm$, $v_\pm$ is the tangent vector $c_\pm z_\pm$, and
$\psi_\pm(b)(c_\pm t)=\exp_P(tP(b))$; so $d\psi_\pm(b)(v_\pm)=d\exp_P(z_\pm P(b))$. By
Lemma~\ref{cv-cm-point-trace}(1), $\sum_k\tau(b_k^\vee)P(b_ku_a)=P(u_a)_\tau$, which gives
$v_a=d\exp_P(z_+P(u_a)_\tau)$, and in the same way $w_i=d\exp_P(z_-P(u_{p+i})_{\bar\tau})$. These are
bases of $\operatorname{Lie}_\tau(y)$ and $\operatorname{Lie}_{\bar\tau}(y)$. Indeed, for
$x=\sum_mc_mu_m$ with $c_m\in E\otimes\mathbf R$ one has $P(x)_\tau=\sum_{a\le p}\tau(c_a)P(u_a)_\tau$,
because $P(u_{p+i})_\tau=0$; as $P(L\otimes\mathbf R)=\mathbf C^g$, the $p$ vectors $P(u_a)_\tau$ span the
$p$-dimensional $\tau$-part of $\mathbf C^g$. Moreover $z_+\ne0$, and $d\exp_P$ is an isomorphism
compatible with the actions of $\mathcal O_E$, which carries the $\tau$-part to
$\operatorname{Lie}_\tau(y)$. The same argument applies to $w$.

(b) Let $s:K\to\mathbf C$ be an embedding. The complex curve of $C_+$ at $s$ is that of $Z_+$ at the
restriction of $s$ to $M_+$, so the squared Hodge norm of $v_+$ there is $n_+(s)$; and $w_0$ acts on
$v_+$ by $s(\frac{1+\theta}2)=s\tau(w_0)$. By Lemma~\ref{cv-cm-point-descent}(1) and (2) and
Lemma~\ref{cv-cm-point-trace}(3), the Hodge Gram matrix at $s(y)$ of the vectors $d\psi_+(b_ku_a)(v_+)$
is $\big(n_+(s)\,s\tau(h_+(b_ku_a,b_{k'}u_{a'}))\big)$. The vectors $v_a$ are their combinations with
coefficients $s\tau(b_k^\vee)$, so by Lemma~\ref{cv-cm-point-trace}(3) and (2), with $\sigma=s\tau$,
the Hodge Gram matrix of the $v_a$ is
$\big(n_+(s)\,\varepsilon_+\,s\tau(\sqrt{-d}\langle u_a,u_{a'}\rangle_E)\big)$. The same computation with
$\bar\tau$, $Z_-$ and $\varepsilon_-$ gives the matrix for $w$.

(c) Let $j:K\to\mathbf C_\ell$ be an embedding. The ring $V_j$ contains $\mathcal O_K$, so $B$ gives an
abelian, in particular semi-abelian, scheme over $V_j$ extending $A_y$ with its $\mathcal O_E$-action and
polarization. Let $\gamma_1,\dots,\gamma_f$ and $\gamma'_k$ be as in Lemma~\ref{cv-cm-point-eigen}(1).
The $g=f(p+q)$ elements $\gamma_ku_m$ lie in $L_+(P)$ for $m\le p$ and in $L_-(P)$ for $m>p$, and their
$\mathcal O$-span contains $N'L$ for an integer $N'$ prime to $\ell$, because $u$ is a local basis of
$L$ over $\mathcal O_E$ at $\ell$ and the $\gamma_k$ are a local basis of $\mathcal O_E$ over $\mathcal O$
at $\ell$. Write $\lambda(\gamma_ku_m)$ for the image of $v_\pm$ under $d\psi_\pm(\gamma_ku_m)$. By
Lemma~\ref{cv-cm-point-descent}(3), these vectors form a basis over $V_j$ of the Lie algebra of $B$ over
$V_j$. Let $U$ be the Lie algebra of $A_y$ over $K$ with the action $\rho$ of $\mathcal O_E$, and
$\sigma:E\to K$ the embedding $\tau$. Then $\lambda(\gamma_ku_m)=\rho(\gamma_k)\lambda(u_m)$, and
$\sqrt{-d}$ acts on $\lambda(u_a)$ by $\sigma(\sqrt{-d})=\theta$ and on $\lambda(u_{p+i})$ by $-\theta$,
since it acts on $v_\pm$ by $\pm\theta$ and the $\psi_\pm$ commute with $\mathcal O$. By
Lemma~\ref{cv-cm-point-trace}(1), $v_a=e_\sigma(\lambda(u_a))$. Take $\delta_0$ as in
Lemma~\ref{cv-cm-point-eigen}(2) and $\delta=\sigma(\delta_0)$; then
$\log|j(\delta)|=\log|\mathfrak d_E|_{j\tau}-\log|j\tau(\sqrt{-d})|$, and
Lemma~\ref{cv-cm-point-eigen}(3) shows that $\delta v$ is a basis over $V_j$ of the intersection of
$U_\sigma$ with the Lie algebra of $B$ over $V_j$, which is its $\tau$-eigenspace. For $w$, apply the
same argument to $\sigma'=\bar\tau$, with the roles of the $u_a$ and the $u_{p+i}$ exchanged, and
$\delta'=\sigma'(\delta_0)=\sigma(\bar\delta_0)$; by Lemma~\ref{cv-cm-point-eigen}(2),
$|j(\delta')|=|j(\delta)|$.
% Lean: (1) is Ar.Fl.exists_curve0_hasPeriod from CM.EllExists.exists_ell (exists_ell_pos,
%   chowla_selberg_lie_lines, cmCurve_of_uniformisation, chowla_selberg_of_mean_log,
%   EllConj.exists_conj_action); (2) is Ar.Fl.exists_cmPoint_locBasis_frames: Step 1 is
%   DSplit.exists_split_period with pure_family_decomposition, Step 2 is CMGiven.exists_pack_given,
%   and v, w come from cm_pack_pure_frame; (a) CoordPart.isDerivAt0_of_isCMVec; (b)
%   GInf.hodgeGram_cmVec (gram_family, dual_basis_hermitian_sum); (c) GFin.intFrames_fin
%   (cm_pack_product_frame, eigen_local_basis_log_norm, different_twisted_generator_exists).
\end{proof}

\subsection{Models and theta-level objects}\label{subsection-cv-models-theta-level}

\begin{situation}\label{cv-function-situation}
Let a PEL datum of signature $(p,q)$ over a CM field $E$, a level, its moduli scheme $X$ over its
coefficient field $F_c\subset\mathbf C$ and the universal abelian scheme $A$ over $X$ be given. A
prime $d\equiv7\pmod{16}$, an element $\sqrt{-d}\in E$ with square $-d$ and a square root $\theta$ of
$-d$ in $\mathbf C$ are part of the datum, and no finite prime of the totally real subfield of $E$
ramifies in $E$; its CM type $\Phi$ is the set of embeddings $\tau:E\to\mathbf C$
with $\tau(\sqrt{-d})=\theta$, and $F_c$ contains $\sigma(E)$ for every embedding $\sigma$. Let $Z$
be a connected component of $X$, with function field $F_c(Z)$, and let $\tau\in\Phi$. For a point $y$
of $X$ with values in a field $k\supset F_c$, let $\operatorname{Lie}_\tau(y)$ and
$\operatorname{Lie}_{\bar\tau}(y)$ be the subspaces of the Lie algebra of $A_y$ on which each
$b\in\mathcal O_E$ acts by $\tau(b)$ and by $\bar\tau(b)$. For a basis $v=(v_1,\dots,v_n)$ of a
vector space write $\wedge v=v_1\wedge\dots\wedge v_n$.
\end{situation}

\begin{definition}\label{cv-function-framed}
In Situation~\ref{cv-function-situation}:
\begin{enumerate}
\item A \emph{framed point} of $X$ is a triple $(y,v,w)$ of a point $y$ of $X$ with values in a field
  $k\supset F_c$, a basis $v=(v_1,\dots,v_p)$ of $\operatorname{Lie}_\tau(y)$ and a basis
  $w=(w_1,\dots,w_q)$ of $\operatorname{Lie}_{\bar\tau}(y)$. It is a framed point of $Z$ if $y$ is a
  point of $Z$, and it is complex over an embedding $\sigma:F_c\to\mathbf C$ if $k=\mathbf C$ and
  $F_c\to k$ is $\sigma$.
\item A \emph{frame function} on $Z$ assigns to each framed point $(y,v,w)$ of $Z$ with values in $k$
  an element $\varphi(y,v,w)\in k^\times$, such that
  \[
   \varphi(y,v',w')\det g=\varphi(y,v,w)\det h
  \]
  whenever $v'_a=\sum_bg_{ab}v_b$ and $w'_i=\sum_jh_{ij}w_j$ for invertible matrices $g$ and $h$ over
  $k$, and such that for an extension of fields $k\to k'$ the value of $\varphi$ at the image of
  $(y,v,w)$ is the image of $\varphi(y,v,w)$.
\end{enumerate}
\end{definition}
\leanref{Development/AutomorphicSpectrum/Definitions/GlobalGeometry/IsConstMul.lean}{GlobalTheory.CompanionTheory.Ar.Fl.IsFramed}
\leanref{Development/AutomorphicSpectrum/Definitions/GlobalGeometry/IsConstMul.lean}{GlobalTheory.CompanionTheory.Ar.Fl.IsLineFn}

\begin{definition}\label{cv-model}
In Situation~\ref{cv-function-situation}, let $O\subset F_c$ be a discrete valuation ring with
fraction field $F_c$. A \emph{model} of $X$ over $O$ is a triple $(\mathcal X,\mathcal A,j)$ of a
normal scheme $\mathcal X$, projective and flat over $O$; a semi-abelian scheme $\mathcal A$ over
$\mathcal X$ with an action of $\mathcal O_E$ and a symmetric cubical line bundle for which complex
conjugation on $\mathcal O_E$ is the Rosati involution, normalized as in
Definition~\ref{cv-model-toroidal-semiabelian}(2); and an open immersion
$j:X\to\mathcal X_{F_c}=\mathcal X\times_OF_c$ with dense image, such that $j^*\mathcal A$ is $A$
with its action and polarization, and $j(X)$ is the set of points of $\mathcal X_{F_c}$ over which
$\mathcal A$ is an abelian scheme.
\end{definition}
\leanref{ReciprocityTheory/ReductiveGroups/ArithmeticGeometry/BirationalGeometry/ToroidalCompactification/PushInvertible.lean}{ReciprocityTheory.ChartRestriction.G.Geo.Tor.Model}

\begin{theorem}[Grothendieck \cite{Gro72}]\label{cv-semistable-reduction}
Let $p$ be a prime, let $K$ be a finite extension of $\mathbf Q_p$ in an algebraic closure
$\bar{\mathbf Q}_p$, and let $A$ be a polarized abelian variety over $K$. Then there are a finite
extension $K'$ of $K$ in $\bar{\mathbf Q}_p$ and a semi-abelian scheme $G$ over the ring of integers
of $K'$ whose generic fibre is isomorphic, as a group scheme, to $A\times_KK'$.
\end{theorem}
\leanref{Development/AbelianVarietyTheory/PublishedTheorems/SemistableReduction₂/SemistableReductionSt.lean}{GlobalTheory.Published.SemistableReductionSt}
\leanref{Development/AbelianVarietyTheory/Proofs/DegenerationProofs/SemistableReductionPpavPadicField.lean}{GlobalTheory.SS.semistableReductionSt}

This is the semistable reduction theorem \cite[Expos\'e~IX, Th\'eor\`eme~3.6]{Gro72}: $A$ acquires
semistable reduction over a finite extension $K'$, and $G$ is the identity component of the N\'eron
model of $A\times_KK'$.

\begin{definition}\label{cv-model-toroidal-semiabelian}
In Situation~\ref{cv-function-situation}, let $\Lambda$ be the lattice of the PEL datum, a projective
$\mathcal O_E$-module of rank $p+q$ with a perfect alternating form.
\begin{enumerate}
\item A \emph{level} of the datum is a triple $l=(M,\mathfrak M,N)$ of an odd integer $M\ge3$, an ideal
  $\mathfrak M\subset\mathcal O_E$ generated by elements of $\mathbf Z[(1+\sqrt{-d})/2]$ with
  $\mathfrak M\bar{\mathfrak M}=M\mathcal O_E$ and $\mathfrak M+\bar{\mathfrak M}=\mathcal O_E$, and an
  integer $N\ge1$, the \emph{full level} of $l$. The moduli scheme $X_l$ at level $l$ is a smooth
  quasi-projective scheme over $K_l=\mathbf Q(\sqrt{-d},\zeta_N)\subset\mathbf C$, where
  $\zeta_N=e^{2\pi i/N}$. Its universal abelian scheme $\mathcal A_l$ carries an action of
  $\mathcal O_E$, a principal polarization, a level structure $\eta$ that identifies
  $\Lambda/\mathfrak M\Lambda$ with the $\mathfrak M$-torsion, and a full level structure $\eta'$ that
  identifies $\Lambda/N\Lambda$ with the $N$-torsion. For integers $S,N\ge1$ let
  $R_{S,N}=\mathbf Z[(1+\sqrt{-d})/2,1/S,\zeta_N]\subset\mathbf C$.
\item An \emph{$\mathcal O_E$-semi-abelian scheme} over a scheme $S$ is a semi-abelian scheme $G$ over
  $S$ with an action of $\mathcal O_E$ and a symmetric cubical line bundle $\mathcal L$ on $G$ for which
  complex conjugation on $\mathcal O_E$ is the Rosati involution and which is \emph{normalized}: for
  points $y_1,y_2$ of $G$ with values in a scheme $T$ over $S$ let
  $\Lambda_{\mathcal L}(y_1,y_2)=(y_1+y_2)^*\mathcal L\otimes y_1^*\mathcal L^{-1}\otimes
  y_2^*\mathcal L^{-1}\otimes0^*\mathcal L$, a line bundle on $T$; then
  $\mathcal L\cong\Lambda_{\mathcal L}(y,w_0y)$ on $G$, where $y$ is the identity of $G$ and
  $w_0=(1+\sqrt{-d})/2\in\mathcal O_E$. A \emph{morphism} $G'\to G$ of
  $\mathcal O_E$-semi-abelian schemes over a morphism $f:S'\to S$ is a homomorphism over $f$ that
  identifies $G'$ with $f^*G$, compatibly with the actions and the line bundles. The \emph{abelian
  locus} of $G$ is the set of points of $S$ over which $G$ is an abelian scheme.
\item An $\mathcal O_E$-semi-abelian scheme $G$ over $S$ is \emph{theta-split} if:
  \begin{enumerate}
  \item its abelian locus is open, and $G$ is proper over it;
  \item at every point $x$ of $S$, the fibre $G_x$ becomes, over a finite separable extension of the
    residue field of $x$, an extension of an abelian variety by a split torus;
  \item every point of $S$ has an open neighbourhood $V$ with a section of $\mathcal L$ over $G_V$ that
    vanishes identically on no fibre of $G_V\to V$.
  \end{enumerate}
\item Let $O$ be a Dedekind domain of characteristic zero with fraction field $F$, let $U$ be a scheme
  over $F$, and let $A$ be an $\mathcal O_E$-semi-abelian scheme over $U$. A \emph{compactification} of
  $(U,A)$ over $O$ is a triple $(\mathcal C,\mathcal G,j)$ of a normal scheme $\mathcal C$, projective
  and flat over $O$, an $\mathcal O_E$-semi-abelian scheme $\mathcal G$ over $\mathcal C$, and an open
  immersion $j:U\to\mathcal C_F=\mathcal C\otimes_OF$ over $F$ with dense image, together with a
  morphism $A\to\mathcal G$ of $\mathcal O_E$-semi-abelian schemes over $j$, such that $j(U)$ is the
  intersection of $\mathcal C_F$ with the abelian locus of $\mathcal G$.
\end{enumerate}
\end{definition}
\leanref{ReciprocityTheory/ReductiveGroups/ArithmeticGeometry/BirationalGeometry/ToroidalCompactification/TwoWeightEqOneAddRoot.lean}{ReciprocityTheory.ChartRestriction.G.Geo.Tor.Level}
\leanref{ReciprocityTheory/ReductiveGroups/ArithmeticGeometry/BirationalGeometry/ToroidalCompactification/TwoWeightEqOneAddRoot.lean}{ReciprocityTheory.ChartRestriction.G.Geo.Tor.SAB}
\leanref{ReciprocityTheory/ReductiveGroups/ArithmeticGeometry/BirationalGeometry/ToroidalCompactification/TwoWeightEqOneAddRoot.lean}{ReciprocityTheory.ChartRestriction.G.Geo.Tor.SABMap}
\leanref{Development/AbelianVarietyTheory/Definitions/ModuliSpaces/AuxiliaryInertOddPrime.lean}{GlobalTheory.Models.Good}
\leanref{Development/AbelianVarietyTheory/Definitions/ModuliSpaces/AuxiliaryInertOddPrime.lean}{GlobalTheory.Models.Cpt}

A model of $X$ over a discrete valuation ring $O\subset F_c$ with fraction field $F_c$
(Definition~\ref{cv-model}) is a compactification of $X\otimes F_c$, with its universal abelian
scheme, over $O$.

\begin{lemma}\label{cv-model-toroidal-auxiliary}
In Situation~\ref{cv-function-situation}, call a prime $n$ \emph{auxiliary} if it is odd, does not
divide the discriminant of $E$, and $-d$ is not a square modulo $n$. There are two distinct auxiliary
primes, and an auxiliary prime does not divide the integer $M$ of any level
(Definition~\ref{cv-model-toroidal-semiabelian}).
\end{lemma}
\leanref{Development/AbelianVarietyTheory/Definitions/ModuliSpaces/AuxiliaryInertOddPrime.lean}{GlobalTheory.Models.AuxiliaryInertOddPrime}
\leanref{Development/AbelianVarietyTheory/Proofs/ModuliProofs/CoprimeLevel.lean}{GlobalTheory.Models.exists_two_auxiliary_inert_odd_primes_coprime_level}
\begin{proof}
Existence. Let $b\ge0$ be an integer. Since $d\equiv3\pmod4$, Dirichlet's theorem on primes in
arithmetic progressions gives a prime $n>\max(b,d,|d_E|)$ with $n\equiv-1\pmod{4d}$, where $d_E\ne0$ is
the discriminant of $E$. Then $n\equiv3\pmod4$ and $n\equiv-1\pmod d$. By quadratic reciprocity for the
distinct primes $n$ and $d$, both congruent to $3$ modulo $4$, $d$ is a square modulo $n$ if and only if
$n$ is not a square modulo $d$; and $n\equiv-1$ is not a square modulo $d$, since $d\equiv3\pmod4$. So
$d$ is a nonzero square modulo $n$. As $n\equiv3\pmod4$, $-1$ is not a square modulo $n$, and hence
neither is $-d$. The prime $n$ is odd, and it does not divide $d_E$ since $n>|d_E|$. So $n$ is
auxiliary. Taking $b=0$ gives an auxiliary prime $n$, and then $b=n$ gives a second one $n'>n$.

Coprimality. Let $n$ be auxiliary and $(M,\mathfrak M,N)$ a level, and suppose that $n$ divides $M$.
First, if $\mathfrak p$ is a prime ideal of $\mathcal O_E$ with $n\in\mathfrak p$ and
$a\in\mathfrak p\cap\mathbf Z[w_0]$, where $w_0=(1+\sqrt{-d})/2$, then $\bar a\in\mathfrak p$. Indeed, an
integer $m$ in $\mathfrak p$ is divisible by $n$: otherwise $1=sn+tm$ for integers $s,t$, and
$1\in\mathfrak p$. Write $a=u+vw_0$ with $u,v\in\mathbf Z$. Since $2w_0=1+\sqrt{-d}$ and
$2\bar a=2u+v-v\sqrt{-d}$,
\[
 (2u+v)^2+dv^2=2\big(2u+v-v\sqrt{-d}\big)\,a\in\mathfrak p,
\]
so $n$ divides $(2u+v)^2+dv^2$. If $n\nmid v$, then $-d\equiv\big((2u+v)/v\big)^2\pmod n$, which is
excluded; so $n\mid v$, then $n\mid(2u)^2$, and $n\mid u$ because $n$ is odd. Hence
$\bar a=n(u'+v'\bar w_0)$ with $u=nu'$ and $v=nv'$, and $\bar a\in\mathfrak p$.

Now $n\mathcal O_E\ne\mathcal O_E$, since its norm $n^{[E:\mathbf Q]}$ is not $1$, so $n$ lies in a
maximal ideal $\mathfrak p$ of $\mathcal O_E$. It contains $\mathfrak M\bar{\mathfrak M}=M\mathcal O_E$,
so $\mathfrak M\subset\mathfrak p$ or $\bar{\mathfrak M}\subset\mathfrak p$. In the second case
$\mathfrak M\subset\bar{\mathfrak p}$, the conjugate prime, which also contains $n$. So
$\mathfrak M\subset\mathfrak q$ for a prime $\mathfrak q$ containing $n$. The ideal $\mathfrak M$ is
generated by elements of $\mathfrak M\cap\mathbf Z[w_0]$, whose conjugates lie in $\mathfrak q$ by the
first step, and these conjugates generate $\bar{\mathfrak M}$. So
$\mathcal O_E=\mathfrak M+\bar{\mathfrak M}\subset\mathfrak q$, a contradiction.
\end{proof}
% Lean: exists_large_auxiliary_prime (exists_prime_gt_neg_not_isSquare: Nat.forall_exists_prime_gt_and_eq_mod
% with q = 4d and a = -1, ZMod.exists_sq_eq_prime_iff_of_mod_four_eq_three, ZMod.exists_sq_eq_neg_one_iff),
% and AuxiliaryInertOddPrime.coprime_level (not_Mb_le with l.fromK and l.coprime, and bar_mem_of_mem for
% the first step; CoprimeLevel.lean).

\begin{theorem}[Faltings--Chai \cite{FC90}]\label{cv-model-toroidal-extension}
Let $T$ be a discrete valuation ring with fraction field $K$, and let $G$ and $G'$ be semi-abelian
schemes over $T$. Then every homomorphism $G_K\to G'_K$ extends uniquely to a homomorphism $G\to G'$.
\end{theorem}
\leanref{Development/AbelianVarietyTheory/Definitions/Degenerations/HomExtSt.lean}{GlobalTheory.LN.HomExtSt}
\leanref{Development/AbelianVarietyTheory/Proofs/ModuliProofs/HomExtensionSemiAbelianDiscreteValuationRingStatement.lean}{GlobalTheory.T.homExtension_semiAbelian_discreteValuationRing}

This is the case of a discrete valuation ring of \cite[Chapter~I, Proposition~2.7]{FC90}, which
extends homomorphisms of semi-abelian schemes from a dense open subscheme of a normal noetherian base.

\begin{definition}\label{cv-toroidal-family}
In Situation~\ref{cv-function-situation}, let $\Lambda$ be the lattice of the PEL datum with its
alternating form $\langle\,,\rangle$, let $n$ be a prime, let $N\ge1$, and let $R_{n,N}\subset\mathbf C$
be the ring of Definition~\ref{cv-model-toroidal-semiabelian}(1), which contains $\theta$ and
$\zeta_N=e^{2\pi i/N}$.
\begin{enumerate}
\item Let $S$ be a scheme over $R_{n,N}$. A \emph{full-level object} over $S$ is a pair $(G,\lambda)$ of
  an $\mathcal O_E$-semi-abelian scheme $G$ over $S$ (Definition~\ref{cv-model-toroidal-semiabelian}(2))
  with line bundle $\mathcal L$ and a map $\lambda:\Lambda\to G(S)$, such that $G$ is proper over $S$,
  $\mathcal L$ is relatively ample and defines a principal polarization, $\lambda$ is additive and
  vanishes on $N\Lambda$, and for every algebraically closed field $k$ of characteristic zero and every
  point $x:\operatorname{Spec}k\to S$, with $\beta:R_{n,N}\to k$ the induced homomorphism:
  \begin{enumerate}
  \item for every ring homomorphism $\tau:\mathcal O_E\to k$, the subspace of the Lie algebra of $G_x$
    on which each $b\in\mathcal O_E$ acts by $\tau(b)$ has dimension $p$ if
    $\tau(\sqrt{-d})=\beta(\theta)$, and $q$ otherwise;
  \item there are roots of unity $\zeta'_m\in k$ for $m\ge1$, primitive of order $m$, with
    $(\zeta'_{mm'})^{m'}=\zeta'_m$ and $\zeta'_N=\beta(\zeta_N)$, and $\mathcal O_E$-linear
    homomorphisms $\iota_m:\Lambda\to G_x(k)$ with kernel $m\Lambda$ and image the $m$-torsion, such that
    $m'\iota_{mm'}=\iota_m$, $e_m(\iota_m(a),\iota_m(b))=(\zeta'_m)^{\langle a,b\rangle}$ for the
    commutator pairing $e_m$ of the theta group of $\mathcal L^m$, and $\iota_N(a)=\lambda(a)(x)$ for
    all $a,b\in\Lambda$.
  \end{enumerate}
\item A \emph{full-level family} over $R_{n,N}$ is a triple $(\mathcal M,G,\lambda)$ of a reduced
  scheme $\mathcal M$, flat and quasi-projective over $R_{n,N}$, whose fibre over the fraction field of
  $R_{n,N}$ is smooth, and a full-level object $(G,\lambda)$ over $\mathcal M$, such that:
  \begin{enumerate}
  \item for every reduced, locally noetherian scheme $S$, flat over $R_{n,N}$, and every full-level
    object $(G',\lambda')$ over $S$, there is a unique morphism $f:S\to\mathcal M$ over $R_{n,N}$ for
    which there is a morphism $G'\to G$ of $\mathcal O_E$-semi-abelian schemes over $f$ carrying
    $\lambda'$ to $\lambda$;
  \item for every field $k$ and every point $x:\operatorname{Spec}k[\epsilon]\to\mathcal M$, where
    $\epsilon^2=0$, let $x_0$ be the composite of $\operatorname{Spec}k[\epsilon]\to\operatorname{Spec}k
    \to\operatorname{Spec}k[\epsilon]$ with $x$. If $x$ and $x_0$ have the same composite to
    $\operatorname{Spec}R_{n,N}$, and an $\mathcal O_E$-semi-abelian scheme over $k[\epsilon]$ with a map
    $\lambda':\Lambda\to G'(k[\epsilon])$ has morphisms to $G$ over $x$ and over $x_0$ that carry
    $\lambda'$ to $\lambda$, then $x=x_0$.
  \end{enumerate}
\end{enumerate}
\end{definition}
\leanref{Development/AbelianVarietyTheory/Definitions/ModuliSpaces/TopAmpExists.lean}{GlobalTheory.Models.FullObj}
\leanref{Development/AbelianVarietyTheory/Definitions/ModuliSpaces/TopAmpExists.lean}{GlobalTheory.Models.FamCore}

\begin{definition}\label{cv-toroidal-theta}
Let $g\ge1$ and $N'\ge1$ be integers, let $I$ be a finite set, and let $o$ be a ring with a unit $\zeta$
such that $N'$ is invertible in $o$ and every homomorphism from $o$ to a field maps $\zeta$ to a
primitive root of unity of order $N'^2$. Let $\omega$ be the standard symplectic form on
$\mathbf Z^{2g}$: $\omega(e_i,e_{i+g})=1$ and $\omega(e_{i+g},e_i)=-1$ for $i\le g$, and the other
$\omega(e_i,e_{i'})$ vanish.
\begin{enumerate}
\item Let $S$ be a scheme over $o$. A \emph{theta-level object} over $S$ is a tuple
  $t=(A,L,(x_i),(\alpha_h))$ of an abelian scheme $A$ over $S$ of relative dimension $g$; a line bundle
  $L$ on $A$, trivialized along the zero section, isomorphic to $[-1]^*L$ locally on $S$, relatively
  ample and defining a principal polarization; points $x_1,\dots,x_{2g}\in A(S)$ killed by $N'^2$ such
  that at every geometric point every point of order dividing $N'^2$ is $\sum_ic_ix_i$ for a unique
  $c\in(\mathbf Z/N'^2)^{2g}$, and with $e_{N'^2}(x_i,x_{i'})=\zeta^{\omega(e_i,e_{i'})}$ for the
  commutator pairing of the theta group of $L^{N'^2}$; and endomorphisms $\alpha_h$, $h\in I$, of the
  group scheme $A$. For $c\in(\mathbf Z/N'^2)^{2g}$ write $x_c=\sum_ic_ix_i$. A \emph{morphism}
  $t'\to t$ over a morphism $f:S'\to S$ is a morphism $A'\to A$ over $f$ that identifies $A'$ with
  $f^*A$ as group schemes, $L'$ with the pullback of $L$, $x'_i$ with $f^*x_i$ and $\alpha'_h$ with
  $f^*\alpha_h$.
\item Let $t$ be a theta-level object over $S$ and $y:\operatorname{Spec}R\to S$ a point.
  A family $v=(v_c)_{c\in(\mathbf Z/N'^2)^{2g}}$ in $R$ is a family of \emph{theta constants} of $t$
  at $y$ if $N'$ is odd, $N'$ is invertible in $R$ and, locally on $\operatorname{Spec}R$, there are a generator
  $\vartheta$ of the sections of $y^*L$ over $y^*A$, an isomorphism $[N']^*y^*L\cong y^*L^{N'^2}$, which
  makes $[N']^*\vartheta$ a section $\vartheta'$ of $y^*L^{N'^2}$, symmetric lifts $U_c$ of the points
  $x_c$ to the theta group of $y^*L^{N'^2}$ with $U_c^{N'^2}=1$, a trivialization of $y^*L^{N'^2}$
  along the zero section and a unit $u$, such that for every $c$, $v_c$ is $u$ times the value of
  $U_c\vartheta'$ at the origin.
\item Let $\mathrm{em}:o\to\mathbf C$ be a homomorphism and $\mathcal C$ a class of theta-level objects
  over $\operatorname{Spec}\mathbf C$, regarded over $o$ through $\mathrm{em}$. The \emph{theta ideal}
  of $\mathcal C$ is the ideal of polynomials $P$ in variables $X_c$,
  $c\in(\mathbf Z/N'^2)^{2g}$, with coefficients in $o$ such that $\mathrm{em}(P)(v)=0$ for every
  family $v$ of theta constants at the identity of every member of $\mathcal C$. Let $n\ge0$ be an
  integer, and let $N'_n$ be the integer part of $N'/n$ if $n\ge1$, and $N'_0=0$. A family $v$ in a ring
  $R$, over a homomorphism
  $\beta:o\to R$, is \emph{admissible} for $(\mathcal C,n)$ if $\beta(P)(v)=0$ for every $P$ in the
  theta ideal and, for every prime ideal $\mathfrak q$ of $R$ and every $c$, there is $c'$ with
  $v_{c+N'_nc'}\notin\mathfrak q$.
\item A theta-level object $t$ over $S$ is a $\mathcal C$-object at a point $y:\operatorname{Spec}K\to S$,
  $K$ a field, if there are a field $K_0$ with homomorphisms $o\to K_0$, $K_0\to K$ and
  $K_0\to\mathbf C$, compatible with $\mathrm{em}$ and with the structure morphism of $y$, and a
  theta-level object over $K_0$ of which $y^*t$ and a member of $\mathcal C$ are the base changes to $K$
  and to $\mathbf C$.
\item A \emph{theta family} for $(\mathcal C,n)$ is a pair $(\mathcal M,t)$ of a reduced scheme
  $\mathcal M$, flat and quasi-projective over $o$, and a theta-level object $t$ over $\mathcal M$, such
  that:
  \begin{enumerate}
  \item the points of $\mathcal M$ at which $t$ is a $\mathcal C$-object are dense;
  \item for every local ring $R$ and every point $y:\operatorname{Spec}R\to\mathcal M$, $t$ has at $y$
    a family of theta constants that is admissible for $(\mathcal C,n)$;
  \item for every ring $R$, every homomorphism $\beta:o\to R$ and every family $v$ in $R$ admissible
    for $(\mathcal C,n)$ over $\beta$, there is a unique point $y:\operatorname{Spec}R\to\mathcal M$
    over $\operatorname{Spec}\beta$ at which $v$ is a family of theta constants of $t$;
  \item for every reduced scheme $S$ over $o$ and every theta-level object $t'$ over $S$ whose
    $\mathcal C$-points are dense, there is a unique morphism $f:S\to\mathcal M$ over $o$ with a morphism
    $t'\to t$ over $f$;
  \item for every field $K$ and every point $y:\operatorname{Spec}K[\epsilon]\to\mathcal M$ with the
    same composite to $\operatorname{Spec}o$ as the constant point $y_0$ (as in
    Definition~\ref{cv-toroidal-family}(2)(b)), if a theta-level object over $K[\epsilon]$ has morphisms to
    $t$ over $y$ and over $y_0$, then $y=y_0$.
  \end{enumerate}
\item In Situation~\ref{cv-function-situation}, let $\Lambda$ be the lattice of the PEL datum with its
  alternating form $\langle\,,\rangle$, and let $e_1,\dots,e_{2g}$ be a basis of $\Lambda$ over $\mathbf Z$.
  The basis is \emph{symplectic} if $\langle e_i,e_{i'}\rangle=\omega(e_i,e_{i'})$ for all $i$ and $i'$.
  Let $a:I\to\mathcal O_E$ be a map and $\rho\in o$. A
  \emph{PEL structure} on a theta-level object $t$ over a scheme $S$ over $o$ is an
  $\mathcal O_E$-semi-abelian scheme $G$ over $S$ (Definition~\ref{cv-model-toroidal-semiabelian}(2))
  with an isomorphism of group schemes $G\to A$ that carries the line bundle of $G$ to $L$ and the action
  of $a(h)$ to $\alpha_h$ for every $h$, such that for every algebraically closed field $k$ of
  characteristic zero and every point $x:\operatorname{Spec}k\to S$, with $\beta:o\to k$ the induced
  homomorphism, condition (a) of Definition~\ref{cv-toroidal-family}(1) holds with $\beta(\rho)$ in
  place of $\beta(\theta)$, and there are $\zeta'_m$ and $\iota_m$ as in condition (b) there with
  $\zeta'_{N'^2}=\beta(\zeta)$ and $\iota_{N'^2}(e_i)=x_i(x)$ for every $i$. The \emph{PEL class} of
  $((e_i),a,\rho,\mathrm{em})$ is the class of theta-level objects over $\operatorname{Spec}\mathbf C$,
  through $\mathrm{em}$, that carry a PEL structure.
\item In Situation~\ref{cv-function-situation}, let $n$ be an odd prime, $j\ge1$, $N'=n^j$ and
  $N=N'^2$. A \emph{PEL theta family} of level $N'$ is a triple $(\mathcal M,t,G)$ such that, for some
  symplectic basis $(e_i)$ of $\Lambda$ and some map $a:I\to\mathcal O_E$ from a finite set whose image
  generates $\mathcal O_E$ as a ring and contains $w_0=(1+\sqrt{-d})/2$ and $\bar w_0$, the pair
  $(\mathcal M,t)$ is a theta family for $n$ and the PEL class of $((e_i),a,\theta,\mathrm{em})$, with
  $o=R_{n,N}$, $\zeta=\zeta_N$ and $\mathrm{em}$ the inclusion; $G$ is a PEL structure on $t$; and the
  fibre of $\mathcal M$ over the fraction field of $R_{n,N}$ is smooth of relative dimension
  $\frac12[E:\mathbf Q]pq$.
\end{enumerate}
\end{definition}
\leanref{Development/AbelianVarietyTheory/Definitions/ModuliSpaces/GoodClass.lean}{GlobalTheory.ZM.TObj}
\leanref{Development/AbelianVarietyTheory/Definitions/ModuliSpaces/GoodClass.lean}{GlobalTheory.ZM.TMap}
\leanref{Development/AbelianVarietyTheory/Definitions/ModuliSpaces/GoodClass.lean}{GlobalTheory.ZM.TObj.HasTh}
\leanref{ReciprocityTheory/ReductiveGroups/ArithmeticGeometry/ExtensionDescent/ThetaPt.lean}{ReciprocityTheory.ChartRestriction.G.Out.IsThetaPt}
\leanref{Development/AbelianVarietyTheory/Definitions/ModuliSpaces/GoodClass.lean}{GlobalTheory.ZM.thIdeal}
\leanref{Development/AbelianVarietyTheory/Definitions/ModuliSpaces/GoodClass.lean}{GlobalTheory.ZM.InMM}
\leanref{Development/AbelianVarietyTheory/Definitions/ModuliSpaces/GoodClass.lean}{GlobalTheory.ZM.FromC}
\leanref{Development/AbelianVarietyTheory/Definitions/ModuliSpaces/GoodClass.lean}{GlobalTheory.ZM.ThetaFam}
\leanref{Development/AbelianVarietyTheory/Definitions/ModuliSpaces/ExistsArBaseAway.lean}{GlobalTheory.ZM.IsSympBasis}
\leanref{Development/AbelianVarietyTheory/Definitions/ModuliSpaces/ExistsArBaseAway.lean}{GlobalTheory.ZM.IsPEL}
\leanref{Development/AbelianVarietyTheory/Definitions/ModuliSpaces/ExistsArBaseAway.lean}{GlobalTheory.ZM.PELClass}

\begin{theorem}[Mumford \cite{Mum72}, Faltings--Chai \cite{FC90}]\label{cv-toroidal-mumford}
Let $V$ be a complete discrete valuation ring of characteristic zero with maximal ideal $\mathfrak m$
and fraction field $K$, and let $A$ be an abelian variety over $K$ with a line bundle $L$ that is
trivialized at the origin, symmetric, ample and defines a principal polarization. Write
$\Lambda(L)=m^*L\otimes\operatorname{pr}_1^*L^{-1}\otimes\operatorname{pr}_2^*L^{-1}$ on $A\times A$.
\begin{enumerate}
\item Suppose that $A$ has endomorphisms $w$ and $w'$ with $w+w'=1$, such that $L$ is isomorphic to the
  pullback of $\Lambda(L)$ along $(1,w):A\to A\times A$ and the pullbacks of $\Lambda(L)$ along
  $w\times1$ and $1\times w'$ are isomorphic. Let $G$ be a semi-abelian scheme over $V$ whose fibre over
  the closed point is an extension of an abelian variety by a split torus, and let $\gamma:G_K\to A$ be
  an isomorphism of group schemes. Then there are degeneration data over $V$ and a uniformization of
  $(A,G)$ by Mumford's construction from them that extends $\gamma$.
\item Let $G$ be a semi-abelian scheme over $V$ with degeneration data over $V$ and a uniformization of
  $(A,G)$ by them. Then for every algebraically closed field $F\supset K$ and every $m\ge1$ there is a
  map from the $m$-division points of the degeneration data over $F$ onto the $m$-torsion of $A(F)$,
  given by the uniformization, whose fibres are the classes of the equivalence of division points, and
  which is compatible with addition and carries the pairing of the data to the Weil pairing of $L$.
\end{enumerate}
Here degeneration data consist of an abelian scheme $B$ over $V$ with a principal symmetric line bundle
trivialized at the origin, a free $\mathbf Z$-module $X$ of finite rank with a homomorphism
$c:X\to B(V)$, and trivializations $\psi(\chi)$ over $K$ of the pullbacks along $c(\chi)$ of the
inverse of the line bundle, such that $\psi$ is even, $\psi(\chi+\chi')\psi(\chi)^{-1}\psi(\chi')^{-1}$
is bimultiplicative, $\psi(\chi)$ is integral and lies in $\mathfrak m$ for $\chi\ne0$, and for every
$e$ only finitely many $\psi(\chi)$ lie outside $\mathfrak m^e$. A uniformization by them identifies the
reduction of $G$ modulo each $\mathfrak m^{e+1}$, compatibly with the line bundles, with an extension
of the reduction of $B$ by the split torus with character group $X$, and for every $m\ge1$ it
identifies the sections of $L^m$ with their Fourier expansions in terms of the data.
\end{theorem}
\leanref{Development/AbelianVarietyTheory/Statements/SemistableTorsion/UTorsSt.lean}{GlobalTheory.ZH.ThmUSt}
\leanref{Development/AbelianVarietyTheory/Statements/SemistableTorsion/UTorsSt.lean}{GlobalTheory.ZH.UTorsSt}
\leanref{Development/AbelianVarietyTheory/Proofs/ModuliProofs/UniformisationExistenceAndTorsionStatements.lean}{GlobalTheory.ZH.thmU_and_uTors}

This is Mumford's construction of degenerating abelian varieties over complete bases \cite{Mum72} and
its converse, the theory of degenerations of Faltings and Chai \cite[Chapters~II and~III]{FC90}, by
which every split degeneration of a principally polarized abelian variety comes from degeneration
data, with its torsion points described by them.

\subsection{Theta constants of theta-level objects}\label{subsection-cv-theta-constants}

\begin{situation}\label{cv-toroidal-reduction-situation}
Let $g\ge1$ and $N'\ge1$ be integers, let $I$ be a finite set, and let $o$ be a ring with a unit $\zeta$
such that $N'$ is invertible in $o$ and every homomorphism from $o$ to a field maps $\zeta$ to a primitive
root of unity of order $N'^2$. Theta-level objects, their morphisms and their theta constants are those
of Definition~\ref{cv-toroidal-theta} for these data, and $x_c=\sum_ic_ix_i$ for
$c\in(\mathbf Z/N'^2)^{2g}$.
\end{situation}

\begin{lemma}\label{cv-toroidal-reduction-semistable}
Let $V$ be a complete discrete valuation ring with finite residue field and with fraction field $K$ of
characteristic zero, and let $A$ be an abelian variety over $K$ with a line bundle that is trivialized at
the origin, symmetric, ample and defines a principal polarization. Then there are a complete discrete
valuation ring $W$, finite over $V$ through a local homomorphism $V\to W$, an extension $K\to L$ of
$V\to W$ to the fraction field $L$ of $W$, and a semi-abelian scheme $G$ over $W$, such that $G_L$ is
isomorphic to $A_L$ as a group scheme and the closed fibre of $G$ is an extension of an abelian variety by
a split torus.
\end{lemma}
\leanref{Development/AbelianVarietyTheory/Statements/SemistableTorsion/SemistableFin.lean}{GlobalTheory.SS.SemistablePol}
\leanref{Development/AbelianVarietyTheory/Proofs/ModuliProofs/SemistableReductionSplitOfPadicStatement.lean}{GlobalTheory.T.semistableReductionSplit_of_padic}
\leanref{Development/AbelianVarietyTheory/Proofs/ModuliProofs/SemistableReductionSplitOfPadic.lean}{GlobalTheory.SS.semistablePol_of_padic}
\leanref{Development/AbelianVarietyTheory/Proofs/DegenerationProofs/ExistsSemiAbelianModelSplitSpecialFibre.lean}{GlobalTheory.exists_semiAbelian_model_split_specialFibre}
\begin{proof}
Let $p$ be the characteristic of the residue field of $V$. As $V$ is complete with finite residue field
and $K$ has characteristic zero, $K$ is a finite extension of $\mathbf Q_p$ with ring of integers $V$
\cite{Ser79}. So there are a finite extension $K_0$ of $\mathbf Q_p$ in an algebraic closure
$\bar{\mathbf Q}_p$ and an isomorphism $K\to K_0$ that maps $V$ onto the integral closure $O_{K_0}$ of
$\mathbf Z_p$ in $K_0$; transport $A$ and its line bundle to $K_0$ along it. Theorem~\ref{cv-semistable-reduction}
gives a finite extension $K_1$ of $K_0$ in $\bar{\mathbf Q}_p$ and a semi-abelian scheme $G_1$ over
$O_{K_1}$ whose generic fibre is isomorphic to $A_{K_1}$ as a group scheme.

Let $\bar k$ be the residue field of the integral closure of $\mathbf Z_p$ in $\bar{\mathbf Q}_p$, an
algebraic closure of the residue field of $O_{K_1}$. Since $G_1$ is semi-abelian, its closed fibre becomes
over $\bar k$ an extension of an abelian variety by a split torus, and this extension is defined over
the subfield generated by a finite subset $s$ of $\bar k$: for every local ring $R'$ with a local
homomorphism $O_{K_1}\to R'$ and an embedding of the residue field of $R'$ into $\bar k$, compatible with
the one of the residue field of $O_{K_1}$, whose image contains $s$, the closed fibre of
$G_1\otimes_{O_{K_1}}R'$ is an extension of an abelian variety by a split torus. Let $K'$ be the
extension of $K_1$ in $\bar{\mathbf Q}_p$ generated by lifts of the elements of $s$ to integral elements
of $\bar{\mathbf Q}_p$, a finite extension. Then $W=O_{K'}$ is a complete discrete valuation ring with
fraction field $L=K'$, finite over $O_{K_0}$, hence over $V$, through a local homomorphism, and the
residue field of $W$ contains $s$. So $G=G_1\otimes_{O_{K_1}}W$ is semi-abelian, its closed fibre is an
extension of an abelian variety by a split torus, and $G_L$ is isomorphic to $A_L$ as a group scheme.
\end{proof}
% Lean: semistablePol_of_padic: completeDVR_eq_integralClosure_padicInt (V and K as O_{K_0} and K_0
% inside PadicAlgCl p, and the converse for K'), polBaseChange_holds (transport of A), and
% exists_semiAbelian_model_split_specialFibre (SemistableReductionSt, then
% exists_finite_torusExtension_specialFibre with FieldModel.exists_finite_torusExtension_range for s,
% and exists_subset_range for K').

\begin{lemma}\label{cv-toroidal-reduction-base-change}
In Situation~\ref{cv-toroidal-reduction-situation}, let $S$ and $S'$ be schemes over $o$, let
$f:S'\to S$ be a morphism over $o$, and let $t=(A,L,(x_i),(\alpha_h))$ be a theta-level object over $S$.
\begin{enumerate}
\item There is a theta-level object $t'$ over $S'$ with a morphism $t'\to t$ over $f$.
\item Let $t'=(A',L',(x'_i),(\alpha'_h))$ be a theta-level object over $S'$ with a morphism $t'\to t$
  over $f$, and let $w,w'\in I$. With $\Lambda(L)$ as in Theorem~\ref{cv-toroidal-mumford}, suppose that
  $\alpha_w+\alpha_{w'}=1$, that $L$ is isomorphic to the pullback of $\Lambda(L)$ along
  $(1,\alpha_w):A\to A\times_SA$, and that the pullbacks of $\Lambda(L)$ along $\alpha_w\times1$ and
  $1\times\alpha_{w'}$ are isomorphic. Then the same holds for $L'$, $\alpha'_w$ and $\alpha'_{w'}$.
\item Suppose that $N'$ is odd, and let $t'$ be a theta-level object over $S'$ with a morphism $t'\to t$
  over $f$. Let $R$ be a ring and $y:\operatorname{Spec}R\to S'$ a point. A family $v$ in $R$ is a
  family of theta constants of $t'$ at $y$ if and only if it is a family of theta constants of $t$ at
  $f\circ y$.
\item Suppose that $N'$ is odd, let $v$ be a family of theta constants of $t$ at a point
  $y:\operatorname{Spec}R\to S$, and let $\sigma:R\to R'$ be a ring homomorphism. Then
  $\sigma(v)=(\sigma(v_c))_c$ is a family of theta constants of $t$ at $y\circ\operatorname{Spec}\sigma$.
\item Suppose that $N'$ is odd, let $v$ be a family of theta constants of $t$ at a point
  $y:\operatorname{Spec}R\to S$, and let $r$ be a unit of $R$. Then $rv=(rv_c)_c$ is a family of theta
  constants of $t$ at $y$.
\end{enumerate}
\end{lemma}
\leanref{Development/AbelianVarietyTheory/Proofs/ModuliProofs/ExistsThetaLevelAbelianSchemePullback.lean}{GlobalTheory.tObj_pull}
\leanref{Development/AbelianVarietyTheory/Proofs/ModuliProofs/ThetaLevelShapePrimitivityPullbackDescent.lean}{GlobalTheory.shape_prim_pull_and_descend}
\leanref{Development/AbelianVarietyTheory/Proofs/ModuliProofs/HasThetaConstantsIffOfThetaLevelPullback.lean}{GlobalTheory.hasTh_iff_of_tmap}
\leanref{Development/AbelianVarietyTheory/Proofs/ModuliProofs/HasThetaConstantsRingBaseChangeStatement.lean}{GlobalTheory.T.hasThetaConstants_ring_baseChange}
\leanref{Development/AbelianVarietyTheory/Proofs/ModuliProofs/HasThetaConstantsRingBaseChange.lean}{GlobalTheory.hasTh_baseChange}
\leanref{Development/AbelianVarietyTheory/Proofs/ModuliProofs/GoodReductionOfThetaPointUnits.lean}{GlobalTheory.ThetaCrit.hasTh_unit_mul}
\begin{proof}
(1) Let $A'=A\times_SS'$, let $\varphi:A'\to A$ be the projection, and pull back $L$, the points $x_i$
and the endomorphisms $\alpha_h$ along $\varphi$. Base change preserves the relative dimension, the
trivialization along the zero section, the local isomorphism with $[-1]^*L$, relative ampleness and the
principal polarization, and the pulled-back points are killed by $N'^2$. At a geometric point $y$ of
$S'$, $\varphi$ identifies the fibre of $A'$ at $y$, with its points, with the fibre of $A$ at $f\circ y$;
so the points of order dividing $N'^2$ have the required description. The commutator pairing of the
theta group of $L^{N'^2}$ is compatible with base change, and $\zeta$ is the same element of $o$; so the
pairing condition holds. Then $\varphi$ is a morphism $t'\to t$ over $f$.

(2) The morphism identifies $A'$ with $f^*A$, $\alpha'_h$ with $f^*\alpha_h$ and $L'$ with the pullback
of $L$. So $\alpha'_w+\alpha'_{w'}$ is the pullback of $\alpha_w+\alpha_{w'}=1$, and $\Lambda(L')$ is
the pullback of $\Lambda(L)$ along $\varphi\times\varphi$. The pullbacks of the two given isomorphisms of
line bundles are the required ones.

(3) The data of Definition~\ref{cv-toroidal-theta}(2) for $t'$ at $y$ live on $y^*A'$, with $y^*L'$ and
the points $y^*x'_c$. The morphism $t'\to t$ identifies these with $(f\circ y)^*A$, $(f\circ y)^*L$ and
the points $(f\circ y)^*x_c$, since $x'_c$ is the pullback of $x_c$. So data for either object at its
point give data for the other, with the same values at the origin.

(4) Pull back the data of Definition~\ref{cv-toroidal-theta}(2) along
$\operatorname{Spec}R'\to\operatorname{Spec}R$. The pullback of the generator $\vartheta$ still
generates the sections of the pullback of $L$, since the formation of the sections of a line bundle
that defines a principal polarization commutes with base change. The other data pull back, $N'$ stays
invertible, and the values at the origin are carried to their images under $\sigma$.

(5) Replace the unit $u$ of Definition~\ref{cv-toroidal-theta}(2) by $ru$.
\end{proof}
% Lean: (1) tObj_pull (pol_pull, rigidified_pull, sym_pull, pushHom_injective for the torsion, t.bij at
% x >> f, weilIs_grpPullback_of_comm with CommPairing.weilIs_pull and unitS for the pairing). (2) the
% second conjunct of shape_prim_pull_and_descend (eq_up with push_mul, iso_up with iso_cubeLB_pull).
% (3) hasTh_iff_of_tmap (ThetaPtUp.isThetaPt_iff with lab_over). (4) hasTh_baseChange
% (ThetaBaseChange.isThetaPt_pull with SecBaseChangeAnySt of Out.Laws2). (5) hasTh_unit_mul.

\begin{lemma}[Alp\"oge]\label{cv-toroidal-reduction-local}
In Situation~\ref{cv-toroidal-reduction-situation}, suppose that $N'$ is odd, let $S$ be a scheme over
$o$, and let $t=(A,L,(x_i),(\alpha_h))$ be a theta-level object over $S$. Then every point of $S$ lies in
the image of an open immersion $y:\operatorname{Spec}R\to S$ at which $t$ has a family of theta
constants.
\end{lemma}
\leanref{Development/AbelianVarietyTheory/Proofs/ModuliProofs/ExistsThetaConstantsLocallyOnBase.lean}{GlobalTheory.ThetaLocal.exists_hasTh}
\begin{proof}
Let $p$ be a point of $S$. The integer $N'$ is invertible on $S$, as it is in $o$. Since $L$ is
trivialized along the zero section and isomorphic to $[-1]^*L$ locally on $S$, the theorem of the cube
\cite{Mum70} gives an isomorphism $[N']^*L\cong L^{N'^2}$ on $A$, and the trivialization of $L$ along the zero section
gives one of $L^{N'^2}$. Since $L$ defines a principal polarization, there is an open neighbourhood $V_0$
of $p$ over which the sections of $L$ are generated by one section whose pullback under every base
change of $V_0$ still generates the sections. Since $N'$ is odd and invertible on $S$, each point $x_c$,
which is killed by $N'^2$, has over an open neighbourhood $V_c$ of $p$ a symmetric lift $U_c$ to the
theta group of $L^{N'^2}$ with $U_c^{N'^2}=1$ \cite{Mum66}. Let $y:\operatorname{Spec}R\to S$ be the
inclusion of an affine open neighbourhood of $p$ contained in $V_0$ and in every $V_c$. Over it these
data, with the unit $u=1$, satisfy Definition~\ref{cv-toroidal-theta}(2), and the values at the origin
of the sections $U_c\vartheta'$ are a family of theta constants of $t$ at $y$.
\end{proof}
% Lean: exists_hasTh = exists_affine_isThetaPt: MulPull.mulPull_of_cube_of_symm (from the theorem of
% the cube AbLaw.CubeAssembly.cube_gen and exists_gen_negLB) for [N']^*L = L^{N'^2}, the rigidification
% for the trivialization, SecBaseChangeAnySt with exists_generator for the generator, symLift_odd for
% the symmetric lifts, and exists_isThetaPt, exists_isThetaPtOver with u = 1.

\begin{lemma}[Alp\"oge]\label{cv-toroidal-reduction-coset}
In Situation~\ref{cv-toroidal-reduction-situation}, suppose that $N'$ is odd, and let $n\ge2$ be a
divisor of $N'$. Let $S$ be a scheme over $o$, let $t$ be a theta-level object over $S$, and let $v$ be a
family of theta constants of $t$ at a point $y:\operatorname{Spec}R\to S$. Then for every prime ideal
$\mathfrak q$ of $R$ and every $c\in(\mathbf Z/N'^2)^{2g}$ there is $c'$ with
$v_{c+(N'/n)c'}\notin\mathfrak q$.
\end{lemma}
\leanref{Development/AbelianVarietyTheory/Proofs/ModuliProofs/ThetaCosetNonvanishingOfHasThStatement.lean}{GlobalTheory.T.thetaCosetNonvanishing_of_hasTh}
\leanref{Development/AbelianVarietyTheory/Proofs/ModuliProofs/ThetaCosetNonvanishingOfHasTh.lean}{GlobalTheory.coset_of_hasTh}
\begin{proof}
Let $F$ be an algebraic closure of the residue field of $\mathfrak q$, and $\sigma:R\to F$ the induced
homomorphism. By Lemma~\ref{cv-toroidal-reduction-base-change}(4), $\sigma(v)$ is a family of theta
constants of $t$ at $y\circ\operatorname{Spec}\sigma$, and $\sigma(v_{c''})\ne0$ implies
$v_{c''}\notin\mathfrak q$. So we may assume that $R=F$ is an algebraically closed field and
$\mathfrak q=0$. Let $(A,L,(x_i))$ be the fibre of $t$ at $y$. As $\operatorname{Spec}F$ is a point, the
data of Definition~\ref{cv-toroidal-theta}(2) exist over all of it: a generator $\vartheta$ of the
sections of $L$, which is nonzero; the section $\vartheta'$ of $L^{N'^2}$ given by $[N']^*\vartheta$;
symmetric lifts $U_c$ of the points $x_c$; a trivialization of $L^{N'^2}$ at the origin; and a unit $u$,
with $v_c$ equal to $u$ times the value of $U_c\vartheta'$ at the origin. As $U_c$ lies over $x_c$, this
value is nonzero exactly when $\vartheta'$ does not vanish at $x_c$, that is, when $\vartheta$ does not
vanish at $N'x_c$.

The integer $n$ divides $N'$, which is invertible in $F$. So there is a point $z\in A(F)$ with $nz=0$
and $\vartheta(N'x_c+z)\ne0$: the zero locus of $\vartheta$ contains no translate $b+A[n]$. Indeed,
suppose that it contains $b+A[n]$, and choose $a_0$ with $na_0=b$. The sections of $L^{n^2}$ that vanish
on $a_0+A[n^2]$ form a subspace stable under the theta group of $L^{n^2}$, whose elements lie over points
of $A[n^2]$. By the irreducibility of the action of the theta group \cite{Mum66}, this subspace is $0$ or
all the sections. It is not all of them, as $L^{n^2}$ has a section that does not vanish at $a_0$
\cite{Mum70}. The isomorphism $[n]^*L\cong L^{n^2}$ makes $[n]^*\vartheta$ a section of $L^{n^2}$, and it
vanishes on $a_0+A[n^2]$, since $n(a_0+w)=b+nw\in b+A[n]$ when $n^2w=0$. So $[n]^*\vartheta=0$, and
$\vartheta=0$ as $[n]$ is flat and surjective; a contradiction.

The point $z$ is killed by $N'^2$, so $z=x_e$ for a unique $e\in(\mathbf Z/N'^2)^{2g}$, and
$x_{ne}=nz=0=x_0$ gives $ne=0$. As $N'^2=n\cdot N'(N'/n)$, $e=N'(N'/n)c'$ for some $c'$. Then
$N'x_{c+(N'/n)c'}=N'x_c+x_e=N'x_c+z$, so $\vartheta$ does not vanish at $N'x_{c+(N'/n)c'}$ and
$v_{c+(N'/n)c'}\ne0$.
\end{proof}
% Lean: coset_of_hasTh (hasThetaConstants_ring_baseChange to an algebraic closure of the residue field
% of q) and ThetaCoset.coset_bot: isThetaPtOver_of_subsingleton, exists_ne_zero_over (the point z from
% UT.ExactTerm.exists_good_point, which is the law GoodPtSt of Out.Laws, applied to N' x_c; the value
% identities sv_act, sv_mulPull and sv_of_isValueAt), ne_zero_of_isGenSec, and exists_label with
% exists_eq_smul for c'. The nonvanishing on a coset of A[n] is the law GoodPtSt of Out.Laws, a
% hypothesis of coset_of_hasTh that the archive proves (GlobalTheory.abelianSchemeLaws_all, from
% AbLaw.AllLaws.laws_of_h0_eq_one; the goodPt field is GoodPoint.goodPt_of_irred_of_mulPull), so the
% statement needs no further hypothesis. Its proof there is the argument given here: the subspace
% vanishing P.L B for B = b + A[n^2] is stable under the theta group (IrredSt gives 0 or everything), the
% base-point freeness exists_sec_value_isUnit excludes everything, and [n]^* th (Laws.mulPull) lies in it,
% so th = 0 by sec_eq_zero_of_flat.

\begin{lemma}[Alp\"oge]\label{cv-toroidal-reduction-unit}
In Situation~\ref{cv-toroidal-reduction-situation}, suppose that $N'$ is odd, let $S$ be a scheme over
$o$, let $t$ be a theta-level object over $S$, and let $v$ and $v'$ be families of theta constants of $t$
at a point $y:\operatorname{Spec}R\to S$. If the $v_c$ generate the unit ideal of $R$, then there is a
unit $u$ of $R$ with $v'_c=uv_c$ for every $c$.
\end{lemma}
\leanref{Development/AbelianVarietyTheory/Proofs/ModuliProofs/ThetaConstantsUniqueUpToUnitStatement.lean}{GlobalTheory.T.thetaConstants_unique_up_to_unit}
\leanref{Development/AbelianVarietyTheory/Proofs/ModuliProofs/ThetaConstantsUniqueUpToUnit.lean}{GlobalTheory.hasTh_unique_unit}
\begin{proof}
First let $R$ be local. The data of Definition~\ref{cv-toroidal-theta}(2) for $v$ exist on an open
subset of $\operatorname{Spec}R$ containing the closed point, hence on all of it; so do those for $v'$.
Let $\vartheta,\iota,(U_c),\tau,u_0$ and $\vartheta_1,\iota_1,(U^1_c),\tau_1,u_1$ be these data, where
$\iota$ and $\iota_1$ are the isomorphisms $[N']^*y^*L\cong y^*L^{N'^2}$ and $\tau,\tau_1$ the
trivializations at the origin. The generators $\vartheta$ and $\vartheta_1$ of the sections of $y^*L$
differ by a unit of $R$: if $\vartheta_1=r\vartheta$ and $\vartheta=r'\vartheta_1$ with $r$ not a unit,
then $1-rr'$ is a unit and $(1-rr')\vartheta=0$, so both vanish and any unit will do. The quotient of
$\iota_1$ by $\iota$ is an invertible function on $y^*A$, hence a unit of $R$, since the global functions
on an abelian scheme are those of its base; and $\tau_1$ and $\tau$ differ by a unit of $R$. The symmetric
lifts agree, $U^1_c=U_c$, because $N'$ is odd and a point killed by $N'^2$ has only one symmetric lift
$U$ with $U^{N'^2}=1$. So the values at the origin of $U^1_c\vartheta'_1$ are those of $U_c\vartheta'$
times one unit of $R$, the same for all $c$, and $v'_c=uv_c$ for a unit $u$.

In general, let $\mathfrak m$ be a maximal ideal of $R$. By
Lemma~\ref{cv-toroidal-reduction-base-change}(4), the images of $v$ and $v'$ in $R_{\mathfrak m}$ are
families of theta constants of $t$ at the composite of $\operatorname{Spec}R_{\mathfrak m}\to
\operatorname{Spec}R$ with $y$, so there is a unit $u_{\mathfrak m}$ of $R_{\mathfrak m}$ with
$v'_c=u_{\mathfrak m}v_c$ in $R_{\mathfrak m}$ for every $c$. Write $1=\sum_cr_cv_c$ with $r_c\in R$, and
put $u=\sum_cr_cv'_c$. In each $R_{\mathfrak m}$, $u=u_{\mathfrak m}\sum_cr_cv_c=u_{\mathfrak m}$, a
unit. So $u$ lies in no maximal ideal and is a unit of $R$, and $v'_c=uv_c$ holds in every
$R_{\mathfrak m}$, hence in $R$.
\end{proof}
% Lean: hasTh_unique_unit: ThetaUniq.exists_unit_local at each maximal ideal (isThetaPtOver_of_local;
% exists_unit_over with exists_unit_of_isGenSec_local, exists_ratio_base and SymLiftOdd.surj_fst for the
% units, SymLiftOdd.symLift_unique for the lifts), glued by ThetaUniq.exists_unit_of_local.

\begin{lemma}[Alp\"oge]\label{cv-toroidal-reduction-generic}
In Situation~\ref{cv-toroidal-reduction-situation}, suppose that $N'$ is odd, and let $n\ge2$ be a
divisor of $N'$. Let $W$ be a discrete valuation ring with fraction field $K$, with
$\operatorname{Spec}W$ and $\operatorname{Spec}K$ schemes over $o$, let $t'$ be a theta-level object over
$W$ and $t$ a theta-level object over $K$, with a morphism $t\to t'$ over
$\operatorname{Spec}K\to\operatorname{Spec}W$. Let $v$ be a family in $W$ whose image in $K$ is a family
of theta constants of $t$ at the identity, and suppose that some $v_c$ is a unit. Then $v$ is a family of
theta constants of $t'$ at the identity.
\end{lemma}
\leanref{Development/AbelianVarietyTheory/Proofs/ModuliProofs/GoodReductionOfThetaPointUnits.lean}{GlobalTheory.ThetaCrit.hasTh_of_generic}
\begin{proof}
Let $p$ be the closed point of $\operatorname{Spec}W$. Lemma~\ref{cv-toroidal-reduction-local} gives an
open immersion $y:\operatorname{Spec}R\to\operatorname{Spec}W$ whose image contains $p$, and a family
$v^0$ of theta constants of $t'$ at $y$. Every point of $\operatorname{Spec}W$ specializes to $p$ and the
image of $y$ is open, so $y$ is surjective, hence an isomorphism; let $\sigma:R\to W$ be the homomorphism
with $y\circ\operatorname{Spec}\sigma=\mathrm{id}$. By Lemma~\ref{cv-toroidal-reduction-base-change}(4),
$\sigma(v^0)$ is a family of theta constants of $t'$ at the identity, and its image in $K$ is one at
$\operatorname{Spec}K\to\operatorname{Spec}W$. By Lemma~\ref{cv-toroidal-reduction-base-change}(3),
applied to the morphism $t\to t'$, the image of $v$ in $K$ is also a family of theta constants of $t'$ at
$\operatorname{Spec}K\to\operatorname{Spec}W$.

Lemma~\ref{cv-toroidal-reduction-coset}, for the maximal ideal of $W$ and $c=0$, gives
$c_1=(N'/n)c'$ such that $e=\sigma(v^0_{c_1})$ is a unit of $W$. So the image of $\sigma(v^0)$ in $K$
generates the unit ideal, and Lemma~\ref{cv-toroidal-reduction-unit} gives $u\in K^\times$ with
$v_c=u\,\sigma(v^0_c)$ in $K$ for every $c$. Then $v_c=(v_{c_1}e^{-1})\,\sigma(v^0_c)$ in $W$ for every
$c$, since both sides have the same image in $K$:
$v_{c_1}e^{-1}\sigma(v^0_c)=u\,\sigma(v^0_{c_1})e^{-1}\sigma(v^0_c)=u\,\sigma(v^0_c)$. As some $v_c$ is a
unit, $r=v_{c_1}e^{-1}$ divides a unit, so it is a unit of $W$. By
Lemma~\ref{cv-toroidal-reduction-base-change}(5), $v=r\,\sigma(v^0)$ is a family of theta constants of
$t'$ at the identity.
\end{proof}

\begin{lemma}[Alp\"oge]\label{cv-toroidal-reduction-model}
In Situation~\ref{cv-toroidal-reduction-situation}, let $W$ be a discrete valuation ring with fraction
field $K$, let $\beta:o\to W$ be a homomorphism, and let $t=(A,L,(x_i),(\alpha_h))$ be a theta-level object
over $K$, regarded over $o$ through $\beta$. Let $P$ be an abelian scheme over $W$ with a line bundle $M$
that is trivialized along the zero section, isomorphic to $[-1]^*M$ locally on $W$, relatively ample and
defines a principal polarization, and let $\chi:A\to P$ be a morphism over
$\operatorname{Spec}K\to\operatorname{Spec}W$ that identifies $A$ with $P_K$ as group schemes, with $L$
isomorphic to $\chi^*M$. Then there is a theta-level object $t'$ over $W$, regarded over $o$ through
$\beta$, with a morphism $t\to t'$ over $\operatorname{Spec}K\to\operatorname{Spec}W$.
\end{lemma}
\leanref{Development/AbelianVarietyTheory/Proofs/ModuliProofs/GoodReductionOfThetaPointUnits.lean}{GlobalTheory.ThetaCrit.exists_obj_model}
\leanref{Development/AbelianVarietyTheory/Proofs/ModuliProofs/ThetaLevelEndomorphismsSpreadOverBasicOpen.lean}{GlobalTheory.SpreadLevel.exists_extension}
\begin{proof}
The abelian scheme $P$ is proper over $W$, so by the valuative criterion of properness each point
$\chi\circ x_i$ of $P(K)$ extends to a point $p_i\in P(W)$. By Theorem~\ref{cv-model-toroidal-extension},
applied to $P$ and $P$, each endomorphism of $P_K$ that corresponds to $\alpha_h$ under $\chi$ extends to
an endomorphism $a_h$ of $P$, so that $\chi\circ\alpha_h=a_h\circ\chi$. The relative dimension of $P$ is
$g$, as that of its generic fibre $A$ is $g$ and $\operatorname{Spec}W$ is connected. We show that
$t'=(P,M,(p_i),(a_h))$ is a theta-level object over $W$; then $\chi$ is a morphism $t\to t'$ over
$\operatorname{Spec}K\to\operatorname{Spec}W$.

The integer $N'$ is invertible in $W$, as it is in $o$, and $\beta(\zeta)^{N'^2}=1$ in $W$, since its
image in $K$ is a root of unity of order $N'^2$ and $W\to K$ is injective. The points $N'^2p_i$ and $0$
of $P(W)$ agree over the dense open subscheme $\operatorname{Spec}K$, so they agree, $P$ being
separated. The subgroup scheme $P[N'^2]$ is finite and \'etale over $W$
(Lemma~\ref{cv-toroidal-moduli-torsion}), and the points $p_c=\sum_ic_ip_i$, for
$c\in(\mathbf Z/N'^2)^{2g}$, are sections of it, hence open and closed immersions. At a geometric point
over the generic point, every point of $P[N'^2]$ lies on exactly one of them, by the corresponding
property of $t$. The points of $P[N'^2]$ that lie on none of the $p_c$, or on two of them, form an open
subset, whose image in $\operatorname{Spec}W$ is open, as $P[N'^2]$ is \'etale, and misses the generic
point; so it is empty. Hence at every geometric point of $\operatorname{Spec}W$ every point of order
dividing $N'^2$ is $p_c$ for a unique $c$. For each $i$ and
$i'$, the commutator pairing of $p_i$ and $p_{i'}$ and $\beta(\zeta)^{\omega(e_i,e_{i'})}$ are roots of
unity of order dividing $N'^2$, so they agree on an open and closed subset of $\operatorname{Spec}W$. It
contains the generic point, by the pairing condition for $t$, so it is all of the connected scheme
$\operatorname{Spec}W$.
\end{proof}
% Lean: exists_obj_model: ValuativeCriterion.Existence for the points,
% abelianScheme_hom_extends_dedekind_of_isGrpPullback with HomExtSt (T.semiAbelianHomExtension_of_laws)
% for the endomorphisms, CM.smoothOfRelativeDimension_of_fibre for g, and SpreadLevel.exists_extension:
% ExtUnique.hom_ext_of_dominant for the torsion, torsion_basis_of_dense
% (exists_finiteEtale_latticeTorsionHom, existsUnique_section_of_dense) for the basis condition,
% exists_clopen_weilIs_iff for the pairing, ExtUnique.isMonHom_of_isGrpPullback for the homomorphisms.

\begin{theorem}[Mumford \cite{Mum72}, Faltings--Chai \cite{FC90}]\label{cv-toroidal-reduction-abelian}
Let $W$ be a complete discrete valuation ring with fraction field $K$, let $A$ be an abelian variety over
$K$ with a line bundle $L$ that is trivialized at the origin, symmetric, ample and defines a principal
polarization, and let $G$ be a commutative semi-abelian scheme over $W$. Suppose that there are
degeneration data over $W$ whose free $\mathbf Z$-module $X$ is zero, with a uniformization of $(A,G)$ by
them (Theorem~\ref{cv-toroidal-mumford}). Then there are an abelian scheme $P$ over $W$ with a line
bundle $M$ that is trivialized along the zero section, isomorphic to $[-1]^*M$ locally on $W$, relatively
ample and defines a principal polarization, and a morphism $\chi:A\to P$ over
$\operatorname{Spec}K\to\operatorname{Spec}W$ that identifies $A$ with $P_K$ as group schemes, with $L$
isomorphic to $\chi^*M$.
\end{theorem}
\leanref{Development/AbelianVarietyTheory/Proofs/ModuliProofs/GoodReductionOfThetaPointUnits.lean}{GlobalTheory.ThetaCrit.exists_pol_model}

This is the case $X=0$ of the theory of degenerations. The torus of Theorem~\ref{cv-toroidal-mumford} is
then trivial, and the degeneration data reduce to the abelian scheme $B$ over $W$ with its principal
symmetric line bundle. Mumford's construction from such data returns $B$ \cite{Mum72}, and a
degeneration is determined by its degeneration data \cite[Chapters~II and~III]{FC90}; so $A$, with $L$,
is the generic fibre of $P=B$.

\subsection{Fourier expansions on a degeneration}\label{subsection-cv-fourier-expansions}

\begin{theorem}[Neukirch \cite{Neu99}]\label{cv-toroidal-torus-valuation}
Let $W$ be a discrete valuation ring with uniformizer $\varpi$ and fraction field $K$, and let $F$ be an
algebraic closure of $K$. Then there is a valuation $\nu:F\to\mathbf Q\cup\{\infty\}$ with
$\nu(\varpi^nu)=n$ for every $n\in\mathbf Z$ and every unit $u$ of $W$.
\end{theorem}
\leanref{Development/AbelianVarietyTheory/Proofs/UniformizationProofs/AdicBaseEmbedsValuationAlgebraicClosure.lean}{GlobalTheory.adic_base_embeds_valuation_algebraic_closure}

For complete $W$ the valuation of $K$ extends uniquely to $F$ \cite[Chapter~II, Theorem~4.8]{Neu99}, and
its values on a finite extension $E$ of $K$ are those of the norm from $E$ to $K$ divided by $[E:K]$, so
they lie in $\mathbf Q\cup\{\infty\}$. In general $F$ embeds over $K$ into an algebraic closure of the
completion of $K$, and $\nu$ is the restriction of the valuation there.

\begin{situation}\label{cv-toroidal-torus-situation}
Let $W$ be a complete discrete valuation ring of characteristic zero, with uniformizer $\varpi$, maximal
ideal $\mathfrak m$ and fraction field $K$. Let degeneration data over $W$ be given
(Theorem~\ref{cv-toroidal-mumford}): an abelian scheme $B$ over $W$ with a principal symmetric line bundle
$M$, trivialized at the origin and with its cubical structure, a free $\mathbf Z$-module $X$ of finite
rank, a homomorphism $c:X\to B(W)$ and trivializations $\psi(\chi)$ over $K$ of $c(\chi)^*M^{-1}$. With
$\Lambda(M)$ as in Theorem~\ref{cv-toroidal-mumford}, let $\tau(\chi,\chi')$, for $\chi,\chi'\in X$, be
the trivialization over $K$ of the fibre of $\Lambda(M)^{-1}$ at $(c(\chi),c(\chi'))$ given by
$\psi(\chi+\chi')\psi(\chi)^{-1}\psi(\chi')^{-1}$ and the trivialization of $M$ at the origin. Let $F$ be
an algebraically closed field containing $K$, and let $\nu:F\to\mathbf Q\cup\{\infty\}$ be a valuation
with $\nu(\varpi^nu)=n$ for $n\in\mathbf Z$ and $u\in W^\times$. Let $A$ be an abelian variety over $K$
with a line bundle $L$ that is trivialized at the origin, symmetric, ample and defines a principal
polarization, and let $G$ be a commutative semi-abelian scheme over $W$ with a uniformization of $(A,G)$
by the data (Theorem~\ref{cv-toroidal-mumford}). It includes an isomorphism of group schemes
$\gamma:G_K\to A$ and a line bundle $L_G$ on $G$, trivialized along the zero section, such that
$\gamma^*L$ is isomorphic to the restriction of $L_G$ to $G_K$.
\end{situation}
\leanref{Development/AbelianVarietyTheory/Definitions/Uniformization/AllDV.lean}{GlobalTheory.UT.AllUV}
\leanref{ReciprocityTheory/ReductiveGroups/ArithmeticGeometry/DiscreteValuationRings/ValuationMultiplicationInvariance.lean}{ReciprocityTheory.ChartRestriction.G.DVR.ValOver}
\leanref{ReciprocityTheory/ReductiveGroups/ArithmeticGeometry/ExtensionDescent/ThickTransComp.lean}{ReciprocityTheory.ChartRestriction.G.Out.DegPre.tau}

\begin{lemma}\label{cv-toroidal-torus-quadratic}
In Situation~\ref{cv-toroidal-torus-situation}, for every $\chi\in X$ there is a unique integer $Q(\chi)$
such that $\psi(\chi)=\varpi^{Q(\chi)}ue$ for a unit $u$ of $W$ and a generator $e$ over $W$ of
$c(\chi)^*M^{-1}$. Moreover $Q(0)=0$, and $Q(\chi)\ge1$ for $\chi\ne0$.
\end{lemma}
\leanref{ReciprocityTheory/ReductiveGroups/ArithmeticGeometry/DiscreteValuationRings/WchartLePreimage.lean}{ReciprocityTheory.ChartRestriction.G.DVR.QForm.QVal}
\leanref{ReciprocityTheory/ReductiveGroups/ArithmeticGeometry/DiscreteValuationRings/ExistsRosatiGen.lean}{ReciprocityTheory.ChartRestriction.G.DVR.QForm.exists_QVal}
\leanref{ReciprocityTheory/ReductiveGroups/ArithmeticGeometry/DiscreteValuationRings/ExistsRosatiGen.lean}{ReciprocityTheory.ChartRestriction.G.DVR.QForm.QVal_unique}
\leanref{ReciprocityTheory/ReductiveGroups/ArithmeticGeometry/DiscreteValuationRings/ExistsRosatiGen.lean}{ReciprocityTheory.ChartRestriction.G.DVR.QForm.Q_zero}
\leanref{ReciprocityTheory/ReductiveGroups/ArithmeticGeometry/DiscreteValuationRings/ExistsRosatiGen.lean}{ReciprocityTheory.ChartRestriction.G.DVR.QForm.Q_pos}
\begin{proof}
As $W$ is local, the line bundle $c(\chi)^*M^{-1}$ on $\operatorname{Spec}W$ has a generator $e$, and
$\psi(\chi)=ae$ for some $a\in K^\times$. Write $a=\varpi^nu$ with $u\in W^\times$. Any other generator is
$u'e$ with $u'\in W^\times$, so $n$ does not depend on $e$; this gives existence and uniqueness of $Q(\chi)$.

Let $b(\chi,\chi')=Q(\chi+\chi')-Q(\chi)-Q(\chi')$. The trivialization of $M$ at the origin is a
generator over $W$, so $\tau(\chi,\chi')$ is $\varpi^{b(\chi,\chi')}$ times a generator over $W$. The
bimultiplicativity of $\tau$ says that $\tau(\chi+\chi',\chi'')$, multiplied by the value of the cubical
structure of $M$ at the point $(c(\chi),c(\chi'),c(\chi''))$ of $B^3(W)$, corresponds to
$\tau(\chi,\chi'')\tau(\chi',\chi'')$. That value is a generator over $W$, so
$b(\chi+\chi',\chi'')=b(\chi,\chi'')+b(\chi',\chi'')$. For $\chi=\chi'=\chi''=0$ this gives $b(0,0)=0$,
that is, $Q(0)=0$. For $\chi\ne0$, $\psi(\chi)$ is integral and lies in $\mathfrak m$, so $Q(\chi)\ge1$.
\end{proof}
% Lean: exists_QVal = exists_genHasVal; QVal_unique = genHasVal_unique; Q_zero from bOf_add_left 0 0 0,
% which comes from tau_val (genHasVal_mulGen, genHasVal_invGen, rigK_val), the cubical structure
% (genHasVal_pullGen_of D.cub.t) and D.tau_bimult; Q_pos from genHasVal_ge_of_in and D.psi_pos.

\begin{definition}\label{cv-toroidal-torus-points}
In Situation~\ref{cv-toroidal-torus-situation}, let $Q$ be the function of
Lemma~\ref{cv-toroidal-torus-quadratic}.
\begin{enumerate}
\item A \emph{point of the Raynaud extension} over $F$ is a pair $g=(b,s)$ of a point $b\in B(F)$ and
  trivializations $s(\chi)$, for $\chi\in X$, of the fibres of $\Lambda(M)^{-1}$ at $(c(\chi),b)$, such
  that $s(0)$ is the canonical trivialization of the fibre at $(0,b)$ and $s(\chi+\chi')$ corresponds to
  $s(\chi)s(\chi')$ under the cubical structure of $M$ at $(c(\chi),c(\chi'),b)$. The sum of $(b,s)$ and
  $(b',s')$ is $(b+b',s'')$, where $s''(\chi)$ corresponds to $s(\chi)s'(\chi)$ under the cubical
  structure at $(c(\chi),b,b')$. With this addition the points form an abelian group $\tilde G(F)$, the
  group of points over $F$ of the extension of $B$ by the split torus with character group $X$ that $c$
  defines \cite[Chapter~II]{FC90}, and $\iota(y)=(c(y),\tau(y,\cdot))$ is a homomorphism
  $\iota:X\to\tilde G(F)$.
\item Let $k\ge1$. A \emph{$k$-division point} is a pair $p=(y,g)$ of $y\in X$ and $g\in\tilde G(F)$ with
  $kg=\iota(y)$. The \emph{sum} of $k$-division points $(y,g)$ and $(y',g')$ is $(y+y',g+g')$, and they are
  \emph{equivalent} if $y'=y+kz$ and $g'=g+\iota(z)$ for some $z\in X$.
\item Let $k\ge1$, and for $\chi\in X$ let $M_{k,\chi}$ be the line bundle $M^{k-1}\otimes T_{c(\chi)}^*M$
  on $B$, where $T_{c(\chi)}$ is the translation by $c(\chi)$, tensored with the inverse of the fibre of $M$
  at $c(\chi)$ and with its fibre at the origin. A family of \emph{Fourier coefficients} of weight $k$ over
  $F$ is a family $\sigma=(\sigma_\chi)_{\chi\in X}$ of sections $\sigma_\chi$ of $M_{k,\chi}$ over $B_F$
  such that, for all $y,\chi\in X$, $\sigma_{\chi+ky}$ corresponds to $T_{c(y)}^*\sigma_\chi$ under the
  isomorphism $M_{k,\chi+ky}\cong T_{c(y)}^*M_{k,\chi}$ whose value at the origin is given by
  $\psi(y)^k\tau(y,\chi)$ and the trivialization of $M$ at the origin \cite{Mum72}.
\item For such $\sigma$, $\chi\in X$ and $r\in\mathbf Q$, write $\nu(\sigma_\chi)\ge r$ if
  $\sigma_\chi=\sum_ia_i\sigma_i$ for finitely many sections $\sigma_i$ of $M_{k,\chi}$ over $B$ and
  elements $a_i\in F$ with $\nu(a_i)\ge r$. The family $\sigma$ is \emph{primitive} if
  $\nu(\sigma_\chi)\ge Q(\chi)/k$ for every $\chi$, and there is no rational $r>0$ with
  $\nu(\sigma_\chi)\ge r+Q(\chi)/k$ for every $\chi$.
\end{enumerate}
\end{definition}
\leanref{ReciprocityTheory/ReductiveGroups/ArithmeticGeometry/ExtensionDescent/GtPtBridgeSt.lean}{ReciprocityTheory.ChartRestriction.G.Out.GtPtF}
\leanref{Development/AbelianVarietyTheory/Proofs/ThetaShiftProofs/IsNsmulIffEq.lean}{GlobalTheory.TorusExt.gtGroup}
\leanref{Development/AbelianVarietyTheory/Proofs/ThetaShiftProofs/IsNsmulIffEq.lean}{GlobalTheory.TorusExt.iota}
\leanref{ReciprocityTheory/ReductiveGroups/ArithmeticGeometry/ExtensionDescent/GtPtBridgeSt.lean}{ReciprocityTheory.ChartRestriction.G.Out.QkPt}
\leanref{ReciprocityTheory/ReductiveGroups/ArithmeticGeometry/ExtensionDescent/GtPtBridgeSt.lean}{ReciprocityTheory.ChartRestriction.G.Out.QkPt.Equiv}
\leanref{Development/AbelianVarietyTheory/Proofs/ThetaShiftProofs/RelationsAreFunctions.lean}{GlobalTheory.TorusExt.qkAdd}
\leanref{ReciprocityTheory/ReductiveGroups/ArithmeticGeometry/ExtensionDescent/ThickTransComp.lean}{ReciprocityTheory.ChartRestriction.G.Out.Fam}
\leanref{ReciprocityTheory/ReductiveGroups/ArithmeticGeometry/DiscreteValuationRings/ValuationMultiplicationInvariance.lean}{ReciprocityTheory.ChartRestriction.G.DVR.VGe}
\leanref{ReciprocityTheory/ReductiveGroups/ArithmeticGeometry/DiscreteValuationRings/ValuationMultiplicationInvariance.lean}{ReciprocityTheory.ChartRestriction.G.DVR.NuGe}
\leanref{Development/AbelianVarietyTheory/Definitions/Uniformization/AllDV.lean}{GlobalTheory.UT.Prim}

\begin{lemma}\label{cv-toroidal-torus-extension}
In Situation~\ref{cv-toroidal-torus-situation}, let $\tilde G(F)$ be the group of
Definition~\ref{cv-toroidal-torus-points}(1).
\begin{enumerate}
\item Every point of $B(F)$ is the first component of a point of $\tilde G(F)$.
\item The fibres of $\Lambda(M)^{-1}$ at the points $(c(\chi),0)$ are canonically trivial. Through these
  trivializations, the points of $\tilde G(F)$ with first component $0$ are the pairs
  $(0,s)$ with $s:X\to F^\times$ a homomorphism, and the sum of $(0,s)$ and $(0,s')$ is $(0,ss')$.
\item For every $n\ge1$ and every $g\in\tilde G(F)$ there is $g'\in\tilde G(F)$ with $ng'=g$.
\end{enumerate}
\end{lemma}
\leanref{Development/AbelianVarietyTheory/Proofs/UniformizationProofs/GtOntoDivisibleKernel.lean}{GlobalTheory.UT.gt_onto_divisible_kernel}
\begin{proof}
(1) Let $b\in B(F)$. Choose generators of $M$ near the points $c(\chi)+b$, $c(\chi)$, $b$ and $0$ of
$B_F$. In them, trivializations $s(\chi)$ become elements $f(\chi)\in F^\times$, and the value of the
cubical structure at $(c(\chi),c(\chi'),b)$ becomes an element $\kappa(\chi,\chi')\in F^\times$. The
conditions on $s$ become $f(0)=1$ and $f(\chi+\chi')=\kappa(\chi,\chi')f(\chi)f(\chi')$. The axioms of a
cubical structure make $\kappa$ a symmetric $2$-cocycle on $X$, normalized at $0$. Since $X$ is free, a
symmetric $2$-cocycle on $X$ with values in $F^\times$ is a coboundary, which gives $f$.

(2) The value of the cubical structure at $(c(\chi),c(\chi'),0)$ is $1$. So for $b=0$ the condition on
$s$ says that $s(\chi+\chi')=s(\chi)s(\chi')$, and the sum of $(0,s)$ and $(0,s')$ is $(0,ss')$.

(3) Let $g=(b,s)$. Multiplication by $n$ on $B$ is finite and surjective \cite{Mum70}, and $F$ is
algebraically closed, so $b=nb'$ for some $b'\in B(F)$. By (1) there is $g_1\in\tilde G(F)$ with first
component $b'$. The first component of $g-ng_1$ is $0$, so by (2) $g-ng_1=(0,t)$ for a homomorphism
$t:X\to F^\times$. As $X$ is free and $F^\times$ is divisible, $t=s_0^n$ for a homomorphism
$s_0:X\to F^\times$. Then $g'=g_1+(0,s_0)$ satisfies $ng'=ng_1+(0,t)=g$.
\end{proof}
% Lean: gt_onto_divisible_kernel: (1) TorusExt.exists_lift (exists_good_chart, IsSymmCocycle with IsCub.face,
% s01 and coc, exists_coboundary, exists_gt); (2) torusPt, isTorusPtF_iff_eq, exists_isTorusPtF_of_b,
% IsTorusPtF.add, with IsCub.face3; (3) TorusExt.exists_nsmul_eq (exists_pow_eq with the law MulNSt,
% exists_nsmul_eq_of_pow_eq with exists_lift, exists_torusPt_eq and exists_pow_eq_hom).

\begin{definition}\label{cv-toroidal-torus-expansion}
In Situation~\ref{cv-toroidal-torus-situation}, let $k\ge1$.
\begin{enumerate}
\item Let $\rho$ be a trivialization of $L$ at the origin. A \emph{normalization} for $\rho$ is an
  isomorphism $w_0$ from $\gamma^*L$ to the restriction of $L_G$ to $G_K$ that carries $\rho$ to the
  trivialization of $L_G$ at the origin.
\item Let $w_0$ be an isomorphism from $\gamma^*L$ to the restriction of $L_G$ to $G_K$. A family $\sigma$
  of Fourier coefficients of weight $k$ over $F$ (Definition~\ref{cv-toroidal-torus-points}(3)) is the
  \emph{Fourier expansion} of a section $t$ of $L^k$ over $A_F$ with respect to $w_0$ if there are sections
  $t_1,\dots,t_r$ of $L_G^k$ over $G$ and $a_1,\dots,a_r\in F$ such that $w_0$ and $\gamma$ carry $t$ to
  $\sum_ia_it_i$ and $\sigma=\sum_ia_i\sigma^{(i)}$, where $\sigma^{(i)}$ is the family of Fourier
  coefficients of $t_i$ given by the uniformization (Theorem~\ref{cv-toroidal-mumford}).
\end{enumerate}
\end{definition}
\leanref{Development/AbelianVarietyTheory/Definitions/Uniformization/TXFrame.lean}{GlobalTheory.UT.NormW}
\leanref{ReciprocityTheory/ReductiveGroups/ArithmeticGeometry/ExtensionDescent/DegExpF.lean}{ReciprocityTheory.ChartRestriction.G.Out.UnifPre.ExpandsF}

\begin{lemma}\label{cv-toroidal-torus-normalization}
In Situation~\ref{cv-toroidal-torus-situation}, let $\rho$ be a trivialization of $L$ at the origin. Then
there is a normalization for $\rho$ (Definition~\ref{cv-toroidal-torus-expansion}(1)).
\end{lemma}
\leanref{Development/AbelianVarietyTheory/Proofs/UniformizationProofs/ValueAtOrigin.lean}{GlobalTheory.UT.EvalZero.value_at_origin}
\leanref{Development/AbelianVarietyTheory/Proofs/UniformizationProofs/ValueAtOrigin.lean}{GlobalTheory.UT.EvalZero.exists_normW}
\begin{proof}
Let $w$ be an isomorphism from $\gamma^*L$ to the restriction of $L_G$ to $G_K$. At the origin, the image
of $\rho$ under $w$ and the trivialization of $L_G$ are two trivializations of the same line over $K$, so
the second is $a$ times the first for some $a\in K^\times$. Then $aw$ is a normalization for $\rho$.
\end{proof}

\begin{theorem}[Mumford \cite{Mum70}]\label{cv-toroidal-torus-sections}
In Situation~\ref{cv-toroidal-torus-situation}, let $g$ be the relative dimension of $B$, let $k\ge1$ and
$\chi\in X$, and let $M_{k,\chi}$ be the line bundle on $B$ of
Definition~\ref{cv-toroidal-torus-points}(3). Then there are sections $\sigma_1,\dots,\sigma_{k^g}$ of
$M_{k,\chi}$ over $B$ such that, for every ring homomorphism $W\to R'$, their pullbacks form a basis of the
$R'$-module of sections over $B_{R'}$ of the pullback of $M_{k,\chi}$.
\end{theorem}
\leanref{ReciprocityTheory/ReductiveGroups/ArithmeticGeometry/DiscreteValuationRings/TwistBCPf.lean}{ReciprocityTheory.ChartRestriction.G.DVR.TwistBaseChangeStatement}
\leanref{Development/AbelianVarietyTheory/Proofs/UniformizationProofs/TwistedPrincipalPolarisationSectionsBaseChange.lean}{GlobalTheory.twistBaseChange}

On every geometric fibre of $B$ the line bundle $T_{c(\chi)}^*M$ is algebraically equivalent to $M$, so
$M_{k,\chi}$ is algebraically equivalent to $M^k$. As $M$ defines a principal polarization, $M_{k,\chi}$ is
ample, its higher cohomology vanishes and its space of sections has dimension $k^g$ \cite[\S16]{Mum70}. So
the direct image of $M_{k,\chi}$ is locally free of rank $k^g$ and commutes with base change
\cite[\S5]{Mum70}, and over the local ring $W$ it is free.

\begin{lemma}\label{cv-toroidal-torus-coordinates}
In Situation~\ref{cv-toroidal-torus-situation}, let $g$ be the relative dimension of $B$, and let $k\ge1$
and $\chi\in X$. Then there are sections $\sigma_1,\dots,\sigma_{k^g}$ of $M_{k,\chi}$ over $B$ with the
following properties. Every section $s$ of $M_{k,\chi}$ over $B_F$ is $\sum_ia_i\sigma_i$ for unique
$a_1,\dots,a_{k^g}\in F$. For every rational $r$, $s$ is a finite sum $\sum_ja'_j\sigma'_j$ of sections
$\sigma'_j$ of $M_{k,\chi}$ over $B$ with $a'_j\in F$ and $\nu(a'_j)\ge r$ if and only if $\nu(a_i)\ge r$
for every $i$.
\end{lemma}
\leanref{Development/AbelianVarietyTheory/Proofs/RingsOfIntegers/ExistsBasisOfEq.lean}{GlobalTheory.IntegralCoord.main}
\begin{proof}
Let $\sigma_1,\dots,\sigma_{k^g}$ be the sections of Theorem~\ref{cv-toroidal-torus-sections}. For the
homomorphism $W\to F$ their pullbacks form a basis of the sections of $M_{k,\chi}$ over $B_F$, which gives
the $a_i$. If $\nu(a_i)\ge r$ for every $i$, then $s=\sum_ia_i\sigma_i$ is a sum of the required form.
Conversely, let $s=\sum_ja'_j\sigma'_j$ with sections $\sigma'_j$ over $B$ and $\nu(a'_j)\ge r$. For the
identity of $W$, Theorem~\ref{cv-toroidal-torus-sections} gives $d_{ji}\in W$ with
$\sigma'_j=\sum_id_{ji}\sigma_i$. Then $s=\sum_i\bigl(\sum_ja'_jd_{ji}\bigr)\sigma_i$, so
$a_i=\sum_ja'_jd_{ji}$ by uniqueness. As $\nu(d)\ge0$ for every $d\in W$, $\nu(a_i)\ge r$.
\end{proof}

\begin{lemma}\label{cv-toroidal-torus-translation}
In Situation~\ref{cv-toroidal-torus-situation}, let $Q$ be the function of
Lemma~\ref{cv-toroidal-torus-quadratic}, and let $b(\chi,\chi')=Q(\chi+\chi')-Q(\chi)-Q(\chi')$ for
$\chi,\chi'\in X$. Let $E$ be a field containing $K$ and let $k\ge1$. For $\mu\in X$ let $T_\mu$ be the
translation by $c(\mu)$ on $B$.
\begin{enumerate}
\item For $\mu,\chi\in X$ there is a unique isomorphism $u_{\mu,\chi}:T_\mu^*M_{k,\chi}\to M_{k,\chi+k\mu}$
  over $B_E$ whose value at the origin is given by $\psi(\mu)^k\tau(\mu,\chi)$ and the trivialization of
  $M$ at the origin.
\item For $\mu,\mu',\chi\in X$, $u_{\mu+\mu',\chi}=u_{\mu',\chi+k\mu}\circ T_{\mu'}^*u_{\mu,\chi}$, and
  $u_{0,\chi}$ is the identity.
\item For $\mu,\chi\in X$ there are an isomorphism $u'_{\mu,\chi}:T_\mu^*M_{k,\chi}\to M_{k,\chi+k\mu}$ over
  $B$ and an element $a$ of $K$, equal to $\varpi^{kQ(\mu)+b(\mu,\chi)}$ times a unit of $W$, such that
  $u_{\mu,\chi}$ is $a$ times the pullback of $u'_{\mu,\chi}$ to $B_E$.
\end{enumerate}
\end{lemma}
\leanref{ReciprocityTheory/ReductiveGroups/ArithmeticGeometry/DiscreteValuationRings/PhiGenPf.lean}{ReciprocityTheory.ChartRestriction.G.DVR.gen_fe}
\leanref{ReciprocityTheory/ReductiveGroups/ArithmeticGeometry/DiscreteValuationRings/PhiGenPf.lean}{ReciprocityTheory.ChartRestriction.G.DVR.gen_u}
\leanref{ReciprocityTheory/ReductiveGroups/ArithmeticGeometry/BirationalGeometry/ToroidalCompactification/CubeVal.lean}{ReciprocityTheory.ChartRestriction.G.Geo.Tor.exists_cubN}
\leanref{ReciprocityTheory/ReductiveGroups/ArithmeticGeometry/DiscreteValuationRings/ValPowGen.lean}{ReciprocityTheory.ChartRestriction.G.DVR.UUA.fe_cocycle}
\leanref{ReciprocityTheory/ReductiveGroups/ArithmeticGeometry/DiscreteValuationRings/ExistsValGenHasVal.lean}{ReciprocityTheory.ChartRestriction.G.DVR.SectionScalarValuation.famMult_data}
\leanref{ReciprocityTheory/ReductiveGroups/ArithmeticGeometry/DiscreteValuationRings/ProdNeg.lean}{ReciprocityTheory.ChartRestriction.G.DVR.SectionScalarValuation.exists_eB}
\begin{proof}
(1) For a point $x$ of $B_E$, the product of the values of the cubical structure of $M$ at the points
$(c(\chi)+ic(\mu),c(\mu),x)$, for $0\le i<k$, divided by its value at $(c(\chi),x,c(\mu))$, gives an
isomorphism from the line bundle of homomorphisms from $T_\mu^*M_{k,\chi}$ to $M_{k,\chi+k\mu}$ onto the
constant line bundle given by its fibre at the origin. At the origin it is the identity, since the cubical
structure is $1$ when one of its arguments is the origin. Composing its inverse with the trivialization of
the fibre at the origin given by $\psi(\mu)^k\tau(\mu,\chi)$ gives $u_{\mu,\chi}$. If $u$ is a second such
isomorphism, then $u/u_{\mu,\chi}$ is an invertible function on $B_E$, hence a constant, because the global
functions on $B_E$ are those of $E$; it is $1$ at the origin, so $u=u_{\mu,\chi}$.

(2) Both sides are isomorphisms from $T_{\mu+\mu'}^*M_{k,\chi}$ to $M_{k,\chi+k(\mu+\mu')}$. Their values at
a point of $B_E$ are given by (1) in terms of the cubical structure of $M$ at the points $c(\chi)$,
$c(\mu)$, $c(\mu')$ and that point, and of $\psi$ and $\tau$. They agree by the cocycle and symmetry
identities of the cubical structure, the bimultiplicativity and symmetry of $\tau$, and
$\psi(\mu+\mu')=\psi(\mu)\psi(\mu')\tau(\mu,\mu')$, which is the definition of $\tau$. For $\mu=\mu'=0$ the
identity says that $u_{0,\chi}=u_{0,\chi}\circ u_{0,\chi}$, so $u_{0,\chi}$ is the identity.

(3) By Lemma~\ref{cv-toroidal-torus-quadratic} and its proof, $\psi(\mu)$ is $\varpi^{Q(\mu)}$ times a
generator over $W$, and $\tau(\mu,\chi)$ is $\varpi^{b(\mu,\chi)}$ times a generator over $W$; the
trivialization of $M$ at the origin is a generator over $W$. So the trivialization of the fibre at the
origin given by $\psi(\mu)^k\tau(\mu,\chi)$ is $ae$, with $a$ as in the statement and $e$ a generator over
$W$. The construction of (1), carried out over $B$ with $e$ in place of that trivialization, gives an
isomorphism $u'_{\mu,\chi}$ over $B$ whose value at the origin is $e$, as $c(\chi)$ and $c(\mu)$ lie in
$B(W)$. Its pullback times $a$ has the value at the origin required in (1), so it is $u_{\mu,\chi}$.
\end{proof}

\begin{lemma}\label{cv-toroidal-torus-classes}
In Situation~\ref{cv-toroidal-torus-situation}, let $E$ be a field containing $K$, let $k\ge1$ and let
$\alpha\in X$. Families of Fourier coefficients of weight $k$ over $E$ are defined as in
Definition~\ref{cv-toroidal-torus-points}(3), with $E$ in place of $F$.
\begin{enumerate}
\item Let $\sigma$ and $\sigma'$ be families of Fourier coefficients of weight $k$ over $E$ with
  $\sigma_\alpha=\sigma'_\alpha$. Then $\sigma_{\alpha+k\mu}=\sigma'_{\alpha+k\mu}$ for every $\mu\in X$.
\item Let $s$ be a section of $M_{k,\alpha}$ over $B_E$. Then there is a family $\sigma$ of Fourier
  coefficients of weight $k$ over $E$ with $\sigma_\alpha=s$ and $\sigma_\beta=0$ for every
  $\beta\notin\alpha+kX$.
\end{enumerate}
\end{lemma}
\leanref{Development/AbelianVarietyTheory/Proofs/RingsOfIntegers/FamEqOnClassOfEq.lean}{GlobalTheory.fam_eq_on_class_of_eq}
\leanref{ReciprocityTheory/ReductiveGroups/ArithmeticGeometry/DiscreteValuationRings/ExistsFamRep.lean}{ReciprocityTheory.ChartRestriction.G.DVR.UUA.exists_fam_rep}
\begin{proof}
With $T_\mu$ and $u_{\mu,\chi}$ as in Lemma~\ref{cv-toroidal-torus-translation}, a family $\sigma$ of
Fourier coefficients of weight $k$ over $E$ satisfies $\sigma_{\chi+k\mu}=u_{\mu,\chi}(T_\mu^*\sigma_\chi)$
for all $\mu,\chi\in X$, by Definition~\ref{cv-toroidal-torus-points}(3) and
Lemma~\ref{cv-toroidal-torus-translation}(1). For $\chi=\alpha$ this gives (1).

(2) As $X$ is free, $\mu\mapsto\alpha+k\mu$ is injective. Put $\sigma_{\alpha+k\mu}=u_{\mu,\alpha}(T_\mu^*s)$
for $\mu\in X$, and $\sigma_\beta=0$ for $\beta\notin\alpha+kX$. Then $\sigma_\alpha=s$, as $u_{0,\alpha}$ is
the identity by Lemma~\ref{cv-toroidal-torus-translation}(2). Let $\mu',\beta\in X$. If
$\beta=\alpha+k\mu$, then $T_{\mu+\mu'}=T_\mu\circ T_{\mu'}$, and Lemma~\ref{cv-toroidal-torus-translation}(2)
gives
\[
\sigma_{\beta+k\mu'}=u_{\mu+\mu',\alpha}(T_{\mu+\mu'}^*s)
=u_{\mu',\beta}\bigl(T_{\mu'}^*(u_{\mu,\alpha}(T_\mu^*s))\bigr)=u_{\mu',\beta}(T_{\mu'}^*\sigma_\beta).
\]
If $\beta\notin\alpha+kX$, then $\beta+k\mu'\notin\alpha+kX$, and both sides vanish. So $\sigma$ satisfies
the functional equation of Definition~\ref{cv-toroidal-torus-points}(3).
\end{proof}

\begin{lemma}\label{cv-toroidal-torus-shift}
In Situation~\ref{cv-toroidal-torus-situation}, let $Q$ be the function of
Lemma~\ref{cv-toroidal-torus-quadratic}, let $b(\chi,\chi')=Q(\chi+\chi')-Q(\chi)-Q(\chi')$ for
$\chi,\chi'\in X$, and let $k\ge1$. Let $\sigma$ be a family of Fourier coefficients of weight $k$ over $F$.
\begin{enumerate}
\item For $\alpha,\mu\in X$, $Q(\alpha+k\mu)=Q(\alpha)+k\bigl(kQ(\mu)+b(\mu,\alpha)\bigr)$.
\item Let $\alpha,\mu\in X$, and let $r$ be rational with $\nu(\sigma_\alpha)\ge r$
  (Definition~\ref{cv-toroidal-torus-points}(4)). Then $\nu(\sigma_{\alpha+k\mu})\ge r+kQ(\mu)+b(\mu,\alpha)$.
\item Let $S$ be a finite subset of $X$ that meets every class of $X$ modulo $kX$, and let $r$ be rational.
  Then $\nu(\sigma_\chi)\ge r+Q(\chi)/k$ for every $\chi\in X$ if and only if this holds for every
  $\chi\in S$.
\end{enumerate}
\end{lemma}
\leanref{ReciprocityTheory/ReductiveGroups/ArithmeticGeometry/DiscreteValuationRings/BOfNsmulLeft.lean}{ReciprocityTheory.ChartRestriction.G.DVR.q_shift}
\leanref{ReciprocityTheory/ReductiveGroups/ArithmeticGeometry/DiscreteValuationRings/GsScalJFK.lean}{ReciprocityTheory.ChartRestriction.G.DVR.SectionScalarValuation.fe_move}
\leanref{ReciprocityTheory/ReductiveGroups/ArithmeticGeometry/DiscreteValuationRings/ValuationMultiplicationInvariance.lean}{ReciprocityTheory.ChartRestriction.G.DVR.ValuationMinRepresentatives}
\leanref{ReciprocityTheory/ReductiveGroups/ArithmeticGeometry/DiscreteValuationRings/BOfNsmulLeft.lean}{ReciprocityTheory.ChartRestriction.G.DVR.n1a}
\begin{proof}
(1) The proof of Lemma~\ref{cv-toroidal-torus-quadratic} shows that $b$ is additive in its first argument,
and $b$ is symmetric. As $\psi$ is even, $\psi(-\chi)$ corresponds to $\psi(\chi)$ under the symmetry of
$M$, which is a generator over $W$, so $Q(-\chi)=Q(\chi)$. Then $b(0,\chi)=0$ and
$b(-\chi,\chi)=Q(0)-Q(-\chi)-Q(\chi)=-2Q(\chi)$, so additivity gives $b(\chi,\chi)=2Q(\chi)$. By induction
on $n\ge0$, $b(n\mu,\chi)=nb(\mu,\chi)$ and
$Q((n+1)\mu)=Q(n\mu)+Q(\mu)+b(n\mu,\mu)=(n^2+1+2n)Q(\mu)$, so $Q(n\mu)=n^2Q(\mu)$. Hence
$Q(\alpha+k\mu)=Q(\alpha)+Q(k\mu)+b(k\mu,\alpha)=Q(\alpha)+k^2Q(\mu)+kb(\mu,\alpha)$.

(2) Let $T_\mu$, $u_{\mu,\alpha}$, $u'_{\mu,\alpha}$ and $a$ be as in
Lemma~\ref{cv-toroidal-torus-translation}. By Definition~\ref{cv-toroidal-torus-points}(3) and
Lemma~\ref{cv-toroidal-torus-translation}(1) and~(3),
$\sigma_{\alpha+k\mu}=u_{\mu,\alpha}(T_\mu^*\sigma_\alpha)=a\,u'_{\mu,\alpha}(T_\mu^*\sigma_\alpha)$. Write
$\sigma_\alpha=\sum_ia_i\sigma_i$ with sections $\sigma_i$ of $M_{k,\alpha}$ over $B$ and $\nu(a_i)\ge r$. As
$c(\mu)\in B(W)$ and $u'_{\mu,\alpha}$ is defined over $B$, each $u'_{\mu,\alpha}(T_\mu^*\sigma_i)$ is a
section of $M_{k,\alpha+k\mu}$ over $B$. So
$\sigma_{\alpha+k\mu}=\sum_i(aa_i)\,u'_{\mu,\alpha}(T_\mu^*\sigma_i)$, and
$\nu(aa_i)\ge r+kQ(\mu)+b(\mu,\alpha)$, since $\nu(\varpi^nv)=n$ for every unit $v$ of $W$.

(3) The condition for every $\chi\in X$ implies it for $\chi\in S$. Conversely, let $\chi\in X$, and write
$\chi=\alpha+k\mu$ with $\alpha\in S$. By (2), applied to $r+Q(\alpha)/k$, and by (1),
$\nu(\sigma_\chi)\ge r+Q(\alpha)/k+kQ(\mu)+b(\mu,\alpha)=r+Q(\chi)/k$.
\end{proof}

\begin{lemma}\label{cv-toroidal-torus-rational}
In Situation~\ref{cv-toroidal-torus-situation}, let $q$ be rational. Then there is $t\in F^\times$ with
$\nu(t)=q$, and every $a\in F$ with $\nu(a)\ge q$ is $ct$ for some $c\in F$ with $\nu(c)\ge0$.
\end{lemma}
\leanref{ReciprocityTheory/ReductiveGroups/ArithmeticGeometry/DiscreteValuationRings/BOfNsmulLeft.lean}{ReciprocityTheory.ChartRestriction.G.DVR.exists_val_eq}
\leanref{Development/AbelianVarietyTheory/Proofs/RingsOfIntegers/ExistsLeastOfDiscrete.lean}{GlobalTheory.ValuedPrincipal.main}
\begin{proof}
Write $q=n/d$ with $n\in\mathbf Z$ and $d\ge1$. As $F$ is algebraically closed, there is $t\in F$ with
$t^d=\varpi^n$. Then $t\ne0$ and $d\nu(t)=\nu(\varpi^n)=n$, so $\nu(t)=q$. If $\nu(a)\ge q$, then $c=at^{-1}$
satisfies $a=ct$ and $\nu(c)=\nu(a)-q\ge0$.
\end{proof}

\begin{lemma}\label{cv-toroidal-torus-ball}
In Situation~\ref{cv-toroidal-torus-situation}, let $Q$ be the function of
Lemma~\ref{cv-toroidal-torus-quadratic}, let $r$ be the rank of $X$ and $g$ the relative dimension of $B$,
and let $k\ge1$. The families of Fourier coefficients of weight $k$ over $F$ form an $F$-vector space under
componentwise addition and multiplication by scalars. There are such families
$\eta_1,\dots,\eta_{k^{r+g}}$, linearly independent over $F$, with $\nu(\eta_{i,\chi})\ge Q(\chi)/k$ for
every $i$ and every $\chi\in X$, such that a family $\sigma$ of Fourier coefficients of weight $k$ over $F$
satisfies $\nu(\sigma_\chi)\ge Q(\chi)/k$ for every $\chi\in X$ if and only if $\sigma=\sum_ic_i\eta_i$ with
$c_i\in F$ and $\nu(c_i)\ge0$.
\end{lemma}
\leanref{ReciprocityTheory/ReductiveGroups/ArithmeticGeometry/DiscreteValuationRings/ValuationMultiplicationInvariance.lean}{ReciprocityTheory.ChartRestriction.G.DVR.UUStatement}
\leanref{Development/AbelianVarietyTheory/Proofs/RingsOfIntegers/ExistsBasisOfCoord.lean}{GlobalTheory.UnitBallFree.main}
\leanref{Development/AbelianVarietyTheory/Proofs/RingsOfIntegers/ExistsBasisOfCoord.lean}{GlobalTheory.UnitBallFree.exists_basis_of_coord}
\leanref{Development/AbelianVarietyTheory/Proofs/RingsOfIntegers/ExistsRepresentativesModNsmul.lean}{GlobalTheory.exists_representatives_mod_nsmul}
\begin{proof}
The isomorphisms $u_{\mu,\chi}$ of Lemma~\ref{cv-toroidal-torus-translation} are linear, so sums and
multiples of families satisfy the functional equation of Definition~\ref{cv-toroidal-torus-points}(3).

Choose a basis of $X$. The elements $e_1,\dots,e_{k^r}$ of $X$ whose coordinates lie in
$\{0,\dots,k-1\}$ meet every class of $X$ modulo $kX$ once. For each $i$,
Lemma~\ref{cv-toroidal-torus-coordinates} gives sections $\sigma_{i,1},\dots,\sigma_{i,k^g}$ of
$M_{k,e_i}$ over $B$. For a family $\sigma$, let $a_{i,j}(\sigma)\in F$ be the coefficients of
$\sigma_{e_i}$ in terms of them. The map $\phi:\sigma\mapsto(a_{i,j}(\sigma))_{i,j}$ to $F^{k^r\times k^g}$ is
$F$-linear, and it is injective by Lemma~\ref{cv-toroidal-torus-classes}(1). By
Lemma~\ref{cv-toroidal-torus-classes}(2) there are families $\sigma^{i,j}$ with $\sigma^{i,j}_{e_i}=\sigma_{i,j}$
and $\sigma^{i,j}_\beta=0$ for $\beta\notin e_i+kX$; as the $e_{i'}$ lie in distinct classes,
$\phi(\sigma^{i,j})$ is the $(i,j)$-th unit vector. By Lemma~\ref{cv-toroidal-torus-shift}(3), with
$S=\{e_1,\dots,e_{k^r}\}$, and Lemma~\ref{cv-toroidal-torus-coordinates}, a family $\sigma$ satisfies
$\nu(\sigma_\chi)\ge Q(\chi)/k$ for every $\chi$ if and only if $\nu(a_{i,j}(\sigma))\ge Q(e_i)/k$ for all
$i$ and $j$.

Lemma~\ref{cv-toroidal-torus-rational} gives $t_{i,j}\in F^\times$ with $\nu(t_{i,j})=Q(e_i)/k$. Put
$\eta_{i,j}=t_{i,j}\sigma^{i,j}$; there are $k^r\cdot k^g=k^{r+g}$ of them. Then $\phi(\eta_{i,j})$ is
$t_{i,j}$ times the $(i,j)$-th unit vector, so the $\eta_{i,j}$ are linearly independent and satisfy the
bound. If $\sigma=\sum c_{i,j}\eta_{i,j}$ with $\nu(c_{i,j})\ge0$, then $a_{i,j}(\sigma)=c_{i,j}t_{i,j}$ has
valuation at least $Q(e_i)/k$. Conversely, if $\sigma$ satisfies the bound, then
Lemma~\ref{cv-toroidal-torus-rational} gives $c_{i,j}$ with $\nu(c_{i,j})\ge0$ and
$a_{i,j}(\sigma)=c_{i,j}t_{i,j}$. Then $\phi(\sigma)=\phi\bigl(\sum c_{i,j}\eta_{i,j}\bigr)$, and
$\sigma=\sum c_{i,j}\eta_{i,j}$ as $\phi$ is injective.
\end{proof}

\begin{lemma}\label{cv-toroidal-torus-base-change}
In Situation~\ref{cv-toroidal-torus-situation}, the following hold.
\begin{enumerate}
\item Let $Y$ be a noetherian scheme over $K$ and $N$ a line bundle on $Y$. Then the map
  $F\otimes_K\Gamma(Y,N)\to\Gamma(Y_F,N)$ given by pullback is an isomorphism.
\item Let $k\ge1$. Then the families of Fourier coefficients of weight $k$ over $K$
  (Lemma~\ref{cv-toroidal-torus-classes}) form a $K$-vector space, and pullback gives an isomorphism from its
  tensor product with $F$ over $K$ onto the space of families of Fourier coefficients of weight $k$ over
  $F$.
\end{enumerate}
\end{lemma}
\leanref{Development/AbelianVarietyTheory/Proofs/ThetaShiftProofs/IsExtensionPullOver.lean}{GlobalTheory.FieldBaseChange.isExtension_pullOver}
\leanref{Development/AbelianVarietyTheory/Proofs/ThetaShiftProofs/ExistsFinsetRepresentatives.lean}{GlobalTheory.FieldBaseChange.famPull}
\leanref{Development/AbelianVarietyTheory/Proofs/ThetaShiftProofs/ExistsFinsetRepresentatives.lean}{GlobalTheory.FieldBaseChange.isExtension_famPull}
\begin{proof}
(1) This is flat base change for the quasi-compact and quasi-separated scheme $Y$ \cite[Tag 02KH]{Sta}.

(2) The families over $K$ form a vector space as in the proof of Lemma~\ref{cv-toroidal-torus-ball}. The
pullback to $B_F$ of the isomorphism $u_{\mu,\chi}$ of Lemma~\ref{cv-toroidal-torus-translation}(1) over
$B_K$ has the value at the origin required there over $B_F$, so it is $u_{\mu,\chi}$ over $B_F$. Hence the
pullback of a family over $K$ is a family over $F$. Let $(\beta_l)$ be a basis of $F$ over $K$.

Injectivity. Let $s_l$ be families over $K$, almost all zero, with $\sum_l\beta_ls_l=0$ over $F$. For each
$\chi$, $\sum_l\beta_ls_{l,\chi}=0$ in $\Gamma(B_F,M_{k,\chi})$, so $s_{l,\chi}=0$ for every $l$ by (1),
applied to $B_K$.

Surjectivity. Let $\sigma$ be a family over $F$. For each $\chi$, (1) writes
$\sigma_\chi=\sum_l\beta_ls_{l,\chi}$ with unique sections $s_{l,\chi}$ of $M_{k,\chi}$ over $B_K$, almost all
zero. As the $u_{\mu,\chi}$ over $F$ are pullbacks of those over $K$, the uniqueness of these expressions
shows that each $s_l=(s_{l,\chi})_\chi$ is a family over $K$. Let $T$ be a finite subset of $X$ meeting
every class modulo $kX$. Only finitely many $l$ have $s_{l,\chi}\ne0$ for some $\chi\in T$; for every other
$l$, $s_l$ vanishes on $T$, so $s_l=0$ by Lemma~\ref{cv-toroidal-torus-classes}(1). So $\sigma=\sum_l\beta_ls_l$
is a finite sum.
\end{proof}

\begin{lemma}\label{cv-toroidal-torus-bijection}
In Situation~\ref{cv-toroidal-torus-situation}, let $k\ge1$ and let $w_0$ be an isomorphism from
$\gamma^*L$ to the restriction of $L_G$ to $G_K$. Then every section of $L^k$ over $A_F$ has exactly one
Fourier expansion with respect to $w_0$ (Definition~\ref{cv-toroidal-torus-expansion}(2)), every family of
Fourier coefficients of weight $k$ over $F$ is the Fourier expansion of exactly one section of $L^k$ over
$A_F$, and the map that takes a section to its Fourier expansion is $F$-linear.
\end{lemma}
\leanref{Development/AbelianVarietyTheory/Proofs/ThetaShiftProofs/ExpansionBijectiveOverFields.lean}{GlobalTheory.exists_linearEquiv_expandsF}
\leanref{Development/AbelianVarietyTheory/Proofs/ThetaShiftProofs/ExpansionBijectiveOverFields.lean}{GlobalTheory.FieldBaseChange.expandsF_bijective}
\leanref{Development/AbelianVarietyTheory/Proofs/ThetaShiftProofs/IsExtensionOfIsPushout.lean}{GlobalTheory.FieldBaseChange.graph_bijective}
\begin{proof}
Through $w_0$ and $\gamma$, the sections of $L^k$ over $A_F$ correspond $F$-linearly and bijectively to the
sections of $L_G^k$ over $G_F$. Under this correspondence, $\sigma$ is the Fourier expansion of a section of
$L^k$ if and only if the corresponding section $v$ is $\sum_ia_it_i$ for sections $t_i$ of $L_G^k$ over
$G$ and $a_i\in F$ with $\sigma=\sum_ia_i\sigma^{(i)}$, where $\sigma^{(i)}$ is the family of Fourier
coefficients of $t_i$. Call such a $\sigma$ an expansion of $v$.

Over $K$. By the uniformization (Theorem~\ref{cv-toroidal-mumford}), for every section $v$ of $L_G^k$ over
$G_K$ there is exactly one family $e(v)$ of weight $k$ over $K$ such that, for some nonzero $r\in W$, $rv$
extends to a section $t$ of $L_G^k$ over $G$ whose family of Fourier coefficients is $re(v)$; and every
family over $K$ is $e(v)$ for exactly one $v$. The family of Fourier coefficients of a $W$-linear
combination of sections over $G$ is the same combination of their families.

Existence. Let $v$ be a section of $L_G^k$ over $G_F$. As $G_K$ is isomorphic to $A$, it is noetherian, and
Lemma~\ref{cv-toroidal-torus-base-change}(1) writes $v=\sum_ic_iv_i$ with $c_i\in F$ and sections $v_i$ over
$G_K$. For each $i$ choose $r_i$ and $t_i$ for $v_i$ as above. Then $v=\sum_i(c_i/r_i)t_i$, and
$\sum_i(c_i/r_i)\,r_ie(v_i)=\sum_ic_ie(v_i)$, so $\sum_ic_ie(v_i)$ is an expansion of $v$. In the same way,
Lemma~\ref{cv-toroidal-torus-base-change}(2) writes a family $\sigma$ over $F$ as $\sum_ic_i\sigma_i$ with
families $\sigma_i$ over $K$, and $\sigma$ is an expansion of $\sum_ic_ie^{-1}(\sigma_i)$.

Uniqueness. Let $\sigma$ be an expansion of $v$, with $v=\sum_ia_it_i$ and $\sigma=\sum_ia_i\sigma^{(i)}$. Let
$\ell:F\to K$ be $K$-linear, and let $\Lambda_\ell$ be the map $\ell\otimes1$ on sections over $G_F$ and on
families over $F$, through the isomorphisms of Lemma~\ref{cv-toroidal-torus-base-change}. Then
$\Lambda_\ell(v)=\sum_i\ell(a_i)t_i$ over $G_K$ and $\Lambda_\ell(\sigma)=\sum_i\ell(a_i)\sigma^{(i)}$. Choose a
nonzero $r\in W$ with $r\ell(a_i)\in W$ for every $i$. Then $r\Lambda_\ell(v)$ extends to the section
$\sum_ir\ell(a_i)t_i$ over $G$, whose family of Fourier coefficients is $r\Lambda_\ell(\sigma)$, so
$\Lambda_\ell(\sigma)=e(\Lambda_\ell(v))$. For a basis $(\beta_l)$ of $F$ over $K$ with coordinate functions
$\beta_l^*$, $\sigma=\sum_l\beta_l\Lambda_{\beta_l^*}(\sigma)$, so $\sigma$ is determined by $v$. In the same
way $v=\sum_l\beta_l\Lambda_{\beta_l^*}(v)$ is determined by $\sigma$, through $e^{-1}$.

Linearity. If $\sigma$ and $\sigma'$ are expansions of $v$ and $v'$ and $c\in F$, joining the two sums
shows that $c\sigma+\sigma'$ is an expansion of $cv+v'$.
\end{proof}

\begin{lemma}[Alp\"oge]\label{cv-toroidal-torus-primitive}
In Situation~\ref{cv-toroidal-torus-situation}, let $w_0$ be an isomorphism from $\gamma^*L$ to the
restriction of $L_G$ to $G_K$. Then there is a section $\vartheta$ of $L$ over $A_F$ that generates the
sections of $L$ over $A_F$ and whose Fourier expansion with respect to $w_0$
(Definition~\ref{cv-toroidal-torus-expansion}(2)) is primitive (Definition~\ref{cv-toroidal-torus-points}(4)).
\end{lemma}
\leanref{Development/AbelianVarietyTheory/Proofs/UniformizationProofs/ExistsPrimGenerator.lean}{GlobalTheory.UT.exists_prim_generator}
\begin{proof}
Let $Q$ be the function of Lemma~\ref{cv-toroidal-torus-quadratic}, $r$ the rank of $X$ and $g$ the
relative dimension of $B$. Families of Fourier coefficients have weight $1$ throughout. If
$\nu(\sigma_\chi)\ge s$ and $\nu(\zeta)\ge q$ for $\zeta\in F$, then $\nu(\zeta\sigma_\chi)\ge s+q$, by
multiplying the coefficients in Definition~\ref{cv-toroidal-torus-points}(4) by $\zeta$.

The primitive family. As $1^{r+g}=1$, Lemma~\ref{cv-toroidal-torus-ball} for $k=1$ gives one family $\eta$,
nonzero since it is linearly independent, with $\nu(\eta_\chi)\ge Q(\chi)$ for every $\chi$, such that the
families $\sigma$ with $\nu(\sigma_\chi)\ge Q(\chi)$ for every $\chi$ are the $c\eta$ with $\nu(c)\ge0$. The
family $\eta$ is primitive. Suppose that $r'>0$ is rational and $\nu(\eta_\chi)\ge r'+Q(\chi)$ for every
$\chi$. Lemma~\ref{cv-toroidal-torus-rational} gives $\zeta\in F^\times$ with $\nu(\zeta)=r'$. Then
$\nu(\zeta^{-1}\eta_\chi)\ge Q(\chi)$ for every $\chi$, so $\zeta^{-1}\eta=c\eta$ with $\nu(c)\ge0$. As
$\eta\ne0$, $\zeta c=1$; but $\nu(\zeta c)=r'+\nu(c)>0=\nu(1)$, a contradiction.

Multiples of $\eta$. Let $\sigma$ be a family with $\nu(\sigma_0)\ge s$ for some rational $s$, and let
$s'=|s|+1$. As $Q(0)=0$ and $-s'\le s$, $\nu(\sigma_0)\ge-s'+Q(0)$. Every element of $X$ lies in the class
of $0$ modulo $1\cdot X$, so Lemma~\ref{cv-toroidal-torus-shift}(3), for $k=1$ and $S=\{0\}$, gives
$\nu(\sigma_\chi)\ge-s'+Q(\chi)$ for every $\chi$. Lemma~\ref{cv-toroidal-torus-rational} gives
$\zeta\in F^\times$ with $\nu(\zeta)=s'$. Then $\nu(\zeta\sigma_\chi)\ge Q(\chi)$ for every $\chi$, so
$\zeta\sigma=c\eta$ for some $c$, and $\sigma=\zeta^{-1}c\eta$.

The section. By Lemma~\ref{cv-toroidal-torus-bijection} for $k=1$, taking Fourier expansions with respect
to $w_0$ is an $F$-linear bijection $e$ from the sections of $L$ over $A_F$ onto the families of weight
$1$ over $F$. Let $\vartheta=e^{-1}(\eta)$; its Fourier expansion $\eta$ is primitive. Let $t$ be a section
of $L$ over $A_F$. By Definition~\ref{cv-toroidal-torus-expansion}(2), $e(t)=\sum_ia_i\sigma^{(i)}$ with
$a_i\in F$ and families $\sigma^{(i)}$ of Fourier coefficients of sections of $L_G$ over $G$, whose
components are sections over $B$. So $\nu(e(t)_0)\ge s$ for a rational $s$ at most every $\nu(a_i)$, and
$e(t)=c\eta$ for some $c\in F$ by the previous step. Then $t=c\vartheta$, since $e$ is linear and injective.
\end{proof}
% Lean: exists_prim_generator: PrimGen.exists_prim_span (from UUStatement at k = 1, which is
% UnitBallFree.main, with n1c and exists_val_eq for primitivity and n1a with NuGe.smul for the multiples),
% exists_linearEquiv_expandsF (expandsF_bijective with U.expand) for e, and PrimGen.exists_vge with
% exists_lower for the bound on e(t)_0. It takes LL, TwistBCSpanStatement, VRatioGtStatement, TWStatement and
% UUStatement as hypotheses, which the archive proves (lawsAndExpansions).

\begin{definition}\label{cv-toroidal-torus-normalized}
Let $K$ be a field, let $A$ be an abelian variety over $K$ with a line bundle $L$ that is trivialized at the
origin, symmetric, ample and defines a principal polarization, let $\rho$ be a trivialization of $L$ at the
origin, let $F$ be a field containing $K$ and let $m$ be an odd integer. Let $\vartheta$ be a section of $L$
over $A_F$ and $x\in A(F)$ with $m^2x=0$. An element $c\in F$ is a \emph{normalized theta constant} of
$\vartheta$ at $x$, with a section $t$ of $L^{m^2}$ over $A_F$, if there are an isomorphism
$\alpha:[m]^*L\to L^{m^2}$ over $A_F$ that carries $\rho$ to $\rho^{m^2}$ at the origin and a symmetric lift
$\theta$ of $x$ to the theta group of $L^{m^2}$ over $A_F$ with $\theta^{m^2}=1$, such that
$t=\theta\,\alpha([m]^*\vartheta)$ and the value of $t$ at the origin is $c\rho^{m^2}$.
\end{definition}
\leanref{Development/AbelianVarietyTheory/Definitions/Uniformization/AllDV.lean}{GlobalTheory.UT.ThC}

\begin{lemma}[Alp\"oge]\label{cv-toroidal-torus-constant}
Let $K$ be a field of characteristic zero, let $A$ be an abelian variety over $K$ with a line bundle $L$
that is trivialized at the origin, symmetric, ample and defines a principal polarization, and let $\rho$ be
a trivialization of $L$ at the origin. Let $F$ be a field containing $K$, let $m$ be odd, let $\vartheta$ be
a section of $L$ over $A_F$ and let $x\in A(F)$ with $m^2x=0$. Then there is a unique pair $(t,c)$ such that
$c$ is a normalized theta constant of $\vartheta$ at $x$ with $t$
(Definition~\ref{cv-toroidal-torus-normalized}).
\end{lemma}
\leanref{Development/AbelianVarietyTheory/Proofs/UniformizationProofs/ExistsUniqueThetaConst.lean}{GlobalTheory.existsUnique_thetaConst}
\leanref{Development/AbelianVarietyTheory/Proofs/UniformizationProofs/ExistsUniqueThetaConstStatement.lean}{GlobalTheory.T.existsUnique_theta_const_of_odd}
\begin{proof}
Existence. By the theorem of the cube there is an isomorphism $[m]^*L\cong L^{m^2}$ on $A$
\cite{Mum70}. Over $A_F$, multiply it by a constant so that it carries $\rho$ to $\rho^{m^2}$ at the
origin, and call the result $\alpha$. Since $m$ is odd and $m^2$ is invertible in $F$, $x$ has a symmetric
lift $\theta$ to the theta group of $L^{m^2}$ with $\theta^{m^2}=1$ \cite{Mum66}; it exists locally on
$\operatorname{Spec}F$, which is a point. Put $t=\theta\,\alpha([m]^*\vartheta)$, and let $c$ be the value of
$t$ at the origin divided by $\rho^{m^2}$.

Uniqueness. Let $(\alpha',\theta')$ give a second pair $(t',c')$. The quotient $\alpha'/\alpha$ is an
invertible function on $A_F$, hence a constant, as the global functions on $A_F$ are those of $F$; it is $1$
at the origin, so $\alpha'=\alpha$. Since $m$ is odd, $x$ has only one symmetric lift $\theta$ with
$\theta^{m^2}=1$, so $\theta'=\theta$. Hence $t'=t$, and $c'=c$ by comparing values at the origin.
\end{proof}
% Lean: existsUnique_thetaConst: Laws.mulPull with exists_gen and exists_normalised for alpha,
% ThetaPtExist.exists_isTwistPull for alpha([m]^* th), Laws.symLift with GlobalLift.exists_section and
% exists_symLift_pull for theta, exists_value for c; uniqueness by gen_unique (with CommPairing.surj_fst
% of Laws.globalFns) and the uniqueness part of Laws.symLift.

\begin{lemma}[Alp\"oge]\label{cv-toroidal-torus-family}
Let $K$ be a field, let $A$ be an abelian variety over $K$ with a line bundle $L$ that is trivialized at the
origin, symmetric, ample and defines a principal polarization, and let $\rho$ be a trivialization of $L$ at
the origin. Let $F$ be a field containing $K$, let $m$ be an odd integer that is nonzero in $F$, and let
$\vartheta$ be a section of $L$ over $A_F$ that generates the sections of $L$ over $A_F$. Let $I$ be a
nonempty set, and for $i\in I$ let $x_i\in A(F)$ with $m^2x_i=0$ and let $c_i$ be a normalized theta
constant of $\vartheta$ at $x_i$ (Definition~\ref{cv-toroidal-torus-normalized}). Then $(c_i)_{i\in I}$ is a
family of theta constants of the points $x_i$: the conditions of Definition~\ref{cv-toroidal-theta}(2),
with $N'=m$ and the points $x_i$ in place of the $x_c$, hold over $\operatorname{Spec}F$ with $v_i=c_i$.
\end{lemma}
\leanref{Development/AbelianVarietyTheory/Proofs/ModuliProofs/GoodReductionOfThetaPointUnits.lean}{GlobalTheory.ThetaCrit.isThetaPt_of_normalised}
\begin{proof}
For each $i$ let $\alpha_i$, $\theta_i$ and $t_i$ be the data of Definition~\ref{cv-toroidal-torus-normalized}
for $c_i$. For $i,i'\in I$, the quotient $\alpha_{i'}/\alpha_i$ is an invertible function on $A_F$, hence a
constant, and it is $1$ at the origin; so all the $\alpha_i$ are one isomorphism $\alpha$, and
$\alpha([m]^*\vartheta)$ does not depend on $i$. Then the generator $\vartheta$, the isomorphism $\alpha$,
the symmetric lifts $\theta_i$, the trivialization $\rho^{m^2}$ of $L^{m^2}$ at the origin and the unit $1$
satisfy Definition~\ref{cv-toroidal-theta}(2) over all of $\operatorname{Spec}F$, with $v_i=c_i$; and $m$ is
odd and invertible in $F$.
\end{proof}
% Lean: isThetaPt_of_normalised: gen_unique (with CommPairing.surj_fst of GlobalFnsSt) for the alpha_i,
% then isThetaPtOver_of_eq with u = 1 and ThetaCoord.isValueAt_of_sv, over the open set Spec F.

\subsection{Theta constants as sums of theta terms}\label{subsection-cv-theta-terms}

\begin{definition}\label{cv-toroidal-torus-terms}
In Situation~\ref{cv-toroidal-torus-situation}, let $O$ be the ring of elements of $F$ of nonnegative
valuation, let $e$ be the trivialization of $M$ at the origin, and for $x\in B(F)$ let $M(x)$ be the fibre
of $M$ at $x$.
\begin{enumerate}
\item Let $x_1,\dots,x_n\in B(F)$ and $\epsilon_1,\dots,\epsilon_n\in\mathbf Z$, let $t$ be a trivialization
  of $\bigotimes_iM(x_i)^{\epsilon_i}$, and let $r$ be rational. Then $t$ \emph{has valuation} $r$ if there
  are points $x_i^0\in B(O)$ that restrict to the $x_i$, a trivialization $t_0$ over $O$ of
  $\bigotimes_iM(x_i^0)^{\epsilon_i}$ and $a\in F$ with $\nu(a)=r$ and $t=at_0$. A point $g=(b,s)$ of
  $\tilde G(F)$ (Definition~\ref{cv-toroidal-torus-points}(1)) \emph{has valuations} $l:X\to\mathbf Q$ if
  $s(\chi)$ has valuation $l(\chi)$ for every $\chi\in X$.
\item Let $k\ge1$ and $b\in B(F)$. A \emph{frame} of weight $k$ at $b$ is a trivialization of
  $M(0)^k\otimes M(b)^{-k}$.
\item Let $k\ge1$, let $\sigma$ be a family of Fourier coefficients of weight $k$ over $F$, let
  $g=(b,s)\in\tilde G(F)$ and let $w$ be a frame of weight $k$ at $b$. For $\chi\in X$, the \emph{theta term}
  $a_\chi(\sigma,g,w)\in F$ is given by $\sigma_\chi(b)\otimes s(\chi)\otimes w=a_\chi(\sigma,g,w)\,e^k$, where
  $\sigma_\chi(b)$ is the value of $\sigma_\chi$ at $b$, in the fibre of $M_{k,\chi}$ at $b$.
\item Let $(x_i)_{i\in I}$ be a family of elements of $F$ and let $c\in F$. Then $c=\sum_ix_i$
  \emph{$\nu$-adically} if for every rational $q$ there is a finite set $S\subset I$ with
  $\nu\bigl(c-\sum_{i\in T}x_i\bigr)\ge q$ for every finite $T$ with $S\subset T\subset I$. And
  $c=\sum_ix_i$ \emph{$\mathfrak m$-adically} if there is a finitely generated $W$-submodule $N$ of $F$
  containing $c$ and every $x_i$ such that for every $n\ge0$ there is a finite set $S\subset I$ with
  $c-\sum_{i\in T}x_i\in\mathfrak m^nN$ for every finite $T$ with $S\subset T\subset I$.
\end{enumerate}
\end{definition}
\leanref{Development/AbelianVarietyTheory/Definitions/Uniformization/AllDV.lean}{GlobalTheory.UT.GenV}
\leanref{Development/AbelianVarietyTheory/Definitions/Uniformization/AllDV.lean}{GlobalTheory.UT.GtV}
\leanref{Development/AbelianVarietyTheory/Definitions/Uniformization/TXFrame.lean}{GlobalTheory.UT.Frm}
\leanref{Development/AbelianVarietyTheory/Definitions/Uniformization/TXFrame.lean}{GlobalTheory.UT.IsCoef}
\leanref{Development/AbelianVarietyTheory/Definitions/Uniformization/TXFrame.lean}{GlobalTheory.UT.coef}
\leanref{Development/AbelianVarietyTheory/Definitions/Uniformization/AllDV.lean}{GlobalTheory.UT.SumV}
\leanref{Development/AbelianVarietyTheory/Definitions/Uniformization/TXFrame.lean}{GlobalTheory.UT.SumI}
\leanref{Development/AbelianVarietyTheory/Definitions/Uniformization/TXFrame.lean}{GlobalTheory.UT.SumF}

\begin{lemma}\label{cv-toroidal-torus-sums}
In Situation~\ref{cv-toroidal-torus-situation}, let $(x_i)_{i\in I}$ be a family of elements of $F$, let
$c\in F$ and let $q$ be rational.
\begin{enumerate}
\item Suppose that $c=\sum_ix_i$ $\nu$-adically (Definition~\ref{cv-toroidal-torus-terms}(4)) and that
  $\nu(x_i)\ge q$ for every $i$. Then $\nu(c)\ge q$.
\item Suppose that $c=\sum_ix_i$ $\nu$-adically, that $q>0$, and that there is $i_0\in I$ with
  $\nu(x_{i_0})=0$ and $\nu(x_i)\ge q$ for every $i\ne i_0$. Then $\nu(c)=0$.
\item If $c=\sum_ix_i$ $\mathfrak m$-adically, then $c=\sum_ix_i$ $\nu$-adically.
\end{enumerate}
\end{lemma}
\leanref{Development/AbelianVarietyTheory/Proofs/UniformizationProofs/ThetaConstantValuationsUnitsEveryCoset.lean}{GlobalTheory.UT.ThetaConstVal.le_of_sumV}
\leanref{Development/AbelianVarietyTheory/Proofs/UniformizationProofs/ThetaConstantValuationsUnitsEveryCoset.lean}{GlobalTheory.UT.ThetaConstVal.eq_zero_of_sumV}
\leanref{Development/AbelianVarietyTheory/Proofs/UniformizationProofs/ExistsLeOnFinset.lean}{GlobalTheory.UT.AdicValSum.main}
\begin{proof}
(1) Let $S$ be the finite set given for $q$. Then $\nu(c-\sum_{i\in S}x_i)\ge q$ and
$\nu(\sum_{i\in S}x_i)\ge q$, so $\nu(c)\ge q$.

(2) Let $S$ be the finite set given for $q$, and $T=S\cup\{i_0\}$. Then
$c=x_{i_0}+\bigl(c-\sum_{i\in T}x_i\bigr)+\sum_{i\in T\setminus\{i_0\}}x_i$, where the last two terms have
valuation at least $q>0=\nu(x_{i_0})$. So $\nu(c)=0$.

(3) Let $N$ be generated by $z_1,\dots,z_l$, and let $c_0$ be a rational number at most every $\nu(z_j)$.
As $\nu(a)\ge0$ for $a\in W$ and $\nu(a)\ge n$ for $a\in\mathfrak m^n$, every element of $\mathfrak m^nN$
has valuation at least $n+c_0$. Given $q$, choose an integer $n\ge q-c_0$, and let $S$ be the finite set
given for $n$. Then $\nu(c-\sum_{i\in T}x_i)\ge n+c_0\ge q$ for every finite $T\supset S$.
\end{proof}

\begin{definition}\label{cv-toroidal-torus-symmetric}
In Situation~\ref{cv-toroidal-torus-situation}, let $N\ge1$, let $e$ be the trivialization of $M$ at the
origin, let $M(x)$ be the fibre of $M$ at a point $x$ of $B$, and write $M(0)$ also for the constant line
bundle on $B$ with fibre $M(0)$.
\begin{enumerate}
\item A trivialization $\kappa_N$ of $M^{N^2}\otimes[N]^*M^{-1}\otimes M(0)^{1-N^2}$ over $B$, or $\nu_N$ of
  $[N+1]^*M\otimes[N]^*M^{-1}\otimes M^{-(2N+1)}\otimes M(0)^{2N+1}$ over $B$, is \emph{normalized} if its
  value at the origin is the canonical trivialization of the fibre there.
\item Let $\kappa_N$ and $\nu_N$ be normalized, and let $p=(y,g)$ be an $N$-division point
  (Definition~\ref{cv-toroidal-torus-points}(2)) with $g=(b,s)$. A \emph{symmetric frame} for $p$ is a frame
  $w$ of weight $N$ at $b$ (Definition~\ref{cv-toroidal-torus-terms}(2)) such that
  $w^{\otimes2}=s(y)\otimes\nu_N(b)$ and $w^{\otimes N}\otimes\kappa_N(b)=\psi(y)\otimes e$.
\end{enumerate}
\end{definition}
\leanref{Development/AbelianVarietyTheory/Definitions/Uniformization/TXFrame.lean}{GlobalTheory.UT.Normd}
\leanref{Development/AbelianVarietyTheory/Definitions/Uniformization/TXFrame.lean}{GlobalTheory.UT.kapLB}
\leanref{Development/AbelianVarietyTheory/Proofs/UniformizationProofs/UnivIdentity.lean}{GlobalTheory.UT.SymFrame.nuLB}
\leanref{Development/AbelianVarietyTheory/Definitions/Uniformization/TXFrame.lean}{GlobalTheory.UT.TXFrame}

Here $Nb=c(y)$, as $Ng=\iota(y)$, so $s(y)$ trivializes $M(Nb+b)^{-1}\otimes M(Nb)\otimes M(b)\otimes M(0)^{-1}$
and both conditions compare trivializations of the same lines.

\begin{lemma}[Alp\"oge]\label{cv-toroidal-torus-symmetric-frame}
In Situation~\ref{cv-toroidal-torus-situation}, let $N\ge1$.
\begin{enumerate}
\item There are unique normalized trivializations $\kappa_N$ and $\nu_N$
  (Definition~\ref{cv-toroidal-torus-symmetric}(1)).
\item Suppose that $N$ is odd, and let $p$ be an $N$-division point. Then there is a symmetric frame for $p$,
  and any two symmetric frames for $p$ are equal.
\end{enumerate}
\end{lemma}
\leanref{Development/AbelianVarietyTheory/Proofs/UniformizationProofs/UnivIdentity.lean}{GlobalTheory.UT.SymFrame.exists_kap}
\leanref{Development/AbelianVarietyTheory/Proofs/UniformizationProofs/UnivIdentity.lean}{GlobalTheory.UT.SymFrame.exists_nu}
\leanref{Development/AbelianVarietyTheory/Proofs/UniformizationProofs/UnivIdentity.lean}{GlobalTheory.UT.SymFrame.main}
\leanref{ReciprocityTheory/ReductiveGroups/ArithmeticGeometry/DiscreteValuationRings/PhiGenPf.lean}{ReciprocityTheory.ChartRestriction.G.DVR.chain_val}
\leanref{Development/AbelianVarietyTheory/Proofs/UniformizationProofs/UnivIdentity.lean}{GlobalTheory.UT.SymFrame.tau_self}
\leanref{Development/AbelianVarietyTheory/Proofs/UniformizationProofs/UnivIdentity.lean}{GlobalTheory.UT.SymFrame.univ_identity}
\begin{proof}
(1) As $M$ is symmetric, the theorem of the cube gives an isomorphism $[N]^*M\cong M^{N^2}$ over $B$
\cite[\S6]{Mum70}. Dividing it by its value at the origin gives $\kappa_N$, and
$\kappa_N\otimes\kappa_{N+1}^{-1}$ is a normalized $\nu_N$. Two normalized trivializations of one line bundle
differ by an invertible function on $B$, which is a unit of $W$ since the global functions on $B$ are those
of $W$, and which is $1$ at the origin.

(2) Let $p=(y,g)$ with $g=(b,s)$, and write $N=2k+1$. Let $\theta$ be the cubical structure of $M$ and
$\varsigma$ the symmetry $[-1]^*M\cong M$, and let $C$ be the product of $\theta(Nb,-Nb,Nb)$,
$\varsigma(Nb)^{-1}$ and $\theta(ib,b,Nb)$ for $0\le i<N$. Two computations with the cubical structure give
\[
C\otimes(\psi(y)\otimes e)^{\otimes2}=s(y)^{\otimes N},\qquad C\otimes\nu_N(b)^{\otimes N}\otimes\kappa_N(b)^{\otimes2}=1.
\]
The first comes from $Ng=\iota(y)$, which gives
$\tau(y,y)\otimes\prod_{0\le i<N}\theta(ib,b,Nb)=s(y)^{\otimes N}$ by induction along $g,2g,\dots,Ng$, and from
$\tau(y,y)=\theta(c(y),-c(y),c(y))\otimes\varsigma(c(y))^{-1}\otimes(\psi(y)\otimes e)^{\otimes2}$, which follows
from the bimultiplicativity of $\tau$ and the evenness of $\psi$. The second holds because the corresponding
trivialization over $B$, at the points $Nx$, $-Nx$ and $ix$ for $x$ in $B$, is $1$ at the origin and hence
constant, as the global functions on $B$ are those of $W$. Put
$w=(\psi(y)\otimes e)\otimes\kappa_N(b)^{-1}\otimes(s(y)\otimes\nu_N(b))^{-k}$. The two identities give
$(\psi(y)\otimes e)^{\otimes2}=(s(y)\otimes\nu_N(b))^{\otimes N}\otimes\kappa_N(b)^{\otimes2}$, from which
$w^{\otimes2}=s(y)\otimes\nu_N(b)$ and $w^{\otimes N}\otimes\kappa_N(b)=\psi(y)\otimes e$ follow, as $N=2k+1$.
If $w'$ is a second symmetric frame, then $w'=uw$ with $u\in F^\times$, $u^2=1$ and $u^N=1$; as $N$ is odd,
$u=u^N(u^2)^{-k}=1$.
\end{proof}

\begin{definition}\label{cv-toroidal-torus-theta-sum-phi}
In Situation~\ref{cv-toroidal-torus-situation}, let $k\ge1$, let $p=(y,g)$ be a $k$-division point
(Definition~\ref{cv-toroidal-torus-points}(2)) with $g=(b,s)$, let $w$ be a frame of weight $k$ at $b$
(Definition~\ref{cv-toroidal-torus-terms}(2)), and let $T_b$ be the translation by $b$ on $B_F$. For
$\alpha\in X$, the fibre at the origin of the line bundle of homomorphisms from $T_b^*M_{k,\alpha}$ to
$M_{k,\alpha+y}$ is identified with the line trivialized by $s(\alpha)\otimes w$. Let $\sigma$ and $\sigma'$ be
families of Fourier coefficients of weight $k$ over $F$. Then $\sigma'=\Phi_{p,w}\sigma$ if, for every
$\alpha\in X$ and every isomorphism $u:T_b^*M_{k,\alpha}\to M_{k,\alpha+y}$ over $B_F$ whose value at the origin is
given by $s(\alpha)\otimes w$, $\sigma'_{\alpha+y}=u(T_b^*\sigma_\alpha)$.
\end{definition}
\leanref{ReciprocityTheory/ReductiveGroups/ArithmeticGeometry/ExtensionDescent/GtPtBridgeSt.lean}{ReciprocityTheory.ChartRestriction.G.Out.PhiRel}

\begin{lemma}\label{cv-toroidal-torus-theta-sum-operator}
In Situation~\ref{cv-toroidal-torus-situation}, let $k\ge1$, let $p=(y,g)$ be a $k$-division point with
$g=(b,s)$, let $w$ be a frame of weight $k$ at $b$, and let $T_b$ be the translation by $b$ on $B_F$.
\begin{enumerate}
\item For every $\alpha\in X$ there is exactly one isomorphism $u_\alpha:T_b^*M_{k,\alpha}\to M_{k,\alpha+y}$ over
  $B_F$ whose value at the origin is given by $s(\alpha)\otimes w$.
\item Let $\sigma$ and $\sigma'$ be families of Fourier coefficients of weight $k$ over $F$ with
  $\sigma'=\Phi_{p,w}\sigma$ (Definition~\ref{cv-toroidal-torus-theta-sum-phi}), let $0$ be the zero point of
  $\tilde G(F)$, and let $w_1$ be the frame of weight $k$ at the origin of $B$ given by the canonical
  trivialization of $M(0)^k\otimes M(0)^{-k}$. Then $a_{\alpha+y}(\sigma',0,w_1)=a_\alpha(\sigma,g,w)$ for every
  $\alpha\in X$ (Definition~\ref{cv-toroidal-torus-terms}(3)).
\end{enumerate}
Part (2) is the case of the zero point of a more general identity, in which an arbitrary point $g_0$ of
$\tilde G(F)$ replaces $0$, the point $g_0+g$ replaces $g$, and the frames change accordingly; only the case of the
zero point is stated here.
\end{lemma}
\leanref{ReciprocityTheory/ReductiveGroups/ArithmeticGeometry/DiscreteValuationRings/PhiGenPf.lean}{ReciprocityTheory.ChartRestriction.G.DVR.gen_phi}
\leanref{ReciprocityTheory/ReductiveGroups/ArithmeticGeometry/DiscreteValuationRings/PhiGenPf.lean}{ReciprocityTheory.ChartRestriction.G.DVR.gen_u}
\leanref{Development/AbelianVarietyTheory/Proofs/UniformizationProofs/CubNMul.lean}{GlobalTheory.UT.CoefPhi.coef_phi}
\begin{proof}
(1) As $kg=\iota(y)$, $kb=c(y)$, so $c(\alpha+y)=c(\alpha)+kb$. For a point $x$ of $B_F$, the product of the
values of the cubical structure of $M$ at the points $(c(\alpha)+ib,b,x)$, for $0\le i<k$, divided by its
value at $(c(\alpha),x,b)$, gives an isomorphism from the line bundle of homomorphisms from $T_b^*M_{k,\alpha}$ to
$M_{k,\alpha+y}$ onto the constant line bundle given by its fibre at the origin, as in the proof of
Lemma~\ref{cv-toroidal-torus-translation}(1) with $b$ in place of $c(\mu)$. At the origin it is the identity.
Composing its inverse with the trivialization $s(\alpha)\otimes w$ of the fibre at the origin gives $u_\alpha$.
If $u$ is a second such isomorphism, then $u/u_\alpha$ is an invertible function on $B_F$, hence a constant, as
the global functions on $B_F$ are those of $F$; it is $1$ at the origin, so $u=u_\alpha$.

(2) Let $e$ be the trivialization of $M$ at the origin. The trivializations of the zero point are the
canonical ones (Lemma~\ref{cv-toroidal-torus-extension}(2)), and so is $w_1$. So, by
Definition~\ref{cv-toroidal-torus-terms}(3), $a_{\alpha+y}(\sigma',0,w_1)e^k$ is the value $\sigma'_{\alpha+y}(0)$ of
$\sigma'_{\alpha+y}$ at the origin. By (1) and Definition~\ref{cv-toroidal-torus-theta-sum-phi},
$\sigma'_{\alpha+y}=u_\alpha(T_b^*\sigma_\alpha)$. The value of $T_b^*\sigma_\alpha$ at the origin is $\sigma_\alpha(b)$, and
the value of $u_\alpha$ at the origin carries it to $\sigma_\alpha(b)\otimes s(\alpha)\otimes w$. By
Definition~\ref{cv-toroidal-torus-terms}(3) this is $a_\alpha(\sigma,g,w)e^k$.
\end{proof}
% Lean: (1) gen_phi = gen_u with the point p.g.b and c(alpha + y) = c(alpha) + k b (cPoint_add, cPoint_y);
% gen_u builds the isomorphism from exists_cubN (the value uVal, a product of cube values) and gets
% uniqueness from gen_unique with surj_appTop (GlobalFnsSt). (2) coef_phi with g_0 the zero point
% (zero_isAdd, so g_1 = g), n = 1 and N = k: the Lean proves the identity for any g_0, with the frame
% omega_1 of exists_frame and the cube factor cubN(theta, k, 1, b, g_0.b), and through phiRel_iff (gen_phi)
% for the values; for g_0 = 0 the cube factors are 1 (N_one, face2) and omega_1 = w, which is the case
% written here and the one used.

\begin{lemma}\label{cv-toroidal-torus-theta-sum-origin}
In Situation~\ref{cv-toroidal-torus-situation}, let $\rho$ be a trivialization of $L$ at the origin and $w_0$
a normalization for $\rho$ (Definition~\ref{cv-toroidal-torus-expansion}(1)). Let $k\ge1$, let $t$ be a section
of $L^k$ over $A_F$ with Fourier expansion $\sigma$ with respect to $w_0$, and let $c\in F$ be such that the value
of $t$ at the origin is $c\rho^k$. Let $0$ be the zero point of $\tilde G(F)$, and let $w_1$ be the frame of
weight $k$ at the origin of $B$ given by the canonical trivialization of $M(0)^k\otimes M(0)^{-k}$. Then
$c=\sum_{\beta\in X}a_\beta(\sigma,0,w_1)$ $\mathfrak m$-adically (Definition~\ref{cv-toroidal-torus-terms}).
\end{lemma}
\leanref{Development/AbelianVarietyTheory/Proofs/UniformizationProofs/ValueAtOrigin.lean}{GlobalTheory.UT.EvalZero.value_at_origin}
\leanref{Development/AbelianVarietyTheory/Proofs/RingsOfIntegers/ExpansionValueZero.lean}{GlobalTheory.ExpansionValueZero.value_sub_sum_mem}
\begin{proof}
Let $e$ be the trivialization of $M$ at the origin and $e_G$ the trivialization of $L_G$ along the zero
section. By Definition~\ref{cv-toroidal-torus-expansion}(2) there are sections $t_1,\dots,t_n$ of $L_G^k$ over
$G$ and $c_1,\dots,c_n\in F$ such that $w_0$ and $\gamma$ carry $t$ to $\sum_ic_it_i$ and
$\sigma=\sum_ic_i\sigma^{(i)}$, where $\sigma^{(i)}$ is the family of Fourier coefficients of $t_i$. As $W$ is
local, the value of $t_i$ at the origin of $G$ is $x_ie_G^k$ with $x_i\in W$, and the value of the section
$\sigma^{(i)}_\beta$ of $M_{k,\beta}$ over $B$ at the origin, a point of the fibre $M(0)^k$, is $a_{i\beta}e^k$ with
$a_{i\beta}\in W$. As $w_0$ carries $\rho$ to $e_G$ at the origin, $c=\sum_ic_ix_i$. The trivializations of the
zero point are the canonical ones (Lemma~\ref{cv-toroidal-torus-extension}(2)), and so is $w_1$; so, by
Definition~\ref{cv-toroidal-torus-terms}(3), $a_\beta(\sigma,0,w_1)=\sum_ic_ia_{i\beta}$.

Let $e'\ge0$ and $i\le n$. By the uniformization (Theorem~\ref{cv-toroidal-mumford}), there is a finite set
$S_i\subset X$ such that $\sigma^{(i)}_\beta$ vanishes modulo $\mathfrak m^{e'+1}$ for $\beta\notin S_i$ and, modulo
$\mathfrak m^{e'+1}$, $t_i$ corresponds to the sum over $\beta\in S_i$ of $\sigma^{(i)}_\beta$ times the
trivialization of the character $\beta$ of the torus, compatibly with the trivializations at the origin. The
trivializations of the characters are $1$ along the unit section. Evaluating at the origin gives
$x_i-\sum_{\beta\in T}a_{i\beta}\in\mathfrak m^{e'+1}$ for every finite $T\supset S_i$. Let $N$ be the $W$-submodule
of $F$ generated by $c_1,\dots,c_n$. It contains $c$ and every $a_\beta(\sigma,0,w_1)$, and for every finite $T$
containing $S_1\cup\dots\cup S_n$,
\[
c-\sum_{\beta\in T}a_\beta(\sigma,0,w_1)=\sum_ic_i\Bigl(x_i-\sum_{\beta\in T}a_{i\beta}\Bigr)\in\mathfrak m^{e'+1}N.
\]
So $c=\sum_\beta a_\beta(\sigma,0,w_1)$ $\mathfrak m$-adically.
\end{proof}
% Lean: the third part of EvalZero.value_at_origin: ExpandsF unfolded (sigma', the IsoRelM through w_0, and
% ExpandsG: n, c, sigma, f with ExpandsAt e for every e), exists_value and exists_coeff for a and x,
% value_comb for c = sum c_i x_i (the normalization w_0), coef_comb at the zero point with FrmOne, and
% sumF_comb with sumF_map and ExpansionValueZero.value_sub_sum_mem (expandsAt_val at the unit section unitE
% of the thickening, t_unit for the characters and w_rig for the trivializations at the origin).

\begin{definition}\label{cv-toroidal-torus-expansion-coefficients}
In Situation~\ref{cv-toroidal-torus-situation}, let $e$ be the trivialization of $M$ at the origin, $e_G$ the
trivialization of $L_G$ along the zero section, $\theta$ the cubical structure of $M$, and $\Lambda(M)$ the line
bundle on $B\times B$ of Theorem~\ref{cv-toroidal-mumford}. For $n\ge0$ let $W_n=W/\mathfrak m^{n+1}$, and let
$G_n$ and $B_n$ be the reductions of $G$ and $B$ over $W_n$. The uniformization (Theorem~\ref{cv-toroidal-mumford})
gives, for every $n\ge0$, a surjective homomorphism $q_n:G_n\to B_n$ whose kernel is the split torus $T_n$ over
$W_n$ with character group $X$, an isomorphism $w_n$ from the restriction of $L_G$ to $G_n$ onto $q_n^*M$ that
carries $e_G$ to $e$ at the origin, and, for every $\chi\in X$, a trivialization $t_{n,\chi}$ over $G_n$ of the
pullback of $\Lambda(M)^{-1}$ along $x\mapsto(c(\chi),q_n(x))$. These have the following properties: $t_{n,\chi}$
is $1$ at the origin; $t_{n,\chi}(\tau x)=\chi(\tau)\,t_{n,\chi}(x)$ for points $\tau$ of $T_n$ and $x$ of $G_n$;
and $t_{n,\chi+\chi'}(x)$, multiplied by $\theta(c(\chi),c(\chi'),q_n(x))$, corresponds to
$t_{n,\chi}(x)\,t_{n,\chi'}(x)$.

Let $k\ge1$ and let $\sigma$ be a section of $L_G^k$ over $G$. A family $f=(f_\chi)_{\chi\in X}$ of sections
$f_\chi$ over $B$ of the line bundles $M_{k,\chi}$ of Definition~\ref{cv-toroidal-torus-points}(3) is a
\emph{family of Fourier coefficients} of $\sigma$ if for every $n\ge0$ there is a finite set $S\subset X$ such that
$f_\chi$ vanishes on $B_n$ for $\chi\notin S$ and, on $G_n$,
\[
w_n^{\otimes k}(\sigma)=\sum_{\chi\in S}q_n^*f_\chi\otimes t_{n,\chi}.
\]
\end{definition}
\leanref{ReciprocityTheory/ReductiveGroups/ArithmeticGeometry/ExtensionDescent/DegUnif.lean}{ReciprocityTheory.ChartRestriction.G.Out.GtCoord}
\leanref{ReciprocityTheory/ReductiveGroups/ArithmeticGeometry/ExtensionDescent/DegExp.lean}{ReciprocityTheory.ChartRestriction.G.Out.UnifPre.ExpandsAt}

The fibre of $M_{k,\chi}$ at a point $x$ of $B$ is $M(x)^{k-1}\otimes M(x+c(\chi))\otimes M(c(\chi))^{-1}\otimes
M(0)$, so $q_n^*f_\chi\otimes t_{n,\chi}$ is a section of $q_n^*M^k$ once $M(0)$ is trivialized by $e$. The
families of Fourier coefficients of sections over $G$ in Definition~\ref{cv-toroidal-torus-expansion}(2) are
families in this sense.

\begin{definition}\label{cv-toroidal-torus-expansion-endomorphism}
In Situation~\ref{cv-toroidal-torus-situation}, with the notation of
Definition~\ref{cv-toroidal-torus-expansion-coefficients}, an \emph{endomorphism of the uniformization} is a tuple
$E=(a_G,a_A,a_B,h,u)$ of homomorphisms of group schemes $a_G:G\to G$, $a_A:A\to A$ and $a_B:B\to B$, a
homomorphism $h:X\to X$ and, for every $\chi\in X$, a trivialization $u(\chi)$ over $B$ of the pullback of
$\Lambda(M)$ along $x\mapsto(c(h(\chi)),x)$ tensored with the pullback of $\Lambda(M)^{-1}$ along
$x\mapsto(c(\chi),a_B(x))$, equal to the canonical trivialization at the origin, with the following properties:
$\gamma\circ a_G=a_A\circ\gamma$ on $G_K$; and for every $n\ge0$, $q_n\circ a_G=a_B\circ q_n$ on $G_n$, $a_G$
restricts on $T_n$ to the endomorphism given by $h$, and $t_{n,\chi}(a_G(x))=t_{n,h(\chi)}(x)\otimes
u(\chi)(q_n(x))$ for every $\chi\in X$ and every point $x$ of $G_n$. For points $g=(b,s)$ and $g'=(b',s')$ of
$\tilde G(F)$ (Definition~\ref{cv-toroidal-torus-points}(1)), $E$ \emph{carries} $g$ to $g'$ if $b'=a_B(b)$ and
$s'(\chi)=s(h(\chi))\otimes u(\chi)(b)$ for every $\chi\in X$.
\end{definition}
\leanref{Development/AbelianVarietyTheory/Definitions/Uniformization/TXFrame.lean}{GlobalTheory.UT.EndU}
\leanref{Development/AbelianVarietyTheory/Definitions/Uniformization/TXFrame.lean}{GlobalTheory.UT.EndU.Act}

\begin{lemma}\label{cv-toroidal-torus-expansion-multiples}
In Situation~\ref{cv-toroidal-torus-situation}, with the notation of
Definition~\ref{cv-toroidal-torus-expansion-coefficients}, the following hold.
\begin{enumerate}
\item Let $m\ge0$. Multiplication by $m$ on $G$, $A$, $B$ and $X$, with the trivializations $u(\chi)$ whose value
  at a point $x$ of $B$ is $\prod_{0\le i<m}\theta(ic(\chi),c(\chi),x)\,\theta(c(\chi),ix,x)^{-1}$, is an
  endomorphism $E_m$ of the uniformization (Definition~\ref{cv-toroidal-torus-expansion-endomorphism}). For
  $g,g'\in\tilde G(F)$, $E_m$ carries $g$ to $g'$ if and only if $g'=mg$.
\item Multiplication by $-1$ on $G$, $A$, $B$ and $X$, with the trivializations $u(\chi)$ whose value at a point
  $x$ of $B$ is $\theta(-c(\chi),c(\chi),x)^{-1}\,\theta(c(\chi),-x,x)$, is an endomorphism $E_{-1}$ of the
  uniformization. For $g,g'\in\tilde G(F)$, $E_{-1}$ carries $g$ to $g'$ if and only if $g+g'=0$.
\end{enumerate}
\end{lemma}
\leanref{Development/AbelianVarietyTheory/Proofs/UniformizationProofs/ExistsNsmulAndNeg.lean}{GlobalTheory.IntEnd.exists_nsmul_and_neg}
\leanref{Development/AbelianVarietyTheory/Proofs/UniformizationProofs/ExistsNsmulAndNeg.lean}{GlobalTheory.IntEnd.exists_isCmp}
\leanref{Development/AbelianVarietyTheory/Proofs/UniformizationProofs/ExistsNsmulAndNeg.lean}{GlobalTheory.IntEnd.intEnd_act}
\leanref{Development/AbelianVarietyTheory/Proofs/ThetaShiftProofs/ExistsRatioPoints.lean}{GlobalTheory.TorusCoord.val_t_mul_points}
\begin{proof}
Let $n$ be $m$ in (1) and $-1$ in (2), and let $l\ge0$. Multiplication by $n$ is an endomorphism of each of the
commutative group schemes $G$, $A$, $B$ and $T_l$, so it commutes with $\gamma$ and with $q_l$, and on $T_l$ it is the
endomorphism given by multiplication by $n$ on $X$. The trivialization $u(\chi)$ is a product of pullbacks of
$\theta$ along morphisms $B\to B^3$, and its value at the origin is canonical, since $\theta$ is $1$ when one of
its arguments is the origin.

The two remaining conditions come from one identity. Let $x$ be a point of $B$, let $s_0$ be a trivialization, and
let $(a_j)_{j\in\mathbf Z}$ and $(b_j)_{j\in\mathbf Z}$ be families of trivializations with $a_0=b_0=1$,
$a_{j+1}\,\theta(jc(\chi),c(\chi),x)=a_js_0$ and $b_{j+1}\,\theta(c(\chi),jx,x)=b_js_0$ for every $j$. Then
$b_n=a_n\,u(\chi)(x)$. For $n=m\ge0$, $a_m$ and $b_m$ are $s_0^m$ divided by the products over $j<m$ of the values
of $\theta$ in the two recursions. For $n=-1$, the recursions at $j=-1$ give
$a_{-1}=s_0^{-1}\theta(-c(\chi),c(\chi),x)$ and $b_{-1}=s_0^{-1}\theta(c(\chi),-x,x)$.

The torus. For points $x$ and $y$ of $G_l$,
$t_{l,\chi}(x+y)\,\theta(c(\chi),q_l(x),q_l(y))=t_{l,\chi}(x)\,t_{l,\chi}(y)$. Indeed, the ratio of the two sides is
an invertible function on $G_l\times G_l$. By the rule for $t_{l,\chi}(\tau x)$ it does not change when $x$ is moved
by a point of $T_l$, so it comes from a function on $B_l\times G_l$, which is the pullback of a function on $G_l$,
as the global functions on the abelian scheme $B_l$ are those of $W_l$. At $x=0$ the ratio is $1$, since
$t_{l,\chi}$ is $1$ at the origin and $\theta(c(\chi),0,q_l(y))=1$. Now fix a point $x$ of $G_l$ and put
$a_j=t_{l,j\chi}(x)$, $b_j=t_{l,\chi}(jx)$ and $s_0=t_{l,\chi}(x)$, with $q_l(x)$ in place of $x$ in the
recursions. The multiplicativity of $t_{l,\chi}$ in $\chi$ and the identity just shown are the two recursions, so
$t_{l,\chi}(nx)=t_{l,n\chi}(x)\otimes u(\chi)(q_l(x))$.

Points of $\tilde G(F)$. Let $g=(b,s)$, and write $s_{jg}$ for the trivializations of the point $jg$. Put
$a_j=s(j\chi)$, $b_j=s_{jg}(\chi)$ and $s_0=s(\chi)$, with $b$ in place of $x$. By
Definition~\ref{cv-toroidal-torus-points}(1), the condition on $s$ and the sum $(j+1)g=jg+g$ are the two
recursions, so $s_{ng}(\chi)=s(n\chi)\otimes u(\chi)(b)$. As the first component of $ng$ is $nb$, $E_n$ carries
$g$ to $ng$. A point of $\tilde G(F)$ is determined by its first component and its trivializations, so $E_n$
carries $g$ to $g'$ if and only if $g'=ng$; for $n=-1$ this says that $g+g'=0$.
\end{proof}
% Lean: IntEnd.intEnd: ext, q_lin and i_lin from map_id_zpow, id_zpow_comm and torusEndOf_zpow; u from
% exists_isCmp (IsCmp is the identity b_n = a_n u for all such sequences, proved for every n in Z by induction
% with exists_step), IsCmp.normd for the value at the origin; t_lin from val_t_zpow (val_t_add, which is t_mul, and
% TorusCoord.val_t_mul_points, whose proof is the ratio argument: exists_univ_ratio, val_torus, indep_fst,
% cancel_unit); the action from intEnd_act (charVal_zsmul with charVal_add, add_val and gt_ext_val). The text writes
% out the identity for n = m >= 0 and n = -1, the cases exists_nsmul_and_neg uses. Size: ExistsNsmulAndNeg.lean,
% 276 lines, and ExistsRatioPoints.lean.

\begin{remark}\label{cv-toroidal-torus-descent-structures}
In Situation~\ref{cv-toroidal-torus-situation}, with the notation of
Definition~\ref{cv-toroidal-torus-expansion-coefficients}, the uniformization also includes the following structures,
with the following properties. For a point $x$ of $G$ or of $B$, write $L_G(x)$ and $M(x)$ for the fibres of $L_G$ and
$M$ at $x$. For points
$x_1,x_2,x_3$ of a commutative group scheme and $I\subset\{1,2,3\}$, write $x_I=\sum_{j\in I}x_j$ and
$\delta_I=(-1)^{|I|+1}$; tensor products over $I$ run over the subsets of $\{1,2,3\}$.
\begin{enumerate}
\item A cubical structure $\theta_G$ of $L_G$: a trivialization over $G\times_WG\times_WG$ of the line bundle whose
  fibre at $(x_1,x_2,x_3)$ is $\bigotimes_IL_G(x_I)^{\delta_I}$; it is $1$ when one of its arguments is the origin.
  The cubical structure $\theta$ of $M$ is the same kind of trivialization over $B\times_WB\times_WB$.
\item A symmetry $\varsigma_G$ of $L_G$: a trivialization over $G$ of $[-1]^*L_G\otimes L_G^{-1}$ whose value at the
  origin is the canonical trivialization. The symmetry $\varsigma$ of $M$ is a trivialization over $B$ of
  $[-1]^*M\otimes M^{-1}$.
\item For every $n\ge0$, the morphism $q_n$ is flat and affine, and $w_n$ respects these structures: for points
  $x_1,x_2,x_3$ of $G_n$ the isomorphism $\bigotimes_Iw_n(x_I)^{\otimes\delta_I}$ carries $\theta_G(x_1,x_2,x_3)$ to
  $\theta(q_n(x_1),q_n(x_2),q_n(x_3))$, and for a point $x$ of $G_n$ the isomorphism
  $w_n(-x)\otimes w_n(x)^{\otimes(-1)}$ carries $\varsigma_G(x)$ to $\varsigma(q_n(x))$.
\end{enumerate}
\end{remark}
\leanref{ReciprocityTheory/ReductiveGroups/ArithmeticGeometry/ExtensionDescent/DegUnif.lean}{ReciprocityTheory.ChartRestriction.G.Out.UnifPre}
\leanref{Development/AbelianVarietyTheory/Proofs/UniformizationProofs/ComparisonTorusInvariant.lean}{GlobalTheory.TorusInvariance.cube_val}
\leanref{Development/AbelianVarietyTheory/Proofs/UniformizationProofs/ComparisonTorusInvariant.lean}{GlobalTheory.TorusInvariance.even_val}

\begin{lemma}\label{cv-toroidal-torus-descent-functions}
In Situation~\ref{cv-toroidal-torus-situation}, every global function on $G$, and every global function on
$G\times_WG\times_WG$, is the pullback of an element of $W$.
\end{lemma}
\leanref{Development/AbelianVarietyTheory/Proofs/UniformizationProofs/DescentPowerIsoEndomorphismAbelianPartUniformiza.lean}{GlobalTheory.UT.PullDescent.surj_base}
\leanref{Development/AbelianVarietyTheory/Proofs/UniformizationProofs/RigidifiedIsoRespectsCubeAndSymmetry.lean}{GlobalTheory.RigidIsoCube.surjective_snd_of_generic}
\begin{proof}
Let $T$ be a scheme flat over $W$, let $p:G\times_WT\to T$ be the projection, and let $e:T\to G\times_WT$ be the
section $x\mapsto(0,x)$. Then every global function $f$ on $G\times_WT$ is $p^*e^*f$. Indeed, $G$ is smooth over
$W$, so $G\times_WT$ is flat over the domain $W$, and a global function on it is determined by its restriction to
the generic fibre. Through $\gamma$, the generic fibre is $A\times_KT_K$, an abelian scheme over $T_K$, whose global
functions are those of $T_K$. So the restriction of $f$ is the pullback of a function $c$ on $T_K$, and pulling back
along $e$ shows that $c$ is the restriction of $e^*f$. Hence $f$ and $p^*e^*f$ have the same restriction to the
generic fibre, and they are equal. For $T=\operatorname{Spec}W$ this is the statement for $G$. For $T=G$ and
$T=G\times_WG$ it shows that the global functions on $G\times_WG$ come from $G$ and those on $G\times_WG\times_WG$
come from $G\times_WG$, which gives the statement for $G\times_WG\times_WG$.
\end{proof}
% Lean: surj_base: the first part is RigidIsoCube.surjective_appTop_of_baseChange for the pullback square of U.gen
% (G_K = A) with the zero section and GlobalFnsSt for A; the second composes surjective_snd_of_generic for T = G and
% T = G x G (surjective_appTop_of_baseChange over T with the section (1, id), and surj_snd for the abelian scheme
% A x T_K). Injectivity on the generic fibre is IsSchemeTheoreticallyDominant.of_isPullback with app_injective.
% Size: RigidifiedIsoRespectsCubeAndSymmetry.lean, 127 lines.

\begin{lemma}[Alp\"oge]\label{cv-toroidal-torus-descent-cube}
In Situation~\ref{cv-toroidal-torus-situation}, with the notation of
Definition~\ref{cv-toroidal-torus-expansion-coefficients} and Remark~\ref{cv-toroidal-torus-descent-structures}, let
$a_1,\dots,a_r$ be endomorphisms of the group scheme $G$, also written $a_i$ for their base changes to the $G_n$,
and let $\varepsilon_1,\dots,\varepsilon_r$ be integers with $\varepsilon_1+\dots+\varepsilon_r=0$. Let $\lambda$ be
a trivialization over $G$ of $\bigotimes_ia_i^*L_G^{\varepsilon_i}$ whose value at the origin is the canonical
trivialization of $\bigotimes_iL_G(0)^{\varepsilon_i}$.
\begin{enumerate}
\item For points $x_1,x_2,x_3$ of $G$,
  $\bigotimes_I\lambda(x_I)^{\delta_I}=\bigotimes_i\theta_G(a_i(x_1),a_i(x_2),a_i(x_3))^{\varepsilon_i}$.
\item For a point $x$ of $G$, $\lambda(-x)=\lambda(x)\otimes\bigotimes_i\varsigma_G(a_i(x))^{\varepsilon_i}$.
\item Let $n\ge0$, and for a point $x$ of $G_n$ let $\lambda_n(x)$ be the image of $\lambda(x)$, the value of
  $\lambda$ at the image of $x$ in $G$, under $\bigotimes_iw_n(a_i(x))^{\otimes\varepsilon_i}$; it is a
  trivialization of $\bigotimes_iM(q_n(a_i(x)))^{\varepsilon_i}$. Then for points $x_1,x_2,x_3$ of $G_n$,
  \[
  \bigotimes_I\lambda_n(x_I)^{\delta_I}=\bigotimes_i\theta\bigl(q_n(a_i(x_1)),q_n(a_i(x_2)),q_n(a_i(x_3))\bigr)^{\varepsilon_i},
  \]
  and for a point $x$ of $G_n$, $\lambda_n(-x)=\lambda_n(x)\otimes\bigotimes_i\varsigma(q_n(a_i(x)))^{\varepsilon_i}$.
\end{enumerate}
\end{lemma}
\leanref{Development/AbelianVarietyTheory/Proofs/UniformizationProofs/DescentPowerIsoEndomorphismAbelianPartUniformiza.lean}{GlobalTheory.UT.PullDescent.gen_cube}
\leanref{Development/AbelianVarietyTheory/Proofs/UniformizationProofs/DescentPowerIsoEndomorphismAbelianPartUniformiza.lean}{GlobalTheory.UT.PullDescent.gen_even}
\leanref{Development/AbelianVarietyTheory/Proofs/UniformizationProofs/DescentPowerIsoEndomorphismAbelianPartUniformiza.lean}{GlobalTheory.UT.PullDescent.aV_cube}
\leanref{Development/AbelianVarietyTheory/Proofs/UniformizationProofs/DescentPowerIsoEndomorphismAbelianPartUniformiza.lean}{GlobalTheory.UT.PullDescent.aV_even}
\leanref{Development/AbelianVarietyTheory/Proofs/UniformizationProofs/EqTopOfThick.lean}{GlobalTheory.UT.PullDescent.prod_comm_val}
\begin{proof}
(1) As $(x_1,x_2,x_3)$ varies, the two sides are trivializations over $G^3=G\times_WG\times_WG$ of one line bundle:
its fibre is $\bigotimes_I\bigotimes_iL_G(a_i(x_I))^{\delta_I\varepsilon_i}$ for the left side, and the same for the
right side, as $a_i(x_I)=\sum_{j\in I}a_i(x_j)$. Their ratio is an invertible function on $G^3$, so it is a unit of
$W$ by Lemma~\ref{cv-toroidal-torus-descent-functions}. At $x_1=x_2=x_3=0$ both sides are the canonical
trivialization, as $\lambda(0)$ is and $\theta_G(0,0,0)=1$. So the ratio is $1$.

(2) In the same way, the two sides are trivializations over $G$ of one line bundle, with fibre
$\bigotimes_iL_G(-a_i(x))^{\varepsilon_i}$, as $a_i(-x)=-a_i(x)$ and $\varsigma_G(a_i(x))$ trivializes
$L_G(-a_i(x))\otimes L_G(a_i(x))^{-1}$. Their ratio is a unit of $W$ by
Lemma~\ref{cv-toroidal-torus-descent-functions}, and it is $1$ at $x=0$, where both sides are the canonical
trivialization, as $\lambda(0)$ and $\varsigma_G(0)$ are.

(3) Let $x_1,x_2,x_3$ be points of $G_n$, and apply (1) to their images in $G$. For each $i$, the partial sums of the
points $a_i(x_1),a_i(x_2),a_i(x_3)$ of $G_n$ are the $a_i(x_I)$, so by
Remark~\ref{cv-toroidal-torus-descent-structures}(3) the isomorphism $\bigotimes_Iw_n(a_i(x_I))^{\otimes\delta_I}$
carries $\theta_G(a_i(x_1),a_i(x_2),a_i(x_3))$ to $\theta(q_n(a_i(x_1)),q_n(a_i(x_2)),q_n(a_i(x_3)))$. So the
isomorphism $\bigotimes_I\bigotimes_iw_n(a_i(x_I))^{\otimes\delta_I\varepsilon_i}$ carries the left side of (1) to
$\bigotimes_I\lambda_n(x_I)^{\delta_I}$ and the right side of (1) to the right side of the first identity in (3).
For the second identity, let $x$ be a point of $G_n$, and apply (2) to its image in $G$. By
Remark~\ref{cv-toroidal-torus-descent-structures}(3) at the points $a_i(x)$, the isomorphism
$w_n(-a_i(x))\otimes w_n(a_i(x))^{\otimes(-1)}$ carries $\varsigma_G(a_i(x))$ to $\varsigma(q_n(a_i(x)))$. As
$a_i(-x)=-a_i(x)$, the isomorphism $\bigotimes_iw_n(a_i(-x))^{\otimes\varepsilon_i}$ carries $\lambda(-x)$ to
$\lambda_n(-x)$, and it carries $\lambda(x)\otimes\bigotimes_i\varsigma_G(a_i(x))^{\varepsilon_i}$ to
$\lambda_n(x)\otimes\bigotimes_i\varsigma(q_n(a_i(x)))^{\varepsilon_i}$.
\end{proof}
% Lean: (1) gen_cube and (2) gen_even are prod_comm_val: exists_prod for the two sides, exists_gen for their ratio
% rho, a pullback from the base by surj_base (Lemma cv-toroidal-torus-descent-functions) and 1 at the unit section by
% n_beta and val_cub_one (IsCub.face), resp. symG_rig; psum_comp and EvenPull.inv_left_comp give a_i(x_I) and
% a_i(-x). (3) aV U a beta e z J N x is lambda_n(x) read in charts; aV_cube is gen_cube with cube_val (U.w_cub) and
% prod_swap, aV_even is gen_even with even_val (U.w_even). Size: about 170 lines of
% DescentPowerIsoEndomorphismAbelianPartUniformiza.lean and 45 lines of EqTopOfThick.lean.

\begin{lemma}[Demazure--Grothendieck \cite{DG70b}]\label{cv-toroidal-torus-descent-character}
Let $S$ be a scheme, $X$ an abelian group, $T$ the diagonalizable group over $S$ with character group $X$, and $Y$ a
scheme over $S$ whose underlying space is connected. Let $\rho$ be an invertible function on $T\times_SY$ such that
$\rho(\tau\tau',y)=\rho(\tau,y)\rho(\tau',y)$ for all points $\tau,\tau'$ of $T$ and $y$ of $Y$ with values in a
scheme over $S$. Then there is a unique $\chi\in X$ such that $\rho(\tau,y)=\chi(\tau)$ for all such points.
\end{lemma}
\leanref{Development/AbelianVarietyTheory/Proofs/AbelianSchemeLawProofs/ExistsUniqueCharacterOfMul.lean}{GlobalTheory.existsUnique_character_of_mul}

The function $\rho$ is a homomorphism of group schemes over $Y$ from $T\times_SY$, the diagonalizable group over $Y$
with character group $X$, to the multiplicative group. These homomorphisms correspond to the locally constant
functions from $Y$ to $X$ \cite[Expos\'e~VIII, \S1]{DG70b}, and on the connected scheme $Y$ such a function is
constant.

\begin{lemma}[Alp\"oge]\label{cv-toroidal-torus-descent-ratio}
Let $S$ be a connected scheme, $X$ a free $\mathbf Z$-module, $T$ the split torus over $S$ with character group $X$,
$H$ a commutative group scheme over $S$ whose underlying space is connected, $i:T\to H$ a homomorphism, and $\rho$ an
invertible function on $T\times_SH$. Suppose that for every scheme $Z$ over $S$ and every point $z$ of $Z$ there is a
map $a$ from $H(Z)$ to the units of the local ring $\mathcal O_{Z,z}$ with the following properties, for all
$\tau,\tau'\in T(Z)$ and $y\in H(Z)$, where $\rho(\tau,y)$ is the germ at $z$ of the pullback of $\rho$ along
$(\tau,y):Z\to T\times_SH$.
\begin{enumerate}
\item $a(i(\tau)+y)=\rho(\tau,y)\,a(y)$.
\item $\partial a(i(\tau),i(\tau'),y)=\partial a(0,0,y)$, where for $x,x',y\in H(Z)$
  \[
  \partial a(x,x',y)=a(x+x'+y)\,a(x+x')^{-1}a(x+y)^{-1}a(x'+y)^{-1}a(x)\,a(x')\,a(y)\,a(0)^{-1}.
  \]
\item $a(-i(\tau))=a(i(\tau))$.
\end{enumerate}
Then $\rho=1$.
\end{lemma}
\leanref{Development/AbelianVarietyTheory/Proofs/UniformizationProofs/ComparisonTorusInvariant.lean}{GlobalTheory.TorusInvariance.ratio_eq_one}
\leanref{Development/AbelianVarietyTheory/Proofs/UniformizationProofs/ComparisonTorusInvariant.lean}{GlobalTheory.TorusInvariance.pullTop_eq_one}
\begin{proof}
Let $Z$ be a scheme over $S$, $z$ a point of $Z$ and $a$ as in the hypothesis, and let $\tau,\tau'\in T(Z)$ and
$y\in H(Z)$; the group law of $T$ is written multiplicatively. By (1) with $y=0$, $a(i(\tau))=\rho(\tau,0)\,a(0)$.
Three identities follow.
\begin{enumerate}
\item[(A)] $\rho(\tau\tau',0)=\rho(\tau,i(\tau'))\,\rho(\tau',0)$. Indeed, $a(i(\tau\tau'))$ is $\rho(\tau\tau',0)\,a(0)$,
  and by (1), as $i(\tau\tau')=i(\tau)+i(\tau')$, it is
  $\rho(\tau,i(\tau'))\,a(i(\tau'))=\rho(\tau,i(\tau'))\,\rho(\tau',0)\,a(0)$.
\item[(B)] $\rho(\tau\tau',y)\,\rho(\tau,0)\,\rho(\tau',0)=\rho(\tau,y)\,\rho(\tau',y)\,\rho(\tau\tau',0)$. Indeed,
  writing each of the eight values of $a$ in $\partial a(i(\tau),i(\tau'),y)$ by (1) and the formula for $a(i(\tau))$
  gives
  \[
  \partial a(i(\tau),i(\tau'),y)=\rho(\tau\tau',y)\,\rho(\tau\tau',0)^{-1}\rho(\tau,y)^{-1}\rho(\tau',y)^{-1}
  \rho(\tau,0)\,\rho(\tau',0)\,\partial a(0,0,y),
  \]
  and by (2) the product of the six values of $\rho$ is $1$.
\item[(C)] $\rho(\tau^{-1},0)=\rho(\tau,0)$. Indeed, $a(-i(\tau))=a(i(\tau^{-1}))=\rho(\tau^{-1},0)\,a(0)$, and this is
  $a(i(\tau))=\rho(\tau,0)\,a(0)$ by (3).
\end{enumerate}
These are identities between germs at $z$ of invertible functions on $Z$. As $z$ is arbitrary, they hold as
identities of functions on $Z$, for every scheme $Z$ over $S$ and all $\tau,\tau'\in T(Z)$ and $y\in H(Z)$.

Let $\rho'(\tau,y)=\rho(\tau,y)\,\rho(\tau,0)^{-1}$, an invertible function on $T\times_SH$. By (B),
$\rho'(\tau\tau',y)=\rho'(\tau,y)\,\rho'(\tau',y)$, so by Lemma~\ref{cv-toroidal-torus-descent-character} there is
$\chi\in X$ with $\rho'(\tau,y)=\chi(\tau)$. Then $\rho'(\tau,y)=\chi(\tau)=\rho'(\tau,0)=1$, that is,
$\rho(\tau,y)=\rho(\tau,0)$. With (A) this gives
$\rho(\tau\tau',y)=\rho(\tau\tau',0)=\rho(\tau,i(\tau'))\,\rho(\tau',0)=\rho(\tau,y)\,\rho(\tau',y)$, so by
Lemma~\ref{cv-toroidal-torus-descent-character} again there is $\psi\in X$ with $\rho(\tau,y)=\psi(\tau)$. The
multiplicativity gives $\rho(1,0)=\rho(1,0)^2$, so $\rho(1,0)=1$, and with (C)
$1=\rho(\tau\tau^{-1},0)=\rho(\tau,0)\,\rho(\tau^{-1},0)=\rho(\tau,0)^2=\psi(\tau)^2$ for every point $\tau$ of $T$.
So the characters $2\psi$ and $0$ are the same function on $T$. Distinct characters of $T$ are distinct functions on
$T$, as $S$ is not empty, so $2\psi=0$, and $\psi=0$ as $X$ is free. Hence $\rho=1$.
\end{proof}
% Lean: ratio_eq_one: o (the formula for a(i(tau))), then pullTop_eq_one with h1, h2 (prod_psum, the identities k
% and o, linear_combination) and h3 (i_1, i_2: 1 composed with the inverse is 1), each proved at every point z and
% assembled by unit_ext. pullTop_eq_one: exists_char (existsUnique_character_of_mul) for rho times the inverse of
% rho(., 0) (hb), giving ind; then mul from h1 and exists_char for rho; one, sq with h3, SplitTorusHom.charUnit_add,
% charUnit_injective and smul_eq_zero for psi = 0. Size: about 110 lines of ComparisonTorusInvariant.lean.

\begin{lemma}[Alp\"oge]\label{cv-toroidal-torus-descent-thickenings}
Let $W$ be a complete discrete valuation ring with maximal ideal $\mathfrak m$, and for $n\ge0$ let
$W_n=W/\mathfrak m^{n+1}$. Let $B$ be an abelian scheme over $W$ that is projective over $W$, let $B_n$ be its
reduction over $W_n$, and let $N$ be a line bundle on $B$. Suppose that $N$ has a trivialization along the zero
section of $B$, and a trivialization over $B_n$ for every $n\ge0$. Then $N$ has a trivialization over $B$.
\end{lemma}
\leanref{Development/AbelianVarietyTheory/Proofs/UniformizationProofs/EqTopOfThick.lean}{GlobalTheory.UT.PullDescent.triv_of_thick}
\begin{proof}
Let $c$ be the trivialization of $N$ along the zero section, and let $0_n$ be the zero section of $B_n$.

Normalization. Let $n\ge0$ and let $t$ be a trivialization of $N$ over $B_n$. Its value along $0_n$ and the
restriction of $c$ are trivializations of the same line bundle on $\operatorname{Spec}W_n$, so the second is $u$
times the first for a unit $u$ of $W_n$. Then $t_n=ut$ is a trivialization of $N$ over $B_n$ whose value along $0_n$
is the restriction of $c$.

Compatibility. The restriction of $t_{n+1}$ to $B_n$ and $t_n$ are trivializations of $N$ over $B_n$ with the same
value along $0_n$. Their ratio is an invertible function on $B_n$, which is a unit of $W_n$, as the global functions
on the abelian scheme $B_n$ are those of $W_n$, and it is $1$ along $0_n$. So $t_{n+1}$ restricts to $t_n$.

Algebraization. As $B$ is proper over $W$, the comparison theorem \cite[Th\'eor\`eme~4.1.5]{Gro61} shows that
restriction is a bijection from the sections of $N$ over $B$ to the inverse limit of the sections of $N$ over the
$B_n$, since the sections over $B$ form a finite module over the complete ring $W$. So there is a section $s$ of $N$
over $B$ whose restriction to every $B_n$ is $t_n$.

Generator. The restriction of $s$ to $B_0$ is $t_0$, which generates $N$ at every point of $B_0$. So the open set $U$
of $B$ where $s$ generates $N$ contains $B_0$. The image in $\operatorname{Spec}W$ of the complement of $U$ is closed,
as $B$ is proper over $W$, and it does not contain the closed point. Every nonempty closed subset of
$\operatorname{Spec}W$ contains the closed point, so $U=B$, and $s$ is a trivialization of $N$ over $B$.
\end{proof}
% Lean: triv_of_thick: AbTriv.exists_normalised (t times the pullback of c over the base, divided by the pullback of
% t along the zero section), hcompat (hratio, then rho = 1 by GlobalFnsSt for the abelian scheme thickObj and the
% section sigma), SecAlg.existsUnique_sec, eq_top_of_thick (universally closed, I in the Jacobson radical) for the
% nonzero locus of s, and AbTriv.triv_of_pull_id. existsUnique_sec is the Lean's own proof of the algebraization, for
% B smooth and projective: a finite complex of free modules computing the sections after every base change
% (T.nonempty_perfectSectionsComplex_of_smooth_proj) and the I-adic limit of the coefficient vectors (exists_lim).
% Size: about 130 lines of EqTopOfThick.lean, ExistsUniqueSec.lean 161 lines, EqTopOfClosedFibre.lean 118 lines.

\begin{lemma}[Alp\"oge]\label{cv-toroidal-torus-expansion-descent}
In Situation~\ref{cv-toroidal-torus-situation}, with the notation of
Definition~\ref{cv-toroidal-torus-expansion-coefficients}, let $E=(a_G,a_A,a_B,h,u)$ be an endomorphism of the
uniformization (Definition~\ref{cv-toroidal-torus-expansion-endomorphism}) and let $k,k'\ge1$. Write $L_G(0)$ and
$M(0)$ also for the constant line bundles on $G$ and $B$ with the fibres of $L_G$ and $M$ at the origin. Let
$\lambda$ be a trivialization over $G$ of $L_G^{k'}\otimes a_G^*L_G^{-k}\otimes L_G(0)^{k-k'}$ whose value at the
origin is the canonical trivialization. Then there is a trivialization $\kappa$ over $B$ of
$M^{k'}\otimes a_B^*M^{-k}\otimes M(0)^{k-k'}$ whose value at the origin is the canonical trivialization, such that
for every $n\ge0$ and every point $x$ of $G_n$ the isomorphism
$w_n(x)^{\otimes k'}\otimes w_n(a_G(x))^{\otimes(-k)}\otimes w_n(0)^{\otimes(k-k')}$ carries $\lambda(x)$ to
$\kappa(q_n(x))$.
\end{lemma}
\leanref{Development/AbelianVarietyTheory/Proofs/UniformizationProofs/DescentPowerIsoEndomorphismAbelianPartUniformizaStatement.lean}{GlobalTheory.T.descent_powerIso_endomorphism_abelianPart_uniformization}
\leanref{Development/AbelianVarietyTheory/Proofs/UniformizationProofs/DescentPowerIsoEndomorphismAbelianPartUniformiza.lean}{GlobalTheory.UT.PullDescent.main}
\begin{proof}
Let $N=M^{k'}\otimes a_B^*M^{-k}\otimes M(0)^{k-k'}$, and let $n\ge0$. As $q_n\circ a_G=a_B\circ q_n$, the
isomorphisms $w_n$ carry the restriction of $\lambda$ to $G_n$ to a trivialization $\lambda_n$ of $q_n^*N$ over
$G_n$. It is canonical at the origin, since $\lambda$ is and $w_n$ carries $e_G$ to $e$. Put
$(a_1,a_2,a_3)=(1,a_G,0)$, $(b_1,b_2,b_3)=(1,a_B,0)$ and $(\varepsilon_1,\varepsilon_2,\varepsilon_3)=(k',-k,k-k')$.
Then $\lambda$ trivializes $\bigotimes_ia_i^*L_G^{\varepsilon_i}$, $N=\bigotimes_ib_i^*M^{\varepsilon_i}$,
$q_n\circ a_i=b_i\circ q_n$ for every $i$, $\sum_i\varepsilon_i=0$, and $\lambda_n(x)$ is the image of $\lambda(x)$
under $\bigotimes_iw_n(a_i(x))^{\otimes\varepsilon_i}$, as in Lemma~\ref{cv-toroidal-torus-descent-cube}(3).

Invariance under $T_n$. For points $\tau$ of $T_n$ and $y$ of $G_n$, $q_n(\tau+y)=q_n(y)$, as $T_n$ is the kernel of
$q_n$. So the pullbacks of $\lambda_n$ along the maps $(\tau,y)\mapsto\tau+y$ and $(\tau,y)\mapsto y$ from
$T_n\times_{W_n}G_n$ to $G_n$ trivialize the same line bundle, and their ratio is an invertible function $\rho$ on
$T_n\times_{W_n}G_n$ with $\lambda_n(\tau+y)=\rho(\tau,y)\lambda_n(y)$. Lemma~\ref{cv-toroidal-torus-descent-ratio}
applies to $S=\operatorname{Spec}W_n$, $H=G_n$, the inclusion $i$ of $T_n$, and $\rho$. Indeed,
$\operatorname{Spec}W_n$ is a point, $G_n$ is connected as the fibres of $G$ are connected, and $X$ is free. For the
maps $a$, let $Z$ be a scheme over $W_n$ and $z$ a point of $Z$, and choose, for every point $b$ of $B$ with values in
$Z$, a generator of $M$ near the image of $z$ under $b$. For a point $x$ of $G_n$ with values in $Z$, $\lambda_n(x)$
trivializes $\bigotimes_iM(b_i(q_n(x)))^{\varepsilon_i}$; written in the generators chosen at the points
$b_i(q_n(x))$, its germ at $z$ is a unit $a(x)$ of $\mathcal O_{Z,z}$. Points of $G_n$ with the same image under
$q_n$ get the same generators, so $a(\tau+y)=\rho(\tau,y)\,a(y)$, which is condition (1). By
Lemma~\ref{cv-toroidal-torus-descent-cube}(3), $\partial a(x_1,x_2,x_3)$ is the product over $i$ of the
$\varepsilon_i$-th powers of the values of $\theta$ at $(b_i(q_n(x_1)),b_i(q_n(x_2)),b_i(q_n(x_3)))$, written in the
chosen generators. For $(x_1,x_2,x_3)$ equal to $(\tau,\tau',y)$ and to $(0,0,y)$ these points are the same, as
$q_n(\tau)=q_n(\tau')=0$, which gives condition (2). In the same way $a(-x)=a(x)\,v(x)$, where $v(x)$ is the product
of the $\varepsilon_i$-th powers of the values of $\varsigma$ at the points $b_i(q_n(x))$, and $v(\tau)=v(0)$. As
$-0=0$, $v(0)=1$, so $a(-\tau)=a(\tau)$, which is condition (3). Hence $\rho=1$, that is,
$\lambda_n(\tau+y)=\lambda_n(y)$.

Descent to $B_n$. Let $x$ and $x'$ be points of $G_n$ with $q_n(x)=q_n(x')$. Then $x-x'$ is a point of the kernel
$T_n$ of $q_n$, so $\lambda_n(x)=\lambda_n((x-x')+x')=\lambda_n(x')$. Thus the pullbacks of $\lambda_n$ along the two
projections from $G_n\times_{B_n}G_n$ to $G_n$ agree. As $q_n$ is flat, surjective and affine
(Remark~\ref{cv-toroidal-torus-descent-structures}(3)), faithfully flat descent
\cite[Expos\'e~VIII, Th\'eor\`eme~1.1]{Gro71} gives a trivialization $s_n$ of $N$ over $B_n$ with
$s_n(q_n(x))=\lambda_n(x)$ for every point $x$ of $G_n$.

A trivialization over $B$. As $b_i(0)=0$ for every $i$ and $\sum_i\varepsilon_i=0$, the fibre of $N$ at the origin is
$\bigotimes_iM(0)^{\varepsilon_i}$ with its canonical trivialization, so $N$ has a trivialization along the zero
section. As $M$ is ample, $B$ is projective over $W$. So Lemma~\ref{cv-toroidal-torus-descent-thickenings}, applied
to $N$ with this trivialization and the trivializations $s_n$, gives a trivialization $\kappa_0$ of $N$ over $B$.
% Lean: PullDescent.exists_level for each n: exists_transport (U.w_gen) for lambda_n, with val lambda_n = aV (hv');
% PullDescent.exists_ratio for rho, with the maps (fst i) * snd and snd from T_n x G_n (h1, isKernel.w);
% TorusInvariance.ratio_eq_one (Lemma cv-toroidal-torus-descent-ratio) with a = aV U a beta e z J N for charts J and N
% chosen by exists_chart at every morphism, its three conditions from key, from aV_cube with tri_congr (h2), and from
% aV_even with pG_congr; connectedness from EndTorus.connectedSpace_thickening and
% GroupConn.connectedSpace_of_section (hS.geomConnected); exists_descend along q (q_flat, q_surj, q_affine), chart
% by chart through LevelComparison.exists_unit_descend (FunDescent.exists_app_eq, the Lean's own descent of
% functions), with k_1 k_2^-1 lifted to the torus by isKernel.lift. PullDescent.exists_descent: t_0 by exists_gen
% from oneGen with h_eps, proj from Unif.relAmple_iff_veryAmpleBy and P.ample, kappa_0 by triv_of_thick (Lemma
% cv-toroidal-torus-descent-thickenings) with triv_pull of t_0 and of the s_e. In the Lean, 1 : G -> G is the
% constant map to the origin. Size: DescentPowerIsoEndomorphismAbelianPartUniformiza.lean 388 lines,
% ComparisonTorusInvariant.lean 470, EqTopOfThick.lean 415.
Let $\kappa$ be $\kappa_0$ divided by its value at the origin, a trivialization of $N$ that is canonical at the
origin. For $n\ge0$, the restriction of $\kappa$ to $B_n$ and $s_n$ are trivializations of $N$ over $B_n$ that
are canonical at the origin, as $s_n(0)=\lambda_n(0)$. Their ratio is an invertible function on $B_n$, which is a
unit of $W_n$ since the global functions on $B_n$ are those of $W_n$, and it is $1$ at the origin. So
$\kappa(q_n(x))=s_n(q_n(x))=\lambda_n(x)$ for every point $x$ of $G_n$.
\end{proof}
% Lean: exists_descent, last part: kappa = kappa_0 times the inverse of its pullback along the constant map 1 (nk),
% then gen_unique on each thickening (with the values at the origin from vs, a_0 and q_0); main specializes to
% kapLB with the maps 1, E.aB, 1 and the exponents k', -k, k - k'.

\begin{lemma}[Alp\"oge]\label{cv-toroidal-torus-expansion-integral}
In Situation~\ref{cv-toroidal-torus-situation}, with the notation of
Definition~\ref{cv-toroidal-torus-expansion-coefficients}, let $E=(a_G,a_A,a_B,h,u)$ be an endomorphism of the
uniformization (Definition~\ref{cv-toroidal-torus-expansion-endomorphism}) with $h$ injective, and let $k,k'\ge1$.
Let $\lambda$ be a trivialization over $G$ of $L_G^{k'}\otimes a_G^*L_G^{-k}\otimes L_G(0)^{k-k'}$ and $\kappa$ a
trivialization over $B$ of $M^{k'}\otimes a_B^*M^{-k}\otimes M(0)^{k-k'}$, where $L_G(0)$ and $M(0)$ are the
constant line bundles with the fibres at the origin, such that for every $n\ge0$ and every point $x$ of $G_n$ the
isomorphism $w_n(x)^{\otimes k'}\otimes w_n(a_G(x))^{\otimes(-k)}\otimes w_n(0)^{\otimes(k-k')}$ carries
$\lambda(x)$ to $\kappa(q_n(x))$. Trivialize $L_G(0)$ and $M(0)$ by $e_G$ and $e$. Let $\sigma$ be a section of
$L_G^k$ over $G$ with a family $f$ of Fourier coefficients. Then the section $\sigma'=a_G^*\sigma\otimes\lambda$ of
$L_G^{k'}$ over $G$ has the family $f'$ of Fourier coefficients given by
$f'_{h(\chi)}=a_B^*f_\chi\otimes\kappa\otimes u(\chi)$ for $\chi\in X$ and $f'_\beta=0$ for $\beta\notin h(X)$.
\end{lemma}
\leanref{Development/AbelianVarietyTheory/Proofs/UniformizationProofs/PullbackSectionEndomorphismExpansionStatement.lean}{GlobalTheory.T.pullback_section_endomorphism_expansion}
\leanref{Development/AbelianVarietyTheory/Proofs/UniformizationProofs/PullbackSectionEndomorphismExpansion.lean}{GlobalTheory.PullInt.main}
\begin{proof}
As $h$ is injective, $f'$ is well defined. Let $n\ge0$ and let $S$ be the finite set given for $f$ and $n$. If
$\beta\notin h(S)$, then $f'_\beta$ vanishes on $B_n$: it is $0$ if $\beta\notin h(X)$, and if $\beta=h(\chi)$ with
$\chi\notin S$, then $f_\chi$ vanishes on $B_n$, and so does $a_B^*f_\chi$. Let $x$ be a point of $G_n$. The
expansion of $\sigma$ at the point $a_G(x)$, the identity $q_n\circ a_G=a_B\circ q_n$ and the last condition of
Definition~\ref{cv-toroidal-torus-expansion-endomorphism} give
\[
w_n(a_G(x))^{\otimes k}\bigl(\sigma(a_G(x))\bigr)=\sum_{\chi\in S}f_\chi(a_B(q_n(x)))\otimes t_{n,h(\chi)}(x)\otimes
u(\chi)(q_n(x)).
\]
By hypothesis, $w_n(x)^{\otimes k'}\otimes w_n(a_G(x))^{\otimes(-k)}\otimes w_n(0)^{\otimes(k-k')}$ carries
$\lambda(x)$ to $\kappa(q_n(x))$, and $w_n(0)$ carries $e_G$ to $e$. Tensoring the identity with these gives
\[
w_n(x)^{\otimes k'}\bigl(\sigma'(x)\bigr)=\sum_{\chi\in S}f'_{h(\chi)}(q_n(x))\otimes t_{n,h(\chi)}(x)
=\sum_{\beta\in h(S)}f'_\beta(q_n(x))\otimes t_{n,\beta}(x),
\]
the last step because $h$ is injective. So $h(S)$ is a finite set as in
Definition~\ref{cv-toroidal-torus-expansion-coefficients} for $f'$ and $n$.
\end{proof}
% Lean: PullInt.main: exists_secG for sigma' (pullMulSec with lambda and rig^(k'-k), exists_zpow), exists_secB for
% f'_{h chi} (with kappa, rig^(k'-k) and E.u chi), the zero section off the range of h; ExpandsAt through
% expandsAt_val and expandsAt_of_val with the set S.image E.h: the value at a_G x (act_pGl, act_fGl from E.q_lin),
% the descent relation H at x, rig_val (w_rig: w_n(0) carries e_G to e), E.t_lin, and the algebra alg. Size:
% PullbackSectionEndomorphismExpansion.lean, 212 lines.

\begin{lemma}[Alp\"oge]\label{cv-toroidal-torus-expansion-twist}
In Situation~\ref{cv-toroidal-torus-situation}, let $E=(a_G,a_A,a_B,h,u)$ be an endomorphism of the uniformization
(Definition~\ref{cv-toroidal-torus-expansion-endomorphism}), let $k,k'\ge1$, and let $\lambda$ be a trivialization
over $G$ of $L_G^{k'}\otimes a_G^*L_G^{-k}\otimes L_G(0)^{k-k'}$, where $L_G(0)$ is the constant line bundle with
the fibre of $L_G$ at the origin, whose value at the origin is the canonical trivialization. Trivialize $L_G(0)$ by
the trivialization $e_G$ of $L_G$ at the origin. Let $\rho$ be a trivialization of $L$ at the origin, $w_0$ a
normalization for $\rho$ (Definition~\ref{cv-toroidal-torus-expansion}(1)), and $\alpha$ a trivialization over
$A_F$ of $L^{k'}\otimes a_A^*L^{-k}$ whose value at the origin is $\rho^{k'-k}$. Let $t$ be a section of $L^k$ over
$A_F$, let $t'=\alpha\otimes a_A^*t$, and let $v$ and $v'$ be the sections of $L_G^k$ and $L_G^{k'}$ over $G_F$
that correspond to $t$ and $t'$ through $w_0$ and $\gamma$. Then $v'=a_G^*v\otimes\lambda$.
\end{lemma}
\leanref{Development/AbelianVarietyTheory/Proofs/UniformizationProofs/TwistThroughIsoStatement.lean}{GlobalTheory.T.endU_twist_of_iso}
\leanref{Development/AbelianVarietyTheory/Proofs/UniformizationProofs/TwistThroughIso.lean}{GlobalTheory.UT.TwistCompare.twist_through_iso}
\begin{proof}
Through $e_G$, $\lambda$ is a trivialization $\lambda'$ of $L_G^{k'}\otimes a_G^*L_G^{-k}$ over $G$, whose value at
the origin is $e_G^{k'-k}$. As $\gamma\circ a_G=a_A\circ\gamma$, the isomorphisms $\gamma$ and $w_0$ carry $\alpha$
to a trivialization $\alpha'$ of the same line bundle over $G_F$, the image of $\gamma^*\alpha$ under
$w_0^{\otimes k'}\otimes(a_G^*w_0)^{\otimes(-k)}$. Its value at the origin is the image of $\rho^{k'-k}$ under
$w_0^{\otimes(k'-k)}$, which is $e_G^{k'-k}$, as $w_0$ carries $\rho$ to $e_G$. So $\alpha'/\lambda'$ is an
invertible function on $G_F$ that is $1$ at the origin. Through $\gamma$, $G_F$ is the abelian variety $A_F$, whose
global functions are those of $F$, so $\alpha'=\lambda'$. Hence
\[
v'=w_0^{\otimes k'}(\gamma^*t')=\alpha'\otimes a_G^*\bigl(w_0^{\otimes k}(\gamma^*t)\bigr)=\lambda'\otimes a_G^*v,
\]
which is $a_G^*v\otimes\lambda$.
\end{proof}
% Lean: TwistCompare.twist_through_iso: exists_two for lambda' (lambda times rig^(k'-k)); val_through_iso
% (exists_unpull, H1.exists_gen_move, gen_unique over T_F x A with the origin ptSecT 1 and surj_fst), with the
% values at the origin from the normalizations of lambda and alpha and EvalZero.normW_ext (alg_origin); then
% isoRel_val for t and t', twist_val for t' and dualNumber_unit_mul_decomposition. Size: TwistThroughIso.lean,
% 188 lines.

\begin{lemma}[Alp\"oge]\label{cv-toroidal-torus-expansion-pullback}
In Situation~\ref{cv-toroidal-torus-situation}, with the notation of
Definition~\ref{cv-toroidal-torus-expansion-coefficients}, let $E=(a_G,a_A,a_B,h,u)$ be an endomorphism of the
uniformization (Definition~\ref{cv-toroidal-torus-expansion-endomorphism}) with $h$ injective, and let
$k,k'\ge1$. Suppose that $L_G^{k'}\otimes a_G^*L_G^{-k}$ is trivial over $G$.
\begin{enumerate}
\item There is a trivialization of $M^{k'}\otimes a_B^*M^{-k}\otimes M(0)^{k-k'}$ over $B$ whose value at the
  origin is the canonical trivialization, where $M(0)$ is the constant line bundle with the fibre of $M$ at the
  origin.
\item Let $\rho$ be a trivialization of $L$ at the origin, $w_0$ a normalization for $\rho$
  (Definition~\ref{cv-toroidal-torus-expansion}(1)), and $\alpha$ a trivialization over $A_F$ of
  $L^{k'}\otimes a_A^*L^{-k}$ whose value at the origin is $\rho^{k'-k}$. Let $t$ be a section of $L^k$ over $A_F$,
  let $t'=\alpha\otimes a_A^*t$, and let $\sigma$ and $\sigma'$ be the Fourier expansions of $t$ and $t'$ with
  respect to $w_0$ (Lemma~\ref{cv-toroidal-torus-bijection}). Let $\kappa$ be a trivialization as in (1), let
  $g=(b,s)$ and $g'=(b',s')$ be points of $\tilde G(F)$ such that $E$ carries $g$ to $g'$, and let $\omega'$ be a
  frame of weight $k'$ at $b$ and $\omega$ a frame of weight $k$ at $b'$
  (Definition~\ref{cv-toroidal-torus-terms}(2)) with $\omega'\otimes\kappa(b)=\omega$. Then
  $a_{h(\chi)}(\sigma',g,\omega')=a_\chi(\sigma,g',\omega)$ for every $\chi\in X$, and $a_\beta(\sigma',g,\omega')=0$
  for every $\beta\in X$ outside $h(X)$ (Definition~\ref{cv-toroidal-torus-terms}(3)).
\end{enumerate}
\end{lemma}
\leanref{Development/AbelianVarietyTheory/Proofs/UniformizationProofs/ExpansionPullbackSectionEndomorphismStatement.lean}{GlobalTheory.T.expansion_pullback_section_endomorphism}
\leanref{Development/AbelianVarietyTheory/Proofs/UniformizationProofs/ExpansionPullbackSectionEndomorphism.lean}{GlobalTheory.UT.PullExp.main}
\begin{proof}
Let $\lambda_0$ be a trivialization of $L_G^{k'}\otimes a_G^*L_G^{-k}$ over $G$, and let $\lambda$ be $\lambda_0$
divided by its value at the origin, a trivialization of $L_G^{k'}\otimes a_G^*L_G^{-k}\otimes L_G(0)^{k-k'}$ that
is canonical at the origin.

(1) Lemma~\ref{cv-toroidal-torus-expansion-descent} for $\lambda$ gives a trivialization $\kappa_0$ that is
canonical at the origin and related to $\lambda$ as there.

(2) The trivializations $\kappa$ and $\kappa_0$ are both canonical at the origin. Their ratio is an invertible
function on $B$, a unit of $W$, and it is $1$ at the origin, so $\kappa=\kappa_0$. By
Definition~\ref{cv-toroidal-torus-expansion}(2) there are sections $\sigma_1,\dots,\sigma_r$ of $L_G^k$ over $G$
with families $f^{(1)},\dots,f^{(r)}$ of Fourier coefficients
(Definition~\ref{cv-toroidal-torus-expansion-coefficients}) and $c_1,\dots,c_r\in F$, such that the section $v$
of $L_G^k$ over $G_F$ that corresponds to $t$ through $w_0$ and $\gamma$ is $\sum_ic_i\sigma_i$ and
$\sigma=\sum_ic_if^{(i)}$. Lemma~\ref{cv-toroidal-torus-expansion-integral}, for $\lambda$ and $\kappa$, gives the
families $f'^{(i)}$ of Fourier coefficients of the sections $\sigma'_i=a_G^*\sigma_i\otimes\lambda$ of $L_G^{k'}$
over $G$, with $f'^{(i)}_{h(\chi)}=a_B^*f^{(i)}_\chi\otimes\kappa\otimes u(\chi)$ and $f'^{(i)}_\beta=0$ for
$\beta\notin h(X)$. By Lemma~\ref{cv-toroidal-torus-expansion-twist}, the section of $L_G^{k'}$ over $G_F$ that
corresponds to $t'$ is $a_G^*v\otimes\lambda=\sum_ic_i\sigma'_i$. This section is also the combination of sections
over $G$ given by Definition~\ref{cv-toroidal-torus-expansion}(2) for $\sigma'$. The argument of the uniqueness
part of the proof of Lemma~\ref{cv-toroidal-torus-bijection}, applied to each index $\beta\in X$ and to these two
combinations, gives $\sigma'_\beta=\sum_ic_if'^{(i)}_\beta$ for every $\beta$.

Let $\chi\in X$, and trivialize $M(0)$ by $e$. As $E$ carries $g$ to $g'$, $a_B(b)=b'$ and
$s(h(\chi))\otimes u(\chi)(b)=s'(\chi)$; and $\omega'\otimes\kappa(b)=\omega$. So
\[
\sigma'_{h(\chi)}(b)\otimes s(h(\chi))\otimes\omega'
=\sum_ic_if^{(i)}_\chi(b')\otimes\kappa(b)\otimes u(\chi)(b)\otimes s(h(\chi))\otimes\omega'
=\sigma_\chi(b')\otimes s'(\chi)\otimes\omega.
\]
By Definition~\ref{cv-toroidal-torus-terms}(3) the left side is $a_{h(\chi)}(\sigma',g,\omega')e^{k'}$ and the right
side is $a_\chi(\sigma,g',\omega)e^k$, which is $a_\chi(\sigma,g',\omega)e^{k'}$ once $M(0)^{k'-k}$ is trivialized
by $e$. If $\beta\notin h(X)$, then $\sigma'_\beta=\sum_ic_if'^{(i)}_\beta=0$, so $a_\beta(\sigma',g,\omega')=0$.
\end{proof}
% Lean: PullExp.main: exists_normalised for lambda; T.descent_powerIso_endomorphism_abelianPart_uniformization for
% kappa_0 (part (1)); normd_val_eq for kappa = kappa_0; ExpandsF unfolded (v, n, c, sigma, f);
% T.pullback_section_endomorphism_expansion for each sigma_i; T.endU_twist_of_iso for v'; secLin_of_expandsG
% (comb_eq_zero with term_eq_zero, the argument of the uniqueness half of graph_bijective in Lemma
% cv-toroidal-torus-bijection) for sigma' = sum c_i f'_i; then CoefBasic.coef_eq with sval_comb, isCoef_coef,
% E.Act and the relation of the frames (the factor rig^(k'-k) is e^(k'-k)). Size:
% ExpansionPullbackSectionEndomorphism.lean, 221 lines.

\begin{lemma}\label{cv-toroidal-torus-expansion-phi-map}
In Situation~\ref{cv-toroidal-torus-situation}, let $k\ge1$, let $p=(y,g)$ be a $k$-division point with $g=(b,s)$,
and let $w$ be a frame of weight $k$ at $b$ (Definition~\ref{cv-toroidal-torus-terms}(2)).
\begin{enumerate}
\item For every family $\sigma$ of Fourier coefficients of weight $k$ over $F$ there is exactly one family
  $\sigma'$ with $\sigma'=\Phi_{p,w}\sigma$ (Definition~\ref{cv-toroidal-torus-theta-sum-phi}), and the map
  $\Phi_{p,w}:\sigma\mapsto\sigma'$ is $F$-linear.
\item Let $w'$ be a second frame of weight $k$ at $b$. Then $w'=cw$ for some $c\in F^\times$, and
  $\Phi_{p,w'}\sigma=c\,\Phi_{p,w}\sigma$ for every $\sigma$.
\end{enumerate}
\end{lemma}
\leanref{ReciprocityTheory/ReductiveGroups/ArithmeticGeometry/DiscreteValuationRings/PhiGenPf.lean}{ReciprocityTheory.ChartRestriction.G.DVR.phiFun_of}
\leanref{ReciprocityTheory/ReductiveGroups/ArithmeticGeometry/DiscreteValuationRings/PhiGenPf.lean}{ReciprocityTheory.ChartRestriction.G.DVR.units_comm}
\leanref{Development/AbelianVarietyTheory/Proofs/UniformizationProofs/SymmetricLiftOddOrderActsAsExpansionOperator.lean}{GlobalTheory.UT.SymLiftPhi.exists_ratio}
\begin{proof}
(1) For $\alpha\in X$ let $u_\alpha$ be the isomorphism of Lemma~\ref{cv-toroidal-torus-theta-sum-operator}(1),
let $T_b$ be the translation by $b$ on $B_F$, and for $\mu\in X$ let $T_\mu$ be the translation by $c(\mu)$. By
Definition~\ref{cv-toroidal-torus-theta-sum-phi}, $\sigma'=\Phi_{p,w}\sigma$ says that
$\sigma'_{\alpha+y}=u_\alpha(T_b^*\sigma_\alpha)$ for every $\alpha$. Every element of $X$ is $\alpha+y$ for
exactly one $\alpha$, so these formulas define the only candidate $\sigma'$, and it remains to see that it
satisfies the functional equation of Definition~\ref{cv-toroidal-torus-points}(3). With $u_{\mu,\chi}$ as in
Lemma~\ref{cv-toroidal-torus-translation}(1) over $F$, $\sigma_{\alpha+k\mu}=u_{\mu,\alpha}(T_\mu^*\sigma_\alpha)$,
and the functional equation asks that $\sigma'_{\alpha+y+k\mu}=u_{\mu,\alpha+y}(T_\mu^*\sigma'_{\alpha+y})$. As $T_b$
and $T_\mu$ commute, it follows from the identity
\[
u_{\alpha+k\mu}\circ T_b^*u_{\mu,\alpha}=u_{\mu,\alpha+y}\circ T_\mu^*u_\alpha
\]
of isomorphisms from $T_b^*T_\mu^*M_{k,\alpha}$ to $M_{k,\alpha+y+k\mu}$. By the proofs of
Lemmas~\ref{cv-toroidal-torus-translation}(1) and~\ref{cv-toroidal-torus-theta-sum-operator}(1), the values of both
sides at a point of $B_F$ are products of values of the cubical structure of $M$ with $s(\alpha+k\mu)\otimes w$
and $\psi(\mu)^k\tau(\mu,\alpha)$ on the left, and with $\psi(\mu)^k\tau(\mu,\alpha+y)$ and $s(\alpha)\otimes w$ on
the right. They agree by the relation between $s(\alpha+k\mu)$, $s(\alpha)$ and $s(\mu)$ of
Definition~\ref{cv-toroidal-torus-points}(1), the relation $kg=\iota(y)$, the bimultiplicativity of $\tau$ and an
identity of the cubical structure. The formula for $\sigma'$ is $F$-linear in $\sigma$.

(2) The frames $w$ and $w'$ trivialize the same line over $F$, so $w'=cw$ with $c\in F^\times$. The isomorphism
$T_b^*M_{k,\alpha}\to M_{k,\alpha+y}$ whose value at the origin is given by $s(\alpha)\otimes w'$ is then
$c\,u_\alpha$, so $\Phi_{p,w'}\sigma=c\,\Phi_{p,w}\sigma$.
\end{proof}
% Lean: phiFun_of = phiFun_of_globalFns: exists_phi (gen_phi for u_alpha, exists_sec with pullMulSec) and phi_fe
% (gen_fe for the functional equation, with units_comm: gt_nsmul, tau_mul, chain_val and IsCub.comm), phi_unique,
% sval_add and sval_smul for linearity, phiRel_iff. (2) is SymLiftPhi.exists_ratio (H1.exists_const for c, phiRel_iff).

\begin{lemma}[Alp\"oge]\label{cv-toroidal-torus-expansion-phi-order}
In Situation~\ref{cv-toroidal-torus-situation}, let $N$ be odd, let $p=(y,g)$ be an $N$-division point with
$g=(b,s)$, and let $w$ be a symmetric frame for $p$ (Definition~\ref{cv-toroidal-torus-symmetric}(2)). Then
$\Phi_{p,w}^N\sigma=\sigma$ for every family $\sigma$ of Fourier coefficients of weight $N$ over $F$
(Lemma~\ref{cv-toroidal-torus-expansion-phi-map}).
\end{lemma}
\leanref{Development/AbelianVarietyTheory/Proofs/UniformizationProofs/SymmetricFramePhiPowOrderEqIdStatement.lean}{GlobalTheory.T.symmetricFrame_phi_pow_order_eq_id}
\leanref{Development/AbelianVarietyTheory/Proofs/UniformizationProofs/SymmetricFramePhiPowOrderEqId.lean}{GlobalTheory.UT.PhiOrder.main}
\leanref{ReciprocityTheory/ReductiveGroups/ArithmeticGeometry/FirstCohomology/PhiPow.lean}{ReciprocityTheory.ChartRestriction.G.H1.const_at_zero}
\leanref{Development/AbelianVarietyTheory/Proofs/UniformizationProofs/SymmetricFramePhiPowOrderEqId.lean}{GlobalTheory.UT.PhiOrder.const_one}
\begin{proof}
Let $\sigma^{(0)}=\sigma$ and $\sigma^{(i+1)}=\Phi_{p,w}\sigma^{(i)}$. Let $\alpha\in X$, let $u_\beta$, for
$\beta\in X$, be the isomorphisms of Lemma~\ref{cv-toroidal-torus-theta-sum-operator}(1) for the frame $w$, and for
$x\in B(F)$ let $T_x$ be the translation by $x$. By Definition~\ref{cv-toroidal-torus-theta-sum-phi},
$\sigma^{(i+1)}_{\alpha+(i+1)y}=u_{\alpha+iy}(T_b^*\sigma^{(i)}_{\alpha+iy})$, so by induction on $i$,
$\sigma^{(N)}_{\alpha+Ny}=U_\alpha(T_{Nb}^*\sigma_\alpha)$, where $U_\alpha$ is the composite of the isomorphisms
$T_{(N-1-i)b}^*u_{\alpha+iy}$ for $i=0,\dots,N-1$, from $T_{Nb}^*M_{N,\alpha}$ to $M_{N,\alpha+Ny}$. As
$Ng=\iota(y)$, $Nb=c(y)$, and Lemma~\ref{cv-toroidal-torus-translation}(1), over $F$ and with $k=N$, gives
$\sigma_{\alpha+Ny}=u_{y,\alpha}(T_{Nb}^*\sigma_\alpha)$. So $U_\alpha$ and $u_{y,\alpha}$ are isomorphisms between
the same line bundles on $B_F$, and $U_\alpha=a(\alpha)u_{y,\alpha}$ for an invertible function $a(\alpha)$ on
$B_F$, which is a constant, as the global functions on $B_F$ are those of $F$.

At the origin, the value of $U_\alpha$ is the product of the values of the $u_{\alpha+iy}$ at the points
$(N-1-i)b$, which the proof of Lemma~\ref{cv-toroidal-torus-theta-sum-operator}(1) gives through the cubical
structure of $M$ and $s(\alpha+iy)\otimes w$; the value of $u_{y,\alpha}$ is given by $\psi(y)^N\tau(y,\alpha)$. A
computation with the cubical structure shows that $a(\alpha)=1$. It uses the relations among the $s(iy)$ and the
$s(\alpha+iy)$ from Definition~\ref{cv-toroidal-torus-points}(1), $Ng=\iota(y)$, the bimultiplicativity of $\tau$,
the evenness of $\psi$, and the condition $w^{\otimes N}\otimes\kappa_N(b)=\psi(y)\otimes e$ of
Definition~\ref{cv-toroidal-torus-symmetric}(2), with $\kappa_N(b)$ written through the cubical structure. Hence
$\sigma^{(N)}_{\alpha+Ny}=\sigma_{\alpha+Ny}$. Every element of $X$ is $\alpha+Ny$ for some $\alpha$, so
$\sigma^{(N)}=\sigma$.
\end{proof}
% Lean: PhiOrder.phiPow_eq_one: gen_phi for the u, gen_fe for v = u_{y,alpha}, H1.exists_gen_prod and
% H1.exists_const (surj_appTop) for the constant a(alpha), with H1.tel_formal; H1.const_at_zero for its value at the
% origin and const_one for a(alpha) = 1 (exists_kap, the formula for kappa_N through pot, gt_nsmul, chain_val,
% tau_diag from TauBimult and PsiEven, and const_eq_one); H1.iter for the iterate, then fam_ext. Only the second
% condition of TXFrame is used. Size: SymmetricFramePhiPowOrderEqId.lean, 387 lines, and PhiPow.lean, 254 lines.

\begin{definition}\label{cv-toroidal-torus-expansion-assigned}
In Situation~\ref{cv-toroidal-torus-situation}, let $k\ge1$, let $p=(y,g)$ be a $k$-division point with $g=(b,s)$,
and let $x\in A(F)$. The uniformization \emph{assigns} $x$ to $p$ if $kx=0$ and, for every isomorphism $w_0$ from
$\gamma^*L$ to the restriction of $L_G$ to $G_K$ and every element $U$ of the theta group of $L^k$ over $A_F$ lying
over $x$ \cite{Mum66}, there is a frame $w$ of weight $k$ at $b$ (Definition~\ref{cv-toroidal-torus-terms}(2)) with
the following property: for all sections $t$ and $t'$ of $L^k$ over $A_F$ with $t'=Ut$, the Fourier expansions
$\sigma$ and $\sigma'$ of $t$ and $t'$ with respect to $w_0$ (Definition~\ref{cv-toroidal-torus-expansion}(2))
satisfy $\sigma'=\Phi_{p,w}\sigma$ (Definition~\ref{cv-toroidal-torus-theta-sum-phi}).
\end{definition}
\leanref{ReciprocityTheory/ReductiveGroups/ArithmeticGeometry/ExtensionDescent/QkToAkU.lean}{ReciprocityTheory.ChartRestriction.G.Out.QkToAkU}
\leanref{ReciprocityTheory/ReductiveGroups/ArithmeticGeometry/ExtensionDescent/GtPtBridgeSt.lean}{ReciprocityTheory.ChartRestriction.G.Out.QkToAk}

The map of Theorem~\ref{cv-toroidal-mumford}(2) carries each $k$-division point $p$ to the one point of $A(F)$ that
the uniformization assigns to $p$ in this sense.

\begin{lemma}[Alp\"oge]\label{cv-toroidal-torus-expansion-frame-square}
In Situation~\ref{cv-toroidal-torus-situation}, let $N\ge1$, let $p=(y,g)$ be an $N$-division point with
$g=(b,s)$, let $\nu_N$ be the normalized trivialization of Lemma~\ref{cv-toroidal-torus-symmetric-frame}(1), and let
$w_0$ be an isomorphism from $\gamma^*L$ to the restriction of $L_G$ to $G_K$. Let $U$ be a symmetric element of the
theta group of $L^N$ over $A_F$ \cite{Mum66}, and let $w$ be a frame of weight $N$ at $b$ such that, for all
sections $t$ and $t'$ of $L^N$ over $A_F$ with $t'=Ut$, the Fourier expansions $\sigma$ and $\sigma'$ of $t$ and
$t'$ with respect to $w_0$ satisfy $\sigma'=\Phi_{p,w}\sigma$. Then $w^{\otimes2}=s(y)\otimes\nu_N(b)$.
\end{lemma}
\leanref{Development/AbelianVarietyTheory/Proofs/UniformizationProofs/FrameSqEqOfSymmetricThetaElementStatement.lean}{GlobalTheory.T.frame_sq_eq_of_symmetric_thetaElement}
\leanref{Development/AbelianVarietyTheory/Proofs/UniformizationProofs/FrameSqEqOfSymmetricThetaElement.lean}{GlobalTheory.SymFrame.sq_frame_of_symmetric}
\leanref{Development/AbelianVarietyTheory/Proofs/UniformizationProofs/FrameSqEqOfSymmetricThetaElement.lean}{GlobalTheory.SymFrame.canon_identity}
\begin{proof}
Let $\theta$ be the cubical structure of $M$. Let $\rho$ be the
trivialization of $L$ at the origin that $w_0$ carries to the trivialization of $L_G$ there; then $w_0$ is a
normalization for $\rho$. Lemma~\ref{cv-toroidal-torus-expansion-multiples}(2) gives the endomorphism $E_{-1}$ of
the uniformization, multiplication by $-1$, which carries each point of $\tilde G(F)$ to its negative; the map
$-1$ on $X$ is injective. As $U$ is symmetric \cite{Mum66}, there is a trivialization $\alpha$ of
$L^N\otimes[-1]^*L^{-N}$ over $A_F$, canonical at the origin, such that $U(\alpha\otimes[-1]^*(Ut))=\alpha\otimes
[-1]^*t$ for every section $t$ of $L^N$ over $A_F$. The line bundle $L_G$ is symmetric
(Theorem~\ref{cv-toroidal-mumford}), so $L_G^N\otimes[-1]^*L_G^{-N}$ is trivial over $G$, and
Lemma~\ref{cv-toroidal-torus-expansion-pullback} applies to $E_{-1}$ with $k=k'=N$. Let $\kappa$ be a
trivialization of $M^N\otimes[-1]^*M^{-N}$ over $B$ that is canonical at the origin, as in its part (1).

The negation formula. Let $t$ be a section of $L^N$ over $A_F$, let $t'=\alpha\otimes[-1]^*t$, and let $\sigma$ and
$\sigma'$ be the Fourier expansions of $t$ and $t'$ with respect to $w_0$. Let $\chi,\chi'\in X$ with
$\chi+\chi'=0$, and let $x\in B(F)$. Then
\[
\sigma'_{\chi'}(x)=\sigma_\chi(-x)\otimes\kappa(x)\otimes\theta(c(\chi),x,-x)\otimes\theta(c(\chi),c(\chi'),x)^{-1}.
\tag{$*$}
\]
Indeed, Lemma~\ref{cv-toroidal-torus-extension}(1) gives a point $g_x$ of $\tilde G(F)$ with first component $x$;
write $s_x$ for its trivializations. By Definition~\ref{cv-toroidal-torus-points}(1), $-g_x$ has first component
$-x$ and trivializations $s_x(\chi)^{-1}\otimes\theta(c(\chi),x,-x)$, and $s_x(\chi)\otimes
s_x(\chi')=\theta(c(\chi),c(\chi'),x)$, as $s_x(0)$ is canonical. Let $\omega$ be a frame of weight $N$ at $x$, so
that $\omega\otimes\kappa(x)$ is a frame of weight $N$ at $-x$. As $E_{-1}$ carries $g_x$ to $-g_x$,
Lemma~\ref{cv-toroidal-torus-expansion-pullback}(2) gives
$a_{\chi'}(\sigma',g_x,\omega)=a_\chi(\sigma,-g_x,\omega\otimes\kappa(x))$, which by
Definition~\ref{cv-toroidal-torus-terms}(3) says that $\sigma'_{\chi'}(x)\otimes s_x(\chi')\otimes\omega$ is
$\sigma_\chi(-x)\otimes s_x(\chi)^{-1}\otimes\theta(c(\chi),x,-x)\otimes\omega\otimes\kappa(x)$. This gives $(*)$.

The operator $\Phi_{p,w}$ at a point. If $\sigma'=\Phi_{p,w}\sigma$, then for $\beta\in X$ and $z\in B(F)$,
\[
\sigma'_{\beta+y}(z)=\sigma_\beta(z+b)\otimes\prod_{0\le i<N}\theta(c(\beta)+ib,b,z)\otimes\theta(c(\beta),z,b)^{-1}
\otimes s(\beta)\otimes w,
\tag{$**$}
\]
by Definition~\ref{cv-toroidal-torus-theta-sum-phi} and the value of $u_\beta$ in the proof of
Lemma~\ref{cv-toroidal-torus-theta-sum-operator}(1).

The comparison. By Theorem~\ref{cv-toroidal-torus-sections}, $M_{N,-y}$ has a nonzero section over $B_F$, and
Lemma~\ref{cv-toroidal-torus-classes}(2) gives a family $\sigma$ of weight $N$ over $F$ whose component at $-y$ is
this section and whose components outside $-y+NX$ vanish. As $F$ is algebraically closed and $B_F$ is reduced,
there is $x\in B(F)$ with $\sigma_{-y}(-x)\ne0$. Let $t$ be the section of $L^N$ over $A_F$ with Fourier expansion
$\sigma$ (Lemma~\ref{cv-toroidal-torus-bijection}), let $t_1=Ut$, $t_2=\alpha\otimes[-1]^*t_1$ and $t_3=Ut_2$, and
let $\sigma_1$, $\sigma_2$ and $\sigma_3$ be their Fourier expansions. By hypothesis $\sigma_1=\Phi_{p,w}\sigma$ and
$\sigma_3=\Phi_{p,w}\sigma_2$, and $t_3=\alpha\otimes[-1]^*t$ by the choice of $\alpha$. Note that $c(y)=Nb$, as
$Ng=\iota(y)$, that $s(0)$ is canonical, and that $\theta$ is $1$ when one of its arguments is the origin. Then
$(**)$ for $\sigma_3$ at $\beta=0$ and $z=x$, $(*)$ for $\sigma_2$ at $\chi=\chi'=0$ and the point $x+b$, and
$(**)$ for $\sigma_1$ at $\beta=-y$ and $z=-(x+b)$ give
\[
(\sigma_3)_y(x)=\sigma_{-y}(-x)\otimes\kappa(x+b)\otimes s(-y)\otimes w^{\otimes2}\otimes C,\quad
C=\prod_{0\le i<N}\theta(ib,b,x)\,\theta(-Nb+ib,b,-(x+b))\otimes\theta(-Nb,-(x+b),b)^{-1}.
\]
On the other hand, $(*)$ for $t_3=\alpha\otimes[-1]^*t$ at $\chi=-y$, $\chi'=y$ and the point $x$ gives
$(\sigma_3)_y(x)=\sigma_{-y}(-x)\otimes\kappa(x)\otimes\theta(-Nb,x,-x)\otimes\theta(-Nb,Nb,x)^{-1}$. As
$\sigma_{-y}(-x)\ne0$,
\[
w^{\otimes2}\otimes s(-y)\otimes\kappa(x+b)\otimes C=\kappa(x)\otimes\theta(-Nb,x,-x)\otimes\theta(-Nb,Nb,x)^{-1}.
\]
Moreover
\[
\nu_N(b)\otimes\kappa(x+b)\otimes C\otimes\theta(Nb,-Nb,b)=\kappa(x)\otimes\theta(-Nb,x,-x)\otimes\theta(Nb,-Nb,x)^{-1}.
\]
Indeed, as functions of the pair $(b,x)$, both sides are trivializations of one line bundle on $B\times B$, built
from $\nu_N$, $\kappa$ and $\theta$, and both are canonical at the origin. Their ratio is an invertible function on
$B\times B$, hence a unit of $W$, as the global functions on $B\times B$ are those of $W$, and it is $1$ at the
origin. As $\theta$ is symmetric in its first two
arguments, the last two displays give $w^{\otimes2}\otimes s(-y)=\nu_N(b)\otimes\theta(Nb,-Nb,b)$. Finally
$s(y)\otimes s(-y)=\theta(Nb,-Nb,b)$ by Definition~\ref{cv-toroidal-torus-points}(1), as $y+(-y)=0$ and $s(0)$ is
canonical. So $w^{\otimes2}=s(y)\otimes\nu_N(b)$.
\end{proof}
% Lean: SymFrame.sq_frame_of_symmetric: exists_normW for rho; IntEnd.exists_nsmul_and_neg (the endomorphism for -1);
% exists_norm_sym for alpha; exists_gen with powGen of U.symG^-1 for the trivialization of L_G^N (-1)^*L_G^-N;
% T.expansion_pullback_section_endomorphism for kappa and the formula neg (gt_onto_divisible_kernel for g_x,
% TorusExt.exists_neg, CoefBasic.coef_basic and exists_gen for the frames, isCoef_coef, charVal_add, charVal_zero);
% FamFree.exists_sec_ne_zero (SecTwistSt) and exists_fam_single (UUA.exists_fam_rep) for sigma; exists_sval_ne_zero
% for x; exists_twist for t_1, t_2, t_3 and isTwistPull_of_sym; phi_at for (**) (e_1, e_3), neg (e_2, e_4);
% canon_identity (HasGen.unique over B x B, surj_snd) and C.s01. Size: FrameSqEqOfSymmetricThetaElement.lean,
% 434 lines.

\begin{lemma}[Alp\"oge]\label{cv-toroidal-torus-expansion-symmetric-lift}
In Situation~\ref{cv-toroidal-torus-situation}, let $N$ be odd, let $p=(y,g)$ be an $N$-division point with
$g=(b,s)$, and let $x\in A(F)$ be a point that the uniformization assigns to $p$
(Definition~\ref{cv-toroidal-torus-expansion-assigned}). Let $U$ be a symmetric element of the theta group of $L^N$
over $A_F$ lying over $x$ with $U^N=1$ \cite{Mum66}, let $w_0$ be an isomorphism from $\gamma^*L$ to the restriction
of $L_G$ to $G_K$, and let $w$ be a symmetric frame for $p$ (Definition~\ref{cv-toroidal-torus-symmetric}(2)). Let
$t$ and $t'$ be sections of $L^N$ over $A_F$ with $t'=Ut$, and let $\sigma$ and $\sigma'$ be their Fourier
expansions with respect to $w_0$. Then $\sigma'=\Phi_{p,w}\sigma$
(Definition~\ref{cv-toroidal-torus-theta-sum-phi}).
\end{lemma}
\leanref{Development/AbelianVarietyTheory/Proofs/UniformizationProofs/SymmetricLiftOddOrderActsAsExpansionOperatorStatement.lean}{GlobalTheory.T.symmetricLift_oddOrder_acts_as_expansionOperator}
\leanref{Development/AbelianVarietyTheory/Proofs/UniformizationProofs/SymmetricLiftOddOrderActsAsExpansionOperator.lean}{GlobalTheory.UT.symLift_acts_as_phi}
\leanref{Development/AbelianVarietyTheory/Proofs/RingsOfIntegers/ThetaGrpActIterateEqOfPowEqOne.lean}{GlobalTheory.thetaGrp_act_iterate_eq_of_pow_eq_one}
\begin{proof}
By Definition~\ref{cv-toroidal-torus-expansion-assigned} there is a frame $w'$ of weight $N$ at $b$ such that, for
all sections $t_1$ and $t_2$ of $L^N$ over $A_F$ with $t_2=Ut_1$, the Fourier expansions $\sigma_1$ and $\sigma_2$
satisfy $\sigma_2=\Phi_{p,w'}\sigma_1$. By Lemma~\ref{cv-toroidal-torus-expansion-phi-map}, $\Phi=\Phi_{p,w}$ and
$\Phi'=\Phi_{p,w'}$ are $F$-linear maps on the families of weight $N$ over $F$, and there is $c\in F^\times$ with
$w'=cw$ and $\Phi'=c\Phi$. By Lemma~\ref{cv-toroidal-torus-expansion-phi-order}, $\Phi^N$ is the identity. So is
$\Phi'^N$: for a family $\tau$, let $t_0$ be the section with Fourier expansion $\tau$
(Lemma~\ref{cv-toroidal-torus-bijection}); by induction on $i$, $\Phi'^i\tau$ is the Fourier expansion of $U^it_0$,
and $U^Nt_0=t_0$, as $U^N=1$ and $Nx=0$.

Suppose that $c^N=1$. Then $w'^{\otimes N}=w^{\otimes N}$, so $w'$ satisfies the second condition of
Definition~\ref{cv-toroidal-torus-symmetric}(2), as $w$ does; and $w'^{\otimes2}=s(y)\otimes\nu_N(b)$ by
Lemma~\ref{cv-toroidal-torus-expansion-frame-square}, applied to $U$ and $w'$. So $w'$ is a symmetric frame for $p$,
and $w'=w$ by Lemma~\ref{cv-toroidal-torus-symmetric-frame}(2). Hence $\sigma'=\Phi_{p,w'}\sigma=\Phi_{p,w}\sigma$.

Suppose that $c^N\ne1$. For every family $\tau$ of weight $N$ over $F$,
$\tau=\Phi'^N\tau=c^N\Phi^N\tau=c^N\tau$, so $\tau=0$. Then $\sigma'=0=\Phi_{p,w}\sigma$.
\end{proof}
% Lean: symLift_acts_as_phi: the field of QkToAkU for w' (hx w_0 hw_0, then hx' theta), phiFun_of for Phi and
% Phi', SymLiftPhi.exists_ratio for c, T.symmetricFrame_phi_pow_order_eq_id for Phi^N, ThetaNu.exists_isActOf and
% thetaGrp_act_iterate_eq_of_pow_eq_one (with expBijF_and_expandsF_lin) for Phi'^N; for c^N = 1,
% T.frame_sq_eq_of_symmetric_thetaElement and SymFrame.main (uniqueness), then phiRel_of_val_eq; otherwise every
% family is 0. Size: SymmetricLiftOddOrderActsAsExpansionOperator.lean, 119 lines.

\begin{lemma}[Alp\"oge]\label{cv-toroidal-torus-theta-sum-expansion}
In Situation~\ref{cv-toroidal-torus-situation}, let $\rho$ be a trivialization of $L$ at the origin and $w_0$
a normalization for $\rho$ (Definition~\ref{cv-toroidal-torus-expansion}(1)), let $\vartheta$ be a section of
$L$ over $A_F$, and let $\eta$ be its Fourier expansion with respect to $w_0$. Let $m$ be odd, let $p=(y,g)$ be
an $m^2$-division point with $g=(b,s)$, let $x_p\in A(F)$ be the point that the uniformization assigns to $p$
(Theorem~\ref{cv-toroidal-mumford}(2) and Definition~\ref{cv-toroidal-torus-expansion-assigned}), and let $t$ be
the section of $L^{m^2}$ over $A_F$ with which the
normalized theta constant of $\vartheta$ at $x_p$ is defined (Definition~\ref{cv-toroidal-torus-normalized} and
Lemma~\ref{cv-toroidal-torus-constant}). Let $\sigma$ be the Fourier expansion of $t$ with respect to $w_0$, let
$w$ be the symmetric frame for $p$ and $\kappa_m$ the normalized trivialization
(Lemma~\ref{cv-toroidal-torus-symmetric-frame}, for $N=m^2$ and $N=m$), let $w'=w\otimes\kappa_m(b)$, let $0$
be the zero point of $\tilde G(F)$, and let $w_1$ be the frame of weight $m^2$ at the origin of $B$ given by the
canonical trivialization of $M(0)^{m^2}\otimes M(0)^{-m^2}$. Then $a_{m\chi+y}(\sigma,0,w_1)=a_\chi(\eta,mg,w')$
for every $\chi\in X$, and $a_\beta(\sigma,0,w_1)=0$ for every $\beta\in X$ outside $y+mX$.
\end{lemma}
\leanref{Development/AbelianVarietyTheory/Proofs/UniformizationProofs/ThetaConstantEqSumThetaTermsOddMultipleStatement.lean}{GlobalTheory.T.thetaConstant_eq_sum_thetaTerms_odd_multiple}
\leanref{Development/AbelianVarietyTheory/Proofs/UniformizationProofs/ThetaConstantEqSumThetaTermsOddMultiple.lean}{GlobalTheory.UT.ThetaSum.theta_const_sum}
\begin{proof}
Let $\alpha_m$, $\theta$ and $s_1$ be the data of Definition~\ref{cv-toroidal-torus-normalized} for $t$:
$\alpha_m:[m]^*L\to L^{m^2}$ carries $\rho$ to $\rho^{m^2}$ at the origin, $\theta$ is the symmetric lift of
$x_p$ to the theta group of $L^{m^2}$ with $\theta^{m^2}=1$, $s_1=\alpha_m([m]^*\vartheta)$ and $t=\theta s_1$. Let
$\sigma_1$ be the Fourier expansion of $s_1$ with respect to $w_0$ (Lemma~\ref{cv-toroidal-torus-bijection}, for
$k=m^2$). Two facts relate the expansions.
\begin{enumerate}
\item[(i)] $a_{m\chi}(\sigma_1,g,w)=a_\chi(\eta,mg,w')$ for every $\chi\in X$, and $a_\beta(\sigma_1,g,w)=0$ for
  every $\beta\in X$ outside $mX$.
\item[(ii)] $\sigma=\Phi_{p,w}\sigma_1$ (Definition~\ref{cv-toroidal-torus-theta-sum-phi}).
\end{enumerate}

Proof of (i). Lemma~\ref{cv-toroidal-torus-expansion-multiples}(1) gives the endomorphism $E_m$ of the
uniformization, multiplication by $m$ on $G$, $A$, $B$ and $X$, which carries $g$ to $mg$. As $X$ is free and
$m\ge1$, multiplication by $m$ on $X$ is injective. The line bundle $L_G$ of the uniformization is symmetric and has
a cubical structure (Theorem~\ref{cv-toroidal-mumford}), so $[m]^*L_G\cong L_G^{m^2}$ by the argument of
\cite[\S6]{Mum70}, which uses only the cubical structure, the trivialization at the origin and the symmetry; hence
$L_G^{m^2}\otimes[m]^*L_G^{-1}$ is trivial over $G$. The isomorphism $\alpha_m$ is a trivialization of
$L^{m^2}\otimes[m]^*L^{-1}$ over $A_F$ whose value at the origin is $\rho^{m^2-1}$, and
$s_1=\alpha_m\otimes[m]^*\vartheta$. The trivialization $\kappa_m$ of $M^{m^2}\otimes[m]^*M^{-1}\otimes
M(0)^{1-m^2}$ is canonical at the origin, and $w\otimes\kappa_m(b)=w'$. So
Lemma~\ref{cv-toroidal-torus-expansion-pullback}(2) applies with $E=E_m$, $k=1$, $k'=m^2$, $\alpha=\alpha_m$,
the sections $\vartheta$ and $s_1$, the trivialization $\kappa_m$, the points $g$ and $mg$, and the frames $w$ of
weight $m^2$ at $b$ and $w'$ of weight $1$ at $mb$. It gives (i).

Proof of (ii). By Definition~\ref{cv-toroidal-torus-normalized}, $\theta$ is a symmetric element of the theta group
of $L^{m^2}$ over $A_F$ lying over $x_p$ with $\theta^{m^2}=1$, and $x_p$ is assigned to $p$ in the sense of
Definition~\ref{cv-toroidal-torus-expansion-assigned}. As $m^2$ is odd,
Lemma~\ref{cv-toroidal-torus-expansion-symmetric-lift}, for $N=m^2$, the element $\theta$ and the sections $s_1$ and
$t=\theta s_1$, gives $\sigma=\Phi_{p,w}\sigma_1$.
% Lean: (i) is T.expansion_pullback_section_endomorphism = UT.PullExp.main (Lemma
% cv-toroidal-torus-expansion-pullback), applied to the endomorphism of IntEnd.exists_nsmul_and_neg for m (Lemma
% cv-toroidal-torus-expansion-multiples(1)), with k = 1, k' = m^2, the trivialization exists_gen_pow (MulPull.hasGen_pow,
% a two-step induction from the cube, the rigidification and the symmetry of L_G, the argument of Mumford, Section 6),
% kappa = kappa_m and the frames w and w'; its conclusions are q_1 and q_2. (ii) is
% T.symmetricLift_oddOrder_acts_as_expansionOperator = UT.symLift_acts_as_phi (Lemma
% cv-toroidal-torus-expansion-symmetric-lift) for N = m^2, with IsActOf theta s_1 t. The descent in Lemma
% cv-toroidal-torus-expansion-descent is proved; nothing below them remains open.

By (ii) and Lemma~\ref{cv-toroidal-torus-theta-sum-operator}(2), $a_{\beta+y}(\sigma,0,w_1)=a_\beta(\sigma_1,g,w)$
for every $\beta\in X$. For $\beta=m\chi$, (i) gives $a_{m\chi+y}(\sigma,0,w_1)=a_\chi(\eta,mg,w')$. If
$\beta\notin y+mX$, then $\beta-y\notin mX$, and (i) gives $a_\beta(\sigma,0,w_1)=a_{\beta-y}(\sigma_1,g,w)=0$.
\end{proof}
% Lean: ThetaSum.theta_const_sum, second part (main): hbij.1 s_1 for sigma_1, hpull with q_1 and q_2 for (i),
% hPhi for (ii), CoefPhi.coef_phi with the zero point (Lemma cv-toroidal-torus-theta-sum-operator(2)) for
% key, then q_1 and q_2 through hEh.

\begin{lemma}[Alp\"oge]\label{cv-toroidal-torus-theta-sum}
In Situation~\ref{cv-toroidal-torus-situation}, let $\rho$ be a trivialization of $L$ at the origin and
$w_0$ a normalization for $\rho$ (Definition~\ref{cv-toroidal-torus-expansion}(1)), let $\vartheta$ be a
section of $L$ over $A_F$, and let $\eta$ be its Fourier expansion with respect to $w_0$. Let $m$ be odd, let
$p=(y,g)$ be an $m^2$-division point with $g=(b,s)$, let $x_p\in A(F)$ be the point that the uniformization
assigns to $p$ (Theorem~\ref{cv-toroidal-mumford}(2)), and let $c$ be the normalized theta constant of
$\vartheta$ at $x_p$ (Definition~\ref{cv-toroidal-torus-normalized} and
Lemma~\ref{cv-toroidal-torus-constant}). Let $w$ be the symmetric frame for $p$ and $\kappa_m$ the normalized
trivialization (Lemma~\ref{cv-toroidal-torus-symmetric-frame}, for $N=m^2$ and $N=m$), and let
$w'=w\otimes\kappa_m(b)$, a frame of weight $1$ at $mb$. Then $c=\sum_{\chi\in X}a_\chi(\eta,mg,w')$
$\mathfrak m$-adically (Definition~\ref{cv-toroidal-torus-terms}).
\end{lemma}
\leanref{Development/AbelianVarietyTheory/Proofs/UniformizationProofs/ThetaConstantEqSumThetaTermsOddMultipleStatement.lean}{GlobalTheory.T.thetaConstant_eq_sum_thetaTerms_odd_multiple}
\leanref{Development/AbelianVarietyTheory/Proofs/UniformizationProofs/ThetaConstantEqSumThetaTermsOddMultiple.lean}{GlobalTheory.UT.ThetaSum.theta_const_sum}
\begin{proof}
Let $t$ be the section of $L^{m^2}$ over $A_F$ with which $c$ is the normalized theta constant of $\vartheta$ at
$x_p$ (Lemma~\ref{cv-toroidal-torus-constant}), so that the value of $t$ at the origin is $c\rho^{m^2}$. Let
$\sigma$ be the Fourier expansion of $t$ with respect to $w_0$ (Lemma~\ref{cv-toroidal-torus-bijection}, for
$k=m^2$), let $0$ be the zero point of $\tilde G(F)$, and let $w_1$ be the frame of weight $m^2$ at the origin
of $B$ given by the canonical trivialization of $M(0)^{m^2}\otimes M(0)^{-m^2}$. Write $a_\beta$ for
$a_\beta(\sigma,0,w_1)$. By Lemma~\ref{cv-toroidal-torus-theta-sum-origin}, $c=\sum_{\beta\in X}a_\beta$
$\mathfrak m$-adically. By Lemma~\ref{cv-toroidal-torus-theta-sum-expansion}, $a_{m\chi+y}=a_\chi(\eta,mg,w')$ for
every $\chi\in X$, and $a_\beta=0$ for $\beta\notin y+mX$.

As $X$ is free and $m\ge1$, the map $\chi\mapsto m\chi+y$ is injective. Let $N$ be a finitely generated
$W$-submodule of $F$ as in Definition~\ref{cv-toroidal-torus-terms}(4) for the sum $c=\sum_\beta a_\beta$; it
contains $c$ and every $a_\chi(\eta,mg,w')=a_{m\chi+y}$. Let $n\ge0$, and let $S$ be the finite subset of $X$
given for $n$. Let $S'$ be the finite set of $\chi$ with $m\chi+y\in S$, and let $T\supset S'$ be finite. The set
$T''=S\cup\{m\chi+y:\chi\in T\}$ contains $S$, and $\sum_{\beta\in T''}a_\beta=\sum_{\chi\in T}a_{m\chi+y}$, since
the terms with $\beta\notin y+mX$ vanish and every $\beta=m\chi+y$ in $S$ has $\chi\in S'\subset T$. So
$c-\sum_{\chi\in T}a_\chi(\eta,mg,w')\in\mathfrak m^nN$. Hence $c=\sum_{\chi\in X}a_\chi(\eta,mg,w')$
$\mathfrak m$-adically.
\end{proof}
% Lean: ThetaSum.theta_const_sum, first part: hbij.1 t (expBijF_and_expandsF_lin, Lemma
% cv-toroidal-torus-bijection) for sigma, exists_zero_gt and coef_basic for the zero point o and the frame
% w_1 (FrmOne), main (the second part, Lemma cv-toroidal-torus-theta-sum-expansion), the third part of
% EvalZero.value_at_origin (Lemma cv-toroidal-torus-theta-sum-origin), and the third part of
% AdicSum.sumF_calculus (sumI_comp_iff, the re-indexing along the injection chi -> m chi + y). The descent
% in Lemma cv-toroidal-torus-expansion-descent, the last open part below this lemma, is now proved.

\subsection{Valuations of theta terms and good reduction}\label{subsection-cv-theta-term-valuations}

\begin{lemma}\label{cv-toroidal-torus-term-rules}
In Situation~\ref{cv-toroidal-torus-situation}, let $O$ be the ring of elements of $F$ of nonnegative
valuation, and for a point $x$ of $B$ let $M(x)$ be the fibre of $M$ at $x$. Valuations of trivializations are
those of Definition~\ref{cv-toroidal-torus-terms}(1), and a tensor product of fibres $M(x)$ at points of $B(F)$
is identified with every other in which each point occurs with the same total exponent.
\begin{enumerate}
\item Every trivialization of a tensor product of fibres of $M$ at points of $B(F)$ has exactly one valuation.
\item If trivializations $t$ and $t'$ have valuations $r$ and $r'$, then $t\otimes t'$ has valuation $r+r'$.
\item Let $Y$ be a scheme that is proper over $W$, let $f_1,\dots,f_n:Y\to B$ be morphisms over $W$ and
  $\epsilon_1,\dots,\epsilon_n\in\mathbf Z$, and let $\tau$ be a trivialization over $Y$ of
  $\bigotimes_if_i^*M^{\epsilon_i}$. Then for every $y\in Y(F)$ the value $\tau(y)$, a trivialization of
  $\bigotimes_iM(f_i(y))^{\epsilon_i}$, has valuation $0$.
\item Let $k\ge1$ and $\chi\in X$, let $\sigma$ be a section over $B$ of the line bundle $M_{k,\chi}$ of
  Definition~\ref{cv-toroidal-torus-points}(3), and let $b\in B(F)$. Then $\sigma(b)=at$ for some $a\in F$
  with $\nu(a)\ge0$ and a trivialization $t$ of valuation $0$ of the fibre of $M_{k,\chi}$ at $b$.
\end{enumerate}
\end{lemma}
\leanref{Development/AbelianVarietyTheory/Proofs/UniformizationProofs/ExistsGenV.lean}{GlobalTheory.UT.CoefVal.exists_genV}
\leanref{Development/AbelianVarietyTheory/Proofs/UniformizationProofs/ValuationsSymmetricFrameAndTerms.lean}{GlobalTheory.UT.TxVal.genV_unique}
\leanref{Development/AbelianVarietyTheory/Proofs/UniformizationProofs/ValuationsSymmetricFrameAndTerms.lean}{GlobalTheory.UT.TxVal.genV_mul}
\leanref{Development/AbelianVarietyTheory/Proofs/UniformizationProofs/ValuationsSymmetricFrameAndTerms.lean}{GlobalTheory.UT.TxVal.genV_pull}
\leanref{Development/AbelianVarietyTheory/Proofs/UniformizationProofs/ExistsGenV.lean}{GlobalTheory.UT.CoefVal.hv_sec}
\begin{proof}
As $\nu\ge0$ on $W$, $O$ contains the image of $W$. Every $x\in B(F)$ is the restriction of exactly one
point $x^0\in B(O)$: $B$ is proper over $W$, so the valuative criterion of properness gives one, and $B$ is
separated, so the valuative criterion of separatedness gives uniqueness.

(1) Let $t$ be a trivialization of $\bigotimes_iM(x_i)^{\epsilon_i}$. As $O$ is local,
$\bigotimes_iM(x_i^0)^{\epsilon_i}$ has a trivialization $t_0$ over $O$, and $t=at_0$ with $a\in F^\times$; so
$t$ has the valuation $\nu(a)$, which is rational. Let $t=a't'_0$ be a second such expression, with $t'_0$ a
trivialization over $O$ of $\bigotimes_jM(y_j^0)^{\delta_j}$, where each point of $B(F)$ occurs in
$\bigotimes_jM(y_j)^{\delta_j}$ with the same total exponent as in $\bigotimes_iM(x_i)^{\epsilon_i}$. As the
points $x^0$ are unique, each point of $B(O)$ also occurs in the two tensor products over $O$ with the same
total exponent. So $t'_0=ut_0$ for a unit $u$ of $O$, $a=a'u$ and $\nu(a)=\nu(a')$.

(2) If $t=at_0$ and $t'=a't'_0$ as in Definition~\ref{cv-toroidal-torus-terms}(1), then
$t\otimes t'=aa'(t_0\otimes t'_0)$, and $t_0\otimes t'_0$ is a trivialization over $O$. So $t\otimes t'$ has
valuation $\nu(a)+\nu(a')=r+r'$.

(3) As $Y$ is proper over $W$, $y$ is the restriction of a point $y_0\in Y(O)$, and $f_i(y_0)\in B(O)$
restricts to $f_i(y)$. The value $\tau(y_0)$ is a trivialization over $O$ of
$\bigotimes_iM(f_i(y_0))^{\epsilon_i}$, and $\tau(y)=1\cdot\tau(y_0)$ is its restriction. So $\tau(y)$ has
valuation $\nu(1)=0$.

(4) Let $z=b^0$. The fibre of $M_{k,\chi}$ at $z$ is a free $O$-module of rank one, as $O$ is local; let
$t_0$ be a generator, and $t$ its restriction to $b$. Then $\sigma(z)=at_0$ with $a\in O$, so $\sigma(b)=at$
with $\nu(a)\ge0$, and $t$ has valuation $0$.
\end{proof}
% Lean: (1) CoefVal.exists_genV (exists_liftO through lift_left, exists_gen_of_local, exists_ratio,
% wt_coe_of_ne) and genV_unique (CoefVal's for one combination, TxVal's for formally equal combinations
% FormEq with equal values; both through HV.cmp, HV.le and sec_int, with iota_inj, the valuative criterion of
% separatedness); (2) TxVal.genV_mul; (3) TxVal.genV_pull and CoefVal.hv_gen; (4) CoefVal.hv_sec.

\begin{lemma}\label{cv-toroidal-torus-term-points}
In Situation~\ref{cv-toroidal-torus-situation}, let $\tilde G(F)$ be the group of
Definition~\ref{cv-toroidal-torus-points}(1).
\begin{enumerate}
\item For every $g\in\tilde G(F)$ there is exactly one function $l_g:X\to\mathbf Q$ such that $g$ has
  valuations $l_g$ (Definition~\ref{cv-toroidal-torus-terms}(1)), and $l_g$ is additive.
\item For $g,g'\in\tilde G(F)$, $l_{g+g'}=l_g+l_{g'}$.
\end{enumerate}
\end{lemma}
\leanref{Development/AbelianVarietyTheory/Proofs/UniformizationProofs/ValuationsTorusPointsFramesFamilyTermsStatement.lean}{GlobalTheory.T.valuations_torusPoints_frames_familyTerms}
\leanref{Development/AbelianVarietyTheory/Proofs/UniformizationProofs/ValuationsTorusPointsFramesFamilyTerms.lean}{GlobalTheory.UT.CoefVal.main}
\leanref{Development/AbelianVarietyTheory/Proofs/UniformizationProofs/ValuationsTorusPointsFramesFamilyTerms.lean}{GlobalTheory.UT.CoefVal.exists_gtV}
\leanref{Development/AbelianVarietyTheory/Proofs/UniformizationProofs/ValuationsTorusPointsFramesFamilyTerms.lean}{GlobalTheory.UT.CoefVal.gtV_add}
\leanref{Development/AbelianVarietyTheory/Proofs/UniformizationProofs/ValuationsTorusPointsFramesFamilyTerms.lean}{GlobalTheory.UT.CoefVal.add_of_cub}
\leanref{Development/AbelianVarietyTheory/Proofs/UniformizationProofs/ValuationsTorusPointsFramesFamilyTerms.lean}{GlobalTheory.UT.CoefVal.hv_cub}
\begin{proof}
Valuations are taken in the sense of Lemma~\ref{cv-toroidal-torus-term-rules}. For a point $x$ of $B$ let
$M(x)$ be the fibre of $M$ at $x$. Let $\theta$ be the cubical structure of $M$, a trivialization over $B^3$ of
the line bundle whose fibre at $(x,x',x'')$ is
$M(x+x'+x'')\otimes M(x+x')^{-1}\otimes M(x+x'')^{-1}\otimes M(x'+x'')^{-1}\otimes M(x)\otimes M(x')\otimes
M(x'')\otimes M(0)^{-1}$. As $B^3$ is proper over $W$, its value at every point of $B^3(F)$ has valuation $0$,
by Lemma~\ref{cv-toroidal-torus-term-rules}(3).

(1) Let $g=(b,s)$. For $\chi\in X$, $s(\chi)$ is a trivialization of the fibre of $\Lambda(M)^{-1}$ at
$(c(\chi),b)$, a tensor product of fibres of $M$ at points of $B(F)$. By
Lemma~\ref{cv-toroidal-torus-term-rules}(1) it has exactly one valuation $l_g(\chi)$. Let $\chi,\chi'\in X$. By
Definition~\ref{cv-toroidal-torus-points}(1), $s(\chi+\chi')\otimes\theta(c(\chi),c(\chi'),b)=s(\chi)\otimes s(\chi')$.
By Lemma~\ref{cv-toroidal-torus-term-rules}(1) and~(2), $l_g(\chi+\chi')+0=l_g(\chi)+l_g(\chi')$.

(2) Let $g=(b,s)$, $g'=(b',s')$ and $g+g'=(b+b',s'')$. By Definition~\ref{cv-toroidal-torus-points}(1),
$s''(\chi)\otimes\theta(c(\chi),b,b')=s(\chi)\otimes s'(\chi)$ for every $\chi\in X$. By
Lemma~\ref{cv-toroidal-torus-term-rules}(1) and~(2), $l_{g+g'}(\chi)+0=l_g(\chi)+l_{g'}(\chi)$.
\end{proof}
% Lean: (1) is CV's first part, exists_gtV: exists_genV_gt, add_of_cub with hv_cub (the cubical structure at
% points over O, hv_gen) and gt_mul, uniqueness by genV_unique. (2) is CV's fourth part, gtV_add, with
% add_of_cub at (c(chi), g.b, g'.b). CV (T.valuations_torusPoints_frames_familyTerms = UT.CoefVal.main)
% also contains Lemmas cv-toroidal-torus-term-rules(1), cv-toroidal-torus-term-division and
% cv-toroidal-torus-term-bound, and two parts not used here (genV_frmOne, gtV_iota) and a null-sequence
% statement (finite_sublevel).

\begin{lemma}\label{cv-toroidal-torus-term-division}
In Situation~\ref{cv-toroidal-torus-situation}, let $Q$ be the function of
Lemma~\ref{cv-toroidal-torus-quadratic} and $b(\chi,\chi')=Q(\chi+\chi')-Q(\chi)-Q(\chi')$. Let $k\ge1$ and let
$p=(y,g)$ be a $k$-division point (Definition~\ref{cv-toroidal-torus-points}(2)). Then $g$ has valuations
$\chi\mapsto b(y,\chi)/k$ (Definition~\ref{cv-toroidal-torus-terms}(1)).
\end{lemma}
\leanref{Development/AbelianVarietyTheory/Proofs/UniformizationProofs/ValuationsTorusPointsFramesFamilyTerms.lean}{GlobalTheory.UT.CoefVal.gtV_qk}
\leanref{ReciprocityTheory/ReductiveGroups/ArithmeticGeometry/DiscreteValuationRings/VRatioGtStatement.lean}{ReciprocityTheory.ChartRestriction.G.DVR.TWM.VRatioGtStatement}
\leanref{Development/AbelianVarietyTheory/Proofs/RingsOfIntegers/ValuationThetaRatioEqBilinear.lean}{GlobalTheory.RatioVal.main}
\leanref{Development/AbelianVarietyTheory/Proofs/RingsOfIntegers/ValuationThetaRatioEqBilinear.lean}{GlobalTheory.RatioVal.v_ratio_gt}
\leanref{ReciprocityTheory/ReductiveGroups/ArithmeticGeometry/DiscreteValuationRings/PhiGenPf.lean}{ReciprocityTheory.ChartRestriction.G.DVR.chain_val}
\begin{proof}
Valuations are taken in the sense of Lemma~\ref{cv-toroidal-torus-term-rules}. Let $g=(b,s)$, let $\theta$ be
the cubical structure of $M$, as in the proof of Lemma~\ref{cv-toroidal-torus-term-points}, and let
$\alpha\in X$. As $kg=\iota(y)$, $kb=c(y)$ and the trivializations of the point $kg$ are the $\tau(y,\cdot)$.
For $i\ge0$ let $s_i$ be the trivializations of the point $ig$. A computation along $g,2g,\dots,kg$ with
Definition~\ref{cv-toroidal-torus-points}(1) and the symmetry of the cubical structure gives
$s_i(\alpha)\otimes\prod_{0\le j<i}\theta(jb,b,c(\alpha))=s(\alpha)^{\otimes i}$ by induction on $i$; for $i=k$,
\[
\tau(y,\alpha)\otimes\prod_{0\le j<k}\theta(jb,b,c(\alpha))=s(\alpha)^{\otimes k}.
\]
By Lemma~\ref{cv-toroidal-torus-term-rules}(1), $s(\alpha)$ has a valuation $l(\alpha)$, so the right side has
valuation $kl(\alpha)$ by Lemma~\ref{cv-toroidal-torus-term-rules}(2). By the proof of
Lemma~\ref{cv-toroidal-torus-quadratic}, $\tau(y,\alpha)$ is $\varpi^{b(y,\alpha)}$ times a generator over $W$, so
it has valuation $b(y,\alpha)$. Each $\theta(jb,b,c(\alpha))$ has valuation $0$ by
Lemma~\ref{cv-toroidal-torus-term-rules}(3), for $Y=B^3$. So the left side has valuation $b(y,\alpha)$, and
Lemma~\ref{cv-toroidal-torus-term-rules}(1) gives $kl(\alpha)=b(y,\alpha)$.
\end{proof}
% Lean: gtV_qk (CV's fifth part) writes s(alpha) = t G_0 with G_0 a generator over O at the extension m of b
% (exists_liftO, exists_gen_of_local, exists_ratio) and applies the statement VRatioGtStatement, which the
% archive proves by RatioVal.main = v_ratio_gt: m^k = c(y) over O (liftO_unique), the identity chain_val
% (the computation above; the same one is used in Lemma cv-toroidal-torus-symmetric-frame(2)), the valuation
% kQ(y) + b(y, alpha) of the trivialization psi(y)^k tau(y, alpha) e^k (famMult_data) together with those of
% psi(y)^k and e^k (genHasVal_powGen, rigK_val), which is the valuation b(y, alpha) of tau(y, alpha) used in
% the text, exists_cubN for the cubical factor over O, H1.exists_const for the unit, and unit_calc.

\begin{lemma}\label{cv-toroidal-torus-term-bound}
In Situation~\ref{cv-toroidal-torus-situation}, let $Q$ be the function of
Lemma~\ref{cv-toroidal-torus-quadratic} and let $k\ge1$. Let $\sigma$ be a family of Fourier coefficients of
weight $k$ over $F$ and $r$ a rational number with $\nu(\sigma_\chi)\ge r+Q(\chi)/k$ for every $\chi\in X$
(Definition~\ref{cv-toroidal-torus-points}(4)). Let $g=(b,s)\in\tilde G(F)$ have valuations $l$, and let $w$
be a frame of weight $k$ at $b$ of valuation $\omega$ (Definition~\ref{cv-toroidal-torus-terms}). Then
$\nu(a_\chi(\sigma,g,w))\ge r+Q(\chi)/k+l(\chi)+\omega$ for every $\chi\in X$.
\end{lemma}
\leanref{Development/AbelianVarietyTheory/Proofs/UniformizationProofs/ValuationsTorusPointsFramesFamilyTerms.lean}{GlobalTheory.UT.CoefVal.le_v_coef}
\leanref{Development/AbelianVarietyTheory/Proofs/UniformizationProofs/ValuationsTorusPointsFramesFamilyTerms.lean}{GlobalTheory.UT.CoefVal.le_v_int}
\begin{proof}
Valuations are taken in the sense of Lemma~\ref{cv-toroidal-torus-term-rules}. Let $e$ be the trivialization
of $M$ at the origin; it is a generator over $W$, so $e^k$ has valuation $0$. Let $\chi\in X$, and write
$\sigma_\chi=\sum_ia_i\sigma_i$ with sections $\sigma_i$ of $M_{k,\chi}$ over $B$ and $\nu(a_i)\ge r+Q(\chi)/k$.
For each $i$ let $x_i\in F$ be given by $\sigma_i(b)\otimes s(\chi)\otimes w=x_ie^k$. By
Definition~\ref{cv-toroidal-torus-terms}(3), $a_\chi(\sigma,g,w)e^k=\sum_ia_i\,\sigma_i(b)\otimes s(\chi)\otimes w$,
so $a_\chi(\sigma,g,w)=\sum_ia_ix_i$. By Lemma~\ref{cv-toroidal-torus-term-rules}(4), $\sigma_i(b)=c_it_i$ with
$\nu(c_i)\ge0$ and $t_i$ of valuation $0$. If $c_i=0$, then $x_i=0$. Otherwise $c_it_i\otimes s(\chi)\otimes w$
has valuation $\nu(c_i)+l(\chi)+\omega$ by Definition~\ref{cv-toroidal-torus-terms}(1) and
Lemma~\ref{cv-toroidal-torus-term-rules}(2), and $x_ie^k$ has valuation $\nu(x_i)$; in the tensor product of
fibres of $M$ that they trivialize, each point occurs with the same total exponent as in $M(0)^k$, so
$\nu(x_i)=\nu(c_i)+l(\chi)+\omega\ge l(\chi)+\omega$ by Lemma~\ref{cv-toroidal-torus-term-rules}(1). So every
$a_ix_i$ has valuation at least $r+Q(\chi)/k+l(\chi)+\omega$, and so does their sum.
\end{proof}
% Lean: le_v_coef (the decomposition VGe, exists_isCoef and coef_eq for the sum, AddValuation.map_le_sum) and
% le_v_int (hv_sec for sigma_i(b), hv_of_genV for s(chi) and w, hv_gen of the rigidification for e^k, and
% HV.cmp). This is CV's last part.

\begin{lemma}[Alp\"oge]\label{cv-toroidal-torus-term-valuations}
In Situation~\ref{cv-toroidal-torus-situation}, let $Q$ be the function of
Lemma~\ref{cv-toroidal-torus-quadratic} and $b(\chi,\chi')=Q(\chi+\chi')-Q(\chi)-Q(\chi')$. Let $m\ge1$, let
$p=(y,g)$ be an $m^2$-division point with $g=(b,s)$, let $w$ be a symmetric frame for $p$ and $\kappa_m$ the
normalized trivialization (Definition~\ref{cv-toroidal-torus-symmetric}), and let $w'=w\otimes\kappa_m(b)$.
Then $w'$ has valuation $Q(y)/m^2$, the point $mg$ has valuations $\chi\mapsto b(y,\chi)/m$
(Definition~\ref{cv-toroidal-torus-terms}(1)), and for every family $\sigma$ of Fourier coefficients of
weight $1$ over $F$ with $\nu(\sigma_\chi)\ge Q(\chi)$ for every $\chi\in X$,
$\nu(a_\chi(\sigma,mg,w'))\ge Q(m\chi+y)/m^2$ for every $\chi\in X$.
\end{lemma}
\leanref{Development/AbelianVarietyTheory/Proofs/UniformizationProofs/ValuationsSymmetricFrameAndTermsStatement.lean}{GlobalTheory.T.valuations_symmetricFrame_and_terms}
\leanref{Development/AbelianVarietyTheory/Proofs/UniformizationProofs/ValuationsSymmetricFrameAndTerms.lean}{GlobalTheory.UT.TxVal.main}
\begin{proof}
Valuations of trivializations are taken in the sense of Lemma~\ref{cv-toroidal-torus-term-rules}, and
$\nu_N$ is the normalized trivialization of Definition~\ref{cv-toroidal-torus-symmetric}(1).

The frame $w$. Let $N=m^2$. By Lemma~\ref{cv-toroidal-torus-term-rules}(1), $w$ has a valuation $r$, and
$w^{\otimes2}$ has valuation $2r$ by Lemma~\ref{cv-toroidal-torus-term-rules}(2). By
Lemma~\ref{cv-toroidal-torus-term-division}, $s(y)$ has valuation $b(y,y)/N$, which is $2Q(y)/N$ by the
proof of Lemma~\ref{cv-toroidal-torus-shift}(1). By Lemma~\ref{cv-toroidal-torus-term-rules}(3), for $Y=B$,
$\nu_N(b)$ has valuation $0$. So $s(y)\otimes\nu_N(b)$ has valuation $2Q(y)/N$. As $w$ is a symmetric frame,
$w^{\otimes2}=s(y)\otimes\nu_N(b)$, and Lemma~\ref{cv-toroidal-torus-term-rules}(1) gives $2r=2Q(y)/N$. So $w$
has valuation $Q(y)/m^2$.

The frame $w'$. By Lemma~\ref{cv-toroidal-torus-term-rules}(3), $\kappa_m(b)$ has valuation $0$, so
$w'=w\otimes\kappa_m(b)$ has valuation $Q(y)/m^2$ by Lemma~\ref{cv-toroidal-torus-term-rules}(2).

The point $mg$. For $h\in\tilde G(F)$ let $l_h$ be the function of
Lemma~\ref{cv-toroidal-torus-term-points}(1). As $0+0=0$, Lemma~\ref{cv-toroidal-torus-term-points}(2) gives
$l_0=l_0+l_0$, so $l_0=0$. By Lemma~\ref{cv-toroidal-torus-term-division}, for $k=m^2$,
$l_g(\chi)=b(y,\chi)/m^2$. By Lemma~\ref{cv-toroidal-torus-term-points}(2) and induction on $i\ge0$,
$l_{ig}(\chi)=ib(y,\chi)/m^2$. For $i=m$, $mg$ has valuations $\chi\mapsto b(y,\chi)/m$.

The theta terms. Let $\sigma$ be a family of Fourier coefficients of weight $1$ over $F$ with
$\nu(\sigma_\chi)\ge Q(\chi)$ for every $\chi$. Lemma~\ref{cv-toroidal-torus-term-bound}, for $k=1$, $r=0$,
the point $mg$ with $l=l_{mg}$ and the frame $w'$ with $\omega=Q(y)/m^2$, gives
\[
\nu(a_\chi(\sigma,mg,w'))\ge Q(\chi)+\frac{b(y,\chi)}m+\frac{Q(y)}{m^2}
=\frac{m^2Q(\chi)+mb(\chi,y)+Q(y)}{m^2}=\frac{Q(m\chi+y)}{m^2},
\]
as $b$ is symmetric and $Q(y+m\chi)=Q(y)+m\bigl(mQ(\chi)+b(\chi,y)\bigr)$ by
Lemma~\ref{cv-toroidal-torus-shift}(1).
\end{proof}
% Lean: TxVal.main. The first part (part1: the valuation Q(y)/N of a frame w with TXFrame N p w, of which
% only w.IsGen and the square condition are used) by CV's second part (exists_genV_frm), exists_nu,
% exists_liftO, genV_pull, genV_mul, genV_unique, CV's fifth part (gtV_qk) and QForm.two_Q_eq. A_1 by part1,
% genV_pull for kappa and genV_unique. A_2 by CV's first and fourth parts (exists_gtV, gtV_add) along the
% chain of IsNsmul m p.g g', with the zero point handled by uniqueness (the cube face identity face3 gives
% 0 + 0 = 0). The bound by CV's last part (le_v_coef) for k = 1, r = 0, with q_1, q_2, q_3 (bHom, Q_nsmul,
% two_Q_eq) for Q(m chi + y). CV is T.valuations_torusPoints_frames_familyTerms = UT.CoefVal.main, whose
% parts used here are Lemmas cv-toroidal-torus-term-rules to cv-toroidal-torus-term-bound. Size: about 3,300
% lines of Lean below it, counting the theorems it depends on below the laws LL and the statements the book
% cites.

\begin{lemma}[Alp\"oge]\label{cv-toroidal-torus-series}
In Situation~\ref{cv-toroidal-torus-situation}, let $Q$ be the function of
Lemma~\ref{cv-toroidal-torus-quadratic} and $b(\chi,\chi')=Q(\chi+\chi')-Q(\chi)-Q(\chi')$. Let $\rho$ be a
trivialization of $L$ at the origin and $w_0$ a normalization for $\rho$, let $\vartheta$ be a section of $L$
over $A_F$ whose Fourier expansion $\eta$ with respect to $w_0$ is primitive, and let $m$ be odd. Let
$p=(y,g)$ be an $m^2$-division point, $x_p$ the point that the uniformization assigns to $p$, and $c$ the
normalized theta constant of $\vartheta$ at $x_p$. Then there is a frame $w'$ of weight $1$ at the first
component of $mg$ such that $c=\sum_{\chi\in X}a_\chi(\eta,mg,w')$ $\nu$-adically, $w'$ has valuation
$Q(y)/m^2$, $mg$ has valuations $\chi\mapsto b(y,\chi)/m$, and $\nu(a_\chi(\eta,mg,w'))\ge Q(m\chi+y)/m^2$ for
every $\chi\in X$.
\end{lemma}
\leanref{Development/AbelianVarietyTheory/Proofs/UniformizationProofs/ThetaConstantValuationsUnitsEveryCoset.lean}{GlobalTheory.UT.ThetaConstVal.exists_terms}
\leanref{Development/AbelianVarietyTheory/Proofs/UniformizationProofs/ThetaConstantValuationsUnitsEveryCoset.lean}{GlobalTheory.UT.ThetaConstVal.exists_frm}
\begin{proof}
Let $g=(b,s)$. As $m^2$ is odd, Lemma~\ref{cv-toroidal-torus-symmetric-frame} gives the symmetric frame $w$
for $p$, for $N=m^2$, and the normalized trivialization $\kappa_m$. The tensor product
$w'=w\otimes\kappa_m(b)$ trivializes
$M(0)^{m^2}\otimes M(b)^{-m^2}\otimes M(b)^{m^2}\otimes M(mb)^{-1}\otimes M(0)^{1-m^2}=M(0)\otimes M(mb)^{-1}$, so
it is a frame of weight $1$ at $mb$, the first component of $mg$. By
Lemma~\ref{cv-toroidal-torus-theta-sum}, $c=\sum_\chi a_\chi(\eta,mg,w')$ $\mathfrak m$-adically, hence
$\nu$-adically by Lemma~\ref{cv-toroidal-torus-sums}(3). As $\eta$ is primitive,
$\nu(\eta_\chi)\ge Q(\chi)$ for every $\chi$ (Definition~\ref{cv-toroidal-torus-points}(4)), and
Lemma~\ref{cv-toroidal-torus-term-valuations} gives the valuations of $w'$ and of $mg$ and the bounds on the
theta terms.
\end{proof}

\begin{lemma}\label{cv-toroidal-torus-exact-section}
In Situation~\ref{cv-toroidal-torus-situation}, let $\chi\in X$, and let $M_{1,\chi}$ be the line bundle on $B$
of Definition~\ref{cv-toroidal-torus-points}(3). Then there is a section $\beta$ of $M_{1,\chi}$ over $B$ with
the following properties.
\begin{enumerate}
\item Every section of $M_{1,\chi}$ over $B$ is $c\beta$ for some $c\in W$.
\item Let $\kappa$ be an algebraically closed field with a homomorphism $W\to\kappa$, let $n\ge2$ be an integer
  that is nonzero in $\kappa$, and let $z\in B(\kappa)$. Then there is $x_0\in B(\kappa)$ with $nx_0=0$ such
  that $\beta$ does not vanish at $z+x_0$.
\end{enumerate}
\end{lemma}
\leanref{Development/AbelianVarietyTheory/Proofs/UniformizationProofs/ExistsGeneratorNeZero.lean}{GlobalTheory.UT.ExactTerm.exists_secB_good}
\leanref{Development/AbelianVarietyTheory/Proofs/UniformizationProofs/ExistsGeneratorNeZero.lean}{GlobalTheory.UT.ExactTerm.exists_generator_ne_zero}
\begin{proof}
For $x\in B(W)$ let $M(x)$ be the fibre of $M$ at $x$, a free $W$-module of rank one. As $c(0)=0$, the line
bundle $M_{1,0}$ is $M$. So Theorem~\ref{cv-toroidal-torus-sections}, for $k=1$ and $\chi=0$, gives a section
$\beta_0$ of $M$ over $B$ whose pullback along every homomorphism $W\to R'$ is a basis of the $R'$-module of
sections of $M$ over $B_{R'}$. Let $\xi$ be a generator of $M(c(\chi))^{-1}\otimes M(0)$, and let $T$ be the
translation by $c(\chi)$ on $B$. As $M_{1,\chi}=T^*M\otimes M(c(\chi))^{-1}\otimes M(0)$, the product
$\beta=T^*\beta_0\otimes\xi$ is a section of $M_{1,\chi}$ over $B$.

(1) Let $\sigma$ be a section of $M_{1,\chi}$ over $B$. Then $(T^{-1})^*\sigma\otimes\xi^{-1}$ is a section of
$M$ over $B$, so it is $c\beta_0$ for some $c\in W$. Pulling back along $T$ gives
$\sigma\otimes\xi^{-1}=cT^*\beta_0$, so $\sigma=c\beta$.

(2) The pullback $\beta_{0,\kappa}$ of $\beta_0$ to $B_\kappa$ is a basis of the sections of $M$ over
$B_\kappa$, so it is nonzero. The fibre $B_\kappa$ is an abelian variety over $\kappa$, on which $M$ restricts
to a symmetric line bundle that defines a principal polarization. Lemma~\ref{cv-volume-twisting-coset}, for
$\beta_{0,\kappa}$, $n$ and the point $z+c(\chi)$, gives $x_0\in B(\kappa)$ with $nx_0=0$ such that $\beta_0$ does
not vanish at $z+c(\chi)+x_0$. The image of the generator $\xi$ in $\kappa$ is nonzero, so $\beta$ does not
vanish at $z+x_0$.
\end{proof}
% Lean: exists_secB_good: exists_generator_ne_zero for beta_0 (from the law SecBaseChangeSt for n = 1, the
% same base change statement for M itself; the text cites Theorem cv-toroidal-torus-sections for k = 1 and
% chi = 0), exists_gen_of_local for xi, VVRank.exists_gen_pullTo and exists_sec with pullMulSec for beta;
% (1) by exists_sec with pullMulSec along tau' (translation by -c(chi)); (2) by exists_good_point (the law
% GoodPtSt, Lemma cv-volume-twisting-coset, with PolBaseChangeSt for the fibre) at a + c(chi).

\begin{lemma}\label{cv-toroidal-torus-exact-lift}
Let $F$ be an algebraically closed field, let $O\subset F$ be a valuation ring with fraction field $F$, maximal
ideal $\mathfrak m_O$ and residue field $\kappa$, and let $H$ be a scheme that is finite and flat over $O$.
Then every morphism $p:\operatorname{Spec}\kappa\to H$ over $O$ is the restriction to $\operatorname{Spec}\kappa$
of a section $q:\operatorname{Spec}O\to H$ of $H$ over $O$.
\end{lemma}
\leanref{Development/AbelianVarietyTheory/Proofs/UniformizationProofs/ExistsTorsionLift.lean}{GlobalTheory.UT.ExactTerm.exists_section_lift}
\leanref{Development/AbelianVarietyTheory/Proofs/UniformizationProofs/ExistsTorsionLift.lean}{GlobalTheory.UT.ExactTerm.exists_ringHom_lift}
\leanref{Development/AbelianVarietyTheory/Proofs/UniformizationProofs/ExistsTorsionLift.lean}{GlobalTheory.UT.ExactTerm.exists_torsion_lift}
\begin{proof}
As $H$ is finite over $O$, $H=\operatorname{Spec}A$ for a finite and flat $O$-algebra $A$, and $p$ is a
homomorphism $\varphi:A\to\kappa$ whose restriction to $O$ is the residue map. Its kernel is a prime ideal of
$A$ lying over $\mathfrak m_O$. As $A$ is flat over $O$, going down \cite[Theorem~9.5]{Mat86} gives a prime
ideal $P\subset\ker\varphi$ of $A$ lying over $0$. Then $O\to A/P$ is injective, so the fraction field of
$A/P$ contains $F$. Let $a\in A/P$. As $A$ is finite over $O$, $a$ is integral over $O$, so it is algebraic
over $F$, and $a\in F$ as $F$ is algebraically closed. As the valuation ring $O$ is integrally closed in $F$,
$a\in O$. So $O\to A/P$ is an isomorphism, and its inverse composed with $A\to A/P$ is a homomorphism
$\psi:A\to O$ of $O$-algebras. For $a\in A$, $a-\psi(a)$ lies in $P\subset\ker\varphi$, so $\varphi(a)$ is the
residue of $\psi(a)$. So $q=\operatorname{Spec}\psi$ is a section of $H$ over $O$ whose restriction to
$\operatorname{Spec}\kappa$ is $p$.
\end{proof}
% Lean: exists_section_lift reduces to exists_ringHom_lift (Ideal.exists_ideal_le_liesOver_of_le for going
% down, IsAlgClosed.algebraMap_bijective_of_isIntegral and IsIntegrallyClosed.isIntegral_iff);
% exists_torsion_lift applies it to the pullback of the n-th power map along the zero section over O.

\begin{lemma}\label{cv-toroidal-torus-exact-value}
In Situation~\ref{cv-toroidal-torus-situation}, let $O$ be the ring of elements of $F$ of nonnegative valuation
and $\kappa$ its residue field. Let $\chi\in X$, let $\beta$ be a section over $B$ of the line bundle
$M_{1,\chi}$ of Definition~\ref{cv-toroidal-torus-points}(3), let $\sigma$ be a family of Fourier coefficients
of weight $1$ over $F$, and let $a\in F$ with $\sigma_\chi=a\beta$ over $B_F$. Let $g=(b,s)\in\tilde G(F)$, and
suppose that $b$ is the restriction of a point $z\in B(O)$ such that $\beta$ does not vanish at the image of
$z$ in $B(\kappa)$. Let $w$ be a frame of weight $1$ at $b$, and let $l$ and $\omega$ be rational numbers such
that $s(\chi)$ has valuation $l$ and $w$ has valuation $\omega$ (Definition~\ref{cv-toroidal-torus-terms}).
Then $\nu(a_\chi(\sigma,g,w))=\nu(a)+l+\omega$.
\end{lemma}
\leanref{Development/AbelianVarietyTheory/Proofs/UniformizationProofs/ExactTermAfterTorsion.lean}{GlobalTheory.UT.ExactTerm.v_coef}
\leanref{Development/AbelianVarietyTheory/Proofs/UniformizationProofs/ExactTermAfterTorsion.lean}{GlobalTheory.UT.ExactTerm.exists_unit_const}
\begin{proof}
Let $e$ be the trivialization of $M$ at the origin, and for a point $x$ of $B$ let $M(x)$ be the fibre of $M$
at $x$. As $B$ is separated over $W$, a point of $B(F)$ is the restriction of at most one point of $B(O)$, by
the valuative criterion of separatedness. So, by Definition~\ref{cv-toroidal-torus-terms}(1), $s(\chi)=x_1t_1$
with $\nu(x_1)=l$ and $t_1$ a trivialization over $O$ of
$M(c(\chi)+z)^{-1}\otimes M(c(\chi))\otimes M(z)\otimes M(0)^{-1}$, and $w=x_2t_2$ with $\nu(x_2)=\omega$ and $t_2$
a trivialization over $O$ of $M(0)\otimes M(z)^{-1}$. The value $\beta(z)$ is a section over $O$ of
$M(z+c(\chi))\otimes M(c(\chi))^{-1}\otimes M(0)$ that does not vanish at the closed point; as $O$ is local, it
is a trivialization. So $\beta(z)\otimes t_1\otimes t_2\otimes e^{-1}$ is a trivialization over $O$ of a tensor
product of fibres of $M$ in which each fibre occurs as often with exponent $1$ as with exponent $-1$. That
tensor product is canonically trivial, and $\beta(z)\otimes t_1\otimes t_2\otimes e^{-1}=u$ for a unit $u$ of
$O$; so $\nu(u)=0$. By Definition~\ref{cv-toroidal-torus-terms}(3),
\[
a_\chi(\sigma,g,w)\,e=\sigma_\chi(b)\otimes s(\chi)\otimes w=a\,x_1x_2\,\beta(z)\otimes t_1\otimes t_2=a\,x_1x_2u\,e.
\]
So $a_\chi(\sigma,g,w)=ax_1x_2u$, and $\nu(a_\chi(\sigma,g,w))=\nu(a)+l+\omega$.
\end{proof}
% Lean: v_coef: exists_gen_of_isUnit for the generator beta(z) = t_0, the GenV data of s(chi) and w, and
% exists_unit_const (H1.exists_const for a generator of a formally trivial LinComb over Spec O, with
% comp_generic_injective, the uniqueness of extensions over O, to compare the points) for u; then
% CoefBasic.coef_eq.

\begin{lemma}[Alp\"oge]\label{cv-toroidal-torus-exact}
In Situation~\ref{cv-toroidal-torus-situation}, let $Q$ be the function of
Lemma~\ref{cv-toroidal-torus-quadratic}, let $\eta$ be a primitive family of Fourier coefficients of weight $1$
over $F$, let $n\ge2$ be an integer that is a unit of $W$, let $g\in\tilde G(F)$ and let $\chi\in X$. Then there
is $h\in\tilde G(F)$ with $nh=0$ such that, for $g'=g+h$, every frame $w$ of weight $1$ at the first component
of $g'$, every $l:X\to\mathbf Q$ and every rational $\omega$ such that $g'$ has valuations $l$ and $w$ has
valuation $\omega$ (Definition~\ref{cv-toroidal-torus-terms}),
$\nu(a_\chi(\eta,g',w))=Q(\chi)+l(\chi)+\omega$.
\end{lemma}
\leanref{Development/AbelianVarietyTheory/Proofs/UniformizationProofs/ExactTermAfterTorsionStatement.lean}{GlobalTheory.T.exists_nsmul_add_decomp_with_coef_valuation}
\leanref{Development/AbelianVarietyTheory/Proofs/UniformizationProofs/ExactTermAfterTorsion.lean}{GlobalTheory.UT.ExactTerm.exact_term_after_torsion}
\begin{proof}
Let $O$ be the ring of elements of $F$ of nonnegative valuation, a valuation ring with fraction field $F$,
and $\kappa$ its residue field, which is algebraically closed as $F$ is. As $\nu\ge0$ on $W$, $O$ contains
the image of $W$. For $z\in B(O)$ let $\bar z$ be its image in $B(\kappa)$.

The coefficient. Let $\beta$ be the section of $M_{1,\chi}$ over $B$ of
Lemma~\ref{cv-toroidal-torus-exact-section}. As $\eta$ is primitive, $\nu(\eta_\chi)\ge Q(\chi)$
(Definition~\ref{cv-toroidal-torus-points}(4)): $\eta_\chi=\sum_ia_i\sigma_i$ with sections $\sigma_i$ of
$M_{1,\chi}$ over $B$ and $\nu(a_i)\ge Q(\chi)$. By Lemma~\ref{cv-toroidal-torus-exact-section}(1),
$\sigma_i=c_i\beta$ with $c_i\in W$. So $\eta_\chi=a\beta$ with $a=\sum_ia_ic_i$, and $\nu(a)\ge Q(\chi)$, as
$\nu(c_i)\ge0$. Suppose that $\nu(a)>Q(\chi)$, and choose a rational $q$ with $Q(\chi)<q\le\nu(a)$. Then
$\nu(\eta_\chi)\ge q$, as $\eta_\chi=a\beta$. The set $\{\chi\}$ meets the only class of $X$ modulo $1\cdot X$,
so Lemma~\ref{cv-toroidal-torus-shift}(3), for $k=1$ and $r=q-Q(\chi)$, gives
$\nu(\eta_{\chi'})\ge r+Q(\chi')$ for every $\chi'\in X$. As $r>0$, this contradicts the primitivity of
$\eta$. So $\nu(a)=Q(\chi)$.

The torsion point. Let $g=(b,s)$. As $B$ is proper over $W$, the valuative criterion of properness extends
$b$ to a point $z_0\in B(O)$. As $n$ is a unit of $W$, it is nonzero in $\kappa$.
Lemma~\ref{cv-toroidal-torus-exact-section}(2), for $\kappa$, $n$ and $\bar z_0$, gives $x_0\in B(\kappa)$ with
$nx_0=0$ such that $\beta$ does not vanish at $\bar z_0+x_0$. Multiplication by $n$ on $B$ is finite and
flat: $B$ is proper over $W$, and on each fibre multiplication by $n$ is an isogeny \cite{Mum70}, so it is
flat by the fibrewise criterion of flatness, and it is proper and quasi-finite, hence finite. So its kernel
$B[n]$, the fibre product of multiplication by $n$ with the zero section, is finite and flat over $W$, and
its base change to $O$ is finite and flat over $O$. Lemma~\ref{cv-toroidal-torus-exact-lift}, applied to
this base change and the point $x_0$ of $B[n]$, gives $x\in B(O)$ with $nx=0$ and $\bar x=x_0$. Let
$x_F\in B(F)$ be the restriction of $x$.

By Lemma~\ref{cv-toroidal-torus-extension}(1) there is $g_1\in\tilde G(F)$ with first component $x_F$. The
first component of $-ng_1$ is $-nx_F=0$, so by Lemma~\ref{cv-toroidal-torus-extension}(2), $-ng_1=(0,t)$ for a
homomorphism $t:X\to F^\times$. As $X$ is free and $F^\times$ is divisible, $t=s_0^n$ for a homomorphism
$s_0:X\to F^\times$. Put $h=g_1+(0,s_0)$. Its first component is $x_F$, and $nh=ng_1+(0,s_0^n)=ng_1+(0,t)=0$.

The term. Let $g'=g+h=(b',s')$. Its first component $b'=b+x_F$ is the restriction of $z=z_0+x\in B(O)$,
and $\beta$ does not vanish at $\bar z=\bar z_0+x_0$. Let $w$ be a frame of weight $1$ at $b'$, and let $l$
and $\omega$ be such that $g'$ has valuations $l$ and $w$ has valuation $\omega$. Then $s'(\chi)$ has
valuation $l(\chi)$, and Lemma~\ref{cv-toroidal-torus-exact-value}, for $\sigma=\eta$, the element $a$ and the
point $g'$, gives $\nu(a_\chi(\eta,g',w))=\nu(a)+l(\chi)+\omega=Q(\chi)+l(\chi)+\omega$.
\end{proof}
% Lean: ExactTerm.exact_term_after_torsion: exists_secB_good (with Laws) for beta; vge_iff (with the first
% part of Prim) for a with v(a) >= Q(chi), and the second part of Prim with n1a (Lemma
% cv-toroidal-torus-shift(3) for k = 1, S = {chi}) and WithTop.lt_iff_exists_coe_btwn for v(a) = Q(chi);
% isAlgClosed_residueField_of_valuationRing_of_isAlgClosed for kappa; exists_liftO (TWO.exists_lift, the
% valuative criterion) for z_0; exists_torsion_lift with the law mulN (finite and flat; MulNSt, the text
% cites Mumford and the fibrewise criterion as in Lemma cv-toroidal-moduli-torsion) for x;
% TorusExt.exists_nsmul_eq_of_pow_eq (exists_lift, exists_torusPt_eq, exists_pow_eq_hom) for h; and v_coef
% for the term. Size: about 3,200 lines of Lean below it, counting the theorems it depends on below the laws
% LL and the statements the book cites.

\begin{lemma}[Alp\"oge]\label{cv-toroidal-torus-valuations}
In Situation~\ref{cv-toroidal-torus-situation}, let $Q$ be the function of
Lemma~\ref{cv-toroidal-torus-quadratic}, let $\rho$ be a trivialization of $L$ at the origin and $w_0$ a
normalization for $\rho$, and let $\vartheta$ be a section of $L$ over $A_F$ whose Fourier expansion with
respect to $w_0$ is primitive (Definitions~\ref{cv-toroidal-torus-points}
and~\ref{cv-toroidal-torus-expansion}). Let $m$ be an odd integer that is a unit of $W$. For an
$m^2$-division point $p$ let $x_p\in A(F)$ be the point that the uniformization assigns to $p$
(Theorem~\ref{cv-toroidal-mumford}(2)), and let $c_p$ be the normalized theta constant of $\vartheta$ at
$x_p$ (Definition~\ref{cv-toroidal-torus-normalized} and Lemma~\ref{cv-toroidal-torus-constant}).
\begin{enumerate}
\item Let $p=(y,g)$ be an $m^2$-division point. Then $\nu(c_p)\ge r$ for every rational $r$ with
  $r\le Q(m\chi+y)/m^2$ for all $\chi\in X$, and if $\nu(c_p)=0$, then $y\in mX$.
\item Let $n\ge2$ divide $m$, and let $p=(0,g)$ be an $m^2$-division point. Then there is an $m^2$-division
  point $p'=(0,g')$ with $nm(x_{p'}-x_p)=0$ and $\nu(c_{p'})=0$.
\end{enumerate}
\end{lemma}
\leanref{Development/AbelianVarietyTheory/Proofs/UniformizationProofs/ThetaConstantValuationsUnitsEveryCosetStatement.lean}{GlobalTheory.T.theta_constant_valuations_units_every_coset}
\leanref{Development/AbelianVarietyTheory/Proofs/UniformizationProofs/ThetaConstantValuationsUnitsEveryCoset.lean}{GlobalTheory.UT.ThetaConstVal.main}
\begin{proof}
Let $\eta$ be the Fourier expansion of $\vartheta$ with respect to $w_0$, and let $f$ be the map of
Theorem~\ref{cv-toroidal-mumford}(2) for $m^2$, so that $x_p=f(p)$ and $f$ carries sums of division points
to sums. By Lemma~\ref{cv-toroidal-torus-quadratic}, $Q(\chi)/m^2\ge1/m^2$ for every $\chi\ne0$.

(1) Lemma~\ref{cv-toroidal-torus-series} gives a frame $w'$ with $c_p=\sum_\chi a_\chi(\eta,mg,w')$
$\nu$-adically and $\nu(a_\chi(\eta,mg,w'))\ge Q(m\chi+y)/m^2$ for every $\chi$. If $r\le Q(m\chi+y)/m^2$ for all
$\chi$, then every term has valuation at least $r$, and $\nu(c_p)\ge r$ by
Lemma~\ref{cv-toroidal-torus-sums}(1). Suppose that $\nu(c_p)=0$ and $y\notin mX$. Then $m\chi+y\ne0$ for
every $\chi$, so $r=1/m^2$ is allowed, and $\nu(c_p)\ge1/m^2>0$, a contradiction.

(2) Write $m=nd$. As $m$ is a unit of $W$, so is $n$. Lemma~\ref{cv-toroidal-torus-exact}, for $\eta$, $n$,
the point $mg$ and $\chi=0$, gives $h\in\tilde G(F)$ with $nh=0$ such that, for $g_2=mg+h$,
$\nu(a_0(\eta,g_2,w))=Q(0)+l(0)+\omega$ whenever $g_2$ has valuations $l$ and the frame $w$ of weight $1$ at the
first component of $g_2$ has valuation $\omega$. By Lemma~\ref{cv-toroidal-torus-extension}(3) there is
$h'\in\tilde G(F)$ with $mh'=h$. Then $nmh'=nh=0$ and $m^2h'=d(nmh')=0=\iota(0)$, so $p_1=(0,h')$ is an
$m^2$-division point, and so is the sum $p'=(0,g+h')$ of $p$ and $p_1$.

Let $p_2=(0,g_0)$ be an $m^2$-division point and $j\ge1$ with $jg_0=0$. The pairs $(0,ig_0)$, for $i\ge0$,
are $m^2$-division points, and $(0,(i+i')g_0)$ is the sum of $(0,ig_0)$ and $(0,i'g_0)$. So
$f((0,(i+i')g_0))=f((0,ig_0))+f((0,i'g_0))$, which gives $f((0,0))=0$ and $f((0,ig_0))=if(p_2)$ by
induction; as $(0,jg_0)=(0,0)$, $jf(p_2)=0$. For $p_2=p_1$ and $j=nm$ this gives $nmf(p_1)=0$; for $p_2=p'$
and $j=m^2$, as $m^2(g+h')=\iota(0)=0$, it gives $m^2x_{p'}=0$. As $x_{p'}=f(p)+f(p_1)=x_p+f(p_1)$,
$nm(x_{p'}-x_p)=0$, and
Lemma~\ref{cv-toroidal-torus-constant} gives the normalized theta constant $c_{p'}$ of $\vartheta$ at
$x_{p'}$.

Lemma~\ref{cv-toroidal-torus-series}, for $p'$, gives a frame $w'$ at the first component of
$m(g+h')=mg+h=g_2$ with $c_{p'}=\sum_\chi a_\chi(\eta,g_2,w')$ $\nu$-adically, such that $w'$ has valuation
$Q(0)/m^2=0$, $g_2$ has valuations $\chi\mapsto b(0,\chi)/m=0$, and
$\nu(a_\chi(\eta,g_2,w'))\ge Q(m\chi)/m^2$ for every $\chi$. By the choice of $h$,
$\nu(a_0(\eta,g_2,w'))=Q(0)=0$. Let $\chi\ne0$. If $m\chi=0$, then
$0=Q(m\chi)=m^2Q(\chi)\ge m^2$ by Lemma~\ref{cv-toroidal-torus-shift}(1) and
Lemma~\ref{cv-toroidal-torus-quadratic}, which is impossible. So $m\chi\ne0$ and
$\nu(a_\chi(\eta,g_2,w'))\ge1/m^2$. Lemma~\ref{cv-toroidal-torus-sums}(2), with $q=1/m^2$ and $i_0=0$, gives
$\nu(c_{p'})=0$.
\end{proof}
% Lean: ThetaConstVal.main: (1) from ThetaConstVal.exists_terms with le_of_sumV and Q_pos; (2) from
% T.exists_nsmul_add_decomp_with_coef_valuation, gt_onto_divisible_kernel, TorsionIso.torsion_eq (the
% Lean's own proof of the torsion description, about 11,500 lines, cf. Theorem cv-toroidal-mumford(2); the
% text uses the theorem), pow_eq_one_of_nsmul, T.existsUnique_theta_const_of_odd, Q_nsmul and
% eq_zero_of_sumV. Lemma cv-toroidal-torus-expansion-descent, below facts (i) and (ii) of Lemma
% cv-toroidal-torus-theta-sum-expansion, is now proved.
% It takes LL (Out.Laws to Laws5), TwistBCSpanStatement, VRatioGtStatement, TWStatement
% and UUStatement as hypotheses; the archive proves them (lawsAndExpansions; the expansion statements are
% twistBCSpan_of with twistBaseChange_of_finrank_sections_one, RatioVal.main, TWM.tw_of and
% UnitBallFree.main). Size: about 27,000 lines of Lean below it and below those statements, beyond the
% laws of abelian schemes (Out.Laws, Out.Laws2) and the other statements cited here; about 20,000 of them
% also serve the Lean's proof of Theorem cv-toroidal-mumford (thmU_and_uTors: Laws3 to Laws5, the
% expansions and the torsion description), and about 7,000 lie outside it. Text count following _root_
% references, stopping at statements the book cites; the earlier count of 26,000 lines for the whole
% proof of Lemma cv-toroidal-reduction-torus did not follow the _root_ references.

\begin{lemma}\label{cv-toroidal-torus-free}
Let $X$ be a free $\mathbf Z$-module of finite rank, and let $d\ge2$ be an integer such that every element
of $X$ is $d$ times an element of $X$. Then $X=0$.
\end{lemma}
\leanref{Development/AbelianVarietyTheory/Proofs/ModuliProofs/GoodReductionOfThetaPointUnits.lean}{GlobalTheory.ThetaCrit.subsingleton_of_divisible}
\begin{proof}
Let $(e_l)$ be a basis of $X$. If it had an element $e_l$, then $e_l=dy$ for some $y\in X$, and the
$l$-th coordinate $a$ of $y$ would satisfy $da=1$ in $\mathbf Z$, which is impossible for $d\ge2$. So the
basis is empty, and $X=0$.
\end{proof}

\begin{lemma}[Alp\"oge]\label{cv-toroidal-reduction-torus}
In Situation~\ref{cv-toroidal-reduction-situation}, let $n\ge2$ be odd and $j\ge2$, and suppose that
$N'=n^j$. Let $W$ be a complete discrete valuation ring of characteristic zero in which $N'$ is
invertible, with fraction field $K$, and let $\operatorname{Spec}K$ be a scheme over $o$. Let
$t=(A,L,(x_i),(\alpha_h))$ be a theta-level object over $K$, and let $G$ be a commutative semi-abelian
scheme over $W$, with degeneration data over $W$, with free $\mathbf Z$-module $X$, and a uniformization
of $(A,G)$ by them (Theorem~\ref{cv-toroidal-mumford}). Let $v$ be a family in $W$ whose image in $K$ is
a family of theta constants of $t$ at the identity, such that for every $c$ there is $c'$ for which
$v_{c+n^{j-1}c'}$ is a unit of $W$. Then $X=0$.
\end{lemma}
\leanref{Development/AbelianVarietyTheory/Proofs/ModuliProofs/GoodReductionOfThetaPointUnits.lean}{GlobalTheory.ThetaCrit.subsingleton_characters}
\begin{proof}
Write $j=i+1$, so that $i\ge1$, $N'=n^{i+1}$ and $N'/n=n^i$, and let $N=N'^2$. Let $\varpi$ be a
uniformizer of $W$, $F$ an algebraic closure of $K$ and $\nu$ the valuation of
Theorem~\ref{cv-toroidal-torus-valuation}. With $A$ and $L$ from $t$, the data are those of
Situation~\ref{cv-toroidal-torus-situation}; let $Q$ be the function of
Lemma~\ref{cv-toroidal-torus-quadratic}. As $L$ is trivialized along the zero section, it has a
trivialization $\rho$ at the origin. Lemma~\ref{cv-toroidal-torus-normalization} gives a normalization
$w_0$ for $\rho$, and Lemma~\ref{cv-toroidal-torus-primitive} gives a section $\vartheta$ of $L$ over $A_F$
that generates the sections of $L$ over $A_F$ and whose Fourier expansion with respect to $w_0$ is
primitive. Theorem~\ref{cv-toroidal-mumford}(2), for $m=N$, gives a map $f$ from the $N$-division points
onto the $N$-torsion of $A(F)$ that carries sums of division points to sums; $f(p)$ is the point $x_p$ of
Lemma~\ref{cv-toroidal-torus-valuations}.

The theta constants over $F$. For $a\in(\mathbf Z/N)^{2g}$ write $x_a$ also for the image in $A(F)$ of the
point $x_a$ of $t$. Then $x_{a+a'}=x_a+x_{a'}$, and every point of $A(F)$ killed by $N$ is $x_a$ for some
$a$, by Definition~\ref{cv-toroidal-theta}(1). Lemma~\ref{cv-toroidal-torus-constant}, for $m=N'$, gives
the normalized theta constant $q_a$ of $\vartheta$ at $x_a$. As $N'$ is a unit of $W$, it is nonzero in
$F$, so by Lemma~\ref{cv-toroidal-torus-family} the family $(q_a)$ is a family of theta constants of $t$ at
$\operatorname{Spec}F\to\operatorname{Spec}K$. By Lemma~\ref{cv-toroidal-reduction-base-change}(4), the
image of $v$ in $F$ is also one. The hypothesis for $c=0$ gives $b_0$ with $v_{n^ib_0}$ a unit of $W$, so
the image of $v$ generates the unit ideal of $F$, and Lemma~\ref{cv-toroidal-reduction-unit} gives
$u\in F^\times$ with $q_a=uv_a$ for every $a$. Here $\nu(v_a)\ge0$ for every $a$, as $v_a\in W$, and
$\nu(v_a)=0$ when $v_a$ is a unit of $W$.

The valuation of $u$ is $0$. Lemma~\ref{cv-toroidal-torus-valuations}(2), for $m=N'$, the divisor $n$ of
$N'$ and the division point $(0,0)$, gives an $N$-division point $p'$ such that the normalized theta
constant of $\vartheta$ at $f(p')$ has valuation $0$. As $f(p')$ is killed by $N$, $f(p')=x_a$ for some
$a$, and that constant is $q_a$ by the uniqueness in Lemma~\ref{cv-toroidal-torus-constant}. So
$0=\nu(q_a)=\nu(u)+\nu(v_a)\ge\nu(u)$. Conversely, let $p$ be an $N$-division point with
$f(p)=x_{n^ib_0}$; the normalized theta constant of $\vartheta$ at $f(p)$ is $q_{n^ib_0}$. As
$Q(\chi)\ge0$ for every $\chi$ (Lemma~\ref{cv-toroidal-torus-quadratic}),
Lemma~\ref{cv-toroidal-torus-valuations}(1), with $r=0$, gives $\nu(q_{n^ib_0})\ge0$, so
$\nu(u)=\nu(q_{n^ib_0})-\nu(v_{n^ib_0})\ge0$.

Every element of $X$ is divisible by $n^i$. Let $\chi\in X$. By
Lemma~\ref{cv-toroidal-torus-extension}(3) there is $g_0\in\tilde G(F)$ with $Ng_0=\iota(\chi)$, so
$p=(\chi,g_0)$ is an $N$-division point. As $f(p)$ is killed by $N$, $f(p)=x_a$ for some $a$, and by
hypothesis there is $b$ with $v_{a+n^ib}$ a unit of $W$. Let $p'=(y',g')$ be an $N$-division point with
$f(p')=x_b$, and let $r$ be the sum of $p$ and $n^i$ copies of $p'$, whose first component is
$\chi+n^iy'$. Then $f(r)=x_a+n^ix_b=x_{a+n^ib}$, so the normalized theta constant of $\vartheta$ at $f(r)$
is $q_{a+n^ib}$, and $\nu(q_{a+n^ib})=\nu(u)+\nu(v_{a+n^ib})=0$. By
Lemma~\ref{cv-toroidal-torus-valuations}(1) there is $y\in X$ with $\chi+n^iy'=N'y=n^{i+1}y$. So
$\chi=n^i(ny-y')$.

As $n\ge2$ and $i\ge1$, $n^i\ge2$, and Lemma~\ref{cv-toroidal-torus-free} shows that $X=0$.
\end{proof}
% Lean: subsingleton_characters: T.uniformization_laws_and_expansion_statements (the laws LL and the
% expansion statements, proved in the archive), exists_irreducible, DVR.QForm.exists_QVal (Q),
% adic_base_embeds_valuation_algebraic_closure (nu), t.pol.rigid (rho), UT.EvalZero.value_at_origin (w0),
% UT.exists_prim_generator, T.theta_constant_valuations_units_every_coset (ua1, ua2) with m = n^(i+1),
% UTorsSt with k = N (f, f1, f2, f4), ThetaCoord.lab (the x_a, with lab_add, lab_nsmul, t.bij),
% existsUnique_thetaConst (q), isThetaPt_of_normalised (hc1), hasTh_baseChange (hc2),
% T.thetaConstants_unique_up_to_unit (u, with q_a = u v_a), DVR.GC.v_alg_unit and v_alg_nonneg,
% u0 from ua2 with TorusExt.qkZero and from ua1 with Q_zero and Q_pos, comb (TorusExt.qkAdd,
% qkAdd_qkAdd) for r, UT.gt_onto_divisible_kernel and TorusExt.qkMk for p, and subsingleton_of_divisible
% with d = n^i. Lemma cv-toroidal-torus-primitive is proved; the open steps are in
% cv-toroidal-torus-theta-sum, -term-valuations and -exact.

\begin{lemma}[Alp\"oge]\label{cv-toroidal-theta-reduction}
Let $g\ge1$, let $I$ be a finite set with elements $w$ and $w'$, and let $n$ be an odd prime. There is
$j_0$ such that the following holds for every $j\ge j_0$, with $N'=n^j$. Let $o$ be a ring with a unit
$\zeta$ such that $N'$ is invertible in $o$ and every homomorphism from $o$ to a field maps $\zeta$ to a
primitive root of unity of order $N'^2$, let $V$ be a complete discrete valuation ring with finite residue
field and with fraction field $K$ of characteristic zero, and let $\beta:o\to V$ be a homomorphism. Let
$t=(A,L,(x_i),(\alpha_h))$ be a theta-level object over $K$ (Definition~\ref{cv-toroidal-theta}(1))
with $\alpha_w+\alpha_{w'}=1$, such that, with $\Lambda(L)$ as in Theorem~\ref{cv-toroidal-mumford}, $L$
is isomorphic to the pullback of $\Lambda(L)$ along $(1,\alpha_w)$ and the pullbacks of $\Lambda(L)$
along $\alpha_w\times1$ and $1\times\alpha_{w'}$ are isomorphic. Let $v$ be a family of theta constants
of $t$ at the identity with values in $V$ such that for every $c$ there is $c'$ for which
$v_{c+N'_nc'}$ is a unit of $V$, with $N'_n$ as in Definition~\ref{cv-toroidal-theta}(3). Then there
are a local homomorphism $\phi:V\to V'$ to a discrete valuation ring, an extension $\psi:K\to K'$ of
$\phi$ to the fraction field $K'$ of $V'$, a theta-level object $t'$ over $V'$ and a theta-level object
over $K'$ with morphisms to $t$ over $\operatorname{Spec}\psi$ and to $t'$ over
$\operatorname{Spec}K'\to\operatorname{Spec}V'$, such that $\phi(v)$ is a family of theta constants of
$t'$ at the identity.
\end{lemma}
\leanref{Development/AbelianVarietyTheory/Proofs/ModuliProofs/GoodReductionOfThetaPointUnits.lean}{GlobalTheory.ThetaCrit.good_reduction_of_theta_point}
\begin{proof}
Take $j_0=2$, let $j\ge2$ and $N'=n^j$, which is odd, and let $o$, $\zeta$, $V$, $K$, $\beta$, $t$ and
$v$ be given. The theta-level objects below are those of Situation~\ref{cv-toroidal-reduction-situation}
for $g$, $N'$, $I$, $o$ and $\zeta$.

The model. Lemma~\ref{cv-toroidal-reduction-semistable}, applied to $A$ with $L$, gives a complete
discrete valuation ring $W$ with a local homomorphism $\phi:V\to W$, an extension $\psi:K\to K'$ of $\phi$
to the fraction field $K'$ of $W$, and a semi-abelian scheme $G$ over $W$ whose closed fibre is an
extension of an abelian variety by a split torus, with an isomorphism of group schemes
$\gamma:G_{K'}\to A_{K'}$. By Lemma~\ref{cv-toroidal-reduction-base-change}(1) there is a theta-level
object $t_{K'}$ over $K'$, regarded over $o$ through $\beta$, $\phi$ and $W\to K'$, with a morphism
$t_{K'}\to t$ over $\operatorname{Spec}\psi$; its abelian scheme is $A_{K'}$, with the pullback $L_{K'}$
of $L$. As $\psi$ is injective, $K'$ and $W$ have characteristic zero.

The degeneration data. By Lemma~\ref{cv-toroidal-reduction-base-change}(2), the endomorphisms of
$A_{K'}$ given by $w$ and $w'$ satisfy the hypotheses on $w$ and $w'$ in
Theorem~\ref{cv-toroidal-mumford}(1). So that theorem, applied to $W$, $A_{K'}$ with $L_{K'}$, $G$ and
$\gamma$, gives degeneration data over $W$, with free $\mathbf Z$-module $X$, and a uniformization of
$(A_{K'},G)$ by them.

The theta constants over $K'$. By Lemma~\ref{cv-toroidal-reduction-base-change}(4), applied to $\psi$,
$\psi(v)$ is a family of theta constants of $t$ at $\operatorname{Spec}\psi$, so by
Lemma~\ref{cv-toroidal-reduction-base-change}(3) it is a family of theta constants of $t_{K'}$ at the
identity. It is the image in $K'$ of $\phi(v)$, since $\psi$ extends $\phi$. As $\phi$ is local, it maps
units of $V$ to units of $W$; so for every $c$ there is $c'$ with $\phi(v_{c+N'_nc'})$ a unit of $W$,
where $N'_n=n^{j-1}$. The integer $N'$ is invertible in $o$, hence in $W$.

The reduction is good. Lemma~\ref{cv-toroidal-reduction-torus}, applied to $W$, $t_{K'}$, $G$, the
degeneration data and $\phi(v)$, gives $X=0$. So Theorem~\ref{cv-toroidal-reduction-abelian} gives an
abelian scheme $P$ over $W$ with a line bundle $M$ as there, and a morphism $\chi:A_{K'}\to P$ over
$\operatorname{Spec}K'\to\operatorname{Spec}W$ that identifies $A_{K'}$ with $P_{K'}$ as group schemes,
with $L_{K'}$ isomorphic to $\chi^*M$.

The theta-level object over $W$. Lemma~\ref{cv-toroidal-reduction-model}, applied to $W$, $\phi\circ\beta$,
$t_{K'}$, $P$ and $\chi$, gives a theta-level object $t'$ over $W$, regarded over $o$ through
$\phi\circ\beta$, with a morphism $t_{K'}\to t'$ over $\operatorname{Spec}K'\to\operatorname{Spec}W$. By
the unit condition for $c=0$, some $\phi(v_c)$ is a unit of $W$, and $n\ge2$ divides $N'$. So
Lemma~\ref{cv-toroidal-reduction-generic}, applied to $t'$, $t_{K'}$ and $\phi(v)$, shows that $\phi(v)$
is a family of theta constants of $t'$ at the identity. With $V'=W$ and the object $t_{K'}$ over $K'$,
this proves the lemma.
\end{proof}
% Lean: good_reduction_of_theta_point with j0 = 2: exists_model (SemistablePol, here
% Lemma cv-toroidal-reduction-semistable, and tObj_pull), shape_prim_pull_and_descend and HypS, ThmUSt,
% hasTh_baseChange with hasTh_iff_of_tmap (hvK), the local map phi (hC'), subsingleton_characters,
% exists_pol_model, exists_obj_model, and hasTh_of_generic with hC' 0. The one step left open is the
% proof of Lemma cv-toroidal-reduction-torus.

\subsection{PEL structures on theta-level objects}\label{subsection-cv-pel-structures}

\begin{definition}\label{cv-toroidal-theta-family-good}
In the setting of Definition~\ref{cv-toroidal-theta}, let $\mathrm{em}:o\to\mathbf C$ be a
homomorphism, and let $\mathcal C$ be a class of theta-level objects over $\operatorname{Spec}\mathbf C$, regarded
over $o$ through $\mathrm{em}$. Let $(M_h)_{h\in I}$ be integer matrices of size $2g$, let $\kappa\ge0$ and
$\delta\ge0$ be integers, and let $w,w'\in I$. The class $\mathcal C$ is \emph{good} for
$((M_h),\kappa,w,w',\delta)$ if:
\begin{enumerate}
\item a theta-level object over $\operatorname{Spec}\mathbf C$ with a morphism over the identity to a member of
  $\mathcal C$ belongs to $\mathcal C$;
\item every member $t=(A,L,(x_i),(\alpha_h))$ of $\mathcal C$ satisfies $\alpha_h(x_i)=\sum_{i'}(M_h)_{i'i}x_{i'}$
  for all $h\in I$ and all $i$, and $L^\kappa\otimes\alpha_h^*L^{-1}$ is ample for every $h$; moreover
  $\alpha_w+\alpha_{w'}=1$, $L$ is isomorphic to the pullback of $\Lambda(L)$ along $(1,\alpha_w)$, and the
  pullbacks of $\Lambda(L)$ along $\alpha_w\times1$ and $1\times\alpha_{w'}$ are isomorphic, with $\Lambda(L)$ as
  in Theorem~\ref{cv-toroidal-mumford};
\item there are open subsets $B_1,\dots,B_r$ of $\mathbf C^\delta$, each connected or empty, and holomorphic maps
  $\vartheta_k:B_k\to\mathbf C^{(\mathbf Z/N'^2)^{2g}}$ such that for every $k$ and every $z\in B_k$,
  $\vartheta_k(z)$ is a family of theta constants at the identity of a member of $\mathcal C$, and only countably
  many $z'\in B_k$ have $\vartheta_k(z')\in\mathbf C\,\vartheta_k(z)$; and every member of $\mathcal C$ has a
  family of theta constants at the identity of the form $\vartheta_k(z)$ with $z\in B_k$.
\end{enumerate}
\end{definition}
\leanref{Development/AbelianVarietyTheory/Definitions/ModuliSpaces/GoodClass.lean}{GlobalTheory.ZM.GoodClass}
\leanref{Development/AbelianVarietyTheory/Definitions/ModuliSpaces/GoodClass.lean}{GlobalTheory.ZM.TObj.Shape}
\leanref{Development/AbelianVarietyTheory/Definitions/ModuliSpaces/GoodClass.lean}{GlobalTheory.ZM.TObj.Prim}

\begin{lemma}[Alp\"oge]\label{cv-toroidal-theta-family-pel-good}
In Situation~\ref{cv-function-situation}, suppose $p+q\ge1$, let $g=\frac12[E:\mathbf Q](p+q)$ and
$\delta=\frac12[E:\mathbf Q]pq$, let $e_1,\dots,e_{2g}$ be a symplectic basis of the lattice $\Lambda$ of the datum,
and let $a:I\to\mathcal O_E$ be a map from a finite set whose image generates $\mathcal O_E$ as a ring, with
elements $w,w'\in I$ such that $a(w)=w_0=(1+\sqrt{-d})/2$ and $a(w')=\bar w_0$; the hypothesis that the image of
$a$ generates $\mathcal O_E$ is not used in the proof. Then there are integer matrices $M_h$, $h\in I$, of size
$2g$ and an integer $\kappa\ge0$ such that for every odd prime $n$ there is $j_0$ with the following property: for
every $j\ge j_0$, with $N'=n^j$, every ring $o$ with a unit $\zeta$ as in Definition~\ref{cv-toroidal-theta}, every
homomorphism $\mathrm{em}:o\to\mathbf C$ and every $\rho\in o$ with $\mathrm{em}(\rho)=\theta$, the PEL class of
$((e_i),a,\rho,\mathrm{em})$ is good for $((M_h),\kappa,w,w',\delta)$
(Definition~\ref{cv-toroidal-theta-family-good}).
\end{lemma}
\leanref{Development/AbelianVarietyTheory/Proofs/ModuliProofs/IsGoodThetaComplexClassUnitaryFullLevelStatement.lean}{GlobalTheory.T.isGoodThetaComplexClass_unitaryFullLevel}
\leanref{Development/AbelianVarietyTheory/Proofs/ModuliProofs/IsGoodThetaComplexClassUnitaryFullLevel.lean}{GlobalTheory.PelGood.goodClass_pelClass}
\begin{proof}
Let $M_h$ be the matrix of multiplication by $a(h)$ in the basis $(e_i)$, so that
$a(h)e_i=\sum_{i'}(M_h)_{i'i}e_{i'}$, and let $\kappa$ be an integer with $|\tau(a(h))|^2<\kappa$ for every $h\in I$
and every embedding $\tau:E\to\mathbf C$; it exists because both sets are finite.
\openstep{That the PEL class is good for these data once $j$ is large, from the complex uniformization of its
members by lattices with Hermitian forms, the finitely many marked lattices in the genus of $\Lambda$ and the
domains of their period maps, is not yet written out.}
% Lean: PelGood.goodClass_pelClass: Mat j i' i = bas.symm (a j • bas (single i 1)) i' and kappa from the finite
% range of ‖tau (a j)‖^2 (the two steps written above); then exists_isTypeFn_of_rank (ty),
% T.goodThetaClass_of_lattice_families (j0), MarkedGenus.finite_marked_data_in_genus (the marked data Di, fi,
% nu_i), exists_domain (the domains B i and the period maps P i), pelClass_prim_and_iso (part (1), and the
% conditions on w and w' in part (2)), ConjBound.exists_action_lt (the ampleness of L^kappa (alpha_h^* L)^{-1}),
% T.isPELObject_of_primitive_periodMap_uniformised, T.exists_periods_complexPELObject and
% UnifTransport.transport_uniformisation (part (3)). It takes Out.Laws, Out.Laws2, the theorem of the cube
% (T.theorem_of_cube_polAbelianScheme) and ThetaUnifAV (T.theta_uniformisation_complex_abelian_varieties) as
% hypotheses. Sizes in this section are counted through the declarations reached in the imported files, beyond the
% laws of abelian schemes (abelianSchemeLaws_all, allSixLaws, firstFiveLaws): here about 69,600 lines, about 32,200
% of them not shared with the other inputs of Lemma cv-toroidal-theta-family.
\end{proof}

\begin{lemma}[Alp\"oge]\label{cv-toroidal-theta-family-exists}
Let $g\ge1$, let $I$ be a finite set with elements $w$ and $w'$, let $(M_h)_{h\in I}$ be integer matrices of size
$2g$, let $\kappa\ge0$ be an integer, and let $n$ be an odd prime. There is $j_0$ such that for every $j\ge j_0$,
with $N'=n^j$, the following holds. Let $o$ be a ring with a unit $\zeta$ as in Definition~\ref{cv-toroidal-theta},
let $\mathrm{em}:o\to\mathbf C$ be an injective homomorphism such that $o$ is a Dedekind domain of finite type over
$\mathbf Z$, let $\delta\ge0$, and let $\mathcal C$ be a class of theta-level objects over $\operatorname{Spec}\mathbf C$
that is good for $((M_h),\kappa,w,w',\delta)$ (Definition~\ref{cv-toroidal-theta-family-good}). Then there is a
theta family for $\mathcal C$ and $n$ (Definition~\ref{cv-toroidal-theta}(5)).
\end{lemma}
\leanref{Development/AbelianVarietyTheory/Proofs/ModuliProofs/NonemptyThetaObjectFamilyOpenClosureStatement.lean}{GlobalTheory.T.nonempty_thetaObject_family_open_closure}
\leanref{Development/AbelianVarietyTheory/Proofs/ModuliProofs/NonemptyThetaObjectFamilyOpenClosure.lean}{GlobalTheory.ZM.ThetaFamily.nonempty_thetaFam}
\begin{proof}
\openstep{The construction of the theta family, as the open part of the scheme cut out by the theta ideal of
$\mathcal C$ on which the coset condition holds, with its universal theta-level object, and the proof of
properties (a) to (e) of Definition~\ref{cv-toroidal-theta}(5) for it, with
Lemma~\ref{cv-toroidal-theta-reduction} as an input, are not yet written out.}
% Lean: ThetaFamily.nonempty_thetaFam, with hcrit (the statement of Lemma cv-toroidal-theta-reduction) as a
% hypothesis: T.thetaPointsVanishingIdeal_radical_homogeneous (the theta ideal is radical, homogeneous and
% o-torsion-free), ZM.nonempty_MMPre (the scheme Y over o with its coordinates),
% T.exists_thetaFamily_over_thetaPointScheme (hglue, which uses hcrit: the universal object t over Y, with dense
% C-points and the coordinates as theta constants), T.thetaIdealFamily_universal_property (property (d)),
% T.tangentVector_zero_of_object_constant_along (property (e)), exists_coord_of_isLocalRing with Y.loc (property
% (b)) and Y.pts (property (c)); j0 = max j1 j2. Size: about 45,100 lines, about 10,100 of them not shared with the
% other inputs of Lemma cv-toroidal-theta-family.
\end{proof}

\begin{lemma}[Alp\"oge]\label{cv-toroidal-theta-family-pel-fibre-line}
Let $i:K_0\to K$ be a homomorphism of fields, let $X_0$ be a scheme over $K_0$ that is smooth, proper and geometrically
connected, and let $\varphi:X\to X_0$ be its base change to $K$. A line bundle $M$ on $X_0$ with $\varphi^*M$ trivial is
trivial, and two line bundles $M,N$ on $X_0$ with $\varphi^*M\cong\varphi^*N$ are isomorphic.
\end{lemma}
\leanref{Development/AbelianVarietyTheory/Proofs/ModuliProofs/LineBundleTrivialIsoDescendFieldStatement.lean}{GlobalTheory.T.lineBundle_trivial_iso_descend_field}
\begin{proof}
Let $t$ be a generator of $\varphi^*M$. The scheme $X$ is smooth, proper and geometrically connected over $K$, so its
global functions are the constants $K$ \cite[Tag 0E0L]{Sta}. By flat base change \cite[Tag 02KH]{Sta}, $t$ is a
$K$-linear combination of pullbacks $\varphi^*s_a$ of sections $s_a$ of $M$ over $X_0$. Each $\varphi^*s_a$ is $u_at$ for a
global function $u_a\in K$, and some $u_a$ is nonzero, since otherwise $t=0$. Then $\varphi^*s_a$ generates
$\varphi^*M$, so $s_a$ vanishes at no point of $X_0$, as $\varphi$ is surjective. So $s_a$ generates $M$. For the second
claim apply the first to $N\otimes M^{-1}$.
\end{proof}
% Lean: lb_gen_iso_descend_field: FieldDescentGen.exists_gen (bijective_appTop_of_smooth_of_isProper_of_geometrically-
% Connected, SectionsDescent.main for the sections, exists_ratio, isUnit_of_surjective), applied to LB.div N M.

\begin{lemma}[Alp\"oge]\label{cv-toroidal-theta-family-pel-fibre-descent}
In Situation~\ref{cv-function-situation}, let $e_1,\dots,e_{2g}$ be a basis of the lattice $\Lambda$ of the datum over
$\mathbf Z$, let $a:I\to\mathcal O_E$ be a map whose image generates $\mathcal O_E$ as a ring, let $N'\ge1$, let $o$ be a
ring with a unit $\zeta$ as in Definition~\ref{cv-toroidal-theta}, and let $\rho\in o$. Let $i:K_0\to K$ be a
homomorphism of fields and $\beta_0:o\to K_0$ a homomorphism, let $t_0$ be a theta-level object over $K_0$, regarded over
$o$ through $\beta_0$, and let $t$ be a theta-level object over $K$, regarded over $o$ through $i\circ\beta_0$, with a
morphism $t\to t_0$ over $\operatorname{Spec}i$. If $t$ has a PEL structure for $((e_i),a,\rho)$
(Definition~\ref{cv-toroidal-theta}(6)), then so has $t_0$.
\end{lemma}
\leanref{Development/AbelianVarietyTheory/Proofs/ModuliProofs/PelStructureDescendsFieldExtensionStatement.lean}{GlobalTheory.T.pelStructure_descends_field_extension}
\leanref{Development/AbelianVarietyTheory/Proofs/ModuliProofs/PelStructureDescendsFieldExtension.lean}{GlobalTheory.PELDesc.isPEL_descends}
\begin{proof}
Write $t_0=(A_0,L_0,(x_{0,j}),(\alpha_{0,h}))$ and $t=(A,L,(x_j),(\alpha_h))$, let $\psi:A\to A_0$ be the morphism of
$t\to t_0$, and let $G$ be the PEL structure on $t$, with line bundle $\mathcal L$ and isomorphism $\varphi:G\to A$. The
composite $\chi=\psi\circ\varphi$ identifies $G$ with the base change of $A_0$ to $K$, it carries $\mathcal L$ to the
pullback of $L_0$, and $\chi\circ a(h)=\alpha_{0,h}\circ\chi$ for every $h$.

The action. The morphism $\chi$ is a base change of the faithfully flat and quasi-compact morphism
$\operatorname{Spec}K\to\operatorname{Spec}K_0$, so it is an epimorphism of schemes: two morphisms from $A_0$ that agree
after composing with $\chi$ are equal. The elements $c\in\mathcal O_E$ for which there is an endomorphism $r(c)$ of the
group scheme $A_0$ with $\chi\circ c=r(c)\circ\chi$ form a subring: if $\chi\circ c=f\circ\chi$ and
$\chi\circ c'=f'\circ\chi$, then $\chi\circ(c+c')=(f+f')\circ\chi$ and $\chi\circ(-c)=(-f)\circ\chi$, as $\chi$ is a
homomorphism over $\operatorname{Spec}i$, and $\chi\circ cc'=ff'\circ\chi$. It contains the $a(h)$, with
$r(a(h))=\alpha_{0,h}$, so it is $\mathcal O_E$. As $\chi$ is an epimorphism, $r(c)$ is unique, and the identities
$r(1)=1$, $r(cc')=r(c)r(c')$ and $r(c+c')=r(c)+r(c')$ hold, because they hold after composing with $\chi$. So $r$ is an
action of $\mathcal O_E$ on $A_0$, and $\chi$ carries the action on $G$ to $r$.

The $\mathcal O_E$-semi-abelian scheme. The abelian varieties $A_0$, $A_0\times_{K_0}A_0$ and $A_0^3$ are smooth, proper
and geometrically connected over $K_0$, and through $\chi$ their base changes to $K$ are $G$, $G\times_KG$ and $G^3$.
Write $\Lambda$ and $\Theta$ as in Definition~\ref{cv-germs-chart}(1), and $y$, $y_1,y_2$ for the identity of a group and
the projections of its square. The cubical structure of $\mathcal L$ is a generator of $\Theta$ for $\mathcal L$ on
$G^3$, which is the pullback along $\chi^3$ of $\Theta$ for $L_0$ on $A_0^3$; by
Lemma~\ref{cv-toroidal-theta-family-pel-fibre-line}, the latter is trivial. For $c\in\mathcal O_E$ the isomorphism
$\Lambda_{\mathcal L}(cy_1,y_2)\cong\Lambda_{\mathcal L}(y_1,\bar cy_2)$ on $G\times_KG$ is an isomorphism between the
pullbacks along $\chi\times\chi$ of $\Lambda_{L_0}(r(c)y_1,y_2)$ and $\Lambda_{L_0}(y_1,r(\bar c)y_2)$, so these are
isomorphic on $A_0\times_{K_0}A_0$ by the same lemma. Likewise $\mathcal L\cong\Lambda_{\mathcal L}(y,w_0y)$ on $G$
gives $L_0\cong\Lambda_{L_0}(y,r(w_0)y)$ on $A_0$. Lemma~\ref{cv-toroidal-theta-family-pel-cubical}, applied to $A_0$,
$r$ and $L_0$, gives a symmetric line bundle $P$ on $A_0$, trivialized along the zero section, with a cubical structure,
such that $P\cong\Lambda_{L_0}(y,r(w_0)y)\cong L_0$, $P\cong\Lambda_P(y,r(w_0)y)$ and
$\Lambda_P(r(c)y_1,y_2)\cong\Lambda_P(y_1,r(\bar c)y_2)$ for every $c$. So $A_0$, an abelian and hence semi-abelian
scheme, with $r$ and $P$ is an $\mathcal O_E$-semi-abelian scheme $G_0$ over $K_0$
(Definition~\ref{cv-model-toroidal-semiabelian}(2)). The identity $G_0\to A_0$ is an isomorphism of group schemes that
carries $P$ to $L_0$, up to isomorphism, and $r(a(h))$ to $\alpha_{0,h}$.

The conditions at geometric points. Let $k$ be an algebraically closed field of characteristic zero and
$x:\operatorname{Spec}k\to\operatorname{Spec}K_0$ a point, given by $\kappa:K_0\to k$, with $\beta=\kappa\circ\beta_0$.
The ring $K\otimes_{K_0}k$ is nonzero; let $\Omega$ be an algebraic closure of its quotient by a maximal ideal. It is an
algebraically closed field of characteristic zero, with embeddings $\iota_K:K\to\Omega$ and $\iota_k:k\to\Omega$ such
that $\iota_K\circ i=\iota_k\circ\kappa$. Let $y:\operatorname{Spec}\Omega\to\operatorname{Spec}K$ be the point of
$\iota_K$; its homomorphism from $o$ is $\iota_K\circ i\circ\beta_0=\iota_k\circ\beta$. The morphism $\chi$ identifies
the fibre $G_y$ with the base change of $A_0$ along $\kappa$ and then $\iota_k$, compatibly with the actions of
$\mathcal O_E$ and the line bundles, hence with the eigenspaces of the Lie algebras and the commutator pairings of the
theta groups. As $G$ is a PEL structure on $t$, condition (a) of Definition~\ref{cv-toroidal-family}(1) holds for $G_0$
at $x\circ\operatorname{Spec}\iota_k$ with $\iota_k(\beta(\rho))$, and there are roots of unity $\zeta'_m\in\Omega$ and
maps $\iota_m$ into the points of $A_0$ over $\Omega$ as in condition (b) there, with
$\zeta'_{N'^2}=\iota_k(\beta(\zeta))$ and $\chi(\iota^G_{N'^2}(e_j))=\psi(x_j(y))$, the value of $x_{0,j}$ at
$x\circ\operatorname{Spec}\iota_k$. These data are now descended from $\Omega$ to $k$. Write $A_{0,x}$ for the fibre of
$A_0$ at $x$, an abelian variety over $k$ of dimension $g$, and $L_x$ for the restriction of $L_0$ to it; as $P\cong L_0$,
the commutator pairings of the theta groups of the powers of $P$ are those of $L_0$.

Condition (a) at $x$. By Lemma~\ref{cv-origin-reference-lie-bundle-base}, applied to $A_0$ and $\iota_k$, composition with
$\operatorname{Spec}\iota_k[\varepsilon]$ identifies $\Omega\otimes_k\mathfrak g$, where $\mathfrak g$ is the Lie algebra
of $A_0$ at $x$, with the Lie algebra at $x\circ\operatorname{Spec}\iota_k$, and it commutes with the action of
$\mathcal O_E$. So for a ring homomorphism $\tau:\mathcal O_E\to k$ the subspace on which each $c$ acts by
$\iota_k(\tau(c))$ has the dimension of the subspace of $\mathfrak g$ on which each $c$ acts by $\tau(c)$. Every ring
homomorphism $\mathcal O_E\to\Omega$ is $\iota_k\circ\tau$ for some $\tau$, as its values are integral over $\mathbf Z$
and $k$ is algebraically closed; and $\iota_k(\tau(\sqrt{-d}))=\iota_k(\beta(\rho))$ if and only if
$\tau(\sqrt{-d})=\beta(\rho)$. So condition (a) holds at $x$ with $\beta(\rho)$.

The torsion points. Composition with $\operatorname{Spec}\iota_k$ is an injective homomorphism $\pi$ from $A_{0,x}(k)$ to
the group of points of $A_0$ over $x\circ\operatorname{Spec}\iota_k$, as $\operatorname{Spec}\Omega\to\operatorname{Spec}k$
is faithfully flat, hence an epimorphism of schemes; and $\pi$ commutes with the action of $\mathcal O_E$. For $m\ge1$ the
points killed by $m$ are $m^{2g}$ in number on both sides, as $k$ and $\Omega$ are algebraically closed of characteristic
zero \cite[\S6, Application~3]{Mum70}. So $\pi$ maps the points of $A_{0,x}(k)$ killed by $m$ bijectively onto those
over $\Omega$.

The pairings. Let $m\ge1$, let $z,z'\in A_{0,x}(k)$ be killed by $m$, and let $c\in k^\times$. Then $e_m(z,z')=c$, for
the commutator pairing $e_m$ of the theta group of $L_x^m$, if and only if $e_m(\pi(z),\pi(z'))=\iota_k(c)$. A lift of $z$
to the theta group is a generator of $t_z^*L_x^m\otimes L_x^{-m}$, where $t_z$ is the translation by $z$, and
$e_m(z,z')=c$ says that $z$ and $z'$ have lifts and that the commutator of any two lifts is $c$, an equality between two
generators of a line bundle on $A_{0,x}$. Lifts and commutators pull back along $\operatorname{Spec}\iota_k$, which
gives one implication. Conversely, let $e_m(\pi(z),\pi(z'))=\iota_k(c)$. Then $\pi(z)$ and $\pi(z')$ have lifts, so the
line bundles $t_z^*L_x^m\otimes L_x^{-m}$ and $t_{z'}^*L_x^m\otimes L_x^{-m}$ on $A_{0,x}$, which is smooth, proper and
geometrically connected over $k$, become trivial over $\Omega$; by Lemma~\ref{cv-toroidal-theta-family-pel-fibre-line}
they are trivial, and $z$ and $z'$ have lifts. Let two lifts be given. Their pullbacks are lifts of $\pi(z)$ and
$\pi(z')$, whose commutator is $\iota_k(c)$, so the two generators that the relation compares become equal over
$\Omega$. Pulling back sections along the faithfully flat morphism $A_{0,x}\times_k\operatorname{Spec}\Omega\to A_{0,x}$
is injective; so the two generators are equal, and $e_m(z,z')=c$.

Condition (b) at $x$. As $k$ is algebraically closed of characteristic zero, every root of unity of order $m$ in
$\Omega$ is $\iota_k$ of a root of unity of order $m$ in $k$. So $\zeta'_m=\iota_k(\zeta''_m)$ for primitive roots of
unity $\zeta''_m\in k$ of order $m$, and, by the bijection on torsion points, $\iota_m=\pi\circ\iota''_m$ for maps
$\iota''_m:\Lambda\to A_{0,x}(k)$. As $\iota_k$ and $\pi$ are injective homomorphisms, $\pi$ commutes with the action of
$\mathcal O_E$ and is bijective on the points killed by each $m$, and the pairings correspond, the $\zeta''_m$ and
$\iota''_m$ satisfy condition (b) of Definition~\ref{cv-toroidal-family}(1) at $x$, with $\zeta''_{N'^2}=\beta(\zeta)$
and $\iota''_{N'^2}(e_j)=x_{0,j}(x)$. Hence $G_0$ with the identity of $A_0$ is a PEL structure on $t_0$.
\end{proof}
% Lean: PELDesc.isPEL_descends (with Out.Laws): hpb and hLc (chi = phi ≫ psi), oAct_descend_of_generators (the action;
% ActDescent.exists_oAct and over_hom_ext), exists_triple_of_fieldExtension (T.lineBundle_trivial_iso_descend_field on
% A0, A0^2 and A0^3, and PrimNorm.exists_primitive_normalisation), exists_algClosed_amalgam_of_charZero (Omega),
% detCond_pull and detCond_iff_of_field_extension (condition (a): DetExtend.lieMult_extend through universal charts,
% here Lemma cv-origin-reference-lie-bundle-base, and DetExtend.exists_factor), tate_pull and
% TateIsoFieldExt.tateIso_fieldExt (condition (b)): T.torsion_points_descend_field_extension (comp_injective, and
% card_torsion from the law Out.Laws.mulN, the count cited to Mumford), T.weilPairing_fieldExtension_iff =
% weilIs_fieldExt_iff (weilIs_pull; weilIs_iff, exists_lift_descend with T.lineBundle_trivial_iso_descend_field,
% commIsOf_descend with corr_of_val_epi), exists_unit_of_pow_eq_one and TateIsoFieldExt.down; then tate_pull along the
% isomorphism of G_0 with A_0. The pairing step is about 1,900 of the 2,090 lines of tateIso_fieldExt
% (tools/lean_size.py), mostly the theta group's lifts and their pullbacks.

\begin{lemma}[Alp\"oge]\label{cv-toroidal-theta-family-pel-fibre-pull}
In Situation~\ref{cv-function-situation}, let $e_1,\dots,e_{2g}$ be a basis of the lattice $\Lambda$ of the datum over
$\mathbf Z$, let $a:I\to\mathcal O_E$ be a map, let $N'\ge1$, let $o$ be a ring with a unit $\zeta$ as in
Definition~\ref{cv-toroidal-theta}, and let $\rho\in o$. Let $f:S'\to S$ be a morphism of schemes over $o$, let $t$ be a
theta-level object over $S$ with a PEL structure for $((e_i),a,\rho)$ (Definition~\ref{cv-toroidal-theta}(6)), and let
$t'$ be a theta-level object over $S'$ with a morphism $t'\to t$ over $f$. Then $t'$ has a PEL structure for
$((e_i),a,\rho)$.
\end{lemma}
\leanref{Development/AbelianVarietyTheory/Proofs/ModuliProofs/IsPELPull.lean}{GlobalTheory.PELPull.isPEL_pull}
\begin{proof}
Write $t=(A,L,(x_j),(\alpha_h))$ and $t'=(A',L',(x'_j),(\alpha'_h))$, let $\psi:A'\to A$ be the morphism of $t'\to t$,
and let $G$ be the PEL structure on $t$, with line bundle $\mathcal L$ and isomorphism $\varphi:G\to A$. By
Lemma~\ref{cv-origin-reference-pullback}(1) the pullback $G'=f^*G$, with the pullbacks of the action and of the line
bundle $\mathcal L'$, is an $\mathcal O_E$-semi-abelian scheme over $S'$, and the projection $\nu:G'\to G$ is a morphism
over $f$. Both $\varphi\circ\nu$ and $\psi$ identify their sources with $f^*A$, so there is a unique isomorphism of group
schemes $\eta:G'\to A'$ over $S'$ with $\psi\circ\eta=\varphi\circ\nu$. A morphism to $A'$ over $S'$ is determined by its
composite with $\psi$.

The line bundle $\mathcal L'$ is the pullback of $\mathcal L$ along $\nu$, hence of $L$ along $\varphi\circ\nu=\psi\circ
\eta$, and $L'$ is the pullback of $L$ along $\psi$; so $\eta$ carries $\mathcal L'$ to $L'$. For $h\in I$,
$\psi\circ\eta\circ a(h)=\varphi\circ a(h)\circ\nu=\alpha_h\circ\psi\circ\eta=\psi\circ\alpha'_h\circ\eta$, so
$\eta\circ a(h)=\alpha'_h\circ\eta$.

Let $k$ be an algebraically closed field of characteristic zero and $x:\operatorname{Spec}k\to S'$ a point; its
homomorphism $\beta:o\to k$ is also that of $f\circ x$. The morphism $\nu$ identifies the fibre $G'_x$ with $G_{f\circ
x}$, compatibly with the actions of $\mathcal O_E$ and the line bundles, hence with the eigenspaces of the Lie algebras and
with the commutator pairings of the theta groups. So condition (a) of Definition~\ref{cv-toroidal-family}(1) for $G'$ at
$x$, with $\beta(\rho)$, is that for $G$ at $f\circ x$; and the roots of unity $\zeta'_m$ and the maps $\iota_m$ of
condition (b) for $G$ at $f\circ x$, composed with the inverse of this identification, satisfy condition (b) for $G'$ at
$x$, with $\zeta'_{N'^2}=\beta(\zeta)$. They satisfy $\eta(\iota_{N'^2}(e_j))=x'_j(x)$, because
$\psi(\eta(\iota_{N'^2}(e_j)))=\varphi(\iota^G_{N'^2}(e_j))=x_j(f\circ x)=\psi(x'_j(x))$, as $x'_j$ is the pullback of
$x_j$. So $G'$ with $\eta$ is a PEL structure on $t'$.
\end{proof}
% Lean: PELPull.isPEL_pull: sabPull with Out.CubePf.cubePull (G' and nu), grpPullback_iso (eta), m.lin and lbIso_pull (the
% line bundle), hψ.1.hom_ext with m.act and hφa (the actions), LiePull.lieMult_pull (condition (a)), ObjPull.tateMove with
% ObjPull.pushEquiv and weilIs_grpPullback_of_comm (condition (b)).

\begin{lemma}[Alp\"oge]\label{cv-toroidal-theta-family-pel-fibre}
In Situation~\ref{cv-function-situation}, let $e_1,\dots,e_{2g}$ be a basis of the lattice $\Lambda$ of the datum over
$\mathbf Z$, let $a:I\to\mathcal O_E$ be a map whose image generates $\mathcal O_E$ as a ring, let $N'\ge1$, let $o$ be a
ring with a unit $\zeta$ as in Definition~\ref{cv-toroidal-theta}, let $\mathrm{em}:o\to\mathbf C$ be a homomorphism and
let $\rho\in o$. Let $S$ be a scheme over $o$, let $t$ be a theta-level object over $S$, and let $p$ be a point of $S$ at
which $t$ is a $\mathcal C$-object (Definition~\ref{cv-toroidal-theta}(4)), for the PEL class $\mathcal C$ of
$((e_i),a,\rho,\mathrm{em})$ (Definition~\ref{cv-toroidal-theta}(6)). Then there are a field $K$ of characteristic zero, a point $x:\operatorname{Spec}K\to S$
whose image is $p$, and a theta-level object $t_K$ over $K$ with a morphism $t_K\to t$ over $x$ and a PEL structure for
$((e_i),a,\rho)$ (Definition~\ref{cv-toroidal-theta}(6)).
\end{lemma}
\leanref{Development/AbelianVarietyTheory/Proofs/ModuliProofs/ExistsIsPELTripleOfDenseFibres.lean}{GlobalTheory.PELFamily.exists_fibre}
\leanref{Development/AbelianVarietyTheory/Proofs/ModuliProofs/ExistsIsPELTripleOfDenseFibres.lean}{GlobalTheory.PELFamily.exists_fibre_map}
\begin{proof}
By Definition~\ref{cv-toroidal-theta}(4) there are a field $K$, a point $x:\operatorname{Spec}K\to S$ with image $p$,
a field $K_0$ with homomorphisms $\beta_0:o\to K_0$, $i:K_0\to K$ and $j:K_0\to\mathbf C$, with $j\circ\beta_0=\mathrm{em}$
and $i\circ\beta_0$ the homomorphism of $x$, a theta-level object $t_0$ over $K_0$, a theta-level object $t_K$ over $K$
with morphisms $t_K\to t$ over $x$ and $t_K\to t_0$ over $\operatorname{Spec}i$, and a member $t_{\mathbf C}$ of
$\mathcal C$ with a morphism $t_{\mathbf C}\to t_0$ over $\operatorname{Spec}j$. The member $t_{\mathbf C}$ has a PEL
structure. The field $K_0$ embeds in $\mathbf C$ through $j$, so it has characteristic zero, and so has $K$.

Lemma~\ref{cv-toroidal-theta-family-pel-fibre-descent}, applied to $j$, $\beta_0$, $t_0$ and $t_{\mathbf C}$ with its PEL
structure, gives a PEL structure on $t_0$; it uses that the image of $a$ generates $\mathcal O_E$. Then
Lemma~\ref{cv-toroidal-theta-family-pel-fibre-pull}, applied to $\operatorname{Spec}i$, $t_0$ and $t_K$ with its morphism
to $t_0$, gives a PEL structure on $t_K$. With $x$ and the morphism $t_K\to t$, this proves the lemma.
\end{proof}
% Lean: exists_fibre (with Out.Laws): the data of fromSet, T.pelStructure_descends_field_extension (the descent from C
% to K0) and PELPull.isPEL_pull (the pullback to K); CharZero K from j.charZero and i. exists_fibre_map composes the
% isomorphism of the PEL structure with the morphism tx -> t.

\begin{lemma}[Alp\"oge]\label{cv-toroidal-theta-family-pel-action}
Let $O$ be a commutative ring, let $a:I\to O$ be a map whose image generates $O$ as a ring, let $S$ be a reduced
scheme, let $A$ be an abelian scheme over $S$, and let $\alpha_h$, $h\in I$, be endomorphisms of the group scheme $A$. Suppose that the
points $p$ of $S$ with the following property are dense: there are a field $K$, a point $x:\operatorname{Spec}K\to S$
whose image contains $p$, a group scheme $G$ over $K$ with a morphism $\varphi:G\to A$ over $x$ that identifies $G$
with $x^*A$, and an action of $O$ on $G$ under which $\varphi\circ a(h)=\alpha_h\circ\varphi$ for every $h$. Then there
is an action $r$ of $O$ on $A$ with $r(a(h))=\alpha_h$ for every $h$, such that for every morphism $x:T\to S$, every
group scheme $G$ over $T$ with a morphism $\varphi:G\to A$ over $x$ that identifies $G$ with $x^*A$, and every action of
$O$ on $G$ with $\varphi\circ a(h)=\alpha_h\circ\varphi$ for every $h$, $\varphi\circ c=r(c)\circ\varphi$ for every
$c\in O$.
\end{lemma}
\leanref{Development/AbelianVarietyTheory/Proofs/ModuliProofs/OActSpreadFromDenseFibresAbelianSchemeStatement.lean}{GlobalTheory.T.oAct_spread_from_dense_fibres_abelianScheme}
\begin{proof}
Call an endomorphism $f$ of the group scheme $A$ \emph{compatible} with $c\in O$ if $\varphi\circ c=f\circ\varphi$ for
every morphism $x:T\to S$, every group scheme $G$ over $T$ with a morphism $\varphi:G\to A$ over $x$ that identifies $G$
with $x^*A$, and every action of $O$ on $G$ with $\varphi\circ a(h)=\alpha_h\circ\varphi$ for every $h$.

Compatible endomorphisms. Let $x$, $G$ and $\varphi$ be as there, $u,v$ endomorphisms of $G$ and $f,f'$ endomorphisms of
$A$ with $\varphi\circ u=f\circ\varphi$ and $\varphi\circ v=f'\circ\varphi$. As $\varphi$ is a homomorphism of group
schemes over $x$, $\varphi\circ(u+v)=\varphi\circ u+\varphi\circ v=(f+f')\circ\varphi$ and $\varphi\circ(-u)=(-f)\circ
\varphi$; also $\varphi\circ uv=ff'\circ\varphi$, $\varphi\circ0=0\circ\varphi$ and $\varphi\circ1=1\circ\varphi$. So the
elements of $O$ with which some endomorphism of $A$ is compatible form a subring. By hypothesis it contains each $a(h)$,
with which $\alpha_h$ is compatible, so it is $O$. For each $c\in O$ choose an endomorphism $r(c)$ of $A$ compatible
with $c$.

Endomorphisms agreeing on the dense set. Let $\mathcal D$ be the dense set of the hypothesis, and let $f,f'$ be
endomorphisms of $A$ with $f\circ\varphi=f'\circ\varphi$ for all data $K$, $x$, $G$, $\varphi$ attached to the points of
$\mathcal D$. Then $f=f'$. Indeed, $A$ is smooth over the reduced scheme $S$, so $A$ is reduced. As $A$ is flat and
locally of finite presentation over $S$, the morphism $A\to S$ is open \cite[Th\'eor\`eme~2.4.6]{Gro65}, so the
preimage of $\mathcal D$ in $A$ is dense. Let $y$ be a point of $A$ over $p\in\mathcal D$, with $K$, $x$, $G$, $\varphi$
attached to $p$. As $p$ lies in the image of $x$ and $\varphi$ identifies $G$ with $x^*A$, $y=\varphi(z)$ for a point
$z$ of $G$, and $(f,f')(y)=(f\varphi(z),f'\varphi(z))$ lies on the diagonal of $A\times_SA$. The diagonal is a closed
subscheme, as $A$ is separated over $S$. So $(f,f'):A\to A\times_SA$ maps a dense set of points into it, and, as $A$ is
reduced, factors through it. Hence $f=f'$.

The action. For the data at the points of $\mathcal D$, compatibility gives $r(1)\circ\varphi=\varphi$,
$r(cc')\circ\varphi=\varphi\circ cc'=r(c)r(c')\circ\varphi$, $r(c+c')\circ\varphi=(r(c)+r(c'))\circ\varphi$ and
$r(a(h))\circ\varphi=\varphi\circ a(h)=\alpha_h\circ\varphi$. By the previous paragraph, $r(1)=1$, $r(cc')=r(c)r(c')$,
$r(c+c')=r(c)+r(c')$ and $r(a(h))=\alpha_h$. So $r$ is an action of $O$ on $A$ by endomorphisms of the group scheme, and
$\varphi\circ c=r(c)\circ\varphi$ for all data as in the statement, as $r(c)$ is compatible with $c$.
\end{proof}
% Lean: oAct_spread_of_dense_fibres: ActSpread.exists_universal (the compatible endomorphisms, by
% Subring.closure_induction with compat_mul, compat_inv, compat_comp, compat_one and compat_id), ActSpread.ext_of_dense
% (isReduced_of_smooth_of_isReduced, dense_fibres_ext: the preimage of a dense set under the open map A -> S, and
% exists_lift_closedImmersion_of_dense for the diagonal), and the four identities by ext.

\begin{lemma}[Alp\"oge]\label{cv-toroidal-theta-family-pel-bundles-seesaw}
Let $S_0$ be a scheme, let $B$ be an abelian scheme over $S_0$ with a relatively ample line bundle $M$, let $S$ be a
reduced scheme with a morphism $S\to S_0$, let $X=B\times_{S_0}S$ with projection $f:X\to S$, and let $N$ be a line bundle
on $X$. Suppose that the points $p$ of $S$ with the following property are dense: there are a field $K$ and a point
$x:\operatorname{Spec}K\to S$ whose image contains $p$ such that the pullback of $N$ to $X\times_S\operatorname{Spec}K$
is trivial. If $\sigma^*N$ is trivial for a section $\sigma$ of $f$, then $N$ is trivial.
\end{lemma}
\leanref{Development/AbelianVarietyTheory/Proofs/ModuliProofs/SeesawReducedBaseAbelianSchemeOfDenseFibresStatement.lean}{GlobalTheory.T.seesaw_reduced_base_abelianScheme_of_dense_fibres}
\begin{proof}
Local triviality. Every point of $S$ has an open neighbourhood $U$ such that $N$ is trivial on $f^{-1}(U)$.
\openstep{That $N$ is trivial over a neighbourhood of each point of the reduced base, from its triviality on a dense set
of fibres, is not yet written out.}
% Lean: SeesawLocal.triv_locally_of_dense: an affine open Spec R of S, which is reduced;
% AmpleBaseChange.isProjBy_smooth_isProper_geometricallyIntegral (X over Spec R is projective, smooth and proper with
% geometrically integral fibres, from M; 1,968 lines); SeesawReduced.triv_locally_of_dense over Spec R with dense_affine
% (an r with the point in D(r) and N trivial over D(r); 695 lines: perfect complexes for N and N^{-1} by
% T.nonempty_perfectSectionsComplex_of_smooth_proj, TrivLocus.isClosed_trivLocus, AbLaw.Seesaw.seesaw, and the
% reducedness of R). Size (tools/lean_size.py): 2,738 lines.

Gluing. For every open $U$ of $S$, $f^{-1}(U)$ is an abelian scheme over $U$, so its global functions are those of $U$
\cite[Tag 0E0L]{Sta}. Let $N'=N\otimes f^*\sigma^*N^{-1}$. As $f\circ\sigma=1$, $\sigma^*N'$ has a canonical generator.
For $U$ as in the first step, $N'$ is trivial on $f^{-1}(U)$, since $\sigma^*N$ is trivial on $U$. On $f^{-1}(U)$ choose a
generator of $N'$ whose restriction along $\sigma$ is the canonical one: any generator restricts to a unit of
$\Gamma(U,\mathcal O_U)$ times it, and dividing by the pullback of that unit along $f$ gives one. Two such generators,
on $f^{-1}(U)$ and $f^{-1}(U')$, differ on $f^{-1}(U\cap U')$ by an invertible function, which comes from $U\cap U'$ and is
$1$ along $\sigma$, so it is $1$. So these generators glue to a generator of $N'$, and $N\cong f^*\sigma^*N$, which is
trivial as $\sigma^*N$ is.
\end{proof}
% Lean: seesaw_reduced_of_dense (with Out.GlobalFnsSt): SeesawLocal.triv_locally_of_dense (the first step),
% BaseChangeFns.app_surjective (the global functions) and exists_gen_of_local_of_section (exists_gen_of_local_on_base for
% N' = LinComb N ![1, f ≫ sigma] ![1, -1], then BoxTheta.exists_gen_of_section).

\begin{lemma}[Alp\"oge]\label{cv-toroidal-theta-family-pel-bundles}
Let $S$ be a reduced scheme, let $A$ be an abelian scheme over $S$ and let $L$ be a relatively ample line bundle on $A$.
For points $y_1,y_2$ of $A$ write $\Lambda_L(y_1,y_2)$ as in Definition~\ref{cv-model-toroidal-semiabelian}(2), and let
$y$ be the identity of $A$ and $y_1,y_2$ the projections of $A\times_SA$.
\begin{enumerate}
\item Let $u$ and $v$ be endomorphisms of the group scheme $A$. Suppose that the points $p$ of $S$ with the following
  property are dense: there are a field $K$, a point $x:\operatorname{Spec}K\to S$ whose image contains $p$, a
  commutative group scheme $G$ over $K$ with a morphism $\varphi:G\to A$ over $x$ that identifies $G$ with $x^*A$, a line
  bundle $L'$ on $G$ isomorphic to $\varphi^*L$, and endomorphisms $u',v'$ of $G$ with $\varphi u'=u\varphi$,
  $\varphi v'=v\varphi$ and $\Lambda_{L'}(u'y_1,y_2)\cong\Lambda_{L'}(y_1,v'y_2)$ on $G\times_KG$. Then
  $\Lambda_L(uy_1,y_2)\cong\Lambda_L(y_1,vy_2)$ on $A\times_SA$.
\item Suppose that $L$ is trivialized along the zero section, and let $u$ be an endomorphism of the group scheme $A$.
  Suppose that the points $p$ of $S$ with the following property are dense: there are $K$, $x$, $G$, $\varphi$ and
  $L'\cong\varphi^*L$ as in (1) and an endomorphism $u'$ of $G$ with $\varphi u'=u\varphi$ and
  $L'\cong\Lambda_{L'}(y,u'y)$ on $G$. Then $L\cong\Lambda_L(y,uy)$ on $A$.
\end{enumerate}
\end{lemma}
\leanref{Development/AbelianVarietyTheory/Proofs/ModuliProofs/AdjointEndomorphismsCubeLBOfDenseFibresStatement.lean}{GlobalTheory.T.adjoint_endomorphisms_cubeLB_of_dense_fibres}
\leanref{Development/AbelianVarietyTheory/Proofs/ModuliProofs/LbIsoCubeGraphOfDenseFibresStatement.lean}{GlobalTheory.T.lbIso_cube_graph_of_dense_fibres}
\begin{proof}
Let $\mathcal D$ be the dense set of the hypothesis. For data $K$, $x$, $G$, $\varphi$ as there and points $b_1,b_2$ of
$G$, $\Lambda_{\varphi^*L}(b_1,b_2)=\Lambda_L(\varphi b_1,\varphi b_2)$, as $\varphi$ is a homomorphism over $x$.

(1) The scheme $A$ is smooth over the reduced scheme $S$, so it is reduced. Apply
Lemma~\ref{cv-toroidal-theta-family-pel-bundles-seesaw} with $S_0=S$, $B=A$, $M=L$, the base $A$ over $S$, so that
$X=A\times_SA$ and $f=y_1$, the line bundle $N=\Lambda_L(y_1,vy_2)\otimes\Lambda_L(uy_1,y_2)^{-1}$, and the section
$\sigma=(1,0)$ of $y_1$. Then $\sigma^*N=\Lambda_L(y,0)\otimes\Lambda_L(uy,0)^{-1}$ is trivial, as $\Lambda_L(b,0)\cong
\mathcal O$ canonically. As $A$ is flat and locally of finite presentation over $S$, $A\to S$ is open
\cite[Th\'eor\`eme~2.4.6]{Gro65}, so the points of $A$ over $\mathcal D$ are dense. Let $a$ be such a point, over
$p\in\mathcal D$ with data $K$, $x$, $G$, $\varphi$, $L'$, $u'$, $v'$. As $p$ lies in the image of $x$, $a=\varphi(z)$
for a point $z$ of $G$; let $x'$ be the point of $A$ with values in the residue field $\kappa(z)$ given by $z$ and
$\varphi$. The base change of $X$ along $x'$ is the base change along $z$ of $G\times_KG$ over its first factor, mapped
to $A\times_SA$ by $\varphi\times\varphi$. On $G\times_KG$, as $\varphi u'=u\varphi$ and $\varphi v'=v\varphi$, the
pullback of $N$ is $\Lambda_{\varphi^*L}(y_1,v'y_2)\otimes\Lambda_{\varphi^*L}(u'y_1,y_2)^{-1}\cong
\Lambda_{L'}(y_1,v'y_2)\otimes\Lambda_{L'}(u'y_1,y_2)^{-1}$, which is trivial by hypothesis. So the pullback of $N$ along
the fibre at $x'$ is trivial, and the lemma shows that $N$ is trivial: $\Lambda_L(uy_1,y_2)\cong\Lambda_L(y_1,vy_2)$.

(2) Apply Lemma~\ref{cv-toroidal-theta-family-pel-bundles-seesaw} with $S_0=S$, $B=A$, $M=L$, the base $S$ itself, so
that $X=A$ and $f$ is the structure morphism, the line bundle $N=\Lambda_L(y,uy)\otimes L^{-1}$ and the zero section
$\sigma$. Then $\sigma^*L$ is trivial by hypothesis and $\sigma^*\Lambda_L(y,uy)=\Lambda_L(0,0)$ is trivial. For
$p\in\mathcal D$ with data $K$, $x$, $G$, $\varphi$, $L'$, $u'$, the morphism $\varphi$ identifies $G$ with
$X\times_S\operatorname{Spec}K$, and $\varphi^*N\cong\Lambda_{L'}(y,u'y)\otimes L'^{-1}$, as $\varphi u'=u\varphi$ and
$L'\cong\varphi^*L$; it is trivial by hypothesis. So $N$ is trivial, and $L\cong\Lambda_L(y,uy)$.
\end{proof}
% Lean: (1) adjoint_cubeLB_of_dense_fibres (with Out.GlobalFnsSt): T.seesaw_reduced_base_abelianScheme_of_dense_fibres
% with S := A, X := A ⊗ A, f := fst, pi := snd, N = cubeLB(fst, snd ≫ v) / cubeLB(fst ≫ u, snd), sigma = lift 1 1;
% the dense set through A.hom.isOpenMap and the residue field of z (CechPic.lam_pull_of_isGrpPullback, cls_cube_two,
% lam_one_right). (2) lbIso_cube_of_dense: the same lemma with S := S, X := A, N = cubeLB(1, u) / L, sigma = eSec A
% (PrimDense.exists_gen_cube_eSec and the rigidification of L).

\begin{lemma}[Alp\"oge]\label{cv-toroidal-theta-family-pel-cubical}
In Situation~\ref{cv-function-situation}, let $S$ be a scheme, let $B$ be an abelian scheme over $S$ with an action of
$\mathcal O_E$, and let $M$ be a line bundle on $B$ for which the line bundle $\Theta$ of the theorem of the cube on
$B\times_SB\times_SB$ has a generator. Write $\Lambda_M$ as in Definition~\ref{cv-model-toroidal-semiabelian}(2), with
$y$ the identity of $B$ and $y_1,y_2$ the projections of $B\times_SB$. Suppose that
$\Lambda_M(cy_1,y_2)\cong\Lambda_M(y_1,\bar cy_2)$ on $B\times_SB$ for every $c\in\mathcal O_E$. Then there is a line
bundle $P$ on $B$, trivialized along the zero section, with a cubical structure, and symmetric, such that
$P\cong\Lambda_M(y,w_0y)$ on $B$, $\Lambda_P(y_1,y_2)\cong\Lambda_M(y_1,y_2)$ on $B\times_SB$,
$P\cong\Lambda_P(y,w_0y)$, and $\Lambda_P(cy_1,y_2)\cong\Lambda_P(y_1,\bar cy_2)$ for every $c\in\mathcal O_E$.
\end{lemma}
\leanref{Development/AbelianVarietyTheory/Proofs/ModuliProofs/ExistsPrimitiveNormalisation.lean}{GlobalTheory.PrimNorm.exists_primitive_normalisation}
\begin{proof}
Write $w=w_0$, so that $\bar w=1-w$, and let $P=N=\Lambda_M(y,wy)$, a line bundle on $B$.

Classes. For a scheme $T$ over $S$ and a point $b$ of $B$ with values in $T$, let $m(b)\in\operatorname{Pic}T$ be the
class of $b^*M$, and for points $b_1,b_2$ let $\lambda(b_1,b_2)=m(b_1+b_2)-m(b_1)-m(b_2)+m(0)$, the class of
$\Lambda_M(b_1,b_2)$. Pulling back a generator of $\Theta$ along $(b_1,b_2,b_3)$ gives
$m(b_1+b_2+b_3)-m(b_1+b_2)-m(b_1+b_3)-m(b_2+b_3)+m(b_1)+m(b_2)+m(b_3)-m(0)=0$, so $\lambda$ is additive in each variable.
It is symmetric, and additivity gives $\lambda(-b_1,-b_2)=\lambda(b_1,b_2)$. Pulling back the hypothesis along
$(b_1,b_2)$ gives $\lambda(cb_1,b_2)=\lambda(b_1,\bar cb_2)$ for $c\in\mathcal O_E$.

The class of $N$. The pullback of $N$ along $b$ has class $n(b)=\lambda(b,wb)$. By additivity,
$n(b_1+b_2)-n(b_1)-n(b_2)+n(0)=\lambda(b_1,wb_2)+\lambda(b_2,wb_1)$, and $\lambda(b_2,wb_1)=\lambda(wb_1,b_2)=
\lambda(b_1,\bar wb_2)$. So this is $\lambda(b_1,(w+\bar w)b_2)=\lambda(b_1,b_2)$: $\Lambda_N(b_1,b_2)$ and
$\Lambda_M(b_1,b_2)$ have the same class. As $\lambda$ is additive, the alternating sum of $n$ that corresponds to
$\Theta$ for $N$ vanishes, so $\Theta$ for $N$ is trivial on $B\times_SB\times_SB$.

The structures. As $n(0)=\lambda(0,0)=0$, $0^*N$ is trivial; fix a trivialization. Let $\theta$ be a generator of
$\Theta$ for $N$, divided by the pullback along $B^3\to S$ of its restriction to the zero section, so that $\theta$ is $1$
along the zero section. Each identity required of a cubical structure (Definition~\ref{cv-germs-chart}(1)) says that two
generators of one line bundle on $B^2$, $B^3$ or $B^4$ built from $\theta$ are equal; both are $1$ along the zero section.
Their ratio is an invertible function on the abelian scheme $B^n$ over $S$, hence comes from $S$
\cite[Tag 0E0L]{Sta}, and it is $1$ along the zero section, so it is $1$. So $\theta$ is a cubical structure on $N$. As
$n(-b)=\lambda(-b,-wb)=n(b)$, the line bundle $[-1]^*N\otimes N^{-1}$ is trivial; dividing a generator by the pullback of
its restriction to the zero section, where $[-1]$ is the identity, gives an isomorphism $[-1]^*N\cong N$ that is the
identity along the zero section, so $N$ is symmetric.

The conclusions. $N=\Lambda_M(y,wy)$ by definition, and $\Lambda_N(y_1,y_2)\cong\Lambda_M(y_1,y_2)$ on $B\times_SB$ by the
class of $N$, for the two projections. The class of $\Lambda_N(y,wy)$ is $\lambda(y,wy)=n(y)$, that of $N$. The class
of $\Lambda_N(cy_1,y_2)$ is $\lambda(cy_1,y_2)=\lambda(y_1,\bar cy_2)$, that of $\Lambda_N(y_1,\bar cy_2)$.
\end{proof}
% Lean: exists_primitive_normalisation (with Out.GlobalFnsSt) applies exists_normalised_bundle with w = act w0 and
% w' = act (bar w0) (cmHost_bar_halfIntegerGenerator: bar w0 = 1 - w0). Classes are PullClass (cls, hasGen_iff,
% lbIso_iff), the subgroup of Pic T generated by the pullbacks of M; m is ptCls, lambda is diff2, the cube is diff3_ptCls;
% the class of N is normCls (normCls_left), and diff2_quadratic, diff3_quadratic and diff2_inv_inv are the three
% computations. The cubical structure is AbLaw.Cube.exists_normalised and cubicalOfNormalised (GlobalFnsSt, through
% corr_of_val_section); the symmetry is isSymmetric_of_hasGen. Lemma cv-germs-abelian-normal states the case of an
% abelian variety over an algebraically closed field, with the same computation.

\begin{lemma}[Alp\"oge]\label{cv-toroidal-abelian-dvr-eigen}
In Situation~\ref{cv-function-situation}, let $m_+,m_-\ge0$ be integers, let $R$ be a reduced ring with an
element $r$ such that $r^2=-d$, let $n\ge0$, and let $M:\mathcal O_E\to M_n(R)$ be a ring homomorphism. For a
homomorphism $f:R\to k$ to an algebraically closed field of characteristic zero, let
$M_f:\mathcal O_E\to M_n(k)$ be the composite of $M$ with $f$ applied to the entries, and say that $M$ has
\emph{signature} $(m_+,m_-)$ at $f$ if for every ring homomorphism $\tau:\mathcal O_E\to k$ the subspace of
$k^n$ on which each $b\in\mathcal O_E$ acts through $M_f$ by $\tau(b)$ has dimension $m_+$ if
$\tau(\sqrt{-d})=f(r)$ and $m_-$ otherwise. If the kernels of the homomorphisms $f$ at which $M$ has
signature $(m_+,m_-)$ are dense in $\operatorname{Spec}R$, then $M$ has signature $(m_+,m_-)$ at every such
homomorphism $f$.
\end{lemma}
\leanref{Development/AbelianVarietyTheory/Proofs/ModuliProofs/CmTypeEigenMultiplicityOfDensePrimesStatement.lean}{GlobalTheory.T.cmType_eigenMultiplicity_of_dense_primes}
\leanref{Development/AbelianVarietyTheory/Proofs/ModuliProofs/CmTypeEigenMultiplicityOfDensePrimes.lean}{GlobalTheory.eigenMult_of_dense}
\begin{proof}
Let $f:R\to k$ be as in the statement, let $W=k^n$ with the action of $\mathcal O_E$ through $M_f$, and for
$\tau:\mathcal O_E\to k$ let $W_\tau$ be the subspace on which each $b$ acts by $\tau(b)$. As $E$ is
separable over $\mathbf Q$ and $k$ is algebraically closed of characteristic zero,
$\mathcal O_E\otimes k$ is a product of copies of $k$ indexed by the $[E:\mathbf Q]$ homomorphisms $\tau$. So
$W$ is the direct sum of the $W_\tau$, and $\operatorname{tr}M_f(b)=\sum_\tau\dim W_\tau\,\tau(b)$. Each
$\tau$ maps $\sqrt{-d}$ to $f(r)$ or to $-f(r)$, as both square to $-d$, and $f(r)\ne0$. Let $m_\tau$ be $m_+$
if $\tau(\sqrt{-d})=f(r)$ and $m_-$ otherwise; then
$2f(r)\,m_\tau\tau(b)=(m_++m_-)f(r)\tau(b)+(m_+-m_-)\tau(\sqrt{-d}\,b)$. The distinct homomorphisms $\tau$
are linearly independent over $k$ \cite[Chapter~VI, Theorem~4.1]{Lan02}. So $M$ has signature $(m_+,m_-)$
at $f$ if and only if $\operatorname{tr}M_f(b)=\sum_\tau m_\tau\tau(b)$ for all $b\in\mathcal O_E$, that is,
since $\sum_\tau\tau(c)=\operatorname{Tr}_{E/\mathbf Q}(c)$ for $c\in\mathcal O_E$ and $2f(r)$ is invertible
in $k$, if and only if $f(X_b)=0$ for all $b$, where
\[
 X_b=2r\operatorname{tr}M(b)-(m_++m_-)\,r\operatorname{Tr}_{E/\mathbf Q}(b)
 -(m_+-m_-)\operatorname{Tr}_{E/\mathbf Q}(\sqrt{-d}\,b)\in R.
\]
By hypothesis, each $X_b$ lies in a dense set of prime ideals of $R$, hence in every prime ideal, as those
containing $X_b$ form a closed set. So $X_b$ is nilpotent, and $X_b=0$ because $R$ is reduced. Hence
$f(X_b)=0$ for every $f$ and every $b$.
\end{proof}
% Lean: eigenMult_of_dense: TraceSpread.finrank_iff_trace (the criterion by traces) and TraceSpread.map_trace.

\begin{lemma}[Alp\"oge]\label{cv-toroidal-theta-family-pel-lie}
In Situation~\ref{cv-function-situation}, let $o$ be a ring with an element $\rho$ such that $\rho^2=-d$, let $S$ be a
reduced scheme over $o$, and let $G$ be a commutative group scheme over $S$, smooth of relative dimension $g$, with an
action of $\mathcal O_E$. For an algebraically closed field $k$ of characteristic zero and a point
$x:\operatorname{Spec}k\to S$ over $\beta:o\to k$, say that $G$ satisfies the Lie condition at $x$ if condition (a) of
Definition~\ref{cv-toroidal-family}(1) holds for $G$ at $x$ with $\beta(\rho)$ in place of $\beta(\theta)$. Suppose that
the points $p$ of $S$ that lie in the image of a point $x$ at which $G$ satisfies the Lie condition are dense. Then $G$
satisfies the Lie condition at every such point $x$.
\end{lemma}
\leanref{Development/AbelianVarietyTheory/Proofs/ModuliProofs/DeterminantConditionSpreadsOfDensePointsStatement.lean}{GlobalTheory.T.determinantCondition_spreads_of_dense_points}
\begin{proof}
This is Lemma~\ref{cv-toroidal-abelian-dvr-determinant}(1), for $G$, the structure morphism
$S\to\operatorname{Spec}o$ and $\rho$: the Lie condition at $x$ is the signature $(p,q)$ there.
\end{proof}
% Lean: the same declaration T.determinantCondition_spreads_of_dense_points, whose proof
% (DetSpread.detCond_of_dense) Lemma cv-toroidal-abelian-dvr-determinant(1) follows.

\begin{lemma}[Alp\"oge]\label{cv-toroidal-theta-family-pel-level}
In Situation~\ref{cv-function-situation}, let $o$ be a ring with an element $z$, let $S$ be a scheme over $o$, let $A$
be an abelian scheme over $S$ of relative dimension $g$ with a line bundle $L$ that is relatively ample, trivialized
along the zero section, isomorphic to $[-1]^*L$ locally on $S$ and defines a principal polarization, and let
$\alpha(c)$, $c\in\mathcal O_E$, be endomorphisms of $A$. Let $(\lambda_i)_{i\in\kappa}$ be a family in $\Lambda$ and
$(y_i)_{i\in\kappa}$ a family of points of $A$ over $S$, and let $n\ge1$ and $m\ge0$. For a field $K$, a point
$x:\operatorname{Spec}K\to S$ over $\beta:o\to K$ with $n\ne0$ in $K$, call a \emph{level datum} at $x$ a primitive
$n$-th root of unity $\zeta''\in K$ and an additive map $\iota:\Lambda\to A_x(K)$ with $n\iota=0$,
$\iota(c\lambda)=\alpha(c)\iota(\lambda)$, $e_n(\iota(\lambda),\iota(\lambda'))=\zeta''^{\langle\lambda,\lambda'\rangle}$
for the commutator pairing $e_n$ of the theta group of $L^n$, $\zeta''^m=\beta(z)$ and $m\iota(\lambda_i)=y_i(x)$.
Suppose that the points $p$ of $S$ that lie in the image of a point $x$ with a level datum are dense. Then there is a
level datum at every point $x:\operatorname{Spec}K\to S$ with $K$ algebraically closed and $n\ne0$ in $K$.
\end{lemma}
\leanref{Development/AbelianVarietyTheory/Proofs/ModuliProofs/ExistsLevelIsometryOfDenseGeometricPointsStatement.lean}{GlobalTheory.T.exists_level_isometry_of_dense_geometric_points}
\begin{proof}
This is Lemma~\ref{cv-toroidal-abelian-dvr-isometry}, for $P=A$ with $\mathcal L=L$, $\alpha_b=\alpha(b)$, the
structure morphism $S\to\operatorname{Spec}o$, $z$, $n$, $m$, $a_i=\lambda_i$ and $p_i=y_i$: a level datum at $x$ is an
$n$-isometry at $x$ there.
\end{proof}
% Lean: the same declaration T.exists_level_isometry_of_dense_geometric_points (with Out.Laws), whose proof
% (exists_levelIsometry_of_dense) Lemma cv-toroidal-abelian-dvr-isometry follows.

\begin{lemma}[Alp\"oge]\label{cv-toroidal-theta-family-pel-tate}
In Situation~\ref{cv-function-situation}, let $K$ be a field and $Q$ a group with maps $c\cdot{}:Q\to Q$ for
$c\in\mathcal O_E$, and for each $n\ge1$ let $e_n$ be a relation between pairs of elements of $Q$ and elements of
$K^\times$.
\begin{enumerate}
\item Let $n\ge1$, let $\zeta''\in K$ be a primitive $n$-th root of unity, and let $\iota:\Lambda\to Q$ be a homomorphism
  with $\iota(\lambda)^n=1$ and $e_n(\iota(\lambda),\iota(\lambda'),\zeta''^{\langle\lambda,\lambda'\rangle})$ for all
  $\lambda,\lambda'\in\Lambda$. Suppose that the elements of $Q$ killed by $n$ are $n^{\operatorname{rk}\Lambda}$ in
  number, and that $e_n(1,y,c)$ implies $c=1$ when $y^n=1$. Then the kernel of $\iota$ is $n\Lambda$, and its image is the set
  of elements killed by $n$.
\item Let $k\ge1$, $z\in K$, $(\lambda_i)_{i\in\kappa}$ a family in $\Lambda$ and $(y_i)_{i\in\kappa}$ a family in $Q$.
  Suppose that only finitely many elements of $Q$ are killed by each $n\ge1$; that $e_{nm}(y,y',c)$ implies
  $e_n(y^m,y'^m,c^m)$ for $n,m\ge1$ and $y,y'$ killed by $nm$; and that for every $m\ge1$ there are a primitive $km$-th
  root of unity $\zeta''\in K$ and a homomorphism $\iota:\Lambda\to Q$ with $\iota(c\lambda)=c\cdot\iota(\lambda)$,
  kernel $km\Lambda$ and image the elements killed by $km$, with
  $e_{km}(\iota(\lambda),\iota(\lambda'),\zeta''^{\langle\lambda,\lambda'\rangle})$, $\zeta''^m=z$ and
  $\iota(\lambda_i)^m=y_i$. Then there are primitive $n$-th roots of unity $\zeta'_n\in K$ and homomorphisms
  $\iota_n:\Lambda\to Q$, for $n\ge1$, with $(\zeta'_{nn'})^{n'}=\zeta'_n$, $\iota_n(c\lambda)=c\cdot\iota_n(\lambda)$,
  kernel $n\Lambda$, image the elements killed by $n$, $\iota_{nn'}(\lambda)^{n'}=\iota_n(\lambda)$ and
  $e_n(\iota_n(\lambda),\iota_n(\lambda'),(\zeta'_n)^{\langle\lambda,\lambda'\rangle})$, such that $\zeta'_k=z$ and
  $\iota_k(\lambda_i)=y_i$ for every $i$.
\end{enumerate}
\end{lemma}
\leanref{Development/AbelianVarietyTheory/Proofs/ModuliProofs/TorsionIsometryPrimitiveRootBijectiveModNStatement.lean}{GlobalTheory.T.torsionIsometry_primitiveRoot_bijective_mod_n}
\leanref{Development/AbelianVarietyTheory/Proofs/ModuliProofs/TateIsometryOfFiniteLevelIsometriesStatement.lean}{GlobalTheory.T.tateIsometry_of_finiteLevel_isometries}
\begin{proof}
(1) If $\lambda\in n\Lambda$ then $\iota(\lambda)=1$, as $\iota$ is killed by $n$. Conversely let $\iota(\lambda)=1$.
For every $\lambda'$, $e_n(1,\iota(\lambda'),\zeta''^{\langle\lambda,\lambda'\rangle})$ holds, so
$\zeta''^{\langle\lambda,\lambda'\rangle}=1$ and $n$ divides $\langle\lambda,\lambda'\rangle$. As the alternating form
is perfect, the linear form $\lambda'\mapsto\langle\lambda,\lambda'\rangle/n$ is $\langle\mu,\cdot\rangle$ for some
$\mu\in\Lambda$, and $\lambda=n\mu$. So $\iota$ induces an injective map from $\Lambda/n\Lambda$, a set of
$n^{\operatorname{rk}\Lambda}$ elements as $\Lambda$ is a free $\mathbf Z$-module, to the set of elements of $Q$ killed by
$n$, which has as many elements. So this map is onto.

(2) For $n\ge1$ call a \emph{level-$n$ isometry} a pair $(\zeta'',\iota)$ of a primitive $n$-th root of unity
$\zeta''\in K$ and a homomorphism $\iota:\Lambda\to Q$ with $\iota(c\lambda)=c\cdot\iota(\lambda)$, kernel $n\Lambda$,
image the elements killed by $n$, and $e_n(\iota(\lambda),\iota(\lambda'),\zeta''^{\langle\lambda,\lambda'\rangle})$.

If $(\zeta'',\iota)$ is a level-$nr$ isometry, then $(\zeta''^r,\iota^r)$, with $\iota^r(\lambda)=\iota(\lambda)^r=
\iota(r\lambda)$, is a level-$n$ isometry. The root $\zeta''^r$ is primitive of order $n$, and $\iota^r$ is additive and
compatible with the maps $c\cdot{}$. Its kernel: $\iota(r\lambda)=1$ exactly when $r\lambda\in nr\Lambda$, that is
$\lambda\in n\Lambda$, as $\Lambda$ has no torsion. Its image: if $y^n=1$, then $y^{nr}=1$, so $y=\iota(\mu)$ for some
$\mu$; then $\iota(n\mu)=1$, so $n\mu\in nr\Lambda$, $\mu=r\mu'$ and $y=\iota^r(\mu')$; conversely
$\iota^r(\lambda)^n=\iota(nr\lambda)=1$. The pairing: the hypothesis on the $e_n$ turns
$e_{nr}(\iota(\lambda),\iota(\lambda'),\zeta''^{\langle\lambda,\lambda'\rangle})$ into
$e_n(\iota(\lambda)^r,\iota(\lambda')^r,(\zeta''^r)^{\langle\lambda,\lambda'\rangle})$.

There are finitely many level-$n$ isometries: $\zeta''$ is an $n$-th root of unity, and $\iota$ is determined by its
values on a finite set of generators of the finitely generated group $\Lambda$, which lie in the finite set of elements
killed by $n$.

For $j\ge1$ let $p_j$ be the level-$kj!$ isometry of the hypothesis for $m=j!$, so that its root has $j!$-th power $z$ and
its map sends $\lambda_i$ to an element whose $j!$-th power is $y_i$. Let $\mathcal U$ be an ultrafilter on the positive
integers that contains the complements of finite sets. Let $n\ge1$. For $j\ge n$, $n$ divides $kj!$, so
$p_j^{kj!/n}$ is a level-$n$ isometry. As these are finitely many, there is one, $s_n$, with $p_j^{kj!/n}=s_n$ for all
$j$ in a set in $\mathcal U$. For $n,m\ge1$, some $j$ lies in the sets for $nm$ and for $n$, and then
$s_{nm}^m=p_j^{kj!/n}=s_n$. Write $s_n=(\zeta'_n,\iota_n)$. These are primitive roots and homomorphisms with the
required properties, and $(\zeta'_{nn'})^{n'}=\zeta'_n$ and $\iota_{nn'}(\lambda)^{n'}=\iota_n(\lambda)$. Finally,
for some $j$, $s_k=p_j^{j!}$, so $\zeta'_k=z$ and $\iota_k(\lambda_i)=y_i$.
\end{proof}
% Lean: (1) is TorsIsometry.ker_and_onto (exists_eq_nsmul_of_dvd with D.unimod; the injection from (Z/n)^{rank}
% through a Z-basis of La and Fintype.bijective_iff_injective_and_card); (2) is TateTower.tateIso_of_levels
% (= T.tateIsometry_of_finiteLevel_isometries): LevelIso, LevelIso.pow (the power of an isometry, with nsmul_cancel),
% finite_levelIso (through a finite generating set of La), the ultrafilter Ultrafilter.of Filter.atTop on the levels
% k * j!, Ultrafilter.eq_pure_of_finite_mem (s n), htow (the tower), and hj for the level k.

\begin{lemma}[Alp\"oge]\label{cv-toroidal-theta-family-pel}
In Situation~\ref{cv-function-situation}, suppose $p+q\ge1$, let $g=\frac12[E:\mathbf Q](p+q)$, let
$e_1,\dots,e_{2g}$ be a symplectic basis of the lattice $\Lambda$ of the datum, and let $a:I\to\mathcal O_E$ be a
map whose image generates $\mathcal O_E$ as a ring, with elements $w,w'\in I$ such that $a(w)=w_0$ and
$a(w')=\bar w_0$; the hypotheses that $p+q\ge1$, that the basis is symplectic, that $a(w)=w_0$ and that
$a(w')=\bar w_0$ are not used in the proof. Let $N'\ge1$, let $o$ be a ring with a unit $\zeta$ as in Definition~\ref{cv-toroidal-theta},
let $\mathrm{em}:o\to\mathbf C$ be an injective homomorphism and let $\rho\in o$ with $\mathrm{em}(\rho)=\theta$. Let
$S$ be a reduced, locally noetherian scheme over $o$, and let $t$ be a theta-level object over $S$ such that the
points of $S$ at which $t$ is a $\mathcal C$-object (Definition~\ref{cv-toroidal-theta}(4)), for the PEL class
$\mathcal C$ of $((e_i),a,\rho,\mathrm{em})$, are dense. Then $t$ has a PEL structure for $((e_i),a,\rho)$.
\end{lemma}
\leanref{Development/AbelianVarietyTheory/Proofs/ModuliProofs/ExistsIsPELTripleOfDenseFibresStatement.lean}{GlobalTheory.T.exists_isPELTriple_of_dense_fibres}
\leanref{Development/AbelianVarietyTheory/Proofs/ModuliProofs/ExistsIsPELTripleOfDenseFibres.lean}{GlobalTheory.PELFamily.exists_isPEL_of_dense}
\begin{proof}
Write $t=(A,L,(x_i),(\alpha_h))$, and let $\mathcal D$ be the dense set of points of $S$ at which $t$ is a
$\mathcal C$-object.

The fibres. Let $p\in\mathcal D$. Lemma~\ref{cv-toroidal-theta-family-pel-fibre} gives a field $K$ of characteristic
zero, a point $x:\operatorname{Spec}K\to S$ with image $p$, and a theta-level object over $K$ with a morphism to $t$
over $x$ and a PEL structure. Composing the isomorphism of the PEL structure with this morphism gives an
$\mathcal O_E$-semi-abelian scheme $G_x$ over $K$ and a morphism $\psi_x:G_x\to A$ over $x$ that identifies $G_x$ with
$x^*A$, carries the line bundle of $G_x$ to $L$ and $a(h)$ to $\alpha_h$ for every $h$, and satisfies the conditions
at the points of characteristic zero of Definition~\ref{cv-toroidal-theta}(6), with the points $x_i$ of $t$.

The action. As the image of $a$ generates $\mathcal O_E$ and $S$ is reduced,
Lemma~\ref{cv-toroidal-theta-family-pel-action}, applied to the $\alpha_h$ and the fibres $G_x$, gives an action $r$
of $\mathcal O_E$ on $A$ with $r(a(h))=\alpha_h$ for every $h$, such that every $\psi_x$ carries the action on $G_x$ to
$r$.

The line bundle. Let $c\in\mathcal O_E$. On each $G_x$, complex conjugation is the Rosati involution, and the line
bundle is normalized. Through $\psi_x$, Lemma~\ref{cv-toroidal-theta-family-pel-bundles}(1), applied to $r(c)$ and
$r(\bar c)$, gives $\Lambda_L(r(c)y_1,y_2)\cong\Lambda_L(y_1,r(\bar c)y_2)$ on $A\times_SA$, and part (2), applied to
$r(w_0)$ and using that $L$ is trivialized along the zero section, gives $L\cong\Lambda_L(y,r(w_0)y)$ on $A$. By the
theorem of the cube, the line bundle $\Theta$ of $L$ has a generator. So
Lemma~\ref{cv-toroidal-theta-family-pel-cubical}, applied to $A$, $r$ and $L$, gives a symmetric line bundle $P$ on $A$,
trivialized along the zero section, with a cubical structure, such that $P\cong\Lambda_L(y,r(w_0)y)$, hence $P\cong L$;
$P\cong\Lambda_P(y,r(w_0)y)$; and $\Lambda_P(r(c)y_1,y_2)\cong\Lambda_P(y_1,r(\bar c)y_2)$ for every $c$. The abelian
scheme $A$ is semi-abelian, so $A$ with $r$ and $P$ is an $\mathcal O_E$-semi-abelian scheme $G$ over $S$
(Definition~\ref{cv-model-toroidal-semiabelian}(2)). The identity $G\to A$ is an isomorphism of group schemes that
carries the line bundle $P$ of $G$ to $L$ and the action of $a(h)$ to $\alpha_h$.

The Lie algebra. As $\mathrm{em}(\rho)^2=\theta^2=-d$ and $\mathrm{em}$ is injective, $\rho^2=-d$ in $o$. Let
$p\in\mathcal D$, with $x$ and $G_x$ as above, and let $\bar x$ be the composite of $x$ with
$\operatorname{Spec}\bar K\to\operatorname{Spec}K$ for an algebraic closure $\bar K$ of $K$. The condition on the Lie
algebra holds for $G_x$ at $\operatorname{Spec}\bar K$, and $\psi_x$ identifies $G_x$, with its action, with the base
change of $A$ with $r$; so it holds for $A$ with $r$ at $\bar x$, a point with image $p$. The scheme $A$ is smooth of
relative dimension $g$, so Lemma~\ref{cv-toroidal-theta-family-pel-lie} gives condition (a) of
Definition~\ref{cv-toroidal-family}(1) for $G$, with $\beta(\rho)$, at every point $x:\operatorname{Spec}k\to S$ with $k$
algebraically closed of characteristic zero and $\beta:o\to k$ the induced homomorphism.

The Tate systems. Let $x:\operatorname{Spec}k\to S$ and $\beta$ be as above. For $n\ge1$ the points of $A_x(k)$ killed
by $n$ are $n^{2g}$ in number, and the commutator pairing of the theta group of $L^n$ takes the value $1$ when its first
argument is the origin. For $m\ge1$, the systems of roots of unity and maps of condition (b) at the points $\bar x$ above,
carried to $A$ by $\psi_x$, give at $\bar x$ a primitive $N'^2m$-th root of unity $\zeta''$ and a map
$\iota:\Lambda\to A_{\bar x}(\bar K)$ as in Lemma~\ref{cv-toroidal-theta-family-pel-level}, for $n=N'^2m$, $m$,
$\alpha=r$, $z=\zeta$, the family $(e_i)$ and the points $(x_i)$: $\zeta''$ and $\iota$ are those of level $N'^2m$, and
$\zeta''^m=\beta(\zeta)$ and $m\iota(e_i)=x_i$ by the compatibilities of the system. By that lemma, such data exist at
$x$. By Lemma~\ref{cv-toroidal-theta-family-pel-tate}(1), their map $\iota$ has kernel $N'^2m\Lambda$ and image the
$N'^2m$-torsion. The commutator pairings of the theta groups of the powers of $L$ satisfy
$e_{nm}(y,y',c)\Rightarrow e_n(my,my',c^m)$ for points killed by $nm$. So
Lemma~\ref{cv-toroidal-theta-family-pel-tate}(2), with $Q=A_x(k)$, $k=N'^2$ there and $z=\beta(\zeta)$, gives roots of
unity $\zeta'_n$ and maps $\iota_n$ as in condition (b) of Definition~\ref{cv-toroidal-family}(1), with
$\zeta'_{N'^2}=\beta(\zeta)$ and $\iota_{N'^2}(e_i)=x_i(x)$, for the pairings of $L$. As $P\cong L$, the pairings of $P$
are those of $L$. So $G$, with the identity $G\to A$, is a PEL structure on $t$ for $((e_i),a,\rho)$.
\end{proof}
% Lean: PELFamily.exists_isPEL_of_dense (for theta-level objects of any level k; the text takes k = N'^2): exists_fibre_map
% (the fibres); T.oAct_spread_from_dense_fibres_abelianScheme (r, hra, hr); T.adjoint_endomorphisms_cubeLB_of_dense_fibres
% (hadj, from tau.adj); T.lbIso_cube_graph_of_dense_fibres (hprim, from tau.prim); PrimNorm.exists_primitive_normalisation
% with hcube (the theorem of the cube), and iso' (P ≅ L); the SAB on t.A with Iso.refl; T.determinantCondition_spreads_of_
% dense_points with exists_geom; for the Tate systems hcard (hL.mulN), T.tateIsometry_of_finiteLevel_isometries with
% weilIs_lower (hL.weil), T.exists_level_isometry_of_dense_geometric_points (from exists_geom: I.zeta (k*m), I.ev (k*m),
% I.tower, I.zeta_pow) and T.torsionIsometry_primitiveRoot_bijective_mod_n with eq_one_of_weilIs_one; weilIs_of_lbIso.
% The hypotheses 1 <= p + q, IsSympBasis and a w' = bar w0 are introduced as _, and a w = w0 is named hw but not used;
% Out.Laws2 is not used. Size: about 9,200 lines, about 3,900 of them not shared with the other inputs of Lemma
% cv-toroidal-theta-family; about 600 beyond Lemmas cv-toroidal-theta-family-pel-fibre to -tate.

\subsection{Smooth theta families}\label{subsection-cv-smooth-theta-families}

\begin{theorem}[Osgood {\cite[Theorem~1.2.1]{Leb26}}]\label{cv-toroidal-theta-family-trdeg-osgood}
Let $m\ge0$, let $U\subset\mathbf C^m$ be open, and let $f:U\to\mathbf C$ be complex differentiable at every point of
$U$. Then $f$ is analytic on $U$: near every point of $U$ it is the sum of a convergent power series.
\end{theorem}
\leanref{ReciprocityTheory/AnalyticForms/AnalyticCharts/Osgood.lean}{ReciprocityTheory.AnForm.Osgood.osgood}

\begin{lemma}[Alp\"oge]\label{cv-toroidal-theta-family-trdeg-general}
Let $o$ be a countable ring with a homomorphism $\mathrm{em}:o\to\mathbf C$, let $I$ be a countable set, let $m\ge0$,
let $B\subset\mathbf C^m$ be a nonempty connected open set, and let $\vartheta:B\to\mathbf C^I$ be a map whose
coordinates are complex differentiable on $B$. Then there are $z\in B$ and $c\in\mathbf C^\times$ such that every
polynomial $F$ in the variables $X_a$, $a\in I$, with coefficients in $o$ and with $\mathrm{em}(F)(c\vartheta(z))=0$
satisfies $\mathrm{em}(F)(c'\vartheta(z'))=0$ for all $z'\in B$ and $c'\in\mathbf C^\times$.
\end{lemma}
\leanref{Development/AbelianVarietyTheory/Proofs/ModuliProofs/ExistsVeryGeneralPointOfHolomorphicFamily.lean}{GlobalTheory.exists_veryGeneral_point_of_holomorphic_family}
\begin{proof}
Let $\Omega=\mathbf C^\times\times B\subset\mathbf C^{1+m}$, a nonempty connected open set. For a polynomial $F$ let
$f_F(c',z')=\mathrm{em}(F)(c'\vartheta(z'))$ for $(c',z')\in\Omega$. It involves finitely many coordinates of
$\vartheta$, so it is complex differentiable on $\Omega$, hence analytic there by
Theorem~\ref{cv-toroidal-theta-family-trdeg-osgood}. Let $W_F=\mathbf C^{1+m}$ if $f_F$ vanishes on $\Omega$; otherwise
let $W_F$ be the union of $\{p\in\Omega:f_F(p)\ne0\}$ and the complement of the closure of $\Omega$. Each $W_F$ is open.
It is dense: if $f_F$ does not vanish on $\Omega$, its zeros in the connected open set $\Omega$ have empty interior by
the identity theorem \cite[Theorem~1.2.7]{Leb26}, so every point of $\Omega$, and hence of its closure, lies in the
closure of $\{p\in\Omega:f_F(p)\ne0\}$.

The polynomial ring over the countable ring $o$ in countably many variables is countable. By the Baire category
theorem, the intersection of the $W_F$ is dense in $\mathbf C^{1+m}$, so it meets the nonempty open set $\Omega$ in a
point $(c,z)$. If $\mathrm{em}(F)(c\vartheta(z))=0$, then $(c,z)$ lies neither in $\{p\in\Omega:f_F(p)\ne0\}$ nor
outside the closure of $\Omega$; so $W_F=\mathbf C^{1+m}$, and $f_F$ vanishes on $\Omega$.
\end{proof}
% Lean: exists_veryGeneral_point_of_holomorphic_family: Omega = {p : Fin.tail p in B, p 0 ≠ 0} (connected as the image
% of B × (C \ {0}), isConnected_compl_singleton_of_one_lt_rank), f F p = eval₂ em (p 0 • θ (Fin.tail p)) F, analytic
% by AnForm.Osgood.osgood; W F as in the text, dense by AnalyticOnNhd.eqOn_zero_of_preconnected_of_eventuallyEq_zero;
% dense_iInter_of_isOpen (Baire) over the countable MvPolynomial I o.

\begin{lemma}\label{cv-toroidal-theta-family-trdeg-rank}
Let $E$ and $F$ be finite-dimensional complex vector spaces, let $U\subset E$ be a nonempty open set, and let
$f:U\to F$ be continuously complex differentiable at every point of $U$. Suppose that for every $z\in U$ only countably
many $z'\in U$ have $f(z')=f(z)$. Then the complex derivative of $f$ is injective at some point of $U$.
\end{lemma}
\leanref{Development/AbelianVarietyTheory/Proofs/ModuliProofs/MvPolynomialEqZeroOfEventually.lean}{GlobalTheory.CountableFibres.exists_injective_fderiv}
\begin{proof}
Write $Df(z)$ for the derivative of $f$ at $z$, and choose $z_0\in U$ at which the rank $r$ of $Df(z_0)$ is largest.
Suppose that $L=Df(z_0)$ is not injective. Let $W$ be the image of $L$ and $P:F\to W$ a linear projection onto $W$.
Then $P\circ L:E\to W$ is surjective, with kernel $K=\ker L\ne0$. As $Df$ is continuous at $z_0$, $P\circ Df(z)$ is
surjective for $z$ near $z_0$. For such $z$, $Df(z)$ has rank at least $r$, hence equal to $r$, and the inclusion
$\ker Df(z)\subset\ker(P\circ Df(z))$ is an equality, both having dimension $\dim E-r$.

Let $\pi$ be a linear projection of $E$ onto $K$. The implicit function theorem \cite[Theorem~1.3.8]{Leb26}, applied
to $P\circ f$ at $z_0$, gives a ball $S$ around $0$ in $K$ and a complex differentiable map $\psi:S\to U$ with
$\psi(0)=z_0$, $P(f(\psi(t)))=P(f(z_0))$ and $\pi(\psi(t)-z_0)=t$ for $t\in S$; in particular $\psi$ is injective.
Shrinking $S$, the points $\psi(t)$ lie in the neighbourhood of $z_0$ above. For $t\in S$ and $v\in K$, the derivative
of $P\circ f\circ\psi$ vanishes, so $D\psi(t)v$ lies in $\ker(P\circ Df(\psi(t)))=\ker Df(\psi(t))$. So $f\circ\psi$ has
derivative $0$ on the connected set $S$, and it is constant, equal to $f(z_0)$. Then $\psi$ maps the uncountable set
$S$ injectively into the countable set of $z'\in U$ with $f(z')=f(z_0)$, which is absurd.
\end{proof}
% Lean: exists_injective_fderiv: z0 maximizes rk z = finrank (range (fderiv f z)) (Nat.sSup_mem); P a left inverse
% of W.subtype; hker (the kernels agree near z0, from the units being open and the maximality of the rank);
% HasStrictFDerivAt.implicitToOpenPartialHomeomorph for P ∘ f, with psi t = Psi.symm (P (f z0), t); hderiv and
% IsOpen.is_const_of_fderiv_eq_zero (f ∘ psi constant on a ball); Set.Countable.dense_compl for the contradiction.

\begin{lemma}[Alp\"oge]\label{cv-toroidal-theta-family-trdeg-ratios}
Let $d\ge0$, let $I$ be a finite set, let $B\subset\mathbf C^d$ be a nonempty connected open set, and let
$\vartheta:B\to\mathbf C^I$ be a map whose coordinates are complex differentiable on $B$. Suppose that for every
$z\in B$, $\vartheta(z)\ne0$ and only countably many $z'\in B$ have $\vartheta(z')\in\mathbf C\,\vartheta(z)$. The
hypothesis that $B$ is connected is not used in the proof. Then there are $a_0,a_1,\dots,a_d\in I$ such that
$\vartheta_{a_0}(z)\ne0$ for some $z\in B$, and such that the only polynomial $F\in\mathbf C[Y_1,\dots,Y_d]$ with
$F(\vartheta_{a_1}(z)/\vartheta_{a_0}(z),\dots,\vartheta_{a_d}(z)/\vartheta_{a_0}(z))=0$ for every $z\in B$ with
$\vartheta_{a_0}(z)\ne0$ is $F=0$.
\end{lemma}
\leanref{Development/AbelianVarietyTheory/Proofs/ModuliProofs/MvPolynomialEqZeroOfEventually.lean}{GlobalTheory.CountableFibres.exists_independent_ratios}
\begin{proof}
By Theorem~\ref{cv-toroidal-theta-family-trdeg-osgood} the coordinates $\vartheta_a$ are analytic on $B$. Choose
$z_1\in B$, and $a_0$ with $\vartheta_{a_0}(z_1)\ne0$, which exists as $\vartheta(z_1)\ne0$. Let $U$ be the open set of
$z\in B$ with $\vartheta_{a_0}(z)\ne0$, which contains $z_1$, and let $\rho:U\to\mathbf C^I$ be
$\rho(z)=\vartheta(z)/\vartheta_{a_0}(z)$, a continuously complex differentiable map. If $\rho(z')=\rho(z)$, then
$\vartheta(z')=(\vartheta_{a_0}(z')/\vartheta_{a_0}(z))\vartheta(z)$; so each fibre of $\rho$ is countable. By
Lemma~\ref{cv-toroidal-theta-family-trdeg-rank} there is $z_2\in U$ at which the derivative $D\rho(z_2):\mathbf C^d\to
\mathbf C^I$ is injective. The linear forms $\ell_a$, $a\in I$, given by the coordinates of $D\rho(z_2)$ then span the
dual of $\mathbf C^d$, so $d$ of them, $\ell_{a_1},\dots,\ell_{a_d}$, form a basis. With $\pi:\mathbf C^I\to\mathbf C^d$
the projection to the coordinates $a_1,\dots,a_d$, the derivative of $\pi\circ\rho$ at $z_2$ is invertible.

By the inverse function theorem, $\pi\circ\rho$ maps $U$ onto a neighbourhood of $\pi(\rho(z_2))$, which contains a
product $D_1\times\dots\times D_d$ of open discs. Let $F$ vanish at $\pi(\rho(z))$ for every $z\in U$. Each disc $D_i$ is
infinite, so it contains a finite set $S_i$ with more elements than the degree of $F$ in $Y_i$. Then $F$ vanishes on
$S_1\times\dots\times S_d$, and so $F=0$: by induction on $d$, the coefficients of $F$ as a polynomial in $Y_d$ vanish on
$S_1\times\dots\times S_{d-1}$, since for each point there the polynomial in $Y_d$ has more roots than its degree.
\end{proof}
% Lean: exists_independent_ratios (the hypothesis IsPreconnected B is introduced as _): AnForm.Osgood.osgood,
% exists_injective_fderiv for rho = theta / theta a0 on U, exists_linearIndependent' for the rows a, hsurj and
% HasStrictFDerivAt.map_nhds_eq_of_surj (the inverse function theorem), Metric.mem_nhds_iff and ball_pi (the
% polydisc), Set.Infinite.exists_subset_card_eq, and MvPolynomial.eq_zero_of_eval_zero_at_prod_finset.

\begin{definition}\label{cv-toroidal-theta-family-trdeg-cone}
Let $o$ be a ring with a homomorphism $\mathrm{em}:o\to\mathbf C$, let $I$ be a set, let $m\ge0$, let
$B\subset\mathbf C^m$, and let $\vartheta:B\to\mathbf C^I$ be a map. The \emph{cone ideal} of $(B,\vartheta)$ is the
ideal of the polynomials $F$ in the variables $X_a$, $a\in I$, with coefficients in $o$ such that
$\mathrm{em}(F)(c\vartheta(z))=0$ for all $z\in B$ and $c\in\mathbf C^\times$.
\end{definition}
\leanref{Development/AbelianVarietyTheory/Proofs/ModuliProofs/ThetaIdealPrimeAvoidingVariable.lean}{GlobalTheory.ThetaPrime.coneIdeal}

\begin{lemma}[Alp\"oge]\label{cv-toroidal-theta-family-trdeg-cone-ideal}
Let $o$ be a ring with a homomorphism $\mathrm{em}:o\to\mathbf C$, let $I$ be a set, let $m\ge0$, let
$B\subset\mathbf C^m$ be open, let $\vartheta:B\to\mathbf C^I$ be a map whose coordinates are complex differentiable on
$B$, and let $\mathfrak P$ be the cone ideal of $(B,\vartheta)$ (Definition~\ref{cv-toroidal-theta-family-trdeg-cone}).
\begin{enumerate}
\item If $o$ and $I$ are countable and $B$ is nonempty and connected, then $\mathfrak P$ is a prime ideal.
\item Suppose that $\mathrm{em}$ is injective, that $I$ is finite, that for every $z\in B$ only countably many
  $z'\in B$ have $\vartheta(z')\in\mathbf C\,\vartheta(z)$, and that $X_b\notin\mathfrak P$ for some $b\in I$. Then
  there are $a_0\in I$ with $X_{a_0}\notin\mathfrak P$ and $a_1,\dots,a_m\in I$ such that for every field $K$, every
  homomorphism $\beta:o\to K$ and every family $v\in K^I$ for which the kernel of the evaluation at $v$ over $\beta$ is
  $\mathfrak P$, the only polynomial $G\in o[Y_1,\dots,Y_m]$ with $\beta(G)(v_{a_1}/v_{a_0},\dots,v_{a_m}/v_{a_0})=0$
  is $G=0$.
\end{enumerate}
\end{lemma}
\leanref{Development/AbelianVarietyTheory/Proofs/ModuliProofs/ThetaIdealPrimeAvoidingVariable.lean}{GlobalTheory.ThetaPrime.coneIdeal_isPrime}
\leanref{Development/AbelianVarietyTheory/Proofs/ModuliProofs/ThetaIdealPrimeAvoidingVariable.lean}{GlobalTheory.ThetaPrime.exists_ratios}
\begin{proof}
(1) Lemma~\ref{cv-toroidal-theta-family-trdeg-general} gives $z\in B$ and $c\in\mathbf C^\times$ such that a polynomial
$F$ lies in $\mathfrak P$ exactly when $\mathrm{em}(F)(c\vartheta(z))=0$. So $\mathfrak P$ is the kernel of a homomorphism
to the field $\mathbf C$, and it is prime.

(2) As $X_b\notin\mathfrak P$, there are $z_1\in B$ and $c_1\in\mathbf C^\times$ with $c_1\vartheta_b(z_1)\ne0$. By
continuity there is an open ball $D\subset B$ around $z_1$ on which $\vartheta_b$ does not vanish. On $D$, $\vartheta$
does not vanish, and only countably many $z'\in D$ have $\vartheta(z')\in\mathbf C\,\vartheta(z)$, for $z\in D$. So
Lemma~\ref{cv-toroidal-theta-family-trdeg-ratios}, applied to $D$, gives $a_0,a_1,\dots,a_m$ and $z_0\in D$ with
$\vartheta_{a_0}(z_0)\ne0$, such that a polynomial $F\in\mathbf C[Y_1,\dots,Y_m]$ vanishing at
$(\vartheta_{a_i}(z)/\vartheta_{a_0}(z))_i$ for every $z\in D$ with $\vartheta_{a_0}(z)\ne0$ is $0$. Evaluation at
$\vartheta(z_0)$ shows $X_{a_0}\notin\mathfrak P$.

Let $K$, $\beta$ and $v$ be as in (2). Then $v_{a_0}\ne0$, as $X_{a_0}\notin\mathfrak P$. Let $G\in o[Y_1,\dots,Y_m]$
satisfy $\beta(G)(v_{a_1}/v_{a_0},\dots,v_{a_m}/v_{a_0})=0$. Let $e$ be the total degree of $G$ and
$H=X_{a_0}^eG(X_{a_1}/X_{a_0},\dots,X_{a_m}/X_{a_0})$, a polynomial in the $X_a$ with coefficients in $o$. For every
field $K'$, homomorphism $\beta':o\to K'$ and $w\in K'^I$ with $w_{a_0}\ne0$,
$\beta'(H)(w)=w_{a_0}^e\,\beta'(G)(w_{a_1}/w_{a_0},\dots,w_{a_m}/w_{a_0})$. So $\beta(H)(v)=0$, and $H\in\mathfrak P$.
For $z\in D$ with $\vartheta_{a_0}(z)\ne0$ this gives
$0=\mathrm{em}(H)(\vartheta(z))=\vartheta_{a_0}(z)^e\,\mathrm{em}(G)(\vartheta_{a_1}(z)/\vartheta_{a_0}(z),\dots)$, so
$\mathrm{em}(G)$ vanishes there. Hence $\mathrm{em}(G)=0$, and $G=0$ as $\mathrm{em}$ is injective.
\end{proof}
% Lean: coneIdeal_isPrime (exists_veryGeneral_point_of_holomorphic_family; the cone ideal is the kernel of evaluation
% at c • θ z); exists_ratios: mem_coneIdeal (z1, c1), Metric.mem_nhds_iff (the ball), exists_independent_ratios on the
% ball, exists_homogenisation (H, built by induction on G with X a0 ^ n; the text takes n the total degree) and
% MvPolynomial.map_injective.

\begin{lemma}[Alp\"oge]\label{cv-toroidal-theta-family-trdeg-prime}
In the setting of Definition~\ref{cv-toroidal-theta-family-good}, suppose that $\mathrm{em}$ is injective and that $o$
is a Dedekind domain of finite type over $\mathbf Z$, and let $\mathcal C$ be good for $((M_h),\kappa,w,w',\delta)$.
The hypothesis that $o$ is a Dedekind domain and conditions (1) and (2) of
Definition~\ref{cv-toroidal-theta-family-good} are not used in the proof. Let $\mathfrak p$ be a prime ideal of the
polynomial ring $o[X_c]$, $c\in(\mathbf Z/N'^2)^{2g}$, that contains the theta ideal of $\mathcal C$, and let $b$ be an
index with $X_b\notin\mathfrak p$. Then there are a prime ideal $\mathfrak P\subset\mathfrak p$ that contains the theta
ideal, an index $a_0$ with $X_{a_0}\notin\mathfrak P$ and indices $a_1,\dots,a_\delta$ such that for every field $K$,
every homomorphism $\beta:o\to K$ and every family $v$ in $K$ for which the kernel of the evaluation $o[X_c]\to K$ at
$v$ over $\beta$ is $\mathfrak P$, the only polynomial $G\in o[Y_1,\dots,Y_\delta]$ with
$\beta(G)(v_{a_1}/v_{a_0},\dots,v_{a_\delta}/v_{a_0})=0$ is $G=0$.
\end{lemma}
\leanref{Development/AbelianVarietyTheory/Proofs/ModuliProofs/ThetaIdealPrimeAvoidingVariable.lean}{GlobalTheory.ThetaPrime.exists_prime_independent_ratios}
\begin{proof}
The ring $o$ is countable, being a quotient of a polynomial ring in finitely many variables over $\mathbf Z$. Let
$B_1,\dots,B_r$ and $\vartheta_1,\dots,\vartheta_r$ be as in condition (3) of
Definition~\ref{cv-toroidal-theta-family-good}, and let $\mathfrak P_k$ be the cone ideal of $(B_k,\vartheta_k)$
(Definition~\ref{cv-toroidal-theta-family-trdeg-cone}).

The intersection of the $\mathfrak P_k$ lies in the theta ideal. Let $F$ lie in every $\mathfrak P_k$, let $t$ be a
member of $\mathcal C$ and let $v$ be a family of theta constants of $t$ at the identity; so $N'$ is odd
(Definition~\ref{cv-toroidal-theta}(2)). By condition (3) there are $k$
and $z\in B_k$ such that $\vartheta_k(z)$ is also one. Then $v=c\,\vartheta_k(z)$ for some $c\in\mathbf C^\times$. If
$\vartheta_k(z)\ne0$, its entries generate the unit ideal of $\mathbf C$, and Lemma~\ref{cv-toroidal-reduction-unit}
gives $c$. If $\vartheta_k(z)=0$, then $v=0$, since otherwise the same lemma, applied with $v$ in the first place,
would make $\vartheta_k(z)$ a nonzero multiple of $v$; take $c=1$. So $\mathrm{em}(F)(v)=0$, as $F\in\mathfrak P_k$.

The prime $\mathfrak p$ therefore contains the finite intersection of the $\mathfrak P_k$, and hence one of them,
$\mathfrak P=\mathfrak P_k$. Then $X_b\notin\mathfrak P$, so $B_k$ is nonempty, hence connected. By
Lemma~\ref{cv-toroidal-theta-family-trdeg-cone-ideal}(1), $\mathfrak P$ is prime, as $o$ and the set of indices are
countable. Part (2) of that lemma, with the countability in condition (3) and the injectivity of $\mathrm{em}$, gives
$a_0$ with $X_{a_0}\notin\mathfrak P$ and $a_1,\dots,a_\delta$ with the stated property.

Finally, $\mathfrak P$ contains the theta ideal. Let $F$ lie in the theta ideal, and let $z\in B_k$ and
$c\in\mathbf C^\times$. By condition (3), $\vartheta_k(z)$ is a family of theta constants at the identity of a member of
$\mathcal C$, so $N'$ is odd. The homogeneous components $F_e$ of $F$ lie in the theta ideal
(Lemma~\ref{cv-toroidal-theta-envelope-homogeneous}), so they vanish at $\vartheta_k(z)$, and
$\mathrm{em}(F)(c\vartheta_k(z))=\sum_ec^e\,\mathrm{em}(F_e)(\vartheta_k(z))=0$.
\end{proof}
% Lean: exists_prime_independent_ratios: ArBase gives only finite type and injectivity (⟨-, hft, hem⟩), and only
% hG.fam is used; countable_of_finiteType; hinf (the intersection of the coneIdeal of the B i, through hcov and
% hasTh_unique_unit, both cases of theta i z = 0); Ideal.IsPrime.inf_le' (one of them in p); coneIdeal_isPrime,
% exists_ratios; the theta ideal in the cone ideal by thIdeal_radical_homogeneous_torsionFree (homogeneity) and
% eval₂_smul_of_isHomogeneous.

\begin{lemma}[Alp\"oge]\label{cv-toroidal-theta-family-trdeg}
In the setting of Definition~\ref{cv-toroidal-theta-family-good}, let $n\ge0$ be an integer, suppose that $\mathrm{em}$ is
injective and that $o$ is a Dedekind domain of finite type over $\mathbf Z$, let $\mathcal C$ be good for
$((M_h),\kappa,w,w',\delta)$, and let $(\mathcal M,t)$ be a theta family for $\mathcal C$ and $n$. Let
$y:\operatorname{Spec}R\to\mathcal M$ be an open immersion over $o$, and let $\mathfrak q$ be a minimal prime of $R$.
The hypothesis that $o$ is a Dedekind domain and conditions (1) and (2) of
Definition~\ref{cv-toroidal-theta-family-good} are not used in the proof. Then the fraction field of $R/\mathfrak q$
contains $\delta$ elements that are algebraically independent over $o$.
\end{lemma}
\leanref{Development/AbelianVarietyTheory/Proofs/ModuliProofs/ThetaClosureComponentTrdegGeStatement.lean}{GlobalTheory.T.thetaClosure_component_trdeg_ge}
\leanref{Development/AbelianVarietyTheory/Proofs/ModuliProofs/ThetaClosureComponentTrdegGe.lean}{GlobalTheory.ZM.GenericTrdeg.exists_algebraicIndependent}
\begin{proof}
Let $K$ be the fraction field of $R/\mathfrak q$ and $\varphi:R\to K$ the induced homomorphism, with kernel
$\mathfrak q$; every element of $K$ is a quotient $\varphi(r)/\varphi(r')$. Let $\eta=y\circ\operatorname{Spec}\varphi$,
a point of $\mathcal M$ over $o\to K$.

A prime ideal. By property (b) of Definition~\ref{cv-toroidal-theta}(5), $t$ has at $\eta$ a family $v$ of theta
constants that is admissible for $(\mathcal C,n)$ over $o\to K$; so $N'$ is odd (Definition~\ref{cv-toroidal-theta}(2)),
as the lemmas on theta constants used below require. Let $e:o[X_c]\to K$ be the evaluation at $v$ and
$\mathfrak p$ its kernel, a prime ideal. It contains the theta ideal, as $v$ is admissible. The coset condition at the
prime $0$ of $K$, for $c=0$, gives an index $b$ with $v_b\ne0$, so $X_b\notin\mathfrak p$.
Lemma~\ref{cv-toroidal-theta-family-trdeg-prime} gives a prime $\mathfrak P\subset\mathfrak p$ that contains the theta
ideal and indices $a_0,a_1,\dots,a_\delta$ as there.

A point with values in a local domain. Let $R'$ be the localization of $o[X_c]/\mathfrak P$ at the prime
$\mathfrak p/\mathfrak P$, a local domain, let $\rho:o[X_c]\to R'$ be the induced homomorphism, with kernel
$\mathfrak P$, and let $\psi:R'\to K$ be the homomorphism induced by $e$, so that $\psi\circ\rho=e$. An element $r$ of
$R'$ with $\psi(r)\ne0$ is a unit. Let $\beta'$ be the restriction of $\rho$ to $o$ and $v'_c=\rho(X_c)$; then
$\psi(v'_c)=v_c$, and $\psi\circ\beta'$ is the structure map $o\to K$. The family $v'$ is admissible for $(\mathcal C,n)$
over $\beta'$. If $P$ lies in the theta ideal, then $\beta'(P)(v')=\rho(P)=0$, as $P\in\mathfrak P$. For every $c$ the
coset condition for $v$ at the prime $0$ gives $c'$ with $\psi(v'_{c+N'_nc'})=v_{c+N'_nc'}\ne0$, so $v'_{c+N'_nc'}$ is
a unit of $R'$ and lies in no prime ideal. Property (c) of Definition~\ref{cv-toroidal-theta}(5) gives a point
$x':\operatorname{Spec}R'\to\mathcal M$ over $\operatorname{Spec}\beta'$ at which $v'$ is a family of theta constants of
$t$. By Lemma~\ref{cv-toroidal-reduction-base-change}(4), $v=\psi(v')$ is one at $x'\circ\operatorname{Spec}\psi$, a
point over $o\to K$. By the uniqueness in property (c) for $v$, $x'\circ\operatorname{Spec}\psi=\eta$.

The point $x'$ factors through $y$. Every prime ideal of $R'$ is contained in $\psi^{-1}(0)$, as the elements outside
$\psi^{-1}(0)$ are units. So every point of $\operatorname{Spec}R'$ specializes to the point $\psi^{-1}(0)$, whose image
under $x'$ is the image of $\eta$, a point of the open image of $y$. As open sets are stable under generization, $x'$
maps into the image of $y$. So $x'=y\circ\operatorname{Spec}\chi$ for a homomorphism $\chi:R\to R'$. Then
$y\circ\operatorname{Spec}(\psi\circ\chi)=x'\circ\operatorname{Spec}\psi=y\circ\operatorname{Spec}\varphi$, and
$\psi\circ\chi=\varphi$, as $y$ is a monomorphism.

Comparison of the two families. Let $K'$ be the fraction field of $R'$ and $\lambda:R'\to K'$ the inclusion. The kernel
of $\lambda\circ\chi$ is a prime ideal of $R$ contained in the kernel $\mathfrak q$ of $\psi\circ\chi=\varphi$, which
is minimal; so it equals $\mathfrak q$. As $K$ is the fraction field of $R/\mathfrak q$, there is a homomorphism
$j:K\to K'$ with $j\circ\varphi=\lambda\circ\chi$. Then $\eta\circ\operatorname{Spec}j=x'\circ\operatorname{Spec}\lambda$,
and by Lemma~\ref{cv-toroidal-reduction-base-change}(4) both $j(v)$ and $\lambda(v')$ are families of theta constants
of $t$ at this point. As $j(v_b)\ne0$, the entries of $j(v)$ generate the unit ideal of $K'$, and
Lemma~\ref{cv-toroidal-reduction-unit} gives $u\in K'^\times$ with $\lambda(v'_c)=u\,j(v_c)$ for every $c$. The point
lies over $\operatorname{Spec}o$ through $j\circ(o\to K)$ and through $\lambda\circ\beta'$, so these two homomorphisms
agree. The kernel of the evaluation $o[X_c]\to K'$ at $\lambda(v')$ over $\lambda\circ\beta'$ is the kernel of
$\lambda\circ\rho$, which is $\mathfrak P$.

The independent elements. As $X_{a_0}\notin\mathfrak P$, $\lambda(v'_{a_0})\ne0$; so $j(v_{a_0})\ne0$ and
$v_{a_0}\ne0$. Let $s_i=v_{a_i}/v_{a_0}\in K$ for $i=1,\dots,\delta$, and let $G\in o[Y_1,\dots,Y_\delta]$ vanish at
$(s_1,\dots,s_\delta)$. Applying $j$, $G$ vanishes, over $\lambda\circ\beta'$, at the elements
$j(s_i)=\lambda(v'_{a_i})/\lambda(v'_{a_0})$ of $K'$. By the property of $a_0,\dots,a_\delta$ in
Lemma~\ref{cv-toroidal-theta-family-trdeg-prime}, applied to $K'$, $\lambda\circ\beta'$ and $\lambda(v')$, $G=0$. So
$s_1,\dots,s_\delta$ are algebraically independent over $o$.
\end{proof}
% Lean: ZM.GenericTrdeg.exists_algebraicIndependent: F.loc K at spMap phi ≫ x (v), hvI.2 at ⊥ for c = 0 (a1),
% ThetaPrime.exists_prime_independent_ratios (P, a0, a), exists_local_domain (R', rho, psi), hI' (admissibility of v'),
% F.pts R' (x'), uniqueness in F.pts K with hasTh_baseChange (hcl), exists_factor (chi, through
% IsOpenImmersion.lift), hker with the minimality of ker phi, exists_hom_of_ker_eq (j), hasTh_baseChange twice and
% hasTh_unique_unit (u), hjβ, hkP, and hind for K'. ArBase is used only through exists_prime_independent_ratios, and
% GoodClass only through its field fam.

\begin{lemma}[Alp\"oge]\label{cv-toroidal-theta-family-criterion}
Let $o$ be a domain with fraction field $F$ and an injective homomorphism $\mathrm{em}:o\to\mathbf C$, let $R$ be a
finitely generated $o$-algebra, and let $\delta\ge0$. Suppose that for every minimal prime $\mathfrak q$ of $R$ the
fraction field of $R/\mathfrak q$ contains $\delta$ elements algebraically independent over $o$, and that
$\dim_{\mathbf C}\operatorname{Der}_o(R,\mathbf C)\le\delta$ for every homomorphism of $o$-algebras $R\to\mathbf C$,
where $\mathbf C$ is an $o$-algebra through $\mathrm{em}$. Then $\operatorname{Spec}(F\otimes_oR)$ is smooth of
relative dimension $\delta$ over $F$.
\end{lemma}
\leanref{Development/AbelianVarietyTheory/Proofs/ModuliProofs/SmoothGenericFibreOfIndependentAndTangentBoundStatement.lean}{GlobalTheory.T.smooth_genericFibre_of_independent_and_tangent_bound}
\leanref{Development/AbelianVarietyTheory/Proofs/ModuliProofs/SmoothGenericFibreOfIndependentAndTangentBound.lean}{GlobalTheory.smoothOfRelativeDimension_genericFibre_of_independent_of_finrank_derivation_le}
\leanref{Development/AbelianVarietyTheory/Proofs/ModuliProofs/SmoothOfRelativeDimensionOfDimOfFinrankDerivatio.lean}{GlobalTheory.smoothOfRelativeDimension_of_dim_of_finrank_derivation_le}
\begin{proof}
Let $e:F\to\mathbf C$ be the extension of $\mathrm{em}$; then $F$ has characteristic zero. The ring
$S=F\otimes_oR$ is the localization of $R$ at the nonzero elements of $o$, and it is a finitely generated
$F$-algebra.

The components. Let $Q$ be a minimal prime of $S$. Its inverse image $\mathfrak q$ in $R$ is a minimal prime of $R$,
as $S$ is a localization of $R$, and $S/Q$ is a finitely generated $F$-algebra and a domain whose fraction field is
that of $R/\mathfrak q$. The $\delta$ elements of this field that are algebraically independent over $o$ are
algebraically independent over $F$, so its transcendence degree over $F$ is at least $\delta$, and so is the Krull
dimension of $S/Q$, which equals that transcendence degree.

The tangent spaces. Let $S\to\mathbf C$ be a homomorphism of $F$-algebras, $\mathbf C$ being an $F$-algebra through
$e$, and regard $\mathbf C$ as an $R$-algebra through $R\to S$. Restriction along $R\to S$ is an injective map
$\operatorname{Der}_F(S,\mathbf C)\to\operatorname{Der}_o(R,\mathbf C)$: if a derivation $D$ vanishes on the image
of $R$ and $s\in S$ satisfies $bs=a$ with $a$ the image of an element of $R$ and $b$ that of a nonzero element of
$o$, then $bD(s)=D(a)=0$, and $b$ is invertible in $\mathbf C$. So
$\dim_{\mathbf C}\operatorname{Der}_F(S,\mathbf C)\le\delta$.

Smoothness. Let $\mathfrak m$ be a maximal ideal of $S$. Some minimal prime $Q\subset\mathfrak m$ has
$\dim S/Q\ge\delta$, and in the finitely generated domain $S/Q$ over $F$ the maximal ideal $\mathfrak m/Q$ has height
$\dim S/Q$; so $\mathfrak m$ has height at least $\delta$. The field $S/\mathfrak m$ is finite over $F$, so $e$
extends to an embedding $S/\mathfrak m\to\mathbf C$; with the composite $S\to S/\mathfrak m\to\mathbf C$, the dual of
$\mathfrak m/\mathfrak m^2$ over $S/\mathfrak m$ injects into $\operatorname{Der}_F(S,\mathbf C)$, as $S/\mathfrak m$ is
separable over $F$. So $\dim\mathfrak m/\mathfrak m^2\le\delta\le\dim S_{\mathfrak m}$, and $S_{\mathfrak m}$ is a
regular local ring of dimension $\delta$. As $F$ is perfect, $S$ is smooth over $F$ at $\mathfrak m$, with
$\Omega_{S_{\mathfrak m}/F}$ free of rank $\delta$, so $S_f$ is smooth of relative dimension $\delta$ over $F$ for some
$f\notin\mathfrak m$. These $\operatorname{Spec}S_f$ cover $\operatorname{Spec}S$.
\end{proof}
% Lean: smoothOfRelativeDimension_genericFibre_of_independent_of_finrank_derivation_le (e = IsFractionRing.lift hem,
% IsLocalization.tensorRight), GenericSmooth.le_ringKrullDim_of_mem_minimalPrimes (the components, with
% exists_ringKrullDim_eq_and_trdeg_eq), GenericSmooth.finrank_derivation_le (the tangent spaces), and
% smoothOfRelativeDimension_of_dim_of_finrank_derivation_le: SmoothTangent.le_height_of_forall_minimalPrimes (with
% CA.height_eq_of_isMaximal_of_trdeg_eq_field), SmoothTangent.finrank_cotangent_le_finrank_derivation
% (kerCotangentToTensor_injective), isSmoothAt_of_finrank_cotangent_le_of_perfectField and
% exists_isStandardSmoothOfRelativeDimension_of_isSmoothAt. About 530 lines.

\begin{lemma}[Alp\"oge]\label{cv-toroidal-theta-family-smooth}
In the setting of Lemma~\ref{cv-toroidal-theta-family-trdeg}, suppose that
$\dim_{\mathbf C}\operatorname{Der}_o(R,\mathbf C)\le\delta$ for every open immersion
$y:\operatorname{Spec}R\to\mathcal M$ over $o$ and every homomorphism of $o$-algebras $R\to\mathbf C$, where
$\mathbf C$ is an $o$-algebra through $\mathrm{em}$. Then the fibre of $\mathcal M$ over the fraction field $F$ of
$o$ is smooth of relative dimension $\delta$.
\end{lemma}
\leanref{Development/AbelianVarietyTheory/Proofs/ModuliProofs/FiniteTypeOfIsQProjOverAffineOpen.lean}{GlobalTheory.ZM.ThetaFam.smooth_genericFibre}
\begin{proof}
Being smooth of relative dimension $\delta$ is local on the source, and the fibres over $F$ of the open subschemes
$\operatorname{Spec}R$ of $\mathcal M$ cover the fibre of $\mathcal M$; that of $\operatorname{Spec}R$ is
$\operatorname{Spec}(F\otimes_oR)$. So it suffices that $F\otimes_oR$ is smooth of relative dimension $\delta$ over
$F$ for every open immersion $y:\operatorname{Spec}R\to\mathcal M$ over $o$. The ring $o$ is a domain, and $R$ is a
finitely generated $o$-algebra, because $\mathcal M$ is quasi-projective, hence locally of finite type, over $o$.
By Lemma~\ref{cv-toroidal-theta-family-trdeg}, for every minimal prime $\mathfrak q$ of $R$ the fraction field of
$R/\mathfrak q$ contains $\delta$ elements algebraically independent over $o$, and by hypothesis
$\dim_{\mathbf C}\operatorname{Der}_o(R,\mathbf C)\le\delta$ for every homomorphism of $o$-algebras $R\to\mathbf C$.
So Lemma~\ref{cv-toroidal-theta-family-criterion} applies.
\end{proof}
% Lean: ThetaFam.smooth_genericFibre: T.smooth_genericFibre_of_affine_opens (smoothOfRelativeDimension_genericFibre_of_affine,
% by IsZariskiLocalAtSource over the affine open cover), T.smooth_genericFibre_of_independent_and_tangent_bound,
% finiteType_of_isQProjOver_affineOpen and T.thetaClosure_component_trdeg_ge. Size: about 2,500 lines with Lemmas
% cv-toroidal-theta-family-trdeg and -criterion.

\begin{definition}\label{cv-toroidal-theta-family-tangent-good}
Let $A_0$ be an abelian variety over $\mathbf C$ of dimension $g$ with a line bundle $L_0$, and let $u:\mathbf C^g\to
A_0(\mathbf C)$ be a uniformization of $A_0$, with lattice $\Lambda_u$ and trivializing functions $\theta_\alpha$ for
$L_0$, with Appell--Humbert data $(H,\chi)$ (Definition~\ref{cv-origin-hodge-norm-gram}(1), with $T=\operatorname{Spec}
\mathbf C$, $G=A_0$ and $\varphi$ the identity). The uniformization is \emph{good} if:
\begin{enumerate}
\item every morphism $v:\operatorname{Spec}\mathbf C[\epsilon]\to A_0$ over $\mathbf C$ that restricts to the origin
  has a unique differential $w\in\mathbf C^g$, and every $w\in\mathbf C^g$ is the differential of a unique such $v$;
\item for every endomorphism $f$ of the group scheme $A_0$ there is a $\mathbf C$-linear map $F$ of $\mathbf C^g$ with
  $u\circ F=f\circ u$;
\item for every $n\ge0$ and every holomorphic function $\vartheta$ on $\mathbf C^g$ with
  $\vartheta(z+l)=\big(\chi(l)\exp(\pi z^{\mathsf T}H\bar l+\tfrac\pi2l^{\mathsf T}H\bar l)\big)^n\vartheta(z)$ for
  $l\in\Lambda_u$, there is a section $s$ of $L_0^n$ such that $s\circ u=\vartheta\,\theta_\alpha^n$ on each chart.
\end{enumerate}
\end{definition}
\leanref{Development/AbelianVarietyTheory/Definitions/ModuliSpaces/IsSymm.lean}{GlobalTheory.KS.GoodU}

\begin{lemma}[Alp\"oge]\label{cv-toroidal-theta-family-tangent-ks}
Let $g\ge1$ and $N'\ge1$, let $o$ be a ring with a unit $\zeta$ as in Definition~\ref{cv-toroidal-theta}, let $R$ be an
$o$-algebra with a homomorphism of $o$-algebras $R\to\mathbf C$, for some structure of $o$-algebra on $\mathbf C$, and
let $t=(A,L,(x_i),(\alpha_h))$ be a theta-level object over $\operatorname{Spec}R$. Suppose that for every point
$x:\operatorname{Spec}\mathbf C[\epsilon]\to\operatorname{Spec}R$ with the same composite to $\operatorname{Spec}o$ as
the constant point $x_0$, if a theta-level object over $\mathbf C[\epsilon]$ has morphisms to $t$ over $x$ and over
$x_0$, then $x=x_0$. Then there are:
\begin{enumerate}
\item a theta-level object $t_0=(A_0,L_0,(x^0_i),(\alpha^0_h))$ over $\mathbf C$ with a morphism $m:t_0\to t$ over the
  point $\operatorname{Spec}\mathbf C\to\operatorname{Spec}R$;
\item a good uniformization $u$ of $A_0$ with Appell--Humbert data $(H,\chi)$
  (Definition~\ref{cv-toroidal-theta-family-tangent-good});
\item an injective $\mathbf C$-linear map $\kappa:\operatorname{Der}_o(R,\mathbf C)\to M_g(\mathbf C)$ such that for every
  $D$ the matrix $\Theta=\kappa(D)$ satisfies $h(\Theta\bar z,w)=h(\Theta\bar w,z)$ for all $z,w\in\mathbf C^g$, where
  $h(z,w)=z^{\mathsf T}H\bar w$, and $F(\Theta\bar z)=\Theta\,\overline{F(z)}$ for every morphism $f:A\to A$ over $R$,
  every morphism $f_0:A_0\to A_0$ over $\mathbf C$ with $m\circ f_0=f\circ m$ and every $\mathbf C$-linear $F$ with
  $u\circ F=f_0\circ u$.
\end{enumerate}
\end{lemma}
\leanref{Development/AbelianVarietyTheory/Proofs/ModuliProofs/TangentInjectsSymmetricAntilinearCommutingStatement.lean}{GlobalTheory.T.tangent_injects_symmetric_antilinear_commuting}
\begin{proof}
\openstep{The construction of the Kodaira--Spencer map $\kappa$, from the deformations of the uniformized complex fibre
along the derivations of $R$, its injectivity, from the hypothesis on points with values in $\mathbf C[\epsilon]$, and
the symmetry and commutation of its values, are not yet written out.}
% Lean: T.tangent_injects_symmetric_antilinear_commuting, with Out.Laws, Out.Laws2 and ThetaUnifAV (the analytic
% uniformization of principally polarized complex abelian varieties, T.theta_uniformisation_complex_abelian_varieties) as
% hypotheses; KSRel is the Lean's relation between a derivation and its matrix. Size, counted through the declarations
% reached in the imported files beyond the laws of abelian schemes: about 8,900 lines.
\end{proof}

\begin{lemma}[Alp\"oge]\label{cv-toroidal-theta-family-tangent-fibre}
In Situation~\ref{cv-function-situation}, let $e_1,\dots,e_{2g}$ be a basis of the lattice $\Lambda$ of the datum over
$\mathbf Z$, let $a:I\to\mathcal O_E$ be a map, let $N'\ge1$, let $o$ be a ring with a unit $\zeta$ as in
Definition~\ref{cv-toroidal-theta}, and let $\rho\in o$. Let $S$ be a scheme over $o$, let $t=(A,L,(x_i),(\alpha_h))$ be
a theta-level object over $S$ with a PEL structure for $((e_i),a,\rho)$ (Definition~\ref{cv-toroidal-theta}(6)), let
$x:\operatorname{Spec}\mathbf C\to S$ be a point over $\beta:o\to\mathbf C$, and let $A_0$ be a group scheme over
$\mathbf C$ with a morphism $m:A_0\to A$ over $x$ that identifies $A_0$ with $x^*A$. Then there are actions of
$\mathcal O_E$ on $A$ and on $A_0$, compatible through $m$, such that condition (a) of
Definition~\ref{cv-toroidal-family}(1) holds for $A_0$ with this action at the identity point of
$\operatorname{Spec}\mathbf C$, with $\beta(\rho)$ in place of $\beta(\theta)$.
\end{lemma}
\leanref{Development/AbelianVarietyTheory/Proofs/ModuliProofs/ExistsFibreAct.lean}{GlobalTheory.FibreAct.exists_fibre_act}
\begin{proof}
Let $G$ be the PEL structure, with its isomorphism of group schemes $\varphi:G\to A$. If a group scheme $G'$ over a scheme
$S'$ has a morphism $\mu:G'\to G''$ over $f:S'\to S''$ that identifies it with $f^*G''$, then an action of
$\mathcal O_E$ on $G''$ gives one on $G'=f^*G''$, by base change of the endomorphisms, with $\mu\circ c=c\circ\mu$; the
relations of an action hold, as base change preserves sums and composites of homomorphisms. Applied to
$\varphi^{-1}:A\to G$ over the identity of $S$, this gives an action of $\mathcal O_E$ on $A$ compatible with that on $G$
through $\varphi^{-1}$; applied to $m$, an action on $A_0$ compatible with that on $A$ through $m$.

The Lie algebra of $A_0$ at the identity point is identified, through $m$ and then $\varphi^{-1}$, with the Lie algebra of
$G$ at $x$, compatibly with the actions of $\mathcal O_E$; so for every ring homomorphism $\tau:\mathcal O_E\to\mathbf C$
the two subspaces on which each $c$ acts by $\tau(c)$ have the same dimension. As $\mathbf C$ is algebraically closed of
characteristic zero, the PEL structure satisfies condition (a) of Definition~\ref{cv-toroidal-family}(1) at $x$ with
$\beta(\rho)$, and so does $A_0$ at the identity point.
\end{proof}
% Lean: FibreAct.exists_fibre_act: isGrpPullback_of_iso (phi^{-1} over the identity), exists_lift_act twice (the
% actions, through LiePull.exists_pushEquiv), LiePull.lieMult_pull twice (the Lie multiplicities), and the condition of
% IsPEL at x. About 180 lines, counted as for Lemma cv-toroidal-theta-family-tangent-ks.

\begin{lemma}[Alp\"oge]\label{cv-toroidal-theta-family-tangent-lift}
Let $A_0$ be an abelian variety over $\mathbf C$ of dimension $g$ with a line bundle $L_0$ that is ample, trivialized at
the origin, symmetric and defines a principal polarization, and let $u$ be a good uniformization of $A_0$ with
Appell--Humbert data $(H,\chi)$ (Definition~\ref{cv-toroidal-theta-family-tangent-good}).
\begin{enumerate}
\item The linear map $F$ of Definition~\ref{cv-toroidal-theta-family-tangent-good}(2) is unique. For every commutative
  ring $O$ acting on $A_0$ there is a ring homomorphism $F:O\to\operatorname{End}_{\mathbf C}(\mathbf C^g)$ with
  $u\circ F(c)=c\circ u$ for every $c\in O$.
\item Let $O$ be a commutative ring acting on $A_0$, let $F:O\to\operatorname{End}_{\mathbf C}(\mathbf C^g)$ be a map with
  $u\circ F(c)=c\circ u$ for every $c$, let $\tau:O\to\mathbf C$ be a ring homomorphism and $n\ge0$, and suppose that the
  subspace of the Lie algebra of $A_0$ on which each $c\in O$ acts by $\tau(c)$ has dimension $n$. Then the subspace of
  $v\in\mathbf C^g$ with $F(c)v=\tau(c)v$ for every $c$ has dimension $n$.
\end{enumerate}
\end{lemma}
\leanref{Development/AbelianVarietyTheory/Proofs/ModuliProofs/EqZeroOfLineSubsetDiscrete.lean}{GlobalTheory.KS.exists_linear_lift_action}
\leanref{Development/AbelianVarietyTheory/Proofs/ModuliProofs/FinrankEigenspaceOfLieMultiplicityStatement.lean}{GlobalTheory.T.finrank_eigenspace_of_lieMultiplicity}
\begin{proof}
Let $\Lambda_u$ be the kernel of $u$, a discrete subgroup of $\mathbf C^g$.

(1) Let $F,F'$ be linear maps with $u\circ F=f\circ u=u\circ F'$ for a morphism $f:A_0\to A_0$. For $z\in\mathbf C^g$,
$u(Fz)=u(F'z)$, so $Fz-F'z\in\Lambda_u$, as $u$ is a homomorphism; and then $t(Fz-F'z)=F(tz)-F'(tz)\in\Lambda_u$ for every
$t\in\mathbf C$. The discrete subgroup $\Lambda_u$ is countable; if $Fz\ne F'z$, then $t\mapsto t(Fz-F'z)$ would
be an injective map from the uncountable set $\mathbf C$ to $\Lambda_u$. So $F=F'$. Now let $O$ act on $A_0$. For $c\in O$,
condition (2) of Definition~\ref{cv-toroidal-theta-family-tangent-good}, applied to the endomorphism $c$, gives a linear
map $F(c)$ with $u\circ F(c)=c\circ u$. By uniqueness, $F(1)=1$, $F(cc')=F(c)F(c')$, as $u\circ F(c)F(c')=c\circ
u\circ F(c')=cc'\circ u$, and $F(c+c')=F(c)+F(c')$, as $u(F(c)z+F(c')z)=u(F(c)z)+u(F(c')z)=(c+c')(u(z))$. So $F$ is a
ring homomorphism.

(2) Tangent vectors of $A_0$ at the origin are morphisms $v:\operatorname{Spec}\mathbf C[\varepsilon]\to A_0$ over
$\mathbf C$ that restrict to the origin; they form the Lie algebra, with the sum given by the group law and $s\in\mathbf C$
acting through $\varepsilon\mapsto s\varepsilon$. By condition (1) of
Definition~\ref{cv-toroidal-theta-family-tangent-good}, the map that sends $v$ to its differential
(Definition~\ref{cv-origin-hodge-norm-gram}(1)) is a bijection from tangent vectors to $\mathbf C^g$. It is linear. For $s\in\mathbf C$, the coefficient of $\varepsilon$
in $f\circ v$ is multiplied by $s$ when $\varepsilon$ is replaced by $s\varepsilon$, so $sv$ has differential $s$ times
that of $v$. For tangent vectors $v,v'$, the map that sends a function $f$ near the origin of $A_0\times A_0$ to the
coefficient of $\varepsilon$ in $f\circ(v,v')$ is the sum of those for $(v,0)$ and $(0,v')$, the tangent space of a
product being the direct sum; applied to $f\circ m$, where $m$ is the group law, with $m\circ(v,0)=v$ and
$m\circ(0,v')=v'$, this shows that the differential of $v+v'$ is the sum of those of $v$ and $v'$.

Let $v$ be a tangent vector with differential $z$, and $c\in O$. The tangent vector $c\circ v$ has differential $F(c)z$:
for a function $f$ near the origin, the coefficient of $\varepsilon$ in $f\circ c\circ v$ is the derivative at $0$ of
$f\circ c\circ u=f\circ u\circ F(c)$ in the direction $z$, which is the derivative of $f\circ u$ in the direction $F(c)z$,
as $F(c)$ is linear. The tangent vector $\tau(c)v$ has differential $\tau(c)z$. By the uniqueness in condition (1),
$c\circ v=\tau(c)v$ if and only if $F(c)z=\tau(c)z$. So the linear bijection maps the subspace of the Lie algebra on which
each $c$ acts by $\tau(c)$ onto the subspace of the $z$ with $F(c)z=\tau(c)z$ for every $c$, which therefore has
dimension $n$.
\end{proof}
% Lean: (1) KS.exists_linear_lift_action: linear_eq_of_ex_eq (eq_zero_of_line_subset_discrete: a discrete subgroup is
% countable), then hone, hmul and hadd. (2) UnifEigenRank.finrank_eigenspace_of_lieMult: LiePull.basis_transport (the
% basis v of LieMult), key (with UnifLieBasis.isDiff_comp, isDiff_smul, isTangent_comp and the uniqueness of GoodU.tan1
% and tan2), isDiff_list_prod (the differential of a sum) and finrank_of_existsUnique.

\begin{lemma}[Alp\"oge]\label{cv-toroidal-theta-family-tangent-dimension}
In Situation~\ref{cv-function-situation}, let $g,p,q\ge0$, let $\rho_0\in\mathbf C$ with $\rho_0^2=-d$, let $H\in
M_g(\mathbf C)$ be a Hermitian matrix with $\operatorname{Re}h(z,z)>0$ for $z\ne0$, where $h(z,w)=z^{\mathsf T}H\bar w$,
and let $F:\mathcal O_E\to\operatorname{End}_{\mathbf C}(\mathbf C^g)$ be a ring homomorphism. For a ring homomorphism
$\tau:\mathcal O_E\to\mathbf C$ let $V_\tau$ be the space of $v\in\mathbf C^g$ with $F(c)v=\tau(c)v$ for every $c$, and
suppose that $\dim V_\tau=p$ if $\tau(\sqrt{-d})=\rho_0$ and $\dim V_\tau=q$ otherwise. Let $S$ be a space of matrices
$\Theta\in M_g(\mathbf C)$ with $h(\Theta\bar z,w)=h(\Theta\bar w,z)$ and $F(c)(\Theta\bar z)=\Theta\,\overline{F(c)z}$
for all $z,w\in\mathbf C^g$ and $c\in\mathcal O_E$. The hypothesis that $H$ is Hermitian is not used in the proof. Then
$\dim_{\mathbf C}S\le\frac12[E:\mathbf Q]pq$.
\end{lemma}
\leanref{Development/AbelianVarietyTheory/Proofs/ModuliProofs/FinrankSymmAntilinearCommCMLeStatement.lean}{GlobalTheory.T.finrank_symm_antilinear_comm_CM_le}
\leanref{Development/AbelianVarietyTheory/Proofs/ModuliProofs/FinrankSymmAntilinearCommCMLe.lean}{GlobalTheory.KS.DimPEL.finrank_symm_comm_le}
\begin{proof}
Ring homomorphisms $\tau:\mathcal O_E\to\mathbf C$ are the restrictions of the embeddings of $E$. For each of them
$\tau(\sqrt{-d})^2=-d<0$, so $\tau(\sqrt{-d})$ is purely imaginary and $\bar\tau(\sqrt{-d})=-\tau(\sqrt{-d})$; and
$\rho_0\ne0$. Let $\Phi$ be the set of $\tau$ with $\tau(\sqrt{-d})=\rho_0$. If $\tau\in\Phi$, then
$\bar\tau(\sqrt{-d})=-\rho_0\ne\rho_0$, so $\bar\tau\notin\Phi$; if $\tau\notin\Phi$, then $\tau(\sqrt{-d})=-\rho_0$ and
$\bar\tau\in\Phi$. Complex conjugation thus maps $\Phi$ injectively into its complement, and $\Phi$ has at most
$\frac12[E:\mathbf Q]$ elements.

For $\Theta\in S$ let $\psi_\Theta(z)=\Theta\bar z$, an antilinear map. If $v\in V_\tau$, then
$F(c)\psi_\Theta(v)=\psi_\Theta(\tau(c)v)=\overline{\tau(c)}\psi_\Theta(v)$, so $\psi_\Theta(V_\tau)\subset V_{\bar\tau}$.
For $\tau\in\Phi$ choose a basis $b_{\tau,1},\dots,b_{\tau,p}$ of $V_\tau$. The linear map from $S$ to the product over
$\tau\in\Phi$ of $V_{\bar\tau}^p$, given by $\Theta\mapsto(\psi_\Theta(b_{\tau,i}))_{\tau,i}$, has a target of dimension
$|\Phi|pq\le\frac12[E:\mathbf Q]pq$. It is injective. Let $\Theta$ map to $0$ and $\psi=\psi_\Theta$. Then $\psi$
vanishes on $V_\tau$ for $\tau\in\Phi$, as it is antilinear. For $\tau\notin\Phi$ and $v\in V_\tau$, $\psi(v)$ lies in
$V_{\bar\tau}$ with $\bar\tau\in\Phi$, so $\psi(\psi(v))=0$, and $h(\psi(v),\psi(v))=h(\psi(\psi(v)),v)=0$ by the
symmetry of $\Theta$; so $\psi(v)=0$. Finally every $v\in\mathbf C^g$ is a sum of elements of the $V_\tau$: through a
$\mathbf Z$-basis of $\mathcal O_E$, $F$ extends to a $\mathbf C$-linear homomorphism from
$\mathbf C\otimes\mathcal O_E\cong\mathbf C^{\operatorname{Hom}(E,\mathbf C)}$, and the images of its standard idempotents
project $v$ to the $V_\tau$. So $\psi=0$, and $\Theta=0$.
\end{proof}
% Lean: DimPEL.finrank_symm_comm_le (the hypothesis Hm.IsHermitian is named _hH and not used): conj_eq_neg, hne, heq and
% hcard (the set T of tau with tau(rt) = rho0), theta_mem (psi maps Eig tau to Eig (conjugate tau)), Psi through the bases
% b, hinj with k1, k2 (hpos and the symmetry) and k3 (exists_sum_eigen, through canonicalEmbedding.latticeBasis), and
% LinearMap.finrank_le_finrank_of_injective.

\begin{lemma}[Alp\"oge]\label{cv-toroidal-theta-family-tangent}
In Situation~\ref{cv-function-situation}, suppose $p+q\ge1$, let $g=\frac12[E:\mathbf Q](p+q)$ and
$\delta=\frac12[E:\mathbf Q]pq$, let $e_1,\dots,e_{2g}$ be a symplectic basis of the lattice $\Lambda$ of the datum,
and let $a:I\to\mathcal O_E$ be a map from a finite set whose image generates $\mathcal O_E$ as a ring, with
elements $w,w'\in I$ such that $a(w)=w_0$ and $a(w')=\bar w_0$; these hypotheses on $p+q$, the basis and $a$ are
not used in the proof. Let $N'\ge1$, let $o$ be a ring with a unit $\zeta$ as in
Definition~\ref{cv-toroidal-theta}, and let $\rho\in o$ with $\rho^2=-d$. Let $R$ be a finitely generated
$o$-algebra with a homomorphism of $o$-algebras $R\to\mathbf C$, for some structure of $o$-algebra on $\mathbf C$,
and let $t$ be a theta-level object over $\operatorname{Spec}R$ with a PEL structure for $((e_i),a,\rho)$; that $R$
is finitely generated is not used in the proof either. Suppose
that for every point $x:\operatorname{Spec}\mathbf C[\epsilon]\to\operatorname{Spec}R$ with the same composite to
$\operatorname{Spec}o$ as the constant point $x_0$ (as in Definition~\ref{cv-toroidal-family}(2)(b)), if a
theta-level object over $\mathbf C[\epsilon]$ has morphisms to $t$ over $x$ and over $x_0$, then $x=x_0$. Then
$\dim_{\mathbf C}\operatorname{Der}_o(R,\mathbf C)\le\delta$.
\end{lemma}
\leanref{Development/AbelianVarietyTheory/Proofs/ModuliProofs/FinrankTangentPELFamilyLeSymmetricDomainStatement.lean}{GlobalTheory.T.finrank_tangent_PEL_family_le_symmetricDomain}
\leanref{Development/AbelianVarietyTheory/Proofs/ModuliProofs/FinrankTangentPELFamilyLeSymmetricDomain.lean}{GlobalTheory.PelTan.finrank_derivation_le}
\begin{proof}
Write $t=(A,L,(x_i),(\alpha_h))$. Lemma~\ref{cv-toroidal-theta-family-tangent-ks} gives a theta-level object
$t_0=(A_0,L_0,(x^0_i),(\alpha^0_h))$ over $\mathbf C$ with a morphism $m:t_0\to t$ over the point
$x:\operatorname{Spec}\mathbf C\to\operatorname{Spec}R$, a good uniformization $u$ of $A_0$ with Appell--Humbert data
$(H,\chi)$, and an injective linear map $\kappa:\operatorname{Der}_o(R,\mathbf C)\to M_g(\mathbf C)$ as there. The point
$x$ lies over $\beta:o\to\mathbf C$, the structure map of $\mathbf C$, as $R\to\mathbf C$ is a homomorphism of
$o$-algebras.

Lemma~\ref{cv-toroidal-theta-family-tangent-fibre}, applied to the PEL structure on $t$, $x$, $A_0$ and $m$, gives actions
of $\mathcal O_E$ on $A$ and on $A_0$, compatible through $m$, such that condition (a) of
Definition~\ref{cv-toroidal-family}(1) holds for $A_0$ at the identity point with $\beta(\rho)$. As $A_0$ with $L_0$ is a
principally polarized abelian variety, Lemma~\ref{cv-toroidal-theta-family-tangent-lift}(1) gives a ring homomorphism
$F:\mathcal O_E\to\operatorname{End}_{\mathbf C}(\mathbf C^g)$ with $u\circ F(c)=c\circ u$ for every $c$, and part (2)
of that lemma shows that for every ring homomorphism $\tau:\mathcal O_E\to\mathbf C$ the space of $v$ with
$F(c)v=\tau(c)v$ for every $c$ has dimension $p$ if $\tau(\sqrt{-d})=\beta(\rho)$ and $q$ otherwise.

Let $S$ be the image of $\kappa$. For $D$ a derivation and $\Theta=\kappa(D)$, $h(\Theta\bar z,w)=h(\Theta\bar w,z)$, and,
applying Lemma~\ref{cv-toroidal-theta-family-tangent-ks}(3) to the action of $c$ on $A$, its action on $A_0$, which is
compatible through $m$, and $F(c)$, $F(c)(\Theta\bar z)=\Theta\,\overline{F(c)z}$ for every $c\in\mathcal O_E$. The
Appell--Humbert matrix $H$ is Hermitian with $\operatorname{Re}h(z,z)>0$ for $z\ne0$, and $\beta(\rho)^2=\beta(-d)=-d$.
So Lemma~\ref{cv-toroidal-theta-family-tangent-dimension}, with $\rho_0=\beta(\rho)$, gives
$\dim S\le\frac12[E:\mathbf Q]pq=\delta$. As $\kappa$ is injective,
$\dim_{\mathbf C}\operatorname{Der}_o(R,\mathbf C)=\dim S\le\delta$.
\end{proof}
% Lean: PelTan.finrank_derivation_le: T.tangent_injects_symmetric_antilinear_commuting (t0, m, u, Hm, chi, kappa),
% FibreAct.exists_fibre_act (act, act0 with hc and hdet), KS.exists_linear_lift_action (F), T.finrank_symm_antilinear_comm_CM_le
% with hG.ah (Hermitian and positive), T.finrank_eigenspace_of_lieMultiplicity (the dimensions p and q) and hD for
% the range of kappa; LinearMap.finrank_range_of_inj. The theorem of the cube and Algebra.FiniteType o R are
% introduced as _ and not used, as are 1 <= p + q, IsSympBasis, a w = w0, a w' = bar w0 and the closure of the
% range of a. Size: about 9,900 lines, about 7,900 of them not shared with the other inputs of Lemma
% cv-toroidal-theta-family; about 8,900 are those of Lemma cv-toroidal-theta-family-tangent-ks.

\begin{lemma}[Alp\"oge]\label{cv-toroidal-theta-family}
In Situation~\ref{cv-function-situation}, suppose $p+q\ge1$, let $g=\frac12[E:\mathbf Q](p+q)$, let
$e_1,\dots,e_{2g}$ be a symplectic basis of the lattice $\Lambda$ of the datum, let $a:I\to\mathcal O_E$
be a map from a finite set whose image generates $\mathcal O_E$ as a ring, with elements $w,w'\in I$
such that $a(w)=w_0=(1+\sqrt{-d})/2$ and $a(w')=\bar w_0$, and let $n$ be an odd prime. There is $j_0$
such that for every $j\ge j_0$, with $N'=n^j$, the following holds. Let $o$ be a ring with a unit
$\zeta$ such that $N'$ is invertible in $o$ and every homomorphism from $o$ to a field maps $\zeta$ to a
primitive root of unity of order $N'^2$, let $\mathrm{em}:o\to\mathbf C$ be an injective homomorphism
such that $o$ is a Dedekind domain of finite type over $\mathbf Z$, and let $\rho\in o$ with
$\mathrm{em}(\rho)=\theta$. Then there are a theta family $(\mathcal M,t)$ for $n$ and the PEL class of
$((e_i),a,\rho,\mathrm{em})$ (Definition~\ref{cv-toroidal-theta}) and a PEL structure on $t$, and the
fibre of $\mathcal M$ over the fraction field of $o$ is smooth of relative dimension
$\frac12[E:\mathbf Q]pq$.
\end{lemma}
\leanref{Development/AbelianVarietyTheory/Proofs/ModuliProofs/ExistsThetaFamilyPelSmoothGenericFibreStatement.lean}{GlobalTheory.T.exists_thetaFamily_pel_smooth_generic_fibre}
\leanref{Development/AbelianVarietyTheory/Proofs/ModuliProofs/ExistsThetaFamilyPelSmoothGenericFibre.lean}{GlobalTheory.ZM.exists_thetaFam_isPEL_smooth}
\begin{proof}
Let $\delta=\frac12[E:\mathbf Q]pq$. Lemma~\ref{cv-toroidal-theta-family-pel-good} gives integer matrices $(M_h)$
and an integer $\kappa$, and for the prime $n$ an integer $j_1$. Lemma~\ref{cv-toroidal-theta-family-exists}, for
$g$, $I$, $w$, $w'$, $(M_h)$, $\kappa$ and $n$, gives an integer $j_2$. Let $j_0=\max(j_1,j_2)$, let $j\ge j_0$ and
$N'=n^j$, let $o$, $\zeta$, $\mathrm{em}$ and $\rho$ be as in the statement, and let $\mathcal C$ be the PEL class of
$((e_i),a,\rho,\mathrm{em})$. By the choice of $j_1$, $\mathcal C$ is good for $((M_h),\kappa,w,w',\delta)$, so by
the choice of $j_2$ there is a theta family $(\mathcal M,t)$ for $\mathcal C$ and $n$.

The PEL structure. The scheme $\mathcal M$ is reduced, and it is quasi-projective, hence locally of finite type,
over the noetherian ring $o$; so it is locally noetherian. By property (a) of Definition~\ref{cv-toroidal-theta}(5)
the points at which $t$ is a $\mathcal C$-object are dense, and $\mathrm{em}$ is injective. So
Lemma~\ref{cv-toroidal-theta-family-pel} gives a PEL structure $G$ on $t$ for $((e_i),a,\rho)$.

The generic fibre. By Lemma~\ref{cv-toroidal-theta-family-smooth}, it suffices to show that
$\dim_{\mathbf C}\operatorname{Der}_o(R,\mathbf C)\le\delta$ for every open immersion
$y:\operatorname{Spec}R\to\mathcal M$ over $o$ and every homomorphism of $o$-algebras $R\to\mathbf C$, where
$\mathbf C$ is an $o$-algebra through $\mathrm{em}$. Lemma~\ref{cv-toroidal-reduction-base-change}(1) gives a
theta-level object $t_R$ over $\operatorname{Spec}R$ with a morphism $t_R\to t$ over $y$. The pullback of $G$ along
$y$, with the isomorphism to the abelian scheme of $t_R$ induced by that of $G$, is a PEL structure on $t_R$: the
action of $\mathcal O_E$ and the line bundle pull back, and the conditions at points of characteristic zero concern
the fibres at those points, which the pullback identifies with fibres of $G$ and of $t$. In $o$, $\rho^2=-d$, since
$\mathrm{em}(\rho^2)=\theta^2=-d$ and $\mathrm{em}$ is injective. The $o$-algebra $R$ is finitely generated, as
$\mathcal M$ is locally of finite type over $o$. Finally, let $x:\operatorname{Spec}\mathbf C[\epsilon]\to
\operatorname{Spec}R$ have the same composite to $\operatorname{Spec}o$ as the constant point $x_0$, and let a
theta-level object $t'$ over $\mathbf C[\epsilon]$ have morphisms to $t_R$ over $x$ and over $x_0$. Composed with
$t_R\to t$, they give morphisms $t'\to t$ over $y\circ x$ and over $y\circ x_0$, which is the constant point of
$y\circ x$ and has the same composite to $\operatorname{Spec}o$. Property (e) of
Definition~\ref{cv-toroidal-theta}(5) gives $y\circ x=y\circ x_0$, and so $x=x_0$, as the open immersion $y$ is a
monomorphism. Lemma~\ref{cv-toroidal-theta-family-tangent}, applied to $o$, $\rho$, $R$, $t_R$ and this PEL
structure, gives $\dim_{\mathbf C}\operatorname{Der}_o(R,\mathbf C)\le\delta$.
\end{proof}
% Lean: ZM.exists_thetaFam_isPEL_smooth: T.isGoodThetaComplexClass_unitaryFullLevel (Mat, kappa, j1),
% T.nonempty_thetaObject_family_open_closure with hcrit (j2), j0 = max j1 j2, thetaFam_isPEL (F.reduced, F.qproj
% through isProper_str for LocallyOfFiniteType, LocallyOfFiniteType.isLocallyNoetherian, ArBase for the injectivity of
% em, then T.exists_isPELTriple_of_dense_fibres with F.dense), and ThetaFam.smooth_genericFibre with the tangent bound:
% tObj_pull (t_R), PELPull.isPEL_pull (the PEL structure on t_R: sabPull, grpPullback_iso, LiePull.lieMult_pull and
% ObjPull.tateMove), hρ2 from hb.2.2, finiteType from isQProjOver_iff_locallyOfFiniteType_and_exists_isAmpleLB, and
% T.finrank_tangent_PEL_family_le_symmetricDomain with ThetaFam.unram_of_mono (F.unram after composing with x, by
% thetaLevelPullback_id_comp_cancel and cancel_mono). Inputs Out.Laws, Out.Laws2, hcrit (Lemma
% cv-toroidal-theta-reduction), the theorem of the cube and ThetaUnifAV are hypotheses. Size of the whole, counted
% through the declarations reached in the imported files and beyond the laws of abelian schemes: about 93,800 lines,
% about 82,000 of them outside the proof of Lemma cv-toroidal-theta-reduction; Lemma cv-toroidal-theta-family-pel-good
% accounts for about 69,600 and -exists for about 45,100 (they share much).

\subsection{Universal families and theta envelopes}\label{subsection-cv-theta-envelopes}

\begin{lemma}[Alp\"oge]\label{cv-toroidal-universal-cubical}
In Situation~\ref{cv-function-situation}, let $f:X'\to X$ be a morphism of schemes, let $G$ and $G'$ be
$\mathcal O_E$-semi-abelian schemes over $X$ and $X'$ with line bundles $\mathcal L$ and $\mathcal L'$
(Definition~\ref{cv-model-toroidal-semiabelian}(2)), and let $\varphi:G'\to G$ be a homomorphism over $f$
that identifies $G'$ with $f^*G$ as group schemes, commutes with the actions of $\mathcal O_E$, and
satisfies $\mathcal L'\cong\varphi^*\mathcal L$. Suppose that $X'$ is reduced and that its generic points lie
in the abelian locus of $G'$. Then $\varphi$ is a morphism of $\mathcal O_E$-semi-abelian schemes over $f$.
\end{lemma}
\leanref{Development/AbelianVarietyTheory/Proofs/ModuliProofs/ExistsSabMapOfLbIso.lean}{GlobalTheory.exists_sabMap_of_lbIso}
\begin{proof}
It remains to see that $\varphi$ respects the cubical structures. Choose the isomorphism
$\mathcal L'\cong\varphi^*\mathcal L$ compatible with the trivializations along the zero sections.
Through it, the cubical structure of $\mathcal L$ gives a second cubical structure on $\mathcal L'$, and
the two differ by a unit $c$ on $G'^3=G'\times_{X'}G'\times_{X'}G'$ that is $1$ along the zero section.
The scheme $G'^3$ is smooth over $X'$ with geometrically integral fibres, and over a generic point of $X'$
its fibre is an abelian variety, whose regular functions are constant; so $c=1$ on those fibres. As $X'$
is reduced and $G'^3$ is smooth over it, $G'^3$ is reduced, and its generic points lie over generic points
of $X'$; so $c=1$.
\end{proof}
% Lean: exists_sabMap_of_lbIso: TripleMap.exists_rigidified (w'), TripleMap.pmaps_cubical (the unit on
% G'^3), smooth_geometricallyIntegral_abLocus_tensorObj and eq_one_of_appLE_section_eq_one (c = 1).

\begin{definition}\label{cv-toroidal-universal-matched}
In Situation~\ref{cv-function-situation}, let $n$ be a prime, $N'\ge1$ and $N=N'^2$, let
$e_1,\dots,e_{2g}$ be a basis of the lattice $\Lambda$ of the datum over $\mathbf Z$, and let $a:I\to\mathcal O_E$
be a map from a finite set. Let $S$ be a scheme over $R_{n,N}$, let $t=(A,L,(x_i),(\alpha_h))$ be a
theta-level object over $S$ for $o=R_{n,N}$ and $\zeta=\zeta_N$ (Definition~\ref{cv-toroidal-theta}(1)),
and let $G$ be an $\mathcal O_E$-semi-abelian scheme over $S$ with a map $\lambda:\Lambda\to G(S)$. An
isomorphism $\varphi:G\to A$ of group schemes over $S$ \emph{matches} $(G,\lambda)$ with $t$ if
$\varphi^*L$ is isomorphic to the line bundle of $G$, $\varphi$ carries the action of $a(h)$ to $\alpha_h$
for every $h\in I$, and $\varphi$ carries $\lambda(e_i)$ to $x_i$ for every $i$.
\end{definition}
\leanref{Development/AbelianVarietyTheory/Proofs/ModuliProofs/ActCommOfClosure.lean}{GlobalTheory.TripleTheta.Under}

\begin{lemma}[Alp\"oge]\label{cv-toroidal-universal-structures}
In Situation~\ref{cv-function-situation}, let $n$ be a prime, $N'\ge1$ and $N=N'^2$, let
$e_1,\dots,e_{2g}$ be a basis of the lattice $\Lambda$ of the datum over $\mathbf Z$, let $a:I\to\mathcal O_E$
be a map from a finite set, and consider theta-level objects for $o=R_{n,N}$ and $\zeta=\zeta_N$.
\begin{enumerate}
\item Let $t$ be a theta-level object over a scheme $S$ over $R_{n,N}$ with a PEL structure $G$ for
  $((e_i),a,\theta)$ (Definition~\ref{cv-toroidal-theta}(6)). Then there is $\lambda:\Lambda\to G(S)$
  such that $(G,\lambda)$ is a full-level object over $S$ (Definition~\ref{cv-toroidal-family}(1)), and the
  isomorphism of the PEL structure matches $(G,\lambda)$ with $t$
  (Definition~\ref{cv-toroidal-universal-matched}).
\item Suppose that the basis $(e_i)$ is symplectic (Definition~\ref{cv-toroidal-theta}(6)). Let $S$ be a
  scheme flat over $R_{n,N}$ and let $(G,\lambda)$ be a full-level object over $S$. Then
  $t=(G,\mathcal L,(\lambda(e_i)),(a(h)))$, with $\mathcal L$ the line bundle of $G$, is a theta-level
  object over $S$, the identity matches $(G,\lambda)$ with $t$, and $G$ with the identity is a PEL
  structure on $t$ for $((e_i),a,\theta)$.
\item Let $(G,\lambda)$ over $S$ and $(G',\lambda')$ over $S'$ be matched with theta-level objects $t$ and
  $t'$, with $\lambda$ and $\lambda'$ additive, and let $f:S'\to S$ be a morphism over $R_{n,N}$. A morphism
  $G'\to G$ of $\mathcal O_E$-semi-abelian schemes over $f$ that carries $\lambda'$ to $\lambda$ gives a
  morphism $t'\to t$ over $f$. Conversely, if $S'$ is reduced, $G'$ is proper over $S'$ and the image of
  $a$ generates $\mathcal O_E$ as a ring, then every morphism $t'\to t$ over $f$ comes in this way from a
  morphism $G'\to G$ of $\mathcal O_E$-semi-abelian schemes over $f$ that carries $\lambda'$ to $\lambda$.
\item Let $(G,\lambda)$ over $S$ be matched with $t$, let $t'$ be a theta-level object over $S'$, let $G'$
  be an $\mathcal O_E$-semi-abelian scheme over $S'$ with a map $\lambda':\Lambda\to G'(S')$, and let
  $f:S'\to S$ be a morphism with a morphism $t'\to t$ over $f$ and a morphism $G'\to G$ of
  $\mathcal O_E$-semi-abelian schemes over $f$ that carries $\lambda'$ to $\lambda$. Then $(G',\lambda')$ is
  matched with $t'$.
\end{enumerate}
\end{lemma}
\leanref{Development/AbelianVarietyTheory/Proofs/ModuliProofs/ExistsUniversalFullLevelOfPelThetaFamily.lean}{GlobalTheory.TripleTheta.exists_level}
\leanref{Development/AbelianVarietyTheory/Proofs/ModuliProofs/ExistsUniversalFullLevelOfPelThetaFamily.lean}{GlobalTheory.TripleTheta.exists_under}
\leanref{Development/AbelianVarietyTheory/Proofs/ModuliProofs/ActCommOfClosure.lean}{GlobalTheory.TripleTheta.tmap_of_sabMap}
\leanref{Development/AbelianVarietyTheory/Proofs/ModuliProofs/ActCommOfClosure.lean}{GlobalTheory.TripleTheta.sabMap_of_tmap}
\leanref{Development/AbelianVarietyTheory/Proofs/ModuliProofs/ActCommOfClosure.lean}{GlobalTheory.TripleTheta.under_of_tmap}
\begin{proof}
(1) Let $\varphi:G\to A$ be the isomorphism of the PEL structure, and let
$\lambda(\sum_ic_ie_i)=\sum_ic_i\varphi^{-1}(x_i)$. It is additive, $N\lambda=0$ because $N x_i=0$, and
$\varphi$ carries $\lambda(e_i)$ to $x_i$. Through $\varphi$, $G$ is proper over $S$, and its line bundle,
isomorphic to $\varphi^*L$, is relatively ample and defines a principal polarization. At a point $x$ of $S$ with
values in an algebraically closed field of characteristic zero, the PEL structure gives condition (a) and a
Tate system with $\zeta'_N=\beta(\zeta_N)$ and $\iota_N(e_i)=x_i(x)$ through $\varphi$, that is,
$\iota_N(e_i)=\lambda(e_i)(x)$; both sides are additive in $e_i$, so $\iota_N(b)=\lambda(b)(x)$ for all
$b\in\Lambda$, which is condition (b).

(2) As $G$ is proper over $S$, it is an abelian scheme. Since $S$ is flat over the domain $R_{n,N}$ of
characteristic zero, every point of $S$ is a specialization of a point over the generic point of
$\operatorname{Spec}R_{n,N}$, the image of a point $x$ of $S$ with values in an algebraically closed field
$k$ of characteristic zero. At such a point, condition (a) gives the Lie algebra
of the fibre the dimension $\sum_\tau(p\text{ or }q)=\frac12[E:\mathbf Q](p+q)=g$, the sum over the
embeddings $\tau:\mathcal O_E\to k$; as $G$ is smooth with connected fibres, its relative dimension is $g$.
Its line bundle is trivialized along the zero section, symmetric, relatively ample and defines a principal
polarization. The points $x_i=\lambda(e_i)$ are killed by $N$, which is invertible on $S$. At a point of
characteristic zero as above, condition (b) gives $\iota_N$ with kernel $N\Lambda$, image the $N$-torsion
and $\iota_N=\lambda(\cdot)(x)$, so every point of order dividing $N$ is $\sum_ic_ix_i$ for a unique
$c\in(\mathbf Z/N)^{2g}$, and $e_N(x_i,x_{i'})=\beta(\zeta_N)^{\langle e_i,e_{i'}\rangle}=\beta(\zeta_N)^{\omega(e_i,e_{i'})}$,
as the basis is symplectic. Both properties
hold on open and closed subsets of $S$: the first because the $N$-torsion of $G$ is finite and \'etale over
$S$ (Lemma~\ref{cv-toroidal-moduli-torsion}), the second because moreover the commutator pairing on it is
locally constant. These subsets contain a generization of every point, so they are all of $S$. So $t$ is
a theta-level object, the identity matches $(G,\lambda)$ with it, and the conditions at points of
characteristic zero for a PEL structure are conditions (a) and (b) for $(G,\lambda)$, with
$\iota_N(e_i)=\lambda(e_i)(x)=x_i(x)$.

(3) Let $\varphi$ and $\varphi'$ match $(G,\lambda)$ with $t=(A,L,(x_i),(\alpha_h))$ and $(G',\lambda')$ with
$t'=(A',L',(x'_i),(\alpha'_h))$. Given a morphism $\mu:G'\to G$ over $f$, the composite
$\varphi\circ\mu\circ\varphi'^{-1}:A'\to A$ identifies $A'$ with $f^*A$, $L'$ with the pullback of $L$, and
$x'_i$ and $\alpha'_h$ with the pullbacks of $x_i$ and $\alpha_h$; it is a morphism $t'\to t$. Conversely,
given a morphism $\chi:A'\to A$ of theta-level objects over $f$, the composite
$\mu=\varphi^{-1}\circ\chi\circ\varphi':G'\to G$ identifies $G'$ with $f^*G$, its line bundle with the
pullback of that of $G$, and commutes with the action of $a(h)$ for every $h$. The elements of
$\mathcal O_E$ whose actions commute with $\mu$ form a subring, which contains the image of $a$, so it is
$\mathcal O_E$. As $S'$ is reduced and $G'$ is proper, Lemma~\ref{cv-toroidal-universal-cubical} shows that
$\mu$ is a morphism of $\mathcal O_E$-semi-abelian schemes. It carries $\lambda'(e_i)$ to
$\lambda(e_i)$, hence $\lambda'$ to $\lambda$, both being additive.

(4) The morphisms $G'\to G$ and $t'\to t$ identify $G'$ with $f^*G$ and $A'$ with $f^*A$, so $\varphi$
induces an isomorphism $\varphi':G'\to A'$ over $S'$ compatible with them. It carries the line bundle,
the actions of the $a(h)$ and the points $\lambda'(e_i)$ to those of $t'$, since $\varphi$ does so for
$(G,\lambda)$ and $t$, and these properties can be checked after composing with the projections to $G$ and
$A$.
\end{proof}
% Lean: (1) exists_level (lv from q i = T.pt i >> phi.inv, pol_pull_of_relAmple_local, eq_of_basis);
% (2) exists_under (exists_generization, smoothOfRelativeDimension_of_generization with
% ObjDim.dim_eq_of_detCond, torsion_basis_of_generization via torsion_basis_of_dense, weilIs_of_generization
% via exists_clopen_weilIs_iff, which uses Out.Laws, and existsUnique_coeff); (3) tmap_of_sabMap and
% sabMap_of_tmap (act_comm_of_closure, exists_sabMap_of_lbIso, eq_of_basis); (4) under_of_tmap
% (grpPullback_iso).

\begin{lemma}[Alp\"oge]\label{cv-toroidal-universal-complex}
In Situation~\ref{cv-function-situation}, suppose $p+q\ge1$, let $N'\ge1$, let $e_1,\dots,e_{2g}$ be a
symplectic basis of the lattice $\Lambda$ of the datum, and let $a:I\to\mathcal O_E$ be a map from a finite
set whose image generates $\mathcal O_E$ as a ring and contains $w_0=(1+\sqrt{-d})/2$ and $\bar w_0$. Let
$o$ be a ring with a unit $\zeta$ as in Definition~\ref{cv-toroidal-theta}, let $\mathrm{em}:o\to\mathbf C$
be an injective homomorphism such that $o$ is a Dedekind domain of finite type over $\mathbf Z$, and let
$\rho\in o$ with $\mathrm{em}(\rho)=\theta$. Let $t$ be a theta-level object over a scheme $S$ over $o$
with a PEL structure for $((e_i),a,\rho)$ (Definition~\ref{cv-toroidal-theta}(6)). The hypotheses that
$p+q\ge1$, that the basis is symplectic, that the image of $a$ contains $w_0$ and $\bar w_0$, and that
$\mathrm{em}(\rho)=\theta$ are not used in the proof. Then for every field $K$ of characteristic zero and every
point $y:\operatorname{Spec}K\to S$, $t$ is a $\mathcal C$-object at $y$ (Definition~\ref{cv-toroidal-theta}(4)),
for the PEL class $\mathcal C$ of $((e_i),a,\rho,\mathrm{em})$.
\end{lemma}
\leanref{Development/AbelianVarietyTheory/Proofs/ModuliProofs/PelObjectFibreFromComplexCharZero.lean}{GlobalTheory.PELFrom.fromC_of_isPEL}
\begin{proof}
Let $\beta:o\to K$ be the homomorphism of $y$. By Lemma~\ref{cv-toroidal-reduction-base-change}(1) there is a
theta-level object $y^*t$ over $K$ with a morphism $y^*t\to t$ over $y$, and by
Lemma~\ref{cv-toroidal-theta-family-pel-fibre-pull} it has a PEL structure for $((e_i),a,\rho)$.

Step 1. There are a subfield $K_0\subset K$, generated over $\beta(o)$ by finitely many elements, a
theta-level object $t_0$ over $K_0$, regarded over $o$ through $\beta$ with its values in $K_0$, written
$\beta_0:o\to K_0$, and a morphism $y^*t\to t_0$ over $\operatorname{Spec}K\to\operatorname{Spec}K_0$. This is
the descent of finitely presented schemes, morphisms, properties of morphisms and line bundles along the filtered
union of such subfields \cite[\S8]{Gro66}. As $o$ is of finite type over $\mathbf Z$, $K_0$ is countable.

Step 2. Lemma~\ref{cv-toroidal-theta-family-pel-fibre-descent}, applied to the inclusion $K_0\to K$, $\beta_0$,
$t_0$ and $y^*t$ with its morphism to $t_0$ and its PEL structure, gives a PEL structure on $t_0$ for
$((e_i),a,\rho)$; it uses that the image of $a$ generates $\mathcal O_E$.

Step 3. Since $K_0$ is countable of characteristic zero and $o$ is a Dedekind domain of finite type over
$\mathbf Z$ with $\mathrm{em}$ injective, there is an embedding $j:K_0\to\mathbf C$ with
$j\circ\beta_0=\mathrm{em}$: extend $\mathrm{em}$ to the fraction field of $o$, map a transcendence basis of
$K_0$ over it to elements of $\mathbf C$ algebraically independent over the image, and extend to the
algebraic part.

By Lemma~\ref{cv-toroidal-reduction-base-change}(1) there is a theta-level object $t_{\mathbf C}$ over
$\mathbf C$, regarded over $o$ through $j\circ\beta_0=\mathrm{em}$, with a morphism $t_{\mathbf C}\to t_0$ over
$\operatorname{Spec}j$, and by Lemma~\ref{cv-toroidal-theta-family-pel-fibre-pull} it has a PEL structure. So
$t_{\mathbf C}$ belongs to $\mathcal C$, and $K_0$, $t_0$, $y^*t$ and $t_{\mathbf C}$ show that $t$ is a
$\mathcal C$-object at $y$.
\end{proof}
% Lean: fromC_of_isPEL: T.exists_thetaLevelAbelianScheme_pullback (Lemma cv-toroidal-reduction-base-change(1)) and
% PELPull.isPEL_pull (Lemma cv-toroidal-theta-family-pel-fibre-pull) for the pullbacks to K and to C,
% T.thetaLevel_object_descends_fg_subfield = ThetaDescent.descends_to_fg_subfield (Step 1),
% T.pelStructure_descends_field_extension (Lemma cv-toroidal-theta-family-pel-fibre-descent, Step 2; uses Out.Laws),
% T.exists_ringHom_complex_comp_eq_of_countable_field (Step 3). The hypotheses p + q >= 1, IsSympBasis,
% a w = w0, a w' = bar w0, em rho = rt0 and Out.Laws2 are introduced as _ and not used.

\begin{lemma}[Alp\"oge]\label{cv-toroidal-theta-universal}
In Situation~\ref{cv-function-situation}, suppose $p+q\ge1$, let $n$ be an odd prime, $j\ge1$, $N'=n^j$
and $N=N'^2$, and let $(\mathcal M,t,G)$ be a PEL theta family of level $N'$
(Definition~\ref{cv-toroidal-theta}(7)). Then there is a map $\lambda:\Lambda\to G(\mathcal M)$ such
that $(\mathcal M,G,\lambda)$ is a full-level family over $R_{n,N}$
(Definition~\ref{cv-toroidal-family}).
\end{lemma}
\leanref{Development/AbelianVarietyTheory/Proofs/ModuliProofs/ExistsUniversalFullLevelOfPelThetaFamilyStatement.lean}{GlobalTheory.T.exists_universal_fullLevel_of_pelThetaFamily}
\leanref{Development/AbelianVarietyTheory/Proofs/ModuliProofs/ExistsUniversalFullLevelOfPelThetaFamily.lean}{GlobalTheory.exists_famCore_of_thetaFam}
\begin{proof}
Let $(e_i)$ and $a:I\to\mathcal O_E$ be as in Definition~\ref{cv-toroidal-theta}(7) for $(\mathcal M,t,G)$,
let $\varphi:G\to A$ be the isomorphism of the PEL structure, and let $s_{\mathcal M}$ be the structure
morphism of $\mathcal M$. By Lemma~\ref{cv-toroidal-universal-structures}(1) there is
$\lambda:\Lambda\to G(\mathcal M)$ such that $(G,\lambda)$ is a full-level object over $\mathcal M$, matched
with $t$ by $\varphi$. The scheme $\mathcal M$ is reduced, flat and quasi-projective over $R_{n,N}$, and its
fibre over the fraction field of $R_{n,N}$ is smooth. The ring $R_{n,N}$ is a Dedekind domain of finite
type over $\mathbf Z$, contained in $\mathbf C$.

Property (a) of Definition~\ref{cv-toroidal-family}(2). Let $S$ be a reduced, locally noetherian scheme,
flat over $R_{n,N}$, and $(G',\lambda')$ a full-level object over $S$. By
Lemma~\ref{cv-toroidal-universal-structures}(2) there is a theta-level object $t'$ over $S$ matched with
$(G',\lambda')$ by the identity, with the PEL structure $G'$. By Lemma~\ref{cv-toroidal-universal-complex},
$t'$ is a $\mathcal C$-object, for the PEL class $\mathcal C$, at every point of $S$ with values in a field of
characteristic zero. As $S$ is flat over $R_{n,N}$, every point of $S$ is a specialization of such a point,
so the points at which $t'$ is a $\mathcal C$-object are dense. Property (d) of
Definition~\ref{cv-toroidal-theta}(5) gives a unique $f:S\to\mathcal M$ over $R_{n,N}$ with a morphism
$t'\to t$ over $f$. As $S$ is reduced, $G'$ is proper and the image of $a$ generates $\mathcal O_E$,
Lemma~\ref{cv-toroidal-universal-structures}(3) turns it into a morphism $G'\to G$ over $f$ that carries
$\lambda'$ to $\lambda$. If $f':S\to\mathcal M$ over $R_{n,N}$ has a morphism $G'\to G$ over $f'$ that carries
$\lambda'$ to $\lambda$, then part (3) of that lemma gives a morphism $t'\to t$ over $f'$, and $f'=f$.

Property (b). Let $x:\operatorname{Spec}k[\epsilon]\to\mathcal M$ and $x_0$ be as there, with the same
composite to $\operatorname{Spec}R_{n,N}$, and let $G'$ over $k[\epsilon]$ with $\lambda'$ have morphisms
$\mu$ and $\mu_0$ to $G$ over $x$ and $x_0$ that carry $\lambda'$ to $\lambda$. The pullback $t'=x^*t$ is a
theta-level object over $k[\epsilon]$ with a morphism $t'\to t$ over $x$. By
Lemma~\ref{cv-toroidal-universal-structures}(4), applied to $\mu$, $(G',\lambda')$ is matched with $t'$, and
by part (3) of that lemma, $\mu_0$ gives a morphism $t'\to t$ over $x_0$. Property (e) of
Definition~\ref{cv-toroidal-theta}(5) gives $x=x_0$.
\end{proof}
% Lean: exists_famCore_of_thetaFam: arBase_RZ, exists_level, F.reduced, F.flat, F.qproj and the smooth
% generic fibre; rep from exists_under, F.rep with exists_generization and PELFrom.fromC_of_isPEL (density),
% sabMap_of_tmap and tmap_of_sabMap; unram from tObj_pull, under_of_tmap, tmap_of_sabMap and F.unram.

\begin{definition}\label{cv-toroidal-envelope}
In Situation~\ref{cv-function-situation}:
\begin{enumerate}
\item Let $G$ be an $\mathcal O_E$-semi-abelian scheme over a scheme $Z$, with line bundle $\mathcal L$,
  and let $N\ge0$. An \emph{ample theta presentation} of level $N$ of $G$ consists of an ample line
  bundle $\mathcal A$ on $Z$, an open cover $(U_i)$ of $Z$ on which $\mathcal A$ is trivialized,
  an integer $n_0$, two integers $r$ and $s$ (the \emph{integers} of the presentation), and for each
  $i$ a section $\vartheta_i$ of $\mathcal L$ over $G_{U_i}$
  and sections $b_{i,1},\dots,b_{i,n_0}$ of $\mathcal L^N$ over $G_{U_i}$, such that:
  \begin{enumerate}
  \item on $U_i\cap U_{i'}$ there are a unit $c$ and an invertible matrix $M$ of size $n_0$ with
    $\vartheta_i=c\,\vartheta_{i'}$ and $b_{i,k}=\sum_{k'}M_{kk'}b_{i',k'}$, and the transition function
    of $\mathcal A$ on $U_i\cap U_{i'}$ is $c^r\det(M)^s$;
  \item for every complete discrete valuation ring $V$ of characteristic zero and every morphism
    $\chi:\operatorname{Spec}V\to U_i$ that maps the generic point into the abelian locus of $G$, the
    pullbacks of $\vartheta_i$ and of $b_{i,1},\dots,b_{i,n_0}$ are bases over $V$ of the sections of
    $\mathcal L$ and of $\mathcal L^N$ over $G\times_Z\operatorname{Spec}V$.
  \end{enumerate}
\item Let $O$ be a ring, $\mathcal M$ a scheme over $O$ and $G_{\mathcal M}$ an
  $\mathcal O_E$-semi-abelian scheme over $\mathcal M$. A \emph{theta envelope} of level $N$ of
  $(\mathcal M,G_{\mathcal M})$ over $O$ is a tuple $(Z,G,W,j)$ of a normal scheme $Z$, projective and
  flat over $O$; a theta-split $\mathcal O_E$-semi-abelian scheme $G$ over $Z$
  (Definition~\ref{cv-model-toroidal-semiabelian}(3)), quasi-projective over $O$, with an ample theta
  presentation of level $N$; a dense open subscheme $W$ of $\mathcal M$; and an open immersion
  $j:W\to Z$ over $O$ with dense image; such that some $\mathcal O_E$-semi-abelian scheme $G_W$ over $W$
  has morphisms to $G_{\mathcal M}$ over the inclusion of $W$ and to $G$ over $j$.
\end{enumerate}
\end{definition}
\leanref{Development/AbelianVarietyTheory/Definitions/Degenerations/AdmissiblePadicIntMaximalIdeal.lean}{GlobalTheory.LN.AmpFrames}
\leanref{Development/AbelianVarietyTheory/Definitions/ModuliSpaces/TopAmpExists.lean}{GlobalTheory.Models.ZOn}

\begin{lemma}\label{cv-toroidal-theta-envelope-chart}
Let $o$ be a noetherian domain, let $X$ be a reduced scheme, locally of finite type and flat over $o$, let $\xi$
be a generic point of $X$, and let $x:\operatorname{Spec}R_0\to X$ be an open immersion whose image contains
$\xi$. Then there are a field $\Omega$ with an injective homomorphism $o\to\Omega$, a finitely generated
$o$-subalgebra $R$ of $\Omega$ such that every element of $\Omega$ is a quotient $a/c$ with $a,c\in R$, and a
homomorphism $\varphi:R_0\to R$ such that $x\circ\operatorname{Spec}\varphi$ is an open immersion over $o$ whose
image contains $\xi$.
\end{lemma}
\leanref{Development/AbelianVarietyTheory/Proofs/ModuliProofs/NonemptyZOnOfComponents.lean}{GlobalTheory.FamilyModel.exists_integral_chart}
\leanref{Development/AbelianVarietyTheory/Proofs/ModuliProofs/NonemptyZOnOfComponents.lean}{GlobalTheory.FamilyModel.exists_isDomain_away}
\begin{proof}
The composite of $x$ with the structure morphism of $X$ is $\operatorname{Spec}\beta$ for a homomorphism
$\beta:o\to R_0$, which is of finite type as $X$ is locally of finite type over $o$. So $R_0$ is noetherian, and it
is reduced, as $X$ is. Let $\mathfrak p$ be the prime of $R_0$ with $x(\mathfrak p)=\xi$. It is a minimal prime: if
$\mathfrak q\subset\mathfrak p$ is prime, then $x(\mathfrak q)$ specializes to $\xi$, so $x(\mathfrak q)=\xi$ as
$\xi$ is a generic point, and $\mathfrak q=\mathfrak p$ as $x$ is injective.

There is $f\in R_0\setminus\mathfrak p$ such that $(R_0)_f$ is a domain. Indeed, $R_0$ has finitely many minimal
primes, and $\mathfrak p$ contains none of the others, so it does not contain their intersection, which therefore
has an element $f\notin\mathfrak p$. For $a\in\mathfrak p$, $fa$ lies in every minimal prime, hence is nilpotent,
hence $0$, as $R_0$ is reduced. So the kernel of $R_0\to(R_0)_f$ is $\mathfrak p$: it contains $\mathfrak p$, and
$f^ka=0$ implies $a\in\mathfrak p$ as $f\notin\mathfrak p$. Hence $(R_0)_f$ is nonzero, and it has no zero
divisors: if $(a/f^k)(b/f^l)=0$, then $ab\in\mathfrak p$, so $a$ or $b$ lies in $\mathfrak p$ and maps to $0$.

Let $\Omega$ be the fraction field of $(R_0)_f$, let $R\subset\Omega$ be the image of $(R_0)_f$, which is
isomorphic to it, and let $\varphi:R_0\to R$ be the induced homomorphism; $R$ is an $o$-algebra through
$\varphi\circ\beta$. Then $\operatorname{Spec}\varphi$ is the inclusion of the open subset of the primes not
containing $f$, which contains $\mathfrak p$; so $x\circ\operatorname{Spec}\varphi$ is an open immersion over $o$
whose image contains $\xi$. The $o$-algebra $R\cong R_0[1/f]$ is finitely generated, and every element of $\Omega$
is a quotient of elements of $R$. Finally, $\operatorname{Spec}R\to X\to\operatorname{Spec}o$ is flat, so $R$ is a
flat, hence torsion-free, $o$-module; as $1\ne0$ in $R$, $o\to R$ is injective, and so is $o\to\Omega$.
\end{proof}
% Lean: exists_integral_chart: Spec.map_surjective (beta), Algebra.FiniteType.isNoetherianRing,
% isReduced_of_isOpenImmersion, eq_of_specializes_of_mem_genericPoints (minimality), exists_isDomain_away (f, from
% minimalPrimes.finite_of_isNoetherianRing, Ideal.IsPrime.inf_le' and Ideal.sInf_minimalPrimes), Omega =
% FractionRing (Localization.Away f), R the range of the algebra map, and injective_of_flat. About 100 lines,
% counted through the declarations reached in the imported files.

\begin{lemma}\label{cv-toroidal-theta-envelope-homogeneous}
In the setting of Definition~\ref{cv-toroidal-theta}(3), let $N'$ be odd, and suppose that $\mathrm{em}$ is
injective; this hypothesis is not used in the proof. If $P$ lies in the theta ideal of $\mathcal C$, then so does
every homogeneous component of $P$.
\end{lemma}
\leanref{Development/AbelianVarietyTheory/Proofs/ModuliProofs/ThetaPointsVanishingIdealRadicalHomogeneous.lean}{GlobalTheory.ZM.ThIdeal.thIdeal_radical_homogeneous_torsionFree}
\leanref{Development/AbelianVarietyTheory/Proofs/ModuliProofs/ThetaPointsVanishingIdealRadicalHomogeneous.lean}{GlobalTheory.ZM.ThIdeal.component_eq_zero}
\begin{proof}
Write $P=\sum_eP_e$ with $P_e$ homogeneous of degree $e$. Let $t$ be a member of $\mathcal C$ and $v$ a family of
theta constants of $t$ at the identity. For every $c\in\mathbf C^\times$, $cv=(cv_a)_a$ is also one, by
Lemma~\ref{cv-toroidal-reduction-base-change}(5), so $\mathrm{em}(P)(cv)=0$. The polynomial
$\mathrm{em}(P)(Tv)=\sum_e\mathrm{em}(P_e)(v)\,T^e$ in one variable $T$ over $\mathbf C$ thus has infinitely many
roots, so it is $0$, and $\mathrm{em}(P_e)(v)=0$ for every $e$. As $t$ and $v$ were arbitrary, each $P_e$ lies in
the theta ideal.
\end{proof}
% Lean: the second part of thIdeal_radical_homogeneous_torsionFree (mem_thIdeal, component_eq_zero with
% coeff_scale and Polynomial.eq_zero_of_infinite_isRoot, hasTh_unit_mul); the other two parts (the ideal is radical,
% and C c * F in it with c nonzero implies F in it, which uses that em is injective) are not needed here. About 160
% lines with what it reaches.

\begin{lemma}[Alp\"oge]\label{cv-toroidal-theta-envelope-ratios}
In the setting of Definition~\ref{cv-toroidal-theta}(3), let $N'$ be odd, suppose that $\mathrm{em}$ is injective,
and let $(\mathcal M,t)$ be a theta family for $\mathcal C$ and $n$ (Definition~\ref{cv-toroidal-theta}(5)). Let
$\Omega$ be a field over $o$, let $R$ be an $o$-subalgebra of $\Omega$ such that every element of $\Omega$ is a
quotient of elements of $R$, let $i:\operatorname{Spec}R\to\mathcal M$ be an open immersion over $o$, and let $v$ be
a family of theta constants of $t$ at $i$. Then $\Omega$ is generated as a field by the image of $o$ and the
quotients $v_c/v_{c'}$ with $v_{c'}\ne0$.
\end{lemma}
\leanref{Development/AbelianVarietyTheory/Proofs/ModuliProofs/NonemptyZOnOfComponents.lean}{GlobalTheory.FamilyModel.closure_ratios_eq_top}
\leanref{Development/AbelianVarietyTheory/Proofs/ModuliProofs/NonemptyZOnOfComponents.lean}{GlobalTheory.FamilyModel.eval₂_mul_eq_zero}
\begin{proof}
Let $K\subset\Omega$ be the subfield generated by the image of $o$ and these quotients, let $\beta:o\to K$ be the
induced homomorphism, and let $y_\Omega:\operatorname{Spec}\Omega\to\mathcal M$ be the composite of $i$ with the
morphism given by the inclusion $R\subset\Omega$; it lies over $o\to\Omega$. By
Lemma~\ref{cv-toroidal-reduction-base-change}(4), $v$, regarded in $\Omega$, is a family of theta constants of $t$
at $y_\Omega$. Property (b) of Definition~\ref{cv-toroidal-theta}(5), for the field $\Omega$, gives a family $v'$ of
theta constants of $t$ at $y_\Omega$ that is admissible for $(\mathcal C,n)$ over $o\to\Omega$. Its coset
condition at the prime $0$ of $\Omega$ gives $a_0$ with $v'_{a_0}\ne0$; so the $v'_c$ generate the unit ideal of
$\Omega$, and Lemma~\ref{cv-toroidal-reduction-unit} gives $u\in\Omega^\times$ with $v_c=uv'_c$ for every $c$. In
particular $v_{a_0}\ne0$, and $w_c=v_c/v_{a_0}$ is an element of $K$ with $w_c=v'^{-1}_{a_0}v'_c$ in $\Omega$.

The family $v''=(v'^{-1}_{a_0}v'_c)_c$ is admissible for $(\mathcal C,n)$ over $o\to\Omega$. Indeed, let $P$ be in
the theta ideal, with homogeneous components $P_e$. They lie in the theta ideal
(Lemma~\ref{cv-toroidal-theta-envelope-homogeneous}), so they vanish at $v'$, and
$P(v'')=\sum_ev'^{-e}_{a_0}P_e(v')=0$. The coset condition need only be checked at the prime $0$ of $\Omega$, where
it holds for $v''$ because it holds for $v'$. As $K\to\Omega$ is injective and carries $w$ to $v''$, the family
$w$ is admissible for $(\mathcal C,n)$ over $\beta$: the polynomial conditions hold after applying $K\to\Omega$,
and so does the coset condition at the prime $0$ of $K$.

Property (c) gives a point $y:\operatorname{Spec}K\to\mathcal M$ over $\operatorname{Spec}\beta$ at which $w$ is a
family of theta constants. By Lemma~\ref{cv-toroidal-reduction-base-change}(4), the composite of $y$ with the
morphism given by $K\subset\Omega$ is a point over $o\to\Omega$ at which $v''$ is a family of theta constants; by
Lemma~\ref{cv-toroidal-reduction-base-change}(5), applied to $v'$ and the unit $v'^{-1}_{a_0}$, so is $y_\Omega$. By
the uniqueness in property (c) for $v''$, these two points are equal. So the image under $y$ of the point of
$\operatorname{Spec}K$ is the image under $y_\Omega$ of the point of $\operatorname{Spec}\Omega$, which lies in the
image of $i$. As $i$ is an open immersion, $y=i\circ\operatorname{Spec}\psi$ for a homomorphism $\psi:R\to K$. Then
$i$ composed with the morphism given by $R\xrightarrow{\psi}K\subset\Omega$ is $y_\Omega$, and as $i$ is a
monomorphism, $R\xrightarrow{\psi}K\subset\Omega$ is the inclusion of $R$. Hence $R\subset K$, and $K=\Omega$, since
every element of $\Omega$ is a quotient of elements of $R$.
\end{proof}
% Lean: closure_ratios_eq_top: thIdeal_radical_homogeneous_torsionFree (homogeneity), hasThetaConstants_ring_baseChange
% (hvOmega), F.loc at Omega, ThetaUniq.exists_unit_local (the local case of Lemma cv-toroidal-reduction-unit, which the
% Lean applies directly, writing v' = u v), w a = v a / v a0, eval₂_mul_eq_zero (the polynomial conditions for
% (v' a0)^{-1} v'), InMM over Omega and over K, F.pts K beta w (the point y), the uniqueness in F.pts Omega with
% hasTh_unit_mul (hxy), IsOpenImmersion.lift, Spec.map_injective and K.div_mem. About 940 lines, counted through the
% declarations reached in the imported files beyond the laws of abelian schemes, about 230 of them outside Lemmas
% cv-toroidal-theta-envelope-component and -glue.

\begin{lemma}\label{cv-toroidal-theta-envelope-pull}
In Situation~\ref{cv-function-situation}, let $t=(A,L,(x_i),(\alpha_h))$ be a theta-level object over a scheme $S$
over $o$ (Definition~\ref{cv-toroidal-theta}(1)), let $G$ be an $\mathcal O_E$-semi-abelian scheme over $S$
(Definition~\ref{cv-model-toroidal-semiabelian}(2)) with an isomorphism of group schemes $\varphi:G\to A$ such that
the line bundle of $G$ is isomorphic to $\varphi^*L$, and let $f:S'\to S$ be a morphism, with $S'$ regarded over $o$
through $f$. Then there are a theta-level object $u=(A',L',(x'_i),(\alpha'_h))$ over $S'$ with a morphism $u\to t$
over $f$, an $\mathcal O_E$-semi-abelian scheme $G'$ over $S'$ with a morphism $G'\to G$ over $f$, and an
isomorphism of group schemes $e:G'\to A'$ such that the line bundle of $G'$ is isomorphic to $e^*L'$.
\end{lemma}
\leanref{Development/AbelianVarietyTheory/Proofs/ModuliProofs/NonemptyZOnOfComponents.lean}{GlobalTheory.FamilyModel.exists_pull}
\begin{proof}
Lemma~\ref{cv-toroidal-reduction-base-change}(1) gives $u$ with a morphism $\chi:A'\to A$ over $f$, which identifies
$A'$ with $f^*A$ as group schemes and $L'$ with $\chi^*L$. Let $G'=f^*G$, with the pulled-back action of
$\mathcal O_E$ and line bundle, and $\mu:G'\to G$ the projection, as in the proof of
Lemma~\ref{cv-toroidal-abelian-pullback}(1); it is a morphism of $\mathcal O_E$-semi-abelian schemes over $f$. Both
$\varphi\circ\mu$ and $\chi$ identify their sources with $f^*A$ as group schemes, so there is a unique isomorphism
of group schemes $e:G'\to A'$ with $\chi\circ e=\varphi\circ\mu$. The line bundle of $G'$ is isomorphic to the
pullback along $\mu$ of that of $G$, hence to $\mu^*\varphi^*L=e^*\chi^*L$, hence to $e^*L'$.
\end{proof}
% Lean: exists_pull: tObj_pull (u), sabPull Out.CubePf.cubePull (G' and m), grpPullback_iso (e), m.lin and
% GenSpread.lbIso_pull_of_lbIso (the line bundles). About 600 lines with what it reaches beyond the laws of abelian
% schemes, about 130 of them outside Lemmas cv-toroidal-theta-envelope-component and -glue.

\begin{lemma}[Alp\"oge]\label{cv-toroidal-theta-envelope-component}
In Situation~\ref{cv-function-situation}, let $g\ge1$ and let $n$ be an odd prime. There are integers $c\ge0$ and
$j_0$ such that for every $j\ge j_0$, with $N'=n^j$, the following holds. Let $o$ be a Dedekind domain of finite
type over $\mathbf Z$ with a unit $\zeta$ as in Definition~\ref{cv-toroidal-theta}, let $\Omega$ be a field of
characteristic zero with an injective homomorphism $o\to\Omega$, and let $R$ be a finitely generated $o$-subalgebra
of $\Omega$ such that every element of $\Omega$ is a quotient of elements of $R$. Let $u=(A,L,(x_i),(\alpha_h))$ be
a theta-level object over $\operatorname{Spec}R$, let $G_0$ be an $\mathcal O_E$-semi-abelian scheme over
$\operatorname{Spec}R$ with an isomorphism of group schemes $e:G_0\to A$ such that the line bundle of $G_0$ is
isomorphic to $e^*L$, and let $v$ be a family of theta constants of $u$ at the identity such that $\Omega$ is
finite over the subfield generated by the image of $o$ and the quotients $v_c/v_{c'}$ with $v_{c'}\ne0$. Then
$(\operatorname{Spec}R,G_0)$ has a theta envelope of level $n^{2c}$ over $o$
(Definition~\ref{cv-toroidal-envelope}(2)) whose scheme $Z$ is integral and whose ample theta presentation has
integers $r_0$ and $1$, for some $r_0$, such that for every integer $r\ge r_0$ its $\mathcal O_E$-semi-abelian scheme
has an ample theta presentation of level $n^{2c}$ with integers $r$ and $1$.
\end{lemma}
\leanref{Development/AbelianVarietyTheory/Proofs/ModuliProofs/ExistsModelTripleAmpleFramesComponentStatement.lean}{GlobalTheory.T.exists_model_triple_ampleFrames_component}
\leanref{Development/AbelianVarietyTheory/Proofs/ModuliProofs/ExistsModelTripleAmpleFramesComponent.lean}{GlobalTheory.exists_zon_ampFrames_component}
\begin{proof}
\openstep{The construction of the theta envelope over such a chart, as an open part of the closure of the abelian
scheme in a projective space of theta functions of a large level, with the semistable reduction theorem and
Mumford's construction (Theorems~\ref{cv-semistable-reduction} and~\ref{cv-toroidal-mumford}) controlling the theta
constants near the boundary, is not yet written out.}
% Lean: exists_zon_ampFrames_component: exists_big_thetaSys_with_action (k, j0 and the theta system S with its real
% point rho), T.thetaSystem_hasTests_of_real_point (from SemistableReductionSt and ThmUSt, then
% thetaConstant_bounded_by_shortestVector, exists_norm_torsionSum_shortestVector_eq and their companions),
% ZH.exists_minor_model (the normal model Z of the maximal minors), the closure statement ZH.ZHGeoSt (the closure Y;
% a hypothesis here, proved by T.zarhinGeometric_of_semistableReduction = zhGeo_of_semistableReduction),
% isOpen_setOf_TT_eq_bot (the abelian locus is open), exists_ampFrames_exponents (r0) and
% T.exists_birationalModel_semiAbelian_ampleFrames. Stating ZHGeoSt (a ThetaSys with a realization Real, a RatModel
% with a lattice Lat satisfying Over0 and IsEE, conclusion ZHOut) as its own lemma is the next division of this
% step. Size, counted through the declarations reached in the imported files beyond the laws of abelian schemes:
% about 62,800 lines, about 60,600 of them not shared with the other inputs of Lemma cv-toroidal-theta-envelope, and
% about 37,500 beyond the proof of Theorem cv-toroidal-mumford (thmU_and_uTors), which it reaches. The proof of
% ZHGeoSt adds about 92,600 lines beyond these.
\end{proof}

\begin{lemma}[Alp\"oge]\label{cv-toroidal-theta-envelope-glue-dense}
In Situation~\ref{cv-function-situation}, let $O$ be a ring, let $\mathcal M$ be a scheme over $O$, and let
$G_{\mathcal M}$ be an $\mathcal O_E$-semi-abelian scheme over $\mathcal M$. Let $i:V\to\mathcal M$ be an open immersion
with dense image, with $V$ regarded over $O$ through $i$, and let $G_V$ be an $\mathcal O_E$-semi-abelian scheme over $V$
with a morphism $G_V\to G_{\mathcal M}$ over $i$. Let $m\ge0$, $n_0\ge0$, $r$ and $s$ be integers.
\begin{enumerate}
\item If $(V,G_V)$ has a theta envelope of level $m^2$ over $O$ (Definition~\ref{cv-toroidal-envelope}(2)) whose ample
  theta presentation has $n_0$ sections and integers $r$ and $s$, then so has $(\mathcal M,G_{\mathcal M})$.
\item If $(\mathcal M,G_{\mathcal M})$ has a theta envelope of level $m^2$ over $O$ whose ample theta presentation has
  $n_0$ sections and integers $r$ and $s$, then so has $(V,G_V)$.
\end{enumerate}
\end{lemma}
\leanref{Development/AbelianVarietyTheory/Proofs/ModuliProofs/GenericPointsOpenIntegralCover.lean}{GlobalTheory.ModelSum.HasModel.grow}
\leanref{Development/AbelianVarietyTheory/Proofs/ModuliProofs/GenericPointsOpenIntegralCover.lean}{GlobalTheory.ModelSum.HasModel.shrink}
\begin{proof}
A theta envelope may be given with a dense open immersion $a:Y\to\mathcal M$ in place of a dense open subscheme $W$:
if $j:Y\to Z$ is an open immersion over $O$ with dense image and some $\mathcal O_E$-semi-abelian scheme $G_Y$ over $Y$
has morphisms to $G_{\mathcal M}$ over $a$ and to $G$ over $j$, then $W=a(Y)$, $j\circ a^{-1}$ and $G_Y$, regarded
over $W$ through $a$, satisfy Definition~\ref{cv-toroidal-envelope}(2).

(1) Let $(Z,G,W,j)$ be the theta envelope of $(V,G_V)$, with $G_W$ and its morphisms to $G_V$ over the inclusion and to
$G$ over $j$. The composite of $W\subset V$ with $i$ is an open immersion with dense image, as $W$ is dense in $V$ and
$i(V)$ is dense in $\mathcal M$, and the composite $G_W\to G_V\to G_{\mathcal M}$ is a morphism over it. With $Z$, $G$,
its presentation and $j$, this is a theta envelope of $(\mathcal M,G_{\mathcal M})$ in the form above.

(2) Let $(Z,G,W,j)$ be the theta envelope of $(\mathcal M,G_{\mathcal M})$, with $G_W$ and its morphisms to
$G_{\mathcal M}$ and to $G$. Let $Y=i^{-1}(W)$, an open subscheme of $V$. It is dense in $V$: a nonempty open subset of
$V$ has open image in $\mathcal M$, which meets the dense open set $W$. Its image $i(Y)=i(V)\cap W$ is dense in $W$, as
$i(V)$ is dense and open in $\mathcal M$; so the composite $Y\to W\to Z$ of $i$ and $j$ is an open immersion over $O$
with dense image. Let $G_Y$ be the pullback of $G_W$ to $Y$. Its morphism to $G_{\mathcal M}$, over $Y\to V\to\mathcal M$,
factors through a morphism $G_Y\to G_V$ over the inclusion $Y\subset V$, since $G_V\to G_{\mathcal M}$ identifies
$G_V$ with the pullback of $G_{\mathcal M}$; and $G_Y\to G_W\to G$ is a morphism over $Y\to Z$. With $Z$, $G$ and its
presentation, this is a theta envelope of $(V,G_V)$ in the form above.
\end{proof}
% Lean: HasModel.grow and HasModel.shrink, through hasModel_of_maps (the form above, with Y and the isomorphism
% a.isoOpensRange), map_comp (ObjMaps.exists_sab_comp) and map_cancel (ObjMaps.exists_sab_cancel), exists_pull
% (sabPull) and isDominant_pullback_fst_of_isOpenMap. About 580 and 450 lines, counted through the declarations reached
% in the imported files beyond the laws of abelian schemes.

\begin{lemma}[Alp\"oge]\label{cv-toroidal-theta-envelope-glue-rank}
In Situation~\ref{cv-function-situation}, let $\mathcal M$ be a scheme and $G_{\mathcal M}$ an $\mathcal O_E$-semi-abelian
scheme over $\mathcal M$, proper and smooth of relative dimension $g$ over $\mathcal M$, whose line bundle is relatively
ample and defines a principal polarization. Let $O$ be a ring, let $V$ be a scheme over $O$ with a morphism
$i:V\to\mathcal M$, and let $G_V$ be an $\mathcal O_E$-semi-abelian scheme over $V$ with a morphism $G_V\to G_{\mathcal M}$
over $i$. Let $m\ge0$, let $(Z,G,W,j)$ be a theta envelope of level $m^2$ of $(V,G_V)$ over $O$
(Definition~\ref{cv-toroidal-envelope}(2)) with $Z$ integral, and suppose that $G$ has an ample theta presentation of
level $m^2$ with integers $r$ and $s$. Then $(V,G_V)$ has a theta envelope of level $m^2$ over $O$ whose ample theta
presentation has $\max(1,m^{2g})$ sections and integers $r$ and $s$.
\end{lemma}
\leanref{Development/AbelianVarietyTheory/Proofs/ModuliProofs/GenericPointsOpenIntegralCover.lean}{GlobalTheory.ModelSum.hasModel_rank}
\leanref{Development/AbelianVarietyTheory/Proofs/ModuliProofs/GenericPointsOpenIntegralCover.lean}{GlobalTheory.ModelSum.rank_eq}
\begin{proof}
Let the given presentation have line bundle $\mathcal A$, charts $U_a$, sections $\vartheta_a$ and $b_{a,1},\dots,
b_{a,n_0}$, and integers $r$ and $s$, and let $L$ be the line bundle of $G$. Put $n_1=\max(1,m^{2g})$. It suffices to
give a presentation of $G$ with the same $\mathcal A$, charts, $\vartheta_a$, $r$ and $s$ and with $n_1$ sections: with
it in place of the presentation of the theta envelope, $(Z,G,W,j)$ is the required envelope.

Suppose first that no morphism $\operatorname{Spec}R\to Z$ with $R$ a complete discrete valuation ring of characteristic
zero exists. Then condition (b) of Definition~\ref{cv-toroidal-envelope}(1) holds for any sections. Take the sections
$b'_{a,k}=0$ for $k=1,\dots,n_1$. On $U_a\cap U_{a'}$, let $c$ and $M$ be as in condition (a), and let $M'$ be the
diagonal matrix of size $n_1\ge1$ with entries $\det M,1,\dots,1$. Then $b'_{a,k}=\sum_{k'}M'_{kk'}b'_{a',k'}$, and
$c^r\det(M')^s=c^r\det(M)^s$ is the transition function of $\mathcal A$.

Otherwise, let $\operatorname{Spec}R_0\to Z$ be such a morphism. Let $\xi$ be the generic point of $Z$ and $\kappa$ its
residue field, the function field of $Z$. The image of the generic point of $\operatorname{Spec}R_0$ has a residue
field of characteristic zero, so the integers are units in its local ring, and $\kappa$ has characteristic zero. So
$R=\kappa[[\tau]]$ is a complete discrete valuation ring of characteristic zero, and the composite
$\chi:\operatorname{Spec}R\to\operatorname{Spec}\kappa\to Z$ maps both points of $\operatorname{Spec}R$ to $\xi$. The
point $\xi$ lies in the dense open set $j(W)$ and in some chart $U_a$, so $\chi$ factors through both, as
$\chi=j\circ\chi'$. Through the morphisms $G_W\to G_V\to G_{\mathcal M}$ and $G_W\to G$, the base change
$G\times_Z\operatorname{Spec}R$, with $L$, is identified with the pullback $A_R$ of $G_{\mathcal M}$ along
$i\circ\chi'$, with the pullback of its line bundle. So $A_R$ is an abelian scheme of relative dimension $g$ over $R$,
with a line bundle that is relatively ample, trivialized along the zero section, isomorphic to its pullback under $-1$
and defines a principal polarization; and the generic point of $\operatorname{Spec}R$ maps into the abelian locus of
$G$. Condition (b) for $\chi$ says that the pullbacks of $b_{a,1},\dots,b_{a,n_0}$ form a basis over $R$ of the sections
of $L^{m^2}$ over $A_R$. If $m\ge1$, Theorem~\ref{cv-volume-key-sections}, with $S=\operatorname{Spec}R$ and $n=m^2$,
shows that these sections also have a basis of $m^{2g}=n_1$ elements, as every open neighbourhood of the closed point of
$\operatorname{Spec}R$ is all of it. If $m=0$, $L^0$ is the trivial line bundle, whose sections over $A_R$ are the global
functions on the abelian scheme $A_R$, which are those of $R$ \cite[Tag 0E0L]{Sta}; they have the basis $1$ of $n_1=1$
element. Two bases of a free module over a nonzero commutative ring have the same number of elements, so $n_0=n_1$, and
the given presentation is the required one.
\end{proof}
% Lean: hasModel_rank: the case without complete discrete valuation points is
% T.hasTransitionCompatibleFrames_of_no_completeDVR_points (ModelSum.hasFrames_rank_of_isEmpty: noTraitZero and
% noTrait_exists_gl_det_eq); otherwise T.exists_completeDVR_point_genericPoint (GenTrait.exists_trait_genericPoint:
% R = PowerSeries of the function field, charZero_functionField), exists_chart, the factorization through C.j, A.trait
% at a, and rank_eq: pulledSectionsFormBasis_id_of_pulledSectionsFormBasis, exists_sectionPower_transport_cartesian_square,
% relAmple_pull and kTriv_pull, and T.card_pulledSectionsFormBasis_eq_pow_dim (LN.SecRank.card_basis_eq_pow) for N >= 1;
% T.card_basis_polarisation_power_zero_eq_one for N = 0, where the Lean's count max 1 ((m^2)^g) is 1. About 1,260
% lines, counted as above.

\begin{lemma}[Alp\"oge]\label{cv-toroidal-theta-envelope-glue-triple}
In Situation~\ref{cv-function-situation}, let $X$ be a scheme, let $\psi_1:X_1\to X$ and $\psi_2:X_2\to X$ be open
immersions whose images cover $X$, and let $X_0=X_1\times_XX_2$, with projections $\varphi_1:X_0\to X_1$ and
$\varphi_2:X_0\to X_2$. For $k=0,1,2$ let $G_k$ be an $\mathcal O_E$-semi-abelian scheme over $X_k$
(Definition~\ref{cv-model-toroidal-semiabelian}(2)), with morphisms $G_0\to G_1$ over $\varphi_1$ and $G_0\to G_2$ over
$\varphi_2$. Suppose that $X_1$ and $X_2$ are reduced and that their generic points lie in the abelian loci of $G_1$ and
$G_2$. Then there is an $\mathcal O_E$-semi-abelian scheme $G$ over $X$ with morphisms $G_1\to G$ over $\psi_1$ and
$G_2\to G$ over $\psi_2$.
\end{lemma}
\leanref{Development/AbelianVarietyTheory/Proofs/ModuliProofs/ExistsAbelianTripleGlueTwoOpenImmersionsStatement.lean}{GlobalTheory.T.exists_abelianTriple_glue_two_open_immersions}
\leanref{Development/AbelianVarietyTheory/Proofs/ModuliProofs/ExistsAbelianTripleGlueTwoOpenImmersions.lean}{GlobalTheory.TripleGlueMap.exists_sab_glue_of_isPullback}
\begin{proof}
Reduction. The open immersion $\psi_k$ identifies $X_k$ with an open subscheme $U_k$ of $X$, and $X_0$ with
$U_1\cap U_2$; transporting $G_0$, $G_1$, $G_2$ and their morphisms along these isomorphisms, we may assume that
$X_1=U_1$ and $X_2=U_2$ are open subschemes of $X$ with $U_1\cup U_2=X$, that $X_0=U_1\cap U_2$, and that $\psi_k$ and
$\varphi_k$ are the inclusions. Then $G_0\to G_k$ identifies $G_0$ with the restriction of $G_k$ to $U_1\cap U_2$.

The group scheme. Glue $G_1$ and $G_2$ along the isomorphism between their restrictions to $U_1\cap U_2$ given by $G_0$.
This gives a scheme $G$ over $X$ with open immersions $\iota_k:G_k\to G$ that identify $G_k$ with the restriction of $G$
to $U_k$. The group laws, unit sections, inverses and actions of $\mathcal O_E$ of $G_1$ and $G_2$ agree over
$U_1\cap U_2$, as the morphisms from $G_0$ respect them, so they glue: $G$ is a commutative group scheme over $X$ with an
action of $\mathcal O_E$, and it is semi-abelian, since this is local on $X$.

The line bundle. Let $\mathcal L_k$ be the line bundle of $G_k$. The isomorphisms of the pullbacks of $\mathcal L_1$ and
$\mathcal L_2$ to $G_0$ with its line bundle give an isomorphism between the restrictions of $\mathcal L_1$ and
$\mathcal L_2$ over $U_1\cap U_2$; gluing gives a line bundle $\mathcal L$ on $G$ with $\iota_k^*\mathcal L\cong\mathcal
L_k$. Let $\pi:G\to X$ be the structure morphism, $0$ the zero section, and $\mathcal M=\mathcal L\otimes\pi^*0^*\mathcal
L^{-1}$. Then $0^*\mathcal M$ is trivial; fix a trivialization. As $\mathcal L_k$ is trivialized along the zero section,
$\iota_k^*\mathcal M\cong\mathcal L_k$.

Rigidity. The scheme $X$ is reduced, as $U_1$ and $U_2$ are, and its generic points lie in $U_1$ or $U_2$, hence in the
abelian locus of $G$. Let $n\ge1$ and $G^n=G\times_X\dots\times_XG$, smooth over $X$ with geometrically integral fibres.
For an open $W\subset X$, an invertible function on the preimage of $W$ in $G^n$ that is $1$ along the zero section is
$1$, by the argument of the proof of Lemma~\ref{cv-toroidal-universal-cubical}: over a generic point of $W$ the fibre
of $G^n$ is an abelian variety, whose regular functions are constant, and the preimage is reduced with its generic
points over those of $W$. Hence a line bundle on $G^n$ that is trivial over $U_1$, over $U_2$ and along the zero section
is trivial: choose generators over $U_1$ and over $U_2$ that agree with a generator along the zero section, after
dividing each by the pullback of a unit of the base; their ratio over $U_1\cap U_2$ is then $1$, and they glue.

The identities. Each of the following line bundles is trivial over $U_1$ and over $U_2$, because there $\mathcal M$ is
isomorphic to $\mathcal L_k$, compatibly with the actions of $\mathcal O_E$, and $\mathcal L_k$ satisfies the
corresponding condition of Definition~\ref{cv-model-toroidal-semiabelian}(2); and it is trivial along the zero section,
since $0^*\mathcal M$ is: the line bundle $\Theta$ of $\mathcal M$ on $G^3$ (Definition~\ref{cv-germs-chart}(1));
$[-1]^*\mathcal M\otimes\mathcal M^{-1}$ on $G$; $\Lambda_{\mathcal M}(cy_1,y_2)\otimes\Lambda_{\mathcal
M}(y_1,\bar cy_2)^{-1}$ on $G^2$ for each $c\in\mathcal O_E$, with $y_1,y_2$ the projections; and $\Lambda_{\mathcal
M}(y,w_0y)\otimes\mathcal M^{-1}$ on $G$, with $y$ the identity. By rigidity they are trivial.

The $\mathcal O_E$-semi-abelian scheme. Let $\theta$ be a generator of $\Theta$ for $\mathcal M$, divided by the pullback
of its restriction to the zero section. Each identity required of a cubical structure in
Definition~\ref{cv-germs-chart}(1) compares two generators of a line bundle on $G^2$, $G^3$ or $G^4$ that agree along the
zero section; by rigidity they are equal, so $\theta$ is a cubical structure. Likewise an isomorphism $[-1]^*\mathcal
M\cong\mathcal M$, divided by its restriction to the zero section, makes $\mathcal M$ symmetric. So $G$ with the action
and $\mathcal M$ is an $\mathcal O_E$-semi-abelian scheme over $X$.

The morphisms. For $k=1,2$, $\iota_k$ identifies $G_k$ with the restriction of $G$ to $U_k$, commutes with the actions
and satisfies $\mathcal L_k\cong\iota_k^*\mathcal M$, and $U_k$ is reduced with its generic points in the abelian locus
of $G_k$. So Lemma~\ref{cv-toroidal-universal-cubical} shows that $\iota_k$ is a morphism of $\mathcal O_E$-semi-abelian
schemes over the inclusion of $U_k$.
\end{proof}
% Lean: TripleGlueMap.exists_sab_glue_of_isPullback: exists_transport, generic_transport, the iso e0 of X0 with the
% intersection, ObjMaps.exists_sab_comp (the reduction), then TripleGlue.exists_sab_glue: GroupGlue.exists_glue_two (G, a,
% psi_1, psi_2), GlueTwoLB.glue (L), cls_surj and pull_rigidify (M = L - pull (toUnit ≫ eSec) L), isReduced_of_cover and
% abLocus_pullback (X), smooth_geometricallyIntegral_abLocus_tensorObj (G^2, G^3), eq_zero_of_cover
% (RigidTriv.exists_gen_of_rigid with eq_one_of_appLE_section_eq_one) applied to the classes given by local_identities,
% CubeGen.cubical_and_symmetric (the cubical structure and the symmetry) and exists_sabMap_of_lbIso (Lemma
% cv-toroidal-universal-cubical) for n_1, n_2.

\begin{lemma}[Alp\"oge]\label{cv-toroidal-theta-envelope-glue-sum}
In Situation~\ref{cv-function-situation}, let $O$ be a noetherian ring, let $\mathcal M$ be a reduced scheme over $O$,
and let $G_{\mathcal M}$ be an $\mathcal O_E$-semi-abelian scheme over $\mathcal M$, proper over $\mathcal M$. Let
$i_1:V_1\to\mathcal M$ and $i_2:V_2\to\mathcal M$ be open immersions with disjoint images whose union is dense in
$\mathcal M$, with $V_k$ regarded over $O$ through $i_k$, and for $k=1,2$ let $G_k$ be an $\mathcal O_E$-semi-abelian
scheme over $V_k$ with a morphism $G_k\to G_{\mathcal M}$ over $i_k$. Let $m\ge0$, $n_0\ge0$, $r$ and $s$ be integers.
If $(V_1,G_1)$ and $(V_2,G_2)$ have theta envelopes of level $m^2$ over $O$ (Definition~\ref{cv-toroidal-envelope}(2))
whose ample theta presentations have $n_0$ sections and integers $r$ and $s$, then so has $(\mathcal M,G_{\mathcal M})$.
\end{lemma}
\leanref{Development/AbelianVarietyTheory/Proofs/ModuliProofs/GenericPointsOpenIntegralCover.lean}{GlobalTheory.ModelSum.HasModel.sum}
\begin{proof}
For $k=1,2$ let $(Z_k,G'_k,W_k,j_k)$ be the theta envelope of $(V_k,G_k)$, with $\mathcal O_E$-semi-abelian scheme
$G_{W_k}$ over $W_k$, its morphisms to $G_k$ over the inclusion and to $G'_k$ over $j_k$, and the presentation $A_k$
of $G'_k$.

The open part. The morphism $V_1\sqcup V_2\to\mathcal M$ given by $i_1$ and $i_2$ is an open immersion, as the images
are disjoint, and its image is dense. Let $Y=W_1\sqcup W_2$, and let $a:Y\to\mathcal M$ be the composite of
$Y\subset V_1\sqcup V_2$ with it, an open immersion with dense image; $Y$ is reduced, as $\mathcal M$ is. Let $G_Y$ be
the pullback of $G_{\mathcal M}$ along $a$, with its morphism to $G_{\mathcal M}$; it is proper over $Y$. For $k=1,2$
the composite $G_{W_k}\to G_k\to G_{\mathcal M}$ is a morphism over $W_k\to\mathcal M$, which factors through $a$; as
$G_Y$ is the pullback of $G_{\mathcal M}$, it factors through a morphism $G_{W_k}\to G_Y$ over the inclusion of $W_k$
in $Y$.

The scheme. Let $Z=Z_1\sqcup Z_2$, over $O$ through the structure morphisms of $Z_1$ and $Z_2$. It is normal and flat
over $O$. It is projective over $O$. Let $Z_k$ be a closed subscheme of $\mathbf P^{e_k}_O$ over $O$, and embed
$\mathbf P^{e_1}_O$ and $\mathbf P^{e_2}_O$ in $\mathbf P^{e_1+e_2+1}_O$ as the linear subspaces on which the last
$e_2+1$, respectively the first $e_1+1$, homogeneous coordinates vanish. These closed subschemes are disjoint: a point of
the first lies in the chart where one of the first $e_1+1$ coordinates is invertible, and that coordinate vanishes on
the second. A morphism from $Z_1\sqcup Z_2$ that is a closed immersion on each $Z_k$, with disjoint images, is a closed
immersion; so $Z_1\sqcup Z_2\to\mathbf P^{e_1+e_2+1}_O$ is a closed immersion over $O$. So $Z$ is noetherian, and it is
reduced. The morphism $j=j_1\sqcup j_2:Y\to Z$ is an open
immersion over $O$ with dense image; it is the composite of $j_1\sqcup\mathrm{id}:Y\to Z_1\sqcup W_2$ and
$\mathrm{id}\sqcup j_2:Z_1\sqcup W_2\to Z$, both open immersions with dense image.

The semi-abelian scheme. The generic points of the noetherian scheme $Z_k$ lie in the dense open image of $j_k$. Over
it, $G'_k$ is identified through $G_{W_k}\to G'_k$ with $G_{W_k}$, and so with the restriction of the proper $G_Y$;
so these generic points lie in the abelian locus of $G'_k$. Lemma~\ref{cv-toroidal-theta-envelope-glue-triple},
applied to the open immersions $j_1\sqcup\mathrm{id}:Y\to Z_1\sqcup W_2$ and $Z_1\to Z_1\sqcup W_2$, whose fibre product
is $W_1$, and to $G_{W_1}$ with its morphisms to $G_Y$ and to $G'_1$, gives an $\mathcal O_E$-semi-abelian scheme $G''$
over $Z_1\sqcup W_2$ with morphisms $G_Y\to G''$ over $j_1\sqcup\mathrm{id}$ and $G'_1\to G''$ over the inclusion. The
scheme $Z_1\sqcup W_2$ is reduced, as an open subscheme of $Z$, and its generic points lie in the image of $Y$, hence
in the abelian locus of $G''$. The same lemma, applied to $\mathrm{id}\sqcup j_2:Z_1\sqcup W_2\to Z$ and $Z_2\to Z$,
whose fibre product is $W_2$, and to $G_{W_2}$ with its morphisms $G_{W_2}\to G_Y\to G''$ and $G_{W_2}\to G'_2$, gives
an $\mathcal O_E$-semi-abelian scheme $G$ over $Z$ with morphisms $G''\to G$ over $\mathrm{id}\sqcup j_2$ and
$G'_2\to G$ over the inclusion of $Z_2$. Composing, there are morphisms $G'_1\to G$ over the inclusion of $Z_1$ and
$G_Y\to G$ over $j$.

The conditions. The morphisms $G'_k\to G$ identify $G'_k$ with the restriction of $G$ to $Z_k$; so the open subschemes
$G'_1$ and $G'_2$ of $G$ are disjoint and cover $G$.

The scheme $G$ is quasi-projective over $O$. Let $G'_k\to\mathbf P^{f_k}_O$ be an immersion over $O$, which is
quasi-compact. Then $G=G'_1\sqcup G'_2\to\mathbf P^{f_1}_O\sqcup\mathbf P^{f_2}_O$ is an immersion, and composed with the
closed immersion $\mathbf P^{f_1}_O\sqcup\mathbf P^{f_2}_O\to\mathbf P^{f_1+f_2+1}_O$ of the previous paragraph but one
it is an immersion over $O$. It is quasi-compact, as $G$ is the union of the two quasi-compact spaces $G'_k$ and the
projective space is separated over $O$.

The scheme $G$ is theta-split. Each of the three conditions of Definition~\ref{cv-model-toroidal-semiabelian}(3) holds for
$G$ because it holds for $G'_1$ and $G'_2$ over the open sets $Z_1$ and $Z_2$, which cover $Z$: (a) the abelian locus
of $G$ meets $Z_k$ in that of $G'_k$, so it is open, and properness over it is local on the base; (b) is a condition at
each point of $Z$; (c) asks for a neighbourhood of each point.

The presentations $A_1$ and $A_2$ give an ample theta presentation of $G$ of level $m^2$ with $n_0$ sections and integers
$r$ and $s$. Its line bundle is the line bundle $\mathcal A$ on $Z$ that is $\mathcal A_k$ on $Z_k$. It is ample: a
point $z$ of $Z_k$ lies in the nonvanishing locus, an affine open set, of a section of $\mathcal A_k^d$ for some $d\ge1$,
and this section, extended by $0$ on the other part, is a section of $\mathcal A^d$ with the same nonvanishing locus. Its
charts are those of $A_1$ and of $A_2$, open subsets of $Z$ that cover it and on which $\mathcal A$ is trivialized, with
the sections $\vartheta_a$ and $b_{a,k'}$ carried to $G$ by the morphisms $G'_k\to G$. Condition (a) of
Definition~\ref{cv-toroidal-envelope}(1) holds for two charts in the same $Z_k$ as it does for $A_k$, as $G$ over $Z_k$
is $G'_k$ with its line bundle, and two charts in different $Z_k$ do not meet. Condition (b) holds for a chart in $Z_k$
as it does for $A_k$, as a morphism $\operatorname{Spec}V\to Z_k$ maps the generic point into the abelian locus of $G$
only if it maps it into that of $G'_k$. So $Z$,
$G$ with this presentation, $j$, $a$ and $G_Y$ form a theta envelope of $(\mathcal M,G_{\mathcal M})$ in the form of the
proof of Lemma~\ref{cv-toroidal-theta-envelope-glue-dense}.
\end{proof}
% Lean: HasModel.sum: isOpenImmersion_desc and range_desc (V_1 ⊔ V_2 -> M), a = coprod.map W1.ι W2.ι ≫ coprod.desc,
% isNormalScheme_coprod, ProjSpace.isProjOver_coprodDesc, noetherianSpace_of_isProjOver, exists_pull (tY), map_comp
% and map_cancel (p1, p2), two applications of T.exists_abelianTriple_glue_two_open_immersions (with
% isPullback_inl_inl_coprodMap and isPullback_inr_inr_coprodMap, and genericPoints_subset_abLocus), amp_of_two
% (T.transitionFrames_push_openImmersion_image, T.isAmpleCechLineBundle_disjointUnion_pushforward,
% T.hasFrames_disjointUnionCechData), qproj_of_two (T.isQProjOver_of_isCompl_openImmersions), good_of_two
% (T.abelianLocus_isOpen_of_openImmersion_cover and its two companions), flatness by IsZariskiLocalAtSource, and
% hasModel_of_maps. About 4,830 lines, counted as above, of which about 2,350 are those of Lemma
% cv-toroidal-theta-envelope-glue-triple. The four facts on disjoint unions follow the Lean's proofs, which have been
% read: projectivity is ProjSpace.isProjOver_coprodDesc = isProjOver_sigmaDesc (the linear embeddings linearι with
% disjoint sets of coordinates, disjoint by exists_mem_chart and linearι_preimage_chart_of_notMem, and
% isClosedImmersion_sigmaDesc); quasi-projectivity is qproj_of_two = ProjSpace.isQProjOver_of_isCompl (the two parts are
% the preimages of Z_1 and Z_2, the coproduct of the immersions followed by the closed immersion of
% isProjOver_coprodDesc, quasi-compactness from the two parts); theta-splitting is good_of_two
% (AbLocusLocal.abOpen_of_cover, splitFib_of_openCover, ThetaLocal.hasTheta_of_cover); the presentation is amp_of_two
% (ModelSum.hasFrames_pushLB for the charts of each A_k with its sections carried by the cartesian square,
% UnionAmple.sumLB_isLB_ampleFam for the line bundle and its ampleness, ModelSum.hasFrames_sumLB for charts in different
% parts). Sizes (tools/lean_size.py): amp_of_two 1,826 lines, good_of_two 822, isProjOver_coprodDesc 186, qproj_of_two
% 110; 2,710 together, mostly general facts on line bundles and sections under base change along open immersions.

\begin{lemma}[Alp\"oge]\label{cv-toroidal-theta-envelope-glue}
In Situation~\ref{cv-function-situation}, let $O$ be a noetherian ring, let $\mathcal M$ be a reduced scheme,
locally of finite type and quasi-compact over $O$, and let $G_{\mathcal M}$ be an $\mathcal O_E$-semi-abelian scheme
over $\mathcal M$, proper and smooth of some relative dimension over $\mathcal M$, whose line bundle is relatively
ample and defines a principal polarization. Let $m\ge0$. Suppose that every generic point $\xi$ of $\mathcal M$ lies
in the image of an open immersion $i:V\to\mathcal M$, with $V$ integral, for which there are an
$\mathcal O_E$-semi-abelian scheme $G_V$ over $V$ with a morphism $G_V\to G_{\mathcal M}$ over $i$ and a theta
envelope of level $m^2$ of $(V,G_V)$ over $O$ (Definition~\ref{cv-toroidal-envelope}(2)) whose scheme is integral and
whose ample theta presentation has integers $r_\xi$ and $1$, for some $r_\xi$, such that for every integer
$r\ge r_\xi$ its $\mathcal O_E$-semi-abelian scheme has an ample theta presentation of level $m^2$ with integers $r$
and $1$. That the second integer of the presentation of each of these theta envelopes is $1$ is not used in the
proof. Then $(\mathcal M,G_{\mathcal M})$ has a theta envelope of level $m^2$ over $O$.
\end{lemma}
\leanref{Development/AbelianVarietyTheory/Proofs/ModuliProofs/GenericPointsOpenIntegralCoverStatement.lean}{GlobalTheory.T.generic_points_open_integral_cover}
\leanref{Development/AbelianVarietyTheory/Proofs/ModuliProofs/GenericPointsOpenIntegralCover.lean}{GlobalTheory.Models.nonempty_zOn_of_components}
\begin{proof}
As $\mathcal M$ is locally of finite type over the noetherian ring $O$ and quasi-compact, it is a noetherian
scheme, so it has finitely many generic points. For each generic point $\xi$ choose $i_\xi:V_\xi\to\mathcal M$,
$G_{V_\xi}$ and a theta envelope as in the hypothesis, and let $r$ be an integer with $r\ge r_\xi$ for every $\xi$.
Every point of $V_\xi$ maps to a specialization of $\xi$, as $V_\xi$ is integral and its generic point maps to
$\xi$. So the images $U_\xi$ of the $i_\xi$ are pairwise disjoint: if $i_\xi(x)$ lies in the open set $U_\zeta$, then
so does its generization $\xi$, so $\xi=i_\zeta(x')$ is a specialization of $\zeta$, and $\xi=\zeta$ as $\xi$ is a
generic point. Their union $U$ contains every generic point, so it is dense in $\mathcal M$.

The pieces. Let $g$ be the relative dimension of $G_{\mathcal M}$ and $n_0=\max(1,m^{2g})$. Fix $\xi$. As $r\ge r_\xi$, the
$\mathcal O_E$-semi-abelian scheme of the chosen theta envelope of $(V_\xi,G_{V_\xi})$ has an ample theta presentation
of level $m^2$ with integers $r$ and $1$, and the scheme of the envelope is integral. So
Lemma~\ref{cv-toroidal-theta-envelope-glue-rank}, applied to $i_\xi$, gives a theta envelope of level $m^2$ of
$(V_\xi,G_{V_\xi})$ over $O$ whose presentation has $n_0$ sections and integers $r$ and $1$. Let $G'$ be an
$\mathcal O_E$-semi-abelian scheme over $U_\xi$ with a morphism $G'\to G_{\mathcal M}$ over the inclusion. As
$G_{V_\xi}\to G_{\mathcal M}$ identifies $G_{V_\xi}$ with the pullback of $G_{\mathcal M}$, the morphism
$G'\to G_{\mathcal M}$ factors through a morphism $G'\to G_{V_\xi}$ over the isomorphism $U_\xi\to V_\xi$ inverse to
$i_\xi$. By Lemma~\ref{cv-toroidal-theta-envelope-glue-dense}(2), applied to this isomorphism, $(U_\xi,G')$ has a theta
envelope of level $m^2$ over $O$ whose presentation has $n_0$ sections and integers $r$ and $1$.

The union. For a set $\Xi$ of generic points let $U_\Xi$ be the union of the $U_\xi$ with $\xi\in\Xi$, an open
subscheme of $\mathcal M$, reduced as $\mathcal M$ is. By induction on $\Xi$: for every
$\mathcal O_E$-semi-abelian scheme $G'$ over $U_\Xi$ with a morphism $G'\to G_{\mathcal M}$ over the inclusion,
$(U_\Xi,G')$ has a theta envelope of level $m^2$ over $O$ whose presentation has $n_0$ sections and integers $r$ and
$1$. For $\Xi$ empty, $U_\Xi$ is empty, and the empty scheme, which is normal and projective and flat over $O$, with the
empty $G'$, the empty open subscheme and the presentation with no charts, is such an envelope. Let
$\Xi=\Xi'\cup\{\xi\}$ with $\xi\notin\Xi'$. The open subschemes $U_\xi$ and $U_{\Xi'}$ of $U_\Xi$ are disjoint and cover
it. Pulling back $G'$ to them gives $\mathcal O_E$-semi-abelian schemes with morphisms to $G'$, and hence to
$G_{\mathcal M}$; by the previous paragraph and by induction, both have theta envelopes with $n_0$ sections and integers
$r$ and $1$. The scheme $G'$ is proper over $U_\Xi$, as $G_{\mathcal M}$ is proper over $\mathcal M$. So
Lemma~\ref{cv-toroidal-theta-envelope-glue-sum}, applied to the reduced scheme $U_\Xi$ over $O$, $G'$ and the two
inclusions, gives such an envelope of $(U_\Xi,G')$.

The conclusion. Let $G_U$ be the pullback of $G_{\mathcal M}$ to $U$, with its morphism to $G_{\mathcal M}$. The case of
the set of all generic points gives a theta envelope of level $m^2$ of $(U,G_U)$ over $O$. As $U$ is dense in
$\mathcal M$, Lemma~\ref{cv-toroidal-theta-envelope-glue-dense}(1), applied to the inclusion of $U$, gives one of
$(\mathcal M,G_{\mathcal M})$.
\end{proof}
% Lean: Models.nonempty_zOn_of_components: genericPoints.finite, choose (the hypothesis hs : C.amp.s = 1 is not used),
% the bound r of the C.amp.r, gen with eq_of_specializes_of_mem_genericPoints (the disjointness), then hasModel_sup with
% n = max 1 ((m^2)^g) (the induction: hasModel_of_isEmpty and HasModel.sum), whose hypothesis for each U_xi is hC at r,
% map_cancel with isoOpensRange_inv_comp, hasModel_rank and HasModel.shrink; exists_pull to the union, isDominant_ι
% with genericPoints.closure, and HasModel.grow. Size, counted through the declarations reached in the imported files
% beyond the laws of abelian schemes: about 5,300 lines, of which about 2,350 are those of Lemma
% cv-toroidal-theta-envelope-glue-triple.

\begin{lemma}[Alp\"oge]\label{cv-toroidal-theta-envelope}
In Situation~\ref{cv-function-situation}, suppose $p+q\ge1$, and let $n$ be an odd prime. There are
integers $c\ge0$ and $j_0$ such that for every $j\ge j_0$, with $N'=n^j$ and $N=N'^2$, and every PEL
theta family $(\mathcal M,t,G)$ of level $N'$ (Definition~\ref{cv-toroidal-theta}(7)), the pair
$(\mathcal M,G)$ has a theta envelope of level $n^{2c}$ over $R_{n,N}$
(Definition~\ref{cv-toroidal-envelope}). Of the PEL theta family, the proof uses only that $(\mathcal M,t)$ is a
theta family and the isomorphism of group schemes from $G$ to the abelian scheme of $t$ that carries the line bundle
of $G$ to that of $t$. The hypothesis $p+q\ge1$ and the other conditions of a PEL theta family are not used in the
proof: that the basis $(e_i)$ is symplectic, that the image of $a$ generates $\mathcal O_E$ and contains $w_0$ and
$\bar w_0$, that $\zeta=\zeta_N$, that the fibre of $\mathcal M$ over the fraction field of $R_{n,N}$ is smooth, and
the conditions of the PEL structure on the action of $\mathcal O_E$ and at the points of characteristic zero.
\end{lemma}
\leanref{Development/AbelianVarietyTheory/Proofs/ModuliProofs/NonemptyAmpleFramedModelPELThetaFamilyStatement.lean}{GlobalTheory.T.nonempty_ampleFramedModel_PELThetaFamily}
\leanref{Development/AbelianVarietyTheory/Proofs/ModuliProofs/NonemptyAmpleFramedModelPELThetaFamily.lean}{GlobalTheory.nonempty_zOn_thetaFam}
\begin{proof}
Let $2g$ be the rank of $\Lambda$ over $\mathbf Z$. Lemma~\ref{cv-toroidal-theta-envelope-component}, for $g$ and
$n$, gives integers $c\ge0$ and $j_0$. Let $j\ge j_0$, $N'=n^j$ and $N=N'^2$, and let $(\mathcal M,t,G)$ be a PEL
theta family of level $N'$, with $t=(A,L,(x_i),(\alpha_h))$ and structure morphism
$s:\mathcal M\to\operatorname{Spec}R_{n,N}$. Let $\varphi:G\to A$ be the isomorphism of the PEL structure; the line
bundle of $G$ is isomorphic to $\varphi^*L$. The symplectic basis of
Definition~\ref{cv-toroidal-theta}(7) has $2g$ elements, so $A$ has relative dimension $g$.

The ring $R_{n,N}$ is a Dedekind domain of finite type over $\mathbf Z$, so it is noetherian, and $s$ is locally of
finite type and quasi-compact, as $\mathcal M$ is quasi-projective over $R_{n,N}$. Through $\varphi$, $G$ is proper
and smooth of relative dimension $g$ over $\mathcal M$, and its line bundle is relatively ample and defines a
principal polarization, because $A$ and $L$ have these properties. So
Lemma~\ref{cv-toroidal-theta-envelope-glue}, for $O=R_{n,N}$, $\mathcal M$, $G$ and $m=n^c$, gives a theta envelope
of level $n^{2c}$ of $(\mathcal M,G)$ over $R_{n,N}$, once its hypothesis is checked at each generic point $\xi$ of
$\mathcal M$.

Lemma~\ref{cv-toroidal-reduction-local} gives an open immersion $x:\operatorname{Spec}R_0\to\mathcal M$ whose image
contains $\xi$, and a family $v_0$ of theta constants of $t$ at $x$. As $\mathcal M$ is reduced, flat and locally of
finite type over $R_{n,N}$, Lemma~\ref{cv-toroidal-theta-envelope-chart} gives a field $\Omega$, an
$R_{n,N}$-subalgebra $R\subset\Omega$ as there and a homomorphism $\psi:R_0\to R$ such that
$i=x\circ\operatorname{Spec}\psi$ is an open immersion over $R_{n,N}$ whose image contains $\xi$, with
$R_{n,N}\to\Omega$ injective; so $\Omega$ has characteristic zero. By
Lemma~\ref{cv-toroidal-reduction-base-change}(4), $\psi(v_0)$ is a family of theta constants of $t$ at $i$.
Lemma~\ref{cv-toroidal-theta-envelope-pull}, applied to $t$, $G$, $\varphi$ and $i$, gives a theta-level object $u$
over $\operatorname{Spec}R$ with a morphism $u\to t$ over $i$, an $\mathcal O_E$-semi-abelian scheme $G_0$ over
$\operatorname{Spec}R$ with a morphism $G_0\to G$ over $i$, and an isomorphism $e$ from $G_0$ to the abelian scheme
of $u$ as there. By Lemma~\ref{cv-toroidal-reduction-base-change}(3), $\psi(v_0)$ is a family of theta constants of
$u$ at the identity. By Lemma~\ref{cv-toroidal-theta-envelope-ratios}, for the injective inclusion
$R_{n,N}\subset\mathbf C$, $\Omega$ is generated as a field by $R_{n,N}$ and the quotients of the entries of
$\psi(v_0)$, so it is finite over that subfield. Lemma~\ref{cv-toroidal-theta-envelope-component}, applied to
$R_{n,N}$, $\zeta_N$, $\Omega$, $R$, $u$, $G_0$, $e$ and $\psi(v_0)$, gives a theta envelope of level $n^{2c}$ of
$(\operatorname{Spec}R,G_0)$ over $R_{n,N}$ with the properties required in
Lemma~\ref{cv-toroidal-theta-envelope-glue}. As $R$ is a subring of a field, $\operatorname{Spec}R$ is integral. So
$V=\operatorname{Spec}R$, $i$, $G_0$ and this envelope satisfy the hypothesis at $\xi$.
\end{proof}
% Lean: nonempty_zOn_thetaFam = FamilyModel.nonempty_zOn_of_components applied to
% T.exists_model_triple_ampleFrames_component (hz, for g = finrank La / 2; obtain rfl from bas.finrank_eq). Of the
% hypotheses only F, tau and the first three parts phi, hphi, hL of IsPEL are used (rintro ... - - - - ⟨φ, hφ, hL, -⟩
% -); 1 <= p + q, IsSympBasis, the closure of the range of a, w and w', the value of b.ζ and the smoothness of the
% generic fibre are discarded. BaseRing.isDedekindDomain_RZ and BaseRing.finiteType_RZ (R_{n,N}), F.qproj with
% isProper_str (finite type, quasi-compact), pol_pull (proper, RelAmple, KTriv) and the smoothness of F.univ through
% isGrpPullback_of_iso; then T.generic_points_open_integral_cover with, at xi, ThetaLocal.exists_hasTh (Lemma
% cv-toroidal-reduction-local), exists_integral_chart, T.hasThetaConstants_ring_baseChange, ArithBase, exists_pull,
% closure_ratios_eq_top, hasTh_iff_of_tmap and hz. Inputs Out.Laws, Out.Laws2, SemistableReductionSt (Theorem
% cv-semistable-reduction), ThmUSt and UTorsSt (Theorem cv-toroidal-mumford) and ZHGeoSt (proved by
% T.zarhinGeometric_of_semistableReduction) are hypotheses. Size of the whole, counted through the declarations
% reached in the imported files: about 67,000 lines beyond the laws of abelian schemes, about 41,600 beyond them and the
% proof of Theorem cv-toroidal-mumford; Lemma cv-toroidal-theta-envelope-component accounts for about 62,800 and -glue
% for about 5,300. The proof of ZHGeoSt adds about 92,600 lines; the earlier estimate of 105,000 lines included it.

\subsection{Level structures and flat kernels}\label{subsection-cv-flat-kernels}

\begin{lemma}[Alp\"oge]\label{cv-toroidal-abelian-pullback}
In Situation~\ref{cv-function-situation}, let $n$ be a prime, $N\ge1$, and $R_{n,N}$ as in
Definition~\ref{cv-toroidal-family}. Let $S$ be a scheme over $R_{n,N}$, let $G$ be an
$\mathcal O_E$-semi-abelian scheme over $S$ (Definition~\ref{cv-model-toroidal-semiabelian}(2)), let
$\lambda:\Lambda\to G(S)$ be a map, and let $f:S'\to S$ be a morphism of schemes, with $S'$ regarded as
a scheme over $R_{n,N}$ through $f$.
\begin{enumerate}
\item There are an $\mathcal O_E$-semi-abelian scheme $G'$ over $S'$, a morphism $\nu:G'\to G$ over $f$
  and a map $\lambda':\Lambda\to G'(S')$ with $\nu\circ\lambda'(a)=\lambda(a)\circ f$ for every
  $a\in\Lambda$.
\item Let $G'$, $\nu$ and $\lambda'$ be as in (1). If $(G,\lambda)$ is a full-level object over $S$
  (Definition~\ref{cv-toroidal-family}(1)), then $(G',\lambda')$ is a full-level object over $S'$.
\end{enumerate}
\end{lemma}
\leanref{Development/AbelianVarietyTheory/Proofs/ModuliProofs/FullLevelObjectPullbackExistsAndFullStatement.lean}{GlobalTheory.T.fullLevelObject_pullback_exists_and_full}
\leanref{Development/AbelianVarietyTheory/Proofs/ModuliProofs/FullLevelObjectPullbackExistsAndFull.lean}{GlobalTheory.fullObj_pull}
\begin{proof}
For (1), let $G'=f^*G$ with the pulled-back action of $\mathcal O_E$ and line bundle, let $\nu$ be the
projection, and let $\lambda'(a)$ be the section of $G'$ given by $\lambda(a)\circ f$.

For (2), $\nu$ identifies $G'$ with $f^*G$, and the line bundle $\mathcal L'$ of $G'$ with the pullback of
the line bundle $\mathcal L$ of $G$; so $\lambda'(a)$ is the section given by $\lambda(a)\circ f$. Hence
$G'$ is proper over $S'$, $\mathcal L'$ is relatively ample \cite[Tag 0893]{Sta}, and $\lambda'$ is
additive and vanishes on $N\Lambda$. The polarization of $\mathcal L'$ is principal: a point $x$ of $G'$
over an $S'$-scheme $T$ such that $t_x^*\mathcal L'\otimes\mathcal L'^{-1}$ is trivial locally on $T$,
where $t_x$ is the translation by $x$ on $G'_T$, is through $\nu$ a point of $G$ over $T$ with the same
property, hence zero. Let $k$ be an algebraically closed field of characteristic zero and
$x:\operatorname{Spec}k\to S'$ a point. Then $\nu$ identifies $G'_x$ with $G_{f\circ x}$, compatibly with
the actions of $\mathcal O_E$ and the line bundles, hence with the eigenspaces of the Lie algebras and
with the commutator pairings of the theta groups, and it carries $\lambda'(a)(x)$ to
$\lambda(a)(f\circ x)$; the homomorphisms $R_{n,N}\to k$ induced by $x$ and by $f\circ x$ are equal. So
conditions (a) and (b) of Definition~\ref{cv-toroidal-family}(1) at $x$ follow from the same conditions
at $f\circ x$, with the same roots of unity $\zeta'_m$ and with the maps $\iota_m$ composed with this
identification.
\end{proof}
% Lean: fullObj_pull: sabPull with Out.CubePf.cubePull (the pullback), ObjPull.secPull,
% PolPull.pol_pull_of_relAmple_local (RelAmple and KTriv pull back), LiePull.lieMult_pull, and
% ObjPull.tateMove with ObjPull.pushEquiv and weilIs_grpPullback_of_comm.

\begin{lemma}[Alp\"oge]\label{cv-toroidal-abelian-torsion}
Let $S$ be a locally noetherian scheme whose local rings are integrally closed domains, let $j:W\to S$
be an open immersion with dense image, and let $N\ge1$ be an integer that is invertible on $S$. Let $G$
be a commutative group scheme over $S$, proper over $S$, let $G'$ be a group scheme over $W$ with a
morphism $\varphi:G'\to G$ over $j$ that identifies $G'$ with $j^*G$ as group schemes, let $\Lambda$ be an
abelian group, and let $\lambda':\Lambda\to G'(W)$ be a homomorphism with $\lambda'(N\Lambda)=0$. Then
there is a homomorphism $\lambda:\Lambda\to G(S)$ with $\lambda(N\Lambda)=0$ and
$\varphi\circ\lambda'(a)=\lambda(a)\circ j$ for every $a\in\Lambda$.
\end{lemma}
\leanref{Development/AbelianVarietyTheory/Proofs/ModuliProofs/TorsionLevelExtendsFromDenseOpenNormalStatement.lean}{GlobalTheory.T.torsionLevel_extends_from_dense_open_normal}
\leanref{Development/AbelianVarietyTheory/Proofs/ModuliProofs/TorsionLevelExtendsFromDenseOpenNormal.lean}{GlobalTheory.exists_torsion_level_extend}
\begin{proof}
The scheme $S$ is reduced, as its local rings are domains. Let $[N]$ be multiplication by $N$ on $G$,
and $e:S\to G$ the zero section. The morphism $[N]$ is formally unramified. Indeed, let $g_1,g_2$ be
points of $G$ over an affine $S$-scheme $T$ that agree on a closed subscheme $T_0$ defined by an ideal of
square zero, with $Ng_1=Ng_2$. Then $g=g_1-g_2$ is trivial on $T_0$ and $Ng=0$. The points of $G$ over
$T$ that are trivial on $T_0$ multiply by adding the derivations they define, so $Ng=0$ gives $g=0$, as
$N$ is invertible. Since $G$ is proper over $S$, $[N]$ is proper; being locally of finite type and
formally unramified, it is locally quasi-finite. So $[N]$ is finite \cite[Tag 02LS]{Sta}.

Let $a\in\Lambda$ and $\sigma'=\varphi\circ\lambda'(a):W\to G$. Then $[N]\circ\sigma'=e\circ j$, as
$N\lambda'(a)=\lambda'(Na)=0$. There is a unique $\sigma:S\to G$ with $\sigma\circ j=\sigma'$ and
$[N]\circ\sigma=e$. Uniqueness holds because $W$ is dense, $S$ is reduced and $[N]$ is separated. For
existence, the irreducible components of $S$ are open and disjoint, since $S$ is locally noetherian and
its local rings are domains, and each is an integral normal scheme $S_i$; it suffices to construct the
lift over each $S_i$. The open subscheme $W_i=j^{-1}(S_i)$ is dense in $S_i$, hence integral with the
function field $K_i$ of $S_i$. As $S_i$ is normal, it is the normalization of itself in
$\operatorname{Spec}K_i$. Since $[N]$ is finite, hence integral, the universal property of relative
normalization \cite[Tag 035I]{Sta}, applied to $\operatorname{Spec}K_i\to W_i\to G$ over $e$, gives
$\sigma_i:S_i\to G$ with $[N]\circ\sigma_i=e$ that agrees with $\sigma'$ at the generic point of $W_i$.
Then $\sigma_i$ and $\sigma'$ agree on $W_i$, both being lifts of $e\circ j$ through the separated
morphism $[N]$ that agree on a dense subscheme of the reduced scheme $W_i$. The $\sigma_i$ glue to
$\sigma$, and $\sigma$ is a section of $G$ because $e$ is.

Put $\lambda(a)=\sigma$. For $a,c\in\Lambda$, the sections $\lambda(a+c)$ and $\lambda(a)+\lambda(c)$
agree after composition with $j$, because $\varphi$ is a homomorphism and $\lambda'$ is additive. As $W$
is dense, $S$ reduced and $G$ separated over $S$, they are equal. In the same way $\lambda(Na)=0$.
\end{proof}
% Lean: exists_torsion_level_extend with exists_section_of_pow_eq_one (formallyUnramified_mulA, from
% GlobalTheory.formallyUnramified_mulN and eq_one_of_pow_eq_one; IsFinite.of_isProper_of_locallyQuasiFinite)
% and existsUnique_lift_of_densePart (IsNormalScheme.component, existsUnique_lift_of_birational with
% isIso_fromNormalization_of_birational, lift_unique_of_isDominant), and ext_of_isDominant_of_isSeparated.

\begin{lemma}[Alp\"oge]\label{cv-toroidal-abelian-cover}
In Situation~\ref{cv-function-situation}, let $S$ be a scheme and $G$ an $\mathcal O_E$-semi-abelian
scheme over $S$ (Definition~\ref{cv-model-toroidal-semiabelian}(2)), proper over $S$, with line bundle
$\mathcal L$. Say that a morphism $G'\to G$ of $\mathcal O_E$-semi-abelian schemes over a morphism
$f:S'\to S$ \emph{covers} a point $p$ of $S$ if $p\in f(S')$.
\begin{enumerate}
\item Let $S$ be locally noetherian. Suppose that for every ring $A$ and every morphism
  $\operatorname{Spec}A\to S$ the base change of $G$ to $\operatorname{Spec}A$ is projective over $A$, and
  that every point of $S$ is covered by some $G'\to G$ for which the line bundle of $G'$ is relatively
  ample and defines a principal polarization. Then $\mathcal L$ is relatively ample.
\item Let $S$ be locally noetherian and $\mathcal L$ relatively ample. Suppose that every point of $S$
  is covered by some $G'\to G$ for which the line bundle of $G'$ defines a principal polarization. Then
  $\mathcal L$ defines a principal polarization.
\item Let $n$ be a prime, $N\ge1$, and let $S$ be a scheme over $R_{n,N}$
  (Definition~\ref{cv-toroidal-family}). Suppose that $\mathcal L$ is relatively ample and defines a
  principal polarization, and let $\lambda:\Lambda\to G(S)$ be additive with $\lambda(N\Lambda)=0$.
  Suppose that every point of $S$ is covered by some $\nu:G'\to G$ over $f:S'\to S$ with a map
  $\lambda':\Lambda\to G'(S')$ such that $\nu\circ\lambda'(a)=\lambda(a)\circ f$ for every $a\in\Lambda$
  and $(G',\lambda')$ is a full-level object over $S'$ (Definition~\ref{cv-toroidal-family}(1)). Then
  $(G,\lambda)$ is a full-level object over $S$.
\end{enumerate}
\end{lemma}
\leanref{Development/AbelianVarietyTheory/Proofs/ModuliProofs/RelAmpleOfPolarisedAtEveryPointStatement.lean}{GlobalTheory.T.relAmple_of_polarised_at_every_point}
\leanref{Development/AbelianVarietyTheory/Proofs/ModuliProofs/KernelTrivialOfEveryPointCoveredStatement.lean}{GlobalTheory.T.kernelTrivial_of_every_point_covered}
\leanref{Development/AbelianVarietyTheory/Proofs/ModuliProofs/FullLevelObjectOfCoveringPullbacksStatement.lean}{GlobalTheory.T.fullLevelObject_of_covering_pullbacks}
\begin{proof}
(1) Let $p\in S$, covered by $G'\to G$ over $f:S'\to S$, with $p=f(p')$. The fibre $G'_{p'}$ is the base
change of $G_p$ along $\kappa(p)\to\kappa(p')$, with the pullback of $\mathcal L_p$, which is ample
\cite[Tag 0893]{Sta}. As $\operatorname{Spec}\kappa(p')\to\operatorname{Spec}\kappa(p)$ is an fpqc
covering, $\mathcal L_p$ is ample on $G_p$ \cite[Tag 0D2P]{Sta}. Let $V$ be a noetherian affine open
neighbourhood of $p$. Since $G_V$ is proper over $V$, there is an open neighbourhood $V'\subset V$ of $p$
over which $\mathcal L$ is relatively ample \cite[Tag 0D2N]{Sta}. Relative ampleness is local on $S$, so
$\mathcal L$ is relatively ample.

(2) Let $p\in S$, covered by $G'\to G$ over $f:S'\to S$, with $p=f(q)$, and let $K$ be the residue field
of $S'$ at $q$. The base change of $G$ along $\operatorname{Spec}K\to S'\to S$ is the base change of $G'$,
so its line bundle defines a principal polarization. Let $V=\operatorname{Spec}A$ be an affine open
neighbourhood of $p$, with $A$ noetherian, and $\chi:A\to K$ the induced homomorphism. The points $x$ of
$G_V$ over $V$-schemes $T$ for which $t_x^*\mathcal L\otimes\mathcal L^{-1}$ is trivial locally on $T$
are the points of a closed subscheme $\operatorname{Spec}B$ of $G_V$, finite over $A$, since $\mathcal L$
is relatively ample \cite[Chapter~6]{MFK94}. Let $\alpha:A\to B$ be the structure map, $\sigma:B\to A$
the zero section, and $I=\ker\sigma$, a direct summand of $B$ as an $A$-module, hence finite over $A$.
Every point of $\operatorname{Spec}B$ over a $K$-algebra is zero, since the polarization over $K$ is
principal. Applied to the point of $\operatorname{Spec}B$ over $K\otimes_AB$, this gives
$K\otimes_AI=0$. So the prime $\ker\chi$ is not in the support of $I$, and there is $r\in A$ with
$\chi(r)\ne0$ and $rI=0$. Over the neighbourhood $D(r)$ of $p$, $\operatorname{Spec}B$ is the zero
section, so the polarization of $\mathcal L$ is principal over $D(r)$. Principality is local on $S$.

(3) The conditions on $G$, $\mathcal L$ and $\lambda$ other than (a) and (b) of
Definition~\ref{cv-toroidal-family}(1) hold by hypothesis. Let $k$ be an algebraically closed field of
characteristic zero, $x:\operatorname{Spec}k\to S$ a point and $\beta:R_{n,N}\to k$ the induced
homomorphism. Let $\nu:G'\to G$ over $f:S'\to S$, with $\lambda'$, cover the image of $x$. A point of
$\operatorname{Spec}k\times_SS'$ and an algebraic closure $\Omega$ of its residue field give an embedding
$i:k\to\Omega$ and a point $y:\operatorname{Spec}\Omega\to S'$ with
$f\circ y=x\circ\operatorname{Spec}i$. Conditions (a) and (b) for $(G',\lambda')$ at $y$ give, through
$\nu$, the same conditions for $(G,\lambda)$ at $x\circ\operatorname{Spec}i$, over $i\circ\beta$.
Condition (a) at $x$ follows, as the dimensions of the eigenspaces of the Lie algebra do not change under
$k\to\Omega$. The fibre $G_x$ is an abelian variety over $k$, and for $m\ge1$ its points of order
dividing $m$ over $k$ map bijectively to those over $\Omega$: the map is injective, and in characteristic zero
both groups have $m^{2g}$ elements, where $g$ is the dimension of $G_x$ \cite[\S6, Application~3]{Mum70}; the commutator
pairings correspond. So the roots of unity $\zeta'_m$ and the maps $\iota_m$ at
$x\circ\operatorname{Spec}i$ come from roots of unity in $k$ and maps into $G_x(k)$, which satisfy (b)
at $x$, with $\iota_N(a)=\lambda(a)(x)$ because $\nu$ carries $\lambda'(a)(y)$ to $\lambda(a)(x)$.
\end{proof}
% Lean: (1) AmpleGlobal.relAmple_of_polarised_at_points = relAmple_of_pointwise:
% FibreAmple.relAmple_of_factor and relAmple_of_fieldExtension (the fibre), AmpleNear.relAmple_near_fibre
% (a neighbourhood, using the projectivity of G over affine bases), AmpleLocal.relAmple_of_locally.
% (2) KTrivGlob.kTriv_of_forall_point: kTriv_transfer, T.exists_basicOpen_kTriv_of_fibre =
% KTrivSpread.exists_kTriv_of_isUnit (exists_finite_rep, exists_mul_sub_eq_zero, exists_smul_eq_zero),
% CoverGlue.kTriv_of_cover. (3) FullDescent.fullObj_of_cover: exists_lift, detCond_iff_of_sabMap,
% detCond_iff_of_field_extension, tate_push, torsion_ext (torsion_points_field_extension, from Laws),
% TateIsoFieldExt.down and weilIs_fieldExt_iff.

\begin{lemma}[Alp\"oge]\label{cv-toroidal-abelian-dvr-closed}
Let $R$ be a ring, let $B$ be a commutative group scheme over $R$ that is an abelian scheme and projective
over $R$, and let $\mathcal L$ and $\mathcal O$ be line bundles on $B$. For a scheme $T$ over $R$ and a point
$x\in B(T)$, let $\mathcal L_x=t_x^*\mathcal L_T\otimes\mathcal L_T^{-1}\otimes\mathcal O_T^{-1}$ on
$B_T=T\times_RB$, where $\mathcal L_T$ and $\mathcal O_T$ are the pullbacks of $\mathcal L$ and $\mathcal O$
and $t_x$ is the translation by $x$ on $B_T$.
\begin{enumerate}
\item There is a closed subscheme $H$ of $B$ such that a point $x\in B(T)$ factors through $H$ if and only
  if $\mathcal L_x$ is trivial locally on $T$, that is, $T$ has an open cover by subschemes $U$ with
  $\mathcal L_x$ trivial on $B_U$.
\item If $\mathcal L$ is relatively ample, then $H$ is finite over $R$.
\end{enumerate}
For $\mathcal O=\mathcal O_B$, write $K(\mathcal L)$ for the group of points of $H$; then $\mathcal L$
defines a principal polarization when every point of $K(\mathcal L)$ is zero.
\end{lemma}
\leanref{Development/AbelianVarietyTheory/Proofs/UniformizationProofs/ExistsClosedSubschemeTranslationKernelStatement.lean}{GlobalTheory.T.exists_closed_subscheme_translation_kernel}
\leanref{Development/AbelianVarietyTheory/Proofs/UniformizationProofs/IsAffineOfQuasiAffine.lean}{GlobalTheory.AmpleStabFinite.isFinite_of_relAmple}
\begin{proof}
(1) Let $f:B\times_RB\to B$ be the first projection and let $\mathcal F$ be the line bundle
$\mathcal L_{x_0}$ on $B\times_RB$ for the point $x_0\in B(B)$ given by the identity of $B$. For $x\in B(T)$,
the base change of $f$ along $x:T\to B$ is the projection $p:B_T\to T$, and the pullback of $\mathcal F$ to
$B_T$ is $\mathcal L_x$. The morphism $f$ is flat, proper and of finite presentation, and its geometric fibres
are abelian varieties, so $\mathcal O_T\to p_*\mathcal O_{B_T}$ is an isomorphism for every such $T$
\cite[Tag 0E0L]{Sta}. Hence there is a closed subscheme $H$ of $B$ such that $x$ factors through $H$ if and
only if $\mathcal L_x\cong p^*\mathcal N$ for an invertible $\mathcal O_T$-module $\mathcal N$
\cite[Tags 0BDP and 0BF0]{Sta}. This holds if and only if $\mathcal L_x$ is trivial locally on $T$. Indeed,
$p^*\mathcal N$ is trivial over the open sets on which $\mathcal N$ is. Conversely, if $\mathcal L_x$ is
trivial locally on $T$, then $\mathcal N=p_*\mathcal L_x$ is invertible, since $p_*\mathcal O_{B_T}=\mathcal
O_T$, and the natural map $p^*\mathcal N\to\mathcal L_x$ is an isomorphism, because it is one locally on $T$.

(2) First, let $z\in B(T)$ factor through $H$ and let $w\in B(T)$. Then
\[
 (z+w)^*\mathcal L\otimes w^*\mathcal L^{-1}\otimes w^*\mathcal O^{-1}\cong
 z^*\mathcal L\otimes0^*\mathcal L^{-1}\otimes0^*\mathcal O^{-1}
\]
on $T$. Indeed, $\mathcal L_z\cong p^*\mathcal N$ as in (1), and the two sides are the pullbacks of
$\mathcal L_z$ along the sections $(\mathrm{id},w)$ and $(\mathrm{id},0)$ of $p$, both isomorphic to
$\mathcal N$. If $x,y\in B(T)$ factor through $H$, applying this to $(z,w)=(x,-x-y)$ and to $(y,-x-y)$ and
dividing gives
\[
 x^*\mathcal L\otimes(-x)^*\mathcal L\cong y^*\mathcal L\otimes(-y)^*\mathcal L.
\]

The scheme $H$ is closed in $B$, which is proper over $R$, so $H$ is proper over $R$, and it suffices to show
that its fibres are finite \cite[Tag 02LS]{Sta}. Let $h$ be a point of $H$ with residue field $\kappa$, let
$P=\operatorname{Spec}\kappa$, regarded over $R$ through $h$, and let $H_P=H\times_RP$. Let $x\in B(H_P)$ be the
composite $H_P\to H\to B$ and $y\in B(H_P)$ the composite $H_P\to P\to H\to B$, where $P\to H$ is the point $h$.
Both factor through $H$, and $y$ and $-y$ factor through $\operatorname{Spec}\kappa$, over which every line
bundle is trivial. By the display, $x^*\mathcal L\otimes(-x)^*\mathcal L$ is trivial on $H_P$. Since
$\operatorname{Spec}R$ is affine and $\mathcal L$ is relatively ample, there are $d\ge1$ and sections
$s_0,\dots,s_m$ of $\mathcal L^d$ over $B$ whose nonvanishing loci $B_{s_i}$ are affine and cover $B$. The
morphisms $x$ and $-x$ are affine, as $H_P\to H$ is a base change of $P\to\operatorname{Spec}R$, $H\to B$ is a
closed immersion and $-1$ is an automorphism of $B$. For each $i$ and $j$, the product of $x^*s_i$ and $(-x)^*s_j$ is a section of the trivial bundle
$x^*\mathcal L^d\otimes(-x)^*\mathcal L^d$, that is, a function on $H_P$, whose nonvanishing locus
$x^{-1}(B_{s_i})\cap(-x)^{-1}(B_{s_j})$ is affine. These loci cover $H_P$, which is quasi-compact, so
$\mathcal O_{H_P}$ is ample and $H_P$ is quasi-affine \cite[Tag 01QD]{Sta}. The open immersion of $H_P$ into
$\operatorname{Spec}\Gamma(H_P,\mathcal O_{H_P})$ has dense image, and closed image, since $H_P$ is proper over
$\kappa$; so it is an isomorphism, and $H_P$ is affine and proper over $\kappa$, hence finite over $\kappa$
\cite[Tag 01WN]{Sta}. Every point of $H$ over the image $r$ of $h$ in $\operatorname{Spec}R$ is the image of a
point of $H_P$, as $P$ lies over $r$. So the fibre of $H$ over $r$ is finite.
\end{proof}
% Lean: (1) exists_closed_twist_kernel: T.exists_closedSubscheme_locallyTrivial_product_lineBundle
% (= TrivSubscheme.exists_closed_locTriv) for F = translDiff B L (id) / O on B x B, and cls_twist_pull. The Lean
% proves the representability itself (about 7,200 lines, with Grothendieck's complex for cohomology and base
% change); the text cites Stacks. The hypothesis GlobalFnsSt of the Lean (O_T = p_* O, a law of abelian
% schemes) is Stacks Tag 0E0L. (2) isFinite_of_relAmple: Unif.relAmple_iff_veryAmpleBy, add_inv_eq (from
% transl_const), pull_spec_local, quasiAffine_of_pull_add (with cls_pw), isAffine_of_quasiAffine,
% IsFinite.iff_isProper_and_isAffineHom and Scheme.Pullback.exists_preimage_pullback.

\begin{lemma}[Alp\"oge]\label{cv-toroidal-abelian-dvr-flat-criterion}
Let $V$ be a discrete valuation ring with uniformizer $\varpi$, let $X$ be a scheme over $V$, and let $j:H\to X$ be a
closed immersion. Suppose that for every point $h\in H$ at which the image of $\varpi$ in
$A=\mathcal O_{X,j(h)}$ is not a unit, there is an ideal $I\subset A$ such that $\varpi$ is a nonzerodivisor on
$A/I$ and such that for every local ring $R$ and homomorphism $\phi:A\to R$, the morphism
$\operatorname{Spec}R\to\operatorname{Spec}A\to X$ factors through $H$ if and only if $I\subset\ker\phi$. Then $H$ is
flat over $V$.
\end{lemma}
\leanref{Development/AbelianVarietyTheory/Proofs/UniformizationProofs/LocalizationOfIsStandardSmooth.lean}{GlobalTheory.StabFlat.flat_of_ideals}
\begin{proof}
It suffices to show that $\varpi$ is a nonzerodivisor on each local ring $\mathcal O_{H,h}$: then
$\mathcal O_{H,h}$ is torsion free, hence flat, over $V$ \cite[Tag 0539]{Sta}, and so over the localization of
$V$ at the image of $h$. If the image of $\varpi$ in $A=\mathcal O_{X,j(h)}$ is a unit, this is clear. Otherwise
let $I$ be as in the statement. The homomorphism $\pi:A\to\mathcal O_{H,h}$ induced by $j$ is surjective, and
$\operatorname{Spec}\mathcal O_{H,h}\to\operatorname{Spec}A\to X$ factors through $H$, so $I\subset\ker\pi$. The
ring $A/I$ is local, and $\operatorname{Spec}(A/I)\to X$ factors through $H$, by a morphism that maps the closed
point to $h$, as $j$ is injective; this gives $\psi:\mathcal O_{H,h}\to A/I$ with $\psi\circ\pi$ the quotient map.
Let $b\in\mathcal O_{H,h}$ with $\varpi b=0$, and $b=\pi(a)$. Then $\varpi$ kills the image of $a$ in $A/I$, so
$a\in I\subset\ker\pi$, and $b=0$.
\end{proof}
% Lean: flat_of_ideals: Flat.of_stalkMap, j.stalkMap_surjective, the factorization of Spec(A/I) through H (with
% stalkSpecializes and stalkClosedPointTo for psi), and flat_of_regular for the flatness over V.

\begin{lemma}[Alp\"oge]\label{cv-toroidal-abelian-dvr-flat-stalk}
Let $V$ be a discrete valuation ring with uniformizer $\varpi$ and residue field $\kappa$, let $X$ be smooth of
relative dimension $g$ over $V$, and let $x$ be a closed point of $X$ over the closed point of
$\operatorname{Spec}V$. Then $A=\mathcal O_{X,x}$ is noetherian, the image of $\varpi$ is a nonzerodivisor of $A$,
and $A/\varpi A$ is a regular local ring of dimension $g$.
\end{lemma}
\leanref{Development/AbelianVarietyTheory/Proofs/UniformizationProofs/LocalizationOfIsStandardSmooth.lean}{GlobalTheory.StabFlat.stalk_of_smooth}
\begin{proof}
The ring $A$ is noetherian, as $X$ is locally of finite type over $V$, and $\varpi$ is a nonzerodivisor on $A$, as
$A$ is flat over $V$. The ring $A/\varpi A$ is the local ring of the special fibre $X_\kappa$ at $x$. The scheme
$X_\kappa$ is smooth over $\kappa$, so this local ring is regular \cite[Tag 056S]{Sta}; and $x$ is a closed point
of $X_\kappa$, at which $X_\kappa$ has dimension $g$, so the local ring has dimension $g$ \cite[Tag 0A21]{Sta}.
\end{proof}
% Lean: stalk_of_smooth: exists_isStandardSmoothOfRelativeDimension and localization_of_isStandardSmooth (with
% height_of_etale_mvPolynomial), about 700 lines. The text cites Stacks for these facts, which the Lean proves
% itself through a standard smooth presentation near x.

\begin{definition}\label{cv-toroidal-abelian-dvr-flat-complex-sections}
Let $A$ be a ring, let $X$ be a scheme over $A$, and let $\mathcal M$ be a line bundle on $X$. For an $A$-algebra $R$
let $X_R=X\times_A\operatorname{Spec}R$, and write $\mathcal M_R$, $W_R$, $\epsilon_R$ and $u_R$ for the pullbacks to
$X_R$ of $\mathcal M$, of an open subscheme $W$ of $X$, of a section $\epsilon$ of $\mathcal M$ and of a function $u$.
A \emph{section complex} for $\mathcal M$ consists of integers $a,b,c\ge0$, matrices $D_0\in M_{b\times a}(A)$ and
$D_1\in M_{c\times b}(A)$ with $D_1D_0=0$, an open cover $(W_i)_{i\in J}$ of $X$ with generators $\epsilon_i$ of
$\mathcal M$ over $W_i$, and functions $u_{k,i}\in\Gamma(W_i,\mathcal O_X)$ for $1\le k\le a$ and $i\in J$, such that
for every $A$-algebra $R$:
\begin{enumerate}
\item for every $v\in R^a$ with $D_0v=0$ there is a section $s_v$ of $\mathcal M_R$ over $X_R$ with
  $s_v=\ell_i(v)\epsilon_{i,R}$ on $W_{i,R}$ for every $i$, where $\ell_i(v)=\sum_kv_ku_{k,i,R}$, and $v\mapsto s_v$ is
  a bijection from the kernel of $D_0$ on $R^a$ to the set of sections of $\mathcal M_R$ over $X_R$;
\item if every \v{C}ech $1$-cocycle of $\mathcal M_R$ for an open cover of $X_R$ is a coboundary, then every
  $w\in R^b$ with $D_1w=0$ is $D_0v$ for some $v\in R^a$.
\end{enumerate}
Here a \v{C}ech $1$-cocycle of $\mathcal M_R$ for an open cover $(V_\lambda)_{\lambda\in\Lambda}$ of $X_R$ is a family of
sections $\theta_{\lambda\mu}$ of $\mathcal M_R$ over $V_\lambda\cap V_\mu$ with $\theta_{\lambda\mu}+\theta_{\mu\nu}=
\theta_{\lambda\nu}$ on $V_\lambda\cap V_\mu\cap V_\nu$ for all $\lambda,\mu,\nu\in\Lambda$, and it is a coboundary if
there are sections $\psi_\lambda$ of $\mathcal M_R$ over $V_\lambda$ with $\theta_{\lambda\mu}=\psi_\lambda-\psi_\mu$ on
$V_\lambda\cap V_\mu$ for all $\lambda,\mu$.
\end{definition}
\leanref{Development/AbelianVarietyTheory/Definitions/AbelianSchemeLaws/FirstCohomologyVanishesSpecOfIsLocalRing.lean}{GlobalTheory.AbLaw.Perf}
% Lean: the hypothesis of exact is FirstCohomologyVanishes (M.pull pi), the condition of (2) for the open covers
% subordinate to the charts of M.pull pi, with sections written in the charts; CechOne.h1Zero_iff shows that it is
% equivalent to CechOne.CechHOneVanishesOnOpen (M.pull pi) top, the condition of (2) for every open cover (IsCocycle,
% IsCoboundary).

\begin{lemma}[Mumford {\cite[\S5]{Mum70}}]\label{cv-toroidal-abelian-dvr-flat-complex-model}
Let $A$ be a ring, let $X$ be a scheme smooth over $A$, and let $\mathcal H$ be a line bundle on $X$ with sections
$s_0,\dots,s_m$ that define a closed immersion of $X$ into $\mathbf P^m_A$, so that the nonvanishing loci $X_i$ of the
$s_i$ cover $X$ and all their finite intersections $X_I=\bigcap_{i\in I}X_i$ are affine. Let $\mathcal M$ be a line
bundle on $X$ and let $\check C$ be the \v{C}ech complex of $\mathcal M$ for the cover $(X_i)$, with
$\check C^p=\prod_{|I|=p+1}\Gamma(X_I,\mathcal M)$. Then there are a complex $K$ of finite projective $A$-modules, zero in
negative degrees, and a map of complexes $\varphi:K\to\check C$ such that $\varphi\otimes_AR$ is a quasi-isomorphism for
every $A$-algebra $R$.
\end{lemma}
\leanref{Development/AbelianVarietyTheory/Proofs/CoherentCohomologyProofs/ExistsFiniteProjectiveCechComplexLineBundleStatement.lean}{GlobalTheory.T.exists_finite_projective_cechComplex_lineBundle}
\leanref{Development/AbelianVarietyTheory/Proofs/CoherentCohomologyProofs/ExistsFiniteProjectiveCechComplexLineBundle.lean}{GlobalTheory.Cohom.exists_finite_projective_complex_cechLB}
% Lean: Cohom.exists_finite_projective_complex_cechLB: exists_noetherian_model_grCube (a model of the graded Cech cube
% over a noetherian subring R_0 of A, through exists_dehom_chart and secOn_finitePresentation_projective),
% exists_grCube_twistedSections_hom_secCube_bijective, secCubeIso and exists_finite_projective_complex_of_noetherian_model.
% Mumford states the result for a noetherian base; the Lean proves it for any ring A through the noetherian model.
% Size: about 3,600 lines.

\begin{lemma}[Alp\"oge]\label{cv-toroidal-abelian-dvr-flat-complex-matrices}
Let $A$ be a ring and let $P_0\xrightarrow{d_0}P_1\xrightarrow{d_1}P_2$ be homomorphisms of finite projective $A$-modules
with $d_1d_0=0$. There are integers $a,b,c\ge0$, matrices $D_0\in M_{b\times a}(A)$ and $D_1\in M_{c\times b}(A)$ with
$D_1D_0=0$, and elements $u_1,\dots,u_a\in P_0$ such that for every $A$-algebra $R$:
\begin{enumerate}
\item $v\mapsto\sum_kv_k\otimes u_k$ is a bijection from the kernel of $D_0$ on $R^a$ to the kernel of $d_0\otimes R$;
\item if $P_0\otimes_AR\to P_1\otimes_AR\to P_2\otimes_AR$ is exact, then every $w\in R^b$ with $D_1w=0$ is $D_0v$ for
  some $v\in R^a$.
\end{enumerate}
\end{lemma}
\leanref{Development/AbelianVarietyTheory/Proofs/CoherentCohomologyProofs/ExistsMatricesOfProjectiveComplex.lean}{GlobalTheory.exists_matrices_of_projective_complex}
\leanref{Development/AbelianVarietyTheory/Proofs/CoherentCohomologyProofs/ExistsMatricesOfProjectiveComplex.lean}{GlobalTheory.projMatrix_abstract}
\begin{proof}
Choose $F_i=A^{n_i}$ with $p_i:F_i\to P_i$ and $s_i:P_i\to F_i$ such that $p_is_i=\mathrm{id}$. Let $a=n_0$,
$b=n_1+n_0$, $c=n_2+n_1+n_0$, let $D_0:F_0\to F_1\oplus F_0$ and $D_1:F_1\oplus F_0\to F_2\oplus F_1\oplus F_0$ be
\[
 D_0(x)=(s_1d_0p_0x,\;x-s_0p_0x),\qquad D_1(y,x)=(s_2d_1p_1y,\;y-s_1p_1y,\;s_0p_0x),
\]
and let $u_k=p_0(e_k)$ for the standard basis $e_k$ of $F_0$, so that $p_0(v)=\sum_kv_ku_k$. Then $D_1D_0=0$, as $d_1d_0=0$
and $p_is_i=\mathrm{id}$. Let $R$ be an $A$-algebra; all maps below are tensored with $R$.

(1) If $D_0v=0$, then $v=s_0p_0v$ and $s_1d_0p_0v=0$, so $d_0p_0v=0$ (apply $p_1$); and $p_0$ is injective on such
$v$, as $v=s_0p_0v$. If $d_0x=0$, then $D_0(s_0x)=(s_1d_0x,s_0x-s_0p_0s_0x)=0$ and $p_0s_0x=x$.

(2) Let $D_1(y,x)=0$. Then $d_1p_1y=0$ (apply $p_2$), $y=s_1p_1y$, and $p_0x=0$ (apply $p_0$). By exactness
$p_1y=d_0z$ for some $z$. Let $v=s_0z+x$; then $p_0v=z$, and $D_0v=(s_1d_0z,\,s_0z+x-s_0z)=(s_1p_1y,x)=(y,x)$.
\end{proof}
% Lean: exists_matrices_of_projective_complex (Module.Finite.exists_comp_eq_id_of_projective for p_i, s_i; f_0 = D_0
% and f_1 = D_1 through the equivalences E and E'), projMatrix_abstract (the argument over R), projMatrix_mulVec and
% projMatrix_sum (the matrices and the sum over the basis after base change).

\begin{lemma}[Alp\"oge]\label{cv-toroidal-abelian-dvr-flat-complex}
Let $A$ be a ring, let $X$ be a scheme smooth and projective over $A$, and let $\mathcal M$ be a line bundle on
$X$. Then $\mathcal M$ has a section complex $(D_0,D_1,(W_i,\epsilon_i),(u_{k,i}))$
(Definition~\ref{cv-toroidal-abelian-dvr-flat-complex-sections}).
\end{lemma}
\leanref{Development/AbelianVarietyTheory/Proofs/AbelianSchemeLawProofs/NonemptyPerfectSectionsComplexOfSmoothProjStatement.lean}{GlobalTheory.T.nonempty_perfectSectionsComplex_of_smooth_proj}
\leanref{Development/AbelianVarietyTheory/Proofs/AbelianSchemeLawProofs/NonemptyPerfectSectionsComplexOfSmoothProj.lean}{GlobalTheory.AbLaw.nonempty_perf_of_smooth_isProjBy}
\begin{proof}
As $X$ is projective over $A$, there are a line bundle $\mathcal H$ on $X$ and sections $s_0,\dots,s_m$ as in
Lemma~\ref{cv-toroidal-abelian-dvr-flat-complex-model}, with nonvanishing loci $X_i$ whose finite intersections $X_I$
are affine. Let $\check C$ be the \v{C}ech complex of $\mathcal M$ for this cover, and let $\varphi:K\to\check C$ be
given by Lemma~\ref{cv-toroidal-abelian-dvr-flat-complex-model}. Lemma~\ref{cv-toroidal-abelian-dvr-flat-complex-matrices},
applied to $K^0\to K^1\to K^2$, gives $a,b,c$, $D_0$, $D_1$ and $u_1,\dots,u_a\in K^0$. Let $(U_j,\epsilon_j)$ be open
subschemes covering $X$ with generators $\epsilon_j$ of $\mathcal M$ over $U_j$. The open cover is
$W_{(i,j)}=X_i\cap U_j$, with the generators $\epsilon_j$, and $u_{k,(i,j)}$ is the function on $X_i\cap U_j$ with
$\varphi^0(u_k)_i=u_{k,(i,j)}\epsilon_j$ there, where $\varphi^0(u_k)_i$ is the component of $\varphi^0(u_k)$ in
$\Gamma(X_i,\mathcal M)$.

Let $R$ be an $A$-algebra. As the $X_I$ are affine, $\Gamma(X_{I,R},\mathcal M_R)=\Gamma(X_I,\mathcal M)\otimes_AR$, so
$\check C\otimes_AR$ is the \v{C}ech complex of $\mathcal M_R$ for the cover $(X_{i,R})$.

(1) The kernel of $d_0\otimes R$ on $K^0\otimes_AR$ maps bijectively, through $\varphi\otimes R$, onto the kernel of
$\check C^0\otimes_AR\to\check C^1\otimes_AR$, as $\varphi\otimes R$ is a quasi-isomorphism; and that kernel is the set of
sections of $\mathcal M_R$ over $X_R$, by the sheaf property. Composed with
Lemma~\ref{cv-toroidal-abelian-dvr-flat-complex-matrices}(1), this is a bijection $v\mapsto s_v$ from the kernel of $D_0$
on $R^a$ to the sections of $\mathcal M_R$, with $s_v$ the section whose component on $X_{i,R}$ is
$\sum_kv_k\varphi^0(u_k)_i$; on $W_{(i,j),R}$ this is $\bigl(\sum_kv_ku_{k,(i,j),R}\bigr)\epsilon_{j,R}$.

(2) Suppose that every \v{C}ech $1$-cocycle of $\mathcal M_R$ for an open cover of $X_R$ is a coboundary. Order the
indices of the cover $(X_i)$. The \v{C}ech complex $\check C\otimes_AR$ of $\mathcal M_R$ for the cover $(X_{i,R})$ is
exact in degree $1$. Indeed, let $x\in\check C^1\otimes_AR$, with components $x_{\alpha\beta}$ over
$X_{\alpha,R}\cap X_{\beta,R}$ for $\alpha<\beta$, be killed by the differential, that is,
$x_{\beta\gamma}-x_{\alpha\gamma}+x_{\alpha\beta}=0$ for $\alpha<\beta<\gamma$. Let $\theta_{\alpha\beta}=x_{\alpha\beta}$
and $\theta_{\beta\alpha}=-x_{\alpha\beta}$ for $\alpha<\beta$, and $\theta_{\alpha\alpha}=0$. By these relations
$(\theta_{\alpha\beta})$ is a \v{C}ech $1$-cocycle (Definition~\ref{cv-toroidal-abelian-dvr-flat-complex-sections}), so
there are sections $\psi_\alpha$ of $\mathcal M_R$ over $X_{\alpha,R}$ with $\theta_{\alpha\beta}=\psi_\alpha-\psi_\beta$.
The differential maps $(-\psi_\alpha)_\alpha\in\check C^0\otimes_AR$ to the family with components
$-\psi_\beta+\psi_\alpha=x_{\alpha\beta}$ for $\alpha<\beta$, that is, to $x$. So $K\otimes_AR$ is exact in degree $1$,
as $\varphi\otimes R$ is a quasi-isomorphism, and Lemma~\ref{cv-toroidal-abelian-dvr-flat-complex-matrices}(2) gives the
claim.
\end{proof}
% Lean: AbLaw.nonempty_perf_of_smooth_isProjBy: IsProjBy (H, s), T.exists_finite_projective_cechComplex_lineBundle
% (Lemma cv-toroidal-abelian-dvr-flat-complex-model), T.veryAmple_compact_quasiSeparated_affine_charts (the X_I are
% affine) and T.threeTermComplex_of_projective_cech_model = AbLaw.perf_of_projective_model:
% exists_matrices_of_projective_complex, W = D_i inf M.U j, u = SecOnR.coord of phi^0(u_k); sec through
% PerfOfModel.termMap_bijOn (secTermBaseChange_bijective), bijOn_ker_and_exact_of_quasiIso_baseChange and
% CechOne.secKerEquiv, with coord_termMap_sum for the chart formula; exact through Cohom.exact_cech_one_of_h1Zero for
% the cover pi^-1 D_i of X_R (from FirstCohomologyVanishes (M.pull pi), through CechOne.h1Zero_iff, which passes to
% every open cover, and CechOne.exactAt_of_isCoboundary, the argument of (2): phi from x with add_eq_of_ordered for
% the cocycle condition, and y = -psi through restrict_d_one), PerfOfModel.exact_baseChange and
% bijOn_ker_and_exact_of_quasiIso_baseChange. Size: about 5,500 lines, of which about 3,600 are the model.

\begin{lemma}[Alp\"oge]\label{cv-toroidal-abelian-dvr-flat-cyclic-minor}
Let $R$ be a ring, let $g_0:R\to K$ be a homomorphism to a field, let $D\in M_{b\times a}(R)$, and let $v\in K^a$ be
nonzero such that every $w\in K^a$ with $g_0(D)w=0$ is a multiple of $v$. Then there are $\delta\in R$ with
$g_0(\delta)\ne0$ and $v_0\in R^a$ such that for every homomorphism $g:R\to R'$ with $g(\delta)$ a unit, every
$w\in R'^a$ with $g(D)w=0$ is $c\,g(v_0)$ for some $c\in R'$.
\end{lemma}
\leanref{Development/AbelianVarietyTheory/Proofs/AbelianSchemeLawProofs/ExistsKernelVector.lean}{GlobalTheory.AbLaw.Seesaw.exists_kernel_vector}
\leanref{Development/AbelianVarietyTheory/Proofs/AbelianSchemeLawProofs/ExistsKernelVector.lean}{GlobalTheory.AbLaw.Seesaw.exists_minor}
\begin{proof}
First, if a matrix $N\in M_{m\times n}(K)$ is injective on $K^n$, then for some map $\rho:\{1,\dots,n\}\to\{1,\dots,m\}$
the square matrix $N_\rho$ with rows $N_{\rho(1)},\dots,N_{\rho(n)}$ has nonzero determinant. Indeed, $N$ has a left
inverse $L$, and as the $i$-th row of $LN$ is $\sum_jL_{ij}N_j$, multilinearity of the determinant in the rows gives
$1=\det(LN)=\sum_\rho\bigl(\prod_iL_{i\rho(i)}\bigr)\det N_\rho$.

Choose $k$ with $v_k\ne0$, and let $D'$ be the matrix $D$ with the row $e_k^{\mathsf T}$ added, $e_k$ being the $k$-th
standard basis vector. If $g(D)w=0$, then $g(D')w$ is $w_k$ in the added row and $0$ elsewhere. The matrix $g_0(D')$
is injective: if $g_0(D')w=0$, then $w=cv$ and $cv_k=w_k=0$, so $c=0$. By the first paragraph there are rows
$\rho(1),\dots,\rho(a)$ of $D'$ whose matrix $N$ has $g_0(\det N)\ne0$. Let $\delta=\det N$, let $e\in R^a$ have
$e_j=1$ if $\rho(j)$ is the added row and $e_j=0$ otherwise, and let $v_0=\operatorname{adj}(N)e$. If $g(\delta)$ is a
unit and $g(D)w=0$, then $g(N)w=w_kg(e)$, so $g(\delta)w=\operatorname{adj}(g(N))g(N)w=w_kg(v_0)$, and
$w=g(\delta)^{-1}w_k\,g(v_0)$.
\end{proof}
% Lean: exists_kernel_vector (D' with the row Pi.single k 1, key) and exists_minor (left inverse, Cauchy-Binet
% expansion through detRowAlternating).

\begin{lemma}[Alp\"oge]\label{cv-toroidal-abelian-dvr-flat-cyclic-coker}
Let $R$ be a ring, let $D\in M_{b\times a}(R)$, and let $v_0\in R^a$ have a coordinate $v_{0,k}$ that is a unit.
Suppose that for every $R$-algebra $R'$, every $w\in R'^a$ with $Dw=0$ is a multiple of the image of $v_0$. Then the
cokernel of $D^{\mathsf T}:R^b\to R^a$ is isomorphic to $R/I$, where $I$ is the ideal generated by the entries of
$Dv_0$.
\end{lemma}
\leanref{Development/AbelianVarietyTheory/Proofs/UniformizationProofs/CyclicIdeal.lean}{GlobalTheory.UU.SecCyclic.coker_equiv}
\begin{proof}
Let $N=R^a/D^{\mathsf T}(R^b)$, and let $e'_j$ be the class of the $j$-th standard basis vector. For each $i$,
$\sum_jD_{ij}e'_j$ is the class of $D^{\mathsf T}e_i$, which is $0$. Let $R'=R\oplus N$ with $N^2=0$, and let
$w\in R'^a$ have coordinates $e'_j\in N\subset R'$. Then $(Dw)_i=\sum_jD_{ij}e'_j=0$, so $w=c\,v_0$ for some
$c=(c_1,n)\in R'$, and comparing components in $N$ gives $e'_j=v_{0,j}n$. Hence every $x\in R^a$ has class
$(x\cdot v_0)n$ in $N$. Moreover $In=0$: $(Dv_0)_in=\sum_jD_{ij}v_{0,j}n=\sum_jD_{ij}e'_j=0$.

Let $\varphi:R^a\to R/I$, $x\mapsto x\cdot v_0+I$. It is surjective, as $\varphi(rv_{0,k}^{-1}e_k)=r+I$. Its kernel is
$D^{\mathsf T}(R^b)$: on the one hand $D^{\mathsf T}y\cdot v_0=y\cdot Dv_0\in I$; on the other, if $x\cdot v_0\in I$,
then the class of $x$ is $(x\cdot v_0)n=0$. So $\varphi$ induces $N\cong R/I$.
\end{proof}
% Lean: coker_equiv, with R' = TrivSqZeroExt R N.

\begin{lemma}[Alp\"oge]\label{cv-toroidal-abelian-dvr-flat-cyclic}
Let $R$ be a local ring, let $f:X\to\operatorname{Spec}R$ be universally closed and surjective, and let $\mathcal M$
be a line bundle on $X$ with a section complex $(D_0,D_1,(W_i,\epsilon_i),(u_{k,i}))$
(Definition~\ref{cv-toroidal-abelian-dvr-flat-complex-sections}). Let $R\to K_0$ be a local homomorphism to a field
such that the functions on $X_{K_0}$ are the constants $K_0$ and $\mathcal M_{K_0}$ is trivial. Then there is an ideal
$I$ of $R$ such that the cokernel of the transpose $D_0^{\mathsf T}:R^b\to R^a$ is isomorphic to $R/I$ and, for every
$R$-algebra $R'$, the line bundle $\mathcal M_{R'}$ on $X_{R'}$ is trivial if and only if $I\subset\ker(R\to R')$.
\end{lemma}
\leanref{Development/AbelianVarietyTheory/Proofs/UniformizationProofs/CyclicIdeal.lean}{GlobalTheory.UU.SecCyclic.cyclic_ideal}
\begin{proof}
For an $R$-algebra $R'$ and $v\in R'^a$ let $\ell_i(v)=\sum_kv_ku_{k,i,R'}$ on $W_{i,R'}$, as in
Definition~\ref{cv-toroidal-abelian-dvr-flat-complex-sections}. Then $\ell_i(cv)=c\ell_i(v)$, and for $v\in R^a$,
$\ell_i$ of the image of $v$ is the pullback of $\ell_i(v)$. Three facts on these functions come first.

(a) If $\mathcal M_{R'}$ is trivial, there is $v\in R'^a$ with $D_0v=0$ such that every $\ell_i(v)$ is a unit; if
moreover every function on $X_{R'}$ is a constant from $R'$, every $w\in R'^a$ with $D_0w=0$ is a multiple of $v$.
Indeed, a generator of $\mathcal M_{R'}$ is $s_v$ for some $v$, and its functions $\ell_i(v)$ are units. For $w$, the
section $s_w$ is $\rho s_v$ for a function $\rho$ on $X_{R'}$, which is a constant $c\in R'$; so $s_w=s_{cv}$, and
$w=cv$ by the bijectivity of $v\mapsto s_v$.

(b) If $c\in R'$ and $w\in R'^a$ are such that every $\ell_i(cw)$ is a unit, then $c$ is a unit. Otherwise $c$ lies
in a maximal ideal $\mathfrak n$ of $R'$. As $f$ is surjective, so is its base change $X_{R'}\to\operatorname{Spec}R'$,
and some $y\in X_{R'}$ lies over $\mathfrak n$, say $y\in W_{i,R'}$. Then $\ell_i(cw)=c\ell_i(w)$ does not vanish at
$y$, so $c\notin\mathfrak n$.

(c) If $v\in R'^a$ has $D_0v=0$ and every point of $X_{R'}$ lies in some $W_{i,R'}$ at which $\ell_i(v)$ does not
vanish, then $s_v$ vanishes nowhere, and $\mathcal M_{R'}$ is trivial.

By (a) for $K_0$, there is $v\in K_0^a$ with $D_0v=0$ and every $\ell_i(v)$ a unit, of which every vector in the kernel
of $D_0$ on $K_0^a$ is a multiple. Then $v\ne0$, since otherwise (b) for $c=0$ would make $0$ a unit of $K_0$.
Lemma~\ref{cv-toroidal-abelian-dvr-flat-cyclic-minor} gives $\delta\in R$ and $v_0\in R^a$. As $R\to K_0$ is local and
the image of $\delta$ is nonzero, $\delta$ is a unit of $R$. So for every $R$-algebra $R'$, every $w\in R'^a$ with
$D_0w=0$ is a multiple of the image $v'_0$ of $v_0$. In particular $v$ is $c_0$ times the image of $v_0$ for some
$c_0\in K_0$, and for $k$ with $v_k\ne0$ the image of $v_{0,k}$ is nonzero, so $v_{0,k}$ is a unit of $R$.
Lemma~\ref{cv-toroidal-abelian-dvr-flat-cyclic-coker} shows that $\operatorname{coker}D_0^{\mathsf T}\cong R/I$, with
$I$ the ideal generated by the entries of $D_0v_0$. For an $R$-algebra $R'$, $I\subset\ker(R\to R')$ if and only if
$D_0v'_0=0$.

For each $i$ let $G_i$ be the open set of points of $W_i$ at which $\ell_i(v_0)$ does not vanish, and let $G$ be the
union of the $G_i$. As $f$ is closed, $Z=f(X\setminus G)$ is closed. The image $p_0$ of $\operatorname{Spec}K_0$ is not in $Z$: if $x\notin G$ and
$f(x)=p_0$, some $y\in X_{K_0}$ lies over $x$, say in $W_{i,K_0}$, and $\ell_i(v)$, which is $c_0$ times the pullback of
$\ell_i(v_0)$, is a unit at $y$; so $\ell_i(v_0)$ does not vanish at $x$, and $x\in G$. So there is $r_1\in R$ with
$p_0\in D(r_1)$ and $D(r_1)\cap Z=\emptyset$. The image of $r_1$ in $K_0$ is nonzero, so $r_1$ is a unit, $D(r_1)$ is
all of $\operatorname{Spec}R$, and $G=X$.

Let $R'$ be an $R$-algebra. If $\mathcal M_{R'}$ is trivial, (a) gives $w$ with $D_0w=0$ and every $\ell_i(w)$ a unit;
then $w=cv'_0$, and $c$ is a unit by (b), so $D_0v'_0=c^{-1}D_0w=0$. Conversely, let $D_0v'_0=0$. Every $y\in X_{R'}$
maps to a point of $G$, which lies in some $W_i$ at which $\ell_i(v_0)$ does not vanish; so $\ell_i(v'_0)$, its pullback,
does not vanish at $y$. By (c), $\mathcal M_{R'}$ is trivial.
\end{proof}
% Lean: cyclic_ideal: loc_of_triv (a, with sec_ratio), unit_cst (b), triv_of_loc (c), exists_kernel_vector (Lemma
% cv-toroidal-abelian-dvr-flat-cyclic-minor), coker_equiv (Lemma cv-toroidal-abelian-dvr-flat-cyclic-coker), the open
% G = union of the basic opens of loc_0 P v_0 i, f.isClosedMap and isTopologicalBasis_basic_opens for r_1. Size: about
% 350 lines with the helpers of ExistsKernelVector.lean.

\begin{lemma}[Mumford {\cite[\S5]{Mum70}}]\label{cv-toroidal-abelian-dvr-flat-free-model}
Let $A$ be a noetherian ring and let $\check C$ be a complex of flat $A$-modules, zero in negative degrees and in
degrees larger than $n$, whose cohomology modules are finite over $A$. Then there are a complex $C$ of finite projective
$A$-modules, zero in negative degrees and in degrees larger than $n$, and a map of complexes $\psi:C\to\check C$ such
that $\psi\otimes_AR$ is a quasi-isomorphism for every $A$-algebra $R$.
\end{lemma}
\leanref{Development/AbelianVarietyTheory/Proofs/CoherentCohomologyProofs/ExistsFiniteProjectiveComplexQuasiIso.lean}{GlobalTheory.exists_finite_projective_complex_quasiIso}
% Lean: exists_finite_projective_complex_quasiIso, which also gives C free in positive degrees and psi tensored with
% every module a quasi-isomorphism. Mumford's construction (Section 5, Lemma 1 and the second theorem) is the one cited.

\begin{lemma}[Alp\"oge]\label{cv-toroidal-abelian-dvr-flat-free-yoneda}
Let $A$ be a ring and let $D\in M_{b\times a}(A)$ and $D'\in M_{b'\times a'}(A)$. Suppose that for every $A$-algebra $R$
there is a bijection $e_R$ from the kernel of $D$ on $R^a$ to the kernel of $D'$ on $R^{a'}$, with $e_A(0)=0$, such that
$e_S(h\circ v)=h\circ e_R(v)$ for every homomorphism $h:R\to S$ of $A$-algebras. Then the cokernels of
$D^{\mathsf T}:A^b\to A^a$ and of $D'^{\mathsf T}:A^{b'}\to A^{a'}$ are isomorphic.
\end{lemma}
\leanref{Development/AbelianVarietyTheory/Proofs/UniformizationProofs/ExistsPerfectFreeFiniteComplex.lean}{GlobalTheory.KernelYoneda.nonempty_equiv}
\begin{proof}
Let $Q=\operatorname{coker}D^{\mathsf T}$ and $Q'=\operatorname{coker}D'^{\mathsf T}$, with the classes $[e_k]$ and
$[e'_l]$ of the standard basis vectors. For an $A$-module $N$ let $A\oplus N$ be the ring with $N^2=0$. A vector
$(n_k)\in N^a\subset(A\oplus N)^a$ lies in the kernel of $D$ if and only if $\sum_kD_{jk}n_k=0$ for all $j$, that is, if
and only if there is a homomorphism $Q\to N$ with $[e_k]\mapsto n_k$; and similarly for $D'$.

Apply $e_{A\oplus Q}$ to the vector $([e_k])_k\in Q^a$, and let $w$ be its image. The projection $A\oplus Q\to A$ maps
$([e_k])$ to $0$, so it maps $w$ to $e_A(0)=0$; hence $w\in Q^{a'}$, and $w$ defines a homomorphism $\alpha:Q'\to Q$
with $\alpha([e'_l])=w_l$. In the same way the inverses of the $e_R$ give $\beta:Q\to Q'$ with
$e_{A\oplus Q'}((\beta[e_k])_k)=([e'_l])_l$. The homomorphism $A\oplus Q'\to A\oplus Q$ induced by $\alpha$ carries this
identity, by the compatibility with base change, to $e_{A\oplus Q}((\alpha\beta[e_k])_k)=(\alpha[e'_l])_l=w=
e_{A\oplus Q}(([e_k])_k)$. As $e_{A\oplus Q}$ is injective, $\alpha\beta[e_k]=[e_k]$ for all $k$, and
$\alpha\beta=\mathrm{id}$. Likewise $\beta\alpha=\mathrm{id}$.
\end{proof}
% Lean: KernelYoneda.nonempty_equiv (with a relation Rel in place of the bijections e_R): exists_hom_of_rel (alpha and
% beta, through mulVec_inr_iff and exists_hom on TrivSqZeroExt A (Coker D)) and comp_eq_id (through TrivSqZeroExt.map).

\begin{lemma}[Alp\"oge]\label{cv-toroidal-abelian-dvr-flat-free}
Let $A$ be a noetherian local ring and $B$ a commutative group scheme over $A$ that is an abelian scheme,
projective and smooth of relative dimension $g$ over $A$. Let $\mathcal F$ be a line bundle on $B$ with
$\Lambda(\mathcal F)$ trivial on $B\times_AB$, and let $(D_0,D_1,(W_i,\epsilon_i),(u_{k,i}))$ be a section complex for
$\mathcal F$ (Definition~\ref{cv-toroidal-abelian-dvr-flat-complex-sections}). There is a complex $C^0\to C^1\to\cdots$ of finite free
$A$-modules, zero in large degrees, such that the cokernel of the dual $(C^1)^\vee\to(C^0)^\vee$ of its first
differential is isomorphic to the cokernel of $D_0^{\mathsf T}:A^b\to A^a$, and such that for every algebraically
closed field $K$ and homomorphism $A\to K$, the complex $C\otimes_AK$ is exact in degree $i$ whenever $i>g$ or the
pullback of $\mathcal F$ to $B_K$ is nontrivial.
\end{lemma}
\leanref{Development/AbelianVarietyTheory/Proofs/UniformizationProofs/ExistsPerfectFreeFiniteComplexStatement.lean}{GlobalTheory.T.exists_perfect_free_finite_complex}
\begin{proof}
As $B$ is projective over $A$, there are a line bundle $\mathcal H$ on $B$ and sections $s_0,\dots,s_m$ as in
Lemma~\ref{cv-toroidal-abelian-dvr-flat-complex-model}, with nonvanishing loci $X_i$ whose finite intersections $X_I$
are affine. Let $\check C$ be the \v{C}ech complex of $\mathcal F$ for the cover $(X_i)$. It is zero in degrees larger
than $m$; its terms are flat over $A$, as $B$ is flat over $A$ and the $X_I$ are affine; and its cohomology modules
$H^p(B,\mathcal F)$ \cite[Tag 01XD]{Sta} are finite over $A$, as $B$ is proper over the noetherian ring $A$
\cite[Tag 02O5]{Sta}. Lemma~\ref{cv-toroidal-abelian-dvr-flat-free-model} gives a complex $C$ of finite projective
$A$-modules, zero in large degrees, and $\psi:C\to\check C$ with $\psi\otimes R$ a quasi-isomorphism for every
$A$-algebra $R$. As $A$ is local, the $C^p$ are free \cite[Theorem~2.5]{Mat86}.

The cokernel. Choose bases of $C^0$ and $C^1$, of sizes $a'$ and $b'$, and let $M_0\in M_{b'\times a'}(A)$ be the
matrix of the first differential $d_0$. With the dual bases, the cokernel of $(C^1)^\vee\to(C^0)^\vee$ is that of
$M_0^{\mathsf T}$. By the argument of part (1) of the proof of Lemma~\ref{cv-toroidal-abelian-dvr-flat-complex}, applied
to $\psi$ and the bases, for every $A$-algebra $R$ the kernel of $M_0$ on $R^{a'}$ maps bijectively to the sections of $\mathcal F_R$, compatibly with base change, a vector $v$
going to the section whose functions on the charts $X_i\cap U_j$ are given by $\psi^0$ of the basis of $C^0$ and the
coordinates of $v$. The section complex gives another such bijection from the kernel of $D_0$, with the functions
$\sum_kv_ku_{k,i}$ on the $W_i$. On the common refinement of the two covers by the sets $W_i\cap X_{i'}\cap U_j$, a
section is determined by its functions there, so composing the first bijection with the inverse of the second gives
bijections $e_R$ as in Lemma~\ref{cv-toroidal-abelian-dvr-flat-free-yoneda}, compatible with base change and with
$e_A(0)=0$. That lemma shows that the cokernels of $D_0^{\mathsf T}$ and of $M_0^{\mathsf T}$ are isomorphic.

The fibres. Let $K$ be an algebraically closed field with a homomorphism $A\to K$, and let $i$ be such that $i>g$ or
$\mathcal F_K$ is nontrivial. The complex $\check C\otimes_AK$ is the \v{C}ech complex of $\mathcal F_K$ for the affine
cover $(X_{i,K})$ of $B_K$, and $\psi\otimes K$ is a quasi-isomorphism, so it suffices that $H^i(B_K,\mathcal F_K)=0$
\cite[Tag 01XD]{Sta}. If $i>g$, this holds as $B_K$ has dimension $g$ \cite[Tag 02UZ]{Sta}. If $\mathcal F_K$ is
nontrivial, then, as $\Lambda(\mathcal F_K)$ is trivial, $\mathcal F_K$ lies in $\operatorname{Pic}^0(B_K)$, and all the
cohomology of a nontrivial line bundle in $\operatorname{Pic}^0$ of an abelian variety vanishes \cite[\S8]{Mum70}.
\end{proof}
% Lean: AbLaw.exists_free_complex_of_primitive: IsProjBy (H, s), VeryAmpleCharts.compact_quasiSeparated_isAffineOpen,
% exists_finite_projective_complex_quasiIso (Lemma cv-toroidal-abelian-dvr-flat-free-model) with isZero_cechLB_X,
% flat_cechLB_X and cechLB_homology_finite_of_affine_cover (finiteness of the cohomology, proved in the Lean through
% T.finite_cechLineBundle_homology), Module.free_of_flat_of_isLocalRing; the cokernel through
% SecKernel.nonempty_equiv_of_model (KernelYoneda.nonempty_dual_equiv, SecKernel.computes_of_model with bijOn_basis,
% computes_perf, Computes.refine and Computes.nonempty_equiv, which applies KernelYoneda.nonempty_equiv); the fibres
% through cechLBBaseChangeIso, cechLB_exactAt_of_dim_lt (a cover of the fibre by g + 1 affine opens,
% AffineCover.exists_cover_of_isProjBy, and the independence of the affine cover; the text cites Stacks Tag 02UZ) and
% T.cechComplex_acyclic_primitive_nontrivial_algClosed = PrimVan.acyclic_cechLB_of_primitive (the text cites Mumford,
% Section 8). The Lean proves these cited results itself. Size: about 7,000 lines.

\begin{lemma}[Alp\"oge]\label{cv-toroidal-abelian-dvr-flat-regular-split}
Let $A$ be a local ring with residue field $\kappa$, let $M$ be an $A$-module, let $P$ and $P'$ be finite free
$A$-modules, and let $f:M\to P$ and $h:P\to P'$ be homomorphisms with $hf=0$ such that
$M\otimes_A\kappa\to P\otimes_A\kappa\to P'\otimes_A\kappa$ is exact. Then there are homomorphisms $s:P\to M$ and
$t:P'\to P$ with $fs+th=\mathrm{id}_P$.
\end{lemma}
\leanref{Development/AbelianVarietyTheory/Proofs/CoherentCohomologyProofs/ExactAndKerBaseChangeOfForallResidueField.lean}{GlobalTheory.exists_homotopy_of_forall_exact_residueField}
\begin{proof}
First, a homomorphism $u:N\to F$ from a finite $A$-module to a finite free one for which $u\otimes_A\kappa$ is
injective has a left inverse. Indeed, choose $n_1,\dots,n_r\in N$ whose images form a basis of $N\otimes_A\kappa$,
and $f_{r+1},\dots,f_m\in F$ whose images, with those of the $u(n_i)$, form a basis of $F\otimes_A\kappa$, where $m$
is the rank of $F$. By Nakayama's lemma \cite[Theorem~2.2]{Mat86} the $n_i$ generate $N$, and the $u(n_i)$ and $f_j$
generate $F$, so they form a basis of $F$ \cite[Theorem~2.4]{Mat86}. The homomorphism $F\to N$ with
$u(n_i)\mapsto n_i$ and $f_j\mapsto0$ is a left inverse of $u$.

The homomorphism $h$ vanishes on $f(M)$ and induces $\bar h:P/f(M)\to P'$. Its reduction
$\operatorname{coker}(f\otimes_A\kappa)\to P'\otimes_A\kappa$ is injective, by exactness, and $P/f(M)$ is finite;
so $\bar h$ has a left inverse $r$. As $P'$ is free, $r$ lifts to $t:P'\to P$, with $t(y)$ in the class $r(y)$. For
$x\in P$, the class of $x-th(x)$ in $P/f(M)$ is $[x]-r\bar h[x]=0$, so $x-th(x)\in f(M)$. As $P$ is free, the
homomorphism $\mathrm{id}-th:P\to f(M)$ lifts along $f$ to $s:P\to M$, and $fs=\mathrm{id}-th$.
\end{proof}
% Lean: exists_homotopy_of_forall_exact_residueField (any ring, B and C finite projective, exactness at every maximal
% ideal): lTensor_liftQ_injective_of_exact, split_injective_of_forall_lTensor_residueField_injective (through
% localization at maximal ideals and Mathlib's IsLocalRing.split_injective_iff_lTensor_residueField_injective, whose
% local case the first paragraph proves) and Module.projective_lifting_property twice.

\begin{lemma}[Alp\"oge]\label{cv-toroidal-abelian-dvr-flat-regular-homotopy}
Let $A$ be a ring, let $\mathfrak p$ be a prime ideal of $A$ with residue field $\kappa(\mathfrak p)$, and let
$C^0\to C^1\to\cdots$ be a complex of finite free $A$-modules, zero in large degrees, such that
$C\otimes_A\kappa(\mathfrak p)$ is exact in every degree. Write $d$ for its differentials and set $C^{-1}=0$. Then for
every $i\ge0$ there are $t\in A\setminus\mathfrak p$ and homomorphisms $k_0:C^i\to C^{i-1}$ and
$k_1:C^{i+1}\to C^i$ with $dk_0+k_1d=t$ on $C^i$.
\end{lemma}
\leanref{Development/AbelianVarietyTheory/Proofs/UniformizationProofs/IsSMulRegularCokerDualMapOfExactFibres.lean}{GlobalTheory.DualCokernel.exists_homotopy}
\leanref{Development/AbelianVarietyTheory/Proofs/UniformizationProofs/IsSMulRegularCokerDualMapOfExactFibres.lean}{GlobalTheory.DualCokernel.exists_homotopy_step}
\begin{proof}
By descending induction on $i$. For $i$ large, $C^i=0$, and $t=1$ and $k_0=k_1=0$ will do. Otherwise, by induction
there are $t'\notin\mathfrak p$, $h_1:C^{i+1}\to C^i$ and $h_2:C^{i+2}\to C^{i+1}$ with $dh_1+h_2d=t'$ on $C^{i+1}$. Write
$a:C^{i-1}\to C^i$ and $b:C^i\to C^{i+1}$ for the differentials. As $db=0$, $bh_1b=t'b$. Let $\pi=t'-h_1b$ on $C^i$;
then $b\pi=0$. Let $N=a(C^{i-1})+h_1b(C^i)\subset C^i$ and let $q:C^i\to C^i/N$ be the quotient map, so that $qa=0$
and $q\pi=t'q$.

First, $(C^i/N)\otimes_A\kappa(\mathfrak p)=0$. Let $x\in C^i\otimes_A\kappa(\mathfrak p)$. Then $(b\otimes1)(\pi\otimes
1)(x)=0$, so by exactness $(\pi\otimes1)(x)=(a\otimes1)(z)$ for some $z$. Applying $q\otimes1$ gives
$t'(q\otimes1)(x)=0$, and $t'$ is a unit in $\kappa(\mathfrak p)$. As $C^i/N$ is finite, Nakayama's lemma over
$A_{\mathfrak p}$ gives $(C^i/N)_{\mathfrak p}=0$, so there is $t''\notin\mathfrak p$ with $t''C^i\subset N$.

For $x\in C^i$ write $t''x=a(u)+h_1b(v)$. Then $t''\pi(x)=\pi(a(u))+\pi(h_1b(v))=t'a(u)+t'h_1b(v)-h_1bh_1b(v)=
t'a(u)$, as $ba=0$ and $bh_1b=t'b$. So $t''\pi$ maps $C^i$ into $a(C^{i-1})$, and as $C^i$ is free it lifts along $a$ to
$k:C^i\to C^{i-1}$ with $ak=t''\pi$. Then $ak+t''h_1b=t''t'$ on $C^i$; take $t=t''t'$, $k_0=k$ and $k_1=t''h_1$.
\end{proof}
% Lean: exists_homotopy (induction on n + 1 - i) and exists_homotopy_step (pi = s - h1 b, N = range a + range (h1 b),
% exists_smul_eq_zero_of_fibre for t'', Module.projective_lifting_property for k).

\begin{lemma}[Alp\"oge]\label{cv-toroidal-abelian-dvr-flat-regular-zero}
Let $A$ be a noetherian local ring with maximal ideal $\mathfrak m$ and residue field $\kappa$, let $M$ be an
$A$-module with $\operatorname{Hom}_A(\kappa,M)=0$, and let $x\in M$. If for every prime ideal
$\mathfrak p\ne\mathfrak m$ there is $s\in A\setminus\mathfrak p$ with $sx=0$, then $x=0$.
\end{lemma}
\leanref{Development/AbelianVarietyTheory/Proofs/UniformizationProofs/IsSMulRegularCokerDualMapOfExactFibres.lean}{GlobalTheory.DualCokernel.eq_zero_of_extZero}
\begin{proof}
Suppose $x\ne0$. Among the annihilators of the nonzero elements of $Ax$, all of which contain the annihilator of
$x$, a maximal one $P=\operatorname{Ann}(y)$ is a prime ideal \cite[Theorem~6.1]{Mat86}. If $P=\mathfrak m$, then
$1\mapsto y$ is a nonzero homomorphism $\kappa\to M$. Otherwise some $s\notin P$ has $sx=0$, so $s\in P$. Both are
impossible.
\end{proof}
% Lean: eq_zero_of_extZero, with exists_le_isAssociatedPrime_of_isNoetherianRing and
% isAssociatedPrime_iff_exists_injective_linearMap.

\begin{lemma}[Alp\"oge]\label{cv-toroidal-abelian-dvr-flat-regular-step}
Let $A$ be a noetherian local ring with maximal ideal $\mathfrak m$ and residue field $\kappa$, let
$\varpi\in\mathfrak m$, and let $n\ge0$. Let $\beta:D_2\to D_1$ and $\alpha:D_1\to D_0$ be homomorphisms of
$A$-modules with $\alpha\beta=0$ and $D_1$ finite, and let $Q=D_1/\beta(D_2)$. Suppose that:
\begin{enumerate}
\item $\varpi$ is a nonzerodivisor on $D_0$ and on $Q$;
\item $\operatorname{Ext}^j_A(\kappa,D_0)=0$ for $j\le n$, and $\operatorname{Ext}^j_A(\kappa,Q)=0$ for $j\le n+1$;
\item for every prime ideal $\mathfrak p\ne\mathfrak m$ containing $\varpi$ there is $s\notin\mathfrak p$ with
  $s\ker\alpha\subset\beta(D_2)$;
\item for every prime ideal $\mathfrak p\ne\mathfrak m$ containing $\varpi$ there is $s\notin\mathfrak p$ with
  $sx\in\alpha(D_1)$ for every $x\in D_0$ such that $\varpi x\in\alpha(D_1)$.
\end{enumerate}
Then $\varpi$ is a nonzerodivisor on $D_0/\alpha(D_1)$, and $\operatorname{Ext}^j_A(\kappa,D_0/\alpha(D_1))=0$ for
$j\le n$.
\end{lemma}
\leanref{Development/AbelianVarietyTheory/Proofs/UniformizationProofs/IsSMulRegularCokerDualMapOfExactFibres.lean}{GlobalTheory.DualCokernel.coker_step}
\begin{proof}
First, $\operatorname{Hom}_A(\kappa,Q/\varpi Q)=0$: the sequence $0\to Q\xrightarrow{\varpi}Q\to Q/\varpi Q\to0$ is
exact by (1), and in its long exact sequence $\operatorname{Hom}_A(\kappa,Q)\to\operatorname{Hom}_A(\kappa,Q/\varpi Q)
\to\operatorname{Ext}^1_A(\kappa,Q)$ the outer terms vanish by (2).

Second, $\ker\alpha=\beta(D_2)$. One inclusion is $\alpha\beta=0$. The module $\ker\alpha$ is finite, as $D_1$ is
and $A$ is noetherian, so by Nakayama's lemma it suffices to show $\ker\alpha\subset\beta(D_2)+\mathfrak m\ker\alpha$.
Let $x\in\ker\alpha$, with image $\bar x$ in $Q/\varpi Q$. For a prime $\mathfrak p\ne\mathfrak m$ containing
$\varpi$, (3) gives $s\notin\mathfrak p$ with $s\bar x=0$; for a prime not containing $\varpi$, $\varpi\bar x=0$.
By the first point and Lemma~\ref{cv-toroidal-abelian-dvr-flat-regular-zero}, $\bar x=0$, so $x=\varpi z+\beta(y)$
with $z\in D_1$ and $y\in D_2$. Then $\varpi\alpha(z)=\alpha(x)-\alpha\beta(y)=0$, so $\alpha(z)=0$ by (1), and
$x\in\beta(D_2)+\mathfrak m\ker\alpha$.

Third, by the second point $\alpha$ induces an injection $Q\to D_0$ with cokernel $D_0/\alpha(D_1)$. In the long exact
sequence $\operatorname{Ext}^j_A(\kappa,D_0)\to\operatorname{Ext}^j_A(\kappa,D_0/\alpha(D_1))\to
\operatorname{Ext}^{j+1}_A(\kappa,Q)$ the outer terms vanish for $j\le n$ by (2).

Fourth, let $x\in D_0$ with $\varpi x\in\alpha(D_1)$, and let $\bar x$ be its class in $D_0/\alpha(D_1)$. For a prime
$\mathfrak p\ne\mathfrak m$ containing $\varpi$, (4) gives $s\notin\mathfrak p$ with $s\bar x=0$; for a prime not
containing $\varpi$, $\varpi\bar x=0$. By the third point with $j=0$ and
Lemma~\ref{cv-toroidal-abelian-dvr-flat-regular-zero}, $\bar x=0$.
\end{proof}
% Lean: coker_step: extZero_coker for Q/varpi Q, Submodule.le_of_le_smul_of_le_jacobson_bot (Nakayama) with
% eq_zero_of_extZero for ker alpha = range beta, extZero_coker for D_0/range alpha, and eq_zero_of_extZero again.

\begin{lemma}[Alp\"oge]\label{cv-toroidal-abelian-dvr-flat-regular}
Let $A$ be a noetherian local ring, and let $\varpi\in A$ be a nonzerodivisor such that $A/\varpi A$ is a regular
local ring of dimension $g$. Let $C^0\to C^1\to\cdots$ be a complex of finite free $A$-modules, zero in large
degrees. Suppose that for every prime ideal $\mathfrak p$ of $A$ containing $\varpi$ there are a field $K$ and a
homomorphism $A\to K$ with kernel $\mathfrak p$ such that $C\otimes_AK$ is exact in degree $i$ whenever
$\mathfrak p$ is not the maximal ideal or $i>g$. Then $\varpi$ is a nonzerodivisor on the cokernel of the dual
$(C^1)^\vee\to(C^0)^\vee$ of the first differential.
\end{lemma}
\leanref{Development/AbelianVarietyTheory/Proofs/UniformizationProofs/IsSMulRegularCokerDualMapOfExactFibres.lean}{GlobalTheory.isSMulRegular_coker_dualMap_of_exact_fibres}
\begin{proof}
Let $\mathfrak m$ be the maximal ideal of $A$ and $\kappa$ its residue field. Then $\varpi\in\mathfrak m$, as
$A/\varpi A$ is a local ring. Write $d_i:C^i\to C^{i+1}$ for the differentials, $D^i=\operatorname{Hom}_A(C^i,A)$,
$d_i^\vee:D^{i+1}\to D^i$, $f\mapsto f\circ d_i$, with $D^{-1}=0$, and $Q_i=D^i/d_i^\vee(D^{i+1})$. The lemma says that
$\varpi$ is a nonzerodivisor on $Q_0$.

(i) The element $\varpi$ is a nonzerodivisor on each $D^i$, as on $A$, since $D^i$ is finite free.

(ii) $\operatorname{Ext}^j_A(\kappa,D^i)=0$ for $j\le g$. Indeed, the maximal ideal of $A/\varpi A$ is generated by $g$
elements, which form a regular sequence \cite[Tag 00NQ]{Sta}; with lifts $x_1,\dots,x_g\in\mathfrak m$ of them,
$\varpi,x_1,\dots,x_g$ is a regular sequence on $A$ in $\mathfrak m$. So $\operatorname{Ext}^j_A(\kappa,A)=0$ for
$j\le g$ \cite[Tag 00LW]{Sta}, and the same holds for the finite free module $D^i$.

(iii) For every prime $\mathfrak p\ne\mathfrak m$ containing $\varpi$ and every $i$, there is $s\notin\mathfrak p$ such
that $sf\in d_i^\vee(D^{i+1})$ for every $f\in D^i$ with $d_{i-1}^\vee(f)=0$. Indeed, by hypothesis there is a field
$K$ with a homomorphism $A\to K$ of kernel $\mathfrak p$ for which $C\otimes_AK$ is exact in every degree; as $K$ is
faithfully flat over $\kappa(\mathfrak p)$, $C\otimes_A\kappa(\mathfrak p)$ is exact in every degree.
Lemma~\ref{cv-toroidal-abelian-dvr-flat-regular-homotopy} gives $s\notin\mathfrak p$, $k_0$ and $k_1$ with
$d_{i-1}k_0+k_1d_i=s$ on $C^i$. If $f\circ d_{i-1}=0$, then $sf=f\circ k_1\circ d_i=d_i^\vee(f\circ k_1)$.

(iv) We show by descending induction on $i$, for $0\le i\le g$, that $\varpi$ is a nonzerodivisor on $Q_i$ and that
$\operatorname{Ext}^j_A(\kappa,Q_i)=0$ for $j\le i$. The case $i=0$ is the lemma.

For $i=g$: by hypothesis for $\mathfrak p=\mathfrak m$ there is a field $K$ with a homomorphism $A\to K$ of kernel
$\mathfrak m$ for which $C\otimes_AK$ is exact in degree $g+1$, and so $C\otimes_A\kappa$ is exact in degree $g+1$.
Lemma~\ref{cv-toroidal-abelian-dvr-flat-regular-split}, applied to $d_g$ and $d_{g+1}$, gives $s:C^{g+1}\to C^g$ and
$t$ with $d_gs+td_{g+1}=\mathrm{id}$; as $d_{g+1}d_g=0$, $d_gsd_g=d_g$. The map $\iota:Q_g\to D^g$,
$[f]\mapsto f-f\circ s\circ d_g$, is well defined, since for $f=h\circ d_g$ it gives $hd_g-hd_gsd_g=0$; and the class
of $\iota([f])$ is $[f]$, as $f\circ s\circ d_g=d_g^\vee(f\circ s)$. So $Q_g$ is a direct summand of $D^g$, and (i) and
(ii) give the claim.

For $i<g$: apply Lemma~\ref{cv-toroidal-abelian-dvr-flat-regular-step} with $n=i$, $D_0=D^i$, $D_1=D^{i+1}$,
$D_2=D^{i+2}$, $\alpha=d_i^\vee$ and $\beta=d_{i+1}^\vee$, so that $Q=Q_{i+1}$. Its conditions (1) and (2) hold by (i),
(ii) and the induction hypothesis, and (3) is (iii) for $i+1$. For (4), let $\varpi x=d_i^\vee(y)$; then
$\varpi d_{i-1}^\vee(x)=d_{i-1}^\vee d_i^\vee(y)=0$, so $d_{i-1}^\vee(x)=0$ by (i), and (iii) for $i$ gives
$s\notin\mathfrak p$ with $sx\in d_i^\vee(D^{i+1})$.
\end{proof}
% Lean: isSMulRegular_coker_dualMap_of_exact_fibres: hm, hDreg (i), hDext (ii) with extZero_self (the regular sequence
% varpi :: lifts of a regular system of parameters, ModuleCat.subsingleton_ext_of_exists_isRegular) and
% FiniteLengthComplex.subsingleton_ext_free, hloc (iii) with exact_fibre (exact_lTensor_iff_residueField) and
% exists_homotopy, and the induction key (iv): the base i = g through exists_homotopy_of_forall_exact_residueField
% and the retraction iota (extZero_retract), the step through coker_step. Size: about 330 lines.

\begin{lemma}[Alp\"oge]\label{cv-toroidal-abelian-dvr-flat-ideal}
Let $V$ be a ring and $B$ a commutative group scheme over $V$ that is an abelian scheme, projective and smooth of
relative dimension $g$ over $V$. Let $\mathcal L$ and $\mathcal O$ be line bundles on $B$ such that the following two
line bundles are trivial: the line bundle $\Theta(\mathcal L)$ on $B\times_VB\times_VB$ whose fibre at
$(x_1,x_2,x_3)$ is $\bigotimes_I\mathcal L(x_I)^{\otimes\delta_I}$, where $I$ runs over the subsets of $\{1,2,3\}$,
$x_I=\sum_{i\in I}x_i$ and $\delta_I=(-1)^{|I|+1}$; and the line bundle $\Lambda(\mathcal O)=m^*\mathcal O\otimes
\operatorname{pr}_1^*\mathcal O^{-1}\otimes\operatorname{pr}_2^*\mathcal O^{-1}\otimes0^*\mathcal O$ on $B\times_VB$. Let
$H\to B$ be a
morphism such that a point $x\in B(T)$ of a scheme $T$ over $V$ factors through $H$ if and only if
$t_x^*\mathcal L_T\otimes\mathcal L_T^{-1}\otimes\mathcal O_T^{-1}$ is trivial locally on $T$. Let $A$ be a noetherian
local ring with a morphism $\alpha:\operatorname{Spec}A\to B$, and let $\varpi\in A$ be a nonzerodivisor such that
$A/\varpi A$ is a regular local ring of dimension $g$. Suppose that there is a local homomorphism $A\to K_0$ to a field
for which $\operatorname{Spec}K_0\to\operatorname{Spec}A\to B$ factors through $H$, and that every homomorphism
$\phi:A\to K$ to a field with $\phi(\varpi)=0$ for which $\operatorname{Spec}K\to\operatorname{Spec}A\to B$ factors
through $H$ has kernel the maximal ideal. Then there is an ideal $I$ of $A$ such that $\varpi$ is a nonzerodivisor on
$A/I$ and, for every local ring $R$ and homomorphism $\phi:A\to R$, the morphism
$\operatorname{Spec}R\to\operatorname{Spec}A\to B$ factors through $H$ if and only if $I\subset\ker\phi$.
\end{lemma}
\leanref{Development/AbelianVarietyTheory/Proofs/UniformizationProofs/FlatKernelPolarisationCosetsOverDVR.lean}{GlobalTheory.StabFlat.exists_ideal}
\leanref{Development/AbelianVarietyTheory/Proofs/UniformizationProofs/FlatKernelPolarisationCosetsOverDVR.lean}{GlobalTheory.StabFlat.exists_family}
\begin{proof}
Let $B_A=B\times_V\operatorname{Spec}A$, through $\alpha$ and the structure morphism, an abelian scheme over $A$,
projective and smooth of relative dimension $g$; let $\mathcal L_A$ and $\mathcal O_A$ be the pullbacks of
$\mathcal L$ and $\mathcal O$, and let $x\in B_A(A)$ be the point given by $\alpha$. Let
$\mathcal F=t_x^*\mathcal L_A\otimes\mathcal L_A^{-1}\otimes\mathcal O_A^{-1}$ on $B_A$.

First, $\Lambda(\mathcal F)$ is trivial. Indeed, $\Lambda$ is multiplicative in the line bundle, $\Lambda(\mathcal O_A)$
is trivial, and the pullback of a generator of $\Theta(\mathcal L_A)$ along $(x,y_1,y_2)$, where $y_1$ and $y_2$
are the projections of $B_A\times_AB_A$, is an isomorphism $\Lambda(t_x^*\mathcal L_A)\cong\Lambda(\mathcal L_A)$,
as the fibres of the two sides at $(y_1,y_2)$ are those of $\Theta(\mathcal L_A)$ at $(x,y_1,y_2)$ rearranged.

Second, for a local ring $R$ and a homomorphism $\phi:A\to R$, the pullback $\mathcal F_R$ of $\mathcal F$ to $B_R$
is trivial if and only if $\operatorname{Spec}R\to\operatorname{Spec}A\to B$ factors through $H$. Indeed,
$\mathcal F_R$ is $t_{x_R}^*\mathcal L_R\otimes\mathcal L_R^{-1}\otimes\mathcal O_R^{-1}$ for the point $x_R$ of
$B$ given by $\operatorname{Spec}R\to\operatorname{Spec}A\to B$, and as $R$ is local, a line bundle on $B_R$ is
trivial locally on $\operatorname{Spec}R$ if and only if it is trivial.

By Lemma~\ref{cv-toroidal-abelian-dvr-flat-complex}, $\mathcal F$ has a section complex with matrices $D_0$ and $D_1$. The
morphism $B_A\to\operatorname{Spec}A$ is proper and surjective, the functions on $B_{K_0}$ are the constants
\cite[Tag 0E0L]{Sta}, and $\mathcal F_{K_0}$ is trivial by the second point. By
Lemma~\ref{cv-toroidal-abelian-dvr-flat-cyclic}, there is an ideal $I$ of $A$ with
$\operatorname{coker}D_0^{\mathsf T}\cong A/I$ such that, by the second point, a morphism
$\operatorname{Spec}R\to\operatorname{Spec}A\to B$ with $R$ local factors through $H$ if and only if
$I\subset\ker\phi$.

By the first point and Lemma~\ref{cv-toroidal-abelian-dvr-flat-free}, there is a complex $C$ of finite free
$A$-modules, zero in large degrees, with $\operatorname{coker}((C^1)^\vee\to(C^0)^\vee)\cong A/I$, such that
$C\otimes_AK$ is exact in degree $i$ for every algebraically closed field $K$ and homomorphism $A\to K$ with $i>g$
or $\mathcal F_K$ nontrivial. Let $\mathfrak p$ be a prime of $A$ containing $\varpi$, let $K$ be an algebraic
closure of its residue field, and $\phi:A\to K$ the homomorphism, with kernel $\mathfrak p$. If $\mathfrak p$ is not
the maximal ideal, then $\mathcal F_K$ is nontrivial: otherwise $\operatorname{Spec}K\to B$ factors through $H$ by
the second point, and the hypothesis gives that $\mathfrak p=\ker\phi$ is the maximal ideal. So $C\otimes_AK$ is
exact in degree $i$ whenever $\mathfrak p$ is not maximal or $i>g$, and
Lemma~\ref{cv-toroidal-abelian-dvr-flat-regular} shows that $\varpi$ is a nonzerodivisor on $A/I$.
\end{proof}
% Lean: exists_ideal (with GlobalFnsSt): exists_family (F with the class dcl of x and b, and its primitivity
% from dcl_primitive, that is lam_mul_left and gen_lam_iff), TLinPf.isAb_pullback and isProjBy_baseChange (B_A),
% test (with cls_eq_dcl and loc_triv_iff), T.nonempty_perfectSectionsComplex_of_smooth_proj,
% SecCyclic.cyclic_ideal with CubeOpen.bijective_appTop (the law GlobalFnsSt; Stacks Tag 0E0L),
% T.exists_perfect_free_finite_complex and isSMulRegular_coker_dualMap_of_exact_fibres. In the Lean V is any ring
% and the morphism H -> B need not be a closed immersion.

\begin{lemma}[Alp\"oge]\label{cv-toroidal-abelian-dvr-flat}
Let $V$ be a discrete valuation ring and $B$ a commutative group scheme over $V$ that is an abelian scheme and
projective over $V$. Let $\mathcal L$ be a line bundle on $B$, trivialized along the zero section, with a
cubical structure, and let $\mathcal O$ be a line bundle on $B$ such that
$\Lambda(\mathcal O)=m^*\mathcal O\otimes\operatorname{pr}_1^*\mathcal O^{-1}\otimes\operatorname{pr}_2^*
\mathcal O^{-1}\otimes0^*\mathcal O$ is trivial on $B\times_VB$. Let $H$ be the closed subscheme of $B$ of
Lemma~\ref{cv-toroidal-abelian-dvr-closed}(1) for $\mathcal L$ and $\mathcal O$. If $H$ is finite over $V$,
then $H$ is flat over $V$.
\end{lemma}
\leanref{Development/AbelianVarietyTheory/Proofs/UniformizationProofs/FlatKernelPolarisationCosetsOverDVRStatement.lean}{GlobalTheory.T.flat_kernelPolarisation_cosets_over_DVR}
\begin{proof}
Let $\varpi$ be a uniformizer of $V$, $\kappa$ its residue field, and $g$ the relative dimension of $B$, which is
smooth over $V$. Let $j:H\to B$ be the closed immersion. By Lemma~\ref{cv-toroidal-abelian-dvr-flat-criterion}, it
suffices to find, for each point $h\in H$ at which the image of $\varpi$ in $A=\mathcal O_{B,j(h)}$ is not a unit, an
ideal $I$ of $A$ with the properties stated there. Such a point lies over the closed point of
$\operatorname{Spec}V$. As $H$ is finite over $V$, its fibre over the closed point is a finite discrete closed
subset of $H$; so its points are closed in $H$, and their images are closed in $B$. In particular $j(h)$ is a closed
point of $B$ over the closed point of $\operatorname{Spec}V$, and by Lemma~\ref{cv-toroidal-abelian-dvr-flat-stalk},
$A$ is noetherian, $\varpi$ is a nonzerodivisor of $A$, and $A/\varpi A$ is a regular local ring of dimension $g$.

We apply Lemma~\ref{cv-toroidal-abelian-dvr-flat-ideal} to $B$, $\mathcal L$, $\mathcal O$, $j$, $A$ with the
canonical morphism $\operatorname{Spec}A\to B$, and $\varpi$. The line bundle $\Theta(\mathcal L)$ is trivial, by the
cubical structure, and so is $\Lambda(\mathcal O)$; the morphism $j$ has the stated property by
Lemma~\ref{cv-toroidal-abelian-dvr-closed}(1). Let $K_0$ be the residue field of $H$ at $h$; then
$\operatorname{Spec}K_0\to\operatorname{Spec}A\to B$ is the point $h$ of $H$. Let $\phi:A\to K$ be a homomorphism
to a field with $\phi(\varpi)=0$ for which $\operatorname{Spec}K\to\operatorname{Spec}A\to B$ factors through $H$,
and let $\mathfrak p=\ker\phi$, a point of $\operatorname{Spec}A$. Its image $z$ in $B$ lies in $j(H)$ over the
closed point of $\operatorname{Spec}V$, so it is closed in $B$; and $z$ is a generization of $j(h)$, as are all
points in the image of $\operatorname{Spec}A$. So $z=j(h)$, and $\mathfrak p$ is the maximal ideal of $A$.
Lemma~\ref{cv-toroidal-abelian-dvr-flat-ideal} now gives the ideal $I$, and $H$ is flat over $V$.
\end{proof}
% Lean: flat_of_isFinite (with GlobalFnsSt): AbelianDim.exists_relDim, flat_of_ideals (Lemma
% cv-toroidal-abelian-dvr-flat-criterion), isClosed_of_isFinite (the points of the special fibre of H are closed),
% stalk_of_smooth (Lemma cv-toroidal-abelian-dvr-flat-stalk) and exists_ideal (Lemma
% cv-toroidal-abelian-dvr-flat-ideal, with K_0 the residue field of H at h, and hE from fromSpecStalk and
% range_fromSpecStalk). Lemmas cv-toroidal-abelian-dvr-flat-complex, -flat-cyclic, -flat-free and -flat-regular are
% written at the floor; the theorem on cohomology and base change and four facts on coherent cohomology and abelian
% varieties are cited there.

\subsection{Abelian schemes over a discrete valuation ring}\label{subsection-cv-abelian-schemes-dvr}

\begin{lemma}[Alp\"oge]\label{cv-toroidal-abelian-dvr-kernel}
Let $V$ be a discrete valuation ring with fraction field $K$, let $B$ be a commutative group scheme over $V$
that is an abelian scheme, and let $\mathcal L$ be a line bundle on $B$, trivialized along the zero section,
with a cubical structure, that is relatively ample. Let $B'$ be a group scheme over $K$ with a line bundle
$\mathcal L'$, and let $\varphi:B'\to B$ be a morphism over $\operatorname{Spec}K\to\operatorname{Spec}V$
that identifies $B'$ with $B\times_V\operatorname{Spec}K$ as group schemes, with
$\mathcal L'\cong\varphi^*\mathcal L$. If $\mathcal L'$ defines a principal polarization, then so does
$\mathcal L$.
\end{lemma}
\leanref{Development/AbelianVarietyTheory/Proofs/ModuliProofs/ExistsPrincipallyPolarizedCMExtensionDvr.lean}{GlobalTheory.GoodPol.kTriv_of_generic}
\leanref{Development/AbelianVarietyTheory/Proofs/ModuliProofs/ThetaGroupKTrivialOfPrincipalGenericFibreDvrStatement.lean}{GlobalTheory.T.thetaGroupK_trivial_of_principal_genericFibre_dvr}
\begin{proof}
As $\operatorname{Spec}V$ is affine and $\mathcal L$ is relatively ample, sections of a power of $\mathcal L$
define a closed immersion of $B$ into a projective space over $V$, so $B$ is projective over $V$. Let $H$ be
the closed subscheme of Lemma~\ref{cv-toroidal-abelian-dvr-closed}(1) for $\mathcal L$ and
$\mathcal O=\mathcal O_B$, with inclusion $i:H\to B$, so that the points of $H$ are those of
$K(\mathcal L)$. By Lemma~\ref{cv-toroidal-abelian-dvr-closed}(2), $H$ is finite over $V$. As
$\Lambda(\mathcal O_B)$ is trivial and $\mathcal L$ has a cubical structure,
Lemma~\ref{cv-toroidal-abelian-dvr-flat} shows that $H$ is flat over $V$.

Let $H_K=H\times_V\operatorname{Spec}K$, with projection $\pi:H_K\to H$. As $H$ is finite and flat over $V$, it
is the spectrum of a finite flat $V$-algebra $C$, which is torsion free, so $C\to C\otimes_VK$ is injective
and $\pi$ is scheme-theoretically dominant. The point $i\circ\pi$ of $B$ lies over
$\operatorname{Spec}K\to\operatorname{Spec}V$, so it is $\varphi\circ x$ for a point $x\in B'(H_K)$. Through
$\mathcal L'\cong\varphi^*\mathcal L$, the line bundle $\mathcal L'_x$ on $B'_{H_K}$ is the pullback of
$\mathcal L_{i\circ\pi}$, which is trivial locally on $H_K$ because $i\circ\pi$ factors through $H$. So $x$ is a
point of $K(\mathcal L')$, and $x=0$ since $\mathcal L'$ defines a principal polarization. Hence $i\circ\pi$
is the composite of $\pi$, the structure morphism of $H$ and the zero section of $B$. As $\pi$ is
scheme-theoretically dominant and $B$ is separated over $V$, $i$ itself is this composite of the structure
morphism and the zero section. Every point of $K(\mathcal L)$ factors through $i$, so it is zero.
\end{proof}
% Lean: GoodPol.kTriv_of_generic: Unif.relAmple_iff_veryAmpleBy (IsProjBy), TrivTwist.inK_iff_trivial_twist
% (pw L 0 is the trivial bundle, with its generator of abelian2), TwistKer.exists_closed_twist_kernel,
% AmpleStabFinite.isFinite_of_relAmple, T.flat_kernelPolarisation_cosets_over_DVR,
% schemeTheoreticallyDominant_fractionField, eq_one_of_inK and over_ext_of_schemeTheoreticallyDominant.
% T.thetaGroupK_trivial_of_principal_genericFibre_dvr = KTrivSpread.kTriv_of_generic_fibre is the same argument
% (AmpleBaseChange.isProjBy_smooth_isProper_geometricallyIntegral for IsProjBy, kTriv_of_finite_flat for the
% last paragraph).

\begin{lemma}[Alp\"oge]\label{cv-toroidal-abelian-dvr-specialize}
Let $S$ be an integral, locally noetherian scheme, let $K$ be a field and $x:\operatorname{Spec}K\to S$ a
morphism whose image is the generic point of $S$. Let $P$ be an abelian scheme over $S$ with a line bundle
that is trivialized along the zero section, isomorphic to its pullback by $-1$ locally on $S$, relatively
ample and defines a principal polarization, and let $\mathcal M$ be a line bundle on $P$. If the pullback of
$\mathcal M$ to $P\times_S\operatorname{Spec}K$ is relatively ample, then $\mathcal M$ is relatively ample.
\end{lemma}
\leanref{Development/AbelianVarietyTheory/Proofs/ModuliProofs/ExistsSecNeZeroOfLbIso.lean}{GlobalTheory.AmpleSpecialisation.relAmple_of_genericFibre}
\begin{proof}
Let $\mathcal L$ be the line bundle of $P$, and let $P_K=P\times_S\operatorname{Spec}K$, with the pullbacks
$\mathcal L_K$ and $\mathcal M_K$. For a point $y:\operatorname{Spec}K'\to S$ write $P_y$, $\mathcal L_y$ and
$\mathcal M_y$ for the base changes along $y$.

Step 1. The scheme $P_K$ is proper over $K$ and nonempty, and $\mathcal M_K$ is ample \cite[Tag 01VH]{Sta}.
So for some $a\ge1$ the line bundle $\mathcal M_K^a\otimes\mathcal L_K^{-1}$ is generated by its global
sections \cite[Tag 01Q3]{Sta}, and it has a nonzero section. Let
$\mathcal N=\mathcal M^a\otimes\mathcal L^{-1}$ on $P$.

Step 2. For every point $y:\operatorname{Spec}K'\to S$, $\mathcal N_y$ has a nonzero section. Indeed, $P$ is
proper and of finite presentation over $S$, and $\mathcal N$ is flat over $S$, so the function
$s\mapsto\dim_{\kappa(s)}H^0(P_s,\mathcal N_s)$ is upper semicontinuous, and its formation commutes with base
change \cite[Tag 0BDN]{Sta}. By Step 1 and the base change from the residue field of the generic point $\eta$
to $K$, it is at least $1$ at $\eta$; so it is at least $1$ on the closure of $\eta$, which is $S$, and by base
change $H^0(P_y,\mathcal N_y)\ne0$.

Step 3. Let $y:\operatorname{Spec}K'\to S$ with $K'$ algebraically closed. Then $P_y$ is an abelian variety
over $K'$, and $\mathcal N_y$ has a nonzero section $\sigma$ by Step 2. The line bundle $\mathcal N_y^3$ is
generated by its global sections. Indeed, the nonvanishing locus $D$ of $\sigma$ is a dense open subset of the
integral scheme $P_y$. Let $w$ be a closed point of $P_y$. There are points $u,v\in P_y(K')$ such that $w+u$,
$w+v$ and $w-u-v$ lie in $D$, since nonempty open subsets of $P_y$ meet and their $K'$-points are dense. By the
theorem of the square \cite[\S6, Corollary~4]{Mum70}, applied twice,
$t_u^*\mathcal N_y\otimes t_v^*\mathcal N_y\otimes t_{-u-v}^*\mathcal N_y\cong\mathcal N_y^3$, so
$t_u^*\sigma\otimes t_v^*\sigma\otimes t_{-u-v}^*\sigma$ is a section of $\mathcal N_y^3$ that does not vanish
at $w$. The nonvanishing loci of these sections form an open set that contains every closed point of the
scheme $P_y$, of finite type over $K'$, so it is $P_y$.

Step 4. With $y$ as in Step 3, $\mathcal L_y$ is ample, as $\mathcal L$ is relatively ample
\cite[Tags 0893 and 01VH]{Sta}. So $\mathcal L_y^3$ is ample, and by Step 3 and \cite[Tag 0890]{Sta},
$\mathcal L_y^3\otimes\mathcal N_y^3=\mathcal M_y^{3a}$ is ample; hence $\mathcal M_y$ is ample
\cite[Tag 01PT]{Sta}.

Step 5. Let $s\in S$ and let $y:\operatorname{Spec}K'\to S$ be the point given by an algebraic closure $K'$ of
$\kappa(s)$. By Step 4 and descent along $\operatorname{Spec}K'\to\operatorname{Spec}\kappa(s)$
\cite[Tag 0D2P]{Sta}, $\mathcal M_s$ is ample on $P_s$. As $P$ is proper over $S$, $\mathcal M$ is relatively
ample over an open neighbourhood of $s$ \cite[Tag 0D2S]{Sta}. Relative ampleness is local on $S$
\cite[Tag 01VJ]{Sta}, so $\mathcal M$ is relatively ample.
\end{proof}
% Lean: relAmple_of_genericFibre: T.exists_sec_relAmplePow_div_ne_zero (Step 1; AmpleDivSec, about 900 lines),
% lbIso_pull_div_lbPow and exists_sec_ne_zero_of_lbIso (transport of the section), then
% T.relAmple_of_ample_geometricFibres_abelianScheme (Step 5; FibreAmple.relAmple_of_geometricFibres, about 9,100,
% with Out.Laws) applied to T.exists_section_ne_zero_fibre_of_generic_section (Step 2;
% GenericSection.exists_sec_ne_zero_of_generic, about 6,400, for any integral S),
% T.exists_sections_cover_cube_effective_abelianVariety (Step 3; EffectiveCube.exists_sections_cube_cover with
% the theorem of the cube for Theta, NoBasePoint.exists_two_points and SepPoints.exists_gen_three) and
% T.relAmple_of_cube_div_generated_field (Step 4; about 1,000). The Lean proves the facts cited to Stacks in
% Steps 1, 2, 4 and 5 itself.

\begin{lemma}[Alp\"oge]\label{cv-toroidal-abelian-dvr-isogenous}
Let $R$ be a Dedekind domain with fraction field $K$, and let $P$ be an abelian scheme over $R$ with a line
bundle that is trivialized along the zero section, isomorphic to its pullback by $-1$ locally on
$\operatorname{Spec}R$, relatively ample and defines a principal polarization. Let $B$ be a group scheme over
$R$ that is an abelian scheme, let $\Psi:P\to B$ be a morphism over $R$ that is finite, flat and surjective,
and let $\mathcal M$ be a line bundle on $B$. Let $B'$ be a group scheme over $K$, let $\varphi:B'\to B$ be a
morphism over $\operatorname{Spec}K\to\operatorname{Spec}R$ that identifies $B'$ with
$B\times_R\operatorname{Spec}K$ as group schemes, and let $\mathcal M'$ be a line bundle on $B'$ with
$\mathcal M'\cong\varphi^*\mathcal M$. If $\mathcal M'$ is relatively ample, then $\mathcal M$ is relatively
ample.
\end{lemma}
\leanref{Development/AbelianVarietyTheory/Proofs/CubicalLineBundles/RelAmpleOfGenericIsogenousDedekindStatement.lean}{GlobalTheory.T.relAmple_of_generic_isogenous_dedekind}
\begin{proof}
The base $\operatorname{Spec}K$ is affine, so $\mathcal M'$ is ample, and hence so is $\varphi^*\mathcal M$
\cite[Tag 01VH]{Sta}. Let $P_K=P\times_R\operatorname{Spec}K$, and let $g:P_K\to B'$ be the morphism through
which the base change of $\Psi$ factors, so that $\varphi\circ g=\Psi\circ q$ for the projection
$q:P_K\to P$. It is a base change of $\Psi$, hence finite, so $g^*\varphi^*\mathcal M$ is ample
\cite[Tag 0892]{Sta}; it is the pullback of $\Psi^*\mathcal M$ to $P_K$. The ring $R$ is noetherian and
integral, and $\operatorname{Spec}K\to\operatorname{Spec}R$ has image the generic point, so
Lemma~\ref{cv-toroidal-abelian-dvr-specialize}, applied to $P$ and $\Psi^*\mathcal M$, shows that
$\Psi^*\mathcal M$ is relatively ample. Let $U=\operatorname{Spec}A$ be an affine open subscheme of
$\operatorname{Spec}R$. Over $U$, $\Psi$ is a finite surjective morphism between schemes proper over the
noetherian ring $A$, and the pullback of $\mathcal M$ to $P_U$ is ample; so $\mathcal M$ is ample on $B_U$
\cite[Tag 0B5V]{Sta}. Hence $\mathcal M$ is relatively ample \cite[Tag 01VJ]{Sta}.
\end{proof}
% Lean: AmpleDedekind.relAmple_of_generic_of_isogenous: relAmple_iff_isAmpleLB_of_isAffine (a0),
% IsAmpleLB.pull_of_isFinite with FiniteAmple.lbIso_pull_pull (a1), AmpleSpecialisation.relAmple_of_genericFibre
% (a2) and relAmple_finite_pow_iso (3) (FiniteAmple.isAmpleLB_of_pull, which also uses that Psi is flat). The
% text cites Stacks for these three facts on ample line bundles, which the Lean proves itself (about 2,500
% lines).

\begin{lemma}[Alp\"oge]\label{cv-toroidal-abelian-dvr-cubical-prime}
Let $X$ be an integral, locally noetherian scheme with function field $F$, let $t\in\Gamma(X,\mathcal O_X)$, and
let $\eta\in X$. Write $X_t$ for the open subscheme on which $t$ is invertible. Suppose that a point of $X$ lies
outside $X_t$ if and only if it is a specialization of $\eta$, and that the image of $t$ in $\Gamma(W,\mathcal O_X)$
is a prime element for every affine open subscheme $W$ containing $\eta$. Let $\mathcal N$ be a line bundle on $X$.
\begin{enumerate}
\item The local ring $\mathcal O_{X,\eta}$ is a discrete valuation ring, with uniformizer the image of $t$.
\item Let $r$ be a nonzero section of $\mathcal N$ over $X_t$. There is exactly one integer $n$ for which there is
  a section $s$ of $\mathcal N$ over $X$ with $s|_{X_t}=t^nr$ that does not vanish at some point outside $X_t$.
\item If moreover $r$ vanishes nowhere on $X_t$, then the section $s$ of (2) vanishes nowhere on $X$.
\end{enumerate}
\end{lemma}
\leanref{Development/AbelianVarietyTheory/Proofs/FoundationProofs/EqZeroOfForallPowDvdGerm.lean}{GlobalTheory.PrimeSection.isDiscreteValuationRing_stalk}
\leanref{Development/AbelianVarietyTheory/Proofs/FoundationProofs/EqZeroOfForallPowDvdGerm.lean}{GlobalTheory.PrimeSection.existsUnique_zpow_mul}
\leanref{Development/AbelianVarietyTheory/Proofs/FoundationProofs/EqZeroOfForallPowDvdGerm.lean}{GlobalTheory.PrimeSection.exists_zpow_mul_isUnit}
\begin{proof}
As $t$ is a prime element of some ring of functions, $t\ne0$, so the generic point $\xi$ of $X$ lies in $X_t$.
Every open subscheme of $X$ that contains a point outside $X_t$ contains $\eta$, a generization of that point.

(1) Let $W$ be an affine open subscheme containing $\eta$, let $A=\Gamma(W,\mathcal O_X)$, a noetherian domain,
and let $\mathfrak p$ be the prime ideal of $A$ at $\eta$. Then $t\in\mathfrak p$, as $\eta\notin X_t$. The prime
ideal $tA$ is the prime of a point of $W$ outside $X_t$, which is a specialization of $\eta$; so
$\mathfrak p\subset tA$, and $\mathfrak p=tA$. Hence $\mathcal O_{X,\eta}=A_{tA}$ is a noetherian local domain whose
maximal ideal is generated by the nonzero element $t$, that is, a discrete valuation ring \cite[Tag 00PD]{Sta}.

Two facts on functions follow. Let $W\subset X$ be open. First, a function $f$ on $W\cap X_t$ extends to $W$ if
$\eta\notin W$, or if $\eta\in W$ and the image of $f$ in $F$ lies in $\mathcal O_{X,\eta}$. If $\eta\notin W$, then
$W\subset X_t$. Otherwise, let $x\in W$ lie outside $X_t$, and let $V\subset W$ be an affine open neighbourhood of
$x$, with $B=\Gamma(V,\mathcal O_X)$; then $\eta\in V$ and $t$ is a prime element of $B$. The function $f$ on
$V\cap X_t$ lies in $B_t$, and its image in $F$ lies in $B_{tB}=\mathcal O_{X,\eta}$. Write it as $b/t^m$ with
$b\in B$ and $m\ge0$ least. If $m\ge1$, then $t$ does not divide $b$; but $ub=t^mb'$ for some $u\in B\setminus tB$
and $b'\in B$, and as $t$ is prime, $t$ divides $b$. So $m=0$ and $f$ extends to $V$. These extensions agree with
each other and with $f$ at $\xi$, so on overlaps, as $X$ is integral, and they glue. Second, a function $s$ on $W$
is a unit if it is a unit on $W\cap X_t$ and, in case $\eta\in W$, its germ at $\eta$ is a unit. Indeed, for
$x$ and $V$ as above, the image of $s^{-1}$ in $F$ lies in $\mathcal O_{X,\eta}$, so by the first fact $s^{-1}$
extends from $V\cap X_t$ to $V$, and $s$ is a unit on $V$.

(2) Choose an open cover $(U_a)$ of $X$ with generators $\epsilon_a$ of $\mathcal N$ over $U_a$, and write
$r=r_a\epsilon_a$ on $U_a\cap X_t$. Let $\eta\in U_{a_0}$. As $r\ne0$ and $X$ is integral, the image of $r_{a_0}$
in $F$ is nonzero, so by (1) it is $t^{-n}u$ for exactly one integer $n$ and a unit $u$ of $\mathcal O_{X,\eta}$.
For each $a$ with $\eta\in U_a$, the image of $t^nr_a$ in $F$ is $u$ times the germ at $\eta$ of the unit
$\epsilon_{a_0}/\epsilon_a$, so it lies in $\mathcal O_{X,\eta}$. By the first fact, each $t^nr_a$ extends to a
function $s_a$ on $U_a$. The sections $s_a\epsilon_a$ agree on overlaps, as they agree at $\xi$, and they glue to
a section $s$ of $\mathcal N$ over $X$ with $s|_{X_t}=t^nr$. The germ of $s_{a_0}$ at $\eta$ is $u$, so $s$ does
not vanish at $\eta$. Conversely, let $s'$ be a section with $s'|_{X_t}=t^{n'}r$ that does not vanish at a point
$x\notin X_t$. Its nonvanishing locus is open and contains $x$, so it contains $\eta$. Write $s'=s'_{a_0}
\epsilon_{a_0}$ on $U_{a_0}$; the germ of $s'_{a_0}$ at $\eta$ is a unit with image $t^{n'-n}u$ in $F$, so
$n'=n$.

(3) On $U_a\cap X_t$, the function $s_a$ is $t^nr_a$, a unit. If $\eta\in U_a$, the germ of $s_a$ at $\eta$ is a
unit, as $s$ does not vanish at $\eta$. By the second fact, $s_a$ is a unit on $U_a$, so $s$ vanishes nowhere.
\end{proof}
% Lean: (1) isDiscreteValuationRing_stalk (primeIdealOf_eq_span, PrimeElement.isDiscreteValuationRing).
% The two facts are exists_res_eq (exists_section_restrict_eq_of_germ_eq, PrimeElement.exists_algebraMap_eq with
% range_inf_range) and isUnit_of_isUnit_res. (2) existsUnique_zpow_mul (PrimeElement.existsUnique_zpow_mul_eq_unit,
% compatible_of_res_eq, mem_basicOpen_of_compatible). (3) exists_zpow_mul_isUnit. In the Lean the line bundle is
% given by its charts and transition units, and sections by their functions in the charts.

\begin{lemma}[Alp\"oge]\label{cv-toroidal-abelian-dvr-cubical-dvr}
Let $V$ be a discrete valuation ring with uniformizer $\varpi$, fraction field $K$ and residue field $\kappa$. Let
$H$ be a scheme flat and locally of finite type over $V$ whose fibres $H_K$ and $H_\kappa$ are integral, and let
$\mathcal N$ be a line bundle on $H$.
\begin{enumerate}
\item The scheme $H$ is integral, and $H$ with $t=\varpi$ and the generic point $\eta$ of $H_\kappa$ satisfies the
  hypotheses of Lemma~\ref{cv-toroidal-abelian-dvr-cubical-prime}, with $H_\varpi=H_K$.
\item Let $\sigma$ be a nonzero section of $\mathcal N$ over $H_K$. There is exactly one integer $n$ such that
  $\varpi^n\sigma$ extends to a section of $\mathcal N$ over $H$ that does not vanish at some point of $H_\kappa$.
\item Let $e:\operatorname{Spec}V\to H$ be a section, let $\tau$ be a generator of $\mathcal N$ over $H_K$, and let
  $\rho$ be a generator of $e^*\mathcal N$ whose restriction to $\operatorname{Spec}K$ is the pullback of $\tau$
  along the section $e_K$ of $H_K$ induced by $e$. Then $\tau$ extends uniquely to a generator of $\mathcal N$ over
  $H$.
\end{enumerate}
\end{lemma}
\leanref{Development/AbelianVarietyTheory/Proofs/FoundationProofs/IsDiscreteValuationRingStalkGeneric.lean}{GlobalTheory.DvrModel.isIntegral}
\leanref{Development/AbelianVarietyTheory/Proofs/FoundationProofs/IsDiscreteValuationRingStalkGeneric.lean}{GlobalTheory.DvrModel.existsUnique_zpow_mul_sec}
\leanref{Development/AbelianVarietyTheory/Proofs/FoundationProofs/IsDiscreteValuationRingStalkGeneric.lean}{GlobalTheory.DvrModel.existsUnique_gen}
\leanref{Development/AbelianVarietyTheory/Proofs/AbelianSchemeLawProofs/ExistsUniqueGeneratorExtendDiscreteValuationRing.lean}{GlobalTheory.GenExtend.existsUnique_gen_extend}
\begin{proof}
(1) The generic fibre $H_K$ is the open subscheme $H_\varpi$, as $K=V[1/\varpi]$, and the points outside it are
those of $H_\kappa$, the specializations of $\eta$. As $H$ is flat over $V$, $\varpi$ is a nonzerodivisor on
$\Gamma(W,\mathcal O_H)$ for every affine open $W$. So every nonempty open subscheme of $H$ meets $H_\varpi$, since
$\varpi$ is not nilpotent on it, and a nilpotent function on $W$ is $0$, since it is $0$ on the reduced scheme
$W\cap H_\varpi$ and hence killed by a power of $\varpi$. Thus $H$ is reduced and irreducible, as $H_\varpi$ is;
so $H$ is integral. If $W$ is affine and contains $\eta$, then $W\cap H_\kappa$ is a nonempty affine open
subscheme of the integral scheme $H_\kappa$, with ring $\Gamma(W,\mathcal O_H)/\varpi$; so this ring is a domain
and $\varpi$ is a prime element of $\Gamma(W,\mathcal O_H)$.

(2) This is Lemma~\ref{cv-toroidal-abelian-dvr-cubical-prime}(2), by (1).

(3) Two generators that agree on the nonempty open subscheme $H_K$ of the integral scheme $H$ are equal. By (1)
and Lemma~\ref{cv-toroidal-abelian-dvr-cubical-prime}(2) and (3), there are an integer $n$ and a generator $s$ of
$\mathcal N$ over $H$ with $s|_{H_K}=\varpi^n\tau$. We compare $s$ with $\rho$ in a chart. Let $Q$ be an open
subscheme of $H$ that contains the image under $e$ of the closed point, with a generator $\epsilon$ of $\mathcal N$
over $Q$. Let $Z=e^{-1}(Q)$, an open subscheme of $\operatorname{Spec}V$ that contains the closed point, and let
$Z'$ be its inverse image in $\operatorname{Spec}K$, which is nonempty. Write $s=\sigma\epsilon$ on $Q$, with
$\sigma$ a unit; $\tau=r\epsilon$ on $Q\cap H_K$; and $\rho=u\,e^*\epsilon$ on $Z$, with $u$ a unit of
$\Gamma(Z,\mathcal O)$. Then $\sigma=\varpi^nr$ on $Q\cap H_K$, and by the hypothesis on $\rho$ the functions $u$
and $e_K^*r$ have the same image in $\Gamma(Z',\mathcal O)$. So the unit $e^*\sigma$ of $\Gamma(Z,\mathcal O)$ and
$\varpi^nu$ have the same image in $\Gamma(Z',\mathcal O)$. The restriction
$\Gamma(Z,\mathcal O)\to\Gamma(Z',\mathcal O)$ is injective, as $\operatorname{Spec}V$ is integral and $Z'$ is
nonempty. The image of $\varpi$ in $\Gamma(Z,\mathcal O)$ is not a unit, as it vanishes at the closed point. If
$n>0$, then $e^*\sigma=\varpi^nu$ in $\Gamma(Z,\mathcal O)$, and $\varpi$ is a unit there; if $n<0$, then
$\varpi^{-n}e^*\sigma=u$, and again $\varpi$ is a unit. So $n=0$, and $s$ extends $\tau$.
\end{proof}
% Lean: (1) DvrModel.isIntegral (isIntegral_of_basicOpen with res_mem_nonZeroDivisors), opensRange_generic,
% notMem_basicOpen_iff_specializes and prime_res_generic (ker_app_eq_span). (2) existsUnique_zpow_mul_sec. (3)
% existsUnique_gen and, for a linear combination of pullbacks of a line bundle, GenExtend.existsUnique_gen_extend,
% both through existsUnique_isUnit_appLE_eq: uniqueness by appLE_injective; existence by
% PrimeSection.exists_isUnit_appLE_eq: exists_transport (the functions r_a on the charts of H_K), exists_zpow_mul_isUnit
% (n and s), and zpow_eq_zero for n = 0, in the chart a_0 at e(closed point) with A = e^{-1}(U_{a_0}) and A' its
% inverse image in Spec K (appLE_injective for A' -> A, and the non-unit p = e^*(t) on A).

\begin{lemma}[Alp\"oge]\label{cv-toroidal-abelian-dvr-cubical-normal}
Let $S$ be an integral, locally noetherian, normal scheme and let $j:U\to S$ be an open immersion with $U$
nonempty. Let $p:H\to S$ be a smooth morphism with geometrically integral fibres and with a section $e$, let
$H_U=H\times_SU$, and let $e_U$ be the section of $H_U$ over $U$ induced by $e$. Let $\mathcal N$ be a line bundle
on $H$, let $\tau$ be a generator of $\mathcal N$ over $H_U$, and let $\rho$ be a generator of $e^*\mathcal N$
whose restriction to $U$ is $e_U^*\tau$. Then $\tau$ extends uniquely to a generator of $\mathcal N$ over $H$.
\end{lemma}
\leanref{Development/AbelianVarietyTheory/Proofs/DegenerationProofs/ExistsUniqueRigidifiedGeneratorExtendNormalStatement.lean}{GlobalTheory.T.existsUnique_rigidified_generator_extend_normal}
\leanref{Development/AbelianVarietyTheory/Proofs/DegenerationProofs/ExistsUniqueRigidifiedGeneratorExtendNormal.lean}{GlobalTheory.GenExtendNormal.existsUnique_gen_extend}
\leanref{Development/AbelianVarietyTheory/Proofs/CoherentCohomologyProofs/ExistsUniqueLbIsoPullEqOfCoheightLeOne.lean}{GlobalTheory.exists_germToFunctionField_eq_of_coheight_eq_one}
\begin{proof}
The scheme $H$ is locally noetherian and normal \cite[Tag 033C]{Sta}. The morphism $p$ is open
\cite[Tag 01UA]{Sta} and surjective, with connected fibres, and $S$ is connected, so $H$ is connected
\cite[Tag 0377]{Sta}; hence $H$ is integral \cite[Tag 033N]{Sta}. Its open subscheme $H_U$ is nonempty, so two
generators that agree on $H_U$ are equal. This gives uniqueness.

For existence, let $Q\subset H$ be open with a generator $\epsilon$ of $\mathcal N$ over $Q$, and write
$\tau=f\epsilon$ on $Q\cap H_U$, with $f$ a unit there. It suffices to show that $f$ extends to a unit on $Q$: the
sections $f\epsilon$ so obtained then agree on overlaps, as they agree on $H_U$, and they glue to the extension of
$\tau$. As $H$ is integral, normal and locally noetherian, a rational function on $Q$ that lies in
$\mathcal O_{H,x}$ for every point $x\in Q$ of codimension one is a function on $Q$ \cite[Tag 031T]{Sta}. We apply
this to $f$ and $f^{-1}$.

Let $x\in Q$ have codimension one. If $p(x)\in U$, then $x\in Q\cap H_U$. Otherwise let $s=p(x)$. Generalizations
lift along the flat morphism $p$ \cite[Tag 03HV]{Sta}, so $s$ has codimension at most one; it is not the generic
point of $S$, which lies in $U$, so it has codimension one, and $V=\mathcal O_{S,s}$ is a discrete valuation ring
\cite[Tag 00PD]{Sta}, as $S$ is normal. Let $K$ be its fraction field, $v:\operatorname{Spec}V\to S$ the canonical
morphism, $H_V=H\times_S\operatorname{Spec}V$ with projection $g:H_V\to H$, and $e_V$ the section of $H_V$ induced
by $e$. The scheme $H_V$ is smooth over $V$ with integral fibres. The morphism $\operatorname{Spec}K\to S$ factors
through $U$, as its image is the generic point; so the pullback of $\tau$ to the generic fibre of $H_V$ is a
generator, and the pullback of $\rho$ along $v$ is a generator of $e_V^*g^*\mathcal N$ that agrees with it along
the section over $K$. By Lemma~\ref{cv-toroidal-abelian-dvr-cubical-dvr}(3), the pullback of $\tau$ extends to a
generator $c$ of $g^*\mathcal N$ over $H_V$. The morphism $g$ is a base change of $v$, so it is injective on
points and induces surjections on local rings, and it maps the generic point of $H_V$ to that of $H$. Let $y\in H_V$
with $g(y)=x$, and write $c=f'g^*\epsilon$ near $y$, with $f'$ a unit there. The germ of $f'$ at $y$ is the image of
an element of $\mathcal O_{H,x}$, and its image in the function field is that of $f$, since $c$ restricts to the
pullback of $\tau$. So $f\in\mathcal O_{H,x}$, and in the same way $f^{-1}\in\mathcal O_{H,x}$.
\end{proof}
% Lean: GenExtendNormal.existsUnique_gen_extend: IsNormalScheme.of_pred with
% forall_isDomain_and_isIntegrallyClosed_stalk_of_smooth (H normal), GeometricallyIntegral.isIntegral_of_isLocallyNoetherian
% (H integral), exists_germToFunctionField_eq_of_coheight_eq_one (the normality step), and mem_range_of_notMem:
% coheight_apply_le_one (with Flat.generalizingMap), IsNormalScheme.isDiscreteValuationRing_stalk, exists_gen_germ
% (= GenExtend.existsUnique_gen_extend over the stalk) and mem_range_of_preimmersion; then
% unit_ext_of_isOpenImmersion for the gluing and the uniqueness.

\begin{lemma}[Alp\"oge]\label{cv-toroidal-abelian-dvr-cubical-iso}
Let $S$ be an integral, locally noetherian, normal scheme and let $j:U\to S$ be an open immersion with $U$
nonempty. Let $p:H\to S$ be a smooth morphism with geometrically integral fibres and with a section $e$, and let
$\mathcal L$ and $\mathcal M$ be line bundles on $H$ whose restrictions to $H\times_SU$ are isomorphic.
\begin{enumerate}
\item If $e^*\mathcal L\cong e^*\mathcal M$, then $\mathcal L\cong\mathcal M$.
\item Every point of $S$ has an open neighbourhood $N$ such that $\mathcal L$ and $\mathcal M$ are isomorphic over
  $p^{-1}(N)$.
\end{enumerate}
\end{lemma}
\leanref{Development/AbelianVarietyTheory/Proofs/DegenerationProofs/LineBundleIsoExtendOfNormalBaseStatement.lean}{GlobalTheory.T.lineBundleIso_extend_of_normal_base}
\begin{proof}
Let $\mathcal D=\mathcal M\otimes\mathcal L^{-1}$, with a generator $t$ over $H_U=H\times_SU$, and let
$\mathcal D'=\mathcal D\otimes p^*e^*\mathcal D^{-1}$. With $p_U$ and $e_U$ the base changes of $p$ and $e$, the
generator $t\otimes p_U^*e_U^*t^{-1}$ of $\mathcal D'$ over $H_U$ pulls back along $e_U$ to the canonical generator
of $e^*\mathcal D'=e^*\mathcal D\otimes e^*\mathcal D^{-1}$ over $U$. By
Lemma~\ref{cv-toroidal-abelian-dvr-cubical-normal}, with $\rho$ the canonical generator, $\mathcal D'$ has a
generator, so $\mathcal D\cong p^*e^*\mathcal D$. (1) If $e^*\mathcal L\cong e^*\mathcal M$, then $e^*\mathcal D$ is
trivial, hence so is $\mathcal D$. (2) The line bundle $e^*\mathcal D$ is trivial on an open neighbourhood $N$ of
the point, so $\mathcal D$ is trivial on $p^{-1}(N)$.
\end{proof}
% Lean: lbIso_extend_of_normal: lbIso_pull_iff_gen (D = M / L), exists_gen_diff (the generator of D tensor
% (p >> e)^*D^(-1), through GenExtendNormal.existsUnique_gen_extend), exists_gen_of_diff and exists_gen_of_chart.
% TripleExt.pic_ext restates (1) for classes in Pic(H).

\begin{lemma}[Alp\"oge]\label{cv-toroidal-abelian-dvr-cubical-extension}
Let $S$ be an integral, locally noetherian, normal scheme and let $j:U\to S$ be an open immersion with $U$
nonempty. Let $G$ be a commutative group scheme over $S$, smooth with geometrically connected fibres, let $A$ be a
commutative group scheme over $U$, and let $a:A\to G$ be a morphism over $j$ that identifies $A$ with $j^*G$ as
group schemes. Let $\mathcal L_U$ be a line bundle on $A$, trivialized along the zero section, with a cubical
structure $\theta_U$, and let $\mathcal M$ be a line bundle on $G$ with $a^*\mathcal M\cong\mathcal L_U$. Then there
are a line bundle $\mathcal L$ on $G$, trivialized along the zero section, with a cubical structure $\theta$, and an
isomorphism $w:\mathcal L_U\to a^*\mathcal L$ that respects the trivializations along the zero sections and carries
$\theta_U$ to the pullback of $\theta$. If $\mathcal L_U$ is symmetric, so is $\mathcal L$.
\end{lemma}
\leanref{Development/AbelianVarietyTheory/Proofs/DegenerationProofs/ExistsRigidifiedCubicalSymmetricExtensionNormalBStatement.lean}{GlobalTheory.T.exists_rigidified_cubical_symmetric_extension_normalBase}
\leanref{Development/AbelianVarietyTheory/Proofs/DegenerationProofs/ExistsRigidifiedCubicalSymmetricExtensionNormalB.lean}{GlobalTheory.NormalCubeExt.exists_cubical_extension}
\begin{proof}
For $n\ge1$ let $G^n$ be the $n$-fold fibre product of $G$ over $S$, with its zero section. It is smooth over $S$
with geometrically connected fibres, and its fibres over algebraically closed fields are regular
\cite[Tag 056S]{Sta}, hence integral \cite[Tag 033N]{Sta}; so $G^n$ has geometrically integral fibres, and it is
integral by the first paragraph of the proof of Lemma~\ref{cv-toroidal-abelian-dvr-cubical-normal}. Its base change
to $U$ is $A^n$, through $a$. Write $\pi:G\to S$ for the structure morphism and $0$ for zero sections. For a line
bundle $\mathcal L$ on $G$ let $\Theta(\mathcal L)$ be the line bundle on $G^3$ whose fibre at $(x_1,x_2,x_3)$ is
$\bigotimes_I\mathcal L(x_I)^{\otimes\delta_I}$, over the subsets $I\subset\{1,2,3\}$, with $x_I=\sum_{i\in I}x_i$
and $\delta_I=(-1)^{|I|+1}$.

The line bundle. Let $\mathcal L=\mathcal M\otimes\pi^*0^*\mathcal M^{-1}$. Its pullback along $0$ is
$0^*\mathcal M\otimes0^*\mathcal M^{-1}$, which is trivialized canonically. As $a$ commutes with the zero sections,
$a^*\mathcal L\cong a^*\mathcal M\otimes\pi_A^*0^*a^*\mathcal M^{-1}\cong\mathcal L_U\otimes\pi_A^*0^*
\mathcal L_U^{-1}$, with $\pi_A:A\to U$ the structure morphism, and $0^*\mathcal L_U$ is trivial; so
$\mathcal L_U\cong a^*\mathcal L$. Let $w$ be such an isomorphism. The trivialization of $0^*\mathcal L_U$ and the
pullback through $w$ of that of $0^*\mathcal L$ differ by a unit $c$ of $\Gamma(U,\mathcal O_U)$; replacing $w$ by
$\pi_A^*(c)w$, or by $\pi_A^*(c^{-1})w$, makes $w$ respect the trivializations.

The cubical structure. Through $w$, $\theta_U$ is a generator of $\Theta(\mathcal L)$ over $A^3=G^3\times_SU$. Its
pullback along the zero section of $A^3$ is $1$, as $\theta_U$ is $1$ when one of its arguments is the origin, and
this is the restriction to $U$ of the canonical generator of $0^*\Theta(\mathcal L)$, the tensor product of copies
of $0^*\mathcal L$ with exponents of sum zero. By Lemma~\ref{cv-toroidal-abelian-dvr-cubical-normal} for
$G^3\to S$ and its zero section, $\theta_U$ extends to a generator $\theta$ of $\Theta(\mathcal L)$. The conditions
for $\theta$ to be a cubical structure, that it is $1$ when one argument is the origin, that it is invariant under
exchanging the first two and the last two arguments, and the cocycle identity
$\theta(x_1+x_2,x_3,x_4)\theta(x_1,x_2,x_4)=\theta(x_1,x_2+x_3,x_4)\theta(x_2,x_3,x_4)$, are equalities between two
generators of one line bundle on $G^2$, $G^3$ or $G^4$. They hold over $U$, where they are the same conditions for
$\theta_U$, transported through $w$. Two generators of a line bundle on an integral scheme that agree on a nonempty
open subscheme are equal, as their ratio is a function that is $1$ there. So $\theta$ is a cubical structure, and
$w$ carries $\theta_U$ to the pullback of $\theta$.

The symmetry. Suppose that $\mathcal L_U$ is symmetric: there is an isomorphism $[-1]^*\mathcal L_U\cong
\mathcal L_U$ that respects the trivialization along the zero section, that is, a generator of $[-1]^*\mathcal L_U
\otimes\mathcal L_U^{-1}$ whose pullback along $0$ is $1$. Through $w$ it is a generator of $[-1]^*\mathcal L\otimes
\mathcal L^{-1}$ over $A$ whose pullback along $0$ is the restriction to $U$ of the canonical generator, and
Lemma~\ref{cv-toroidal-abelian-dvr-cubical-normal} for $G\to S$ extends it to $G$. This gives an isomorphism
$[-1]^*\mathcal L\cong\mathcal L$, which, adjusted by a unit of $\Gamma(S,\mathcal O_S)$ as above, respects the
trivialization along the zero section; so $\mathcal L$ is symmetric.
\end{proof}
% Lean: exists_cubical_extension: exists_rigidified (L' = M tensor (G.hom >> eSec)^*M^(-1), through classes in
% Pic), Unif.RigidifyIso.exists_rigidified_iso (exists_ratio, exists_scaled: the unit c), exists_cubical
% (exists_extension, which applies T.existsUnique_rigidified_generator_extend_normal to the generator transported
% by exists_up, with CubeExtension.formEq_origin for the canonical generator at the origin; then
% exists_cubical_of_open with corr_of_open and val_eq_of_open for the conditions face, swaps and cocycle), maps_iff
% for the compatibility of w with the cubical structures, and isSymmetric (isSymmetric_iff_exists_gen,
% exists_extension for G, exists_rigidified_iso).

\begin{lemma}[Alp\"oge]\label{cv-toroidal-abelian-dvr-cubical}
In Situation~\ref{cv-function-situation}, let $S$ be an integral, locally noetherian, normal scheme, let
$j:U\to S$ be an open immersion with $U$ nonempty, and let $G$ be a commutative semi-abelian scheme over $S$
with an action of $\mathcal O_E$. Let $G_U$ be an $\mathcal O_E$-semi-abelian scheme over $U$
(Definition~\ref{cv-model-toroidal-semiabelian}(2)) with line bundle $\mathcal L_U$, and let $a:G_U\to G$ be a
morphism over $j$ that identifies $G_U$ with $j^*G$ as group schemes and commutes with the actions of
$\mathcal O_E$. Let $\mathcal M$ be a line bundle on $G$ with $a^*\mathcal M\cong\mathcal L_U$. Then $G$ has a
line bundle $\mathcal L$, trivialized along the zero section, with a cubical structure, such that $G$ with its
action and $\mathcal L$ is an $\mathcal O_E$-semi-abelian scheme and $a$ is a morphism of
$\mathcal O_E$-semi-abelian schemes over $j$.
\end{lemma}
\leanref{Development/AbelianVarietyTheory/Proofs/DegenerationProofs/ExistsSabExtension.lean}{GlobalTheory.TripleExt.exists_sab_extension}
\begin{proof}
The semi-abelian scheme $G$ is smooth over $S$ with geometrically connected fibres. Let $\mathcal L_U$ have the
trivialization along the zero section and the cubical structure $\theta_U$ of $G_U$.
Lemma~\ref{cv-toroidal-abelian-dvr-cubical-extension}, applied to $G$, $A=G_U$, $a$, $\mathcal L_U$ and
$\mathcal M$, gives a line bundle $\mathcal L$ on $G$, trivialized along the zero section, with a cubical
structure $\theta$, and an isomorphism $w:\mathcal L_U\to a^*\mathcal L$ that respects the trivializations and
carries $\theta_U$ to the pullback of $\theta$; and $\mathcal L$ is symmetric, as $\mathcal L_U$ is. Then $a$
identifies $G_U$ with $j^*G$ and commutes with the actions, and through $w$ it respects the line bundles with
their trivializations and cubical structures.

It remains to show that complex conjugation is the Rosati involution for $\mathcal L$, and that $\mathcal L$ is
normalized: for $b\in\mathcal O_E$, $\Lambda_{\mathcal L}(by_1,y_2)\cong\Lambda_{\mathcal L}(y_1,\bar by_2)$ on
$G\times_SG$, where $y_1$ and $y_2$ are the two projections, and $\mathcal L\cong\Lambda_{\mathcal L}(y,w_0y)$ on
$G$, where $y$ is the identity. The morphisms $G\times_SG\to S$ and $G\to S$ are smooth with geometrically
integral fibres (as in the proof of Lemma~\ref{cv-toroidal-abelian-dvr-cubical-extension}), with zero sections,
and their base changes to $U$ are $G_U\times_UG_U$ and $G_U$, through $a$. Through $a$ and $w$, the pullbacks to
$G_U\times_UG_U$ and to $G_U$ of the two sides of each isomorphism are the corresponding line bundles for
$\mathcal L_U$, as $a$ commutes with the actions and the group laws; these are isomorphic, as $G_U$ is an
$\mathcal O_E$-semi-abelian scheme. Along the zero sections all of these line bundles are trivial: the actions fix
the zero section, so the pullback of $\Lambda_{\mathcal L}(by_1,y_2)$ or $\Lambda_{\mathcal L}(y_1,\bar by_2)$ is
a tensor product of copies of $0^*\mathcal L$ with exponents of sum zero, and that of $\mathcal L$ or of
$\Lambda_{\mathcal L}(y,w_0y)$ is a tensor power of $0^*\mathcal L$, which is trivial. So
Lemma~\ref{cv-toroidal-abelian-dvr-cubical-iso}(1) gives both isomorphisms, and $G$ with its action and $\mathcal L$
is an $\mathcal O_E$-semi-abelian scheme, with $a$ a morphism of $\mathcal O_E$-semi-abelian schemes over $j$.
\end{proof}
% Lean: exists_sab_extension: T.exists_rigidified_cubical_symmetric_extension_normalBase (Lemma
% cv-toroidal-abelian-dvr-cubical-extension) for the line bundle, its cubical structure and its symmetry; then, for
% adj on G x G and prim on G, pic_ext (T.lineBundleIso_extend_of_normal_base, Lemma cv-toroidal-abelian-dvr-cubical-iso)
% with geometricallyIntegral_of_smooth_of_geometricallyConnected_fibres, the classes cls_cube, pull_cubeCls_pre and
% pull_cubeCls_push (the comparison over U through t0.adj and t0.prim) and cubeCls_one (the zero section).

\begin{lemma}[Alp\"oge]\label{cv-toroidal-abelian-dvr-ample-square}
Let $O$ be a valuation ring whose fraction field $F$ is algebraically closed, and let $B$ be a commutative group
scheme over $O$ that is an abelian scheme. Let $\mathcal N$ be a line bundle on $B$ such that the line bundle
$\Theta(\mathcal N)$ on $B\times_OB\times_OB$ whose fibre at $(x_1,x_2,x_3)$ is
$\bigotimes_I\mathcal N(x_I)^{\otimes\delta_I}$ is trivial, where $I$ runs over the subsets of $\{1,2,3\}$,
$x_I=\sum_{i\in I}x_i$ and $\delta_I=(-1)^{|I|+1}$. Let $s$ be a section of $\mathcal N$ that does not vanish at
some point of the closed fibre of $B$. Then there are sections $u_0,\dots,u_m$ of $\mathcal N^2$ whose nonvanishing
loci cover $B$.
\end{lemma}
\leanref{Development/AbelianVarietyTheory/Proofs/ModuliProofs/ExistsSectionsGenerateSquareBundle.lean}{GlobalTheory.SquareGenerated.exists_secs_sq_generate}
\leanref{Development/AbelianVarietyTheory/Proofs/ModuliProofs/ValuationRingHenselianAndSmoothPointsLift.lean}{GlobalTheory.valuationRing_henselian_and_lift_of_isAlgClosed}
\begin{proof}
Let $k$ be the residue field of $O$. The ring $O$ is a normal local domain with separably closed fraction field,
so it is henselian \cite[Tag 0BSQ]{Sta}. The field $k$ is algebraically closed: a monic polynomial over $k$ lifts
to a monic polynomial over $O$, which has a root in $F$; the root lies in $O$, as $O$ is integrally closed, and its
image in $k$ is a root. So every point with values in $k$ of a scheme smooth over $O$ lifts to a point with values
in $O$ \cite[Tag 0H74]{Sta}.

Let $W$ be the union of the nonvanishing loci of the sections of $\mathcal N^2$, and suppose that $W\ne B$. As $B$
is proper over $O$, the image of the nonempty closed set $B\setminus W$ in $\operatorname{Spec}O$ is closed, so it
contains the closed point. The closed fibre $B_k$ is of finite type over $k$, so the nonempty closed subset
$B_k\setminus W$ contains a closed point of $B_k$, which is a point $x\in B(k)$ as $k$ is algebraically closed. The
fibre $B_k$ is integral, being smooth with geometrically connected fibres (as in the proof of
Lemma~\ref{cv-toroidal-abelian-dvr-cubical-extension}), and the nonvanishing locus $\Omega$ of $s$ meets it. So
the sets of points $a$ of $B_k$ with $a+x\in\Omega$ and with $-a+x\in\Omega$ are nonempty open subsets of $B_k$;
they meet, and as the points with values in $k$ are dense, there is $a\in B(k)$ with $a+x$ and $-a+x$ in $\Omega$.
Lift $a$ to $b\in B(O)$. Pulling back a generator of $\Theta(\mathcal N)$ along $(b,-b,y)$, where $y$ is the
identity of $B$, gives $t_b^*\mathcal N\otimes t_{-b}^*\mathcal N\cong\mathcal N^2\otimes\pi^*\mathcal P$, with
$\pi:B\to\operatorname{Spec}O$ the structure morphism and $\mathcal P=b^*\mathcal N\otimes(-b)^*\mathcal N\otimes
0^*\mathcal N^{-2}$, which is trivial as $O$ is local. Through this isomorphism $t_b^*s\otimes t_{-b}^*s$ is a
section of $\mathcal N^2$ that does not vanish at $x$, since $s$ does not vanish at $a+x$ and at $-a+x$. This
contradicts $x\notin W$. So $W=B$, and as $B$ is quasi-compact, finitely many sections suffice.
\end{proof}
% Lean: exists_secs_sq_generate: T.valuationRing_henselian_and_smooth_points_lift
% (henselianLocalRing_of_valuationRing_of_isAlgClosed, isAlgClosed_residueField_of_valuationRing_of_isAlgClosed),
% forall_geometricallyIntegral_of_smooth_of_geometricallyConnected, nonempty_inter_closedPoints and
% pointOfClosedPoint (the point x), NoBasePoint.exists_two_points (the point a), exists_sec_sq (the section of
% N^2 from Theta(N) trivial) and isCompact_univ.elim_finite_subcover.

\begin{lemma}[Alp\"oge]\label{cv-toroidal-abelian-dvr-ample-descent}
Let $R\to R'$ be a flat ring homomorphism, let $X$ be a quasi-compact and quasi-separated scheme over $R$, let
$X'=X\times_{\operatorname{Spec}R}\operatorname{Spec}R'$, and suppose that the projection $\pi:X'\to X$ is
surjective. Let $\mathcal M$ be a line bundle on $X$. If there are sections of $\pi^*\mathcal M$ whose nonvanishing
loci cover $X'$, then there are finitely many sections of $\mathcal M$ whose nonvanishing loci cover $X$.
\end{lemma}
\leanref{Development/AbelianVarietyTheory/Proofs/UniformizationProofs/ExistsSecPullNeZero.lean}{GlobalTheory.AmpleCrit.exists_cover_of_pull}
\begin{proof}
As $X'$ is quasi-compact, finitely many of the given sections $t_i$ suffice. By flat base change
\cite[Tag 02KH]{Sta}, $\Gamma(X',\pi^*\mathcal M)=\Gamma(X,\mathcal M)\otimes_RR'$, so each $t_i$ is a finite sum
$\sum_jc_{ij}\pi^*u_{ij}$ with $c_{ij}\in R'$ and $u_{ij}\in\Gamma(X,\mathcal M)$. Let $x\in X$, and let $x'\in X'$
lie over $x$. Some $t_i$ does not vanish at $x'$, so in a trivialization near $x'$ its germ at $x'$ is a unit of
the local ring; a sum of nonunits of a local ring is a nonunit, so some $c_{ij}\pi^*u_{ij}$ does not vanish at
$x'$, and then $u_{ij}$ does not vanish at $x$. So the sections $u_{ij}$ have nonvanishing loci covering $X$.
\end{proof}
% Lean: exists_cover_of_pull: FlatSec.exists_comb_of_flat (each section over X' is an R'-combination of pullbacks;
% through exists_finite_affine_refinement and exists_equiv_of_isPushout on affine opens) and
% IsLocalRing.exists_of_isUnit_sum.

\begin{lemma}[Alp\"oge]\label{cv-toroidal-abelian-dvr-ample-generating}
Let $V$ be a discrete valuation ring with fraction field $K$, let $A$ be a commutative group scheme over $V$ that
is an abelian scheme, and let $\mathcal L$ be a line bundle on $A$, trivialized along the zero section, with a
cubical structure. If the pullback of $\mathcal L$ to $A_K=A\times_V\operatorname{Spec}K$ is relatively ample,
then there are an integer $d\ge1$ and sections $u_0,\dots,u_m$ of $\mathcal L^d$ whose nonvanishing loci cover $A$.
\end{lemma}
\leanref{Development/AbelianVarietyTheory/Proofs/ModuliProofs/ExistsPrincipallyPolarizedCMExtensionDvr.lean}{GlobalTheory.GoodPol.exists_generating}
\begin{proof}
The scheme $A_K$ is nonempty and the pullback $\mathcal L_K$ of $\mathcal L$ to it is ample, as $\operatorname{Spec}K$
is affine \cite[Tag 01VH]{Sta}; so for some $a\ge1$ there is a nonzero section $\sigma$ of $\mathcal L_K^a$
\cite[Tag 01PS]{Sta}. The scheme $A$ is flat and of finite type over $V$, with integral fibres (as in the proof of
Lemma~\ref{cv-toroidal-abelian-dvr-cubical-extension}). By
Lemma~\ref{cv-toroidal-abelian-dvr-cubical-dvr}(2), there is an integer $n$ such that $\varpi^n\sigma$, for a
uniformizer $\varpi$, extends to a section $t$ of $\mathcal L^a$ that does not vanish at some point $x$ of the
special fibre.

Let $F$ be an algebraic closure of $K$, and let $O$ be a valuation ring of $F$ that dominates $V$
\cite[Tag 00IA]{Sta}. The homomorphism $V\to O$ is injective and local; $O$ is torsion free over $V$, hence flat
\cite[Tag 0539]{Sta}, and $\operatorname{Spec}O\to\operatorname{Spec}V$ is surjective, as it maps the closed point
to the closed point and the generic point to the generic point. Let $B=A\times_V\operatorname{Spec}O$, an abelian
scheme over $O$, with projection $\pi:B\to A$, which is surjective \cite[Tag 01S1]{Sta}. As $V\to O$ is local,
there is a point of $B$ over $x$ and over the closed point of $\operatorname{Spec}O$ \cite[Tag 01JT]{Sta}, and
$\pi^*t$ does not vanish at it. The cubical structure of $\mathcal L$
pulls back to a generator of $\Theta(\pi^*\mathcal L)$, and $\Theta(\pi^*\mathcal L^a)=\Theta(\pi^*\mathcal L)^{\otimes a}$,
so $\Theta(\pi^*\mathcal L^a)$ is trivial. By Lemma~\ref{cv-toroidal-abelian-dvr-ample-square}, there are sections
of $\pi^*\mathcal L^{2a}$ whose nonvanishing loci cover $B$, and by
Lemma~\ref{cv-toroidal-abelian-dvr-ample-descent}, applied to $V\to O$, there are such sections of
$\mathcal L^{2a}$ over $A$. Let $d=2a$.
\end{proof}
% Lean: exists_generating: Unif.relAmple_iff_veryAmpleBy over K (the nonzero section sigma),
% DvrModel.existsUnique_zpow_mul_sec (Lemma cv-toroidal-abelian-dvr-cubical-dvr(2)), the flatness and surjectivity
% of V -> O, Out.CubePf.cubePull and triv_theta_pw (Theta of the pullback of L^d is trivial),
% SecPullIso.section_pull, SquareGenerated.exists_secs_sq_generate and AmpleCrit.exists_cover_of_pull. The valuation
% ring O and its injective local homomorphism g : V -> O are hypotheses of the Lean statement; the Lean lemma gives
% d = 2a.

\begin{lemma}[Alp\"oge]\label{cv-toroidal-abelian-dvr-ample-effective-finite}
Let $K$ be an algebraically closed field, let $B$ be a commutative group scheme over $K$ that is an abelian scheme,
and let $n$ be a prime that is nonzero in $K$. Let $C$ be a subgroup of $B(K)$ and $W$ a closed subset of $B$ such
that a point $x\in B(K)$ lies in $C$ if and only if its image lies in $W$. If $0$ is the only element of $C$ killed by
$n$, then $C$ is finite.
\end{lemma}
\leanref{Development/AbelianVarietyTheory/Proofs/UnitaryExistence/ExistsLiftOfRangeSubset.lean}{GlobalTheory.AbelianSubgroupFinite.finite_of_no_torsion}
\begin{proof}
Let $N$ be $W$ with its reduced structure. A morphism to $B$ from a reduced scheme $T$ locally of finite type over $K$
factors through $N$ if it maps the points of $T$ with values in $K$ into $C$, as these points are dense in $T$. Applied
to the origin, to the group law on $N\times_KN$ and to the inverse on $N$, this makes $N$ a closed subgroup scheme of
$B$ whose points with values in $K$ are those of $C$. It is reduced and locally of finite type over $K$, and $K$ is
perfect, so $N$ is geometrically reduced over $K$ \cite[Tag 020I]{Sta}, and its smooth locus $U$ is open and dense in $N$
\cite[Tag 056V]{Sta}.

The scheme $N$ is smooth over $K$. Indeed, suppose that $U\ne N$. As $N$ is locally of finite type over $K$, it is
Jacobson \cite[Tag 02J6]{Sta}, so the nonempty closed set $N\setminus U$ contains a closed point $x$ of $N$, and the
open set $U$, which is nonempty as it is dense in $N$ and $N$ contains the origin, contains a closed point $y$. Their
residue fields are finite over $K$ \cite[Tag 01TB]{Sta}, hence
equal to $K$, so $x$ and $y$ are points with values in $K$. The translation $z\mapsto z-x+y$ is an automorphism of $N$
over $K$, so the inverse image of $U$ under it is $U$; as it maps $x$ to $y\in U$, $x$ lies in $U$, a contradiction.

The connected component $N_0$ of the origin is open and closed in $N$, as
$N$ is noetherian; it is a connected smooth closed subgroup scheme of $B$, proper over $K$, so an abelian variety, of
some dimension $g$. Its points killed by $n$ lie in $C$, so there is only one; as there are $n^{2g}$ of them
\cite[\S6, Application~3]{Mum70}, $g=0$. Then $N_0$ is \'etale over $K$, so the origin is an open point of $N$, and by
translation so is every point with values in $K$. As $N$ is quasi-compact, there are finitely many, and $C$ is finite.
\end{proof}
% Lean: AbelianSubgroupFinite.finite_of_no_torsion, with the count of n-torsion points (the law Out.Laws.mulN) as
% hypothesis: ReducedSubgroup.exists_smooth_closed_subgroup (N, smooth, with range W), exists_neutralComponent (N_0
% open and closed, geometrically connected, smooth of relative dimension g, and N(K) finite when g = 0), and
% exists_lift_of_range_subset. The Lean proves these facts on algebraic groups itself (about 360 lines). In
% exists_smooth_closed_subgroup, N is the subscheme of vanishingIdeal Z (isReduced_subscheme_vanishingIdeal), the
% group law comes from hlift and exists_grpObj_of_mono, and the smoothness is geometricallyReduced_of_perfectField
% (Stacks 020I in the text) with Mathlib's AlgebraicGeometry.smooth_of_grpObj. That lemma descends smoothness from
% the base change to an algebraic closure Omega of K (here K -> Omega is an isomorphism) and there applies
% smooth_of_grpObj_of_isAlgClosed, the second paragraph of the proof: LocallyOfFiniteType.jacobsonSpace,
% nonempty_inter_closedPoints (x and y), Scheme.Hom.dense_smoothLocus_of_perfectField (Stacks 056V in the text),
% pointEquivClosedPoint, the translation alpha = (GrpObj.mulRight x)^(-1) then GrpObj.mulRight y, and
% Scheme.Hom.preimage_smoothLocus_eq.

\begin{lemma}[Alp\"oge]\label{cv-toroidal-abelian-dvr-ample-effective-rigid}
Let $K$ be an algebraically closed field, let $B$ be a commutative group scheme over $K$ that is an abelian scheme,
let $Z$ be a scheme separated and locally of finite type over $K$, and let $\varphi:B\to Z$ be a morphism over $K$. If
only finitely many $x\in B(K)$ satisfy $\varphi\circ t_x=\varphi$, then $\varphi$ is locally quasi-finite.
\end{lemma}
\leanref{Development/AbelianVarietyTheory/Proofs/UniformizationProofs/LocallyQuasiFiniteOfFiniteStabiliser.lean}{GlobalTheory.FiniteStabiliser.locallyQuasiFinite_of_finite_stabiliser}
\leanref{Development/AbelianVarietyTheory/Proofs/UniformizationProofs/LocallyQuasiFiniteOfFiniteStabiliser.lean}{GlobalTheory.FiniteStabiliser.finite_closedPoints_of_contracted}
\leanref{Development/AbelianVarietyTheory/Proofs/UniformizationProofs/LocallyQuasiFiniteOfFiniteStabiliser.lean}{GlobalTheory.FiniteStabiliser.slice_eq}
\begin{proof}
The scheme $B$ is integral, noetherian and proper over $K$, so $\varphi$ is proper. It suffices to show that for every
closed point $y$ of $B$ the fibre $\varphi^{-1}(\varphi(y))$, a closed subset, has finitely many closed points. It is a
finite union of irreducible closed subsets; let $S$ be one of them, mapped by $\varphi$ to the point $z$, and suppose
that $S$ has a closed point. Let $F$ be $S$ with its reduced structure, an integral scheme proper over $K$, and let
$u_0,u\in F(K)$. The morphism $g:F\times_KB\to Z$, $(f,b)\mapsto\varphi(f+b)$, maps $F\times\{0\}$ to $z$. By the
rigidity lemma \cite[\S4]{Mum70}, $g(u,b)=g(u_0,b)$ for every $b$, that is, $\varphi\circ t_u=\varphi\circ t_{u_0}$, so
$\varphi\circ t_{u-u_0}=\varphi$. So the closed points of $S$, which are points with values in $K$, lie in the
translate by $u_0$ of the finite set of the statement.
\end{proof}
% Lean: locallyQuasiFinite_of_finite_stabiliser: locallyQuasiFinite_of_finite_fibres_closedPoints,
% NoetherianSpace.exists_finite_set_isClosed_irreducible and finite_closedPoints_of_contracted, whose step
% slice_eq is the rigidity lemma (through an affine open W around z, the closed image of the complement of g^{-1}(W),
% and apply_eq_of_range_subset_affine: a morphism from a proper integral K-scheme into an affine scheme is constant).
% The text cites Mumford for the rigidity lemma.

\begin{lemma}[Alp\"oge]\label{cv-toroidal-abelian-dvr-ample-effective-quasifinite}
Let $K$ be an algebraically closed field, let $B$ be a commutative group scheme over $K$ that is an abelian scheme,
and let $\mathcal L$ be a line bundle on $B$, trivialized along the zero section, with a cubical structure. Let $r\ge1$
and let $n$ be a prime that is nonzero in $K$ and does not divide $r$, such that every $x\in B(K)$ with $nx=0$ for
which $t_x^*\mathcal L\otimes\mathcal L^{-1}$ is trivial is zero. Let $Z$ be a scheme separated and locally of finite
type over $K$, let $\varphi:B\to Z$ be a morphism over $K$, and let $\mathcal H$ be a line bundle on $Z$ with
$\mathcal L^r\cong\varphi^*\mathcal H$. Then $\varphi$ is locally quasi-finite.
\end{lemma}
\leanref{Development/AbelianVarietyTheory/Proofs/UniformizationProofs/LocallyQuasiFiniteOfCubicalPowIsoPullStatement.lean}{GlobalTheory.T.locallyQuasiFinite_of_cubical_pow_iso_pull}
\leanref{Development/AbelianVarietyTheory/Proofs/UniformizationProofs/LocallyQuasiFiniteOfCubicalPowIsoPull.lean}{GlobalTheory.CubicalQuasiFinite.locallyQuasiFinite_of_pow_iso_pull}
\leanref{Development/AbelianVarietyTheory/Proofs/UniformizationProofs/ExistsClosedStabiliser.lean}{GlobalTheory.Stabiliser.exists_closed_stabiliser}
\begin{proof}
Let $C$ be the set of $x\in B(K)$ with $\varphi\circ t_x=\varphi$, a subgroup of $B(K)$. As $B$ is reduced and its
points with values in $K$ are dense, $x\in C$ if and only if $\varphi(x+z)=\varphi(z)$ for every $z\in B(K)$. For each
$z$ these $x$ are the points of the closed subscheme of $B$ on which $b\mapsto\varphi(b+z)$ agrees with the constant
$\varphi(z)$, closed as $Z$ is separated; so $x\in C$ if and only if $x$ lies in the intersection $W$ of these closed
subsets.

Let $x\in C$ with $nx=0$. Then $t_x^*\mathcal L^r\otimes\mathcal L^{-r}\cong t_x^*\varphi^*\mathcal H\otimes
\varphi^*\mathcal H^{-1}$ is trivial, as $\varphi\circ t_x=\varphi$. By the theorem of the square
\cite[\S6, Corollary~4]{Mum70}, $t_x^*\mathcal L^r\otimes\mathcal L^{-r}\cong t_{rx}^*\mathcal L\otimes\mathcal L^{-1}$,
so the latter is trivial; as $n(rx)=0$, the hypothesis gives $rx=0$. As $nx=0$ and $n$ is a prime not dividing $r$,
$x=0$. By Lemma~\ref{cv-toroidal-abelian-dvr-ample-effective-finite}, $C$ is finite, and by
Lemma~\ref{cv-toroidal-abelian-dvr-ample-effective-rigid}, $\varphi$ is locally quasi-finite.
\end{proof}
% Lean: CubicalQuasiFinite.locallyQuasiFinite_of_pow_iso_pull: Stabiliser.exists_closed_stabiliser (C = stab phi and
% W, through comp_eq_iff_range_subset and ComplexPoints.hom_ext_of_points), triv_translDiff_of_fixes,
% Square.square_of_cubical (the theorem of the square from the cubical structure, through triv_square and triv_pow;
% the text cites Mumford), pow_gcd_eq_one, AbelianSubgroupFinite.finite_of_no_torsion and
% FiniteStabiliser.locallyQuasiFinite_of_finite_stabiliser.

\begin{lemma}[Alp\"oge]\label{cv-toroidal-abelian-dvr-ample-effective-cover}
Let $K$ be an algebraically closed field, let $B$ be a commutative group scheme over $K$ that is an abelian scheme,
and let $\mathcal L$ be a line bundle on $B$, trivialized along the zero section, with a cubical structure. Let $m\ge0$
and let $s$ be a nonzero section of $\mathcal L^m$. Then there are sections $t_0,\dots,t_N$ of $\mathcal L^{2m}$ whose
nonvanishing loci cover $B$.
\end{lemma}
\leanref{Development/AbelianVarietyTheory/Proofs/UniformizationProofs/ExistsCoverNonzeroLocus.lean}{GlobalTheory.CubeFree.exists_cover_nonzeroLocus}
\begin{proof}
The scheme $B$ is integral, being smooth with geometrically connected fibres over $K$ (as in the proof of
Lemma~\ref{cv-toroidal-abelian-dvr-cubical-extension}). As $B$ is reduced and $s\ne0$, the nonvanishing locus $D$ of
$s$ is a nonempty open subset of $B$. As $B$ is locally of finite type over $K$, it is Jacobson \cite[Tag 02J6]{Sta},
and the residue field of each of its closed points is finite over $K$ \cite[Tag 01TB]{Sta}, hence equal to $K$; so
every nonempty locally closed subset of $B$ contains a point with values in $K$.

First, every $x\in B(K)$ lies in the nonvanishing locus of a section of $\mathcal L^{2m}$. Indeed, the sets of points
$a$ of $B$ with $a+x\in D$ and with $-a+x\in D$ are the images of $D$ under the automorphisms $z\mapsto z-x$ and
$z\mapsto x-z$ of the scheme $B$, so they are nonempty and open. They meet, as $B$ is irreducible, so there is
$a\in B(K)$ with $a+x$ and $-a+x$ in $D$. Then $t_a^*s$ and $t_{-a}^*s$ are sections of $t_a^*\mathcal L^m$ and
$t_{-a}^*\mathcal L^m$ that do not vanish at $x$. The cubical structure is a generator of the line bundle
$\Theta(\mathcal L)$ on $B\times_KB\times_KB$ whose fibre at $(x_1,x_2,x_3)$ is
$\bigotimes_I\mathcal L(x_I)^{\otimes\delta_I}$, where $I$ runs over the subsets of $\{1,2,3\}$, $x_I=\sum_{i\in I}x_i$
and $\delta_I=(-1)^{|I|+1}$. Its pullback along $(a,-a,y)$, where $y$ is the identity of $B$, is a generator of
$\mathcal L^2\otimes t_a^*\mathcal L^{-1}\otimes t_{-a}^*\mathcal L^{-1}$ tensored with the pullback of a line bundle on
$\operatorname{Spec}K$, which is trivial. This gives a generator $\rho$ of $\mathcal L^2\otimes t_a^*\mathcal L^{-1}
\otimes t_{-a}^*\mathcal L^{-1}$ (the theorem of the square). Then $t_a^*s\otimes t_{-a}^*s\otimes\rho^m$ is a section of
$\mathcal L^{2m}$ that does not vanish at $x$.

Let $W$ be the union of the nonvanishing loci of the sections of $\mathcal L^{2m}$, and suppose that $W\ne B$. The
nonempty closed set $B\setminus W$ contains a point $z\in B(K)$, and by the first point $z\in W$, a contradiction. So
$W=B$, and as $B$ is quasi-compact, finitely many sections suffice.
\end{proof}
% Lean: CubeFree.exists_cover_nonzeroLocus: forall_geometricallyIntegral_of_smooth_of_geometricallyConnected (B
% integral), basicOpen_eq_bot_iff for a nonzero chart function s.f j (D nonempty), then key:
% NoBasePoint.exists_two_points (the point a, through exists_section_mem), unitAt_of_pw (t_a^* s and t_{-a}^* s),
% Square.triv_square c a a^-1 with unitAt_of_triv (rho^m; triv_square pulls back the cubical generator along
% (a, -a, y) and uses triv_const for the line bundles from Spec K), exists_sec_of_unitAt (the section of L^(2m)); then
% top: LocallyOfFiniteType.jacobsonSpace, nonempty_inter_closedPoints and pointOfClosedPoint (the point z); and
% isCompact_univ.elim_finite_subcover. The Lean's exponent is k.

\begin{lemma}[Alp\"oge]\label{cv-toroidal-abelian-dvr-ample-effective}
Let $k$ be a field and let $B$ be a commutative group scheme over $k$ that is an abelian scheme. Let $\mathcal L$ be
a line bundle on $B$, trivialized along the zero section, with a cubical structure, let $m\ge1$, and let $s$ be a
nonzero section of $\mathcal L^m$. Let $\Omega$ be an algebraically closed field with a homomorphism
$k\to\Omega$, and let $n$ be a prime that is nonzero in $\Omega$ and does not divide $2m$. Suppose that every point
$x\in B(\Omega)$ with $nx=0$ for which $t_x^*\mathcal L_\Omega\otimes\mathcal L_\Omega^{-1}$ is trivial on
$B_\Omega$ is zero. Then $\mathcal L$ is ample.
\end{lemma}
\leanref{Development/AbelianVarietyTheory/Proofs/UniformizationProofs/ExistsSecPullNeZero.lean}{GlobalTheory.AmpleCrit.relAmple_of_effective}
\begin{proof}
Let $B_\Omega=B\times_k\operatorname{Spec}\Omega$, with projection $\pi:B_\Omega\to B$, which is flat and surjective,
and let $\mathcal L_\Omega=\pi^*\mathcal L$, with the pullbacks of the trivialization and of the cubical structure.

Step 1. The pullback $\pi^*s$ is a nonzero section of $\mathcal L_\Omega^m$, as $\pi$ is flat and surjective.

Step 2. There are sections of $\mathcal L_\Omega^{2m}$ whose nonvanishing loci cover $B_\Omega$, by
Lemma~\ref{cv-toroidal-abelian-dvr-ample-effective-cover} applied to the abelian scheme $B_\Omega$ over $\Omega$, to
$\mathcal L_\Omega$ with the pullbacks of the trivialization and of the cubical structure, and to the section $\pi^*s$
of $\mathcal L_\Omega^m$.

Step 3. By Lemma~\ref{cv-toroidal-abelian-dvr-ample-descent}, applied to $k\to\Omega$, there are sections
$t_0,\dots,t_M$ of $\mathcal L^{2m}$ whose nonvanishing loci $B_{t_i}$ cover $B$.

Step 4. These sections define a morphism $\varphi:B\to P=\mathbf P^M_k$ over $k$ with
$\varphi^*\mathcal O(1)\cong\mathcal L^{2m}$ and $B_{t_i}=\varphi^{-1}(D_+(T_i))$ \cite[Tag 01NE]{Sta}. The morphism
$\varphi$ is proper, as $B$ is proper over $k$.

Step 5. The base change $\varphi_\Omega:B_\Omega\to P_\Omega$ is locally quasi-finite. Indeed,
$\mathcal L_\Omega^{2m}\cong\varphi_\Omega^*\mathcal O(1)$, and $P_\Omega$ is separated and of finite type over
$\Omega$. The points of $B_\Omega$ with values in $\Omega$ are the points $x\in B(\Omega)$, and
$t_x^*\mathcal L_\Omega\otimes\mathcal L_\Omega^{-1}$ is the same line bundle for both; so by hypothesis every such $x$
with $nx=0$ and $t_x^*\mathcal L_\Omega\otimes\mathcal L_\Omega^{-1}$ trivial is zero. As $n$ does not divide $2m$,
Lemma~\ref{cv-toroidal-abelian-dvr-ample-effective-quasifinite}, with $r=2m$, applies.

Step 6. The morphism $\varphi$ is finite. Every fibre of $\varphi$ is the image of a fibre of $\varphi_\Omega$, as
$P_\Omega\to P$ is surjective, and the fibres of $\varphi_\Omega$ are finite, as $\varphi_\Omega$ is locally
quasi-finite and quasi-compact. So $\varphi$ is proper with finite fibres, hence finite \cite[Tag 02OG]{Sta}.

Step 7. The nonvanishing loci $B_{t_i}=\varphi^{-1}(D_+(T_i))$ are affine, as $\varphi$ is finite, and they cover $B$.
So $\mathcal L^{2m}$ is ample \cite[Tag 01PS]{Sta}, and so is $\mathcal L$ \cite[Tag 01PT]{Sta}.
\end{proof}
% Lean: AmpleCrit.relAmple_of_effective: exists_sec_pull_ne_zero (Step 1), CubeFree.exists_cover_nonzeroLocus (Step 2,
% Lemma cv-toroidal-abelian-dvr-ample-effective-cover, applied to B' with (cubeMapOf h).rigLB and the pulled back
% cubical structure c' from CubePf.cubePull), exists_cover_of_pull (Step 3),
% exists_map (Step 4, through ProjQuot.exists_veryAmpleBy_of_ideal and exists_hom_of_veryAmpleBy),
% T.locallyQuasiFinite_of_cubical_pow_iso_pull with pushHom and triv_translDiff_push (Step 5),
% Unif.locallyQuasiFinite_of_isPullback_surjective and IsFinite.of_isProper_of_locallyQuasiFinite (Step 6),
% veryAmpleBy_pw_of_affine_loci and Unif.relAmple_iff_veryAmpleBy (Step 7).
% The Lean takes as hypothesis the count of n-torsion points over algebraically closed fields (n^(2g) when n is
% nonzero in the field), the law Out.Laws.mulN, which is [Mum70, Section 6, Application 3]. Size: about 5,400 lines.

\begin{definition}\label{cv-toroidal-abelian-dvr-ample-torsion-weil}
Let $S$ be a scheme, let $A$ be a commutative group scheme over $S$, and let $\mathcal L$ be a line bundle on $A$,
trivialized along the zero section, with a cubical structure $\theta$, a generator of $\Theta(\mathcal L)$ on
$A\times_SA\times_SA$. For a scheme $T$ over $S$ and $a,m,x\in A(T)$, write $\theta(a,m,x)$ for the pullback of
$\theta$ along $(a,m,x)$, a generator of the line bundle $\Theta(\mathcal L)(a,m,x)$ on $T$, and for $k\ge0$ let
$c_k(a;m,x)=\prod_{i=0}^{k-1}\theta(a+im,m,x)$. For $x,y\in A(T)$ let $\Lambda_{\mathcal L}(x,y)=(x+y)^*\mathcal L
\otimes x^*\mathcal L^{-1}\otimes y^*\mathcal L^{-1}\otimes0^*\mathcal L$. Let $N\ge1$ and let $x,y\in A(T)$ with
$Nx=Ny=0$. The tensor product of the line bundles $\Theta(\mathcal L)(ix,x,y)$ for $0\le i<N$ is that of
$\Lambda_{\mathcal L}(x,y)^{-N}$ and of $(Nx+y)^*\mathcal L\otimes(Nx)^*\mathcal L^{-1}\otimes y^*\mathcal L^{-1}\otimes
0^*\mathcal L$, which is trivial as $Nx=0$; so $c_N(x,y)=c_N(0;x,y)$ is a generator of $\Lambda_{\mathcal L}(x,y)^{-N}$.
As $\Lambda_{\mathcal L}(x,y)=\Lambda_{\mathcal L}(y,x)$, there is a unique unit $e_N(x,y)$ of $\Gamma(T,\mathcal O_T)$
with $e_N(x,y)c_N(x,y)=c_N(y,x)$, the \emph{cubical Weil pairing} of $x$ and $y$.
\end{definition}
\leanref{Development/AbelianVarietyTheory/Definitions/Uniformization/PointEqOfSubsingleton.lean}{GlobalTheory.UU.CubWeil}
\leanref{Development/AbelianVarietyTheory/Proofs/ThetaShiftProofs/WeilIsGrpPullbackIff.lean}{GlobalTheory.WeilCube.exists_gen_torsion}
\leanref{Development/AbelianVarietyTheory/Definitions/Uniformization/PointEqOfSubsingleton.lean}{GlobalTheory.UU.existsUnique_cubWeil}

\begin{lemma}[Alp\"oge]\label{cv-toroidal-abelian-dvr-ample-torsion-basic}
Let $S$ be a scheme, let $A$ be a commutative group scheme over $S$, let $\mathcal L$ be a line bundle on $A$,
trivialized along the zero section, with a cubical structure $\theta$, let $T$ be a scheme over $S$, and let $N\ge1$.
Let $e_N$ be the cubical Weil pairing (Definition~\ref{cv-toroidal-abelian-dvr-ample-torsion-weil}), and let
$x,x',y\in A(T)$ be killed by $N$.
\begin{enumerate}
\item $e_N(x+x',y)=e_N(x,y)e_N(x',y)$.
\item $e_N(y,x)=e_N(x,y)^{-1}$ and $e_N(x,y)^N=1$.
\item For a morphism $\rho:T'\to T$ over $S$, $e_N(x\rho,y\rho)$ is the pullback of $e_N(x,y)$.
\end{enumerate}
\end{lemma}
\leanref{Development/AbelianVarietyTheory/Proofs/UniformizationProofs/CubWeilOneLeft.lean}{GlobalTheory.UU.cubWeil_basic}
\begin{proof}
The cubical structure is $1$ when one argument is $0$, is invariant under exchanging the first two and the last two
arguments, and satisfies $\theta(x_1+x_2,x_3,y)\theta(x_1,x_2,y)=\theta(x_1,x_2+x_3,y)\theta(x_2,x_3,y)$. From these
identities, by induction on $k$, for points $a,m,x,b,x'$ of $A$ with values in $T$:
\begin{align*}
 c_k(0;x+x',y)\,\theta(x,x',y)^k&=\theta(kx,kx',y)\,c_k(0;x,y)\,c_k(0;x',y),\tag{i}\\
 c_k(a;m,x+b)\,\theta(a,x,b)\,\theta(m,x,b)^k&=\theta(a+km,x,b)\,c_k(a;m,x)\,c_k(a;m,b).\tag{ii}
\end{align*}

(1) By (i) with $k=N$, as $Nx=Nx'=0$, $c_N(x+x',y)\theta(x,x',y)^N=c_N(x,y)c_N(x',y)$. By (ii) with $a=0$, $m=y$ and
$k=N$, as $Ny=0$, $c_N(y,x+x')\theta(y,x,x')^N=c_N(y,x)c_N(y,x')$. As $\theta(x,x',y)=\theta(y,x,x')$, dividing the
second identity by the first gives (1).

(2) The first assertion is the definition. By (1), $e_N(0,y)=e_N(0,y)^2$, so $e_N(0,y)=1$, and
$e_N(x,y)^N=e_N(Nx,y)=e_N(0,y)=1$.

(3) The generators $c_N(x,y)$ and $c_N(y,x)$ pull back along $\rho$ to those for $x\rho$ and $y\rho$.
\end{proof}
% Lean: cubWeil_basic: existsUnique_cubWeil (WeilCube.exists_isRatio, CubWeil.unique), CubWeil.mul_left (cubN_mul_left,
% which is (i), and IsCub.N_add, which is (ii)), CubWeil.symm, CubWeil.pow_eq_one (cubWeil_one_left, pow_left) and
% CubWeil.comp. The last conjunct of the Lean statement (pairings with the image of a homomorphism carrying a
% bi-trivialization) is not used here.

\begin{lemma}[Alp\"oge]\label{cv-toroidal-abelian-dvr-ample-torsion-kernel}
Let $S$ be a scheme, let $B$ be a commutative group scheme over $S$ that is an abelian scheme, let $\mathcal L$ be a
line bundle on $B$, trivialized along the zero section, with a cubical structure $\theta$, and let $e_N$ be its
cubical Weil pairing (Definition~\ref{cv-toroidal-abelian-dvr-ample-torsion-weil}). Let $N\ge1$, let $T$ be a scheme
over $S$, and let $x,y\in B(T)$ with $Nx=Ny=0$. If $t_x^*\mathcal L_T\otimes\mathcal L_T^{-1}$ is trivial on
$B_T=T\times_SB$, then $e_N(x,y)=1$.
\end{lemma}
\leanref{Development/AbelianVarietyTheory/Proofs/UniformizationProofs/CubWeilOfGen.lean}{GlobalTheory.UU.KerPair.cubWeil_of_gen}
\begin{proof}
Let $u$ be a generator of $t_x^*\mathcal L_T\otimes\mathcal L_T^{-1}$, and for $b\in B(T)$ let $U(b)$ be its pullback
along $(\mathrm{id},b):T\to B_T$, a generator of $(x+b)^*\mathcal L\otimes b^*\mathcal L^{-1}$. First, for
$a,b\in B(T)$,
\[
 U(a+b)\,U(0)=\theta(x,a,b)\,U(a)\,U(b)\qquad\text{and}\qquad c_N(x,b)\,U(b)^N=U(0)^N.
\]
For the first, replace $b$ by the tautological point $q$ of $B_T$ and $x$, $a$ by their pullbacks to $B_T$: the two
sides become generators of one line bundle on $B_T$. Their ratio is a function on $B_T$, which comes from $T$
\cite[Tag 0E0L]{Sta}; at the zero section $q=0$ the two sides agree, as $\theta(x,a,0)=1$; so they agree everywhere,
and pulling back along $(\mathrm{id},b)$ gives the identity. The second is proved in the same way, as $c_N(x,0)=1$.

By induction on $n$, the first identity with $a=ny$ and $b=y$, and $\theta(x,ny,y)=\theta(ny,y,x)$, give
$U(ny)U(0)^n=U(0)\,c_n(0;y,x)\,U(y)^n$. For $n=N$, as $Ny=0$, this is $U(0)^N=c_N(y,x)U(y)^N$. With the second identity
for $b=y$, $c_N(x,y)U(y)^N=c_N(y,x)U(y)^N$, so $c_N(x,y)=c_N(y,x)$ and $e_N(x,y)=1$.
\end{proof}
% Lean: KerPair.cubWeil_of_gen (with GlobalFnsSt): val_mul_pt and val_pow_pt (the two identities, through gen_unique'
% and WeilCube.surj_fst, the surjectivity of Gamma(T) -> Gamma(T x B) from GlobalFnsSt), and the induction b.

\begin{lemma}[Alp\"oge]\label{cv-toroidal-abelian-dvr-ample-torsion-descent}
Let $S$ be a scheme, let $G$ be a commutative group scheme over $S$ that is smooth over $S$, let $K$ be an
algebraically closed field with a morphism $\operatorname{Spec}K\to S$, and let $N\ge1$ be an integer with $N\ne0$ in
$K$ such that multiplication by $N$ on $G$ is finite, flat and surjective. Let $G_K=G\times_S\operatorname{Spec}K$, let
$m$ be multiplication by $N$ on $G_K$, let $V$ be an open subscheme of $G_K$, and let $f$ be a function on $m^{-1}(V)$
such that $\tau_y^*f=f$ for every $y\in G(K)$ with $Ny=0$, where $\tau_y$ is the translation by $y$, which preserves
$m^{-1}(V)$. Then $f=m^*g$ for a unique function $g$ on $V$.
\end{lemma}
\leanref{Development/AbelianVarietyTheory/Proofs/UnramifiedDepth/TorsionInvariantDescendsFibre.lean}{GlobalTheory.torsion_invariant_descends_fibre}
\leanref{Development/AbelianVarietyTheory/Proofs/AbelianSchemeLawProofs/ExistsUniqueDescendOfTorsionInvariant.lean}{GlobalTheory.existsUnique_descend_of_torsion_invariant}
\begin{proof}
The morphism $m$ is the base change of multiplication by $N$ on $G$, so it is finite, flat and surjective. It is
\'etale, as $G_K$ is smooth over $K$ and $N\ne0$ in $K$, so that its differential is multiplication by $N$ on the
tangent spaces. Its restriction $m^{-1}(V)\to V$ is flat, quasi-compact and surjective, so it is an effective
epimorphism of schemes \cite[Tag 023Q]{Sta}. Regard $f$ as a morphism from $m^{-1}(V)$ to the affine line. It
suffices to show that $f\circ a=f\circ b$ for any two morphisms $a,b:Z\to m^{-1}(V)$ with $ma=mb$. The difference
$d=a-b$ is killed by $N$, so it factors through the kernel of $m$, which is \'etale over the algebraically closed
field $K$. Every point of that kernel lies in the image of a section over $\operatorname{Spec}K$, which is an open
immersion, as a section of an \'etale morphism. So $Z$ is covered by open subschemes on each of which $d$ is a point
$y\in G(K)$ with $Ny=0$, and there $a=\tau_y\circ b$ and $f\circ a=f\circ b$.
\end{proof}
% Lean: torsion_invariant_descends_fibre transports the statement to A = G x_S Spec K (TorsionLift.ptsEquiv) and
% applies existsUnique_descend_of_torsion_invariant: GrpTangent.etale_id_pow (m etale), exists_local_rational_torsion
% (exists_section_of_etale and isOpenImmersion_of_section_of_etale) and existsUnique_appTop_eq_of_flat (EffectiveEpi
% for a flat, quasi-compact, surjective morphism).

\begin{lemma}[Alp\"oge]\label{cv-toroidal-abelian-dvr-ample-torsion-nondegenerate}
Let $S$ be a scheme, let $A$ be a commutative group scheme over $S$ that is an abelian scheme, let $\mathcal L$ be a
line bundle on $A$, trivialized along the zero section, with a cubical structure $\theta$, and let $e_N$ be its cubical
Weil pairing (Definition~\ref{cv-toroidal-abelian-dvr-ample-torsion-weil}). Let $K$ be an algebraically closed field
with a morphism $t:\operatorname{Spec}K\to S$, and let $N\ge1$ be an integer with $N\ne0$ in $K$ such that
multiplication by $N$ on $A$ is finite, flat and surjective. Let $x\in A(K)$, a point over $t$, with $Nx=0$ and with
$e_N(x,y)=1$ for every $y\in A(K)$ over $t$ with $Ny=0$. Then $t_x^*\mathcal L_K\otimes\mathcal L_K^{-1}$ is trivial
on $A_K=A\times_S\operatorname{Spec}K$.
\end{lemma}
\leanref{Development/AbelianVarietyTheory/Proofs/ModuliProofs/RatioTranslate.lean}{GlobalTheory.CubNondeg.main}
\begin{proof}
Let $m$ be multiplication by $N$ on $A_K$, let $a\in A(A_K)$ be the projection, and write $x$ also for the pullback of
$x$ to $A_K$. Let $\Lambda=\Lambda_{\mathcal L}(a,x)$ on $A_K$. As $Nx=0$, its pullback $m^*\Lambda$ is
$(Na+x)^*\mathcal L\otimes(Na)^*\mathcal L^{-1}\otimes x^*\mathcal L^{-1}\otimes0^*\mathcal L$.

First, $\sigma=c_N(0;a,x)\,c_N(x,a)^{-1}$ is a generator of $m^*\Lambda$. Indeed, $c_N(0;a,x)$ is a generator of the
tensor product of $(Na+x)^*\mathcal L\otimes(Na)^*\mathcal L^{-1}\otimes x^*\mathcal L^{-1}\otimes0^*\mathcal L$ and
$\Lambda_{\mathcal L}(a,x)^{-N}$, by the computation of Definition~\ref{cv-toroidal-abelian-dvr-ample-torsion-weil},
and $c_N(x,a)$ is a generator of $\Lambda_{\mathcal L}(x,a)^{-N}$, as $Nx=0$.

Second, $\sigma$ is invariant under the translations $\tau_y$ by the points $y\in A(K)$ with $Ny=0$, which satisfy
$m\tau_y=m$: its pullback $\tau_y^*\sigma=c_N(0;a+y,x)c_N(x,a+y)^{-1}$ is $\sigma$. Indeed, identities (i) and (ii) of
the proof of Lemma~\ref{cv-toroidal-abelian-dvr-ample-torsion-basic}, with $Nx=Ny=0$, give
$c_N(0;a+y,x)\theta(a,y,x)^N=c_N(0;a,x)c_N(y,x)$ and $c_N(x,a+y)\theta(x,a,y)^N=c_N(x,a)c_N(x,y)$. Moreover
$\theta(a,y,x)=\theta(x,a,y)$, and $c_N(x,y)=c_N(y,x)$ as $e_N(x,y)=1$.

Third, $\Lambda$ has a generator. On an open subscheme $W$ of $A_K$ with a generator $\epsilon$ of $\Lambda$, write
$\sigma=g\,m^*\epsilon$ with $g$ a unit on $m^{-1}(W)$. By the second point $\tau_y^*g=g$, so by
Lemma~\ref{cv-toroidal-abelian-dvr-ample-torsion-descent}, applied to $g$ and to $g^{-1}$, $g=m^*h$ for a unit $h$ on $W$.
The generators $h\epsilon$ agree on overlaps, by the uniqueness in that lemma, and glue to a generator of $\Lambda$.

Finally, $x^*\mathcal L\otimes0^*\mathcal L^{-1}$ is a line bundle on $\operatorname{Spec}K$, hence trivial, and the
tensor product of $\Lambda$ with its pullback to $A_K$ is $(a+x)^*\mathcal L\otimes a^*\mathcal L^{-1}=
t_x^*\mathcal L_K\otimes\mathcal L_K^{-1}$. So the latter is trivial.
\end{proof}
% Lean: CubNondeg.main: exists_sigma (sigma from exists_cubN and WeilCube.exists_gen_torsion), ratio_translate (the
% invariance, from cubN_mul_left and IsCub.N_add), descend_linComb with torsion_invariant_descends_fibre (the third
% point), DVR.exists_gen_of_local (the line bundle on Spec K) and exists_gen for the swap. The Lean also assumes the
% laws Out.Laws, which this proof does not use.

\begin{lemma}[Alp\"oge]\label{cv-toroidal-abelian-dvr-ample-torsion}
Let $S$ be a scheme, let $A$ be a commutative group scheme over $S$ that is an abelian scheme, and let $\mathcal L$
be a line bundle on $A$, trivialized along the zero section, with a cubical structure. Let $O$ be a valuation ring
whose fraction field $F$ is algebraically closed, with residue field $k$, let $q:\operatorname{Spec}O\to S$ be a
morphism, and let $n\ge1$ be an integer that is invertible on $S$. Suppose that every point $y$ of $A$ over
$\operatorname{Spec}F\to\operatorname{Spec}O\to S$ for which $t_y^*\mathcal L\otimes\mathcal L^{-1}$ is trivial is
zero. Then every point $x$ of $A$ over $\operatorname{Spec}k\to\operatorname{Spec}O\to S$ with $nx=0$ for which
$t_x^*\mathcal L\otimes\mathcal L^{-1}$ is trivial is zero.
\end{lemma}
\leanref{Development/AbelianVarietyTheory/Proofs/ModuliProofs/EqOneOfTriv.lean}{GlobalTheory.GoodPol.eq_one_of_triv}
\begin{proof}
Let $\theta$ be the cubical structure and $e_n$ its cubical Weil pairing
(Definition~\ref{cv-toroidal-abelian-dvr-ample-torsion-weil}). As $n$ is invertible on $S$, it is a unit of $O$ and
nonzero in $F$. Points of $A$ are taken over the morphisms to $S$ through $q$.

First, a point of $A$ over $\operatorname{Spec}O$ is determined by its restriction to $\operatorname{Spec}F$. Indeed,
$A$ is separated over $S$, so two such points agree on a closed subscheme of $\operatorname{Spec}O$, through which
$\operatorname{Spec}F\to\operatorname{Spec}O$ factors when their restrictions agree; as $O\to F$ is injective, that
closed subscheme is $\operatorname{Spec}O$.

Second, $x$ lifts to a point $X$ of $A$ over $\operatorname{Spec}O$ with $nX=0$. Indeed, the kernel $A[n]$ of
multiplication by $n$ is finite and \'etale over $S$, as in the proof of Lemma~\ref{cv-toroidal-moduli-torsion}, and
$x$ is a point of it. The ring $O$ is henselian with algebraically closed residue field, and points with values in $k$
of schemes smooth over $O$ lift to points with values in $O$, as in the proof of
Lemma~\ref{cv-toroidal-abelian-dvr-ample-square}; apply this to $A[n]\times_S\operatorname{Spec}O$.

Third, let $X_F$ be the restriction of $X$ to $\operatorname{Spec}F$; then $X_F=0$. By hypothesis it suffices to show
that $t_{X_F}^*\mathcal L\otimes\mathcal L^{-1}$ is trivial. Multiplication by $n$ on $A$ is finite, \'etale and
surjective, as in the proof of Lemma~\ref{cv-toroidal-moduli-torsion}, and $nX_F=0$; by
Lemma~\ref{cv-toroidal-abelian-dvr-ample-torsion-nondegenerate} it suffices to show that $e_n(X_F,y)=1$ for every point
$y$ of $A$ over $\operatorname{Spec}F$ with $ny=0$. As $A$ is proper over $S$, $y$ extends to a point $Y$ over
$\operatorname{Spec}O$ \cite[Tag 0BX5]{Sta}, and $nY=0$ by the first point, as its restriction to $\operatorname{Spec}F$
is $ny=0$. Let $\zeta=e_n(X,Y)$, a unit of $O$. By Lemma~\ref{cv-toroidal-abelian-dvr-ample-torsion-basic}(3), its image
in $k$ is $e_n(x,Y_k)$, where $Y_k$ is the restriction of $Y$ to $\operatorname{Spec}k$, and $e_n(x,Y_k)=1$ by
Lemma~\ref{cv-toroidal-abelian-dvr-ample-torsion-kernel}, as $t_x^*\mathcal L\otimes\mathcal L^{-1}$ is trivial. By
Lemma~\ref{cv-toroidal-abelian-dvr-ample-torsion-basic}(2), $\zeta^n=1$. So
$0=(\zeta-1)(1+\zeta+\cdots+\zeta^{n-1})$, and the second factor is congruent to $n$ modulo the maximal ideal, hence a
unit; as $O$ is a domain, $\zeta=1$. By Lemma~\ref{cv-toroidal-abelian-dvr-ample-torsion-basic}(3),
$e_n(X_F,y)=1$.

By the third point and the first, $X=0$, so $x$, its restriction to $\operatorname{Spec}k$, is $0$.
\end{proof}
% Lean: eq_one_of_triv (with Out.Laws): GrpTangent.finiteEtale_mulN_torsion and
% valuationRing_henselian_and_lift_of_isAlgClosed (the lift X of x, through LatticeTorsion.existsUnique_torsion);
% then CubNondeg.main (nondegeneracy) for X over F, checked with the cubical Weil pairing UU.CubWeil: cubWeil_basic
% (existence, uniqueness, powers and base change), the valuative criterion (each y over F killed by n extends to
% Y over O), KerPair.cubWeil_of_gen (the pairing of x is 1 since its translate is trivial) and unit_eq_one (a
% root of unity of order n that is 1 modulo the maximal ideal is 1); over_points_ext_of_injective for X = 0.
% Size: about 1,800 lines.

\begin{lemma}[Alp\"oge]\label{cv-toroidal-abelian-dvr-ample}
Let $V$ be a discrete valuation ring with fraction field $K$, let $A$ be a commutative group scheme over $V$
that is an abelian scheme, and let $\mathcal L$ be a line bundle on $A$, trivialized along the zero section,
with a cubical structure. Let $A'$ be a group scheme over $K$ with a line bundle $\mathcal L'$, and let
$\varphi:A'\to A$ be a morphism over $\operatorname{Spec}K\to\operatorname{Spec}V$ that identifies $A'$ with
$A\times_V\operatorname{Spec}K$ as group schemes, with $\mathcal L'\cong\varphi^*\mathcal L$. If
$\mathcal L'$ is relatively ample and defines a principal polarization, then $\mathcal L$ is relatively
ample.
\end{lemma}
\leanref{Development/AbelianVarietyTheory/Proofs/ModuliProofs/ExistsPrincipallyPolarizedCMExtensionDvr.lean}{GlobalTheory.GoodPol.relAmple_of_generic}
\begin{proof}
Let $F$ be an algebraic closure of $K$, let $O$ be a valuation ring of $F$ that dominates $V$
\cite[Tag 00IA]{Sta}, and let $k$ be its residue field, which is algebraically closed (as in the proof of
Lemma~\ref{cv-toroidal-abelian-dvr-ample-square}). Let $\beta:V\to O\to k$ be the composite homomorphism, whose
kernel is the maximal ideal of $V$.

Step 1. Through $\varphi$, the line bundle $\mathcal L'$ on $A'$ is the pullback of $\mathcal L$ to
$A\times_V\operatorname{Spec}K$, which is therefore relatively ample.

Step 2. By Step 1 and Lemma~\ref{cv-toroidal-abelian-dvr-ample-generating}, there are an integer $d\ge1$ and
sections $u_0,\dots,u_m$ of $\mathcal L^d$ whose nonvanishing loci cover $A$.

Step 3. Let $B=A\times_V\operatorname{Spec}k$, through $\beta$, with the pullbacks $\mathcal L_B$ of $\mathcal L$, of
its trivialization along the zero section and of its cubical structure. Then $\mathcal L_B$ is ample. Indeed,
some $u_i$ does not vanish at the image of the origin of $B$, so its pullback is a nonzero section of
$\mathcal L_B^d$. Let $n$ be a prime larger than $2d$ and than the characteristic of $k$; then $n\ne0$ in $k$ and
$n$ does not divide $2d$. By Lemma~\ref{cv-toroidal-abelian-dvr-ample-effective}, with $\Omega=k$, it suffices to
show that every $x\in B(k)$ with $nx=0$ for which $t_x^*\mathcal L_B\otimes\mathcal L_B^{-1}$ is trivial is zero.
As $\beta(n)\ne0$, $n$ is a unit of $V$. By Lemma~\ref{cv-toroidal-abelian-dvr-ample-torsion}, applied to $A$ over
$\operatorname{Spec}V$, $O$ and $\operatorname{Spec}O\to\operatorname{Spec}V$, it suffices to show that every point
$y$ of $A$ over $\operatorname{Spec}F\to\operatorname{Spec}V$ for which $t_y^*\mathcal L\otimes\mathcal L^{-1}$ is
trivial is zero. Such a point lies over $\operatorname{Spec}F\to\operatorname{Spec}K$, so it is $\varphi\circ y'$ for
a point $y'$ of $A'$, and through $\mathcal L'\cong\varphi^*\mathcal L$ the line bundle
$t_{y'}^*\mathcal L'\otimes\mathcal L'^{-1}$ is trivial. So $y'$ is a point of $K(\mathcal L')$
(Lemma~\ref{cv-toroidal-abelian-dvr-closed}), and $y'=0$, as $\mathcal L'$ defines a principal polarization. Hence
$y=0$.

Step 4. The line bundle $\mathcal L$ is ample on the fibre of $A$ over the generic point of
$\operatorname{Spec}V$, by Step 1, and on the fibre over the closed point, by Step 3 and descent along the
extension $k$ of its residue field \cite[Tag 0D2P]{Sta}. As $A$ is proper over $V$, $\mathcal L$ is relatively
ample over an open neighbourhood of each point of $\operatorname{Spec}V$ \cite[Tag 0D2S]{Sta}, hence relatively
ample \cite[Tag 01VJ]{Sta}.
\end{proof}
% Lean: GoodPol.relAmple_of_generic (with Out.Laws): relAmple_pull (Step 1), IsLocalRing.exists_factor_valuationRing
% and valuationRing_henselian_and_lift_of_isAlgClosed (O and k), exists_generating (Step 2), exists_ne_zero,
% AmpleCrit.relAmple_of_effective with abMulN_holds, eq_one_of_triv, eq_one_of_inK and inK_of_triv (Step 3). For
% Step 4 the Lean uses the sections u_i: T.isAffineOpen_nonzeroLocus_of_fibrewise_ample (their nonvanishing loci
% are affine, from ampleness on the fibres over both points; about 600 lines) and
% AbLaw.veryAmpleBy_pw_of_affine_loci (a power of L^d is then very ample by sections; about 700), then
% Unif.relAmple_iff_veryAmpleBy. The text cites Stacks for this step, which the Lean proves itself.
% Lemmas cv-toroidal-abelian-dvr-ample-effective and -ample-torsion are written at the floor.

\begin{lemma}[Alp\"oge]\label{cv-toroidal-abelian-dvr-structure}
In Situation~\ref{cv-function-situation}, let $V$ be a discrete valuation ring with fraction field $K$, let
$A$ be a group scheme over $V$ that is an abelian scheme, let $G_K$ be an $\mathcal O_E$-semi-abelian scheme
over $K$ (Definition~\ref{cv-model-toroidal-semiabelian}(2)) with line bundle $\mathcal L_K$, and let
$\varphi:G_K\to A$ be a morphism over $\operatorname{Spec}K\to\operatorname{Spec}V$ that identifies $G_K$
with $A\times_V\operatorname{Spec}K$ as group schemes. Suppose that $\mathcal L_K$ is relatively ample and
defines a principal polarization. Then there are an $\mathcal O_E$-semi-abelian scheme $G'$ over $V$ with line
bundle $\mathcal L'$, a morphism $\nu':G_K\to G'$ of $\mathcal O_E$-semi-abelian schemes over
$\operatorname{Spec}K\to\operatorname{Spec}V$ and an isomorphism of group schemes $e:G'\to A$ over $V$ with
$e\circ\nu'=\varphi$, such that $\mathcal L'$ is relatively ample and defines a principal polarization.
\end{lemma}
\leanref{Development/AbelianVarietyTheory/Proofs/ModuliProofs/ExistsPrincipallyPolarizedCMExtensionDvrStatement.lean}{GlobalTheory.T.exists_principallyPolarizedCM_extension_dvr}
\begin{proof}
The group scheme $A$ is commutative, smooth and proper over $V$, with geometrically connected fibres.

The action. For $b\in\mathcal O_E$, the endomorphism of $A\times_V\operatorname{Spec}K$ that corresponds
through $\varphi$ to the action of $b$ on $G_K$ extends uniquely to an endomorphism $[b]$ of $A$, by
Theorem~\ref{cv-model-toroidal-extension} applied to the abelian schemes $A$ and $A$. By the uniqueness,
$[1]=1$, $[bb']=[b][b']$ and $[b+b']=[b]+[b']$, since both sides of each identity extend the same
endomorphism of the generic fibre. So $b\mapsto[b]$ is an action of $\mathcal O_E$ on $A$, and $\varphi$
commutes with the actions.

The line bundle. The scheme $A$ is smooth over the regular ring $V$, so its local rings are regular
\cite[Tag 07NF]{Sta}, hence factorial \cite[Tag 0AG0]{Sta}. The generic fibre $A\times_V\operatorname{Spec}K$
is open in $A$, so the line bundle on it that corresponds to $\mathcal L_K$ through $\varphi$ extends to a line
bundle $\mathcal M$ on $A$ \cite[Tag 0BD9]{Sta}. Then $\varphi^*\mathcal M\cong\mathcal L_K$.

The structure. The morphism $\operatorname{Spec}K\to\operatorname{Spec}V$ is an open immersion, and
$\operatorname{Spec}V$ is integral, noetherian and normal. Lemma~\ref{cv-toroidal-abelian-dvr-cubical},
applied to $U=\operatorname{Spec}K$, $G=A$ with the action above, $G_U=G_K$, $a=\varphi$ and $\mathcal M$,
gives a line bundle $\mathcal L'$ on $A$, trivialized along the zero section, with a cubical structure, that
makes $A$ an $\mathcal O_E$-semi-abelian scheme $G'$, and for which $\varphi$ is a morphism $\nu':G_K\to G'$
of $\mathcal O_E$-semi-abelian schemes. In particular $\mathcal L_K\cong\nu'^*\mathcal L'$. Let $e$ be the
identity of $A$.

Ampleness and principality. Lemma~\ref{cv-toroidal-abelian-dvr-ample}, applied to $A$ with the line bundle
$\mathcal L'$, to $A'=G_K$ with $\mathcal L_K$, and to $\varphi$, shows that $\mathcal L'$ is relatively
ample. Then Lemma~\ref{cv-toroidal-abelian-dvr-kernel}, applied to $B=A$ with $\mathcal L'$, to $B'=G_K$ with
$\mathcal L_K$, and to $\varphi$, shows that $\mathcal L'$ defines a principal polarization.
\end{proof}
% Lean: T.exists_principallyPolarizedCM_extension_dvr = GoodPol.exists_sab_of_generic:
% CM.exists_oAct_extend_dedekind with HomExtSt (Theorem cv-model-toroidal-extension) for the action;
% gen_chartPull and the law Out.Laws.extend (Semi.extend_existsSt: a rigidified line bundle on the generic fibre
% of a smooth group over a discrete valuation ring extends), which the text derives from Stacks; CechPic for
% LBIso (M.pull phi) tK.PP.L; ExtendRig.isOpenImmersion_etaK, Published.SemiAbelianHomExtension.isNormalScheme_dvr
% and TripleExt.exists_sab_extension; then relAmple_of_generic and kTriv_of_generic. The isomorphism e is the
% identity of A.

\begin{lemma}[Alp\"oge]\label{cv-toroidal-abelian-dvr-determinant}
In Situation~\ref{cv-function-situation}, let $S$ be a scheme, let $G$ be a group scheme over $S$ with an
action of $\mathcal O_E$, let $o$ be a ring with a morphism $s:S\to\operatorname{Spec}o$, and let $\rho\in o$
with $\rho^2=-d$. For an algebraically closed field $k$ of characteristic zero and a point
$x:\operatorname{Spec}k\to S$, with $\beta:o\to k$ the homomorphism of $s\circ x$, say that $G$ has
\emph{signature} $(p,q)$ at $x$ if for every ring homomorphism $\tau:\mathcal O_E\to k$ the subspace of the Lie
algebra of $G_x$ on which each $b\in\mathcal O_E$ acts by $\tau(b)$ has dimension $p$ if
$\tau(\sqrt{-d})=\beta(\rho)$ and $q$ otherwise.
\begin{enumerate}
\item Suppose that $S$ is reduced and that $G$ is commutative and smooth of relative dimension $g$ over $S$. If
  the points of $S$ that are images of points $x$ at which $G$ has signature $(p,q)$ are dense in $S$, then $G$
  has signature $(p,q)$ at every point $x:\operatorname{Spec}k\to S$ with $k$ algebraically closed of
  characteristic zero.
\item If $G$ has signature $(p,q)$ at $x$ and $G_x$ is smooth of dimension $n$ over $k$, then
  $n=\frac12[E:\mathbf Q](p+q)$.
\end{enumerate}
\end{lemma}
\leanref{Development/AbelianVarietyTheory/Proofs/ModuliProofs/DeterminantConditionSpreadsOfDensePointsStatement.lean}{GlobalTheory.T.determinantCondition_spreads_of_dense_points}
\leanref{Development/AbelianVarietyTheory/Proofs/ModuliProofs/ObjIsAbelianSchemeDimRigidifiedSym.lean}{GlobalTheory.ObjDim.dim_eq_of_detCond}
\begin{proof}
(1) Let $x:\operatorname{Spec}k\to S$ be such a point, with image $y$, and let $\operatorname{Spec}A$ be an
affine open subscheme of $S$ containing $y$. As $G$ is smooth of relative dimension $g$, there is $a\in A$
with $y\in D(a)$ over which the Lie algebra of $G$ is free: there are points $v_1,\dots,v_g$ of $G$ with values
in $A_a[\epsilon]$, $\epsilon^2=0$, that reduce to $0$ modulo $\epsilon$, whose images form a basis of the Lie
algebra of the fibre of $G$ at every point $\operatorname{Spec}k'\to\operatorname{Spec}A_a$ with $k'$ a field.
In this basis, the action of $b\in\mathcal O_E$ on the Lie algebra of $G$ over $A_a$ is a matrix
$M(b)\in M_g(A_a)$, the map $b\mapsto M(b)$ is a ring homomorphism, and the action of $b$ on the Lie algebra
of the fibre at a point given by $f:A_a\to k'$ is $f(M(b))$. So $G$ has signature $(p,q)$ at such a point,
with $k'$ algebraically closed of characteristic zero, if and only if $M$ has signature $(p,q)$ at $f$ in the
sense of Lemma~\ref{cv-toroidal-abelian-dvr-eigen}, for the image $r$ of $\rho$ in $A_a$. The ring $A_a$ is
reduced, as $S$ is. The open set $D(a)$ meets the dense set of the hypothesis in a dense subset of
$\operatorname{Spec}A_a$, and each of its points is the kernel of a homomorphism $f$ at which $M$ has
signature $(p,q)$, since a point $\operatorname{Spec}k'\to S$ with image in $D(a)$ factors through
$\operatorname{Spec}A_a$. The point $x$ factors through $\operatorname{Spec}A_a$ as well, and
Lemma~\ref{cv-toroidal-abelian-dvr-eigen} shows that $G$ has signature $(p,q)$ at $x$.

(2) The Lie algebra $\mathfrak g$ of $G_x$ has dimension $n$. As in the proof of
Lemma~\ref{cv-toroidal-abelian-dvr-eigen}, $\mathfrak g$ is the direct sum of its subspaces
$\mathfrak g_\tau$ on which $\mathcal O_E$ acts through $\tau$, over the $[E:\mathbf Q]$ homomorphisms
$\tau:\mathcal O_E\to k$. Complex conjugation $c$ of the CM field $E$ maps $\sqrt{-d}$ to $-\sqrt{-d}$, and
$\beta(\rho)\ne-\beta(\rho)$; so $\tau\mapsto\tau\circ c$ exchanges the $\tau$ with
$\tau(\sqrt{-d})=\beta(\rho)$ and the $\tau$ with $\tau(\sqrt{-d})=-\beta(\rho)$, and there are
$\frac12[E:\mathbf Q]$ of each. Hence $n=\frac12[E:\mathbf Q]p+\frac12[E:\mathbf Q]q$.
\end{proof}
% Lean: (1) T.determinantCondition_spreads_of_dense_points = DetSpread.detCond_of_dense: LieLoc.lie_locally_free
% (a universal chart of the smooth G), exists_factor, DetSpread.detCond_of_basis with
% exists_lie_coordinates_lieMult_iff (the matrices M and their base change; about 1,300 lines),
% T.cmType_eigenMultiplicity_of_dense_primes and LiePull.lieMult_pull. (2) dim_eq_of_detCond:
% univChart_of_relativeDimension, tangentAction, Chart.finrank_of_lieMult and finrank_eq_of_eigenSub (with
% EigenSplit.eigenEquiv and sum_embeddings).

\begin{lemma}[Alp\"oge]\label{cv-toroidal-abelian-dvr-pairing}
Let $S$ be a scheme and $P$ an abelian scheme over $S$ with a line bundle $\mathcal L$ that is trivialized
along the zero section, isomorphic to its pullback by $-1$ locally on $S$, relatively ample and defines a
principal polarization. Let $n\ge1$, let $T$ be a scheme over $S$ on which $n$ is invertible, let
$x,y\in P(T)$ be killed by $n$, and let $c$ be a unit of $\Gamma(T,\mathcal O_T)$ with $c^n=1$. Then there is
an open and closed subset $U$ of $T$ such that a morphism $h:T'\to T$ over $S$ satisfies
$e_n(x\circ h,y\circ h)=h^*c$, for the commutator pairing $e_n$ of the theta group of $\mathcal L^n$ over
$T'$, if and only if $h(T')\subset U$.
\end{lemma}
\leanref{Development/AbelianVarietyTheory/Proofs/ModuliProofs/ExistsClopenWeilIsIff.lean}{GlobalTheory.exists_clopen_weilIs_iff}
\begin{proof}
For a point $x$ killed by $n$, $t_x^*\mathcal L^n\cong\mathcal L^n$ locally on $T$, by the theorem of the
square over $T$, whose form over an algebraically closed field is \cite[\S6, Corollary~4]{Mum70}; so locally on
$T$, $x$ and $y$ lift to the theta group of
$\mathcal L^n$. The commutator of two lifts is an automorphism of the pullback of $\mathcal L^n$, that is, a
unit of the global functions of $P_T$ over the open set, which are those of $T$ \cite[Tag 0E0L]{Sta}; it
does not depend on the lifts. These units glue to a unit $c_0$ of $\Gamma(T,\mathcal O_T)$ with
$e_n(x,y)=c_0$, and $e_n(x\circ h,y\circ h)=h^*c_0$ for every $h$, as lifts pull back to lifts
\cite[\S23]{Mum70}. By bilinearity, $c_0^n=e_n(nx,y)=e_n(0,y)=1$. Let $u=c_0c^{-1}$, so that $u^n=1$, and let
$Q(u)=1+u+\dots+u^{n-1}$. Then $(u-1)Q(u)=0$, and $Q(u)\equiv n$ modulo $u-1$, a unit, so the closed subsets
$V(u-1)$ and $V(Q(u))$ of $T$ are disjoint and cover $T$; both are open. Let $U=V(u-1)$. On the open
subscheme $U$, $Q(u)$ is a unit, so $u=1$ there. Hence if $h(T')\subset U$, then $h^*u=1$, that is,
$e_n(x\circ h,y\circ h)=h^*c$; conversely, if $h^*u=1$, then $h(T')\subset V(u-1)=U$.
\end{proof}
% Lean: exists_clopen_weilIs_iff = PairingLocus.exists_clopen_of_lift with GlobalFnsSt and kerPow_of_square
% (AbSquare.square_of_cube, from the theorem of the cube) for the lifts: CommutatorPairing.existsUnique_weilIs
% (exists_weilIs, weilIs_unique), PairingLocus.exists_clopen_of_weilIs (pow_eq_one, CommPairing.weilIs_pull,
% exists_clopen_eq). The Lean proves these properties of the commutator pairing itself (about 1,200 lines).

\begin{lemma}[Alp\"oge]\label{cv-toroidal-abelian-dvr-isometry}
In Situation~\ref{cv-function-situation}, let $\Lambda$ be the lattice of the datum with its alternating form
$\langle\,,\rangle$. Let $S$ be a scheme, let $P$ be an abelian scheme over $S$ with a line bundle
$\mathcal L$ as in Lemma~\ref{cv-toroidal-abelian-dvr-pairing}, and let $\alpha_b$, for $b\in\mathcal O_E$, be
endomorphisms of $P$. Let $o$ be a ring with a morphism $s:S\to\operatorname{Spec}o$, let $z\in o$, let
$n\ge1$ and $m\ge0$ be integers, and let $(a_i)_{i\in I}$ be a family in $\Lambda$ and $(p_i)_{i\in I}$ a
family in $P(S)$. For a field $K$ in which $n\ne0$ and a point $x:\operatorname{Spec}K\to S$, with
$\beta:o\to K$ the homomorphism of $s\circ x$, an \emph{$n$-isometry} at $x$ is a pair $(\zeta,\iota)$ of a
primitive $n$-th root of unity $\zeta\in K$ and an additive map $\iota:\Lambda\to P_x(K)$ such that
$n\iota=0$, $\iota(ba)=\alpha_b(\iota(a))$, $e_n(\iota(a),\iota(a'))=\zeta^{\langle a,a'\rangle}$ for the
commutator pairing of the theta group of $\mathcal L^n$, $\zeta^m=\beta(z)$ and $m\iota(a_i)=p_i(x)$, for
all $a,a'\in\Lambda$, $b\in\mathcal O_E$ and $i\in I$. If the points of $S$ that are images of points at which
there is an $n$-isometry are dense in $S$, then there is an $n$-isometry at every point
$x:\operatorname{Spec}K\to S$ with $K$ algebraically closed and $n\ne0$ in $K$.
\end{lemma}
\leanref{Development/AbelianVarietyTheory/Proofs/ModuliProofs/ExistsLevelIsometryOfDenseGeometricPointsStatement.lean}{GlobalTheory.T.exists_level_isometry_of_dense_geometric_points}
\begin{proof}
Let $S_n$ be the open subscheme of $S$ on which $n$ is invertible; every point $x$ as in the statement factors
through it. Lemma~\ref{cv-toroidal-moduli-torsion}(1), for $P\times_SS_n$, $n$ and $\Lambda$, gives a scheme
$Z_0$ finite over $S_n$ and an additive map $\epsilon_0:\Lambda\to P(Z_0)$ with $n\epsilon_0=0$ that is
universal among such maps. Let $Z=Z_0\times_{\mathbf Z}\operatorname{Spec}\mathbf Z[t]/(\Phi_n(t))$, where
$\Phi_n$ is the $n$-th cyclotomic polynomial, let $\epsilon:\Lambda\to P(Z)$ be the pullback of $\epsilon_0$,
and let $c\in\Gamma(Z,\mathcal O_Z)$ be the image of $t$, a unit with $c^n=1$. The morphism $\pi:Z\to S$ is
the composite of the finite morphisms $Z\to Z_0\to S_n$ with the open immersion $S_n\to S$. At every point
$h:\operatorname{Spec}K\to Z$, $h^*c$ is a primitive $n$-th root of unity, as $\Phi_n(h^*c)=0$ and $n\ne0$ in
$K$. Conversely, for a field $K$ with $n\ne0$ in $K$, a point $x:\operatorname{Spec}K\to S$, an additive map
$\iota:\Lambda\to P_x(K)$ with $n\iota=0$ and a primitive $n$-th root of unity $\zeta\in K$, there is a point
$h:\operatorname{Spec}K\to Z$ over $x$ with $\epsilon(a)\circ h=\iota(a)$ for all $a$ and $h^*c=\zeta$, by the
universal property of $Z_0$ and since $\Phi_n(\zeta)=0$.

Let $C\subset Z$ be the intersection of the following subsets: for $b\in\mathcal O_E$ and $a\in\Lambda$, the
closed subset on which the morphisms $\epsilon(ba)$ and $\alpha_b\circ\epsilon(a)$ from $Z$ to $P$ agree,
which is closed as $P$ is separated over $S$; for $a,a'\in\Lambda$, the open and closed subset of
Lemma~\ref{cv-toroidal-abelian-dvr-pairing} for $T=Z$, $x=\epsilon(a)$, $y=\epsilon(a')$ and
$c^{\langle a,a'\rangle}$; the closed subset $V(c^m-z)$; and for $i\in I$, the closed subset on which
$m\epsilon(a_i)$ and $p_i\circ\pi$ agree. So $C$ is closed. For a field $K$ and a point
$h:\operatorname{Spec}K\to Z$ over $x$, the pair $(h^*c,\epsilon\circ h)$ is an $n$-isometry at $x$ if and only
if the image of $h$ lies in $C$, as $\operatorname{Spec}K$ is reduced, so that $h$ factors through a closed
subscheme if and only if its image lies in it.

By the previous paragraphs, every point of $S$ at which there is an $n$-isometry lies in $\pi(C)$, so
$\pi(C)$ is dense in $S$. As $Z$ is finite over $S_n$, $\pi(C)$ is closed in $S_n$, so it contains $S_n$. Let
$K$ be algebraically closed with $n\ne0$ in $K$, and $x:\operatorname{Spec}K\to S$ with image $y\in S_n$.
Choose $w\in C$ with $\pi(w)=y$. The scheme $\operatorname{Spec}K\times_SZ$ is finite over $K$ and has a point
over $w$; its residue field is $K$, as $K$ is algebraically closed. This gives a point
$h:\operatorname{Spec}K\to Z$ over $x$ with image $w$, and $(h^*c,\epsilon\circ h)$ is an $n$-isometry at
$x$.
\end{proof}
% Lean: exists_levelIsometry_of_dense: LevCover.datum_free, T.exists_finite_scheme_latticeTorsionHom_primitiveRoot
% (= exists_latticeTorsionHom_primitiveRoot_param = TorsRootParam.main_lemma_3: exists_finiteEtale_latticeTorsionHom
% over the basic open of n, the ring CycRing n = Z[t]/(Phi_n), isPrimitiveRoot_map, exists_map_eq and
% exists_lift_of_range_subset_closure_image_finite), exists_closed_eq (C1, C4), exists_clopen_weilIs_iff (C2),
% the basic open of c^m - z (C3), and pow_eq_iff.

\begin{lemma}[Alp\"oge]\label{cv-toroidal-abelian-dvr-tate}
In Situation~\ref{cv-function-situation}, let $\Lambda$ be the lattice of the datum with its alternating form
$\langle\,,\rangle$, of rank $r$ over $\mathbf Z$. Let $K$ be a field, let $Q$ be a group, written additively,
with maps $\alpha_b:Q\to Q$ for $b\in\mathcal O_E$, and for each $m\ge1$ let $e_m$ be a relation between pairs
$y,y'\in Q$ and units $c\in K^\times$, written $e_m(y,y')=c$. Write $Q[m]$ for the elements killed by $m$. For
$m\ge1$, a \emph{level-$m$ isometry} is a pair $(\zeta,\iota)$ of a primitive $m$-th root of unity
$\zeta\in K$ and an additive map $\iota:\Lambda\to Q$ with kernel $m\Lambda$ and image $Q[m]$, such that
$\iota(ba)=\alpha_b(\iota(a))$ and $e_m(\iota(a),\iota(a'))=\zeta^{\langle a,a'\rangle}$ for all
$a,a'\in\Lambda$ and $b\in\mathcal O_E$.
\begin{enumerate}
\item Let $m\ge1$ be such that $Q[m]$ has $m^r$ elements and $e_m(0,y)=c$ implies $c=1$ for $y\in Q[m]$.
  Let $\zeta\in K$ be a primitive $m$-th root of unity and $\iota:\Lambda\to Q$ an additive map with
  $m\iota=0$ and $e_m(\iota(a),\iota(a'))=\zeta^{\langle a,a'\rangle}$ for all $a,a'$. Then $\iota$ has kernel
  $m\Lambda$ and image $Q[m]$.
\item Suppose that $Q[m]$ is finite for every $m\ge1$, and that $e_{mm'}(y,y')=c$ implies
  $e_m(m'y,m'y')=c^{m'}$ for $y,y'\in Q[mm']$. Let $l\ge1$, $z\in K$, and let $(a_i)_{i\in I}$ in $\Lambda$ and
  $(y_i)_{i\in I}$ in $Q$ be families. If for every $m'\ge1$ there is a level-$lm'$ isometry $(\zeta,\iota)$
  with $\zeta^{m'}=z$ and $m'\iota(a_i)=y_i$ for all $i$, then there are level-$m$ isometries
  $(\zeta_m,\iota_m)$ for all $m\ge1$ such that $\zeta_{mm'}^{m'}=\zeta_m$ and $m'\iota_{mm'}=\iota_m$ for all
  $m,m'\ge1$, $\zeta_l=z$, and $\iota_l(a_i)=y_i$ for all $i$.
\end{enumerate}
\end{lemma}
\leanref{Development/AbelianVarietyTheory/Proofs/ModuliProofs/TorsionIsometryPrimitiveRootBijectiveModNStatement.lean}{GlobalTheory.T.torsionIsometry_primitiveRoot_bijective_mod_n}
\leanref{Development/AbelianVarietyTheory/Proofs/ModuliProofs/TateIsometryOfFiniteLevelIsometries.lean}{GlobalTheory.TateTower.tateIso_of_levels}
\begin{proof}
(1) If $\iota(a)=0$, then $\zeta^{\langle a,a'\rangle}=1$ for every $a'$, by the hypothesis on $e_m(0,\cdot)$;
so $m$ divides $\langle a,a'\rangle$ for every $a'$, and $a\in m\Lambda$ as the form is perfect. Conversely
$\iota(m\Lambda)=0$. So $\iota$ induces an injective map $\Lambda/m\Lambda\to Q[m]$, which is bijective as
both sets have $m^r$ elements.

(2) For a pair $(\zeta,\iota)$ and $t\ge1$ write $(\zeta,\iota)^t=(\zeta^t,t\iota)$. If $(\zeta,\iota)$ is a
level-$mt$ isometry, then $(\zeta,\iota)^t$ is a level-$m$ isometry. Indeed, $\zeta^t$ is a primitive $m$-th
root of unity, and $t\iota$ is additive and commutes with the $\alpha_b$. As $\Lambda$ is torsion free,
$t\iota(a)=\iota(ta)=0$ if and only if $ta\in mt\Lambda$, that is, $a\in m\Lambda$. If $my=0$, then
$y\in Q[mt]$, so $y=\iota(a)$, and $\iota(ma)=0$ gives $ma\in mt\Lambda$, so $a=ta'$ and $y=t\iota(a')$;
conversely $mt\iota(a)=0$. Finally $e_m(t\iota(a),t\iota(a'))=\zeta^{t\langle a,a'\rangle}$ by the hypothesis
on the relations. For each $m$ there are finitely many level-$m$ isometries, since $\zeta$ is an $m$-th root
of unity and $\iota$ is determined by its values on a finite set of generators of $\Lambda$, which lie in the
finite set $Q[m]$.

For $j\ge1$ let $(\xi_j,\kappa_j)$ be a level-$l\,j!$ isometry with $\xi_j^{j!}=z$ and $j!\kappa_j(a_i)=y_i$
for all $i$. Let $\mathcal U$ be an ultrafilter on the positive integers that contains every set
$\{j:j\ge j_0\}$. For $m\ge1$, $m$ divides $l\,j!$ for all $j\ge m$, and then $(\xi_j,\kappa_j)^{l\,j!/m}$ is a
level-$m$ isometry. As there are finitely many of these, one of them, $(\zeta_m,\iota_m)$, is equal to
$(\xi_j,\kappa_j)^{l\,j!/m}$ for all $j$ in a set that belongs to $\mathcal U$. For $m,m'\ge1$, choose
$j\ge mm'$ in the sets of $\mathcal U$ for $mm'$ and for $m$; then
$(\zeta_{mm'},\iota_{mm'})^{m'}=(\xi_j,\kappa_j)^{l\,j!/m}=(\zeta_m,\iota_m)$. For $m=l$, there is $j$ with
$(\zeta_l,\iota_l)=(\xi_j,\kappa_j)^{j!}$, so $\zeta_l=z$ and $\iota_l(a_i)=y_i$.
\end{proof}
% Lean: (1) TorsIsometry.ker_and_onto (exists_eq_nsmul_of_dvd with D.unimod; the count through
% Module.finBasis). (2) tateIso_of_levels: LevelIso, LevelIso.pow, finite_levelIso, Ultrafilter.of atTop and
% Ultrafilter.eq_pure_of_finite_mem.

\begin{lemma}[Alp\"oge]\label{cv-toroidal-abelian-dvr-spread}
In Situation~\ref{cv-function-situation}, let $n$ be a prime and $N\ge1$. For a scheme $S$ over $R_{n,N}$
(Definition~\ref{cv-toroidal-family}), an $\mathcal O_E$-semi-abelian scheme $G$ over $S$
(Definition~\ref{cv-model-toroidal-semiabelian}(2)), a map $\lambda:\Lambda\to G(S)$ and a point
$x:\operatorname{Spec}k\to S$ with $k$ an algebraically closed field of characteristic zero, say that
$(G,\lambda)$ satisfies (a) and (b) at $x$ if conditions (a) and (b) of
Definition~\ref{cv-toroidal-family}(1) hold at $x$.
\begin{enumerate}
\item Let $f:S'\to S$ be a morphism of schemes over $R_{n,N}$, let $\nu:G'\to G$ be a morphism of
  $\mathcal O_E$-semi-abelian schemes over $f$, and let $\lambda:\Lambda\to G(S)$ and
  $\lambda':\Lambda\to G'(S')$ be maps with $\nu\circ\lambda'(a)=\lambda(a)\circ f$ for all $a\in\Lambda$. If
  $(G',\lambda')$ satisfies (a) and (b) at a point $y:\operatorname{Spec}k\to S'$, then $(G,\lambda)$
  satisfies them at $f\circ y$.
\item Let $S$ be a locally noetherian, reduced and connected scheme over $R_{n,N}$, let $G$ be an
  $\mathcal O_E$-semi-abelian scheme over $S$, proper over $S$, whose line bundle is relatively ample and
  defines a principal polarization, and let $\lambda:\Lambda\to G(S)$ be a map. If the points of $S$ that are
  images of points at which $(G,\lambda)$ satisfies (a) and (b) are dense in $S$, then $(G,\lambda)$ satisfies
  (a) and (b) at every point $x:\operatorname{Spec}k\to S$ with $k$ an algebraically closed field of
  characteristic zero.
\end{enumerate}
\end{lemma}
\leanref{Development/AbelianVarietyTheory/Proofs/ModuliProofs/FullObjectOfGenericFibreDvr.lean}{GlobalTheory.FullSpread.GeomAt.push}
\leanref{Development/AbelianVarietyTheory/Proofs/ModuliProofs/FullObjectOfGenericFibreDvr.lean}{GlobalTheory.FullSpread.geomAt_of_dense}
\begin{proof}
(1) The morphism $\nu$ identifies the fibre of $G'$ at $y$ with the fibre of $G$ at $f\circ y$, compatibly
with the actions of $\mathcal O_E$ and the line bundles, and the homomorphisms $R_{n,N}\to k$ of $y$ and of
$f\circ y$ agree. So the eigenspaces of the Lie algebras correspond, which gives (a). A family
$(\zeta'_m,\iota_m)$ as in (b) for $(G',\lambda')$ at $y$, composed with this identification, is one for
$(G,\lambda)$ at $f\circ y$, since the commutator pairings of the theta groups correspond and $\nu$ carries
$\lambda'(a)(y)$ to $\lambda(a)(f\circ y)$.

(2) Let $\mathcal L$ be the line bundle of $G$, which is trivialized along the zero section and symmetric. As
$G$ is proper and smooth over $S$ with connected fibres, it is an abelian scheme; as $G\to S$ is open with
connected fibres and $S$ is connected, $G$ is connected, so its relative dimension is a constant $g$. Let
$x:\operatorname{Spec}k\to S$ be as in (2), with $\beta:R_{n,N}\to k$ its homomorphism.

Condition (a) at a point is the signature $(p,q)$ of Lemma~\ref{cv-toroidal-abelian-dvr-determinant} for
$o=R_{n,N}$ and $\rho=\theta$. By part (1) of that lemma it holds at $x$, and by part (2),
$g=\frac12[E:\mathbf Q](p+q)$. As $\Lambda$ is a projective $\mathcal O_E$-module of rank $p+q$, its rank
over $\mathbf Z$ is $[E:\mathbf Q](p+q)=2g$. For $m\ge1$, the points of $G_x(k)$ killed by $m$ are $m^{2g}$ in
number, as multiplication by $m$ on the abelian variety $G_x$ is \'etale of degree $m^{2g}$ in characteristic
zero and $k$ is algebraically closed \cite[\S6, Application~3]{Mum70}.

For condition (b), let $Q=G_x(k)$ with the action $\alpha_b$ of $\mathcal O_E$, and let $e_m$ be the commutator
pairing of the theta group of $\mathcal L_x^m$. If $e_{mm'}(y,y')=c$ for $y,y'\in Q[mm']$, then
$e_{mm'}(y,m'y')=c^{m'}$, by bilinearity and skew-symmetry, and $e_{mm'}(y,m'y')=e_m(m'y,m'y')$ as
$m'y'\in Q[m]$, by property (4) of the pairings $e^{\mathcal L}$ in \cite[\S23, p.~228]{Mum70} and
skew-symmetry; so $e_m(m'y,m'y')=c^{m'}$. By bilinearity, $e_m(0,y)=c$ implies $c=1$. Let
$m'\ge1$. At each point $x'$ of the dense set of the hypothesis, with homomorphism $\beta'$, the family
$(\zeta'_j,\iota_j)$ of condition (b) gives the $Nm'$-isometry $(\zeta'_{Nm'},\iota_{Nm'})$ of
Lemma~\ref{cv-toroidal-abelian-dvr-isometry}, for $P=G$, the action of $\mathcal O_E$, $o=R_{n,N}$,
$z=\zeta_N$, the integers $Nm'$ and $m'$, and the families $(a)_{a\in\Lambda}$ and
$(\lambda(a))_{a\in\Lambda}$: indeed $(\zeta'_{Nm'})^{m'}=\zeta'_N=\beta'(\zeta_N)$,
$m'\iota_{Nm'}(a)=\iota_N(a)=\lambda(a)(x')$, and $Nm'\ne0$ in characteristic zero. So that lemma gives an
$Nm'$-isometry $(\zeta,\iota)$ at $x$. By Lemma~\ref{cv-toroidal-abelian-dvr-tate}(1), applied to $Q$, whose
elements killed by $Nm'$ are $(Nm')^{2g}$ in number, $(\zeta,\iota)$ is a level-$Nm'$ isometry; and
$\zeta^{m'}=\beta(\zeta_N)$ and $m'\iota(a)=\lambda(a)(x)$. Lemma~\ref{cv-toroidal-abelian-dvr-tate}(2), with
$l=N$, $z=\beta(\zeta_N)$ and the families $(a)_{a\in\Lambda}$ and $(\lambda(a)(x))_{a\in\Lambda}$, gives the
roots of unity $\zeta'_m$ and the maps $\iota_m$ of condition (b) at $x$, with $\zeta'_N=\beta(\zeta_N)$ and
$\iota_N(a)=\lambda(a)(x)$.
\end{proof}
% Lean: (1) GeomAt.push: detCond_iff_of_sabMap, ObjPull.tateMove with ObjPull.pushEquiv,
% weilIs_iff_of_isGrpPullback_comm (with the LBIso of nu.lin) and toUnit_left'. (2) geomAt_of_dense:
% exists_dim_of_preconnected, the PolAV built from t (ObjDim.rigidified, ObjDim.sym),
% T.determinantCondition_spreads_of_dense_points, ObjDim.dim_eq_of_detCond, ZM.finrank_int_lattice, the law
% Out.Laws.mulN for the count, TateTower.tateIso_of_levels with weilIs_lower (Out.WeilSt: weilIs_symm,
% weilIs_pow_left and the compatibility of levels), T.exists_level_isometry_of_dense_geometric_points,
% T.torsionIsometry_primitiveRoot_bijective_mod_n and eq_one_of_weilIs_one. The count of torsion points and the
% properties of the commutator pairing are laws of abelian schemes in the Lean (Out.Laws), cited to Mumford.

\begin{lemma}[Alp\"oge]\label{cv-toroidal-abelian-dvr}
In Situation~\ref{cv-function-situation}, let $V$ be a discrete valuation ring with fraction field $K$,
let $G$ be an $\mathcal O_E$-semi-abelian scheme over $V$, proper over $V$, with line bundle $\mathcal L$,
and let $\nu:G_K\to G$ be a morphism of $\mathcal O_E$-semi-abelian schemes over
$\operatorname{Spec}K\to\operatorname{Spec}V$ (Definition~\ref{cv-model-toroidal-semiabelian}(2)).
\begin{enumerate}
\item If the line bundle of $G_K$ is ample and defines a principal polarization, then $\mathcal L$ is
  relatively ample and defines a principal polarization.
\item Let $n$ be a prime, $N\ge1$, and let $V$ be an algebra over $R_{n,N}$
  (Definition~\ref{cv-toroidal-family}). Suppose that $\mathcal L$ is relatively ample and defines a
  principal polarization. Let $\lambda:\Lambda\to G(V)$ be additive with $\lambda(N\Lambda)=0$, and let
  $\lambda_K:\Lambda\to G_K(K)$ be a map with $\nu\circ\lambda_K(a)=\lambda(a)|_{\operatorname{Spec}K}$
  for every $a\in\Lambda$. If $(G_K,\lambda_K)$ is a full-level object over $K$
  (Definition~\ref{cv-toroidal-family}(1)), then $(G,\lambda)$ is a full-level object over $V$.
\end{enumerate}
\end{lemma}
\leanref{Development/AbelianVarietyTheory/Proofs/ModuliProofs/RelAmplePrincipalOfGenericFibreDvrStatement.lean}{GlobalTheory.T.relAmple_principal_of_genericFibre_dvr}
\leanref{Development/AbelianVarietyTheory/Proofs/ModuliProofs/FullObjectOfGenericFibreDvrStatement.lean}{GlobalTheory.T.fullObject_of_genericFibre_dvr}
\begin{proof}
(1) The group scheme $G$ is commutative, smooth with geometrically connected fibres, and proper over $V$, so
it is an abelian scheme. As $\nu$ is a morphism of $\mathcal O_E$-semi-abelian schemes, it identifies $G_K$
with $G\times_V\operatorname{Spec}K$ as group schemes and the line bundle $\mathcal L_K$ of $G_K$ with
$\nu^*\mathcal L$. Lemma~\ref{cv-toroidal-abelian-dvr-structure}, applied to $A=G$ and $\varphi=\nu$, gives an
$\mathcal O_E$-semi-abelian scheme $G'$ over $V$ with line bundle $\mathcal L'$ and an isomorphism of group
schemes $e:G'\to G$, with $\mathcal L'$ relatively ample and defining a principal polarization. As the line
bundle of an $\mathcal O_E$-semi-abelian scheme, $\mathcal L'$ is trivialized along the zero section and
symmetric, and $e$ is finite, flat and surjective. Lemma~\ref{cv-toroidal-abelian-dvr-isogenous}, applied to
$P=G'$, $B=G$, $\Psi=e$, $\mathcal M=\mathcal L$, $B'=G_K$, $\varphi=\nu$ and $\mathcal M'=\mathcal L_K$,
shows that $\mathcal L$ is relatively ample. The line bundle $\mathcal L$ is trivialized along the zero
section and has a cubical structure, and $\mathcal L_K$ defines a principal polarization; so
Lemma~\ref{cv-toroidal-abelian-dvr-kernel}, applied to $B=G$, $B'=G_K$ and $\varphi=\nu$, shows that
$\mathcal L$ defines a principal polarization.

(2) The conditions on $G$, $\mathcal L$ and $\lambda$ other than (a) and (b) of
Definition~\ref{cv-toroidal-family}(1) hold by hypothesis. Let $k$ be an algebraically closed field of
characteristic zero and $x:\operatorname{Spec}k\to\operatorname{Spec}V$ a point. Then $V$ has characteristic
zero, as it maps to $k$, and so has $K$. Let $\Omega$ be an algebraic closure of $K$ and $y$ the point
$\operatorname{Spec}\Omega\to\operatorname{Spec}K$. As $(G_K,\lambda_K)$ is a full-level object over $K$, it
satisfies (a) and (b) at $y$, and by Lemma~\ref{cv-toroidal-abelian-dvr-spread}(1), applied to $\nu$ over
$\operatorname{Spec}K\to\operatorname{Spec}V$, $(G,\lambda)$ satisfies them at the composite point
$\operatorname{Spec}\Omega\to\operatorname{Spec}V$. Its image is the generic point of
$\operatorname{Spec}V$, which is dense. The scheme $\operatorname{Spec}V$ is noetherian, reduced and
connected, and $G$ is proper over $V$ with $\mathcal L$ relatively ample and defining a principal
polarization; so Lemma~\ref{cv-toroidal-abelian-dvr-spread}(2) shows that $(G,\lambda)$ satisfies (a) and (b)
at $x$.
\end{proof}
% Lean: (1) PolDvr.relAmple_kTriv_of_generic_fibre: T.exists_principallyPolarizedCM_extension_dvr (with
% HomExtSt), exists_polAV (G' as an Out.PolAV, with e' the identity), T.relAmple_of_generic_isogenous_dedekind
% for Psi = e' >> e, then exists_polAV for tK (G_K is proper by base change), SymmetricBundle.gen_of_isSymmetric
% and T.abelianScheme_cubical_bundle_exists_and_maps (the cubical package of L, whose fields other than the
% cubical structure are not used) and T.thetaGroupK_trivial_of_principal_genericFibre_dvr.
% (2) FullSpread.fullObj_of_genericFibre: GeomAt.push (o.geom at AlgebraicClosure K) and geomAt_of_dense with
% the dense singleton of the generic point.
% The lemmas cv-toroidal-abelian-dvr-* are written at the floor; no open step of this proof remains.

\subsection{The abelian locus and the moduli scheme}\label{subsection-cv-abelian-locus}

\begin{lemma}[Alp\"oge]\label{cv-toroidal-abelian-level}
In Situation~\ref{cv-function-situation}, let $n$ be a prime and $N\ge1$ an integer that is invertible
in $R_{n,N}$ (Definition~\ref{cv-toroidal-family}). Let $S$ be a locally noetherian scheme over
$R_{n,N}$ whose local rings are integrally closed domains, let $G$ be an $\mathcal O_E$-semi-abelian
scheme over $S$ (Definition~\ref{cv-model-toroidal-semiabelian}(2)) that is proper over $S$ and
quasi-projective over $R_{n,N}$, let $j:W\to S$ be an open immersion with dense image, let $G_W$ be an
$\mathcal O_E$-semi-abelian scheme over $W$ with a morphism $\nu:G_W\to G$ over $j$, and let
$\lambda_W:\Lambda\to G_W(W)$ be a map such that $(G_W,\lambda_W)$ is a full-level object over $W$
(Definition~\ref{cv-toroidal-family}(1)). Then there is a map $\lambda:\Lambda\to G(S)$ such that
$(G,\lambda)$ is a full-level object over $S$ and $\nu\circ\lambda_W(a)=\lambda(a)\circ j$ for every
$a\in\Lambda$.
\end{lemma}
\leanref{Development/AbelianVarietyTheory/Proofs/ModuliProofs/AbelianLocusOpenFamilyOverLevel.lean}{GlobalTheory.AbelianLocusNormalisation.exists_full}
\leanref{Development/AbelianVarietyTheory/Proofs/ModuliProofs/AbelianLocusOpenFamilyOverLevel.lean}{GlobalTheory.AbelianLocusNormalisation.exists_full_through}
\leanref{Development/AbelianVarietyTheory/Proofs/ModuliProofs/IsProjOverBaseChangeOfProperQuasiProjectiveStatement.lean}{GlobalTheory.T.isProjOver_baseChange_of_proper_quasiProjective}
\begin{proof}
The integer $N$ is invertible on $S$. The morphism $\nu$ identifies $G_W$ with $j^*G$ as group schemes,
and $\lambda_W$ is additive and vanishes on $N\Lambda$, as $(G_W,\lambda_W)$ is a full-level object. So
Lemma~\ref{cv-toroidal-abelian-torsion} gives an additive $\lambda:\Lambda\to G(S)$ vanishing on
$N\Lambda$ with $\nu\circ\lambda_W(a)=\lambda(a)\circ j$.

Covers. Every point $p$ of $S$ is covered, in the sense of Lemma~\ref{cv-toroidal-abelian-cover}, by a
morphism $\nu':G'\to G$ over some $f:S'\to S$ with a map $\lambda'$ such that
$\nu'\circ\lambda'(a)=\lambda(a)\circ f$ and $(G',\lambda')$ is a full-level object over $S'$. Indeed,
let $\xi$ be a generic point of an irreducible component of $S$ that specializes to $p$. Since $j$ has
dense image, $\xi$ lies in $j(W)$. If $\xi=p$, take $f=j$, $\nu'=\nu$ and $\lambda'=\lambda_W$. Otherwise
there are a complete discrete valuation ring $V$, with fraction field $K$, and a morphism
$f:\operatorname{Spec}V\to S$ that maps the generic point to $\xi$ and the closed point to $p$
\cite[Tag 054F]{Sta}, after replacing the discrete valuation ring given there by its completion. The
point $\operatorname{Spec}K\to S$ lands in the open set $j(W)$, so it is $j\circ x$ for a morphism
$x:\operatorname{Spec}K\to W$. Let $G_V$ be the pullback of $G$ along $f$, with the pullback $\lambda_V$
of $\lambda$, and let $G_K$ be the pullback of $G_V$ to $\operatorname{Spec}K$. Through $G_K\to G_V\to G$,
which lies over $j\circ x$, and $\nu$, the scheme $G_K$ is also the pullback of $G_W$ along $x$; with the
pullback $\lambda_K$ of $\lambda_W$ it is a full-level object over $K$, by
Lemma~\ref{cv-toroidal-abelian-pullback}. The group scheme $G_V$ is proper over $V$, so
Lemma~\ref{cv-toroidal-abelian-dvr}(1) shows that its line bundle is relatively ample and defines a
principal polarization. The map $\lambda_V$ is additive and vanishes on $N\Lambda$, and $G_K\to G_V$
carries $\lambda_K$ to $\lambda_V$, since $\nu\circ\lambda_W(a)=\lambda(a)\circ j$. So
Lemma~\ref{cv-toroidal-abelian-dvr}(2) shows that $(G_V,\lambda_V)$ is a full-level object over $V$, and
$G_V\to G$ over $f$ covers $p$.

Projectivity. Let $A$ be a ring and $y:\operatorname{Spec}A\to S$ a morphism, with base change $G_A$ of
$G$. The morphism $G_A=G\times_S\operatorname{Spec}A\to G\times_{R_{n,N}}\operatorname{Spec}A$ is an
immersion, and $G\times_{R_{n,N}}\operatorname{Spec}A$ has a quasi-compact immersion into a projective
space over $A$, since $G$ is quasi-projective over $R_{n,N}$. As $G_A$ is proper over $A$, the composite
immersion is closed, so $G_A$ is projective over $A$.

Lemma~\ref{cv-toroidal-abelian-cover}(1) now shows that the line bundle of $G$ is relatively ample,
Lemma~\ref{cv-toroidal-abelian-cover}(2) that it defines a principal polarization, and
Lemma~\ref{cv-toroidal-abelian-cover}(3) that $(G,\lambda)$ is a full-level object over $S$.
\end{proof}
% Lean: exists_full: T.torsionLevel_extends_from_dense_open_normal, exists_full_through (the covers:
% exists_mem_genericPoints_specializes, exists_complete_discreteValuationRing_of_specializes, exists_pull,
% exists_cancel, T.fullLevelObject_pullback_exists_and_full, T.relAmple_principal_of_genericFibre_dvr and
% T.fullObject_of_genericFibre_dvr), T.isProjOver_baseChange_of_proper_quasiProjective
% (isImmersion_pullback_lift_of_isPullback, isProjOver_of_isProper_of_isImmersion),
% T.relAmple_of_polarised_at_every_point, T.kernelTrivial_of_every_point_covered and
% T.fullLevelObject_of_covering_pullbacks.

\begin{lemma}[Alp\"oge]\label{cv-toroidal-abelian-valuative}
Let $R$ be a noetherian ring, let $Y$ be a locally noetherian scheme, separated over $R$, and let
$f:X\to Y$ be a morphism of schemes such that $X$ is quasi-projective over $R$. Suppose that for every
complete discrete valuation ring $V$ with fraction field $K$, every morphism
$x:\operatorname{Spec}K\to X$ whose image is the generic point of an irreducible component of $X$, and
every morphism $y:\operatorname{Spec}V\to Y$ with $f\circ x=y|_{\operatorname{Spec}K}$, there is
$x':\operatorname{Spec}V\to X$ with $x'|_{\operatorname{Spec}K}=x$. Then $f$ is proper.
\end{lemma}
\leanref{Development/AbelianVarietyTheory/Proofs/ModuliProofs/IsProperOfCompleteDVRLiftGenericStatement.lean}{GlobalTheory.T.isProper_of_completeDVR_lift_generic}
\leanref{Development/AbelianVarietyTheory/Proofs/ModuliProofs/IsProperOfCompleteDVRLiftGeneric.lean}{GlobalTheory.ValuativeGeneric.isProper_of_lift_generic}
\begin{proof}
Let $i:X\to\mathbf P^m_R$ be a quasi-compact immersion over $R$, and $P=Y\times_R\mathbf P^m_R$, which
is proper over $Y$, locally noetherian and separated. The morphism $h=(f,i):X\to P$ is a quasi-compact
immersion, because its composite with the projection to $\mathbf P^m_R$ is. Its image is closed once it
is stable under specialization \cite[Tag 05JL]{Sta}; then $h$ is a closed immersion, and $f$, the
composite of $h$ with the projection $P\to Y$, is proper.

Let $z$ be a specialization of a point $h(x_1)$. The point $x_1$ is a specialization of the generic
point $x_0$ of an irreducible component of $X$, so $h(x_0)$ specializes to $z$. If $h(x_0)=z$, there is
nothing to prove. Otherwise there are a complete discrete valuation ring $V$, with fraction field $K$,
and $g:\operatorname{Spec}V\to P$ mapping the generic point to $h(x_0)$ and the closed point to $z$
\cite[Tag 054F]{Sta}, after completion. As $h$ is an immersion, it induces an isomorphism of the residue
fields at $x_0$ and $h(x_0)$, so $g|_{\operatorname{Spec}K}=h\circ x$ for a morphism
$x:\operatorname{Spec}K\to X$ with image $x_0$. The hypothesis, applied to $x$ and to the composite $y$
of $g$ with the projection to $Y$, gives $x':\operatorname{Spec}V\to X$ extending $x$. The morphisms
$h\circ x'$ and $g$ agree on $\operatorname{Spec}K$, which is dense in the reduced scheme
$\operatorname{Spec}V$, and $P$ is separated; so $h\circ x'=g$, and $z=h(x'(\mathfrak m))$, where
$\mathfrak m$ is the closed point.
\end{proof}
% Lean: ValuativeGeneric.isProper_of_lift_generic: isClosedImmersion_of_lift
% (isClosedMap_iff_specializingMap, exists_mem_genericPoints_specializes,
% exists_complete_discreteValuationRing_of_specializes, exists_lift_of_surjectiveOnStalks,
% ext_of_isDominant).

\begin{lemma}[Alp\"oge]\label{cv-toroidal-abelian-proper}
In Situation~\ref{cv-function-situation}, let $R$ be a noetherian ring, let $Z$ be a scheme projective
over $R$ with an $\mathcal O_E$-semi-abelian scheme $G$ (Definition~\ref{cv-model-toroidal-semiabelian}(2)),
and let $M$ be a scheme quasi-projective over $R$ with an $\mathcal O_E$-semi-abelian scheme $G_M$ that
is proper over $M$. Let $U$ be an open subscheme of $Z$ that contains the abelian locus of $G$, let
$u:U\to M$ be a morphism over $R$, and let $G_U$ be an $\mathcal O_E$-semi-abelian scheme over $U$ with
morphisms $G_U\to G$ over the inclusion and $G_U\to G_M$ over $u$. Then $u$ is proper.
\end{lemma}
\leanref{Development/AbelianVarietyTheory/Proofs/ModuliProofs/AbelianLocusOpenFamilyOverLevel.lean}{GlobalTheory.AbelianLocusNormalisation.isProper_of_abLocus}
\leanref{Development/AbelianVarietyTheory/Proofs/ModuliProofs/IsAbelianSchemeOfCommonGenericFibre.lean}{GlobalTheory.isAbelianScheme_of_common_generic_fibre}
\begin{proof}
The scheme $Z$ is proper over $R$, and $M$ is separated and of finite type over $R$, so both are locally
noetherian. The scheme $U$ is quasi-projective over $R$, being a quasi-compact open subscheme of $Z$. We
apply Lemma~\ref{cv-toroidal-abelian-valuative} to $u$. Let $V$, $K$, $x:\operatorname{Spec}K\to U$ and
$y:\operatorname{Spec}V\to M$ be as there. As $Z$ is proper over $R$, the valuative criterion
\cite[Tag 0BX5]{Sta} gives $z:\operatorname{Spec}V\to Z$ that extends $\operatorname{Spec}K\to U\to Z$
and has the same composite to $\operatorname{Spec}R$ as $y$. Let $G_z$ and $G_y$ be the pullbacks of $G$
along $z$ and of $G_M$ along $y$, and $G_x$ the pullback of $G_U$ along $x$. Then $G_y$ is an abelian
scheme over $V$, as $G_M$ is proper over $M$, and through $G_U\to G$ and $G_U\to G_M$ the group scheme
$G_x$ is the generic fibre of both $G_z$ and $G_y$.

The resulting isomorphism of the generic fibres of $G_y$ and $G_z$ extends to a homomorphism
$F:G_y\to G_z$ by Theorem~\ref{cv-model-toroidal-extension}. As $G_y$ is proper over $V$ and $G_z$ is
separated, $F$ is proper, so its image is closed. It contains the generic fibre of $G_z$, and every point
of $G_z$ is a specialization of a point of the generic fibre, as $G_z$ is flat over $V$
\cite[Tag 03HV]{Sta}. So $F$ is surjective, and $G_z$ is universally closed over $V$, being the image of
the proper $G_y$. Hence $G_z$ is proper over $V$, an abelian scheme, and $z$ maps into the abelian locus
of $G$, since the abelian locus of a pullback is the preimage of the abelian locus. So $z$ factors
through $U$, which gives the required extension of $x$.
\end{proof}
% Lean: isProper_of_abLocus: T.isProper_of_completeDVR_lift_generic, valuative_lift_of_isProper,
% exists_pull, exists_comp, exists_cancel, isAbelianScheme_of_common_generic_fibre (with HomExtSt;
% grpPullback_iso and isAbelianScheme_of_generically_surjective, through
% surjective_of_generic_fibre_subset_range and UniversallyClosed.of_comp_surjective) and
% range_subset_abLocus_of_isProper.

\begin{lemma}[Alp\"oge]\label{cv-toroidal-abelian-fibres}
In Situation~\ref{cv-function-situation}, let $Z$ be a locally noetherian scheme, flat over
$\mathbf Z$, whose local rings are integrally closed domains, and let $G$ be an
$\mathcal O_E$-semi-abelian scheme over $Z$ (Definition~\ref{cv-model-toroidal-semiabelian}(2)), proper
over $Z$, whose line bundle $\mathcal L$ is relatively ample and defines a principal polarization. Let
$N\ge1$, and suppose that $G$ has an ample theta presentation of level $N$
(Definition~\ref{cv-toroidal-envelope}(1)). Let $f:Z\to Y$ be a proper morphism and $G\to G'$ a morphism
of $\mathcal O_E$-semi-abelian schemes over $f$. Then $f$ is finite.
\end{lemma}
\leanref{Development/AbelianVarietyTheory/Proofs/ModuliProofs/IsFiniteOfAmpleDeterminantFramesPullbackStatement.lean}{GlobalTheory.T.isFinite_of_ampleDeterminantFrames_pullback}
\leanref{Development/AbelianVarietyTheory/Proofs/ModuliProofs/IsFiniteOfAmpleDeterminantFramesPullback.lean}{GlobalTheory.FrameFinite.isFinite_of_ampFrames}
\begin{proof}
Let the presentation consist of $\mathcal A$, $(U_i)$, $n_0$, the integers $r$ and $s$, and the sections
$\vartheta_i$ and $b_{i,1},\dots,b_{i,n_0}$, with transition units $c_{ii'}$, matrices $M_{ii'}$ and
transition functions $g_{ii'}=c_{ii'}^r\det(M_{ii'})^s$ of $\mathcal A$.

Step 1. For every $i$, every ring $A$ and every morphism $\chi:\operatorname{Spec}A\to U_i$, the
pullbacks along $\chi$ of $\vartheta_i$ and of $b_{i,1},\dots,b_{i,n_0}$ are bases of the $A$-modules of
sections of $\mathcal L$ and of $\mathcal L^N$ over $G\times_Z\operatorname{Spec}A$. First, for every
$p\in Z$ and $a\ge1$ there are an open neighbourhood $V$ of $p$ and sections $e_1,\dots,e_{m'}$ of
$\mathcal L^a$ over $G_V$ whose pullbacks along every $\operatorname{Spec}A\to V$ form a basis of the
sections of $\mathcal L^a$ over $G_A$: over a connected affine neighbourhood $G$ is an abelian scheme of
some relative dimension $g$ with a principal polarization, and the direct image of $\mathcal L^a$ is
locally free of rank $a^g$ and commutes with base change \cite[\S\S5 and~16]{Mum70}. Second, let
$p\in U_i$. The local ring $\mathcal O_{Z,p}$ is a noetherian domain of characteristic zero, as $Z$ is
flat over $\mathbf Z$. If it is a field, let $W=\mathcal O_{Z,p}[[t]]$; otherwise let $W$ be the
completion of a discrete valuation ring that dominates it \cite[Tag 00PH]{Sta}. Then $W$ is a complete
discrete valuation ring of characteristic zero with a morphism $\operatorname{Spec}W\to U_i$ mapping the
closed point to $p$, and the generic point lies in the abelian locus, which is all of $Z$. By condition
(b) of Definition~\ref{cv-toroidal-envelope}(1), the pullbacks of $\vartheta_i$ and of the $b_{i,k}$ to
$W$ are bases. Write the pullbacks of $\vartheta_i$ (or of the $b_{i,k}$) to an affine neighbourhood
$V'\subset V\cap U_i$ of $p$ in terms of the $e_l$ for $a=1$ (or $a=N$), with a matrix $P$. Over $W$ both
are bases, so $P$ is square and $\det P$ maps to a unit of $W$; as $\operatorname{Spec}W\to V'$ maps the
closed point to $p$, $\det P$ does not vanish at $p$. Over the neighbourhood $D(\det P)$ of $p$ the
pullbacks are therefore bases after every base change. Being a basis is local on $\operatorname{Spec}A$,
which proves Step 1.

Step 2. Let $q:C\to Z$ be a morphism, $i_0$ an index and $q_0:C\to U_{i_0}$ a morphism with
$f\circ q_0=f\circ q$. Then $q^*\mathcal A$ has a nowhere vanishing section. Let $T$ be an affine scheme
and $t:T\to C$ a morphism such that $q\circ t$ maps into $U_i$. The morphism $G\to G'$ identifies the
pullbacks of $G$ along $q\circ t$ and along $q_0\circ t$, with their line bundles, with the pullback of
$G'$ along $f\circ q\circ t=f\circ q_0\circ t$. By Step 1 the pullbacks of $\vartheta_i$ and of
$\vartheta_{i_0}$ are then two bases of the same module of rank one, and those of the $b_{i,k}$ and of
the $b_{i_0,k}$ two bases of the same module of rank $n_0$. Let $c$ be the unit and $M'$ the invertible
matrix relating them, and $d_i(t)=c^r\det(M')^s$, a unit on $T$. These units are compatible with
pullback along morphisms $T'\to T$, and on $q^{-1}(U_i\cap U_{i'})$ condition (a) of
Definition~\ref{cv-toroidal-envelope}(1) gives $d_i=g_{ii'}\,d_{i'}$. So the $d_i$ glue to a nowhere
vanishing section of $q^*\mathcal A$.

Step 3. Let $y\in Y$, $\Omega$ an algebraic closure of its residue field and $Z_\Omega$ the fibre of $f$
over $\operatorname{Spec}\Omega\to Y$, a noetherian scheme proper over $\Omega$. Let $C$ be an
irreducible component of $Z_\Omega$ with its reduced structure, and $q:C\to Z$ the projection. As $C$ is
of finite type over $\Omega$, it has a point $c:\operatorname{Spec}\Omega\to C$; let $U_{i_0}$ contain
$q(c)$, and let $q_0$ be the composite of $C\to\operatorname{Spec}\Omega$ with $q\circ c$. Then
$f\circ q_0=f\circ q$, so $q^*\mathcal A$ is trivial by Step 2. As $\mathcal A$ is ample, every point of
$Z$ lies in the nonvanishing locus, which is affine, of a section of some $\mathcal A^{d}$ with $d\ge1$.
On $C$ these sections become functions $h$ whose nonvanishing loci $D(h)$ cover $C$. Each $D(h)$ is
affine, as the preimage of an affine open subscheme under the affine morphism $C\to Z$, and it is open
and closed in $C$, as $C$ is proper over $\Omega$. So each $D(h)$ is a closed
affine subscheme of a proper scheme over $\Omega$, hence finite over $\Omega$, and $C$ is finite. So the
fibre of $f$ over $y$ is finite. Thus $f$ is locally quasi-finite and proper, hence finite
\cite[Tag 02LS]{Sta}.
\end{proof}
% Lean: FrameFinite.isFinite_of_ampFrames: FrameBasis.basisAlong_of_ampFrames (Step 1: exists_trait with
% exists_trait_ring, exists_stableBas_local from Laws2.secBaseChangeAny, stableBas_of_traits),
% exists_gen_pull (Step 2: exists_frames, exists_det, Det.mul, exists_det_trans and exists_gen_of_affine)
% and AmpleFibre.locallyQuasiFinite_of_trivial_on_fibres (Step 3: exists_integral_closedImmersion,
% exists_fn, MulFinite.locallyQuasiFinite_of_cover), with IsFinite.of_isProper_of_locallyQuasiFinite.

\begin{lemma}[Alp\"oge]\label{cv-toroidal-abelian-finite}
In Situation~\ref{cv-function-situation}, let $R$ be a ring and $Z$ a locally noetherian scheme,
separated over $R$ and flat over $\mathbf Z$, whose local rings are integrally closed domains, with an
$\mathcal O_E$-semi-abelian scheme $G$ (Definition~\ref{cv-model-toroidal-semiabelian}(2)) that has an
ample theta presentation of level $N\ge1$ (Definition~\ref{cv-toroidal-envelope}(1)). Let $M$ be a
scheme over $R$ with an $\mathcal O_E$-semi-abelian scheme $G_M$, proper over $M$, whose line bundle is
relatively ample and defines a principal polarization. Let $U$ be an open subscheme of $Z$, let
$u:U\to M$ be a proper morphism over $R$, and let $G_U$ be an $\mathcal O_E$-semi-abelian scheme over $U$
with morphisms $G_U\to G$ over the inclusion and $G_U\to G_M$ over $u$. Then $u$ is finite.
\end{lemma}
\leanref{Development/AbelianVarietyTheory/Proofs/ModuliProofs/AbelianLocusOpenFamilyOverLevel.lean}{GlobalTheory.AbelianLocusNormalisation.isFinite_of_isProper}
\leanref{Development/AbelianVarietyTheory/Proofs/ModuliProofs/AmpleDeterminantFramesPullbackOfIsAffineHomStatement.lean}{GlobalTheory.T.ampleDeterminantFrames_pullback_of_isAffineHom}
\begin{proof}
A proper morphism is finite if it is affine, so it suffices that $u^{-1}(V)\to V$ is finite for every
affine open $V$ of $M$. The inclusion $g:u^{-1}(V)\to Z$ is an affine morphism: it is the composite of
the morphism $u^{-1}(V)\to V\times_RZ$ with components the restriction of $u$ and $g$, which is a proper
immersion, hence a closed immersion, with the projection $V\times_RZ\to Z$, a base change of the affine
morphism $V\to\operatorname{Spec}R$.

Let $G_g$ be the pullback of $G$ along $g$, with morphisms to $G$ over $g$ and, through $G_U$, to $G_M$
over $u^{-1}(V)\to V\to M$. The ample theta presentation pulls back along the affine morphism $g$: the
line bundle, its trivializations, the sections and the transition data pull back; the nonvanishing locus
of the pullback of a section is the preimage of its nonvanishing locus, hence affine; and condition (b)
holds, since a morphism from $\operatorname{Spec}W$ to $u^{-1}(V)$ whose generic point maps into the
abelian locus of $G_g$ gives one to $Z$ whose generic point maps into the abelian locus of $G$. As $G_g$
is also the pullback of $G_M$, it is proper over $u^{-1}(V)$, and its line bundle is relatively ample
\cite[Tag 0893]{Sta} and defines a principal polarization, as in the proof of
Lemma~\ref{cv-toroidal-abelian-pullback}. The scheme $u^{-1}(V)$ is open in $Z$, so it is locally
noetherian, flat over $\mathbf Z$ and has integrally closed domains as local rings. So
Lemma~\ref{cv-toroidal-abelian-fibres}, applied to $u^{-1}(V)\to V$ and the morphism from $G_g$ to the
pullback of $G_M$ to $V$, shows that $u^{-1}(V)\to V$ is finite.
\end{proof}
% Lean: isFinite_of_isProper: IsFinite.iff_isProper_and_isAffineHom, isAffineHom_of_isProper,
% T.ampleDeterminantFrames_pullback_of_isAffineHom (AmpPull.ampFrames_pull = exists_pull: chartPullLinComb,
% nonzeroLocus_chartSec, exists_rel_down, lin_down, basis_down, abLocus_pullback), pol_of_map and
% T.isFinite_of_ampleDeterminantFrames_pullback.

\begin{lemma}[Alp\"oge]\label{cv-toroidal-abelian-birational}
Let $M$ be a locally noetherian scheme, let $u:U\to M$ be a finite morphism with $U$ reduced, and let
$W$ be a dense open subscheme of $M$ such that $u^{-1}(W)$ is dense in $U$ and $u^{-1}(W)\to W$ is an
isomorphism. Let $x\in M$ be a point whose local ring $\mathcal O_{M,x}$ is an integrally closed domain.
Then $x$ has an open neighbourhood $V$ such that $u^{-1}(V)\to V$ is an isomorphism.
\end{lemma}
\leanref{Development/AbelianVarietyTheory/Proofs/ModuliProofs/IsIsoNearNormalPointOfFiniteBirationalStatement.lean}{GlobalTheory.T.isIso_near_normal_point_of_finite_birational}
\leanref{Development/AbelianVarietyTheory/Proofs/ModuliProofs/IsIsoNearNormalPointOfFiniteBirational.lean}{GlobalTheory.FiniteBirationalNear.main}
\begin{proof}
First, $x$ has an integral open neighbourhood. Let $\operatorname{Spec}A$ be an affine open
neighbourhood, with $A$ noetherian and $x$ corresponding to the prime $\mathfrak p$. The kernel of
$A\to A_{\mathfrak p}$ is finitely generated, so some $a\notin\mathfrak p$ kills it; then $A_a$ embeds in
the domain $A_{\mathfrak p}$, and $D(a)$ is integral.

So let $V\ni x$ be integral, and let $u_V:U_V=u^{-1}(V)\to V$ be the restriction, which is finite. The
open set $W\cap V$ is nonempty, as $W$ is dense, and $u^{-1}(W\cap V)\cong W\cap V$ is integral. It is
dense in $U_V$, since $u^{-1}(W)$ is dense in $U$; so $U_V$ is irreducible, and, being reduced, integral.
Thus $u_V$ is a finite morphism of integral schemes that is an isomorphism over a dense open subset, so
it is dominant and induces an isomorphism of function fields.

Let $\operatorname{Spec}A\subset V$ be an affine open neighbourhood of $x$, with $x$ corresponding to
$\mathfrak p$, and $B$ the ring of $u^{-1}(\operatorname{Spec}A)$, a finite $A$-module. The map $A\to B$
is injective and $B$ embeds in the common fraction field of $A$ and $B$. Each $b\in B$ is integral over
$A$, hence over the integrally closed domain $A_{\mathfrak p}$, so $b\in A_{\mathfrak p}$. Applying this to
finitely many generators of the $A$-module $B$ gives $r\notin\mathfrak p$ with $rB\subset A$. Then
$A_r\to B_r$ is bijective, so $u$ is an isomorphism over the neighbourhood $D(r)$ of $x$.
\end{proof}
% Lean: FiniteBirationalNear.main: exists_isIntegral_nhds (exists_away_isDomain), birational_of_isIso_morphismRestrict,
% exists_isIso_of_isIntegral (exists_mul_mem_range, exists_common_smul_mem).

\begin{lemma}[Alp\"oge]\label{cv-toroidal-abelian-locus}
In Situation~\ref{cv-function-situation}, let $n$ be a prime, $N$ a power of $n$ and $m\ge1$, let
$(\mathcal M,G_{\mathcal M},\lambda)$ be a full-level family over $R_{n,N}$
(Definition~\ref{cv-toroidal-family}), let $(Z,G,W,j)$ be a theta envelope of level $m^2$ of
$(\mathcal M,G_{\mathcal M})$ over $R_{n,N}$ (Definition~\ref{cv-toroidal-envelope}), and let $U$ be the
abelian locus of $G$, an open subscheme of $Z$. Then there are a finite surjective morphism
$u:U\to\mathcal M$ over $R_{n,N}$, an $\mathcal O_E$-semi-abelian scheme $G_U$ over $U$ with morphisms
$G_U\to G$ over the inclusion and $G_U\to G_{\mathcal M}$ over $u$, and an open subscheme $W_0$ of
$\mathcal M$ that contains every point over the generic point of $\operatorname{Spec}R_{n,N}$ and over
which $u$ is an isomorphism. Moreover, $u\circ j'$ is the inclusion of $W$ in $\mathcal M$ for every
morphism $j':W\to U$ whose composite with the inclusion of $U$ is $j$.
\end{lemma}
\leanref{Development/AbelianVarietyTheory/Proofs/ModuliProofs/AbelianLocusOpenFamilyOverLevelStatement.lean}{GlobalTheory.T.abelian_locus_open_family_over_level}
\leanref{Development/AbelianVarietyTheory/Proofs/ModuliProofs/AbelianLocusOpenFamilyOverLevel.lean}{GlobalTheory.AbelianLocusNormalisation.exists_finite_map}
\begin{proof}
Write $R=R_{n,N}$, with structure morphisms $s:Z\to\operatorname{Spec}R$ and
$s_{\mathcal M}:\mathcal M\to\operatorname{Spec}R$, let $\iota_U$ and $\iota_W$ be the inclusions, and let
$G_W$ be the $\mathcal O_E$-semi-abelian scheme over $W$ of Definition~\ref{cv-toroidal-envelope}(2),
with morphisms $\mu_1:G_W\to G_{\mathcal M}$ over $\iota_W$ and $\mu_2:G_W\to G$ over $j$. The ring $R$
is of finite type over $\mathbf Z$, hence noetherian, and $n$, hence $N$, is invertible in $R$. The
scheme $Z$ is projective over $R$, so proper over $R$ and noetherian, and $\mathcal M$ is
quasi-projective over $R$, so separated over $R$ and locally noetherian. The open subscheme $U$ of the
normal scheme $Z$ is normal, hence reduced.

Let $G_U$ be the pullback of $G$ to $U$, with its morphism $G_U\to G$ over $\iota_U$. By condition (a)
of Definition~\ref{cv-model-toroidal-semiabelian}(3), $G$ is proper over its abelian locus $U$, so $G_U$
is proper over $U$. It is quasi-projective over $R$, since $G$ is and $G_U\to G$ is a quasi-compact open
immersion, as a base change of $\iota_U$.

The group scheme $G_W$ is proper over $W$, as the pullback of $G_{\mathcal M}$, which is proper because
$(G_{\mathcal M},\lambda)$ is a full-level object. Through $\mu_2$ it is also the pullback of $G$ along
$j$, and the abelian locus of a pullback is the preimage of the abelian locus. So $j(W)\subset U$, and
$j=\iota_U\circ j'$ for an open immersion $j':W\to U$ with dense image. As $G_U\to G$ identifies $G_U$
with the pullback of $G$, $\mu_2$ is the composite of a morphism $\nu:G_W\to G_U$ over $j'$ with
$G_U\to G$.

Let $\lambda_W:\Lambda\to G_W(W)$ be the pullback of $\lambda$ along $\mu_1$, so that
$\mu_1\circ\lambda_W(a)=\lambda(a)\circ\iota_W$. By Lemma~\ref{cv-toroidal-abelian-pullback},
$(G_W,\lambda_W)$ is a full-level object over $W$, regarded over $R$ through
$s_{\mathcal M}\circ\iota_W=s\circ j$. Lemma~\ref{cv-toroidal-abelian-level}, applied to $U$ over $R$
through $s\circ\iota_U$, to $G_U$, $j'$, $\nu$ and $\lambda_W$, gives $\lambda_U:\Lambda\to G_U(U)$ such
that $(G_U,\lambda_U)$ is a full-level object over $U$ and $\nu\circ\lambda_W(a)=\lambda_U(a)\circ j'$ for
every $a\in\Lambda$.

The scheme $U$ is reduced, locally noetherian and flat over $R$, being open in $Z$. By
Definition~\ref{cv-toroidal-family}(2)(a) there is a unique morphism $u:U\to\mathcal M$ over $R$ with a
morphism $\mu:G_U\to G_{\mathcal M}$ over $u$ carrying $\lambda_U$ to $\lambda$. Then $u\circ j'=\iota_W$.
Indeed, $W$ is reduced, locally noetherian and flat over $R$, being open in $\mathcal M$; the morphisms
$\iota_W$, with $\mu_1$, and $u\circ j'$, with $\mu\circ\nu$, both carry $\lambda_W$ to $\lambda$; and
$(G_W,\lambda_W)$ is a full-level object, so the uniqueness in Definition~\ref{cv-toroidal-family}(2)(a)
applies.

Lemma~\ref{cv-toroidal-abelian-proper}, with $G_U\to G$ and $\mu$, shows that $u$ is proper. The scheme
$Z$ is separated over $R$, and flat over $\mathbf Z$, since it is flat over $R$ and $R\subset\mathbf C$
is torsion-free. The line bundle of $G_{\mathcal M}$ is relatively ample and defines a principal
polarization, and $G$ has an ample theta presentation of level $m^2\ge1$. So
Lemma~\ref{cv-toroidal-abelian-finite} shows that $u$ is finite.

Over $W$. Let $u_W:u^{-1}(W)\to W$ be the restriction of $u$. Since $u\circ j'=\iota_W$, $j'$ factors
through the open subscheme $u^{-1}(W)$ as a morphism $\sigma:W\to u^{-1}(W)$ with $u_W\circ\sigma=1$, and
$\sigma$ has dense image. The morphisms $\sigma\circ u_W$ and $1$ of $u^{-1}(W)$ have the same composite
with the separated morphism $u_W$ and agree after composition with $\sigma$; as $u^{-1}(W)$ is reduced,
they are equal. So $u_W$ is an isomorphism with inverse $\sigma$. Moreover $u^{-1}(W)$ contains $j'(W)$,
so it is dense in $U$.

Surjectivity. The image of $u$ is closed, as $u$ is proper, and contains the dense subset $W$ of
$\mathcal M$. So $u$ is surjective.

The open subscheme $W_0$. Let $F$ be the fraction field of $R$ and $x$ a point of $\mathcal M$ over the
generic point of $\operatorname{Spec}R$. The projection $\mathcal M\otimes_RF\to\mathcal M$ is flat and
surjective on local rings, so it identifies the local ring of $\mathcal M$ at $x$ with a local ring of
$\mathcal M\otimes_RF$, which is smooth over $F$. So this local ring is regular \cite[Tag 056S]{Sta},
hence an integrally closed domain. Lemma~\ref{cv-toroidal-abelian-birational}, with $u$ finite, $U$
reduced and $u_W$ an isomorphism, gives an open neighbourhood of $x$ over which $u$ is an isomorphism.
Let $W_0$ be the union of all open subsets of $\mathcal M$ over which $u$ is an isomorphism. Being an
isomorphism is local on the target, so $u$ is an isomorphism over $W_0$, and $W_0$ contains every point
over the generic point of $\operatorname{Spec}R$.

Finally, let $j'':W\to U$ satisfy $\iota_U\circ j''=j$. Then $j''=j'$, as $\iota_U$ is a monomorphism, so
$u\circ j''=\iota_W$.
\end{proof}
% Lean: AbelianLocusNormalisation.exists_finite_map: exists_pull (G_U), C.good.1 (G_U proper),
% range_subset_abLocus_of_isProper (j(W) in U), exists_cancel (nu), fullLevelObject_pullback_exists_and_full
% (oW), exists_full (Lemma cv-toroidal-abelian-level), F.rep (u and hju), isProper_of_abLocus,
% isFinite_of_isProper, isIso_restrict_of_section, the closed dense image, normal_stalk_of_generic
% (isReduced_isNormalScheme_isLocallyNoetherian_of_smooth_over_field and
% isIso_stalkMap_of_flat_of_surjectiveOnStalks), T.isIso_near_normal_point_of_finite_birational and
% exists_isIso_restrict.

\begin{definition}\label{cv-toroidal-moduli-object}
In Situation~\ref{cv-function-situation}, let $\Lambda$ be the lattice of the PEL datum with its
alternating form $\langle\,,\rangle$, and let $l=(M,\mathfrak M,N)$ be a level
(Definition~\ref{cv-model-toroidal-semiabelian}(1)).
\begin{enumerate}
\item A \emph{family of similitudes} of $\Lambda$ is a family of $\mathcal O_E$-linear maps
  $\omega_m:\Lambda\to\Lambda$ and integers $\nu_m$, $m\ge1$, such that for all $m,m'\ge1$ the map
  $\omega_{mm'}$ is congruent to $\omega_m$ modulo $m\Lambda$, $\omega_m$ induces a bijection of
  $\Lambda/m\Lambda$, and $\langle\omega_m(a),\omega_m(b)\rangle\equiv\nu_m\langle a,b\rangle$ modulo $m$
  for all $a,b\in\Lambda$.
\item Let $G$ be an $\mathcal O_E$-semi-abelian scheme over a scheme $S$
  (Definition~\ref{cv-model-toroidal-semiabelian}(2)), with line bundle $\mathcal L$, and let
  $x:\operatorname{Spec}k\to S$ be a point with $k$ an algebraically closed field. A \emph{Tate system} of
  $G$ at $x$ is a family of roots of unity $\zeta'_m\in k$ and $\mathcal O_E$-linear homomorphisms
  $\iota_m:\Lambda\to G_x(k)$, $m\ge1$, such that $\zeta'_m$ is primitive of order $m$,
  $(\zeta'_{mm'})^{m'}=\zeta'_m$, $\iota_m$ has kernel $m\Lambda$ and image the $m$-torsion,
  $m'\iota_{mm'}=\iota_m$, and $e_m(\iota_m(a),\iota_m(b))=(\zeta'_m)^{\langle a,b\rangle}$ for the
  commutator pairing $e_m$ of the theta group of $\mathcal L^m$. These are the data of condition (b) of
  Definition~\ref{cv-toroidal-family}(1).
\item Let $G$ be an $\mathcal O_E$-semi-abelian scheme over a scheme $S$ and let $T$ be a scheme over $S$.
  An \emph{$\mathfrak M$-level structure} of $G$ over $T$ is an additive map $\chi:\Lambda\to G(T)$ that
  vanishes on $\mathfrak M\Lambda$, such that at every point $z:\operatorname{Spec}k\to T$ with $k$
  algebraically closed, $\chi(ca)(z)=c\,\chi(a)(z)$ for $c\in\mathcal O_E$ and $a\in\Lambda$, and
  $\chi(a)(z)=0$ only if $a\in\mathfrak M\Lambda$.
\item Let $S$ be a scheme over $K_l$. An \emph{object at level $l$} over $S$ is a triple $(G,\eta,\eta')$ of
  an $\mathcal O_E$-semi-abelian scheme $G$ over $S$, proper over $S$, whose line bundle is relatively
  ample and defines a principal polarization, and additive maps $\eta,\eta':\Lambda\to G(S)$ that vanish
  on $\mathfrak M\Lambda$ and on $N\Lambda$, such that at every point $x:\operatorname{Spec}k\to S$ with $k$
  algebraically closed, with $\beta:K_l\to k$ the induced homomorphism, condition (a) of
  Definition~\ref{cv-toroidal-family}(1) holds and $G$ has a Tate system at $x$ with
  $\zeta'_N=\beta(\zeta_N)$, $\iota_N(a)=\eta'(a)(x)$ for $a\in\Lambda$, and $\iota_M(a)=\eta(a)(x)$ for
  $a\in\bar{\mathfrak M}\Lambda$. A \emph{morphism} of objects at level $l$ over a morphism $f:S'\to S$ is a
  morphism of $\mathcal O_E$-semi-abelian schemes over $f$ that carries $\eta$ and $\eta'$ to $\eta$ and
  $\eta'$.
\item The moduli scheme $X_l$, with structure morphism $s_l:X_l\to\operatorname{Spec}K_l$, carries its
  universal object $(\mathcal A_l,\eta,\eta')$ at level $l$, and it represents objects at level $l$ on
  reduced schemes: for every reduced scheme $S$ over $K_l$ and every object at level $l$ over $S$ there is
  a unique morphism $f:S\to X_l$ over $K_l$ with a morphism of objects to $(\mathcal A_l,\eta,\eta')$
  over $f$.
\end{enumerate}
\end{definition}
\leanref{ReciprocityTheory/ReductiveGroups/ArithmeticGeometry/BirationalGeometry/ToroidalCompactification/HeckeKer.lean}{ReciprocityTheory.ChartRestriction.G.Geo.Tor.Sim}
\leanref{ReciprocityTheory/ReductiveGroups/ArithmeticGeometry/BirationalGeometry/ToroidalCompactification/TwoWeightEqOneAddRoot.lean}{ReciprocityTheory.ChartRestriction.G.Geo.Tor.TateIso}
\leanref{Development/AbelianVarietyTheory/Proofs/ModuliProofs/ModuliFiniteSurjectiveOverGenericFamily.lean}{GlobalTheory.ModFam.IsLev}
\leanref{ReciprocityTheory/ReductiveGroups/ArithmeticGeometry/BirationalGeometry/ToroidalCompactification/TwoWeightEqOneAddRoot.lean}{ReciprocityTheory.ChartRestriction.G.Geo.Tor.Obj}
\leanref{ReciprocityTheory/ReductiveGroups/ArithmeticGeometry/BirationalGeometry/ToroidalCompactification/TwoWeightEqOneAddRoot.lean}{ReciprocityTheory.ChartRestriction.G.Geo.Tor.ObjMap}

\begin{lemma}[Alp\"oge]\label{cv-toroidal-moduli-lift}
In Situation~\ref{cv-function-situation}, let $\Lambda$ be the lattice of the PEL datum with its
alternating form $\langle\,,\rangle$, let $n\ge1$, let $f:\Lambda\to\Lambda$ be an $\mathcal O_E$-linear
map, and let $\nu$ be an integer prime to $n$ with $\langle f(a),f(b)\rangle\equiv\nu\langle a,b\rangle$
modulo $n$ for all $a,b\in\Lambda$. Then there is a family of similitudes $(\omega_m,\nu_m)$ of $\Lambda$
(Definition~\ref{cv-toroidal-moduli-object}(1)) with $\omega_n\equiv f$ modulo $n\Lambda$ and
$\nu_n\equiv\nu$ modulo $n$.
\end{lemma}
\leanref{Development/AbelianVarietyTheory/Proofs/ModuliProofs/ModNOfDvdModulus.lean}{GlobalTheory.exists_sim_lift}
\begin{proof}
For $m\ge1$, call a pair $(f,\nu)$ of an $\mathcal O_E$-linear map $f:\Lambda\to\Lambda$ and an integer
$\nu$ a similitude modulo $m$ if $\nu$ is prime to $m$ and $\langle f(a),f(b)\rangle\equiv\nu\langle
a,b\rangle$ modulo $m$ for all $a,b$. Let $w_0=(1+\sqrt{-d})/2$, so that $\bar w_0=1-w_0$.

Step 1: a similitude $(f,\nu)$ modulo $m$ is congruent modulo $m\Lambda$ to a similitude $(f',\nu)$ modulo
$m^2$. Write $\langle f(a),f(b)\rangle-\nu\langle a,b\rangle=mB(a,b)$. Then $B$ is biadditive, skew, and
$B(ca,b)=B(a,\bar cb)$ for $c\in\mathcal O_E$. As the form is perfect, $B(a,b)=\langle u(a),b\rangle$ for
an $\mathcal O_E$-linear map $u$, and $\langle a,u(b)\rangle=\langle u(a),b\rangle$ because $B$ is skew.
Let $v=w_0u$. Then $\langle a,v(b)\rangle=\langle \bar w_0a,u(b)\rangle=\langle(1-w_0)u(a),b\rangle$, so
$\langle v(a),b\rangle+\langle a,v(b)\rangle=B(a,b)$. Choose integers $\kappa,t$ with $\kappa\nu+tm=1$ and
let $f'=f-m\kappa\,f\circ v$. Modulo $m^2$, using $\langle f(x),f(y)\rangle\equiv\nu\langle x,y\rangle$
modulo $m$ in the cross terms,
\[
 \langle f'(a),f'(b)\rangle\equiv\nu\langle a,b\rangle+mB(a,b)-m\kappa\nu B(a,b)
 =\nu\langle a,b\rangle+m^2tB(a,b)\equiv\nu\langle a,b\rangle.
\]

Step 2: if $r$ is prime to $m$, a similitude $(f,\nu)$ modulo $m$ is congruent modulo $m$ to a similitude
modulo $mr$. Choose integers $x,y$ with $xm+yr=1$ and put $f'=yrf+xm$ and $\nu'=yr\nu+xm$. Modulo $m$,
$(f',\nu')$ is $(f,\nu)$, and modulo $r$ it is $(1,1)$; both are similitudes, so $(f',\nu')$ is a
similitude modulo $mr$.

Step 3: by induction on the number of prime factors of $r$, using Step 1 for a prime factor that divides
$m$ (then $m^2$ is divisible by $m$ times that prime) and Step 2 otherwise, every similitude $(f,\nu)$
modulo $m$ is congruent modulo $m$, with $\nu$ modulo $m$, to a similitude modulo $mr$, for every
$r\ge1$.

Step 4: starting from $(f,\nu)$, Step 3 gives similitudes $(g_k,\mu_k)$ modulo $n\cdot k!$ for $k\ge0$,
with $(g_0,\mu_0)=(f,\nu)$ and $(g_{k+1},\mu_{k+1})$ congruent to $(g_k,\mu_k)$ modulo $n\cdot k!$. Put
$\omega_m=g_m$ and $\nu_m=\mu_m$. Since $m$ divides $n\cdot m!$, $(\omega_m,\nu_m)$ is a similitude
modulo $m$, and $\omega_{mm'}\equiv\omega_m$ modulo $n\cdot m!$, hence modulo $m$. Also
$\omega_n\equiv f$ modulo $n$ and $\nu_n\equiv\nu$ modulo $n$.

Finally, a similitude $(f,\nu)$ modulo $m$ induces a bijection of $\Lambda/m\Lambda$. Let $u$ be the
$\mathcal O_E$-linear map with $\langle u(a),b\rangle=\langle a,f(b)\rangle$, and let $\kappa\nu\equiv1$
modulo $m$. Then $\langle\kappa u(f(a)),b\rangle=\kappa\langle f(a),f(b)\rangle\equiv\langle a,b\rangle$
modulo $m$ for all $b$, so $\kappa u(f(a))\equiv a$ modulo $m\Lambda$, as the form is perfect. So $f$ is
injective on the finite group $\Lambda/m\Lambda$, hence bijective.
\end{proof}
% Lean: exists_sim_lift: SimLift.step_sq (Step 1, with exists_rep for u and v = w0 . u), step_coprime
% (Step 2), step (Step 3, induction_on_primes), the sequence s k modulo N * k! (Step 4), and exists_inv
% for the bijection.

\begin{lemma}[Alp\"oge]\label{cv-toroidal-moduli-similitude}
In Situation~\ref{cv-function-situation}, suppose $p+q\ge1$, let $\Lambda$ be the lattice of the PEL datum
with its alternating form $\langle\,,\rangle$, and let $l=(M,\mathfrak M,N)$ be a level.
\begin{enumerate}
\item Let $(\omega_m,\nu_m)$ be a family of similitudes of $\Lambda$
  (Definition~\ref{cv-toroidal-moduli-object}(1)). Then $\nu_m$ is prime to $m$ and
  $\nu_{mm'}\equiv\nu_m$ modulo $m$ for all $m,m'\ge1$. If $((\zeta'_m),(\iota_m))$ is a Tate system of an
  $\mathcal O_E$-semi-abelian scheme at a point (Definition~\ref{cv-toroidal-moduli-object}(2)), then
  so is $(((\zeta'_m)^{\nu_m}),(\iota_m\circ\omega_m))$.
\item Let $A\ge1$ be an integer prime to $M$, and let $f:\Lambda\to\Lambda$ be an $\mathcal O_E$-linear map
  such that $f(a)\in\mathfrak M\Lambda$ only if $a\in\mathfrak M\Lambda$. Then there is a family of
  similitudes $(\omega_m,\nu_m)$ of $\Lambda$ with $\omega_M(a)-f(a)\in\mathfrak M\Lambda$ and
  $\omega_A(a)-a\in A\Lambda$ for all $a\in\Lambda$, and $\nu_A\equiv1$ modulo $A$.
\end{enumerate}
\end{lemma}
\leanref{Development/AbelianVarietyTheory/Proofs/ModuliProofs/ExistsFormEqOne.lean}{GlobalTheory.SimAlg.sim_group}
\leanref{Development/AbelianVarietyTheory/Proofs/ModuliProofs/ExistsObjTwistBySimilitude.lean}{GlobalTheory.LevelTwist.twist}
\leanref{Development/AbelianVarietyTheory/Proofs/ModuliProofs/ExistsSimilitudeOfAutModLevelIdealStatement.lean}{GlobalTheory.T.exists_similitude_of_aut_mod_levelIdeal}
\leanref{Development/AbelianVarietyTheory/Proofs/ModuliProofs/ExistsSimilitudeOfAutModLevelIdeal.lean}{GlobalTheory.exists_sim_of_aut_mod}
\begin{proof}
(1) As $p+q\ge1$, $\Lambda\ne0$, and since the form is perfect there are $a_0,b_0\in\Lambda$ with
$\langle a_0,b_0\rangle=1$. Then $\nu_{mm'}\equiv\langle\omega_{mm'}(a_0),\omega_{mm'}(b_0)\rangle
\equiv\langle\omega_m(a_0),\omega_m(b_0)\rangle\equiv\nu_m$ modulo $m$. If $h$ inverts $\omega_m$ modulo
$m\Lambda$, then $1=\langle a_0,b_0\rangle\equiv\langle\omega_mh(a_0),\omega_mh(b_0)\rangle
\equiv\nu_m\langle h(a_0),h(b_0)\rangle$ modulo $m$, so $\nu_m$ is prime to $m$. Hence
$(\zeta'_m)^{\nu_m}$ is primitive of order $m$, and
$(\zeta'_{mm'})^{m'\nu_{mm'}}=(\zeta'_m)^{\nu_{mm'}}=(\zeta'_m)^{\nu_m}$. The map $\iota_m\circ\omega_m$ is
$\mathcal O_E$-linear with kernel $m\Lambda$ and image the $m$-torsion, since $\omega_m$ induces a
bijection of $\Lambda/m\Lambda$; $m'\iota_{mm'}\omega_{mm'}=\iota_m\omega_{mm'}=\iota_m\omega_m$; and
$e_m(\iota_m\omega_m(a),\iota_m\omega_m(b))=(\zeta'_m)^{\langle\omega_m(a),\omega_m(b)\rangle}
=((\zeta'_m)^{\nu_m})^{\langle a,b\rangle}$.

(2) Step 1. The ideal $\mathfrak M$ is nonzero, since $\mathfrak M\bar{\mathfrak M}=M\mathcal O_E$, so
$\Lambda/\mathfrak M\Lambda$ is finite. By hypothesis $f$ induces an injective, hence bijective,
endomorphism of it. As $\Lambda$ is projective, there is an $\mathcal O_E$-linear $h:\Lambda\to\Lambda$
with $h(f(a))-a\in\mathfrak M\Lambda$ for all $a$.

Step 2. As the form is perfect, there is a map $u:\Lambda\to\Lambda$ with
$\langle u(a),b\rangle=\langle a,h(b)\rangle$. It is $\mathcal O_E$-linear:
$\langle u(ca),b\rangle=\langle a,\bar c\,h(b)\rangle=\langle a,h(\bar cb)\rangle
=\langle u(a),\bar cb\rangle=\langle cu(a),b\rangle$.

Step 3. Since $\mathfrak M+\bar{\mathfrak M}=\mathcal O_E$, write $1=e+e'$ with $e\in\mathfrak M$ and
$e'\in\bar{\mathfrak M}$, and let $f_1=e'f+eu$. Then $\langle f_1(a),f_1(b)\rangle\equiv\langle a,b\rangle$
modulo $M$. Indeed, $\langle ex,ey\rangle=\langle x,\bar eey\rangle$ and $\langle e'x,e'y\rangle$ are
divisible by $M$, since $\bar ee$ and $\bar e'e'$ lie in $\mathfrak M\bar{\mathfrak M}=M\mathcal O_E$.
Next, $\langle eu(a),e'f(b)\rangle=\langle a,\bar ee'h(f(b))\rangle$, and
$\bar e(h(f(b))-b)\in\bar{\mathfrak M}\mathfrak M\Lambda=M\Lambda$, so
$\langle eu(a),e'f(b)\rangle\equiv\langle ea,e'b\rangle$ modulo $M$; by skew-symmetry also
$\langle e'f(a),eu(b)\rangle\equiv\langle e'a,eb\rangle$. Expanding $\langle f_1(a),f_1(b)\rangle$ and
$\langle a,b\rangle=\langle ea+e'a,eb+e'b\rangle$ gives the claim.

Step 4. Choose integers $s,t$ with $sM+tA=1$, and let $f_0=sM+tAf_1$. Writing
$\langle f_1(a),f_1(b)\rangle=\langle a,b\rangle+Mk$,
\[
 \langle f_0(a),f_0(b)\rangle-\langle a,b\rangle
 =MA\big(st(\langle a,f_1(b)\rangle+\langle f_1(a),b\rangle-2\langle a,b\rangle)+t^2Ak\big),
\]
so $(f_0,1)$ is a similitude modulo $MA$ in the sense of the proof of
Lemma~\ref{cv-toroidal-moduli-lift}.

Step 5. Lemma~\ref{cv-toroidal-moduli-lift} gives a family of similitudes $(\omega_m,\nu_m)$ with
$\omega_{MA}\equiv f_0$ modulo $MA\Lambda$ and $\nu_{MA}\equiv1$ modulo $MA$. Modulo $M\Lambda$,
$\omega_M\equiv\omega_{MA}\equiv f_0$, and $f_0-f_1=sM(1-f_1)$. Also $f_1-f=e(u-f)$ has values in
$\mathfrak M\Lambda$, and $M\Lambda\subset\mathfrak M\Lambda$. So $\omega_M(a)-f(a)\in\mathfrak M\Lambda$.
Modulo $A\Lambda$, $\omega_A\equiv\omega_{AM}\equiv f_0$ and $f_0-1=tA(f_1-1)$, so
$\omega_A(a)-a\in A\Lambda$. By (1), $\nu_A\equiv\nu_{AM}\equiv1$ modulo $A$.
\end{proof}
% Lean: (1) SimAlg.sim_group (exists_form_eq_one, dvd_mult_tower, dvd_mult_inv) and LevelTwist.twist;
% (2) exists_sim_of_aut_mod: SimMb.exists_left_inv (Step 1), SimLift.exists_rep (Step 2), hcross and hF
% (Step 3), hform (Step 4), exists_sim_lift and the three congruences (Step 5).

\begin{lemma}[Alp\"oge]\label{cv-toroidal-moduli-tate}
In Situation~\ref{cv-function-situation}, let $\Lambda$ be the lattice of the PEL datum, let
$l=(M,\mathfrak M,N)$ be a level with $N$ prime to $M$, let $G$ be an $\mathcal O_E$-semi-abelian scheme
over a scheme $S$, let $x:\operatorname{Spec}k\to S$ be a point with $k$ algebraically closed, and let
$((\zeta'_m),(\iota_m))$ be a Tate system of $G$ at $x$ (Definition~\ref{cv-toroidal-moduli-object}(2)).
Let $\alpha:\Lambda\to G_x(k)$ be an additive, $\mathcal O_E$-linear map with kernel $\mathfrak M\Lambda$.
Then $G$ has a Tate system $((\zeta''_m),(\iota'_m))$ at $x$ with $\zeta''_N=\zeta'_N$, $\iota'_N=\iota_N$,
and $\iota'_M(a)=\alpha(a)$ for $a\in\bar{\mathfrak M}\Lambda$.
\end{lemma}
\leanref{Development/AbelianVarietyTheory/Proofs/ModuliProofs/ModuliFiniteSurjectiveOverGenericFamily.lean}{GlobalTheory.ModFam.exists_tate_of_lin}
\begin{proof}
If $p+q=0$, then $\Lambda=0$ and the given system will do. Otherwise, $M\in\mathfrak M$, so
$\alpha(Ma)=0$ and $\alpha(a)$ is $M$-torsion; hence $\alpha(a)=\iota_M(s(a))$ for an element $s(a)$ of
$\Lambda$ that is unique modulo $M\Lambda$. As $\iota_M$ is additive and $\mathcal O_E$-linear with
kernel $M\Lambda$, $a\mapsto s(a)$ is an $\mathcal O_E$-linear map $\Lambda\to\Lambda/M\Lambda$, and as
$\Lambda$ is projective it lifts to an $\mathcal O_E$-linear $f:\Lambda\to\Lambda$ with
$\alpha(a)=\iota_M(f(a))$.

Write $1=e+e'$ with $e\in\bar{\mathfrak M}$ and $e'\in\mathfrak M$. Then
$\mathfrak M\Lambda\cap\bar{\mathfrak M}\Lambda\subset M\Lambda$: an element $y$ of both is $ey+e'y$, with
$ey$ and $e'y$ in $\bar{\mathfrak M}\mathfrak M\Lambda=M\Lambda$. If $f(a)\in\mathfrak M\Lambda$, then
$ef(a)\in M\Lambda$, so $\alpha(ea)=\iota_M(ef(a))=0$, hence $ea\in\mathfrak M\Lambda$ and
$a=ea+e'a\in\mathfrak M\Lambda$. So Lemma~\ref{cv-toroidal-moduli-similitude}(2) applies with $A=N$ and
gives a family of similitudes $(\omega_m,\nu_m)$. By part (1) of that lemma,
$\zeta''_m=(\zeta'_m)^{\nu_m}$ and $\iota'_m=\iota_m\circ\omega_m$ form a Tate system of $G$ at $x$. Since
$\nu_N\equiv1$ modulo $N$, $\zeta''_N=\zeta'_N$, and since $\omega_N(a)-a\in N\Lambda$,
$\iota'_N=\iota_N$. For $a\in\bar{\mathfrak M}\Lambda$, the element $\omega_M(a)-f(a)$ lies in
$\mathfrak M\Lambda$, and in $\bar{\mathfrak M}\Lambda$ because $\omega_M$ and $f$ are
$\mathcal O_E$-linear, so it lies in $M\Lambda$, and $\iota'_M(a)=\iota_M(f(a))=\alpha(a)$.
\end{proof}
% Lean: exists_tate_of_lin: projective_lifting_property (f), exists_idem (e), the hypothesis of
% T.exists_similitude_of_aut_mod_levelIdeal, LevelTwist.twist with SimAlg.sim_group, twist_zeta and
% zeta_zpow_congr, ev_congr, modN_of_mem_both.

\begin{lemma}[Alp\"oge]\label{cv-toroidal-moduli-objects}
In Situation~\ref{cv-function-situation}, let $n$ be a prime, let $l=(M,\mathfrak M,N)$ be a level with $N$
prime to $M$, and let $\beta:R_{n,N}\to K_l$ be the inclusion.
\begin{enumerate}
\item Let $f:S'\to S$ be a morphism of schemes over $R_{n,N}$, let $(G,\lambda)$ be a full-level object
  over $S$ (Definition~\ref{cv-toroidal-family}(1)), and let $G'\to G$ be a morphism of
  $\mathcal O_E$-semi-abelian schemes over $f$. Let $\lambda'(a)$ be the section of $G'$ that it carries
  to $\lambda(a)\circ f$. Then $(G',\lambda')$ is a full-level object over $S'$.
\item Let $S$ be a scheme over $K_l$ with structure morphism $s$, and let $(G,\eta,\eta')$ be an object at
  level $l$ over $S$ (Definition~\ref{cv-toroidal-moduli-object}(4)). Then $(G,\eta')$ is a full-level
  object over $S$ regarded over $R_{n,N}$ through $\operatorname{Spec}\beta\circ s$, and $\eta$ is an
  $\mathfrak M$-level structure of $G$ over $S$.
\item Let $S$ be a scheme over $K_l$ with structure morphism $s$, let $(G,\lambda)$ be a full-level object
  over $S$ regarded over $R_{n,N}$ through $\operatorname{Spec}\beta\circ s$, and let $\chi$ be an
  $\mathfrak M$-level structure of $G$ over $S$. Then $(G,\chi,\lambda)$ is an object at level $l$ over
  $S$.
\item Let $k$ be an algebraically closed field, let $S=\operatorname{Spec}k$ be a scheme over $K_l$ with
  structure morphism $s$, and let $(G,\lambda)$ be a full-level object over $S$ regarded over $R_{n,N}$
  through $\operatorname{Spec}\beta\circ s$. Then $G$ has an $\mathfrak M$-level structure over $S$.
\end{enumerate}
\end{lemma}
\leanref{Development/AbelianVarietyTheory/Proofs/ModuliProofs/ModuliFiniteSurjectiveOverGenericFamily.lean}{GlobalTheory.ModFam.fullObj_pull}
\leanref{Development/AbelianVarietyTheory/Proofs/ModuliProofs/FullObjOfObj.lean}{GlobalTheory.Models.fullObj_of_obj}
\leanref{Development/AbelianVarietyTheory/Proofs/ModuliProofs/ModuliFiniteSurjectiveOverGenericFamily.lean}{GlobalTheory.ModFam.isLev_obj}
\leanref{Development/AbelianVarietyTheory/Proofs/ModuliProofs/ModuliFiniteSurjectiveOverGenericFamily.lean}{GlobalTheory.ModFam.objOfFull}
\leanref{Development/AbelianVarietyTheory/Proofs/ModuliProofs/ModuliFiniteSurjectiveOverGenericFamily.lean}{GlobalTheory.ModFam.exists_isLev}
\begin{proof}
Write $1=e+e'$ with $e\in\bar{\mathfrak M}$ and $e'\in\mathfrak M$; then $ea\in M\Lambda$ for
$a\in\mathfrak M\Lambda$, and $a\in\mathfrak M\Lambda$ whenever $ea\in M\Lambda$, since $a=ea+e'a$ and
$M\Lambda\subset\mathfrak M\Lambda$.

(1) The morphism identifies $G'$ with $f^*G$, with its action and line bundle, so $G'$ is proper over
$S'$ and its line bundle is relatively ample and defines a principal polarization. The map $\lambda'$ is
additive and vanishes on $N\Lambda$. At a point $x$ of $S'$ with values in an algebraically closed field
of characteristic zero, the morphism identifies $G'_x$ with $G_{f(x)}$, compatibly with the actions on
the Lie algebras and with the commutator pairings, and carries $\lambda'(a)(x)$ to $\lambda(a)(f(x))$; so
conditions (a) and (b) at $x$ follow from those for $(G,\lambda)$ at $f\circ x$.

(2) At a point $x$ of $S$ with values in an algebraically closed field $k$, with $\delta:K_l\to k$, the
conditions on $(G,\eta,\eta')$ at $x$ give condition (a) with $\delta(\beta(\theta))=\delta(\theta)$ and a
Tate system with $\zeta'_N=\delta(\beta(\zeta_N))$ and $\iota_N(a)=\eta'(a)(x)$, which is condition (b)
for $\operatorname{Spec}\beta\circ s$. For $\eta$: it is additive and vanishes on
$\mathfrak M\Lambda$. Since $e'a\in\mathfrak M\Lambda$, $\eta(a)=\eta(ea)$, and $ea\in\bar{\mathfrak
M}\Lambda$, so $\eta(a)(x)=\iota_M(ea)$. Hence $\eta(ca)(x)=\iota_M(cea)=c\,\iota_M(ea)=c\,\eta(a)(x)$, and
if $\eta(a)(x)=0$ then $ea\in M\Lambda$ and $a\in\mathfrak M\Lambda$.

(3) The conditions on $G$, $\chi$ and $\lambda$ other than those at points are given. At a point $x$ of $S$
with values in an algebraically closed field $k$, with $\delta:K_l\to k$, the field $k$ has
characteristic zero, and the conditions for $(G,\lambda)$ at $x$, over $\delta\circ\beta$, give condition
(a) and a Tate system with $\zeta'_N=\delta(\zeta_N)$ and $\iota_N(a)=\lambda(a)(x)$. The map
$a\mapsto\chi(a)(x)$ is additive and $\mathcal O_E$-linear with kernel $\mathfrak M\Lambda$, so
Lemma~\ref{cv-toroidal-moduli-tate} gives a Tate system with the same $\zeta'_N$ and $\iota_N$ whose
$\iota'_M$ agrees with $\chi(\cdot)(x)$ on $\bar{\mathfrak M}\Lambda$.

(4) Let $((\zeta'_m),(\iota_m))$ be the Tate system that condition (b) gives at the identity of $S$, and
let $\chi(a)=\iota_M(ea)$. It is additive, and it vanishes on $\mathfrak M\Lambda$ because $ea\in
M\Lambda$ there. At a point $z:\operatorname{Spec}k'\to S$ with $k'$ algebraically closed,
$\chi(ca)(z)=c\,\chi(a)(z)$ because $\iota_M$ is $\mathcal O_E$-linear. If $\chi(a)(z)=0$, then
$\chi(a)=0$, since $z$ is faithfully flat; so $ea\in M\Lambda$ and $a\in\mathfrak M\Lambda$.
\end{proof}
% Lean: (1) fullObj_pull (pol_pull_of_relAmple_local, lieMult_pull, weilIs_grpPullback_of_comm, tateMove);
% (2) fullObj_of_obj and isLev_obj (hidem with exists_idem, mem_of_idem); (3) objOfFull with
% exists_tate_of_lin; (4) exists_isLev (Flat.epi_of_flat_of_surjective for the last condition).

\begin{lemma}\label{cv-toroidal-moduli-torsion}
Let $G$ be an abelian scheme over a scheme $S$, let $M\ge1$ be an integer invertible on $S$, and let
$\Lambda$ be a free $\mathbf Z$-module of finite rank $r$.
\begin{enumerate}
\item There are a scheme $Z$, finite and \'etale over $S$, and an additive map $\epsilon:\Lambda\to G(Z)$
  with $M\epsilon=0$, such that for every scheme $T$ over $S$ and every additive map
  $\chi:\Lambda\to G(T)$ with $M\chi=0$ there is a unique morphism $h:T\to Z$ over $S$ with
  $\epsilon(a)\circ h=\chi(a)$ for all $a\in\Lambda$.
\item Let $T$ be a scheme over $S$ and let $y,y'\in G(T)$ be killed by $M$. Then there is an open and
  closed subset $U$ of $T$ such that a morphism $g:T'\to T$ over $S$ satisfies $y\circ g=y'\circ g$ if and
  only if $g(T')\subset U$.
\end{enumerate}
\end{lemma}
\leanref{Development/AbelianVarietyTheory/Proofs/ModuliProofs/ExistsFiniteEtaleLatticeTorsionHom.lean}{GlobalTheory.exists_finiteEtale_latticeTorsionHom}
\leanref{Development/AbelianVarietyTheory/Proofs/ModuliProofs/ModuliFiniteSurjectiveOverGenericFamily.lean}{GlobalTheory.ModFam.clopenCond_pt}
\begin{proof}
Multiplication by $M$ on $G$ is finite and \'etale. Indeed, $G$ is smooth and proper over $S$, and on
each fibre multiplication by $M$ is an isogeny of degree $M^{2g}$, where $g$ is the relative dimension,
which is prime to the characteristic \cite{Mum70}; so it is flat by the fibrewise criterion of flatness,
proper and quasi-finite, hence finite, and \'etale. So its kernel $G[M]$ is finite and \'etale over $S$.

(1) Let $e_1,\dots,e_r$ be a basis of $\Lambda$, let $Z$ be the $r$-fold fibre product of $G[M]$ over $S$
with projections $x_1,\dots,x_r$, and let $\epsilon(\sum_ic_ie_i)=\sum_ic_ix_i$. An additive map
$\chi:\Lambda\to G(T)$ with $M\chi=0$ is determined by the points $\chi(e_i)$ of $G[M](T)$, and these are
the components of a unique morphism $T\to Z$.

(2) The points $y,y'$ factor through $G[M]$. As $G[M]$ is finite and \'etale over $S$, its diagonal is an
open and closed immersion, and $U$ is the inverse image of the diagonal under
$(y,y'):T\to G[M]\times_SG[M]$.
\end{proof}
% Lean: exists_finiteEtale_latticeTorsionHom (GrpTangent.finiteEtale_mulN_torsion, the power powObj of
% torsion G n, existsUnique_torsion, eq_prod_zpow); clopenCond_pt (RepTors.clopenCond for the lattice Z,
% clopenCond_eq with FinEt.agreeι).

\begin{lemma}[Alp\"oge]\label{cv-toroidal-moduli-closed}
In Situation~\ref{cv-function-situation}, let $n$ be a prime, let $l=(M,\mathfrak M,N)$ be a level with $N$
prime to $M$, and let $\beta:R_{n,N}\to K_l$ be the inclusion. Let $S$ be a scheme with a separated
morphism $s:S\to\operatorname{Spec}K_l$, let $(G,\lambda)$ be a full-level object over $S$ regarded over
$R_{n,N}$ through $\operatorname{Spec}\beta\circ s$ (Definition~\ref{cv-toroidal-family}(1)), and let
$a:X_l\to S$ be a morphism with $s\circ a=s_l$ and $\nu:\mathcal A_l\to G$ a morphism of
$\mathcal O_E$-semi-abelian schemes over $a$ that carries $\eta'$ to $\lambda$. Then $a$ is finite.
\end{lemma}
\leanref{Development/AbelianVarietyTheory/Proofs/ModuliProofs/ModuliFiniteSurjectiveOverGenericFamily.lean}{GlobalTheory.ModFam.isFinite_of_full}
\begin{proof}
The scheme $X_l$ is smooth over $K_l$, hence reduced. As $G$ is proper over $S$, it is an abelian scheme,
and $M$ is invertible on $S$. Let $Z$ and $\epsilon$ be given by
Lemma~\ref{cv-toroidal-moduli-torsion}(1) for $G$, $M$ and $\Lambda$, with structure morphism
$\pi_Z:Z\to S$.

The morphism $\nu$ identifies $\mathcal A_l$ with $a^*G$, so for $b\in\Lambda$ it carries $\eta(b)$ to a
section $\chi(b)\in G(X_l)$ of $G$ over $X_l$, and $\chi$ is an $\mathfrak M$-level structure of $G$ over $X_l$,
because $\eta$ is one (Lemma~\ref{cv-toroidal-moduli-objects}(2)) and $\nu$ identifies the fibres. As
$M\Lambda\subset\mathfrak M\Lambda$, $M\chi=0$, so $\chi$ is classified by a morphism $h:X_l\to Z$ over
$S$. Let $V$ be the closure of $h(X_l)$ in $Z$ with its reduced structure, $i:V\to Z$ the closed immersion
and $j:X_l\to V$ the morphism with $i\circ j=h$.

The restriction $\epsilon_V=\epsilon\circ i$ is an $\mathfrak M$-level structure of $G$ over $V$. Each
condition holds on an open and closed subset of $Z$ by Lemma~\ref{cv-toroidal-moduli-torsion}(2). For
$b\in\mathfrak M\Lambda$, the set where $\epsilon(b)=0$ contains $h(X_l)$, hence its closure. For
$c\in\mathcal O_E$ and $b\in\Lambda$, the set where $\epsilon(cb)=c\,\epsilon(b)$ contains the image of
every geometric point of $X_l$, since $\chi$ is $\mathcal O_E$-linear there, so it contains $h(X_l)$ and
$V$. For $b\notin\mathfrak M\Lambda$, the set where $\epsilon(b)=0$ meets no point of $h(X_l)$; its
complement is closed, so it contains $V$.

Let $u=\pi_Z\circ i:V\to S$, let $G_V=u^*G$ with the pullbacks $\lambda_V$ of $\lambda$ and $\chi_V$ of
$\epsilon_V$. By Lemma~\ref{cv-toroidal-moduli-objects}(1) and~(3), $(G_V,\chi_V,\lambda_V)$ is an object
at level $l$ over $V$, regarded over $K_l$ through $s\circ u$. As $V$ is reduced, the representing
property of $X_l$ (Definition~\ref{cv-toroidal-moduli-object}(5)) gives $g:V\to X_l$ over $K_l$ with a
morphism $p$ of objects to $(\mathcal A_l,\eta,\eta')$ over $g$. Since $u\circ j=a$, $\nu$ factors through
a morphism $q:\mathcal A_l\to G_V$ over $j$, which carries $\eta$ to $\chi_V$ and $\eta'$ to $\lambda_V$.
Then $p\circ q$ is a morphism of objects from $(\mathcal A_l,\eta,\eta')$ to itself over $g\circ j$, and by
the uniqueness in the representing property, $g\circ j$ is the identity.

The morphism $s\circ u=s_l\circ g$ is separated, since $i$ is a closed immersion, $\pi_Z$ is finite and $s$
is separated; so $g$ is separated. As $g\circ j$ is a closed immersion and $g$ is separated, $j$ is a
closed immersion. Hence $a=\pi_Z\circ i\circ j$ is finite.
\end{proof}
% Lean: isFinite_of_full: exists_finiteEtale_latticeTorsionHom, IsLev.push (chi), exists_closure (V, i, j),
% IsLev.closure (with clopenCond_pt and range_subset_of_geom), exists_sabPull and IsLev.pull, fullObj_pull,
% objOfFull, Xo.rep (g, p), exists_cancel and sec_cancel (q), sabComp and the uniqueness of Xo.rep (hjg),
% IsSeparated.of_comp and IsClosedImmersion.of_comp.

\begin{lemma}[Alp\"oge]\label{cv-toroidal-moduli-finite}
In Situation~\ref{cv-function-situation}, let $n$ be an odd prime, $j\ge1$ and $N=n^j$, let
$(\mathcal M,G_{\mathcal M},\lambda)$ be a full-level family over $R_{n,N}$
(Definition~\ref{cv-toroidal-family}), and let $l=(M,\mathfrak M,N)$ be a level of full level $N$
(Definition~\ref{cv-model-toroidal-semiabelian}) with $M$ prime to $N$. Let $\beta:R_{n,N}\to K_l$ be the
inclusion. Then there are a morphism $a:X_l\to\mathcal M$ whose composite with the structure morphism of
$\mathcal M$ is $\operatorname{Spec}\beta$ composed with that of $X_l$, and a morphism
$\mathcal A_l\to G_{\mathcal M}$ of $\mathcal O_E$-semi-abelian schemes over $a$ that carries $\eta'$ to
$\lambda$, such that the induced morphism $X_l\to\mathcal M\otimes_{R_{n,N}}K_l$ is finite and
surjective.
\end{lemma}
\leanref{Development/AbelianVarietyTheory/Proofs/ModuliProofs/ModuliFiniteSurjectiveOverGenericFamilyStatement.lean}{GlobalTheory.T.moduli_finite_surjective_over_generic_family}
\leanref{Development/AbelianVarietyTheory/Proofs/ModuliProofs/ModuliFiniteSurjectiveOverGenericFamily.lean}{GlobalTheory.moduli_finite_surjective_over_family}
\begin{proof}
Let $s_l:X_l\to\operatorname{Spec}K_l$ be the structure morphism. The scheme $X_l$ is smooth over $K_l$,
so it is reduced and locally noetherian. The ring $R_{n,N}$ is a Dedekind domain with fraction field
$K_l$, so $\operatorname{Spec}\beta$ is flat, and $X_l$ is flat over $R_{n,N}$ through
$\operatorname{Spec}\beta\circ s_l$. By Lemma~\ref{cv-toroidal-moduli-objects}(2),
$(\mathcal A_l,\eta')$ is a full-level object over $X_l$ regarded over $R_{n,N}$ through
$\operatorname{Spec}\beta\circ s_l$. Property (a) of Definition~\ref{cv-toroidal-family}(2) gives
$a:X_l\to\mathcal M$ with $s_{\mathcal M}\circ a=\operatorname{Spec}\beta\circ s_l$, where $s_{\mathcal M}$
is the structure morphism of $\mathcal M$, and a morphism $\nu:\mathcal A_l\to G_{\mathcal M}$ over $a$ that
carries $\eta'$ to $\lambda$.

Let $Y=\mathcal M\otimes_{R_{n,N}}K_l$, with projections $p_1:Y\to\mathcal M$ and
$p_2:Y\to\operatorname{Spec}K_l$, and let $b:X_l\to Y$ be the morphism with $p_1\circ b=a$ and
$p_2\circ b=s_l$. Let $G_Y=p_1^*G_{\mathcal M}$, with the pullback $\lambda_Y$ of $\lambda$. By
Lemma~\ref{cv-toroidal-moduli-objects}(1), $(G_Y,\lambda_Y)$ is a full-level object over $Y$, regarded
over $R_{n,N}$ through $s_{\mathcal M}\circ p_1=\operatorname{Spec}\beta\circ p_2$.

Finiteness. The morphism $\nu$ factors through a morphism $\mathcal A_l\to G_Y$ over $b$ that carries
$\eta'$ to $\lambda_Y$. The morphism $p_2$ is separated, because $\mathcal M$ is quasi-projective over
$R_{n,N}$. So Lemma~\ref{cv-toroidal-moduli-closed} applies with $S=Y$, $s=p_2$ and $a=b$, and $b$ is
finite.

Surjectivity. Let $y\in Y$, and let $z:\operatorname{Spec}k\to Y$ be a point with image $y$ and $k$
algebraically closed; $k$ has characteristic zero, as $Y$ lies over $K_l$. Let $(G_k,\lambda_k)$ be the
pullback of $(G_Y,\lambda_Y)$ along $z$, a full-level object over $\operatorname{Spec}k$ by
Lemma~\ref{cv-toroidal-moduli-objects}(1). By parts (4) and (3) of that lemma, $G_k$ has an
$\mathfrak M$-level structure $\chi$, and $(G_k,\chi,\lambda_k)$ is an object at level $l$ over
$\operatorname{Spec}k$, regarded over $K_l$ through $p_2\circ z$. By the representing property of $X_l$
(Definition~\ref{cv-toroidal-moduli-object}(5)) there is $f:\operatorname{Spec}k\to X_l$ with
$s_l\circ f=p_2\circ z$ and a morphism of objects $G_k\to\mathcal A_l$ over $f$. Its composite with $\nu$
is a morphism $G_k\to G_{\mathcal M}$ over $a\circ f$ that carries $\lambda_k$ to $\lambda$, and so is the
morphism $G_k\to G_{\mathcal M}$ over $p_1\circ z$ given by the pullbacks. The scheme
$\operatorname{Spec}k$ is reduced, noetherian and flat over $R_{n,N}$, so the uniqueness in property (a)
of Definition~\ref{cv-toroidal-family}(2) gives $a\circ f=p_1\circ z$. Hence $b\circ f=z$, and $y$ lies in
the image of $b$.
\end{proof}
% Lean: moduli_finite_surjective_over_family: isReduced_isNormalScheme_isLocallyNoetherian_of_smooth_over_field,
% BaseRing.isDedFrac_RZ (Spec beta flat), fullObj_of_obj and F.rep (a and nu); finiteness by exists_sabPull,
% fullObj_pull, exists_cancel and isFinite_of_full on the fibre product; surjectivity by
% FinEt.exists_algClosed_point, exists_isLev, objOfFull, Xo.rep, sabComp and the uniqueness of F.rep.

\begin{lemma}[Alp\"oge]\label{cv-toroidal-level-change}
In Situation~\ref{cv-function-situation}, suppose $p+q\ge1$, let $n$ be an odd prime, $j\ge1$ and
$N=n^j$, let $(\mathcal M,G_{\mathcal M},\lambda)$ be a full-level family over $R_{n,N}$ with structure
morphism $s_{\mathcal M}$ (Definition~\ref{cv-toroidal-family}), let $l$ be a level of full level $N$
(Definition~\ref{cv-model-toroidal-semiabelian}), with structure morphism
$s_l:X_l\to\operatorname{Spec}K_l$, and let $\beta:R_{n,N}\to K_l$ be the inclusion. Let
$a:X_l\to\mathcal M$ be a morphism with $s_{\mathcal M}\circ a=\operatorname{Spec}\beta\circ s_l$ and
$\nu_0:\mathcal A_l\to G_{\mathcal M}$ a morphism of $\mathcal O_E$-semi-abelian schemes over $a$
carrying $\eta'$ to $\lambda$, such that the induced morphism $X_l\to\mathcal M\otimes_{R_{n,N}}K_l$ is
surjective. Let $\sigma$ be an automorphism of $K_l$ with $\sigma(\sqrt{-d})=\sqrt{-d}$, $r$ an
automorphism of $X_l$ with $s_l\circ r=\operatorname{Spec}\sigma\circ s_l$, $g$ an
$\mathcal O_E$-linear endomorphism of $\Lambda$, and $\nu:\mathcal A_l\to\mathcal A_l$ a morphism of
$\mathcal O_E$-semi-abelian schemes over $r$ with $\nu\circ\eta'(g(y))=\eta'(y)\circ r$ for all
$y\in\Lambda$. Then there are an automorphism $\rho$ of $\mathcal M$ with a morphism
$G_{\mathcal M}\to G_{\mathcal M}$ of $\mathcal O_E$-semi-abelian schemes over $\rho$, and an
automorphism $\sigma_0$ of $R_{n,N}$ with $\sigma_0(\theta)=\theta$, such that $\rho\circ a=a\circ r$
and $s_{\mathcal M}\circ\rho=\operatorname{Spec}\sigma_0\circ s_{\mathcal M}$.
\end{lemma}
\leanref{Development/AbelianVarietyTheory/Proofs/ModuliProofs/FullLevelChangeActsOnUnitaryFamilyStatement.lean}{GlobalTheory.T.fullLevelChange_acts_on_unitaryFamily}
\leanref{Development/AbelianVarietyTheory/Proofs/ModuliProofs/FullLevelChangeActsOnUnitaryFamily.lean}{GlobalTheory.famCore_exists_aut}
\begin{proof}
If $X_l$ is empty, take identities. Otherwise, let $\pi=\operatorname{Spec}\beta\circ s_l$. Since
$(\mathcal A_l,\eta')$ satisfies at every geometric point of $X_l$ the conditions of the moduli problem
at level $l$, it is a full-level object over $X_l$ regarded over $R_{n,N}$ through $\pi$. The scheme
$X_l$ is smooth over $K_l$, hence reduced and locally noetherian, and it is flat over $R_{n,N}$, since
$K_l$ is a localization of $R_{n,N}$.

The ring $R_{n,N}$ is generated by $(1+\theta)/2$, $1/n$ and $\zeta_N$. The automorphism $\sigma$ fixes
the first two and maps $\zeta_N$ to a power of itself, and so does $\sigma^{-1}$. So $\sigma$ restricts
to an automorphism $\sigma_0$ of $R_{n,N}$ with $\beta\circ\sigma_0=\sigma\circ\beta$ and
$\sigma_0(\theta)=\theta$. Then
$s_{\mathcal M}\circ a\circ r=\operatorname{Spec}\sigma_0\circ s_{\mathcal M}\circ a$.

Every point $x:\operatorname{Spec}k\to\mathcal M$ with $k$ algebraically closed of characteristic zero
is $a\circ y$ for a point $y$ of $X_l$. Indeed, the homomorphism $R_{n,N}\to k$ given by $x$ is
injective, since a nonzero element of its kernel would be algebraic over $\mathbf Z$ and the kernel
would then meet $\mathbf Z$ in a nonzero ideal, which is impossible in characteristic zero. So it
extends to the fraction field $K_l$, and $x$ with this extension is a point of
$\mathcal M\otimes_{R_{n,N}}K_l$, which lifts to $X_l$ because the morphism from $X_l$ is surjective and
locally of finite type.

The pair $(\mathcal A_l,\eta'\circ g)$ is a full-level object over $X_l$, regarded over $R_{n,N}$
through $s_{\mathcal M}\circ a\circ r$. Additivity and vanishing on $N\Lambda$ follow from those of
$\eta'$, as $g$ is linear. At a geometric point $y$ of $X_l$, conditions (a) and (b) of
Definition~\ref{cv-toroidal-family}(1) for $(\mathcal A_l,\eta'\circ g)$ at $y$ are transported by $\nu$
from those for $(\mathcal A_l,\eta')$ at $r\circ y$. In the same way, the pair
$(G_{\mathcal M},\lambda\circ g)$ is a full-level object over $\mathcal M$ regarded over $R_{n,N}$
through $\operatorname{Spec}\sigma_0\circ s_{\mathcal M}$: at a geometric point $a\circ y$ of
$\mathcal M$, the conditions are transported by $\nu_0$ and $\nu$ from those for $(G_{\mathcal M},\lambda)$
at $a\circ r\circ y$.

Fix a geometric point $y$ of $X_l$. The homomorphisms $\iota_N$ of condition (b) for $\eta'$ and for
$\eta'\circ g$ at $y$ induce bijections of $\Lambda/N\Lambda$ with the $N$-torsion of $\mathcal A_{l,y}$,
with $\iota_N(g(z))=\iota'_N(z)$ for $z\in\Lambda$. So there is a map $g':\Lambda\to\Lambda$, additive
modulo $N\Lambda$, with $g\circ g'$ and $g'\circ g$ congruent to the identity modulo $N\Lambda$. As $\nu$
identifies $\mathcal A_l$ with $r^*\mathcal A_l$ and $r$ is an automorphism, $\nu$ has an inverse over
$r^{-1}$, and $\nu^{-1}\circ\eta'(g'(z))=\eta'(z)\circ r^{-1}$ for $z\in\Lambda$ because
$\eta'\circ g\circ g'=\eta'$, as $\eta'$ vanishes on $N\Lambda$. As above,
$(G_{\mathcal M},\lambda\circ g')$ is a full-level object over $\mathcal M$ regarded over $R_{n,N}$ through
$\operatorname{Spec}\sigma_0^{-1}\circ s_{\mathcal M}$.

The scheme $\mathcal M$ is reduced, locally noetherian and flat over $R_{n,N}$ through
$\operatorname{Spec}\sigma_0^{\pm1}\circ s_{\mathcal M}$. By property (a) of
Definition~\ref{cv-toroidal-family}(2), there are $\rho:\mathcal M\to\mathcal M$ with
$s_{\mathcal M}\circ\rho=\operatorname{Spec}\sigma_0\circ s_{\mathcal M}$ and a morphism
$\mu:G_{\mathcal M}\to G_{\mathcal M}$ over $\rho$ carrying $\lambda\circ g$ to $\lambda$, and
$\rho'$ with $s_{\mathcal M}\circ\rho'=\operatorname{Spec}\sigma_0^{-1}\circ s_{\mathcal M}$ and $\mu'$
over $\rho'$ carrying $\lambda\circ g'$ to $\lambda$. Then $\mu'\circ\mu$ lies over $\rho'\circ\rho$, which
lies over the identity of $\operatorname{Spec}R_{n,N}$, and carries $\lambda$ to $\lambda$, because
$\lambda\circ g\circ g'=\lambda$ as $\lambda$ vanishes on $N\Lambda$. By the uniqueness in property (a),
applied to $(G_{\mathcal M},\lambda)$, $\rho'\circ\rho$ is the identity, and so is $\rho\circ\rho'$. So
$\rho$ is an automorphism.

Finally, $a\circ r$ and $\rho\circ a$ both lie over $s_{\mathcal M}\circ a\circ r$, and
$\nu_0\circ\nu$ over $a\circ r$ and $\mu\circ\nu_0$ over $\rho\circ a$ both carry $\eta'\circ g$ to
$\lambda$. The scheme $X_l$ is flat over $R_{n,N}$ through $s_{\mathcal M}\circ a\circ r$, which is
$\pi\circ r$. By the uniqueness in property (a), applied to the full-level object
$(\mathcal A_l,\eta'\circ g)$, $\rho\circ a=a\circ r$.
\end{proof}
% Lean: famCore_exists_aut: exists_restrict (sigma_0), exists_extend and ObjDesc.exists_geomPt_lift
% (geometric points lift), fullObj_of_obj and geomAt_iff (the full-level objects), exists_inv_mod (g'),
% sab_cast and sabCancel (the inverse of nu), exists_endo (rho, rho'), eq_id, and FamCore.rep for the
% last step. The hypotheses Laws, Laws2 and 1 <= p + q are not used.

\subsection{Construction of the model}\label{subsection-cv-model-construction}

\begin{lemma}\label{cv-toroidal-extend-lift}
Let $T$ be a normal, locally noetherian scheme and let $j:W\to T$ be a flat, dominant, quasi-compact
morphism that is a homeomorphism onto its image and surjective on stalks. Let $q:Y\to X$ be an integral
morphism, and let $f:T\to X$ and $g:W\to Y$ be morphisms with $q\circ g=f\circ j$. Then there is a unique
morphism $h:T\to Y$ with $h\circ j=g$ and $q\circ h=f$.
\end{lemma}
\leanref{Development/AbelianVarietyTheory/Proofs/ModuliProofs/ExistsUniqueLiftOfDensePart.lean}{GlobalTheory.existsUnique_lift_of_densePart}
\leanref{Development/AbelianVarietyTheory/Proofs/ModuliProofs/IsIsoStalkMapOfFlatOfSurjectiveOnStalks.lean}{GlobalTheory.densePart_basic}
\begin{proof}
The local homomorphisms $\mathcal O_{T,j(w)}\to\mathcal O_{W,w}$ are surjective and flat, hence
isomorphisms. As $j$ is flat, $j(W)$ is stable under generization, and as $j$ is dominant and
quasi-compact, it contains the generic points of $T$. The scheme $T$ is normal and locally noetherian, so
its irreducible components $T_i$ are open and integral, and $W_i=j^{-1}(T_i)$ is integral with the same
function field. The base change $Y\times_XT_i\to T_i$ is integral, and $g$ gives a point of it over the
generic point of $T_i$, with the same residue field. Its closure, with the reduced structure, is integral
and birational over the normal scheme $T_i$, hence isomorphic to $T_i$; this gives $h$ on $T_i$ with
$q\circ h=f$. The morphisms $h\circ j$ and $g$ from $W_i$ to $Y$ lie over the same morphism to $X$ and
agree at the generic point of the integral scheme $W_i$; as $q$ is separated, they are equal. The
morphisms on the $T_i$ glue. Uniqueness follows in the same way, since $T$ is reduced, $j$ is dominant
and $q$ is separated.
\end{proof}
% Lean: existsUnique_lift_of_densePart: isIntegral_of_densePart and isIso_stalkMap_of_flat_of_surjectiveOnStalks
% on each component (IsNormalScheme.component), existsUnique_lift_of_birational (through the normalization,
% isIso_fromNormalization_of_birational), Cofan.IsColimit.desc to glue, lift_unique_of_isDominant.

\begin{lemma}[Alp\"oge]\label{cv-toroidal-extend-trivial}
Let $X$ be a normal, locally noetherian scheme and let $j:W\to X$ be a flat, dominant, quasi-compact
morphism that is a homeomorphism onto its image and surjective on stalks. Let $p:P\to X$ be smooth with
geometrically integral fibres and with a section $e$, and let $\mathcal N$ be a line bundle on $P$ such
that $e^*\mathcal N$ and the pullback of $\mathcal N$ to $P_W=P\times_XW$ are trivial. Then $\mathcal N$ is
trivial.
\end{lemma}
\leanref{Development/AbelianVarietyTheory/Proofs/ModuliProofs/ExistsLineBundleGeneratorExtendingDensePartStatement.lean}{GlobalTheory.T.exists_lineBundleGenerator_extending_densePart}
\leanref{Development/AbelianVarietyTheory/Proofs/ModuliProofs/ExistsLineBundleGeneratorExtendingDensePart.lean}{GlobalTheory.GenericExtend.exists_isGen_of_densePart}
\begin{proof}
Let $\rho$ generate $e^*\mathcal N$ and $t$ generate the pullback of $\mathcal N$ to $P_W$. The section
$e_W$ of $P_W$ over $W$ induced by $e$ carries $t$ to a generator of the pullback of $e^*\mathcal N$ to $W$,
which is $v\rho$ for a unit $v$ on $W$. Replacing $t$ by $t$ divided by the pullback of $v$, we may assume
that $t$ and $\rho$ agree along $e_W$.

As in the proof of Lemma~\ref{cv-toroidal-extend-lift}, the projection $P_W\to P$ induces isomorphisms
of local rings, and its image is stable under generization and contains the generic points of $P$. The
scheme $P$ is normal, being smooth over $X$. Let $\sigma$ generate $\mathcal N$ on an open subset $Q$ of
$P$. On the inverse image of $Q$ in $P_W$, $t=f\sigma$ for a unit $f$; as the local rings of $P_W$ are those
of $P$, $f$ is a rational function on $Q$ that is regular and invertible at the points of the image of
$P_W$. Since $P$ is normal, $f$ and $f^{-1}$ are regular on $Q$ once they are
regular at every point $y$ of codimension one of $Q$ outside that image. Such a point lies over a point
$x=p(y)$ that is not in $j(W)$ and has codimension at most one, as $p$ is flat; so $x$ has codimension one,
and $V=\mathcal O_{X,x}$ is a discrete valuation ring, with fraction field $K$. The fibre $P_x$ is integral
with generic point $y$, and $\operatorname{Spec}K\to X$ factors through $W$. Choose $\sigma$ near $e(x)$, so
that $Q$ contains $e(x)$ and $y$. On $Q\times_X\operatorname{Spec}V$, which is normal with integral special fibre,
$f$ is invertible over $K$, so its divisor is a multiple of the special fibre, and $f=\varpi^mf'$ with
$\varpi$ a uniformizer of $V$ and $f'$ a unit near $e(x)$ and $y$. Pulled back along $e$, $f$ becomes $\rho$ divided by $e^*\sigma$, a unit of
$V$; so $m=0$, and $f$ is a unit at $y$. The generators $t$ thus extend to generators of $\mathcal N$ on
each $Q$, and they agree on overlaps, so $\mathcal N$ is trivial.
\end{proof}
% Lean: exists_isGen_of_densePart: exists_compat, exists_normalised (t agrees with r along e_W),
% coheight_apply_le_one and exists_specializes_mem_range (y over a point of coheight one),
% extAt_of_coheight_eq_one with exists_unit_dvr = DvrModel.existsUnique_isUnit_appLE_eq (the discrete
% valuation ring step), exists_app_eq (normality) and the glued generator sigma.

\begin{lemma}\label{cv-toroidal-extend-homomorphisms-torsion}
Let $k$ be an algebraically closed field, let $G$ be a semi-abelian scheme over $k$, and let $n\ge1$ be an
integer. Then only finitely many points of $G$ with values in $k$ have order dividing $n$.
\end{lemma}
\leanref{Development/AbelianVarietyTheory/Proofs/ModuliProofs/FiniteTorsionSemiAbelian.lean}{GlobalTheory.finite_torsion_semiAbelian}
\leanref{Development/AbelianVarietyTheory/Proofs/ModuliProofs/FiniteTorsionSemiAbelianOfAbelian.lean}{GlobalTheory.TorsionFinite.finite_torsion_semiAbelian_of_abelian}
\begin{proof}
As $k$ is algebraically closed, $G$ is an extension of an abelian variety $B$ by a split torus $T$: there are
homomorphisms $\iota:T\to G$ and $\pi:G\to B$ such that $\iota$ is a kernel of $\pi$, and $T$ is the torus with
character group $\Xi$, a free abelian group of finite rank. On points with values in $k$, $\iota$ is injective
and its image is the kernel of $\pi$. A point of $T$ with values in $k$ is a homomorphism $\Xi\to k^\times$. If
its order divides $n$, it takes values in the finite group of $n$-th roots of unity of $k$, and it is
determined by its values on a finite set of generators of $\Xi$; so $T$ has finitely many such points.
Multiplication by $n$ on $B$ is finite \cite[Tag 0BFG]{Sta}, so its fibre over the origin is finite, and $B$
has finitely many points with values in $k$ of order dividing $n$. If $g$ and $g'$ are points of $G$ of order
dividing $n$ with $\pi(g)=\pi(g')$, then $g'g^{-1}=\iota(t)$ for a point $t$ of $T$, and $t^n=1$ because $G$ is
commutative and $\iota$ is injective. So the points of $G$ of order dividing $n$ have finitely many images in
$B$, and finitely many of them lie over each image.
\end{proof}
% Lean: finite_torsion_semiAbelian = TorsionFinite.finite_torsion_semiAbelian_of_abelian applied to
% AbTorsionFinite.finite_torsion. The extension is the TorusExtension of IsSemiAbelian.fibres; the count is
% semiAbelian_torsion_finite_and_closedPoint_finiteOrder(1): finite_pow_eq_one_of_extension, finite_chars_pow_eq_one
% and splitTorusPointsEquiv (about 410 lines). AbTorsionFinite.finite_torsion is finite_fibre_points with
% T.isFinite_mulByNat_abelianVariety_algClosed (MulFiniteAny.isFinite_mulA): the Lean's own proof that [n] is
% finite on an abelian variety over an algebraically closed field, about 15,700 lines, cited as Stacks Tag 0BFG.

\begin{lemma}\label{cv-toroidal-extend-homomorphisms-points}
Let $i:W\to Z$ be a flat morphism of schemes that is a homeomorphism onto its image and surjective on stalks.
\begin{enumerate}
\item Every morphism $x:T\to Z$ whose image lies in $i(W)$ is $i\circ y$ for a morphism $y:T\to W$.
\item Suppose moreover that $Z$ is locally noetherian and that $i$ is dominant and quasi-compact, and let $z$
  be a point of $Z$ outside $i(W)$. Then there are a discrete valuation ring $V$ with fraction field $K$, a
  morphism $a:\operatorname{Spec}V\to Z$ that maps the closed point to $z$, and a morphism
  $c:\operatorname{Spec}K\to W$ with $i\circ c=a\circ\operatorname{Spec}(V\subset K)$.
\end{enumerate}
\end{lemma}
\leanref{Development/AbelianVarietyTheory/Proofs/ModuliProofs/IsIsoOfDensePartOfIsIntegral.lean}{GlobalTheory.HomExt.exists_lift_of_range_subset}
\leanref{Development/AbelianVarietyTheory/Proofs/ModuliProofs/ExistsUniqueHomExtensionSemiAbelianNormal.lean}{GlobalTheory.HomExt.mem_range_or_exists_dvr}
\begin{proof}
(1) Let $P=T\times_ZW$, with projections $\pi:P\to T$ and $\pi':P\to W$. The morphism $\pi$ is a base change
of $i$, so it is flat, a homeomorphism onto its image and surjective on stalks, as these properties are stable
under base change. Its image is the inverse image of $i(W)$ under $x$, which is all of $T$. As in the proof of
Lemma~\ref{cv-toroidal-extend-lift}, its local homomorphisms are surjective and flat, hence isomorphisms. So
$\pi$ is a homeomorphism that induces isomorphisms of local rings, hence an isomorphism, and
$y=\pi'\circ\pi^{-1}$ satisfies $i\circ y=x\circ\pi\circ\pi^{-1}=x$.

(2) The point $z$ lies in an irreducible component of $Z$; let $\xi$ be its generic point, which specializes
to $z$. As in the proof of Lemma~\ref{cv-toroidal-extend-lift}, $i(W)$ contains the generic points of $Z$, so
$\xi\in i(W)$ and $\xi\ne z$. As $Z$ is locally noetherian, there are a discrete valuation ring $V$ and a
morphism $a:\operatorname{Spec}V\to Z$ that maps the generic point to $\xi$ and the closed point to $z$
\cite[Tag 054F]{Sta}. The image of $a\circ\operatorname{Spec}(V\subset K)$ is $\{\xi\}$, which lies in $i(W)$,
so (1) gives $c$.
\end{proof}
% Lean: exists_lift_of_range_subset: isIso_iff_isIso_stalkMap, isIso_stalkMap_of_flat_of_surjectiveOnStalks and
% TopCat.isIso_of_bijective_of_isOpenMap. mem_range_or_exists_dvr: exists_mem_genericPoints_specializes,
% genericPoints_subset_range_of_isDominant, exists_complete_discreteValuationRing_of_specializes (whose V is also
% complete, which is not used) and exists_lift_of_range_subset.

\begin{lemma}\label{cv-toroidal-extend-homomorphisms-closure}
Let $S$ be a scheme, let $f_1:Z_1\to S$ and $f_2:Z_2\to S$ be morphisms with $f_2$ quasi-separated, and let
$u:U\to Z_1$ and $\varphi:U\to Z_2$ be morphisms with $f_1\circ u=f_2\circ\varphi$ and $u$ quasi-compact. Call a
morphism $d:T'\to T$ \emph{scheme-theoretically dominant} if $\mathcal O_T\to d_*\mathcal O_{T'}$ is injective.
Let $\Gamma\subset Z_1\times_SZ_2$ be the scheme-theoretic image of $(u,\varphi):U\to Z_1\times_SZ_2$, the
\emph{closure of the graph} of $\varphi$, with the projections $p:\Gamma\to Z_1$ and $q:\Gamma\to Z_2$ and the
morphism $i:U\to\Gamma$ through which $(u,\varphi)$ factors, so that $p\circ i=u$ and $q\circ i=\varphi$.
\begin{enumerate}
\item Two morphisms $a,b:T\to\Gamma$ with $p\circ a=p\circ b$ and $q\circ a=q\circ b$ are equal.
\item The morphism $i$ is quasi-compact and scheme-theoretically dominant, and $\Gamma$ is reduced if $U$ is.
\item If $f_2$ is separated, locally of finite type or quasi-compact, then so is $p$.
\item Let $t_1:T\to Z_1$ and $t_2:T\to Z_2$ be morphisms with $f_1\circ t_1=f_2\circ t_2$, let $d:T'\to T$ be
  scheme-theoretically dominant, and let $w:T'\to U$ satisfy $t_1\circ d=u\circ w$ and $t_2\circ d=\varphi\circ w$.
  Then there is a unique morphism $l:T\to\Gamma$ with $p\circ l=t_1$ and $q\circ l=t_2$.
\item If $u$ is flat and a monomorphism and $f_2$ is separated, then $i$ and the identity of $U$ identify $U$
  with the fibre product $\Gamma\times_{Z_1}U$ of $p$ and $u$.
\item If $p$ is an isomorphism, then $\psi=q\circ p^{-1}:Z_1\to Z_2$ satisfies $f_2\circ\psi=f_1$ and
  $\psi\circ u=\varphi$.
\end{enumerate}
\end{lemma}
\leanref{Development/AbelianVarietyTheory/Definitions/ShimuraVarieties/FlatImageOfOverFractionFieldOfValuationRing.lean}{GlobalTheory.graphClosure}
\leanref{Development/AbelianVarietyTheory/Definitions/ShimuraVarieties/FlatImageOfOverFractionFieldOfValuationRing.lean}{GlobalTheory.graphClosure.hom_ext}
\leanref{Development/AbelianVarietyTheory/Definitions/ShimuraVarieties/FlatImageOfOverFractionFieldOfValuationRing.lean}{GlobalTheory.graphClosure.existsUnique_lift}
\leanref{Development/AbelianVarietyTheory/Definitions/ShimuraVarieties/FlatImageOfOverFractionFieldOfValuationRing.lean}{GlobalTheory.graphClosure.isPullback_incl}
\leanref{Development/AbelianVarietyTheory/Definitions/ShimuraVarieties/FlatImageOfOverFractionFieldOfValuationRing.lean}{GlobalTheory.graphClosure.exists_extension_of_isIso_fst}
\begin{proof}
(1) The inclusion of $\Gamma$ in $Z_1\times_SZ_2$ is a closed immersion, hence a monomorphism, and it maps $a$
and $b$ to the same morphism $(p\circ a,q\circ a)=(p\circ b,q\circ b)$.

(2) The projection $Z_1\times_SZ_2\to Z_1$ is quasi-separated, as a base change of $f_2$, and its composite with
$(u,\varphi)$ is the quasi-compact morphism $u$; so $(u,\varphi)$ is quasi-compact. Hence the ideal of $\Gamma$ is
the kernel of $\mathcal O\to(u,\varphi)_*\mathcal O_U$ on $Z_1\times_SZ_2$, so $\mathcal O_\Gamma\to i_*\mathcal O_U$ is
injective. The morphism $i$ is quasi-compact because its composite with the closed immersion of $\Gamma$ is. If $U$
is reduced, the rings of sections of $i_*\mathcal O_U$ are reduced, and so are those of $\mathcal O_\Gamma$.

(3) The morphism $p$ is the composite of the closed immersion of $\Gamma$ in $Z_1\times_SZ_2$, which is separated,
of finite type and quasi-compact, with the projection to $Z_1$, a base change of $f_2$.

(4) Let $t=(t_1,t_2):T\to Z_1\times_SZ_2$, so that $t\circ d=(u,\varphi)\circ w$. A section of the ideal of $\Gamma$
over an open subset of $Z_1\times_SZ_2$ pulls back to $0$ along $(u,\varphi)$, hence along $t\circ d$. As $d$ is
scheme-theoretically dominant, it pulls back to $0$ along $t$. So $t$ factors through the closed subscheme
$\Gamma$, which gives $l$, and $l$ is unique by (1).

(5) Let $u'=u\times_SZ_2:U\times_SZ_2\to Z_1\times_SZ_2$, a base change of $u$, hence flat and a monomorphism. The
morphism $c=(\mathrm{id},\varphi):U\to U\times_SZ_2$ is a section of the projection to $U$, which is separated as a
base change of $f_2$; so $c$ is a closed immersion, and $u'\circ c=(u,\varphi)$. As $u'$ is a monomorphism, the
base change of $(u,\varphi)$ along $u'$ is $c$, and the scheme-theoretic image of the closed immersion $c$ is $c$
itself. Taking scheme-theoretic images of quasi-compact morphisms commutes with flat base change
\cite[Tag 081I]{Sta}, so the inverse image of $\Gamma$ under $u'$ is $U$, embedded by $c$, with $i$ as its map to
$\Gamma$. As $U\times_SZ_2$ is the fibre product of $u$ and the projection $Z_1\times_SZ_2\to Z_1$, this inverse
image is $\Gamma\times_{Z_1}U$, which is therefore identified with $U$ through $i$ and the identity.

(6) As $f_1\circ p=f_2\circ q$, $f_2\circ\psi=f_1\circ p\circ p^{-1}=f_1$, and $\psi\circ u=q\circ p^{-1}\circ p\circ
i=q\circ i=\varphi$.
\end{proof}
% Lean: graphClosure = (graphMap ...).image, the scheme-theoretic image, with incl = toImage, fst and snd; (1) is
% hom_ext; (2) is schemeDominant_incl (schemeDominant_toImage), quasiCompact_incl and isReduced
% (IsSchemeTheoreticallyDominant.isReduced); (3) is isSeparated_fst, locallyOfFiniteType_fst and quasiCompact_fst;
% (4) is existsUnique_lift through lift = closureLift and ker_comp_of_schemeDominant; (5) is isPullback_incl through
% isPullback_toImage_of_isClosedImmersion and ker_eq_comap_of_flat (the flat base change of the kernel); (6) is
% exists_extension_of_isIso_fst.

\begin{lemma}[Alp\"oge]\label{cv-toroidal-extend-homomorphisms-isomorphism}
Let $Y$ be a normal, locally noetherian scheme, let $\Gamma$ be a reduced scheme, and let $p:\Gamma\to Y$ be a
separated, quasi-compact and surjective morphism that is locally of finite type and locally quasi-finite. Let
$i:W\to\Gamma$ be a flat, dominant morphism that is a homeomorphism onto its image and surjective on stalks, such
that $p\circ i$ is flat, dominant, quasi-compact, a homeomorphism onto its image and surjective on stalks. Then
$p$ is an isomorphism.
\end{lemma}
\leanref{Development/AbelianVarietyTheory/Proofs/ModuliProofs/IsIsoOfDensePartOfIsIntegral.lean}{GlobalTheory.HomExt.isIso_of_densePart}
\leanref{Development/AbelianVarietyTheory/Proofs/ModuliProofs/IsIsoOfDensePartOfIsIntegral.lean}{GlobalTheory.HomExt.isIso_of_densePart_of_isIntegral}
\leanref{Development/AbelianVarietyTheory/Proofs/FoundationProofs/BijectiveAlgebraMapIntegralClosureOfBirational.lean}{GlobalTheory.isOpenImmersion_of_quasiFinite_of_isIso_stalkMap}
\begin{proof}
As in the proof of Lemma~\ref{cv-toroidal-extend-lift}, the irreducible components of $Y$ are open and integral.
Being an isomorphism is local on $Y$, and the hypotheses are preserved when $Y$ is replaced by one of its
components $Y_0$, $\Gamma$ by $p^{-1}(Y_0)$ and $W$ by $(p\circ i)^{-1}(Y_0)$. So we may assume that $Y$ is
integral, with generic point $\eta$.

As in the proof of Lemma~\ref{cv-toroidal-extend-lift}, the image of $p\circ i$ contains $\eta$; let $w_0\in W$
with $p(i(w_0))=\eta$. For $w\in W$, $\eta$ specializes to $p(i(w))$, and as $p\circ i$ is a homeomorphism onto
its image, $w_0$ specializes to $w$; so $i(w)$ lies in the closure of $\{i(w_0)\}$. As $i$ is dominant, this
closure is $\Gamma$. So $\Gamma$ is irreducible with generic point $\theta=i(w_0)$, and it is integral, being
reduced. The local homomorphisms of $p\circ i$ and of $i$ at $w_0$ are flat and surjective, hence isomorphisms, as
in the proof of Lemma~\ref{cv-toroidal-extend-lift}; therefore so is that of $p$ at $\theta$, and $p$ identifies
the function field $K$ of $Y$ with that of $\Gamma$.

The morphism $p$ is quasi-finite, being quasi-compact, locally of finite type and locally quasi-finite. By
Zariski's main theorem \cite[Tag 02LR]{Sta}, $p$ is the composite of an open immersion $\Gamma\to Y'$ with the
morphism $\nu:Y'\to Y$, where $Y'$ is the normalization of $Y$ in $\Gamma$. Over an affine open subset
$\operatorname{Spec}A$ of $Y$, $Y'$ is the spectrum of the integral closure of $A$ in the ring of functions on the
inverse image of $\operatorname{Spec}A$ in $\Gamma$, a nonempty open subset of the integral scheme $\Gamma$; that
ring is a subring of $K$. As $Y$ is normal and integral, $A$ is integrally closed in $K$, so this integral closure
is $A$, and $\nu$ is an isomorphism. So $p$ is an open immersion, and as it is surjective, it is an isomorphism.
\end{proof}
% Lean: isIso_of_densePart restricts to the open components (IsNormalScheme.isOpen_of_mem_irreducibleComponents,
% IsZariskiLocalAtTarget) and applies isIso_of_densePart_of_isIntegral: genericPoints_subset_range_of_isDominant,
% the generic point hgen, isIso_stalkMap_of_flat_of_surjectiveOnStalks for i o p and i, then
% isOpenImmersion_of_quasiFinite_of_isIso_stalkMap: Mathlib's Zariski main theorem
% (exists_isIso_morphismRestrict_toNormalization) and isIso_fromNormalization_of_birational
% (bijective_algebraMap_integralClosure_of_isIntegrallyClosed).

\begin{lemma}[Alp\"oge]\label{cv-toroidal-extend-homomorphisms-quasi-finite}
Let $X$ be a scheme, locally of finite type over $\operatorname{Spec}\mathbf Z$, let $G$ be a group scheme over
$X$, locally of finite type and quasi-compact over $X$, and let $G'$ be a semi-abelian scheme over $X$. Let
$\Gamma$ be a scheme over $X$ and let $p:\Gamma\to G$ and $q:\Gamma\to G'$ be morphisms over $X$, with $p$ locally
of finite type. For a morphism $c:T\to\Gamma$, regard $T$ as a scheme over $X$ through $\Gamma$, so that $p\circ c$
and $q\circ c$ are points of $G$ and of $G'$ with values in $T$, elements of groups. Suppose:
\begin{enumerate}
\item[(a)] two morphisms $a,b:T\to\Gamma$ with $p\circ a=p\circ b$ and $q\circ a=q\circ b$ are equal;
\item[(b)] for every field $k$, every morphism $c:\operatorname{Spec}k\to\Gamma$ and every integer $n\ge1$, if
  $(p\circ c)^n=1$, then $(q\circ c)^n=1$.
\end{enumerate}
Then $p$ is locally quasi-finite.
\end{lemma}
\leanref{Development/AbelianVarietyTheory/Proofs/ModuliProofs/ExistsUniqueHomExtensionSemiAbelianNormal.lean}{GlobalTheory.HomExt.locallyQuasiFinite_of_torsion}
\begin{proof}
Step 1: closed points. The schemes $G$ and $\Gamma$ are locally of finite type over $\operatorname{Spec}\mathbf Z$,
which is Jacobson, so they are Jacobson \cite[Tag 02J5]{Sta}. The points of $\Gamma$ at which $p$ is quasi-finite
form an open subset \cite[Tag 01TI]{Sta}. In a Jacobson scheme the closed points are dense in every closed
subset, so an open subset that contains every closed point is the whole scheme. So it suffices that for every
closed point $w$ of $\Gamma$ the fibre of $p$ through $w$ is finite \cite[Tag 02NG]{Sta}.

Step 2: a finite group. Let $w$ be a closed point of $\Gamma$ and $\sigma=p(w)$. As $p$ is locally of finite type
and $G$ is Jacobson, $\sigma$ is a closed point of $G$; and its residue field $\kappa$ is finite, being finite over
the residue field of a closed point of $\operatorname{Spec}\mathbf Z$ \cite[Tag 01TB]{Sta}. Let
$x_0:\operatorname{Spec}\kappa\to G$ be the canonical morphism and $t_0$ its composite with the structure morphism
of $G$. The points of $G$ with values in $\operatorname{Spec}\kappa$ over $t_0$ form a finite group. Indeed, the
image of $t_0$ lies in an affine open subset $U$ of $X$; the inverse image of $U$ in $G$ is quasi-compact, so it
is covered by finitely many affine open subsets $\operatorname{Spec}A_l$, with each $A_l$ a finitely generated
$\mathbf Z$-algebra; and such a point lies in one of them and is given by a ring homomorphism $A_l\to\kappa$, of
which there are finitely many as $\kappa$ is finite. So $x_0^N=1$ for some $N\ge1$.

Step 3: the fibre. Let $\Omega$ be an algebraic closure of $\kappa$ and $y:\operatorname{Spec}\Omega\to G$ the
composite of $\operatorname{Spec}\Omega\to\operatorname{Spec}\kappa$ with $x_0$; its image is $\sigma$, and
$y^N=1$, as composition with $\operatorname{Spec}\Omega\to\operatorname{Spec}\kappa$ is a homomorphism of groups of
points. Let $x:\operatorname{Spec}\Omega\to\Gamma$ satisfy $p\circ x=y$. Then $(p\circ x)^N=1$, so
$(q\circ x)^N=1$ by (b). Now $q\circ x$ is a point with values in $\Omega$ of the base change of $G'$ along
$\operatorname{Spec}\Omega\to X$, a semi-abelian scheme over $\Omega$, and by
Lemma~\ref{cv-toroidal-extend-homomorphisms-torsion} only finitely many such points have order dividing $N$. By
(a), $x$ is determined by $q\circ x$, as $p\circ x=y$ is fixed. So there are finitely many morphisms $x$ with
$p\circ x=y$. They are the points with values in $\Omega$ of $P=\Gamma\times_G\operatorname{Spec}\Omega$ over
$\Omega$. As $P$ is locally of finite type over $\Omega$, it is Jacobson and its closed points are its points with
values in $\Omega$ \cite[Tag 01TB]{Sta}, which are finitely many; a finite set of closed points is closed, so by
the density of the closed points it is all of $P$. So $P$ is finite, and so is its image in $\Gamma$, which is the
fibre $p^{-1}(\sigma)$.
\end{proof}
% Lean: locallyQuasiFinite_of_torsion: jacobsonSpace_of_locallyOfFiniteType_int, locallyQuasiFinite_of_finite_fibres_closedPoints
% (quasiFiniteLocus open, subset_of_isOpen_of_closedPoints), closePoints_subset_preimage_closedPoints,
% jacobson_closedPoints_tools_int(2) (finite residue fields), SemiAbTorsion.finite_points_over (through
% finite_hom_range_subset), pow_card_eq_one', finite_torsion_points (Lemma
% cv-toroidal-extend-homomorphisms-torsion through TorsionLift.ptsEquiv), finite_preimage_of_finite_lifts with
% finite_of_finite_sections (pointEquivClosedPoint and JacobsonSpace.closure_inter_closedPoints). The hypothesis (b)
% is the Lean's key, stated as c o (p^n) = c o 1 for all n; for n = 0 it holds trivially.

\begin{theorem}\label{cv-toroidal-extend-homomorphisms}
Let $X$ be a normal scheme, locally of finite type and flat over $\operatorname{Spec}\mathbf Z$ (the
flatness is not used in the proof), and let
$j:W\to X$ be a flat, dominant, quasi-compact morphism that is a homeomorphism onto its image and
surjective on stalks. Let $G$ and $G'$ be semi-abelian schemes over $X$. Then every homomorphism
$j^*G\to j^*G'$ of group schemes over $W$ extends uniquely to a homomorphism $G\to G'$.
\end{theorem}
\leanref{Development/AbelianVarietyTheory/Proofs/ModuliProofs/ExistsUniqueHomExtensionSemiAbelianNormalStatement.lean}{GlobalTheory.T.existsUnique_hom_extension_semiAbelian_normal_base_alt}

When $j$ is the inclusion of a dense open subscheme, this is \cite[Chapter~I, Proposition~2.7]{FC90}.
The morphisms $j$ here are more general, and Alp\"oge derives the form stated here from its case of a
discrete valuation ring, Theorem~\ref{cv-model-toroidal-extension}.

\begin{proof}
Let $f:j^*G\to j^*G'$ be a homomorphism over $W$. Write $G_W=G\times_XW$, let $u:G_W\to G$ be the projection,
and let $\varphi:G_W\to G'$ be the composite of $f$ with the projection $G'\times_XW\to G'$; then $u$ and
$\varphi$ have the same composite to $X$. For a scheme $T$ over $X$, the points of $G$ and of $G'$ with values in
$T$ over $X$ form groups, and composition with a morphism $T'\to T$ over $X$ is a homomorphism. Write
$\iota=\operatorname{Spec}(V\subset K)$ for a discrete valuation ring $V$ with fraction field $K$; it is
scheme-theoretically dominant (Lemma~\ref{cv-toroidal-extend-homomorphisms-closure}), as $V\to K$ is injective.

Step 1: the base. As $G$ is smooth over the normal scheme $X$, it is normal \cite[Tag 033C]{Sta}; it is locally
noetherian, being locally of finite type over $\operatorname{Spec}\mathbf Z$. The morphism $u$ is a base change
of $j$, so it is flat, quasi-compact, a homeomorphism onto its image and surjective on stalks; in particular it
is a monomorphism. Its image is the inverse image of $j(W)$ in $G$, which is dense because $G\to X$ is flat and
locally of finite presentation, hence open; so $u$ is dominant. Its local homomorphisms are isomorphisms, as in
the proof of Lemma~\ref{cv-toroidal-extend-lift}, so $G_W$ is normal, and in particular reduced. The scheme $X$ is
reduced and $j$ is quasi-compact and dominant, so $j$ is scheme-theoretically dominant: its scheme-theoretic image
is a closed subscheme of $X$ whose underlying set is the closure of $j(W)$, which is $X$.

Step 2: the closure of the graph. Let $\Gamma$, $p$, $q$ and $i$ be as in
Lemma~\ref{cv-toroidal-extend-homomorphisms-closure}, for $S=X$, $Z_1=G$, $Z_2=G'$, $u$ and $\varphi$. As $G'$ is
separated, quasi-compact and locally of finite type over $X$, so is $p$, by part (3) of that lemma. By part (2),
$i$ is quasi-compact and scheme-theoretically dominant, hence dominant, and $\Gamma$ is reduced. By part (5), $i$
is a base change of $u$, so it is flat, a homeomorphism onto its image and surjective on stalks. The scheme
$\Gamma$ is locally of finite type over $G$, hence locally noetherian.

Step 3: discrete valuation rings. Let $V$ be a discrete valuation ring with fraction field $K$, and let
$a:\operatorname{Spec}V\to G$ and $c:\operatorname{Spec}K\to G_W$ satisfy $u\circ c=a\circ\iota$; regard
$\operatorname{Spec}V$ as a scheme over $X$ through $a$. Then there is a point $b$ of $G'$ with values in
$\operatorname{Spec}V$ such that $b\circ\iota=\varphi\circ c$ and, for every morphism $e:T\to\operatorname{Spec}V$ and
every $n\ge1$, $(b\circ e)^n=1$ whenever $(a\circ e)^n=1$. Indeed, let $g:\operatorname{Spec}V\to X$ be the composite
of $a$ with the structure morphism of $G$, and let $G_V$ and $G'_V$ be the base changes of $G$ and $G'$ along $g$,
semi-abelian schemes over $V$. The composite $w'$ of $c$ with the projection $G_W\to W$ satisfies
$j\circ w'=g\circ\iota$, so the base change of $f$ along $w'$ is a homomorphism $f_K$ between the generic fibres
of $G_V$ and $G'_V$. By Theorem~\ref{cv-model-toroidal-extension}, $f_K$ extends to a homomorphism
$F:G_V\to G'_V$. Let $s:\operatorname{Spec}V\to G_V$ be the section given by $a$, and let $b$ be the composite of
$F\circ s$ with the projection $G'_V\to G'$. Over $K$, $s$ is the point of the generic fibre of $G_V$ given by $c$,
which $f_K$ maps to the point given by $f\circ c$; so $b\circ\iota=\varphi\circ c$. The points of $G$ over $X$ with
values in $T$ through $a\circ e$ correspond to the points of $G_V$ with values in $T$, compatibly with the group
laws, and $a\circ e$ corresponds to $s\circ e$; in the same way $b\circ e$ corresponds to $F\circ s\circ e$. As $F$
is a homomorphism, $(s\circ e)^n=1$ implies $(F\circ s\circ e)^n=1$.

Step 4: $p$ is surjective. Let $\sigma\in G$. If $\sigma=u(w)$ for some $w\in G_W$, then $\sigma=p(i(w))$.
Otherwise Lemma~\ref{cv-toroidal-extend-homomorphisms-points}(2), applied to $u$, gives $V$, $K$, a morphism
$a:\operatorname{Spec}V\to G$ that maps the closed point to $\sigma$, and $c:\operatorname{Spec}K\to G_W$ with
$u\circ c=a\circ\iota$. Let $b$ be as in Step 3. Then $a$ and $b$ have the same composite to $X$,
$a\circ\iota=u\circ c$ and $b\circ\iota=\varphi\circ c$, so by Lemma~\ref{cv-toroidal-extend-homomorphisms-closure}(4)
there is $l:\operatorname{Spec}V\to\Gamma$ with $p\circ l=a$. So $\sigma$ is the image under $p$ of the image of the
closed point under $l$.

Step 5: $p$ is locally quasi-finite. Lemma~\ref{cv-toroidal-extend-homomorphisms-quasi-finite} applies to $X$, $G$,
$G'$, $\Gamma$, $p$ and $q$: condition (a) is part (1) of Lemma~\ref{cv-toroidal-extend-homomorphisms-closure}. For
condition (b), let $c:\operatorname{Spec}k\to\Gamma$ satisfy $(p\circ c)^n=1$, and let $x$ be the image of $c$.

If $x\in i(G_W)$, then by Lemma~\ref{cv-toroidal-extend-homomorphisms-points}(1), applied to $i$, $c=i\circ c'$
for some $c':\operatorname{Spec}k\to G_W$, so that $p\circ c=u\circ c'$ and $q\circ c=\varphi\circ c'$. Let $w'$ be the
composite of $c'$ with $G_W\to W$. The points of $G$ with values in $\operatorname{Spec}k$ over $j\circ w'$ correspond
to the points of $j^*G$ over $w'$, compatibly with the group laws, and $u\circ c'$ corresponds to $c'$; in the same
way $\varphi\circ c'$ corresponds to $f\circ c'$. As $f$ is a homomorphism and $(c')^n=1$, $(f\circ c')^n=1$, so
$(q\circ c)^n=1$.

If $x\notin i(G_W)$, then Lemma~\ref{cv-toroidal-extend-homomorphisms-points}(2), applied to $i$, which is
possible by Step 2, gives $V$, $K$, a morphism $a:\operatorname{Spec}V\to\Gamma$ that maps the closed point to $x$,
and $c':\operatorname{Spec}K\to G_W$ with $i\circ c'=a\circ\iota$. As $u\circ c'=p\circ i\circ c'=p\circ a\circ\iota$,
Step 3 applies to $p\circ a$ and $c'$; let $b$ be the point it gives. Then
$b\circ\iota=\varphi\circ c'=q\circ i\circ c'=q\circ a\circ\iota$. The morphisms $b$ and $q\circ a$ lie over the same
morphism to $X$, and $G'$ is separated over $X$, so the locus where they agree is a closed subscheme of
$\operatorname{Spec}V$ through which $\iota$ factors; as $\iota$ is scheme-theoretically dominant, it is all of
$\operatorname{Spec}V$, and $b=q\circ a$. Let $r$ be the inclusion of the closed point of $\operatorname{Spec}V$. The
morphisms $c$ and $a\circ r$ both factor through the canonical morphism $\operatorname{Spec}\kappa(x)\to\Gamma$,
through morphisms of spectra of fields, which are faithfully flat and hence epimorphisms. So two morphisms from
$\Gamma$ to $G$, or to $G'$, agree after composition with $c$ if and only if they agree after composition with
$a\circ r$. Applied to the $n$-th power of $p$, as a point of $G$ with values in $\Gamma$, and to $1$, this gives
$(p\circ a\circ r)^n=1$. By Step 3 with $e=r$, $(q\circ a\circ r)^n=(b\circ r)^n=1$, and by the same equivalence,
applied to $q$, $(q\circ c)^n=1$.

Step 6: $p$ is an isomorphism. Lemma~\ref{cv-toroidal-extend-homomorphisms-isomorphism} applies with $Y=G$, which
is normal and locally noetherian by Step 1: $\Gamma$ is reduced and $p$ is separated, quasi-compact and locally of
finite type by Step 2, surjective by Step 4 and locally quasi-finite by Step 5; $i$ has the required properties
by Step 2, and $p\circ i=u$ by Step 1.

Step 7: the extension. By Lemma~\ref{cv-toroidal-extend-homomorphisms-closure}(6), $\psi=q\circ p^{-1}:G\to G'$ is a
morphism over $X$ with $\psi\circ u=\varphi$. Its base change $j^*\psi:G_W\to G'\times_XW$ is the morphism over $W$
whose composite with the projection to $G'$ is $\psi\circ u=\varphi$, which is also the composite of $f$ with that
projection; so $j^*\psi=f$.

Step 8: uniqueness and the homomorphism property. Let $T$ be a scheme flat over $X$, and let $\psi_1,\psi_2:T\to G'$
be morphisms over $X$ whose base changes along $j$ agree. The projection $T\times_XW\to T$ is the base change of
the quasi-compact, scheme-theoretically dominant morphism $j$ along the flat morphism $T\to X$, so it is
scheme-theoretically dominant \cite[Tag 081I]{Sta}, and $\psi_1$ and $\psi_2$ agree after composition with it. As
$G'$ is separated over $X$, the locus where $\psi_1$ and $\psi_2$ agree is a closed subscheme of $T$ through which
$T\times_XW\to T$ factors, so it is $T$, and $\psi_1=\psi_2$. For $T=G$ this is the uniqueness of the extension. The
schemes $X$ and $G\times_XG$ are flat over $X$, and the identities that make $\psi$ a homomorphism, for the unit
sections and for the multiplications, are equalities of morphisms from them to $G'$ over $X$ whose base changes
along $j$ are the same identities for $f$, which hold. So $\psi$ is a homomorphism.
\end{proof}
\leanref{Development/AbelianVarietyTheory/Proofs/ModuliProofs/ExistsUniqueHomExtensionSemiAbelianNormal.lean}{GlobalTheory.existsUnique_hom_extension_semiAbelian_normal}
\leanref{Development/AbelianVarietyTheory/Proofs/ModuliProofs/CompPowEqOneOfPullbackHom.lean}{GlobalTheory.HomExt.exists_hom_along_dvr}
\leanref{Development/AbelianVarietyTheory/Proofs/ModuliProofs/CompPowEqOneOfPullbackHom.lean}{GlobalTheory.HomExt.comp_pow_eq_one_of_pullback_hom}
\leanref{Development/AbelianVarietyTheory/Proofs/UnitaryExistence/ExistsLiftOfSchemeTheoreticallyDominant.lean}{GlobalTheory.pullback_map_injective_of_flat}
\leanref{Development/AbelianVarietyTheory/Proofs/UnitaryExistence/ExistsLiftOfSchemeTheoreticallyDominant.lean}{GlobalTheory.isMonHom_of_pullback_of_flat}
% Lean: existsUnique_hom_extension_semiAbelian_normal, with HomExtSt (Theorem cv-model-toroidal-extension) as its first
% hypothesis; the flatness of X over Spec Z is discarded in its first line. Step 1: forall_isDomain_and_isIntegrallyClosed_stalk_of_smooth,
% isDominant_pullback_fst_of_isOpenMap, IsNormalScheme.of_isIso_stalkMap, IsSchemeTheoreticallyDominant.of_isDominant.
% Step 2: graphClosure G.hom G'.hom u phi with isPullback_incl, schemeDominant_incl, isReduced, locallyOfFiniteType_fst,
% quasiCompact_fst, isSeparated_fst. Step 3 is exists_hom_along_dvr (through exists_iso_pullback_pullback and
% TorsionLift.ptsEquiv) with comp_pow_eq_one_of_pullback_hom for the powers. Step 4: mem_range_or_exists_dvr,
% exists_hom_along_dvr and graphClosure.existsUnique_lift with schemeTheoreticallyDominant_fractionField. Step 5:
% locallyQuasiFinite_of_torsion with graphClosure.hom_ext; the first case is exists_lift_of_range_subset and
% comp_pow_eq_one_of_pullback_hom with g = j; the second is mem_range_or_exists_dvr for i,
% over_ext_of_schemeTheoreticallyDominant (b = q o a, hb) and field_point_comp_eq_iff. Step 6: isIso_of_densePart.
% Step 7: graphClosure.exists_extension_of_isIso_fst. Step 8: pullback_map_injective_of_flat (through
% over_ext_of_schemeTheoreticallyDominant) and isMonHom_of_pullback_of_flat. Size: about 1,900 lines of Lean below the
% theorem besides Lemma cv-toroidal-extend-homomorphisms-torsion's citation.

\begin{lemma}[Alp\"oge]\label{cv-toroidal-extend-identify}
In Situation~\ref{cv-function-situation}, let $X$ be a normal scheme, locally of finite type and flat over
$\operatorname{Spec}\mathbf Z$, and let $j:W\to X$ be a flat, dominant, quasi-compact morphism that is a
homeomorphism onto its image and surjective on stalks. Let $G_0$ be an $\mathcal O_E$-semi-abelian scheme
over $W$, proper over $W$, and let $G$ and $G'$ be $\mathcal O_E$-semi-abelian schemes over $X$ with
morphisms $a:G_0\to G$ and $a':G_0\to G'$ over $j$ (Definition~\ref{cv-model-toroidal-semiabelian}(2)).
Then there is a morphism $b:G\to G'$ of $\mathcal O_E$-semi-abelian schemes over the identity of $X$ with
$b\circ a=a'$.
\end{lemma}
\leanref{Development/AbelianVarietyTheory/Proofs/ModuliProofs/ExistsSabMapIdOfCommonGeneric.lean}{GlobalTheory.exists_sabMap_id_of_common_generic}
\begin{proof}
As in the proof of Lemma~\ref{cv-toroidal-extend-lift}, the generic points of $X$ lie in $j(W)$. The
scheme $X$ is reduced, being normal, and locally noetherian. The morphisms $a$ and $a'$ identify $j^*G$
and $j^*G'$ with $G_0$. By Theorem~\ref{cv-toroidal-extend-homomorphisms}, the resulting isomorphism
$j^*G\to j^*G'$ extends to a homomorphism $F:G\to G'$ with $F\circ a=a'$, and its inverse to
$F':G'\to G$ with $F'\circ a'=a$. By the uniqueness in that theorem, $F'\circ F$ and $F\circ F'$ are the
identities, so $F$ is an isomorphism. For $c\in\mathcal O_E$, the homomorphisms $F\circ c$ and $c\circ F$
agree over $W$, as $a$ and $a'$ commute with $c$, so they are equal.

Let $\mathcal L$ and $\mathcal L'$ be the line bundles of $G$ and $G'$, and let
$\mathcal N=F^*\mathcal L'\otimes\mathcal L^{-1}$. Both are trivialized along the zero section, so
$\mathcal N$ is trivial along the zero section. Over $G_0$, both $\mathcal L$ and $F^*\mathcal L'$ pull back
to the line bundle of $G_0$, so the pullback of $\mathcal N$ to $G\times_XW$ is trivial. The scheme $G$ is
smooth over $X$ with geometrically integral fibres, being a smooth group scheme with geometrically
connected fibres. By Lemma~\ref{cv-toroidal-extend-trivial}, $\mathcal N$ is trivial, so
$\mathcal L\cong F^*\mathcal L'$. The generic points of $X$ lie in $j(W)$, over which $G$ is the proper
$G_0$, so they lie in the abelian locus of $G$. By Lemma~\ref{cv-toroidal-universal-cubical}, $F$ is a
morphism of $\mathcal O_E$-semi-abelian schemes over the identity, and $b=F$ will do.
\end{proof}
% Lean: exists_sabMap_id_of_common_generic: densePart_basic, existsUnique_hom_extension_semiAbelian_normal_base_alt
% with TripleExt.exists_hom (F and F'), ExtUnique.isIso_of_isPullback, ExtUnique.over_hom_ext_of_isPullback
% (actions), exists_quot_of_isos and exists_quot_of_gens (the two trivializations of N),
% exists_lineBundleGenerator_extending_densePart, exists_sabMap_of_lbIso.

\begin{lemma}[Alp\"oge]\label{cv-toroidal-extend-bundles}
In Situation~\ref{cv-function-situation}, let $Z$ be a normal, locally noetherian scheme, flat over
$\operatorname{Spec}\mathbf Z$, and let $G$ be an $\mathcal O_E$-semi-abelian scheme over $Z$ whose abelian
locus contains the generic points of $Z$. For $k=1,2$, let $G_k$ be an $\mathcal O_E$-semi-abelian scheme
over a scheme $Z_k$ with an ample theta presentation of level $N$ with ample line bundle $\mathcal A_k$
(Definition~\ref{cv-toroidal-envelope}(1)), whose integers $r$ and $s$ do not depend on $k$, and let
$f_k:Z\to Z_k$ be a morphism with a morphism $G\to G_k$ of $\mathcal O_E$-semi-abelian schemes over
$f_k$. Then $f_1^*\mathcal A_1\cong f_2^*\mathcal A_2$.
\end{lemma}
\leanref{Development/AbelianVarietyTheory/Proofs/ModuliProofs/LineBundleIsoOfAmpleDeterminantFramesStatement.lean}{GlobalTheory.T.lineBundleIso_of_ampleDeterminantFrames}
\leanref{Development/AbelianVarietyTheory/Proofs/ModuliProofs/LineBundleIsoOfAmpleDeterminantFrames.lean}{GlobalTheory.FrameCmp.lbIso_of_ampFrames}
\begin{proof}
Let $\mathcal L$ be the line bundle of $G$. The morphism $G\to G_k$ identifies $G$ with $f_k^*G_k$ and
$\mathcal L$ with the pullback of the line bundle of $G_k$. So the sections of the presentation of $G_k$
over a set $U$ of its cover pull back to a section $\vartheta_U$ of $\mathcal L$ and sections
$b_{U,1},\dots,b_{U,n_k}$ of $\mathcal L^N$ over $G$ restricted to $f_k^{-1}(U)$; write $b_U$ for the
column of the $b_{U,m}$. For sets $U$ of the cover of $Z_1$ and $U'$ of that of $Z_2$, let
$V=f_1^{-1}(U)\cap f_2^{-1}(U')$.

Step 1. Let $W$ be a complete discrete valuation ring of characteristic zero and
$\chi:\operatorname{Spec}W\to V$ a morphism that maps the generic point to a generic point of $Z$. That
point lies in the abelian locus of $G$, so $f_k\circ\chi$ maps it into the abelian locus of $G_k$, and
condition (b) of Definition~\ref{cv-toroidal-envelope}(1) says that $\chi^*\vartheta_U$ and
$\chi^*\vartheta_{U'}$ are bases of the sections of $\mathcal L$ over $G\times_Z\operatorname{Spec}W$, and
$\chi^*b_U$ and $\chi^*b_{U'}$ bases of those of $\mathcal L^N$. Hence $\vartheta_U=c\,\vartheta_{U'}$ and
$b_U=P\,b_{U'}$ over $W$, with $c\in W^\times$ and $P$ an invertible matrix over $W$; in particular
$n_1=n_2$ if $V$ is nonempty. For a field $F$ of characteristic zero and a morphism
$\operatorname{Spec}F\to V$ onto a generic point of $Z$, the same relations hold over $F$: apply the
above to $W=F[[t]]$ and set $t=0$. The sections in each family are linearly independent over the field,
so these relations are unique.

Step 2. Let $\Omega\subset V$ be an affine open subset whose ring $B$ is a domain. It is noetherian and
normal, and it has characteristic zero because $Z$ is flat over $\mathbf Z$; let $F$ be its fraction
field. By Step 1 over $F$, $\vartheta_U=c\,\vartheta_{U'}$ and $b_U=P\,b_{U'}$ with $c\in F^\times$ and $P$
invertible over $F$. Let $\mathfrak p$ be a prime of height one of $B$ and $W$ the completion of the
discrete valuation ring $B_{\mathfrak p}$. Step 1 over $W$ gives such relations over $W$, which agree with
those over $F$ in the fraction field of $W$, by uniqueness. So $c$ and $\det P$ lie in $F\cap W=B_{\mathfrak
p}$. As $B$ is noetherian and normal, $c$ and $\det P$ lie in $B$, and by symmetry so do their inverses.

Step 3. As $Z$ is normal and locally noetherian, $V$ is covered by affine open subsets as in Step 2. The
units of Step 2 agree on overlaps, since they agree at the generic points and $Z$ is reduced. So there are
units $c_{U,U'}$ and $d_{U,U'}$ on $V$ such that at every generic point of $V$, $\vartheta_U=
c_{U,U'}\vartheta_{U'}$ and $b_U=P\,b_{U'}$ with $\det P=d_{U,U'}$.

Step 4. The units $c_{U,U'}^r d_{U,U'}^s$ on the sets $V$ define an isomorphism
$f_1^*\mathcal A_1\cong f_2^*\mathcal A_2$. Let $(U_1,U'_1)$ and $(U_2,U'_2)$ be two pairs of sets. By
Definition~\ref{cv-toroidal-envelope}(1)(a), $\vartheta_{U_1}=c_1\vartheta_{U_2}$ and $b_{U_1}=M_1b_{U_2}$
on $f_1^{-1}(U_1\cap U_2)$, where the transition function of $\mathcal A_1$ is $c_1^r\det(M_1)^s$, and
likewise $c_2$, $M_2$ for $U'_1$, $U'_2$ and $\mathcal A_2$. At a generic point of $Z$ in all four sets,
expressing $\vartheta_{U_1}$ and $b_{U_1}$ through $\vartheta_{U'_2}$ and $b_{U'_2}$ in two ways gives, by
uniqueness, $c_{U_1,U'_1}c_2=c_1c_{U_2,U'_2}$ and $d_{U_1,U'_1}\det M_2=\det M_1\,d_{U_2,U'_2}$. Hence
\[
 c_{U_1,U'_1}^rd_{U_1,U'_1}^s\cdot c_2^r\det(M_2)^s=c_1^r\det(M_1)^s\cdot c_{U_2,U'_2}^rd_{U_2,U'_2}^s
\]
at every generic point, hence everywhere, as $Z$ is reduced. This is the compatibility of the local
isomorphisms with the transition functions of $f_1^*\mathcal A_1$ and $f_2^*\mathcal A_2$.
\end{proof}
% Lean: FrameCmp.lbIso_of_ampFrames: exists_values (theta and b pulled back as probe values, Bas from
% the trait condition by bas_of_basisAlong, Tr from A.trans by rel_of_lin); GeoCmp, Cmp.of_powerSeries and
% TraitCmp.field (Step 1); TraitCmp.det_mem with mem_of_completion and TraitCmp.exists_unit (Step 2);
% exists_domain_chart, GeoCmp.exists_unit with glue_unit and unit_ext_genPt (Step 3); lbIso_of_cmp with
% det_square (Step 4).

\begin{lemma}[Alp\"oge]\label{cv-toroidal-extend-finite}
Let $S$ be an affine scheme, let $Z_1$ and $Z_2$ be schemes over $S$ with $Z_2$ proper over $S$, and let
$\mathcal A_1$ and $\mathcal A_2$ be line bundles on $Z_1$ and $Z_2$ such that every point of $Z_2$ lies in
the nonvanishing locus, an affine open subset, of a section of $\mathcal A_2^d$ for some $d\ge1$. Let
$\pi:\tilde Z\to Z_1\times_SZ_2$ be a finite morphism, let $p_1$ and $p_2$ be its composites with the
projections, and suppose that $p_1^*\mathcal A_1\cong p_2^*\mathcal A_2$. Then $p_1$ is finite.
\end{lemma}
\leanref{Development/AbelianVarietyTheory/Proofs/ModuliProofs/ExistsHomBirationalModelsAmpleFrames.lean}{GlobalTheory.AmpUniq.isFinite_fst}
\leanref{Development/AbelianVarietyTheory/Proofs/ModuliProofs/LocallyQuasiFiniteOfTrivialOnFibres.lean}{GlobalTheory.AmpleFibre.locallyQuasiFinite_of_trivial_on_fibres}
\begin{proof}
The morphism $p_1$ is proper, as $\pi$ is finite and $Z_1\times_SZ_2\to Z_1$ is proper. A proper, locally
quasi-finite morphism is finite, so it suffices that $p_1$ has finite fibres. Let $y\in Z_1$, let $U$ be an
affine open neighbourhood of $y$, let $\bar k$ be an algebraic closure of the residue field of $y$, and let
$\Phi=\tilde Z\times_{Z_1}\operatorname{Spec}\bar k$, a scheme proper over $\bar k$. The morphism
$q:\Phi\to Z_2$ induced by $p_2$ is affine, because $\pi$ is finite and $U\times_SZ_2\to Z_2$ is affine.
The pullback of $\mathcal A_1$ to $\operatorname{Spec}\bar k$ is trivial, so $q^*\mathcal A_2$ is trivial,
with a generator $\tau$. Let $C$ be an irreducible component of $\Phi$, with its reduced structure. For a
section $\sigma$ of $\mathcal A_2^d$ with affine nonvanishing locus, $q^*\sigma=a\tau^d$ with $a$ a regular
function on $C$, and the locus $D(a)$ where $a$ does not vanish is the inverse image of that of $\sigma$,
an affine scheme. The regular functions on $C$ form a finite-dimensional algebra over $\bar k$, as $C$ is
proper, so $a$ is integral over $\bar k$ and $D(a)$ is also closed. So $D(a)$ is proper and affine over
$\bar k$, hence finite. These sets cover $C$, so $C$ is finite, and so is $\Phi$, which has finitely many
irreducible components.
\end{proof}
% Lean: isFinite_fst: IsFinite.of_isProper_of_locallyQuasiFinite; locally quasi-finite on an affine cover of
% Z1 by locallyQuasiFinite_of_trivial_on_fibres (the affine morphism q, DVR.exists_gen_of_local and
% SiegelDescent.exists_gen_of_lbIso for the generator), with exists_fn, exists_integral_closedImmersion and
% MulFinite.locallyQuasiFinite_of_cover (isClosed_basicOpen).

\begin{lemma}[Alp\"oge]\label{cv-toroidal-extend}
In Situation~\ref{cv-function-situation}, let $R$ be a Dedekind domain of characteristic zero and of
finite type over $\mathbf Z$, let $Z_1$ and $Z_2$ be schemes projective over $R$, with $Z_1$ normal and
flat over $R$, and let $G_1$ and $G_2$ be $\mathcal O_E$-semi-abelian schemes over $Z_1$ and $Z_2$ with
ample theta presentations of the same level and with the same integers $r$ and $s$
(Definition~\ref{cv-toroidal-envelope}(1)). Let $W$ be a scheme with an $\mathcal O_E$-semi-abelian
scheme $G_W$ that is proper over $W$, and let $j_1:W\to Z_1$ and $j_2:W\to Z_2$ be flat, dominant,
quasi-compact morphisms over $R$, each a homeomorphism onto its image and surjective on stalks, with
morphisms $G_W\to G_1$ over $j_1$ and $G_W\to G_2$ over $j_2$. Then there is a morphism
$\rho:Z_1\to Z_2$ with $\rho\circ j_1=j_2$.
\end{lemma}
\leanref{Development/AbelianVarietyTheory/Proofs/ModuliProofs/ExistsHomBirationalModelsAmpleFramesStatement.lean}{GlobalTheory.T.exists_hom_birational_models_ampleFrames}
\leanref{Development/AbelianVarietyTheory/Proofs/ModuliProofs/ExistsHomBirationalModelsAmpleFrames.lean}{GlobalTheory.AmpUniq.exists_hom_of_ampFrames}
\begin{proof}
The schemes $Z_1$ and $Z_2$ are proper over $R$, and $Z_1$ is noetherian. As in the proof of
Lemma~\ref{cv-toroidal-extend-lift}, $j_1$ induces isomorphisms of local rings, so $W$ is normal and
locally noetherian, and $W$ is flat over $R$ because $Z_1$ is. Let $f:W\to Z_1\times_RZ_2$ be the morphism
with components $j_1$ and $j_2$. Its composite with the first projection is $j_1$, so $f$ is
quasi-compact, a homeomorphism onto its image and surjective on stalks.

Let $\tilde Z$ be the normalization of $Z_1\times_RZ_2$ in $W$, with $u:W\to\tilde Z$ and the integral
morphism $\pi:\tilde Z\to Z_1\times_RZ_2$, and let $p_1,p_2$ be the composites of $\pi$ with the
projections, so that $p_1\circ u=j_1$ and $p_2\circ u=j_2$. The generic points of $W$ lie over the
generic point of $\operatorname{Spec}R$, so their residue fields have characteristic zero; they are the
residue fields of their images, as $f$ is surjective on stalks; and $R$ is of finite type over $\mathbf Z$.
So $\pi$ is finite \cite[Scholie~7.8.3]{Gro65}. The scheme $\tilde Z$ is normal
and locally noetherian, and it is flat over the Dedekind domain $R$, since it has no torsion. The morphism
$u$ is quasi-compact, a homeomorphism onto its image and surjective on stalks, as $p_1\circ u=j_1$ is;
it is dominant, and flat, since its local homomorphisms are surjective and, as $\tilde Z$ is the
normalization in $W$, injective.

Let $G'_k=p_k^*G_k$ for $k=1,2$. The morphisms $G_W\to G_k$ over $j_k=p_k\circ u$ factor through morphisms
$G_W\to G'_k$ over $u$. By Lemma~\ref{cv-toroidal-extend-identify}, applied to $u$, there is a morphism
$G'_1\to G'_2$ over the identity of $\tilde Z$; composed with $G'_2\to G_2$, it gives a morphism
$G'_1\to G_2$ over $p_2$. The generic points of $\tilde Z$ lie in $u(W)$, over which $G'_1$ is the proper
$G_W$, so they lie in the abelian locus of $G'_1$. Lemma~\ref{cv-toroidal-extend-bundles}, with $Z=\tilde Z$,
$G=G'_1$, $f_1=p_1$ and $f_2=p_2$, gives $p_1^*\mathcal A_1\cong p_2^*\mathcal A_2$. By
Lemma~\ref{cv-toroidal-extend-finite}, with $S=\operatorname{Spec}R$, $p_1$ is finite.

Finally, Lemma~\ref{cv-toroidal-extend-lift}, with $T=X=Z_1$, $j=j_1$, $q=p_1$, $f$ the identity and $g=u$,
gives $h:Z_1\to\tilde Z$ with $h\circ j_1=u$ and $p_1\circ h$ the identity. Then $\rho=p_2\circ h$ satisfies
$\rho\circ j_1=p_2\circ u=j_2$.
\end{proof}
% Lean: AmpUniq.exists_hom_of_ampFrames: isIso_stalkMap_of_flat_of_surjectiveOnStalks and
% isLocallyNoetherian_of_isPreimmersion (W), pullback.lift j j' (f), isFinite_fromNormalization_graph (pi),
% flat_of_isPreimmersion (u), flat_normalization, sabPull and ObjMaps.exists_sab_cancel (G'_k),
% exists_sabMap_id_of_common_generic (with HomExtSt), ObjMaps.exists_sab_comp,
% T.lineBundleIso_of_ampleDeterminantFrames with genericPoints_subset_abLocus, isFinite_fst,
% existsUnique_lift_of_densePart.

\begin{lemma}[Alp\"oge]\label{cv-model-toroidal-top}
In Situation~\ref{cv-function-situation}, suppose $p+q\ge1$, and let $n$ be an auxiliary prime
(Lemma~\ref{cv-model-toroidal-auxiliary}). Then there are an integer $J\ge1$ and, for $N=n^J$, a
normal scheme $Z$, projective and flat over $R_{n,N}$, with a theta-split $\mathcal O_E$-semi-abelian
scheme $G$ that is quasi-projective over $R_{n,N}$ and abelian at the generic points of $Z$
(Definition~\ref{cv-model-toroidal-semiabelian}), with the following properties for every level $l$
of full level $N$.
\begin{enumerate}
\item There are a morphism $a_l:X_l\to Z$ and a morphism $\mathcal A_l\to G$ of
  $\mathcal O_E$-semi-abelian schemes over $a_l$. The morphism $a_l$ is the composite of a finite
  surjective morphism over $K_l$, from $X_l$ onto the open subscheme of $Z\otimes_{R_{n,N}}K_l$ over
  which $G$ is abelian, with the projection to $Z$.
\item Let $\sigma$ be an automorphism of $K_l$ with $\sigma(\sqrt{-d})=\sqrt{-d}$, let $r$ be an
  automorphism of $X_l$ over $\sigma$, and let $g$ be an $\mathcal O_E$-linear endomorphism of
  $\Lambda$. If $\mathcal A_l$ has an automorphism $\nu$ over $r$, as an $\mathcal O_E$-semi-abelian
  scheme, with $\nu\circ\eta'(g(y))=\eta'(y)\circ r$ for all $y\in\Lambda$, then there is an
  automorphism $\rho$ of $Z$ with $\rho\circ a_l=a_l\circ r$.
\end{enumerate}
\end{lemma}
\leanref{Development/AbelianVarietyTheory/Definitions/ModuliSpaces/AuxiliaryInertOddPrime.lean}{GlobalTheory.Models.Top}
\leanref{Development/AbelianVarietyTheory/Definitions/ModuliSpaces/TopAmpExists.lean}{GlobalTheory.Models.TopAmpExists}
\leanref{Development/AbelianVarietyTheory/Proofs/ModuliProofs/ExistsFullLevelAmpleModelOfSemistableReduction.lean}{GlobalTheory.topAmpExists_of_semistableReduction}
\leanref{Development/AbelianVarietyTheory/Proofs/ModuliProofs/ExistsFullLevelModelAmpleFrames.lean}{GlobalTheory.Models.topAmpExists}
\begin{proof}
The construction is Alp\"oge's, from Mumford's algebraic theta functions \cite{Mum66}. Let
$g=\frac12[E:\mathbf Q](p+q)\ge1$, so that $\Lambda$ is a free $\mathbf Z$-module of rank $2g$ with a
perfect alternating form. Choose a symplectic basis $e_1,\dots,e_{2g}$ of $\Lambda$ and a map from a
finite set $I$ to $\mathcal O_E$ whose image generates $\mathcal O_E$ as a ring and contains
$w_0=(1+\sqrt{-d})/2$ and $\bar w_0$. Lemma~\ref{cv-toroidal-theta-family} gives $j_1$, and
Lemma~\ref{cv-toroidal-theta-envelope} gives $c$ and $j_2$, such that their conclusions hold for every
$j\ge j_1$ and every $j\ge j_2$. Let $j=\max(j_1,j_2,1)$, $N'=n^j$ and $N=N'^2$, and let $J=2j$, so that
$N=n^J$ and $J\ge1$.

The ring $R_{n,N}$ is a Dedekind domain of finite type over $\mathbf Z$, contained in $\mathbf C$ and
containing $\theta$. The integer $N$ is invertible in it, and $\zeta_N$ maps to a primitive $N$-th root
of unity under every homomorphism to a field, since it is a root of the $N$-th cyclotomic polynomial and
$N$ is invertible there. So Lemma~\ref{cv-toroidal-theta-family} applies with $o=R_{n,N}$,
$\zeta=\zeta_N$, $\mathrm{em}$ the inclusion and $\rho=\theta$. It gives a PEL theta family
$(\mathcal M,t,G_{\mathcal M})$ of level $N'$ (Definition~\ref{cv-toroidal-theta}(7)), with structure
morphism $s_{\mathcal M}:\mathcal M\to\operatorname{Spec}R_{n,N}$.
Lemma~\ref{cv-toroidal-theta-envelope} gives a theta envelope $(Z,G,W,j)$ of
$(\mathcal M,G_{\mathcal M})$ of level $n^{2c}$ over $R_{n,N}$, with $G_W$ over $W$ as in
Definition~\ref{cv-toroidal-envelope}(2), and Lemma~\ref{cv-toroidal-theta-universal} gives
$\lambda:\Lambda\to G_{\mathcal M}(\mathcal M)$ such that $(\mathcal M,G_{\mathcal M},\lambda)$ is a
full-level family over $R_{n,N}$.

The scheme $Z$ is normal, projective and flat over $R_{n,N}$, and $G$ is theta-split and
quasi-projective over $R_{n,N}$. So $Z$ is reduced, separated and noetherian. Let $U$ be the abelian
locus of $G$, which is open by condition (a) of Definition~\ref{cv-model-toroidal-semiabelian}(3). Over
$W$, the pullbacks of $G$ along $j$ and of $G_{\mathcal M}$ along the inclusion are both isomorphic to
$G_W$, and $G_{\mathcal M}$ is proper; so $j(W)\subset U$, and $j$ factors through a morphism
$j':W\to U$. Since $j$ is dominant and $Z$ is noetherian, the generic points of $Z$ lie in $j(W)$, so
$G$ is abelian at the generic points of $Z$.

Lemma~\ref{cv-toroidal-abelian-locus} gives a finite surjective morphism $u:U\to\mathcal M$ over
$R_{n,N}$, with $G_U$ over $U$ and morphisms $G_U\to G$ over the inclusion $\iota_U$ and
$G_U\to G_{\mathcal M}$ over $u$, and an open subscheme $W_0$ of $\mathcal M$, with inclusion
$\iota_{W_0}$, containing every point over the generic point of $\operatorname{Spec}R_{n,N}$, over
which $u$ is an isomorphism; moreover $u\circ j'=\iota_W$, the inclusion of $W$. Let $i:W_0\to U$ be the
inverse of $u^{-1}(W_0)\to W_0$ followed by the inclusion of $u^{-1}(W_0)$. Then $i$ is an open
immersion with $u\circ i=\iota_{W_0}$, and:
\begin{enumerate}
\item[$(*)$] if $h:T\to U$ and $b:T\to W_0$ satisfy $u\circ h=\iota_{W_0}\circ b$, then $h=i\circ b$,
\end{enumerate}
since $h$ then factors through $u^{-1}(W_0)$, which $u$ maps isomorphically onto $W_0$.

Let $l$ be a level of full level $N$, with integer $M$, so that $K_l=\mathbf Q(\sqrt{-d},\zeta_N)$
contains $R_{n,N}$; let $\beta:R_{n,N}\to K_l$ be the inclusion, which is injective and flat. By
Lemma~\ref{cv-model-toroidal-auxiliary}, $n$ does not divide $M$, so $N$ is prime to $M$.
Lemma~\ref{cv-toroidal-moduli-finite} gives $a:X_l\to\mathcal M$ with
$s_{\mathcal M}\circ a=\operatorname{Spec}\beta\circ s_l$ and a morphism
$\nu_0:\mathcal A_l\to G_{\mathcal M}$ over $a$ carrying $\eta'$ to $\lambda$, such that
$X_l\to\mathcal M\otimes_{R_{n,N}}K_l$ is finite and surjective. As $\beta$ is injective, every point
of $X_l$ maps to a point of $\mathcal M$ over the generic point of $\operatorname{Spec}R_{n,N}$, which
lies in $W_0$; so $a=\iota_{W_0}\circ b$ for a morphism $b:X_l\to W_0$. Put
$a_l=\iota_U\circ i\circ b$. Since $u\circ i\circ b=a$ and $G_U\to G_{\mathcal M}$ identifies $G_U$ with
the pullback of $G_{\mathcal M}$ along $u$, $\nu_0$ factors through a morphism $\mathcal A_l\to G_U$ over
$i\circ b$, and its composite with $G_U\to G$ is a morphism $\mathcal A_l\to G$ over $a_l$.

Let $P$ be the open subscheme of $Z\otimes_{R_{n,N}}K_l$ over which $G$ is abelian, the preimage of $U$.
The image of $\mathcal M\otimes_{R_{n,N}}K_l$ in $\mathcal M$ lies in $W_0$, by the same argument. So the
morphism $\mathcal M\otimes_{R_{n,N}}K_l\to P$ given by $\iota_U\circ i$ on the first factor and the
identity on the second is defined; it is an isomorphism, with inverse given by $u$ on the first factor,
by $u\circ i=\iota_{W_0}$ and $(*)$. Its composite with $X_l\to\mathcal M\otimes_{R_{n,N}}K_l$ is a
finite surjective morphism $X_l\to P$ over $K_l$, whose composite with the projection to $Z$ is $a_l$.
This proves (1).

For (2), let $\sigma$, $r$, $g$ and $\nu$ be as there. Lemma~\ref{cv-toroidal-level-change}, applied to
$a$ and $\nu_0$, gives an automorphism $\rho_{\mathcal M}$ of $\mathcal M$ with a morphism
$\mu:G_{\mathcal M}\to G_{\mathcal M}$ over it, and an automorphism $\sigma_0$ of $R_{n,N}$, with
$\rho_{\mathcal M}\circ a=a\circ r$ and
$s_{\mathcal M}\circ\rho_{\mathcal M}=\operatorname{Spec}\sigma_0\circ s_{\mathcal M}$.

Let $W'$ be the fibre product of $\rho_{\mathcal M}\circ\iota_W$ and $\iota_W$, with projections
$f_1,f_2:W'\to W$, so that $\rho_{\mathcal M}\circ\iota_W\circ f_1=\iota_W\circ f_2$. Both are open
immersions, and $j\circ f_1$ and $j\circ f_2$ are dominant, because $W$ is dense in $\mathcal M$,
$\rho_{\mathcal M}$ is an automorphism and $j$ is dominant. Let $G'$ be the pullback of $G_W$ along
$f_1$, which is proper over $W'$. Composing $G'\to G_W\to G$ gives a morphism $G'\to G$ over
$j\circ f_1$. Composing $G'\to G_W\to G_{\mathcal M}$ with $\mu$ gives a morphism over
$\rho_{\mathcal M}\circ\iota_W\circ f_1=\iota_W\circ f_2$; it factors through $G_W$ over $f_2$, and
composing with $G_W\to G$ gives a morphism $G'\to G$ over $j\circ f_2$. With
$s:Z\to\operatorname{Spec}R_{n,N}$ the structure morphism,
$s\circ j\circ f_2=\operatorname{Spec}\sigma_0\circ s\circ j\circ f_1$. Apply
Lemma~\ref{cv-toroidal-extend} with $R=R_{n,N}$, with $Z_1=Z$ over $R$ through $s$ and $Z_2=Z$ over $R$
through $\operatorname{Spec}\sigma_0^{-1}\circ s$, which is projective, with $G$ and its ample theta
presentation on both, and with $W'$, $G'$, $j_1=j\circ f_1$ and $j_2=j\circ f_2$, which are open
immersions with dense image. It gives $\tau:Z\to Z$ with $\tau\circ j\circ f_1=j\circ f_2$. With
$\sigma_0$ in place of $\sigma_0^{-1}$ and $f_1$ and $f_2$ exchanged, it gives $\tau'$ with
$\tau'\circ j\circ f_2=j\circ f_1$. The composites $\tau'\circ\tau$ and $\tau\circ\tau'$ are the
identity on the dense images of $j\circ f_1$ and $j\circ f_2$; as $Z$ is reduced and separated, $\tau$
is an automorphism with inverse $\tau'$.

It remains to show $\tau\circ a_l=a_l\circ r$. Let $D$ be the fibre product of
$\rho_{\mathcal M}\circ\iota_{W_0}$ and $\iota_{W_0}$, with projections $d_1,d_2:D\to W_0$, so that
$\rho_{\mathcal M}\circ\iota_{W_0}\circ d_1=\iota_{W_0}\circ d_2$. Through the open immersion
$\iota_U\circ i\circ d_1$, $D$ is open in $Z$, hence reduced. Let $Q$ be the fibre product of
$j\circ f_1$ and $\iota_U\circ i\circ d_1$, with projections $q_1:Q\to W'$ and $q_2:Q\to D$; then $q_2$
is dominant, as a base change of the dominant morphism $j\circ f_1$ along an open immersion. Since
$\iota_U$ is a monomorphism, $j'\circ f_1\circ q_1=i\circ d_1\circ q_2$. Hence
\[
 u\circ j'\circ f_2\circ q_1=\iota_W\circ f_2\circ q_1=\rho_{\mathcal M}\circ\iota_W\circ f_1\circ q_1
 =\rho_{\mathcal M}\circ u\circ j'\circ f_1\circ q_1=\rho_{\mathcal M}\circ u\circ i\circ d_1\circ q_2
 =\rho_{\mathcal M}\circ\iota_{W_0}\circ d_1\circ q_2=\iota_{W_0}\circ d_2\circ q_2,
\]
and $(*)$ gives $j'\circ f_2\circ q_1=i\circ d_2\circ q_2$. So
$\tau\circ\iota_U\circ i\circ d_1\circ q_2=\tau\circ j\circ f_1\circ q_1=j\circ f_2\circ q_1
=\iota_U\circ i\circ d_2\circ q_2$, and since $q_2$ is dominant, $D$ is reduced and $Z$ is separated,
$\tau\circ\iota_U\circ i\circ d_1=\iota_U\circ i\circ d_2$. Finally,
$\rho_{\mathcal M}\circ a=a\circ r$ says
$\rho_{\mathcal M}\circ\iota_{W_0}\circ b=\iota_{W_0}\circ b\circ r$, so $(b,b\circ r)$ is a morphism
$X_l\to D$, and
$\tau\circ a_l=\tau\circ\iota_U\circ i\circ d_1\circ(b,b\circ r)=\iota_U\circ i\circ d_2\circ(b,b\circ r)
=\iota_U\circ i\circ b\circ r=a_l\circ r$. This proves (2) with $\rho=\tau$.
\end{proof}
% Lean: topAmpExists_of_semistableReduction = T.exists_fullLevel_model_ample_frames =
% Models.topAmpExists: ZM.exists_isSympBasis, exists_ringGenerators, the three lemmas
% (exists_thetaFamily_pel_smooth_generic_fibre, nonempty_ampleFramedModel_PELThetaFamily,
% exists_universal_fullLevel_of_pelThetaFamily), exists_tBase_RZ and arBase_RZ, and
% T.exists_fullLevel_model_of_family_birationalModel = TopOfFamily.exists_top: C.good.1 (U),
% genericPoints_subset and range_j_subset, abelian_locus_open_family_over_level, exists_section (i and
% (*)), exists_inc, exists_two_auxiliary_inert_odd_primes_coprime_level,
% moduli_finite_surjective_over_generic_family, cancel_of_eq and comp_of_eq (the morphism over a_l),
% exists_finite_lift (fin), fullLevelChange_acts_on_unitaryFamily, exists_iso (with
% exists_hom_birational_models_ampleFrames twice and isProjOver_comp_iso), comp_eq_of_dense (sym). The
% Siegel family (T.exists_universal_quasiProjective_siegelFamily_large_level,
% T.cmFamily_finite_unramified_over_siegelFamily) is used only for the field sieg of Top, and
% eq_of_level only for compat; the lemma states neither.

\begin{lemma}\label{cv-model-toroidal-normalization}
Let $O$ be a Dedekind domain of characteristic zero with fraction field $F$, let $Y$ be a scheme
projective over $O$, let $W$ be an open subscheme of $Y_F=Y\otimes_OF$, and let $f:U\to W$ be a finite
morphism from a normal scheme $U$.
\begin{enumerate}
\item There are a normal scheme $Y'$, projective and flat over $O$, a finite morphism $\pi:Y'\to Y$
  over $O$, and an open immersion $j:U\to Y'_F$ over $F$ with dense image, such that $\pi\circ j=f$ and
  $j(U)=Y'_F\cap\pi^{-1}(W)$.
\item In Situation~\ref{cv-function-situation}, let $G$ be an $\mathcal O_E$-semi-abelian scheme over
  $Y$ whose abelian locus meets $Y_F$ in $W$, and let $A$ be an $\mathcal O_E$-semi-abelian scheme over
  $U$ with a morphism $A\to G$ over $f$. Then $(Y',\pi^*G,j)$ is a compactification of $(U,A)$ over $O$
  (Definition~\ref{cv-model-toroidal-semiabelian}). If $G$ is theta-split, or quasi-projective over
  $O$, so is $\pi^*G$.
\end{enumerate}
\end{lemma}
\leanref{Development/AbelianVarietyTheory/Proofs/ModuliProofs/ExistsNormalFiniteModelOfFiniteOverOpen.lean}{GlobalTheory.exists_normal_finite_model_of_finite_over_open}
\leanref{Development/AbelianVarietyTheory/Proofs/ModuliProofs/ExistsNormalModelAboveFiniteCover.lean}{GlobalTheory.Models.Cpt.exists_above}
\leanref{Development/AbelianVarietyTheory/Proofs/ModuliProofs/ExistsCompactModelAboveFullLevel.lean}{GlobalTheory.Models.Top.exists_cpt_above}
\begin{proof}
(1) Let $Y'$ be the normalization of $Y$ in $U$, the relative spectrum of the integral closure of
$\mathcal O_Y$ in the direct image of $\mathcal O_U$. A Dedekind domain of characteristic zero is
excellent, so $\pi:Y'\to Y$ is finite; hence $Y'$ is normal and projective over $O$. Every component
of $Y'$ is integral and is the closure of a component of $U$, so it dominates $\operatorname{Spec}O$,
and $Y'$ is flat over $O$. Since $f$ is finite and $U$ is normal, the normalization of $W$ in $U$ is
$U$, which gives $j$ with $j(U)=Y'_F\cap\pi^{-1}(W)$, dense in $Y'_F$.

(2) The pullback $\pi^*G$ is an $\mathcal O_E$-semi-abelian scheme, and the morphism $A\to G$ over
$f=\pi\circ j$ is a morphism $A\to\pi^*G$ over $j$. The abelian locus of $\pi^*G$ is the preimage of
that of $G$, so it meets $Y'_F$ in $\pi^{-1}(W)\cap Y'_F=j(U)$. The three conditions of being
theta-split are stable under pullback: the preimage of an open set is open and properness is stable
under base change, fibres are fibres of $G$, and sections pull back. Quasi-projectivity is preserved
because $\pi$ is finite.
\end{proof}
% Lean: exists_normal_finite_model_of_finite_over_open; the semi-abelian scheme is the pullback by sabPull
% (with sabCancel for the map from A), theta-split pulls back by good_pullback.

\begin{lemma}\label{cv-model-toroidal-quotient}
In Situation~\ref{cv-function-situation}, suppose $p+q\ge1$. Let $l_0=(M_0,\mathfrak M_0,1)$ and
$l_1=(M_0,\mathfrak M_0,N)$ be levels (Definition~\ref{cv-model-toroidal-semiabelian}), and let
$r:X_{l_1}\to X_{l_0}$ be a morphism over the inclusion $K_{l_0}\subset K_{l_1}$ with a morphism
$\nu:\mathcal A_{l_1}\to\mathcal A_{l_0}$ of $\mathcal O_E$-semi-abelian schemes over $r$ that carries
the level structure $\eta$ of $\mathcal A_{l_1}$ to that of $\mathcal A_{l_0}$. Then there is a finite
group $\Gamma$ that acts faithfully on $X_{l_1}$ by automorphisms over $X_{l_0}$, and on
$\mathcal A_{l_1}$ by automorphisms of $\mathcal O_E$-semi-abelian schemes over them that commute with
$\nu$, such that:
\begin{enumerate}
\item $r$ is affine, and for every open $V\subset X_{l_0}$ the functions on $V$ are the
  $\Gamma$-invariant functions on $r^{-1}(V)$;
\item every $\gamma\in\Gamma$ preserves $\eta$, and there are an $\mathcal O_E$-linear endomorphism $g$
  of $\Lambda$ and an automorphism $\sigma$ of $K_{l_1}$ with $\sigma(\sqrt{-d})=\sqrt{-d}$ such that
  $\gamma$ acts on $X_{l_1}$ over $\sigma$ and carries $\eta'(g(y))$ to $\eta'(y)$ for every
  $y\in\Lambda$.
\end{enumerate}
\end{lemma}
\leanref{Development/AbelianVarietyTheory/Proofs/ModuliProofs/ModuliLevelQuotientByFiniteGroup.lean}{GlobalTheory.LevQuot.exists_level_quotient}
\begin{proof}
The morphism $r$ is finite, \'etale and surjective, and it is the unique morphism over
$K_{l_0}\subset K_{l_1}$ that classifies $\mathcal A_{l_1}$ with its action, polarization and level
structure $\eta$. Since $\nu$ identifies $\mathcal A_{l_1}$ with $r^*\mathcal A_{l_0}$, every
automorphism of $X_{l_1}$ over $X_{l_0}$ lifts uniquely to an automorphism of $\mathcal A_{l_1}$ that
commutes with $\nu$. The automorphisms whose lifts change $\eta'$ by an $\mathcal O_E$-linear
endomorphism of $\Lambda$ and act on $K_{l_1}$ by an automorphism fixing $\sqrt{-d}$ form a finite
monoid, and $\Gamma$ is a group contained in it. Two geometric points of $X_{l_1}$ with the same image
in $X_{l_0}$ carry the same abelian variety with the same $\eta$, and differ in their full level
structures and in the chosen primitive $N$-th root of unity; so an element of $\Gamma$ carries one to
the other. Since $r$ is finite and \'etale and $\Gamma$ is transitive on its geometric fibres, the
invariant functions on $r^{-1}(V)$ are the functions on $V$. Every element of $\Gamma$ preserves $\eta$
because its lift commutes with $\nu$.
\end{proof}
% Lean: LevelMap.exists_levelMap (finite, etale, surjective), LevUniq.map_to_univ_unique, exists_liftAut,
% LevelCover.sab_cancel, twistMonoid with finite_twistMonoid and exists_mem_twistMonoid,
% exists_subgroup_of_finite, invariant_descends.

\begin{lemma}\label{cv-model-toroidal-stabilizer}
In Situation~\ref{cv-function-situation}, let $l=(M,\mathfrak M,N)$ be a level, let $S\ge1$ be an
integer with $M\mid S$ and $N\mid S^e$ for some $e$, and let $(\mathcal C,\mathcal G,j)$ be a
compactification of $(X_l,\mathcal A_l)$ over $R_{S,N}$
(Definition~\ref{cv-model-toroidal-semiabelian}). Let $\sigma$ be an automorphism of $K_l$ with
$\sigma(\sqrt{-d})=\sqrt{-d}$, let $r$ be an automorphism of $X_l$ over $\sigma$ with an automorphism
$\nu$ of the $\mathcal O_E$-semi-abelian scheme $\mathcal A_l$ over $r$, and let $\rho$ be an
automorphism of $\mathcal C$ with an automorphism $\mu$ of the $\mathcal O_E$-semi-abelian scheme
$\mathcal G$ over $\rho$, such that $\rho\circ j=j\circ r$ and $\mu$ extends $\nu$. Let $x$ be a
geometric point of $\mathcal C$ with $\rho(x)=x$.
\begin{enumerate}
\item The automorphism $\sigma$ is the identity.
\item If $\nu$ preserves the level structure $\eta$ and $\mu$ has finite order, then $\mu$ acts
  trivially on the fibre $\mathcal G_x$.
\item If $\nu$ preserves $\eta$, $\mu$ acts trivially on $\mathcal G_x$, and
  $\nu\circ\eta'(g(y))=\eta'(y)\circ r$ for all $y\in\Lambda$ and an $\mathcal O_E$-linear
  endomorphism $g$ of $\Lambda$, then $r^N$ is the identity.
\end{enumerate}
\end{lemma}
\leanref{Development/AbelianVarietyTheory/Proofs/ModuliProofs/IsDominantPullbackFstOfFlatOfInjective.lean}{GlobalTheory.galoisPart_eq_refl_of_fixes_point}
\leanref{Development/AbelianVarietyTheory/Proofs/ModuliProofs/FiniteOrderAutomorphismFixingLevelTrivialOnFixed.lean}{GlobalTheory.fibre_fixed_of_finite_order_of_level_fixed}
\leanref{Development/AbelianVarietyTheory/Proofs/ModuliProofs/IterateEqIdOfTrivialOnFixedFibre.lean}{GlobalTheory.StabTame.iterate_eq_id_of_fixes_fibre}
\begin{proof}
(1) Let $k$ be the field of $x$. Since $N$ is invertible on $R_{S,N}$, the $N$-th roots of unity of
$R_{S,N}$ map injectively to $k$. The automorphism $\rho$ lies over the automorphism of $R_{S,N}$
induced by $\sigma$, and it fixes $x$, so $\sigma$ fixes the image of $\zeta_N$ in $k$; hence
$\sigma(\zeta_N)=\zeta_N$, and $\sigma$ is the identity on $K_l=\mathbf Q(\sqrt{-d},\zeta_N)$.

(2) By (1), $r$ lies over the identity of $K_l$. The group scheme $\mathcal G[M]$ is quasi-finite and
\'etale over the normal scheme $\mathcal C$, since $M$ is invertible. After adjoining a primitive
$M$-th root of unity, there are finitely many $M$-torsion sections of $\mathcal A_l$, determined by
$\eta$ and the root of unity, whose values are all the $M$-torsion points at every geometric point; they
are fixed by $\nu$, which preserves $\eta$ and lies over the identity. By normality, $\mu$ then fixes
every $M$-torsion point of $\mathcal G_x$. So $\mu$ acts on the semi-abelian variety $\mathcal G_x$ by
an automorphism of finite order that fixes its $M$-torsion, with $M$ odd, at least $3$ and invertible
in $k$. Such an automorphism is the identity, by Serre's lemma \cite{Ser61}, whose argument applies to
semi-abelian varieties.

(3) Let $P\subset\Lambda$ be the set of $y$ for which the section $\eta'(y)$ extends to a section of
$\mathcal G$ over a neighbourhood of $x$. The group scheme $\mathcal G[N]$ is \'etale over the normal
scheme $\mathcal C$, so each of its components through a point over $x$ is a section over an open
neighbourhood of $x$; hence taking values at $x$ is a bijection from $P/N\Lambda$ onto the $N$-torsion
of $\mathcal G_x$, which has $N^{2h-t}$ elements, where $h=\dim\mathcal G_x$ and $t$ is the dimension of
its torus. For $y\in P$, the sections $\eta'(g(y))$ and $\eta'(y)$ have the same value at $x$, since
$\mu$ acts trivially on $\mathcal G_x$; so $g(y)\equiv y$ modulo $N\Lambda$. The automorphism $\nu$
preserves the Weil pairing, so $g$ preserves the alternating form of $\Lambda$ modulo $N$. The elements
of $P$ whose values at $x$ lie in the torus, $N^t$ of them modulo $N\Lambda$, are orthogonal to $P$
modulo $N$; since the form is perfect, the elements orthogonal to $P$ modulo $N$ are also $N^t$ in
number modulo $N\Lambda$, so they all lie in $P$. For $y\in\Lambda$ and $w\in P$, the pairing of
$g(y)-y$ with $w$ is that of $g(y)$ with $w-g(w)$ modulo $N$, which is divisible by $N$; so
$g(y)-y\in P$, and $g^N(y)-y\equiv N(g(y)-y)\equiv0$ modulo $N\Lambda$. Hence $r^N$, with $\nu^N$,
preserves $\eta$, the full level structure $\eta'$ and the structure map to $K_l$, and it is the
identity because $X_l$ represents its moduli problem.
\end{proof}
% Lean: (1) galoisPart_eq_refl_of_fixes_point; (2) fibre_fixed_of_finite_order_of_level_fixed via
% torsion_fixed (torsion_point_fixed_of_generic_sections_fixed) and fibre_fixed
% (serre_lemma_semiAbelian); (3) iterate_eq_id_of_fixes_fibre: modN_of_extends, isometry_modN,
% T.levelLabel_extends_of_orthogonal_to_extended (exists_toric_torsion_cubWeil_one and a count),
% iterate_modN, exists_sabMap_iterate, LiftAut.eq_id.

\begin{lemma}\label{cv-model-toroidal-pushforward}
In Situation~\ref{cv-function-situation}, let $O$ be a Dedekind domain of characteristic zero with
fraction field $F$, let $l$ be a level, let $U$ be a scheme over $F$, and let $A$ be an abelian scheme
over $U$ with the structures of the moduli problem at level $l$
(Definition~\ref{cv-model-toroidal-semiabelian}), viewed as an $\mathcal O_E$-semi-abelian scheme. Let
$(\mathcal C,\mathcal G,j)$ be a compactification of $(U,A)$ over $O$ that satisfies condition (c) of
Definition~\ref{cv-model-toroidal-semiabelian}(3). Then every point of $\mathcal C$ has an open
neighbourhood $V$ with a section $\theta$ of the line bundle $\mathcal L$ of $\mathcal G$ over
$\mathcal G_V$ that vanishes identically on no fibre, such that for every open $W\subset V$ each section
of $\mathcal L$ over $\mathcal G_W$ is $f\theta$ for a unique function $f$ on $W$.
\end{lemma}
\leanref{Development/AbelianVarietyTheory/Proofs/ModuliProofs/PushforwardPolarisationBundleInvertibleOnModel.lean}{GlobalTheory.PushInv.pushInvertible_of_model}
\leanref{Development/AbelianVarietyTheory/Proofs/ModuliProofs/PushInvertibleOfDensePart.lean}{GlobalTheory.PushInv.pushInvertible_of_densePart}
\begin{proof}
Let $V$ and $\theta$ be given by condition (c), and let $s$ be a section of $\mathcal L$ over
$\mathcal G_W$. Over $j(U)$, $\mathcal L$ is the line bundle of the principal polarization of $A$, so
its direct image is invertible and $\theta$ generates it; thus $s=f\theta$ over $W\cap j(U)$ for a
unique function $f$. The scheme $\mathcal C$ is normal, so $f$ extends to $W$ once it is regular at the
points $y$ of codimension one of $W$. The fibre $\mathcal G_y$ is integral, since $\mathcal G$ is smooth
with geometrically connected fibres, and $\theta$ does not vanish identically on it; so $\theta$
generates $\mathcal L$ at the generic point of $\mathcal G_y$, and $f=s/\theta$ is regular there. The
local ring of $\mathcal G$ at that generic point is faithfully flat over the local ring of $\mathcal C$
at $y$, so $f$ is regular at $y$.
\end{proof}
% Lean: pushInvertible_of_model from pushInvertible_of_densePart (normal base, smooth geometrically
% integral G) and pushInvertible_any (principally polarized abelian schemes, exists_basisOn_any).

\begin{lemma}\label{cv-model-toroidal-descent}
In Situation~\ref{cv-function-situation}, let $O'\to O$ be a finite injective homomorphism of Dedekind
domains of characteristic zero, with fraction fields $F'\to F$. Let $U$ be a scheme over $F$ with an
$\mathcal O_E$-semi-abelian scheme $A$ proper over $U$, let $U'$ be a scheme over $F'$ with an
$\mathcal O_E$-semi-abelian scheme $A'$, and let $r:U\to U'$ be a morphism over $F'\to F$ with a
morphism $\nu:A\to A'$ of $\mathcal O_E$-semi-abelian schemes over $r$. Let a finite group $\Gamma$ act
on $U$, and on $A$ by automorphisms over those of $U$ that commute with $\nu$, such that $r$ is affine
and for every open $V\subset U'$ the functions on $V$ are the $\Gamma$-invariant functions on
$r^{-1}(V)$. Let $(\mathcal C,\mathcal G,j)$ be a compactification of $(U,A)$ over $O$
(Definition~\ref{cv-model-toroidal-semiabelian}), with $\mathcal G$ theta-split and quasi-projective
over $O$, which satisfies the conclusion of Lemma~\ref{cv-model-toroidal-pushforward}. Let $\Gamma$ act
on $\mathcal C$, and on $\mathcal G$ by automorphisms of $\mathcal O_E$-semi-abelian schemes over those
of $\mathcal C$, compatibly with $j$ and with its action on $A$. Suppose that, for every geometric point
$x$ of $\mathcal C$, the stabilizer of $x$ in $\Gamma$ has order invertible in the field of $x$ and acts
trivially on the fibre $\mathcal G_x$. Then $(U',A')$ has a compactification
$(\mathcal C',\mathcal G',j')$ over $O'$ with $\mathcal G'$ theta-split, together with a morphism
$\pi:\mathcal C\to\mathcal C'$ over $O'\to O$, finite as a morphism to $\mathcal C'\otimes_{O'}O$, and
a morphism $\mathcal G\to\mathcal G'$ of $\mathcal O_E$-semi-abelian schemes over $\pi$, which extend
$r$ and $\nu$.
\end{lemma}
\leanref{Development/AbelianVarietyTheory/Proofs/ModuliProofs/ExistsQuotientModelByTameFiniteGroup.lean}{GlobalTheory.Models.exists_quotient_model}
\leanref{Development/AbelianVarietyTheory/Proofs/ModuliProofs/ExistsSemiAbelianDescentTameQuotient.lean}{GlobalTheory.exists_sab_descent_of_tame_quotient}
\begin{proof}
Let $\mathcal C'$ be the quotient of $\mathcal C$ by $\Gamma$, which exists because $\mathcal C$ is
projective over $O$ \cite[Expos\'e~V, Proposition~1.8]{Gro71}, with $\pi:\mathcal C\to\mathcal C'$. It
is normal, projective and flat over $O'$, and $\pi$ is finite. Since $U'$ is the quotient of $U$ by
$\Gamma$, the open immersion $j$ descends to an open immersion $j':U'\to\mathcal C'_{F'}$ with dense
image, and $\pi^{-1}(j'(U'))=j(U)$.

The quotient $\mathcal G'=\mathcal G/\Gamma$ exists because $\mathcal G$ is quasi-projective, and it is
a scheme over $\mathcal C'$. The morphism $\mathcal G\to\mathcal G'\times_{\mathcal C'}\mathcal C$ is
an isomorphism: at a geometric point $x$ of $\mathcal C$, the stabilizer of $x$ acts trivially on
$\mathcal G_x$ and has invertible order, so forming invariants commutes with passing to the fibre. Hence
$\mathcal G'$ is smooth over $\mathcal C'$ with the fibres of $\mathcal G$, a semi-abelian scheme, and
the group law and the action of $\mathcal O_E$ descend. The line bundle $\mathcal L$, rigidified along
the unit section, carries an action of $\Gamma$ over that on $\mathcal G$; since the stabilizers act
trivially on its fibres at the points of $\mathcal G$ over $x$ and have invertible order, it descends
to a line bundle on $\mathcal G'$, with its cubical structure and its symmetry.

The three conditions of being theta-split descend: the abelian locus of $\mathcal G'$ is the image of
that of $\mathcal G$, which is open and stable under $\Gamma$, and $\mathcal G'$ is proper over it;
the fibre of $\mathcal G'$ at a point of $\mathcal C'$ becomes the fibre of $\mathcal G$ at a point
over it after a residue field extension that is separable, because the stabilizers have order prime to
the residue characteristic; and averaging over $\Gamma$ the generators $\theta$ of the direct image of
$\mathcal L$ gives invariant sections, one of which vanishes identically on no fibre near a given
point, since the order of its stabilizer is invertible there.
\end{proof}
% Lean: exists_quotient_model with exists_quotient_model_scheme (the quotient C' with pi, j') and
% exists_sab_descent_of_tame_quotient: quotient_flat_baseChange, quotient_descend_grpObj,
% SemiAbelianDescent.main, RigidLinearisation.exists_linearisation, LinDescent.exists_descent_of_tame,
% T.exists_semiAbelianTriple_of_finiteQuotient; theta-split: splitFibres_descend_of_tame_quotient,
% exists_secOn_nonvanishing_of_tame_quotient.

\begin{lemma}\label{cv-model-toroidal-base}
In Situation~\ref{cv-function-situation}, suppose $p+q\ge1$. Let $n$ be an auxiliary prime
(Lemma~\ref{cv-model-toroidal-auxiliary}) and let $l_0=(M_0,\mathfrak M_0,1)$ be a level of full level
$1$ (Definition~\ref{cv-model-toroidal-semiabelian}). Then $(X_{l_0},\mathcal A_{l_0})$ has a
compactification over $R_{nM_0,1}=\mathbf Z[(1+\sqrt{-d})/2,1/(nM_0)]$ whose
$\mathcal O_E$-semi-abelian scheme is theta-split.
\end{lemma}
\leanref{Development/AbelianVarietyTheory/Proofs/ModuliProofs/ExistsBaseModelSemiAbelian.lean}{GlobalTheory.BaseModelEx.exists_baseModel}
\leanref{Development/AbelianVarietyTheory/Definitions/ModuliSpaces/AuxiliaryInertOddPrime.lean}{GlobalTheory.Models.BaseModel}
\begin{proof}
Let $J$, $N=n^J$, $Z$ and $G$ be given by Lemma~\ref{cv-model-toroidal-top}, let
$l_1=(M_0,\mathfrak M_0,N)$, and let $R_0=R_{nM_0,1}$ and $R_1=R_{nM_0,N}$. These are Dedekind domains
with fraction fields $K_{l_0}=\mathbf Q(\sqrt{-d})$ and $K_{l_1}=\mathbf Q(\sqrt{-d},\zeta_N)$, the
ring $R_1$ is finite over $R_0$, and $nM_0$ is invertible on both.

By Lemma~\ref{cv-model-toroidal-top}(1), $X_{l_1}$ maps finitely onto the open subscheme of
$Z\otimes_{R_{n,N}}K_{l_1}$ over which $G$ is abelian. Lemma~\ref{cv-model-toroidal-normalization},
applied to $Z\otimes_{R_{n,N}}R_1$, gives a compactification $(\mathcal C_1,\mathcal G_1,j_1)$ of
$(X_{l_1},\mathcal A_{l_1})$ over $R_1$ with a finite morphism $q:\mathcal C_1\to
Z\otimes_{R_{n,N}}R_1$ and $\mathcal G_1=q^*G$, theta-split and quasi-projective over $R_1$.

Let $r:X_{l_1}\to X_{l_0}$ be the morphism that forgets the full level structure, with
$\nu:\mathcal A_{l_1}\to\mathcal A_{l_0}$ over it; it is finite, \'etale and surjective, and
Lemma~\ref{cv-model-toroidal-quotient} gives a finite group $\Gamma$. Let $\gamma\in\Gamma$, acting by
$r_\gamma$ on $X_{l_1}$ and by $\nu_\gamma$ on $\mathcal A_{l_1}$, with $g$ and $\sigma$ as in
Lemma~\ref{cv-model-toroidal-quotient}(2). Lemma~\ref{cv-model-toroidal-top}(2) gives an automorphism
of $Z$ compatible with $r_\gamma$, which preserves the open subscheme $Z\otimes_{R_{n,N}}R_1$ on which
$M_0$ is invertible. Since $\mathcal C_1$ is
the normalization of $Z\otimes_{R_{n,N}}R_1$ in $X_{l_1}$, $r_\gamma$ extends to an automorphism
$\rho_\gamma$ of $\mathcal C_1$, and by Theorem~\ref{cv-model-toroidal-extension}, applied at the points
of codimension one of the normal scheme $\mathcal C_1$, $\nu_\gamma$ extends to an automorphism
$\mu_\gamma$ of $\mathcal G_1$ over $\rho_\gamma$. This is an action of $\Gamma$ on
$(\mathcal C_1,\mathcal G_1)$ compatible with $j_1$.

Let $x$ be a geometric point of $\mathcal C_1$ and $\gamma$ an element of its stabilizer. Apply
Lemma~\ref{cv-model-toroidal-stabilizer} with $l=l_1$ and $S=nM_0$, so that $M_0\mid S$ and $N\mid S^J$.
By part (1), $\sigma$ is the identity. By part (2), since $\nu_\gamma$ preserves $\eta$ and
$\mu_\gamma$ has finite order, $\gamma$ acts trivially on $\mathcal G_{1,x}$. By part (3),
$r_\gamma^N$ is the identity, so $\gamma^N=1$ because $\Gamma$ acts faithfully. So the stabilizer of
$x$ is a group of order a power of $n$, which is invertible in the field of $x$.

By Lemma~\ref{cv-model-toroidal-pushforward}, applied to $X_{l_1}$ with $\mathcal A_{l_1}$ and to the
theta-split $\mathcal G_1$, the compactification $\mathcal C_1$ satisfies the conclusion of that
lemma. Lemma~\ref{cv-model-toroidal-descent}, applied to $R_0\subset R_1$, $r$, $\nu$ and the action
of $\Gamma$, gives the compactification of $(X_{l_0},\mathcal A_{l_0})$ over $R_0$ with a theta-split
$\mathcal O_E$-semi-abelian scheme.
\end{proof}
% Lean: exists_baseModel: isDedFrac_RZ, cmModelCoefficient_inclusions_etale_openImmersion,
% Top.exists_cpt_above (C_1), exists_finite_etale_surjective_levelMap, exists_level_quotient,
% exists_action_model (with T.sym, exists_restrict, Models.exists_sabMap_extending),
% galoisPart_eq_refl_of_fixes_point, finiteOrderAutomorphism_fixing_level_trivial_on_fixedFibre,
% iterate_eq_id_of_trivial_on_fixed_fibre, card_stab_ne_zero, good_pullback,
% pushforward_polarisationBundle_invertible_on_model, exists_quotient_model_by_tame_finite_group.

\begin{theorem}[Faltings--Chai \cite{FC90}, Lan \cite{Lan13}]\label{cv-model-toroidal-degeneration}
In Situation~\ref{cv-function-situation}, let $R$ be a complete noetherian normal local domain of
characteristic zero, and let $G$ be an $\mathcal O_E$-semi-abelian scheme over $R$
(Definition~\ref{cv-model-toroidal-semiabelian}) with line bundle $\mathcal L$. Suppose that the
generic fibre of $G$ is an abelian variety on which $\mathcal L$ is ample and defines a principal
polarization, that the torus of $G$ over the closed point is split, and that $\mathcal L$ has a section
that vanishes identically on no fibre of $G$. Then:
\begin{enumerate}
\item $G$, with its action of $\mathcal O_E$, its line bundle and the polarization of its generic fibre,
  is given by Mumford's construction from degeneration data: a split extension of an abelian scheme by
  a torus over $R$, a group of periods, and trivializations $\psi$ of the associated line bundles over
  the fraction field of $R$;
\item for every algebraically closed field $k$ containing the fraction field of $R$ and every
  $m\ge1$, the $m$-torsion points of the generic fibre with values in $k$ are described by the
  degeneration data, compatibly with addition, the action of $\mathcal O_E$ and the Weil pairing;
\item if the abelian locus of $G$ is open, the trivializations $\psi$ are units over it.
\end{enumerate}
\end{theorem}
\leanref{Development/AbelianVarietyTheory/Proofs/DegenerationProofs/ExistsUniformisationCompleteLocalNormalPrime.lean}{GlobalTheory.LN.UnifPrime.exists_pelUnif}
\leanref{Development/AbelianVarietyTheory/Proofs/DegenerationProofs/PelTrivialisationIsUnitOnAbelianLocus.lean}{GlobalTheory.psiUnitOn_abLocus_pel}

This is the theory of degenerations of polarized abelian varieties over complete normal bases, with
Mumford's construction \cite{Mum72} as its inverse \cite[Chapters~II and~III]{FC90}, in the form with
an action of $\mathcal O_E$ and with level structures \cite[Chapters~4 and~5]{Lan13}.

\begin{theorem}\label{cv-model-toroidal}
In Situation~\ref{cv-function-situation}, suppose $p+q\ge1$. Let $\ell$ be a prime for which the
prime-to-$\ell$ part of the level is at least $3$, let $\iota:F_c\to\mathbf C_\ell$ be an embedding,
and let $O=\{b\in F_c:|\iota(b)|\le1\}$, a discrete valuation ring with fraction field $F_c$. Then
there is a model $(\mathcal X,\mathcal A,j)$ of $X$ over $O$ (Definition~\ref{cv-model}) such that:
\begin{enumerate}
\item at every point $x$ of $\mathcal X$, the fibre $\mathcal A_x$ becomes, over a finite separable
  extension of the residue field of $x$, an extension of an abelian variety by a split torus;
\item for every complete noetherian normal local domain $R$ of characteristic zero and every
  morphism $\operatorname{Spec}R\to\mathcal X$ that maps some point into $j(X)$ and over whose closed
  point the torus of $\mathcal A$ is split, the pullback of $\mathcal A$ to $R$, with its action of
  $\mathcal O_E$, its polarization and its level structure over the points that map into $j(X)$, is
  given by Mumford's construction from a split extension of an abelian scheme by a torus over $R$ and
  a group of periods.
\end{enumerate}
\end{theorem}
\leanref{Development/AutomorphicSpectrum/Proofs/GeometryProofs/GoodReductionDvrModelExistsStatement.lean}{GlobalTheory.T.good_reduction_dvr_model_exists}
\leanref{Development/WeilBound/CoreStatements/LocalModelCharts/ExistsLocalModelDvr.lean}{PositivityTheory.Ar2.exists_local_model_dvr}
\leanref{ReciprocityTheory/ReductiveGroups/ArithmeticGeometry/BirationalGeometry/ToroidalCompactification/UniformisationConclusion.lean}{ReciprocityTheory.ChartRestriction.G.Geo.Tor.Model.SplitFibres}
\leanref{ReciprocityTheory/ReductiveGroups/ArithmeticGeometry/BirationalGeometry/ToroidalCompactification/UniformisationConclusion.lean}{ReciprocityTheory.ChartRestriction.G.Geo.Tor.Model.N3}

\begin{proof}
Write the level of $X$ as $l=(M,\mathfrak M,N)$ (Definition~\ref{cv-model-toroidal-semiabelian}), so
that $X$ is $X_l$, over $K_l\subset F_c$, and $A$ is $\mathcal A_l$. Let $M_0$ be the prime-to-$\ell$
part of $M$, which is at least $3$ by hypothesis, and let $l_0=(M_0,\mathfrak M+M_0\mathcal O_E,1)$,
the prime-to-$\ell$ part of the level, which is a level. By
Lemma~\ref{cv-model-toroidal-auxiliary} there are two auxiliary primes; let $n$ be one that is
different from $\ell$. By Lemma~\ref{cv-model-toroidal-base}, $(X_{l_0},\mathcal A_{l_0})$ has a
compactification $(\mathcal C_0,\mathcal G_0,j_0)$ over $R_0=R_{nM_0,1}$ with $\mathcal G_0$
theta-split.

The ring $O$ is a discrete valuation ring with fraction field $F_c$. Since $\ell$ divides neither $n$
nor $M_0$, it contains $R_0$, compatibly with $K_{l_0}=\mathbf Q(\sqrt{-d})\subset F_c$. The level map
$X\otimes_{K_l}F_c\to X_{l_0}$, under which $M/M_0$ times the level structure $\eta$ of $A$
is that of $\mathcal A_{l_0}$, comes with a morphism from $A$ to $\mathcal A_{l_0}$ of
$\mathcal O_E$-semi-abelian schemes over it, and the induced morphism
$X\otimes_{K_l}F_c\to X_{l_0}\otimes_{K_{l_0}}F_c$ is finite. The scheme $X\otimes_{K_l}F_c$ is
smooth over $F_c$, hence normal. Lemma~\ref{cv-model-toroidal-normalization}, applied to
$\mathcal C_0\otimes_{R_0}O$, to the open subscheme $j_0(X_{l_0})\otimes_{K_{l_0}}F_c$ of its generic
fibre and to $\mathcal G_0\otimes_{R_0}O$, gives a compactification $(\mathcal X,\mathcal A,j)$ of
$(X\otimes_{K_l}F_c,A)$ over $O$ with $\mathcal A$ theta-split. It is a model of $X$ over $O$
(Definition~\ref{cv-model}).

Property (1) is condition (b) of Definition~\ref{cv-model-toroidal-semiabelian}(3).

For (2), let $R$ and $\xi:\operatorname{Spec}R\to\mathcal X$ be as in (2), and let $G=\xi^*\mathcal A$,
an $\mathcal O_E$-semi-abelian scheme over $R$. The abelian locus of $\mathcal A$ is open, and its
preimage under $\xi$ is the abelian locus of $G$; it is open and nonempty, so it contains the generic
point of $\operatorname{Spec}R$. Since $R$ has characteristic zero, the generic point maps to
$\mathcal X_{F_c}$, where the abelian locus is $j(X)$. So the generic fibre of $G$ is the pullback of
$A$ by a point of $X$, an abelian variety on which the line bundle of $G$ is ample and
defines a principal polarization. By Lemma~\ref{cv-model-toroidal-pushforward}, applied to
$(\mathcal X,\mathcal A,j)$, the image of the closed point of $\operatorname{Spec}R$ has a neighbourhood
$V$ with a section of the line bundle of $\mathcal A$ over $\mathcal A_V$ that vanishes identically on
no fibre; $\xi$ maps $\operatorname{Spec}R$ into $V$, and the pullback of this section to $G$ vanishes
identically on no fibre of $G$. The torus of $G$ over the closed point is split by hypothesis. Then
Theorem~\ref{cv-model-toroidal-degeneration} applies to $G$: it is given by Mumford's construction from
a split extension of an abelian scheme by a torus over $R$ and a group of periods, its torsion points,
and with them the level structure at the generic point, are described by the degeneration data, and
the trivializations of the data are units on the preimage of the abelian locus of $\mathcal A$. This is
(2).
\end{proof}
% Lean: good_reduction_dvr_model_exists = PositivityTheory.Ar2.exists_local_model_dvr takes the model
% Y p W hp hW i of the system of Ar.Top.partI_holds = partI_of_semistableReduction
% (partI_of_inputs, ModelsSt from modelsSt_of_semistableReduction = Models.exists_localSystems_ample).
% The model is built in SystemAmple.exists_system_ampFrames: exists_two_auxiliary_inert_odd_primes_coprime_level
% (n /= p), topAmpExists (Top), exists_base_model_semiAbelian (BaseModel for the prime-to-p part l_0,
% primeToPart_spec), LocalModelExists.exists_model_raises (LevelMapField.exists_levelMap and
% exists_normal_model_above_finite_cover). (1): System.levelStructureConditions_and_splitFibres
% (good_pullback). (2): model_uniformisation_of_abelianLocus_directImage = model_uniformisation_of_clauses
% with LN.generic_point_of_meets_abLocus, ThetaLocal.thetaNZ_of_pushInvertible (from
% pushforward_polarisationBundle_invertible_on_model), isAdmissibleCompleteLocalDomainPrime_maximalIdeal,
% LN.UnifPrime.exists_pelUnif and psiUnitOn_abLocus_pel.

\begin{remark}\label{cv-model-literature}
Toroidal compactifications give models of the same kind in many cases. For the Siegel moduli schemes
of principally polarized abelian varieties with a level $n\ge3$ prime to $\ell$, Faltings and Chai
construct projective toroidal compactifications over $\mathbf Z[\zeta_n,1/n]$ that carry a
semi-abelian scheme extending the universal one \cite[Chapters~IV and~V]{FC90}; where the PEL datum
has good reduction at $\ell$, Lan constructs toroidal compactifications of its moduli schemes with the
additional structures \cite{Lan13}. Property (1) holds for every semi-abelian scheme, since a torus
over a field splits over a finite separable extension. Property (2) is the theory of degenerations of
Faltings and Chai, by which a polarized semi-abelian scheme over a complete noetherian normal local
domain, abelian at the generic point and with split torus over the closed point, comes from Mumford's
construction \cite{Mum72}, \cite[Chapters~II and~III]{FC90}. These sources do not state the theorem
in its generality, for every level whose prime-to-$\ell$ part is at least $3$ and with a symmetric
cubical line bundle on the semi-abelian scheme over the whole model.
\end{remark}

\subsection{The frame function of the periods}\label{subsection-cv-period-frame-function}

\begin{lemma}\label{cv-model-eigenbundles}
In Situation~\ref{cv-function-situation}, let $(\mathcal X,\mathcal A,j)$ be a model of $X$ over a
discrete valuation ring $O\subset F_c$ with fraction field $F_c$ (Definition~\ref{cv-model}). Then
$Z$ is open and closed in $X$, and the closure $\bar Z$ of $j(Z)$ in $\mathcal X_{F_c}$ is a normal
integral scheme with function field $F_c(Z)$. The subsheaves $\operatorname{Lie}_\tau(\mathcal A)$
and $\operatorname{Lie}_{\bar\tau}(\mathcal A)$ of the Lie algebra of $\mathcal A$ on $\bar Z$, on
which each $b\in\mathcal O_E$ acts by $\tau(b)$ and by $\bar\tau(b)$, are vector bundles of ranks $p$
and $q$, with fibres $\operatorname{Lie}_\tau(y)$ and $\operatorname{Lie}_{\bar\tau}(y)$ at the
points $y$ of $Z$.
\end{lemma}
\leanref{Development/AutomorphicSpectrum/Proofs/GeometryProofs/IrreducibleComponentClopenStatement.lean}{GlobalTheory.T.irreducible_component_clopen}
\leanref{Development/AutomorphicSpectrum/Proofs/GeometryProofs/EigenLatticeVectorBundleStatement.lean}{GlobalTheory.T.eigen_lattice_vector_bundle}
\begin{proof}
A connected component of the noetherian scheme $X$ is open and closed. The scheme $\mathcal X$ is
flat over $O$, so its generic fibre is dense in it, and $j(X)$ is dense in $\mathcal X_{F_c}$. So the
generic point of $j(Z)$ is a generic point of the normal scheme $\mathcal X$, and its closures in
$\mathcal X$ and in $\mathcal X_{F_c}$ are normal and integral, with function field $F_c(Z)$.

The group scheme $\mathcal A$ is smooth, so its Lie algebra is a vector bundle on $\mathcal X$, with
an action of $\mathcal O_E$. On $\mathcal X_{F_c}$ the discriminant of $E$ is invertible and every
embedding of $E$ takes values in $F_c$, so $\mathcal O_E\otimes\mathcal O_{\bar Z}$ is the product of
copies of $\mathcal O_{\bar Z}$, one for each embedding. The Lie algebra on $\bar Z$ is therefore the
direct sum of its eigensubsheaves, one for each embedding. Each is a direct summand of a vector
bundle, hence a vector bundle, and its rank is the dimension of its fibre at a point of $j(Z)$, which
is $p$ for $\tau$ and $q$ for $\bar\tau$.
\end{proof}
% Lean: cptOver_any (nonempty_cptOver_of_fraction) and isVBOn_eig_genFib (lat_isVB_of_smooth,
% lat_stable, EigSplit.isVBOn_eigSp with hdisc, finrank_eigSp); normality is
% GlobalTheory.CompanionTheory.Ar.XR.normalCpt_of_isCptGen.

\begin{lemma}\label{cv-periods-function}
In Situation~\ref{cv-function-situation}, let $\sigma$ be a complex embedding of $F_c$ with
$\sigma(\tau(\sqrt{-d}))=\theta$, and put $\tau_\sigma=\sigma\circ\tau$. Suppose that the
uniformization of $X$ at $\sigma$ presents each complex point of $X$ over $\sigma$ as $y_P$ for a
period $P$, of one type $(\sigma_k)_{k\le g}$ of signature $(p,q)$ for $\theta$, of one of finitely
many finitely generated projective $\mathcal O_E$-modules $L_1,\dots,L_n$ of rank $p+q$ with perfect
alternating forms (Definition~\ref{cv-cm-point-periods}), the index being determined by the point,
with $A_{y_P}=\mathbf C^g/P(L_i)$ compatibly with the actions of $\mathcal O_E$ and the polarizations.
For each $i$ fix a family $e^i=(e^i_1,\dots,e^i_{p+q})$ in $L_i$, linearly independent over
$\mathcal O_E$, and put $S_i=\tau_\sigma(G(e^i))$, with $\tau_\sigma$ applied to each entry of the
Gram matrix.

Let $(y,v,w)$ be a complex framed point of $X$ over $\sigma$ with $y=y_P$ for a period $P$ of $L_i$;
then $\operatorname{Lie}_\tau(y)$ and $\operatorname{Lie}_{\bar\tau}(y)$ consist of the vectors of
$\mathbf C^g$ supported on the coordinates $k$ with $\sigma_k=\tau_\sigma$ and with
$\sigma_k=\bar\tau_\sigma$. Let $C$ be a $p\times(p+q)$ matrix with
$v_a=\sum_mC_{am}P(e^i_m)_{\tau_\sigma}$ for every $a$, let $B$ be the $q\times(p+q)$ matrix with
$P(e^i_n)_{\bar\tau_\sigma}=\sum_jB_{jn}w_j$ for every $n$, and let $F=BS_i^{-1}$. Then
$\det S_i\ne0$, the determinant $\delta(y,v,w)$ of the square matrix whose rows are those of $C$
followed by those of $F$ is nonzero and depends only on $(y,v,w)$, and the function
\[
 \Phi_\sigma(y,v,w)=\big((2\pi)^{p-q}|\det S_i|\big)^{-1/2}\,\delta(y,v,w)^{-1}
\]
on the complex framed points of $X$ over $\sigma$ has the following properties.
\begin{enumerate}
\item It satisfies the transformation rule of a frame function (Definition~\ref{cv-function-framed}).
\item $|\Phi_\sigma(y,v,w)|^2\det\Gamma(v)=\det\Delta(w)$, where $\Gamma(v)$ and $\Delta(w)$ are the
  Hodge Gram matrices of $v$ and $w$ (Definition~\ref{cv-cm-point-periods}).
\item It is holomorphic: if complex framed points $(y(b),v(b),w(b))$ over $\sigma$ depend
  holomorphically on $b$ in an open subset of some $\mathbf C^m$, then
  $\Phi_\sigma(y(b),v(b),w(b))$ is a holomorphic function of $b$.
\end{enumerate}
\end{lemma}
\leanref{Development/AutomorphicSpectrum/Proofs/GeometryProofs/ExistsUnitAnalyticFunctionStatement.lean}{GlobalTheory.T.exists_unit_analytic_function}
\leanref{Development/AutomorphicSpectrum/Definitions/GlobalGeometry/IsFlatAnIffIsConstMul.lean}{GlobalTheory.CompanionTheory.Ar.Flat.IsUnitAn}
\leanref{Development/AutomorphicSpectrum/Definitions/GlobalGeometry/IsFlatAnIffIsConstMul.lean}{GlobalTheory.CompanionTheory.Ar.Flat.IsFlatAn}
\begin{proof}
The map $P$ extends to an $\mathbf R$-linear isomorphism $L_i\otimes\mathbf R\to\mathbf C^g$, and
$e^i$ is a basis of $L_i\otimes\mathbf Q$ over $E$, so the vectors $P(e^i_m)_{\tau_\sigma}$ span
$\operatorname{Lie}_\tau(y)$ and $C$ exists; $B$ exists and is unique because $w$ is a basis. By
Lemma~\ref{cv-cm-point-gram}(2), $\det S_i\ne0$. The relation
$\operatorname{Im}H(P(x),P(x'))=\langle x,x'\rangle$ expresses the pairing between the
$\tau_\sigma$- and the $\bar\tau_\sigma$-components of the periods through $S_i$, and a computation
in linear algebra with it shows that $\delta$ is nonzero and does not depend on the choice of $C$.
Two periods of $L_i$ that present $y$ differ by a linear automorphism of $\mathbf C^g$ and an
automorphism of $L_i$ that preserves the form and is congruent to $1$ modulo the level; the
determinant of the latter over $\mathcal O_E$ is $1$, so $\delta$ depends only on $(y,v,w)$.

(1) If $v'=gv$ and $w'=hw$, then $gC$ serves for $v'$, while $B$ and $F$ are replaced by
$(h^{\mathsf t})^{-1}B$ and $(h^{\mathsf t})^{-1}F$. So $\delta(y,v',w')=\delta(y,v,w)\det g/\det h$,
which is the transformation rule for $\Phi_\sigma$.

(2) The Hodge Gram matrices are $\Gamma(v)=(\pi H(v_a,v_b))_{a,b}$ and
$\Delta(w)=(\pi H(w_j,w_{j'}))_{j,j'}$. Expressing $v$ through $C$ and $w$ through $B$, the same
relation gives
\[
 \det\Gamma(v)=(2\pi)^{p-q}|\det S_i|\,|\delta(y,v,w)|^2\det\Delta(w),
\]
which is (2).

(3) Locally in $b$, the points $y(b)$ are presented by periods $P_b$ of one $L_i$ that depend
holomorphically on $b$, because the uniformization is holomorphic. Then $C$, $B$ and $F$ can be
chosen to depend holomorphically on $b$, and so does $\delta$, which does not vanish.
\end{proof}
% Lean: anUnit, with rho_i = (t_i)^{-1/2}. The existence of C, F and delta and the independence of C
% are GlobalTheory.T.uniformised_eigen_frame_one = anSolve (period_eigen_frame,
% period_hermitian_type_compatibility, hermitian_period_linear_algebra); the independence of the
% period is PositivityTheory.Xgn.analytic_type_definition (familyUnif_exp_unique, per_det_one,
% cond_transport); (1) is Unit.chg_of_isFlatAn with Unit.det_chg; (2) is
% GlobalTheory.T.uniformised_eigen_frame_two = hodgeNorm_flat (hodgeGram_iff_period,
% riemann_tau_component, norm_core), with t = (2 pi)^p (2 pi)^{-q} |det S|; (3) is
% GlobalTheory.T.uniformised_eigen_frames_two = anHol.

\begin{situation}\label{cv-boundary-situation}
In Situation~\ref{cv-function-situation}, suppose that $p+q\ge1$ and that $F_c$ is algebraically
closed in $F_c(Z)$. Let $(\mathcal X,\mathcal A,j)$ be a model of $X$ over a discrete valuation ring
$O\subset F_c$ with fraction field $F_c$ (Definition~\ref{cv-model}), let $\bar Z$ be the closure of
$j(Z)$ in $\mathcal X_{F_c}$, and let
$\mathcal L=\det\operatorname{Lie}_\tau(\mathcal A)\otimes(\det\operatorname{Lie}_{\bar\tau}(\mathcal A))^{-1}$,
a line bundle on $\bar Z$ by Lemma~\ref{cv-model-eigenbundles}. Let $\sigma$ be a complex embedding of
$F_c$ with $\sigma(\tau(\sqrt{-d}))=\theta$, and let $\Phi_\sigma$ be the function of
Lemma~\ref{cv-periods-function} for the uniformization of $X$ at $\sigma$. Points of $X$ are
identified with their images under $j$.
\end{situation}

\begin{definition}\label{cv-period-matrix}
Let $n\ge a\ge0$ and $g\ge0$ be integers. Let $\Delta^n\subset\mathbf C^n$ be the polydisc of the $z$
with $|z_k|<1$ for every $k$, let $\Delta^n_a\subset\Delta^n$ be the open subset where $z_k\ne0$ for
every $k\le a$, and let $\mathcal H_a\subset\mathbf C^n$ be the set of the $\xi$ with
$\operatorname{Re}\xi_k<0$ for $k\le a$ and $|\xi_k|<1$ for $k>a$. Let
$\exp_a:\mathcal H_a\to\Delta^n_a$ and $\log_a:\Delta^n_a\to\mathcal H_a$ be the maps given in the
coordinates $k\le a$ by $\xi_k\mapsto e^{\xi_k}$ and by the principal branch of $z_k\mapsto\log z_k$,
and in the coordinates $k>a$ by the identity; so $\exp_a\circ\log_a$ is the identity. For
$z\in\Delta^n_a$ put $\lambda_a(z)=1+\sum_{k\le a}\big|\log|z_k|\big|$, and let $\varepsilon_k$ be
the $k$-th unit vector of $\mathbf C^n$.

A \emph{period matrix over $\Delta^n_a$} of size $g$ is a pair $(Q,B)$ of a holomorphic map $Q$ from
$\mathcal H_a$ to the symmetric complex $g\times g$ matrices, such that the real matrix
$\operatorname{Im}Q(\xi)$ of the imaginary parts of the entries is positive definite for every
$\xi\in\mathcal H_a$, and a family $B=(B_1,\dots,B_a)$ of symmetric $g\times g$ matrices with
integral entries such that $Q(\xi+2\pi i\varepsilon_k)=Q(\xi)+B_k$ for every $k\le a$ and every
$\xi\in\mathcal H_a$. Its \emph{reduced matrix} is
\[
 Q_0(\xi)=Q(\xi)-\frac1{2\pi i}\sum_{k\le a}\xi_kB_k,
\]
which satisfies $Q_0(\xi+2\pi i\varepsilon_k)=Q_0(\xi)$ for every $k\le a$.
\end{definition}
\leanref{GroundTheory/HeightsAndTheta/Heights/Regularisation/Statements.lean}{GroundTheory.Heights₂.Heights.Regular.PeriodPolydisc}
\leanref{GroundTheory/HeightsAndTheta/Heights/Regularisation/Statements.lean}{GroundTheory.Heights₂.Heights.Regular.PeriodPolydisc.tau0}
\leanref{GroundTheory/HeightsAndTheta/Heights/Regularisation/Statements.lean}{GroundTheory.Heights₂.Heights.Regular.logCover}
\leanref{GroundTheory/HeightsAndTheta/Heights/Regularisation/Statements.lean}{GroundTheory.Heights₂.Heights.Regular.expCover}
\leanref{GroundTheory/HeightsAndTheta/Heights/Regularisation/Statements.lean}{GroundTheory.Heights₂.Heights.Regular.logLift}

\begin{lemma}\label{cv-period-riemann-slices}
Let $n\ge0$, let $\Delta^n\subset\mathbf C^n$ be the polydisc of the $z$ with $|z_k|<1$ for every
$k$, and let $G$ be a holomorphic function on the open subset of $\Delta^n$ where
$z_1\cdots z_n\ne0$. Suppose that for every $z\in\Delta^n$ and every $i\le n$ with $z_j\ne0$ for all
$j\ne i$, the function $x\mapsto G(z_1,\dots,z_{i-1},x,z_{i+1},\dots,z_n)$ has a limit as $x$ tends
to $0$ with $x\ne0$. Then there is a holomorphic function $G'$ on $\Delta^n$ that is equal to $G$
where $z_1\cdots z_n\ne0$.
\end{lemma}
\leanref{GroundTheory/HeightsAndTheta/Heights/Regularisation/StRiemannSlices.lean}{GroundTheory.Heights₂.Heights.Regular.RiemannExtensionSlices.Statement}
\begin{proof}
Let $0<\rho<1$, and let $T_\rho\subset\mathbf C^n$ be the torus of the $\zeta$ with $|\zeta_k|=\rho$
for every $k$. We prove by induction on $n$ that
\[
 G(z)=\frac1{(2\pi i)^n}\int_{T_\rho}\frac{G(\zeta)\,d\zeta_1\cdots d\zeta_n}
 {(\zeta_1-z_1)\cdots(\zeta_n-z_n)}
\]
for every $z$ with $0<|z_k|<\rho$ for every $k$. For $n=0$ both sides are $G(z)$. Let $n\ge1$, and
write $z=(z_1,z')$ and $\zeta=(\zeta_1,\zeta')$. For $|\zeta_1|=\rho$, the function
$\zeta'\mapsto G(\zeta_1,\zeta')$ satisfies the hypotheses of the lemma in $n-1$ variables, since
$\zeta_1\ne0$, so by induction $G(\zeta_1,z')$ is given by the formula in the variables $\zeta'$. The
function $x\mapsto G(x,z')$ is holomorphic for $0<|x|<1$ and has a limit at $0$, since no coordinate
of $z'$ vanishes. Extended by this limit, it is continuous on $|x|\le\rho$ and holomorphic for
$0<|x|<\rho$, so Cauchy's formula on the circle $|\zeta_1|=\rho$ gives $G(z_1,z')$ as the integral of
$G(\zeta_1,z')/(\zeta_1-z_1)$. Substituting the first formula into the second gives the formula for
$G(z)$, as an iterated integral of a continuous function on $T_\rho$.

The right side of the formula is a holomorphic function $G_\rho$ of $z$ on the polydisc $|z_k|<\rho$,
since the kernel is holomorphic in $z$ there and $G$ is continuous on $T_\rho$. For
$\rho<\rho'<1$, the functions $G_\rho$ and $G_{\rho'}$ are equal to $G$ on the dense subset of
$|z_k|<\rho$ where $z_1\cdots z_n\ne0$, so they agree on $|z_k|<\rho$. Together they define a
holomorphic function $G'$ on $\Delta^n$, equal to $G$ where $z_1\cdots z_n\ne0$.
\end{proof}
\leanref{GroundTheory/HeightsAndTheta/Heights/Regularisation/PRiemannExtensionSlices.lean}{GroundTheory.Heights₂.Heights.Regular.RiemannExtensionSlices.proof}

\begin{lemma}\label{cv-period-punctured-order}
Let $R>0$, let $\psi$ be a holomorphic function on the punctured disc $0<|x|<R$, and let
$k\in\mathbf Z$ and $M\in\mathbf R$ be such that $|e^{\psi(x)}x^k|\le M$ for $0<|x|<R$. Then $k\ge0$,
and $e^{\psi(x)}$ has a nonzero limit as $x$ tends to $0$.
\end{lemma}
\leanref{GroundTheory/HeightsAndTheta/Heights/Regularisation/StPuncturedExpOrder.lean}{GroundTheory.Heights₂.Heights.Regular.PuncturedExpOrder.Statement}
\begin{proof}
The function $F(x)=e^{\psi(x)}x^k$ is holomorphic and bounded on $0<|x|<R$, so it extends
holomorphically to $|x|<R$ by Riemann's theorem on removable singularities. It has no zeros for
$x\ne0$, so near $0$ it is $x^mu(x)$ with an integer $m\ge0$ and a holomorphic function $u$ with
$u(0)\ne0$. Choose $0<\delta<R$ such that $u$ is holomorphic without zeros on $|x|<\delta$ and
$F=x^mu$ there. Then $e^{\psi(x)}=x^{m-k}u(x)$ for $0<|x|<\delta$, and taking logarithmic
derivatives,
\[
 \psi'(x)=\frac{m-k}{x}+\frac{u'(x)}{u(x)}.
\]
Integrate over the circle $|x|=\delta/2$. The integral of $\psi'$ is $0$, since $\psi'$ has the
primitive $\psi$; the integral of $u'/u$ is $0$ by Cauchy's theorem, since $u'/u$ is holomorphic on
$|x|<\delta$; and the integral of $(m-k)/x$ is $2\pi i(m-k)$. So $k=m\ge0$, and
$e^{\psi(x)}=u(x)$ tends to $u(0)\ne0$.
\end{proof}
\leanref{GroundTheory/HeightsAndTheta/Heights/Regularisation/PPuncturedExpOrder.lean}{GroundTheory.Heights₂.Heights.Regular.PuncturedExpOrder.proof}

\begin{lemma}\label{cv-period-log-descent}
In the notation of Definition~\ref{cv-period-matrix}, let $\varphi$ be a holomorphic function on
$\mathcal H_a$ with $\varphi(\xi+2\pi i\varepsilon_k)=\varphi(\xi)$ for every $k\le a$ and every
$\xi\in\mathcal H_a$. Then $\varphi\circ\log_a$ is holomorphic on $\Delta^n_a$, and
$\varphi(\log_a(\exp_a\xi))=\varphi(\xi)$ for every $\xi\in\mathcal H_a$.
\end{lemma}
\leanref{GroundTheory/HeightsAndTheta/Heights/Regularisation/StLogCoverDescent.lean}{GroundTheory.Heights₂.Heights.Regular.LogCoverDescent.Statement}
\begin{proof}
The translation by $2\pi i\varepsilon_k$, for $k\le a$, and its inverse map $\mathcal H_a$ to
itself. So $\varphi(\xi')=\varphi(\xi)$ whenever $\xi,\xi'\in\mathcal H_a$ and
$\exp_a\xi'=\exp_a\xi$, since then $\xi'-\xi=2\pi i\sum_{k\le a}m_k\varepsilon_k$ with integers
$m_k$. As $\exp_a(\log_a(\exp_a\xi))=\exp_a\xi$, this gives the second assertion. Let
$z^0\in\Delta^n_a$. For $z$ near $z^0$, let $L(z)$ be the point with the coordinates
$\log z^0_k+\log(z_k/z^0_k)$ for $k\le a$, with principal branches, and $z_k$ for $k>a$. Then $L$ is
holomorphic near $z^0$, with values in $\mathcal H_a$ as $\operatorname{Re}L(z)_k=\log|z_k|<0$, and
$\exp_aL(z)=z=\exp_a\log_az$. So $\varphi(\log_az)=\varphi(L(z))$ near $z^0$, and this is a
holomorphic function of $z$.
\end{proof}
\leanref{GroundTheory/HeightsAndTheta/Heights/Regularisation/PLogCoverDescent.lean}{GroundTheory.Heights₂.Heights.Regular.LogCoverDescent.proof}

\begin{lemma}\label{cv-period-exp-monomial}
Let $(Q,B)$ be a period matrix over $\Delta^n_a$ of size $g$ (Definition~\ref{cv-period-matrix}),
and let $v\in\mathbf Z^g$. Then $v^{\mathsf t}B_kv\ge0$ for every $k\le a$, and there is a
holomorphic function $h$ on $\Delta^n$ without zeros such that for every $\xi\in\mathcal H_a$, with
$z=\exp_a\xi$,
\[
 \exp\big(2\pi i\,v^{\mathsf t}Q(\xi)v\big)=h(z)\prod_{k\le a}z_k^{v^{\mathsf t}B_kv}.
\]
\end{lemma}
\leanref{GroundTheory/HeightsAndTheta/Heights/Regularisation/Statements.lean}{GroundTheory.Heights₂.Heights.Regular.ExpQuadMonomial.Statement}
\begin{proof}
Put $b_k=v^{\mathsf t}B_kv\in\mathbf Z$ and $\varphi(\xi)=2\pi i\,v^{\mathsf t}Q_0(\xi)v$, where
$Q_0$ is the reduced matrix. The function $\varphi$ is holomorphic on $\mathcal H_a$ and invariant
under the translations by $2\pi i\varepsilon_k$, so by Lemma~\ref{cv-period-log-descent} the function
$\psi=\varphi\circ\log_a$ is holomorphic on $\Delta^n_a$, with $\psi(\exp_a\xi)=\varphi(\xi)$. By the
definition of $Q_0$,
\[
 \exp\big(2\pi i\,v^{\mathsf t}Q(\xi)v\big)=e^{\varphi(\xi)}\prod_{k\le a}e^{b_k\xi_k},
\]
and the left side has absolute value $\exp(-2\pi\,v^{\mathsf t}\operatorname{Im}Q(\xi)v)\le1$.

Let $z\in\Delta^n$ and $i\le a$ be such that $z_j\ne0$ for every $j\le a$ other than $i$, and let
$z(x)$ be $z$ with its $i$-th coordinate replaced by $x$. For $0<|x|<1$ the point $z(x)$ lies in
$\Delta^n_a$, and the identity above at $\xi=\log_az(x)$ gives
\[
 \big|e^{\psi(z(x))}x^{b_i}\big|\le|c_z|^{-1},\qquad c_z=\prod_{j\le a,\,j\ne i}z_j^{b_j}\ne0 .
\]
By Lemma~\ref{cv-period-punctured-order}, $b_i\ge0$, and $e^{\psi(z(x))}$ has a nonzero limit as $x$
tends to $0$, so $e^{-\psi(z(x))}$ has a limit too. Along a slice in a coordinate $m>a$, through a
point whose other coordinates are nonzero, the point of the slice with $m$-th coordinate $0$ lies in
$\Delta^n_a$, where $\psi$ is holomorphic, so $e^{\psi}$ and $e^{-\psi}$ have limits along the slice
as well. Lemma~\ref{cv-period-riemann-slices}, applied to the restrictions of $e^{\psi}$ and
$e^{-\psi}$ to the points with no zero coordinate, gives holomorphic functions $h$ and $h'$ on
$\Delta^n$ equal to $e^{\psi}$ and $e^{-\psi}$ there. Then $hh'=1$ on a dense subset of $\Delta^n$,
hence everywhere, and $h$ has no zeros. Both $h$ and $e^\psi$ are continuous on $\Delta^n_a$, so
$h=e^{\psi}$ there. For $\xi\in\mathcal H_a$ and $z=\exp_a\xi$ this gives $e^{\varphi(\xi)}=h(z)$,
and $e^{b_k\xi_k}=z_k^{b_k}$, which proves the formula.
\end{proof}
\leanref{GroundTheory/HeightsAndTheta/Heights/Regularisation/PExpExtension.lean}{GroundTheory.Heights₂.Heights.Regular.ExpQuadMonomial.proof}

\begin{lemma}\label{cv-period-reduced-extends}
Let $(Q,B)$ be a period matrix over $\Delta^n_a$ of size $g$ (Definition~\ref{cv-period-matrix}),
with reduced matrix $Q_0$. Then $B_1,\dots,B_a$ are positive semidefinite, and there is a
holomorphic map $T_0$ from $\Delta^n$ to the complex $g\times g$ matrices with
$Q_0(\xi)=T_0(\exp_a\xi)$ for every $\xi\in\mathcal H_a$.
\end{lemma}
\leanref{GroundTheory/HeightsAndTheta/Heights/Regularisation/Statements.lean}{GroundTheory.Heights₂.Heights.Regular.Tau0Holomorphic.Statement}
\begin{proof}
By Lemma~\ref{cv-period-exp-monomial}, $v^{\mathsf t}B_kv\ge0$ for every $v\in\mathbf Z^g$. Let
$x\in\mathbf R^g$, and let $x_N=N^{-1}\lfloor Nx\rfloor$ for integers $N\ge1$, with the integer part
taken in each coordinate. Then $x_N^{\mathsf t}B_kx_N=N^{-2}\lfloor Nx\rfloor^{\mathsf t}B_k\lfloor
Nx\rfloor\ge0$ and $x_N\to x$, so $x^{\mathsf t}B_kx\ge0$, and $B_k$ is positive semidefinite.

Let $v\in\mathbf Z^g$, put $b_k=v^{\mathsf t}B_kv$, and let $h$ be as in
Lemma~\ref{cv-period-exp-monomial}. The polydisc $\Delta^n$ is convex and $h$ has no zeros on it, so
$h=e^G$ for a holomorphic function $G$ on $\Delta^n$. Let $\varphi(\xi)=2\pi i\,v^{\mathsf t}Q_0(\xi)v$.
By the definition of $Q_0$,
$\exp(2\pi i\,v^{\mathsf t}Q(\xi)v)=e^{\varphi(\xi)}\prod_{k\le a}(e^{\xi_k})^{b_k}$, and comparing with
Lemma~\ref{cv-period-exp-monomial} gives $e^{\varphi(\xi)}=h(\exp_a\xi)=e^{G(\exp_a\xi)}$. So
$\varphi-G\circ\exp_a$ is a continuous function on the convex set $\mathcal H_a$ with values in
$2\pi i\mathbf Z$, hence a constant $c_v$, and $v^{\mathsf t}Q_0(\xi)v=q_v(\exp_a\xi)$ with
$q_v=(G+c_v)/(2\pi i)$ holomorphic on $\Delta^n$. The matrix $Q_0(\xi)$ is symmetric, as $Q(\xi)$ and
the $B_k$ are. With $u_1,\dots,u_g$ the unit vectors of $\mathbf Z^g$, its entries are therefore
\[
 Q_0(\xi)_{pq}=\tfrac12\big(q_{u_p+u_q}-q_{u_p}-q_{u_q}\big)(\exp_a\xi),
\]
and the right sides are the entries of the required $T_0$.
\end{proof}
\leanref{GroundTheory/HeightsAndTheta/Heights/Regularisation/PosSemidefOfIntQuadNonneg.lean}{GroundTheory.Heights₂.Heights.Regular.Tau0Holomorphic.proof}
\leanref{GroundTheory/HeightsAndTheta/Heights/Regularisation/Assembly.lean}{GroundTheory.Heights₂.Heights.Regular.Tau0Holomorphic.cond}

\begin{lemma}\label{cv-period-kernel-positive}
Let $(Q,B)$ be a period matrix over $\Delta^n_a$ of size $g$ (Definition~\ref{cv-period-matrix}),
with reduced matrix $Q_0$, and let $T_0$ be a holomorphic map from $\Delta^n$ to the complex
$g\times g$ matrices with $Q_0(\xi)=T_0(\exp_a\xi)$ for every $\xi\in\mathcal H_a$. Let
$v\in\mathbf R^g$ be nonzero with $B_kv=0$ for every $k\le a$. Then
$v^{\mathsf t}\operatorname{Im}T_0(z)\,v>0$ for every $z\in\Delta^n$.
\end{lemma}
\leanref{GroundTheory/HeightsAndTheta/Heights/Regularisation/Statements.lean}{GroundTheory.Heights₂.Heights.Regular.KernelPositive.Statement}
\begin{proof}
Let $f(z)=v^{\mathsf t}T_0(z)v$, a holomorphic function on $\Delta^n$ whose imaginary part is
$v^{\mathsf t}\operatorname{Im}T_0(z)\,v$. Let $z\in\Delta^n$ have no zero coordinate. Then
$\xi=\log_az$ lies in $\mathcal H_a$ with $\exp_a\xi=z$, and since $B_kv=0$ for every $k$,
$f(z)=v^{\mathsf t}Q_0(\xi)v=v^{\mathsf t}Q(\xi)v$, whose imaginary part
$v^{\mathsf t}\operatorname{Im}Q(\xi)v$ is positive.

Now let $z\in\Delta^n$ be arbitrary, let $\mathbf 1=(1,\dots,1)$, and let $F(t)=f(z+t\mathbf 1)$ for
$t\in\mathbf C$ near $0$. There is a disc $V$ around $0$ such that for $t\in V$ with $t\ne0$ the
point $z+t\mathbf 1$ lies in $\Delta^n$ and has no zero coordinate, so that
$\operatorname{Im}F(t)>0$. If $F$ is constant on $V$, then $\operatorname{Im}F(0)>0$. Otherwise,
by the open mapping theorem, $F(V)$ contains a neighbourhood of $F(0)$. If
$\operatorname{Im}F(0)\le0$, this neighbourhood contains a point $F(0)-i\delta$ with $\delta>0$, which
is $F(t)$ for some $t\in V$ with $t\ne0$ and has negative imaginary part; this is impossible. So
$\operatorname{Im}f(z)=\operatorname{Im}F(0)>0$.
\end{proof}
\leanref{GroundTheory/HeightsAndTheta/Heights/Regularisation/PKernelPositive.lean}{GroundTheory.Heights₂.Heights.Regular.kernelPositive_statement}

\begin{proposition}\label{cv-period-matrix-bounds}
Let $(Q,B)$ be a period matrix over $\Delta^n_a$ of size $g$ (Definition~\ref{cv-period-matrix}).
Then there are real numbers $0<\rho<1$, $c>0$ and $C$ such that
\[
 c\,v^{\mathsf t}v\le v^{\mathsf t}\operatorname{Im}Q(\xi)\,v\le C\lambda_a(z)\,v^{\mathsf t}v
\]
for every $v\in\mathbf R^g$ and every $\xi\in\mathcal H_a$ whose image $z=\exp_a\xi$ satisfies
$|z_k|\le\rho$ for every $k$.
\end{proposition}
\leanref{GroundTheory/HeightsAndTheta/Heights/Regularisation/Statements.lean}{GroundTheory.Heights₂.Heights.Regular.ImTauBounds.Statement}
\leanref{GroundTheory/HeightsAndTheta/Heights/Regularisation/Assembly.lean}{GroundTheory.Heights₂.Heights.Regular.ImTauBounds.cond}
\begin{proof}
By Lemma~\ref{cv-period-reduced-extends}, the $B_k$ are positive semidefinite and
$Q_0(\xi)=T_0(\exp_a\xi)$ for a holomorphic map $T_0$ on $\Delta^n$, where $Q_0$ is the reduced
matrix. Let $M(z)=\operatorname{Im}T_0(z)$, a continuous function of $z\in\Delta^n$ with values in
the real symmetric matrices. Taking imaginary parts in the definition of $Q_0$ gives, with
$z=\exp_a\xi$,
\[
 \operatorname{Im}Q(\xi)=M(z)+\sum_{k\le a}\ell_kB_k,\qquad
 \ell_k=-\frac{\operatorname{Re}\xi_k}{2\pi}=\frac{\big|\log|z_k|\big|}{2\pi}.
\]

Let $S=B_1+\dots+B_a$, which is positive semidefinite, and let $K\subset\Delta^n$ be the compact
polydisc $|z_k|\le\frac12$. If $v\ne0$ and $Sv=0$, then $v^{\mathsf t}B_kv=0$ for every $k$, so
$B_kv=0$ as $B_k$ is positive semidefinite, and $v^{\mathsf t}M(z)v>0$ for every $z\in\Delta^n$ by
Lemma~\ref{cv-period-kernel-positive}. It follows that there are an integer $L\ge0$ and $c>0$
with
\[
 c\,v^{\mathsf t}v\le L\,v^{\mathsf t}Sv+v^{\mathsf t}M(z)v\qquad(z\in K,\ v\in\mathbf R^g).
\]
Indeed, for an integer $m\ge0$ let $U_m$ be the set of $(z,u)\in\Delta^n\times\mathbf R^g$ with
$m\,u^{\mathsf t}Su+u^{\mathsf t}M(z)u>1/(m+1)$. The sets $U_m$ are open and increase with $m$, as
$u^{\mathsf t}Su\ge0$. They cover $K\times\Sigma$, where $\Sigma$ is the unit sphere of
$\mathbf R^g$: if $Su=0$ then $u^{\mathsf t}M(z)u>0$, and otherwise $u^{\mathsf t}Su>0$. As
$K\times\Sigma$ is compact, it lies in one $U_m$, and $L=m$ and $c=1/(m+1)$ satisfy the
inequality, which is homogeneous in $v$. By compactness again, there are $A\ge0$ and
$\beta_1,\dots,\beta_a\ge0$ with $v^{\mathsf t}M(z)v\le A\,v^{\mathsf t}v$ for $z\in K$ and
$v^{\mathsf t}B_kv\le\beta_k\,v^{\mathsf t}v$, for every $v$.

Let $\rho=\min\{\frac12,e^{-2\pi L}\}$ and $C=A+\beta_1+\dots+\beta_a$, and let
$\xi\in\mathcal H_a$ be such that $z=\exp_a\xi$ satisfies $|z_k|\le\rho$ for every $k$. Then
$z\in K$, and $e^{\operatorname{Re}\xi_k}=|z_k|\le e^{-2\pi L}$ gives $\ell_k\ge L$ for $k\le a$. As
$v^{\mathsf t}B_kv\ge0$,
\[
 v^{\mathsf t}\operatorname{Im}Q(\xi)v\ge v^{\mathsf t}M(z)v+L\,v^{\mathsf t}Sv\ge c\,v^{\mathsf t}v .
\]
On the other side, $\ell_k\le\lambda_a(z)$ and $\lambda_a(z)\ge1$, so
\[
 v^{\mathsf t}\operatorname{Im}Q(\xi)v\le A\,v^{\mathsf t}v+\lambda_a(z)\sum_{k\le a}\beta_k\,v^{\mathsf t}v
 \le C\lambda_a(z)\,v^{\mathsf t}v .
\]
\end{proof}
\leanref{GroundTheory/HeightsAndTheta/Heights/Regularisation/PImTauBounds.lean}{GroundTheory.Heights₂.Heights.Regular.ImTauBounds.proof}

\begin{lemma}\label{cv-boundary-hodge-growth}
In Situation~\ref{cv-boundary-situation}, let $\pi:X'\to\mathcal X_{F_c}$ be a proper morphism from a
scheme $X'$ smooth of dimension $D$ over $F_c$, that restricts to an isomorphism from
$\pi^{-1}(j(X))$ onto $j(X)$, with $X'\setminus\pi^{-1}(j(X))$ cut out locally by coordinate
functions as in Lemma~\ref{cv-resolution-coordinates}(3). Let $x$ be a complex point of $X'$ over
$\sigma$, and let $z_1,\dots,z_D$ be holomorphic coordinates that identify a neighbourhood
$\mathcal N$ of $x$ in the complex points of $X'$ over $\sigma$ with the polydisc $|z_k|<1$, with $x$
at $0$ and with $\mathcal N\cap\pi^{-1}(j(X))$ the set where $z_1\cdots z_a\ne0$. Let
$e_1,\dots,e_m$ be sections of the Lie algebra of $\mathcal A$ on an affine open neighbourhood of
$\pi(x)$ in $\mathcal X$ whose values at $\pi(x)$ are linearly independent. Then there are
$0<r\le1$, $K>0$ and an integer $N\ge0$ such that at every point $y$ of
$\mathcal N\cap\pi^{-1}(j(X))$ with $|z_k(y)|<r$ for all $k$, the Hodge Gram matrix $\Gamma(y)$ of
the values of $e_1,\dots,e_m$ at $\pi(y)$ satisfies
\[
 \det\Gamma(y)\le K,\qquad K\lambda(y)^N\det\Gamma(y)\ge1,\qquad
 \lambda(y)=1+\sum_{k\le a}\big|\log|z_k(y)|\big|.
\]
\end{lemma}
\leanref{Development/AutomorphicSpectrum/Proofs/GeometryProofs/HodgeGramBoundNearBoundaryStatement.lean}{GlobalTheory.T.hodge_gram_bound_near_boundary}
\begin{proof}
Shrink $\mathcal N$. Over $\mathcal N\cap\pi^{-1}(j(X))$ the kernels of the exponential maps of the
abelian varieties at the points form lattices of rank $2g$; continued along $\log z_k$ for $k\le a$,
they become one lattice with the Riemann form of the polarization, and the periods that are
unchanged by these continuations extend holomorphically over $\mathcal N$ and span the Lie algebra
at $\pi(x)$. In a symplectic basis of the lattice and a holomorphic basis of the Lie algebra of
$\mathcal A$ near $\pi(x)$ that is a basis at $\pi(x)$, the periods are given by a symmetric matrix
$Q(\xi)$ with positive definite imaginary part, holomorphic in
$\xi=(\log z_1,\dots,\log z_a,z_{a+1},\dots,z_D)$, with
$Q(\xi+2\pi i\varepsilon_k)=Q(\xi)+B_k$ for integral symmetric matrices $B_k$, where
$\varepsilon_k$ is the $k$-th unit vector; and the Hodge Gram matrix of that basis is
$\pi(\operatorname{Im}Q)^{-1}$. So $(Q,(B_k)_{k\le a})$ is a period matrix over $\Delta^D_a$ in the
coordinates $z_1,\dots,z_D$ (Definition~\ref{cv-period-matrix}), and
$\lambda(y)=\lambda_a(z_1(y),\dots,z_D(y))$. By Proposition~\ref{cv-period-matrix-bounds}, there are
$c,C>0$ with $c\le\operatorname{Im}Q(\xi)\le C\lambda(y)$ as hermitian matrices for $y$ near $x$.
% Lean: ImTauBounds.cond RiemannExtensionSlices.proof PuncturedExpOrder.proof
% LogCoverDescent.proof, applied in logPeriods_subframe_det_bounds with chartRadius_le_log_sum.
Write the values of $e_1,\dots,e_m$
in the holomorphic basis as the rows of an $m\times g$ matrix $M(y)$, continuous near $x$ and of rank
$m$ at $x$. Then $\Gamma(y)=\pi M(\operatorname{Im}Q)^{-1}\bar M^{\mathsf t}$, so
\[
 \pi^m(C\lambda(y))^{-m}\det(M\bar M^{\mathsf t})\le\det\Gamma(y)\le\pi^mc^{-m}\det(M\bar M^{\mathsf t}).
\]
For $r$ small, $\det(M\bar M^{\mathsf t})$ lies between two positive constants, and the bounds hold
with $N=m$.
\end{proof}
% Lean: hodgeGram_det_log_bounds; the periods are exists_logPeriods_isRiem (LogPeriods, through
% exists_family_exponential_locally, period_lattice_discrete_free, period_holomorphic_continuation);
% the period matrix is period_matrix_local_trivialization = exists_periodPolydisc_frame
% (PeriodPolydisc); the bounds on Im tau are ImTauBounds.cond; the determinant bounds are
% logPeriods_subframe_det_bounds (subframe_det_log_growth).

\begin{lemma}\label{cv-boundary-unit}
In Situation~\ref{cv-boundary-situation}, let $\pi:X'\to\mathcal X_{F_c}$ be a proper morphism from a
scheme $X'$ smooth over $F_c$, that restricts to an isomorphism from $\pi^{-1}(j(X))$ onto $j(X)$,
with $X'\setminus\pi^{-1}(j(X))$ cut out locally by coordinate functions as in
Lemma~\ref{cv-resolution-coordinates}(3). Let $x$ be a complex point of $X'$ over $\sigma$, and let
$e_1=(e_{1,1},\dots,e_{1,p})$ and $e_2=(e_{2,1},\dots,e_{2,q})$ be families of sections of the Lie
algebra of $\mathcal A$ on an affine open neighbourhood $U$ of $\pi(x)$ in $\mathcal X$ whose values
at every complex point of $U$ over $\sigma$ are bases of the subspaces on which $\mathcal O_E$ acts
through $\sigma\tau$ and through $\sigma\bar\tau$. Then there are a neighbourhood $\mathcal N$ of $x$
in the complex points of $X'$ over $\sigma$ and a holomorphic function $h$ on $\mathcal N$ without
zeros such that $h(y)=\Phi_\sigma(\pi(y),e_1(\pi(y)),e_2(\pi(y)))$ at every
$y\in\mathcal N\cap\pi^{-1}(j(X))$.
\end{lemma}
\leanref{Development/AutomorphicSpectrum/Proofs/GeometryProofs/AffineChartBundleSectionsStatement.lean}{GlobalTheory.T.affine_chart_bundle_sections}
\leanref{Development/AutomorphicSpectrum/Proofs/GeometryProofs/BoundaryChartFrameGrowthStatement.lean}{GlobalTheory.T.boundary_chart_frame_growth}
\begin{proof}
By the coordinates of Lemma~\ref{cv-resolution-coordinates}(3) and the holomorphic inverse function
theorem, there are holomorphic coordinates $z_1,\dots,z_D$ that identify a neighbourhood
$\mathcal N_0$ of $x$, mapped by $\pi$ into $U$, with the polydisc $|z_k|<1$, with $x$ at $0$ and
with $\mathcal N_0\cap\pi^{-1}(j(X))$ the set where $z_1\cdots z_a\ne0$. On this set
$\phi(y)=\Phi_\sigma(\pi(y),e_1(\pi(y)),e_2(\pi(y)))$ is holomorphic by
Lemma~\ref{cv-periods-function}(3), since $e_1$ and $e_2$ are algebraic sections, and it does not
vanish. Lemma~\ref{cv-boundary-hodge-growth}, applied to $e_1$ and to $e_2$, gives $r$, $K$ and $N$
with $\det\Gamma\le K$ and $K\lambda^N\det\Gamma\ge1$ for the Hodge Gram matrices $\Gamma$ of $e_1$ and
of $e_2$, where $|z_k|<r$ and $z_1\cdots z_a\ne0$. By Lemma~\ref{cv-periods-function}(2),
$|\phi|^2$ is the quotient of these two determinants, so $|\phi|\le K\lambda^{N/2}$ and
$|\phi|^{-1}\le K\lambda^{N/2}$ there.

A holomorphic function $\phi$ with these bounds extends to a holomorphic function without zeros on
the polydisc $|z_k|<r$. Indeed, $\eta=z_1\cdots z_a\phi$ is bounded, since
$\epsilon\big(1+a|\log\epsilon|\big)^{N/2}$ is bounded for $0<\epsilon<1$, so $\eta$ extends
holomorphically to the polydisc by Riemann's extension theorem. The same bound tends to $0$ with
$\epsilon$, applied to the least of the $|z_k|$, $k\le a$, so $\eta$ vanishes on each hyperplane
$z_k=0$, and $\eta$ is divisible by $z_1\cdots z_a$; the quotient $h$ extends $\phi$. The same
argument extends $1/\phi$ to a holomorphic $h'$, and $hh'=1$ by continuity, so $h$ has no zeros.
\end{proof}
% Lean: flat_unit_along_families (holomorphic_affine_polydisc_chart, ChartFamily.exists_model_chart,
% ChartFamily.chart_unit) and exists_unit_extension_of_bdryChart (hodge_gram_bound_near_boundary,
% IsUnitAn.norm, norm_bounds_of_gram, exists_unit_extension_of_log_growth).

\begin{lemma}\label{cv-boundary-gaga}
Let $K$ be a field of characteristic zero with an embedding $\sigma:K\to\mathbf C$, let $Y$ be a
smooth proper scheme over $K$, and let $Y_0$ be an irreducible component of $Y$ whose function field
$K(Y_0)$ contains no element algebraic over $K$ outside $K$. Let $\mathcal M$ be an invertible
$\mathcal O_{Y_0}$-submodule of the constant sheaf $K(Y_0)$ on $Y_0$, let $W\subset Y_0$ be a dense
open subscheme with $\mathcal M|_W=\mathcal O_W$, and let $F$ be a complex function on the complex
points of $W$ over $\sigma$. Suppose that every complex point $x$ of $Y_0$ over $\sigma$ has a
neighbourhood $\mathcal N_x$ in the complex points of $Y_0$ over $\sigma$, an element $t_x\in K(Y_0)$
that generates $\mathcal M$ on a Zariski neighbourhood of $x$ containing $\mathcal N_x$, and a
holomorphic function $h_x$ on $\mathcal N_x$ without zeros, such that $F(y)=h_x(y)t_x(y)$ at every
$y\in\mathcal N_x\cap W$. Then there are $f\in K(Y_0)$ that generates $\mathcal M$ at every point of
$Y_0$ and $\kappa\in\mathbf C^\times$ with $f(y)=\kappa F(y)$ at every complex point $y$ of $W$ over
$\sigma$.
\end{lemma}
\leanref{Development/AutomorphicSpectrum/Proofs/GeometryProofs/GagaLineGeneratorEmbeddingStatement.lean}{GlobalTheory.T.gaga_line_generator_embedding}
\begin{proof}
Let $Y_{0,\mathbf C}=Y_0\times_{K,\sigma}\mathbf C$. Since $K$ is algebraically closed in $K(Y_0)$ and
has characteristic zero, $Y_{0,\mathbf C}$ is integral, it is a connected component of the
smooth proper scheme $Y\times_{K,\sigma}\mathbf C$, and $\mathbf C$ is algebraically closed in its
function field; its complex points are the complex points of $Y_0$ over $\sigma$. On a Zariski open
set $U$ on which some $t_x$ generates $\mathcal M$, the function $F/t_x$ on the complex points of
$U\cap W$ extends holomorphically to all complex points of $U$, as $h_{x'}t_{x'}/t_x$ near each of
them; for two such open sets these extensions satisfy the gluing condition of Theorem~\ref{cv-gaga}
with the transition functions $t_x/t_{x'}$, because they do on the dense set $W$. So
Theorem~\ref{cv-gaga} gives a global section $f'$ of the pullback $\mathcal M_{\mathbf C}$ of
$\mathcal M$ to $Y_{0,\mathbf C}$ with $f'=F$ on $W$, and $f'$ generates $\mathcal M_{\mathbf C}$ at
every point because the $h_x$ have no zeros. Then $\mathcal M_{\mathbf C}$ is trivial, and its global
sections are $\mathbf Cf'$, because the global functions on the proper integral scheme
$Y_{0,\mathbf C}$ are constant. By flat base change the global sections of $\mathcal M$ form a
$K$-vector space of dimension one, spanned by some $f$ with $f=cf'$ for some $c\in\mathbf C^\times$.
The section $f$ generates $\mathcal M$ at every point, since it does after the faithfully flat base
change to $\mathbf C$. Put $\kappa=c$.
\end{proof}
% Lean: gaga_rat: cpt_base_change, GlobalTheory.T.gaga_line_generator_complex = gaga_rat_C
% (Theorem cv-gaga over C), and GlobalTheory.T.proper_model_generic_fibre_compare =
% exists_gen_of_extension (the descent of the generator).

\begin{lemma}\label{cv-boundary-section}
In Situation~\ref{cv-boundary-situation}, let $b_1$ and $b_2$ be bases of
$\operatorname{Lie}_\tau(\mathcal A)$ and $\operatorname{Lie}_{\bar\tau}(\mathcal A)$ at the generic
point of $\bar Z$. Then there are $f\in F_c(Z)$, $\kappa\in\mathbf C^\times$ and an open neighbourhood
$W$ of the generic point of $Z$ in $X$ such that $f\wedge b_1\otimes(\wedge b_2)^{-1}$ generates
$\mathcal L$ at every point of $\bar Z$, and such that at every complex point $y$ of $Z\cap W$ over
$\sigma$ the values $v$ and $w$ of $b_1$ and $b_2$ are defined and are bases of
$\operatorname{Lie}_\tau(y)$ and $\operatorname{Lie}_{\bar\tau}(y)$, $f$ is regular at $y$, and
$f(y)=\kappa\Phi_\sigma(y,v,w)$.
\end{lemma}
\leanref{Development/AutomorphicSpectrum/Proofs/GeometryProofs/EigenFramesAnalyticStatement.lean}{GlobalTheory.T.eigen_frames_analytic}
\leanref{Development/AutomorphicSpectrum/Proofs/GeometryProofs/EigenFrameFunctionCompareStatement.lean}{GlobalTheory.T.eigen_frame_function_compare}
\leanref{Development/AutomorphicSpectrum/Proofs/GeometryProofs/ProperPullbackLineBundleStatement.lean}{GlobalTheory.T.proper_pullback_line_bundle}
\begin{proof}
By Lemma~\ref{cv-model-eigenbundles}, $b_1$ and $b_2$ are bases of the two eigenbundles on an open
neighbourhood $W_0$ of the generic point of $\bar Z$, so $s_0=\wedge b_1\otimes(\wedge b_2)^{-1}$
generates $\mathcal L$ on $W_0$; through $s_0\mapsto1$, $\mathcal L$ is an invertible submodule of the
constant sheaf $F_c(Z)$ on $\bar Z$. Let $W=Z\cap W_0$. At the points of $W$, the values of $b_1$ and
$b_2$ are bases of $\operatorname{Lie}_\tau(y)$ and $\operatorname{Lie}_{\bar\tau}(y)$ by
Lemma~\ref{cv-model-eigenbundles}.

The generic fibre $\mathcal X_{F_c}$ is projective over $F_c$ and reduced, and $j(X)$ is a dense open
subscheme, smooth over $F_c$. Lemma~\ref{cv-resolution-coordinates} gives a proper morphism
$\pi:X'\to\mathcal X_{F_c}$ from a smooth projective scheme $X'$ over $F_c$ that restricts to an
isomorphism over $j(X)$, with $\pi^{-1}(j(X))$ dense and its complement cut out locally by
coordinate functions. Let $X'_Z$ be the closure of $\pi^{-1}(j(Z))$, an irreducible component of $X'$
with function field $F_c(Z)$; being proper with dense image, $\pi$ maps $X'_Z$ onto $\bar Z$.

For a complex point $y$ of $\pi^{-1}(W)$ over $\sigma$, let $F(y)=\Phi_\sigma(\pi(y),v,w)$, with $v$
and $w$ the values of $b_1$ and $b_2$ at $\pi(y)$. Let $x$ be a complex point of $X'_Z$ over $\sigma$.
Choose an affine open neighbourhood $U$ of $\pi(x)$ in $\mathcal X$ with $U_{F_c}\subset\bar Z$,
which is open in the normal scheme $\mathcal X_{F_c}$, and sections $e_1$ and $e_2$ of the Lie
algebra of $\mathcal A$ on $U$ whose values at the points of $U_{F_c}$ are bases of the two
eigenspaces (Lemma~\ref{cv-model-eigenbundles}, with sections over $U_{F_c}$ multiplied by a power
of a uniformizer of $O$). At the generic point, $e_1=g_1b_1$ and $e_2=g_2b_2$ with invertible
matrices $g_1$ and $g_2$ over $F_c(Z)$, and $t_x=\det g_1/\det g_2$, the element corresponding to
$\wedge e_1\otimes(\wedge e_2)^{-1}$, generates $\mathcal L$ near $\pi(x)$. By
Lemma~\ref{cv-periods-function}(1), $F(y)=\Phi_\sigma(\pi(y),e_1,e_2)\,t_x(\pi(y))$ near $x$, and by
Lemma~\ref{cv-boundary-unit} the first factor is a holomorphic function without zeros near $x$.

So Lemma~\ref{cv-boundary-gaga} applies to $Y_0=X'_Z$, the pullback of $\mathcal L$ and $F$, and gives
$f\in F_c(Z)$ that generates the pullback of $\mathcal L$ at every point of $X'_Z$ and $\kappa$ with
$f=\kappa F$ on $\pi^{-1}(W)$. Then $f$ generates $\mathcal L$ at every point of $\bar Z$. Let $t$
generate $\mathcal L$ on an open set $U'\subset\bar Z$. At a point $u$ of codimension one in $U'$, the
local ring of the normal scheme $\bar Z$ is a discrete valuation ring; $f/t$ is a unit at a point of
$X'_Z$ over $u$, and a local homomorphism of local rings maps nonunits to nonunits, so $f/t$ has
valuation $0$ at $u$. On the normal scheme $U'$, $f/t$ and $t/f$ are then regular, and $f$ generates
$\mathcal L$ on $U'$. So $f\wedge b_1\otimes(\wedge b_2)^{-1}=fs_0$ generates $\mathcal L$ on $\bar Z$,
and $f(y)=\kappa F(y)=\kappa\Phi_\sigma(y,v,w)$ at the complex points $y$ of $W$ over $\sigma$.
\end{proof}
% Lean: exists_gen_of_unitAn: frame_near_generic_point, wedge_lattice_generator and
% wedge_line_lattice_vector_bundle (s_0); exists_resolution_generic_fibre (Lemma
% cv-resolution-coordinates); generic_of_smooth_surjective (X'_Z); point_frame_specialization (F);
% extendsUnitAt_resolution (frame_at_component_point, pair_frames_on_set,
% affine_chart_bundle_sections, pullback_frames_complex_point); gaga_line_generator_embedding;
% proper_pullback_line_bundle = PushLoc.gen_descends.

\begin{lemma}\label{cv-function-generator}
In Situation~\ref{cv-boundary-situation}, let $s$ be a rational section of $\mathcal L$ that generates
$\mathcal L$ at every point of $\bar Z$.
\begin{enumerate}
\item For a framed point $(y,v,w)$ of $Z$ with values in $k$, the fibre of $\mathcal L$ at $y$ is
  $\wedge^p\operatorname{Lie}_\tau(y)\otimes(\wedge^q\operatorname{Lie}_{\bar\tau}(y))^{-1}$; let
  $\varphi_s(y,v,w)\in k^\times$ be the element with
  $s(y)=\varphi_s(y,v,w)\wedge v\otimes(\wedge w)^{-1}$. Then $\varphi_s$ is a frame function on $Z$
  (Definition~\ref{cv-function-framed}).
\item If $s'$ is another rational section of $\mathcal L$ that generates it at every point of
  $\bar Z$, then $\varphi_{s'}=x\varphi_s$ for some $x\in F_c^\times$.
\item There is $\kappa_\sigma\in\mathbf C^\times$ with $\varphi_s(y,v,w)=\kappa_\sigma\Phi_\sigma(y,v,w)$
  at every complex framed point $(y,v,w)$ of $Z$ over $\sigma$.
\item Let $(\mathcal X',\mathcal A',j')$ be a model of $X$ over a discrete valuation ring
  $O'\subset F_c$ with fraction field $F_c$, and let $s'$ be a rational section of
  $\det\operatorname{Lie}_\tau(\mathcal A')\otimes(\det\operatorname{Lie}_{\bar\tau}(\mathcal A'))^{-1}$
  that generates it at every point of the closure of $j'(Z)$ in $\mathcal X'_{F_c}$. Then
  $\varphi_{s'}=x\varphi_s$ for some $x\in F_c^\times$.
\end{enumerate}
\end{lemma}
\leanref{Development/AutomorphicSpectrum/Proofs/GeometryProofs/WedgeEigenLineGeneratorStatement.lean}{GlobalTheory.T.wedge_eigen_line_generator}
\leanref{Development/AutomorphicSpectrum/Definitions/GlobalGeometry/IsFlatAnIffIsConstMul.lean}{GlobalTheory.CompanionTheory.Ar.Flat.FlatOnModel}
\leanref{Development/AutomorphicSpectrum/Proofs/GeometryProofs/FrameFunctionsFlatAgreeStatement.lean}{GlobalTheory.T.frame_functions_flat_agree}
\leanref{Development/AutomorphicSpectrum/Proofs/GeometryProofs/FrameFunctionAnalyticNormStatement.lean}{GlobalTheory.T.frame_function_analytic_norm}
\leanref{Development/AutomorphicSpectrum/Proofs/GeometryProofs/FrameFunctionTwoModelsRelayStatement.lean}{GlobalTheory.T.frame_function_two_models_relay}
\begin{proof}
(1) By Lemma~\ref{cv-model-eigenbundles} the fibres of the two eigenbundles at $y$ are
$\operatorname{Lie}_\tau(y)$ and $\operatorname{Lie}_{\bar\tau}(y)$, and $s(y)\ne0$ because $s$
generates $\mathcal L$ near $y$. If $v'=gv$ and $w'=hw$, then
$\wedge v'\otimes(\wedge w')^{-1}=(\det g/\det h)\wedge v\otimes(\wedge w)^{-1}$, which gives the
transformation rule; the compatibility with extensions of fields holds because the fibre of
$\mathcal L$ is compatible with them.

(2) The quotient $s'/s$ is a unit at every point of $\bar Z$, so it is a global function on $\bar Z$,
which is integral and proper over $F_c$ as a closed subscheme of $\mathcal X_{F_c}$. Such a function
is algebraic over $F_c$, so it lies in $F_c$, which is algebraically closed in $F_c(Z)$. Put
$x=s'/s$.

(3) Let $b_1$ and $b_2$ be bases at the generic point, and let $f$, $\kappa$ and $W$ be given by
Lemma~\ref{cv-boundary-section}. For $s''=f\wedge b_1\otimes(\wedge b_2)^{-1}$, (1) gives
$\varphi_{s''}(y,v,w)=f(y)=\kappa\Phi_\sigma(y,v,w)$ when $y\in W$ and $v$ and $w$ are the values of
$b_1$ and $b_2$, and then at all complex framed points of $W$ over $\sigma$ by the transformation rules
(part (1) and Lemma~\ref{cv-periods-function}(1)). By (2), $\varphi_s=x\varphi_{s''}$ with
$x\in F_c^\times$, so $\varphi_s=\sigma(x)\kappa\Phi_\sigma$ on these framed points. Near a complex
point $y_0$ of $Z$ over $\sigma$, choose holomorphic coordinates on a connected neighbourhood and
sections of the two eigenbundles that are bases near $y_0$. In these bases $\varphi_s$ and
$\Phi_\sigma$ are holomorphic functions, by Lemma~\ref{cv-periods-function}(3), and $\Phi_\sigma$ has
no zeros. Their quotient equals $\sigma(x)\kappa$ off the proper Zariski closed subset $Z\setminus W$,
which is nowhere dense, so it equals $\sigma(x)\kappa$ near $y_0$. Put $\kappa_\sigma=\sigma(x)\kappa$.

(4) Let $\sigma_0$ be the inclusion $F_c\subset\mathbf C$; then $\sigma_0(\tau(\sqrt{-d}))=\theta$
because $\tau\in\Phi$. By (3) for both models at $\sigma_0$, $\varphi_{s'}/\varphi_s$ is a constant
$c\in\mathbf C^\times$ at the complex framed points of $Z$ over $\sigma_0$. Over $Z$ both line bundles
are $\det\operatorname{Lie}_\tau(A)\otimes(\det\operatorname{Lie}_{\bar\tau}(A))^{-1}$, so
$\varphi_{s'}/\varphi_s$ is the value of the rational function $r=s'/s\in F_c(Z)$, which is regular on
$Z$. Every embedding of $F_c(Z)$ into $\mathbf C$ over $\sigma_0$ is evaluation at a complex point of
$Z$, as in the proof of Lemma~\ref{cv-constant-sections-proportional}, so all of them send $r$ to $c$.
Hence $r$ is algebraic over $F_c$, since otherwise an embedding could send it to any transcendental
number, and $r=x\in F_c^\times$.
\end{proof}
% Lean: (1) wedge_eigen_line_generator (Represents, FlatOnModel.of_gen); (2) frame_functions_flat_agree;
% (3) frame_function_analytic_norm = flat_eq_const_mul_unitAn (eigen_frames_analytic and
% Cmp.ext_frames: component_open_stalks_domain, component_point_value_holomorphic,
% holomorphic_chart_near_point); (4) frame_function_two_models_relay = line_function_model_ratio
% (flat_eq_const_mul_unitAn at algebraMap c.F C for both models, frame_ratio_constant_rational).

\begin{lemma}\label{cv-norms-zariski}
Let $O$ be a discrete valuation ring with uniformizer $\varpi$ and fraction field $K$, and let
$\mathcal Z$ be a normal integral scheme, proper over $O$, whose generic point lies over the generic
point of $\operatorname{Spec}O$ and whose function field contains no element algebraic over $K$
outside $K$. Let $\mathcal M$ be a line bundle on $\mathcal Z$ and $s$ a rational section of
$\mathcal M$ that generates $\mathcal M$ at every point of the generic fibre. Suppose that the
pullback of $\mathcal M$ to every proper normal curve over a field that maps into the special fibre of
$\mathcal Z$ has degree $0$. Then there are integers $M>0$ and $a$ such that $s^M\varpi^{-a}$ generates
$\mathcal M^{\otimes M}$ at every point of $\mathcal Z$.
\end{lemma}
\leanref{Development/AutomorphicSpectrum/Proofs/GeometryProofs/LineDegreeSpecialFibreStatement.lean}{GlobalTheory.T.line_degree_special_fibre}
\leanref{Development/AutomorphicSpectrum/Proofs/GeometryProofs/IsConnectedZeroSetOfIsProperStatement.lean}{GlobalTheory.T.isConnected_zeroSet_uniformizer}
\begin{proof}
The global functions on $\mathcal Z$ form a finite $O$-algebra inside its function field, integral
over $O$, hence equal to $O$; by Zariski's connectedness theorem \cite[Corollaire~4.3.2]{Gro61} the
special fibre $\mathcal Z_0$ is connected. Since $\mathcal Z$ is normal and $s$ generates
$\mathcal M$ on the generic fibre, the divisor of $s$ is $\sum_in_iD_i$ over the irreducible
components $D_i$ of $\mathcal Z_0$, and the divisor of $\varpi$ is $\sum_im_iD_i$ with $m_i>0$. Let
$a/M$, with $M>0$, be the least of the $n_i/m_i$. Then $t=s^M\varpi^{-a}$ has an effective divisor, so
it is a regular section of $\mathcal M^{\otimes M}$, and it generates $\mathcal M^{\otimes M}$ at the
generic point of a component $D_i$ with $Mn_i=am_i$.

Let $V$ be the open set on which $t$ generates $\mathcal M^{\otimes M}$; it contains the generic fibre
and meets $\mathcal Z_0$. Suppose $\mathcal Z_0\not\subset V$. Since $\mathcal Z_0$ is connected,
$V\cap\mathcal Z_0$ is not closed in it, so some point $y\in V\cap\mathcal Z_0$ specializes to a point
outside $V$, and then to a closed point $x$ of $\mathcal Z_0$ outside $V$. In the closure of $y$ there
is an integral curve through $x$ whose generic point lies in $V$: prime avoidance in the local ring
of the closure of $y$ at $x$ gives a prime of dimension one that does not contain the ideal of the
complement of $V$. On the normalization of this curve, which is proper and normal and maps into
$\mathcal Z_0$, the pullback of $t$ is a regular section of the pullback of $\mathcal M^{\otimes M}$,
nonzero at the generic point and zero at a point over $x$, so that pullback has positive degree. This
contradicts the hypothesis. So $t$ generates $\mathcal M^{\otimes M}$ at every point of $\mathcal Z_0$,
and also on the generic fibre, where $s$ generates $\mathcal M$ and $\varpi$ is a unit.
\end{proof}
% Lean: propZ_numerical = PositivityTheory.Xgn.normal_proper_line_order (exists_regular_section_generator,
% subset_of_specializes, special_fibre_closed_specialization, degree_zero_curve_positive); the
% connectedness is isConnected_zeroSet_of_isProper. The curves are FibreCurve and the degree is
% FibreCurve.HasDeg.

\subsection{Formal charts and germs}\label{subsection-cv-formal-charts-germs}

\begin{definition}\label{cv-hodge-lattice-germ}
Let $E$ be a CM field, with complex conjugation $b\mapsto\bar b$ on $\mathcal O_E$, let $R$ be a
commutative ring and let $n\ge0$. The letters $x,y,z,w$ stand for families of $n$ variables, into
which $n$-tuples of power series without constant term are substituted, and $M^{\mathsf t}$ is the
transpose of a matrix $M$.
\begin{enumerate}
\item A \emph{formal germ} of dimension $n$ over $R$ is a tuple $B=(F,[\,\cdot\,],\Phi_1,\Phi_2,S,P)$
  of: a commutative formal group law $F$ of dimension $n$ over $R$, whose sum is written $x+_Fy$; for
  each $a\in\mathcal O_E$ an $n$-tuple $[a]x$ of power series over $R$ without constant term, with
  $[a](x+_Fy)=[a]x+_F[a]y$, $[a+b]x=[a]x+_F[b]x$, $[ab]x=[a]([b]x)$ and $[1]x=x$; and power series
  $\Phi_1(x,y,z)$, $\Phi_2(x,y,z)$, $S(x,y)$ and $P_a(x,y)$, for $a\in\mathcal O_E$, over $R$, each
  equal to $1$ when one of its arguments is $0$, such that for all $a,b\in\mathcal O_E$
  \begin{align*}
   &\Phi_1(x,y,w)\Phi_1(x+_Fy,z,w)=\Phi_1(y,z,w)\Phi_1(x,y+_Fz,w),\qquad\Phi_1(y,x,z)=\Phi_1(x,y,z),\\
   &\Phi_2(x,y,z)\Phi_2(x,y+_Fz,w)=\Phi_2(x,z,w)\Phi_2(x,y,z+_Fw),\qquad\Phi_2(x,z,y)=\Phi_2(x,y,z),\\
   &\Phi_1(x,y,z)\Phi_1(x,y,w)\Phi_2(x+_Fy,z,w)=\Phi_2(x,z,w)\Phi_2(y,z,w)\Phi_1(x,y,z+_Fw),\\
   &S(x,y)S(y,x)=1,\qquad\Phi_2(x,y,z)S(x,y)S(x,z)=S(x,y+_Fz)\Phi_1(y,z,x),\\
   &\Phi_1([a]x,[a]y,z)P_a(x+_Fy,z)=P_a(x,z)P_a(y,z)\Phi_1(x,y,[\bar a]z),\\
   &\Phi_2([a]x,y,z)P_a(x,y+_Fz)=P_a(x,y)P_a(x,z)\Phi_2(x,[\bar a]y,[\bar a]z),\\
   &P_1=1,\qquad P_{ab}(x,y)=P_a([b]x,y)P_b(x,[\bar a]y),\\
   &\Phi_1([a]x,[b]x,y)P_{a+b}(x,y)=P_a(x,y)P_b(x,y)\Phi_2(x,[\bar a]y,[\bar b]y),\\
   &P_a(x,y)P_{\bar a}(y,x)S([a]x,y)=S(x,[\bar a]y).
  \end{align*}
\item For a formal germ $B$, let $\rho_B(a)$ be the $n\times n$ matrix whose $(i,j)$ entry is the
  coefficient of $x_j$ in the $i$-th component of $[a]x$, and let $\beta_a$ and $S_{11}$ be the
  matrices whose $(i,j)$ entries are the coefficients of $x_iy_j$ in $P_a$ and in $S$. The
  \emph{Hodge module} of $B$ is $H(B)=R^n\oplus R^n$ with the endomorphisms
  \[
   \rho^H_B(a)(x,y)=\big(\rho_B(a)x,\ \rho_B(\bar a)^{\mathsf t}y-\beta_a^{\mathsf t}x\big)
   \qquad(a\in\mathcal O_E)
  \]
  and the pairing $\langle(x,y),(x',y')\rangle_B=y'\cdot x-y\cdot x'+x\cdot S_{11}x'$.
\item For a ring homomorphism $f:R\to R'$, the formal germ $B\otimes R'$ over $R'$ has the images
  under $f$ of the coefficients of $B$.
\item Let $B'$ be another formal germ of dimension $n$ over $R$, with its data written with primes.
  An \emph{isomorphism} from $B'$ to $B$ is a triple $(\phi,\psi,\lambda)$ of $n$-tuples $\phi$ and
  $\psi$ of power series over $R$ without constant term, with $\phi(\psi(x))=x=\psi(\phi(x))$, and a
  power series $\lambda(x,y)$ over $R$ equal to $1$ when $x=0$ or $y=0$, such that for all
  $a\in\mathcal O_E$
  \begin{align*}
   &\phi(x+_{F'}y)=\phi(x)+_F\phi(y),\qquad\phi([a]'x)=[a]\phi(x),\\
   &\Phi_1'(x,y,z)\lambda(x+_{F'}y,z)=\lambda(x,z)\lambda(y,z)\Phi_1(\phi(x),\phi(y),\phi(z)),\\
   &\Phi_2'(x,y,z)\lambda(x,y+_{F'}z)=\lambda(x,y)\lambda(x,z)\Phi_2(\phi(x),\phi(y),\phi(z)),\\
   &S'(x,y)\lambda(x,y)=S(\phi(x),\phi(y))\lambda(y,x),\qquad
   P_a'(x,y)\lambda(x,[\bar a]'y)=\lambda([a]'x,y)P_a(\phi(x),\phi(y)).
  \end{align*}
  Its \emph{transport map} is the $R$-linear map $T_{\phi,\psi,\lambda}:H(B')\to H(B)$,
  $(x,y)\mapsto\big(\Phi x,\Psi^{\mathsf t}(y+\Lambda^{\mathsf t}x)\big)$, where $\Phi$ and $\Psi$ are
  the matrices of the linear coefficients of $\phi$ and $\psi$, and $\Lambda$ is the matrix of the
  coefficients of $x_iy_j$ in $\lambda$.
\item Suppose that $\ell=0$ in $R$ for a prime $\ell$. The \emph{conjugate maps} of $B$ are the
  $R$-linear maps $\alpha_B:R^n\to H(B)$ and $\beta_B:H(B)\to R^n$,
  \[
   \alpha_B(u)=(Wu,D^{\mathsf t}u),\qquad\beta_B(x,y)=W^{\mathsf t}y-(D-W^{\mathsf t}S_{11})x,
  \]
  where the $(i,j)$ entry of $W$ is the coefficient of $x_j^\ell$ in the $i$-th component of
  $[\ell]x$, and that of $D$ is the coefficient of $x_i^\ell y_j$ in
  $\prod_{m=0}^{\ell-1}\Phi_1([m]x,x,y)$.
\end{enumerate}
\end{definition}
\leanref{ReciprocityTheory/ReductiveGroups/ArithmeticGeometry/BirationalGeometry/ToroidalCompactification/CocycleRightZeroThd.lean}{ReciprocityTheory.ChartRestriction.G.Geo.Tor.Germ.FB}
\leanref{Development/AutomorphicSpectrum/Definitions/CommutativeAlgebra/GermComp1ZeroPointLeft.lean}{GlobalTheory.CompanionTheory.Ar.Conj.rhoH}
\leanref{Development/AutomorphicSpectrum/Definitions/CommutativeAlgebra/GermComp1ZeroPointLeft.lean}{GlobalTheory.CompanionTheory.Ar.Conj.pairH}
\leanref{Development/AutomorphicSpectrum/Definitions/CommutativeAlgebra/GermComp1ZeroPointLeft.lean}{GlobalTheory.CompanionTheory.Ar.Conj.IsMapOf}
\leanref{Development/AutomorphicSpectrum/Definitions/CommutativeAlgebra/GermComp1ZeroPointLeft.lean}{GlobalTheory.CompanionTheory.Ar.Conj.IsIso}
\leanref{Development/AutomorphicSpectrum/Definitions/CommutativeAlgebra/GermComp1ZeroPointLeft.lean}{GlobalTheory.CompanionTheory.Ar.Conj.transH}
\leanref{Development/AutomorphicSpectrum/Definitions/CommutativeAlgebra/GermComp1ZeroPointLeft.lean}{GlobalTheory.CompanionTheory.Ar.Conj.alpha}
\leanref{Development/AutomorphicSpectrum/Definitions/CommutativeAlgebra/GermComp1ZeroPointLeft.lean}{GlobalTheory.CompanionTheory.Ar.Conj.beta}

\begin{lemma}\label{cv-hodge-lattice-transport}
Let $E$ be a CM field, let $R$ be a commutative ring, and let $B$ be a formal germ of dimension $n$
over $R$ (Definition~\ref{cv-hodge-lattice-germ}).
\begin{enumerate}
\item The map $a\mapsto\rho^H_B(a)$ is a ring homomorphism from $\mathcal O_E$ to
  $\operatorname{End}_RH(B)$, $S_{11}^{\mathsf t}=-S_{11}$, and
  $\langle\rho^H_B(a)v,v'\rangle_B=\langle v,\rho^H_B(\bar a)v'\rangle_B$ for $v,v'\in H(B)$.
\item Let $(\phi,\psi,\lambda)$ be an isomorphism from a formal germ $B'$ of dimension $n$ over $R$ to
  $B$, with transport map $T$. Then $\Phi\Psi=\Psi\Phi=1$ and $\rho_B(a)\Phi=\Phi\rho_{B'}(a)$, the
  map $T$ is bijective, $\rho^H_B(a)\circ T=T\circ\rho^H_{B'}(a)$, and
  $\langle Tv,Tv'\rangle_B=\langle v,v'\rangle_{B'}$.
\item Suppose that $\ell=0$ in $R$ for a prime $\ell$, and let $M^{(\ell)}$ be the matrix of the
  $\ell$-th powers of the entries of a matrix $M$ over $R$. Then $u\cdot\beta_B(v)=\langle\alpha_B(u),
  v\rangle_B$, $\beta_B(\alpha_B(u))=0$, $\alpha_B(\rho_B(a)^{(\ell)}u)=\rho^H_B(a)\alpha_B(u)$ and
  $\beta_B(\rho^H_B(a)v)=(\rho_B(\bar a)^{\mathsf t})^{(\ell)}\beta_B(v)$. In the situation of (2),
  $\alpha_B(\Phi^{(\ell)}u)=T(\alpha_{B'}(u))$ and $\beta_B(Tv)=(\Psi^{\mathsf t})^{(\ell)}\beta_{B'}(v)$.
\end{enumerate}
\end{lemma}
\leanref{Development/AutomorphicSpectrum/Proofs/AlgebraProofs/FormalHermitianActionPairingStatement.lean}{GlobalTheory.T.formal_hermitian_action_pairing}
\leanref{Development/AutomorphicSpectrum/Proofs/AlgebraProofs/FormalIsomorphismTransportStatement.lean}{GlobalTheory.T.formal_isomorphism_transport}
\leanref{Development/AutomorphicSpectrum/Proofs/AlgebraProofs/AlphaBetaAdjointExactStatement.lean}{GlobalTheory.T.alpha_beta_adjoint_exact}
\leanref{Development/AutomorphicSpectrum/Proofs/AlgebraProofs/AlphaEquivariantFrobeniusStatement.lean}{GlobalTheory.T.alpha_equivariant_frobenius}
\leanref{Development/AutomorphicSpectrum/Proofs/AlgebraProofs/BetaEquivariantFrobeniusStatement.lean}{GlobalTheory.T.beta_equivariant_frobenius}
\leanref{Development/AutomorphicSpectrum/Proofs/AlgebraProofs/AlphaBetaIsomorphismTransportStatement.lean}{GlobalTheory.T.alpha_beta_isomorphism_transport}
\begin{proof}
Each of $\Phi_1$, $\Phi_2$, $S$, $P_a$ and $\lambda$ is $1$ when one of its arguments is $0$, so a
product of such series has as terms of bidegree $(1,1)$ in two families of variables the sum of
those of its factors.

(1) Since $F(x,y)$ is $x+y$ plus terms of degree at least $2$, the linear parts of the identities for
$[a+b]$, $[ab]$ and $[1]$ show that $\rho_B$ is a ring homomorphism. In the identity for $P_{a+b}$ the
factors $\Phi_1([a]x,[b]x,y)$ and $\Phi_2(x,[\bar a]y,[\bar b]y)$ have no terms of bidegree $(1,1)$
in $(x,y)$, so that identity gives $\beta_{a+b}=\beta_a+\beta_b$. The identities for $P_{ab}$ and
$P_1$ give $\beta_{ab}=\rho_B(b)^{\mathsf t}\beta_a+\beta_b\rho_B(\bar a)$ and $\beta_1=0$. With
$\rho_B(\overline{ab})=\rho_B(\bar a)\rho_B(\bar b)$ these say that $\rho^H_B$ is additive,
multiplicative and unital. The terms of bidegree $(1,1)$ of $S(x,y)S(y,x)=1$ give
$S_{11}+S_{11}^{\mathsf t}=0$. Those of the last identity of a formal germ give
$\beta_a+\beta_{\bar a}^{\mathsf t}+\rho_B(a)^{\mathsf t}S_{11}=S_{11}\rho_B(\bar a)$, which is the
adjointness.

(2) The linear parts of $\phi(\psi(x))=x=\psi(\phi(x))$ and of $\phi([a]'x)=[a]\phi(x)$ give
$\Phi\Psi=\Psi\Phi=1$ and $\rho_B(a)\Phi=\Phi\rho_{B'}(a)$. The terms of bidegree $(1,1)$ of the
conditions on $S'$ and on $P'_a$ give
\[
 S'_{11}+\Lambda=\Phi^{\mathsf t}S_{11}\Phi+\Lambda^{\mathsf t},\qquad
 \beta'_a+\Lambda\rho_{B'}(\bar a)=\rho_{B'}(a)^{\mathsf t}\Lambda+\Phi^{\mathsf t}\beta_a\Phi.
\]
The map $(x,y)\mapsto(\Psi x,\Phi^{\mathsf t}y-\Lambda^{\mathsf t}\Psi x)$ is inverse to $T$. The first
components of $\rho^H_B(a)T(x,y)$ and $T\rho^H_{B'}(a)(x,y)$ are $\rho_B(a)\Phi x=\Phi\rho_{B'}(a)x$.
Their second components agree by the second identity displayed, transposed and multiplied by
$\Psi^{\mathsf t}$, together with $\Psi\rho_B(\bar a)=\rho_{B'}(\bar a)\Psi$. Expanding
$\langle T(x,y),T(x',y')\rangle_B$ with $\Psi\Phi=1$ gives
$y'\cdot x-y\cdot x'+x\cdot(\Lambda^{\mathsf t}-\Lambda+\Phi^{\mathsf t}S_{11}\Phi)x'$, which is
$\langle(x,y),(x',y')\rangle_{B'}$ by the first identity displayed.

(3) The first formula follows from the definitions:
$\langle\alpha_B(u),(x,y)\rangle_B=y\cdot Wu-D^{\mathsf t}u\cdot x+Wu\cdot S_{11}x=u\cdot\beta_B(x,y)$.
The others are computations with the coefficients of degree $\ell$ in the identities of the formal
germs and of the isomorphism, in which $\ell=0$ in $R$ is used; only their results are given here.
\end{proof}
% Lean: (1) rhoH_hom_adjoint; (2) Ar.Conj.Iso.rho_pairing; (3) Isotropy.beta_comp_alpha,
% AlphaAct.alpha_comm_act, beta_equivariant_frobenius, alpha_beta_of_isIso.

\begin{definition}\label{cv-germs-chart}
Let $E$ be a CM field, let $G$ be an $\mathcal O_E$-semi-abelian scheme over a scheme $X$
(Definition~\ref{cv-model-toroidal-semiabelian}) with line bundle $\mathcal L$, let $R$ be a
commutative ring, let $x_0:\operatorname{Spec}R\to X$ be a morphism, and let $d\ge0$.
\begin{enumerate}
\item For points $y_1,y_2,y_3$ of $G$ with values in a scheme $T$ over $X$, let
  \[
   \Lambda(y_1,y_2)=(y_1+y_2)^*\mathcal L\otimes y_1^*\mathcal L^{-1}\otimes y_2^*\mathcal L^{-1}
   \otimes0^*\mathcal L,\qquad
   \Theta(y_1,y_2,y_3)=\bigotimes_{K\subset\{1,2,3\}}\Big(\sum_{i\in K}y_i\Big)^*\mathcal L^{(-1)^{3-|K|}},
  \]
  line bundles on $T$. There are canonical isomorphisms $\Lambda(y_1,y_2)\cong\Lambda(y_2,y_1)$,
  $\Lambda(0,y_2)\cong\mathcal O_T\cong\Lambda(y_1,0)$ and
  \[
   \Theta(y_1,y_2,y_3)\cong\Lambda(y_1+y_2,y_3)\otimes\Lambda(y_1,y_3)^{-1}\otimes\Lambda(y_2,y_3)^{-1}
   \cong\Lambda(y_1,y_2+y_3)\otimes\Lambda(y_1,y_2)^{-1}\otimes\Lambda(y_1,y_3)^{-1},
  \]
  through which generators of these line bundles are compared below. The cubical structure of
  $\mathcal L$ is a generator of $\Theta$ for the three projections of $G\times_XG\times_XG$, and
  $\theta(y_1,y_2,y_3)$ is its pullback by $(y_1,y_2,y_3)$. It is $1$ when one of its arguments is
  $0$, it is invariant under permutations of its arguments, and
  $\theta(y_1+y_2,y_3,y_4)\,\theta(y_1,y_2,y_4)=\theta(y_1,y_2+y_3,y_4)\,\theta(y_2,y_3,y_4)$
  \cite[Chapitre~I]{MB85}.
\item For $k,m\ge0$ let $J^k_m=R[t^{(1)},\dots,t^{(k)}]/I^{m+1}$, where $t^{(1)},\dots,t^{(k)}$ are
  families of $d$ variables and $I$ is the ideal that they generate, so that
  $\operatorname{Spec}J^k_m$ is a scheme over $X$ through $x_0$; write $t,t',t''$ for
  $t^{(1)},t^{(2)},t^{(3)}$. For a point $x$ of $G$ with values in $J^1_m$ and a $d$-tuple $u$ of power
  series over $R$ without constant term in the variables $t^{(i)}$, $x(u)$ is the point with values in
  $J^k_m$ obtained by substituting $u$ for $t$. Generators of line bundles on
  $\operatorname{Spec}J^2_m$ are evaluated at pairs $(u,u')$ in the same way, and power series are read
  modulo $I^{m+1}$. A \emph{tangent vector} at $x_0$ is a point of $G$ with values in $R[\varepsilon]$,
  $\varepsilon^2=0$, over $x_0$ that is the unit modulo $\varepsilon$. The tangent vectors form an
  $R$-module, with the sum given by the group law and $c\in R$ acting through
  $\varepsilon\mapsto c\varepsilon$.
\item A \emph{formal chart} of dimension $d$ at $x_0$ consists of a formal germ
  $B=(F,[\,\cdot\,],\Phi_1,\Phi_2,S,P)$ of dimension $d$ over $R$
  (Definition~\ref{cv-hodge-lattice-germ}), points $x_m$ of $G$ with values in $J^1_m$ and generators
  $e_m$ of $\Lambda(x_m(t),x_m(t'))$ on $\operatorname{Spec}J^2_m$, for $m\ge0$, and tangent vectors
  $v_1,\dots,v_d$ at $x_0$, such that:
  \begin{enumerate}
  \item $x_m(0)$ is the unit, $x_{m+1}$ reduces to $x_m$ modulo $I^{m+1}$,
    $x_m(t+_Ft')=x_m(t)+x_m(t')$, and $a\,x_m(t)=x_m([a]t)$ for $a\in\mathcal O_E$;
  \item $v_l$ is the image of $x_1$ under $t_j\mapsto\delta_{jl}\varepsilon$, and every tangent vector
    at $x_0$ is $\sum_lc_lv_l$ for a unique $c\in R^d$;
  \item $e_{m+1}$ reduces to $e_m$, and $e_m(0,t)=1=e_m(t,0)$;
  \item for every $m$, on $\operatorname{Spec}J^3_m$ and on $\operatorname{Spec}J^2_m$,
    \begin{align*}
     \theta(x_m(t),x_m(t'),x_m(t''))&=\Phi_1(t,t',t'')\,e_m(t+_Ft',t'')\,e_m(t,t'')^{-1}e_m(t',t'')^{-1},\\
     \theta(x_m(t'),x_m(t''),x_m(t))&=\Phi_2(t,t',t'')\,e_m(t,t'+_Ft'')\,e_m(t,t')^{-1}e_m(t,t'')^{-1},\\
     e_m(t,t')&=S(t,t')^{-1}e_m(t',t);
    \end{align*}
  \item for every $a\in\mathcal O_E$ there is a generator $w_a$ of
    $\Lambda(y,\bar ay')\otimes\Lambda(ay,y')^{-1}$, where $y$ and $y'$ are the two projections of
    $G\times_XG$, such that $w_a(x_m(t),x_m(t'))\,e_m([a]t,t')=P_a(t,t')\,e_m(t,[\bar a]t')$ for every
    $m$.
  \end{enumerate}
\end{enumerate}
\end{definition}
\leanref{Development/AutomorphicSpectrum/Definitions/GlobalGeometry/ToFirst.lean}{GlobalTheory.CompanionTheory.Ar.Conj.GermAt}
\leanref{Formalism/LocalToGlobal/ProjectiveLattices/IsChartIff.lean}{Formalism.L2.Pj.IsChartG}
\leanref{Formalism/LocalToGlobal/ProjectiveLattices/IsChartIff.lean}{Formalism.L2.Pj.IsFrameG}
\leanref{Formalism/SharpRestriction/Divisibility/GeneratorPowers/GermCut.lean}{Formalism.MainNotions.ModuliIntegralCharts.H1.FrTrunc}
\leanref{Formalism/SharpRestriction/Divisibility/GeneratorPowers/GermCut.lean}{Formalism.MainNotions.ModuliIntegralCharts.H1.FrAxes}
\leanref{ReciprocityTheory/ReductiveGroups/ArithmeticGeometry/SemiabelianSchemes/CubicalStructures/CubeSymmetricSwapFirst.lean}{ReciprocityTheory.ChartRestriction.G.Semi.cubeLB}
\leanref{ReciprocityTheory/ReductiveGroups/ArithmeticGeometry/SemiabelianSchemes/CubicalStructures/CubeSymmetricSwapFirst.lean}{ReciprocityTheory.ChartRestriction.G.Semi.CubicalStr}

\begin{theorem}\label{cv-germs-formal}
Let $R$ be a commutative ring, let $G$ be a smooth commutative group scheme over
$\operatorname{Spec}R$, let $d\ge0$, and let the rings $J^1_m$ and the tangent vectors of $G$ at the
identity of $\operatorname{Spec}R$ be as in Definition~\ref{cv-germs-chart}.
\begin{enumerate}
\item Let $v_1,\dots,v_d$ be tangent vectors such that every tangent vector is $\sum_lc_lv_l$ for a
  unique $c\in R^d$. Then there are points $x_m$ of $G$ with values in $J^1_m$, for $m\ge0$, such that
  $x_m(0)$ is the unit, $x_{m+1}$ reduces to $x_m$, and $v_l$ is the image of $x_1$ under
  $t_j\mapsto\delta_{jl}\varepsilon$ for every $l$.
\item Let $x_m$, for $m\ge0$, be points as in (1) for some tangent vectors $v_1,\dots,v_d$ as in (1).
  Let $m\ge0$, let $C$ be an $R$-algebra with an ideal $N$ such that $N^{m+1}=0$, and let $y$ be a point
  of $G$ with values in $C$ that is the unit modulo $N$. Then $y=x_m\circ\operatorname{Spec}\chi$ for a
  unique homomorphism of $R$-algebras $\chi:J^1_m\to C$.
\end{enumerate}
\end{theorem}
\leanref{Development/AbelianVarietyTheory/Proofs/CuspFormEvaluations/ExistsIsChartGOfBasis.lean}{GlobalTheory.ChartData.exists_isChartG_of_basis}
\leanref{Development/AbelianVarietyTheory/Proofs/UnitaryExistence/DifBijectiveOfParam.lean}{GlobalTheory.ChartCompare.param_of_dif_bijective}

The points $x_m$ are coordinates of $G$ along its unit section. The unit section of the smooth group
scheme $G$ is a quasi-regular immersion \cite[Th\'eor\`eme~17.12.1]{Gro67}, and the completion of $G$
along it is therefore, locally on $\operatorname{Spec}R$, an algebra of formal power series over the base,
augmented by the unit section \cite[Corollaire~16.9.9]{Gro67}.

\begin{lemma}\label{cv-germs-local}
Let $E$ be a CM field, and let $G$ be an $\mathcal O_E$-semi-abelian scheme over a scheme $X$
(Definition~\ref{cv-model-toroidal-semiabelian}) such that, for every $m\ge1$, every function on the
$m$-fold fibre product $G^m$ of $G$ over $X$ is the pullback of a function on $X$. Then every point
$z$ of $X$ has an affine open neighbourhood $\operatorname{Spec}A\subset X$ with a formal chart
(Definition~\ref{cv-germs-chart}) at the inclusion $\operatorname{Spec}A\to X$.
\end{lemma}
\leanref{Development/WeilBound/CoreStatements/TorusUniformisation/ChowMahlerFamilies/EqMulInversePow.lean}{PositivityTheory.Xgn.GermSystem.germ_aff}
\leanref{Development/WeilBound/CoreStatements/HodgeTatePoints/CdChart.lean}{PositivityTheory.HtPt.Cd.chart_affine}
\leanref{Development/WeilBound/CoreStatements/HodgeTatePoints/CdFrame.lean}{PositivityTheory.HtPt.Cd.frame_affine}
\leanref{Formalism/SharpRestriction/Divisibility/GeneratorPowers/FBIdAd.lean}{Formalism.MainNotions.ModuliIntegralCharts.H1.FBIdAd.fbId}
\begin{proof}
Coordinates. The line bundle $\mathcal L$ is trivial on an open subset $U$ of $G$ that contains the
unit over $z$. Let $\operatorname{Spec}A'\subset X$ be an affine open neighbourhood of $z$ whose image
under the unit section lies in $U$. The group scheme $G$ is smooth, so its module of tangent vectors over
$A'$ is finitely generated and projective, and its formation commutes with base change. So it is free
over $A=A'_r$ for some $r\in A'$ outside the prime ideal of $z$; let $v_1,\dots,v_d$ be a basis.
Theorem~\ref{cv-germs-formal}(1), applied to the pullback of $G$ to $\operatorname{Spec}A$, gives points
$x_m$ with values in $J^1_m$ such that $x_m(0)$ is the unit, $x_{m+1}$ reduces to $x_m$, and $v_l$ is
the image of $x_1$ under $t_j\mapsto\delta_{jl}\varepsilon$.

The series $F$ and $[a]$. The point $x_m(t)+x_m(t')$ with values in $J^2_m$ is the unit modulo the ideal
$I$ generated by $t$ and $t'$, and $I^{m+1}=0$ in $J^2_m$. By Theorem~\ref{cv-germs-formal}(2) it is
$x_m\circ\operatorname{Spec}\chi_m$ for a unique homomorphism $\chi_m:J^1_m\to J^2_m$. By the uniqueness,
the $\chi_m(t)$ reduce to one another as $m$ grows, so they are the reductions of a $d$-tuple $F(t,t')$
of power series, and its constant terms vanish by the uniqueness applied modulo $I$. In the same way
$a\,x_m(t)=x_m([a]t)$ defines $[a]t$. By the uniqueness in Theorem~\ref{cv-germs-formal}(2), a $d$-tuple
$u$ of power series without constant term is determined by the points $x_m(u)$, $m\ge0$. So $F$ is a
commutative formal group law, and the $[a]$ satisfy the identities of
Definition~\ref{cv-hodge-lattice-germ}(1) that involve only $F$ and $[\,\cdot\,]$: the two sides of each
give the same points, for example $x_m((t+_Ft')+_Ft'')=x_m(t)+x_m(t')+x_m(t'')=x_m(t+_F(t'+_Ft''))$ and
$x_m([ab]t)=a\,b\,x_m(t)=x_m([a]([b]t))$.

Trivializations. Each point $x_m(u)$, and each sum of such points, is the unit modulo a nilpotent ideal,
so it maps $\operatorname{Spec}J^k_m$ into $U$. The trivialization of $\mathcal L$ on $U$ therefore
trivializes $\Lambda$ and $\Theta$ at these points, and it gives generators $e_m$ of
$\Lambda(x_m(t),x_m(t'))$ that reduce to one another and are $1$ where $t=0$ or $t'=0$. As complex
conjugation is the Rosati involution, $\Lambda(ay,y')\cong\Lambda(y,\bar ay')$ on $G\times_XG$; let $w_a$
be a generator of $\Lambda(y,\bar ay')\otimes\Lambda(ay,y')^{-1}$, divided by its value on the unit
section, which is a unit on $X$, so that $w_a(0,0)=1$. In each relation of
Definition~\ref{cv-germs-chart}(3)(d) and~(e), the two sides without the power series are generators of
the same line bundle on $\operatorname{Spec}J^k_m$, so their ratio is a unit of $J^k_m$. These units
reduce to one another as $m$ grows, so they are the reductions of power series $\Phi_1$, $\Phi_2$, $S$
and $P_a$ over $A$. Where $t=t'=0$ every point involved is the unit and every generator is $1$, so $P_a$
has constant term $1$.

Identities. It remains to show that $B=(F,[\,\cdot\,],\Phi_1,\Phi_2,S,P)$ is a formal germ. A power series
over $A$ is determined by its reductions modulo $I^{m+1}$ for all $m$, so each identity of
Definition~\ref{cv-hodge-lattice-germ}(1) may be checked in $J^k_m$, with $k\le4$, for every $m$.
Substituting into the relations of Definition~\ref{cv-germs-chart}(3)(d) and~(e) writes each side as a
product of values of $\theta$, of $e_m$ and of the $w_a$ at points $x_m(u)$, which are added and
multiplied by elements of $\mathcal O_E$ through $x_m(u+_Fu')=x_m(u)+x_m(u')$ and
$x_m([a]u)=a\,x_m(u)$. The factors $e_m$ cancel, and the identity becomes one between values of $\theta$
and of the $w_a$. For the identities that involve only $\Phi_1$, $\Phi_2$ and $S$, these follow from the
conditions on $\theta$ in Definition~\ref{cv-germs-chart}(1) and from $e_m=1$ on the axes; they are the
axioms of the symmetric biextension of $G\times_XG$ by the multiplicative group that the cubical
structure defines on $\Lambda$, written in the generators $e_m$ \cite[Chapitre~I]{MB85},
\cite[Expos\'e~VII]{Gro72}. For the identities that involve $P$, they are the following, for points
$y,y',y''$ of $G$ with values in a scheme over $X$ and for $a,b\in\mathcal O_E$:
\begin{align*}
 &w_a(y+y',y'')\,\theta(ay,ay',y'')=w_a(y,y'')\,w_a(y',y'')\,\theta(y,y',\bar ay''),\\
 &w_a(y,y'+y'')\,\theta(ay,y',y'')=w_a(y,y')\,w_a(y,y'')\,\theta(y,\bar ay',\bar ay''),\\
 &w_{ab}(y,y')=w_a(by,y')\,w_b(y,\bar ay'),\qquad
  w_{a+b}(y,y')\,\theta(ay,by,y')=w_a(y,y')\,w_b(y,y')\,\theta(y,\bar ay',\bar by'),\\
 &w_a(y,y')\,w_{\bar a}(y',y)\,w_1(y,y')=w_1(y,y').
\end{align*}
For the universal points, the projections of $G^3$, the two sides of each are generators of the same
line bundle, so their ratio is a unit on $G^3$. By hypothesis it is the pullback of a unit on $X$, and on
the unit section it is $1$, because $w_a(0,0)=1$ and $\theta$ is $1$ when one of its arguments is $0$.
So each of them holds, and with $a=b=1$ the third gives $w_1=1$.
\end{proof}

\begin{lemma}\label{cv-germs-compare}
Let $E$ be a CM field, and let $G$ be an $\mathcal O_E$-semi-abelian scheme over a scheme $X$
(Definition~\ref{cv-model-toroidal-semiabelian}) such that, for every $m\ge1$, every function on the
$m$-fold fibre product $G^m$ of $G$ over $X$ is the pullback of a function on $X$. Let $R$ be a
commutative ring, let $x_0:\operatorname{Spec}R\to X$ be a morphism, and let $d\ge0$. Formal charts are
those of Definition~\ref{cv-germs-chart}, and the data of a formal chart $g'$ are written with primes.
\begin{enumerate}
\item Let $f:R\to R'$ be a ring homomorphism and $g$ a formal chart of dimension $d$ at $x_0$. The
  images under $f$ of the data of $g$ form a formal chart $g\otimes R'$ at
  $x_0\circ\operatorname{Spec}f$, with formal germ $B\otimes R'$, and $(g\otimes R')\otimes R''=g\otimes
  R''$ for a ring homomorphism $R'\to R''$.
\item If $R$ is local, formal charts at $x_0$ of dimensions $d$ and $d'$ have $d=d'$.
\item Let $g$ and $g'$ be formal charts of dimension $d$ at $x_0$. There are $d$-tuples $\phi$ and
  $\psi$ of power series over $R$ without constant term and a power series $\lambda$ over $R$ in two
  families of $d$ variables such that, for every $m$,
  \[
   x'_m(t)=x_m(\phi(t)),\qquad x_m(t)=x'_m(\psi(t)),\qquad
   e'_m(t,t')=\lambda(t,t')\,e_m(\phi(t),\phi(t')).
  \]
  A triple $(\phi,\psi,\lambda)$ with these properties is a \emph{comparison} of $g'$ with $g$, and the
  $d$-tuples $\phi$ and $\psi$ are the same for all comparisons of $g'$ with $g$. Every comparison is an
  isomorphism from $B'$ to $B$ (Definition~\ref{cv-hodge-lattice-germ}), and $v'_l=\sum_j\Phi_{jl}v_j$
  for every $l$. For a ring homomorphism $f:R\to R'$, the image of a comparison of $g'$ with $g$ under
  $f$ is a comparison of $g'\otimes R'$ with $g\otimes R'$.
\item Let $g$, $g'$ and $g''$ be formal charts of dimension $d$ at $x_0$, and let
  $(\phi,\psi,\lambda)$, $(\phi',\psi',\lambda')$ and $(\phi'',\psi'',\lambda'')$ be comparisons of
  $g'$ with $g$, of $g''$ with $g'$ and of $g''$ with $g$. Then
  $T_{\phi'',\psi'',\lambda''}=T_{\phi,\psi,\lambda}\circ T_{\phi',\psi',\lambda'}$.
\end{enumerate}
\end{lemma}
\leanref{Development/WeilBound/CoreStatements/TorusUniformisation/ChowMahlerFamilies/EqMulInversePow.lean}{PositivityTheory.Xgn.GermSystem.exists_bc}
\leanref{Development/WeilBound/CoreStatements/HodgeTatePoints/CdUniv.lean}{PositivityTheory.HtPt.Cd.germ_univ}
\leanref{Development/WeilBound/CoreStatements/TorusUniformisation/ChowMahlerFamilies/EqMulInversePow.lean}{PositivityTheory.Xgn.GermSystem.exists_cmp}
\leanref{Development/WeilBound/CoreStatements/TorusUniformisation/ChowMahlerFamilies/ActGeneratorZero.lean}{PositivityTheory.Xgn.GermSystem.PP_inj}
\leanref{Development/WeilBound/CoreStatements/TorusUniformisation/ChowMahlerFamilies/EqMulInversePow.lean}{PositivityTheory.Xgn.GermSystem.Cmp.iso}
\leanref{Development/WeilBound/CoreStatements/TorusUniformisation/ChowMahlerFamilies/ActGeneratorZero.lean}{PositivityTheory.Xgn.GermSystem.coord_tan}
\leanref{Development/WeilBound/CoreStatements/TorusUniformisation/ChowMahlerFamilies/EqMulInversePow.lean}{PositivityTheory.Xgn.GermSystem.Cmp.trans}
\begin{proof}
(1) The relations of Definition~\ref{cv-germs-chart}(3) for $g$ pull back along the morphisms
$\operatorname{Spec}(J^k_m\otimes_RR')\to\operatorname{Spec}J^k_m$, and the generators $w_a$ serve again.
The module of tangent vectors of the smooth group scheme $G$ commutes with base change, so the images
of $v_1,\dots,v_d$ form a basis of the tangent vectors at $x_0\circ\operatorname{Spec}f$. Taking images
under $f$ and then under $R'\to R''$ is taking images under the composite.

(2) By Definition~\ref{cv-germs-chart}(3)(b), the tangent vectors at $x_0$ form a free $R$-module with a
basis of $d$ elements and one of $d'$ elements, so $d=d'$.

(3) Theorem~\ref{cv-germs-formal}(2), for the pullback of $G$ to $\operatorname{Spec}R$, applies to the
points $x_m$ of $g$, by Definition~\ref{cv-germs-chart}(3)(a) and~(b). Applied to $y=x'_m$ and the ideal generated by $t$ in
$J^1_m$, it gives $x'_m=x_m\circ\operatorname{Spec}\chi_m$ for a unique homomorphism
$\chi_m:J^1_m\to J^1_m$. By the uniqueness, the $\chi_m(t)$ reduce to one another, so they are the
reductions of a unique $d$-tuple $\phi$ of power series, without constant term by the uniqueness modulo
$t$. Exchanging $g$ and $g'$ gives $\psi$. The generators $e'_m(t,t')$ and $e_m(\phi(t),\phi(t'))$ of
$\Lambda(x'_m(t),x'_m(t'))$ differ by a unit of $J^2_m$, and these units reduce to one another, so they
are the reductions of a power series $\lambda$. If $(\tilde\phi,\tilde\psi,\tilde\lambda)$ is another
comparison, then $x_m(\tilde\phi(t))=x'_m(t)=x_m(\phi(t))$ for every $m$, so $\tilde\phi=\phi$ by the
uniqueness in Theorem~\ref{cv-germs-formal}(2) applied modulo $I^{m+1}$ for every $m$; likewise
$\tilde\psi=\psi$.

The isomorphism. By the uniqueness in Theorem~\ref{cv-germs-formal}(2), a $d$-tuple $u$ of power series
without constant term is determined by the points $x_m(u)$. As
$x_m(\phi(\psi(t)))=x'_m(\psi(t))=x_m(t)$, $\phi(\psi(x))=x$, and likewise $\psi(\phi(x))=x$. As
$x_m(\phi(t+_{F'}t'))=x'_m(t)+x'_m(t')=x_m(\phi(t)+_F\phi(t'))$ and
$x_m(\phi([a]'t))=a\,x'_m(t)=x_m([a]\phi(t))$, the map $\phi$ is compatible with the group laws and the
actions of $\mathcal O_E$. As $e_m$ and $e'_m$ are $1$ where one argument is $0$, so is $\lambda$. The
conditions on $\Phi'_1$, $\Phi'_2$, $S'$ and $P'_a$ are checked in $J^3_m$ or $J^2_m$ for every $m$:
substituting $x'_m=x_m\circ\phi$ and $e'_m=\lambda\,e_m(\phi,\phi)$ into the relations of
Definition~\ref{cv-germs-chart}(3)(d) and~(e) for $g'$ and comparing them with those for $g$ at the points
$x_m(\phi(\cdot))$, the values of $\theta$ and of $e_m$ cancel and the conditions remain. For the condition
on $P'_a$ this uses $w'_a=w_a$. Both are generators of the same line bundle on $G\times_XG$, and both are
$1$ on the unit section, by the relations of Definition~\ref{cv-germs-chart}(3)(e) at $t=t'=0$, as
$P_a(0,0)=1=P'_a(0,0)$ and the generators $e_m$ and $e'_m$ are $1$ there. So their ratio is a unit on
$G\times_XG$, by hypothesis the pullback of a unit on $X$, and that unit is $1$.

Tangent vectors and base change. For $c\in R^d$, $\phi(\varepsilon c)=\varepsilon\Phi c$ as
$\varepsilon^2=0$. The map $c\mapsto x_1(\varepsilon c)$ to the tangent vectors is additive, as
$x_1(\varepsilon c)+x_1(\varepsilon c')=x_1(\varepsilon c+_F\varepsilon c')=x_1(\varepsilon(c+c'))$, and
compatible with the action of $R$, so $x_1(\varepsilon c)=\sum_jc_jv_j$. Hence
$v'_l=x'_1(\varepsilon\delta_l)=x_1(\varepsilon\Phi\delta_l)=\sum_j\Phi_{jl}v_j$, where $\delta_l$ is the
$l$-th unit vector. The images of $\phi$, $\psi$ and $\lambda$ under $f$ satisfy the relations that define
a comparison of $g'\otimes R'$ with $g\otimes R'$, since the data of $g\otimes R'$ and $g'\otimes R'$ are
the images of those of $g$ and $g'$.

(4) By the uniqueness in (3), $x''_m(t)=x'_m(\phi'(t))=x_m(\phi(\phi'(t)))$ gives
$\phi''=\phi\circ\phi'$, and likewise $\psi''=\psi'\circ\psi$. In $J^2_m$, both $\lambda''(t,t')$ and
$\lambda'(t,t')\,\lambda(\phi'(t),\phi'(t'))$ are the ratio of the generators $e''_m(t,t')$ and
$e_m(\phi''(t),\phi''(t'))$ of one line bundle, by the relations for the three comparisons. A power
series is determined by its images in the rings $J^2_m$, so
$\lambda''(x,y)=\lambda'(x,y)\,\lambda(\phi'(x),\phi'(y))$. Since $\lambda$ and $\lambda'$ are $1$ when one
argument is $0$, $\Phi''=\Phi\Phi'$, $\Psi''=\Psi'\Psi$ and
$\Lambda''=\Lambda'+\Phi'^{\mathsf t}\Lambda\Phi'$. Then
\[
 T_{\phi'',\psi'',\lambda''}(x,y)=\big(\Phi\Phi'x,\ \Psi^{\mathsf t}\Psi'^{\mathsf t}(y+\Lambda'^{\mathsf t}x
 +\Phi'^{\mathsf t}\Lambda^{\mathsf t}\Phi'x)\big),
\]
which is $T_{\phi,\psi,\lambda}\big(\Phi'x,\Psi'^{\mathsf t}(y+\Lambda'^{\mathsf t}x)\big)$ because
$\Psi'^{\mathsf t}\Phi'^{\mathsf t}=(\Phi'\Psi')^{\mathsf t}=1$ by
Lemma~\ref{cv-hodge-lattice-transport}(2).
\end{proof}

\begin{remark}\label{cv-germs-unique-lambda}
The power series $\lambda$ of a comparison in Lemma~\ref{cv-germs-compare}(3) is unique as well. In
$J^2_m$ its image is the ratio of the generators $e'_m(t,t')$ and $e_m(\phi(t),\phi(t'))$ of one line
bundle, and $\phi$ is unique; a power series is determined by its images in the rings $J^2_m$. This is
the argument for $\lambda''$ in the proof of Lemma~\ref{cv-germs-compare}(4). Nothing below uses the
uniqueness of $\lambda$: a comparison is chosen where one is needed.
\end{remark}

\begin{definition}\label{cv-germs-principal}
Let $E$ be a CM field, let $G$ be an $\mathcal O_E$-semi-abelian scheme over a scheme $X$
(Definition~\ref{cv-model-toroidal-semiabelian}) with line bundle $\mathcal L$, and let $k$ be a field.
\begin{enumerate}
\item For a scheme $T$ over $X$ and a point $b$ of $G$ with values in $T$, let
  $\mathcal L_b=(b+y)^*\mathcal L\otimes y^*\mathcal L^{-1}$, a line bundle on $T\times_XG$, where $y$ is
  the second projection and $b$ also stands for its composite with the first projection. The same
  notation is used for any commutative group scheme over $X$ with a line bundle.
\item Let $B$ be an abelian variety over $k$ with a line bundle $\mathcal M$. The group $K(\mathcal M)$
  is \emph{trivial} if a point $b$ of $B$ with values in a commutative $k$-algebra $R$ is the unit
  whenever $\operatorname{Spec}R$ is covered by open sets $V$ such that $\mathcal M_b$ has a generator
  over the preimage of $V$ in $\operatorname{Spec}R\times_kB$.
\item Let $x:\operatorname{Spec}k\to X$ be a morphism, and let $G_x$ be the fibre of $G$ at $x$, a group
  scheme over $k$. Then $G$ has a \emph{principal fibre} at $x$ if there are a free abelian group $\Xi$
  of finite rank, an abelian variety $B$ over $k$, a flat, surjective and quasi-compact homomorphism
  $q:G_x\to B$, a homomorphism $i:T_\Xi\to G_x$ from the split torus over $k$ with character group $\Xi$
  that identifies $T_\Xi$ with the kernel of $q$, and a line bundle $\mathcal M$ on $B$ with
  $K(\mathcal M)$ trivial, such that the restriction of $\mathcal L$ to $G_x$ is isomorphic to
  $q^*\mathcal M$.
\item $G$ is \emph{principal} at a point $z$ of $X$ if it has a principal fibre at some morphism
  $\operatorname{Spec}k\to X$ with image $z$ and with $k$ algebraically closed.
\item Let $x:\operatorname{Spec}k\to X$ be a morphism, let $p$ be a prime, and for $m\ge0$ let
  $k_m=k[s]/(s^{m+1})$. Then $G$ is \emph{tangentially principal} at $x$ for $p$ if the following
  holds for every point $y$ of $G$ with values in $k_p$ over $x$ that is the unit modulo $s$ and for
  which $py$ is the unit: if $y_1$ is the image of $y$ in $G(k_1)$ and $\mathcal L_{y_1}$ has a
  generator on $\operatorname{Spec}k_1\times_XG$, then $y_1$ is the unit.
\end{enumerate}
\end{definition}
\leanref{ReciprocityTheory/ReductiveGroups/ArithmeticGeometry/ExtensionDescent/NonzeroLocus.lean}{ReciprocityTheory.ChartRestriction.G.Out.translDiff}
\leanref{ReciprocityTheory/ReductiveGroups/ArithmeticGeometry/ExtensionDescent/NonzeroLocus.lean}{ReciprocityTheory.ChartRestriction.G.Out.KTriv}
\leanref{ReciprocityTheory/ReductiveGroups/ArithmeticGeometry/SemiabelianSchemes/AbelianSchemeQuotients/IsAbelianScheme.lean}{ReciprocityTheory.ChartRestriction.G.Semi.TorusExtension}
\leanref{Development/AutomorphicSpectrum/Definitions/GlobalGeometry/ToFirst.lean}{GlobalTheory.CompanionTheory.Ar.Conj.GeomPF}
\leanref{Development/AutomorphicSpectrum/Definitions/GlobalGeometry/ToFirst.lean}{GlobalTheory.CompanionTheory.Ar.Conj.FibPF}
\leanref{Development/AutomorphicSpectrum/Definitions/GlobalGeometry/ToFirst.lean}{GlobalTheory.CompanionTheory.Ar.Conj.TanPF}

\begin{lemma}\label{cv-germs-principal-base-change}
Let $E$ be a CM field, and let $G$ be an $\mathcal O_E$-semi-abelian scheme over a scheme $X$
(Definition~\ref{cv-model-toroidal-semiabelian}).
\begin{enumerate}
\item If $G$ has a principal fibre at a morphism $x:\operatorname{Spec}k\to X$
  (Definition~\ref{cv-germs-principal}), with $k$ a field, and $k\to k'$ is a homomorphism of fields,
  then $G$ has a principal fibre at the composite $\operatorname{Spec}k'\to\operatorname{Spec}k\to X$.
\item Let $G'$ be an $\mathcal O_E$-semi-abelian scheme over a scheme $X'$ with a morphism $G'\to G$ of
  $\mathcal O_E$-semi-abelian schemes over a morphism $f:X'\to X$, and let $x:\operatorname{Spec}k\to X'$
  be a morphism, with $k$ a field. Then $G'$ has a principal fibre at $x$ if and only if $G$ has a
  principal fibre at $f\circ x$.
\end{enumerate}
\end{lemma}
\leanref{Development/AutomorphicSpectrum/Proofs/GeometryProofs/GeometricPowerFreeFunctorialStatement.lean}{GlobalTheory.T.geometric_power_free_functorial}
\leanref{Development/WeilBound/CoreStatements/HodgeTatePoints/LbIsoPullPull.lean}{PositivityTheory.HtPt.Cf.geompf_bc}
\leanref{ReciprocityTheory/ReductiveGroups/ArithmeticGeometry/BirationalGeometry/ToroidalCompactification/AmplePt.lean}{ReciprocityTheory.ChartRestriction.G.Geo.Tor.kTriv_pull}
\begin{proof}
(1) Let $\Xi$, $B$, $q$, $i$ and $\mathcal M$ be data as in Definition~\ref{cv-germs-principal}(3) for
$x$, and let $x'$ be the composite. The fibre $G_{x'}$ is $G_x\otimes_kk'$. The base changes
$B\otimes_kk'$, $q\otimes k'$, $i\otimes k'$ and $\mathcal M\otimes k'$ are data of the same kind for
$x'$: flatness, surjectivity and quasi-compactness are stable under base change, the kernel of
$q\otimes k'$ is the base change of the kernel of $q$, and the isomorphism of line bundles pulls back.
The group $K(\mathcal M\otimes k')$ is trivial: a point $b$ of $B\otimes_kk'$ with values in a
$k'$-algebra $R$ is a point of $B$ with values in $R$, regarded as a $k$-algebra, and the line bundle
$\mathcal M_b$ on $\operatorname{Spec}R\times_kB=\operatorname{Spec}R\times_{k'}(B\otimes_kk')$ is the
same for both.

(2) The morphism $G'\to G$ identifies $G'$ with $f^*G$ and the line bundle of $G'$ with the pullback of
that of $G$. So it identifies the fibre $G'_x$ with $G_{f\circ x}$, as group schemes over $k$ with line
bundles, and data as in Definition~\ref{cv-germs-principal}(3) for one fibre are data for the other.
\end{proof}

\begin{lemma}[Grothendieck {\cite[Proposition~7.1.7]{Gro61b}}, {\cite[Chapitre~0, Proposition~10.3.1]{Gro61}}]\label{cv-germs-dvr}
Let $A$ be a noetherian local domain that is not a field, with residue field $\kappa$, and let $k$ be a
finite extension of $\kappa$. Then there are a complete discrete valuation ring $R$, with residue field
$\kappa_R$, an injective local homomorphism $g:A\to R$, and a homomorphism $s:k\to\kappa_R$ whose
restriction to $\kappa$ is the homomorphism $\kappa\to\kappa_R$ induced by $g$.
\end{lemma}
\leanref{Development/AutomorphicSpectrum/Proofs/GeometryProofs/ModelFibresPowerFree.lean}{GlobalTheory.CompanionTheory.Ar.ModelFib.exists_complete_dvr}
\leanref{GroundTheory/HeightsAndTheta/Heights/IntegralModels/Assembly.lean}{GroundTheory.Heights₂.Heights.Model.DominatingDVR.full}
\leanref{Development/AutomorphicSpectrum/Proofs/AlgebraProofs/CompleteDvrDominatingValuationStatement.lean}{GlobalTheory.T.complete_dvr_dominating_valuation}
\begin{proof}
By \cite[Proposition~7.1.7]{Gro61b}, the fraction field of $A$ contains a discrete valuation ring $V$
that dominates $A$, so that the inclusion $A\to V$ is an injective local homomorphism. Let $\kappa_V$ be
the residue field of $V$, choose a homomorphism $\sigma$ over $\kappa$ from $k$ to an algebraic closure
of $\kappa_V$, and let $k'$ be the field generated over $\kappa_V$ by $\sigma(k)$, a finite extension of
$\kappa_V$. By \cite[Chapitre~0, Proposition~10.3.1]{Gro61} there is a flat local noetherian
$V$-algebra $V'$ whose maximal ideal is generated by a uniformizer $\varpi$ of $V$ and whose residue
field is $k'$ over $\kappa_V$. By flatness $\varpi$ is not a zero divisor on $V'$, so $V'$ is a discrete
valuation ring, and its completion $R$ is a complete discrete valuation ring, faithfully flat over $V$,
with residue field $k'$. So $V\to R$ is an injective local homomorphism. Let $g$ be the composite
$A\to V\to R$ and $s$ the composite of $\sigma$ with $k'=\kappa_R$.
\end{proof}

\begin{theorem}[Mumford \cite{Mum72}, Faltings--Chai \cite{FC90}]\label{cv-germs-raynaud}
In Situation~\ref{cv-function-situation}, let $R$ be a local domain with maximal ideal $\mathfrak m$ and
residue field $\kappa$, and let $G$ be an $\mathcal O_E$-semi-abelian scheme over $R$
(Definition~\ref{cv-model-toroidal-semiabelian}) with line bundle $\mathcal L$ that is given by Mumford's
construction from degeneration data with respect to $\mathfrak m$, as in
Theorem~\ref{cv-model-toroidal-degeneration}(1). Then $G$ has a principal fibre at the closed point
$\operatorname{Spec}\kappa\to\operatorname{Spec}R$ (Definition~\ref{cv-germs-principal}): the
degeneration data contain a free abelian group $\Xi$ of finite rank and an abelian scheme
$\mathcal B$ over $R$ with a line bundle $\mathcal M$ whose restriction $\mathcal M_\kappa$ to the closed
fibre has $K(\mathcal M_\kappa)$ trivial; over $R/\mathfrak m$ they give a flat surjective homomorphism
$q:G_\kappa\to\mathcal B_\kappa$ whose kernel is the split torus with character group $\Xi$, and an
isomorphism of the restriction of $\mathcal L$ to $G_\kappa$ with $q^*\mathcal M_\kappa$.
\end{theorem}
\leanref{Development/AutomorphicSpectrum/Proofs/GeometryProofs/PelUniformisedGeometricFreeStatement.lean}{GlobalTheory.T.pel_uniformised_geometric_free}
\leanref{Development/WeilBound/CoreStatements/HodgeTatePoints/CgUnif.lean}{PositivityTheory.HtPt.Cg.unif_pf}
\leanref{ReciprocityTheory/ReductiveGroups/ArithmeticGeometry/BirationalGeometry/ToroidalCompactification/UniformisationConclusion.lean}{ReciprocityTheory.ChartRestriction.G.Geo.Tor.PELUnif}

The extension $q$ is the closed fibre of the Raynaud extension of $G$. In the theory of degenerations,
the formal completion of $G$ along its closed fibre is that of an extension $\tilde G$ of the abelian
scheme $\mathcal B$ by a torus, the line bundle of $G$ corresponds to the pullback to $\tilde G$ of an
ample line bundle $\mathcal M$ on $\mathcal B$, and a principal polarization of $G$ induces one on
$\mathcal B$, so that $K(\mathcal M)$ is trivial \cite[Chapters~II and~III]{FC90}. Mumford's construction
recovers $G$ from these data and the group of periods \cite{Mum72}.

\begin{lemma}\label{cv-germs-principal-fibres}
In Situation~\ref{cv-function-situation}, let $(\mathcal X,\mathcal A,j)$ be a model of $X$ over a
discrete valuation ring $O\subset F_c$ with fraction field $F_c$ (Definition~\ref{cv-model}) that has
properties (1) and (2) of Theorem~\ref{cv-model-toroidal}, and suppose that $j(X)$ is dense in $\mathcal X$.
Let $p$ be a prime and $z$ a point of $\mathcal X$ whose residue field $\kappa(z)$ has characteristic $p$.
Then $\mathcal A$ is principal at $z$ (Definition~\ref{cv-germs-principal}).
\end{lemma}
\leanref{Development/AutomorphicSpectrum/Proofs/GeometryProofs/ModelFibresPowerFreeStatement.lean}{GlobalTheory.T.model_fibres_power_free}
\leanref{Development/AutomorphicSpectrum/Proofs/GeometryProofs/ModelFibresPowerFree.lean}{GlobalTheory.CompanionTheory.Ar.ModelFib.model_fibPF}
\begin{proof}
The local ring. Let $A$ be the local ring of $\mathcal X$ at $z$. The scheme $\mathcal X$ is normal and
locally noetherian, so $A$ is a noetherian normal domain, and the generic point of
$\operatorname{Spec}A$ maps to a generic point $\gamma$ of $\mathcal X$. As $j(X)$ is dense in $\mathcal X$,
$\gamma$ specializes to a point $u$ of $j(X)$. The points of $j(X)$ lie over the generic point $\eta$ of
$\operatorname{Spec}O$, and, again as $j(X)$ is dense, every point of $\mathcal X$ lies over a specialization
of $\eta$. The image of $\gamma$ is therefore a specialization of $\eta$ that specializes to the image $\eta$
of $u$, so it is $\eta$. So $\gamma$ lies in $\mathcal X_{F_c}$, where it is a generic point; as $j(X)$ is
open and dense in $\mathcal X_{F_c}$, $\gamma\in j(X)$. The ring $A$ maps to the local ring of $\mathcal X$ at $\gamma$,
which contains $F_c$, so $A$ has characteristic $0$. Then $p$ is a nonzero element of the maximal ideal
of $A$, and $A$ is not a field.

The complete discrete valuation ring. By property (1) of Theorem~\ref{cv-model-toroidal}, there is a
finite separable extension $k$ of $\kappa(z)$ such that the fibre of $\mathcal A$ at
$\operatorname{Spec}k\to\mathcal X$ is an extension of an abelian variety by a split torus. By
Lemma~\ref{cv-germs-dvr} there are a complete discrete valuation ring $R$, an injective local
homomorphism $h:A\to R$ and a homomorphism $s:k\to\kappa_R$ whose restriction to $\kappa(z)$ is induced
by $h$. Let $\xi:\operatorname{Spec}R\to\mathcal X$ be the composite of $\operatorname{Spec}h$ with
$\operatorname{Spec}A\to\mathcal X$, and let $G=\xi^*\mathcal A$, an $\mathcal O_E$-semi-abelian scheme
over $R$ with a morphism $G\to\mathcal A$ over $\xi$.

The degeneration. The ring $R$ is a complete noetherian normal local domain, of characteristic $0$ as it
contains $A$. As $h$ is injective, $\xi$ maps the generic point of $\operatorname{Spec}R$ to $\gamma$,
a point of $j(X)$. It maps the closed point to $z$, and the morphism
$\operatorname{Spec}\kappa_R\to\mathcal X$ is the composite of $\operatorname{Spec}s$ with
$\operatorname{Spec}k\to\mathcal X$, by the compatibility of $s$ with $h$. So the fibre of $G$ over the
closed point is the base change along $s$ of an extension of an abelian variety by a split torus, and
its torus is split. By property (2) of Theorem~\ref{cv-model-toroidal}, $G$ is given by Mumford's
construction from degeneration data with respect to the maximal ideal of $R$. By
Theorem~\ref{cv-germs-raynaud}, $G$ has a principal fibre at
$\operatorname{Spec}\kappa_R\to\operatorname{Spec}R$.

Conclusion. By Lemma~\ref{cv-germs-principal-base-change}(2), applied to $G\to\mathcal A$, the scheme
$\mathcal A$ has a principal fibre at $\operatorname{Spec}\kappa_R\to\mathcal X$, and by
Lemma~\ref{cv-germs-principal-base-change}(1) at the composite
$\operatorname{Spec}\Omega\to\operatorname{Spec}\kappa_R\to\mathcal X$ for an algebraic closure $\Omega$
of $\kappa_R$. Its image is the image of the closed point of $\operatorname{Spec}R$, which is $z$. So
$\mathcal A$ is principal at $z$.
\end{proof}

\begin{lemma}\label{cv-germs-torus-lift}
Let $k$ be a field of characteristic $p>0$, let $k_m=k[s]/(s^{m+1})$ for $m\ge0$, let $\Xi$ be a free
abelian group of finite rank, and let $T$ be the split torus over $k$ with character group $\Xi$. Let
$y_1$ be a point of $T$ with values in $k_1$ that is the unit modulo $s$. Then there is a point $y$ of
$T$ with values in $k_p$ that is the unit modulo $s$, has image $y_1$ in $T(k_1)$, and is such that
$y_1$ is the unit if $py$ is the unit.
\end{lemma}
\leanref{Development/WeilBound/CoreStatements/HodgeTatePoints/IsPullbackWhisker.lean}{PositivityTheory.HtPt.Ce.torus_tan}
\begin{proof}
A point of $T$ with values in a $k$-algebra $C$ is a homomorphism $\Xi\to C^\times$. An element of $k_p$
whose image in $k_1$ is a unit is a unit, so the reduction $k_p^\times\to k_1^\times$ is onto. As $\Xi$ is
free, $y_1:\Xi\to k_1^\times$ lifts to a homomorphism $y:\Xi\to k_p^\times$. Its image in $T(k_1)$ is
$y_1$, and modulo $s$ it agrees with $y_1$, so it is the unit. Suppose that $py$ is the unit, that is,
$y(\chi)^p=1$ for every $\chi\in\Xi$. Then $a=y(\chi)-1$ satisfies $a^p=y(\chi)^p-1=0$, since $k_p$ has
characteristic $p$. Write $a\equiv a_0+a_1s$ modulo $s^2$ with $a_0,a_1\in k$. Then
$a^p=a_0^p+a_1^ps^p$ in $k_p$, because $(a-a_0-a_1s)^p$ lies in $(s^{2p})$ and $2p\ge p+1$. So
$a_0=a_1=0$, and $y_1(\chi)=1$ for every $\chi$.
\end{proof}

\begin{lemma}\label{cv-germs-frobenius-series}
Let $E$ be a CM field, let $p$ be a prime, let $R$ be a commutative ring in which $p=0$, and let $B$ be a
formal germ of dimension $n$ over $R$ (Definition~\ref{cv-hodge-lattice-germ}). Let $W$ be the matrix of
Definition~\ref{cv-hodge-lattice-germ}(5) for $\ell=p$. Then $[p]x=V(x_1^p,\dots,x_n^p)$ for an $n$-tuple
$V$ of power series over $R$ without constant term whose matrix of linear coefficients is $W$.
\end{lemma}
\leanref{Development/WeilBound/CoreStatements/HodgeTatePoints/CbFrob.lean}{PositivityTheory.HtPt.Cb.mulp_V}
\leanref{Development/WeilBound/CoreStatements/HodgeTatePoints/CbFrob.lean}{PositivityTheory.HtPt.Cb.d_act_p}
\begin{proof}
The matrix of linear coefficients of $[p]x$ is $\rho_B(p)=p\cdot1=0$, by the linear parts of the
identities $[a+b]x=[a]x+_F[b]x$ and $[1]x=x$. Differentiate $[p](x+_Fy)=[p]x+_F[p]y$ in $y$ at $y=0$. The
right side gives $0$, as the linear coefficients of $[p]y$ vanish; the left side gives the product of the
Jacobian matrix of $[p]x$ with that of $x+_Fy$ in $y$ at $y=0$, which is invertible, as its constant term is
$1$. So all partial derivatives of the components of $[p]x$ vanish. An integer prime to $p$ is a unit in
$R$, so these components have nonzero coefficients only at monomials whose exponents are all divisible by
$p$; that is, $[p]x=V(x_1^p,\dots,x_n^p)$. The series $V$ has no constant term since $[p]x$ has none, and
its linear coefficients are the coefficients of the $x_j^p$ in $[p]x$, which form $W$.
\end{proof}

\begin{lemma}\label{cv-germs-frobenius-points}
Let $E$ be a CM field, let $G$ be an $\mathcal O_E$-semi-abelian scheme over a scheme $X$
(Definition~\ref{cv-model-toroidal-semiabelian}), let $k$ be a field of characteristic $p>0$, let
$x:\operatorname{Spec}k\to X$ be a morphism, and let $k_m=k[s]/(s^{m+1})$ for $m\ge0$. Points of $G$ with
values in $k$-algebras are taken over $x$.
\begin{enumerate}
\item Let $g$ be a formal chart of dimension $n$ at $x$ (Definition~\ref{cv-germs-chart}), with points
  $x_m$ and formal germ $B$, and let $W$ be the matrix of Definition~\ref{cv-hodge-lattice-germ}(5) for
  $\ell=p$. For a homomorphism of $k$-algebras $\chi:J^1_p\to C$, the point
  $p\,(x_p\circ\operatorname{Spec}\chi)$ is $x_p\circ\operatorname{Spec}\psi$, where $\psi:J^1_p\to C$ is
  the homomorphism with $\psi(t_i)=\sum_jW_{ij}\chi(t_j)^p$ for every $i$.
\item In (1), for $w\in k^n$ let $y_w$ be the point of $G$ with values in $k_p$ obtained from $x_p$ by
  $t_j\mapsto w_js$. Then $y_w$ is the unit modulo $s$; if $\sum_jW_{ij}w_j^p=0$ for every $i$, then $py_w$
  is the unit; and if the image of $y_w$ in $G(k_1)$ is the unit, then $w=0$.
\item Suppose that, for every $m\ge1$, every function on the $m$-fold fibre product of $G$ over $X$ is the
  pullback of a function on $X$. Let $y$ and $y'$ be points of $G$ with values in $k_p$ that are the unit
  modulo $s$ and have the same image in $G(k_1)$. Then $py=py'$.
\end{enumerate}
\end{lemma}
\leanref{Development/WeilBound/CoreStatements/HodgeTatePoints/Comp1PointPowerExpand.lean}{PositivityTheory.HtPt.Ce.pow_coord}
\leanref{Development/WeilBound/CoreStatements/HodgeTatePoints/IsPullbackWhisker.lean}{PositivityTheory.HtPt.Ce.pt_of_vec}
\leanref{Development/WeilBound/CoreStatements/HodgeTatePoints/IsPullbackWhisker.lean}{PositivityTheory.HtPt.Ce.pow_first}
\begin{proof}
(1) By Definition~\ref{cv-germs-chart}(3)(a) for $a=p$, $p\,x_p(t)=x_p([p]t)$, so
$p\,(x_p\circ\operatorname{Spec}\chi)=x_p\circ\operatorname{Spec}\psi$ with
$\psi(t_i)=\chi([p]_it)$, where $[p]_i$ is the $i$-th component of $[p]$. By
Lemma~\ref{cv-germs-frobenius-series}, $[p]_it=V_i(t_1^p,\dots,t_n^p)$, which is
$\sum_jW_{ij}t_j^p$ modulo $I^{2p}$, as $V$ has no constant term. As $I^{p+1}=0$ in $J^1_p$ and
$2p\ge p+1$, $\psi(t_i)=\sum_jW_{ij}\chi(t_j)^p$.

(2) Let $\chi_w:J^1_p\to k_p$ be the homomorphism with $\chi_w(t_j)=w_js$, so that
$y_w=x_p\circ\operatorname{Spec}\chi_w$. Modulo $s$, $\chi_w$ becomes $t\mapsto0$, and $x_p(0)$ is the
unit by Definition~\ref{cv-germs-chart}(3)(a). If $\sum_jW_{ij}w_j^p=0$ for every $i$, then (1) gives
$py_w=x_p\circ\operatorname{Spec}\psi$ with $\psi(t_i)=\sum_jW_{ij}w_j^ps^p=0$, so $py_w=x_p(0)$ is the
unit. The image of $y_w$ in $G(k_1)$ is $x_1\circ\operatorname{Spec}\bar\chi_w$, where
$\bar\chi_w:J^1_1\to k_1$ maps $t_j$ to $w_js$, as $x_p$ reduces to $x_1$. If it is the unit
$x_1(0)$, the uniqueness in Theorem~\ref{cv-germs-formal}(2), for the pullback of $G$ to
$\operatorname{Spec}k$, with $C=k_1$ and $N=(s)$, gives $\bar\chi_w(t_j)=0$, so $w_js=0$ in $k_1$ and
$w_j=0$ for every $j$.

(3) By Lemma~\ref{cv-germs-local}, $x$ factors through an affine open subset of $X$ with a formal chart,
and by Lemma~\ref{cv-germs-compare}(1) its base change is a formal chart $g$ at $x$, of some dimension
$n$, with points $x_m$. By Theorem~\ref{cv-germs-formal}(2) with $C=k_p$ and $N=(s)$, $y=x_p\circ
\operatorname{Spec}\chi$ and $y'=x_p\circ\operatorname{Spec}\chi'$ for homomorphisms
$\chi,\chi':J^1_p\to k_p$, which map each $t_j$ into $(s)$. Their reductions modulo $s^2$ factor through
homomorphisms $\bar\chi,\bar\chi':J^1_1\to k_1$, and the images of $y$ and $y'$ in $G(k_1)$ are
$x_1\circ\operatorname{Spec}\bar\chi$ and $x_1\circ\operatorname{Spec}\bar\chi'$. These are equal, so
$\bar\chi=\bar\chi'$ by the uniqueness in Theorem~\ref{cv-germs-formal}(2): $\chi(t_j)\equiv\chi'(t_j)$
modulo $s^2$ for every $j$. For $a,b\in k_p$ with $a\equiv b$ modulo $s^2$, $a^p-b^p=(a-b)^p$ lies in
$(s^{2p})=0$. So $\chi(t_j)^p=\chi'(t_j)^p$ for every $j$, and by (1), $py$ and $py'$ are
$x_p\circ\operatorname{Spec}\psi$ for the same $\psi$.
\end{proof}

\begin{lemma}\label{cv-germs-descent}
Let $k$ be a field, let $k_1=k[s]/(s^2)$, and let $G$ be a commutative group scheme over $k$. Let $\Xi$ be
a free abelian group of finite rank, $B$ an abelian variety over $k$, $q:G\to B$ a flat, surjective and
quasi-compact homomorphism, and $i:T_\Xi\to G$ a homomorphism from the split torus with character group
$\Xi$ that identifies $T_\Xi$ with the kernel of $q$. Let $\mathcal M$ be a line bundle on $B$, and let
$y$ be a point of $G$ with values in $k_1$ that is the unit modulo $s$. If the line bundle
$(q^*\mathcal M)_y$ on $\operatorname{Spec}k_1\times_kG$ (Definition~\ref{cv-germs-principal}(1)) has a
generator, then $\mathcal M_{q(y)}$ has a generator on $\operatorname{Spec}k_1\times_kB$.
\end{lemma}
\leanref{Development/WeilBound/CoreStatements/HodgeTatePoints/IsPullbackWhisker.lean}{PositivityTheory.HtPt.Ce.triv_descend}
\leanref{Development/WeilBound/CoreStatements/HodgeTatePoints/IsPullbackWhisker.lean}{PositivityTheory.HtPt.Ce.desc_first}
\leanref{Development/WeilBound/CoreStatements/HodgeTatePoints/IsPullbackWhisker.lean}{PositivityTheory.HtPt.Ce.cobound}
\begin{proof}
Let $S=\operatorname{Spec}k_1$, $Y=S\times_kG$, $Y_B=S\times_kB$ and $Q=1\times q:Y\to Y_B$, which is flat,
surjective and quasi-compact, and let $N=\mathcal M_{q(y)}$ on $Y_B$. As $q$ is a homomorphism,
$(q^*\mathcal M)_y=Q^*N$, so $Q^*N$ has a generator $c_0$.

Normalization. Let $z_0:G\to Y$ be the section over the closed point $s=0$ of $S$. As $y$ is the unit modulo
$s$, the pullback of $N$ along $Q\circ z_0$ is $\mathcal M\otimes\mathcal M^{-1}$, which is canonically
trivial. Dividing $c_0$ by the pullback, along the projection $Y\to G$, of its pullback along $z_0$ gives
a generator $c$ of $Q^*N$ that is $1$ along $z_0$.

The cocycle. The torus acts on $Y$ by $a(\tau,(\sigma,u))=(\sigma,i(\tau)+u)$, and $Q\circ a=Q\circ\pi$,
where $\pi:T_\Xi\times_kY\to Y$ is the projection, since $i(\tau)$ lies in the kernel of $q$. So $a^*c$ and
$\pi^*c$ are generators of one line bundle on $T_\Xi\times_kY$, and $a^*c=\rho\,\pi^*c$ for a unit $\rho$
on $T_\Xi\times_kY$. A point $v$ of $Y$ at which $s$ vanishes factors through $z_0$, and so does
$a(\tau,v)$; hence $\rho(\tau,v)=1$ for such $v$. As $a(\tau\tau',v)=a(\tau,a(\tau',v))$,
\[
 \rho(\tau\tau',v)=\rho(\tau,a(\tau',v))\,\rho(\tau',v).
\]

The coboundary. For a point $v$ of $Y$ with values in a ring $C$, the function $\tau\mapsto\rho(\tau,v)$ on
the torus over $C$ is an element of the group algebra $C[\Xi]$; let $f(v)$ be its coefficient at $0$,
minus $1$. This is compatible with homomorphisms $C\to C'$, so $f$ is a function on $Y$. Since $\rho=1$
where $s=0$, applied over $C/sC$, both $\rho(\tau,v)-1$ and $f(v)$ are multiples of $s$, and $f^2=0$ as
$s^2=0$. Evaluate the cocycle identity at the universal point $\tau$ of the torus over $C[\Xi]$ and at a
point $\tau'$ with values in $C$. The left side is the element of $C[\Xi]$ attached to $v$, with the
coefficient of each $\chi$ multiplied by $\chi(\tau')$, which leaves the coefficient at $0$ unchanged;
the right side is the element attached to $a(\tau',v)$ times $\rho(\tau',v)$. Products of two multiples
of $s$ vanish, so the coefficients at $0$ give $f(v)=f(a(\tau',v))+\rho(\tau',v)-1$. Hence $u=1+f$ is a unit on $Y$, with inverse $1-f$,
and $u(v)=u(a(\tau',v))\,\rho(\tau',v)$.

Descent. The generator $c'=uc$ of $Q^*N$ satisfies $a^*c'=a^*u\cdot\rho\cdot\pi^*c=\pi^*u\cdot\pi^*c=
\pi^*c'$. Two points of $Y$ with the same image under $Q$ are $a(\tau,v)$ and $v$ for a point $\tau$ of the
torus, as $T_\Xi$ is the kernel of $q$, so the two pullbacks of $c'$ to $Y\times_{Y_B}Y$ agree. Since $Q$
is faithfully flat and quasi-compact, $c'$ is the pullback of a section of $N$
\cite[Expos\'e~VIII, Th\'eor\`eme~1.1]{Gro71}, which is a generator as $c'$ is one.
\end{proof}

\begin{lemma}\label{cv-germs-tangent}
Let $E$ be a CM field, and let $G$ be an $\mathcal O_E$-semi-abelian scheme over a scheme $X$
(Definition~\ref{cv-model-toroidal-semiabelian}) such that, for every $m\ge1$, every function on the
$m$-fold fibre product of $G$ over $X$ is the pullback of a function on $X$. Let $k$ be a field of
characteristic $p>0$ and $x:\operatorname{Spec}k\to X$ a morphism. If $G$ has a principal fibre at $x$,
then $G$ is tangentially principal at $x$ for $p$ (Definition~\ref{cv-germs-principal}).
\end{lemma}
\leanref{Development/AutomorphicSpectrum/Proofs/GeometryProofs/GeometricToTangentPowerFreeStatement.lean}{GlobalTheory.T.geometric_to_tangent_power_free}
\leanref{Development/WeilBound/CoreStatements/HodgeTatePoints/IsPullbackWhisker.lean}{PositivityTheory.HtPt.Ce.geompf_tan}
\begin{proof}
Let $\Xi$, $B$, $q$, $i$ and $\mathcal M$ be data as in Definition~\ref{cv-germs-principal}(3) for $x$,
and let $\mathcal L$ be the line bundle of $G$. Let $y$ be a point of $G$ with values in $k_p$ over $x$
that is the unit modulo $s$, with $py$ the unit, and such that $\mathcal L_{y_1}$ has a generator, where
$y_1$ is the image of $y$ in $G(k_1)$. Points of $G$ over $x$ with values in a $k$-algebra are points of
the fibre $G_x$, so $y$ and $y_1$ are points of $G_x$. The restriction of the generator to
$\operatorname{Spec}k_1\times_kG_x$ is a generator of $(\mathcal L|_{G_x})_{y_1}$, and through the
isomorphism $\mathcal L|_{G_x}\cong q^*\mathcal M$ it gives a generator of $(q^*\mathcal M)_{y_1}$. By
Lemma~\ref{cv-germs-descent}, $\mathcal M_{q(y_1)}$ has a generator on
$\operatorname{Spec}k_1\times_kB$. As $K(\mathcal M)$ is trivial, $q(y_1)$ is the unit, so
$y_1=i(\tau_1)$ for a point $\tau_1$ of $T_\Xi$ with values in $k_1$, which is the unit modulo $s$ since
$i$ is a monomorphism. Let $\tau$ be the point of $T_\Xi$ given by Lemma~\ref{cv-germs-torus-lift} for
$\tau_1$. The points $i(\tau)$ and $y$ of $G$ with values in $k_p$ are the unit modulo $s$ and have the
same image $y_1$ in $G(k_1)$, so $p\,i(\tau)=py$ by Lemma~\ref{cv-germs-frobenius-points}(3), and this is
the unit. As $i$ is a monomorphism, $p\tau$ is the unit, so $\tau_1$ is the unit by
Lemma~\ref{cv-germs-torus-lift}, and so is $y_1=i(\tau_1)$.
\end{proof}

\begin{lemma}\label{cv-germs-conj-injective}
Let $E$ be a CM field, let $G$ be an $\mathcal O_E$-semi-abelian scheme over a scheme $X$
(Definition~\ref{cv-model-toroidal-semiabelian}) with line bundle $\mathcal L$, let $k$ be an
algebraically closed field of characteristic $p>0$, let $x:\operatorname{Spec}k\to X$ be a morphism, and
let $g$ be a formal chart of dimension $n$ at $x$ (Definition~\ref{cv-germs-chart}) with formal germ $B$.
For $w\in k^n$ let $y_w$ be the point of Lemma~\ref{cv-germs-frobenius-points}(2) and $y_{w,1}$ its image
in $G(k_1)$. Suppose that $G$ is tangentially principal at $x$ for $p$
(Definition~\ref{cv-germs-principal}(5)), and that $\mathcal L_{y_{w,1}}$ has a generator on
$\operatorname{Spec}k_1\times_XG$ for every $w\in k^n$ with $\alpha_B(w_1^p,\dots,w_n^p)=0$. Then the
conjugate map $\alpha_B$ (Definition~\ref{cv-hodge-lattice-germ}(5), with $\ell=p$) is injective.
\end{lemma}
\leanref{Development/WeilBound/CoreStatements/HodgeTatePoints/IsPullbackWhisker.lean}{PositivityTheory.HtPt.Ce.tanpf_inj_of}
\begin{proof}
Let $u\in k^n$ with $\alpha_B(u)=0$. As $k$ is algebraically closed, $u=(w_1^p,\dots,w_n^p)$ for some
$w\in k^n$. The first component of $\alpha_B(u)$ is $Wu$, so $\sum_jW_{ij}w_j^p=0$ for every $i$, and by
Lemma~\ref{cv-germs-frobenius-points}(2), $y_w$ is the unit modulo $s$ and $py_w$ is the unit. By
hypothesis $\mathcal L_{y_{w,1}}$ has a generator, so $y_{w,1}$ is the unit, as $G$ is tangentially
principal at $x$. Then $w=0$ by Lemma~\ref{cv-germs-frobenius-points}(2), and $u=0$.
\end{proof}

\begin{lemma}\label{cv-germs-kernel-bound}
Let $k$ be a field and $n,n'\ge0$. For $n\times n$ matrices $\rho$, $\bar\rho$ and $A$ over $k$, let
$T(\rho,\bar\rho,A)$ be the endomorphism $(x,y)\mapsto(\rho x,\bar\rho^{\mathsf t}y-A^{\mathsf t}x)$ of
$k^n\oplus k^n$, and likewise for $n'\times n'$ matrices.
\begin{enumerate}
\item For every $n\times n$ matrix $Z$ over $k$, the kernels of $T(\rho,\bar\rho,A)$ and of
  $T(\rho,\bar\rho,A+\rho^{\mathsf t}Z-Z\bar\rho)$ have the same dimension.
\item Let $L$ be a finite-dimensional $k$-vector space with endomorphisms $S$ and $\bar S$, let $P$ be an
  $n'\times n$ matrix and $j:L\to k^n$ an injective linear map such that $P:k^n\to k^{n'}$ is onto with
  kernel $j(L)$, and let $\rho'$, $\bar\rho'$ and $A'$ be $n'\times n'$ matrices with $P\rho=\rho'P$,
  $P\bar\rho=\bar\rho'P$, $A=P^{\mathsf t}A'P$, $jS=\rho j$ and $j\bar S=\bar\rho j$. Then
  \[
   \dim\ker T(\rho,\bar\rho,A)\le\dim\ker\bar S+\dim\ker T(\rho',\bar\rho',A')+\dim\ker S.
  \]
\end{enumerate}
\end{lemma}
\leanref{Development/AutomorphicSpectrum/Proofs/GeometryProofs/HeightCharFreeGeometric.lean}{GlobalTheory.CompanionTheory.Ar.Flat.LemFGeo.finrank_ker_gauge}
\leanref{Development/AutomorphicSpectrum/Proofs/GeometryProofs/HeightCharFreeGeometric.lean}{GlobalTheory.CompanionTheory.Ar.Flat.LemFGeo.ker_bound}
\begin{proof}
(1) The automorphism $\theta(x,y)=(x,y+Z^{\mathsf t}x)$ of $k^n\oplus k^n$ satisfies
$T(\rho,\bar\rho,A)\circ\theta=\theta\circ T(\rho,\bar\rho,A+\rho^{\mathsf t}Z-Z\bar\rho)$, as both send
$(x,y)$ to $(\rho x,\bar\rho^{\mathsf t}y+\bar\rho^{\mathsf t}Z^{\mathsf t}x-A^{\mathsf t}x)$. So $\theta$
maps the second kernel onto the first.

(2) Let $K=\ker T(\rho,\bar\rho,A)$; for $(x,y)\in K$, $\rho x=0$ and $\bar\rho^{\mathsf t}y=A^{\mathsf t}x$.
Choose an $n\times n'$ matrix $Q$ with $PQ=1$. Let $c:k^n\oplus k^n\to L^\vee$ map $(x,y)$ to
$v\mapsto y\cdot j(v)$, and let $m(x,y)=(Px,Q^{\mathsf t}y)$. As $Pj=0$, also $Aj=0$.
\begin{itemize}
\item $c(K)\subset\ker\bar S^\vee$: for $(x,y)\in K$,
  $y\cdot j(\bar Sv)=\bar\rho^{\mathsf t}y\cdot j(v)=x\cdot Aj(v)=0$.
\item If $c(x,y)=0$, then $y$ vanishes on $j(L)=\ker P$, so $y=P^{\mathsf t}Q^{\mathsf t}y$. For
  $(x,y)\in K\cap\ker c$ this gives $m(x,y)\in\ker T(\rho',\bar\rho',A')$: $\rho'Px=P\rho x=0$, and
  $Q^{\mathsf t}$ applied to $\bar\rho^{\mathsf t}y-A^{\mathsf t}x=P^{\mathsf t}(\bar\rho'^{\mathsf t}
  Q^{\mathsf t}y-A'^{\mathsf t}Px)$ gives $\bar\rho'^{\mathsf t}Q^{\mathsf t}y-A'^{\mathsf t}Px=0$, as
  $Q^{\mathsf t}P^{\mathsf t}=1$.
\item If $(x,y)\in K\cap\ker c\cap\ker m$, then $Px=0$, so $x=j(v)$ for some $v\in L$, and
  $y=P^{\mathsf t}Q^{\mathsf t}y=0$; as $j(Sv)=\rho x=0$ and $j$ is injective, $Sv=0$. So this space is
  the image of a subspace of $\ker S$.
\end{itemize}
Hence $\dim K\le\dim c(K)+\dim m(K\cap\ker c)+\dim(K\cap\ker c\cap\ker m)$ is at most
$\dim\ker\bar S^\vee+\dim\ker T(\rho',\bar\rho',A')+\dim\ker S$, and $\dim\ker\bar S^\vee=\dim\ker\bar S$.
\end{proof}

\begin{theorem}[Mumford \cite{Mum70}]\label{cv-germs-kernel-degree}
Let $E$ be a CM field of degree $e$ over $\mathbf Q$, let $k$ be an algebraically closed field, and let
$B$ be an $\mathcal O_E$-semi-abelian scheme over $\operatorname{Spec}k$
(Definition~\ref{cv-model-toroidal-semiabelian}) that is proper over $k$ and has a formal chart of
dimension $n_0$ at the identity of $\operatorname{Spec}k$ (Definition~\ref{cv-germs-chart}). Let
$a\in\mathcal O_E$ be nonzero, with norm $N(a)$. Then the kernel of the action of $a$ on $B$, the fibre
product of $a:B\to B$ with the unit section, is $\operatorname{Spec}A_a$ for a finite $k$-algebra
$A_a$, and $(\dim_kA_a)^e=|N(a)|^{2n_0}$.
\end{theorem}
\leanref{Development/AutomorphicSpectrum/Proofs/GeometryProofs/ExistsFiniteKernelNonzeroStatement.lean}{GlobalTheory.T.exists_finite_kernel_nonzero}
\leanref{Development/WeilBound/CoreStatements/TorusUniformisation/ChowMahlerFamilies/DegreeMultiplicationNorm.lean}{PositivityTheory.Xgn.degree_multiplication_norm}

The scheme $B$ is an abelian variety of dimension $n_0$, since its tangent vectors at the identity form
a $k$-vector space with a basis of $n_0$ elements (Definition~\ref{cv-germs-chart}(3)(b)). As $a$
divides the nonzero integer $N(a)$ in $\mathcal O_E$, the action of $a$ is an isogeny, and $\dim_kA_a$
is its degree. The degree of an endomorphism of $B$ is the value at $0$ of its characteristic polynomial
on the Tate module $V_qB$ for a prime $q$ different from the characteristic of $k$, and this polynomial
has rational coefficients \cite[\S19]{Mum70}. Over an algebraic closure of $\mathbf Q_q$, each
$b\in E$ acts on $V_qB$ with the eigenvalues $\sigma(b)$, for the embeddings $\sigma$ of $E$, with
multiplicities that do not depend on $b$. For a primitive element $b$ of $E$ the characteristic
polynomial is rational, so these multiplicities are all $2n_0/e$, and the degree of $a$ is
$\prod_\sigma\sigma(a)^{2n_0/e}=N(a)^{2n_0/e}$, so $(\dim_kA_a)^e=|N(a)|^{2n_0}$.

\begin{lemma}\label{cv-germs-kernel-pairing}
Let $E$ be a CM field, let $k$ be a field, let $B$ be a formal germ of dimension $n$ over $k$
(Definition~\ref{cv-hodge-lattice-germ}), and let $a\in\mathcal O_E$. Let $W_a$ be the space of
$u\in\ker\rho_B(a)$ such that $z\cdot\beta_a^{\mathsf t}u=0$ for every $z\in\ker\rho_B(\bar a)$. Then
\[
 \dim\ker\rho^H_B(a)\le\dim\ker\rho_B(\bar a)+\dim W_a.
\]
\end{lemma}
\leanref{Development/AutomorphicSpectrum/Proofs/GeometryProofs/ProperMultiplicationKernel.lean}{GlobalTheory.CompanionTheory.Ar.Flat.KerBd.finrank_ker_rhoH_le}
\begin{proof}
Let $(x,y)\in\ker\rho^H_B(a)$, so that $\rho_B(a)x=0$ and $\rho_B(\bar a)^{\mathsf t}y=\beta_a^{\mathsf t}x$.
For $z\in\ker\rho_B(\bar a)$, $z\cdot\beta_a^{\mathsf t}x=z\cdot\rho_B(\bar a)^{\mathsf t}y=\rho_B(\bar a)z\cdot
y=0$, so $x\in W_a$. The first projection therefore maps $\ker\rho^H_B(a)$ to $W_a$. Its kernel consists
of the $(0,y)$ with $\rho_B(\bar a)^{\mathsf t}y=0$, a space of dimension
$\dim\ker\rho_B(\bar a)^{\mathsf t}=\dim\ker\rho_B(\bar a)$.
\end{proof}

\begin{lemma}\label{cv-germs-lattice-kernel}
Let $E$ be a number field, let $k$ be a field of characteristic $p>0$, and let $\Lambda$ be an
$\mathcal O_E$-module that is free of finite rank over $\mathbf Z$, of rank $\rho_0$ over
$\mathcal O_E$, so that $E\otimes_{\mathcal O_E}\Lambda$ has dimension $\rho_0$ over $E$. Let
$a\in\mathcal O_E$ be nonzero, with norm $N(a)$, and let $v_p$ be the $p$-adic valuation. Then the
characteristic polynomial of the endomorphism $1\otimes a$ of $k\otimes_{\mathbf Z}\Lambda$ is the
$\rho_0$-th power of that of multiplication by $a$ on $\mathcal O_E$, and
$\dim_k\ker(1\otimes a)\le\rho_0\,v_p(N(a))$.
\end{lemma}
\leanref{Development/AutomorphicSpectrum/Proofs/AlgebraProofs/LatticeKernelRankBoundNormStatement.lean}{GlobalTheory.T.lattice_kernel_rank_bound_norm}
\leanref{Development/WeilBound/CoreStatements/TorusUniformisation/ChowMahlerFamilies/KernelRankNormBound.lean}{PositivityTheory.Xgn.kernel_rank_norm_bound}
\begin{proof}
The determinant. There are $v_1,\dots,v_{\rho_0}\in\Lambda$ that are linearly independent over
$\mathcal O_E$, and $L:\mathcal O_E^{\rho_0}\to\Lambda$, $c\mapsto\sum_ic_iv_i$, is injective. Its
image has rank $\rho_0$, so the quotient $\Lambda/L(\mathcal O_E^{\rho_0})$ has rank $0$ over
$\mathcal O_E$: for $x\in\Lambda$, $bx\in L(\mathcal O_E^{\rho_0})$ for some nonzero $b\in\mathcal O_E$.
As $b$ is integral
over $\mathbf Z$, $bc$ is a nonzero integer $m$ for some $c\in\mathcal O_E$, and
$mx=c\,bx\in L(\mathcal O_E^{\rho_0})$. So $L(\mathcal O_E^{\rho_0})$ has finite index in $\Lambda$, and
the free $\mathbf Z$-modules $\mathcal O_E^{\rho_0}$ and $\Lambda$ have the same rank. If $f$ and $g$ are
endomorphisms of free $\mathbf Z$-modules $M$ and $M'$ of the same finite rank, and $\iota:M'\to M$ is an
injective homomorphism with $f\iota=\iota g$, then $\det f=\det g$: after an identification of $M'$
with $M$, $\iota$ becomes an endomorphism with nonzero determinant, and
$\det f\det\iota=\det\iota\det g$. Multiplication by $a$ commutes with $L$, so its determinant on
$\Lambda$ is its determinant on $\mathcal O_E^{\rho_0}$, which is $N(a)^{\rho_0}$.

The characteristic polynomial. For $t\in\mathbf Z$, the value at $t$ of the characteristic polynomial
of $a$ on $\Lambda$ is the determinant of $t-a$ on $\Lambda$. By the preceding paragraph it is
$N(t-a)^{\rho_0}$, the value at $t$ of the $\rho_0$-th power of the characteristic polynomial of $a$ on
$\mathcal O_E$. Two polynomials over $\mathbf Z$ that agree at every integer are equal. The
characteristic polynomial of $1\otimes a$ on $k\otimes_{\mathbf Z}\Lambda$ is the image of that of $a$
on $\Lambda$.

The kernel. Let $f$ be multiplication by $a$ on $\Lambda$. Its determinant $N(a)^{\rho_0}$ is nonzero,
so $f$ is injective. By the theorem on elementary divisors \cite[Chapter~III, \S7]{Lan02}, there are a
basis $t_1,\dots,t_m$ of $\Lambda$ and nonzero integers $d_1,\dots,d_m$ such that
$d_1t_1,\dots,d_mt_m$ is a basis of $f(\Lambda)$. Let $Q$ be the matrix of the isomorphism
$\Lambda\to f(\Lambda)$ induced by $f$ in the bases $(t_i)$ and $(d_it_i)$; its determinant is a unit
of $\mathbf Z$. The matrix of $f$ in the basis $(t_i)$ is $\operatorname{diag}(d_1,\dots,d_m)\,Q$. So
$d_1\cdots d_m$ divides $\det f$, and the matrix of $1\otimes f$ on $k\otimes_{\mathbf Z}\Lambda$ is the
product of the diagonal matrix of the images of the $d_i$ in $k$ with an invertible matrix. Its kernel
therefore has dimension the number $c$ of indices $i$ with $p\mid d_i$. Then $p^c$ divides
$d_1\cdots d_m$, hence $\det f=N(a)^{\rho_0}$, so $c\le v_p(N(a)^{\rho_0})=\rho_0\,v_p(N(a))$.
\end{proof}
% Lean: det_lsmul (det_eq_of_intertwine), charpoly_lsmul, corank_le (smithNormalFormBotBasis), in
% PositivityTheory.Xgn.AmalgamKernelBound.

\subsection{Hopf algebras and kernels of endomorphisms}\label{subsection-cv-hopf-kernels}

\begin{lemma}\label{cv-germs-free-local}
Let $k$ be a field, let $A$ be a finite-dimensional commutative $k$-algebra, let $U$ be an $A$-module
that is finite-dimensional over $k$, and let $r\ge0$. Let $t,\pi\in A$, and let $\chi_i:A\to k$, for $i$
in a finite set $I$, be homomorphisms of $k$-algebras. Suppose that the characteristic polynomial of
multiplication by $t$ on $U$ is the $r$-th power of that on $A$, that $\pi t$ is nilpotent, that the
values $\chi_i(t)$ are nonzero and pairwise distinct, and that the kernel of multiplication by $\pi$ on
$U$ has dimension at most $r|I|$. Then for every $j\in I$ there are $x_1,\dots,x_r\in U$ with
$U=Ax_1+\dots+Ax_r+\ker(\chi_j)\,U$.
\end{lemma}
\leanref{Development/WeilBound/CoreStatements/TorusUniformisation/ChowMahlerFamilies/KernelRankAlgebraicBound.lean}{PositivityTheory.Xgn.KernelRankAlgebraic.core}
\begin{proof}
Fitting's lemma. Let $u\in A$. As $A$ is artinian, the chain of ideals $u^mA$ is stationary, so
$u^n=u^{n+1}y$ for some $n\ge0$ and $y\in A$. Then $u^n(uy)^m=u^n$ for every $m\ge0$. Let
$\varepsilon=1-(uy)^n$. Then $u^n\varepsilon=0$, and for every $M\ge0$,
$1-\varepsilon=u^ny^n=u^n(uy)^My^n\in Au^M$.

Eigenspaces. For $i\in I$ let $\alpha_i=\chi_i(t)$, and let $G_i\subset U$ and $G^A_i\subset A$ be the
generalized eigenspaces of $t$ for $\alpha_i$: the elements killed by a power of $t-\alpha_i$. Let
$d_i=\dim G^A_i$. Multiplication by $\pi$ commutes with $t$, so it preserves $G_i$; let $f_i$ be its
restriction to $G_i$. The dimension of a generalized eigenspace is the multiplicity of the eigenvalue
as a root of the characteristic polynomial, so $\dim G_i=r\,d_i$.

Lower bound. Fix $i$, and let $\varepsilon$ and $n$ be as above for $u=t-\alpha_i$. Then
$\varepsilon\in G^A_i$, and $\varepsilon\ne0$: otherwise $1=cu$ for some $c\in A$, and applying
$\chi_i$ gives $1=0$. So $d_i\ge1$. Let $(\pi t)^N=0$. As $\alpha_i\ne0$, the elements $t$ and
$t-\alpha_i$ generate the unit ideal of $A$, and so do $t^N$ and $(t-\alpha_i)^m$ for every $m$, say
$at^N+b(t-\alpha_i)^m=1$. For $v\in G^A_i$ killed by $(t-\alpha_i)^m$ this gives
$\pi^Nv=a(\pi t)^Nv+b\pi^N(t-\alpha_i)^mv=0$. So $\pi$ acts nilpotently on the $d_i$-dimensional space
$G^A_i$, and $\pi^{d_i}G^A_i=0$. Let $v\in G_i$ be killed by $(t-\alpha_i)^M$. As
$1-\varepsilon\in A(t-\alpha_i)^M$, $\varepsilon v=v$, and $\pi^{d_i}v=(\pi^{d_i}\varepsilon)v=0$. So
$f_i^{d_i}=0$. If $\phi$ is an endomorphism of a finite-dimensional vector space $V$ with $\phi^e=0$,
then $\dim V\le e\dim\ker\phi$: the space $V/\phi(V)$ has dimension $s=\dim\ker\phi$, so
$V=kx_1+\dots+kx_s+\phi(V)$ for some $x_j$; the span $S$ of the $\phi^mx_j$, $m<e$, is stable under
$\phi$, and $V=S+\phi(V)=S+\phi^2(V)=\dots=S+\phi^e(V)=S$. Hence $r\,d_i=\dim G_i\le d_i\dim\ker f_i$,
and $\dim\ker f_i\ge r$.

Upper bound. The generalized eigenspaces $G_i$ belong to the distinct eigenvalues $\alpha_i$, so they
are independent, and $\ker f_i$ lies in $G_i$ and in the kernel of $\pi$ on $U$. So
$\sum_i\dim\ker f_i$ is at most the dimension of that kernel, hence at most $r|I|$, and with the lower
bounds, $\dim\ker f_j\le r$ for every $j\in I$.

Generators. Fix $j\in I$. As $\dim G_j/\pi G_j=\dim\ker f_j\le r$, there are $x_1,\dots,x_r\in G_j$ with
$G_j=kx_1+\dots+kx_r+\pi G_j$. Let $\varepsilon$ and $n$ be as above for $u=t-\alpha_j$, so that
$1-\varepsilon=cu$ for some $c\in A$. The value $\chi_j(\pi)\chi_j(t)=\chi_j(\pi t)$ is nilpotent in
$k$, so it is $0$, and $\chi_j(t)\ne0$; hence $\chi_j(\pi)=0$. Also $\chi_j(1-\varepsilon)=\chi_j(c)\chi_j(u)
=0$. Let $w\in U$. Then $\varepsilon w\in G_j$, as $u^n\varepsilon=0$, so $\varepsilon w=s+\pi v$ with
$s\in kx_1+\dots+kx_r$ and $v\in G_j$. So $w=s+\pi v+(1-\varepsilon)w$ lies in
$Ax_1+\dots+Ax_r+\ker(\chi_j)\,U$.
\end{proof}
% Lean: fitting, mem_G, mapsTo_G, pow_finrank_kill, finrank_le_of_nil (JordanBasis.exists_span_sup_range,
% JordanBasis.span_eq_top), finrank_G, sum_finrank_le, in PositivityTheory.Xgn.KernelRankAlgebraic.

\begin{lemma}\label{cv-germs-free-primes}
Let $E$ be a number field, let $p$ be a prime, and let $\mathfrak w$ be a maximal ideal of
$\mathcal O_E$ that contains $p$, with residue field $F=\mathcal O_E/\mathfrak w$ of order $p^f$. Let
$N(a)$ be the norm of $a\in\mathcal O_E$ and $v_p$ the $p$-adic valuation.
\begin{enumerate}
\item There is a nonzero $\pi\in\mathfrak w$ with $v_p(N(\pi))=f$.
\item There is $\theta\in\mathcal O_E$ whose image in $F$ generates the multiplicative group of $F$, such
  that for every $\pi\in\mathfrak w$ some power of $\pi\theta$ lies in $p\mathcal O_E$.
\item If $k$ is an algebraically closed field of characteristic $p$, there are $f$ ring
  homomorphisms from $F$ to $k$.
\end{enumerate}
\end{lemma}
\leanref{Development/WeilBound/CoreStatements/TorusUniformisation/ChowMahlerFamilies/KernelRankAlgebraicBound.lean}{PositivityTheory.Xgn.KernelRankAlgebraic.arith_pi}
\leanref{Development/WeilBound/CoreStatements/TorusUniformisation/ChowMahlerFamilies/KernelRankAlgebraicBound.lean}{PositivityTheory.Xgn.KernelRankAlgebraic.arith_theta}
\leanref{Development/WeilBound/CoreStatements/TorusUniformisation/ChowMahlerFamilies/KernelRankAlgebraicBound.lean}{PositivityTheory.Xgn.KernelRankAlgebraic.count_hom}
\begin{proof}
(1) As $\mathcal O_E$ is a Dedekind domain and $p\mathfrak w\subset\mathfrak w$ is a nonzero ideal, there
is $\pi\in\mathcal O_E$ with $p\mathfrak w+\pi\mathcal O_E=\mathfrak w$. Then $\pi\mathcal O_E\subset
\mathfrak w$, so $\pi\mathcal O_E=\mathfrak w\mathfrak a$ for an ideal $\mathfrak a$, and
$\mathfrak w(p\mathcal O_E+\mathfrak a)=\mathfrak w$ gives $p\mathcal O_E+\mathfrak a=\mathcal O_E$. So
$\mathfrak a\ne0$, as $p$ is not a unit, and $\pi\ne0$. The order of $\mathcal O_E/\mathfrak a$ is prime to
$p$, since $p$ is invertible modulo $\mathfrak a$. So $|N(\pi)|=|\mathcal O_E/\mathfrak w|\,
|\mathcal O_E/\mathfrak a|$ has $p$-adic valuation $f$.

(2) Write $p\mathcal O_E=\mathfrak w^c\mathfrak q$ with $\mathfrak w+\mathfrak q=\mathcal O_E$. Let
$g\in\mathcal O_E$ map to a generator of the cyclic group $F^\times$, choose $z\in\mathfrak q$ with
$z\equiv1$ modulo $\mathfrak w$, and let $\theta=gz$, which maps to the same generator. For
$\pi\in\mathfrak w$, $(\pi\theta)^{c+1}=\pi^c\cdot\pi g^{c+1}z^c\cdot z$ lies in
$\mathfrak w^c\mathfrak q=p\mathcal O_E$.

(3) The field $F$ has characteristic $p$, so the ring homomorphisms $F\to k$ are the homomorphisms of
$\mathbf F_p$-algebras. As $F$ is a separable extension of $\mathbf F_p$ of degree $f$ and $k$ is
algebraically closed, there are $f$ of them.
\end{proof}

\begin{lemma}\label{cv-germs-free-criterion}
Let $E$ be a number field of degree $e$ over $\mathbf Q$, with a $\mathbf Z$-basis $b_1,\dots,b_e$ of
$\mathcal O_E$, let $p$ be a prime, let $k$ be an algebraically closed field of characteristic $p$, and
let $r\ge0$. Let $N(a)$ be the norm of $a\in\mathcal O_E$ and $v_p$ the $p$-adic valuation. Let $U$ be a
finite-dimensional $k$-vector space and $\rho:\mathcal O_E\to\operatorname{End}_kU$ a ring homomorphism
such that, for every $a\in\mathcal O_E$, the characteristic polynomial of $\rho(a)$ is the $r$-th power
of that of multiplication by $a$ on $\mathcal O_E$ and, if $a\ne0$, $\dim\ker\rho(a)\le r\,v_p(N(a))$.
Then there are $x_1,\dots,x_r\in U$ such that the $re$ vectors $\rho(b_q)x_i$ form a basis of $U$. In
particular $U$ is free of rank $r$ over $\mathcal O_E\otimes k$.
\end{lemma}
\leanref{Development/AutomorphicSpectrum/Proofs/AlgebraProofs/KernelRankBoundNormValuationStatement.lean}{GlobalTheory.T.exists_free_basis_of_kernel_rank_le_norm_valuation}
\leanref{Development/WeilBound/CoreStatements/TorusUniformisation/ChowMahlerFamilies/KernelRankAlgebraicBound.lean}{PositivityTheory.Xgn.kernel_rank_algebraic_bound}
\begin{proof}
Let $A=k\otimes_{\mathbf Z}\mathcal O_E$, a commutative $k$-algebra of dimension $e$, and make $U$ an
$A$-module, with $c\otimes a$ acting as $c\rho(a)$. The ring $A$ is artinian, so it has finitely many
maximal ideals, and they are pairwise coprime.

Local generators. Let $\mathfrak m$ be a maximal ideal of $A$. Its residue field is finite over the
algebraically closed field $k$, so $\mathfrak m$ is the kernel of a homomorphism of $k$-algebras
$\chi_0:A\to k$. The kernel $\mathfrak w$ of $a\mapsto\chi_0(1\otimes a)$ is a prime ideal of
$\mathcal O_E$ that contains $p$, so it is nonzero and maximal; let $F=\mathcal O_E/\mathfrak w$, of order
$p^f$, and
let $\pi$ and $\theta$ be as in Lemma~\ref{cv-germs-free-primes}(1) and~(2). For each of the $f$ ring
homomorphisms $\iota:F\to k$ (Lemma~\ref{cv-germs-free-primes}(3)) let $\chi_\iota:A\to k$ be the
homomorphism of $k$-algebras $c\otimes a\mapsto c\,\iota(a\bmod\mathfrak w)$. For the homomorphism
$\iota_0:F\to k$ induced by $a\mapsto\chi_0(1\otimes a)$, $\chi_{\iota_0}=\chi_0$.
Lemma~\ref{cv-germs-free-local} applies to $t=1\otimes\theta$, to $1\otimes\pi$ and to these $f$
characters, with $j=\iota_0$:
\begin{itemize}
\item the characteristic polynomial of $t$ on $U$ is that of $\rho(\theta)$, the $r$-th power of that of
  $\theta$ on $\mathcal O_E$, which is the characteristic polynomial of $t$ on $A$;
\item $(1\otimes\pi)t=1\otimes\pi\theta$ is nilpotent, as a power of $\pi\theta$ lies in $p\mathcal O_E$
  and $p=0$ in $A$;
\item $\chi_\iota(t)=\iota(\theta\bmod\mathfrak w)$ is nonzero, and it determines $\iota$, as
  $\theta\bmod\mathfrak w$ generates $F^\times$; so these values are pairwise distinct;
\item the kernel of $1\otimes\pi$ on $U$ is $\ker\rho(\pi)$, of dimension at most
  $r\,v_p(N(\pi))=rf$.
\end{itemize}
So $U=Ay^{\mathfrak m}_1+\dots+Ay^{\mathfrak m}_r+\mathfrak mU$ for some $y^{\mathfrak m}_i\in U$.

Gluing. By the Chinese remainder theorem there are $c_{\mathfrak m}\in A$, for the maximal ideals
$\mathfrak m$, with $c_{\mathfrak m}-1\in\mathfrak m$ and $c_{\mathfrak m}\in\mathfrak m'$ for
$\mathfrak m'\ne\mathfrak m$. Let $x_i=\sum_{\mathfrak m}c_{\mathfrak m}y^{\mathfrak m}_i$ and
$N=Ax_1+\dots+Ax_r$. For each $\mathfrak m$, $y^{\mathfrak m}_i=x_i-\sum_{\mathfrak m'}
(c_{\mathfrak m'}-\delta_{\mathfrak m\mathfrak m'})y^{\mathfrak m'}_i$ lies in $N+\mathfrak mU$, so
$U=N+\mathfrak mU$. By Nakayama's lemma there is then $a\in A$ with $a-1\in\mathfrak m$ and $aU\subset N$.
So the ideal of the $a\in A$ with $aU\subset N$ lies in no maximal ideal; it is $A$, and $U=N$.

The basis. Every $a\in\mathcal O_E$ is an integral combination of the $b_q$, so $\rho(a)x_i$ lies in the
$k$-span $S$ of the vectors $\rho(b_q)x_i$, and so does $(c\otimes a)x_i=c\rho(a)x_i$. Hence $U=N\subset S$.
The characteristic polynomial of $\rho(0)$ has degree $\dim U$ and is the $r$-th power of a polynomial of
degree $e$, so $\dim U=re$. So the $re$ vectors $\rho(b_q)x_i$ form a basis, and
$(c_i)\mapsto\sum_ic_ix_i$ is an isomorphism $(\mathcal O_E\otimes k)^r\to U$.
\end{proof}
% Lean: kernel_rank_algebraic_bound: the module structure is Module.compHom through
% Algebra.TensorProduct.lift; locmax (exists_char, arith_pi, arith_theta, count_hom, core); glue
% (Ideal.pi_quotient_surjective, Submodule.exists_sub_one_mem_and_smul_le_of_fg_of_le_sup).

\begin{definition}\label{cv-germs-hopf}
Let $k$ be a field and let $C$ be a finite-dimensional commutative and cocommutative Hopf algebra over
$k$, with comultiplication $\Delta$, counit $\varepsilon$ and antipode $S$.
\begin{enumerate}
\item An element $f$ of $C$ is \emph{primitive} if $\Delta(f)=f\otimes1+1\otimes f$, and
  \emph{grouplike} if $\Delta(f)=f\otimes f$ and $\varepsilon(f)=1$.
\item A \emph{derivation at the counit} of $C$ is a linear form $v$ on $C$ with
  $v(fh)=v(f)\varepsilon(h)+\varepsilon(f)v(h)$ for all $f,h\in C$.
\item The \emph{Cartier dual} $C^\vee$ is the dual space $\operatorname{Hom}_k(C,k)$ with the product
  $(ww')(f)=(w\otimes w')(\Delta(f))$ and the unit $\varepsilon$, the comultiplication dual to the
  multiplication of $C$, the counit $w\mapsto w(1)$ and the antipode $w\mapsto w\circ S$.
\end{enumerate}
\end{definition}
\leanref{Development/AbelianVarietyTheory/Definitions/ShimuraVarieties/MemPrimitiveTensorsIffIsPrimitiveElem.lean}{GlobalTheory.CartierDual}
\leanref{Development/AbelianVarietyTheory/Definitions/ShimuraVarieties/MemPrimitiveTensorsIffIsPrimitiveElem.lean}{GlobalTheory.primitiveElems}

\begin{lemma}\label{cv-germs-hopf-dual}
Let $k$ be a field and let $C$ be a finite-dimensional commutative and cocommutative Hopf algebra over
$k$, with comultiplication $\Delta$, counit $\varepsilon$ and antipode $S$
(Definition~\ref{cv-germs-hopf}).
\begin{enumerate}
\item The Cartier dual $C^\vee$ is a finite-dimensional commutative and cocommutative Hopf algebra over
  $k$, and $f\mapsto(w\mapsto w(f))$ is an isomorphism of Hopf algebras $C\to(C^\vee)^\vee$.
\item An element of $C^\vee$ is primitive if and only if it is a derivation at the counit of $C$.
\item If $f\in C$ is primitive, then $\varepsilon(f)=0$, $S(f)=-f$, and
  $(ww')(f)=w(f)w'(1)+w(1)w'(f)$ for $w,w'\in C^\vee$.
\item Every homomorphism of $k$-algebras $\chi:C^\vee\to k$ is $w\mapsto w(g)$ for a grouplike element
  $g$ of $C$, and $g$ is a unit of $C$.
\item Let $I$ be an ideal of $C$ with $\varepsilon(I)=0$, $S(I)\subset I$ and
  $\Delta(I)\subset I\otimes C+C\otimes I$. Then $C/I$, with the induced comultiplication, counit and
  antipode, is a commutative and cocommutative Hopf algebra over $k$.
\end{enumerate}
\end{lemma}
\leanref{Development/AbelianVarietyTheory/Definitions/ShimuraVarieties/MemPrimitiveTensorsIffIsPrimitiveElem.lean}{GlobalTheory.CartierDual.bidual}
\leanref{Development/AbelianVarietyTheory/Definitions/ShimuraVarieties/MemPrimitiveTensorsIffIsPrimitiveElem.lean}{GlobalTheory.CartierDual.comul_eq_iff}
\leanref{Development/AutomorphicSpectrum/Proofs/GeometryProofs/PrimitiveDerivationPairing.lean}{GlobalTheory.CompanionTheory.Ar.HopfCount.mul_apply_prim}
\leanref{Development/AutomorphicSpectrum/Proofs/GeometryProofs/PrimitiveDerivationPairing.lean}{GlobalTheory.CompanionTheory.Ar.HopfCount.antipode_prim}
\leanref{Development/AutomorphicSpectrum/Proofs/GeometryProofs/PrimitiveDerivationPairing.lean}{GlobalTheory.CompanionTheory.Ar.HopfCount.point_eval}
\leanref{Development/AutomorphicSpectrum/Proofs/GeometryProofs/PrimitiveDerivationPairing.lean}{GlobalTheory.CompanionTheory.Ar.HopfCount.isCocomm_quot}
\begin{proof}
(1) and (2) are Cartier duality \cite[\S2.4]{Wat79}. The comultiplication of $C^\vee$ maps $w$ to
$w\circ m$, where $m:C\otimes C\to C$ is the multiplication, so $w$ is primitive if and only if
$w\circ m=w\otimes\varepsilon+\varepsilon\otimes w$, which is (2). Part (5) is \cite[\S2.1]{Wat79}; the
quotient is cocommutative because $C$ is.

(3) Applying $\varepsilon\otimes1$ to $\Delta(f)=f\otimes1+1\otimes f$ gives $f=\varepsilon(f)+f$, so
$\varepsilon(f)=0$. The identity $m\circ(S\otimes1)\circ\Delta=\varepsilon$ gives $S(f)+f=\varepsilon(f)=0$.
The formula for $(ww')(f)$ is the definition of the product.

(4) By (1) there is $g\in C$ with $\chi(w)=w(g)$ for every $w\in C^\vee$. As $\chi$ is multiplicative,
$(w\otimes w')(\Delta(g))=w(g)w'(g)$ for all $w,w'\in C^\vee$, so $\Delta(g)=g\otimes g$; and
$\varepsilon(g)=\chi(1)=1$. The identity $m\circ(S\otimes1)\circ\Delta=\varepsilon$ gives $S(g)g=1$.
\end{proof}

\begin{lemma}\label{cv-germs-hopf-frobenius}
Let $k$ be an algebraically closed field of characteristic $p>0$, let $A$ be a finite-dimensional
commutative $k$-algebra, and let $V\ne0$ be a subspace of $A$ with $f^p\in V$ for every $f\in V$. Then
there is a nonzero $f\in V$ with $f^p=0$ or $f^p=f$.
\end{lemma}
\leanref{Development/AutomorphicSpectrum/Proofs/GeometryProofs/FrobeniusStableSubspaceElementStatement.lean}{GlobalTheory.T.frobenius_stable_subspace_element}
\begin{proof}
Let $f\in V$ be nonzero. The proof is by induction on $\dim_kfA$.

If $f$ is nilpotent, let $m$ be the least integer with $f^{p^m}=0$. Then $m\ge1$, and
$g=f^{p^{m-1}}$ lies in $V$, is nonzero and has $g^p=0$.

Otherwise $f$ lies outside some prime ideal of $A$. As $A$ is artinian, this ideal is maximal; its residue
field is finite over $k$, hence equal to $k$. So there is a homomorphism of $k$-algebras $\chi:A\to k$
with $\chi(f)\ne0$. Let $f'=\chi(f)^{-1}f$, an element of $V$ with $\chi(f')=1$, and let $g=f'^p-f'\in V$.
If $g=0$, then $f'^p=f'$. Otherwise $\chi(g)=0$, and $g=f\cdot\chi(f)^{-1}(f'^{p-1}-1)$, so $gA\subset fA$.
As $\chi$ vanishes on $gA$ and not at $f$, $f\notin gA$. So $\dim_kgA<\dim_kfA$, and the induction
hypothesis applies to $g$.
\end{proof}
% Lean: HopfFrob.exists_frobenius_fixed_or_killed (exists_pow_eq_zero_of_isNilpotent,
% exists_algHom_ne_zero_of_not_isNilpotent; the induction is on the rank of multiplication by f).

\begin{lemma}\label{cv-germs-hopf-unit}
Let $k$ be an algebraically closed field, let $R$ be a finite-dimensional commutative $k$-algebra, and let
$U$ be a subspace of $R$ such that for every homomorphism of $k$-algebras $\chi:R\to k$ there is $t\in U$
with $\chi(t)\ne0$. Then $U$ contains a unit of $R$.
\end{lemma}
\leanref{Development/AutomorphicSpectrum/Proofs/GeometryProofs/SubspaceContainsUnitStatement.lean}{GlobalTheory.T.subspace_contains_unit}
\begin{proof}
The ring $R$ is artinian, so it has finitely many maximal ideals. Each is the kernel of a homomorphism of
$k$-algebras $\chi:R\to k$, as its residue field is finite over $k$. By hypothesis, for each maximal
ideal $\mathfrak m$ the subspace $U\cap\mathfrak m$ of $U$ is proper. As $k$ is infinite, a vector space
over $k$ is not a finite union of proper subspaces, so some $w\in U$ lies in no maximal ideal of $R$.
Then $w$ is a unit.
\end{proof}

\begin{lemma}\label{cv-germs-hopf-powers}
Let $k$ be a field of characteristic $p>0$, let $C$ be a finite-dimensional commutative and
cocommutative Hopf algebra over $k$, and let $f\ne0$ be a primitive element of $C$
(Definition~\ref{cv-germs-hopf}). Then $1,f,\dots,f^{p-1}$ are linearly independent over $k$.
\end{lemma}
\leanref{Development/AutomorphicSpectrum/Proofs/GeometryProofs/PrimitiveElementPowersIndepStatement.lean}{GlobalTheory.T.primitive_element_powers_indep}
\begin{proof}
A linear form. By Lemma~\ref{cv-germs-hopf-dual}(3), $\varepsilon(f)=0$. So $f$ is not a multiple of $1$:
if $f=c\cdot1$, then $c=\varepsilon(f)=0$. Hence there is a linear form $y$ on $C$ with $y(1)=0$ and
$y(f)=1$.

Convolution powers. Let $y_0=\varepsilon$ and $y_{m+1}(h)=(y\otimes y_m)(\Delta(h))$. As $\Delta$ is
multiplicative, $\Delta(f^i)=\sum_{l\le i}\binom ilf^l\otimes f^{i-l}$, so
$y_{m+1}(f^i)=\sum_{l\le i}\binom ily(f^l)\,y_m(f^{i-l})$. By induction on $m$, $y_m(f^i)=0$ for $i<m$
and $y_m(f^m)=m!$. Indeed, in the sum for $y_{m+1}(f^i)$ with $i\le m+1$, the term $l=0$ vanishes as
$y(1)=0$, the terms with $l\ge2$ vanish as $i-l<m$, and the term $l=1$ is $i\,y_m(f^{i-1})$, which is $0$
for $i\le m$ and $(m+1)\,m!$ for $i=m+1$.

Independence. Suppose that $\sum_{i<p}c_if^i=0$. For $m<p$, if $c_i=0$ for $m<i<p$, then $c_m=0$.
Indeed, applying $y_m$ to the sum kills the terms with $i<m$, by the preceding step, and those with
$i>m$, by assumption; so $c_m\,m!=0$. As $m<p$, $m!\ne0$ in $k$, and $c_m=0$. By descending induction on
$m$, starting at $m=p-1$, every $c_m$ is $0$.
\end{proof}
% Lean: HopfPrim.powers_indep_of_primitive (exists_form_of_primitive, counit_eq_zero_of_primitive, convPow,
% comul_pow_of_primitive, convPow_pow); the step is key, and the descending induction is hall, by
% induction on n for the indices i with p - 1 - i < n.

\begin{lemma}\label{cv-germs-hopf-derivation}
Let $k$ be a field of characteristic $p>0$, let $R$ be a finite-dimensional commutative $k$-algebra with a
$k$-linear derivation $D:R\to R$, and let $\chi:R\to k$ be a homomorphism of $k$-algebras. Suppose that
one of the following holds:
\begin{enumerate}
\item[(a)] $D^p=0$, and $D^{p-1}(t)$ is a unit of $R$ for some $t\in R$;
\item[(b)] $k$ is algebraically closed, $D^p=D$, and $D(z_0)=z_0$ for some unit $z_0$ of $R$.
\end{enumerate}
Then $\dim_kR=p\dim_k\ker D$, and there is $z\in R$ with $\chi(D(z))\ne0$ such that every linear form
$u$ on $R$ with $u(ab)=u(a)\chi(b)+\chi(a)u(b)$ for all $a,b\in R$, $u(\ker D)=0$ and $u(z)=0$ is zero.
\end{lemma}
\leanref{Development/AutomorphicSpectrum/Proofs/GeometryProofs/NilpotentDerivationRankStatement.lean}{GlobalTheory.T.nilpotent_derivation_rank}
\leanref{Development/AutomorphicSpectrum/Proofs/GeometryProofs/MultiplicativeDerivationRankStatement.lean}{GlobalTheory.T.multiplicative_derivation_rank}
\begin{proof}
In both cases, as shown below, there is $z\in R$ with $\chi(D(z))\ne0$ such that
$(c_0,\dots,c_{p-1})\mapsto\sum_ic_iz^i$ is a bijection $(\ker D)^p\to R$. This gives the dimension. Let
$u$ be as in the statement. Then $u(1)=0$, as $1\in\ker D$, and
$u(z^{i+1})=u(z^i)\chi(z)+\chi(z^i)u(z)$, so $u(z^i)=0$ for all $i$ by induction. Then
$u(cz^i)=u(c)\chi(z^i)+\chi(c)u(z^i)=0$ for $c\in\ker D$, and $u=0$ by the bijection.

(a) Let $w=D^{p-1}(t)$. Then $D(w)=D^p(t)=0$, and $D(w^{-1})=0$, from $D(ww^{-1})=0$. Let
$z=w^{-1}D^{p-2}(t)$. Then $D(z)=1$ and $\chi(D(z))=1$. For $c\in\ker D$ and $i,n\ge0$,
$D^n(cz^i)=i(i-1)\cdots(i-n+1)\,cz^{i-n}$. Injectivity: if $\sum_{i<m}c_iz^i=0$ with $m\le p$, applying
$D^{m-1}$ gives $(m-1)!\,c_{m-1}=0$, so $c_{m-1}=0$ as $(m-1)!\ne0$ in $k$, and all $c_i$ vanish by
descending induction. Surjectivity: by induction on $j\le p$, every $a$ with $D^j(a)=0$ is
$\sum_{i<j}c_iz^i$ with $c_i\in\ker D$. If $D^{j+1}(a)=0$, then $b=D^j(a)$ lies in $\ker D$; with
$e=b/j!$, $D^j(a-ez^j)=b-j!\,e=0$, and $a-ez^j$ is such a sum. As $D^p=0$, this covers all of $R$.

(b) Let $z=z_0$ and $w=z^{-1}$. Then $D(z^n)=nz^n$ and $D(w^n)=-nw^n$, and $\chi(D(z))=\chi(z)\ne0$. The
polynomial $X^p-X$ is separable and kills $D$, so $D$ is diagonalizable with eigenvalues among its roots,
which are the images in $k$ of the integers $0,1,\dots,p-1$; let $R_i$ be the eigenspace of $D$ for $i$.
For $c\in\ker D$, $cz^i\in R_i$. Injectivity: if $\sum_{i<p}c_iz^i=0$, then each $c_iz^i$ is $0$, as
eigenvectors for distinct eigenvalues are independent, and $c_i=c_iz^iw^i=0$. Surjectivity: for
$v\in R_i$, $D(vw^i)=ivw^i-ivw^i=0$, and $v=(vw^i)z^i$.
\end{proof}
% Lean: (a) is Ggeo.DerivFree.finrank_and_vanishing (exists_deriv_eq_one, iter_mul_pow, coeff_zero,
% exists_coeff, coeffMap_bijective); (b) is PowerDerivation.finrank_eq_and_derivation_eq_zero
% (iSup_eigenspace_eq_top, exists_natCast_of_pow_eq, powerSum_injective, powerSum_surjective,
% map_powerSum_eq_zero).

\begin{lemma}\label{cv-germs-hopf-quotient}
Let $k$ be an algebraically closed field of characteristic $p>0$, let $C$ be a finite-dimensional
commutative and cocommutative Hopf algebra over $k$, and let $f\ne0$ be a primitive element of $C$
(Definition~\ref{cv-germs-hopf}) with $f^p=0$ or $f^p=f$. Then $C/fC$, with the induced structure, is a
commutative and cocommutative Hopf algebra over $k$, $\dim_kC=p\dim_kC/fC$, and every primitive element
of $C$ that lies in $fC$ is $\lambda f$ for some $\lambda\in k$.
\end{lemma}
\leanref{Development/AutomorphicSpectrum/Proofs/GeometryProofs/PrimitiveDerivationPairing.lean}{GlobalTheory.CompanionTheory.Ar.HopfCount.key}
\leanref{Development/AutomorphicSpectrum/Proofs/GeometryProofs/PrimitiveDerivationPairing.lean}{GlobalTheory.CompanionTheory.Ar.HopfCount.isHopfIdeal_span}
\begin{proof}
The quotient. Let $\varepsilon$ and $S$ be the counit and antipode of $C$. For $c\in C$,
$\varepsilon(cf)=\varepsilon(c)\varepsilon(f)=0$, $\Delta(cf)=\Delta(c)(f\otimes1+1\otimes f)$ lies in
$fC\otimes C+C\otimes fC$, and $S(cf)=-S(c)f$, by Lemma~\ref{cv-germs-hopf-dual}(3). So $C/fC$ is a Hopf
algebra as stated, by Lemma~\ref{cv-germs-hopf-dual}(5).

The derivation. Let $D$ be the linear map of $C^\vee$ with $D(w)(h)=w(fh)$. It is a derivation, as
$D(ww')(h)=(w\otimes w')((f\otimes1+1\otimes f)\Delta(h))=(D(w)w'+wD(w'))(h)$, and by induction
$D^n(w)(h)=w(f^nh)$. Let $\chi:C^\vee\to k$ be the counit $w\mapsto w(1)$ of $C^\vee$, a homomorphism of
$k$-algebras; then $\chi(D(w))=w(f)$.

Case $f^p=0$. Then $D^p=0$, and $f^{p-1}\ne0$ by Lemma~\ref{cv-germs-hopf-powers}. The image of
$D^{p-1}$ satisfies the hypothesis of Lemma~\ref{cv-germs-hopf-unit}. Indeed, a homomorphism of
$k$-algebras $C^\vee\to k$ is $w\mapsto w(g)$ for a unit $g$ of $C$, by
Lemma~\ref{cv-germs-hopf-dual}(4); as $f^{p-1}g\ne0$, there is $w$ with $w(f^{p-1}g)\ne0$, and this is the
value of $D^{p-1}(w)$ at that homomorphism. So $D^{p-1}(t)$ is a unit for some $t\in C^\vee$, and
Lemma~\ref{cv-germs-hopf-derivation}(a) applies to $D$ and $\chi$.

Case $f^p=f$. Then $D^p=D$. Let $u_0=\sum_{i=1}^{p-1}f^i$. Then $fu_0=u_0$, as $f^p=f$, and $u_0\ne0$ by
Lemma~\ref{cv-germs-hopf-powers}. The kernel of $D-1$ satisfies the hypothesis of
Lemma~\ref{cv-germs-hopf-unit}: for a unit $g$ as before there is $w$ with $w(u_0g)\ne0$, and
$w'(h)=w(u_0h)$ defines $w'\in C^\vee$ with $D(w')=w'$, as $fu_0=u_0$, whose value at the homomorphism is
$w(u_0g)\ne0$. So $\ker(D-1)$ contains a unit, and Lemma~\ref{cv-germs-hopf-derivation}(b) applies to $D$
and $\chi$.

The dimension. In both cases $\dim_kC^\vee=p\dim_k\ker D$, and there is $z\in C^\vee$ with
$z(f)=\chi(D(z))\ne0$ and the vanishing property of Lemma~\ref{cv-germs-hopf-derivation}. An element $w$
of $C^\vee$ lies in $\ker D$ if and only if it vanishes on $fC$, so
$\dim_k\ker D=\dim_kC-\dim_kfC=\dim_kC/fC$. As $\dim_kC^\vee=\dim_kC$, this gives $\dim_kC=p\dim_kC/fC$.

Primitive elements. Let $bf$ be primitive, let $\lambda=z(bf)/z(f)$, and let $u(w)=w(bf)-\lambda w(f)$ for
$w\in C^\vee$. As $bf$ and $f$ are primitive, $u(ww')=u(w)\chi(w')+\chi(w)u(w')$ by
Lemma~\ref{cv-germs-hopf-dual}(3). For $w\in\ker D$, $w(f)=w(bf)=0$, so $u(w)=0$; and $u(z)=0$ by the
choice of $\lambda$. So $u=0$: every $w\in C^\vee$ vanishes at $bf-\lambda f$, and $bf=\lambda f$.
\end{proof}
% Lean: key, with dualDer (dualDer_apply, dualDer_iter), point_eval, exists_dual_ne, mul_apply_prim, and
% Subspace.finrank_add_finrank_dualAnnihilator_eq for the dimension; the quotient is isHopfIdeal_span with
% isCocomm_quot.

\begin{lemma}\label{cv-germs-hopf-count}
Let $k$ be an algebraically closed field of characteristic $p>0$, let $C$ be a finite-dimensional
commutative and cocommutative Hopf algebra over $k$, let $Y$ be a space of primitive elements of $C$, and
let $Y'$ be a space of derivations at the counit of $C$ (Definition~\ref{cv-germs-hopf}) such that
$v(f)=0$ for $v\in Y'$ and $f\in Y$. Then $p^{\dim Y+\dim Y'}$ divides $\dim_kC$.
\end{lemma}
\leanref{Development/AutomorphicSpectrum/Proofs/GeometryProofs/PrimitiveDerivationPairingStatement.lean}{GlobalTheory.T.primitive_derivation_pairing}
\leanref{Development/AutomorphicSpectrum/Proofs/GeometryProofs/PrimitiveDerivationPairing.lean}{GlobalTheory.CompanionTheory.Ar.HopfCount.primitive_tangent_dvd_finrank}
\begin{proof}
The proof is by induction on $\dim_kC$. Let $\varepsilon$ be the counit of $C$.

Case $Y\ne0$. Let $Q$ be the space of primitive elements $f$ of $C$ with $v(f)=0$ for every $v\in Y'$.
Then $Y\subset Q$, so $Q\ne0$. The space $Q$ is stable under $f\mapsto f^p$: for $f\in Q$,
$\Delta(f^p)=(f\otimes1+1\otimes f)^p=f^p\otimes1+1\otimes f^p$, and for $v\in Y'$,
$v(f^p)=v(f^{p-1})\varepsilon(f)+\varepsilon(f^{p-1})v(f)=0$, as $\varepsilon(f)=0$ by
Lemma~\ref{cv-germs-hopf-dual}(3). By Lemma~\ref{cv-germs-hopf-frobenius} there is a nonzero $f\in Q$
with $f^p=0$ or $f^p=f$. By Lemma~\ref{cv-germs-hopf-quotient}, $\bar C=C/fC$ is a commutative and
cocommutative Hopf algebra with $\dim_kC=p\dim_k\bar C$, so that $\dim_k\bar C<\dim_kC$, and
$Q\cap fC=kf$. So the image $\bar Y$ of $Q$ in $\bar C$ is a space of primitive elements of dimension
$\dim Q-1$. Every $v\in Y'$ vanishes on $fC$, as $v(cf)=v(c)\varepsilon(f)+\varepsilon(c)v(f)=0$; so $v$ is
the composite of $C\to\bar C$ with a derivation $\bar v$ at the counit of $\bar C$, and $\bar v$ vanishes
on $\bar Y$. These $\bar v$ form a space $\bar Y'$ of dimension $\dim Y'$. By the induction hypothesis
for $\bar C$, $\bar Y$ and $\bar Y'$, $p^{\dim Q-1+\dim Y'}$ divides $\dim_k\bar C$. So
$p^{\dim Q+\dim Y'}$ divides $\dim_kC$, and so does $p^{\dim Y+\dim Y'}$.

Case $Y=0$ and $Y'\ne0$. By Lemma~\ref{cv-germs-hopf-dual}(1) and~(2), $C^\vee$ is a commutative and
cocommutative Hopf algebra of dimension $\dim_kC$, and $Y'$ is a space of primitive elements of $C^\vee$.
The argument of the preceding case, applied to $C^\vee$ with the spaces $Y'$ and $0$, uses the induction
hypothesis only for Hopf algebras of dimension less than $\dim_kC^\vee=\dim_kC$. It shows that
$p^{\dim Y'}$ divides $\dim_kC$.

Case $Y=Y'=0$. There is nothing to prove.
\end{proof}
% Lean: primitive_tangent_dvd_finrank = norm_index_count_bound (strong induction on the dimension), with stepP
% for the case Y nonzero (Q, frobenius_stable_subspace_element, key, the quotient Hopf algebra, PB and WB)
% and, for Y = 0, stepP applied to CartierDual k A with P' the image of W (CartierDual.comul_eq_iff).
% PrimitiveDerivationPairing.lean and the five lemmas it uses have about 1,100 lines.

\begin{lemma}\label{cv-germs-weil-flat}
Let $k$ be an algebraically closed field, let $G$ be a group scheme over $k$ that is integral and locally
of finite type over $k$, and let $\varphi:G\to G$ be a homomorphism of group schemes over $k$ that is a
finite morphism. Then $\varphi$ is flat and surjective.
\end{lemma}
\leanref{Development/AbelianVarietyTheory/Proofs/AbelianSchemeLawProofs/FlatAndSurjectiveOfIsFiniteAbelianEnd.lean}{GlobalTheory.flat_and_surjective_of_isFinite_grpEnd}
\begin{proof}
The scheme $G$ is locally noetherian and Jacobson, as it is locally of finite type over $k$. Let $\xi$ be its
generic point and $o$ its unit, a closed point with $\varphi(o)=o$.

The generic point. As $\varphi$ is finite, its fibres are discrete. So if a point $y$ of $G$ specializes to a
point $z\ne y$, then $\varphi(y)$ specializes to $\varphi(z)\ne\varphi(y)$: otherwise $y$ and $z$ lie in one
fibre, in which $y$ specializes to $z$, and $y=z$. For a point $z$ of $G$ that specializes to $o$, let $h(z)$
be the supremum of the lengths $n$ of chains $o=z_0,z_1,\dots,z_n=z$ in which each $z_{i+1}$ specializes to
$z_i\ne z_{i+1}$. It is at most the dimension of the noetherian local ring of $G$ at $o$, so it is finite.
The image under $\varphi$ of such a chain ending at $\xi$ is such a chain ending at $\varphi(\xi)$, as
$\varphi(o)=o$; so $h(\xi)\le h(\varphi(\xi))$. The point $\xi$ specializes to $\varphi(\xi)$. If
$\varphi(\xi)\ne\xi$, a chain ending at $\varphi(\xi)$ extends by $\xi$, and $h(\xi)>h(\varphi(\xi))$. So
$\varphi(\xi)=\xi$.

A flat open set. Let $U=\operatorname{Spec}A$ be a nonempty affine open subset of $G$. Then $A$ is a
noetherian domain, and $\varphi^{-1}(U)=\operatorname{Spec}B$ for a finite $A$-algebra $B$. Let $K$ be the
fraction field of $A$, and let $b_1,\dots,b_n\in B$ map to a basis of $K\otimes_AB$. The kernel and the
cokernel of the map $A^n\to B$ they define are finitely generated $A$-modules that vanish after tensoring
with $K$, so they are killed by some $r\ne0$ in $A$, and $B_r$ is free over $A_r$. Hence $\varphi$ is flat
at every point of $W=\varphi^{-1}(\operatorname{Spec}A_r)$. This open set contains $\xi$, as $\varphi(\xi)=\xi$ is the
generic point of $\operatorname{Spec}A$ and $r\ne0$.

Closed points. As $G$ is Jacobson, $W$ contains a closed point $y$. Every closed point of $G$ is a $k$-point,
as $k$ is algebraically closed. Let $x$ be a closed point, let $\alpha$ be the right translation of $G$ by
$y^{-1}x$, and let $\beta$ be the right translation by $\varphi(y^{-1}x)$. They are isomorphisms,
$\alpha(y)=x$, and $\varphi\circ\alpha=\beta\circ\varphi$, as $\varphi$ is a homomorphism. As $\varphi$ is
flat at $y$, so is $\beta\circ\varphi=\varphi\circ\alpha$; hence $\varphi$ is flat at $\alpha(y)=x$.

All points. A point $z$ of $G$ specializes to a closed point $x$, as $G$ is Jacobson. The local ring of $G$
at $z$ is a localization of that at $x$, so it is flat over the local ring of $G$ at $\varphi(x)$, and hence
over its localization at $\varphi(z)$. So $\varphi$ is flat at $z$.

Surjectivity. As $\varphi$ is finite, its image is closed. It contains $\xi=\varphi(\xi)$, so it is $G$.
\end{proof}
% Lean: flat_and_surjective_of_isFinite_grpEnd, with FiniteEndFlat.strictMono_of_locallyQuasiFinite and
% apply_genericPoint (heights in the generizations of o, bounded by ringKrullDim of the stalk),
% exists_flat_away (Module.FinitePresentation.exists_free_localizedModule_powers) and exists_open_flat, and
% UnramifiedEndFlat.flat_of_closedPoints with flat_stalkMap_of_specializes. About 150 lines.

\begin{lemma}\label{cv-germs-weil-isogeny}
Let $k$ be an algebraically closed field, let $G$ be a group scheme over $k$ that is proper, smooth and
geometrically connected over $k$, and let $\varphi:G\to G$ be a homomorphism of group schemes over $k$ whose
kernel, the fibre product of $\varphi$ with the unit section, is $\operatorname{Spec}A$ for a
finite-dimensional $k$-algebra $A$. Then $\varphi$ is finite, flat and surjective.
\end{lemma}
\leanref{Development/WeilBound/CoreStatements/TorusUniformisation/ChowMahlerFamilies/KernelRankMultiplicationCharP.lean}{PositivityTheory.Xgn.JetKernelRank.flat_surj}
\leanref{Development/AbelianVarietyTheory/Proofs/UnitaryExistence/LocallyQuasiFiniteOfFiniteGeometricPoints.lean}{GlobalTheory.locallyQuasiFinite_of_finite_geometric_ker}
\begin{proof}
Properness. As $G$ is proper over $k$, so is $\varphi$.

Finite fibres. Let $\Omega$ be an algebraically closed field and $y$ a point of $G$ with values in $\Omega$.
If $x_0$ is a point of $G$ with values in $\Omega$ and $\varphi(x_0)=y$, then $x\mapsto x_0^{-1}x$ maps the
points $x$ with $\varphi(x)=y$ injectively to the points of the kernel with values in $\Omega$. These are the
homomorphisms of $k$-algebras $A\to\Omega$, finitely many as $A$ is finite-dimensional. So the fibre product
$P$ of $\varphi$ with $y$, a scheme locally of finite type over $\Omega$, has finitely many points with values
in $\Omega$. These are its closed points, as $\Omega$ is algebraically closed. Their union is closed, and as
$P$ is Jacobson it contains every point. So $P$ is finite. For a point $z$ of $G$, taking for $y$ the point
with values in an algebraic closure of the residue field at $z$ shows that the fibre of $\varphi$ over $z$,
the image of $P$, is finite. So $\varphi$ is locally quasi-finite.

Conclusion. As $\varphi$ is proper and locally quasi-finite, it is finite
\cite[Th\'eor\`eme~8.11.1]{Gro66}. As $G$ is smooth and connected over $k$, its local rings are regular,
hence normal domains, and $G$ is integral. Lemma~\ref{cv-germs-weil-flat} shows that $\varphi$ is flat and
surjective.
\end{proof}
% Lean: flat_surj (IsProper.of_comp, IsFinite.of_isProper_of_locallyQuasiFinite,
% grpObj_geometricallyIntegral_of_smooth); the fibres by locallyQuasiFinite_of_finite_geometric_ker with
% finite_fiber_of_finite_ker, finite_of_finite_sections and locallyQuasiFinite_of_finite_geometric_points, and
% pt_surj for the points of the kernel.

\begin{lemma}\label{cv-germs-weil}
Let $E$ be a CM field, let $k$ be an algebraically closed field, and let $G$ be an
$\mathcal O_E$-semi-abelian scheme over $\operatorname{Spec}k$
(Definition~\ref{cv-model-toroidal-semiabelian}) that is proper over $k$, with line bundle $\mathcal L$.
Let $g$ be a formal chart of dimension $n$ at the identity of $\operatorname{Spec}k$
(Definition~\ref{cv-germs-chart}), with formal germ $B=(F,[\,\cdot\,],\Phi_1,\Phi_2,S,P)$, points $x_m$,
generators $e_m$ and, for $a\in\mathcal O_E$, the generators $w_a$ of
Definition~\ref{cv-germs-chart}(3)(e). Fix $a\in\mathcal O_E$. Points of $G$ have values in schemes over
$k$, and the group law is written additively. For points $x$ and $y$ of $G$ with values in a scheme $T$,
with $ax=0$ and $\bar ay=0$, the line bundle
$\Lambda(x,\bar ay)\otimes\Lambda(ax,y)^{-1}=\Lambda(x,0)\otimes\Lambda(0,y)^{-1}$ is trivial by
Definition~\ref{cv-germs-chart}(1), and through this trivialization $w_a(x,y)$ is a unit $e_a(x,y)$ of
$\Gamma(T,\mathcal O_T)$, whose formation commutes with base change in $T$.
\begin{enumerate}
\item For such $x$ and $y$, $e_a(x,0)=1=e_a(0,y)$.
\item For such $x$ and $y$ and every point $x'$ of $G$ with values in $T$,
  $w_a(x'+x,y)=e_a(x,y)\,w_a(x',y)$, both read as generators of $\Lambda(ax',y)^{-1}$ through
  $\Lambda(x'+x,0)\cong\mathcal O_T\cong\Lambda(x',0)$. In particular $e_a(x'+x,y)=e_a(x',y)\,e_a(x,y)$
  if $ax'=0$.
\item Suppose that the kernel of $a$ on $G$ is finite over $k$ and that $K(\mathcal L)$ is trivial
  (Definition~\ref{cv-germs-principal}(2)). Let $R$ be a local $k$-algebra and $y$ a point of $G$ with
  values in $R$ with $\bar ay=0$, such that $e_a(x,y_T)=1$ for every scheme $T$ over $\operatorname{Spec}R$
  and every point $x$ of $G$ with values in $T$ with $ax=0$, where $y_T$ is the image of $y$. Then $y=0$.
\item Let $R$ be a $k$-algebra. For $a'\in R^n$ with $a'_ia'_j=0$ for all $i,j$, let $x(a')$ be the point
  of $G$ with values in $R$ obtained from $x_m$ by $t\mapsto a'$, for any $m\ge1$; then $x(0)=0$. If
  $a',b'\in R^n$ satisfy $a'_ia'_j=0=b'_ib'_j$ for all $i,j$, $\rho_B(a)a'=0$ and $\rho_B(\bar a)b'=0$,
  then $a\,x(a')=0$, $\bar a\,x(b')=0$ and $e_a(x(a'),x(b'))=P_a(a',b')$.
\end{enumerate}
\end{lemma}
\leanref{Development/WeilBound/CoreStatements/TorusUniformisation/ChowMahlerFamilies/KernelRankMultiplicationCharP.lean}{PositivityTheory.Xgn.JetKernelRank.chi_exists}
\leanref{Development/WeilBound/CoreStatements/TorusUniformisation/ChowMahlerFamilies/KernelRankMultiplicationCharP.lean}{PositivityTheory.Xgn.JetKernelRank.chi_one}
\leanref{Development/WeilBound/CoreStatements/TorusUniformisation/ChowMahlerFamilies/KernelRankMultiplicationCharP.lean}{PositivityTheory.Xgn.JetKernelRank.pair_mul}
\leanref{Development/WeilBound/CoreStatements/TorusUniformisation/ChowMahlerFamilies/KernelRankMultiplicationCharP.lean}{PositivityTheory.Xgn.JetKernelRank.pair_nondeg}
\leanref{Development/WeilBound/CoreStatements/TorusUniformisation/ChowMahlerFamilies/KernelRankMultiplicationCharP.lean}{PositivityTheory.Xgn.JetKernelRank.chi_chart}
\begin{proof}
(4) The point $x(a')$ does not depend on $m$, as $x_{m+1}$ reduces to $x_m$, and $x(0)$ is the unit, as
$x_m(0)$ is. Products of two of the $a'_i$, or of two of the $b'_j$, vanish, so $t\mapsto a'$,
$t'\mapsto b'$ defines a homomorphism $J^2_2\to R$, and substituting $a'$ into an $n$-tuple of power series
without constant term leaves its linear part. So $[a]a'=\rho_B(a)a'=0$ and $[\bar a]b'=0$, and
$a\,x(a')=x([a]a')=x(0)=0$ by Definition~\ref{cv-germs-chart}(3)(a); likewise $\bar a\,x(b')=0$.
Substituting into the identity of Definition~\ref{cv-germs-chart}(3)(e) for $m=2$ gives
$w_a(x(a'),x(b'))\,e_2(0,b')=P_a(a',b')\,e_2(a',0)$, and $e_2(0,b')=1=e_2(a',0)$ by
Definition~\ref{cv-germs-chart}(3)(c). So $e_a(x(a'),x(b'))=P_a(a',b')$.

Rigidity. As $G$ is a semi-abelian scheme proper over $k$, it is an abelian variety, so
$\Gamma(G,\mathcal O_G)=k$, and for every scheme $T$ over $k$ the functions on $T\times_kG$ are pullbacks of
functions on $T$. So two generators of a line bundle on $T\times_kG$ that agree along $T\times\{0\}$ are
equal: their ratio is a unit pulled back from $T$, and it is $1$ along $T\times\{0\}$.

(1) Let $y'$ be the universal point of $G$, its identity morphism. The generators $w_a(y',0)$ and
$w_a(0,0)$ of the line bundle
$\Lambda(y',0)\otimes\Lambda(ay',0)^{-1}\cong\mathcal O_G$ agree along the unit section, so they are equal
by rigidity with $T=\operatorname{Spec}k$. Pulling back along $x$ gives $e_a(x,0)=e_a(0,0)$, and
likewise $e_a(0,y)=e_a(0,0)$. By (4) with $a'=b'=0$, $e_a(0,0)=P_a(0,0)=1$, as $P_a$ is $1$ when one of its
arguments is $0$.

(2) On $T\times_kG$, with second projection $y'$, the generators $w_a(y'+x,y)\,e_a(0,y)$ and
$w_a(y',y)\,e_a(x,y)$ of $\Lambda(ay',y)^{-1}$ agree along $y'=0$, so they are equal by rigidity. Pulling
back along the point $(1,x')$ of $T\times_kG$ and using $e_a(0,y)=1$ from (1) gives (2).

(3) The abelian variety $G$ is proper, smooth and geometrically connected over $k$, and the kernel of $a$
is $\operatorname{Spec}A$ for a finite-dimensional $k$-algebra $A$. So $a:G\to G$ is finite, flat and
surjective, by Lemma~\ref{cv-germs-weil-isogeny}. Let $T=\operatorname{Spec}R$, let $y'$ be the second projection of
$T\times_kG$, and let $\pi=1\times a$, a flat, surjective and quasi-compact endomorphism of $T\times_kG$.
Let $\mathcal N=\Lambda(y',y)$ on $T\times_kG$. Then $\pi^*\mathcal N=\Lambda(ay',y)$, and
$\sigma=w_a(y',y)^{-1}$, read through $\Lambda(y',\bar ay)=\Lambda(y',0)\cong\mathcal O$, is a generator of
it. Two points $(t,y_1)$ and $(t,y_2)$ of $T\times_kG$ with values in a scheme $T'$ and with the same image
under $\pi$ have $y_1=y_2+x$ with $ax=0$, so $w_a(y_1,y)=e_a(x,y)\,w_a(y_2,y)=w_a(y_2,y)$ by (2) and the
hypothesis. So the two pullbacks of $\sigma$ to $(T\times_kG)\times_{\pi,T\times_kG,\pi}(T\times_kG)$
agree, and by faithfully flat descent \cite[Expos\'e~VIII, Th\'eor\`eme~1.1]{Gro71} $\sigma$ is the
pullback of a generator $u$ of $\mathcal N$. As $R$ is local, the line bundle
$y^*\mathcal L\otimes0^*\mathcal L^{-1}$ on $T$ has a generator $c$. Then $u$ and the pullback of $c$ give
a generator of $\mathcal N\otimes y^*\mathcal L\otimes0^*\mathcal L^{-1}=(y+y')^*\mathcal L\otimes
y'^*\mathcal L^{-1}$, the line bundle $\mathcal L_y$ of Definition~\ref{cv-germs-principal}(1), on all of
$T\times_kG$. As $K(\mathcal L)$ is trivial, $y=0$.
\end{proof}
% Lean: the value is Chi w x v c (chi_exists from pair_exists; Chi.unique, Chi.nat). (4) is chi_chart with
% cp, cp_act, cp_zero and cp_eq (val_left, val_right, second_jet_point). (1) is pair_ax and chi_one, with
% Chi w 1 1 1 proved as h11 in KerGeom.ker_data (chi_chart at 0 and the constant coefficient cc); the
% rigidity is CubeDescent.val_eq_of_section with globalFns_holds. (2) is pair_mul and chi_mul (ratio_const,
% surj_fst). (3) is pair_nondeg (exists_gen, CubeDescent.exists_gen_descend, DVR.exists_gen_of_local and
% Out.KTriv), with flat_surj for the isogeny (Lemma cv-germs-weil-isogeny). All in
% KernelRankMultiplicationCharP.lean.

\begin{lemma}\label{cv-germs-kernel-hopf}
Let $E$ be a CM field, let $k$ be an algebraically closed field of characteristic $p>0$, and let $G$ be an
$\mathcal O_E$-semi-abelian scheme over $\operatorname{Spec}k$
(Definition~\ref{cv-model-toroidal-semiabelian}) that is proper over $k$, whose line bundle $\mathcal L$
has $K(\mathcal L)$ trivial (Definition~\ref{cv-germs-principal}(2)). Let $g$ be a formal chart of
dimension $n$ at the identity of $\operatorname{Spec}k$ (Definition~\ref{cv-germs-chart}) with formal germ
$B$, and for $a\in\mathcal O_E$ let $\beta_a$ be the matrix of Definition~\ref{cv-hodge-lattice-germ}(2)
for $B$. Let $a\in\mathcal O_E$, and suppose that the kernel of the action of $a$ on $G$, the fibre
product of $a:G\to G$ with the unit section, is $\operatorname{Spec}A$ for a finite-dimensional
$k$-algebra $A$. Then there are a finite-dimensional commutative and cocommutative Hopf algebra $C$ over
$k$ with $\dim_kC=\dim_kA$, an injective linear map $\delta$ from $\ker\rho_B(a)$ to the derivations at
the counit of $C$, and an injective linear map $\tau$ from $\ker\rho_B(\bar a)$ to the primitive elements
of $C$ (Definition~\ref{cv-germs-hopf}), such that $\delta(u)(\tau(z))=0$ whenever
$z\cdot\beta_a^{\mathsf t}u=0$.
\end{lemma}
\leanref{Development/AutomorphicSpectrum/Proofs/GeometryProofs/ProperGroupGermFiniteSubschemeStatement.lean}{GlobalTheory.T.proper_group_germ_finite_subscheme}
\leanref{Development/AutomorphicSpectrum/Proofs/GeometryProofs/ProperGroupGermFiniteSubscheme.lean}{GlobalTheory.CompanionTheory.Ar.Flat.KerGeom.ker_data}
\leanref{Development/AutomorphicSpectrum/Proofs/GeometryProofs/ProperGroupGermFiniteSubscheme.lean}{GlobalTheory.CompanionTheory.Ar.Flat.KerGeom.exists_hopf}
\begin{proof}
The Hopf algebra. The scheme $\operatorname{Spec}A$ represents the functor of points $x$ of $G$ with
$ax=0$, with a universal point $\kappa$, and this is a functor of commutative groups. So $A$ is a
commutative and cocommutative Hopf algebra over $k$ \cite[\S1.4]{Wat79}, whose comultiplication $\Delta$,
counit $e$ and antipode correspond to the sum of points, the unit and the negative: $\Delta$ is the
homomorphism $A\to A\otimes A$ through which $\kappa_1+\kappa_2$ factors, where $\kappa_1$ and $\kappa_2$
are the images of $\kappa$ under the two inclusions $A\to A\otimes A$. Let $C=A$.

Notation. Let $u_1,\dots,u_s$ be a basis of $\ker\rho_B(a)$ and $v_1,\dots,v_r$ one of
$\ker\rho_B(\bar a)$. Let $R_s=k[\varepsilon_1,\dots,\varepsilon_s]$ with $\varepsilon_i\varepsilon_j=0$
for all $i,j$, so that $R_s=k\oplus k\varepsilon_1\oplus\dots\oplus k\varepsilon_s$, and define $R_r$
likewise. Let $a'=\sum_ju_j\varepsilon_j\in R_s^n$ and $b'=\sum_iv_i\varepsilon_i\in R_r^n$. Products of two
entries of $a'$ vanish, and $\rho_B(a)a'=\sum_j\varepsilon_j\rho_B(a)u_j=0$; likewise for $b'$. By
Lemma~\ref{cv-germs-weil}(4), $a\,x(a')=0$ and $\bar a\,x(b')=0$. For $w\in k^n$, if the point
$x(\varepsilon w)$ with values in $k[\varepsilon]$ is the unit, then $w=0$: by the uniqueness in
Theorem~\ref{cv-germs-formal}(2), applied to the $k$-algebra $k[\varepsilon]$ and its ideal
$(\varepsilon)$, the homomorphisms $J^1_1\to k[\varepsilon]$ with $t\mapsto\varepsilon w$ and $t\mapsto0$
coincide.

Derivations. The point $x(a')$ is the image of $\kappa$ under a homomorphism $\psi:A\to R_s$. Modulo the
$\varepsilon_j$ it is $x(0)=0$, so $\psi$ reduces to $e$. Let $D_j(h)$ be the coefficient of
$\varepsilon_j$ in $\psi(h)$. As $\psi$ is multiplicative, $D_j(hh')=D_j(h)e(h')+e(h)D_j(h')$, so $D_j$ is
a derivation at the counit. The $D_j$ are linearly independent: if $\sum_jl_jD_j=0$, the composite of
$\psi$ with $R_s\to k[\varepsilon]$, $\varepsilon_j\mapsto l_j\varepsilon$, is $e$, so the point
$x(\varepsilon\sum_jl_ju_j)$ is the unit; hence $\sum_jl_ju_j=0$ and $l=0$.

Primitive elements. Over $A\otimes R_r$ let $c=e_a(\kappa,x(b'))$ (Lemma~\ref{cv-germs-weil}), where
$\kappa$ and $x(b')$ stand for their images. Modulo the $\varepsilon_i$, $x(b')$ becomes $x(0)=0$ and $c$
becomes $e_a(\kappa,0)=1$, by Lemma~\ref{cv-germs-weil}(1). So $c=1+\sum_if_i\otimes\varepsilon_i$ with
$f_i\in A$. The image of $c$ under $\Delta\otimes1$ is $e_a(\kappa_1+\kappa_2,x(b'))$, which is
$e_a(\kappa_1,x(b'))\,e_a(\kappa_2,x(b'))$ by Lemma~\ref{cv-germs-weil}(2), that is,
$(1+\sum_if_i\otimes1\otimes\varepsilon_i)(1+\sum_i1\otimes f_i\otimes\varepsilon_i)$. The coefficients of
$\varepsilon_i$ give $\Delta(f_i)=f_i\otimes1+1\otimes f_i$.

Independence. Suppose that $\sum_il_if_i=0$, and let $\lambda:R_r\to k[\varepsilon]$ map $\varepsilon_i$
to $l_i\varepsilon$. The point $y=x(\varepsilon\sum_il_iv_i)$ with values in the local ring $k[\varepsilon]$
is the image of $x(b')$, so $\bar ay=0$. For a scheme $T$ over $\operatorname{Spec}k[\varepsilon]$ and a
point $x$ of $G$ with values in $T$ with $ax=0$, $x$ is the image of $\kappa$ under a morphism
$T\to\operatorname{Spec}A$, so $e_a(x,y_T)$ is the image of $(1\otimes\lambda)(c)=1+(\sum_il_if_i)\otimes
\varepsilon=1$. As the kernel of $a$ is finite, Lemma~\ref{cv-germs-weil}(3) gives $y=0$. So
$\sum_il_iv_i=0$ and $l=0$.

The values. Under $\psi\otimes1:A\otimes R_r\to R_s\otimes R_r$, the image of $c$ is
$e_a(x(a'\otimes1),x(1\otimes b'))$, which is $P_a(a'\otimes1,1\otimes b')$ by
Lemma~\ref{cv-germs-weil}(4), as products of two entries of $a'\otimes1$, or of $1\otimes b'$, vanish. The
coefficient of $\varepsilon_j\otimes\varepsilon_i$ in $1+\sum_i\psi(f_i)\otimes\varepsilon_i$ is
$D_j(f_i)$. In $P_a(a'\otimes1,1\otimes b')$ only the terms of $P_a$ of bidegree $(1,1)$ contribute to this
coefficient, and they give $\sum_{m,m'}(\beta_a)_{mm'}(u_j)_m(v_i)_{m'}$. So $D_j(f_i)=u_j\cdot\beta_av_i$.

The maps. Let $\delta(\sum_jc_ju_j)=\sum_jc_jD_j$ and $\tau(\sum_id_iv_i)=\sum_id_if_i$. They are injective
and take values in the derivations at the counit and in the primitive elements, and
$\delta(u)(\tau(z))=u\cdot\beta_az=z\cdot\beta_a^{\mathsf t}u$ by bilinearity. So $\delta(u)(\tau(z))=0$
whenever $z\cdot\beta_a^{\mathsf t}u=0$.
\end{proof}
% Lean: kerHopf_geom = proper_group_germ_finite_subscheme. ker_data: exists_ker (kappa), g.frame.ad a (w_a),
% AdSum hW, h11, kill, dirv with cp_inj (jet_inj, germ_chart_universal), z = cp g (sq v), psi from pt_surj,
% e, Delta, c from chi_exists, c1, fpsi; f_i = co i c, D_j = proj_j o snd o psi; primitivity by chi_mul and
% co_mul; independence of the f_i by flat_surj, pair_nondeg, exists_tensor and chi_one; that of the D_j by
% dirv; the values by chi_chart and second_jet_expansion. exists_hopf builds the Hopf algebra by pt_inj
% (Bialgebra.ofAlgHom, HopfAlgebra.ofAlgHom with the inverse point) and assemble builds delta and tau.

\subsection{Torus extensions and cocycles}\label{subsection-cv-torus-extensions-cocycles}

\begin{definition}\label{cv-germs-abelian-extension}
Let $k$ be a field and let $G$ be a commutative group scheme over $k$. A \emph{torus extension} of $G$
consists of a free abelian group $\Xi$ of finite rank, an abelian variety $B$ over $k$, a flat, surjective
and quasi-compact homomorphism $q:G\to B$, and a homomorphism $i:T_\Xi\to G$ from the split torus over $k$
with character group $\Xi$ that identifies $T_\Xi$ with the kernel of $q$. For a commutative group scheme
$H$ over $k$ with a line bundle $N$, and points $h_1,h_2$ of $H$ with values in a scheme $T$ over $k$, let
$\Lambda_N(h_1,h_2)=(h_1+h_2)^*N\otimes h_1^*N^{-1}\otimes h_2^*N^{-1}\otimes0^*N$, a line bundle on $T$;
$\Lambda_N$ alone stands for $\Lambda_N(h_1,h_2)$ on $H\times_kH$, with $h_1,h_2$ the two projections.
\end{definition}
\leanref{ReciprocityTheory/ReductiveGroups/ArithmeticGeometry/SemiabelianSchemes/AbelianSchemeQuotients/IsAbelianScheme.lean}{ReciprocityTheory.ChartRestriction.G.Semi.TorusExtension}

\begin{lemma}\label{cv-germs-abelian-units}
Let $S$ be a scheme, let $\Xi$ be a free abelian group, let $T_\Xi$ be the split torus over $S$ with
character group $\Xi$, and let $Z$ be a reduced scheme over $S$ whose underlying space is connected. Let
$\rho$ be an invertible function on $T_\Xi\times_SZ$. Then there is $\chi\in\Xi$ such that
$\rho(\tau,z)=\chi(\tau)\,\rho(1,z)$ for all points $\tau$ of $T_\Xi$ and $z$ of $Z$ with values in a scheme
over $S$.
\end{lemma}
\leanref{Development/AutomorphicSpectrum/Proofs/GeometryProofs/EqSingleOfIsUnitCoeff.lean}{GlobalTheory.CompanionTheory.TorusMono.torus_unit_eq_char_mul}
\begin{proof}
Domains. Let $R$ be a domain. Every unit $u$ of the group algebra $R[\Xi]$ is $c\chi$ for a unit $c$ of $R$
and some $\chi\in\Xi$. Indeed, choose a total order on the free abelian group $\Xi$ compatible with its
group law, and let $uv=1$. As $R$ is a domain, the product of the largest terms of $u$ and $v$ is the
largest term of $uv=1$, and likewise for the smallest terms; so the largest and the smallest terms of $u$
have the same character, and $u$ has one term.

Near a point. Let $\operatorname{Spec}R$ be an affine open subset of $Z$. Over it $T_\Xi\times_SZ$ is
$\operatorname{Spec}R[\Xi]$, and $\rho$ is a unit $\sum_\chi c_\chi\chi$ of $R[\Xi]$, with $c_\chi\in R$. Let
$\mathfrak p$ be a prime ideal of $R$. As the image of $\rho$ in $(R/\mathfrak p)[\Xi]$ is a unit, some
$c_\chi$ is not in $\mathfrak p$. The ring $R'=R[1/c_\chi]$ is reduced. For every prime ideal $\mathfrak p'$
of $R'$, the image of $\rho$ in $(R'/\mathfrak p')[\Xi]$ is a unit whose coefficient at $\chi$ is nonzero,
so by the case of domains it is a multiple of $\chi$. So every coefficient $c_{\chi'}$ with $\chi'\ne\chi$
lies in every prime ideal of $R'$; it is nilpotent, hence $0$ in $R'$. Thus $\rho=c_\chi\chi$ over
$\operatorname{Spec}R'$, an open neighbourhood of $\mathfrak p$.

Conclusion. The character $\chi$ found near a point $z$ of $Z$ is the only one with a nonzero coefficient
in the image of $\rho$ over the residue field at $z$, so it depends only on $z$, and it is locally constant
in $z$. As $Z$ is connected, it is constant. So $\rho(\tau,z)=\chi(\tau)c(z)$ locally on $Z$, and
$c(z)=\rho(1,z)$.
\end{proof}
% Lean: TorusMono.torus_unit_eq_char_mul = eq_char_mul: exists_isMonoNear (exists_coeff_notMem,
% eq_single_of_isUnit_coeff with Rosenlicht.isUnit_eq_single, which uses TwoUniqueSums for the free group),
% IsMonoNear.unique and IsLocallyConstant. EqSingleOfIsUnitCoeff.lean, 177 lines.

\begin{lemma}\label{cv-germs-abelian-seesaw}
Let $k$ be an algebraically closed field, let $G$ be a semi-abelian scheme over $\operatorname{Spec}k$ with a
torus extension $(\Xi,B,q,i)$ (Definition~\ref{cv-germs-abelian-extension}), and let $Y$ be a scheme that is
smooth and geometrically connected over $k$, with a point $y_0\in Y(k)$. Let $N$ be a line bundle on
$B\times_kY$ such that $(q\times1)^*N$ is trivial on $G\times_kY$ and the restriction of $N$ to
$B\times\{y_0\}$ is trivial. Then $N$ is trivial.
\end{lemma}
\leanref{Development/AutomorphicSpectrum/Proofs/GeometryProofs/PicardClassTrivialBySeesawStatement.lean}{GlobalTheory.T.picard_class_trivial_by_seesaw}
\leanref{Development/AutomorphicSpectrum/Proofs/GeometryProofs/PicardClassTrivialBySeesaw.lean}{GlobalTheory.CompanionTheory.TorusPic.pic_inj}
\begin{proof}
Let $\pi=q\times1:G\times_kY\to B\times_kY$, a flat, surjective and quasi-compact morphism, as $q$ is one.
Let $u$ be a generator of $\pi^*N$ and $s$ a generator of the restriction of $N$ to $B\times\{y_0\}$.

The weight. The torus acts on $G\times_kY$ by $\tau\cdot(g,y)=(i(\tau)+g,y)$, and $\pi(\tau\cdot(g,y))=\pi(g,y)$
as $q\circ i=0$. So $u(i(\tau)+g,y)$ and $u(g,y)$ are generators of one line bundle on
$T_\Xi\times_kG\times_kY$, and their ratio is an invertible function $\rho(\tau,g,y)$. The scheme
$G\times_kY$ is smooth and geometrically connected over $k$, so it is reduced and connected. By
Lemma~\ref{cv-germs-abelian-units} there is $\chi\in\Xi$ with $\rho(\tau,g,y)=\chi(\tau)\rho(0,g,y)$, and
$\rho(0,g,y)=1$. So $u(i(\tau)+g,y)=\chi(\tau)\,u(g,y)$.

A unit of the same weight. The generators $u(g,y_0)$ and $s(q(g))$ of the pullback of $N$ along
$g\mapsto(q(g),y_0)$ have as ratio an invertible function $f$ on $G$: $u(g,y_0)=f(g)\,s(q(g))$. At the point
$i(\tau)+g$, where $s(q(i(\tau)+g))=s(q(g))$, this gives $f(i(\tau)+g)=\chi(\tau)f(g)$.

Descent. Let $u'(g,y)=f(g)^{-1}u(g,y)$, a generator of $\pi^*N$ with $u'(i(\tau)+g,y)=u'(g,y)$. Two points
of $G\times_kY$ with values in a scheme $T$ and with the same image under $\pi$ differ by the action of a
point of $T_\Xi$, as $i$ identifies $T_\Xi$ with the kernel of $q$. So the two pullbacks of $u'$ to
$(G\times_kY)\times_{B\times_kY}(G\times_kY)$ agree, and by faithfully flat descent
\cite[Expos\'e~VIII, Th\'eor\`eme~1.1]{Gro71} $u'$ is the pullback of a generator of $N$.
\end{proof}
% Lean: TorusPic.pic_inj: exists_weight (exists_unit_ratio for u along act and snd, then
% TorusMono.torus_unit_eq_char_mul on G x Y), exists_unit_weight (the ratio f along sec G y0), and
% exists_gen_of_weight (scaleGen by f^-1, then exists_gen_descend_val along proj q Y, flat, surjective and
% quasi-compact by CosetQuotient.*_tensorHom). PicardClassTrivialBySeesaw.lean, 232 lines.

\begin{lemma}\label{cv-germs-abelian-descent}
Let $k$ be an algebraically closed field, and let $G$ be a commutative group scheme over $k$ with a torus
extension $(\Xi,B,q,i)$ (Definition~\ref{cv-germs-abelian-extension}).
\begin{enumerate}
\item Let $Y$ be a group scheme over $k$ and $f:G\to Y$ a homomorphism with $f\circ i$ trivial. Then
  $f=f'\circ q$ for a unique morphism $f':B\to Y$.
\item For every homomorphism $f$ from $G$ to an abelian variety over $k$, $f\circ i$ is trivial.
\item Let a commutative ring $O$ act on $G$ by endomorphisms of group schemes. Then there is a unique action
  of $O$ on $B$ by endomorphisms of group schemes such that $q\circ a=a\circ q$ for every $a\in O$.
\end{enumerate}
\end{lemma}
\leanref{Development/AutomorphicSpectrum/Proofs/GeometryProofs/AbelianQuotientOfPowerFns.lean}{GlobalTheory.CompanionTheory.Ar.AbQ.exists_desc}
\leanref{Development/AutomorphicSpectrum/Proofs/GeometryProofs/AbelianQuotientOfPowerFns.lean}{GlobalTheory.CompanionTheory.Ar.AbQ.i_comp_eq_one}
\leanref{Development/AutomorphicSpectrum/Proofs/GeometryProofs/AbelianQuotientOfPowerFns.lean}{GlobalTheory.CompanionTheory.Ar.AbQ.exists_act}
\leanref{Development/AbelianVarietyTheory/Proofs/DegenerationProofs/SplitTorusToAbelianVarietyConst.lean}{GlobalTheory.splitTorus_to_abelian_const}
\begin{proof}
(1) As $q$ is faithfully flat and quasi-compact, it is an effective epimorphism of schemes
\cite[Expos\'e~VIII, Corollaire~5.3]{Gro71}: a morphism $f:G\to Y$ factors through $q$, and then uniquely, if
$f(g)=f(g')$ for all points $g,g'$ of $G$ with values in a scheme $T$ such that $q(g)=q(g')$. For such
points $q(g-g')=0$, so $g-g'=i(\tau)$ for a point $\tau$ of $T_\Xi$, as $i$ identifies $T_\Xi$ with the
kernel of $q$. As $f$ is a homomorphism and $f\circ i$ is trivial, $f(g)=f(i(\tau)+g')=f(g')$.

(2) Every morphism from a rational variety, such as the split torus $T_\Xi$, to an abelian variety is
constant \cite[Corollary~3.9]{Mil86}. The morphism $f\circ i$ maps the unit to the unit, so it is trivial.

(3) Let $a\in O$. By (2), $q\circ a\circ i$ is trivial, so by (1) $q\circ a=\alpha_a\circ q$ for a unique
morphism $\alpha_a:B\to B$. It maps the unit to the unit, as $q$ and $a$ do, so it is a homomorphism by
rigidity \cite[\S4, Corollary~1]{Mum70}. As $q$ is an epimorphism, composing with $q$ gives
$\alpha_1=1$, $\alpha_{ab}=\alpha_a\circ\alpha_b$ and $\alpha_{a+b}=\alpha_a+\alpha_b$ from the same
identities for the action on $G$, and shows that $\alpha_a$ is the only endomorphism with
$q\circ a=\alpha_a\circ q$.
\end{proof}
% Lean: exists_desc (CosetQuotient.effectiveEpi_of_left, EffectiveEpi.desc, T.isKernel.lift); i_comp_eq_one
% (splitTorus_to_abelian_const = TorusConst.eq_const, 478 lines, by extension of a morphism from the
% multiplicative group to the projective line); exists_act (abelianSchemeLaws_all.1.rigidity, cancel_epi).

\begin{lemma}\label{cv-germs-abelian-rosati}
Let $E$ be a CM field and $k$ an algebraically closed field. Let $G$ be an $\mathcal O_E$-semi-abelian scheme
over $\operatorname{Spec}k$ (Definition~\ref{cv-model-toroidal-semiabelian}) with line bundle $\mathcal L$, let
$(\Xi,B,q,i)$ be a torus extension of $G$ (Definition~\ref{cv-germs-abelian-extension}), and let $M$ be a
line bundle on $B$ with $\mathcal L\cong q^*M$. Let $\mathcal O_E$ act on $B$ by endomorphisms, with
$q\circ a=a\circ q$ for $a\in\mathcal O_E$. Then $\Lambda_M(ab_1,b_2)\cong\Lambda_M(b_1,\bar ab_2)$ on
$B\times_kB$ for every $a\in\mathcal O_E$, where $b_1$ and $b_2$ are the two projections.
\end{lemma}
\leanref{Development/AutomorphicSpectrum/Proofs/GeometryProofs/AbelianQuotientOfPowerFns.lean}{GlobalTheory.CompanionTheory.Ar.AbQ.adj_quot}
\begin{proof}
Fix $a$. For points $b,b'$ of $B$ with values in a scheme $T$ over $k$, let $D(b,b')$ be the class of
$\Lambda_M(ab,b')\otimes\Lambda_M(b,\bar ab')^{-1}$ in $\operatorname{Pic}T$. Its formation commutes with
pullback along morphisms $T'\to T$. By the canonical isomorphisms $\Lambda_M(0,b')\cong\mathcal O_T\cong
\Lambda_M(b,0)$, $D(0,b')=0=D(b,0)$.

On the image of $q$. Let $g,g'$ be points of $G$. In Definition~\ref{cv-model-toroidal-semiabelian}(2),
complex conjugation is the Rosati involution of $\mathcal L$:
$\Lambda_{\mathcal L}(ag,g')\cong\Lambda_{\mathcal L}(g,\bar ag')$. As $\mathcal L\cong q^*M$ and $q$
commutes with the actions, $\Lambda_M(aq(g),q(g'))=\Lambda_M(q(ag),q(g'))\cong\Lambda_{\mathcal L}(ag,g')$
and $\Lambda_M(q(g),\bar aq(g'))\cong\Lambda_{\mathcal L}(g,\bar ag')$. So $D(q(g),q(g'))=0$.

Two applications of Lemma~\ref{cv-germs-abelian-seesaw}. The scheme $G$ is semi-abelian, and $G$ and $B$
are smooth and geometrically connected over $k$. First take $Y=G$, $y_0=0$, and the class
$D(b,q(g'))$ on $B\times_kG$, with $b$ and $g'$ the projections: its pullback along $q\times1$ is
$D(q(g),q(g'))=0$, and its restriction to $B\times\{0\}$ is $D(b,0)=0$. So $D(b,q(g'))=0$ for all points.
Next take $Y=B$, $y_0=0$, and the class $D(b',b)$ on $B\times_kB$, with $b$ the first projection and $b'$
the second: its pullback along $q\times1$ is $D(b',q(g))=0$, and its restriction to $B\times\{0\}$ is
$D(0,b)=0$. So $D(b',b)=0$, and pulling back along the exchange of the factors, $D(b_1,b_2)=0$ for the two
projections. This is the claim.
\end{proof}
% Lean: adj_quot: D x y = lam m (x act a) y - lam m x (y act (bar a)); Dpull, Dq (from tK.adj, cls_cube_two,
% lam_pull), D1l, D1r; c1 and c2 by T.picard_class_trivial_by_seesaw with Y = tK.G and Y = T.B; c3 by
% pulling back along the swap; cls_eq_iff converts back to an isomorphism.

\begin{lemma}\label{cv-germs-abelian-normal}
Let $E$ be a CM field and let $w_0\in\mathcal O_E$ with $w_0+\bar w_0=1$. Let $k$ be an algebraically closed
field, let $B$ be an abelian variety over $k$ with an action of $\mathcal O_E$ by endomorphisms, and let $M$
be a line bundle on $B$ with $\Lambda_M(ab_1,b_2)\cong\Lambda_M(b_1,\bar ab_2)$ on $B\times_kB$ for every
$a\in\mathcal O_E$, with the notation of Definition~\ref{cv-germs-abelian-extension}. Let
$N=\Lambda_M(b,w_0b)$, a line bundle on $B$, where $b$ is the identity of $B$.
\begin{enumerate}
\item The line bundle $0^*N$ is trivial; there is an isomorphism $[-1]^*N\cong N$ that is the identity on
  $0^*N$; and $N$ has a cubical structure, a generator of the line bundle $\Theta$ of
  Definition~\ref{cv-germs-chart}(1), for $N$ and the three projections of $B\times_kB\times_kB$, with the
  properties listed there.
\item $\Lambda_N\cong\Lambda_M$ on $B\times_kB$, and $N\cong\Lambda_N(b,w_0b)$.
\item $\Lambda_N(ab_1,b_2)\cong\Lambda_N(b_1,\bar ab_2)$ on $B\times_kB$ for every $a\in\mathcal O_E$.
\end{enumerate}
\end{lemma}
\leanref{Development/AbelianVarietyTheory/Proofs/ModuliProofs/ExistsPrimitiveNormalisation.lean}{GlobalTheory.PrimNorm.exists_primitive_normalisation}
\leanref{Development/AbelianVarietyTheory/Proofs/ModuliProofs/ExistsPrimitiveNormalisation.lean}{GlobalTheory.PrimNorm.exists_normalised_bundle}
\leanref{Development/AbelianVarietyTheory/Proofs/CoherentCohomologyProofs/ExistsGenThetaOfIsAbelianScheme.lean}{GlobalTheory.exists_gen_theta_of_isAbelianScheme}
\begin{proof}
Classes. For a scheme $T$ over $k$ and a point $b'$ of $B$ with values in $T$, let $m(b')\in\operatorname{Pic}T$
be the class of $b'^*M$, and let $\lambda(b_1,b_2)=m(b_1+b_2)-m(b_1)-m(b_2)+m(0)$, the class of
$\Lambda_M(b_1,b_2)$. By the theorem of the cube for the abelian variety $B$, the line bundle $\Theta$ for
$M$ is trivial on $B\times_kB\times_kB$. So $m$ satisfies
$m(b_1+b_2+b_3)-m(b_1+b_2)-m(b_1+b_3)-m(b_2+b_3)+m(b_1)+m(b_2)+m(b_3)-m(0)=0$, and $\lambda$ is additive in
each variable. It is symmetric, and additivity gives $\lambda(-b_1,-b_2)=\lambda(b_1,b_2)$. The hypothesis
gives $\lambda(ab_1,b_2)=\lambda(b_1,\bar ab_2)$.

The class of $N$. The class of $b'^*N$ is $n(b')=\lambda(b',w_0b')$. By additivity,
\[
n(b_1+b_2)-n(b_1)-n(b_2)+n(0)=\lambda(b_1,w_0b_2)+\lambda(b_2,w_0b_1),
\]
and by symmetry and the hypothesis, $\lambda(b_2,w_0b_1)=\lambda(w_0b_1,b_2)=\lambda(b_1,\bar w_0b_2)$. So
the left side is $\lambda(b_1,(w_0+\bar w_0)b_2)=\lambda(b_1,b_2)$: $\Lambda_N(b_1,b_2)$ and
$\Lambda_M(b_1,b_2)$ have the same class.

(2) For the two projections of $B\times_kB$ this gives $\Lambda_N\cong\Lambda_M$. The class of
$\Lambda_N(b,w_0b)$ is $\lambda(b,w_0b)=n(b)$, the class of $N$.

(3) The class of $\Lambda_N(ab_1,b_2)$ is $\lambda(ab_1,b_2)=\lambda(b_1,\bar ab_2)$, that of
$\Lambda_N(b_1,\bar ab_2)$.

(1) As $n(0)=\lambda(0,0)=0$, $0^*N$ is trivial. As $n(-b)=\lambda(-b,-w_0b)=n(b)$, there is an isomorphism
$[-1]^*N\cong N$; along the unit section, where $[-1]$ is the identity, it is multiplication by a unit of
$k$, and dividing by it gives the identity on $0^*N$. As the class of $\Lambda_N(b_1,b_2)$ is
$\lambda(b_1,b_2)$, additive in each variable, the line bundle $\Theta$ for $N$ is trivial on
$B\times_kB\times_kB$. Let $\theta$ be a generator, divided by its value at the origin. The properties of a
cubical structure in Definition~\ref{cv-germs-chart}(1) say that certain invertible functions on products
of two, three and four copies of $B$, built from $\theta$, are $1$. These functions are constant, as every
global function on a product of copies of the abelian variety $B$ is, and they are $1$ at the origin. So
$\theta$ is a cubical structure.
\end{proof}
% Lean: exists_primitive_normalisation applies exists_normalised_bundle with w = act w0, w' = act (bar w0)
% (cmHost_bar_halfIntegerGenerator: bar w0 = 1 - w0). Classes are PullClass (cls, hasGen_iff, lbIso_iff);
% m is ptCls, lambda is diff2, the cube is diff3_ptCls from exists_gen_theta_of_isAbelianScheme; the class
% of N is normCls (normCls_left), and diff2_quadratic, diff3_quadratic and diff2_inv_inv are the three
% computations. The cubical structure is AbLaw.Cube.exists_normalised and cubicalOfNormalised (GlobalFnsSt);
% the symmetry is isSymmetric_of_hasGen. ExistsPrimitiveNormalisation.lean, 397 lines.

\begin{lemma}\label{cv-germs-abelian-ktriv}
Let $B$ be an abelian variety over a field $k$ with line bundles $M$ and $N$ such that $\Lambda_N\cong\Lambda_M$
on $B\times_kB$ (Definition~\ref{cv-germs-abelian-extension}). If $K(M)$ is trivial
(Definition~\ref{cv-germs-principal}(2)), then $K(N)$ is trivial.
\end{lemma}
\leanref{Development/AbelianVarietyTheory/Proofs/CubicalLineBundles/KTrivOfLamEq.lean}{GlobalTheory.CM.kTriv_of_lamEq}
\begin{proof}
Let $b$ be a point of $B$ with values in a commutative $k$-algebra $R$, and suppose that
$\operatorname{Spec}R$ is covered by open sets $V$ such that $N_b$ has a generator over the preimage of $V$
in $\operatorname{Spec}R\times_kB$. Let $y$ be the second projection of $\operatorname{Spec}R\times_kB$ and
$D=N\otimes M^{-1}$. Then
$N_b\otimes M_b^{-1}=(b+y)^*D\otimes y^*D^{-1}=\Lambda_D(b,y)\otimes b^*D\otimes0^*D^{-1}$, where $b^*D$ and
$0^*D$ are pulled back from $\operatorname{Spec}R$. The line bundle
$\Lambda_D=\Lambda_N\otimes\Lambda_M^{-1}$ on $B\times_kB$ is trivial, so its pullback $\Lambda_D(b,y)$ has a
generator, and $b^*D\otimes0^*D^{-1}$ has generators over the sets of an open cover of
$\operatorname{Spec}R$. Refining the cover, $M_b$ has a generator over the preimage of each of its sets. As
$K(M)$ is trivial, $b$ is the unit.
\end{proof}
% Lean: kTriv_of_lamEq, stated for a commutative group over any base: gen_chartPull with
% DescendedPrincipal.chartPullDiv, UU.KTwist.inK_of_lam and inK_div_iff.

\begin{lemma}\label{cv-germs-abelian-tangent}
Let $k$ be a field, and let $G$ be a commutative group scheme over $k$ with a torus extension
$(\Xi,B,q,i)$ (Definition~\ref{cv-germs-abelian-extension}). For a group scheme $H$ over $k$ let
$\operatorname{Lie}H$ be the $k$-vector space of tangent vectors of $H$ at the identity of
$\operatorname{Spec}k$ (Definition~\ref{cv-germs-chart}(2)); a homomorphism $f$ of group schemes maps
tangent vectors to tangent vectors, by a $k$-linear map $df$.
\begin{enumerate}
\item The map that sends $\delta\in\operatorname{Hom}(\Xi,k)$ to the point of $T_\Xi$ with values in
  $k[\varepsilon]$ at which each $\chi\in\Xi$ takes the value $1+\delta(\chi)\varepsilon$ is an isomorphism
  $\operatorname{Hom}(\Xi,k)\to\operatorname{Lie}T_\Xi$.
\item The morphism $q$ is smooth.
\item The sequence $0\to\operatorname{Lie}T_\Xi\xrightarrow{di}\operatorname{Lie}G\xrightarrow{dq}
  \operatorname{Lie}B\to0$ is exact.
\end{enumerate}
\end{lemma}
\leanref{Development/AutomorphicSpectrum/Proofs/GeometryProofs/ExistsCocharacterLieModelStatement.lean}{GlobalTheory.T.exists_cocharacter_lie_model}
\leanref{Development/AutomorphicSpectrum/Proofs/GeometryProofs/SmoothQ.lean}{GlobalTheory.CompanionTheory.Ar.LieSeq.smooth_q}
\leanref{Development/AutomorphicSpectrum/Proofs/GeometryProofs/AbelianQuotientLattice.lean}{GlobalTheory.CompanionTheory.Ar.LieSeq.lie_kernel}
\leanref{Development/AutomorphicSpectrum/Proofs/GeometryProofs/AbelianQuotientLattice.lean}{GlobalTheory.CompanionTheory.Ar.LieSeq.lie_surjective}
\begin{proof}
(1) A point of $T_\Xi$ with values in $k[\varepsilon]$ is a homomorphism $\Xi\to k[\varepsilon]^\times$. It
is the unit modulo $\varepsilon$ if and only if it is $\chi\mapsto1+\delta(\chi)\varepsilon$ for a map
$\delta:\Xi\to k$, which is then additive. The sum of two tangent vectors is the product of these
homomorphisms, and $(1+\delta(\chi)\varepsilon)(1+\delta'(\chi)\varepsilon)=1+(\delta+\delta')(\chi)\varepsilon$;
and $c\in k$ acts through $\varepsilon\mapsto c\varepsilon$, which maps $1+\delta(\chi)\varepsilon$ to
$1+c\delta(\chi)\varepsilon$.

(2) As $i$ identifies $T_\Xi$ with the kernel of $q$, the morphism $T_\Xi\times_kG\to G\times_BG$,
$(\tau,g)\mapsto(i(\tau)+g,g)$, is an isomorphism. So the base change of $q$ along $q$ is the projection
$T_\Xi\times_kG\to G$, which is smooth as $T_\Xi$ is. As $q$ is faithfully flat and quasi-compact, $q$ is
smooth \cite[Corollaire~17.7.3(ii)]{Gro67}.

(3) The homomorphism $i$ is a monomorphism, so $di$ is injective; and $dq\circ di=d(q\circ i)=0$. Let $v$ be
a tangent vector of $G$ with $dq(v)=0$. Then $q(v)$ is the unit, so $v=i(\tau)$ for a point $\tau$ of $T_\Xi$
with values in $k[\varepsilon]$; modulo $\varepsilon$, $i(\tau)$ is the unit, so $\tau$ is, and $v=di(\tau)$.
Let $w$ be a tangent vector of $B$. Modulo $\varepsilon$ it is the unit of $B$, the image of the unit of
$G$. As $q$ is smooth by (2), it has the infinitesimal lifting property
\cite[D\'efinitions~17.1.1 and~17.3.1]{Gro67}: there is a point $v$ of $G$ with values in $k[\varepsilon]$
that is the unit modulo $\varepsilon$ and has $q(v)=w$. Then $v$ is a tangent vector with $dq(v)=w$.
\end{proof}
% Lean: tangent vectors are LieModel (pt, tan, inj, onto, add, smul); (1) is
% TorusLie.exists_lieModel_splitTorus with IsCocharLie; (2) is LieSeq.smooth_q (H1.torsor_pullback,
% H1.torus_dim, MorphismProperty.of_isPullback_of_descendsAlong); (3) is lie_kernel and lie_surjective
% (exists_dual_lift through Algebra.FormallySmooth.liftOfSurjective on the stalk). df is
% lieModel_exists_unique_diff.

\begin{lemma}\label{cv-germs-abelian-characters}
Let $k$ be a field, let $\Xi$ be a free abelian group of finite rank, and let $T_\Xi$ be the split torus over
$k$ with character group $\Xi$. For an endomorphism $u$ of $\Xi$ let $T(u)$ be the endomorphism of $T_\Xi$
with $\chi(T(u)\tau)=u(\chi)(\tau)$ for $\chi\in\Xi$.
\begin{enumerate}
\item Every endomorphism of the group scheme $T_\Xi$ is $T(u)$ for a unique endomorphism $u$ of $\Xi$.
\item Let $G$ and $B$ be group schemes over $k$, $q:G\to B$ a homomorphism and $i:T_\Xi\to G$ a homomorphism
  that identifies $T_\Xi$ with the kernel of $q$, and let a commutative ring $O$ act by endomorphisms on $G$
  and on $B$, with $q\circ a=a\circ q$ for $a\in O$. Then there is a homomorphism of rings
  $h:O\to\operatorname{End}(\Xi)$ with $a\circ i=i\circ T(h(a))$ for every $a\in O$.
\end{enumerate}
\end{lemma}
\leanref{Development/AutomorphicSpectrum/Proofs/GeometryProofs/AbelianQuotientLattice.lean}{GlobalTheory.CompanionTheory.Ar.LieSeq.exists_diagMap}
\leanref{Development/AutomorphicSpectrum/Proofs/GeometryProofs/AbelianQuotientLattice.lean}{GlobalTheory.CompanionTheory.Ar.LieSeq.exists_charAct}
\begin{proof}
(1) Let $t$ be an endomorphism of $T_\Xi$ and $\chi\in\Xi$. The function $\chi\circ t$ is invertible on
$T_\Xi$, so by Lemma~\ref{cv-germs-abelian-units}, with $Z=\operatorname{Spec}k$, it is
$\chi'\cdot(\chi\circ t)(1)$ for some $\chi'\in\Xi$, and $(\chi\circ t)(1)=1$ as $t$ is a homomorphism. The
character $\chi'$ is unique, as distinct characters are distinct functions; let $u(\chi)=\chi'$. As
$\chi\mapsto\chi\circ t$ is multiplicative, $u$ is additive. A point of $T_\Xi$ is determined by the values
of the characters at it, so $t=T(u)$, and $u$ is unique.

(2) Let $a\in O$. As $q\circ a\circ i=a\circ q\circ i$ is trivial, $a\circ i$ factors through the kernel of
$q$: $a\circ i=i\circ t_a$ for a unique morphism $t_a:T_\Xi\to T_\Xi$, which is a homomorphism as $i$ is a
monomorphism and $a\circ i$ is a homomorphism. By (1), $t_a=T(h(a))$ for a unique $h(a)$. As $i$ is a
monomorphism, the uniqueness in (1) gives $h(1)=1$, $h(ab)=h(a)h(b)$ and $h(a+b)=h(a)+h(b)$ from the same
identities on $G$, as $T(u+u')(\tau)=T(u)(\tau)\,T(u')(\tau)$.
\end{proof}
% Lean: exists_diagMap (torus_hom_coord, Diag.charVal_one, chi_injective, coord_comp_diagMap_single) and
% exists_charAct (ker.lift, isMonHom_of_comp, diagMap_injective, diagMap_id, diagMap_comp, diagMap_add).

\begin{lemma}\label{cv-germs-fibre-taylor}
Let $E$ be a CM field, let $G$ be an $\mathcal O_E$-semi-abelian scheme over a scheme $X$
(Definition~\ref{cv-model-toroidal-semiabelian}), let $k$ be an algebraically closed field of characteristic
$p>0$, let $x:\operatorname{Spec}k\to X$ be a morphism, and let $g$ be a formal chart of dimension $n$ at $x$
(Definition~\ref{cv-germs-chart}) with points $x_m$. Let $G_x$ be the fibre of $G$ at $x$, of which the $x_m$
are points with values in $J^1_m$. Then $G_x$ is integral and normal, its function field has
characteristic $p$, and for every open subset $U$ of $G_x$ that contains the unit and every function $f$ on
$U$ there is a power series $F$ over $k$ in $t_1,\dots,t_n$ such that:
\begin{enumerate}
\item $f(x_m)$ is the image of $F$ in $J^1_m$, for every $m$;
\item if $F=0$, then $f=0$;
\item if every partial derivative of $F$ vanishes, then $f$ is a $p$-th power in the function field of
  $G_x$.
\end{enumerate}
\end{lemma}
\leanref{Development/AutomorphicSpectrum/Proofs/GeometryProofs/FormalGermJetLiftStatement.lean}{GlobalTheory.T.formal_germ_jet_lift}
\leanref{Development/WeilBound/CoreStatements/HodgeTatePoints/CgTaylor.lean}{PositivityTheory.HtPt.Cg.fibre_taylor}
\leanref{Development/WeilBound/CoreStatements/HodgeTatePoints/JetOfMono.lean}{PositivityTheory.HtPt.Cg.taylor}
\begin{proof}
The fibre. The fibre $G_x$ is semi-abelian over $k$, so it is smooth and geometrically connected over $k$.
Its local rings are regular, hence normal domains, and it is connected; so it is integral and normal. Its
function field contains $k$, so it has characteristic $p$.

The local ring at the unit. Let $A$ be the local ring of $G_x$ at the unit, a noetherian local $k$-algebra
with maximal ideal $\mathfrak m$, and let $I$ be the ideal of $J^1_m$ generated by $t$. As $\operatorname{Spec}J^1_m$
has one point, which $x_m$ maps to the unit, $x_m$ factors through $\operatorname{Spec}A$, by a homomorphism
$\tau_m:A\to J^1_m$ of $k$-algebras, with $\tau_m(\mathfrak m)\subset I$. As $I^{m+1}=0$, $\tau_m$ factors
through a homomorphism $\tau'_m:A/\mathfrak m^{m+1}\to J^1_m$. The point of $G_x$ with values in
$A/\mathfrak m^{m+1}$ given by $\operatorname{Spec}A/\mathfrak m^{m+1}\to G_x$ is the unit modulo
$\mathfrak m$, so by Theorem~\ref{cv-germs-formal}(2) it is $x_m\circ\operatorname{Spec}\chi_m$ for a unique
homomorphism $\chi_m:J^1_m\to A/\mathfrak m^{m+1}$ with $\chi_m(I)\subset\mathfrak m/\mathfrak m^{m+1}$. Then
$\chi_m\circ\tau_m$ is the quotient map, as $\operatorname{Spec}A\to G_x$ is a monomorphism, and
$\tau'_m\circ\chi_m$ is the identity of $J^1_m$ by the uniqueness in Theorem~\ref{cv-germs-formal}(2),
applied to $x_m$ itself. So $\tau_m$ is onto with kernel $\mathfrak m^{m+1}$. As $x_{m+1}$ reduces to $x_m$, $\tau_{m+1}$ reduces to $\tau_m$, and the $\tau_m$ define a
homomorphism $j:A\to k[[t_1,\dots,t_n]]$ with $j^{-1}(I^{m+1})=\mathfrak m^{m+1}$ for every $m$.

Conclusion. Let $F=j(f)$, for the image of $f$ in $A$. Then (1) holds by construction. If $F=0$, the image
of $f$ in $A$ lies in $\bigcap_m\mathfrak m^m$, which is $0$ by Krull's intersection theorem
\cite[Theorem~8.10]{Mat86}; as $G_x$ is integral, $f=0$.

The $p$-th power. Suppose that every partial derivative of $F$ vanishes. Then the coefficient of $F$ at a
monomial is $0$ unless all its exponents are divisible by $p$; as $k$ is perfect, $F=g^p$ for the power
series $g$ whose coefficient at $t^e$ is the $p$-th root of that of $F$ at $t^{pe}$. Let $m\ge0$. As $\tau_m$
is onto, there is $a\in A$ with $j(a)-g\in I^{m+1}$; then $j(a^p-f)=(j(a)-g)^p\in I^{m+1}$, so
$a^p-f\in\mathfrak m^{m+1}$. Let $d:A\to\Omega_{A/k}$ be the module of K\"ahler differentials, finitely
generated over $A$. As $d(a^p)=pa^{p-1}da=0$, $df=-d(a^p-f)$ lies in $\mathfrak m^m\Omega_{A/k}$. This holds
for every $m$, so $df=0$ by Krull's intersection theorem. Let $K$ be the function field of $G_x$, the
fraction field of $A$, and $D$ a derivation of $K$. It vanishes on $k$, as every element of $k$ is a $p$-th
power and $D(c^p)=pc^{p-1}D(c)=0$; so its restriction to $A$ is a $k$-derivation, which factors through $d$,
and $D(f)=0$. An element of a field of characteristic $p$ on which every derivation vanishes is a $p$-th
power \cite[\S26]{Mat86}. So $f$ is a $p$-th power in $K$.
\end{proof}
% Lean: formal_germ_jet_lift = FibreTaylor.fibre_taylor_expansion = Cg.fibre_taylor (isIntegral_of_smooth_field_
% of_connectedSpace, forall_isDomain_and_isIntegrallyClosed_stalk_of_smooth, Cd.coord_aug for the universal
% property) and Cg.taylor (stalk_jet, exists_tj, Ideal.iInf_pow_eq_bot_of_isLocalRing, germ_injective_of_
% isIntegral); (3) is CubeFib.exists_pow_eq_fractionRing_of_pderiv_eq_zero (exists_pow_eq_of_pderiv_eq_zero,
% derivation_mem_pow_smul_top, Ideal.iInf_pow_smul_eq_bot_of_isLocalRing, derivation_eq_zero_of_kaehler,
% T.exists_pth_root_of_derivation_zero). FormalGermJetLift.lean, CgTaylor.lean, JetOfMono.lean and
% ExistsPowEqFractionRingOfPderivEqZero.lean have about 400 lines.

\begin{lemma}\label{cv-germs-abelian-bilinear}
Let $R$ be a commutative ring. For a power series $f(t,t')$ over $R$ in two families of $n$ variables, let
$\operatorname{bil}(f)$ be the $n\times n$ matrix whose $(i,j)$ entry is the coefficient of $t_it'_j$ in $f$;
for an $n_0$-tuple $\varphi$ of power series in $t_1,\dots,t_n$ without constant term, let
$\operatorname{lin}(\varphi)$ be the $n_0\times n$ matrix whose $(i,j)$ entry is the coefficient of $t_j$ in
$\varphi_i$. Call $f$ \emph{normalized} if $f(0,t')=1=f(t,0)$.
\begin{enumerate}
\item If $f$ and $f'$ are normalized, so is $ff'$, and $\operatorname{bil}(ff')=\operatorname{bil}(f)+
  \operatorname{bil}(f')$.
\item If $f$ is a normalized power series in two families of $n_0$ variables and $\varphi,\psi$ are
  $n_0$-tuples of power series in $t_1,\dots,t_n$ without constant term, then $f(\varphi(t),\psi(t'))$ is
  normalized and $\operatorname{bil}(f(\varphi(t),\psi(t')))=\operatorname{lin}(\varphi)^{\mathsf t}
  \operatorname{bil}(f)\operatorname{lin}(\psi)$.
\end{enumerate}
\end{lemma}
\leanref{Development/AutomorphicSpectrum/Proofs/AlgebraProofs/FormalBilinearTermProduct.lean}{GlobalTheory.CompanionTheory.Ar.Conj.BilCalc.bil_calc}
\leanref{Development/AutomorphicSpectrum/Proofs/GeometryProofs/SubstitutionBilinearFormTransform.lean}{GlobalTheory.CompanionTheory.Ar.Flat.FrGauge.BilQ.bil_axOne_subQ}
\begin{proof}
A normalized series is $1$ plus terms of positive degree both in $t$ and in $t'$. (1) Then
$ff'=1+(f-1)+(f'-1)+(f-1)(f'-1)$, and the last product has degree at least $2$ in $t$ and in $t'$. (2)
Substituting $\varphi(t)$ and $\psi(t')$ into a monomial of degrees $a$ in $t$ and $b$ in $t'$ gives terms
of degrees at least $a$ and $b$, so only the terms $\operatorname{bil}(f)_{ij}t_it'_j$ contribute terms
$t_lt'_{l'}$, and they give $\sum_{i,j}\operatorname{bil}(f)_{ij}\operatorname{lin}(\varphi)_{il}
\operatorname{lin}(\psi)_{jl'}$. Setting $t=0$ or $t'=0$ makes $\varphi(t)$ or $\psi(t')$ vanish.
\end{proof}
% Lean: Ca.bil_calc (formal_bilinear_term_product) for (1) and for (2) with n_0 = n, and BilQ.bil_axOne_subQ
% (substitution_bilinear_form_transform, 265 lines) for (2) with phi = psi = Q of another dimension. The Lean
% states further identities that are not used here.

\begin{lemma}\label{cv-germs-abelian-quotient}
Let $E$ be a CM field, and let $G$ be an $\mathcal O_E$-semi-abelian scheme over a scheme $X$
(Definition~\ref{cv-model-toroidal-semiabelian}) with line bundle $\mathcal L$, such that, for every
$m\ge1$, every function on the $m$-fold fibre product of $G$ over $X$ is the pullback of a function on $X$.
Let $k$ be an algebraically closed field of characteristic $p>0$, and let $x:\operatorname{Spec}k\to X$ be a
morphism at which $G$ has a principal fibre (Definition~\ref{cv-germs-principal}(3)), with fibre $G_x$.
\begin{enumerate}
\item There are an $\mathcal O_E$-semi-abelian scheme $G_0$ over $\operatorname{Spec}k$, proper over $k$,
  whose line bundle $\mathcal L_0$ has $K(\mathcal L_0)$ trivial, a free abelian group $\Xi$ of finite
  rank, a flat, surjective and quasi-compact homomorphism $q:G_x\to G_0$ compatible with the actions of
  $\mathcal O_E$, and a homomorphism $T_\Xi\to G_x$ that identifies the split torus $T_\Xi$ with the kernel
  of $q$, such that $\mathcal L$ restricts to $q^*\mathcal L_0$ on $G_x$.
\item Let $G_0$ and $q$ be as in (1), let $g$ be a formal chart of $G$ of dimension $n$ at $x$ with formal
  germ $B$, and let $g_0$ be a formal chart of $G_0$ of dimension $n_0$ at the identity of
  $\operatorname{Spec}k$ with formal germ $B_0$ (Definition~\ref{cv-germs-chart}). Then there are an
  $\mathcal O_E$-module $\Lambda$, free of finite rank over $\mathbf Z$, an $n_0\times n$ matrix $P$ over
  $k$ and an injective $k$-linear map $j:k\otimes\Lambda\to k^n$ such that $q$ maps the $l$-th tangent
  vector of $g$ to $\sum_iP_{il}$ times the $i$-th tangent vector of $g_0$, $P:k^n\to k^{n_0}$ is onto with
  kernel $j(k\otimes\Lambda)$, $P\rho_B(a)=\rho_{B_0}(a)P$, and $j((1\otimes a)v)=\rho_B(a)j(v)$ for every
  $a\in\mathcal O_E$.
\item In (2), let $P$ be any $n_0\times n$ matrix over $k$ such that $q$ maps the $l$-th tangent vector of
  $g$ to $\sum_iP_{il}$ times the $i$-th tangent vector of $g_0$. Then there is an $n\times n$ matrix $Z$
  over $k$ with $\beta_a+\rho_B(a)^{\mathsf t}Z-Z\rho_B(\bar a)=P^{\mathsf t}\beta^0_aP$ for every
  $a\in\mathcal O_E$, where $\beta_a$ and $\beta^0_a$ are the matrices of
  Definition~\ref{cv-hodge-lattice-germ}(2) for $B$ and $B_0$.
\end{enumerate}
\end{lemma}
\leanref{Development/AutomorphicSpectrum/Proofs/GeometryProofs/AbelianQuotientOfPowerFnsStatement.lean}{GlobalTheory.T.abelian_quotient_of_power_fns}
\leanref{Development/AutomorphicSpectrum/Proofs/GeometryProofs/AbelianQuotientLatticeStatement.lean}{GlobalTheory.T.abelian_quotient_lattice}
\leanref{Development/AutomorphicSpectrum/Proofs/GeometryProofs/AbelianQuotientBilinearStatement.lean}{GlobalTheory.T.abelian_quotient_bilinear}
\begin{proof}
The fibre. Let $G_x$ be the fibre of $G$ at $x$. By Lemma~\ref{cv-origin-reference-pullback}(1) it is an
$\mathcal O_E$-semi-abelian scheme over $\operatorname{Spec}k$, whose line bundle $\mathcal L_x$ is the
restriction of $\mathcal L$. As $G$ has a principal fibre at $x$, there are a torus extension
$(\Xi,B,q,i)$ of $G_x$ (Definition~\ref{cv-germs-abelian-extension}) and a line bundle $M$ on $B$ with
$K(M)$ trivial and $\mathcal L_x\cong q^*M$. The hypothesis that functions on the fibre powers of $G$ come
from $X$ is not used in the proofs of (1), (2) and (3).

(1) The action. By Lemma~\ref{cv-germs-abelian-descent}(3), $\mathcal O_E$ acts on $B$ by endomorphisms,
with $q\circ a=a\circ q$ for $a\in\mathcal O_E$. By Lemma~\ref{cv-germs-abelian-rosati},
$\Lambda_M(ab_1,b_2)\cong\Lambda_M(b_1,\bar ab_2)$ on $B\times_kB$ for every $a$.

The line bundle. Let $w_0=(1+\sqrt{-d})/2$, an element of $\mathcal O_E$ with $w_0+\bar w_0=1$, where
$\sqrt{-d}$ is as in Situation~\ref{cv-function-situation}. Let $\mathcal L_0=\Lambda_M(b,w_0b)$, for the
identity $b$ of $B$. By Lemma~\ref{cv-germs-abelian-normal}, $\mathcal L_0$ is trivial along the unit
section, symmetric and cubical, $\Lambda_{\mathcal L_0}(ab_1,b_2)\cong\Lambda_{\mathcal L_0}(b_1,\bar ab_2)$
for every $a$, $\Lambda_{\mathcal L_0}\cong\Lambda_M$, and $\mathcal L_0\cong\Lambda_{\mathcal L_0}(b,w_0b)$.
So $B$, with this action of $\mathcal O_E$ and the line bundle $\mathcal L_0$, normalized by the last
isomorphism, is an $\mathcal O_E$-semi-abelian scheme $G_0$ over $\operatorname{Spec}k$
(Definition~\ref{cv-model-toroidal-semiabelian}(2)), as an abelian variety is semi-abelian, and it is
proper over $k$. By Lemma~\ref{cv-germs-abelian-ktriv}, $K(\mathcal L_0)$ is
trivial, as $K(M)$ is.

The restriction of $\mathcal L$. The line bundle of the $\mathcal O_E$-semi-abelian scheme $G_x$ is
normalized (Definition~\ref{cv-model-toroidal-semiabelian}(2)): $\mathcal L_x\cong\Lambda_{\mathcal L_x}(g,w_0g)$,
for the identity $g$ of $G_x$. As $\mathcal L_x\cong q^*M$ and $q\circ w_0=w_0\circ q$,
\[
q^*\mathcal L_0=\Lambda_M(q(g),w_0q(g))=\Lambda_M(q(g),q(w_0g))\cong\Lambda_{\mathcal L_x}(g,w_0g)
\cong\mathcal L_x.
\]
Together with $\Xi$ and $i$, this gives (1).
% The normalization of L_x is the field prim of the SAB of G_x (tK.prim, from sabPull).
% Lean: exists_abQuot, with sabPull for G_x, exists_act,
% adj_quot, PrimNorm.exists_primitive_normalisation (with exists_gen_theta_of_isAbelianScheme), hpull
% (tK.prim, cls_cube_two, lam_pull, hact) and CM.kTriv_of_lamEq. The hypothesis PowFns is not used.

(2) Tangent spaces. Let $\operatorname{Lie}$ and $d$ be as in Lemma~\ref{cv-germs-abelian-tangent}. A tangent
vector of $G$ at $x$ is a point of $G$ with values in $k[\varepsilon]$ over $x$, that is, a tangent vector of
$G_x$ at the identity of $\operatorname{Spec}k$. So by Definition~\ref{cv-germs-chart}(3)(b) the tangent
vectors $v_1,\dots,v_n$ of $g$ form a basis of $\operatorname{Lie}G_x$, and the tangent vectors
$v^0_1,\dots,v^0_{n_0}$ of $g_0$ form one of $\operatorname{Lie}G_0$. Let $x_m$ be the points of $g$. For
$c\in k^n$, $\sum_lc_lv_l=x_1(\varepsilon c)$, as in the proof of Lemma~\ref{cv-germs-compare}(3), and for
$a\in\mathcal O_E$, $a\,x_1(\varepsilon c)=x_1([a](\varepsilon c))=x_1(\varepsilon\rho_B(a)c)$ by
Definition~\ref{cv-germs-chart}(3)(a). So $da$ has the matrix $\rho_B(a)$ in the basis $(v_l)$, and likewise
the matrix $\rho_{B_0}(a)$ in the basis $(v^0_i)$.

The lattice. Let $h:\mathcal O_E\to\operatorname{End}(\Xi)$ be the homomorphism of
Lemma~\ref{cv-germs-abelian-characters}(2), for the actions on $G_x$ and $B$, and let $T(u)$ be as there. Let
$\Lambda=\operatorname{Hom}(\Xi,\mathbf Z)$, on which $a\in\mathcal O_E$ acts by $l\mapsto l\circ h(a)$. It
is an $\mathcal O_E$-module, free of finite rank over $\mathbf Z$ as $\Xi$ is, and the map
$k\otimes\Lambda\to\operatorname{Hom}(\Xi,k)$, $c\otimes l\mapsto cl$, is bijective.

The maps. Let $P$ be the matrix of $dq$ in the bases $(v_l)$ and $(v^0_i)$: $q$ maps $v_l$ to
$\sum_iP_{il}v^0_i$. Let $j$ be the composite of $k\otimes\Lambda\to\operatorname{Hom}(\Xi,k)$, the
isomorphism $\operatorname{Hom}(\Xi,k)\to\operatorname{Lie}T_\Xi$ of
Lemma~\ref{cv-germs-abelian-tangent}(1), $di$, and the coordinates $\operatorname{Lie}G_x\to k^n$ in the
basis $(v_l)$. By Lemma~\ref{cv-germs-abelian-tangent}(3), $P:k^n\to k^{n_0}$ is onto, $j$ is injective,
and the kernel of $P$ is the image of $j$. As $q\circ a=a\circ q$, $dq\circ da=da\circ dq$, that is,
$P\rho_B(a)=\rho_{B_0}(a)P$. Let $\delta\in\operatorname{Hom}(\Xi,k)$. The endomorphism $T(h(a))$ of
$T_\Xi$ maps the tangent vector of $\delta$ to that of $\delta\circ h(a)$, as the value of $\chi$ at the
image is $1+\delta(h(a)\chi)\varepsilon$. As $a\circ i=i\circ T(h(a))$, $di$ maps the tangent vector of
$\delta\circ h(a)$ to $a$ times $di$ of the tangent vector of $\delta$, whose coordinates are $\rho_B(a)$
times those of the latter. As $(1\otimes a)(c\otimes l)=c\otimes(l\circ h(a))$ maps to $c\,l\circ h(a)$,
$j((1\otimes a)v)=\rho_B(a)j(v)$ for $v\in k\otimes\Lambda$.
% Lean: lie_seq: lieModel_exists for G_x and G_0, exists_cocharacter_lie_model for the torus,
% LiePull.lieModel_pullback for the tangent vectors at x, lieModel_exists_unique_act and _diff for da and dq,
% lie_kernel and lie_surjective (with smooth_q), exists_charAct, formal_germ_lie_basis = Cd.germ_lie for the
% bases (germ_coord), tenMap_bijective, lamEnd for the action on Lambda = abelian2 Q.Ch, pt_comp_diagMap and
% tenMap_end for the equivariance of j.

(3) Notation. Let $x_m$, $e_m$ and $w_a$ be the points and generators of $g$ in
Definition~\ref{cv-germs-chart}(3), with $[a]$ and $P_a$ the series of $B$, and let $x^0_m$, $e^0_m$,
$w^0_a$, $[a]_0$ and $P^0_a$ be those of $g_0$ and $B_0$. Fix an isomorphism
$\mathcal L_x\cong q^*\mathcal L_0$, as in (1). Through it,
$\Lambda_{\mathcal L_x}(y_1,y_2)\cong\Lambda_{\mathcal L_0}(q(y_1),q(y_2))$ for points $y_1,y_2$ of $G_x$,
compatibly with the canonical trivializations when $y_1$ or $y_2$ is $0$, and generators of the first are
identified with generators of the second.

The series $Q$. For each $m$, $q(x_m)$ is a point of $G_0$ with values in $J^1_m$ that is the unit modulo
$t$. By Theorem~\ref{cv-germs-formal}(2) for $g_0$, $q(x_m)=x^0_m\circ\operatorname{Spec}\chi_m$ for a unique
homomorphism $\chi_m$ from the ring $J^1_m$ of $g_0$, in $n_0$ variables, to that of $g$ that maps the
variables into the ideal generated by $t$. By uniqueness, $\chi_{m+1}$ reduces to $\chi_m$. So there is an
$n_0$-tuple $Q$ of power series in $t$ without constant term with $q(x_m(t))=x^0_m(Q(t))$ for every $m$.

Its linear part. Let $\operatorname{lin}$ and $\operatorname{bil}$ be as in
Lemma~\ref{cv-germs-abelian-bilinear}. By Definition~\ref{cv-germs-chart}(3)(b), $v_l=x_1(\varepsilon e_l)$
for the $l$-th unit vector $e_l$, and $x^0_1(\varepsilon c)=\sum_ic_iv^0_i$ for $c\in k^{n_0}$, as in (2).
So $q(v_l)=x^0_1(Q(\varepsilon e_l))=x^0_1(\varepsilon\operatorname{lin}(Q)e_l)=\sum_i\operatorname{lin}(Q)_{il}
v^0_i$. As $q(v_l)=\sum_iP_{il}v^0_i$ and coordinates are unique, $\operatorname{lin}(Q)=P$.

The actions. Let $a\in\mathcal O_E$. By Definition~\ref{cv-germs-chart}(3)(a) and $q\circ a=a\circ q$,
$x^0_m(Q([a]t))=q(x_m([a]t))=q(a\,x_m(t))=a\,x^0_m(Q(t))=x^0_m([a]_0Q(t))$. By the uniqueness in
Theorem~\ref{cv-germs-formal}(2), $Q([a]t)=[a]_0Q(t)$.

The gauge. On $\operatorname{Spec}J^2_m$, $e^0_m(Q(t),Q(t'))$ is a generator of
$\Lambda_{\mathcal L_0}(x^0_m(Q(t)),x^0_m(Q(t')))=\Lambda_{\mathcal L_0}(q(x_m(t)),q(x_m(t')))$, which is
identified with $\Lambda_{\mathcal L_x}(x_m(t),x_m(t'))$. So $e^0_m(Q(t),Q(t'))=\lambda_m(t,t')\,e_m(t,t')$ for
a unit $\lambda_m$ of $J^2_m$. As $e_{m+1}$ and $e^0_{m+1}$ reduce to $e_m$ and $e^0_m$, $\lambda_{m+1}$ reduces
to $\lambda_m$, and there is a power series $\lambda(t,t')$ whose image in each $J^2_m$ is $\lambda_m$. As $e_m$
and $e^0_m$ are $1$ when one argument is $0$ (Definition~\ref{cv-germs-chart}(3)(c)), and the identification
respects the canonical trivializations, $\lambda$ is normalized in the sense of
Lemma~\ref{cv-germs-abelian-bilinear}.

The generators $w_a$. Let $y_1,y_2$ be the projections of $G_x\times_kG_x$. As $q$ commutes with $a$ and
$\bar a$, $w^0_a(q(y_1),q(y_2))$ is a generator of the line bundle identified with
$\Lambda_{\mathcal L_x}(y_1,\bar ay_2)\otimes\Lambda_{\mathcal L_x}(ay_1,y_2)^{-1}$, of which $w_a(y_1,y_2)$ is
also a generator. So $w^0_a(q(y_1),q(y_2))=u(y_1,y_2)\,w_a(y_1,y_2)$ for an invertible function $u$ on
$G_x\times_kG_x$. By Definition~\ref{cv-germs-chart}(3)(e) at $t=0$, as $e_m(0,t')=1$ and $P_a(0,t')=1$,
$w_a(0,x_m(t'))=1$, read through the canonical trivializations; likewise $w^0_a(0,x^0_m(s))=1$ for the
variables $s$ of $g_0$, and so at $s=Q(t')$. So $u(0,x_m(t'))=1$ for every $m$, and by
Lemma~\ref{cv-germs-fibre-taylor}(2), applied to the function $u(0,y)-1$ on $G_x$, $u(0,y)=1$. In the same
way $u(y,0)=1$.

The function $u$ is $1$. The scheme $G_x\times_kG_x$ is smooth and connected over $k$, so it is integral. Let
$\rho(\tau,y_1,y_2)=u(i(\tau)+y_1,y_2)\,u(y_1,y_2)^{-1}$, an invertible function on
$T_\Xi\times_kG_x\times_kG_x$ with $\rho(1,y_1,y_2)=1$. By Lemma~\ref{cv-germs-abelian-units},
$\rho(\tau,y_1,y_2)=\chi(\tau)$ for some $\chi\in\Xi$, and $\chi(\tau)=\rho(\tau,0,0)=u(i(\tau),0)\,u(0,0)^{-1}
=1$. So $u(i(\tau)+y_1,y_2)=u(y_1,y_2)$, and in the same way $u(y_1,i(\tau)+y_2)=u(y_1,y_2)$. Two points of
$G_x\times_kG_x$ with the same image under $q\times q$ differ by a point of $T_\Xi\times_kT_\Xi$, and
$q\times q$ is flat, surjective and quasi-compact. By faithfully flat descent
\cite[Expos\'e~VIII, Th\'eor\`eme~1.1]{Gro71}, $u$ is the pullback of a function on $B\times_kB$. That function
is constant, as $B\times_kB$ is an abelian variety, and its value is $u(0,0)=1$. So $u=1$.

The relation. Pull back the identity of Definition~\ref{cv-germs-chart}(3)(e) for $g_0$ along $s=Q(t)$,
$s'=Q(t')$, and write $w_a$ for $w_a(x_m(t),x_m(t'))$. As $u=1$, the first factor becomes $w_a$; as
$Q([a]t)=[a]_0Q(t)$ and $Q([\bar a]t')=[\bar a]_0Q(t')$, the gauge gives
$e^0_m([a]_0Q(t),Q(t'))=\lambda([a]t,t')e_m([a]t,t')$ and $e^0_m(Q(t),[\bar a]_0Q(t'))=\lambda(t,[\bar a]t')
e_m(t,[\bar a]t')$. So
\[
w_a\,\lambda([a]t,t')\,e_m([a]t,t')=P^0_a(Q(t),Q(t'))\,\lambda(t,[\bar a]t')\,e_m(t,[\bar a]t').
\]
Dividing by the identity $w_a\,e_m([a]t,t')=P_a(t,t')\,e_m(t,[\bar a]t')$ for $g$ gives, in every $J^2_m$
and hence as power series,
\[
\lambda([a]t,t')\,P_a(t,t')=P^0_a(Q(t),Q(t'))\,\lambda(t,[\bar a]t').
\]

Bilinear terms. Let $Z=\operatorname{bil}(\lambda)$. The series $\lambda$, $P_a$ and $P^0_a$ are normalized,
as $P_a$ and $P^0_a$ are $1$ when one of their arguments is $0$; and $\operatorname{lin}([a])=\rho_B(a)$,
$\operatorname{lin}(Q)=P$, $\operatorname{bil}(P_a)=\beta_a$ and $\operatorname{bil}(P^0_a)=\beta^0_a$
(Definition~\ref{cv-hodge-lattice-germ}(2)). By Lemma~\ref{cv-germs-abelian-bilinear}, the matrices
$\operatorname{bil}$ of the two sides are $\rho_B(a)^{\mathsf t}Z+\beta_a$ and $P^{\mathsf t}\beta^0_aP+
Z\rho_B(\bar a)$. So $\beta_a+\rho_B(a)^{\mathsf t}Z-Z\rho_B(\bar a)=P^{\mathsf t}\beta^0_aP$ for every
$a\in\mathcal O_E$.
% Lean: frame_gauge: Q.bun (w), ChartMap.main (Q, by formal_germ_dimension_unique and exists_ps_of_compat),
% linQ_eq_diff (tanPt, comb_tanPt), substitution_action_commutes, Lamex.exists_lam (exists_U, U_trunc,
% unit_first_axis_eq_one, unit_second_axis_eq_one), AdTr.adjunction_ratio (u), UAxis.unit_one_on_axes
% (WAxis.adAxes, eq_one_of_jets through formal_germ_jet_lift), FibUnit.unit_eq_one_of_axes (pb_act_eq with
% BiextTriv.exists_char, exists_descent by FlatDescent.existsUnique_appTop_eq_iff, globalFns_holds),
% AdRelPf.adRel_series, and the bilinear terms by bil_calc and bil_axOne_subQ.
% Lean: (1) is abelian_quotient_of_power_fns = Ar.AbQ.exists_abQuot, the structure Ar.Flat.AbQuot
% (AbelianQuotientOfPowerFns.lean, 227 lines: exists_desc, i_comp_eq_one, exists_act, adj_quot; with sabPull,
% 186 lines, PrimNorm.exists_primitive_normalisation, 397 lines, CM.kTriv_of_lamEq and
% picard_class_trivial_by_seesaw, 232 lines): about 1,100 lines. (2) is abelian_quotient_lattice =
% Ar.LieSeq.lie_seq (AbelianQuotientLattice.lean, 522 lines); the tangent condition is AbQuot.IsDiff. (3) is
% abelian_quotient_bilinear = Ar.Flat.FrGauge.frame_gauge (60 lines over abelian_quotient_germ_power_series,
% _differential_matrix, _gauge_transform, _line_bundle_gauge, _gauge_unit, _gauge_relation,
% unit_trivial_on_both_axes, substitution_bilinear_form_transform and formal_bilinear_term_product,
% whose proofs have about 2,150 lines).
\end{proof}

\begin{lemma}\label{cv-germs-cocycle-powers}
Let $E$ be a CM field, let $R$ be a commutative ring, and let $B=(F,[\,\cdot\,],\Phi_1,\Phi_2,S,P)$ be a
formal germ of dimension $n$ over $R$ (Definition~\ref{cv-hodge-lattice-germ}). For $m\ge0$ let
\[
 \Pi_m(x,y)=\prod_{i=0}^{m-1}\Phi_1([i]x,x,y),
\]
where $[i]$ is the series of the integer $i\in\mathcal O_E$. Then:
\begin{enumerate}
\item $\Pi_m(0,y)=1=\Pi_m(x,0)$;
\item $\Pi_m(x+_Fy,z)\,\Phi_1(x,y,z)^m=\Pi_m(x,z)\,\Pi_m(y,z)\,\Phi_1([m]x,[m]y,z)$;
\item $\Pi_m(x,y)\,\Pi_m(x,z)\,\Phi_2([m]x,y,z)=\Phi_2(x,y,z)^m\,\Pi_m(x,y+_Fz)$.
\end{enumerate}
\end{lemma}
\leanref{Development/WeilBound/CoreStatements/HodgeTatePoints/FirstPowerEqProdAt.lean}{PositivityTheory.HtPt.Cb.first_power_zero_left}
\leanref{Development/WeilBound/CoreStatements/HodgeTatePoints/FirstPowerEqProdAt.lean}{PositivityTheory.HtPt.Cb.first_power_zero_right}
\leanref{Development/WeilBound/CoreStatements/HodgeTatePoints/FirstPowerEqProdAt.lean}{PositivityTheory.HtPt.Cb.first_power_add}
\leanref{Development/AutomorphicSpectrum/Proofs/AlgebraProofs/AlphaKernelDerivativeVanishes.lean}{GlobalTheory.CompanionTheory.Ar.Kill.inter2}
\begin{proof}
The identities $[a+b]x=[a]x+_F[b]x$ and $[1]x=x$ of a formal germ give $[0]x=0$ and $[i+1]x=[i]x+_Fx$, and
$[i](x+_Fy)=[i]x+_F[i]y$ is another of its identities. The series $\Phi_1$ and $\Phi_2$ are $1$ when one of
their arguments is $0$.

(1) Each factor of $\Pi_m(0,y)$ is $\Phi_1(0,0,y)=1$, and each factor of $\Pi_m(x,0)$ is
$\Phi_1([i]x,x,0)=1$.

(2) A four-term identity. Let $a,b,c,e$ be $n$-tuples of power series without constant term, and write
$\Phi_1(a,b)$ for $\Phi_1(a,b,z)$. Then
\[
 \Phi_1(a,b)\,\Phi_1(c,e)\,\Phi_1(a+_Fb,c+_Fe)=\Phi_1(a,c)\,\Phi_1(b,e)\,\Phi_1(a+_Fc,b+_Fe).
\]
Indeed, the identity $\Phi_1(x,y,w)\Phi_1(x+_Fy,z,w)=\Phi_1(y,z,w)\Phi_1(x,y+_Fz,w)$ of a formal germ, with
$(x,y,z,w)$ replaced by $(a+_Fb,c,e,z)$ and then by $(a,b,c,z)$, turns the left side into
$\Phi_1(a,b)\Phi_1(a+_Fb,c)\Phi_1(a+_Fb+_Fc,e)$ and then into $\Phi_1(b,c)\Phi_1(a,b+_Fc)\Phi_1(a+_Fb+_Fc,e)$.
With $(x,y,z,w)$ replaced by $(a+_Fc,b,e,z)$ and then by $(a,c,b,z)$, it turns the right side into
$\Phi_1(c,b)\Phi_1(a,c+_Fb)\Phi_1(a+_Fb+_Fc,e)$. The two agree, as $\Phi_1(c,b)=\Phi_1(b,c)$.

Induction on $m$. For $m=0$ both sides of (2) are $1$, as $\Phi_1(0,0,z)=1$. Suppose (2) holds for $m$.
Then $\Pi_{m+1}(x+_Fy,z)\Phi_1(x,y,z)^{m+1}$ is
\[
 \Pi_m(x+_Fy,z)\Phi_1(x,y,z)^m\cdot\Phi_1([m]x+_F[m]y,x+_Fy,z)\,\Phi_1(x,y,z)
 =\Pi_m(x,z)\Pi_m(y,z)\cdot\Phi_1([m]x,[m]y,z)\,\Phi_1(x,y,z)\,\Phi_1([m]x+_F[m]y,x+_Fy,z).
\]
By the four-term identity with $(a,b,c,e)=([m]x,[m]y,x,y)$, the last three factors are
$\Phi_1([m]x,x,z)\Phi_1([m]y,y,z)\Phi_1([m+1]x,[m+1]y,z)$, which gives (2) for $m+1$.

(3) For $m=0$ both sides are $1$, as $\Phi_2(0,y,z)=1$. The identity
$\Phi_1(x,y,z)\Phi_1(x,y,w)\Phi_2(x+_Fy,z,w)=\Phi_2(x,z,w)\Phi_2(y,z,w)\Phi_1(x,y,z+_Fw)$ of a formal germ, with
$(x,y,z,w)$ replaced by $([m]x,x,y,z)$, gives
\[
 \Phi_1([m]x,x,y)\,\Phi_1([m]x,x,z)\,\Phi_2([m+1]x,y,z)=\Phi_2([m]x,y,z)\,\Phi_2(x,y,z)\,\Phi_1([m]x,x,y+_Fz).
\]
Multiplying by $\Pi_m(x,y)\Pi_m(x,z)$ and using (3) for $m$ gives (3) for $m+1$.
\end{proof}
% Lean: (1) first_power_zero_left, first_power_zero_right (first_power_eq_prod_at, first_form_zero_left and
% _right); (2) first_power_add = first_power_add_at, by induction with first_form_four_term (four instances
% of first_form_assoc_at, from cocycleLeft_assoc, and first_form_comm_at); FirstPowerEqProdAt.lean, 179 lines.
% (3) Ar.Kill.inter2 (induction with B.interchange and cocycleRight_zero_fst; formal_module_action_laws for
% [0]x = 0 and [m+1]x = [m]x + x). Pi_m is pow1 B m.

\begin{lemma}\label{cv-germs-cocycle-series}
Let $E$ be a CM field, let $p$ be a prime, let $R$ be a commutative ring in which $p=0$, let $B$ be a
formal germ of dimension $n$ over $R$ (Definition~\ref{cv-hodge-lattice-germ}), let $\Pi_p$ be the series of
Lemma~\ref{cv-germs-cocycle-powers} for $B$, and let $D$ be the matrix of
Definition~\ref{cv-hodge-lattice-germ}(5) for $\ell=p$. Then there is a unique power series $\varphi(u,u')$
over $R$ in two families of $n$ variables such that $\Pi_p(x,y)=\varphi(x_1^p,\dots,x_n^p,y)$. It
satisfies $\varphi(0,u')=1=\varphi(u,0)$, and the coefficient of $u_iu'_j$ in $\varphi$ is $D_{ij}$.
\end{lemma}
\leanref{Development/AutomorphicSpectrum/Proofs/AlgebraProofs/PowOneSeriesPFormulaStatement.lean}{GlobalTheory.T.pow_one_series_p_formula}
\leanref{Development/WeilBound/CoreStatements/HodgeTatePoints/CbFrob.lean}{PositivityTheory.HtPt.Cb.first_power_char_p}
\leanref{Development/WeilBound/CoreStatements/HodgeTatePoints/CbFrob.lean}{PositivityTheory.HtPt.Cb.deriv_first_power_char_p}
\leanref{Development/AutomorphicSpectrum/Proofs/AlgebraProofs/AlphaKernelDerivativeVanishes.lean}{GlobalTheory.CompanionTheory.Ar.Kill.inj2}
\begin{proof}
For $e\in\mathbf Z_{\ge0}^n$ write $x^e=\prod_ix_i^{e_i}$ and $x^{pe}=\prod_ix_i^{pe_i}$.

Uniqueness. The coefficient of $u^eu'^f$ in $\varphi$ is the coefficient of $x^{pe}y^f$ in
$\varphi(x_1^p,\dots,x_n^p,y)$.

The derivatives in $x$. Let $F$ be the formal group law of $B$, and let $M(x)$ be the $n\times n$ matrix with
entries $M_{ck}(x)=(\partial F_c/\partial y_k)(x,0)$. Fix $k$, differentiate the identity of
Lemma~\ref{cv-germs-cocycle-powers}(2) for $m=p$ in $y_k$, and set $y=0$.
\begin{itemize}
\item On the left, the derivative of $\Phi_1(x,y,z)^p$ is $p\Phi_1^{p-1}$ times a series, which is $0$, and
  $\Phi_1(x,0,z)=1$. So the left side gives $\sum_c(\partial\Pi_p/\partial x_c)(x,z)\,M_{ck}(x)$.
\item On the right, $\Pi_p(x,z)$ does not involve $y$. The derivative of $\Pi_p(y,z)$ in $y_k$ vanishes at
  $y=0$: $\Pi_p(y,z)$ is a product of the series $\Phi_1([i]y,y,z)$, which are $1$ at $y=0$, and the
  derivative of each in $y_k$ at $y=0$ is a combination of $(\partial\Phi_1/\partial x_l)(0,0,z)$ and
  $(\partial\Phi_1/\partial y_k)(0,0,z)$, which vanish because $\Phi_1(x,0,z)=1=\Phi_1(0,y,z)$. By
  Lemma~\ref{cv-germs-frobenius-series}, $[p]y$ is a series in $y_1^p,\dots,y_n^p$, so its partial
  derivatives vanish, and so does the derivative of $\Phi_1([p]x,[p]y,z)$ in $y_k$. So the right side gives
  $0$.
\end{itemize}
Hence $\sum_c(\partial\Pi_p/\partial x_c)\,M_{ck}=0$ for every $k$. As $F(x,y)$ is $x+y$ plus terms of degree
at least $2$, $M(0)=1$, so $M$ is invertible, and $\partial\Pi_p/\partial x_c=0$ for every $c$.

The series $\varphi$. Write $\Pi_p=\sum_{e,f}a_{e,f}x^ey^f$. The coefficient of $x^ey^f$ in
$\partial\Pi_p/\partial x_c$ is $(e_c+1)a_{e+\delta_c,f}$, where $\delta_c$ is the $c$-th unit vector. An
integer prime to $p$ is a unit in $R$, so $a_{e,f}=0$ unless $p$ divides every $e_c$. Then
$\varphi=\sum_{e,f}a_{pe,f}u^eu'^f$ satisfies $\Pi_p(x,y)=\varphi(x_1^p,\dots,x_n^p,y)$.

Its properties. By Lemma~\ref{cv-germs-cocycle-powers}(1), $\varphi(0,y)=\Pi_p(0,y)=1$, and
$\varphi(x_1^p,\dots,x_n^p,0)=\Pi_p(x,0)=1$, so $\varphi(u,0)=1$ by the uniqueness of coefficients. The
coefficient of $u_iu'_j$ in $\varphi$ is that of $x_i^py_j$ in $\Pi_p$, which is $D_{ij}$.
\end{proof}
% Lean: pow_one_series_p_formula = Conj.pow_one_series_in_p = Cb.first_power_char_p: deriv_first_power_char_p
% (from first_power_add, with jac_cancel for M, d_act_p for [p] and ev_zero_pderiv_pow1_eq_zero with
% deriv_first_first_form_zero and deriv_second_first_form_zero), series_in_p for phi, first_power_zero_left
% and _right with ev_pow_inj for the axes, bil_ev_pow_two for the coefficients (CbFrob.lean, 254 lines).
% Uniqueness is Ar.Kill.inj2 (expand_injective').

\begin{lemma}\label{cv-germs-cocycle-kernel}
Let $E$ be a CM field, let $p$ be a prime, let $R$ be a commutative ring in which $p=0$, let $B$ be a
formal germ of dimension $n$ over $R$ (Definition~\ref{cv-hodge-lattice-germ}), and let $\varphi(u,u')$ be the
series of Lemma~\ref{cv-germs-cocycle-series} for $B$. Let $c\in R^n$ with $\alpha_B(c)=0$, where $\alpha_B$
is the conjugate map of Definition~\ref{cv-hodge-lattice-germ}(5) for $\ell=p$, and let
\[
 \delta_c(t)=\sum_{i=1}^nc_i\,\frac{\partial\varphi}{\partial u_i}(0,t),
\]
a power series in $t_1,\dots,t_n$. Then every partial derivative of $\delta_c$ vanishes.
\end{lemma}
\leanref{Development/AutomorphicSpectrum/Proofs/AlgebraProofs/AlphaKernelDerivativeVanishesStatement.lean}{GlobalTheory.T.alpha_kernel_derivative_vanishes}
\leanref{Development/AutomorphicSpectrum/Proofs/AlgebraProofs/AlphaKernelDerivativeVanishes.lean}{GlobalTheory.CompanionTheory.Ar.Kill.germKill}
\begin{proof}
Let $F$, $\Phi_2$, $W$ and $D$ be the data of $B$, so that $\alpha_B(c)=(Wc,D^{\mathsf t}c)$, and let $\Pi_p$ be
as in Lemma~\ref{cv-germs-cocycle-powers}. For a power series $f$ write $f^{(p)}$ for the series whose
coefficients are the $p$-th powers of those of $f$, and $y^p$ for $(y_1^p,\dots,y_n^p)$.

An identity. By Lemma~\ref{cv-germs-frobenius-series}, $[p]x=V(x_1^p,\dots,x_n^p)$ for an $n$-tuple $V$
without constant term whose matrix of linear coefficients is $W$. Then
\[
 \varphi(x,y)\,\varphi(x,z)\,\Phi_2(V(x),y,z)=\Phi_2^{(p)}(x,y^p,z^p)\,\varphi(x,y+_Fz).\tag{$*$}
\]
Indeed, substituting $x^p$ for $x$ turns the left side into $\Pi_p(x,y)\Pi_p(x,z)\Phi_2([p]x,y,z)$, and the
right side into $\Phi_2^{(p)}(x^p,y^p,z^p)\Pi_p(x,y+_Fz)=\Phi_2(x,y,z)^p\,\Pi_p(x,y+_Fz)$, as
$f(x,y,z)^p=f^{(p)}(x^p,y^p,z^p)$ for every power series $f$ over $R$. These agree by
Lemma~\ref{cv-germs-cocycle-powers}(3) for $m=p$. Substituting $x^p$ for $x$ maps distinct monomials to distinct
monomials, so $(*)$ holds.

A derivation. For a power series $f(x,y,z)$ let $\partial_cf(y,z)=\sum_ic_i(\partial f/\partial x_i)(0,y,z)$.
Then $\partial_c(ff')=f(0,y,z)\,\partial_cf'+f'(0,y,z)\,\partial_cf$. Every factor in $(*)$ is $1$ at $x=0$:
$\varphi(0,u')=1$, $V(0)=0$ and $\Phi_2(0,y,z)=1$, hence also $\Phi_2^{(p)}(0,y,z)=1$. Moreover
$\partial_c$ maps $\varphi(x,y)$, $\varphi(x,z)$ and $\varphi(x,y+_Fz)$ to $\delta_c(y)$, $\delta_c(z)$ and
$\delta_c(y+_Fz)$; it maps $\Phi_2(V(x),y,z)$ to $\sum_l(\partial\Phi_2/\partial x_l)(0,y,z)\,(Wc)_l=0$, as $Wc=0$;
and it maps $\Phi_2^{(p)}(x,y^p,z^p)$ to $\gamma(y^p,z^p)$, where
$\gamma(y,z)=\sum_ic_i(\partial\Phi_2^{(p)}/\partial x_i)(0,y,z)$. Applying $\partial_c$ to $(*)$ gives
\[
 \delta_c(y)+\delta_c(z)=\gamma(y^p,z^p)+\delta_c(y+_Fz).
\]

The derivatives. Differentiate in $z_j$ and set $z=0$. The derivative of $\gamma(y^p,z^p)$ vanishes, as
$p=0$ in $R$. The derivative of $\delta_c(z)$ at $z=0$ is the coefficient of $t_j$ in $\delta_c$, which is
$\sum_ic_iD_{ij}=(D^{\mathsf t}c)_j$ by Lemma~\ref{cv-germs-cocycle-series}, and $D^{\mathsf t}c=0$. So
\[
 \sum_k\frac{\partial\delta_c}{\partial t_k}(y)\,\frac{\partial F_k}{\partial z_j}(y,0)=0
\]
for every $j$. The matrix $\big((\partial F_k/\partial z_j)(y,0)\big)$ is invertible, as its constant term is $1$.
So $\partial\delta_c/\partial t_k=0$ for every $k$.
\end{proof}
% Lean: alpha_kernel_derivative_vanishes = Ar.Kill.germKill (AlphaKernelDerivativeVanishes.lean, 652 lines):
% multiplication_by_p_factors_frobenius for V, pow_one_series_p_formula and inj2 for phi; identity is (*)
% (inter2, inj3, power_series_pow_p_frobenius = Conj.pow_eq_ev_frobenius); additive applies the derivation
% (Dl, leib, dxFirst, dxV with lin V 1 = W, dxFrob); dz differentiates in z; linc and hbil give the
% coefficient of t_j; inv inverts the matrix. delta_c is dSer g c.

\begin{lemma}\label{cv-germs-cocycle-nfold}
Let $E$ be a CM field, let $G$ be an $\mathcal O_E$-semi-abelian scheme over a scheme $X$
(Definition~\ref{cv-model-toroidal-semiabelian}) with line bundle $\mathcal L$, let $\Lambda$ and $\theta$ be as
in Definition~\ref{cv-germs-chart}(1), and let $n\ge0$. For points $y_1,y_2$ of $G$ with values in a scheme
over $X$ let
\[
 \kappa_n(y_1,y_2)=\prod_{i=0}^{n-1}\theta(iy_1,y_1,y_2).
\]
\begin{enumerate}
\item For the two projections $y_1,y_2$ of $G\times_XG$, $\kappa_n(y_1,y_2)$ is a generator of
  $\Lambda(ny_1,y_2)\otimes\Lambda(y_1,y_2)^{-n}$, and $\kappa_n(0,y_2)=1$.
\item Let $g$ be a formal chart of dimension $d$ at a morphism $x_0:\operatorname{Spec}R\to X$
  (Definition~\ref{cv-germs-chart}), with points $x_m$, generators $e_m$ and formal germ $B$, and let $\Pi_n$
  be as in Lemma~\ref{cv-germs-cocycle-powers} for $B$. Then for every $m\ge0$, on
  $\operatorname{Spec}J^2_m$,
  \[
   \kappa_n(x_m(t),x_m(t'))\,e_m(t,t')^n=\Pi_n(t,t')\,e_m([n]t,t').
  \]
\end{enumerate}
\end{lemma}
\leanref{Development/AutomorphicSpectrum/Proofs/GeometryProofs/LineBundleGeneratorCorrespondenceStatement.lean}{GlobalTheory.T.line_bundle_generator_correspondence}
\leanref{Development/WeilBound/CoreStatements/HodgeTatePoints/CaNfold.lean}{PositivityTheory.HtPt.Ca.nfold_law}
\begin{proof}
(1) By Definition~\ref{cv-germs-chart}(1), $\theta(iy_1,y_1,y_2)$ is a generator of
$\Theta(iy_1,y_1,y_2)\cong\Lambda((i+1)y_1,y_2)\otimes\Lambda(iy_1,y_2)^{-1}\otimes\Lambda(y_1,y_2)^{-1}$. The tensor
product of these line bundles over $i<n$ is $\Lambda(ny_1,y_2)\otimes\Lambda(0,y_2)^{-1}\otimes\Lambda(y_1,y_2)^{-n}$, and
$\Lambda(0,y_2)\cong\mathcal O$. As $\theta$ is $1$ when one of its arguments is $0$,
$\kappa_n(0,y_2)=\prod_i\theta(0,0,y_2)=1$.

(2) By Definition~\ref{cv-germs-chart}(3)(a), $x_m([i]t)=i\,x_m(t)$, and $[i]t+_Ft=[i+1]t$. So the first
relation of Definition~\ref{cv-germs-chart}(3)(d), with $([i]t,t,t')$ substituted for $(t,t',t'')$, reads
\[
 \theta(i\,x_m(t),x_m(t),x_m(t'))=\Phi_1([i]t,t,t')\,e_m([i+1]t,t')\,e_m([i]t,t')^{-1}e_m(t,t')^{-1}.
\]
The product over $i<n$ is $\kappa_n(x_m(t),x_m(t'))=\Pi_n(t,t')\,e_m([n]t,t')\,e_m(0,t')^{-1}e_m(t,t')^{-n}$, and
$e_m(0,t')=1$ by Definition~\ref{cv-germs-chart}(3)(c).
\end{proof}
% Lean: line_bundle_generator_correspondence = Nfold.nfold_first_law = Ca.nfold_law (CaNfold.lean, 121 lines):
% exists_kappa (exists_cubN, with cubN theta n 1 y1 y2 = prod_i theta(i y1, y1, y2)) for (1), face_2 for
% kappa(0, y) = 1, value_first_power (induction with frame_biext_dictionary_first_partial_law, PP_nat) for (2).

\begin{lemma}\label{cv-germs-cocycle-generator}
Let $E$ be a CM field, let $G$ be an $\mathcal O_E$-semi-abelian scheme over a scheme $X$
(Definition~\ref{cv-model-toroidal-semiabelian}) with line bundle $\mathcal L$, let $k$ be a field of
characteristic $p>0$, let $x:\operatorname{Spec}k\to X$ be a morphism with fibre $G_x$, and let $g$ be a formal
chart of dimension $n$ at $x$ (Definition~\ref{cv-germs-chart}) with points $x_m$ and formal germ $B$. For
$m\ge0$ let $k_m=k[s]/(s^{m+1})$, let $Y_m=\operatorname{Spec}k_m\times_XG$, and let $h:Y_m\to G$ be the
projection. Let $\varphi$ be the series of Lemma~\ref{cv-germs-cocycle-series} for $B$, let $w\in k^n$, let
$D_w(t)=\sum_iw_i^p\,(\partial\varphi/\partial u_i)(0,t)$, and let $y=y_w$ be the point of $G$ with values in
$k_p$ of Lemma~\ref{cv-germs-frobenius-points}(2). Suppose that $py$ is the unit. Let $U$ be an open subset of
$G$ that contains the unit of $G_x$ and over which $\mathcal L$ has a generator $\ell$; through $\ell$,
generators of line bundles built from $\mathcal L$ at points of $G$ that lie in $U$ are read as invertible
functions. Then there is a generator $C$ of $\mathcal L_y^{\otimes p}$ on $Y_p$
(Definition~\ref{cv-germs-principal}(1)) such that:
\begin{enumerate}
\item $C$ is $1$ along $s=0$, where $y$ is the unit and $\mathcal L_y^{\otimes p}$ is canonically trivial;
\item for every $m\ge0$ there is $v_m\in J^1_m$ such that the pullback of $C$ along the point $(s,x_m)$ of $Y_p$
  with values in $J^1_m[s]/(s^{p+1})$, read through $\ell$, is $1+s^p\big(v_m^p-D_w(t)\big)$, with $D_w(t)$ read
  in $J^1_m$.
\end{enumerate}
\end{lemma}
\leanref{Development/AutomorphicSpectrum/Proofs/GeometryProofs/UnitsMulRelationRearrangeCommRing.lean}{GlobalTheory.CompanionTheory.Ar.Coc.Pf.pfold}
\begin{proof}
Write $0$ for the unit of $G$ with values in any scheme over $X$.

The generator. Let $\kappa_p$ be the generator of Lemma~\ref{cv-germs-cocycle-nfold}(1) for $n=p$. The
points $y$ and $0$ of $G$ with values in $k_p$ agree modulo $s$, and $\operatorname{Spec}k_p$ has one point,
so both map $\operatorname{Spec}k_p$ into an open subset $V$ of $G$ over which $\mathcal L$ has a generator
$\ell_V$. Then $\tau=y^*\ell_V\otimes0^*\ell_V^{-1}$ is a generator of $y^*\mathcal L\otimes0^*\mathcal L^{-1}$
on $\operatorname{Spec}k_p$ that is $1$ modulo $s$. Let $\kappa=\kappa_p(y,h)$, the pullback of $\kappa_p$ along
$(y,h):Y_p\to G\times_XG$, a generator of $\Lambda(py,h)\otimes\Lambda(y,h)^{-p}$. As $py=0$,
$\Lambda(py,h)\cong\mathcal O$; and
$\Lambda(y,h)=\mathcal L_y\otimes(y^*\mathcal L\otimes0^*\mathcal L^{-1})^{-1}$, the second factor pulled back from
$\operatorname{Spec}k_p$. So $C=\kappa^{-1}\tau^p$ is a generator of $\mathcal L_y^{\otimes p}$.

(1) Along $s=0$, $y$ is the unit, $\kappa_p(0,h)=1$ by Lemma~\ref{cv-germs-cocycle-nfold}(1), and $\tau=1$. So
$C=1$ there.

(2) Fix $m$, let $M=p+m+1$ and $A=J^1_m[s]/(s^{p+1})$. In $A$ the ideal generated by $s$ and $t_1,\dots,t_n$
has vanishing $(M+1)$-th power, as $s^{p+1}=0$ and $t$ generates an ideal whose $(m+1)$-th power is $0$. So
there is a homomorphism $j:J^2_M\to A$ with $j(t)=ws$ and $j(t')=t$. As $x_M$ reduces to $x_p$ and to
$x_m$, the composite of $(s,x_m)$ with $(y,h)$ is $(x_M(ws),x_M(t))$, the image under $j$ of the pair
$(x_M(t),x_M(t'))$. By Lemma~\ref{cv-germs-cocycle-nfold}(2), applied in $J^2_M$ and mapped by $j$,
\[
 \kappa(s,x_m)\,e_M(ws,t)^p=j(\Pi_p)\,e_M([p](ws),t).
\]
\begin{itemize}
\item By Lemma~\ref{cv-germs-cocycle-series}, $j(\Pi_p)=\varphi(w_1^ps^p,\dots,w_n^ps^p,t)$. As
  $(s^p)^2=0$ and $\varphi(0,t)=1$, this is $1+s^pD_w(t)$.
\item By Definition~\ref{cv-germs-chart}(3)(a), $x_p([p](ws))=p\,y$, which is the unit $x_p(0)$. By the
  uniqueness in Theorem~\ref{cv-germs-formal}(2), with $C=k_p$ and $N=(s)$, $[p](ws)=0$ in $k_p$. So
  $e_M([p](ws),t)=e_M(0,t)=1$.
\item The points that occur map $\operatorname{Spec}A$ into $U$, as they are the unit modulo a nilpotent
  ideal. Let $\varepsilon\in A^\times$ be $e_M(ws,t)$ read through $\ell$; modulo $s$ it is $e_M(0,t)=1$.
  Let $\tau_0\in k_p^\times$ be $\tau$ read through $\ell$; modulo $s$ it is $1$.
\end{itemize}
So $C(s,x_m)=\kappa(s,x_m)^{-1}\tau_0^p=(1+s^pD_w(t))^{-1}(\varepsilon\tau_0)^p$. Write $\varepsilon\tau_0=1+sf$
with $f\in A$, and let $v_m\in J^1_m$ be the image of $f$ modulo $s$. As $p=0$ in $A$ and $s^{2p}=0$,
$(\varepsilon\tau_0)^p=1+s^pf^p=1+s^pv_m^p$. Hence $C(s,x_m)=(1-s^pD_w(t))(1+s^pv_m^p)=1+s^p(v_m^p-D_w(t))$.
\end{proof}
% Lean: Ar.Coc.Pf.pfold (UnitsMulRelationRearrangeCommRing.lean, 106 lines): kappa from
% line_bundle_generator_correspondence (Lemma cv-germs-cocycle-nfold); tau from
% line_bundle_section_difference_generator = jetPair_gen (85 lines); C from frobenius_torsion_generator_relation
% = exists_gen_pfold (73 lines); (1) is frobenius_generator_correspondence = Pf.red_one (55 lines); in (2), j is
% jq, the composite is germ_torsion_lift_compat = TestPair.main, j(Pi_p) is germ_jet_quotient_formula =
% jq_psJet_ptPow (184 lines), [p](ws) = 0 is germ_frobenius_torsion_section = Pf.peG_act_one with
% germ_frobenius_torsion_jet = actp_vanish (245 lines), epsilon and tau_0 are germ_jet_unit_on_chart =
% frame_test_unit (137 lines) and jet_section_generator_value = Pf.red_unit (60 lines), v_m is
% jet_augmentation_one_pth_power = QAlg.pow_eq_one_add (55 lines), and the last step is
% units_mul_relation_rearrange_commRing. About 1,100 lines with the files they use.

\begin{lemma}\label{cv-germs-cocycle-first-order}
Let $X$ be a scheme, let $G$ be a scheme over $X$ with a line bundle $\mathcal L$, let $k$ be a field of
characteristic $p>0$, and let $x:\operatorname{Spec}k\to X$ be a morphism whose fibre $G_x$ is integral and
normal. For $m\ge0$ let $k_m=k[s]/(s^{m+1})$, let $Y_m=\operatorname{Spec}k_m\times_XG=
\operatorname{Spec}k_m\times_kG_x$, let $h:Y_m\to G$ be the projection, and let $r:Y_1\to Y_p$ be the morphism
induced by $k_p\to k_1$. Let $a:Y_p\to G$ be a morphism over $X$ that agrees with $h$ along $s=0$, and let $C$
be a generator of $a^*\mathcal L^{\otimes p}\otimes h^*\mathcal L^{\otimes-p}$ on $Y_p$ that is $1$ along $s=0$.
Let $U$ be an open subset of $G$ that meets $G_x$ and over which $\mathcal L$ has a generator $\ell$. On
$h^{-1}(U)=\operatorname{Spec}k_p\times_k(U\cap G_x)$, which $a$ also maps into $U$, write
$C=f\,(a^*\ell)^{\otimes p}\otimes(h^*\ell)^{\otimes-p}$ with $f=\sum_{i\le p}c_is^i$ and functions $c_i$ on
$U\cap G_x$. If $c_p$ is a $p$-th power in the function field of $G_x$, then
$(a\circ r)^*\mathcal L\otimes h^*\mathcal L^{-1}$ has a generator on $Y_1$.
\end{lemma}
\leanref{Development/WeilBound/CoreStatements/HodgeTatePoints/QamCore.lean}{PositivityTheory.HtPt.Qam.core_firstOrder}
\begin{proof}
Points and functions. The scheme $Y_p$ has the underlying space of its closed subscheme $s=0$, which is $G_x$,
and there $a$ agrees with $h$; so $a$ and $h$ agree on underlying points, and $h^{-1}(V)\subset a^{-1}(V)$ for
every open subset $V$ of $G$. As $k_m$ is free over $k$ with basis $1,s,\dots,s^m$, every function on
$\operatorname{Spec}k_m\times_kW$, for an open subset $W$ of $G_x$, is $\sum_{i\le m}b_is^i$ for unique
functions $b_i$ on $W$.

Expansions. Cover $G$ by open subsets $U_j$ over which $\mathcal L$ has generators $\ell_j$, one of which is
$(U_{j_0},\ell_{j_0})=(U,\ell)$, let $W_j=U_j\cap G_x$, and on $U_j\cap U_{j'}$ write $\ell_{j'}=g_{jj'}\ell_j$
with $g_{jj'}$ invertible. On $h^{-1}(U_j)$ write $C=f_j\,(a^*\ell_j)^{\otimes p}\otimes(h^*\ell_j)^{\otimes-p}$ with
$f_j=\sum_ic_{j,i}s^i$; then $c_{j,0}=1$, as $C$ is $1$ along $s=0$, and $c_{j_0,i}=c_i$. On
$h^{-1}(U_j\cap U_{j'})$ write $a^*g_{jj'}/h^*g_{jj'}=\sum_iu_{jj',i}s^i$; then $u_{jj',0}=1$, as $a$ agrees with
$h$ along $s=0$. Comparing the two expressions for $C$ gives $f_j=(\sum_iu_{jj',i}s^i)^pf_{j'}$. As $p=0$ and
$s^{p+1}=0$, $(\sum_iu_{jj',i}s^i)^p=1+u_{jj',1}^ps^p$. The coefficients of $s^p$ give, on $W_j\cap W_{j'}$,
\[
 c_{j,p}=c_{j',p}+u_{jj',1}^p.
\]

Roots. Let $K$ be the function field of $G_x$, of characteristic $p$, and let $\rho\in K$ with $\rho^p=c_{j_0,p}$.
Let $j$ be such that $W_j$ is not empty. As $G_x$ is irreducible, $W_j$ meets $W_{j_0}$, and
$\sigma_j=\rho+u_{jj_0,1}\in K$ satisfies $\sigma_j^p=c_{j_0,p}+u_{jj_0,1}^p=c_{j,p}$. So $\sigma_j$ is integral
over the local ring of $G_x$ at each point of $W_j$; these local rings are integrally closed, as $G_x$ is
normal, so $\sigma_j$ is a function on $W_j$. If $W_j$ is empty let $\sigma_j=0$. On $W_j\cap W_{j'}$,
$(\sigma_j-\sigma_{j'})^p=c_{j,p}-c_{j',p}=u_{jj',1}^p$, and the $p$-th power map of $K$ is injective, so
$u_{jj',1}=\sigma_j-\sigma_{j'}$.

The generator. On $Y_1$, where $s^2=0$, $(a\circ r)^*g_{jj'}/h^*g_{jj'}=1+u_{jj',1}s$. Let $t_j=1+\sigma_js$, a unit
on $\operatorname{Spec}k_1\times_kW_j$. Then $t_j=(1+u_{jj',1}s)\,t_{j'}$ on the overlaps, as
$u_{jj',1}=\sigma_j-\sigma_{j'}$ and $s^2=0$. So the generators $t_j\,(a\circ r)^*\ell_j\otimes h^*\ell_j^{-1}$ of
$(a\circ r)^*\mathcal L\otimes h^*\mathcal L^{-1}$ over the preimages of the $U_j$ agree on the overlaps, and they
glue to a generator on $Y_1$.
\end{proof}
% Lean: Qam.core_firstOrder (QamCore.lean, 391 lines): base_eq and le_pre (zero_section_surjective = secZ_surj)
% for the points, power_series_expansion_in_delta_unique = Thick.sections_thick (303 lines) for the
% expansions; chart_coef (c_{j,i}, c_{j,0} = 1), tr_coef (u_{jj',i}, u_{jj',0} = 1), overlap with
% truncated_expansion_pth_power_coeff = CoefAlg.coeff_comparison (87 lines); artin_schreier_cocycle_root =
% CechP.coboundary_of_pow_coboundary with CechRoot.exists_root and exists_section_of_pow (100 lines) for
% the roots; second (tU, tr_sum, ring_step) with line_bundle_cocycle_pullback_transform = exists_gen_of_units
% (71 lines) for the generator. The Lean's charts are the atlas of L, and j_0 is a chart with nonempty trace on
% the fibre.

\begin{lemma}\label{cv-germs-cocycle}
Let $E$ be a CM field, let $G$ be an $\mathcal O_E$-semi-abelian scheme over a scheme $X$
(Definition~\ref{cv-model-toroidal-semiabelian}) with line bundle $\mathcal L$, let $k$ be an algebraically
closed field of characteristic $p>0$, let $x:\operatorname{Spec}k\to X$ be a morphism, and let $g$ be a
formal chart of dimension $n$ at $x$ (Definition~\ref{cv-germs-chart}) with formal germ $B$. Let $w\in k^n$
with $\alpha_B(w_1^p,\dots,w_n^p)=0$, where $\alpha_B$ is the conjugate map of
Definition~\ref{cv-hodge-lattice-germ}(5) with $\ell=p$; let $y_w$ be the point of
Lemma~\ref{cv-germs-frobenius-points}(2), and $y_{w,1}$ its image in $G(k_1)$. Then the line bundle
$\mathcal L_{y_{w,1}}$ (Definition~\ref{cv-germs-principal}(1)) has a generator on
$\operatorname{Spec}k_1\times_XG$.
\end{lemma}
\leanref{Development/WeilBound/CoreStatements/HodgeTatePoints/CocpMain.lean}{PositivityTheory.HtPt.Cocp.cocycle_p}
\begin{proof}
Notation. For $m\ge0$ let $k_m=k[s]/(s^{m+1})$, let $Y_m=\operatorname{Spec}k_m\times_XG=
\operatorname{Spec}k_m\times_kG_x$, where $G_x$ is the fibre of $G$ at $x$, and let $h:Y_m\to G$ be the
projection. A point $y$ of $G$ with values in $k_m$ over $x$ gives the line bundle
$\mathcal L_y=(y+h)^*\mathcal L\otimes h^*\mathcal L^{-1}$ on $Y_m$ (Definition~\ref{cv-germs-principal}(1)).
Let $y=y_w$; then $\mathcal L_{y_{w,1}}$ is the restriction of $\mathcal L_y$ to $Y_1$. Let $U$ be an open
subset of $G$ that contains the unit of $G_x$ and over which $\mathcal L$ has a generator. Through it,
generators of line bundles built from $\mathcal L$ over the preimage of $U$ are read as invertible functions;
a function on $\operatorname{Spec}k_p\times_k(U\cap G_x)$ is $\sum_{i\le p}c_is^i$ with functions $c_i$ on
$U\cap G_x$. Let $x_m$ be the points of $g$ and $I$ the ideal of $J^1_m$ generated by $t$.

The point $y$. By Lemma~\ref{cv-germs-frobenius-points}(2), $y$ is the unit modulo $s$. The first component
of $\alpha_B(w_1^p,\dots,w_n^p)$ is $\sum_jW_{ij}w_j^p$ (Definition~\ref{cv-hodge-lattice-germ}(5)), so
it vanishes, and $py$ is the unit by Lemma~\ref{cv-germs-frobenius-points}(2).

The series $\varphi$. Let $\varphi(u,u')$ be the series of Lemma~\ref{cv-germs-cocycle-series} for $B$, so that
$\prod_{i=0}^{p-1}\Phi_1([i]u,u,u')=\varphi(u_1^p,\dots,u_n^p,u')$ and $\varphi$ is $1$ when $u=0$ or $u'=0$,
and let $D_w(t)=\sum_iw_i^p\,(\partial\varphi/\partial u_i)(0,t)$. As $\alpha_B(w_1^p,\dots,w_n^p)=0$,
Lemma~\ref{cv-germs-cocycle-kernel}, applied with $c=(w_1^p,\dots,w_n^p)$, shows that every partial
derivative of $D_w$ vanishes.

The generator $C$. As $py$ is the unit, Lemma~\ref{cv-germs-cocycle-generator}, applied with a generator of
$\mathcal L$ over $U$, gives a generator $C$ of $\mathcal L_y^{\otimes p}$ on $Y_p$ that is $1$ along $s=0$ and
such that for every $m$ there is $v_m\in J^1_m$ for which the pullback of $C$ along the point $(s,x_m)$ of $Y_p$
with values in $J^1_m[s]/(s^{p+1})$ is $1+s^p\bigl(v_m^p-D_w(t)\bigr)$, with $D_w(t)$ read in $J^1_m$.
% Lean: phi is gs from pow_one_series_p_formula (Lemma cv-germs-cocycle-series), D_w = dSer gs (w^p), and the
% derivatives vanish by alpha_kernel_derivative_vanishes = Ar.Kill.germKill (Lemma cv-germs-cocycle-kernel,
% AlphaKernelDerivativeVanishes.lean, 652 lines); C is Ar.Coc.Pf.pfold (Lemma cv-germs-cocycle-generator,
% 106 lines with line_bundle_generator_correspondence = Nfold.nfold_first_law).

The top coefficient. Let $C$ be as above, and let $C=\sum_{i\le p}c_is^i$ on
$\operatorname{Spec}k_p\times_k(U\cap G_x)$. Let $F$ be the power series of $c_p$ given by
Lemma~\ref{cv-germs-fibre-taylor}. The pullback of $C$ along $(s,x_m)$ is $\sum_ic_i(x_m)s^i$, so by the choice
of $C$, $c_p(x_m)=v_m^p-D_w(t)$ in $J^1_m$, and by Lemma~\ref{cv-germs-fibre-taylor}(1) the image of $F$ in $J^1_m$ is
$v_m^p-D_w(t)$. Let $V$ be a polynomial with image $v_m$. Then $F-V^p+D_w$ has no terms of degree at most
$m$, so its partial derivatives have no terms of degree less than $m$. As $\partial(V^p)/\partial t_j=
pV^{p-1}\partial V/\partial t_j=0$ and the partial derivatives of $D_w$ vanish, the partial derivatives
of $F$ have no terms of degree less than $m$. This holds for every $m$, so they vanish. By
Lemma~\ref{cv-germs-fibre-taylor}(3), $c_p$ is a $p$-th power in the function field of $G_x$.

Conclusion. By Lemma~\ref{cv-germs-fibre-taylor}, $G_x$ is integral and normal, and $U$ meets $G_x$ at its
unit. As $y$ is the unit modulo $s$, $y+h$ agrees with $h$ along $s=0$, and $\mathcal L_y=(y+h)^*\mathcal L\otimes
h^*\mathcal L^{-1}$. Let $r:Y_1\to Y_p$ be induced by $k_p\to k_1$; then $(y+h)\circ r=y_{w,1}+h$. By
Lemma~\ref{cv-germs-cocycle-first-order}, applied to $a=y+h$, to $C$ and to $U$, the line bundle
$\mathcal L_{y_{w,1}}=(y_{w,1}+h)^*\mathcal L\otimes h^*\mathcal L^{-1}$ has a generator on
$Y_1=\operatorname{Spec}k_1\times_XG$.
% Lean: cocycle_p: pt_of_vec (hy0, hy1), alpha_apply for the first component, the series, the generator and
% Qam.core_firstOrder (Lemma cv-germs-cocycle-first-order, with secZ_pa for a = y + h along s = 0 and
% trY_pa, trY_pb for the restriction to Y_1), formal_germ_jet_lift for F (hT, with hI, hN, hCh for the
% fibre), delta_expansion_restriction for c_p(x_m) and pderiv_zero_of_jet_pth_power_diff for the derivatives
% of F; U is the chart t.PP.L.U j0 containing the unit.
% Lean: PositivityTheory.HtPt.Cocp.cocycle_p (CocpMain.lean, 88 lines) uses Ce.pt_of_vec,
% pow_one_series_p_formula, alpha_kernel_derivative_vanishes (652 lines), formal_germ_jet_lift,
% Ar.Coc.Pf.pfold (106 lines), Qam.core_firstOrder (QamCore.lean, 391 lines), delta_expansion_restriction
% (136 lines) and pderiv_zero_of_jet_pth_power_diff (100 lines), with parts of IsPullbackWhisker.lean:
% about 1,600 lines in all. The steps it uses are Lemmas cv-germs-cocycle-powers to cv-germs-cocycle-first-order.
\end{proof}

\subsection{The Hodge lattice and the flat degree}\label{subsection-cv-hodge-lattice-flat-degree}

\begin{lemma}\label{cv-hodge-lattice-model-germs}
In Situation~\ref{cv-boundary-situation}, suppose that $O=\{b\in F_c:|\iota(b)|\le1\}$ for a prime
$\ell$ and an embedding $\iota:F_c\to\mathbf C_\ell$, and that the model has properties (1) and (2) of
Theorem~\ref{cv-model-toroidal}. Let $\mathcal Z$ be the closure of $j(Z)$ in $\mathcal X$, let $V$ be
the Lie algebra of $\mathcal A$ at the generic point of $\mathcal Z$, and for a point $z$ of
$\mathcal Z$ let $\mathcal O_z\subset F_c(Z)$ be its local ring. Then there are $n\ge0$, an
isomorphism $V\cong F_c(Z)^n$, a formal germ $B_\eta$ of dimension $n$ over $F_c(Z)$
(Definition~\ref{cv-hodge-lattice-germ}), and for every point $z$ of $\mathcal Z$ a formal germ $B_z$
of dimension $n$ over $\mathcal O_z$ and an isomorphism $(\phi_z,\psi_z,\lambda_z)$ from
$B_z\otimes F_c(Z)$ to $B_\eta$, with transport map $T_z$, such that:
\begin{enumerate}
\item each $b\in\mathcal O_E$ acts on $V=F_c(Z)^n$ by $\rho_{B_\eta}(b)$, and the Lie algebra of
  $\mathcal A$ at $z$ is $\Phi_z(\mathcal O_z^n)$;
\item if $z$ is a specialization of a point $z'$, so that $\mathcal O_z\subset\mathcal O_{z'}$, there
  is an isomorphism $(\phi',\psi',\lambda')$ from $B_z\otimes\mathcal O_{z'}$ to $B_{z'}$ with
  $T_z=T_{z'}\circ T_{\phi',\psi',\lambda'}$;
\item every point $z$ has an open neighbourhood $U$ such that
  $T_z(\mathcal O_{z'}^{2n})=T_{z'}(\mathcal O_{z'}^{2n})$ for every point $z'$ of $\mathcal Z\cap U$;
\item let $r\ge0$; if for every $a\in\mathcal O_E$ the characteristic polynomial of
  $\rho^H_{B_z}(a)$ is the $r$-th power of that of multiplication by $a$ on $\mathcal O_E$, then
  $H(B_z)$ is free of rank $r$ over $\mathcal O_E\otimes\mathcal O_z$;
\item if the residue field of $z$ has characteristic $\ell$, then for every homomorphism
  $\mathcal O_z\to k$ to a field that vanishes on the maximal ideal, the conjugate map
  $\alpha_{B_z\otimes k}$ is injective.
\end{enumerate}
\end{lemma}
\leanref{Development/AutomorphicSpectrum/Proofs/GeometryProofs/ExistsGermSystemStatement.lean}{GlobalTheory.T.exists_germ_system}
\leanref{Development/AutomorphicSpectrum/Definitions/GlobalGeometry/GermSys.lean}{GlobalTheory.CompanionTheory.Ar.Flat.GermSys}
\leanref{Development/AutomorphicSpectrum/Proofs/GeometryProofs/GermSystemHeightFreeStatement.lean}{GlobalTheory.T.germ_system_height_free}
\leanref{Development/AutomorphicSpectrum/Proofs/GeometryProofs/FibrePowerFreeConjInjectiveStatement.lean}{GlobalTheory.T.fibre_power_free_conj_injective}
\begin{proof}
The construction of the formal germs uses the following property of $\mathcal A$. The scheme
$\mathcal X$ is normal, so reduced, and locally noetherian, and its generic points lie in $j(X)$, over
which $\mathcal A$ is an abelian scheme. For $m\ge1$ let $\mathcal A^m$ be the $m$-fold fibre product
of $\mathcal A$ over $\mathcal X$. Then every function $f$ on $\mathcal A^m$ is the pullback of a
function on $\mathcal X$, namely of its pullback $g$ along the unit section. Over a generic point of
$\mathcal X$ the fibre of $\mathcal A^m$ is proper, geometrically connected and reduced, so $f$ is
constant on it and equal to the pullback of $g$ there. The scheme $\mathcal A^m$ is smooth over
$\mathcal X$, so it is reduced and its generic points lie over those of $\mathcal X$; hence $f$ is
the pullback of $g$.

The construction. By the preceding paragraph, Lemmas~\ref{cv-germs-local} and~\ref{cv-germs-compare}
apply to $\mathcal A$ over $\mathcal X$. The scheme $\mathcal X$ is normal, so its local ring at a point
$z$ of $\mathcal Z$ is a domain; it is the local ring $\mathcal O_z\subset F_c(Z)$, and the generic point
$\eta:\operatorname{Spec}F_c(Z)\to\mathcal X$ of $\mathcal Z$ factors through
$\operatorname{Spec}\mathcal O_z$. By Lemma~\ref{cv-germs-local}, $\eta$ factors through an affine open
subset of $\mathcal X$ with a formal chart, and by Lemma~\ref{cv-germs-compare}(1) the base change of that
chart is a formal chart $g_\eta$ at $\eta$, of some dimension $n$. Let $B_\eta$ be its formal germ and
$v^\eta_1,\dots,v^\eta_n$ its tangent vectors, and identify $F_c(Z)^n$ with $V$ by
$c\mapsto\sum_lc_lv^\eta_l$, a bijection by Definition~\ref{cv-germs-chart}(3)(b).

For each point $z$ of $\mathcal Z$, Lemma~\ref{cv-germs-local} gives an affine open neighbourhood
$U_z=\operatorname{Spec}A_z$ of $z$ with a formal chart $h_z$. For a point $z'$ of $\mathcal Z\cap U_z$,
the morphism $\operatorname{Spec}\mathcal O_{z'}\to\mathcal X$ factors through $U_z$, by a homomorphism
$A_z\to\mathcal O_{z'}$; since $U_z\to\mathcal X$ is a monomorphism, the composite of this homomorphism
with $\mathcal O_{z'}\subset F_c(Z)$ is the homomorphism $A_z\to F_c(Z)$ through which $\eta$ factors. By
Lemma~\ref{cv-germs-compare}(1), $h_z\otimes\mathcal O_{z'}$ is a formal chart at
$\operatorname{Spec}\mathcal O_{z'}\to\mathcal X$ whose base change to $F_c(Z)$ is $h_z\otimes F_c(Z)$, a
formal chart at $\eta$; by Lemma~\ref{cv-germs-compare}(2), applied at $\eta$, $h_z$ has dimension $n$. Let
$g_z=h_z\otimes\mathcal O_z$, let $B_z$ be its formal germ, and let $(\phi_z,\psi_z,\lambda_z)$ be a
comparison of $g_z\otimes F_c(Z)=h_z\otimes F_c(Z)$ with $g_\eta$ (Lemma~\ref{cv-germs-compare}(3)), an
isomorphism from $B_z\otimes F_c(Z)$ to $B_\eta$ with transport map $T_z$.

Property (1). Let $x_m$ be the points of $g_\eta$. By Definition~\ref{cv-germs-chart}(3)(a),
$b\,x_1(\varepsilon c)=x_1([b](\varepsilon c))=x_1(\varepsilon\rho_{B_\eta}(b)c)$ for $c\in F_c(Z)^n$, and
$x_1(\varepsilon c)=\sum_lc_lv^\eta_l$, as in the proof of Lemma~\ref{cv-germs-compare}(3). So each
$b\in\mathcal O_E$ acts on $V=F_c(Z)^n$ by $\rho_{B_\eta}(b)$. The Lie algebra of $\mathcal A$ at $z$ is
the $\mathcal O_z$-lattice in $V$ spanned by the images of the tangent vectors at
$\operatorname{Spec}\mathcal O_z\to\mathcal X$. By Definition~\ref{cv-germs-chart}(3)(b) for $g_z$, these
tangent vectors are the $\mathcal O_z$-combinations of those of $g_z$, whose images in $V$ are the tangent
vectors of $g_z\otimes F_c(Z)$. By Lemma~\ref{cv-germs-compare}(3), the $l$-th of these is
$\sum_j(\Phi_z)_{jl}v^\eta_j$, the $l$-th column of $\Phi_z$. So the Lie algebra of $\mathcal A$ at $z$ is
$\Phi_z(\mathcal O_z^n)$.

Property (2). Let $z$ be a specialization of $z'$, so that $\mathcal O_z\subset\mathcal O_{z'}$. By
Lemma~\ref{cv-germs-compare}(1), $g'=g_z\otimes\mathcal O_{z'}$ is a formal chart at
$\operatorname{Spec}\mathcal O_{z'}\to\mathcal X$ with formal germ $B_z\otimes\mathcal O_{z'}$, and
$g'\otimes F_c(Z)=g_z\otimes F_c(Z)$. Let $(\phi',\psi',\lambda')$ be a comparison of $g'$ with
$g_{z'}$, an isomorphism from $B_z\otimes\mathcal O_{z'}$ to $B_{z'}$. By
Lemma~\ref{cv-germs-compare}(3), its image over $F_c(Z)$ is a comparison of $g_z\otimes F_c(Z)$ with
$g_{z'}\otimes F_c(Z)$. Lemma~\ref{cv-germs-compare}(4), applied to $g_\eta$, $g_{z'}\otimes F_c(Z)$ and
$g_z\otimes F_c(Z)$, gives $T_z=T_{z'}\circ T_{\phi',\psi',\lambda'}$, the last transport map taken over
$F_c(Z)$.

Property (3). Let $z$ be a point of $\mathcal Z$ and $z'$ a point of $\mathcal Z\cap U_z$. The formal
chart $h_z\otimes\mathcal O_{z'}$ at $\operatorname{Spec}\mathcal O_{z'}\to\mathcal X$ has base change
$g_z\otimes F_c(Z)$ to $F_c(Z)$; let $(\phi',\psi',\lambda')$ be a comparison of it with $g_{z'}$. As in (2),
$T_z=T_{z'}\circ T'$, where $T'$ is the transport map of the image of $(\phi',\psi',\lambda')$ over
$F_c(Z)$. The matrices $\Phi'$, $\Psi'$ and $\Lambda'$ have entries in $\mathcal O_{z'}$ and
$\Phi'\Psi'=1$ (Lemma~\ref{cv-hodge-lattice-transport}(2)), so $T'$ and its inverse
$(x,y)\mapsto(\Psi'x,\Phi'^{\mathsf t}y-\Lambda'^{\mathsf t}\Psi'x)$ preserve $\mathcal O_{z'}^{2n}$.
Hence $T_z(\mathcal O_{z'}^{2n})=T_{z'}(\mathcal O_{z'}^{2n})$ for every point $z'$ of $\mathcal Z\cap U_z$.

Property (4) at a point $z$ whose residue field $k$ has characteristic $0$. The characteristic
polynomials of the $\rho^H_{B_z}(a)\otimes k$ are those of multiplication by $a$ on $\mathcal O_E$,
raised to the power $r$. The algebra $E\otimes k=\mathcal O_E\otimes k$ is a product of fields $k_i$,
and $H(B_z)\otimes k$ is the sum of $k_i$-vector spaces $H_i$. For $a$ generating $E$, the
characteristic polynomial of $a$ on $H(B_z)\otimes k$ is the product of the powers
$f_i^{\dim_{k_i}H_i}$ of the distinct irreducible factors $f_i$ of the minimal polynomial of $a$ over
$k$, and that of $a$ on $E\otimes k$ is the product of the $f_i$. So every $\dim_{k_i}H_i$ is $r$, and
$H(B_z)\otimes k$ is free of rank $r$ over $\mathcal O_E\otimes k$. Lift a basis to
$h_1,\dots,h_r\in H(B_z)$. For $b$ in a $\mathbf Z$-basis of $\mathcal O_E$ and $j\le r$, the vectors
$\rho^H_{B_z}(b)h_j$ span $H(B_z)=\mathcal O_z^{2n}$ by Nakayama's lemma, and there are
$r[E:\mathbf Q]=2n$ of them, by the degrees of the characteristic polynomials; so they form a basis.

The Hodge module at a principal point. Let $k$ be an algebraically closed field of characteristic
$\ell$, let $x:\operatorname{Spec}k\to\mathcal X$ be a morphism at which $\mathcal A$ has a principal fibre
(Definition~\ref{cv-germs-principal}), and let $g$ be a formal chart of dimension $n$ at $x$ whose formal
germ $B$ satisfies the hypothesis of (4) for some $r\ge0$. For $a\in\mathcal O_E$ let $N(a)$ be its norm
and $v_\ell$ the $\ell$-adic valuation, let $e=[E:\mathbf Q]$, and let $\beta_a$ be the matrix of
Definition~\ref{cv-hodge-lattice-germ}(2) for $B$. By the first paragraph of this proof,
Lemma~\ref{cv-germs-abelian-quotient} applies to $\mathcal A$ at $x$; let $\mathcal A_0$, with line bundle
$\mathcal L_0$, and $q:\mathcal A_x\to\mathcal A_0$ be as in its part (1). A formal chart $g_0$ of
$\mathcal A_0$ at the identity of $\operatorname{Spec}k$, of some dimension $n_0$ with formal germ $B_0$,
exists by Lemmas~\ref{cv-germs-local} and~\ref{cv-germs-compare}(1), since every function on a fibre power
of the proper, connected and reduced scheme $\mathcal A_0$ over the algebraically closed field $k$ is
constant. Let $\Lambda$, $P$ and $j$ be as in Lemma~\ref{cv-germs-abelian-quotient}(2), let $Z$ be as in
its part (3) for this $P$, and let $\beta^0_a$ be the matrix of Definition~\ref{cv-hodge-lattice-germ}(2)
for $B_0$. Let $\rho_0=\operatorname{rk}_{\mathcal O_E}(\Lambda)$, and let
$d_\Lambda=\dim_k(k\otimes\Lambda)$. The degrees of the characteristic polynomials of $1$ give
$d_\Lambda=e\rho_0$ by Lemma~\ref{cv-germs-lattice-kernel} and $2n=er$ by the hypothesis of (4), and
$n=n_0+d_\Lambda$ by Lemma~\ref{cv-germs-abelian-quotient}(2). Let $a\ne0$, let $\nu=v_\ell(N(a))=v_\ell(N(\bar a))$, and let
$d_a=\dim\ker\rho^H_{B_0}(a)$. In the notation of Lemma~\ref{cv-germs-kernel-bound},
$\rho^H_B(a)=T(\rho_B(a),\rho_B(\bar a),\beta_a)$ and $\rho^H_{B_0}(a)=T(\rho_{B_0}(a),\rho_{B_0}(\bar
a),\beta^0_a)$. By Lemma~\ref{cv-germs-kernel-bound}(1) with the matrix $Z$, and (2) with
$L=k\otimes\Lambda$, $S=1\otimes a$ and $\bar S=1\otimes\bar a$, and by
Lemma~\ref{cv-germs-lattice-kernel},
\[
 \dim\ker\rho^H_B(a)\le\dim\ker(1\otimes\bar a)+d_a+\dim\ker(1\otimes a)\le2\rho_0\nu+d_a.
\]
By Lemma~\ref{cv-germs-kernel-pairing} for $B_0$, $d_a\le\dim\ker\rho_{B_0}(\bar a)+\dim W_a$, where $W_a$
is the space of $u\in\ker\rho_{B_0}(a)$ with $z\cdot(\beta^0_a)^{\mathsf t}u=0$ for every
$z\in\ker\rho_{B_0}(\bar a)$. By Theorem~\ref{cv-germs-kernel-degree}, the kernel of the action of $a$
on $\mathcal A_0$ is $\operatorname{Spec}A_a$ for a finite $k$-algebra $A_a$. So
Lemma~\ref{cv-germs-kernel-hopf}, applied to $\mathcal A_0$, $g_0$ and $a$, gives a commutative and
cocommutative Hopf algebra $C_a$ with $\dim_kC_a=\dim_kA_a$ and maps $\delta_a$ and $\tau_a$. The spaces
$Y=\tau_a(\ker\rho_{B_0}(\bar a))$ and $Y'=\delta_a(W_a)$ have these dimensions and satisfy the hypothesis
of Lemma~\ref{cv-germs-hopf-count}, so $\ell^{d_a}$ divides $\dim_kC_a=\dim_kA_a$. By
Theorem~\ref{cv-germs-kernel-degree}, $(\dim_kA_a)^e=|N(a)|^{2n_0}$, so
$e\,d_a\le e\,v_\ell(\dim_kA_a)=2n_0\nu$. So
$e\dim\ker\rho^H_B(a)\le(2d_\Lambda+2n_0)\nu=2n\nu=er\nu$, and $\dim\ker\rho^H_B(a)\le r\nu$. As
$\rho^H_B$ is a ring homomorphism (Lemma~\ref{cv-hodge-lattice-transport}(1)),
Lemma~\ref{cv-germs-free-criterion} for $U=H(B)$, with the characteristic polynomials of the hypothesis
of (4), shows that $H(B)$ is free of rank $r$ over $\mathcal O_E\otimes k$.

Property (4) at a point $z$ whose residue field $\kappa$ has characteristic $\ell$. The set $j(X)$ is dense in
$\mathcal X$, as it contains the generic points of the locally noetherian scheme $\mathcal X$. By
Lemma~\ref{cv-germs-principal-fibres}, $\mathcal A$ is principal at $z$: there is a morphism
$\operatorname{Spec}k\to\mathcal X$ with image $z$, with $k$ algebraically closed, at which $\mathcal A$
has a principal fibre. It factors through $\operatorname{Spec}\mathcal O_z$; let $\mathcal O_z\to k$ be
the corresponding local homomorphism. By Lemma~\ref{cv-germs-compare}(1), $g_z\otimes k$ is a formal chart
at this point with formal germ $B_z\otimes k$, whose endomorphisms $\rho^H_{B_z\otimes k}(a)$ are the
images of the $\rho^H_{B_z}(a)$ and so have the characteristic polynomials of (4). By the preceding
paragraph, $H(B_z\otimes k)=H(B_z)\otimes k$ is free of rank $r$ over $\mathcal O_E\otimes k$. Then
$H_\kappa=H(B_z)\otimes\kappa$ is free of rank $r$ over the finite $\kappa$-algebra
$D=\mathcal O_E\otimes\kappa$: the module $H_\kappa\otimes_\kappa k$ is free of rank $r$ over
$D\otimes_\kappa k$, which is faithfully flat over $D$, so $H_\kappa$ is flat over $D$ and of rank $r$
at every prime of $D$, and a finitely generated flat module of constant rank $r$ over the semilocal ring
$D$ is free of rank $r$. Lift a basis of $H_\kappa$ to $h_1,\dots,h_r\in H(B_z)$. The vectors $\rho^H_{B_z}(b)h_j$, for
$b$ in a $\mathbf Z$-basis of $\mathcal O_E$ and $j\le r$, reduce to a basis of $\kappa^{2n}$, so the
matrix they form is square with determinant a unit of $\mathcal O_z$, and they form a basis of
$H(B_z)=\mathcal O_z^{2n}$.

Property (5). Let $z$ be a point whose residue field has characteristic $\ell$, and let
$\mathcal O_z\to k$ be a homomorphism to a field that vanishes on the maximal ideal, so that $k$ has
characteristic $\ell$. By Lemma~\ref{cv-germs-compare}(1), $g_z\otimes k$ is a formal chart at the point
$x:\operatorname{Spec}k\to\mathcal X$ at $z$, with formal germ $B_z\otimes k$. By
Lemma~\ref{cv-germs-principal-fibres}, whose density hypothesis holds as in (4), $\mathcal A$ has a principal
fibre at a morphism
$x':\operatorname{Spec}k'\to\mathcal X$ with image $z$ and $k'$ algebraically closed. As $x$ and $x'$
have the same image, $\operatorname{Spec}k\times_{\mathcal X}\operatorname{Spec}k'$ is not empty; an
algebraic closure $\Omega$ of the residue field at one of its points comes with homomorphisms $k\to\Omega$
and $k'\to\Omega$ through which $x$ and $x'$ give the same morphism $y:\operatorname{Spec}\Omega\to
\mathcal X$. By Lemma~\ref{cv-germs-principal-base-change}(1), $\mathcal A$ has a principal fibre at $y$,
and by Lemma~\ref{cv-germs-tangent}, which applies by the first paragraph of this proof, it is
tangentially principal at $y$ for $\ell$. By Lemma~\ref{cv-germs-compare}(1), $(g_z\otimes k)\otimes
\Omega$ is a formal chart at $y$ with formal germ $B_z\otimes\Omega$, so $\alpha_{B_z\otimes\Omega}$ is
injective by Lemma~\ref{cv-germs-conj-injective}, whose hypothesis on the line bundles
$\mathcal L_{y_{w,1}}$ holds by Lemma~\ref{cv-germs-cocycle}. The matrices $W$ and $D$ of $B_z\otimes\Omega$ are the images of those of
$B_z\otimes k$, so the matrix of $\alpha_{B_z\otimes\Omega}$ is the image of that of
$\alpha_{B_z\otimes k}$; a matrix over $k$ that is injective over $\Omega$ is injective over $k$. So
$\alpha_{B_z\otimes k}$ is injective.
% Lean: at a principal point, the abelian quotient, its lattice and Z are Lemma cv-germs-abelian-quotient
% (abelian_quotient_of_power_fns, abelian_quotient_lattice, abelian_quotient_bilinear); the chart g_0 is
% exists_formal_germ_local with proper_has_power_fns; the divisibility of dim A_a by l^{d_a} is
% proper_multiplication_kernel = Ar.Flat.lemF_ker_order_dvd, from proper_group_germ_finite_subscheme
% (Lemma cv-germs-kernel-hopf), primitive_derivation_pairing (Lemma cv-germs-hopf-count) and
% finrank_ker_rhoH_le (Lemma cv-germs-kernel-pairing). In (5), Lemma cv-germs-conj-injective with
% Lemma cv-germs-cocycle is tangent_power_free_conj_injective = Cocp.tanpf_inj = tanpf_inj_of cocycle_p.
\end{proof}
% Lean: exists_germ_system = PositivityTheory.Xgn.GermSystem.main: exists_germ (g_eta, by germ_aff,
% exists_factor and exists_bc), exists_loc with loc_aux0, st_factor and eq_mul_inverse_pow (h_z, its base
% changes and the dimension through germ_univ), exists_cmp and Cmp.iso (phi_z, psi_z, lambda_z); (1) is
% lieM (gd, gd_bij), lieM_act (gd_act), lieM_dir, key_dir (coord_tan) and GermHodgeLattice.lat_eq
% (comb_map); (2) is key_gen (key, Cmp.bc, Cmp.trans); (3) is key_loc (key, stdLat_comp, transH_surj).
% The first paragraph is powFns_of_abelian_generic = PositivityTheory.HtPt.Cf.pow_fns
% (bijective_of_generic). Property (4) is germ_system_height_free = HodgeFreeStalk.at_point: residue
% characteristic 0 is hfree_char0 (characteristic_polynomial_free_generators,
% representation_basis_linear_independent); residue characteristic l is formal_germ_base_change, hchar_map,
% height_char_free_geometric and hfree_of_map = free_module_nakayama_descent (int_descent,
% abstract_descent, descent, lift_basis). Property (5) is GermFam.exists_residue_germ
% (formal_germ_base_change, fibre_power_free_conj_injective).
% At residue characteristic l: the principal point is the hypothesis FibPF of germ_system_height_free,
% supplied by model_fibres_power_free (Lemma cv-germs-principal-fibres); the paragraph on the Hodge module
% at a principal point is height_char_free_geometric = Ar.Flat.hfree_geom (with LemFGeo.numeric and
% natDegree_charpoly_pow for the counts); Property (5) is fibre_power_free_conj_injective =
% PositivityTheory.HtPt.Cf.fib_inj = fib_inj_of (Cf.exists_common, geompf_bc, Ce.germ_bc, Cc.conjInj_map),
% with Ce.geompf_tan (Lemma cv-germs-tangent) and Cocp.tanpf_inj = tanpf_inj_of cocycle_p.

\begin{lemma}\label{cv-hodge-lattice-curve-degree}
Let $E$ be a number field with an involutive ring automorphism $a\mapsto\bar a$ of $\mathcal O_E$, let
$\ell$ be a prime, let $C$ be a proper normal curve over a field $k_0$
whose function field $\kappa$ has characteristic $\ell$, and let $n,r\ge0$. A vector bundle on $C$
with generic fibre a $\kappa$-vector space $W$ is given by its lattices over the local rings
$\mathcal O_c$ of $C$ at the closed points $c$. For a vector or a matrix $M$ over $\kappa$ let
$M^{(\ell)}$ be the vector or matrix of the $\ell$-th powers of its entries, and for a vector bundle
$\mathcal W$ with generic fibre $\kappa^n$ let $F^*\mathcal W$ be the vector bundle whose lattice at
$c$ is spanned by the $u^{(\ell)}$, $u\in\mathcal W_c$. Let $\mathcal L$ and $\mathcal H$ be vector
bundles on $C$ with generic fibres $\kappa^n$ and $\kappa^n\oplus\kappa^n$, let $\mathcal L^\vee$ be
the dual of $\mathcal L$ under the dot product, and let $\rho:\mathcal O_E\to M_n(\kappa)$ and
$\rho^H:\mathcal O_E\to\operatorname{End}_\kappa(\kappa^n\oplus\kappa^n)$ be ring homomorphisms such
that each $\rho^H(a)$ preserves $\mathcal H$, the first component of $\rho^H(a)(x,y)$ is
$\rho(a)x$, and $\rho^H(a)(0,y)=(0,\rho(\bar a)^{\mathsf t}y)$. Suppose that:
\begin{enumerate}
\item the inclusion of the second summand and the first projection give an exact sequence of vector
  bundles $0\to\mathcal L^\vee\to\mathcal H\to\mathcal L\to0$;
\item there are $\kappa$-linear maps $\alpha:\kappa^n\to\kappa^n\oplus\kappa^n$ and
  $\beta:\kappa^n\oplus\kappa^n\to\kappa^n$ that give an exact sequence of vector bundles
  $0\to F^*\mathcal L\to\mathcal H\to F^*(\mathcal L^\vee)\to0$, with
  $\alpha(\rho(a)^{(\ell)}u)=\rho^H(a)\alpha(u)$ and
  $\beta(\rho^H(a)v)=(\rho(\bar a)^{\mathsf t})^{(\ell)}\beta(v)$;
\item for every closed point $c$ there are $h_1,\dots,h_r\in\mathcal H_c$ such that the
  $\rho^H(b)h_j$, for $b$ in a $\mathbf Z$-basis of $\mathcal O_E$ and $j\le r$, form an
  $\mathcal O_c$-basis of $\mathcal H_c$;
\item for every $a\in\mathcal O_E$, the characteristic polynomial of multiplication by $a$ on
  $\mathcal O_E$ splits into linear factors over $\kappa$. (This hypothesis is not used in the proof.)
\end{enumerate}
Let $\chi:\mathcal O_E\to\kappa$ be a ring homomorphism, let $W_\chi$ be the sum of the images of
the $\rho^H(a)-\chi(a)$, and let $\mathcal H_\chi$ be the quotient bundle of $\mathcal H$ with generic
fibre $(\kappa^n\oplus\kappa^n)/W_\chi$. Then $\deg\mathcal H_\chi=0$.
\end{lemma}
\leanref{Development/AutomorphicSpectrum/Proofs/AlgebraProofs/SplitHermitianGeneratorsExistStatement.lean}{GlobalTheory.T.split_hermitian_generators_exist}
\leanref{Development/AutomorphicSpectrum/Definitions/CommutativeAlgebra/IsSplitGaussChars.lean}{GlobalTheory.CompanionTheory.Ar.Crv.ConjData}
\begin{proof}
For a vector bundle $\mathcal W$ with generic fibre $W$ and a subspace $W'$ of $W$, let
$\deg(\mathcal W,W')$ be the degree of the subbundle of $\mathcal W$ with generic fibre $W'$, whose
lattices are the $\mathcal W_c\cap W'$.

The image of $\chi$ is a finite domain, since $\mathcal O_E$ is a finitely generated abelian group
and $\ell\kappa=0$; so it lies in a finite subfield $k$ of $\kappa$. The nonzero elements of $k$ are
roots of unity, so they are units in every $\mathcal O_c$. The finite ring $A=k\otimes\mathcal O_E$
acts on $\kappa^n\oplus\kappa^n$ by $\rho^H_A(t\otimes a)=t\rho^H(a)$, which preserves $\mathcal H$,
and on $\kappa^n$ by $\rho_A(t\otimes a)=t\rho(a)$ and by
$\bar\rho_A(t\otimes a)=t\rho(\bar a)^{\mathsf t}$; and $\chi$ gives
$\chi_A:A\to k$, $t\otimes a\mapsto t\chi(a)$. A finite commutative ring is a product of local rings,
so there is an idempotent $e\in A$ with $\chi_A(e)=1$ such that every element of $e\ker\chi_A$ is
nilpotent.

Step 1: $\deg(\mathcal H,\rho^H_A(e)(\kappa^{2n}))=m\deg\mathcal H_\chi$ for some $m\ge1$. The ring
$eA$ is local, with nilpotent maximal ideal $e\ker\chi_A$ and residue field $k$. So there is a chain
of ideals $0=J_0\subsetneq J_1\subsetneq\dots\subsetneq J_m=eA$ of $A$, with $m\ge1$, such that
$J_i=J_{i-1}+At_i$ and $t_i\ker\chi_A\subset J_{i-1}$. For an ideal $J$ of $A$ let $W_J$ be the sum
of the images of the $\rho^H_A(t)$, $t\in J$; then $W_{eA}=\rho^H_A(e)(\kappa^{2n})$ and
$W_{\ker\chi_A}=W_\chi$. By (3), $\mathcal H_c$ is free over
$A\otimes_k\mathcal O_c=\mathcal O_E\otimes\mathcal O_c$, so $\mathcal H_c\cap W_J$ is
$J\mathcal H_c$, and $\rho^H_A(t_i)$ induces an isomorphism from $\mathcal H_\chi$ onto the quotient
of the subbundles of $\mathcal H$ with generic fibres $W_{J_i}$ and $W_{J_{i-1}}$. Hence
$\deg(\mathcal H,W_{J_i})=\deg(\mathcal H,W_{J_{i-1}})+\deg\mathcal H_\chi$, and summing over $i$
gives Step 1.

Step 2: $\deg(\mathcal H,\rho^H_A(x)(\kappa^{2n}))=0$ for every idempotent $x\in A$. Let
$F_A:A\to A$ be the ring endomorphism $t\otimes a\mapsto t^\ell\otimes a$. The idempotents of $A$ act
compatibly on the terms of the sequence (1), through $\bar\rho_A$, $\rho^H_A$ and $\rho_A$, and the
images of compatible idempotents in an exact sequence of vector bundles form an exact sequence of
subbundles. So
\[
 \deg(\mathcal H,\rho^H_A(x)(\kappa^{2n}))=\deg(\mathcal L^\vee,\bar\rho_A(x)(\kappa^n))
 +\deg(\mathcal L,\rho_A(x)(\kappa^n)).
\]
In the sequence (2), $A$ acts on the outer terms by $t\otimes a\mapsto t\rho(a)^{(\ell)}$ and
$t\otimes a\mapsto t(\rho(\bar a)^{\mathsf t})^{(\ell)}$, and there $F_A(x)$ acts by
$\rho_A(x)^{(\ell)}$ and $\bar\rho_A(x)^{(\ell)}$. The image of $M^{(\ell)}$ is spanned by the
$u^{(\ell)}$ for $u$ in the image of $M$, so the subbundles of $F^*\mathcal L$ and
$F^*(\mathcal L^\vee)$ that $F_A(x)$ cuts out are the pullbacks by $F$ of those of $\mathcal L$ and
$\mathcal L^\vee$ that $x$ cuts out, and pulling back by $F$ multiplies degrees by $\ell$. So the
sequence (2) gives
\[
 \deg(\mathcal H,\rho^H_A(F_A(x))(\kappa^{2n}))=\ell\deg(\mathcal L,\rho_A(x)(\kappa^n))
 +\ell\deg(\mathcal L^\vee,\bar\rho_A(x)(\kappa^n))=\ell\deg(\mathcal H,\rho^H_A(x)(\kappa^{2n})).
\]
If $k$ has $\ell^N$ elements, $F_A^N$ is the identity, so
$\deg(\mathcal H,\rho^H_A(x)(\kappa^{2n}))=\ell^N\deg(\mathcal H,\rho^H_A(x)(\kappa^{2n}))$, which
forces it to be $0$.

By Steps 1 and 2 for $x=e$, $m\deg\mathcal H_\chi=0$ with $m\ge1$.
\end{proof}
% Lean: deg_quot_eq_zero: exists_finite_subfield, v_le_one_of_finite, exists_idem; Step 1 is deg_count
% (ideal_descending_one_dim_quotient, algebra_action_bundle_filtration_step,
% module_basis_lattice_bundle_free) with idealV_ker; Step 2 is vanish: hodge_part and conj_part
% (filtered_exact_sequence_projector_bundles), range_twist and SubFrob.subDeg_frob in step, and
% frobT_iterate.

\begin{lemma}\label{cv-norms-hodge-lattice}
In Situation~\ref{cv-boundary-situation}, suppose that $O=\{b\in F_c:|\iota(b)|\le1\}$ for a prime
$\ell$ and an embedding $\iota:F_c\to\mathbf C_\ell$, and that the model has properties (1) and (2) of
Theorem~\ref{cv-model-toroidal}. Let $\varpi$ be a uniformizer of $O$, let $\mathcal Z$ be the closure
of $j(Z)$ in $\mathcal X$, and let $V$ be the Lie algebra of $\mathcal A$ at the generic point of
$\mathcal Z$, an $F_c(Z)$-vector space with an action of $\mathcal O_E$, with dual $V^\vee$. Then
there are an action of $\mathcal O_E$ on $V\oplus V^\vee$ that preserves $V^\vee$, acts on it by
$b\cdot\lambda=\lambda\circ\bar b$ and induces the given action on the quotient $V$, and a locally free
$\mathcal O_{\mathcal Z}$-submodule $\mathcal H$ of the constant sheaf $V\oplus V^\vee$ on $\mathcal Z$,
stable under $\mathcal O_E$, such that:
\begin{enumerate}
\item the image of $\mathcal H$ in $V$ is the Lie algebra of $\mathcal A$ on $\mathcal Z$, and the
  intersection of $\mathcal H$ with $V^\vee$ is its dual;
\item every point of $\mathcal Z$ has a neighbourhood on which $\mathcal H$ is free of rank $p+q$
  over $\mathcal O_E\otimes\mathcal O_{\mathcal Z}$;
\item for every embedding $\sigma':E\to F_c(Z)$, let $\mathcal H_{\sigma'}$ be the image of
  $\mathcal H$ under the projection of $V\oplus V^\vee$ onto the subspace on which $\mathcal O_E$
  acts through $\sigma'$; then the pullback of $\det\mathcal H_{\sigma'}$ to every proper normal curve
  over a field that maps into the special fibre of $\mathcal Z$ has degree $0$.
\end{enumerate}
\end{lemma}
\leanref{Development/AutomorphicSpectrum/Proofs/GeometryProofs/EigenLinePAdicBoundStatement.lean}{GlobalTheory.T.eigen_line_p_adic_bound}
\leanref{Development/AutomorphicSpectrum/Definitions/GlobalGeometry/LinearIndependentEigProj.lean}{GlobalTheory.CompanionTheory.Ar.Flat.HodgeLat}
\leanref{Development/AutomorphicSpectrum/Definitions/GlobalGeometry/LinearIndependentEigProj.lean}{GlobalTheory.CompanionTheory.Ar.Flat.HodgeLat.DegZero}
\begin{proof}
Let $n$, the isomorphism $V\cong F_c(Z)^n$, the formal germs $B_\eta$ and $B_z$, and the isomorphisms
$(\phi_z,\psi_z,\lambda_z)$ from $B_z\otimes F_c(Z)$ to $B_\eta$, with transport maps $T_z$, be given
by Lemma~\ref{cv-hodge-lattice-model-germs}, and let $r=p+q$. Identify $V^\vee$ with $F_c(Z)^n$
through the dot product, so that $V\oplus V^\vee=H(B_\eta)$, and let $\mathcal O_E$ act on it by
$\rho^H_{B_\eta}$, a ring homomorphism by Lemma~\ref{cv-hodge-lattice-transport}(1). By its
definition and Lemma~\ref{cv-hodge-lattice-model-germs}(1), this action induces the given action on
the quotient $V$, preserves $V^\vee$, and acts on it by $b\cdot\lambda=\lambda\circ\bar b$, since
$\rho_{B_\eta}(\bar b)^{\mathsf t}y$ represents $\lambda\circ\bar b$ when $y$ represents
$\lambda$. For a point $z$ of $\mathcal Z$ let $\mathcal H_z=T_z(\mathcal O_z^{2n})$, an
$\mathcal O_z$-lattice in $V\oplus V^\vee$ that is stable under $\mathcal O_E$ by
Lemma~\ref{cv-hodge-lattice-transport}(2).

Freeness at the points. The embeddings $s$ of $E$ into $F_c(Z)$ take values in $F_c$, and $V$ is the
sum of the subspaces $V_s$ on which $\mathcal O_E$ acts through $s$. By
Lemma~\ref{cv-model-eigenbundles}, applied to each embedding of the CM type,
$\dim V_s+\dim V_{\bar s}=p+q$, where $\bar s(b)=s(\bar b)$. The matrix of $\rho^H_{B_\eta}(a)$ is
block triangular with diagonal blocks $\rho_{B_\eta}(a)$ and $\rho_{B_\eta}(\bar a)^{\mathsf t}$, so
its characteristic polynomial is
\[
 \prod_s(X-s(a))^{\dim V_s}\prod_s(X-s(\bar a))^{\dim V_s}=\prod_s(X-s(a))^{\dim V_s+\dim V_{\bar s}}
 =\prod_s(X-s(a))^r,
\]
the $r$-th power of the characteristic polynomial of multiplication by $a$ on $\mathcal O_E$. By
Lemma~\ref{cv-hodge-lattice-transport}(2), $T_z$ conjugates $\rho^H_{B_z}(a)\otimes F_c(Z)$ into
$\rho^H_{B_\eta}(a)$, so $\rho^H_{B_z}(a)$ has the same characteristic polynomial, and $H(B_z)$ is
free of rank $r$ over $\mathcal O_E\otimes\mathcal O_z$ by
Lemma~\ref{cv-hodge-lattice-model-germs}(4). Choose a basis $h^z_1,\dots,h^z_r$ of it over
$\mathcal O_E\otimes\mathcal O_z$.

Properties (1) and (2). The image of $\mathcal H_z$ in $V$ is $\Phi_z(\mathcal O_z^n)$, the Lie algebra
of $\mathcal A$ at $z$ by Lemma~\ref{cv-hodge-lattice-model-germs}(1). The vector $T_z(x,y)$ lies in
$V^\vee$ when $\Phi_zx=0$, that is $x=0$; so $\mathcal H_z\cap V^\vee=\Psi_z^{\mathsf t}(\mathcal O_z^n)$,
the dual of $\Phi_z(\mathcal O_z^n)$ since $\Psi_z=\Phi_z^{-1}$. By
Lemma~\ref{cv-hodge-lattice-model-germs}(3), every point $z$ has a neighbourhood $U$ such that
$\mathcal H_{z'}=T_z(\mathcal O_{z'}^{2n})$ for the points $z'$ of $\mathcal Z\cap U$. So the
$\mathcal H_z$ are the stalks of a locally free $\mathcal O_{\mathcal Z}$-submodule $\mathcal H$ of
the constant sheaf, and the vectors
$\rho^H_{B_\eta}(b)T_z(h^z_j)=T_z(\rho^H_{B_z}(b)h^z_j)$, for $b$ in a $\mathbf Z$-basis of
$\mathcal O_E$ and $j\le r$, form an $\mathcal O_{z'}$-basis of $\mathcal H_{z'}$ for every $z'$ in
$\mathcal Z\cap U$. This proves (1) and (2).

Property (3): the eigenspace quotients. Let $\sigma':E\to F_c(Z)$ be an embedding, let
$W_{\sigma'}\subset V\oplus V^\vee$ be the subspace on which $\mathcal O_E$ acts through $\sigma'$, and
let $\pi$ be the projection onto it, which commutes with $\mathcal O_E$ and kills the images of the
$\rho^H_{B_\eta}(b)-\sigma'(b)$. The values $\sigma'(b)$, $b\in\mathcal O_E$, are integral over
$\mathbf Z$ and lie in $F_c$, so they lie in $O$ and in every $\mathcal O_z$. At a point $z$, the
lattice $\mathcal H_{\sigma'}$ is spanned by the vectors $\pi(\rho^H_{B_\eta}(b)T_z(h^z_j))=
\sigma'(b)\pi(T_z(h^z_j))$, so by the $\pi(T_z(h^z_j))$, $j\le r$; these form a basis of
$W_{\sigma'}$, whose dimension is $r$ by the characteristic polynomial. Fix a basis
$u_1,\dots,u_r$ of $W_{\sigma'}$, and let $Q^z$ be the $r\times r$ matrix with
$\pi(T_z(h^z_j))=\sum_lQ^z_{jl}u_l$. Then $\det Q^z\,u_1\wedge\dots\wedge u_r$ generates
$\det\mathcal H_{\sigma'}$ at $z$.

Let $C$ be a proper normal curve over a field $k_0$ with a morphism to $\mathcal Z$ that maps it into
the special fibre, let $\kappa$ be its function field, let $z_0$ be the image of its generic point,
and for a closed point $c$ of $C$, with local ring $\mathcal O_c$, let $z_c$ be its image. Then
$\varpi$, and with it $\ell$, is $0$ in $\kappa$. Each $z_c$ is a specialization of $z_0$, so
$\mathcal O_{z_c}\subset\mathcal O_{z_0}$, and the morphism gives local homomorphisms
$\mathcal O_{z_c}\to\mathcal O_c$ and $\mathcal O_{z_0}\to\kappa$ that agree on $\mathcal O_{z_c}$.

The bundles on $C$. By Lemma~\ref{cv-hodge-lattice-model-germs}(2), for each $c$ there is an
isomorphism $(\phi'_c,\psi'_c,\lambda'_c)$ from $B_{z_c}\otimes\mathcal O_{z_0}$ to $B_{z_0}$ with
transport map $T'_c$ and $T_{z_c}=T_{z_0}\circ T'_c$. Let $B=B_{z_0}\otimes\kappa$ and
$B_c=B_{z_c}\otimes\mathcal O_c$, and let $(\phi_c,\psi_c,\lambda_c)$ be the image of
$(\phi'_c,\psi'_c,\lambda'_c)$ over $\kappa$, an isomorphism from $B_c\otimes\kappa$ to $B$ with
transport map $T_c$. Let $\mathcal L$ and $\mathcal H_C$ have the lattices
$\Phi_c(\mathcal O_c^n)\subset\kappa^n$ and $T_c(\mathcal O_c^{2n})\subset\kappa^n\oplus\kappa^n$ at
the closed points $c$. By Lemma~\ref{cv-hodge-lattice-model-germs}(3) at $z_0$, there is a
neighbourhood $U$ of $z_0$ such that $T_{z_0}(\mathcal O_{z_c}^{2n})=T_{z_c}(\mathcal O_{z_c}^{2n})=
T_{z_0}(T'_c(\mathcal O_{z_c}^{2n}))$ when $z_c\in U$; then $T'_c$ preserves $\mathcal O_{z_c}^{2n}$,
so $\Phi'_c$, $\Psi'_c$ and $\Lambda'_c$ have entries in $\mathcal O_{z_c}$, and $\Phi_c$, $\Psi_c$
and $\Lambda_c$ have entries in $\mathcal O_c$. Only finitely many closed points of $C$ map outside
$U$, so $\mathcal L$ and $\mathcal H_C$ are vector bundles on $C$.

They satisfy the hypotheses of Lemma~\ref{cv-hodge-lattice-curve-degree} with $\rho=\rho_B$,
$\rho^H=\rho^H_B$, $\alpha=\alpha_B$ and $\beta=\beta_B$. By
Lemma~\ref{cv-hodge-lattice-transport}(2), $\rho^H_B(a)T_c=T_c\rho^H_{B_c}(a)$, so $\rho^H_B(a)$
preserves $\mathcal H_C$; the conditions on the first component and on $(0,y)$ hold by the
definition of $\rho^H_B$. The sequence (1) is exact at each $c$ as for $\mathcal H_z$ above, since
$\Psi_c=\Phi_c^{-1}$. For (2), Lemma~\ref{cv-hodge-lattice-transport}(3) gives the relations of
$\alpha_B$ and $\beta_B$ with the action, and $\alpha_B(\Phi_c^{(\ell)}u)=T_c(\alpha_{B_c}(u))$ and
$\beta_B(T_cv)=(\Psi_c^{\mathsf t})^{(\ell)}\beta_{B_c}(v)$; the lattices of $F^*\mathcal L$ and
$F^*(\mathcal L^\vee)$ at $c$ are $\Phi_c^{(\ell)}(\mathcal O_c^n)$ and
$(\Psi_c^{\mathsf t})^{(\ell)}(\mathcal O_c^n)$. So it suffices that
$0\to\mathcal O_c^n\to\mathcal O_c^{2n}\to\mathcal O_c^n\to0$, with the maps $\alpha_{B_c}$ and
$\beta_{B_c}$, is exact. It is a complex, by Lemma~\ref{cv-hodge-lattice-transport}(3). The residue
field $k(c)$ of $\mathcal O_c$ has characteristic $\ell$, and so does that of $z_c$; by
Lemma~\ref{cv-hodge-lattice-model-germs}(5), applied to $\mathcal O_{z_c}\to\mathcal O_c\to k(c)$,
$\alpha_{B_c}$ is injective modulo the maximal ideal, so its image is a direct summand of rank $n$.
The pairing $\langle\,,\rangle_{B_c}$ is perfect: its matrix has determinant $1$. Given
$w\in\mathcal O_c^n$, the linear form $\alpha_{B_c}(u)\mapsto u\cdot w$ on that direct summand extends
to $\mathcal O_c^{2n}$ and is $\langle\,\cdot\,,v\rangle_{B_c}$ for some $v$, and then
$\beta_{B_c}(v)=w$ by Lemma~\ref{cv-hodge-lattice-transport}(3). The kernel of $\beta_{B_c}$ is the
orthogonal of the image of $\alpha_{B_c}$, a direct summand of rank $n$ that contains that image,
which is a direct summand of the same rank; so the two are equal. For (3), the images of the
$h^{z_c}_j$ in $\mathcal O_c^{2n}$ form a basis of $H(B_c)$ over $\mathcal O_E\otimes\mathcal O_c$,
and $T_c$ carries them to vectors as required. For (4), the characteristic polynomial of
multiplication by $a\in\mathcal O_E$ on $\mathcal O_E$ is $\prod_\sigma(X-\sigma(a))$ over the
embeddings $\sigma:E\to F_c$, whose values on $\mathcal O_E$ lie in $O\subset\mathcal O_{z_0}$; so it
splits over $\kappa$.

Let $\chi:\mathcal O_E\to\kappa$ be $\sigma'$ followed by $\mathcal O_{z_0}\to\kappa$. By
Lemma~\ref{cv-hodge-lattice-curve-degree}, the quotient bundle $\mathcal H_{C,\chi}$ has degree $0$.

The comparison with $\det\mathcal H_{\sigma'}$. The module $H(B_{z_0})$, divided by the
$\mathcal O_{z_0}$-span of the images of the $\rho^H_{B_{z_0}}(b)-\sigma'(b)$, is free over
$\mathcal O_{z_0}$ with basis the images of the $h^{z_0}_l$. So there is an $r\times r$ matrix $M^c$
over $\mathcal O_{z_0}$ such that $T'_c(h^{z_c}_j)-\sum_lM^c_{jl}h^{z_0}_l$ lies in that span.
Applying $\pi\circ T_{z_0}$, which kills it, gives $Q^{z_c}=M^cQ^{z_0}$. Let $N^c$ be the image of
$M^c$ over $\kappa$. Then $T_c(\bar h^c_j)-\sum_lN^c_{jl}\bar h_l\in W_\chi$, where $\bar h_l$ and
$\bar h^c_j$ are the images of $h^{z_0}_l$ in $\kappa^{2n}$ and of $h^{z_c}_j$ in
$\mathcal O_c^{2n}$. Since $\mathcal O_E$ acts on $\kappa^{2n}/W_\chi$ through $\chi$, the images of
the $\bar h_l$ form a basis of $\kappa^{2n}/W_\chi$ and those of the $T_c(\bar h^c_j)$ a basis of the
lattice of $\mathcal H_{C,\chi}$ at $c$. So $\det N^c\ne0$ and
\[
 0=\deg\mathcal H_{C,\chi}=-\sum_c[k(c):k_0]\operatorname{ord}_c(\det N^c).
\]
As $\det N^c\ne0$ and $\mathcal O_{z_0}\to\kappa$ is local, $\det M^c$ is a unit in $\mathcal O_{z_0}$.
So the pullback of $\det\mathcal H_{\sigma'}$ to $C$ is generated at the generic point by the image
of $\det Q^{z_0}\,u_1\wedge\dots\wedge u_r$, and at $c$ by the image of
$\det Q^{z_c}\,u_1\wedge\dots\wedge u_r$, which is $\det N^c$ times the first; its degree is the same
sum, which is $0$.
\end{proof}
% Lean: hodgeLat_model. Lemma cv-hodge-lattice-model-germs is exists_germ_system (with the LieModel
% transported to V by lie_model_existsUnique_equiv and hodge_lattice_transport_equiv),
% model_fibres_power_free, powFns_of_abelian_generic and germ_system_height_free; the characteristic
% polynomial is hermitian_signature_character = hodge_character_generic, with the ranks from
% proper_normal_compactification. (1) and (2) are exists_hodge_lattice_germ_system = germ_gives_hodgeLat
% (GermHodgeLattice.lat_eq, fst_stdLat, inr_stdLat, exists_finite, exists_open, S.loc). (3) is
% hodge_lattice_dvr_degree = germ_degZero: hodge_lattice_germ_eigen and hodge_lattice_span_independent
% (the matrices Q^z), fibre_curve_germ_powers = GermFam.germ_family (the germs on C, with
% exists_residue_germ and fibre_curve_degree_compare), exists_trans (M^c), and valuation_bounded_frame
% = DegProd.curve_deg: family_hodge_lattice_data = family_conj_H_lattices (hodge_family, conj_family,
% locSplit with local_matrix_coprime_decomposition), split_hermitian_generators_exist and
% quot_deg_lemma_4 (exists_quot_basis, isLatOf_quot, bundle_degree_eq_sum_valuations).

\begin{lemma}\label{cv-norms-padic}
In Situation~\ref{cv-boundary-situation}, suppose that $O=\{b\in F_c:|\iota(b)|\le1\}$ for a prime
$\ell$ and an embedding $\iota:F_c\to\mathbf C_\ell$, and that the model has properties (1) and (2) of
Theorem~\ref{cv-model-toroidal}. Let $s$ be a rational section of $\mathcal L$ that generates it at
every point of $\bar Z$, and let $\varphi_s$ be its frame function
(Lemma~\ref{cv-function-generator}). Then there is $t\in\mathbf R$ with
$\log|\iota_K(\varphi_s(y,v,w))|=t$ for every discrete valuation ring $V_K$ with fraction field $K$,
every embedding $\iota_K:K\to\mathbf C_\ell$ that extends $\iota$ on $F_c\subset K$ and has
$V_K=\{b\in K:|\iota_K(b)|\le1\}$, every framed point $(y,v,w)$ of $Z$ with values in $K$, and every
semi-abelian scheme $T$ over $V_K$ with an action of $\mathcal O_E$ and a polarization whose generic
fibre is $A_y$ with its action and polarization, such that $v$ and $w$ are bases over $V_K$ of the
submodules of the Lie algebra of $T$ on which $\mathcal O_E$ acts by $\tau$ and by $\bar\tau$.
\end{lemma}
\leanref{Development/AutomorphicSpectrum/Proofs/GeometryProofs/FiniteConstantAllHoldsStatement.lean}{GlobalTheory.T.finite_constant_all_holds}
\leanref{Development/AutomorphicSpectrum/Proofs/GeometryProofs/HodgeLatticeValuationAgreeStatement.lean}{GlobalTheory.T.hodge_lattice_valuation_agree}
\leanref{Development/AutomorphicSpectrum/Proofs/GeometryProofs/EigenLinePAdicFrameStatement.lean}{GlobalTheory.T.eigen_line_p_adic_frame}
\begin{proof}
Let $\varpi$ be a uniformizer of $O$ and $\mathcal Z$ the closure of $j(Z)$ in $\mathcal X$. It is
normal, integral and proper over $O$, with generic point over the generic point of
$\operatorname{Spec}O$, and $F_c$ is algebraically closed in its function field $F_c(Z)$. Let $V$
be the Lie algebra of $\mathcal A$ at the generic point of $\mathcal Z$, let
$\mathcal H\subset V\oplus V^\vee$ be the lattice of Lemma~\ref{cv-norms-hodge-lattice}, and let
$W_\tau$ be the subspace of $V\oplus V^\vee$ on which $\mathcal O_E$ acts through $\tau$; by property
(2) there, $\det\mathcal H_\tau$ is a line bundle on $\mathcal Z$.

Choose bases $v^0$ and $w^0$ of $\operatorname{Lie}_\tau(\mathcal A)$ and
$\operatorname{Lie}_{\bar\tau}(\mathcal A)$ at the generic point, and write
$s=r\wedge v^0\otimes(\wedge w^0)^{-1}$ with $r\in F_c(Z)$. Let $u$ be a basis of $W_\tau$ adapted to
them: $u_1,\dots,u_p$ map to $v^0_1,\dots,v^0_p$ in $V$, and $u_{p+1},\dots,u_{p+q}$ is the basis of
the dual of $\operatorname{Lie}_{\bar\tau}(\mathcal A)$ at the generic point dual to $w^0$, extended
by $0$ on the other eigenspaces, which lies in $W_\tau$ because $b\cdot\lambda=\lambda\circ\bar b$.
At every point of the generic fibre, $\mathcal H_\tau$ is the sum of
$\operatorname{Lie}_\tau(\mathcal A)$ and the dual of $\operatorname{Lie}_{\bar\tau}(\mathcal A)$, by
property (1) of Lemma~\ref{cv-norms-hodge-lattice} and the splitting of $\mathcal O_E\otimes F_c$;
so $r\wedge u$ generates $\det\mathcal H_\tau$ there. By property (3) of
Lemma~\ref{cv-norms-hodge-lattice} and Lemma~\ref{cv-norms-zariski}, there are $M>0$ and $a$ such
that $r^M\varpi^{-a}(\wedge u)^{\otimes M}$ generates $(\det\mathcal H_\tau)^{\otimes M}$ at every
point of $\mathcal Z$.

Now let $V_K$, $K$, $\iota_K$, $(y,v,w)$ and $T$ be as in the statement. Since $\mathcal Z$ is proper
over $O$ and $O\subset V_K$, the point $y$ extends to a morphism
$z:\operatorname{Spec}V_K\to\mathcal Z$. The semi-abelian schemes $T$ and $z^*\mathcal A$ over $V_K$
have the same generic fibre, so they are isomorphic \cite[Chapter~I, Proposition~2.7]{FC90}, and $v$
and $w$ are bases over $V_K$ of the submodules of $z^*\operatorname{Lie}(\mathcal A)$ on which
$\mathcal O_E$ acts by $\tau$ and by $\bar\tau$. By property (2) of
Lemma~\ref{cv-norms-hodge-lattice}, $z^*\mathcal H$ is free of rank $p+q$ over
$\mathcal O_E\otimes V_K$, with image $z^*\operatorname{Lie}(\mathcal A)$ in $V\otimes K$ and with its
dual as intersection with $V^\vee\otimes K$. A computation with such modules after extending scalars
to the ring of integers of $\mathbf C_\ell$ shows: the matrix that expresses a basis of the image of
$z^*\mathcal H$ in $W_\tau\otimes K$ in terms of the basis of $W_\tau\otimes K$ adapted to $v$ and $w$
has determinant $\gamma\varepsilon$, where $\varepsilon$ is a unit of the ring of integers of
$\mathbf C_\ell$ and $\gamma\in\mathbf C_\ell^\times$ depends only on $E$, $\iota\tau$ and $(p,q)$.
Comparing with the generator $r^M\varpi^{-a}(\wedge u)^{\otimes M}$ at the closed point of $z$ gives
$\iota_K(\varphi_s(y,v,w))=\iota_K(c)\,\varepsilon\gamma$ for some $c\in K$ with $c^M=\varpi^a$ times
a unit of $V_K$. Hence
\[
 \log|\iota_K(\varphi_s(y,v,w))|=\frac aM\log|\iota(\varpi)|+\log|\gamma|,
\]
which does not depend on the point.
\end{proof}
% Lean: finConstAll: proper_normal_compactification, constClosed_of_isGeom, eigen_line_p_adic_bound
% (Lemma cv-norms-hodge-lattice), hermitian_signature_basis (the adapted basis, IsAdapted),
% hodge_lattice_valuation_agree = COrder.order_on_component (GenAgree.generic_agree and Lemma
% cv-norms-zariski), and eigen_line_p_adic_frame: dvr_point_eigen_frame_reduction (the extension z),
% dvr_specialization_eigen_frame = IntFrameSpan.span_eq_eigSub (semiabelian_extension_over_dvr),
% cm_type_valuation_embedding (gamma) with residue_hodge_lattice_adapted_basis,
% order_valuation_compare and log_end.

\begin{theorem}\label{cv-flat-degree}
Let $E$ be an abelian CM field and let $d\equiv7\pmod{16}$ be a prime with
$\sqrt{-d}\in E$, such that $E$ is unramified over its totally real subfield and every
prime $w$ of $E$ dividing the different of $E/\mathbf Q$ is moved by complex conjugation
and lies either over a prime $\ell\ne d$ for which $-d$ is a square modulo $\ell$, or over
$d$ with $w^3\nmid d$. Let $\chi$ be the Kronecker character of $\mathbf Q(\sqrt{-d})$.
Let a PEL datum of signature $(p,q)$ over $E$ with $p+q\ge1$, a level of full level $1$
(Definition~\ref{cv-model-toroidal-semiabelian}(1)), its moduli scheme $X$ and coefficient field $F_c$ be
given, with uniformizations at the complex embeddings of $F_c$,
normalized by a real number, that satisfy the compatibility conditions of the moduli problem;
among them, the prime-to-$\ell$ part of the level is at least $3$ for every prime $\ell$, and
$F_c$ is algebraically closed in the function field of every connected component. Then for every connected component $Z$ of $X$ and every
$\tau$ in the CM type $\Phi$ of the datum, flat sections of type $\tau$ on $Z$ exist
(Definition~\ref{cv-constant-sections}, with the prime $d$), and every such section has degree
\[
 \frac{p-q}{2}\log\Omega_E(\chi)+q\big(\log d-\log\operatorname{rd}_E\big)
\]
(Definition~\ref{cv-log-omega}).
\end{theorem}
\leanref{Development/AbelianVarietyTheory/Definitions/Divisors/HtValOnSubset.lean}{GlobalTheory.FlatCore.FlatDegreeSt}
\leanref{Development/AutomorphicArithmetic/ThetaCorrespondence/ThetaPairLifts/ThirdPlaceSeries/FlatDegreeStatement.lean}{GlobalTheory.CovTop.FlatDegProof.flatDeg_alt_alt}
\leanref{Development/AutomorphicSpectrum/Proofs/DivisibilityProofs/FlatDegreeStOfResKGaga.lean}{GlobalTheory.CompanionTheory.Ar.Br.flatDegreeSt_of_resK_gaga}
\leanref{Development/AutomorphicSpectrum/Proofs/GeometryProofs/FlatDegreeHolds.lean}{GlobalTheory.CompanionTheory.Ar.Fl.flatDegree}
\begin{proof}
Fix $Z$ and $\tau$.

Step 1: a section with constant norms. The data of the theorem are in
Situation~\ref{cv-function-situation}, with $p+q\ge1$ and $F_c$ algebraically closed in $F_c(Z)$.
Sections of $\omega^{\lambda_\tau}$ on $Z$ are described by frame functions on $Z$
(Definition~\ref{cv-function-framed}), through the determinants of the two eigenspaces.
\begin{enumerate}
\item[(1a)] Choose an embedding $\iota_2:F_c\to\mathbf C_2$. Since the prime-to-$2$ part of the level
  is at least $3$, Theorem~\ref{cv-model-toroidal} gives a model $(\mathcal X,\mathcal A,j)$ of $X$
  over $O=\{b\in F_c:|\iota_2(b)|\le1\}$. By Lemma~\ref{cv-model-eigenbundles}, $Z$ is open and
  closed in $X$, and
  $\mathcal L=\det\operatorname{Lie}_\tau(\mathcal A)\otimes(\det\operatorname{Lie}_{\bar\tau}(\mathcal A))^{-1}$
  is a line bundle on the closure $\bar Z$ of $j(Z)$ in $\mathcal X_{F_c}$; let $b_1$ and $b_2$ be
  bases of the two eigenbundles at its generic point.
\item[(1b)] Let $\sigma_0$ be the inclusion of $F_c$ in $\mathbf C$. Then
  $\sigma_0(\tau(\sqrt{-d}))=\theta$ because $\tau\in\Phi$, and the uniformization at $\sigma_0$
  gives the function $\Phi_{\sigma_0}$ of Lemma~\ref{cv-periods-function}, built from the periods.
  These data form Situation~\ref{cv-boundary-situation} with $\sigma=\sigma_0$.
  Lemma~\ref{cv-boundary-section} gives $f\in F_c(Z)$ and $\kappa_0\in\mathbf C^\times$ such that
  $s=f\wedge b_1\otimes(\wedge b_2)^{-1}$ generates $\mathcal L$ at every point of $\bar Z$, and
  $f(y)=\kappa_0\Phi_{\sigma_0}(y,v,w)$ at the complex points $y$ over $\sigma_0$ of a dense open subset
  of $Z$, with $v$ and $w$ the values of $b_1$ and $b_2$. Let $\varphi=\varphi_s$ be the frame
  function of Lemma~\ref{cv-function-generator}(1): $\varphi(y,v,w)$ is the value at $y$ of
  $\alpha/\beta$, where $f\wedge b_1=\alpha\wedge v$ and $\wedge b_2=\beta\wedge w$. By
  Lemma~\ref{cv-function-generator}(2), every frame function given in this way by a generator of
  $\mathcal L$ over $\bar Z$ is a multiple of $\varphi$ by an element of $F_c^\times$.
\item[(1c)] Let $\sigma$ be a complex embedding of $F_c$. Since $\sigma(\tau(\sqrt{-d}))^2=-d$ and
  $\bar\theta=-\theta$, either $\sigma(\tau(\sqrt{-d}))=\theta$, or the complex conjugate embedding
  $\bar\sigma$ has this property. In the first case, Lemma~\ref{cv-function-generator}(3) at $\sigma$
  gives $\kappa_\sigma\in\mathbf C^\times$ with $\varphi=\kappa_\sigma\Phi_\sigma$ at every complex
  framed point of $Z$ over $\sigma$, and by Lemma~\ref{cv-periods-function}(2),
  \[
   |\varphi(y,v,w)|\big(\det\Gamma(v)/\det\Delta(w)\big)^{1/2}=|\kappa_\sigma|
  \]
  at all of them, with $\Gamma(v)$ and $\Delta(w)$ the Hodge Gram matrices: the norm of $\varphi$ at
  $\sigma$ is the constant $e^{t_\sigma}=|\kappa_\sigma|$ on $Z$. In the second case, complex
  conjugation carries the complex framed points of $Z$ over $\bar\sigma$ to those over $\sigma$; it
  conjugates the values of $\varphi$, which is compatible with extensions of fields, and the Hodge
  Gram matrices, so the norm of $\varphi$ at $\sigma$ is the constant $e^{t_\sigma}$ with
  $t_\sigma=t_{\bar\sigma}$.
\item[(1d)] Let $\ell$ be a prime and $\kappa:F_c\to\mathbf C_\ell$ an embedding. The prime-to-$\ell$
  part of the level is at least $3$, so Theorem~\ref{cv-model-toroidal} gives a model of $X$ over
  $\{b\in F_c:|\kappa(b)|\le1\}$ with properties (1) and (2) there, and
  Lemma~\ref{cv-boundary-section} on it gives a generator $s_\ell$ of the corresponding line bundle
  on the closure of $j(Z)$ in its generic fibre. By Lemma~\ref{cv-norms-padic},
  $\log|j(\varphi_{s_\ell}(y,v,w))|$ is the same at all integral framed points $(y,v,w)$ of $Z$ at
  embeddings $j$ that extend $\kappa$, and by Lemma~\ref{cv-function-generator}(4),
  $\varphi=x\varphi_{s_\ell}$ with $x\in F_c^\times$. So the $\ell$-adic norm of $\varphi$ is constant
  on these integral framed points. Two embeddings $\kappa$ and $\kappa'$ that induce the same finite
  place $w$ of $F_c$ satisfy $\kappa'=g\circ\kappa$ for an isometric ring homomorphism
  $g:\mathbf C_\ell\to\mathbf C_\ell$, which carries the integral framed points at embeddings over
  $\kappa$ to those over $\kappa'$; so the constant depends only on $w$, and there are real numbers
  $\rho(w)$, one for each finite place $w$ of $F_c$, with that constant equal to
  $|\kappa(\pi_w)|^{\rho(w)}$ for a uniformizer $\pi_w$ of $w$.
\end{enumerate}
% Lean, Step 1 (GlobalTheory.CompanionTheory.Ar.Br.CovEx.exists_constSec):
% (1a) GlobalTheory.T.good_reduction_dvr_model_exists (Theorem cv-model-toroidal),
%   irreducible_component_clopen, geometric_component_alg_closed, eigen_lattice_vector_bundle
%   (Lemma cv-model-eigenbundles);
% (1b) GlobalTheory.CompanionTheory.Ar.Flat.A0.flat_zero_of_resolution_gaga: Phi is
%   exists_unit_analytic_function (Lemma cv-periods-function), f and kappa_0 come from
%   eigen_frames_analytic (Lemma cv-boundary-section), phi is wedge_eigen_line_generator and the
%   uniqueness is frame_functions_flat_agree (Lemma cv-function-generator);
% (1c) GlobalTheory.T.flat_first_from_resolution_and_gaga: exists_infNorm_of_rt,
%   frame_function_analytic_norm (Lemma cv-function-generator(3)) and inf_norm_conjugate_invariant;
% (1d) frame_function_finite_norm = FinGood.finGood: good_reduction_dvr_model_exists at p,
%   finite_constant_all_holds (Lemma cv-norms-padic), frame_function_two_models_relay (Lemma
%   cv-function-generator(4)); finite_norm_local_bound = FinPlace.exists_rho_at with
%   padic_embeddings_conjugate_of_units and hasFinNorm_of_isometry. The reading of phi as a rational
%   section is flat_on_model_exists_rul_sections, line_fn_reads_local_height_sum and
%   unf_section_norm_bound_at_height.
So $\varphi$ defines a rational section of $\omega^{\lambda_\tau}$ on $Z$, with an inverse section of
$\omega^{-\lambda_\tau}$, whose norm in the uniformization at $\sigma$ is the constant $e^{t_\sigma}$ on
$Z$ and whose valuation at each finite place $w$ is the constant $\rho(w)$.

Step 2: the degree at a CM point. Put $\theta=\tau(\sqrt{-d})$, let $\bar\tau$ be the complex
conjugate of $\tau$, and let $x=\gamma+\log(2\pi/\sqrt d)-\frac{L'}{L}(1,\chi)$. In the notation of
Definition~\ref{cv-cm-point-periods}, Step 1 gives the following.
\begin{enumerate}
\item[(i)] At $\sigma_0$ the component $Z$ is uniformized by the periods of a lattice $L$ of a type
  of signature $(p,q)$ for $\theta$, as in Lemma~\ref{cv-cm-point}(2), and the point $y_P$ of every
  period $P$ lies in $Z$.
\item[(ii)] Step (1b), written in terms of the periods, gives a family
  $e=(e_1,\dots,e_{p+q})$ in $L$, linearly independent over $\mathcal O_E$, and a number
  $\mu\in\mathbf C$ with the following property. Let $P$ be a period, let $v$ and $w$ be bases of
  $\operatorname{Lie}_\tau(y_P)$ and $\operatorname{Lie}_{\bar\tau}(y_P)$ with $v_a=d\exp_P(V_a)$ and
  $w_i=d\exp_P(W_i)$, let $R$ be a $p\times(p+q)$ matrix with $\sum_mR_{am}P(e_m)_\tau=V_a$ for
  every $a$, let $B$ be the $q\times(p+q)$ matrix with $P(e_n)_{\bar\tau}=\sum_iB_{in}W_i$ for every
  $n$, and let $S=B\,\tau(G(e))^{-1}$, where $\tau$ is applied to each entry. Then
  \[
   \varphi(y_P,v,w)\det\begin{pmatrix}R\\ S\end{pmatrix}=\mu .
  \]
\item[(iii)] For every complex embedding $\sigma$ of $F_c$, by (1c),
  \[
   \log|\varphi(y,v,w)|+\tfrac12\big(\log|\det\Gamma(v)|-\log|\det\Delta(w)|\big)=t_\sigma
  \]
  at every complex framed point $(y,v,w)$ of $Z$ over $\sigma$, where $\Gamma(v)$ and $\Delta(w)$
  are the Hodge Gram matrices of the frames.
\item[(iv)] For every prime $\ell$ and every embedding $\kappa:F_c\to\mathbf C_\ell$, by (1d),
  there is $\nu(\kappa)\in\mathbf R$ with $\log|j(\varphi(y,v,w))|=\nu(\kappa)$ at every
  \emph{integral framed point at $j$}: a framed point $(y,v,w)$ of $Z$ over a field $K$ with an
  embedding $j:K\to\mathbf C_\ell$ restricting to $\kappa$ on $F_c$, such that the abelian variety at
  $y$, with its $\mathcal O_E$-action and polarization, extends to a semi-abelian scheme over the
  valuation ring $V_j=\{b\in K:|j(b)|\le1\}$, and $v$ and $w$ are bases of the $\tau$- and
  $\bar\tau$-eigenspaces of its Lie algebra. Summed over the embeddings $\kappa$ above a finite place
  $w$ of $F_c$, the $\nu(\kappa)$ give $-\rho(w)\log Nw$.
\end{enumerate}
% Lean: (i) and (ii) are IsFlatGen U tau phi FU i (mem from family_uniformization_type_relay; e and
% mu from IsConstMul, the field flat of FlatLineGeneratorComponent); (iii) and (iv) are HasInfNorm
% and HasFinNorm. The degree is LogNorms.deg, minus the mean of the logarithms.
So $\deg\varphi=-\frac1{[F_c:\mathbf Q]}\big(\sum_\sigma t_\sigma+\sum_\ell\sum_\kappa\nu(\kappa)\big)$.
Let $N_\infty(\sigma)=t_\sigma$, and for $\kappa:F_c\to\mathbf C_\ell$ let
\[
 N(\kappa)=\nu(\kappa)-q\big(\log|\mathfrak d_E|_{\kappa\tau}-\log|d|_\ell\big),
\]
which removes the logarithm of the local defect $\delta_\tau$ (Lemma~\ref{cv-cm-point-gram} for the
notation; $\kappa\tau$ is the composite $E\to F_c\to\mathbf C_\ell$).

Let $Z_+$ and $Z_-$ be CM curves with roots $\theta$ and $-\theta$ given by
Lemma~\ref{cv-cm-point}(1), over number fields $M_\pm$, with numbers $z_\pm$ and squared Hodge
norms $n_\pm$, so that $-\frac1{[M_\pm:\mathbf Q]}\sum_s\log n_\pm(s)=x$. For a prime $\ell$ and a
number field $k_0$, a \emph{CM point at $\ell$ over $k_0$} is the data $P$, $u$, $K$, $y$, $v$, $w$
of Lemma~\ref{cv-cm-point}(2) for $Z_\pm$, $\ell$ and $k_0$, together with the matrix $T$ of
Lemma~\ref{cv-cm-point-gram}(3) for $u$ and $e$; then $\det T\ne0$ and
$\det G(u)=\det G_+(u)\det G_-(u)\ne0$ by Lemma~\ref{cv-cm-point-gram}(2), (3) and (5).
\begin{enumerate}
\item[(2a)] The value at a CM frame. At every CM point,
  \[
   \varphi(y,v,w)\,\tau(\det T)\,z_+^p=\mu\,z_-^q\,\tau(\det G_-(u)).\tag{$*$}
  \]
  Indeed, under $K\subset\mathbf C$ the frame $(y,v,w)$ becomes $(y_P,v,w)$, and $\varphi$ is
  compatible with this extension of fields. By Lemma~\ref{cv-cm-point}(2)(a), $V_a=z_+P(u_a)_\tau$
  and $W_i=z_-P(u_{p+i})_{\bar\tau}$. By Lemma~\ref{cv-cm-point-gram}(3), $u_m=\sum_nT_{mn}e_n$ in
  $L\otimes\mathbf Q$, so $P(u_m)_\sigma=\sum_n\sigma(T_{mn})P(e_n)_\sigma$ for
  $\sigma=\tau,\bar\tau$. Hence $R=(z_+I_p\ \ 0)\,\tau(T)$ satisfies the condition of (ii). For
  $a\le p$, $P(u_a)_{\bar\tau}=0$, because $u_a\in L_+(P)$ and $\bar\tau(\sqrt{-d})=-\theta$; and
  $P(u_{p+i})_{\bar\tau}=z_-^{-1}W_i$. So $B\,\bar\tau(T)^{\mathsf t}=(0\ \ z_-^{-1}I_q)$. By
  Lemma~\ref{cv-cm-point-gram}(3) and (5),
  $\tau(T)\,\tau(G(e))\,\bar\tau(T)^{\mathsf t}=\tau(G(u))$ is block diagonal with blocks
  $\tau(G_+(u))$ and $\tau(G_-(u))$, so
  $S\,\tau(T)^{-1}\tau(G(u))=B\,\bar\tau(T)^{\mathsf t}$ gives
  $S=(0\ \ z_-^{-1}\tau(G_-(u))^{-1})\,\tau(T)$. Therefore
  \[
   \det\begin{pmatrix}R\\ S\end{pmatrix}
   =\det\begin{pmatrix}z_+I_p&0\\0&z_-^{-1}\tau(G_-(u))^{-1}\end{pmatrix}\tau(\det T)
   =\frac{z_+^p\,\tau(\det T)}{z_-^q\,\tau(\det G_-(u))},
  \]
  and (ii) gives $(*)$.

  Take a CM point at $2$ over $\mathbf Q$, and let $k$ be its field $K$, which contains $F_c$,
  $M_+$ and $M_-$. Put $\varrho=\varphi(y,v,w)\tau(\det T)/\tau(\det G_-(u))$, a nonzero element of
  $k$, since $\varphi$ does not vanish on framed points. By $(*)$, $\varrho z_+^p=\mu z_-^q$. Since
  $z_+\ne0$, $\varrho$ is determined by $\mu$ and $z_\pm$, and $(*)$ says that at every CM point
  \[
   \varphi(y,v,w)\,\tau(\det T)=\varrho\,\tau(\det G_-(u)).\tag{$**$}
  \]
\item[(2b)] The archimedean places. Let $s:k\to\mathbf C$ and $\sigma=s|_{F_c}$. Take a CM point at
  $2$ over $k$ and extend $s$ to $s':K\to\mathbf C$. By (i), $y$ lies in $Z$, which is open and
  closed, so $(s'(y),s'(v),s'(w))$ is a complex framed point of $Z$ over $\sigma$, and (iii) applies
  to it, with $\log|s'(\varphi(y,v,w))|$ in the first term. Its Hodge Gram matrices are given by
  Lemma~\ref{cv-cm-point}(2)(b). Since $|\varepsilon_\pm|=1$, $|s'\tau(\sqrt{-d})|=\sqrt d$, and
  $|s'(\bar\tau(b))|=|s'(\tau(b))|$ for $b\in E$ ($s'\tau$ is an embedding of the CM field $E$),
  \[
   \log|\det\Gamma|=p\Big(\log n_+(s)+\frac{\log d}2\Big)+\log|s'\tau(\det G_+(u))|,\quad
   \log|\det\Delta|=q\Big(\log n_-(s)+\frac{\log d}2\Big)+\log|s'\tau(\det G_-(u))|.
  \]
  By $(**)$,
  $\log|s'(\varphi(y,v,w))|=\log|s(\varrho)|+\log|s'\tau(\det G_-(u))|-\log|s'\tau(\det T)|$. Put
  $a=\det G(e)\in E^\times$. By Lemma~\ref{cv-cm-point-gram}(3) and (5),
  $\det G_+(u)\det G_-(u)=\det T\cdot a\cdot\overline{\det T}$, so
  $\log|s'\tau(\det G_+(u))|+\log|s'\tau(\det G_-(u))|=2\log|s'\tau(\det T)|+\log|s\tau(a)|$.
  Substituting these into (iii),
  \[
   N_\infty(\sigma)=\log|s(\varrho)|+\frac p2\Big(\log n_+(s)+\frac{\log d}2\Big)
   -\frac q2\Big(\log n_-(s)+\frac{\log d}2\Big)+\frac12\log|s\tau(a)|,
  \]
  where $n_\pm(s)$ is taken at $s|_{M_\pm}$.
\item[(2c)] The finite places. Let $\ell$ be a prime, take a CM point at $\ell$ over $k$, and put
  $t_\ell=\det T\in E^\times$. By Lemma~\ref{cv-cm-point-gram}(3) and (4),
  \[
   \log|\iota(t_\ell)|+\log|\iota(\bar t_\ell)|+\log|\iota(a)|=\log|\iota(\det G(u))|
   =-(p+q)\log|\mathfrak d_E|_\iota
  \]
  for every embedding $\iota:E\to\mathbf C_\ell$. If $e$ is a local basis at $\ell$, then
  $|\iota(t_\ell)|=1$ for all $\iota$ by Lemma~\ref{cv-cm-point-gram}(4); by
  Lemma~\ref{cv-cm-point-gram}(2) this holds for all but finitely many $\ell$. Now let
  $j:k\to\mathbf C_\ell$, let $\kappa=j|_{F_c}$, and extend $j$ to $j':K\to\mathbf C_\ell$. By
  Lemma~\ref{cv-cm-point}(2)(c), $(y,\delta v,\delta'w)$ is an integral framed point at $j'$, so
  $\nu(\kappa)=\log|j'(\varphi(y,\delta v,\delta'w))|$. The transformation rule of $\varphi$ with
  $g=\delta I_p$ and $h=\delta'I_q$ gives
  $\varphi(y,\delta v,\delta'w)\,\delta^p=\varphi(y,v,w)\,\delta'^q$, so
  \[
   \nu(\kappa)=\log|j'(\varphi(y,v,w))|+(q-p)\big(\log|\mathfrak d_E|_{j\tau}-\log|j\tau(\sqrt{-d})|\big).
  \]
  By $(**)$ and Lemma~\ref{cv-cm-point-gram}(5),
  $\log|j'(\varphi(y,v,w))|=\log|j(\varrho)|-\log|j\tau(t_\ell)|-q\log|\mathfrak d_E|_{j\tau}$.
  Hence
  \[
   \nu(\kappa)=\log|j(\varrho)|-\log|j\tau(t_\ell)|-q\log|\mathfrak d_E|_{j\tau}
   +(q-p)\big(\log|\mathfrak d_E|_{j\tau}-\log|j\tau(\sqrt{-d})|\big).
  \]
  For $\ell\nmid d\,d_E$ at which $\varrho$ and $t_\ell$ are units, every term of $N(\kappa)$
  vanishes. So $N(\kappa)=0$ for all but finitely many $\ell$; in particular $\rho(w)=0$ for almost
  all $w$.
\item[(2d)] The sum. Extend $N_\infty$ and $N$ to the embeddings of $k$ by restriction to $F_c$.
  Each embedding of $F_c$ has $[k:F_c]$ extensions, so
  $\frac1{[k:\mathbf Q]}\big(\sum_sN_\infty(s)+\sum_\ell\sum_jN(j)\big)$ is the same sum for $F_c$.
  Substitute (2b) and (2c), with
  $N(j)=\nu(j|_{F_c})-q(\log|\mathfrak d_E|_{j\tau}-\log|d|_\ell)$, and divide by $[k:\mathbf Q]$:
  the terms in $\varrho$ cancel by the product formula; the composites $j\tau$ run over the
  embeddings $\iota$ of $E$, each $[k:E]$ times, and $\iota\mapsto\iota\circ(b\mapsto\bar b)$
  permutes them, so by (2c) and the product formula for $a$ the terms in $t_\ell$ and $a$ add up to
  $-\frac{p+q}2\log\operatorname{rd}_E$, since the norm of $\mathfrak d_E$ is $|d_E|$; the remaining
  terms in $\mathfrak d_E$ give $(p+q)\log\operatorname{rd}_E$; the finite terms in $d$ give
  $-\frac{p+q}2\log d$; and the archimedean terms in $n_\pm$ and $d$ give
  $-\frac{p-q}2x+\frac{p-q}4\log d$, because the mean of $\log n_\pm$ over the embeddings of $k$ is
  $-x$. The degree is minus the total:
  \[
   \frac{p-q}2\Big(x-\log\operatorname{rd}_E+\frac12\log d\Big)+q\big(\log d-\log\operatorname{rd}_E\big)
   =\frac{p-q}2\log\Omega_E(\chi)+q\big(\log d-\log\operatorname{rd}_E\big).
  \]
\end{enumerate}
% Lean: Step 2 is Ar.Fl.flatDegree. (2a) is frame_function_cm_value, with cm_pack_pure_frame_open
% (CMValue.value_on_cmFrame, CMValue.det_rows) and transition_matrix_period_frame
% (constMul_change_generator); the CM point at 2 is exists_cm_period_embedding with k0 = tau(E).
% (2b) is frame_function_cm_number_field (InfPlace.hasInfNorm_value, log_norm_det_scaled).
% (2c) is frame_function_cm_height (flatGen_finNorm), with gram_det_ne_zero_local_basis_cofinite
% for the finiteness. (2d) is hermitian_metric_log_norm_formula (deg_eq_of_local_two_curves),
% flatDegree_mean_ext and logNorms_ext_deg_isAdelic; the periods are chowla_selberg_period.
Removing the shifts, the unshifted logarithms of $\varphi$ have degree
$\frac{p-q}2\log\Omega_E(\chi)$, by the displayed sum of the defect.

Step 3: existence. By Steps 1 and 2, $\varphi$ is a constant section of weight $\lambda_\tau$
and defect $0$ with $\deg\varphi=\frac{p-q}2\log\Omega_E(\chi)$. The same section, with
$\rho_\varphi$ replaced by $\rho_\varphi-\delta_\tau$, is a flat section $t$ of type $\tau$. By the
definition of the degree and the displayed sum of the defect,
$\deg t=\deg\varphi-q(\log\operatorname{rd}_E-\log d)$, which is the stated value.

Step 4: every flat section. If $t'$ is another flat section of type $\tau$ on $Z$, then
$\deg t'=\deg t$ by Lemma~\ref{cv-constant-sections-proportional}.
\end{proof}
\leanref{Development/AutomorphicSpectrum/Proofs/AdelicProofs/HermitianMetricLogNormFormulaStatement.lean}{GlobalTheory.T.hermitian_metric_log_norm_formula}
\leanref{Development/AutomorphicSpectrum/Proofs/GeometryProofs/FrameFunctionCmValueStatement.lean}{GlobalTheory.T.frame_function_cm_value}
\leanref{Development/AutomorphicSpectrum/Proofs/GeometryProofs/CmPackPureFrameOpenStatement.lean}{GlobalTheory.T.cm_pack_pure_frame_open}
\leanref{Development/AutomorphicSpectrum/Proofs/GeometryProofs/TransitionMatrixPeriodFrameStatement.lean}{GlobalTheory.T.transition_matrix_period_frame}
\leanref{Development/AutomorphicSpectrum/Proofs/GeometryProofs/FrameFunctionCmNumberFieldStatement.lean}{GlobalTheory.T.frame_function_cm_number_field}
\leanref{Development/AutomorphicSpectrum/Proofs/GeometryProofs/FrameFunctionCmHeightStatement.lean}{GlobalTheory.T.frame_function_cm_height}

\begin{remark}\label{cv-where-chowla-selberg-enters}
The two parts of Theorem~\ref{cv-chowla-selberg} enter on the two sides of the comparison.
Part~(1) enters the covolume bound through Theorem~\ref{cv-flat-degree}: the main term
of Theorem~\ref{cv-covolume-bound} is a multiple of $\log\Omega_E(\chi)$ because flat
sections have the degree computed there. Part~(2) enters on the definite side: it bounds
the archimedean size of the algebraic numbers obtained by dividing the special values of
Hecke $L$-functions of the host field by powers of the Chowla--Selberg period. Those
values are the Petersson norms of theta lifts given by the Rallis inner product formula
(Theorem~\ref{th-rallis-inner-product}), and the bound is used in
Chapter~\ref{chapter-integral-structure}.
\end{remark}
\leanref{Development/AbelianVarietyTheory/Proofs/DivisorProofs/BetaConjugateMeanLogAbsBoundOfTransport.lean}{GlobalTheory.SqTop.betaArch_of_transport}
\leanref{Development/AutomorphicSpectrum/Proofs/GeometryProofs/ChowlaSelbergPeriodStatement.lean}{GlobalTheory.T.chowla_selberg_period}

\section{Arithmetic volumes}\label{section-cv-volumes}

\subsection{Weights and flag heights}\label{subsection-cv-weights-flag-heights}

\begin{theorem}\label{cv-functorial-resolution}
For every field $k$ of characteristic zero and every reduced scheme $X$ that is locally of
finite type and quasi-compact over $k$, there is a finite sequence of blow-ups
$X_r\to\cdots\to X_0=X$, each along a closed subscheme, such that:
\begin{enumerate}
\item $X_r$ is smooth over $k$ and $X_r\to X$ is proper;
\item $X_r\to X$ is an isomorphism over the smooth locus $X_{\mathrm{sm}}$ of $X$, and the
  inverse image of $X_{\mathrm{sm}}$ is dense in $X_r$;
\item every centre maps into $X\setminus X_{\mathrm{sm}}$;
\item the reduced inverse image of $X\setminus X_{\mathrm{sm}}$ is a divisor with strict
  normal crossings on $X_r$.
\end{enumerate}
The sequence can be chosen for all such $k$ and $X$ at once so that it commutes with smooth
morphisms $X'\to X$ over $k$, after omitting the blow-ups whose pulled-back centres are
empty, and with extensions $k\subset k'$ of fields of characteristic zero for which the
base change of $X$ is reduced.
\end{theorem}
\leanref{Development/ResolutionOfSingularities/HasSimpleNormalCrossingsOfInStrictNormalCrossing.lean}{FunctorialResolution.functorialResolutionCharZero}
\leanref{Development/ResolutionOfSingularities/Statements.lean}{FunctorialResolution.FunctorialResolutionCharZero}

This is Hironaka's theorem \cite{Hir64} in the functorial form of
\cite[Theorem 3.36]{Kol07}. For integral separated schemes of finite type, the resolution by
blow-ups with centres over the singular locus that is an isomorphism over the smooth locus is
proved in \cite[Tag~R02G]{Res26}, and the compatibility of principalization with smooth
morphisms in \cite[Tag~R022]{Res26}.

\begin{theorem}\label{cv-moriwaki-bigness}
Let $\mathcal X$ be an arithmetic variety of dimension $n+1$, given as the closure of a
point of a product of projective spaces over $\mathbf Z$ whose function field it
generates. Let $\overline A$ be a Hermitian line bundle on $\mathcal X$ with smooth,
conjugation-invariant metric that is semipositive, nef on the fibres of positive
characteristic, and of nonnegative degree on every curve of the generic fibre. Let
$\widehat{\deg}$ be an arithmetic intersection theory on such varieties. Then
\[
 \widehat{\deg}\big(\hat c_1(\overline A)^{n+1}\big)\le\widehat{\operatorname{vol}}(\overline A),
 \qquad
 \widehat{\operatorname{vol}}(\overline A)=\limsup_{\ell\to\infty}
 \frac{\log\#\{s\in H^0(\mathcal X,A^{\ell}):\|s\|_{\sup}\le1\}}{\ell^{n+1}/(n+1)!}.
\]
\end{theorem}
\leanref{Development/WeilBound/CoreStatements/AutomorphicKernels/MoriwakiBigness/MoriwakiBigTheorem/Outright.lean}{PositivityTheory.Pf1.HilbertSamuelMoriwaki.moriwakiBig_outright}
\leanref{ReciprocityTheory/ReductiveGroups/ArithmeticGeometry/ArakelovGeometry/MoriwakiSt.lean}{ReciprocityTheory.ChartRestriction.G.ArakelovGeometry.MoriwakiBigStatement}

For an ample line bundle with positive metric this is the arithmetic Hilbert--Samuel
formula, and for nef bundles it follows from the continuity of the arithmetic volume.
See \cite{Yua08} and \cite{Mor09}.

\begin{definition}\label{cv-reconstruction-functionals}
Let $\Psi$ be a finite nonempty set and let $j<N$ be positive integers. Let
$\mathbf R[x]=\mathbf R[x_{\tau,i}:\tau\in\Psi,\ 1\le i\le N]$ and $d=|\Psi|N(N-1)/2$.
\begin{enumerate}
\item A \emph{weight} is a family $\lambda=(\lambda_{\tau,i})$ in $(\mathbf R^N)^\Psi$. Put
  $\|\lambda\|=\max_{\tau,i}|\lambda_{\tau,i}|$, $\|\lambda\|_1=\sum_{\tau,i}|\lambda_{\tau,i}|$,
  $\operatorname{spread}(\lambda)=\sum_\tau(\max_i\lambda_{\tau,i}-\min_i\lambda_{\tau,i})$,
  $\operatorname{gap}_j(\lambda,\tau)=\lambda_{\tau,j}-\lambda_{\tau,j+1}$ and
  $\ell(\lambda)=-\sum_{\tau,i}\lambda_{\tau,i}x_{\tau,i}$. The weight is \emph{dominant} if
  $\lambda_{\tau,1}\ge\dots\ge\lambda_{\tau,N}$ for every $\tau$, and \emph{regular dominant} if these
  inequalities are strict. A regular dominant weight $u$ has \emph{gaps at least} $g>0$ if
  $\operatorname{gap}_j(u,\tau)\ge g$ for every $\tau$, and $|u_{\tau,a}-u_{\tau,b}|\ge g$ for every
  $\tau$, every $a$ in the \emph{small block} and every $b\ne a$; the small block is
  $\{1,\dots,j\}$ if $j\le N-j$ and $\{j+1,\dots,N\}$ otherwise.
\item The \emph{degree on the flag variety} is the linear form $\deg_{\mathrm{fl}}$ on $\mathbf R[x]$
  that sends $F$ to $(-1)^d\sum_\sigma\prod_\tau\operatorname{sgn}(\sigma_\tau)\,c_\sigma(F)$, where
  $\sigma=(\sigma_\tau)_{\tau\in\Psi}$ runs over the families of permutations of $\{1,\dots,N\}$ and
  $c_\sigma(F)$ is the coefficient in $F$ of $\prod_{\tau,i}x_{\tau,i}^{N-\sigma_\tau(i)}$.
\item For $\tau\in\Psi$ and $k\ge0$ let $p_{\tau,k}=\sum_{i=1}^Nx_{\tau,i}^k$ and
  $p^{\le j}_{\tau,k}=\sum_{i\le j}x_{\tau,i}^k$, and for $F\in\mathbf R[x]$ put
  $\mathrm{Sh}_\tau(F)=\deg_{\mathrm{fl}}\big(\sum_{i\le j}\partial F/\partial x_{\tau,i}\big)$ and
  $\mathrm{Tot}_\tau(F)=\deg_{\mathrm{fl}}\big(\sum_{i=1}^N\partial F/\partial x_{\tau,i}\big)$.
\item Let $\Delta=\prod_\tau\prod_{a<b}(x_{\tau,a}-x_{\tau,b})$. A \emph{block transposition} exchanges
  $x_{\tau,a}$ and $x_{\tau,a+1}$ for some $\tau$ and some $a\ne j$. A linear form is \emph{block
  invariant} if it is invariant under every permutation of the variables that fixes the index
  $\tau$ of each variable and preserves the blocks $\{i\le j\}$ and $\{i>j\}$.
\item Let $c_0>0$ and $\tau_0\in\Psi$. A weight $\lambda$ is \emph{on the wall at $\tau_0$ with margin
  $c_0$} if $\|\lambda\|\le\frac12$, $\operatorname{gap}_j(\lambda,\tau_0)=0$,
  $\operatorname{gap}_j(\lambda,\tau)\ge c_0$ for $\tau\ne\tau_0$, and
  $\lambda_{\tau,i}-\lambda_{\tau,i+1}\ge c_0$ for every $\tau$ and every $i<N$ with $i\ne j$.
\item Let $b\in\mathbf R$. A weight $\mu$ is \emph{beyond the threshold at $b$} if
  $4|\Psi|\,\|\mu\|\le(b-2)\operatorname{gap}_j(\mu,\tau)$ for every $\tau\in\Psi$.
\end{enumerate}
\end{definition}
\leanref{Development/AutomorphicSpectrum/Definitions/DivisibilityTowers/CompactGapsGE.lean}{GlobalTheory.CompanionTheory.Ar.Top.linRoot}
\leanref{Development/AutomorphicSpectrum/Definitions/DivisibilityTowers/CompactGapsGE.lean}{GlobalTheory.CompanionTheory.Ar.Top.gap}
\leanref{Development/AutomorphicSpectrum/Definitions/DivisibilityTowers/CompactGapsGE.lean}{GlobalTheory.CompanionTheory.Ar.Top.Dominant}
\leanref{Development/AutomorphicSpectrum/Definitions/DivisibilityTowers/CompactGapsGE.lean}{GlobalTheory.CompanionTheory.Ar.Top.RegDominant}
\leanref{Development/AutomorphicSpectrum/Definitions/DivisibilityTowers/CompactGapsGE.lean}{GlobalTheory.CompanionTheory.Ar.Top.WeightHyp}
\leanref{Development/AutomorphicSpectrum/Definitions/DivisibilityTowers/CompactGapsGE.lean}{GlobalTheory.CompanionTheory.Ar.Top.spread}
\leanref{Development/AutomorphicSpectrum/Definitions/DivisibilityTowers/CompactGapsGE.lean}{GlobalTheory.CompanionTheory.Ar.Top.degFl}
\leanref{Development/AutomorphicSpectrum/Definitions/DivisibilityTowers/CompactGapsGE.lean}{GlobalTheory.CompanionTheory.Ar.Top.pk}
\leanref{Development/AutomorphicSpectrum/Definitions/DivisibilityTowers/CompactGapsGE.lean}{GlobalTheory.CompanionTheory.Ar.Top.pkLow}
\leanref{Development/AutomorphicSpectrum/Definitions/DivisibilityTowers/CompactGapsGE.lean}{GlobalTheory.CompanionTheory.Ar.Top.Sh}
\leanref{Development/AutomorphicSpectrum/Definitions/DivisibilityTowers/CompactGapsGE.lean}{GlobalTheory.CompanionTheory.Ar.Top.Tot}
\leanref{GroundTheory/HeightsAndTheta/Heights/Antisymmetric/Statements.lean}{GroundTheory.Heights₂.Heights.antisymDelta}
\leanref{Development/AutomorphicSpectrum/Definitions/DivisibilityTowers/HeckeSectionsIsotropicCommutatorPairingPower.lean}{GlobalTheory.CompanionTheory.Ar.Top.InWall}
\leanref{Development/AutomorphicSpectrum/Definitions/DivisibilityTowers/CompactGapsGE.lean}{GlobalTheory.CompanionTheory.Ar.Top.NefThreshold}

\begin{lemma}\label{cv-reconstruction-flag-degree}
In the notation of Definition~\ref{cv-reconstruction-functionals}:
\begin{enumerate}
\item for every weight $\lambda$,
  \[
   \deg_{\mathrm{fl}}(\ell(\lambda)^d)=d!\prod_\tau\prod_{a<b}\frac{\lambda_{\tau,a}-\lambda_{\tau,b}}{b-a},
  \]
  a positive multiple of $\Delta(\lambda)$ that is positive if $\lambda$ is regular dominant and
  nonnegative if $\lambda$ is dominant; and there is $C\ge0$ with
  $\deg_{\mathrm{fl}}(\ell(\lambda)^d)\le C\operatorname{gap}_j(\lambda,\tau_0)$ for every $\tau_0\in\Psi$
  and every dominant weight $\lambda$ with $\|\lambda\|\le1$;
\item for every weight $\lambda$ and every $\tau\in\Psi$,
  \[
   \mathrm{Sh}_\tau(\ell(\lambda)^{d+1})=-(d+1)\deg_{\mathrm{fl}}(\ell(\lambda)^d)\sum_{i\le j}\lambda_{\tau,i},
   \qquad
   \mathrm{Tot}_\tau(\ell(\lambda)^{d+1})=-(d+1)\deg_{\mathrm{fl}}(\ell(\lambda)^d)\sum_{i=1}^N\lambda_{\tau,i};
  \]
\item for $\sigma,\tau\in\Psi$, $k\ge1$ and $Q\in\mathbf R[x]$, $\mathrm{Sh}_\sigma(p_{\tau,k}Q)$ is
  $k\deg_{\mathrm{fl}}(p^{\le j}_{\tau,k-1}Q)$ if $\sigma=\tau$ and $0$ otherwise, and
  $\mathrm{Tot}_\sigma(p_{\tau,k}Q)$ is $N\deg_{\mathrm{fl}}(Q)$ if $\sigma=\tau$ and $k=1$, and $0$
  otherwise;
\item for a permutation $\pi$ of $\Psi$ and $F\in\mathbf R[x]$, let $F^\pi$ be the image of $F$ under
  $x_{\tau,i}\mapsto x_{\pi(\tau),i}$; then $\deg_{\mathrm{fl}}(F^\pi)=\deg_{\mathrm{fl}}(F)$,
  $\mathrm{Sh}_{\pi(\tau)}(F^\pi)=\mathrm{Sh}_\tau(F)$ and $\mathrm{Tot}_{\pi(\tau)}(F^\pi)=\mathrm{Tot}_\tau(F)$.
\end{enumerate}
\end{lemma}
\leanref{Development/AutomorphicSpectrum/Proofs/DivisibilityProofs/FlagDegreeWeightFormula.lean}{GlobalTheory.CompanionTheory.Ar.Top.flagDeg_eq_weyl}
\leanref{Development/AutomorphicSpectrum/Proofs/DivisibilityProofs/AntisymmetricPerturbationBound.lean}{GlobalTheory.CompanionTheory.Ar.Top.flagDeg_eq_delta}
\leanref{Development/AutomorphicSpectrum/Proofs/DivisibilityProofs/FlagDegreeLeGapStatement.lean}{GlobalTheory.T.flag_degree_le_gap}
\leanref{Development/AutomorphicSpectrum/Proofs/DivisibilityProofs/FlagShiftTotalLinRootPowerStatement.lean}{GlobalTheory.T.flag_shift_total_lin_root_power}
\leanref{Development/AutomorphicSpectrum/Proofs/DivisibilityProofs/FlagShiftTotalPowerSumStatement.lean}{GlobalTheory.T.flag_shift_total_power_sum}
\leanref{Development/AutomorphicSpectrum/Proofs/DivisibilityProofs/FlagDegreeSymmetriesStatement.lean}{GlobalTheory.T.flag_degree_symmetries}

These are self-contained computations. Part (1) is the leading term of Weyl's dimension formula,
and the bound by the gap comes from the factor $(a,b)=(j,j+1)$ at $\tau_0$, the other factors being
bounded on $\|\lambda\|\le1$. Part (2) follows from $\partial\ell(\lambda)/\partial x_{\tau,i}=-\lambda_{\tau,i}$.
For (3), $\deg_{\mathrm{fl}}$ is an alternating sum over the permutations of the variables of each
$\tau$, so it vanishes on $SQ$ when $S$ is symmetric in $x_{\tau,1},\dots,x_{\tau,N}$ of positive
degree; in the derivatives of $p_{\tau,k}Q$ only the terms $kx_{\tau,i}^{k-1}Q$ remain, and
$p_{\tau,0}=N$. Part (4) holds because relabelling $\Psi$ permutes the families $\sigma$.

\begin{lemma}\label{cv-reconstruction-polarization}
Let $J$ be a finite set, let $D\ge0$, and let $\Phi$ be a linear form on $\mathbf R[x_y:y\in J]$. Then
$P_\Phi(c)=\Phi\big((-\sum_yc_yx_y)^D\big)$ is a homogeneous polynomial of degree $D$ in
$c=(c_y)_{y\in J}$. Its coefficient at a monomial $c^m$ of degree $D$ is $(-1)^D\binom Dm\Phi(x^m)$,
where $\binom Dm$ is the multinomial coefficient, so that $|\Phi(x^m)|$ is at most the absolute value
of that coefficient.
\end{lemma}
\leanref{Development/AutomorphicSpectrum/Proofs/DivisibilityProofs/PolarizationPolynomialCoeffBoundStatement.lean}{GlobalTheory.T.polarization_polynomial_coeff_bound}

This is the multinomial theorem, applied to $(-\sum_yc_yx_y)^D$ before $\Phi$.

\begin{lemma}\label{cv-wall-lipschitz}
Let $\Psi$, $N$ and $\mathbf R[x]$ be as in Definition~\ref{cv-reconstruction-functionals}, let $a$ be
a further variable, and let $D\ge0$. There is $C\ge0$ with the following property. Let $\Phi$ be a
linear form on $\mathbf R[x,a]$ and $b\ge0$ such that $|\Phi(m)|\le b$ for every monomial $m$. Then
the function $(\lambda,s)\mapsto\Phi\big((\ell(\lambda)+sa)^D\big)$ on $(\mathbf R^N)^\Psi\times\mathbf R$ is
continuous, and
\[
 \big|\Phi\big((\ell(\lambda)+sa)^D\big)-\Phi\big((\ell(\mu)+s'a)^D\big)\big|\le Cb\,\big(\|\lambda-\mu\|+|s-s'|\big)
\]
for all weights $\lambda,\mu$ with $\|\lambda\|,\|\mu\|\le1$ and all real $s,s'$ with $|s|,|s'|\le1$.
\end{lemma}
\leanref{Development/AutomorphicSpectrum/Proofs/DivisibilityProofs/ArithFunctionalLipschitzWeightsStatement.lean}{GlobalTheory.T.arith_functional_lipschitz_weights}

By the multinomial theorem, $\Phi((\ell(\lambda)+sa)^D)$ is a polynomial in the entries of $\lambda$
and in $s$ whose coefficients are at most $b$ times multinomial coefficients. Its partial derivatives
are therefore bounded by $Cb$ on the box $\|\lambda\|\le1$, $|s|\le1$, with $C$ depending only on $D$,
$|\Psi|$ and $N$.

\begin{theorem}\label{cv-wall-kleiman}
Let $X$ be a projective variety of dimension $n$ over a field, and let $L_1,\dots,L_n$ be line
bundles on $X$ that are nef, that is, of nonnegative degree on every curve in $X$. Then the
intersection number $(L_1\cdots L_n)$ is nonnegative.
\end{theorem}
\leanref{ReciprocityTheory/ReductiveGroups/ArithmeticGeometry/ArakelovGeometry/KleimanNefness/KleimanSt.lean}{ReciprocityTheory.ChartRestriction.G.ArakelovGeometry.KH.KleimanNefFStatement}
\leanref{Development/AutomorphicSpectrum/Proofs/AlgebraProofs/KleimanNefCriteriaHoldStatement.lean}{GlobalTheory.T.kleiman_nef_criteria_hold}

This is Kleiman's theorem \cite[Chapter~III]{Kle66}: nef line bundles are limits of ample
$\mathbf Q$-line bundles, whose intersection numbers are positive.

\begin{theorem}\label{cv-wall-successive-minima}
Let $X$ be an arithmetic subvariety of dimension $D+1$ of a product of projective spaces over
$\mathbf Z$, let $(\cdot\mid\cdot)$ be an arithmetic intersection system, and let $\bar A$ be a
metrized line bundle that is semipositive on $X$, invariant under conjugation and nef on the fibres
of $X$ (Definition~\ref{lit-multiprojective}). For an arithmetic subvariety $W$ of dimension $1$ of
$X$, a \emph{point} of $X$, let $h(W)=(\bar A\mid W)/\deg W$, where $\deg W$ is the number of complex
points of $W$, and let $\deg_A(X)$ be the degree of $A$ on the generic fibre of $X$. Let
$e\in\mathbf R$.
\begin{enumerate}
\item If $h(W)\ge e$ for every point $W$ of $X$, then $(\bar A^{D+1}\mid X)\ge(D+1)\,e\deg_A(X)$.
\item If $h(W)\le e$ for every point $W$ in a Zariski-dense set of points of $X$, then
  $(\bar A^{D+1}\mid X)\le(D+1)\,e\deg_A(X)$.
\end{enumerate}
\end{theorem}
\leanref{ReciprocityTheory/ReductiveGroups/ArithmeticGeometry/BirationalGeometry/ArithmeticHilbertSamuelLowerBound.lean}{ReciprocityTheory.ChartRestriction.G.Geo.SuccessiveMinimaLowerBound}
\leanref{ReciprocityTheory/ReductiveGroups/ArithmeticGeometry/BirationalGeometry/ArithmeticHilbertSamuelLowerBound.lean}{ReciprocityTheory.ChartRestriction.G.Geo.SuccessiveMinimaUpperBound}
\leanref{GroundTheory/HeightsAndTheta/Heights/SemipositiveHeights/OfMoriwakiHS.lean}{GroundTheory.Heights₂.Heights.ForSecondHalf.Z1.of_moriwakiHS}
\leanref{GroundTheory/HeightsAndTheta/Heights/SemipositiveHeights/ImportsBis.lean}{GroundTheory.Heights₂.Heights.ForSecondHalf.Z2.proof}

\begin{proof}
For ample $A$ these are the two inequalities of Zhang's theorem on successive minima
\cite[Theorem~5.2]{Zha95}, which bound the height $(\bar A^{D+1}\mid X)/((D+1)\deg_A(X))$ of $X$ between
the absolute minimum and the essential minimum of the heights of its points. For the bundles here,
which are only nef on the fibres, Zhang's argument goes through. (1) follows from the arithmetic
Hilbert--Samuel formula (Theorem~\ref{lit-arithmetic-hilbert-samuel}) applied after scaling the metric
of $\bar A$ by $e^{e}$. (2) follows from the bigness theorem (Theorem~\ref{cv-moriwaki-bigness}): if
(2) failed, then for some $e'>e$ the bundle with its metric scaled by $e^{e'}$ would have positive
self-intersection number, hence a nonzero small section of some power, and the points outside the
divisor of that section would have height at least $e'$.
\end{proof}

\begin{lemma}\label{cv-wall-perturbed-weights}
In the notation of Definition~\ref{cv-reconstruction-functionals}, let $\nu$ be a weight and let
$c_0,\theta_0>0$. There is $B_1$ such that the following holds for every integer $B\ge B_1$, every
$\tau_0\in\Psi$ and every weight $\lambda_0$ on the wall at $\tau_0$ with margin $c_0$. Put
$\theta=5|\Psi|/B$, let $\lambda$ be the weight obtained from $\lambda_0$ by adding $2\theta$ to its
entries at $(\tau_0,i)$ with $i\le j$, and put $\lambda'=\lambda-\theta^2\nu$. Then $B\ge3$,
$0<\theta\le\min(\theta_0,1)$, $\|\lambda\|\le1$ and $\|\lambda'\|\le1$;
$\operatorname{gap}_j(\lambda,\tau)\ge\theta$ for every $\tau$ and
$\operatorname{gap}_j(\lambda,\tau_0)=2\theta$; $\lambda_{\tau,i}-\lambda_{\tau,i+1}\ge c_0$ for every
$\tau$ and every $i<N$ with $i\ne j$; $\lambda$ and $\lambda'$ are dominant; $\lambda'$ is beyond the
threshold at $B$; and $\|\lambda'-\lambda_0\|\le2\theta+\|\nu\|\theta^2$.
\end{lemma}
\leanref{Development/AutomorphicSpectrum/Proofs/DivisibilityProofs/PerturbedWeightGapBoundsStatement.lean}{GlobalTheory.T.perturbed_weight_gap_bounds}

This is a self-contained computation. Every entry moves by at most $2\theta+\|\nu\|\theta^2$; the gap
of $\lambda$ at $(\tau_0,j)$ is $2\theta$ and its other gaps are those of $\lambda_0$. For the
threshold, $(B-2)\big(2\theta-\theta^2|\operatorname{gap}_j(\nu,\tau_0)|\big)$ tends to $10|\Psi|$ as
$B\to\infty$, while $4|\Psi|\,\|\lambda'\|\le4|\Psi|(\frac12+2\theta+\|\nu\|\theta^2)$ tends to
$2|\Psi|$; at $\tau\ne\tau_0$ the gap of $\lambda'$ is at least $c_0-\theta^2|\operatorname{gap}_j(\nu,\tau)|$,
so the right side of the threshold condition tends to infinity.

\begin{lemma}\label{cv-wall-rational-weights}
In the notation of Definition~\ref{cv-reconstruction-functionals}, let $c_0>0$, $b\in\mathbf R$ and
$\tau_0\in\Psi$, and let $\Phi$ be a continuous function on $(\mathbf R^N)^\Psi$ with
$|\Phi(\lambda)|\le b$ for every weight $\lambda$ with rational entries on the wall at $\tau_0$ with
margin $c_0/2$. Then $|\Phi(\lambda)|\le b$ for every weight $\lambda$ on the wall at $\tau_0$ with
margin $c_0$.
\end{lemma}
\leanref{Development/AutomorphicSpectrum/Proofs/DivisibilityProofs/ContinuousWallBoundExtendStatement.lean}{GlobalTheory.T.continuous_wall_bound_extend}

Let $\lambda$ be on the wall at $\tau_0$ with margin $c_0$. Approximate $\lambda$ by rational weights
$\mu_n$ with $\|\mu_n-\lambda\|\le1/n$ and $\mu_{n,\tau_0,j}=\mu_{n,\tau_0,j+1}$, and put
$\lambda_n=(1-2/n)\mu_n$. Then $\|\lambda_n\|\le\frac12$, the gap of $\lambda_n$ at $(\tau_0,j)$ is $0$,
its other gaps are at least $(1-2/n)(c_0-2/n)\ge c_0/2$ for large $n$, and $\lambda_n\to\lambda$; so
$|\Phi(\lambda)|\le b$ by continuity.

\begin{lemma}\label{cv-reconstruction-antisymmetric}
In the notation of Definition~\ref{cv-reconstruction-functionals}, for $\tau\in\Psi$ let $H_\tau$ be
the hyperplane $x_{\tau,j}=x_{\tau,j+1}$ in $(\mathbf R^N)^\Psi$, and let $K_\tau\subset H_\tau$ be a
compact set that contains a nonempty open subset of $H_\tau$. Let $h\in\mathbf R[x]$ vanish on every
$K_\tau$. Then there is $C\ge0$ with the following property. Let $P\in\mathbf R[x]$ and
$v,\varepsilon\in\mathbf R$ be such that $P-vh$ is homogeneous of degree $d+1$ and changes sign under
every block transposition, and $|P|\le\varepsilon$ on every $K_\tau$. Then there is a block invariant
linear form $l$ such that every coefficient of $P-vh-\Delta l$ has absolute value at most
$C\varepsilon$.
\end{lemma}
\leanref{GroundTheory/HeightsAndTheta/Heights/Antisymmetric/Statements.lean}{GroundTheory.Heights₂.Heights.AntisymPerturbation.Statement}
\leanref{Development/AutomorphicSpectrum/Proofs/GeometryProofs/AntisymmetricPerturbationHoldsStatement.lean}{GlobalTheory.T.antisymmetric_perturbation_holds}
\begin{proof}
Let $A$ be the space of homogeneous polynomials of degree $d+1$ that change sign under every block
transposition, and let $R\in A$. A transposition of $x_{\tau,a}$ and $x_{\tau,b}$ with $a<b$ in one
block is a product of an odd number of block transpositions, so it changes the sign of $R$, and $R$
vanishes on the hyperplane $x_{\tau,a}=x_{\tau,b}$. Every permutation of the variables that fixes the
index $\tau$ of each variable maps $R$ to $R$ or $-R$.

Divisibility. Suppose that $R\in A$ vanishes on a nonempty open subset of each hyperplane
$x_{\tau,a}=x_{\tau,b}$ with $a\le j<b$. Then $R$ vanishes on these hyperplanes, and by the previous
paragraph on those with $a<b$ in one block. So each of the pairwise coprime linear forms
$x_{\tau,a}-x_{\tau,b}$ divides $R$, and $R=\Delta l$ with $l$ linear, as $\Delta$ has degree $d$. Both
$R$ and $\Delta$ change sign under every block transposition, so $l$ is block invariant.

Stability. For $\tau\in\Psi$ and $a\le j<b$ let $\pi_{\tau,a,b}$ be the permutation of the variables
that exchanges $x_{\tau,a}$ with $x_{\tau,j}$ and $x_{\tau,b}$ with $x_{\tau,j+1}$, and let
$K_{\tau,a,b}$ be the image of $K_\tau$ under the corresponding permutation of coordinates, a compact
subset of the hyperplane $x_{\tau,a}=x_{\tau,b}$ that contains a nonempty open subset of it. Let
$\Omega$ be the union of the $K_{\tau,a,b}$. Restriction to $\Omega$ is a linear map $\rho$ from the
finite-dimensional space $A$ to the continuous functions on $\Omega$ with the supremum norm. Choose a
linear map $g$ from the image of $\rho$ to $A$ with $\rho\circ g=1$. Its image is finite-dimensional,
so there is $C$ such that every coefficient of $g(f)$ is at most $C\sup_\Omega|f|$ in absolute value.
For $R\in A$, the polynomial $R-g(\rho(R))$ lies in $A$ and vanishes on $\Omega$, so it is $\Delta l$
for a block invariant $l$ by the divisibility step.

Now let $R=P-vh$. Since $\pi_{\tau,a,b}$ maps $R$ to $\pm R$, the supremum of $|R|$ on $K_{\tau,a,b}$ is
its supremum on $K_\tau$, where $R=P$ because $h$ vanishes there; this is at most $\varepsilon$. So
$R-\Delta l=g(\rho(R))$ has coefficients at most $C\varepsilon$ in absolute value.
\end{proof}
\leanref{GroundTheory/HeightsAndTheta/Heights/ElementarySymmetric/PAntisymDivisible.lean}{GroundTheory.Heights₂.Heights.AntisymDivisible.proof}
\leanref{GroundTheory/HeightsAndTheta/Heights/ElementarySymmetric/PAntisymStability.lean}{GroundTheory.Heights₂.Heights.AntisymStability.proof}
\leanref{GroundTheory/HeightsAndTheta/Heights/ElementarySymmetric/PAntisymPerturbation.lean}{GroundTheory.Heights₂.Heights.AntisymPerturbation.proof}

\begin{lemma}\label{cv-reconstruction-coefficients}
In the notation of Definition~\ref{cv-reconstruction-functionals}, let $\Phi$ be a linear form on
$\mathbf R[x]$, let $v>0$ and $\varepsilon\ge0$, and let $P$ be a homogeneous polynomial of degree
$d+1$ in the entries of a weight with $P(\lambda)=\Phi(\ell(\lambda)^{d+1})$ for every weight
$\lambda$. Let $l$ be a block invariant linear form such that every coefficient of $P-\Delta l$ has
absolute value at most $\varepsilon$. Then there are real numbers $\alpha_\tau,\beta_\tau$, for
$\tau\in\Psi$, with
\[
 \Big|\Phi(m)+v\sum_\tau\big((\alpha_\tau-\beta_\tau)\mathrm{Sh}_\tau(m)+\beta_\tau\mathrm{Tot}_\tau(m)\big)\Big|\le\varepsilon
\]
for every monomial $m$ of degree $d+1$.
\end{lemma}
\leanref{Development/AutomorphicSpectrum/Proofs/DivisibilityProofs/AntisymmetricPerturbationBoundStatement.lean}{GlobalTheory.T.antisymmetric_perturbation_bound}
\begin{proof}
A block invariant linear form is
$l=\sum_\tau\big(a_\tau\sum_{i\le j}x_{\tau,i}+b_\tau\sum_ix_{\tau,i}\big)$ for some real $a_\tau,b_\tau$.
By Lemma~\ref{cv-reconstruction-flag-degree}(1), $\deg_{\mathrm{fl}}(\ell(\lambda)^d)=c\Delta(\lambda)$
with $c>0$. Put $K=v(d+1)c$, $\alpha_\tau=(a_\tau+b_\tau)/K$ and $\beta_\tau=b_\tau/K$, and let
$\Phi'=\Phi+v\sum_\tau((\alpha_\tau-\beta_\tau)\mathrm{Sh}_\tau+\beta_\tau\mathrm{Tot}_\tau)$. By
Lemma~\ref{cv-reconstruction-flag-degree}(2), for every weight $\lambda$,
\[
 \Phi'(\ell(\lambda)^{d+1})=P(\lambda)-\Delta(\lambda)\sum_\tau\Big(a_\tau\sum_{i\le j}\lambda_{\tau,i}+b_\tau\sum_i\lambda_{\tau,i}\Big)=(P-\Delta l)(\lambda).
\]
So $P-\Delta l$ is the polynomial $P_{\Phi'}$ of Lemma~\ref{cv-reconstruction-polarization} with
$D=d+1$, and $|\Phi'(m)|$ is at most the absolute value of a coefficient of $P-\Delta l$.
\end{proof}

\begin{lemma}\label{cv-reconstruction-flag-height}
Let $\Psi$, $N\ge2$, $d$ and $\mathbf R[x]$ be as in Definition~\ref{cv-reconstruction-functionals},
and let $(\cdot\mid\cdot)$ be a full arithmetic intersection system
(Definition~\ref{cv-origin-intersection-full}). Let $\mathrm{Fl}$ be the product over $\Psi$ of the schemes over
$\mathbf Z$ of complete flags $0\subset V_1\subset\dots\subset V_{N-1}\subset\mathbf Z^N$, embedded in a
product of projective spaces by the Plücker embeddings of the $V_q$. Let $\bar{\mathcal O}_{\tau,q}$ be
the pullback of $\mathcal O(1)$ from the Plücker factor of $V_q$ at $\tau$, with the Fubini--Study
metric, and let $\bar L_{\tau,i}=\bar{\mathcal O}_{\tau,i}\otimes\bar{\mathcal O}_{\tau,i-1}^{-1}$, where
$\bar{\mathcal O}_{\tau,0}$ and $\bar{\mathcal O}_{\tau,N}$ are trivial. The \emph{flag height} is the
linear form $h$ on the homogeneous polynomials of degree $d+1$ in $\mathbf R[x]$ that sends
$x_{w_1}\cdots x_{w_{d+1}}$ to $(-1)^{d+1}(\bar L_{w_1}\cdots\bar L_{w_{d+1}}\mid\mathrm{Fl})$. Then:
\begin{enumerate}
\item $h(\ell(\lambda)^{d+1})=0$ for every dominant weight $\lambda$ with
  $\lambda_{\tau,a}=\lambda_{\tau,a+1}$ for some $\tau$ and some $a<N$;
\item $h(\ell(\lambda)^{d+1})\ge0$ for every dominant weight $\lambda$;
\item $h(p_{\tau,1}Q)=0$ for every $\tau\in\Psi$ and every homogeneous $Q$ of degree $d$;
\item $h(F^\pi)=h(F)$ for every permutation $\pi$ of $\Psi$, with $F^\pi$ as in
  Lemma~\ref{cv-reconstruction-flag-degree}(4).
\end{enumerate}
\end{lemma}
\leanref{Development/AutomorphicSpectrum/Definitions/DivisibilityTowers/HasBaseDegreeSetting.lean}{GlobalTheory.CompanionTheory.Ar.Top.hhFl}
\leanref{Development/AutomorphicSpectrum/Proofs/DivisibilityProofs/FlagHeightI9PropertiesStatement.lean}{GlobalTheory.T.flag_height_I9_properties}
\leanref{Development/AutomorphicSpectrum/Proofs/MetricProofs/FlagHeightStatement.lean}{GlobalTheory.T.flag_height}
\leanref{Development/AutomorphicSpectrum/Proofs/DivisibilityProofs/FlagHeightPowerSumOneZeroStatement.lean}{GlobalTheory.T.flag_height_power_sum_one_zero}
\leanref{Development/AutomorphicSpectrum/Proofs/DivisibilityProofs/FlagHeightLabelInvariantStatement.lean}{GlobalTheory.T.flag_height_label_invariant}
\begin{proof}
The intersection numbers are multilinear in the multidegrees of the Fubini--Study bundles. As
$\sum_i\lambda_{\tau,i}x_{\tau,i}$ corresponds to
$\sum_i\lambda_{\tau,i}(\bar{\mathcal O}_{\tau,i}-\bar{\mathcal O}_{\tau,i-1})=\sum_{q=1}^{N-1}(\lambda_{\tau,q}-\lambda_{\tau,q+1})\bar{\mathcal O}_{\tau,q}$,
\[
 h(\ell(\lambda)^{d+1})=\sum_w a_w\prod_{r=1}^{d+1}g_{w_r},\qquad
 a_w=(\bar{\mathcal O}_{w_1}\cdots\bar{\mathcal O}_{w_{d+1}}\mid\mathrm{Fl}),
\]
where $w$ runs over the words of length $d+1$ in the pairs $(\tau,q)$ with $1\le q<N$, and
$g_{\tau,q}=\lambda_{\tau,q}-\lambda_{\tau,q+1}$.

(1) If a word $w$ misses a pair $(\tau,q)$, the bundles $\bar{\mathcal O}_{w_r}$ are pulled back from
the image of $\mathrm{Fl}$ in the product of the Plücker spaces other than that of $V_q$ at $\tau$.
This image has dimension less than $d+1$, so $a_w=0$ by the projection formula. Every word with
$a_w\ne0$ therefore contains every pair, and its term contains the factor $g_{\tau,a}=0$.

(2) If the $g_{\tau,q}$ are nonnegative integers, $\sum g_{\tau,q}\bar{\mathcal O}_{\tau,q}$ is a
Fubini--Study bundle generated by monomials in the Plücker coordinates, which have norm at most $1$;
it is semipositive, nef on the fibres, and of nonnegative degree on every horizontal curve, so its
self-intersection number is nonnegative by Proposition~\ref{lit-hilbert-samuel-upper}(1). The
polynomial $\sum_wa_w\prod_rg_{w_r}$ is homogeneous, so it is nonnegative at nonnegative rational
$g$, and by continuity at nonnegative real $g$. A dominant weight has nonnegative $g$.

(3) The polynomial $p_{\tau,1}=\sum_ix_{\tau,i}$ corresponds to
$\sum_i(\bar{\mathcal O}_{\tau,i}-\bar{\mathcal O}_{\tau,i-1})=0$, the trivial bundle with the trivial
metric, so every intersection number with it as a factor vanishes by multilinearity.

(4) A permutation $\pi$ of $\Psi$ permutes the factors of $\mathrm{Fl}$ and of the ambient product of
projective spaces, carries $\bar L_{\tau,i}$ to $\bar L_{\pi(\tau),i}$, and preserves the metrics. The
projection formula for this isomorphism gives (4).
\end{proof}

\begin{lemma}\label{cv-transport-power-sums}
In the notation of Definition~\ref{cv-reconstruction-functionals}, every homogeneous $F\in\mathbf R[x]$
of degree $d+1$ can be written $F=\sum_\tau\sum_{k=1}^Np_{\tau,k}Q_{\tau,k}$ with $Q_{\tau,k}$
homogeneous of degree $d+1-k$.
\end{lemma}
\leanref{Development/AutomorphicSpectrum/Proofs/DivisibilityProofs/HomogeneousPowerSumDecompositionStatement.lean}{GlobalTheory.T.homogeneous_power_sum_decomposition}

For each $\tau$, the polynomials in $x_{\tau,1},\dots,x_{\tau,N}$ are spanned over the symmetric ones by
the monomials $\prod_ix_{\tau,i}^{e_i}$ with $e_i<i$, of degree at most $N(N-1)/2$. So every
monomial of degree greater than $d$ lies in the ideal generated by the symmetric polynomials of
positive degree in the variables of one $\tau$, and in characteristic $0$ these are polynomials in
$p_{\tau,1},\dots,p_{\tau,N}$ by Newton's identities.

\begin{lemma}\label{cv-transport-local-terms}
In the notation of Definition~\ref{cv-reconstruction-functionals}, for $\sigma\in\Psi$ and $k\ge0$ let
$\pi_k(\sigma,+)=p^{\le j}_{\sigma,k}$ and $\pi_k(\sigma,-)=\sum_{i>j}(-x_{\sigma,i})^k$. For every prime
$p$ and every $\tau\in\Psi$ let $(\phi_{p,\tau,n})_{n\ge0}$ be a sequence in $\Psi\times\{+,-\}$, such
that the sign of $\phi_{p,\tau,n}$ does not depend on $\tau$. For $\tau\in\Psi$, $k\ge1$ and
$Q\in\mathbf R[x]$ put
\[
 \delta_p(\tau,k,Q)=k\log p\Big(\sum_{n\ge0}p^{-k(n+1)}\deg_{\mathrm{fl}}\big(\pi_{k-1}(\phi_{p,\tau,n})Q\big)
 -\frac{\deg_{\mathrm{fl}}(p^{\le j}_{\tau,k-1}Q)}{p-1}\Big),
\]
and let $r_p(\tau,k,Q)$ be $k\log p\sum_{n\ge0}p^{-k(n+1)}\deg_{\mathrm{fl}}(\pi_{k-1}(\phi_{p,\tau,n})Q)$ if
$k\ge2$ and $0$ if $k=1$. Then there are real numbers $\kappa_p$ with
\[
 \delta_p(\tau,k,Q)=-\frac{\log p}{p-1}\sum_\sigma\mathrm{Sh}_\sigma(p_{\tau,k}Q)
 +\kappa_p\sum_\sigma\mathrm{Tot}_\sigma(p_{\tau,k}Q)+r_p(\tau,k,Q)
\]
for every prime $p$, every $\tau$, every $k\ge1$ and every $Q$. For fixed $\tau$, $k$ and $Q$ there is
$c$ with $|r_p(\tau,k,Q)|\le c\log p/p^2$ for every prime $p$; so for every finite set $S$ of primes and
every $B_0\ge1$ the sum $R_{B_0}(\tau,k,Q)$ of the $r_p(\tau,k,Q)$ over the primes $p>B_0$ outside $S$
converges.
\end{lemma}
\leanref{Development/AutomorphicSpectrum/Definitions/DivisibilityTowers/CompactGapsGE.lean}{GlobalTheory.CompanionTheory.Ar.Top.deltaTerm}
\leanref{Development/AutomorphicSpectrum/Definitions/DivisibilityTowers/HeckeSectionsIsotropicCommutatorPairingPower.lean}{GlobalTheory.CompanionTheory.Ar.Top.restTerm}
\leanref{Development/AutomorphicSpectrum/Definitions/DivisibilityTowers/HeckeSectionsIsotropicCommutatorPairingPower.lean}{GlobalTheory.CompanionTheory.Ar.Top.restSum}
\leanref{Development/AutomorphicSpectrum/Proofs/DivisibilityProofs/PowerSumDeltaDecompositionStatement.lean}{GlobalTheory.T.power_sum_delta_decomposition}
\leanref{Development/AutomorphicSpectrum/Proofs/DivisibilityProofs/RemainderTermSummableStatement.lean}{GlobalTheory.T.remainder_term_summable}

This is a self-contained computation with Lemma~\ref{cv-reconstruction-flag-degree}(3). For $k\ge2$ the
first term of $\delta_p$ is $r_p$, the second is the term in $\mathrm{Sh}$, and the $\mathrm{Tot}$ term
vanishes. For $k=1$,
$\pi_0(\sigma,+)=j$ and $\pi_0(\sigma,-)=N-j$, so the series is $\deg_{\mathrm{fl}}(Q)$ times a number that
depends only on $p$, because the signs do not depend on $\tau$; this number, divided by $N$ and
multiplied by $\log p$, is $\kappa_p$. For $k\ge2$ the series defining $r_p$ is at most
$2p^{-2}$ times the largest of the finitely many $|\deg_{\mathrm{fl}}(\pi_{k-1}(\sigma,\pm)Q)|$.

\begin{definition}\label{cv-transport-test-polynomial}
In the notation of Definition~\ref{cv-reconstruction-functionals}, put $D_j=|\Psi|\,j(N-j)$ and
$m_j=\min(j,N-j)$. Let $\mu_j=\sum_\tau\sum_{i>j}x_{\tau,i}$ if $N-j\le j$, and
$\mu_j=-\sum_\tau\sum_{i\le j}x_{\tau,i}$ otherwise. Let $\eta_j$ be the product of the
$x_{\tau,b}-x_{\tau,a}$ over all $\tau$ and all pairs $a<b$ that lie both in $\{1,\dots,j\}$ or both in
$\{j+1,\dots,N\}$. The \emph{test polynomial} of type $j$ is $t_j=\mu_j^{D_j+1}\eta_j$, and its
\emph{test degree} is $T_j=\deg_{\mathrm{fl}}(\mu_j^{D_j}\eta_j)$.
\end{definition}
\leanref{Development/AutomorphicSpectrum/Definitions/DivisibilityTowers/HeckeSectionsIsotropicCommutatorPairingPower.lean}{GlobalTheory.CompanionTheory.Ar.Top.muSmall}
\leanref{Development/AutomorphicSpectrum/Definitions/DivisibilityTowers/HeckeSectionsIsotropicCommutatorPairingPower.lean}{GlobalTheory.CompanionTheory.Ar.Top.fibClass}
\leanref{Development/AutomorphicSpectrum/Definitions/DivisibilityTowers/HeckeSectionsIsotropicCommutatorPairingPower.lean}{GlobalTheory.CompanionTheory.Ar.Top.dimX}
\leanref{Development/AutomorphicSpectrum/Definitions/DivisibilityTowers/HeckeSectionsIsotropicCommutatorPairingPower.lean}{GlobalTheory.CompanionTheory.Ar.Top.testPoly}
\leanref{Development/AutomorphicSpectrum/Definitions/DivisibilityTowers/HeckeSectionsIsotropicCommutatorPairingPower.lean}{GlobalTheory.CompanionTheory.Ar.Top.testDeg}

\begin{lemma}\label{cv-transport-test-facts}
In the notation of Definition~\ref{cv-transport-test-polynomial}, $t_j$ is homogeneous of degree $d+1$
and $T_j>0$. For every $\tau\in\Psi$: if $N-j\le j$, then $\mathrm{Sh}_\tau(t_j)=0$ and
$\mathrm{Tot}_\tau(t_j)=(D_j+1)(N-j)\,T_j$; otherwise
$\mathrm{Sh}_\tau(t_j)=\mathrm{Tot}_\tau(t_j)=-(D_j+1)\,j\,T_j$.
\end{lemma}
\leanref{Development/AutomorphicSpectrum/Proofs/DivisibilityProofs/TestPolynomialShiftTotalStatement.lean}{GlobalTheory.T.test_polynomial_shift_total}

This is a self-contained computation. The degree of $\eta_j$ is
$|\Psi|\big(j(j-1)+(N-j)(N-j-1)\big)/2=d-D_j$. The operators $\sum_{i\le j}\partial/\partial x_{\tau,i}$ and
$\sum_i\partial/\partial x_{\tau,i}$ kill $\eta_j$, and they send $\mu_j^{D_j+1}$ to $(D_j+1)\mu_j^{D_j}$
times $0$ and $N-j$ in the first case, and times $-j$ and $-j$ in the second. The positivity of $T_j$
is a computation of the coefficients of $\mu_j^{D_j}\eta_j$.

\begin{lemma}\label{cv-transport-remainders}
In the notation of Definitions~\ref{cv-reconstruction-functionals} and~\ref{cv-transport-test-polynomial}
and Lemma~\ref{cv-transport-local-terms}, let $S$ be a finite set of primes and $B\ge1$ an integer.
For a decomposition $F=\sum_\tau\sum_{k=1}^Np_{\tau,k}Q_{\tau,k}$ as in
Lemma~\ref{cv-transport-power-sums}, put $R_B(F)=\sum_\tau\sum_{k=1}^NR_B(\tau,k,Q_{\tau,k})$.
\begin{enumerate}
\item If $u$ is regular dominant with gaps at least $g>0$, $B+1>4|\Psi|N^2(1+\|u\|/g)$ and
  $F=\ell(u)^{d+1}$, then
  $|R_B(F)|\le2(d+1)\deg_{\mathrm{fl}}(\ell(u)^d)\,m_j\operatorname{spread}(u)$.
\item If $B\ge4|\Psi|N^2$ and $F=t_j$, then $|R_B(F)|\le1.6\,(D_j+1)\,T_j$.
\end{enumerate}
\end{lemma}
\leanref{Development/AutomorphicSpectrum/Proofs/DivisibilityProofs/LinRootRemainderBoundStatement.lean}{GlobalTheory.T.lin_root_remainder_bound}
\leanref{Development/AutomorphicSpectrum/Proofs/DivisibilityProofs/TestPolynomialRemainderBoundStatement.lean}{GlobalTheory.T.test_polynomial_remainder_bound}

These are self-contained computations, uniform in the sequences $\phi_{p,\tau,n}$ and in $S$. Only
the terms with $k\ge2$ contribute. In (1), with $A=N(1+2\|u\|/g)$, the numbers
$|k\deg_{\mathrm{fl}}(\pi_{k-1}(\sigma,\pm)Q_{\tau,k})|$ are at most
$(d+1)\deg_{\mathrm{fl}}(\ell(u)^d)\,m_j\,2^{k-1}A^{k-1}(\max_iu_{\tau,i}-\min_iu_{\tau,i})/k!$, and the sum
over the primes $p>B$ of $\log p\sum_np^{-k(n+1)}$ is at most $5.56\,B'^{\,1-k}$ with
$B'=4|\Psi|N^2(1+\|u\|/g)<B+1$. Since $2A\le B'/2$, the sum over $k\ge2$ is at most twice
$(d+1)\deg_{\mathrm{fl}}(\ell(u)^d)\,m_j(\max_iu_{\tau,i}-\min_iu_{\tau,i})$, and summing over $\tau$ gives
(1). Part (2) is the same computation for $t_j$, with $B\ge4|\Psi|N^2$.

\begin{lemma}\label{cv-transport-grassmannian}
In the notation of Definition~\ref{cv-transport-test-polynomial}, let $h$ be the flag height of
Lemma~\ref{cv-reconstruction-flag-height}. Then
\[
 h(t_j)\le(D_j+1)\,T_j\,|\Psi|\,\frac{m_j}2\log(N+1).
\]
\end{lemma}
\leanref{Development/AutomorphicSpectrum/Proofs/DivisibilityProofs/FlagHeightTestBoundStatement.lean}{GlobalTheory.T.flag_height_test_bound}
\leanref{GroundTheory/HeightsAndTheta/Heights/HodgeLowerBound/FlagIndices.lean}{GroundTheory.Heights₂.Heights.LowerBoundFree.GrassmannianHeightZTwo.Statement}
\begin{proof}
Let $G$ be the product over $\Psi$ of the Grassmannians of $j$-dimensional subspaces of $\mathbf Z^N$
in their Plücker embeddings, of dimension $D_j$, with the Fubini--Study bundle
$\bar{\mathcal O}_G=\sum_\tau\bar{\mathcal O}_{\tau,j}$, and let $\deg(G)$ be its degree on the generic
fibre. In the correspondence of the proof of Lemma~\ref{cv-reconstruction-flag-height} between
polynomials and bundles, $\mu_j$ corresponds to $-\bar{\mathcal O}_G$ pulled back to $\mathrm{Fl}$, up to
multiples of the $p_{\tau,1}$, which do not contribute to $h$ by
Lemma~\ref{cv-reconstruction-flag-height}(3). The projection $\mathrm{Fl}\to G$ has as fibres products
of the flag varieties of the two blocks, and $\eta_j$ is a polynomial in the classes of these blocks.
The projection formula for $\mathrm{Fl}\to G$ therefore gives
\[
 h(t_j)\deg(G)=T_j\,(\bar{\mathcal O}_G^{D_j+1}\mid G).
\]
For an $(N-j)\times j$ matrix $Z$ whose entries are roots of unity, let $V_Z$ be the column space of
the $N\times j$ matrix with blocks $1_j$ and $Z$. One of its Plücker coordinates is $1$ and all are
algebraic integers, and by the Cauchy--Binet formula the square of the Euclidean norm of the Plücker
vector at each complex embedding is $\det(1+ZZ^*)$, where $Z$ is replaced by its conjugate. The
matrix $ZZ^*$ has rank at most $m_j$ and trace $j(N-j)$, so
$\det(1+ZZ^*)\le(1+j(N-j)/m_j)^{m_j}\le(N+1)^{m_j}$, and the height of $V_Z$ is at most
$\frac{m_j}2\log(N+1)$. The points of $G$ whose components are such $V_Z$ are Zariski dense and have
height at most $|\Psi|\frac{m_j}2\log(N+1)$, so
Theorem~\ref{cv-wall-successive-minima}(2) bounds $(\bar{\mathcal O}_G^{D_j+1}\mid G)$ by
$(D_j+1)|\Psi|\frac{m_j}2\log(N+1)\deg(G)$. Dividing by $\deg(G)>0$ gives the bound.
\end{proof}
\leanref{Development/AutomorphicSpectrum/Proofs/DivisibilityProofs/FlagHeightTestBound.lean}{GlobalTheory.CompanionTheory.Ar.Top.HhTestProof.hhTestBound_hhFl}
\leanref{GroundTheory/HeightsAndTheta/Heights/HodgeLowerBound/HFl.lean}{GroundTheory.Heights₂.Heights.LowerBoundFree.GrassmannianHeightZTwo.proof}

\subsection{Analytic systems and proportionality}\label{subsection-cv-analytic-systems-proportionality}

\begin{definition}\label{cv-condition-flag-model}
In the notation of Definition~\ref{cv-reconstruction-functionals}, let $\mathbf R[x,a]$ be the
polynomial ring over $\mathbf R[x]$ in one more variable $a$, and for a weight $\lambda$ and
$s\in\mathbf R$ put $\ell(\lambda,s)=\ell(\lambda)+s\,a$.

A \emph{flag model} $M$ of type $(\Psi,N)$ consists of:
\begin{enumerate}
\item an arithmetic subvariety $X$ of dimension $d+1$ of a product
  $\prod_{k\in K}\mathbf P^{I_k}_{\mathbf Z}$ (Definition~\ref{lit-multiprojective}), the kernel of
  a point $y$ with coordinates in a field $\Omega$ that is its function field;
\item for each $t>0$ and each $(\tau,i)$ a smooth metrized line bundle $\bar x_{\tau,i}(t)$,
  invariant under conjugation, whose multidegree does not depend on $t$;
\item a smooth metrized line bundle $\bar A$ invariant under conjugation, and integers $k_A\ge1$
  and $\delta\ge1$.
\end{enumerate}
Let $(\cdot\mid\cdot)$ be an arithmetic intersection system. For $t>0$ let $\mathrm{PP}^t_M$ be the
linear form on $\mathbf R[x,a]$ that vanishes on the monomials of degree other than $d+1$ and
sends a monomial $a^e\prod_{\tau,i}x_{\tau,i}^{e_{\tau,i}}$ of that degree
to $(\bar A^{e}\prod_{\tau,i}\bar x_{\tau,i}(t)^{e_{\tau,i}}\mid X)/(\delta k_A^{e})$. When the limits
exist, $\mathrm{PP}_M(F)=\lim_{t\to0^+}\mathrm{PP}^t_M(F)$ is the \emph{arithmetic intersection
number} of $F$. The \emph{geometric degree} $\deg_M$ is the linear form on $\mathbf R[x,a]$ that
vanishes on the monomials of degree other than $d$ and sends a monomial
$a^e\prod_{\tau,i}x_{\tau,i}^{e_{\tau,i}}$ of that degree to the integral over the complex points of
$X$ of the product of the first Chern forms of $\bar A$ ($e$ times) and of the
$\bar x_{\tau,i}(1)$ ($e_{\tau,i}$ times), divided by $\delta k_A^{e}$.
\end{definition}
\leanref{Development/AutomorphicSpectrum/Definitions/DivisibilityTowers/Birational.lean}{GlobalTheory.CompanionTheory.Ar.Top.RootModel}
\leanref{Development/AutomorphicSpectrum/Definitions/DivisibilityTowers/Birational.lean}{GlobalTheory.CompanionTheory.Ar.Top.RootModel.arFun}
\leanref{Development/AutomorphicSpectrum/Definitions/DivisibilityTowers/Birational.lean}{GlobalTheory.CompanionTheory.Ar.Top.RootModel.PP}
\leanref{Development/AutomorphicSpectrum/Definitions/DivisibilityTowers/Birational.lean}{GlobalTheory.CompanionTheory.Ar.Top.RootModel.geoFun}
\leanref{Development/AutomorphicSpectrum/Definitions/DivisibilityTowers/CompactGapsGE.lean}{GlobalTheory.CompanionTheory.Ar.Top.linA}

\begin{definition}\label{cv-proportionality-analytic-system}
Let $\Psi$ be a finite set, let $N\ge0$, put $d=|\Psi|\,N(N-1)/2$ and
$\Sigma=(\Psi\times\{1,\dots,N\})\sqcup\{A\}$.
\begin{enumerate}
\item An \emph{analytic system} $S$ of type $(\Psi,N)$ consists of a compact complex manifold $S$
  of dimension $d$ with a K\"ahler form $\omega$; a divisor $\partial S$ with
  simple normal crossings, the zero set of a section $\sigma$ of norm at most $1$ of a Hermitian
  line bundle; numbers $C\ge0$ and $0<t_0<1$; and for each $b\in\Sigma$ a holomorphic line bundle
  $\bar L_b$ on $S$ with a smooth Hermitian metric, a smooth function $\varphi_b$ on
  $S\setminus\partial S$ and smooth functions $\varphi_b(t)$ on $S$ for $0<t\le t_0$, with
  $\varphi_A=\varphi_A(t)=0$. Put $r=\exp(1)-2\log\|\sigma\|$ and
  $\varepsilon(t)=C\log\log(1/t)/\log(1/t)$, and write $\bar L_b(t)$ for $\bar L_b$ with its
  metric multiplied by $e^{-\varphi_b(t)}$. These satisfy, off $\partial S$ where $\varphi_b$
  appears: the metric of $\bar L_b$ multiplied by $e^{-\varphi_b}$ is semipositive and
  $-C\log r\le\varphi_b\le C$; $\varphi_b\le\varphi_b(t)\le C$,
  $\varphi_b(t)\ge-C\log\log(1/t)$ and $\varphi_b(t)\le\varphi_b+t^{1/3}+\varepsilon(t)r$, with
  $\varphi_b(t)\le\varphi_b+t^{1/3}+\varepsilon(t)$ where $r\le\log(1/t)$;
  $\varphi_b(t)\le\varphi_b(t')+C\varepsilon(t')$ for $t\le t'\le t_0$; the first Chern form of
  $\bar L_b(t)$ is at least $-\varepsilon(t)\omega$; and $\varphi_b(t)$ tends to $\varphi_b$ with
  its first and second derivatives, uniformly on compact subsets of $S\setminus\partial S$, as
  $t\to0$. The \emph{degree polynomial} of $S$ is
  $\deg_S(c)=\int_S\big(\sum_{b\in\Sigma}c_b\,c_1(\bar L_b)\big)^d$ for
  $c\in\mathbf R^\Sigma$.
\item Let $0\le j\le N$, let $h_j$ be the Hermitian form on $\mathbf C^N$ with matrix
  $\operatorname{diag}(-1,\dots,-1,1,\dots,1)$ with $N-j$ entries $-1$, and for
  $g\in GL_N(\mathbf C)$ and $0\le k\le N$ let $G_k(g)$ be half the logarithm of the absolute
  value of the Gram determinant for $h_j$ of the first $k$ columns of $g$. The \emph{flag domain}
  $D_j$ is the set of $g$ whose first $N-j$ columns span a negative definite subspace, and on it
  put $\psi_{j,i}=G_{N-j+i}-G_{N-j}$ for $i\le j$ and $\psi_{j,i}=G_N-G_{N-j}+G_{i-j}$ for $i>j$.
  The system $S$ is \emph{uniformized of type $j$} if every point of $S\setminus\partial S$ has a
  neighbourhood $V$ with holomorphic maps $g_\tau:V\to D_j$ ($\tau\in\Psi$) such that, for every
  $(\tau,i)$, a local potential of the metric of $\bar L_{\tau,i}$ multiplied by
  $e^{-\varphi_{\tau,i}}$ equals $\psi_{j,i}\circ g_\tau$ plus the real part of a holomorphic
  function on $V$, and if at one point of $S\setminus\partial S$ the derivative of $(g_\tau)_\tau$
  has no nonzero vector $v$ with $g_\tau^{-1}\,dg_\tau(v)$ upper triangular for every $\tau$.
\item A \emph{realization} of a flag model $M$ of type $(\Psi,N)$ on $S$ with multiplicity $n\ge1$
  (Definition~\ref{cv-condition-flag-model}) consists of a compact complex manifold $\tilde X$ with
  a holomorphic map $\tilde X\to X(\mathbf C)$ to the complex points of $X$ that is bijective and
  immersive over the complement of the zero set of a multihomogeneous ideal not contained in the
  ideal of $X$; a holomorphic map $\pi:\tilde X\to S$ of degree $\deg\pi\ge1$ off a closed set of
  measure zero, with $\delta=n\deg\pi$; and for each $b\in\Sigma$ a smooth metrized line bundle
  $\bar y_b$ on $X$, with $\bar y_A=\bar A$, whose first Chern form on $\tilde X$ is the pullback by
  $\pi$ of that of $\bar L_b$. Write $\bar y_b(t)$ for $\bar y_b$ with its metric
  multiplied by $e^{-\varphi_b(t)\circ\pi}$ on $\tilde X$. For $0<t\le t_0$ the metrized line bundle
  $-\sum_{k\le i}\bar x_{\tau,k}(t)$ is $\bar y_{\tau,i}(t)$, and $\bar y_A(t)=\bar A$.
\end{enumerate}
\end{definition}
\leanref{Development/AutomorphicSpectrum/Definitions/MetricPeriods/IsTwiceDifferentiableAt.lean}{GlobalTheory.CompanionTheory.Ar.PB.AnSys}
\leanref{Development/AutomorphicSpectrum/Definitions/MetricPeriods/IsTwiceDifferentiableAt.lean}{GlobalTheory.CompanionTheory.Ar.PB.IsUniformised}
\leanref{Development/AutomorphicSpectrum/Definitions/MetricPeriods/MetricsProjectionFormulaNef.lean}{GlobalTheory.CompanionTheory.Ar.PB.Realis}
\leanref{Development/AutomorphicSpectrum/Definitions/DivisibilityTowers/RefLinked.lean}{GlobalTheory.CompanionTheory.Ar.Top.RealTie}

\begin{definition}\label{cv-period-monge-ampere}
Let $d\ge0$, let $W\subset\mathbf C^d$ be open and let $y\in W$.
\begin{enumerate}
\item Let $\varphi$ be a real function on $W$ that is twice differentiable at $y$. Its
  \emph{complex Hessian} at $y$ is the Hermitian matrix
  $H_\varphi(y)=\big(\partial^2\varphi/\partial z_a\partial\bar z_b(y)\big)_{a,b}$, and its
  \emph{Levi form} at $y$ is the Hermitian form
  $L_\varphi(y)(v)=\sum_{a,b}\frac{\partial^2\varphi}{\partial z_a\partial\bar z_b}(y)\,v_a\bar v_b$
  on $\mathbf C^d$.
\item Let $\varphi_1,\dots,\varphi_d$ be real functions on $W$ that are twice differentiable at $y$.
  Their \emph{Monge--Amp\`ere density} at $y$ is
  \[
   \operatorname{MA}(\varphi_1,\dots,\varphi_d)(y)=\Big(\frac2\pi\Big)^d\sum_{T\subset\{1,\dots,d\}}
   (-1)^{d-|T|}\det\Big(\sum_{k\in T}H_{\varphi_k}(y)\Big).
  \]
  The sum is $d!$ times the mixed discriminant of the $H_{\varphi_k}(y)$, so
  $\operatorname{MA}(\varphi,\dots,\varphi)(y)=d!\,(2/\pi)^d\det H_\varphi(y)$, and
  $\operatorname{MA}(\varphi_1,\dots,\varphi_d)$ is the density of
  $dd^c\varphi_1\wedge\dots\wedge dd^c\varphi_d$ with respect to Lebesgue measure, for
  $dd^c=(i/\pi)\partial\bar\partial$. It does not change when a $\varphi_k$ is changed by the real
  part of a holomorphic function, whose complex Hessian is $0$.
\item Let $S$ be a compact complex manifold of dimension $d$, with a finite atlas of holomorphic
  charts $\kappa_k:V_k\to\mathbf C^d$ and smooth functions $\chi_k\ge0$ on $S$ with support in $V_k$
  and $\sum_k\chi_k=1$. Let $\vartheta_1,\dots,\vartheta_d$ be local potentials of metrics on line
  bundles on $S$, each determined near every point up to the real part of a holomorphic function.
  The \emph{integral over $S$ of their Monge--Amp\`ere measure} is
  \[
   \sum_k\int_{\kappa_k(V_k)}\chi_k\big(\kappa_k^{-1}(y)\big)\,
   \operatorname{MA}(\vartheta_1\circ\kappa_k^{-1},\dots,\vartheta_d\circ\kappa_k^{-1})(y)\,dy,
  \]
  where at each $y$ the potentials are taken near $\kappa_k^{-1}(y)$.
\end{enumerate}
\end{definition}
\leanref{ReciprocityTheory/AnalyticForms/MaIntEmpty.lean}{ReciprocityTheory.AnForm.levi}
\leanref{ReciprocityTheory/AnalyticForms/MaIntEmpty.lean}{ReciprocityTheory.AnForm.hessC}
\leanref{ReciprocityTheory/AnalyticForms/MaIntEmpty.lean}{ReciprocityTheory.AnForm.mdisc}
\leanref{ReciprocityTheory/AnalyticForms/MaIntEmpty.lean}{ReciprocityTheory.AnForm.maDens}
\leanref{ReciprocityTheory/AnalyticForms/PotentialTheory/IsPotentialOfFun.lean}{ReciprocityTheory.AnForm.M.locDens}
\leanref{ReciprocityTheory/AnalyticForms/PotentialTheory/IsPotentialOfFun.lean}{ReciprocityTheory.AnForm.M.ChartPartition}
\leanref{ReciprocityTheory/AnalyticForms/PotentialTheory/IsPotentialOfFun.lean}{ReciprocityTheory.AnForm.M.mInt}

\begin{lemma}\label{cv-period-flag-domain}
Let $0\le j\le N$, and let $h_j$, the flag domain $D_j$ and the functions $G_k$ be as in
Definition~\ref{cv-proportionality-analytic-system}(2). For a complex $N\times N$ matrix $g$ and
$0\le k\le N$ let $\Gamma_k(g)$ be the upper left $k\times k$ block of $g^*h_jg$, the Gram matrix for
$h_j$ of the first $k$ columns of $g$.
\begin{enumerate}
\item The set $D_j$ is open, and $\det\Gamma_k(g)\neq0$ for every $g\in D_j$ and every
  $0\le k\le N$.
\item For every $g\in D_j$ there is $u\in GL_N(\mathbf C)$ with $u^*h_ju=h_j$ such that $ug$ is
  upper triangular.
\item Let $g$ be a complex $N\times N$ matrix and $0\le k\le N$. If $u^*h_ju=h_j$, then
  $G_k(ug)=G_k(g)$. If $b$ is upper triangular and invertible and $\det\Gamma_k(g)\neq0$, then
  $G_k(gb)=G_k(g)+\sum_{a\le k}\log|b_{aa}|$.
\end{enumerate}
\end{lemma}
\leanref{Development/AutomorphicSpectrum/Proofs/MetricProofs/GramLogUnitaryTriangularInvarianceStatement.lean}{GlobalTheory.T.isOpen_flag_domain_nondegenerate_minors}
\leanref{Development/AutomorphicSpectrum/Proofs/MetricProofs/IsOpenFlagDomainNondegenerateMinorsStatement.lean}{GlobalTheory.T.flag_domain_unitary_triangular}
\leanref{Development/AutomorphicSpectrum/Proofs/MetricProofs/LogDetComparableGramBoundStatement.lean}{GlobalTheory.T.gram_log_unitary_triangular_invariance}
\begin{proof}
Let $g\in D_j$, write $h=h_j$ and $G(x,y)=(gx)^*h(gy)$ for $x,y\in\mathbf C^N$, and let $E$ be the
span of $e_1,\dots,e_{N-j}$, on which $G$ is negative definite. Every subspace $W\subset\mathbf C^N$
of dimension $N-j+1$ contains an $x$ with $G(x,x)>0$: the linear map $W\to\mathbf C^{N-j}$ taking
$x$ to the first $N-j$ coordinates of $gx$ has a nonzero kernel vector $x$, and
$G(x,x)=\sum_{a>N-j}|(gx)_a|^2$ is positive since $gx\neq0$.

(1) The two conditions defining $D_j$, that $\det g\neq0$ and that $-\Gamma_{N-j}(g)$ is positive
definite, are open. Suppose that $\Gamma_k(g)$ is singular. Then a nonzero $x$ in the span of
$e_1,\dots,e_k$ is $G$-orthogonal to $e_1,\dots,e_k$, and $G(x,x)=0$. If $k\le N-j$, this
contradicts the definiteness of $G$ on $E$. If $k>N-j$, then $x\notin E$ for the same reason, $x$
is $G$-orthogonal to $E$, and $G(y+\alpha x,y+\alpha x)=G(y,y)\le0$ for $y\in E$ and $\alpha\in\mathbf C$, which
contradicts the previous paragraph for $W=E+\mathbf Cx$.

(2) By (1) the leading principal minors of $g^*hg$ are nonzero, so Gram--Schmidt gives an
invertible upper triangular $c$ with $c^*g^*hgc=\Delta$ diagonal. The first $N-j$ columns of $c$
span $E$, so $\Delta_{aa}=G(ce_a,ce_a)<0$ for $a\le N-j$. For $a>N-j$ the vector $ce_a$ is
$G$-orthogonal to $E$, so $\Delta_{aa}\le0$ would make $G$ negative semidefinite on
$E+\mathbf Cce_a$, of dimension $N-j+1$; hence $\Delta_{aa}>0$. Replacing $c$ by $c|\Delta|^{-1/2}$
gives $(gc)^*h(gc)=h$. Then $u=(gc)^{-1}$ satisfies $u^*hu=h$, and $ug=c^{-1}$ is upper triangular.

(3) Here $g$ is any matrix. If $u^*hu=h$, then $(ug)^*h(ug)=g^*hg$. If $b$ is upper triangular,
the first $k$ columns of $gb$ are those of $g$ times the upper left $k\times k$ block $b_k$ of $b$, so
$\Gamma_k(gb)=b_k^*\Gamma_k(g)b_k$ and $\det\Gamma_k(gb)=|\det b_k|^2\det\Gamma_k(g)$, with
$\det b_k=\prod_{a\le k}b_{aa}$. Half the logarithm of the absolute value gives the claim.
\end{proof}

\begin{lemma}\label{cv-period-big-cell}
Let $d\ge0$, let $y_0\in\mathbf C^d$, and let $f$ be a holomorphic map from a neighbourhood of
$y_0$ to the complex $N\times N$ matrices such that $f(y_0)$ is upper triangular and invertible.
Then there are a neighbourhood $V$ of $y_0$ and holomorphic maps $Z$ and $b$ on $V$, with $Z(y)$
strictly lower triangular, $Z(y_0)=0$ and $b(y)$ upper triangular and invertible, such that
$f(y)=\exp(Z(y))\,b(y)$ for $y\in V$. For $v\in\mathbf C^d$, $dZ(y_0)(v)=0$ if and only if
$f(y_0)^{-1}\,df(y_0)(v)$ is upper triangular.
\end{lemma}
\leanref{Development/AutomorphicSpectrum/Proofs/MetricProofs/LocalBruhatCellAnalyticCoordinatesStatement.lean}{GlobalTheory.T.local_bruhat_cell_analytic_coordinates}
\leanref{Development/WeilBound/CoreStatements/PluriharmonicCharts/XPerCell.lean}{PositivityTheory.Xgp.t_per_cell}
\begin{proof}
Let $\mathfrak n$ be the space of strictly lower triangular matrices and $\mathfrak b$ that of upper
triangular ones, so that $\dim\mathfrak n+\dim\mathfrak b=N^2$. Put $b_0=f(y_0)$. The map
$F(Z,b)=\exp(Z)\,b$ from $\mathfrak n\times\mathfrak b$ to the $N\times N$ matrices is holomorphic,
$F(0,b_0)=b_0$, and its derivative at $(0,b_0)$ is $(\dot Z,\dot b)\mapsto\dot Zb_0+\dot b$. This
derivative is injective: if $\dot Zb_0+\dot b=0$, then $\dot Z=-\dot bb_0^{-1}$ is both strictly
lower and upper triangular, so $\dot Z=0$ and $\dot b=0$. By the holomorphic inverse function
theorem, $F$ has a holomorphic inverse near $b_0$ taking $b_0$ to $(0,b_0)$. Composing it with $f$
gives $Z$ and $b$, and $b(y)$ is invertible near $y_0$. Differentiating $f=\exp(Z)\,b$ at $y_0$ gives
\[
 f(y_0)^{-1}df(y_0)(v)=b_0^{-1}\,dZ(y_0)(v)\,b_0+b_0^{-1}\,db(y_0)(v).
\]
The second term is upper triangular. Since $b_0$ is upper triangular, the first is upper triangular
if and only if $dZ(y_0)(v)$ is, and a strictly lower triangular matrix is upper triangular only if
it is $0$.
\end{proof}

\begin{lemma}\label{cv-period-cell-coordinates}
Let $0\le j\le N$, let $D_j$ and $\psi_{j,i}$ be as in
Definition~\ref{cv-proportionality-analytic-system}(2), let $d\ge0$ and $y_0\in\mathbf C^d$, and let
$g$ be a holomorphic map from a neighbourhood of $y_0$ to $D_j$. Then there is a holomorphic map $Z$
from a neighbourhood of $y_0$ to the strictly lower triangular $N\times N$ matrices, with
$Z(y_0)=0$, such that:
\begin{enumerate}
\item for $1\le i\le N$, $\psi_{j,i}\circ g-\psi_{j,i}\circ\exp(Z)$ is the real part of a
  holomorphic function near $y_0$;
\item for $v\in\mathbf C^d$, $dZ(y_0)(v)=0$ if and only if $g(y_0)^{-1}\,dg(y_0)(v)$ is upper
  triangular.
\end{enumerate}
\end{lemma}
\leanref{Development/WeilBound/CoreStatements/PluriharmonicCharts/LeviDiagonalOfIsPeriodChart.lean}{PositivityTheory.Xgp.per_diag.PerDiag.exists_cell_coords}
\leanref{Development/WeilBound/CoreStatements/PluriharmonicCharts/LeviDiagonalOfIsPeriodChart.lean}{PositivityTheory.Xgp.per_diag.PerDiag.stepPot_unitary_mul}
\leanref{Development/WeilBound/CoreStatements/PluriharmonicCharts/LeviDiagonalOfIsPeriodChart.lean}{PositivityTheory.Xgp.per_diag.PerDiag.stepPot_mul_triangular}
\begin{proof}
By Lemma~\ref{cv-period-flag-domain}(2) there is $u$ with $u^*h_ju=h_j$ such that $ug(y_0)$ is upper
triangular, and it is invertible. Lemma~\ref{cv-period-big-cell} applied to $f=ug$ gives
$ug=\exp(Z)\,b$ near $y_0$, with $Z(y_0)=0$ and $b$ upper triangular and invertible. The matrix
$\exp(Z(y))=ug(y)b(y)^{-1}$ lies in $D_j$: right multiplication by the upper triangular
$b(y)^{-1}$ does not change the span of the first $N-j$ columns, and $u$ preserves $h_j$. So
$\det\Gamma_k(\exp Z(y))\neq0$ by Lemma~\ref{cv-period-flag-domain}(1), and
Lemma~\ref{cv-period-flag-domain}(3) gives
$G_k(g(y))=G_k(ug(y))=G_k(\exp Z(y))+\sum_{a\le k}\log|b_{aa}(y)|$ for every $k$. Hence
$\psi_{j,i}\circ g-\psi_{j,i}\circ\exp(Z)$ is a sum of terms $\pm\log|b_{aa}|$. Each is the real
part of a logarithm of $b_{aa}$, which is holomorphic and does not vanish near $y_0$. For (2),
$df(y_0)(v)=u\,dg(y_0)(v)$, so $f(y_0)^{-1}df(y_0)(v)=g(y_0)^{-1}dg(y_0)(v)$, and
Lemma~\ref{cv-period-big-cell} gives the claim.
\end{proof}

\begin{lemma}\label{cv-period-cell-levi}
Let $0\le j\le N$ and $1\le i\le N$, and let $\psi_{j,i}$ be as in
Definition~\ref{cv-proportionality-analytic-system}(2). Let $\rho$ be the permutation of
$\{1,\dots,N\}$ with $\rho(a)=a+j$ for $a\le N-j$ and $\rho(a)=a-(N-j)$ for $a>N-j$. On the space of
strictly lower triangular $N\times N$ matrices $Z$, with coordinates $w_{ba}=Z_{ba}$ for $a<b$, the
function $Z\mapsto\psi_{j,i}(\exp Z)$ is twice differentiable at $0$, and its Levi form at $0$ is
\[
 w\mapsto\frac12\sum_{\substack{a<b\\ \min(\rho a,\rho b)\le i<\max(\rho a,\rho b)}}|w_{ba}|^2.
\]
\end{lemma}
\leanref{Development/AutomorphicSpectrum/Proofs/MetricProofs/FlagPotentialLeviFormCellStatement.lean}{GlobalTheory.T.flag_potential_levi_form_cell}
\leanref{ReciprocityTheory/AnalyticForms/LeviForms/HomDens.lean}{ReciprocityTheory.AnForm.P.detGram_deriv}
\leanref{ReciprocityTheory/AnalyticForms/LeviForms/HomDens.lean}{ReciprocityTheory.AnForm.P.detGram_levi}
\leanref{ReciprocityTheory/AnalyticForms/LeviForms/HomDens.lean}{ReciprocityTheory.AnForm.P.leviHalfLogAbs}

This is a self-contained computation. Write $h=h_j=\operatorname{diag}(h_1,\dots,h_N)$, let
$(\cdot)_k$ be the upper left $k\times k$ block, and put $f_k(Z)=\det\big((\exp Z)^*h\exp Z\big)_k$,
so that $G_k(\exp Z)=\frac12\log|f_k(Z)|$ and $f_k(0)=h_1\cdots h_k\neq0$. In
$(\exp Z)^*h\exp Z=h+hZ+Z^*h+Z^*hZ+\frac12(hZ^2+Z^{*2}h)+\cdots$ the terms of degree one have zero
diagonal, so $df_k(0)=0$. Expanding the determinant to second order, the terms of $f_k(Z)/f_k(0)$
of type $w\bar w$ are
\[
 \operatorname{tr}\big(h_k^{-1}(Z^*hZ)_k\big)-\operatorname{tr}\big(h_k^{-1}(Z^*h)_kh_k^{-1}(hZ)_k\big)
 =\sum_{a\le k,\ b>a}\frac{h_b}{h_a}|w_{ba}|^2-\sum_{a<b\le k}\frac{h_b}{h_a}|w_{ba}|^2
 =\sum_{a\le k<b}\frac{h_b}{h_a}|w_{ba}|^2,
\]
and the other terms of degree two are holomorphic or antiholomorphic. Since $df_k(0)=0$, the Levi
form of $\frac12\log|f_k|$ at $0$ is that of $f_k$ divided by $2f_k(0)$, which is
$\frac12\sum_{a\le k<b}(h_b/h_a)|w_{ba}|^2$. The claim follows by taking the combinations of the
$G_k$ that define $\psi_{j,i}$, separately for the pairs with $b\le N-j$, with $a\le N-j<b$ and with
$a>N-j$. For instance, for $i\le j$ a pair with $a\le N-j<b$ has $h_b/h_a=-1$; it contributes
$\frac12$ through $-G_{N-j}$, and $-\frac12$ through $G_{N-j+i}$ when $b>N-j+i$. So its coefficient
is $\frac12$ if $\rho b=b-(N-j)\le i$ and $0$ otherwise, while $\rho a=a+j>i$.

\begin{definition}\label{cv-period-chart}
Let $\Psi$ be a finite set and $0\le j\le N$, let $D_j$ and $\psi_{j,i}$ be as in
Definition~\ref{cv-proportionality-analytic-system}(2), let $S$ be a complex manifold, let $x\in S$,
and for $\tau\in\Psi$ and $1\le i\le N$ let $\vartheta_{\tau,i}$ be a real function on a
neighbourhood of $x$. A \emph{period chart of type $j$} for the $\vartheta_{\tau,i}$ at $x$ is an
open neighbourhood $V$ of $x$ with holomorphic maps $g_\tau:V\to D_j$ ($\tau\in\Psi$) such that
$\vartheta_{\tau,i}-\psi_{j,i}\circ g_\tau$ is the real part of a holomorphic function on $V$ for
every $(\tau,i)$. It is a \emph{flag immersion} at $x$ if for every holomorphic chart $\kappa$ of $S$
at $x$, with $y=\kappa(x)$, no nonzero tangent vector $v$ at $y$ makes
$g_\tau(x)^{-1}\,d(g_\tau\circ\kappa^{-1})(y)(v)$ upper triangular for every $\tau$.
\end{definition}
\leanref{Development/AutomorphicSpectrum/Definitions/MetricPeriods/IsTwiceDifferentiableAt.lean}{GlobalTheory.CompanionTheory.Ar.PB.IsPeriodChart}
\leanref{Development/AutomorphicSpectrum/Definitions/MetricPeriods/IsTwiceDifferentiableAt.lean}{GlobalTheory.CompanionTheory.Ar.PB.IsFlagImmersionAt}

\begin{lemma}\label{cv-period-levi-diagonal}
Let $\Psi$ be a finite set and $0\le j\le N$, and let $\mathcal P$ be the set of pairs $(a,b)$ with
$1\le a<b\le N$. Let $S$ be a complex manifold of dimension $d\ge0$, let $x\in S$, let
$\vartheta_{\tau,i}$ ($\tau\in\Psi$, $1\le i\le N$) be real functions on a neighbourhood of $x$ with
a period chart $(V,(g_\tau))$ of type $j$ at $x$ (Definition~\ref{cv-period-chart}), and let
$\kappa$ be a holomorphic chart of $S$ at $x$, with $y=\kappa(x)$. Then there is a complex linear
map $\ell:\mathbf C^d\to\mathbf C^{\Psi\times\mathcal P}$ such that every
$\vartheta_{\tau,i}\circ\kappa^{-1}$ is twice differentiable at $y$, with Levi form
(Definition~\ref{cv-period-monge-ampere})
\[
 L_{\vartheta_{\tau,i}\circ\kappa^{-1}}(y)(v)=\frac12\sum_{(a,b)\in\mathcal P,\ a\le i<b}
 |\ell(v)_{\tau,a,b}|^2 .
\]
If the period chart is a flag immersion at $x$, then $\ell$ is injective.
\end{lemma}
\leanref{Development/WeilBound/CoreStatements/PluriharmonicCharts/LeviDiagonalOfIsPeriodChart.lean}{PositivityTheory.Xgp.per_diag.levi_diagonal_of_isPeriodChart}
\leanref{Development/AutomorphicSpectrum/Proofs/MetricProofs/PeriodChartPotentialLeviFormStatement.lean}{GlobalTheory.T.period_chart_potential_levi_form}
\begin{proof}
For each $\tau$, Lemma~\ref{cv-period-cell-coordinates} applied to $g_\tau\circ\kappa^{-1}$ at $y$
gives $Z_\tau$. Then $\vartheta_{\tau,i}\circ\kappa^{-1}-\psi_{j,i}\circ\exp\circ Z_\tau$ is the real
part of a holomorphic function near $y$, and its Levi form vanishes. Since $Z_\tau$ is holomorphic
with $Z_\tau(y)=0$, the Levi form of $\psi_{j,i}\circ\exp\circ Z_\tau$ at $y$ is that of
$\psi_{j,i}\circ\exp$ at $0$ evaluated at $dZ_\tau(y)(v)$. By Lemma~\ref{cv-period-cell-levi},
$\vartheta_{\tau,i}\circ\kappa^{-1}$ is twice differentiable at $y$, and its Levi form at $y$ is
$\frac12$ times the sum of $|dZ_\tau(y)(v)_{ba}|^2$ over the pairs $a<b$ with
$\min(\rho a,\rho b)\le i<\max(\rho a,\rho b)$. The map
$(a,b)\mapsto(\min(\rho a,\rho b),\max(\rho a,\rho b))$ is a bijection of $\mathcal P$, and
$\ell(v)_{\tau,\min(\rho a,\rho b),\max(\rho a,\rho b)}=dZ_\tau(y)(v)_{ba}$ defines $\ell$ with these
Levi forms. If $\ell(v)=0$, then $dZ_\tau(y)(v)=0$ for every $\tau$, so
$g_\tau(x)^{-1}d(g_\tau\circ\kappa^{-1})(y)(v)$ is upper triangular for every $\tau$ by
Lemma~\ref{cv-period-cell-coordinates}(2). At a flag immersion this forces $v=0$.
\end{proof}
\leanref{Development/WeilBound/CoreStatements/PluriharmonicCharts/LeviDiagonalOfIsPeriodChart.lean}{PositivityTheory.Xgp.per_diag.PerDiag.levi_of_cell}

\begin{lemma}\label{cv-period-density-diagonal}
Let $d\ge0$, let $\varphi$ be a real function on an open subset of $\mathbf C^d$ that is twice
differentiable at $y$, let $M$ be a complex $d\times d$ matrix and $\delta_1,\dots,\delta_d$ real
numbers, and suppose that $L_\varphi(y)(v)=\sum_{r=1}^d\delta_r|(Mv)_r|^2$ for all
$v\in\mathbf C^d$ (Definition~\ref{cv-period-monge-ampere}). Then
\[
 \operatorname{MA}(\varphi,\dots,\varphi)(y)=d!\,\Big(\frac2\pi\Big)^d\,\delta_1\cdots\delta_d\,
 |\det M|^2.
\]
\end{lemma}
\leanref{Development/AutomorphicSpectrum/Proofs/MetricProofs/MongeAmpereDensityDiagonalLeviStatement.lean}{GlobalTheory.T.monge_ampere_density_diagonal_levi}
\begin{proof}
A Hermitian form is determined by its values on the diagonal, so the transpose of $H_\varphi(y)$ is
$M^*\operatorname{diag}(\delta_1,\dots,\delta_d)M$, whose determinant is
$\delta_1\cdots\delta_d|\det M|^2$. By Definition~\ref{cv-period-monge-ampere},
$\operatorname{MA}(\varphi,\dots,\varphi)(y)=d!\,(2/\pi)^d\det H_\varphi(y)$.
\end{proof}

\begin{lemma}\label{cv-period-pointwise}
Let $\Psi$ be a finite set and $0\le j\le N$, put $d=|\Psi|N(N-1)/2$, and let
$S$ be a complex manifold of dimension $d$. Let $x\in S$, let $\vartheta_{\tau,i}$ ($\tau\in\Psi$,
$1\le i\le N$) be real functions on a neighbourhood of $x$ with a period chart of type $j$ at $x$
(Definition~\ref{cv-period-chart}), and let $\kappa$ be a holomorphic chart of $S$ at $x$, with
$y=\kappa(x)$. For $c\in\mathbf R^{\Psi\times\{1,\dots,N\}}$ put
$\vartheta_c=\sum_{\tau,i}c_{\tau,i}\vartheta_{\tau,i}$. Then there is $\nu\ge0$ such that for every
$c$
\[
 \operatorname{MA}(\vartheta_c\circ\kappa^{-1},\dots,\vartheta_c\circ\kappa^{-1})(y)
 =\nu\prod_{\tau\in\Psi}\prod_{1\le a<b\le N}\ \sum_{m=a}^{b-1}c_{\tau,m}
\]
(Definition~\ref{cv-period-monge-ampere}), and $\nu>0$ if the period chart is a flag immersion at
$x$.
\end{lemma}
\leanref{Development/WeilBound/CoreStatements/PluriharmonicCharts/XPerPoint.lean}{PositivityTheory.Xgp.t_per_point}
\begin{proof}
Let $\mathcal P$ be the set of pairs $(a,b)$ with $1\le a<b\le N$, and let $\ell$ be the linear map
of Lemma~\ref{cv-period-levi-diagonal}. The Levi form at $y$ is linear in the function, among
functions twice differentiable at $y$, so
\[
 L_{\vartheta_c\circ\kappa^{-1}}(y)(v)=\sum_{\tau,i}c_{\tau,i}\,\frac12\sum_{a\le i<b}
 |\ell(v)_{\tau,a,b}|^2
 =\sum_\tau\sum_{(a,b)\in\mathcal P}\Big(\frac12\sum_{m=a}^{b-1}c_{\tau,m}\Big)
 |\ell(v)_{\tau,a,b}|^2.
\]
The set $\Psi\times\mathcal P$ has $d$ elements. Numbering it, $\ell$ becomes a $d\times d$ matrix
$M$, and Lemma~\ref{cv-period-density-diagonal} with
$\delta_{\tau,a,b}=\frac12\sum_{m=a}^{b-1}c_{\tau,m}$ gives the identity with
$\nu=d!\,\pi^{-d}|\det M|^2\ge0$, which does not depend on $c$. If the period chart is a flag
immersion at $x$, then $\ell$ is injective by Lemma~\ref{cv-period-levi-diagonal}, so
$\det M\neq0$ and $\nu>0$.
\end{proof}

\begin{theorem}\label{cv-proportionality-mumford}
Let $\Psi$ be a finite set and $0\le j\le N$, put $d=|\Psi|N(N-1)/2$, and
let $S$ be an analytic system of type $(\Psi,N)$ that is uniformized of type $j$
(Definition~\ref{cv-proportionality-analytic-system}). Then there is $v_S>0$ such that for every
$\lambda\in(\mathbf R^N)^\Psi$, with $\lambda_{\tau,N+1}=0$,
\[
 \int_S\Big(\sum_{\tau,i}(\lambda_{\tau,i}-\lambda_{\tau,i+1})\,c_1(\bar L_{\tau,i})\Big)^{d}
 =v_S\prod_\tau\prod_{1\le i<k\le N}(\lambda_{\tau,i}-\lambda_{\tau,k}).
\]
\end{theorem}
\leanref{Development/AutomorphicSpectrum/Proofs/MetricProofs/UniformisedDegreePolynomialAntisymmetricStatement.lean}{GlobalTheory.T.uniformised_degree_polynomial_antisymmetric}
\leanref{Development/AutomorphicSpectrum/Proofs/MetricProofs/UniformisedDegreePolynomialAntisymmetric.lean}{GlobalTheory.CompanionTheory.Ar.PB.anSys_proportional}
\leanref{Development/AutomorphicSpectrum/Statements/MetricStatements/AgreementIffGuarded.lean}{GlobalTheory.CompanionTheory.Ar.PB.Proportional}
\leanref{Development/AutomorphicSpectrum/Proofs/MetricProofs/ProportionalStatement.lean}{GlobalTheory.T.proportional'}

\begin{proof}
For a word $b_1,\dots,b_d$ in $\Psi\times\{1,\dots,N\}$, the integral over $S$ of
$c_1(\bar L_{b_1})\wedge\dots\wedge c_1(\bar L_{b_d})$ is the integral over $S\setminus\partial S$ of
the mixed Monge--Amp\`ere measure of the singular metrics, the metrics of the $\bar L_b$ multiplied by
$e^{-\varphi_b}$, and that measure is integrable. This comes from the regularizations: the integral
of $c_1(\bar L_{b_1}(t))\wedge\dots\wedge c_1(\bar L_{b_d}(t))$ is the same for every $t$, by Stokes's
theorem, and under the bounds of Definition~\ref{cv-proportionality-analytic-system} it converges, as
$t\to0$, to the integral of the mixed Monge--Amp\`ere measure of the singular metrics.
Expanding the power on the left by multilinearity, the left side is therefore the integral over
$S\setminus\partial S$ of the $d$-th power of
$\sum_{\tau,i}(\lambda_{\tau,i}-\lambda_{\tau,i+1})\,dd^c(\psi_{j,i}\circ g_\tau)$ in the period
charts of the uniformization.

Let $\vartheta_{\tau,i}$ be the local potentials of the metrics of the $\bar L_{\tau,i}$ multiplied by
$e^{-\varphi_{\tau,i}}$. They are smooth on $S\setminus\partial S$, since the metrics of the
$\bar L_b$ and the functions $\varphi_b$ are. Fix a finite holomorphic atlas $(\kappa_k)$ of $S$ and
functions $(\chi_k)$ as in Definition~\ref{cv-period-monge-ampere}(3). For
$c\in\mathbf R^{\Psi\times\{1,\dots,N\}}$ let $\vartheta_c=\sum_{\tau,i}c_{\tau,i}\vartheta_{\tau,i}$,
let $I(c)$ be the integral over $S$ of the Monge--Amp\`ere measure of
$(\vartheta_c,\dots,\vartheta_c)$, and put
\[
 P(c)=\prod_\tau\prod_{1\le a<b\le N}\ \sum_{m=a}^{b-1}c_{\tau,m}.
\]
The Monge--Amp\`ere density is multilinear, so $I(c)$ is the sum over the words $b_1,\dots,b_d$ of
$c_{b_1}\cdots c_{b_d}$ times the integral for $(\vartheta_{b_1},\dots,\vartheta_{b_d})$, which is
finite by the first paragraph. For $c_{\tau,i}=\lambda_{\tau,i}-\lambda_{\tau,i+1}$ this is the left
side, and the inner sums in $P(c)$ telescope to $\lambda_{\tau,a}-\lambda_{\tau,b}$. So it suffices
to find $v_S>0$ with $I(c)=v_S\,P(c)$ for every $c$.

Let $y$ be a point of the image of $\kappa_k$ with $x=\kappa_k^{-1}(y)\notin\partial S$. The
uniformization gives a period chart of type $j$ for the $\vartheta_{\tau,i}$ at $x$
(Definition~\ref{cv-period-chart}). Lemma~\ref{cv-period-pointwise} then gives $\nu_k(y)\ge0$ with
$\operatorname{MA}(\vartheta_c\circ\kappa_k^{-1},\dots,\vartheta_c\circ\kappa_k^{-1})(y)=\nu_k(y)P(c)$
for every $c$. Let $\mathbf 1$ be the vector with all entries $1$, so that
$P(\mathbf 1)=\prod_\tau\prod_{a<b}(b-a)>0$. Since $\partial S$ has measure zero in each chart, the
integrand of $I(c)$ on the image of $\kappa_k$ is almost everywhere $P(c)/P(\mathbf 1)$ times that of
$I(\mathbf 1)$. Hence $I(c)=v_S\,P(c)$ with $v_S=I(\mathbf 1)/P(\mathbf 1)$.

It remains to show that $I(\mathbf 1)>0$. Its integrands
$\chi_k(\kappa_k^{-1}(y))\,\nu_k(y)P(\mathbf 1)$ are nonnegative almost everywhere. By the
uniformization some $x\in S\setminus\partial S$ has a period chart that is a flag immersion at $x$.
Choose $k$ with $\chi_k(x)>0$ and put $y=\kappa_k(x)$. The Monge--Amp\`ere density of
$\vartheta_{\mathbf 1}\circ\kappa_k^{-1}$ is continuous at $y$, since it is a polynomial in the
second derivatives of the $\vartheta_{\tau,i}$, which are smooth near $x$. At $y$ it equals
$\nu_k(y)P(\mathbf 1)>0$ by Lemma~\ref{cv-period-pointwise}. So the integrand over the image of
$\kappa_k$ is positive on a neighbourhood of $y$, its integral is positive, and $I(\mathbf 1)>0$.
\end{proof}
\leanref{Development/AutomorphicSpectrum/Proofs/MetricProofs/UniformisedPotentialIntegralFormulaStatement.lean}{GlobalTheory.T.uniformised_potential_integral_formula}
\leanref{Development/WeilBound/CoreStatements/PluriharmonicCharts/XPerInt.lean}{PositivityTheory.Xgp.t_per_int}
\leanref{Development/AutomorphicSpectrum/Proofs/MetricProofs/MongeAmpereIntegralMultilinearExpansionStatement.lean}{GlobalTheory.T.monge_ampere_integral_multilinear_expansion}
\leanref{Development/WeilBound/CoreStatements/PluriharmonicCharts/XPolyDelta.lean}{PositivityTheory.Xgp.t_poly_delta}

\begin{remark}\label{cv-proportionality-literature}
Theorem~\ref{cv-proportionality-mumford} has the form of the proportionality theorem of Hirzebruch and
Mumford \cite{Mum77}. For the automorphic vector bundles on a toroidal compactification of a locally
symmetric variety $\Gamma\backslash\mathcal D$, with their invariant metrics, Mumford shows that the
Chern numbers are the volume of $\Gamma\backslash\mathcal D$ times the Chern numbers of the
corresponding homogeneous bundles on the compact dual. The analytic systems here are given only by
period charts to the flag domain and by the bounds of
Definition~\ref{cv-proportionality-analytic-system}, and Mumford's theorem is not stated for them.
\end{remark}

\begin{lemma}\label{cv-proportionality-realization}
Let $\Psi$ be a finite nonempty set and $N\ge2$, put $d=|\Psi|N(N-1)/2$, and let $M$ be a flag
model of type $(\Psi,N)$ (Definition~\ref{cv-condition-flag-model}) with a realization of
multiplicity $n$ on an analytic system $S$ of type $(\Psi,N)$
(Definition~\ref{cv-proportionality-analytic-system}). For a weight
$\lambda$ and $s\in\mathbf R$ let $c(\lambda,s)\in\mathbf R^\Sigma$ be given by
$c_{\tau,i}=\lambda_{\tau,i}-\lambda_{\tau,i+1}$, with $\lambda_{\tau,N+1}=0$, and $c_A=s/k_A$.
Then
\[
 n\deg_M\big(\ell(\lambda,s)^{d}\big)=\deg_S\big(c(\lambda,s)\big).
\]
In particular $\deg_M\big(\ell(\lambda,s)^{d}\big)$ depends only on $S$, $n$, $k_A$,
$\lambda$ and $s$.
\end{lemma}
\begin{proof}
Summation by parts gives
$\ell(\lambda)=\sum_{\tau,i}(\lambda_{\tau,i}-\lambda_{\tau,i+1})y_{\tau,i}$ with
$y_{\tau,i}=-\sum_{k\le i}x_{\tau,k}$, and $s\,a$ stands for $(s/k_A)\bar A$. The integral over the
complex points of $X$ of a product of first Chern forms of smooth metrized line bundles depends only
on their multidegrees, by Stokes's theorem. So, by multilinearity, $\delta\deg_M(\ell(\lambda,s)^{d})$
is the integral of $\big(\sum_bc_b\,c_1(\bar y_b)\big)^{d}$ over the complex points
of $X$. The map $\tilde X\to X(\mathbf C)$ is bijective over a dense open set, so this is the
integral over $\tilde X$ of the pullback, which is $\pi^*\big(\sum_bc_b\,c_1(\bar L_b)\big)^{d}$.
Since $\pi$ has degree $\deg\pi$ off a set of measure zero, the integral is $\deg\pi\cdot\deg_S(c)$.
Dividing by $\delta=n\deg\pi$ gives the claim.
\end{proof}
\leanref{Development/AutomorphicSpectrum/Proofs/DivisibilityProofs/RealTowerGeometricAgreementStatement.lean}{GlobalTheory.T.real_tower_geometric_agreement}

\subsection{Tail bounds and convergence}\label{subsection-cv-tail-bounds}

\begin{lemma}\label{cv-tail-shells}
Let $\mu$ be a measure on a measurable space $Y$, let $f$ be a measurable real function on $Y$, and
let $K\ge0$. Suppose that $\mu\{f\ge R\}\le K\log R/R$ for every $R\ge\exp(1)$. Then for every
$R\ge\exp(1)$
\[
 \int_{\{f\ge R\}}\log f\,d\mu\le16K\,\frac{(\log R)^2}{R}.
\]
\end{lemma}
\leanref{ReciprocityTheory/AnalyticForms/LeviForms/TailBound.lean}{ReciprocityTheory.AnForm.P.tail_shells}
\begin{proof}
Let $y\in Y$ with $f(y)\ge R$, and let $N(y)\ge1$ be the number of integers $j\ge0$ with
$f(y)\ge2^jR$. Then $f(y)<2^{N(y)}R$, so $\log f(y)\le\log R+N(y)\log2$. Integrating over
$\{f\ge R\}$, and writing $N(y)$ as the sum over $j\ge0$ of the indicator functions of the sets
$\{f\ge2^jR\}$, gives
\[
 \int_{\{f\ge R\}}\log f\,d\mu\le\log R\cdot\mu\{f\ge R\}+\log2\sum_{j\ge0}\mu\{f\ge2^jR\}.
\]
Since $\log R\ge1\ge\log2$, one has $\log(2^jR)\le(j+1)\log R$, so the hypothesis at $2^jR$ gives
$\mu\{f\ge2^jR\}\le K(j+1)2^{-j}\log R/R$. As $\sum_{j\ge0}(j+1)2^{-j}=4$, the right side is at
most $K(\log R)^2/R+4K\log2\cdot\log R/R$. Since $\log2\le\log R$ and $K\ge0$, the second term is at
most $4K(\log R)^2/R$. So the right side is at most $5K(\log R)^2/R$, which is at most
$16K(\log R)^2/R$ because $K(\log R)^2/R\ge0$.
\end{proof}

\begin{lemma}\label{cv-tail-cutoff}
Let $S$ be a complex manifold, let $\theta$ be a closed real $(1,1)$-form on $S$ with $\theta\ge0$,
let $u$ be a smooth function on $S$ with $u\le0$ and $\theta+dd^cu\ge0$, and let $M\ge0$. Then there
is a smooth function $v$ on $S$ with $\theta+dd^cv\ge0$ and $-M-1\le v\le0$ such that $v=u-1/2$ at
every point where $u\ge-M+1$.
\end{lemma}
\leanref{ReciprocityTheory/AnalyticForms/LeviForms/TailCut.lean}{ReciprocityTheory.AnForm.P.cutoff}
\leanref{ReciprocityTheory/AnalyticForms/PlurisubharmonicFunctions/IsPshLocallyIntegrableStatement.lean}{ReciprocityTheory.AnForm.A.regMax}
\leanref{ReciprocityTheory/AnalyticForms/PlurisubharmonicFunctions/PshRegMax.lean}{ReciprocityTheory.AnForm.A.regMax_thm}
\begin{proof}
Let $\mathcal M$ be Demailly's regularized maximum of two real variables with both parameters equal
to $1/2$ \cite[Chapter~I, \S5]{Dem12}. It is smooth, convex and nondecreasing in each variable, with
$\max(t_1,t_2)\le\mathcal M(t_1,t_2)\le\max(t_1,t_2)+1/2$, with $\mathcal M(t_1,t_2)=t_1$ if
$t_2+1\le t_1$, and with $\mathcal M(t_1+c,t_2+c)=\mathcal M(t_1,t_2)+c$; and
$\mathcal M(\psi_1,\psi_2)$ is plurisubharmonic when $\psi_1$ and $\psi_2$ are. Put
$v=\mathcal M(u,-M)-1/2$, a smooth function. Near each point write $\theta=dd^c\varphi$ with
$\varphi$ smooth. Then $\varphi+v=\mathcal M(\varphi+u,\varphi-M)-1/2$, and $\varphi+u$ and
$\varphi-M$ are plurisubharmonic since $\theta+dd^cu\ge0$ and $\theta\ge0$; so $\theta+dd^cv\ge0$.
Since $u\le0$ and $M\ge0$, one has $-M\le\max(u,-M)\le0$, so $-M-1/2\le v\le0$. Where
$u\ge-M+1$, $\mathcal M(u,-M)=u$ and $v=u-1/2$.
\end{proof}

\begin{situation}\label{cv-tail-situation}
Let $S$ be a compact complex manifold of dimension $d\ge1$. Let $\bar L_\sigma$ be a holomorphic line
bundle on $S$ with a smooth Hermitian metric, and let $\sigma$ be a holomorphic section of
$\bar L_\sigma$ with $\|\sigma\|\le1$ whose divisor is $\sum_km_kW_k$, where the $W_k$ are finitely
many smooth hypersurfaces crossing normally and the $m_k\ge0$ are integers. Let $\partial S$ be the
zero set of $\sigma$, put $r=\exp(1)-2\log\|\sigma\|$ on $S\setminus\partial S$, so that
$r\ge\exp(1)$, and for $R>0$ write $\{r\ge R\}$ for the set of points of $S\setminus\partial S$ at
which $r\ge R$. The operator $dd^c$ is normalized so that a metric whose norm is $|\cdot|e^{-\varphi}$
in a local trivialization has first Chern form $dd^c\varphi$; thus multiplying a metric by $e^{-u}$
adds $dd^cu$ to its first Chern form. For closed real $(1,1)$-forms $\omega_1,\dots,\omega_d$ on $S$,
the product $\omega_1\wedge\dots\wedge\omega_d$ is a signed measure with a smooth density, and a
positive measure if every $\omega_i$ is semipositive.
\end{situation}

\begin{theorem}[Poincar\'e--Lelong]\label{cv-tail-poincare-lelong}
In Situation~\ref{cv-tail-situation}, let $v$ be a smooth function on $S$, let
$\Phi_1,\dots,\Phi_{d-1}$ be closed real $(1,1)$-forms on $S$, and write
$\Phi=\Phi_1\wedge\dots\wedge\Phi_{d-1}$. Then $\log\|\sigma\|$ is integrable against
$\omega_1\wedge\dots\wedge\omega_d$ for all closed real $(1,1)$-forms $\omega_1,\dots,\omega_d$ on
$S$, and
\[
 \int_S\log\|\sigma\|\,dd^cv\wedge\Phi=\sum_km_k\int_{W_k}v\,\Phi
 -\int_Sv\,c_1(\bar L_\sigma)\wedge\Phi.
\]
\end{theorem}
\leanref{ReciprocityTheory/AnalyticForms/PotentialTheory/PLLocalStatement.lean}{ReciprocityTheory.AnForm.M.poincareLelong}
\leanref{ReciprocityTheory/AnalyticForms/PotentialTheory/LogInt.lean}{ReciprocityTheory.AnForm.M.logNormIntegrable}

The integrability holds because near each point $\log\|\sigma\|$ is a smooth function plus a sum of
multiples of $\log|z_i|$ for coordinates $z_i$. The formula is the Lelong--Poincar\'e equation
$dd^c\log\|\sigma\|=\sum_km_k[W_k]-c_1(\bar L_\sigma)$ of currents, paired with the closed form
$v\,\Phi$ \cite[Chapter~III, \S2]{Dem12}.

\begin{lemma}\label{cv-tail-energy}
In Situation~\ref{cv-tail-situation}, let $\theta_1,\dots,\theta_d$ be positive definite closed real
$(1,1)$-forms on $S$. There is a constant $C_E$ with the following property. Let $M\ge0$, and let
$v_1,\dots,v_d$ be smooth functions on $S$ with $\theta_i+dd^cv_i\ge0$ and $-M-1\le v_i\le0$ for
every $i$. Then
\[
 \int_{S\setminus\partial S}r\,(\theta_1+dd^cv_1)\wedge\dots\wedge(\theta_d+dd^cv_d)\le C_E(M+1).
\]
\end{lemma}
\leanref{ReciprocityTheory/AnalyticForms/LeviForms/TailBound.lean}{ReciprocityTheory.AnForm.P.EnergyBoundStatement}
\leanref{Development/AutomorphicSpectrum/Proofs/MetricProofs/TailBoundStatementHolds.lean}{GlobalTheory.CompanionTheory.Ar.Giv.tailEnergyBound}
\leanref{Development/AutomorphicSpectrum/Proofs/MetricProofs/TailBoundStatementHolds.lean}{GlobalTheory.CompanionTheory.Ar.Giv.tailTelescope}
\leanref{Development/AutomorphicSpectrum/Proofs/MetricProofs/TailBoundStatementHolds.lean}{GlobalTheory.CompanionTheory.Ar.Giv.tailEnergyStep}
\begin{proof}
Write $\mu_v=(\theta_1+dd^cv_1)\wedge\dots\wedge(\theta_d+dd^cv_d)$, a positive measure, and
$\mu_0=\theta_1\wedge\dots\wedge\theta_d$. By Stokes's theorem $\mu_v(S)=\mu_0(S)$. The set
$\partial S$ has measure $0$ for $\mu_v$, $r=\exp(1)-2\log\|\sigma\|$, and $\log\|\sigma\|$ is
$\mu_v$-integrable by Theorem~\ref{cv-tail-poincare-lelong}, so
\[
 \int_{S\setminus\partial S}r\,\mu_v=\exp(1)\,\mu_0(S)-2\int_S\log\|\sigma\|\,\mu_v.
\]
It therefore suffices to find $K$, not depending on $M$ or on the $v_i$, with
$\int_S\log\|\sigma\|\,\mu_v\ge\int_S\log\|\sigma\|\,\mu_0-(M+1)K$: then
$C_E=|\exp(1)\mu_0(S)-2\int_S\log\|\sigma\|\,\mu_0|+2|K|$ has the property, since $M+1\ge1$.

For each $j$ choose $C_j\ge0$ with $C_j\theta_j+c_1(\bar L_\sigma)\ge0$; it exists since $\theta_j$
is positive definite and $S$ is compact. Replace the factors one at a time: for $0\le j\le d$ let
$\mu^{(j)}$ be the product of the forms $\theta_i+dd^cv_i$ for $i\le j$ and $\theta_i$ for $i>j$, so
that $\mu^{(0)}=\mu_0$ and $\mu^{(d)}=\mu_v$. Fix $1\le j\le d$, and let $\Phi$ be the product of the
$d-1$ factors of $\mu^{(j)}$ other than the $j$-th, all of them semipositive. By multilinearity and
Theorem~\ref{cv-tail-poincare-lelong},
\[
 \int_S\log\|\sigma\|\,\mu^{(j)}-\int_S\log\|\sigma\|\,\mu^{(j-1)}
 =\int_S\log\|\sigma\|\,dd^cv_j\wedge\Phi
 =\sum_km_k\int_{W_k}v_j\,\Phi-\int_Sv_j\,c_1(\bar L_\sigma)\wedge\Phi.
\]
The restriction of $\Phi$ to $W_k$ is a positive measure and $v_j\ge-M-1$, so the sum is at least
$-(M+1)\sum_km_k\int_{W_k}\Phi$. Write $c_1(\bar L_\sigma)=(C_j\theta_j+c_1(\bar L_\sigma))-C_j\theta_j$.
Since $v_j\le0$ and $(C_j\theta_j+c_1(\bar L_\sigma))\wedge\Phi$ is a positive measure, and since
$-v_j\le M+1$ and $\theta_j\wedge\Phi$ is a positive measure, the last integral is at most
$C_j(M+1)\int_S\theta_j\wedge\Phi$. By Stokes's theorem on $W_k$ and on $S$, the integrals
$\int_{W_k}\Phi$ and $\int_S\theta_j\wedge\Phi$ do not change when the terms $dd^cv_i$ are removed
from the factors of $\Phi$. Hence the step from $j-1$ to $j$ lowers the integral of $\log\|\sigma\|$
by at most $(M+1)K_j$, where
\[
 K_j=\sum_km_k\int_{W_k}\bigwedge_{i\ne j}\theta_i+C_j\int_S\theta_1\wedge\dots\wedge\theta_d,
\]
and $K=K_1+\dots+K_d$ has the required property.
\end{proof}

\begin{lemma}\label{cv-tail-bound}
In Situation~\ref{cv-tail-situation}, let $\theta_1,\dots,\theta_d$ be positive definite closed real
$(1,1)$-forms on $S$. There is a constant $C$ with the following property. Let $C_1\ge0$, let
$u_1,\dots,u_d$ be smooth functions on $S$ with $\theta_i+dd^cu_i\ge0$ on $S$ and
$-C_1\log r\le u_i\le0$ on $S\setminus\partial S$, and let $\mu$ be the measure
$(\theta_1+dd^cu_1)\wedge\dots\wedge(\theta_d+dd^cu_d)$. Then for every $R\ge\exp(1)$
\[
 \mu\big(\partial S\cup\{r\ge R\}\big)\le C(C_1+3)\,\frac{\log R}{R},\qquad
 \int_{\{r\ge R\}}\log r\,d\mu\le C(C_1+3)\,\frac{(\log R)^2}{R}.
\]
\end{lemma}
\leanref{ReciprocityTheory/AnalyticForms/LeviForms/IsRegFamilyInhabitStatement.lean}{ReciprocityTheory.AnForm.P.TailBoundStatement}
\leanref{Development/AutomorphicSpectrum/Proofs/MetricProofs/TailBoundStatementHolds.lean}{GlobalTheory.CompanionTheory.Ar.Giv.tailBound_holds}
\begin{proof}
Let $C_E$ be a constant of Lemma~\ref{cv-tail-energy} for $\theta_1,\dots,\theta_d$; enlarging it, we
may assume $C_E\ge0$. Put $C=16C_E$. Since $\partial S$ has empty interior and the $u_i$ are
continuous, $u_i\le0$ on all of $S$.

Let $R'\ge\exp(1)$ and put $M=C_1\log R'+2$. Let $v_i$ be the function that
Lemma~\ref{cv-tail-cutoff} gives for $\theta_i$, $u_i$ and $M$, and let
$\mu'=(\theta_1+dd^cv_1)\wedge\dots\wedge(\theta_d+dd^cv_d)$. On the open set
$V=(S\setminus\partial S)\setminus\{r\ge R'\}$ one has $u_i\ge-C_1\log R'=-M+2$, so $v_i=u_i-1/2$
there, and $\mu$ and $\mu'$ agree on $V$. By Stokes's theorem both have total mass
$\int_S\theta_1\wedge\dots\wedge\theta_d$, so they also agree on the complement
$\partial S\cup\{r\ge R'\}$ of $V$. The measure $\mu'$ has a smooth density, so
$\mu'(\partial S)=0$. Since $-M-1\le v_i\le0$, Lemma~\ref{cv-tail-energy} gives
\[
 R'\,\mu'\{r\ge R'\}\le\int_{S\setminus\partial S}r\,d\mu'\le C_E(M+1)
 =C_E(C_1\log R'+3)\le C_E(C_1+3)\log R',
\]
the last because $\log R'\ge1$. Hence
$\mu(\partial S\cup\{r\ge R'\})\le C_E(C_1+3)\log R'/R'$ for every $R'\ge\exp(1)$, which gives the
first bound since $C_E\le C$. Lemma~\ref{cv-tail-shells}, applied to the restriction of $\mu$ to
$S\setminus\partial S$, to $f=r$ and to $K=C_E(C_1+3)$, gives the second.
\end{proof}

\begin{lemma}\label{cv-tail-uniform}
Let $\Psi$ be a finite nonempty set and $N\ge2$, put $d=|\Psi|N(N-1)/2$, and let $S$ be an analytic
system of type $(\Psi,N)$ (Definition~\ref{cv-proportionality-analytic-system}), with K\"ahler form
$\omega$ and function $r$; for $R>0$ write $\{r\ge R\}$ for the set of points of
$S\setminus\partial S$ at which $r\ge R$. There are $K$ and $0<t_1\le t_0$ with the following
property. For $i=1,\dots,d$ let $\omega_i$ be either $\omega$ or $c_1(\bar L_b(\tau))+\omega$ for
some $b\in\Sigma$ and $0<\tau\le t_1$, and let $\mu=\omega_1\wedge\dots\wedge\omega_d$. Then the
$\omega_i$ are semipositive, and for every $R\ge\exp(1)$
\[
 \mu\big(\partial S\cup\{r\ge R\}\big)\le K\,\frac{\log R}{R},\qquad
 \int_{\{r\ge R\}}\log r\,d\mu\le K\,\frac{(\log R)^2}{R}.
\]
\end{lemma}
\leanref{Development/AutomorphicSpectrum/Proofs/MetricProofs/MaMeasureEqOfConst.lean}{GlobalTheory.CompanionTheory.Ar.Giv.tail_uniform}
\leanref{Development/AutomorphicSpectrum/Proofs/MetricProofs/FixedMemberBase.lean}{GlobalTheory.CompanionTheory.Ar.Giv.fixedMember_base}
\leanref{Development/AutomorphicSpectrum/Proofs/MetricProofs/FixedMemberBase.lean}{GlobalTheory.CompanionTheory.Ar.Giv.epsT_small}
\begin{proof}
The manifold $S$, the section $\sigma$, the divisor $\partial S$ and the function $r$ of the
analytic system satisfy the hypotheses of Situation~\ref{cv-tail-situation}. Let $C$, $t_0$ and
$\varepsilon(t)$ be as in Definition~\ref{cv-proportionality-analytic-system}. Since
$\varepsilon(t)\to0$ as $t\to0$, there is $0<t_1\le t_0$ with $|\varepsilon(t)|\le1/2$ and
$C\varepsilon(t)\le1$ for $0<t\le t_1$. For such $t$ and $b\in\Sigma$ one has
$c_1(\bar L_b(t))\ge-\varepsilon(t)\omega\ge-\omega/2$, so $c_1(\bar L_b(t))+\omega\ge\omega/2$.
This gives the semipositivity, and shows that $\theta_b=c_1(\bar L_b(t_1))+\omega$ is positive
definite.

Let $b\in\Sigma$ and $0<\tau\le t_1$, and put $u=\varphi_b(\tau)-\varphi_b(t_1)-C\varepsilon(t_1)$, a
smooth function on $S$. Since $dd^cu=c_1(\bar L_b(\tau))-c_1(\bar L_b(t_1))$, one has
$\theta_b+dd^cu=c_1(\bar L_b(\tau))+\omega\ge0$. The bound
$\varphi_b(\tau)\le\varphi_b(t_1)+C\varepsilon(t_1)$ gives $u\le0$. Off $\partial S$ one has
$\varphi_b(\tau)\ge\varphi_b\ge-C\log r$, $\varphi_b(t_1)\le C$ and $C\varepsilon(t_1)\le1$, so
$u\ge-C\log r-C-1\ge-(2C+1)\log r$, as $\log r\ge1$.

Given $\omega_1,\dots,\omega_d$ as in the statement, put $\theta_i=\omega$ and $u_i=0$ if
$\omega_i=\omega$, and $\theta_i=\theta_b$ and $u_i=u$ as above if
$\omega_i=c_1(\bar L_b(\tau))+\omega$. Then $\omega_i=\theta_i+dd^cu_i$, and
Lemma~\ref{cv-tail-bound} applies with $C_1=2C+1$. The tuple $(\theta_1,\dots,\theta_d)$ takes
finitely many values. If $C_{\max}\ge0$ is at least the constants of Lemma~\ref{cv-tail-bound} for
all of them, then $K=(2C+4)C_{\max}$ has the property.
\end{proof}

\begin{lemma}\label{cv-tail-rate}
Let $\Psi$ be a finite nonempty set and $N\ge2$, put $d=|\Psi|N(N-1)/2$, let $S$ be an analytic
system of type $(\Psi,N)$ (Definition~\ref{cv-proportionality-analytic-system}), and put
$\rho(t)=(\log\log(1/t))^2/\log(1/t)$. There are $C'$ and $0<t_1\le t_0$ with the following
property. Let $0<t'\le t\le t_1$, let $a,b_1,\dots,b_d\in\Sigma$, and let
$\tau_1,\dots,\tau_d\in\{t,t'\}$. Then
\[
 \Big|\int_S\big(\varphi_a(t)-\varphi_a(t')\big)\,
 c_1(\bar L_{b_1}(\tau_1))\wedge\dots\wedge c_1(\bar L_{b_d}(\tau_d))\Big|\le C'\rho(t).
\]
\end{lemma}
\leanref{Development/AutomorphicSpectrum/Proofs/MetricProofs/MongeAmpereRegularizationLogLogBoundStatement.lean}{GlobalTheory.T.monge_ampere_regularization_log_log_bound}
\leanref{Development/AutomorphicSpectrum/Proofs/MetricProofs/MongeAmpereRegularizationLogLogBound.lean}{GlobalTheory.CompanionTheory.Ar.PB.signed_rate}
\leanref{Development/AutomorphicSpectrum/Proofs/MetricProofs/ConvergenceRateStatementHoldsStatement.lean}{GlobalTheory.T.convergence_rate_statement_holds}
\leanref{Development/AutomorphicSpectrum/Proofs/MetricProofs/ConvergenceRateStatementHolds.lean}{GlobalTheory.CompanionTheory.Ar.Giv.convergenceRate_holds}
\begin{proof}
Let $\omega$, $C$, $\varepsilon(t)$ and $r$ be as in
Definition~\ref{cv-proportionality-analytic-system}. Take $K$ and $t_1$ from
Lemma~\ref{cv-tail-uniform}, and reduce $t_1$ so that $t_1\le\exp(-\exp(1))$. Let $t'$, $t$, $a$,
the $b_i$ and the $\tau_i$ be as in the statement, and put $g=\varphi_a(t)-\varphi_a(t')$,
$\ell=\log(1/t)$ and $q=(\log\ell)^2/\ell=\rho(t)$. Then $\ell\ge\exp(1)$ and $\log\ell\ge1$, so
$1/\ell\le q$ and $\log\ell/\ell\le q$.

Write each $c_1(\bar L_{b_i}(\tau_i))$ as $(c_1(\bar L_{b_i}(\tau_i))+\omega)-\omega$. By
multilinearity the integral is a sum of $2^d$ terms $\pm\int_Sg\,\mu$, where
$\mu=\omega_1\wedge\dots\wedge\omega_d$ with each $\omega_i$ equal to $c_1(\bar L_{b_i}(\tau_i))+\omega$
or to $\omega$. By Lemma~\ref{cv-tail-uniform} each such $\mu$ is a positive measure, so
$|\int_Sg\,\mu|\le\int_S|g|\,d\mu$. It suffices to bound $\int_S|g|\,d\mu$ by $C''q$ with $C''$
independent of $t$, $t'$ and the choices; then $C'=2^dC''$.

Let $y\in S\setminus\partial S$ and $u\in\{t,t'\}$. If $r(y)<\ell$, then $r(y)\le\log(1/u)$, so
$\varphi_a\le\varphi_a(u)\le\varphi_a+u^{1/3}+\varepsilon(u)$ at $y$. Here
$u^{1/3}\le t^{1/3}=e^{-\ell/3}\le3/\ell$, and $\varepsilon(u)\le C\log\ell/\ell$, because
$\log(1/u)\ge\ell$ and $x\mapsto\log x/x$ decreases on $[\exp(1),\infty)$. Hence
$|g(y)|\le(3+C)q$. If $r(y)\ge\ell$, then $-C\log r\le\varphi_a\le\varphi_a(u)\le C$ at $y$, and
$\log r(y)\ge1$, so $|g(y)|\le2C\log r(y)$.

The measure $\mu$ has a smooth density, so $\mu(\partial S)=0$. By Stokes's theorem $\mu(S)$ is the
integral over $S$ of the product of the forms $c_1(\bar L_{b_i})+\omega$ or $\omega$, which takes
finitely many values; let $M_0$ be the largest. By Lemma~\ref{cv-tail-uniform} with $R=\ell$,
\[
 \int_S|g|\,d\mu\le(3+C)q\,\mu(S)+2C\int_{\{r\ge\ell\}}\log r\,d\mu\le\big((3+C)M_0+2CK\big)q.
\]
\end{proof}

\begin{lemma}\label{cv-condition-convergence}
Let $\Psi$ be a finite nonempty set and $N\ge2$, put $d=|\Psi|N(N-1)/2$, let $S$ be an analytic
system of type $(\Psi,N)$ (Definition~\ref{cv-proportionality-analytic-system}), let
$(\cdot\mid\cdot)$ be a full arithmetic intersection system (Definition~\ref{cv-origin-intersection-full}), and
let $b\ge0$. There are $C\ge0$
and $s_0>0$ with the following property. Let $M$ be a flag model of type $(\Psi,N)$ with a realization on $S$. Then the arithmetic
intersection numbers $\mathrm{PP}_M(F)$ exist for all $F\in\mathbf R[x,a]$
(Definition~\ref{cv-condition-flag-model}), and
\[
 \Big|\mathrm{PP}^{t}_M\big(\ell(\lambda,s')^{d+1}\big)
 -\mathrm{PP}_M\big(\ell(\lambda,s')^{d+1}\big)\Big|\le Cs,\qquad t=e^{-1/s^2},
\]
for $0<s\le s_0$ and all weights $\lambda$ and $s'\in\mathbf R$ with
$\sum_{\tau,i}|\lambda_{\tau,i}-\lambda_{\tau,i+1}|+|s'|/k_A\le b$, where $\lambda_{\tau,N+1}=0$.
\end{lemma}
\begin{proof}
Let $n$ be the multiplicity of the realization and $\pi$ its map to $S$. As in the proof of
Lemma~\ref{cv-proportionality-realization}, multilinearity gives
$n\,\mathrm{PP}^t_M(\ell(\lambda,s')^{d+1})=(\deg\pi)^{-1}\big(\bar E(t)^{d+1}\mid X\big)$
with $\bar E(t)=\sum_bc_b\bar y_b(t)$ and $c=c(\lambda,s')$, so it suffices to bound the arithmetic
intersection numbers $\nu(t)=(\deg\pi)^{-1}(\bar y_{b_0}(t)\cdots\bar y_{b_d}(t)\mid X)$ for each word
$b_0,\dots,b_d$ in $\Sigma$.

Let $0<t'\le t$. Change the metrics from $t$ to $t'$ one factor at a time. By the recursion of the
intersection system, each change alters the arithmetic intersection number by the integral over the
complex points of $X$ of $(\varphi_b(t)-\varphi_b(t'))\circ\pi$ against a product of first Chern forms
of the other bundles, each at $t$ or at $t'$. Pulled back to $\tilde X$ and pushed down by $\pi$, this
is $\deg\pi$ times an integral over $S$ of $\varphi_b(t)-\varphi_b(t')$ against the Monge--Amp\`ere measure
of the metrics of the bundles $\bar L_{b'}(t)$ or $\bar L_{b'}(t')$. By
Lemma~\ref{cv-tail-rate}, these integrals are at most $C'\rho(t)$ in absolute value, with
$\rho(t)=(\log\log(1/t))^2/\log(1/t)$, for $t\le t_1$, where $C'$ and $t_1>0$ depend only on $S$.
Hence $|\nu(t)-\nu(t')|\le(d+1)C'\rho(t)$ for $t'\le t\le t_1$. The numbers $\nu(t)$ are therefore
Cauchy as $t\to0$, their limits exist, and $|\nu(t)-\lim\nu|\le(d+1)C'\rho(t)$.

For $t=e^{-1/s^2}$ and $0<s\le1$ one has $\rho(t)=4s^2\log^2(1/s)\le4s$. Expanding $\bar E(t)$ by
multilinearity and using $\sum_b|c_b|\le b$ gives
$n\,|\mathrm{PP}^t_M(\ell(\lambda,s')^{d+1})-\mathrm{PP}_M(\ell(\lambda,s')^{d+1})|\le b^{d+1}(d+1)4C's$,
which is the claim with $C=4(d+1)b^{d+1}C'$, since $n\ge1$, and $s_0=1$ reduced so that
$e^{-1/s_0^2}\le t_1$.
\end{proof}
\leanref{Development/AutomorphicSpectrum/Statements/GeometryStatements/AgreementPos.lean}{GlobalTheory.CompanionTheory.Ar.PB.AgreementPos}
\leanref{Development/AutomorphicSpectrum/Proofs/MetricProofs/AgreementPositiveStatement.lean}{GlobalTheory.T.agreement_positive}
\leanref{Development/AutomorphicSpectrum/Proofs/MetricProofs/AgreementPositive.lean}{GlobalTheory.CompanionTheory.Ar.PB.agree_cauchy}
\leanref{Development/AutomorphicSpectrum/Proofs/DivisibilityProofs/RealTowerArithAgreementStatement.lean}{GlobalTheory.T.real_tower_arith_agreement}

\begin{lemma}\label{cv-condition-comparison}
Let $\Psi$ be a finite nonempty set and $N\ge2$, let $S$ be an analytic system of type $(\Psi,N)$
(Definition~\ref{cv-proportionality-analytic-system}), let $\bar L_H$ be a smooth Hermitian line
bundle on $S$, and let $(\cdot\mid\cdot)$ be a full arithmetic intersection system
(Definition~\ref{cv-origin-intersection-full}). There is $C_2\ge0$ with the following property. Let $X$ be an arithmetic subvariety of dimension $d+1$,
where $d=|\Psi|N(N-1)/2$, of a product of projective spaces, and let $\bar y_b(t)$ and
$\bar y'_b(t)$ ($b\in\Sigma$, $0<t\le t_0$) be the bundles of two realizations on $S$ of flag
models whose arithmetic variety is $X$, whose maps to $S$ have the same degree $\deg\pi$. Let $T$
be a finite set of primes, let $\alpha_p\ge0$ for $p\in T$, and let $\bar H$ be a smooth metrized
line bundle on $X$ whose first Chern form on the manifold $\tilde X$ of the first realization is the
pullback of that of $\bar L_H$. Suppose that:
\begin{enumerate}
\item for every $b$ and $t$ there are multihomogeneous $f,g$ not in the ideal of $X$ with
  $\deg f-\deg g$ equal to the multidegree of $\bar y_b(t)-\bar y'_b(t)$ and
  $|f|\,h'_b(t)=|g|\,h_b(t)$ on the complex points, where $h_b(t)$ and $h'_b(t)$ are the metrics,
  whose divisors on $X$ agree away from the fibres over $T$ and whose multiplicities at a
  component of the fibre over $p\in T$ differ by at most $\alpha_p$ times that of $p$;
\item $\bar H$, $\bar y_b+\bar H$ and $\bar y'_b+\bar H$ are nef on the fibres over the primes in
  $T$, for every $b$.
\end{enumerate}
Then for every word $b_0,\dots,b_d$ in $\Sigma$ and $0<t\le t_0$
\[
 (\deg\pi)^{-1}\Big|\big(\bar y_{b_0}(t)\cdots\bar y_{b_d}(t)\mid X\big)
 -\big(\bar y'_{b_0}(t)\cdots\bar y'_{b_d}(t)\mid X\big)\Big|\le C_2\sum_{p\in T}\alpha_p\log p.
\]
Moreover, the coefficients of the polynomials in $(\lambda,u)$
\[
 Q(\lambda,u)=\lim_{t\to0^+}(\deg\pi)^{-1}\Big(\big(\textstyle\sum_bc_b\bar y_b(t)\big)^{d+1}\,\Big|\,X\Big),
 \qquad c_{\tau,i}=\lambda_{\tau,i}-\lambda_{\tau,i+1},\quad c_A=u,
\]
and $Q'$, defined in the same way with the $\bar y'_b(t)$, differ by at most
$C_2'\sum_{p\in T}\alpha_p\log p$, where $C_2'$ depends only on $C_2$, $|\Psi|$ and $N$.
\end{lemma}
\begin{proof}
Replace the factors $\bar y_{b_k}(t)$ by $\bar y'_{b_k}(t)$ one at a time. By multilinearity each
step is a difference $(\bar E_1\cdots\bar E_d\,\bar y_b(t)\mid X)-(\bar E_1\cdots\bar E_d\,\bar y'_b(t)\mid X)$,
and writing each $\bar E_k$ as $(\bar E_k+\bar H)-\bar H$ expands it into at most $2^d$ such
differences in which every $\bar E_k$ is $\bar H$ or of the form $\bar y_{b'}(t)+\bar H$ or
$\bar y'_{b'}(t)+\bar H$, all nef on the fibres over $T$. For such factors the recursion of the
intersection system, applied to the section $f/g$, writes the difference as
$\sum_{p\in T}\log p$ times the degree of the $\bar E_k$ on the difference of the vertical parts of
$\operatorname{div}f$ and $\operatorname{div}g$ over $p$, which is at most $\alpha_p$ times their
degree on the fibre over $p$, by nefness; the archimedean terms cancel since
$|f|\,h'_b(t)=|g|\,h_b(t)$. The degree on a fibre equals the degree on the generic fibre, an integral over
the complex points of $X$ of first Chern forms; pulled back to $\tilde X$ it is $\deg\pi$ times an integral
over $S$ of products of first Chern forms of the $\bar L_{b'}$ and $\bar L_H$, bounded by a constant
depending only on $S$ and $\bar L_H$. Summing gives the first bound. The limits as $t\to0$ exist by
Lemma~\ref{cv-condition-convergence}. The coefficient of $\prod_bc_b^{e_b}$ in the limit
polynomial in the variables $c_b$ is a multinomial coefficient times the limit for one word, and the
substitution $c_{\tau,i}=\lambda_{\tau,i}-\lambda_{\tau,i+1}$, $c_A=u$ is a fixed linear change of
variables, which gives the second bound.
\end{proof}
\leanref{Development/AutomorphicSpectrum/Statements/MetricStatements/AgreementIffGuarded.lean}{GlobalTheory.CompanionTheory.Ar.PB.Coefficients}
\leanref{Development/AutomorphicSpectrum/Proofs/MetricProofs/CoefficientsStatement.lean}{GlobalTheory.T.coefficients'}
\leanref{Development/AutomorphicSpectrum/Proofs/MetricProofs/RealisationHeightIndependence.lean}{GlobalTheory.CompanionTheory.Ar.PB.abs_num_sub_num_le}

\begin{lemma}\label{cv-nef-positivity}
Let $\Psi$ be a finite nonempty set and $N\ge2$, let $S$ be an analytic system of type $(\Psi,N)$
(Definition~\ref{cv-proportionality-analytic-system}), and let $\Sigma_0\subset\Sigma$ be a set of
indices $b$ for which the first Chern form of $\bar L_b(t)$ is $0$ for every $t$. Let $M$ be a flag
model of type $(\Psi,N)$ (Definition~\ref{cv-condition-flag-model}) with a realization on $S$ whose
map $\tilde X\to X(\mathbf C)$ is bijective and immersive everywhere, and let $\bar y_b$
($b\in\Sigma$) be the bundles of the realization. Let $n\in\mathbf Z^\Sigma$ with $n_b\ge0$ for
$b\notin\Sigma_0$. Then the line bundle $\sum_bn_by_b$ has nonnegative degree on every curve in the
generic fibre of $X$.
\end{lemma}
\leanref{Development/AutomorphicSpectrum/Statements/MetricStatements/AgreementIffGuarded.lean}{GlobalTheory.CompanionTheory.Ar.PB.NefPol}
\leanref{Development/AutomorphicSpectrum/Proofs/MetricProofs/NefPolarization.lean}{GlobalTheory.CompanionTheory.Ar.PB.nefPol}
\leanref{Development/WeilBound/CoreStatements/PluriharmonicCharts/XNefCurve.lean}{PositivityTheory.Xgp.t_nef_curve}
\begin{proof}
Let $W$ be a curve in the generic fibre of $X$. The degree on $W$ is additive in the line bundle, so
it suffices to show that $\deg_W\bar y_b\ge0$ for every $b$ and that $\deg_W\bar y_b=0$ for
$b\in\Sigma_0$. Let $\tilde W\to W(\mathbf C)$ be a resolution of the complex points of $W$, a compact
Riemann surface. Since $\tilde X\to X(\mathbf C)$ is bijective and immersive, this map lifts to a
holomorphic map $h:\tilde W\to\tilde X$; put $g=\pi\circ h$. The degree of $\bar y_b$ on $W$ is the
integral over $\tilde W$ of the pullback of its first Chern form, and by the realization this is
$\int_{\tilde W}g^*c_1(\bar L_b)$. For $0<t\le t_0$ the metrics of $\bar L_b(t)$ and $\bar L_b$ differ
by the factor $e^{-\varphi_b(t)}$, with $\varphi_b(t)$ smooth, so by Stokes's theorem on $\tilde W$
\[
 \deg_W\bar y_b=\int_{\tilde W}g^*c_1(\bar L_b(t))\ge-\varepsilon(t)\int_{\tilde W}g^*\omega,
\]
by the lower bound for $c_1(\bar L_b(t))$ in Definition~\ref{cv-proportionality-analytic-system}. As
$\varepsilon(t)\to0$ when $t\to0$, $\deg_W\bar y_b\ge0$. If $b\in\Sigma_0$, the same identity at
$t=t_0$ gives $\deg_W\bar y_b=0$.
\end{proof}
% Lean: nefPol takes the degree from GlobalTheory.T.geometric_degree_nonnegative_psh
% (= PositivityTheory.Xgp.t_nef_curve: exists_anLogRes for the resolution, t_res_lift for the lift,
% mass_sys for the two signs) and GlobalTheory.T.generic_intersection_equals_points_integral.

\subsection{Towers and the volume theorem}\label{subsection-cv-towers-volume-theorem}

\begin{definition}\label{cv-origin-hodge-component}
Let $E$ be a CM field with a host structure as in Theorem~\ref{cv-flat-degree}, with its prime $d$
and $\sqrt{-d}\in E$, let $\theta$ be a square root of $-d$ in $\mathbf C$, and let $\Psi$ be the set
of embeddings $\tau:E\to\mathbf C$ with $\tau(\sqrt{-d})=\theta$. Let a PEL datum of signature $(p,q)$
over $E$, a level, its moduli scheme $X$ over a coefficient field $F_c$
(Situation~\ref{cv-function-situation}) and a model $(\mathcal X,\mathcal A,j)$ of $X$ over the ring
of integers $\mathcal O_c$ of $F_c$ be given, defined as in Definition~\ref{cv-model} with
$\mathcal O_c$ in place of $O$. Put $N=p+q$. The \emph{labels} are the embeddings $\tau$ and
$\bar\tau$ of $\mathcal O_E$ into $F_c$, for $\tau\in\Psi$, where $\bar\tau$ is $\tau$ composed with
complex conjugation; put $r_\tau=p$ and $r_{\bar\tau}=q$. Let $S$ be a finite set of primes, and let
$y=(y_{k,i})$ be a point of a product $\prod_{k\in K}\mathbf P^{I_k}$ with values in a field $\Omega$
that is its function field (Definition~\ref{lit-multiprojective}(4)). A \emph{chart} is a choice
$\kappa=(\kappa_k)_{k\in K}$ of indices $\kappa_k\in I_k$ with $y_{k,\kappa_k}\neq0$ for every $k$.
\begin{enumerate}
\item A \emph{Hodge component} of $\mathcal X$ at $S$ with coordinates $y$ consists of: a normal
  integral scheme $Z$ with function field $\Omega$, covered by affine open subschemes $Z_\kappa$, one
  for each chart $\kappa$, where $Z_\kappa$ is the set of points at which the ratios
  $y_{k,i}/y_{k,\kappa_k}$ are regular and has ring $\mathbf Z[y_{k,i}/y_{k,\kappa_k}:k\in K,\ i\in I_k]$;
  a morphism
  $f:Z\to\mathcal X$ over $\mathcal O_c$ that maps the generic point of $Z$ to the generic point of an
  irreducible component of $\mathcal X$, whose residue field $\Omega_0$ it embeds in $\Omega$; a
  semi-abelian scheme $\mathcal A_Z$ over $Z$ with an action of $\mathcal O_E$ and a symmetric cubical
  line bundle, identified with $f^*\mathcal A$ with these structures; and an integer $R_p\ge4$ for
  each $p\in S$, such that every point of order dividing $p^{R_p}$ of $\mathcal A_Z$ over a field
  extension of $\Omega$ is defined over $\Omega$. Let $\mathrm{Lie}$ be the Lie algebra of
  $\mathcal A_Z$ at the generic point and, for a label $s$, let $\mathrm{Lie}_s$ be the subspace on
  which each $b\in\mathcal O_E$ acts by $s(b)$. For $p\in S$, the \emph{Hodge--Tate periods} are the
  invariant differentials $d\log$ of the functions on $\mathcal A_Z$ given by the Weil pairing with
  the points of order dividing $p^{R_p}$. At a point of $Z$, let $H_s$ be the module of
  $e\in\mathrm{Lie}_s$ at which all Hodge--Tate periods take regular values, and put
  $\mathcal V^p_{\bar\tau}=H_{\bar\tau}$ and
  $\mathcal V^p_\tau=\{e\in\mathrm{Lie}_\tau:\tau(\delta)e\in H_\tau$ for every $\delta$ in the
  different of $E\}$. The condition is that for each label $s$ the modules $\mathcal V_s$, given at a
  point over no prime of $S$ by the vectors of $\mathrm{Lie}_s$ that lie in the Lie algebra of
  $\mathcal A_Z$ at that point, and at a point over primes of $S$ by the intersection of the
  $\mathcal V^p_s$ over those primes $p$, form a vector bundle of rank $r_s$ on $Z$.
\item A \emph{flag} of degree $\deg\ge1$ of the Hodge component is a family of bases
  $b_s=(b_{s,1},\dots,b_{s,r_s})$ of the $\mathrm{Lie}_s$ such that near every point of $Z$ the vector
  bundle $\mathcal V_s$ has a basis $b'$ with
  $\langle b'_1,\dots,b'_i\rangle=\langle b_{s,1},\dots,b_{s,i}\rangle$ for every $i$, and such that
  there are bases $(c_{s,l})_l$ over $\Omega_0$ of the eigenspaces of the Lie algebra of $\mathcal A$
  at the generic point of $f(Z)$ and elements $u_{s,lk}\in\Omega$ for $l>k$, algebraically
  independent over $\Omega_0$, with $\Omega$ of degree $\deg$ over $\Omega_0(u)$ and
  $\langle b_{s,1},\dots,b_{s,i}\rangle$ spanned by the vectors $c_{s,k}+\sum_{l>k}u_{s,lk}c_{s,l}$
  with $k\le i$.
\item For $(\tau,i)\in\Psi\times\{1,\dots,N\}$, the \emph{root line} $\Lambda_{\tau,i}$ of a flag is
  the line bundle on $Z$ given by the $i$-th graded piece of the flag of $\mathcal V_\tau$ if $i\le p$,
  and by the dual of the $(N+1-i)$-th graded piece of the flag of $\mathcal V_{\bar\tau}$ if $i>p$;
  it is a module in $\Omega$ through the coefficient of $b_{\tau,i}$, or the inverse of the
  coefficient of $b_{\bar\tau,N+1-i}$. The flag has \emph{root lines of multidegrees}
  $a_{\tau,i}\in\mathbf Z^K$ \emph{with generators} $e_{\tau,i}\in\Omega^\times$ if, for every chart
  $\kappa$, $\Lambda_{\tau,i}$ is generated at every point of $Z_\kappa$ by
  $e_{\tau,i}\prod_ky_{k,\kappa_k}^{a_{\tau,i,k}}$.
\item Metrized line bundles $\bar\sigma_{\tau,i}=(a_{\tau,i},h_{\tau,i})$
  (Definition~\ref{lit-multiprojective}(2)) are \emph{Hodge metrics} on root lines with generators
  $e_{\tau,i}$ if at every complex point of $Z$ over which $\mathcal A_Z$ is abelian, with
  coordinates $w$, and for every chart $\kappa$ with all $w_{k,\kappa_k}\neq0$, the generator
  $e_{\tau,i}\prod_ky_{k,\kappa_k}^{a_{\tau,i,k}}$ of $\Lambda_{\tau,i}$ has norm
  $\prod_k|w_{k,\kappa_k}|^{a_{\tau,i,k}}/h_{\tau,i}(w)$ for the Hodge form of the abelian variety at
  that point. Here the $i$-th graded piece of a flag in a space with a positive definite Hermitian
  form gives the class of the $i$-th vector of a basis adapted to the flag the norm
  $(\det\Gamma_i/\det\Gamma_{i-1})^{1/2}$, where $\Gamma_i$ is the Gram matrix of the first $i$
  vectors, and a dual has the dual norm.
\end{enumerate}
\end{definition}
\leanref{Development/AutomorphicSpectrum/Definitions/GlobalGeometry/IsHodgeLine.lean}{GlobalTheory.CompanionTheory.Ar.HodgeCpt}
\leanref{Development/AutomorphicSpectrum/Definitions/GlobalGeometry/IsHodgeLine.lean}{GlobalTheory.CompanionTheory.Ar.HodgeCpt.Flag}
\leanref{Development/AutomorphicSpectrum/Definitions/GlobalGeometry/IsGenericFlag.lean}{GlobalTheory.CompanionTheory.Ar.IsGenericFlag}
\leanref{Development/AutomorphicSpectrum/Definitions/GlobalGeometry/IsHodgeLine.lean}{GlobalTheory.CompanionTheory.Ar.HodgeCpt.Flag.RootLines}
\leanref{Development/AutomorphicSpectrum/Definitions/GlobalGeometry/IsHodgeLine.lean}{GlobalTheory.CompanionTheory.Ar.HodgeCpt.Flag.IsHodgeLine}
\leanref{Development/AutomorphicSpectrum/Definitions/GlobalGeometry/CoordsAt.lean}{GlobalTheory.CompanionTheory.Ar.CoordTie}

\begin{definition}\label{cv-origin-tower}
Let $\Psi$ be a finite nonempty set and let $N\ge2$.
\begin{enumerate}
\item A \emph{tower} of type $(\Psi,N)$ consists of an index $j$ with $1\le j<N$; a nonempty finite
  set of \emph{components}; a finite set of primes containing $2$, the \emph{bad primes}; a weight
  $\nu$ with rational entries (Definition~\ref{cv-reconstruction-functionals}); and for each integer
  $B$ and each component $c$ a flag model $M(B,c)$ of type $(\Psi,N)$
  (Definition~\ref{cv-condition-flag-model}). For an integer $B$, $S_B$ is the set of primes $p\le B$
  together with the bad primes.
\item In the setting of Definition~\ref{cv-origin-hodge-component}, with $\Psi$ the set of
  embeddings there and $N=p+q$, an \emph{origin} of a tower $T$ of type $(\Psi,N)$ in $\mathcal X$
  consists of the following data.
  \begin{enumerate}
  \item[(a)] The index of $T$ is $j=p$, and the bad primes are the primes dividing $2d_EM'$, where
    $d_E$ is the discriminant of $E$ and $M'$ is the least common multiple of the two integers of the
    level.
  \item[(b)] For each component $c$: an analytic system $S_c$ of type $(\Psi,N)$
    (Definition~\ref{cv-proportionality-analytic-system}) with $\varphi_A=\varphi_A(t)=0$ whose
    bundle $\bar L_A$ has a positive metric; a point $\xi_c$ of $\mathcal X$ with values in a field
    $\Omega_c$, over the generic point $\eta_c$ of an irreducible component of $\mathcal X$, such that
    $c\mapsto\eta_c$ is a bijection from the components of $T$ to the irreducible components of
    $\mathcal X$; vectors $b^c_{s,i}$, for the labels $s$ and $i\le r_s$, of the Lie algebra of
    $\mathcal A$ at $\xi_c$; elements $y^c_0,\dots,y^c_m$ of $\Omega_c$ whose ratios generate the field
    $\Omega_c$; a function $h_c$ on $\mathbf C^{m+1}$; a map
    $\mathrm{pt}_c:S_c\to\mathbf C^{m+1}\setminus\{0\}$ such that $\mathrm{pt}_c(x)$ and
    $\mathrm{pt}_c(x')$ are proportional only if $x=x'$; integers $k_{A,c},q_c\ge1$ such that the
    numbers $n^c_{\tau,i}=k_{A,c}q_c\nu_{\tau,i}$ are integers; and an element $s_c\neq0$ of
    $\Omega_c$.
  \item[(c)] For each integer $B$ and each component $c$, with $M=M(B,c)$, its point $y$ with values
    in $\Omega$, its arithmetic variety $X_M$, the multidegrees $a_{\tau,i}$ of its bundles
    $\bar x_{\tau,i}(t)$, and its $\bar A$, $k_A$ and $\delta$: a Hodge component $Z$ of $\mathcal X$
    at $S_B$ with coordinates $y$ (Definition~\ref{cv-origin-hodge-component}); a flag of $Z$ of some
    degree $\deg$, with root lines of multidegrees $a_{\tau,i}$ with generators $e_{\tau,i}$, and
    Hodge metrics $\bar\sigma_{\tau,i}$ on them; a field embedding $g:\Omega_c\to\Omega$ of degree
    $\deg$ through which the generic point of $Z$ maps to $\xi_c$, and which carries each $b^c_{s,i}$
    to $b_{s,i}$ and the elements of $\Omega_c$ algebraic over $\mathbf Q$ into $F_c$; an index
    $k_0\in K$ with a bijection $i\mapsto i'$ from $\{0,\dots,m\}$ to $I_{k_0}$ such that
    $y_{k_0,i'}=g(y^c_i)$; and a compact complex manifold $\tilde X$ with a holomorphic map
    $x\mapsto w(x)$ to the complex points of $X_M$, bijective and immersive off the zero set of a
    multihomogeneous ideal $J$ not contained in the ideal of $X_M$, a holomorphic map
    $\pi:\tilde X\to S_c$ of degree $\deg$ off a closed set of measure zero, and smooth metrized line
    bundles $\bar y_b$ on $X_M$, for $b\in\Sigma$, whose first Chern forms on $\tilde X$ are the
    pullbacks by $\pi$ of those of the $\bar L_b$, with $\bar y_b(t)$ as in
    Definition~\ref{cv-proportionality-analytic-system}(3). These satisfy:
    \begin{enumerate}
    \item[(i)] $\delta=[F_c:\mathbf Q]\deg$ and $k_A=k_{A,c}$;
    \item[(ii)] the multidegree of $\bar A$ is $1$ at $k_0$ and $0$ at the other indices, and the
      metric of $\bar A$ at a point $w$ is $h_c\big((w_{k_0,i'})_i\big)$;
    \item[(iii)] at every point $P$ of $Z$ over no prime, if $y_{k_0,i'}\neq0$, the ratios
      $y_{k_0,l}/y_{k_0,i'}$ are regular at $P$ and $t$ generates the line bundle
      $\bigotimes_{\tau,l}\Lambda_{\tau,l}^{-n^c_{\tau,l}}$ at $P$, then
      $g(s_c)\,y_{k_0,i'}^{q_c}/t$ is regular at $P$, and it is not a unit at $P$ if
      $\mathcal A_Z$ is not abelian at $P$;
    \item[(iv)] for $x\in\tilde X$ the coordinates $(w(x)_{k_0,i'})_i$ are proportional to
      $\mathrm{pt}_c(\pi(x))$, and a complex point of $Z$ with coordinates $w(x)$ is one over which
      $\mathcal A_Z$ is abelian if and only if $\pi(x)\notin\partial S_c$;
    \item[(v)] $\bar y_{\tau,i}$ has multidegree $-\sum_{k\le i}a_{\tau,k}$; at the points $w(x)$
      with $\pi(x)\notin\partial S_c$, the metric of $\bar y_{\tau,i}$ multiplied by
      $e^{-\varphi_{\tau,i}(\pi(x))}$ is that of $-\sum_{k\le i}\bar\sigma_{\tau,k}$; for
      $0<t\le t_0$, $\bar y_{\tau,i}(t)$ is $-\sum_{k\le i}\bar x_{\tau,k}(t)$; and $\bar y_A=\bar A$.
    \end{enumerate}
  \end{enumerate}
\item The origin is \emph{good} if moreover:
  \begin{enumerate}
  \item[(a)] for every $B$, every component $c$ and every $p\in S_B$, the Hodge component at $(B,c)$
    has $R_p=4$, and the module spanned by the Hodge--Tate periods of its points of order dividing
    $p^4$, and its images in the duals of the $\mathrm{Lie}_s$, are vector bundles of ranks
    $\dim\mathrm{Lie}$ and $r_s$ on the fibre of $Z$ over $p$;
  \item[(b)] the maps $\tilde X\to X_M(\mathbf C)$ at all $(B,c)$ are bijective and immersive
    everywhere, and the generic fibres of the arithmetic varieties of the $M(B,c)$ are smooth;
  \item[(c)] for every component $c$, every $\tau\in\Psi$ and every $t$, the regularized metric of
    $\bar L_{\tau,N}(t)$ on $S_c$ has pluriharmonic local potentials;
  \item[(d)] for every component $c$, the map $S_c\to\mathbf P^m(\mathbf C)$,
    $x\mapsto[\mathrm{pt}_c(x)]$, is a holomorphic immersion;
  \item[(e)] for every component $c$ there are a flag model $R_c$ of type $(\Psi,N)$ and data as in
    (2c) for $R_c$ in place of $M(B,c)$ and the empty set in place of $S_B$, with $\deg=1$ and with
    $\tilde X\to X_{R_c}(\mathbf C)$ bijective and immersive everywhere, such that the multidegree
    $a_A$ of the bundle $\bar A$ of $R_c$ is ample on the generic fibre: for some $m'\ge1$, finitely
    many elements of $\tilde H^0(X_{R_c},\bar L)$, with $\bar L$ of multidegree $m'a_A-(1,\dots,1)$,
    have no common zero on the generic fibre of the arithmetic variety $X_{R_c}$ of $R_c$
    (Definition~\ref{lit-multiprojective}(4)).
  \end{enumerate}
\end{enumerate}
\end{definition}
\leanref{Development/AutomorphicSpectrum/Definitions/DivisibilityTowers/Proportional.lean}{GlobalTheory.CompanionTheory.Ar.Top.RootTower}
\leanref{Development/AutomorphicSpectrum/Definitions/GlobalGeometry/TowerOrigin.lean}{GlobalTheory.CompanionTheory.Ar.TowerOrigin}
\leanref{Development/AutomorphicSpectrum/Definitions/GlobalGeometry/TowerOrigin.lean}{GlobalTheory.CompanionTheory.Ar.FixedData}
\leanref{Development/AutomorphicSpectrum/Definitions/GlobalGeometry/TowerOrigin.lean}{GlobalTheory.CompanionTheory.Ar.RootOrigin}
\leanref{Development/AutomorphicSpectrum/Definitions/DivisibilityTowers/GoodOrigin.lean}{GlobalTheory.CompanionTheory.Ar.Top.GoodOrigin}
\leanref{Development/AutomorphicSpectrum/Statements/TowerExistence/GoodTowerExistsX.lean}{GlobalTheory.CompanionTheory.Ar.Top.GoodOrigin.Extra}
\leanref{Development/AutomorphicSpectrum/Definitions/GlobalGeometry/SplitComps.lean}{GlobalTheory.CompanionTheory.Ar.HodgeCpt.OmVB}
\leanref{Development/AutomorphicSpectrum/Definitions/DivisibilityTowers/RefLinked.lean}{GlobalTheory.CompanionTheory.Ar.Top.RealTower.Smooth}
\leanref{Development/AutomorphicSpectrum/Definitions/DivisibilityTowers/RefLinked.lean}{GlobalTheory.CompanionTheory.Ar.Top.RealTower.Flat}

\begin{lemma}\label{cv-origin-realization}
Let $T$ be a tower of type $(\Psi,N)$ with an origin in $\mathcal X$
(Definition~\ref{cv-origin-tower}), and put $n=[F_c:\mathbf Q]$. For every integer $B$ and every
component $c$, the manifold $\tilde X$, the map $\pi$ and the bundles $\bar y_b$ of the origin at
$(B,c)$ form a realization of $M(B,c)$ on $S_c$ with multiplicity $n$
(Definition~\ref{cv-proportionality-analytic-system}(3)), and the integer $k_A$ of $M(B,c)$ is
$k_{A,c}$, which does not depend on $B$.
\end{lemma}
\leanref{Development/AutomorphicSpectrum/Definitions/DivisibilityTowers/GoodOrigin.lean}{GlobalTheory.CompanionTheory.Ar.Top.realOf}
\leanref{Development/AutomorphicSpectrum/Definitions/DivisibilityTowers/RefLinked.lean}{GlobalTheory.CompanionTheory.Ar.Top.RealTie}
\leanref{Development/AutomorphicSpectrum/Definitions/DivisibilityTowers/RefLinked.lean}{GlobalTheory.CompanionTheory.Ar.Top.RealTower}
\begin{proof}
The map $x\mapsto w(x)$, the map $\pi$ of degree $\deg$ and the bundles $\bar y_b$, with
$\bar y_A=\bar A$ by Definition~\ref{cv-origin-tower}(2c)(v), have the properties that
Definition~\ref{cv-proportionality-analytic-system}(3) asks of a realization. By (2c)(i),
$\delta=n\deg=n\deg\pi$. By (2c)(v), $\bar y_{\tau,i}(t)=-\sum_{k\le i}\bar x_{\tau,k}(t)$ for
$0<t\le t_0$, and $\bar y_A(t)=\bar y_A=\bar A$ since $\varphi_A(t)=0$ by (2b). Finally
$k_A=k_{A,c}$ by (2c)(i).
\end{proof}

\begin{lemma}\label{cv-origin-effective}
Let $T$ be a tower of type $(\Psi,N)$ with an origin in $\mathcal X$
(Definition~\ref{cv-origin-tower}), let $B$ be an integer and $c$ a component, put $M=M(B,c)$, and let
$y$, $\Omega$, $X_M$, $a_{\tau,i}$, $\bar A$ and $k_A$ be the point, its function field, the arithmetic
variety, the multidegrees of the bundles $\bar x_{\tau,i}(t)$, the bundle and the integer of $M$; let
$a_A$ be the multidegree of $\bar A$. Then there are an integer $q\ge1$, such that the $q\nu_{\tau,i}$ and
$q/k_A$ are integers, and an element $s\neq0$ of $\Omega$ such that, with
\[
 \alpha=-q\sum_{\tau,i}\nu_{\tau,i}a_{\tau,i}-\frac q{k_A}\,a_A\in\mathbf Z^K,
\]
$s\prod_ky_{k,\kappa_k}^{-\alpha_k}$ is integral over $\mathbf Z[y_{k,i}/y_{k,\kappa_k}:k\in K,\ i\in I_k]$
for every chart $\kappa$, that is, every choice of indices with all $y_{k,\kappa_k}\neq0$. So $s$ is a
nonzero element of $\tilde H^0(X_M,\bar L)$
(Definition~\ref{lit-multiprojective}(4)) for every metrized line bundle $\bar L$ of multidegree
$\alpha$, which is the multidegree of the class $q\,\ell(\nu,-1)$ of
Definition~\ref{cv-condition-flag-model},
with $x_{\tau,i}$ standing for $\bar x_{\tau,i}(t)$ and $a$ for $\bar A/k_A$.
\end{lemma}
\leanref{Development/AutomorphicSpectrum/Proofs/MetricProofs/RootOriginWeightIdentityStatement.lean}{GlobalTheory.T.root_origin_weight_identity}
\leanref{Development/WeilBound/CoreStatements/FlagImmersions/GlobalModelTowerBounds.lean}{RootFieldNumerics.global_model_tower_bounds}
\leanref{Development/AutomorphicSpectrum/Definitions/DivisibilityTowers/Birational.lean}{GlobalTheory.CompanionTheory.Ar.Top.RootModel.EffDiff}
\begin{proof}
Let $k_0$, $q_c$, $s_c$, $g$ and the root lines $\Lambda_{\tau,i}$ with generators $e_{\tau,i}$ be
those of the origin at $(B,c)$ (Definition~\ref{cv-origin-tower}(2)), and put
$n_{\tau,i}=n^c_{\tau,i}=k_{A,c}q_c\nu_{\tau,i}$; here $k_A=k_{A,c}$ by (2c)(i). Put $q=k_Aq_c$, so
that $q\nu_{\tau,i}=n_{\tau,i}$ and $q/k_A=q_c$ are integers, and
\[
 s_0=g(s_c)\prod_{\tau,i}e_{\tau,i}^{n_{\tau,i}}\in\Omega,
\]
which is nonzero. By (2c)(ii), $a_A$ is $1$ at $k_0$ and $0$ elsewhere, and
$\alpha=-\sum_{\tau,i}n_{\tau,i}a_{\tau,i}-q_ca_A$.

Fix a chart $\kappa$, write $Y_k=y_{k,\kappa_k}$, and let $i$ be the index with $i'=\kappa_{k_0}$. Put
$t=\prod_{\tau,l}\big(e_{\tau,l}\prod_kY_k^{a_{\tau,l,k}}\big)^{-n_{\tau,l}}$. Expanding the products
gives
\[
 f_\kappa:=s_0\prod_kY_k^{-\alpha_k}=\frac{g(s_c)\,Y_{k_0}^{q_c}}t.
\]
Let $Z$ be the Hodge component of the origin at $(B,c)$ and $P$ a point of its open set $Z_\kappa$
over no prime. Then $Y_{k_0}=y_{k_0,i'}\neq0$ and the ratios $y_{k_0,l}/Y_{k_0}$ are regular at $P$.
Each $e_{\tau,l}\prod_kY_k^{a_{\tau,l,k}}$ generates $\Lambda_{\tau,l}$ at $P$, because the flag has
root lines with these generators, so $t$ generates $\bigotimes_{\tau,l}\Lambda_{\tau,l}^{-n_{\tau,l}}$
at $P$. By (2c)(iii), $f_\kappa$ is regular at $P$.

The open set $Z_\kappa$ is affine with ring $R_\kappa=\mathbf Z[y_{k,l}/Y_k]$, so its points over no
prime are the prime ideals of the domain $R_\kappa\otimes\mathbf Q$, with the localizations of
$R_\kappa\otimes\mathbf Q$ as local rings. An element of the fraction field that lies in all of these
localizations lies in $R_\kappa\otimes\mathbf Q$. So $f_\kappa\in R_\kappa\otimes\mathbf Q$, and
$N_\kappa f_\kappa\in R_\kappa$ for some integer $N_\kappa\ge1$. Let $N_0$ be the product of the
$N_\kappa$ over the finitely many charts. Then $s=N_0s_0$ is nonzero, and
$s\prod_ky_{k,\kappa_k}^{-\alpha_k}=N_0f_\kappa$ lies in $R_\kappa$ for every chart $\kappa$.
\end{proof}
\leanref{Development/AutomorphicSpectrum/Proofs/MetricProofs/RootOriginWeightIdentity.lean}{GlobalTheory.CompanionTheory.Ar.Small.OriginEps.section_on_model}
\leanref{Development/AutomorphicSpectrum/Proofs/MetricProofs/RootOriginWeightIdentity.lean}{GlobalTheory.CompanionTheory.Ar.Small.OriginEps.mem_adjoin_rat_of_regular}
\leanref{Development/AutomorphicSpectrum/Proofs/AlgebraProofs/IntegralSectionLatticeMultipleStatement.lean}{GlobalTheory.T.integral_section_lattice_multiple}

\begin{definition}\label{cv-volume-global-model}
In Situation~\ref{cv-function-situation}, let $l=(M,\mathfrak M,N)$ be the level of $X$, let the moduli
schemes $X_{l'}$ of the PEL datum at all levels $l'$ be given
(Definition~\ref{cv-model-toroidal-semiabelian}(1)), with $X=X_l$, and let $(\mathcal Y,\mathcal A,j)$ be a
model of $X$ over the ring of integers $\mathcal O_c$ of $F_c$, defined as in
Definition~\ref{cv-model} with $\mathcal O_c$ in place of $O$. For a prime $p$, a \emph{place over} $p$ is
a valuation ring $W$ of $\mathbf C$ whose maximal ideal contains $p$; for a number field
$F'\subset\mathbf C$ put $O_{F',W}=F'\cap W$. An \emph{index at} $p$ is a pair $i=(l_i,F_i)$ of a level
$l_i=(M_i,\mathfrak M_i,N_i)$ such that the prime-to-$p$ part of $M_i$ is at least $3$, and a coefficient
field $F_i\subset\mathbf C$ for $l_i$: a number field that contains $K_{l_i}$, the images of all
embeddings of $E$ and the roots of unity of order the least common multiple of $M_i$ and $N_i$. A
\emph{local system of models at} $(p,W)$ is a family of models $(\mathcal Y_i,\mathcal A_i,j_i)$ of
$X_{l_i}\otimes F_i$ over $O_{F_i,W}$, one for each index $i$ at $p$. The model $\mathcal Y$ is
\emph{global} if:
\begin{enumerate}
\item $N=1$, and for every prime $p$ the prime-to-$p$ part of $M$ is at least $3$;
\item $\mathcal Y$ \emph{descends}: there are a scheme $X_K$ over $K_l$, an open immersion $X\to X_K$
  over $K_l$ and a group scheme $G_K$ over $X_K$ whose pullback to $X$ is $A$, such that the generic
  fibre $\mathcal Y\otimes_{\mathcal O_c}F_c$ is $X_K\otimes_{K_l}F_c$, compatibly with $j$, and the
  pullbacks of $\mathcal A$ and of $G_K$ to it are identified, compatibly with their identifications
  with $A$ over $X\otimes_{K_l}F_c$;
\item for every prime $p$ and every place $W$ over $p$ there is a local system $(\mathcal Y_i)$ of models
  at $(p,W)$ with the following properties.
  \begin{enumerate}
  \item[(a)] For each $i$: every automorphism of $X_{l_i}\otimes F_i$ over $F_i$, with an automorphism
    over it of the universal abelian scheme that changes the level structures $\eta$ and $\eta'$ by
    $\mathcal O_E$-linear endomorphisms of $\Lambda$, extends to an automorphism of $\mathcal Y_i$ over
    $O_{F_i,W}$ with an automorphism of $\mathcal A_i$ over it; if $N_i=1$, $\mathcal Y_i$ descends as in
    (2); the conclusion of Lemma~\ref{cv-model-toroidal-pushforward} holds for $\mathcal A_i$; for
    $a\in\Lambda$ and $n$ prime to $p$, each of the sections $\eta(a)$ and $\eta'(a)$ whose order divides
    $n$ extends to a section of $\mathcal A_i$ over its abelian locus; and, for $V$ a complete discrete
    valuation ring with fraction field $K$, or the valuation ring of a complete algebraically closed field
    $K$ of characteristic zero in which $|p|<1$, every $K$-point of $X_{l_i}\otimes F_i$ over a
    homomorphism $O_{F_i,W}\to V$, at which the abelian variety extends to an abelian scheme over $V$
    with a relatively ample line bundle compatible with its polarization, extends to a $V$-point of
    $\mathcal Y_i$ in the abelian locus of $\mathcal A_i$.
  \item[(b)] If $l_i$ divides $l_{i_2}$ and $F_i\subset F_{i_2}$, there are a morphism
    $\pi\colon\mathcal Y_{i_2}\to\mathcal Y_i$, finite as a morphism to
    $\mathcal Y_i\otimes O_{F_{i_2},W}$, and a morphism $\mathcal A_{i_2}\to\mathcal A_i$ over it, that
    extend a morphism $X_{l_{i_2}}\otimes F_{i_2}\to X_{l_i}\otimes F_i$ under which $M_{i_2}/M_i$ times
    $\eta$ and $N_{i_2}/N_i$ times $\eta'$ of $\mathcal A_{i_2}$ become $\eta$ and $\eta'$ of $\mathcal A_i$.
  \item[(c)] There are an odd prime $n\ne p$, prime to the integer $M$ of every level, an integer $J$ with
    $n^J\ge3$, and a Siegel moduli scheme $\mathcal S$, smooth of relative dimension $g(g+1)/2$ over
    $\mathbf Z[\zeta_{n^J},1/m:p\nmid m]$ with $g=[E:\mathbf Q](p+q)/2$, carrying a universal principally
    polarized abelian scheme with a symplectic level-$n^J$ structure, unramified over its base, and
    through which every such abelian scheme over the valuation ring of a complete algebraically closed
    field over $p$ with ample reduction factors, with the following property. For every index $i$, let
    $i_1$ be the index with the same $M_i$ and $\mathfrak M_i$, full level the least common multiple of
    $N_i$ and $n^J$, and field $F_i(\zeta_{n^J})$. Then $\pi$ of (b) is finite, \'etale and surjective
    from the abelian locus of $\mathcal Y_{i_1}$ to the base change of that of $\mathcal Y_i$; for a
    symplectic basis of $\Lambda$, the abelian locus of $\mathcal Y_{i_1}$ has a finite morphism to
    $\mathcal S$ that classifies $\mathcal A_{i_1}$ there with its polarization and the level structure
    $(N_{i_1}/n^J)\eta'$; and if $p$ does not divide the least common multiple of $M_i$ and $N_i$, this
    morphism is the composite of a finite surjective morphism onto a reduced scheme, an isomorphism over a
    dense open subscheme, that carries an abelian scheme with the level structures, and a finite, flat and
    unramified morphism to $\mathcal S$.
  \item[(d)] Each $\mathcal Y_i$ has properties (1) and (2) of Theorem~\ref{cv-model-toroidal}, and the
    following arc property. Let $O'$ be a domain of characteristic zero with a homomorphism
    $O_{F_i,W}\to O'$, which is either a field, with $\varpi=1$, or a complete discrete valuation ring
    with uniformizer $\varpi$, and let an arc over $O'$ be given: degeneration data over $O'$ for
    Mumford's construction, with an action of $\mathcal O_E$, a positive definite symmetric form on its
    group of periods and the level structures of the moduli problem at $l_i$. Then there is a unique
    morphism $\chi\colon\operatorname{Spec}O'[[x]]\to\mathcal Y_i$ over $O_{F_i,W}$ whose preimage of the
    abelian locus is where $x\ne0$, which maps the locus where $\varpi x\ne0$ into
    $j_i(X_{l_i}\otimes F_i)$, and along which $\mathcal A_i$ is given by Mumford's construction from the
    arc with parameter $x$; and if $\varpi$ is a uniformizer, every $O'$-point of $\mathcal Y_i$ over
    $O_{F_i,W}\to O'$ whose generic point lies in $j_i(X_{l_i}\otimes F_i)$ and along which $\mathcal A_i$
    is given by the arc with parameter $\varpi$ is $\chi$ composed with $x\mapsto\varpi$.
  \end{enumerate}
  Moreover $\mathcal Y$ \emph{dominates} $\mathcal Y_{(l,F_c)}$: there is a morphism $f$ from
  $\mathcal Y\otimes_{\mathcal O_c}O_{F_c,W}$ to $\mathcal Y_{(l,F_c)}$ over $O_{F_c,W}$ that is the
  identity on $X\otimes F_c$, with an identification of the pullback of $\mathcal A$ with
  $f^*\mathcal A_{(l,F_c)}$ that is compatible with the identifications of both with $A$.
\end{enumerate}
\end{definition}
\leanref{Development/AutomorphicSpectrum/Definitions/DivisibilityTowers/GoodOrigin.lean}{GlobalTheory.CompanionTheory.Ar.Top.IsGlobal}
\leanref{ReciprocityTheory/ReductiveGroups/ArithmeticGeometry/BirationalGeometry/ToroidalCompactification/Clauses.lean}{ReciprocityTheory.ChartRestriction.G.Geo.Tor.LocalSystem}
\leanref{ReciprocityTheory/ReductiveGroups/ArithmeticGeometry/BirationalGeometry/ToroidalCompactification/Clauses.lean}{ReciprocityTheory.ChartRestriction.G.Geo.Tor.N2}
\leanref{ReciprocityTheory/ReductiveGroups/ArithmeticGeometry/BirationalGeometry/ToroidalCompactification/Clauses.lean}{ReciprocityTheory.ChartRestriction.G.Geo.Tor.Model.LevelStructureConditions}
\leanref{ReciprocityTheory/ReductiveGroups/ArithmeticGeometry/BirationalGeometry/ToroidalCompactification/Clauses.lean}{ReciprocityTheory.ChartRestriction.G.Geo.Tor.SiegelCompare}
\leanref{ReciprocityTheory/ReductiveGroups/ArithmeticGeometry/BirationalGeometry/ToroidalCompactification/UniformisationConclusion.lean}{ReciprocityTheory.ChartRestriction.G.Geo.Tor.Model.N4}
\leanref{ReciprocityTheory/ReductiveGroups/ArithmeticGeometry/BirationalGeometry/ToroidalCompactification/PushInvertible.lean}{ReciprocityTheory.ChartRestriction.G.Geo.Tor.Model.Descends}
\leanref{ReciprocityTheory/ReductiveGroups/ArithmeticGeometry/BirationalGeometry/ToroidalCompactification/PushInvertible.lean}{ReciprocityTheory.ChartRestriction.G.Geo.Tor.Dominates}

\begin{definition}\label{cv-origin-trivial-flag}
Let $K$, $(I_k)$ and $R$ be as in Definition~\ref{lit-multiprojective}, let $\Omega$ be a field, let
$y\in\Omega^{\sqcup_kI_k}$, and let $o\subset\Omega$ be a subring. A \emph{chart} of $y$ is a choice
$\kappa=(\kappa_k)_{k\in K}$ of indices $\kappa_k\in I_k$ with $y_{k,\kappa_k}\neq0$ for every $k$. For
a chart $\kappa$ and $a\in\mathbf Z^K$ put $Y_\kappa=\prod_ky_{k,\kappa_k}$ and
$y_\kappa^a=\prod_ky_{k,\kappa_k}^{a_k}$, and let $o_\kappa=o[y_{k,i}/y_{k,\kappa_k}:k\in K,\ i\in I_k]$
be the subring of $\Omega$ generated over $o$ by these ratios. Let $\varpi\in o$, let $N\ge1$ and let
$b_1,\dots,b_N\in\mathbf Z^K$. The family $(b_i)$ has a \emph{trivial flag modulo $\varpi$} on $y$ over
$o$ if there are $N\times N$ matrices $g_\kappa$ with entries in $\Omega$, one for each chart $\kappa$,
such that:
\begin{enumerate}
\item the entries of $g_\kappa$ lie in $o_\kappa$, and $\det(g_\kappa)\,h=1+\varpi r$ for some
  $h,r\in o_\kappa$;
\item for all charts $\kappa$ and $\kappa'$, with $\phi=Y_\kappa/Y_{\kappa'}$, there are an integer
  $m\ge0$ and $N\times N$ matrices $U$ and $Z$ with entries in $o_{\kappa'}$, with $U$ upper triangular
  and $U_{ii}=\phi^m\,y_\kappa^{b_i}/y_{\kappa'}^{b_i}$ for every $i$, such that $\phi^mg_\kappa$ has
  entries in $o_{\kappa'}$ and $\phi^mg_\kappa=g_{\kappa'}U+\varpi Z$.
\end{enumerate}
\end{definition}
\leanref{Development/AutomorphicSpectrum/Definitions/CommutativeAlgebra/TrivFlagMod.lean}{GlobalTheory.CompanionTheory.Ar.Num.TrivFlagMod}
\leanref{ReciprocityTheory/ReductiveGroups/ArithmeticGeometry/BirationalGeometry/Bundle.lean}{ReciprocityTheory.ChartRestriction.G.Geo.chartRing}

\begin{lemma}\label{cv-origin-trivial-flag-divisor}
Let $\Omega$, $y$, $o$, $N$ and $b_1,\dots,b_N$ be as in Definition~\ref{cv-origin-trivial-flag}, and let
$\varpi,\varpi'\in o$ with $\varpi'$ dividing $\varpi$ in $o$. If $(b_i)$ has a trivial flag modulo
$\varpi$ on $y$ over $o$, then it has a trivial flag modulo $\varpi'$ on $y$ over $o$.
\end{lemma}
\leanref{Development/AutomorphicSpectrum/Proofs/GeometryProofs/TrivFlagModOfDvd.lean}{GlobalTheory.CompanionTheory.Ar.HT.trivFlagMod_of_dvd}
\begin{proof}
Write $\varpi=\varpi'u$ with $u\in o\subset o_\kappa$. The same matrices $g_\kappa$ serve, with $ur$ in
place of $r$ in (1) and $uZ$ in place of $Z$ in (2).
\end{proof}

\begin{lemma}\label{cv-origin-trivial-flag-section}
Let $K$, $(I_k)$ and $R$ be as in Definition~\ref{lit-multiprojective}, let $\Omega$ be a field of
characteristic zero, let $y\in\Omega^{\sqcup_kI_k}$, and let $P\subset R$ be the kernel of
Definition~\ref{lit-multiprojective}(4). Let $p$ be a prime, let $o\subset\Omega$ be the ring of
fractions $a/b$ with $a,b\in\mathbf Z$ and $p\nmid b$, and use the charts $\kappa$ of $y$, the elements
$Y_\kappa$ and $y_\kappa^a$ and the rings $o_\kappa$ of Definition~\ref{cv-origin-trivial-flag}. Let
$W\supset P$ be a relevant multihomogeneous prime of $R$ with $W\cap\mathbf Z=p\mathbf Z$ and
$\dim R/W=1+|K|$. Let $a\in\mathbf Z^K$, and let $r_\kappa\in o_\kappa$ for each chart $\kappa$. Assume:
\begin{enumerate}
\item[(i)] for all charts $\kappa$ and $\kappa'$ there are $n\ge0$ and $z\in o_{\kappa'}$ with
  $(Y_\kappa/Y_{\kappa'})^n\big(r_\kappa\,y_\kappa^a/y_{\kappa'}^a-r_{\kappa'}\big)=pz$;
\item[(ii)] there are indices $\kappa_k\in I_k$ with $x_{k,\kappa_k}\notin W$ for every $k$, integers
  $n\ge1$ prime to $p$ and $d\ge0$, and a multihomogeneous $g\in R\setminus W$ of multidegree
  $(d,\dots,d)$ with $n\,r_\kappa Y_\kappa^d=g(y)$.
\end{enumerate}
Then $\mathcal O(a)$ has nonnegative degree on $W$.
\end{lemma}
\leanref{Development/WeilBound/CoreStatements/HodgeTatePoints/GfCurve.lean}{PositivityTheory.HtPt.Gf.curve_sec_nef}
\leanref{Development/WeilBound/CoreStatements/HodgeTatePoints/GeoNumNonnegOfRegular.lean}{PositivityTheory.HtPt.Gf.geoNum_nonneg_of_regular}
\leanref{Development/WeilBound/CoreStatements/HodgeTatePoints/GfForms.lean}{PositivityTheory.HtPt.Gf.exists_rep}
\begin{proof}
For a choice of indices $\kappa=(\kappa_k)$ and $c\in\mathbf N^K$ write
$x_\kappa^c=\prod_kx_{k,\kappa_k}^{c_k}\in R$, and write $(e)$ for the constant family $(e,\dots,e)$. Two
facts are used. A multihomogeneous $F\in R$ with $F(y)=0$ lies in $P$, since $P$ is the kernel of
$x_{k,i}\mapsto y_{k,i}t_k$ and $F$ maps to $F(y)$ times a monomial in the $t_k$. For a chart $\kappa'$,
every $x\in o_{\kappa'}$ can be written $n'xY_{\kappa'}^e=g'(y)$ with $n'\ge1$ prime to $p$, $e\ge0$ and
$g'\in R$ multihomogeneous of multidegree $(e)$: write $x$ as a polynomial in the ratios
$y_{k,i}/y_{k,\kappa'_k}=y_{k,i}\prod_{k'\neq k}y_{k',\kappa'_{k'}}/Y_{\kappa'}$ with coefficients in
$o$, and clear the denominators.

Let $\kappa$, $n$, $d$ and $g$ be as in (ii); $\kappa$ is a chart, since $x_{k,\kappa_k}\notin W\supset P$.
Write $a=a^+-a^-$ with $a^\pm\in\mathbf N^K$, and put $f=g\,x_\kappa^{a^+}$ and
$f'=x_\kappa^{(d)}x_\kappa^{a^-}$. They are multihomogeneous, the difference of their multidegrees is
$a$, and neither lies in the prime $W$, since $g\notin W$ and no $x_{k,\kappa_k}$ lies in $W$.

Let $\kappa'$ be indices with $x_{k,\kappa'_k}\notin W$ for every $k$; $\kappa'$ is a chart. By (i) for
$(\kappa,\kappa')$ there are $m\ge0$ and $z\in o_{\kappa'}$ with
$(Y_\kappa/Y_{\kappa'})^m(r_\kappa y_\kappa^a/y_{\kappa'}^a-r_{\kappa'})=pz$. Write
$n_1r_{\kappa'}Y_{\kappa'}^{e_1}=g_1(y)$ and $n_2zY_{\kappa'}^{e_2}=g_2(y)$ as above. Multiplying the
relation by $nn_1n_2Y_{\kappa'}^{m+e_1+e_2}Y_\kappa^dy_{\kappa'}^{a^+}y_\kappa^{a^-}$ shows that the
multihomogeneous element
\[
 x_\kappa^{(m)}\big(n_1n_2\,x_{\kappa'}^{(e_1+e_2)}x_{\kappa'}^{a^-}f-nn_2\,g_1x_{\kappa'}^{(e_2)}x_{\kappa'}^{a^+}f'\big)
 -p\,nn_1\,g_2\,x_{\kappa'}^{(e_1+m)}x_{\kappa'}^{a^+}f'
\]
vanishes at $y$, so it lies in $P\subset W$. As $p\in W$ and $x_\kappa^{(m)}\notin W$, the element
$n_1n_2\,x_{\kappa'}^{(e_1+e_2)+a^-}f-v f'$ lies in $W$, where $v=nn_2\,g_1x_{\kappa'}^{(e_2)+a^+}$, and
$n_1n_2$ is prime to $p$.

Let $C$ be the integral projective curve over $\mathbf F_p$ defined by $W$; its points are the relevant
multihomogeneous primes $W'\supset W$. The quotient $f/f'$ is a nonzero rational section of
$\mathcal O(a)$ on $C$. It is regular at every point $W'$: as $W'$ is relevant there are indices
$\kappa'$ with $x_{k,\kappa'_k}\notin W'$, hence $x_{k,\kappa'_k}\notin W$, for every $k$, and with
$n_1n_2$, $e_1+e_2$ and $v$ as above, $f/f'=v/(n_1n_2\,x_{\kappa'}^{(e_1+e_2)+a^-})$ near $W'$, whose
denominator is a unit there. A line bundle
with a nonzero regular section on an integral projective curve has nonnegative degree, the length of the
quotient of the bundle by the subsheaf the section generates. So $\mathcal O(a)$ has nonnegative degree
on $W$.
\end{proof}

\begin{lemma}\label{cv-origin-trivial-flag-degrees}
Let $\Omega$, $y$, $P$, $p$, $o$ and $W$ be as in Lemma~\ref{cv-origin-trivial-flag-section}, let $N\ge1$,
and let $b_1,\dots,b_N\in\mathbf Z^K$ have a trivial flag modulo $p$ on $y$ over $o$
(Definition~\ref{cv-origin-trivial-flag}). Then $\sum_{l\le i}\deg_W\mathcal O(b_l)\le0$ for every
$0\le i\le N$, and $\sum_{l=1}^N\deg_W\mathcal O(b_l)=0$.
\end{lemma}
\leanref{Development/WeilBound/CoreStatements/HodgeTatePoints/ExistsVanishingIdeal.lean}{PositivityTheory.HtPt.Gf.trivflag_deg}
\leanref{Development/WeilBound/CoreStatements/HodgeTatePoints/ExistsVanishingIdeal.lean}{PositivityTheory.HtPt.Gf.exists_vanishing_ideal}
\leanref{Development/WeilBound/CoreStatements/HodgeTatePoints/ExistsMinorNeZero.lean}{PositivityTheory.HtPt.Gf.lead_minor}
\begin{proof}
The degree on $W$ is additive in the multidegree. Let $(g_\kappa)$ be as in
Definition~\ref{cv-origin-trivial-flag}, with $h_\kappa$ and $\rho_\kappa$ in $o_\kappa$ such that
$\det(g_\kappa)h_\kappa=1+p\rho_\kappa$. As $W$ is relevant there are indices $\kappa_0$ with
$x_{k,\kappa_{0,k}}\notin W$ for every $k$; $\kappa_0$ is a chart. We use the notation and the two facts of
the proof of Lemma~\ref{cv-origin-trivial-flag-section}.

The ideal at $\kappa_0$. Let $I$ be the set of $x\in o_{\kappa_0}$ with $nxY_{\kappa_0}^d=g(y)$ for some
$n\ge1$ prime to $p$, some $d\ge0$ and some multihomogeneous $g\in W$ of multidegree $(d)$. It is an
ideal: if $nxY_{\kappa_0}^d=g(y)$ and $n'x'Y_{\kappa_0}^{d'}=g'(y)$, then
$nn'(x+x')Y_{\kappa_0}^{d+d'}$ is the value at $y$ of $n'gx_{\kappa_0}^{(d')}+nx_{\kappa_0}^{(d)}g'$, and for
$u\in o_{\kappa_0}$ with $n''uY_{\kappa_0}^{e}=g''(y)$ the product $n''n\,uxY_{\kappa_0}^{e+d}$ is the value of
$g''g$. It contains $p$, the value of $g=p$. It does not contain $1$: if $nY_{\kappa_0}^d=g(y)$ with
$g\in W$, then $nx_{\kappa_0}^{(d)}-g$ vanishes at $y$, so it lies in $P\subset W$ and
$nx_{\kappa_0}^{(d)}\in W$; but $n\notin W$, since $n$ is prime to $p\in W$ and $1\notin W$, and
$x_{\kappa_0}^{(d)}\notin W$. Finally every $x\in o_{\kappa_0}\setminus I$ has a representation
$nxY_{\kappa_0}^d=g(y)$ with $p\nmid n$ and $g$ multihomogeneous of multidegree $(d)$, and then
$g\notin W$.

The minors. For $0\le i\le N$, a map $J_0\colon\{1,\dots,i\}\to\{1,\dots,N\}$ and a chart $\kappa$, let
$\Delta_\kappa$ be the determinant of the $i\times i$ matrix with entries $(g_\kappa)_{J_0(a),b}$ for
$a,b\le i$, and put $c=b_1+\dots+b_i$. Let $\kappa,\kappa'$ be charts, $\phi=Y_\kappa/Y_{\kappa'}$, and
$m$, $U$, $Z$ as in Definition~\ref{cv-origin-trivial-flag}(2). Since $U$ is upper triangular, the
submatrix of $g_{\kappa'}U$ on the rows $J_0$ and the first $i$ columns is the product of the same
submatrix of $g_{\kappa'}$ with the upper left $i\times i$ block of $U$, whose determinant is
$\prod_{l\le i}U_{ll}=\phi^{mi}y_\kappa^c/y_{\kappa'}^c$. Adding $pZ$, whose entries lie in
$o_{\kappa'}$, changes the determinant by $p$ times an element of $o_{\kappa'}$. So
\[
 \phi^{mi}\big(\Delta_\kappa-\Delta_{\kappa'}\,y_\kappa^c/y_{\kappa'}^c\big)\in p\,o_{\kappa'}.\tag{$*$}
\]
For every $u\in\mathbf Z^K$ and every $x\in o_\kappa$, some power of $\phi$ times $y_\kappa^u/y_{\kappa'}^u$,
and some power of $\phi$ times $x$, lie in $o_{\kappa'}$: indeed
$\phi^n\,y_\kappa^u/y_{\kappa'}^u=\prod_k(y_{k,\kappa_k}/y_{k,\kappa'_k})^{n+u_k}$, and
$\phi\,y_{k,i}/y_{k,\kappa_k}=(y_{k,i}/y_{k,\kappa'_k})\prod_{k'\neq k}y_{k',\kappa_{k'}}/y_{k',\kappa'_{k'}}$.
Also $\phi\in o_{\kappa'}$.

The partial sums. Let $0\le i\le N$. The image of $\det(g_{\kappa_0})$ in $o_{\kappa_0}/I$ is a unit, since
$\det(g_{\kappa_0})h_{\kappa_0}=1+p\rho_{\kappa_0}$ and $p\in I$. Over a commutative ring, the $i\times i$
minors on the first $i$ columns of a square matrix with unit determinant generate the unit ideal: by the
Cauchy--Binet formula, the product of the first $i$ rows of the inverse matrix with the first $i$ columns
of the matrix is the identity, and its determinant is a combination of these minors. As
$o_{\kappa_0}/I\neq0$, there is a map $J_0$ with $\Delta_{\kappa_0}\notin I$. Apply
Lemma~\ref{cv-origin-trivial-flag-section} with $-c$ in place of $a$ and $r_\kappa=\Delta_\kappa$. Its
condition (i) follows from $(*)$, multiplied by a power of $\phi$ times $y_{\kappa'}^c/y_\kappa^c$ that lies
in $o_{\kappa'}$; its condition (ii) holds at $\kappa_0$, since $\Delta_{\kappa_0}\notin I$. So
$\deg_W\mathcal O(-c)\ge0$, that is, $\sum_{l\le i}\deg_W\mathcal O(b_l)\le0$.

The total. Put $c=b_1+\dots+b_N$ and $\gamma=y_\kappa^c/y_{\kappa'}^c$. By $(*)$ with $i=N$ and $J_0$ the
identity, $\phi^{mN}(\det g_\kappa-\det g_{\kappa'}\,\gamma)=pz$ with $z\in o_{\kappa'}$. From
$\det(g_\kappa)h_\kappa=1+p\rho_\kappa$ and $\det(g_{\kappa'})h_{\kappa'}=1+p\rho_{\kappa'}$,
\[
 h_\kappa\gamma-h_{\kappa'}=-h_\kappa h_{\kappa'}\big(\det g_\kappa-\det g_{\kappa'}\,\gamma\big)
 +p\big(h_{\kappa'}\rho_\kappa-h_\kappa\gamma\rho_{\kappa'}\big),
\]
and multiplying by a suitable power of $\phi$ puts $h_\kappa$, $\gamma$ and $\rho_\kappa$ into $o_{\kappa'}$,
so some power of $\phi$ times $h_\kappa\gamma-h_{\kappa'}$ lies in $p\,o_{\kappa'}$. Moreover
$h_{\kappa_0}\notin I$, since otherwise $1=\det(g_{\kappa_0})h_{\kappa_0}-p\rho_{\kappa_0}$ would lie in
$I$. So Lemma~\ref{cv-origin-trivial-flag-section} with $a=c$ and $r_\kappa=h_\kappa$ gives
$\sum_l\deg_W\mathcal O(b_l)\ge0$. With the partial sum for $i=N$ the total is $0$.
\end{proof}

\begin{lemma}\label{cv-origin-trivial-flag-nef}
Let $M$ be a flag model of type $(\Psi,N)$ (Definition~\ref{cv-condition-flag-model}) whose field
$\Omega$ has characteristic zero, with point $y$, arithmetic variety $X$ and multidegrees $a_{\tau,i}$ of
its bundles $\bar x_{\tau,i}(t)$. Let $p$ be a prime and $o\subset\Omega$ the ring of fractions $a/b$ with
$a,b\in\mathbf Z$ and $p\nmid b$. Assume that for every $\tau\in\Psi$ the family
$(a_{\tau,1},\dots,a_{\tau,N})$ has a trivial flag modulo $p$ on $y$ over $o$
(Definition~\ref{cv-origin-trivial-flag}). Then for every dominant weight $\mu$ and every vertical curve
$W$ in $X$ with $W\cap\mathbf Z=p\mathbf Z$,
\[
 \deg_W\ell(\mu)=-\sum_{\tau,i}\mu_{\tau,i}\deg_W\mathcal O(a_{\tau,i})\ge0.
\]
\end{lemma}
\leanref{Development/AutomorphicSpectrum/Proofs/GeometryProofs/TrivflagNefFibre.lean}{GlobalTheory.CompanionTheory.Ar.HT.trivflag_nef_fibre}
\leanref{Development/WeilBound/CoreStatements/HodgeTatePoints/TrivflagNef.lean}{PositivityTheory.HtPt.Gf.trivflag_nef}
\leanref{Development/WeilBound/CoreStatements/HodgeTatePoints/ExistsMinorNeZero.lean}{PositivityTheory.HtPt.Gf.abel_sum}
\begin{proof}
Fix $\tau$ and put $u_i=-\deg_W\mathcal O(a_{\tau,i})$ and $U_i=\sum_{l\le i}u_l$. By
Lemma~\ref{cv-origin-trivial-flag-degrees}, $U_i\ge0$ for every $i$ and $U_N=0$. As $\mu$ is dominant,
$\mu_{\tau,1}\ge\dots\ge\mu_{\tau,N}$, so by Abel summation
\[
 \sum_{i=1}^N\mu_{\tau,i}u_i=\sum_{i<N}(\mu_{\tau,i}-\mu_{\tau,i+1})U_i+\mu_{\tau,N}U_N\ge0.
\]
Summing over $\tau$ gives the claim.
\end{proof}

\begin{theorem}\label{cv-volume-theorem}
There is a constant $C\ge0$ with the following property. Let $E$ be an abelian CM field
with a host structure as in Theorem~\ref{cv-flat-degree}, let $\Phi$ be its set of CM
types, and let a PEL unitary datum of signature $(p,q)$ over $E$ be given, with
$q\ge1$ and $3q+1\le N:=p+q$, together with its moduli schemes at all levels, a level, a
coefficient field $F_c$, a global model $\mathcal Y$ of the moduli scheme over the ring of
integers of $F_c$ (Definition~\ref{cv-volume-global-model}), a tower of integral models of the
flag bundle over the moduli scheme, that is, a tower of type $(\Phi,N)$ with a good origin in
$\mathcal Y$ (Definition~\ref{cv-origin-tower}), whose set of bad primes has at most $4|\Phi|N^2$
elements, and a full arithmetic intersection system (Definition~\ref{cv-origin-intersection-full}; one exists by
Lemma~\ref{cv-origin-intersection-exists}). Write $\dim_{\mathrm{fl}}=|\Phi|\,N(N-1)/2$. Then:
\begin{enumerate}
\item the geometric degrees of tautological monomials on each component $c$ are
  $\operatorname{vol}(c)$ times their degrees on the flag variety, with
  $\operatorname{vol}(c)>0$ independent of $c$;
\item for every regular dominant weight $u\in(\mathbf R^N)^\Phi$ with gaps at least $g>0$
  and every $B$ with $B+1>4|\Phi|N^2(1+\|u\|/g)$, small perturbations of the tautological
  bundle of weight $u$ and of its metric are nef on the fibres and semipositive, and their
  arithmetic self-intersection numbers and slopes converge, as the perturbation tends to
  $0$, to the arithmetic intersection number and the slope of weight $u$ at level $B$;
\item the arithmetic intersection numbers of tautological monomials at level $B$ agree,
  up to $O(\log B/B)$, with $\sum_c\operatorname{vol}(c)$ times an explicit linear form in
  the monomial whose coefficients depend on $B$ only through two constants for each
  $\tau\in\Phi$;
\item for $u$ and $B$ as in (2) with $B\ge128$ and at least the number of bad primes, the
  arithmetic top intersection number of weight $u$ and the mean over components of the
  arithmetic slopes of weight $u$ are bounded below; the second bound reads
  \[
   -C\Big(\big(\log(B\operatorname{rd}_E)+\big|\tfrac{L'}{L}(1,\chi)\big|\big)\|u\|_1
   +\min(j,N-j)\operatorname{spread}(u)\Big),
  \]
  where $j$ is the type of the flag, $\|u\|_1=\sum_{\tau,i}|u_{\tau,i}|$ and
  $\operatorname{spread}(u)=\sum_\tau(\max_iu_{\tau,i}-\min_iu_{\tau,i})$.
\end{enumerate}
\end{theorem}
\leanref{Development/AutomorphicSpectrum/Statements/TowerExistence/GoodTowerExistsX.lean}{GlobalTheory.CompanionTheory.Ar.Top.VolumeTheoremX}
\leanref{Development/WeilBound/CoreStatements/AutomorphicKernels/MoriwakiBigness/HilbertSamuelApplications/GoodTowerVolumeTheoremBis.lean}{GlobalTheory.CompanionTheory.Ar.Link.good_tower_volume_theorem_bis}
\begin{proof}
Write $N=p+q$ and $\Psi$ for the set of CM types. The tower consists of an index $j$ with
$1\le j<N$, a finite set of components $c$, a finite set of bad primes containing $2$, a rational
weight $\nu$, and for each level $B\ge2$ and component $c$ an integral model $M(B,c)$ of the flag
bundle. For a polynomial $F$ in the tautological classes $x_{\tau,i}$ ($\tau\in\Psi$, $1\le i\le N$)
let $\deg_{B,c}(F)$ be its geometric degree on $M(B,c)$ and, in degree $\dim_{\mathrm{fl}}+1$,
$\mathrm{PP}_{B,c}(F)$ its arithmetic intersection number; put
$\mathrm{PP}_B=\sum_c\mathrm{PP}_{B,c}$, $\operatorname{vol}_T=\sum_c\operatorname{vol}(c)$, and for a
weight $u$ let $\ell(u)=-\sum_{\tau,i}u_{\tau,i}x_{\tau,i}$. Let $\deg_{\mathrm{fl}}$ be the degree on
the flag variety (Definition~\ref{cv-reconstruction-functionals}) and $h$ the flag height, the
functional on polynomials of degree $\dim_{\mathrm{fl}}+1$ given by the arithmetic intersection
numbers of the flag variety over $\mathbf Z$ with its standard metrics
(Lemma~\ref{cv-reconstruction-flag-height}). The proof uses, besides the inputs of
Theorem~\ref{cv-covolume-bound-from-literature}, the upper and lower bounds for successive minima, a
lemma on antisymmetric perturbations, and a lower bound for heights on Grassmannians.

Step 1: proportionality and volumes. Use Definition~\ref{cv-reconstruction-functionals} for
$\Psi$, $N$ and the index $j$ of the tower, with $d=\dim_{\mathrm{fl}}$. Each $M(B,c)$ is a flag model
of type $(\Psi,N)$ (Definition~\ref{cv-condition-flag-model}) with $\deg_{B,c}=\deg_{M(B,c)}$ and
$\mathrm{PP}_{B,c}=\mathrm{PP}_{M(B,c)}$, and
$\operatorname{vol}(c)=\deg_{2,c}(\ell(\rho)^d)/\deg_{\mathrm{fl}}(\ell(\rho)^d)$
for the weight $\rho_{\tau,i}=N-i$. The origin of the tower in the moduli scheme
(Definition~\ref{cv-origin-tower}) gives, for each component $c$, an analytic system $S_c$ of type
$(\Psi,N)$ (Definition~\ref{cv-proportionality-analytic-system}): the complex points of the toroidal
compactification of the flag bundle over $c$, with the invariant metrics of the tautological line
bundles, singular along the boundary, and their regularizations; by
Definition~\ref{cv-origin-tower}(2c)(v) the singular metrics are induced by the Hodge metrics of the
root lines.
For every $B$ it gives a realization of $M(B,c)$ on $S_c$ whose multiplicity $n$ is the degree over
$\mathbf Q$ of the field of coefficients, and the integer $k_A$ of $M(B,c)$ does not depend on $B$
(Lemma~\ref{cv-origin-realization}). Each $S_c$ is uniformized of type $j$, with $j=p$: the period
maps of the component to the flag domain are the charts.

Fix $c$. By Lemma~\ref{cv-proportionality-realization} with $s=0$ and
Theorem~\ref{cv-proportionality-mumford},
$n\deg_{B,c}(\ell(\lambda)^d)=v_{S_c}\Delta(\lambda)$ for every $B$ and every weight $\lambda$. By
Lemma~\ref{cv-reconstruction-flag-degree}(1) the right side is
$n\kappa_c\deg_{\mathrm{fl}}(\ell(\lambda)^d)$ with
$\kappa_c=v_{S_c}\big(\prod_{a<b}(b-a)\big)^{|\Psi|}/(n\,d!)>0$. Taking $B=2$ and the regular
dominant weight $\lambda=\rho$ gives $\operatorname{vol}(c)=\kappa_c>0$. So the linear form
$F\mapsto\deg_{B,c}(F)-\operatorname{vol}(c)\deg_{\mathrm{fl}}(F)$ on $\mathbf R[x]$ vanishes on the
powers $\ell(\lambda)^d$, and on the monomials of degree other than $d$, since every monomial in the
definitions of $\deg_{B,c}$ and $\deg_{\mathrm{fl}}$ has degree $d$. By
Lemma~\ref{cv-reconstruction-polarization} it vanishes, that is,
$\deg_{B,c}(F)=\operatorname{vol}(c)\deg_{\mathrm{fl}}(F)$ for all $B$ and $F$. In the same way the
forms $\deg_{B,c}$ and $\deg_{B',c}$ on $\mathbf R[x,a]$ agree on every power $\ell(\lambda,s)^d$ by
Lemma~\ref{cv-proportionality-realization}, so they agree by
Lemma~\ref{cv-reconstruction-polarization}. This is the proportionality theorem of Hirzebruch
and Mumford \cite{Mum77} in the form used here.

For the equality of volumes, the complex points of the moduli scheme are a disjoint union of
quotients $\Gamma_k\backslash\mathcal B$ of the symmetric domain $\mathcal B$ of $U(p,q)^\Psi$, where
$\Gamma_k$ is the level subgroup of the unitary group of a Hermitian $\mathcal O_E$-lattice $L_k$ of
signature $(p,q)$ at every CM type. Computing the geometric degree of a component on its quotient
gives $c_0\operatorname{vol}(c)=\kappa v_c$, where $v_c$ is the invariant volume of a fundamental set
of the group of $c$, and $c_0,\kappa>0$ depend only on $|\Psi|$, $p$ and $q$. The volumes $v_c$
agree. The lattices $L_k$ are self-dual for the trace form attached to a fixed totally imaginary
element of $E$, so they lie in one genus, and their Hermitian spaces are isometric by Landherr's
theorem \cite{Lan36} and the Hasse norm theorem for the quadratic extension $E/E^+$. An isometry
carries the level subgroup of one lattice to a subgroup of the unitary group of the other with the
same index as the level subgroup of the other. Subgroups of the same index in a common group acting
on $\mathcal B$ have fundamental sets of the same invariant volume. So
$\operatorname{vol}(c)=\operatorname{vol}(c')$ for all components. This is (1).

Step 2: the conditions on the tower. There are constants, each written $C$, and a constant
$C_0\ge0$ such that, with $C'=4$:
\begin{enumerate}
\item[(i)] $|\mathrm{PP}_{B,c}(m)|\le C\log B$ for every monomial $m$;
\item[(ii)] for weights $\lambda$ of bounded $\ell^1$-norm and small $s>0$ with $\lambda-s\nu$
  dominant, the arithmetic self-intersection of the bundle of weight $\lambda-s\nu$ with its metric
  scaled by $e^{-1/s^2}$ differs from $\mathrm{PP}_{B,c}$ of the corresponding class by at most $Cs$,
  and that metric is semipositive;
\item[(iii)] the models are nef on the fibres at the primes $p\le B$ and at the bad primes for
  dominant weights, nef at the other primes for dominant weights beyond a threshold depending on $j$
  and $p$, and the tautological bundle is nef on all fibres;
\item[(iv)] for rational weights $\lambda$ with $\|\lambda\|\le1$ whose gaps at $j$ are at least
  $\theta$ and whose other gaps are bounded below, and for $s=\theta^2$ small, every point of the
  perturbed bundle has height at least $-C\log(B/s)$, and there is a Zariski-dense set of points of
  height at most $C\log B$;
\item[(v)] the models are nef on the generic fibre for dominant weights, and $\nu$ satisfies the
  effectivity condition of the perturbation;
\item[(vi)] exchanging the entries $a$ and $a+1$ of a weight $\lambda$ at one CM type, for $a\ne j$,
  changes the sign of $\mathrm{PP}_{B,c}(\ell(\lambda)^{\dim_{\mathrm{fl}}+1})-\operatorname{vol}(c)\,h(\ell(\lambda)^{\dim_{\mathrm{fl}}+1})$;
\item[(vii)] for $B_0\le B$, a CM type $\tau$, the power sum $p_{\tau,k}=\sum_ix_{\tau,i}^k$ and $Q$
  homogeneous of degree $\dim_{\mathrm{fl}}+1-k$,
  $\mathrm{PP}_{B_0,c}(p_{\tau,k}Q)=\mathrm{PP}_{B,c}(p_{\tau,k}Q)+\operatorname{vol}(c)\sum_p\delta_p(\tau,k,Q)$,
  over the primes $B_0<p\le B$ that are not bad, where $\delta_p$ is an explicit local term built from
  the Frobenius at $p$;
\item[(viii)] a group of permutations of $\Psi$ acting transitively preserves $\mathrm{PP}_B$;
\item[(ix)] $\mathrm{PP}_{B,c}(p_{\tau,1}Q)=h_{B,c,\tau}\operatorname{vol}(c)\deg_{\mathrm{fl}}(Q)$ with
  $|h_{B,c,\tau}|\le C'N\big(\log(B\operatorname{rd}_E)+|\frac{L'}{L}(1,\chi)|\big)$ for $B\ge128$ and $B$
  at least the number of bad primes; this rests on Theorem~\ref{cv-flat-degree};
\item[(x)] for $B$ as in (ix) with $B\ge4|\Psi|N^2$, the arithmetic intersection number
  $\mathrm{PP}_B$ of an explicit test polynomial of the type $j$ is at least
  $-C_0|\Psi|\min(j,N-j)\big(\log(B\operatorname{rd}_E)+|\frac{L'}{L}(1,\chi)|\big)$ times
  $\operatorname{vol}_T$ and an explicit degree; this is the lower bound at the base of the tower.
\end{enumerate}
% Lean: the steps are combined in GlobalTheory.CompanionTheory.Ar.Cusp.volumeTheoremX_of_inputs and
% GlobalTheory.CompanionTheory.Ar.Top.chain_tower (GlobalTheory.T.tower_weil_inputs_combined).
% Step 1 is real_tower_proportional (Ar.Top.rootTower_proportional_of_unif) and
% good_tower_intersection_extra. The conditions are the definitions RootTower.I1 to I8, A1,
% ModelNefStatement, M2, ANef, Prop6, PropA, PropF and SlopeLower.

The proofs of (i) to (x) below use Step 1, the realizations of Step 1, and properties of the origin
of the tower. Write $S_B$ for the set of primes $p\le B$ together with the bad primes, and
$y_{\tau,i}=-\sum_{k\le i}x_{\tau,k}$, so that
$\ell(\lambda)=\sum_{\tau,i}(\lambda_{\tau,i}-\lambda_{\tau,i+1})y_{\tau,i}$ with
$\lambda_{\tau,N+1}=0$.

Proof of (i). The origin gives, for each component $c$, a smooth Hermitian line bundle $\bar L_H$ on
$S_c$, a constant $C_l$ and a polynomial $P_c$ in the variables $(\lambda,u)$, and for each $B\ge2$
two realizations of $S_c$ on a common arithmetic variety, of the same degree, satisfying the
hypotheses of Lemma~\ref{cv-condition-comparison} with a set of primes $T$ and numbers $\alpha_p$
such that $\sum_{p\in T}\alpha_p\log p\le C_l\log B$; the limit polynomial of the first is
$P_{B,c}(\lambda,u)=n\,\mathrm{PP}_{B,c}(\ell(\lambda,k_Au)^{d+1})$, and that of the second is
$P_c$. By Lemma~\ref{cv-condition-comparison} the coefficients of $P_{B,c}$ and $P_c$ differ by at
most $C_2'C_l\log B$. For a monomial $m=a^e\prod x_{\tau,i}^{e_{\tau,i}}$ of degree $d+1$, the
coefficient of $u^e\prod\lambda_{\tau,i}^{e_{\tau,i}}$ in $P_{B,c}$ is $nk_A^e$ times the
coefficient of Lemma~\ref{cv-reconstruction-polarization} for the form $\mathrm{PP}_{B,c}$ and the
variables $x_{\tau,i}$ and $-a$, so $|\mathrm{PP}_{B,c}(m)|$ is at most its absolute value. With $A_0$
the largest absolute value of a coefficient of the finitely many $P_c$, this gives
$|\mathrm{PP}_{B,c}(m)|\le A_0+C_2'C_l\log B\le(A_0/\log2+C_2'C_l)\log B$ for $B\ge2$. Monomials of
other degrees give $0$.

Proof of (ii). Let $\|\lambda\|_1\le b_1$ and $0<s\le1$, and put $\mu=\lambda-s\nu$. The coefficients
$c(\mu,s)$ of Lemma~\ref{cv-proportionality-realization} satisfy
$\sum_b|c_b|\le2\|\mu\|_1+s\le b:=2(b_1+\|\nu\|_1)+1$. Lemma~\ref{cv-condition-convergence}, applied
to the finitely many $S_c$ with this $b$, gives $C$ and $s_0$, independent of $B$, with
$|\mathrm{PP}^{t}_{B,c}(\ell(\mu,s)^{d+1})-\mathrm{PP}_{B,c}(\ell(\mu,s)^{d+1})|\le Cs$ for
$t=e^{-1/s^2}$ and $s\le s_0$; the left side is the difference in (ii). For the semipositivity, the
origin is good, which gives three facts (Definition~\ref{cv-origin-tower}(3b), (3c) and (2b)): the
realizations are isomorphisms onto smooth generic fibres; the regularized metrics of the
determinant bundles $\bar L_{\tau,N}(t)$ have pluriharmonic potentials, so first Chern form $0$; and
the metric of $\bar L_A$ is positive, so, $S_c$ being compact, its first Chern form is at least
$c_0\omega$ for some $c_0>0$. If $\mu$ is
dominant, then $c_{\tau,i}\ge0$ for $i<N$ and $c_A=s/k_A>0$. On the complex points of $M(B,c)$ the
metric of $\ell(\mu,s)$ at $t$ has as first Chern form the pullback of that of
$\sum_bc_b\bar L_b(t)$, which is at least
$\big(sc_0/k_A-\varepsilon(t)\sum_\tau\sum_{i<N}c_{\tau,i}\big)\omega\ge(sc_0/k_A-b\,\varepsilon(t))\omega$
by Definition~\ref{cv-proportionality-analytic-system}. At $t=e^{-1/s^2}$ one has
$\varepsilon(t)=2C_cs^2\log(1/s)$, where $C_c$ is the constant of $S_c$, so this is nonnegative for
$s$ small, uniformly in $B$ because $k_A$ does not depend on $B$.

Proof of (iii). The bundle $\bar A$ of $M(B,c)$ has multidegree $1$ in one factor of the product of
projective spaces and $0$ in the others (Definition~\ref{cv-origin-tower}(2c)(ii)), so it has
nonnegative degree on every curve, and the
weight $\ell(0,1)=a$ is nef on all fibres. Let $p\in S_B$, and let $Z$ be the Hodge component of the
origin at $(B,c)$, with coordinates the point $y$ of $M(B,c)$ and a flag whose root lines have the
multidegrees $a_{\tau,i}$ of the bundles $\bar x_{\tau,i}(t)$ of $M(B,c)$. Its field $\Omega$ has
characteristic zero: it contains the residue field of a generic point of $\mathcal Y$, which lies over the
generic point of $\operatorname{Spec}\mathcal O_c$ as $\mathcal Y$ is flat. By (3a), the Hodge--Tate
modules of $Z$ at $p$ are vector bundles. From this and the exactness of the Hodge--Tate sequence
(Theorem~\ref{lit-hodge-tate-sequence}), for every $\tau$ the family $(a_{\tau,1},\dots,a_{\tau,N})$ has a
trivial flag modulo $p^{R_p-2}$ on $y$ over the ring $\mathbf Z_{(p)}\subset\Omega$ of fractions with
denominators prime to $p$ (Definition~\ref{cv-origin-trivial-flag}). As $R_p\ge4$, $p$ divides
$p^{R_p-2}$, so the family has a trivial flag modulo $p$ by
Lemma~\ref{cv-origin-trivial-flag-divisor}, and Lemma~\ref{cv-origin-trivial-flag-nef} shows that
$\ell(\mu)$ is nef on the fibre at $p$ for every dominant $\mu$. At a prime $p\notin S_B$, which is odd
and prime to the discriminant of $E$, the nefness comes from the good reduction of the moduli scheme at
$p$, for dominant $\mu$ beyond the threshold at $p$ (Definition~\ref{cv-reconstruction-functionals}(6)).
% Lean: at p in S_B, good_tower_M1_nef = Top.or_m1nef applies root_lines_weight_bound_of_free
% (= HT.rootModel_nefAt_of_hodgeCpt) with GoodOrigin.omVB (HTFree): trivFlagMod_rootLines_model gives
% TrivFlagMod modulo p^(R p - 2) for o = zpLocSubring, trivFlagMod_of_dvd passes to p, and
% trivflag_nef_fibre concludes.

Proof of (iv). The origin gives, for each component, a reference model of the flag bundle with an
ample $\bar A$ and the same analytic system, and at each prime of $S_B$ a Hodge--Tate comparison
between $M(B,c)$ and the reference model. The lower bound for heights of points and the Zariski-dense
set of points of small height are proved on the reference model and transported along these
comparisons, at a cost of $C\log B$.

Proof of (v). By the facts from the origin used in (ii), the realization of $M(B,c)$ on $S_c$ is
an isomorphism onto the smooth generic fibre, so its map $\tilde X\to X(\mathbf C)$ is bijective and
immersive everywhere, and the first Chern forms of the $\bar L_{\tau,N}(t)$ are $0$. So
Lemma~\ref{cv-nef-positivity} applies with $\Sigma_0=\{(\tau,N):\tau\in\Psi\}$. Let $W$ be a curve
in the generic fibre of $M(B,c)$. For $n$ equal to $1$ at one $b\in\Sigma$ and $0$ elsewhere, the
lemma gives $\deg_W\bar y_b\ge0$. For $n$ equal to $-1$ at one $(\tau,N)$ and $0$ elsewhere, it gives
$\deg_W\bar y_{\tau,N}\le0$, so $\deg_W\bar y_{\tau,N}=0$. For dominant $\mu$,
$\deg_W\ell(\mu)=\sum_{\tau}\sum_{i<N}(\mu_{\tau,i}-\mu_{\tau,i+1})\deg_W\bar y_{\tau,i}\ge0$, and
$\deg_W\ell(0,1)=\deg_W\bar y_A/k_A\ge0$, since $\bar y_A=\bar A$. The effectivity condition, that some positive multiple of
$\ell(\nu,-1)$ has a nonzero integral section on each $M(B,c)$, is Lemma~\ref{cv-origin-effective}.

Proof of (vi). For $a\ne j$, forgetting the $a$-th member of the flag at $\tau$ makes the flag
bundle a $\mathbf P^1$-bundle, and the block transposition of $x_{\tau,a}$ and $x_{\tau,a+1}$ is the
reflection of this $\mathbf P^1$-bundle. The arithmetic intersection theory of this
$\mathbf P^1$-bundle with its invariant metrics changes the sign of the self-intersection numbers up
to a term from the fibres, which is $\operatorname{vol}(c)$ times the corresponding term of the flag
height $h$ of Lemma~\ref{cv-reconstruction-flag-height}.

Proof of (vii). The origin gives, for every prime $p$, a permutation $\phi_p$ of
$\Psi\times\{+,-\}$, the Frobenius at $p$, such that the sign of $\phi_p^{n}(\tau,+)$ does not depend
on $\tau$; the term $\delta_p(\tau,k,Q)$ of (vii) is that of Lemma~\ref{cv-transport-local-terms} for
the sequences $\phi_{p,\tau,n}=\phi_p^{n+1}(\tau,+)$. For $2\le B_0\le B$ and a component $c$, the
origin gives a flag model $M'$ with the arithmetic intersection numbers of $M(B_0,c)$ on $\mathbf R[x]$
and its geometric degrees, and a change of $M'$ at the set $T$ of primes in $(B_0,B]$ that are not
bad, whose changed model has the arithmetic intersection numbers of $M(B,c)$ on $\mathbf R[x]$. At
each parameter $t$ the change alters $\mathrm{PP}^t(p_{\tau,k}Q)$ by
$\sum_{p\in T}\log p\,\deg_{M'}(\beta_{p}Q)$, where $\beta_p$ is $k$ times a combination of the
polynomials $\pi_{k-1}(\sigma,\pm)$ of Lemma~\ref{cv-transport-local-terms} with coefficients given by
$\phi_p$; this passes to the limit $t\to0$. By Step 1,
$\deg_{M'}(\beta_pQ)=\operatorname{vol}(c)\deg_{\mathrm{fl}}(\beta_pQ)$, and a computation with
$\deg_{\mathrm{fl}}$ gives $\log p\,\deg_{\mathrm{fl}}(\beta_pQ)=\delta_p(\tau,k,Q)$.

Proof of (viii). A group of automorphisms of the host field $E$ permutes the CM types
transitively. Each of its elements twists the universal object of the moduli problem, and the
twist is isomorphic to the universal object by an automorphism of the moduli scheme; this
automorphism permutes the components of the tower and the tautological bundles of the CM types.
The arithmetic intersection numbers, which exist by (ii), are preserved, so $\mathrm{PP}_B$, the
sum over the components, is invariant under the induced permutations of $\Psi$.

Proof of (ix). For each component $c$ and $\tau\in\Psi$, the origin gives a Hermitian line bundle
$\bar A_\tau$ on $M(2,c)$ with a flat section over a number field and with the multidegree of
$\sum_ix_{\tau,i}$, such that the metric of $\sum_i\bar x_{\tau,i}(t)$ is that of $\bar A_\tau$
multiplied by $e^{-v_t}$ with $\int v_t$ against products of first Chern forms tending to $0$. By the
recursion of the intersection system applied to the flat section,
$\mathrm{PP}_{2,c}(p_{\tau,1}Q)=\operatorname{ht}(\bar A_\tau)\deg_{2,c}(Q)
=\operatorname{ht}(\bar A_\tau)\operatorname{vol}(c)\deg_{\mathrm{fl}}(Q)$ by Step 1, where
$\operatorname{ht}(\bar A_\tau)$ is the height of the flat section. By
Theorem~\ref{cv-flat-degree} this height has the form $h_0+\sum_{p\text{ bad}}d_p$ with
$|h_0|\le N(\frac32\log\operatorname{rd}_E+\frac12|\frac{L'}L(1,\chi)|)$,
$|d_p|\le(N/(p-1)+Nw_p)\log p$, $w_p\ge0$ and $\sum_pw_p\log p\le\log\operatorname{rd}_E$. For $k=1$
the term $\delta_p(\tau,1,Q)$ of (vii) is $e_p\log p\deg_{\mathrm{fl}}(Q)$ with $|e_p|\le N/(p-1)$,
since $\pi_0(\sigma,+)=j$ and $\pi_0(\sigma,-)=N-j$. So (vii) with $B_0=2$ gives (ix) with
$h_{B,c,\tau}=\operatorname{ht}(\bar A_\tau)-\sum_pe_p\log p$, the sum over the primes in $(2,B]$
that are not bad. For the bound, $\sum_{p\in S}\log p/(p-1)\le3\log B+1$ for every set $S$ of primes
each at most $B$ or bad, when $B\ge128$ and there are at most $B$ bad primes: the primes $p\le B$
contribute at most $\log B+17/5$, by an estimate of Mertens's type, the at most $B$ others at most
$\log(B+1)$, and $\log B\ge4.85$. Adding up,
$|h_{B,c,\tau}|\le N\big(\frac32\log\operatorname{rd}_E+\frac12|\frac{L'}L(1,\chi)|\big)
+N(3\log B+1)+N\log\operatorname{rd}_E\le4N\big(\log(B\operatorname{rd}_E)+|\frac{L'}L(1,\chi)|\big)$.

Proof of (x). Let $t_j$, $T_j$, $D_j$ and $m_j$ be the test polynomial, the test degree and the
numbers of Definition~\ref{cv-transport-test-polynomial}; $t_j$ is the test polynomial of (x), and
the degree there is $(D_j+1)T_j$. Write
$\Lambda_B=\log(B\operatorname{rd}_E)+|\frac{L'}L(1,\chi)|$. For each $B$ as in (x) and each
component $c$, the origin gives a setting for the test class: an arithmetic subvariety $Y$ of
dimension $n+1$ of a product of projective spaces, and Hermitian line bundles $\bar A_{r,t}$ on $Y$
for rational $r>0$ and small $t>0$, which perturb the bundle of the test class by $r$ times a
boundary bundle and regularize its metric at $t$. These are semipositive, invariant under
conjugation and nef on the fibres, and by the bounds for heights in the interior and near the
boundary every point of $Y$ has height at least $-C_h|\Psi|m_j\Lambda_B-rK$ for $\bar A_{r,t}$,
with a constant $K$. By Theorem~\ref{cv-wall-successive-minima}(1) the normalized heights
$(\bar A_{r,t}^{n+1}\mid Y)/((n+1)\deg_{A_{r,t}}(Y))$ satisfy the same bound. The setting also gives
$\mathrm{PP}_{B,c}(t_j)=\sigma_c(D_j+1)\operatorname{vol}(c)T_j$, where $\sigma_c$ is the limit, by
Stokes's formula, of these normalized heights as $r$ and $t$ tend to $0$, corrected by a mean over
the cusps that is at least $-\frac{C_u}2|\Psi|m_j\Lambda_B$. Hence
$\sigma_c\ge-(C_h+C_u/2)|\Psi|m_j\Lambda_B$. Since $\operatorname{vol}(c)>0$ and $T_j>0$
(Lemma~\ref{cv-transport-test-facts}), summing over $c$ gives (x) with $C_0=\max(C_h+C_u/2,0)$.
\openstep{The derivation from the good origin (Definition~\ref{cv-origin-tower}) of the other
properties used above is not yet written out: the uniformization of $S_c$ (Step 1); the computation
of $\operatorname{vol}(c)$ on the ball quotient (Step 1); the reference realizations of (i); in (iii),
the trivial flag modulo $p^{R_p-2}$ of the root lines at a prime $p\in S_B$, from the Hodge--Tate
modules, and the nefness at the primes outside $S_B$; the reference models and Hodge--Tate
comparisons of (iv); the $\mathbf P^1$-bundle antisymmetry of (vi); the change of model at the primes
of $(B_0,B]$ and its effect on $\mathrm{PP}^t$ in (vii); the action on the tower in (viii); the flat
bundles $\bar A_\tau$, the pushforward identity and the form of their heights in (ix); and the
settings of (x).}
% Lean, for the next writer. Written out: the origin (Definitions cv-origin-hodge-component and
% cv-origin-tower), realOf o.toTowerOrigin (Lemma cv-origin-realization), the facts
% GoodOrigin.smooth, .flat and FixedData.pos used in (ii), RootOrigin.A_deg with root_model_nef_fibres
% in (iii), the passage from TrivFlagMod to nefness at p in S_B in (iii) (Definition
% cv-origin-trivial-flag, Lemmas cv-origin-trivial-flag-divisor, -section, -degrees and -nef), and
% global_model_tower_bounds (root_origin_weight_identity = Lemma cv-origin-effective)
% in (v). Open: Step 1: GlobalTheory.T.good_tower_real_agreement
% (RootFieldNumerics.global_model_real_tower: Un.unif_of from the reference origin refOr, ref_top,
% ref_deg, ample, with Biholo, LocPer, Transport, FrameVals, PerMat, FlagImm, CountGen, PotLoc and
% PtsFam, all in WeilBound/CoreStatements/FlagImmersions) for the uniformization;
% SheetDecomposition.volEq_of with DegPush, degMu_eq and BallVolEq (famVol_of: MoebVol,
% StdModel, ArithIdx = arithIdx, LevelDict; WeilBound/CoreStatements/VolumeEquality) for the
% volumes. (i) real_tower_I1 =
% TowerOverReals.tower_I1_of_reference_linked, with RefLinked from good_tower_real_I6
% (RootFieldNumerics.global_model_real_tower_smooth: val_bound, OrRef.exists_comparison,
% OrRef.sum_bound, the sandwiches below and level, from ResExistsStatement). (iii) at p in S_B:
% HT.trivFlagMod_rootLines_model (= PositivityTheory.HtPt.Ge.triv_flag_mod_of_hodge_tate_free_of_root_lines,
% from FarguesHodgeTate through Fargues.pointHT, HT.hodgeTate_family_exists, Famproj.family_proj,
% Flagb.chart_flag_bases, chart_sequence, hodge_component_integral_at_prime_in_S,
% hodge_chart_transition_integral and trivflag_charts); outside S_B: good_tower_M2
% (hodge_cpt_global_nef_bound = Conj.hodge_cpt_global_nefness: global_model_power_free,
% root_lines_unramified_primes = HcptTie.drTie_of_hodgeCpt, rootLineDeRhamModel_holds = Conj.m2Model_holds;
% it uses IsGlobal). (iv) good_tower_I5
% (TopGood.tower_good_of_sandwich_split, exists_sandwich_of_hodge_tate_free). (vi)
% full_intersection_tower_I8_identity (SquareStatements.good_tower_antisymmetry). (vii)
% good_tower_global_conditions (Link.goodOrigin_linked: link_step, with R p = 4 from o.Extra, LinkedTower,
% FrobShape) and linked_tower_power_sum (power_sum_height_change_model, log_mul_degFl_bracket_mul).
% (viii) good_tower_propA_of_global_extra (TowerPropA.propA_of_origin: hodge_tie_galois with
% TwistAut and CmpPair, the bijection comp_inj, comp_onto with the irreducible components, and R p = 4).
% (ix) PropFTop_holds (Flat.propFTop, root_flat_pushforward_geoFun,
% propFTopE_spec, sum_log_prime_div_pred_bound, hB_bound) with good_tower_real_extra (NearRecord,
% from FlatDegree). (x) slope_lower_base_settings (ppc_lower_of_setting, slX_lower_of_stokesAt,
% ZhangLowerX) with base_setting_of_good_tower (Cusp.TopBase.base_setting_of_good_origin).

Step 3: nef perturbations, and (2). Write $d=\dim_{\mathrm{fl}}$, and use
Definition~\ref{cv-reconstruction-functionals} for $\Psi$, $N$ and the index $j$ of the tower. Let $a$ be
the class of the bundle that (iii) calls tautological, so that the bundle of weight $\lambda-s\nu$ in
(ii) has class $\ell(\lambda-s\nu)+sa$; extend $\deg_{B,c}$ and $\mathrm{PP}_{B,c}$ to polynomials in
the $x_{\tau,i}$ and $a$. For a weight $\lambda$ and $s>0$ let $\bar m_{B,c}(\lambda,s)$ be the
arithmetic self-intersection number of the bundle of weight $\lambda-s\nu$ on $M(B,c)$ with its metric
scaled by $e^{-1/s^2}$, and let $m_{B,c}(\lambda,s)=\deg_{B,c}((\ell(\lambda-s\nu)+sa)^d)$. So (ii) bounds
$|\bar m_{B,c}(\lambda,s)-\mathrm{PP}_{B,c}((\ell(\lambda-s\nu)+sa)^{d+1})|$ by $Cs$. The threshold of
(iii) at a prime $p$ is the one of Definition~\ref{cv-reconstruction-functionals}(6) at $b=p$.

(3a) For every $A\ge0$ there is $s_0>0$ such that the following holds for every $B\ge2$, every
component $c$, every weight $\lambda$ with $\|\lambda\|_1\le A$ and every $0<s\le s_0$ for which
$\lambda-s\nu$ is dominant and beyond the threshold at $B$: the bundle of weight $\lambda-s\nu$ is nef on
every fibre of $M(B,c)$ and semipositive. Take $s_0$ from (ii), which gives the semipositivity. Let
$p$ be a prime. If $p\le B$ or $p$ is bad, a bundle of class $\ell(\lambda-s\nu)$ is nef on the fibre at
$p$ by (iii). Otherwise $p>B$; as $\operatorname{gap}_j(\lambda-s\nu,\tau)\ge0$ by dominance, the right
side of the threshold condition grows with $b$, so $\lambda-s\nu$ is beyond the threshold at $p$ and
(iii) applies again. The class $a$ is nef on the fibre at $p$ by (iii), and degrees on curves are
linear in the class, so $\ell(\lambda-s\nu)+sa$ is nef on the fibre at $p$.
% Lean: GlobalTheory.T.tower_I3_of_semipositive_nef (RootTower.I3), with
% RootTower.nefFib_of_nefThreshold and GlobalTheory.T.curve_degree_linear_nef_cone.

(3b) Let $u$ be regular dominant with gaps at least $g>0$, and let $B+1>4|\Psi|N^2(1+\|u\|/g)$. Then
$B\ge2$, since $4|\Psi|N^2\ge16$. For $s$ near $0$ the weight $u-s\nu$ is dominant, since the
inequalities between the entries of $u$ are strict. It is also beyond the threshold at $B$ for $s$
near $0$: both sides of the threshold condition are continuous in $s$, and at $s=0$
\[
 (B-2)\operatorname{gap}_j(u,\tau)\ge(B-2)g>4|\Psi|N^2g+4|\Psi|N^2\|u\|-3g>4|\Psi|\,\|u\|.
\]
By (3a), the bundle of weight $u-s\nu$ is nef on the fibres and semipositive for small $s>0$. By (ii),
$\bar m_{B,c}(u,s)-\mathrm{PP}_{B,c}((\ell(u-s\nu)+sa)^{d+1})$ tends to $0$ as $s\to0$. As
$\ell(u-s\nu)+sa=\ell(u)+s(a-\ell(\nu))$, the second term is a polynomial in $s$ with value
$\mathrm{PP}_{B,c}(\ell(u)^{d+1})$ at $s=0$. So $\bar m_{B,c}(u,s)$ tends to
$\mathrm{PP}_{B,c}(\ell(u)^{d+1})$. In the same way $m_{B,c}(u,s)$ tends to
$\deg_{B,c}(\ell(u)^d)=\operatorname{vol}(c)\deg_{\mathrm{fl}}(\ell(u)^d)$, by Step 1, and this limit is
positive by Lemma~\ref{cv-reconstruction-flag-degree}(1). Hence the slopes
$\bar m_{B,c}(u,s)/((d+1)m_{B,c}(u,s))$ tend to
$\mathrm{PP}_{B,c}(\ell(u)^{d+1})/((d+1)\operatorname{vol}(c)\deg_{\mathrm{fl}}(\ell(u)^d))$, the slope of
weight $u$ at level $B$. This is (2).
% Lean: GlobalTheory.T.tower_T1_limit (RootTower.T1), through LimitVolumeRatio.eventually_dominant,
% LimitVolumeRatio.eventually_threshold and LimitVolumeRatio.tendsto_lin.

Step 4: walls. Fix $c_0>0$. The claim is that there are $C_W$ and $B_1$ with
$|\mathrm{PP}_B(\ell(\lambda)^{d+1})|\le C_W\log B/B$ for every $B\ge B_1$ and every weight $\lambda$ on
the wall at some $\tau_0$ with margin $c_0$ (Definition~\ref{cv-reconstruction-functionals}(5)).

(4a) Heights. There are $\theta_0>0$ and $C_s\ge0$ such that the following holds for every $B\ge2$,
every $c$, every weight $\lambda$ with rational entries and $\|\lambda\|\le1$, and every
$0<\theta\le\theta_0$ with $s=\theta^2$ rational, such that $\operatorname{gap}_j(\lambda,\tau)\ge\theta$
for every $\tau$, $\lambda_{\tau,i}-\lambda_{\tau,i+1}\ge c_0/2$ for every $\tau$ and every $i\ne j$, the
weight $\lambda-s\nu$ is dominant and beyond the threshold at $B$, and $m_{B,c}(\lambda,s)\ge0$:
\[
 |\bar m_{B,c}(\lambda,s)|\le C_s\log(B/s)\,m_{B,c}(\lambda,s).
\]
Take $\theta_0\le1$, at most the bound of (iv) for the lower bound $c_0/2$ on the other gaps, and at
most the $s_0$ of (3a) for $A=|\Psi|N$. By (3a) the bundle of weight $\lambda-s\nu$ is nef on the
fibres and semipositive, since $\|\lambda\|_1\le|\Psi|N$ and $s\le\theta\le s_0$; as $\lambda$, $\nu$
and $s$ are rational, a positive multiple of it is a metrized line bundle. By (iv), every point has
height at least $-C\log(B/s)$, and the points of a Zariski-dense set have height at most $C\log B$.
Theorem~\ref{cv-wall-successive-minima} gives
\[
 -(d+1)C\log(B/s)\,m_{B,c}(\lambda,s)\le\bar m_{B,c}(\lambda,s)\le(d+1)C\log B\;m_{B,c}(\lambda,s),
\]
and $0\le\log B\le\log(B/s)$ since $s\le1$. Take $C_s=2(d+1)\max(C,0)$.
% Lean: GlobalTheory.T.tower_gap_height_uniform_bound (mbarTop_squeeze), from
% GlobalTheory.T.zhang_lower_bound_rational and GlobalTheory.T.zhang_upper_bound_dense_points.

(4b) Degrees. For a dominant weight $\lambda$ with rational entries and rational $s\ge0$ with
$\lambda-s\nu$ dominant, $0\le m_{B,c}(\lambda,s)\le\operatorname{vol}(c)\deg_{\mathrm{fl}}(\ell(\lambda)^d)$.
By (v), on the generic fibre of $M(B,c)$ the classes $\ell(\mu)$ for dominant $\mu$ and the class $a$
are nef, and a positive multiple of $\ell(\nu)-a$ is the class of an effective divisor. So
$\ell(\lambda-s\nu)+sa$ is nef there, and its $d$-th power has nonnegative degree by
Theorem~\ref{cv-wall-kleiman}. Since $\ell(\lambda)-(\ell(\lambda-s\nu)+sa)=s(\ell(\nu)-a)$,
\[
 \ell(\lambda)^d-(\ell(\lambda-s\nu)+sa)^d
 =\sum_{i=0}^{d-1}\ell(\lambda)^i(\ell(\lambda-s\nu)+sa)^{d-1-i}\,s(\ell(\nu)-a),
\]
and each term has nonnegative degree, by Theorem~\ref{cv-wall-kleiman} on the components of the
effective divisor. Finally $\deg_{B,c}(\ell(\lambda)^d)=\operatorname{vol}(c)\deg_{\mathrm{fl}}(\ell(\lambda)^d)$ by
Step 1.
% Lean: GlobalTheory.T.tower_component_height_bound (RootFieldNumerics.tower_rational_weight_bound).

(4c) Rational weights. Let $C_1$ be the constant of (i), let $C_2$ and $s_2$ be the constant and the
bound for $s$ in (ii) for weights of $\ell^1$-norm at most $|\Psi|N$, let $C_8$ be the constant of
Lemma~\ref{cv-reconstruction-flag-degree}(1), let $C_p$ be the constant of
Lemma~\ref{cv-wall-lipschitz} for $D=d+1$, and let $B_1$ be the bound of
Lemma~\ref{cv-wall-perturbed-weights} for the margin $c_0/2$ and for $\min(\theta_0,s_2)$. Let $B\ge B_1$
and let $\lambda_0$ be a weight with rational entries on the wall at $\tau_0$ with margin $c_0/2$. Put
$\theta=5|\Psi|/B$, let $\lambda$ be $\lambda_0$ with $2\theta$ added to its entries at $(\tau_0,i)$ for
$i\le j$, and put $s=\theta^2$; these are rational. Lemma~\ref{cv-wall-perturbed-weights} gives the
hypotheses of (4a) and (4b) for $\lambda$ and $s$, with $m_{B,c}(\lambda,s)\ge0$ from (4b), and also
$B\ge3$, $\operatorname{gap}_j(\lambda,\tau_0)=2\theta$ and
$\|\lambda-s\nu-\lambda_0\|\le2\theta+\|\nu\|\theta^2$. For each $c$ write
$q_c=\mathrm{PP}_{B,c}(\ell(\lambda_0)^{d+1})$ and $p_c=\mathrm{PP}_{B,c}((\ell(\lambda-s\nu)+sa)^{d+1})$.
Since $B\ge3$, $\log B\ge1$; also $\theta\le1$ and $\log(B/s)=\log(B^3/(5|\Psi|)^2)\le3\log B$. Then:
\begin{itemize}
\item $|\bar m_{B,c}(\lambda,s)|\le C_s\cdot3\log B\cdot\operatorname{vol}(c)C_8\cdot2\theta$, by (4a),
  (4b) and Lemma~\ref{cv-reconstruction-flag-degree}(1);
\item $|\bar m_{B,c}(\lambda,s)-p_c|\le C_2s\le|C_2|\,\theta\log B$, by (ii);
\item $|p_c-q_c|\le C_pC_1\log B\,(\|\lambda-s\nu-\lambda_0\|+s)\le C_pC_1(3+\|\nu\|)\,\theta\log B$, by
  Lemma~\ref{cv-wall-lipschitz} with $b=C_1\log B$, which (i) allows.
\end{itemize}
Hence $|q_c|\le(|C_2|+6C_sC_8\operatorname{vol}(c)+C_pC_1(3+\|\nu\|))\,\theta\log B$. Summing over $c$
and inserting $\theta=5|\Psi|/B$ gives $|\mathrm{PP}_B(\ell(\lambda_0)^{d+1})|\le C_W\log B/B$ with
$C_W=5|\Psi|\sum_c(|C_2|+6C_sC_8\operatorname{vol}(c)+C_pC_1(3+\|\nu\|))$.

(4d) All weights. The function $\lambda\mapsto\mathrm{PP}_B(\ell(\lambda)^{d+1})$ is continuous by
Lemma~\ref{cv-wall-lipschitz}, so Lemma~\ref{cv-wall-rational-weights} extends the bound of (4c) from
the weights with rational entries on walls with margin $c_0/2$ to all weights on walls with margin
$c_0$. This proves the claim.
% Lean: GlobalTheory.T.tower_weil_theorem (Ar.Top.thmW_tower, RootTower.ThmW), with
% GlobalTheory.T.perturbed_weight_gap_bounds, GlobalTheory.T.arith_functional_lipschitz_weights,
% GlobalTheory.T.flag_degree_le_gap and GlobalTheory.T.continuous_wall_bound_extend.

Step 5: (3). For $\tau\in\Psi$ let $x^{(\tau)}$ be the weight with entries
$x^{(\tau)}_{\sigma,i}=-(i-1-\epsilon_{\sigma,i})/(4(N+1))$, where $\epsilon_{\sigma,i}$ is $1$ if
$\sigma=\tau$ and $i>j$, and $0$ otherwise. Let $K_\tau$ be the set of weights $\lambda$ with
$\lambda_{\tau,j}=\lambda_{\tau,j+1}$ and $\|\lambda-x^{(\tau)}\|\le1/(16(N+1))$, a compact subset of the
hyperplane $H_\tau$ of Lemma~\ref{cv-reconstruction-antisymmetric} that contains a nonempty open subset
of it. The entries of $x^{(\tau)}$ decrease in steps of $1/(4(N+1))$, except from $(\tau,j)$ to
$(\tau,j+1)$, where they are equal, and $\|x^{(\tau)}\|\le\frac14$. So every $\lambda\in K_\tau$ is
dominant and on the wall at $\tau$ with margin $c_0=1/(8(N+1))$. By
Lemma~\ref{cv-reconstruction-flag-height}(1), $h(\ell(\lambda)^{d+1})=0$ for $\lambda\in K_\tau$, and by
Step 4 there are $C_W$ and $B_1\ge2$ with $|\mathrm{PP}_B(\ell(\lambda)^{d+1})|\le C_W\log B/B$ for
$\lambda\in K_\tau$ and $B\ge B_1$.

Let $B\ge B_1$. Let $P_h$ and $P_B$ be the polynomials that Lemma~\ref{cv-reconstruction-polarization}
attaches to $h$ and to $\mathrm{PP}_B$ with $D=d+1$, so that $P_h(\lambda)=h(\ell(\lambda)^{d+1})$ and
$P_B(\lambda)=\mathrm{PP}_B(\ell(\lambda)^{d+1})$. Then $P_h$ vanishes on every $K_\tau$, and
$|P_B|\le C_W\log B/B$ there. The polynomial $P_B-\operatorname{vol}_TP_h$ is homogeneous of degree
$d+1$, and it changes sign under every block transposition: its value at $\lambda$ is
$\sum_c\big(\mathrm{PP}_{B,c}(\ell(\lambda)^{d+1})-\operatorname{vol}(c)h(\ell(\lambda)^{d+1})\big)$, and by
(vi) each term changes sign when $\lambda_{\tau,a}$ and $\lambda_{\tau,a+1}$ are exchanged, for every
$\tau$ and every $a\ne j$. Lemma~\ref{cv-reconstruction-antisymmetric}, applied with $h=P_h$,
$P=P_B$, $v=\operatorname{vol}_T$ and $\varepsilon=C_W\log B/B$, gives a constant $C_A$ that does not
depend on $B$, and a block invariant linear form $l$ such that every coefficient of
$P_B-\operatorname{vol}_TP_h-\Delta l$ is at most $C_AC_W\log B/B$ in absolute value. The polynomial
$P_B-\operatorname{vol}_TP_h$ represents the linear form $\mathrm{PP}_B-\operatorname{vol}_Th$ in the sense
of Lemma~\ref{cv-reconstruction-coefficients}, which, with $v=\operatorname{vol}_T>0$, gives
$\alpha_\tau,\beta_\tau$ with
\[
 \Big|\mathrm{PP}_B(m)-\operatorname{vol}_T\Big(h(m)-\sum_\tau(\alpha_\tau-\beta_\tau)\,\mathrm{Sh}_\tau(m)
 -\sum_\tau\beta_\tau\,\mathrm{Tot}_\tau(m)\Big)\Big|\le C_AC_W\frac{\log B}{B}
\]
for every monomial $m$ of degree $d+1$, with $\mathrm{Sh}_\tau$ and $\mathrm{Tot}_\tau$ as in
Definition~\ref{cv-reconstruction-functionals}(3). This is (3).
% Lean: GlobalTheory.T.tower_theorem_B (Ar.Top.thmB_functional, RootFieldNumerics.antisymmetry_tower_theorem,
% RootTower.ThmB), with GlobalTheory.T.flag_height_I9_properties and
% GlobalTheory.T.antisymmetric_perturbation_holds.

The $\alpha_\tau$ and $\beta_\tau$ can be taken independent of $\tau$. By (viii) there is a group
$\Gamma$ of permutations of $\Psi$ that acts transitively and satisfies $\mathrm{PP}_B(F^\pi)=\mathrm{PP}_B(F)$
for $\pi\in\Gamma$, with $F^\pi$ as in Lemma~\ref{cv-reconstruction-flag-degree}(4); and $h(F^\pi)=h(F)$ by
Lemma~\ref{cv-reconstruction-flag-height}(4). For $\pi\in\Gamma$, the estimate for the monomial $m^\pi$
becomes, by Lemma~\ref{cv-reconstruction-flag-degree}(4), the estimate for $m$ with $\alpha_{\pi(\tau)}$
and $\beta_{\pi(\tau)}$ in place of $\alpha_\tau$ and $\beta_\tau$. The expression inside the absolute
value is affine in $(\alpha,\beta)$, so the estimate also holds with the averages over $\pi\in\Gamma$.
As $\Gamma$ is transitive, these averages are the means $\alpha_B$ and $\beta_B$ of the $\alpha_\tau$ and
the $\beta_\tau$ over $\Psi$. Thus, with $C_5=C_AC_W$, for every $B\ge B_1$ there are $\alpha_B,\beta_B$
with
\[
 \Big|\mathrm{PP}_B(m)-\operatorname{vol}_T\Big(h(m)-(\alpha_B-\beta_B)\sum_\tau\mathrm{Sh}_\tau(m)
 -\beta_B\sum_\tau\mathrm{Tot}_\tau(m)\Big)\Big|\le C_5\frac{\log B}{B}
\]
for every monomial $m$ of degree $d+1$.
% Lean: GlobalTheory.T.tower_theorem_B_averaged (thmBA_of_thmB_propA, RootTower.ThmBA), with
% GlobalTheory.T.flag_height_label_invariant.

Step 6: (4). Put $\Lambda_B=\log(B\operatorname{rd}_E)+|\frac{L'}{L}(1,\chi)|$; as
$\operatorname{rd}_E\ge1$, $\Lambda_B\ge\log B$. Let $t_j$, $D_j$, $m_j$ and $T_j$ be as in
Definition~\ref{cv-transport-test-polynomial}, so that (ix) and (x) read $|h_{B,c,\tau}|\le C'N\Lambda_B$
and $\mathrm{PP}_B(t_j)\ge-C_0|\Psi|m_j\Lambda_B(D_j+1)\operatorname{vol}_TT_j$. The local terms $\delta_p$
of (vii) are those of Lemma~\ref{cv-transport-local-terms}, for Frobenius data whose signs do not
depend on $\tau$. Let $r_p$ and $R_{B_0}$ be as in that lemma, with $S$ the set of bad primes, and
$R_{B_0}(F)$ as in Lemma~\ref{cv-transport-remainders}.

(6a) The identity at a base level. Let $B_0\ge2$. There are real numbers $A_0,b_0$ such that
\[
 \mathrm{PP}_{B_0}(p_{\tau,k}Q)=\operatorname{vol}_T\Big(h(p_{\tau,k}Q)-A_0\sum_\sigma\mathrm{Sh}_\sigma(p_{\tau,k}Q)
 -b_0\sum_\sigma\mathrm{Tot}_\sigma(p_{\tau,k}Q)+R_{B_0}(\tau,k,Q)\Big)
\]
for every $\tau$, every $1\le k\le N$ and every homogeneous $Q$ of degree $d+1-k$. Let
$B\ge\max(B_0,B_1)$, with $B_1$ and $\alpha_B,\beta_B$ from Step 5. Summing (vii) over the components
gives $\mathrm{PP}_{B_0}(p_{\tau,k}Q)=\mathrm{PP}_B(p_{\tau,k}Q)+\operatorname{vol}_T\sum_p\delta_p(\tau,k,Q)$, the
sum over the good primes $B_0<p\le B$. Insert the decomposition of $\delta_p$ from
Lemma~\ref{cv-transport-local-terms}, and the estimate of Step 5 applied to the monomials of
$p_{\tau,k}Q$. With $a_B=\alpha_B-\beta_B+\sum_p\frac{\log p}{p-1}$ and $b_B=\beta_B-\sum_p\kappa_p$, both
sums over the good primes in $(B_0,B]$, this gives
\[
 \mathrm{PP}_{B_0}(p_{\tau,k}Q)=\operatorname{vol}_T\Big(h(p_{\tau,k}Q)-a_B\sum_\sigma\mathrm{Sh}_\sigma(p_{\tau,k}Q)
 -b_B\sum_\sigma\mathrm{Tot}_\sigma(p_{\tau,k}Q)+\sum_pr_p(\tau,k,Q)\Big)+E_B,
\]
where $|E_B|$ is at most $C_5\log B/B$ times the sum of the absolute values of the coefficients of
$p_{\tau,k}Q$, so that $E_B\to0$ as $B\to\infty$. Fix $\tau_0\in\Psi$, a homogeneous $Q_1$ of degree $d$
with $\deg_{\mathrm{fl}}(Q_1)=1$, and a homogeneous $Q_2$ of degree $d-1$ with
$\deg_{\mathrm{fl}}(p^{\le j}_{\tau_0,1}Q_2)=1$; they exist because $\deg_{\mathrm{fl}}$ and
$Q\mapsto\deg_{\mathrm{fl}}(p^{\le j}_{\tau_0,1}Q)$ are nonzero linear forms. By
Lemma~\ref{cv-reconstruction-flag-degree}(3), $\sum_\sigma\mathrm{Sh}_\sigma(p_{\tau_0,2}Q_2)=2$,
$\sum_\sigma\mathrm{Tot}_\sigma(p_{\tau_0,2}Q_2)=0$ and $\sum_\sigma\mathrm{Tot}_\sigma(p_{\tau_0,1}Q_1)=N$.
The partial sums of the $r_p$ converge by Lemma~\ref{cv-transport-local-terms}, and
$\operatorname{vol}_T>0$. So the identity for $(\tau_0,2,Q_2)$ shows that $a_B$ converges as
$B\to\infty$, to $A_0$ say, and then the identity for $(\tau_0,1,Q_1)$ shows that $b_B$ converges, to
$b_0$ say. Letting $B\to\infty$ in the identity for any $(\tau,k,Q)$ proves (6a).
% Lean: GlobalTheory.T.tower_identity_at_cutoff (exists_identAt, RootTower.IdentAt), with
% exists_finite_identity, GlobalTheory.T.exists_flag_degree_one_polys, tendsto_coeffs_of_tests and
% eq_of_tendsto_coeffs.

(6b) Let $A_0,b_0$ be as in (6a). For every homogeneous $F$ of degree $d+1$ and every decomposition
$F=\sum_{\tau,k}p_{\tau,k}Q_{\tau,k}$ as in Lemma~\ref{cv-transport-power-sums}, the identities of (6a) add
up to
\[
 \mathrm{PP}_{B_0}(F)=\operatorname{vol}_T\Big(h(F)-A_0\sum_\sigma\mathrm{Sh}_\sigma(F)-b_0\sum_\sigma\mathrm{Tot}_\sigma(F)+R_{B_0}(F)\Big).
\]
% Lean: GlobalTheory.T.tower_height_identity_expansion (identOnDecomposedForm).

(6c) Let $B_0\ge128$ be at least the number of bad primes, and let $A_0,b_0$ be as in (6a). Then
$|jA_0+Nb_0|\le C'N\Lambda_{B_0}$. Take $\tau_0$ and $Q_1$ as in (6a). By
Lemma~\ref{cv-reconstruction-flag-height}(3), $h(p_{\tau_0,1}Q_1)=0$; by
Lemma~\ref{cv-reconstruction-flag-degree}(3), $\sum_\sigma\mathrm{Sh}_\sigma(p_{\tau_0,1}Q_1)=
\deg_{\mathrm{fl}}(p^{\le j}_{\tau_0,0}Q_1)=j$ and $\sum_\sigma\mathrm{Tot}_\sigma(p_{\tau_0,1}Q_1)=N$; and
$R_{B_0}(\tau_0,1,Q_1)=0$, since $r_p$ vanishes for $k=1$. So (6a) reads
$\mathrm{PP}_{B_0}(p_{\tau_0,1}Q_1)=-\operatorname{vol}_T(jA_0+Nb_0)$. By (ix) the left side is
$\sum_ch_{B_0,c,\tau_0}\operatorname{vol}(c)$, whose absolute value is at most
$C'N\Lambda_{B_0}\operatorname{vol}_T$.
% Lean: GlobalTheory.T.cutoff_coefficients_bound (flat_85_bound).

(6d) The lower bound at the base. Let $B_0\ge128$ with $B_0\ge4|\Psi|N^2$ and $B_0$ at least the number
of bad primes, and let $A_0,b_0$ be as in (6a). Then $A_0\ge-(2(C_0+C')+3)\Lambda_{B_0}$. Put $f=|\Psi|$,
$K=(D_j+1)T_j$, which is positive by Lemma~\ref{cv-transport-test-facts}, and $l=\log(N+1)$. By (x),
(6b) for $F=t_j$, Lemma~\ref{cv-transport-grassmannian}, which gives $h(t_j)\le Kf\frac{m_j}2l$, and
Lemma~\ref{cv-transport-remainders}(2), which gives $|R_{B_0}(t_j)|\le1.6K$,
\[
 -C_0fm_j\Lambda_{B_0}K\le Kf\frac{m_j}2l-A_0\sum_\sigma\mathrm{Sh}_\sigma(t_j)-b_0\sum_\sigma\mathrm{Tot}_\sigma(t_j)+1.6K.
\]
If $N-j\le j$, then $m_j=N-j$, and by Lemma~\ref{cv-transport-test-facts} the sums of $\mathrm{Sh}$ and
$\mathrm{Tot}$ are $0$ and $fKm_j$. Dividing by $fKm_j\ge K$ gives
$b_0\le C_0\Lambda_{B_0}+l/2+1.6$. By (6c), $jA_0\ge-Nb_0-C'N\Lambda_{B_0}$, and $N\le2j$, so
$A_0\ge-2(C_0\Lambda_{B_0}+l/2+1.6+C'\Lambda_{B_0})$. If $N-j>j$, then $m_j=j$, both sums are
$-fKm_j$, and the same division gives $-(A_0+b_0)\le C_0\Lambda_{B_0}+l/2+1.6$. By (6c),
$(N-j)A_0=N(A_0+b_0)-(jA_0+Nb_0)\ge-N(C_0\Lambda_{B_0}+l/2+1.6+C'\Lambda_{B_0})$, and $N\le2(N-j)$, which
gives the same bound for $A_0$. Finally $N+1\le B_0$ and $128\le B_0$ give
$l\le\log B_0\le\Lambda_{B_0}$ and $1.6\le\log B_0\le\Lambda_{B_0}$, so
$A_0\ge-(2C_0+2C'+3)\Lambda_{B_0}$.
% Lean: GlobalTheory.T.test_height_bound (test_height_bound_slope_flat, RootTower.LowerBoundAtBase),
% with GlobalTheory.T.flag_height_test_bound, GlobalTheory.T.test_polynomial_shift_total and
% GlobalTheory.T.test_polynomial_remainder_bound.

(6e) The first bound of (4). Let $u$ be regular dominant with gaps at least $g$, and let $B\ge128$ be
at least the number of bad primes with $B+1>4|\Psi|N^2(1+\|u\|/g)$. Then $B\ge4|\Psi|N^2$, and (6a),
(6c) and (6d) apply at $B_0=B$: with $K_0=2(C_0+C')+3$, there are $A_0,b_0$ with
$A_0\ge-K_0\Lambda_B$ and $|jA_0+Nb_0|\le C'N\Lambda_B$. Put
$E=(d+1)\deg_{\mathrm{fl}}(\ell(u)^d)$, which is positive, and for $\sigma\in\Psi$ let
$s_\sigma=\sum_{i\le j}u_{\sigma,i}$ and $t_\sigma=\sum_iu_{\sigma,i}$. By
Lemma~\ref{cv-reconstruction-flag-degree}(2),
\[
 -A_0\sum_\sigma\mathrm{Sh}_\sigma(\ell(u)^{d+1})-b_0\sum_\sigma\mathrm{Tot}_\sigma(\ell(u)^{d+1})
 =E\sum_\sigma(A_0s_\sigma+b_0t_\sigma),
\]
and $A_0s_\sigma+b_0t_\sigma=A_0(s_\sigma-jt_\sigma/N)+(jA_0+Nb_0)t_\sigma/N$. As $u$ is dominant, the
mean of its first $j$ entries at $\sigma$ is at least the mean of all of them, so
$0\le s_\sigma-jt_\sigma/N\le\sum_i|u_{\sigma,i}|$; and $|t_\sigma|\le\sum_i|u_{\sigma,i}|$. Hence
$A_0s_\sigma+b_0t_\sigma\ge-(K_0+C')\Lambda_B\sum_i|u_{\sigma,i}|$, and the sum over $\sigma$ is at least
$-(K_0+C')\Lambda_B\|u\|_1$. Now apply (6b) to $F=\ell(u)^{d+1}$, with $h(\ell(u)^{d+1})\ge0$ by
Lemma~\ref{cv-reconstruction-flag-height}(2) and $|R_B(F)|\le2Em_j\operatorname{spread}(u)$ by
Lemma~\ref{cv-transport-remainders}(1):
\[
 \mathrm{PP}_B(\ell(u)^{d+1})\ge-\operatorname{vol}_TE\big((K_0+C')\Lambda_B\|u\|_1+2m_j\operatorname{spread}(u)\big)
 \ge-(K_0+C'+2)(d+1)\operatorname{vol}_T\deg_{\mathrm{fl}}(\ell(u)^d)\big(\Lambda_B\|u\|_1+m_j\operatorname{spread}(u)\big).
\]
Here $K_0+C'+2=2C_0+3C'+5$, and $m_j=\min(j,N-j)$. This is the lower bound for the arithmetic top
intersection number in (4).
% Lean: GlobalTheory.T.tower_T2 (tower_T2_of_inputs, RootTower.T2), with
% GlobalTheory.T.homogeneous_power_sum_decomposition, GlobalTheory.T.lin_root_remainder_bound and
% GlobalTheory.T.dominant_partial_sum_bounds.

(6f) The mean slope. By Step 1 all the $\operatorname{vol}(c)$ are equal, to $\operatorname{vol}_T/n$, where
$n$ is the number of components. So the mean over $c$ of the slopes
$\mathrm{PP}_{B,c}(\ell(u)^{d+1})/((d+1)\operatorname{vol}(c)\deg_{\mathrm{fl}}(\ell(u)^d))$ is
$\mathrm{PP}_B(\ell(u)^{d+1})/((d+1)\operatorname{vol}_T\deg_{\mathrm{fl}}(\ell(u)^d))$, which by (6e) is at
least $-(2C_0+3C'+5)(\Lambda_B\|u\|_1+m_j\operatorname{spread}(u))$. This is the second bound of (4),
with $C=2C_0+3C'+5$.
% Lean: GlobalTheory.T.mean_slope_proportional_volume (RootTower.MeanSlope). Steps 3 to 6 are
% combined in GlobalTheory.CompanionTheory.Ar.Top.chain_tower (GlobalTheory.T.tower_weil_inputs_combined).
\end{proof}

\subsection{Reference geometries and lattice families}\label{subsection-cv-reference-geometries}

\begin{lemma}\label{cv-origin-components}
Let a PEL datum of signature $(p,q)$ over a CM field $E$ with $p+q\ge1$, a level, its moduli scheme $X$
over a coefficient field $F_c$ and a model $(\mathcal Y,\mathcal A,j)$ of $X$ over the ring of integers
$\mathcal O_c$ of $F_c$ be given, defined as in Definition~\ref{cv-model} with $\mathcal O_c$ in place of
$O$. Then the points of $\mathcal Y$ that have no generization other than themselves form a finite
nonempty set, and each of them is the generic point of an irreducible component of $\mathcal Y$.
\end{lemma}
\leanref{Development/AutomorphicSpectrum/Proofs/GeometryProofs/GenericPointsFiniteNonemptyStatement.lean}{GlobalTheory.T.generic_points_finite_nonempty}
\leanref{Development/AutomorphicSpectrum/Proofs/GeometryProofs/GenericPointsFiniteNonempty.lean}{GlobalTheory.CompanionTheory.Ar.Comps.finite_nonempty_maximal}
\begin{proof}
The scheme $\mathcal Y$ is projective over the Noetherian ring $\mathcal O_c$, so its underlying space
is Noetherian and has finitely many irreducible components. Let $x$ be a point with no generization
other than itself, and let $Z$ be an irreducible component that contains $x$, with generic point
$\zeta$. Then $\zeta$ generizes $x$, so $x=\zeta$ and $\overline{\{x\}}=Z$. Distinct
points of a scheme have distinct closures, so $x\mapsto\overline{\{x\}}$ is injective, and the set is
finite. It is nonempty: the moduli scheme $X$ has complex points (Theorem~\ref{cv-uniformization-family}),
a point of $X$ gives a point of $X\otimes F_c$ and $j$ a point $x$ of $\mathcal Y$, and the generic point
$\zeta$ of an irreducible component $Z$ containing $x$ has no other generization: if $x'$ generizes
$\zeta$, then $Z\subset\overline{\{x'\}}$, which is irreducible, so $\overline{\{x'\}}=Z$ and $x'=\zeta$.
\end{proof}
% Lean: finite_maximal_points (the injection x -> closure {x} into irreducibleComponents, finite in a
% Noetherian space) and nonempty_maximal_points; Noetherian by ModelSum.noetherianSpace_of_isProjOver;
% nonempty by model_nonempty and moduli_nonempty (N1c).

\begin{definition}\label{cv-origin-reference}
In the setting of Definition~\ref{cv-origin-hodge-component}, put $D=|\Psi|N(N-1)/2$. A \emph{reference
geometry} of $\mathcal X$ consists of:
\begin{enumerate}
\item a field $\Omega$ of characteristic zero, finite sets $K$ and $I_k$ for $k\in K$, a point
  $y\in\Omega^{\sqcup_kI_k}$, and a Hodge component $Z$ of $\mathcal X$ at the empty set with coordinates
  $y$ (Definition~\ref{cv-origin-hodge-component}(1));
\item a flag of $Z$ of degree $1$ whose root lines have multidegrees $a_{\tau,i}\in\mathbf Z^K$ with
  generators $e_{\tau,i}\in\Omega^\times$ (Definition~\ref{cv-origin-hodge-component}(2) and (3));
\item an ideal $J_\partial$ of $R=\mathbf Z[x_{k,i}:k\in K,\ i\in I_k]$, the
  \emph{boundary ideal}, such that for every chart $\kappa$ of $y$ there are elements $f_1,\dots,f_m$ and
  $a_1,\dots,a_m$ of $A_\kappa=\mathbf Q[y_{k,i}/y_{k,\kappa_k}:k\in K,\ i\in I_k]$ with
  $\sum_la_lf_l=1$, and for each $l$ with $f_l\neq0$ the algebra $A_\kappa[1/f_l]$ is smooth over
  $\mathbf Q$ and has elements $z_1,\dots,z_D$, whose differentials form a basis of its module of
  differentials over $\mathbf Q$, and an integer $n\le D$, such that a prime ideal of $A_\kappa[1/f_l]$
  contains the values $g(y_{k,i}/y_{k,\kappa_k})$ of all $g\in J_\partial$ if and only if it contains one
  of $z_1,\dots,z_n$;
\item the complex points of $Z$ whose coordinates lie in the zero set of $J_\partial$ in $\hat C$
  (Definition~\ref{lit-multiprojective}) are those over which $\mathcal A_Z$ is not abelian; here
  $w\in\hat C$ are \emph{coordinates} of a complex point $P$ of $Z$ if, for each $k$ and each $c\in I_k$
  with $w_{k,c}\neq0$, $y_{k,c}\neq0$ and every ratio $y_{k,i}/y_{k,c}$ is regular at $P$ with value
  $w_{k,i}/w_{k,c}$;
\item an index $k_0\in K$ whose block \emph{embeds the generic fibre}: every point $P$ of $Z$ over no
  prime lies in a \emph{chart} $i\in I_{k_0}$ of the block, that is, $y_{k_0,i}\neq0$ and the ratios
  $y_{k_0,j}/y_{k_0,i}$ are regular at $P$; for every such $i$, every element $a$ of the local ring of
  $Z$ at $P$ has an element $s$ of $B_i=\mathbf Q[y_{k_0,j}/y_{k_0,i}:j\in I_{k_0}]$ that is a unit at $P$
  with $sa\in B_i$; and two points of $Z$ over no prime that lie in the chart $i$, at which the same
  elements of $B_i$ are units, are equal.
\end{enumerate}
\end{definition}
\leanref{Development/AutomorphicSpectrum/Definitions/GlobalGeometry/EmbedsGeneric.lean}{GlobalTheory.CompanionTheory.Ar.RefGeo}
\leanref{Development/AutomorphicSpectrum/Definitions/GlobalGeometry/EmbedsGeneric.lean}{GlobalTheory.CompanionTheory.Ar.EmbedsGeneric}
\leanref{Development/AutomorphicSpectrum/Definitions/GlobalGeometry/IsSncModel.lean}{GlobalTheory.CompanionTheory.Ar.Cmp.IsSncModel}
\leanref{Development/AutomorphicSpectrum/Definitions/GlobalGeometry/CoordsAt.lean}{GlobalTheory.CompanionTheory.Ar.CoordsAt}

\begin{lemma}\label{cv-origin-reference-char}
In the setting of Definition~\ref{cv-origin-hodge-component}, let $\Omega_0$ be a field and
$\xi_0:\operatorname{Spec}\Omega_0\to\mathcal X$ a morphism such that the local ring of $\mathcal X$ at the
image $x_0$ of $\xi_0$ maps bijectively onto $\Omega_0$, and the local ring at every specialization of
$x_0$ maps injectively into $\Omega_0$. Then $\Omega_0$ has characteristic zero.
\end{lemma}
\leanref{Development/AutomorphicSpectrum/Proofs/GeometryProofs/ComponentGenericCharZeroStatement.lean}{GlobalTheory.T.component_generic_char_zero}
\leanref{Development/AutomorphicSpectrum/Proofs/GeometryProofs/ComponentGenericCharZero.lean}{GlobalTheory.CompanionTheory.Ar.CharZeroModel.charZero_of_flat}
\begin{proof}
Let $u$ be the image of $x_0$ in $\operatorname{Spec}\mathcal O_c$. The ring $\mathcal O_c$ is a domain of
characteristic zero, so a nonzero integer $n$ is nonzero in the local ring $\mathcal O_{c,u}$, and it is a
nonzerodivisor there. The model $\mathcal X$ is flat over $\mathcal O_c$
(Definition~\ref{cv-model}), so $\mathcal O_{\mathcal X,x_0}$ is flat over $\mathcal O_{c,u}$ and $n$ is
a nonzerodivisor on $\mathcal O_{\mathcal X,x_0}$; in particular $n\neq0$ there. The map
$\mathcal O_{\mathcal X,x_0}\to\Omega_0$ is injective, so $n\neq0$ in $\Omega_0$.
\end{proof}

\begin{definition}\label{cv-origin-reference-lattices}
Let $T$ be a scheme, $\Omega$ a field and $\xi:\operatorname{Spec}\Omega\to T$ a morphism with image $\eta$. For a
point $P$ that specializes $\eta$, let $\mathcal O_P\subset\Omega$ be the ring of elements regular at $P$
(Definition~\ref{cv-origin-segre-coordinates}); it is the local ring of $T$ at $P$ when $\mathcal O_{T,P}\to\Omega$ is
injective.
\begin{enumerate}
\item Let $G$ be a commutative group scheme over $T$, $R$ a ring and $x:\operatorname{Spec}R\to T$ a morphism, and put
  $R[\varepsilon]=R[\varepsilon]/(\varepsilon^2)$. A \emph{tangent vector} of $G$ at $x$ is a morphism
  $v:\operatorname{Spec}R[\varepsilon]\to G$ over $T$, where $\operatorname{Spec}R[\varepsilon]$ lies over $T$ through $x$,
  whose restriction to $\operatorname{Spec}R$ is the unit section at $x$. The \emph{Lie algebra} of $G$ at $x$ is the set of
  tangent vectors, in which the sum of $v$ and $v'$ is their product $v\cdot v'$ for the group law and
  $av=v\circ\operatorname{Spec}(\varepsilon\mapsto a\varepsilon)$ for $a\in R$. Let $\mathrm{Lie}$ be the Lie algebra of $G$
  at $\xi$. A vector $e\in\mathrm{Lie}$ \emph{extends} to a point $P$ that specializes $\eta$ if
  $e=v\circ\operatorname{Spec}(\mathcal O_{T,P}[\varepsilon]\to\Omega[\varepsilon])$ for a tangent vector $v$ of $G$ at
  $\operatorname{Spec}\mathcal O_{T,P}\to T$. The \emph{Lie algebra of $G$ at $P$} is the $\mathcal O_P$-submodule of
  $\mathrm{Lie}$ generated by the vectors that extend to $P$.
\item Let $V$ be an $\Omega$-vector space. A \emph{lattice family} $L$ in $V$ assigns to each point $P$ that specializes
  $\eta$ an $\mathcal O_P$-submodule $L(P)$ of $V$; for a subspace $W$, $L\cap W$ is the lattice family
  $P\mapsto L(P)\cap W$. Let $S$ be a set of points that specialize $\eta$, $W\subset V$ a subspace and $r\ge0$. The family
  $L$ is a \emph{vector bundle of rank $r$ on $S$ with generic fibre $W$} if every $P\in S$ has an open neighbourhood $U$
  with a basis $b_1,\dots,b_r$ of $W$ such that $L(Q)=\sum_i\mathcal O_Qb_i$ for every $Q\in S\cap U$; it is a
  \emph{vector bundle of rank $r$ with generic fibre $W$} if $S$ is the set of all points that specialize $\eta$. For
  $b^0_1,\dots,b^0_r\in V$, $L$ is a \emph{flag bundle on $S$ adapted to $b^0$} if every $P\in S$ has an open
  neighbourhood $U$ with linearly independent $b_1,\dots,b_r$ such that
  $\langle b_1,\dots,b_i\rangle=\langle b^0_1,\dots,b^0_i\rangle$ for every $i$ and $L(Q)=\sum_i\mathcal O_Qb_i$ for every
  $Q\in S\cap U$.
\item Let $T'$ be a scheme with a morphism $\xi':\operatorname{Spec}\Omega'\to T'$ from the spectrum of a field, $f:T'\to T$ a
  morphism, $V'$ an $\Omega'$-vector space and $j:V\to V'$ an additive map. The \emph{pullback} $f^*_jL$ of a lattice
  family $L$ in $V$ is the lattice family in $V'$ for which $(f^*_jL)(P')$ is the $\mathcal O_{P'}$-submodule generated by
  $j(L(f(P')))$ if $f(P')$ specializes $\eta$, and is $0$ otherwise; $f^*L$ is $f^*_jL$ for $j$ the identity.
\item Let $\iota:\Omega\to\Omega'$ be a field embedding, $V'$ an $\Omega'$-vector space, $j:V\to V'$ an additive map and
  $d\ge1$, and for $s$ in a finite set let $W_s\subset V$ be a subspace and $b_s=(b_{s,1},\dots,b_{s,r_s})$ a family in
  $V'$. The families $b_s$ form a \emph{generic flag of degree $d$ along $j$ over the $W_s$} if $j(ae)=\iota(a)j(e)$ for
  $a\in\Omega$, each $b_s$ is linearly independent, and there are bases $c_s=(c_{s,1},\dots,c_{s,r_s})$ of the $W_s$, with
  each family $j(c_s)$ linearly independent, and elements $u_{s,lk}\in\Omega'$ for $l>k$, algebraically independent over
  $\iota(\Omega)$, such that $\Omega'$ has degree $d$ over the field $\iota(\Omega)(u)$ and
  $\langle b_{s,1},\dots,b_{s,i}\rangle$ is spanned by the vectors $j(c_{s,k})+\sum_{l>k}u_{s,lk}j(c_{s,l})$ with $k\le i$,
  for every $s$ and $i$.
\item Let $\iota:\Omega\to\Omega'$ be a field embedding, $T'$ a scheme with a morphism $\xi':\operatorname{Spec}\Omega'\to T'$
  and $f:T'\to T$ a morphism with $f\circ\xi'=\xi\circ\operatorname{Spec}\iota$, and let $P$ be a point that specializes
  $\eta$, and $\tilde A$ the integral closure of $\iota(\mathcal O_P)$ in $\Omega'$. Then $f$ is a \emph{normalization over}
  $P$ if, for every point $P'$ that specializes the image of $\xi'$ with $f(P')=P$, $\mathcal O_{P'}$ is the localization of
  $\tilde A$ at a maximal ideal, and for every maximal ideal $\mathfrak m$ of $\tilde A$ there is exactly one such point
  $P'$ with $\mathcal O_{P'}=\tilde A_{\mathfrak m}$.
\end{enumerate}
\end{definition}
\leanref{Development/AutomorphicSpectrum/Definitions/GlobalGeometry/LieModel.lean}{GlobalTheory.CompanionTheory.Ar.LieModel}
\leanref{Development/AutomorphicSpectrum/Definitions/GlobalGeometry/EigLat.lean}{GlobalTheory.CompanionTheory.Ar.LieModel.lat}
\leanref{Development/AutomorphicSpectrum/Definitions/GlobalGeometry/EtaleLocalizationAtPrimeTensorOfLowHeight.lean}{GlobalTheory.CompanionTheory.Ar.Lat}
\leanref{Development/AutomorphicSpectrum/Definitions/GlobalGeometry/EtaleLocalizationAtPrimeTensorOfLowHeight.lean}{GlobalTheory.CompanionTheory.Ar.Lat.IsVBOn}
\leanref{Development/AutomorphicSpectrum/Definitions/GlobalGeometry/IsGenericFlag.lean}{GlobalTheory.CompanionTheory.Ar.Lat.IsFlagVBOn}
\leanref{Development/AutomorphicSpectrum/Definitions/GlobalGeometry/EtaleLocalizationAtPrimeTensorOfLowHeight.lean}{GlobalTheory.CompanionTheory.Ar.Lat.pull}
\leanref{Development/AutomorphicSpectrum/Definitions/GlobalGeometry/IsGenericFlag.lean}{GlobalTheory.CompanionTheory.Ar.IsGenericFlag}
\leanref{Development/AutomorphicSpectrum/Definitions/GlobalGeometry/GlueApplyOfNotInFibre.lean}{GlobalTheory.CompanionTheory.Ar.NormAt}

\begin{lemma}\label{cv-origin-reference-pullback}
In the setting of Definition~\ref{cv-origin-hodge-component}, let $f:T'\to T$ and $g:T''\to T'$ be morphisms of
schemes and $G$ an $\mathcal O_E$-semi-abelian scheme over $T$ (Definition~\ref{cv-model-toroidal-semiabelian}(2)).
\begin{enumerate}
\item The fibre product $f^*G=G\times_TT'$, with the pullbacks of the action of $\mathcal O_E$ and of the line bundle with
  its cubical structure, is an $\mathcal O_E$-semi-abelian scheme over $T'$, and the projection $f^*G\to G$ is a morphism
  over $f$.
\item If $G'$ and $G''$ are $\mathcal O_E$-semi-abelian schemes over $T'$ and $T''$ with morphisms $G'\to G$ over $f$ and
  $G''\to G'$ over $g$, then the composite $G''\to G$ is a morphism over $f\circ g$.
\end{enumerate}
\end{lemma}
\leanref{ReciprocityTheory/ReductiveGroups/ArithmeticGeometry/BirationalGeometry/ToroidalCompactification/SabPull.lean}{ReciprocityTheory.ChartRestriction.G.Geo.Tor.sabPull}
\leanref{ReciprocityTheory/ReductiveGroups/ArithmeticGeometry/BirationalGeometry/ToroidalCompactification/SabComp.lean}{ReciprocityTheory.ChartRestriction.G.Geo.Tor.sabComp}
\begin{proof}
(1) The fibre product is a commutative group scheme over $T'$, with the group law pulled back. It is smooth and
separated over $T'$, and its fibre at a point $t'$ is the base change of the fibre of $G$ at $f(t')$ to the residue field
of $t'$, an extension of an abelian variety by a torus; so it is semi-abelian. The action of $\mathcal O_E$ by
homomorphisms pulls back to an action by homomorphisms. The line bundle pulls back with its trivialization along the unit
section. Its cubical structure is an isomorphism of line bundles on $G\times_TG\times_TG$ subject to identities of symmetry
and of compatibility with the group law, and its pullback to $f^*G\times_{T'}f^*G\times_{T'}f^*G$ is a cubical structure on
the pulled-back line bundle, symmetric if the line bundle is. The isomorphisms that make complex conjugation the Rosati
involution, and the isomorphism of the normalization, pull back in the same way. The projection identifies $f^*G$ with the
pullback of $G$, compatibly with these
structures, so it is a morphism over $f$.

(2) The composite of two cartesian squares is cartesian, and the compatibilities with the actions, the line bundles and
the cubical structures compose.
\end{proof}
% Lean: sabPull (IsSemiAbelian.pullback, OAct.pullback, Out.CubePullSt = Out.proved_CubePullSt for the cubical
% structure and its symmetry, LBIso.chartPull for adj and prim); sabComp (IsGrpPullback.comp, and
% PullCubCompSt = pullCubComp for the cubical structures).

\begin{lemma}\label{cv-origin-reference-lie}
Let $T$ be a scheme, $G$ and $G'$ commutative group schemes over $T$, $R$ a ring and $x:\operatorname{Spec}R\to T$ a
morphism.
\begin{enumerate}
\item The Lie algebra of $G$ at $x$ (Definition~\ref{cv-origin-reference-lattices}(1)) is an $R$-module.
\item For a homomorphism $\varphi:G\to G'$ of group schemes over $T$, $v\mapsto\varphi\circ v$ is an $R$-linear map from
  the Lie algebra of $G$ at $x$ to that of $G'$.
\item Let $O$ be a commutative ring and $\alpha(a)$, for $a\in O$, endomorphisms of the group scheme $G$ with
  $\alpha(1)$ the identity, $\alpha(ab)=\alpha(a)\circ\alpha(b)$ and $\alpha(a+b)=\alpha(a)\cdot\alpha(b)$ for the group
  law. Then $\rho(a)(v)=\alpha(a)\circ v$ defines a ring homomorphism $\rho$ from $O$ to the ring of endomorphisms of the
  $R$-module of (1).
\end{enumerate}
\end{lemma}
\leanref{Development/AutomorphicSpectrum/Proofs/GeometryProofs/ExistsLieModel.lean}{GlobalTheory.CompanionTheory.Ar.lieModel_exists}
\leanref{Development/AutomorphicSpectrum/Proofs/GeometryProofs/LieModelFunctorialMap.lean}{GlobalTheory.CompanionTheory.Ar.lieModel_exists_unique_diff}
\leanref{Development/AutomorphicSpectrum/Proofs/GeometryProofs/LieModelExistsUniqueAct.lean}{GlobalTheory.CompanionTheory.Ar.lieModel_exists_unique_act}
\begin{proof}
(1) This is the module structure of the Lie algebra of a group scheme \cite[Expos\'e~II]{DG70}.

(2) As $\varphi$ is a homomorphism, $\varphi\circ v$ restricts to the unit section,
$\varphi\circ(v\cdot v')=(\varphi\circ v)\cdot(\varphi\circ v')$, and $\varphi\circ v\circ\operatorname{Spec}(\varepsilon\mapsto
a\varepsilon)$ is $a(\varphi\circ v)$.

(3) Each $\rho(a)$ is $R$-linear by (2). The identities $\rho(1)=1$ and $\rho(ab)=\rho(a)\rho(b)$ come from those of
$\alpha$, and $\rho(a+b)=\rho(a)+\rho(b)$ holds because composition with $v$ carries the product $\alpha(a)\cdot\alpha(b)$
to $(\alpha(a)\circ v)\cdot(\alpha(b)\circ v)$.
\end{proof}
% Lean: lieModel_exists takes the tangent vectors themselves as the model (LieEx.tanModule, whose distributivity
% (a+b)v = av + bv is LieEx.hom_add); lieModel_exists_unique_diff (hadd from MonObj.mul_comp, hsmul from M.smul) and
% lieModel_exists_unique_act (hone, hmul', hadd'). The Lean's LieModel is a module E with a bijection pt onto the
% tangent vectors; its uniqueness clauses say that the linear map, and rho, are determined by pt.

\begin{lemma}[Demazure--Grothendieck \cite{DG70}]\label{cv-origin-reference-lie-bundle-base}
Let $T$ be a scheme, $G$ a commutative group scheme over $T$, smooth over $T$, $R$ a ring and $x:\operatorname{Spec}R\to T$ a
morphism, and let $L_x$ be the Lie algebra of $G$ at $x$ (Definition~\ref{cv-origin-reference-lattices}(1)). Then $L_x$ is a finite
projective $R$-module. For a ring map $\psi:R\to A$, let $\psi[\varepsilon]:R[\varepsilon]\to A[\varepsilon]$ be $\psi$ with
$\varepsilon\mapsto\varepsilon$; then $v\mapsto v\circ\operatorname{Spec}\psi[\varepsilon]$ is a $\psi$-semilinear map from $L_x$ to the
Lie algebra $L_{x\circ\operatorname{Spec}\psi}$ of $G$ at $x\circ\operatorname{Spec}\psi$, and the $A$-linear map
$A\otimes_RL_x\to L_{x\circ\operatorname{Spec}\psi}$ it induces is an isomorphism.
\end{lemma}
\leanref{Development/AutomorphicSpectrum/Proofs/GeometryProofs/LieModelFiniteBaseChangeStatement.lean}{GlobalTheory.T.lie_model_finite_base_change}
\leanref{Development/AutomorphicSpectrum/Proofs/GeometryProofs/LieModelExistsUniqueSemilinearStatement.lean}{GlobalTheory.T.lie_model_existsUnique_semilinear}
\begin{proof}
By \cite[Expos\'e~II, 4.11]{DG70}, the functor $\mathrm{Lie}(G/T)$, whose value at $x$ is $L_x$, is represented by the vector
bundle $\mathbf V(\omega^1_{G/T})$, where $\omega^1_{G/T}$ is the pullback of $\Omega^1_{G/T}$ along the unit section; it is locally
free of finite type, as $G$ is smooth over $T$, and the construction commutes with base change. So $L_x$ is the dual of the finite
projective $R$-module of sections of the pullback of $\omega^1_{G/T}$ along $x$, and its formation commutes with $R\to A$.
\end{proof}
% Lean: lie_model_finite_base_change = LieProj.lie_finite_projective_baseChange = HtPt.Htat.lie_proj (Lie.exists_cover,
% lieModel_finite_projective, Lie.isBaseChange_lie); the semilinear map is the unique j of lie_model_existsUnique_semilinear.
% The Lean proves the published result in its own way.

\begin{lemma}\label{cv-origin-reference-lie-pull}
In the setting of Definition~\ref{cv-origin-hodge-component}, let $T$ and $T'$ be schemes with morphisms
$\xi:\operatorname{Spec}\Omega\to T$ and $\xi':\operatorname{Spec}\Omega'\to T'$ from the spectra of fields, such that the
local rings at the images of $\xi$ and $\xi'$ map bijectively onto $\Omega$ and $\Omega'$ and those at their
specializations map injectively. Let $\iota:\Omega\to\Omega'$ be a field embedding and $f:T'\to T$ a morphism with
$f\circ\xi'=\xi\circ\operatorname{Spec}\iota$, let $G$ and $G'$ be $\mathcal O_E$-semi-abelian schemes over $T$ and $T'$
with a morphism $m:G'\to G$ over $f$, and let $\mathrm{Lie}$ and $\mathrm{Lie}'$ be their Lie algebras at $\xi$ and $\xi'$,
with the actions $\rho$ and $\rho'$ of $\mathcal O_E$ of Lemma~\ref{cv-origin-reference-lie}(3). Let
$\iota[\varepsilon]:\Omega[\varepsilon]\to\Omega'[\varepsilon]$ be $\iota$ with $\varepsilon\mapsto\varepsilon$.
\begin{enumerate}
\item There is an additive map $j:\mathrm{Lie}\to\mathrm{Lie}'$ with $m\circ j(e)=e\circ\operatorname{Spec}\iota[\varepsilon]$
  for every $e$, such that $j(ae)=\iota(a)j(e)$ for $a\in\Omega$, $j$ carries linearly independent families to linearly
  independent families, its image spans $\mathrm{Lie}'$, and the Lie algebras of $G'$ at the points that specialize the
  image of $\xi'$ form the pullback $f^*_j$ of the lattice family of the Lie algebras of $G$
  (Definition~\ref{cv-origin-reference-lattices}).
\item Every additive map $j:\mathrm{Lie}\to\mathrm{Lie}'$ with $m\circ j(e)=e\circ\operatorname{Spec}\iota[\varepsilon]$ for
  every $e$ satisfies $j(\rho(b)e)=\rho'(b)j(e)$ for $b\in\mathcal O_E$.
\end{enumerate}
\end{lemma}
\leanref{Development/AutomorphicSpectrum/Proofs/GeometryProofs/LieModelDominantHodgeStatement.lean}{GlobalTheory.T.lie_model_dominant_hodge}
\leanref{Development/AutomorphicSpectrum/Proofs/GeometryProofs/LieModelPullbackActionRingStatement.lean}{GlobalTheory.T.lie_model_pullback_action_ring}
\begin{proof}
(1) As $m$ identifies $G'$ with $f^*G$, for a ring $A$ and a morphism $z:\operatorname{Spec}A\to T'$, composition with $m$ is a
bijection from the tangent vectors of $G'$ at $z$ onto those of $G$ at $f\circ z$, compatible with sums and with the action of
$A$; in particular a tangent vector of $G'$ is determined by its composite with $m$. The group scheme $G$ is smooth over $T$,
being semi-abelian. For a point $Q$ that specializes the image of $\xi$, or of $\xi'$, let $g_Q$ be the injective map from the
local ring at $Q$ to $\Omega$, or to $\Omega'$, through the local ring at that image.

\emph{The map $j$.} As $f\circ\xi'=\xi\circ\operatorname{Spec}\iota$, Lemma~\ref{cv-origin-reference-lie-bundle-base}, applied to $G$,
$\xi$ and $\iota$, shows that $e\mapsto e\circ\operatorname{Spec}\iota[\varepsilon]$ is an $\iota$-semilinear map from $\mathrm{Lie}$ to
the Lie algebra of $G$ at $f\circ\xi'$ that carries a basis to a basis. Let $j(e)$ be the tangent vector of $G'$ at $\xi'$ with
$m\circ j(e)=e\circ\operatorname{Spec}\iota[\varepsilon]$. Then $j$ is additive, $j(ae)=\iota(a)j(e)$, and $j$ carries a basis of
$\mathrm{Lie}$ to a basis of $\mathrm{Lie}'$. So its image spans $\mathrm{Lie}'$, and, as every linearly independent family extends to
a basis, $j$ carries linearly independent families to linearly independent families.

\emph{Lifting extensions.} Let $P'$ be a point that specializes the image of $\xi'$, and $P=f(P')$, which specializes the image of
$\xi$. Let $\lambda:\mathcal O_{T,P}\to\mathcal O_{T',P'}$ be the local homomorphism of $f$; then $g_{P'}\circ\lambda=\iota\circ g_P$, as
$f\circ\xi'=\xi\circ\operatorname{Spec}\iota$. Let $e\in\mathrm{Lie}$ extend to $P$ through a tangent vector $v$ of $G$ at
$\operatorname{Spec}\mathcal O_{T,P}\to T$, so that $e=v\circ\operatorname{Spec}g_P[\varepsilon]$. Then $v\circ\operatorname{Spec}\lambda[\varepsilon]$
is a tangent vector of $G$ at $f\circ(\operatorname{Spec}\mathcal O_{T',P'}\to T')$, so it is $m\circ v'$ for a tangent vector $v'$ of
$G'$ at $\operatorname{Spec}\mathcal O_{T',P'}\to T'$. Both $v'\circ\operatorname{Spec}g_{P'}[\varepsilon]$ and $j(e)$ have the composite
$v\circ\operatorname{Spec}(g_{P'}\circ\lambda)[\varepsilon]=e\circ\operatorname{Spec}\iota[\varepsilon]$ with $m$, so they are equal: $j(e)$
extends to $P'$ through $v'$.

\emph{The lattices.} The Lie algebra $\Lambda(P)$ of $G$ at $P$ is generated by the vectors that extend to $P$, whose images under $j$
extend to $P'$; as $j$ is semilinear and $\iota(\mathcal O_P)\subset\mathcal O_{P'}$, the module $(f^*_j\Lambda)(P')$ generated by
$j(\Lambda(P))$ lies in the Lie algebra $\Lambda'(P')$ of $G'$ at $P'$. Conversely, the vectors of $\mathrm{Lie}$ that extend to $P$
form an $\mathcal O_{T,P}$-module through $g_P$, which is identified with the Lie algebra of $G$ at
$\operatorname{Spec}\mathcal O_{T,P}\to T$, as $g_P$ is injective and a tangent vector at $\operatorname{Spec}\mathcal O_{T,P}$ is therefore
determined by its composite with $\operatorname{Spec}g_P[\varepsilon]$. By Lemma~\ref{cv-origin-reference-lie-bundle-base} it is finite
projective over the local ring $\mathcal O_{T,P}$, hence free \cite[Tag 00NX]{Sta}, with a basis $b_1,\dots,b_n$ extending through
tangent vectors $v_1,\dots,v_n$. By the previous paragraph, each $j(b_k)$ extends to $P'$ through a tangent vector $u_k$ of $G'$ with
$m\circ u_k=v_k\circ\operatorname{Spec}\lambda[\varepsilon]$. By Lemma~\ref{cv-origin-reference-lie-bundle-base}, applied to $\lambda$, and the
bijection given by $m$, the $u_k$ form a basis over $\mathcal O_{T',P'}$ of the Lie algebra of $G'$ at
$\operatorname{Spec}\mathcal O_{T',P'}\to T'$. Let $e'\in\mathrm{Lie}'$ extend to $P'$ through $w=\sum_kc_ku_k$, with
$c_k\in\mathcal O_{T',P'}$. Then $e'=w\circ\operatorname{Spec}g_{P'}[\varepsilon]=\sum_kg_{P'}(c_k)j(b_k)$, which lies in
$(f^*_j\Lambda)(P')$. As $\Lambda'(P')$ is generated by such $e'$, $\Lambda'(P')=(f^*_j\Lambda)(P')$.
% Lean: T.lie_model_dominant_hodge = LiePullSab.lie_pull_lat: exists_bc (T.lie_model_base_change_map, with the model of
% G' pushed to G by LieMove.pushModel m.cart) gives j, and T.lie_lattice_dominant_compat = LiePull.main the rest:
% LiePull.bc (T.lie_model_base_change_map with Smooth G.hom, t.semi.smooth) and semilin for j; indep_of_bases (extend to a
% basis); RegAtDom.genMap_stalkMap for g_P' o rho = iota o g_P; lieLat_carrier (the lattice is the set of extendable vectors);
% ext_up (LieMove.liftEps through the cartesian square m.cart) for the first inclusion; core for the second (extModel, the
% extendable vectors as a Lie model at Spec O_(T,P), using IsCptGen.inj; T.lie_model_free_over_local_ring; bc for the u_k).

(2) Let $e\in\mathrm{Lie}$ and $b\in\mathcal O_E$, and let $\alpha$ and $\alpha'$ be the actions of $\mathcal O_E$ on $G$ and
$G'$. The tangent vectors $j(\rho(b)e)$ and $\rho'(b)j(e)=\alpha'(b)\circ j(e)$ of $G'$ at $\xi'$ both lie over $\xi'$. Their
composites with $m$ are $\alpha(b)\circ e\circ\operatorname{Spec}\iota[\varepsilon]$ and
$m\circ\alpha'(b)\circ j(e)=\alpha(b)\circ m\circ j(e)$, which are equal, since $m$ commutes with the actions. As $G'$ is the
fibre product of $G$ and $T'$ over $T$, the two tangent vectors are equal.
\end{proof}
% Lean: (2) is lie_model_pullback_action_ring = FlagStep.lieMap_act_pull (hcanc from m.cart, m.act).

\begin{lemma}\label{cv-origin-reference-lie-same}
In the setting of Lemma~\ref{cv-origin-reference-lie-pull}, suppose that $\Omega'=\Omega$ and that $\iota$ is the
identity. Then $v'\mapsto m\circ v'$ is an isomorphism of $\Omega$-vector spaces from $\mathrm{Lie}'$ onto $\mathrm{Lie}$
that commutes with the actions of $\mathcal O_E$. With $\mathrm{Lie}'$ identified with $\mathrm{Lie}$ through it, the Lie
algebras of $G'$ at the points that specialize the image of $\xi'$ form the pullback $f^*$ of the lattice family of the
Lie algebras of $G$.
\end{lemma}
\leanref{Development/AutomorphicSpectrum/Proofs/GeometryProofs/GenericCoordsDescendSubfield.lean}{GlobalTheory.CompanionTheory.Ar.SatModel.exists_lieModel_pull}
\begin{proof}
For $v'\in\mathrm{Lie}'$, $m\circ v'$ is a tangent vector of $G$ at $f\circ\xi'=\xi$. Let $j$ be as in
Lemma~\ref{cv-origin-reference-lie-pull}(1), so that $m\circ j(e)=e$ for every $e$. It is $\Omega$-linear. A nonzero vector
is a linearly independent family, so $j$ is injective; its image is a subspace that spans $\mathrm{Lie}'$, so $j$ is
surjective. Hence $v'\mapsto m\circ v'$ is the inverse of $j$. By Lemma~\ref{cv-origin-reference-lie-pull}(2), $j$ commutes
with the actions. The Lie algebras of $G'$ at the points form the pullback $f^*_j$ of those of $G$, which is $f^*$ once
$\mathrm{Lie}'$ is identified with $\mathrm{Lie}$ through $j$.
\end{proof}
% Lean: exists_lieModel_pull applies lie_model_dominant_hodge to a fresh model M1, transports it to the model M2 on Li
% along j (hinj, hsurj), applies lie_model_dominant_hodge again to M and M2, and shows that the second map is the
% identity by hcanc (a tangent vector of t' is determined by its composite with m); the action is checked the same way.

\begin{lemma}\label{cv-origin-reference-lie-bundle}
Let $T$ be a scheme, $\Omega$ a field and $\xi:\operatorname{Spec}\Omega\to T$ a morphism such that the local ring at the
image of $\xi$ maps bijectively onto $\Omega$ and those at its specializations injectively, and let $G$ be a commutative
group scheme over $T$, smooth over $T$, with Lie algebra $\mathrm{Lie}$ at $\xi$. Then the lattice family $\Lambda$ of the Lie
algebras of $G$ at the points that specialize the image of $\xi$ is a vector bundle of rank $\dim_\Omega\mathrm{Lie}$ with
generic fibre $\mathrm{Lie}$ (Definition~\ref{cv-origin-reference-lattices}), and the lattice family
$P\mapsto\{\varphi\in\mathrm{Lie}^\vee:\varphi(\Lambda(P))\subset\mathcal O_P\}$ in the dual space is a vector bundle of the
same rank with generic fibre $\mathrm{Lie}^\vee$.
\end{lemma}
\leanref{Development/AutomorphicSpectrum/Proofs/GeometryProofs/LieOmegaLatticeVectorBundleStatement.lean}{GlobalTheory.T.lie_omega_lattice_vector_bundle}
\begin{proof}
Let $\eta$ be the image of $\xi$, and $P$ a point that specializes $\eta$. For a point $Q$ that specializes $\eta$, let
$g_Q:\mathcal O_{T,Q}\to\Omega$ be the map through $\mathcal O_{T,\eta}$, so that $\mathcal O_Q$ is its image. Composition with
$\operatorname{Spec}\psi[\varepsilon]$, for a ring map $\psi$, is $\psi$-semilinear on tangent vectors (Lemma~\ref{cv-origin-reference-lie-bundle-base}).

\emph{A chart.} Let $U$ be an affine open neighbourhood of $P$ with ring $R_0$, and $x_0:\operatorname{Spec}R_0\to T$ its inclusion. By
Lemma~\ref{cv-origin-reference-lie-bundle-base}, the Lie algebra of $G$ at $x_0$ is a finite projective $R_0$-module, so there is
$f\in R_0$ with $P$ in the standard open $D(f)$ such that its localization at $f$ is free \cite[Tag 00NX]{Sta}. Let $R=(R_0)_f$ and
$x:\operatorname{Spec}R\to T$ the open immersion with image $D(f)$. By the same lemma, the Lie algebra of $G$ at $x$ is the base
change of that at $x_0$, so it is free, with a basis $b_1,\dots,b_n$.

\emph{The generic point.} The open set $D(f)$ contains $P$, so it contains its generization $\eta$, and $\xi$ factors through it:
$\xi=x\circ\operatorname{Spec}\varphi$ for a ring map $\varphi:R\to\Omega$ \cite[Tag 01I1]{Sta}. By
Lemma~\ref{cv-origin-reference-lie-bundle-base}, applied to $\varphi$, the vectors $c_k=b_k\circ\operatorname{Spec}\varphi[\varepsilon]$
form a basis of $\mathrm{Lie}$; so $n=\dim_\Omega\mathrm{Lie}$.

\emph{The other points.} Let $Q\in D(f)$ specialize $\eta$. Every point of $\operatorname{Spec}\mathcal O_{T,Q}$ maps to a
generization of $Q$, which lies in the open set $D(f)$; so the canonical morphism $\operatorname{Spec}\mathcal O_{T,Q}\to T$ is
$x\circ\operatorname{Spec}\psi$ for a ring map $\psi:R\to\mathcal O_{T,Q}$. Then $x\circ\operatorname{Spec}(g_Q\circ\psi)=\xi=
x\circ\operatorname{Spec}\varphi$, and $x$ is a monomorphism, so $g_Q\circ\psi=\varphi$. By
Lemma~\ref{cv-origin-reference-lie-bundle-base}, applied to $\psi$, the vectors $\beta_k=b_k\circ\operatorname{Spec}\psi[\varepsilon]$
form a basis over $\mathcal O_{T,Q}$ of the Lie algebra of $G$ at $\operatorname{Spec}\mathcal O_{T,Q}\to T$, and
$\beta_k\circ\operatorname{Spec}g_Q[\varepsilon]=c_k$. A vector $e\in\mathrm{Lie}$ extends to $Q$ if and only if
$e=v\circ\operatorname{Spec}g_Q[\varepsilon]$ for a tangent vector $v=\sum_ka_k\beta_k$ with $a_k\in\mathcal O_{T,Q}$, that is, if and
only if $e=\sum_kg_Q(a_k)c_k$. So $\Lambda(Q)=\sum_k\mathcal O_Qc_k$ for every $Q\in D(f)$ that specializes $\eta$, and $\Lambda$ is a
vector bundle of rank $n$ with generic fibre $\mathrm{Lie}$.

\emph{The duals.} Let $(c^\vee_k)$ be the dual basis of $\mathrm{Lie}^\vee$. For $Q$ as above, a linear form $\varphi$ maps
$\Lambda(Q)=\sum_k\mathcal O_Qc_k$ into $\mathcal O_Q$ if and only if every $\varphi(c_k)$ lies in $\mathcal O_Q$, that is, if and only if
$\varphi=\sum_k\varphi(c_k)c^\vee_k$ has coefficients in $\mathcal O_Q$. So the dual lattice family is generated by $(c^\vee_k)$ at these
points, and it is a vector bundle of rank $n$ with generic fibre $\mathrm{Lie}^\vee$.
\end{proof}
% Lean: T.lie_omega_lattice_vector_bundle = lieLat_isVB_of_smooth: LieVB.exists_chart (an affine open from isBasis_affineOpens,
% T.lie_model_finite_base_change for the finite projective Lie algebra, Module.FinitePresentation.exists_free_localizedModule_powers
% for f, and the base change to Localization.Away f), LieVB.exists_lift (IsOpenImmersion.lift for the factorizations through x),
% LieVB.frame_of_basis (basis_tr for c and for beta, hcomp for beta_k composed with g_Q, mem_lat_iff and mem_span_range for the
% frame), and cancel_mono x for g_Q o psi = phi. The duals are Lat.isVBOn_dual_and_bidual. The hypothesis on the local rings
% (IsCptGen xi) is not used.

\begin{lemma}\label{cv-origin-reference-integral}
Let $\Omega$ be a field of characteristic zero, $T$ a scheme with a morphism $\xi:\operatorname{Spec}\Omega\to T$, covered by
the affine open subschemes of the charts of a point $y\in\Omega^{\sqcup_kI_k}$ as in
Definition~\ref{cv-origin-hodge-component}(1), and let $J$ be an ideal of $\mathbf Z[x_{k,i}]$ that satisfies condition (3)
of Definition~\ref{cv-origin-reference} for $y$ with an integer $e$ in place of $D$. Then the local ring of $T$ at every
point that specializes the image of $\xi$ and lies over no prime is integrally closed.
\end{lemma}
\leanref{Development/AutomorphicSpectrum/Proofs/GeometryProofs/ExistsCoordsForComponent.lean}{GlobalTheory.CompanionTheory.Ar.EigBase.isIntegrallyClosed_of_snc}
\begin{proof}
Let $P$ be such a point, in the chart $\kappa$ of $y$ (Definition~\ref{cv-origin-segre-coordinates}(2)). Each ratio
$y_{k,i}/y_{k,\kappa_k}$ is regular at $P$, being the Segre ratio for the chart $\kappa$ with $\kappa_k$ replaced by $i$. A
nonzero integer is $\pm1$ times a product of primes, which are units at $P$; so every rational number is regular at $P$,
and $A=\mathbf Q[y_{k,i}/y_{k,\kappa_k}]$ lies in $\mathcal O_{T,P}$. Condition (3) at $\kappa$ gives $f_l,a_l\in A$ with
$\sum_la_lf_l=1$. As $\mathcal O_{T,P}$ is local, some $f_l$ is a unit there, and $S=A[1/f_l]\subset\mathcal O_{T,P}$ is
smooth over $\mathbf Q$, hence normal \cite[Tag 033C]{Sta}. Let $\mathfrak p$ be the prime of the elements of $S$ that are
not units at $P$. By Lemma~\ref{cv-origin-snc-charts}(1), every element of $\mathcal O_{T,P}$ is $a/u$ with $a$ and $u$ in
$\mathbf Z[y_{k,i}/y_{k,\kappa_k}]\subset S$ and $u$ a unit at $P$. So $\mathcal O_{T,P}=S_{\mathfrak p}$, the localization of a
normal ring at a prime, an integrally closed domain.
\end{proof}
% Lean: isIntegrallyClosed_of_snc: htie.cover, the ratios r p (hr), hQ through SatModel.unitAt_int, the unit f l
% (maximal ideal argument), GlobalTheory.isIntegrallyClosed_of_smooth, the localization IsLocalization.AtPrime via
% regAt_iff_chart_localization (Lemma cv-origin-snc-charts(1)), and isDomain_and_isIntegrallyClosed_atPrime_of_smooth.

\begin{lemma}\label{cv-origin-reference-eigen-split}
Let $T$, $\Omega$ and $\xi$ be as in Lemma~\ref{cv-origin-reference-lie-bundle}, with $\Omega$ of characteristic zero. Let
$F$ be a number field, $V$ an $\Omega$-vector space with a ring homomorphism $\rho:\mathcal O_F\to\operatorname{End}_\Omega V$,
and $\sigma_1,\dots,\sigma_n:\mathcal O_F\to\Omega$ distinct ring homomorphisms, where $n=[F:\mathbf Q]$; let $V_i$ be the
subspace of $V$ on which each $\rho(a)$ acts by $\sigma_i(a)$. Let $S$ be a set of points that specialize the image of $\xi$,
and $L$ a lattice family in $V$ that is a vector bundle of rank $g$ on $S$ with generic fibre $V$
(Definition~\ref{cv-origin-reference-lattices}). Suppose that at every $P\in S$, $L(P)$ is stable under $\rho(\mathcal O_F)$,
the values of the $\sigma_i$ are regular, and the discriminant of $F$ is a unit. Then each $L\cap V_i$ is a vector bundle of
rank $\dim V_i$ on $S$ with generic fibre $V_i$, and $L(P)=\sum_i(L(P)\cap V_i)$ for every $P\in S$.
\end{lemma}
\leanref{Development/AutomorphicSpectrum/Proofs/GeometryProofs/LatticeEigenVectorBundleStatement.lean}{GlobalTheory.T.lattice_eigen_vector_bundle}
\begin{proof}
Each $\sigma_i$ extends to an embedding $\tau_i:F\to\Omega$, and the $\tau_i$ are distinct. Composed with an embedding of
$\Omega$ in an algebraic closure, distinct embeddings of $F$ in $\Omega$ stay distinct, so there are at most $[F:\mathbf Q]$ of
them, and the $\tau_i$ are all of them. Let $(a_k)$ be a $\mathbf Z$-basis of $\mathcal O_F$ and $(\delta_k)$ the dual basis of
$F$ for the trace form, and put $\pi_i=\sum_k\tau_i(\delta_k)\rho(a_k)$. A computation with the dual basis, from the expansion
$z=\sum_l\operatorname{Tr}_{F/\mathbf Q}(z\delta_l)a_l$ of $z\in\mathcal O_F$, the identity $\sum_k\delta_ka_k=1$ and
$\sum_i\tau_i(\delta_k)=\operatorname{Tr}_{F/\mathbf Q}(\delta_k)$, shows that $\rho(a)\pi_i=\tau_i(a)\pi_i$ for $a\in\mathcal O_F$,
that $\pi_i$ is the identity on $V_i$, and that $\sum_i\pi_i=1$. So $\pi_i$ is an idempotent with image $V_i$.

The discriminant $d$ of $F$ lies in the different of $F$, which multiplies the $\mathbf Z$-span of the $\delta_k$ into
$\mathcal O_F$. So $d\delta_k\in\mathcal O_F$, and $\tau_i(\delta_k)=d^{-1}\sigma_i(d\delta_k)$ is regular at every $P\in S$.
Hence $\pi_i$ preserves $L(P)$, and $L(P)=\sum_i\pi_i(L(P))\subset\sum_i(L(P)\cap V_i)$; the other inclusion is clear.

Let $P\in S$, with $U$ and $b$ as in the definition of a vector bundle. Over the local ring $\mathcal O_P$, $L(P)$ is the
direct sum of $\pi_i(L(P))=L(P)\cap V_i$ and $(1-\pi_i)(L(P))$, which are free, as direct summands of a free module over a
local ring \cite[Theorem~2.5]{Mat86}. Let $f$ and $f'$ be bases of the two summands; $f$ is linearly independent and spans
$\pi_i(V)=V_i$. The families $b$ and $(f,f')$ are combinations of each other with coefficients in $\mathcal O_P$. These
finitely many coefficients are regular on an open neighbourhood $U'\subset U$ of $P$, by
Lemma~\ref{cv-origin-segre-regular}(1), so $(f,f')$ generates $L(Q)$ for $Q\in S\cap U'$, and $f$ generates
$\pi_i(L(Q))=L(Q)\cap V_i$.
\end{proof}
% Lean: EigSplit.isVBOn_eigSp: tau s = IsFractionRing.lift, hbij (counting with AlgHom.card), HT.eigProj (pi_i),
% HT.sum_eigProj_eq_one, HT.ProjEig.eigProj_spec, hcoef from HT.exists_different_mul_dualInt and discr_mem_differentIdeal,
% hstable, and isVBOn_inter_range (exists_frame_of_proj over the local ring, exists_open_mem_span for the neighbourhood).

\begin{lemma}\label{cv-origin-reference-eigen}
In the setting of Definition~\ref{cv-origin-hodge-component}, let $T$, $\Omega$ and $\xi$ be as in
Lemma~\ref{cv-origin-reference-lie-bundle}, with $\Omega$ of characteristic zero, let $F_c\to\Omega$ be a field embedding,
and let $G$ be an $\mathcal O_E$-semi-abelian scheme over $T$ with Lie algebra $\mathrm{Lie}$ at $\xi$ and its action of
$\mathcal O_E$. For a label $s$ let $\mathrm{Lie}_s$ be the subspace on which each $b\in\mathcal O_E$ acts by $s(b)$, through
$F_c\to\Omega$, and let $\Lambda$ be the lattice family of the Lie algebras of $G$ at the points that specialize the image of
$\xi$. Suppose that the local ring of $T$ at every such point over no prime is integrally closed. Then:
\begin{enumerate}
\item the subspaces $\mathrm{Lie}_s$ are independent;
\item $\Lambda(P)=\sum_s(\Lambda(P)\cap\mathrm{Lie}_s)$ at every such point $P$ over no prime;
\item each $\Lambda\cap\mathrm{Lie}_s$ is a vector bundle of rank $\dim\mathrm{Lie}_s$ with generic fibre $\mathrm{Lie}_s$ on the
  set of these points over no prime.
\end{enumerate}
\end{lemma}
\leanref{Development/AutomorphicSpectrum/Proofs/GeometryProofs/GenericCoordsDescendSubfield.lean}{GlobalTheory.CompanionTheory.Ar.SatModel.lat_eq_iSup_of_char0}
\leanref{Development/AutomorphicSpectrum/Proofs/GeometryProofs/ExistsCoordsForComponent.lean}{GlobalTheory.CompanionTheory.Ar.EigBase.inter_isVBOn_char0}
\begin{proof}
The $2|\Psi|=[E:\mathbf Q]$ labels are distinct embeddings of $\mathcal O_E$ into $F_c$, so their composites $\sigma_s$ with
$F_c\to\Omega$ are $[E:\mathbf Q]$ distinct ring homomorphisms $\mathcal O_E\to\Omega$. Eigenspaces for distinct characters
are independent, which is (1). For (2) and (3), apply Lemma~\ref{cv-origin-reference-eigen-split} with $F=E$, $L=\Lambda$
and $S$ the set of points over no prime; $\Lambda$ is a vector bundle of rank $\dim\mathrm{Lie}$ with generic fibre
$\mathrm{Lie}$ by Lemma~\ref{cv-origin-reference-lie-bundle}, as $G$ is smooth over $T$. Let $P\in S$. If $e$ extends to $P$
through a tangent vector $v$, then $\rho(b)e$ extends to $P$ through $\alpha(b)\circ v$, for the action $\alpha$ of
$\mathcal O_E$ on $G$; so $\Lambda(P)$ is stable under $\mathcal O_E$. The values $\sigma_s(b)$ are integral over $\mathbf Z$,
hence over $\mathcal O_P$, whose field of fractions is $\Omega$; as $\mathcal O_P$ is integrally closed, they lie in it. The
discriminant of $E$ is a nonzero integer, $\pm1$ times a product of primes, which are units at $P$.
\end{proof}
% Lean: lat_eq_iSup_of_char0 and inter_isVBOn_char0: T.label_embedding_injective (labels_all_once) for hinj and hcard,
% T.eigenspace_decomposition_internal for (1), T.lie_omega_lattice_vector_bundle, T.lattice_eigen_vector_bundle with
% SatModel.lat_stable, integrality through regAt_isLocalRing_isFractionRing_stalk, and SatModel.unitAt_int.

\begin{lemma}\label{cv-origin-reference-normalization}
Let $T$, $\Omega$ and $\xi$ be as in Lemma~\ref{cv-origin-reference-lie-bundle}, let $T'$ be a scheme with a morphism
$\xi':\operatorname{Spec}\Omega\to T'$, and $f:T'\to T$ a morphism with $f\circ\xi'=\xi$. Let $P$ be a point that specializes
the image of $\xi$, with $\mathcal O_{T,P}$ integrally closed, such that $f$ is a normalization over $P$
(Definition~\ref{cv-origin-reference-lattices}(5), with $\iota$ the identity). Then every point $P'$ that specializes the
image of $\xi'$ with $f(P')=P$ has $\mathcal O_{P'}=\mathcal O_{T,P}$ as subrings of $\Omega$, and there is at most one such
point.
\end{lemma}
\leanref{Development/AutomorphicSpectrum/Proofs/GeometryProofs/GenericCoordsDescendSubfield.lean}{GlobalTheory.CompanionTheory.Ar.SatModel.normAt_id_elim}
\begin{proof}
The field of fractions of $\mathcal O_{T,P}$ is the local ring at the image of $\xi$, which is $\Omega$. As $\mathcal O_{T,P}$
is integrally closed, its integral closure $\tilde A$ in $\Omega$ is $\mathcal O_{T,P}$ itself, a local ring with one maximal
ideal $\mathfrak m$, and $\tilde A_{\mathfrak m}=\mathcal O_{T,P}$. By the definition of a normalization over $P$, every such
$P'$ has $\mathcal O_{P'}=\tilde A_{\mathfrak m}$, and there is only one point with this property.
\end{proof}
% Lean: normAt_id_elim: regAt_isLocalRing_isFractionRing_stalk, hA (the integral closure is regAt p), NormAt.exists_max
% and NormAt.existsUnique with IsLocalRing.eq_maximalIdeal.

\begin{lemma}\label{cv-origin-reference-refinement-plucker}
Let $\Omega$ be a field, $W$ an $\Omega$-vector space of finite dimension $r$ with a basis $\beta$, $z_1,\dots,z_n$ vectors of
$W$ and $A\subset\Omega$ a subring. For a map $c:\{1,\dots,r\}\to\{1,\dots,n\}$ let $\Delta(c)$ be the determinant of
$(z_{c(1)},\dots,z_{c(r)})$ in the basis $\beta$.
\begin{enumerate}
\item If $\Delta(c_0)\neq0$ and $\Delta(c)/\Delta(c_0)\in A$ for every $c$, then $z_{c_0(1)},\dots,z_{c_0(r)}$ is a basis of
  $W$ and $\sum_kAz_k=\sum_iAz_{c_0(i)}$.
\item If $A$ is local and $\sum_kAz_k=\sum_iAf_i$ for a linearly independent family $f_1,\dots,f_r$ of $W$, then there is
  $c_0$ with $\Delta(c_0)\neq0$ and $\Delta(c)/\Delta(c_0)\in A$ for every $c$.
\end{enumerate}
\end{lemma}
\leanref{Development/WeilBound/CoreStatements/HodgeTatePoints/McCramer.lean}{PositivityTheory.HtPt.Mc.cramer_sub}
\leanref{Development/WeilBound/CoreStatements/HodgeTatePoints/McCramer.lean}{PositivityTheory.HtPt.Mc.frame_plucker}
\begin{proof}
(1) As $\Delta(c_0)\neq0$, the $z_{c_0(i)}$ form a basis of $W$. By Cramer's rule, the $i$-th coordinate of $z_k$ in this basis is
$\Delta(c)/\Delta(c_0)$, where $c$ is $c_0$ with the value at $i$ replaced by $k$; it lies in $A$. So every $z_k$ lies in
$\sum_iAz_{c_0(i)}$, and the other inclusion is clear.

(2) As $r=\dim W$, the $f_i$ form a basis of $W$. Let $\delta(c)$ be the determinant of $(z_{c(1)},\dots,z_{c(r)})$ in this basis,
so that $\Delta(c)=\Delta_f\delta(c)$, where $\Delta_f\neq0$ is the determinant of $(f_1,\dots,f_r)$ in the basis $\beta$. Every
$z_k$ lies in $\sum_iAf_i$, so its coordinates in the basis $f$ lie in $A$, and $\delta(c)\in A$. Every $f_i$ lies in $\sum_kAz_k$;
expanding the determinant $1$ of $(f_1,\dots,f_r)$ in the basis $f$ by multilinearity gives $1\in\sum_cA\delta(c)$. As $1$ is not
in the maximal ideal of $A$, some $\delta(c_0)$ is a unit of $A$. Then $\Delta(c_0)\neq0$, and
$\Delta(c)/\Delta(c_0)=\delta(c)\delta(c_0)^{-1}\in A$ for every $c$.
\end{proof}
% Lean: cramer_sub (Basis.det_smul_mk_coord_eq_det_update for the coordinates as ratios of determinants) and frame_plucker
% (det_mem for delta(c) in A, det_mem_span for 1 in the span of the delta(c), and the maximal ideal of the local ring A).

\begin{lemma}\label{cv-origin-reference-refinement-twist}
Let $X$ be a scheme, $\Omega$ a field and $\xi:\operatorname{Spec}\Omega\to X$ a morphism of whose image every point of $X$ is
a specialization, and let $x=(x_j)_{j\in J}$ be a finite family in $\Omega$ such that $X$ has the coordinates $x$
(Definition~\ref{cv-origin-segre-coordinates}). Let $V$ be an $\Omega$-vector space and $L$ a lattice family in $V$
(Definition~\ref{cv-origin-reference-lattices}) such that every point of $X$ has an open neighbourhood $U$ and finitely many
vectors that generate $L(Q)$ over $\mathcal O_Q$ for every $Q\in U$. Then there are an integer $d\ge0$ and vectors
$z_1,\dots,z_n$ of $V$ such that $L(P)=\sum_k\mathcal O_Px_c^{-d}z_k$ for every index $c$ and every point $P$ in the chart $c$
of $x$.
\end{lemma}
\leanref{Development/WeilBound/CoreStatements/HodgeTatePoints/MaSerre.lean}{PositivityTheory.HtPt.Ma.serre}
\leanref{Development/WeilBound/CoreStatements/HodgeTatePoints/MaSerre.lean}{PositivityTheory.HtPt.Ma.coh_chart}
\leanref{Development/WeilBound/CoreStatements/HodgeTatePoints/MaSerre.lean}{PositivityTheory.HtPt.Ma.twist_mem}
\leanref{Development/AutomorphicSpectrum/Proofs/GeometryProofs/ExistsGenOnChart.lean}{GlobalTheory.CompanionTheory.Ar.ChCoh.finite_gen}
\leanref{Development/AutomorphicSpectrum/Proofs/GeometryProofs/ExistsGenOnChart.lean}{GlobalTheory.CompanionTheory.Ar.ChCoh.pow_smul_mem}
\begin{proof}
The chart $c$ is empty if $x_c=0$. For an index $c$ with $x_c\neq0$ let $X_c$ and $A_c=\mathbf Z[x_j/x_c:j\in J]$ be as in
Lemma~\ref{cv-origin-reference-descent-roots-piece}. By Lemma~\ref{cv-origin-snc-charts}, applied to $x$ as a point with one
block, an element of $\Omega$ is regular at a point $Q$ of $X_c$ if and only if it is $u/s$ with $u,s\in A_c$ and $s$ a unit at
$Q$, and for every prime $\mathfrak p$ of $A_c$ there is a point of $X_c$ at which the elements of $A_c$ that are units are those
outside $\mathfrak p$.

\emph{Saturation.} Let $x_c\neq0$, let $g,f\in A_c$, let $t$ be a finite set of vectors with $L(Q)=\sum_{e'\in t}\mathcal O_Qe'$
at every point $Q$ of $X_c$ at which $g$ is a unit, and let $e\in V$ lie in $L(Q)$ at every point $Q$ of $X_c$ at which $f$ is a
unit. Then there is $m\ge0$ such that $f^me\in L(Q)$ at every point $Q$ of $X_c$ at which $g$ is a unit. Indeed, let $I$ be the
ideal of the $a\in A_c$ with $g^kae\in\sum_{e'\in t}A_ce'$ for some $k\ge0$. Let $\mathfrak p$ be a prime of $A_c$ that contains
$I$, and suppose that $gf\notin\mathfrak p$. Let $Q$ be the point of $X_c$ at which the elements of $A_c$ that are units are those
outside $\mathfrak p$. Then $g$ and $f$ are units at $Q$, so $e\in\sum_{e'\in t}\mathcal O_Qe'$. Writing the coefficients as
fractions with denominators in $A_c$ that are units at $Q$, and multiplying by the product $s$ of these denominators, gives
$se\in\sum_{e'\in t}A_ce'$; so $s\in I\subset\mathfrak p$, although $s$ is a unit at $Q$. Hence $gf$ lies in every prime that
contains $I$, and $(gf)^m\in I$ for some $m$ \cite[Tag 00E0]{Sta}: $g^k(gf)^me\in\sum_{e'\in t}A_ce'$ for some $k$. At a point $Q$
of $X_c$ at which $g$ is a unit, $f^me=g^{-k-m}\cdot g^k(gf)^me$ lies in $\sum_{e'\in t}\mathcal O_Qe'=L(Q)$.

\emph{Generators on a chart.} Let $x_c\neq0$. For every point $P$ of $X_c$ there are, by hypothesis and by
Lemma~\ref{cv-origin-reference-descent-roots-piece}(1), a nonzero $g_P\in A_c$ that is a unit at $P$ and a finite set $t_P$ of
vectors with $L(Q)=\sum_{e\in t_P}\mathcal O_Qe$ at every point $Q$ of $X_c$ at which $g_P$ is a unit. The $g_P$ generate the unit
ideal of $A_c$: a maximal ideal that contained them all would be the set of the elements of $A_c$ that are not units at some point
$P$ of $X_c$, but $g_P$ is a unit at $P$. So $1=\sum_{l=1}^Na_lg_{P_l}$ for some points $P_1,\dots,P_N$ and some $a_l\in A_c$, and
at every point $Q$ of $X_c$ some $g_{P_l}$ is a unit, as the elements of $A_c$ that are not units at $Q$ lie in the maximal ideal of
$\mathcal O_Q$. For $l\le N$ and $e\in t_{P_l}$, the saturation step, applied with $g=g_{P_{l'}}$, $t=t_{P_{l'}}$ and $f=g_{P_l}$ for
each $l'\le N$, gives $m_{l,e}\ge0$ with $g_{P_l}^{m_{l,e}}e\in L(Q)$ at every point $Q$ of $X_c$, as $g_{P_l}\in A_c$ is regular at
$Q$. Let $t_c$ be the finite set of these vectors $g_{P_l}^{m_{l,e}}e$. At a point $Q$ of $X_c$ at which $g_{P_l}$ is a unit, $L(Q)$
is generated by the vectors $e\in t_{P_l}$, each of which is $g_{P_l}^{-m_{l,e}}$ times an element of $t_c$; so
$L(Q)=\sum_{e\in t_c}\mathcal O_Qe$.

\emph{Twisting.} Let $x_c,x_j\neq0$ and $e\in t_c$. A point of $X_j$ at which $x_c/x_j$ is a unit lies in the chart $c$, as
$x_i/x_c=(x_i/x_j)(x_c/x_j)^{-1}$, so $e$ lies in $L$ there. So the saturation step on the chart $j$, with $g=1$, $t=t_j$ and
$f=x_c/x_j$, gives $m_{c,j,e}\ge0$ with $(x_c/x_j)^{m_{c,j,e}}e\in L(P)$ for every $P\in X_j$. Let $d$ be the largest
$m_{c,j,e}$, and let $z$ be the family of the vectors $x_c^de$ with $x_c\neq0$ and $e\in t_c$. Let $P$ be a point in the chart
$j$ of $x$. Then $L(P)=\sum_{e\in t_j}\mathcal O_Pe$, and each $e\in t_j$ is $x_j^{-d}$ times the vector $x_j^de$ of $z$.
Conversely, $x_j^{-d}x_c^de=(x_c/x_j)^{d-m}\cdot(x_c/x_j)^me\in L(P)$ with $m=m_{c,j,e}$, as $x_c/x_j$ is regular at $P$. So
$L(P)=\sum_k\mathcal O_Px_j^{-d}z_k$.
\end{proof}
% Lean: serre: coh_chart (ChCoh.finite_gen, with Ma.basic for g_P, Ma.prime_pt for the unit ideal and the maximal ideal of
% regAt q for some g_P a unit at q; pow_smul_mem for the saturation step, through satIdeal, clear_denoms and
% Ideal.radical_eq_sInf), twist_mem (the saturation step with g = 1, satIdeal (span t) 1 e), and the family
% x c^d * b c k with d an upper bound of the m (s, j).

\begin{lemma}\label{cv-origin-reference-refinement-flatten}
Let $X$, $\Omega$, $\xi$ and $x$ be as in Lemma~\ref{cv-origin-reference-refinement-twist}, let $V$ be an $\Omega$-vector space
with a finite-dimensional subspace $W$, and let $L$ be a lattice family in $V$ such that every point of $X$ has an open
neighbourhood $U$ and finitely many vectors that generate $L(Q)$ over $\mathcal O_Q$ for every $Q\in U$, and such that $L(P)$
spans $W$ for every point $P$. Then there is a finite family $v$ in $\Omega$ with a nonzero entry such that:
\begin{enumerate}
\item every point $P$ of $X$ at which $L(P)$ is generated over $\mathcal O_P$ by linearly independent vectors lies in a chart of
  $v$;
\item for every scheme $X'$ with a morphism $\xi':\operatorname{Spec}\Omega\to X'$ that satisfies condition (1) of
  Definition~\ref{cv-origin-segre-coordinates}, and every morphism $g:X'\to X$ with $g\circ\xi'=\xi$, the pullback $g^*L$ is a
  vector bundle of rank $\dim W$ with generic fibre $W$ on the set of points of $X'$ that lie in a chart of $v$
  (Definition~\ref{cv-origin-reference-lattices}).
\end{enumerate}
\end{lemma}
\leanref{Development/WeilBound/CoreStatements/HodgeTatePoints/MbLFlatten.lean}{PositivityTheory.HtPt.Mb.flatten}
\leanref{Development/WeilBound/CoreStatements/HodgeTatePoints/MbLPullEq.lean}{PositivityTheory.HtPt.Mb.pull_eq}
\begin{proof}
Let $\eta$ be the image of $\xi$. Lemma~\ref{cv-origin-reference-refinement-twist} gives $d\ge0$ and $z_1,\dots,z_n\in V$ with
$L(P)=\sum_k\mathcal O_Px_c^{-d}z_k$ for every point $P$ in a chart $c$ of $x$. The point $\eta$ lies in a chart $c_0$, and
$\mathcal O_\eta=\Omega$; so $L(\eta)$ is the subspace spanned by the $x_{c_0}^{-d}z_k$, and as it spans $W$, the $z_k$ lie in $W$
and span it. Let $r=\dim W$, fix a basis $\beta$ of $W$, and for a map $c:\{1,\dots,r\}\to\{1,\dots,n\}$ let $v_c$ be the
determinant of $(z_{c(1)},\dots,z_{c(r)})$ in the basis $\beta$. Some $r$ of the $z_k$ form a basis of $W$, and the corresponding
$v_c$ is nonzero. For $t\in\Omega^\times$, the determinant of $(tz_{c(1)},\dots,tz_{c(r)})$ is $t^rv_c$, so the family $tz$ has the
same ratios $v_c/v_{c'}$ of determinants as $z$.

(1) Let $L(P)=\sum_i\mathcal O_Pf_i$ with $f_1,\dots,f_s$ linearly independent, and let $P$ lie in the chart $c$ of $x$. The $f_i$
span $W$, as $L(P)$ does, so $s=r$, and $\sum_k\mathcal O_Px_c^{-d}z_k=\sum_i\mathcal O_Pf_i$. The ring $\mathcal O_P$, the image of
a local ring, is local; so Lemma~\ref{cv-origin-reference-refinement-plucker}(2), applied to the family $x_c^{-d}z$, gives $c_0$ with
$v_{c_0}\neq0$ and $v_c/v_{c_0}\in\mathcal O_P$ for every $c$: $P$ lies in the chart $c_0$ of $v$.

(2) Let $\eta'$ be the image of $\xi'$, and $P'$ a point of $X'$ in the chart $l$ of $v$. By Lemma~\ref{cv-origin-segre-regular}(3),
$g(P')$ specializes $\eta$; let it lie in the chart $c$ of $x$. By Lemma~\ref{cv-origin-segre-regular}(1), $P'$ has an open
neighbourhood $U$ at every point of which that specializes $\eta'$ the ratios $v_{c'}/v_l$ and $x_j/x_c$ are regular. Put
$b_i=x_c^{-d}z_{l(i)}$ for $i\le r$; as $v_l\neq0$, $b_1,\dots,b_r$ is a basis of $W$. Let $Q'\in U$ lie in a chart of $v$. By
Lemma~\ref{cv-origin-segre-regular}(3), $g(Q')$ specializes $\eta$, the elements regular at $g(Q')$ are regular at $Q'$, and, as
$\mathcal O_{X',Q'}\to\Omega$ is injective, an element regular at $g(Q')$ that is a unit at $Q'$ is a unit at $g(Q')$. Let $g(Q')$ lie
in the chart $c_1$ of $x$. The element $x_c/x_{c_1}$ is regular at $g(Q')$, and its inverse is regular at $Q'$; so it is a unit at
$Q'$, hence at $g(Q')$, and every $x_j/x_c=(x_j/x_{c_1})(x_c/x_{c_1})^{-1}$ is regular at $g(Q')$: $g(Q')$ lies in the chart $c$.
Hence $L(g(Q'))=\sum_k\mathcal O_{g(Q')}x_c^{-d}z_k$, and, as $\mathcal O_{g(Q')}\subset\mathcal O_{Q'}$,
$(g^*L)(Q')=\sum_k\mathcal O_{Q'}x_c^{-d}z_k$. The ratios of the determinants of the family $x_c^{-d}z$ are the $v_{c'}/v_l$, which lie
in $\mathcal O_{Q'}$; so Lemma~\ref{cv-origin-reference-refinement-plucker}(1) gives $(g^*L)(Q')=\sum_i\mathcal O_{Q'}b_i$.
\end{proof}
% Lean: flatten = LFlatten.Flatten.plucker_block: Ma.serre (Lemma cv-origin-reference-refinement-twist), hzV and hspanE for
% the span, w l = b0.det (zV o l) with Module.finBasis, exists_det_ne, dominant_minor (frame_plucker with det_smul_cols) for
% (1), and vb_on for (2): reg_near for U, ext_basis, the centre q0 of q' (Mb.centre = exists_centre, whose chart clause is the
% unit argument above), pullloc_centre and cramer_sub. The Lean states (2) for the local pullback pullLoc along a field
% embedding i and an additive map j with IsExtOver i j, on the points of ctrSet; Mb.pull_eq identifies it with the pullback
% g^*L along a morphism g, the case used here.

\begin{lemma}\label{cv-origin-reference-refinement-join}
Let $X$, $\Omega$, $\xi$ and $x$ be as in Lemma~\ref{cv-origin-reference-refinement-twist}, and let $w_a$, for $a$ in a finite
set $A$, be finite families in $\Omega$, each with a nonzero entry. Then there are a finite family $x'$ in $\Omega$, a scheme $X'$
with a morphism $\xi':\operatorname{Spec}\Omega\to X'$ of whose image every point is a specialization, with the coordinates $x'$,
and a morphism $f:X'\to X$ with $f\circ\xi'=\xi$, such that every point of $X'$ lies in a chart of each $w_a$, and such that for
every point $P$ of $X$ that lies in a chart of each $w_a$ there is exactly one point $P'$ of $X'$ with $f(P')=P$, and
$\mathcal O_{P'}=\mathcal O_P$.
\end{lemma}
\leanref{Development/WeilBound/CoreStatements/HodgeTatePoints/MbLJoinMany.lean}{PositivityTheory.HtPt.Mb.join_many}
\leanref{Development/WeilBound/CoreStatements/HodgeTatePoints/MbLJoinIso.lean}{PositivityTheory.HtPt.Mb.join_iso}
\begin{proof}
\emph{One family.} Let $w$ be the family of the products $w_c=\prod_aw_{a,c_a}$, indexed by the tuples $c=(c_a)_{a\in A}$ of
indices of the $w_a$. Some $w_c$ is nonzero, as each $w_a$ has a nonzero entry. Let $Q$ be a point of a scheme with a morphism
from $\operatorname{Spec}\Omega$ that specializes its image. If $Q$ lies in the chart $c_a$ of each $w_a$, then $w_c\neq0$ and
every ratio $w_{c'}/w_c=\prod_aw_{a,c'_a}/w_{a,c_a}$ is regular at $Q$, so $Q$ lies in the chart $c$ of $w$. Conversely, if $Q$
lies in the chart $c$ of $w$, then each factor $w_{a,c_a}$ of $w_c$ is nonzero, and each $w_{a,l}/w_{a,c_a}$ is the ratio
$w_{c'}/w_c$ for the tuple $c'$ obtained from $c$ by replacing $c_a$ with $l$; so $Q$ lies in the chart $c_a$ of $w_a$.

\emph{The model.} Lemma~\ref{cv-origin-segre-join}, applied to $x$ and $w$, gives a scheme $X'$ with a morphism
$\xi':\operatorname{Spec}\Omega\to X'$ of whose image every point is a specialization, with the coordinates
$x'=(x_jw_c)_{(j,c)}$, and a morphism $f:X'\to X$ with $f\circ\xi'=\xi$. A point $P'$ of $X'$ lies in a chart $(j,c)$ of $x'$; then
$w_c\neq0$, and every $w_{c'}/w_c=x_jw_{c'}/(x_jw_c)$ is regular at $P'$. So $P'$ lies in the chart $c$ of $w$, hence in a chart of
each $w_a$.

\emph{A point over $P$.} Let $P$ be a point of $X$ in a chart of each $w_a$, hence in a chart $l$ of $w$, and in a chart $j_0$ of
$x$. Then $x'_{(j_0,l)}\neq0$, and the ratios $x'_{(j,c)}/x'_{(j_0,l)}=(x_j/x_{j_0})(w_c/w_l)$ are regular at $P$, so the ring $A'$
they generate lies in $\mathcal O_P$. Let $\mathfrak p$ be the prime of the elements of $A'$ that are not units at $P$. By
Lemma~\ref{cv-origin-snc-charts}, applied to $X'$ and $x'$ as a point with one block, there is a point $P'$ of $X'$ in the chart
$(j_0,l)$ of $x'$ at which the elements of $A'$ that are units are those outside $\mathfrak p$, and every element of
$\mathcal O_{P'}$ is $u/s$ with $u,s\in A'$ and $s$ a unit at $P'$. Then $s\notin\mathfrak p$ is a unit at $P$, so
$\mathcal O_{P'}\subset\mathcal O_P$; and if $u/s$ is a unit at $P$, then so is $u$, which is then a unit at $P'$, and $u/s$ is a
unit at $P'$. Let $Q=f(P')$. By Lemma~\ref{cv-origin-segre-regular}(3), $\mathcal O_Q\subset\mathcal O_{P'}$, and an element of
$\mathcal O_Q$ that is a unit at $P'$ is a unit at $Q$. The ratios $x_j/x_{j_0}=x'_{(j,l)}/x'_{(j_0,l)}$ are regular at $P'$. If $Q$
lies in the chart $j_1$ of $x$, then $x_{j_0}/x_{j_1}$ is regular at $Q$, and at $P'$ it is a unit, as its inverse is regular there;
so it is a unit at $Q$, and $Q$ lies in the chart $j_0$ of $x$. For $u\in A_{j_0}=\mathbf Z[x_j/x_{j_0}]$, which lies in
$\mathcal O_Q\subset\mathcal O_{P'}\subset\mathcal O_P$, the three conditions that $u$ is a unit at $Q$, at $P'$ and at $P$ are
therefore equivalent. So $P$ and $Q$ are points of the chart $j_0$ of $x$ at which the same elements of $A_{j_0}$ are units, and
$Q=P$ by Lemma~\ref{cv-origin-snc-charts}(2). Hence $f(P')=P$, and
$\mathcal O_P=\mathcal O_Q\subset\mathcal O_{P'}\subset\mathcal O_P$, so $\mathcal O_{P'}=\mathcal O_P$.

\emph{Uniqueness.} Let $P''$ be a point of $X'$ with $f(P'')=P$. By Lemma~\ref{cv-origin-segre-regular}(3),
$\mathcal O_P\subset\mathcal O_{P''}$, and an element of $\mathcal O_P$ that is a unit at $P''$ is a unit at $P$. The ratios of $x'$
at $(j_0,l)$ lie in $\mathcal O_P$, so $P''$ lies in the chart $(j_0,l)$, and for $u\in A'\subset\mathcal O_P=\mathcal O_{P'}$, $u$ is
a unit at $P''$ if and only if it is a unit at $P$, that is, at $P'$. By Lemma~\ref{cv-origin-snc-charts}(2), $P''=P'$.
\end{proof}
% Lean: join_many = LJoinMany.join_finitely_many: the product family over (forall a, kappa a), join_exists (Lemma
% cv-origin-coords-join, here with the identity of Omega, which is Lemma cv-origin-segre-join), join_charts and
% JoinMany.inChart_of_inChart_prod for the charts, JoinMany.inChart_prod_of_forall and join_iso = JoinIso.isoAt_of_chart
% (Mb.centre for P', with hle and hun; dom_local for f(P'); hu and hqc for the chart j0; eq_of_units, through reg_le_spec
% and the antisymmetry of specialization, where the book uses Lemma cv-origin-snc-charts(2); uniq for P''). The Lean's
% conclusion is IsoAt (existsUnique and regAt_eq).

\begin{lemma}\label{cv-origin-reference-refinement}
Let $\Omega$ be a field, $T$ a scheme with a morphism $\xi:\operatorname{Spec}\Omega\to T$ of whose image every point is a
specialization, and $x$ a finite family in $\Omega$ such that $T$ has the coordinates $x$
(Definition~\ref{cv-origin-segre-coordinates}). Let $w_a$, for $a$ in a finite set, be finite families in $\Omega$, each
with a nonzero entry. For $b$ in a finite set let $V_b$ be an $\Omega$-vector space with a finite-dimensional subspace
$W_b$, and $L_b$ a lattice family in $V_b$ (Definition~\ref{cv-origin-reference-lattices}) such that every point of $T$ has an
open neighbourhood $U$ and finitely many vectors that generate $L_b(Q)$ over $\mathcal O_Q$ for every $Q\in U$, and such that
$L_b(P)$ spans $W_b$ for every point $P$. Then there are a finite family $x'$ in $\Omega$, a scheme $T'$ with a morphism
$\xi':\operatorname{Spec}\Omega\to T'$ of whose image every point is a specialization, with the coordinates $x'$, and a
morphism $f:T'\to T$ with $f\circ\xi'=\xi$, such that:
\begin{enumerate}
\item the local rings of $T'$ are integrally closed in $\Omega$;
\item every point of $T'$ lies in a chart of each family $w_a$;
\item each pullback $f^*L_b$ is a vector bundle of rank $\dim W_b$ with generic fibre $W_b$;
\item $f$ is a normalization over every point $P$ of $T$ that lies in a chart of each $w_a$ and at which each $L_b(P)$ is
  generated over $\mathcal O_{T,P}$ by linearly independent vectors.
\end{enumerate}
\end{lemma}
\leanref{Development/AutomorphicSpectrum/Proofs/GeometryProofs/FlagModelCommonRefinementStatement.lean}{GlobalTheory.T.flag_model_common_refinement}
\begin{proof}
For each $b$, Lemma~\ref{cv-origin-reference-refinement-flatten}, applied to $T$, $x$ and $L_b$ with $W_b$, gives a finite family
$v_b$ in $\Omega$ with a nonzero entry. Lemma~\ref{cv-origin-reference-refinement-join}, applied to $T$, $x$ and the families $w_a$
and $v_b$, gives a finite family $x_1$, a scheme $T_1$ with a morphism $\xi_1:\operatorname{Spec}\Omega\to T_1$ of whose image
every point is a specialization, with the coordinates $x_1$, and a morphism $f_1:T_1\to T$ with $f_1\circ\xi_1=\xi$, such that
every point of $T_1$ lies in a chart of each $w_a$ and of each $v_b$, and such that over every point $P$ of $T$ in a chart of each
$w_a$ and of each $v_b$ there is exactly one point $P_1$ of $T_1$, with $\mathcal O_{P_1}=\mathcal O_P$.
Lemma~\ref{cv-origin-coords-normal}, applied to $T_1$ and $x_1$, gives a finite family $x'$, a scheme $T'$ with a morphism
$\xi':\operatorname{Spec}\Omega\to T'$ of whose image every point is a specialization, with the coordinates $x'$ and with local
rings integrally closed in $\Omega$, and a morphism $g:T'\to T_1$ with $g\circ\xi'=\xi_1$ that is a normalization over every
point of $T_1$ (Definition~\ref{cv-origin-reference-lattices}(5), with $\iota$ the identity). Put $f=f_1\circ g$. Then
$f\circ\xi'=\xi$, and (1) holds.

(2) Let $P'$ be a point of $T'$. The point $g(P')$ lies in a chart $l$ of $w_a$, and the ratios of $w_a$ at $l$, which are regular
at $g(P')$, are regular at $P'$ by Lemma~\ref{cv-origin-segre-regular}(3); so $P'$ lies in the chart $l$ of $w_a$. In the same way
every point of $T'$ lies in a chart of each $v_b$.

(3) The scheme $T'$ has coordinates, so it satisfies condition (1) of Definition~\ref{cv-origin-segre-coordinates}, and by the
previous paragraph every point of $T'$ lies in a chart of $v_b$. So Lemma~\ref{cv-origin-reference-refinement-flatten}(2), applied
to $f$, shows that $f^*L_b$ is a vector bundle of rank $\dim W_b$ with generic fibre $W_b$.

(4) Let $P$ be a point of $T$ in a chart of each $w_a$ at which each $L_b(P)$ is generated by linearly independent vectors. By
Lemma~\ref{cv-origin-reference-refinement-flatten}(1), $P$ lies in a chart of each $v_b$. So there is exactly one point $P_1$ of
$T_1$ with $f_1(P_1)=P$, and $\mathcal O_{P_1}=\mathcal O_P$; in particular $\mathcal O_P$ and $\mathcal O_{P_1}$ have the same
integral closure $\tilde A$ in $\Omega$. The points $P'$ of $T'$ with $f(P')=P$ are those with $g(P')=P_1$. As $g$ is a
normalization over $P_1$, the local ring at each of them is the localization of $\tilde A$ at a maximal ideal, and for every
maximal ideal $\mathfrak m$ of $\tilde A$ exactly one of them has the local ring $\tilde A_{\mathfrak m}$. So $f$ is a normalization
over $P$.
\end{proof}
% Lean: T.flag_model_common_refinement = Relay.flag_model_common_refinement = PositivityTheory.HtPt.Mb.top_common =
% LTopCommon.exists_common_normal_model: Mb.flatten for each b (Lemma cv-origin-reference-refinement-flatten), Mb.join_many
% for the families w a and v b indexed by A + B (Lemma cv-origin-reference-refinement-join), Mb.norm_exists with
% Omega' = Omega (Lemma cv-origin-coords-normal), dom_local for the charts of the points of T' (hch2), Mb.pull_eq with the
% clause hv2 of flatten for (3), and the first part of Mb.iso_comp (an isomorphism over P followed by a normalization over the
% point above it) with the clause hv1 of flatten for (4).

\subsection{Coordinates over a base ring}\label{subsection-cv-coordinates-base-ring}

\begin{definition}\label{cv-origin-reference-base-over}
Let $X$ be a scheme, $\Omega$ a field, $\xi:\operatorname{Spec}\Omega\to X$ a morphism and $x=(x_j)_{j\in J}$ a finite family in
$\Omega$, and let $o$ be a ring with a homomorphism $o\to\Omega$. The scheme $X$ \emph{has the coordinates $x$ over $o$} if it
satisfies Definition~\ref{cv-origin-segre-coordinates} with the $o$-algebra $o[x_i/x_j:i\in J]\subset\Omega$ generated by the ratios
in place of $\mathbf Z[x_i/x_j:i\in J]$ in condition (4). Over $o=\mathbf Z$ this is Definition~\ref{cv-origin-segre-coordinates}.
\end{definition}
\leanref{Development/AutomorphicSpectrum/Definitions/GlobalGeometry/CoordsAt.lean}{GlobalTheory.CompanionTheory.Ar.CoordTie}

\begin{lemma}\label{cv-origin-reference-base-integral}
Let $X$ be a scheme, $\Omega$ a field and $\xi:\operatorname{Spec}\Omega\to X$ a morphism of whose image every point of $X$ is a
specialization. Let $o$ be a ring with a homomorphism $o\to\Omega$, generated as an abelian group by elements $g_1,\dots,g_n$, and
suppose that $X$ has the coordinates $x=(x_j)_{j\in J}$ over $o$ (Definition~\ref{cv-origin-reference-base-over}). Let $w_0=1$ and
let $w_m$ be the image of $g_m$ in $\Omega$ for $1\le m\le n$. Then $X$ has the coordinates $(w_mx_j)_{(j,m)}$ over $\mathbf Z$.
\end{lemma}
\leanref{Development/AutomorphicSpectrum/Proofs/GeometryProofs/CoordOverIntOfFiniteOrderStatement.lean}{GlobalTheory.T.coord_over_int_of_finite_order}
\leanref{Development/AutomorphicSpectrum/Proofs/GeometryProofs/CoordOverIntOfFiniteOrder.lean}{GlobalTheory.CompanionTheory.Ar.coordTie_int_of_finite}
\begin{proof}
Put $x'_{(j,m)}=w_mx_j$, and for $x_c\neq0$ let $X_c$ be the affine open subscheme of the chart $c$ of $x$ and $A_c$ the
$o$-algebra generated by the ratios $x_j/x_c$. The $w_m$ lie in $A_c$, so they are regular at every point of $X_c$. Condition (1) of
Definition~\ref{cv-origin-segre-coordinates} does not depend on the family.

\emph{Charts.} A point $P$ lies in the chart $(c,m)$ of $x'$ if and only if it lies in the chart $c$ of $x$ and $w_m$ is a unit at $P$.
Indeed, if $P$ lies in the chart $(c,m)$, then $x_c\neq0$, every $x_{c'}/x_c=w_mx_{c'}/(w_mx_c)$ is regular at $P$, and so is
$w_m^{-1}=w_0x_c/(w_mx_c)$. Conversely, the ratios of $x'$ at $(c,m)$ are $w_{m'}(x_{c'}/x_c)w_m^{-1}$. As $w_0=1$, every point lies in
a chart $(c,0)$ of $x'$, which is condition (2). By the proof of Lemma~\ref{cv-origin-reference-descent-roots-piece}(2), with $A_c$ in
place of $\mathbf Z[x_j/x_c]$, the points of $X_c$ at which $w_m$ is a unit form a standard open subset of the affine scheme $X_c$, so
an affine open subscheme; this is condition (3).

\emph{Rings.} Let $x_c,w_m\neq0$ and let $S$ be the ring generated by the ratios $w_{m'}x_{c'}/(w_mx_c)$, whose elements are regular at
every point of the chart $(c,m)$. It contains $w_m^{-1}=w_0x_c/(w_mx_c)$. As $o$ is a finitely generated abelian group, its elements
are integral over $\mathbf Z$ \cite[Tag 00GK]{Sta}, and a nonzero element integral over $\mathbf Z$ whose inverse lies in $S$ lies in
$S$, as in the proof of Lemma~\ref{cv-origin-coords-normal}. So $w_m\in S$, every $w_{m'}=(w_{m'}x_c/(w_mx_c))w_m$ lies in $S$, the
image of $o$, which is generated by the $w_{m'}$ as an abelian group, lies in $S$, and every $x_{c'}/x_c=w_mx_{c'}/(w_mx_c)$ lies in
$S$; so $A_c\subset S$. Let $a$ be regular at every point of the chart $(c,m)$. By the proof of
Lemma~\ref{cv-origin-reference-descent-roots-piece}(2), with $A_c$ in place of $\mathbf Z[x_j/x_c]$, $w_m^ka\in A_c$ for some $k$; so
$a=w_m^ka\cdot(w_m^{-1})^k\in S$. This is condition (4).
\end{proof}
% Lean: coordTie_int_of_finite: Module.Finite.exists_fin for g, w = Option.elim 1 g, CoordBase.coordTie_mul (inChart_mul_iff
% for the charts, data and mem_adjoin_iff for the section f with value w m, basicOpen for (3), mem_of_inv_mem for w m in S,
% hgen for the image of o, exists_pow_mul_mem for w m^k a), and coordTie_reindex for the numbering by Fin (N + 1).

\begin{lemma}\label{cv-origin-reference-base-rings}
Let $o$ be a domain with field of fractions $F$, $\Omega$ a field with a homomorphism $F\to\Omega$, and $t$ a finite family in
$\Omega$. Let $X$ and $X'$ be schemes with morphisms $\xi:\operatorname{Spec}\Omega\to X$ and $\xi':\operatorname{Spec}\Omega\to X'$ such
that $X$ has the coordinates $t$ over $o$ and $X'$ the coordinates $t$ over $F$ (Definition~\ref{cv-origin-reference-base-over}).
Then every local ring $\mathcal O_{P'}$ of $X'$, at a point $P'$ that specializes the image of $\xi'$, is a local ring $\mathcal O_P$ of
$X$, as subrings of $\Omega$; and if $P$ is a point of $X$ that specializes the image of $\xi$ at which every nonzero element of $o$
is a unit, then $\mathcal O_P$ is a local ring of $X'$.
\end{lemma}
\leanref{Development/AutomorphicSpectrum/Proofs/GeometryProofs/CoordChartsMatchRegularStatement.lean}{GlobalTheory.T.coord_charts_match_regular}
\leanref{Development/AutomorphicSpectrum/Proofs/GeometryProofs/CoordChartsMatchRegular.lean}{GlobalTheory.CompanionTheory.Ar.SncBase.FracPts.regAt_frac}
\begin{proof}
For a chart $c$ of $t$ let $A_c$ and $B_c$ be the $o$-algebra and the $F$-algebra generated by the ratios $t_j/t_c$. Every element
of $B_c$ is $b/d$ with $b\in A_c$ and $d\in o$ nonzero, by clearing the denominators of its coefficients. Lemma~\ref{cv-origin-snc-charts},
with $A_c$ or $B_c$ in place of $\mathbf Z[t_j/t_c]$, holds with the same proof.

\emph{Equal units give equal rings.} Let $P\in X$ and $P'\in X'$ lie in the chart $c$, and suppose that the elements of $A_c$ that are
units at $P$ are those that are units at $P'$. Then $\mathcal O_P=\mathcal O_{P'}$. Indeed, by Lemma~\ref{cv-origin-snc-charts}(1),
an element of $\mathcal O_P$ is $u/s$ with $u,s\in A_c$ and $s$ a unit at $P$, hence at $P'$, and $u/s\in\mathcal O_{P'}$, as
$A_c\subset B_c\subset\mathcal O_{P'}$. Conversely, an element $a$ of $\mathcal O_{P'}$ has $sa\in B_c$ for some $s\in B_c$ that is a unit
at $P'$; writing $s=b/d$ and $sa=b'/d'$ with $b,b'\in A_c$ and nonzero $d,d'\in o$, which are units at $P'$ as $F\subset\mathcal O_{P'}$,
the element $d'b=dd's\in A_c$ is a unit at $P'$, hence at $P$, and $d'ba=d\,b'\in A_c$; so $a\in\mathcal O_P$.

\emph{The two statements.} Let $P'$ be a point of $X'$, in the chart $c$. As $A_c\subset\mathcal O_{P'}$, the elements of $A_c$ that are
not units at $P'$ form a prime, and Lemma~\ref{cv-origin-snc-charts}(2) gives a point $P$ of $X$ in the chart $c$ at which the same
elements of $A_c$ are units; so $\mathcal O_P=\mathcal O_{P'}$. Let $P$ be a point of $X$, in the chart $c$, at which every nonzero
element of $o$ is a unit. Then $B_c\subset\mathcal O_P$, as $b/d$ with $b\in A_c$ and $d\neq0$ is regular at $P$. The elements of $B_c$
that are not units at $P$ form a prime, and Lemma~\ref{cv-origin-snc-charts}(2) gives a point $P'$ of $X'$ in the chart $c$ at which the
same elements of $B_c$, hence of $A_c$, are units; so $\mathcal O_P=\mathcal O_{P'}$.
\end{proof}
% Lean: regAt_frac: exists_prime_of_subset, prime_chart_existsUnique_point over o and over F (the Lean's form of Lemma
% cv-origin-snc-charts(2) for any coefficient ring), regAt_eq_of_units (regAt_iff_chart_localization, exists_mul_mem_adjoin for
% clearing denominators, unitAt_algebraMap).

\begin{lemma}\label{cv-origin-reference-base-projective}
Let $X$ be a scheme, $\Omega$ a field and $\xi:\operatorname{Spec}\Omega\to X$ a morphism with image $\eta$ such that the local ring
of $X$ at $\eta$ maps bijectively onto $\Omega$ and those at its specializations map injectively.
\begin{enumerate}
\item If $X$ is locally noetherian, the set of the specializations of $\eta$ is open and closed in $X$.
\item Let $o$ be a noetherian ring with a homomorphism $o\to\Omega$, and $i:X\to\mathbf P^n_o$ a closed immersion over $o$ such that
  $\xi$ lies over $o$. Let $i(\eta)$ lie in the standard open subset $D_+(T_{j_0})$, and let $t_j\in\Omega$ be the value at $\xi$ of the
  function $T_j/T_{j_0}$, for $0\le j\le n$. Then $X$ has the coordinates $t=(t_0,\dots,t_n)$ over $o$
  (Definition~\ref{cv-origin-reference-base-over}).
\end{enumerate}
\end{lemma}
\leanref{Development/AutomorphicSpectrum/Proofs/GeometryProofs/IsClopenSpecializationsGenPtStatement.lean}{GlobalTheory.T.isClopen_specializations_genPt}
\leanref{Development/AutomorphicSpectrum/Proofs/GeometryProofs/ProjectiveOverCoordsExistStatement.lean}{GlobalTheory.T.projective_over_coords_exist}
\leanref{Development/WeilBound/CoreStatements/HodgeTatePoints/MfCoords.lean}{PositivityTheory.HtPt.Mf.coordTie_of_all}
\begin{proof}
(1) The set $Z$ of the specializations of $\eta$ is its closure. Let $p\in Z$. An affine open neighbourhood $W$ of $p$ is the spectrum
of a noetherian ring, so it has finitely many irreducible components \cite[Tag 0052]{Sta}, and the complement $V$ in $W$ of those
that do not contain $p$ is an open neighbourhood of $p$. For $q\in V$, the generic point $z$ of an irreducible component of $W$ that
contains $p$ and $q$ generizes $p$. Every generization of $p$ is the image of a prime of the local ring at $p$ under the canonical
morphism to $X$ \cite[Tag 01J7]{Sta}. The generization $\eta$ is the image of the preimage of the maximal ideal of the local ring at
$\eta$, a field; this preimage is the kernel of the map to $\Omega$, which is $0$. As $0$ is contained in every prime and the canonical
morphism is continuous, every generization of $p$, in particular $z$, specializes $\eta$, and so does $q$: $V\subset Z$.

(2) The scheme $X$ is of finite type over the noetherian ring $o$, so it is locally noetherian, and by (1) the set $Z$ of the
specializations of $\eta$ is an open and closed subscheme; the composite $Z\to X\to\mathbf P^n_o$ is a closed immersion, and $\xi$
factors through $Z$. If $i(\eta)\in D_+(T_c)$, then on $D_+(T_c)\cap D_+(T_{j_0})$ the function $T_c/T_{j_0}$ is invertible and
$T_j/T_c=(T_j/T_{j_0})/(T_c/T_{j_0})$; so $t_c\neq0$, and the value of $T_j/T_c$ at $\xi$ is $t_j/t_c$.

\emph{Charts.} A point $q$ of $Z$ lies in the chart $c$ of $t$ if and only if $i(q)\in D_+(T_c)$. If $i(q)\in D_+(T_d)$, then
$i(\eta)\in D_+(T_d)$, as $\eta$ generizes $q$; so $t_d\neq0$, and the functions $T_j/T_d$ give sections over $i^{-1}(D_+(T_d))$ whose
values at $\xi$ are the $t_j/t_d$, which are therefore regular at $q$. So if $i(q)\in D_+(T_c)$, $q$ lies in the chart $c$. Conversely, let $q$
lie in the chart $c$ and $i(q)\in D_+(T_d)$. Then $t_c/t_d$, the value of the section $T_c/T_d$, is a unit at $q$; as the local ring at
$q$ embeds in $\Omega$, the germ of $T_c/T_d$ is a unit, so $i(q)$ lies in $D_+(T_d)\cap D_+(T_c)$. Every point of $Z$ lies in some
$D_+(T_d)$, which gives condition (2) of Definition~\ref{cv-origin-segre-coordinates}; condition (1) is the hypothesis on $\xi$. The
points of $Z$ in the chart $c$ form $Z\cap i^{-1}(D_+(T_c))$, the preimage of an affine open subscheme under a closed immersion, which is
affine; this is condition (3).

\emph{Rings.} The ring of $D_+(T_c)$ is $o[T_j/T_c]$, and as $i$ restricted to $Z$ is a closed immersion, the ring map to the sections of
$Z\cap i^{-1}(D_+(T_c))$ is surjective. So the values at $\xi$ of these sections are the elements of $o[t_j/t_c]$. Every point of $Z$
specializes $\eta$, so by Lemma~\ref{cv-origin-segre-regular}(2) these values are the elements of $\Omega$ regular at every point of
the chart; this is condition (4). As $Z$ is open in $X$, the local rings of $X$ at the points of $Z$ are those of $Z$, and $X$ has the
coordinates $t$ over $o$.
\end{proof}
% Lean: (1) is isClopen_setOf_genPt_specializes (the closure, exists_isOpen_common_generization in the noetherian affine
% neighbourhood, genPt_specializes_of_specializes). (2) is Relay.projective_over_coords_exist: the open U of the
% specializations, IsOpenImmersion.lift, IsClosedImmersion.of_isPreimmersion, isCptGen_of_comp, HtPt.Mf.exists_point (the
% coordinates t of the point, with t i = 1), HtPt.Mf.coordTie_of_all (inChart_iff, valR_point_chart_aeval_ratio,
% exists_valR_eq_iff_mem_adjoin_ratio with app_surjective of the closed immersion, and
% exists_valR_eq_iff_forall_inChart_mem_regAt, which argues through the function field of the integral scheme and
% PrimeSpectrum.iInf_localization_eq_bot where the book uses Lemma cv-origin-segre-regular(2)), and HtPt.Mf.coordTie_comp for
% the open immersion of U into X.

\begin{lemma}\label{cv-origin-reference-base-product}
Let $X$ be a scheme, $\Omega$ a field and $\xi:\operatorname{Spec}\Omega\to X$ a morphism, let $o$ be a ring with a homomorphism
$o\to\Omega$, and suppose that $X$ has the coordinates $x=(x_c)_{c\in J}$ over $o$ (Definition~\ref{cv-origin-reference-base-over}).
Let $s=(s_i)_{i\in I}$ be a finite family in $\Omega$ such that every point of $X$ that specializes the image of $\xi$ lies in a chart
of $s$. Then there is an integer $n\ge0$ such that $X$ has the coordinates
$z=(x_{a_1}\cdots x_{a_n}x_cs_i)_{(a,c,i)\in J^n\times J\times I}$ over $o$.
\end{lemma}
\leanref{Development/AutomorphicSpectrum/Proofs/GeometryProofs/CoordsProductRefinementStatement.lean}{GlobalTheory.T.coords_product_refinement}
\leanref{Development/AutomorphicSpectrum/Proofs/GeometryProofs/CoordsProductRefinement.lean}{GlobalTheory.CompanionTheory.Ar.coordTie_add_block}
\begin{proof}
Take $n=0$. The products $x_{a_1}\cdots x_{a_n}$ are then $1$, and $z=(x_cs_i)_{(c,i)\in J\times I}$. For $x_c\neq0$ let $X_c$ be the
affine open subscheme of the chart $c$ of $x$, whose points specialize the image of $\xi$, and $A_c$ the $o$-algebra generated by
the $x_j/x_c$; by Lemma~\ref{cv-origin-segre-regular}(2), applied to $X_c$, and condition (4), the ring of $X_c$ maps
isomorphically onto $A_c$, and we identify the two. Lemmas~\ref{cv-origin-snc-charts}
and~\ref{cv-origin-reference-descent-roots-piece} hold over $o$, with $A_c$ in place of $\mathbf Z[x_j/x_c]$ and the same proofs,
which use only the points of $X_c$. Condition (1) of Definition~\ref{cv-origin-segre-coordinates} does not depend on the family.

\emph{Charts.} A point $P$ lies in the chart $(c,i)$ of $z$ if and only if it lies in the chart $c$ of $x$ and in the chart $i$ of $s$.
If it lies in both, the ratios $(x_{c'}/x_c)(s_j/s_i)$ are regular at $P$. Conversely, let $P$ lie in the chart $(c,i)$; then
$x_c,s_i\neq0$, and $P$ lies in some chart $c_0$ of $x$ and $i_0$ of $s$. The element $w=(x_{c_0}/x_c)(s_{i_0}/s_i)$ is regular at
$P$, so are $x_{c_0}/x_c=w\cdot s_i/s_{i_0}$ and $s_{i_0}/s_i=w\cdot x_c/x_{c_0}$, and then every $x_{c'}/x_c=(x_{c'}/x_{c_0})(x_{c_0}/x_c)$
and every $s_j/s_i$. As every point that specializes the image of $\xi$ lies in charts of $x$ and of $s$, this gives condition (2).

\emph{Affine charts.} Let $x_c,s_i\neq0$. Let $D$ be the set of the $g\in A_c$ with $gs_j/s_i\in A_c$ for every $j$, and $W$ the union
of the standard open subsets $D(g)$ of $X_c$, $g\in D$. A point $P$ of $X_c$ lies in $W$ if and only if it lies in the chart $i$ of $s$:
on $D(g)$ every $s_j/s_i=(gs_j/s_i)g^{-1}$ is regular; conversely, if the $s_j/s_i$ are regular at $P$, Lemma~\ref{cv-origin-snc-charts}(1)
gives a common denominator $g\in A_c$, a unit at $P$, with all $gs_j/s_i\in A_c$, and $P\in D(g)$. So $W$ is the set of the points in the
chart $(c,i)$ of $z$. Let $E$ be the set of the elements of $A_c$ of the form $hs_j$ with $h\in\Omega$ and $hs_l\in A_c$ for every $l$.
They generate the unit ideal of $A_c$: otherwise they lie in the ideal of the elements that are not units at some point $P$ of $X_c$
(Lemma~\ref{cv-origin-snc-charts}(2)); $P$ lies in a chart $j$ of $s$, and a common denominator $g$ of the $s_l/s_j$ at $P$, as above,
is $hs_j$ with $h=g/s_j$, an element of $E$ that is a unit at $P$. An element $hs_j$ of $E$ is $a\cdot s_j/s_i$ with $a=hs_i$, which lies
in $D$. As the $s_j/s_i$ are regular at every point of $W$, they are the values of sections $\tau_j$ over $W$, by
Lemma~\ref{cv-origin-segre-regular}(2); so the restriction to $W$ of $hs_j$ is $a\tau_j$, and the restrictions of the elements of $D$
generate the unit ideal of the ring of $W$. For $a\in D$, the subset of $W$ where $a$ is invertible is $D(a)$, a standard open subset of
the affine $X_c$. So $W$ is affine \cite[Tag 01QF]{Sta}; this is condition (3).

\emph{Rings.} Let $S$ be the $o$-algebra generated by the ratios of $z$ at $(c,i)$; its elements are regular at every point of $W$. It
contains $x_{c'}/x_c=z_{(c',i)}/z_{(c,i)}$ and $s_j/s_i=z_{(c,j)}/z_{(c,i)}$, so $A_c\subset S$. Let $a$ be regular at every point of $W$,
and let $I$ be the ideal of the $r\in A_c$ with $ra\in S$. Every element $e=b\cdot s_j/s_i$ of $E$, with $b\in D$, lies in the radical of
$I$: this is clear if $b=0$; otherwise $a$ is regular at every point of $D(b)\subset W$, so $b^Na\in A_c$ for some $N$, by
Lemma~\ref{cv-origin-reference-descent-roots-piece}(2), and $e^Na=(b^Na)(s_j/s_i)^N\in S$. As $E$ generates the unit ideal, $1\in I$, and
$a\in S$. This is condition (4).
\end{proof}
% Lean: T.coords_product_refinement = coordTie_add_block, whose proof takes n = 0 (the product over Fin 0 -> kappa is 1) and
% applies AddBlock.coordTie_mul: inChart_of_mul and hchart' for the charts, denoms and chartOpen with mem_chartOpen
% (exists_common_den) for W, chartOpen_affine (span_trace, exists_denom_of_trace, exists_val_eq for tau,
% isAffineOpen_of_isAffineOpen_basicOpen) for (3), and mem_of_regular (the ideal c, exists_pow, Ideal.radical_eq_top) for (4).

\begin{definition}\label{cv-origin-reference-base-patch}
Let $\Omega$ be a field of characteristic zero, $\mathcal P$ a set of subrings of $\Omega$, $\mathcal B\subset\mathcal P$ and $e\ge0$. For
a subring $A\subset\Omega$ and a prime $\mathfrak q$ of $A$ let $A_{\mathfrak q}$ be the ring of the $a\in\Omega$ with $sa\in A$ for some
$s\in A\setminus\mathfrak q$. A \emph{normal crossings patch} of dimension $e$ for $(\mathcal P,\mathcal B)$ is a finitely generated
$\mathbf Q$-subalgebra $A$ of $\Omega$, smooth over $\mathbf Q$, with elements $z_1,\dots,z_e$ whose differentials form a basis of
$\Omega_{A/\mathbf Q}$ and an integer $n\le e$, such that for every prime $\mathfrak q$ of $A$, $A_{\mathfrak q}\in\mathcal P$, and
$A_{\mathfrak q}\in\mathcal B$ if and only if $\mathfrak q$ contains one of $z_1,\dots,z_n$.
\end{definition}
\leanref{Development/AutomorphicSpectrum/Proofs/GeometryProofs/IsSncPatchRat.lean}{GlobalTheory.CompanionTheory.Ar.SncBase.IsSncPatch}
\leanref{Development/AutomorphicSpectrum/Proofs/GeometryProofs/IsSncPatchRat.lean}{GlobalTheory.CompanionTheory.Ar.SncBase.IsLocAt}

\begin{lemma}\label{cv-origin-reference-base-patches}
Let $F$ be a number field, $\Omega$ a field with a homomorphism $F\to\Omega$, $X'$ a scheme smooth of relative dimension $e$ over $F$,
and $\xi':\operatorname{Spec}\Omega\to X'$ a morphism over $F$ such that the local ring of $X'$ at its image $\eta'$ maps bijectively
onto $\Omega$ and those at the specializations of $\eta'$ map injectively, and such that distinct specializations of $\eta'$ have
distinct local rings $\mathcal O_{P'}\subset\Omega$. Let $T_\partial\subset X'$, and suppose that $X'$ is covered by finitely many affine
open subschemes $V$, each with functions $z_1,\dots,z_e$ whose differentials form a basis of $\Omega_{V/F}$ and an integer $k\le e$ such
that $V\cap T_\partial$ is the set of the points of $V$ at which one of $z_1,\dots,z_k$ vanishes. Let $\mathcal P$ be the set of the local
rings $\mathcal O_{P'}$ at the specializations $P'$ of $\eta'$, and $\mathcal B$ the set of those with $P'\in T_\partial$. Then for every
specialization $P'$ of $\eta'$ there are a normal crossings patch $A$ of dimension $e$ for $(\mathcal P,\mathcal B)$
(Definition~\ref{cv-origin-reference-base-patch}) and a prime $\mathfrak q$ of $A$ with $A_{\mathfrak q}=\mathcal O_{P'}$.
\end{lemma}
\leanref{Development/AutomorphicSpectrum/Proofs/GeometryProofs/SncPatchAtEachComponentStatement.lean}{GlobalTheory.T.snc_patch_at_each_component}
\leanref{Development/AutomorphicSpectrum/Proofs/GeometryProofs/SncPatchAtEachComponent.lean}{GlobalTheory.CompanionTheory.Ar.SncBase.Patch.exists_patch}
\begin{proof}
Let $V$ be an affine open subscheme of the cover that contains $P'$, with its $z_1,\dots,z_e$ and $k$; it contains the generization
$\eta'$ of $P'$. Its ring $\Gamma(V)$ is smooth over $F$ \cite[Tag 01V6]{Sta}, hence finitely generated over a field and noetherian.
Let $\mathrm{ev}:\Gamma(V)\to\Omega$ be the value at $\xi'$.

\emph{An injective chart.} The kernel of $\mathrm{ev}$ is generated by finitely many elements. The germ at $P'$ of each of them maps to
$0$ in $\Omega$, so it is $0$, as $\mathcal O_{P'}\to\Omega$ is injective; so each is killed by an element of $\Gamma(V)$ outside the
prime of $P'$. Let $g$ be the product of these elements. Then $P'\in D(g)$, $g$ kills the kernel of $\mathrm{ev}$, and $\mathrm{ev}$ is
injective on $\Gamma(D(g))=\Gamma(V)[1/g]$: if $a/g^m$ has value $0$, then $\mathrm{ev}(a)=0$, so $ga=0$ and $a/g^m=0$. Let $A\subset\Omega$
be the image of $\Gamma(D(g))$, isomorphic to it; it is a finitely generated $\mathbf Q$-algebra.

\emph{Smoothness and differentials.} The ring $\Gamma(D(g))$ is smooth over $F$ \cite[Tag 01V6]{Sta}. The field $F$ is a finite
separable extension of $\mathbf Q$, so it is \'etale \cite[Tag 00U3]{Sta}, hence formally \'etale \cite[Tag 00UR]{Sta}, over
$\mathbf Q$. Being formally smooth and of finite presentation, $F$ is smooth over $\mathbf Q$ \cite[Tag 00TN]{Sta}; so
$\Gamma(D(g))$ is smooth over $\mathbf Q$ \cite[Tag 00TD]{Sta}, and so is $A$, which is isomorphic to it. Being formally unramified,
$F$ has $\Omega_{F/\mathbf Q}=0$ \cite[Tag 00UO]{Sta}. So in the exact sequence
$\Gamma(V)\otimes_F\Omega_{F/\mathbf Q}\to\Omega_{\Gamma(V)/\mathbf Q}\to\Omega_{\Gamma(V)/F}\to0$ \cite[Tag 00RS]{Sta} the second map,
which sends $dz_i$ to $dz_i$, is an isomorphism, and the differentials of $z_1,\dots,z_e$ form a basis of
$\Omega_{\Gamma(V)/\mathbf Q}$. For a formally \'etale homomorphism $R\to S$ of $\mathbf Q$-algebras, the map
$S\otimes_R\Omega_{R/\mathbf Q}\to\Omega_{S/\mathbf Q}$, $s\otimes dr\mapsto s\,dr$, is an isomorphism \cite[Tag 0H93]{Sta}; so if the
differentials of some $r_1,\dots,r_e\in R$ form a basis of $\Omega_{R/\mathbf Q}$, those of their images form a basis of
$\Omega_{S/\mathbf Q}$. The localization $\Gamma(V)\to\Gamma(V)[1/g]=\Gamma(D(g))$ is formally \'etale \cite[Tag 04EG]{Sta}, and so
is the isomorphism $\Gamma(D(g))\to A$. Applying this twice, the differentials of the values $z_1,\dots,z_e\in A$ of the functions
$z_i$ form a basis of $\Omega_{A/\mathbf Q}$. Put $n=k$.

\emph{Primes.} Let $\mathfrak q$ be a prime of $A$, and $Q'$ the point of the affine scheme $D(g)$ whose prime is the preimage of
$\mathfrak q$. As $\mathrm{ev}$ is injective on $\Gamma(D(g))$, the prime of $\eta'$ is $0$, which $\mathfrak q$ contains, so $Q'$
specializes $\eta'$. The local ring of $D(g)$ at $Q'$ is the localization of $\Gamma(D(g))$ at the prime of $Q'$, so
$\mathcal O_{Q'}=A_{\mathfrak q}$, which lies in $\mathcal P$. By hypothesis, $A_{\mathfrak q}\in\mathcal B$ if and only if $Q'\in T_\partial$, as
distinct specializations of $\eta'$ have distinct local rings; that is, as $Q'\in V$, if and only if one of $z_1,\dots,z_k$ vanishes at
$Q'$, that is, lies in $\mathfrak q$. So $A$ is a normal crossings patch for $(\mathcal P,\mathcal B)$, and the prime $\mathfrak q$ of $P'$
has $A_{\mathfrak q}=\mathcal O_{P'}$.
\end{proof}
% Lean: Patch.exists_patch: the cover from HasSncCoords, exists_kill (IsNoetherian of the kernel of ev, the localization
% at the prime of p'), ev_basicOpen_injective, and patch_core: smooth_secAlgebra' (Smooth.smooth_appLE), the smoothness of F
% over Q (Algebra.FormallyEtale.of_isSeparable, finite presentation, Algebra.Smooth.comp), coords_rat on Gamma(V)
% (KaehlerDifferential.exact_mapBaseChange_map with Omega_(F/Q) = 0) and SncRes.coords_of_formallyEtale
% (KaehlerDifferential.tensorKaehlerEquivOfFormallyEtale, Stacks 0H93) twice, for the localization at g
% (FormallyEtale.of_isLocalization, Stacks 04EG) and the isomorphism onto A (FormallyEtale.of_equiv), specializes_fromSpec,
% isLocAt_of_point (stalk_eq_germ_div_of_affine), and the hypothesis hreg for the boundary clause.

\begin{situation}\label{cv-origin-reference-base-situation}
Let $\Omega$ be a field of characteristic zero, $X$ a scheme, $\xi:\operatorname{Spec}\Omega\to X$ a morphism with image $\eta$,
and $x=(x_0,\dots,x_N)$ a family in $\Omega$ such that $X$ has the coordinates $x$ (Definition~\ref{cv-origin-segre-coordinates}).
For a point $P$ that specializes $\eta$, let $\mathcal O_P\subset\Omega$ be the ring of the elements regular at $P$; the point
lies \emph{over no prime} if every prime number is a unit at it. For an index $c$ with $x_c\neq0$ put
$A_c=\mathbf Q[x_j/x_c:0\le j\le N]\subset\Omega$, and for $g\in\mathbf Z[X_0,\dots,X_N]$ write
$g(x/x_c)=g(x_0/x_c,\dots,x_N/x_c)\in A_c$. For a $\mathbf Q$-subalgebra $A$ of $\Omega$ and a nonzero $h\in A$, let $A[1/h]$ be
the $\mathbf Q$-subalgebra of $\Omega$ generated by $A$ and $h^{-1}$, and for a prime $\mathfrak q$ of $A$ let $A_{\mathfrak q}$ be
as in Definition~\ref{cv-origin-reference-base-patch}. Let $e\ge0$, let $B$ be a set of points of $X$ over no prime, $\mathcal P$
the set of the rings $\mathcal O_P$ at the points $P$ over no prime, and $\mathcal B$ the set of the $\mathcal O_P$ with $P\in B$.
A \emph{boundary form} is a homogeneous $g\in\mathbf Z[X_0,\dots,X_N]$ such that, for every index $c$, $g(x/x_c)$ is not a unit at
any point of $B$ in the chart $c$ of $x$. Let $J_B$ be the ideal of $\mathbf Z[X_0,\dots,X_N]$ generated by the boundary forms.
\end{situation}
\leanref{Development/AutomorphicSpectrum/Proofs/GeometryProofs/UnitAtOfMulRight.lean}{GlobalTheory.CompanionTheory.Ar.SncBase.Mdl.chartQ}
\leanref{Development/AutomorphicSpectrum/Proofs/GeometryProofs/UnitAtOfMulRight.lean}{GlobalTheory.CompanionTheory.Ar.SncBase.Mdl.invAdj}
\leanref{Development/AutomorphicSpectrum/Proofs/GeometryProofs/UnitAtOfMulRight.lean}{GlobalTheory.CompanionTheory.Ar.SncBase.Mdl.ratEval}
\leanref{Development/AutomorphicSpectrum/Proofs/GeometryProofs/UnitAtOfMulRight.lean}{GlobalTheory.CompanionTheory.Ar.SncBase.Mdl.bdGens}
\leanref{Development/AutomorphicSpectrum/Proofs/GeometryProofs/IsSncPatchRat.lean}{GlobalTheory.CompanionTheory.Ar.SncBase.Char0}
\leanref{Development/AutomorphicSpectrum/Proofs/GeometryProofs/IsSncPatchRat.lean}{GlobalTheory.CompanionTheory.Ar.SncBase.Pts}

\begin{lemma}\label{cv-origin-reference-base-localization}
Let $\Omega$ be a field of characteristic zero, $A$ and $A'$ finitely generated $\mathbf Q$-subalgebras of $\Omega$, and
$\mathfrak q\subset A$ and $\mathfrak q'\subset A'$ primes with $A_{\mathfrak q}=A'_{\mathfrak q'}$
(Definition~\ref{cv-origin-reference-base-patch}). For a nonzero $g\in\Omega$ let $A[1/g]$ be the $\mathbf Q$-subalgebra of
$\Omega$ generated by $A$ and $g^{-1}$, and $A'[1/g]$ that generated by $A'$ and $g^{-1}$. Then there are $g\in A\setminus\mathfrak q$
and $h\in A'\setminus\mathfrak q'$ with $A[1/g]=A'[1/h]$.
\end{lemma}
\leanref{Development/AutomorphicSpectrum/Proofs/GeometryProofs/LocalizationsEqualBasicOpensStatement.lean}{GlobalTheory.T.localizations_equal_basic_opens}
\leanref{Development/AutomorphicSpectrum/Proofs/GeometryProofs/LocalizationsEqualBasicOpens.lean}{GlobalTheory.CompanionTheory.Ar.SncBase.CommonLoc.exists_common_loc}
\begin{proof}
Put $R=A_{\mathfrak q}=A'_{\mathfrak q'}$. Every element of $A$ lies in $R$. An element $s\in A\setminus\mathfrak q$ is nonzero, and
$s^{-1}\in R$, as $s\cdot s^{-1}=1\in A$. Conversely, a nonzero $b\in A$ with $b^{-1}\in R$ lies outside $\mathfrak q$: there is
$s\in A\setminus\mathfrak q$ with $sb^{-1}\in A$, and $s=(sb^{-1})b$. The same holds for $A'$ and $\mathfrak q'$. For a nonzero
$g\in A$, every $a\in A[1/g]$ has $ag^k\in A$ for some $k\ge0$, as the elements with this property form a $\mathbf Q$-subalgebra
that contains $A$ and $g^{-1}$; the same holds for $A'$.

\emph{Denominators.} Let $b_1,\dots,b_m$ generate $A'$. They lie in $R$, so there are $s_i\in A\setminus\mathfrak q$ with
$s_ib_i\in A$. Then $g_0=\prod_is_i\notin\mathfrak q$, and $b_i=(s_ib_i)\prod_{l\neq i}s_l\cdot g_0^{-1}\in A[1/g_0]$; so
$A'\subset A[1/g_0]$. In the same way there is $h_0\in A'\setminus\mathfrak q'$ with $A\subset A'[1/h_0]$.

\emph{Clearing.} As $g_0\in A'[1/h_0]$, $a=g_0h_0^m$ lies in $A'$ for some $m$, and $a^{-1}=g_0^{-1}h_0^{-m}\in R$; so
$a\notin\mathfrak q'$, and $h=h_0a$ lies in $A'\setminus\mathfrak q'$. As $h\in A'\subset A[1/g_0]$, $a'=hg_0^l$ lies in $A$ for some
$l$, and $a'^{-1}=h^{-1}g_0^{-l}\in R$, so $a'\notin\mathfrak q$; put $g=g_0a'\in A\setminus\mathfrak q$. Then $h=h_0^{m+1}g_0$ and $g=g_0^{l+1}h$.

\emph{Equality.} The elements $h_0^{-1}=ah^{-1}$ and $g_0^{-1}=h_0^{m+1}h^{-1}$ lie in $A'[1/h]$; so
$A\subset A'[1/h_0]\subset A'[1/h]$, and $g^{-1}=g_0^{-(l+1)}h^{-1}\in A'[1/h]$. Hence $A[1/g]\subset A'[1/h]$. Conversely,
$g_0^{-1}=a'g^{-1}$, so $A'\subset A[1/g_0]\subset A[1/g]$, and $h^{-1}=g_0^{l+1}g^{-1}\in A[1/g]$. Hence $A'[1/h]\subset A[1/g]$.
\end{proof}
% Lean: exists_common_loc: exists_den twice (the denominators g0 and h0, by induction over the finite generating set),
% exists_clear twice (a and a', with exists_mul_pow_mem, unit_of_not_mem and not_mem_of_unit), invAdj_le_mul and invAdj_le_of
% for the two inclusions.

\begin{lemma}\label{cv-origin-reference-base-shrink}
Let $\Omega$ be a field of characteristic zero, $A$ a $\mathbf Q$-subalgebra of $\Omega$, $g\in A$ nonzero, and $A[1/g]$ the
$\mathbf Q$-subalgebra of $\Omega$ generated by $A$ and $g^{-1}$.
\begin{enumerate}
\item Every $s\in A[1/g]$ has $sg^k\in A$ for some $k\ge0$. Let $\mathfrak q$ be a prime of $A[1/g]$. Then $g\notin\mathfrak q$,
  $\mathfrak q\cap A$ is a prime of $A$, an element $s$ of $A[1/g]$ with $sg^k\in A$ lies in $\mathfrak q$ if and only if $sg^k$
  does, and $A[1/g]_{\mathfrak q}=A_{\mathfrak q\cap A}$ (Definition~\ref{cv-origin-reference-base-patch}).
\item Let $\mathcal P$ be a set of subrings of $\Omega$, $\mathcal B\subset\mathcal P$ and $e\ge0$. If $A$ is a normal crossings
  patch of dimension $e$ for $(\mathcal P,\mathcal B)$ with elements $z_1,\dots,z_e$ and an integer $n$, then so is $A[1/g]$, with
  the same $z_1,\dots,z_e$ and $n$.
\end{enumerate}
\end{lemma}
\leanref{Development/AutomorphicSpectrum/Proofs/GeometryProofs/UnitAtOfMulRight.lean}{GlobalTheory.CompanionTheory.Ar.SncBase.Mdl.exists_mul_pow_mem}
\leanref{Development/AutomorphicSpectrum/Proofs/GeometryProofs/UnitAtOfMulRight.lean}{GlobalTheory.CompanionTheory.Ar.SncBase.Mdl.mem_iff_mul_pow}
\leanref{Development/AutomorphicSpectrum/Proofs/GeometryProofs/UnitAtOfMulRight.lean}{GlobalTheory.CompanionTheory.Ar.SncBase.Mdl.isLocAt_invAdj}
\leanref{Development/AutomorphicSpectrum/Proofs/GeometryProofs/ExistsSncModelOfPatches.lean}{GlobalTheory.CompanionTheory.Ar.SncBase.Mdl.patch_invAdj}
\leanref{Development/AutomorphicSpectrum/Proofs/GeometryProofs/SmoothAdjoinInverseStatement.lean}{GlobalTheory.T.smooth_adjoin_inverse}
\leanref{Development/AutomorphicSpectrum/Proofs/GeometryProofs/KaehlerBasisAdjoinInverseStatement.lean}{GlobalTheory.T.kaehler_basis_adjoin_inverse}
\begin{proof}
(1) The elements $s$ of $\Omega$ with $sg^k\in A$ for some $k$ form a $\mathbf Q$-subalgebra that contains $A$ and $g^{-1}$, so it
contains $A[1/g]$. As $g\cdot g^{-1}=1$, $g\notin\mathfrak q$, and $\mathfrak q\cap A$, the preimage of $\mathfrak q$, is a prime of
$A$. Let $t=sg^k\in A$. If $s\in\mathfrak q$, then $t\in\mathfrak q$; if $t\in\mathfrak q$, then $s\in\mathfrak q$ or $g\in\mathfrak q$,
so $s\in\mathfrak q$. Let $a\in\Omega$. If $sa\in A$ for some $s\in A\setminus\mathfrak q$, then $s\in A[1/g]\setminus\mathfrak q$ and
$sa\in A[1/g]$. Conversely, let $s\in A[1/g]\setminus\mathfrak q$ with $sa\in A[1/g]$, and choose $k$ and $l$ with $sg^k\in A$ and
$sag^l\in A$. Then $t=sg^{k+l}$ lies in $A$ and outside $\mathfrak q$, and $ta=(sag^l)g^k\in A$. So
$A[1/g]_{\mathfrak q}=A_{\mathfrak q\cap A}$.

(2) The algebra $A[1/g]$ is generated by $g^{-1}$ and the generators of $A$, so it is finitely generated. As $A$ is a domain, by (1)
$A[1/g]$ is the localization of $A$ at $g$, which is $A[T]/(gT-1)$, a standard smooth $A$-algebra; so it is smooth over $A$
\cite[Tag 00T7]{Sta}, and over $\mathbf Q$ \cite[Tag 00TD]{Sta}. The localization $A\to A[1/g]$ is formally \'etale
\cite[Tag 04EG]{Sta}, so the map $A[1/g]\otimes_A\Omega_{A/\mathbf Q}\to\Omega_{A[1/g]/\mathbf Q}$, $s\otimes da\mapsto s\,da$, is an
isomorphism \cite[Tag 0H93]{Sta}. It carries the basis $1\otimes dz_1,\dots,1\otimes dz_e$ of the source to
$dz_1,\dots,dz_e$, which therefore form a basis of $\Omega_{A[1/g]/\mathbf Q}$. Let $\mathfrak q$ be a prime of
$A[1/g]$. As $A$ is a normal crossings patch, $A_{\mathfrak q\cap A}$ lies in $\mathcal P$, and it lies in $\mathcal B$ if and only if
$\mathfrak q\cap A$ contains one of $z_1,\dots,z_n$, that is, as the $z_i$ lie in $A$, if and only if $\mathfrak q$ contains one of
them. By (1), $A[1/g]_{\mathfrak q}=A_{\mathfrak q\cap A}$.
\end{proof}
% Lean: (1) is exists_mul_pow_mem, not_mem_of_isPrime, primeUnder (the preimage), mem_iff_mul_pow and isLocAt_invAdj. (2) is
% patch_invAdj: hA.fg with g^-1 added, T.smooth_adjoin_inverse (smooth_invAdj: isLocalizationAway_invAdj,
% Algebra.Smooth.of_isLocalization_Away and Algebra.Smooth.comp), T.kaehler_basis_adjoin_inverse (differentials_basis_invAdj:
% FormallyEtale.of_isLocalization, Stacks 04EG, and SncRes.coords_of_formallyEtale through
% KaehlerDifferential.tensorKaehlerEquivOfFormallyEtale, Stacks 0H93), and hA.pts at primeUnder with isLocAt_invAdj.

\begin{lemma}\label{cv-origin-reference-base-chart}
In Situation~\ref{cv-origin-reference-base-situation}, let $c$ be an index with $x_c\neq0$.
\begin{enumerate}
\item Let $P$ be a point over no prime in the chart $c$ of $x$. Then $A_c\subset\mathcal O_P$, the elements of $A_c$ that are not
  units at $P$ form a prime $I$ of $A_c$, and $(A_c)_I=\mathcal O_P$.
\item For every prime $I$ of $A_c$ there is a point $P$ over no prime in the chart $c$ at which the elements of $A_c$ that are units
  are those outside $I$, and $(A_c)_I=\mathcal O_P$.
\end{enumerate}
\end{lemma}
\leanref{Development/AutomorphicSpectrum/Proofs/GeometryProofs/ComponentGivesChartPrimeStatement.lean}{GlobalTheory.T.component_gives_chart_prime}
\leanref{Development/AutomorphicSpectrum/Proofs/GeometryProofs/ChartPrimeGivesComponentStatement.lean}{GlobalTheory.T.chart_prime_gives_component}
\begin{proof}
Let $A^{\mathbf Z}_c=\mathbf Z[x_j/x_c]\subset A_c$. For a point $P$ that specializes $\eta$, the ring $\mathcal O_P$ is isomorphic to
the local ring of $X$ at $P$, by condition (1) of Definition~\ref{cv-origin-segre-coordinates}; so it is a local ring, whose units
are the elements that are units at $P$. Lemma~\ref{cv-origin-snc-charts} applies to $x$ as a point with one block.

(1) The ratios $x_j/x_c$ are regular at $P$. A rational number is $m/d$ with integers $m$ and $d\ge1$, and $d$ is a product of
primes, which are units at $P$; so $A_c\subset\mathcal O_P$. The elements of $A_c$ that are not units at $P$ are those that map into
the maximal ideal of $\mathcal O_P$, so they form a prime $I$. If $a\in\mathcal O_P$, Lemma~\ref{cv-origin-snc-charts}(1) gives
$s\in A^{\mathbf Z}_c$, a unit at $P$, with $sa\in A^{\mathbf Z}_c$; so $s\in A_c\setminus I$ and $sa\in A_c$. Conversely, if
$s\in A_c\setminus I$ and $sa\in A_c$, then $a=s^{-1}(sa)\in\mathcal O_P$. So $(A_c)_I=\mathcal O_P$.

(2) The intersection $I\cap A^{\mathbf Z}_c$ is a prime of $A^{\mathbf Z}_c$, and Lemma~\ref{cv-origin-snc-charts}(2) gives a point
$P$ in the chart $c$ at which the elements of $A^{\mathbf Z}_c$ that are units are those outside $I$. A prime number $r$ is a unit of
$A_c$, so $r\notin I$, and $r$ is a unit at $P$; so $P$ lies over no prime, and $A_c\subset\mathcal O_P$ as in (1). Let $s\in A_c$.
Clearing the denominators of its coefficients gives an integer $m\ge1$ with $ms\in A^{\mathbf Z}_c$. As $m$ is a unit of $A_c$,
$ms\notin I$ if and only if $s\notin I$; as $m$ is a unit at $P$, $ms$ is a unit at $P$ if and only if $s$ is. So the elements of $A_c$
that are units at $P$ are those outside $I$, and $(A_c)_I=\mathcal O_P$ as in (1).
\end{proof}
% Lean: (1) is Mdl.exists_prime_chartQ (exists_prime_of_chart: the preimage of the maximal ideal of regAt p, which is local by
% regAt_isLocalRing_isFractionRing_stalk; chartQ_le with Char0.algebraMap_mem; regAt_iff_chart_localization, which is Lemma
% cv-origin-snc-charts(1)). (2) is Mdl.chartPrime_point = chart_prime_component_point (prime_chart_existsUnique_point, which is
% Lemma cv-origin-snc-charts(2), for the preimage in chartZ; char_zero_of_chart_inclusion, exists_natCast_mul_mem,
% notMem_iff_unitAt and isLocAt_of_unit).

\begin{lemma}\label{cv-origin-reference-base-cover}
In Situation~\ref{cv-origin-reference-base-situation}, suppose that at every point $P$ over no prime there are a normal crossings
patch $A$ of dimension $e$ for $(\mathcal P,\mathcal B)$ (Definition~\ref{cv-origin-reference-base-patch}) and a prime $\mathfrak q$ of
$A$ with $A_{\mathfrak q}=\mathcal O_P$. Let $P$ be a point over no prime in the chart $c$ of $x$. Then there is $h\in A_c$, a unit at
$P$, such that $A_c[1/h]$ is a normal crossings patch of dimension $e$ for $(\mathcal P,\mathcal B)$.
\end{lemma}
\leanref{Development/AutomorphicSpectrum/Proofs/GeometryProofs/ExistsSncModelOfPatches.lean}{GlobalTheory.CompanionTheory.Ar.SncBase.Mdl.exists_patch}
\begin{proof}
Let $A$ and $\mathfrak q$ be as in the hypothesis, for $P$. By Lemma~\ref{cv-origin-reference-base-chart}(1), the elements of $A_c$
that are not units at $P$ form a prime $I$ with $(A_c)_I=\mathcal O_P$. Both $A$ and $A_c$ are finitely generated $\mathbf Q$-subalgebras
of $\Omega$, so Lemma~\ref{cv-origin-reference-base-localization} gives $g\in A\setminus\mathfrak q$ and $h\in A_c\setminus I$ with
$A[1/g]=A_c[1/h]$. The element $g$ is nonzero, so by Lemma~\ref{cv-origin-reference-base-shrink}(2) $A[1/g]$ is a normal crossings
patch of dimension $e$ for $(\mathcal P,\mathcal B)$. As $h\notin I$, $h$ is a unit at $P$.
\end{proof}
% Lean: exists_patch: hP at p, T.component_gives_chart_prime, T.localizations_equal_basic_opens with chartQ_fg, and
% patch_invAdj.

\begin{lemma}\label{cv-origin-reference-base-points}
In Situation~\ref{cv-origin-reference-base-situation}, let $c$ be an index with $x_c\neq0$ and $h\in A_c$ nonzero.
\begin{enumerate}
\item For every prime $\mathfrak q$ of $A_c[1/h]$ there is a point $P$ over no prime in the chart $c$, at which $h$ is a unit, such
  that the elements of $A_c[1/h]$ that are units at $P$ are those outside $\mathfrak q$, and $A_c[1/h]_{\mathfrak q}=\mathcal O_P$.
\item Let $P$ be a point over no prime in the chart $c$ at which $h$ is a unit. Then $A_c[1/h]\subset\mathcal O_P$, and the elements
  of $A_c[1/h]$ that are not units at $P$ form a prime $\mathfrak q$ with $A_c[1/h]_{\mathfrak q}=\mathcal O_P$.
\item Let $\mathfrak q$ be a prime of $A_c[1/h]$ and $P$ a point over no prime with $A_c[1/h]_{\mathfrak q}=\mathcal O_P$, and suppose
  that $A_c[1/h]$ is a normal crossings patch of dimension $e$ for $(\mathcal P,\mathcal B)$ with elements $z_1,\dots,z_e$ and an
  integer $n$. Then $P\in B$ if and only if $\mathfrak q$ contains one of $z_1,\dots,z_n$.
\end{enumerate}
\end{lemma}
\leanref{Development/AutomorphicSpectrum/Proofs/GeometryProofs/ExistsSncModelOfPatches.lean}{GlobalTheory.CompanionTheory.Ar.SncBase.Mdl.primePoint}
\leanref{Development/AutomorphicSpectrum/Proofs/GeometryProofs/ExistsSncModelOfPatches.lean}{GlobalTheory.CompanionTheory.Ar.SncBase.Mdl.pointPrime}
\leanref{Development/AutomorphicSpectrum/Proofs/GeometryProofs/ExistsSncModelOfPatches.lean}{GlobalTheory.CompanionTheory.Ar.SncBase.Mdl.patchClause}
\leanref{Development/AutomorphicSpectrum/Proofs/GeometryProofs/RegularRingDeterminesComponentStatement.lean}{GlobalTheory.T.regular_ring_determines_component}
\begin{proof}
(1) By Lemma~\ref{cv-origin-reference-base-shrink}(1), $h\notin\mathfrak q$, and $I=\mathfrak q\cap A_c$ is a prime of $A_c$.
Lemma~\ref{cv-origin-reference-base-chart}(2) gives a point $P$ over no prime in the chart $c$ at which the elements of $A_c$ that are
units are those outside $I$, with $(A_c)_I=\mathcal O_P$. As $h\notin I$, $h$ is a unit at $P$; as $A_c\subset\mathcal O_P$ by
Lemma~\ref{cv-origin-reference-base-chart}(1), $A_c[1/h]\subset\mathcal O_P$. Let $s\in A_c[1/h]$, and $k$ with $sh^k\in A_c$
(Lemma~\ref{cv-origin-reference-base-shrink}(1)). Then $s\notin\mathfrak q$ if and only if $sh^k\notin I$, if and only if $sh^k$ is a
unit at $P$, if and only if $s$ is, as $h$ is a unit at $P$. By Lemma~\ref{cv-origin-reference-base-shrink}(1),
$A_c[1/h]_{\mathfrak q}=(A_c)_I=\mathcal O_P$.

(2) By Lemma~\ref{cv-origin-reference-base-chart}(1), $A_c\subset\mathcal O_P$, and $h^{-1}\in\mathcal O_P$; so
$A_c[1/h]\subset\mathcal O_P$. As $\mathcal O_P$ is a local ring whose units are the units at $P$ (proof of
Lemma~\ref{cv-origin-reference-base-chart}), the elements of $A_c[1/h]$ that are not units at $P$ form a prime $\mathfrak q$, the
preimage of the maximal ideal. Its intersection with $A_c$ is the prime $I$ of Lemma~\ref{cv-origin-reference-base-chart}(1), with
$(A_c)_I=\mathcal O_P$, and by Lemma~\ref{cv-origin-reference-base-shrink}(1), $A_c[1/h]_{\mathfrak q}=(A_c)_I=\mathcal O_P$.

(3) Two points $P$ and $P'$ that specialize $\eta$ with $\mathcal O_P=\mathcal O_{P'}$ are equal: $P$ lies in a chart $c'$ of $x$, and
so does $P'$, as the ratios $x_j/x_{c'}$ lie in $\mathcal O_P=\mathcal O_{P'}$; the same elements of $\mathbf Z[x_j/x_{c'}]$ are units at
$P$ and at $P'$, so $P=P'$ by Lemma~\ref{cv-origin-snc-charts}(2). As $A_c[1/h]$ is a normal crossings patch,
$A_c[1/h]_{\mathfrak q}=\mathcal O_P$ lies in $\mathcal B$ if and only if $\mathfrak q$ contains one of $z_1,\dots,z_n$. By the first
sentence, $\mathcal O_P\in\mathcal B$ if and only if $P\in B$.
\end{proof}
% Lean: (1) is primePoint (T.chart_prime_gives_component for primeUnder, not_mem_primeUnder, invAdj_le, exists_mul_pow_mem,
% mem_iff_mul_pow, unitAt_of_mul_left and unitAt_mul, isLocAt_invAdj). (2) is pointPrime (the preimage of the maximal ideal of
% regAt p, T.component_gives_chart_prime, and isLocAt_invAdj). (3) is patchClause (hS.pts, and
% T.regular_ring_determines_component = eq_of_regAt_eq_of_coordTie through prime_chart_existsUnique_point).

\begin{lemma}\label{cv-origin-reference-base-separate}
In Situation~\ref{cv-origin-reference-base-situation}, let $c$ be an index, $h\in A_c$ nonzero, $z_1,\dots,z_e\in A_c[1/h]$ and
$n\ge0$, and suppose that a point $Q$ over no prime in the chart $c$ at which $h$ is a unit lies in $B$ if and only if one of the
$z_i$ with $i\le n$ is not a unit at $Q$. Let $P\notin B$ be a point over no prime in the chart $c$ at which $h$ is a unit. Then there
is $a\in A_c$ that is a unit at $P$ and is not a unit at any point of $B$ in the chart $c$.
\end{lemma}
\leanref{Development/AutomorphicSpectrum/Proofs/GeometryProofs/SncPatchBoundaryUnitsStatement.lean}{GlobalTheory.T.snc_patch_boundary_units}
\begin{proof}
Let $\zeta$ be the product of the $z_i$ with $i\le n$. By Lemma~\ref{cv-origin-reference-base-shrink}(1) there is $m$ with
$\zeta h^m\in A_c$; put $a=\zeta h^{m+1}\in A_c$. As $P\notin B$, each $z_i$ with $i\le n$ is a unit at $P$, and so is $h$; so $a$ is a
unit at $P$. Let $Q\in B$ lie in the chart $c$, and suppose that $a$ is a unit at $Q$. The point $Q$ lies over no prime, so
$A_c\subset\mathcal O_Q$, a local ring whose units are the units at $Q$ (Lemma~\ref{cv-origin-reference-base-chart}(1) and its
proof). In a local ring, a product is a unit only if each factor is. If $h$ is a unit at $Q$, then $A_c[1/h]\subset\mathcal O_Q$, so
$\zeta$ is a unit at $Q$, and then each $z_i$ with $i\le n$; this contradicts $Q\in B$. Otherwise $a=(\zeta h^m)h$, with both factors
in $A_c\subset\mathcal O_Q$, shows that $h$ is a unit at $Q$, a contradiction.
\end{proof}
% Lean: snc_patch_boundary_units = Mdl.exists_chartQ_unit_sep: zeta the product over the i < n, exists_mul_pow_mem for m,
% unitAt_finsetProd for P, and for Q the two cases on unitAt Q h (chartQ_le, invAdj_le, unitAt_of_mul_left,
% unitAt_of_finsetProd, unitAt_of_mul_right).

\begin{lemma}\label{cv-origin-reference-base-form}
In Situation~\ref{cv-origin-reference-base-situation}, let $P$ be a point over no prime in the chart $c$ of $x$, and $a\in A_c$ a unit
at $P$ that is not a unit at any point of $B$ in the chart $c$. Then there is a boundary form $g$ such that $g(x/x_c)$ is a unit at
$P$.
\end{lemma}
\leanref{Development/AutomorphicSpectrum/Proofs/GeometryProofs/BoundaryGeneratorUnitStatement.lean}{GlobalTheory.T.boundary_generator_unit}
\begin{proof}
\emph{A form.} There are integers $m\ge1$ and $d\ge0$ and a homogeneous $H\in\mathbf Z[X_0,\dots,X_N]$ of degree $d$ with
$H(x/x_c)=ma$. This holds for the generators of $A_c$: for $x_j/x_c$ take $m=1$ and $H=X_j$, and for a rational number $u/v$ with
integers $u$ and $v\ge1$ take $m=v$ and the constant $H=u$. If it holds for $a_1$ and $a_2$, with $m_1,d_1,H_1$ and $m_2,d_2,H_2$,
then it holds for $a_1a_2$ with $m_1m_2$, $d_1+d_2$ and $H_1H_2$, and for $a_1+a_2$ with $m_1m_2$, $d_1+d_2$ and
$m_2H_1X_c^{d_2}+m_1H_2X_c^{d_1}$, as the value of $X_c$ at $x/x_c$ is $1$.

\emph{A boundary form.} Put $g=X_cH$, homogeneous of degree $d+1$. Then $g(x/x_c)=H(x/x_c)=ma$, which is a unit at $P$, as $P$ lies
over no prime. Let $c'$ be an index and $Q\in B$ a point in the chart $c'$, and suppose that $g(x/x_{c'})=(x_c/x_{c'})H(x/x_{c'})$ is
a unit at $Q$. Both factors are regular at $Q$, and $\mathcal O_Q$ is local (proof of Lemma~\ref{cv-origin-reference-base-chart}), so
$x_c/x_{c'}$ is a unit at $Q$; then every $x_j/x_c=(x_j/x_{c'})(x_c/x_{c'})^{-1}$ is regular at $Q$, and $Q$ lies in the chart $c$. As
$g$ is homogeneous of degree $d+1$, $g(x/x_{c'})=(x_c/x_{c'})^{d+1}g(x/x_c)=(x_c/x_{c'})^{d+1}ma$. The point $Q$ lies over no prime,
so $a\in A_c\subset\mathcal O_Q$ (Lemma~\ref{cv-origin-reference-base-chart}(1)) and $m$ is a unit at $Q$; so $ma$, and then $a$, is a unit at $Q$, which contradicts the hypothesis on
$a$. So $g$ is a boundary form.
\end{proof}
% Lean: boundary_generator_unit = Mdl.Sep.exists_bdGens_unitAt: exists_form (induction on the adjoin), exists_form_sep with
% G = X c * H (isHomogeneous_X_mul), not_unitAt_form (ratEval_X_mul, unitAt_of_mul_left, inChart_of_unitAt,
% ratEval_eq_pow_mul, Char0.unitAt_natCast, unitAt_of_mul_right).

\begin{lemma}\label{cv-origin-reference-base-zeros}
In Situation~\ref{cv-origin-reference-base-situation}, suppose that for every index $c$ and every point $P\notin B$ over no prime in
the chart $c$ there is a boundary form $g$ with $g(x/x_c)$ a unit at $P$. Let $z$ be a complex point of $X$ with coordinates $w$
for $x$, as a point with one block (Definition~\ref{cv-origin-reference}(4)). Then $w$ lies in the zero set of $J_B$ in $\hat C$
(Definition~\ref{lit-multiprojective}) if and only if the centre of $z$, the image of its closed point, lies in $B$.
\end{lemma}
\leanref{Development/AutomorphicSpectrum/Proofs/GeometryProofs/BoundaryPointsCoordsStatement.lean}{GlobalTheory.T.boundary_points_coords}
\begin{proof}
Let $P$ be the centre of $z$. As $w\in\hat C$, some $w_c$ is nonzero; then $x_c\neq0$, and the ratios $x_j/x_c$ are regular at $P$
with values $w_j/w_c$. In particular $P$ specializes $\eta$ and lies in the chart $c$. Let $\varphi:\mathcal O_P\to\mathbf C$ be the
composite of the isomorphism of $\mathcal O_P$ with the local ring of $X$ at $P$ and the local homomorphism of $z$; it sends an
element regular at $P$ to its value, and an element $s$ of $\mathcal O_P$ is a unit at $P$ if and only if $\varphi(s)\neq0$. In
particular every prime number is a unit at $P$, and $P$ lies over no prime. For a homogeneous $g$ of degree $d$,
\[
g(w)=w_c^dg(w_0/w_c,\dots,w_N/w_c)=w_c^d\varphi\big(g(x/x_c)\big),
\]
so $g(w)=0$ if and only if $g(x/x_c)$ is not a unit at $P$. If $w$ lies in the zero set of $J_B$ and $P\notin B$, the hypothesis
gives a boundary form $g$ with $g(x/x_c)$ a unit at $P$, and then $g(w)\neq0$, a contradiction. Conversely, if $P\in B$, then every
boundary form $g$ has $g(x/x_c)$ not a unit at $P$, as $P$ lies in the chart $c$, so $g(w)=0$; the boundary forms generate $J_B$, so
every element of $J_B$ vanishes at $w$, which lies in $\hat C$.
\end{proof}
% Lean: boundary_points_coords = Mdl.mem_mcPts_span_bdGens_iff: ValHom.exists_valAt_hom for phi (with isUnit_iff_unitAt),
% char_zero_centre_complex_point, eval_map_mul_of_isHomogeneous and exists_ratEval_val for g(w), and span_induction on the
% ideal. The Lean assumes only condition (1) of the coordinates (IsCptGen xi), for any field Omega.

\begin{lemma}\label{cv-origin-reference-base-model}
In Situation~\ref{cv-origin-reference-base-situation}, let $J$ be an ideal of $\mathbf Z[X_0,\dots,X_N]$, and suppose that for every
index $c$ and every point $P$ over no prime in the chart $c$ of $x$:
\begin{enumerate}
\item there is $h\in A_c$, a unit at $P$, such that $A_c[1/h]$ is a normal crossings patch of dimension $e$ for
  $(\mathcal P,\mathcal B)$ (Definition~\ref{cv-origin-reference-base-patch});
\item $P\in B$ if and only if $g(x/x_c)$ is not a unit at $P$ for any $g\in J$.
\end{enumerate}
Then $J$ satisfies condition (3) of Definition~\ref{cv-origin-reference} for $x$, as a point with one block, with $e$ in place of $D$.
\end{lemma}
\leanref{Development/AutomorphicSpectrum/Proofs/GeometryProofs/SncModelFromLocalPatchesStatement.lean}{GlobalTheory.T.snc_model_from_local_patches}
\begin{proof}
For one block, a chart of $x$ is an index $c$ with $x_c\neq0$, the algebra of condition (3) at $c$ is $A_c$, and the values there of
$g\in J$ are the $g(x/x_c)$. Let $c$ be such an index, and $\mathcal T$ the set of the nonzero $t\in A_c$ for which $A_c[1/t]$ is a
normal crossings patch of dimension $e$ for $(\mathcal P,\mathcal B)$.

\emph{A partition of unity.} The set $\mathcal T$ generates the unit ideal of $A_c$. Otherwise it lies in a maximal ideal $I$;
Lemma~\ref{cv-origin-reference-base-chart}(2) gives a point $P$ over no prime in the chart $c$ at which the elements of $A_c$ that are
units are those outside $I$, and (1) gives $h\in A_c$, a unit at $P$, with $A_c[1/h]$ a normal crossings patch; then
$h\in\mathcal T\subset I$, and $h$ is not a unit at $P$, a contradiction. So there are $f_1,\dots,f_m\in\mathcal T$ and
$a_1,\dots,a_m\in A_c$ with $\sum_la_lf_l=1$.

\emph{The boundary.} Let $1\le l\le m$. As a normal crossings patch, $A_c[1/f_l]$ is smooth over $\mathbf Q$ and has elements
$z_1,\dots,z_e$, whose differentials form a basis of its module of differentials over $\mathbf Q$, and an integer $n\le e$. Let
$\mathfrak q$ be a prime of $A_c[1/f_l]$. Lemma~\ref{cv-origin-reference-base-points}(1) gives a point $P$ over no prime in the chart
$c$ at which the elements of $A_c[1/f_l]$ that are units are those outside $\mathfrak q$, with $A_c[1/f_l]_{\mathfrak q}=\mathcal O_P$;
by Lemma~\ref{cv-origin-reference-base-points}(3), $P\in B$ if and only if $\mathfrak q$ contains one of $z_1,\dots,z_n$. By (2),
$P\in B$ if and only if no $g(x/x_c)$ with $g\in J$ is a unit at $P$, that is, if and only if $\mathfrak q$ contains all the
$g(x/x_c)$ with $g\in J$. This is condition (3) at the chart $c$.
\end{proof}
% Lean: snc_model_from_local_patches = Mdl.isSncModel_of_chart_patches = chart_cover_core: the set T, Ideal.exists_le_maximal,
% exists_sum_eq_one_of_span_eq_top, and the boundary clause from h3 and h4. Its hypotheses h2 and h3 are, in the book, Lemma
% cv-origin-reference-base-chart(2) and Lemma cv-origin-reference-base-points(1) and (3), which the Lean supplies in
% exists_isSncModel_of_patches (T.chart_prime_gives_component, primePoint and patchClause).

\begin{lemma}\label{cv-origin-reference-base-ideal}
In Situation~\ref{cv-origin-reference-base-situation}, suppose that at every point $P$ over no prime there are a normal crossings
patch $A$ of dimension $e$ for $(\mathcal P,\mathcal B)$ (Definition~\ref{cv-origin-reference-base-patch}) and a prime $\mathfrak q$ of
$A$ with $A_{\mathfrak q}=\mathcal O_P$. Then the ideal $J_B$ generated by the boundary forms satisfies condition (3) of
Definition~\ref{cv-origin-reference} for $x$, as a point with one block, with $e$ in place of $D$, and a complex point of $X$ with
coordinates $w$ for $x$ has its centre in $B$ if and only if $w$ lies in the zero set of $J_B$ in $\hat C$.
\end{lemma}
\leanref{Development/AutomorphicSpectrum/Proofs/GeometryProofs/ExistsSncModelOfPatchesStatement.lean}{GlobalTheory.T.exists_snc_model_of_patches}
\leanref{Development/AutomorphicSpectrum/Proofs/GeometryProofs/ExistsSncModelOfPatches.lean}{GlobalTheory.CompanionTheory.Ar.SncBase.exists_isSncModel_of_patches}
\begin{proof}
\emph{Patches on the charts.} By Lemma~\ref{cv-origin-reference-base-cover}, for every index $c$ and every point $P$ over no prime in
the chart $c$ there is $h\in A_c$, a unit at $P$, such that $A_c[1/h]$ is a normal crossings patch of dimension $e$ for
$(\mathcal P,\mathcal B)$.

\emph{Separation.} Let $P\notin B$ be a point over no prime in the chart $c$, and let $h$ be as in the previous paragraph, with
elements $z_1,\dots,z_e$ and an integer $n$. Let $Q$ be a point over no prime in the chart $c$ at which $h$ is a unit. By
Lemma~\ref{cv-origin-reference-base-points}(2), the elements of $A_c[1/h]$ that are not units at $Q$ form a prime $\mathfrak q$ with
$A_c[1/h]_{\mathfrak q}=\mathcal O_Q$, and by Lemma~\ref{cv-origin-reference-base-points}(3), $Q\in B$ if and only if $\mathfrak q$
contains one of $z_1,\dots,z_n$, that is, one of them is not a unit at $Q$. So Lemma~\ref{cv-origin-reference-base-separate} gives
$a\in A_c$, a unit at $P$, that is not a unit at any point of $B$ in the chart $c$, and Lemma~\ref{cv-origin-reference-base-form}
gives a boundary form $g$ with $g(x/x_c)$ a unit at $P$.

\emph{The ideal.} Let $P$ be a point over no prime in the chart $c$. If $P\in B$, then $g(x/x_c)$ is not a unit at $P$ for any
$g\in J_B$: such a $g$ is $\sum_lr_lg_l$ with boundary forms $g_l$ and polynomials $r_l$, the values $r_l(x/x_c)$ and
$g_l(x/x_c)$ are regular at $P$, the $g_l(x/x_c)$ are not units at $P$, and the non-units of the local ring $\mathcal O_P$ form an
ideal. If $P\notin B$, the separation gives a boundary form $g$ with $g(x/x_c)$ a unit at $P$. So
Lemma~\ref{cv-origin-reference-base-model} applies to $J_B$, with the first paragraph as its hypothesis (1), and $J_B$ satisfies
condition (3).

\emph{Complex points.} By the separation, Lemma~\ref{cv-origin-reference-base-zeros} applies, and it gives the last assertion.
\end{proof}
% Lean: exists_isSncModel_of_patches: hcov (exists_patch), hsep (pointPrime and patchClause for hcl,
% T.snc_patch_boundary_units, T.boundary_generator_unit), the ideal Ideal.span (bdGens xi x B), T.snc_model_from_local_patches
% with T.chart_prime_gives_component, primePoint, patchClause and, for its last hypothesis, not_unitAt_of_mem_span (span
% induction with not_unitAt_add and unitAt_of_mul_right) and hsep; and T.boundary_points_coords with hsep. The Lean states
% that some ideal exists; its proof takes the ideal generated by bdGens, which the book states. The lemmas of this section from
% Situation cv-origin-reference-base-situation on are about 1,400 lines of Lean.

\begin{lemma}\label{cv-origin-reference-base}
In the setting of Definition~\ref{cv-origin-hodge-component}, let $\Omega_0$ be a field of characteristic zero and
$\xi_0:\operatorname{Spec}\Omega_0\to\mathcal X$ a morphism as in Lemma~\ref{cv-origin-reference-char}. Then there are a
nonempty finite set $K$, finite sets $I_k$, a point $y\in\Omega_0^{\sqcup_kI_k}$, a scheme $T$ with a morphism
$\xi:\operatorname{Spec}\Omega_0\to T$ of whose image every point is a specialization, covered by the affine open subschemes of
the charts of $y$ as in Definition~\ref{cv-origin-hodge-component}(1), a morphism $f:T\to\mathcal X$ with $f\circ\xi=\xi_0$,
and an ideal $J$ of $\mathbf Z[x_{k,i}]$ that satisfies condition (3) of Definition~\ref{cv-origin-reference} for $y$ with
$|\Psi|pq$ in place of $D$, such that the complex points of $T$ whose coordinates lie in the zero set of $J$ are those whose
image under $f$ lies outside the abelian locus of $\mathcal A$.
\end{lemma}
\leanref{Development/AutomorphicSpectrum/Proofs/GeometryProofs/ExistsCoordsOfGenericPointStatement.lean}{GlobalTheory.T.exists_coords_of_generic_point}
\begin{proof}
Write $F=F_c$, $\mathcal X_F=\mathcal X\times_{\mathcal O_c}F$ with its projection $\pi:\mathcal X_F\to\mathcal X$, $U=j(X)$, and
$\Lambda$ for the abelian locus of $\mathcal A$, the set of points of $\mathcal X$ over which $\mathcal A$ is abelian. By
Definition~\ref{cv-model}, $U$ is open and dense in $\mathcal X_F$ and $\pi^{-1}(\Lambda)=U$, and $U\cong X$ is smooth of relative
dimension $D=|\Psi|pq$ over $F$. Coordinates over a ring are as in Definition~\ref{cv-origin-reference-base-over}. Let $\eta_0$ be
the image of $\xi_0$.

\emph{The generic fibre.} On an affine open subscheme $\operatorname{Spec}R$ of $\mathcal X$, $\pi$ is
$\operatorname{Spec}$ of the localization of $R$ at the nonzero elements of $\mathcal O_c$, so $\pi$ identifies the local rings of
$\mathcal X_F$ with those of $\mathcal X$. Hence $\mathcal X_F$ is reduced, as $\mathcal X$ is normal, and locally noetherian; and
it is projective over $F$, as $\mathcal X$ is projective over $\mathcal O_c$, with a closed immersion into some $\mathbf P^M_F$. As
$\Omega_0$ has characteristic zero, the homomorphism $\mathcal O_c\to\Omega_0$ of $\xi_0$ is injective and extends to $F$; so
$\xi_0=\pi\circ\xi_F$ for a morphism $\xi_F:\operatorname{Spec}\Omega_0\to\mathcal X_F$ over $F$. As $\pi$ identifies the local
rings, the local ring of $\mathcal X_F$ at the image $\eta_F$ of $\xi_F$ maps bijectively onto $\Omega_0$, and those at its
specializations map injectively.

\emph{The point $\eta_F$ lies in $U$.} By Lemma~\ref{cv-origin-reference-base-projective}(1), the set $Z$ of the specializations of
$\eta_F$ is open. As $U$ is dense, it meets the nonempty open set $Z$ at a point that specializes $\eta_F$; as $U$ is open,
$\eta_F\in U$.

\emph{A resolution.} Lemma~\ref{cv-resolution-coordinates}, applied to $K=F$, to $\mathcal X_F\subset\mathbf P^M_F$, to $U$ and to
$D$, gives $r\ge0$, a closed subscheme $X'\subset\mathbf P^M_F\times\mathbf P^r_F$ and the restriction $f':X'\to\mathcal X_F$ of
the first projection, with $X'$ smooth of relative dimension $D$ over $F$, $f'$ an isomorphism over $U$ with $f'^{-1}(U)$ dense,
and $X'$ covered by finitely many affine open subschemes with coordinates $z_1,\dots,z_D$ in which $X'\setminus f'^{-1}(U)$ is
the set of points at which one of $z_1,\dots,z_k$ vanishes. Through the Segre embedding, $X'$ is projective over $F$. Its affine
open subschemes are spectra of smooth $F$-algebras, which are normal \cite[Tag 033C]{Sta}; so the local rings of $X'$ are domains.
As $f'$ is an isomorphism over $U\ni\eta_F$, $\xi_F=f'\circ\xi'$ for a morphism $\xi':\operatorname{Spec}\Omega_0\to X'$ whose
image $\eta'$ lies in $f'^{-1}(U)$, and the local ring of $X'$ at $\eta'$ is that of $\mathcal X_F$ at $\eta_F$, which maps
bijectively onto $\Omega_0$. A generization of $\eta'$ lies in the open set $f'^{-1}(U)$ and maps to a generization of $\eta_F$,
which is $\eta_F$, as the local ring there is a field; as $f'$ is injective on $f'^{-1}(U)$, $\eta'$ has no generization other
than itself. Let $p'$ specialize $\eta'$. The local ring at $p'$ is a domain, its zero ideal corresponds to the generic point of
the only irreducible component through $p'$, which is $\eta'$, and its map to the local ring at $\eta'$ is the localization at
the zero ideal, which is injective.

\emph{Coordinates.} Let $g=\pi\circ f':X'\to\mathcal X$; then $g\circ\xi'=\xi_0$. By
Lemma~\ref{cv-origin-reference-base-projective}(2), applied to a closed immersion of $\mathcal X$ into a projective space over
$\mathcal O_c$ and to one of $X'$ into a projective space over $F$, $\mathcal X$ has coordinates $u$ over $\mathcal O_c$ with respect
to $\xi_0$, and $X'$ has coordinates $x'$ over $F$ with respect to $\xi'$. Let
$p'$ be a point of $X'$ that specializes $\eta'$. By Lemma~\ref{cv-origin-segre-regular}(3), $g(p')$ specializes $\eta_0$; it lies
in a chart $c$ of $u$, and the ratios of $u$ at $c$, which are regular at $g(p')$, are regular at $p'$. So every point of $X'$
that specializes $\eta'$ lies in a chart of $u$, and Lemma~\ref{cv-origin-reference-base-product} gives $n\ge0$ such that $X'$ has
the coordinates $t=(x'_{a_1}\cdots x'_{a_n}x'_\delta u_\varepsilon)_{(a,\delta,\varepsilon)}$ over $F$, where $a=(a_1,\dots,a_n)$
runs over the $n$-tuples of indices of $x'$.

\emph{A model over $\mathcal O_c$.} The point $\eta'$ lies in a chart $c_0$ of $t$. Every element of $\Omega_0$ is regular at
$\eta'$, so, as in Lemma~\ref{cv-origin-snc-charts}(1), it is a quotient $a/s$ of elements of the $F$-algebra generated by the
ratios $t_j/t_{c_0}$; multiplying $a$ and $s$ by a common denominator in $\mathcal O_c$ of their coefficients writes it as a
quotient of two elements of the $\mathcal O_c$-algebra generated by these ratios. Lemma~\ref{cv-origin-segre-model}, with the
same proof over the image of $\mathcal O_c$ in place of the image of $\mathbf Z$, gives a scheme $T$ with a morphism
$\xi:\operatorname{Spec}\Omega_0\to T$ of whose image $\eta$ every point is a specialization, such that $T$ has the coordinates $t$
over $\mathcal O_c$. A point of $T$ lies in a chart $(a,\delta,\varepsilon)$ of $t$; then $u_\varepsilon\neq0$, and every
$u_c/u_\varepsilon=t_{(a,\delta,c)}/t_{(a,\delta,\varepsilon)}$ is regular there, so it lies in a chart of $u$. The
elements of $\mathcal O_c$ are regular at every point of $T$. So Lemma~\ref{cv-origin-coords-join-map}(2), with the same proof
over $\mathcal O_c$, gives a morphism $f:T\to\mathcal X$ with $f\circ\xi=\xi_0$.

\emph{Points over no prime.} A point $p$ of $T$ lies over no prime if and only if every nonzero $d\in\mathcal O_c$ is a unit at
$p$: if $p$ lies over no prime and $d\neq0$, the norm of $d$ is a nonzero integer $dd'$ with $d'\in\mathcal O_c$, a unit at $p$, and
$d^{-1}=d'(dd')^{-1}$ is regular at $p$. The ring $\mathcal O_c$ is generated as an abelian group by finitely many elements, so
Lemma~\ref{cv-origin-reference-base-integral} shows that $T$ has coordinates $x=(x_0,\dots,x_N)$ over $\mathbf Z$, after numbering
the family. By Lemma~\ref{cv-origin-reference-base-rings}, applied to $\mathcal O_c$, $F$, $t$, $T$ and $X'$, every local ring of
$X'$ is a local ring of $T$, and every local ring of $T$ at a point over no prime is a local ring of $X'$. Let $p\in T$ and
$p'\in X'$ have $\mathcal O_p=\mathcal O_{p'}$. Then $f(p)=g(p')$. Indeed, by Lemma~\ref{cv-origin-segre-regular}(3),
$\mathcal O_{f(p)}$ and $\mathcal O_{g(p')}$ lie in $\mathcal O_p$, and an element of either that is a unit at $p$
is a unit at $f(p)$, or at $g(p')$. If $f(p)$ lies in the chart $c_1$ of $u$ and $g(p')$ in the chart $c$, then $u_c/u_{c_1}$ is
regular at $f(p)$ and a unit at $p$, as its inverse is regular at $g(p')$; so it is a unit at $f(p)$, and $f(p)$ lies in the chart
$c$. An element of the $\mathcal O_c$-algebra generated by the ratios $u_j/u_c$ is then a unit at $f(p)$ if and only if it is a
unit at $p$, if and only if it is a unit at $g(p')$; so $f(p)=g(p')$ by Lemma~\ref{cv-origin-snc-charts}(2), over $\mathcal O_c$.
Hence $f(p)\notin\Lambda$ if and only if $f'(p')\notin\pi^{-1}(\Lambda)=U$.

\emph{The boundary.} Let $B$ be the set of the points $p$ of $T$ over no prime with $f(p)\notin\Lambda$. By the previous
paragraph, the local rings of $T$ at the points of $B$ are those of $X'$ at the points of $X'\setminus f'^{-1}(U)$, and the local
rings of $T$ at the points over no prime are those of $X'$. Distinct points of $X'$ have distinct local rings: two points with the
same local ring lie in the same chart of $x'$, where the same elements of the chart ring are units, so they are equal by
Lemma~\ref{cv-origin-snc-charts}(2), over $F$. Let $\mathcal P$ be the set of the local rings of $T$ at the points over no prime, and
$\mathcal B$ the set of those at the points of $B$; by the above, they are also the sets of the local rings of $X'$ at all its points
and at the points of $T_\partial=X'\setminus f'^{-1}(U)$. Lemma~\ref{cv-origin-reference-base-patches}, applied to $X'$, $\xi'$ and
$T_\partial$ with the cover of Lemma~\ref{cv-resolution-coordinates}(3), gives at every point $p'$ of $X'$ a normal crossings patch $A$ of
dimension $D$ for $(\mathcal P,\mathcal B)$ and a prime $\mathfrak q$ of $A$ with $A_{\mathfrak q}=\mathcal O_{p'}$. As every local ring of
$T$ at a point over no prime is one of $X'$, at every point $p$ of $T$ over no prime there are such $A$ and $\mathfrak q$ with
$A_{\mathfrak q}=\mathcal O_p$. So Lemma~\ref{cv-origin-reference-base-ideal} applies to $\Omega_0$, $T$, $\xi$, $x$, $e=D$ and $B$
(Situation~\ref{cv-origin-reference-base-situation}): the ideal $J$ of $\mathbf Z[X_0,\dots,X_N]$ generated by the boundary forms
satisfies condition (3) of Definition~\ref{cv-origin-reference} for $x$, as a point with one block, with $D$,
and the complex points of $T$ whose coordinates lie in the zero set of $J$ are those whose centre lies in $B$. The centre of a
complex point lies over no prime: the local homomorphism from its local ring to
$\mathbf C$ sends every prime to a unit. So these complex points are those whose image under $f$ lies outside $\Lambda$.

\emph{One block.} Let $K$ have one element $k$, with $I_k=\{0,\dots,N\}$ and $y=x$, so that $\mathbf Z[x_{k,i}]$ is
$\mathbf Z[X_0,\dots,X_N]$ with $x_{k,i}=X_i$. The Segre products of $y$ are the entries of $x$, so $T$ is covered by the affine open
subschemes of the charts of $y$.
% Lean: T.exists_coords_of_generic_point = SncBase.exists_sncBase, under ResKStatement: exists_model_of_res with K = Unit and
% T.coords_single_block_segre (= SncBase.Reindex.coordTie_one_block). exists_model_of_res, step by step:
% T.generic_fibre_projective_embedding (genericFibre_reduced_noetherian_closed), T.generic_point_generic_fibre
% (GenFibre.exists_point: injective_of_charZero, isIso_stalkMap_fst, isCptGen_of_isIso_stalkMap),
% T.generic_point_mem_dense_open (genPt_mem_of_dense_open, with isClopen_setOf_genPt_specializes through
% exists_isOpen_common_generization), hRes (Lemma cv-resolution-coordinates), T.closed_in_product_projective,
% T.smooth_scheme_stalks_domains, T.generic_point_lift_along_iso (exists_lift_isCptGen, isCptGen_of_maximal_of_surjective),
% T.every_component_in_some_chart, T.coords_generate_function_field (ChartFrac.chartRing_frac), T.exists_proper_integral_model
% with o = O_F, T.dominates_of_coords_regular with o = O_F, T.char_zero_iff_nonzero_units (CharTie, Ideal.absNorm),
% T.regular_rings_equal_images_equal (image_eq_of_regAt_eq, through prime_chart_existsUnique_point over O_F), hgen, hBd and
% hPts, T.coord_over_int_of_finite_order (Lemma cv-origin-reference-base-integral), T.coord_charts_match_regular (Lemma
% cv-origin-reference-base-rings), T.projective_over_coords_exist (Lemma cv-origin-reference-base-projective), and
% T.coords_product_refinement (Lemma cv-origin-reference-base-product, whose proof takes n = 0), and
% T.snc_patch_at_each_component (Lemma cv-origin-reference-base-patches) with T.regular_ring_determines_component and hPts,
% hBd, T.exists_snc_model_of_patches (Lemma cv-origin-reference-base-ideal, with B = {p | Char0 p and f p not in Ab}, whose
% hypotheses come from R2, hPts, hBd and hpatch), and T.char_zero_centre_of_complex.
\end{proof}

\begin{lemma}\label{cv-origin-reference-rank}
In the setting of Definition~\ref{cv-origin-hodge-component}, let $\Omega$ be a field of characteristic zero, $T$ a scheme
with a morphism $\xi:\operatorname{Spec}\Omega\to T$, $f:T\to\mathcal X$ a morphism with $f\circ\xi=j\circ y_0$ for a point
$y_0$ of $X$ with values in $\Omega$, and $F_c\to\Omega$ an embedding through which $f\circ\xi$ lies over $\mathcal O_c$. Let
$G$ be an $\mathcal O_E$-semi-abelian scheme over $T$ with a morphism $G\to\mathcal A$ over $f$, and $\mathrm{Lie}$ its Lie
algebra at $\xi$ with its action of $\mathcal O_E$. Then $\mathrm{Lie}$ is finite-dimensional, and for every label $s$ the
subspace on which each $b\in\mathcal O_E$ acts by $s(b)$, through $F_c\to\Omega$, has dimension $r_s$.
\end{lemma}
\leanref{Development/AutomorphicSpectrum/Proofs/GeometryProofs/EigenLatticeConstantRankStatement.lean}{GlobalTheory.T.eigen_lattice_constant_rank}
\begin{proof}
Let $\sigma:F_c\to\Omega$ be the given embedding and $\rho$ the action of $\mathcal O_E$ on $\mathrm{Lie}$. For a homomorphism
$\chi:\mathcal O_E\to K$ to a field and a $K$-vector space $V$ with an action of $\mathcal O_E$, let $V^\chi$ be the subspace on which
each $b$ acts by $\chi(b)$. The group scheme $G$ is smooth over $T$, being semi-abelian, and so is the universal abelian scheme $A$
over $X$.

\emph{Finiteness.} By Lemma~\ref{cv-origin-reference-lie-bundle-base}, $\mathrm{Lie}$ is a finite projective $\Omega$-module, that is,
a finite-dimensional vector space.

\emph{Transport to $X$.} The morphism $G\to\mathcal A$ over $f$ identifies $G$ with $f^*\mathcal A$, and $j^*\mathcal A$ is $A$
(Definition~\ref{cv-model}), through a morphism $A\to\mathcal A$ over $j$ compatible with the actions. As $f\circ\xi=j\circ y_0$,
composition with these morphisms identifies the tangent vectors of $G$ at $\xi$ and those of $A$ at $y_0$ with the tangent vectors
of $\mathcal A$ at $f\circ\xi$, compatibly with sums, scalars and the actions of $\mathcal O_E$. So $\mathrm{Lie}$, with $\rho$, is
isomorphic to the Lie algebra of $A$ at $y_0$ with its action.

\emph{An algebraic closure.} Let $\bar\Omega$ be an algebraic closure of $\Omega$, and $x=y_0\circ\operatorname{Spec}(\bar\Omega\supset\Omega)$.
By Lemma~\ref{cv-origin-reference-lie-bundle-base}, composition with $\operatorname{Spec}(\bar\Omega[\varepsilon]\supset\Omega[\varepsilon])$
induces an isomorphism from $\bar\Omega\otimes_\Omega\mathrm{Lie}$ onto the Lie algebra $L_x$ of $A$ at $x$, which is that of the fibre
$A_x$; it commutes with the actions of $\mathcal O_E$, which act by composition on the other side. Let $s$ be a label, $\chi=\sigma\circ s$,
and $\bar\chi$ its composite with $\Omega\subset\bar\Omega$. For a $\mathbf Z$-basis $b_1,\dots,b_n$ of $\mathcal O_E$, $\mathrm{Lie}^\chi$
is the kernel of the linear map $e\mapsto(\rho(b_i)e-\chi(b_i)e)_i$ from $\mathrm{Lie}$ to $\mathrm{Lie}^n$, and the dimension of the
kernel of a linear map does not change under extension of the field. So $\dim_\Omega\mathrm{Lie}^\chi=\dim_{\bar\Omega}L_x^{\bar\chi}$.

\emph{The signature.} The point $y_0$ of $X$ gives an embedding $F_c\to\Omega$; it agrees with $\sigma$ on $\mathcal O_c$, as
$j\circ y_0=f\circ\xi$ lies over $\mathcal O_c$ through $\sigma$, and so it is $\sigma$, as $F_c$ is the field of fractions of
$\mathcal O_c$. The moduli scheme $X$ is $X_l\otimes_{K_l}F_c$, and $A$ is the pullback along the projection of the universal
object over $X_l$ (Definition~\ref{cv-toroidal-moduli-object}(5)). So condition (a) of Definition~\ref{cv-toroidal-family}(1) holds at the point $x$
of $X_l$, with values in the algebraically closed field $\bar\Omega$ of characteristic zero, where $\beta:K_l\to\bar\Omega$ is the
restriction of $F_c\to\Omega\subset\bar\Omega$: $L_x^\tau$ has dimension $p$ if $\tau(\sqrt{-d})=\beta(\theta)$, and $q$ otherwise. Apply
this to $\tau=\bar\chi$. If $s$ is a label $\tau_0\in\Psi$, then $s(\sqrt{-d})=\theta$, so $\bar\chi(\sqrt{-d})=\beta(\theta)$ and
$\dim\mathrm{Lie}^\chi=p=r_s$. If $s=\bar\tau_0$, then $s(\sqrt{-d})=\bar\theta=-\theta$, as $\theta^2=-d$ is negative; so
$\bar\chi(\sqrt{-d})=-\beta(\theta)$, which differs from $\beta(\theta)$, as $\beta(\theta)^2=-d$ is nonzero in a field of
characteristic zero. Then $\dim\mathrm{Lie}^\chi=q=r_s$.
% Lean: T.eigen_lattice_constant_rank = EigRank.eig_finrank_label: T.lie_model_finite_projective (finiteness);
% EigRank.lieModel_transport (LiePull.lieModel_pullback along m.cart and Y.m.cart, with m.act and Y.m.act) to the point y0 of
% Xo.XF c and the object Y.oF; T.lie_model_base_change_map (with Smooth Y.oF.t.G.hom) to AlgebraicClosure Omega, with
% lieModel_exists_unique_act for rho'; EigRank.finrank_eigSp_baseChange (Chart.finrank_eig_bc, the eigenspace in coordinates);
% EigRank.structure_of_model_point (h5: the embedding of F_c is cst, by IsLocalization.ringHom_ext); the condition Y.oF.geom,
% DetCond with rho = rt0, whose LieMult is the Lean's form of the dimension (EigRank.finrank_eigSp_of_lieMult); and
% EigRank.labelEmb_rt with rt0_conj_neg_ne_zero for the values at sqrt(-d), and hne (d is prime) for beta(theta) nonzero.
\end{proof}

\begin{lemma}\label{cv-origin-reference-coords}
In the setting of Definition~\ref{cv-origin-hodge-component}, let $\Omega_0$ and $\xi_0$ be as in
Lemma~\ref{cv-origin-reference-char}. Then there are finite sets $K$ and $I_k$, a point
$y\in\Omega_0^{\sqcup_kI_k}$, a scheme $X_1$ with a point $\xi_1:\operatorname{Spec}\Omega_0\to X_1$ of which
every point of $X_1$ is a specialization, covered by the affine open subschemes of the charts of $y$ as
in Definition~\ref{cv-origin-hodge-component}(1), a morphism $f_1:X_1\to\mathcal X$ with
$f_1\circ\xi_1=\xi_0$, an embedding $F_c\to\Omega_0$ through which $f_1\circ\xi_1$ lies over
$\mathcal O_c$, a semi-abelian scheme $\mathcal A_1$ over $X_1$ with an action of $\mathcal O_E$ and a
symmetric cubical line bundle, identified with $f_1^*\mathcal A$ with these structures, and for each
label $s$ an $\mathcal O_{X_1,P}$-submodule $L_s(P)$, at each point $P$ of $X_1$, of the Lie algebra
$\mathrm{Lie}_1$ of $\mathcal A_1$ at $\xi_1$, and an ideal $J_\partial$ of
$\mathbf Z[x_{k,i}]$, such that:
\begin{enumerate}
\item the $L_s$ form a vector bundle of rank $r_s$ on $X_1$ whose generic fibre is the subspace of
  $\mathrm{Lie}_1$ on which each $b\in\mathcal O_E$ acts by $s(b)$;
\item at every point $P$ over no prime, $L_s(P)$ is the set of vectors of that subspace that lie in the
  Lie algebra of $\mathcal A_1$ at $P$;
\item $J_\partial$ satisfies condition (3) of Definition~\ref{cv-origin-reference} for $y$, with
  $|\Psi|pq$ in place of $D$;
\item the complex points of $X_1$ whose coordinates lie in the zero set of $J_\partial$ are those whose
  image under $f_1$ lies outside the abelian locus of $\mathcal A$.
\end{enumerate}
\end{lemma}
\leanref{Development/AutomorphicSpectrum/Proofs/GeometryProofs/ExistsCoordsForComponentStatement.lean}{GlobalTheory.T.exists_coords_for_component}
\begin{proof}
By Lemma~\ref{cv-origin-reference-char}, $\Omega_0$ has characteristic zero.

\emph{A first model.} Lemma~\ref{cv-origin-reference-base} gives finite sets $K_0$ and $I^0_k$, a point
$y_0\in\Omega_0^{\sqcup_kI^0_k}$, a scheme $T_0$ with a morphism $\xi:\operatorname{Spec}\Omega_0\to T_0$ of whose image $\eta$
every point is a specialization, covered by the affine open subschemes of the charts of $y_0$, a morphism
$f_0:T_0\to\mathcal X$ with $f_0\circ\xi=\xi_0$, and an ideal $J_0$ that satisfies condition (3) of
Definition~\ref{cv-origin-reference} for $y_0$ with $|\Psi|pq$ in place of $D$, such that the complex points of $T_0$ whose
coordinates lie in the zero set of $J_0$ are those whose image under $f_0$ lies outside the abelian locus of $\mathcal A$. By
condition (1) of Definition~\ref{cv-origin-segre-coordinates}, the local ring of $T_0$ at $\eta$ maps bijectively onto
$\Omega_0$ and those at the other points injectively. The point $\eta$ lies over no prime, as every nonzero element of
$\Omega_0$ is a unit there.

\emph{Constants and Lie algebras.} The local ring of $\mathcal X$ at the image $x_0$ of $\xi_0$ is a field, as it maps
bijectively onto $\Omega_0$. Its primes correspond to the generizations of $x_0$ \cite[Tag 01J7]{Sta}, so $x_0$ has no
generization other than itself, and it is the generic point of an irreducible component of $\mathcal X$. By
Lemma~\ref{cv-origin-change-factor}, $\xi_0=j\circ y_p$ for a point $y_p$ of $X$ with values in $\Omega_0$. As $X$ lies over
$F_c$, $y_p$ gives an embedding $F_c\to\Omega_0$, through which $\xi_0=f_0\circ\xi$ lies over $\mathcal O_c$. Let
$\mathcal A_0=f_0^*\mathcal A$, with its morphism to $\mathcal A$ over $f_0$ (Lemma~\ref{cv-origin-reference-pullback}(1)), let
$\mathrm{Lie}$ be its Lie algebra at $\xi$ with the action $\rho$ of $\mathcal O_E$ (Lemma~\ref{cv-origin-reference-lie}), and
for a label $s$ let $\mathrm{Lie}_s$ be the subspace on which each $b\in\mathcal O_E$ acts by $s(b)$, through $F_c\to\Omega_0$.
By Lemma~\ref{cv-origin-reference-rank}, $\mathrm{Lie}$ is finite-dimensional and $\dim\mathrm{Lie}_s=r_s$.

\emph{The lattices.} By Lemma~\ref{cv-origin-reference-integral}, applied to $J_0$, the local ring of $T_0$ at every point
over no prime is integrally closed. Let $\Lambda$ be the lattice family of the Lie algebras of $\mathcal A_0$ at the points of
$T_0$ (Definition~\ref{cv-origin-reference-lattices}). As $\mathcal A_0$ is smooth over $T_0$, $\Lambda$ is a vector bundle
of rank $\dim\mathrm{Lie}$ with generic fibre $\mathrm{Lie}$ (Lemma~\ref{cv-origin-reference-lie-bundle}): near every point
it is generated by a basis of $\mathrm{Lie}$, and $\Lambda(P)$ spans $\mathrm{Lie}$ for every $P$. By
Lemma~\ref{cv-origin-reference-eigen}, the $\mathrm{Lie}_s$ are independent and
$\Lambda(P)=\sum_s(\Lambda(P)\cap\mathrm{Lie}_s)$ at every point $P$ over no prime. At $P=\eta$, where $\Lambda(\eta)$ spans
$\mathrm{Lie}$, this gives $\mathrm{Lie}=\bigoplus_s\mathrm{Lie}_s$. Let $\pi_s$ be the projection onto $\mathrm{Lie}_s$ along
the other summands, and $L^0_s$ the lattice family $P\mapsto\pi_s(\Lambda(P))$. Near every point, $L^0_s$ is generated by the
images of a basis of $\mathrm{Lie}$, and $L^0_s(P)$ spans $\pi_s(\mathrm{Lie})=\mathrm{Lie}_s$ for every $P$. At a point $P$
over no prime, $\pi_s$ is the identity on $\Lambda(P)\cap\mathrm{Lie}_s$ and kills $\Lambda(P)\cap\mathrm{Lie}_{s'}$ for
$s'\neq s$, so $L^0_s(P)=\Lambda(P)\cap\mathrm{Lie}_s$.

\emph{A refinement.} Let $x$ be the Segre products of $y_0$, the coordinates of $T_0$
(Definition~\ref{cv-origin-segre-coordinates}); $\eta$ lies in a chart of $x$, so some entry of $x$ is nonzero.
Lemma~\ref{cv-origin-reference-refinement}, applied to $T_0$, $x$, the single family $w=x$, and the lattice families $L^0_s$
with the subspaces $\mathrm{Lie}_s$, gives a finite family $x'$, a scheme $X_1$ with a morphism
$\xi_1:\operatorname{Spec}\Omega_0\to X_1$ of whose image every point is a specialization, with the coordinates $x'$, and a
morphism $g:X_1\to T_0$ with $g\circ\xi_1=\xi$, such that every point of $X_1$ lies in a chart of $x$, each $g^*L^0_s$ is a
vector bundle of rank $r_s$ with generic fibre $\mathrm{Lie}_s$, and $g$ is a normalization over every point of $T_0$ in a
chart of $x$ at which each $L^0_s$ is generated by linearly independent vectors. A point in the chart $c$ of $x$ lies in the
chart $c_k$ of each block $k$ of $y_0$, as $y_{0,k,i}/y_{0,k,c_k}$ is the Segre ratio for the chart $c$ with $c_k$ replaced
by $i$. So Lemma~\ref{cv-origin-reference-segre}, applied to $X_1$, with $x'$ as a point with one block, and to the blocks of
$y_0$, shows that $X_1$ is covered by the affine open subschemes of the charts of the point $y$ with the blocks of $y_0$ and
the block $x'$. Its blocks are $K=K_0\sqcup\{*\}$, with $I_k=I^0_k$ and $y_{k,i}=y_{0,k,i}$ for $k\in K_0$.

\emph{The semi-abelian scheme.} Lemma~\ref{cv-origin-reference-pullback} gives $\mathcal A_1=g^*\mathcal A_0$ with a morphism
to $\mathcal A_0$ over $g$, and its composite with $\mathcal A_0\to\mathcal A$, a morphism $\mathcal A_1\to\mathcal A$ over
$f_1=f_0\circ g$. By Lemma~\ref{cv-origin-reference-lie-same}, the Lie algebra $\mathrm{Lie}_1$ of $\mathcal A_1$ at $\xi_1$ is
identified with $\mathrm{Lie}$, compatibly with the actions of $\mathcal O_E$, and the Lie algebras of $\mathcal A_1$ at the
points of $X_1$ form $g^*\Lambda$. Put $L_s=g^*L^0_s$, and let $J_\partial$ be the image of $J_0$ under the renaming of the
blocks of $y_0$ into those of $y$.

\emph{Points over no prime.} Let $P'$ be a point of $X_1$ over no prime, and $P=g(P')$. Every prime is regular at $P$ and a
unit at $P'$, so it is a unit at $P$ by Lemma~\ref{cv-origin-segre-regular}(3): $P$ lies over no prime. It lies in a chart of
$x$, and each $L^0_s(P)=\Lambda(P)\cap\mathrm{Lie}_s$ is generated by linearly independent vectors, by
Lemma~\ref{cv-origin-reference-eigen}(3). So $g$ is a normalization over $P$, and, as $\mathcal O_{T_0,P}$ is integrally
closed, Lemma~\ref{cv-origin-reference-normalization} shows that the local ring of $X_1$ at $P'$ equals that of $T_0$ at $P$
and that $P'$ is the only point of $X_1$ over $P$.

\emph{The conclusions.} The composite $f_1\circ\xi_1=f_0\circ\xi=\xi_0$ lies over $\mathcal O_c$ through $F_c\to\Omega_0$,
and (1) holds by the choice of $L_s$. For (2), let $P'$ be a point of $X_1$ over no prime and $P=g(P')$. As the local rings
at $P'$ and $P$ are equal, $L_s(P')$, the module generated by $L^0_s(P)$, is $L^0_s(P)=\Lambda(P)\cap\mathrm{Lie}_s$; in the
same way the Lie algebra of $\mathcal A_1$ at $P'$ is $(g^*\Lambda)(P')=\Lambda(P)$, so $L_s(P')$ is the set of vectors of
$\mathrm{Lie}_s$ in it. For (3), Lemma~\ref{cv-origin-reference-snc}, applied to $g$, the renaming of blocks, and the local
rings and preimages of the previous paragraph, shows that $J_\partial$ satisfies condition (3) for $y$ with $|\Psi|pq$. For
(4), let $w$ be coordinates for $y$ of a complex point $z$ of $X_1$. By Lemma~\ref{cv-origin-reference-points},
$w_0=(w_{k,i})_{k\in K_0}$ are coordinates for $y_0$ of $g\circ z$. A polynomial of $J_0$ and its renaming take the same value
at $w_0$ and at $w$, and the renamings generate $J_\partial$; so $w$ lies in the zero set of $J_\partial$ if and only if
$w_0$ lies in that of $J_0$, that is, if and only if $f_0(g(z))=f_1(z)$ lies outside the abelian locus of $\mathcal A$.
\end{proof}
% Lean: exists_coords_for_component = exists_eigBase, step by step: T.exists_coords_of_generic_point (Lemma
% cv-origin-reference-base); T.component_generic_is_generic_point (GenPtModel.genPt_mem_genericPoints) and
% T.generic_point_factors_open with EigBase.exists_cst for the constants; sabPull; lieModel_exists and
% lieModel_exists_unique_act; T.eigen_lattice_constant_rank; EigBase.isIntegrallyClosed_of_snc;
% T.lie_omega_lattice_vector_bundle; SatModel.lat_eq_iSup_of_char0 with Flat.ZLat.unitAt_gp (htop); EigBase.exists_proj,
% T.coherent_lattice_closure, span_map_eq_range and map_proj_eq for the L^0_s; T.flag_model_common_refinement;
% T.segre_refine_line_bundles with no line bundles (ExtraBlocks.exists_blocks_of_lines, whose proof with no line
% bundles is the construction join_coords_coordTie of Lemma cv-origin-reference-segre); sabPull, sabComp and
% SatModel.exists_lieModel_pull; inter_isVBOn_char0, RegAtDom.regAt_dominates and SatModel.normAt_id_elim for hreg and
% huniq; EigBase.mem_pull_iff; T.snc_model_reindexed_coords; T.complex_coords_reindexed and EigBase.mem_mcPts_rename.

\subsection{Flag bundles and descent}\label{subsection-cv-flag-bundles-descent}

\begin{lemma}\label{cv-origin-reference-flag-bundle-twist}
Let $X$ be a scheme, $\Omega$ a field and $\xi:\operatorname{Spec}\Omega\to X$ a morphism of whose image every point of $X$ is
a specialization, and let $x=(x_j)_{j\in J}$ be a finite family in $\Omega$ such that $X$ has the coordinates $x$
(Definition~\ref{cv-origin-segre-coordinates}). Let $V$ be an $\Omega$-vector space and $L$ a lattice family in $V$ that is a
vector bundle of rank $r$ with generic fibre a subspace $W$ (Definition~\ref{cv-origin-reference-lattices}). Then there are an
integer $m\ge0$ and linear forms $\varphi_1,\dots,\varphi_n$ on $V$ such that for every index $c$ and every point $P$ in the
chart $c$ of $x$:
\begin{enumerate}
\item $\varphi_j(e)/x_c^m\in\mathcal O_P$ for every $j$ and every $e\in L(P)$;
\item for every linear form $w$ on $V$ with $w(L(P))\subset\mathcal O_P$ there are $g_1,\dots,g_n\in\mathcal O_P$ with
  $w(e)=\sum_jg_j\varphi_j(e)/x_c^m$ for every $e\in L(P)$.
\end{enumerate}
\end{lemma}
\leanref{Development/AutomorphicSpectrum/Proofs/GeometryProofs/ExistsTwistedGeneratorsStatement.lean}{GlobalTheory.T.exists_twisted_generators}
\leanref{Development/AutomorphicSpectrum/Proofs/GeometryProofs/ExistsTwistedGenerators.lean}{GlobalTheory.CompanionTheory.Ar.FlagBun.Twist.exists_twistGen}
\leanref{Development/AutomorphicSpectrum/Proofs/GeometryProofs/CoordTieCompEquiv.lean}{GlobalTheory.CompanionTheory.Ar.FlagBun.IsTwistGen}
\begin{proof}
For an index $c$ with $x_c\neq0$ let $X_c$ and $A_c=\mathbf Z[x_j/x_c:j\in J]$ be as in
Lemma~\ref{cv-origin-reference-descent-roots-piece}; the elements of $A_c$ are regular at the points of $X_c$. By
Lemma~\ref{cv-origin-snc-charts}(2), applied to $x$ as a point with one block, for every prime $\mathfrak p$ of $A_c$ there is
exactly one point of $X_c$ at which the elements of $A_c$ that are units are those outside $\mathfrak p$.

\emph{Frames on a chart.} Let $x_c\neq0$, and let $G$ be the set of the $s\in A_c$ for which there is a linearly independent
family $b$ of $r$ vectors with $L(Q)=\sum_i\mathcal O_Qb_i$ at every point $Q$ of $X_c$ at which $s$ is a unit. The set $G$
generates the unit ideal of $A_c$. Otherwise it lies in a maximal ideal $\mathfrak m$; let $Q$ be the point of $X_c$ at which the
elements of $A_c$ that are not units are those of $\mathfrak m$. As $L$ is a vector bundle, $Q$ has an open neighbourhood $U$ with
a basis $b$ of $W$ such that $L(Q')=\sum_i\mathcal O_{Q'}b_i$ for $Q'\in U$, and
Lemma~\ref{cv-origin-reference-descent-roots-piece}(1) gives $s\in A_c$, a unit at $Q$, such that every point of $X_c$ at which
$s$ is a unit lies in $U$. Then $s\in G$ and $s\notin\mathfrak m$. So a finite subset $S_c$ of $G$ generates the unit ideal, and at
every point $Q$ of $X_c$ some element of $S_c$ is a unit, as the elements of $A_c$ that are not units at $Q$ form a proper ideal.
For $s\in S_c$ fix a family $b_s$ as in the definition of $G$.

\emph{Degrees.} For $s\in A_c$ there is $d_s\ge0$ such that $sx_c^{d'}/x_{c'}^{d'}\in A_{c'}$ for every $d'\ge d_s$ and every $c'$ with
$x_{c'}\neq0$: this holds for the generators $x_j/x_c$ with $d_s=1$, as $(x_j/x_c)(x_c/x_{c'})=x_j/x_{c'}$, and for the integers
with $d_s=0$; the property passes to sums and products, adding the integers $d_s$; and it passes from $d'$ to $d'+1$, as
$x_c/x_{c'}\in A_{c'}$. Let $I$ be the finite set of the pairs $i=(c_i,s_i)$ with $x_{c_i}\neq0$ and $s_i\in S_{c_i}$, put
$b_i=b_{s_i}$, and let $d$ be the largest $d_{s_i}$.

\emph{Dual bases.} For $i\in I$ choose linear forms $\beta_{i,1},\dots,\beta_{i,r}$ on $V$ with $\beta_{i,k}(b_{i,l})=1$ if
$k=l$ and $0$ otherwise. If $Q$ is a point of $X_{c_i}$ at which $s_i$ is a unit and $e\in L(Q)$, then
$e=\sum_l\beta_{i,l}(e)b_{i,l}$ with all $\beta_{i,l}(e)\in\mathcal O_Q$.

\emph{A unit criterion.} Let $i\in I$, let $x_{c'}\neq0$, put $\alpha=s_ix_{c_i}^{d+1}/x_{c'}^{d+1}\in A_{c'}$, and let $Q$ be a
point of $X_{c'}$ at which $\alpha$ is a unit. Then $Q$ lies in the chart $c_i$ and $s_i$ is a unit at $Q$. Indeed, $y=x_{c_i}/x_{c'}$
and $z=s_ix_{c_i}^d/x_{c'}^d$ lie in $A_{c'}$, and $\alpha=zy$, so $y^{-1}=z\alpha^{-1}$ is regular at $Q$, and so is every
$x_j/x_{c_i}=(x_j/x_{c'})y^{-1}$. Then $s_i\in A_{c_i}$ is regular at $Q$, and $s_i^{-1}=\alpha^{-1}y^{d+1}$ is regular at $Q$.

\emph{Comparison.} Let $i,i'\in I$ and $k,l\le r$, put $c'=c_{i'}$, $\alpha=s_ix_{c_i}^{d+1}/x_{c'}^{d+1}$ and $t=\alpha s_{i'}\in A_{c'}$.
Then $t^N\beta_{i,k}(b_{i',l})\in A_{c'}$ for some $N\ge0$. If $t=0$ this holds with $N=1$. Otherwise, let $Q$ be a point of $X_{c'}$
at which $t$ is a unit. Its factors $\alpha$ and $s_{i'}$ are regular at $Q$, so they are units there. By the unit criterion $Q$
lies in the chart $c_i$ and $s_i$ is a unit at $Q$; as $b_{i',l}\in L(Q)$, $\beta_{i,k}(b_{i',l})\in\mathcal O_Q$. So
$\beta_{i,k}(b_{i',l})$ is regular at every point of $X_{c'}$ at which $t$ is a unit, and
Lemma~\ref{cv-origin-reference-descent-roots-piece}(2) gives $N$. Let $M$ be the largest of these integers $N$ over all $i,i',k,l$.

\emph{The forms.} Put $m=(d+1)M$ and $\varphi_{i,k}=(s_ix_{c_i}^{d+1})^M\beta_{i,k}$ for $i\in I$ and $k\le r$. Let $P$ be a point
in the chart $c$ of $x$. Some $j=(c,s')\in I$ has $s'=s_j$ a unit at $P$, and every $e\in L(P)$ is $\sum_l\beta_{j,l}(e)b_{j,l}$
with $\beta_{j,l}(e)\in\mathcal O_P$. For $(i,k)$ put $\alpha=s_ix_{c_i}^{d+1}/x_c^{d+1}\in A_c$. Then
$\varphi_{i,k}(e)/x_c^m=\alpha^M\beta_{i,k}(e)=\sum_l\beta_{j,l}(e)\,\alpha^M\beta_{i,k}(b_{j,l})$, and with the integer $N\le M$ of the
comparison for $(i,j,k,l)$,
\[
\alpha^M\beta_{i,k}(b_{j,l})=\alpha^{M-N}s'^{-N}\cdot(\alpha s')^N\beta_{i,k}(b_{j,l})\in\mathcal O_P,
\]
as $s'^{-1}\in\mathcal O_P$. This gives (1). For (2), let $w(L(P))\subset\mathcal O_P$, and put $g_{j,l}=w(b_{j,l})s'^{-M}$, which
lies in $\mathcal O_P$, and $g_{i,k}=0$ for $i\neq j$. For $e\in L(P)$, $\varphi_{j,l}(e)/x_c^m=s'^M\beta_{j,l}(e)$, as $x_{c_j}=x_c$; so
$\sum_{i,k}g_{i,k}\varphi_{i,k}(e)/x_c^m=\sum_lw(b_{j,l})\beta_{j,l}(e)=w(e)$.
\end{proof}
% Lean: exists_twistGen: frame_cover (chart_prime, chart_nbhd, Ideal.span_eq_top_iff_finite, chart_point), exists_dual,
% exists_isSec (IsSec, induction on the adjoin), d = sup d0 with hd and hd1, compare (chart_denom, the Lean's form of
% Lemma cv-origin-reference-descent-roots-piece(2); unit_chart; frame_coords), M = sup N, phi t = (s x^(d+1))^M beta, and
% twistGen_of_fintype with power_step and twist_eval for (1) and (2). The Lean's IsTwistGen is (1) and (2).

\begin{lemma}\label{cv-origin-reference-flag-bundle-inverse}
Let $X$, $\Omega$, $\xi$, $x$, $V$, $L$, $r$ and $W$ be as in Lemma~\ref{cv-origin-reference-flag-bundle-twist}, and let $m$ and
$\varphi_1,\dots,\varphi_n$ satisfy its conclusions (1) and (2).
\begin{enumerate}
\item If $c_1,\dots,c_r$ is a basis of $W$, then there is an $r\times n$ matrix $B$ over $\Omega$ with
  $B\cdot(\varphi_j(c_l))_{j,l}=1$.
\item If $P$ is a point in the chart $c$ of $x$ and $\beta_1,\dots,\beta_r$ are linearly independent vectors with
  $L(P)=\sum_a\mathcal O_P\beta_a$, then the entries of the $n\times r$ matrix $A=(\varphi_j(\beta_a)/x_c^m)_{j,a}$ lie in
  $\mathcal O_P$, and there is an $r\times n$ matrix $B$ with entries in $\mathcal O_P$ and $BA=1$.
\end{enumerate}
\end{lemma}
\leanref{Development/AutomorphicSpectrum/Proofs/GeometryProofs/FrameViaTwistedGeneratorsStatement.lean}{GlobalTheory.T.frame_via_twisted_generators}
\leanref{Development/AutomorphicSpectrum/Proofs/GeometryProofs/FrameCoeffsRegularInChartStatement.lean}{GlobalTheory.T.frame_coeffs_regular_in_chart}
\begin{proof}
(2) The entries of $A$ lie in $\mathcal O_P$ by (1) of Lemma~\ref{cv-origin-reference-flag-bundle-twist}. Choose linear forms
$w_1,\dots,w_r$ on $V$ with $w_a(\beta_b)=1$ if $a=b$ and $0$ otherwise; they map $L(P)$ into $\mathcal O_P$. By (2) of that
lemma there are $g_{a,j}\in\mathcal O_P$ with $w_a(e)=\sum_jg_{a,j}\varphi_j(e)/x_c^m$ for $e\in L(P)$; for $e=\beta_b$ this says
that $B=(g_{a,j})$ has $BA=1$.

(1) The image $\eta$ of $\xi$ lies in a chart $k$ of $x$, and $\mathcal O_\eta=\Omega$. As $L$ is a vector bundle with generic fibre
$W$, $L(\eta)$ is spanned by a basis of $W$, so $L(\eta)=W$. Choose linear forms $w_1,\dots,w_r$ on $V$ with $w_i(c_l)=1$ if
$i=l$ and $0$ otherwise. Every linear form maps $L(\eta)$ into $\mathcal O_\eta$, so by (2) of
Lemma~\ref{cv-origin-reference-flag-bundle-twist} there are $g_{i,j}\in\Omega$ with $w_i=\sum_jg_{i,j}\varphi_j/x_k^m$ on $W$.
The matrix $B=(g_{i,j}/x_k^m)$ has $B\cdot(\varphi_j(c_l))=(w_i(c_l))=1$.
\end{proof}
% Lean: (1) is exists_leftInverse_flagForms (the dual forms w from Basis.span, IsTwistGen at genPt with regAt_generic_eq_top,
% and LinearMap.eqOn_span' on the frame b at the generic point); (2) is twist_left_inverse.

\begin{lemma}\label{cv-origin-reference-flag-bundle-coords}
Let $\Omega$ be a field and $x=(x_j)_{j\in J}$ a finite family in $\Omega$ with an index $c_0$ such that $x_{c_0}\neq0$ and every
element of $\Omega$ is a quotient of two elements of $\mathbf Z[x_j/x_{c_0}:j\in J]$. For $s$ in a finite set $S$ let $r_s,n_s\ge0$
and let $F^0_s$ be an $n_s\times r_s$ matrix over $\Omega$ with a left inverse. Let $\Omega'$ be the field of rational functions over
$\Omega$ in indeterminates $u_{s,lk}$, for $s\in S$ and $1\le k<l\le r_s$, let $U_s$ be the $r_s\times r_s$ matrix over $\Omega'$
whose entry in the row $l$ and the column $k$ is $u_{s,lk}$ if $l>k$, $1$ if $l=k$ and $0$ if $l<k$, and put $F_s=F^0_sU_s$. For
$1\le k\le r_s$ and a map $\mu:\{1,\dots,k\}\to\{1,\dots,n_s\}$ let $\Delta_{s,k}(\mu)$ be the determinant of the $k\times k$ matrix
of the entries of $F_s$ in the rows $\mu(1),\dots,\mu(k)$ and the columns $1,\dots,k$. For an index $c\in J$ and a family
$\mu=(\mu_{s,k})$ of such maps, one for each $s$ and $k\le r_s$, put
\[
y_{(c,\mu)}=x_c\prod_{s,k}\Delta_{s,k}(\mu_{s,k}).
\]
Then there is a family $\rho$ with $y_{(c_0,\rho)}\neq0$ such that every element of $\Omega'$ is a quotient of two elements of the
ring generated by the ratios $y_{(c,\mu)}/y_{(c_0,\rho)}$.
\end{lemma}
\leanref{Development/AutomorphicSpectrum/Proofs/GeometryProofs/FrameMatrixChartRatiosStatement.lean}{GlobalTheory.T.frame_matrix_chart_ratios}
\leanref{Development/AutomorphicSpectrum/Proofs/GeometryProofs/CoordTieCompEquiv.lean}{GlobalTheory.CompanionTheory.Ar.FlagBun.flagCoord}
\begin{proof}
Let $B_s$ be a left inverse of $F^0_s$; then $B_sF_s=U_s$. For $1\le k\le r_s$ and a map $R:\{1,\dots,k\}\to\{1,\dots,r_s\}$, the
rows $R(1),\dots,R(k)$ of the first $k$ columns of $U_s$ are combinations of the rows of the first $k$ columns of $F_s$ with
coefficients from $B_s$; expanding the determinant by multilinearity in the rows gives
\[
\det\bigl((U_s)_{R(i),k'}\bigr)_{i,k'\le k}=\sum_{f:\{1,\dots,k\}\to\{1,\dots,n_s\}}\Bigl(\prod_{i\le k}(B_s)_{R(i),f(i)}\Bigr)\Delta_{s,k}(f).
\tag{$*$}
\]
For $l\ge k$ let $R_l$ be the map with $R_l(i)=i$ for $i<k$ and $R_l(k)=l$. The matrix on the left of $(*)$ for $R=R_l$ is lower
triangular with diagonal entries $1,\dots,1,(U_s)_{lk}$, so its determinant is $u_{s,lk}$ if $l>k$ and $1$ if $l=k$. By $(*)$ for
$R=R_k$, some $\Delta_{s,k}(f)$ is nonzero; choose such an $f=\rho_{s,k}$. Then $y_{(c_0,\rho)}\neq0$.

Let $L$ be the subfield of $\Omega'$ generated by the ratios $y_{(c,\mu)}/y_{(c_0,\rho)}$. As $y_{(j,\rho)}/y_{(c_0,\rho)}=x_j/x_{c_0}$,
$L$ contains $\mathbf Z[x_j/x_{c_0}]$, and hence $\Omega$. For $s$, $k$ and $f$, replacing $\rho_{s,k}$ by $f$ in $\rho$ gives a family
$\mu$ with $y_{(c_0,\mu)}/y_{(c_0,\rho)}=\Delta_{s,k}(f)/\Delta_{s,k}(\rho_{s,k})$, which therefore lies in $L$. Dividing $(*)$ by
$\Delta_{s,k}(\rho_{s,k})$ shows, as the entries of $B_s$ lie in $\Omega\subset L$, that the left side of $(*)$ divided by
$\Delta_{s,k}(\rho_{s,k})$ lies in $L$ for every $R$; with $R=R_l$ and $R=R_k$, $u_{s,lk}$ is a quotient of two elements of $L$. So
$L=\Omega'$, and as $L$ is the field of fractions of the ring generated by the ratios, every element of $\Omega'$ is a quotient of
two elements of that ring.
\end{proof}
% Lean: CoordGen.flagCoord_generate: exists_minor_ne_zero (det_genU_submatrix through det_mul_eq_sum, and det_genU_rows), the
% index (c0, rho), base_mem, ratio_update for the minors, genU_mem and mem_of_genU_mem, and exists_div_of_mem_closure. The
% Lean's flagCoord indexes the maps by k in Fin r with k + 1 columns; its matrix flagMat is (phi s j (c s l)) * genU.

\begin{lemma}\label{cv-origin-reference-flag-bundle-valuative}
Let $\Omega$, $x$ and $j_0$ be as in Lemma~\ref{cv-origin-segre-model}. Then the scheme $X$ and the morphism $\xi$ of that lemma
can be chosen such that, moreover, $X$ is integral and the morphism $X\to\operatorname{Spec}\mathbf Z$ is proper.
\end{lemma}
\leanref{Development/AutomorphicSpectrum/Proofs/GeometryProofs/ExistsProperIntegralModelStatement.lean}{GlobalTheory.T.exists_proper_integral_model}
\leanref{Development/AutomorphicSpectrum/Proofs/GeometryProofs/ExistsProperIntegralModel.lean}{GlobalTheory.CompanionTheory.Ar.exists_cl_scheme}
\begin{proof}
Let $X=\operatorname{Proj}S$ and $\xi$ be as constructed in the proof of Lemma~\ref{cv-origin-segre-model}, and $\eta$ the image
of $\xi$.

\emph{Integrality.} Every point of $X$ specializes $\eta$, so $X$ is the closure of $\eta$, which is irreducible. The local rings
of $X$ embed in the field $\Omega$, so they are reduced, and $X$ is reduced. So $X$ is integral \cite[Tag 01ON]{Sta}.

\emph{Properness.} The scheme $X$ is covered by the finitely many affine open subschemes $D_+(t_j)$ with $x_j\neq0$, whose rings
$S_{(t_j)}\cong A_j$ are generated over $S_0$ by the $t_i/t_j$; so $\operatorname{Proj}S\to\operatorname{Spec}S_0$ is quasi-compact
and of finite type. The graded ring $S$ is generated over $S_0$ by the finitely many elements $t_j$, so this morphism satisfies the
existence part of the valuative criterion \cite[Tag 01MF]{Sta}; as it is quasi-compact, it is universally closed
\cite[Tag 01KF]{Sta}. The scheme $X$ is separated \cite[Tag 01MC]{Sta}, so the morphism is separated \cite[Tag 01KV]{Sta}, and it is
proper \cite[Tag 01W1]{Sta}. The ring $S_0$ is the image of $\mathbf Z$ in $\Omega$, so $\operatorname{Spec}S_0\to\operatorname{Spec}\mathbf Z$
is a closed immersion \cite[Tag 01IG]{Sta}, which is proper \cite[Tag 01W5]{Sta}. So the composite
$X\to\operatorname{Spec}\mathbf Z$ is proper \cite[Tag 01W3]{Sta}.
\end{proof}
% Lean: exists_cl_scheme (= T.exists_proper_integral_model, here with o = Z) gives IsIntegral X (in ClSch.tie_of_charts:
% irreducible as the closure of the generic point, reduced as the stalks embed in Omega) and IsProper s for
% s = Proj.toSpecZero followed by Spec of zeroMap: Mathlib's properness of Proj of a finitely generated graded algebra
% (finiteType_grading), which the book takes from Stacks 01MF, 01KF, 01MC and 01W1, and the closed immersion of
% zeroMap_surjective (IsClosedImmersion.spec_of_surjective). The existence part of the valuative criterion is taken from the
% properness in Lemma cv-origin-reference-flag-bundle-onto, as in FlagBun.surjective_of_dominates.

\begin{lemma}\label{cv-origin-reference-flag-bundle-onto}
Let $X$, $\Omega$, $\xi$ and $x$ be as in Lemma~\ref{cv-origin-reference-flag-bundle-twist}, let $\iota:\Omega\to\Omega'$ be an
embedding of fields, and let $X'$ be a scheme with a morphism $\xi':\operatorname{Spec}\Omega'\to X'$ such that the morphism
$X'\to\operatorname{Spec}\mathbf Z$ is proper. Then every morphism $f:X'\to X$ with $f\circ\xi'=\xi\circ\operatorname{Spec}\iota$ is
surjective.
\end{lemma}
\leanref{Development/AutomorphicSpectrum/Proofs/GeometryProofs/ProperDominatingSurjectiveStatement.lean}{GlobalTheory.T.proper_dominating_surjective}
\leanref{Development/AutomorphicSpectrum/Proofs/GeometryProofs/ProperDominatingSurjective.lean}{GlobalTheory.CompanionTheory.Ar.FlagBun.surjective_of_dominates}
\begin{proof}
Let $\eta$ be the image of $\xi$ and $Q$ a point of $X$; it specializes $\eta$. The ring $\iota(\mathcal O_Q)$ is a local subring of
$\Omega'$, so there is a valuation ring $A$ of $\Omega'$ with field of fractions $\Omega'$ that dominates it \cite[Tag 00IA]{Sta}:
$\iota(\mathcal O_Q)\subset A$, and $\iota$ maps the non-units of $\mathcal O_Q$ into the maximal ideal of $A$. The morphism
$X'\to\operatorname{Spec}\mathbf Z$ is proper, so it is of finite type, hence quasi-compact, and universally closed
\cite[Tag 01W1]{Sta}; so it satisfies the existence part of the valuative criterion \cite[Tag 01KF]{Sta}. Applied to $\xi'$ and the
morphism $\operatorname{Spec}A\to\operatorname{Spec}\mathbf Z$, it gives a morphism $l:\operatorname{Spec}A\to X'$ with
$\xi'=l\circ\operatorname{Spec}(A\subset\Omega')$. Let $g=f\circ l$, and let $Z$
be the image under $g$ of the closed point of $\operatorname{Spec}A$. Then $g\circ\operatorname{Spec}(A\subset\Omega')=
\xi\circ\operatorname{Spec}\iota$, so $Z$ specializes $\eta$, and the argument of the proof of
Lemma~\ref{cv-origin-coords-join-map}(1), with $g$ and the closed point of $\operatorname{Spec}A$ in place of $f$ and $P'$, shows that
the local homomorphism $\mathcal O_{X,Z}\to A$ of $g$, followed by the inclusion of $A$ in $\Omega'$, is $\iota$ applied to the map
$\mathcal O_{X,Z}\to\Omega$. So for each of the points $P=Q$ and $P=Z$: $\iota(\mathcal O_P)\subset A$, and an element $a$ of
$\mathcal O_P$ is a unit at $P$ if and only if $\iota(a)$ is a unit of $A$.

The point $Q$ lies in a chart of $x$, so some $x_{c'}$ is nonzero. As $A$ is a valuation ring, there is $c$ with $x_c\neq0$ such
that $\iota(x_{c'}/x_c)\in A$ for every $c'$. Let $P$ be $Q$ or $Z$, in a chart $c''$ of $x$. Then $x_c/x_{c''}\in\mathcal O_P$, and
$\iota$ of its inverse lies in $A$, so $\iota(x_c/x_{c''})$ is a unit of $A$ and $x_c/x_{c''}$ is a unit at $P$. So every
$x_{c'}/x_c=(x_{c'}/x_{c''})(x_c/x_{c''})^{-1}$ is regular at $P$, and $P$ lies in the chart $c$. For $a\in\mathbf Z[x_{c'}/x_c]$, $a$ is a
unit at $Q$ if and only if $\iota(a)$ is a unit of $A$, if and only if $a$ is a unit at $Z$. By Lemma~\ref{cv-origin-snc-charts}(2),
applied to $x$ as a point with one block, $Q=Z=f(l(\text{closed point}))$.
\end{proof}
% Lean: surjective_of_dominates: exists_domAt_of_injective_genMap (LocalSubring.exists_le_valuationSubring for A), the lift l
% from the valuative criterion of the proper s', domAt_of_spec (val_comp_spec, through appLE, for the local homomorphism of
% l >> f), Finite.exists_max of the valuations of the i(x c) for the chart c, inChart_of_domAt_of_div_mem with DomAt.inv_mem,
% and DomAt.eq_of_mem_affine (the same basic opens in the affine chart, and the antisymmetry of specialization), where the
% book uses Lemma cv-origin-snc-charts(2).

\begin{lemma}\label{cv-origin-reference-flag-bundle-minors}
Let $K$ be a field, $R\subset K$ a subring, $A$ an $n\times r$ and $B$ an $r\times n$ matrix with entries in $R$ such that $BA=1$, $G$
an $r\times r$ matrix over $K$, $a\in K^\times$ and $F=aAG$. Let $1\le k\le r$ and $\mu:\{1,\dots,k\}\to\{1,\dots,n\}$, let $H$ be the
$k\times k$ matrix of the entries of $F$ in the rows $\mu(1),\dots,\mu(k)$ and the columns $1,\dots,k$, and for $l\le k$ and
$j\le n$ let $H_{l\to j}$ be $H$ with its row $l$ replaced by the entries of $F$ in the row $j$ and the columns $1,\dots,k$. Suppose
that $\det H\neq0$ and $\det H_{l\to j}/\det H\in R$ for all $l$ and $j$. Then there are vectors $D_1,\dots,D_k$ and
$E_1,\dots,E_k$ in $R^r$ with $\sum_iE_{l,i}D_{l',i}=1$ if $l=l'$ and $0$ otherwise, such that every $D_l$ lies in the $K$-span of
the first $k$ columns of $G$.
\end{lemma}
\leanref{Development/AutomorphicSpectrum/Proofs/GeometryProofs/MinorRatioRegularOfInvertibleStatement.lean}{GlobalTheory.T.minor_ratio_regular_of_invertible}
\leanref{Development/AutomorphicSpectrum/Proofs/GeometryProofs/MinorRatioRegularOfInvertible.lean}{GlobalTheory.CompanionTheory.Ar.FlagBun.split_of_regular_ratios}
\begin{proof}
Let $F_k$ and $G_k$ be the matrices of the first $k$ columns of $F$ and of $G$, and put $Q=aG_kH^{-1}$, an $r\times k$ matrix. Then
$AQ=F_kH^{-1}$. By Cramer's rule, the entry of $F_kH^{-1}$ in the row $j$ and the column $l$ is $\det H_{l\to j}/\det H$, which lies in
$R$. So $Q=B(AQ)$ has its entries in $R$. Let $D_l$ be the column $l$ of $Q$, a combination of the columns of $G_k$, and $E_l$ the row
$\mu(l)$ of $A$. Then $\sum_iE_{l,i}D_{l',i}$ is the entry of $AQ$ in the row $\mu(l)$ and the column $l'$, which is the entry of
$HH^{-1}=1$ in the row $l$ and the column $l'$.
\end{proof}
% Lean: split_of_regular_ratios, with mul_inv_apply_eq_det_div (Matrix.cramer_transpose_apply and the adjugate) for the entries
% of F_k H^-1. The Lean indexes k by Fin r, with k + 1 columns.

\begin{lemma}\label{cv-origin-reference-flag-bundle-adapted}
Let $K$ be a field and $R\subset K$ a local subring with residue field $\kappa$, and let $g_1,\dots,g_r\in K^r$. Suppose that for
every $k\le r$ there are $D_1,\dots,D_k$ and $E_1,\dots,E_k$ in $R^r$ with $\sum_iE_{l,i}D_{l',i}=1$ if $l=l'$ and $0$ otherwise,
such that every $D_l$ lies in the $K$-span of $g_1,\dots,g_k$. Then there is an $r\times r$ matrix $P$ with entries in $R$ and
determinant a unit of $R$, such that for every $i$ the first $i$ columns of $P$ and $g_1,\dots,g_i$ span the same subspace of $K^r$.
\end{lemma}
\leanref{Development/AutomorphicSpectrum/Proofs/GeometryProofs/FlagBasisChangeIntegralStatement.lean}{GlobalTheory.T.flag_basis_change_integral}
\leanref{Development/AutomorphicSpectrum/Proofs/GeometryProofs/FlagBasisChangeIntegral.lean}{GlobalTheory.CompanionTheory.Ar.FlagBun.Adapted.exists_adapted}
\begin{proof}
Write $\bar v\in\kappa^r$ for the residue of $v\in R^r$. Let $v_1,\dots,v_t\in R^r$ have linearly independent residues, and let $M$ be
the $r\times t$ matrix with these columns. Its residue $\bar M$ has a left inverse over $\kappa$; lifting its entries to $R$ gives a
$t\times r$ matrix $N$ over $R$ such that the residue of $\det(NM)$ is $1$. So $\det(NM)$ is a unit of $R$, and $v_1,\dots,v_t$ are
linearly independent over $K$; if $t=r$, then $\det M$ is a unit of $R$, as $\det(NM)=\det N\det M$.

We choose $p_1,p_2,\dots\in R^r$ inductively such that, for each $k\le r$, $\bar p_1,\dots,\bar p_k$ are linearly independent and
$p_1,\dots,p_i$ span the same subspace as $g_1,\dots,g_i$ for every $i\le k$. Let $k<r$ and suppose $p_1,\dots,p_k$ are chosen. Take
$D_1,\dots,D_{k+1}$ and $E_1,\dots,E_{k+1}$ as in the hypothesis for $k+1$. If $\sum_lc_l\bar D_l=0$ with $c_l\in\kappa$, then
$c_{l'}=\sum_lc_l\sum_i\bar E_{l',i}\bar D_{l,i}=0$ for every $l'$; so $\bar D_1,\dots,\bar D_{k+1}$ are linearly independent, and one
of them, $\bar D_l$, does not lie in the span of $\bar p_1,\dots,\bar p_k$. Put $p_{k+1}=D_l$. Then $\bar p_1,\dots,\bar p_{k+1}$ are
linearly independent, and by the first paragraph so are $p_1,\dots,p_{k+1}$. They lie in the span of $g_1,\dots,g_{k+1}$, which has
dimension at most $k+1$, so the two spans are equal.

Let $P$ be the matrix with the columns $p_1,\dots,p_r$. By the first paragraph, $\det P$ is a unit of $R$. For $i>r$, both families
span what they span for $i=r$.
\end{proof}
% Lean: exists_adapted: exists_cols (the induction, with red_indep, exists_notMem_span and linearIndependent_of_red, which goes
% through exists_isUnit_det_mul, the lifted left inverse Q), isUnit_det_of_red, and flagStep_of_le for i > r.

\begin{lemma}\label{cv-origin-reference-flag-bundle-frame}
Let $X$, $\Omega$, $\xi$, $x$, $V$, $L$, $r$ and $W$ be as in Lemma~\ref{cv-origin-reference-flag-bundle-twist}, let $m$ and
$\varphi_1,\dots,\varphi_n$ satisfy its conclusions (1) and (2), and let $c_1,\dots,c_r$ be a basis of $W$. Let $P$ be a point in
the chart $c$ of $x$, and $A_c=\mathbf Z[x_j/x_c]$. Then there are $g\in A_c$, a unit at $P$, an invertible $r\times r$ matrix
$T_0$ over $\Omega$, and an $n\times r$ matrix $M$ and an $r\times n$ matrix $N$ with entries in $A_c[1/g]$, such that $NM=1$ and
$(\varphi_j(c_l))_{j,l}=x_c^mMT_0$.
\end{lemma}
\leanref{Development/AutomorphicSpectrum/Proofs/GeometryProofs/TwistedFrameInChartStatement.lean}{GlobalTheory.T.twisted_frame_in_chart}
\leanref{Development/AutomorphicSpectrum/Proofs/GeometryProofs/TwistedFrameInChart.lean}{GlobalTheory.CompanionTheory.Ar.FlagBun.Sm.exists_chart_frame}
\begin{proof}
As $L$ is a vector bundle, $L(P)=\sum_k\mathcal O_Pb_k$ for a basis $b_1,\dots,b_r$ of $W$. Choose linear forms $w_1,\dots,w_r$
on $V$ with $w_k(b_l)=1$ if $k=l$ and $0$ otherwise. Then $c_l=\sum_kw_k(c_l)b_k$, and $T_0=(w_k(c_l))_{k,l}$ is invertible, as
the $c_l$ form a basis of $W$. The matrix $M=(\varphi_j(b_k)/x_c^m)_{j,k}$ has entries in $\mathcal O_P$, by (1) of
Lemma~\ref{cv-origin-reference-flag-bundle-twist}, and $\varphi_j(c_l)=\sum_kw_k(c_l)\varphi_j(b_k)=x_c^m(MT_0)_{jl}$. The forms $w_l$
map $L(P)$ into $\mathcal O_P$, so (2) of that lemma gives $g_{l,j}\in\mathcal O_P$ with $w_l(e)=\sum_jg_{l,j}\varphi_j(e)/x_c^m$ for
$e\in L(P)$; at $e=b_{l'}$ this says that $N=(g_{l,j})$ has $NM=1$. By Lemma~\ref{cv-origin-snc-charts}(1), applied to $x$ as a point
with one block, each of the finitely many entries of $M$ and $N$ is $u/s$ with $u,s\in A_c$ and $s$ a unit at $P$. The product $g$
of these $s$ is a unit at $P$, and all the entries lie in $A_c[1/g]$.
\end{proof}
% Lean: exists_chart_frame: the frame b' of L at p (IsVB), exists_dual_family, the matrix T = (w k (b l)) with
% Matrix.exists_mulVec_eq_zero_iff for det T != 0, M from hTw1, N from hTw2 (hNM), and regAt_iff_chart_localization with
% locSub_le_of_mul and the product g of the denominators.

\begin{lemma}\label{cv-origin-reference-flag-bundle-normalized}
Let $K$ be a field and $O\subset K$ a local subring. For a matrix $Y$ over $K$ with $r$ columns, an integer $k\le r$ and a map $I$
from $\{1,\dots,k\}$ to the rows of $Y$, let $\delta_Y(I)$ be the determinant of the $k\times k$ matrix of the entries of $Y$ in the
rows $I(1),\dots,I(k)$ and the columns $1,\dots,k$, with $\delta_Y(I)=1$ for $k=0$; call $I$ \emph{normalized} for $Y$ if
$\delta_Y(I)\neq0$ and $\delta_Y(I')/\delta_Y(I)\in O$ for every $I'$. Let $A$ be an $n\times r$ and $N$ an $r\times n$ matrix with
entries in $O$ such that $NA=1$, and $Y$ an $r\times r$ matrix over $K$, and suppose that for every $k\le r$ some map
$\{1,\dots,k\}\to\{1,\dots,n\}$ is normalized for $AY$. Then there is a permutation $w$ of $\{1,\dots,r\}$ such that for every
$k\le r$ the map $(w(1),\dots,w(k))$ is normalized for $Y$.
\end{lemma}
\leanref{Development/AutomorphicSpectrum/Proofs/GeometryProofs/FrameMatricesRegularGeneric.lean}{GlobalTheory.CompanionTheory.Ar.FlagBun.Sm.Cell.exists_cell}
\leanref{Development/AutomorphicSpectrum/Proofs/GeometryProofs/ExistsPermNormalizedMinorsStatement.lean}{GlobalTheory.T.exists_perm_normalized_minors}
\leanref{Development/AutomorphicSpectrum/Proofs/GeometryProofs/NormalizedMinorExtendsStatement.lean}{GlobalTheory.T.normalized_minor_extends}
\begin{proof}
For a matrix $Z$ with $r$ columns and a matrix $C$ with as many columns as $Z$ has rows, expanding by multilinearity in the rows
gives
\[
\delta_{CZ}(I)=\sum_J\Bigl(\prod_{a\le k}C_{I(a),J(a)}\Bigr)\delta_Z(J).
\tag{$*$}
\]

(a) For every $k\le r$ some $I$ is normalized for $Y$. For $k=0$ every $I$ is. Let $k\ge1$ and let $\mu$ be normalized for $AY$. As
$Y=N(AY)$, $(*)$ shows that $\delta_Y(I)/\delta_{AY}(\mu)\in O$ for every $I$, and $(*)$ for $AY$ gives
$1=\sum_I\bigl(\prod_aA_{\mu(a),I(a)}\bigr)\delta_Y(I)/\delta_{AY}(\mu)$, a sum of elements of $O$. As $O$ is local, some term is a
unit of $O$. For that $I$, $\delta_Y(I)\neq0$, and $\delta_Y(I')/\delta_Y(I)$ is $\delta_Y(I')/\delta_{AY}(\mu)$ times
$\prod_aA_{\mu(a),I(a)}$ times the inverse of the term, so it lies in $O$.

(b) If $I'$ is normalized for $k+1$ and $I$ for $k$, then $I'\circ f$ is normalized for some injective
$f:\{1,\dots,k\}\to\{1,\dots,k+1\}$. Let $Y_j$ be the matrix of the first $j$ columns of $Y$, let $H'$ and $H$ be the matrices of the
rows $I'$ of $Y_{k+1}$ and the rows $I$ of $Y_k$, and put $Z=Y_{k+1}H'^{-1}$ and $Z'=Y_kH^{-1}$. By Cramer's rule, the entries of $Z$
and $Z'$ are quotients of minors by $\delta_Y(I')$ and by $\delta_Y(I)$, so they lie in $O$. The first $k$ columns of $Y_{k+1}=ZH'$
give $Y_k=Z\cdot(\text{rows }I'\text{ of }Y_k)$, so $Z'=ZZ'_{I'}$, where $Z'_{I'}$ is the matrix of the rows $I'$ of $Z'$. The rows $I$
of $Z'$ form the identity matrix, so $1=Z_IZ'_{I'}$, where $Z_I$ is the matrix of the rows $I$ of $Z$. Taking determinants and
expanding by $(*)$,
\[
1=\sum_{f:\{1,\dots,k\}\to\{1,\dots,k+1\}}\Bigl(\prod_aZ_{I(a),f(a)}\Bigr)\frac{\delta_Y(I'\circ f)}{\delta_Y(I)}.
\]
As $O$ is local, some term is a unit of $O$. Then $\delta_Y(I'\circ f)\neq0$, so $f$ is injective, and for every $I''$,
$\delta_Y(I'')/\delta_Y(I'\circ f)$ is $\delta_Y(I'')/\delta_Y(I)$ times $\prod_aZ_{I(a),f(a)}$ times the inverse of the term, which lies
in $O$.

(c) Reordering the values of $I$ changes $\delta_Y(I)$ by a sign, so it keeps $I$ normalized. We show by descending induction on $k$
that some permutation $w$ has $(w(1),\dots,w(j))$ normalized for every $j\ge k$. For $k=r$, a normalized $I$ given by (a) has
$\delta_Y(I)\neq0$, so it is injective, a permutation. Let $k<r$, and let $w$ be given for $k+1$. By (a) and (b), applied to
$(w(1),\dots,w(k+1))$ and a normalized map for $k$, there is an injective $f:\{1,\dots,k\}\to\{1,\dots,k+1\}$ with
$(w(f(1)),\dots,w(f(k)))$ normalized. Extend $f$ to a permutation $g$ of $\{1,\dots,k+1\}$, and $g$ to $\{1,\dots,r\}$ by the
identity, and put $w'=w\circ g$. Then $(w'(1),\dots,w'(k))$ is normalized, and for $j\ge k+1$, $(w'(1),\dots,w'(j))$ is a reordering
of $(w(1),\dots,w(j))$, so it is normalized.
\end{proof}
% Lean: exists_cell, from T.exists_perm_normalized_minors (exists_perm_normAt: exists_perm_top, perm_step with
% exists_perm_castSucc and normAt_comp_perm), T.normalized_minor_extends (normAt_exchange: normAt_mul_inv_mem through
% mul_adjugate_apply_eq_kMinor, det_mul_expand and exists_unit_of_sum_eq_one) and exists_normAt (kMinor_mul_expand,
% kMinor_mul_div_mem, exists_unit_of_sum_eq_one). The Lean's NormAt O Y k hk I is "normalized"; kMinor is delta.

\begin{lemma}\label{cv-origin-reference-flag-bundle-cell}
Let $\Omega$, $S$, the $r_s$, $\Omega'$ and the $U_s$ be as in Lemma~\ref{cv-origin-reference-flag-bundle-coords}, let $R\subset\Omega$
be a subring and $O\subset\Omega'$ a local subring that contains $R$. For $s\in S$ let $M_s$ be an $n_s\times r_s$ and $N_s$ an
$r_s\times n_s$ matrix with entries in $R$ such that $N_sM_s=1$, let $T_s$ be an invertible $r_s\times r_s$ matrix over $\Omega$, and
put $\Phi_s=M_sT_sU_s$. For $1\le k\le r_s$ let $\mu_{s,k}$ be normalized for $\Phi_s$ with respect to $O$
(Lemma~\ref{cv-origin-reference-flag-bundle-normalized}), and let $A$ be the subring of $\Omega'$ generated by $R$ and the quotients
$\delta_{\Phi_s}(\mu')/\delta_{\Phi_s}(\mu_{s,k})$, for all $s$, $k$ and $\mu':\{1,\dots,k\}\to\{1,\dots,n_s\}$. Then there is a nonzero
$t\in A$ with $t,t^{-1}\in O$ such that $A[1/t]$ is the localization at one element of a ring of polynomials over $R$ in finitely
many variables. In particular, the inclusion of $R$ in $A[1/t]$ is smooth.
\end{lemma}
\leanref{Development/AutomorphicSpectrum/Proofs/GeometryProofs/FrameMatricesRegularGenericStatement.lean}{GlobalTheory.T.frame_matrices_regular_generic}
\leanref{Development/AutomorphicSpectrum/Proofs/GeometryProofs/FrameMatricesRegularGeneric.lean}{GlobalTheory.CompanionTheory.Ar.FlagBun.Sm.Cell.big_cell_smooth}
\begin{proof}
Put $Y_s=T_sU_s$, so that $\Phi_s=M_sY_s$ and $Y_s=N_s\Phi_s$; we use the expansion $(*)$ of the proof of
Lemma~\ref{cv-origin-reference-flag-bundle-normalized} for these products.

\emph{An order of the rows.} For each $s$, Lemma~\ref{cv-origin-reference-flag-bundle-normalized}, applied to $O$, $M_s$, $N_s$, $Y_s$
and the $\mu_{s,k}$, gives a permutation $w_s$ such that $(w_s(1),\dots,w_s(k))$ is normalized for $Y_s$ for every $k$. Let $Y'_s$ be
the matrix whose row $i$ is the row $w_s(i)$ of $Y_s$, and $D_{s,k}=\delta_{Y_s}(w_s(1),\dots,w_s(k))$, the leading $k\times k$
minor of $Y'_s$, which is nonzero.

\emph{An LU decomposition.} As the leading minors of $Y'_s$ are nonzero, there are an upper triangular matrix $C_s$ and a lower
triangular matrix $V_s$ with $1$ on the diagonal such that $Y'_sC_s=V_s$: the column $k$ of $C_s$, with zeros below the row $k$,
solves the system with the leading $k\times k$ block of $Y'_s$ and the $k$-th unit vector as right side. As $C_s$ is upper triangular,
$\delta_{Y'_sC_s}(I)=\delta_{Y'_s}(I)\det(C_s)_{\le k}$ for a map $I$ on $\{1,\dots,k\}$, where $(C_s)_{\le k}$ is the leading
$k\times k$ block; for $I=(1,\dots,k)$ this gives $1=D_{s,k}\det(C_s)_{\le k}$. So $\delta_{Y'_s}(I)/D_{s,k}=\delta_{V_s}(I)$, and for
$l>k$, $(V_s)_{lk}=\delta_{V_s}(1,\dots,k-1,l)=\delta_{Y_s}(w_s(1),\dots,w_s(k-1),w_s(l))/D_{s,k}$.

\emph{The rings.} Let $v$ be the family of the entries $(V_s)_{lk}$ with $l>k$, and $B$ the subring of $\Omega'$ generated by $R$ and
$v$. The entries of $V_s$ lie in $B$, so $\delta_{Y_s}(I)/D_{s,k}\in B$ for every $I$, and by $(*)$ for $\Phi_s=M_sY_s$,
$\delta_{\Phi_s}(\mu')/D_{s,k}\in B$ for every $\mu'$. By $(*)$ for $Y_s=N_s\Phi_s$, $\delta_{Y_s}(I)/\delta_{\Phi_s}(\mu_{s,k})\in A$.
Put $t_{s,k}=D_{s,k}/\delta_{\Phi_s}(\mu_{s,k})$. Then $t_{s,k}\in A$ and $t_{s,k}^{-1}\in B$; and $t_{s,k}^{-1}\in O$, as by $(*)$ it is
a combination, with coefficients from $M_s$, of the quotients $\delta_{Y_s}(I)/D_{s,k}$, which lie in $O$ since
$(w_s(1),\dots,w_s(k))$ is normalized. The generators $\delta_{\Phi_s}(\mu')/\delta_{\Phi_s}(\mu_{s,k})=(\delta_{\Phi_s}(\mu')/D_{s,k})t_{s,k}$
of $A$ lie in the ring generated by $B$ and the $t_{s,k}$, and the generators
$(V_s)_{lk}=\bigl(\delta_{Y_s}(w_s(1),\dots,w_s(k-1),w_s(l))/\delta_{\Phi_s}(\mu_{s,k})\bigr)t_{s,k}^{-1}$ of $B$ lie in the ring generated
by $A$ and the $t_{s,k}^{-1}$. Let $t=\prod_{s,k}t_{s,k}$. Each $t_{s,k}^{-1}$ is $t^{-1}$ times a product of elements of $A$, and each
$t_{s,k}$ is $t$ times a product of elements of $B$; so $A[1/t]=B[1/t^{-1}]$. As $R$ and the generators of $A$ lie in $O$, $A\subset O$;
so $t\in O$, and $t^{-1}\in O$.

\emph{Independence.} Let $T'_s$ be $T_s$ with its rows permuted by $w_s$, so that $Y'_s=T'_sU_s$ and $U_sC_s=T'^{-1}_sV_s$, a matrix
with entries in the subfield $\Omega(v)$ of $\Omega'$ generated by $v$. As above, with $U_s$ in place of $Y'_s$, whose leading minors
are $1$, the entry $u_{s,lk}=\delta_{U_s}(1,\dots,k-1,l)$ is $\delta_{U_sC_s}(1,\dots,k-1,l)/\delta_{U_sC_s}(1,\dots,k)$, which lies in
$\Omega(v)$. So $\Omega(v)$ contains $\Omega(u)=\Omega'$, whose transcendence degree over $\Omega$ is the number of the $u_{s,lk}$,
which is the number of entries of $v$. A transcendence basis of $\Omega'$ over $\Omega$ can be taken among the entries of $v$, and it
has as many elements as $v$; so $v$ is algebraically independent over $\Omega$.

\emph{Smoothness.} As $v$ is algebraically independent over $\Omega$, the homomorphism to $\Omega'$ from the ring of polynomials over $R$
in variables indexed by the entries of $v$, sending the variables to $v$, is injective, and its image is $B$. So $A[1/t]=B[1/t^{-1}]$
is the localization at one element of a ring of polynomials over $R$. A ring of polynomials over $R$ is a standard smooth $R$-algebra
with no equations, so its localization at one element is smooth over $R$ \cite[Tag 00T7]{Sta}.
\end{proof}
% Lean: big_cell_smooth: exists_cell (Lemma cv-origin-reference-flag-bundle-normalized), T.lu_decomposition_of_minors
% (exists_lu_of_leading_minors, the columns a k = (B k)^-1 e_last), ratio_eq_kMinor and kMinor_mul_upper, kMinor_lower_lastRow,
% hyB, hPB, hyA (kMinor_mul_div_mem), tq with htqA, htqB, htqO, hAB and hBA, entries_mem with genU_indep and genU_gen,
% T.algebraic_independent_of_adjoin (transcendence bases), T.localization_closure_equal (locSub_eq_of_inverse_family) and
% T.localization_closure_smooth (evalSub_injective, evalSub_range, Algebra.Smooth.of_isLocalization_Away and the smoothness
% of MvPolynomial).

\begin{lemma}\label{cv-origin-reference-flag-bundle}
Let $\Omega$ be a field, $T$ a scheme with a morphism $\xi:\operatorname{Spec}\Omega\to T$ of whose image every point is a
specialization, with coordinates $x$ (Definition~\ref{cv-origin-segre-coordinates}), and $V$ an $\Omega$-vector space. For
$s$ in a finite set let $L_s$ be a lattice family in $V$ that is a vector bundle of rank $r_s$ with generic fibre $W_s$
(Definition~\ref{cv-origin-reference-lattices}). Then there are a field $\Omega'\supset\Omega$, a scheme $T'$ with a morphism
$\xi':\operatorname{Spec}\Omega'\to T'$ of whose image every point is a specialization, with coordinates
$x'=(x'_0,\dots,x'_N)$, a smooth surjective morphism $h:T'\to T$ with $h\circ\xi'=\xi\circ\operatorname{Spec}(\Omega'\supset\Omega)$,
and families $b_s$ of $r_s$ vectors of $\Omega'\otimes_\Omega V$ that form a generic flag of degree $1$ along $j:e\mapsto1\otimes e$
over the $W_s$, such that each pullback $h^*_jL_s$ is a flag bundle on $T'$ adapted to $b_s$.
\end{lemma}
\leanref{Development/AutomorphicSpectrum/Proofs/GeometryProofs/ExistsModelVectorBundlesStatement.lean}{GlobalTheory.T.exists_model_vector_bundles}
\begin{proof}
Let $\eta$ be the image of $\xi$, and $S$ the finite set of the indices $s$.

\emph{Bases and twisted forms.} Let $s\in S$. As $L_s$ is a vector bundle with generic fibre $W_s$, it is generated near $\eta$ by
a basis $c_s=(c_{s,1},\dots,c_{s,r_s})$ of $W_s$. Lemma~\ref{cv-origin-reference-flag-bundle-twist}, applied to $L_s$, gives an
integer $m_s\ge0$ and linear forms $\varphi_{s,1},\dots,\varphi_{s,n_s}$ on $V$, and by
Lemma~\ref{cv-origin-reference-flag-bundle-inverse}(1) the $n_s\times r_s$ matrix $F^0_s=(\varphi_{s,j}(c_{s,l}))_{j,l}$ has a left
inverse.

\emph{The generic flag.} Let $\Omega'$, the $u_{s,lk}$ and the matrices $U_s$ be as in
Lemma~\ref{cv-origin-reference-flag-bundle-coords}, let $\iota$ be the inclusion of $\Omega$ in $\Omega'$, and put
\[
b_{s,k}=\sum_l(U_s)_{lk}(1\otimes c_{s,l})=1\otimes c_{s,k}+\sum_{l>k}u_{s,lk}(1\otimes c_{s,l})\in\Omega'\otimes_\Omega V.
\]
The map $j:e\mapsto1\otimes e$ has $j(ae)=aj(e)$ for $a\in\Omega$, and it carries linearly independent families to linearly
independent families, as a basis of $V$ gives a basis of $\Omega'\otimes_\Omega V$. Each family $b_s$ is linearly independent, as
$U_s$ has determinant $1$. The $u_{s,lk}$ are algebraically independent over $\Omega$ and generate $\Omega'$ over $\Omega$. So the
$b_s$ form a generic flag of degree $1$ along $j$ over the $W_s$ (Definition~\ref{cv-origin-reference-lattices}(4)), with the
bases $c_s$ and the elements $u_{s,lk}$. Extending each $\varphi_{s,j}$ to an $\Omega'$-linear form on $\Omega'\otimes_\Omega V$,
the matrix $F_s=F^0_sU_s$ of Lemma~\ref{cv-origin-reference-flag-bundle-coords} is $(\varphi_{s,j}(b_{s,k}))_{j,k}$.

\emph{The model.} The point $\eta$ lies in a chart $c_0$ of $x$. Every element of $\Omega$ is regular at $\eta$, so it is a quotient
of two elements of $\mathbf Z[x_j/x_{c_0}]$, by Lemma~\ref{cv-origin-snc-charts}(1) applied to $x$ as a point with one block. Let
$y=(y_{(c,\mu)})$ be the family of Lemma~\ref{cv-origin-reference-flag-bundle-coords} for $x$, $c_0$ and the $F^0_s$, with its
family $\rho$. Lemmas~\ref{cv-origin-segre-model} and~\ref{cv-origin-reference-flag-bundle-valuative}, applied to $y$ and the index
$(c_0,\rho)$, give a scheme $T'$ with a morphism $\xi':\operatorname{Spec}\Omega'\to T'$ of whose image every point is a
specialization, such that $T'$ has the coordinates $y$ and the morphism $T'\to\operatorname{Spec}\mathbf Z$ is proper. Numbering the indices of $y$ by $0,\dots,N$ gives the coordinates $x'$.

\emph{Charts.} Let $P'$ be a point of $T'$, in the chart $(c_1,\mu)$ of $y$. Then $y_{(c_1,\mu)}\neq0$, so $x_{c_1}\neq0$ and every
$\Delta_{s,k}(\mu_{s,k})$ is nonzero. The ratios $\iota(x_j)/\iota(x_{c_1})=y_{(j,\mu)}/y_{(c_1,\mu)}$ are regular at $P'$, so $P'$
lies in the chart $c_1$ of the family $\iota(x)$. For $s$, $k\le r_s$, $l\le k$ and $j\le n_s$, let $\mu'$ be $\mu$ with $\mu_{s,k}$
replaced by the map $\mu_{s,k}$ with its value at $l$ replaced by $j$; then
$y_{(c_1,\mu')}/y_{(c_1,\mu)}=\Delta_{s,k}(\mu'_{s,k})/\Delta_{s,k}(\mu_{s,k})$, which is therefore regular at $P'$.

\emph{The morphism.} By the previous paragraph every point of $T'$ lies in a chart of $\iota(x)$, and $T'$, which has coordinates,
satisfies condition (1) of Definition~\ref{cv-origin-segre-coordinates}. So Lemma~\ref{cv-origin-coords-join-map}(2) gives a
morphism $h:T'\to T$ with $h\circ\xi'=\xi\circ\operatorname{Spec}\iota$, such that a point $P'$ lies in the chart $c$ of $\iota(x)$
if and only if $h(P')$ lies in the chart $c$ of $x$. By Lemma~\ref{cv-origin-reference-flag-bundle-onto}, $h$ is surjective.

\emph{Smoothness.} Let $P'$ be a point of $T'$, in the chart $(c_2,\mu)$ of $y$. By the paragraphs on charts and on the morphism,
every point of $T'$ in this chart lies in the chart $c_2$ of $\iota(x)$, and its image under $h$ in the chart $c_2$ of $x$; in
particular $P=h(P')$ lies in the chart $c_2$ of $x$. Let $A_{c_2}=\mathbf Z[x_j/x_{c_2}]$, and $A'$ the ring generated by the ratios
$y_{(c,\mu')}/y_{(c_2,\mu)}$. For each $s$, Lemma~\ref{cv-origin-reference-flag-bundle-frame}, applied at $P$ to $L_s$, the
$\varphi_{s,j}$ and $c_s$, gives $g_s\in A_{c_2}$, a unit at $P$, and matrices $T_s$, $M_s$ and $N_s$ with $N_sM_s=1$ and
$F^0_s=x_{c_2}^{m_s}M_sT_s$. Let $g=\prod_sg_s$ and $R=A_{c_2}[1/g]$; then $M_s$ and $N_s$ have entries in $R$, and $R\subset\mathcal O_P$,
so $R\subset\mathcal O_{P'}$ by Lemma~\ref{cv-origin-coords-join-map}(1). With $\Phi_s=M_sT_sU_s$, $F_s=x_{c_2}^{m_s}\Phi_s$, so
$\Delta_{s,k}(\mu')=x_{c_2}^{km_s}\delta_{\Phi_s}(\mu')$ in the notation of Lemma~\ref{cv-origin-reference-flag-bundle-normalized}, and
\[
\frac{y_{(c,\mu')}}{y_{(c_2,\mu)}}=\frac{x_c}{x_{c_2}}\prod_{s,k}\frac{\delta_{\Phi_s}(\mu'_{s,k})}{\delta_{\Phi_s}(\mu_{s,k})}.
\]
Replacing in $\mu$ only the map $\mu_{s,k}$ by $\mu'$, and taking $c=c_2$, shows that $\delta_{\Phi_s}(\mu')/\delta_{\Phi_s}(\mu_{s,k})$ lies in
$A'$, hence in $\mathcal O_{P'}$, for every $\mu'$: each $\mu_{s,k}$ is normalized for $\Phi_s$ with respect to $\mathcal O_{P'}$. Let $A$ be
the ring of Lemma~\ref{cv-origin-reference-flag-bundle-cell} for $R$, the $M_s$, $N_s$, $T_s$ and $\mu_{s,k}$. Then $A=A'[1/g]$: $A'$
contains the $x_c/x_{c_2}=y_{(c,\mu)}/y_{(c_2,\mu)}$, so $A_{c_2}$ and $g$, and the generators of $A$; conversely the generators of $A'$
are products of elements of $A_{c_2}$ and of generators of $A$, and $g^{-1}\in R$. Lemma~\ref{cv-origin-reference-flag-bundle-cell}, with
$O=\mathcal O_{P'}$, gives a nonzero $t\in A$ with $t$ and $t^{-1}$ regular at $P'$ and the inclusion of $R$ in $A[1/t]$ smooth. Write
$t=a/g^e$ with $a\in A'$. Let $U_1$ be the set of the points of the chart $c_2$ of $x$ at which $g$ is a unit, and $W$ the set of the
points of the chart $(c_2,\mu)$ of $y$ at which $ga$ is a unit. By Lemma~\ref{cv-origin-reference-descent-roots-piece}, applied to $T$
and $x$ and to $T'$ and $y$, these are standard open subsets of affine schemes, hence affine, and by
Lemma~\ref{cv-origin-segre-regular}(2) their rings map isomorphically onto $R$ and onto $A'[1/ga]=A[1/t]$. The point $P$ lies in $U_1$
and $P'$ in $W$, as $g$ is a unit at $P$, hence at $P'$, and $t$ is a unit at $P'$. For $Q'\in W$, $h(Q')$ lies in the chart $c_2$ of
$x$, and $g$, a unit at $Q'$, is a unit at $h(Q')$ by Lemma~\ref{cv-origin-coords-join-map}(1); so $h(W)\subset U_1$. As
$h\circ\xi'=\xi\circ\operatorname{Spec}\iota$, the ring map of $h:W\to U_1$ is, through these isomorphisms, the inclusion of $R$ in
$A[1/t]$, which is smooth. So $h$ is smooth at $P'$ \cite[Tag 01V5]{Sta}, and as $P'$ was arbitrary, $h$ is smooth.
% Lean: T.flag_frames_generic_model_right = FlagBun.Sm.flag_smooth: T.dominating_pieces_smooth (smooth_of_pieces, through
% HasRingHomProperty.iff_exists_appLE and val_appLE_of_dominates), hover for the charts, T.twisted_frame_in_chart (Lemma
% cv-origin-reference-flag-bundle-frame) for each s, g0 and R = locSub (chartSub x c0) g0, flagCoord_eq and prod_update_div
% for hbetaJ, T.frame_matrices_regular_generic (Lemma cv-origin-reference-flag-bundle-cell), T.chart_frame_matrix_regular
% (chartRing_eq_ratioRing: A'[1/g] = ratioRing R (flagB M T) J), and T.piece_localization_at_element three times (U1, W1 and
% W2); the book reaches W2 in one step, as the points of the chart at which g a is a unit.

\emph{The flag bundles.} Let $s\in S$, write $r=r_s$, $n=n_s$, $m=m_s$, $\varphi_j=\varphi_{s,j}$ and $c_l=c_{s,l}$, and let $P'$ be a
point of $T'$, in the chart $(c_1,\mu)$ of $y$. By the paragraphs on charts and on the morphism, $P=h(P')$ lies in the chart $c_1$ of
$x$. As $L_s$
is a vector bundle, $P$ has an open neighbourhood $U$ with a basis $\beta_1,\dots,\beta_r$ of $W_s$ such that
$L_s(Q)=\sum_a\mathcal O_Q\beta_a$ for $Q\in U$. Lemma~\ref{cv-origin-reference-flag-bundle-inverse}(2) gives the matrix
$A_0=(\varphi_j(\beta_a)/x_{c_1}^m)_{j,a}$ and a matrix $B_0$, both with entries in $\mathcal O_P$, such that $B_0A_0=1$; by
Lemma~\ref{cv-origin-coords-join-map}(1) their entries are regular at $P'$. Write $c_l=\sum_aC_{la}\beta_a$ with $C_{la}\in\Omega$,
and let $G$ be the $r\times r$ matrix over $\Omega'$ with the entries $G_{ak}=\sum_lC_{la}(U_s)_{lk}$, so that
$b_{s,k}=\sum_aG_{ak}(1\otimes\beta_a)$. Then $\varphi_j(b_{s,k})=x_{c_1}^m(A_0G)_{jk}$, that is, $F_s=x_{c_1}^mA_0G$. For $1\le k\le r$,
the minor $\Delta_{s,k}(\mu_{s,k})$ of $F_s$ is nonzero, and its quotients into the minors with one of the rows $\mu_{s,k}(l)$
replaced are regular at $P'$, by the paragraph on charts. So Lemma~\ref{cv-origin-reference-flag-bundle-minors}, applied with
$R=\mathcal O_{P'}$, $A_0$, $B_0$, $G$, $a=x_{c_1}^m$ and $\mu_{s,k}$, gives the vectors $D_l$ and $E_l$ for each $k$, and
Lemma~\ref{cv-origin-reference-flag-bundle-adapted}, applied to the local ring $\mathcal O_{P'}$ and the columns of $G$, gives an
$r\times r$ matrix $P_s$ with entries in $\mathcal O_{P'}$ and determinant a unit at $P'$, whose first $i$ columns span the same
subspace as the first $i$ columns of $G$, for every $i$. Put $b'_k=\sum_a(P_s)_{ak}(1\otimes\beta_a)$. The linear map
$v\mapsto\sum_av_a(1\otimes\beta_a)$ carries the first $i$ columns of $G$ to $b_{s,1},\dots,b_{s,i}$ and those of $P_s$ to
$b'_1,\dots,b'_i$, so $\langle b'_1,\dots,b'_i\rangle=\langle b_{s,1},\dots,b_{s,i}\rangle$ for every $i$. By
Lemma~\ref{cv-origin-segre-regular}(1), $P'$ has an open neighbourhood $U'\subset h^{-1}(U)$ at every point of which the entries of
$P_s$ and the inverse of $\det P_s$ are regular. Let $Q'\in U'$. Then $h(Q')\in U$, so $(h^*_jL_s)(Q')$ is generated over
$\mathcal O_{Q'}$ by the linearly independent vectors $1\otimes\beta_a$, as $\iota(\mathcal O_{h(Q')})\subset\mathcal O_{Q'}$. As $P_s$
and its inverse $\det(P_s)^{-1}\operatorname{adj}(P_s)$ have entries in $\mathcal O_{Q'}$, the $b'_k$ are linearly independent and
generate the same $\mathcal O_{Q'}$-module. So $h^*_jL_s$ is a flag bundle on $T'$ adapted to $b_s$.
\end{proof}
% Lean: T.exists_model_vector_bundles = exists_flagBundle: the frames c s at genPt, T.exists_twisted_generators (Lemma
% cv-origin-reference-flag-bundle-twist), T.frame_via_twisted_generators (Lemma cv-origin-reference-flag-bundle-inverse(1)),
% T.frame_matrix_chart_ratios with T.function_field_ratio_generators (Lemma cv-origin-reference-flag-bundle-coords),
% T.exists_proper_integral_model on GenField Omega r and flagCoord (Lemmas cv-origin-segre-model and
% cv-origin-reference-flag-bundle-valuative), T.dominates_of_coords_regular (Lemma cv-origin-coords-join-map(2)) with the
% charts from the cover of flagCoord (inChart_flagCoord_ratios), T.proper_dominating_surjective (Lemma
% cv-origin-reference-flag-bundle-onto), T.flag_frames_generic_model_right (the paragraph on smoothness), T.generic_flag_over_generic_field
% (genB_isGenericFlag, with genField_adjoin_eq_top and genB_linearIndependent) and T.flag_frames_generic_model_left
% (Flags.pull_isFlagVB_genB: T.flag_chart_generic_field, T.frame_coeffs_regular_in_chart, flagMat = a (A G),
% T.minor_ratio_regular_of_invertible, T.flag_basis_change_integral, FlagPull.pull_isFrameAt and
% T.frame_change_by_unit_matrix), and coordTie_comp_equiv for the numbering of the coordinates.

\begin{lemma}\label{cv-origin-reference-flag-linear}
Let $\Omega_0\subset\Omega$ be fields.
\begin{enumerate}
\item Let $V$ be an $\Omega_0$-vector space, $V'$ an $\Omega$-vector space and $j:V\to V'$ an additive map with $j(ae)=aj(e)$
  for $a\in\Omega_0$, which carries linearly independent families to linearly independent families and whose image spans
  $V'$. Then there is an isomorphism of $\Omega$-vector spaces $\Phi:\Omega\otimes_{\Omega_0}V\to V'$ with
  $\Phi(1\otimes e)=j(e)$.
\item Let $a:V_0\to V$ be an isomorphism of $\Omega_0$-vector spaces, $j:V\to V'$ an additive map to an $\Omega$-vector space,
  $\Phi:V'\to V''$ an isomorphism of $\Omega$-vector spaces, and for $s$ in a finite set let $W_s\subset V$ be a subspace and
  $W^0_s=a^{-1}(W_s)$. If families $b_s$ in $V'$ form a generic flag of degree $d$ along $j$ over the $W_s$
  (Definition~\ref{cv-origin-reference-lattices}(4)), then the families $\Phi(b_s)$ form a generic flag of degree $d$ along
  $\Phi\circ j\circ a$ over the $W^0_s$.
\item Let $T$ and $T'$ be schemes with morphisms $\xi:\operatorname{Spec}\Omega_0\to T$ and $\xi':\operatorname{Spec}\Omega\to T'$,
  $f:T'\to T$ a morphism, $L$ a lattice family in an $\Omega_0$-vector space $V$, $j:V\to V'$ an additive map to an
  $\Omega$-vector space, $\Phi:V'\to V''$ an isomorphism of $\Omega$-vector spaces, $S$ a set of points of $T'$ that specialize
  the image of $\xi'$, and $b$ a finite family in $V'$. If $f^*_jL$ is a flag bundle on $S$ adapted to $b$, then
  $f^*_{\Phi\circ j}L$ is a flag bundle on $S$ adapted to $\Phi(b)$.
\item Let $u_{s,lk}\in\Omega$, for $s$ in a set and $1\le k,l\le r_s$, with $u_{s,kk}=1$ and $u_{s,lk}=0$ for $l<k$. If
  $\Omega$ has degree $1$ over the field $\Omega_0(u)$ generated by all of them, then $\Omega$ is generated over $\Omega_0$ by
  the $u_{s,lk}$ with $l>k$.
\end{enumerate}
\end{lemma}
\leanref{Development/AutomorphicSpectrum/Proofs/GeometryProofs/TensorEquivOfBasisExtensionStatement.lean}{GlobalTheory.T.tensor_equiv_of_basis_extension}
\leanref{Development/AutomorphicSpectrum/Proofs/GeometryProofs/GenericFlagUnderEquivStatement.lean}{GlobalTheory.T.generic_flag_under_equiv}
\leanref{Development/AutomorphicSpectrum/Proofs/GeometryProofs/FlagBundleUnderEquivStatement.lean}{GlobalTheory.T.flag_bundle_under_equiv}
\leanref{Development/AutomorphicSpectrum/Proofs/GeometryProofs/UnipotentEntriesFieldDegreeStatement.lean}{GlobalTheory.T.unipotent_entries_field_degree}
\begin{proof}
(1) Let $\Phi$ be the $\Omega$-linear map with $\Phi(c\otimes e)=cj(e)$. It is surjective, as the image of $j$ spans $V'$.
For a basis $(e_i)$ of $V$, the $1\otimes e_i$ form a basis of $\Omega\otimes_{\Omega_0}V$, and their images $j(e_i)$ are
linearly independent, as every finite subfamily of the $e_i$ is carried to a linearly independent family. So $\Phi$ is
injective.

(2) Let $c_s$ and $u_{s,lk}$ be as in the definition for the $b_s$. Then the $a^{-1}(c_s)$ are bases of the $W^0_s$, the map
$\Phi\circ j\circ a$ is semilinear, it carries $a^{-1}(c_s)$ to the linearly independent family $\Phi(j(c_s))$, the $\Phi(b_s)$
are linearly independent, and $\langle\Phi(b_{s,1}),\dots,\Phi(b_{s,i})\rangle=\Phi(\langle b_{s,1},\dots,b_{s,i}\rangle)$ is
spanned by the vectors $\Phi(j(c_{s,k}))+\sum_{l>k}u_{s,lk}\Phi(j(c_{s,l}))$ with $k\le i$. The $u_{s,lk}$ are unchanged.

(3) Let $P'\in S$, with $U$ and linearly independent $b'$ as in the definition for $f^*_jL$. Then $\Phi(b')$ is linearly
independent, $\langle\Phi(b'_1),\dots,\Phi(b'_i)\rangle=\Phi(\langle b_1,\dots,b_i\rangle)$ for every $i$, and for $Q'\in S\cap U$
the module $(f^*_{\Phi\circ j}L)(Q')$, generated by $\Phi(j(L(f(Q'))))$, is $\Phi((f^*_jL)(Q'))=\sum_i\mathcal O_{Q'}\Phi(b'_i)$.

(4) As the degree is $1$, $\Omega_0(u)=\Omega$, and the generators $u_{s,lk}$ with $l\le k$ are $0$ or $1$.
\end{proof}
% Lean: FlagStep.exists_tensor_equiv_of_semilinear (liftBaseChange, Algebra.TensorProduct.basis), FlagStep.Transport.
% isGenericFlag_transport (flagStep_equiv), FlagStep.isFlagVBOn_pull_map (flagStep_linearEquiv_image) and
% FlagStep.AdjoinLower.adjoin_lower_eq_top.

\begin{lemma}\label{cv-origin-reference-flag-coords}
Let $\Omega_0\subset\Omega$ be fields, $T$ a scheme with a morphism $\xi:\operatorname{Spec}\Omega_0\to T$ of whose image every
point is a specialization, covered by the affine open subschemes of the charts of a point $y\in\Omega_0^{\sqcup_{k\in K}I_k}$ as
in Definition~\ref{cv-origin-hodge-component}(1), and $T'$ a scheme with a morphism $\xi':\operatorname{Spec}\Omega\to T'$ and a
morphism $h:T'\to T$ with $h\circ\xi'=\xi\circ\operatorname{Spec}(\Omega\supset\Omega_0)$. Let $x'=(x'_0,\dots,x'_N)\in\Omega^{N+1}$,
and let $y'$ be the point with the blocks of $y$ and the block $x'$.
\begin{enumerate}
\item If $T'$ has the coordinates $x'$ (Definition~\ref{cv-origin-segre-coordinates}), then $T'$ is covered by the affine open
  subschemes of the charts of $y'$.
\item Suppose that the local ring of $T'$ at the image of $\xi'$ maps bijectively onto $\Omega$ and those at its
  specializations injectively, and that every point of $T'$ specializes the image of $\xi'$. If $w$ are coordinates for $y'$ of
  a complex point $z$ of $T'$, then $(w_{k,i})_{k\in K}$ are coordinates for $y$ of $h\circ z$.
\end{enumerate}
\end{lemma}
\leanref{Development/AutomorphicSpectrum/Proofs/GeometryProofs/CoordsAddedProjectiveStatement.lean}{GlobalTheory.T.coords_added_projective}
\leanref{Development/AutomorphicSpectrum/Proofs/GeometryProofs/ComplexCoordsAddedStatement.lean}{GlobalTheory.T.complex_coords_added}
\begin{proof}
(1) Let $P'$ be a point of $T'$ that specializes the image of $\xi'$. Its image $h(P')$ lies in a chart $c$ of $y$, so for
each block $k$, $y_{k,c_k}\neq0$ and every $y_{k,i}/y_{k,c_k}$, a quotient of two Segre ratios, is regular at $h(P')$. By
Lemma~\ref{cv-origin-model-pullback}(1), these ratios are regular at $P'$: $P'$ lies in the chart $c_k$ of each block $k$ of
$y$. Lemma~\ref{cv-origin-reference-segre}, applied to $T'$, with $x'$ as a point with one block, and to the blocks of $y$,
gives (1).

(2) This is Lemma~\ref{cv-origin-model-pullback}(2), applied to $h$ and to the blocks of $y$ in $y'$ with the identity maps.
\end{proof}
% Lean: (1) is FlagStep.TieBlocks.tie_blocks (CoordBase.coordTie_reindex for the one block, RegAtDom.regAt_dominates with
% block_ratio_mem, and T.coords_join_blocks, Lemma cv-origin-reference-segre). (2) is FlagStep.CoordsRestrict.coordsAt_restrict,
% which repeats for the blocks of y the argument of T.dominating_coords_reindexed (Lemma cv-origin-model-pullback(2)).

\begin{lemma}\label{cv-origin-reference-flag-snc-chart}
Let $\Omega_0\subset\Omega$ be fields, let $X$ and $X'$ be schemes with morphisms $\xi:\operatorname{Spec}\Omega_0\to X$ and
$\xi':\operatorname{Spec}\Omega\to X'$ of whose images every point is a specialization, and let $x=(x_j)_{j\in J}$ in $\Omega_0$ and
$x'=(x'_l)_{l\in L}$ in $\Omega$ be finite families such that $X$ has the coordinates $x$ and $X'$ the coordinates $x'$
(Definition~\ref{cv-origin-segre-coordinates}). Let $h:X'\to X$ be a smooth morphism with
$h\circ\xi'=\xi\circ\operatorname{Spec}(\Omega\supset\Omega_0)$, and let $c$ and $c'$ be indices with $x_c\neq0$ and $x'_{c'}\neq0$
such that every $x_j/x_c$ lies in $A'_{c'}=\mathbf Z[x'_l/x'_{c'}:l\in L]$. Then $A_c=\mathbf Z[x_j/x_c:j\in J]$ is contained in
$A'_{c'}$, and the inclusion $A_c\to A'_{c'}$ is a smooth ring map.
\end{lemma}
\leanref{Development/AutomorphicSpectrum/Proofs/GeometryProofs/DominantSmoothMapCoordsStatement.lean}{GlobalTheory.T.dominant_smooth_map_coords}
\begin{proof}
Let $X_c$ and $X'_{c'}$ be the affine open subschemes of the charts $c$ and $c'$. First, $h$ maps $X'_{c'}$ into $X_c$. Let
$P'\in X'_{c'}$. By Lemma~\ref{cv-origin-segre-regular}(3), $h(P')$ specializes the image of $\xi$, so it lies in a chart $d$ of $x$,
and $x_c/x_d$ is regular at $h(P')$. Its inverse $x_d/x_c$ lies in $A'_{c'}$, so it is regular at $P'$. By
Lemma~\ref{cv-origin-model-pullback}(1), $x_c/x_d$ is then a unit at $h(P')$, and every $x_j/x_c=(x_j/x_d)(x_c/x_d)^{-1}$ is regular at
$h(P')$: $h(P')$ lies in $X_c$.

So $h$ restricts to a morphism $X'_{c'}\to X_c$ of affine schemes, which is smooth, and the ring map
$\Gamma(X_c,\mathcal O_X)\to\Gamma(X'_{c'},\mathcal O_{X'})$ is smooth \cite[Tag 01V6]{Sta}. By Lemma~\ref{cv-origin-segre-regular}(2)
and condition (4) of Definition~\ref{cv-origin-segre-coordinates}, the maps of these rings to $\Omega_0$ and to $\Omega$ through the
local rings at the images of $\xi$ and $\xi'$ are injective, with images $A_c$ and $A'_{c'}$. As
$h\circ\xi'=\xi\circ\operatorname{Spec}(\Omega\supset\Omega_0)$, the ring map is compatible with these embeddings and the inclusion
$\Omega_0\subset\Omega$. So $A_c\subset A'_{c'}$, and the inclusion is smooth.
\end{proof}
% Lean: ChartSmooth.smooth_chartRing: T.affine for U and U', inChart_map (RegAtDom.regAt_dominates, the hypothesis hi for the
% inverse ratio) for hle, exists_secEquiv (mem_adjoin_iff, secVal_injective) for the two rings, Smooth.smooth_appLE, and
% secVal_app with secVal_res for the compatibility. The Lean states it for any field embedding i.

\begin{lemma}\label{cv-origin-reference-flag-snc}
Let $\Omega_0\subset\Omega$ be fields of characteristic zero, let $T$, $\xi$, $y$, $T'$, $\xi'$, $h$, $x'$ and $y'$ be as in
Lemma~\ref{cv-origin-reference-flag-coords}, and suppose that every point of $T'$ specializes the image of $\xi'$, that $T'$ is
covered by the affine open subschemes of the charts of $y'$, and that $h$ is smooth. Let $u_1,\dots,u_n\in\Omega$ be
algebraically independent over $\Omega_0$ and generate $\Omega$ over $\Omega_0$. If an ideal $J$ of $\mathbf Z[x_{k,i}]$ satisfies
condition (3) of Definition~\ref{cv-origin-reference} for $y$ with an integer $e$ in place of $D$, then its image under the
renaming of the blocks of $y$ into those of $y'$ satisfies it for $y'$ with $e+n$.
\end{lemma}
\leanref{Development/AutomorphicSpectrum/Proofs/GeometryProofs/SncModelAddedCoordsStatement.lean}{GlobalTheory.T.snc_model_added_coords}
\begin{proof}
Let $J'$ be the image of $J$. For a chart $\kappa$ of $y$ write $A^{\mathbf Z}_\kappa=\mathbf Z[y_{k,i}/y_{k,\kappa_k}]$ and
$A_\kappa=\mathbf Q[y_{k,i}/y_{k,\kappa_k}]$, and likewise $A'^{\mathbf Z}_{\kappa'}$ and $A'_{\kappa'}$ for a chart $\kappa'$ of $y'$.
Every element of $\Omega_0$ is a quotient of two elements of $A^{\mathbf Z}_\kappa$, by Lemma~\ref{cv-origin-snc-charts}(1) at the image
of $\xi$, where every element is regular; likewise every element of $\Omega$ is a quotient of two elements of $A'^{\mathbf Z}_{\kappa'}$.

\emph{The differentials of $\Omega$.} The field $\Omega$ is the field of fractions of the polynomial ring $\Omega_0[u_1,\dots,u_n]$,
so $\Omega_{\Omega/\Omega_0}$ is $\Omega\otimes\Omega_{\Omega_0[u]/\Omega_0}$ \cite[Tag 00RT]{Sta}, and it has the basis
$du_1,\dots,du_n$ \cite[Tag 00RX]{Sta}.

\emph{A chart.} Let $\kappa'$ be a chart of $y'$, and $\kappa$ its restriction to the blocks of $y$; then $y_{k,\kappa_k}=y'_{k',\kappa'_{k'}}\neq0$,
so $\kappa$ is a chart of $y$, and $A_\kappa\subset A'_{\kappa'}$. The ratios of the Segre products of $y$ at $\kappa$ are products of
the ratios $y_{k,i}/y_{k,\kappa_k}$, which lie in $A'^{\mathbf Z}_{\kappa'}$. So Lemma~\ref{cv-origin-reference-flag-snc-chart}, applied
to $h$ and the Segre products of $y$ and of $y'$, shows that the inclusion $A^{\mathbf Z}_\kappa\to A'^{\mathbf Z}_{\kappa'}$ is smooth.
Condition (3) for $y$ at $\kappa$ gives $f_l,a_l\in A_\kappa$ with $\sum_la_lf_l=1$, and for $f_l\neq0$ elements
$z_{l,1},\dots,z_{l,e}$ of $A_\kappa[1/f_l]$ and an integer $n_l\le e$.

\emph{A local statement.} Let $\mathfrak P$ be a prime of $A'_{\kappa'}$. Then there is $t\in A'_{\kappa'}$ outside $\mathfrak P$ such
that $S'=A'_{\kappa'}[1/t]$ is smooth over $\mathbf Q$ and has elements $z'_1,\dots,z'_{e+n}$, whose differentials form a basis of
$\Omega_{S'/\mathbf Q}$, and an integer $n'\le e+n$, such that a prime of $S'$ contains the values of $J'$ if and only if it contains
one of $z'_1,\dots,z'_{n'}$. Indeed, as $\sum_la_lf_l=1$ in $A'_{\kappa'}$, some $f_l$ is outside $\mathfrak P$; it is nonzero. Put
$f=f_l$, $S=A_\kappa[1/f]$, with its smooth structure over $\mathbf Q$, $z_i=z_{l,i}$ and $n_0=n_l$, and $T_1=A'_{\kappa'}[1/f]$.
\begin{enumerate}
\item[(a)] Some integer $r\ge1$ has $rf\in A^{\mathbf Z}_\kappa$. Then $S$ is the localization of $A^{\mathbf Z}_\kappa$ at the products
  of nonzero integers and powers of $rf$, and $T_1$ is the localization of $A'^{\mathbf Z}_{\kappa'}$ at the same products. So
  $T_1=S\otimes_{A^{\mathbf Z}_\kappa}A'^{\mathbf Z}_{\kappa'}$ \cite[Tag 00DK]{Sta}, and $S\to T_1$ is smooth, by base change
  \cite[Tag 00T4]{Sta}.
\item[(b)] The ideal $\mathfrak PT_1$ is prime, as $f\notin\mathfrak P$. There is $g\in T_1$ outside $\mathfrak PT_1$ such that
  $T_2=T_1[1/g]$ is standard smooth over $S$ \cite[Tag 00TA]{Sta}, so $\Omega_{T_2/S}$ is free on the differentials of some
  $w_1,\dots,w_m\in T_2$ \cite[Tag 00T7]{Sta}, and $T_2$ is smooth over $\mathbf Q$, as $S$ is \cite[Tag 00TD]{Sta}. The fields of
  fractions of $S$ and $T_2$ are $\Omega_0$ and $\Omega$. So $\Omega_{\Omega/S}=\Omega\otimes_{T_2}\Omega_{T_2/S}$ is free of rank $m$,
  and it is $\Omega_{\Omega/\Omega_0}$, as the nonzero elements of $S$ are units in $\Omega$ \cite[Tag 00RT]{Sta}. Hence $m=n$.
\item[(c)] As $S\to T_2$ is smooth, it is formally smooth \cite[Tag 00TN]{Sta}, and
  $0\to T_2\otimes_S\Omega_{S/\mathbf Q}\to\Omega_{T_2/\mathbf Q}\to\Omega_{T_2/S}\to0$ is split exact \cite[Tag 031K]{Sta}. So the
  differentials of $z_1,\dots,z_e,w_1,\dots,w_n$ form a basis of $\Omega_{T_2/\mathbf Q}$. Let $z'$ be this family and $n'=n_0$.
\item[(d)] Write $g=g_0/f^k$ with $g_0\in A'_{\kappa'}$. Then $g_0\notin\mathfrak P$, as $g\notin\mathfrak PT_1$ and $f\notin\mathfrak P$. Put
  $t=g_0f$, outside $\mathfrak P$. Then $S'=A'_{\kappa'}[1/t]=T_1[1/g]=T_2$.
\item[(e)] For $F\in J$ with image $F'\in J'$, $F'$ takes the same value at the ratios of $y'$ at $\kappa'$ as $F$ at the ratios of
  $y$ at $\kappa$, and the values of $J'$ are the combinations of these values with coefficients in $S'$. So a prime $\mathfrak q$ of
  $S'$ contains the values of $J'$ if and only if $\mathfrak q\cap S$ contains those of $J$, which by (3) for $y$ holds if and only
  if $\mathfrak q\cap S$, or equivalently $\mathfrak q$, contains one of $z_1,\dots,z_{n_0}$.
\end{enumerate}

\emph{Gluing.} The elements $t$ given by the local statement for all primes $\mathfrak P$ of $A'_{\kappa'}$ generate the unit ideal: an
ideal that contains them all lies in no maximal ideal. So $1=\sum_la'_lt_l$ for finitely many of them, with $a'_l\in A'_{\kappa'}$,
and each $t_l$ is nonzero, with the $z'$ and $n'$ of the local statement. This is condition (3) for $y'$ at $\kappa'$ with $e+n$.
\end{proof}
% Lean: SncPull.isSncModel_pull: the chart c of y from c1 (hc), T.dominant_smooth_map_coords (Lemma
% cv-origin-reference-flag-snc-chart) for phi, T.kaehler_rank_transcendence_degree (= finrank_kaehlerDifferential_of_adjoin_eq_top:
% Omega is a localization of MvPolynomial I Omega0, isBaseChange_of_formallyEtale) for n, T.fraction_field_coords_generated for
% the fields of fractions, and hloc for each prime P: the unit f l, adjoin_insert_eq, T.smooth_map_localization_compat (=
% LocalisedSmooth.smooth_localised: denomMonoid, isLocalization_adjoin, isPushout_of_isLocalization and
% RingHom.Smooth.isStableUnderBaseChange) for (a), and T.kaehler_rank_localization_bound (= LocalStep.local_step:
% exists_away_basis with Algebra.IsSmoothAt.exists_notMem_isStandardSmooth, Algebra.Smooth.comp,
% T.kaehler_finrank_fraction_field and T.kaehler_basis_tower_alt for (b) and (c); IsLocalization.Away.surj for g0 and t = g0 h1
% in (d); T.adjoin_inverse_is_localization_away; T.localization_kaehler_prime_ideal for (e)). The cover is SncPull.exists_cover.

\begin{lemma}\label{cv-origin-reference-flag-constants}
In the setting of Definition~\ref{cv-origin-hodge-component}, let $\Omega_0\subset\Omega$ be fields of characteristic zero, $T$
and $T'$ schemes with morphisms $\xi:\operatorname{Spec}\Omega_0\to T$ and $\xi':\operatorname{Spec}\Omega\to T'$ such that the local
rings at their images map bijectively onto $\Omega_0$ and $\Omega$ and those at their specializations injectively, and $h:T'\to T$
a morphism with $h\circ\xi'=\xi\circ\operatorname{Spec}(\Omega\supset\Omega_0)$. Let $G$ and $G'$ be $\mathcal O_E$-semi-abelian
schemes over $T$ and $T'$ with Lie algebras $\mathrm{Lie}$ and $\mathrm{Lie}'$ at $\xi$ and $\xi'$ and their actions of
$\mathcal O_E$, and $\Lambda$ and $\Lambda'$ the lattice families of their Lie algebras at the points
(Definition~\ref{cv-origin-reference-lattices}). Let $F_c\to\Omega_0$ be an embedding, and for a label $s$ let $\mathrm{Lie}_s$ and
$\mathrm{Lie}'_s$ be the subspaces on which each $b\in\mathcal O_E$ acts by $s(b)$, through $F_c\to\Omega_0\subset\Omega$. Let
$j:\mathrm{Lie}\to\mathrm{Lie}'$ be an additive map, and for each label $s$ let $L_s$ be a lattice family in $\mathrm{Lie}$ and $b_s$ a
finite family in $\mathrm{Lie}'$. Suppose that:
\begin{enumerate}
\item[(a)] the local rings of $T$ and of $T'$ at the points over no prime are integrally closed;
\item[(b)] $L_s(P)=\Lambda(P)\cap\mathrm{Lie}_s$ for every point $P$ of $T$ over no prime;
\item[(c)] $\Lambda'=h^*_j\Lambda$;
\item[(d)] $j(\mathrm{Lie}_s)\subset\mathrm{Lie}'_s$ for every $s$;
\item[(e)] each $h^*_jL_s$ is a flag bundle on $T'$ adapted to $b_s$.
\end{enumerate}
Then each $\Lambda'\cap\mathrm{Lie}'_s$ is a flag bundle adapted to $b_s$ on the set of points of $T'$ over no prime.
\end{lemma}
\leanref{Development/AutomorphicSpectrum/Proofs/GeometryProofs/LieModelConstantsExtensionStatement.lean}{GlobalTheory.T.lie_model_constants_extension}
\begin{proof}
Let $Q$ be a point of $T'$ over no prime, and $P=h(Q)$. Every prime is regular at $P$ and a unit at $Q$, so it is a unit at
$P$ by Lemma~\ref{cv-origin-model-pullback}(1): $P$ lies over no prime. By Lemma~\ref{cv-origin-reference-eigen}, applied to $G$
and to $G'$ with (a), $\Lambda(P)=\sum_u(\Lambda(P)\cap\mathrm{Lie}_u)=\sum_uL_u(P)$, by (b), and the subspaces $\mathrm{Lie}'_u$
are independent. For each label $u$, $(h^*_jL_u)(Q)$ is generated by the $j(e)$ with $e\in L_u(P)=\Lambda(P)\cap\mathrm{Lie}_u$,
which lie in $(h^*_j\Lambda)(Q)=\Lambda'(Q)$ and, by (d), in $\mathrm{Lie}'_u$; so
$(h^*_jL_u)(Q)\subset\Lambda'(Q)\cap\mathrm{Lie}'_u$. Conversely $\Lambda'(Q)$ is generated by the $j(e)$ with
$e\in\Lambda(P)=\sum_uL_u(P)$, so $\Lambda'(Q)\subset\sum_u(h^*_jL_u)(Q)$. Let $v\in\Lambda'(Q)\cap\mathrm{Lie}'_s$, and write
$v=\sum_uv_u$ with $v_u\in(h^*_jL_u)(Q)\subset\mathrm{Lie}'_u$. By the independence of the $\mathrm{Lie}'_u$, $v=v_s$. So
$(\Lambda'\cap\mathrm{Lie}'_s)(Q)=(h^*_jL_s)(Q)$ at every point $Q$ over no prime, and the local frames of the flag bundle
$h^*_jL_s$ of (e) are local frames of $\Lambda'\cap\mathrm{Lie}'_s$ on the points over no prime.
\end{proof}
% Lean: FlagStep.eigLat_isFlagVBOn_of_pull_char0: RegAtDom.regAt_dominates for the units, SatModel.lat_eq_iSup_of_char0 for t
% (hsum) and for t' (hind), hle, hall and mem_of_mem_iSup_of_indep for the key equality, then the frames of hfl.

\begin{lemma}\label{cv-origin-reference-flag}
In the setting of Definition~\ref{cv-origin-hodge-component}, let $\Omega_0$ and $\xi_0$ be as in
Lemma~\ref{cv-origin-reference-char}, and let $K$, $I_k$, $y$, $X_1$, $\xi_1$, $f_1$, $F_c\to\Omega_0$,
$\mathcal A_1$, $L_s$ and $J_\partial$ be as in the conclusion of
Lemma~\ref{cv-origin-reference-coords}. Then there are a field $\Omega\supset\Omega_0$ of characteristic
zero, finite sets $K'$ and $I'_k$, a point $y'\in\Omega^{\sqcup_kI'_k}$, a scheme $X_2$ with a point
$\xi_2:\operatorname{Spec}\Omega\to X_2$ of which every point of $X_2$ is a specialization, covered by the
affine open subschemes of the charts of $y'$, a morphism $f_2:X_2\to\mathcal X$ with
$f_2\circ\xi_2=\xi_0\circ\operatorname{Spec}(\Omega\supset\Omega_0)$, an embedding $F_c\to\Omega$ through
which $f_2\circ\xi_2$ lies over $\mathcal O_c$, a semi-abelian scheme $\mathcal A_2$ over $X_2$ with an action
of $\mathcal O_E$ and a symmetric cubical line bundle, identified with $f_2^*\mathcal A$, with Lie algebra
$\mathrm{Lie}_2$ at $\xi_2$, for each label $s$ a basis $b_s=(b_{s,1},\dots,b_{s,r_s})$ of the subspace
$\mathrm{Lie}_{2,s}$ of $\mathrm{Lie}_2$ on which $\mathcal O_E$ acts through $s$, and an
ideal $J'_\partial$ of $\mathbf Z[x_{k,i}:k\in K',\ i\in I'_k]$, such that:
\begin{enumerate}
\item the local rings of $X_2$ at the points over no prime are integrally closed;
\item with the Lie algebra of $\mathcal A$ at $\xi_0$, which maps to $\mathrm{Lie}_2$ compatibly with the
  maps to $\mathcal A$, the bases $b_s$ satisfy the genericity condition of
  Definition~\ref{cv-origin-hodge-component}(2) over $\Omega_0$ with $\deg=1$;
\item near every point over no prime, the vectors of $\mathrm{Lie}_{2,s}$ in the Lie algebra of
  $\mathcal A_2$ form a vector bundle with a basis $b'$ with
  $\langle b'_1,\dots,b'_i\rangle=\langle b_{s,1},\dots,b_{s,i}\rangle$ for every $i$;
\item $J'_\partial$ satisfies condition (3) of Definition~\ref{cv-origin-reference} for $y'$, with
  $D=|\Psi|N(N-1)/2$;
\item the complex points of $X_2$ whose coordinates lie in the zero set of $J'_\partial$ are those whose
  image under $f_2$ lies outside the abelian locus of $\mathcal A$.
\end{enumerate}
\end{lemma}
\leanref{Development/AutomorphicSpectrum/Proofs/GeometryProofs/CoordsWithConstantsMapStatement.lean}{GlobalTheory.T.coords_with_constants_map}
\begin{proof}
By Lemma~\ref{cv-origin-reference-char}, $\Omega_0$ has characteristic zero. Let $\mathrm{Lie}_1$ be the Lie algebra of
$\mathcal A_1$ at $\xi_1$ and, for a label $s$, $\mathrm{Lie}_{1,s}$ the subspace on which $\mathcal O_E$ acts through $s$; by (1)
of Lemma~\ref{cv-origin-reference-coords}, $L_s$ is a vector bundle of rank $r_s$ with generic fibre $\mathrm{Lie}_{1,s}$.

\emph{The flag bundle.} Let $x$ be the Segre products of $y$, the coordinates of $X_1$.
Lemma~\ref{cv-origin-reference-flag-bundle}, applied to $X_1$, $x$ and the $L_s$, gives a field $\Omega\supset\Omega_0$, a
scheme $X_2$ with a morphism $\xi_2:\operatorname{Spec}\Omega\to X_2$ of whose image every point is a specialization, with
coordinates $x'=(x'_0,\dots,x'_N)$, a smooth morphism $h:X_2\to X_1$ with
$h\circ\xi_2=\xi_1\circ\operatorname{Spec}(\Omega\supset\Omega_0)$, and families $b^1_s$ of $r_s$ vectors of
$\Omega\otimes_{\Omega_0}\mathrm{Lie}_1$ that form a generic flag of degree $1$ along $e\mapsto1\otimes e$ over the
$\mathrm{Lie}_{1,s}$, such that the pullback of each $L_s$ along $e\mapsto1\otimes e$ is a flag bundle on $X_2$ adapted to
$b^1_s$. The field $\Omega$ has characteristic zero, as it contains $\Omega_0$.

\emph{Lie algebras.} Lemma~\ref{cv-origin-reference-pullback} gives $\mathcal A_2=h^*\mathcal A_1$ with a morphism to
$\mathcal A_1$ over $h$, and its composite, a morphism $\mathcal A_2\to\mathcal A$ over $f_2=f_1\circ h$. Let $\mathrm{Lie}_2$ be
the Lie algebra of $\mathcal A_2$ at $\xi_2$, with its action of $\mathcal O_E$. Lemma~\ref{cv-origin-reference-lie-pull}(1),
applied to $h$, gives an additive map $j':\mathrm{Lie}_1\to\mathrm{Lie}_2$, compatible with the maps to $\mathcal A_1$, with
$j'(ae)=aj'(e)$ for $a\in\Omega_0$, which carries linearly independent families to linearly independent families and whose image
spans $\mathrm{Lie}_2$, such that the Lie algebras of $\mathcal A_2$ at the points of $X_2$ form the pullback along $j'$ of those
of $\mathcal A_1$. Let $\mathrm{Lie}_0$ be the Lie algebra of $\mathcal A$ at $\xi_0$. As $f_1\circ\xi_1=\xi_0$, the same lemma,
applied to $f_1$ with the identity of $\Omega_0$, gives an $\Omega_0$-linear map $a:\mathrm{Lie}_0\to\mathrm{Lie}_1$ compatible
with the maps to $\mathcal A$, which carries nonzero vectors to nonzero vectors and whose image spans; so $a$ is an isomorphism.
By Lemma~\ref{cv-origin-reference-flag-linear}(1) there is an isomorphism $\Phi:\Omega\otimes_{\Omega_0}\mathrm{Lie}_1\to\mathrm{Lie}_2$
with $\Phi(1\otimes e)=j'(e)$; in particular $j'$ is injective. By Lemma~\ref{cv-origin-reference-lie-pull}(2), $j'$ and $a$
commute with the actions of $\mathcal O_E$. Let $\mathrm{Lie}_{2,s}$ be the subspace of $\mathrm{Lie}_2$ on which $\mathcal O_E$
acts through $s$ and $F_c\to\Omega_0\subset\Omega$. As $j'$ is injective and semilinear, $j'(e)\in\mathrm{Lie}_{2,s}$ if and only
if $e\in\mathrm{Lie}_{1,s}$; and $V^0_s=a^{-1}(\mathrm{Lie}_{1,s})$ is the subspace of $\mathrm{Lie}_0$ on which $\mathcal O_E$ acts
through $s$.

\emph{The generic flag.} Put $b_s=\Phi(b^1_s)$. By Lemma~\ref{cv-origin-reference-flag-linear}(2), applied to $a$,
$e\mapsto1\otimes e$ and $\Phi$, the $b_s$ form a generic flag of degree $1$ along $\Phi\circ(1\otimes-)\circ a=j'\circ a$ over
the $V^0_s$. Its elements $u_{s,lk}$, for $l>k$, are algebraically independent over $\Omega_0$, and by
Lemma~\ref{cv-origin-reference-flag-linear}(4), as the degree is $1$, they generate $\Omega$ over $\Omega_0$. Since $j'\circ a$ is
compatible with the maps to $\mathcal A$, this is (2).

\emph{Coordinates and the boundary.} Let $y'$ be the point with the blocks of $y$, with values in $\Omega_0\subset\Omega$, and the
block $x'$, and let $J'_\partial$ be the image of $J_\partial$ under the renaming of the blocks of $y$ into those of $y'$. By
Lemma~\ref{cv-origin-reference-flag-coords}(1), $X_2$ is covered by the affine open subschemes of the charts of $y'$. By
Lemma~\ref{cv-origin-reference-flag-snc}, applied to $h$ and the $u_{s,lk}$, $J'_\partial$ satisfies condition (3) of
Definition~\ref{cv-origin-reference} for $y'$ with $|\Psi|pq+n$ in place of $D$, where $n$ is the number of the $u_{s,lk}$.
There are $|\Psi|$ labels with $r_s=p$ and $|\Psi|$ with $r_s=q$, so $n=|\Psi|(p(p-1)/2+q(q-1)/2)$ and
$|\Psi|pq+n=|\Psi|N(N-1)/2$, as $N=p+q$. This is (4). By Lemma~\ref{cv-origin-reference-integral}, applied to $J'_\partial$, the
local rings of $X_2$ at the points over no prime are integrally closed: this is (1). For (5), let $w$ be coordinates for $y'$
of a complex point $z$ of $X_2$. By Lemma~\ref{cv-origin-reference-flag-coords}(2), $w_1=(w_{k,i})_{k\in K}$ are coordinates for
$y$ of $h\circ z$. As in the proof of Lemma~\ref{cv-origin-reference-coords}, $w$ lies in the zero set of $J'_\partial$ if and
only if $w_1$ lies in that of $J_\partial$, which by (4) of Lemma~\ref{cv-origin-reference-coords} holds if and only if
$f_1(h(z))=f_2(z)$ lies outside the abelian locus of $\mathcal A$.

\emph{Flags.} By Lemma~\ref{cv-origin-reference-flag-linear}(3), the pullback of each $L_s$ along $j'=\Phi\circ(1\otimes-)$ is a
flag bundle on $X_2$ adapted to $b_s$. Lemma~\ref{cv-origin-reference-flag-constants}, applied to $h$, $\mathcal A_1$,
$\mathcal A_2$, $j'$, the $L_s$ and the $b_s$, with (1) for $X_2$, Lemma~\ref{cv-origin-reference-integral} for $X_1$, (2) of
Lemma~\ref{cv-origin-reference-coords}, the pullback of the Lie algebras along $j'$ and $j'(\mathrm{Lie}_{1,s})\subset\mathrm{Lie}_{2,s}$,
shows that the vectors of $\mathrm{Lie}_{2,s}$ in the Lie algebras of $\mathcal A_2$ form a flag bundle adapted to $b_s$ on the
points of $X_2$ over no prime: this is (3). The image of $\xi_2$ lies over no prime, and the Lie algebra of $\mathcal A_2$ there
is $\mathrm{Lie}_2$, by Lemma~\ref{cv-origin-reference-lie-bundle}, as the local ring there is $\Omega$. So by (3) the linearly
independent family $b_s$ spans $\mathrm{Lie}_{2,s}$, and it is a basis of it.

Finally $f_2\circ\xi_2=f_1\circ\xi_1\circ\operatorname{Spec}(\Omega\supset\Omega_0)=\xi_0\circ\operatorname{Spec}(\Omega\supset\Omega_0)$,
which lies over $\mathcal O_c$ through $F_c\to\Omega_0\subset\Omega$.
\end{proof}
% Lean: coords_with_constants_map = exists_flagStep, step by step: T.exists_model_vector_bundles (Lemma
% cv-origin-reference-flag-bundle); sabPull and sabComp; lieModel_exists and lieModel_exists_unique_act;
% T.lie_model_dominant_hodge for f' (j') and for f (a0, hainj and hasurj for the isomorphism a);
% T.tensor_equiv_of_basis_extension (Phi); T.lie_model_pullback_action_ring and heig; V0 = comap a (eigSp);
% T.generic_flag_under_equiv; T.unipotent_entries_field_degree (htop); T.coords_added_projective (tie1);
% T.snc_model_added_coords with T.dimension_plus_flag_count (FlagStep.dim_add_card_strictLower, the count of n);
% EigBase.isIntegrallyClosed_of_snc; T.flag_bundle_under_equiv and T.lie_model_constants_extension; T.complex_coords_added
% with EigBase.mem_mcPts_rename. The Lean does not state that b_s spans Lie_{2,s}; the last sentence of (3) derives it, as
% SatModel.span_eq_of_isFlagVBOn does in the proof of Lemma cv-origin-reference-descent.

\begin{lemma}\label{cv-origin-reference-descent-quotients}
Let $T$ be a scheme, $\Omega$ a field and $\xi:\operatorname{Spec}\Omega\to T$ a morphism, $V$ an $\Omega$-vector space, $S$ a set of
points that specialize the image of $\xi$, and $L$ a lattice family in $V$ that is a vector bundle of rank $g$ on $S$ with
generic fibre $W$ (Definition~\ref{cv-origin-reference-lattices}). For a subspace $V_0\subset V$, let $L/V_0$ be the lattice family
of the images of the $L(P)$ in $V/V_0$.
\begin{enumerate}
\item Let $b_1,\dots,b_r\in W$ be linearly independent, and suppose that for every $n$, $L/\langle b_1,\dots,b_n\rangle$ is a
  vector bundle on $S$ of some rank with some generic fibre. Then $L\cap\langle b_1,\dots,b_r\rangle$ is a flag bundle on $S$
  adapted to $b$.
\item Let $V_0\subset W$ be a subspace such that $L/V_0$ is a vector bundle on $S$ of some rank with some generic fibre. Then
  $L\cap V_0$ is a vector bundle of rank $\dim V_0$ on $S$ with generic fibre $V_0$, and at every $P\in S$ the
  $\mathcal O_P$-module $(L\cap V_0)(P)$ has a complement in $L(P)$.
\end{enumerate}
\end{lemma}
\leanref{Development/AutomorphicSpectrum/Proofs/GeometryProofs/LatticeFlagVectorBundleStatement.lean}{GlobalTheory.T.lattice_flag_vector_bundle}
\leanref{Development/AutomorphicSpectrum/Proofs/GeometryProofs/LatticeSubQuotientSplitStatement.lean}{GlobalTheory.T.lattice_sub_quotient_split}
\begin{proof}
For $P\in S$, the ring $\mathcal O_P$ is local, as the image of the local ring of $T$ at $P$. A finite projective module over
$\mathcal O_P$ is free \cite[Tag 00NX]{Sta}, and two finite families that generate the same $\mathcal O_P$-module are combinations of
each other with coefficients in $\mathcal O_P$.

(2) Let $P\in S$. By the definition of a vector bundle there are an open neighbourhood $U$ of $P$ and a basis $b$ of $W$ with
$L(Q)=\sum_i\mathcal O_Qb_i$ for $Q\in S\cap U$, and a linearly independent family $c=(c_1,\dots,c_s)$ in $V/V_0$ that generates the
image of $L(P)$. Lift each $c_i$ to $e_i\in L(P)$, and put $C=\sum_i\mathcal O_Pe_i$ and $N=L(P)\cap V_0$. An element $x$ of $L(P)$
has image $\sum_ia_ic_i$ with $a_i\in\mathcal O_P$, and $x-\sum_ia_ie_i$ lies in $N$; an element $\sum_ia_ie_i$ of $C$ that lies in
$V_0$ has $\sum_ia_ic_i=0$, so it is $0$. So $L(P)=C\oplus N$. The module $L(P)$ is free with basis $b$, so $N$, a direct summand,
is finite projective, hence free, with a basis $n=(n_1,\dots,n_k)$; the $e_i$ are independent over $\mathcal O_P$, as their images
are, so $(e,n)$ is a basis of $L(P)$ and $s+k=g$. The family $(e,n)$ spans the same $\Omega$-subspace as $L(P)$, which is $W$, and
it has $g=\dim W$ members, so it is linearly independent. It spans $V_0$ through $n$: an element of $V_0\subset W$ is $y+z$ with
$y\in\langle e\rangle$ and $z\in\langle n\rangle\subset V_0$, and then $y$ lies in $V_0$, so the images of its coefficients on the
$c_i$ vanish and $y=0$. Hence $k=\dim V_0$.

The bases $b$ and $(e,n)$ of $L(P)$ are combinations of each other with coefficients in $\mathcal O_P$. These finitely many
coefficients are regular on an open neighbourhood $U'\subset U$ of $P$, by Lemma~\ref{cv-origin-segre-regular}(1), so $(e,n)$
generates $L(Q)$ for every $Q\in S\cap U'$. For such $Q$, an element $\sum_ia_ie_i+\sum_ja'_jn_j$ of $L(Q)$ lies in $V_0$ only if
$\sum_ia_ie_i$ does, that is, only if every $a_i$ is $0$; so $(L\cap V_0)(Q)=\sum_j\mathcal O_Qn_j$. So $L\cap V_0$ is a vector
bundle of rank $\dim V_0$ on $S$ with generic fibre $V_0$, and $C$ is a complement of $(L\cap V_0)(P)$ in $L(P)$.

(1) For $n\ge0$ put $F_n=\langle b_1,\dots,b_{\min(n,r)}\rangle\subset W$. By (2), applied to $V_0=F_n$, each $L\cap F_n$ is a vector
bundle on $S$ of rank $\dim F_n$ with generic fibre $F_n$, and at every $P\in S$ the module $M_n=(L\cap F_n)(P)$ has a complement
$C_n$ in $L(P)$; $M_n$ is free of rank $\dim F_n$, generated by a basis of $F_n$. Fix $P\in S$. By induction on $n\le r$ there are
vectors $c_1,\dots,c_n$ with $\langle c_1,\dots,c_i\rangle=F_i$ for $i\le n$ and $M_n=\sum_{i\le n}\mathcal O_Pc_i$. For $n=0$ there
is nothing to prove. Let $n<r$, with $c_1,\dots,c_n$ given. Put $D=C_n\cap M_{n+1}$. An element $x$ of $M_{n+1}\subset C_n\oplus M_n$ is $y+z$ with
$y\in C_n$ and $z\in M_n\subset M_{n+1}$, so $y\in D$; hence $M_{n+1}=D\oplus M_n$. So $D$ is finite projective, hence free, and of
rank $(n+1)-n=1$. Let $c_{n+1}$ generate it. Then $M_{n+1}=M_n+\mathcal O_Pc_{n+1}$, so
$\langle c_1,\dots,c_{n+1}\rangle=F_n+\Omega c_{n+1}$ is contained in $F_{n+1}$ and contains the $\Omega$-span $F_{n+1}$ of
$M_{n+1}$.

For $n=r$ this gives $c=(c_1,\dots,c_r)$, with $\langle c_1,\dots,c_i\rangle=\langle b_1,\dots,b_i\rangle$ for every $i$, which
generates $(L\cap F_r)(P)$. It spans $F_r$, of dimension $r$, so it is linearly independent. Near $P$, the vector bundle $L\cap F_r$
is generated at the points of $S$ by a basis of $F_r$, and $c$ and this basis are combinations of each other with coefficients in
$\mathcal O_P$; as in (2), $c$ generates $(L\cap F_r)(Q)$ at every point $Q$ of $S$ in an open neighbourhood of $P$. So
$L\cap\langle b_1,\dots,b_r\rangle$ is a flag bundle on $S$ adapted to $b$.
\end{proof}
% Lean: (2) is inter_isVBOn_of_map: T.coherent_lattice_closure (a vector bundle is locally generated by finitely many
% vectors), T.regAt_local_order_fraction (regAt p is local), SatSub.exists_frame (the lifts e, hdec and hdisj, N free by
% exists_basis_of_isCompl with Module.free_of_flat_of_isLocalRing, the count hcard, hcatOmega, and the span of n), then
% T.frame_open_neighbourhood (= HtPt.Mc.frame_spread, the mutual combinations near_span_eq) and SatSub.sup_inf_eq. (1) is
% isFlagVBOn_inter_of_steps: (2) for every flagStep n, FlagVB.exists_adapted (the induction, with
% FlagVB.exists_gen_of_summand for D of rank 1), hlib, and T.frame_open_neighbourhood for L.inter (span b0).

\begin{lemma}\label{cv-origin-reference-descent-model}
Let $\Omega$ be a field, $T$ a scheme with a morphism $\xi:\operatorname{Spec}\Omega\to T$ of whose image every point is a
specialization, and $x$ a finite family in $\Omega$ such that $T$ has the coordinates $x$
(Definition~\ref{cv-origin-segre-coordinates}). Let $w_a$, for $a$ in a finite set, be finite families in $\Omega$, each with a
nonzero entry, and $V$ an $\Omega$-vector space. For $s$ in a finite set let $L_s$ be a lattice family in $V$ that is a vector
bundle of rank $g_s$ with generic fibre $W_s$ (Definition~\ref{cv-origin-reference-lattices}), and
$b_s=(b_{s,1},\dots,b_{s,r_s})$ a linearly independent family in $W_s$. Then there are a finite family $x'$ in $\Omega$, a scheme
$T'$ with a morphism $\xi':\operatorname{Spec}\Omega\to T'$ of whose image every point is a specialization, with the coordinates
$x'$, and a morphism $f:T'\to T$ with $f\circ\xi'=\xi$, such that:
\begin{enumerate}
\item the local rings of $T'$ are integrally closed in $\Omega$;
\item every point of $T'$ lies in a chart of each family $w_a$;
\item each $f^*L_s\cap\langle b_s\rangle$ is a flag bundle on $T'$ adapted to $b_s$;
\item for every $s$, every $n\ge0$ and every point $P'$ of $T'$, the $\mathcal O_{P'}$-module
  $(f^*L_s)(P')\cap\langle b_{s,1},\dots,b_{s,n}\rangle$ has a complement in $(f^*L_s)(P')$;
\item $f$ is a normalization over every point $P$ of $T$ that lies in a chart of each $w_a$ and at which, for every $s$ and
  $n$, the image of $L_s(P)$ in $V/\langle b_{s,1},\dots,b_{s,n}\rangle$ is generated over $\mathcal O_{T,P}$ by linearly independent
  vectors.
\end{enumerate}
\end{lemma}
\leanref{Development/AutomorphicSpectrum/Proofs/GeometryProofs/FlagModelStatementHoldsStatement.lean}{GlobalTheory.T.FlagModelStatement_holds}
\leanref{Development/AutomorphicSpectrum/Statements/GeometryStatements/FlagModelStatement.lean}{GlobalTheory.CompanionTheory.Ar.FlagModelStatement}
\begin{proof}
Each $W_s$ is spanned by a local basis of $L_s$, so it is finite-dimensional. For $s$ and $0\le n\le r_s$ let $q_{s,n}$ be the
projection $V\to V/\langle b_{s,1},\dots,b_{s,n}\rangle$. Apply Lemma~\ref{cv-origin-reference-refinement} to $T$, $x$, the families
$w_a$, and the lattice families $L_s/\langle b_{s,1},\dots,b_{s,n}\rangle$ with the subspaces $q_{s,n}(W_s)$: near every point,
$L_s/\langle b_{s,1},\dots,b_{s,n}\rangle$ is generated by the image of a local basis of $L_s$, and its value at every point spans
$q_{s,n}(W_s)$. This gives $x'$, $T'$, $\xi'$ and $f$ with (1) and (2), such that each
$f^*(L_s/\langle b_{s,1},\dots,b_{s,n}\rangle)=(f^*L_s)/\langle b_{s,1},\dots,b_{s,n}\rangle$ is a vector bundle, and such that $f$ is a
normalization over every point $P$ in a chart of each $w_a$ at which each image of $L_s(P)$ in a quotient by
$\langle b_{s,1},\dots,b_{s,n}\rangle$ with $n\le r_s$ is generated by linearly independent vectors. This gives (5). As
$\langle b_{s,1},\dots,b_{s,n}\rangle=\langle b_s\rangle$ for $n\ge r_s$, each $(f^*L_s)/\langle b_{s,1},\dots,b_{s,n}\rangle$ with
$n\ge0$ is a vector bundle. A local basis of $L_s$ near $f(P')$ generates
$f^*L_s$ near $P'$, as $f(P')$ specializes the image of $\xi$ and $\mathcal O_{f(P')}\subset\mathcal O_{P'}$ by
Lemma~\ref{cv-origin-segre-regular}(3); so $f^*L_s$ is a vector bundle of rank $g_s$ with generic fibre $W_s$.
Lemma~\ref{cv-origin-reference-descent-quotients}(1), applied to $f^*L_s$ and $b_s$, gives (3), and (2) of that lemma, applied
to $f^*L_s$ and the subspaces $\langle b_{s,1},\dots,b_{s,n}\rangle\subset W_s$, gives (4).
\end{proof}
% Lean: FlagModel.flagModel: T.flag_model_common_refinement for the quotient lattices (L s).map (flagStep n).mkQ, indexed
% by (s, n) with n <= r s (coh_map, span_map), pull_vb, pull_map with flagStep_min for (5), then T.lattice_flag_vector_bundle
% for (3) and T.lattice_sub_quotient_split for (4).

\begin{lemma}\label{cv-origin-reference-descent-roots-piece}
Let $X$ be a scheme, $\Omega$ a field and $\xi:\operatorname{Spec}\Omega\to X$ a morphism of whose image every point of $X$ is a
specialization, and let $x=(x_j)_{j\in J}$ be a finite family in $\Omega$ such that $X$ has the coordinates $x$
(Definition~\ref{cv-origin-segre-coordinates}). Let $c$ be an index with $x_c\neq0$, let $X_c$ be the affine open subscheme of the
points in the chart $c$, and put $A_c=\mathbf Z[x_j/x_c:j\in J]$.
\begin{enumerate}
\item For every point $P$ of $X_c$ and every open neighbourhood $V$ of $P$ there is a nonzero $g\in A_c$ that is a unit at $P$,
  such that every point of $X_c$ at which $g$ is a unit lies in $V$.
\item For a nonzero $g\in A_c$, the points of $X_c$ at which $g$ is a unit form an open subset $X_{c,g}$ of $X$, and an element
  $a\in\Omega$ is regular at every point of $X_{c,g}$ if and only if $g^ka\in A_c$ for some $k\ge0$.
\end{enumerate}
\end{lemma}
\leanref{Development/AutomorphicSpectrum/Proofs/GeometryProofs/LineBundleGlobalGenerators.lean}{GlobalTheory.CompanionTheory.Ar.LineGen.exists_basic}
\leanref{Development/AutomorphicSpectrum/Proofs/GeometryProofs/PieceLocalizationAtElementStatement.lean}{GlobalTheory.T.piece_localization_at_element}
\begin{proof}
Let $\eta$ be the image of $\xi$. By condition (1) of Definition~\ref{cv-origin-segre-coordinates} every element of $\Omega$ is
regular at $\eta$, so $\eta$ lies in $X_c$. By Lemma~\ref{cv-origin-segre-regular}(2) and condition (4), the map
$\Gamma(X_c,\mathcal O_X)\to\Omega$ is injective with image $A_c$. For a section $f$ over $X_c$ with image $g$ and a point $Q$ of
$X_c$, the germ of $f$ is a unit in $\mathcal O_{X,Q}$ if and only if $g$ is a unit at $Q$: if $g^{-1}$ is the image of
$t\in\mathcal O_{X,Q}$, the germ of $f$ times $t$ maps to $1$, so it is $1$, as $\mathcal O_{X,Q}\to\Omega$ is injective.

(1) There is a section $f$ over $X_c$ whose standard open $D(f)$ contains $P$ and lies in $V$ \cite[Tag 00E0]{Sta}. Its image
$g$ lies in $A_c$ and is a unit at $P$, so it is nonzero, and a point of $X_c$ at which $g$ is a unit lies in $D(f)\subset V$.

(2) Let $f$ be the section over $X_c$ with image $g$. By the first paragraph, $X_{c,g}$ is the standard open $D(f)$, and it
contains $\eta$. As $X_c$ is affine, $\Gamma(D(f),\mathcal O_X)$ is the localization of $\Gamma(X_c,\mathcal O_X)$ at $f$
\cite[Tag 01HV]{Sta}, so its image in $\Omega$ is $A_c[1/g]$. By Lemma~\ref{cv-origin-segre-regular}(2), applied to $D(f)$, the
elements regular at every point of $D(f)$ are those of $A_c[1/g]$, the $a$ with $g^ka\in A_c$ for some $k$.
\end{proof}
% Lean: (1) is LineGen.exists_basic (IsAffineOpen.exists_basicOpen_le, the germ g at the generic point, and IsCptGen.inj
% for the units); (2) is piece_localization_at_element = FlagBun.Sm.exists_basicOpen_piece (germ_genPt_injective_range_regAt,
% IsAffineOpen.basicOpen, isLocalization_basicOpen), with locSub R a and FlagBun.Sm.mem_locSub_iff for the powers of g.

\begin{lemma}\label{cv-origin-reference-descent-roots-twist}
Let $X$, $\Omega$, $\xi$ and $x$ be as in Lemma~\ref{cv-origin-reference-descent-roots-piece}, and let $\Lambda$ be a lattice
family in $\Omega$ that is a vector bundle of rank $1$ with generic fibre $\Omega$ (Definition~\ref{cv-origin-reference-lattices}).
Then there are an integer $n\ge0$ and a finite family $(s_l)$ in $\Omega$ such that, for every point $P$ of $X$ and every chart $c$
of $x$ that contains $P$, every $s_l/x_c^n$ lies in $\Lambda(P)$ and some $s_l/x_c^n$ generates $\Lambda(P)$.
\end{lemma}
\leanref{Development/AutomorphicSpectrum/Proofs/GeometryProofs/LineBundleGlobalGeneratorsStatement.lean}{GlobalTheory.T.line_bundle_global_generators}
\begin{proof}
For an index $c$ with $x_c\neq0$ let $X_c$ and $A_c$ be as in Lemma~\ref{cv-origin-reference-descent-roots-piece}; the elements of
$A_c$ are regular at the points of $X_c$.

\emph{Local generators.} Let $x_c\neq0$ and $P\in X_c$. As $\Lambda$ is a vector bundle of rank $1$, there are an open
neighbourhood $V$ of $P$ and a nonzero $b\in\Omega$ with $\Lambda(Q)=\mathcal O_Qb$ for every $Q\in V$. By
Lemma~\ref{cv-origin-reference-descent-roots-piece}(1) there is a nonzero $g\in A_c$, a unit at $P$, with $X_{c,g}\subset V$. The
open sets $X_{c,g}$ so obtained cover $X_c$, which is quasi-compact, being affine \cite[Tag 00E8]{Sta}. So there is a finite set
$T_c$ of pairs $(g,b)$ of nonzero elements of $\Omega$ with $g\in A_c$, such that $\Lambda(Q)=\mathcal O_Qb$ at every point $Q$ of
$X_c$ at which $g$ is a unit, and every point of $X_c$ has a pair in $T_c$ whose $g$ is a unit there. Let $I$ be the finite set of
the triples $i=(c_i,g_i,b_i)$ with $x_{c_i}\neq0$ and $(g_i,b_i)\in T_{c_i}$.

\emph{Degrees.} For $i,j\in I$ put $u_{ij}=x_{c_i}/x_{c_j}\in A_{c_j}$. As $u_{ij}\cdot x_l/x_{c_i}=x_l/x_{c_j}$, multiplication by
$u_{ij}^d$ carries a monomial of degree $d$ in the ratios $x_l/x_{c_i}$ into $A_{c_j}$; so there is $e_{ij}$ with
$u_{ij}^{e'}g_i\in A_{c_j}$ for every $e'\ge e_{ij}$. Let $e$ be the largest $e_{ij}$.

\emph{Pairs.} Let $i,j\in I$, write $u=u_{ij}$, and put $w=g_ju^{e+1}g_i$, a nonzero element of $A_{c_j}$. Let $Q$ be a point of
$X_{c_j}$ at which $w$ is a unit. Its factors $g_j$, $u$ and $u^eg_i$ in $A_{c_j}$ are regular at $Q$, so they are units there:
the inverse of each is $w^{-1}$ times the other two.
Then $Q$ lies in the chart $c_i$, as $x_l/x_{c_i}=(x_l/x_{c_j})u^{-1}$ is regular at $Q$, and $g_i=u^{-e}(u^eg_i)$ is a unit at
$Q$. So $\mathcal O_Qb_i=\Lambda(Q)=\mathcal O_Qb_j$, and $b_i/b_j$ is regular at $Q$. By
Lemma~\ref{cv-origin-reference-descent-roots-piece}(2), $w^{k_{ij}}b_i/b_j\in A_{c_j}$ for some $k_{ij}\ge0$. Let $M$ be the
largest $k_{ij}$. Let $Q$ be a point in the chart $c_j$ at which $g_j$ is a unit. Then $\Lambda(Q)=\mathcal O_Qb_j$, and
\[
u^{(e+1)M}g_i^Mb_i=g_j^{-M}\cdot w^{M-k_{ij}}\cdot\bigl(w^{k_{ij}}b_i/b_j\bigr)\cdot b_j
\]
with $g_j^{-M}$ regular at $Q$ and the middle factors in $A_{c_j}$, so $u^{(e+1)M}g_i^Mb_i\in\Lambda(Q)$.

\emph{Generators.} Put $n=(e+1)M$ and $s_i=x_{c_i}^ng_i^Mb_i$ for $i\in I$. Let $P$ be a point in the chart $c$ of $x$. Then
$x_c\neq0$ and $P\in X_c$, so there is $j\in I$ with $c_j=c$ and $g_j$ a unit at $P$. For every $i$,
$s_i/x_c^n=u_{ij}^ng_i^Mb_i$ lies in $\Lambda(P)$ by the previous paragraph, and $s_j/x_c^n=g_j^Mb_j$ generates
$\Lambda(P)=\mathcal O_Pb_j$, as $g_j^M$ is a unit at $P$.
\end{proof}
% Lean: line_twist_generated: piece_chart (the chart opens as IsPiece), LineGen.exists_cover (exists_basic,
% piece_localization_at_element, compactness of the affine open), LineGen.exists_deg (induction on the adjoin), LineGen.pair
% (hf: b/b' in locSub through hU1.ring, then mem_locSub_iff), and the sections T i = x_(cc i)^n * gg i^M * bb i. The Lean
% indexes the family by Option I, with the entry 0 at none, to make it a Fin (m+1).

\begin{lemma}\label{cv-origin-reference-descent-roots}
In the setting of Definition~\ref{cv-origin-hodge-component}, let $Z$ be a Hodge component of $\mathcal X$ at a finite set of
primes with coordinates $y\in\Omega^{\sqcup_{k\in K}I_k}$, and let a flag of $Z$ be given. Then there are finite sets $K'$ and
$I'_k$ and a point $y'\in\Omega^{\sqcup_{k\in K'}I'_k}$, with an injection $k\mapsto k'$ from $K$ to $K'$ and bijections
$i\mapsto i'$ from $I_k$ to $I'_{k'}$ with $y'_{k',i'}=y_{k,i}$, such that $Z$ is covered by the affine open subschemes of the
charts of $y'$ as in Definition~\ref{cv-origin-hodge-component}(1), and the flag has root lines of some multidegrees
$a_{\tau,i}\in\mathbf Z^{K'}$ with generators $e_{\tau,i}$ for the coordinates $y'$
(Definition~\ref{cv-origin-hodge-component}(3)).
\end{lemma}
\leanref{Development/AutomorphicSpectrum/Proofs/GeometryProofs/HodgeCptBlockRefinementStatement.lean}{GlobalTheory.T.hodge_cpt_block_refinement}
\begin{proof}
Let $\mathcal V_s$ be the vector bundles of Definition~\ref{cv-origin-hodge-component}(1) and $b_s$ the bases of the flag. Let $x$ be
the Segre products of $y$, so that $Z$ has the coordinates $x$ (Definition~\ref{cv-origin-segre-coordinates}).

\emph{The root lines are line bundles.} Let $s$ be a label, $1\le i\le r_s$, and $P$ a point of $Z$. By
Definition~\ref{cv-origin-hodge-component}(2) there are an open neighbourhood $U$ of $P$ and a basis $b'$ with
$\langle b'_1,\dots,b'_k\rangle=\langle b_{s,1},\dots,b_{s,k}\rangle$ for every $k$ and $\mathcal V_s(Q)=\sum_k\mathcal O_Qb'_k$ for
every $Q\in U$. Write $b'_i=tb_{s,i}+v$ with $v\in\langle b_{s,1},\dots,b_{s,i-1}\rangle$; then $t\neq0$, as $b'_i$ is not in
$\langle b'_1,\dots,b'_{i-1}\rangle$. Let $Q\in U$. An element of $\mathcal V_s(Q)$ lies in
$\langle b_{s,1},\dots,b_{s,i}\rangle=\langle b'_1,\dots,b'_i\rangle$ if and only if it is $\sum_{k\le i}a_kb'_k$ with
$a_k\in\mathcal O_Q$, as $b'$ is linearly independent, and then its coefficient of $b_{s,i}$ is $a_it$. So the $i$-th graded
piece at $Q$ is $\mathcal O_Qt$. Each root line $\Lambda_{\tau,i}$ is a graded piece or its inverse, so it is generated by $t$ or
$t^{-1}$ at every point of $U$: it is a vector bundle of rank $1$ with generic fibre $\Omega$
(Definition~\ref{cv-origin-reference-lattices}).

\emph{New blocks.} For each $(\tau,i)\in\Psi\times\{1,\dots,N\}$, Lemma~\ref{cv-origin-reference-descent-roots-twist}, applied to $Z$,
$x$ and $\Lambda_{\tau,i}$, gives an integer $n_{\tau,i}\ge0$ and a finite family $s_{\tau,i}=(s_{\tau,i,l})_{l\in L_{\tau,i}}$. Let
$K'=K\sqcup(\Psi\times\{1,\dots,N\})$, and let $y'$ be the point with the blocks $y_k$ of $y$ and the blocks $s_{\tau,i}$, with the
inclusion of $K$ and the identity maps of the $I_k$. Let $P$ be a point of $Z$, in the chart $c$ of $x$. Some
$s_{\tau,i,l}/x_c^{n_{\tau,i}}$ generates $\Lambda_{\tau,i}(P)$, so $s_{\tau,i,l}\neq0$, and every $s_{\tau,i,l'}/x_c^{n_{\tau,i}}$ lies
in $\Lambda_{\tau,i}(P)$, so $s_{\tau,i,l'}/s_{\tau,i,l}$ is regular at $P$. So $P$ lies in the chart $l$ of the block
$s_{\tau,i}$, and Lemma~\ref{cv-origin-reference-segre} shows that $Z$ is covered by the affine open subschemes of the charts of
$y'$.

\emph{Generators.} Let $\kappa'$ be a chart of $y'$ and $P$ a point of $Z_{\kappa'}$. The restriction $\kappa$ of $\kappa'$ to $K$ is a
chart of $y$ whose ratios are regular at $P$, so $P$ lies in the chart $\kappa$ of $x$, as each ratio of $x$ there is a product of
ratios of the blocks of $y$. Let $l=\kappa'_{(\tau,i)}$. Some $s_{\tau,i,l'}/x_\kappa^{n_{\tau,i}}$ generates $\Lambda_{\tau,i}(P)$, and
it is $s_{\tau,i,l'}/s_{\tau,i,l}$, which is regular at $P$, times $s_{\tau,i,l}/x_\kappa^{n_{\tau,i}}$, which lies in
$\Lambda_{\tau,i}(P)$. So $\Lambda_{\tau,i}(P)$ is generated by
\[
s_{\tau,i,l}\prod_{k\in K}y_{k,\kappa_k}^{-n_{\tau,i}}=\prod_{k\in K'}y'^{\,a_{\tau,i,k}}_{k,\kappa'_k},
\]
where $a_{\tau,i}\in\mathbf Z^{K'}$ is $-n_{\tau,i}$ on $K$, $1$ at the block $(\tau,i)$ and $0$ at the other new blocks. So the
flag has root lines of multidegrees $a_{\tau,i}$ with generators $e_{\tau,i}=1$ for the coordinates $y'$.
\end{proof}
% Lean: T.hodge_cpt_block_refinement = HodgeCpt.Flag.exists_recoord_rootLines: hlin from Fl.generic, Lat.IsFlagVBOn.isVBOn_gr_inter
% (exists_diag and gr_eq: the graded pieces are line bundles, generated by t), T.lattice_tensor_power_line_bundle =
% LatTens.tens_isVBOn (Fl.rootLat r is Lat.tens of a graded piece with exponent rootSgn = 1 or -1, generated by t^(+-1)), and
% T.exists_block_refinement_lines = exists_addBlocks_isDegOn: T.line_bundle_global_generators for each root line, the blocks
% s_2 indexed by I = CMType H x Fin (p+q), exists_chart_of_gen, T.coords_join_blocks, block_chart and segre_chart,
% isGen_of_chart, and the multidegree Sum.elim (-n i) (indicator of i) with e0 = 1.

\begin{lemma}\label{cv-origin-reference-descent}
In the setting of Definition~\ref{cv-origin-hodge-component}, let $\Omega_0$ and $\xi_0$ be as in
Lemma~\ref{cv-origin-reference-char}, let $\deg\ge1$, and let $\Omega$, $y'$, $X_2$, $\xi_2$, $f_2$,
$F_c\to\Omega$, $\mathcal A_2$ and $b_s$ satisfy the conclusions of Lemma~\ref{cv-origin-reference-flag}
other than (4) and (5), with $\deg$ in place of $1$ in (2). Then there are finite sets $K''$ and $I''_k$,
a point $y''\in\Omega^{\sqcup_kI''_k}$ with an injection $k\mapsto k''$ from the index set $K'$ of $y'$ to
$K''$ and bijections $i\mapsto i''$ from $I'_k$ to $I''_{k''}$ with $y''_{k'',i''}=y'_{k,i}$, a Hodge
component $Z$ of $\mathcal X$ at the empty set with coordinates $y''$, lying over the irreducible
component of $\mathcal X$ whose generic point is the image of $\xi_0$, a morphism $g:Z\to X_2$ that maps the generic point
of $Z$ to $\xi_2$, with $f=f_2\circ g$ for the morphism $f:Z\to\mathcal X$ of $Z$, and a flag of $Z$ of
degree $\deg$, whose bases have the same images as the $b_{s,i}$ in the Lie algebra of $\mathcal A$, with
root lines of some multidegrees with generators, such that at every point $P$ of $Z$ over no prime the
local ring of $Z$ at $P$ equals that of $X_2$ at $g(P)$ as subrings of $\Omega$, and $P$ is the only point of
$Z$ with image $g(P)$.
\end{lemma}
\leanref{Development/AutomorphicSpectrum/Proofs/GeometryProofs/GenericCoordsDescendSubfieldStatement.lean}{GlobalTheory.T.generic_coords_descend_subfield}
\begin{proof}
Let $\mathrm{Lie}_2$ be the Lie algebra of $\mathcal A_2$ at $\xi_2$, with the subspaces $\mathrm{Lie}_{2,s}$, and $\Lambda$ the
lattice family of the Lie algebras of $\mathcal A_2$ at the points of $X_2$
(Definition~\ref{cv-origin-reference-lattices}). As $\mathcal A_2$ is smooth over $X_2$, $\Lambda$ is a vector bundle of rank
$\dim\mathrm{Lie}_2$ with generic fibre $\mathrm{Lie}_2$ (Lemma~\ref{cv-origin-reference-lie-bundle}). The families $b_s$ are
linearly independent, by (2). The image $\eta_2$ of $\xi_2$ lies over no prime, and $\Lambda(\eta_2)=\mathrm{Lie}_2$, as the
local ring there is $\Omega$; by (3), $\mathrm{Lie}_{2,s}=\Lambda(\eta_2)\cap\mathrm{Lie}_{2,s}$ has a basis with the same flag as
$b_s$, so $\langle b_s\rangle=\mathrm{Lie}_{2,s}$.

\emph{A model of the flags.} Let $x$ be the Segre products of $y'$, the coordinates of $X_2$; as $\eta_2$ lies in a chart of
$x$, some entry of $x$ is nonzero. Lemma~\ref{cv-origin-reference-descent-model}, applied to $X_2$, $x$, the single family $w=x$,
and for each label $s$ the lattice family $\Lambda$ with the family $b_s$, gives a finite family $x'$, a scheme $Z$ with a
morphism $\xi_Z:\operatorname{Spec}\Omega\to Z$ of whose image every point is a specialization, with the coordinates $x'$, and a
morphism $g:Z\to X_2$ with $g\circ\xi_Z=\xi_2$, such that the local rings of $Z$ are integrally closed in $\Omega$, every point of
$Z$ lies in a chart of $x$, each $g^*\Lambda\cap\mathrm{Lie}_{2,s}$ is a flag bundle on $Z$ adapted to $b_s$, and $g$ is a
normalization over every point $P$ of $X_2$ in a chart of $x$ at which every image of $\Lambda(P)$ in a quotient
$\mathrm{Lie}_2/\langle b_{s,1},\dots,b_{s,n}\rangle$ is generated by linearly independent vectors. As in the proof of
Lemma~\ref{cv-origin-reference-coords}, Lemma~\ref{cv-origin-reference-segre} shows that $Z$ is covered by the affine open
subschemes of the charts of the point $y''_1$ with the blocks of $y'$ and the block $x'$.

\emph{The Hodge component.} Lemma~\ref{cv-origin-reference-pullback} gives $\mathcal A_Z=g^*\mathcal A_2$ with a morphism to
$\mathcal A_2$ over $g$, and its composite, a morphism $\mathcal A_Z\to\mathcal A$ over $f=f_2\circ g$. By
Lemma~\ref{cv-origin-reference-lie-same}, the Lie algebra of $\mathcal A_Z$ at $\xi_Z$ is identified with $\mathrm{Lie}_2$,
compatibly with the actions of $\mathcal O_E$ and the maps to $\mathcal A_2$, and the Lie algebras of $\mathcal A_Z$ at the points
of $Z$ form $g^*\Lambda$. So the vectors of $\mathrm{Lie}_{2,s}$ in them form the flag bundle $g^*\Lambda\cap\mathrm{Lie}_{2,s}$
adapted to $b_s$, in particular a vector bundle of rank $r_s$ with generic fibre $\mathrm{Lie}_{2,s}$. The local rings of $Z$ are
domains with field of fractions $\Omega$, integrally closed in $\Omega$, so $Z$ is normal; it is irreducible, as every point
specializes the image of $\xi_Z$. The morphism $f$ maps the image of $\xi_Z$ to the image of $\xi_0$, the generic point of an
irreducible component of $\mathcal X$, with $f\circ\xi_Z=f_2\circ\xi_2=\xi_0\circ\operatorname{Spec}(\Omega\supset\Omega_0)$, and
$f\circ\xi_Z$ lies over $\mathcal O_c$ through $F_c\to\Omega$. At the empty set of primes, the modules $\mathcal V_s$ of
Definition~\ref{cv-origin-hodge-component}(1) are the vectors of $\mathrm{Lie}_{2,s}$ in the Lie algebras of $\mathcal A_Z$. So
these data form a Hodge component of $\mathcal X$ at the empty set with coordinates $y''_1$. The bases $b_s$, with the Lie
algebra of $\mathcal A$ at $\xi_0$ and its map to $\mathrm{Lie}_2$ of (2), form a flag of degree $\deg$: the $\mathcal V_s$ have
bases adapted to $b_s$ near every point, and the genericity condition is (2). Each $b_{s,i}$ has the same image in the Lie
algebra of $\mathcal A$ as before, since $\mathcal A_Z\to\mathcal A$ is the composite $\mathcal A_Z\to\mathcal A_2\to\mathcal A$.

\emph{Points over no prime.} By Lemma~\ref{cv-origin-reference-eigen}, applied to $\mathcal A_2$ with (1), the
$\mathrm{Lie}_{2,s}$ are independent and $\Lambda(P)=\sum_s(\Lambda(P)\cap\mathrm{Lie}_{2,s})$ at every point $P$ of $X_2$ over no
prime. Let $P$ be such a point. By (3), each $\Lambda(P)\cap\mathrm{Lie}_{2,t}$ is generated by linearly independent vectors
$B_{t,1},\dots,B_{t,r_t}$ with the same flag as $b_t$. Together the $B_{t,i}$ are linearly independent, as the $\mathrm{Lie}_{2,t}$
are independent, and they generate $\Lambda(P)$. For a label $s$ and $n\ge0$, $\langle b_{s,1},\dots,b_{s,n}\rangle$ is spanned by
the $B_{s,i}$ with $i\le n$, and the images of the other $B_{t,i}$ in $\mathrm{Lie}_2/\langle b_{s,1},\dots,b_{s,n}\rangle$ are
linearly independent and generate the image of $\Lambda(P)$. As $P$ lies in a chart of $x$, $g$ is a normalization over $P$. As
in the proof of Lemma~\ref{cv-origin-reference-coords}, every point $P'$ of $Z$ over no prime lies over a point $P=g(P')$ over
no prime, and Lemma~\ref{cv-origin-reference-normalization}, with (1), shows that the local ring of $Z$ at $P'$ is that of $X_2$
at $P$ and that $P'$ is the only point of $Z$ over $P$.

\emph{Root lines.} Lemma~\ref{cv-origin-reference-descent-roots}, applied to this Hodge component and its flag, gives a point
$y''$ that contains $y''_1$ as some of its blocks, through an injection of blocks and bijections, such that $Z$ is covered by the
affine open subschemes of the charts of $y''$ and the flag has root lines of some multidegrees with generators for $y''$. With
the coordinates $y''$, the same data form a Hodge component of $\mathcal X$ at the empty set with the same flag. The blocks of
$y'$ are carried into those of $y''$ through $y''_1$.
\end{proof}
% Lean: generic_coords_descend_subfield = exists_sat_model = SatModel.exists_flagged_model followed by
% T.hodge_cpt_block_refinement (Lemma cv-origin-reference-descent-roots), with the blocks composed (hk.comp hbk). In
% exists_flagged_model: T.lie_omega_lattice_vector_bundle, SatModel.span_eq_of_isFlagVBOn (hspan), T.FlagModelStatement_holds
% with L = M.lat for every label and W = top, T.segre_refine_line_bundles with no line bundles (Lemma cv-origin-reference-segre,
% as in exists_eigBase), sabPull, sabComp, SatModel.exists_lieModel_pull, hsub, SatModel.isNormalScheme_of_normalCpt, the
% HodgeCpt with S empty (Lat.patch_empty, R = 0, isVBOn_of_isFlagVBOn for vb) and its Flag (j_pt from hn, hMl and
% spMap_dualMap_id), SatModel.lat_eq_iSup_of_char0, SatModel.exists_frame_map_mkQ for hnormAt, RegAtDom.regAt_dominates for
% hdown, and SatModel.normAt_id_elim for hreg and huniq.

\subsection{Existence of reference geometries}\label{subsection-cv-reference-geometries-exist}

\begin{lemma}\label{cv-origin-snc-charts}
Let $X$ be a scheme, $\Omega$ a field and $\xi:\operatorname{Spec}\Omega\to X$ a morphism, let
$y\in\Omega^{\sqcup_{k\in K}I_k}$, and suppose that $X$ is covered by the affine open subschemes of the
charts of $y$ as in Definition~\ref{cv-origin-hodge-component}(1). Let $\kappa$ be a chart of $y$, with
affine open subscheme $X_\kappa$, and put $A^{\mathbf Z}_\kappa=\mathbf Z[y_{k,i}/y_{k,\kappa_k}:k\in K,\ i\in I_k]$.
\begin{enumerate}
\item Let $P$ be a point of $X_\kappa$. An element $a\in\Omega$ is regular at $P$ if and only if $sa\in A^{\mathbf Z}_\kappa$
  for some $s\in A^{\mathbf Z}_\kappa$ that is a unit at $P$.
\item For every prime $\mathfrak p$ of $A^{\mathbf Z}_\kappa$ there is exactly one point $P$ of $X_\kappa$ at which the
  elements of $A^{\mathbf Z}_\kappa$ that are units are those outside $\mathfrak p$.
\end{enumerate}
\end{lemma}
\leanref{Development/AutomorphicSpectrum/Proofs/GeometryProofs/RegAtIffChartLocalizationStatement.lean}{GlobalTheory.T.regAt_iff_chart_localization}
\leanref{Development/AutomorphicSpectrum/Proofs/GeometryProofs/PrimeChartExistsUniquePointStatement.lean}{GlobalTheory.T.prime_chart_existsUnique_point}
\begin{proof}
Let $B$ be the ring of $X_\kappa$ and $\varphi:B\to\Omega$ the homomorphism through the local ring at the image
of $\xi$, which specializes to every point of $X_\kappa$. For a point $P$ of $X_\kappa$ with prime
$\mathfrak q\subset B$, the local ring of $X$ at $P$ is $B_{\mathfrak q}$ \cite[Tag 01HV]{Sta}, and its embedding in
$\Omega$ extends $\varphi$. The elements of $\varphi(B)$ are regular at every point of $X_\kappa$, so
$\varphi(B)\subset A^{\mathbf Z}_\kappa$ by Definition~\ref{cv-origin-hodge-component}(1).

(1) If $a$ is regular at $P$, then $a=\varphi(t)/\varphi(r)$ with $t,r\in B$ and $r\notin\mathfrak q$. Then
$s=\varphi(r)$ lies in $A^{\mathbf Z}_\kappa$ and is a unit at $P$, and $sa=\varphi(t)\in A^{\mathbf Z}_\kappa$. Conversely, the
elements of $A^{\mathbf Z}_\kappa$ are regular at $P$, so $a=s^{-1}(sa)$ is.

(2) First, $\varphi(B)=A^{\mathbf Z}_\kappa$. Let $a\in A^{\mathbf Z}_\kappa$, and let $I$ be the ideal of the $r\in B$ with
$\varphi(r)a\in\varphi(B)$. If $I$ were proper, it would lie in a maximal ideal $\mathfrak m$, the prime of a point
$Q$ of $X_\kappa$. As $a$ is regular at $Q$, $a=\varphi(t)/\varphi(r)$ with $r\notin\mathfrak m$, and then
$r\in I\subset\mathfrak m$. So $1\in I$ and $a\in\varphi(B)$. For $r\in B$, $r\notin\mathfrak q$ if and only if
the germ of $r$ is a unit at $P$, if and only if $\varphi(r)$ is a unit at $P$, since the local ring embeds in
$\Omega$. So the point with prime $\varphi^{-1}(\mathfrak p)$ has the property, and every point with the property has
the prime $\varphi^{-1}(\mathfrak p)$.
\end{proof}
% Lean: (1) is regAt_iff_exists_chart_fraction (the stalk as a localization of Gamma(U), isLocalization_stalk);
% (2) is chart_point_of_prime, with ChartPrimes.exists_secVal_eq for phi(B) = A_kappa and key for the units.

\begin{definition}\label{cv-origin-segre-coordinates}
Let $X$ be a scheme, $\Omega$ a field and $\xi:\operatorname{Spec}\Omega\to X$ a morphism with image $\eta$. For a
point $P$ of $X$ that specializes $\eta$, an element of $\Omega$ is \emph{regular at} $P$ if it is the image of an
element of the local ring $\mathcal O_{X,P}$ under $\mathcal O_{X,P}\to\mathcal O_{X,\eta}\to\Omega$, and it is a
\emph{unit at} $P$ if it is nonzero and it and its inverse are regular at $P$. Let $x=(x_j)_{j\in J}$ be a finite
family in $\Omega$. A point $P$ that specializes $\eta$ lies in the \emph{chart} $j$ of $x$ if $x_j\neq0$ and every
$x_i/x_j$ is regular at $P$. The scheme $X$ \emph{has the coordinates} $x$ if:
\begin{enumerate}
\item the map $\mathcal O_{X,\eta}\to\Omega$ is bijective, and $\mathcal O_{X,P}\to\Omega$ is injective for every
  point $P$ that specializes $\eta$;
\item every point that specializes $\eta$ lies in a chart of $x$;
\item for every $j$ with $x_j\neq0$, the points that specialize $\eta$ and lie in the chart $j$ are the points of an
  affine open subscheme $X_j$ of $X$;
\item for every $j$ with $x_j\neq0$, the elements of $\Omega$ that are regular at every point of $X_j$ are those of
  $\mathbf Z[x_i/x_j:i\in J]$.
\end{enumerate}
For a point $y\in\Omega^{\sqcup_{k\in K}I_k}$, the \emph{Segre products} of $y$ are the family
$(\prod_{k\in K}y_{k,c_k})_{c\in\prod_kI_k}$. Its charts are the charts of $y$, and its ratios at a chart $\kappa$
generate the ring $\mathbf Z[y_{k,i}/y_{k,\kappa_k}:k\in K,\ i\in I_k]$. The scheme $X$ is \emph{covered by the affine
open subschemes of the charts} of $y$, as in Definition~\ref{cv-origin-hodge-component}(1), if it has the Segre
products of $y$ as coordinates.
\end{definition}
\leanref{Development/AutomorphicSpectrum/Definitions/GlobalGeometry/CoordsAt.lean}{GlobalTheory.CompanionTheory.Ar.CoordTie}
\leanref{Development/AutomorphicSpectrum/Definitions/GlobalGeometry/CoordsAt.lean}{GlobalTheory.CompanionTheory.Ar.InChart}
\leanref{Development/AutomorphicSpectrum/Definitions/GlobalGeometry/IsCptGen.lean}{GlobalTheory.CompanionTheory.Ar.IsCptGen}
\leanref{Development/AutomorphicSpectrum/Definitions/GlobalGeometry/IsCptGen.lean}{GlobalTheory.CompanionTheory.Ar.unitAt}

\begin{lemma}\label{cv-origin-segre-regular}
Let $X$ be a scheme, $\Omega$ a field and $\xi:\operatorname{Spec}\Omega\to X$ a morphism with image $\eta$.
\begin{enumerate}
\item If $a\in\Omega$ is regular at a point $P$ that specializes $\eta$, then $P$ has an open neighbourhood $O$ such
  that $a$ is regular at every point of $O$ that specializes $\eta$. If every point of $X$ specializes $\eta$, then
  for every finite family $x$ in $\Omega$ and every index $j$ the points in the chart $j$ of $x$ form an open subset
  of $X$.
\item Suppose that every point of $X$ specializes $\eta$ and that $X$ satisfies condition (1) of
  Definition~\ref{cv-origin-segre-coordinates}, and let $V$ be an open subset of $X$ that contains $\eta$. Then the
  map $\Gamma(V,\mathcal O_X)\to\mathcal O_{X,\eta}\to\Omega$ is injective, and its image is the set of elements of
  $\Omega$ that are regular at every point of $V$.
\item Let $X'$ be a scheme with a morphism $\xi':\operatorname{Spec}\Omega\to X'$, let $f:X'\to X$ be a morphism with
  $f\circ\xi'=\xi$, and let $P'$ be a point of $X'$ that specializes the image of $\xi'$. Then $f(P')$ specializes
  $\eta$, the map $\mathcal O_{X,f(P')}\to\mathcal O_{X',P'}\to\Omega$ is the map $\mathcal O_{X,f(P')}\to\Omega$ of
  Definition~\ref{cv-origin-segre-coordinates}, and so every element regular at $f(P')$ is regular at $P'$. If
  $\mathcal O_{X',P'}\to\Omega$ is injective, then every element regular at $f(P')$ that is a unit at $P'$ is a unit
  at $f(P')$.
\end{enumerate}
\end{lemma}
\leanref{Development/AutomorphicSpectrum/Proofs/GeometryProofs/GenMapGermEqGenMapGerm.lean}{GlobalTheory.CompanionTheory.Ar.exists_open_mem_regAt}
\leanref{Development/AutomorphicSpectrum/Proofs/GeometryProofs/DominatesOfCoordsRegular.lean}{GlobalTheory.CompanionTheory.Ar.CoordMap.isOpen_chart}
\leanref{Development/AutomorphicSpectrum/Proofs/GeometryProofs/DominatesOfCoordsRegular.lean}{GlobalTheory.CompanionTheory.Ar.CoordMap.val_injective}
\leanref{Development/AutomorphicSpectrum/Proofs/GeometryProofs/DominatesOfCoordsRegular.lean}{GlobalTheory.CompanionTheory.Ar.CoordMap.exists_val_eq}
\leanref{Development/AutomorphicSpectrum/Proofs/GeometryProofs/GenMapStalkMap.lean}{GlobalTheory.CompanionTheory.Ar.RegAtDom.regAt_dominates}
\begin{proof}
(1) The element $a$ is the image of a germ at $P$ of a section $s$ of $\mathcal O_X$ over an open neighbourhood $O$ of
$P$. As $O$ contains a specialization of $\eta$, it contains $\eta$. For a point $Q$ of $O$ that specializes $\eta$,
the germs of $s$ at $Q$ and at $P$ have the same image as the germ of $s$ at $\eta$ in $\mathcal O_{X,\eta}$, so $a$ is
regular at $Q$. If every point specializes $\eta$, the chart $j$ is empty if $x_j=0$, and otherwise it is the
intersection of the finitely many open sets on which the ratios $x_i/x_j$ are regular.

(2) For $s\in\Gamma(V,\mathcal O_X)$ with image $a$, the germ of $s$ at every point $P$ of $V$ maps to $a$, so $a$ is
regular at every point of $V$; and if $a=0$, every germ of $s$ is $0$, as $\mathcal O_{X,P}\to\Omega$ is injective,
so $s=0$. Conversely, let $a$ be regular at every point of $V$. By (1), $V$ is covered by open sets $O_P$ with
sections $s_P$ whose images in $\Omega$ are $a$. On $O_P\cap O_Q$, if it is nonempty, the restrictions of $s_P$ and
$s_Q$ have the same image, so they are equal by the injectivity just shown. They glue to a section over $V$ with
image $a$.

(3) The morphism $f$ maps the image $\eta'$ of $\xi'$ to $\eta$, as $f\circ\xi'=\xi$, and preserves specialization, so
$f(P')$ specializes $\eta$. As $\xi=f\circ\xi'$, the map $\mathcal O_{X,\eta}\to\Omega$ is
$\mathcal O_{X,\eta}\to\mathcal O_{X',\eta'}\to\Omega$, and the maps of $f$ on local rings commute with
$\mathcal O_{X,f(P')}\to\mathcal O_{X,\eta}$ and $\mathcal O_{X',P'}\to\mathcal O_{X',\eta'}$; this gives the
compatibility, and with it the transfer of regular elements. Let $a$ be the image of $s\in\mathcal O_{X,f(P')}$, and a
unit at $P'$, with $a^{-1}$ the image of $t\in\mathcal O_{X',P'}$. Then $f^\sharp(s)t$ maps to $1$, so
$f^\sharp(s)t=1$ by injectivity, and $f^\sharp(s)$ is a unit. As $f^\sharp$ is a local homomorphism, $s$ is a unit,
and $a^{-1}$ is the image of $s^{-1}$.
\end{proof}

\begin{lemma}\label{cv-origin-segre-model}
Let $\Omega$ be a field and $x=(x_j)_{j\in J}$ a finite family in $\Omega$ with an index $j_0$ such that
$x_{j_0}\neq0$ and every element of $\Omega$ is a quotient of two elements of $\mathbf Z[x_j/x_{j_0}:j\in J]$. Then
there are a scheme $X$ and a morphism $\xi:\operatorname{Spec}\Omega\to X$ of whose image every point of $X$ is a
specialization, such that $X$ has the coordinates $x$ (Definition~\ref{cv-origin-segre-coordinates}).
\end{lemma}
\leanref{Development/AutomorphicSpectrum/Proofs/GeometryProofs/ProjectiveModelNonemptyStatement.lean}{GlobalTheory.T.projective_model_nonempty}
\begin{proof}
For $j$ with $x_j\neq0$ write $A_j=\mathbf Z[x_i/x_j:i\in J]\subset\Omega$. Let $V\subset\Omega$ be the subgroup
generated by the $x_j$, let $V^n$ be the subgroup generated by the products of $n$ elements of $V$, with
$V^0=\mathbf Z\cdot1$, and let $S=\bigoplus_{n\ge0}V^nT^n\subset\Omega[T]$, a graded ring. Put $t_j=x_jT\in S_1$. The
ring $S_0$ is the image of $\mathbf Z$, and the $t_j$ generate $S$ over $S_0$, since $V^nT^n$ is spanned by the
products of $n$ of them. Let $\mathrm{ev}:S\to\Omega$ be the evaluation at $T=1$; it is injective on each $S_n$.

Let $j$ be an index with $x_j\neq0$. As $\mathrm{ev}(t_j)=x_j\neq0$, the evaluation extends to $S[1/t_j]$, and its
restriction to the degree-zero part $S_{(t_j)}$ is $\theta_j(a/t_j^n)=\mathrm{ev}(a)/x_j^n$ for $a\in S_n$. The map
$\theta_j$ is injective, as $\mathrm{ev}$ is injective on $S_n$, and its image is $A_j$: $\mathrm{ev}(a)\in V^n$ and
$V^n/x_j^n\subset A_j$, while $x_i/x_j=\theta_j(t_i/t_j)$. In the same way $\theta_{jk}(a/(t_jt_k)^n)=\mathrm{ev}(a)/(x_jx_k)^n$
embeds $S_{(t_jt_k)}$ in $\Omega$ when $x_jx_k\neq0$.

Let $X=\operatorname{Proj}S$. For $j$ with $x_j\neq0$, the open subset $D_+(t_j)$ is an affine open subscheme
\cite[Tag 01MB]{Sta} with ring of sections $S_{(t_j)}$ \cite[Tag 01M7]{Sta}, so it is isomorphic to
$\operatorname{Spec}S_{(t_j)}$; write $\iota_j:\operatorname{Spec}A_j\to X$ for the
open immersion given by $\theta_j$. If $x_j=0$ then $t_j=0$. The $t_j$ generate an ideal that contains $S_+$, so the
$D_+(t_j)$ with $x_j\neq0$ cover $X$ \cite[Tag 00JP]{Sta}. Let $\xi$ be $\iota_{j_0}$ composed with the morphism
$\operatorname{Spec}\Omega\to\operatorname{Spec}A_{j_0}$ given by the inclusion of $A_{j_0}$ in $\Omega$.

Let $j$ and $k$ be indices with $x_j,x_k\neq0$; then $D_+(t_j)\cap D_+(t_k)=D_+(t_jt_k)$ \cite[Tag 00JP]{Sta}. The
restriction map $S_{(t_j)}\to S_{(t_jt_k)}$ sends $a/t_j^n$ to $at_k^n/(t_jt_k)^n$ \cite[Tag 01M4]{Sta}, so it is
compatible with $\theta_j$ and $\theta_{jk}$. Hence $\iota_j$ and $\iota_k$, composed with the morphisms from
$\operatorname{Spec}\Omega$ given by $A_j,A_k\subset\Omega$, both equal $\operatorname{Spec}\Omega\to\operatorname{Spec}S_{(t_jt_k)}\to X$,
and so equal $\xi$. Moreover, inside $D_+(t_k)$ the open subset $D_+(t_jt_k)$ is the standard open of
$t_jt_k/t_k^2=t_j/t_k$ \cite[Tag 01M4]{Sta}, and $\theta_k(t_j/t_k)=x_j/x_k$: a point $\iota_k(\mathfrak p)$ with
$x_j/x_k\notin\mathfrak p$ lies in the image of $\iota_j$.

\emph{Points and local rings.} Let $x_j\neq0$, let $\mathfrak p$ be a prime of $A_j$ and $P=\iota_j(\mathfrak p)$. The
image $\eta$ of $\xi$ is $\iota_j$ of the zero ideal, which $\mathfrak p$ contains, so $P$ specializes $\eta$; as the
images of the $\iota_j$ cover $X$, every point of $X$ specializes $\eta$. The local ring of $X$ at $P$ is
$(A_j)_{\mathfrak p}$ \cite[Tag 01HV]{Sta}, and its map to $\Omega$ through $\mathcal O_{X,\eta}$ is the inclusion of
$(A_j)_{\mathfrak p}$ in $\Omega$. So it is injective, the elements regular at $P$ are the $a\in\Omega$ with $sa\in A_j$
for some $s\in A_j\setminus\mathfrak p$, and every $s\in A_j\setminus\mathfrak p$ is a unit at $P$. The local ring at
$\eta$, the zero ideal of $A_{j_0}$, is the field of fractions of $A_{j_0}$, which is $\Omega$ by hypothesis. This
gives condition (1) of Definition~\ref{cv-origin-segre-coordinates}.

\emph{Charts.} Every point $\iota_j(\mathfrak p)$ lies in the chart $j$ of $x$, as the ratios $x_i/x_j$ lie in $A_j$.
Conversely, let a point $Q=\iota_k(\mathfrak p)$ lie in the chart $j$. As $x_k/x_j$ is regular at $Q$, there is
$s'\in A_k\setminus\mathfrak p$ with $u=s'x_k/x_j\in A_k$. Then $(x_j/x_k)u=s'\notin\mathfrak p$, so
$x_j/x_k\notin\mathfrak p$, and $Q$ lies in the image of $\iota_j$. So the points in the chart $j$ are those of the
affine open subscheme $\iota_j(\operatorname{Spec}A_j)$, and every point lies in a chart: conditions (2) and (3).

\emph{Rings of ratios.} The elements of $A_j$ are regular at every point of $\iota_j(\operatorname{Spec}A_j)$.
Conversely, let $a\in\Omega$ be regular at all these points, and let $I$ be the ideal of the $b\in A_j$ with
$ba\in A_j$. For a maximal ideal $\mathfrak m$ of $A_j$, regularity at $\iota_j(\mathfrak m)$ gives $s\notin\mathfrak m$
with $sa\in A_j$, so $I\not\subset\mathfrak m$. Hence $I=A_j$ and $a\in A_j$: condition (4).
\end{proof}
% Lean: projective_model_nonempty = nonempty_pModel_of_chart, from GlobalTheory.T.exists_proper_integral_model =
% exists_cl_scheme (Proj of ClSch.grading, the Rees-type algebra rees o x; chartVal, chartVal_injective,
% mem_adjoin_iff_mem_range, chart_compat, exists_chart_point, chart_overlap), verified by ClSch.tie_of_charts (spec,
% chart_stalk, inj, chart1, chart2, gen, and the ideal J for the ring condition). exists_cl_scheme also shows that X is
% integral and proper over o, which is not used here. Size: about 770 lines.

\begin{lemma}\label{cv-origin-segre-map}
Let $\Omega$ be a field, let $X$ and $X'$ be schemes with morphisms $\xi:\operatorname{Spec}\Omega\to X$ and
$\xi':\operatorname{Spec}\Omega\to X'$, and let $x=(x_j)_{j\in J}$ be a finite family in $\Omega$ such that $X$ has the
coordinates $x$ (Definition~\ref{cv-origin-segre-coordinates}). Suppose that every point of $X'$ specializes the image
of $\xi'$, that $X'$ satisfies condition (1) of Definition~\ref{cv-origin-segre-coordinates}, and that every point of
$X'$ lies in a chart of $x$. Then there is a morphism $f:X'\to X$ with $f\circ\xi'=\xi$, and a point $P'$ of $X'$ lies
in the chart $j$ of $x$ if and only if $f(P')$ does.
\end{lemma}
\leanref{Development/AutomorphicSpectrum/Proofs/GeometryProofs/DominatesOfCoordsRegularStatement.lean}{GlobalTheory.T.dominates_of_coords_regular}
\begin{proof}
Let $\eta$ and $\eta'$ be the images of $\xi$ and $\xi'$. By condition (1) every element of $\Omega$ is regular at
$\eta$, so $\eta$ lies in the chart $j$ of $x$ for every $j$ with $x_j\neq0$, and so does $\eta'$. For such $j$, let
$U_j$ be the affine open subscheme of $X$ of condition (3), and $V_j$ the set of points of $X'$ in the chart $j$,
which is open by Lemma~\ref{cv-origin-segre-regular}(1). As $U_j$ and $V_j$ contain $\eta$ and $\eta'$, the morphisms
$\xi$ and $\xi'$ factor through them.

\emph{A morphism on each chart.} Let $x_j\neq0$. By Lemma~\ref{cv-origin-segre-regular}(2), applied to $X'$ and
$V_j$, the map $\Gamma(V_j,\mathcal O_{X'})\to\Omega$ is injective, and its image consists of the elements regular at
every point of $V_j$. The image of $\Gamma(U_j,\mathcal O_X)\to\Omega$ consists of elements regular at every point of
$U_j$, so it lies in $\mathbf Z[x_i/x_j]$ by condition (4), and the elements of this ring are regular at every point of
$V_j$. So $\Gamma(U_j,\mathcal O_X)\to\Omega$ factors through a ring homomorphism
$\Gamma(U_j,\mathcal O_X)\to\Gamma(V_j,\mathcal O_{X'})$, which, as $U_j$ is affine, is the map on global sections of a
morphism $g_j:V_j\to U_j$ \cite[Tag 01I1]{Sta}. The morphisms $g_j\circ\xi'$ and $\xi$ from
$\operatorname{Spec}\Omega$ to $U_j$ have the same map on global sections, so they are equal, by the same reference. Let
$G_j:V_j\to X$ be $g_j$ followed by the inclusion of $U_j$.

\emph{Agreement on overlaps.} Let $x_j,x_k\neq0$ and $P'\in V_j\cap V_k$. The point $G_k(P')$ lies in $U_k$, so
$x_j/x_k$ is regular at $G_k(P')$; at $P'$, both $x_j/x_k$ and its inverse are regular. By
Lemma~\ref{cv-origin-segre-regular}(3), applied to $G_k$ and $\xi'$, the element $x_j/x_k$ is a unit at $G_k(P')$, so
every $x_i/x_j=(x_i/x_k)(x_j/x_k)^{-1}$ is regular there, and $G_k(P')\in U_j$. So the restrictions of $G_j$ and $G_k$
to $W=V_j\cap V_k$ are two morphisms $W\to U_j$, which agree after composition with $\xi':\operatorname{Spec}\Omega\to W$.
The scheme $W$ is reduced, as its local rings embed in $\Omega$, and $U_j$ is affine, hence separated. The locus where
the two morphisms agree is a closed subscheme of $W$ \cite[Tag 01KM]{Sta} through which $\xi'$ factors, so it contains
$\eta'$ and hence every point of $W$; as $W$ is reduced, the identity of $W$ factors through it
\cite[Tag 0356]{Sta}, and the two morphisms are equal.

\emph{Gluing.} Every point of $X'$ lies in a chart of $x$, so the $V_j$ cover $X'$, and the $G_j$ glue to a morphism
$f:X'\to X$ \cite[Tag 01JB]{Sta}. As $\xi'$ factors through some $V_j$, $f\circ\xi'=G_j\circ\xi'=\xi$. If $P'$ lies
in the chart $j$, then $f(P')=G_j(P')\in U_j$ lies in the chart $j$. Conversely, if $f(P')$ lies in the chart $j$, the
ratios $x_i/x_j$ are regular at $f(P')$, hence at $P'$ by Lemma~\ref{cv-origin-segre-regular}(3).
\end{proof}
% Lean: CoordMap.exists_dominating_map: the opens U c (tie.affine) and V c (isOpen_chart), the lifts (exists_lift),
% exists_chart_map = exists_map_affine (val_injective, exists_val_eq), glue_compat (inChart_of_inChart, ext_of_land with
% Mathlib's ext_of_isDominant), Scheme.OpenCover.glueMorphisms, and inChart_pull for the charts. The Lean states it for
% a field embedding i and a coefficient ring o; the book states the case i = id, o = Z. Size: about 370 lines.

\begin{lemma}\label{cv-origin-segre-unique}
Let $\Omega$ be a field, let $X$ and $X'$ be schemes with morphisms $\xi:\operatorname{Spec}\Omega\to X$ and
$\xi':\operatorname{Spec}\Omega\to X'$ of whose images every point is a specialization, and let $x$ and $x'$ be finite
families in $\Omega$ such that $X$ has the coordinates $x$ and $X'$ the coordinates $x'$
(Definition~\ref{cv-origin-segre-coordinates}). If every point of $X'$ lies in a chart of $x$, then there is exactly one
morphism $f:X'\to X$ with $f\circ\xi'=\xi$.
\end{lemma}
\leanref{Development/AutomorphicSpectrum/Proofs/GeometryProofs/ExistsUniqueDominatingMapStatement.lean}{GlobalTheory.T.existsUnique_dominating_map}
\begin{proof}
A morphism exists by Lemma~\ref{cv-origin-segre-map}, as $X'$ satisfies condition (1) of
Definition~\ref{cv-origin-segre-coordinates}. Let $f$ and $g$ be two such morphisms, and let $P'$ be a point of $X'$.
By Lemma~\ref{cv-origin-segre-regular}(3), $f(P')$ and $g(P')$ specialize the image of $\xi$; let them lie in the charts
$c$ and $c'$ of $x$.

First, $g(P')$ specializes $f(P')$. The element $x_c/x_{c'}$ is regular at $g(P')$. It is a unit at $P'$: it is
regular there by Lemma~\ref{cv-origin-segre-regular}(3) applied to $g$, and so is its inverse, which is regular at
$f(P')$, by the same lemma applied to $f$. By Lemma~\ref{cv-origin-segre-regular}(3) for $g$, $x_c/x_{c'}$ is a unit at
$g(P')$, so every $x_i/x_c=(x_i/x_{c'})(x_c/x_{c'})^{-1}$ is regular there, and $g(P')$ lies in the chart $c$. Let $O$
be an open neighbourhood of $f(P')$, and $X_c$ the affine open subscheme of the chart $c$. There is a section $a$ of
$\mathcal O_X$ over $X_c$ whose standard open contains $f(P')$ and lies in $O$ \cite[Tag 00E0]{Sta}. A point $Q$ of
$X_c$ lies in this standard open if and only if the germ of $a$ is a unit in $\mathcal O_{X,Q}$, that is, as
$\mathcal O_{X,Q}$ embeds in $\Omega$, if and only if the image $\alpha$ of $a$ in $\Omega$ is a unit at $Q$. By
condition (4), $\alpha$ lies in $\mathbf Z[x_i/x_c]$, so it is regular at $g(P')$. It is a unit at $f(P')$, so at $P'$,
by Lemma~\ref{cv-origin-segre-regular}(3) for $f$, and then at $g(P')$, by the same lemma for $g$. So $g(P')$ lies in
$O$. As $O$ was arbitrary, $g(P')$ specializes $f(P')$.

By symmetry $f(P')$ specializes $g(P')$, so $f(P')=g(P')$, and $f$ and $g$ agree as maps of sets. Let $s$ be a section
of $\mathcal O_X$ over an open subset $O$. At each point $P'$ of $f^{-1}(O)$, the germs of $f^\sharp(s)$ and
$g^\sharp(s)$ both map to the image of $s$ in $\Omega$, by Lemma~\ref{cv-origin-segre-regular}(3); as
$\mathcal O_{X',P'}\to\Omega$ is injective, the germs are equal. So $f^\sharp(s)=g^\sharp(s)$, and $f=g$.
\end{proof}
% Lean: ChMorph.existsUnique_dominates: existence from dominates_of_coords_regular; uniqueness from
% specializes_of_dominates (RegAtDom.regAt_dominates, chart_unit_basic_open = HtPt.Ma.basic with
% IsAffineOpen.exists_basicOpen_le, Specializes.antisymm) and hom_eq_of_points (section_ext with IsCptGen.inj). The Lean
% states it for a field embedding i and coefficient rings o and o'; the book states the case i = id, o = o' = Z.

\begin{lemma}\label{cv-origin-segre-join}
Let $X$ be a scheme, $\Omega$ a field and $\xi:\operatorname{Spec}\Omega\to X$ a morphism of whose image every point of
$X$ is a specialization, and let $x=(x_j)_{j\in J}$ and $w=(w_l)_{l\in L}$ be finite families in $\Omega$ such that $X$
has the coordinates $x$ (Definition~\ref{cv-origin-segre-coordinates}) and some $w_l$ is nonzero. Then there are a
scheme $X'$ with a morphism $\xi':\operatorname{Spec}\Omega\to X'$ of whose image every point of $X'$ is a
specialization, such that $X'$ has the coordinates $(x_jw_l)_{(j,l)\in J\times L}$, and a morphism $f:X'\to X$ with
$f\circ\xi'=\xi$.
\end{lemma}
\leanref{Development/AutomorphicSpectrum/Proofs/GeometryProofs/ExistsJoinModelDominatesStatement.lean}{GlobalTheory.T.exists_join_model_dominates}
\begin{proof}
Let $\eta$ be the image of $\xi$. It lies in a chart $c$ of $x$, by condition (2), and some $w_l$ is nonzero, so
$x_cw_l\neq0$. Let $A$ be the ring generated by the ratios $x_jw_m/(x_cw_l)$; it contains every
$x_j/x_c=x_jw_l/(x_cw_l)$. Every element of $\Omega$ is regular at $\eta$, by condition (1), so it is a quotient of two
elements of $\mathbf Z[x_j/x_c]\subset A$, by Lemma~\ref{cv-origin-snc-charts}(1) applied to $x$ as a point with one
block. Lemma~\ref{cv-origin-segre-model}, applied to the family $(x_jw_m)$ and the index $(c,l)$, gives a scheme $X'$
and a morphism $\xi'$ of whose image every point of $X'$ is a specialization, such that $X'$ has the coordinates
$(x_jw_m)$. Every point $P'$ of $X'$ lies in a chart $(j,m)$ of this family; then $x_j\neq0$, and every
$x_i/x_j=x_iw_m/(x_jw_m)$ is regular at $P'$, so $P'$ lies in the chart $j$ of $x$. Lemma~\ref{cv-origin-segre-unique}
gives a morphism $f:X'\to X$ with $f\circ\xi'=\xi$.
\end{proof}
% Lean: exists_join_model_dominates = join_exists, for a field embedding i : Omega -> Omega' and a coefficient ring o; the
% book states the case i = id, o = Z, in which the hypothesis that i(Omega) and the ratios of w generate Omega' holds.
% hfrac uses regAt_chart_localization and regAt_local_order_fraction at the generic point; then projective_model_nonempty,
% join_coords_centre_chart = HtPt.Mb.join_charts (every point of X' lies in a chart of x), and
% existsUnique_dominating_map for f.

\begin{lemma}\label{cv-origin-reference-segre}
Let $X$ be a scheme, $\Omega$ a field and $\xi:\operatorname{Spec}\Omega\to X$ a morphism, let
$y\in\Omega^{\sqcup_{k\in K}I_k}$, and suppose that $X$ is covered by the affine open subschemes of the
charts of $y$ as in Definition~\ref{cv-origin-hodge-component}(1). Let $s\in\Omega^{\sqcup_{k\in K'}I'_k}$
be a further family of blocks such that every point of $X$ that specializes the image of $\xi$ lies in a chart of
each block of $s$: for each $k\in K'$ there is $i$ with $s_{k,i}\neq0$ and every $s_{k,j}/s_{k,i}$ regular at the
point. Then $X$
is covered in the same way by the affine open subschemes of the charts of the point
$(y,s)\in\Omega^{\sqcup_{k\in K\sqcup K'}}$.
\end{lemma}
\leanref{Development/AutomorphicSpectrum/Proofs/GeometryProofs/CoordsJoinBlocksStatement.lean}{GlobalTheory.T.coords_join_blocks}
\begin{proof}
Let $W$ be the union of the open subsets of $X$ all of whose points specialize the image of $\xi$. It contains the
affine open subschemes of the charts of $y$, and the open subscheme $W$ is covered by them in the same way, with the same
local rings; so we may replace $X$ by $W$ and assume that every point of $X$ specializes the image of $\xi$.

Let $x$ and $w$ be the Segre products of $y$ and of $s$ (Definition~\ref{cv-origin-segre-coordinates}). The Segre
products of $(y,s)$ are the products $x_cw_d$, indexed by the pairs $(c,d)$, so we must show that $X$ has the
coordinates $(x_cw_d)_{(c,d)}$. Every point of $X$ lies in a chart of $w$: if it lies in the chart $i_k$ of each block
$k$ of $s$, then $w_{(i_k)_k}\neq0$, and every ratio of $w$ at $(i_k)_k$ is a product of ratios $s_{k,j}/s_{k,i_k}$,
which are regular at the point. In particular some $w_d$ is nonzero, and Lemma~\ref{cv-origin-segre-join} gives a
scheme $X'$ with a morphism $\xi':\operatorname{Spec}\Omega\to X'$ of whose image every point is a specialization, with
the coordinates $(x_cw_d)$, and a morphism $f:X'\to X$ with $f\circ\xi'=\xi$. Every point of $X$ lies in a chart $c$
of $x$ and a chart $d$ of $w$, hence in the chart $(c,d)$ of $(x_cw_d)$, as each ratio $x_{c'}w_{d'}/(x_cw_d)$ is the
product of two regular ratios. So Lemma~\ref{cv-origin-segre-unique} gives a morphism $g:X\to X'$ with
$g\circ\xi=\xi'$. The same lemma, applied with $X$ in place of both schemes and then with $X'$ in place of both, shows
that the identity is the only endomorphism of $X$ compatible with $\xi$, and the only one of $X'$ compatible with
$\xi'$. So $f\circ g$ and $g\circ f$ are the identities, $f$ is an isomorphism with $f\circ\xi'=\xi$, and $X$, like
$X'$, has the coordinates $(x_cw_d)$.
\end{proof}
% Lean: GlobalTheory.T.coords_join_blocks = AddBlockMulti.add_block_multi: the Segre products of s have a chart at every
% point (hw), coordTie_join with exists_openModel (the open W, coordTie_of_push, coordTie_push) and
% join_coords_coordTie = Retie.coordTie_join_of_charts (exists_join_model_dominates, existsUnique_dominating_map for g
% with mem_ctrSet_join, and for the identities with mem_ctrSet_self, then coordTie_comp_iso), then coordTie_comp_equiv
% (Equiv.sumPiEquivProdPi). The three facts are Lemmas cv-origin-segre-join, cv-origin-segre-map and
% cv-origin-segre-unique.

\begin{lemma}\label{cv-origin-reference-snc}
Let $\Omega$ be a field of characteristic zero, let $X$ and $X'$ be schemes with morphisms
$\xi:\operatorname{Spec}\Omega\to X$ and $\xi':\operatorname{Spec}\Omega\to X'$ of whose images every point
is a specialization, and let $y\in\Omega^{\sqcup_{k\in K}I_k}$ and $y'\in\Omega^{\sqcup_{k\in K'}I'_k}$ be
points such that $X$ and $X'$ are covered by the affine open subschemes of the charts of $y$ and of $y'$.
Let $g:X'\to X$ be a morphism with $g\circ\xi'=\xi$, let $k\mapsto k'$ be a map from $K$ to $K'$ and
$i\mapsto i'$ bijections from $I_k$ to $I'_{k'}$ with $y'_{k',i'}=y_{k,i}$. Suppose that at every point
$P$ of $X'$ over no prime the local ring of $X'$ at $P$ equals that of $X$ at $g(P)$ as subrings of
$\Omega$, and that a point of $X'$ over no prime is the only point of $X'$ with its image under $g$. If an
ideal $J_\partial$ satisfies condition (3) of Definition~\ref{cv-origin-reference} for $y$ with an integer $e$
in place of $D$, then its image under the renaming $x_{k,i}\mapsto x_{k',i'}$ satisfies it for $y'$ with the
same $e$. (The hypotheses that every point of $X'$ is a specialization of the image of $\xi'$, that
$g\circ\xi'=\xi$ and that $g$ separates a point over no prime from the other points are not used in the
proof.)
\end{lemma}
\leanref{Development/AutomorphicSpectrum/Proofs/GeometryProofs/SncModelReindexedCoordsStatement.lean}{GlobalTheory.T.snc_model_reindexed_coords}
\begin{proof}
For a chart $\kappa$ of $y$ write $A_\kappa=\mathbf Q[y_{k,i}/y_{k,\kappa_k}]$, as in
Definition~\ref{cv-origin-reference}(3), and $A^{\mathbf Z}_\kappa$ as in Lemma~\ref{cv-origin-snc-charts}; likewise
$A'_{\kappa'}$ and $A'^{\mathbf Z}_{\kappa'}$ for a chart $\kappa'$ of $y'$. Let $\kappa'$ be a chart of $y'$, and let
$\kappa$ be given by $\kappa_k'=\kappa'_{k'}$. Then $y_{k,\kappa_k}=y'_{k',\kappa'_{k'}}\neq0$, so $\kappa$ is a chart of
$y$, and $y_{k,i}/y_{k,\kappa_k}=y'_{k',i'}/y'_{k',\kappa'_{k'}}$. Hence $A^{\mathbf Z}_\kappa\subset A'^{\mathbf Z}_{\kappa'}$ and
$A_\kappa\subset A'_{\kappa'}$. Condition (3) for $y$ at $\kappa$ gives $f_l,a_l\in A_\kappa$ with $\sum_la_lf_l=1$. Fix $l$
with $f_l\neq0$, let $S=A_\kappa[1/f_l]$ with its elements $z_1,\dots,z_e$ and integer $n\le e$, and put
$S'=A'_{\kappa'}[1/f_l]\supset S$.

\emph{Denominators.} Let $\mathfrak q$ be a prime of $S'$ and $t\in S'$; then $t=v/u$ with $u,v\in S$ and
$u\notin\mathfrak q$. By Lemma~\ref{cv-origin-snc-charts}(2), applied to $X'$ and the prime
$\mathfrak q\cap A'^{\mathbf Z}_{\kappa'}$, there is a point $P'$ of $X'_{\kappa'}$ at which the elements of $A'^{\mathbf Z}_{\kappa'}$
outside $\mathfrak q$ are the units. A positive integer is a unit in $S'$, so it lies outside $\mathfrak q$ and is a
unit at $P'$: $P'$ lies over no prime. Some integer $N\ge1$ has $Nf_l\in A'^{\mathbf Z}_{\kappa'}$, and $Nf_l$ is a unit
in $S'$, so it is a unit at $P'$ and $f_l^{-1}$ is regular at $P'$. Hence $S'$ lies in the local ring of $X'$ at
$P'$, which by hypothesis is that of $X$ at $g(P')$. The ratios of the chart $\kappa$ lie in $A'^{\mathbf Z}_{\kappa'}$, so
$g(P')$ lies in $X_\kappa$, and Lemma~\ref{cv-origin-snc-charts}(1) gives $u\in A^{\mathbf Z}_\kappa$, a unit at $g(P')$, with
$v=ut\in A^{\mathbf Z}_\kappa$. As the two local rings are equal, $u$ is a unit at $P'$, so $u\notin\mathfrak q$.

\emph{Smoothness and differentials.} Let $\mathfrak q$ be a prime of $S'$ and $\mathfrak p=\mathfrak q\cap S$. By the
previous paragraph $S_{\mathfrak p}\to S'_{\mathfrak q}$ is surjective, and it is injective since both are subrings of
$\Omega$. The ring $S$ is noetherian and $S'$ is of finite type over $S$, so $S\to S'$ is smooth at $\mathfrak q$
\cite[Tag 07VH]{Sta}, an isomorphism being formally smooth; and the localization of $\Omega_{S'/S}$ at $\mathfrak q$ is
$\Omega_{S_{\mathfrak p}/S}=0$ \cite[Tag 00RT]{Sta}. Hence $S\to S'$ is smooth \cite[Tag 00TC]{Sta}, so formally smooth
\cite[Tag 00TN]{Sta}, and formally unramified as $\Omega_{S'/S}=0$ \cite[Tag 00UO]{Sta}; it is formally \'etale. So
$S'$ is smooth over $\mathbf Q$ \cite[Tag 00TD]{Sta}, and $S'\otimes_S\Omega_{S/\mathbf Q}\to\Omega_{S'/\mathbf Q}$ is an
isomorphism \cite[Tag 0H93]{Sta}: the differentials of $z_1,\dots,z_e$, a basis of $\Omega_{S/\mathbf Q}$, form a basis
of $\Omega_{S'/\mathbf Q}$.

\emph{The boundary.} For $h\in J_\partial$ with image $h'$ under the renaming,
$h'(y'_{k,i}/y'_{k,\kappa'_k})=h(y_{k,i}/y_{k,\kappa_k})$. The image ideal is generated by these $h'$, and the value of
$\sum_mr_mh'_m$ is $\sum_mr_m(y'/y'_{\kappa'})h'_m(y'/y'_{\kappa'})$ with $r_m(y'/y'_{\kappa'})\in S'$. So a prime
$\mathfrak q$ of $S'$ contains the values of the image ideal if and only if $\mathfrak q\cap S$ contains those of
$J_\partial$, which by (3) for $y$ holds if and only if $\mathfrak q\cap S$, or equivalently $\mathfrak q$, contains one
of $z_1,\dots,z_n$. As $l$ was arbitrary, the $f_l$ and $a_l$, with the $z_i$ and $n$ of each $l$, give (3) for $y'$ at
$\kappa'$ with the same $e$.
\end{proof}
% Lean: SncRefine.isSncModel_refine. The chart c is (bx k).symm (c' (bk k)); exists_denominator is the paragraph on
% denominators (prime_chart_existsUnique_point, regAt_iff_chart_localization, hreg); smooth_coords_of_denominators with
% formallyEtale_of_denominators (isLocalization_atPrime, FormallyEtale.of_isLocalization, the local criteria
% formallyUnramified_iff_forall and smoothLocus_eq_univ_iff) and SncRes.coords_of_formallyEtale
% (tensorKaehlerEquivOfFormallyEtale) is the paragraph on smoothness; the last clause is the span induction on bd.map.

\begin{lemma}\label{cv-origin-reference-points}
Let $\Omega$, $X$, $X'$, $\xi$, $\xi'$, $g$, $y$, $y'$, $k\mapsto k'$ and $i\mapsto i'$ be as in
Lemma~\ref{cv-origin-reference-snc}, without the hypotheses on charts and on injectivity, and suppose
that the local ring of $X'$ at the image of $\xi'$ maps bijectively onto $\Omega$ and the local rings at
its specializations map injectively into $\Omega$. If $w\in\hat C$ are coordinates for $y'$ of a complex
point $P$ of $X'$ (Definition~\ref{cv-origin-reference}(4)), then $(w_{k',i'})_{k\in K,i\in I_k}$ are
coordinates for $y$ of the complex point $g(P)$ of $X$.
\end{lemma}
\leanref{Development/AutomorphicSpectrum/Proofs/GeometryProofs/ComplexCoordsReindexedStatement.lean}{GlobalTheory.T.complex_coords_reindexed}
\leanref{Development/AutomorphicSpectrum/Proofs/GeometryProofs/ComplexCoordsReindexed.lean}{GlobalTheory.CompanionTheory.Ar.CoordsDown.coordsAt_down}
\begin{proof}
For $k\in K$, some $w_{k',c'}$ is nonzero, since $w\in\hat C$, so the restricted point lies in $\hat C$.
Let $c\in I_k$ with $w_{k',c'}\neq0$. Then $y_{k,c}=y'_{k',c'}\neq0$, and for $i\in I_k$ the ratio
$y_{k,i}/y_{k,c}=y'_{k',i'}/y'_{k',c'}$ is regular at $P$ with value $w_{k',i'}/w_{k',c'}$. The point $P$
is complex, so the local homomorphism from the local ring at $P$ to $\mathbf C$ sends every prime to a unit,
every prime is a unit at $P$, and $P$ lies over no prime. Hence the local ring of $X'$ at
$P$ equals that of $X$ at $g(P)$, and the ratio is regular at $g(P)$. Its value at $g(P)$ is the image of
its germ under $\mathcal O_{X,g(P)}\to\mathcal O_{X',P}\to\mathbf C$; the germs at $g(P)$ and at $P$ have
the same image in $\Omega$, and $\mathcal O_{X',P}\to\Omega$ is injective, so they correspond, and the value
is $w_{k',i'}/w_{k',c'}$.
\end{proof}

\begin{lemma}\label{cv-origin-reference-embeds}
Let $X$ be a scheme, $\Omega$ a field of characteristic zero and $\xi:\operatorname{Spec}\Omega\to X$ a
morphism, and let $x=(x_i)_{i\in I}$ be a finite family in $\Omega$. Suppose that the local ring of $X$
at the image of $\xi$ maps bijectively onto $\Omega$ and the local rings at its specializations map
injectively into $\Omega$, that every point of $X$ that specializes the image of $\xi$ lies in a chart
$i$ of $x$ ($x_i\neq0$ and all $x_j/x_i$ regular at the point), and that for each $i$ with $x_i\neq0$ the
points in the chart $i$ form an affine open subscheme with ring $\mathbf Z[x_j/x_i:j\in I]$. Then the
block $x$ embeds the generic fibre in the sense of Definition~\ref{cv-origin-reference}(5).
\end{lemma}
\leanref{Development/AutomorphicSpectrum/Proofs/GeometryProofs/EmbedsGenericOfCharZeroStatement.lean}{GlobalTheory.T.embeds_generic_of_char_zero}
\leanref{Development/AutomorphicSpectrum/Proofs/GeometryProofs/EmbedsGenericOfCharZero.lean}{GlobalTheory.CompanionTheory.Ar.EmbedsOfTie.embedsGeneric_of_coordTie}
\begin{proof}
Let $P$ be a point over no prime. It lies in a chart $i$ by hypothesis. Let $i$ be such a chart, let
$U_i$ be its affine open subscheme with ring $A_i=\mathbf Z[x_j/x_i]$, and let $\mathfrak p\subset A_i$ be
the prime of $P$. The local ring of $X$ at $P$ is the localization $(A_i)_{\mathfrak p}$, so every element
$a$ of it is $r/s$ with $r,s\in A_i$ and $s\notin\mathfrak p$; then $s$ is a unit at $P$, and $s$ and
$sa=r$ lie in $A_i\subset B_i=\mathbf Q[x_j/x_i]$. For $s\in A_i$, $s$ is a unit at $P$ if and only if
$s\notin\mathfrak p$. So if $P'$ is a second point over no prime in the chart $i$ at which the same elements
of $B_i$ are units, then $P$ and $P'$ have the same prime in $A_i$, and $P=P'$ since $U_i$ is affine.
\end{proof}

\begin{lemma}\label{cv-origin-reference-exists}
In the setting of Definition~\ref{cv-origin-hodge-component}, let $\Omega_0$ be a field and
$\xi_0:\operatorname{Spec}\Omega_0\to\mathcal X$ a morphism as in Lemma~\ref{cv-origin-reference-char}.
Then there is a reference geometry of $\mathcal X$ (Definition~\ref{cv-origin-reference}) whose Hodge
component lies over the irreducible component of $\mathcal X$ with generic point the image of $\xi_0$.
\end{lemma}
\leanref{Development/AutomorphicSpectrum/Proofs/GeometryProofs/ExistsReferenceGeometrySameGenericStatement.lean}{GlobalTheory.T.exists_reference_geometry_same_generic}
\leanref{Development/AutomorphicSpectrum/Proofs/GeometryProofs/ExistsReferenceGeometrySameGeneric.lean}{GlobalTheory.CompanionTheory.Ar.exists_refGeo}
\begin{proof}
By Lemma~\ref{cv-origin-reference-char}, $\Omega_0$ has characteristic zero. Lemma~\ref{cv-origin-reference-coords}
gives coordinates $y$, a model $X_1$ of the component with a semi-abelian scheme, its Lie lattices
$L_s$ and a boundary ideal. Lemma~\ref{cv-origin-reference-flag} applied to these gives a field
$\Omega\supset\Omega_0$, coordinates $y'$, a scheme $X_2$ with $f_2:X_2\to\mathcal X$, generic bases
$b_s$ of degree $1$ and an ideal $J'_\partial$ with its properties (4) and (5). Lemma~\ref{cv-origin-reference-descent}
with $\deg=1$ gives coordinates $y''$ containing $y'$ as some of their blocks, a Hodge component $Z$ at
the empty set with coordinates $y''$ over the component of $\xi_0$, a morphism $g:Z\to X_2$ with
$f=f_2\circ g$, and a flag of $Z$ of degree $1$ with root lines of multidegrees $a''_{\tau,i}$ and
generators $e_{\tau,i}$.

Let $K=K''\sqcup\{k_0\}$, let $I_{k_0}=\prod_{k\in K''}I''_k$, and let $y\in\Omega^{\sqcup_{k\in K}I_k}$ be
$y''$ on the blocks of $K''$ and $y_{k_0,c}=\prod_{k\in K''}y''_{k,c_k}$ on the new block. A point of $Z$
lies in a chart $c$ of $y''$, and then in the chart $c$ of the block $k_0$, since
$y_{k_0,c'}/y_{k_0,c}=\prod_ky''_{k,c'_k}/y''_{k,c_k}$ is regular there. So
Lemma~\ref{cv-origin-reference-segre} shows that $Z$ is covered by the affine open subschemes of the charts
of $y$, and $Z$, with the same scheme, map, semi-abelian scheme, Lie algebra and integers, is a Hodge
component at the empty set with coordinates $y$, with the same flag. A chart $\kappa$ of $y$ restricts to
the chart $(\kappa_k)_{k\in K''}$ of $y''$, and the generator
$e_{\tau,i}\prod_{k\in K''}y_{k,\kappa_k}^{a''_{\tau,i,k}}$ does not involve the block $k_0$, so the root
lines have multidegrees $a_{\tau,i}=(a''_{\tau,i},0)$ with generators $e_{\tau,i}$. This gives (1) and (2)
of Definition~\ref{cv-origin-reference}.

For (3), let $J_\partial$ be the image of $J'_\partial$ under the renaming of the blocks of $y'$ into those
of $y$. Lemma~\ref{cv-origin-reference-snc}, applied to $g:Z\to X_2$, the local rings of
Lemma~\ref{cv-origin-reference-descent} and condition (4) of Lemma~\ref{cv-origin-reference-flag}, shows
that $J_\partial$ satisfies (3) for $y$ with $D=|\Psi|N(N-1)/2$. For (4), let $w$ be coordinates of a
complex point $P$ of $Z$. By Lemma~\ref{cv-origin-reference-points}, the restriction $w'$ of $w$ to the
blocks of $y'$ gives coordinates of $g(P)$. The polynomials of $J_\partial$ only involve the variables of
these blocks, so $w$ lies in the zero set of $J_\partial$ if and only if $w'$ lies in that of $J'_\partial$,
which by (5) of Lemma~\ref{cv-origin-reference-flag} holds if and only if $f_2(g(P))=f(P)$ lies outside the
abelian locus of $\mathcal A$. As $\mathcal A_Z=f^*\mathcal A$, this is the case if and only if
$\mathcal A_Z$ is not abelian at $P$.

For (5), the coordinates of the block $k_0$ are the products $\prod_ky''_{k,c_k}$, and the charts of
this block are the charts of $y''$. By Definition~\ref{cv-origin-hodge-component}(1) for $y''$, every point of
$Z$ lies in such a chart, and the chart $c$ is an affine open subscheme with ring
$\mathbf Z[y''_{k,i}/y''_{k,c_k}]$, which is the ring generated by the ratios $y_{k_0,c'}/y_{k_0,c}$: these
ratios are products of the $y''_{k,i}/y''_{k,c_k}$, and $y''_{k,i}/y''_{k,c_k}$ is the ratio for $c'$ equal
to $c$ with $c_k$ replaced by $i$. Lemma~\ref{cv-origin-reference-embeds}, with $x$ the coordinates of the
block $k_0$, gives (5).
\end{proof}
% Lean: exists_refGeo: component_generic_char_zero, exists_coords_for_component (with ResKStatement),
% coords_with_constants_map, generic_coords_descend_subfield with d = 1, the Segre block segY y' with
% k0 = Sum.inr () (segTie through coords_join_blocks; segDeg for the root lines), bd.map (rename v) with
% snc_model_reindexed_coords, bd_iff from complex_coords_reindexed, mem_mcPts_map and abLocus_pullback, and
% embeds_generic_of_char_zero applied to the CoordTie of the Segre products of y'.

\subsection{Analytic data and Hodge norms}\label{subsection-cv-analytic-data-hodge-norms}

\begin{definition}\label{cv-origin-analytic}
In the setting of Definition~\ref{cv-origin-reference}, let a reference geometry of $\mathcal X$ be given,
with its notation, and put $\Sigma=(\Psi\times\{1,\dots,N\})\sqcup\{A\}$. \emph{Analytic data} on it
consist of:
\begin{enumerate}
\item metrized line bundles $\bar\sigma_{\tau,i}=(a_{\tau,i},h_{\tau,i})$ that are Hodge metrics on the
  root lines with generators $e_{\tau,i}$ (Definition~\ref{cv-origin-hodge-component}(4));
\item an analytic system $S$ of type $(\Psi,N)$ (Definition~\ref{cv-proportionality-analytic-system}(1));
\item a compact complex manifold $\tilde X$ of dimension $D$ with, for each $x\in\tilde X$, an open
  neighbourhood $U_x$ and a holomorphic map $\lambda_x:U_x\to\hat X(\mathbf C)$, where the kernel of $y$ is
  an arithmetic subvariety $X$ of dimension $D+1$ (Definition~\ref{lit-multiprojective}), such that any two
  of these maps agree up to the action of $(\mathbf C^\times)^K$ where both are defined, every point of
  $\hat X(\mathbf C)$ lies in the orbit of some $\lambda_x(x)$, $\lambda_x(x)$ and $\lambda_{x'}(x')$ lie in
  one orbit only if $x=x'$, and the affine coordinates of $\lambda_x$ in every chart have injective
  derivative at $x$;
\item a bijective holomorphic map $\pi:\tilde X\to S$ of degree $1$ off a closed set of measure zero, as
  in Definition~\ref{cv-proportionality-analytic-system}(3);
\item for each $b\in\Sigma$ a smooth metrized line bundle $\bar y_b$ on $X$ whose first Chern form on
  $\tilde X$ is the pullback by $\pi$ of that of $\bar L_b$, and for each real $t$ a metrized line bundle
  $\bar y_b(t)$ of the multidegree of $\bar y_b$, smooth and invariant under conjugation, with
  $\log h_{\bar y_b(t)}(\lambda_x(x))=\log h_{\bar y_b}(\lambda_x(x))+\varphi_b(t)(\pi(x))$ for
  $0<t\le t_0$ and $x\in\tilde X$, where $h_{\bar L}$ is the function of a metrized line bundle $\bar L$.
\end{enumerate}
These satisfy:
\begin{enumerate}
\item[(i)] for $x\in\tilde X$, a complex point of $Z$ with coordinates $\lambda_x(x)$ is one over which
  $\mathcal A_Z$ is abelian if and only if $\pi(x)\notin\partial S$;
\item[(ii)] $\bar y_{\tau,i}$ has multidegree $-\sum_{k\le i}a_{\tau,k}$, and
  $\log h_{\bar y_{\tau,i}}(\lambda_x(x))+\varphi_{\tau,i}(\pi(x))=-\sum_{k\le i}\log h_{\tau,k}(\lambda_x(x))$
  for $x$ with $\pi(x)\notin\partial S$;
\item[(iii)] $\bar y_A$ has multidegree $1$ at $k_0$ and $0$ at the other indices, and
  $h_{\bar y_A}(z)=\big(\sum_{i\in I_{k_0}}|z_{k_0,i}|^2\big)^{1/2}$;
\item[(iv)] the bundle $\bar L_A$ of $S$ has a positive metric, and $\varphi_A=\varphi_A(t)=0$;
\item[(v)] condition (3c) of Definition~\ref{cv-origin-tower} holds for $S$: for every $\tau\in\Psi$ and every
  $t$, the regularized metric of $\bar L_{\tau,N}(t)$ has pluriharmonic local potentials.
\end{enumerate}
\end{definition}
\leanref{Development/AutomorphicSpectrum/Definitions/GlobalGeometry/EmbedsGeneric.lean}{GlobalTheory.CompanionTheory.Ar.RefGeo.An}
\leanref{Development/AutomorphicSpectrum/Definitions/MetricPeriods/MetricsProjectionFormulaNef.lean}{GlobalTheory.CompanionTheory.Ar.PB.Realis}
\leanref{Development/AutomorphicSpectrum/Definitions/GlobalGeometry/IsSncModel.lean}{GlobalTheory.CompanionTheory.Ar.Cmp.IsEmbedded}

\begin{definition}\label{cv-origin-hodge-norm-gram}
Let $T$ be a scheme, $G$ a commutative group scheme over $T$, and $\mathcal L$ a line bundle on $G$, trivialized on an
open cover $(U_\alpha)$ of $G$ with transition functions $g_{\alpha\beta}$. Let $z:\operatorname{Spec}\mathbf C\to T$ be a
complex point, and let tangent vectors be as in Definition~\ref{cv-origin-reference-lattices}(1).
\begin{enumerate}
\item A \emph{uniformization} of the fibre of $G$ at $z$ consists of a commutative group scheme $A$ over $\mathbf C$ with a
  homomorphism $\varphi:A\to G$ over $z$ that identifies $A$ with $G\times_T\operatorname{Spec}\mathbf C$; an integer
  $g\ge0$ and a surjective homomorphism $u:\mathbf C^g\to A(\mathbf C)$ whose kernel $\Lambda$ is discrete and spans
  $\mathbf C^g$ over $\mathbf R$, such that for every open $U\subset A$ and every regular function $f$ on $U$ the set
  $u^{-1}(U)$ is open and $f\circ u$ is holomorphic on it; and holomorphic functions $\theta_\alpha$ without zeros on
  $u^{-1}(\varphi^{-1}(U_\alpha))$ with $\theta_\alpha=(g_{\alpha\beta}\circ\varphi\circ u)\,\theta_\beta$ where both are
  defined, the \emph{trivializing functions} of $u$. It has \emph{Appell--Humbert data} $(H,\chi)$, for a Hermitian matrix $H\in M_g(\mathbf C)$ with
  $x^{\mathsf T}H\bar x>0$ for $x\neq0$ and a map $\chi:\Lambda\to\mathbf C$ with $|\chi|=1$, if
  \[
  \theta_\alpha(x)=\chi(l)\exp\big(\pi x^{\mathsf T}H\bar l+\tfrac\pi2\,l^{\mathsf T}H\bar l\big)\,\theta_\alpha(x+l)
  \]
  for every $\alpha$, $l\in\Lambda$ and $x\in u^{-1}(\varphi^{-1}(U_\alpha))$. A morphism
  $v':\operatorname{Spec}\mathbf C[\varepsilon]\to A$ over $\mathbf C$ has \emph{differential} $a\in\mathbf C^g$ if, for every
  open $U\subset A$ containing its image and every regular function $f$ on $U$, the coefficient of $\varepsilon$ in
  $f\circ v'$ is the derivative of $f\circ u$ at $0$ in the direction $a$.
\item Tangent vectors $v$ and $w$ of $G$ at $z$ have \emph{Hodge inner product} $c\in\mathbf C$ if there are a
  uniformization of the fibre of $G$ at $z$ with Appell--Humbert data $(H,\chi)$ and morphisms
  $v',w':\operatorname{Spec}\mathbf C[\varepsilon]\to A$ over $\mathbf C$ with differentials $a$ and $b$, such that
  $\varphi\circ v'=v$, $\varphi\circ w'=w$ and $c=\pi\,b^{\mathsf T}H\bar a$.
\item Let $\Omega$ be a field, $\xi:\operatorname{Spec}\Omega\to T$ a morphism with image $\eta$ and $\mathrm{Lie}$ the Lie
  algebra of $G$ at $\xi$, and suppose that the centre $P$ of $z$ specializes $\eta$. A tangent vector $v$ of $G$ at $z$ is
  the \emph{value at $z$} of $e\in\mathrm{Lie}$ if $e$ extends to $P$ through a tangent vector $v_P$ of $G$ at
  $\operatorname{Spec}\mathcal O_{T,P}\to T$ and $v=v_P\circ\operatorname{Spec}\psi[\varepsilon]$, where
  $\psi:\mathcal O_{T,P}\to\mathbf C$ is the homomorphism of $z$ and $\psi[\varepsilon]$ is $\psi$ with
  $\varepsilon\mapsto\varepsilon$. A family $b=(b_1,\dots,b_r)$ in $\mathrm{Lie}$ has \emph{Hodge Gram matrix}
  $\Gamma\in M_r(\mathbf C)$ at $z$ if there are values $v_p$ at $z$ of the $b_p$ such that $v_p$ and $v_q$ have Hodge
  inner product $\Gamma_{pq}$ for all $p$ and $q$.
\end{enumerate}
\end{definition}
\leanref{ReciprocityTheory/ReductiveGroups/ArithmeticGeometry/BirationalGeometry/ToroidalCompactification/IsLieBasis.lean}{ReciprocityTheory.ChartRestriction.G.Geo.Tor.HolUnif}
\leanref{ReciprocityTheory/ReductiveGroups/ArithmeticGeometry/BirationalGeometry/ToroidalCompactification/IsLieBasis.lean}{ReciprocityTheory.ChartRestriction.G.Geo.Tor.HolUnif.IsAH}
\leanref{ReciprocityTheory/ReductiveGroups/ArithmeticGeometry/BirationalGeometry/ToroidalCompactification/IsLieBasis.lean}{ReciprocityTheory.ChartRestriction.G.Geo.Tor.HolUnif.IsDiff}
\leanref{Development/AutomorphicSpectrum/Definitions/GlobalGeometry/HodgeIPIffHodgeFormConj.lean}{GlobalTheory.CompanionTheory.Ar.HodgeIP}
\leanref{Development/AutomorphicSpectrum/Definitions/GlobalGeometry/EigLat.lean}{GlobalTheory.CompanionTheory.Ar.LieModel.ValAt}
\leanref{Development/AutomorphicSpectrum/Definitions/GlobalGeometry/HodgeIPIffHodgeFormConj.lean}{GlobalTheory.CompanionTheory.Ar.LieModel.GramAt}
% Name (STYLE.md rule 13): the theta_alpha of (1), the Lean's HolUnif.fr, are the *trivializing functions* of u, as for
% the same functions in Lemma cv-uniformization-theta-half-sign-glue; "frame" is reserved for Chapter 5.

\begin{lemma}\label{cv-origin-hodge-norm-eigen}
In the setting of Definition~\ref{cv-origin-hodge-component}, let $Z$ be a Hodge component of $\mathcal X$ at $S$ with
coordinates $y$ with values in $\Omega$, with the action $\rho$ of $\mathcal O_E$ on $\mathrm{Lie}$, and for a point $P$ of
$Z$ let $L(P)$ be the Lie algebra of $\mathcal A_Z$ at $P$ (Definition~\ref{cv-origin-reference-lattices}(1)).
\begin{enumerate}
\item The field $\Omega$ has characteristic zero, the labels are distinct as embeddings of $\mathcal O_E$ into $\Omega$
  through $F_c\subset\Omega$, $\mathrm{Lie}$ is the direct sum of the $\mathrm{Lie}_s$ over the labels $s$, and $s(a)$ is
  regular at every point of $Z$ for every label $s$ and every $a\in\mathcal O_E$.
\item At every point $P$ of $Z$ over no prime, $\mathcal V_s(P)=L(P)\cap\mathrm{Lie}_s$ for every label $s$, and
  $L(P)=\sum_s\mathcal V_s(P)$.
\end{enumerate}
\end{lemma}
\leanref{Development/AutomorphicSpectrum/Proofs/GeometryProofs/HodgeCptInteriorGenericStatement.lean}{GlobalTheory.T.hodge_cpt_interior_generic}
\leanref{Development/AutomorphicSpectrum/Proofs/GeometryProofs/HodgeCptAwayFromPrimesStatement.lean}{GlobalTheory.T.hodge_cpt_away_from_primes}
\begin{proof}
(1) The embedding of $F_c$ in $\Omega$ is injective, so $\Omega$ has characteristic zero. The labels $\tau$ and $\bar\tau$,
for $\tau\in\Psi$, are distinct embeddings of $\mathcal O_E$ into $F_c$, and every embedding of $E$ into $\mathbf C$ is $\tau$
or $\bar\tau$ for exactly one $\tau\in\Psi$, as exactly one of the two maps $\sqrt{-d}$ to $\theta$. So there are
$[E:\mathbf Q]$ labels, and they remain distinct when composed with $F_c\to\Omega$. Hence $E\otimes_{\mathbf Q}\Omega$ is
the product of $[E:\mathbf Q]$ copies of $\Omega$, one for each label, and the module $\mathrm{Lie}$ over it is the direct
sum of the eigenspaces $\mathrm{Lie}_s$. For $a\in\mathcal O_E$, $s(a)$ is integral over $\mathbf Z$. The local rings of the
normal scheme $Z$ are integrally closed in its function field $\Omega$, so $s(a)$ is regular at every point of $Z$.

(2) The first assertion is the definition of $\mathcal V_s$ at a point over no prime of $S$
(Definition~\ref{cv-origin-hodge-component}(1)). Every prime number is a unit at $P$, hence so is every nonzero integer.
For labels $s\neq s'$ choose $a\in\mathcal O_E$ with $s(a)\neq s'(a)$. Then $x=s(a)-s'(a)$ is a nonzero algebraic integer,
regular at $P$ by (1); its inverse is a polynomial in $x$ with rational coefficients, whose denominators are units at $P$,
so $x^{-1}$ is regular at $P$. The module $L(P)$ is stable under $\mathcal O_E$: if $e$ extends to $P$ through $v$, then
$\rho(a)e$ extends to $P$ through $\alpha(a)\circ v$, for the action $\alpha$ of $\mathcal O_E$ on $\mathcal A_Z$
(Lemma~\ref{cv-origin-reference-lie}(3)).

Let $\Sigma$ be a set of labels and $f_s\in\mathrm{Lie}_s$, for $s\in\Sigma$, with $e=\sum_{s\in\Sigma}f_s\in L(P)$. We
show that every $f_s$ lies in $L(P)$, by induction on $|\Sigma|$. Let $\Sigma=\Sigma'\sqcup\{s'\}$ and $s\in\Sigma'$, and
choose $a$ with $s(a)\neq s'(a)$. The element
$\rho(a)e-s'(a)e=\sum_{u\in\Sigma'}(u(a)-s'(a))f_u$ lies in $L(P)$, as $s'(a)$ is regular at $P$, and
$(u(a)-s'(a))f_u\in\mathrm{Lie}_u$. By induction $(s(a)-s'(a))f_s\in L(P)$, and so $f_s\in L(P)$, as
$(s(a)-s'(a))^{-1}$ is regular at $P$. Then $f_{s'}=e-\sum_{s\in\Sigma'}f_s\in L(P)$ too. Applied to the decomposition of
(1) of an element of $L(P)$, this gives $L(P)\subset\sum_s(L(P)\cap\mathrm{Lie}_s)=\sum_s\mathcal V_s(P)$; the other
inclusion is clear.
\end{proof}
% Lean: (1) is part of HodgeCpt.basic (charZero_of_injective_ringHom, HodgeBasic.labelEmb_injective,
% HodgeBasic.labelHom_bijective with Embeddings.card, T.eigenspace_decomposition_internal, HodgeBasic.mem_regAt_of_isIntegral).
% (2) is HodgeChar0.lattices_of_char_zero: hV by Lat.patch_of_not, hlat by mem_of_sum_mem (the induction of the text) with
% lat_stable, unitAt_natCast and inv_mem_regAt (Subalgebra.inv_mem_of_algebraic over Q). Size: about 540 and 930 lines.

\begin{lemma}\label{cv-origin-hodge-norm-values}
Let $T$ be a scheme, $\Omega$ a field and $\xi:\operatorname{Spec}\Omega\to T$ a morphism with image $\eta$, such that the
local ring at $\eta$ maps bijectively onto $\Omega$ and the local rings at the specializations of $\eta$ map injectively into
$\Omega$. Let $G$ be a commutative group scheme over $T$, smooth over $T$, let $\mathrm{Lie}$ be its Lie algebra at $\xi$, and
for a point $P$ that specializes $\eta$ let $L(P)\subset\mathrm{Lie}$ be the Lie algebra of $G$ at $P$ and
$\mathcal O_P\subset\Omega$ the ring of elements regular at $P$ (Definition~\ref{cv-origin-reference-lattices}). Let $z$ be a
complex point of $T$ whose centre $P$ specializes $\eta$, let $\psi:\mathcal O_P\to\mathbf C$ be the value at $z$, and let
$T_z$ be the Lie algebra of $G$ at $z$. Then there is an additive map $e\mapsto e(z)$ from $L(P)$ to $T_z$ such that:
\begin{enumerate}
\item for $e\in\mathrm{Lie}$ and $v\in T_z$, $v$ is the value of $e$ at $z$ (Definition~\ref{cv-origin-hodge-norm-gram}(3)) if
  and only if $e\in L(P)$ and $v=e(z)$;
\item $(ae)(z)=\psi(a)e(z)$ for $a\in\mathcal O_P$ and $e\in L(P)$;
\item the $e(z)$, for $e\in L(P)$, span $T_z$ over $\mathbf C$;
\item for every $\Omega$-linear form $\pi$ on $\mathrm{Lie}$ with $\pi(L(P))\subset\mathcal O_P$ there is a $\mathbf C$-linear
  form $g$ on $T_z$ with $g(e(z))=\psi(\pi(e))$ for every $e\in L(P)$.
\end{enumerate}
\end{lemma}
\leanref{Development/AutomorphicSpectrum/Proofs/GeometryProofs/LieLatticeEvalAtLocalRingStatement.lean}{GlobalTheory.T.lie_lattice_eval_at_local_ring}
\leanref{Development/AutomorphicSpectrum/Proofs/GeometryProofs/LatticeRegularAtLocalPointStatement.lean}{GlobalTheory.T.lattice_regular_at_local_point}
\begin{proof}
Let $B=\mathcal O_{T,P}$, identified with $\mathcal O_P$, let $x:\operatorname{Spec}B\to T$ be the canonical morphism and
$L_x$ the Lie algebra of $G$ at $x$. The morphisms $\xi$ and $z$ factor through $x$, by the homomorphisms $\iota:B\to\Omega$ and
$\psi:B\to\mathbf C$. By Lemma~\ref{cv-origin-reference-lie-bundle-base}, $L_x$ is a finite projective, hence free, $B$-module,
the maps $i:L_x\to\mathrm{Lie}$, $v\mapsto v\circ\operatorname{Spec}\iota[\varepsilon]$, and $j:L_x\to T_z$,
$v\mapsto v\circ\operatorname{Spec}\psi[\varepsilon]$, are $\iota$-semilinear and $\psi$-semilinear, and
$\mathbf C\otimes_BL_x\to T_z$ is an isomorphism. The map $i$ is injective: two tangent vectors at $x$ with the same image
in $\mathrm{Lie}$ agree after $B[\varepsilon]\to\Omega[\varepsilon]$ and both restrict to the unit section over
$\operatorname{Spec}B$, and $B[\varepsilon]\to B\times\Omega[\varepsilon]$ is injective, so they are equal. By
Definition~\ref{cv-origin-reference-lattices}(1), the vectors of $\mathrm{Lie}$ that extend to $P$ are the elements of
$i(L_x)$, an $\mathcal O_P$-submodule, so $L(P)=i(L_x)$ and $i:L_x\to L(P)$ is bijective. Put $e(z)=j(i^{-1}(e))$.

(1) is the definition of the value at $z$, as the tangent vector $v_P$ through which $e$ extends is $i^{-1}(e)$. (2) holds
because $i$ and $j$ are semilinear. (3) The images under $j$ of a basis of $L_x$ span $T_z$, as
$\mathbf C\otimes_BL_x\to T_z$ is an isomorphism. (4) The map $\pi\circ i:L_x\to\mathcal O_P=B$ is $B$-linear, so
$\psi\circ\pi\circ i$ induces a $\mathbf C$-linear form on $\mathbf C\otimes_BL_x$; through the isomorphism with $T_z$ it is
the form $g$, and $g(j(v))=\psi(\pi(i(v)))$ for $v\in L_x$.
\end{proof}
% Lean: (1)-(3) are lieLat_val, for any local ring R in place of C: lieModel_exists at the stalk, LieVal.exists_base
% (T.lie_model_base_change_map) for i and j, hinj through Extends.unique (LieCarrier.tan_ext, LieEx.joint_epi), hsurj and
% the AddEquiv q, and T.lie_model_free_over_local_ring with the spanning clause of exists_base. (4) is
% RootHt.exists_lieFunctional: the stalk-linear map j, T.lie_model_isBaseChange, and IsBaseChange.lift of f'. The base
% change results are those of Lemma cv-origin-reference-lie-bundle-base. Size: about 4,400 and 4,600 lines (overlapping).

\begin{definition}\label{cv-origin-hodge-norm-form-ample}
Let $S$ be a scheme, $\pi:G\to S$ a morphism of schemes and $\mathcal L$ a line bundle on $G$. Then $\mathcal L$ is
\emph{relatively ample by sections} if for every affine open $U\subset S$ there are an integer $N\ge1$ and sections
$s_0,\dots,s_M$ of $\mathcal L^{\otimes N}$ over $\pi^{-1}(U)$ such that the open sets $G_{s_i}\subset\pi^{-1}(U)$ on which
$s_i$ does not vanish cover $\pi^{-1}(U)$, and for each $i$ the scheme $G_{s_i}$ is affine and its ring is generated over
$\Gamma(U,\mathcal O_S)$ by regular functions $r_{i0},\dots,r_{iM}$ on $G_{s_i}$ with $s_j=r_{ij}s_i$ on $G_{s_i}$ for every $j$.
\end{definition}
\leanref{ReciprocityTheory/ReductiveGroups/ArithmeticGeometry/ExtensionDescent/NonzeroLocus.lean}{ReciprocityTheory.ChartRestriction.G.Out.RelAmple}

\begin{lemma}\label{cv-origin-hodge-norm-form-descent-ample-affine}
Let $B$ be a ring, $X$ a scheme with a universally closed morphism $f:X\to\operatorname{Spec}B$, and $H$ a line bundle on $X$.
For a section $a$ of $H$ over $X$, let $X_a$ be the open set where $a$ does not vanish. Let $s_1,\dots,s_m$ be sections of
$H$ over $X$ such that the sets $X_{s_j}$ are affine and cover $X$. Then $X_t$ is affine for every section $t$ of $H$ over
$X$.
\end{lemma}
\leanref{Development/AbelianVarietyTheory/Proofs/UniformizationProofs/IsAffineOpenNonzeroLocusOfUniversallyClosed.lean}{GlobalTheory.isAffineOpen_nonzeroLocus_of_universallyClosed}
\leanref{Development/AbelianVarietyTheory/Proofs/ThetaShiftProofs/PolarizedAbelianEmbeddingBySections.lean}{GlobalTheory.AmpleSpan.ker_sup_span_eq_top}
\begin{proof}
For sections $a$ and $b$ of $H$ there is a unique function $a/b$ on $X_b$ with $a=(a/b)\,b$ on $X_b$: in a trivialization of
$H$ over an open set $V$, $b$ is a function that is a unit on $V\cap X_b$, and the quotients of $a$ by $b$ on these sets glue.
In such a trivialization $a/b$ is a unit at a point of $X_b$ if and only if $a$ is, so the set of points of $X_b$ where $a/b$
does not vanish is $X_a\cap X_b$.

\emph{Quasi-compactness.} The set $X_t$ is the union of the sets $X_{s_j}\cap X_t$. Each of them is the set of points of the
affine scheme $X_{s_j}$ where $t/s_j$ does not vanish, a basic open subset, so it is affine and quasi-compact. Hence $X_t$ is
quasi-compact.

\emph{A unit ideal.} Let $r_j=s_j/t$, functions on $X_t$. Let $P=B[T_1,\dots,T_m]$, and let $\psi:P\to\Gamma(X_t,\mathcal O_X)$
be the homomorphism that is $f^\#$ on $B$ and maps $T_j$ to $r_j$. We show that $\ker\psi+(T_1,\dots,T_m)=P$. Let
$q:X_t\to\operatorname{Spec}P$ be the morphism given by $\psi$, so that the restriction of $f$ to $X_t$ is the composite of $q$
with $\operatorname{Spec}P\to\operatorname{Spec}B$. Let $X_P=X\times_{\operatorname{Spec}B}\operatorname{Spec}P$, with
projections $p_1$ and $p_2$, and let $d_j=p_1^*s_j-T_j\,p_1^*t$, sections of $p_1^*H$. Let $O\subset X_P$ be the union of the
sets where the $d_j$ do not vanish, and $C=p_2(X_P\setminus O)$. As $p_2$ is a base change of $f$, it is closed, and so $C$ is
closed in $\operatorname{Spec}P$.

First, $q(X_t)\subset C$. Let $\gamma:X_t\to X_P$ be the morphism with components the inclusion of $X_t$ and $q$. Then
$\gamma^*d_j=s_j-r_jt=0$ on $X_t$, so for $z\in X_t$ the point $\gamma(z)$ is not in $O$, and $q(z)=p_2(\gamma(z))$ lies in
$C$.

Suppose that $\ker\psi+(T_1,\dots,T_m)$ is contained in a maximal ideal $\mathfrak m$ of $P$, and let $x$ be its point of
$\operatorname{Spec}P$. If $x\notin C$, there is $a\in P$ with $x\in D(a)$ and $D(a)\cap C=\emptyset$. The set where $\psi(a)$
does not vanish on $X_t$ is $q^{-1}(D(a))$, which is empty by the previous paragraph. A function on the quasi-compact scheme
$X_t$ that vanishes at every point is nilpotent, so $\psi(a^e)=0$ for some $e\ge1$; then $a^e\in\mathfrak m$, so
$a\in\mathfrak m$, which contradicts $x\in D(a)$. Hence $x\in C$: there is a point $z$ of $X_P\setminus O$ with $p_2(z)=x$.
The point $p_1(z)$ lies in some $X_{s_j}$. In a trivialization of $H$ near $p_1(z)$, $p_1^*s_j=d_j+T_j\,p_1^*t$. Here $d_j$ is
not a unit at $z$, as $z\notin O$, and $T_j$ vanishes at $z$, as $T_j\in\mathfrak m$ and $p_2(z)=x$. So $p_1^*s_j$ is not a
unit at $z$, that is, $p_1(z)\notin X_{s_j}$, a contradiction. So $\ker\psi+(T_1,\dots,T_m)=P$.

\emph{Affineness.} Write $1=k+\sum_jc_jT_j$ with $k\in\ker\psi$ and $c_j\in P$. Applying $\psi$ gives
$1=\sum_j\psi(c_j)\,r_j$, so $r_1,\dots,r_m$ generate the unit ideal of $\Gamma(X_t,\mathcal O_X)$. The set of points of $X_t$
where $r_j$ does not vanish is $X_{s_j}\cap X_t$, which is affine by the first step. So $X_t$ is affine
\cite[Tag 01QF]{Sta}.
\end{proof}
% Lean: isAffineOpen_nonzeroLocus_of_universallyClosed: AmpleSpan.exists_isRatio, IsRatio.unique and IsRatio.basicOpen for
% the ratios; AmpleSpan.isCompact_nonzeroLocus; AmpleSpan.ker_sup_span_eq_top with t = s and c the identity matrix (the
% Lean allows s_j = sum_l c_jl t_l with c_jl in B; the text takes c the identity, which is the only use): P, psi, q from
% exists_hom_spec, the fibre product of f and Spec P -> Spec B, the sections d, O and the closed set C (isClosedMap of the
% universally closed p2), h1 through pullback.lift, the maximal ideal m, Scheme.isNilpotent_iff_basicOpen_eq_bot_of_isCompact
% (Mathlib), and the contradiction key at a point of C over m. Then Ideal.map of the identity under psi and
% isAffineOpen_of_isAffineOpen_basicOpen (Mathlib; Stacks, Tag 01QF, checked). About 320 lines with the ratios.

\begin{lemma}\label{cv-origin-hodge-norm-form-descent-ample-embedding}
Let $R$ be a ring, $X$ a scheme proper over $\operatorname{Spec}R$ and $\mathcal L$ a line bundle on $X$. For a section $a$ of a
power of $\mathcal L$ over $X$, let $X_a$ be the open set where $a$ does not vanish. Suppose that every point of $X$ lies in an
affine set $X_s$ for a section $s$ of some $\mathcal L^{\otimes n}$ with $n\ge1$. Then there are $n\ge1$ and sections
$t_0,\dots,t_M$ of $\mathcal L^{\otimes n}$ over $X$ such that the sets $X_{t_l}$ cover $X$, and for each $l$ the scheme
$X_{t_l}$ is affine and its ring is generated over $R$ by the functions $t_j/t_l$, $0\le j\le M$, where $t_j/t_l$ is the
function on $X_{t_l}$ with $t_j=(t_j/t_l)\,t_l$ there.
\end{lemma}
\leanref{Development/AbelianVarietyTheory/Proofs/CovolumeBounds/IsProjOverIffIsProperAndExistsIsAmpleLB.lean}{GlobalTheory.IsAmpleLB.exists_veryAmpleBy}
\leanref{Development/AbelianVarietyTheory/Proofs/CovolumeBounds/IsProjOverIffIsProperAndExistsIsAmpleLB.lean}{GlobalTheory.IsAmpleLB.exists_isImmersion_mul}
\begin{proof}
The functions $t_j/t_l$ exist and are unique, as in the proof of Lemma~\ref{cv-origin-hodge-norm-form-descent-ample-affine}.

\emph{A cover.} The scheme $X$ is quasi-compact, as it is proper over $\operatorname{Spec}R$. So finitely many of the affine
sets of the hypothesis cover $X$, given by sections $s'_i$ of $\mathcal L^{\otimes n_i}$, $0\le i\le c$. Let $n=\prod_in_i$ and
$s_i=s_i'^{\,n/n_i}$, sections of $\mathcal L^{\otimes n}$ with $X_{s_i}=X_{s'_i}$. The scheme $X$ is quasi-separated: for affine
open sets $U$ and $V$ of $X$, $U\cap V$ is the union of the sets $(X_{s_i}\cap U)\cap(X_{s_i}\cap V)$, where $X_{s_i}\cap U$ and
$X_{s_i}\cap V$ are affine \cite[Tag 01PV]{Sta} open subsets of the affine scheme $X_{s_i}$, so that their intersection is
quasi-compact.

\emph{Generators.} For each $i$, the ring of $X_{s_i}$ is a finitely generated $R$-algebra \cite[Tag 01T2]{Sta}; let
$g_{i1},\dots,g_{im_i}$ generate it. As $X$ is quasi-compact and quasi-separated, for each $i$ and $j$ there is $N_{ij}$ such
that for every $k\ge N_{ij}$ the section $g_{ij}s_i^k$ of $\mathcal L^{\otimes nk}$ over $X_{s_i}$ is the restriction of a section of
$\mathcal L^{\otimes nk}$ over $X$ \cite[Tag 01PW]{Sta}. Fix $k\ge1$ with $k\ge N_{ij}$ for all $i$ and $j$, and let $u_{ij}$ be
a section of $\mathcal L^{\otimes nk}$ with $u_{ij}=g_{ij}s_i^k$ on $X_{s_i}$.

\emph{The morphism to projective space.} Let $t_0,\dots,t_M$ be the sections $s_i^k$ and $u_{ij}$ of $\mathcal L^{\otimes nk}$,
in some order, and let $l(i)$ be the index with $t_{l(i)}=s_i^k$. As $X_{s_i^k}=X_{s_i}$, the sets $X_{t_l}$ cover $X$. Let
$\varphi:X\to\mathbf P^M_R$ be the morphism over $R$ obtained by gluing the morphisms
$X_{t_l}\to D_+(T_l)=\operatorname{Spec}R[T_0/T_l,\dots,T_M/T_l]$ given by $T_j/T_l\mapsto t_j/t_l$; it is the morphism
given by the generating sections $t_0,\dots,t_M$ \cite[Tag 01NE]{Sta}. A point $y\in X_{t_{l'}}$ lies in
$\varphi^{-1}(D_+(T_l))$ if and only if $t_l/t_{l'}$ is a unit at $y$, that is, $y\in X_{t_l}$; so
$\varphi^{-1}(D_+(T_l))=X_{t_l}$.

\emph{An immersion.} Over $D_+(T_{l(i)})$, $\varphi$ is the morphism $X_{s_i}\to\operatorname{Spec}R[T_j/T_{l(i)}]$ given by
$T_j/T_{l(i)}\mapsto t_j/s_i^k$. Its ring map is surjective, as its image contains the functions $u_{ij}/s_i^k=g_{ij}$, so it
is a closed immersion. The union $W$ of the sets $D_+(T_{l(i)})$ is open in $\mathbf P^M_R$ with $\varphi^{-1}(W)=X$, so
$\varphi:X\to W$ is a closed immersion \cite[Tag 01QO]{Sta}, and $\varphi$ is an immersion.

\emph{A closed immersion.} As $X$ is proper over $\operatorname{Spec}R$ and $\mathbf P^M_R$ is separated over
$\operatorname{Spec}R$, $\varphi$ is proper and $\varphi(X)$ is closed \cite[Tag 01W6]{Sta}. So $\varphi$ is a closed immersion
\cite[Tag 01IQ]{Sta}. By \cite[Tag 01QO]{Sta}, for every $l$ the inverse image $X_{t_l}$ of the affine open set $D_+(T_l)$ is
affine and the map $R[T_j/T_l]\to\Gamma(X_{t_l},\mathcal O_X)$, $T_j/T_l\mapsto t_j/t_l$, is surjective. With $nk$ in place of
$n$, this is the claim.
\end{proof}
% Lean: IsAmpleLB.exists_veryAmpleBy = exists_veryAmpleBy_mul: IsAmpleLB.exists_isImmersion_mul (h.compactSpace;
% IsAmpleLB.quasiSeparatedSpace with isAffineOpen_nonzeroLocus_inf, which is Stacks, Tag 01PV, checked; IsAmpleLB.exists_cover
% with exists_common_pow; exists_generators, Stacks, Tag 01T2, checked; exists_sec_eq_fn_mul_pow = exists_sec_eq_pow_mul with
% m = 0, which is Stacks, Tag 01PW, checked, in the form "for all k >= N", proved its own way on a finite affine refinement;
% isImmersion_toProj with toProj = IsCoordFamily.toHom, toProj_preimage_chart, surjective_appLE_chart_iff and
% IsZariskiLocalAtTarget.of_range_subset_iSup), then isClosedImmersion_of_isProper (IsProper.of_comp and
% IsClosedImmersion.of_isPreimmersion; Stacks, Tags 01W6 and 01IQ, checked) and veryAmpleBy_iff through
% isClosedImmersion_toProj_iff = isClosedImmersion_iff_pullRatio (Stacks, Tag 01QO, checked). Stacks, Tag 01NE (checked) is
% the published form of toProj; the Lean builds it by gluing, as the text does. About 900 lines with
% IsClosedImmersionIffPullRatio and ExistsSecNonzeroLocusEqBasicOpen.

\begin{lemma}\label{cv-origin-hodge-norm-form-descent-ample}
Let $R$ be a ring, $G$ a scheme proper over $\operatorname{Spec}R$ and $\mathcal L$ a line bundle on $G$. For a section $s$ of
a power of $\mathcal L$ over $G$, let $G_s$ be the open set where $s$ does not vanish. Say that $\mathcal L$ is ample if for every
point $y$ of $G$ there are $n\ge1$ and a section $s$ of $\mathcal L^{\otimes n}$ over $G$ with $y\in G_s$ and $G_s$ affine
\cite[Tag 01PS]{Sta}.
\begin{enumerate}
\item $\mathcal L$ is relatively ample by sections (Definition~\ref{cv-origin-hodge-norm-form-ample}) if and only if it is
  ample.
\item If $\mathcal L$ is ample, then $G_s$ is affine for every $n\ge1$ and every section $s$ of $\mathcal L^{\otimes n}$ over
  $G$.
\end{enumerate}
\end{lemma}
\leanref{Development/AbelianVarietyTheory/Proofs/CovolumeBounds/IsProjOverIffIsProperAndExistsIsAmpleLB.lean}{GlobalTheory.relAmple_iff_isAmpleLB_of_isAffine}
\leanref{Development/AbelianVarietyTheory/Definitions/ThetaMetrics/IsAmpleLBOfRelAmpleOfIsAffine.lean}{GlobalTheory.isAmpleLB_of_relAmple_of_isAffine}
\leanref{Development/AbelianVarietyTheory/Proofs/UniformizationProofs/IsAffineOpenNonzeroLocusOfUniversallyClosed.lean}{GlobalTheory.IsAmpleLB.isAffineOpen_nonzeroLocus}
\begin{proof}
If $\mathcal L$ is relatively ample by sections, then Definition~\ref{cv-origin-hodge-norm-form-ample} for the affine open
$\operatorname{Spec}R$ gives $N\ge1$ and sections $s_i$ of $\mathcal L^{\otimes N}$ over $G$ whose sets $G_{s_i}$ are affine
and cover $G$; so $\mathcal L$ is ample.

Conversely, let $\mathcal L$ be ample, let $\pi:G\to\operatorname{Spec}R$ be the structure morphism, $U\subset\operatorname{Spec}R$
an affine open set and $G_U=\pi^{-1}(U)$. The inclusion of $G_U$ in $G$ is a base change of $U\to\operatorname{Spec}R$, a
morphism of affine schemes, so it is affine. Hence for a section $s$ of $\mathcal L^{\otimes n}$ over $G$ with $G_s$ affine, the
set $G_s\cap G_U$ where the restriction of $s$ to $G_U$ does not vanish is affine, and every point of $G_U$ lies in such a set.
The scheme $G_U$ is proper over $U=\operatorname{Spec}\Gamma(U,\mathcal O)$. So
Lemma~\ref{cv-origin-hodge-norm-form-descent-ample-embedding}, applied to the ring $\Gamma(U,\mathcal O)$, the scheme $G_U$ and the
restriction of $\mathcal L$, gives $N\ge1$ and sections $s_0,\dots,s_M$ of $\mathcal L^{\otimes N}$ over $G_U$ whose sets $G_{s_i}$
cover $G_U$ and are affine, with rings generated over $\Gamma(U,\mathcal O)$ by the functions $r_{ij}=s_j/s_i$. These satisfy
$s_j=r_{ij}s_i$ on $G_{s_i}$. So $\mathcal L$ is relatively ample by sections. This proves (1).

For (2), let $\mathcal L$ be ample, $n\ge1$ and $s$ a section of $\mathcal L^{\otimes n}$ over $G$. As in the first step of the
proof of Lemma~\ref{cv-origin-hodge-norm-form-descent-ample-embedding}, there are $N\ge1$ and sections $u_0,\dots,u_M$ of
$\mathcal L^{\otimes N}$ over $G$ whose sets $G_{u_i}$ are affine and cover $G$. The sections $u_i^n$ and $s^N$ of
$\mathcal L^{\otimes Nn}$ have $G_{u_i^n}=G_{u_i}$ and $G_{s^N}=G_s$. As $G$ is proper over $\operatorname{Spec}R$, the
morphism $G\to\operatorname{Spec}R$ is universally closed, and Lemma~\ref{cv-origin-hodge-norm-form-descent-ample-affine},
applied to $H=\mathcal L^{\otimes Nn}$, the sections $u_i^n$ and $t=s^N$, shows that $G_s$ is affine.
% Lean: the converse is relAmple_of_isAmpleLB (through relAmple_iff_isAmpleLB_of_isAffine, with IsAmpleLB.pull_of_isPullback
% = pull_of_isAffineHom for each affine open of the base, the base change of U.iota being affine): IsAmpleLB.exists_veryAmpleBy
% (Lemma cv-origin-hodge-norm-form-descent-ample-embedding) for A.hom restricted to U composed with hU.isoSpec, then
% exists_veryAmpleBy_of_lbIso with lbIso_lbPow_finPow (the Lean's L^N over the inverse image of U is powOver) and
% veryAmpleBy_comp_iff (the isomorphism of U with Spec of its ring). (2) is IsAmpleLB.isAffineOpen_nonzeroLocus for G
% universally closed over an affine base: IsAmpleLB.exists_cover_mul with k = n, secPow s N, and
% isAffineOpen_nonzeroLocus_of_universallyClosed (Lemma cv-origin-hodge-norm-form-descent-ample-affine). Size: about 1,600
% lines with the two lemmas before this one (IsProjOverIffIsProperAndExistsIsAmpleLB, IsAmpleLBOfRelAmpleOfIsAffine,
% IsAffineOpenNonzeroLocusOfUniversallyClosed, IsClosedImmersionIffPullRatio and part of PolarizedAbelianEmbeddingBySections).
\end{proof}
% Lean: the first direction is isAmpleLB_of_relAmple_of_isAffine (isAmpleLB_of_relAmple for the affine open top, and the
% isomorphism of the inverse image of top with A). The text's "ample" is the Lean's IsAmpleLB (with CompactSpace, which holds
% as G is proper over an affine base).

\begin{lemma}\label{cv-origin-hodge-norm-form-descent}
Let $j:X'\to X$ be a morphism of schemes, and let $T'$ and $T$ be $\mathcal O_E$-semi-abelian schemes over $X'$ and $X$
(Definition~\ref{cv-model-toroidal-semiabelian}(2)) with a morphism $T'\to T$ over $j$. Suppose that $T'$ is proper over
$X'$ and that its line bundle is relatively ample by sections (Definition~\ref{cv-origin-hodge-norm-form-ample}) and defines
a principal polarization. Let $K$ be a field, $x:\operatorname{Spec}K\to X$ a morphism whose image lies in $j(X')$, and $T_x$
an $\mathcal O_E$-semi-abelian scheme over $\operatorname{Spec}K$ with a morphism $T_x\to T$ over $x$. Then $T_x$ is proper
over $K$, and its line bundle is relatively ample by sections and defines a principal polarization.
\end{lemma}
\leanref{Development/AbelianVarietyTheory/Proofs/UniformizationProofs/ExistsPolAVGenericFibre.lean}{GlobalTheory.GenericPol.proper_ample_kTriv}
\leanref{Development/AbelianVarietyTheory/Proofs/UniformizationProofs/ProperRelAmpleKTrivOfFieldExtension.lean}{GlobalTheory.AmpleDescent.proper_relAmple_kTriv_of_fieldExtension}
\leanref{Development/AbelianVarietyTheory/Proofs/UniformizationProofs/ProperRelAmpleKTrivOfFieldExtension.lean}{GlobalTheory.AmpleDescent.isAmpleLB_of_pull}
\leanref{ReciprocityTheory/ReductiveGroups/ArithmeticGeometry/BirationalGeometry/ToroidalCompactification/AmplePt.lean}{ReciprocityTheory.ChartRestriction.G.Geo.Tor.relAmple_pull}
\leanref{ReciprocityTheory/ReductiveGroups/ArithmeticGeometry/BirationalGeometry/ToroidalCompactification/AmplePt.lean}{ReciprocityTheory.ChartRestriction.G.Geo.Tor.kTriv_pull}
\leanref{Development/AbelianVarietyTheory/Proofs/ModuliProofs/RelAmpleKTrivOfFiniteEtaleCover.lean}{GlobalTheory.PolDescent.kTriv_of_pull}
\begin{proof}
The image of $x$ lies in $j(X')$, so the fibre product $\operatorname{Spec}K\times_XX'$ has a point. Let $K'$ be its residue
field. This gives a field extension $i:K\to K'$ and a morphism $x':\operatorname{Spec}K'\to X'$ with
$j\circ x'=x\circ\operatorname{Spec}i$. Let $T_{K'}$ be the base change of $T_x$ along $\operatorname{Spec}i$, an
$\mathcal O_E$-semi-abelian scheme over $K'$ with a morphism $T_{K'}\to T_x$ over $\operatorname{Spec}i$. Composed with
$T_x\to T$, it is a morphism $T_{K'}\to T$ over $j\circ x'$. As $T'\to T$ identifies $T'$ with $j^*T$, compatibly with the
actions and the line bundles, this morphism factors through a morphism $T_{K'}\to T'$ over $x'$, by the universal property
of the fibre product, and $T_{K'}\to T'$ identifies $T_{K'}$ with $x'^*T'$. So $T_{K'}$ is proper over $K'$, as a base change
of $T'$.

Here the line bundle $\mathcal L$ of an $\mathcal O_E$-semi-abelian scheme $G$ over a scheme $S$ defines a principal
polarization in the following sense: for every ring $R$, morphism $\operatorname{Spec}R\to S$ and point $y$ of $G$ over
$\operatorname{Spec}R$, if the line bundle $m_y^*\mathcal L\otimes\operatorname{pr}_2^*\mathcal L^{-1}$ on
$\operatorname{Spec}R\times_SG$, where $m_y(t,a)=y(t)+a$, is trivial over the inverse images of the sets of an open cover of
$\operatorname{Spec}R$, then $y=0$.

\emph{Ampleness over $K'$.} Let $U\subset X'$ be an affine open containing the point $x'$, and let $N$, the sections $s_i$ of
$\mathcal L'^{\otimes N}$ over the inverse image $T'_U$ of $U$ and the functions $r_{ij}$ be given by
Definition~\ref{cv-origin-hodge-norm-form-ample} for $T'$ and $U$. The morphism $\psi:T_{K'}\to T'_U$ is a base change of
$\operatorname{Spec}K'\to U$, so it is affine. The pullbacks $\psi^*s_i$ are sections of the $N$-th power of the line bundle of
$T_{K'}$; the set where $\psi^*s_i$ does not vanish is $\psi^{-1}$ of the set $T'_{s_i}$ where $s_i$ does not vanish, an affine
scheme whose ring is $\Gamma(T'_{s_i},\mathcal O)\otimes_{\Gamma(U,\mathcal O)}K'$; these sets cover $T_{K'}$, the ring is generated
over $K'$ by the pullbacks of the $r_{ij}$, and $\psi^*s_j=\psi^*(r_{ij})\,\psi^*s_i$. So the line bundle of $T_{K'}$ is
relatively ample by sections.

\emph{The polarization over $K'$.} Let $y$ be a point of $T_{K'}$ over $\operatorname{Spec}R\to\operatorname{Spec}K'$. Through
$\psi$, it is a point of $T'$ over $\operatorname{Spec}R\to X'$, and its line bundle $m_y^*\mathcal L\otimes\operatorname{pr}_2^*\mathcal L^{-1}$ is
the pullback of the corresponding line bundle for $T'$ under the isomorphism
$\operatorname{Spec}R\times_{K'}T_{K'}\cong\operatorname{Spec}R\times_{X'}T'$. If one is trivial over the inverse images of an
open cover of $\operatorname{Spec}R$, so is the other. Then $\psi\circ y=0$, as $\mathcal L'$ defines a principal polarization,
and $y=0$, as $T_{K'}$ is the fibre product of $\operatorname{Spec}K'$ and $T'$ over $X'$.

\emph{Descent to $K$: properness.} The scheme $T_x$ is separated, smooth and quasi-compact over $K$, as $T_x$ is semi-abelian.
Its base change $T_{K'}$ is universally closed over $K'$, and $\operatorname{Spec}K'\to\operatorname{Spec}K$ is flat, surjective
and quasi-compact, so $T_x$ is universally closed over $K$ \cite[Tag 02KS]{Sta}, hence proper.

\emph{Descent to $K$: the polarization.} Let $y$ be a point of $T_x$ over $\operatorname{Spec}R\to\operatorname{Spec}K$ whose line
bundle $m_y^*\mathcal L_x\otimes\operatorname{pr}_2^*\mathcal L_x^{-1}$ is trivial over the inverse images of an open cover
$(V_k)$ of $\operatorname{Spec}R$. Let $R'=R\otimes_KK'$ and $q:\operatorname{Spec}R'\to\operatorname{Spec}R$ the projection, which is
flat and surjective, hence an epimorphism of schemes. The point $y\circ q$ lifts to a point $y'$ of $T_{K'}$ over
$\operatorname{Spec}R'\to\operatorname{Spec}K'$, whose line bundle is the pullback of that of $y$, hence trivial over the inverse
images of the $q^{-1}(V_k)$. By the previous step for $T_{K'}$, $y'=0$, so $y\circ q=0$, and $y=0$ as $q$ is an epimorphism.

\emph{Descent to $K$: ampleness.} Let $\mathcal L_x$ and $\mathcal L_{K'}$ be the line bundles of $T_x$ and $T_{K'}$, so that
$\mathcal L_{K'}$ is the pullback of $\mathcal L_x$ along the projection $\pi:T_{K'}\to T_x$. As $T_{K'}$ and $T_x$ are proper
over $K'$ and $K$, Lemma~\ref{cv-origin-hodge-norm-form-descent-ample}(1) shows that $\mathcal L_{K'}$ is ample, in the sense
of that lemma, and that it suffices to show that $\mathcal L_x$ is ample. Let $y$ be a point of $T_x$. The morphism $\pi$ is
surjective, as a base change of $\operatorname{Spec}K'\to\operatorname{Spec}K$; let $z$ be a point of $T_{K'}$ with
$\pi(z)=y$. There are $n\ge1$ and a section $\sigma$ of $\mathcal L_{K'}^{\otimes n}$ that does not vanish at $z$. As $T_x$ is
quasi-compact and quasi-separated and $K\to K'$ is flat,
$\Gamma(T_{K'},\mathcal L_{K'}^{\otimes n})=\Gamma(T_x,\mathcal L_x^{\otimes n})\otimes_KK'$ \cite[Tag 02KH]{Sta}, so
$\sigma=\sum_jc_j\,\pi^*u_j$ with $c_j\in K'$ and sections $u_j$ of $\mathcal L_x^{\otimes n}$. In a trivialization of
$\mathcal L_x$ near $y$, the value of $\sigma$ at $z$ is a unit of the local ring of $T_{K'}$ at $z$, and it is
$\sum_jc_j\,\pi^\#(u_j(y))$, where $u_j(y)$ is the value of $u_j$ in the local ring of $T_x$ at $y$. A sum of nonunits of a
local ring is a nonunit, so some $c_j\,\pi^\#(u_j(y))$ is a unit; as $\pi^\#$ is a local homomorphism, $u_j(y)$ is a unit, that
is, $u_j$ does not vanish at $y$. The inverse image under $\pi$ of the set $T_{x,u_j}$ where $u_j$ does not vanish is the set
where $\pi^*u_j$ does not vanish, which is affine by Lemma~\ref{cv-origin-hodge-norm-form-descent-ample}(2) for $T_{K'}$. It
is the base change of $T_{x,u_j}\to\operatorname{Spec}K$ along the flat, surjective and quasi-compact morphism
$\operatorname{Spec}K'\to\operatorname{Spec}K$, so $T_{x,u_j}\to\operatorname{Spec}K$ is affine \cite[Tag 02L5]{Sta}, and
$T_{x,u_j}$ is affine. So $\mathcal L_x$ is ample, and by Lemma~\ref{cv-origin-hodge-norm-form-descent-ample}(1) it is
relatively ample by sections.
% Lean: GenericPol.proper_ample_kTriv: exists_lift (the residue field of a point of the pullback), sabPull with
% CubePf.cubePull, sabComp, sabCancel (the lift through the cartesian square of mj), MorphismProperty.of_isPullback for
% IsProper. KTriv (Out.KTriv, with translDiff) is the polarization condition of the text. Ampleness over K':
% relAmple_pull on an affine open of X' containing the image of x' (exists_isAffineOpen_mem_and_subset; affine psi by
% MorphismProperty.of_isPullback, sec_pull_down, the nonzero loci, isIso_pushoutSection_of_isAffineOpen for the generators);
% about 560 lines. The polarization over K': kTriv_pull (pushPt and its injectivity, exists_up and gen_chartPull). The
% field descent is AmpleDescent.proper_relAmple_kTriv_of_fieldExtension: properness by
% MorphismProperty.of_isPullback_of_descendsAlong (universally closed, along Spec K' -> Spec K) with smooth, separated and
% quasi-compact; the polarization by PolDescent.kTriv_of_pull (R' = Gamma of the fibre product, Flat.epi_of_flat_of_surjective).
% Ampleness over K: relAmple_iff_isAmpleLB_of_isAffine (Lemma cv-origin-hodge-norm-form-descent-ample(1)),
% isAmpleLB_of_relAmple_of_isAffine for T_K' with the LBIso to the pullback, and AmpleDescent.isAmpleLB_of_pull: pi
% surjective by MorphismProperty.of_isPullback, the lift z, h.exists_affine, exists_sec_nonzeroLocus_eq_of_lbIso,
% FlatSec.exists_comb_of_flat (the flat base change of sections over the quasi-compact quasi-separated T_x, through
% exists_finite_affine_refinement, exists_equiv_of_isPushout and exists_comb_ker; Stacks, Tag 02KH, checked), exists_chart,
% IsLocalRing.exists_of_isUnit_sum, mem_nonzeroLocus_iff_pull, IsAmpleLB.isAffineOpen_nonzeroLocus for T_K' (Lemma
% cv-origin-hodge-norm-form-descent-ample(2)) and isAffine_of_isPullback_of_flat_surjective (affineness descends along
% fpqc morphisms; Stacks, Tag 02L5, checked). Size: about 450 lines besides Lemma cv-origin-hodge-norm-form-descent-ample.
\end{proof}

\begin{lemma}\label{cv-origin-hodge-norm-form-fibre}
In the setting of Definition~\ref{cv-origin-hodge-component}, let $Z$ be a Hodge component of $\mathcal X$ at $S$, and let $z$
be a complex point of $Z$ over which $\mathcal A_Z$ is abelian, that is, whose centre lies in the abelian locus of
$\mathcal A_Z$ (Definition~\ref{cv-model-toroidal-semiabelian}(2)). Let $A$ be the fibre of $\mathcal A_Z$ at $z$, the base
change of $\mathcal A_Z$ along $z$ with its group law, its action of $\mathcal O_E$ and its symmetric cubical line bundle
$L_A$, and let $\varphi:A\to\mathcal A_Z$ be the projection, a morphism of $\mathcal O_E$-semi-abelian schemes over $z$. Then
$A$ is proper over $\mathbf C$, and $L_A$ is relatively ample by sections (Definition~\ref{cv-origin-hodge-norm-form-ample}).
\end{lemma}
\leanref{Development/AutomorphicSpectrum/Proofs/GeometryProofs/HodgeCptInteriorFlag.lean}{GlobalTheory.CompanionTheory.Ar.GrNorm.interior_pol}
\leanref{ReciprocityTheory/ReductiveGroups/ArithmeticGeometry/BirationalGeometry/ToroidalCompactification/OneArc.lean}{ReciprocityTheory.ChartRestriction.G.Geo.Tor.abLocus_pullback}
\leanref{Development/AutomorphicSpectrum/Proofs/GeometryProofs/HodgeCptInteriorFlag.lean}{GlobalTheory.CompanionTheory.Ar.GrNorm.injective_of_algebraic}
\begin{proof}
Let $f:Z\to\mathcal X$ be the morphism of the Hodge component, and $\mathcal A$ the semi-abelian scheme of the model, so that
$\mathcal A_Z\to\mathcal A$ is a morphism of $\mathcal O_E$-semi-abelian schemes over $f$
(Definition~\ref{cv-origin-hodge-component}(1)). Composed with $\varphi$, it gives a morphism $A\to\mathcal A$ over $f\circ z$.
Let $P$ be the centre of $z$.

\emph{The point $f(P)$ lies in the abelian locus of $\mathcal A$.} A semi-abelian scheme is abelian over a point $Q$ of its
base when its fibres at the points with values in algebraically closed fields at $Q$ are proper. Let $\bar\kappa$ be an
algebraic closure of the residue field of $P$, and $t:\operatorname{Spec}\bar\kappa\to Z$ the corresponding point. As $P$ lies
in the abelian locus of $\mathcal A_Z$, the fibre of $\mathcal A_Z$ at $t$ is proper; as $\mathcal A_Z$ is the pullback of
$\mathcal A$, it is the fibre of $\mathcal A$ at $f\circ t$. Let
$t':\operatorname{Spec}K\to\mathcal X$ be a point at $f(P)$ with $K$ algebraically closed. The fibre product of
$\operatorname{Spec}\bar\kappa$ and $\operatorname{Spec}K$ over $\mathcal X$ has a point; let $K''$ be its residue field. The
fibres of $\mathcal A$ at $f\circ t$ and at $t'$ become isomorphic over $K''$, so the base change to $K''$ of the fibre at
$t'$ is universally closed. As $\operatorname{Spec}K''\to\operatorname{Spec}K$ is an fpqc covering, the fibre at $t'$ is
universally closed \cite[Tag 02KS]{Sta}; it is separated and smooth, as $\mathcal A$ is semi-abelian, so it is proper.

\emph{The point $f(P)$ lies over the generic point of $\operatorname{Spec}\mathcal O_c$.} Let $\psi:\mathcal O_c\to\mathbf C$ be
the homomorphism given by $z$ and the structure morphism $\mathcal X\to\operatorname{Spec}\mathcal O_c$. It is injective. Let
$o\in\mathcal O_c$ be nonzero. As $F_c$ is algebraic over $\mathbf Q$, $o$ is a root of a nonzero polynomial with integer
coefficients; dividing it by the highest power of the variable that divides it gives a polynomial $Q$ with $Q(o)=0$, as
$o\neq0$ in the field $F_c$, and with nonzero constant term $Q(0)$. If $\psi(o)=0$, then $0=\psi(Q(o))=Q(\psi(o))=Q(0)$, a
contradiction. So the image of $f(P)$ in $\operatorname{Spec}\mathcal O_c$ is the prime $\ker\psi=0$, which is also the image
of the point of $\operatorname{Spec}F_c$. Hence $f(P)$ is the image of a point $w$ of the generic fibre
$\mathcal X_{F_c}=\mathcal X\times_{\mathcal O_c}F_c$. The image of $w$ in $\mathcal X$ is $f(P)$, which lies in the abelian
locus of $\mathcal A$; so $w$ lies in $j(X)$, by Definition~\ref{cv-model} for the model over $\mathcal O_c$. Thus $f(P)$ lies
in the image of the composite $X\to\mathcal X_{F_c}\to\mathcal X$.

\emph{Conclusion.} The universal abelian scheme of the moduli scheme $X$ is proper over $X$, and its line bundle is
relatively ample by sections and defines a principal polarization; by Definition~\ref{cv-model} it has a morphism of
$\mathcal O_E$-semi-abelian schemes to $\mathcal A$ over $X\to\mathcal X$. Lemma~\ref{cv-origin-hodge-norm-form-descent}, applied
to this morphism, to $K=\mathbf C$, the point $x=f\circ z$, whose image $f(P)$ lies in the image of $X$, and to
$T_x=A$ with its morphism to $\mathcal A$, shows that $A$ is proper over $\mathbf C$ and $L_A$ is relatively ample by
sections.
\end{proof}
% Lean: GrNorm.interior_pol: sabPull with Out.CubePf.cubePull for the fibre t' and m1, sabComp (pullCubComp) for n : SABMap t'
% Y.t (z >> f); abLocus_pullback with abLocusOne (AbLocusOneSt: the common field extension from the pullback of the two
% points and MorphismProperty.of_pullback_snd_of_descendsAlong for UniversallyClosed); injective_of_algebraic (Q from
% Polynomial.exists_eq_pow_rootMultiplicity_mul_and_not_dvd at 0); spec_injective_point;
% Scheme.Pullback.exists_preimage_pullback; Y.ab_gen (the preimage of the abelian locus in the generic fibre is the range of
% jQ); GenericPol.proper_ample_kTriv (Lemma cv-origin-hodge-norm-form-descent) with Y.m and the fields proper, ample and
% ktriv of Y.oF. The Lean states the existence of a fibre t' with these properties; the fibre it constructs is A. Size:
% about 200 lines besides Lemma cv-origin-hodge-norm-form-descent.

\begin{lemma}\label{cv-origin-hodge-norm-form-torus-exp-log}
Let $V$ be a finite-dimensional complex normed vector space, and let $m$ be a holomorphic map from a neighbourhood of $(0,0)$
in $V\times V$ to $V$ such that $m(x,0)=x$, $m(x,y)=m(y,x)$ and $m(m(x,y),z)=m(x,m(y,z))$ for $x$, $y$ and $z$ near $0$. Then
there is a holomorphic map $\lambda$ from a neighbourhood of $0$ in $V$ to $V$ with $\lambda(0)=0$, derivative the identity
at $0$, and $\lambda(m(x,y))=\lambda(x)+\lambda(y)$ for $x$ and $y$ near $0$.
\end{lemma}
\leanref{Development/AbelianVarietyTheory/Proofs/HomogeneousIdeals/IsLocalHomeomorphAddMonoidHomOfEventuallyEq.lean}{GlobalTheory.IsLocalCommGroupLaw.of_unit}
\leanref{Development/AbelianVarietyTheory/Proofs/HomogeneousIdeals/IsLocalHomeomorphAddMonoidHomOfEventuallyEq.lean}{GlobalTheory.IsLocalCommGroupLaw.log}
\leanref{Development/AbelianVarietyTheory/Proofs/HomogeneousIdeals/IsLocalHomeomorphAddMonoidHomOfEventuallyEq.lean}{GlobalTheory.IsLocalCommGroupLaw.log_law}
\leanref{Development/AbelianVarietyTheory/Proofs/HomogeneousIdeals/IsLocalHomeomorphAddMonoidHomOfEventuallyEq.lean}{GlobalTheory.IsLocalCommGroupLaw.hasFDerivAt_log}
\leanref{Development/AbelianVarietyTheory/Proofs/HomogeneousIdeals/IsLocalHomeomorphAddMonoidHomOfEventuallyEq.lean}{GlobalTheory.IsLocalCommGroupLaw.analyticAt_log}
\leanref{Development/AbelianVarietyTheory/Proofs/HomogeneousIdeals/TendstoUniformlyOnDoublingLogOfContraction.lean}{GlobalTheory.tendsto_doublingLog_of_contraction}
\begin{proof}
Use the maximum norm $|(x,y)|=\max(|x|,|y|)$ on $V\times V$.

\emph{The linear part.} The unit law at $x=0$ gives $m(0,0)=0$. By $m(x,0)=x$ and $m(0,y)=m(y,0)=y$, the derivative of $m$
at $(0,0)$ is $(x,y)\mapsto x+y$. As $m$ is holomorphic, Taylor's formula gives $D\ge0$ with $|m(x,y)-x-y|\le D|(x,y)|^2$ for
$(x,y)$ near $(0,0)$.

\emph{Halving.} The doubling map $d(x)=m(x,x)$ is holomorphic near $0$, with $d(0)=0$ and derivative $2$ at $0$. By the
holomorphic inverse function theorem it has a holomorphic local inverse $\varepsilon$ near $0$: $\varepsilon(0)=0$, and
$\varepsilon(d(x))=x$ and $d(\varepsilon(x))=x$ for $x$ near $0$. Its derivative at $0$ is $\frac12$, so Taylor's formula
gives $C\ge0$ with $|\varepsilon(x)-x/2|\le C|x|^2$ for $x$ near $0$. For $x$ and $y$ near $0$, associativity and
commutativity give
\[
d(m(x,y))=m(x,m(y,m(x,y)))=m(x,m(y,m(y,x)))=m(x,m(d(y),x))=m(x,m(x,d(y)))=m(d(x),d(y)).
\]
At $\varepsilon(x)$ and $\varepsilon(y)$ this reads $d(m(\varepsilon(x),\varepsilon(y)))=m(x,y)$, and applying $\varepsilon$
gives $\varepsilon(m(x,y))=m(\varepsilon(x),\varepsilon(y))$ for $x$ and $y$ near $0$.

\emph{The logarithm.} Choose $r>0$ such that $\varepsilon$ is holomorphic on the ball $B(0,r)$, with
$|\varepsilon(x)-x/2|\le C|x|^2\le|x|/8$ there. Then $|\varepsilon(x)|\le\frac58|x|$ on $B(0,r)$, so $\varepsilon$ maps
$B(0,r)$ into itself and $|\varepsilon^j(x)|\le(\frac58)^j|x|$ for $j\ge0$, where $\varepsilon^j$ is the $j$-th iterate. For
$x\in B(0,r)$,
\[
\big|2^{j+1}\varepsilon^{j+1}(x)-2^j\varepsilon^j(x)\big|=2^{j+1}\big|\varepsilon(\varepsilon^j(x))-\tfrac12\varepsilon^j(x)\big|
\le2^{j+1}C\big(\tfrac58\big)^{2j}|x|^2=2C|x|^2\big(\tfrac{25}{32}\big)^j.
\]
So the holomorphic functions $2^j\varepsilon^j$ converge uniformly on $B(0,r)$ to a function $\lambda$, holomorphic as a
uniform limit of holomorphic functions, with $|\lambda(x)-2^j\varepsilon^j(x)|\le\frac{64}7C|x|^2(\frac{25}{32})^j$. For $j=0$
this gives $|\lambda(x)-x|\le\frac{64}7C|x|^2$, so $\lambda(0)=0$ and the derivative of $\lambda$ at $0$ is the identity.

\emph{The law.} Choose $s\in(0,r]$ such that for $|(x,y)|<s$ the bound $|m(x,y)-x-y|\le D|(x,y)|^2$ and the relation
$\varepsilon(m(x,y))=m(\varepsilon(x),\varepsilon(y))$ hold and $|m(x,y)|<r$. Let $|(x,y)|<s$. As $x,y\in B(0,r)$,
$|(\varepsilon^j(x),\varepsilon^j(y))|\le(\frac58)^j|(x,y)|<s$ for every $j$, so by induction
$\varepsilon^j(m(x,y))=m(\varepsilon^j(x),\varepsilon^j(y))$, and
\[
\big|2^j\varepsilon^j(m(x,y))-2^j\varepsilon^j(x)-2^j\varepsilon^j(y)\big|
=2^j\big|m(\varepsilon^j(x),\varepsilon^j(y))-\varepsilon^j(x)-\varepsilon^j(y)\big|\le D|(x,y)|^2\big(\tfrac{25}{32}\big)^j,
\]
which tends to $0$. As $x$, $y$ and $m(x,y)$ lie in $B(0,r)$, the three sequences converge to $\lambda(m(x,y))$, $\lambda(x)$
and $\lambda(y)$, so $\lambda(m(x,y))=\lambda(x)+\lambda(y)$.
\end{proof}
% Lean: IsLocalCommGroupLaw (analyticAt, isBigO, comm, assoc) from IsLocalCommGroupLaw.of_unit (the derivative from hu and
% hu', isBigO_sub_linear_sq_of_analyticAt); map_zero, hasFDerivAt, doubling_law (the chain of the text: assoc, comm, assoc,
% comm, assoc), hasStrictFDerivAt_doubling with twoEquiv, halving = localInverse, eventuallyEq_halving, doubling_inverse,
% analyticAt_inverse, isBigO_inverse and inverse_law; log = doublingLog halving, the limit of 2^j eps^[j]
% (exists_ball_contraction with the factor 5/8, norm_iterate_le_of_contraction, dist_two_pow_smul_iterate_succ_le,
% tendsto_doublingLog_of_contraction, dist_doublingLog_le_of_contraction with 64/7, isBigO_doublingLog_sub), log_zero and
% hasFDerivAt_log by hasFDerivAt_of_isBigO_sq, analyticAt_log (analyticAt_doublingLog through
% tendstoUniformlyOn_doublingLog_of_contraction and SCV.differentiableOn_of_tendstoLocallyUniformlyOn), and log_law
% (doublingLog_law: hit, hq, hqs and the bound D |p|^2 (25/32)^j). Size: about 450 lines, besides the archive's proof that a
% locally uniform limit of holomorphic functions of several variables is holomorphic.

\begin{lemma}\label{cv-origin-hodge-norm-form-torus-exp}
Let $\Gamma$ be a commutative topological group, written additively, and $V$ a finite-dimensional complex vector space. Let
$\phi$ be a homeomorphism from an open neighbourhood of $0$ in $\Gamma$ onto an open subset of $V$ with $\phi(0)=0$, such that
the law $m(x,y)=\phi(\phi^{-1}(x)+\phi^{-1}(y))$ is holomorphic near $(0,0)$. Then there is a homomorphism $e:V\to\Gamma$ such
that:
\begin{enumerate}
\item $e$ is a local homeomorphism, $v\mapsto\phi(e(v))$ is holomorphic near $0$ with derivative the identity at $0$, and the
  kernel of $e$ is discrete;
\item if $\Gamma$ is compact, the kernel of $e$ spans $V$ over $\mathbf R$;
\item if $\Gamma$ is compact and $n\Gamma=\Gamma$ for every integer $n\ge1$, then $e$ is surjective.
\end{enumerate}
\end{lemma}
\leanref{Development/AbelianVarietyTheory/Proofs/HomogeneousIdeals/IsLocalHomeomorphAddMonoidHomOfEventuallyEq.lean}{GlobalTheory.IsAnalyticGroupChart.exp}
\leanref{Development/AbelianVarietyTheory/Proofs/HomogeneousIdeals/IsLocalHomeomorphAddMonoidHomOfEventuallyEq.lean}{GlobalTheory.IsAnalyticGroupChart.exp_surjective}
\leanref{Development/AbelianVarietyTheory/Proofs/HomogeneousIdeals/IsLocalHomeomorphAddMonoidHomOfEventuallyEq.lean}{GlobalTheory.IsAnalyticGroupChart.isZLattice_periodLattice}
\leanref{Development/AbelianVarietyTheory/Proofs/HomogeneousIdeals/IsLocalHomeomorphAddMonoidHomOfEventuallyEq.lean}{GlobalTheory.IsAnalyticGroupChart.hasFDerivAt_chart_exp}
\begin{proof}
\emph{The logarithm.} For $x$ near $0$ in $V$, $\phi^{-1}(x)$ is near $0$ in $\Gamma$, so $m(x,0)=\phi(\phi^{-1}(x))=x$, and
$m(x,y)=m(y,x)$ as $\Gamma$ is commutative. For $x$ and $y$ near $0$, the sum $\phi^{-1}(x)+\phi^{-1}(y)$ is near $0$ in
$\Gamma$, by the continuity of the addition, so it lies in the source of $\phi$ and $\phi^{-1}(m(x,y))=\phi^{-1}(x)+\phi^{-1}(y)$.
Hence $m(m(x,y),z)=\phi(\phi^{-1}(x)+\phi^{-1}(y)+\phi^{-1}(z))=m(x,m(y,z))$ for $x$, $y$ and $z$ near $0$.
Lemma~\ref{cv-origin-hodge-norm-form-torus-exp-log} gives a holomorphic map $\lambda$ near $0$ in $V$ with $\lambda(0)=0$,
derivative the identity at $0$, and $\lambda(m(x,y))=\lambda(x)+\lambda(y)$ for $x$ and $y$ near $0$.

\emph{A local exponential.} By the inverse function theorem, $\lambda$ restricts to a biholomorphism between neighbourhoods of
$0$ in $V$; let $\psi=\phi^{-1}\circ\lambda^{-1}$, a homeomorphism from a neighbourhood of $0$ in $V$ onto one of $0$ in $\Gamma$.
For $v$ and $w$ near $0$, with $x=\lambda^{-1}(v)$ and $y=\lambda^{-1}(w)$, $\lambda(m(x,y))=v+w$, so
$\psi(v+w)=\phi^{-1}(m(x,y))=\psi(v)+\psi(w)$.

\emph{The homomorphism.} Let $\rho>0$ be such that $\psi(v+w)=\psi(v)+\psi(w)$ for $|v|,|w|<\rho$. For $v\in V$ choose $j$ with
$|2^{-j}v|<\rho$ and put $e(v)=2^j\psi(2^{-j}v)$, as in the proof of Lemma~\ref{cv-origin-hodge-norm-form-lift-real}: this does not
depend on $j$, $e$ is a homomorphism, and $e=\psi$ on $B(0,\rho)$.

(1) Near any $v$, $e(w)=e(v)+\psi(w-v)$, so $e$ is a local homeomorphism; in particular it is continuous and open. Near $0$,
$\phi\circ e=\lambda^{-1}$, which is holomorphic with derivative the identity at $0$. As $e$ is injective on a neighbourhood of
$0$, this neighbourhood meets the kernel only in $0$.

(2) The image $e(V)$ is an open subgroup of $\Gamma$, hence closed, hence compact. The open sets $e(B(0,n))$, for $n\ge1$,
increase and cover it, so $e(V)=e(B(0,n))$ for some $n$. Let $v\in V$, $\eta>0$ and $t=(n+1)/\eta$. Then $e(tv)=e(b)$ for some $b$
with $|b|<n$, so $tv-b$ lies in the kernel, and $|v-t^{-1}(tv-b)|=|b|/t<\eta$. So $v$ lies in the closure of the real span of the
kernel, which is closed, as $V$ is finite-dimensional.

(3) The open subgroup $e(V)$ of the compact group $\Gamma$ has finite index $N$. For $a\in\Gamma$ choose $b$ with $Nb=a$; as $N$
kills $\Gamma/e(V)$, $a=Nb$ lies in $e(V)$.
\end{proof}
% Lean: IsAnalyticGroupChart, with chartLaw and isLocalCommGroupLaw (IsLocalCommGroupLaw.of_unit with the unit law from
% eventually_right_inverse, comm from add_comm, and assoc through symm_chartLaw); the logarithm lambda is
% IsLocalCommGroupLaw.log (Lemma cv-origin-hodge-norm-form-torus-exp-log), with logChart from HasStrictFDerivAt; localExp,
% localExp_add (logChart_symm_add and symm_chartLaw), exp by exists_addMonoidHom_eventuallyEq, exp_eventuallyEq,
% isLocalHomeomorph_exp (isLocalHomeomorph_addMonoidHom_of_eventuallyEq), discreteTopology_ker_of_isLocalHomeomorph,
% hasFDerivAt_chart_exp, span_ker_eq_top_of_isOpenMap_of_compact for (2) and
% surjective_of_isOpenMap_of_compact_of_divisible for (3). Size: about 250 lines besides Lemma
% cv-origin-hodge-norm-form-torus-exp-log.

\begin{lemma}\label{cv-origin-hodge-norm-form-torus-points}
Let $X$ be a scheme over $\mathbf C$. Give the set $X(\mathbf C)$ of its points over $\mathbf C$ the topology generated by the
sets of points $x\in U$ with $f(x)\in O$, for $U\subset X$ open, $f$ regular on $U$ and $O\subset\mathbf C$ open.
\begin{enumerate}
\item Let $Z$ be a topological space and $c:Z\to X(\mathbf C)$ a map such that every $z\in Z$ has an affine open $P\subset X$
  containing $c(z)$ for which the set $N$ of $z'\in Z$ with $c(z')\in P$ is open and $z'\mapsto g(c(z'))$ is continuous on $N$
  for every $g\in\Gamma(P,\mathcal O_X)$. Then $c$ is continuous.
\item If $X$ is a group scheme over $\mathbf C$, then $X(\mathbf C)$ is a topological group.
\item If $X$ has a line bundle that is relatively ample by sections (Definition~\ref{cv-origin-hodge-norm-form-ample}), then
  $X(\mathbf C)$ is compact.
\end{enumerate}
\end{lemma}
\leanref{Development/AbelianVarietyTheory/Proofs/ToricGeometry/IsTopologicalGroupComplexPoints.lean}{GlobalTheory.ComplexPointsGroup.isContFamily_of_affine}
\leanref{Development/AbelianVarietyTheory/Proofs/ToricGeometry/IsTopologicalGroupComplexPoints.lean}{GlobalTheory.isTopologicalGroup_complexPoints}
\leanref{Development/AbelianVarietyTheory/Proofs/ToricGeometry/CompactSpacePointsPolAV.lean}{GlobalTheory.PointsCompact.compact_of_sections}
\leanref{Development/AutomorphicSpectrum/Proofs/GeometryProofs/ExistsExpChartCoordinates.lean}{GlobalTheory.CompanionTheory.Ar.AVExp.compactSpace_points}
\begin{proof}
(1) Let $U\subset X$ be open, $f$ regular on $U$, $O\subset\mathbf C$ open, and $z\in Z$ with $c(z)\in U$ and $f(c(z))\in O$;
let $P$ and $N$ be as in the hypothesis for $z$. There is $s\in\Gamma(P,\mathcal O_X)$ such that the set $D(s)$ where $s$ does
not vanish contains $c(z)$ and lies in $U\cap P$, and $f=h/s^n$ on $D(s)$ with $h\in\Gamma(P,\mathcal O_X)$. The set $N'$ of
$z'\in N$ with $s(c(z'))\neq0$ is open, as $s\circ c$ is continuous on $N$, and for $z'\in N'$, $c(z')\in U$ and
$f(c(z'))=h(c(z'))/s(c(z'))^n$, a continuous function of $z'$. So the set of $z'\in N'$ with $f(c(z'))\in O$ is an open
neighbourhood of $z$ that $c$ maps into the set of $x\in U$ with $f(x)\in O$. Hence $c$ is continuous.

(2) By the definition of the topology, a map $c:Z\to X(\mathbf C)$ is continuous if and only if, for every open $U\subset X$
and every regular $f$ on $U$, the set $c^{-1}(U)$ is open and $f\circ c$ is continuous on it. So for a morphism $q:X\to Y$ of
schemes over $\mathbf C$ and a continuous $c$, the map $q\circ c$ is continuous, as $(q\circ c)^{-1}(U)=c^{-1}(q^{-1}(U))$ and
$f\circ q$ is regular on $q^{-1}(U)$. Let $a:Z\to X(\mathbf C)$ and $b:Z\to Y(\mathbf C)$ be continuous; we show that
$(a,b):Z\to(X\times_{\mathbf C}Y)(\mathbf C)$ is continuous. Let $z\in Z$, and let $V_X\subset X$ and $V_Y\subset Y$ be affine
opens containing $a(z)$ and $b(z)$. Then $P=V_X\times_{\mathbf C}V_Y$ is an affine open of $X\times_{\mathbf C}Y$ with ring
$\Gamma(V_X,\mathcal O_X)\otimes_{\mathbf C}\Gamma(V_Y,\mathcal O_Y)$, and the set $N$ of $z'$ with $(a(z'),b(z'))\in P$ is
$a^{-1}(V_X)\cap b^{-1}(V_Y)$, which is open. The regular functions $g$ on $P$ for which $z'\mapsto g(a(z'),b(z'))$ is continuous
on $N$ form a subring, which contains the functions $f\circ p$ and $f'\circ p'$ for the projections $p$ and $p'$, as
$(f\circ p)(a,b)=f\circ a$ and $(f'\circ p')(a,b)=f'\circ b$. These generate the ring of $P$, so (1) applies. For a group
scheme $X$, the group law of $X(\mathbf C)$ is the composite of the continuous map $X(\mathbf C)\times X(\mathbf C)\to
(X\times_{\mathbf C}X)(\mathbf C)$ given by the two projections with the multiplication morphism, and the inverse is the
composite of the identity with the inverse morphism; both are continuous.

(3) Let $\mathcal L$ be the line bundle, and let $N\ge1$, the sections $s_0,\dots,s_M$ of $\mathcal L^{\otimes N}$ and the
functions $r_{ij}$ be given by Definition~\ref{cv-origin-hodge-norm-form-ample} for the affine open $\operatorname{Spec}\mathbf C$.
For each $i$, the ring of the affine scheme $X_{s_i}$ is the quotient of $\mathbf C[y_0,\dots,y_M]$ by the kernel $I_i$ of
$y_j\mapsto r_{ij}$. Let $K_i$ be the set of $z\in\mathbf C^{M+1}$ with $|z_j|\le1$ for every $j$ and $q(z)=0$ for every
$q\in I_i$, a compact set. For $z\in K_i$, evaluation at $z$ induces a homomorphism of $\mathbf C$-algebras from the ring of
$X_{s_i}$ to $\mathbf C$, so a point $\sigma_i(z)\in X_{s_i}(\mathbf C)$ with $r_{ij}(\sigma_i(z))=z_j$. Every regular function
on $X_{s_i}$ is a polynomial in the $r_{ij}$, so its composite with $\sigma_i$ is a polynomial function of $z$. By (1),
$\sigma_i$ is continuous, and $C_i=\sigma_i(K_i)$ is compact. Let $x\in X(\mathbf C)$. Choose an open set containing $x$ on
which $\mathcal L$, and so $\mathcal L^{\otimes N}$, is trivialized, and let $a_j$ be the value at $x$ of $s_j$ in this
trivialization. As the $X_{s_i}$ cover $X$, some $a_i$ is nonzero; choose $i_0$ with $|a_{i_0}|$ maximal. Then $a_{i_0}\neq0$,
so $x\in X_{s_{i_0}}$, and $a_j=r_{i_0j}(x)a_{i_0}$, as $s_j=r_{i_0j}s_{i_0}$ on $X_{s_{i_0}}$. So $|r_{i_0j}(x)|\le1$ for
every $j$, the point $z=(r_{i_0j}(x))_j$ lies in $K_{i_0}$, and $x=\sigma_{i_0}(z)\in C_{i_0}$, as a point of the affine scheme
$X_{s_{i_0}}$ is determined by its values at the generators $r_{i_0j}$. So $X(\mathbf C)$ is the union of the finitely many
compact sets $C_i$, hence compact.
\end{proof}
% Lean: (1) is ComplexPointsGroup.isContFamily_of_affine (exists_eq_pow_mul_of_isAffineOpen for f = h / s^n), with
% PointsCompact.continuous_of_affine, the case of one affine open, used in (3). (2) is isTopologicalGroup_complexPoints:
% continuous_iff_isContFamily, IsContFamily.comp, isContFamily_lift through isContFamily_of_isPullback
% (isIso_pushoutSection_of_isAffineOpen and closure_range_union_range_eq_top_of_isPushout for the generators, and the subring
% by closure_induction), with mu and the inverse of A. (3) is AVExp.compactSpace_points: the RelAmple data for the affine
% open top, restricted to Y over the open top of Sp C, and PointsCompact.compact_of_sections: exists_compact_of_affine (K the
% polydisc intersected with the zero set of the kernel of eval2Hom c r, sigma from AffinePoints, continuous_of_affine), and
% the cover by Finite.exists_max of the values a j. The points are the Lean's morphisms 1 -> A over Sp C. Size: about 450
% lines.

\begin{lemma}\label{cv-origin-hodge-norm-form-torus-chart-etale}
Let $R$ be a $\mathbf C$-algebra, $\phi:\mathbf C[x_1,\dots,x_g]\to R$ an \'etale homomorphism of $\mathbf C$-algebras, and
$e:R\to\mathbf C$ a homomorphism of $\mathbf C$-algebras with $e(\phi(x_i))=0$ for every $i$. Then there are an open set
$W\ni0$ in $\mathbf C^g$ and homomorphisms of $\mathbf C$-algebras $\sigma_w:R\to\mathbf C$, for $w\in W$, such that:
\begin{enumerate}
\item $\sigma_0=e$, $\sigma_w(\phi(x_i))=w_i$ for $w\in W$ and every $i$, and $w\mapsto\sigma_w(a)$ is holomorphic on $W$ for
  every $a\in R$;
\item there are $f_1,\dots,f_n\in R$ and $\eta>0$ such that every homomorphism of $\mathbf C$-algebras $\xi:R\to\mathbf C$
  with $|\xi(f_k)-e(f_k)|<\eta$ for every $k$ is $\sigma_w$ for some $w\in W$.
\end{enumerate}
\end{lemma}
\leanref{Development/AbelianVarietyTheory/Proofs/HomogeneousIdeals/ExistsOpenPartialHomeomorphAnalytic.lean}{GlobalTheory.AffinePoints.exists_chart_of_etale}
\leanref{Development/AbelianVarietyTheory/Proofs/ToricGeometry/ExistsTestFunctionsSubsetOfMemNhds.lean}{GlobalTheory.exists_holomorphic_section_of_etale}
\leanref{Development/AbelianVarietyTheory/Proofs/ToricGeometry/ExistsTestFunctionsSubsetOfMemNhds.lean}{GlobalTheory.exists_testFunctions_subset_of_mem_nhds}
\begin{proof}
Write $P=\mathbf C[x_1,\dots,x_g]$. As $R$ is \'etale over $P$, it has a presentation $R=P[y_1,\dots,y_m]/(F_1,\dots,F_m)$
such that the image $\delta$ of $\det(\partial F_r/\partial y_l)$ in $R$ is invertible \cite[Tag 00U9]{Sta}. Let $\bar y_l$ be
the image of $y_l$ in $R$, and for $(w,v)\in\mathbf C^g\times\mathbf C^m$ and $q\in P[y_1,\dots,y_m]$ let $q(w,v)$ be the
value of $q$ at $x=w$ and $y=v$.

\emph{Points as solutions.} For a homomorphism of $\mathbf C$-algebras $\xi:R\to\mathbf C$ put
$j(\xi)=((\xi(\phi(x_i)))_i,(\xi(\bar y_l))_l)$. Then $\xi$ maps the image of $q$ to $q(j(\xi))$, and $j$ is injective, as $R$
is generated by the $\phi(x_i)$ and the $\bar y_l$. Conversely, if $F_r(w,v)=0$ for every $r$, then $(w,v)=j(\xi)$ for the
homomorphism $\xi$ induced by evaluation at $(w,v)$, as the $F_r$ generate the kernel of $P[y_1,\dots,y_m]\to R$.

\emph{The inverse function theorem.} Let $\Psi(w,v)=(w,(F_r(w,v))_r)$, a polynomial map from $\mathbf C^g\times\mathbf C^m$ to
itself, and $p_0=j(e)$. If the derivative of $\Psi$ at $p_0$ kills $(u,v)$, then $u=0$ and
$\sum_l\frac{\partial F_r}{\partial y_l}(p_0)v_l=0$ for every $r$; the determinant of this system is $e(\delta)\neq0$, as
$\delta$ is a unit, so $v=0$. By the holomorphic inverse function theorem, $\Psi$ restricts to a biholomorphism $\Phi$ from an
open neighbourhood $S$ of $p_0$ onto an open set $T$. As $\Psi(j(\xi))=((\xi(\phi(x_i)))_i,0)$ for every $\xi$,
$\Phi(p_0)=(0,0)$.

\emph{The sections.} Let $W$ be the set of $w\in\mathbf C^g$ with $(w,0)\in T$, an open set containing $0$. For $w\in W$,
$\Phi^{-1}(w,0)=(w',v)$ satisfies $\Psi(w',v)=(w,0)$, so $w'=w$ and $F_r(w,v)=0$ for every $r$; let $\sigma_w$ be the
homomorphism with $j(\sigma_w)=\Phi^{-1}(w,0)$. Then $\sigma_w(\phi(x_i))=w_i$, and for $a\in R$, the image of a polynomial
$q$, $\sigma_w(a)=q(\Phi^{-1}(w,0))$ is holomorphic in $w$. As $\Phi(p_0)=(0,0)$, $j(\sigma_0)=p_0=j(e)$, so $\sigma_0=e$. This
is (1).

(2) Let the $f_k$ be the $\phi(x_i)$ and the $\bar y_l$, and choose $\eta>0$ such that $S$ contains every point whose
coordinates differ from those of $p_0$ by less than $\eta$. If $|\xi(f_k)-e(f_k)|<\eta$ for every $k$, then
$j(\xi)\in S$, so $\Phi(j(\xi))=(w,0)$ with $w=(\xi(\phi(x_i)))_i\in W$, and $j(\xi)=\Phi^{-1}(w,0)=j(\sigma_w)$; so
$\xi=\sigma_w$.
\end{proof}
% Lean: exists_holomorphic_section_of_etale: AffinePoints.exists_chart_of_etale for the etale g = phi and the point e, an
% OpenPartialHomeomorph c from the C-points of A (topology of pointwise convergence) to C^d with c = coords g and c.symm
% analytic: the presentation P from Algebra.Etale.iff_isStandardSmoothOfRelativeDimension_zero (hcard and hbij: as many
% relations as variables), j, ev, hevj, hker, hex, hjinj, Psi, hPsi_an (analyticAt_aevalTower), hinj (hasDerivAt_aevalTower_line,
% the jacobian P.jacobian_isUnit mapped by p, eq_zero_of_vecMul_eq_zero), exists_openPartialHomeomorph_analytic (the
% holomorphic inverse function theorem), hex' and han; then exists_testFunctions_subset_of_mem_nhds, which takes generators
% of a finite presentation of A over C and the closed embedding isClosedEmbedding_coords (the text takes the generators of the
% etale presentation). The Lean's sigma is c.symm on c.target, the text's W. Size: about 500 lines.

\begin{lemma}\label{cv-origin-hodge-norm-form-torus-chart}
Let $X$ be a scheme over $\mathbf C$, smooth of relative dimension $g$, with the topology on $X(\mathbf C)$ of
Lemma~\ref{cv-origin-hodge-norm-form-torus-points}, and let $o\in X(\mathbf C)$. Then there are an open set $W\ni0$ in
$\mathbf C^g$, a map $\sigma:W\to X(\mathbf C)$ with $\sigma(0)=o$, an affine open $U\subset X$ and regular functions
$t_1,\dots,t_g$ on $U$ such that:
\begin{enumerate}
\item for every open $V\subset X$ and every regular function $f$ on $V$, the set of $w\in W$ with $\sigma(w)\in V$ is open and
  $w\mapsto f(\sigma(w))$ is holomorphic on it;
\item $\sigma(w)\in U$ and $t_i(\sigma(w))=w_i$ for $w\in W$ and every $i$, and $\sigma(W)$ is open in $X(\mathbf C)$;
\item every additive map $D:\Gamma(U,\mathcal O_X)\to\mathbf C$ with $D(c)=0$ for the constants $c\in\mathbf C$,
  $D(ff')=f(o)D(f')+f'(o)D(f)$ for all $f,f'$, and $D(t_i)=0$ for all $i$, is zero.
\end{enumerate}
\end{lemma}
\leanref{Development/AbelianVarietyTheory/Proofs/ToricGeometry/ExistsEtaleCoordsVanishing.lean}{GlobalTheory.HolChart.exists_holomorphic_chart}
\leanref{Development/AbelianVarietyTheory/Proofs/ToricGeometry/ExistsEtaleCoordsVanishing.lean}{GlobalTheory.HolChart.exists_etale_coords_vanishing}
\leanref{Development/AbelianVarietyTheory/Proofs/ToricGeometry/ExistsEtaleCoordsVanishing.lean}{GlobalTheory.HolChart.derivation_eq_zero}
\leanref{Development/AbelianVarietyTheory/Proofs/ToricGeometry/ExistsEtaleCoordsVanishing.lean}{GlobalTheory.HolChart.hol_of_affine}
\begin{proof}
\emph{\'Etale coordinates.} As $X$ is smooth of relative dimension $g$ over $\mathbf C$, there are an affine open $U$
containing $o$ and an \'etale homomorphism of $\mathbf C$-algebras $\phi_0:P=\mathbf C[x_1,\dots,x_g]\to R=\Gamma(U,\mathcal O_X)$
\cite[Tags 00TA and 054L]{Sta}. Let $e:R\to\mathbf C$ be the homomorphism of the point $o$, and let $\phi$ be the composite of
$\phi_0$ with the automorphism of $P$ that sends $x_i$ to $x_i-e(\phi_0(x_i))$. It is \'etale, and $e(\phi(x_i))=0$. Put
$t_i=\phi(x_i)$.

(3) Let $D$ be as in (3). By the product rule and induction on polynomials, $D$ vanishes on $\phi(P)$. So $F(f)=e(f)+\varepsilon
D(f)$ and $F_0(f)=e(f)$ are homomorphisms of $P$-algebras from $R$ to $\mathbf C[\varepsilon]/(\varepsilon^2)$, with $P$ acting
through $e\circ\phi$, and they agree modulo the square-zero ideal $(\varepsilon)$. As $R$ is \'etale over $P$,
$\Omega_{R/P}=0$ \cite[Tag 00U1]{Sta}, so $R$ is formally unramified over $P$ \cite[Tag 00UO]{Sta}, and $F=F_0$. So $D=0$.

\emph{The section.} Lemma~\ref{cv-origin-hodge-norm-form-torus-chart-etale}, applied to $R$, $\phi$ and $e$, gives an open set
$W_0\ni0$, homomorphisms $\sigma_w$ for $w\in W_0$, elements $f_1,\dots,f_n$ of $R$ and $\eta>0$. As $U$ is affine,
each $\sigma_w$ is the homomorphism of a unique point $\sigma(w)\in U(\mathbf C)$, so that $f(\sigma(w))=\sigma_w(f)$ for
$f\in R$, and $\sigma(0)=o$. Let $W$ be the set of $w\in W_0$ with $|\sigma_w(f_k)-e(f_k)|<\eta$ for every $k$, an open
set containing $0$. Then $t_i(\sigma(w))=w_i$ for $w\in W$, and $w\mapsto f(\sigma(w))$ is holomorphic on $W$ for every
$f\in R$. By the step ``Affine charts suffice'' of the proof of Lemma~\ref{cv-origin-hodge-norm-form-appell-pair}, applied to
$\sigma$ on $W$ with the affine open $U$, (1) holds. The image $\sigma(W)$ is the set of points $x\in U$ with
$|f_k(x)-f_k(o)|<\eta$ for every $k$: such a point is $\sigma(w)$ for some $w\in W_0$, by
Lemma~\ref{cv-origin-hodge-norm-form-torus-chart-etale}(2), and then $w\in W$. It is a finite intersection of the open sets
that generate the topology of $X(\mathbf C)$. This is (2).
\end{proof}
% Lean: HolChart.exists_holomorphic_chart: SmoothOfRelativeDimension.exists_isStandardSmoothOfRelativeDimension at the point
% (U0 = top, as Sp C has one point), RingHom.IsStandardSmoothOfRelativeDimension for c, exists_etale_coords_vanishing
% (IsStandardSmoothOfRelativeDimension.exists_etale_mvPolynomial and the translation T), derivation_eq_zero (hP by
% MvPolynomial.induction_on, F and F0 into DualNumber C, Algebra.FormallyUnramified.ext'), existsUnique_point_of_ringHom
% for the points pt of U, exists_holomorphic_section_of_etale (Lemma cv-origin-hodge-norm-form-torus-chart-etale), W', hmem,
% hol_of_affine (the affine criterion of Lemma cv-origin-hodge-norm-form-appell-pair) and himg for the open image. The
% published counterparts of the first two steps are the local structure of smooth morphisms (Stacks, Tags 00TA and 054L,
% checked); the Lean takes them from Mathlib. Size: about 300 lines besides Lemma cv-origin-hodge-norm-form-torus-chart-etale.

\begin{lemma}\label{cv-origin-hodge-norm-form-torus-divisible}
Let $A$ be an abelian scheme over $\operatorname{Spec}\mathbf C$ and $r\ge1$ an integer. Then every $a\in A(\mathbf C)$ is $rb$
for some $b\in A(\mathbf C)$.
\end{lemma}
\leanref{Development/AutomorphicSpectrum/Proofs/GeometryProofs/ExistsExpChartCoordinates.lean}{GlobalTheory.CompanionTheory.Ar.AVExp.exists_div}
\leanref{Development/AutomorphicSpectrum/Proofs/GeometryProofs/ExistsExpChartCoordinates.lean}{GlobalTheory.CompanionTheory.Ar.AVExp.surjective_mulA}
\leanref{ReciprocityTheory/ReductiveGroups/ArithmeticGeometry/FirstCohomology/SectionDivision.lean}{ReciprocityTheory.ChartRestriction.G.H1.SectionDivision.exists_section}
\leanref{Development/AbelianVarietyTheory/Proofs/DivisorProofs/IsFiniteAndFlatTorsion.lean}{GlobalTheory.AbTorsionAnyBase.isFinite_and_flat_mulA}
\begin{proof}
Multiplication by $r$ on $A$ is a morphism $[r]:A\to A$ that is finite and flat, as it is an isogeny of the abelian variety
$A$ \cite{Mum70}. It is locally of finite presentation, as $A$ is of finite type over $\mathbf C$, so it is an open map
\cite[Tag 01UA]{Sta}; it is finite, hence integral and universally closed \cite[Tag 01WM]{Sta}, so it is a closed map. As $A$
is connected, the image of $[r]$, which is open, closed and nonempty, is all of $A$.

Let $a\in A(\mathbf C)$, and let $F=\operatorname{Spec}\mathbf C\times_{a,A,[r]}A$ be the fibre of $[r]$ over $a$. It is finite
over $\operatorname{Spec}\mathbf C$, as a base change of $[r]$, and nonempty, as the image of $a$ lies in the image of $[r]$
\cite[Tag 01JT]{Sta}. So $F=\operatorname{Spec}B$ for a nonzero finite-dimensional $\mathbf C$-algebra $B$. Let $\mathfrak m$ be
a maximal ideal of $B$. The field $B/\mathfrak m$ is finite over $\mathbf C$, hence equal to $\mathbf C$, as $\mathbf C$ is
algebraically closed. The homomorphism $B\to B/\mathfrak m=\mathbf C$ gives a section of $F\to\operatorname{Spec}\mathbf C$;
its image $b$ in $A(\mathbf C)$ satisfies $rb=[r]\circ b=a$.
\end{proof}
% Lean: AVExp.exists_div: GlobalTheory.AbTorsionAnyBase.isFinite_and_flat_mulA (for an abelian scheme over any base),
% surjective_mulA (LocallyOfFinitePresentation through TangentCrit.locallyOfFinitePresentation_of_comp, isClosedMap and
% isOpenMap, ConnectedSpace from GeometricallyConnected, IsClopen.eq_univ), and H1.SectionDivision.exists_section for the
% pullback of mulA along a0 (a maximal ideal of the finite C-algebra, IsAlgClosed.algebraMap_bijective_of_isIntegral). The
% finiteness and flatness of multiplication by r is published (Mumford, Abelian Varieties: it is an isogeny); the Lean proves
% it for abelian schemes over any base, with about 16,000 lines by the archive's dependency count. Size: about 80 lines
% besides isFinite_and_flat_mulA.

\begin{lemma}\label{cv-origin-hodge-norm-form-torus}
Let $A$ be an abelian scheme over $\operatorname{Spec}\mathbf C$, smooth of relative dimension $g$, with a line bundle that is
relatively ample by sections (Definition~\ref{cv-origin-hodge-norm-form-ample}). Then there is a uniformization $u$ of $A$ for
the trivial line bundle, with $\theta=1$ (Definition~\ref{cv-origin-hodge-norm-gram}(1), with $T=\operatorname{Spec}\mathbf C$,
$G=A$ and $\varphi$ the identity): a surjective homomorphism $u:\mathbf C^g\to A(\mathbf C)$ whose kernel $\Lambda$ is
discrete and spans $\mathbf C^g$ over $\mathbf R$, such that for every open $U\subset A$ and every regular function $f$ on $U$
the set $u^{-1}(U)$ is open and $f\circ u$ is holomorphic on it. Moreover there are an affine open $U_0\subset A$ containing
the origin $u(0)$ and regular functions $t_1,\dots,t_g$ on $U_0$ such that:
\begin{enumerate}
\item[(a)] every additive map $D:\Gamma(U_0,\mathcal O_A)\to\mathbf C$ with $D(c)=0$ for the constants $c\in\mathbf C$,
  $D(ff')=f(u(0))D(f')+f'(u(0))D(f)$ for all $f,f'$, and $D(t_i)=0$ for all $i$, is zero;
\item[(b)] the map $x\mapsto(t_i(u(x)))_i$, defined near $0$, has derivative the identity at $0$.
\end{enumerate}
\end{lemma}
\leanref{Development/AutomorphicSpectrum/Proofs/GeometryProofs/ExistsExpChartCoordinatesStatement.lean}{GlobalTheory.T.exists_exp_chart_coordinates}
\leanref{Development/AutomorphicSpectrum/Proofs/GeometryProofs/ExistsExpChartCoordinates.lean}{GlobalTheory.CompanionTheory.Ar.AVExp.exists_holUnif}
\begin{proof}
Give $A(\mathbf C)$ the topology of Lemma~\ref{cv-origin-hodge-norm-form-torus-points}. By (2) and (3) of that lemma, it is a
compact topological group, commutative as $A$ is an abelian scheme. Call a map from an open subset of $\mathbf C^n$ to
$A(\mathbf C)$ holomorphic if it is so in the sense of Lemma~\ref{cv-origin-hodge-norm-form-appell-pair}. A holomorphic map
$c$ on an open set $W'$ is continuous: the set of $w\in W'$ with $c(w)\in U$ and $f(c(w))\in O$ is open, as $f\circ c$ is
continuous on the open set of $w$ with $c(w)\in U$.

\emph{The chart.} Lemma~\ref{cv-origin-hodge-norm-form-torus-chart}, applied to $A$ and its origin, gives an open set
$W\ni0$ in $\mathbf C^g$, a holomorphic map $\sigma:W\to A(\mathbf C)$ with $\sigma(0)$ the origin and $\sigma(W)$ open, an
affine open $U_0\subset A$ and regular functions $t_1,\dots,t_g$ on $U_0$, such that $\sigma(w)\in U_0$ and
$t_i(\sigma(w))=w_i$ for $w\in W$, and such that (a) holds. Let $\phi:\sigma(W)\to W$ be the map $a\mapsto(t_i(a))_i$. It is
continuous, as the $t_i$ are continuous on the points of $U_0$ by the definition of the topology, and it is inverse to the
continuous map $\sigma$. So $\phi$ is a homeomorphism from the open neighbourhood $\sigma(W)$ of the origin onto $W$, with
$\phi$ of the origin equal to $0$. The map $(x,y)\mapsto\sigma(x)+\sigma(y)$ is holomorphic on $W\times W$, by
Lemma~\ref{cv-origin-hodge-norm-form-appell-pair} and the group law of $A$, and it maps $(0,0)$ to the origin, which lies in
$U_0$; so the law $m(x,y)=\phi(\sigma(x)+\sigma(y))$, whose coordinates are the $t_i$ at $\sigma(x)+\sigma(y)$, is holomorphic
near $(0,0)$.

\emph{The uniformization.} By Lemma~\ref{cv-origin-hodge-norm-form-torus-divisible}, $rA(\mathbf C)=A(\mathbf C)$ for every
integer $r\ge1$. Lemma~\ref{cv-origin-hodge-norm-form-torus-exp}, applied to $\Gamma=A(\mathbf C)$, $V=\mathbf C^g$ and the
chart $\phi$, gives a homomorphism $u=e:\mathbf C^g\to A(\mathbf C)$. It is surjective and its kernel $\Lambda$ is discrete and
spans $\mathbf C^g$ over $\mathbf R$, by (1)--(3) of that lemma, as $A(\mathbf C)$ is compact and divisible. As $u$ is
continuous and $\sigma(W)$ is open, $u(v)\in\sigma(W)$ for $v$ near $0$, and then $u(v)=\sigma(h(v))$ with $h=\phi\circ u$,
which is holomorphic near $0$ with values in $W$ by (1) of that lemma. So near any $z\in\mathbf C^g$,
$u(v)=u(z)+\sigma(h(v-z))$ is the sum of a constant point and a holomorphic map, hence holomorphic
by Lemma~\ref{cv-origin-hodge-norm-form-appell-pair} and the group law. As holomorphy is a local condition, for every open
$U\subset A$ and regular $f$ on $U$ the set $u^{-1}(U)$ is open and $f\circ u$ is holomorphic on it. With the trivial line bundle
and $\theta=1$, $u$ is a uniformization of $A$ (Definition~\ref{cv-origin-hodge-norm-gram}(1)).

\emph{The coordinates.} The origin $u(0)=\sigma(0)$ lies in $U_0$, and (a) holds. Near $0$, $(t_i(u(x)))_i=\phi(u(x))$, which has
derivative the identity at $0$ by Lemma~\ref{cv-origin-hodge-norm-form-torus-exp}(1). This is (b).
\end{proof}
% Lean: T.exists_exp_chart_coordinates = AVExp.exists_holUnif: compactSpace_points and isTopologicalGroup_complexPoints
% (Lemma cv-origin-hodge-norm-form-torus-points), HolChart.exists_holomorphic_chart (Lemma cv-origin-hodge-norm-form-torus-chart)
% for W, sigma, U and t, the chart (AbelianExp.chart: continuousOn_value for phi, HolOn.continuousAt for sigma) and
% IsAnalyticGroupChart (hmulHol, HolOn.post of holFamily_lift and the multiplication, HolOn.analyticAt), exp_surjective with
% exists_div (Lemma cv-origin-hodge-norm-form-torus-divisible), periodLattice and isZLattice_periodLattice, HolOn.of_local
% for the holomorphy (large2 = logChart.symm on W', with exp = sigma o large2 near 0 by exp_eventuallyEq; the text's h is
% phi o u, equal to large2 near 0), and hasFDerivAt_chart_exp for (b). The published counterpart is that A(C) is a complex
% torus C^g / Lambda (Mumford, Abelian Varieties, Section 1; Birkenhake-Lange, Chapter 1); not cited, as the locators are not
% checked. Size of the top level: about 150 lines; about 2,200 lines with Lemmas cv-origin-hodge-norm-form-torus-exp-log,
% -torus-exp, -torus-points, -torus-chart-etale, -torus-chart and -torus-divisible, besides the finiteness and flatness of
% multiplication by r on an abelian scheme (about 16,000 lines by the archive's dependency count), which the text cites.

\begin{lemma}\label{cv-origin-hodge-norm-form-tangent}
Let $A$ be a group scheme over $\mathbf C$ with a line bundle trivialized on an open cover, and let $u:\mathbf C^g\to A(\mathbf C)$
with trivializing functions $(\theta_\alpha)$ be a uniformization of $A$ for this line bundle
(Definition~\ref{cv-origin-hodge-norm-gram}(1), with $T=\operatorname{Spec}\mathbf C$, $G=A$ and $\varphi$ the identity).
Suppose that there are an affine open $U_0\subset A$
containing $u(0)$ and regular functions $t_1,\dots,t_g$ on $U_0$ with the properties (a) and (b) of
Lemma~\ref{cv-origin-hodge-norm-form-torus}. Then, with differentials as in Definition~\ref{cv-origin-hodge-norm-gram}(1), every
tangent vector of $A$ at the origin (Definition~\ref{cv-origin-reference-lattices}(1)) has exactly one differential, and every
$a\in\mathbf C^g$ is the differential of exactly one tangent vector of $A$ at the origin.
\end{lemma}
\leanref{Development/AbelianVarietyTheory/Proofs/ToricGeometry/TangentBijectionOfValues.lean}{GlobalTheory.TangentOrigin.tangent_bijection_of_values}
\begin{proof}
As $u$ is a homomorphism, $u(0)$ is the origin. A morphism $v:\operatorname{Spec}\mathbf C[\varepsilon]\to A$ over $\mathbf C$
maps the single point of $\operatorname{Spec}\mathbf C[\varepsilon]$ to a point of $A$; if $v$ is a tangent vector at the
origin, this point is the origin, and for an open $W$ containing it let $v^*_W:\Gamma(W,\mathcal O_A)\to\mathbf C[\varepsilon]$
be the homomorphism of $v$. Its constant part is $f\mapsto f(u(0))$, and $v^*_W(c)=c$ for constants $c$, as $v$ is over
$\mathbf C$.

\emph{Jets.} For $a\in\mathbf C^g$ and an open $W$ containing the origin, let
$J_{a,W}:\Gamma(W,\mathcal O_A)\to\mathbf C[\varepsilon]$ send $f$ to $f(u(0))+\varepsilon\,d(f\circ u)_0(a)$. It is a ring
homomorphism, by the rules for the derivative of a sum and a product, and $J_{a,W}(c)=c$ for constants, whose composites
with $u$ are constant near $0$. It is compatible with restriction to smaller opens. A tangent vector $v$ at the origin has
differential $a$ if and only if $v^*_{U_0}=J_{a,U_0}$. Indeed, if $v$ has differential $a$, then $v^*_W=J_{a,W}$ for every
$W$. Conversely, suppose that $v^*_{U_0}=J_{a,U_0}$, and let $W$ be open containing the origin and $f$ regular on $W$. As
$U_0$ is affine, there is $s\in\Gamma(U_0,\mathcal O_A)$ such that the open set $D(s)$ where $s$ does not vanish contains the
origin and lies in $W$. Its ring is the localization of $\Gamma(U_0,\mathcal O_A)$ at $s$, so the ring homomorphisms $v^*_{D(s)}$
and $J_{a,D(s)}$, which agree on the image of $\Gamma(U_0,\mathcal O_A)$, are equal. By the compatibility with restriction,
$v^*_W(f)$ and $J_{a,W}(f)$ are both the common value at the restriction of $f$ to $D(s)$. So $v$ has differential $a$.

\emph{Tangent vectors have one differential.} Let $v$ be a tangent vector at the origin, and let $a_i$ be the coefficient of
$\varepsilon$ in $v^*_{U_0}(t_i)$. Let $D(f)$ be the coefficient of $\varepsilon$ in $v^*_{U_0}(f)-J_{a,U_0}(f)$. Then $D$ is
additive, $D(c)=0$ for constants, and, as both homomorphisms have constant part $f\mapsto f(u(0))$,
$D(ff')=f(u(0))D(f')+f'(u(0))D(f)$. By (b), the coefficient of $\varepsilon$ in $J_{a,U_0}(t_i)$ is $a_i$, so $D(t_i)=0$. By
(a), $D=0$, so $v^*_{U_0}=J_{a,U_0}$ and $v$ has differential $a$. If $v$ has differentials $a$ and $a'$, then
$J_{a,U_0}(t_i)=J_{a',U_0}(t_i)$ for every $i$, and $a=a'$ by (b).

\emph{Differentials have one tangent vector.} Let $a\in\mathbf C^g$. As $U_0$ is affine, the ring homomorphism $J_{a,U_0}$ is
the homomorphism of a unique morphism $x:\operatorname{Spec}\mathbf C[\varepsilon]\to U_0\subset A$. It lies over $\mathbf C$,
as $J_{a,U_0}(c)=c$ for constants, and its restriction to $\operatorname{Spec}\mathbf C$ has the homomorphism
$f\mapsto f(u(0))$, so it is the origin. So $x$ is a tangent vector at the origin, with differential $a$ by the jet
criterion. If $v$ is a tangent vector at the origin with differential $a$, then $v^*_{U_0}=J_{a,U_0}=x^*_{U_0}$, and $v=x$
since $U_0$ is affine.
\end{proof}
% Lean: TangentOrigin.tangent_bijection_of_values = tangent_bijection: evalPt, jet (with fderiv_val_const, fderiv_val_add,
% fderiv_val_mul, jet_res), tangent_mem_iff and tangent_fst, isDiff_iff (the jet criterion, through
% IsAffineOpen.exists_basicOpen_le and IsLocalization.ringHom_ext), the derivation D for existence, hjet (from (b)) for
% uniqueness, exists_evalPt_eq and ext_evalPt for the converse. Size: about 360 lines.

\subsection{Appell--Humbert data and Hodge metrics}\label{subsection-cv-appell-humbert-hodge-metrics}

\begin{definition}\label{cv-origin-hodge-norm-form-appell-datum}
Let $g\ge0$ and let $\Lambda\subset\mathbf C^g$ be a discrete subgroup that spans $\mathbf C^g$ over $\mathbf R$. For
$x=(x_0,x_1,x_2)\in(\mathbf C^g)^3$ and $I\subset\{0,1,2\}$ put $x_I=\sum_{a\in I}x_a$ and $\varepsilon_I=(-1)^{3-|I|}$. A
\emph{periodic cube datum} for $\Lambda$ consists of:
\begin{enumerate}
\item open sets $W_i\subset\mathbf C^g$, for $i$ in a set $\iota$, that cover $\mathbf C^g$, with $W_i+l=W_i$ for $l\in\Lambda$;
\item holomorphic functions $c_{ij}$ without zeros on $W_i\cap W_j$, with $c_{ij}(z+l)=c_{ij}(z)$ for $l\in\Lambda$, and
  $c_{ij}c_{jk}=c_{ik}$ on $W_i\cap W_j\cap W_k$;
\item for every map $j$ from the set of subsets of $\{0,1,2\}$ to $\iota$, a holomorphic function $T_j$ without zeros on the
  open set $C_j$ of the $x\in(\mathbf C^g)^3$ with $x_I\in W_{j(I)}$ for every $I$, such that $T_j(x+l)=T_j(x)$ for
  $l\in\Lambda^3$ and
  \[
  T_j(x)=\prod_Ic_{j(I)j'(I)}(x_I)^{\varepsilon_I}\,T_{j'}(x)\qquad(x\in C_j\cap C_{j'}).
  \]
\end{enumerate}
\end{definition}
\leanref{Development/AbelianVarietyTheory/Proofs/ToricGeometry/CubeTrivialisationOnUniformisingSpace.lean}{GlobalTheory.cube_trivialisation_on_uniformising_space}

\begin{lemma}\label{cv-origin-hodge-norm-form-appell-pair}
Let $X$ and $Y$ be schemes over $\mathbf C$, $E$ a complex normed vector space and $V\subset E$ a subset. Call a map
$a:E\to X(\mathbf C)$ \emph{holomorphic on $V$} if for every open $U\subset X$ and every regular function $f$ on $U$ the set of
$w\in V$ with $a(w)\in U$ is open and $w\mapsto f(a(w))$ is holomorphic on it. If $a:E\to X(\mathbf C)$ and
$b:E\to Y(\mathbf C)$ are holomorphic on $V$, then so is $(a,b):E\to(X\times_{\mathbf C}Y)(\mathbf C)$.
\end{lemma}
\leanref{Development/AbelianVarietyTheory/Proofs/ToricGeometry/HolFamilyOfIsPullback.lean}{GlobalTheory.holFamily_lift}
\leanref{Development/AbelianVarietyTheory/Proofs/ToricGeometry/HolFamilyOfIsPullback.lean}{GlobalTheory.holFamily_of_affine}
\begin{proof}
Let $c=(a,b)$, and let $p$ and $q$ be the projections of $X\times_{\mathbf C}Y$.

\emph{Affine charts suffice.} Let $Z$ be a scheme over $\mathbf C$ and $c:E\to Z(\mathbf C)$ a map such that every $w\in V$ has
an affine open $V'\subset Z$ containing $c(w)$, for which the set $N$ of $w'\in V$ with $c(w')\in V'$ is open and $g\circ c$ is
holomorphic on $N$ for every $g\in\Gamma(V',\mathcal O_Z)$. Then $c$ is holomorphic on $V$. Indeed, let $U\subset Z$ be open, $f$
regular on $U$, and $w\in V$ with $c(w)\in U$. There is $s\in\Gamma(V',\mathcal O_Z)$ such that the set $D(s)$ where $s$ does
not vanish contains $c(w)$ and lies in $U\cap V'$, and $f=h/s^n$ on $D(s)$ with $h\in\Gamma(V',\mathcal O_Z)$. The set of
$w'\in N$ with $s(c(w'))\neq0$ is an open neighbourhood of $w$ on which $c(w')\in U$ and $f\circ c=(h\circ c)/(s\circ c)^n$ is
holomorphic.

\emph{The charts of the product.} Let $w\in V$, and let $V_X\subset X$ and $V_Y\subset Y$ be affine opens containing $a(w)$ and
$b(w)$. Then $V'=p^{-1}(V_X)\cap q^{-1}(V_Y)=V_X\times_{\mathbf C}V_Y$ is affine, with ring
$\Gamma(V_X,\mathcal O_X)\otimes_{\mathbf C}\Gamma(V_Y,\mathcal O_Y)$. The set of $w'\in V$ with $c(w')\in V'$ is the
intersection of the sets where $a(w')\in V_X$ and $b(w')\in V_Y$, so it is open. The regular functions $g$ on $V'$ with $g\circ c$
holomorphic there form a subring, which contains the functions $f\circ p$ and $f'\circ q$, as $(f\circ p)\circ c=f\circ a$ and
$(f'\circ q)\circ c=f'\circ b$. These generate the ring of $V'$, so the subring is the whole ring, and the previous paragraph
applies.
\end{proof}
% Lean: holFamily_lift through holFamily_of_isPullback for the cartesian square of X x Y over Sp C: holFamily_of_affine
% (exists_basicOpen_le and IsLocalization.surj for f = h / s^n), exists_isAffineOpen_mem_and_subset for VX and VY,
% isIso_pushoutSection_of_isAffineOpen (the ring of VX x VY is the tensor product), the subring T and the induced map from
% the pushout. The rules for values at complex points used throughout are valAt'_rules. Size: about 370 lines with
% valAt'_rules.

\begin{lemma}\label{cv-origin-hodge-norm-form-appell-analytic}
Let $A$ be a commutative group scheme over $\mathbf C$ and $\mathcal L$ a line bundle on $A$, trivialized on an open cover
$(U_i)_{i\in\iota}$ with transition functions $g_{ij}$, with a cubical structure (Definition~\ref{cv-germs-chart}(1)). Let
$u:\mathbf C^g\to A(\mathbf C)$ be a uniformization of $A$ for the trivial line bundle, with $\theta=1$
(Definition~\ref{cv-origin-hodge-norm-gram}(1), with $T=\operatorname{Spec}\mathbf C$, $G=A$ and $\varphi$ the identity), with
kernel $\Lambda$. For a map $j$ from the subsets of $\{0,1,2\}$ to $\iota$, let $t_j$ be the regular function without zeros
that gives the cubical structure on the open set of $A^3$ where the sum of the coordinates in $I$ lies in $U_{j(I)}$ for every
$I$, in the trivialization of $\Theta$ by those of $\mathcal L$. Then the sets $W_i=u^{-1}(U_i)$, the functions
$c_{ij}=g_{ij}\circ u$ and the functions $T_j(x)=t_j(u(x_0),u(x_1),u(x_2))$ form a periodic cube datum for $\Lambda$
(Definition~\ref{cv-origin-hodge-norm-form-appell-datum}).
\end{lemma}
\leanref{Development/AbelianVarietyTheory/Proofs/ToricGeometry/CubeTrivialisationOnUniformisingSpace.lean}{GlobalTheory.cube_trivialisation_on_uniformising_space}
\begin{proof}
(1) The sets $W_i$ are open by Definition~\ref{cv-origin-hodge-norm-gram}(1), and they cover $\mathbf C^g$, as the $U_i$ cover $A$.
As $u$ is a homomorphism with kernel $\Lambda$, $u(z+l)=u(z)$ for $l\in\Lambda$, so $W_i+l=W_i$. (2) The functions $c_{ij}$ are
holomorphic by Definition~\ref{cv-origin-hodge-norm-gram}(1), without zeros as the $g_{ij}$ are units, invariant under $\Lambda$,
and they satisfy the cocycle relation of the $g_{ij}$. (3) By Lemma~\ref{cv-origin-hodge-norm-form-appell-pair}, applied twice,
$x\mapsto y(x)=(u(x_0),u(x_1),u(x_2))$ is holomorphic on $(\mathbf C^g)^3$ as a map to $A^3(\mathbf C)$. As $u$ is a
homomorphism, the sum of the coordinates of $y(x)$ in $I$ is $u(x_I)$, so $y$ maps $C_j$ into the domain of $t_j$. Hence $T_j$
is holomorphic on $C_j$ and has no zeros. The cubical structure is a nowhere vanishing section of $\Theta$, so
$t_j=\prod_I(g_{j(I)j'(I)}\circ m_I)^{\varepsilon_I}\,t_{j'}$, where $m_I$ is the sum of the coordinates in $I$; at $y(x)$ this is
the relation between $T_j$ and $T_{j'}$. Finally $y(x+l)=y(x)$ for $l\in\Lambda^3$.
\end{proof}
% Lean: cube_trivialisation_on_uniformising_space, with valAt'_rules (the rules ValueRules for values at complex points) and
% holFamily_lift (Lemma cv-origin-hodge-norm-form-appell-pair) as hypotheses: u.hol, the cover from hL.cover, val_unit_ne_zero,
% val_cocycle, ex_add_period; T j x is the value of t.s at pt3 u x, with tripleProductPoint_comp_partialSum,
% pt3_preimage_theta_whole_iff, val_gen and valAt_theta_transition_eq_prod_zpow (val_linComb) for the transformation law and
% tripleProductPoint_add_period for the periodicity. Only the section t of Theta and its IsGen are used from the cubical
% structure. Size: about 250 lines.

\begin{lemma}\label{cv-origin-hodge-norm-form-appell-faces}
Let $(W_i,c_{ij},T_j)$ be a periodic cube datum for $\Lambda\subset\mathbf C^g$
(Definition~\ref{cv-origin-hodge-norm-form-appell-datum}), let $i_0$ be an index with $0\in W_{i_0}$, and let $F=T_{j_0}$ for the
constant map $j_0$ with value $i_0$. Then for $x\in(\mathbf C^g)^3$ near $0$:
\begin{enumerate}
\item $F(x)=F(0)$ if $x_a=0$ for some $a\in\{0,1,2\}$;
\item $F(2x)\prod_IF(x_I,x_I,-x_I)^{\varepsilon_I}\,F(0)^3=F(x)^3F(-x)$.
\end{enumerate}
\end{lemma}
\leanref{Development/AbelianVarietyTheory/Proofs/ToricGeometry/CubeSectionFacesAndDoubling.lean}{GlobalTheory.cube_section_faces_and_doubling}
\leanref{Development/AbelianVarietyTheory/Proofs/CovolumeBounds/EntireApplyEqOfPeriodicOfSpanReal.lean}{GlobalTheory.entire_apply_eq_of_periodic_of_span_real}
\begin{proof}
\emph{Periodic functions.} Let $U_k$, for $k$ in a set, be open sets covering $(\mathbf C^g)^3$ with $U_k+l=U_k$ for
$l\in\Lambda^3$, and $R_k$ holomorphic functions on $U_k$ that agree on overlaps, with $R_k(y+l)=R_k(y)$. They glue to a
holomorphic function $R$ on $(\mathbf C^g)^3$ invariant under $\Lambda^3$. As $\Lambda^3$ spans $(\mathbf C^g)^3$ over $\mathbf R$,
$|R|$ is bounded by its maximum on a compact parallelepiped, and $R$ is constant by Liouville's theorem.

(1) Let $a\in\{0,1,2\}$, and for $y\in(\mathbf C^g)^3$ let $y^a$ be $y$ with $y_a$ replaced by $0$, so that
$y^a_I=y^a_{I\setminus\{a\}}$. For a map $j$ from the subsets of $\{0,1,2\}$ to $\iota$, let $j^a(I)=j(I\setminus\{a\})$, let $U_j$ be
the set of $y$ with $y^a\in C_{j^a}$, and put $R_j(y)=T_{j^a}(y^a)$. On $U_j\cap U_k$, the transformation law gives
$R_j(y)/R_k(y)=\prod_Ic_{j(I\setminus\{a\})k(I\setminus\{a\})}(y^a_{I\setminus\{a\}})^{\varepsilon_I}=1$, as the factors for $I$ and
$I\cup\{a\}$, for $a\notin I$, are equal with $\varepsilon_{I\cup\{a\}}=-\varepsilon_I$. The $U_j$ cover, as $j(J)$ can be chosen
with $y^a_J\in W_{j(J)}$; they are stable under $\Lambda^3$, and the $R_j$ are invariant, as $(y+l)^a=y^a+l^a$. So the $R_j$ glue
to a constant. If $x_a=0$ and $x\in C_{j_0}$, then $x=x^a$, and $x$ and $0$ lie in $U_{j_0}$, so $F(x)=R_{j_0}(x)=R_{j_0}(0)=F(0)$.

(2) For $y\in(\mathbf C^g)^3$ and $I\subset\{0,1,2\}$ let $d_I(y)=(y_I,y_I,-y_I)$. Its partial sums are $0$ for $J=\emptyset$,
$\{0,2\}$, $\{1,2\}$, they are $y_I$ for $J=\{0\}$, $\{1\}$, $\{0,1,2\}$, $-y_I$ for $J=\{2\}$ and $2y_I$ for $J=\{0,1\}$. For a
triple $k=(k_1,k_2,k_3)$ of maps from the subsets to $\iota$, let $U_k$ be the set of $y$ with $y\in C_{k_1}$, $-y\in C_{k_2}$
and $2y\in C_{k_3}$. For $y\in U_k$, $d_I(y)\in C_\delta$ for the map $\delta$ with value $i_0$ at the first three sets,
$k_1(I)$ at the next three, $k_2(I)$ at $\{2\}$ and $k_3(I)$ at $\{0,1\}$; put $D_k(I,y)=T_\delta(d_I(y))$ and
\[
R_k(y)=T_{k_3}(2y)\prod_ID_k(I,y)^{\varepsilon_I}\big/\big(T_{k_1}(y)^3T_{k_2}(-y)\big).
\]
Changing $k$ to $k'$ multiplies $D_k(I,y)$ by $c(y_I)^3c(-y_I)c(2y_I)^{-1}$, with the transition functions $c$ between the
charts of $k$ and $k'$ at $I$, since the factor at the partial sums $0$ is $c_{i_0i_0}(0)=1$. So it multiplies
$\prod_ID_k(I,y)^{\varepsilon_I}$ by the cube of the factor of $T_{k_1}(y)$, times the factor of $T_{k_2}(-y)$, divided by the
factor of $T_{k_3}(2y)$, and $R_k=R_{k'}$ on $U_k\cap U_{k'}$. The $U_k$ cover, are stable under $\Lambda^3$, and the $R_k$ are
invariant. So they glue to a constant. Let $x$ be near $0$, with $x$, $-x$ and $2x$ in $C_{j_0}$, and $k=(j_0,j_0,j_0)$. Then
$D_k(I,x)=F(d_I(x))$, and $R_k(x)=R_k(0)=F(0)\prod_IF(0)^{\varepsilon_I}/F(0)^4=F(0)^{-3}$, as $\sum_I\varepsilon_I=0$. This is (2).
\end{proof}
% Lean: cube_section_faces_and_doubling, through CubeFace.face and CubeFace.doubling: glue_const (exists of a global section,
% entire_apply_eq_of_periodic_of_span_real, span_triples), prod_alt_erase for the cancellation in (1); diagChart, sum_diag,
% diag_mem, diag_prod and hone for the charts of d_I(y) in (2). Size: about 400 lines.

\begin{lemma}\label{cv-origin-hodge-norm-form-appell-local}
Let $g\ge0$, and let $F$ be a holomorphic function near $0$ in $(\mathbf C^g)^3$ with $F(0)\neq0$ that satisfies (1) and (2) of
Lemma~\ref{cv-origin-hodge-norm-form-appell-faces} for $x$ near $0$, with the notation of
Definition~\ref{cv-origin-hodge-norm-form-appell-datum}. Then there is a holomorphic function $k$ near $0$ in $\mathbf C^g$ with
$k(0)=0$ and $F(x)=F(0)\prod_I\exp(k(x_I))^{\varepsilon_I}$ for $x$ near $0$.
\end{lemma}
\leanref{Development/AbelianVarietyTheory/Proofs/ToricGeometry/SeriesApplyEqZeroOfFace.lean}{GlobalTheory.LocalCube.exists_third_difference}
\leanref{Development/AbelianVarietyTheory/Proofs/ToricGeometry/DifferentiableOnDoublingIterate.lean}{GlobalTheory.exists_doubling_solution_pi}
\leanref{Development/AbelianVarietyTheory/Proofs/ToricGeometry/EventuallyEqZeroOfIsBigOCubeOfDoubling.lean}{GlobalTheory.eventually_eq_zero_of_isBigO_cube_of_doubling}
\begin{proof}
Use the maximum norm on $(\mathbf C^g)^3$. For a function $w$ near $0$ in $\mathbf C^g$ put $\Delta w(x)=\sum_I\varepsilon_Iw(x_I)$;
note that $\sum_I\varepsilon_I=0$.

\emph{A cubic bound.} Let $f(x)=\log(F(x)/F(0))$ with the principal branch, a holomorphic function near $0$ with $F=F(0)e^f$
there. By (1), $f(x)=0$ near $0$ when some $x_a=0$. Then $|f(x)|\le C|x|^3$ near $0$. Indeed, choose $R>0$ such that on the ball
$|x|<R$, $f$ is holomorphic, $|f|<1$, and $f(x)=0$ when some $x_a=0$. Let $0<|y|<R$, put $\rho=R/|y|>1$ and
$G(t_0,t_1,t_2)=f(t_0y_0,t_1y_1,t_2y_2)$ for $|t_a|<\rho$. Then $G$ is holomorphic in each variable, $|G|\le1$, and $G=0$ when
some $t_a=0$. The Schwarz lemma in $t_0$ gives $|G|\le|t_0|/\rho$; applied in $t_1$ to $G/t_0$ it gives $|G|\le|t_0t_1|/\rho^2$,
and in $t_2$, $|G|\le|t_0t_1t_2|/\rho^3$. At $t=(1,1,1)$, $|f(y)|\le\rho^{-3}=|y|^3/R^3$.

\emph{The doubling equation.} Let $h$ be holomorphic on a ball $B(0,r)\subset\mathbf C^g$ with $|h(z)|\le C|z|^3$ there. Put
$h_0(z)=h(z/2)$ and $h_{n+1}(z)=3h_n(z/2)+h_n(-z/2)$ on $B(0,r)$. By induction $|h_n(z)|\le2^{-n}C|z|^3$, as
$4\cdot2^{-n}C|z/2|^3=2^{-n-1}C|z|^3$. So $k=\sum_{n\ge0}h_n$ converges uniformly on $B(0,r)$ to a holomorphic function with
$|k(z)|\le2C|z|^3$, and for $|z|<r/2$,
$k(2z)=h(z)+\sum_nh_{n+1}(2z)=h(z)+3k(z)+k(-z)$.

\emph{A vanishing criterion.} Let $G$ be a function on a ball $|z|<\delta$ of a complex normed space with $|G(z)|\le C|z|^3$ and
$G(2z)=3G(z)+G(-z)$ when $|z|<\delta$. Then $G=0$ on the ball. Indeed $M(w)=|G(w)|+|G(-w)|$ satisfies $M(2w)\le4M(w)$, so
$M(z)\le4^nM(2^{-n}z)\le4^n\cdot2C\,2^{-3n}|z|^3$ for every $n$, which tends to $0$.

\emph{The logarithm of (2).} Let $\varphi(s)=f(s,s,-s)$ for $s\in\mathbf C^g$ near $0$; by the cubic bound
$|\varphi(s)|\le C|s|^3$. Write the two sides of (2) at $x$ near $0$ with $F=F(0)e^f$. As $\sum_I\varepsilon_I=0$, the left side is
$F(0)^4\exp(f(2x)+\Delta\varphi(x))$ and the right side is $F(0)^4\exp(3f(x)+f(-x))$. So
$f(2x)+\Delta\varphi(x)-3f(x)-f(-x)$ lies in $2\pi i\mathbf Z$. It tends to $0$ as $x\to0$, so it vanishes near $0$.

\emph{Conclusion.} Let $k$ be the solution of the doubling equation for $h=-\varphi$: $k(2z)-3k(z)-k(-z)=-\varphi(z)$ near $0$ and
$|k(z)|\le C'|z|^3$, so that $k(0)=0$. As $(2x)_I=2x_I$ and $(-x)_I=-x_I$,
$\Delta k(2x)-3\Delta k(x)-\Delta k(-x)=\sum_I\varepsilon_I\big(k(2x_I)-3k(x_I)-k(-x_I)\big)=-\Delta\varphi(x)$. So $G=f-\Delta k$
satisfies $G(2x)=3G(x)+G(-x)$ near $0$ by the previous paragraph, and $|G(x)|\le C''|x|^3$, as $|x_I|\le3|x|$. By the vanishing
criterion $G=0$ near $0$, and $F(x)=F(0)e^{\Delta k(x)}=F(0)\prod_I\exp(k(x_I))^{\varepsilon_I}$.
\end{proof}
% Lean: LocalCube.exists_third_difference, with exists_doubling_solution_pi (doublingIterate, norm_doublingIterate_le,
% hasSum_geometric_two, SCV.differentiableOn_tsum_of_locally_normal) and eventually_eq_zero_of_isBigO_cube_of_doubling
% (eq_zero_of_cube_bound_of_doubling) as hypotheses: f = log(F/F 0) (AnalyticAt.clog), hfax, the private
% isBigO_norm_cube_of_analytic_vanishing_coords (three applications of the Schwarz lemma,
% Complex.dist_le_div_mul_dist_of_mapsTo_ball), phi = f o M with M s = (s, s, -s), k from the doubling solution for -phi,
% hEt and hE (exp_eq_exp_iff_exists_int, the integer is 0 by the bound 2 pi), hGeq, and Gf = f - D k vanishing. Size: about
% 550 lines with the two inputs.

\begin{lemma}\label{cv-origin-hodge-norm-form-appell-global}
Let $(W_i,c_{ij},T_j)$ be a periodic cube datum for $\Lambda\subset\mathbf C^g$
(Definition~\ref{cv-origin-hodge-norm-form-appell-datum}). Then there are holomorphic functions $f_i$ without zeros on $W_i$, with
$f_i=c_{ij}f_j$ on $W_i\cap W_j$, and $\kappa\in\mathbf C$, such that $T_j(x)=\kappa\prod_If_{j(I)}(x_I)^{\varepsilon_I}$ for every
$j$ and $x\in C_j$.
\end{lemma}
\leanref{Development/AbelianVarietyTheory/Proofs/ToricGeometry/ExistsLocalFrameCompatibleCube.lean}{GlobalTheory.exists_localFrame_compatible_cube}
\leanref{Development/AbelianVarietyTheory/Proofs/ToricGeometry/ExistsGlobalFrameOfLocal.lean}{GlobalTheory.CubeFrame.exists_global_frame}
\leanref{Development/AbelianVarietyTheory/Proofs/ToricGeometry/ExistsGlobalFrameOfLocal.lean}{GlobalTheory.CubeFrame.exists_global_frame_of_local}
\begin{proof}
For $\rho>0$ let $B_\rho\subset\mathbf C^g$ be the open ball of radius $\rho$ about $0$, and $U_\rho$ the set of
$x\in(\mathbf C^g)^3$ with $x_I\in B_\rho$ for every $I$, a convex open set containing $0$. Say that a family $(f_i)$ satisfies the
\emph{cube identity on $B_\rho$} if each $f_i$ is holomorphic without zeros on $W_i\cap B_\rho$, $f_i=c_{ij}f_j$ on
$W_i\cap W_j\cap B_\rho$, and $T_j(x)=\kappa\prod_If_{j(I)}(x_I)^{\varepsilon_I}$ for every $j$ and $x\in C_j\cap U_\rho$.

\emph{Near $0$.} Choose $i_0$ with $0\in W_{i_0}$, and let $F=T_{j_0}$ for the constant map $j_0$ with value $i_0$, holomorphic
near $0$ with $F(0)\neq0$. By Lemma~\ref{cv-origin-hodge-norm-form-appell-faces}, $F$ satisfies the hypotheses of
Lemma~\ref{cv-origin-hodge-norm-form-appell-local}, which gives $k$ with $F(x)=F(0)\prod_I\exp(k(x_I))^{\varepsilon_I}$ near $0$.
Choose $r>0$ with $k$ holomorphic on $B_r\subset W_{i_0}$, and put $\kappa=F(0)$ and $f_i=c_{ii_0}e^k$ on $W_i\cap B_r$. These are
holomorphic without zeros, $f_i=c_{ij}f_j$ by the cocycle relation, and for $x\in C_j$ near $0$ the transformation law gives
$T_j(x)=\prod_Ic_{j(I)i_0}(x_I)^{\varepsilon_I}F(x)=\kappa\prod_If_{j(I)}(x_I)^{\varepsilon_I}$.

\emph{From near $0$ to $U_\rho$.} Let $(f_i)$ satisfy the first two conditions on $B_\rho$, and the identity for $x$ near $0$. On
$C_j\cap U_\rho$ put $Q_j(x)=T_j(x)/\prod_If_{j(I)}(x_I)^{\varepsilon_I}$. By the transformation laws of the $T_j$ and of the
$f_i$, $Q_j=Q_{j'}$ on overlaps, and the $C_j$ cover, so the $Q_j$ glue to a holomorphic function on $U_\rho$. It is $\kappa$ near
$0$, hence on the connected open set $U_\rho$, by the identity theorem. So $(f_i)$ satisfies the cube identity on $B_\rho$.

\emph{Doubling the radius.} Let $(f_i)$ satisfy the cube identity on $B_\rho$. Then $\kappa\neq0$, as $T_{j_0}(0)\neq0$. For
$w\in\mathbf C^g$ let $h(w)=(w/2,w/2,-w/2)$ and $D=\{0,1\}$; then $h(w)_D=w$, and $|h(w)_I|\le|w|/2$ for $I\neq D$. For
$w\in W_i\cap B_{2\rho}$ and $j$ with $h(w)\in C_j$ put
\[
f'_i(w)=c_{ij(D)}(w)\,\kappa\prod_{I\neq D}f_{j(I)}(h(w)_I)^{\varepsilon_I}\Big/T_j(h(w)).
\]
By the transformation laws and the cocycle relation, this does not depend on $j$. So $f'_i$ is holomorphic without zeros on
$W_i\cap B_{2\rho}$, and $f'_i=c_{ii'}f'_{i'}$. For $w\in B_\rho$ all $h(w)_I$ lie in $B_\rho$, and the cube identity at $h(w)$,
where $\varepsilon_D=-1$, gives $f'_i(w)=c_{ij(D)}(w)f_{j(D)}(w)=f_i(w)$. So $(f'_i)$ satisfies the identity near $0$, and by the
previous paragraph the cube identity on $B_{2\rho}$.

\emph{Conclusion.} Starting from $B_r$, doubling gives families satisfying the cube identity on $B_{2^nr}$ for every $n$, each
agreeing with the previous one on the smaller ball. They glue to holomorphic functions $f_i$ on $W_i$ with the required
properties, as every $x\in(\mathbf C^g)^3$ lies in some $U_{2^nr}$.
\end{proof}
% Lean: exists_localFrame_compatible_cube with cube_section_faces_and_doubling and LocalCube.exists_third_difference as
% hypotheses (F = T (fun _ => i0), SCV.analyticAt_of_differentiableOn, fr i = c i i0 * exp k on the ball of radius min eps
% delta, kappa = F 0); CubeFrame.exists_global_frame through exists_global_frame_of_local: cube_eq_on_ball (exists_glue and
% SCV.eqOn_of_eventuallyEq on the convex set U), exists_double (halfPoint, sum_halfPoint, norm_sum_halfPoint_le, hsplit),
% and the sequence s n on the balls of radius 2^n r glued by exists_glue. Size: about 400 lines.

\begin{lemma}\label{cv-origin-hodge-norm-form-appell-factor}
Let $(W_i,c_{ij},T_j)$ be a periodic cube datum for $\Lambda\subset\mathbf C^g$
(Definition~\ref{cv-origin-hodge-norm-form-appell-datum}), and let $(f_i)$ and $\kappa$ be as in
Lemma~\ref{cv-origin-hodge-norm-form-appell-global}. Then there are a map $a:\Lambda\to\mathbf C^\times$ and an additive map
$l\mapsto B_l$ from $\Lambda$ to the $\mathbf C$-linear forms on $\mathbf C^g$, such that $a(l+m)=a(l)a(m)\exp(B_l(m))$ for
$l,m\in\Lambda$ and $f_i(z)=a(l)\exp(B_l(z))f_i(z+l)$ for every $i$, $l\in\Lambda$ and $z\in W_i$.
\end{lemma}
\leanref{Development/AbelianVarietyTheory/Proofs/ToricGeometry/AutomorphyFactorExpLinearOfCube.lean}{GlobalTheory.automorphyFactor_exp_linear_of_cube}
\leanref{Development/AbelianVarietyTheory/Proofs/ToricGeometry/AutomorphyFactorExpLinearOfCube.lean}{GlobalTheory.entire_hom_eq_cexp_fderiv}
\begin{proof}
\emph{The factor of automorphy.} For $l\in\Lambda$ and $z\in W_i$ put $\alpha_l(z)=f_i(z)/f_i(z+l)$; it is defined, as
$W_i+l=W_i$. It does not depend on $i$, as $c_{ij}(z+l)=c_{ij}(z)$. So $\alpha_l$ is a holomorphic function on $\mathbf C^g$
without zeros, $\alpha_0=1$, and $\alpha_{l+m}(z)=\alpha_l(z)\alpha_m(z+l)$.

\emph{It is exponential.} Let $l\in\Lambda$ and $y,z\in\mathbf C^g$. Put $x=(0,y,z)$ and $L=(l,0,0)\in\Lambda^3$, and choose $j$
with $x_I\in W_{j(I)}$ for every $I$; then $(x+L)_I=x_I+L_I\in W_{j(I)}$ too. The identity of
Lemma~\ref{cv-origin-hodge-norm-form-appell-global} at $x$ and $x+L$, the periodicity $T_j(x+L)=T_j(x)$ and $\kappa\neq0$ give
$\prod_I\alpha_{L_I}(x_I)^{\varepsilon_I}=1$. Only the four sets $I$ containing $0$ contribute, with $x_I=0$, $y$, $z$, $y+z$ and
$\varepsilon_I=1,-1,-1,1$, so $\alpha_l(0)\alpha_l(y+z)=\alpha_l(y)\alpha_l(z)$. Hence $e(z)=\alpha_l(z)/\alpha_l(0)$ is a holomorphic
homomorphism from $\mathbf C^g$ to $\mathbf C^\times$. Let $B_l$ be its derivative at $0$. As $e(w+z)=e(w)e(z)$, the derivative of
$e$ at $w$ is $e(w)B_l$, so $e\exp(-B_l)$ has derivative $0$ and is constant, equal to $1$: $e(z)=\exp(B_l(z))$.

\emph{Conclusion.} Put $a(l)=\alpha_l(0)$. Then $f_i(z)=\alpha_l(z)f_i(z+l)=a(l)\exp(B_l(z))f_i(z+l)$, and
$a(l+m)=\alpha_{m+l}(0)=\alpha_m(0)\alpha_l(m)=a(l)a(m)\exp(B_l(m))$. With this, $\alpha_{m+l}(z)=\alpha_m(z)\alpha_l(z+m)$ gives
$\exp(B_{l+m}(z))=\exp(B_l(z)+B_m(z))$ for every $z$; two linear forms with the same exponential have the same derivative at $0$,
so $B_{l+m}=B_l+B_m$.
\end{proof}
% Lean: automorphyFactor_exp_linear_of_cube: alpha l z = fr (sigma z) z / fr (sigma z) (z + l), halpha, hαc, hcubeα (with
% x = ![0, y, z], L = ![l, 0, 0], prod_finset_fin_three), entire_hom_eq_cexp_fderiv for B0 l, ha, and clm_eq_of_cexp_eq for the
% additivity. Size: about 200 lines.

\begin{lemma}\label{cv-origin-hodge-norm-form-appell-normal}
Let $\Lambda\subset\mathbf C^g$ be a discrete subgroup that spans $\mathbf C^g$ over $\mathbf R$, $a:\Lambda\to\mathbf C^\times$ a
map and $l\mapsto B_l$ an additive map from $\Lambda$ to the $\mathbf C$-linear forms on $\mathbf C^g$, with
$a(l+m)=a(l)a(m)\exp(B_l(m))$ for $l,m\in\Lambda$. Then there are a Hermitian matrix $H\in M_g(\mathbf C)$, a map
$\chi:\Lambda\to\mathbf C$ with $|\chi|=1$ and a holomorphic function $Q$ on $\mathbf C^g$ with $Q(0)=0$ such that
\[
a(l)\exp\big(B_l(z)+Q(z)-Q(z+l)\big)=\chi(l)\exp\big(\pi z^{\mathsf T}H\bar l+\tfrac\pi2\,l^{\mathsf T}H\bar l\big)
\]
for $l\in\Lambda$ and $z\in\mathbf C^g$. Moreover $\operatorname{Im}l^{\mathsf T}H\bar m\in\mathbf Z$ for $l,m\in\Lambda$.
\end{lemma}
\leanref{Development/AbelianVarietyTheory/Proofs/ToricGeometry/AppellHumbertNormalFormOfLinearExponent.lean}{GlobalTheory.appellHumbert_normal_form_of_linear_exponent}
\begin{proof}
A $\mathbf Z$-basis of $\Lambda$ is an $\mathbf R$-basis of $\mathbf C^g$, so every additive map from $\Lambda$ to a real vector
space extends to an $\mathbf R$-linear map on $\mathbf C^g$.

\emph{The matrices $M$ and $N$.} Let $e_1,\dots,e_g$ be the standard basis of $\mathbf C^g$. Extend the additive map
$l\mapsto(B_l(e_j))_j$ from $\Lambda$ to $\mathbf C^g$ to an $\mathbf R$-linear map on $\mathbf C^g$, and for $u\in\mathbf C^g$ let
$B_u$ be the $\mathbf C$-linear form whose value at $e_j$ is the $j$-th coordinate of its image; so $u\mapsto B_u$ is
$\mathbf R$-linear and extends $l\mapsto B_l$. As $a(l+m)=a(m+l)$, $\exp(B_l(m))=\exp(B_m(l))$, so
$B_l(m)-B_m(l)\in2\pi i\mathbf Z$ for $l,m\in\Lambda$. By $\mathbf R$-linearity in $u$ and then in $v$, $S(u,v)=B_u(v)-B_v(u)$
has real part $0$ for all $u,v\in\mathbf C^g$. Put $\beta_{kj}=B_{e_k}(e_j)$ and $\gamma_{kj}=B_{ie_k}(e_j)$. The real parts of
$S(e_k,e_j)=\beta_{kj}-\beta_{jk}$, $S(ie_k,ie_j)=i(\gamma_{kj}-\gamma_{jk})$ and $S(ie_k,e_j)=\gamma_{kj}-i\beta_{jk}$ vanish:
\[
\operatorname{Re}\beta_{kj}=\operatorname{Re}\beta_{jk},\qquad\operatorname{Im}\gamma_{kj}=\operatorname{Im}\gamma_{jk},\qquad
\operatorname{Re}\gamma_{kj}=-\operatorname{Im}\beta_{jk}.
\]
Let $M$ and $N$ be the matrices with entries $M_{kj}=\frac12(\beta_{kj}-i\gamma_{kj})$ and $N_{kj}=\frac12(\beta_{kj}+i\gamma_{kj})$.
Writing $u=\sum_k(\operatorname{Re}u_k\,e_k+\operatorname{Im}u_k\,ie_k)$ gives
$B_u(e_j)=\sum_k(\operatorname{Re}u_k\,\beta_{kj}+\operatorname{Im}u_k\,\gamma_{kj})=\sum_k(u_kM_{kj}+\bar u_kN_{kj})$, so
\[
B_u(z)=u^{\mathsf T}Mz+\bar u^{\mathsf T}Nz\qquad(u,z\in\mathbf C^g).
\]
By the three identities, $\operatorname{Re}M_{kj}=\frac12(\operatorname{Re}\beta_{kj}+\operatorname{Im}\gamma_{kj})$ and
$\operatorname{Im}M_{kj}=\frac12(\operatorname{Im}\beta_{kj}+\operatorname{Im}\beta_{jk})$ are symmetric in $(k,j)$, so $M$ is
symmetric; $\operatorname{Re}N_{kj}=\frac12(\operatorname{Re}\beta_{kj}-\operatorname{Im}\gamma_{kj})$ is symmetric and
$\operatorname{Im}N_{kj}=\frac12(\operatorname{Im}\beta_{kj}-\operatorname{Im}\beta_{jk})$ is antisymmetric, so $N_{jk}=\overline{N_{kj}}$.

\emph{The forms $\mathcal A$ and $\mathcal K$.} Put $\mathcal A(u,v)=u^{\mathsf T}Mv$ and $\mathcal K(z,u)=\bar u^{\mathsf T}Nz$, so
that $B_u(z)=\mathcal A(u,z)+\mathcal K(z,u)$. As $M$ is symmetric, $\mathcal A$ is a symmetric complex bilinear form. As
$\bar N=N^{\mathsf T}$, $\overline{\mathcal K(u,v)}=v^{\mathsf T}\bar N\bar u=\bar u^{\mathsf T}Nv=\mathcal K(v,u)$. Let
$H=\pi^{-1}N^{\mathsf T}$, a Hermitian matrix with $\pi z^{\mathsf T}H\bar u=\mathcal K(z,u)$. For $l,m\in\Lambda$,
$B_l(m)-B_m(l)=\mathcal K(m,l)-\mathcal K(l,m)=-2i\operatorname{Im}\mathcal K(l,m)$, as $\mathcal A$ is symmetric; this lies in
$2\pi i\mathbf Z$, so $\operatorname{Im}l^{\mathsf T}H\bar m=\pi^{-1}\operatorname{Im}\mathcal K(l,m)\in\mathbf Z$.

\emph{The character.} The relation at $l=m=0$ gives $a(0)=1$. The map
$\phi_0(l)=\log|a(l)|-\frac12\operatorname{Re}\mathcal A(l,l)-\frac12\operatorname{Re}\mathcal K(l,l)$ is additive on $\Lambda$: by
the relation, $\log|a(l+m)|=\log|a(l)|+\log|a(m)|+\operatorname{Re}B_l(m)$, and
$\operatorname{Re}(\mathcal A(l+m,l+m)+\mathcal K(l+m,l+m))$ exceeds the sum of the terms for $l$ and $m$ by
$2\operatorname{Re}\mathcal A(l,m)+\operatorname{Re}(\mathcal K(l,m)+\mathcal K(m,l))=2\operatorname{Re}B_l(m)$. Extend $\phi_0$ to an
$\mathbf R$-linear form, and write it as $u\mapsto\operatorname{Re}(c^{\mathsf T}u)$ with $c\in\mathbf C^g$. Put
\[
\chi(l)=a(l)\exp\big(-\tfrac12\mathcal A(l,l)-\tfrac12\mathcal K(l,l)-c^{\mathsf T}l\big),\qquad
Q(z)=\tfrac12\mathcal A(z,z)+c^{\mathsf T}z.
\]
Then $|\chi(l)|=|a(l)|\exp(-\log|a(l)|)=1$, and $Q$ is holomorphic with $Q(0)=0$. As $\mathcal A$ is symmetric,
$Q(z)-Q(z+l)=-\mathcal A(l,z)-\frac12\mathcal A(l,l)-c^{\mathsf T}l$, so $B_l(z)+Q(z)-Q(z+l)=\mathcal K(z,l)-\frac12\mathcal A(l,l)-c^{\mathsf T}l$
and $a(l)\exp(B_l(z)+Q(z)-Q(z+l))=\chi(l)\exp(\mathcal K(z,l)+\frac12\mathcal K(l,l))$. This is the claim, as
$\mathcal K(z,l)=\pi z^{\mathsf T}H\bar l$.
\end{proof}
% Lean: appellHumbert_normal_form_of_linear_exponent: AppellHumbert.exists_real_extension (a Z-basis of Lambda, IsZLattice)
% for F with B l z = F l . z (the text's B_u is z -> F u . z), hskew, hlin, hD1 and hD; the matrices M and N
% (M k j = (F e_k j - I F (I e_k) j) / 2, N with +) with F u = u M + conj(u) N (hFMN), hsym (hD at (e_k, e_j), (I e_k, I e_j)
% and (I e_k, e_j)), hMsym and hNherm; A and K, hBF, hKherm, Hm = pi^{-1} N^T, hint and Ef; ha0, phi0,
% exists_real_extension, c, chi and Q. Size: about 270 lines.

\begin{lemma}\label{cv-origin-hodge-norm-form-appell}
Let $A$ be a commutative group scheme over $\mathbf C$ with a line bundle $\mathcal L$, trivialized at the origin, with a
cubical structure, and trivialized on an open cover $(U_\alpha)$ of $A$ with transition functions $g_{\alpha\beta}$. Let
$u:\mathbf C^g\to A(\mathbf C)$ be a uniformization of $A$ for the trivial line bundle, with $\theta=1$
(Definition~\ref{cv-origin-hodge-norm-gram}(1), with $T=\operatorname{Spec}\mathbf C$, $G=A$ and $\varphi$ the identity), with
kernel $\Lambda$. Then there are trivializing functions $\theta_\alpha$ that make the same $u$ a uniformization of $A$ for
$\mathcal L$, a
Hermitian matrix $H\in M_g(\mathbf C)$ and a map $\chi:\Lambda\to\mathbf C$ with $|\chi|=1$, such that
\[
\theta_\alpha(x)=\chi(l)\exp\big(\pi x^{\mathsf T}H\bar l+\tfrac\pi2\,l^{\mathsf T}H\bar l\big)\,\theta_\alpha(x+l)
\]
for every $\alpha$, $l\in\Lambda$ and $x\in u^{-1}(U_\alpha)$.
\end{lemma}
\leanref{Development/AutomorphicSpectrum/Proofs/GeometryProofs/ThetaMultiplierNormalFormStatement.lean}{GlobalTheory.T.theta_multiplier_normal_form}
\leanref{Development/AutomorphicSpectrum/Proofs/GeometryProofs/ThetaMultiplierNormalForm.lean}{GlobalTheory.CompanionTheory.Ar.exists_appellHumbert_frame_of_cubical}
\begin{proof}
\emph{The periodic cube datum.} By Lemma~\ref{cv-origin-hodge-norm-form-appell-analytic}, applied to the cover $(U_\alpha)$ and
the cubical structure of $\mathcal L$ (Definition~\ref{cv-germs-chart}(1)), the sets $u^{-1}(U_\alpha)$, the functions
$g_{\alpha\beta}\circ u$ and the values of the cubical structure along $u$ form a periodic cube datum for $\Lambda$
(Definition~\ref{cv-origin-hodge-norm-form-appell-datum}). Lemma~\ref{cv-origin-hodge-norm-form-appell-global} gives holomorphic
functions $f_\alpha$ without zeros on $u^{-1}(U_\alpha)$, with $f_\alpha=(g_{\alpha\beta}\circ u)f_\beta$, and $\kappa\in\mathbf C$
with the cube identity. Lemma~\ref{cv-origin-hodge-norm-form-appell-factor} then gives $a:\Lambda\to\mathbf C^\times$ and an additive
map $l\mapsto B_l$ to the linear forms, with $a(l+m)=a(l)a(m)\exp(B_l(m))$ and $f_\alpha(x)=a(l)\exp(B_l(x))f_\alpha(x+l)$.

\emph{The normal form.} Lemma~\ref{cv-origin-hodge-norm-form-appell-normal} gives a Hermitian $H$, a map $\chi$ with $|\chi|=1$ and
a holomorphic $Q$ on $\mathbf C^g$. Put $\theta_\alpha(x)=f_\alpha(x)e^{Q(x)}$ on $u^{-1}(U_\alpha)$. These functions are
holomorphic without zeros and $\theta_\alpha=(g_{\alpha\beta}\circ u)\theta_\beta$ where both are defined, so with them $u$ is a
uniformization of $A$ for $\mathcal L$ (Definition~\ref{cv-origin-hodge-norm-gram}(1)). For $l\in\Lambda$ and
$x\in u^{-1}(U_\alpha)$,
\[
\theta_\alpha(x)=a(l)\exp\big(B_l(x)+Q(x)-Q(x+l)\big)\,\theta_\alpha(x+l)
=\chi(l)\exp\big(\pi x^{\mathsf T}H\bar l+\tfrac\pi2\,l^{\mathsf T}H\bar l\big)\,\theta_\alpha(x+l).
\]
\end{proof}
% Lean: T.theta_multiplier_normal_form = exists_appellHumbert_frame_of_cubical: valAt'_rules and holFamily_lift (Lemma
% cv-origin-hodge-norm-form-appell-pair), cube_trivialisation_on_uniformising_space (-appell-analytic) with cs.t and cs.isGen,
% cube_section_faces_and_doubling (-appell-faces), LocalCube.exists_third_difference with exists_doubling_solution_pi and
% eventually_eq_zero_of_isBigO_cube_of_doubling (-appell-local), exists_localFrame_compatible_cube and
% CubeFrame.exists_global_frame (-appell-global), automorphyFactor_exp_linear_of_cube (-appell-factor) and
% appellHumbert_normal_form_of_linear_exponent (-appell-normal); the trivializing functions are fr i z * exp(Q z), with the
% same ex as u0.
% The rigidification of L and the face, symmetry and cocycle conditions of the cubical structure are not used. The published
% counterpart is the Appell-Humbert theorem (Birkenhake-Lange, Chapter 2; Mumford, Abelian Varieties, Section 2); not cited,
% as the locators are not checked. Size: about 2,800 lines with the lemmas before this one.

\begin{lemma}\label{cv-origin-hodge-norm-form-theta}
Let $A$ be an abelian scheme over $\operatorname{Spec}\mathbf C$ with a line bundle $\mathcal L$, trivialized on an open cover
$(U_\alpha)$ with transition functions $g_{\alpha\beta}$, that is relatively ample by sections
(Definition~\ref{cv-origin-hodge-norm-form-ample}). Let $u:\mathbf C^g\to A(\mathbf C)$ with trivializing functions
$(\theta_\alpha)$ be a
uniformization of $A$ for $\mathcal L$ (Definition~\ref{cv-origin-hodge-norm-gram}(1), with $T=\operatorname{Spec}\mathbf C$,
$G=A$ and $\varphi$ the identity), with kernel $\Lambda$. Let $H\in M_g(\mathbf C)$ and $\chi:\Lambda\to\mathbf C$, put
$e_l(x)=\chi(l)\exp(\pi x^{\mathsf T}H\bar l+\frac\pi2l^{\mathsf T}H\bar l)$ for $l\in\Lambda$ and $x\in\mathbf C^g$, and
suppose that $\theta_\alpha(x)=e_l(x)\theta_\alpha(x+l)$ for every $\alpha$, $l\in\Lambda$ and $x\in u^{-1}(U_\alpha)$. Then
there are $N\ge1$ and holomorphic functions $\Theta_j$ on $\mathbf C^g$, for $j$ in a set $J$, such that:
\begin{enumerate}
\item $\Theta_j(x+l)=e_l(x)^N\Theta_j(x)$ for every $j$, $l\in\Lambda$ and $x$;
\item some $\Theta_j$ is not identically zero;
\item for every nonzero $w\in\mathbf C^g$ there are $j$, $t\in\mathbf C$ and $x$ with $\Theta_j(x+tw)\neq\Theta_j(x)$.
\end{enumerate}
\end{lemma}
\leanref{Development/AutomorphicSpectrum/Proofs/GeometryProofs/ExistsThetaBasisOfFactorStatement.lean}{GlobalTheory.T.exists_theta_basis_of_factor}
\begin{proof}
Let $N\ge1$ and sections $s_0,\dots,s_M$ of $\mathcal L^{\otimes N}$ over $A$, with the functions $r_{ij}$ on the affine open
sets $A_{s_i}$, be given by Definition~\ref{cv-origin-hodge-norm-form-ample} for the affine open $\operatorname{Spec}\mathbf C$.
Trivialize $\mathcal L^{\otimes N}$ on $U_\alpha$ by the $N$-th power of the trivialization of $\mathcal L$, so that $s_i$ is
given on $U_\alpha$ by a regular function $s_{i,\alpha}$, with $s_{i,\alpha}=g_{\alpha\beta}^Ns_{i,\beta}$ on $U_\alpha\cap U_\beta$.

\emph{The theta functions.} For $x\in\mathbf C^g$ choose $\alpha$ with $u(x)\in U_\alpha$ and put
$\Theta_i(x)=s_{i,\alpha}(u(x))/\theta_\alpha(x)^N$. This does not depend on $\alpha$, as
$\theta_\alpha=(g_{\alpha\beta}\circ u)\theta_\beta$ where both are defined. Near $x$ it is a quotient of holomorphic functions
with a denominator without zeros, so $\Theta_i$ is holomorphic. As $u(x+l)=u(x)$ for $l\in\Lambda$,
$\Theta_i(x+l)=s_{i,\alpha}(u(x))/\theta_\alpha(x+l)^N=e_l(x)^N\Theta_i(x)$, which is (1). Also $\Theta_i(x)\neq0$ if and only if
$u(x)\in A_{s_i}$, and then $\Theta_j(x)=r_{ij}(u(x))\Theta_i(x)$ for every $j$, as $s_j=r_{ij}s_i$ on $A_{s_i}$.

(2) The origin $u(0)$ lies in some $A_{s_i}$, so $\Theta_i(0)\neq0$.

(3) Let $w\in\mathbf C^g$ be such that $\Theta_j(x+tw)=\Theta_j(x)$ for every $j$, $t$ and $x$; we show that $w=0$. Then
$\Theta_j(tw)=\Theta_j(0)$ for every $j$ and $t$. Choose $i$ with $\Theta_i(0)\neq0$. For every $t\in\mathbf C$,
$\Theta_i(tw)=\Theta_i(0)\neq0$, so $u(tw)$ and $u(0)$ lie in $A_{s_i}$, and
$r_{ij}(u(tw))=\Theta_j(tw)/\Theta_i(tw)=\Theta_j(0)/\Theta_i(0)=r_{ij}(u(0))$ for every $j$. The points $u(tw)$ and $u(0)$ of
the affine scheme $A_{s_i}$ over $\mathbf C$ are given by homomorphisms from its ring to $\mathbf C$ over $\mathbf C$, which
agree on the generators $r_{ij}$, so $u(tw)=u(0)$ and $tw\in\Lambda$ for every $t\in\mathbf C$. As $\Lambda$ is discrete, there
is $\varepsilon>0$ such that $0$ is the only element of $\Lambda$ of norm less than $\varepsilon$. If $w\neq0$, then
$t=\varepsilon/(2|w|)$ gives an element $tw$ of $\Lambda$ of norm $\varepsilon/2$, a contradiction. So $w=0$.
\end{proof}
% Lean: T.exists_theta_basis_of_factor = SecTheta.exists_theta_family: the RelAmple data for U = top, exists_theta (the
% quotient by the product of the trivializing functions over an N-tuple of charts, well defined by quot_eq), theta_differentiable,
% theta_add for (1), theta_ne_zero_iff, val_ratio, eq_of_generators (the generators of the ring of the affine nonzero
% locus), lp_injective, and eq_zero_of_smul_mem for the discreteness. The Lean uses N-tuples of charts of L for the
% charts of L^N; the book takes the diagonal tuples, which suffice. The hypothesis that A is abelian is not used. Size:
% about 450 lines.

\begin{lemma}\label{cv-origin-hodge-norm-form-positive}
Let $\Lambda\subset\mathbf C^g$ be a discrete subgroup that spans $\mathbf C^g$ over $\mathbf R$, let $H\in M_g(\mathbf C)$ be
Hermitian, $\chi:\Lambda\to\mathbf C$ a map with $|\chi|=1$ and $N\ge1$, and put
$e_l(x)=\chi(l)\exp(\pi x^{\mathsf T}H\bar l+\frac\pi2l^{\mathsf T}H\bar l)$ for $l\in\Lambda$ and $x\in\mathbf C^g$. Let
$\Theta_j$, for $j$ in a set $J$, be holomorphic functions on $\mathbf C^g$ with $\Theta_j(x+l)=e_l(x)^N\Theta_j(x)$ for every
$j$, $l$ and $x$, not all identically zero, such that for every nonzero $w\in\mathbf C^g$ there are $j$, $t\in\mathbf C$ and $x$
with $\Theta_j(x+tw)\neq\Theta_j(x)$. Then $x^{\mathsf T}H\bar x>0$ for every nonzero $x\in\mathbf C^g$.
\end{lemma}
\leanref{Development/AutomorphicSpectrum/Proofs/GeometryProofs/ThetaFunctionsFiniteDimStatement.lean}{GlobalTheory.T.theta_functions_finite_dim}
\leanref{Development/AbelianVarietyTheory/Proofs/ToricGeometry/ReFormPosOfThetaNeZero.lean}{GlobalTheory.ThetaPositivity.exists_norm_le}
\leanref{Development/AbelianVarietyTheory/Proofs/ToricGeometry/ReFormPosOfThetaNeZero.lean}{GlobalTheory.ThetaPositivity.eq_zero_of_neg}
\begin{proof}
Put $H'=NH$, a Hermitian matrix, and $q(x)=\operatorname{Re}x^{\mathsf T}H'\bar x$. Then
$e_l(x)^N=\chi(l)^N\exp(\pi x^{\mathsf T}H'\bar l+\frac\pi2l^{\mathsf T}H'\bar l)$ with $|\chi(l)^N|=1$. As $H'$ is Hermitian,
$y^{\mathsf T}H'\bar x=\overline{x^{\mathsf T}H'\bar y}$, so
\[
q(x+y)=q(x)+2\operatorname{Re}x^{\mathsf T}H'\bar y+q(y)\qquad(x,y\in\mathbf C^g).\tag{$*$}
\]

\emph{Growth.} Every $\Theta_j$ satisfies $|\Theta_j(x)|\le C_j\exp(\pi q(x)/2)$ for some $C_j$. Indeed, the function
$F(x)=|\Theta_j(x)|\exp(-\pi q(x)/2)$ is continuous, and, as
$|e_l(x)^N|=\exp(\pi\operatorname{Re}x^{\mathsf T}H'\bar l+\frac\pi2q(l))$, $(*)$ gives $F(x+l)=F(x)$ for $l\in\Lambda$. As
$\Lambda$ is discrete and spans $\mathbf C^g$ over $\mathbf R$, it has a $\mathbf Z$-basis that is an $\mathbf R$-basis of
$\mathbf C^g$, and every $x$ is $l+y$ with $l\in\Lambda$ and $y$ in the compact parallelepiped $P$ spanned by this basis. So
$F$ is bounded by its maximum $C_j$ on $P$.

\emph{Vanishing.} Let $\theta$ be a holomorphic function on $\mathbf C^g$ with $|\theta(x)|\le C\exp(\pi q(x)/2)$, and suppose that
$q(w)<0$ for some $w$. Then $\theta=0$. Indeed, fix $x$ and put $f(t)=\theta(x+tw)$ for $t\in\mathbf C$. By $(*)$ and
$q(tw)=|t|^2q(w)$,
\[
q(x+tw)=q(x)+2\operatorname{Re}\big(\bar t\,x^{\mathsf T}H'\bar w\big)+|t|^2q(w)\le q(x)+2B|t|-\alpha|t|^2,
\]
with $B=|x^{\mathsf T}H'\bar w|$ and $\alpha=-q(w)>0$. So $|f(t)|\le C\exp(\pi(q(x)+2B|t|-\alpha|t|^2)/2)$, which tends to $0$ as
$|t|\to\infty$. The holomorphic function $f$ is bounded, hence constant by Liouville's theorem, with limit $0$ at infinity;
so $\theta(x)=f(0)=0$.

\emph{Semidefiniteness.} Some $\Theta_{j_0}$ is not identically zero, so by the two previous paragraphs $q(w)\ge0$ for every
$w$, that is, $\operatorname{Re}w^{\mathsf T}H\bar w\ge0$. As $H$ is Hermitian, $w^{\mathsf T}H\bar w$ is real. For
$v\in\mathbf C^g$ and $w=\bar v$ this gives $v^*Hv=w^{\mathsf T}H\bar w\ge0$, so the Hermitian matrix $H$ is positive
semidefinite.

\emph{Definiteness.} Let $x\neq0$ with $x^{\mathsf T}H\bar x=0$, that is, $v^*Hv=0$ for $v=\bar x$. As $H$ is positive
semidefinite, $Hv=0$, so $H'\overline{tx}=\bar tNH\bar x=0$ for every $t\in\mathbf C$. By $(*)$, for every $y\in\mathbf C^g$
and $t\in\mathbf C$, $q(y+tx)=q(y)+2\operatorname{Re}y^{\mathsf T}H'\overline{tx}+\operatorname{Re}(tx)^{\mathsf T}H'\overline{tx}=q(y)$.
So for each $j$ the holomorphic function $t\mapsto\Theta_j(y+tx)$ is bounded by $C_j\exp(\pi q(y)/2)$, hence constant by
Liouville's theorem: $\Theta_j(y+tx)=\Theta_j(y)$ for every $j$, $t$ and $y$. This contradicts the hypothesis for $w=x$. So
$x^{\mathsf T}H\bar x>0$ for every nonzero $x$.
\end{proof}
% Lean: T.theta_functions_finite_dim = HodgeFormAna.pos_of_theta_power: H' = N Hm, htr', ThetaPositivity.exists_norm_le
% (the periodic function F on the closure of ZSpan.fundamentalDomain, with form_add_self), ThetaPositivity.eq_zero_of_neg
% (re_form_add_smul_of_isHermitian, the bound tending to 0 on the cocompact filter, Differentiable.apply_eq_of_tendsto_cocompact),
% hnn and him, Matrix.posSemidef_iff_dotProduct_mulVec, PosSemidef.dotProduct_mulVec_zero_iff for hker, form_add_self with
% hker' for hq, and Differentiable.apply_eq_apply_of_bounded. Size: about 210 lines.

\begin{lemma}\label{cv-origin-hodge-norm-form-lie}
Let $T$ be a scheme, $G$ a commutative group scheme over $T$, $z$ a complex point of $T$ and $T_z$ the Lie algebra of $G$ at $z$
(Definition~\ref{cv-origin-reference-lattices}(1)). Let $A$ be a group scheme over $\mathbf C$ with a homomorphism
$\varphi:A\to G$ over $z$ that identifies $A$ with $G\times_T\operatorname{Spec}\mathbf C$, with a line bundle trivialized on an
open cover, and let $u:\mathbf C^g\to A(\mathbf C)$ with trivializing functions $(\theta_\alpha)$ be a uniformization of $A$
for this line bundle
(Definition~\ref{cv-origin-hodge-norm-gram}(1), with $T=\operatorname{Spec}\mathbf C$, $G=A$ and the identity). Suppose that
every tangent vector of $A$ at the origin has exactly one differential, and every $a\in\mathbf C^g$ is the differential of
exactly one tangent vector of $A$ at the origin. Then there is a $\mathbf C$-linear isomorphism $D:T_z\to\mathbf C^g$ such that
every $v\in T_z$ is $\varphi\circ x$ for a tangent vector $x$ of $A$ at the origin with differential $D(v)$.
\end{lemma}
\leanref{Development/AutomorphicSpectrum/Proofs/GeometryProofs/TangentSpaceEquivHolUnifStatement.lean}{GlobalTheory.T.tangent_space_equiv_hol_unif}
\leanref{Development/AbelianVarietyTheory/Proofs/ThetaLawProofs/SubsingletonPrimeSpectrum.lean}{GlobalTheory.UnifDifferential.isDiff_mul}
\leanref{Development/AbelianVarietyTheory/Proofs/ThetaLawProofs/SubsingletonPrimeSpectrum.lean}{GlobalTheory.UnifDifferential.isDiff_scale}
\begin{proof}
\emph{Tangent vectors of $A$ and of $G$.} Morphisms $x:\operatorname{Spec}\mathbf C[\varepsilon]\to A$ over $\mathbf C$
correspond, through $x\mapsto\varphi\circ x$, to morphisms $\operatorname{Spec}\mathbf C[\varepsilon]\to G$ over $T$, where
$\operatorname{Spec}\mathbf C[\varepsilon]$ lies over $T$ through $z$, as $\varphi$ identifies $A$ with
$G\times_T\operatorname{Spec}\mathbf C$. This correspondence respects the products for the group laws, as $\varphi$ is a
homomorphism, and it maps the origin of $A$ to the unit section of $G$ at $z$. So $x$ is a tangent vector of $A$ at the origin
if and only if $\varphi\circ x\in T_z$. For $v\in T_z$ let $x_v$ be the tangent vector of $A$ with $\varphi\circ x_v=v$, and
$D(v)$ its differential. Then $x_{v+v'}=x_v\cdot x_{v'}$ and $x_{cv}=x_v\circ\operatorname{Spec}(\varepsilon\mapsto c\varepsilon)$
for $c\in\mathbf C$, by Definition~\ref{cv-origin-reference-lattices}(1).

\emph{Differentials add and scale.} Let $x$ and $x'$ be tangent vectors at the origin with differentials $a$ and $a'$. Let
$R=\mathbf C[\varepsilon_1,\varepsilon_2]/(\varepsilon_1,\varepsilon_2)^2$, let $\rho_k:\mathbf C[\varepsilon]\to R$ send
$\varepsilon$ to $\varepsilon_k$, and let $p:\operatorname{Spec}R\to A$ be the product of $x\circ\operatorname{Spec}\rho_1$ and
$x'\circ\operatorname{Spec}\rho_2$. Let $\sigma,\sigma_1,\sigma_2:R\to\mathbf C[\varepsilon]$ send $(\varepsilon_1,\varepsilon_2)$
to $(\varepsilon,\varepsilon)$, $(\varepsilon,0)$ and $(0,\varepsilon)$. Then $p\circ\operatorname{Spec}\sigma=x\cdot x'$,
$p\circ\operatorname{Spec}\sigma_1=x$ and $p\circ\operatorname{Spec}\sigma_2=x'$, since $x$ and $x'$ restrict to the origin at
$\varepsilon=0$. For an open $U$ containing the origin and $f$ regular on $U$, write $f\circ p=c+d_1\varepsilon_1+d_2\varepsilon_2$.
The coefficients of $\varepsilon$ in $f\circ(x\cdot x')$, $f\circ x$ and $f\circ x'$ are $d_1+d_2$, $d_1$ and $d_2$, so the first is
$d(f\circ u)_0(a)+d(f\circ u)_0(a')=d(f\circ u)_0(a+a')$, and $x\cdot x'$ has differential $a+a'$. The coefficient of $\varepsilon$ in
$f\circ x\circ\operatorname{Spec}(\varepsilon\mapsto c\varepsilon)$ is $c$ times that in $f\circ x$, so this tangent vector has
differential $ca$.

\emph{The isomorphism.} By the previous paragraphs, $x_{v+v'}$ has differentials $D(v+v')$ and $D(v)+D(v')$, and $x_{cv}$ has
$D(cv)$ and $cD(v)$; by the uniqueness of differentials, $D$ is $\mathbf C$-linear. If $D(v)=D(v')$, then $x_v$ and $x_{v'}$
are tangent vectors with the same differential, hence equal, and $v=\varphi\circ x_v=v'$. For $a\in\mathbf C^g$ let $y$ be the
tangent vector at the origin with differential $a$; then $v=\varphi\circ y\in T_z$, $x_v=y$ and $D(v)=a$. So $D$ is an
isomorphism, and $v=\varphi\circ x_v$ with $x_v$ of differential $D(v)$.
\end{proof}
% Lean: T.tangent_space_equiv_hol_unif = HodgeLin.exists_linearEquiv: tanEquiv (LiePull.pts, a multiplicative bijection),
% isTangent_iff with unit_left, xadd (M.add), xsmul (M.smul), UnifDifferential.isDiff_mul (snd_pullFn_mul with
% TrivSqZeroExt K (Fin 2 -> K), the maps q, and epsTo_zero_comp) and isDiff_scale, inj and surj, LinearEquiv.ofBijective.
% The book's T_z is the Lean's LieModel V of G at z, with M.pt the tangent vector of v. Size: about 380 lines.

\begin{lemma}\label{cv-origin-hodge-norm-form-lift-open}
Let $V$ and $W$ be finite-dimensional complex vector spaces, $\Gamma$ an abelian group, $u:W\to\Gamma$ a surjective homomorphism
and $G:V\to\Gamma$ a homomorphism. Let $S\subset\Gamma$ be a subset containing $0$, let $Y$ be a Hausdorff space, and let
$f_i:\Gamma\to Y$, for $i$ in a set $I$, be maps that separate the points of $S$: if $a,b\in S$ and $f_i(a)=f_i(b)$ for every
$i$, then $a=b$. Suppose that $u^{-1}(S)$ and $G^{-1}(S)$ are open, and that for every $i$ the map $f_i\circ u$ is continuous on
$u^{-1}(S)$ and $f_i\circ G$ is continuous on $G^{-1}(S)$. Then for every neighbourhood $N$ of $0$ in $W$ there is a
neighbourhood $N'$ of $0$ in $V$ with $G(N')\subset u(N)$.
\end{lemma}
\leanref{Development/AbelianVarietyTheory/Proofs/CubicalLineBundles/ExistsUniqueLinearLiftOfHolomorphicAbelianVariet.lean}{GlobalTheory.tendsto_map_nhds_of_separating}
\begin{proof}
Fix norms on $V$ and $W$. As $u^{-1}(S)$ is an open set containing $0$, there is $R>0$ such that the closed ball $\bar B(0,R)$
of $W$ lies in $N\cap u^{-1}(S)$. Put $r=R/4$, and let $T$ be the set of $z\in V$ with $G(z)=u(w)$ for some $w\in\bar B(0,2r)$.
As $\bar B(0,2r)\subset u^{-1}(S)$, $T\subset G^{-1}(S)$.

\emph{A closure property of $T$.} Let $z$ be in the closure of $T$ with $G(z)\in S$; we show that $z\in T$. Choose $z_n\in T$
with $z_n\to z$, and $w_n\in\bar B(0,2r)$ with $u(w_n)=G(z_n)$. As the ball is compact, a subsequence of $(w_n)$ converges to
some $w_0\in\bar B(0,2r)$. Let $i\in I$. As $f_i\circ u$ is continuous on the open set $u^{-1}(S)$, which contains $w_0$, the
values $f_i(u(w_n))$ converge to $f_i(u(w_0))$ along the subsequence. As $f_i\circ G$ is continuous on the open set
$G^{-1}(S)$, which contains $z$, the values $f_i(G(z_n))=f_i(u(w_n))$ converge to $f_i(G(z))$. As $Y$ is Hausdorff,
$f_i(u(w_0))=f_i(G(z))$. Both $u(w_0)$ and $G(z)$ lie in $S$, so $u(w_0)=G(z)$, and $z\in T$.

\emph{A translate of $T$ with interior.} Let $(c_j)_{j\ge0}$ be a dense sequence in $W$, and let $B_j$ be the set of $z\in V$
with $G(z)=u(w)$ for some $w$ in the open ball $B(c_j,r)$. Every $z\in V$ lies in some $B_j$: as $u$ is surjective, $G(z)=u(w)$
for some $w$, and $w\in B(c_j,r)$ for some $j$. So $V$ is the union of the closures $\bar B_j$, and by the Baire category
theorem some $\bar B_j$ has a nonempty interior. This interior is a nonempty open subset of $\bar B_j$, so it contains a point
$z_0\in B_j$; let $w_0\in B(c_j,r)$ with $u(w_0)=G(z_0)$. For $z\in B_j$, with $w\in B(c_j,r)$ and $u(w)=G(z)$, we have
$|w-w_0|<2r$ and $u(w-w_0)=G(z-z_0)$, so $z-z_0\in T$. As $z\mapsto z-z_0$ is continuous, $z-z_0$ lies in the closure of $T$
for every $z\in\bar B_j$.

\emph{Conclusion.} As $z_0$ is an interior point of $\bar B_j$ and $G^{-1}(S)$ is an open set containing $0$, there is a
neighbourhood $N'$ of $0$ in $V$ such that $z+z_0\in\bar B_j$ and $G(z)\in S$ for $z\in N'$. For $z\in N'$, $z$ lies in the
closure of $T$ by the previous paragraph, and so $z\in T$ by the closure property: $G(z)=u(w)$ with $|w|\le2r<R$, hence
$w\in N$.
\end{proof}
% Lean: tendsto_map_nhds_of_separating, for V complete and W proper normed groups: hN' and nhds_basis_closedBall for R; T,
% hKS and hTS; hcl (mem_closure_iff_seq_limit, isCompact_closedBall.tendsto_subseq, tendsto_nhds_unique) for the closure
% property; exists_dense_seq, B, hB and nonempty_interior_of_iUnion_of_closed (the Baire category theorem), z0, w0 and hBT;
% h1 and the final filter_upwards. Size: about 55 lines.

\begin{lemma}\label{cv-origin-hodge-norm-form-lift-real}
Let $E$ and $F$ be finite-dimensional real vector spaces, $\Gamma$ an abelian group, and $G:E\to\Gamma$ and $u:F\to\Gamma$
homomorphisms such that the kernel $\Lambda$ of $u$ is discrete. Let $\gamma:E\to F$ be a map, continuous at $0$, with
$u(\gamma(z))=G(z)$ for all $z$ in a neighbourhood of $0$. Then there is exactly one $\mathbf R$-linear map $L:E\to F$ with
$u\circ L=G$.
\end{lemma}
\leanref{Development/AbelianVarietyTheory/Proofs/HomogeneousIdeals/ExistsUniqueRealLinearOfEventuallyAddModDiscrete.lean}{GlobalTheory.existsUnique_realLinear_lift_of_germ}
\leanref{Development/AbelianVarietyTheory/Proofs/HomogeneousIdeals/ExistsUniqueRealLinearOfEventuallyAddModDiscrete.lean}{GlobalTheory.existsUnique_realLinear_of_eventually_add_mod_discrete}
\leanref{Development/AbelianVarietyTheory/Proofs/HomogeneousIdeals/ExistsUniqueRealLinearOfEventuallyAddModDiscrete.lean}{GlobalTheory.realLinear_eq_of_sub_mem_discrete}
\leanref{Development/AbelianVarietyTheory/Proofs/HomogeneousIdeals/TendstoUniformlyOnDoublingLogOfContraction.lean}{GlobalTheory.exists_addMonoidHom_eventuallyEq}
\begin{proof}
Fix norms on $E$ and $F$. As $\Lambda$ is discrete, there is a neighbourhood of $0$ in $F$ that meets $\Lambda$ only in $0$. So a
map to $F$ with values in $\Lambda$ that tends to $0$ along a filter is eventually $0$ along it.

\emph{Additivity near $0$.} As $u(\gamma(0))=G(0)=0$, the map $\gamma_0(z)=\gamma(z)-\gamma(0)$ is continuous at $0$, with
$\gamma_0(0)=0$ and $u(\gamma_0(z))=G(z)$ for $z$ near $0$. For $x$ and $y$ near $0$, the vectors $x$, $y$ and $x+y$ are near
$0$, so $u(\gamma_0(x+y)-\gamma_0(x)-\gamma_0(y))=G(x+y)-G(x)-G(y)=0$: the difference
$\delta(x,y)=\gamma_0(x+y)-\gamma_0(x)-\gamma_0(y)$ lies in $\Lambda$. As $\delta(x,y)\to0$ when $(x,y)\to(0,0)$, $\delta$ vanishes
near $(0,0)$: there is $\rho>0$ with $\gamma_0(x+y)=\gamma_0(x)+\gamma_0(y)$ for $|x|,|y|<\rho$.

\emph{Extension.} For $v\in E$ choose $j\ge0$ with $|2^{-j}v|<\rho$, and put $L(v)=2^j\gamma_0(2^{-j}v)$. This does not depend
on $j$: if $|2^{-j}v|<\rho$, then $\gamma_0(2^{-j}v)=2\gamma_0(2^{-j-1}v)$ by the additivity near $0$. For $v,w\in E$, choose $j$
with $2^{-j}v$, $2^{-j}w$ and $2^{-j}(v+w)$ of norm less than $\rho$; the additivity near $0$ gives $L(v+w)=L(v)+L(w)$. With
$j=0$, $L=\gamma_0$ on the ball $B(0,\rho)$, so $L$ is continuous at $0$. An additive map continuous at $0$ is continuous, and a
continuous additive map between real vector spaces is $\mathbf R$-linear, as it is $\mathbf Q$-linear.

\emph{The lift.} For $z$ near $0$, $u(L(z))=u(\gamma_0(z))=G(z)$, so the homomorphisms $u\circ L$ and $G$ agree on a ball
$B(0,\rho')$. Two homomorphisms $E\to\Gamma$ that agree on a ball around $0$ are equal, as every $v$ is $2^j(2^{-j}v)$ with
$2^{-j}v$ in the ball for large $j$. So $u\circ L=G$.

\emph{Uniqueness.} Let $L'$ be $\mathbf R$-linear with $u\circ L'=G$, let $x\in E$ and $d=L(x)-L'(x)$. For $t\in\mathbf R$,
$td=L(tx)-L'(tx)$ lies in $\Lambda$, and $td\to0$ as $t\to0$; so $td=0$ for some $t\neq0$, and $d=0$.
\end{proof}
% Lean: existsUnique_realLinear_lift_of_germ, for normed real spaces E and F, with uniqueness among continuous real-linear
% maps: h0, and existsUnique_realLinear_of_eventually_add_mod_discrete for z -> g z - g 0 (eventually_eq_zero_of_tendsto_of_
% mem_discrete for the additivity near 0, exists_addMonoidHom_eventuallyEq for the extension by halving, with the least j
% where the text takes any j, continuous_of_continuousAt_zero and AddMonoidHom.toRealLinearMap); hL' by
% addMonoidHom_eq_of_eventuallyEq; realLinear_eq_of_sub_mem_discrete for the uniqueness. Size: about 120 lines.

\begin{lemma}\label{cv-origin-hodge-norm-form-lift}
Let $A$ be an abelian scheme over $\operatorname{Spec}\mathbf C$ with a line bundle trivialized on an open cover, and let
$u:\mathbf C^g\to A(\mathbf C)$ with trivializing functions $(\theta_\alpha)$ be a uniformization of $A$ for this line bundle
(Definition~\ref{cv-origin-hodge-norm-gram}(1), with $T=\operatorname{Spec}\mathbf C$, $G=A$ and the identity), such that every
tangent vector of $A$ at the origin has exactly one differential and every $a\in\mathbf C^g$ is the differential of exactly one
tangent vector of $A$ at the origin. Let $m\ge0$ and let $G:\mathbf C^m\to A(\mathbf C)$ be a homomorphism such that for every
open $U\subset A$ and every regular function $f$ on $U$ the set $G^{-1}(U)$ is open and $f\circ G$ is holomorphic on it. Then
there is exactly one $\mathbf C$-linear map $F:\mathbf C^m\to\mathbf C^g$ with $u\circ F=G$.
\end{lemma}
\leanref{Development/AbelianVarietyTheory/Proofs/CubicalLineBundles/ExistsUniqueLinearLiftOfHolomorphicAbelianVariet.lean}{GlobalTheory.existsUnique_linear_lift_of_isHolTo}
\leanref{Development/AbelianVarietyTheory/Proofs/CubicalLineBundles/ExistsUniqueLinearLiftOfHolomorphicAbelianVariet.lean}{GlobalTheory.existsUnique_linear_lift_of_separating}
\begin{proof}
Let $\Lambda$ be the kernel of $u$. The group $A(\mathbf C)$ is commutative, as it is the image of $\mathbf C^g$ under the
homomorphism $u$. Let $U\subset A$ be an affine open set containing the origin $u(0)$, let $S\subset A(\mathbf C)$ be the set of
points in $U$, and for $s\in\Gamma(U,\mathcal O_A)$ let $f_s:A(\mathbf C)\to\mathbf C$ be the value of $s$ at the points of $U$,
and $0$ elsewhere.

\emph{The functions $f_s$.} (i) The origin lies in $S$. (ii) The $f_s$ separate the points of $S$: a point $a$ of $A(\mathbf C)$
in the affine scheme $U$ is determined by its homomorphism $\Gamma(U,\mathcal O_A)\to\mathbf C$ over $\mathbf C$, which is
$s\mapsto f_s(a)$. (iii) By Definition~\ref{cv-origin-hodge-norm-gram}(1) and the hypothesis on $G$, the sets
$u^{-1}(S)=u^{-1}(U)$ and $G^{-1}(S)=G^{-1}(U)$ are open, and $f_s\circ u$ and $f_s\circ G$ are holomorphic on them. (iv) If
$a\in\mathbf C^g$ and $d(f_s\circ u)_0(a)=0$ for every $s$, then $a=0$. Indeed, let $x$ be the tangent vector of $A$ at the
origin with differential $a$. By the jet criterion in the proof of Lemma~\ref{cv-origin-hodge-norm-form-tangent}, whose
argument applies to every affine open set containing the origin, the homomorphism of $x$ on $U$ is
$s\mapsto f_s(u(0))+\varepsilon\,d(f_s\circ u)_0(a)=f_s(u(0))$. This is also the jet of $0$, so by the same criterion $x$ has
differential $0$, and $a=0$ by the uniqueness of differentials.

\emph{A continuous local lift.} As $\Lambda$ is discrete, there is a neighbourhood $N$ of $0$ in $\mathbf C^g$ such that $x-y$
lies in $\Lambda$ for $x,y\in N$ only if $x=y$; so $u$ is injective on $N$. Let $\sigma:A(\mathbf C)\to\mathbf C^g$ be a map with
$\sigma(u(x))=x$ for $x\in N$. Lemma~\ref{cv-origin-hodge-norm-form-lift-open}, applied to $u$, $G$, $S$ and the $f_s$, shows that
for every neighbourhood $M$ of $0$ in $\mathbf C^g$ and every $z$ near $0$, $G(z)=u(w)$ for some $w\in M\cap N$, and then
$\sigma(G(z))=w\in M$ and $u(\sigma(G(z)))=G(z)$. So $\gamma=\sigma\circ G$ is continuous at $0$, with $u(\gamma(z))=G(z)$ for $z$
near $0$.

\emph{The real lift.} Lemma~\ref{cv-origin-hodge-norm-form-lift-real}, applied to the real vector spaces underlying $\mathbf C^m$
and $\mathbf C^g$, to $G$, $u$ and $\gamma$, gives exactly one $\mathbf R$-linear map $F:\mathbf C^m\to\mathbf C^g$ with
$u\circ F=G$.

\emph{Complex linearity.} Let $s\in\Gamma(U,\mathcal O_A)$. As $u^{-1}(U)$ and $G^{-1}(U)$ are open neighbourhoods of $0$,
$f_s\circ u$ and $f_s\circ G$ are holomorphic near $0$, and $(f_s\circ u)\circ F=f_s\circ G$, as $u\circ F=G$. As $F$ is
$\mathbf R$-linear with $F(0)=0$, the chain rule gives $d(f_s\circ u)_0\circ F=d(f_s\circ G)_0$, and both derivatives are
$\mathbf C$-linear. So for $c\in\mathbf C$ and $z\in\mathbf C^m$,
\[
d(f_s\circ u)_0(F(cz))=d(f_s\circ G)_0(cz)=c\,d(f_s\circ G)_0(z)=d(f_s\circ u)_0(cF(z)).
\]
By (iv), applied to $a=F(cz)-cF(z)$, $F(cz)=cF(z)$. So $F$ is $\mathbf C$-linear with $u\circ F=G$. A $\mathbf C$-linear map
$F'$ with $u\circ F'=G$ is $\mathbf R$-linear, so $F'=F$ by the uniqueness in Lemma~\ref{cv-origin-hodge-norm-form-lift-real}.
\end{proof}
% Lean: existsUnique_linear_lift_of_isHolTo: the CommGroup structure on the points from u.ex_onto and u.ex_add; U from
% exists_isAffineOpen_mem_and_subset; existsUnique_linear_lift_of_separating with S the points in U and f s the valAt' of s,
% with hsep from ext_evalPt and valAt'_eq_evalPt, u.hol and the IsHolTo hypothesis, and hD from tangent_mem_iff, isDiff_iff
% (the jet criterion), jet_snd and the uniqueness h1. existsUnique_linear_lift_of_separating: tendsto_map_nhds_of_separating
% (Lemma cv-origin-hodge-norm-form-lift-open), N from isDiscrete_iff_forall_mem_exists_isOpen and exists_nhds_half_neg,
% Function.invFunOn for sigma, hc, existsUnique_realLinear_lift_of_germ (Lemma cv-origin-hodge-norm-form-lift-real), hI
% (HasFDerivAt.comp with L.hasFDerivAt and HasFDerivAt.unique) and realLinear_eq_of_sub_mem_discrete. The hypothesis that
% A is abelian, and the trivializing functions, are not used. The published counterpart is the lifting of holomorphic homomorphisms of
% complex tori to their universal covers (Birkenhake-Lange, Chapter 1); not cited, as the locator is not checked. Size:
% about 360 lines with Lemmas cv-origin-hodge-norm-form-lift-open and -lift-real.

\begin{lemma}\label{cv-origin-hodge-norm-form-pullback-trivial}
Let $\Lambda\subset\mathbf C^g$ be a discrete subgroup that spans $\mathbf C^g$ over $\mathbf R$, $B\in M_g(\mathbf C)$ a Hermitian
matrix, $c:\Lambda\to\mathbf C$ a map with $|c|=1$, and $r$ a holomorphic function on $\mathbf C^g$ without zeros such that
\[
r(x+l)=c(l)\exp\big(\pi x^{\mathsf T}B\bar l+\tfrac\pi2\,l^{\mathsf T}B\bar l\big)\,r(x)
\]
for every $l\in\Lambda$ and $x\in\mathbf C^g$. Then $B=0$, $c(l)=1$ for every $l\in\Lambda$, and $r$ is constant.
\end{lemma}
\leanref{Development/AbelianVarietyTheory/Proofs/CubicalLineBundles/AppellHumbertTrivialOfNowhereVanishing.lean}{GlobalTheory.appellHumbert_trivial_of_nowhere_vanishing}
\begin{proof}
Put $q(x)=\operatorname{Re}x^{\mathsf T}B\bar x$. As $B$ is Hermitian, $y^{\mathsf T}B\bar x=\overline{x^{\mathsf T}B\bar y}$, so
\[
q(x+y)=q(x)+2\operatorname{Re}x^{\mathsf T}B\bar y+q(y)\qquad(x,y\in\mathbf C^g),\tag{$*$}
\]
and $q(tx)=|t|^2q(x)$ for $t\in\mathbf C$.

\emph{Bounds.} Let $\phi(x)=|r(x)|\exp(-\pi q(x)/2)$, a continuous positive function. By the functional equation,
$|r(x+l)|=\exp(\pi\operatorname{Re}x^{\mathsf T}B\bar l+\frac\pi2q(l))\,|r(x)|$, so $(*)$ gives $\phi(x+l)=\phi(x)$ for
$l\in\Lambda$. As $\Lambda$ is discrete and spans $\mathbf C^g$ over $\mathbf R$, every $x$ is $l+y$ with $l\in\Lambda$ and $y$ in
a compact parallelepiped, so $\phi$ and $1/\phi$ are bounded, by $M$ and $M'$. Hence $|r(x)|\le M\exp(\pi q(x)/2)$ and
$|r(x)|^{-1}\le M'\exp(-\pi q(x)/2)$ for every $x$.

\emph{The function $q$ vanishes.} Let $v\in\mathbf C^g$. Suppose that $q(v)<0$. The holomorphic function $t\mapsto r(tv)$ on
$\mathbf C$ satisfies $|r(tv)|\le M\exp(\pi|t|^2q(v)/2)$, so it is bounded, hence constant by Liouville's theorem, and it tends
to $0$ as $|t|\to\infty$; so $r(0)=0$, a contradiction. If $q(v)>0$, the same argument, applied to $t\mapsto r(tv)^{-1}$ with
the second bound, gives $r(0)^{-1}=0$, a contradiction. So $q(v)=0$.

\emph{The matrix $B$ vanishes.} By $(*)$ and $q=0$, $\operatorname{Re}x^{\mathsf T}B\bar y=0$ for all $x$ and $y$. Applied to $x$
and $iy$, as $x^{\mathsf T}B\overline{iy}=-i\,x^{\mathsf T}B\bar y$, it gives $\operatorname{Im}x^{\mathsf T}B\bar y=0$. So
$x^{\mathsf T}B\bar y=0$ for all $x$ and $y$, and $B=0$, by taking for $x$ and $y$ vectors of the standard basis.

\emph{The function $r$ and the map $c$.} Now $|r|\le M$, so $r$ is constant by Liouville's theorem. The functional equation at
$x=0$ gives $r(0)=r(l)=c(l)r(0)$, and $c(l)=1$ as $r(0)\neq0$.
\end{proof}
% Lean: appellHumbert_trivial_of_nowhere_vanishing: hnorm and hphiper (with re_hermForm_add_self), exists_bound_of_periodic
% (IsZLattice.isCompact_range_of_periodic) for phi and its inverse, hup and hlow, hscale, hq0 by trichotomy with
% eq_zero_of_gaussian_decay (Differentiable.apply_eq_apply_of_bounded and the limit along t = n), hH from
% re_hermForm_add_self at (x, y) and (x, I y), hB0 at the basis vectors Pi.single, hconst, and hper at 0. Size: about 200
% lines.

\begin{lemma}\label{cv-origin-hodge-norm-form-pullback}
Let $A$ and $A'$ be group schemes over $\mathbf C$ with line bundles $\mathcal L$ and $\mathcal L'$, each trivialized on an open
cover, and let $u:\mathbf C^g\to A(\mathbf C)$ and $u':\mathbf C^{g'}\to A'(\mathbf C)$, with their trivializing functions, be
uniformizations of
$A$ for $\mathcal L$ and of $A'$ for $\mathcal L'$ (Definition~\ref{cv-origin-hodge-norm-gram}(1), with
$T=\operatorname{Spec}\mathbf C$ and the identity), with kernels $\Lambda$ and $\Lambda'$ and Appell--Humbert data $(H,\chi)$ and
$(H',\chi')$. Let $f:A\to A'$ be a morphism over $\mathbf C$, $F:\mathbf C^g\to\mathbf C^{g'}$ a $\mathbf C$-linear map with
$F(\Lambda)\subset\Lambda'$ and $u'(F(x))=f(u(x))$ for every $x$, and $n\ge1$ with $\mathcal L^{\otimes n}\cong
f^*\mathcal L'^{\otimes n}$. Then $x^{\mathsf T}H\bar y=F(x)^{\mathsf T}H'\overline{F(y)}$ for all $x,y\in\mathbf C^g$, and
$\chi(l)^n=\chi'(F(l))^n$ for every $l\in\Lambda$.
\end{lemma}
\leanref{Development/AbelianVarietyTheory/Proofs/ModuliProofs/AppellHumbertDataEqOfLbIsoPowPull.lean}{GlobalTheory.appellHumbert_data_eq_of_lbIso_pow_pull}
\leanref{Development/AbelianVarietyTheory/Proofs/CubicalLineBundles/HermFormPullEqOfLbIsoPow.lean}{GlobalTheory.hermForm_pull_eq_of_lbIso_pow}
\leanref{Development/AbelianVarietyTheory/Proofs/ModuliProofs/AppellHumbertDataEqOfLbIsoPowPull.lean}{GlobalTheory.semicharacter_pull_eq_of_lbIso_pow}
\begin{proof}
Let $(U_i)$ and $(U'_j)$ be the open covers on which $\mathcal L$ and $\mathcal L'$ are trivialized, with transition functions
$g_{ii'}$ and $g'_{jj'}$, and let $(\theta_i)$ and $(\theta'_j)$ be the trivializing functions of $u$ and $u'$. In these
trivializations, the
isomorphism of $\mathcal L^{\otimes n}$ with $f^*\mathcal L'^{\otimes n}$ is given by regular functions $w_{ij}$ without zeros
on $U_i\cap f^{-1}(U'_j)$ with $w_{ij}\,(g'_{jj'}\circ f)^n=g_{ii'}^n\,w_{i'j'}$ where all are defined.

\emph{A function without zeros.} For each pair $(i,j)$ let $V_{ij}$ be the set of $x\in\mathbf C^g$ with $u(x)\in U_i$ and
$u'(F(x))\in U'_j$, an open set, and put
\[
r_{ij}(x)=w_{ij}(u(x))\,\theta'_j(F(x))^n\,\theta_i(x)^{-n}\qquad(x\in V_{ij}).
\]
As $u'(F(x))=f(u(x))$, the point $u(x)$ lies in $U_i\cap f^{-1}(U'_j)$ for $x\in V_{ij}$, so $r_{ij}$ is holomorphic on $V_{ij}$
and has no zeros there. On $V_{ij}\cap V_{i'j'}$, $\theta_i=(g_{ii'}\circ u)\,\theta_{i'}$ and
$\theta'_j(F(x))=g'_{jj'}(f(u(x)))\,\theta'_{j'}(F(x))$, so the relation between $w_{ij}$ and $w_{i'j'}$ gives $r_{ij}=r_{i'j'}$.
Every point of $A$ and of $A'$ lies in some chart, so the $V_{ij}$ cover $\mathbf C^g$, and the $r_{ij}$ glue to a holomorphic
function $r$ on $\mathbf C^g$ without zeros.

\emph{Its factor of automorphy.} Let $P$ be the matrix of $F$, so that $F(x)^{\mathsf T}H'\overline{F(y)}=x^{\mathsf T}H_2\bar y$
for the Hermitian matrix $H_2=P^{\mathsf T}H'\bar P$. Put $B=n(H-H_2)$, a Hermitian matrix, and
$c(l)=\chi(l)^n\chi'(F(l))^{-n}$ for $l\in\Lambda$, so that $|c|=1$. Write $e_l(x)=\exp(\pi x^{\mathsf T}H\bar l+\frac\pi2l^{\mathsf T}H\bar l)$,
and $e'_{l'}(x')$ for the same expression with $H'$. Let $l\in\Lambda$ and $x\in V_{ij}$. As $u(x+l)=u(x)$ and
$u'(F(x+l))=u'(F(x)+F(l))=u'(F(x))$, since $F(l)\in\Lambda'$, the point $x+l$ lies in $V_{ij}$. The functional equations of the
trivializing functions give $\theta_i(x+l)=\theta_i(x)/(\chi(l)e_l(x))$ and $\theta'_j(F(x)+F(l))=\theta'_j(F(x))/(\chi'(F(l))e'_{F(l)}(F(x)))$. So
\[
r(x+l)=c(l)\Big(\frac{e_l(x)}{e'_{F(l)}(F(x))}\Big)^n r(x)=c(l)\exp\big(\pi x^{\mathsf T}B\bar l+\tfrac\pi2\,l^{\mathsf T}B\bar l\big)\,r(x).
\]

\emph{Conclusion.} Lemma~\ref{cv-origin-hodge-norm-form-pullback-trivial}, applied to $\Lambda$, $B$, $c$ and $r$, gives $B=0$.
As $n\ge1$, $H=H_2$, that is, $x^{\mathsf T}H\bar y=F(x)^{\mathsf T}H'\overline{F(y)}$ for all $x,y\in\mathbf C^g$. Then
$r(x+l)=c(l)r(x)$ for $l\in\Lambda$, and the same lemma, applied to $\Lambda$, the zero matrix, $c$ and $r$, gives $c(l)=1$, that
is, $\chi(l)^n=\chi'(F(l))^n$.
\end{proof}
% Lean: appellHumbert_data_eq_of_lbIso_pow_pull: hermForm_pull_eq_of_lbIso_pow with m = n and the symmetric LBIso, then
% mul_left_cancel by n; semicharacter_pull_eq_of_lbIso_pow for chi. hermForm_pull_eq_of_lbIso_pow: hFl, hLambda, hper,
% exists_glue for r with the agreement from PullHermitian.val_iso_pow (the relation of the w_ij, from unitVal_iso and
% unitVal_linComb), u.fr_compat and u'.fr_compat; H2 = P^T Hm' conj(P) (hermForm_pull) with B = n Hm - m H2 (hermForm_sub_smul),
% and appellHumbert_trivial_of_nowhere_vanishing (Lemma cv-origin-hodge-norm-form-pullback-trivial). The second application,
% with B = 0, is semicharacter_pull_eq_of_lbIso_pow, which glues the same r and uses the agreement of the Hermitian forms.
% Only the Hermitian forms are used below. Size: about 440 lines with Lemma cv-origin-hodge-norm-form-pullback-trivial.

\begin{lemma}\label{cv-origin-hodge-norm-form-compare}
Let $T$, $G$, $\mathcal L$ and $z$ be as in Definition~\ref{cv-origin-hodge-norm-gram}, and let $(A,\varphi,u)$ and
$(A',\varphi',u')$, with their trivializing functions, be uniformizations of the fibre of $G$ at $z$ of dimensions $g$ and
$g'$, with
Appell--Humbert data $(H,\chi)$ and $(H',\chi')$. Suppose that $A$ is an abelian scheme over $\operatorname{Spec}\mathbf C$, and
that every tangent vector of $A$ at the origin has exactly one differential for $u$ and every $a\in\mathbf C^g$ is the
differential of exactly one tangent vector of $A$ at the origin. Let $x$ and $y$ be tangent vectors of $A$ at the origin with
differentials $a$ and $b$, and let $x',y':\operatorname{Spec}\mathbf C[\varepsilon]\to A'$ be morphisms over $\mathbf C$ with
$\varphi'\circ x'=\varphi\circ x$ and $\varphi'\circ y'=\varphi\circ y$, with differentials $a'$ and $b'$. Then
$a'^{\mathsf T}H'\bar b'=a^{\mathsf T}H\bar b$.
\end{lemma}
\leanref{Development/AutomorphicSpectrum/Proofs/GeometryProofs/HolUnifCompareIsogenyStatement.lean}{GlobalTheory.T.hol_unif_compare_isogeny}
\begin{proof}
\emph{The isomorphism of the fibres.} Both $\varphi$ and $\varphi'$ identify their sources with $G\times_T\operatorname{Spec}\mathbf C$,
so there is an isomorphism $e:A'\to A$ over $\mathbf C$ with $\varphi\circ e=\varphi'$; it is a homomorphism, as $\varphi$ and
$\varphi'$ are. A morphism to $A$ over $\mathbf C$ is determined by its composite with $\varphi$.

\emph{The linear lift.} The map $w\mapsto e(u'(w))$ is a homomorphism $\mathbf C^{g'}\to A(\mathbf C)$, and for an open
$U\subset A$ and $f$ regular on $U$, the set of $w$ with $e(u'(w))\in U$ is $u'^{-1}(e^{-1}(U))$, which is open, and
$f\circ e\circ u'$ is holomorphic on it, as $f\circ e$ is regular on $e^{-1}(U)$. By Lemma~\ref{cv-origin-hodge-norm-form-lift}
there is a $\mathbf C$-linear $F:\mathbf C^{g'}\to\mathbf C^g$ with $u\circ F=e\circ u'$. For $l\in\Lambda'$, $u'(l)$ is the
origin of $A'$, so $u(F(l))=e(u'(l))$ is the origin of $A$, and $F(\Lambda')\subset\Lambda$.

\emph{The Hermitian forms.} The line bundle of $u'$ is $\varphi'^*\mathcal L\cong e^*\varphi^*\mathcal L$. So
Lemma~\ref{cv-origin-hodge-norm-form-pullback}, applied to $e:A'\to A$, $F$ and $n=1$, gives
$c^{\mathsf T}H'\bar d=F(c)^{\mathsf T}H\overline{F(d)}$ for all $c,d\in\mathbf C^{g'}$.

\emph{The differentials.} The morphism $e\circ x'$ satisfies $\varphi\circ e\circ x'=\varphi'\circ x'=\varphi\circ x$, so
$e\circ x'=x$. Let $U\subset A$ be an open set containing the origin and $f$ regular on $U$. As $x'$ has differential $a'$
and $f\circ e$ is regular on $e^{-1}(U)$, the coefficient of $\varepsilon$ in $f\circ x=(f\circ e)\circ x'$ is the derivative
of $f\circ e\circ u'=f\circ u\circ F$ at $0$ in the direction $a'$, which is $d(f\circ u)_0(F(a'))$ by the chain rule. So $x$ has
differential $F(a')$, and $F(a')=a$ by the uniqueness of differentials. In the same way $F(b')=b$. Hence
$a'^{\mathsf T}H'\bar b'=F(a')^{\mathsf T}H\overline{F(b')}=a^{\mathsf T}H\bar b$.
\end{proof}
% Lean: T.hol_unif_compare_isogeny = HodgeUniq.hermForm_eq_of_pullbacks: grpPullback_iso (e with IsMonHom and
% e.hom.left >> phi = phi'), HolTorsion.isHolTo_comp, existsUnique_linear_lift_of_isHolTo (Lemma
% cv-origin-hodge-norm-form-lift), hLambda, GenSpread.lbIso_pull_pull with lbIso_lbPow, appellHumbert_data_eq_of_lbIso_pow_pull
% with n = 1 (Lemma cv-origin-hodge-norm-form-pullback), and coord: the cartesian square gives x' >> e = x, and isDiff_comp
% (the tangent vector with differential F a', identified with x' >> e by VerticalDeriv.isDerivAt0_unique from
% isDerivAt0_comp_linMap and FibreDeriv.isDerivAt0_comp, which is the chain rule of the text), with the uniqueness d1.
% Size: about 1,250 lines with Lemmas cv-origin-hodge-norm-form-lift and cv-origin-hodge-norm-form-pullback.

\begin{lemma}\label{cv-origin-hodge-norm-form}
In the setting of Definition~\ref{cv-origin-hodge-component}, let $Z$ be a Hodge component of $\mathcal X$ at $S$, let $z$ be a
complex point of $Z$ over which $\mathcal A_Z$ is abelian, and let $T_z$ be the Lie algebra of $\mathcal A_Z$ at $z$. Then
there is a map $h:T_z\times T_z\to\mathbf C$, conjugate-linear in the first variable and linear in the second, with
$h(w,v)=\overline{h(v,w)}$ for all $v,w$ and $h(v,v)>0$ for $v\neq0$, such that tangent vectors $v,w\in T_z$ have Hodge
inner product $c$ for the line bundle of $\mathcal A_Z$ (Definition~\ref{cv-origin-hodge-norm-gram}(2)) if and only if
$c=h(v,w)$.
\end{lemma}
\leanref{Development/AutomorphicSpectrum/Proofs/GeometryProofs/HodgeHermitianFormExistsStatement.lean}{GlobalTheory.T.hodge_hermitian_form_exists}
\leanref{Development/AutomorphicSpectrum/Proofs/GeometryProofs/HodgeCptInteriorFlag.lean}{GlobalTheory.CompanionTheory.Ar.GrNorm.interior_pol}
\begin{proof}
Let $\mathcal L$ be the line bundle of $\mathcal A_Z$, trivialized on an open cover $(V_\beta)$ with transition functions
$g'_{\beta\beta'}$, and let $A$, its line bundle $L_A$ and $\varphi:A\to\mathcal A_Z$ be the fibre of $\mathcal A_Z$ at $z$ as in
Lemma~\ref{cv-origin-hodge-norm-form-fibre}. Trivialize $L_A$ on an open cover $(U_\alpha)$ of $A$ with transition functions
$g_{\alpha\alpha'}$.

\emph{The fibre.} By Lemma~\ref{cv-origin-hodge-norm-form-fibre}, $A$ is proper over $\mathbf C$ and $L_A$ is relatively ample
by sections. As $\mathcal A_Z$ is semi-abelian, $A$ is commutative, smooth and has connected fibre; being proper, it is an
abelian scheme over $\operatorname{Spec}\mathbf C$. As $\operatorname{Spec}\mathbf C$ is connected, $A$ is smooth of some
relative dimension $g$.

\emph{A uniformization of $A$.} Lemma~\ref{cv-origin-hodge-norm-form-torus} gives a uniformization $u:\mathbf C^g\to A(\mathbf C)$ of
$A$ for the trivial line bundle, with kernel $\Lambda$, and an affine open $U_0$ with coordinates $t_1,\dots,t_g$ satisfying
its (a) and (b). Lemma~\ref{cv-origin-hodge-norm-form-appell}, applied to $L_A$ with its trivialization at the origin and its
cubical structure, gives trivializing functions $\theta_\alpha$ that make $u$ a uniformization of $A$ for $L_A$, a Hermitian
$H$ and $\chi$
with $|\chi|=1$ such that $\theta_\alpha(x)=\chi(l)\exp(\pi x^{\mathsf T}H\bar l+\frac\pi2l^{\mathsf T}H\bar l)\theta_\alpha(x+l)$.
Lemma~\ref{cv-origin-hodge-norm-form-theta} gives $N\ge1$ and holomorphic functions $\Theta_j$ with its properties (1)--(3), and
then Lemma~\ref{cv-origin-hodge-norm-form-positive} shows that $x^{\mathsf T}H\bar x>0$ for $x\neq0$. By
Lemma~\ref{cv-origin-hodge-norm-form-tangent}, every tangent vector of $A$ at the origin has exactly one differential for $u$,
and every $a\in\mathbf C^g$ is the differential of exactly one tangent vector of $A$ at the origin.

\emph{A uniformization of the fibre of $\mathcal A_Z$.} As $\varphi$ is a morphism of $\mathcal O_E$-semi-abelian schemes,
there is an isomorphism $\iota:L_A\to\varphi^*\mathcal L$. In the trivializations, $\iota$ is multiplication by regular functions
$\iota_{\beta\alpha}$ without zeros on $U_\alpha\cap\varphi^{-1}(V_\beta)$, with
$\iota_{\beta\alpha}=(g'_{\beta\beta'}\circ\varphi)\,\iota_{\beta'\alpha'}\,g_{\alpha'\alpha}$ where all are defined. For each
$\beta$ and $x\in u^{-1}(\varphi^{-1}(V_\beta))$ put $\theta'_\beta(x)=\iota_{\beta\alpha}(u(x))\theta_\alpha(x)$ for any $\alpha$
with $u(x)\in U_\alpha$. By the relation between the $\iota_{\beta\alpha}$ and $\theta_\alpha=(g_{\alpha\alpha'}\circ u)\theta_{\alpha'}$,
this does not depend on $\alpha$, and $\theta'_\beta=(g'_{\beta\beta'}\circ\varphi\circ u)\theta'_{\beta'}$. Near each point,
$\theta'_\beta$ is a product of holomorphic functions without zeros. As $u(x+l)=u(x)$ for $l\in\Lambda$, the functional
equation of the $\theta_\alpha$ gives
$\theta'_\beta(x)=\chi(l)\exp(\pi x^{\mathsf T}H\bar l+\frac\pi2l^{\mathsf T}H\bar l)\theta'_\beta(x+l)$. So
$(A,\varphi,u,(\theta'_\beta))$ is a uniformization of the fibre of $\mathcal A_Z$ at $z$ for $\mathcal L$
(Definition~\ref{cv-origin-hodge-norm-gram}(1)), with Appell--Humbert data $(H,\chi)$.

\emph{The form.} Lemma~\ref{cv-origin-hodge-norm-form-lie}, applied to $\varphi$ and this uniformization, gives a
$\mathbf C$-linear isomorphism $D:T_z\to\mathbf C^g$ and, for each $v\in T_z$, a tangent vector $x_v$ of $A$ at the origin
with $\varphi\circ x_v=v$ and differential $D(v)$. Put
\[
h(v,w)=\pi\,D(w)^{\mathsf T}H\,\overline{D(v)}.
\]
It is conjugate-linear in $v$ and linear in $w$. As $H$ is Hermitian, $\overline{b^{\mathsf T}H\bar a}=a^{\mathsf T}H\bar b$ for
$a,b\in\mathbf C^g$, so $h(w,v)=\overline{h(v,w)}$. For $v\neq0$, $D(v)\neq0$ and $h(v,v)=\pi D(v)^{\mathsf T}H\overline{D(v)}>0$.

If $c=h(v,w)$, then the uniformization $(A,\varphi,u,(\theta'_\beta))$ with Appell--Humbert data $(H,\chi)$, and the morphisms
$x_v$ and $x_w$ with differentials $D(v)$ and $D(w)$, show that $v$ and $w$ have Hodge inner product $c$
(Definition~\ref{cv-origin-hodge-norm-gram}(2)). Conversely, suppose that $v$ and $w$ have Hodge inner product $c$, through a
uniformization $(A',\varphi',u')$ with Appell--Humbert data $(H',\chi')$ and morphisms $v',w'$ with $\varphi'\circ v'=v$ and
$\varphi'\circ w'=w$, with differentials $a'$ and $b'$, so that $c=\pi\,b'^{\mathsf T}H'\bar a'$. Lemma~\ref{cv-origin-hodge-norm-form-compare},
applied to the two uniformizations, to $x=x_w$ and $y=x_v$ with differentials $D(w)$ and $D(v)$, and to $x'=w'$ and $y'=v'$,
gives $b'^{\mathsf T}H'\bar a'=D(w)^{\mathsf T}H\overline{D(v)}$. So $c=h(v,w)$.
\end{proof}
% Lean: T.hodge_hermitian_form_exists = exists_hodgeInnerProduct_form, for any SAB t over X and z, under the hypothesis that
% the fibre at z (an SAB t' over Sp C with SABMap t' t z) is proper with RelAmple line bundle, which GrNorm.interior_pol
% (Lemma cv-origin-hodge-norm-form-fibre) gives at the points of the abelian locus: isAbelianScheme_of_proper,
% AbelianDim.exists_relDim, T.exists_exp_chart_coordinates (Lemma cv-origin-hodge-norm-form-torus),
% T.theta_multiplier_normal_form (Lemma cv-origin-hodge-norm-form-appell), T.exists_theta_basis_of_factor and
% T.theta_functions_finite_dim (Lemmas cv-origin-hodge-norm-form-theta and -positive), TangentOrigin.tangent_bijection_of_values
% (Lemma cv-origin-hodge-norm-form-tangent), ArchMetricCompare.carry and carry_isAH (the trivializing functions theta'_beta, through
% exists_frame and val_iso, from the LBIso w0 of m.lin), T.tangent_space_equiv_hol_unif (Lemma cv-origin-hodge-norm-form-lie)
% and T.hol_unif_compare_isogeny (Lemma cv-origin-hodge-norm-form-compare); h v w = pi * hermForm Hm (D w) (D v), Hermitian by
% conj_hermForm. Size of the top level with carry: about 300 lines. Nothing below this lemma remains open: Lemmas -torus
% (with -torus-exp-log, -torus-exp, -torus-points, -torus-chart-etale, -torus-chart and -torus-divisible, citing the
% finiteness and flatness of multiplication by r), -descent, -descent-ample (with -descent-ample-affine and
% -descent-ample-embedding: relAmple_of_isAmpleLB and IsAmpleLB.isAffineOpen_nonzeroLocus, about 1,600 lines), -appell,
% -lift and -pullback are written out.

\begin{lemma}\label{cv-origin-hodge-norm-adapted}
Let $\Omega$ be a field, $A\subset\Omega$ a subring and $\varphi:A\to\mathbf C$ a ring homomorphism such that every $a\in A$
with $\varphi(a)\neq0$ is a unit of $A$. Let $W$ be an $\Omega$-vector space, $L\subset W$ an $A$-submodule, $V$ a complex
vector space, $\mathrm{val}:L\to V$ an additive map with $\mathrm{val}(ae)=\varphi(a)\mathrm{val}(e)$ for $a\in A$, and
$h:V\times V\to\mathbf C$ a map, conjugate-linear in the first variable and linear in the second, with
$h(w,v)=\overline{h(v,w)}$ and $h(v,v)>0$ for $v\neq0$. For a family $c=(c_1,\dots,c_r)$ in $L$ let
$\Gamma(c)=(h(\mathrm{val}(c_p),\mathrm{val}(c_q)))_{p,q}$, with upper left $l\times l$ blocks $\Gamma(c)_l$, and for a family
$c$ in $W$ let $\langle c\rangle_l$ be the $\Omega$-span of $c_1,\dots,c_l$. Let $Q\subset L$ be an $A$-submodule and
$b^0=(b^0_1,\dots,b^0_r)$ a linearly independent family in $W$. Call a family $c$ in $W$ an \emph{adapted basis} of $Q$ if it
is linearly independent over $\Omega$, $Q=\sum_kAc_k$ and $\langle c\rangle_l=\langle b^0\rangle_l$ for every $l$, and
suppose that the $\mathrm{val}(c_k)$ are linearly independent over $\mathbf C$ for every adapted basis $c$. Let $b$ be an
adapted basis, $1\le m\le r$ and $t\in\Omega^\times$ with $b_m-tb^0_m\in\langle b^0\rangle_{m-1}$. Then there is a real
$\nu>0$ such that for $t'\in\Omega$ and $\nu'\in\mathbf R$ the following are equivalent:
\begin{enumerate}
\item[(a)] there is an adapted basis $b'$ of $Q$ with $b'_m-t'b^0_m\in\langle b^0\rangle_{m-1}$, $\Gamma(b')$ positive
  definite and $\nu'=(\det\Gamma(b')_m/\det\Gamma(b')_{m-1})^{1/2}$;
\item[(b)] $t'/t\in A$, $\varphi(t'/t)\neq0$ and $\nu'=|\varphi(t'/t)|\nu$.
\end{enumerate}
\end{lemma}
\leanref{Development/AutomorphicSpectrum/Proofs/GeometryProofs/HodgeCptInteriorFlag.lean}{GlobalTheory.CompanionTheory.Ar.GrNorm.core}
\begin{proof}
\emph{Gram matrices.} Let $c$ be an adapted basis. For $x\in\mathbf C^r$,
$\bar x^{\mathsf T}\Gamma(c)x=h\big(\sum_px_p\mathrm{val}(c_p),\sum_qx_q\mathrm{val}(c_q)\big)$, which is positive for
$x\neq0$, as the $\mathrm{val}(c_k)$ are linearly independent. So $\Gamma(c)$ is positive definite, and so are its blocks
$\Gamma(c)_l$, whose determinants are positive; $\det\Gamma(c)_0=1$. Put
$\nu=(\det\Gamma(b)_m/\det\Gamma(b)_{m-1})^{1/2}>0$.

\emph{Change of basis.} Let $b'$ be an adapted basis with $b'_m-t'b^0_m\in\langle b^0\rangle_{m-1}$. As $b'_k\in Q=\sum_lAb_l$,
$b'_k=\sum_lP_{lk}b_l$ with $P_{lk}\in A$; as $b'_k\in\langle b'\rangle_k=\langle b^0\rangle_k=\langle b\rangle_k$ and $b$ is
linearly independent, $P_{lk}=0$ for $l>k$. In the same way $b_k=\sum_lR_{lk}b'_l$ with $R_{lk}\in A$ and $R_{lk}=0$ for
$l>k$. Then $PR=1$, as $b$ is linearly independent, and its diagonal gives $P_{kk}R_{kk}=1$: the $P_{kk}$ are units of $A$.
Let $C=(\varphi(P_{lk}))$. As $\mathrm{val}$ is additive and $\mathrm{val}(ae)=\varphi(a)\mathrm{val}(e)$,
$\mathrm{val}(b'_k)=\sum_lC_{lk}\mathrm{val}(b_l)$, so $\Gamma(b')=\bar C^{\mathsf T}\Gamma(b)C$. As $C$ is upper triangular,
the block $\Gamma(b')_l$ is $\bar C_l^{\mathsf T}\Gamma(b)_lC_l$ for the upper left block $C_l$ of $C$, and
$\det\Gamma(b')_l=\prod_{k\le l}|C_{kk}|^2\det\Gamma(b)_l$. Hence
$(\det\Gamma(b')_m/\det\Gamma(b')_{m-1})^{1/2}=|\varphi(P_{mm})|\nu$. Finally
$b'_m-P_{mm}b_m=\sum_{l<m}P_{lm}b_l\in\langle b\rangle_{m-1}=\langle b^0\rangle_{m-1}$, so
\[
(t'-P_{mm}t)b^0_m=(b'_m-P_{mm}b_m)-(b'_m-t'b^0_m)+P_{mm}(b_m-tb^0_m)\in\langle b^0\rangle_{m-1},
\]
and $t'=P_{mm}t$, as $b^0$ is linearly independent.

\emph{The equivalence.} If (a) holds with $b'$, the previous paragraph gives $t'/t=P_{mm}$, a unit of $A$, so
$\varphi(t'/t)\neq0$, and $\nu'=|\varphi(t'/t)|\nu$. Conversely, suppose (b), and put $a=t'/t$. As $\varphi(a)\neq0$, $a$ is a
unit of $A$. Let $b'$ be $b$ with $b_m$ replaced by $ab_m$. It is linearly independent, it spans $Q$ over $A$ and has the
same spans $\langle b'\rangle_l=\langle b\rangle_l$, so it is an adapted basis, and
$b'_m-t'b^0_m=a(b_m-tb^0_m)\in\langle b^0\rangle_{m-1}$. Its Gram matrix is positive definite by the first paragraph, and the
previous paragraph gives $P_{mm}t=t'$, so $P_{mm}=a$ and $(\det\Gamma(b')_m/\det\Gamma(b')_{m-1})^{1/2}=|\varphi(a)|\nu=\nu'$.
So (a) holds with $b'$.
\end{proof}
% Lean: GrNorm.core, with the Hermitian form h : V ->l* V ->l C (hh, hp) and hfr for the adapted bases: gram_posDef and
% leadMinor_pos_of_posDef for nu > 0; compare, through frame_change (exists_tri both ways and the product of the two
% triangular matrices), val_sum, the conjugated Gram matrix hG, leadMinor_conj and sqrt_ratio_conj, and the coefficient
% argument m1, m2 with exists_coeff_of_mem_flagStep; scale_frame with Function.update 1 i w for the converse. leadMinor is
% the real part of the determinant, which is the determinant for a positive definite matrix. Size: about 450 lines.

\begin{lemma}\label{cv-origin-analytic-hodge-points}
In the setting of Definition~\ref{cv-origin-hodge-component}, let $Z$ be a Hodge component of $\mathcal X$ at $S$ with
coordinates $y$ with values in $\Omega$, and let $z$ be a complex point of $Z$ with centre $P$ and coordinates $w$
(Definition~\ref{cv-origin-reference}(4)).
\begin{enumerate}
\item If $z'$ is a complex point of $Z$ with coordinates $w$, then $z'=z$.
\item For $\kappa\in(\mathbf C^\times)^K$, the point $\kappa w=(\kappa_kw_{k,i})$ gives coordinates of $z$.
\item Let $\bar z=z\circ\operatorname{Spec}(\bar{\ })$, for complex conjugation $x\mapsto\bar x$. Then $\bar z$ has centre $P$;
  an element of $\Omega$ regular at $z$ with value $u$ is regular at $\bar z$ with value $\bar u$; $\bar w$ are coordinates
  of $\bar z$; and if a tangent vector $v$ of $\mathcal A_Z$ at $z$ is the value at $z$ of $e\in\mathrm{Lie}$
  (Definition~\ref{cv-origin-hodge-norm-gram}(3)), then $\bar v=v\circ\operatorname{Spec}(\bar{\ }[\varepsilon])$ is the
  value of $e$ at $\bar z$, where $\bar{\ }[\varepsilon]$ is complex conjugation with $\varepsilon\mapsto\varepsilon$.
\end{enumerate}
\end{lemma}
\leanref{Development/AutomorphicSpectrum/Proofs/GeometryProofs/LocalRingPointsExistsUnique.lean}{GlobalTheory.CompanionTheory.Ar.TiePoints.existsUnique_point}
\leanref{Development/AutomorphicSpectrum/Proofs/GeometryProofs/LocalRingPointsExistsUnique.lean}{GlobalTheory.CompanionTheory.Ar.TiePoints.coords_iff}
\leanref{Development/AutomorphicSpectrum/Proofs/GeometryProofs/HodgeCptFlagXdeg.lean}{GlobalTheory.CompanionTheory.Ar.HodgeLine.centre_cj}
\leanref{Development/AutomorphicSpectrum/Proofs/GeometryProofs/HodgeCptFlagXdeg.lean}{GlobalTheory.CompanionTheory.Ar.HodgeLine.valAt_conj}
\leanref{Development/AutomorphicSpectrum/Proofs/GeometryProofs/HodgeCptFlagXdeg.lean}{GlobalTheory.CompanionTheory.Ar.HodgeLine.coordsAt_conj}
\leanref{Development/AutomorphicSpectrum/Proofs/GeometryProofs/HodgeCptFlagXdeg.lean}{GlobalTheory.CompanionTheory.Ar.HodgeLine.lieVal_conj}
\begin{proof}
(1) Choose $c_k\in I_k$ with $w_{k,c_k}\neq0$ for every $k$. By the coordinates, $y_{k,c_k}\neq0$ and the ratios
$y_{k,i}/y_{k,c_k}$ are regular at the centres $P$ of $z$ and $P'$ of $z'$, with the value $w_{k,i}/w_{k,c_k}$ at both points.
So $P$ and $P'$ lie in the affine open subscheme $Z_c$ of the chart $c$, with ring
$A_c=\mathbf Z[y_{k,i}/y_{k,c_k}]$ (Definition~\ref{cv-origin-hodge-component}(1)), and the homomorphisms $A_c\to\mathbf C$
given by $z$ and $z'$ agree on the generators, hence everywhere. As the homomorphisms from the local rings to $\mathbf C$ are
local, an element of $A_c$ is a unit at $P$, or at $P'$, if and only if its value is nonzero. So $P=P'$ by
Lemma~\ref{cv-origin-snc-charts}(2), applied to the prime of the elements of $A_c$ with value $0$. By
Lemma~\ref{cv-origin-snc-charts}(1), every element regular at $P$ is $r/s$ with $r,s\in A_c$ and $s$ a unit at $P$, so the
two homomorphisms from the local ring at $P$ to $\mathbf C$ agree. A complex point is determined by its centre and this
homomorphism \cite[Tag 01J6]{Sta}, so $z'=z$.

(2) The conditions of Definition~\ref{cv-origin-reference}(4) involve $w$ only through the entries that are nonzero and the
ratios $w_{k,i}/w_{k,c}$ of two entries of one block, which $\kappa$ does not change; and $\kappa w\in\hat C$.

(3) The point $z$ is given by $P$ and a local homomorphism $\mathcal O_{Z,P}\to\mathbf C$, and $\bar z$ by $P$ and its
composite with complex conjugation, which is again local \cite[Tag 01J6]{Sta}. So $\bar z$ has centre $P$, and the values at
$\bar z$ are the conjugates of the values at $z$. Conjugation preserves the nonzero entries and the ratios of the coordinates,
so $\bar w$ are coordinates of $\bar z$. If $e$ extends to $P$ through $v_P$ and $v=v_P\circ\operatorname{Spec}\psi[\varepsilon]$,
where $\psi$ is the homomorphism of $z$, then the homomorphism of $\bar z$ is $\bar\psi$, and
$v_P\circ\operatorname{Spec}\bar\psi[\varepsilon]=v_P\circ\operatorname{Spec}\psi[\varepsilon]\circ\operatorname{Spec}(\bar{\ }[\varepsilon])=\bar v$.
\end{proof}
% Lean: (1) is the uniqueness clause of TiePoints.existsUnique_point (the third part of
% T.local_ring_points_existsUnique), for R = C: the chart c with unit coordinates, T.prime_chart_existsUnique_point for the
% centre (hpp), T.regAt_iff_chart_localization for the local homomorphism, and T.existsUnique_centre_local_hom. (2) is the
% reverse implication of TiePoints.coords_iff. (3) is HodgeLine.centre_cj (through conj_transport and
% SpecToEquivOfLocalRing), valAt_conj, coordsAt_conj and lieVal_conj, with cj z = Spec(conj) >> z and cjTan. Size: about 500
% lines.

\begin{lemma}\label{cv-origin-analytic-hodge-conjugate}
\begin{enumerate}
\item Let $T$, $G$, $\mathcal L$ and $z$ be as in Definition~\ref{cv-origin-hodge-norm-gram}, let
  $\bar z=z\circ\operatorname{Spec}(\bar{\ })$, and for a tangent vector $v$ of $G$ at $z$ let
  $\bar v=v\circ\operatorname{Spec}(\bar{\ }[\varepsilon])$, a tangent vector of $G$ at $\bar z$. If $v$ and $w$ have Hodge
  inner product $c$, then $\bar v$ and $\bar w$ have Hodge inner product $\bar c$.
\item In the setting of Definition~\ref{cv-origin-hodge-component}, let $Z$ be a Hodge component of $\mathcal X$ at $S$, let a
  flag of $Z$ have bases $b_s$, let $s$ be a label and $1\le m\le r_s$, and let $z$ be a complex point of $Z$ over which
  $\mathcal A_Z$ is abelian. If $t'\in\Omega$ generates the $m$-th graded piece of the flag of $\mathcal V_s$ at $z$ with Hodge
  norm $\nu'$, in the sense of Lemma~\ref{cv-origin-model-metric-norm}, then it generates it at $\bar z$ with Hodge norm
  $\nu'$.
\end{enumerate}
\end{lemma}
\leanref{Development/AutomorphicSpectrum/Proofs/GeometryProofs/HodgeInnerProductConjugateStatement.lean}{GlobalTheory.T.hodge_inner_product_conjugate}
\leanref{Development/AutomorphicSpectrum/Proofs/GeometryProofs/HodgeCptFlagXdeg.lean}{GlobalTheory.CompanionTheory.Ar.HodgeLine.isGrNorm_conj}
\begin{proof}
(1) Let $A$, $\varphi$, $u$, $\Lambda$, $(\theta_\alpha)$, $(H,\chi)$, $v'$, $w'$, $a$ and $b$ be data for $v$, $w$ and $c$ as in
Definition~\ref{cv-origin-hodge-norm-gram}(2). Let $\bar A$ be the scheme $A$ over $\operatorname{Spec}\mathbf C$ through its
structure morphism followed by $\operatorname{Spec}(\bar{\ })$. The same homomorphism $\varphi$ identifies $\bar A$ with the
fibre of $G$ at $\bar z$. The complex points of $\bar A$ are the $s\circ\operatorname{Spec}(\bar{\ })$ for the complex points
$s$ of $A$, and $\bar u(x)=u(\bar x)\circ\operatorname{Spec}(\bar{\ })$ is a surjective homomorphism
$\mathbf C^g\to\bar A(\mathbf C)$ with kernel $\bar\Lambda$, which is discrete and spans $\mathbf C^g$ over $\mathbf R$. A
regular function takes at $\bar u(x)$ the conjugate of its value at $u(\bar x)$, so $f\circ\bar u$ is holomorphic where it is
defined, and so are the functions $\bar\theta_\alpha(x)=\overline{\theta_\alpha(\bar x)}$, without zeros and with
$\bar\theta_\alpha=(g_{\alpha\beta}\circ\varphi\circ\bar u)\bar\theta_\beta$. For $l\in\bar\Lambda$, the functional equation
of $\theta_\alpha$ at $\bar x$ and $\bar l\in\Lambda$, conjugated, is
\[
\bar\theta_\alpha(x)=\overline{\chi(\bar l)}\exp\big(\pi x^{\mathsf T}\bar H\bar l+\tfrac\pi2\,l^{\mathsf T}\bar H\bar l\big)\,
\bar\theta_\alpha(x+l),
\]
and $\bar H$ is Hermitian with $x^{\mathsf T}\bar H\bar x=\overline{\bar x^{\mathsf T}Hx}>0$ for $x\neq0$. So $\bar A$, $\varphi$,
$\bar u$ and the $\bar\theta_\alpha$ form a uniformization of the fibre at $\bar z$ with Appell--Humbert data
$(\bar H,l\mapsto\overline{\chi(\bar l)})$. The morphisms $\bar v'=v'\circ\operatorname{Spec}(\bar{\ }[\varepsilon])$ and
$\bar w'=w'\circ\operatorname{Spec}(\bar{\ }[\varepsilon])$ to $\bar A$ satisfy $\varphi\circ\bar v'=\bar v$ and
$\varphi\circ\bar w'=\bar w$, and they have differentials $\bar a$ and $\bar b$: the coefficient of $\varepsilon$ in $f\circ\bar v'$
is the conjugate of that in $f\circ v'$, and the derivative of $f\circ\bar u$ at $0$ in the direction $\bar a$ is the conjugate of
the derivative of $f\circ u$ in the direction $a$. Finally $\pi\,\bar b^{\mathsf T}\bar H\,\overline{\bar a}=
\overline{\pi\,b^{\mathsf T}H\bar a}=\bar c$.

(2) Let $P$ be the centre of $z$, and let $b'$ be a family as in the definition of the Hodge norm at $z$, with Hodge Gram
matrix $\Gamma$ at $z$. By Lemma~\ref{cv-origin-analytic-hodge-points}(3), $\bar z$ also has centre $P$, so $b'$ satisfies the
conditions at $P$, and the conjugates $\bar v_p$ of values $v_p$ of the $b'_p$ at $z$ are values of the $b'_p$ at $\bar z$. By
(1), $\bar v_p$ and $\bar v_q$ have Hodge inner product $\bar\Gamma_{pq}$. So $b'$ has Hodge Gram matrix
$\bar\Gamma=\Gamma^{\mathsf T}$ at $\bar z$, as $\Gamma$ is Hermitian. It is positive definite, and its upper left blocks have
the determinants $\overline{\det\Gamma_l}=\det\Gamma_l$, which are real. So $t'$ generates the $m$-th graded piece at $\bar z$
with Hodge norm $\nu'$.
\end{proof}
% Lean: (1) is T.hodge_inner_product_conjugate = HodgeConj.hodgeInnerProduct_conj: tw A with isGrpPullback_tw, cjU
% (ex z = pc (u.ex (star z)), Lambda = comap star, trivializing functions conj (fr i (star z)), conjugate_uniformisation_holomorphic,
% conjugate_frame_holomorphic), isAH_cjU (Hm.map conj, chi(star l) conj, form_conj), isTan_conj, cjV and isDiff_cj
% (snd_val_cj, fderiv_cj), and form_conj for the value. (2) is HodgeLine.isGrNorm_conj: the same frame, lieVal_conj, the
% Gram matrix Gamma.map conj (with Gamma^T = Gamma.map conj from IsHermitian), PosDef.transpose and leadMinor_conj. Size:
% about 360 and 60 lines.

\begin{lemma}\label{cv-origin-analytic-hodge}
In the setting of Definition~\ref{cv-origin-hodge-component}, let $Z$ be a Hodge component of $\mathcal X$
at $S$ with coordinates $y$, and let a flag of $Z$ have root lines of multidegrees $a_{\tau,i}$ with
generators $e_{\tau,i}$. Then for each $(\tau,i)$ there is a metrized line bundle
$\bar\sigma_{\tau,i}=(a_{\tau,i},h_{\tau,i})$, invariant under conjugation, that is a Hodge metric on the
root line $\Lambda_{\tau,i}$ with generator $e_{\tau,i}$.
\end{lemma}
\leanref{Development/AutomorphicSpectrum/Proofs/GeometryProofs/HodgeCptFlagXdegStatement.lean}{GlobalTheory.T.hodge_cpt_flag_xdeg}
\leanref{Development/AutomorphicSpectrum/Proofs/GeometryProofs/HodgeCptFlagXdeg.lean}{GlobalTheory.CompanionTheory.Ar.HodgeLine.exists_hodgeLine}
\begin{proof}
Fix $(\tau,i)$ and write $a=a_{\tau,i}$ and $e=e_{\tau,i}$. Let $\sigma=1$ and $(s,m)=(\tau,i)$ if $i\le p$, and $\sigma=-1$ and
$(s,m)=(\bar\tau,N+1-i)$ if $i>p$, so that $\Lambda_{\tau,i}$ is the $m$-th graded piece of the flag of $\mathcal V_s$ or its
dual (Definition~\ref{cv-origin-hodge-component}(3)). With the Hodge norms of Lemma~\ref{cv-origin-model-metric-norm}, a
metrized line bundle $(a,h)$ is a Hodge metric on $\Lambda_{\tau,i}$ with generator $e$ if, for every complex point $z$ of $Z$
over which $\mathcal A_Z$ is abelian, with coordinates $w$, every $t\in\Omega$ that generates the $m$-th graded piece of the
flag of $\mathcal V_s$ at $z$ with Hodge norm $\nu$, and every chart $c$ with $w_{k,c_k}\neq0$ for all $k$, the element
$t^\sigma/(e\prod_ky_{k,c_k}^{a_k})$ is regular at $z$ with a value $U$ such that
$\nu^\sigma h(w)=|U|\prod_k|w_{k,c_k}|^{a_k}$. Call $(z,w,h_0)$, for a complex point $z$ of $Z$ over which $\mathcal A_Z$ is
abelian, coordinates $w$ of $z$ and a real $h_0>0$, a \emph{good triple} if this condition holds at $z$ and $w$ with $h_0$ in
place of $h(w)$.

\emph{Good triples exist.} Let $z$ be a complex point over which $\mathcal A_Z$ is abelian, with centre $P$ and coordinates
$w$. Choose $c^0$ with $w_{k,c^0_k}\neq0$ for all $k$. By the coordinates, the ratios $y_{k,j}/y_{k,c^0_k}$ are regular at
$P$, so $P$ lies in the open subscheme of the chart $c^0$, and $\Lambda_{\tau,i}$ is generated at $P$ by
$g=e\prod_ky_{k,c^0_k}^{a_k}\neq0$. The $m$-th graded piece is generated at $P$ by some $t_0$ (the first step of the proof of
Lemma~\ref{cv-origin-reference-descent-roots}), so $\Lambda_{\tau,i}(P)$ is generated by $t_0^\sigma$, and $g/t_0^\sigma$ is a
unit at $P$, with a nonzero value $u_0$ at $z$. Lemma~\ref{cv-origin-model-metric-norm}, for $z$ and $t_0$, gives $\nu_0>0$
such that $t$ generates the $m$-th graded piece at $z$ with Hodge norm $\nu$ if and only if $t/t_0$ is regular at $z$ with a
nonzero value $u_1$ and $\nu=|u_1|\nu_0$; in particular $t_0$ generates it with Hodge norm $\nu_0$. Put
$h_0=\prod_k|w_{k,c^0_k}|^{a_k}/(|u_0|\nu_0^\sigma)$. Let $t$ generate the graded piece with Hodge norm $\nu=|u_1|\nu_0$, and
let $c$ be a chart with $w_{k,c_k}\neq0$ for all $k$. As in the proof of Lemma~\ref{cv-origin-model-metric-chart},
$\mu=\prod_k(y_{k,c_k}/y_{k,c^0_k})^{a_k}$ is regular at $z$ with a nonzero value $\mu_0$, and
$|\mu_0|\prod_k|w_{k,c^0_k}|^{a_k}=\prod_k|w_{k,c_k}|^{a_k}$. So
\[
\frac{t^\sigma}{e\prod_ky_{k,c_k}^{a_k}}=\Big(\frac t{t_0}\Big)^\sigma\Big(\frac g{t_0^\sigma}\Big)^{-1}\mu^{-1}
\]
is regular at $z$ with value $U=u_1^\sigma u_0^{-1}\mu_0^{-1}$, and
$\nu^\sigma h_0=|u_1|^\sigma|u_0|^{-1}\prod_k|w_{k,c^0_k}|^{a_k}=|U|\prod_k|w_{k,c_k}|^{a_k}$. So $(z,w,h_0)$ is a good triple.

\emph{Uniqueness.} Let $(z,w,h_0)$ and $(z',w,h'_0)$ be good triples. Then $z'=z$ by
Lemma~\ref{cv-origin-analytic-hodge-points}(1). The previous paragraph gives $t_0$ that generates the $m$-th graded piece at
$z$ with Hodge norm $\nu_0>0$. For a chart $c$ with $w_{k,c_k}\neq0$ for all $k$, the two triples give values $U$ and $U'$ at $z$
of the same element with $\nu_0^\sigma h_0=|U|\prod_k|w_{k,c_k}|^{a_k}$ and $\nu_0^\sigma h'_0=|U'|\prod_k|w_{k,c_k}|^{a_k}$. A
value at $z$ is unique, as the local ring at the centre embeds in $\Omega$, so $U=U'$ and $h_0=h'_0$.

\emph{The metric.} For $w\in\hat C$ put $h_{\tau,i}(w)=h_0$ if $(z,w,h_0)$ is a good triple for some $z$, and
$h_{\tau,i}(w)=\prod_k(\sum_{j\in I_k}|w_{k,j}|^2)^{a_k/2}$ otherwise; this is well defined by the previous paragraph, and
positive. Let $\kappa\in(\mathbf C^\times)^K$. If $(z,w,h_0)$ is a good triple, so is $(z,\kappa w,\prod_k|\kappa_k|^{a_k}h_0)$:
$\kappa w$ are coordinates of $z$ by Lemma~\ref{cv-origin-analytic-hodge-points}(2), the charts $c$ with $\kappa_kw_{k,c_k}\neq0$
are those with $w_{k,c_k}\neq0$, and both sides of the identity are multiplied by $\prod_k|\kappa_k|^{a_k}$. Applied to
$\kappa^{-1}$, this shows that no good triple has $\kappa w$ if none has $w$. The second formula is also homogeneous of
multidegree $a$, so $h_{\tau,i}(\kappa w)=\prod_k|\kappa_k|^{a_k}h_{\tau,i}(w)$, and $\bar\sigma_{\tau,i}=(a,h_{\tau,i})$ is a
metrized line bundle. It is a Hodge metric on $\Lambda_{\tau,i}$ with generator $e$: at a complex point $z$ over which
$\mathcal A_Z$ is abelian, with coordinates $w$, the first step gives a good triple $(z,w,h_0)$, and $h_{\tau,i}(w)=h_0$.

\emph{Conjugation.} Let $(z,w,h_0)$ be a good triple, and $\bar z$ the conjugate point. By
Lemma~\ref{cv-origin-analytic-hodge-points}(3), $\bar w$ are coordinates of $\bar z$, and $\bar z$ has the centre of $z$, so
$\mathcal A_Z$ is abelian over $\bar z$. Let $t$ generate the $m$-th graded piece at $\bar z$ with Hodge norm $\nu$. By
Lemma~\ref{cv-origin-analytic-hodge-conjugate}(2), applied to $\bar z$, whose conjugate point is $z$, $t$ generates it at $z$
with Hodge norm $\nu$. So for a chart $c$ with $\bar w_{k,c_k}\neq0$ for all $k$, $t^\sigma/(e\prod_ky_{k,c_k}^{a_k})$ is regular
at $z$ with a value $U$ such that $\nu^\sigma h_0=|U|\prod_k|w_{k,c_k}|^{a_k}$, and its value at $\bar z$ is $\bar U$, of the same
absolute value. So $(\bar z,\bar w,h_0)$ is a good triple, and $h_{\tau,i}(\bar w)=h_{\tau,i}(w)$. If no good triple has $w$,
none has $\bar w$, by the same argument applied to $\bar w$, and the second formula is invariant under conjugation. So
$\bar\sigma_{\tau,i}$ is invariant under conjugation.
\end{proof}
% Lean: GlobalTheory.T.hodge_cpt_flag_xdeg = HodgeLine.exists_hodgeLine, for a model over any O with O -> F_c injective; the
% metric is the HermLine with a = xdeg r and h = metricF, with IsConjInvariantH. Good is the good triple (with ChartNorm);
% exists_metric is the first step (inChart_of_coords, exists_gen, T.hodge_cpt_interior_flag = Lemma
% cv-origin-model-metric-norm, exists_valAt_of_unit, valAt_ratio, valAt_zpow, valAt_inv); hu is the uniqueness
% (T.local_ring_points_existsUnique, valAt_unique); metricF takes refNorm (the Fubini-Study function) off the good triples,
% metricF_eq and metricF_of_not; Good.smul with ChartNorm.smul and refNorm_smul give the homogeneity; Good.conj (coordsAt_conj,
% centre_cj, isGrNorm_conj, ChartNorm.conj) and hneg give the conjugation invariance. Size: about 28,000 lines with its
% inputs, of which the Hodge norm (Lemma cv-origin-model-metric-norm and the lemmas before this one) is most. Nothing under
% the Hodge norm remains open: Lemmas cv-origin-hodge-norm-form-descent, cv-origin-hodge-norm-form-descent-ample (with
% -descent-ample-affine and -descent-ample-embedding) and cv-origin-hodge-norm-form-torus, below Lemma
% cv-origin-hodge-norm-form, are written out.

\subsection{Tame functions and normal crossings}\label{subsection-cv-tame-functions-crossings}

\begin{lemma}\label{cv-origin-analytic-local-tame}
Let $n\ge0$ and $0\le a\le n$, and let functions tame for $a$ be as in
Definition~\ref{cv-origin-analytic-system-weight}(2) and Levi forms as in Definition~\ref{cv-period-monge-ampere}(1).
\begin{enumerate}
\item If $f$ and $g$ are tame for $a$ and $c\in\mathbf R$, then $f+g$ and $cf$ are tame for $a$. If $f$ is tame for $a$
  and $g$ is a real function with $g(t)=f(t)$ for all $t$ near $0$ with $t_j\neq0$ for $j\le a$, then $g$ is tame for $a$.
\item If $f$ and $g$ are real functions of class $C^2$ near $0$ with nonnegative Levi forms at $0$, and $c\ge0$, then
  $f+g$ and $cf$ are of class $C^2$ near $0$ with nonnegative Levi forms at $0$, and so is every real function equal to $f$
  near $0$.
\item Let $F$ be a holomorphic function near $0$ with $F(0)\neq0$. Then $\log|F|$ is tame for $a$, and for all real $c$
  and $d$ the function $c\log|F|+d$ is of class $C^2$ near $0$ with nonnegative Levi form at $0$.
\end{enumerate}
\end{lemma}
\leanref{Development/WeilBound/CoreStatements/CentralCharts/PshAt0LogNorm.lean}{AnalyticReference.An.Tame.add}
\leanref{Development/WeilBound/CoreStatements/CentralCharts/PshAt0LogNorm.lean}{AnalyticReference.An.Tame.const_mul}
\leanref{Development/WeilBound/CoreStatements/CentralCharts/PshAt0LogNorm.lean}{AnalyticReference.An.Tame.congr}
\leanref{Development/WeilBound/CoreStatements/CentralCharts/PshAt0LogNorm.lean}{AnalyticReference.An.tame_log_norm}
\leanref{Development/WeilBound/CoreStatements/CentralCharts/PshAt0LogNorm.lean}{AnalyticReference.An.PshAt0.add}
\leanref{Development/WeilBound/CoreStatements/CentralCharts/PshAt0LogNorm.lean}{AnalyticReference.An.PshAt0.const_mul}
\leanref{Development/WeilBound/CoreStatements/CentralCharts/PshAt0LogNorm.lean}{AnalyticReference.An.PshAt0.congr}
\leanref{Development/WeilBound/CoreStatements/CentralCharts/PshAt0LogNorm.lean}{AnalyticReference.An.pshAt0_log_norm}
\begin{proof}
Put $\ell_a(t)=\exp(1)+\sum_{i\le a}\log(1/|t_i|^2)$. If $\max_j|t_j|<1$ and $t_j\neq0$ for $j\le a$, every term
$\log(1/|t_i|^2)$ is positive, so $\ell_a(t)\ge\exp(1)$ and $\log\ell_a(t)\ge1$. Hence if $f$ is smooth at these $t$ with
$\max_j|t_j|<\varepsilon$, and $|f(t)|\le C_1\log\ell_a(t)$ and $\|df(t)\|\le C_2\ell_a(t)^ee^{\ell_a(t)/2}$ there for some
$C_1,C_2,e\ge0$, then $f$ is tame for $a$, with $\min(\varepsilon,1)$ and $C=\max(C_1,C_2,e)$, as
$\ell_a(t)^e\le\ell_a(t)^C$.

(1) Let $f$ and $g$ be tame with constants $(\varepsilon_f,C_f)$ and $(\varepsilon_g,C_g)$. For $\max_j|t_j|<\min(\varepsilon_f,\varepsilon_g)$
with $t_j\neq0$ for $j\le a$, $f+g$ is smooth at $t$, $|f(t)+g(t)|\le(C_f+C_g)\log\ell_a(t)$ and
\[
\|d(f+g)(t)\|\le C_f\ell_a(t)^{C_f}e^{\ell_a(t)/2}+C_g\ell_a(t)^{C_g}e^{\ell_a(t)/2}\le(C_f+C_g)\ell_a(t)^{C_f+C_g}e^{\ell_a(t)/2}.
\]
For $cf$ the bounds hold with $|c|C_f$ in place of $C_1$ and $C_2$ and with $e=C_f$. For the last assertion, let $\delta>0$
be such that $g(t)=f(t)$ when $\max_j|t_j|<\delta$ and $t_j\neq0$ for $j\le a$. These $t$ form an open set, so $g$ agrees
with $f$ near each of them, and $g$ is tame with $\min(\varepsilon_f,\delta)$ and $C_f$.

(2) The class $C^2$ near $0$ is kept by sums and constant multiples, the second derivative at $0$ of $f+g$ is the sum of
those of $f$ and $g$, and that of $cf$ is $c$ times that of $f$. So the Levi forms at $0$ add and scale in the same way.
A function equal to $f$ near $0$ has the second derivative of $f$ at $0$.

(3) Near $0$, $F$ is holomorphic and nonzero, so $\log|F|$ is smooth there. By continuity at $0$, there is $\delta>0$ such that
$|\log|F(t)||$ and $\|d\log|F|(t)\|$ are at most $K=\max(|\log|F(0)||,\|d\log|F|(0)\|)+1$ for $\max_j|t_j|<\delta$. As
$\log\ell_a\ge1$, $\ell_a^K\ge1$ and $e^{\ell_a/2}\ge1$ at the relevant $t$, the first paragraph applies with
$C_1=C_2=e=K$. Next, $F(t)/F(0)$ is near $1$ for $t$ near $0$, so $G(t)=c\log(F(t)/F(0))$, with the principal branch of
the logarithm, is holomorphic near $0$, and $c\log|F(t)|+d=\operatorname{Re}G(t)+c\log|F(0)|+d$. The second derivative at
$0$ of $\operatorname{Re}G$ is $(v,w)\mapsto\operatorname{Re}d^2G(0)(v,w)$, which is complex bilinear inside the real part, so
$\operatorname{Re}d^2G(0)(iw,iw)=-\operatorname{Re}d^2G(0)(w,w)$ and the Levi form of $\operatorname{Re}G$ at $0$ is $0$. A
constant has Levi form $0$, so $c\log|F|+d$ is of class $C^2$ near $0$ with Levi form $0$ at $0$.
\end{proof}
% Lean: the AnalyticReference.An.H lemmas of ContDiffAtLogNorm.lean: tame_of (the first paragraph, with one_le_logR and
% one_le_log_logR), add, const_mul, congr (isOpen_off), tame_of_smooth and contDiffAt_log_norm for tame_log_norm;
% padd_lemma and padd, pconst_mul, pcongr; pre (the Levi form of the real part of a holomorphic function, through
% fderiv_clm_apply) and plog_norm (Complex.log of F/F(0), analytic as F(0)/F(0) = 1 is in the slit plane). The Lean also
% proves Tame.sum and PshAt0.sum by induction from add. Size: about 490 lines.

\begin{lemma}\label{cv-origin-analytic-local-roots}
Let $\mathcal X$, $Z$ with coordinates $y$, the flag with root lines of multidegrees $a_{\tau,i}$ and generators
$e_{\tau,i}$, the Hodge metrics $h_{\tau,i}$, $B$, $n$, $a$, $\Delta$ and $\gamma$ satisfy the hypotheses of
Lemma~\ref{cv-origin-analytic-local}, except possibly the condition on $F_c$ and $p+q\ge1$. For
$(\tau,i)\in\Psi\times\{1,\dots,N\}$ put $(\sigma,s,m)=(1,\tau,i)$ if $i\le p$ and $(\sigma,s,m)=(-1,\bar\tau,N+1-i)$ if
$i>p$. Then there are real functions $M_{s,m}$ on $\mathbf C^n$, for each label $s$ and $0\le m\le r_s$, and for each
$(\tau,i)$ a function $F_{\tau,i}$ holomorphic near $0$ with $F_{\tau,i}(0)\neq0$, such that:
\begin{enumerate}
\item $M_{s,0}=0$, every $M_{s,m}$ is tame for $a$ (Definition~\ref{cv-origin-analytic-system-weight}(2)), and if $a=0$
  every $M_{s,m}$ is of class $C^2$ near $0$ with nonnegative Levi form at $0$;
\item for every $(\tau,i)$, with $(\sigma,s,m)$ as above, and every $t$ near $0$ with $t_j\neq0$ for $j\le a$,
  \[
  \log h_{\tau,i}(\gamma(t))=\log|F_{\tau,i}(t)|-\tfrac\sigma2\big(M_{s,m}(t)-M_{s,m-1}(t)\big).
  \]
\end{enumerate}
\end{lemma}
\leanref{Development/WeilBound/CoreStatements/CentralCharts/ChartAnalytic.lean}{AnalyticReference.An.root_data}
\begin{proof}
\openstep{The construction of the potentials $M_{s,m}$, as the logarithms of the leading principal minors of the Hodge
Gram matrices of frames of the $\mathcal V_s$ along $\gamma$, with their bounds and plurisubharmonicity from the period
maps near the boundary, and of the functions $F_{\tau,i}$ from the root lines, is not yet written out.}
% Lean: AnalyticReference.An.root_data. For each label s, label_data: Fl.sub and exists_basic give a frame b of V s at the
% centre z0 of gamma(0) adapted to the flag; T.global_pair_sum_gram_bound gives the period Gram matrices
% periodGram L U (uncentre 0 rho pi t), and M m t = log leadMinor m of them; tame_logdet and pshAt0_logdet give (1);
% T.frame_domain_in_segre_chart, RootGram.exists_gram, T.gram_matrix_change_frame and T.interior_point_polarised identify
% them with the Hodge Gram matrices of b at the interior points. For each root r: GrNorm.exists_sub_mem_flagStep,
% T.root_lines_flag_degree, exists_val_ne and exists_val_fam give F r = F1 * prod_k gamma_(k, c k) ^ xdeg r k, and (2)
% comes from the Hodge metric (hhod, SecNorm) with leadMinor_pos_of_posDef and Real.log_sqrt. Size: about 77,000 lines by
% the archive's dependency count, which includes the Hodge norm of Lemma cv-origin-model-metric-norm and much of the
% next lemma; the period bound global_pair_sum_gram_bound is most of the rest.
\end{proof}

\begin{lemma}\label{cv-origin-analytic-local-central}
Let $\mathcal X$, $Z$ with coordinates $y$, the flag with root lines of multidegrees $a_{\tau,i}$ and generators
$e_{\tau,i}$, the Hodge metrics $h_{\tau,i}$, $B$, $n$, $a$, $\Delta$ and $\gamma$ satisfy the hypotheses of
Lemma~\ref{cv-origin-analytic-local}, with $a=0$, and let $\tau\in\Psi$. Then there are a function $F$ holomorphic near
$0$ with $F(0)\neq0$ and a real number $d$ such that $\sum_{i=1}^N\log h_{\tau,i}(\gamma(t))=\log|F(t)|+d$ for all $t$
near $0$.
\end{lemma}
\leanref{Development/WeilBound/CoreStatements/CentralCharts/ReferenceGeometryHolomorphicInverse.lean}{AnalyticReference.An.central_chart}
\begin{proof}
\openstep{The proof that the product of the Hodge metrics of the root lines of $\tau$ is, along $\gamma$, the absolute
value of a holomorphic function without zeros times a constant, from the flat generators of
Lemma~\ref{cv-function-generator} and Theorem~\ref{cv-flat-degree}, is not yet written out.}
% Lean: AnalyticReference.An.central_chart, under Flat.FlatLineGenerator, Flat.FlatLineGeneratorComponent, SplitComps
% and 1 <= p + q: the point z0 of gamma(0) (T.complex_points_iff_centres), An.central (c1, the exponents a and the
% constants tt s: the value of c1 / (prod_k y(k, c k) ^ a k * prod_i e(tau, i)) is nonzero, and at interior points its
% absolute value times prod_k |w(k, c k)| ^ a k is exp(tt s) prod_i h(tau, i)(w)), exists_val_fam for F1 and
% exists_emb_fam for the embedding s of F_c; F = F1 * prod_k gamma_(k, c k) ^ a k and d = -tt s. Size: about 36,000
% lines by the archive's dependency count, through An.central (the flat line generators, the comparison of the flag with
% the CM type lattices and the degree of the flag).
\end{proof}

\begin{lemma}\label{cv-origin-analytic-local}
In the setting of Definition~\ref{cv-origin-hodge-component}, let the moduli schemes of the PEL datum at
all levels be given, let $\mathcal X$ be a global model (Definition~\ref{cv-volume-global-model}) such that
$F_c$ is algebraically closed in the function field of every irreducible component of $\mathcal X$, and
suppose $p+q\ge1$. Let $Z$ be a Hodge component of $\mathcal X$ at $S$ with coordinates $y$, a flag of $Z$
with root lines of multidegrees $a_{\tau,i}$ with generators $e_{\tau,i}$, and Hodge metrics
$\bar\sigma_{\tau,i}=(a_{\tau,i},h_{\tau,i})$ on them. Let $B\subset\hat C$ be a set such that a complex point
of $Z$ with coordinates $w$ is one over which $\mathcal A_Z$ is not abelian if and only if $w\in B$. Let
$n\ge0$, let $0\le a\le n$, let $\Delta=\{t\in\mathbf C^n:\max_i|t_i|<1\}$, and let $\gamma:\Delta\to\hat X(\mathbf C)$
be holomorphic, where $X$ is the kernel of $y$, such that no nonzero tangent vector at $t$ is mapped
into the tangent space of the orbit of $\gamma(t)$ under $(\mathbf C^\times)^K$, $\gamma(t)$ and $\gamma(t')$
lie in one orbit only if $t=t'$, and $\gamma(t)\in B$ if and only if $t_i=0$ for some $i\le a$. Put
$\ell(t)=\exp(1)+\sum_{i\le a}\log(1/|t_i|^2)$, and for $\tau\in\Psi$ and $1\le i\le N$ put
$u_{\tau,i}(t)=-\sum_{k\le i}\log h_{\tau,k}(\gamma(t))$. Then:
\begin{enumerate}
\item there are $0<\varepsilon\le1$ and $C\ge0$ such that, for $\max_j|t_j|<\varepsilon$ with $t_j\neq0$ for
  $j\le a$, $u_{\tau,i}$ is smooth at $t$, $|u_{\tau,i}(t)|\le C\log\ell(t)$ and
  $\|du_{\tau,i}(t)\|\le C\ell(t)^Ce^{\ell(t)/2}$;
\item if $a=0$, then $u_{\tau,i}$ is of class $C^2$ near $0$ with nonnegative Levi form at $0$, and so is
  $-u_{\tau,N}$.
\end{enumerate}
\end{lemma}
\leanref{Development/WeilBound/CoreStatements/CentralCharts/ReferenceGeometryHolomorphicInverse.lean}{AnalyticReference.An.local_bounds}
\begin{proof}
For $(\tau,i)\in\Psi\times\{1,\dots,N\}$ let $(\sigma,s,m)$ be as in Lemma~\ref{cv-origin-analytic-local-roots}, and let
$M_{s,m}$ and $F_{\tau,i}$ be the functions that it gives.

(1) Let $(\tau,i)$ be given. By Lemma~\ref{cv-origin-analytic-local-roots}(2), for $t$ near $0$ with $t_j\neq0$ for $j\le a$,
$\log h_{\tau,i}(\gamma(t))$ equals $\log|F_{\tau,i}(t)|-\frac\sigma2(M_{s,m}(t)-M_{s,m-1}(t))$. This function is tame for $a$
by Lemma~\ref{cv-origin-analytic-local-roots}(1) and Lemma~\ref{cv-origin-analytic-local-tame}(1) and (3). So
$t\mapsto\log h_{\tau,i}(\gamma(t))$ is tame for $a$ by Lemma~\ref{cv-origin-analytic-local-tame}(1), and so is
$u_{\tau,i}=-\sum_{k\le i}\log h_{\tau,k}\circ\gamma$. Being tame for $a$ is the assertion (1)
(Definition~\ref{cv-origin-analytic-system-weight}(2)).

(2) Let $a=0$, so that Lemma~\ref{cv-origin-analytic-local-roots}(2) holds for all $t$ near $0$, and let $\tau\in\Psi$ and
$1\le i\le N$. By Lemma~\ref{cv-origin-analytic-local-tame}(3), every $c\log|F_{\tau,k}|+d$ is of class $C^2$ near $0$ with
nonnegative Levi form at $0$; by Lemma~\ref{cv-origin-analytic-local-roots}(1) so is every $M_{s,m}$.

If $i\le p$, every $k\le i$ has $(\sigma,s,m)=(1,\tau,k)$. The sum telescopes, and $M_{\tau,0}=0$, so near $0$
\[
u_{\tau,i}(t)=-\sum_{k\le i}\Big(\log|F_{\tau,k}(t)|-\tfrac12\big(M_{\tau,k}(t)-M_{\tau,k-1}(t)\big)\Big)
=\sum_{k\le i}\big(-\log|F_{\tau,k}(t)|\big)+\tfrac12M_{\tau,i}(t),
\]
with $i\le p=r_\tau$. By Lemma~\ref{cv-origin-analytic-local-tame}(2), $u_{\tau,i}$ is of class $C^2$ near $0$ with
nonnegative Levi form at $0$.

If $i>p$, every $k>i$ has $(\sigma,s,m)=(-1,\bar\tau,N+1-k)$. As $k$ runs from $i+1$ to $N$, $N+1-k$ runs from $N-i$ down to
$1$, so the sum telescopes, and $M_{\bar\tau,0}=0$:
\[
\sum_{k>i}\log h_{\tau,k}(\gamma(t))=\sum_{k>i}\Big(\log|F_{\tau,k}(t)|+\tfrac12\big(M_{\bar\tau,N+1-k}(t)-M_{\bar\tau,N-k}(t)\big)\Big)
=\sum_{k>i}\log|F_{\tau,k}(t)|+\tfrac12M_{\bar\tau,N-i}(t),
\]
with $N-i<q=r_{\bar\tau}$. Lemma~\ref{cv-origin-analytic-local-central} gives $F$ holomorphic near $0$ with $F(0)\neq0$ and
$d\in\mathbf R$ with $\sum_{k=1}^N\log h_{\tau,k}(\gamma(t))=\log|F(t)|+d$ near $0$. Subtracting,
\[
u_{\tau,i}(t)=-\sum_{k\le i}\log h_{\tau,k}(\gamma(t))=\sum_{k>i}\log|F_{\tau,k}(t)|+\tfrac12M_{\bar\tau,N-i}(t)+\big(-\log|F(t)|-d\big)
\]
near $0$, and $u_{\tau,i}$ is of class $C^2$ near $0$ with nonnegative Levi form at $0$ by
Lemma~\ref{cv-origin-analytic-local-tame}(2) and (3).

Finally $-u_{\tau,N}(t)=\sum_{k=1}^N\log h_{\tau,k}(\gamma(t))=\log|F(t)|+d$ near $0$, for $F$ and $d$ given by
Lemma~\ref{cv-origin-analytic-local-central}, so $-u_{\tau,N}$ is of class $C^2$ near $0$ with nonnegative Levi form at $0$
by Lemma~\ref{cv-origin-analytic-local-tame}(2) and (3).
\end{proof}
% Lean: AnalyticReference.An.local_bounds: root_data (Lemma cv-origin-analytic-local-roots), hr from tame_log_norm, Tame.add
% and Tame.const_mul with Tame.congr, then Tame.sum; for a = 0, central_chart (Lemma cv-origin-analytic-local-central), the
% case i < p (zero-based) with PshAt0.sum, pshAt0_log_norm (-1) 0 and (hM (inl tau) (i+1)).2 with const_mul (1/2) and
% sum_Iic_tele; the case i >= p with the complement sum over Finset.Ioi i, pshAt0_log_norm 1 0 and (-1) (-d),
% (hM (inr tau) (p+q-(i+1))).2 and sum_Ioi_tele; and the last root with pshAt0_log_norm 1 d. Size of the top level: about
% 120 lines; the open parts are Lemmas cv-origin-analytic-local-roots and cv-origin-analytic-local-central.

\begin{definition}\label{cv-origin-analytic-system-pullback}
Let $\Omega$ be a field, let $y\in\Omega^{\sqcup_{k\in K}I_k}$, and let $\tilde X$ with maps
$\lambda_x:U_x\to\hat X(\mathbf C)$ be as in Definition~\ref{cv-origin-analytic}(3) for $y$. For $z\in\hat C$ and
$k\in K$ put $\|z_k\|^2=\sum_{i\in I_k}|z_{k,i}|^2$.
\begin{enumerate}
\item For $d\in\mathbf Z^K$, the \emph{Fubini--Study bundle} of multidegree $d$ is the metrized line bundle
  $\bar{\mathcal O}(d)=(d,h_d)$ (Definition~\ref{lit-multiprojective}(2)) with $h_d(z)=\prod_k\|z_k\|^{d_k}$. It is
  smooth and invariant under conjugation. Write $e_k\in\mathbf Z^K$ for the multidegree that is $1$ at $k$ and $0$
  elsewhere, and $\mathbf 1$ for the multidegree that is $1$ everywhere.
\item The \emph{pullback} to $\tilde X$ of a smooth metrized line bundle $\bar A=(a,h)$ is the family of functions
  $\vartheta_{\bar A,x}(x')=\log h(\lambda_x(x'))$ on $U_x$, for $x\in\tilde X$. It is \emph{positive} if, for every
  $x\in\tilde X$ and every chart $\chi$ of $\tilde X$, the Levi form of $\vartheta_{\bar A,x}\circ\chi^{-1}$
  (Definition~\ref{cv-period-monge-ampere}(1)) is positive definite at every point of $\chi(U_x)$.
\end{enumerate}
\end{definition}
\leanref{GroundTheory/HeightsAndTheta/Heights/HodgeLowerBound/FlagIndices/InstFintypeFlIdx.lean}{GroundTheory.Heights₂.Heights.LowerBoundFree.fsHerm}
\leanref{ReciprocityTheory/ReductiveGroups/ArithmeticGeometry/BirationalGeometry/InterRes.lean}{ReciprocityTheory.ChartRestriction.G.Geo.potOf}
\leanref{ReciprocityTheory/AnalyticForms/PotentialTheory/IsPotentialOfFun.lean}{ReciprocityTheory.AnForm.M.IsPosPot}

\begin{lemma}\label{cv-origin-analytic-system-arith}
In the setting of Definition~\ref{cv-origin-reference}, suppose $D\ge1$, and let a reference geometry of
$\mathcal X$ be given, with its notation. Then the kernel of $y$ is an arithmetic subvariety of dimension $D+1$
(Definition~\ref{lit-multiprojective}).
\end{lemma}
\leanref{Development/AutomorphicSpectrum/Proofs/MetricProofs/IsArithVarOfSncModelStatement.lean}{GlobalTheory.T.isArithVar_of_snc_model}
\begin{proof}
Let $\kappa$ be a chart of $y$ whose affine open subscheme $Z_\kappa$ contains the generic point of $Z$, and put
$A_\kappa=\mathbf Q[y_{k,i}/y_{k,\kappa_k}]$. Let $f_l$ and $a_l$ be as in Definition~\ref{cv-origin-reference}(3) at
$\kappa$. As $\sum_la_lf_l=1$, some $f_l$ is nonzero, and then $S=A_\kappa[1/f_l]$ has elements $z_1,\dots,z_D$ whose
differentials form a basis of $\Omega_{S/\mathbf Q}$. Every element of $\Omega$ is regular at the generic point of $Z$,
so by Lemma~\ref{cv-origin-snc-charts}(1) it is a quotient of two elements of
$\mathbf Z[y_{k,i}/y_{k,\kappa_k}]\subset S$. Hence $\Omega$ is the field of fractions of $S$. It is a finitely
generated field extension of $\mathbf Q$, and $\Omega_{\Omega/\mathbf Q}=\Omega\otimes_S\Omega_{S/\mathbf Q}$
\cite[Tag 00RT]{Sta} has dimension $D$ over $\Omega$.

The set $K$ is not empty: otherwise $\mathbf Z[y_{k,i}/y_{k,\kappa_k}]=\mathbf Z$, so $\Omega=\mathbf Q$ and
$\Omega_{\Omega/\mathbf Q}=0$, while $D\ge1$. As $\Omega$ has characteristic zero, it is formally smooth over
$\mathbf Q$ \cite[Tag 0321]{Sta}, so $H_1(L_{\Omega/\mathbf Q})=0$ \cite[Tag 031Z]{Sta}, and Cartier's equality
\cite[Tag 07E1]{Sta} gives $\operatorname{trdeg}_{\mathbf Q}\Omega=\dim_\Omega\Omega_{\Omega/\mathbf Q}=D$. The
paragraphs ``The function field'' and ``The dimension'' of the proof of Lemma~\ref{cv-origin-weight-birational} use
only that $Z$ is a Hodge component with coordinates $y$, that $\Omega$ has characteristic zero and transcendence
degree $D$ over $\mathbf Q$, and that $K$ is not empty. So they show that the kernel of $y$ is an arithmetic
subvariety of dimension $D+1$.
\end{proof}
% Lean: PositivityTheory.Xgp.t_c_bm_exists_arith (= Relay.isArithVar_of_snc_model): T.cover for the chart c of the
% generic point, the IsSncModel data at c with f l != 0, IsFractionRing S Omega from FracChart.exists_eq_div_adjoin,
% finrank_kaehler_fractionRing_eq_of_basis, hK (K nonempty), rank_kaehler_eq_trdeg (the text cites Cartier's
% equality), trdeg_int_eq_rat and birational_of_transcendence_degree (the last two paragraphs of the proof of
% Lemma cv-origin-weight-birational). Size: about 300 lines.

\begin{lemma}\label{cv-origin-analytic-system-boundary}
In the setting of Definition~\ref{cv-origin-reference}, let a reference geometry of $\mathcal X$ be given, with its
notation, let $X$ be the kernel of $y$, and let $B\subset\hat X(\mathbf C)$ be the set of points at which every
element of $J_\partial$ vanishes. Then $B$ is closed in $\hat X(\mathbf C)$, and $cz\in B$ and $\bar z\in B$ for
every $z\in B$ and every $c\in(\mathbf C^\times)^K$, acting blockwise.
\end{lemma}
\leanref{Development/AutomorphicSpectrum/Proofs/GeometryProofs/BoundaryPointsClosedStatement.lean}{GlobalTheory.T.boundary_points_closed}
\begin{proof}
First, a point $w\in\hat C$ lies in $\hat X(\mathbf C)$ if and only if $w$ are coordinates of a complex point of $Z$
(Definition~\ref{cv-origin-reference}(4)). Let $w\in\hat X(\mathbf C)$, and choose $\kappa$ with
$w_{k,\kappa_k}\neq0$ for every $k$. Then $y_{k,\kappa_k}\neq0$, as otherwise $x_{k,\kappa_k}$ lies in the kernel
and vanishes at $w$; so $\kappa$ is a chart, and Lemma~\ref{cv-origin-boundary-points}(2) with $L=\mathbf C$ gives a
complex point of $Z$ with coordinates $w$. Conversely, let $w$ be coordinates of a complex point $\zeta$ of $Z$,
choose $\kappa$ with $w_{k,\kappa_k}\neq0$ for every $k$, and let $f$ be in the kernel of $y$. The kernel is
multihomogeneous, so each multihomogeneous component $f_d$ of $f$, of multidegree $d$, lies in it. As
$f_d(y_{k,i}T_k)=\prod_kT_k^{d_k}f_d(y)$, $f_d(y)=0$, and $f_d(y_{k,i}/y_{k,\kappa_k})=\prod_ky_{k,\kappa_k}^{-d_k}f_d(y)=0$.
The ratios $y_{k,i}/y_{k,\kappa_k}$ are
regular at the image of $\zeta$ with values $w_{k,i}/w_{k,\kappa_k}$, so $f_d(w_{k,i}/w_{k,\kappa_k})=0$ and
$f_d(w)=\prod_kw_{k,\kappa_k}^{d_k}f_d(w_{k,i}/w_{k,\kappa_k})=0$. So $f(w)=0$.

The set $B$ is the set of points of $\hat X(\mathbf C)$ at which the polynomials of $J_\partial$ vanish, so it is
closed. Let $z\in B$; by the first paragraph, $z$ are coordinates of a complex point $\zeta$ of $Z$. For
$c\in(\mathbf C^\times)^K$, $cz$ are
coordinates of $\zeta$ as well, since the ratios $z_{k,i}/z_{k,j}$ do not change. The complex conjugate $\bar z$ are
coordinates of $\zeta$ composed with the morphism of $\operatorname{Spec}\mathbf C$ given by complex conjugation,
which has the same image in $Z$. Whether $\mathcal A_Z$ is abelian over a complex point depends only on its image in
$Z$. So by Definition~\ref{cv-origin-reference}(4), $cz$ and $\bar z$ lie in the zero set of $J_\partial$, as $z$
does.
\end{proof}
% Lean: BdSym.boundary_closed_torus_conj: complex_points_iff_centres (= mem_mcPts_iff_exists_coordsAt, from
% local_ring_points_existsUnique, Lemma cv-origin-boundary-points), the intersection of the zero sets for closedness,
% CoordsAt of torusSmul c w for the torus, and HodgeLine.coordsAt_conj with centre_cj for the conjugate; bd_iff is
% Definition cv-origin-reference(4) with T the abelian locus. Size: about 230 lines.

\begin{lemma}\label{cv-origin-analytic-system-ideal}
Let $\Omega$ be a field of characteristic zero, let $e\ge0$, let $y\in\Omega^{\sqcup_{k\in K}I_k}$ be a point whose
kernel is an arithmetic subvariety $X$ of dimension $e+1$, and let $J$ be an ideal of $R=\mathbf Z[x_{k,i}]$ that
satisfies condition (3) of Definition~\ref{cv-origin-reference} for $y$ with $e$ in place of $D$. Let
$B\subset\hat X(\mathbf C)$ be the set of points at which every element of $J$ vanishes, and suppose that $cz\in B$
for every $z\in B$ and every $c\in(\mathbf C^\times)^K$. For $z\in\hat C$ let $I_z$ be the kernel of
$R\to\mathbf C[T_k:k\in K]$, $x_{k,i}\mapsto z_{k,i}T_k$, and put $I=\bigcap_{z\in B}I_z$. Then $I$ is
multihomogeneous and not contained in the kernel of $y$, it satisfies condition (3) of
Definition~\ref{cv-origin-reference} for $y$ with $e$ in place of $D$, and $B$ is the set of points of
$\hat X(\mathbf C)$ at which every element of $I$ vanishes.
\end{lemma}
\leanref{Development/WeilBound/CoreStatements/PluriharmonicCharts/LeMCoordIdealOfTorus.lean}{PositivityTheory.Xgp.c_bm_exists_boundary.exists_homogIdeal}
\begin{proof}
For a point $z$ and a polynomial $f$, $f\in I_z$ if and only if $f(cz)=0$ for every $c\in(\mathbf C^\times)^K$, as
$(\mathbf C^\times)^K$ is Zariski dense in $\mathbf C^K$. So $J\subset I$, by the hypothesis on $B$. Each $I_z$ is
multihomogeneous, being the kernel of a homomorphism compatible with the gradings by multidegree, and hence so is
$I$. If $f\in I$ and $z\in B$, then $f(z)=0$; with $J\subset I$ this shows that $B$ is the set of points of
$\hat X(\mathbf C)$ at which $I$ vanishes.

As the kernel of $y$ is relevant, for each $k$ some $x_{k,i}$ is not in it, so $y_{k,i}\neq0$, and $y$ has a chart
$\kappa$. Let $f_l$ and $a_l$ be as in condition (3) for $J$ at $\kappa$. As $\sum_la_lf_l=1$, some $f_l\neq0$; let
$S=A_\kappa[1/f_l]$, with its elements $z_1,\dots,z_e$ and integer $n$. The zero ideal of $S\subset\Omega$ is prime,
and it contains none of the $z_j$, whose differentials are members of a basis. So by condition (3) it does not
contain all the values of $J$: some $g\in J$ has $g(y_{k,i}/y_{k,\kappa_k})\neq0$. If $I$ were contained
in the kernel of $y$, then so would be $g$, and $g(y_{k,i}/y_{k,\kappa_k})=0$ as in the proof of
Lemma~\ref{cv-origin-analytic-system-boundary}. So $I$ is not contained in the kernel.

For condition (3) for $I$, take at each chart $\kappa$ the elements $f_l$, $a_l$, $z_j$ and integers $n$ of condition
(3) for $J$. Let $S=A_\kappa[1/f_l]$ with $f_l\neq0$, and let $\mathfrak q$ be a prime of $S$. If $\mathfrak q$
contains the values $g(y_{k,i}/y_{k,\kappa_k})$ of all $g\in I$, it contains those of all $g\in J$. Conversely,
suppose that $\mathfrak q$ contains the values of all $g\in J$. The ring $S$ is a finitely generated
$\mathbf Q$-algebra, so $\mathfrak q$ is the intersection of the maximal ideals $\mathfrak m\supset\mathfrak q$, and
each $S/\mathfrak m$ is finite over $\mathbf Q$ \cite[Tag 00FV]{Sta}. Let $\mathfrak m$ be one of them; as
$S/\mathfrak m$ embeds in $\mathbf C$, there is a homomorphism $\mathrm{ev}:S\to\mathbf C$ with kernel $\mathfrak m$. Put
$w_{k,i}=\mathrm{ev}(y_{k,i}/y_{k,\kappa_k})$. Then $w_{k,\kappa_k}=1$; every element of the kernel of $y$ vanishes
at $w$, as it vanishes on the ratios; and every $g\in J$ vanishes at $w$, as $g(w)=\mathrm{ev}(g(y_{k,i}/y_{k,\kappa_k}))$
and the value lies in $\mathfrak m$. So $w\in B$. For $g\in I$ this gives $g\in I_w$, so
$\mathrm{ev}(g(y_{k,i}/y_{k,\kappa_k}))=g(w)=0$, and the value of $g$ lies in $\mathfrak m$. Hence it lies in
$\mathfrak q$.
\end{proof}
% Lean: exists_homogIdeal: I = iInf over z in mcPts P meet mcPts bd of mCoordIdeal z; hle (bd <= I) from
% le_mCoordIdeal_of_torus; hhom from CoordinateConeArithVar.component_mem; I not <= P from the IsSncModel data at c
% (hq at the zero ideal, aeval_ratio_eq_zero); IsSncModel for I from cond_of_bd (isJacobsonRing_of_finiteType,
% exists_algHom_ker_eq, aeval_ratio_eq_zero); the zero sets from aeval_eq_zero_of_mem. Size: about 350 lines.

\begin{lemma}\label{cv-origin-analytic-system-crossings-dehomogenize}
Let $K$ be a finite set with finite sets $I_k$ for $k\in K$, let $Q$ be a prime ideal of
$\mathbf C[x_{k,i}:k\in K,\ i\in I_k]$ generated by multihomogeneous elements, and let $\kappa=(\kappa_k)$ be indices
$\kappa_k\in I_k$ with $x_{k,\kappa_k}\notin Q$ for every $k$. Let $Q_\kappa$ be the ideal generated by $Q$ and the
$x_{k,\kappa_k}-1$, and for a polynomial $f$ let $f^\kappa$ be the polynomial obtained from $f$ by putting
$x_{k,\kappa_k}=1$ for every $k$. Then $Q_\kappa$ is prime, and a multihomogeneous polynomial $f$ lies in $Q$ if and
only if $f^\kappa\in Q_\kappa$.
\end{lemma}
\leanref{Development/AutomorphicSpectrum/Proofs/MetricProofs/AffineChartIdealPrimeAndSaturatedStatement.lean}{GlobalTheory.T.affine_chart_ideal_prime_and_saturated}
\begin{proof}
Let $E$ be the field of fractions of $\mathbf C[x_{k,i}]/Q$, write $\bar x_{k,i}$ for the images of the variables, so
that $\bar x_{k,\kappa_k}\neq0$, and let $\theta:\mathbf C[x_{k,i}]\to E$ be the homomorphism with
$\theta(x_{k,i})=\bar x_{k,i}/\bar x_{k,\kappa_k}$. If $f$ is multihomogeneous of multidegree $a$, then
$f(\bar x)=\prod_k\bar x_{k,\kappa_k}^{a_k}\,\theta(f)$, so $\theta(f)=0$ if and only if $f\in Q$. Hence $\theta$ kills
the multihomogeneous generators of $Q$ and the $x_{k,\kappa_k}-1$, and $Q_\kappa\subset\ker\theta$.

Conversely, every polynomial $g$ is congruent modulo $Q_\kappa$ to a multihomogeneous polynomial. This holds for the
constants. If $g_1\equiv H_1$ and $g_2\equiv H_2$ with $H_1$ and $H_2$ multihomogeneous of multidegrees $d_1$ and $d_2$,
then $g_1+g_2\equiv H_1\prod_kx_{k,\kappa_k}^{d_{2,k}}+H_2\prod_kx_{k,\kappa_k}^{d_{1,k}}$, which is multihomogeneous of
multidegree $d_1+d_2$, and $g_1x_{k,i}\equiv H_1x_{k,i}$. Let $g\in\ker\theta$ and $g\equiv H$ with $H$
multihomogeneous. Then $\theta(H)=\theta(g)=0$, so $H\in Q$ and $g\in Q_\kappa$. So $Q_\kappa=\ker\theta$, which is
prime, as $E$ is a field. Finally $f-f^\kappa$ lies in the ideal generated by the $x_{k,\kappa_k}-1$, so
$\theta(f^\kappa)=\theta(f)$, and $f^\kappa\in Q_\kappa$ if and only if $\theta(f)=0$; for multihomogeneous $f$ this
means $f\in Q$.
\end{proof}
% Lean: SncRes.Dict.ideal_isPrime_and_mem_iff: theta into the FractionRing of the quotient, phi_eq_prod_mul_theta,
% theta_eq_zero_iff, exists_form (the congruence to a multihomogeneous H), ker_theta, and theta_symm_bind1 for f^kappa,
% which the Lean writes as bind1 (SncRes.lift p) f in renamed variables. Size: about 150 lines.

\begin{lemma}\label{cv-origin-analytic-system-crossings-regular}
Let $\Omega$ be a field of characteristic zero, let $e\ge0$, let $y\in\Omega^{\sqcup_{k\in K}I_k}$ be a point whose
kernel $P$ is an arithmetic subvariety of dimension $e+1$, and let $J$ be an ideal of $R=\mathbf Z[x_{k,i}]$ that
satisfies condition (3) of Definition~\ref{cv-origin-reference} for $y$ with $e$ in place of $D$. Let $\kappa$ be a
chart of $y$, and let $I_\kappa$ be the ideal of $\mathbf C[x_{k,i}]$ generated by $P$ and the $x_{k,\kappa_k}-1$ for
$k\in K$. Then at every common zero $w$ of $I_\kappa$ the local ring $\mathcal O_w$ of $\mathbf C[x_{k,i}]/I_\kappa$ at
$w$ is regular of dimension $e$.
\end{lemma}
\leanref{Development/AutomorphicSpectrum/Proofs/MetricProofs/SncResolutionRegularLocalRingStatement.lean}{GlobalTheory.T.snc_resolution_regular_local_ring}
\begin{proof}
Put $u=(y_{k,i}/y_{k,\kappa_k})_{k,i}$ and $A_\kappa=\mathbf Q[u]$. The proof of
Lemma~\ref{cv-origin-model-resolution-regular}, with $e$ in place of $D$, uses the smoothness of $A_\kappa$ only to
bound the dimension of the space $N$ of the vectors $v\in\mathbf C^{\sqcup_kI_k}$ with $\partial_vh(w)=0$ for every
$h\in I_\kappa$, which is the dimension of the cotangent space of $\mathcal O_w$. Its other steps show that
$\mathcal O_w$ has dimension $e$, and that every polynomial $g$ over $\mathbf Q$ with $g(u)=0$ lies in $I_\kappa$, so
that $g(w)=0$ and $\partial_vg(w)=0$ for $v\in N$. It remains to show that $\dim N\le e$.

Let $f_l=F_l(u)$ and $a_l=A_l(u)$ be as in condition (3) at $\kappa$, with polynomials $F_l$ and $A_l$ over
$\mathbf Q$. As $\sum_lA_lF_l-1$ vanishes at $u$, it vanishes at $w$, so $F_l(w)\neq0$ for some $l$; then $f_l\neq0$,
as otherwise $F_l(w)=0$. Let $S=A_\kappa[1/f_l]$, with its elements $z_1,\dots,z_e$, whose differentials span
$\Omega_{S/\mathbf Q}$. Let $C$ be $\mathbf C$ or $\mathbf C[\varepsilon]/(\varepsilon^2)$, and let
$t\in C^{\sqcup_kI_k}$ be such that $g(t)=0$ for every polynomial $g$ over $\mathbf Q$ with $g(u)=0$, and $F_l(t)$ is a
unit. Then there is exactly one homomorphism $\Phi:S\to C$ with $\Phi(u)=t$: every element of $S$ is
$g(u)/F_l(u)^n$, and $g(u)/F_l(u)^n\mapsto g(t)/F_l(t)^n$ is well defined, since a relation
$g(u)F_l(u)^{n'}=g'(u)F_l(u)^n$ gives $g(t)F_l(t)^{n'}=g'(t)F_l(t)^n$. For $C=\mathbf C$ and $t=w$ this gives
$\psi:S\to\mathbf C$. Write $du_q=\sum_ic_{q,i}\,dz_i$ in $\Omega_{S/\mathbf Q}$, and let
$b_i=(\psi(c_{q,i}))_q\in\mathbf C^{\sqcup_kI_k}$. For $v\in N$ take $C=\mathbf C[\varepsilon]/(\varepsilon^2)$ and
$t=w+\varepsilon v$: for $g$ over $\mathbf Q$ with $g(u)=0$, $g(w+\varepsilon v)=g(w)+\varepsilon\,\partial_vg(w)=0$,
and $F_l(w+\varepsilon v)$ is a unit. The $\varepsilon$-part of the resulting $\Phi$ is a $\mathbf Q$-derivation
$\delta:S\to\mathbf C$ along $\psi$, with $\delta(u_q)=v_q$. So $v_q=\sum_i\psi(c_{q,i})\delta(z_i)$, and
$v=\sum_i\delta(z_i)b_i$ lies in the span of $b_1,\dots,b_e$. Hence $\dim N\le e$.
\end{proof}
% Lean: SncRes.Reg.isRegularLocalRing_chart, the proof of SmoothMan.isRegularLocalRing_chart (Lemma
% cv-origin-model-resolution-regular) with card_le_finrank_jacSpanAt in place of SmoothMan.card_le_of_smooth: the
% IsSncModel data at p, F l with eval w (F l) != 0, card_le_of_kaehler (exists_unique_lift for psi and for Phi into
% TrivSqZeroExt C C, aeval_dual, exists_span_of_derivations for the b_i, card_le_finrank_add). Size: about 2,000 lines,
% most of them shared with Lemma cv-origin-model-resolution-regular.

\begin{lemma}\label{cv-origin-analytic-system-crossings-params}
In the setting of Lemma~\ref{cv-origin-analytic-system-crossings-regular}, let $w$ be a common zero of $I_\kappa$, let
$\mathfrak m_w$ be the ideal of the polynomials that vanish at $w$, and let $J_\kappa$ be the ideal of
$\mathbf C[x_{k,i}]$ generated by $P$, $J$ and the $x_{k,\kappa_k}-1$. Then there are polynomials $r_1,\dots,r_e$ with
complex coefficients that vanish at $w$, whose images generate the maximal ideal of $\mathcal O_w$, and a set
$T\subset\{1,\dots,e\}$, such that a prime ideal $\mathfrak q$ of $\mathbf C[x_{k,i}]$ with
$I_\kappa\subset\mathfrak q\subset\mathfrak m_w$ contains $J_\kappa$ if and only if it contains $r_i$ for some $i\in T$.
\end{lemma}
\leanref{Development/AutomorphicSpectrum/Proofs/MetricProofs/SncChartParametersAtConePointStatement.lean}{GlobalTheory.T.snc_chart_parameters_at_cone_point}
\begin{proof}
Put $u=(y_{k,i}/y_{k,\kappa_k})_{k,i}$ and $A_\kappa=\mathbf Q[u]$. As in the proof of
Lemma~\ref{cv-origin-analytic-system-crossings-regular}, every polynomial $g$ over $\mathbf Q$ with $g(u)=0$ lies in
$I_\kappa$, and condition (3) at $\kappa$ gives a polynomial $F$ over $\mathbf Q$ with $F(u)\neq0$ and $F(w)\neq0$, the
algebra $S=A_\kappa[1/F(u)]$, its elements $z_1,\dots,z_e$, whose differentials span $\Omega_{S/\mathbf Q}$, and an
integer $n\le e$ such that a prime of $S$ contains the values $g(u)$ of all $g\in J$ if and only if it contains one of
$z_1,\dots,z_n$. Write $z_i=Z_i(u)/F(u)^{N_i}$ with polynomials $Z_i$ over $\mathbf Q$ and integers $N_i\ge0$, put
$c_i=Z_i(w)/F(w)^{N_i}$ and $r_i=Z_i-c_iF^{N_i}$, and let $T$ be the set of the $i\le n$ with $c_i=0$. Then
$r_i(w)=0$.

\emph{The parameters.} Let $v\in\mathbf C^{\sqcup_kI_k}$ with $\partial_vh(w)=0$ for every $h\in I_\kappa$ and for
$h=r_1,\dots,r_e$. As in the proof of Lemma~\ref{cv-origin-analytic-system-crossings-regular}, there is a homomorphism
$\Phi:S\to\mathbf C[\varepsilon]/(\varepsilon^2)$ with $\Phi(g(u))=g(w)+\varepsilon\,\partial_vg(w)$ for polynomials $g$
over $\mathbf Q$. As $r_i(w)=0$ and $\partial_vr_i(w)=0$, $\Phi(Z_i(u))=c_i\Phi(F(u))^{N_i}$, so $\Phi(z_i)=c_i$, whose
$\varepsilon$-part is $0$. The $\varepsilon$-part of $\Phi$ is a $\mathbf Q$-derivation of $S$ along the homomorphism
given by the first part. It vanishes on the $z_i$, whose differentials span $\Omega_{S/\mathbf Q}$, hence on $S$, and
its value at $u_q$ is $v_q$. So $v=0$. Hence the gradients at $w$ of the elements of $I_\kappa$ and of the $r_i$ span
the dual of $\mathbf C^{\sqcup_kI_k}$, and each $x_q-w_q$ is congruent modulo $\mathfrak m_w^2$ to a linear combination
of such elements, as a polynomial that vanishes at $w$ with gradient $0$ there lies in $\mathfrak m_w^2$. The maximal
ideal $M$ of $\mathcal O_w$ is generated by the image of $\mathfrak m_w$, and the elements of $I_\kappa$ map to $0$; so
$M\subset(r_1,\dots,r_e)+M^2$, and the images of the $r_i$ generate $M$ by Nakayama's lemma \cite[Tag 00DV]{Sta}.

\emph{The primes.} Let $\mathfrak q$ be a prime with $I_\kappa\subset\mathfrak q\subset\mathfrak m_w$, let $E$ be the
field of fractions of $\mathbf C[x_{k,i}]/\mathfrak q$, and let $\pi$ be the map to $E$. As $I_\kappa\subset\mathfrak q$,
$g(u)\mapsto\pi(g)$ is a well defined homomorphism on $A_\kappa$, and as $F\notin\mathfrak q$, it extends to a
homomorphism $\varphi:S\to E$, whose kernel $\mathfrak q'$ is prime. The value $g(u)$ of $g\in J$ maps to $\pi(g)$, and
$J_\kappa$ is generated by $I_\kappa$ and $J$. So $J_\kappa\subset\mathfrak q$ if and only if $\mathfrak q'$ contains the
values of all $g\in J$, that is, by condition (3), if and only if $z_i\in\mathfrak q'$ for some $i\le n$. As
$\varphi(z_i)\varphi(F(u))^{N_i}=\pi(Z_i)$ and $\varphi(F(u))\neq0$, this means that $Z_i\in\mathfrak q$ for some
$i\le n$. If $Z_i\in\mathfrak q\subset\mathfrak m_w$, then $Z_i(w)=0$, so $c_i=0$, $i\in T$ and $r_i=Z_i$; and
$r_i=Z_i$ for $i\in T$. So $J_\kappa\subset\mathfrak q$ if and only if $r_i\in\mathfrak q$ for some $i\in T$.
\end{proof}
% Lean: SncRes.Par.exists_params: exists_basicOpen_of_isSncModel (F, S, z and n), exists_mul_pow_eq_aeval_of_mem_adjoin_inv
% (Z and N), r0 = Z - C c * F^N and Jf = {i < n : c i = 0}; span_algebraMap_eq_map_ptIdeal (the Jacobian criterion with
% Nakayama's lemma, ptI_le_of_span_grad_eq_top) with Ktan.tangent_eq_zero (snd_eq_zero_of_span); the regularity, which
% the Lean also concludes here, from Kreg.isRegularLocalRing_locRing_of_span and ringKrullDim_locAt_chart; the primes
% from PrimeRead.boundary_on_primes (read_in_field, Kext.exists_algHom_adjoin_inv). Size: about 2,250 lines, partly
% shared with Lemma cv-origin-analytic-system-crossings-regular.

\begin{lemma}\label{cv-origin-analytic-system-crossings-local}
Let $\mathfrak a\subset\mathfrak b$ be ideals of a polynomial ring $\mathbf C[x_1,\dots,x_n]$, let $w$ be a common zero
of $\mathfrak a$, let $\mathfrak m_w$ be the ideal of the polynomials that vanish at $w$, and suppose that the local ring
$\mathcal O_w$ of $\mathbf C[x_1,\dots,x_n]/\mathfrak a$ at $w$ is regular of dimension $e$. Let $r_1,\dots,r_e$ be
polynomials whose images $\rho_1,\dots,\rho_e$ generate the maximal ideal of $\mathcal O_w$, and let
$T\subset\{1,\dots,e\}$ be a set such that a prime $\mathfrak q$ with $\mathfrak a\subset\mathfrak q\subset\mathfrak m_w$
contains $\mathfrak b$ if and only if it contains $r_i$ for some $i\in T$. For each $i$ let $\mathfrak q_i$ be the set of
polynomials whose image in $\mathcal O_w$ lies in $\rho_i\mathcal O_w$. Then:
\begin{enumerate}
\item for $i\in T$, $\mathfrak q_i$ is a prime ideal with $\mathfrak a\subset\mathfrak q_i\subset\mathfrak m_w$ that
  contains $\mathfrak b$ and $r_i$ and none of the $r_j$ with $j\in T\setminus\{i\}$, and it is contained in every prime
  $\mathfrak q$ with $\mathfrak a\subset\mathfrak q\subset\mathfrak m_w$ and $r_i\in\mathfrak q$;
\item there is a polynomial $h$ with $h(w)\neq0$ such that $h\prod_{i\in T}r_i$ lies in the radical $\sqrt{\mathfrak b}$
  of $\mathfrak b$, and $hg\in(\prod_{i\in T}r_i)+\mathfrak a$ for every $g\in\sqrt{\mathfrak b}$.
\end{enumerate}
\end{lemma}
\leanref{Development/AutomorphicSpectrum/Proofs/MetricProofs/RegularSystemOfParametersLocalRingAltBis.lean}{GlobalTheory.T.regular_system_of_parameters_local_ring_alt}
\leanref{Development/AutomorphicSpectrum/Proofs/MetricProofs/RegularSystemOfParametersLocalRingStatement.lean}{GlobalTheory.T.exists_radical_locally_eq_span_prod}
\begin{proof}
A polynomial $g$ maps to a unit of $\mathcal O_w$ if and only if $g(w)\neq0$; every element of $\mathcal O_w$ is the
image of a quotient $a/b$ of polynomials with $b(w)\neq0$; and if a polynomial $a$ maps to $0$ in $\mathcal O_w$, then
$sa\in\mathfrak a$ for some $s$ with $s(w)\neq0$. Let $M$ be the maximal ideal of $\mathcal O_w$. No $\rho_j$ lies in the
ideal generated by the other $\rho_i$ and $M^2$: otherwise, by Nakayama's lemma \cite[Tag 00DV]{Sta}, $M$ would be
generated by $e-1$ elements, while $\mathcal O_w$ is regular of dimension $e$. So $\rho_1,\dots,\rho_e$ is a minimal set
of generators of $M$, each $\mathcal O_w/\rho_i\mathcal O_w$ is regular \cite[Tag 00NQ]{Sta}, hence a domain
\cite[Tag 00NP]{Sta}, and each $\rho_i$ is prime. For $i\neq j$, $\rho_i$ and $\rho_j$ are not associated, as
$\rho_j\notin\rho_i\mathcal O_w$. The ring $\mathcal O_w$ is a domain \cite[Tag 00NP]{Sta}.

(1) The ideal $\mathfrak q_i$ is the preimage of the prime $\rho_i\mathcal O_w$, so it is prime, and it contains
$\mathfrak a$ and $r_i$. If $g\in\mathfrak q_i$ had $g(w)\neq0$, its image would be a unit in $\rho_i\mathcal O_w\subset M$;
so $\mathfrak q_i\subset\mathfrak m_w$, and by the hypothesis on $T$ it contains $\mathfrak b$. If $r_j\in\mathfrak q_i$
with $j\neq i$, then $\rho_i$ divides the prime $\rho_j$, and the two are associated, which is excluded. Let
$\mathfrak q$ be a prime with $\mathfrak a\subset\mathfrak q\subset\mathfrak m_w$ and $r_i\in\mathfrak q$, and let
$g\in\mathfrak q_i$, with image $\rho_ia_0/b_0$, where $b_0(w)\neq0$. Then $b_0g-a_0r_i$ maps to $0$, so
$s(b_0g-a_0r_i)\in\mathfrak a\subset\mathfrak q$ for some $s$ with $s(w)\neq0$. So $sb_0g\in\mathfrak q$, and
$g\in\mathfrak q$, as $(sb_0)(w)\neq0$ and $\mathfrak q\subset\mathfrak m_w$.

(2) The preimage of a prime $\mathfrak P$ of $\mathcal O_w$ is a prime $\mathfrak q$ with
$\mathfrak a\subset\mathfrak q\subset\mathfrak m_w$. So $\mathfrak b\mathcal O_w\subset\mathfrak P$ if and only if
$\mathfrak b\subset\mathfrak q$, if and only if $r_i\in\mathfrak q$ for some $i\in T$, if and only if
$\rho=\prod_{i\in T}\rho_i$ lies in $\mathfrak P$. As the radical of an ideal is the intersection of the primes that
contain it, $\mathfrak b\mathcal O_w$ and $\rho\mathcal O_w$ have the same radical. The ideal $\rho\mathcal O_w$ is its
own radical: if $g^N\in\rho\mathcal O_w$, then each prime $\rho_i$ with $i\in T$ divides $g$, and as they are pairwise
not associated, their product divides $g$, by induction on the number of factors. So
$\sqrt{\mathfrak b\mathcal O_w}=\rho\mathcal O_w$.

Hence $\rho^N\in\mathfrak b\mathcal O_w$ for some $N$. An element of $\mathfrak b\mathcal O_w$ is the image of a
quotient $a/b$ with $a\in\mathfrak b$ and $b(w)\neq0$; so $b(\prod_{i\in T}r_i)^N-a$ maps to $0$, and
$s(b(\prod_{i\in T}r_i)^N-a)\in\mathfrak a\subset\mathfrak b$ for some $s$ with $s(w)\neq0$. So $h'=sb$ has
$h'(w)\neq0$ and $h'(\prod_{i\in T}r_i)^N\in\mathfrak b$, and
$(h'\prod_{i\in T}r_i)^{N+1}=h'^N\prod_{i\in T}r_i\cdot h'(\prod_{i\in T}r_i)^N\in\mathfrak b$. For
$g\in\sqrt{\mathfrak b}$, the image of $g$ lies in $\sqrt{\mathfrak b\mathcal O_w}=\rho\mathcal O_w$, so it is
$\rho a/b$ with $b(w)\neq0$, and $s(bg-a\prod_{i\in T}r_i)\in\mathfrak a$ for some $s$ with $s(w)\neq0$; that is,
$(sb)g\in(\prod_{i\in T}r_i)+\mathfrak a$. The ideal $\sqrt{\mathfrak b}$ is generated by finitely many elements
$g_1,\dots,g_p$; let $H$ be the product of the corresponding factors $sb$. Then
$Hg\in(\prod_{i\in T}r_i)+\mathfrak a$ for every $g\in\sqrt{\mathfrak b}$, and $h=h'H$ has the properties of (2).
\end{proof}
% Lean: (1) is SncRes.Par.primes_below_params, with q_i = chPrime I t (r i), prime_of_regular_params and
% regular_params_spec (the Nakayama argument, with prime_of_not_mem_sq for the primality) and chPrime_le_of_mem; (2) is
% BdRad.exists_den_prod_params: hpr2, radical_span_prod_eq, exists_den_of_mem_map for h', hone, and the finite generation
% of the radical. Size: about 1,800 lines.

\begin{lemma}\label{cv-origin-analytic-system-crossings-components}
Let $K$ be a finite set with finite sets $I_k$, let $\mathfrak p$ and $\mathfrak d$ be ideals of $\mathbf C[x_{k,i}]$
generated by multihomogeneous elements, and let $Q_1,\dots,Q_m$ be the minimal primes over $\mathfrak p+\mathfrak d$.
For indices $\kappa=(\kappa_k)$ and an ideal $Q$ write $Q_\kappa$ for the ideal generated by $Q$ and the
$x_{k,\kappa_k}-1$, and for $z\in\hat C$ with every $z_{k,\kappa_k}\neq0$ put $a_\kappa(z)=(z_{k,i}/z_{k,\kappa_k})_{k,i}$.
Then each $Q_j$ is generated by multihomogeneous elements. Let $\kappa$ be indices, let $z\in\hat C$ be a common zero of
$\mathfrak p$ with $z_{k,\kappa_k}\neq0$ for every $k$, and put $w=a_\kappa(z)$. Suppose that the ideals
$\mathfrak a=\mathfrak p_\kappa\subset\mathfrak b=(\mathfrak p+\mathfrak d)_\kappa$, the point $w$, polynomials
$r_1,\dots,r_e$ and a set $T$ satisfy the hypotheses of Lemma~\ref{cv-origin-analytic-system-crossings-local}, and let
$\mathfrak q_i$ be its primes. Then:
\begin{enumerate}
\item there is an injective map $\nu:T\to\{1,\dots,m\}$ such that, for $i\in T$, $(Q_{\nu(i)})_\kappa=\mathfrak q_i$ and
  no $x_{k,\kappa_k}$ lies in $Q_{\nu(i)}$, and such that every $Q_j$ whose elements vanish at $z$ is $Q_{\nu(i)}$ for
  some $i\in T$;
\item there is a polynomial $F$ with $F(w)\neq0$ such that, for every common zero $z'\in\hat C$ of $\mathfrak p$ with
  every $z'_{k,\kappa_k}\neq0$ and $F(a_\kappa(z'))\neq0$, and every $j$, the elements of $Q_j$ vanish at $z'$ if and
  only if $j=\nu(i)$ for some $i\in T$ with $r_i(a_\kappa(z'))=0$.
\end{enumerate}
\end{lemma}
\leanref{Development/AutomorphicSpectrum/Proofs/MetricProofs/MinimalPrimesOfMultihomogSupStatement.lean}{GlobalTheory.T.minimal_primes_of_multihomog_sup}
\leanref{Development/AutomorphicSpectrum/Proofs/MetricProofs/SncChartRegularAtConePointStatement.lean}{GlobalTheory.T.snc_chart_regular_at_cone_point}
\leanref{Development/AutomorphicSpectrum/Proofs/MetricProofs/SncChartRegularWithBoundaryChoiceStatement.lean}{GlobalTheory.T.snc_chart_regular_with_boundary_choice}
\begin{proof}
For a multihomogeneous $f$ of multidegree $a$ and $z'\in\hat C$ with every $z'_{k,\kappa_k}\neq0$,
$f(z')=\prod_kz'^{a_k}_{k,\kappa_k}\,f(a_\kappa(z'))$, and the coordinates of $a_\kappa(z')$ at the $(k,\kappa_k)$ are
$1$. So for an ideal $Q$ generated by multihomogeneous elements, the elements of $Q$ vanish at $z'$ if and only if those
of $Q_\kappa$ vanish at $a_\kappa(z')$. Let $\delta$ be the substitution $x_{k,\kappa_k}\mapsto1$; then $f-\delta(f)$
lies in the ideal generated by the $x_{k,\kappa_k}-1$, for every $f$.

\emph{Homogeneity.} For each $k$, grade $\mathbf C[x_{k,i}]$ by the degree in the variables of the block $k$. The ideal
$\mathfrak p+\mathfrak d$ is homogeneous for each of these gradings, so its minimal primes are \cite[Tag 00JU]{Sta}. An
ideal that is homogeneous for each of them contains the multihomogeneous components of its elements, which are obtained
by taking homogeneous components block by block; so it is generated by multihomogeneous elements.

(1) Let $i\in T$. By Lemma~\ref{cv-origin-analytic-system-crossings-local}(1), $\mathfrak q_i$ is a prime that contains
$\mathfrak b$. For $f\in\mathfrak p+\mathfrak d$, $\delta(f)\in\mathfrak b\subset\mathfrak q_i$; so
$\delta^{-1}(\mathfrak q_i)$ is a prime containing $\mathfrak p+\mathfrak d$, and it contains a minimal prime $Q$ over
$\mathfrak p+\mathfrak d$. No $x_{k,\kappa_k}$ lies in $Q$, as $\delta(x_{k,\kappa_k})=1$, and $Q_\kappa\subset\mathfrak q_i$,
as the $x_{k,\kappa_k}-1$ lie in $\mathfrak a\subset\mathfrak q_i$. By
Lemma~\ref{cv-origin-analytic-system-crossings-dehomogenize}, $Q_\kappa$ is prime; it contains $\mathfrak b$, hence
$\mathfrak a$, and $Q_\kappa\subset\mathfrak q_i\subset\mathfrak m_w$. By the hypothesis on $T$, $r_j\in Q_\kappa\subset\mathfrak q_i$
for some $j\in T$, so $j=i$ by Lemma~\ref{cv-origin-analytic-system-crossings-local}(1), and then $\mathfrak q_i\subset Q_\kappa$ by
the same lemma. So $Q_\kappa=\mathfrak q_i$; let $\nu(i)$ be the index of $Q$. If $\nu(i)=\nu(i')$, then
$\mathfrak q_i=\mathfrak q_{i'}$ contains $r_{i'}$, so $i=i'$.

Let $Q_j$ be a minimal prime whose elements vanish at $z$. No $x_{k,\kappa_k}$ lies in $Q_j$, as $z_{k,\kappa_k}\neq0$, so
$(Q_j)_\kappa$ is prime by Lemma~\ref{cv-origin-analytic-system-crossings-dehomogenize}. It contains $\mathfrak b$, hence
$\mathfrak a$, and it lies in $\mathfrak m_w$, as its elements vanish at $w$ by the first paragraph. So
$r_i\in(Q_j)_\kappa$ for some $i\in T$, and $(Q_{\nu(i)})_\kappa=\mathfrak q_i\subset(Q_j)_\kappa$. For a multihomogeneous
$f\in Q_{\nu(i)}$, $\delta(f)\in(Q_{\nu(i)})_\kappa\subset(Q_j)_\kappa$, so $f\in Q_j$ by
Lemma~\ref{cv-origin-analytic-system-crossings-dehomogenize}. So $Q_{\nu(i)}\subset Q_j$, and the two are equal, as
$Q_j$ is minimal.

(2) For each $i\in T$ there is a polynomial $s_i$ with $s_i(w)\neq0$ such that every element of $\mathfrak q_i$ vanishes at
every common zero $w'$ of $\mathfrak a$ with $s_i(w')\neq0$ and $r_i(w')=0$. Indeed, let $f_1,\dots,f_p$ generate
$\mathfrak q_i$. The image of $f_l$ in $\mathcal O_w$ is $\rho_ia/b$, where $\rho_i$ is the image of $r_i$ and
$b(w)\neq0$; so $s(bf_l-ar_i)\in\mathfrak a$ for some $s$ with $s(w)\neq0$, and $f_l(w')=0$ at every such $w'$ with
$(sb)(w')\neq0$. Let $s_i$ be the product of these factors $sb$. For $j$ not of the form $\nu(i)$, the elements of $Q_j$ do not all vanish at $z$, by (1), so those of
$(Q_j)_\kappa$ do not all vanish at $w$, and there is $g_j\in(Q_j)_\kappa$ with $g_j(w)\neq0$; for $j=\nu(i)$ put
$g_j=1$. Let $F=\prod_{i\in T}s_i\prod_{j=1}^mg_j$, so that $F(w)\neq0$. Let $z'$ be as in (2), and put $w'=a_\kappa(z')$,
a common zero of $\mathfrak a$ by the first paragraph. If the elements of $Q_j$ vanish at $z'$, then those of
$(Q_j)_\kappa$ vanish at $w'$; as $g_j(w')\neq0$, $j=\nu(i)$ for some $i\in T$, and $r_i\in\mathfrak q_i=(Q_j)_\kappa$
vanishes at $w'$. Conversely, if $j=\nu(i)$ and $r_i(w')=0$, then, as $s_i(w')\neq0$, the elements of
$\mathfrak q_i=(Q_j)_\kappa$ vanish at $w'$, and those of $Q_j$ vanish at $z'$.
\end{proof}
% Lean: MinHom.isMultiHomogIdeal_of_mem_minimalPrimes_sup (cyc_isHomogeneous_of_mem_minimalPrimes, for each block, and
% cyc_isMultiHomogIdeal_of_isHomogeneous); (1) is SncRes.Sig.components_of_params (exists_minimal_below,
% affine_chart_ideal_prime_and_saturated, regular_system_of_parameters_local_ring_alt, exists_compIdeal_eq); (2) is the
% proof of SncRes.NearComp.comp_eq_near (exists_den_chPrime for s_i, the g k, zero_iff), whose open set is the set of
% points of the chart at which F(coord) != 0. Size: about 1,900 lines.

\begin{lemma}\label{cv-origin-analytic-system-crossings-mult}
In the setting of Lemma~\ref{cv-origin-analytic-system-crossings-components}, with its map $\nu$, let $\rho_i$ be the
image of $r_i$ in $\mathcal O_w$, and for a polynomial $f$ let $f^\kappa$ be obtained from $f$ by putting
$x_{k,\kappa_k}=1$ for every $k$. Let $G$ be a multihomogeneous polynomial with complex coefficients such that every
common zero $z'\in\hat C$ of $\mathfrak p$ with $G(z')=0$ is a common zero of $\mathfrak d$. For $1\le j\le m$ let
$\Lambda_j$ be the set of integers $n\ge0$ for which there is a polynomial $s\notin Q_j$ with $sG\in Q_j^n+\mathfrak p$.
Then there are a unit $u$ of $\mathcal O_w$ and integers $m_i\ge0$ for $i\in T$ such that the image of $G^\kappa$ in
$\mathcal O_w$ is $u\prod_{i\in T}\rho_i^{m_i}$, and $\Lambda_{\nu(i)}=\{0,1,\dots,m_i\}$ for every $i\in T$.
\end{lemma}
\leanref{Development/AutomorphicSpectrum/Proofs/MetricProofs/SncChartMultihomogOrderAtPointStatement.lean}{GlobalTheory.T.snc_chart_multihomog_order_at_point}
\leanref{Development/AutomorphicSpectrum/Proofs/MetricProofs/UnitTimesParameterMonomialStatement.lean}{GlobalTheory.T.unit_times_parameter_monomial}
\leanref{Development/AutomorphicSpectrum/Proofs/MetricProofs/SymbolicPowerFromAffineChartStatement.lean}{GlobalTheory.T.symbolic_power_from_affine_chart}
\leanref{Development/AutomorphicSpectrum/Proofs/MetricProofs/SymbolicPowerInAffineChartStatement.lean}{GlobalTheory.T.symbolic_power_in_affine_chart}
\begin{proof}
We use the facts on $\mathcal O_w$ recalled at the start of the proof of
Lemma~\ref{cv-origin-analytic-system-crossings-local}: a polynomial maps to a unit if and only if it does not vanish at
$w$; every element is the image of a quotient $a/b$ of polynomials with $b(w)\neq0$; a polynomial that maps to $0$ is
sent into $\mathfrak a$ by multiplication with some $s$ with $s(w)\neq0$; $\mathcal O_w$ is a domain; and the $\rho_i$ are
primes, pairwise not associated. Let $\mathfrak q_i$ be the primes of that lemma.

\emph{Primes that contain $G^\kappa$.} Let $\mathfrak q$ be a prime with $\mathfrak a\subset\mathfrak q\subset\mathfrak m_w$
and $G^\kappa\in\mathfrak q$. Let $w'$ be a common zero of $\mathfrak q$, and $z'$ the point with coordinates $w'$. As
$\mathfrak a$ contains the $x_{k,\kappa_k}-1$, $w'_{k,\kappa_k}=1$, so $z'\in\hat C$ and $a_\kappa(z')=w'$. By the first
paragraph of the proof of Lemma~\ref{cv-origin-analytic-system-crossings-components}, $z'$ is a common zero of
$\mathfrak p$, and $G(z')=G^\kappa(w')=0$. So $z'$ is a common zero of $\mathfrak d$, and the elements of $\mathfrak b$
vanish at $w'$. By the Nullstellensatz \cite[Tag 00FV]{Sta} they lie in the radical of $\mathfrak q$, which is
$\mathfrak q$. So $\mathfrak b\subset\mathfrak q$, and $r_i\in\mathfrak q$ for some $i\in T$.

\emph{The monomial.} Let $g$ be the image of $G^\kappa$ in $\mathcal O_w$. The preimage of a prime $\mathfrak P$ of
$\mathcal O_w$ is a prime $\mathfrak q$ with $\mathfrak a\subset\mathfrak q\subset\mathfrak m_w$. If $g=0$, the preimage of
the zero ideal contains $G^\kappa$, hence some $r_i$ with $i\in T$, and $\rho_i=0$, which is impossible for a prime. So
$g\neq0$. If $g\in\mathfrak P$, the preimage of $\mathfrak P$ contains $G^\kappa$, hence some $r_i$ with $i\in T$, and
$\prod_{i\in T}\rho_i\in\mathfrak P$. As the radical of $g\mathcal O_w$ is the intersection of the primes that contain it,
$g$ divides $(\prod_{i\in T}\rho_i)^N$ for some $N$. In a domain a divisor of a product of primes is a unit times a
product of some of the factors, by induction on their number: if $gc=p_1\cdots p_n$, then $p_n$ divides $g$ or $c$, and
one cancels $p_n$. So $g=u\prod_{i\in T}\rho_i^{m_i}$ with $u$ a unit and integers $m_i\ge0$.

\emph{The exponents in the chart.} Let $i\in T$ and $n\ge0$. We show that there is $s\notin\mathfrak q_i$ with
$sG^\kappa\in\mathfrak q_i^n+\mathfrak a$ if and only if $n\le m_i$. Let $s$ be such a polynomial. The ideal
$\mathfrak q_i$ maps into $\rho_i\mathcal O_w$, and $\mathfrak a$ to $0$, so the image $\bar s$ of $s$ has
$\bar sg\in\rho_i^n\mathcal O_w$, and $\rho_i$ does not divide $\bar s$, as $s\notin\mathfrak q_i$. Put
$c=\bar su\prod_{j\in T\setminus\{i\}}\rho_j^{m_j}$; then $c\rho_i^{m_i}=\bar sg$. The prime $\rho_i$ does not divide $c$, as
it divides neither $\bar s$, nor the unit $u$, nor the primes $\rho_j$ not associated with it. If $n>m_i$, then
$\rho_i^{m_i+1}$ divides $c\rho_i^{m_i}$, and $\rho_i$ divides $c$ after cancelling $\rho_i^{m_i}$ in the domain
$\mathcal O_w$, a contradiction. So $n\le m_i$. Conversely, let $n\le m_i$, and write $u$ as the image of $a_0/b_0$ with
$b_0(w)\neq0$. Then $b_0G^\kappa-a_0\prod_{j\in T}r_j^{m_j}$ maps to $0$, so
$c'(b_0G^\kappa-a_0\prod_{j\in T}r_j^{m_j})\in\mathfrak a$ for some $c'$ with $c'(w)\neq0$. Put $s=c'b_0$. Its image is a
unit, so $s\notin\mathfrak q_i$, and
$sG^\kappa=c'(b_0G^\kappa-a_0\prod_{j\in T}r_j^{m_j})+c'a_0\prod_{j\in T}r_j^{m_j}\in\mathfrak a+\mathfrak q_i^n$, as
$r_i\in\mathfrak q_i$.

\emph{Homogenization.} Let $i\in T$ and $Q=Q_{\nu(i)}$. By
Lemma~\ref{cv-origin-analytic-system-crossings-components}(1), $Q$ is generated by multihomogeneous elements, no
$x_{k,\kappa_k}$ lies in $Q$, and $Q_\kappa=\mathfrak q_i$; and $\mathfrak p_\kappa=\mathfrak a$. Let $\delta$ be the
substitution $x_{k,\kappa_k}\mapsto1$, so that $\delta(f)=f^\kappa$; as $f-\delta(f)$ lies in the ideal generated by the
$x_{k,\kappa_k}-1$, $\delta$ maps $Q^n+\mathfrak p$ into $Q_\kappa^n+\mathfrak p_\kappa$. Let $a$ be the multidegree of
$G$, and $n\ge0$. We show that there is $s\notin Q$ with $sG\in Q^n+\mathfrak p$ if and only if there is
$s'\notin Q_\kappa$ with $s'G^\kappa\in Q_\kappa^n+\mathfrak p_\kappa$. With the previous paragraph this gives
$\Lambda_{\nu(i)}=\{0,\dots,m_i\}$.

Let $s\notin Q$ with $sG\in Q^n+\mathfrak p$. Some multihomogeneous component $s_b$ of $s$, of multidegree $b$, is not in
$Q$, as $s$ is their sum. The ideal $Q^n+\mathfrak p$ is generated by multihomogeneous elements, so it contains the
component of multidegree $b+a$ of $sG$, which is $s_bG$. So $s_b^\kappa G^\kappa=\delta(s_bG)\in
Q_\kappa^n+\mathfrak p_\kappa$, and $s_b^\kappa\notin Q_\kappa$ by
Lemma~\ref{cv-origin-analytic-system-crossings-dehomogenize}.

Conversely, let $s'\notin Q_\kappa$ with $s'G^\kappa\in Q_\kappa^n+\mathfrak p_\kappa$. Put $\xi=\prod_kx_{k,\kappa_k}$,
let $L=\mathbf C[x_{k,i}][1/\xi]$, and let $\psi:\mathbf C[x_{k,i}]\to L$ be the homomorphism with
$\psi(x_{k,i})=x_{k,i}/x_{k,\kappa_k}$. It kills the $x_{k,\kappa_k}-1$ and maps a multihomogeneous $f$ of multidegree
$b$ to $f\prod_kx_{k,\kappa_k}^{-b_k}$. So $\psi(Q_\kappa)\subset QL$, $\psi(\mathfrak p_\kappa)\subset\mathfrak pL$ and
$\psi(G^\kappa)=\psi(G)=G\prod_kx_{k,\kappa_k}^{-a_k}$, and $\psi(s')G\in(Q^n+\mathfrak p)L$. Write
$\psi(s')=s_0/\xi^e$ with a polynomial $s_0$. Then $s_0G\in(Q^n+\mathfrak p)L$, so $\xi^Ns_0G\in Q^n+\mathfrak p$ for some
$N\ge0$. Put $s=\xi^Ns_0$. As $Q$ is prime and contains no $x_{k,\kappa_k}$, $\xi\notin Q$. If $s_0\in Q$, let
$\Phi:L\to\mathbf C[x_{k,i}]/Q_\kappa$ be the homomorphism that sends each $x_{k,i}$ to its class, which sends $\xi$ to
$1$. Then $\Phi\circ\psi$ is the quotient map, as both send $x_{k,i}$ to its class, so the class of $s'$ is
$\Phi(s_0)\Phi(\xi)^{-e}=0$, and $s'\in Q_\kappa$, a contradiction. So $s\notin Q$, and $sG\in Q^n+\mathfrak p$.
\end{proof}
% Lean: SncRes.Mono.exists_unit_mul_monomial (= snc_chart_multihomog_order_at_point): exists_unit_mul_monomial_of_primes
% (hcontr, hg0 with the zero ideal, hrad, and eq_unit_mul_monomial through exists_le_associated_of_dvd_prod_primes) with
% ideal_sup_le_of_mem for the primes containing G^kappa (vanishingIdeal_zeroLocus_eq_radical, conePts_of_chartIdeal);
% SncRes.Par.exponent_eq_order_polynomial (= unit_times_parameter_monomial; le_of_mul_pow_mem_span_pow,
% exists_mul_mem_of_algebraMap_locRing_eq_zero); SncRes.Dehom.relation_dehom (part, part_mul_mem, map_deh_le,
% affine_chart_ideal_prime_and_saturated) and SncRes.Homog.relation_homogenises (Localization.Away x, chartSub, the lift
% Phi to the quotient by ideal Q p). The Lean states the set Lambda through SncRes.ordAlong, its supremum. Size: about 900
% lines.

\begin{lemma}\label{cv-origin-analytic-system-crossings}
Let $\Omega$ be a field of characteristic zero, let $e\ge0$, let $y\in\Omega^{\sqcup_{k\in K}I_k}$ be a point whose
kernel is an arithmetic subvariety $X$ of dimension $e+1$, and let $I$ be a multihomogeneous ideal of
$R=\mathbf Z[x_{k,i}]$, not contained in the kernel of $y$, that satisfies condition (3) of
Definition~\ref{cv-origin-reference} for $y$ with $e$ in place of $D$. Let $B_I\subset\hat X(\mathbf C)$ be the set of
points at which every element of $I$ vanishes. Then there are a compact complex manifold $\tilde X$ of dimension $e$
with maps $\lambda_x$ as in Definition~\ref{cv-origin-analytic}(3) for $y$, and subsets $W_1,\dots,W_m$ of $\tilde X$,
with the following property. For every $x\in\tilde X$ there are a chart $\psi:V\to\mathbf C^e$ of $\tilde X$ at $x$
with $V\subset U_x$, and a partial map $\sigma$ from $\{1,\dots,m\}$ to $\{1,\dots,e\}$, injective where it is defined,
such that, with $\mu(x')=\prod_l\psi(x')_{\sigma(l)}$, the product over the $l$ at which $\sigma$ is defined:
\begin{enumerate}
\item $W_l\cap V$ is the set of $x'\in V$ with $\psi(x')_{\sigma(l)}=0$ if $\sigma(l)$ is defined, and is empty
  otherwise;
\item for $x'\in V$, $\lambda_x(x')\in B_I$ if and only if $\mu(x')=0$;
\item for every multihomogeneous polynomial $G$ with complex coefficients that vanishes on $B_I$, there is a
  holomorphic function $v$ on $V$ with $G(\lambda_x(x'))=v(x')\mu(x')$ for $x'\in V$;
\item there are such $G$ and $v$ with $v(x)\neq0$.
\end{enumerate}
\end{lemma}
\leanref{Development/WeilBound/CoreStatements/PluriharmonicCharts/EvalFam.lean}{PositivityTheory.Xgp.c_bm_exists_boundary.exists_res}
\begin{proof}
The hypothesis that $I$ is not contained in the kernel of $y$ is not used. Let $P$ be the kernel of $y$, and let
$P_{\mathbf C}$ and $I_{\mathbf C}$ be the ideals of $\mathbf C[x_{k,i}]$ generated by $P$ and by $I$; both are
generated by multihomogeneous elements. For indices $\kappa$ write $Q_\kappa$,
$f^\kappa$ and $a_\kappa$ as in Lemmas~\ref{cv-origin-analytic-system-crossings-dehomogenize}
and~\ref{cv-origin-analytic-system-crossings-components}, so that $(P_{\mathbf C})_\kappa$ is the ideal $I_\kappa$ of
Lemma~\ref{cv-origin-analytic-system-crossings-regular}. Let $Q_1,\dots,Q_m$ be the minimal primes over
$P_{\mathbf C}+I_{\mathbf C}$; by Lemma~\ref{cv-origin-analytic-system-crossings-components} they are generated by
multihomogeneous elements.

\emph{The manifold.} Let $\tilde X$ be the image of $\hat X(\mathbf C)$ in $\prod_k\mathbf P^{I_k}(\mathbf C)$, with the
open sets $\tilde X_\kappa$ and the maps $a_\kappa$ to the affine coordinates of the proof of
Lemma~\ref{cv-origin-model-resolution}. That proof uses its hypothesis only through
Lemma~\ref{cv-origin-model-resolution-regular}, whose conclusion Lemma~\ref{cv-origin-analytic-system-crossings-regular}
gives here. So it shows the following. The space $\tilde X$ is a compact complex manifold of dimension $e$. For
$x\in\tilde X$ choose $\kappa(x)$ with $x\in\tilde X_{\kappa(x)}$, which is then a chart of $y$; with
$U_x=\tilde X_{\kappa(x)}$ and $\lambda_x(x')$ the point with coordinates $a_{\kappa(x)}(x')$, the maps $\lambda_x$ are as
in Definition~\ref{cv-origin-analytic}(3). The map $a_\kappa$ is a homeomorphism from $\tilde X_\kappa$ onto the set of
common zeros of $I_\kappa$, and every polynomial in the coordinates $a_\kappa$ is holomorphic on $\tilde X_\kappa$.
Finally, if $x\in\tilde X_\kappa$, the images of polynomials $r_1,\dots,r_e$ generate the maximal ideal of the local ring
of $\mathbf C[x_{k,i}]/I_\kappa$ at $a_\kappa(x)$, and $O$ is an open neighbourhood of $x$, then
$x'\mapsto(r_i(a_\kappa(x')))_i$ restricts to a chart of $\tilde X$ on an open neighbourhood of $x$ in
$O\cap\tilde X_\kappa$, by Lemma~\ref{cv-origin-model-resolution-coords}.

\emph{The boundary in a chart.} Let $W_l$ be the set of $x\in\tilde X$ such that the elements of $Q_l$ vanish at
$\lambda_x(x)$; as $Q_l$ is generated by multihomogeneous elements, this does not change when $\lambda_x(x)$ is replaced
by a point of its orbit. Let $D_\kappa$ be the radical of $(P_{\mathbf C}+I_{\mathbf C})_\kappa$. For
$x'\in\tilde X_\kappa$, the point $z$ with coordinates $a_\kappa(x')$ is a common zero of $P_{\mathbf C}$ with
$z_{k,\kappa_k}=1$; by the first paragraph of the proof of Lemma~\ref{cv-origin-analytic-system-crossings-components},
$z\in B_I$ if and only if the elements of $(P_{\mathbf C}+I_{\mathbf C})_\kappa$, or equivalently those of $D_\kappa$,
vanish at $a_\kappa(x')$.

\emph{The chart.} Let $x_0\in\tilde X$, $\kappa=\kappa(x_0)$ and $w=a_\kappa(x_0)$. By
Lemma~\ref{cv-origin-analytic-system-crossings-regular}, the local ring $\mathcal O_w$ of $\mathbf C[x_{k,i}]/I_\kappa$ at
$w$ is regular of dimension $e$. Lemma~\ref{cv-origin-analytic-system-crossings-params}, applied with $J=I$, for which
$J_\kappa=(P_{\mathbf C}+I_{\mathbf C})_\kappa$, gives polynomials $r_1,\dots,r_e$ and a set $T$.
Lemma~\ref{cv-origin-analytic-system-crossings-local}, applied to $I_\kappa\subset(P_{\mathbf C}+I_{\mathbf C})_\kappa$,
gives a polynomial $h$ with $h(w)\neq0$, $h\prod_{i\in T}r_i\in D_\kappa$ and $hg\in(\prod_{i\in T}r_i)+I_\kappa$ for
every $g\in D_\kappa$. Lemma~\ref{cv-origin-analytic-system-crossings-components}, applied to $\mathfrak p=P_{\mathbf C}$,
$\mathfrak d=I_{\mathbf C}$ and $z=\lambda_{x_0}(x_0)$, gives $\nu$ and $F$. Let $V\subset\tilde X_\kappa=U_{x_0}$ be an
open neighbourhood of $x_0$ on which $h(a_\kappa(x'))\neq0$ and $F(a_\kappa(x'))\neq0$, and on which
$\psi(x')=(r_i(a_\kappa(x')))_i$ is a chart. Let $\sigma$ be the partial map from $\{1,\dots,m\}$ to $\{1,\dots,e\}$
with $\sigma(\nu(i))=i$ for $i\in T$, undefined outside the image of $\nu$. As $\nu$ is injective, $\sigma$ is well
defined and injective where it is defined, and $\mu(x')=\prod_{i\in T}r_i(a_\kappa(x'))$ for $x'\in V$.

(1) Let $x'\in V$. By Lemma~\ref{cv-origin-analytic-system-crossings-components}(2), $x'\in W_l$ if and only if
$l=\nu(i)$ for some $i\in T$ with $r_i(a_\kappa(x'))=0$, that is, if and only if $\sigma(l)$ is defined and
$\psi(x')_{\sigma(l)}=0$.

(2) Let $x'\in V$ and $t=a_\kappa(x')$. If $\lambda_{x_0}(x')\in B_I$, then $h\prod_{i\in T}r_i\in D_\kappa$ vanishes at
$t$, and $h(t)\neq0$, so $\mu(x')=0$. Conversely, let $\mu(x')=0$ and $g\in D_\kappa$, and write
$hg=c\prod_{i\in T}r_i+b$ with $b\in I_\kappa$. As $t$ is a common zero of $I_\kappa$, $h(t)g(t)=0$, so $g(t)=0$; hence
$\lambda_{x_0}(x')\in B_I$.

(3) Let $G$ be a multihomogeneous polynomial with complex coefficients that vanishes on $B_I$, and let $t$ be a common
zero of $D_\kappa$. It is a common zero of $I_\kappa$, so $t=a_\kappa(x')$ for some $x'\in\tilde X_\kappa$, and the point
with coordinates $t$ lies in $B_I$ by the paragraph on the boundary in a chart. As the coordinates of $t$ at the
$(k,\kappa_k)$ are $1$, $G^\kappa(t)=0$. By the Nullstellensatz \cite[Tag 00FV]{Sta}, the radical ideal $D_\kappa$ is the
intersection of the maximal ideals that contain it, and each of these is the ideal of the polynomials vanishing at one of
its common zeros, as its residue field is $\mathbf C$; so $G^\kappa\in D_\kappa$. Write $hG^\kappa=c\prod_{i\in T}r_i+b$ with $b\in I_\kappa$. For $x'\in V$, the point
$\lambda_{x_0}(x')$ has coordinates $a_\kappa(x')$, a common zero of $I_\kappa$, so
$G(\lambda_{x_0}(x'))=G^\kappa(a_\kappa(x'))=v(x')\mu(x')$ with $v=c(a_\kappa)/h(a_\kappa)$, which is holomorphic on $V$.

(4) Let $T_0$ be larger than the total degree of $h\prod_{i\in T}r_i$, and let $G$ be obtained from
$h\prod_{i\in T}r_i$ by multiplying each monomial by $\prod_kx_{k,\kappa_k}^{T_0-d_k}$, where $d_k$ is the degree of the
monomial in the variables of the block $k$. Then $G$ is multihomogeneous of multidegree $(T_0,\dots,T_0)$, each of its
monomials contains every $x_{k,\kappa_k}$, and $G(z)=\prod_kz_{k,\kappa_k}^{T_0}\cdot(h\prod_{i\in T}r_i)(a_\kappa(z))$
if every $z_{k,\kappa_k}\neq0$. So $G$ vanishes on $B_I$: at the $z\in B_I$ with some $z_{k,\kappa_k}=0$ because of its
monomials, and at the others because the elements of $P_{\mathbf C}+I_{\mathbf C}$ vanish at $z$, so that those of
$(P_{\mathbf C}+I_{\mathbf C})_\kappa$ and of its radical $D_\kappa$, among them $h\prod_{i\in T}r_i$, vanish at
$a_\kappa(z)$. For $x'\in V$, $G(\lambda_{x_0}(x'))=h(a_\kappa(x'))\mu(x')$, and $v=h\circ a_\kappa$ is holomorphic on
$V$ with $v(x_0)=h(w)\neq0$.

\emph{The multiplicities.} The charts also carry the orders of vanishing of multihomogeneous forms along the $W_l$,
which (1) to (4) do not use. Let $G$ be a multihomogeneous polynomial with complex coefficients whose zeros in
$\hat X(\mathbf C)$ lie in $B_I$, and for each $l$ let $\Lambda_l$ be the set of integers $n\ge0$ for which there is a
polynomial $s\notin Q_l$ with $sG\in Q_l^n+P_{\mathbf C}$. Let $x_0\in\tilde X$, with $\kappa$, $w$, the $r_i$, $T$ and
$\nu$ as in the paragraph on the chart. By Lemma~\ref{cv-origin-analytic-system-crossings-mult}, applied with
$\mathfrak p=P_{\mathbf C}$ and $\mathfrak d=I_{\mathbf C}$, each $\Lambda_{\nu(i)}$ with $i\in T$ is $\{0,\dots,m_{\nu(i)}\}$
for an integer $m_{\nu(i)}$, and the image of $G^\kappa$ in $\mathcal O_w$ is a unit times
$\prod_{i\in T}r_i^{m_{\nu(i)}}$. So the exponent at $r_i$ depends only on $G$ and on the component $Q_{\nu(i)}$, and not on
$x_0$.
\end{proof}
% Lean: exists_res. The AlgRes R on M = Pts (pC (mCoordIdeal x)) comes from snc_resolution_compact_embedded,
% snc_resolution_transition_polynomial, snc_resolution_cover_and_cone (with zero_bd, the paragraph on the boundary in a
% chart) and snc_resolution_local_chart_polynomial, as for Lemma cv-origin-model-resolution, with
% snc_resolution_regular_local_ring (Lemma cv-origin-analytic-system-crossings-regular) and
% snc_resolution_boundary_normal_crossing = SncRes.Par.exists_sncData_of_isSncModel (Lemmas
% cv-origin-analytic-system-crossings-params, -local and -components, through snc_chart_parameters_at_cone_point,
% exists_radical_locally_eq_span_prod, snc_chart_regular_at_cone_point and snc_chart_regular_with_boundary_choice) for
% sncAt and comp_eq. The atlas is R0.exists_atlas for the copy R0 with bd = top, a good atlas for R (hA'); Rb has the
% lifts at idx. The chart psi is hA'.2 at the generators r of sncAt on O meet {h != 0}; sigma' is the partial inverse of
% d.sigma, and locMono psi sigma' is the product of the r over d.J (prod_reindex_option). (1) is comp_eq; (2) is zero_bd
% with mem_bd, le_span and coord_mem_zero; (3) is vanishingIdeal_zeroLocus_eq_radical with R.image and le_span; (4) is
% exists_homogenize. The field mult of R is the paragraph on the multiplicities: the second part of
% exists_sncData_of_isSncModel (Lemma cv-origin-analytic-system-crossings-mult), with m k = SncRes.ordAlong (compIdeal k) f;
% its hypothesis that f is not in pC P is not used, and the conclusion of exists_res does not use the field. The
% hypothesis hJP (I not contained in the kernel of y) is passed to snc_resolution_boundary_normal_crossing, whose proof
% exists_sncData_of_isSncModel never uses it (hJn); nondeg of Rb comes from hP. Size: about 6,500 lines.

\begin{lemma}\label{cv-origin-analytic-system-section}
Let $\Omega$ be a field, let $y\in\Omega^{\sqcup_{k\in K}I_k}$, let $\tilde X$ with maps $\lambda_x$ be as in
Definition~\ref{cv-origin-analytic}(3) for $y$, of dimension $e$, and let $B\subset\hat C$ be a set with $cz\in B$ and
$\bar z\in B$ for every $z\in B$ and every $c\in(\mathbf C^\times)^K$. Let $W_1,\dots,W_m$ be subsets of $\tilde X$,
and suppose that every $x\in\tilde X$ has a chart $\psi$ on $V\subset U_x$ and a partial map $\sigma$ with
properties (1) to (4) of Lemma~\ref{cv-origin-analytic-system-crossings}, with $B$ in place of $B_I$. Then there are a
holomorphic line bundle $\bar L$ on $\tilde X$ with a smooth Hermitian metric and a holomorphic section $s$ of
$\bar L$ with $\|s\|\le1$, whose divisor is $W_1+\dots+W_m$ for smooth hypersurfaces $W_l$ crossing normally, such
that $s(x)=0$ if and only if $\lambda_x(x)\in B$, and $\|s(x')\|=\|s(x)\|$ whenever $\lambda_{x'}(x')$ lies in the
orbit of $\overline{\lambda_x(x)}$ under $(\mathbf C^\times)^K$.
\end{lemma}
\leanref{Development/WeilBound/CoreStatements/PluriharmonicCharts/HoloOnLocMono.lean}{PositivityTheory.Xgp.c_bm_exists_boundary.assemble}
\begin{proof}
For each $b\in\tilde X$ choose $\psi_b$ on $V_b$, $\sigma_b$ and $\mu_b$ as in the hypothesis, and $G_b$,
multihomogeneous of some multidegree $a_b\in\mathbf Z_{\ge0}^K$, and $v_b$ as in (4). The open sets
$V'_b=\{x'\in V_b:v_b(x')\neq0\}$ cover $\tilde X$, as $b\in V'_b$, so finitely many of them, for $b$ in a finite
set $T$, cover the compact space $\tilde X$. Let $\mathcal G$ be the finite family of the $G_b$ and of the
polynomials $\bar G_b$ with conjugate coefficients, for $b\in T$, and write $a_G$ for the multidegree of
$G\in\mathcal G$. Each $G\in\mathcal G$ vanishes on $B$, as $\bar G_b(z)=\overline{G_b(\bar z)}$ and $B$ is stable
under conjugation. With $h_a$ as in Definition~\ref{cv-origin-analytic-system-pullback}(1), put
\[
 Q(z)=\sum_{G\in\mathcal G}|G(z)|^2/h_{a_G}(z)^2\qquad(z\in\hat C).
\]
As $G(cz)=\prod_kc_k^{a_{G,k}}G(z)$, $Q(cz)=Q(z)$ for $c\in(\mathbf C^\times)^K$; and $Q(\bar z)=Q(z)$, as
$\mathcal G$ is stable under conjugation of the coefficients.

\emph{Local forms.} For $b\in\tilde X$ and $G\in\mathcal G$, (3) gives a holomorphic function $w_{b,G}$ on $V_b$
with $G(\lambda_b(x'))=w_{b,G}(x')\mu_b(x')$. Put $S_b(x')=\sum_G|w_{b,G}(x')|^2/h_{a_G}(\lambda_b(x'))^2$, a smooth
function on $V_b$ with $Q(\lambda_b(x'))=|\mu_b(x')|^2S_b(x')$. For $b\in T$, $w_{b,G_b}=v_b$: the two agree where
$\mu_b\neq0$, a dense open subset of $V_b$, and both are continuous. So $S_b>0$ on $V'_b$.

\emph{Overlaps.} Let $b,b'\in T$ and $x'\in V'_b\cap V'_{b'}$, and for each $k$ choose $i_k$ with
$\lambda_{b'}(x')_{k,i_k}\neq0$. Near $x'$, $c_k=\lambda_b(\cdot)_{k,i_k}/\lambda_{b'}(\cdot)_{k,i_k}$ is holomorphic
and nonvanishing, and $\lambda_b=c\lambda_{b'}$ by the compatibility of the lifts. Put
$\rho=\prod_kc_k^{a_{b,k}}w_{b',G_b}/v_b$ and $\rho'=\prod_kc_k^{-a_{b',k}}w_{b,G_{b'}}/v_{b'}$, holomorphic near
$x'$. As $v_b\mu_b=G_b(\lambda_b)=\prod_kc_k^{a_{b,k}}G_b(\lambda_{b'})=\prod_kc_k^{a_{b,k}}w_{b',G_b}\mu_{b'}$, we have
$\mu_b=\rho\mu_{b'}$, and likewise $\mu_{b'}=\rho'\mu_b$. So $\rho\rho'=1$ where $\mu_b\neq0$, hence near $x'$ by
continuity, and $\rho$ does not vanish there. As $Q(\lambda_b)=Q(\lambda_{b'})$,
$|\mu_b|^2S_b=|\mu_{b'}|^2S_{b'}$, so $S_{b'}=|\rho|^2S_b$ where $\mu_{b'}\neq0$, and near $x'$ by continuity.

\emph{The bundle and the section.} The function $\nu(x')=Q(\lambda_x(x'))$, for any $x$ with $x'\in U_x$, does not
depend on $x$, by the compatibility of the lifts, and it is continuous. Let $M$ be a bound for it on the compact
space $\tilde X$, and put $c_0=\log(\sqrt M+1)$. For $x\in\tilde X$ choose $b(x)\in T$ with $x\in V'_{b(x)}$, and on
$V'_{b(x)}$ put $\phi_x=c_0-\frac12\log S_{b(x)}$ and $s_x=\mu_{b(x)}$. The $\phi_x$ are smooth, and by the previous
paragraph their differences are near each point of the form $\log|\rho|$, the real part of a branch of $\log\rho$. So
they are the local potentials of a smooth Hermitian metric on a holomorphic line bundle $\bar L$, in which the norm is
$|\cdot|e^{-\phi_x}$ near $x$. Moreover $|s_x|e^{-\phi_x}=e^{-c_0}\nu^{1/2}$ on $V'_{b(x)}$, which does not depend on
$x$; so the holomorphic functions $s_x$ define a holomorphic section $s$ of $\bar L$ with
$\|s\|=e^{-c_0}\nu^{1/2}\le\sqrt M/(\sqrt M+1)\le1$.

\emph{The divisor.} On $V'_{b(x)}$, $s_x=\mu_{b(x)}$ is a product of distinct coordinates of the chart
$\psi_{b(x)}$, and by (1) the zero set there of the coordinate of index $\sigma_{b(x)}(l)$ is the trace of $W_l$. So the
divisor of
$s$ is $W_1+\dots+W_m$, with the $W_l$ smooth hypersurfaces crossing normally. By (2), $s(x)=0$ if and only if
$\lambda_{b(x)}(x)\in B$, that is, $\lambda_x(x)\in B$, as $B$ is stable under $(\mathbf C^\times)^K$. Finally, if
$\lambda_{x'}(x')=c\,\overline{\lambda_x(x)}$, then $\nu(x')=Q(\lambda_{x'}(x'))=Q(\lambda_x(x))=\nu(x)$, so
$\|s(x')\|=\|s(x)\|$.
\end{proof}
% Lean: assemble: the finite subcover t of the V b (isCompact_univ.elim_finite_subcover); GG = G b and its conjugate
% (isMultiHomog_map_conj, eval_map_conj); Qf with nd a = h_a^2 (term_smul, nd_conj); w from hdiv, S, hQS, hwv
% (eqOn_of_coords_ne) and hSpos; Nf continuous with bound M and c0 = log(sqrt M + 1); hover (r and r' with locMono,
% eval_smul_homog and eqOn_of_coords_ne); L with phi = c0 - log(S)/2 (smoothOn_pot, and
% eq_zero_of_mem_all_maximalIdeal_pow for the holomorphic logarithm of r) and s = locMono; IsSNC with multiplicity 1
% on the restricted charts; logNormSec <= 0; bdry from hBc and htor; rW from hQsame and hQconj. Size: about 560 lines.

\begin{lemma}\label{cv-origin-analytic-system-snc}
Let $\Omega$ be a field of characteristic zero, let $e\ge0$, let $y\in\Omega^{\sqcup_{k\in K}I_k}$ be a point whose
kernel is an arithmetic subvariety $X$ of dimension $e+1$, and let $J$ be an ideal of $R=\mathbf Z[x_{k,i}]$ that
satisfies condition (3) of Definition~\ref{cv-origin-reference} for $y$ with $e$ in place of $D$. Let
$B\subset\hat X(\mathbf C)$ be the set of points at which every element of $J$ vanishes, and suppose that $B$ is
closed in $\hat X(\mathbf C)$ and that $cz\in B$ and $\bar z\in B$ for every $z\in B$ and every
$c\in(\mathbf C^\times)^K$. Then there are a compact complex manifold $\tilde X$ of dimension $e$ with maps
$\lambda_x$ as in Definition~\ref{cv-origin-analytic}(3) for $y$, and a holomorphic line bundle $\bar L$ on $\tilde X$
with a smooth Hermitian metric and a holomorphic section $s$ with $\|s\|\le1$, whose divisor is $W_1+\dots+W_m$ for
smooth hypersurfaces $W_l$ crossing normally, such that $s(x)=0$ if and only if $\lambda_x(x)\in B$, and
$\|s(x')\|=\|s(x)\|$ whenever $\lambda_{x'}(x')$ lies in the orbit of $\overline{\lambda_x(x)}$ under
$(\mathbf C^\times)^K$.
\end{lemma}
\leanref{Development/AutomorphicSpectrum/Proofs/MetricProofs/LowerBoundHeightOnBoundaryStatement.lean}{GlobalTheory.T.lowerBound_height_on_boundary}
\begin{proof}
Lemma~\ref{cv-origin-analytic-system-ideal} gives a multihomogeneous ideal $I$, not contained in the kernel of $y$,
that satisfies condition (3) for $y$ with $e$ in place of $D$, and whose set of zeros in $\hat X(\mathbf C)$ is $B$.
Lemma~\ref{cv-origin-analytic-system-crossings}, applied to $I$, gives $\tilde X$ with maps $\lambda_x$, subsets
$W_1,\dots,W_m$ and charts with properties (1) to (4) for $B_I=B$. Lemma~\ref{cv-origin-analytic-system-section},
applied to them and to $B$, gives $\bar L$ and $s$.
\end{proof}
% Lean: Relay.lowerBound_height_on_boundary: exists_homogIdeal, exists_res and assemble, with m = 1. The closedness
% of B is a hypothesis that the proof does not use.

\begin{lemma}\label{cv-origin-analytic-system-positive}
Let $\Omega$ be a field, let $y\in\Omega^{\sqcup_{k\in K}I_k}$, and let $\tilde X$ with maps $\lambda_x$ be as in
Definition~\ref{cv-origin-analytic}(3) for $y$.
\begin{enumerate}
\item For every smooth metrized line bundle $\bar A=(a,h)$, the functions $\vartheta_{\bar A,x}$ of its pullback
  (Definition~\ref{cv-origin-analytic-system-pullback}(2)) are smooth, and near each point of $U_x\cap U_{x'}$ the
  difference $\vartheta_{\bar A,x}-\vartheta_{\bar A,x'}$ is the real part of a holomorphic function.
\item The pullback of $\bar{\mathcal O}(\mathbf 1)$ is positive.
\item Let $k_0\in K$, and suppose that for every $x\in\tilde X$, every $p\in I_{k_0}$ with $\lambda_x(x)_{k_0,p}\neq0$
  and every chart $\chi$ of $\tilde X$ at $x$, the map
  $v\mapsto\big(\lambda_x(\chi^{-1}(v))_{k_0,i}/\lambda_x(\chi^{-1}(v))_{k_0,p}\big)_{i\in I_{k_0}}$ has injective
  derivative at $\chi(x)$. Then the pullback of $\bar{\mathcal O}(e_{k_0})$ is positive.
\end{enumerate}
\end{lemma}
\leanref{Development/AutomorphicSpectrum/Proofs/MetricProofs/LogResolutionFubiniStudyPositive.lean}{GlobalTheory.CompanionTheory.Ar.PB.fs_pos_potential}
\leanref{Development/WeilBound/CoreStatements/CentralCharts/IsPosPotBlock.lean}{AnalyticReference.An.isPosPot_block}
\begin{proof}
(1) The function $h$ is smooth and positive on $\hat C$, and $\lambda_x$ is holomorphic, so $\vartheta_{\bar A,x}$ is
smooth. Let $x''\in U_x\cap U_{x'}$, and for each $k$ choose $i_k$ with $\lambda_{x'}(x'')_{k,i_k}\neq0$. Near $x''$,
$c_k=\lambda_x(\cdot)_{k,i_k}/\lambda_{x'}(\cdot)_{k,i_k}$ is holomorphic and nonvanishing, and $\lambda_x=c\lambda_{x'}$
by the compatibility of the lifts. So $\vartheta_{\bar A,x}-\vartheta_{\bar A,x'}=\sum_ka_k\log|c_k|$, the real part of
$\sum_ka_k\log c_k$ for branches of the logarithms near $x''$.

(2) and (3). For $u\in\mathbf C^I\setminus\{0\}$ and $v\in\mathbf C^I$, the Levi form of
$\frac12\log\|u\|^2=\frac12\log\sum_i|u_i|^2$ at $u$ on $v$ is
$\frac12\big(\|u\|^2\|v\|^2-|\sum_iv_i\bar u_i|^2\big)/\|u\|^4$. By the Cauchy--Schwarz inequality it is nonnegative,
and it is zero only if $v\in\mathbf Cu$. For a holomorphic map $g$ and a smooth function $f$, the Levi form of $f\circ g$
at $t$ on $w$ is that of $f$ at $g(t)$ on $dg_t(w)$, as $g$ is holomorphic.

Let $x\in\tilde X$, let $\chi$ be a chart, let $t=\chi(x')$ with $x'\in U_x$, put $g=\lambda_x\circ\chi^{-1}$ near $t$,
and let $w\neq0$. As $\log h_{\mathbf 1}(z)=\sum_k\frac12\log\|z_k\|^2$, the Levi form of
$\vartheta_{\bar{\mathcal O}(\mathbf 1),x}\circ\chi^{-1}$ at $t$ on $w$ is the sum over $k$ of the Levi forms of
$\frac12\log\|\cdot\|^2$ at $g(t)_k$ on $dg_t(w)_k$. It is nonnegative, and it is zero only if $dg_t(w)=c\,g(t)$ for
some $c\in\mathbf C^K$. In that case choose $p$ with $g(t)_{k,p_k}\neq0$ for every $k$. The affine coordinates
$g_{k,i}/g_{k,p_k}$ have derivative
$\big(dg_t(w)_{k,i}\,g(t)_{k,p_k}-g(t)_{k,i}\,dg_t(w)_{k,p_k}\big)/g(t)_{k,p_k}^2=0$ on $w$. Near $t$ they are the
affine coordinates of $\lambda_{x'}\circ\chi^{-1}$ in the chart $p$, by the compatibility of the lifts, and their
derivative at $t$ is injective by Definition~\ref{cv-origin-analytic}(3). So $w=0$, a contradiction; this proves (2).
For (3), $\log h_{e_{k_0}}(z)=\frac12\log\|z_{k_0}\|^2$, and the same argument for the block $k_0$ alone shows that a
vector $w$ on which the Levi form vanishes has $dg_t(w)_{k_0}=c\,g(t)_{k_0}$ for some $c\in\mathbf C$. Choose $p$ with
$g(t)_{k_0,p}\neq0$; then $\lambda_{x'}(x')_{k_0,p}\neq0$, the ratios $g_{k_0,i}/g_{k_0,p}$ have derivative $0$ on $w$,
and near $t$ they are the ratios of $\lambda_{x'}\circ\chi^{-1}$, whose derivative at $t$ is injective by hypothesis.
So $w=0$.
\end{proof}
% Lean: fs_pos_potential: isPotential_potOf for (1), StepW.fsHerm_isSmoothHerm, and PosFs.isPosPot_fs for (2)
% (levi_comp, levi_fs_pos = FS.levi_fsMulti_pos, fderiv_affCoord_torus and R.immersion through the compatibility of
% the lifts, hev); isPosPot_block for (3) (fs_single_h, the Fubini-Study norm on the block, levi_comp, levi_fs_pos,
% fderiv_affCoord_torus, the hypothesis himm at y' = chart^-1(y) and hev). Sizes: about 920 and 860 lines, mostly
% shared.

\begin{lemma}\label{cv-origin-analytic-system-charts}
Let $\Omega$ be a field, let $y\in\Omega^{\sqcup_{k\in K}I_k}$, and let $\tilde X$ with maps $\lambda_x$ be as in
Definition~\ref{cv-origin-analytic}(3) for $y$, of dimension $e$. Let $s$ be a holomorphic section of a holomorphic
line bundle on $\tilde X$ with a smooth Hermitian metric, with $\|s\|\le1$, whose divisor is $\sum_lm_lW_l$ for smooth
hypersurfaces $W_l$ crossing normally and integers $m_l\ge0$. Let $\partial$ be the zero set of $s$, put
$r=\exp(1)-2\log\|s\|$ on $\tilde X\setminus\partial$, and for $0\le a\le e$ and $t\in\mathbf C^e$ put
$\ell_a(t)=\exp(1)+\sum_{i\le a}\log(1/|t_i|^2)$. Let $\Delta=\{t\in\mathbf C^e:\max_i|t_i|<1\}$.
\begin{enumerate}
\item For every $x\in\tilde X$ and every open set $U\ni x$ there are a chart $\psi$ of $\tilde X$ at $x$ with
  $\psi(x)=0$, an integer $0\le a\le e$ and a number $M$ such that, for every $t\in\Delta$, $\psi^{-1}(t)$ is defined
  and lies in $U$, $\psi^{-1}(t)\in\partial$ if and only if $t_i=0$ for some $i\le a$, and
  $\ell_a(t)\le r(\psi^{-1}(t))+M$ if $\psi^{-1}(t)\notin\partial$.
\item Let $x\in\tilde X$, let $\psi$ and $a$ be as in (1) for $U=U_x$, and let $B\subset\hat C$ be a set such that,
  for $x'\in U_x$, $\lambda_x(x')\in B$ if and only if $x'\in\partial$. Then $\gamma=\lambda_x\circ\psi^{-1}$ is a
  holomorphic map $\Delta\to\hat X(\mathbf C)$, no nonzero tangent vector at $t\in\Delta$ is mapped into the tangent
  space of the orbit of $\gamma(t)$ under $(\mathbf C^\times)^K$, $\gamma(t)$ and $\gamma(t')$ lie in one orbit only if
  $t=t'$, and $\gamma(t)\in B$ if and only if $t_i=0$ for some $i\le a$.
\end{enumerate}
\end{lemma}
\leanref{Development/WeilBound/CoreStatements/CentralCharts/ExistsPermFirst.lean}{AnalyticReference.An.exists_nchart}
\leanref{Development/WeilBound/CoreStatements/CentralCharts/ExistsPermFirst.lean}{AnalyticReference.An.NChart.snc}
\begin{proof}
(1) As the divisor of $s$ has normal crossings, there are a chart $\varphi$ at $x$ with $\varphi(x)=0$ and a
nonvanishing holomorphic function $v$ near $x$ such that, in a local trivialization in which the norm is
$|\cdot|e^{-\phi}$ for a smooth function $\phi$, $s=v\prod_{i\in S}\varphi_i^{\mu_i}$ for a set
$S\subset\{1,\dots,e\}$ and integers $\mu_i\ge1$. After rescaling $\varphi$, the closed polydisc $\Delta_2$ of radius
$2$ lies in its image, and its preimage in the domain of the trivialization. For $t\in\Delta_2$,
$\varphi^{-1}(t)\in\partial$ if and only if $t_i=0$ for some $i\in S$, and off $\partial$
\[
 r(\varphi^{-1}(t))=\exp(1)-2\log|v|-2\sum_{i\in S}\mu_i\log|t_i|+2\phi=\ell'(t)+c(t),
\]
where $\ell'(t)=\exp(1)+\sum_{i\in S}\mu_i\log(4/|t_i|^2)$ and
$c(t)=-2\log|v(\varphi^{-1}(t))|+2\phi(\varphi^{-1}(t))-\sum_{i\in S}\mu_i\log4$ is continuous on $\Delta_2$. Let $M$
be a bound for $|c|$ on $\Delta_2$. Choose $0<\delta\le1$ such that $\varphi^{-1}$ maps the polydisc of radius $\delta$
into $U$, put $a=|S|$, and choose a permutation $\pi$ of $\{1,\dots,e\}$ with $\pi(i)\le a$ if and only if $i\in S$.
Let $\psi$ be the chart with $\psi^{-1}(t)=\varphi^{-1}(L(t))$, where $L(t)_i=\delta t_{\pi(i)}$. For $t\in\Delta$,
$L(t)$ lies in the polydisc of radius $\delta$, so $\psi^{-1}(t)$ is defined and lies in $U$, and
$\psi^{-1}(t)\in\partial$ if and only if $t_{\pi(i)}=0$ for some $i\in S$, that is, $t_j=0$ for some $j\le a$. Off
$\partial$, as $\mu_i\ge1$, $\log(1/|t_j|^2)\ge0$ and $4/(\delta^2|t_j|^2)\ge1/|t_j|^2$,
\[
 \ell_a(t)=\exp(1)+\sum_{i\in S}\log(1/|t_{\pi(i)}|^2)\le\ell'(L(t))\le r(\psi^{-1}(t))+M.
\]

(2) The map $\gamma$ is holomorphic on $\Delta$ with values in $\hat X(\mathbf C)$. Let $t\in\Delta$, put
$x'=\psi^{-1}(t)$, and suppose that $d\gamma_t(v)=c\,\gamma(t)$ for some $c\in\mathbf C^K$. Choose $p$ with
$\lambda_{x'}(x')_{k,p_k}\neq0$ for every $k$; then $\gamma(t)_{k,p_k}\neq0$, as $\gamma(t)$ and $\lambda_{x'}(x')$
lie in one orbit. As in the proof of Lemma~\ref{cv-origin-analytic-system-positive}, the affine coordinates of
$\gamma$ in the chart $p$ have derivative $0$ on $v$ at $t$. Near $t$ they are the affine coordinates of
$\lambda_{x'}\circ\chi^{-1}$, for a chart $\chi$ at $x'$, composed with the change of charts $\chi\circ\psi^{-1}$,
whose derivative is invertible. By Definition~\ref{cv-origin-analytic}(3), $v=0$. If $\gamma(t)$ and $\gamma(t')$ lie
in one orbit, then so do $\lambda_{x_1}(x_1)$ and $\lambda_{x_2}(x_2)$ for $x_1=\psi^{-1}(t)$ and $x_2=\psi^{-1}(t')$,
by the compatibility of the lifts; so $x_1=x_2$ and $t=t'$. Finally, $\gamma(t)\in B$ if and only if
$\psi^{-1}(t)\in\partial$, that is, $t_i=0$ for some $i\le a$.
\end{proof}
% Lean: exists_nchart: exists_snc_adapted_chart_at_point (= LocalForm.exists_local_form: the chart of IsSNC,
% translated and scaled by 3/r, about 230 lines), snc_adapted_chart_section_form (= LocalModel.zero_set_and_bound:
% the zero set and |rW - rLoc S mu| <= M on the closed ball of radius 2, about 170 lines), exists_perm_first and
% scalePerm. (2) is NChart.snc: analyticAt_lift, R.mem, R.immersion with PosFs.fderiv_affCoord_torus and the change
% of charts, R.inj with R.compat. Size: about 520 lines.

\subsection{Singular weights and their regularization}\label{subsection-cv-singular-weights}

\begin{definition}\label{cv-origin-analytic-system-weight}
\begin{enumerate}
\item Let $S$ be a compact complex manifold with a K\"ahler form $\omega$, and let $s$ be a holomorphic section of a
  holomorphic line bundle on $S$ with a smooth Hermitian metric, with $\|s\|\le1$, whose divisor is $\sum_lm_lW_l$ for
  smooth hypersurfaces $W_l$ crossing normally and integers $m_l\ge0$. Let $\partial$ be the zero set of $s$, and
  $r=\exp(1)-2\log\|s\|$. Let $\bar L$ be a holomorphic line bundle on $S$ with a smooth Hermitian metric, and let
  $C\ge0$ and $0<t_0<1$. A \emph{regularized weight} on $\bar L$ with constants $(C,t_0)$ consists of a smooth
  function $\varphi$ on $S\setminus\partial$ and smooth functions $\varphi(t)$ on $S$ for $0<t\le t_0$ that satisfy
  the conditions of Definition~\ref{cv-proportionality-analytic-system}(1) on $\varphi_b$ and $\varphi_b(t)$, with
  $\bar L$, $\varphi$ and $\varphi(t)$ in place of $\bar L_b$, $\varphi_b$ and $\varphi_b(t)$, with $s$ in place of
  $\sigma$, and with these $\omega$, $C$ and $t_0$.
\item Let $n\ge0$ and $0\le a\le n$, and for $t\in\mathbf C^n$ put $\ell_a(t)=\exp(1)+\sum_{i\le a}\log(1/|t_i|^2)$.
  A real function $f$ defined near $0$ in $\mathbf C^n$ is \emph{tame} for $a$ if there are $0<\varepsilon\le1$ and
  $C\ge0$ such that, for every $t$ with $\max_j|t_j|<\varepsilon$ and $t_j\neq0$ for $j\le a$, $f$ is smooth at $t$,
  $|f(t)|\le C\log\ell_a(t)$ and $\|df(t)\|\le C\ell_a(t)^Ce^{\ell_a(t)/2}$.
\end{enumerate}
\end{definition}
\leanref{ReciprocityTheory/AnalyticForms/LeviForms/IsRegFamilyInhabitStatement.lean}{ReciprocityTheory.AnForm.P.IsSingWeight}
\leanref{ReciprocityTheory/AnalyticForms/LeviForms/IsRegFamilyInhabitStatement.lean}{ReciprocityTheory.AnForm.P.IsRegFamily}
\leanref{ReciprocityTheory/AnalyticForms/LeviForms/IsRegFamilyInhabitStatement.lean}{ReciprocityTheory.AnForm.P.IsRegFamilyC2}
\leanref{Development/WeilBound/CoreStatements/CentralCharts/PshAt0.lean}{AnalyticReference.An.Tame}

\begin{lemma}\label{cv-origin-analytic-system-enlarge}
Let $S$, $\omega$, $s$, $\partial$ and $r$ be as in Definition~\ref{cv-origin-analytic-system-weight}(1), and let
$\bar L$ be a holomorphic line bundle on $S$ with a smooth Hermitian metric.
\begin{enumerate}
\item A regularized weight on $\bar L$ with constants $(C,t_0)$, restricted to $0<t\le t'_0$, is a regularized weight
  with constants $(C',t'_0)$ whenever $C\le C'$, $0<t'_0\le t_0$ and $t'_0\le\exp(-1)$.
\item Let $u$ be a smooth function on $S$ such that the metric of $\bar L$ multiplied by $e^{-u}$ is semipositive.
  Then there is $C\ge0$ such that $\varphi=u$ on $S\setminus\partial$ and $\varphi(t)=u$ form a regularized weight on
  $\bar L$ with constants $(C,\exp(-\exp(1)))$.
\end{enumerate}
\end{lemma}
\leanref{Development/AutomorphicSpectrum/Proofs/MetricProofs/SingularWeightMonotoneRegularizationStatement.lean}{GlobalTheory.T.singular_weight_monotone_regularization}
\begin{proof}
Write $\varepsilon_C(t)=C\log\log(1/t)/\log(1/t)$. For $0<t\le\exp(-1)$, $\log(1/t)\ge1$, so
$q(t)=\log\log(1/t)/\log(1/t)\ge0$ and $\varepsilon_C(t)=Cq(t)$. As $\|s\|\le1$, $r\ge\exp(1)$ and $\log r\ge1$.

(1) Let $0<t\le t'\le t'_0$. Then $-C'\log r\le-C\log r\le\varphi\le C\le C'$, $\varphi\le\varphi(t)\le C\le C'$ and
$\varphi(t)\ge-C\log\log(1/t)\ge-C'\log\log(1/t)$. The upper bounds for $\varphi(t)-\varphi$ that involve
$\varepsilon_C(t)$ and $\varepsilon_C(t)r$ remain true for $C'$, as $\varepsilon_C(t)\le\varepsilon_{C'}(t)$ and
$r\ge0$. Next $C\varepsilon_C(t')=C^2q(t')\le C'^2q(t')=C'\varepsilon_{C'}(t')$. The first Chern form of the metric
multiplied by $e^{-\varphi(t)}$ is at least $-\varepsilon_C(t)\omega\ge-\varepsilon_{C'}(t)\omega$, as $\omega$ is
positive. The smoothness and the convergence as $t\to0$ do not involve the constants, and $0<t'_0<1$.

(2) As $S$ is compact, $|u|\le C$ for some $C\ge1$. Then $-C\log r\le-C\le u\le C$. Let
$0<t\le t'\le\exp(-\exp(1))$. Then $\log(1/t)\ge\exp(1)$, so $\log\log(1/t)\ge1$ and $\varepsilon_C(t)\ge0$. So
$u\le\varphi(t)=u\le C$, $\varphi(t)\ge-C\ge-C\log\log(1/t)$, $\varphi(t)\le u+t^{1/3}+\varepsilon_C(t)$,
$\varphi(t)\le u+t^{1/3}+\varepsilon_C(t)r$ and $\varphi(t)\le\varphi(t')+C\varepsilon_C(t')$. The metric of
$\bar L$ multiplied by $e^{-\varphi(t)}$ is semipositive, so its first Chern form is at least
$-\varepsilon_C(t)\omega$. As $\varphi(t)=\varphi$, the convergence as $t\to0$ holds, and
$0<\exp(-\exp(1))<1$.
\end{proof}
% Lean: singular_weight_monotone_regularization = PB.enlarging_constants_smooth_members: AnMono.isSingWeight_mono,
% isRegFamily_enlarge (epsT_mono, mul_epsT_mono, Giv.isPsh_rsmul_mono) and smooth_member (C0 = |M| + 1,
% isRegFamilyC2_const, isPsh_add_rsmul). The Lean keeps the two constants C5 (IsSingWeight) and C (IsRegFamily),
% with |C| <= C' for the second; the book has one constant C >= 0. Size: about 450 lines.

\begin{definition}\label{cv-origin-analytic-system-weights-family}
Let $S$, $\omega$, $s$, $\partial$ and $r$ be as in Definition~\ref{cv-origin-analytic-system-weight}(1), let $\bar L$ be
a holomorphic line bundle on $S$ with a smooth Hermitian metric, let $\varphi$ be a real function on $S\setminus\partial$,
and let $C$ and $0<t_0<1$ be real numbers. Smooth functions $\varphi(t)$ on $S$, for $0<t\le t_0$, \emph{regularize}
$\varphi$ on $\bar L$ with constants $(C,t_0)$ if they satisfy the conditions of
Definition~\ref{cv-proportionality-analytic-system}(1) in which $\varphi_b(t)$ appears, with $\bar L$, $\varphi$ and
$\varphi(t)$ in place of $\bar L_b$, $\varphi_b$ and $\varphi_b(t)$, with $s$ in place of $\sigma$, and with these
$\omega$, $C$ and $t_0$. So, for $C\ge0$, a regularized weight on $\bar L$ with constants $(C,t_0)$ consists of a smooth
function $\varphi$ on $S\setminus\partial$ such that the metric of $\bar L$ multiplied by $e^{-\varphi}$ is semipositive
on $S\setminus\partial$ and $-C\log r\le\varphi\le C$ there, and of functions $\varphi(t)$ that regularize $\varphi$ on
$\bar L$ with constants $(C,t_0)$.
\end{definition}
\leanref{ReciprocityTheory/AnalyticForms/LeviForms/IsRegFamilyInhabitStatement.lean}{ReciprocityTheory.AnForm.P.IsRegFamily}
\leanref{ReciprocityTheory/AnalyticForms/LeviForms/IsRegFamilyInhabitStatement.lean}{ReciprocityTheory.AnForm.P.IsRegFamilyC2}

\begin{lemma}\label{cv-origin-analytic-system-weights-conjugation}
Let $\Omega$ be a field, let $y\in\Omega^{\sqcup_{k\in K}I_k}$, and let $\tilde X$ with maps $\lambda_x$ be as in
Definition~\ref{cv-origin-analytic}(3) for $y$. Then there is a continuous map $c:\tilde X\to\tilde X$ with $c(c(x))=x$,
such that $\lambda_{c(x)}(c(x))$ lies in the orbit of $\overline{\lambda_x(x)}$ under $(\mathbf C^\times)^K$ for every
$x\in\tilde X$. Moreover:
\begin{enumerate}
\item $c$ is \emph{antiholomorphic}: for charts $\chi$ and $\chi'$ of $\tilde X$, the complex conjugate of
  $\chi'\circ c\circ\chi^{-1}$ is holomorphic where it is defined; so for a holomorphic function $g$ on an open set
  $V\subset\tilde X$, the function $\overline{g\circ c}$ is holomorphic on $c^{-1}(V)$;
\item for every smooth metrized line bundle $\bar A$ invariant under conjugation, with pullback $(\vartheta_x)$
  (Definition~\ref{cv-origin-analytic-system-pullback}(2)), every $x\in\tilde X$ has an open neighbourhood $V$ with
  $c(V)\subset U_{c(x)}$ on which $\vartheta_{c(x)}\circ c-\vartheta_x$ is the real part of a holomorphic function.
\end{enumerate}
\end{lemma}
\leanref{Development/AutomorphicSpectrum/Proofs/MetricProofs/LiftSystemComplexConjugationStatement.lean}{GlobalTheory.T.lift_system_complex_conjugation}
\leanref{Development/AutomorphicSpectrum/Proofs/MetricProofs/ConjugationAntiholomorphicChartsStatement.lean}{GlobalTheory.T.conjugation_antiholomorphic_charts}
\leanref{Development/AutomorphicSpectrum/Proofs/MetricProofs/HolomorphicConjugatePullbackStatement.lean}{GlobalTheory.T.holomorphic_conjugate_pullback}
\leanref{Development/AutomorphicSpectrum/Proofs/MetricProofs/LiftSystemRealStructureStatement.lean}{GlobalTheory.T.lift_system_real_structure}
\begin{proof}
Two points $z,z'\in\hat C$ lie in one orbit if and only if $z_{k,a}z'_{k,b}=z_{k,b}z'_{k,a}$ for all $k$, $a$ and $b$:
given these relations, choose $j$ with $z'_{k,j}\neq0$ and put $\gamma_k=z_{k,j}/z'_{k,j}$; then
$z_{k,a}=\gamma_kz'_{k,a}$ for every $a$, and $\gamma_k\neq0$ as $z_k\neq0$. For $x\in\tilde X$, the point $\overline{\lambda_x(x)}$ lies in
$\hat X(\mathbf C)$, as the elements of the kernel of $y$ have integer coefficients; so it lies in the orbit of
$\lambda_{x'}(x')$ for some $x'$, which is unique by Definition~\ref{cv-origin-analytic}(3). Put $c(x)=x'$. As
$\overline{\lambda_{c(x)}(c(x))}$ lies in the orbit of $\lambda_x(x)$, $c(c(x))=x$.

\emph{Continuity.} As $\tilde X$ is compact, it suffices to show that if a net $(x_\alpha)$ converges to $x$ and
$(c(x_\alpha))$ converges to $x'$, then $x'=c(x)$. For large $\alpha$, $x_\alpha\in U_x$ and $c(x_\alpha)\in U_{x'}$, and
by the compatibility of the lifts $\lambda_{x'}(c(x_\alpha))$ lies in the orbit of
$\lambda_{c(x_\alpha)}(c(x_\alpha))$, hence of $\overline{\lambda_{x_\alpha}(x_\alpha)}$, hence of
$\overline{\lambda_x(x_\alpha)}$. So
$\lambda_{x'}(c(x_\alpha))_{k,a}\overline{\lambda_x(x_\alpha)_{k,b}}=\lambda_{x'}(c(x_\alpha))_{k,b}\overline{\lambda_x(x_\alpha)_{k,a}}$,
and by the continuity of $\lambda_{x'}$ and $\lambda_x$ the same relations hold for $\lambda_{x'}(x')$ and
$\overline{\lambda_x(x)}$. So $\lambda_{x'}(x')$ lies in the orbit of $\overline{\lambda_x(x)}$, hence of
$\lambda_{c(x)}(c(x))$, and $x'=c(x)$.

(1) Let $\chi$ and $\chi'$ be charts, and let $x_0$ be a point of the domain of $\chi$ such that $x_1=c(x_0)$ lies in that
of $\chi'$. Choose $p$ with $\lambda_{x_1}(x_1)_{k,p_k}\neq0$ for every $k$; then $\lambda_{x_0}(x_0)_{k,p_k}\neq0$. Let
$\alpha$ and $\beta$ be the affine coordinates in the chart $p$ of $\lambda_{x_0}\circ\chi^{-1}$ near $\chi(x_0)$ and of
$\lambda_{x_1}\circ\chi'^{-1}$ near $\chi'(x_1)$. They are holomorphic, and the derivative of $\beta$ at $\chi'(x_1)$ is
injective by Definition~\ref{cv-origin-analytic}(3). For $y$ near $x_0$, $\lambda_{x_1}(c(y))$ lies in the orbit of
$\overline{\lambda_{x_0}(y)}$, as in the previous paragraph, so the affine coordinates agree:
$\beta(\chi'(c(y)))=\overline{\alpha(\chi(y))}$. Put $G(t)=\overline{\chi'(c(\chi^{-1}(t)))}$, a continuous map, and
$\beta^*(v)=\overline{\beta(\bar v)}$, a holomorphic map with injective derivative at $G(\chi(x_0))$. Then
$\beta^*(G(t))=\alpha(t)$ near $\chi(x_0)$. A linear projection $\pi$ makes $\pi\circ\beta^*$ a biholomorphism near
$G(\chi(x_0))$, by the holomorphic inverse function theorem; as $G$ is continuous, $G=(\pi\circ\beta^*)^{-1}\circ\pi\circ\alpha$
near $\chi(x_0)$, which is holomorphic. For $g$ holomorphic on $V$ and $t$ near $\chi(x_0)$ with $x_1\in V$,
$\overline{g(c(\chi^{-1}(t)))}=\overline{(g\circ\chi'^{-1})(\overline{G(t)})}$, and $v\mapsto\overline{(g\circ\chi'^{-1})(\bar v)}$
is holomorphic; so $\overline{g\circ c}$ is holomorphic on $c^{-1}(V)$.

(2) Let $\bar A=(a,h)$ and $x\in\tilde X$. Choose $i_k$ with $\lambda_x(x)_{k,i_k}\neq0$, and let $V$ be an open
neighbourhood of $x$ in $U_x\cap c^{-1}(U_{c(x)})$ on which every $\lambda_x(y)_{k,i_k}\neq0$. For $y\in V$, as in the
continuity paragraph, $\lambda_{c(x)}(c(y))=\gamma(y)\overline{\lambda_x(y)}$ for some $\gamma(y)\in(\mathbf C^\times)^K$,
acting blockwise. Then $\overline{\gamma_k(y)}=b_k(y)$, where
$b_k(y)=\overline{\lambda_{c(x)}(c(y))_{k,i_k}}/\lambda_x(y)_{k,i_k}$ is holomorphic on $V$ by (1). As
$h(\bar z)=h(z)$ and $h(\gamma z)=\prod_k|\gamma_k|^{a_k}h(z)$,
\[
 \vartheta_{c(x)}(c(y))-\vartheta_x(y)=\sum_ka_k\log|\gamma_k(y)|=\sum_ka_k\log|b_k(y)|.
\]
Shrinking $V$ so that every $b_k(y)/b_k(x)$ lies in the disc of radius $1$ about $1$, this is the real part of
$\sum_ka_k\big(\log(b_k(y)/b_k(x))+\log|b_k(x)|\big)$, with the principal branch of the logarithm, which is holomorphic on
$V$.
\end{proof}
% Lean: (the involution) Cj.exists_conj_of_isLiftSys: c from the surjectivity of the lifts at cj (lift x x)
% (cj_mem_mcPts), continuity by ultrafilters and compactness with sameOrbit_cross and sameOrbit_of_cross, c (c x) = x
% by injectivity. (1) is Cj.conj_analytic, through T.analytic_of_immersion_composition (Rigid.analyticAt_of_comp_immersion,
% the left inverse of a holomorphic immersion), with affCoord_eq_of_orbit and affCoord_cj; its consequence for g is
% T.holomorphic_conjugate_pullback (Cj.holoOn_conj_comp). (2) is Cj.pot_conj_pluriharmonic for a potential th with
% log h(lift x y) - th x y the real part of a holomorphic G; the book takes th the pullback itself, with G = 0, the case
% used. Size: about 500 lines.

\begin{lemma}\label{cv-origin-analytic-system-weights-tame}
Let $\tilde X$, $s$, $\partial$, $r$ and $\ell_a$ be as in Lemma~\ref{cv-origin-analytic-system-charts}, and let $u$ be a
real function on $\tilde X$ such that every $x\in\tilde X$ has $\psi$, $a$ and $M$ as in
Lemma~\ref{cv-origin-analytic-system-charts}(1), for some open set $U\ni x$, with $u\circ\psi^{-1}$ tame for $a$
(Definition~\ref{cv-origin-analytic-system-weight}(2)). Then:
\begin{enumerate}
\item $u$ is smooth on $\tilde X\setminus\partial$;
\item every $x\in\tilde X$ has a chart $\psi$ at $x$, an open neighbourhood $V$ of $x$ in its domain and a number
  $C'\ge0$ such that $|u(x')|\le C'\log r(x')$ and $\|d(u\circ\psi^{-1})(\psi(x'))\|\le C'r(x')^{C'}e^{r(x')/2}$ for
  $x'\in V\setminus\partial$;
\item there is $C_1\ge0$ with $|u|\le C_1\log r$ on $\tilde X\setminus\partial$.
\end{enumerate}
\end{lemma}
\leanref{Development/WeilBound/CoreStatements/CentralCharts/SmoothOnOfTame.lean}{AnalyticReference.An.smoothOn_of_tame}
\leanref{Development/WeilBound/CoreStatements/CentralCharts/SmoothOnOfTame.lean}{AnalyticReference.An.local_of_tame}
\leanref{Development/WeilBound/CoreStatements/CentralCharts/SmoothOnOfTame.lean}{AnalyticReference.An.size_of_tame}
\begin{proof}
(1) Let $x\notin\partial$, with $\psi$, $a$ and $M$ as in the hypothesis. As $\psi^{-1}(0)=x\notin\partial$,
Lemma~\ref{cv-origin-analytic-system-charts}(1) gives $a=0$. The tameness of $u\circ\psi^{-1}$ at $t=0$ then says that
it is smooth at $0$; as the changes of charts are holomorphic, $u$ is smooth at $x$.

(2) Let $x$, $\psi$, $a$ and $M$ be as in the hypothesis, let $\varepsilon$ and $C$ be as in
Definition~\ref{cv-origin-analytic-system-weight}(2), and put $K=1+\max(M,0)$. Let $V$ be the set of $x'$ in the domain
of $\psi$ with $\max_j|\psi(x')_j|<\varepsilon$. Let $x'\in V\setminus\partial$ and $t=\psi(x')$. Then $t_j\neq0$ for
$j\le a$, as $x'\notin\partial$, so $|u(x')|\le C\log\ell_a(t)$ and
$\|d(u\circ\psi^{-1})(t)\|\le C\ell_a(t)^Ce^{\ell_a(t)/2}$. As $\|s\|\le1$, $r\ge\exp(1)$ and $\log r\ge1$; and
$\ell_a(t)\le r(x')+M$, so $\ell_a(t)\le Kr(x')$ and $\ell_a(t)\le r(x')+K$. Hence
$\log\ell_a(t)\le\log r(x')+\log K\le(1+\log K)\log r(x')$, $\ell_a(t)^C\le K^Cr(x')^C$ and
$e^{\ell_a(t)/2}\le e^{K/2}e^{r(x')/2}$. Put $C'=\max\big(C(1+\log K),CK^Ce^{K/2}\big)$; it is at least $C$, so
$r^C\le r^{C'}$, and the bounds of (2) follow.

(3) Finitely many of the sets $V$ of (2) cover the compact space $\tilde X$; let $C_1$ be the sum of their constants $C'$.
\end{proof}
% Lean: smoothOn_of_tame (NChart.zero_of_not_mem for a = 0), local_of_tame (K = 1 + max M 0, c1 = C (1 + log K),
% c2 = C K^C exp(K/2)) and size_of_tame (a finite subcover, with the sum of the |c x|). Size: about 100 lines.

\begin{lemma}\label{cv-origin-analytic-system-weights-line}
Let $\Omega$ be a field, let $y\in\Omega^{\sqcup_{k\in K}I_k}$, let $\tilde X$ with maps $\lambda_x$ be as in
Definition~\ref{cv-origin-analytic}(3) for $y$, and let $c:\tilde X\to\tilde X$ be a map such that
$\lambda_{c(x)}(c(x))$ lies in the orbit of $\overline{\lambda_x(x)}$ under $(\mathbf C^\times)^K$ for every $x$. Let
$d\in\mathbf Z^K$, and let $f$ be a smooth function on $\tilde X$ with $f\circ c=f$. Then there is a smooth metrized line
bundle $\bar A=(d,h_A)$, invariant under conjugation, with $\log h_A(\lambda_x(x))=\log h_d(\lambda_x(x))+f(x)$ for every
$x\in\tilde X$.
\end{lemma}
\leanref{Development/WeilBound/CoreStatements/CentralCharts/CommonMembers.lean}{AnalyticReference.An.exists_line}
\begin{proof}
Lemma~\ref{cv-origin-correction-extend} gives a smooth metrized line bundle $\bar B=(0,h_B)$ with
$\log h_B(\lambda_x(x))=f(x)$ for every $x$. Put $h'(z)=(h_B(z)h_B(\bar z))^{1/2}$ for $z\in\hat C$. It is smooth and
positive, $h'(\gamma z)=h'(z)$ for $\gamma\in(\mathbf C^\times)^K$, as $|\bar\gamma_k|=|\gamma_k|$ and $\bar B$ has
multidegree $0$, and $h'(\bar z)=h'(z)$. As $h_B$ is constant on orbits,
$h_B(\overline{\lambda_x(x)})=h_B(\lambda_{c(x)}(c(x)))=\exp f(c(x))=\exp f(x)$, so
$\log h'(\lambda_x(x))=(f(x)+f(c(x)))/2=f(x)$. Let $h_A=h_dh'$. Then $\bar A=(d,h_A)$ is a smooth metrized line bundle,
invariant under conjugation as $h_d$ and $h'$ are, and $\log h_A(\lambda_x(x))=\log h_d(\lambda_x(x))+f(x)$.
\end{proof}
% Lean: exists_line: exists_hermLine_log_eq_of_smooth (Lemma cv-origin-correction-extend), Cj.conjAvg B (conjAvg_log,
% conjAvg_conj, conjAvg_smooth), FA.deg0_orbit, and the tensor product with fsHerm J d (Cj.fs_conj). Size: about 400
% lines.

\begin{lemma}\label{cv-origin-analytic-system-weights-singular-disc}
Let $r>0$, let $D=\{\zeta\in\mathbf C:|\zeta|<r\}$ and $\bar D=\{\zeta:|\zeta|\le r\}$, and let $F\subset D$ be a finite set
with $0\notin F$. Let $g$ be a continuous real function on $\bar D\setminus F$ such that:
\begin{enumerate}
\item[(a)] for every $\varepsilon>0$ and every $e\in F$, $g(\zeta)+\varepsilon\log|\zeta-e|\to-\infty$ as $\zeta\to e$;
\item[(b)] $g(a)$ is at most the average of $g$ over the circle $|\zeta-a|=s$, whenever $s>0$ and the closed disc
  $|\zeta-a|\le s$ lies in $D\setminus F$.
\end{enumerate}
Then $g(0)$ is at most the average of $g$ over the circle $|\zeta|=r$, and $g(0)\le g(\zeta)$ for some $\zeta$ with
$|\zeta|=r$.
\end{lemma}
\leanref{ReciprocityTheory/AnalyticForms/PlurisubharmonicFunctions/PshExtA.lean}{ReciprocityTheory.AnForm.A.PlurisubharmonicExtension.smg2}
\leanref{ReciprocityTheory/AnalyticForms/PlurisubharmonicFunctions/PshExtA.lean}{ReciprocityTheory.AnForm.A.PlurisubharmonicExtension.exists_sphere_ge}
\begin{proof}
Write $\mu(h;a,s)$ for the average of a function $h$ over the circle $|\zeta-a|=s$. Let $\Phi$ be the solution of the
Dirichlet problem on $D$ with the values of $g$ on the circle $|\zeta|=r$: it is continuous on $\bar D$, equal to $g$ on
that circle and harmonic on $D$ \cite[Chapter~I, \S4.B.2]{Dem12}. So $\Phi(0)=\mu(g;0,r)$, by the formula (4.6) of
\cite[Chapter~I]{Dem12} for the Poisson kernel at the centre, and $\mu(\Phi;a,s)=\Phi(a)$ whenever the closed disc
$|\zeta-a|\le s$ lies in $D$, by the mean value equality of harmonic functions \cite[Chapter~I, \S4.C.1]{Dem12}.

Suppose that $g(0)>\Phi(0)$, and put $\gamma=g(0)-\Phi(0)$. Let $B$ be a bound for $-\Phi$ on $\bar D$. Put
$\Lambda=|F|\log(2r)-\sum_{f\in F}\log|f|$, $\varepsilon=\gamma/(4(|\Lambda|+1))$ and $\delta=\gamma/(4(r^2+1))$, so that
$\varepsilon\Lambda<\gamma/4$ and $\delta r^2<\gamma/4$. On $\bar D\setminus F$ put
\[
 k(\zeta)=g(\zeta)-\Phi(\zeta)+\varepsilon\sum_{f\in F}\log|\zeta-f|+\delta|\zeta|^2,
\]
a continuous function with $k(0)=\gamma+\varepsilon\sum_{f\in F}\log|f|$.
\begin{enumerate}
\item[(i)] If $|\zeta|=r$, then $g(\zeta)=\Phi(\zeta)$ and $|\zeta-f|\le2r$, so
  $k(\zeta)\le\varepsilon|F|\log(2r)+\delta r^2=k(0)-\gamma+\varepsilon\Lambda+\delta r^2<k(0)$.
\item[(ii)] By (a) there is $\eta>0$ such that
  $g(\zeta)+\varepsilon\log|\zeta-e|<k(0)-B-\delta r^2-\varepsilon|F|\,|\log(2r)|$ for $e\in F$ and $\zeta\in\bar D\setminus F$
  with $|\zeta-e|<\eta$. For such $\zeta$, as $\log|\zeta-f|\le|\log(2r)|$ for every $f$, $k(\zeta)<k(0)$.
\item[(iii)] If the closed disc $|\zeta-a|\le s$, with $s>0$, lies in $D\setminus F$, then $\mu(k;a,s)>k(a)$. Indeed
  $\mu(g;a,s)\ge g(a)$ by (b) and $\mu(\Phi;a,s)=\Phi(a)$; $\mu(\log|\zeta-f|;a,s)\ge\log|a-f|$, as $\log|\zeta-f|$ is
  subharmonic \cite[Chapter~I, (5.12) and (4.12)]{Dem12}; and $\mu(|\zeta|^2;a,s)=|a|^2+s^2$, as
  $|a+se^{i\theta}|^2=|a|^2+s^2+2s\operatorname{Re}(\bar ae^{i\theta})$.
\end{enumerate}
Let $K$ be the set of $\zeta\in\bar D$ at distance at least $\eta/2$ from every point of $F$. It is compact, $k$ is
continuous on it, and $0\in K$ by (ii). Let $\zeta_1\in K$ be a point at which $k$ attains its maximum on $K$, so that
$k(\zeta_1)\ge k(0)$. By (i), $|\zeta_1|<r$, and by (ii), $|\zeta_1-e|\ge\eta$ for every $e\in F$. With
$s=\min(r-|\zeta_1|,\eta/2)/2$, the closed disc $|\zeta-\zeta_1|\le s$ lies in $K$ and in $D\setminus F$, so by (iii)
$k(\zeta_1)<\mu(k;\zeta_1,s)\le k(\zeta_1)$, a contradiction. So $g(0)\le\Phi(0)=\mu(g;0,r)$.

If $g(\zeta)<g(0)$ for every $\zeta$ with $|\zeta|=r$, then, $g$ being continuous on that circle, its maximum there is
less than $g(0)$, and so is $\mu(g;0,r)$, a contradiction.
\end{proof}
% Lean: smg2: poisson_ext (PshPoisson.lean: the Poisson integral poissonInt, its continuity up to the circle and its mean
% value property), gamma, BPhi, L, epsilon and delta as in the text, k = g - Phi + pert F 0 epsilon delta, and mp_core for
% the maximum on K = closedBall minus the balls of radius eta/2; pert_submean uses Mathlib's
% circleAverage_log_norm_sub_const_eq_log_radius_add_posLog (the average of log|zeta - f| is log s + log+(|a - f|/s)) and
% circleAverage_norm_sub_sq. exists_sphere_ge is the last paragraph. Size: about 750 lines with PshPoisson.lean.

\begin{lemma}\label{cv-origin-analytic-system-weights-singular-extension}
Let $n\ge0$, let $S\subset\{1,\dots,n\}$, let $U$ be an open subset of the polydisc
$\Delta=\{z\in\mathbf C^n:\max_i|z_i|<1\}$, and put $A=\{z:z_i=0\text{ for some }i\in S\}$. Let $f$ be a continuous real
function on $U\setminus A$, and $C$ a real number, such that:
\begin{enumerate}
\item[(a)] $f(z)$ is at most the average of $f(z+e^{i\theta}w)$ over $0\le\theta\le2\pi$, whenever $z+\zeta w\in U\setminus A$
  for every $\zeta\in\mathbf C$ with $|\zeta|\le1$;
\item[(b)] $f(z)\le C\log\big(\exp(1)+\sum_{i\in S}\log(1/|z_i|)\big)$ for $z\in U\setminus A$.
\end{enumerate}
Then $f$ is bounded above on $K\setminus A$ for every compact set $K\subset U$.
\end{lemma}
\leanref{ReciprocityTheory/AnalyticForms/PlurisubharmonicFunctions/PshExtB.lean}{ReciprocityTheory.AnForm.A.PlurisubharmonicExtension.part_a}
\leanref{ReciprocityTheory/AnalyticForms/PlurisubharmonicFunctions/PshExtB.lean}{ReciprocityTheory.AnForm.A.PlurisubharmonicExtension.local_bound}
\leanref{ReciprocityTheory/AnalyticForms/PlurisubharmonicFunctions/PshExtB.lean}{ReciprocityTheory.AnForm.A.PlurisubharmonicExtension.slice_step}
\begin{proof}
Write $e_j$ for the $j$-th basis vector of $\mathbf C^n$ and $B(p,\rho)=\{z:\max_i|z_i-p_i|<\rho\}$.

\emph{Slices.} Let $x\in U\setminus A$, $j\in S$ and $r>0$ with $|x_j|<r$, such that $x+\zeta e_j\in U$ for every
$|\zeta|\le r$. Then $f(x)\le f(x+\zeta e_j)$ for some $\zeta$ with $|\zeta|=r$. Indeed, as $x\notin A$, the point
$x+\zeta e_j$ lies in $A$ only if $\zeta=-x_j$. Put $g(\zeta)=f(x+\zeta e_j)$ for $|\zeta|\le r$ and $\zeta\neq-x_j$. It
is continuous, and it satisfies (b) of Lemma~\ref{cv-origin-analytic-system-weights-singular-disc}, by (a) applied with
$w=se_j$. As $U\subset\Delta$, every $|x_i|<1$, so $a_0=\sum_{i\in S\setminus\{j\}}\log(1/|x_i|)\ge0$, and by (b), with
$\lambda=\log(1/|\zeta+x_j|)$,
\[
 g(\zeta)+\varepsilon\log|\zeta+x_j|\le C\log(\exp(1)+a_0+\lambda)-\varepsilon\lambda,
\]
which tends to $-\infty$ as $\lambda\to\infty$, that is, as $\zeta\to-x_j$, for every $\varepsilon>0$. As $x_j\neq0$ and
$|x_j|<r$, Lemma~\ref{cv-origin-analytic-system-weights-singular-disc}, applied with $F=\{-x_j\}$, gives $\zeta$.

\emph{Near a point.} Let $p\in U$. Choose $R>0$ with $B(p,R)\subset U$, and $0<r<R/2$ with $r<|p_i|/4$ for every
$i\in S$ with $p_i\neq0$. Let $L$ be the set of $y$ with $|y_i-p_i|\le3r/2$ for every $i$ and $|y_i|\ge r/2$ for every
$i\in S$. It is compact and contained in $U\setminus A$, so $f\le M$ on $L$ for some $M$. Let $F_0$ be the set of $i\in S$
with $p_i=0$. We show by induction on the number of elements of $F'\subset F_0$: if $x\notin A$, $|x_i-p_i|<3r/2$ for
every $i$, $|x_i|>r/2$ for $i\in F_0\setminus F'$ and $|x_i|<r/2$ for $i\in F'$, then $f(x)\le f(y)$ for some $y\in L$.
If $F'$ is empty, then $x\in L$: for $i\in S\setminus F_0$, $|x_i|\ge|p_i|-3r/2>r/2$. Otherwise let $j\in F'$. For
$|\zeta|\le r$, the point $x+\zeta e_j$ lies in $B(p,R)\subset U$, as $|x_j+\zeta-p_j|<r/2+r<R$. By the paragraph on slices
there is $\zeta$ with $|\zeta|=r$ and $f(x)\le f(x')$, where $x'=x+\zeta e_j$. Then $r/2<|x'_j|<3r/2$, the other
coordinates of $x'$ are those of $x$, and $x'\notin A$; so the induction hypothesis for $F'\setminus\{j\}$ applies to $x'$.
Every $z\in B(p,r/2)\setminus A$ satisfies the hypothesis with $F'=F_0$, so $f(z)\le M$ there.

\emph{Compact sets.} Finitely many of the sets $B(p,r/2)$ of the previous paragraph cover $K$; the sum of the absolute values
of the corresponding $M$ bounds $f$ on $K\setminus A$.
\end{proof}
% Lean: part_a (part (1) of PshExtensionStatement, proved with part (2), the extension of max(u, -kappa), as pshExtension_thm,
% which the text does not need): local_bound (the compact set L, and the induction key over the finite set
% S.filter (p i = 0)), slice_step (with smg2 for F = {-x j}, hgrow from loglog_bound, and exists_sphere_ge), and the finite
% subcover. Size: about 400 lines.

\begin{lemma}\label{cv-origin-analytic-system-weights-singular}
Let $S$, $s$, $\partial$ and $r$ be as in Definition~\ref{cv-origin-analytic-system-weight}(1), let $\bar L$ be a
holomorphic line bundle on $S$ with a smooth Hermitian metric, and let $u$ be a smooth function on $S\setminus\partial$
such that the metric of $\bar L$ multiplied by $e^{-u}$ is semipositive on $S\setminus\partial$ and $|u|\le C_1\log r$
there, for some $C_1$. Then there is $C_5\ge0$ with $-C_5\log r\le u\le C_5$ on $S\setminus\partial$.
\end{lemma}
\leanref{Development/AutomorphicSpectrum/Proofs/MetricProofs/SingularPotentialLogGrowthPshStatement.lean}{GlobalTheory.T.singular_potential_log_growth_psh}
\begin{proof}
We show that every $x\in S$ has an open neighbourhood $V_x$ and a number $B_x$ with $u\le B_x$ on
$V_x\setminus\partial$. As $\|s\|\le1$, $r\ge\exp(1)$ and $\log r\ge1$, so $u\le|C_1|\log r$ on $S\setminus\partial$.

\emph{A chart.} Fix $x$, and let $\vartheta_x$ be a local potential of the metric of $\bar L$ near $x$; by hypothesis
$\vartheta_x+u$ has nonnegative Levi form in every chart, off $\partial$. As the divisor of $s$ has normal crossings,
there are a chart $\psi$ at $x$, a smooth function $\phi$ with $\|s\|=|s_x|e^{-\phi}$ on the domain of $\psi$ for a
holomorphic function $s_x$, a holomorphic function $v$ without zeros there, a set $S_0\subset\{1,\dots,e\}$ and integers
$\mu_i\ge0$ with $s_x=v\prod_{i\in S_0}\psi_i^{\mu_i}$, where $e=\dim S$. Put $a=\psi(x)$, and let $S'$ be the set of
$i\in S_0$ with $a_i=0$ and $\mu_i\ge1$. Choose $\rho>0$ such that the closed polydisc
$\bar P=\{w:\max_i|w_i-a_i|\le\rho\}$ lies in the image of $\psi$, $\psi^{-1}(\bar P)$ lies in the domain of
$\vartheta_x$, and $|w_i|>|a_i|/2$ on $\bar P$ for every $i$ with $a_i\neq0$. Let $M_1$, $M_2$ and $M_3$ be bounds for
$|\vartheta_x\circ\psi^{-1}|$, $|\phi\circ\psi^{-1}|$ and $|\log|v\circ\psi^{-1}||$ on the compact set $\bar P$.

\emph{The growth in the chart.} Let $\Delta$ be the open unit polydisc in $\mathbf C^e$, let $A$ be the set of $z$ with
$z_i=0$ for some $i\in S'$, and for $z\in\Delta\setminus A$ put $w=a+\rho z\in\bar P$ and
$\ell(z)=\sum_{i\in S'}\log(1/|z_i|)\ge0$. For $i\in S_0$ with $\mu_i\ge1$ let $B_i=|\mu_i\log\rho|$ if $a_i=0$ and
$B_i=|\mu_i\log(|a_i|/2)|$ if $a_i\neq0$, and let $B_i=0$ if $\mu_i=0$. Then $w_i^{\mu_i}\neq0$ and
$\log|w_i^{\mu_i}|\ge-B_i-\mu_i\ell(z)$ for every $i\in S_0$: if $a_i=0$ and $\mu_i\ge1$, then
$\log|w_i^{\mu_i}|=\mu_i\log\rho-\mu_i\log(1/|z_i|)$; if $a_i\neq0$, then $|w_i|>|a_i|/2$. So $y=\psi^{-1}(w)$ has
$s_x(y)\neq0$, hence $y\notin\partial$, and
\[
 r(y)=\exp(1)-2\log|v(y)|-2\sum_{i\in S_0}\mu_i\log|w_i|+2\phi(y)\le Q+D\ell(z),
\]
with $Q=\exp(1)+2M_3+2M_2+2\sum_iB_i$ and $D=2\sum_i\mu_i$. Put $K=|Q|+|D|+1\ge1$. Then $r(y)\le K(\exp(1)+\ell(z))$, so
$\log r(y)\le\log K+\log(\exp(1)+\ell(z))\le(|\log K|+1)\log(\exp(1)+\ell(z))$, as $\log(\exp(1)+\ell(z))\ge1$. With
$f(z)=\vartheta_x(y)+u(y)$ and $C_c=|M_1|+|C_1|(|\log K|+1)$ this gives
\[
 f(z)\le|M_1|+|C_1|\log r(y)\le C_c\log\Big(\exp(1)+\sum_{i\in S'}\log(1/|z_i|)\Big)\qquad(z\in\Delta\setminus A).
\]

\emph{The bound.} The function $f$ is continuous on $\Delta\setminus A$, and of class $C^2$ there with nonnegative Levi form,
as $z\mapsto\psi^{-1}(a+\rho z)$ is holomorphic and composing with a holomorphic map composes the Levi form with its
derivative. So $f$ is plurisubharmonic on $\Delta\setminus A$ \cite[Chapter~I, \S5.A, before Theorem~5.8]{Dem12}, and in
particular it satisfies the mean value inequality (a) of
Lemma~\ref{cv-origin-analytic-system-weights-singular-extension} on complex discs \cite[Chapter~I, Definition~5.1 and
Theorem~4.12]{Dem12}. That lemma, applied with $U=\Delta$, $S=S'$ and $C=C_c$, gives $B$ with $f\le B$ on the set of
$z\in\Delta\setminus A$ with $\max_i|z_i|\le1/2$. Let $V_x$ be the set of points of the domain of $\psi$ with $\max_i|\psi_i-a_i|<\rho/2$, and let
$y\in V_x\setminus\partial$. Then $z=(\psi(y)-a)/\rho$ has $\max_i|z_i|<1/2$, and $z\notin A$: if $z_i=0$ with $i\in S'$,
then $\psi_i(y)=a_i=0$ and $\mu_i\ge1$, so $s_x(y)=0$ and $y\in\partial$. So $u(y)=f(z)-\vartheta_x(y)\le B+M_1$; put
$B_x=B+M_1$.

\emph{Compactness.} Finitely many of the $V_x$ cover the compact space $S$; let $C_5$ be $|C_1|$ plus the sum of the
corresponding $|B_x|$. As $\log r\ge1$, $-C_5\log r\le-|C_1|\log r\le u$, and $u\le C_5$.
\end{proof}
% Lean: PB.isSingWeight_of_logGrowth: localBound, a finite subcover (CompactSpace.elim_nhds_subcover) and
% C5 = |C1| + the sum of the |B x|. localBound: the IsSNC chart at x (tm k for the factors, grouped here by coordinate),
% S', rho from hnb, M1, M2, M3, the bounds hk (Bk), hD and hoff, hr (logNormSecLocal), log_le_growth, Cc and hFb; pshPull
% (levi_comp and isPsh_iff_isLeviPsh_thm, the C^2 criterion cited from Demailly) on ball 0 1 minus coordDiv S'; and
% pshExtension_thm part (1) (Lemma cv-origin-analytic-system-weights-singular-extension) on closedBall 0 (1/2); the
% bound Bm + M1 on psi^-1(ball a (rho/2)). Size: about 330 lines here, with the two lemmas before it.

\begin{lemma}\label{cv-origin-analytic-system-weights-flat-extend}
Let $n\ge0$, $c\in\mathbf C^n$ and $r>0$, put $B(c,\rho)=\{z\in\mathbf C^n:\max_i|z_i-c_i|<\rho\}$, and let
$Q\subset\mathbf C^n$ be an open set that contains every $z\in B(c,r)$ with $z_i\neq c_i$ for all $i$. Let $F$ be a
smooth real function on $B(c,r)\cap Q$ whose Levi form is $0$ at every point, with $|F|\le M$ there for some $M$. Then
there is a smooth function $G$ on $B(c,r/2)$ with $G=F$ on $B(c,r/2)\cap Q$.
\end{lemma}
\leanref{Development/AutomorphicSpectrum/Proofs/MetricProofs/BoundedPshCorrectionLocallySmooth.lean}{GlobalTheory.CompanionTheory.Ar.PB.Flat.locExt_core}
\begin{proof}
This is the extension theorem for bounded pluriharmonic functions. The set $A=B(c,r)\setminus Q$ is closed in $B(c,r)$
and contained in the union of the hyperplanes $z_i=c_i$, on which the plurisubharmonic function $\sum_i\log|z_i-c_i|$
is $-\infty$, so it is pluripolar \cite[Chapter~I, Definition~5.22]{Dem12}. As the Levi form of $F$ is $0$, $F$ and $-F$
are plurisubharmonic \cite[Chapter~I, \S5.A, before Theorem~5.8]{Dem12}, so $F$ is pluriharmonic on $B(c,r)\setminus A$,
and it is bounded. So it extends to a pluriharmonic function $G$ on $B(c,r)$ \cite[Chapter~I, Theorem~5.24, applied to
$F$ and $-F$ as in the proof of Corollary~5.25]{Dem12}, which is smooth \cite[Chapter~I, \S5.D]{Dem12}.
\end{proof}
% Lean: locExt_core proves this itself, on the half polydisc, with iterated Poisson integrals: levi_flat_separately_harmonic
% (F is harmonic in each variable on the punctured polydisc where every z i != c i), polydisc_poisson_reproduces_harmonic
% (a bounded continuous separately harmonic function there is the Poisson integral polyPoisson (r/2) of its values on the
% torus of radius r/2), polydisc_poisson_smooth, and dense_coordinate_torus with continuity for the points of Q where some
% z i = c i. Size: about 850 lines, in IntegralAngMeas.lean, PoissonIntegralPuncturedDisc.lean and the files of these four
% statements.

\begin{lemma}\label{cv-origin-analytic-system-weights-flat}
Let $S$, $s$, $\partial$ and $r$ be as in Definition~\ref{cv-origin-analytic-system-weight}(1), let $\bar L$ be a
holomorphic line bundle on $S$ with a smooth Hermitian metric with local potentials $\vartheta_x$, and let $u$ be a smooth
function on $S\setminus\partial$ with $|u|\le C_1\log r$ there, for some $C_1$, such that $\vartheta_x+u$ and
$-(\vartheta_x+u)$ are plurisubharmonic on $S\setminus\partial$. Then there is a smooth function $u'$ on $S$ with
$u'=u$ on $S\setminus\partial$ such that $\vartheta_x+u'$ and $-(\vartheta_x+u')$ are plurisubharmonic on $S$.
\end{lemma}
\leanref{Development/AutomorphicSpectrum/Proofs/MetricProofs/FlatWeightExtends.lean}{GlobalTheory.CompanionTheory.Ar.PB.Flat.flat_weight_extends}
\begin{proof}
\emph{A bound.} Lemma~\ref{cv-origin-analytic-system-weights-singular}, applied to $\bar L$ and $u$, gives $C_5$ with
$u\le C_5$ on $S\setminus\partial$. The functions $-\vartheta_x$ are local potentials of the dual metric on the dual
bundle, and the metric of the dual multiplied by $e^{u}$ is semipositive off $\partial$, as $-\vartheta_x-u$ is
plurisubharmonic there; so the same lemma, applied to the dual and to $-u$, gives $C_6$ with $-u\le C_6$. So
$|u|\le M_0=\max(C_5,C_6)$ on $S\setminus\partial$.

\emph{The boundary.} As the divisor of $s$ has normal crossings, every $x\in S$ has a chart $\psi$ and a set $S_x$ of
indices such that a point $y$ of the domain of $\psi$ lies in $\partial$ if and only if $\psi_i(y)=0$ for some $i\in S_x$.
So $\partial$ is closed, and $S\setminus\partial$ is dense in $S$, as the points of $\mathbf C^e$ with no coordinate equal
to $0$ are dense.

\emph{Local extensions.} Let $x\in S$, with $\psi$ and $S_x$ as above, and put $c=\psi(x)$ and
$Q=\{z:z_i\neq0\text{ for }i\in S_x\}$; write $B(c,\rho)$ for the open polydisc of radius $\rho$ about $c$. Choose $r>0$
such that the closed polydisc $\bar B$ of radius $r$ about $c$ lies
in $\psi$ of the intersection of the domains of $\psi$ and of $\vartheta_x$, and such that $z_i\neq0$ on $\bar B$ for every
$i\in S_x$ with $c_i\neq0$. Then every $z\in B(c,r)$ with $z_i\neq c_i$ for all $i$ lies in $Q$, and $\psi^{-1}(z)\notin\partial$
for $z\in B(c,r)\cap Q$. On $B(c,r)\cap Q$ the function $F=(\vartheta_x+u)\circ\psi^{-1}$ is smooth, its Levi form is $0$, as
those of $F$ and $-F$ are nonnegative, and $|F|\le M+M_0$, where $M$ bounds $|\vartheta_x\circ\psi^{-1}|$ on $\bar B$.
Lemma~\ref{cv-origin-analytic-system-weights-flat-extend} gives a smooth $G$ on $B(c,r/2)$ with $G=F$ on
$B(c,r/2)\cap Q$. On $N_x=\psi^{-1}(B(c,r/2))$ put $E_x=G\circ\psi-\vartheta_x$. It is smooth, and $E_x(y)=u(y)$ for
$y\in N_x\setminus\partial$, as then $\psi(y)\in Q$.

\emph{Gluing.} On $N_x\cap N_{x'}$ the continuous functions $E_x$ and $E_{x'}$ agree off $\partial$, hence everywhere, as
$S\setminus\partial$ is dense and $N_x\cap N_{x'}$ is open. So $u'=E_x$ on $N_x$ defines a smooth function on $S$, with
$u'=u$ on $S\setminus\partial$.

\emph{Plurisubharmonicity.} Let $\chi$ be a chart, $t$ a point of $\chi$ of the domain of $\vartheta_x$ and $w$ a vector,
and suppose that the Levi form of $(\vartheta_x+u')\circ\chi^{-1}$ at $t$ on $w$ is negative. It is continuous in $t$, as $\vartheta_x+u'$ is smooth,
so it is negative on an open neighbourhood of $t$, which contains $\chi(y)$ for some $y\notin\partial$. Near $y$,
$\vartheta_x+u'=\vartheta_x+u$, as $S\setminus\partial$ is open, and the Levi form there is nonnegative, a contradiction.
So $\vartheta_x+u'$ is plurisubharmonic on $S$, and so is $-(\vartheta_x+u')$, by the same argument.
\end{proof}
% Lean: flat_weight_extends: singular_potential_log_growth_psh for (th, u) and for (th.rsmul (-1), -u) (isPotentialRsmul),
% with |u| <= max C5 C6 (hM); section_boundary_closed_and_codense (Bdry.bdry_isClosed_dense_compl, from
% snc_boundary_coordinate_chart and dense_coordinate_torus); bounded_psh_correction_locally_smooth (Flat.locExt: the chart of
% snc_boundary_coordinate_chart, Q, S', r from hball, the bound C + M, locExt_core, E = G o psi - th x);
% smooth_extension_across_closed_dense (SmoothGlue.glue, eq_of_dense_compl); and psh_extension_across_closed_dense
% (AnForm.isPsh_of_dense_compl, continuousOn_levi) for th + u' and its negative. Size: about 400 lines besides
% locExt_core and Lemma cv-origin-analytic-system-weights-singular.

\begin{lemma}\label{cv-origin-analytic-system-weights-regularize-atlas}
Let $S$, $s$, $\partial$ and $r$ be as in Definition~\ref{cv-origin-analytic-system-weight}(1), with $e=\dim S$. For
$\rho>0$ let $\Delta_\rho$ and $\bar\Delta_\rho$ be the open and the closed polydisc of radius $\rho$ about $0$ in
$\mathbf C^e$. For a set $T\subset\{1,\dots,e\}$, integers $\mu_i\ge1$ for $i\in T$, and $y\in\mathbf C^e$ with
$y_i\neq0$ for $i\in T$, put $r_{T,\mu}(y)=\exp(1)+\sum_{i\in T}\mu_i\log(4/|y_i|^2)$. Then there are finitely many charts
$\kappa_1,\dots,\kappa_N$ of $S$, sets $T_j\subset\{1,\dots,e\}$ with integers $\mu_{j,i}\ge1$ for $i\in T_j$, and a number
$M$, such that:
\begin{enumerate}
\item the image of every $\kappa_j$ contains $\bar\Delta_2$;
\item every $x\in S$ has some $j$ with $x$ in the domain of $\kappa_j$ and $\kappa_j(x)\in\Delta_1$;
\item for $y$ in the image of $\kappa_j$, $\kappa_j^{-1}(y)\in\partial$ if and only if $y_i=0$ for some $i\in T_j$;
\item $|r(\kappa_j^{-1}(y))-r_{T_j,\mu_j}(y)|\le M$ for $y\in\bar\Delta_2$ with $\kappa_j^{-1}(y)\notin\partial$.
\end{enumerate}
\end{lemma}
\leanref{Development/AutomorphicSpectrum/Proofs/GeometryProofs/AdaptedAtlasSncSectionStatement.lean}{GlobalTheory.T.adapted_atlas_snc_section}
\begin{proof}
Let $x\in S$. As the divisor of $s$ has normal crossings, after translating and rescaling the chart of its local form, as
in the proof of Lemma~\ref{cv-origin-analytic-system-charts}(1), there are a chart $\varphi_x$ with $\varphi_x(x)=0$ whose
image contains $\bar\Delta_2$, a holomorphic function $s_x$ and a smooth function $\phi$ with $\|s\|=|s_x|e^{-\phi}$ on its
domain, a holomorphic function $v$ without zeros there, a set $T_x$ and integers $\mu_i\ge1$ with
$s_x=v\prod_{i\in T_x}\varphi_{x,i}^{\mu_i}$. As $v$ has no zeros and every $\mu_i\ge1$, a point $y$ of the image lies in
$\varphi_x(\partial)$ if and only if $y_i=0$ for some $i\in T_x$. For $y\in\bar\Delta_2$ with $\varphi_x^{-1}(y)\notin\partial$,
\[
 r(\varphi_x^{-1}(y))-r_{T_x,\mu}(y)=-2\log|v(\varphi_x^{-1}(y))|+2\phi(\varphi_x^{-1}(y))-\sum_{i\in T_x}\mu_i\log4,
\]
as $r=\exp(1)-2\log|v|-2\sum_{i\in T_x}\mu_i\log|\varphi_{x,i}|+2\phi$; so its absolute value is at most
$M_x=2C_1+2C_2+\sum_{i\in T_x}\mu_i\log4$, where $C_1$ and $C_2$ bound $|\log|v\circ\varphi_x^{-1}||$ and
$|\phi\circ\varphi_x^{-1}|$ on the compact set $\bar\Delta_2$. The open sets $\varphi_x^{-1}(\Delta_1)$ cover the compact
space $S$, so finitely many of them do; their charts are the $\kappa_j$, and $M$ is the sum of the corresponding $|M_x|$.
\end{proof}
% Lean: Adapt.exists_adaptedAtlas, which gives IsAdaptedAtlas A B L s for a FinHolAtlas B compatible with A: hgood from
% PB.LocalForm.exists_local_form (the IsSNC chart translated and scaled), hzero and hmono for the boundary, hlogN, hloc and
% hdiff for the difference, the bound 2 C1 + 2 C2 + sum of mu i log 4; CompactSpace.elim_nhds_subcover, atlasOfCompat,
% and the sum of the |M y|. Size: about 170 lines, besides the local form shared with Lemma cv-origin-analytic-system-charts.

\begin{lemma}\label{cv-origin-analytic-system-weights-regularize-gradient}
Let $S$, $s$, $\partial$ and $r$ be as in Definition~\ref{cv-origin-analytic-system-weight}(1), and let
$\kappa_1,\dots,\kappa_N$ be charts of $S$ with property (1) of Lemma~\ref{cv-origin-analytic-system-weights-regularize-atlas},
with the polydiscs $\Delta_\rho$ and $\bar\Delta_\rho$ of that lemma. Let $\varphi$ be a smooth function on $S\setminus\partial$ such that every $x\in S$ has a chart $\psi$, an open neighbourhood
$V$ of $x$ in its domain and a number $c$ with $\|d(\varphi\circ\psi^{-1})(\psi(x'))\|\le c\,r(x')^ce^{r(x')/2}$ for
$x'\in V\setminus\partial$. Then there is $C_6$ such that, for every $C\ge C_6$, every $j$ and every $y\in\Delta_{3/2}$ with
$x'=\kappa_j^{-1}(y)\notin\partial$, $\|d(\varphi\circ\kappa_j^{-1})(y)\|\le C\,r(x')^Ce^{r(x')/2}$.
\end{lemma}
\leanref{Development/AutomorphicSpectrum/Proofs/MetricProofs/AdaptedAtlasSmoothPotentialLocalStatement.lean}{GlobalTheory.T.adapted_atlas_smooth_potential_local}
\begin{proof}
Fix $j$ and $y_0\in\bar\Delta_{3/2}$, and let $\psi$, $V$ and $c$ be as in the hypothesis for $x=\kappa_j^{-1}(y_0)$. The
change of charts $\tau=\psi\circ\kappa_j^{-1}$ is holomorphic near $y_0$; put $K=\|d\tau(y_0)\|+1$. There is a
neighbourhood $W$ of $y_0$ on which $\|d\tau\|<K$ and $\kappa_j^{-1}(W)\subset V$. For $y\in W$ with
$x'=\kappa_j^{-1}(y)\notin\partial$, the chain rule gives
\[
 \|d(\varphi\circ\kappa_j^{-1})(y)\|\le\|d(\varphi\circ\psi^{-1})(\tau(y))\|\,\|d\tau(y)\|\le K\,c\,r(x')^ce^{r(x')/2}.
\]
Let $D_{y_0}=\max(K|c|,|c|)$ and $C\ge D_{y_0}$. As $\|s\|\le1$, $r\ge\exp(1)\ge1$. If $c\ge0$, then $Kc\le C$ and
$r^c\le r^C$; if $c<0$, then $Kc\,r^c\le0\le C\,r^C$. In both cases the right side is at most $C\,r(x')^Ce^{r(x')/2}$.
Finitely many of the sets $W$ cover the compact set $\bar\Delta_{3/2}$; let $D_j$ be the sum of the corresponding
$|D_{y_0}|$, and $C_6$ the sum of the $|D_j|$.
\end{proof}
% Lean: PB.gradBound_of_local_bounds = PositivityTheory.Xgp.t_c_grad_atlas (XCGradAtlas.lean): local_bd (tau = psi o
% (B.chart j)^-1, K = ||D tau y0|| + 1, the chain rule hchain, and mono_bd with max (K |c|) |c|), finite subcovers of
% closedBall 0 (3/2) and the sums of the |D y| and of the |Dj j|; the conclusion is IsGradBound B L s ph C for every
% C >= C5. Size: about 200 lines.

\begin{definition}\label{cv-origin-analytic-system-weights-regularize-global-notation}
Let $e\ge0$. On $\mathbf C^e$ use the maximum norm $|y|=\max_i|y_i|$; norms of derivatives of real functions on open
subsets of $\mathbf C^e$ are operator norms for it.
\begin{enumerate}
\item For $z\in\mathbf C^e$ and $\rho>0$, $\Delta(z,\rho)$ and $\bar\Delta(z,\rho)$ are the polydiscs $\{y:|y-z|<\rho\}$ and
  $\{y:|y-z|\le\rho\}$, and $\Delta_\rho=\Delta(0,\rho)$. For $T\subset\{1,\dots,e\}$, $H_T$ is the set of $y$ with $y_i=0$
  for some $i\in T$. For integers $\mu_i\ge0$ ($i\in T$) and $y\notin H_T$, put
  $r_{T,\mu}(y)=\exp(1)+\sum_{i\in T}\mu_i\log(4/|y_i|^2)$ and $|\mu|=\sum_{i\in T}\mu_i$.
\item Let $t>0$. For a real function $u$ bounded above on $\bar\Delta(z,t)\setminus H_T$, $m_T(u;z,t)$ is its supremum
  there; for a real function $w$ bounded above on $\bar\Delta(z,t)$, $m(w;z,t)$ is its supremum on $\bar\Delta(z,t)$.
\item For a real number $C$ and $0<t<1$, $\varepsilon_C(t)=C\log\log(1/t)/\log(1/t)$.
\item A \emph{mollifier} on $\mathbf C^e$ is a smooth function $\varrho\ge0$ on $\mathbf C^e$ with support in $\Delta_1$ and
  integral $1$ for the Lebesgue measure, such that $\varrho(cw)=\varrho(w)$ for every $w$ and every $c\in\mathbf C$ with
  $|c|=1$. For $t>0$ put $\varrho_t(w)=t^{-2e}\varrho(w/t)$ and $(u*\varrho_t)(z)=\int_{\mathbf C^e}\varrho_t(w)u(z-w)\,dw$,
  for a function $u$ for which the integral is defined.
\end{enumerate}
\end{definition}
\leanref{ReciprocityTheory/AnalyticForms/PlurisubharmonicFunctions/IsPshLocallyIntegrableStatement.lean}{ReciprocityTheory.AnForm.A.coordDiv}
\leanref{ReciprocityTheory/AnalyticForms/LeviForms/IsRegFamilyInhabitStatement.lean}{ReciprocityTheory.AnForm.P.rLoc}
\leanref{ReciprocityTheory/AnalyticForms/LeviForms/RegChart.lean}{ReciprocityTheory.AnForm.P.polyMaxOff}
\leanref{ReciprocityTheory/AnalyticForms/PlurisubharmonicFunctions/IsPshLocallyIntegrableStatement.lean}{ReciprocityTheory.AnForm.A.polyMax}
\leanref{ReciprocityTheory/AnalyticForms/LeviForms/IsRegFamilyInhabitStatement.lean}{ReciprocityTheory.AnForm.P.epsT}
\leanref{ReciprocityTheory/AnalyticForms/PlurisubharmonicFunctions/IsPshLocallyIntegrableStatement.lean}{ReciprocityTheory.AnForm.A.Mollifier}
\leanref{ReciprocityTheory/AnalyticForms/PlurisubharmonicFunctions/IsPshLocallyIntegrableStatement.lean}{ReciprocityTheory.AnForm.A.molli}

\begin{lemma}\label{cv-origin-analytic-system-weights-regularize-global-numbers}
Let $0<t\le t'\le\exp(-\exp(1))$, and let $C$ be a real number.
\begin{enumerate}
\item $\log(1/t)\ge\exp(1)$, $\log\log(1/t)\ge1$, and $0\le\varepsilon_1(t)\le\varepsilon_1(t')\le1$.
\item $\varepsilon_C(t)=C\varepsilon_1(t)$; so $\varepsilon_C(t)\le\varepsilon_C(t')$ if $C\ge0$, and
  $\varepsilon_C(t)\le\varepsilon_{C'}(t)$ if $C\le C'$.
\item If $C\ge1$, then $t\le\varepsilon_C(t)$.
\item $\log b/b\le\log a/a$ whenever $\exp(1)\le a\le b$; and $\varepsilon_C(\tau)\to0$ as $\tau\to0$, so that
  $\varepsilon_C(\tau)\le1$ for all small $\tau>0$.
\end{enumerate}
\end{lemma}
\leanref{Development/AutomorphicSpectrum/Proofs/MetricProofs/EpsilonSchedulePropertiesStatement.lean}{GlobalTheory.T.epsilon_schedule_properties}
\begin{proof}
The function $x\mapsto\log x/x$ has derivative $(1-\log x)/x^2\le0$ for $x\ge\exp(1)$, and tends to $0$ as
$x\to\infty$; this gives (4), as $\log(1/\tau)\to\infty$ when $\tau\to0$. Put $\ell=\log(1/t)$ and $\ell'=\log(1/t')$.
Then $\ell\ge\ell'\ge\exp(1)$, so $\log\ell\ge1$, and $\varepsilon_1(t)=\log\ell/\ell$. As $0\le\log\ell\le\ell-1$,
$0\le\varepsilon_1(t)\le1$, and $\varepsilon_1(t)\le\varepsilon_1(t')$ by (4); this gives (1), and (2) follows. For (3),
$\ell+1\le e^\ell=1/t$, so $t\le1/\ell\le\log\ell/\ell\le C\log\ell/\ell$.
\end{proof}
% Lean: epsilon_schedule_properties = Num.epsT_numerics (log_div_self_anti, epsT_mono, le_epsT, epsT_tendsto,
% epsT_le_one; its last clause, a positive lower bound for finitely many positive numbers, gives t0 below), with
% Giv.log_inv_bounds, Giv.epsT_le_epsT, RegFin.quot_mem_unit and RegFin.quot_mono. Size: about 150 lines.

\begin{lemma}\label{cv-origin-analytic-system-weights-regularize-global-transfer}
Let $C_5$ and $M$ be real numbers. There is $C'_5\ge0$ with the following property. Let $r_1,r_2\ge\exp(1)$ with
$|r_1-r_2|\le M$, and let $p$ and $d$ be real numbers. If $-C_5\log r_1\le p\le C_5$, then $-C'_5\log r_2\le p\le C'_5$;
and if $d\le|C_5|\,r_1^{|C_5|}e^{r_1/2}$, then $d\le C'_5\,r_2^{C'_5}e^{r_2/2}$.
\end{lemma}
\leanref{Development/AutomorphicSpectrum/Proofs/MetricProofs/LogBoundTransferBetweenRadiiStatement.lean}{GlobalTheory.T.log_bound_transfer_between_radii}
\begin{proof}
Put $A=|C_5|$, $M'=|M|$ and $C'_5=A(1+M')^Ae^{M'}$; then $C'_5\ge A$ and $C'_5\ge A(1+M')$. As $r_2\ge1$,
$r_1\le r_2+M'\le r_2(1+M')$, and $\log r_2\ge1$. Suppose that $-C_5\log r_1\le p\le C_5$. Then $p\le A\le C'_5$. If
$C_5\le0$, then $p\ge0\ge-C'_5\log r_2$. If $C_5>0$, then
$\log r_1\le\log r_2+\log(1+M')\le\log r_2+M'\le(1+M')\log r_2$, so $p\ge-C_5(1+M')\log r_2\ge-C'_5\log r_2$. Suppose that
$d\le Ar_1^Ae^{r_1/2}$. Then $r_1^A\le r_2^A(1+M')^A$, $e^{r_1/2}\le e^{r_2/2}e^{M'}$ and $r_2^A\le r_2^{C'_5}$, so
$d\le A(1+M')^Ae^{M'}r_2^{C'_5}e^{r_2/2}=C'_5r_2^{C'_5}e^{r_2/2}$.
\end{proof}
% Lean: log_bound_transfer_between_radii = BoundConv.bound_conv (bound_conv_local, with
% K = |C5| (1 + |M|)^|C5| exp |M|). Size: about 80 lines.

\begin{lemma}\label{cv-origin-analytic-system-weights-regularize-global-refine}
Let $S$, $s$, $\partial$ and $r$ be as in Definition~\ref{cv-origin-analytic-system-weight}(1), and let
$\kappa_1,\dots,\kappa_N$, $T_j$, $\mu_{j,i}$ and $M$ be as in Lemma~\ref{cv-origin-analytic-system-weights-regularize-atlas}.
For each $x\in S$ let $U_x$ be an open neighbourhood of $x$. Let $\varphi$ be a smooth function on $S\setminus\partial$, and
$C_5$ a number with $\|d(\varphi\circ\kappa_j^{-1})(y)\|\le C_5\,r(x')^{C_5}e^{r(x')/2}$ for every $j$ and every
$y\in\Delta_{3/2}$ with $x'=\kappa_j^{-1}(y)\notin\partial$. Then there are charts $\kappa'_1,\dots,\kappa'_{N'}$ of $S$,
sets $T'_j$, integers $\mu'_{j,i}\ge1$ for $i\in T'_j$ and a number $M'$ with the properties (1)--(4) of
Lemma~\ref{cv-origin-analytic-system-weights-regularize-atlas}, and points $x_1,\dots,x_{N'}$ of $S$, such that the domain
of $\kappa'_j$ lies in $U_{x_j}$ and $\|d(\varphi\circ(\kappa'_j)^{-1})(y)\|\le|C_5|\,r(x')^{|C_5|}e^{r(x')/2}$ for every $j$
and every $y\in\Delta_{3/2}$ with $x'=(\kappa'_j)^{-1}(y)\notin\partial$.
\end{lemma}
\leanref{Development/AutomorphicSpectrum/Proofs/MetricProofs/RefineAdaptedAtlasWithGradientStatement.lean}{GlobalTheory.T.refine_adapted_atlas_with_gradient}
\begin{proof}
Let $x\in S$. By (2) there is $j$ with $x$ in the domain of $\kappa_j$ and $z_0=\kappa_j(x)\in\Delta_1$. The set of points
$y$ of the image of $\kappa_j$ with $\kappa_j^{-1}(y)\in U_x$ is open and contains $z_0$, so it contains $\Delta(z_0,\delta)$
for some $\delta>0$. Choose $0<\rho\le1/8$ with $3\rho\le\delta$ and $3\rho<|z_{0,i}|$ for every $i$ with $z_{0,i}\neq0$.
Let $\psi_x$ be the chart $y'\mapsto(\kappa_j(y')-z_0)/\rho$ on the set of points $y'$ of the domain of $\kappa_j$ with
$\kappa_j(y')\in\Delta(z_0,3\rho)$. Its image is $\Delta_3$, which contains $\bar\Delta_2$; its domain lies in $U_x$;
$\psi_x(x)=0$; and $\psi_x^{-1}(w)=\kappa_j^{-1}(\rho w+z_0)$.

Let $T'_x$ be the set of $i\in T_j$ with $z_{0,i}=0$, with the integers $\mu_{j,i}$. For $w\in\Delta_3$ and $i\in T_j$ with
$z_{0,i}\neq0$, $|\rho w_i+z_{0,i}|\ge|z_{0,i}|-3\rho>0$; so, by (3) for $\kappa_j$, $\psi_x^{-1}(w)\in\partial$ if and only
if $w_i=0$ for some $i\in T'_x$. Let $w\in\bar\Delta_2$ with $\psi_x^{-1}(w)\notin\partial$. Then $|\rho w+z_0|<1/4+1$, so (4)
for $\kappa_j$ gives $|r(\psi_x^{-1}(w))-r_{T_j,\mu_j}(\rho w+z_0)|\le M$. For $i\in T_j$ with $z_{0,i}=0$,
$\log(4/|\rho w_i|^2)=\log(4/|w_i|^2)+\log(1/\rho^2)$ with $0\le\log(1/\rho^2)\le\log(4/\rho^2)$; for the other $i\in T_j$,
$\rho\le|\rho w_i+z_{0,i}|\le2$, so $0\le\log(4/|\rho w_i+z_{0,i}|^2)\le\log(4/\rho^2)$. Hence
$|r_{T_j,\mu_j}(\rho w+z_0)-r_{T'_x,\mu_j}(w)|\le c_x=\sum_{i\in T_j}\mu_{j,i}\log(4/\rho^2)$, and (4) holds for $\psi_x$
with the bound $M+c_x$. For $w\in\Delta_{3/2}$, $\rho w+z_0\in\Delta_{3/2}$ and
$d(\varphi\circ\psi_x^{-1})(w)=\rho\,d(\varphi\circ\kappa_j^{-1})(\rho w+z_0)$, whose norm is at most that of
$d(\varphi\circ\kappa_j^{-1})(\rho w+z_0)$, as $\rho\le1$. As $\|s\|\le1$, $r\ge\exp(1)\ge1$, so
$C_5r^{C_5}\le|C_5|\,r^{|C_5|}$, which gives the bound on the derivative for $\psi_x$.

The open sets $\psi_x^{-1}(\Delta_1)$ cover the compact space $S$, so finitely many of them do, for points
$x_1,\dots,x_{N'}$. The charts $\kappa'_j=\psi_{x_j}$, with $T'_j=T'_{x_j}$ and $M'$ the sum of the numbers $|M+c_{x_j}|$,
have the properties stated.
\end{proof}
% Lean: refine_adapted_atlas_with_gradient = Giv.refine_adapted_atlas: refine_local_chart (the chart
% (c.restrOpen _ hs).transHomeomorph (pa_aff rho^-1 z0), rho from hfar with rho <= 1/8 and 3 rho <= r, the boundary
% hbd with S.filter (z0 i = 0), rLoc_rescale_bound for c_x, fderiv_rescale_norm_le for the derivative), then
% CompactSpace.elim_nhds_subcover and atlasOfCompat, with the sum of the |Mf|. Size: about 450 lines.

\begin{lemma}\label{cv-origin-analytic-system-weights-regularize-global-extend}
Let $e\ge0$, $T\subset\{1,\dots,e\}$, $R>0$ and $B\ge0$, let $u$ be a continuous plurisubharmonic function on
$\Delta_R\setminus H_T$ with $u\le B$ there, and let $\kappa$ be a real number. Then there is a plurisubharmonic function
$w$ on $\Delta_R$, equal to $\max(u,-\kappa)$ on $\Delta_R\setminus H_T$, such that for every $x\in\Delta_R$ and every
$\delta>0$ some $y\in\Delta_R\setminus H_T$ has $|y-x|<\delta$ and $w(x)\le\max(u(y),-\kappa)+\delta$. Let
$z\in\mathbf C^e$ and $s'>0$ with $\bar\Delta(z,s')\subset\Delta_R$. Then:
\begin{enumerate}
\item the function $\sigma\mapsto m(w;z,e^\sigma)$ is convex and nondecreasing on $\sigma\le\log s'$;
\item for $0<s\le s'$, $m_T(u;z,s)\le m(w;z,s)$, and $m(w;z,s_0)\le\max(m_T(u;z,s),-\kappa)$ for $0<s_0<s$;
\item for $0<s<s'$, $m(w;z,s)\le\max(m_T(u;z,s),-\kappa)$.
\end{enumerate}
\end{lemma}
\leanref{ReciprocityTheory/AnalyticForms/LeviForms/RegChartA.lean}{ReciprocityTheory.AnForm.P.RegularChartScaling.ext_exists}
\leanref{ReciprocityTheory/AnalyticForms/PlurisubharmonicFunctions/PshMoll.lean}{ReciprocityTheory.AnForm.A.polyMaxLogConvex_thm}
\leanref{ReciprocityTheory/AnalyticForms/LeviForms/RegChartB.lean}{ReciprocityTheory.AnForm.P.RegularChartScaling.pmo_le_polyMax}
\leanref{ReciprocityTheory/AnalyticForms/LeviForms/RegChartB.lean}{ReciprocityTheory.AnForm.P.RegularChartScaling.polyMax_le_pmo}
\leanref{Development/AutomorphicSpectrum/Proofs/MetricProofs/LocalPshRegularizationLogBound.lean}{GlobalTheory.CompanionTheory.Ar.Giv.polyMax_le_pmoff_same}
\begin{proof}
The set $H_T$ is closed, and it is pluripolar, as the plurisubharmonic function $\sum_{i\in T}\log|y_i|$ is $-\infty$ on it
\cite[Chapter~I, (5.12) and Definition~5.22]{Dem12}. The function $v=\max(u,-\kappa)$ is plurisubharmonic on
$\Delta_R\setminus H_T$ \cite[Chapter~I, Theorem~5.6]{Dem12} and at most $\max(B,-\kappa)$, so it extends to a
plurisubharmonic function $w$ on $\Delta_R$ \cite[Chapter~I, Theorem~5.24]{Dem12}. Let $x\in\Delta_R$ and $\delta>0$, and
let $\beta$ be a Euclidean ball about $x$ contained in $\Delta(x,\delta)\cap\Delta_R$. As $w$ is subharmonic, $w(x)$ is at
most the mean value of $w$ over $\beta$ \cite[Chapter~I, Theorem~4.12 and (5.3)]{Dem12}. As $H_T$ has measure zero, this is
the mean value of $v$ over $\beta\setminus H_T$, which is at most the supremum of $v$ there; so some $y\in\beta\setminus H_T$
has $w(x)\le v(y)+\delta$.

(1) Choose $R''>s'$ with $\Delta(z,R'')\subset\Delta_R$. By \cite[Chapter~I, Corollary~5.14]{Dem12}, the supremum
$m^\circ(\rho)$ of $w$ over the open polydisc $\Delta(z,\rho)$ is a convex nondecreasing function of $\log\rho$ for
$0<\rho<R''$, hence a continuous one. As $w$ is upper semicontinuous, it is bounded above on compact sets, and
$m^\circ(\rho)\le m(w;z,\rho)\le m^\circ(\rho')$ for $\rho<\rho'<R''$; so $m(w;z,\rho)=m^\circ(\rho)$ for $\rho<R''$, which
gives (1).

(2) As $u\le v=w$ off $H_T$, $m_T(u;z,s)\le m(w;z,s)$. Let $x\in\bar\Delta(z,s_0)$ and $0<\delta\le s-s_0$. The point $y$
given by the first paragraph lies in $\bar\Delta(z,s)\setminus H_T$, so $w(x)\le\max(m_T(u;z,s),-\kappa)+\delta$. Letting
$\delta\to0$ and taking the supremum over $x$ gives the second inequality.

(3) By (2), $m(w;z,e^\sigma)\le\max(m_T(u;z,s),-\kappa)$ for $\sigma<\log s$, and by (1) the left side is continuous at
$\sigma=\log s$, as $\log s<\log s'$.
\end{proof}
% Lean: ext_exists: w = extT S R u kappa, the limsup extension limsupExt of max(u(R y), -kappa) across coordDiv S on the
% unit polydisc, rescaled; it is plurisubharmonic by pshExtension_thm, part (2) (PlurisubharmonicExtension.vext_isPsh),
% which the Lean proves itself, and the limsup gives the points y (frequently_lt_of_lt_limsup). The book cites
% Demailly's extension theorem and identifies the extension with the limsup by the mean value inequality. (1) is
% polyMaxLogConvex_thm, proved for closed polydiscs through polydiscMax_exp_chord_bound; (2) is pmo_le_polyMax and
% polyMax_le_pmo; (3) is polyMax_le_pmoff_same. Size: about 1,700 lines, of which about 1,500 prove the extension.

\begin{lemma}\label{cv-origin-analytic-system-weights-regularize-global-oscillation}
Let $e\ge0$, $T\subset\{1,\dots,e\}$ with integers $\mu_i\ge0$ for $i\in T$, and $5/4<R\le2$. Let $C_5$ be a number and
$u$ a continuous plurisubharmonic function on $\Delta_R\setminus H_T$ with $-C_5\log r_{T,\mu}\le u\le C_5$ there, and let
$C_1\ge1$. Then there are $C\ge1$ and $\ell_0$ such that every $t>0$ with $\log(1/t)\ge\ell_0$ has $\log(1/t)\ge\exp(1)$
and $C_1t\le(R-5/4)/2$, and every $z\in\Delta_{5/4}$ has
\[
 -C\log\log(1/t)\le m_T\big(u;z,t/(8C_1)\big)\quad\text{and}\quad
 m_T(u;z,C_1t)-m_T\big(u;z,t/(4C_1)\big)\le\varepsilon_C(t).
\]
\end{lemma}
\leanref{ReciprocityTheory/AnalyticForms/LeviForms/RegChartC.lean}{ReciprocityTheory.AnForm.P.RegularChartScaling.oscGen}
\begin{proof}
\emph{The sign of $C_5$.} The point $x$ with $x_i=1/2$ for every $i$ lies in $\Delta_R\setminus H_T$, and
$\log r_{T,\mu}(x)\ge1$, as every term of $r_{T,\mu}$ is nonnegative on $\bar\Delta_2$. So $-C_5\log r_{T,\mu}(x)\le C_5$,
that is $C_5(1+\log r_{T,\mu}(x))\ge0$, and $C_5\ge0$.

\emph{Constants.} Put $R'=R-5/4$, $\Lambda=\log C_1+\log(8C_1)>0$, $K_1=\exp(1)+|\mu|(\log4+2\log(8C_1)+2)$,
$C=6C_5\Lambda+2C_5+1\ge1$ and $\ell_0=\max\big(\exp(1),K_1,\log(2C_1/R'),2(\log C_1-\log R')\big)$. Let $t>0$ with
$\ell=\log(1/t)\ge\ell_0$. Then $\ell\ge\exp(1)$, $\log\ell\ge1$ and $C_1t=C_1e^{-\ell}\le R'/2$. Fix $z\in\Delta_{5/4}$;
then $\bar\Delta(z,R')\subset\Delta_R$.

\emph{A lower bound.} Put $s_8=t/(8C_1)$ and $s_4=t/(4C_1)$. Let $x_i=z_i+s_8z_i/|z_i|$ if $z_i\neq0$, and $x_i=s_8$ if
$z_i=0$. Then $x\in\bar\Delta(z,s_8)$ and $|x_i|\ge s_8$ for every $i$, so $x\notin H_T$ and
\[
 r_{T,\mu}(x)\le\exp(1)+|\mu|\log(4/s_8^2)=\exp(1)+|\mu|\big(\log4+2\ell+2\log(8C_1)\big)\le K_1\ell\le\ell^2,
\]
as $\ell\ge1$ and $K_1\le\ell$. Hence
\[
 m_T(u;z,s_8)\ge u(x)\ge-C_5\log r_{T,\mu}(x)\ge-2C_5\log\ell.\tag{$*$}
\]
As $C\ge2C_5$ and $\log\ell\ge0$, this gives the first inequality of the lemma.

\emph{Convexity.} Put $\kappa=2C_5\log\ell+1$, and let $w$ be the function of
Lemma~\ref{cv-origin-analytic-system-weights-regularize-global-extend} for $u$, $B=C_5$ and $\kappa$. Every value of $w$
is at most $\max(C_5,-\kappa)=C_5$, by the points $y$ of that lemma. Let $g(\sigma)=m(w;z,e^\sigma)$, a convex
nondecreasing function for $\sigma\le\log R'$ by part (1) of that lemma. Put $a=\log s_8=-\ell-\log(8C_1)$,
$b=\log(C_1t)=\log C_1-\ell$ and $\sigma_0=\log R'$. Then $b-a=\Lambda$, and $\sigma_0-b\ge\ell/2$, as
$\ell\ge2(\log C_1-\log R')$. By parts (1) and (2) of that lemma and $(*)$, $g(b)\ge g(a)\ge m_T(u;z,s_8)\ge-2C_5\log\ell$; and
$g(\sigma_0)\le C_5$. By convexity,
\[
 \frac{g(b)-g(a)}{\Lambda}\le\frac{g(\sigma_0)-g(b)}{\sigma_0-b}\le\frac{C_5+2C_5\log\ell}{\ell/2}.
\]
As $C_5\le C_5\log\ell$, this gives
$g(b)-g(a)\le2\Lambda(C_5+2C_5\log\ell)/\ell\le6C_5\Lambda\log\ell/\ell\le\varepsilon_C(t)$.

\emph{Conclusion.} By $(*)$ and the monotonicity of $m_T$ in the radius, $m_T(u;z,s_4)\ge-2C_5\log\ell>-\kappa$. By part
(2) of Lemma~\ref{cv-origin-analytic-system-weights-regularize-global-extend}, $m_T(u;z,C_1t)\le m(w;z,C_1t)=g(b)$ and
$g(a)=m(w;z,s_8)\le\max(m_T(u;z,s_4),-\kappa)=m_T(u;z,s_4)$. So
$m_T(u;z,C_1t)-m_T(u;z,s_4)\le g(b)-g(a)\le\varepsilon_C(t)$.
\end{proof}
% Lean: oscGen (RegChartC.lean): fifth_volume_constant_nonneg for C5 >= 0 (exists_far_point at 0 with s = 1/2), the
% constants Lambda, K1, C = 6 C5 Lambda + 2 C5 + 1 and l0 as in the text, pmo_lower with exists_far_point and rLoc_le_far
% for (*), ext_exists with kappa = 2 C5 log l + 1, the convexity of polyMaxLogConvex_thm with
% ConvexOn.slope_mono_adjacent at a < b < sigma, pmo_le_polyMax and polyMax_le_pmo. Size: about 350 lines besides
% Lemma cv-origin-analytic-system-weights-regularize-global-extend.

\begin{lemma}\label{cv-origin-analytic-system-weights-regularize-global-mollifier}
Let $e\ge0$, let $\varrho$ be a mollifier on $\mathbf C^e$ with $\varrho\le b$, let $U\subset\mathbf C^e$ be open, let
$t>0$, and let $U_t$ be the set of $z$ with $\bar\Delta(z,t)\subset U$.
\begin{enumerate}
\item Let $w$ be a real plurisubharmonic function on $U$. Then $U_t$ is open, $w*\varrho_t$ is smooth and
  plurisubharmonic on $U_t$, and $w(z)\le(w*\varrho_t)(z)\le m(w;z,t)$ for $z\in U_t$. If moreover
  $\bar\Delta(z,2t)\subset U$, then
  \[
   (w*\varrho_t)(z)\ge m(w;z,2t)-\pi^e(3/2)^{2e}b\,\big(m(w;z,2t)-m(w;z,t/4)\big).
  \]
\item Let $u$ be a smooth function on $U$ and $K\subset U$ a compact set. Then $K\subset U_t$ for all small $t$, and as
  $t\to0$, $u*\varrho_t$ tends to $u$ with its first and second derivatives uniformly on $K$, with
  $\|d(u*\varrho_t)-du\|\le M't$ on $K$ for some $M'$ and all small $t$.
\end{enumerate}
\end{lemma}
\leanref{ReciprocityTheory/AnalyticForms/PlurisubharmonicFunctions/PshMoll.lean}{ReciprocityTheory.AnForm.A.mollify_thm}
\leanref{GroundTheory/HeightsAndTheta/Heights/Regularisation/PMollifierLowerBound.lean}{GroundTheory.Heights₂.Heights.Regular.MollifiedLowerBound.closed}
\leanref{ReciprocityTheory/AnalyticForms/PlurisubharmonicFunctions/PshMoll.lean}{ReciprocityTheory.AnForm.A.mollification_second_order}
\leanref{Development/AutomorphicSpectrum/Proofs/MetricProofs/MollifyConvergesC2Rate.lean}{GlobalTheory.CompanionTheory.Ar.Giv.mollify_converges_C2_rate}
\begin{proof}
(1) If $z\in U_t$, then $\bar\Delta(z',t)\subset U$ for $z'$ near $z$, by the compactness of $\bar\Delta(z,t)$; so $U_t$ is
open. The function $w$ is locally integrable \cite[Chapter~I, Theorem~4.17 and (5.3)]{Dem12} and bounded above on compact
sets. For $z\in U_t$, $(w*\varrho_t)(z)=\int w(y)\varrho_t(z-y)\,dy$, to which only $y\in\bar\Delta(z,t)$ contribute; so
$w*\varrho_t$ can be differentiated under the integral sign any number of times, and it is smooth on $U_t$. It is
plurisubharmonic there: integrating the mean value inequality \cite[Chapter~I, (5.2)]{Dem12} for the plurisubharmonic
functions $z\mapsto w(z-y)$ against $\varrho_t(y)\,dy$ gives that for $w*\varrho_t$. Let $z\in U_t$. For $c\in\mathbf C$
with $|c|=1$, the substitution $y\mapsto cy$ and $\varrho_t(cy)=\varrho_t(y)$ give
$(w*\varrho_t)(z)=\int\varrho_t(y)w(z-cy)\,dy$. Averaging over $c=e^{i\theta}$, and using the mean value inequality for $w$
on the disc $\{z-\zeta y:|\zeta|\le1\}$, which lies in $\bar\Delta(z,t)$ when $\varrho_t(y)\neq0$, gives
$(w*\varrho_t)(z)\ge\int\varrho_t(y)w(z)\,dy=w(z)$. As $\varrho_t\ge0$ has integral $1$ and support in $\Delta_t$,
$(w*\varrho_t)(z)\le m(w;z,t)$.

Suppose that $\bar\Delta(z,2t)\subset U$, and put $M_2=m(w;z,2t)$ and $M_4=m(w;z,t/4)$. As $w$ is upper semicontinuous,
$w(x)=M_4$ for some $x\in\bar\Delta(z,t/4)$. The polydisc $P=\bar\Delta(x,3t/2)$ contains $\bar\Delta(z,t)$ and lies in
$\bar\Delta(z,2t)$, and its volume is $(\pi(3t/2)^2)^e$. The mean value inequality (5.2) of \cite[Chapter~I]{Dem12} on
the discs $\{x+\zeta v:|\zeta|\le1\}$, integrated over $v\in\bar\Delta(0,3t/2)$, whose measure is invariant under
$v\mapsto cv$ for $|c|=1$, shows that $w(x)$ is at most the mean value of $w$ over $P$; so
$\int_P(M_2-w)\le(\pi(3t/2)^2)^e(M_2-M_4)$. As $\varrho_t\le t^{-2e}b$, and $\varrho_t(y)\neq0$ only if
$z-y\in\bar\Delta(z,t)\subset P$, where $M_2-w\ge0$,
\[
 M_2-(w*\varrho_t)(z)=\int\varrho_t(y)\big(M_2-w(z-y)\big)\,dy\le t^{-2e}b\int_P(M_2-w)\le\pi^e(3/2)^{2e}b\,(M_2-M_4).
\]

(2) Choose $t_1>0$ such that the compact set $K_1$ of points at distance at most $t_1$ from $K$ lies in $U$. For
$t\le t_1$, $K\subset U_t$, and near each point of $K$ the derivatives of order $k\le2$ of $u*\varrho_t$ are
$(d^ku)*\varrho_t$, by differentiation under the integral sign. As $\varrho_t\ge0$ has integral $1$ and support in
$\Delta_t$, for $z\in K$ and $k\le2$ the mean value inequality gives
\[
 \big\|\big((d^ku)*\varrho_t\big)(z)-d^ku(z)\big\|\le\sup_{|y|\le t}\|d^ku(z-y)-d^ku(z)\|\le t\sup_{K_1}\|d^{k+1}u\|,
\]
which gives (2), with $M'=\sup_{K_1}\|d^2u\|$.
\end{proof}
% Lean: (1) is mollify_thm = mollify_main (PshMollB.lean: molli_contDiffAt, molli_submean by circle_fubini, le_molli by
% circle averaging with the circle invariance of the mollifier, molli_le_polyMax) and MollifiedLowerBound.closed
% (= MollifiedLowerBound.proof: the maximum point x of polyMaxLogConvex_thm, subMeanPolydisc_thm on closedBall x (3/2 t),
% volume_real_closedBall_polydisc, the bound t^(-2D) b on the scaled mollifier). (2) is mollification_second_order
% (stated for functions of class C^3, with the rate C t for the value and both derivatives) and
% mollify_converges_C2_rate. Demailly's kernels are radial; the Lean's mollifiers are invariant under the circle acting
% by scalars, and the Lean proves these facts itself. Size: about 1,500 lines.

\begin{lemma}\label{cv-origin-analytic-system-weights-regularize-global-mollify}
Let $e\ge0$, $T\subset\{1,\dots,e\}$ with integers $\mu_i\ge0$ for $i\in T$, let $C_5$ be a number and $u$ a continuous
plurisubharmonic function on $\Delta_{3/2}\setminus H_T$ with $-C_5\log r_{T,\mu}\le u\le C_5$ there, and let $\varrho$ be
a mollifier on $\mathbf C^e$. Then there are $C\ge1$, a number $t_0$ with $0<t_0\le\exp(-\exp(1))$ and $t_0\le1/8$, and,
for $0<t\le t_0$, functions $w_t$ on $\Delta_{3/2}$ equal to $\max\big(u,-C\log\log(1/t)\big)$ on $\Delta_{3/2}\setminus H_T$,
with the following property. For $0<t\le t_0$, the function $u_t=w_t*\varrho_t$ is smooth with nonnegative Levi form on
$\Delta_{5/4}$, $u\le u_t$ on $\Delta_{5/4}\setminus H_T$, and every $z\in\Delta_{5/4}$ has
\[
 -C\log\log(1/t)\le m_T(u;z,t)\quad\text{and}\quad m_T(u;z,t)-\varepsilon_C(t)\le u_t(z)\le m_T(u;z,t).
\]
\end{lemma}
\leanref{Development/AutomorphicSpectrum/Proofs/MetricProofs/LocalPshRegularizationLogBoundStatement.lean}{GlobalTheory.T.local_psh_regularization_log_bound}
\begin{proof}
As in the proof of Lemma~\ref{cv-origin-analytic-system-weights-regularize-global-oscillation}, the point with all
coordinates $1/2$ shows that $C_5\ge0$. That lemma, applied to $u$ with $R=3/2$ and $C_1=2$, gives $C_0\ge1$ and $\ell_0$.
Let $b$ be the maximum of $\varrho$, put $K=\pi^e(3/2)^{2e}b$ and $C=C_0(K+1)$, so that $C\ge C_0$ and $C\ge KC_0$, and let
$t_0=\min(\exp(-\ell_0),\exp(-\exp(1)),1/8)$. Let $0<t\le t_0$, put $\ell=\log(1/t)\ge\ell_0$ and $\kappa=C\log\ell$, and let
$w_t$ be the function $w$ of Lemma~\ref{cv-origin-analytic-system-weights-regularize-global-extend} for $u$, $R=3/2$,
$B=C_5$ and $\kappa$. By Lemma~\ref{cv-origin-analytic-system-weights-regularize-global-oscillation}, $2t\le1/8$, and every
$z\in\Delta_{5/4}$ has
\[
 -C_0\log\ell\le m_T(u;z,t/16)\quad\text{and}\quad m_T(u;z,2t)-m_T(u;z,t/8)\le\varepsilon_{C_0}(t).\tag{$*$}
\]
For $z\in\Delta_{5/4}$, $\bar\Delta(z,1/4)\subset\Delta_{3/2}$. So
Lemma~\ref{cv-origin-analytic-system-weights-regularize-global-mollifier}(1), with $U=\Delta_{3/2}$, shows that $u_t$ is
smooth and plurisubharmonic on $\Delta_{5/4}$, hence has nonnegative Levi form \cite[Chapter~I, \S5.A, before
Theorem~5.8]{Dem12}, and that $u\le w_t\le u_t$ on $\Delta_{5/4}\setminus H_T$.

Fix $z\in\Delta_{5/4}$. For $t/16\le s\le1/4$, $m_T(u;z,s)\ge m_T(u;z,t/16)\ge-C_0\log\ell\ge-\kappa$, by $(*)$. So, by
parts (2) and (3) of Lemma~\ref{cv-origin-analytic-system-weights-regularize-global-extend} with $s'=1/4$,
$m(w_t;z,s)=m_T(u;z,s)$ for $t/16\le s\le1/8$. In particular $-C\log\log(1/t)=-\kappa\le m_T(u;z,t)$, and
$u_t(z)\le m(w_t;z,t)=m_T(u;z,t)$ by Lemma~\ref{cv-origin-analytic-system-weights-regularize-global-mollifier}(1). By the
same lemma,
\[
 u_t(z)\ge m(w_t;z,2t)-K\big(m(w_t;z,2t)-m(w_t;z,t/4)\big).
\]
Here $m(w_t;z,2t)=m_T(u;z,2t)\ge m_T(u;z,t)$, and $m(w_t;z,t/4)\ge m_T(u;z,t/4)\ge m_T(u;z,t/8)$, so the bracket is at
most $\varepsilon_{C_0}(t)$ by $(*)$. As $K\varepsilon_{C_0}(t)\le\varepsilon_C(t)$, $u_t(z)\ge m_T(u;z,t)-\varepsilon_C(t)$.
\end{proof}
% Lean: local_psh_regularization_log_bound = Giv.molli_trunc_near_polyMax (LocalPshRegularizationLogBound.lean):
% oscGen with R = 3/2 and C1 = 2, mollifier_bdd for b, K = pi^D (3/2)^(2D) b, C = C0 (K + 1), the extension extT with
% kappa = C log l, mollify_thm, polyMax_le_pmoff_same (hPM) and pmo_le_polyMax (hMP), and MollifiedLowerBound.closed.
% Size: about 150 lines besides the lemmas it uses.

\begin{lemma}\label{cv-origin-analytic-system-weights-regularize-global-near}
Let $e\ge0$, $T\subset\{1,\dots,e\}$ with integers $\mu_i\ge1$ for $i\in T$, and let $C_5$ and $K$ be real numbers. Let $u$
be a differentiable function on $\Delta_{3/2}\setminus H_T$ with
$\|du(z)\|\le C_5\,r_{T,\mu}(z)^{C_5}e^{r_{T,\mu}(z)/2}$ there. Then there is $t_0>0$ with the following property. Let
$0<t\le t_0$ and $z\in\Delta_{5/4}\setminus H_T$ with $r_{T,\mu}(z)\le\log(1/t)+K$. Then $\bar\Delta(z,t)$ does not meet
$H_T$, $u\le u(z)+t^{1/3}/2$ on $\bar\Delta(z,t)$, and so $m_T(u;z,t)\le u(z)+t^{1/3}/2$.
\end{lemma}
\leanref{Development/AutomorphicSpectrum/Proofs/MetricProofs/LocalRegularizationGradientBoundStatement.lean}{GlobalTheory.T.local_regularization_gradient_bound}
\begin{proof}
Put $a=|C_5|$ and $K'=K+|\mu|\log4$. As $y^ae^{-y/6}\to0$ when $y\to\infty$, there is $t_1>0$ with
$ae^{K'/2}(\log(1/t)+K')^at^{1/2}\le t^{1/3}/2$ for $0<t\le t_1$: write $t^{1/2}=t^{1/3}e^{K'/6}e^{-y/6}$ with
$y=\log(1/t)+K'$. Let $t_0=\min(1/4,e^{-K},t_1)$, and let $t$ and $z$ be as in the statement. For $i\in T$, as $\mu_i\ge1$
and every term of $r_{T,\mu}(z)$ is nonnegative, $\log(4/|z_i|^2)\le\log(1/t)+K\le2\log(1/t)$, as $e^K\le1/t$; so
$|z_i|\ge2t$. For $x\in\bar\Delta(z,t)$ and $i\in T$, $|x_i|\ge|z_i|-t\ge|z_i|/2>0$. So $x\in\Delta_{3/2}\setminus H_T$, and
$\log(4/|x_i|^2)\le\log(4/|z_i|^2)+\log4$, so that $r_{T,\mu}(x)\le r_{T,\mu}(z)+|\mu|\log4\le y$. As $r_{T,\mu}(x)\ge1$,
$\|du(x)\|\le ay^ae^{y/2}$. As $\log t=K'-y$, $ay^ae^{y/2}t=ae^{K'/2}y^at^{1/2}\le t^{1/3}/2$. By the mean value inequality
on the convex set $\bar\Delta(z,t)$, $u(x)\le u(z)+ay^ae^{y/2}t\le u(z)+t^{1/3}/2$ for every $x\in\bar\Delta(z,t)$.
\end{proof}
% Lean: local_regularization_gradient_bound = Giv.polyMaxOff_le_of_rLoc_le (LocalRegularizationGradientBound.lean):
% near_small (tendsto_rpow_mul_exp_neg_mul_atTop_nhds_zero), t0 = min(min (1/4) (exp(-K))) t1, hcoord (|z i| >= 2 t),
% hfar, hmem, hrx, the bound B = a y^a exp(y/2) and Convex.norm_image_sub_le_of_norm_fderiv_le. Size: about 200 lines.

\begin{lemma}\label{cv-origin-analytic-system-weights-regularize-global-chart}
Let $e\ge0$, $T\subset\{1,\dots,e\}$ with integers $\mu_i\ge1$ for $i\in T$, and let $C_5$ and $K$ be real numbers. Let $q$
be a smooth function on $\Delta_2$ and $\varphi$ a smooth function on $\Delta_2\setminus H_T$ such that $q+\varphi$ has
nonnegative Levi form at every point of $\Delta_2\setminus H_T$, $-C_5\log r_{T,\mu}\le\varphi\le C_5$ on
$\Delta_2\setminus H_T$, and $\|d\varphi(z)\|\le C_5\,r_{T,\mu}(z)^{C_5}e^{r_{T,\mu}(z)/2}$ for $z\in\Delta_{3/2}\setminus H_T$.
Then there are $C\ge1$, $0<t_0\le\exp(-\exp(1))$ and smooth functions $\varphi_t$ on $\Delta_{5/4}$, for $0<t\le t_0$, such
that for $0<t\le t_0$:
\begin{enumerate}
\item[(a)] $q+\varphi_t$ has nonnegative Levi form at every point of $\Delta_{5/4}$;
\item[(b)] $-C\log\log(1/t)\le\varphi_t\le C$ on $\Delta_{5/4}$;
\item[(c)] $\varphi\le\varphi_t$ on $\Delta_{5/4}\setminus H_T$, and $\varphi_t(z)\le\varphi(z)+t^{1/3}$ at the points
  $z\in\Delta_{5/4}\setminus H_T$ with $r_{T,\mu}(z)\le\log(1/t)+K$;
\item[(d)] $m_T(q+\varphi;z,t)-\varepsilon_C(t)\le q(z)+\varphi_t(z)\le m_T(q+\varphi;z,t)$ for $z\in\Delta_{5/4}$;
\item[(e)] $\varphi_t\le\varphi_{t'}+\varepsilon_C(t')$ on $\Delta_{5/4}$ for $t\le t'\le t_0$;
\end{enumerate}
and, for every compact set $K'\subset\Delta_{5/4}\setminus H_T$, $\varphi_t$ tends to $\varphi$ with its first and second
derivatives uniformly on $K'$ as $t\to0$, with $\|d\varphi_t-d\varphi\|\le M't$ on $K'$ for some $M'$ and all small $t$.
\end{lemma}
\leanref{Development/AutomorphicSpectrum/Proofs/MetricProofs/LocalLeviPshRegularizationStatement.lean}{GlobalTheory.T.local_levi_psh_regularization}
\begin{proof}
\emph{The function $u=q+\varphi$.} As in the proof of
Lemma~\ref{cv-origin-analytic-system-weights-regularize-global-oscillation}, the point with all coordinates $1/2$ shows that
$C_5\ge0$. As $\bar\Delta_{3/2}$ is a compact subset of $\Delta_2$, there is $Q\ge0$ with $|q|\le Q$ and $\|dq\|\le Q$ on
$\Delta_{3/2}$. On $\Delta_{3/2}\setminus H_T$, $\log r_{T,\mu}\ge1$, and $u$ is continuous and plurisubharmonic, as a
function of class $C^2$ with nonnegative Levi form \cite[Chapter~I, \S5.A, before Theorem~5.8]{Dem12}, with
\[
 -(C_5+Q)\log r_{T,\mu}\le u\le C_5+Q,\qquad
 \|du\|\le Q+C_5r_{T,\mu}^{C_5}e^{r_{T,\mu}/2}\le(C_5+Q)\,r_{T,\mu}^{C_5+Q}e^{r_{T,\mu}/2}.
\]

\emph{Construction.} Fix a mollifier $\varrho$ on $\mathbf C^e$, for example a multiple of a smooth function of
$\sum_i|w_i|^2$ supported in the unit ball. Lemma~\ref{cv-origin-analytic-system-weights-regularize-global-mollify},
applied to $u$ with $C_5+Q$, gives $C_1\ge1$, $t_1$, the functions $w_t$ and $u_t=w_t*\varrho_t$.
Lemma~\ref{cv-origin-analytic-system-weights-regularize-global-near}, applied to $u$, $C_5+Q$ and $K$, gives $t_2$. Put
$\varphi_t=u_t-q$ on $\Delta_{5/4}$, $C=2C_1+2Q+C_5+1$ and $t_0=\min(t_1,t_2)$. Let $0<t\le t_0$. By
Lemma~\ref{cv-origin-analytic-system-weights-regularize-global-numbers}, $\log\log(1/t)\ge1$,
$\varepsilon_{C_1}(t)\le\varepsilon_C(t)$, and $\varepsilon_{C_1}(t)\le C_1\log\log(1/t)$, as $\log(1/t)\ge1$.
\begin{enumerate}
\item[(a)] $q+\varphi_t=u_t$ is smooth with nonnegative Levi form on $\Delta_{5/4}$.
\item[(b)] For $z\in\Delta_{5/4}$, $\bar\Delta(z,t)\subset\Delta_{3/2}$, so $m_T(u;z,t)\le C_5+Q$. Hence
  $\varphi_t(z)\le m_T(u;z,t)+Q\le C_5+2Q\le C$, and
  $\varphi_t(z)\ge m_T(u;z,t)-\varepsilon_{C_1}(t)-Q\ge-2C_1\log\log(1/t)-Q\ge-C\log\log(1/t)$.
\item[(c)] Off $H_T$, $u\le u_t$, so $\varphi\le\varphi_t$. If moreover $r_{T,\mu}(z)\le\log(1/t)+K$, then
  $u_t(z)\le m_T(u;z,t)\le u(z)+t^{1/3}/2$ by Lemma~\ref{cv-origin-analytic-system-weights-regularize-global-near}, so
  $\varphi_t(z)\le\varphi(z)+t^{1/3}$.
\item[(d)] This is the bound of Lemma~\ref{cv-origin-analytic-system-weights-regularize-global-mollify}, as
  $\varepsilon_{C_1}(t)\le\varepsilon_C(t)$.
\item[(e)] For $t\le t'\le t_0$, $u_t(z)\le m_T(u;z,t)\le m_T(u;z,t')\le u_{t'}(z)+\varepsilon_{C_1}(t')$, and
  $\varepsilon_{C_1}(t')\le\varepsilon_C(t')$.
\end{enumerate}

\emph{Convergence.} Let $K'\subset\Delta_{5/4}\setminus H_T$ be compact. As $r_{T,\mu}$ is continuous off $H_T$, it is at most
some $B_0$ on $K'$; let $V$ be the open set of $z\in\Delta_{5/4}\setminus H_T$ with $r_{T,\mu}(z)<B_0+1$, which contains
$K'$. For small $t$, $C_1\log\log(1/t)\ge(C_5+Q)\log(B_0+1)$, and then $u\ge-(C_5+Q)\log r_{T,\mu}\ge-C_1\log\log(1/t)$ on
$V$, so $w_t=u$ on $V$. For small $t$, the polydiscs $\bar\Delta(z,t)$ with $z$ in a neighbourhood of $K'$ lie in $V$, so
$u_t=u*\varrho_t$ there, computed on $V$. By
Lemma~\ref{cv-origin-analytic-system-weights-regularize-global-mollifier}(2), applied to $u$ on $V$, $u_t$ tends to $u$ with
its first and second derivatives uniformly on $K'$, with $\|du_t-du\|\le M't$ there for small $t$. Subtracting $q$ gives
the same for $\varphi_t$ and $\varphi$.
\end{proof}
% Lean: local_levi_psh_regularization = Giv.regularise_one_chart (LocalLeviPshRegularization.lean): oc_q_bounds for Q,
% hubd and hugrad (oc_grad) for u = q + ph, isPsh_iff_isLeviPsh_thm, mollifierInhabit_thm for the mollifier,
% local_psh_regularization_log_bound, local_regularization_gradient_bound, phT t = molli rho t (w t) - q,
% C = 2 C1 + 2 Q + C5 + 1, t0 = min t1 t2; for the convergence, the set V and mollify_converges_C2_rate with
% fderiv_add_and_second_fderiv_add. Size: about 300 lines here, about 4,300 with the lemmas it uses.

\begin{lemma}\label{cv-origin-analytic-system-weights-regularize-global-compare}
Let $S$, $s$, $\partial$ and $r$ be as in Definition~\ref{cv-origin-analytic-system-weight}(1), with $e=\dim S$, and let
$\bar L$ be a holomorphic line bundle on $S$ with a smooth Hermitian metric with local potentials $\vartheta_x$, each defined
on an open neighbourhood $U_x$ of $x$. Let $C_5$ be a number and $\varphi$ a smooth function on $S\setminus\partial$ such
that the metric of $\bar L$ multiplied by $e^{-\varphi}$ is semipositive on $S\setminus\partial$ and
$-C_5\log r\le\varphi\le C_5$ there. Let $\kappa_1,\dots,\kappa_N$ be charts of $S$, with sets $T_j$, integers
$\mu_{j,i}\ge1$ for $i\in T_j$ and a number $M$ satisfying (1), (3) and (4) of
Lemma~\ref{cv-origin-analytic-system-weights-regularize-atlas}, and let $x_1,\dots,x_N\in S$ be points such that the domain
of $\kappa_j$ lies in $U_{x_j}$. For each $j$ put $u_j(y)=\vartheta_{x_j}(\kappa_j^{-1}(y))+\varphi(\kappa_j^{-1}(y))$ for
$y\in\Delta_2\setminus H_{T_j}$, and let $B_j$ be the set of points $x$ of the domain of $\kappa_j$ with
$\kappa_j(x)\in\Delta_{5/4}$. Then there are $C$ and $t_0>0$ such that, for $0<t\le t_0$, all $j$ and $k$ and every
$x\in B_j\cap B_k$,
\[
 \Big|\big(m_{T_j}(u_j;\kappa_j(x),t)-\vartheta_{x_j}(x)\big)-\big(m_{T_k}(u_k;\kappa_k(x),t)-\vartheta_{x_k}(x)\big)\Big|
 \le\varepsilon_C(t).
\]
\end{lemma}
\leanref{Development/AutomorphicSpectrum/Proofs/MetricProofs/GluedLocalRegularizationsOnAtlasStatement.lean}{GlobalTheory.T.glued_local_regularizations_on_atlas}
\begin{proof}
\emph{The functions $u_j$.} Fix $j$. By (3), $\Delta_2\setminus H_{T_j}$ is the set of $y\in\Delta_2$ with
$\kappa_j^{-1}(y)\notin\partial$, and $u_j$ is smooth on it. Near each of its points, $u_j$ is the composite of
$\vartheta_{x_j}+\varphi$, read in a chart of $S$, with a holomorphic change of charts. The former has nonnegative Levi
form, as the metric of $\bar L$ multiplied by $e^{-\varphi}$ is semipositive off $\partial$, and holomorphic maps carry
Levi forms \cite[Chapter~I, (5.10)]{Dem12}; so $u_j$ has nonnegative Levi form, and it is plurisubharmonic. Let $Q_j$ bound
$|\vartheta_{x_j}\circ\kappa_j^{-1}|$ on $\bar\Delta_2$, and put $\Lambda=\log(1+|M|)$. Let $y\in\Delta_2\setminus H_{T_j}$,
and put $\rho_j=r_{T_j,\mu_j}(y)$ and $\rho'=r(\kappa_j^{-1}(y))$, both at least $\exp(1)$. By (4),
$\rho'\le\rho_j+|M|\le\rho_j(1+|M|)$, so $\log\rho'\le\log\rho_j+\Lambda\le(1+\Lambda)\log\rho_j$. Hence
$-C^b_j\log\rho_j\le u_j(y)\le C^b_j$, with $C^b_j=|Q_j|+|C_5|(1+\Lambda)+|C_5|$.

\emph{Changes of charts.} Fix $j$ and $k$. The set $\kappa_j^{-1}(\bar\Delta_{5/4})\cap\kappa_k^{-1}(\bar\Delta_{5/4})$ is
compact, and its image under $\kappa_j$ lies in the open set of points $y$ of the image of $\kappa_j$ with
$\kappa_j^{-1}(y)$ in the domain of $\kappa_k$. On that open set the maps $y\mapsto\kappa_k(\kappa_j^{-1}(y))$ and
$y\mapsto g_{jk}(\kappa_j^{-1}(y))$, with $g_{jk}=\vartheta_{x_j}-\vartheta_{x_k}$, are smooth. By compactness and the mean
value inequality, there are $r_{jk}>0$ and $L_{jk}\ge0$ such that, for $x\in B_j\cap B_k$ and
$y\in\bar\Delta(\kappa_j(x),r_{jk})$, the point $\kappa_j^{-1}(y)$ lies in the domain of $\kappa_k$, and
$|\kappa_k(\kappa_j^{-1}(y))-\kappa_k(x)|$ and $|g_{jk}(\kappa_j^{-1}(y))-g_{jk}(x)|$ are at most $L_{jk}|y-\kappa_j(x)|$. Let
$0<r_0\le1$ be at most every $r_{jk}$, and $C_1=\max(1,\max_{j,k}L_{jk})$.

\emph{Oscillation.} Lemma~\ref{cv-origin-analytic-system-weights-regularize-global-oscillation}, applied to each $u_k$
with $R=2$, the constant $C^b_k$ and $C_1$, gives $C^o_k\ge1$ and $\ell_k$. Let $C_m$ and $\ell_m$ be the largest of the
$C^o_k$ and of the $\ell_k$, put $C=C_1+C_m$ and $t_0=\min(\exp(-\ell_m),r_0)$, and let $0<t\le t_0$, so that
$\log(1/t)\ge\ell_k$ for every $k$. By that lemma, $\log(1/t)\ge\exp(1)$ and $C_1t\le3/8$.

\emph{The comparison.} Let $x\in B_j\cap B_k$, $z=\kappa_j(x)$ and $z'=\kappa_k(x)$. Let $y\in\bar\Delta(z,t)\setminus H_{T_j}$
and $p=\kappa_j^{-1}(y)$, so that $p\notin\partial$. As $t\le r_{jk}$, $p$ lies in the domain of $\kappa_k$, and
$y'=\kappa_k(p)$ has $|y'-z'|\le C_1t$, so $y'\in\Delta_2$, and $y'\notin H_{T_k}$ by (3). Then
\[
 u_j(y)=\vartheta_{x_j}(p)+\varphi(p)=u_k(y')+g_{jk}(p)\le m_{T_k}(u_k;z',C_1t)+g_{jk}(x)+C_1t.
\]
So $m_{T_j}(u_j;z,t)\le m_{T_k}(u_k;z',C_1t)+g_{jk}(x)+C_1t$. By
Lemma~\ref{cv-origin-analytic-system-weights-regularize-global-oscillation} and the monotonicity of $m_{T_k}$ in the radius,
$m_{T_k}(u_k;z',C_1t)\le m_{T_k}(u_k;z',t/(4C_1))+\varepsilon_{C^o_k}(t)\le m_{T_k}(u_k;z',t)+\varepsilon_{C_m}(t)$; and
$C_1t\le\varepsilon_{C_1}(t)$ by Lemma~\ref{cv-origin-analytic-system-weights-regularize-global-numbers}. So
\[
 \big(m_{T_j}(u_j;z,t)-\vartheta_{x_j}(x)\big)-\big(m_{T_k}(u_k;z',t)-\vartheta_{x_k}(x)\big)\le\varepsilon_{C_1}(t)
 +\varepsilon_{C_m}(t)=\varepsilon_C(t).
\]
Exchanging $j$ and $k$ gives the other inequality.
\end{proof}
% Lean: glued_local_regularizations_on_atlas = Giv.polyMaxOff_charts_compare (GluedLocalRegularizationsOnAtlas.lean):
% Compare.chart_data (plurisubharmonicity through compat_local and levi_comp, and the bound
% |Q| + |C5| (1 + log(1 + |M|)) + |C5|), Compare.overlap_pair (exists_lip_near on the compact overlap, for the change
% of charts and th (x j) - th (x k) together), r0 = (max Rinv 1)^-1, C1 = max L0 1, oscGen with R = 2, hkey, and
% Compare.le_log_div for C1 t <= epsT C1 t. Size: about 1,000 lines besides
% Lemma cv-origin-analytic-system-weights-regularize-global-oscillation.

\begin{lemma}\label{cv-origin-analytic-system-weights-regularize-global-max}
Let $I$ be a finite nonempty set, let $\theta\ge0$ be a smooth even function on $\mathbf R$ with support in
$[-1/2,1/2]$ and integral $1$, and for $\eta>0$ and $\tau\in\mathbf R^I$ put
\[
 \mathcal M_\eta(\tau)=\int_{\mathbf R^I}\max_{j\in I}(\tau_j+\eta h_j)\prod_{j\in I}\theta(h_j)\,dh.
\]
Let $\eta>0$. Then:
\begin{enumerate}
\item $\mathcal M_\eta$ is smooth, convex and nondecreasing in each variable,
  $\max_j\tau_j\le\mathcal M_\eta(\tau)\le\max_j\tau_j+\eta/2$, and $\mathcal M_\eta((\tau_j+c)_j)=\mathcal M_\eta(\tau)+c$ for
  $c\in\mathbf R$;
\item if $k_0\in I'\subset I$ and $\tau_j+\eta\le\tau_{k_0}$ for every $j\notin I'$, then
  $\mathcal M_\eta(\tau)=\mathcal M_\eta(\tau')$, where $\tau'_j=\tau_j$ for $j\in I'$ and $\tau'_j=\tau_{k_0}-\eta$ for
  $j\notin I'$;
\item $\partial_j\mathcal M_\eta\ge0$ and $\sum_j\partial_j\mathcal M_\eta=1$, and the second derivative of
  $\mathcal M_\eta$ has norm at most $C_I/\eta$ at every point, for the maximum norm on $\mathbf R^I$ and a constant $C_I$
  depending only on $I$ and $\theta$;
\item let $u_j$ ($j\in I$) be functions of class $C^2$ on an open subset of $\mathbf C^e$, let $y$ be a point of it and $v$
  a vector, and suppose that the Levi form of every $u_j$ at $y$ on $v$ is at least $g$. Then so is the Levi form of
  $\mathcal M_\eta((u_j)_j)$ at $y$ on $v$.
\end{enumerate}
\end{lemma}
\leanref{ReciprocityTheory/AnalyticForms/PlurisubharmonicFunctions/IsPshLocallyIntegrableStatement.lean}{ReciprocityTheory.AnForm.A.regMax}
\leanref{ReciprocityTheory/AnalyticForms/PlurisubharmonicFunctions/PshRegMax.lean}{ReciprocityTheory.AnForm.A.regMax_thm}
\leanref{ReciprocityTheory/AnalyticForms/LeviForms/RegGlue.lean}{ReciprocityTheory.AnForm.P.regMax_dominated}
\leanref{ReciprocityTheory/AnalyticForms/PlurisubharmonicFunctions/PshRegMax.lean}{ReciprocityTheory.AnForm.A.regMaxLevi_thm}
\begin{proof}
Up to the substitution $h_j\mapsto\eta h_j$, $\mathcal M_\eta$ is Demailly's regularized maximum with all parameters
$\eta/2$ and the kernel $h\mapsto\theta(h/2)/2$, supported in $[-1,1]$ \cite[Chapter~I, Lemma~5.18]{Dem12}, and (1) is its
properties (a), (b) and (d).

(2) If $|h_j|<1/2$ for every $j$, then for $j\notin I'$ both $\tau_j+\eta h_j$ and $\tau_{k_0}-\eta+\eta h_j$ are less than
$\tau_{k_0}-\eta/2<\tau_{k_0}+\eta h_{k_0}$; so the maxima for $\tau$ and $\tau'$ under the integral are both the maximum over
$I'$. Otherwise $\prod_j\theta(h_j)=0$.

(3) The first two claims follow from the monotonicity and from differentiating
$\mathcal M_\eta((\tau_j+c)_j)=\mathcal M_\eta(\tau)+c$ in $c$. For the third, $\mathcal M_\eta(\tau)=\eta\mathcal M_1(\tau/\eta)$,
so it suffices to bound the second derivatives of $\mathcal M_1$. Substituting $s=\tau+h$,
$\mathcal M_1(\tau)=\int\max_js_j\prod_j\theta(s_j-\tau_j)\,ds$. At a point $\tau_0$, the integral over $s$ of the second
derivatives in $\tau$ of $\prod_j\theta(s_j-\tau_j)$ is $0$, so the second derivatives of $\mathcal M_1$ at $\tau_0$ are the
integrals of $(\max_js_j-\max_j\tau_{0,j})$ against them. On their support $|s_j-\tau_{0,j}|\le1/2$ for every $j$, so
$|\max_js_j-\max_j\tau_{0,j}|\le1/2$, and the bound depends only on $I$ and $\theta$.

(4) By the chain rule, the Levi form of $\mathcal M_\eta((u_j)_j)$ at $y$ on $v$ is
$\sum_j\partial_j\mathcal M_\eta\cdot L_{u_j}(y)(v)$ plus a quarter of the sum of the second derivative of $\mathcal M_\eta$
on the vectors $(du_j(y)v)_j$ and on $(du_j(y)(iv))_j$. The latter is nonnegative, as $\mathcal M_\eta$ is convex; and by
(3) the former is at least $g$. This is property (e) of \cite[Chapter~I, Lemma~5.18]{Dem12} at one point.
\end{proof}
% Lean: regMax eta t is the integral of tmax (t j + eta h j) against the product of ktheta (h j), ktheta the normed
% ContDiffBump with radii 1/4 and 1/2. regMax_thm (PshRegMax.lean, RegMaxStatement): regMax_contDiff, regMax_convex,
% regMax_mono, regMax_ge and regMax_le, regMax_eq_of_dom, regMax_add_const, regMax_fderiv_sum, and regMax_fderiv2_bound
% (fderiv2_regMax for the scaling, regMax_one_fderiv2_bound with regMax_one_eq_conv_sub and norm_coordMax_sub_le_half);
% regMax_fderiv_nonneg. (2) is regMax_dominated (RegGlue.lean); (4) is regMaxLevi_thm (regMax_comp_fderiv2,
% regMax_fderiv2_nonneg, regMax_fderiv_sum). The Lean proves Demailly's properties itself. Size: about 900 lines.

\begin{lemma}\label{cv-origin-analytic-system-weights-regularize-global-max-limit}
Let $I$, $\theta$ and $\mathcal M_\eta$ be as in Lemma~\ref{cv-origin-analytic-system-weights-regularize-global-max}. Let
$U$ be an open subset of a real normed space, $K\subset U$, and $f$ a function of class $C^2$ on $U$. For $t>0$ let
$\varepsilon(t)$ and $\delta(t)$ be numbers and $u_{t,j}$ ($j\in I$) functions on $U$, such that: for all small $t>0$,
$\varepsilon(t)>0$, the $u_{t,j}$ are of class $C^2$, and $\|du_{t,j}-df\|\le\delta(t)$ on $K$; as $t\to0$,
$\varepsilon(t)\to0$ and $\delta(t)^2/\varepsilon(t)\to0$; and each $u_{t,j}$ tends to $f$ with its first and second
derivatives uniformly on $K$. Then $\mathcal M_{\varepsilon(t)}((u_{t,j})_j)+\varepsilon(t)/2$ tends to $f$ with its first
and second derivatives uniformly on $K$ as $t\to0$.
\end{lemma}
\leanref{GroundTheory/HeightsAndTheta/Heights/AnalyticLimit/RegularizedMaxC2ExampleRealLine.lean}{GroundTheory.Heights₂.Heights.AnalyticLimit.RegularizedMaxConvergenceC2.proof}
\begin{proof}
Let $t$ be small, $\eta=\varepsilon(t)$, and $y\in K$, and let $s_0$, $s_1$ and $s_2$ be the largest over $j$ of
$|u_{t,j}(y)-f(y)|$, $\|du_{t,j}(y)-df(y)\|$ and $\|d^2u_{t,j}(y)-d^2f(y)\|$. Put $b_j=u_{t,j}-f$ and
$G=\mathcal M_\eta((b_j)_j)$, so that $\mathcal M_\eta((u_{t,j})_j)=G+f$ by
Lemma~\ref{cv-origin-analytic-system-weights-regularize-global-max}(1). By the same lemma,
$|G(y)+\eta/2|\le s_0+\eta$, and $dG(y)=\sum_j\partial_j\mathcal M_\eta\,db_j(y)$ has norm at most $s_1$. For a vector $v$,
$d^2G(y)(v,v)$ is the second derivative of $\mathcal M_\eta$ on the vector $(db_j(y)v)_j$, plus
$\sum_j\partial_j\mathcal M_\eta\,d^2b_j(y)(v,v)$; so $|d^2G(y)(v,v)|\le(C_Is_1^2/\eta+s_2)|v|^2$, and, $d^2G(y)$ being
symmetric, $\|d^2G(y)\|\le2(C_Is_1^2/\eta+s_2)$ by polarization. So, on $K$, the differences between
$\mathcal M_\eta((u_{t,j})_j)+\eta/2$ and $f$, and between their first and second derivatives, are at most $s_0+\varepsilon(t)$,
$s_1$ and $2(C_I\delta(t)^2/\varepsilon(t)+s_2)$, and these tend to $0$ uniformly on $K$.
\end{proof}
% Lean: RegularizedMaxConvergenceC2.proof (RegularizedMaxC2ExampleRealLine.lean): norm_fderiv2_regMax_le, and
% regMax_c2_pointwise_bound (b j = u j - f, g = regMax eta b with regMax_add_const, abs_regMax_le,
% abs_fderiv_regMax_apply_le, regMax_comp_fderiv2, norm_le_of_symm_diag for the polarization). Size: about 300 lines
% besides Lemma cv-origin-analytic-system-weights-regularize-global-max.

\begin{lemma}\label{cv-origin-analytic-system-weights-regularize-global-glue}
Let $S$ be a compact complex manifold of dimension $e$ with a K\"ahler form $\omega$, let $\partial\subset S$ be closed,
let $\bar L$ be a holomorphic line bundle on $S$ with a smooth Hermitian metric, and let $\varphi$ be a smooth function on
$S\setminus\partial$. Let $\kappa_1,\dots,\kappa_N$ be charts of $S$ whose images contain $\bar\Delta_2$; for $\rho\le2$ let
$B_j(\rho)$ be the set of points $x$ of the domain of $\kappa_j$ with $\kappa_j(x)\in\Delta_\rho$, and suppose that the
$B_j(1)$ cover $S$. Then there is $K>0$ with the following property. Let $t_0>0$, let $\delta(t)$ be numbers with
$\delta(t)\to0$ as $t\to0$, and for $0<t\le t_0$ and every $j$ let $f_{t,j}$ be a smooth function on $B_j(5/4)$, such that
for $0<t\le t_0$:
\begin{enumerate}
\item[(a)] $t\le\delta(t)\le1$, and the $|f_{t,j}|$ are bounded on $B_j(5/4)$ by a number independent of $j$;
\item[(b)] the metric of $\bar L$ multiplied by $e^{-f_{t,j}}$ is semipositive on $B_j(5/4)$;
\item[(c)] $|f_{t,j}-f_{t,k}|\le\delta(t)$ on $B_j(5/4)\cap B_k(5/4)$;
\end{enumerate}
and such that, for every $j$ and every compact set $K'\subset\Delta_{5/4}$ with $\kappa_j^{-1}(K')\cap\partial=\emptyset$,
$f_{t,j}\circ\kappa_j^{-1}$ tends to $\varphi\circ\kappa_j^{-1}$ with its first and second derivatives uniformly on $K'$ as
$t\to0$, with $\|d(f_{t,j}\circ\kappa_j^{-1})-d(\varphi\circ\kappa_j^{-1})\|\le M't$ on $K'$ for some $M'$ and all small
$t$. Then there are smooth functions $F_t$ on $S$, for $0<t\le t_0$, such that the first Chern form of $\bar L$ with its
metric multiplied by $e^{-F_t}$ is at least $-K\delta(t)\omega$, $f_{t,j}\le F_t\le f_{t,j}+K\delta(t)$ on $B_j(5/4)$, and
$F_t$ tends to $\varphi$ with its first and second derivatives, in every chart and uniformly on compact subsets of
$S\setminus\partial$, as $t\to0$.
\end{lemma}
\leanref{Development/AutomorphicSpectrum/Proofs/MetricProofs/RegularizationGlueC2ExplicitStatement.lean}{GlobalTheory.T.regularization_glue_C2_explicit}
\leanref{ReciprocityTheory/AnalyticForms/LeviForms/RegGlue.lean}{ReciprocityTheory.AnForm.P.glue_cutoff}
\begin{proof}
Let $\vartheta_x$ and $\Theta_x$ be local potentials of the metric of $\bar L$ and of $\omega$. If $N=0$, then $S$ is empty and
there is nothing to prove; suppose $N\ge1$, and let $\mathcal M_\eta$ be as in
Lemma~\ref{cv-origin-analytic-system-weights-regularize-global-max} with $I=\{1,\dots,N\}$.

\emph{Cutoffs.} Let $\chi$ be a smooth function on $\mathbf C^e$ with $0\le\chi\le1$, equal to $1$ on $\Delta_1$ and to $0$
outside $\Delta_{9/8}$. For each $j$ let $\eta_j=\chi\circ\kappa_j-1$ on the domain of $\kappa_j$ and $\eta_j=-1$ elsewhere.
It is smooth, as $\chi\circ\kappa_j$ vanishes outside the compact set $\kappa_j^{-1}(\bar\Delta_{9/8})$; $-1\le\eta_j\le0$,
$\eta_j=0$ on $B_j(1)$, and $\eta_j=-1$ outside $B_j(9/8)$. As $S$ is compact and $\omega$ is positive, there are numbers
$c_j$ such that $\eta_j+c_j\Theta_x$ has nonnegative Levi form in every chart, for every $x$. Put $K_1=\sum_j|c_j|$,
$K_0=4K_1+2$ and $K=2K_0$.

\emph{The construction.} Let $0<t\le t_0$, write $\delta=\delta(t)$ and $f_j=f_{t,j}$, and let $M_b$ bound the $|f_j|$. For
$y\in S$ put $a_j(y)=f_j(y)+4\delta\eta_j(y)$ if $y\in B_j(5/4)$ and $a_j(y)=-M_b-10$ otherwise, and
$F_t(y)=\mathcal M_\delta((a_j(y))_j)+K_0\delta$.

\emph{Local form.} Let $y_0\in S$. Choose $k_0$ with $y_0\in B_{k_0}(1)$, let $I_0$ be the set of $j$ with $y_0\in B_j(5/4)$,
and let $W$ be the intersection of $B_{k_0}(1)$, of the $B_j(5/4)$ for $j\in I_0$, and of the complements of the compact
sets $\kappa_j^{-1}(\bar\Delta_{9/8})$ for $j\notin I_0$. It is an open neighbourhood of $y_0$, and it does not depend on $t$.
Let $y\in W$. Then $a_{k_0}(y)=f_{k_0}(y)$, and $a_j(y)+\delta\le a_{k_0}(y)$ for $j\notin I_0$: if $y\in B_j(5/4)$, then
$y\notin B_j(9/8)$, so $a_j(y)=f_j(y)-4\delta\le f_{k_0}(y)-3\delta$ by (c); otherwise $a_j(y)=-M_b-10$. By
Lemma~\ref{cv-origin-analytic-system-weights-regularize-global-max}(2), on $W$
\[
 \mathcal M_\delta((a_j)_j)=\mathcal M_\delta((b_j)_j),\quad b_j=f_j+4\delta\eta_j\ (j\in I_0),\quad
 b_j=f_{k_0}+4\delta\eta_{k_0}-\delta\ (j\notin I_0),\tag{$*$}
\]
and the $b_j$ are smooth on $W$. So $F_t$ is smooth.

\emph{Bounds.} Let $y\in B_j(5/4)$. For every $j'$, $a_{j'}(y)\le f_j(y)+\delta$: if $y\in B_{j'}(5/4)$, then
$a_{j'}(y)\le f_{j'}(y)\le f_j(y)+\delta$, as $\eta_{j'}\le0$; otherwise $a_{j'}(y)=-M_b-10$. For $k$ with $y\in B_k(1)$,
$a_k(y)=f_k(y)\ge f_j(y)-\delta$. By Lemma~\ref{cv-origin-analytic-system-weights-regularize-global-max}(1),
$|\mathcal M_\delta((a_{j'}(y))_{j'})-f_j(y)|\le2\delta\le K_0\delta$, so $f_j\le F_t\le f_j+K\delta$ on $B_j(5/4)$.

\emph{Positivity.} Near $y_0$, in a chart of $S$, $(*)$ and
Lemma~\ref{cv-origin-analytic-system-weights-regularize-global-max}(1) give
$\vartheta_x+F_t+K\delta\Theta_x=\mathcal M_\delta((v_j)_j)$, with $v_j=P_j+4\delta Q_j+K\delta\Theta_x+K_0\delta$, where
$P_j=\vartheta_x+f_j$ and $Q_j=\eta_j$ for $j\in I_0$, and $P_j=\vartheta_x+f_{k_0}-\delta$ and $Q_j=\eta_{k_0}$ otherwise. The
Levi forms of $P_j$, of $Q_j+K_1\Theta_x$ and of $\Theta_x$ are nonnegative, by (b), by the choice of the $c_j$ and as
$\omega$ is positive; and $K\delta=(8K_1+4)\delta\ge4\delta K_1$. So the Levi form of every $v_j$ is nonnegative, and by
Lemma~\ref{cv-origin-analytic-system-weights-regularize-global-max}(4), with $g=0$, so is that of
$\vartheta_x+F_t+K\delta\Theta_x$. That is, the first Chern form of $\bar L$ with its metric multiplied by $e^{-F_t}$ is at
least $-K\delta\omega$.

\emph{Convergence.} Let $\psi$ be a chart of $S$ and $z_0$ a point of its image with $\psi^{-1}(z_0)\notin\partial$. Take
$y_0=\psi^{-1}(z_0)$, with $k_0$, $I_0$ and $W$ as above, and a closed ball $N$ about $z_0$ in $\psi(W\setminus\partial)$.
Let $\beta(j)=j$ for $j\in I_0$ and $\beta(j)=k_0$ otherwise, and let $Q'_j=4\eta_{\beta(j)}\circ\psi^{-1}$, minus $1$ if
$j\notin I_0$. By $(*)$ and Lemma~\ref{cv-origin-analytic-system-weights-regularize-global-max}(1), near $N$
\[
 F_t\circ\psi^{-1}=\mathcal M_{\delta(t)}((u_{t,j})_j)+\delta(t)/2+(K_0-1/2)\delta(t),\qquad
 u_{t,j}=f_{t,\beta(j)}\circ\psi^{-1}+\delta(t)Q'_j.
\]
As $W\subset B_{\beta(j)}(5/4)$, the change of charts $\kappa_{\beta(j)}\circ\psi^{-1}$ is holomorphic near $N$ and maps $N$
into a compact subset of $\Delta_{5/4}$ whose image under $\kappa_{\beta(j)}^{-1}$ does not meet $\partial$. By the chain
rule, $f_{t,\beta(j)}\circ\psi^{-1}$ tends to $\varphi\circ\psi^{-1}$ with its first and second derivatives uniformly on
$N$, with $\|d(f_{t,\beta(j)}\circ\psi^{-1})-d(\varphi\circ\psi^{-1})\|\le M_jt$ there. As $Q'_j$ is smooth and
$\delta(t)\to0$, the same holds for $u_{t,j}$, with $\|du_{t,j}-d(\varphi\circ\psi^{-1})\|\le M_jt+\delta(t)\sup_N\|dQ'_j\|$,
which is at most $M''\delta(t)$ for a number $M''$, as $t\le\delta(t)$. As $(M''\delta(t))^2/\delta(t)\to0$,
Lemma~\ref{cv-origin-analytic-system-weights-regularize-global-max-limit} shows that
$\mathcal M_{\delta(t)}((u_{t,j})_j)+\delta(t)/2$ tends to $\varphi\circ\psi^{-1}$ with its first and second derivatives
uniformly on $N$, and so does $F_t\circ\psi^{-1}$, as $(K_0-1/2)\delta(t)\to0$. Finitely many such balls $N$ cover a given
compact subset of the image under $\psi$ of the complement of $\partial$.
\end{proof}
% Lean: regularization_glue_C2_explicit = Giv.regularization_glue_C2 (RegularizationGlueC2Explicit.lean): glue_cutoff
% (glueCut, polydiscCutoff), reg_glue_explicit (pshShift for the Cj, K' = sum |Cj|, K0 = 4 K' + 2; glueArg with
% -Mb - 10 off chartBall (5/4); glue_local_unif and glue_loc_rep with regMax_dominated for (*); the bound hbd; the Levi
% form through regMaxLevi_thm with c = 2 K0 delta), F t = regMax (delta t) (glueArg ...) + K0 delta t, K = 2 K0, and
% glue_C2_local (chart_C2_convergence, tendstoC2On_add_small, RegularizedMaxConvergenceC2.proof with M' delta,
% tendstoC2On_congr) with tendstoLocallyUniformlyOn_of_nhds. Size: about 750 lines here, about 2,600 with the lemmas it
% uses.

\begin{lemma}\label{cv-origin-analytic-system-weights-regularize-global-bounds}
Let $S$, $\omega$, $s$, $\partial$ and $r$ be as in Definition~\ref{cv-origin-analytic-system-weight}(1), with $e=\dim S$,
let $\bar L$ be a holomorphic line bundle on $S$ with a smooth Hermitian metric, and let $\varphi$ be a function on
$S\setminus\partial$ with $\varphi\ge-C_5\log r$ there, for a number $C_5$. Let $\kappa_1,\dots,\kappa_N$ be charts of $S$,
let $B_j(\rho)$ be the set of points of the domain of $\kappa_j$ that $\kappa_j$ maps into $\Delta_\rho$, and suppose that
the $B_j(1)$ cover $S$. Let $K>0$, $C_l\ge1$ and $0<t_0\le\exp(-\exp(1))$. For $0<t\le t_0$ let $F_t$ be a smooth function on
$S$ and $f_{t,j}$ a function on $B_j(5/4)$, for every $j$, such that for $0<t\le t'\le t_0$ and every $j$:
\begin{enumerate}
\item[(a)] the first Chern form of $\bar L$ with its metric multiplied by $e^{-F_t}$ is at least $-K\varepsilon_{C_l}(t)\omega$,
  and $f_{t,j}\le F_t\le f_{t,j}+K\varepsilon_{C_l}(t)$ on $B_j(5/4)$;
\item[(b)] $-C_l\log\log(1/t)\le f_{t,j}\le C_l$ on $B_j(5/4)$, and every $y\in B_j(5/4)\setminus\partial$ has
  $\varphi(y)\le f_{t,j}(y)$, and $f_{t,j}(y)\le\varphi(y)+t^{1/3}$ if $r(y)\le\log(1/t)$;
\item[(c)] $f_{t,j}\le f_{t',j}+\varepsilon_{C_l}(t')$ on $B_j(5/4)$.
\end{enumerate}
Then there is $C$ such that the functions $F_t$ satisfy the conditions of
Definition~\ref{cv-proportionality-analytic-system}(1) in which $\varphi_b(t)$ appears, except the convergence as $t\to0$,
with $\bar L$, $\varphi$ and $F_t$ in place of $\bar L_b$, $\varphi_b$ and $\varphi_b(t)$, with $s$ in place of $\sigma$,
and with these $\omega$, $C$ and $t_0$.
\end{lemma}
\leanref{Development/AutomorphicSpectrum/Proofs/MetricProofs/RegularizationGlueBoundsStatement.lean}{GlobalTheory.T.regularization_glue_bounds}
\begin{proof}
Put $C=C_l(1+K)+|C_5|+1\ge1$. Let $0<t\le t_0$, and put $g=\varepsilon_1(t)$, so that $0\le g\le1$ and
$\varepsilon_{C'}(t)=C'g$ for every $C'$, by Lemma~\ref{cv-origin-analytic-system-weights-regularize-global-numbers}. Let
$y\in S$, and choose $j$ with $y\in B_j(1)\subset B_j(5/4)$; write $f=f_{t,j}(y)$.
\begin{enumerate}
\item[(i)] $F_t(y)\le f+KC_lg\le C_l(1+K)\le C$.
\item[(ii)] If $y\notin\partial$, then $\varphi(y)\le f\le F_t(y)$.
\item[(iii)] $F_t(y)\ge f\ge-C_l\log\log(1/t)\ge-C\log\log(1/t)$, as $\log\log(1/t)\ge0$.
\item[(iv)] As $K\varepsilon_{C_l}(t)=KC_lg\le Cg=\varepsilon_C(t)$ and $\omega$ is positive, the first Chern form of $\bar L$
  with its metric multiplied by $e^{-F_t}$ is at least $-\varepsilon_C(t)\omega$.
\item[(v)] If $y\notin\partial$ and $r(y)\le\log(1/t)$, then $F_t(y)\le f+KC_lg\le\varphi(y)+t^{1/3}+Cg$; and, as $r\ge1$,
  $F_t(y)\le\varphi(y)+t^{1/3}+Cg\,r(y)$.
\item[(vi)] Let $y\notin\partial$ with $r(y)>\ell=\log(1/t)$. As $\log x/x$ decreases for $x\ge\exp(1)$,
  $\log r(y)\le g\,r(y)$, and $g\,r(y)=\log\ell\cdot r(y)/\ell\ge1$. So $\varphi(y)\ge-|C_5|\log r(y)\ge-|C_5|g\,r(y)$, and
  $F_t(y)\le C_l(1+K)\le C_l(1+K)g\,r(y)$; hence $F_t(y)-\varphi(y)\le Cg\,r(y)\le t^{1/3}+\varepsilon_C(t)r(y)$.
\item[(vii)] For $t\le t'\le t_0$, with $g'=\varepsilon_1(t')\ge g$: $F_t(y)\le f+KC_lg\le f_{t',j}(y)+C_lg'+KC_lg'\le
  F_{t'}(y)+Cg'\le F_{t'}(y)+C\varepsilon_C(t')$, as $C\ge1$.
\end{enumerate}
Finally $0<t_0<1$.
\end{proof}
% Lean: regularization_glue_bounds = RegFin.isRegFamily_of_glue (RegularizationGlueBounds.lean): C = Cl (1 + K) + |C5| + 1,
% g = log log(1/t) / log(1/t) (quot_mem_unit, quot_mono, epsT_eq_mul), isPsh_rsmul_mono for (iv), the case split on
% rW <= log(1/t) with Real.log_div_self_antitoneOn for (v) and (vi). Size: about 320 lines.

\begin{lemma}\label{cv-origin-analytic-system-weights-regularize-global}
Let $S$, $\omega$, $s$, $\partial$ and $r$ be as in Definition~\ref{cv-origin-analytic-system-weight}(1), and let
$\kappa_1,\dots,\kappa_N$, $T_j$, $\mu_{j,i}$ and $M$ be as in
Lemma~\ref{cv-origin-analytic-system-weights-regularize-atlas}. Let $\bar L$ be a holomorphic line bundle on $S$ with a
smooth Hermitian metric, let $C_5$ be a number, and let $\varphi$ be a smooth function on $S\setminus\partial$ such that
the metric of $\bar L$ multiplied by $e^{-\varphi}$ is semipositive on $S\setminus\partial$, $-C_5\log r\le\varphi\le C_5$
there, and $\|d(\varphi\circ\kappa_j^{-1})(y)\|\le C_5\,r(x')^{C_5}e^{r(x')/2}$ for every $j$ and every $y\in\Delta_{3/2}$ with
$x'=\kappa_j^{-1}(y)\notin\partial$. Then there are numbers $C$ and $0<t_0<1$ and smooth functions $\varphi(t)$ on $S$ that
regularize $\varphi$ on $\bar L$ with constants $(C,t_0)$ (Definition~\ref{cv-origin-analytic-system-weights-family}).
\end{lemma}
\leanref{Development/AutomorphicSpectrum/Proofs/MetricProofs/GlobalRegularizationStatementHoldsStatement.lean}{GlobalTheory.T.global_regularization_statement_holds}
\begin{proof}
\emph{The atlas.} Let $\vartheta_x$ be local potentials of the metric of $\bar L$, each defined on an open neighbourhood
$U_x$ of $x$. Lemma~\ref{cv-origin-analytic-system-weights-regularize-global-refine}, applied to the $\kappa_j$, the
$U_x$, $\varphi$ and $C_5$, gives charts with the properties (1)--(4) of
Lemma~\ref{cv-origin-analytic-system-weights-regularize-atlas}, for sets $T_j$, integers $\mu_{j,i}\ge1$ and a number $M$,
and points $x_j$ such that the domain of the $j$-th chart lies in $U_{x_j}$ and the bound on the derivative of $\varphi$
holds with $|C_5|$. Replace the $\kappa_j$, $T_j$, $\mu_{j,i}$ and $M$ by these. Write $r_j=r_{T_j,\mu_j}$ and $H_j=H_{T_j}$
(Definition~\ref{cv-origin-analytic-system-weights-regularize-global-notation}), and $B_j(\rho)$ for the set of points of
the domain of $\kappa_j$ that $\kappa_j$ maps into $\Delta_\rho$. By (3), $\Delta_2\setminus H_j$ is the set of $y\in\Delta_2$
with $\kappa_j^{-1}(y)\notin\partial$.

\emph{The weight in the charts.} Put $q_j=\vartheta_{x_j}\circ\kappa_j^{-1}$, smooth on $\Delta_2$, and
$\varphi_j=\varphi\circ\kappa_j^{-1}$, smooth on $\Delta_2\setminus H_j$. The function $q_j+\varphi_j$ has nonnegative Levi
form on $\Delta_2\setminus H_j$: near each point it is the composite of $\vartheta_{x_j}+\varphi$, read in a chart of $S$,
which has nonnegative Levi form off $\partial$, with a holomorphic change of charts \cite[Chapter~I, (5.10)]{Dem12}.
Lemma~\ref{cv-origin-analytic-system-weights-regularize-global-transfer}, with $C_5$ and $M$, gives $C'_5\ge0$. On
$\Delta_2\setminus H_j$, $r_j\ge\exp(1)$, as every term of $r_j$ is nonnegative on $\bar\Delta_2$;
$r\circ\kappa_j^{-1}\ge\exp(1)$, as $\|s\|\le1$; and $|r\circ\kappa_j^{-1}-r_j|\le M$ by (4). So
$-C'_5\log r_j\le\varphi_j\le C'_5$ on $\Delta_2\setminus H_j$, and $\|d\varphi_j\|\le C'_5r_j^{C'_5}e^{r_j/2}$ on
$\Delta_{3/2}\setminus H_j$.

\emph{Constants.} Lemma~\ref{cv-origin-analytic-system-weights-regularize-global-chart}, applied for each $j$ to $T_j$,
$\mu_j$, $q_j$, $\varphi_j$, $C'_5$ and $K=M$, gives $C_j\ge1$, $0<t_j\le\exp(-\exp(1))$ and smooth functions $g_{j,t}$ on
$\Delta_{5/4}$ with its properties (a)--(e) and its convergence.
Lemma~\ref{cv-origin-analytic-system-weights-regularize-global-compare}, applied to the charts $\kappa_j$ and the points
$x_j$, whose functions $u_j$ are the $q_j+\varphi_j$, gives $C'$ and $t_c>0$. Put $C_m=1+\sum_j|C_j|$ and $C_l=|C'|+C_m$, so
that $C_j\le C_m\le C_l$ and $C_l\ge1$. By Lemma~\ref{cv-origin-analytic-system-weights-regularize-global-numbers} there
is $a>0$ with $\varepsilon_{C_l}(t)\le1$ for $0<t\le a$. Let $t_0$ be the least of the $t_j$, $t_c$, $a$ and
$\exp(-\exp(1))$.

\emph{The local regularizations.} For $0<t\le t_0$ put $f_{t,j}=g_{j,t}\circ\kappa_j$ on $B_j(5/4)$. Let $0<t\le t_0$.
\begin{enumerate}
\item[(i)] For $y\in B_j(5/4)$, put $m_j(y)=m_{T_j}(q_j+\varphi_j;\kappa_j(y),t)$. By (d) of
  Lemma~\ref{cv-origin-analytic-system-weights-regularize-global-chart},
  $m_j(y)-\varepsilon_{C_j}(t)\le\vartheta_{x_j}(y)+f_{t,j}(y)\le m_j(y)$.
\item[(ii)] By (b) and (c) of that lemma, as $C_j\le C_l$ and $\log\log(1/t)\ge1$,
  $-C_l\log\log(1/t)\le f_{t,j}\le C_l$ on $B_j(5/4)$; and every $y\in B_j(5/4)\setminus\partial$ has
  $\varphi(y)\le f_{t,j}(y)$, and, if $r(y)\le\log(1/t)$, then $r_j(\kappa_j(y))\le r(y)+M\le\log(1/t)+M$, so that
  $f_{t,j}(y)\le\varphi(y)+t^{1/3}$.
\item[(iii)] By (e) of that lemma and
  Lemma~\ref{cv-origin-analytic-system-weights-regularize-global-numbers}(2), $f_{t,j}\le f_{\tau,j}+\varepsilon_{C_l}(\tau)$
  on $B_j(5/4)$ for $t\le\tau\le t_0$.
\item[(iv)] By (i), Lemma~\ref{cv-origin-analytic-system-weights-regularize-global-compare} and
  Lemma~\ref{cv-origin-analytic-system-weights-regularize-global-numbers}(2), on $B_j(5/4)\cap B_k(5/4)$,
  $|f_{t,j}-f_{t,k}|\le\varepsilon_{C'}(t)+\varepsilon_{C_m}(t)\le\varepsilon_{|C'|}(t)+\varepsilon_{C_m}(t)
  =\varepsilon_{C_l}(t)$.
\item[(v)] By (a) of Lemma~\ref{cv-origin-analytic-system-weights-regularize-global-chart}, $f_{t,j}$ is smooth on
  $B_j(5/4)$, and the metric of $\bar L$ multiplied by $e^{-f_{t,j}}$ is semipositive there: near each point, a local
  potential $\vartheta_x$ differs from $\vartheta_{x_j}$ by the real part of a holomorphic function, whose Levi form is $0$,
  and $\vartheta_{x_j}+f_{t,j}=(q_j+g_{j,t})\circ\kappa_j$ has nonnegative Levi form in every chart, by
  \cite[Chapter~I, (5.10)]{Dem12}.
\item[(vi)] Let $K'\subset\Delta_{5/4}$ be compact with $\kappa_j^{-1}(K')\cap\partial=\emptyset$, so that
  $K'\subset\Delta_{5/4}\setminus H_j$. By Lemma~\ref{cv-origin-analytic-system-weights-regularize-global-chart},
  $f_{t,j}\circ\kappa_j^{-1}=g_{j,t}$ tends to $\varphi_j=\varphi\circ\kappa_j^{-1}$ with its first and second derivatives
  uniformly on $K'$ as $t\to0$, with $\|dg_{j,t}-d\varphi_j\|\le M't$ there for small $t$.
\end{enumerate}

\emph{Gluing.} By (2), the $B_j(1)$ cover $S$; by
Lemma~\ref{cv-origin-analytic-system-weights-regularize-global-numbers}, $t\le\varepsilon_{C_l}(t)\le1$ for $0<t\le t_0$, and
$\varepsilon_{C_l}(t)\to0$ as $t\to0$. So, by (ii), (iv), (v) and (vi),
Lemma~\ref{cv-origin-analytic-system-weights-regularize-global-glue} applies to $\partial$, $\varphi$, the charts
$\kappa_j$, $\delta(t)=\varepsilon_{C_l}(t)$ and the $f_{t,j}$, whose absolute values are at most
$C_l+C_l\log\log(1/t)$; it gives $K>0$ and smooth functions $F_t$ on $S$.
Lemma~\ref{cv-origin-analytic-system-weights-regularize-global-bounds}, applied with $K$, $C_l$, $t_0$, the $F_t$ and the
$f_{t,j}$, using (ii), (iii) and $\varphi\ge-C_5\log r$, gives $C$. With the convergence given by
Lemma~\ref{cv-origin-analytic-system-weights-regularize-global-glue}, the functions $\varphi(t)=F_t$ regularize $\varphi$ on
$\bar L$ with constants $(C,t_0)$, and $0<t_0<1$.
\end{proof}
% Lean: ReciprocityTheory.AnForm.P.RegGlobalStatement, proved as Giv.regGlobalStatement_holds
% (GlobalRegularizationStatementHolds.lean): refine_adapted_atlas_with_gradient, levi_nonneg_in_chart_of_psh (hlevi),
% log_bound_transfer_between_radii with rW_ge and rLoc_ge (hbnd, hgrad), local_levi_psh_regularization with K = M (hch),
% glued_local_regularizations_on_atlas (hcmp), Cm = 1 + sum |Cj|, Cl = |C'| + Cm, epsilon_schedule_properties for a and
% tm, t0 = min (min tm t') (min a (exp(-exp 1))), hsand, hloc, hmono', hov, smoothOn_of_chart_levi_nonneg (hback),
% regularization_glue_C2_explicit with delta = epsT Cl (C2_tendsto_pullback_through_chart for the convergence), and
% regularization_glue_bounds. Size: about 190 lines here, about 8,400 with the lemmas above.

\begin{lemma}\label{cv-origin-analytic-system-weights-regularize}
Let $S$, $\omega$, $s$, $\partial$ and $r$ be as in Definition~\ref{cv-origin-analytic-system-weight}(1), let $\bar L$ be a
holomorphic line bundle on $S$ with a smooth Hermitian metric, let $C_5\ge0$, and let $\varphi$ be a smooth function on
$S\setminus\partial$ such that the metric of $\bar L$ multiplied by $e^{-\varphi}$ is semipositive on $S\setminus\partial$
and $-C_5\log r\le\varphi\le C_5$ there. Suppose that every $x\in S$ has a chart $\psi$, an open neighbourhood $V$ of $x$
in its domain and a number $c$ with $\|d(\varphi\circ\psi^{-1})(\psi(x'))\|\le c\,r(x')^ce^{r(x')/2}$ for
$x'\in V\setminus\partial$. Then there are numbers $C$ and $0<t_0<1$ and smooth functions $\varphi(t)$ on $S$ that
regularize $\varphi$ on $\bar L$ with constants $(C,t_0)$ (Definition~\ref{cv-origin-analytic-system-weights-family}).
\end{lemma}
\leanref{Development/AutomorphicSpectrum/Proofs/GeometryProofs/AdaptedAtlasSncSectionStatement.lean}{GlobalTheory.T.adapted_atlas_snc_section}
\leanref{Development/AutomorphicSpectrum/Proofs/MetricProofs/AdaptedAtlasSmoothPotentialLocalStatement.lean}{GlobalTheory.T.adapted_atlas_smooth_potential_local}
\leanref{Development/AutomorphicSpectrum/Proofs/MetricProofs/GlobalRegularizationStatementHoldsStatement.lean}{GlobalTheory.T.global_regularization_statement_holds}
\begin{proof}
Lemma~\ref{cv-origin-analytic-system-weights-regularize-atlas} gives charts $\kappa_1,\dots,\kappa_N$ with its properties,
and Lemma~\ref{cv-origin-analytic-system-weights-regularize-gradient}, applied to them, to $\varphi$ and to the local
bounds of the hypothesis, gives $C_6$. Put $C'=\max(C_5,C_6)$. As $\|s\|\le1$, $\log r\ge1$, so
$-C'\log r\le-C_5\log r\le\varphi\le C_5\le C'$ on $S\setminus\partial$, and the bounds on the derivatives in the charts
$\kappa_j$ hold with $C'$, as $C'\ge C_6$. Lemma~\ref{cv-origin-analytic-system-weights-regularize-global}, applied with
$C'$ in place of $C_5$, gives $C$, $t_0$ and the functions $\varphi(t)$.
\end{proof}
% Lean: the case not cen of exists_member: adapted_atlas_snc_section, adapted_atlas_smooth_potential_local with the local
% bounds of local_of_tame, the singular weight with max C5 C6 (the first clause of enlarging_constants_smooth_members), and
% RegGlobalStatement, which exists_member takes as a hypothesis hR and the Lean proves as
% GlobalTheory.T.global_regularization_statement_holds.

\begin{lemma}\label{cv-origin-analytic-system-weights-symmetric}
Let $S$, $\omega$, $s$, $\partial$ and $r$ be as in Definition~\ref{cv-origin-analytic-system-weight}(1), and let
$c:S\to S$ be a continuous map with $c(c(x))=x$ that is antiholomorphic, in the sense of
Lemma~\ref{cv-origin-analytic-system-weights-conjugation}(1), with $c(x)\in\partial$ if and only if $x\in\partial$ and
$r(c(x))=r(x)$. Let $\bar L$ be a holomorphic line bundle on $S$ with a smooth Hermitian metric with local potentials
$\vartheta_x$, such that every $x$ has a neighbourhood on which $\vartheta_{c(x)}\circ c-\vartheta_x$ is the real part of
a holomorphic function. Let $\varphi$ be a function on $S\setminus\partial$ with $\varphi\circ c=\varphi$, and let smooth
functions $\varphi(t)$ regularize $\varphi$ on $\bar L$ with constants $(C,t_0)$
(Definition~\ref{cv-origin-analytic-system-weights-family}). Then there are $C'\ge0$ and $t'_0$ such that the functions
$(\varphi(t)+\varphi(t)\circ c)/2$ regularize $\varphi$ on $\bar L$ with constants $(C',t'_0)$.
\end{lemma}
\leanref{Development/AutomorphicSpectrum/Proofs/MetricProofs/ConjugationInvariantSingularPotentialStatement.lean}{GlobalTheory.T.conjugation_invariant_singular_potential}
\begin{proof}
Let $\Theta_x$ be local potentials of $\omega$, and write $\varepsilon_C(t)=C\log\log(1/t)/\log(1/t)$. A function $g$
is plurisubharmonic near $c(x)$ if and only if $g\circ c$ is plurisubharmonic near $x$: in charts $c$ is the complex
conjugation composed with a holomorphic map, and the Levi form of $g(\bar v)$ at $v$ on $w$ is that of $g$ at $\bar v$
on $\bar w$. So the $\Theta_{c(x)}\circ c$ are plurisubharmonic, smooth, and their differences are real parts of
holomorphic functions, by Lemma~\ref{cv-origin-analytic-system-weights-conjugation}(1). As $S$ is compact and $\omega$ is
positive, there is $K\ge0$ such that every $K\Theta_x-\Theta_{c(x)}\circ c$ is plurisubharmonic. Put $M=\max(1,K)$,
$C'=|C|M$ and $t'_0=\min(t_0,\exp(-\exp(1)))$. The argument of the proof of
Lemma~\ref{cv-origin-analytic-system-enlarge}(1) for the conditions involving $\varphi(t)$ uses only that the new
constant is at least $|C|$; so, restricted to $0<t\le t'_0$, the $\varphi(t)$ regularize $\varphi$ with constants
$(|C|,t'_0)$ and $(C',t'_0)$.

Let $0<t\le t'_0$ and $\varphi'(t)=(\varphi(t)+\varphi(t)\circ c)/2$, a smooth function. The bounds
$\varphi'(t)\le C'$, $\varphi\le\varphi'(t)$, $\varphi'(t)\ge-C'\log\log(1/t)$,
$\varphi'(t)\le\varphi+t^{1/3}+\varepsilon_{C'}(t)r$, $\varphi'(t)\le\varphi+t^{1/3}+\varepsilon_{C'}(t)$ where
$r\le\log(1/t)$, and $\varphi'(t)\le\varphi'(t')+C'\varepsilon_{C'}(t')$ for $t\le t'\le t'_0$, are the averages of those
for $\varphi(t)$ with $(C',t'_0)$ at $x$ and at $c(x)$, as $\varphi(c(x))=\varphi(x)$, $r(c(x))=r(x)$ and
$c(x)\notin\partial$ for $x\notin\partial$. Put $\varepsilon=\varepsilon_{|C|}(t)$. As $t\le\exp(-\exp(1))$,
$\log\log(1/t)\ge1$ and $\varepsilon\ge0$, and $\varepsilon_{C'}(t)=M\varepsilon$. The functions
$\vartheta_x+\varphi(t)+\varepsilon\Theta_x$ are plurisubharmonic, by the condition on the first Chern form with
$(|C|,t'_0)$, and so are their composites with $c$. Near $x$, with $\vartheta_{c(x)}\circ c-\vartheta_x=\operatorname{Re}G$,
\begin{align*}
 \vartheta_x+\varphi'(t)+M\varepsilon\Theta_x
 ={}&\tfrac12\big(\vartheta_x+\varphi(t)+\varepsilon\Theta_x\big)
   +\tfrac12\big(\vartheta_{c(x)}+\varphi(t)+\varepsilon\Theta_{c(x)}\big)\circ c\\
 &+\tfrac\varepsilon2\big(K\Theta_x-\Theta_{c(x)}\circ c\big)
   +\big(M\varepsilon-\tfrac{(K+1)\varepsilon}2\big)\Theta_x-\tfrac12\operatorname{Re}G.
\end{align*}
Each term is plurisubharmonic, as $M\ge(K+1)/2$, or pluriharmonic. So the first Chern form of the metric of $\bar L$
multiplied by $e^{-\varphi'(t)}$ is at least $-\varepsilon_{C'}(t)\omega$. Finally, $c$ is a diffeomorphism of
$S\setminus\partial$, so $\varphi(t)\circ c$ tends to $\varphi\circ c=\varphi$ with its first and second derivatives,
uniformly on compact subsets of $S\setminus\partial$, as $t\to0$, and so does the average $\varphi'(t)$. As
$0<t'_0<1$, the $\varphi'(t)$ regularize $\varphi$ with constants $(C',t'_0)$.
\end{proof}
% Lean: PositivityTheory.Xgp.t_an_conj: TailBoundAlt.exists_psh_chern_bound for K (with Theta.pullback c, which
% an_conj.isPotential_pull makes a potential), M = max 1 K, C' = |C| M, t1 = min t0 (exp(-exp 1)), Giv.isRegFamily_mono
% twice (F1 with |C|, F2 with C'), smoothOn_pull, the averages of the bounds a2 to a8, the plurisubharmonicity through
% isPsh_pull (levi_cj) and isPsh_of_same with the identity displayed, and regular_family_average with
% regular_family_pullback_involution for the convergence. Size: about 1,200 lines.

\begin{lemma}\label{cv-origin-analytic-system-weights}
Let $\Omega$ be a field, let $y\in\Omega^{\sqcup_{k\in K}I_k}$, and let $\tilde X$ with maps $\lambda_x$ be as in
Definition~\ref{cv-origin-analytic}(3) for $y$, of dimension $e$. Let $s$ be a holomorphic section of a holomorphic
line bundle on $\tilde X$ with a smooth Hermitian metric, with $\|s\|\le1$, whose divisor is $\sum_lm_lW_l$ for smooth
hypersurfaces $W_l$ crossing normally and integers $m_l\ge0$; let $\partial$ be its zero set, and
$r=\exp(1)-2\log\|s\|$. Suppose that whenever $\lambda_{x'}(x')$ lies in the orbit of $\overline{\lambda_x(x)}$ under
$(\mathbf C^\times)^K$, $x'\in\partial$ if and only if $x\in\partial$, and $\|s(x')\|=\|s(x)\|$. Let $\omega$ be the
K\"ahler form whose local potentials are the pullback of $\bar{\mathcal O}(\mathbf 1)$
(Lemma~\ref{cv-origin-analytic-system-positive}(2)). Let $d\in\mathbf Z^K$, write $\vartheta_x$ for the functions of
the pullback of $\bar{\mathcal O}(d)$, and let $u$ be a real function on $\tilde X$ such that:
\begin{enumerate}
\item[(a)] $u(x')=u(x)$ whenever $\lambda_{x'}(x')$ lies in the orbit of $\overline{\lambda_x(x)}$;
\item[(b)] every $x\in\tilde X$ has $\psi$, $a$ and $M$ as in Lemma~\ref{cv-origin-analytic-system-charts}(1) for
  $U=U_x$ such that $u\circ\psi^{-1}$ is tame for $a$ (Definition~\ref{cv-origin-analytic-system-weight}(2));
\item[(c)] every $x\in\tilde X\setminus\partial$ has a chart $\psi$ at $x$ with $\psi(x)=0$ such that
  $f_x(t)=\vartheta_x(\psi^{-1}(t))+u(\psi^{-1}(t))$ is of class $C^2$ near $0$ with nonnegative Levi form at $0$.
\end{enumerate}
Then:
\begin{enumerate}
\item there are a regularized weight $(\varphi,\varphi(t))$ with some constants $(C,t_0)$ on the pullback of
  $\bar{\mathcal O}(d)$ (Definition~\ref{cv-origin-analytic-system-weight}(1)), with $\varphi=u$ on
  $\tilde X\setminus\partial$, and for every real $t$ a smooth metrized line bundle $\bar A(t)=(d,h_{A(t)})$, invariant
  under conjugation, with $\log h_{A(t)}(\lambda_x(x))=\log h_d(\lambda_x(x))+\varphi(t)(x)$ for $0<t\le t_0$ and
  $x\in\tilde X$;
\item if in (c) also $-f_x$ is of class $C^2$ near $0$ with nonnegative Levi form at $0$, then these can be chosen
  such that moreover, for $0<t\le t_0$, the functions $\vartheta_x+\varphi(t)$ are pluriharmonic.
\end{enumerate}
\end{lemma}
\leanref{Development/WeilBound/CoreStatements/CentralCharts/CommonMembers.lean}{AnalyticReference.An.exists_member}
\begin{proof}
Let $c$ be the involution of Lemma~\ref{cv-origin-analytic-system-weights-conjugation}. For every $x$,
$\lambda_{c(x)}(c(x))$ lies in the orbit of $\overline{\lambda_x(x)}$; so by hypothesis $c(x)\in\partial$ if and only if
$x\in\partial$, $r(c(x))=r(x)$, and $u(c(x))=u(x)$ by (a). By Lemma~\ref{cv-origin-analytic-system-positive}(1), the
$\vartheta_x$ are the local potentials of a smooth Hermitian metric on a holomorphic line bundle $\bar L$ on $\tilde X$,
and by Lemma~\ref{cv-origin-analytic-system-weights-conjugation}(2), applied to $\bar{\mathcal O}(d)$, which is invariant
under conjugation, every $x$ has a neighbourhood on which $\vartheta_{c(x)}\circ c-\vartheta_x$ is the real part of a
holomorphic function. With $\omega$ and $s$, $\tilde X$ is as in Definition~\ref{cv-origin-analytic-system-weight}(1).

\emph{The weight $u$.} By (b) and Lemma~\ref{cv-origin-analytic-system-weights-tame}, $u$ is smooth on
$\tilde X\setminus\partial$, $|u|\le C_1\log r$ there for some $C_1$, and $u$ has the bounds on its derivative of
Lemma~\ref{cv-origin-analytic-system-weights-tame}(2). Let $z\in U_x\setminus\partial$, let $\psi$ be the chart at $z$
given by (c), and let $\chi$ be a chart at $z$. Near $z$, $\vartheta_x-\vartheta_z$ is the real part of a holomorphic
function $G$ (Lemma~\ref{cv-origin-analytic-system-positive}(1)), so near $\chi(z)$
\[
 (\vartheta_x+u)\circ\chi^{-1}=f_z\circ(\psi\circ\chi^{-1})+\operatorname{Re}(G\circ\chi^{-1}).
\]
The Levi form of the second term is $0$, and that of the first at $\chi(z)$ on a vector $v$ is the Levi form of $f_z$ at
$0$ on the image of $v$ under the derivative of the holomorphic map $\psi\circ\chi^{-1}$, which is nonnegative by (c).
So the metric of $\bar L$ multiplied by $e^{-u}$ is semipositive on $\tilde X\setminus\partial$, and by
Lemma~\ref{cv-origin-analytic-system-weights-singular} there is $C_5\ge0$ with $-C_5\log r\le u\le C_5$ on
$\tilde X\setminus\partial$.

\emph{The bundles.} For a smooth function $f$ on $\tilde X$ with $f\circ c=f$,
Lemma~\ref{cv-origin-analytic-system-weights-line} gives a smooth metrized line bundle of multidegree $d$, invariant
under conjugation, with $\log h(\lambda_x(x))=\log h_d(\lambda_x(x))+f(x)$ for every $x$.

\emph{Case (2).} Suppose that, for every $x\notin\partial$, $-f_x$ is also of class $C^2$ near $0$ with nonnegative Levi
form at $0$. By the argument above, $-(\vartheta_x+u)$ is then plurisubharmonic on $\tilde X\setminus\partial$, and
Lemma~\ref{cv-origin-analytic-system-weights-flat} gives a smooth function $u'$ on $\tilde X$ with $u'=u$ on
$\tilde X\setminus\partial$ such that $\vartheta_x+u'$ and $-(\vartheta_x+u')$ are plurisubharmonic, that is,
$\vartheta_x+u'$ is pluriharmonic. By Lemma~\ref{cv-origin-analytic-system-enlarge}(2), applied to $u'$, there is
$C_0\ge0$ such that $\varphi=u'=u$ on $\tilde X\setminus\partial$ and $\varphi(t)=u'$ form a regularized weight with
constants $(C_0,\exp(-\exp(1)))$. The set $\partial$ has empty interior, as near each of its points it is a union of
coordinate hyperplanes of a chart; so $\tilde X\setminus\partial$ is dense. The continuous functions $u'\circ c$ and
$u'$ agree on $\tilde X\setminus\partial$, by (a), as $c$ maps $\tilde X\setminus\partial$ to itself; so $u'\circ c=u'$.
For every $t$ let $\bar A(t)$ be the bundle given by the previous paragraph for $f=u'$.

\emph{Case (1).} Otherwise, Lemma~\ref{cv-origin-analytic-system-weights-regularize}, applied to $\bar L$, to $\varphi=u$
on $\tilde X\setminus\partial$ with $C_5$, and to the bounds of Lemma~\ref{cv-origin-analytic-system-weights-tame}(2),
gives $C$, $t_0$ and smooth functions $\varphi(t)$ that regularize $u$ on $\bar L$ with constants $(C,t_0)$.
Lemma~\ref{cv-origin-analytic-system-weights-symmetric}, applied to $c$, gives $C'\ge0$ and $t'_0$ such that the functions
$\varphi'(t)=(\varphi(t)+\varphi(t)\circ c)/2$ regularize $u$ on $\bar L$ with constants $(C',t'_0)$. As $c(c(x))=x$,
$\varphi'(t)\circ c=\varphi'(t)$. For $0<t\le t'_0$ let $\bar A(t)$ be the bundle given by the paragraph on the bundles for
$f=\varphi'(t)$, and for the other $t$ let $\bar A(t)=\bar{\mathcal O}(d)$.

\emph{The constants.} In either case, $\varphi=u$ on $\tilde X\setminus\partial$ satisfies the conditions on $\varphi$ of
a regularized weight with a constant $C_a\ge0$ ($C_0$, or $C_5$), and the functions $\varphi(t)$ ($u'$, or
$\varphi'(t)$) regularize it on $\bar L$ with constants $(C_b,t_b)$ with $C_b\ge0$ ($(C_0,\exp(-\exp(1)))$, or
$(C',t'_0)$). The proof of Lemma~\ref{cv-origin-analytic-system-enlarge}(1) treats the conditions on $\varphi$, with the
first constant, and those involving $\varphi(t)$, with the second, separately. So, with $C=\max(C_a,C_b)$ and
$t_0=\min(t_b,\exp(-1))$, $\varphi$ and the $\varphi(t)$ for $0<t\le t_0$ form a regularized weight on the pullback of
$\bar{\mathcal O}(d)$ with constants $(C,t_0)$, and $\log h_{A(t)}(\lambda_x(x))=\log h_d(\lambda_x(x))+\varphi(t)(x)$ for
$0<t\le t_0$. This proves (1), and in Case (2) also (2), as there $\vartheta_x+\varphi(t)=\vartheta_x+u'$ is
pluriharmonic.
\end{proof}
% Lean: AnalyticReference.An.exists_member, under RegGlobalStatement (proved as
% GlobalTheory.T.global_regularization_statement_holds and carried in PB.Inputs): fs_pos_potential (hA, hfs, hTheta);
% lift_system_complex_conjugation and conjugation_antiholomorphic_charts (Lemma cv-origin-analytic-system-weights-conjugation);
% smoothOn_of_tame and size_of_tame (Lemma cv-origin-analytic-system-weights-tame); pshOn_of_local, with hth and hP,
% for the semipositivity; singular_potential_log_growth_psh (C5); hlines (exists_line, with fsHerm J d for the t outside
% (0, t0]). Case cen: the negative potential by pshOn_of_local, FA.central_regular (Flat.flat_weight_extends and
% enlarging_constants_smooth_members, Lemma cv-origin-analytic-system-enlarge(2)), and hcp by density
% (LeviFormBoundaryTopology.interior_bdry_eq_empty, Continuous.ext_on). Case not cen: adapted_atlas_snc_section,
% adapted_atlas_smooth_potential_local with local_of_tame, enlarging_constants_smooth_members and hR (Lemma
% cv-origin-analytic-system-weights-regularize); lift_system_real_structure (hthc) and
% conjugation_invariant_singular_potential (Lemma cv-origin-analytic-system-weights-symmetric). IsMember keeps C5 and C
% apart; the last paragraph combines them, as the proof of Lemma cv-origin-analytic-system does with common_members.
% Lemma cv-origin-analytic-system-weights-regularize-global (RegGlobalStatement), used by Lemma
% cv-origin-analytic-system-weights-regularize, is proved; no open step remains below this lemma.

\begin{lemma}\label{cv-origin-analytic-system}
In the setting of Definition~\ref{cv-origin-reference}, suppose $p\ge1$ and $q\ge1$, and let a reference
geometry of $\mathcal X$ be given, with metrized line bundles $\bar\sigma_{\tau,i}=(a_{\tau,i},h_{\tau,i})$,
invariant under conjugation, that are Hodge metrics on its root lines. Suppose that for every
$\gamma:\Delta\to\hat X(\mathbf C)$ and $a$ as in Lemma~\ref{cv-origin-analytic-local}, with $n=D$ and $B$ the
zero set of the boundary ideal $J_\partial$, the conclusions (1) and (2) of
Lemma~\ref{cv-origin-analytic-local} hold. Then there are analytic data on the reference geometry
(Definition~\ref{cv-origin-analytic}) and a holomorphic map $\psi:S\to\tilde X$ with $\psi\circ\pi$ the
identity.
\end{lemma}
\leanref{Development/WeilBound/CoreStatements/CentralCharts/OneLeDimFl.lean}{AnalyticReference.An.of_local}
\begin{proof}
Put $N=p+q\ge2$ and $\Sigma=(\Psi\times\{1,\dots,N\})\sqcup\{A\}$. As $\Psi$ has $[E:\mathbf Q]/2\ge1$ elements,
$D=|\Psi|N(N-1)/2\ge1$.

\emph{The manifold and the boundary.} By Lemma~\ref{cv-origin-analytic-system-arith}, the kernel $X$ of $y$ is an
arithmetic subvariety of dimension $D+1$. Let $B\subset\hat X(\mathbf C)$ be the set of points at which every element
of $J_\partial$ vanishes; by Lemma~\ref{cv-origin-analytic-system-boundary} it is closed and stable under
$(\mathbf C^\times)^K$ and under conjugation. As $J_\partial$ satisfies condition (3) of
Definition~\ref{cv-origin-reference}, Lemma~\ref{cv-origin-analytic-system-snc} gives a compact complex manifold
$\tilde X$ of dimension $D$ with maps $\lambda_x$ as in Definition~\ref{cv-origin-analytic}(3), and a holomorphic
section $s$, with $\|s\|\le1$, of a holomorphic line bundle with a smooth Hermitian metric, whose divisor is
$W_1+\dots+W_m$ for smooth hypersurfaces crossing normally, such that $s(x)=0$ if and only if $\lambda_x(x)\in B$,
and $\|s(x')\|=\|s(x)\|$ whenever $\lambda_{x'}(x')$ lies in the orbit of $\overline{\lambda_x(x)}$. Let $\partial$ be
the zero set of $s$, and $r=\exp(1)-2\log\|s\|$. Then:
\begin{enumerate}
\item[(a)] for $x'\in U_x$, $\lambda_x(x')$ lies in the zero set of $J_\partial$ if and only if $x'\in\partial$, as
  $\lambda_x(x')$ and $\lambda_{x'}(x')$ lie in one orbit and $B$ is stable under $(\mathbf C^\times)^K$;
\item[(b)] if $\lambda_{x'}(x')$ lies in the orbit of $\overline{\lambda_x(x)}$, then $x'\in\partial$ if and only if
  $x\in\partial$, as $B$ is stable under conjugation.
\end{enumerate}

\emph{Positivity.} By Lemma~\ref{cv-origin-analytic-system-positive}(1), the pullback of every smooth metrized line
bundle gives the local potentials of a smooth Hermitian metric on $\tilde X$. By part (2) the pullback of
$\bar{\mathcal O}(\mathbf 1)$ is positive; let $\omega$ be its K\"ahler form. By
Lemma~\ref{cv-origin-reference-immersion}, the hypothesis of part (3) holds for the index $k_0$ of
Definition~\ref{cv-origin-reference}(5), so the pullback of $\bar{\mathcal O}(e_{k_0})$ is positive.

\emph{The weights.} For $b=(\tau,i)$ put $d_b=-\sum_{k\le i}a_{\tau,k}$, and for $z\in\hat C$ put
$w_b(z)=-\sum_{k\le i}\log h_{\tau,k}(z)-\log h_{d_b}(z)$. For $c\in(\mathbf C^\times)^K$,
$h_{\tau,k}(cz)=\prod_j|c_j|^{a_{\tau,k,j}}h_{\tau,k}(z)$ and $h_{d_b}(cz)=\prod_j|c_j|^{d_{b,j}}h_{d_b}(z)$, so
$w_b(cz)=w_b(z)$; and $w_b(\bar z)=w_b(z)$, as the $\bar\sigma_{\tau,k}$ and $\bar{\mathcal O}(d_b)$ are invariant
under conjugation. Put $u_b(x)=w_b(\lambda_x(x))$; then $u_b(x')=w_b(\lambda_x(x'))$ for $x'\in U_x$. The hypotheses of
Lemma~\ref{cv-origin-analytic-system-weights} for $d_b$ and $u_b$ hold: those on $s$ by (b) and the choice of $s$,
and (a), (b) and (c) as follows.
\begin{enumerate}
\item[(a)] If $\lambda_{x'}(x')=c\,\overline{\lambda_x(x)}$, then $u_b(x')=w_b(\lambda_x(x))=u_b(x)$.
\item[(b)] Let $x\in\tilde X$, and let $\chi$, $a$ and $M$ be as in Lemma~\ref{cv-origin-analytic-system-charts}(1)
  for $U=U_x$. By (a) above and Lemma~\ref{cv-origin-analytic-system-charts}(2), $\gamma=\lambda_x\circ\chi^{-1}$ is a
  map as in Lemma~\ref{cv-origin-analytic-local}, with $n=D$ and $B$ the zero set of $J_\partial$. By hypothesis,
  conclusion (1) of that lemma holds: $t\mapsto-\sum_{k\le i}\log h_{\tau,k}(\gamma(t))$ is tame for $a$
  (Definition~\ref{cv-origin-analytic-system-weight}(2)). The function $t\mapsto\log h_{d_b}(\gamma(t))$ is smooth on
  $\Delta$, so it and its derivative are bounded on the closed polydisc of radius $\frac12$; as $\ell_a\ge\exp(1)$, it
  is tame for $a$, with $\varepsilon=\frac12$. The negative of a tame function is tame, and a sum of two functions
  that are tame for $a$ is tame for $a$, with the smaller $\varepsilon$ and the sum of the constants, as $\ell_a\ge1$.
  So $u_b\circ\chi^{-1}$, which is the difference of these two functions on $\Delta$, is tame for $a$.
\item[(c)] Let $x\notin\partial$, and let $\chi$ and $a$ be as in Lemma~\ref{cv-origin-analytic-system-charts}(1) for
  $U=U_x$. As $\chi^{-1}(0)=x\notin\partial$, $a=0$. By hypothesis, conclusion (2) of
  Lemma~\ref{cv-origin-analytic-local} holds for $\gamma=\lambda_x\circ\chi^{-1}$: the function
  $\log h_{d_b}(\gamma(t))+u_b(\chi^{-1}(t))=-\sum_{k\le i}\log h_{\tau,k}(\gamma(t))$ is of class $C^2$ near $0$ with
  nonnegative Levi form at $0$, and so is its negative if $i=N$.
\end{enumerate}
So Lemma~\ref{cv-origin-analytic-system-weights}, with its part (2) when $i=N$, gives a regularized weight
$(\varphi_b,\varphi_b(t))$ on the pullback of $\bar{\mathcal O}(d_b)$, with constants $(C_b,t_b)$ and with
$\varphi_b=u_b$ off $\partial$, and smooth metrized line bundles $\bar y_b(t)$ of multidegree $d_b$, invariant under
conjugation, with $\log h_{y_b(t)}(\lambda_x(x))=\log h_{d_b}(\lambda_x(x))+\varphi_b(t)(x)$ for $0<t\le t_b$; for
$i=N$, the functions $\vartheta_{\bar{\mathcal O}(d_b),x}+\varphi_b(t)$ are pluriharmonic. For $b=A$ put
$d_A=e_{k_0}$, $\varphi_A=0$, $\varphi_A(t)=0$ and $\bar y_A(t)=\bar{\mathcal O}(e_{k_0})$ for every $t$. As the
pullback of $\bar{\mathcal O}(e_{k_0})$ is positive, Lemma~\ref{cv-origin-analytic-system-enlarge}(2) shows that this
is a regularized weight with constants $(C_A,t_A)$ for some $C_A$ and $t_A=\exp(-\exp(1))$. Put $C=\max_bC_b$ and
$t_0=\min(\exp(-1),\min_bt_b)$. By Lemma~\ref{cv-origin-analytic-system-enlarge}(1), every
$(\varphi_b,\varphi_b(t))$, for $0<t\le t_0$, is a regularized weight with constants $(C,t_0)$.

\emph{The analytic system.} Let $S$ be $\tilde X$, of dimension $D$, with the K\"ahler form $\omega$, the section $s$
in place of $\sigma$ and the divisor $\partial S=\partial$, the numbers $C$ and $t_0$, and for $b\in\Sigma$ the
pullback $\bar L_b$ of $\bar{\mathcal O}(d_b)$ with $\varphi_b$ and $\varphi_b(t)$. Then $\varphi_A=\varphi_A(t)=0$,
and the conditions of Definition~\ref{cv-proportionality-analytic-system}(1) are those of the regularized weights.
So $S$ is an analytic system of type $(\Psi,N)$.

\emph{The analytic data.} Take the given $\bar\sigma_{\tau,i}$ for (1) of Definition~\ref{cv-origin-analytic}, $S$
for (2), and $\tilde X$ with the maps $\lambda_x$ for (3). For (4), let $\pi$ be the identity of $\tilde X$; it is
bijective and holomorphic, of degree $1$ off the empty set. For (5), put $\bar y_b=\bar{\mathcal O}(d_b)$, which is
smooth; its pullback is the metric of $\bar L_b$, so its first Chern form on $\tilde X$ is that of $\bar L_b$. The
bundles $\bar y_b(t)$ are smooth, invariant under conjugation and of the multidegree of $\bar y_b$, and
$\log h_{y_b(t)}(\lambda_x(x))=\log h_{d_b}(\lambda_x(x))+\varphi_b(t)(\pi(x))$ for $0<t\le t_0$. The conditions
hold:
\begin{enumerate}
\item[(i)] by Definition~\ref{cv-origin-reference}(4), a complex point of $Z$ with coordinates $\lambda_x(x)$ is one
  over which $\mathcal A_Z$ is not abelian if and only if $\lambda_x(x)$ lies in the zero set of $J_\partial$, that is,
  $x\in\partial S$;
\item[(ii)] $\bar y_{\tau,i}$ has multidegree $-\sum_{k\le i}a_{\tau,k}$, and for $x\notin\partial S$ and $b=(\tau,i)$,
  $\log h_{d_b}(\lambda_x(x))+\varphi_b(x)=\log h_{d_b}(\lambda_x(x))+u_b(x)=-\sum_{k\le i}\log h_{\tau,k}(\lambda_x(x))$;
\item[(iii)] $\bar y_A=\bar{\mathcal O}(e_{k_0})$ has multidegree $1$ at $k_0$ and $0$ at the other indices, and
  $h_{e_{k_0}}(z)=\big(\sum_{i\in I_{k_0}}|z_{k_0,i}|^2\big)^{1/2}$;
\item[(iv)] $\bar L_A$ is the pullback of $\bar{\mathcal O}(e_{k_0})$, which is positive, and
  $\varphi_A=\varphi_A(t)=0$;
\item[(v)] for $\tau\in\Psi$ and $0<t\le t_0$, the local potentials of the metric of $\bar L_{\tau,N}(t)$ are
  pluriharmonic, by the choice of $\varphi_{\tau,N}(t)$.
\end{enumerate}
Finally, the identity $\psi:S\to\tilde X$ is holomorphic, and $\psi\circ\pi$ is the identity.
\end{proof}
% Lean: AnalyticReference.An.of_local, under PB.Inputs (for RegGlobalStatement in exists_member): one_le_dimFl;
% isArithVar_of_snc_model (Lemma cv-origin-analytic-system-arith); boundary_points_closed (Lemma
% cv-origin-analytic-system-boundary); lowerBound_height_on_boundary (Lemma cv-origin-analytic-system-snc) for R, L,
% s, nb, m and W; hBd and hcb are (a) and (b); hemb (Cmp.IsEmbedded) from R with J = top; hpos from isPosPot_block
% and holomorphic_atlas_coord_compat; d r, wt r and u r with weight_orbit; hsome from exists_member with huc
% (weight_orbit, Cj.fs_conj, hconj), hT (exists_nchart, NChart.snc, hloc, tame_smooth, Tame.add) and hP
% (NChart.zero_of_not_mem, hloc); common_members with member_zero for none (Pi.single k0 1, hpos, through
% enlarging_constants_smooth_members) and IsMember.mono; fs_pos_potential for Theta; S and Re with J = top, pi = id,
% deg = 1, br empty (FA.degree_id_self), ref o = fsHerm (d o) and ref_pot trivial; the fields of RefGeo.An from
% G.bd_iff and hbd (interior), hdeg (fund_deg), hM (fund_sing, flat), fs_single_h (none_h) and hpos (pos).
% No open step remains below this lemma: Lemma cv-origin-analytic-system-weights-regularize-global
% (RegGlobalStatement) is proved.

\begin{lemma}\label{cv-origin-analytic-exists}
In the setting of Definition~\ref{cv-origin-reference}, let the moduli schemes at all levels be given, let
$\mathcal X$ be a global model such that $F_c$ is algebraically closed in the function field of every
irreducible component of $\mathcal X$, and suppose $p\ge1$ and $q\ge1$. Then every reference geometry of
$\mathcal X$ has analytic data (Definition~\ref{cv-origin-analytic}) with a holomorphic map
$\psi:S\to\tilde X$ such that $\psi\circ\pi$ is the identity.
\end{lemma}
\leanref{Development/AutomorphicSpectrum/Proofs/GeometryProofs/ExistsHolomorphicSectionOfAnalytic.lean}{GlobalTheory.T.exists_holomorphic_section_of_analytic}
\leanref{Development/WeilBound/CoreStatements/CentralCharts/ReferenceGeometryHolomorphicInverse.lean}{AnalyticReference.reference_geometry_holomorphic_inverse}
\begin{proof}
Lemma~\ref{cv-origin-analytic-hodge}, applied to the Hodge component and the flag of the reference
geometry, gives Hodge metrics $\bar\sigma_{\tau,i}=(a_{\tau,i},h_{\tau,i})$, invariant under conjugation, on
its root lines. By Definition~\ref{cv-origin-reference}(4), the zero set $B$ of $J_\partial$ satisfies the
hypothesis of Lemma~\ref{cv-origin-analytic-local}, and $p+q\ge1$. So that lemma gives (1) and (2) for
every chart $\gamma$ with boundary $B$, and Lemma~\ref{cv-origin-analytic-system} gives the analytic data and
$\psi$.
\end{proof}
% Lean: reference_geometry_holomorphic_inverse: hodge_cpt_flag_xdeg, An.of_local with PB.Inputs from
% Top.Given, and An.local_bounds with G.bd_iff for each chart chi.

\subsection{The boundary form and analytic resolutions}\label{subsection-cv-boundary-form-resolutions}

\begin{lemma}\label{cv-origin-weight-birational}
In the setting of Definition~\ref{cv-origin-hodge-component}, let $Z$ be a Hodge component of $\mathcal X$
at $S$ with coordinates $y\in\Omega^{\sqcup_{k\in K}I_k}$, with $K$ nonempty, and let $Z$ have a flag. Then
$\Omega$ is the function field of the kernel of $y$ (Definition~\ref{lit-multiprojective}(4)), and this
kernel is an arithmetic subvariety of dimension $|\Psi|N(N-1)/2+1$.
\end{lemma}
\leanref{Development/AutomorphicSpectrum/Proofs/GeometryProofs/HodgeCptBirationalArithVarStatement.lean}{GlobalTheory.T.hodge_cpt_birational_arith_var}
\begin{proof}
Put $D=|\Psi|N(N-1)/2$. The field $\Omega$ has characteristic zero, since $F_c$ embeds in it. Let $\Omega_0\subset\Omega$
be the residue field of the image $\eta_0$ of the generic point of $Z$, the generic point of an irreducible component
of $\mathcal X$.

\emph{The field $\Omega_0$.} Lemma~\ref{cv-origin-change-factor}, applied to the canonical morphism
$\operatorname{Spec}\Omega_0\to\mathcal X$, factors it through $j$; let $w$ be the point of $X$ that it gives, so that
$j(w)=\eta_0$ and $j$ identifies the residue fields at $w$ and $\eta_0$. As $j$ is injective and continuous, $w$ has no
generization other than itself, and the local ring of $X$ at $w$ has dimension $0$. The moduli scheme $X$ is smooth of
relative dimension $\tfrac12[E:\mathbf Q]pq=|\Psi|pq$ over $F_c$, so its dimension at $w$ is $|\Psi|pq$
\cite[Tags 02G1 and 02G2]{Sta}, and this is the transcendence degree of the residue field at $w$ over $F_c$
\cite[Tag 00P1]{Sta}. As $F_c$ is algebraic over $\mathbf Q$, $\Omega_0$ has transcendence degree $|\Psi|pq$ over
$\mathbf Q$.

\emph{The field $\Omega$.} By Definition~\ref{cv-origin-hodge-component}(2), $\Omega$ is finite over $\Omega_0(u)$ for a
family $u=(u_{s,lk})_{l>k}$ algebraically independent over $\Omega_0$, with $r_s(r_s-1)/2$ elements for the label $s$. The
labels are $\tau$ and $\bar\tau$ for $\tau\in\Psi$, with $r_\tau=p$ and $r_{\bar\tau}=q$, so $\Omega$ has transcendence
degree $|\Psi|\big(p(p-1)+q(q-1)\big)/2$ over $\Omega_0$, and
$|\Psi|\big(pq+p(p-1)/2+q(q-1)/2\big)=D$ over $\mathbf Q$.

\emph{The function field.} The generic point of $Z$ lies in a chart $\kappa$, and every $a\in\Omega$ is regular there.
By Lemma~\ref{cv-origin-snc-charts}(1), $a=r/s$ with $r,s\in\mathbf Z[y_{k,i}/y_{k,\kappa_k}]$ and $s\neq0$. Write $r$ and $s$ as
polynomials in the ratios and multiply both by $\prod_ky_{k,\kappa_k}^{d_k}$ for large $d_k$: then $a=g(y)/h(y)$ with $g$
and $h$ multihomogeneous of the same multidegree and $h(y)\neq0$. So $\Omega$ is the function field of the kernel $P$ of
$y$.

\emph{The dimension.} The ideal $P$ is the kernel of the homomorphism $R\to\Omega[T_k:k\in K]$,
$x_{k,i}\mapsto y_{k,i}T_k$, to a domain, compatible with the gradings by multidegree; so it is a multihomogeneous prime.
It is relevant, as $y_{k,\kappa_k}\neq0$, and $P\cap\mathbf Z=0$, as $\Omega$ has characteristic zero. The ring $R/P$ is the
subring of $\Omega[T_k]$ generated by the $y_{k,i}T_k$. Its fraction field contains the ratios $y_{k,i}/y_{k,\kappa_k}$,
which generate $\Omega$ by the previous paragraph, so it is $\Omega(y_{k,\kappa_k}T_k:k\in K)$, of transcendence degree
$D+|K|$ over $\mathbf Q$. Let $\mathfrak q$ be a maximal ideal of $R/P$. As $\mathbf Z$ is a Jacobson ring, $\mathfrak q$
lies over a prime $p\mathbf Z$ and its residue field is finite \cite[Tags 00G4 and 00GB]{Sta}. The ring $\mathbf Z$ is
regular, hence Cohen--Macaulay and universally catenary \cite[Tag 00NM]{Sta}, so the dimension formula
\cite[Tag 02JU]{Sta} gives $\dim(R/P)_{\mathfrak q}=1+(D+|K|)-0$. Hence $\dim R/P=(D+1)+|K|$, and $P$ is an arithmetic
subvariety of dimension $D+1$.
\end{proof}
% Lean: flagged_model_dim: hodge_cpt_interior_generic (= HodgeCpt.basic: CharZero and the lift x, through mem_range_jQ),
% generic_trdeg_eq_dimension (= CptTrdeg.trdeg_eq_dim: Xo.smooth base-changed, eq_of_specializes_of_fst,
% residueField_fg_trdeg_of_smoothOfRelativeDimension, trdeg_residueField_of_fst), generic_flag_trdeg_formula
% (= FlagTrdeg.trdeg_of_isGenericFlag), trdeg_add_eq and the count harith, and birational_of_transcendence_degree
% (= coordTie_birational_isArithVar: regAt_iff_chart_localization at the generic point, birational_of_ratioField_eq_top
% with rep_of_mem, isArithVar_of_ratioField_eq_top with Cmp.ringKrullDim_coord = IntDomainDim.full (dim = trdeg + 1 for a
% finitely generated domain over Z of characteristic zero) and ConeTrdeg.full). The book gives the case O = O_c of its
% statement and takes the lift x from Lemma cv-origin-change-factor; the Lean argues for any injective O -> F_c.

\begin{lemma}\label{cv-origin-boundary-points}
Let $X$, $\Omega$, $\xi$ and $y\in\Omega^{\sqcup_{k\in K}I_k}$ be as in Lemma~\ref{cv-origin-snc-charts}, let $\kappa$ be a
chart of $y$, with $X_\kappa$ and $A^{\mathbf Z}_\kappa$ as there, and let $L$ be a field. For a morphism
$\zeta:\operatorname{Spec}L\to X$ with image $P$ and an element $a\in\Omega$ regular at $P$, the \emph{value} of $a$ at
$\zeta$ is the image of its germ under the homomorphism $\mathcal O_{X,P}\to L$ given by $\zeta$.
\begin{enumerate}
\item For every ring homomorphism $\psi:A^{\mathbf Z}_\kappa\to L$ there is a morphism $\zeta:\operatorname{Spec}L\to X$
  whose image $P$ lies in $X_\kappa$, such that every $a\in A^{\mathbf Z}_\kappa$ has value $\psi(a)$ at $\zeta$, and the
  elements of $A^{\mathbf Z}_\kappa$ that are units at $P$ are those $a$ with $\psi(a)\neq0$.
\item Let $w\in L^{\sqcup_kI_k}$ with $w_{k,\kappa_k}\neq0$ for every $k$, at which every element of the kernel of $y$
  vanishes. Then there is a ring homomorphism $\psi:A^{\mathbf Z}_\kappa\to L$ with
  $\psi(y_{k,i}/y_{k,\kappa_k})=w_{k,i}/w_{k,\kappa_k}$, and the morphism $\zeta$ of (1) has \emph{coordinates} $w$: for
  every $k$ and every $c\in I_k$ with $w_{k,c}\neq0$, $y_{k,c}\neq0$ and every $y_{k,i}/y_{k,c}$ is regular at $P$ with
  value $w_{k,i}/w_{k,c}$ at $\zeta$.
\end{enumerate}
\end{lemma}
\leanref{Development/WeilBound/CoreStatements/VolumeEquality/ChartMeetsPoint.lean}{SheetDecomposition.ChartForm.exists_point_hom}
\leanref{Development/AutomorphicSpectrum/Proofs/GeometryProofs/LocalRingPointsExistsUniqueStatement.lean}{GlobalTheory.T.local_ring_points_existsUnique}
\begin{proof}
(1) The kernel $\mathfrak p$ of $\psi$ is a prime of $A^{\mathbf Z}_\kappa$. By Lemma~\ref{cv-origin-snc-charts}(2) there is a
point $P$ of $X_\kappa$ at which the elements of $A^{\mathbf Z}_\kappa$ that are units are those outside $\mathfrak p$, and by
Lemma~\ref{cv-origin-snc-charts}(1) the local ring of $X$ at $P$ is the localization of $A^{\mathbf Z}_\kappa$ at
$\mathfrak p$. So $\psi$ extends to a local homomorphism $\mathcal O_{X,P}\to L$, which is a morphism
$\zeta:\operatorname{Spec}L\to X$ with image $P$ \cite[Tag 01J6]{Sta}, and every $a\in A^{\mathbf Z}_\kappa$ has value
$\psi(a)$ at $\zeta$.

(2) Put $t_{k,i}=y_{k,i}/y_{k,\kappa_k}$; these generate $A^{\mathbf Z}_\kappa$. Let $f$ be a polynomial with integer
coefficients in variables $x_{k,i}$ with $f(t)=0$, and let $d$ be at least its degree. Then
$F=\prod_kx_{k,\kappa_k}^d\cdot f(x_{k,i}/x_{k,\kappa_k})$ is a polynomial, and
$F(y_{k,i}T_k)=\prod_k(y_{k,\kappa_k}T_k)^df(t)=0$, so $F$ lies in the kernel of $y$. Hence
$0=F(w)=\prod_kw_{k,\kappa_k}^d\,f(w_{k,i}/w_{k,\kappa_k})$, and $f(w_{k,i}/w_{k,\kappa_k})=0$. So $\psi$ is well defined.
If $w_{k,c}\neq0$, then $\psi(t_{k,c})\neq0$, so $t_{k,c}$ is a unit at $P$ by (1); hence $y_{k,c}\neq0$, and
$y_{k,i}/y_{k,c}=t_{k,i}t_{k,c}^{-1}$ is regular at $P$ with value
$(w_{k,i}/w_{k,\kappa_k})(w_{k,c}/w_{k,\kappa_k})^{-1}=w_{k,i}/w_{k,c}$.
\end{proof}
% Lean: TiePoints.existsUnique_point (exists_homog and aeval_ratio_eq_zero for psi, prime_chart_existsUnique_point,
% exists_local_ext, existsUnique_centre_local_hom, valAt_inv and valAt_mul for the coordinates); exists_point_hom applies
% it to the single block of the Segre products and extends the values from the generators to A_kappa.

\begin{lemma}\label{cv-origin-weight-boundary}
In the setting of Definition~\ref{cv-origin-reference}, let a reference geometry of $\mathcal X$ be given,
with its notation. Then there are $d_\beta\in\mathbf Z_{\ge0}^K$ and a multihomogeneous $\beta\in R$ of
multidegree $d_\beta$, not in the kernel of $y$ and with $\beta(y)\neq0$, such that
$\beta(y)\prod_ky_{k,\kappa_k}^{-d_{\beta,k}}\in\mathbf Z[y_{k,i}/y_{k,\kappa_k}]$ for every chart $\kappa$ of
$y$, and such that at every point $P$ of $Z$ over no prime over which $\mathcal A_Z$ is not abelian, and
every chart $\kappa$ whose open subscheme $Z_\kappa$ contains $P$,
$\beta(y)\prod_ky_{k,\kappa_k}^{-d_{\beta,k}}$ is not a unit at $P$.
\end{lemma}
\leanref{Development/AutomorphicSpectrum/Proofs/GeometryProofs/ReferenceGeometryMultihomogFormStatement.lean}{GlobalTheory.T.reference_geometry_multihomog_form}
\begin{proof}
For a chart $\kappa$ of $y$ let $t^\kappa$ be the point with coordinates $y_{k,i}/y_{k,\kappa_k}$, and let $A_\kappa$ and
$A^{\mathbf Z}_\kappa$ be the algebras over $\mathbf Q$ and $\mathbf Z$ generated by its coordinates. Let $\kappa^0$ be a chart
whose open subscheme contains the generic point of $Z$.

\emph{A form not vanishing at $t^{\kappa^0}$.} Some $h\in J_\partial$ has $h(t^{\kappa^0})\neq0$. Otherwise, let $f_l$ and
$a_l$ be as in Definition~\ref{cv-origin-reference}(3) at $\kappa^0$. As $\sum_la_lf_l=1$, some $f_l\neq0$. The zero ideal
of the domain $S=A_{\kappa^0}[1/f_l]$ is prime and contains the values of $J_\partial$, so by (3) it contains some $z_i$.
But the differential of $z_i$ is a member of a basis of $\Omega_{S/\mathbf Q}$, so $z_i\neq0$.

Write $h=\sum_dh_d$ with $h_d$ multihomogeneous of multidegree $d\in\mathbf Z_{\ge0}^K$, and let $\beta=h_{d_\beta}$ be a
component with $\beta(t^{\kappa^0})\neq0$. For every chart $\kappa$, homogeneity gives
$\beta(y)\prod_ky_{k,\kappa_k}^{-d_{\beta,k}}=\beta(t^\kappa)\in A^{\mathbf Z}_\kappa$. At $\kappa^0$ this is nonzero, so
$\beta(y)\neq0$, and $\beta$ is not in the kernel of $y$, whose elements vanish at $y$.

\emph{The boundary.} Let $P$ be a point of $Z$ over no prime over which $\mathcal A_Z$ is not abelian, and $\kappa$ a chart
with $P\in Z_\kappa$, and suppose that $\beta(t^\kappa)$ is a unit at $P$. Let $\mathfrak p$ be the prime of the elements of
$A^{\mathbf Z}_\kappa$ that are not units at $P$. The ring $A^{\mathbf Z}_\kappa$ is countable and no prime number lies in
$\mathfrak p$, so the fraction field of $A^{\mathbf Z}_\kappa/\mathfrak p$ is a countable field of characteristic zero and
embeds in $\mathbf C$. This gives $\psi:A^{\mathbf Z}_\kappa\to\mathbf C$ with kernel $\mathfrak p$. By
Lemma~\ref{cv-origin-boundary-points}(1) there is a complex point $\zeta$ of $Z$ whose image lies in $Z_\kappa$ and has
the units of $A^{\mathbf Z}_\kappa$ outside $\mathfrak p$, so its image is $P$ by Lemma~\ref{cv-origin-snc-charts}(2). Put
$w_{k,i}=\psi(t^\kappa_{k,i})$, so that $w_{k,\kappa_k}=1$. As in the proof of
Lemma~\ref{cv-origin-boundary-points}(2), $w$ are coordinates of $\zeta$ in the sense of
Definition~\ref{cv-origin-reference}(4), and so are the $(u_kw_{k,i})$ for $u\in(\mathbf C^\times)^K$. As $\mathcal A_Z$
is not abelian at $P$, condition (4) of Definition~\ref{cv-origin-reference} says that every element of $J_\partial$
vanishes at these points: $\sum_dh_d(w)\prod_ku_k^{d_k}=0$ for all $u\in(\mathbf C^\times)^K$. So every $h_d(w)$ is $0$.
But $\beta(w)=\psi(\beta(t^\kappa))\neq0$, since $\beta(t^\kappa)$ is a unit at $P$ and so lies outside $\mathfrak p$.
\end{proof}
% Lean: ExRefEff.BdForm.bd_form: exists_chart_form (G.snc at the chart of the generic point, the zero prime, hli.ne_zero),
% bcomp and exists_aeval_bcomp_ne for the component, dehom_eq and mem_chartLat, and exists_point_at (exists_hom_ker with
% nonempty_hom_complex for psi, exists_point_hom, centre_eq_of_ker, coordsAt_of_val), coordsAt_scale, G.bd_iff and
% aeval_bcomp_eq_zero for the last paragraph.

\begin{lemma}\label{cv-origin-weight-nef-degree}
Let $\Psi$ be a finite nonempty set and $N\ge2$, put $D=|\Psi|N(N-1)/2$, and let $S$ be an analytic system of
type $(\Psi,N)$ (Definition~\ref{cv-proportionality-analytic-system}). Let $K$ and $(I_k)_{k\in K}$ be as in
Definition~\ref{lit-multiprojective}, let $\Omega$ be a field, let $y\in\Omega^{\sqcup_kI_k}$ be a point whose
kernel is an arithmetic subvariety $X$ of dimension $D+1$, and let $\tilde X$, $\lambda_x$, $\pi$ and $\bar y_b$
be as in (3) to (5) of Definition~\ref{cv-origin-analytic} for $y$ and $S$, except that $\pi$ is not assumed
bijective and has some degree $\deg\pi\ge1$ off a closed set of measure zero. Let $n\in\mathbf Z^\Sigma$, let
$a_n=\sum_bn_ba_b$, where $a_b$ is the multidegree of $\bar y_b$, and
let $\bar a_n=(a_n,h_n)$ with $h_n=\prod_bh_{\bar y_b}^{n_b}$, where $h_{\bar y_b}$ is the function of $\bar y_b$.
Then $\bar a_n$ is a smooth metrized line bundle, and the integral over the complex points of $X$ of the $D$-th
power of its first Chern form is $\deg\pi\cdot\deg_S(n)$.
\end{lemma}
\leanref{Development/AutomorphicSpectrum/Proofs/MetricProofs/RealisationPotentialDegreeCompatibilityStatement.lean}{GlobalTheory.T.realisation_potential_degree_compatibility}
\leanref{Development/WeilBound/CoreStatements/PluriharmonicCharts/XRePot.lean}{PositivityTheory.Xgp.t_re_pot}
\begin{proof}
The function $h_n$ is positive and smooth on $\hat C$, as a product of integer powers of positive smooth
functions, and $h_n(cz)=\prod_k|c_k|^{a_{n,k}}h_n(z)$ for $c\in(\mathbf C^\times)^K$, since each $h_{\bar y_b}$
has this property with the exponents $a_b$. So $\bar a_n$ is a smooth metrized line bundle
(Definition~\ref{lit-multiprojective}(2)).

For a smooth metrized line bundle $\bar L=(a,h)$, the functions $\log h\circ\lambda_x$ near the points $x$ are
local potentials on $\tilde X$ of the pullback of the first Chern form of $\bar L$: near a common point, two of
the maps $\lambda_x$ differ by a holomorphic map $c$ to $(\mathbf C^\times)^K$, which changes $\log h$ by
$\sum_ka_k\log|c_k|$, locally the real part of a holomorphic function. By (3) of
Definition~\ref{cv-origin-analytic}, the map from $\tilde X$ to the complex points of $X$ given by the
$\lambda_x$ is bijective and immersive. So the integral over the complex points of $X$ of a product of $D$ first
Chern forms of smooth metrized line bundles is the integral over $\tilde X$ of the Monge--Amp\`ere measure of
these potentials (Definition~\ref{cv-period-monge-ampere}(3)).

The potential $\log h_n\circ\lambda_x$ of $\bar a_n$ is $\sum_bn_b\log h_{\bar y_b}\circ\lambda_x$. The
Monge--Amp\`ere density is multilinear, so the integral of the statement is the sum over the words
$w=(w_1,\dots,w_D)$ in $\Sigma$ of $n_{w_1}\cdots n_{w_D}I_{\tilde X}(w)$, where $I_{\tilde X}(w)$ is the
integral over $\tilde X$ of the Monge--Amp\`ere measure of the potentials of $\bar y_{w_1},\dots,\bar y_{w_D}$.
By (5) of Definition~\ref{cv-origin-analytic}, the first Chern form of $\bar y_b$ on $\tilde X$ is the pullback
by $\pi$ of that of $\bar L_b$, so the potential of $\bar y_b$ is the pullback by $\pi$ of a local potential
$\vartheta_b$ of the metric of $\bar L_b$, up to the real part of a holomorphic function. As $\pi$ has degree
$\deg\pi$ off a closed set of measure zero, the change of variables along $\pi$ gives
$I_{\tilde X}(w)=\deg\pi\cdot I_S(w)$, where $I_S(w)$ is the integral over $S$ of the Monge--Amp\`ere measure of
$(\vartheta_{w_1},\dots,\vartheta_{w_D})$, that is, of $c_1(\bar L_{w_1})\wedge\dots\wedge c_1(\bar L_{w_D})$. By
multilinearity again,
$\sum_wn_{w_1}\cdots n_{w_D}I_S(w)=\int_S\big(\sum_bn_bc_1(\bar L_b)\big)^D=\deg_S(n)$.
\end{proof}
% Lean: clause 4 of t_re_pot (= PB.re_pot), under ResExistsStatement: combLine_smooth_conj_invariant
% (combLine_eq_top) for the smoothness; ptsIntegral_one_eq_mInt (with resBridge) for the integral over the
% manifold R.A of the realization; Real.log_prod and Real.log_zpow (locDens_congr_of_eqOn) for the potential of
% combLine; mInt_sumPot_expand (GlobalTheory.T.monge_ampere_integral_multilinear_expansion) for the expansion;
% re_pot.mass_eq (degreeFormula with Re.degree, and ref_pot) for the factor Re.deg; AnSys.degPoly =
% formPoly D degNum and formPoly_isHomogeneous_and_eval for deg_S. Sizes: ptsIntegral_one_eq_mInt about 1,000
% lines, mass_eq about 2,100, mInt_sumPot_expand about 400.

\begin{lemma}\label{cv-origin-weight-nef-generic-fibre}
Let $K$, $(I_k)_{k\in K}$ and $R$ be as in Definition~\ref{lit-multiprojective}, let $(\cdot\mid\cdot)$ be an
arithmetic intersection system, let $D\ge0$, let $X$ be an arithmetic subvariety of dimension $D+1$ defined by
the prime $P$, let $\bar L_j=(a_j,h_j)$, $j=1,\dots,D$, be smooth metrized line bundles, and let $p$ be a prime
number. Then the integral over the complex points of $X$ of $c_1(\bar L_1)\wedge\dots\wedge c_1(\bar L_D)$ is
\[
 \sum_Vm_V\,(\mathcal O(a_1)\cdots\mathcal O(a_D)\cdot V),
\]
the sum over the relevant minimal primes $V$ of $P+(p)$, where $m_V$ is the multiplicity of $V$ in
$R/(P+(p))$ and $(\mathcal O(a_1)\cdots\mathcal O(a_D)\cdot V)$ is the intersection number on the projective
variety over $\mathbf F_p$ defined by $V$.
\end{lemma}
\leanref{Development/AutomorphicSpectrum/Proofs/DivisibilityProofs/IntersectionSystemDegreeStatement.lean}{GlobalTheory.T.intersection_system_degree}
\leanref{ReciprocityTheory/ReductiveGroups/ArithmeticGeometry/BirationalGeometry/InterSt.lean}{ReciprocityTheory.ChartRestriction.G.Geo.DegreeStatement}
\leanref{GroundTheory/HeightsAndTheta/Heights/Intersections/ConservationDegreeFormOfExistence.lean}{GroundTheory.Heights₂.Heights.Inter.Conservation.statement_holds}
\begin{proof}
Let $\bar A=(0,h)$, where $h$ is the constant function $p$ on $\hat C$. It is a smooth metrized line bundle of
multidegree $0$. Compute $(\bar L_1\cdots\bar L_D\bar A\mid X)$ twice by the recursion of
Definition~\ref{lit-multiprojective}(3), with the sections $s=p$ and $s=1$ of multidegree $0$. Neither lies in
$P$, since $P\cap\mathbf Z=0$.

For $s=p$, every relevant minimal prime $V$ of $P+(p)$ contains $p$, so $V\cap\mathbf Z$ is a prime ideal
containing $p$, and $V\cap\mathbf Z=p\mathbf Z$. So no $V$ is horizontal, and each contributes
$m_V\,(\mathcal O(a_1)\cdots\mathcal O(a_D)\cdot V)\log p$. The function $\log(|p|/h)$ is $0$, so its integral
vanishes. Hence $(\bar L_1\cdots\bar L_D\bar A\mid X)=\log p\sum_Vm_V\,(\mathcal O(a_1)\cdots\mathcal O(a_D)\cdot V)$.

For $s=1$, the ideal $P+(1)$ is $R$, which no prime ideal contains, so the sum in the recursion is empty. The
function $\log(|1|/h)$ is the constant $-\log p$. Hence $(\bar L_1\cdots\bar L_D\bar A\mid X)$ is $\log p$
times the integral over the complex points of $X$ of $c_1(\bar L_1)\wedge\dots\wedge c_1(\bar L_D)$.

As $\log p>0$, comparing the two values gives the claim.
\end{proof}
% Lean: GlobalTheory.T.intersection_system_degree = Top.degreeStatement_of_intersectionSystem =
% Giv.degree_of_isIntersectionSystem: Conservation.conservation_iff_degree_statement,
% conservation_degree_form_of_existence and Conservation.statement_holds. Abar is constLine (-log p), with the
% constant function exp(log p) = p; the recursion is hI.2.2 with s = C p (isMultiHomog_C, C_not_mem;
% not_horizontal, charOf_eq and vertR_C_prime for the sum, logNorm = 0) and with s = C 1 (divSupp empty,
% logNorm = -log p, ptsIntegral_const); mul_left_cancel by log p > 0.

\begin{lemma}\label{cv-origin-weight-nef-generic-conservation}
Let $K$, $(I_k)_{k\in K}$ and $R$ be as in Definition~\ref{lit-multiprojective}, let $D\ge0$, let $X$ be an
arithmetic subvariety of dimension $D+1$ defined by the prime $P$, with generic fibre $X_{\mathbf Q}$, let
$a_1,\dots,a_D\in\mathbf Z^K$, and let $p$ be a prime number. Then
\[
 (\mathcal O(a_1)\cdots\mathcal O(a_D)\cdot X_{\mathbf Q})=\sum_Vm_V\,(\mathcal O(a_1)\cdots\mathcal O(a_D)\cdot V),
\]
the sum over the relevant minimal primes $V$ of $P+(p)$, where $m_V$ is the multiplicity of $V$ in
$R/(P+(p))$ and $(\mathcal O(a_1)\cdots\mathcal O(a_D)\cdot V)$ is the intersection number on the projective
variety over $\mathbf F_p$ defined by $V$.
\end{lemma}
\leanref{GroundTheory/HeightsAndTheta/Heights/OrbitLimits/StatementsDegreeFibre.lean}{GroundTheory.Heights₂.Heights.Targets.FibreDegreeSum.Statement}
\leanref{GroundTheory/HeightsAndTheta/Heights/OrbitLimits/PFibreDegreeSum.lean}{GroundTheory.Heights₂.Heights.Targets.FibreDegreeSum.proof}

This is the conservation of number for the base change $X_{(p)}$ of $X$ to $\mathbf Z_{(p)}$, which is integral,
projective and flat over this discrete valuation ring, of relative dimension $D$, with generic fibre
$X_{\mathbf Q}$. Its special fibre is defined by $P+(p)$, and its fundamental cycle is $\sum_Vm_V[V]$
\cite[\S1.5]{Ful98}. Fulton's specialization map takes $[X_{\mathbf Q}]$ to this cycle \cite[\S20.3]{Ful98}, and
the bundles $\mathcal O(a_j)$ are pseudo-divisors with support $X_{(p)}$ \cite[Definition~2.2.1 and
\S2.5]{Ful98}, so the equality is the case of these pseudo-divisors in \cite[Example~20.3.3]{Ful98}.
% Lean: FibreDegreeSum.proof, its own proof: FibreDegreeSum.intNum_eq_cycEval localizes P at Z_(p) and applies
% NumbersDVR.ConservationOfNumberLE.proof over the DVR (GenericPieceRank and SpecialPieceRank give equal Hilbert
% functions of the generic and special fibres, IntNumAssoc and SpecialCycleEval split the special fibre into its
% components), with IntNumBaseChange.intersectionNumber_baseChange for Q -> Frac Z_(p) and F_p -> the residue
% field and WeightedDegree.cycEval_divCycle_map for the components; Bridge.intNum_map_eq_geoNum with
% FibreCompDimZ.full gives the degrees geoNum of the components. With Targets.ExistsUnique.full, about 4,200 lines.

\begin{definition}\label{cv-origin-intersection-resolution}
Let $K$, $(I_k)_{k\in K}$, $R$ and $\hat C$ be as in Definition~\ref{lit-multiprojective}, let $e\ge0$, let $X$ be an
arithmetic subvariety of dimension $e+1$ defined by the prime $P$, and let $J$ be an ideal of $R$ that is generated by
multihomogeneous elements and is not contained in $P$. Write $Z(J)$ for the set of points of $\hat C$ at which every
element of $J$ vanishes. An \emph{analytic resolution} of $X$ off $J$ consists of a compact complex manifold
$\tilde X$ of dimension $e$ and, for each $x\in\tilde X$, an open neighbourhood $U_x$ of $x$ and a holomorphic map
$\lambda_x:U_x\to\hat X(\mathbf C)$, such that:
\begin{enumerate}
\item at every point of $U_x\cap U_{x'}$ the values of $\lambda_x$ and $\lambda_{x'}$ lie in one orbit of
  $(\mathbf C^\times)^K$;
\item every point of $\hat X(\mathbf C)$ outside $Z(J)$ lies in the orbit of some $\lambda_x(x)$;
\item if $\lambda_x(x)\notin Z(J)$ and $\lambda_{x'}(x')$ lies in the orbit of $\lambda_x(x)$, then $x=x'$;
\item if $\lambda_x(x)\notin Z(J)$, then for every choice of indices $\kappa_k\in I_k$ with
  $\lambda_x(x)_{k,\kappa_k}\neq0$ for all $k$, the affine coordinates $z_{k,i}/z_{k,\kappa_k}$ of $\lambda_x$ have
  injective derivative at $x$ in every chart of $\tilde X$;
\item for every multihomogeneous $f\in R\setminus P$ whose zeros in $\hat X(\mathbf C)$ lie in $Z(J)$, there are
  finitely many subsets $W_1,\dots,W_r$ of $\tilde X$ and integers $m_1,\dots,m_r\ge0$ such that every $x\in\tilde X$
  has a holomorphic chart $\psi$ on a neighbourhood $V\subset U_x$ of $x$ with the following property: each
  $W_l\cap V$ is empty or the zero set of a coordinate $\psi_{i(l)}$, with $i(l)\neq i(l')$ for $l\neq l'$, and
  $f\circ\lambda_x=u\prod_l\psi_{i(l)}^{m_l}$ on $V$, the product over the $l$ with $W_l\cap V$ nonempty, for a
  holomorphic function $u$ on $V$ without zeros;
\item the set of $x\in\tilde X$ with $\lambda_x(x)\in Z(J)$ has measure zero.
\end{enumerate}
\end{definition}
\leanref{ReciprocityTheory/ReductiveGroups/ArithmeticGeometry/BirationalGeometry/InterRes.lean}{ReciprocityTheory.ChartRestriction.G.Geo.AnLogRes}

\begin{lemma}\label{cv-origin-intersection-resolution-exists-arith}
Let $(\kappa,v)$ be a block map (Definition~\ref{cv-origin-intersection-full}(1)), let $\Omega$ be a field of
characteristic zero, and let $y\in\Omega^{\sqcup_{k'}I'_{k'}}$ have a nonzero coordinate in every block. Suppose that the
kernel of $y\circ v=(y_{v(k,i)})_{k,i}$ (Definition~\ref{lit-multiprojective}(4)) is an arithmetic subvariety of dimension
$d$, and that the field generated over $\mathbf Q$ by the ratios $y_{k',i}/y_{k',j}$ is algebraic over the one generated
by the ratios of $y\circ v$. Then the kernel of $y$ is an arithmetic subvariety of dimension $d$.
\end{lemma}
\leanref{Development/AutomorphicSpectrum/Proofs/GeometryProofs/IsArithVarOfBlockAlgebraic.lean}{GlobalTheory.CompanionTheory.Ar.Cmp.isArithVar_of_isAlgebraic_image}
\begin{proof}
The kernel $P$ of $y$ is the kernel of a homomorphism to a domain that respects the gradings, so it is a multihomogeneous
prime; it is relevant, as every block of $y$ has a nonzero coordinate, and $P\cap\mathbf Z=0$, as $\Omega$ has
characteristic zero. Let $T$ and $T_0$ be the transcendence degrees over $\mathbf Q$ of the fields of ratios of $y$ and of
$y\circ v$; they are equal, as the first field is algebraic over the second. As in the proof of
Lemma~\ref{cv-origin-weight-birational}, the dimension of the quotient by the kernel of a point with a nonzero
coordinate in each block is the transcendence degree of its field of ratios plus the number of blocks plus one. So
$d+|K|=T_0+|K|+1$, and $\dim R'/P=T+|K'|+1=d+|K'|$.
\end{proof}
% Lean: isArithVar_of_isAlgebraic_image: ringKrullDim_coord twice and trdeg_eq_of_le_of_isAlgebraic (blockFnField of
% y o v lies in blockFnField y); relevant, horizontal and isMultiHomogIdeal for the kernel of mConeMap y.

\begin{lemma}\label{cv-origin-intersection-resolution-exists-join}
Let $K$, $(I_k)_{k\in K}$ and $R$ be as in Definition~\ref{lit-multiprojective}, let $\Omega$ be a field of
characteristic zero, let $e\ge0$, and let $y\in\Omega^{\sqcup_kI_k}$ be a point whose kernel $P$ is an arithmetic
subvariety $X$ of dimension $e+1$. Let $\beta\in R\setminus P$ be multihomogeneous. For $M\ge1$ and
$y'\in\Omega^M$, write $R_1=R[x_{\ast,1},\dots,x_{\ast,M}]$ for the ring with one more block $\ast$ with
$I_\ast=\{1,\dots,M\}$, $(y,y')$ for the joined point, $P_1$ for its kernel and $\hat C_1$ for the set of points of
Definition~\ref{lit-multiprojective} for $R_1$. Then there are a multihomogeneous $\beta'\in R\setminus P$, an integer
$M\ge1$, a nonzero $y'\in\Omega^M$ and an ideal $J_1$ of $R_1$, generated by multihomogeneous elements and not contained in
$P_1$, such that:
\begin{enumerate}
\item every $z\in\hat X(\mathbf C)$ with $\beta(z)=0$ has $\beta'(z)=0$;
\item the ratios of the coordinates of $(y,y')$ in each block generate the same field as those of $y$;
\item $J_1$ satisfies, for the point $(y,y')$ and with $D=e$, the condition on the boundary ideal in
  Definition~\ref{cv-origin-reference}: in each chart, the ring of ratios is covered by smooth localizations with
  coordinates $z_1,\dots,z_e$ in which $J_1$ cuts out the union of the $\{z_i=0\}$ for $i\le n$;
\item every point of $\hat C_1$ in the zero set of $P_1$ at whose old coordinates $\beta'$ vanishes lies in the zero set
  of $J_1$;
\item for every $z\in\hat X(\mathbf C)$ with $\beta'(z)\neq0$ there are multihomogeneous $G_1,\dots,G_M\in R$ of one
  multidegree with $y'_iG_j(y)=y'_jG_i(y)$ for all $i,j$ and $G_i(z)\neq0$ for some $i$.
\end{enumerate}
\end{lemma}
\leanref{Development/AutomorphicSpectrum/Proofs/MetricProofs/MultihomogeneousLiftToJoinStatement.lean}{GlobalTheory.T.multihomogeneous_lift_to_join}
\begin{proof}
\emph{The construction.} Construct the following as in the proof of Lemma~\ref{cv-origin-coords-snc}, for $y$ and
$\beta$: a chart $c$ of $y$; the Segre products $u_j=\prod_ky_{k,j_k}$ for $j\in\prod_kI_k$, with $u_c\neq0$, and the ring
$A=\mathbf Q[u_j/u_c]=\mathbf Q[y_{k,i}/y_{k,c_k}]$; for a form $H$ in variables $T_j$, the element $\tilde H\in R$ obtained
by evaluating $H$ at the monomials $\prod_kx_{k,j_k}$; a form $G$ with integer coefficients, of degree $d+1$ and divisible
by $T_c$, with $g_0=G(u/u_c)\neq0$; the scheme theoretic image $C\subset\mathbf P^N_{\mathbf Q}$ of the point $p$ given by
$u$, the affine open $W_0=C\cap D_+(T_c)$, whose ring maps onto $A$, and the function $\sigma_0=G/T_c^{d+1}$ on $W_0$, with
value $g_0$, whose standard open is $U=C\cap D_+(G)$; the morphism $f:X'\to C$ of Lemma~\ref{cv-resolution-coordinates},
an isomorphism over $U$; and the lift $\zeta$ of $p$, whose second projection is given by a nonzero $w\in\Omega^{r+1}$. Put
$\beta'=\tilde G$, $M=r+1$, $y'=w$, and let $J_1\subset R_1$ be the ideal generated by $\tilde G$.

The element $\beta'$ is multihomogeneous, of degree $d+1$ in each block, and $\beta'(y)=G(u)=u_c^{d+1}g_0\neq0$; an element of
$R$ of one multidegree lies in $P$ if and only if it vanishes at $y$, so $\beta'\notin P$. As shown there, $J_1$ is
generated by a multihomogeneous element and is not contained in $P_1$. (1) is the third step of that proof: if
$z\in\hat X(\mathbf C)$ and $\beta(z)=0$, then $G(u(z))=0$, where $u(z)_j=\prod_kz_{k,j_k}$, that is, $\beta'(z)=0$. (2) and
(3) are (1) and (2) of Lemma~\ref{cv-origin-coords-snc}. (4) holds as $J_1$ is generated by $\beta'$.

\emph{Local denominators for $w$.} Let $\psi:A\to\mathbf C$ be a ring homomorphism with $\psi(g_0)\neq0$. Then there are an
index $l$ with $w_l\neq0$ and $b\in A$ with $\psi(b)\neq0$ and $(w_i/w_l)\,b\in A$ for all $i$. Indeed, let $B$ be the ring of
$W_0$. The kernel of $\psi$ composed with $B\to A$ is a prime of $B$, the prime of a point $t\in W_0$, and $t$ lies in the
standard open of $a\in B$ if and only if $\psi$ does not kill the value of $a$. As $\psi(g_0)\neq0$, $t\in U$. As $f$ is an
isomorphism over $U$, $t=f(t')$ with $t'\in f^{-1}(U)$; choose $l$ with $T'_l\neq0$ at $t'$. Let $O$ be the image under $f$
of the set of points of $f^{-1}(U)$ where $T'_l\neq0$, an open subset of $C$ containing $t$, on which the functions
$T'_i/T'_l$, transported through $f$, are regular. Choose a standard open $D(b_0)$ of $W_0$ with $t\in D(b_0)\subset O$. The
image of $p$ lies in $D(b_0)$, as it is dense in $C$; so $\zeta$ lies where $T'_l\neq0$, $w_l\neq0$, and the transported
functions take the values $w_i/w_l$ at $p$. On $D(b_0)$ they are $a_i/b_0^n$ with $a_i\in B$ and one $n$. Taking values,
$(w_i/w_l)\,\bar b_0^{\,n}=\bar a_i$, where the bar is the value in $A$; so $b=\bar b_0^{\,n+1}$ has $(w_i/w_l)\,b=\bar a_i\bar
b_0\in A$, and $\psi(b)\neq0$ as $t\in D(b_0)$.

\emph{Homogenizing.} Every polynomial $Q$ over $\mathbf Q$ in the $T_j$ has a nonzero integer $m$, an $n\ge0$ and a form
$H$ with integer coefficients of degree $n$ such that $H(v)=m\,Q(v)$ whenever $v_c=1$, in any field of characteristic zero:
for $T_j$ take $H=T_j$, for a rational number its numerator and denominator, and for sums and products multiply by powers of
$T_c$ to make the degrees agree. The same works for finitely many $Q_i$ at once, with one $m$ and one $n$.

(5) Let $z\in\hat X(\mathbf C)$ with $\beta'(z)\neq0$. As $T_c$ divides $G$, $u(z)_c\neq0$. If $Q$ is a polynomial over
$\mathbf Q$ with $Q(u/u_c)=0$, then $Q(u(z)/u(z)_c)=0$: with $m$, $n$ and $H$ as above for $Q$, the multihomogeneous element
$\tilde H$ vanishes at $y$, as $H(u)=u_c^nm\,Q(u/u_c)$, so it lies in $P$ and vanishes at $z$, and
$H(u(z))=u(z)_c^nm\,Q(u(z)/u(z)_c)$. So $\psi:A\to\mathbf C$, $Q(u/u_c)\mapsto Q(u(z)/u(z)_c)$, is a well defined ring
homomorphism, and $\psi(g_0)=G(u(z))/u(z)_c^{d+1}=\beta'(z)/u(z)_c^{d+1}\neq0$. The previous paragraph gives $l$ and $b$; write
$(w_i/w_l)\,b=Q_i(u/u_c)$ for polynomials $Q_i$ over $\mathbf Q$, and take $m$, $n$ and forms $H_i$ as above for the $Q_i$.
Put $G_i=\tilde H_i\in R$, multihomogeneous of degree $n$ in each block. Then $G_i(y)=H_i(u)=u_c^nm\,(w_i/w_l)\,b$, so
$y'_iG_j(y)=y'_jG_i(y)$ for all $i,j$. And $G_l(z)=H_l(u(z))=u(z)_c^nm\,\psi(Q_l(u/u_c))=u(z)_c^nm\,\psi(b)\neq0$, as
$(w_l/w_l)\,b=b$.
\end{proof}
\leanref{Development/AutomorphicSpectrum/Proofs/MetricProofs/MultihomogeneousLiftToJoin.lean}{GlobalTheory.CompanionTheory.Ar.Giv.exists_snc_model_iso}
% Lean: GlobalTheory.T.multihomogeneous_lift_to_join = Giv.exists_snc_model_iso, under ResKStatement (Lemma
% cv-resolution-coordinates). It repeats the steps of SncRes.exists_snc_model (Lemma cv-origin-coords-snc): exists_smooth,
% exists_form, and snc_of_res_local, a copy of snc_of_res with the extra clause local_of_hom; beta' = segForm E Gz, y' = x',
% J = span of the renamed segForm. The local denominators are local_of_hom (exists_point_of_hom for t, exists_local_quot
% for l, b and the a_i, with exists_basicOpen_le, pullRatio and IsLocalization.exist_integer_multiples_of_finite); (5) is
% exists_forms (ker_ratio, exists_adjoin_hom for psi, exists_homog_poly and exists_homog_family for the homogenizing, and
% aeval_segForm); (1) to (4) come from exists_form, blockFnField_joinPt with adjoin_seg_le, isSncModel_joinPt and
% mem_mcPts_span. Size: MultihomogeneousLiftToJoin.lean, 403 lines, besides the files shared with Lemma
% cv-origin-coords-snc (AdjoinCoordsArithVar, RegularValuesInAdjoinedField, StandardSmoothChartOfNonzero and
% HomogeneousFormNonvanishingPoint, 903 lines, and their imports).

\begin{lemma}\label{cv-origin-intersection-resolution-exists-embedded-regular}
Let $K$, $(I_k)_{k\in K}$ and $R$ be as in Definition~\ref{lit-multiprojective}, let $\Omega$ be a field of
characteristic zero, let $e\ge0$, let $y\in\Omega^{\sqcup_kI_k}$ be a point whose kernel $P$ is an arithmetic
subvariety of dimension $e+1$, and let $J$ be an ideal of $R$ that satisfies for $y$ and with $D=e$ the condition on the
boundary ideal in Definition~\ref{cv-origin-reference}(3). Let $\kappa$ be a chart of $y$, and let $I_\kappa$ be the ideal
of $\mathbf C[x_{k,i}]$ generated by $P$ and the $x_{k,\kappa_k}-1$ for $k\in K$. Then at every common zero $w$ of
$I_\kappa$, the local ring $\mathcal O_w$ of $\mathbf C[x_{k,i}]/I_\kappa$ at $w$ is regular of dimension $e$.
\end{lemma}
\leanref{Development/AutomorphicSpectrum/Proofs/MetricProofs/SncResolutionRegularLocalRingStatement.lean}{GlobalTheory.T.snc_resolution_regular_local_ring}
\begin{proof}
The first paragraph of the proof of Lemma~\ref{cv-origin-model-resolution-regular} uses only that $P$ is an
arithmetic subvariety of dimension $e+1$, and it shows that $\mathcal O_w$ has dimension $e$. As in its second
paragraph, evaluation at $w$ defines a homomorphism $\psi:A_\kappa\to\mathbf C$,
$g(y_{k,i}/y_{k,\kappa_k})\mapsto g(w)$, and each vector $v$ with $\partial_vh(w)=0$ for all $h\in I_\kappa$ defines a
$\mathbf Q$-derivation $A_\kappa\to\mathbf C$ along $\psi$ that determines $v$; these vectors form a space whose dimension
is that of the cotangent space of $\mathcal O_w$. The condition gives $f_1,\dots,f_m$ and $a_1,\dots,a_m$ in $A_\kappa$
with $\sum_la_lf_l=1$. As $\sum_l\psi(a_l)\psi(f_l)=1$, some $\psi(f_l)$ is nonzero; then $f_l\neq0$, and $\psi$ and these
derivations extend to $B=A_\kappa[1/f_l]$. The condition gives $z_1,\dots,z_e\in B$ whose differentials span the module of
differentials of $B$ over $\mathbf Q$, so a derivation of $B$ along $\psi$ is determined by its values at the $z_i$. So the
space of vectors $v$ has dimension at most $e$, the cotangent space of $\mathcal O_w$ has dimension at most
$e=\dim\mathcal O_w$, and $\mathcal O_w$ is regular.
\end{proof}
\leanref{Development/AutomorphicSpectrum/Proofs/MetricProofs/SncResolutionRegularLocalRing.lean}{GlobalTheory.CompanionTheory.Ar.Cmp.SncRes.Reg.isRegularLocalRing_chart}
% Lean: isRegularLocalRing_chart: ringKrullDim_locAt_chart with pcEquidim for the dimension, as for
% SmoothMan.isRegularLocalRing_chart; finrank_cotangentSpace_locAt and card_le_finrank_jacSpanAt, through
% card_le_of_kaehler: exists_unique_lift (psi on Q[u, 1/F(u)] with F(w) nonzero, F the f_l), exists_span_of_derivations
% (the derivations lie in the span of e vectors given by the dz_i) and card_le_finrank_add. Size:
% SncResolutionRegularLocalRing.lean, 429 lines.

\begin{lemma}\label{cv-origin-intersection-resolution-exists-embedded-snc}
Let $K$, $(I_k)_{k\in K}$, $R$, $\Omega$, $e$, $y$, $P$ and $J$ be as in
Lemma~\ref{cv-origin-intersection-resolution-exists-embedded-regular}, with $J$ generated by multihomogeneous elements and
not contained in $P$, and let $X$ be the arithmetic subvariety defined by $P$. Let $\kappa$ be a chart of $y$, let
$I_\kappa$ be as there, let $I^J_\kappa$ be the radical of the ideal of $\mathbf C[x_{k,i}]$ generated by $I_\kappa$ and
$J$, and let $Q_1,\dots,Q_c$ be the minimal primes of the ideal of $\mathbf C[x_{k,i}]$ generated by $P$ and $J$. Let $w$ be
a common zero of $I_\kappa$, and $\mathcal O_w$ the local ring of $\mathbf C[x_{k,i}]/I_\kappa$ at $w$. Then there are
polynomials $r_1,\dots,r_e$ whose images generate the maximal ideal of $\mathcal O_w$, a set $L\subset\{1,\dots,e\}$, an
injective map $\tau:L\to\{1,\dots,c\}$ and a polynomial $h$ with $h(w)\neq0$, such that:
\begin{enumerate}
\item $h\prod_{l\in L}r_l\in I^J_\kappa$, and $hg$ lies in the ideal generated by $I_\kappa$ and $\prod_{l\in L}r_l$ for
  every $g\in I^J_\kappa$;
\item on a neighbourhood of $w$, a common zero of $I_\kappa$ is a zero of $Q_b$ if and only if one of the $r_l$ with
  $\tau(l)=b$ vanishes there;
\item for every multihomogeneous $f\in R\setminus P$ whose zeros in $\hat X(\mathbf C)$ lie in the zero set of $J$, the
  image of $f$ in $\mathcal O_w$ is a unit times $\prod_{l\in L}r_l^{m_{\tau(l)}}$, where $m_b$ is the largest $m$ such
  that $sf\in Q_b^m+P\,\mathbf C[x_{k,i}]$ for some $s\notin Q_b$.
\end{enumerate}
\end{lemma}
\leanref{Development/AutomorphicSpectrum/Proofs/MetricProofs/SncResolutionBoundaryNormalCrossingStatement.lean}{GlobalTheory.T.snc_resolution_boundary_normal_crossing}
\leanref{ReciprocityTheory/AnalyticForms/RegularLocalRings/IsOpenAlt.lean}{ReciprocityTheory.AnForm.R.SncData}
\begin{proof}
Let $P_{\mathbf C}$ and $J_{\mathbf C}$ be the ideals of $\mathbf C[x_{k,i}]$ generated by $P$ and by $J$. Both are
generated by multihomogeneous elements, and $Q_1,\dots,Q_c$ are the minimal primes over $P_{\mathbf C}+J_{\mathbf C}$. In
the notation of Lemma~\ref{cv-origin-analytic-system-crossings-dehomogenize}, $I_\kappa=(P_{\mathbf C})_\kappa$; let
$J_\kappa=(P_{\mathbf C}+J_{\mathbf C})_\kappa$, the ideal generated by $I_\kappa$ and $J$, so that $I^J_\kappa$ is its
radical. As $I_\kappa$ contains the $x_{k,\kappa_k}-1$, a common zero $w'$ of $I_\kappa$ has $w'_{k,\kappa_k}=1$ for every
$k$; so $w'$ is a point of $\hat C$ with $a_\kappa(w')=w'$, in the notation of
Lemma~\ref{cv-origin-analytic-system-crossings-components}, and it is a common zero of $P_{\mathbf C}$ by the first
paragraph of the proof of that lemma. This applies to $w$.

By Lemma~\ref{cv-origin-intersection-resolution-exists-embedded-regular}, $\mathcal O_w$ is regular of dimension $e$.
Lemma~\ref{cv-origin-analytic-system-crossings-params} gives polynomials $r_1,\dots,r_e$ that vanish at $w$, whose images
generate the maximal ideal of $\mathcal O_w$, and a set $L\subset\{1,\dots,e\}$ (its $T$) such that a prime
$\mathfrak q$ with $I_\kappa\subset\mathfrak q\subset\mathfrak m_w$ contains $J_\kappa$ if and only if it contains $r_l$
for some $l\in L$. So $\mathfrak a=I_\kappa\subset\mathfrak b=J_\kappa$, $w$, the $r_l$ and $L$ satisfy the hypotheses
of Lemma~\ref{cv-origin-analytic-system-crossings-local}, and with $\mathfrak p=P_{\mathbf C}$, $\mathfrak d=J_{\mathbf
C}$ and $z=w$ they satisfy those of Lemma~\ref{cv-origin-analytic-system-crossings-components}. Let $h$ be the polynomial
of Lemma~\ref{cv-origin-analytic-system-crossings-local}(2), and $\tau:L\to\{1,\dots,c\}$ the injective map $\nu$ of
Lemma~\ref{cv-origin-analytic-system-crossings-components}(1).

(1) This is Lemma~\ref{cv-origin-analytic-system-crossings-local}(2), as the radical of $\mathfrak b$ is $I^J_\kappa$.

(2) Let $F$ be the polynomial of Lemma~\ref{cv-origin-analytic-system-crossings-components}(2). The set of points at
which $F$ does not vanish is a neighbourhood of $w$. Let $w'$ be a common zero of $I_\kappa$ in it. By the first
paragraph, $w'$ is a common zero of $P_{\mathbf C}$ in $\hat C$ with $w'_{k,\kappa_k}=1$ and $a_\kappa(w')=w'$. So by
Lemma~\ref{cv-origin-analytic-system-crossings-components}(2), the elements of $Q_b$ vanish at $w'$ if and only if
$b=\tau(l)$ for some $l\in L$ with $r_l(w')=0$.

(3) Let $f$ be as in (3), viewed as a polynomial with complex coefficients. The points of $\hat X(\mathbf C)$ are the
common zeros of $P_{\mathbf C}$ in $\hat C$, so every such zero of $f$ is a common zero of $J_{\mathbf C}$, and
Lemma~\ref{cv-origin-analytic-system-crossings-mult} applies to $G=f$. It gives a unit $u$ of $\mathcal O_w$ and integers
$m'_l\ge0$ for $l\in L$ such that the image of $f^\kappa$ in $\mathcal O_w$ is $u\prod_{l\in L}\rho_l^{m'_l}$, where
$\rho_l$ is the image of $r_l$, and such that the integers $m$ for which $sf\in Q_{\tau(l)}^m+P\,\mathbf C[x_{k,i}]$ for some
$s\notin Q_{\tau(l)}$ are those with $m\le m'_l$. So $m'_l=m_{\tau(l)}$. Finally $f-f^\kappa$ lies in the ideal
generated by the $x_{k,\kappa_k}-1$, which is contained in $I_\kappa$, so $f$ and $f^\kappa$ have the same image in
$\mathcal O_w$.
\end{proof}
\leanref{Development/AutomorphicSpectrum/Proofs/MetricProofs/SncResolutionBoundaryNormalCrossing.lean}{GlobalTheory.CompanionTheory.Ar.Cmp.SncRes.Par.exists_sncData_of_isSncModel}
% Lean: snc_resolution_boundary_normal_crossing = SncRes.Par.exists_sncData_of_isSncModel: SncData (r, J, h, sigma with
% span_eq, eval_h, mem_bd, le_span) from GlobalTheory.T.snc_chart_parameters_at_cone_point,
% exists_radical_locally_eq_span_prod (RegularSystemOfParametersLocalRing), snc_chart_regular_at_cone_point and
% snc_chart_regular_with_boundary_choice for (2), snc_chart_multihomog_order_at_point, unit_times_parameter_monomial,
% symbolic_power_in_affine_chart and symbolic_power_from_affine_chart for (3), and minimal_primes_of_multihomog_sup.
% The text states it for f in R outside P; the Lean, for complex f outside the ideal generated by P (pcContract), and
% its proof does not use that f lies outside it. The components are the minimal primes of pC P + pC J (compIdeal), and
% the Lean's (2) is stated on the open set of points of the chart, coordinates coord p x', with F(coord) != 0
% (comp_eq_near), which the text reads on the affine coordinates. The steps are Lemmas
% cv-origin-analytic-system-crossings-params (exists_params), -local (primes_below_params and exists_den_prod_params),
% -components (components_of_params, comp_eq_near, minimal_primes_of_multihomog_sup) and -mult
% (exists_unit_mul_monomial, exponent_eq_order_polynomial, relation_dehom, relation_homogenises).
% Size: SncResolutionBoundaryNormalCrossing.lean, 66 lines, and the ten files it rests on (ChPrime,
% MinimalPrimesOfMultihomogSup, RegularSystemOfParametersLocalRing, SncChartRegularAtConePoint,
% SncChartRegularWithBoundaryChoice, SncChartMultihomogOrderAtPoint, UnitTimesParameterMonomial,
% SymbolicPowerInAffineChart, SymbolicPowerFromAffineChart, SncChartParametersAtConePoint), 1,332 lines, besides their
% own imports.

\begin{lemma}\label{cv-origin-intersection-resolution-exists-embedded}
Let $K$, $(I_k)_{k\in K}$ and $R$ be as in Definition~\ref{lit-multiprojective}, let $\Omega$ be a field of
characteristic zero, let $e\ge0$, let $y\in\Omega^{\sqcup_kI_k}$ be a point whose kernel $P$ is an arithmetic
subvariety $X$ of dimension $e+1$, and let $J$ be an ideal of $R$, generated by multihomogeneous elements and not
contained in $P$, that satisfies for $y$ and with $D=e$ the condition on the boundary ideal in
Definition~\ref{cv-origin-reference}. Then there is an analytic resolution of $X$ off $J$
(Definition~\ref{cv-origin-intersection-resolution}) whose maps satisfy (1) to (4) of that definition at every point,
whether or not its value lies in $Z(J)$.
\end{lemma}
\leanref{Development/AutomorphicSpectrum/Proofs/MetricProofs/SncModelEmbeddedLogResolutionStatement.lean}{GlobalTheory.T.snc_model_embedded_log_resolution}
\begin{proof}
Construct $\tilde X$, the open sets $U_x$ and the maps $\lambda_x$ as in the proof of
Lemma~\ref{cv-origin-model-resolution}, with $D=e$: $\tilde X$ is the image of $\hat X(\mathbf C)$ in
$\prod_k\mathbf P^{I_k}(\mathbf C)$, $U_x$ is a set $\tilde X_\kappa$ of classes with $z_{k,\kappa_k}\neq0$ that contains
$x$, and $\lambda_x$ has the affine coordinates $a_\kappa$, with $1$ at the $(k,\kappa_k)$. That proof uses the smoothness
of the rings $A_\kappa$ only through Lemma~\ref{cv-origin-model-resolution-regular}, whose conclusion holds here by
Lemma~\ref{cv-origin-intersection-resolution-exists-embedded-regular}. So $\tilde X$ is a compact complex manifold of
dimension $e$, a chart at $x\in\tilde X_\kappa$ is $r\circ a_\kappa$ for any polynomials $r_1,\dots,r_e$ whose images
generate the maximal ideal of the local ring at $a_\kappa(x)$ (Lemma~\ref{cv-origin-model-resolution-coords}), and the
maps $\lambda_x$ have the properties of Definition~\ref{cv-origin-analytic}(3). These are (1) to (4) of
Definition~\ref{cv-origin-intersection-resolution} at every point. It remains to prove (5) and (6).

Let $x\in\tilde X$, with $U_x=\tilde X_\kappa$, and let $w=a_\kappa(x)$. Let $r_1,\dots,r_e$, $L$, $\tau$ and $h$ be as in
Lemma~\ref{cv-origin-intersection-resolution-exists-embedded-snc} at $w$, with the minimal primes $Q_1,\dots,Q_c$ named
there. Then $\psi=r\circ a_\kappa$ is a holomorphic chart of $\tilde X$ on a neighbourhood $V\subset U_x$ of $x$, and
$\psi(x)=0$, as the $r_l$ vanish at $w$.

(5) Let $f\in R\setminus P$ be multihomogeneous with its zeros in $\hat X(\mathbf C)$ in $Z(J)$. For $b=1,\dots,c$ let
$W_b$ be the set of classes in $\tilde X$ of zeros of $Q_b$, and let $m_b$ be as in
Lemma~\ref{cv-origin-intersection-resolution-exists-embedded-snc}(3); neither depends on $x$. By (3) of that lemma there
are polynomials $g_1,g_2,g_3$ that do not vanish at $w$ with
$g_3\,(g_2f-g_1\prod_{l\in L}r_l^{m_{\tau(l)}})\in I_\kappa$. Shrinking $V$, the $g_i$ do not vanish on $a_\kappa(V)$, so
$f\circ\lambda_x=u\prod_{l\in L}\psi_l^{m_{\tau(l)}}$ on $V$ with $u=(g_1/g_2)\circ a_\kappa$ holomorphic and without
zeros. By (2) of that lemma, shrinking $V$ again, $W_b\cap V$ is empty if $b$ is not a value of $\tau$, and is the zero set
of $\psi_l$ if $b=\tau(l)$, which contains $x$. As $\tau$ is injective, distinct $W_b$ meeting $V$ are the zero sets of
distinct coordinates, and the product is over the $b$ with $W_b\cap V$ nonempty.

(6) Shrink $V$ so that $h$ does not vanish on $a_\kappa(V)$. If $x'\in V$ and $\lambda_x(x')\in Z(J)$, then the elements of
$I_\kappa$ and of $J$, hence of $I^J_\kappa$, vanish at $a_\kappa(x')$; by (1) of
Lemma~\ref{cv-origin-intersection-resolution-exists-embedded-snc}, $h\prod_{l\in L}r_l$ vanishes there, so some
$\psi_l(x')$ is $0$. So in the chart $\psi$ the set of (6) meets $V$ in a finite union of coordinate hyperplanes, of measure
zero, and its image in any other chart has measure zero, as the change of charts is holomorphic. Finitely many such $V$
cover the compact space $\tilde X$.
\end{proof}
\leanref{Development/AutomorphicSpectrum/Proofs/MetricProofs/SncModelEmbeddedLogResolution.lean}{GlobalTheory.CompanionTheory.Ar.Cmp.SncRes.isEmbedded_of_isSncModel}
\leanref{Development/AutomorphicSpectrum/Proofs/MetricProofs/SncModelEmbeddedLogResolution.lean}{GlobalTheory.CompanionTheory.Ar.Cmp.SncRes.exists_isEmbedded_of_algRes}
\leanref{ReciprocityTheory/AnalyticForms/RegularLocalRings/SNCstep.lean}{ReciprocityTheory.AnForm.R.AlgRes.isSNC_pullback}
\leanref{ReciprocityTheory/AnalyticForms/RegularLocalRings/SNCstep.lean}{ReciprocityTheory.AnForm.R.AlgRes.volume_preimage_eq_zero}
% Lean: GlobalTheory.T.snc_model_embedded_log_resolution = Cmp.SncRes.isEmbedded_of_isSncModel (IsSncModel y e J gives
% AnLogRes (mCoordIdeal y) J e with IsEmbedded): the AlgRes on Pts (pC P) from snc_resolution_compact_embedded,
% snc_resolution_transition_polynomial, snc_resolution_cover_and_cone, snc_resolution_local_chart_polynomial (the
% counterparts, for this case, of the files of SmoothMan.nonempty_algRes used in Lemma cv-origin-model-resolution),
% snc_resolution_regular_local_ring (Lemma cv-origin-intersection-resolution-exists-embedded-regular) and
% snc_resolution_boundary_normal_crossing (Lemma cv-origin-intersection-resolution-exists-embedded-snc);
% exists_isEmbedded_of_algRes: AlgRes.exists_atlas, the fields compat, surj, inj and immersion
% (fderiv_affCoord_ker_eq_bot) at every point, snc by AlgRes.isSNC_pullback (exists_poly_of_locRing_eq gives the text's
% g_1, g_2, g_3; comp_eq and prod_reindex_option for the components), null by AlgRes.volume_preimage_eq_zero
% (mem_bd, zero_bd, addHaar_image_eq_zero_of_differentiableOn_of_addHaar_eq_zero). Size:
% SncModelEmbeddedLogResolution.lean, 221 lines, the six SncResolution files, 1,027 lines with their statements
% (Lemmas -regular and -snc among them), FiniteMinimalPrimes.lean, 89, and SNCstep.lean, 288, besides the RegularLocalRings
% files shared with Lemma cv-origin-model-resolution.

\begin{lemma}\label{cv-origin-intersection-resolution-exists-descent}
Let $K$, $(I_k)_{k\in K}$, $R$, $\Omega$, $e$, $y$, $P$ and $X$ be as in
Lemma~\ref{cv-origin-intersection-resolution-exists-join}, let $M\ge1$ and let $y'\in\Omega^M$ be nonzero, with $R_1$,
$P_1$ and $\hat C_1$ as there, and let $J_1$ be an ideal of $R_1$ such that $P_1$ has an analytic resolution off $J_1$
whose maps satisfy (1) to (4) of Definition~\ref{cv-origin-intersection-resolution} at every point. Let
$\gamma\in R\setminus P$ be multihomogeneous, and suppose that every point of $\hat C_1$ in the zero set of $P_1$ at whose
old coordinates $\gamma$ vanishes lies in the zero set of $J_1$, and that (5) of
Lemma~\ref{cv-origin-intersection-resolution-exists-join} holds with $\gamma$ in place of $\beta'$. Then $X$ has an
analytic resolution off the ideal $(\gamma)$.
\end{lemma}
\leanref{Development/AutomorphicSpectrum/Proofs/MetricProofs/MultihomogeneousNonvanishingOnJoinStatement.lean}{GlobalTheory.T.multihomogeneous_nonvanishing_on_join}
\begin{proof}
For $u\in\hat C_1$ write $u^\circ$ for its old coordinates, the blocks $k\in K$, and $u_\ast$ for its block $\ast$, and
write $\hat X_1(\mathbf C)$ for the zero set of $P_1$ in $\hat C_1$. An element $f\in R$ lies in $P_1$ if and only if it
lies in $P$, since $x_{k,i}\mapsto y_{k,i}t_k$ treats it in the same way for the two points; so $u^\circ\in\hat X(\mathbf
C)$ for every $u\in\hat X_1(\mathbf C)$.

\emph{The forms.} Let $G_1,\dots,G_M\in R$ be multihomogeneous of one multidegree $g$ with $y'_iG_j(y)=y'_jG_i(y)$ for all
$i,j$.
\begin{enumerate}
\item[(a)] For $u\in\hat X_1(\mathbf C)$, $u_{\ast,i}G_j(u^\circ)=u_{\ast,j}G_i(u^\circ)$ for all $i,j$. Indeed
  $x_{\ast,i}G_j-x_{\ast,j}G_i\in P_1$, as $x_{k,i}\mapsto y_{k,i}t_k$, $x_{\ast,i}\mapsto y'_it_\ast$ sends it to
  $(y'_iG_j(y)-y'_jG_i(y))\,t_\ast\prod_kt_k^{g_k}=0$.
\item[(b)] Hence, if $u\in\hat X_1(\mathbf C)$ and $G_l(u^\circ)\neq0$, then $u_\ast=\mu\,(G_i(u^\circ))_i$ with
  $\mu=u_{\ast,l}/G_l(u^\circ)$, and $\mu\neq0$ as $u_\ast\neq0$.
\item[(c)] If $z\in\hat X(\mathbf C)$ and $G_i(z)\neq0$ for some $i$, then the point $(z,(G_i(z))_i)$ lies in
  $\hat X_1(\mathbf C)$. Indeed, let $\sigma:R_1\to R$ be the identity on $R$ with $\sigma(x_{\ast,i})=G_i$, choose $i_0$
  with $y'_{i_0}\neq0$ and put $c=G_{i_0}(y)/y'_{i_0}$, so that $G_i(y)=cy'_i$ for all $i$. Composing
  $x_{k,i}\mapsto y_{k,i}t_k$ with $\sigma$ gives the map of $(y,y')$ followed by $t_k\mapsto t_k$,
  $t_\ast\mapsto c\prod_kt_k^{g_k}$. So $\sigma(P_1)\subset P$, and $F(z,(G_i(z))_i)=\sigma(F)(z)=0$ for $F\in P_1$.
\item[(d)] If $u,u'\in\hat X_1(\mathbf C)$, $u^\circ=c\cdot u'^\circ$ for some $c\in(\mathbf C^\times)^K$ and
  $G_l(u^\circ)\neq0$, then $u$ and $u'$ lie in one orbit of $(\mathbf C^\times)^{K\sqcup\{\ast\}}$. Indeed
  $G_i(u^\circ)=c^gG_i(u'^\circ)$ with $c^g=\prod_kc_k^{g_k}\neq0$, so $G_l(u'^\circ)\neq0$, and by (b)
  $u_\ast=\mu\,G(u^\circ)$ and $u'_\ast=\nu\,G(u'^\circ)$ with $\mu,\nu\neq0$; so $u_\ast=(\mu c^g/\nu)\,u'_\ast$.
\end{enumerate}

\emph{The resolution.} Let $(\tilde X,(U_x),(\lambda_x))$ be the given resolution of $P_1$ off $J_1$. Keep $\tilde X$ and
the $U_x$, and put $\lambda'_x=\lambda_x^\circ$, a holomorphic map to $\hat X(\mathbf C)$ by the first paragraph. The ideal
$(\gamma)$ is generated by a multihomogeneous element outside $P$. We check (1) to (6) of
Definition~\ref{cv-origin-intersection-resolution} for $X$ off $(\gamma)$.

(1) If two points of $\hat C_1$ lie in one orbit, so do their old coordinates.

(2) Let $z\in\hat X(\mathbf C)$ with $\gamma(z)\neq0$. Hypothesis (5) gives forms $G_i$ as above with some $G_i(z)\neq0$,
and by (c), $u=(z,(G_i(z))_i)$ lies in $\hat X_1(\mathbf C)$. By (2) for the given resolution at every point, $u$ lies in
the orbit of some $\lambda_x(x)$, so $z=u^\circ$ lies in the orbit of $\lambda'_x(x)$.

(3) Let $\gamma(\lambda'_x(x))\neq0$, and let $\lambda'_{x'}(x')$ lie in the orbit of $\lambda'_x(x)$. Hypothesis (5) at
$\lambda'_x(x)$ gives forms with $G_l(\lambda'_x(x))\neq0$, and by (d), $\lambda_x(x)$ and $\lambda_{x'}(x')$ lie in one
orbit. By (3) for the given resolution at every point, $x=x'$.

(4) Let $\gamma(\lambda'_x(x))\neq0$, and let $\kappa_k\in I_k$ with $\lambda_x(x)_{k,\kappa_k}\neq0$ for all $k\in K$.
Take forms from (5) with some $G_i(\lambda'_x(x))\neq0$. By (b), $\lambda_x(x)_\ast$ is a nonzero multiple of
$(G_i(\lambda'_x(x)))_i$, so for an $l$ with $\lambda_x(x)_{\ast,l}\neq0$ also $G_l(\lambda'_x(x))\neq0$. In a chart of
$\tilde X$ near $x$, let $a$ be the affine coordinates $z_{k,i}/z_{k,\kappa_k}$ of $\lambda'_x$, and let $a_1$ be those of
$\lambda_x$ for the indices $\kappa$ and $l$. By multihomogeneity, $G_i(\lambda'_x)=\prod_k\lambda_{x,k,\kappa_k}^{g_k}\,
G_i(a)$, so $G_l(a)\neq0$ near $x$, and by (a) the ratios $\lambda_{x,\ast,i}/\lambda_{x,\ast,l}$ are $G_i(a)/G_l(a)$.
So $a_1=\Phi\circ a$ near $x$, where $\Phi(b)=(b,(G_i(b)/G_l(b))_i)$ is holomorphic near the value of $a$ at $x$. The
derivative of $a_1$ at $x$ is injective by (4) for the given resolution at every point, and it is the derivative of
$\Phi$ composed with that of $a$; so the derivative of $a$ at $x$ is injective.

(5) Let $f\in R\setminus P$ be multihomogeneous with its zeros in $\hat X(\mathbf C)$ in the zero set of $\gamma$. In
$R_1$, $f$ is multihomogeneous, of multidegree $0$ at $\ast$, and not in $P_1$. If $u\in\hat X_1(\mathbf C)$ and
$f(u)=0$, then $f(u^\circ)=0$, so $\gamma(u^\circ)=0$, and $u$ lies in the zero set of $J_1$ by hypothesis. So (5) for the
given resolution applies to $f$, and as $f\circ\lambda_x=f\circ\lambda'_x$, its sets $W_l$, multiplicities and charts
serve.

(6) If $\gamma(\lambda'_x(x))=0$, then $\lambda_x(x)$ lies in the zero set of $J_1$ by hypothesis. So the set of (6) is
contained in that of (6) for the given resolution, which has measure zero.
\end{proof}
\leanref{Development/AutomorphicSpectrum/Proofs/MetricProofs/MultihomogeneousNonvanishingOnJoin.lean}{GlobalTheory.CompanionTheory.Ar.Giv.anLogRes_of_embedded_join}
% Lean: GlobalTheory.T.multihomogeneous_nonvanishing_on_join = Giv.anLogRes_of_embedded_join, in ResProj:
% rename_mem_iff and comp_inl_mem (old coordinates), cross_mem and cross_eval (a), proportional (b), subst, subst_mem
% and joinVec_mem (c, through cone_homog), sameOrbit_of_inl (d); the fields: compat by sameOrbit_comp_inl, surj by
% hE.surj and joinVec_mem, inj by hE.inj and sameOrbit_of_inl, immersion by immersion_inl (hEv: Phi = F o Psi near the
% chart point, and fderiv_comp with hE.immersion), snc by R.snc with the degree Sum.elim a 0 and homog_rename_inl, null
% by R.null and the hypothesis on J. Size: MultihomogeneousNonvanishingOnJoin.lean, about 420 lines.

\begin{lemma}\label{cv-origin-intersection-resolution-exists}
Let $K$, $(I_k)_{k\in K}$ and $R$ be as in Definition~\ref{lit-multiprojective}, let $e\ge0$, let $X$ be an arithmetic
subvariety of dimension $e+1$ defined by the prime $P$, and let $J$ be an ideal of $R$ that is generated by
multihomogeneous elements and is not contained in $P$. Then there are an ideal $J'\subset J$, generated by
multihomogeneous elements and not contained in $P$, and an analytic resolution of $X$ off $J'$
(Definition~\ref{cv-origin-intersection-resolution}).
\end{lemma}
\leanref{ReciprocityTheory/ReductiveGroups/ArithmeticGeometry/BirationalGeometry/InterRes.lean}{ReciprocityTheory.ChartRestriction.G.Geo.ResExistsStatement}
\leanref{Development/AutomorphicSpectrum/Proofs/MetricProofs/ResolutionExistsOfResolutionStatementStatement.lean}{GlobalTheory.T.resolution_exists_of_resolution_statement}
\leanref{Development/AutomorphicSpectrum/Proofs/MetricProofs/ResolutionExistsOfResolutionStatement.lean}{GlobalTheory.CompanionTheory.Ar.Giv.resExists_of_resK}
\begin{proof}
The prime $P$ is the kernel of a point $y\in\Omega^{\sqcup_kI_k}$ with function field $\Omega$ (as in the proof of
Lemma~\ref{cv-origin-intersection-contraction-ratios}), and $\Omega$ has characteristic zero, as $P\cap\mathbf Z=0$. As
$J$ is generated by multihomogeneous elements and is not contained in $P$, some multihomogeneous $\beta\in J$ is not in
$P$. Lemma~\ref{cv-origin-intersection-resolution-exists-join} gives $\beta'$, $M$, $y'$ and $J_1$. The joined point
$(y,y')$ has a nonzero coordinate in every block, its restriction to the old blocks is $y$, and by (2) its field of
ratios is that of $y$; so by Lemma~\ref{cv-origin-intersection-resolution-exists-arith}, for the inclusion of the old
blocks, its kernel $P_1$ is an arithmetic subvariety of dimension $e+1$. By (3) and
Lemma~\ref{cv-origin-intersection-resolution-exists-embedded}, $P_1$ has an analytic resolution off $J_1$ whose maps
satisfy (1) to (4) of Definition~\ref{cv-origin-intersection-resolution} everywhere.

Put $\gamma=\beta\beta'$ and $J'=(\gamma)$. Then $\gamma\in J$, it is multihomogeneous, and it is not in the prime $P$, so
$J'\subset J$ is generated by a multihomogeneous element and is not contained in $P$. If $\gamma$ vanishes at the old
coordinates of a point of $\hat C_1$ in the zero set of $P_1$, then $\beta$ or $\beta'$ does; in the first case so does
$\beta'$, by (1), as the old coordinates form a point of $\hat X(\mathbf C)$; so the point lies in the zero set of $J_1$
by (4). Where $\gamma\neq0$, also $\beta'\neq0$, and (5) gives the $G_i$. Lemma~\ref{cv-origin-intersection-resolution-exists-descent}
gives an analytic resolution of $X$ off $J'$.
\end{proof}
% Lean: Giv.resExists_of_resK (= GlobalTheory.T.resolution_exists_of_resolution_statement), under ResKStatement
% (Lemma cv-resolution-coordinates): HilbertSamuel.BirationalModelExists.proof writes P as the kernel of a point y;
% beta in J is multihomogeneous and outside P; GlobalTheory.T.multihomogeneous_lift_to_join gives beta', a further
% block y' of M coordinates and an ideal J_1 with an snc condition; isArithVar_of_isAlgebraic_image makes the joined
% point (y, y') an arithmetic variety of the same dimension; GlobalTheory.T.snc_model_embedded_log_resolution gives
% its resolution; J' is the principal ideal of beta * beta', and GlobalTheory.T.multihomogeneous_nonvanishing_on_join
% gives AnLogRes P J' e (hmul, hzero, hz and hG for the two hypotheses). Size: about 9,500 lines, in Lemmas
% cv-origin-intersection-resolution-exists-join, -embedded and -descent (an estimate from file sizes; the files named in
% their comments have about 4,800 lines, 903 of them shared with Lemma cv-origin-coords-snc, besides further imports
% shared with Lemmas cv-origin-coords-snc and cv-origin-model-resolution).

\subsection{The Poincar\'e--Lelong formula}\label{subsection-cv-poincare-lelong}

\begin{theorem}[Poincar\'e--Lelong]\label{cv-origin-intersection-poincare-lelong}
Let $K$, $(I_k)_{k\in K}$ and $R$ be as in Definition~\ref{lit-multiprojective}, let $n\ge0$, let $X$ be an arithmetic
subvariety of dimension $n+2$ defined by the prime $P$, and let $\bar L_1,\dots,\bar L_n$, $\bar A=(a,h_A)$ and
$\bar B=(0,h_B)$ be smooth metrized line bundles that are invariant under conjugation, with $a\in\mathbf Z_{\ge0}^K$.
Let $s\in R\setminus P$ be multihomogeneous of multidegree $a$, put
$\Phi=c_1(\bar L_1)\wedge\dots\wedge c_1(\bar L_n)$, and for an arithmetic subvariety $Y$ write $\int_{Y(\mathbf C)}$
for the integral over the complex points of $Y$. Then
\[
 \int_{X(\mathbf C)}\log\frac{|s|}{h_A}\;c_1(\bar B)\wedge\Phi
 =\sum_Wm_W\int_{W(\mathbf C)}\log h_B\;\Phi-\int_{X(\mathbf C)}\log h_B\;c_1(\bar A)\wedge\Phi,
\]
the sum over the horizontal relevant minimal primes $W$ of $P+(s)$, where $m_W$ is the multiplicity of $W$ in
$R/(P+(s))$.
\end{theorem}
\leanref{Development/AutomorphicSpectrum/PublishedTheorems/Metrics₂/PoincareLelongInv.lean}{GlobalTheory.CompanionTheory.Hypotheses₂.PoincareLelongInv}

As $\bar B$ has multidegree $0$, $\log h_B$ is a smooth function on the complex points of $X$, and
$c_1(\bar B)=dd^c\log h_B$. So this is the formula of Theorem~\ref{cv-tail-poincare-lelong} with $v=\log h_B$, on the
complex points of $X$, which may be singular. Pulled back to a resolution of singularities of the complex points of
$X$ (Corollary~\ref{lit-resolution}), it is the Lelong--Poincar\'e equation \cite[Chapter~III, \S2]{Dem12}: the
divisor of the pulled-back section pushes forward to the cycle $\sum_Wm_W[W(\mathbf C)]$ \cite[\S1.4]{Ful98}, and
the components that the resolution contracts contribute nothing, since the forms pulled back to them vanish.

\begin{definition}\label{cv-origin-intersection-poincare-lelong-all-twist}
Let $K$, $(I_k)_{k\in K}$, $R$ and $\hat C$ be as in Definition~\ref{lit-multiprojective}, and let $p$ be a prime
number. Let $K^+=K\sqcup\{\ast\}$ with $I_\ast=\{0,1\}$, so that $R^+=R[x_{\ast,0},x_{\ast,1}]$ contains $R$, and let
$\hat C^+$ be the corresponding set of points.
\begin{enumerate}
\item An ideal $W$ of $R$ is \emph{twistable at $p$} if $a^2+pb^2\in W$ with $a,b\in R$ implies $b\in W$. Its
  \emph{twist} is the ideal $W^{(p)}=WR^++(x_{\ast,0}^2+p\,x_{\ast,1}^2)$ of $R^+$.
\item For $z\in\hat C$ and $w\in\mathbf C^2\setminus\{0\}$, write $(z,w)$ for the point of $\hat C^+$ with blocks $z$ and
  $w$, and put $w_\pm=(\pm i\sqrt p,1)$, the two points of the conic $w_0^2+pw_1^2=0$ with $w_1=1$.
\item For $w\in\mathbf C^2\setminus\{0\}$ put
  \[
   \chi_p(w)=\frac12+\frac{p+1}{2\sqrt p}\cdot\frac{\operatorname{Im}(w_0\bar w_1)}{|w_0|^2+|w_1|^2}.
  \]
  The \emph{twist} of a metrized line bundle $\bar A=(a,h)$ is $\bar A^{(p)}=(a^+,h^{(p)})$, where $a^+$ is $a$ on $K$
  and $0$ at $\ast$, and $h^{(p)}(z,w)=h(z)^{\chi_p(w)}\,h(\bar z)^{1-\chi_p(w)}$.
\end{enumerate}
\end{definition}
\leanref{Development/AutomorphicSpectrum/Proofs/MetricProofs/IsHorizontalTwist.lean}{GlobalTheory.CompanionTheory.Ar.Tw.twist}
\leanref{Development/AutomorphicSpectrum/Proofs/MetricProofs/IsHorizontalTwist.lean}{GlobalTheory.CompanionTheory.Ar.Tw.NoRoot}
\leanref{Development/AutomorphicSpectrum/Proofs/MetricProofs/IsHorizontalTwist.lean}{GlobalTheory.CompanionTheory.Ar.Tw.chi}
\leanref{Development/AutomorphicSpectrum/Proofs/MetricProofs/IsHorizontalTwist.lean}{GlobalTheory.CompanionTheory.Ar.Tw.twLine}

\begin{lemma}\label{cv-origin-intersection-poincare-lelong-all-prime}
Let $K$, $(I_k)_{k\in K}$ and $R$ be as in Definition~\ref{lit-multiprojective}, and let $W_1,\dots,W_m$ be prime ideals
of $R$ with $W_j\cap\mathbf Z=0$. Then there is a prime number $p$ at which every $W_j$ is twistable
(Definition~\ref{cv-origin-intersection-poincare-lelong-all-twist}(1)).
\end{lemma}
\leanref{Development/AutomorphicSpectrum/Proofs/MetricProofs/PoincareLelongOfInv.lean}{GlobalTheory.CompanionTheory.Ar.Tw.Main.exists_prime_noRoot}
\begin{proof}
Let $W$ be one of the $W_j$, and let $F$ be the fraction field of the domain $R/W$, of characteristic zero and
finitely generated over $\mathbf Q$. If $W$ is not twistable at $p$, there are $a,b\in R/W$ with $b\neq0$ and
$a^2+pb^2=0$, and $x=a/b\in F$ has $x^2=-p$. The algebraic closure $E$ of $\mathbf Q$ in $F$ is finite over $\mathbf Q$,
as $F$ is finitely generated, and it contains $x$. For distinct primes $p$ and $q$ with square roots $x$ and $y$ of
$-p$ and $-q$ in $E$, the fields $\mathbf Q(x)$ and $\mathbf Q(y)$ differ: otherwise $x=\alpha+\beta y$ with
$\alpha,\beta\in\mathbf Q$, and $x^2=-p$ gives $\alpha^2-q\beta^2+p=0$ and $2\alpha\beta=0$, so either $\alpha=0$ and
$p=q\beta^2$, which the $q$-adic valuation excludes, or $\beta=0$ and $\alpha^2=-p$, which is impossible. As $E$ has
finitely many subfields, $W$ is not twistable at only finitely many primes. There are infinitely many primes, so one
of them is outside the finitely many exceptions for $W_1,\dots,W_m$.
\end{proof}
\leanref{Development/AutomorphicSpectrum/Proofs/MetricProofs/FinitePrimesSumSquaresZeroStatement.lean}{GlobalTheory.T.finite_primes_sum_squares_zero}
\leanref{Development/AutomorphicSpectrum/Proofs/MetricProofs/FinitePrimesSumSquaresZero.lean}{GlobalTheory.CompanionTheory.QuadFin.finite_primes_neg_sq}
% Lean: exists_prime_noRoot with finite_primes_sum_squares_zero = QuadFin.finite_primes_neg_sq:
% finite_primes_of_isFractionRing (x = a/b), finite_primes_field (algebraicClosure Q F is finite dimensional,
% finiteDimensional_algebraicClosure_of_essFiniteType), finite_primes_of_finiteDimensional (finitely many intermediate
% fields, primitive element theorem), exists_eq_add_mul_of_mem_adjoin, false_of_sq_eq and not_sq_eq_of_prime.

\begin{lemma}\label{cv-origin-intersection-poincare-lelong-all-metric}
Let $K$ and $(I_k)_{k\in K}$ be as in Definition~\ref{lit-multiprojective}, let $p$ be a prime number, and let
$\bar A=(a,h)$ be a metrized line bundle with twist $\bar A^{(p)}=(a^+,h^{(p)})$
(Definition~\ref{cv-origin-intersection-poincare-lelong-all-twist}(3)).
\begin{enumerate}
\item $\bar A^{(p)}$ is a metrized line bundle, smooth if $\bar A$ is, and invariant under conjugation.
\item For $z\in\hat C$, $h^{(p)}(z,w_+)=h(z)$ and $h^{(p)}(z,w_-)=h(\bar z)$.
\item The twist of $(-a,1/h)$ is $(-a^+,1/h^{(p)})$.
\end{enumerate}
\end{lemma}
\leanref{Development/AutomorphicSpectrum/Proofs/MetricProofs/TwLineSmoothConjInvariantGlueStatement.lean}{GlobalTheory.T.twLine_smooth_conjInvariant_glue}
\begin{proof}
This is a self-contained computation. The function $\chi_p$ is smooth on $\mathbf C^2\setminus\{0\}$ and does not
change when $w$ is multiplied by a nonzero scalar. For $c\in(\mathbf C^\times)^K$, $h(cz)$ and $h(\bar c\bar z)$ are
both $\prod_k|c_k|^{a_k}$ times $h(z)$ and $h(\bar z)$, so $h^{(p)}$ is homogeneous of multidegree $a^+$, and it is
smooth where $h$ is, as $h^{(p)}=\exp\big(\chi_p\log h(z)+(1-\chi_p)\log h(\bar z)\big)$. Conjugating $w$ changes the sign
of $\operatorname{Im}(w_0\bar w_1)$, so $\chi_p(\bar w)=1-\chi_p(w)$ and $h^{(p)}(\bar z,\bar w)=h^{(p)}(z,w)$. At $w_\pm$,
$w_0\bar w_1=\pm i\sqrt p$ and $|w_0|^2+|w_1|^2=p+1$, so $\chi_p(w_+)=1$ and $\chi_p(w_-)=0$, which gives (2). For (3),
$h^{(p)}$ is multiplicative in $h$.
\end{proof}
\leanref{Development/AutomorphicSpectrum/Proofs/MetricProofs/TwLineSmoothConjInvariantGlue.lean}{GlobalTheory.CompanionTheory.Ar.Tw.Lines.twLine_main}
% Lean: twLine_main: chi_contDiffOn, twLine_smooth, chi_conj and twLine_conj, chi_vPlus and chi_vMinus with
% twLine_vPlus and twLine_vMinus; (3) is logNorm_tw_inv_one in PoincareLelongOfInv.lean (twH of invLine).

\begin{lemma}\label{cv-origin-intersection-poincare-lelong-all-components-prime}
Let $K$, $(I_k)_{k\in K}$ and $R$ be as in Definition~\ref{lit-multiprojective}, let $p$ be a prime number, and let $W$
be a prime ideal of $R$ that is twistable at $p$, with twist $W^{(p)}\subset R^+$
(Definition~\ref{cv-origin-intersection-poincare-lelong-all-twist}(1)). Then $W^{(p)}$ is a prime ideal, and
$\dim R^+/W^{(p)}=\dim R/W+1$.
\end{lemma}
\leanref{Development/AutomorphicSpectrum/Proofs/MetricProofs/TwistIsPrimeStatement.lean}{GlobalTheory.T.twist_isPrime_of_isPrime_of_noRoot}
\leanref{Development/AutomorphicSpectrum/Proofs/MetricProofs/RingKrullDimTwistSuccStatement.lean}{GlobalTheory.T.ringKrullDim_twist_succ}
\begin{proof}
Let $D=R/W$, a noetherian domain, and let $F$ be its fraction field. Write $q_p=x_{\ast,0}^2+p\,x_{\ast,1}^2$.

\emph{Primality.} The polynomial $X^2+p$ has no root in $F$: if $(a/b)^2+p=0$ with $a,b\in D$ and $b\neq0$, then
$a^2+pb^2=0$ in $D$, and as $W$ is twistable at $p$, $b=0$. So $E=F[\sigma]/(\sigma^2+p)$ is a field, and $u+v\sigma=0$
with $u,v\in D$ implies $u=v=0$. Let $\varphi:R^+\to E[t]$ send $x_{k,i}$ to its class in $D$, $x_{\ast,0}$ to $\sigma t$
and $x_{\ast,1}$ to $t$. It kills $W$, and $\varphi(q_p)=(\sigma^2+p)t^2=0$, so it induces
$\varphi':R^+/W^{(p)}\to E[t]$. In $R^+/W^{(p)}$ the relation $x_{\ast,0}^2=-p\,x_{\ast,1}^2$ holds, so, by induction on
monomials, every element is $a(x_{\ast,1})+b(x_{\ast,1})\,x_{\ast,0}$ with $a,b\in D[T]$. Its image is
$a(t)+b(t)\,\sigma t$. If this is $0$, the coefficient of $t^0$ gives $a_0=0$, and the coefficient of $t^{n+1}$ gives
$a_{n+1}+b_n\sigma=0$, so $a_{n+1}=b_n=0$. So $\varphi'$ is injective, $R^+/W^{(p)}$ is a subring of the domain $E[t]$,
and $W^{(p)}$ is prime.

\emph{Dimension.} Let $\beta:D[T]\to R^+/W^{(p)}$ be the homomorphism that is the identity on $D$ and sends $T$ to
$x_{\ast,1}$. Let $\alpha:R^+/W^{(p)}\to D[T][S]/(S^2+pT^2)$ send $x_{k,i}$ to its class, $x_{\ast,0}$ to $S$ and
$x_{\ast,1}$ to $T$; it is well defined, as it kills $W$ and $q_p$. Then $\alpha\circ\beta$ is the inclusion of $D[T]$ in
the free $D[T]$-module with basis $1,S$, so $\beta$ is injective. The ring $R^+/W^{(p)}$ is generated over the image of
$\beta$ by $x_{\ast,0}$, a root of the monic polynomial $S^2+p\,x_{\ast,1}^2$, so it is integral over $D[T]$. An
integral extension has the dimension of its base \cite[Tag 00OK]{Sta}, so $\dim R^+/W^{(p)}=\dim D[T]$, and $\dim D[T]=\dim D+1$ as $D$ is noetherian \cite[Theorem~15.4]{Mat86}.
\end{proof}
\leanref{Development/AutomorphicSpectrum/Proofs/MetricProofs/TwistIsPrime.lean}{GlobalTheory.CompanionTheory.Ar.Tw.twist_isPrime}
\leanref{Development/AutomorphicSpectrum/Proofs/MetricProofs/RingKrullDimTwistSucc.lean}{GlobalTheory.CompanionTheory.Ar.Tw.ringKrullDim_twist}
% Lean: twist_isPrime: exists_root_field (X^2 + p irreducible over FractionRing D, E = AdjoinRoot), phi to
% Polynomial E (xN to C s * X, yN to X), hform (every class is sigma a + sigma b * xN) and the coefficient comparison
% for hinj. ringKrullDim_twist: baseP (D[T] to R+/W^(p), T to yN), baseMap_injective through quotA to
% AdjoinRoot (quadP p W), baseMap_isIntegral, ReciprocityTheory.CA.ringKrullDim_eq_of_isIntegral_of_injective (its own
% proof of the equality of dimensions for an integral extension) and Polynomial.ringKrullDim_of_isNoetherianRing. Size:
% TwistIsPrime.lean and RingKrullDimTwistSucc.lean, about 400 lines.

\begin{lemma}\label{cv-origin-intersection-poincare-lelong-all-components}
Let $K$, $(I_k)_{k\in K}$ and $R$ be as in Definition~\ref{lit-multiprojective}, let $p$ be a prime number, and let $X$
be an arithmetic subvariety of dimension $d$ defined by a prime $P$ that is twistable at $p$
(Definition~\ref{cv-origin-intersection-poincare-lelong-all-twist}(1)).
\begin{enumerate}
\item The twist $P^{(p)}$ is an arithmetic subvariety $X^{(p)}$ of dimension $d$ of
  $\prod_{k\in K}\mathbf P^{I_k}_{\mathbf Z}\times\mathbf P^1_{\mathbf Z}$, and $P^{(p)}\cap R=P$.
\item Let $s\in R\setminus P$ be multihomogeneous, and suppose that every horizontal relevant minimal prime of $P+(s)$ is
  twistable at $p$. Then the horizontal relevant minimal primes of $P^{(p)}+(s)$ are the twists $W^{(p)}$ of those of
  $P+(s)$, and the multiplicity of $W^{(p)}$ in $R^+/(P^{(p)}+(s))$ is that of $W$ in $R/(P+(s))$.
\end{enumerate}
\end{lemma}
\leanref{Development/AutomorphicSpectrum/Proofs/MetricProofs/PoincareLelongOfInv.lean}{GlobalTheory.CompanionTheory.Ar.Tw.Main.isArithVar_twist}
\leanref{Development/AutomorphicSpectrum/Proofs/MetricProofs/HorizontalDivSuppTwistStatement.lean}{GlobalTheory.T.horizontal_divSupp_twist_eq_image}
\leanref{Development/AutomorphicSpectrum/Proofs/MetricProofs/MultAtTwistQuadEqStatement.lean}{GlobalTheory.T.multAt_twist_quad_eq}
\begin{proof}
Write $q_p=x_{\ast,0}^2+p\,x_{\ast,1}^2$, so that $W^{(p)}=WR^++(q_p)$ for every ideal $W$ of $R$. Let $\rho:R^+\to R$
be the homomorphism that is the identity on $R$ and sends $x_{\ast,0}$ and $x_{\ast,1}$ to $0$. It kills $q_p$, so
$\rho(W^{(p)})\subset W$. As $\rho$ is the identity on $R$, this gives $W^{(p)}\cap R\subset W$; the other inclusion is
clear, so $W^{(p)}\cap R=W$. Likewise $W^{(p)}\cap\mathbf Z\subset W\cap\mathbf Z$.

(1) By Lemma~\ref{cv-origin-intersection-poincare-lelong-all-components-prime}, $P^{(p)}$ is a prime ideal and
$\dim R^+/P^{(p)}=\dim R/P+1=d+|K|+1$, which is $d$ plus the number of blocks of $R^+$. It is generated by
multihomogeneous elements: the images of multihomogeneous generators of $P$, whose multidegree at $\ast$ is $0$, and
$q_p$, of multidegree $0$ on $K$ and $2$ at $\ast$. It is relevant. For $k\in K$, some $x_{k,i}$ is not in
$P=P^{(p)}\cap R$. For the block $\ast$, let $\theta:R^+\to(R/P)[T]$ send $x_{k,i}$ to its class, $x_{\ast,0}$ to $T$ and
$x_{\ast,1}$ to $0$. It kills $P$ and sends $q_p$ to $T^2$, so $\theta(P^{(p)})\subset(T^2)$; as
$\theta(x_{\ast,0})=T\notin(T^2)$, $x_{\ast,0}\notin P^{(p)}$. Finally $P^{(p)}\cap\mathbf Z\subset P\cap\mathbf Z=0$.
So $P^{(p)}$ is an arithmetic subvariety of dimension $d$, and $P^{(p)}\cap R=P$ by the first paragraph.

(2) As $P^{(p)}+(s)=(P+(s))R^++(q_p)$, an ideal $Q$ of $R^+$ contains $P^{(p)}+(s)$ if and only if it contains $q_p$
and $Q\cap R\supset P+(s)$; then $W^{(p)}\subset Q$ for every ideal $W\subset Q\cap R$ of $R$.

Let $W'$ be a horizontal relevant minimal prime of $P^{(p)}+(s)$. The prime $W'\cap R$ contains $P+(s)$, so it contains a
minimal prime $W$ of $P+(s)$. If all $x_{k,i}$ with $i\in I_k$ were in $W$, they would be in $W'$; so $W$ is relevant.
As $W\cap\mathbf Z\subset W'\cap\mathbf Z=0$, $W$ is horizontal, so twistable at $p$ by hypothesis. By the first
paragraph of (2), $W^{(p)}\subset W'$, and $W^{(p)}$ is a prime containing $P^{(p)}+(s)$ by
Lemma~\ref{cv-origin-intersection-poincare-lelong-all-components-prime}. By the minimality of $W'$, $W'=W^{(p)}$.

Conversely, let $W$ be a horizontal relevant minimal prime of $P+(s)$. By
Lemma~\ref{cv-origin-intersection-poincare-lelong-all-components-prime}, $W^{(p)}$ is a prime; it contains
$P^{(p)}+(s)$, and it is relevant and horizontal by the argument of (1). If $Q$ is a prime with
$P^{(p)}+(s)\subset Q\subset W^{(p)}$, then $P+(s)\subset Q\cap R\subset W^{(p)}\cap R=W$, so $Q\cap R=W$ by the
minimality of $W$, and $W^{(p)}\subset Q$. So $W^{(p)}$ is a horizontal relevant minimal prime of $P^{(p)}+(s)$. Distinct
$W$ have distinct twists, as $W^{(p)}\cap R=W$.

For the multiplicities, let $B=R^+/(q_p)$ and let $\psi:R\to B$ be the inclusion followed by the quotient map. The ring
$B$ is the quotient of $R[x_{\ast,1}][x_{\ast,0}]$ by a monic polynomial of degree $2$ in $x_{\ast,0}$, so it is free over
$R[x_{\ast,1}]$ with basis $1,x_{\ast,0}$, and $\psi$ is flat. For $W$ as above, $\mathfrak P=W^{(p)}/(q_p)$ is a prime of
$B$ with $\psi^{-1}(\mathfrak P)=W^{(p)}\cap R=W$ and $\mathfrak P=WB$. So $R_W\to B_{\mathfrak P}$ is a flat local
homomorphism, and the maximal ideal of $R_W$ generates that of $B_{\mathfrak P}$. Put $J=P+(s)$. Then
$R^+/(P^{(p)}+(s))=B/JB$, and its localization at $W^{(p)}$ is $(R/J)_W\otimes_{R_W}B_{\mathfrak P}$. As
$R_W\to B_{\mathfrak P}$ is a flat local homomorphism, its length is the length of $(R/J)_W$ times that of the quotient
of $B_{\mathfrak P}$ by the maximal ideal of $R_W$ \cite[Tag 02M1]{Sta}. This quotient is the residue field of $B_{\mathfrak P}$, of length $1$.
\end{proof}
\leanref{Development/AutomorphicSpectrum/Proofs/MetricProofs/HorizontalDivSuppTwist.lean}{GlobalTheory.CompanionTheory.Ar.Tw.horizontal_divSupp_twist}
\leanref{Development/AutomorphicSpectrum/Proofs/MetricProofs/MultAtTwistQuadEq.lean}{GlobalTheory.CompanionTheory.Ar.Tw.multAt_twist}
% Lean: isArithVar_twist from GlobalTheory.T.twist_isPrime_of_isPrime_of_noRoot (Tw.twist_isPrime),
% twist_isMultiHomogIdeal (inc of the generators and isMultiHomog_quad), isRelevant_twist (xN_not_mem_twist through
% theta to Polynomial (R/W), and comap_inc_twist), isHorizontal_twist (twist_le_comap_ret) and ringKrullDim_twist_succ
% (Tw.ringKrullDim_twist); comap_inc_twist (with ret, the text's rho) for P^(p) meet R = P; horizontal_divSupp_twist
% (exists_minimalPrimes_le, twist_le_of_map_le, sup_span_le_comap_inc) and multAt_twist (with the ideal
% (P + (s)).map inc + (quad p), which is P^(p) + (s)): length_cut_of_flat with flat_quad (polyEquiv and
% AdjoinRoot.powerBasis'), Localization.localRingHom and length_quot_map_of_flat (IsLocalRing.length_baseChange, the
% Lean's form of the length formula of Stacks Tag 02M1). Twist_injective is used in Lemma cv-origin-intersection-poincare-lelong-all. Size:
% about 1,000 lines (TwistIsPrime, RingKrullDimTwistSucc, TwistIsMultiHomogIdeal, HorizontalDivSuppTwist,
% MultAtTwistQuadEq and IsHorizontalTwist in Development/AutomorphicSpectrum/Proofs/MetricProofs).

\begin{definition}\label{cv-origin-intersection-poincare-lelong-all-double-projection}
Let $K$, $(I_k)_{k\in K}$, $R$ and $\hat C$ be as in Definition~\ref{lit-multiprojective}, let $e\ge0$, and let $X$ be an
arithmetic subvariety of dimension $e+1$ defined by the prime $P$. Put $I=\prod_{k\in K}I_k$, and let
$\sigma:\mathbf C^{\sqcup_kI_k}\to\mathbf C^I$ be the Segre map, $\sigma(z)_i=\prod_kz_{k,i_k}$.
\begin{enumerate}
\item A \emph{good projection} for $X$ is a $\mathbf C$-linear map $\Psi:\mathbf C^I\to\mathbf C^{e+1}$ with
  $\Psi(\sigma(z))\neq0$ for every $z\in\hat X(\mathbf C)$.
\item Let $\Psi$ be a good projection for $X$ and $y\in\mathbf C^e$. The \emph{fibre} $F_\Psi(y)$ is the set of the
  $\sigma(z)$ with $z\in\hat X(\mathbf C)$ and $\Psi(\sigma(z))=(1,y)$. A \emph{sheet} through $w\in\mathbf C^I$ at $y$ is a
  map $\gamma:\mathbf C^e\to\mathbf C^{\sqcup_kI_k}$ with $\sigma(\gamma(y))=w$ such that, for every $y'$ in a
  neighbourhood of $y$, $\gamma$ is complex differentiable at $y'$, $\gamma(y')\in\hat X(\mathbf C)$ and
  $\Psi(\sigma(\gamma(y')))=(1,y')$.
\item Let $\varphi_1,\dots,\varphi_e$ and $g$ be real functions on $\mathbf C^{\sqcup_kI_k}$. Their \emph{sheet density}
  at $(y,w)$ is $g(\gamma(y))\operatorname{MA}(\varphi_1\circ\gamma,\dots,\varphi_e\circ\gamma)(y)$
  (Definition~\ref{cv-period-monge-ampere}(2)) for a sheet $\gamma$ through $w$ at $y$, chosen once for each $(y,w)$, and
  $0$ if there is none. The \emph{integral through $\Psi$} of $g$ against $\varphi_1,\dots,\varphi_e$ is the integral over
  $y\in\mathbf C^e$ of the sum of the sheet densities at $(y,w)$ over $w\in F_\Psi(y)$.
\item Let $\bar L_j=(a_j,h_j)$, $j=1,\dots,e$, be smooth metrized line bundles and $g$ a real function on
  $\mathbf C^{\sqcup_kI_k}$. The \emph{integral over the complex points of $X$} of
  $g\,c_1(\bar L_1)\wedge\dots\wedge c_1(\bar L_e)$, written $\int_{X(\mathbf C)}g\,c_1(\bar L_1)\wedge\dots\wedge
  c_1(\bar L_e)$, is the integral through a good projection for $X$, chosen once for $X$, of $g$ against
  $\log h_1,\dots,\log h_e$, and $0$ if $X$ has no good projection.
\end{enumerate}
\end{definition}
\leanref{ReciprocityTheory/ReductiveGroups/ArithmeticGeometry/BirationalGeometry/IsIntersectionFamily.lean}{ReciprocityTheory.ChartRestriction.G.Geo.segre}
\leanref{ReciprocityTheory/ReductiveGroups/ArithmeticGeometry/BirationalGeometry/IsIntersectionFamily.lean}{ReciprocityTheory.ChartRestriction.G.Geo.IsGoodProj}
\leanref{ReciprocityTheory/ReductiveGroups/ArithmeticGeometry/BirationalGeometry/IsIntersectionFamily.lean}{ReciprocityTheory.ChartRestriction.G.Geo.fibre}
\leanref{ReciprocityTheory/ReductiveGroups/ArithmeticGeometry/BirationalGeometry/IsIntersectionFamily.lean}{ReciprocityTheory.ChartRestriction.G.Geo.IsSheet}
\leanref{ReciprocityTheory/ReductiveGroups/ArithmeticGeometry/BirationalGeometry/IsIntersectionFamily.lean}{ReciprocityTheory.ChartRestriction.G.Geo.sheetDensity}
\leanref{ReciprocityTheory/ReductiveGroups/ArithmeticGeometry/BirationalGeometry/IsIntersectionFamily.lean}{ReciprocityTheory.ChartRestriction.G.Geo.ptsIntegralAt}
\leanref{ReciprocityTheory/ReductiveGroups/ArithmeticGeometry/BirationalGeometry/IsIntersectionFamily.lean}{ReciprocityTheory.ChartRestriction.G.Geo.ptsIntegral}

\begin{lemma}\label{cv-origin-intersection-poincare-lelong-all-double-good-normalization}
Let $k$ be an infinite field, let $A$ be a $k$-algebra, let $v_1,\dots,v_n\in A$, and let $m\ge0$. Suppose that for
every $r>m$, every family of $r$ $k$-linear combinations of $v_1,\dots,v_n$ is algebraically dependent over $k$. Then
there are $r\le m$ and $k$-linear combinations $\ell_1,\dots,\ell_r$ of $v_1,\dots,v_n$ such that every $v_t$ is a root
of a monic polynomial whose coefficients are polynomials over $k$ in $\ell_1,\dots,\ell_r$.
\end{lemma}
\leanref{ReciprocityTheory/ReductiveGroups/ArithmeticGeometry/BirationalGeometry/InterGoodE.lean}{ReciprocityTheory.ChartRestriction.G.Geo.exists_linear_normalization}
\begin{proof}
Say that $k$-linear combinations $\ell_1,\dots,\ell_r$ of the $v_t$ \emph{normalize} if every $v_t$ is integral over
$k[\ell_1,\dots,\ell_r]$ in the sense of the statement. The family $v_1,\dots,v_n$ normalizes. If $n\le m$, it gives the
lemma. Otherwise it suffices to show the following: if $\ell_0,\dots,\ell_r$ normalize and $r+1>m$, then there are
$\lambda_1,\dots,\lambda_r\in k$ such that the $\ell'_i=\ell_i-\lambda_i\ell_0$, $i=1,\dots,r$, normalize. Applied
$n-m$ times, starting from $v_1,\dots,v_n$, this gives $m$ linear combinations that normalize.

By hypothesis there is a nonzero $F\in k[T_0,\dots,T_r]$ with $F(\ell_0,\dots,\ell_r)=0$. Let $d$ be its total degree and
$F_d\neq0$ its homogeneous component of degree $d$. There is $\lambda\in k^r$ with $F_d(1,\lambda)\neq0$. Otherwise, for
every $y\in k^{r+1}$ with $y_0\neq0$, $F_d(y)=y_0^dF_d(1,y_1/y_0,\dots,y_r/y_0)=0$; so the polynomial $T_0F_d$ vanishes on
$k^{r+1}$, hence is $0$ as $k$ is infinite, and $F_d=0$.

Substitute $T_0=S$ and $T_i=\lambda_iS+T'_i$ in $F$. This gives a polynomial $G$ in $S$ with coefficients in
$k[T'_1,\dots,T'_r]$, of degree at most $d$. A monomial of degree $d$ in the $T_j$ contributes to the coefficient of
$S^d$ its coefficient times the product of the corresponding powers of $1$ and the $\lambda_i$, and a monomial of lower
degree contributes nothing; so the coefficient of $S^d$ is $F_d(1,\lambda)$. So $G/F_d(1,\lambda)$ is monic in $S$, and,
as $\ell_i=\lambda_i\ell_0+\ell'_i$, its value at $S=\ell_0$ and $T'=\ell'$ is $F(\ell_0,\dots,\ell_r)/F_d(1,\lambda)=0$. So
$\ell_0$ is integral over $k[\ell'_1,\dots,\ell'_r]$, hence so is every $\ell_i=\ell'_i+\lambda_i\ell_0$, and
$k[\ell_0,\dots,\ell_r]$ is integral over $k[\ell'_1,\dots,\ell'_r]$. By the transitivity of integrality, every $v_t$
is integral over $k[\ell'_1,\dots,\ell'_r]$.
\end{proof}
\leanref{ReciprocityTheory/ReductiveGroups/ArithmeticGeometry/BirationalGeometry/InterGoodD.lean}{ReciprocityTheory.ChartRestriction.G.Geo.exists_isIntegralElem_linear}
\leanref{ReciprocityTheory/ReductiveGroups/ArithmeticGeometry/BirationalGeometry/InterGoodE.lean}{ReciprocityTheory.ChartRestriction.G.Geo.intOver_trans}
% Lean: exists_linear_normalization (Good r: some r forms over which every v t is IntOver; hstart with the forms v t,
% hstep by exists_isIntegralElem_linear and intOver_trans, hdesc down to m). exists_isIntegralElem_linear: the top
% homogeneous component of F (homogeneousComponent_totalDegree_ne_zero), exists_eval_cons_one_ne_zero (MvPolynomial.funext
% for H * X 0 and eval_smul_of_degree), subst (T_0 to S, T_i to lam_i S + T'_i) with subst_natDegree_le and
% subst_coeff_totalDegree, eval₂_subst. Size: InterGoodD.lean and InterGoodE.lean, 243 lines.

\begin{lemma}\label{cv-origin-intersection-poincare-lelong-all-double-good}
Let $K$, $(I_k)_{k\in K}$ and $R$ be as in Definition~\ref{lit-multiprojective}, let $e\ge0$, and let $X$ be an
arithmetic subvariety of dimension $e+1$ defined by the prime $P$.
\begin{enumerate}
\item $X$ has a good projection (Definition~\ref{cv-origin-intersection-poincare-lelong-all-double-projection}(1)).
\item For every good projection $\Psi$ for $X$ and every $y\in\mathbf C^e$, the fibre $F_\Psi(y)$ is finite.
\end{enumerate}
\end{lemma}
\leanref{ReciprocityTheory/ReductiveGroups/ArithmeticGeometry/BirationalGeometry/InterSt.lean}{ReciprocityTheory.ChartRestriction.G.Geo.GoodProjExistsStatement}
\leanref{ReciprocityTheory/ReductiveGroups/ArithmeticGeometry/BirationalGeometry/InterSt.lean}{ReciprocityTheory.ChartRestriction.G.Geo.FibreFiniteStatement}
\begin{proof}
(1) Let $A=\mathbf Q[x]/P\mathbf Q[x]$, where $\mathbf Q[x]=R\otimes\mathbf Q$; it is a domain, as $P$ is prime and
$P\cap\mathbf Z=0$. For $i\in I$ let $u_i\in A$ be the class of the Segre monomial $m_i=\prod_kx_{k,i_k}$.

First, if $\ell_1,\dots,\ell_r$ are $\mathbf Q$-linear combinations of the $u_i$ that are algebraically independent over
$\mathbf Q$, then $r\le e+1$. Fix $k_0\in K$ and, as $P$ is relevant, $i_0\in I$ with $x_{k,i_{0,k}}\notin P$ for all $k$,
and let $v_k$ be the class of $x_{k,i_{0,k}}$. The $\ell_j$ and the $v_k$ with $k\neq k_0$ are algebraically
independent. Indeed, a monomial $\prod_j\ell_j^{n_j}\prod_kv_k^{m_k}$ is multihomogeneous of a multidegree $d$ with
$d_{k_0}=\sum_jn_j$ and $m_k=d_k-d_{k_0}$. As $P\mathbf Q[x]$ is multihomogeneous, a polynomial relation splits into
relations of fixed multidegree, each the product of $\prod_kv_k^{m_k}$, which is nonzero in the domain $A$, and a polynomial
in the $\ell_j$, which must therefore vanish. So $r+|K|-1$ is at most the transcendence degree over $\mathbf Q$ of the
fraction field of $R/P$, which is $\dim R/P-1=e+|K|$, as in the proof of Lemma~\ref{cv-origin-weight-birational}
\cite[Tag 02JU]{Sta}. So $r\le e+1$.

By the previous paragraph, Lemma~\ref{cv-origin-intersection-poincare-lelong-all-double-good-normalization} applies to
the $u_i$ with $k=\mathbf Q$ and $m=e+1$. It gives $r\le e+1$ and $\mathbf Q$-linear combinations
$\ell_j=\sum_i\ell_{j,i}u_i$, $j=1,\dots,r$, over which every $u_i$ is integral. Let $\Psi:\mathbf C^I\to\mathbf C^{e+1}$ have the
coordinates $w\mapsto\sum_i\ell_{j,i}w_i$ for $j\le r$ and $0$ for the others. Suppose that $z\in\hat X(\mathbf C)$ and
$\Psi(\sigma(z))=0$. Evaluation at $z$ is a homomorphism $A\to\mathbf C$ that sends $u_i$ to $\sigma(z)_i$ and every
$\ell_j$ to $0$. Each $u_i$ is a root of a monic polynomial with coefficients in $\mathbf Q[\ell_1,\dots,\ell_r]$, so
$\sigma(z)_i$ is a root of the monic polynomial $\bar p_i\in\mathbf Q[T]$ obtained by setting the $\ell_j$ to $0$.
Choose $i$ with $\sigma(z)_i\neq0$, which exists as every block of $z$ is nonzero. For $n\ge0$, multiplying the block
$k_0$ of $z$ by $n+1$ gives a point of $\hat X(\mathbf C)$, as $P$ is multihomogeneous, at which $\Psi\circ\sigma$ still
vanishes and whose Segre coordinate at $i$ is $(n+1)\sigma(z)_i$. These are infinitely many roots of $\bar p_i$, which is
impossible. So $\Psi$ is a good projection.

(2) Let $\psi_j=\Psi_j(\sigma(x))\in\mathbf C[x]$ for $j=0,\dots,e$, multihomogeneous of multidegree $(1,\dots,1)$. At a
common zero in $\mathbf C^{\sqcup_kI_k}$ of $P$ and the $\psi_j$ some block vanishes, as $\Psi$ is good, so every $m_i$
vanishes there. By the Nullstellensatz, $m_i^{n+1}\in P\mathbf C[x]+(\psi_0,\dots,\psi_e)$ for some $n\ge0$. Taking the
component of multidegree $(n+1,\dots,n+1)$ gives $m_i^{n+1}=f+\sum_jc_j\psi_j$ with $f\in P\mathbf C[x]$ and $c_j$
multihomogeneous of multidegree $(n,\dots,n)$, which is a polynomial $G_{ij}$ of degree at most $n$ in the Segre monomials.
So every $w\in F_\Psi(y)$ satisfies $w_i^{n+1}=H_i(w)$, where $H_i=\sum_j(1,y)_jG_{ij}$ has degree at most $n$. For such a
system of equations, one for each $i\in I$, every polynomial function on its set of solutions agrees with a combination
of the finitely many monomials whose exponent at each $i$ is at most $n$ (with the $n$ of $i$), by induction on the degree.
So the powers of each $w_i$ satisfy a nonzero linear relation on the solutions, each $w_i$ takes finitely many values, and
$F_\Psi(y)$ is finite.
\end{proof}
\leanref{ReciprocityTheory/ReductiveGroups/ArithmeticGeometry/BirationalGeometry/InterGood.lean}{ReciprocityTheory.ChartRestriction.G.Geo.goodProjExists}
\leanref{ReciprocityTheory/ReductiveGroups/ArithmeticGeometry/BirationalGeometry/InterFib.lean}{ReciprocityTheory.ChartRestriction.G.Geo.fibreFinite}
% Lean: goodProjExists: forms_not_algebraicIndependent (algebraicIndependent_sumElim with the weights blkWt and sWt, and
% succ_le_ringKrullDim_of_algebraicIndependent, which bounds the number of algebraically independent elements by
% dim R/P - 1), exists_linear_normalization (its own proof of Noether normalization by linear forms, for at most e + 1
% forms that need not be independent), evQ, scaleBlock_mem_zeroLocus and Polynomial.finite_setOf_isRoot. fibreFinite:
% segre_rel (vanishingIdeal_zeroLocus_eq_radical, the Nullstellensatz, and diag_of_homog) and fibSys_finite. Size:
% InterGood and InterGoodA to InterGoodG, 994 lines; InterFib, InterFibA and InterFibB, 453 lines.

\begin{lemma}\label{cv-origin-intersection-poincare-lelong-all-double-sheets}
Let $K$, $(I_k)_{k\in K}$ and $R$ be as in Definition~\ref{lit-multiprojective}, let $e\ge0$, let $X$ be an arithmetic
subvariety of dimension $e+1$ defined by the prime $P$, and let $\Psi$ be a good projection for $X$
(Definition~\ref{cv-origin-intersection-poincare-lelong-all-double-projection}). Let $\bar L_j=(a_j,h_j)$,
$j=1,\dots,e$, and $\bar A=(a,h_A)$ be smooth metrized line bundles with $a\in\mathbf Z_{\ge0}^K$, and let $s\in R\setminus
P$ be multihomogeneous of multidegree $a$. Then for almost every $y\in\mathbf C^e$ the following holds: for any two sheets
$\gamma$ and $\gamma'$ through the same $w$ at $y$,
\[
 \log\frac{|s(\gamma(y))|}{h_A(\gamma(y))}\operatorname{MA}(\log h_1\circ\gamma,\dots,\log h_e\circ\gamma)(y)
 =\log\frac{|s(\gamma'(y))|}{h_A(\gamma'(y))}\operatorname{MA}(\log h_1\circ\gamma',\dots,\log h_e\circ\gamma')(y).
\]
\end{lemma}
\leanref{ReciprocityTheory/ReductiveGroups/ArithmeticGeometry/BirationalGeometry/InterSt.lean}{ReciprocityTheory.ChartRestriction.G.Geo.SheetIndepStatement}
\begin{proof}
A sheet is complex differentiable on a neighbourhood of $y$, hence holomorphic there.

\emph{Sheets with equal Segre images.} Suppose that $\sigma\circ\gamma=\sigma\circ\gamma'$ near $y$. Choose $i_0\in I$
with $\gamma(y)_{k,i_{0,k}}\neq0$ for all $k$. Near $y$ these coordinates of $\gamma$ do not vanish, and the equality of
the Segre images gives $\gamma'=c\cdot\gamma$ blockwise, with $c_k=\gamma'_{k,i_{0,k}}/\gamma_{k,i_{0,k}}$ holomorphic and
nonzero. As $s$ and $h_A$ are homogeneous of the same multidegree, $|s|/h_A$ takes the same value at $\gamma(y)$ and at
$\gamma'(y)$. Near $y$, $\log h_j\circ\gamma'=\log h_j\circ\gamma+\sum_ka_{j,k}\log|c_k|$, and
$\sum_ka_{j,k}\log|c_k/c_k(y)|$ is the real part of the holomorphic function $\sum_ka_{j,k}\log(c_k/c_k(y))$. So the
complex Hessians of $\log h_j\circ\gamma$ and $\log h_j\circ\gamma'$ at $y$ agree (Definition~\ref{cv-period-monge-ampere}),
and so do the two sides.

\emph{Sheets through one point have equal Segre images near almost every $y$.} Let $B\subset\mathbf C^e$ be an open ball,
and let $S_B$ be the set of holomorphic maps $G:B\to\mathbf C^I$ with $G(y')\in F_\Psi(y')$ for all $y'\in B$. Two distinct
elements of $S_B$ differ in some coordinate, a holomorphic function on $B$ that does not vanish identically, so they agree
only on a set of measure zero. The set $S_B$ is finite. Indeed, if it has $m$ elements, then off a set of measure zero in
$B$ their values at $y'$ are $m$ distinct points of $F_\Psi(y')$, so $\{y'\in B:|F_\Psi(y')|<m\}$ has measure zero; if
$S_B$ were infinite this would hold for every $m$, and as the fibres are finite
(Lemma~\ref{cv-origin-intersection-poincare-lelong-all-double-good}(2)), $B$ would have measure zero. Now suppose that
$y$ has sheets $\gamma$ and $\gamma'$ through the same $w$ whose Segre images differ near $y$. Both satisfy the conditions
of a sheet on a ball around $y$; let $B$ be a ball inside it that contains $y$, centred at a point of a fixed countable
dense subset of $\mathbf C^e$ and with rational radius. The restrictions of $\sigma\circ\gamma$ and $\sigma\circ\gamma'$ to
$B$ are distinct elements of $S_B$ that agree at $y$. So $y$ lies in the union, over the countably many such balls $B$ and
the pairs of distinct elements of $S_B$, of the sets where they agree, and this union has measure zero.
\end{proof}
\leanref{ReciprocityTheory/ReductiveGroups/ArithmeticGeometry/BirationalGeometry/InterSheet.lean}{ReciprocityTheory.ChartRestriction.G.Geo.sheetIndep}
\leanref{ReciprocityTheory/ReductiveGroups/ArithmeticGeometry/BirationalGeometry/InterSheet.lean}{ReciprocityTheory.ChartRestriction.G.Geo.sheetIndep_of}
% Lean: sheetIndep = sheetIndep_of osgood fibreFinite: sheet_analyticAt (Osgood's theorem, InterOsgood and
% AnalyticCharts/Osgood), density_eq (eq_bsmul_of_segre_eq, logNorm_bsmul, cxHessian_add_re), and badSet_null
% (secSet, coincide_null with volume_zero_set_of_analyticOnNhd, secSet_finite, a countable dense set and rational
% radii). Size: InterSheet, InterSheetA and InterSheetB, 455 lines, with InterOsgood and AnalyticCharts/Osgood, 308.

\begin{lemma}\label{cv-origin-intersection-poincare-lelong-all-double-bridge-ma}
Let $d\ge0$, let $G$ be a holomorphic map from a neighbourhood of $y\in\mathbf C^d$ to $\mathbf C^d$, with complex
derivative $dG(y)$ at $y$, and let $\varphi_1,\dots,\varphi_d$ be real functions that are twice continuously
differentiable near $G(y)$. Then
\[
 \operatorname{MA}(\varphi_1\circ G,\dots,\varphi_d\circ G)(y)
 =\operatorname{MA}(\varphi_1,\dots,\varphi_d)(G(y))\,|\det dG(y)|^2
\]
(Definition~\ref{cv-period-monge-ampere}(2)).
\end{lemma}
\leanref{ReciprocityTheory/AnalyticForms/PlurisubharmonicFunctions/LintegralImageEqLintegralNormSqDetMul.lean}{ReciprocityTheory.AnForm.A.maDens_comp}
\begin{proof}
This is a self-contained computation. Let $T$ be the matrix of $dG(y)$, with entries
$T_{c,a}=\partial G_c/\partial z_a(y)$. As $G$ is holomorphic, $\partial\bar G_c/\partial z_a=0$ and
$\partial^2G_c/\partial z_a\partial\bar z_b=0$, so the chain rule gives
\[
 \frac{\partial^2(\varphi\circ G)}{\partial z_a\partial\bar z_b}(y)
 =\sum_{c,c'}T_{c,a}\,\frac{\partial^2\varphi}{\partial w_c\partial\bar w_{c'}}(G(y))\,\bar T_{c',b},
\]
that is, $H_{\varphi\circ G}(y)=T^{\mathsf T}H_\varphi(G(y))\bar T$. So for every subset $T'$ of $\{1,\dots,d\}$,
$\det\sum_{k\in T'}H_{\varphi_k\circ G}(y)=|\det T|^2\det\sum_{k\in T'}H_{\varphi_k}(G(y))$, and the formula follows from
the definition of $\operatorname{MA}$.
\end{proof}
% Lean: maDens_comp: hessC_comp (the chain rule for the complex Hessian through jacC) and mdisc_transpose_mul_mul_conj.

\begin{lemma}\label{cv-origin-intersection-poincare-lelong-all-double-bridge-area}
Let $S$ be a compact complex manifold of dimension $d$, with charts $\kappa_j:V_j\to\mathbf C^d$, functions $\chi_j$ and
local potentials $\vartheta_1,\dots,\vartheta_d$ as in Definition~\ref{cv-period-monge-ampere}(3), the potentials being
smooth. For $x\in V_j$ let $M_j(x)$ be the Monge--Amp\`ere density at $\kappa_j(x)$ of the potentials taken near $x$,
composed with $\kappa_j^{-1}$. Let $g$ be a real function on $S$ such that $(gM_j)\circ\kappa_j^{-1}$ is integrable on
every compact subset of $\kappa_j(V_j)$, for every $j$. Let $q:S\to\mathbf C^d$ and $E\subset S$ be such that:
\begin{enumerate}
\item $\kappa_j(E\cap V_j)$ has measure zero for every $j$;
\item for every $x\in S\setminus E$ and every $j$ with $x\in V_j$, the map $q\circ\kappa_j^{-1}$ is holomorphic near
  $\kappa_j(x)$, and its Jacobian determinant $\Delta_j(x)$ at $\kappa_j(x)$ is nonzero;
\item $q^{-1}(y)\setminus E$ is finite for every $y\in\mathbf C^d$.
\end{enumerate}
For $x\notin E$ and $j$ with $x\in V_j$ put $p(x)=M_j(x)/|\Delta_j(x)|^2$. Then $p(x)$ does not depend on $j$, the
function $y\mapsto\sum_{x\in q^{-1}(y)\setminus E}g(x)p(x)$ is integrable on $\mathbf C^d$, and its integral is
\[
 \sum_j\int_{\kappa_j(V_j)}\chi_j(\kappa_j^{-1}(u))\,g(\kappa_j^{-1}(u))\,M_j(\kappa_j^{-1}(u))\,du,
\]
the integral over $S$ of $g$ against the Monge--Amp\`ere measure of $\vartheta_1,\dots,\vartheta_d$.
\end{lemma}
\leanref{ReciprocityTheory/ReductiveGroups/ArithmeticGeometry/BirationalGeometry/InterBridgeA.lean}{ReciprocityTheory.ChartRestriction.G.Geo.pushForward_null}
\leanref{GroundTheory/HeightsAndTheta/Heights/SemipositiveHeights/PtsIntegrable.lean}{GroundTheory.Heights₂.Heights.ForSecondHalf.pushForward_integrable}
\begin{proof}
\emph{The density does not depend on the chart.} Let $x\notin E$ lie in $V_j$ and in $V_k$, and let
$\tau=\kappa_k\circ\kappa_j^{-1}$ near $\kappa_j(x)$, a holomorphic map. By
Lemma~\ref{cv-origin-intersection-poincare-lelong-all-double-bridge-ma}, applied to $\tau$ and to the potentials taken
near $x$ composed with $\kappa_k^{-1}$, $M_j(x)=M_k(x)\,|\det d\tau(\kappa_j(x))|^2$. By the chain rule,
$\Delta_j(x)=\Delta_k(x)\det d\tau(\kappa_j(x))$. So $M_j(x)/|\Delta_j(x)|^2=M_k(x)/|\Delta_k(x)|^2$.

\emph{One chart.} Fix $j$, and put $q_j=q\circ\kappa_j^{-1}$ and $h_j=(\chi_jgM_j)\circ\kappa_j^{-1}$ on
$\kappa_j(V_j)$. As $\chi_j$ is continuous, with support a compact subset of $V_j$, $h_j$ is integrable. Let $O_j$ be the
set of the $u\in\kappa_j(V_j)$ such that $q_j$ is holomorphic near $u$ with nonzero Jacobian determinant at $u$. It is
open, and by (2) the rest of $\kappa_j(V_j)$ lies in $\kappa_j(E\cap V_j)$, which has measure zero by (1). By the
holomorphic inverse function theorem every point of $O_j$ has an open neighbourhood on which $q_j$ is injective; as
$\mathbf C^d$ has a countable basis of open sets, $O_j$ is the union of open sets $O_{j,1},O_{j,2},\dots$ on which $q_j$
is injective. The sets $D_{j,n}=O_{j,n}\setminus(O_{j,1}\cup\dots\cup O_{j,n-1})$ are measurable and disjoint, with union
$O_j$, and $q_j$ is injective on each. The real Jacobian determinant of $q_j$ at $u$ is $|\det dq_j(u)|^2$. Let
$f_{j,n}(y)=h_j(u)/|\det dq_j(u)|^2$ if $y=q_j(u)$ with $u\in D_{j,n}$, and $f_{j,n}(y)=0$ if $y\notin q_j(D_{j,n})$. By the
change of variables formula \cite[263D]{Fre01}, $f_{j,n}$ is integrable, with $\int f_{j,n}=\int_{D_{j,n}}h_j$ and
$\int|f_{j,n}|=\int_{D_{j,n}}|h_j|$. As $\sum_n\int_{D_{j,n}}|h_j|=\int_{O_j}|h_j|$ is finite, $\sum_nf_{j,n}$ converges
almost everywhere to an integrable function, whose integral is $\sum_n\int_{D_{j,n}}h_j=\int_{\kappa_j(V_j)}h_j$.

The set $O_j\cap\kappa_j(E\cap V_j)$ has measure zero and $q_j$ is differentiable on it, so its image $N_j$ under $q_j$
has measure zero \cite[263D(ii)]{Fre01}. Let $y\notin N_j$. The points $u\in O_j$ with $q_j(u)=y$ are the $\kappa_j(x)$
with $x\in V_j\cap q^{-1}(y)\setminus E$: for such a $u$, $\kappa_j^{-1}(u)\notin E$ as $y\notin N_j$, and conversely
$\kappa_j(x)\in O_j$ for $x\in V_j\setminus E$ by (2). They are finitely many, by (3). So
\[
 \sum_nf_{j,n}(y)=\sum_{x\in V_j\cap q^{-1}(y)\setminus E}\frac{h_j(\kappa_j(x))}{|\Delta_j(x)|^2}
 =\sum_{x\in q^{-1}(y)\setminus E}\chi_j(x)g(x)p(x),
\]
as $\chi_j$ vanishes off $V_j$.

\emph{All charts.} For $y$ outside the set $\bigcup_jN_j$ of measure zero, summing over $j$ and using $\sum_j\chi_j=1$
gives $\sum_j\sum_nf_{j,n}(y)=\sum_{x\in q^{-1}(y)\setminus E}g(x)p(x)$. So this function of $y$ is integrable, and its
integral is $\sum_j\int_{\kappa_j(V_j)}h_j$.
\end{proof}
\leanref{ReciprocityTheory/ReductiveGroups/ArithmeticGeometry/BirationalGeometry/InterBridgeG.lean}{ReciprocityTheory.ChartRestriction.G.Geo.area_open}
\leanref{ReciprocityTheory/ReductiveGroups/ArithmeticGeometry/BirationalGeometry/InterBridgeH.lean}{ReciprocityTheory.ChartRestriction.G.Geo.fib_term}
% Lean: pushForward_null = pushForward_null' (InterBridgeI) and pushForward_integrable, for an arbitrary ChartPartition:
% goodT (the text's O_j) with isOpen_goodT and target_diff_goodT_subset, integrableOn_chart, area_open (exists_inj_cover
% through a countable basis, disjointed pieces, piece_change through integral_image_eq_integral_abs_det_fderiv_smul and
% abs_det_restrictScalars, integral_tsum_of_summable_integral_norm and tsum_finsum_disjointed),
% addHaar_image_eq_zero_of_differentiableOn_of_addHaar_eq_zero for N_j, goodT_fibre_eq, and fib_term (locDensChart,
% which is maDens_comp for the change of charts, and the chain rule for the Jacobian), with sum_eq_one of the partition.
% The text defines p(x) in any chart; the Lean's pushDens uses the chart chartIdx A x. Size: InterBridgeG, InterBridgeH
% and InterBridgeI, 462 lines, and PtsIntegrable.lean, 81 lines.

\begin{lemma}\label{cv-origin-intersection-poincare-lelong-all-double-bridge-norm}
Let $K$, $(I_k)_{k\in K}$ and $R$ be as in Definition~\ref{lit-multiprojective}, let $e\ge0$, let $X$ be an arithmetic
subvariety of dimension $e+1$ defined by the prime $P$, let $\Psi$ be a good projection for $X$, with Segre map $\sigma$
(Definition~\ref{cv-origin-intersection-poincare-lelong-all-double-projection}), and let $f\in R\setminus P$ be
multihomogeneous. Then there is a nonzero homogeneous polynomial $c$ in $e+1$ variables with complex coefficients such
that $c(\Psi(\sigma(z)))=0$ for every $z\in\hat X(\mathbf C)$ with $f(z)=0$.
\end{lemma}
\leanref{ReciprocityTheory/ReductiveGroups/ArithmeticGeometry/BirationalGeometry/InterBridgeO.lean}{ReciprocityTheory.ChartRestriction.G.Geo.NormPolyStatement}
\leanref{ReciprocityTheory/ReductiveGroups/ArithmeticGeometry/BirationalGeometry/InterBridgeQ.lean}{ReciprocityTheory.ChartRestriction.G.Geo.normPoly_of}
\begin{proof}
Put $I=\prod_kI_k$.

\emph{Monic relations.} Let $\Lambda$ be the ring of the complex functions on $\hat X(\mathbf C)$ of the form
$z\mapsto G(\sigma(z))$ with $G$ a polynomial on $\mathbf C^I$, an algebra over $\mathbf C[Y_0,\dots,Y_e]$ through
$Y_j\mapsto\Psi(\sigma(z))_j$. As in the proof of
Lemma~\ref{cv-origin-intersection-poincare-lelong-all-double-good}(2), for every $i\in I$ there are $r_i\ge1$ and
polynomials $G_{i,0},\dots,G_{i,e}$ of degree less than $r_i$ with
$\sigma(z)_i^{r_i}=\sum_j\Psi(\sigma(z))_jG_{i,j}(\sigma(z))$ for $z\in\hat X(\mathbf C)$. By induction on the degree,
every monomial in $\sigma(z)$ lies in the $\mathbf C[Y]$-submodule of $\Lambda$ generated by the monomials whose
exponent at each $i$ is less than $r_i$: a monomial with exponent at least $r_i$ at some $i$ is, by the relation for $i$,
a combination with coefficients $\Psi(\sigma(z))_j$ of monomials of lower degree times the $G_{i,j}(\sigma(z))$. So
$\Lambda$ is a finite $\mathbf C[Y]$-module, and its elements are integral over $\mathbf C[Y]$ \cite[Tag 00GK]{Sta}. So
for every polynomial $F$ on $\mathbf C^I$ there are $n\ge0$ and polynomials $c_0,\dots,c_n$ in $e+1$ variables with
$c_n=1$ and $\sum_{l=0}^nc_l(\Psi(\sigma(z)))F(\sigma(z))^l=0$ for every $z\in\hat X(\mathbf C)$.

\emph{Cancelling a factor.} As $P$ is relevant, there is $i^0\in I$ with $x_{k,i^0_k}\notin P$ for every $k$; put
$m=\prod_kx_{k,i^0_k}$, which is not in the prime $P$. Let $u\in R\setminus P$, and let $H$ be a polynomial with complex
coefficients such that $uH$ vanishes on $\hat X(\mathbf C)$. Then $H$ vanishes on $\hat X(\mathbf C)$. Indeed, $umH$
vanishes at every common zero of $P$ in $\mathbf C^{\sqcup_kI_k}$: if every $x_{k,i^0_k}$ is nonzero there, it is a
point of $\hat X(\mathbf C)$, and otherwise $m$ vanishes. By the Nullstellensatz \cite[Tag 00FV]{Sta}, $umH$ lies in
every minimal prime $Q$ over $P\,\mathbf C[x]$. Such a $Q$ meets $R$ in $P$: the ideal $P\,\mathbf Q[x]$ is prime and
meets $R$ in $P$, as $P\cap\mathbf Z=0$; $\mathbf C[x]$ is flat over $\mathbf Q[x]$; and if $Q\cap\mathbf Q[x]$ were
larger than $P\,\mathbf Q[x]$, going down \cite[Tag 00HS]{Sta} would give a prime $Q'\subsetneq Q$ with
$Q'\cap\mathbf Q[x]=P\,\mathbf Q[x]$, which contains $P\,\mathbf C[x]$, contradicting the minimality of $Q$. So
$um\notin Q$ and $H\in Q$, for every such $Q$. Hence $H$ lies in the radical of $P\,\mathbf C[x]$, and it vanishes on
$\hat X(\mathbf C)$.

\emph{The polynomial.} Let $a$ be the multidegree of $f$, put $\delta=\sum_ka_k$ and
$f'=f\prod_kx_{k,i^0_k}^{\delta-a_k}$. Then $f'\notin P$, and $f'$ is multihomogeneous of multidegree
$(\delta,\dots,\delta)$, so $f'(z)=F(\sigma(z))$ for a polynomial $F$ on $\mathbf C^I$, as each monomial of $f'$ is a
product of $\delta$ Segre monomials. Take $n$ and the $c_l$ of the first step for $F$, let $l_0$ be the least $l$ with
$c_l\neq0$, and let $H(z)=\sum_{l\ge l_0}c_l(\Psi(\sigma(z)))f'(z)^{l-l_0}$, a polynomial in $z$. Then $f'^{l_0}H$
vanishes on $\hat X(\mathbf C)$, so $H$ does, by the second step with $u=f'^{l_0}$. If $z\in\hat X(\mathbf C)$ and
$f(z)=0$, then $f'(z)=0$, and $0=H(z)=c_{l_0}(\Psi(\sigma(z)))$. So $c_{l_0}$ vanishes on the set $T$ of the
$\Psi(\sigma(z))$ with $z\in\hat X(\mathbf C)$ and $f(z)=0$. Multiplying one block of such a $z$ by $t\neq0$ gives a
point of $\hat X(\mathbf C)$ at which $f$ vanishes, with Segre image $t\sigma(z)$; so $tT=T$. Write
$c_{l_0}=\sum_\mu c^{(\mu)}$ with $c^{(\mu)}$ homogeneous of degree $\mu$. For $w\in T$,
$\sum_\mu t^\mu c^{(\mu)}(w)=c_{l_0}(tw)=0$ for every $t\neq0$, so $c^{(\mu)}(w)=0$ for every $\mu$. A nonzero $c^{(\mu)}$
is the polynomial $c$ of the statement.
\end{proof}
\leanref{ReciprocityTheory/ReductiveGroups/ArithmeticGeometry/BirationalGeometry/InterBridgeP.lean}{ReciprocityTheory.ChartRestriction.G.Geo.exists_monic_rel}
\leanref{ReciprocityTheory/ReductiveGroups/ArithmeticGeometry/BirationalGeometry/InterBridgeQ.lean}{ReciprocityTheory.ChartRestriction.G.Geo.eval_eq_zero_of_mul}
\leanref{ReciprocityTheory/ReductiveGroups/ArithmeticGeometry/BirationalGeometry/InterPCProof.lean}{ReciprocityTheory.ChartRestriction.G.Geo.pcMinContract}
% Lean: normPoly_of pcMinContract: diag_of_homog for F, exists_monic_rel (segre_rel, the module N spanned by the
% monomials of bounded exponents, and IsIntegral.of_mem_of_fg), Nat.find for l_0, eval_eq_zero_of_mul
% (vanishingIdeal_zeroLocus_eq_radical, Ideal.sInf_minimalPrimes, and pcMinContract = extQ_spec with
% comap_of_mem_minimalPrimes_extC, going down for the flat C[x] over Q[x]), and exists_homog_of_cone. Size: InterBridgeP
% and InterBridgeQ, 334 lines, with parts of InterBridgeO and InterPCProof.

\begin{definition}\label{cv-origin-intersection-poincare-lelong-all-double-bridge-map}
Let $K$, $(I_k)_{k\in K}$ and $R$ be as in Definition~\ref{lit-multiprojective}, let $e\ge0$, let $X$ be an arithmetic
subvariety of dimension $e+1$ defined by the prime $P$, let $J$ be an ideal of $R$, let $(\tilde X,(U_x),(\lambda_x))$
be an analytic resolution of $X$ off $J$ (Definition~\ref{cv-origin-intersection-resolution}), and let $\Psi$ be a good
projection for $X$, with $I$ and the Segre map $\sigma$ as in
Definition~\ref{cv-origin-intersection-poincare-lelong-all-double-projection}.
\begin{enumerate}
\item For $x\in\tilde X$ put $v(x)=\Psi(\sigma(\lambda_x(x)))\in\mathbf C^{e+1}$, with coordinates
  $v_0(x),\dots,v_e(x)$.
\item The \emph{affine projection} of $\tilde X$ along $\Psi$ is the map $Q:\tilde X\to\mathbf C^e$ with
  $Q(x)=(v_1(x)/v_0(x),\dots,v_e(x)/v_0(x))$ if $v_0(x)\neq0$, and $Q(x)=0$ otherwise.
\item For $x$ with $v_0(x)\neq0$ put $W(x)=v_0(x)^{-1}\sigma(\lambda_x(x))\in\mathbf C^I$.
\item The \emph{excluded set} $B$ is the set of the $x\in\tilde X$ such that $\lambda_x(x)\in Z(J)$, or $v_0(x)=0$, or
  for some chart $\kappa$ of $\tilde X$ at $x$ the map $Q\circ\kappa^{-1}$ is not holomorphic near $\kappa(x)$ or has
  Jacobian determinant $0$ at $\kappa(x)$.
\end{enumerate}
\end{definition}
\leanref{ReciprocityTheory/ReductiveGroups/ArithmeticGeometry/BirationalGeometry/InterBridgeD.lean}{ReciprocityTheory.ChartRestriction.G.Geo.brQ}
\leanref{ReciprocityTheory/ReductiveGroups/ArithmeticGeometry/BirationalGeometry/InterBridgeD.lean}{ReciprocityTheory.ChartRestriction.G.Geo.brW}
\leanref{ReciprocityTheory/ReductiveGroups/ArithmeticGeometry/BirationalGeometry/InterBridgeD.lean}{ReciprocityTheory.ChartRestriction.G.Geo.brBad}

\begin{lemma}\label{cv-origin-intersection-poincare-lelong-all-double-bridge-fibres}
Let $X$, $J$, $(\tilde X,(U_x),(\lambda_x))$ and $\Psi$ be as in
Definition~\ref{cv-origin-intersection-poincare-lelong-all-double-bridge-map}, with $v$, $Q$, $W$ and $B$ as there.
\begin{enumerate}
\item Let $x\in\tilde X$. For $z\in U_x$, $v(z)$ is a nonzero multiple of $\Psi(\sigma(\lambda_x(z)))$. If
  $v_0(x)\neq0$ and $\kappa$ is a chart of $\tilde X$ at $x$, then $Q\circ\kappa^{-1}$ is holomorphic near $\kappa(x)$.
\item If $v_0(x)\neq0$, then $W(x)$ lies in the fibre $F_\Psi(Q(x))$.
\item $W$ is injective on the set of the $x$ with $\lambda_x(x)\notin Z(J)$ and $v_0(x)\neq0$, and for every
  $y\in\mathbf C^e$ only finitely many of these $x$ have $Q(x)=y$. In particular $Q^{-1}(y)\setminus B$ is finite.
\item Let $y\in\mathbf C^e$. Suppose that $\Psi(\sigma(z))\neq(1,y)$ for every $z\in\hat X(\mathbf C)\cap Z(J)$, and
  that $Q(x)\neq y$ for every $x\in B$ with $v_0(x)\neq0$. Then $W$ restricts to a bijection from $Q^{-1}(y)\setminus B$
  onto $F_\Psi(y)$.
\end{enumerate}
\end{lemma}
\leanref{ReciprocityTheory/ReductiveGroups/ArithmeticGeometry/BirationalGeometry/InterBridgeE.lean}{ReciprocityTheory.ChartRestriction.G.Geo.brQ_eq_local}
\leanref{ReciprocityTheory/ReductiveGroups/ArithmeticGeometry/BirationalGeometry/InterBridgeD.lean}{ReciprocityTheory.ChartRestriction.G.Geo.brW_mem_fibre}
\leanref{ReciprocityTheory/ReductiveGroups/ArithmeticGeometry/BirationalGeometry/InterBridgeD.lean}{ReciprocityTheory.ChartRestriction.G.Geo.brW_injOn}
\leanref{ReciprocityTheory/ReductiveGroups/ArithmeticGeometry/BirationalGeometry/InterBridgeD.lean}{ReciprocityTheory.ChartRestriction.G.Geo.brFib_finite}
\leanref{ReciprocityTheory/ReductiveGroups/ArithmeticGeometry/BirationalGeometry/InterBridgeD.lean}{ReciprocityTheory.ChartRestriction.G.Geo.exists_of_mem_fibre}
\begin{proof}
(1) By (1) of Definition~\ref{cv-origin-intersection-resolution}, $\lambda_z(z)=(c_k\lambda_x(z)_{k,i})_{k,i}$ for some
$c\in(\mathbf C^\times)^K$, so $\sigma(\lambda_z(z))=(\prod_kc_k)\sigma(\lambda_x(z))$ and
$v(z)=(\prod_kc_k)\Psi(\sigma(\lambda_x(z)))$. The map $\Psi\circ\sigma\circ\lambda_x\circ\kappa^{-1}$ is holomorphic near
$\kappa(x)$, as $\lambda_x$ is, and its coordinate $0$ is $v_0(x)\neq0$ at $\kappa(x)$, hence nonzero nearby. There
$Q\circ\kappa^{-1}$ is the vector of the quotients of its coordinates $1,\dots,e$ by its coordinate $0$.

(2) Fix $k_0\in K$, and let $z$ be $\lambda_x(x)$ with its block $k_0$ multiplied by $v_0(x)^{-1}$. It is a point of
$\hat X(\mathbf C)$, as $P$ is generated by multihomogeneous elements, $\sigma(z)=W(x)$, and
$\Psi(W(x))=v(x)/v_0(x)=(1,Q(x))$.

(3) Let $x$ and $x'$ be two such points with $W(x)=W(x')$. Then $\sigma(\lambda_x(x))=\mu\,\sigma(\lambda_{x'}(x'))$ with
$\mu=v_0(x)/v_0(x')\neq0$. Let $z,z'\in\hat C$ with $\sigma(z)=\mu\sigma(z')$, and choose for each $k$ an index $p_k$ with
$z'_{k,p_k}\neq0$. Then $\sigma(z)_p=\mu\sigma(z')_p\neq0$, so every $z_{k,p_k}$ is nonzero, and comparing the Segre
coordinates at $p$ with $p_k$ replaced by $i$ gives $z_{k,i}=c_kz'_{k,i}$ with $c_k=\mu\prod_{l\neq k}z'_{l,p_l}/z_{l,p_l}$.
So $\lambda_x(x)$ and $\lambda_{x'}(x')$ lie in one orbit of $(\mathbf C^\times)^K$, and $x=x'$ by (3) of
Definition~\ref{cv-origin-intersection-resolution}. By (2), $W$ maps the points of the set with $Q(x)=y$ injectively into
$F_\Psi(y)$, which is finite by Lemma~\ref{cv-origin-intersection-poincare-lelong-all-double-good}(2). The points of
$Q^{-1}(y)\setminus B$ lie in the set.

(4) By (2) and (3), $W$ maps $Q^{-1}(y)\setminus B$ injectively into $F_\Psi(y)$. Let $w\in F_\Psi(y)$, so that
$w=\sigma(z)$ with $z\in\hat X(\mathbf C)$ and $\Psi(w)=(1,y)$. By the first hypothesis $z\notin Z(J)$, so by (2) of
Definition~\ref{cv-origin-intersection-resolution} there are $x\in\tilde X$ and $c\in(\mathbf C^\times)^K$ with
$z=(c_k\lambda_x(x)_{k,i})_{k,i}$. Put $\nu=\prod_kc_k$. Then $\sigma(\lambda_x(x))=\nu^{-1}w$, so $v(x)=\nu^{-1}(1,y)$,
$v_0(x)=\nu^{-1}\neq0$, $Q(x)=y$ and $W(x)=w$. By the second hypothesis $x\notin B$.
\end{proof}
% Lean: brQ_eq_local and analyticAt_brQ (InterBridgeE, through R.compat and segre_smul) for (1); brW_mem_fibre
% (scaling the block k0, mem_mcPts_of_sameOrbit) for (2); brW_injOn with sameOrbit_of_segre (InterBridgeB) and R.inj,
% and brFib_finite and brFib_finite' (InterBridgeL) with fibreFinite for (3); exists_of_mem_fibre (R.surj) and the
% equality hfib in the proof of brFibreSum_ae for (4). Size: InterBridgeB, InterBridgeD and InterBridgeE, about 300 lines.

\begin{lemma}\label{cv-origin-intersection-poincare-lelong-all-double-bridge-null}
Let $X$, $P$, $J$, $(\tilde X,(U_x),(\lambda_x))$ and $\Psi$ be as in
Definition~\ref{cv-origin-intersection-poincare-lelong-all-double-bridge-map}, with $v$, $Q$ and $B$ as there.
\begin{enumerate}
\item For every chart $\kappa:V\to\mathbf C^e$ of $\tilde X$, the set $\kappa(B\cap V)$ has measure zero.
\item The set of the $y\in\mathbf C^e$ with $\Psi(\sigma(z))=(1,y)$ for some $z\in\hat X(\mathbf C)\cap Z(J)$ has
  measure zero.
\item The set of the $Q(x)$ with $x\in B$ and $v_0(x)\neq0$ has measure zero.
\end{enumerate}
\end{lemma}
\leanref{ReciprocityTheory/ReductiveGroups/ArithmeticGeometry/BirationalGeometry/InterBridgeA.lean}{ReciprocityTheory.ChartRestriction.G.Geo.brBad_null}
\leanref{ReciprocityTheory/ReductiveGroups/ArithmeticGeometry/BirationalGeometry/InterBridgeA.lean}{ReciprocityTheory.ChartRestriction.G.Geo.projB_null}
\leanref{ReciprocityTheory/ReductiveGroups/ArithmeticGeometry/BirationalGeometry/InterBridgeA.lean}{ReciprocityTheory.ChartRestriction.G.Geo.brQ_image_null}
\begin{proof}
The zero set of a holomorphic function $\phi$ on a connected open subset of $\mathbf C^e$ that does not vanish
identically has measure zero: $\log|\phi|$ is plurisubharmonic \cite[Chapter~I, (5.12)]{Dem12}, hence subharmonic
\cite[Chapter~I, (5.3)]{Dem12} and not identically $-\infty$, hence locally integrable \cite[Chapter~I,
Theorem~4.17]{Dem12}, and it is $-\infty$ on the zero set. So a set of which every point has an open neighbourhood in
which the set lies in such a zero set has measure zero, as countably many of these neighbourhoods cover it.

\emph{No flat pieces.} Let $\kappa:V\to\mathbf C^e$ be a chart, and let $O\subset\kappa(V)$ be a nonempty open set such
that $v_0$ does not vanish on $\kappa^{-1}(O)$ and $\tilde Q=Q\circ\kappa^{-1}$ is holomorphic on $O$. Then the Jacobian
determinant of $\tilde Q$ does not vanish identically on $O$. Suppose it does. Let $O'$ be the set of the $u\in O$ with
$\lambda_x(x)\notin Z(J)$ for $x=\kappa^{-1}(u)$. It is open: if $g\in J$ does not vanish at $\lambda_x(x)$, then
$g(\lambda_x(z))\neq0$ for $z$ near $x$, and $\lambda_z(z)$ lies in the orbit of $\lambda_x(z)$, so
$\lambda_z(z)\notin Z(J)$. It is nonempty, as $O$ has positive measure and $O\setminus O'$ has measure zero by (6) of
Definition~\ref{cv-origin-intersection-resolution}. Let $r$ be the largest rank of $d\tilde Q$ on $O'$, attained at
$u_1$; then $r<e$. Let $\pi$ be a linear projection of $\mathbf C^e$ onto the image $L$ of $d\tilde Q(u_1)$, let $N$ be the
kernel of $\pi\circ d\tilde Q(u_1)$, of dimension $e-r>0$, and let $\rho$ be a linear projection of $\mathbf C^e$ onto
$N$. The map $\Phi=(\pi\circ\tilde Q,\rho):O'\to L\times N$ has invertible derivative at $u_1$, so by the holomorphic
inverse function theorem it has a holomorphic inverse $\Phi^{-1}$ near $\Phi(u_1)$, with values in $O'$. Near $u_1$,
$\pi\circ d\tilde Q$ maps onto $L$, as at $u_1$, and as $d\tilde Q$ has rank at most $r$, the kernels of $d\tilde Q$ and of
$\pi\circ d\tilde Q$ agree. Differentiating $\pi\circ\tilde Q\circ\Phi^{-1}=\mathrm{pr}_L$ shows that $d\Phi^{-1}$ maps
$\{0\}\times N$ into the kernel of $\pi\circ d\tilde Q$, hence of $d\tilde Q$. So the derivative of $\tilde Q\circ\Phi^{-1}$
vanishes on $\{0\}\times N$ near $\Phi(u_1)$, and $\tilde Q\circ\Phi^{-1}$ is constant on the set of the
$(\pi\tilde Q(u_1),\rho(u_1)+n)$ with $n$ in a small ball of $N$. The image of this infinite set under $\Phi^{-1}$ is an
infinite subset of $O'$ on which $\tilde Q$ is constant, and $\kappa^{-1}$ maps it injectively to points $x$ with
$\lambda_x(x)\notin Z(J)$, $v_0(x)\neq0$ and one value of $Q$. This contradicts
Lemma~\ref{cv-origin-intersection-poincare-lelong-all-double-bridge-fibres}(3).

(1) The set $\kappa(B\cap V)$ is contained in the union of three sets. The first, of the $u$ with $\lambda_x(x)\in Z(J)$
for $x=\kappa^{-1}(u)$, has measure zero by (6) of Definition~\ref{cv-origin-intersection-resolution}.

The second is the set of the $u$ with $v_0(x)\neq0$, for $x=\kappa^{-1}(u)$, such that for some chart $\kappa'$ at $x$
the map $Q\circ\kappa'^{-1}$ is not holomorphic near $\kappa'(x)$ or has Jacobian determinant $0$ at $\kappa'(x)$. By
Lemma~\ref{cv-origin-intersection-poincare-lelong-all-double-bridge-fibres}(1) it is holomorphic there, and by the chain
rule its Jacobian determinant at $\kappa'(x)$ times that of the holomorphic change of charts $\kappa'\circ\kappa^{-1}$,
which is nonzero, is the Jacobian determinant $\delta(u)$ of $Q\circ\kappa^{-1}$ at $u$. So this is the set of the $u$
with $v_0(\kappa^{-1}(u))\neq0$ and $\delta(u)=0$. By
Lemma~\ref{cv-origin-intersection-poincare-lelong-all-double-bridge-fibres}(1), each such $u$ has a ball around it on
which $v_0\circ\kappa^{-1}$ does not vanish and $Q\circ\kappa^{-1}$ is holomorphic; there $\delta$ is holomorphic, and by
the previous paragraph it does not vanish identically. So this set has measure zero.

The third is the set of the $u$ with $v_0(\kappa^{-1}(u))=0$. Let $u_0$ be such a point, $x_0=\kappa^{-1}(u_0)$, and
$D_0$ a ball around $u_0$ in $\kappa(V\cap U_{x_0})$. On $D_0$ the function
$\phi(u)=\Psi(\sigma(\lambda_{x_0}(\kappa^{-1}(u))))_0$ is holomorphic, and by
Lemma~\ref{cv-origin-intersection-poincare-lelong-all-double-bridge-fibres}(1) the set meets $D_0$ in the zero set of
$\phi$. It remains to see that $\phi$ does not vanish identically on $D_0$. Suppose it does. As $\Psi$ is a good
projection and $\lambda_{x_0}(x_0)\in\hat X(\mathbf C)$, some coordinate $i_0$ of $\Psi(\sigma(\lambda_{x_0}(x_0)))$ is
nonzero, and $i_0\neq0$. Let $\Psi'$ be $\Psi$ followed by the exchange of the coordinates $0$ and $i_0$, again a good
projection, with $v'$ and $Q'$ defined for it as $v$ and $Q$ are for $\Psi$. The set $O$ of the $u\in D_0$ with
$\Psi'(\sigma(\lambda_{x_0}(\kappa^{-1}(u))))_0\neq0$ is open and contains $u_0$, and by
Lemma~\ref{cv-origin-intersection-poincare-lelong-all-double-bridge-fibres}(1) for $\Psi'$, $v'_0$ does not vanish on
$\kappa^{-1}(O)$, $Q'\circ\kappa^{-1}$ is holomorphic on $O$, and its coordinate $i_0$ is the quotient of $\phi$ by
$\Psi'(\sigma(\lambda_{x_0}\circ\kappa^{-1}))_0$, which is $0$ on $O$. So the Jacobian determinant of $Q'\circ\kappa^{-1}$
vanishes identically on $O$, which the previous paragraph excludes for $\Psi'$.

(2) As $J$ is generated by multihomogeneous elements and is not contained in $P$, some multihomogeneous $f\in J$ is not in
$P$. Let $c$ be the polynomial of Lemma~\ref{cv-origin-intersection-poincare-lelong-all-double-bridge-norm} for $f$,
homogeneous of degree $\mu$. Then $y\mapsto c(1,y)$ is a nonzero polynomial: otherwise
$c(w)=w_0^\mu c(1,w_1/w_0,\dots,w_e/w_0)=0$ whenever $w_0\neq0$, and $c=0$. Its zero set has measure zero, by the first
paragraph. If $z\in\hat X(\mathbf C)\cap Z(J)$ and $\Psi(\sigma(z))=(1,y)$, then $f(z)=0$, so $c(1,y)=0$.

(3) By (1), $\kappa(B\cap V)$ has measure zero for every chart $\kappa:V\to\mathbf C^e$, and by
Lemma~\ref{cv-origin-intersection-poincare-lelong-all-double-bridge-fibres}(1), $Q\circ\kappa^{-1}$ is differentiable at
the points of $\kappa(B\cap V)$ at which $v_0\circ\kappa^{-1}$ does not vanish. So the image of these points under
$Q\circ\kappa^{-1}$ has measure zero \cite[263D(ii)]{Fre01}. Finitely many charts cover $\tilde X$.
\end{proof}
\leanref{ReciprocityTheory/ReductiveGroups/ArithmeticGeometry/BirationalGeometry/InterBridgeL.lean}{ReciprocityTheory.ChartRestriction.G.Geo.not_flat}
\leanref{ReciprocityTheory/ReductiveGroups/ArithmeticGeometry/BirationalGeometry/InterBridgeK.lean}{ReciprocityTheory.ChartRestriction.G.Geo.exists_infinite_fibre}
% Lean: brBad_null from R.null, psi0_null = volume_projection_zeroSet_null (the exchange Psi' by Equiv.swap and not_flat
% for Psi') and ram_null = ram_null' (det_brQ_chart for the change of charts, analyticAt_det_fderiv, not_flat), with
% volume_eq_zero_of_locally_in_analytic_zero_set (AnalyticCharts/NullZero, the Lean's own proof that the zero set of an
% analytic function not identically zero has measure zero); not_flat with isOpen_chart_notB and exists_infinite_fibre
% (InterBridgeK, the constant rank argument: eventually_surjective_clm and a local inverse of (PL o Q, gL));
% projB_null = projB_null_of normPoly (exists_eval_cons_ne_zero, volume_eval_eq_zero_lemma); brQ_image_null =
% brQ_image_null_of (addHaar_image_eq_zero_of_differentiableOn_of_addHaar_eq_zero). Size: InterBridgeK, InterBridgeL and
% InterBridgeM, 422 lines, with parts of InterBridgeA, InterBridgeE and InterBridgeO.

\begin{lemma}\label{cv-origin-intersection-poincare-lelong-all-double-bridge-sheet}
Let $X$, $P$, $J$, $(\tilde X,(U_x),(\lambda_x))$ and $\Psi$ be as in
Definition~\ref{cv-origin-intersection-poincare-lelong-all-double-bridge-map}, with $v$, $Q$, $W$ and $B$ as there. Let
$\bar L_j=(a_j,h_j)$, $j=1,\dots,e$, and $\bar A=(a,h_A)$ be smooth metrized line bundles with $a\in\mathbf Z_{\ge0}^K$,
and let $s\in R\setminus P$ be multihomogeneous of multidegree $a$. Let $x\in\tilde X\setminus B$, let $\kappa$ be a chart
of $\tilde X$ at $x$, let $M(x)$ be the Monge--Amp\`ere density at $\kappa(x)$ of the functions
$\log h_j\circ\lambda_x\circ\kappa^{-1}$, and let $\Delta(x)$ be the Jacobian determinant of $Q\circ\kappa^{-1}$ at
$\kappa(x)$. Then there is a sheet $\gamma$ through $W(x)$ at $Q(x)$
(Definition~\ref{cv-origin-intersection-poincare-lelong-all-double-projection}(2)) with
\[
 \log\frac{|s(\gamma(Q(x)))|}{h_A(\gamma(Q(x)))}\operatorname{MA}(\log h_1\circ\gamma,\dots,\log h_e\circ\gamma)(Q(x))
 =\log\frac{|s(\lambda_x(x))|}{h_A(\lambda_x(x))}\cdot\frac{M(x)}{|\Delta(x)|^2}.
\]
\end{lemma}
\leanref{ReciprocityTheory/ReductiveGroups/ArithmeticGeometry/BirationalGeometry/InterBridgeF.lean}{ReciprocityTheory.ChartRestriction.G.Geo.exists_sheet_brW}
\begin{proof}
Put $u=\kappa(x)$, $\tilde Q=Q\circ\kappa^{-1}$ and $y_0=Q(x)$. As $x\notin B$, $\tilde Q$ is holomorphic near $u$ with
$\Delta(x)\neq0$, so by the holomorphic inverse function theorem it has a holomorphic inverse $G$ near $y_0$ with
$G(y_0)=u$ and $\det dG(y_0)=\Delta(x)^{-1}$. Near $y_0$ let $\gamma_0=\lambda_x\circ\kappa^{-1}\circ G$, holomorphic with
$\gamma_0(y_0)=\lambda_x(x)$, and $\psi(y')=\Psi(\sigma(\gamma_0(y')))_0$, holomorphic with $\psi(y_0)=v_0(x)\neq0$. Fix
$k_0\in K$, and for $y'$ near $y_0$ let $\gamma(y')$ be $\gamma_0(y')$ with its block $k_0$ multiplied by
$\psi(y')^{-1}$; elsewhere let $\gamma$ be arbitrary.

\emph{It is a sheet.} Near $y_0$, $\gamma$ is holomorphic, and $\gamma(y')\in\hat X(\mathbf C)$, as $\gamma_0(y')$ lies in
$\hat X(\mathbf C)$ and $P$ is generated by multihomogeneous elements. Its Segre image is
$\psi(y')^{-1}\sigma(\gamma_0(y'))$. By Lemma~\ref{cv-origin-intersection-poincare-lelong-all-double-bridge-fibres}(1),
$y'=\tilde Q(G(y'))$ is the vector of the quotients of the coordinates $1,\dots,e$ of $\Psi(\sigma(\gamma_0(y')))$ by its
coordinate $0$, which is $\psi(y')$. So $\Psi(\sigma(\gamma(y')))=(1,y')$. Finally
$\sigma(\gamma(y_0))=v_0(x)^{-1}\sigma(\lambda_x(x))=W(x)$.

\emph{The density.} As $s$ and $h_A$ are homogeneous of the same multidegree $a$, $|s|/h_A$ takes the same value at
$\gamma(y_0)$ and at $\gamma_0(y_0)=\lambda_x(x)$. Near $y_0$, $\log h_j\circ\gamma=\log h_j\circ\gamma_0-a_{j,k_0}\log|\psi|$,
and $\log|\psi/\psi(y_0)|$ is the real part of a holomorphic function. So the complex Hessians of $\log h_j\circ\gamma$
and of $\log h_j\circ\gamma_0$ at $y_0$ agree, and so do their Monge--Amp\`ere densities
(Definition~\ref{cv-period-monge-ampere}(2)). By Lemma~\ref{cv-origin-intersection-poincare-lelong-all-double-bridge-ma},
applied to $G$ and to the functions $\log h_j\circ\lambda_x\circ\kappa^{-1}$, which are smooth near $u$, the density of the
$\log h_j\circ\gamma_0$ at $y_0$ is $M(x)\,|\det dG(y_0)|^2=M(x)/|\Delta(x)|^2$.
\end{proof}
% Lean: exists_sheet_brW, with the chart chartIdx R.A x: brGood, toOpenPartialHomeomorph of Q o chart^-1 (its inverse F.symm,
% the text's G, analytic by analyticAt_symm), gamma0 = lift x o chart^-1 o F.symm, psi0, the block k0 scaled by psi0^-1
% (bsmul), brQ_eq_local for the sheet, logNorm_bsmul, cxHessian_congr and cxHessian_add_re for the pluriharmonic term,
% maDens_comp with det_jacC and the inverse derivative for the density. The text states it in any chart, as the density
% M(x)/|Delta(x)|^2 does not depend on it (Lemma cv-origin-intersection-poincare-lelong-all-double-bridge-area). Size:
% InterBridgeF.lean, 183 lines.

\begin{lemma}\label{cv-origin-intersection-poincare-lelong-all-double-bridge}
Let $K$, $(I_k)_{k\in K}$ and $R$ be as in Definition~\ref{lit-multiprojective}, let $e\ge0$, let $X$ be an arithmetic
subvariety of dimension $e+1$ defined by the prime $P$, let $J$ be an ideal of $R$, and let
$(\tilde X,(U_x),(\lambda_x))$ be an analytic resolution of $X$ off $J$
(Definition~\ref{cv-origin-intersection-resolution}). Let $\Psi$ be a good projection for $X$, let $\bar L_j=(a_j,h_j)$,
$j=1,\dots,e$, and $\bar A=(a,h_A)$ be smooth metrized line bundles with $a\in\mathbf Z_{\ge0}^K$, and let
$s\in R\setminus P$ be multihomogeneous of multidegree $a$. Let $D(y)$ be the sum over $w\in F_\Psi(y)$ of the sheet
densities at $(y,w)$ of $\log(|s|/h_A)$ against $\log h_1,\dots,\log h_e$
(Definition~\ref{cv-origin-intersection-poincare-lelong-all-double-projection}(3)). Then $D$ is integrable on
$\mathbf C^e$, and its integral is the integral over $\tilde X$ of the function
$x\mapsto\log\big(|s(\lambda_x(x))|/h_A(\lambda_x(x))\big)$ against the Monge--Amp\`ere measure of the potentials
$\log h_j\circ\lambda_x$, defined as in Definition~\ref{cv-period-monge-ampere}(3) with this function as a factor of each
integrand.
\end{lemma}
\leanref{ReciprocityTheory/ReductiveGroups/ArithmeticGeometry/BirationalGeometry/InterRes.lean}{ReciprocityTheory.ChartRestriction.G.Geo.ResBridgeStatement}
\leanref{ReciprocityTheory/ReductiveGroups/ArithmeticGeometry/BirationalGeometry/InterBridge.lean}{ReciprocityTheory.ChartRestriction.G.Geo.resBridge}
\leanref{GroundTheory/HeightsAndTheta/Heights/SemipositiveHeights/PtsIntegrable.lean}{GroundTheory.Heights₂.Heights.ForSecondHalf.PtsIntegrable.of_resExists}
\begin{proof}
Let $v$, $Q$, $W$ and $B$ be as in Definition~\ref{cv-origin-intersection-poincare-lelong-all-double-bridge-map}, and
for $x\in\tilde X$ let $g(x)=\log\big(|s(\lambda_x(x))|/h_A(\lambda_x(x))\big)$, with $\log0$ read as $0$, as for the sheet
densities. Give $\tilde X$ a finite atlas of holomorphic charts $\kappa_j:V_j\to\mathbf C^e$ and functions $\chi_j$ as in
Definition~\ref{cv-period-monge-ampere}(3). The potentials $\log h_l\circ\lambda_x$ near $x$ are smooth, and by (1) of
Definition~\ref{cv-origin-intersection-resolution} two of them differ near a point by $\sum_ka_{l,k}\log|c_k|$, with
holomorphic functions $c_k$ without zeros, which is the real part of a holomorphic function.

\emph{Integrability on $\tilde X$.} Let $x_0\in V_j$, and let $D_0$ be a ball around $\kappa_j(x_0)$ in
$\kappa_j(V_j\cap U_{x_0})$. For $u\in D_0$, by (1) of Definition~\ref{cv-origin-intersection-resolution},
$\lambda_{\kappa_j^{-1}(u)}(\kappa_j^{-1}(u))$ lies in the orbit of $\lambda_{x_0}(\kappa_j^{-1}(u))$, and as $s$ and $h_A$
are homogeneous of the same multidegree, $g(\kappa_j^{-1}(u))=\log|F(u)|-\log H(u)$, where
$F(u)=s(\lambda_{x_0}(\kappa_j^{-1}(u)))$ is holomorphic and $H(u)=h_A(\lambda_{x_0}(\kappa_j^{-1}(u)))$ is continuous and
positive. If $F$ vanishes identically on $D_0$, so does $\log|F|$. Otherwise $\log|F|$ is plurisubharmonic
\cite[Chapter~I, (5.12)]{Dem12}, hence subharmonic \cite[Chapter~I, (5.3)]{Dem12}, and not identically $-\infty$ on the
ball $D_0$, so it is integrable near $\kappa_j(x_0)$ \cite[Chapter~I, Theorem~4.17]{Dem12}. So $g\circ\kappa_j^{-1}$ is
locally integrable on $\kappa_j(V_j)$. The Monge--Amp\`ere density $M_j$ of the smooth potentials is continuous, so
$(gM_j)\circ\kappa_j^{-1}$ is integrable on every compact subset of $\kappa_j(V_j)$.

\emph{The change of variables.} Lemma~\ref{cv-origin-intersection-poincare-lelong-all-double-bridge-area} applies to
$\tilde X$, these charts, functions and potentials, $g$, $q=Q$ and $E=B$: its hypothesis (1) is
Lemma~\ref{cv-origin-intersection-poincare-lelong-all-double-bridge-null}(1), (2) holds by the definition of $B$, and (3)
is Lemma~\ref{cv-origin-intersection-poincare-lelong-all-double-bridge-fibres}(3). So, with $p$ as there, the function
$y\mapsto\sum_{x\in Q^{-1}(y)\setminus B}g(x)p(x)$ is integrable on $\mathbf C^e$, and its integral is the integral over
$\tilde X$ of the statement.

\emph{The fibre sums.} By Lemmas~\ref{cv-origin-intersection-poincare-lelong-all-double-bridge-null}(2) and (3)
and~\ref{cv-origin-intersection-poincare-lelong-all-double-sheets}, almost every $y\in\mathbf C^e$ has the following
properties: $\Psi(\sigma(z))\neq(1,y)$ for every $z\in\hat X(\mathbf C)\cap Z(J)$; $Q(x)\neq y$ for every $x\in B$ with
$v_0(x)\neq0$; and any two sheets through one point at $y$ give the same density. Let $y$ be such a point. By
Lemma~\ref{cv-origin-intersection-poincare-lelong-all-double-bridge-fibres}(4), $W$ is a bijection from
$Q^{-1}(y)\setminus B$ onto $F_\Psi(y)$. For $x\in Q^{-1}(y)\setminus B$,
Lemma~\ref{cv-origin-intersection-poincare-lelong-all-double-bridge-sheet} gives a sheet through $W(x)$ at $y$ whose
density is $g(x)p(x)$, and the sheet density at $(y,W(x))$, computed with the chosen sheet, has the same value. So
$D(y)=\sum_{x\in Q^{-1}(y)\setminus B}g(x)p(x)$ for almost every $y$, and the lemma follows from the previous paragraph.
\end{proof}
\leanref{ReciprocityTheory/ReductiveGroups/ArithmeticGeometry/BirationalGeometry/InterBridgeA.lean}{ReciprocityTheory.ChartRestriction.G.Geo.brFibreSum_ae}
\leanref{ReciprocityTheory/ReductiveGroups/ArithmeticGeometry/BirationalGeometry/InterBridgeJ.lean}{ReciprocityTheory.ChartRestriction.G.Geo.brIntegrable_of}
% Lean: resBridge: pushForward_null (Lemma -bridge-area) along brQ R Psi off brBad, with potOf_isPotential, brBad_null
% (Lemma -bridge-null (1)), brGood, brFib_finite (Lemma -bridge-fibres (3)) and brIntegrable = brIntegrable_of
% logAbsLocInt (R.compat and logNorm_bsmul for the local form, logAbsLocInt' in InterBridgeN through isPsh_logmax and
% integrableOn_log_norm_closedBall, the Lean's own proof that log |F| is locally integrable for any analytic F, with
% Real.log 0 = 0 where F vanishes identically, and locDensCont); then integral_congr_ae with brFibreSum_ae (projB_null
% and brQ_image_null, Lemma -bridge-null (2) and (3); sheetIndep, Lemma cv-origin-intersection-poincare-lelong-all-double-sheets;
% exists_of_mem_fibre, brW_mem_fibre and brW_injOn, Lemma -bridge-fibres (4); exists_sheet_brW, Lemma -bridge-sheet).
% The integrability of D is PtsIntegrable.of_resExists: pushForward_integrable with the same inputs and brFibreSum_ae.
% The integrability on the manifold does not use the normal crossings of (5) of Definition
% cv-origin-intersection-resolution. Size: InterBridge and InterBridgeA to InterBridgeQ, 2,432 lines, and
% PtsIntegrable.lean, 81 lines, besides the potential-theory files of ReciprocityTheory/AnalyticForms they use
% (LocDens, NullZero, LintegralImageEqLintegralNormSqDetMul), whose size was not measured.

\begin{lemma}\label{cv-origin-intersection-poincare-lelong-all-double-independent}
Let $K$, $(I_k)_{k\in K}$ and $R$ be as in Definition~\ref{lit-multiprojective}, let $e\ge0$, let $X$ be an arithmetic
subvariety of dimension $e+1$ defined by the prime $P$, let $\bar L_j=(a_j,h_j)$, $j=1,\dots,e$, and $\bar A=(a,h_A)$ be
smooth metrized line bundles with $a\in\mathbf Z_{\ge0}^K$, and let $s\in R\setminus P$ be multihomogeneous of multidegree
$a$. For a good projection $\Psi$ for $X$ let $D_\Psi(y)$ be the sum over $w\in F_\Psi(y)$ of the sheet densities at
$(y,w)$ of $\log(|s|/h_A)$ against $\log h_1,\dots,\log h_e$
(Definition~\ref{cv-origin-intersection-poincare-lelong-all-double-projection}).
\begin{enumerate}
\item For every good projection $\Psi$, $D_\Psi$ is integrable on $\mathbf C^e$.
\item For any two good projections $\Psi$ and $\Psi'$, $D_\Psi$ and $D_{\Psi'}$ have the same integral. So
  $\int_{X(\mathbf C)}\log(|s|/h_A)\,c_1(\bar L_1)\wedge\dots\wedge c_1(\bar L_e)$ is the integral of $D_\Psi$ for every good
  projection $\Psi$.
\end{enumerate}
\end{lemma}
\leanref{ReciprocityTheory/ReductiveGroups/ArithmeticGeometry/BirationalGeometry/InterSt.lean}{ReciprocityTheory.ChartRestriction.G.Geo.PtsIntegrableStatement}
\leanref{ReciprocityTheory/ReductiveGroups/ArithmeticGeometry/BirationalGeometry/InterSt.lean}{ReciprocityTheory.ChartRestriction.G.Geo.PtsIndepStatement}
\begin{proof}
As $P$ is relevant, some $x_{k,i}$ is not in $P$, and the ideal $(x_{k,i})$ is generated by a multihomogeneous element and
is not contained in $P$. By Lemma~\ref{cv-origin-intersection-resolution-exists}, $X$ has an analytic resolution off some
ideal. By Lemma~\ref{cv-origin-intersection-poincare-lelong-all-double-bridge}, for every good projection $\Psi$ the
function $D_\Psi$ is integrable, and its integral is an integral over this resolution that does not involve $\Psi$. The
last statement follows, as $X$ has a good projection by
Lemma~\ref{cv-origin-intersection-poincare-lelong-all-double-good}(1).
\end{proof}
\leanref{Development/AutomorphicSpectrum/Proofs/MetricProofs/PointsIntegrableOfResolutionStatement.lean}{GlobalTheory.CompanionTheory.Ar.Giv.ptsIntegrable_of_resK}
\leanref{ReciprocityTheory/ReductiveGroups/ArithmeticGeometry/BirationalGeometry/InterRed.lean}{ReciprocityTheory.ChartRestriction.G.Geo.ptsIndep_of}
% Lean: ptsIntegrable_of_resK = PtsIntegrable.of_resExists (resolution_exists_of_resolution_statement) and ptsIndep_of
% with resBridge: exists_anLogRes takes J = (X k i) for an X k i outside P; ptsIntegral_eq_at for the last statement.

\begin{lemma}\label{cv-origin-intersection-poincare-lelong-all-double-fibre}
Let $K$, $(I_k)_{k\in K}$ and $R$ be as in Definition~\ref{lit-multiprojective}, let $p$ be a prime number, let $e\ge0$,
and let $X$ be an arithmetic subvariety of dimension $e+1$ defined by a prime $P$ that is twistable at $p$, with twist
$X^{(p)}$ (Lemma~\ref{cv-origin-intersection-poincare-lelong-all-components}(1)). Put $I=\prod_kI_k$, so that the Segre
coordinates of $R^+$ form $\mathbf C^{I\times\{0,1\}}=\mathbf C^I\otimes\mathbf C^2$ and the Segre map of $R^+$ is
$(z,w)\mapsto\sigma(z)\otimes w$. Let $\Psi$ be a good projection for $X^{(p)}$
(Definition~\ref{cv-origin-intersection-poincare-lelong-all-double-projection}), and put $\Psi_\pm(u)=\Psi(u\otimes w_\pm)$
for $u\in\mathbf C^I$, with $w_\pm$ as in Definition~\ref{cv-origin-intersection-poincare-lelong-all-twist}(2).
\begin{enumerate}
\item For $z\in\hat X(\mathbf C)$, $(z,w_\pm)$ lies in $\hat X^{(p)}(\mathbf C)$. The maps $\Psi_+$ and $\Psi_-$ are good
  projections for $X$, and $u\mapsto u\otimes w_\pm$ is injective.
\item For every $y\in\mathbf C^e$, the fibre $F_\Psi(y)$ of $X^{(p)}$ is the union of the images of the fibres
  $F_{\Psi_+}(y)$ and $F_{\Psi_-}(y)$ of $X$ under $u\mapsto u\otimes w_+$ and $u\mapsto u\otimes w_-$, and these images
  are disjoint.
\end{enumerate}
\end{lemma}
\leanref{Development/AutomorphicSpectrum/Proofs/MetricProofs/FibreTwistEqUnionEmbeddingsStatement.lean}{GlobalTheory.T.fibre_twist_eq_union_embeddings}
\begin{proof}
(1) Let $z\in\hat X(\mathbf C)$. The point $(z,w_\pm)$ has nonzero blocks, the elements of $P$ vanish at it, and
$w_{\pm,0}^2+p\,w_{\pm,1}^2=(\pm i\sqrt p)^2+p=0$; so it lies in $\hat X^{(p)}(\mathbf C)$, and
$\Psi_\pm(\sigma(z))=\Psi(\sigma(z)\otimes w_\pm)\neq0$. The coordinate of $u\otimes w_\pm$ at $(i,1)$ is $u_i$, as
$w_{\pm,1}=1$, so $u\mapsto u\otimes w_\pm$ is injective.

(2) If $u=\sigma(z)\in F_{\Psi_\pm}(y)$, then $u\otimes w_\pm$ is the Segre image of $(z,w_\pm)\in\hat X^{(p)}(\mathbf C)$
and lies in $F_\Psi(y)$. Conversely, let $v$ be the Segre image of $z^+\in\hat X^{(p)}(\mathbf C)$ with $\Psi(v)=(1,y)$. The
block $z^+_\ast$ is nonzero and satisfies $(z^+_{\ast,0})^2+p\,(z^+_{\ast,1})^2=0$, so $c=z^+_{\ast,1}\neq0$ and
$z^+_{\ast,0}=\pm i\sqrt p\,c$, that is, $z^+_\ast=c\,w_\pm$ for one sign. The old coordinates of $z^+$ form a point of
$\hat X(\mathbf C)$, as $P\subset P^{(p)}$, and multiplying one of their blocks by $c$ gives $z\in\hat X(\mathbf C)$, as $P$
is multihomogeneous, with $v=\sigma(z)\otimes w_\pm$. Then $\Psi_\pm(\sigma(z))=(1,y)$, and $v$ is the image of
$\sigma(z)\in F_{\Psi_\pm}(y)$. Finally, if $u\otimes w_+=u'\otimes w_-$, the coordinates at $(i,1)$ give $u=u'$ and those
at $(i,0)$ give $2i\sqrt p\,u=0$; so $u=0$, which is impossible for $u$ in a fibre, as $\Psi_+(u)=(1,y)\neq0$.
\end{proof}
\leanref{Development/AutomorphicSpectrum/Proofs/MetricProofs/FibreTwistEqUnionEmbeddings.lean}{GlobalTheory.CompanionTheory.Ar.Tw.Fib.fibre_twist}
% Lean: (1) is glue_mem_mcPts, segre_glue and embL_inj, used as good and v1p, v1m in twist_integral; (2) is
% Tw.Fib.fibre_twist (new_coords, exists_segre_eq_embL with mcPts_smul, embL_image_subset), with the disjointness hdis
% proved in twist_integral. Here u tensor w is embL w u. Size: FibreTwistEqUnionEmbeddings.lean, 113 lines, with about 40
% lines of twist_integral.

\begin{lemma}\label{cv-origin-intersection-poincare-lelong-all-double-density}
Let $p$, $e$, $X$, $P$, $X^{(p)}$, $\Psi$ and $\Psi_\pm$ be as in
Lemma~\ref{cv-origin-intersection-poincare-lelong-all-double-fibre}, and fix one of the two signs. Let
$\bar L'_j=(a'_j,h'_j)$, $j=1,\dots,e$, and $\bar A'=(a',h'_A)$ be smooth metrized line bundles for $R^+$, and let
$\bar L_j=(a_j,h_j)$ and $\bar A=(a,h_A)$ be metrized line bundles for $R$, with $a\in\mathbf Z_{\ge0}^K$ and $a'=(a,0)$,
such that $h'_j(z,w_\pm)=h_j(z)$ and $h'_A(z,w_\pm)=h_A(z)$ for $z\in\hat X(\mathbf C)$. Let $s\in R\setminus P$ be
multihomogeneous of multidegree $a$. Then for almost every $y\in\mathbf C^e$ and every $u\in\mathbf C^I$, the sheet
density for $X^{(p)}$ and $\Psi$ at $(y,u\otimes w_\pm)$ of $\log(|s|/h'_A)$ against $\log h'_1,\dots,\log h'_e$ equals
the sheet density for $X$ and $\Psi_\pm$ at $(y,u)$ of $\log(|s|/h_A)$ against $\log h_1,\dots,\log h_e$
(Definition~\ref{cv-origin-intersection-poincare-lelong-all-double-projection}(3)).
\end{lemma}
\leanref{Development/AutomorphicSpectrum/Proofs/MetricProofs/TwistHeightIntegralDecompositionStatement.lean}{GlobalTheory.T.twist_height_integral_decomposition}
\begin{proof}
By Lemma~\ref{cv-origin-intersection-poincare-lelong-all-components}(1), $X^{(p)}$ is an arithmetic subvariety of
dimension $e+1$ and $s\notin P^{(p)}$. Let $y$ lie outside the set of measure zero of
Lemma~\ref{cv-origin-intersection-poincare-lelong-all-double-sheets} for $X^{(p)}$, $\Psi$, the $\bar L'_j$, $\bar A'$
and $s$, and let $u\in\mathbf C^I$. Write $\rho=w_{\pm,0}$, so that $\rho^2=-p$.

If $\gamma$ is a sheet through $u$ at $y$ for $X$ and $\Psi_\pm$, then $\gamma^+(y')=(\gamma(y'),w_\pm)$ is a sheet through
$u\otimes w_\pm$ at $y$ for $X^{(p)}$ and $\Psi$: it is complex differentiable where $\gamma$ is, its values near $y$ lie in
$\hat X^{(p)}(\mathbf C)$ by Lemma~\ref{cv-origin-intersection-poincare-lelong-all-double-fibre}(1), and its Segre image is
$\sigma(\gamma)\otimes w_\pm$. By the choice of $y$, the density for $X^{(p)}$ may be computed with $\gamma^+$. As $s$
involves only the old coordinates and $\gamma(y')\in\hat X(\mathbf C)$ near $y$, $\log(|s|/h'_A)$ takes the same value at
$\gamma^+(y)$ as $\log(|s|/h_A)$ at $\gamma(y)$, and $\log h'_j\circ\gamma^+=\log h_j\circ\gamma$ near $y$. So the two
densities agree.

If there is no such sheet, there is none through $u\otimes w_\pm$ either, and both densities are $0$. Indeed, let
$\gamma'$ be a sheet through $u\otimes w_\pm$ at $y$ for $X^{(p)}$ and $\Psi$. As $\Psi(u\otimes w_\pm)=(1,y)\neq0$, there
is an $i$ with $u_i\neq0$. Writing $\gamma'^\circ$ for the old coordinates of $\gamma'$, the coordinates at $(i,1)$ and
$(i,0)$ of the Segre image of $\gamma'(y)$ give $\sigma(\gamma'^\circ(y))_i\,\gamma'_{\ast,1}(y)=u_i$ and
$\sigma(\gamma'^\circ(y))_i\,\gamma'_{\ast,0}(y)=\rho\,u_i$, so $\gamma'_{\ast,0}(y)=\rho\,\gamma'_{\ast,1}(y)$. Points of
$\hat X^{(p)}(\mathbf C)$ have a nonzero coordinate $\ast,1$, as their block $\ast$ is nonzero and
$x_{\ast,0}^2+p\,x_{\ast,1}^2$ vanishes there. Near $y$,
$(\gamma'_{\ast,0}-\rho\gamma'_{\ast,1})(\gamma'_{\ast,0}+\rho\gamma'_{\ast,1})=\gamma'^2_{\ast,0}+p\,\gamma'^2_{\ast,1}=0$,
and the second factor is $2\rho\,\gamma'_{\ast,1}(y)\neq0$ at $y$, hence nonzero near $y$; so
$\gamma'_\ast=\gamma'_{\ast,1}w_\pm$ near $y$. Let $\gamma(y')$ be $\gamma'^\circ(y')$ with one block $k_0$ multiplied by
$\gamma'_{\ast,1}(y')$. Near $y$, $\gamma$ is complex differentiable, $\gamma(y')\in\hat X(\mathbf C)$, as $P\subset P^{(p)}$ is
multihomogeneous and $\gamma'_{\ast,1}(y')\neq0$, and $\sigma(\gamma)\otimes w_\pm$ is the Segre image of $\gamma'$. So
$\Psi_\pm(\sigma(\gamma(y')))=(1,y')$, and $\sigma(\gamma(y))=u$ by the injectivity in
Lemma~\ref{cv-origin-intersection-poincare-lelong-all-double-fibre}(1). So $\gamma$ is a sheet through $u$ at $y$ for $X$
and $\Psi_\pm$.
\end{proof}
\leanref{Development/AutomorphicSpectrum/Proofs/MetricProofs/TwistHeightIntegralDecomposition.lean}{GlobalTheory.CompanionTheory.Ar.Tw.SheetD.density_twist}
% Lean: density_twist: sheetIndep for the twist, glue_isSheet (the sheet gamma^+), maDensity_congr_of_eventuallyEq and
% eval_inc, and sheet_of_twist_sheet (hseg, hay, hline, second_ne_zero, the rescaled sheet with mcPts_smul and
% segre_smul, embL_injective). The Lean assumes that the twist is an arithmetic subvariety of dimension e + 1, which the
% text derives from the twistability of P. Size: TwistHeightIntegralDecomposition.lean, 229 lines.

\begin{lemma}\label{cv-origin-intersection-poincare-lelong-all-double-conjugate}
Let $K$, $(I_k)_{k\in K}$ and $R$ be as in Definition~\ref{lit-multiprojective}, let $e\ge0$, let $X$ be an arithmetic
subvariety of dimension $e+1$ defined by the prime $P$, let $\bar L_j=(a_j,h_j)$, $j=1,\dots,e$, and $\bar A=(a,h_A)$ be
smooth metrized line bundles with $a\in\mathbf Z_{\ge0}^K$, and let $s\in R\setminus P$ be multihomogeneous of multidegree
$a$. For a metrized line bundle $\bar B=(b,h)$ let $\bar B^c=(b,h^c)$, with $h^c(z)=h(\bar z)$, be its \emph{conjugate}.
Then
\[
 \int_{X(\mathbf C)}\log\frac{|s|}{h_A^c}\;c_1(\bar L_1^c)\wedge\dots\wedge c_1(\bar L_e^c)
 =\int_{X(\mathbf C)}\log\frac{|s|}{h_A}\;c_1(\bar L_1)\wedge\dots\wedge c_1(\bar L_e)
\]
(Definition~\ref{cv-origin-intersection-poincare-lelong-all-double-projection}(4)).
\end{lemma}
\leanref{Development/AutomorphicSpectrum/Proofs/MetricProofs/PtsIntegralConjLineEqStatement.lean}{GlobalTheory.T.ptsIntegral_conjLine_eq}
\leanref{Development/AutomorphicSpectrum/Proofs/MetricProofs/IsHorizontalTwist.lean}{GlobalTheory.CompanionTheory.Ar.Tw.conjLine}
\begin{proof}
The conjugate bundles are smooth metrized line bundles. Complex conjugation maps $\hat X(\mathbf C)$ to itself, as the
elements of $P$ have integer coefficients, and $\sigma(\bar z)=\overline{\sigma(z)}$. As $s$ has integer coefficients,
$|s(\bar z)|=|s(z)|$, so $\log(|s|/h_A^c)(z)=\log(|s|/h_A)(\bar z)$.

Let $\Psi$ be a good projection for $X$ (Lemma~\ref{cv-origin-intersection-poincare-lelong-all-double-good}(1)), and let
$\Psi'(u)=\overline{\Psi(\bar u)}$, a $\mathbf C$-linear map. It is a good projection for $X$: if $\Psi'(\sigma(z))=0$,
then $\Psi(\sigma(\bar z))=0$ with $\bar z\in\hat X(\mathbf C)$. By
Lemma~\ref{cv-origin-intersection-poincare-lelong-all-double-independent}(2), the left side of the statement is the
integral through $\Psi'$ for the conjugate bundles, and the right side is the integral through $\Psi$ for the given ones.

For $y\in\mathbf C^e$, conjugation $u\mapsto\bar u$ is a bijection from $F_{\Psi'}(y)$ onto $F_\Psi(\bar y)$. If $\gamma$
is a sheet through $u$ at $y$ for $\Psi'$, then $\tilde\gamma(y')=\overline{\gamma(\bar y')}$ is a sheet through $\bar u$
at $\bar y$ for $\Psi$, and conversely; so the two sheet densities vanish together. Otherwise, let $\bar y$ lie outside
the set of measure zero of Lemma~\ref{cv-origin-intersection-poincare-lelong-all-double-sheets} for $\Psi$, the
$\bar L_j$, $\bar A$ and $s$; the conjugate of that set also has measure zero. Then the density for $\Psi$ at
$(\bar y,\bar u)$ may be computed with $\tilde\gamma$. Its first factor is $\log(|s|/h_A)(\overline{\gamma(y)})
=\log(|s|/h^c_A)(\gamma(y))$. For the second, $\log h_j\circ\tilde\gamma$ is $y'\mapsto\varphi_j(\bar y')$ with
$\varphi_j=\log h_j^c\circ\gamma$. The complex Hessian of $y'\mapsto\varphi_j(\bar y')$ at $\bar y$ is the complex
conjugate of that of $\varphi_j$ at $y$, and the determinants in Definition~\ref{cv-period-monge-ampere}(2) of the
conjugate Hermitian matrices are the conjugates of real numbers; so the Monge--Amp\`ere densities agree. Hence the sum of
the densities over $F_\Psi(\bar y)$ for the given bundles equals the sum over $F_{\Psi'}(y)$ for the conjugate bundles,
for almost every $y$. As $y\mapsto\bar y$ preserves Lebesgue measure, the two integrals agree.
\end{proof}
\leanref{Development/AutomorphicSpectrum/Proofs/MetricProofs/PtsIntegralConjLineEq.lean}{GlobalTheory.CompanionTheory.Ar.Tw.ptsIntegral_conjLine}
\leanref{Development/WeilBound/CoreStatements/HomogeneousNonvanishing/CArint.lean}{PositivityTheory.UnitarySelfGeneration.c_arint}
\leanref{Development/WeilBound/CoreStatements/HomogeneousNonvanishing/CPtsat.lean}{PositivityTheory.UnitarySelfGeneration.c_ptsat}
\leanref{Development/WeilBound/CoreStatements/HomogeneousNonvanishing/CMadens.lean}{PositivityTheory.UnitarySelfGeneration.c_madens}
% Lean: ptsIntegral_conjLine = Cmp.arInt_conj = c_arint with sheetIndep, PtsIndep and goodProjExists: Psi' is the
% conjugated projection, c_ptsat (ptsat_mem, ptsat_fibre, ptsat_sheet, the measure-preserving star on C^e and
% finsum_mem_eq_of_bijOn), complex_conjugation_herm_compat (logNorm of conjLine) and c_madens (madens_cx: the complex
% Hessian of phi o conj is the conjugate matrix; the mixed discriminant of conjugate matrices is the conjugate). Size:
% PtsIntegralConjLineEq, ArIntConj, CArint, CPtsat and CMadens, about 240 lines.

\begin{lemma}\label{cv-origin-intersection-poincare-lelong-all-double}
Let $K$, $(I_k)_{k\in K}$ and $R$ be as in Definition~\ref{lit-multiprojective}, let $p$ be a prime number, let $e\ge0$,
and let $X$ be an arithmetic subvariety of dimension $e+1$ defined by a prime $P$ that is twistable at $p$, with twist
$X^{(p)}$ (Lemma~\ref{cv-origin-intersection-poincare-lelong-all-components}(1)). Let $\bar L_1,\dots,\bar L_e$ and
$\bar A=(a,h_A)$ be smooth metrized line bundles with $a\in\mathbf Z_{\ge0}^K$, and let $s\in R\setminus P$ be
multihomogeneous of multidegree $a$. Then, with the twists of
Definition~\ref{cv-origin-intersection-poincare-lelong-all-twist}(3),
\[
 \int_{X^{(p)}(\mathbf C)}\log\frac{|s|}{h_A^{(p)}}\;c_1(\bar L_1^{(p)})\wedge\dots\wedge c_1(\bar L_e^{(p)})
 =2\int_{X(\mathbf C)}\log\frac{|s|}{h_A}\;c_1(\bar L_1)\wedge\dots\wedge c_1(\bar L_e).
\]
\end{lemma}
\leanref{Development/AutomorphicSpectrum/Proofs/MetricProofs/PoincareLelongOfInv.lean}{GlobalTheory.CompanionTheory.Ar.Tw.Main.twist_double}
\begin{proof}
Write $\Phi=c_1(\bar L_1)\wedge\dots\wedge c_1(\bar L_e)$, $\Phi^{(p)}$ for the product for the twisted bundles, and
$\Phi^c$ for the product for the conjugate bundles $\bar L_j^c$ of
Lemma~\ref{cv-origin-intersection-poincare-lelong-all-double-conjugate}.

By Lemma~\ref{cv-origin-intersection-poincare-lelong-all-components}(1), $X^{(p)}$ is an arithmetic subvariety of
dimension $e+1$, and $s\notin P^{(p)}$. By Lemma~\ref{cv-origin-intersection-poincare-lelong-all-metric}(1) the twisted
bundles are smooth, and $\bar A^{(p)}$ has multidegree $(a,0)$, that of $s$ in $R^+$. Let $\Psi$ be a good projection for
$X^{(p)}$ (Lemma~\ref{cv-origin-intersection-poincare-lelong-all-double-good}(1)), and let $\Psi_+$ and $\Psi_-$ be the
good projections for $X$ of Lemma~\ref{cv-origin-intersection-poincare-lelong-all-double-fibre}(1). Let $D(y)$ be the sum
over $F_\Psi(y)$ of the sheet densities for $X^{(p)}$ of $\log(|s|/h_A^{(p)})$ against the $\log h_j^{(p)}$, let $D_+(y)$
be the sum over $F_{\Psi_+}(y)$ of those for $X$ of $\log(|s|/h_A)$ against the $\log h_j$, and let $D_-(y)$ be the sum over
$F_{\Psi_-}(y)$ of those for $X$ of $\log(|s|/h_A^c)$ against the $\log h_j^c$. By
Lemma~\ref{cv-origin-intersection-poincare-lelong-all-double-independent}, these functions are integrable, and their
integrals are $\int_{X^{(p)}(\mathbf C)}\log(|s|/h_A^{(p)})\,\Phi^{(p)}$, $\int_{X(\mathbf C)}\log(|s|/h_A)\,\Phi$ and
$\int_{X(\mathbf C)}\log(|s|/h_A^c)\,\Phi^c$.

By Lemma~\ref{cv-origin-intersection-poincare-lelong-all-metric}(2), each twisted metric satisfies
$h^{(p)}(z,w_+)=h(z)$ and $h^{(p)}(z,w_-)=h(\bar z)=h^c(z)$. So
Lemma~\ref{cv-origin-intersection-poincare-lelong-all-double-density} applies with the sign $+$ to the $\bar L_j$ and
$\bar A$, and with the sign $-$ to the $\bar L^c_j$ and $\bar A^c$. By
Lemma~\ref{cv-origin-intersection-poincare-lelong-all-double-fibre}(2), $F_\Psi(y)$ is the disjoint union of the images
of $F_{\Psi_+}(y)$ and $F_{\Psi_-}(y)$ under the injective maps $u\mapsto u\otimes w_\pm$, and these fibres are finite by
Lemma~\ref{cv-origin-intersection-poincare-lelong-all-double-good}(2). So $D(y)=D_+(y)+D_-(y)$ for almost every $y$, and
\[
 \int_{X^{(p)}(\mathbf C)}\log\frac{|s|}{h_A^{(p)}}\;\Phi^{(p)}
 =\int_{X(\mathbf C)}\log\frac{|s|}{h_A}\;\Phi+\int_{X(\mathbf C)}\log\frac{|s|}{h_A^c}\;\Phi^c.
\]
By Lemma~\ref{cv-origin-intersection-poincare-lelong-all-double-conjugate}, the two terms on the right are equal.
\end{proof}
\leanref{Development/AutomorphicSpectrum/Proofs/MetricProofs/PoincareLelongOfInv.lean}{GlobalTheory.CompanionTheory.Ar.Tw.Main.twist_integral}
% Lean: twist_double = twist_integral and GlobalTheory.T.ptsIntegral_conjLine_eq (Tw.ptsIntegral_conjLine), under
% PtsIntegrableStatement (Giv.ptsIntegrable_of_resK) and PtsIndepStatement (ptsIndep_of with
% resolution_exists_of_resolution_statement and resBridge). twist_integral: goodProjExists for the twist, the projections
% pj (vPlus p) and pj (vMinus p) are good for X (glue_mem_mcPts, segre_glue), atOf (PtsIndep: the integral does not
% depend on the good projection), GlobalTheory.T.fibre_twist_eq_union_embeddings (Tw.Fib.fibre_twist, the fibre is the
% disjoint union of the two embedded fibres), GlobalTheory.T.twist_height_integral_decomposition (Tw.SheetD.density_twist,
% the sheet densities agree), and integral_add with PtsIntegrable. The estimate of about 7,000 of the 8,700 lines of
% Lemma cv-origin-intersection-poincare-lelong-all was made from file sizes; the measured files are about 5,400 lines
% (twist files about 480, good projections and fibres 1,447, sheets 763, bridge 2,513, conjugation 240), besides
% unmeasured lemmas of ReciprocityTheory/AnalyticForms. What remains open is the bridge, Lemma
% cv-origin-intersection-poincare-lelong-all-double-bridge.

\begin{lemma}\label{cv-origin-intersection-poincare-lelong-all}
Let $K$, $(I_k)_{k\in K}$ and $R$ be as in Definition~\ref{lit-multiprojective}, let $n\ge0$, let $X$ be an arithmetic
subvariety of dimension $n+2$ defined by the prime $P$, and let $\bar L_1,\dots,\bar L_n$, $\bar A=(a,h_A)$ and
$\bar B=(0,h_B)$ be smooth metrized line bundles, with $a\in\mathbf Z_{\ge0}^K$. Let $s\in R\setminus P$ be
multihomogeneous of multidegree $a$, put $\Phi=c_1(\bar L_1)\wedge\dots\wedge c_1(\bar L_n)$, and for an arithmetic
subvariety $Y$ write $\int_{Y(\mathbf C)}$ for the integral over the complex points of $Y$. Then
\[
 \int_{X(\mathbf C)}\log\frac{|s|}{h_A}\;c_1(\bar B)\wedge\Phi
 =\sum_Wm_W\int_{W(\mathbf C)}\log h_B\;\Phi-\int_{X(\mathbf C)}\log h_B\;c_1(\bar A)\wedge\Phi,
\]
the sum over the horizontal relevant minimal primes $W$ of $P+(s)$, where $m_W$ is the multiplicity of $W$ in
$R/(P+(s))$.
\end{lemma}
\leanref{ReciprocityTheory/ReductiveGroups/ArithmeticGeometry/BirationalGeometry/InterSt.lean}{ReciprocityTheory.ChartRestriction.G.Geo.ArithmeticPoincareLelongFormula}
\leanref{Development/AutomorphicSpectrum/PublishedTheorems/Metrics₂/PoincareLelong.lean}{GlobalTheory.CompanionTheory.Ar.Pub.PoincareLelong}
\leanref{Development/AutomorphicSpectrum/Proofs/MetricProofs/PoincareLelongOfInv.lean}{GlobalTheory.CompanionTheory.Ar.Tw.Main.poincareLelong_of_inv}
\begin{proof}
Let $\mathcal H$ be the set of horizontal relevant minimal primes $W$ of $P+(s)$, a finite set of arithmetic
subvarieties of dimension $n+1$. By Lemma~\ref{cv-origin-intersection-poincare-lelong-all-prime} there is a prime number
$p$ at which $P$ and every $W\in\mathcal H$ are twistable. Use the twists of
Definition~\ref{cv-origin-intersection-poincare-lelong-all-twist} for this $p$, and put
$\Phi^{(p)}=c_1(\bar L_1^{(p)})\wedge\dots\wedge c_1(\bar L_n^{(p)})$. By
Lemma~\ref{cv-origin-intersection-poincare-lelong-all-components}(1), $P^{(p)}$ is an arithmetic subvariety
$X^{(p)}$ of dimension $n+2$, and $s\notin P^{(p)}$. By Lemma~\ref{cv-origin-intersection-poincare-lelong-all-metric}(1)
the twisted bundles are smooth and invariant under conjugation; $\bar A^{(p)}$ has multidegree $a^+$, which is that of
$s$ in $R^+$, and $\bar B^{(p)}$ has multidegree $0$. Theorem~\ref{cv-origin-intersection-poincare-lelong} for
$X^{(p)}$, these bundles and $s$ gives
\[
 \int_{X^{(p)}(\mathbf C)}\log\frac{|s|}{h_A^{(p)}}\;c_1(\bar B^{(p)})\wedge\Phi^{(p)}
 =\sum_{W'}m_{W'}\int_{W'(\mathbf C)}\log h_B^{(p)}\;\Phi^{(p)}
 -\int_{X^{(p)}(\mathbf C)}\log h_B^{(p)}\;c_1(\bar A^{(p)})\wedge\Phi^{(p)},
\]
the sum over the horizontal relevant minimal primes $W'$ of $P^{(p)}+(s)$.

Each term is twice the corresponding one on $X$, by Lemma~\ref{cv-origin-intersection-poincare-lelong-all-double}. For
the left side, apply it to $X$, the bundles $\bar B,\bar L_1,\dots,\bar L_n$ and the section $s$ of $\bar A$. For the
last term, apply it to $X$, the bundles $\bar A,\bar L_1,\dots,\bar L_n$ and the section $1$ of $(0,1/h_B)$, whose twist
is $(0,1/h_B^{(p)})$ by Lemma~\ref{cv-origin-intersection-poincare-lelong-all-metric}(3); here
$\log(|1|\cdot h_B)=\log h_B$. By Lemma~\ref{cv-origin-intersection-poincare-lelong-all-components}(2), the $W'$ are
the twists $W^{(p)}$ of the $W\in\mathcal H$, with $m_{W^{(p)}}=m_W$. Lemma~\ref{cv-origin-intersection-poincare-lelong-all-double}
for $W$, of dimension $n+1$, the bundles $\bar L_1,\dots,\bar L_n$ and the section $1$ of $(0,1/h_B)$ gives
$\int_{W^{(p)}(\mathbf C)}\log h_B^{(p)}\,\Phi^{(p)}=2\int_{W(\mathbf C)}\log h_B\,\Phi$. Dividing by $2$ gives the
formula.
\end{proof}
% Lean: GlobalTheory.CompanionTheory.Ar.Tw.Main.poincareLelong_of_inv, under ResKStatement (Lemma
% cv-resolution-coordinates) and PoincareLelongInv (Theorem cv-origin-intersection-poincare-lelong): PtsIntegrable
% from Giv.ptsIntegrable_of_resK and PtsIndep from resolution_exists_of_resolution_statement and resBridge;
% exists_prime_noRoot chooses p with NoRoot p for P and the horizontal W (no a, b with b outside W and
% a^2 + p b^2 in W); twist p (the ideal W.map inc + (quad p), quad p = xN^2 + p yN^2 in a new block of two
% coordinates), twLine p (the metric interpolating, through chi p, between h and its conjugate at the two points of
% the conic), twLine_smooth_conjInvariant_glue; hInv applied to the twist; twist_double for the three integrals, each
% twice the one on X; horizontal_divSupp_twist_eq_image and multAt for the components (hsup: P^(p) + (s) is
% (P + (s)).map inc + (quad p)); e1, e3 and e2 are the three applications of twist_double, with invLine B, the section 1
% (isSectionOn_inv_one, B.a = 0) and logNorm_inv_one, logNorm_tw_inv_one. Size: about 8,700 lines, besides the 9,500 of
% Lemma cv-origin-intersection-resolution-exists; what remains open is Lemma
% cv-origin-intersection-poincare-lelong-all-double-bridge.

\subsection{Arithmetic intersection systems}\label{subsection-cv-arithmetic-intersection-systems}

\begin{definition}\label{cv-origin-intersection-family-numbers}
Let $K$, $(I_k)_{k\in K}$ and $R$ be as in Definition~\ref{lit-multiprojective}.
\begin{enumerate}
\item Let $I$ assign a real number $I(\bar L_1\cdots\bar L_n\mid W)$ to every arithmetic subvariety $W$ of some
  dimension $n$ and all smooth metrized line bundles $\bar L_1,\dots,\bar L_n$. Let $X$ be an arithmetic subvariety of
  dimension $n+1$ defined by $P$, let $\bar L_1,\dots,\bar L_n$ and $\bar B=(b,h_B)$ be smooth metrized line bundles, and
  let $s\in R\setminus P$ be multihomogeneous of multidegree $b$. Then $\mathcal R_I(\bar L_1\cdots\bar L_n;\bar B,s\mid X)$
  is the right side of the recursion of Definition~\ref{lit-multiprojective}(3): the sum over the relevant minimal
  primes $W$ of $P+(s)$, with their multiplicities, of $I(\bar L_1\cdots\bar L_n\mid W)$ for horizontal $W$ and of
  $\deg_{L_1\cdots L_n}(W)\log p$ for $W\cap\mathbf Z=p\mathbf Z$, minus the integral of $\log(|s|/h_B)$ over
  $\hat X(\mathbf C)$ against $c_1(\bar L_1)\wedge\dots\wedge c_1(\bar L_n)$.
\item For every prime $P$ and every $b\in\mathbf N^K$ fix a multihomogeneous $s_b(P)$ of multidegree $b$, not in $P$
  when $P$ is relevant. The numbers $(\bar L_1\cdots\bar L_n\mid X)_0$ are defined by recursion on $n$: they are $0$ for
  $n=0$, and for $X$ of dimension $n+1$ defined by $P$ and $\bar A=(a,h)$, with $m=\sum_k|a_k|$ and $\mathbf 1$ the
  multidegree that is $1$ everywhere,
  \[
   (\bar L_1\cdots\bar L_n\bar A\mid X)_0=\mathcal R_0\big(\bar L_1\cdots\bar L_n;\bar A\otimes\bar{\mathcal O}(m\mathbf 1),
   s_{a+m\mathbf 1}(P)\mid X\big)-\mathcal R_0\big(\bar L_1\cdots\bar L_n;\bar{\mathcal O}(m\mathbf 1),s_{m\mathbf 1}(P)\mid X\big),
  \]
  where $\mathcal R_0$ is $\mathcal R_I$ for the numbers $(\cdot\mid\cdot)_0$ in dimension $n$ and
  $\bar{\mathcal O}(m\mathbf 1)$ is the Fubini--Study bundle (Definition~\ref{cv-origin-analytic-system-pullback}(1)).
\end{enumerate}
\end{definition}
\leanref{ReciprocityTheory/ReductiveGroups/ArithmeticGeometry/BirationalGeometry/IsIntersectionFamily.lean}{ReciprocityTheory.ChartRestriction.G.Geo.recR}
\leanref{GroundTheory/HeightsAndTheta/Heights/Intersections/ThmIExistsFamily.lean}{GroundTheory.Heights₂.Heights.Inter.ThmIExists.step}
\leanref{GroundTheory/HeightsAndTheta/Heights/Intersections/ThmIExistsFamily.lean}{GroundTheory.Heights₂.Heights.Inter.ThmIExists.fam}
\leanref{GroundTheory/HeightsAndTheta/Heights/Intersections/ThmIExistsFamily.lean}{GroundTheory.Heights₂.Heights.Inter.ThmIExists.famI}

\begin{lemma}\label{cv-origin-intersection-family-curves-points}
Let $K$, $(I_k)_{k\in K}$ and $R$ be as in Definition~\ref{lit-multiprojective}, let $X$ be an arithmetic subvariety of
dimension $1$ defined by $P$, and let $F$ be the subring of the fraction field of $R/P$ formed by the quotients $u/v$ of
multihomogeneous $u,v\in R$ of the same multidegree with $v\notin P$.
\begin{enumerate}
\item There is no arithmetic subvariety of dimension $0$. For multihomogeneous $s\in R\setminus P$, no relevant minimal
  prime of $P+(s)$ is horizontal, and $s(z)\neq0$ for every $z\in\hat X(\mathbf C)$.
\item There is a finite set $T\subset\hat X(\mathbf C)$ that meets every orbit of $(\mathbf C^\times)^K$ on
  $\hat X(\mathbf C)$ in one point, such that $\int_{X(\mathbf C)}g=\sum_{r\in T}g(r)$ for every function $g$ on
  $\mathbf C^{\sqcup_kI_k}$ that is constant on these orbits, the integral being that of
  Definition~\ref{cv-origin-intersection-poincare-lelong-all-double-projection}(4) with no line bundles.
\item $F$ is a field, and for $r\in T$, $\sigma_r(u/v)=u(r)/v(r)$ defines a field embedding $\sigma_r:F\to\mathbf C$. The
  map $r\mapsto\sigma_r$ is a bijection from $T$ onto the set of field embeddings $F\to\mathbf C$.
\end{enumerate}
\end{lemma}
\leanref{GroundTheory/HeightsAndTheta/Heights/Intersections/ArithVarZero.lean}{GroundTheory.Heights₂.Heights.Inter.arithVarZero}
\leanref{GroundTheory/HeightsAndTheta/Heights/Intersections/NotIsHorizontalOfDimOne.lean}{GroundTheory.Heights₂.Heights.Inter.eval_ne_zero_of_dimOne}
\leanref{GroundTheory/HeightsAndTheta/Heights/Intersections/ExistsIsSheetFinZero.lean}{GroundTheory.Heights₂.Heights.Inter.regularityPointsSum_holds.proof}
\leanref{GroundTheory/HeightsAndTheta/Heights/Intersections/NotMemOfNonZeroDivisor.lean}{GroundTheory.Heights₂.Heights.Inter.regularityEmbedding_holds.proof}
\begin{proof}
(1) Let $W$ be a relevant multihomogeneous prime with $W\cap\mathbf Z=0$, and number the blocks $k_1,\dots,k_m$, $m=|K|$.
For $0\le l\le m$ let $Q_l$ be the set of the $f\in R$ such that $f$, with the variables of the blocks
$k_1,\dots,k_l$ set to $0$, lies in $W$. It is a prime ideal, the inverse image of $W$ under a ring homomorphism, and
$Q_0=W$. As $W$ is generated by multihomogeneous elements, it contains the part of each of its elements of degree $0$ in a
given block, so $Q_l\subset Q_{l+1}$; and some $x_{k_{l+1},i}$ is not in $W$, so it lies in $Q_{l+1}\setminus Q_l$. Setting
every variable to $0$ maps $f$ to its constant term, so $Q_m$ is the set of the $f$ with constant term in
$W\cap\mathbf Z=0$, which is strictly contained in the prime ideal of the $f$ with even constant term. This chain of
$m+2$ primes starting at $W$ gives $\dim R/W\ge m+1$, so $W$ is not an arithmetic subvariety of dimension $0$. A horizontal
relevant minimal prime of $P+(s)$ would be one, as for the components in Definition~\ref{lit-multiprojective}(3). If
$s(z)=0$ for some $z\in\hat X(\mathbf C)$, the kernel of $f\mapsto f(z)$ on $R$ is a prime containing $P+(s)$, so it contains
a minimal prime $W$ of $P+(s)$; $W\cap\mathbf Z=0$ as the kernel meets $\mathbf Z$ in $0$, and $W$ is relevant as every
block of $z$ is nonzero. This contradicts what was just shown.

(2) Let $\Psi:\mathbf C^I\to\mathbf C$ be the good projection chosen for $X$, here with $e=0$, so that
$\Psi(\sigma(z))\neq0$ for $z\in\hat X(\mathbf C)$. The fibre $F_\Psi$ over the point $\mathbf C^0$ is the set of the
$\sigma(z)$ with $z\in\hat X(\mathbf C)$ and $\Psi(\sigma(z))=1$; it is finite by
Lemma~\ref{cv-origin-intersection-poincare-lelong-all-double-good}(2). A sheet through $w\in F_\Psi$ is a point
$\gamma_w\in\hat X(\mathbf C)$ with $\sigma(\gamma_w)=w$, which exists, and as the Monge--Amp\`ere density of no functions
is $1$ (Definition~\ref{cv-period-monge-ampere}(2)), the sheet density of $g$ at $w$ is $g(\gamma_w)$, for the chosen
$\gamma_w$. Let $T$ be the set of the chosen $\gamma_w$; the map $w\mapsto\gamma_w$ is injective. As $\mathbf C^0$ is a
point of measure $1$, $\int_{X(\mathbf C)}g=\sum_{w\in F_\Psi}g(\gamma_w)=\sum_{r\in T}g(r)$. Let $z\in\hat X(\mathbf C)$ and
$a=\Psi(\sigma(z))\neq0$. Multiplying one block of $z$ by $a^{-1}$ gives a point of $\hat X(\mathbf C)$, as $P$ is
multihomogeneous, whose Segre image $w=a^{-1}\sigma(z)$ lies in $F_\Psi$. Then $\gamma_w$ and $z$ have proportional Segre
images, so they lie in one orbit, as in the proof of
Lemma~\ref{cv-origin-intersection-poincare-lelong-all-double-bridge-fibres}(3). Conversely, if $z=c\gamma_{w'}$ with
$c\in(\mathbf C^\times)^K$ acting blockwise, then $\sigma(z)=(\prod_kc_k)w'$ and $a=\prod_kc_k$, as $\Psi(w')=1$; so
$w'=a^{-1}\sigma(z)=w$. So $T$ meets the orbit of $z$ only in $\gamma_w$.

(3) By (1), $v(r)\neq0$ for $r\in T$ and multihomogeneous $v\notin P$. If $u/v=u'/v'$, then $uv'-u'v\in P$ vanishes at $r$,
so $\sigma_r$ is well defined; it respects sums and products, as $u/v+u'/v'=(uv'+u'v)/(vv')$. So $F$, a subring of a field
on which $\sigma_r$ is a ring homomorphism, is a domain; its nonzero elements $u/v$ have $u\notin P$ and inverse $v/u$, so
$F$ is a field and $\sigma_r$ an embedding. For each $k$ choose $j_k$ with $x_{k,j_k}\notin P$.

Injectivity: if $\sigma_r=\sigma_{r'}$, then $r_{k,i}/r_{k,j_k}=\sigma_r(x_{k,i}/x_{k,j_k})=r'_{k,i}/r'_{k,j_k}$ for all
$k,i$, so $r=cr'$ blockwise with $c_k=r_{k,j_k}/r'_{k,j_k}$; $r$ and $r'$ lie in one orbit, and $r=r'$ by (2).

Surjectivity: let $\tau:F\to\mathbf C$ be an embedding, and put $z_{k,i}=\tau(x_{k,i}/x_{k,j_k})$. Then $z_{k,j_k}=1$, and
$z\in\hat X(\mathbf C)$: for $g\in P$ multihomogeneous of multidegree $a$, $g(z)=\tau\big(g/\prod_kx_{k,j_k}^{a_k}\big)=0$.
By (2), $z=cr$ for some $r\in T$ and $c\in(\mathbf C^\times)^K$, so
$\sigma_r(x_{k,i}/x_{k,j_k})=r_{k,i}/r_{k,j_k}=z_{k,i}=\tau(x_{k,i}/x_{k,j_k})$. Every element of $F$ is
$(u/\prod_kx_{k,j_k}^{a_k})/(v/\prod_kx_{k,j_k}^{a_k})$, a quotient of polynomials in the $x_{k,i}/x_{k,j_k}$, so
$\sigma_r=\tau$.
\end{proof}
% Lean: (1) arithVarZero (killBlocks, the chain Q n with comap_killBlocks_mono, and Q2 of the even constant terms),
% not_isHorizontal_of_dimOne (divSuppArith) and eval_ne_zero_of_dimOne (Ideal.exists_minimalPrimes_le). (2)
% regularityPointsSum_holds.proof: goodProjExists and fibreFinite for e = 0, pick through exists_isSheet_fin_zero,
% maDensity_fin_zero, integral_unique with volume 1 on Fin 0 -> C, the scaling of the block k0 by a^-1, horb and
% sameOrbit_of_segre. (3) regularityEmbedding_holds.proof: funField P (the Lean's subring generated by the quotients,
% inside FractionRing (MRing / P)), sigmaAt (IsLocalization.lift on the multihomogeneous forms outside P), sigmaAt_spec,
% hinj through hratio and SameOrbit, hsurj through chartHom r and chart_eval, and funField_hom_ext. Size: ArithVarZero.lean
% (178 lines), ExistsIsSheetFinZero.lean (105), NotMemOfNonZeroDivisor.lean (272), ExistsFunFieldMulEq.lean (43) and part
% of NotIsHorizontalOfDimOne.lean, about 650 lines.

\begin{definition}\label{cv-origin-intersection-family-curves-vertical-ring}
Let $K$, $(I_k)_{k\in K}$ and $R$ be as in Definition~\ref{lit-multiprojective}.
\begin{enumerate}
\item For prime ideals $P\subset W$ of $R$, let $\mathcal O_{P,W}$ be the subring of the fraction field of $R/P$ generated
  by the quotients $u/v$ of multihomogeneous $u,v\in R$ of the same multidegree with $v\notin W$, where $u$ and $v$ stand
  for their images in $R/P$.
\item Let $X$ be an arithmetic subvariety of dimension $1$ defined by $P$. A \emph{closed point} of $X$ is a relevant
  multihomogeneous prime $W\supset P$ with $\dim R/W=|K|$. Let $A_X$ be the subring of the field $F$ of
  Lemma~\ref{cv-origin-intersection-family-curves-points} formed by the elements that lie in $\mathcal O_{P,W}$ for
  every closed point $W$ of $X$.
\end{enumerate}
\end{definition}
\leanref{GroundTheory/HeightsAndTheta/Heights/Intersections/LocRing.lean}{GroundTheory.Heights₂.Heights.Inter.locRing}
\leanref{GroundTheory/HeightsAndTheta/Heights/Intersections/StatementsArith.lean}{GroundTheory.Heights₂.Heights.Inter.IsClosedPt}
\leanref{GroundTheory/HeightsAndTheta/Heights/Intersections/TwoMemRegRingSpecZ.lean}{GroundTheory.Heights₂.Heights.Inter.regRing}

\begin{lemma}\label{cv-origin-intersection-family-curves-vertical-local}
Let $K$, $(I_k)_{k\in K}$ and $R$ be as in Definition~\ref{lit-multiprojective}, let $P\subset W$ be prime ideals of
$R$, and write $\mathcal O_W$ for $\mathcal O_{P,W}$ (Definition~\ref{cv-origin-intersection-family-curves-vertical-ring}).
\begin{enumerate}
\item Every element of $\mathcal O_W$ is a quotient $u/v$ of multihomogeneous $u,v\in R$ of the same multidegree with
  $v\notin W$.
\item $\mathcal O_W$ is a local ring, and a quotient $u/v$ as in (1) is a unit of $\mathcal O_W$ if and only if
  $u\notin W$.
\item If $W$ is relevant, then $\mathcal O_W$ is noetherian.
\item Suppose that $P$ and $W$ are relevant and multihomogeneous, and choose $i_k\in I_k$ with $x_{k,i_k}\notin W$ for
  $k\in K$. Let $\mathfrak m_W$ be the maximal ideal of $\mathcal O_W$, and let $B=(R/P)_W$ be the localization of $R/P$
  at the image of $W$. Then the homomorphism of $\mathcal O_W$-algebras $\varphi:\mathcal O_W[T_k:k\in K]\to B$ with
  $\varphi(T_k)=x_{k,i_k}$ is injective, and it identifies $B$ with the localization of $\mathcal O_W[T_k:k\in K]$ at the
  prime ideal $\mathfrak m_W\mathcal O_W[T_k:k\in K]$. In particular $B$ is flat over $\mathcal O_W$, and its maximal
  ideal is $\mathfrak m_WB$.
\item Let $W'\supset P$ be a prime. If $W$ and $W'$ are relevant and multihomogeneous and
  $\mathcal O_{P,W}=\mathcal O_{P,W'}$, then $W=W'$.
\end{enumerate}
\end{lemma}
\leanref{GroundTheory/HeightsAndTheta/Heights/Intersections/PLocRingFrac.lean}{GroundTheory.Heights₂.Heights.Inter.LocRingFrac.proof}
\leanref{GroundTheory/HeightsAndTheta/Heights/Intersections/PLocRingLocal.lean}{GroundTheory.Heights₂.Heights.Inter.LocRingLocal.proof}
\leanref{GroundTheory/HeightsAndTheta/Heights/Intersections/PLocRingNoeth.lean}{GroundTheory.Heights₂.Heights.Inter.LocRingNoeth.proof}
\leanref{GroundTheory/HeightsAndTheta/Heights/Intersections/ImportsBisTer.lean}{GroundTheory.Heights₂.Heights.Inter.regularityLocalFlat_holds.proof}
\leanref{GroundTheory/HeightsAndTheta/Heights/Intersections/IsLocalizationConeVarMap.lean}{GroundTheory.Heights₂.Heights.Inter.regularityLocalFlat_holds.isLocalization_coneVarMap}
\leanref{GroundTheory/HeightsAndTheta/Heights/Intersections/PLocRingInj.lean}{GroundTheory.Heights₂.Heights.Inter.LocRingInj.proof}
\begin{proof}
(1) The quotients of (1) include $0=0/1$ and $1=1/1$, and they are closed under negation, sums and products:
$u/v+u'/v'=(uv'+u'v)/(vv')$ and $(u/v)(u'/v')=uu'/(vv')$, where $vv'\notin W$ as $W$ is prime. So they form the subring
that they generate.

(2) If $u\notin W$, then $v/u\in\mathcal O_W$ is an inverse of $u/v$. Conversely, let $u/v$ have an inverse, written
$u'/v'$ as in (1). Then $uu'-vv'\in P\subset W$; if $u\in W$, then $vv'\in W$, which is false. So the nonunits of
$\mathcal O_W$ are the $u/v$ with $u\in W$, and by the formula for sums in (1) a sum of two nonunits is a nonunit. So
$\mathcal O_W$ is local.

(3) Choose $c_k$ with $x_{k,c_k}\notin W$. Let $B_0\subset\mathcal O_W$ be the subring generated by the finitely many
quotients $x_{k,i}/x_{k,c_k}$, which is noetherian by Hilbert's basis theorem, and let $\theta:R\to B_0$ be the
homomorphism with $\theta(x_{k,i})=x_{k,i}/x_{k,c_k}$. For $u$ multihomogeneous of multidegree $a$,
$u=\theta(u)\prod_kx_{k,c_k}^{a_k}$ in the fraction field of $R/P$. So for $u,v$ as in (1), of multidegree $a$,
$u/v=\theta(u)/\theta(v)$, and $\theta(v)$ is a unit of $\mathcal O_W$, with inverse $\prod_kx_{k,c_k}^{a_k}/v$. Hence
$\mathcal O_W$ is the localization of $B_0$ at the elements of $B_0$ that are units of $\mathcal O_W$, and it is
noetherian.

(4) Write $x^e=\prod_kx_{k,i_k}^{e_k}$ for $e\in\mathbf N^K$, so that $\varphi(\sum_eg_eT^e)=\sum_eg_ex^e$. Let
$G=\sum_eg_eT^e$. By (1), applied to its finitely many coefficients and with a common denominator, there are a
multihomogeneous $V\notin W$ of some multidegree $b$ and multihomogeneous $N_e$ of multidegree $b$ with $g_eV=N_e$ for
every $e$. Then $\varphi(G)V=\sum_eN_ex^e$, a sum of multihomogeneous terms of the distinct multidegrees $b+e$.

If $\varphi(G)=0$, then $\sum_eN_ex^e\in P$, so each $N_ex^e$ lies in $P$, as $P$ is multihomogeneous, and $N_e\in P$, as
$x_{k,i_k}\notin P$. So every $g_e$ is $0$, and $\varphi$ is injective. The ideal $\mathfrak m_W\mathcal O_W[T]$ is
prime, the kernel of the reduction to $(\mathcal O_W/\mathfrak m_W)[T]$. If $G$ is not in it, some $g_e$ is a unit of
$\mathcal O_W$, so $N_e\notin W$ by (2), and $\sum_eN_ex^e\notin W$, as $W$ is multihomogeneous and $x^e\notin W$; so
$\varphi(G)V$ is a unit of $B$, and so is $\varphi(G)$. Every element of $R/P$ is a value of $\varphi$, as
$x_{k,i}=\varphi\big((x_{k,i}/x_{k,i_k})T_k\big)$. If $g\notin W$ and $\varphi(G)=g$, then $G$ is not in
$\mathfrak m_W\mathcal O_W[T]$: otherwise every $N_e$ lies in $W$ by (2), and $gV\equiv\sum_eN_ex^e$ modulo $P\subset W$
lies in $W$, which is false. So every element $f/g$ of $B$ is $\varphi(G_1)/\varphi(G)$ with $G$ outside
$\mathfrak m_W\mathcal O_W[T]$, and $B$ is the localization at this prime. The polynomial ring is free over
$\mathcal O_W$ and localizations are flat, so $B$ is flat over $\mathcal O_W$; the maximal ideal of the localization at
a prime is generated by that prime.

(5) We show $W\subset W'$; the other inclusion follows by symmetry. Choose $i_k$ and $j_k$ with $x_{k,i_k}\notin W$ and
$x_{k,j_k}\notin W'$. Then $x_{k,i_k}/x_{k,j_k}$ lies in $\mathcal O_{P,W'}=\mathcal O_{P,W}$ and is inverse to
$x_{k,j_k}/x_{k,i_k}\in\mathcal O_{P,W}$, so $x_{k,j_k}\notin W$ by (2). As $W$ is generated by multihomogeneous
elements, it suffices to show that a multihomogeneous $f\in W$, of multidegree $a$, lies in $W'$. Put
$t=\prod_kx_{k,j_k}^{a_k}$, which lies neither in $W$ nor in $W'$. By (2), $f/t$ is not a unit of $\mathcal O_{P,W}$; if
$f\notin W'$, it would be a unit of $\mathcal O_{P,W'}=\mathcal O_{P,W}$. So $f\in W'$.
\end{proof}
% Lean: (1) LocRingFrac.proof (Subring.closure_induction); (2) LocRingLocal.proof (key, IsLocalRing.of_nonunits_add);
% (3) LocRingNoeth.proof (B = Subring.closure of the ratios g p, is_noetherian_subring_closure, theta = eval2Hom g, hkey,
% IsLocalization.isNoetherianRing for the units of locRing P W); (4) regularityLocalFlat_holds.proof: coneVarMap (aeval of
% the X (k, v k) into coneLocRing P W hPW, the Lean's (R/P)_W inside the fraction field), isLocalization_coneVarMap
% (exists_clear for the common denominator, coneVarMap_injective through mem_of_clear_mem, isUnit_coneVarMap,
% exists_lift and exists_isUnit_coeff), IsLocalization.flat with Module.Flat.of_free and Module.Flat.trans, and
% IsLocalization.AtPrime.map_eq_maximalIdeal; the local rings come from regularityLocalFlat_holds.isLocalRing and
% IsLocalization.AtPrime.isLocalRing; (5) LocRingInj.proof (le_of_locRing_eq both ways: hunit, hjW, then frac_mem_locRing
% for f/t). Size: PLocRingFrac.lean, PLocRingLocal.lean, PLocRingNoeth.lean, IsLocalizationConeVarMap.lean,
% OneNotMemOfIsPrime.lean, ImportsBisTer.lean and PLocRingInj.lean, about 700 lines.

\begin{lemma}\label{cv-origin-intersection-family-curves-vertical-chain}
Let $K$ and $(I_k)_{k\in K}$ be as in Definition~\ref{lit-multiprojective}, let $A_0$ be $\mathbf Z$ or a field, and let
$R_0=A_0[x_{k,i}:k\in K,\ i\in I_k]$, graded by multidegree. Let $E$ be a field, let $\psi:R_0\to E$ be a ring
homomorphism, and let $j_k\in I_k$ with $\psi(x_{k,j_k})=1$ for $k\in K$. For $S\subset K$ let
$\varphi_S:R_0\to E[t_k:k\in K]$ be the ring homomorphism with $\varphi_S(x_{k,i})=0$ for $k\in S$ and
$\varphi_S(x_{k,i})=\psi(x_{k,i})t_k$ for $k\notin S$, and let $Q_S$ be its kernel, a prime ideal.
\begin{enumerate}
\item $\varphi_\emptyset(g)=\psi(g)\prod_kt_k^{a_k}$ for $g$ multihomogeneous of multidegree $a$.
\item $Q_S\subsetneq Q_{S\cup\{k\}}$ for $k\notin S$.
\item Let $V$ be a multihomogeneous prime of $R_0$ with $\psi(g)=0$ for every multihomogeneous $g\in V$, and suppose that
  $\psi(u)=0$ for some multihomogeneous $u\notin V$. Then $\dim R_0/V\ge|K|+1$.
\item If in (3) moreover $A_0=\mathbf Z$ and $E$ has characteristic zero, then $\dim R_0/V\ge|K|+2$.
\end{enumerate}
\end{lemma}
\leanref{GroundTheory/HeightsAndTheta/Heights/Intersections/FunFieldLeLocRing.lean}{GroundTheory.Heights₂.Heights.Inter.RegularityFunctionFieldNumber.specMap}
\leanref{GroundTheory/HeightsAndTheta/Heights/Intersections/FunFieldLeLocRing.lean}{GroundTheory.Heights₂.Heights.Inter.RegularityFunctionFieldNumber.card_add_two_le}
\leanref{ReciprocityTheory/ReductiveGroups/ArithmeticGeometry/ArithmeticIntersection/MultihomogeneousIdeals/CurvePt.lean}{ReciprocityTheory.ChartRestriction.G.AI.I1.isUnit_mk_chartIdeal_form}
\begin{proof}
(1) For a monomial of multidegree $a$, $\varphi_\emptyset$ multiplies its image under $\psi$ by $\prod_kt_k^{a_k}$; sum
over the monomials of $g$.

(2) The map $\varphi_{S\cup\{k\}}$ is $\varphi_S$ followed by $t_k\mapsto0$, so $Q_S\subset Q_{S\cup\{k\}}$. The element
$x_{k,j_k}$ lies in $Q_{S\cup\{k\}}$ but not in $Q_S$, as $\varphi_S(x_{k,j_k})=t_k$.

(3) The ideal $V$ is generated by multihomogeneous elements $g$, and $\varphi_\emptyset(g)=0$ for these by (1); so
$V\subset Q_\emptyset$, and $u\in Q_\emptyset\setminus V$ by (1). A strict inclusion of primes $Q\subsetneq Q'$ gives
$\dim R_0/Q\ge\dim R_0/Q'+1$. Adding the blocks one at a time, (2) gives a chain
$V\subsetneq Q_\emptyset\subsetneq\dots\subsetneq Q_K$ of length $|K|+1$, so $\dim R_0/V\ge|K|+1+\dim R_0/Q_K$.

(4) The map $\varphi_K$ sends $g$ to the image in $E$ of its constant term. As $E$ has characteristic zero, $Q_K$ is the
set of the $g$ with constant term $0$, which is strictly contained in the prime ideal of the $g$ with even constant
term. So $\dim R_0/Q_K\ge1$, and the chain of (3) gives $\dim R_0/V\ge|K|+2$.
\end{proof}
% Lean: the chain argument appears twice. card_add_two_le is (4) for A_0 = Z, with specMap psi S (aeval of cS S k * C
% (psi (X p))), specMap_form for (1), hPQ, hstep and hlast (Q univ < Qt, the kernel of the even constant terms), chain by
% Finset.induction_on and HorizNotMem.ringKrullDim_quotient_add_one_le; it is used in
% RegularityFunctionFieldNumber.proof. isUnit_mk_chartIdeal_form contains (3) for A_0 = F a field, with phi S = aeval of
% c S q.1 * theta (X q), hVQ (aeval_blockScale_eq_zero), hmono, hin, hout, hbase and chain (cyc_dim_add_one_le); it is
% used in chart_isField. Size: about 100 lines in FunFieldLeLocRing.lean and about 100 in CurvePt.lean.

\begin{lemma}\label{cv-origin-intersection-family-curves-vertical-closed}
Let $K$, $(I_k)_{k\in K}$ and $R$ be as in Definition~\ref{lit-multiprojective}, and let $X$ be an arithmetic subvariety
of dimension $1$ defined by $P$.
\begin{enumerate}
\item For prime ideals $Q\subset Q'$ of $R$, $\dim R/Q=\dim R/Q'+\operatorname{ht}(Q'/Q)$, where $\operatorname{ht}(Q'/Q)$
  is the height of $Q'/Q$ in $R/Q$. If $s\in R\setminus Q$ and $Q'$ is a minimal prime of $Q+(s)$, then
  $\dim R/Q'=\dim R/Q-1$.
\item Every relevant multihomogeneous prime $W$ of $R$ has $\dim R/W\ge|K|$.
\item A relevant multihomogeneous prime $W$ with $P\subsetneq W$ is a closed point of $X$
  (Definition~\ref{cv-origin-intersection-family-curves-vertical-ring}). For every closed point $W$ of $X$ there is a
  prime number $p$ with $W\cap\mathbf Z=p\mathbf Z$.
\item Let $s\in R\setminus P$ be multihomogeneous. The relevant minimal primes of $P+(s)$ are multihomogeneous, and they
  are the closed points of $X$ that contain $s$. There are finitely many of them.
\end{enumerate}
\end{lemma}
\leanref{ReciprocityTheory/ReductiveGroups/ArithmeticGeometry/BirationalGeometry/InterDivSuppA.lean}{ReciprocityTheory.ChartRestriction.G.Geo.mvPolynomialInt_dimFormula}
\leanref{ReciprocityTheory/ReductiveGroups/ArithmeticGeometry/BirationalGeometry/InterDivSupp.lean}{ReciprocityTheory.ChartRestriction.G.Geo.dim_quot_minimal_add_one}
\leanref{GroundTheory/HeightsAndTheta/Heights/Intersections/CardLeRingKrullDim.lean}{GroundTheory.Heights₂.Heights.Inter.RegularityRelativeClosedPoint.card_le_ringKrullDim}
\leanref{GroundTheory/HeightsAndTheta/Heights/Intersections/CardLeRingKrullDim.lean}{GroundTheory.Heights₂.Heights.Inter.RegularityRelativeClosedPoint.proof}
\leanref{GroundTheory/HeightsAndTheta/Heights/Intersections/ImportsBisRest.lean}{GroundTheory.Heights₂.Heights.Inter.RegularityLocalize.exists_prime_over}
\leanref{ReciprocityTheory/ReductiveGroups/ArithmeticGeometry/BirationalGeometry/InterDivSupp.lean}{ReciprocityTheory.ChartRestriction.G.Geo.isMultiHomogIdeal_of_mem_minimalPrimes}
\leanref{GroundTheory/HeightsAndTheta/Heights/Intersections/LeOfHeightMapEqZero.lean}{GroundTheory.Heights₂.Heights.Inter.RegularityDivisorSupportClosed.proof}
\begin{proof}
(1) The ring $R/Q$ is a domain of finite type over $\mathbf Z$. As in the proof of
Lemma~\ref{cv-origin-weight-birational}, the dimension formula \cite[Tags 00NM and 02JU]{Sta} shows that every maximal
ideal of $R/Q$ has height $\dim R/Q$; as $R/Q$ is catenary, this gives the formula. If $Q'$ is a minimal prime of
$Q+(s)$, then $Q'/Q$ is minimal over the image of $s$ in $R/Q$, so its height is at most one \cite[Tag 00KV]{Sta}, and
it is not zero, as $s\notin Q$.

(2) The chain $W=Q_0\subsetneq Q_1\subsetneq\dots\subsetneq Q_{|K|}$ of the proof of
Lemma~\ref{cv-origin-intersection-family-curves-points}(1), which uses only that $W$ is relevant and multihomogeneous,
gives $\dim R/W\ge|K|$.

(3) By (1), $1+|K|=\dim R/P=\dim R/W+\operatorname{ht}(W/P)$ with $\operatorname{ht}(W/P)\ge1$, so $\dim R/W\le|K|$, and
$\dim R/W=|K|$ by (2). For a closed point $W$, $W\cap\mathbf Z$ is a prime ideal of $\mathbf Z$, and it is not $0$, as
there is no arithmetic subvariety of dimension $0$ (Lemma~\ref{cv-origin-intersection-family-curves-points}(1)).

(4) The ideal $P+(s)$ is generated by elements that are homogeneous for the grading by the degree in each block, so its
minimal primes are homogeneous for each of these gradings \cite[Tag 00JU]{Sta}, that is, multihomogeneous. A relevant
minimal prime $W$ of $P+(s)$ has $\dim R/W=\dim R/P-1=|K|$ by (1), so it is a closed point that contains $s$. Conversely,
let $W$ be a closed point that contains $s$, and $Q$ a prime with $P+(s)\subset Q\subsetneq W$. Then $Q\neq P$, as
$s\notin P$, and (1), applied to $P\subset Q$ and to $Q\subset W$, gives
$1+|K|=\dim R/W+\operatorname{ht}(W/Q)+\operatorname{ht}(Q/P)\ge|K|+2$, which is impossible. So $W$ is minimal over
$P+(s)$. The noetherian ring $R/(P+(s))$ has finitely many minimal primes.
\end{proof}
% Lean: (1) mvPolynomialInt_dimFormula (DimF.dim_quot_add_height with DimF.height_eq_of_isMaximal_int, the Lean's own
% proof that the maximal ideals of R/P all have height dim R/P, in InterDivSuppAA, -AB and -AC) and
% dim_quot_minimal_add_one (Ideal.height_le_one_of_isPrincipal_of_mem_minimalPrimes); (2) card_le_ringKrullDim (killBlocks
% chain); (3) RegularityRelativeClosedPoint.proof and exists_prime_over (with arithVarZero); (4)
% isMultiHomogIdeal_of_mem_minimalPrimes (isHomogeneous_of_mem_minimalPrimes for each block grading),
% RegularityDivisorSupportClosed.proof (e1, e2, le_of_height_map_eq_zero, not_add_add) and divSupp_finite. Size:
% InterDivSuppA.lean, InterDivSuppAA.lean, InterDivSuppAB.lean, InterDivSuppAC.lean, parts of InterDivSupp.lean,
% CardLeRingKrullDim.lean and LeOfHeightMapEqZero.lean, about 750 lines.

\begin{lemma}\label{cv-origin-intersection-family-curves-vertical-component}
Let $K$, $(I_k)_{k\in K}$ and $R$ be as in Definition~\ref{lit-multiprojective}, let $X$ be an arithmetic subvariety of
dimension $1$ defined by $P$, let $s\in R\setminus P$ be multihomogeneous of multidegree $a$, let $W$ be a relevant
minimal prime of $P+(s)$, and let $t\in R\setminus W$ be multihomogeneous of multidegree $a$. Write $\mathcal O_W$ for
$\mathcal O_{P,W}$ (Definition~\ref{cv-origin-intersection-family-curves-vertical-ring}), put $y=s/t\in\mathcal O_W$, and
let $p$ be the prime number with $W\cap\mathbf Z=p\mathbf Z$
(Lemma~\ref{cv-origin-intersection-family-curves-vertical-closed}(3) and (4)).
\begin{enumerate}
\item Every element of the maximal ideal of $\mathcal O_W$ has a power in $y\mathcal O_W$.
\item $\mathcal O_W$ is a noetherian local domain of dimension $1$. For $z\in\mathcal O_W\setminus\{0\}$, the
  $\mathcal O_W$-module $\mathcal O_W/z\mathcal O_W$ has finite length $\ell(z)$.
\item Choose $i_k\in I_k$ with $x_{k,i_k}\notin W$, and put $J=W+(x_{k,i_k}-1:k\in K)$. Then $R/J$ is a finite field with
  $p^{\deg(W)}$ elements, where $\deg(W)$ is its dimension over $\mathbf F_p$, the degree of the point
  $W_{\mathbf F_p}$ (Lemma~\ref{cv-origin-intersection-family-sections-degree} with no line bundles). The homomorphism
  $R\to\mathcal O_W$ with $x_{k,i}\mapsto x_{k,i}/x_{k,i_k}$, followed by the reduction modulo the maximal ideal,
  induces an isomorphism of $R/J$ onto the residue field of $\mathcal O_W$.
\item For $z\in\mathcal O_W\setminus\{0\}$, $\mathcal O_W/z\mathcal O_W$ is finite, with $p^{\deg(W)\ell(z)}$ elements.
\item The multiplicity $m_W$ of $W$ in $R/(P+(s))$ is $\ell(y)$. Hence $\log|\mathcal O_W/y\mathcal O_W|=m_W\deg(W)\log p$.
\end{enumerate}
\end{lemma}
\leanref{GroundTheory/HeightsAndTheta/Heights/Intersections/PLocRingPow.lean}{GroundTheory.Heights₂.Heights.Inter.LocRingPow.proof}
\leanref{GroundTheory/HeightsAndTheta/Heights/Intersections/ImportsBisBis.lean}{GroundTheory.Heights₂.Heights.Inter.regularityLocalDimension_holds.proof}
\leanref{GroundTheory/HeightsAndTheta/Heights/Intersections/NotIsUnitOfNumMem.lean}{GroundTheory.Heights₂.Heights.Inter.regularityLocalResidue_holds.proof}
\leanref{GroundTheory/HeightsAndTheta/Heights/Intersections/NotIsUnitOfNumMem.lean}{GroundTheory.Heights₂.Heights.Inter.regularityLocalResidue_holds.chart_field_card}
\leanref{GroundTheory/HeightsAndTheta/Heights/Intersections/ProofOfNoethPow.lean}{GroundTheory.Heights₂.Heights.Inter.RegularityLocalFinite.proof}
\leanref{GroundTheory/HeightsAndTheta/Heights/Intersections/SpanSingletonEqRangeLsmul.lean}{GroundTheory.Heights₂.Heights.Inter.regularityLocalCardinality_holds.proof}
\leanref{GroundTheory/HeightsAndTheta/Heights/Intersections/ImportsLast.lean}{GroundTheory.Heights₂.Heights.Inter.regularityDehomogenization_holds.proof}
\begin{proof}
(1) By Lemma~\ref{cv-origin-intersection-family-curves-vertical-closed}(4), $W$ is multihomogeneous. Let $y'$ be in the
maximal ideal, written $y'=u/v$ with $u,v$ multihomogeneous of multidegree $b$, $v\notin W$ and $u\in W$
(Lemma~\ref{cv-origin-intersection-family-curves-vertical-local}(1) and (2)). As $W$ is a minimal prime of $P+(s)$, the
radical of $(P+(s))R_W$ is $WR_W$, so there are $n\ge0$ and $d\notin W$ with $du^n\in P+(s)$. As $W$ is multihomogeneous,
some multihomogeneous component $d_e$ of $d$, of multidegree $e$, is not in $W$. Write $du^n=q+gs$ with $q\in P$, and put
$D=e+nb$. The component of multidegree $D$ of $du^n$ is $d_eu^n$; that of $gs$ is $Gs$, where $G$ is the component of
$g$ of multidegree $D-a$ if $D\ge a$ and $G=0$ otherwise; and that of $q$ lies in $P$, as $P$ is multihomogeneous. So
$d_eu^n-Gs\in P$, and in the fraction field of $R/P$
\[
 y'^n=\frac{u^n}{v^n}=\frac st\cdot\frac{Gt}{d_ev^n},
\]
where $Gt$ and $d_ev^n\notin W$ are multihomogeneous of multidegree $D$. So $y'^n\in y\mathcal O_W$.

(2) The ring $\mathcal O_W$ is a subring of a field, and it is local and noetherian by
Lemma~\ref{cv-origin-intersection-family-curves-vertical-local}(2) and (3). By (1), every prime ideal of $\mathcal O_W$
that contains $y$ is the maximal ideal $\mathfrak m_W$; so $\mathfrak m_W$ is minimal over $(y)$ and has height at most
one \cite[Tag 00KV]{Sta}. As $s\notin P$, $y\neq0$, and $y$ is not a unit, as $s\in W$
(Lemma~\ref{cv-origin-intersection-family-curves-vertical-local}(2)); so $\mathfrak m_W$ has height one. For $z\neq0$,
$\mathcal O_W/z\mathcal O_W$ is a noetherian local ring of dimension zero, so it has finite length.

(3) By Lemma~\ref{cv-origin-intersection-family-curves-vertical-closed}(4), $W$ is a closed point, so $\dim R/W=|K|$. Let
$\pi:R\to\mathbf F_p[x_{k,i}]$ be the reduction modulo $p$. Its kernel $pR$ lies in $W$, so $\bar W=\pi(W)$ is a relevant
multihomogeneous prime with $\mathbf F_p[x_{k,i}]/\bar W\cong R/W$, of dimension $|K|$, and $R/J$ is isomorphic to
$\mathbf F_p[x_{k,i}]/\bar J$ with $\bar J=\bar W+(x_{k,i_k}-1:k\in K)$. Let $b\notin\bar J$. Multiplying each monomial of
$b$ by a suitable product of powers of the $x_{k,i_k}$ gives, for every large enough $e\in\mathbf N^K$, a multihomogeneous
$b'$ of multidegree $e$ with $b'\equiv b$ modulo $\bar J$; then $b'\notin\bar W$. If $b'$ were not a unit modulo $\bar J$,
it would lie in a maximal ideal $M\supset\bar J$, and the reduction $\psi$ to the field $\mathbf F_p[x_{k,i}]/M$ would
satisfy $\psi(x_{k,i_k})=1$, $\psi(\bar W)=0$ and $\psi(b')=0$; Lemma~\ref{cv-origin-intersection-family-curves-vertical-chain}(3)
would give $\dim\mathbf F_p[x_{k,i}]/\bar W\ge|K|+1$. So $R/J$ is a field. It is a finitely generated $\mathbf F_p$-algebra,
so it is finite over $\mathbf F_p$ \cite[Tag 00FV]{Sta}, and it has $p^{\deg(W)}$ elements.

Let $\chi:R\to\mathcal O_W/\mathfrak m_W$ be the homomorphism of the statement. For $g$ multihomogeneous of multidegree
$c$, the image of $g$ in $\mathcal O_W$ is $g/\prod_kx_{k,i_k}^{c_k}$. So $\chi(x_{k,i_k}-1)=0$, and $\chi$ kills the
multihomogeneous generators of $W$, by Lemma~\ref{cv-origin-intersection-family-curves-vertical-local}(2); so $\chi$
factors through the field $R/J$, and it is injective on it. It is onto: an element $u/v$ of $\mathcal O_W$, with $u,v$ of
multidegree $c$ and $v\notin W$ (Lemma~\ref{cv-origin-intersection-family-curves-vertical-local}(1)), is the image of
$u/\prod_kx_{k,i_k}^{c_k}$ times the inverse of the unit $v/\prod_kx_{k,i_k}^{c_k}$, so its residue is
$\chi(u)\chi(v)^{-1}$.

(4) A module of finite length $\ell$ over a local ring whose residue field has $q$ elements has $q^\ell$ elements, by a
composition series. Apply this to (2) and (3).

(5) Let $B=(R/P)_W$. The multiplicity $m_W$ is the length of $(R/(P+(s)))_W=B/sB$ over $R_W$, which is its length over
$B$, as $R_W\to B$ is onto. As $t$ is a unit of $B$, $sB=yB$. By
Lemma~\ref{cv-origin-intersection-family-curves-vertical-local}(4), $B$ is a flat local $\mathcal O_W$-algebra with
maximal ideal $\mathfrak m_WB$, so $\operatorname{length}_B(B/yB)=\ell(y)\operatorname{length}_B(B/\mathfrak m_WB)=\ell(y)$
\cite[Tag 02M1]{Sta}. The last statement follows from (4).
\end{proof}
% Lean: (1) LocRingPow.proof (radical_map_of_mem_minimalPrimes, exists_component_not_mem, component_mul_right,
% component_mul_of_not_le, component_mem), stated for P prime and multihomogeneous; (2)
% regularityLocalDimension_holds.proof (hx0, hxu, hmin, Ideal.height_le_one_of_isPrincipal_of_mem_minimalPrimes,
% Ideal.height_span_singleton_eq_one_of_mem_nonZeroDivisors) and length_coker_ne_top; (3)
% regularityLocalResidue_holds.proof: geoNum 0 W is by definition the length over Z of R/J for the chart chosen by
% (hWrel k).choose; chart_field_card (the reduction to ZMod p, IsProjVarF Wb 0, chart_isField through exists_homog_chart
% and isUnit_mk_chartIdeal_form, chart_finite through finite_of_finite_type_of_isJacobsonRing), chartMap, chartMap_form,
% not_isUnit_of_num_mem and the bijection eK; (4) RegularityLocalFinite.proof and
% regularityLocalCardinality_holds.proof with card_eq_card_residueField_pow_length (card_eq_pow_of_length_eq); (5)
% regularityDehomogenization_holds.proof: LocalizedQuotientLength.proof for the first two equalities, hspan for sB = yB,
% and key (Algebra.TensorProduct.quotIdealMapEquivTensorQuot and IsLocalRing.length_baseChange, the Lean's form of Stacks
% Tag 02M1). Size: PLocRingPow.lean, ImportsBisBis.lean, NotIsUnitOfNumMem.lean, ProofOfNoethPow.lean,
% SpanSingletonEqRangeLsmul.lean, CardLength.lean, ImportsLast.lean, PLocalizedQuotientLength.lean, part of
% KrullDimQuotientEqOneIffEqBot.lean and part of CurvePt.lean, about 650 lines.

\begin{lemma}\label{cv-origin-intersection-family-curves-vertical-order}
Let $K$, $(I_k)_{k\in K}$ and $R$ be as in Definition~\ref{lit-multiprojective}, let $X$ be an arithmetic subvariety of
dimension $1$ defined by $P$, let $F$ be as in Lemma~\ref{cv-origin-intersection-family-curves-points}, and let
$A=A_X\subset F$ and $\mathcal O_W=\mathcal O_{P,W}$ for the closed points $W$ of $X$ be as in
Definition~\ref{cv-origin-intersection-family-curves-vertical-ring}.
\begin{enumerate}
\item $F$ is a number field.
\item For every $x\in F$ there is a nonzero integer $M$ with $Mx\in A$.
\item Every element of $A$ is integral over $\mathbf Z$, and $A$ is a free $\mathbf Z$-module of rank $[F:\mathbf Q]$.
\item Let $p$ be a prime number, and let $y\in F$ lie in $\mathcal O_W$ for every closed point $W$ of $X$ with $p\in W$.
  Then there is an integer $n$ prime to $p$ with $ny\in A$.
\item For a closed point $W$, the set $\mathfrak m_W$ of the $x\in A$ that are not units of $\mathcal O_W$ is a maximal
  ideal of $A$, and $\mathcal O_W$ is the localization of $A$ at $\mathfrak m_W$: its elements are the $x/s$ with
  $x,s\in A$ and $s\notin\mathfrak m_W$. The map $W\mapsto\mathfrak m_W$ is a bijection from the set of the closed points
  of $X$ onto the set of the maximal ideals of $A$.
\item Let $b\in A$ be nonzero. Then $\mathcal O_W/b\mathcal O_W$ is finite for every closed point $W$, it has one element
  for all but finitely many $W$, and
  \[
   \sum_W\log|\mathcal O_W/b\mathcal O_W|=\log|A/bA|=\sum_\tau\log|\tau(b)|,
  \]
  the sums over the closed points $W$ of $X$ and over the field embeddings $\tau:F\to\mathbf C$.
\end{enumerate}
\end{lemma}
\leanref{GroundTheory/HeightsAndTheta/Heights/Intersections/FunFieldLeLocRing.lean}{GroundTheory.Heights₂.Heights.Inter.RegularityFunctionFieldNumber.proof}
\leanref{GroundTheory/HeightsAndTheta/Heights/Intersections/ImportsBisAlt.lean}{GroundTheory.Heights₂.Heights.Inter.RegularityDenominator.proof}
\leanref{GroundTheory/HeightsAndTheta/Heights/Intersections/KerIsMultiHomogIdeal.lean}{GroundTheory.Heights₂.Heights.Inter.RegularityValuationCentre.proof}
\leanref{GroundTheory/HeightsAndTheta/Heights/Intersections/ImportsBisMore.lean}{GroundTheory.Heights₂.Heights.Inter.RegularityIntegral.proof}
\leanref{GroundTheory/HeightsAndTheta/Heights/Intersections/OrderOfIntegral.lean}{GroundTheory.Heights₂.Heights.Inter.RegularityOrder.proof}
\leanref{GroundTheory/HeightsAndTheta/Heights/Intersections/ImportsBisFurther.lean}{GroundTheory.Heights₂.Heights.Inter.RegularityLocalPrime.proof}
\leanref{GroundTheory/HeightsAndTheta/Heights/Intersections/ImportsBisRest.lean}{GroundTheory.Heights₂.Heights.Inter.RegularityLocalize.proof}
\leanref{GroundTheory/HeightsAndTheta/Heights/Intersections/FunFieldExistsInv.lean}{GroundTheory.Heights₂.Heights.Inter.RegularityMaximalSurjective.proof}
\leanref{GroundTheory/HeightsAndTheta/Heights/Intersections/IntCastMemRegMaxAt.lean}{GroundTheory.Heights₂.Heights.Inter.RegularityMaximalSpectrum.proof}
\leanref{GroundTheory/HeightsAndTheta/Heights/Intersections/ImportsBisExtra.lean}{GroundTheory.Heights₂.Heights.Inter.RegularityLocal.proof}
\leanref{ReciprocityTheory/ReductiveGroups/ArithmeticGeometry/ArithmeticIntersection/CardQuotEqNatAbsNorm.lean}{ReciprocityTheory.ChartRestriction.G.AI.index_thm}
\begin{proof}
(1) Let $u/v\in F$ be nonzero, with $v\notin P$. Then $u\notin P$, and $v/u\in F$ is its inverse; so $F$ is a field, of
characteristic zero as $P\cap\mathbf Z=0$. Choose $j_k$ with $x_{k,j_k}\notin P$, put $r_{k,i}=x_{k,i}/x_{k,j_k}\in F$,
let $D\subset F$ be the $\mathbf Q$-algebra generated by the $r_{k,i}$, and let $\phi:R\to D$ be the homomorphism with
$\phi(x_{k,i})=r_{k,i}$. For $g$ multihomogeneous of multidegree $c$, $g=\phi(g)\prod_kx_{k,j_k}^{c_k}$ in the fraction
field of $R/P$; so $\phi(g)=0$ for multihomogeneous $g\in P$, and $u/v=\phi(u)/\phi(v)$ for $u,v$ of the same
multidegree with $v\notin P$.

We show that $D$ is a field. Let $d\in D$ be nonzero, written $d=u/v$ with $u,v\notin P$
(Lemma~\ref{cv-origin-intersection-family-curves-vertical-local}(1)), and suppose that $d$ is not a unit of $D$. Then
$d$ lies in a maximal ideal $\mathfrak n$ of $D$, and $E=D/\mathfrak n$ is a field of characteristic zero, as the nonzero
integers are units of $D$. Let $\psi:R\to E$ be $\phi$ followed by the reduction. Then $\psi(x_{k,j_k})=1$, $\psi$ kills
the multihomogeneous elements of $P$, and $\psi(u)$ is the reduction of $\phi(u)=d\,\phi(v)\in\mathfrak n$, so
$\psi(u)=0$, while $u\notin P$. Lemma~\ref{cv-origin-intersection-family-curves-vertical-chain}(4) gives
$\dim R/P\ge|K|+2$, against $\dim R/P=|K|+1$. So $D$ is a field; it contains $\phi(u)/\phi(v)=u/v$ for every $u/v\in F$,
so $D=F$. A field that is a finitely generated $\mathbf Q$-algebra is finite over $\mathbf Q$ \cite[Tag 00FV]{Sta}. So $F$
is a number field.

(2) Write $x=u/v$ with $u,v$ of multidegree $a$ and $v\notin P$. If a closed point $W$ does not contain $v$, then
$x\in\mathcal O_W$. The closed points that contain $v$ are the finitely many relevant minimal primes of $P+(v)$
(Lemma~\ref{cv-origin-intersection-family-curves-vertical-closed}(4)). For such a $W$, choose a nonzero integer $n\in W$
(Lemma~\ref{cv-origin-intersection-family-curves-vertical-closed}(3)) and $t=\prod_kx_{k,c_k}^{a_k}$ with
$x_{k,c_k}\notin W$. The integer $n$ lies in the maximal ideal of $\mathcal O_W$
(Lemma~\ref{cv-origin-intersection-family-curves-vertical-local}(2)), so by
Lemma~\ref{cv-origin-intersection-family-curves-vertical-component}(1) there are $N\ge0$ and $z\in\mathcal O_W$ with
$n^N=(v/t)z$; then $n^Nx=(u/t)z\in\mathcal O_W$. The product $M$ of these $n^N$ over the closed points that contain $v$
is a nonzero integer with $Mx\in\mathcal O_W$ for every closed point $W$.

(3) Let $\mathcal O_F$ be the ring of integers of $F$, let $\mathfrak p$ be a nonzero prime of $\mathcal O_F$ with
valuation $v_{\mathfrak p}$ and valuation ring $\mathcal O_{\mathfrak p}$, and let $n$ be a nonzero integer with
$v_{\mathfrak p}(n)>0$. We find a closed point $W$ of $X$ with $\mathcal O_W\subset\mathcal O_{\mathfrak p}$. For each $k$
fix $j^0_k$ with $x_{k,j^0_k}\notin P$, and choose $j_k$, among the $i$ with $x_{k,i}\notin P$, with
$v_{\mathfrak p}(x_{k,i}/x_{k,j^0_k})$ smallest. Then $r_{k,i}=x_{k,i}/x_{k,j_k}$ has $v_{\mathfrak p}(r_{k,i})\ge0$ for
every $i$, and the homomorphism $\phi:R\to F$ with $\phi(x_{k,i})=r_{k,i}$ takes values in $\mathcal O_{\mathfrak p}$. Let
$\psi:R\to E$ be $\phi$ followed by the reduction to the residue field $E$ of $\mathcal O_{\mathfrak p}$, so that
$\psi(x_{k,j_k})=1$, and let $W=Q_\emptyset$ be the kernel of the map $\varphi_\emptyset$ of
Lemma~\ref{cv-origin-intersection-family-curves-vertical-chain} for $\psi$ and the $j_k$. By (1) of that lemma, a
multihomogeneous $g$ lies in $W$ if and only if $v_{\mathfrak p}(\phi(g))>0$. The map $\varphi_\emptyset$ takes the
component of $g$ of degree $d$ in the block $k$ to the component of $\varphi_\emptyset(g)$ of degree $d$ in $t_k$, so the
prime $W$ is multihomogeneous; it is relevant, as $x_{k,j_k}\notin W$; it contains $P$, as $\phi$ kills the
multihomogeneous elements of $P$; and it contains $n$, so $W\neq P$. By
Lemma~\ref{cv-origin-intersection-family-curves-vertical-closed}(3), $W$ is a closed point. Let $z=u/w\in\mathcal O_W$,
with $u,w$ of multidegree $c$ and $w\notin W$. Then $z=\phi(u)/\phi(w)$, $v_{\mathfrak p}(\phi(u))\ge0$, and
$v_{\mathfrak p}(\phi(w))=0$ as $w\notin W$; so $v_{\mathfrak p}(z)\ge0$. Hence every $x\in A$ has $v_{\mathfrak p}(x)\ge0$
for every $\mathfrak p$, and $x\in\mathcal O_F$ \cite[Tag 031T]{Sta}.

So $A\subset\mathcal O_F$, and $A$ is a free $\mathbf Z$-module of rank at most $[F:\mathbf Q]$. Conversely, let
$s_1,\dots,s_r$ be a $\mathbf Z$-basis of $\mathcal O_F$. By (2) there are nonzero integers $m_i$ with $m_is_i\in A$, and
multiplication by $\prod_im_i$ embeds $\mathcal O_F$ in $A$. So $A$ has rank $[F:\mathbf Q]$.

(4) By (2) there is a nonzero integer $N$ with $Ny\in A$; write $N=p^em$ with $p\nmid m$. Let $W$ be a closed point. If
$p\in W$, then $my\in\mathcal O_W$ by hypothesis. If $p\notin W$, then $1/p\in\mathcal O_W$
(Lemma~\ref{cv-origin-intersection-family-curves-vertical-local}(2)), and $my=p^{-e}(Ny)\in\mathcal O_W$. So $my\in A$.

(5) By Lemma~\ref{cv-origin-intersection-family-curves-vertical-local}(2), $\mathfrak m_W$ is the inverse image of the
maximal ideal of the local ring $\mathcal O_W$, a prime ideal of $A$. Let $W\cap\mathbf Z=p\mathbf Z$; then $p\in
\mathfrak m_W$. The ring $A$ is finite over $\mathbf Z$ by (3), so $\mathfrak m_W$ lies in a maximal ideal $\mathfrak m$ of
$A$, both lie over $p\mathbf Z$, and $\mathfrak m_W=\mathfrak m$ \cite[Tag 00GT]{Sta}.

\emph{The localization.} Let $y\in\mathcal O_W$. Let $S$ be the set of the closed points that contain $p$, finite by
Lemma~\ref{cv-origin-intersection-family-curves-vertical-closed}(4) applied to $s=p$, and let $d_0$ be as in
Lemma~\ref{cv-origin-intersection-family-sections-avoid} for $S$. Write $y=u/v$ with $u,v$ multihomogeneous of multidegree
$a+\mathbf 1$ and $v\notin W$, where $\mathbf 1$ is the multidegree that is $1$ everywhere: multiply a representative of
Lemma~\ref{cv-origin-intersection-family-curves-vertical-local}(1) above and below by $\prod_kx_{k,c_k}\notin W$. Put
$m=\sum_kd_{0,k}$, and let $h$ be multihomogeneous of multidegree $(m+1)(a+\mathbf 1)\ge d_0$ and in no element of $S$.
Then $s_0=v^{m+1}/h$ and $x_0=v^mu/h$ lie in $F$ and in $\mathcal O_{W'}$ for every $W'\in S$, and $ys_0=x_0$. By (4)
there are integers $n_1$ and $n_2$ prime to $p$ with $n_1s_0,n_2x_0\in A$. Put $s=n_1n_2s_0$ and $x=n_1n_2x_0$, elements of
$A$ with $ys=x$. As $p\nmid n_1n_2$ and $v\notin W$, $n_1n_2v^{m+1}\notin W$, so $s$ is a unit of $\mathcal O_W$ by
Lemma~\ref{cv-origin-intersection-family-curves-vertical-local}(2), that is, $s\notin\mathfrak m_W$. Conversely $x/s$ lies
in $\mathcal O_W$ for $x\in A$ and $s\in A\setminus\mathfrak m_W$.

\emph{The bijection.} If $\mathfrak m_W=\mathfrak m_{W'}$, then $\mathcal O_{P,W}=\mathcal O_{P,W'}$, both being the
localization of $A$ at this ideal, so $W=W'$ by Lemma~\ref{cv-origin-intersection-family-curves-vertical-local}(5). Let
$\mathfrak m$ be a maximal ideal of $A$. As $A$ is integral over $\mathbf Z$, $\mathfrak m\cap\mathbf Z$ is a maximal ideal
$p\mathbf Z$. Suppose that $\mathfrak m\not\subset\mathfrak m_W$ for every closed point $W$. Let $S$ be the finite set of
the closed points that contain $p$. The prime ideals $0$ and $\mathfrak m_W$ for $W\in S$ do not contain $\mathfrak m$, as
$p\in\mathfrak m$, so there is $a\in\mathfrak m$ in none of them \cite[Tag 00DS]{Sta}. Then $a\neq0$, and $a$ is a unit of
$\mathcal O_W$ for $W\in S$, so $a^{-1}\in F$ lies in $\mathcal O_W$ for every closed point $W$ that contains $p$. By (4)
there is an integer $n$ prime to $p$ with $na^{-1}\in A$, and then $n=a\cdot na^{-1}\in\mathfrak m\cap\mathbf Z=p\mathbf
Z$, a contradiction. So $\mathfrak m\subset\mathfrak m_W$ for some $W$, and $\mathfrak m=\mathfrak m_W$ as $\mathfrak m$
is maximal.

(6) By (3), $A$ is free over $\mathbf Z$ of finite rank, so $A/bA$ is finite: multiplication by $b$ is injective, and
$|A/bA|$ is the absolute value of its determinant. A $\mathbf Z$-basis of $A$ is a $\mathbf Q$-basis of $F$, as $A$ has
rank $[F:\mathbf Q]$, so this determinant is the norm $N_{F/\mathbf Q}(b)=\prod_\tau\tau(b)$. The finite ring $A/bA$ is
artinian, so it is the product of its localizations at its finitely many prime ideals, which are maximal
\cite[Tag 00JA]{Sta}. These are the $\mathfrak m/bA$ for the maximal ideals $\mathfrak m\ni b$ of $A$, which by (5) are the
$\mathfrak m_W$ for the closed points $W$ with $b\in\mathfrak m_W$, each for one $W$. The localization of $A/bA$ at
$\mathfrak m_W$ is $\mathcal O_W/b\mathcal O_W$, by (5). For the other $W$, $b$ is a unit of $\mathcal O_W$ and
$\mathcal O_W/b\mathcal O_W$ has one element. So $\sum_W\log|\mathcal O_W/b\mathcal O_W|=\log|A/bA|
=\log|N_{F/\mathbf Q}(b)|=\sum_\tau\log|\tau(b)|$.
\end{proof}
% Lean: (1) RegularityFunctionFieldNumber.proof (hF, hY for the characteristic, D = Algebra.adjoin Q of the ratios,
% hunit through card_add_two_le with psi to D/m, hFD, finite_of_finite_type_of_isJacobsonRing). (2)
% RegularityDenominator.proof (key through LocRingPow with C n, M the product over divSupp P v); its statement only says
% that some nonzero c in funField P has c and x c in every locRing P W, and RegularityOrder.proof (order_of_integral) and
% RegularityLocalPrime.proof recover a nonzero integer in (c) through the integrality of c
% (Ideal.comap_ne_bot_of_integral_mem), where the text uses the integer M of the proof. (3) RegularityIntegral.proof from
% RegularityValuationCentre.proof (j by Finset.exists_max_image, hrle, chartHom r, psiE to the fraction field of
% v.integer / ltIdeal 1, the kernel of specMap psiE {} with ker_isMultiHomogIdeal and specMap_component) and
% RegularityRelativeClosedPoint, with mem_integers_of_valuation_le_one; order_of_integral (the embedding into 𝓞 F and
% the multiple of 𝓞 F). (4) RegularityLocalPrime.proof. (5) regMaxAt and not_mem_regMaxAt_iff, regMaxAt_isMaximal
% (intCast_mem_regMaxAt, Ideal.isMaximal_of_isIntegral_of_isMaximal_comap), RegularityLocalize.proof (exists_prime_over,
% AvoidRing.exists_form_avoid, h of multidegree (m + 1) (a + 1), exists_funField_mul_eq), RegularityMaximalSpectrum.proof
% (hsub, LocRingInj) and RegularityMaximalSurjective.proof (Ideal.subset_union_prime, funField_exists_inv); the Lean's
% AutoClose.Inter_R1RegMaxSurj is RegularityMaximalSurjective.proof. (6) RegularityLocal.proof
% (Ideal.finiteQuotientOfFreeOfNeBot, isLocalization_of_mem_iff, card_quot_span_eq_card_localization,
% log_card_eq_sum_log_card_localization through PrimeSpectrum.toPiLocalizationEquiv) and index_thm
% (card_quot_eq_natAbs_norm through Submodule.natAbs_det_equiv, norm_int_eq_norm_rat, abs_norm_eq_prod), combined in
% RegularityFormula.proof. Size: FunFieldLeLocRing.lean, ImportsBisAlt.lean, KerIsMultiHomogIdeal.lean,
% ExistsIntValuationLtOne.lean, ImportsBisMore.lean, OrderOfIntegral.lean, ImportsBisFurther.lean, ImportsBisRest.lean,
% FunFieldExistsInv.lean, IntCastMemRegMaxAt.lean, ImportsBisExtra.lean, ImportsBisSide.lean, LocQuotCard.lean and
% CardQuotEqNatAbsNorm.lean, about 1,300 lines.

\begin{lemma}\label{cv-origin-intersection-family-curves-vertical}
Let $K$, $(I_k)_{k\in K}$ and $R$ be as in Definition~\ref{lit-multiprojective}, let $X$ be an arithmetic subvariety of
dimension $1$ defined by $P$, with $F$ as in Lemma~\ref{cv-origin-intersection-family-curves-points}, and let
$f,g\in R\setminus P$ be multihomogeneous of the same multidegree. For multihomogeneous $s\in R\setminus P$ let
$\nu(s)=\sum_Wm_W\deg(W)\log p_W$, the sum over the relevant minimal primes $W$ of $P+(s)$, with their multiplicities
$m_W$ in $R/(P+(s))$, where $W\cap\mathbf Z=p_W\mathbf Z$ and $\deg(W)$ is the degree of
Definition~\ref{lit-multiprojective}(3) with no line bundles. Then
\[
 \nu(f)-\nu(g)=\sum_\tau\log|\tau(f/g)|,
\]
the sum over the field embeddings $\tau:F\to\mathbf C$.
\end{lemma}
\leanref{GroundTheory/HeightsAndTheta/Heights/Intersections/ImportsExtraAlt.lean}{GroundTheory.Heights₂.Heights.Inter.VerticalEmbedding.Statement}
\leanref{GroundTheory/HeightsAndTheta/Heights/Intersections/ImportsBisTail.lean}{GroundTheory.Heights₂.Heights.Inter.VerticalEmbedding.proof}
\begin{proof}
Let $a$ be the multidegree of $f$ and $g$, and $x=f/g\in F$. Write $\mathcal O_W=\mathcal O_{P,W}$ and $A=A_X$
(Definition~\ref{cv-origin-intersection-family-curves-vertical-ring}). For a relevant prime $W$ choose $i_k$ with
$x_{k,i_k}\notin W$ and put $t_W=\prod_kx_{k,i_k}^{a_k}$, multihomogeneous of multidegree $a$ and not in $W$.

\emph{The terms as local quotients.} Let $s\in\{f,g\}$. By
Lemma~\ref{cv-origin-intersection-family-curves-vertical-closed}(4), the relevant minimal primes of $P+(s)$ are the closed
points of $X$ that contain $s$. For such a $W$, Lemma~\ref{cv-origin-intersection-family-curves-vertical-component}(5)
with $t=t_W$ shows that its term in $\nu(s)$ is $m_W\deg(W)\log p_W=\log|\mathcal O_W/y_s(W)\mathcal O_W|$, where
$y_s(W)=s/t_W$. For a closed point $W$ let $\nu_s(W)$ be $\log|\mathcal O_W/y_s(W)\mathcal O_W|$ if $s\in W$ and $0$
otherwise. Then $\nu(s)=\sum_W\nu_s(W)$, the sum over the closed points.

\emph{A common denominator.} By Lemma~\ref{cv-origin-intersection-family-curves-vertical-order}(2) there is a nonzero
integer $c$ with $b=cx\in A$; also $c\in A$, and $b\neq0$. We show that for every closed point $W$
\[
 \nu_f(W)-\nu_g(W)=\log|\mathcal O_W/b\mathcal O_W|-\log|\mathcal O_W/c\mathcal O_W|.
 \tag{$*$}
\]
For nonzero $z_1,z_2\in\mathcal O_W$ with $\mathcal O_W/z_1\mathcal O_W$ and $\mathcal O_W/z_2\mathcal O_W$ finite,
multiplication by $z_1$ gives an exact sequence
$0\to\mathcal O_W/z_2\mathcal O_W\to\mathcal O_W/z_1z_2\mathcal O_W\to\mathcal O_W/z_1\mathcal O_W\to0$, as
$\mathcal O_W$ is a domain; so $|\mathcal O_W/z_1z_2\mathcal O_W|=|\mathcal O_W/z_1\mathcal O_W|\,
|\mathcal O_W/z_2\mathcal O_W|$. Also $z\mathcal O_W=zu\mathcal O_W$ for a unit $u$.

Suppose that $W$ contains $f$ or $g$. Then $t_W\notin W$, $y_f(W)$ and $y_g(W)$ are nonzero elements of $\mathcal O_W$,
and by Lemma~\ref{cv-origin-intersection-family-curves-vertical-component}(4), applied with $s=f$ or $s=g$, the quotients
of $\mathcal O_W$ by its nonzero elements are finite. As $y_f(W)=x\,y_g(W)$, $y_f(W)\,c=b\,y_g(W)$, and so
\[
 \log|\mathcal O_W/y_f(W)\mathcal O_W|+\log|\mathcal O_W/c\mathcal O_W|
 =\log|\mathcal O_W/b\mathcal O_W|+\log|\mathcal O_W/y_g(W)\mathcal O_W|.
\]
If $g\notin W$, then $y_g(W)$ is a unit of $\mathcal O_W$
(Lemma~\ref{cv-origin-intersection-family-curves-vertical-local}(2)), and its term is $0=\nu_g(W)$; likewise for $f$. So
$(*)$ holds. If $W$ contains neither $f$ nor $g$, then $x=f/g$ is a unit of $\mathcal O_W$ by the same lemma, so
$b\mathcal O_W=c\mathcal O_W$, and both sides of $(*)$ are $0$.

\emph{Summation.} By Lemma~\ref{cv-origin-intersection-family-curves-vertical-order}(6) for $b$ and for $c$, and as only
finitely many closed points contain $f$ or $g$, each term of $(*)$ is nonzero for finitely many $W$. Summing $(*)$ over the
closed points and applying Lemma~\ref{cv-origin-intersection-family-curves-vertical-order}(6),
\[
 \nu(f)-\nu(g)=\sum_\tau\log|\tau(b)|-\sum_\tau\log|\tau(c)|=\sum_\tau\log|\tau(x)|,
\]
as $\tau(b)=\tau(x)\tau(c)$ with $\tau(x)$ and $\tau(c)$ nonzero.
\end{proof}
\leanref{GroundTheory/HeightsAndTheta/Heights/Intersections/LogCardLocIdealUnitMul.lean}{GroundTheory.Heights₂.Heights.Inter.RegularityLocalProduct.proof}
% Lean: VerticalEmbedding.of_remaining = VerticalEmbedding.proof (t W the coordinate product, y s W = s / t W, term and
% hvert through regularityDehomogenization_holds and regularityLocalCardinality_holds) with
% RegularityLocalProduct.of_remaining = RegularityLocalProduct.proof (RegularityDenominator for c, key for (*) with
% log_card_locIdeal_mul through card_quot_span_mul, log_card_locIdeal_unit and log_card_locIdeal_unit_mul,
% frac_unit_locRing, RegularityLocalFinite and RegularityDivisorSupportClosed; RegularityFormula for xc and for c, with
% finsum_sub_distrib). The Lean's c is an element of funField P with c and x c in every locRing P W; the text takes the
% integer of Lemma cv-origin-intersection-family-curves-vertical-order (2). Size: ImportsBisTail.lean and
% LogCardLocIdealUnitMul.lean, about 240 lines, besides Lemmas cv-origin-intersection-family-curves-vertical-local,
% -chain, -closed, -component and -order; about 4,700 lines with these.

\begin{lemma}\label{cv-origin-intersection-family-curves}
Let $K$, $(I_k)_{k\in K}$ and $R$ be as in Definition~\ref{lit-multiprojective}, let $X$ be an arithmetic subvariety of
dimension $1$ defined by $P$, let $\bar B=(b,h_B)$ be a smooth metrized line bundle, and let $s,t\in R\setminus P$ be
multihomogeneous of multidegree $b$. Then for every $I$ as in
Definition~\ref{cv-origin-intersection-family-numbers}(1),
$\mathcal R_I(\,;\bar B,s\mid X)=\mathcal R_I(\,;\bar B,t\mid X)$.
\end{lemma}
\leanref{ReciprocityTheory/ReductiveGroups/ArithmeticGeometry/ArithmeticIntersection/RationalSections/ThmISt.lean}{ReciprocityTheory.ChartRestriction.G.AI.I2.DimensionOneRationalSectionIndependence}
\begin{proof}
Let $\nu$ be as in Lemma~\ref{cv-origin-intersection-family-curves-vertical}, and $T$ and the $\sigma_r$ as in
Lemma~\ref{cv-origin-intersection-family-curves-points}. By (1) of that lemma, no relevant minimal prime of $P+(s)$ or of
$P+(t)$ is horizontal, so $\mathcal R_I(\,;\bar B,s\mid X)=\nu(s)-\int_{X(\mathbf C)}\log(|s|/h_B)$, whatever $I$ is, and
likewise for $t$. Also by (1), $s$ and $t$ vanish nowhere on $\hat X(\mathbf C)$. The functions $|s|/h_B$, $|t|/h_B$ and
$|s|/|t|$ are constant on the orbits of $(\mathbf C^\times)^K$, as $s$, $t$ and $h_B$ are homogeneous of the same
multidegree. By (2), the integrals are the sums over $T$, and
\[
 \int_{X(\mathbf C)}\log\frac{|s|}{h_B}-\int_{X(\mathbf C)}\log\frac{|t|}{h_B}=\sum_{r\in T}\log\frac{|s(r)|}{|t(r)|}.
\]
By Lemma~\ref{cv-origin-intersection-family-curves-vertical} and (3),
$\nu(s)-\nu(t)=\sum_\tau\log|\tau(s/t)|=\sum_{r\in T}\log|\sigma_r(s/t)|=\sum_{r\in T}\log(|s(r)|/|t(r)|)$. Subtracting
gives $\mathcal R_I(\,;\bar B,s\mid X)=\mathcal R_I(\,;\bar B,t\mid X)$.
\end{proof}
\leanref{GroundTheory/HeightsAndTheta/Heights/Intersections/NotIsHorizontalOfDimOne.lean}{GroundTheory.Heights₂.Heights.Inter.RegularityIndependence.proof}
\leanref{GroundTheory/HeightsAndTheta/Heights/Intersections/LogNormDivSameOrbit.lean}{GroundTheory.Heights₂.Heights.Inter.RegularityProductFormula.proof}
\leanref{GroundTheory/HeightsAndTheta/Heights/Intersections/ImportsBisLast.lean}{GroundTheory.Heights₂.Heights.Inter.VerticalClosedPoints.proof}
% Lean: AutoClose.Inter_R1Indep = RegularityIndependence.of_remaining = RegularityIndependence.proof arithVarZero
% fibreFinite (RegularityProductFormula.of_remaining ...): cycR_eq_vertR_of_dimOne (no horizontal terms),
% ptsIntegral_add_of_dimOne with eval_ne_zero_of_dimOne (the integral splits as that of log(|s|/|t|) plus that of
% log(|t|/h)), and RegularityProductFormula.proof (regularityPointsSum_holds with log_norm_div_sameOrbit, and
% VerticalClosedPoints.proof: exists_funField_mul_eq for x = s/t, regularityEmbedding_holds for the bijection e, and
% VerticalEmbedding). The text uses (2) for all three integrals, where the Lean splits the integral through the finite
% fibre directly. Size: about 100 lines in NotIsHorizontalOfDimOne, LogNormDivSameOrbit and ImportsBisLast, besides
% Lemmas cv-origin-intersection-family-curves-points and -vertical.

\begin{lemma}\label{cv-origin-intersection-family-product}
Let $K$, $(I_k)_{k\in K}$ and $R$ be as in Definition~\ref{lit-multiprojective}, let $I$ be as in
Definition~\ref{cv-origin-intersection-family-numbers}(1), let $X$ be an arithmetic subvariety of dimension $n+1$ defined
by $P$, let $\bar L_1,\dots,\bar L_n$, $\bar A$ and $\bar B$ be smooth metrized line bundles, and let $s,t\in R\setminus P$
be multihomogeneous of the multidegrees of $\bar A$ and $\bar B$. Then
$\mathcal R_I(\bar L_1\cdots\bar L_n;\bar A\otimes\bar B,st\mid X)=\mathcal R_I(\bar L_1\cdots\bar L_n;\bar A,s\mid X)
+\mathcal R_I(\bar L_1\cdots\bar L_n;\bar B,t\mid X)$.
\end{lemma}
\leanref{ReciprocityTheory/ReductiveGroups/ArithmeticGeometry/ArithmeticIntersection/RationalSections/ThmISt.lean}{ReciprocityTheory.ChartRestriction.G.AI.I2.RMulStatement}
\leanref{GroundTheory/HeightsAndTheta/Heights/Intersections/ArIntMul.lean}{GroundTheory.Heights₂.Heights.Inter.RMul.proof}
\begin{proof}
The function $\log(|st|/(h_Ah_B))$ is $\log(|s|/h_A)+\log(|t|/h_B)$ where $s$ and $t$ do not vanish. Both sides of the
statement are a sum over components minus an integral, and the two parts add separately.

\emph{The components.} For $f\in R\setminus P$ and a prime $W\supset P$ let $m_W(f)$ be the length of the localization
at $W$ of $R/(P+(f))$, the multiplicity of $W$ in $R/(P+(f))$; it is finite if $W$ is minimal over $P+(f)$, and $0$ if
$f\notin W$, as $f$ is then a unit in the localization. Multiplication by $s$ gives an exact sequence
\[
 0\to R/(P+(t))\to R/(P+(st))\to R/(P+(s))\to0:
\]
if $sr=p+str'$ with $p\in P$, then $s(r-tr')\in P$, so $r-tr'\in P$ as $P$ is prime and $s\notin P$. Localizing at $W$
and taking lengths gives $m_W(st)=m_W(s)+m_W(t)$. A relevant minimal prime $W$ over $P+(st)$ contains $s$ or $t$; if it
contains $s$, it is a relevant minimal prime over $P+(s)$, as every prime between $P+(s)$ and $W$ contains $P+(st)$, and
likewise for $t$. Conversely, let $W$ be a relevant minimal prime over $P+(s)$, and $Q$ a prime with
$P+(st)\subset Q\subset W$. If $s\in Q$, then $Q=W$ by minimality. If $t\in Q$, then $Q\neq P$ as $t\notin P$, and no prime
lies strictly between $P$ and $W$ \cite[Tag 00KV]{Sta}, so again $Q=W$. So $W$ is a relevant minimal prime over
$P+(st)$, and likewise for $t$. Hence the sum over the components for $st$, with the multiplicities $m_W(st)$, is the sum
of those for $s$ and for $t$.

\emph{The integrals.} Let $\Psi$ be the good projection chosen for $X$ in
Definition~\ref{cv-origin-intersection-poincare-lelong-all-double-projection}(4); one exists by
Lemma~\ref{cv-origin-intersection-poincare-lelong-all-double-good}(1). By
Lemma~\ref{cv-origin-intersection-resolution-exists}, $X$ has an analytic resolution off some ideal, so by
Lemma~\ref{cv-origin-intersection-poincare-lelong-all-double-bridge} the sums over the fibres of $\Psi$ of the sheet
densities of $\log(|s|/h_A)$ and of $\log(|t|/h_B)$ against $\log h_1,\dots,\log h_n$ are integrable functions of
$y\in\mathbf C^n$. Let $c$ be the polynomial of Lemma~\ref{cv-origin-intersection-poincare-lelong-all-double-bridge-norm}
for $s$, homogeneous of degree $\mu$. As in the proof of
Lemma~\ref{cv-origin-intersection-poincare-lelong-all-double-bridge-null}(2), $y\mapsto c(1,y)$ is a nonzero polynomial
whose zero set has measure zero, and $c(1,y)=0$ whenever $s(z)=0$ for some $z\in\hat X(\mathbf C)$ with
$\Psi(\sigma(z))=(1,y)$. So for almost every $y$, $s$ has no zero at such $z$, and likewise $t$. For such $y$ and
$w\in F_\Psi(y)$, the chosen sheet $\gamma$ through $w$ at $y$ has $\gamma(y)\in\hat X(\mathbf C)$ and
$\Psi(\sigma(\gamma(y)))=(1,y)$, so $s$ and $t$ do not vanish at $\gamma(y)$, and the sheet density of
$\log(|st|/(h_Ah_B))$ at $(y,w)$ is the sum of those of $\log(|s|/h_A)$ and $\log(|t|/h_B)$. So the fibre sums add almost
everywhere, and so do their integrals.
\end{proof}
\leanref{GroundTheory/HeightsAndTheta/Heights/Intersections/NotIsHorizontalOfDimOne.lean}{GroundTheory.Heights₂.Heights.Inter.cycR_mul}
\leanref{GroundTheory/HeightsAndTheta/Heights/Intersections/ArIntMul.lean}{GroundTheory.Heights₂.Heights.Inter.arInt_mul}
% Lean: RMul.proof under PtsIntegrableStatement (ForSecondHalf.PtsIntegrable.of_resExists, the integrability part of
% Lemma cv-origin-intersection-poincare-lelong-all-double-bridge with Lemma cv-origin-intersection-resolution-exists):
% recR_eq; cycR_mul = cycR_eq_sum with divCycle_mul' (locLen_mul through mulQ_injective and mulQ_exact, locLen_quot_eq_zero,
% mem_divSupp_of_mul', and mem_divSupp_mul_left' through eq_of_le_of_mem_minimalPrimes, Krull's principal ideal theorem
% Ideal.height_le_one_of_isPrincipal_of_mem_minimalPrimes); arInt_mul (ptsIntegral, integral_add with the two
% integrability statements, ae_eval_ne_zero, which is projB_null for the ideal (s), and Real.log_mul at the chosen sheets).
% Size: ArIntMul.lean, 59 lines, and NotIsHorizontalOfDimOne.lean and MemDivSuppOfDivCycleNeZero.lean (parts), about
% 200 lines.

\begin{lemma}\label{cv-origin-intersection-family-sections-avoid}
Let $K$, $(I_k)_{k\in K}$ and $R$ be as in Definition~\ref{lit-multiprojective}, and let $S$ be a finite set of relevant
multihomogeneous primes of $R$. There is $d_0\in\mathbf N^K$ such that for every $d\in\mathbf N^K$ with $d\ge d_0$ there
is a multihomogeneous $f\in R$ of multidegree $d$ that lies in no $W\in S$.
\end{lemma}
\leanref{GroundTheory/HeightsAndTheta/Heights/Intersections/PChoice.lean}{GroundTheory.Heights₂.Heights.Inter.Choice.exists_form_avoid}
\leanref{GroundTheory/HeightsAndTheta/Heights/Intersections/AvoidRing.lean}{GroundTheory.Heights₂.Heights.Inter.AvoidRing.exists_form_avoid}
\begin{proof}
Let $T$ be the set of the maximal elements of $S$; every $W\in S$ lies in some $W'\in T$. For distinct $W,W'\in T$,
$W'\not\subset W$, and as $W'$ is generated by multihomogeneous elements, there is a multihomogeneous
$g_{W,W'}\in W'\setminus W$. For $W\in T$ let $a_W$ be the product of the $g_{W,W'}$ over $W'\in T\setminus\{W\}$, of
multidegree $\delta_W$. Then $a_W\notin W$, as $W$ is prime, and $a_W$ lies in every other element of $T$. Put
$d_0=\sum_{W\in T}\delta_W$, and let $d\ge d_0$. For $W\in T$, as $W$ is relevant, there are $i_k$ with $x_{k,i_k}\notin W$;
let $m_W=\prod_kx_{k,i_k}^{d_k-\delta_{W,k}}$, which is not in $W$. Let $f=\sum_{W\in T}a_Wm_W$, multihomogeneous of
multidegree $d$. For $W_0\in T$, every term with $W\neq W_0$ lies in $W_0$, and $a_{W_0}m_{W_0}\notin W_0$, so
$f\notin W_0$. If $W\in S$ lies in $W_0\in T$, then $f\notin W$.
\end{proof}
% Lean: Choice.exists_form_avoid (with R = Z), and its copy AvoidRing.exists_form_avoid used by WellDefined.proof: T the
% maximal elements (exists_le_maximal), gen W W' from the multihomogeneous generators of W', a W the product over
% T.erase W, m W = coordProd with the indices from relevance (coordProd_not_mem), and the sum. Size: about 110 lines
% in each of PChoice.lean and AvoidRing.lean.

\begin{lemma}\label{cv-origin-intersection-family-sections-degree}
Let $K$, $(I_k)_{k\in K}$ and $R$ be as in Definition~\ref{lit-multiprojective}, let $p$ be a prime number, let $n\ge0$,
and let $V$ be a relevant multihomogeneous prime of $R$ with $p\in V$ and $\dim R/V=n+|K|$, so that it defines a projective
variety $V_{\mathbf F_p}$ of dimension $n$ over $\mathbf F_p$. For $a_1,\dots,a_n\in\mathbf Z^K$ let
$\deg_{a_1\cdots a_n}(V)=(\mathcal O(a_1)\cdots\mathcal O(a_n)\cdot V_{\mathbf F_p})$, the degree
$\deg_{L_1\cdots L_n}(V)$ of Definition~\ref{lit-multiprojective}(3) for line bundles $\bar L_j$ of multidegrees $a_j$.
\begin{enumerate}
\item $\deg_{a_1\cdots a_n}(V)$ is additive in each $a_j$ and symmetric in $a_1,\dots,a_n$.
\item Suppose $\dim R/V=n+1+|K|$, and let $t\in R\setminus V$ be multihomogeneous of multidegree $b\in\mathbf N^K$. Then
  every relevant minimal prime $U$ of $V+(t)$ contains $p$ and has $\dim R/U=n+|K|$, and
  \[
   \deg_{a_1\cdots a_nb}(V)=\sum_Um_U\deg_{a_1\cdots a_n}(U),
  \]
  the sum over these $U$, where $m_U$ is the multiplicity of $U$ in $R/(V+(t))$.
\end{enumerate}
\end{lemma}
\leanref{ReciprocityTheory/ReductiveGroups/ArithmeticGeometry/BirationalGeometry/IsIntersectionFamily.lean}{ReciprocityTheory.ChartRestriction.G.Geo.geoNum}
\leanref{GroundTheory/HeightsAndTheta/Heights/Intersections/StatementsSwap.lean}{GroundTheory.Heights₂.Heights.Inter.GeoNumCut.Statement}
\leanref{GroundTheory/HeightsAndTheta/Heights/Intersections/PGeoNumCut.lean}{GroundTheory.Heights₂.Heights.Inter.GeoNumCut.full}
\leanref{GroundTheory/HeightsAndTheta/Heights/Intersections/PRSlotAdd.lean}{GroundTheory.Heights₂.Heights.Inter.geoNum_snoc_tensor}

These are properties of intersection numbers of line bundles on a projective variety over a field: they are additive
in each line bundle \cite[Tag 0BER]{Sta}, symmetric by their definition as coefficients of the Euler characteristic
of $\mathcal O(n_1a_1+\dots+n_na_n)$ \cite[Tag 0BEP]{Sta}, and computed through the effective Cartier divisor of $t$ on
$V_{\mathbf F_p}$, whose components are the $U_{\mathbf F_p}$ with multiplicities $m_U$ \cite[Tags 0BEU and 0BES]{Sta}.
The $U$ contain $V\ni p$, and $\dim R/U=\dim R/V-1$, as for the components in Definition~\ref{lit-multiprojective}(3).
% Lean: geoNum n V a is the Lean's own recursion through the coordinates X (k, i_k) of a chart of V, additive in its
% last argument by definition (geoNum_snoc_tensor); GeoNumCut.full is (2), through intNum_map_eq_geoNum (the reduction
% red p to Z/p), CutRegular.full and cycEval_divCycle_eq for the field intersection numbers intNum over Z/p; the
% symmetry is geoNum_comp_perm (Lemma cv-origin-intersection-family-symmetry), from GeoNumStatement,
% I1.ExistsUniqueStatement and I1.SymmStatement over Z/p.

\begin{lemma}\label{cv-origin-intersection-family-sections-position}
Let $K$, $(I_k)_{k\in K}$ and $R$ be as in Definition~\ref{lit-multiprojective}, let $n\ge0$, and let $I$ be as in
Definition~\ref{cv-origin-intersection-family-numbers}(1), additive in the last line bundle on arithmetic subvarieties of
dimension $n+1$. Let $X$ be an arithmetic subvariety of dimension $n+2$ defined by $P$, let
$\bar L_1,\dots,\bar L_n$, $\bar A$, $\bar C$ and $\bar B=(b,h_B)$ be smooth metrized line bundles, and let
$s\in R\setminus P$ be multihomogeneous of multidegree $b$. Then
\[
 \mathcal R_I(\bar L_1\cdots\bar L_n(\bar A\otimes\bar C);\bar B,s\mid X)
 =\mathcal R_I(\bar L_1\cdots\bar L_n\bar A;\bar B,s\mid X)+\mathcal R_I(\bar L_1\cdots\bar L_n\bar C;\bar B,s\mid X).
\]
\end{lemma}
\leanref{GroundTheory/HeightsAndTheta/Heights/Intersections/StatementsWellDefined.lean}{GroundTheory.Heights₂.Heights.Inter.RPositionAdd.Statement}
\leanref{GroundTheory/HeightsAndTheta/Heights/Intersections/PRSlotAdd.lean}{GroundTheory.Heights₂.Heights.Inter.RPositionAdd.proof}
\begin{proof}
Write $\bar L$ for $\bar L_1\cdots\bar L_n$ and $\bar A\otimes\bar C=(a+c,h_Ah_C)$. The three sides are sums over the
same relevant minimal primes $W$ of $P+(s)$, with the same multiplicities, minus an integral, and the terms add
separately. For horizontal $W$, an arithmetic subvariety of dimension $n+1$,
$I(\bar L(\bar A\otimes\bar C)\mid W)=I(\bar L\bar A\mid W)+I(\bar L\bar C\mid W)$ by hypothesis. For $W$ with
$W\cap\mathbf Z=p\mathbf Z$, the degree is additive in its last argument by
Lemma~\ref{cv-origin-intersection-family-sections-degree}(1).

For the integrals, let $\Psi$ be the good projection chosen for $X$
(Definition~\ref{cv-origin-intersection-poincare-lelong-all-double-projection}(4)). Let $y\in\mathbf C^{n+1}$ and
$w\in F_\Psi(y)$ with a sheet $\gamma$ through $w$ at $y$. Near $y$ every block of $\gamma$ is nonzero, so
$\log((h_Ah_C)\circ\gamma)=\log(h_A\circ\gamma)+\log(h_C\circ\gamma)$ there, and the complex Hessians add. The
Monge--Amp\`ere density, $(n+1)!\,(2/\pi)^{n+1}$ times a mixed discriminant (Definition~\ref{cv-period-monge-ampere}(2)),
is additive in each argument. So the sheet density of $\log(|s|/h_B)$ against
$\log h_1,\dots,\log h_n,\log(h_Ah_C)$ at $(y,w)$ is the sum of those against $\log h_1,\dots,\log h_n,\log h_A$ and
against $\log h_1,\dots,\log h_n,\log h_C$. The fibres are finite
(Lemma~\ref{cv-origin-intersection-poincare-lelong-all-double-good}(2)), so the fibre sums add, and the last two are
integrable by Lemma~\ref{cv-origin-intersection-poincare-lelong-all-double-independent}(1); so the integrals add.
\end{proof}
% Lean: RPositionAdd.proof (with fibreFinite and PtsIntegrableStatement, which RPositionAdd.of_remaining takes from
% PtsIntegrable): termwise, hI (AddAt) on the horizontal components through divSuppArith, geoNum_snoc_tensor on the
% vertical ones, and arInt_snoc_tensor (sheetDensity_snoc_tensor, maDensity_snoc_tensor: contDiffAt_log_metric,
% gaamix_cxHessian_comb and mdisc_update_add, the additivity of the mixed discriminant in one argument; integral_add).
% Size: PRSlotAdd.lean, about 150 lines.

\begin{lemma}\label{cv-origin-intersection-family-sections-cycles-local}
Let $A$ be a noetherian local domain of dimension $2$, and let $a,b\in A$ be nonzero with $A/(a,b)$ of finite length.
For a prime $\mathfrak p$ of $A$ of height one, let $\operatorname{ord}_{\mathfrak p}(a)$ be the length of
$A_{\mathfrak p}/aA_{\mathfrak p}$ and $\operatorname{ord}_{A/\mathfrak p}(b)$ the length of $A/(\mathfrak p+bA)$, and let
$\ker_{A/a}(b)$ be the kernel of multiplication by $b$ on $A/aA$. Then $\ker_{A/a}(b)\cong\ker_{A/b}(a)$, these kernels
have finite length, and
\[
 \ell\big(A/(a,b)\big)=\ell\big(\ker_{A/a}(b)\big)+\sum_{\mathfrak p\ni a}\operatorname{ord}_{\mathfrak p}(a)\,
 \operatorname{ord}_{A/\mathfrak p}(b),
\]
the sum over the primes of height one that contain $a$. In particular
\[
 \sum_{\mathfrak p\ni a}\operatorname{ord}_{\mathfrak p}(a)\,\operatorname{ord}_{A/\mathfrak p}(b)
 =\sum_{\mathfrak p\ni b}\operatorname{ord}_{\mathfrak p}(b)\,\operatorname{ord}_{A/\mathfrak p}(a).
\]
\end{lemma}
\leanref{ReciprocityTheory/ReductiveGroups/ArithmeticGeometry/ArithmeticIntersection/LocalLengthAdditivityTorsionFree.lean}{ReciprocityTheory.ChartRestriction.G.AI.LocalLengthMultiplicativity}
\leanref{ReciprocityTheory/ReductiveGroups/ArithmeticGeometry/ArithmeticIntersection/LengthTwoGeneratorQuotient.lean}{ReciprocityTheory.ChartRestriction.G.AI.length_two_generator_quotient}
\begin{proof}
\emph{The kernels.} Let $S=\{(x,y)\in A^2:bx=ay\}$. The maps $S\to\ker_{A/a}(b)$, $(x,y)\mapsto x\bmod a$, and
$S\to\ker_{A/b}(a)$, $(x,y)\mapsto y\bmod b$, are onto, and they have the same kernel: if $bx=ay$ and $x=az$, then
$ay=abz$ and $y=bz$, as $A$ is a domain, and conversely. So $\ker_{A/a}(b)\cong\ker_{A/b}(a)$. This module is finitely
generated and killed by $a$ and $b$, so it has finite length, as $A/(a,b)$ has.

\emph{The formula.} If $a$ is a unit, all three terms are $0$. Otherwise $A/aA$ is a noetherian local ring of dimension $1$
\cite[Tag 00KW]{Sta}. The multiplicity $e(A/aA,0,b)=\ell(A/(a,b))-\ell(\ker_{A/a}(b))$ of the module $A/aA$ over this ring
is the sum, over the primes $\mathfrak q$ of $A/aA$ with $\dim(A/aA)/\mathfrak q=1$, of $\ell\big((A/aA)_{\mathfrak q}\big)$
times the length of $(A/aA)/(\mathfrak q+bA)$ \cite[Tag 02QF]{Sta}. These $\mathfrak q$ are the $\mathfrak p/aA$ for the
primes $\mathfrak p\ni a$ of $A$ with $\dim A/\mathfrak p=1$, and these are the primes of height one that contain $a$: such a
$\mathfrak p$ has height at least one as $a\neq0$, $\operatorname{ht}\mathfrak p+\dim A/\mathfrak p\le2$, and the maximal
ideal has height $2$. For them $(A/aA)_{\mathfrak q}=A_{\mathfrak p}/aA_{\mathfrak p}$ and
$(A/aA)/(\mathfrak q+bA)=A/(\mathfrak p+bA)$, which gives the formula. Exchanging $a$ and $b$ and subtracting, with the
isomorphism of the kernels and their finite length, gives the last identity.
\end{proof}
% Lean: length_two_generator_quotient: length_quotient_pair_formula (dimOneFormula = LocalLengthAdditivity, proved as
% length_coker_eq_ker_add_dimOneOrderSum by induction on a prime filtration of the module with euler_ses, the Lean's own
% proof of Stacks Tag 02QF for a one-dimensional noetherian local ring; length_quotient_cokernel, length_quotient_kernel,
% length_quotient_sum_ord with height_one_iff_dim_one, length_loc_quot and ordMod_quot), length_ker_swap (the module S of
% the (x, y) with b x = a y, length_eq_of_surj_of_ker_eq) and length_ker_ne_top (length_ne_top_of_killed). ordLoc p a and
% ordMod p b are Ring.ord of A_p at a and of A/p at b. Size: LengthTwoGeneratorQuotient.lean,
% LocalLengthAdditivityTorsionFree.lean, LengthCokerEqKerAddDimOneOrderSum.lean, LengthEqOfSurjOfKerEq.lean,
% LengthEqRangeAddRange.lean and HeightOneIffDimOne.lean, about 1,170 lines.

\begin{lemma}\label{cv-origin-intersection-family-sections-cycles}
Let $K$, $(I_k)_{k\in K}$ and $R$ be as in Definition~\ref{lit-multiprojective}, let $X$ be an arithmetic subvariety
defined by $P$, and let $s,s'\in R\setminus P$ be multihomogeneous, such that no relevant minimal prime of $P+(s)$ is a
relevant minimal prime of $P+(s')$. For a prime $W$ and $f\in R$ and a prime $U$, let $m_U(W,f)$ be the multiplicity of
$U$ in $R/(W+(f))$ if $U$ is a relevant minimal prime of $W+(f)$, and $0$ otherwise. Then for every prime $U$,
\[
 \sum_Wm_W(P,s)\,m_U(W,s')=\sum_{W'}m_{W'}(P,s')\,m_U(W',s),
\]
the sums over the relevant minimal primes $W$ of $P+(s)$ and $W'$ of $P+(s')$.
\end{lemma}
\leanref{GroundTheory/HeightsAndTheta/Heights/Intersections/StatementsSwap.lean}{GroundTheory.Heights₂.Heights.Inter.SwapCycles.Statement}
\leanref{GroundTheory/HeightsAndTheta/Heights/Intersections/PSwapCycles.lean}{GroundTheory.Heights₂.Heights.Inter.SwapCycles.full}
\begin{proof}
Let $d$ be the dimension of $X$, so that $\dim R/P=d+|K|$. If $W$ is a relevant minimal prime of $P+(s)$, then
$s'\notin W$: otherwise $W$ contains $P+(s')$, and it is minimal over it, as a prime containing $P+(s')$ is not $P$ and
no prime lies strictly between $P$ and $W$ \cite[Tag 00KV]{Sta}; so $W$ would be a relevant minimal prime of $P+(s')$.
Likewise $s$ lies in no relevant minimal prime of $P+(s')$. We use
Lemma~\ref{cv-origin-intersection-family-curves-vertical-closed}(1) and (2), which hold for all primes of $R$.

\emph{A nonzero term.} Suppose that $m_W(P,s)\,m_U(W,s')\neq0$, so that $W$ is a relevant minimal prime of $P+(s)$ and
$U$ one of $W+(s')$. Both are multihomogeneous \cite[Tag 00JU]{Sta}, as in
Lemma~\ref{cv-origin-intersection-family-curves-vertical-closed}(4), and by (1) of that lemma, applied twice,
$\dim R/U=d+|K|-2$; by (2), $d\ge2$. Moreover $U$ is a minimal prime of $P+(s)+(s')$. Indeed, let $Q$ be a prime with
$P+(s)+(s')\subset Q\subsetneq U$, choose a minimal prime $W'$ of $P+(s)$ inside $Q$, which is relevant as $U$ is, so that
$s'\notin W'$, and choose a minimal prime $Q'$ of $W'+(s')$ inside $Q$. By (1), applied twice, $\dim R/Q'=d+|K|-2=\dim
R/U$, while $Q'\subset Q\subsetneq U$ gives $\dim R/Q'\ge\dim R/U+1$ by (1). The same holds with $s$ and $s'$ exchanged.
So both sides vanish if $d\le1$, and both vanish if $U$ is not a relevant minimal prime of $P+(s)+(s')$.

\emph{The local ring.} Let $d\ge2$, and let $U$ be a relevant minimal prime of $P+(s)+(s')$. Let $A=(R/P)_U$, the
localization of $R/P$ at the image of $U$, a noetherian local domain. Its primes are the $QA$ for the primes $Q$ with
$P\subset Q\subset U$, and $QA$ has inverse image $Q$ in $R$. As $U$ is minimal over $P+(s)+(s')$, the maximal ideal
$\mathfrak m_A$ is minimal over $(s,s')$; so $A/(s,s')$ has finite length, and $\mathfrak m_A$ has height at most $2$
\cite[Tag 0BBZ]{Sta}. Let $W$ be a minimal prime of $P+(s)$ inside $U$. It is relevant, so $s'\notin W$, and
$0\subsetneq WA\subsetneq\mathfrak m_A$, as $s\in WA$ is nonzero and $s'\in\mathfrak m_A\setminus WA$. So $A$ has
dimension $2$, and $s$ and $s'$ are nonzero in $A$.

\emph{The terms.} Let $\mathfrak p=QA$ be a prime of $A$. Then $\mathfrak p$ has height one and contains $s$ if and only if
$Q$ is a minimal prime of $P+(s)$. Indeed, if $\mathfrak p$ has height one and contains $s$, and $Q_1$ is a prime with
$P+(s)\subset Q_1\subset Q$, then $Q_1A\subset\mathfrak p$ is a nonzero prime, so $Q_1A=\mathfrak p$ and $Q_1=Q$.
Conversely, if $Q$ is minimal over $P+(s)$, then $\mathfrak p$ is minimal over $sA$, so its height is at most one
\cite[Tag 00KV]{Sta}, and it is not zero. Such a $Q$ is relevant, as $Q\subset U$. For it, $A_{\mathfrak p}=(R/P)_Q$, so
$m_Q(P,s)$, the length of $(R/(P+(s)))_Q=A_{\mathfrak p}/sA_{\mathfrak p}$, is $\operatorname{ord}_{\mathfrak p}(s)$ in the
notation of Lemma~\ref{cv-origin-intersection-family-sections-cycles-local}. Also $A/\mathfrak p=(R/Q)_U$, and $U$ is a
relevant minimal prime of $Q+(s')$, as every prime between $Q+(s')$ and $U$ contains $P+(s)+(s')$; so $m_U(Q,s')$, the
length of $(R/(Q+(s')))_U=A/(\mathfrak p+s'A)$, is $\operatorname{ord}_{A/\mathfrak p}(s')$. A relevant minimal prime $W$
of $P+(s)$ that is not contained in $U$ has $m_U(W,s')=0$. Hence
\[
 \sum_Wm_W(P,s)\,m_U(W,s')=\sum_{\mathfrak p\ni s}\operatorname{ord}_{\mathfrak p}(s)\,\operatorname{ord}_{A/\mathfrak p}(s'),
\]
the sum over the primes of $A$ of height one that contain $s$, and likewise with $s$ and $s'$ exchanged. These agree by
Lemma~\ref{cv-origin-intersection-family-sections-cycles-local} for $a=s$ and $b=s'$.
\end{proof}
\leanref{GroundTheory/HeightsAndTheta/Heights/Intersections/PSwapCycles.lean}{GroundTheory.Heights₂.Heights.Inter.SwapCycles.proof}
\leanref{GroundTheory/HeightsAndTheta/Heights/Intersections/PSwapCycles.lean}{GroundTheory.Heights₂.Heights.Inter.SwapCyclesAlt.coeff_eq}
% Lean: SwapCycles.full = SwapCycles.proof length_two_generator_quotient: the case d < 2 through
% divCycle_eq_zero_of_lt_two (dim_quot_minimal_add_one twice and card_le_dim), and for d = m + 2 the coefficientwise
% coeff_eq: sum_eq_zero_of_not with mem_minimalPrimes_of_mem_divSupp_divSupp (not_mem_of_mem_divSupp,
% dim_of_mem_minimalPrimes, DivisorCycleSwap.not_lt_of_dim_eq) when U is not a relevant minimal prime of P + (s) + (s');
% otherwise A = Localization.AtPrime (U.map (mk P)) with isLocData_atPrime_quot and exists_prime_comap_atPrime_quot,
% IsLocData.ringKrullDim_eq_two (Ideal.height_le_card_of_mem_minimalPrimes_span_finset and the chain 0 < p < m),
% IsLocData.maximalIdeal_mem_minimalPrimes with length_quot_ne_top_of_mem_minimalPrimes, sum_eq_finsum
% (IsLocData.height_eq_one_iff, mem_minimalPrimes_sup, cast_multAt, IsLocData.atPrime and IsLocData.quot for the lengths),
% and the third part of LocalLengthMultiplicativity. Size: PSwapCycles.lean and SwapLoc.lean, about 750 lines, besides
% Lemma cv-origin-intersection-family-sections-cycles-local.

\begin{definition}\label{cv-origin-intersection-family-sections-analytic-defect}
Let $K$, $(I_k)_{k\in K}$ and $R$ be as in Definition~\ref{lit-multiprojective}.
\begin{enumerate}
\item A smooth metrized line bundle $\bar M=(m,h_M)$ is \emph{semipositive everywhere} if the Levi form of $\log h_M$
  (Definition~\ref{cv-period-monge-ampere}(1)) is nonnegative at every point of $\hat C$.
\item Let $n\ge0$, let $X$ be an arithmetic subvariety of dimension $n+2$ defined by $P$, let $\bar A=(a,h_A)$ and
  $\bar A'=(a',h_{A'})$ be smooth metrized line bundles, and let $s,s'\in R\setminus P$ be multihomogeneous of
  multidegrees $a$ and $a'$. For smooth metrized line bundles $\bar L_1,\dots,\bar L_n$, let
  $D(\bar L_1,\dots,\bar L_n)$ be the left side minus the right side of the identity of
  Lemma~\ref{cv-origin-intersection-family-sections-analytic}.
\end{enumerate}
\end{definition}
\leanref{GroundTheory/HeightsAndTheta/Heights/Intersections/NoCommonComponentCoordinatesPlane.lean}{GroundTheory.Heights₂.Heights.Inter.IsSemiPosHerm}
\leanref{GroundTheory/HeightsAndTheta/Heights/Intersections/NoCommonComponentCoordinatesPlane.lean}{GroundTheory.Heights₂.Heights.Inter.De}

\begin{lemma}\label{cv-origin-intersection-family-sections-analytic-additive-slot}
Let $K$, $(I_k)_{k\in K}$ and $R$ be as in Definition~\ref{lit-multiprojective}, let $e\ge1$, let $Y$ be an arithmetic
subvariety of dimension $e+1$ defined by the prime $Q$, let $1\le j\le e$, let $\bar L_1,\dots,\bar L_e$, $\bar M$,
$\bar N$ and $\bar B=(b,h_B)$ be smooth metrized line bundles, and let $t\in R\setminus Q$ be multihomogeneous of
multidegree $b$. For a smooth metrized line bundle $\bar X$ let $I(\bar X)$ be the integral
$\int_{Y(\mathbf C)}\log(|t|/h_B)\,c_1(\bar L_1)\wedge\dots\wedge c_1(\bar L_e)$ with $\bar X$ in place of $\bar L_j$.
Then $I(\bar M\otimes\bar N)=I(\bar M)+I(\bar N)$.
\end{lemma}
\leanref{GroundTheory/HeightsAndTheta/Heights/Intersections/PDeSlotAdd.lean}{GroundTheory.Heights₂.Heights.Inter.ptsIntegral_update_tensor}
\begin{proof}
As $Q$ is relevant, some $x_{k,i}$ is not in $Q$. By Lemma~\ref{cv-origin-intersection-resolution-exists}, applied to
the ideal $(x_{k,i})$, $Y$ has an analytic resolution $(\tilde Y,(U_x),(\lambda_x))$ off some ideal
(Definition~\ref{cv-origin-intersection-resolution}). For a smooth metrized line bundle $\bar X=(c,h_X)$ let
$\vartheta_X$ be the local potential $\log h_X\circ\lambda_x$ on $U_x$, for each $x\in\tilde Y$, and let
$g(x)=\log\big(|t(\lambda_x(x))|/h_B(\lambda_x(x))\big)$, with $\log0$ read as $0$. By
Lemma~\ref{cv-origin-intersection-poincare-lelong-all-double-good}(1), $Y$ has a good projection. So by
Lemmas~\ref{cv-origin-intersection-poincare-lelong-all-double-independent}(2)
and~\ref{cv-origin-intersection-poincare-lelong-all-double-bridge}, for smooth metrized line bundles
$\bar X_1,\dots,\bar X_e$ the integral $\int_{Y(\mathbf C)}\log(|t|/h_B)\,c_1(\bar X_1)\wedge\dots\wedge c_1(\bar X_e)$
is the integral over $\tilde Y$ of $g$ against the Monge--Amp\`ere measure of $\vartheta_{X_1},\dots,\vartheta_{X_e}$
(Definition~\ref{cv-period-monge-ampere}(3)).

\emph{The density.} For Hermitian $e\times e$ matrices $H_1,\dots,H_e$, expand each $\det(\sum_{k\in T}H_k)$ in
Definition~\ref{cv-period-monge-ampere}(2) by multilinearity in the columns. The term in which the column $l$ comes from
$H_{f(l)}$, for a map $f:\{1,\dots,e\}\to T$, has total sign $\sum_{T\supset f(\{1,\dots,e\})}(-1)^{e-|T|}$ over the
$T$, which is $1$ if $f$ is onto and $0$ otherwise. So the sum of Definition~\ref{cv-period-monge-ampere}(2) is the sum,
over the permutations $\sigma$ of $\{1,\dots,e\}$, of the determinant of the matrix whose column $l$ is the column $l$ of
$H_{\sigma(l)}$. This is symmetric in $H_1,\dots,H_e$ and additive in each of them. Hence the Monge--Amp\`ere density
is symmetric in its arguments, and
$\operatorname{MA}(\varphi+\varphi',\varphi_2,\dots,\varphi_e)=\operatorname{MA}(\varphi,\varphi_2,\dots,\varphi_e)
+\operatorname{MA}(\varphi',\varphi_2,\dots,\varphi_e)$ at points near which $\varphi$ and $\varphi'$ are twice
continuously differentiable, as their complex Hessians add there.

\emph{The first place.} Let $\bar T_i$ be $\bar L_1$ for $i=j$ and $\bar L_i$ otherwise, for $2\le i\le e$. By the
symmetry, exchanging the places $1$ and $j$, $I(\bar X)$ is the integral over $\tilde Y$ of $g$ against the
Monge--Amp\`ere measure of $\vartheta_X,\vartheta_{T_2},\dots,\vartheta_{T_e}$. Write $\bar M=(m,h_M)$ and
$\bar N=(m',h_N)$, so that $\bar M\otimes\bar N=(m+m',h_Mh_N)$. The values of $\lambda_x$ lie in $\hat C$, where $h_M$
and $h_N$ are positive, so $\vartheta_{M\otimes N}=\vartheta_M+\vartheta_N$ on each $U_x$. These potentials are smooth in
every chart, as the metrics are smooth on $\hat C$ and $\lambda_x$ is holomorphic. So in every chart the density for
$\vartheta_{M\otimes N},\vartheta_{T_2},\dots,\vartheta_{T_e}$ is the sum of those for $\vartheta_M$ and for
$\vartheta_N$ in the first place.

\emph{The integrals.} Let $\kappa_k:V_k\to\mathbf C^e$ and $\chi_k$ be the charts and functions of
Definition~\ref{cv-period-monge-ampere}(3). The first paragraph of the proof of
Lemma~\ref{cv-origin-intersection-poincare-lelong-all-double-bridge} shows that, for any smooth metrized line bundles
in the places, $g\circ\kappa_k^{-1}$ times the Monge--Amp\`ere density is integrable on every compact subset of
$\kappa_k(V_k)$. As $\chi_k$ is continuous with compact support in $V_k$, the integrand of the term $k$ of
Definition~\ref{cv-period-monge-ampere}(3) is integrable, for $\vartheta_M$ and for $\vartheta_N$ in the first place. So
each term for $\vartheta_{M\otimes N}$ is the sum of those for $\vartheta_M$ and for $\vartheta_N$, and summing over $k$
gives $I(\bar M\otimes\bar N)=I(\bar M)+I(\bar N)$.
\end{proof}
\leanref{ReciprocityTheory/AnalyticForms/PotentialTheory/PotentialContDiffAt.lean}{ReciprocityTheory.AnForm.M.mIntSymm}
\leanref{ReciprocityTheory/AnalyticForms/PotentialTheory/PotentialContDiffAt.lean}{ReciprocityTheory.AnForm.M.mIntAdd}
\leanref{ReciprocityTheory/AnalyticForms/PlurisubharmonicFunctions/MdiscSt.lean}{ReciprocityTheory.AnForm.A.mdisc_eq_sum_perm}
\leanref{ReciprocityTheory/ReductiveGroups/ArithmeticGeometry/BirationalGeometry/InterRed.lean}{ReciprocityTheory.ChartRestriction.G.Geo.ptsIntegral_eq_mInt}
% Lean: ptsIntegral_update_tensor: exists_anLogRes (Lemma cv-origin-intersection-resolution-exists with J = (X k i)),
% key from ptsIntegral_eq_mInt' (goodProjExists, ptsIntegral_eq_at and resBridge, Lemmas -double-independent and
% -double-bridge) and mIntSymm with Equiv.swap 0 j (mdisc_comp_perm_potential, from mdisc_eq_sum_perm and
% mdisc_comp_perm); e1 (germ_of_eqOn, Real.log_mul with the positivity of the metrics on R.mem); hI from brIntegrable
% (= brIntegrable_of logAbsLocInt, the first paragraph of the proof of Lemma -double-bridge); mIntAdd (integrableOn_weight
% for the compact support of the partition, locDens_cons_add with hessC_add' and mdisc_cons_add = mdisc_update_add, the
% multilinearity of mdiscML). Size: PDeSlotAdd.lean (ptsIntegral_update_tensor, about 60 lines), PotentialContDiffAt.lean
% (166 lines), MdiscSt.lean (part, about 110 lines) and InterRed.lean, InterRedE.lean and InterRedF.lean (about 330
% lines), besides Lemmas cv-origin-intersection-resolution-exists and -double-bridge.

\begin{lemma}\label{cv-origin-intersection-family-sections-analytic-additive}
Let $K$, $(I_k)_{k\in K}$ and $R$ be as in Definition~\ref{lit-multiprojective}, and let $n$, $X$, $P$, $\bar A$,
$\bar A'$, $s$, $s'$ and $D$ be as in Definition~\ref{cv-origin-intersection-family-sections-analytic-defect}(2), such
that no horizontal relevant minimal prime of $P+(s)$ is one of $P+(s')$; this hypothesis is not used in the proof. Then
$D$ is additive in each argument: for $1\le j\le n$ and smooth metrized line bundles $\bar L_i$ ($i\neq j$), $\bar M$ and
$\bar N$, the value of $D$ with $\bar M\otimes\bar N$ in the place $j$ is the sum of its values with $\bar M$ and with
$\bar N$ there.
\end{lemma}
\leanref{GroundTheory/HeightsAndTheta/Heights/Intersections/NoCommonComponentCoordinatesPlane.lean}{GroundTheory.Heights₂.Heights.Inter.DeSlotAdd.Statement}
\leanref{GroundTheory/HeightsAndTheta/Heights/Intersections/PDeSlotAdd.lean}{GroundTheory.Heights₂.Heights.Inter.DeSlotAdd.proof}
\begin{proof}
By definition, $D$ is a combination of four kinds of integrals: over $W(\mathbf C)$ for the horizontal relevant minimal
primes $W$ of $P+(s)$, of $\log(|s'|/h_{A'})$ against $c_1(\bar L_1)\wedge\dots\wedge c_1(\bar L_n)$; the same with the roles
of $(\bar A,s)$ and $(\bar A',s')$ exchanged; and over $X(\mathbf C)$, of $\log(|s|/h_A)$ against
$c_1(\bar A')\wedge c_1(\bar L_1)\wedge\dots\wedge c_1(\bar L_n)$, and the same exchanged. Each such $W$ is an arithmetic
subvariety of dimension $n+1$, as for the components in Definition~\ref{lit-multiprojective}(3). If $s'\in W$, then
$\log(|s'|/h_{A'})$, with $\log0$ read as $0$ as in the sheet densities, vanishes on $\hat W(\mathbf C)$, and the three
integrals over $W(\mathbf C)$ are $0$; likewise with $s$ and $s'$ exchanged. So it suffices that, for an arithmetic
subvariety $Y$ of
dimension $e+1$, smooth metrized line bundles $\bar B=(b,h_B)$ and $\bar L_1,\dots,\bar L_e$, and a multihomogeneous
$t\notin Y$ of multidegree $b$, the integral of $\log(|t|/h_B)$ over $Y(\mathbf C)$ against
$c_1(\bar L_1)\wedge\dots\wedge c_1(\bar L_e)$ is additive in each $\bar L_j$. This is
Lemma~\ref{cv-origin-intersection-family-sections-analytic-additive-slot}: for the integrals over the $W(\mathbf C)$ with
$e=n$ and the place $j$, and for those over $X(\mathbf C)$ with $e=n+1$, the bundles $\bar A'$ (or $\bar A$),
$\bar L_1,\dots,\bar L_n$ and the place $j+1$.
\end{proof}
% Lean: DeSlotAdd.of_remaining = DeSlotAdd.proof (resExists_of_hironaka): hZ for the integrals over P with the first form
% log h of A or A' (Fin.cons_update and the place j.succ), hW for the sums over the horizontal components (divSuppArith;
% the case t' in W through ptsIntegral_eq_zero_of_vanish, as Real.log 0 = 0), both from
% ptsIntegral_update_tensor (Lemma cv-origin-intersection-family-sections-analytic-additive-slot). The sheet-density
% argument of Lemma cv-origin-intersection-family-sections-position, which the Lean uses for the last place
% (arInt_snoc_tensor), would also give it. Size: PDeSlotAdd.lean, about 100 lines besides ptsIntegral_update_tensor.

\begin{lemma}\label{cv-origin-intersection-family-sections-analytic-twist}
Let $K$, $(I_k)_{k\in K}$ and $R$ be as in Definition~\ref{lit-multiprojective}, and let $\bar M=(m,h_M)$ be a smooth
metrized line bundle. There is an integer $N\ge0$ such that the Fubini--Study bundle $\bar T=\bar{\mathcal O}(N\mathbf 1)$
(Definition~\ref{cv-origin-analytic-system-pullback}(1)) and $\bar M\otimes\bar T$ are semipositive everywhere
(Definition~\ref{cv-origin-intersection-family-sections-analytic-defect}(1)).
\end{lemma}
\leanref{GroundTheory/HeightsAndTheta/Heights/Intersections/NoCommonComponentCoordinatesPlane.lean}{GroundTheory.Heights₂.Heights.Inter.SemiPosTwist.Statement}
\leanref{GroundTheory/HeightsAndTheta/Heights/Intersections/PSemiPosTwist.lean}{GroundTheory.Heights₂.Heights.Inter.SemiPosTwist.proof}
\leanref{GroundTheory/HeightsAndTheta/Heights/Intersections/PTorusInvLevi.lean}{GroundTheory.Heights₂.Heights.Inter.TorusInvLevi.proof}
\begin{proof}
Write $\mathcal L_\phi(z,w)$ for the Levi form of a function $\phi$ at $z$ in the direction $w$. For $k\in K$ let
$F_k(z)=\log\|z_k\|$ on $\hat C$, with $\|z_k\|$ as in Definition~\ref{cv-origin-analytic-system-pullback}. The function
$\log\|\cdot\|$ on $\mathbf C^{I_k}\setminus\{0\}$ is plurisubharmonic, so $\mathcal L_{F_k}(z,w)\ge0$. Put
$u=\log h_M-\sum_km_kF_k$. As $h_M(cz)=\prod_k|c_k|^{m_k}h_M(z)$ and $F_k(cz)=\log|c_k|+F_k(z)$ for $c\in(\mathbf C^\times)^K$
acting blockwise, $u$ is invariant under $(\mathbf C^\times)^K$, and it is twice continuously differentiable on $\hat C$.

\emph{A bound for the Levi form of $u$.} Let $Q\subset\hat C$ be the compact set of the $z$ with $|z_{k,i}|\le1$ for all
$k,i$ and with some $z_{k,i}=1$ for each $k$. By continuity there is $M_0\ge0$ with $|\mathcal L_u(x,w)|\le
M_0\sum_{k,i}|w_{k,i}|^2$ for $x\in Q$. Let $z\in\hat C$, choose for each $k$ an index $i_k$ with $|z_{k,i_k}|$ largest,
and let $\tau(y)_{k,i}=y_{k,i}/y_{k,i_k}$, a holomorphic map near $z$ with $\tau(z)\in Q$. By the invariance, $u=u\circ\tau$
near $z$, so $\mathcal L_u(z,w)=\mathcal L_u(\tau(z),d\tau(z)w)$. The entries of $d\tau(z)w$ in the block $k$ are
$(z_{k,i_k}w_{k,l}-z_{k,l}w_{k,i_k})/z_{k,i_k}^2$, and a computation, using $|z_{k,l}|\le|z_{k,i_k}|$, bounds the sum of
their squared absolute values by $4|I_k|^2\mathcal L_{F_k}(z,w)$. Hence
$|\mathcal L_u(z,w)|\le C\sum_k\mathcal L_{F_k}(z,w)$ with $C=4M_0\sum_k|I_k|^2$.

\emph{The twist.} Let $N$ be an integer with $N\ge|C|+\sum_k|m_k|$, so that $m_k+N\ge C$ for every $k$. The metric of
$\bar T$ has $\log h_T=N\sum_kF_k$, whose Levi form $N\sum_k\mathcal L_{F_k}$ is nonnegative. The metric of
$\bar M\otimes\bar T$ has $\log(h_Mh_T)=u+\sum_k(m_k+N)F_k$, whose Levi form is at least
$-C\sum_k\mathcal L_{F_k}+\sum_k(m_k+N)\mathcal L_{F_k}=\sum_k(m_k+N-C)\mathcal L_{F_k}\ge0$.
\end{proof}
\leanref{GroundTheory/HeightsAndTheta/Heights/Intersections/PTorusInvLevi.lean}{GroundTheory.Heights₂.Heights.Inter.TorusInvariantLevi.chart_bound}
% Lean: SemiPosTwist.full = SemiPosTwist.proof Closed.Inter.TorusInvLevi: u, hinv, hC from TorusInvLevi.proof (Q the
% normalized points, levi_bound_compact, tau with hev and levi_comp, hD for the derivative, chart_bound, the computation
% with leviForm_phi_nonneg and sum_cross_eq_levi, and levi_Phi), N = ceil(|C| + sum |M.a k|), T = sqMetric (fun _ => N),
% and the Laplacian along each complex line through laplacian_line_eq_levi and levi_add_sum. IsSemiPosHerm asks the
% Laplacian of t |-> log h(z + t v) at 0 to be nonnegative, which is four times the Levi form. Size: PSemiPosTwist.lean and
% PTorusInvLevi.lean, about 540 lines.

\begin{definition}\label{cv-origin-intersection-family-sections-analytic-semipositive-regularized}
Let $K$, $(I_k)_{k\in K}$ and $R$ be as in Definition~\ref{lit-multiprojective}, let $\bar E=(a,h_E)$ be a smooth
metrized line bundle, and let $t\in R$ be multihomogeneous of multidegree $a$. For $\varepsilon>0$, the
\emph{regularization} $\bar E_{t,\varepsilon}=(0,h_{t,\varepsilon})$ is the metrized line bundle of multidegree $0$ with
\[
 \log h_{t,\varepsilon}(z)=\tfrac12\log\Big(\frac{|t(z)|^2}{h_E(z)^2}+\varepsilon\Big),\qquad z\in\hat C.
\]
\end{definition}
\leanref{GroundTheory/HeightsAndTheta/Heights/Intersections/RegLine.lean}{GroundTheory.Heights₂.Heights.Inter.regLine}
\leanref{GroundTheory/HeightsAndTheta/Heights/Intersections/RegLine.lean}{GroundTheory.Heights₂.Heights.Inter.isSmoothHerm_regLine}
\leanref{GroundTheory/HeightsAndTheta/Heights/Intersections/RegLine.lean}{GroundTheory.Heights₂.Heights.Inter.logNorm_le_log_regLine_h}
% The function |t|/h_E is invariant under (C^x)^K, as t and h_E have the same multidegree, so h_{t,eps} has
% multidegree 0; it is smooth on the cone as the argument of the logarithm is at least eps (isSmoothHerm_regLine). Where
% t(z) is nonzero, log(|t(z)|/h_E(z)) <= log h_{t,eps}(z) (logNorm_le_log_regLine_h). Size: RegLine.lean, 113 lines.

\begin{lemma}\label{cv-origin-intersection-family-sections-analytic-semipositive-first}
Let $K$, $(I_k)_{k\in K}$ and $R$ be as in Definition~\ref{lit-multiprojective}, let $e\ge0$, let $Y$ be an arithmetic
subvariety of dimension $e+2$ defined by $Q$, let $\bar L_1,\dots,\bar L_e$, $\bar M$, $\bar N$ and $\bar E=(b,h_E)$ be
smooth metrized line bundles, and let $t\in R\setminus Q$ be multihomogeneous of multidegree $b$. Put
$\Phi=c_1(\bar L_1)\wedge\dots\wedge c_1(\bar L_e)$. Then
\[
 \int_{Y(\mathbf C)}\log\frac{|t|}{h_E}\,c_1(\bar M\otimes\bar N)\wedge\Phi
 =\int_{Y(\mathbf C)}\log\frac{|t|}{h_E}\,c_1(\bar M)\wedge\Phi+\int_{Y(\mathbf C)}\log\frac{|t|}{h_E}\,c_1(\bar N)\wedge\Phi,
\]
and the same holds with $\log h_D$ in place of $\log(|t|/h_E)$, for a smooth metrized line bundle $\bar D=(0,h_D)$ of
multidegree $0$.
\end{lemma}
\leanref{GroundTheory/HeightsAndTheta/Heights/Intersections/StPtsAddFirst.lean}{GroundTheory.Heights₂.Heights.Inter.PtsAddFirst.Statement}
\leanref{GroundTheory/HeightsAndTheta/Heights/Intersections/PPtsAddFirst.lean}{GroundTheory.Heights₂.Heights.Inter.PtsAddFirst.proof}
\leanref{GroundTheory/HeightsAndTheta/Heights/Intersections/PDeLeSemiPos.lean}{GroundTheory.Heights₂.Heights.Inter.ptsIntegral_add_first_log}
\begin{proof}
The argument for the integrals in the proof of Lemma~\ref{cv-origin-intersection-family-sections-position}, with the
first place in place of the last, gives the first identity: through the good projection chosen for $Y$, the sheet
densities add, as $\log(h_Mh_N)\circ\gamma=\log h_M\circ\gamma+\log h_N\circ\gamma$ near $y$ for a sheet $\gamma$ and the
Monge--Amp\`ere density is additive in each argument; the fibres are finite; and the two fibre sums on the right are
integrable by Lemma~\ref{cv-origin-intersection-poincare-lelong-all-double-independent}(1). If $Y$ has no good projection,
the three integrals are $0$. For the second identity, $\log h_D=-\log(|1|/h_D)$, where $1$ is multihomogeneous of
multidegree $0$ and not in $Q$; apply the first identity to $\bar D$ and $t=1$.
\end{proof}
% Lean: PtsAddFirst.proof (under PtsIntegrableStatement): sheetDensity_add (maDensity_cons_add, mixedDisc_cons_add,
% cxHessian_add), fibreFinite and integral_add; ptsIntegral_add_first_log through log_h_eq_neg_logNorm_one,
% isSectionOn_one and ptsIntegral_neg. Size: PPtsAddFirst.lean, 108 lines, and about 15 lines in PDeLeSemiPos.lean.

\begin{lemma}\label{cv-origin-intersection-family-sections-analytic-semipositive-curvature}
Let $K$, $(I_k)_{k\in K}$ and $R$ be as in Definition~\ref{lit-multiprojective}, let $\bar E=(a,h_E)$ be a smooth
metrized line bundle, and let $t\in R$ be multihomogeneous of multidegree $a$. There is an integer $N\ge0$ such that for
every $\varepsilon>0$ the metrized line bundle $\bar E\otimes\bar E_{t,\varepsilon}\otimes\bar{\mathcal O}(N\mathbf 1)$
(Definitions~\ref{cv-origin-intersection-family-sections-analytic-semipositive-regularized}
and~\ref{cv-origin-analytic-system-pullback}(1)) is semipositive everywhere
(Definition~\ref{cv-origin-intersection-family-sections-analytic-defect}(1)).
\end{lemma}
\leanref{GroundTheory/HeightsAndTheta/Heights/Intersections/StRegCurvBound.lean}{GroundTheory.Heights₂.Heights.Inter.RegCurvBound.Statement}
\leanref{GroundTheory/HeightsAndTheta/Heights/Intersections/PRegCurvBound.lean}{GroundTheory.Heights₂.Heights.Inter.RegCurvBound.proof}
\leanref{GroundTheory/HeightsAndTheta/Heights/Intersections/RegularizedCurvature.lean}{GroundTheory.Heights₂.Heights.Inter.RegularizedCurvature.levi_logSum}
\begin{proof}
Write $\mathcal L_\phi(z,w)$ for the Levi form, $F_k(z)=\log\|z_k\|$ and $\psi=\log h_E$. As in the proof of
Lemma~\ref{cv-origin-intersection-family-sections-analytic-twist}, $u=\psi-\sum_ka_kF_k$ is invariant under
$(\mathbf C^\times)^K$, and there is $C$ with $|\mathcal L_u(z,w)|\le C\sum_k\mathcal L_{F_k}(z,w)$, where
$\mathcal L_{F_k}\ge0$. Let $N$ be an integer with $N\ge|C|+\sum_k|a_k|$. Then
$\mathcal L_\psi+N\sum_k\mathcal L_{F_k}=\mathcal L_u+\sum_k(a_k+N)\mathcal L_{F_k}\ge\sum_k(a_k+N-C)\mathcal L_{F_k}\ge0$.

Let $\varepsilon>0$, and put $\Phi_\varepsilon=\psi+\log h_{t,\varepsilon}=\psi+\frac12\log(|t|^2e^{-2\psi}+\varepsilon)$.
Fix $z\in\hat C$ and $w$. We show that $\mathcal L_{\Phi_\varepsilon}(z,w)\ge\mu\,\mathcal L_\psi(z,w)$ for some
$0\le\mu\le1$. If $t(z)=0$, then $\Phi_\varepsilon\ge\psi+\frac12\log\varepsilon$ near $z$, with equality at $z$, so the
difference has a minimum at $z$ and $\mathcal L_{\Phi_\varepsilon}(z,w)\ge\mathcal L_\psi(z,w)$; take $\mu=1$. If
$t(z)\neq0$, let $x_0=|t(z)|^2e^{-2\psi(z)}$, $\lambda=x_0/(x_0+\varepsilon)$ and $\mu=\varepsilon/(x_0+\varepsilon)$, so
that $\lambda+\mu=1$. Near $z$, with $x=|t|^2e^{-2\psi}$,
\[
 \Phi_\varepsilon-\mu\psi-\lambda\log|t|=\tfrac12\big(\log(x+\varepsilon)-\lambda\log x\big),
\]
and the function $x\mapsto\log(x+\varepsilon)-\lambda\log x$ on $x>0$ has its minimum at $x=x_0$. So
$\Phi_\varepsilon-\mu\psi-\lambda\log|t|$ has a minimum at $z$, and as $\log|t|$ is pluriharmonic near $z$,
$\mathcal L_{\Phi_\varepsilon}(z,w)\ge\mu\,\mathcal L_\psi(z,w)$.

The metric of $\bar E\otimes\bar E_{t,\varepsilon}\otimes\bar{\mathcal O}(N\mathbf 1)$ has logarithm
$\Phi_\varepsilon+N\sum_kF_k$, whose Levi form at $(z,w)$ is at least $\mu\mathcal L_\psi+N\sum_k\mathcal L_{F_k}$. If
$\mathcal L_\psi(z,w)\ge0$ this is $\ge0$; otherwise it is at least $\mathcal L_\psi+N\sum_k\mathcal L_{F_k}\ge0$, as
$\mu\le1$.
\end{proof}
% Lean: RegCurvBound.proof (TorusInvLevi.proof, as in Lemma -twist, for u = log h_E - sum a_k F_k; N = ceil(|C| + sum
% |a_k|); hLW; levi_logSum, with levi_le_of_le at a zero of t and log_norm_pluriharmonic and the minimum of
% log(x + eps) - l log x otherwise; laplacian_line_eq_levi for IsSemiPosHerm). Size: PRegCurvBound.lean (122 lines) and
% RegularizedCurvature.lean (310 lines), besides PTorusInvLevi.lean.

\begin{lemma}\label{cv-origin-intersection-family-sections-analytic-semipositive-absolute}
Let $K$, $(I_k)_{k\in K}$ and $R$ be as in Definition~\ref{lit-multiprojective}, let $e\ge0$, let $Y$ be an arithmetic
subvariety of dimension $e+1$ defined by $Q$, let $\Psi$ be a good projection for $Y$, let $\bar L_1,\dots,\bar L_e$ and
$\bar E=(b,h_E)$ be smooth metrized line bundles, and let $t\in R\setminus Q$ be multihomogeneous of multidegree $b$.
Then the function of $y\in\mathbf C^e$ that sums, over $w\in F_\Psi(y)$, the absolute values of the sheet densities at
$(y,w)$ of $\log(|t|/h_E)$ against $\log h_1,\dots,\log h_e$
(Definition~\ref{cv-origin-intersection-poincare-lelong-all-double-projection}(3)) is integrable.
\end{lemma}
\leanref{GroundTheory/HeightsAndTheta/Heights/Intersections/StPtsAbsIntegrable.lean}{GroundTheory.Heights₂.Heights.Inter.PtsAbsIntegrable.Statement}
\leanref{GroundTheory/HeightsAndTheta/Heights/Intersections/PPtsAbsIntegrable.lean}{GroundTheory.Heights₂.Heights.Inter.PtsAbsIntegrable.proof}
\begin{proof}
\openstep{The proof is not yet written out. It twists each $\bar L_j$ by the Fubini--Study bundle of
Lemma~\ref{cv-origin-intersection-family-sections-analytic-twist} to bound the absolute value of each sheet density by a
combination of sheet densities of semipositive bundles, with integrands $\log(|t|/h_E)$ and $\log h_{t,1}$, whose fibre
sums are integrable by Lemma~\ref{cv-origin-intersection-poincare-lelong-all-double-independent}(1), and it proves the
measurability through the resolution of Lemma~\ref{cv-origin-intersection-resolution-exists}.}
% Lean: PtsAbsIntegrable.proof (SemiPosTwist.proof, abs_maDensity_le with the bound 3^e, regLine A s 1 for D,
% PtsIntegrableStatement twice) and PtsAbsMeasurable.of_remaining (under FinTypeQE and GFLogStC). Size
% (tools/lean_size.py PtsAbsIntegrable.of_remaining): 5,741 lines.
\end{proof}

\begin{lemma}\label{cv-origin-intersection-family-sections-analytic-semipositive-limit}
Let $K$, $(I_k)_{k\in K}$ and $R$ be as in Definition~\ref{lit-multiprojective}, let $e\ge0$, let $Y$ be an arithmetic
subvariety of dimension $e+1$ defined by $Q$, let $\bar M_1,\dots,\bar M_e$ and $\bar E=(a,h_E)$ be smooth metrized line
bundles, and let $t\in R\setminus Q$ be multihomogeneous of multidegree $a$. Put
$\Phi=c_1(\bar M_1)\wedge\dots\wedge c_1(\bar M_e)$. Then, as $\varepsilon\to0$ with $\varepsilon>0$,
\[
 \int_{Y(\mathbf C)}\log h_{t,\varepsilon}\,\Phi\to\int_{Y(\mathbf C)}\log\frac{|t|}{h_E}\,\Phi.
\]
\end{lemma}
\leanref{GroundTheory/HeightsAndTheta/Heights/Intersections/StPtsRegLimit.lean}{GroundTheory.Heights₂.Heights.Inter.PtsRegLimit.Statement}
\leanref{GroundTheory/HeightsAndTheta/Heights/Intersections/PPtsRegLimit.lean}{GroundTheory.Heights₂.Heights.Inter.PtsRegLimit.proof}
\begin{proof}
If $Y$ has no good projection, all the integrals are $0$. Otherwise let $\Psi$ be the good projection chosen for $Y$, and
write $S_g(y,w)$ for the sheet density at $(y,w)$ of a function $g$ against $\log h_{M_1},\dots,\log h_{M_e}$. As in the
proof of Lemma~\ref{cv-origin-intersection-family-product}, for almost every $y$, $t$ vanishes at no
$z\in\hat Y(\mathbf C)$ with $\Psi(\sigma(z))=(1,y)$. Let $y$ be such a point and $w\in F_\Psi(y)$ with a chosen sheet
$\gamma$, and put $x=|t(\gamma(y))|/h_E(\gamma(y))\neq0$. Then $\log h_{t,\varepsilon}(\gamma(y))=\frac12\log(x^2+\varepsilon)$
tends to $\log x$, so $S_{\log h_{t,\varepsilon}}(y,w)\to S_{\log(|t|/h_E)}(y,w)$, and the fibre $F_\Psi(y)$ is finite
(Lemma~\ref{cv-origin-intersection-poincare-lelong-all-double-good}(2)). For $0<\varepsilon\le1$,
$|\frac12\log(x^2+\varepsilon)|\le|\log x|+\frac12\log(x^2+1)$: if $x^2+\varepsilon\ge1$ it lies between $0$ and
$\frac12\log(x^2+1)$, and otherwise between $\log x$ and $0$. So
$|S_{\log h_{t,\varepsilon}}|\le|S_{\log(|t|/h_E)}|+|S_{\log h_{t,1}}|$ there. The fibre sums of both are integrable by
Lemma~\ref{cv-origin-intersection-family-sections-analytic-semipositive-absolute}, for $\bar E$ and $t$, and for
$\bar E_{t,1}$ and $1$, as $\log h_{t,1}=-\log(|1|/h_{t,1})$. Dominated convergence gives the statement.
\end{proof}
% Lean: PtsRegLimit.proof (under PtsIntegrableStatement and PtsAbsIntegrable): projB_null for the ideal (t), hpt,
% abs_half_log_le and tendsto_half_log, hsdlog for log h = -logNorm D 1, and tendsto_integral_filter_of_dominated_convergence.
% Size: PPtsRegLimit.lean, 131 lines, besides Lemma -absolute.

\begin{lemma}\label{cv-origin-intersection-family-sections-analytic-semipositive-mono}
Let $K$, $(I_k)_{k\in K}$ and $R$ be as in Definition~\ref{lit-multiprojective}, let $e\ge0$, let $Y$ be an arithmetic
subvariety of dimension $e+1$ defined by $Q$, let $\bar M_1,\dots,\bar M_e$ be smooth metrized line bundles that are
semipositive everywhere, let $\bar E=(a,h_E)$ and $\bar D=(0,h_D)$ be smooth metrized line bundles, and let
$t\in R\setminus Q$ be multihomogeneous of multidegree $a$, such that $\log(|t(z)|/h_E(z))\le\log h_D(z)$ for every
$z\in\hat Y(\mathbf C)$ with $t(z)\neq0$. Put $\Phi=c_1(\bar M_1)\wedge\dots\wedge c_1(\bar M_e)$. Then
\[
 \int_{Y(\mathbf C)}\log\frac{|t|}{h_E}\,\Phi\le\int_{Y(\mathbf C)}\log h_D\,\Phi.
\]
\end{lemma}
\leanref{GroundTheory/HeightsAndTheta/Heights/Intersections/StPtsMonoSemiPos.lean}{GroundTheory.Heights₂.Heights.Inter.PtsMonoSemiPos.Statement}
\leanref{GroundTheory/HeightsAndTheta/Heights/Intersections/PPtsMonoSemiPos.lean}{GroundTheory.Heights₂.Heights.Inter.PtsMonoSemiPos.proof}
\leanref{ReciprocityTheory/AnalyticForms/PlurisubharmonicFunctions/MdiscSt.lean}{ReciprocityTheory.AnForm.A.maDens_nonneg_of_posSemidef}
\begin{proof}
If $Y$ has no good projection, both sides are $0$. Otherwise let $\Psi$ be the good projection chosen for $Y$. Both fibre
sums are integrable by Lemma~\ref{cv-origin-intersection-poincare-lelong-all-double-independent}(1), the second as
$\log h_D=-\log(|1|/h_D)$. As in the proof of Lemma~\ref{cv-origin-intersection-family-product}, for almost every $y$,
$t$ vanishes at no $z\in\hat Y(\mathbf C)$ with $\Psi(\sigma(z))=(1,y)$; it suffices that for such $y$ every sheet density
of $\log(|t|/h_E)$ is at most that of $\log h_D$. Let $\gamma$ be the sheet chosen at $(y,w)$. Then $t(\gamma(y))\neq0$,
so $\log(|t|/h_E)(\gamma(y))\le\log h_D(\gamma(y))$, and it remains that
$\operatorname{MA}(\log h_{M_1}\circ\gamma,\dots,\log h_{M_e}\circ\gamma)(y)\ge0$.

The complex Hessian of $\log h_{M_j}\circ\gamma$ at $y$ is positive semidefinite: as $\gamma$ is holomorphic, its Levi
form in the direction $v$ is that of $\log h_{M_j}$ at $\gamma(y)$ in the direction $d\gamma(y)v$, which is $\ge0$. For
positive semidefinite $H_1,\dots,H_e$, write $H_j=\sum_lw_{j,l}w_{j,l}^*$ with column vectors $w_{j,l}$. By the
expansion over permutations in the proof of
Lemma~\ref{cv-origin-intersection-family-sections-analytic-additive-slot}, the sum of
Definition~\ref{cv-period-monge-ampere}(2) is additive in each $H_j$. For rank-one matrices $H_j=u_ju_j^*$, with
$u_{j,l}$ the coordinates of $u_j$, the column $l$ of $H_{\sigma(l)}$ is $\overline{u_{\sigma(l),l}}\,u_{\sigma(l)}$, so the
sum is $\sum_\sigma\prod_l\overline{u_{\sigma(l),l}}\operatorname{sgn}(\sigma)\det(u_1,\dots,u_e)
=|\det(u_1,\dots,u_e)|^2\ge0$. Expanding each $H_j$ into its rank-one terms, the density is a sum of nonnegative terms.
\end{proof}
% Lean: PtsMonoSemiPos.proof (under PtsIntegrableStatement): key for log h_D = -logNorm D 1, projB_null for (t),
% integral_mono_ae, and maDens_nonneg_of_posSemidef with SheetHessPSD.full (the Hessian along a sheet). The
% nonnegativity of the mixed discriminant of positive semidefinite matrices is mixedDiscriminant_nonneg_of_gram
% (MdiscPos.lean, through the Gram decompositions pm_exists_gram_of_posSemidef and pm_mdisc_gram). Size:
% PPtsMonoSemiPos.lean (72 lines), MdiscPos.lean (part, about 150) and the SheetHessPSD files, about 400 lines.

\begin{lemma}\label{cv-origin-intersection-family-sections-analytic-semipositive-stokes}
Let $K$, $(I_k)_{k\in K}$ and $R$ be as in Definition~\ref{lit-multiprojective}, let $n\ge0$, let $X$ be an arithmetic
subvariety of dimension $n+2$, and let $\bar L_1,\dots,\bar L_n$, $\bar B=(0,h_B)$ and $\bar B'=(0,h_{B'})$ be smooth
metrized line bundles. Put $\Phi=c_1(\bar L_1)\wedge\dots\wedge c_1(\bar L_n)$. Then
\[
 \int_{X(\mathbf C)}\log h_B\,c_1(\bar B')\wedge\Phi=\int_{X(\mathbf C)}\log h_{B'}\,c_1(\bar B)\wedge\Phi.
\]
\end{lemma}
\leanref{ReciprocityTheory/ReductiveGroups/ArithmeticGeometry/BirationalGeometry/InterSt.lean}{ReciprocityTheory.ChartRestriction.G.Geo.StokesStatement}
\leanref{ReciprocityTheory/ReductiveGroups/ArithmeticGeometry/BirationalGeometry/InterHironaka.lean}{ReciprocityTheory.ChartRestriction.G.Geo.stokes_of_hironaka}
\begin{proof}
\openstep{The proof is not yet written out. It computes both integrals on an analytic resolution of $X$
(Lemma~\ref{cv-origin-intersection-resolution-exists}), where $\log h_B$ and $\log h_{B'}$ are smooth functions, and
integrates by parts: the integral of $dd^c$ of a smooth function against the other forms vanishes.}
% Lean: stokes_of_hironaka (under GFLogStC and FinTypeQE): ReciprocityTheory.AnForm.M.StokesStatement on the resolution
% (mInt of the potential LocPot.ofFun w is 0), through StokesC.lean and the integration by parts of PartsPiola.lean,
% PartsCalc.lean and PartsIbp.lean. Size (tools/lean_size.py stokes_of_hironaka): 7,869 lines.
\end{proof}

\begin{lemma}\label{cv-origin-intersection-family-sections-analytic-semipositive}
Let $K$, $(I_k)_{k\in K}$ and $R$ be as in Definition~\ref{lit-multiprojective}, and let $n$, $X$, $P$, $\bar A$,
$\bar A'$, $s$, $s'$ and $D$ be as in Definition~\ref{cv-origin-intersection-family-sections-analytic-defect}(2), such
that no horizontal relevant minimal prime of $P+(s)$ is one of $P+(s')$. If the smooth metrized line bundles
$\bar L_1,\dots,\bar L_n$ are semipositive everywhere, then $D(\bar L_1,\dots,\bar L_n)\le0$.
\end{lemma}
\leanref{GroundTheory/HeightsAndTheta/Heights/Intersections/NoCommonComponentCoordinatesPlane.lean}{GroundTheory.Heights₂.Heights.Inter.DeLeSemiPos.Statement}
\leanref{GroundTheory/HeightsAndTheta/Heights/Intersections/PDeLeSemiPos.lean}{GroundTheory.Heights₂.Heights.Inter.DeLeSemiPos.proof}
\begin{proof}
Put $\varphi=\log(|s|/h_A)$, $\varphi'=\log(|s'|/h_{A'})$ and $\Phi=c_1(\bar L_1)\wedge\dots\wedge c_1(\bar L_n)$. For a
smooth metrized line bundle $\bar F$ and a function $g$ write $I(\bar F;g)=\int_{X(\mathbf C)}g\,c_1(\bar F)\wedge\Phi$,
and $I_W(g)=\int_{W(\mathbf C)}g\,\Phi$. Let $\mathcal H$ and $\mathcal H'$ be the sets of the horizontal relevant minimal
primes of $P+(s)$ and of $P+(s')$, with their multiplicities $m_W$. By definition,
\[
 D(\bar L_1,\dots,\bar L_n)=\sum_{W\in\mathcal H}m_WI_W(\varphi')+I(\bar A';\varphi)-\sum_{W\in\mathcal H'}m_WI_W(\varphi)
 -I(\bar A;\varphi').
\]
As in the proof of Lemma~\ref{cv-origin-intersection-family-sections-analytic-additive}, every $W\in\mathcal H$ is an
arithmetic subvariety of dimension $n+1$ with $s'\notin W$, and every $W\in\mathcal H'$ one with $s\notin W$.

For $\varepsilon,\delta>0$ let $\bar B_\varepsilon=\bar A'_{s',\varepsilon}$ and $\bar D_\delta=\bar A_{s,\delta}$
(Definition~\ref{cv-origin-intersection-family-sections-analytic-semipositive-regularized}), with
$\beta_\varepsilon=\log h_{s',\varepsilon}$ and $\gamma_\delta=\log h_{s,\delta}$; they are smooth of multidegree $0$. By
Lemma~\ref{cv-origin-intersection-family-sections-analytic-semipositive-curvature} there is a Fubini--Study bundle
$\bar T=\bar{\mathcal O}(N\mathbf 1)$, which is smooth, such that $\bar F_\varepsilon=\bar A'\otimes\bar B_\varepsilon\otimes\bar T$
is semipositive everywhere for every $\varepsilon>0$.

\emph{The limit in $\varepsilon$.} By Lemma~\ref{cv-origin-intersection-family-sections-analytic-semipositive-first},
applied twice, and Lemma~\ref{cv-origin-intersection-poincare-lelong-all} for $\bar A$, $s$ and $\bar B_\varepsilon$,
\[
 I(\bar F_\varepsilon;\varphi)=I(\bar A';\varphi)+I(\bar B_\varepsilon;\varphi)+I(\bar T;\varphi)
 =I(\bar A';\varphi)+\sum_{W\in\mathcal H}m_WI_W(\beta_\varepsilon)-I(\bar A;\beta_\varepsilon)+I(\bar T;\varphi).
\]
By Lemma~\ref{cv-origin-intersection-family-sections-analytic-semipositive-limit}, on each $W\in\mathcal H$ and on $X$
with the bundles $\bar A,\bar L_1,\dots,\bar L_n$, this tends to
$\lambda=I(\bar A';\varphi)+\sum_{W\in\mathcal H}m_WI_W(\varphi')-I(\bar A;\varphi')+I(\bar T;\varphi)$ as
$\varepsilon\to0$.

\emph{The comparison.} Fix $\delta>0$. Where $s\neq0$, $\varphi\le\gamma_\delta$. As $\bar F_\varepsilon$ and the
$\bar L_j$ are semipositive everywhere,
Lemma~\ref{cv-origin-intersection-family-sections-analytic-semipositive-mono} gives
$I(\bar F_\varepsilon;\varphi)\le I(\bar F_\varepsilon;\gamma_\delta)$. By the second part of
Lemma~\ref{cv-origin-intersection-family-sections-analytic-semipositive-first}, applied twice, and
Lemma~\ref{cv-origin-intersection-family-sections-analytic-semipositive-stokes} for $\bar B_\varepsilon$ and $\bar D_\delta$,
\[
 I(\bar F_\varepsilon;\gamma_\delta)=I(\bar A';\gamma_\delta)+I(\bar B_\varepsilon;\gamma_\delta)+I(\bar T;\gamma_\delta)
 =I(\bar A';\gamma_\delta)+I(\bar D_\delta;\beta_\varepsilon)+I(\bar T;\gamma_\delta).
\]
As $\varepsilon\to0$, $I(\bar D_\delta;\beta_\varepsilon)$ tends to $I(\bar D_\delta;\varphi')$ by
Lemma~\ref{cv-origin-intersection-family-sections-analytic-semipositive-limit} on $X$ with the bundles
$\bar D_\delta,\bar L_1,\dots,\bar L_n$, and $I(\bar D_\delta;\varphi')=\sum_{W\in\mathcal H'}m_WI_W(\gamma_\delta)
-I(\bar A';\gamma_\delta)$ by Lemma~\ref{cv-origin-intersection-poincare-lelong-all} for $\bar A'$, $s'$ and
$\bar D_\delta$. Passing to the limit in the inequality,
$\lambda\le\sum_{W\in\mathcal H'}m_WI_W(\gamma_\delta)+I(\bar T;\gamma_\delta)$.

\emph{The limit in $\delta$.} By Lemma~\ref{cv-origin-intersection-family-sections-analytic-semipositive-limit}, on each
$W\in\mathcal H'$ and on $X$ with the bundles $\bar T,\bar L_1,\dots,\bar L_n$, the right side tends to
$\sum_{W\in\mathcal H'}m_WI_W(\varphi)+I(\bar T;\varphi)$ as $\delta\to0$. So
$\lambda\le\sum_{W\in\mathcal H'}m_WI_W(\varphi)+I(\bar T;\varphi)$, and subtracting $I(\bar T;\varphi)$ from both sides
gives $D(\bar L_1,\dots,\bar L_n)\le0$.
% Lean: DeLeSemiPos.of_remaining = DeLeSemiPos.proof with PoincareLelongAlt.of_remaining (Pub.PoincareLelong, Lemma
% cv-origin-intersection-poincare-lelong-all), stokes_of_hironaka (-stokes), PtsAddFirst (-first), PtsRegLimit
% (-limit), PtsMonoSemiPos (-mono), RegCurvBound (-curvature) and HorizNotMem, under FinTypeQE, GFLogStC, DivPushPoint,
% DivPushResSide and PtsIntegrable; regLine A' s' ht' eps is B_eps and regLine A s ht delta is D_delta; e1 and h1 are
% the limit in eps, h2 the comparison (eR, hlimD, le_of_tendsto_of_tendsto, logNorm_le_log_regLine_h), h3 and h4 the
% limit in delta. Size: PDeLeSemiPos.lean, 190 lines, and PHorizNotMem.lean, 60 lines, besides these lemmas (about
% 7,700 lines in all, besides the resolution of singularities and the integrability statements; the open steps are in
% -absolute and -stokes).
\end{proof}

\begin{lemma}\label{cv-origin-intersection-family-sections-analytic}
Let $K$, $(I_k)_{k\in K}$ and $R$ be as in Definition~\ref{lit-multiprojective}, let $n\ge0$, let $X$ be an arithmetic
subvariety of dimension $n+2$ defined by $P$, let $\bar L_1,\dots,\bar L_n$, $\bar A=(a,h_A)$ and $\bar A'=(a',h_{A'})$ be
smooth metrized line bundles, and let $s,s'\in R\setminus P$ be multihomogeneous of multidegrees $a$ and $a'$, such that
no horizontal relevant minimal prime of $P+(s)$ is one of $P+(s')$. Put $\Phi=c_1(\bar L_1)\wedge\dots\wedge
c_1(\bar L_n)$. Then
\[
 \sum_Wm_W\int_{W(\mathbf C)}\log\frac{|s'|}{h_{A'}}\,\Phi+\int_{X(\mathbf C)}\log\frac{|s|}{h_A}\,c_1(\bar A')\wedge\Phi
 =\sum_{W'}m_{W'}\int_{W'(\mathbf C)}\log\frac{|s|}{h_A}\,\Phi+\int_{X(\mathbf C)}\log\frac{|s'|}{h_{A'}}\,c_1(\bar A)\wedge\Phi,
\]
the sums over the horizontal relevant minimal primes $W$ of $P+(s)$ and $W'$ of $P+(s')$, with their multiplicities
in $R/(P+(s))$ and $R/(P+(s'))$.
\end{lemma}
\leanref{ReciprocityTheory/ReductiveGroups/ArithmeticGeometry/BirationalGeometry/InterSt.lean}{ReciprocityTheory.ChartRestriction.G.Geo.SectionSpaceSymmetry}
\leanref{GroundTheory/HeightsAndTheta/Heights/Intersections/DeEqZeroOfSemiPosOff.lean}{GroundTheory.Heights₂.Heights.Inter.ProjectiveLineTwo.proof}
\begin{proof}
Let $D$ be as in Definition~\ref{cv-origin-intersection-family-sections-analytic-defect}(2); the statement is
$D(\bar L_1,\dots,\bar L_n)=0$. The hypothesis on the components is symmetric in $(\bar A,s)$ and $(\bar A',s')$, and
exchanging these two pairs changes $D$ into $-D$.

We show by induction on $|S|$, for subsets $S\subset\{1,\dots,n\}$, that $D(\bar L_1,\dots,\bar L_n)=0$ for all smooth
$\bar L_1,\dots,\bar L_n$ such that $\bar L_j$ is semipositive everywhere for every $j\notin S$. For $S=\emptyset$, all
$\bar L_j$ are semipositive everywhere, and Lemma~\ref{cv-origin-intersection-family-sections-analytic-semipositive}
gives $D\le0$ and, for the exchanged pairs, $-D\le0$. Now let $S=S_0\cup\{j\}$ with $j\notin S_0$, and suppose the claim
for $S_0$. By Lemma~\ref{cv-origin-intersection-family-sections-analytic-twist} there is a smooth $\bar T$ such that
$\bar T$ and $\bar L_j\otimes\bar T$ are semipositive everywhere. By
Lemma~\ref{cv-origin-intersection-family-sections-analytic-additive}, the value of $D$ with $\bar L_j\otimes\bar T$ in the
place $j$ is the sum of its values with $\bar L_j$ and with $\bar T$ there. The first and the last of these values vanish
by the claim for $S_0$, so the value with $\bar L_j$ vanishes. For $S=\{1,\dots,n\}$ this gives the statement.
\end{proof}
\leanref{GroundTheory/HeightsAndTheta/Heights/Intersections/DeEqZeroOfSemiPosOff.lean}{GroundTheory.Heights₂.Heights.Inter.De_eq_zero_of_semiPos_off}
% Lean: ProjectiveLineTwo.of_remaining = ProjectiveLineTwo.proof (De_eq_zero_of_semiPos_off with Finset.univ; the base
% case De_eq_zero_of_semiPos with De_swap) from DeSlotAdd.of_remaining (Lemma
% cv-origin-intersection-family-sections-analytic-additive), SemiPosTwist.full (-twist) and DeLeSemiPos.of_remaining
% (-semipositive), under FinTypeQE and GFLogStC (Assembly.lean). Size: DeEqZeroOfSemiPosOff.lean, about 60 lines, besides
% these three lemmas.

\begin{lemma}\label{cv-origin-intersection-family-sections-swap}
Let $K$, $(I_k)_{k\in K}$ and $R$ be as in Definition~\ref{lit-multiprojective}, let $n\ge0$, and let $I$ be as in
Definition~\ref{cv-origin-intersection-family-numbers}(1), satisfying the recursion of
Definition~\ref{lit-multiprojective}(3) with every section on arithmetic subvarieties of dimension $n+1$. Let $X$ be an
arithmetic subvariety of dimension $n+2$ defined by $P$, let $\bar L_1,\dots,\bar L_n$, $\bar A=(a,h_A)$ and
$\bar A'=(a',h_{A'})$ be smooth metrized line bundles, and let $s,s'\in R\setminus P$ be multihomogeneous of multidegrees
$a$ and $a'$, such that no relevant minimal prime of $P+(s)$ is a relevant minimal prime of $P+(s')$. Then
\[
 \mathcal R_I(\bar L_1\cdots\bar L_n\bar A';\bar A,s\mid X)=\mathcal R_I(\bar L_1\cdots\bar L_n\bar A;\bar A',s'\mid X).
\]
\end{lemma}
\leanref{GroundTheory/HeightsAndTheta/Heights/Intersections/StatementsArith.lean}{GroundTheory.Heights₂.Heights.Inter.Swap.Statement}
\leanref{GroundTheory/HeightsAndTheta/Heights/Intersections/NotLtLtOfMemMinimalPrimes.lean}{GroundTheory.Heights₂.Heights.Inter.Swap.of_cycles}
\begin{proof}
Write $\bar L$ for $\bar L_1\cdots\bar L_n$ and $\Phi=c_1(\bar L_1)\wedge\dots\wedge c_1(\bar L_n)$, and use the notation
$m_U(W,f)$ of Lemma~\ref{cv-origin-intersection-family-sections-cycles}. For a relevant multihomogeneous prime $U$ let
$c(U)=I(\bar L\mid U)$ if $U$ is horizontal and $c(U)=\deg_{L_1\cdots L_n}(U)\log p$ if $U\cap\mathbf Z=p\mathbf Z$.

First, $s'$ lies in no relevant minimal prime $W$ of $P+(s)$. Otherwise $W$ contains $P+(s')$, and it is minimal over it:
a prime $Q$ with $P+(s')\subset Q\subset W$ is not $P$, as $s'\notin P$, and no prime lies strictly between $P$ and $W$
\cite[Tag 00KV]{Sta}; so $W$ would be a relevant minimal prime of $P+(s')$ too. Likewise $s$ lies in no relevant minimal
prime of $P+(s')$.

Let $W$ be a relevant minimal prime of $P+(s)$. We show
\[
 \text{(term of $W$)}=\sum_Um_U(W,s')\,c(U)-\delta_W\int_{W(\mathbf C)}\log\frac{|s'|}{h_{A'}}\,\Phi,
 \tag{$*$}
\]
where the term of $W$ is $I(\bar L\bar A'\mid W)$ if $W$ is horizontal and $\deg_{L_1\cdots L_nA'}(W)\log p$ if
$W\cap\mathbf Z=p\mathbf Z$, the sum is over the relevant minimal primes $U$ of $W+(s')$, and $\delta_W$ is $1$ if $W$ is
horizontal and $0$ otherwise. If $W$ is horizontal, it is an arithmetic subvariety of dimension $n+1$, and $s'\notin W$, so
by the recursion with the section $s'$, $I(\bar L\bar A'\mid W)=\mathcal R_I(\bar L;\bar A',s'\mid W)$, which is the
right side of $(*)$. If $W\cap\mathbf Z=p\mathbf Z$, then $\dim R/W=n+1+|K|$, and
Lemma~\ref{cv-origin-intersection-family-sections-degree}(2) with $t=s'$ gives
$\deg_{L_1\cdots L_nA'}(W)=\sum_Um_U(W,s')\deg_{L_1\cdots L_n}(U)$. Each such $U$ contains $p$, so $c(U)$ is
$\deg_{L_1\cdots L_n}(U)\log p$, and $(*)$ holds.

The integral in $\mathcal R_I(\bar L\bar A';\bar A,s\mid X)$ is that of $\log(|s|/h_A)$ against $c_1(\bar A')\wedge\Phi$,
as the Monge--Amp\`ere density is symmetric in its arguments (Definition~\ref{cv-period-monge-ampere}(2)). Multiplying
$(*)$ by $m_W(P,s)$ and summing gives
\[
 \mathcal R_I(\bar L\bar A';\bar A,s\mid X)=\sum_U\Big(\sum_Wm_W(P,s)\,m_U(W,s')\Big)c(U)
 -\sum_Wm_W(P,s)\int_{W(\mathbf C)}\log\frac{|s'|}{h_{A'}}\,\Phi-\int_{X(\mathbf C)}\log\frac{|s|}{h_A}\,c_1(\bar A')\wedge\Phi,
\]
the middle sum over the horizontal $W$, and likewise with the roles of $(\bar A,s)$ and $(\bar A',s')$ exchanged. The first
terms agree by Lemma~\ref{cv-origin-intersection-family-sections-cycles}, and the remaining terms by
Lemma~\ref{cv-origin-intersection-family-sections-analytic}.
\end{proof}
\leanref{GroundTheory/HeightsAndTheta/Heights/Intersections/NotLtLtOfMemMinimalPrimes.lean}{GroundTheory.Heights₂.Heights.Inter.Swap.proof}
% Lean: Swap.proof = Swap.of_cycles SwapCycles.full GeoNumCut SectionSpaceSymmetry: hs'W and hsW' (mem_divSupp_of_mem,
% mem_minimalPrimes_of_mem), key (the recursion hR on horizontal W through divSuppArith, GeoNumCut on vertical W with
% comap_eq_of_vertical and charOf), side (hperm: Fin.cons of the forms is a rotation of Fin.snoc, ptsIntegral_comp_perm;
% sum_cycle_comp for the double sum), hCyc and hPL2. Size: NotLtLtOfMemMinimalPrimes.lean, 291 lines, besides
% Lemmas cv-origin-intersection-family-sections-cycles and -analytic.

\begin{lemma}\label{cv-origin-intersection-family-sections}
Let $K$, $(I_k)_{k\in K}$ and $R$ be as in Definition~\ref{lit-multiprojective}, let $n\ge0$, and let $I$ be as in
Definition~\ref{cv-origin-intersection-family-numbers}(1), additive in the last bundle and satisfying the recursion of
Definition~\ref{lit-multiprojective}(3) with every section on arithmetic subvarieties of dimension $n+1$. Let $X$ be an
arithmetic subvariety of dimension $n+2$ defined by $P$, let $\bar L_1,\dots,\bar L_{n+1}$ and $\bar B=(b,h_B)$ be smooth
metrized line bundles, and let $s,t\in R\setminus P$ be multihomogeneous of multidegree $b$. Then
$\mathcal R_I(\bar L_1\cdots\bar L_{n+1};\bar B,s\mid X)=\mathcal R_I(\bar L_1\cdots\bar L_{n+1};\bar B,t\mid X)$.
\end{lemma}
\leanref{GroundTheory/HeightsAndTheta/Heights/Intersections/StatementsWellDefined.lean}{GroundTheory.Heights₂.Heights.Inter.WellDefined.Statement}
\leanref{GroundTheory/HeightsAndTheta/Heights/Intersections/PWellDefined.lean}{GroundTheory.Heights₂.Heights.Inter.WellDefined.proof}
\begin{proof}
Write $\bar L'$ for $\bar L_1\cdots\bar L_n$ and $\bar M=(m,h_M)$ for $\bar L_{n+1}$. Let $S$ be the set consisting of
$P$ and the relevant minimal primes of $P+(s)$ and of $P+(t)$, a finite set of relevant multihomogeneous primes, and
let $d_0$ be as in Lemma~\ref{cv-origin-intersection-family-sections-avoid} for $S$. Put
$k=\sum_l(d_{0,l}+|m_l|)$ and let $\bar H=\bar{\mathcal O}(k\mathbf 1)$, the Fubini--Study bundle
(Definition~\ref{cv-origin-analytic-system-pullback}(1)). As $m+k\mathbf 1\ge d_0$ and $k\mathbf 1\ge d_0$, there are
multihomogeneous $a_1$ of multidegree $m+k\mathbf 1$, the multidegree of $\bar M\otimes\bar H$, and $a_2$ of
multidegree $k\mathbf 1$, that lie in no element of $S$; in particular $a_1,a_2\notin P$. No relevant minimal prime $W$ of
$P+(s)$ or of $P+(t)$ is one of $P+(a_1)$ or of $P+(a_2)$, as $a_1,a_2\notin W$.

By Lemma~\ref{cv-origin-intersection-family-sections-swap}, applied with the recursion of $I$ in dimension $n+1$,
\[
 \mathcal R_I(\bar L'(\bar M\otimes\bar H);\bar B,s\mid X)=\mathcal R_I(\bar L'\bar B;\bar M\otimes\bar H,a_1\mid X)
 =\mathcal R_I(\bar L'(\bar M\otimes\bar H);\bar B,t\mid X),
\]
and likewise $\mathcal R_I(\bar L'\bar H;\bar B,s\mid X)=\mathcal R_I(\bar L'\bar B;\bar H,a_2\mid X)
=\mathcal R_I(\bar L'\bar H;\bar B,t\mid X)$. By Lemma~\ref{cv-origin-intersection-family-sections-position}, applied with
the additivity of $I$ in dimension $n+1$, for $f=s$ and for $f=t$,
\[
 \mathcal R_I(\bar L'(\bar M\otimes\bar H);\bar B,f\mid X)=\mathcal R_I(\bar L'\bar M;\bar B,f\mid X)
 +\mathcal R_I(\bar L'\bar H;\bar B,f\mid X).
\]
Subtracting gives $\mathcal R_I(\bar L'\bar M;\bar B,s\mid X)=\mathcal R_I(\bar L'\bar M;\bar B,t\mid X)$, which is the
statement.
\end{proof}
% Lean: WellDefined.proof (h1 = Swap.Statement, h2 = RPositionAdd.Statement): S = divSupp P s u divSupp P t u {P}
% (isMultiProjPrime_of_mem_divSupp), AvoidRing.exists_form_avoid for d0, k = sum of d0 l + |M.a l|, H = sqMetric (k),
% a1 and a2 from hd0, divSupp_disjoint_of_not_mem, e1s, e1t, e2s, e2t (Swap) and fs, ft (RPositionAdd). In
% ThmIExists.proof, h1 is Swap.proof and h2 is RPositionAdd.of_remaining. Size: PWellDefined.lean, about 90 lines,
% besides Lemmas cv-origin-intersection-family-sections-avoid, -position and -swap.

\begin{lemma}\label{cv-origin-intersection-family-recursion}
The numbers $(\cdot\mid\cdot)_0$ of Definition~\ref{cv-origin-intersection-family-numbers}(2) are additive in the last
line bundle, and they satisfy the recursion of Definition~\ref{lit-multiprojective}(3) with every section.
\end{lemma}
\leanref{GroundTheory/HeightsAndTheta/Heights/Intersections/ThmIExistsFamily.lean}{GroundTheory.Heights₂.Heights.Inter.ThmIExists.isRFamily_famI}
\begin{proof}
Write $(\cdot)_0$ for $(\cdot\mid\cdot)_0$ and $\mathcal R_0$ for $\mathcal R_I$ with these numbers; all line bundles below
are smooth. For $n\ge0$ consider three properties of arithmetic subvarieties $X$ of dimension $n+1$, defined by $P$, and
smooth metrized line bundles $\bar L=(\bar L_1,\dots,\bar L_n)$:
\begin{itemize}
\item[(W$_n$)] $\mathcal R_0(\bar L;\bar B,s\mid X)=\mathcal R_0(\bar L;\bar B,t\mid X)$ for all $\bar B$ and all
  $s,t\in R\setminus P$ multihomogeneous of the multidegree of $\bar B$;
\item[(R$_n$)] $(\bar L\,\bar A\mid X)_0=\mathcal R_0(\bar L;\bar A,s\mid X)$ for all $\bar A$ and all $s\in R\setminus P$
  multihomogeneous of the multidegree of $\bar A$;
\item[(A$_n$)] $(\bar L\,(\bar A\otimes\bar B)\mid X)_0=(\bar L\,\bar A\mid X)_0+(\bar L\,\bar B\mid X)_0$ for all $\bar A$ and
  $\bar B$.
\end{itemize}
For $\bar A$ of multidegree $a$ write $H_{\bar A}=\bar{\mathcal O}(m\mathbf 1)$ with $m=\sum_k|a_k|$, a smooth metrized
line bundle, $t_{\bar A}=s_{a+m\mathbf 1}(P)$ and $u_{\bar A}=s_{m\mathbf 1}(P)$, sections of $\bar A\otimes H_{\bar A}$ and
of $H_{\bar A}$ outside $P$, so that by definition
$(\bar L\,\bar A\mid X)_0=\mathcal R_0(\bar L;\bar A\otimes H_{\bar A},t_{\bar A}\mid X)-\mathcal R_0(\bar L;H_{\bar A},
u_{\bar A}\mid X)$.

\emph{A change of auxiliary bundle.} Assume (W$_n$). Let $\bar C$ be a line bundle, and let $t,c\in R\setminus P$ be
multihomogeneous of the multidegrees of $\bar A\otimes\bar C$ and of $\bar C$. Then
\[
 (\bar L\,\bar A\mid X)_0=\mathcal R_0(\bar L;\bar A\otimes\bar C,t\mid X)-\mathcal R_0(\bar L;\bar C,c\mid X).
\]
Indeed, write $H=H_{\bar A}$. By Lemma~\ref{cv-origin-intersection-family-product},
$\mathcal R_0(\bar L;(\bar A\otimes H)\otimes\bar C,t_{\bar A}c)=\mathcal R_0(\bar L;\bar A\otimes H,t_{\bar A})
+\mathcal R_0(\bar L;\bar C,c)$ and $\mathcal R_0(\bar L;(\bar A\otimes\bar C)\otimes H,tu_{\bar A})
=\mathcal R_0(\bar L;\bar A\otimes\bar C,t)+\mathcal R_0(\bar L;H,u_{\bar A})$, on $X$. The two bundles on the left are equal,
and their sections $t_{\bar A}c$ and $tu_{\bar A}$ are not in $P$, so by (W$_n$) the left sides are equal. Subtracting
$\mathcal R_0(\bar L;H,u_{\bar A})+\mathcal R_0(\bar L;\bar C,c)$ from both sides gives the claim, by the definition.

\emph{(W$_n$) implies (R$_n$) and (A$_n$).} For (R$_n$), take $\bar C=H_{\bar A}$, $c=u_{\bar A}$ and $t=su_{\bar A}$, and
apply Lemma~\ref{cv-origin-intersection-family-product} to $\mathcal R_0(\bar L;\bar A\otimes H_{\bar A},su_{\bar A})$.
For (A$_n$), take $\bar C=H_{\bar A}\otimes H_{\bar B}$, $c=u_{\bar A}u_{\bar B}$ and $t=t_{\bar A}t_{\bar B}$, a section
of $(\bar A\otimes H_{\bar A})\otimes(\bar B\otimes H_{\bar B})=(\bar A\otimes\bar B)\otimes\bar C$. By
Lemma~\ref{cv-origin-intersection-family-product}, applied to both terms, $(\bar L\,(\bar A\otimes\bar B)\mid X)_0$ is
$\mathcal R_0(\bar L;\bar A\otimes H_{\bar A},t_{\bar A})+\mathcal R_0(\bar L;\bar B\otimes H_{\bar B},t_{\bar B})
-\mathcal R_0(\bar L;H_{\bar A},u_{\bar A})-\mathcal R_0(\bar L;H_{\bar B},u_{\bar B})$, which is
$(\bar L\,\bar A\mid X)_0+(\bar L\,\bar B\mid X)_0$ by the definition.

\emph{Induction.} (W$_0$) is Lemma~\ref{cv-origin-intersection-family-curves}. If (W$_n$) holds, then so do (R$_n$) and
(A$_n$), and Lemma~\ref{cv-origin-intersection-family-sections}, applied to the numbers $(\cdot\mid\cdot)_0$, gives
(W$_{n+1}$). So (R$_n$) and (A$_n$) hold for every $n$, which is the statement.
\end{proof}
\leanref{GroundTheory/HeightsAndTheta/Heights/Intersections/ThmIExistsFamily.lean}{GroundTheory.Heights₂.Heights.Inter.ThmIExists.famI_eq_diff}
\leanref{GroundTheory/HeightsAndTheta/Heights/Intersections/ThmIExistsFamily.lean}{GroundTheory.Heights₂.Heights.Inter.ThmIExists.all_dims}
% Lean: isRFamily_famI from all_dims (hasRAt_famI, lastAddAt_famI, wdAt_zero, famI_snoc, famI_eq_diff), with
% RegularityIndependence (AutoClose.Inter_R1Indep), RMul.proof and WellDefined.proof. The text's (W_n), (R_n) and (A_n)
% are WDAt, HasRAt and LastAddAt; the change of auxiliary bundle is famI_eq_diff (e1, e2 from RMul and e3 from WDAt with
% tensor_right_comm); (R_n) is hasRAt_famI and (A_n) is lastAddAt_famI (tensor_tensor_comm); the sections s_b(P) are
% secOf, products of the x_{k,i_k}^{b_k} (exists_form), and H_A is segreLine (shift A), smooth by isSmoothHerm_segreLine.
% The Lean builds the numbers through fam and lift, bookkeeping for the dimension. Size: ThmIExistsFamily.lean, 239 lines.

\begin{lemma}\label{cv-origin-intersection-family-symmetry}
Let $K$, $(I_k)_{k\in K}$ and $R$ be as in Definition~\ref{lit-multiprojective}, and let $I$ be as in
Definition~\ref{cv-origin-intersection-family-numbers}(1), additive in the last line bundle and satisfying the recursion
of Definition~\ref{lit-multiprojective}(3) with every section. Then $I(\bar L_1\cdots\bar L_n\mid X)$ is symmetric in the
$\bar L_j$.
\end{lemma}
\leanref{ReciprocityTheory/ReductiveGroups/ArithmeticGeometry/ArithmeticIntersection/RationalSections/ThmISt.lean}{ReciprocityTheory.ChartRestriction.G.AI.I2.SymmStatement}
\leanref{GroundTheory/HeightsAndTheta/Heights/Intersections/PThmISymm.lean}{GroundTheory.Heights₂.Heights.Inter.ThmISymm.proof}
\begin{proof}
All line bundles below are smooth. For a metrized line bundle $\bar A=(a,h)$ let $\bar H_A=\bar{\mathcal O}(m_A\mathbf 1)$
with $m_A=\sum_k|a_k|$, the Fubini--Study bundle (Definition~\ref{cv-origin-analytic-system-pullback}(1)). For an
arithmetic subvariety $X$ defined by $P$, the multidegrees $a+m_A\mathbf 1$ and $m_A\mathbf 1$ of
$\bar A\otimes\bar H_A$ and $\bar H_A$ are nonnegative, and as $P$ is relevant there are multihomogeneous
$t_A,u_A\in R\setminus P$ of these multidegrees, for instance products of powers of coordinates $x_{k,i_k}\notin P$. By
the additivity in the last bundle and the recursion,
\[
 I(\bar L\bar A\mid X)=I(\bar L(\bar A\otimes\bar H_A)\mid X)-I(\bar L\bar H_A\mid X)
 =\mathcal R_I(\bar L;\bar A\otimes\bar H_A,t_A\mid X)-\mathcal R_I(\bar L;\bar H_A,u_A\mid X)
 \tag{$*$}
\]
for every $\bar L=\bar L_1\cdots\bar L_{n-1}$, if $X$ has dimension $n$.

We show by induction on $n$ that $I(\bar L_{\sigma(1)}\cdots\bar L_{\sigma(n)}\mid X)=I(\bar L_1\cdots\bar L_n\mid X)$ for
every arithmetic subvariety $X$ of dimension $n$ and every permutation $\sigma$. For $n\le1$ there is nothing to show. Let
$n=m+2$, and assume the claim in dimension $m+1$. The permutations of $\{1,\dots,m+2\}$ are products of adjacent
transpositions, so it suffices to treat these. Let $X$, of dimension $m+2$, be defined by $P$, and write $\bar L$ for
$\bar L_1\cdots\bar L_m$.

\emph{A transposition that fixes the last place.} Let $\tau$ permute $\{1,\dots,m+1\}$, and let $\bar A$ be the last
bundle. For a section $f\in R\setminus P$ of a bundle $\bar B$,
$\mathcal R_I(\bar L_{\tau(1)}\cdots\bar L_{\tau(m+1)};\bar B,f\mid X)=\mathcal R_I(\bar L_1\cdots\bar L_{m+1};\bar B,f\mid X)$
term by term: on a horizontal relevant minimal prime $W$ of $P+(f)$, an arithmetic subvariety of dimension $m+1$, by the
induction hypothesis; on one with $W\cap\mathbf Z=p\mathbf Z$, by the symmetry of the degree
(Lemma~\ref{cv-origin-intersection-family-sections-degree}(1)); and the integral by the symmetry of the
Monge--Amp\`ere density (Definition~\ref{cv-period-monge-ampere}(2)). By $(*)$ for both orders, the claim follows.

\emph{The transposition of the last two places.} Here the claim is $I(\bar L\bar X\bar Y\mid X)=I(\bar L\bar Y\bar X\mid X)$
for all $\bar X$ and $\bar Y$. We use four steps.
\begin{enumerate}
\item[(a)] If $x,y\in R\setminus P$ are sections of $\bar X$ and $\bar Y$ such that no relevant minimal prime of $P+(y)$ is
  one of $P+(x)$, then $I(\bar L\bar X\bar Y\mid X)=I(\bar L\bar Y\bar X\mid X)$. Indeed, by the recursion and
  Lemma~\ref{cv-origin-intersection-family-sections-swap}, with the recursion in dimension $m+1$,
  $I(\bar L\bar X\bar Y\mid X)=\mathcal R_I(\bar L\bar X;\bar Y,y\mid X)=\mathcal R_I(\bar L\bar Y;\bar X,x\mid X)
  =I(\bar L\bar Y\bar X\mid X)$.
\item[(b)] If $\bar Y$ has a section $y\in R\setminus P$, then
  $I(\bar L(\bar A\otimes\bar C)\bar Y\mid X)=I(\bar L\bar A\bar Y\mid X)+I(\bar L\bar C\bar Y\mid X)$ for all $\bar A$ and
  $\bar C$: by the recursion with $y$ on the three terms and
  Lemma~\ref{cv-origin-intersection-family-sections-position}, with the additivity in dimension $m+1$. For any $\bar Y$,
  the same additivity follows by writing each term as in $(*)$, through $\bar Y\otimes\bar H_Y$ with $t_Y$ and $\bar H_Y$
  with $u_Y$.
\item[(c)] If $\bar Y$ has a section $y\in R\setminus P$, then $I(\bar L\bar X\bar Y\mid X)=I(\bar L\bar Y\bar X\mid X)$.
  Let $S$ consist of $P$ and the relevant minimal primes of $P+(y)$, let $d_0$ be as in
  Lemma~\ref{cv-origin-intersection-family-sections-avoid} for $S$, put $k=m_X+\sum_ld_{0,l}$ and
  $\bar H=\bar{\mathcal O}(k\mathbf 1)$. There are sections $x$ of $\bar X\otimes\bar H$ and $x'$ of $\bar H$, of
  multidegrees at least $d_0$, that lie in no element of $S$; so no relevant minimal prime of $P+(y)$ is one of $P+(x)$ or
  of $P+(x')$. By (b), (a) twice and the additivity in the last bundle,
  \[
   I(\bar L\bar X\bar Y\mid X)=I(\bar L(\bar X\otimes\bar H)\bar Y\mid X)-I(\bar L\bar H\bar Y\mid X)
   =I(\bar L\bar Y(\bar X\otimes\bar H)\mid X)-I(\bar L\bar Y\bar H\mid X)=I(\bar L\bar Y\bar X\mid X).
  \]
\item[(d)] For any $\bar Y$, by the additivity in the last bundle, (c) for $\bar Y\otimes\bar H_Y$ and for $\bar H_Y$, and
  (b),
  \[
   I(\bar L\bar X\bar Y\mid X)=I(\bar L\bar X(\bar Y\otimes\bar H_Y)\mid X)-I(\bar L\bar X\bar H_Y\mid X)
   =I(\bar L(\bar Y\otimes\bar H_Y)\bar X\mid X)-I(\bar L\bar H_Y\bar X\mid X)=I(\bar L\bar Y\bar X\mid X).
  \]
\end{enumerate}
\end{proof}
\leanref{GroundTheory/HeightsAndTheta/Heights/Intersections/PThmISymm.lean}{GroundTheory.Heights₂.Heights.Inter.ThmISymm.lastFix}
\leanref{GroundTheory/HeightsAndTheta/Heights/Intersections/PThmISymm.lean}{GroundTheory.Heights₂.Heights.Inter.ThmISymm.swap_last}
% Lean: ThmISymm.proof: induction on n, the cases 0 and 1, Equiv.Perm.mclosure_swap_castSucc_succ (adjacent
% transpositions generate) and Submonoid.closure_induction; lastFix (key is (*), with segreLine (shift A), isSectionOn_up
% and isSectionOn_segre; recR_comp_perm: the induction hypothesis on horizontal components through divSuppArith,
% geoNum_comp_perm from FibreCompStatement, I1.ExistsUniqueStatement, I1.SymmStatement and GeoNumStatement over Z/p, and
% ptsIntegral_comp_perm, maDensity_comp_perm); swap_last (d), swap_half (c) with Choice.exists_form_avoid, addX (b),
% addX_sec (b) with RPositionAdd, swap_sec (a) with Swap. ThmIExists.proof takes Swap.proof, RPositionAdd.of_remaining,
% fibreComp, Targets.ExistsUnique.full, Targets.SymmStatement.of_type and geoNumStatement. Size: PThmISymm.lean, 293
% lines, besides Lemmas cv-origin-intersection-family-sections-degree, -position, -swap and -avoid.

\begin{lemma}\label{cv-origin-intersection-family}
There is an arithmetic intersection system (Definition~\ref{lit-multiprojective}(3)).
\end{lemma}
\leanref{ReciprocityTheory/ReductiveGroups/ArithmeticGeometry/ArithmeticIntersection/RationalSections/ThmISt.lean}{ReciprocityTheory.ChartRestriction.G.AI.I2.ExistsStatement}
\leanref{Development/AutomorphicSpectrum/Proofs/MetricProofs/IntersectionExistsOfPoincareLelongStatement.lean}{GlobalTheory.T.intersection_exists_of_poincare_lelong}
\leanref{Development/AutomorphicSpectrum/Proofs/MetricProofs/IntersectionExistsOfPoincareLelong.lean}{GlobalTheory.CompanionTheory.Ar.Pub.thmIExists}
\begin{proof}
Let $(\cdot\mid\cdot)_0$ be the numbers of Definition~\ref{cv-origin-intersection-family-numbers}(2). By
Lemma~\ref{cv-origin-intersection-family-recursion} they are additive in the last line bundle and satisfy the recursion
of Definition~\ref{lit-multiprojective}(3) with every section, and by
Lemma~\ref{cv-origin-intersection-family-symmetry} they are symmetric. A symmetric function that is additive in its
last argument is additive in each argument. So $(\cdot\mid\cdot)_0$ is an arithmetic intersection system.
\end{proof}
\leanref{GroundTheory/HeightsAndTheta/Heights/Intersections/PThmIExists.lean}{GroundTheory.Heights₂.Heights.Inter.ThmIExists.proof}
% Lean: Pub.thmIExists (= GlobalTheory.T.intersection_exists_of_poincare_lelong), under Pub.PoincareLelong (Lemma
% cv-origin-intersection-poincare-lelong-all) and ResExistsStatement (Lemma
% cv-origin-intersection-resolution-exists): Inter.ThmIExists.proof with the family famI (ThmIExists.fam: step,
% recR of A.tensor (segreLine (shift A)) with secOf P (upDeg A) minus recR of segreLine (shift A) with
% secOf P (fun _ => shift A)), isRFamily_famI from RegularityIndependence (AutoClose.Inter_R1Indep) and RMul,
% WellDefined.proof (from Swap and RPositionAdd) for the independence of the sections, and ThmISymm.proof for the
% symmetry; Swap rests on SectionSpaceSymmetry (ProjectiveLineTwo), DeLeSemiPos (from the Poincare-Lelong formula,
% StokesStatement and PtsRegLimit) and the integrability statements PtsIntegrable, PtsAbsIntegrable and PtsBilinear.
% ThmIExists.proof returns IsIntersectionFamily (symmetry, additivity in the last argument, recursion); additivity in
% each argument is multilin_of_isFamily. Size: about 35,000 lines, in Lemmas cv-origin-intersection-family-curves,
% -product, -sections, -recursion and -symmetry.

\subsection{Projections and generic cycles}\label{subsection-cv-projections-generic-cycles}

\begin{definition}\label{cv-origin-intersection-full}
Let $K$, $(I_k)_{k\in K}$ and $R$, and $K'$, $(I'_{k'})_{k'\in K'}$ and $R'=\mathbf Z[x'_{k',i}:k'\in K',\ i\in I'_{k'}]$,
be two sets of data as in Definition~\ref{lit-multiprojective}.
\begin{enumerate}
\item A \emph{block map} $(\kappa,v)$ consists of injective maps $\kappa:K\to K'$ and
  $v:\bigsqcup_kI_k\to\bigsqcup_{k'}I'_{k'}$ such that $v$ maps $I_k$ onto $I'_{\kappa(k)}$ for every $k$. Let
  $v_*:R\to R'$ be the ring homomorphism with $v_*(x_{k,i})=x'_{v(k,i)}$. The \emph{pullback} of a metrized line
  bundle $\bar L=(a,h)$ on $\prod_k\mathbf P^{I_k}_{\mathbf Z}$ is the metrized line bundle $\bar L'=(a',h')$ on
  $\prod_{k'}\mathbf P^{I'_{k'}}_{\mathbf Z}$ with $a'_{\kappa(k)}=a_k$ for $k\in K$, $a'_{k'}=0$ for
  $k'\notin\kappa(K)$, and $h'(z')=h\big((z'_{v(k,i)})_{k,i}\big)$.
\item An arithmetic intersection system $(\cdot\mid\cdot)$ (Definition~\ref{lit-multiprojective}(3)) satisfies the
  \emph{projection formula} if the following holds for every block map $(\kappa,v)$. Let $\Omega$ be a field of
  characteristic zero, $d\ge0$, and $y\in\Omega^{\sqcup_{k'}I'_{k'}}$ a point whose kernel $P'$
  (Definition~\ref{lit-multiprojective}(4)) is an arithmetic subvariety $X'$ of dimension $d$, such that
  $v_*^{-1}(P')$ is an arithmetic subvariety $X$ of dimension $d$. Let $\mathbf Q(y)\subset\Omega$ be the field
  generated by the ratios $y_{k',i}/y_{k',j}$ and $\mathbf Q(y\circ v)\subset\mathbf Q(y)$ the field generated by the
  ratios $y_{v(k,i)}/y_{v(k,j)}$; the extension is finite, as both fields have transcendence degree $d-1$ over
  $\mathbf Q$. Then for smooth metrized line bundles $\bar L_1,\dots,\bar L_d$ on $\prod_k\mathbf P^{I_k}_{\mathbf Z}$
  with pullbacks $\bar L'_1,\dots,\bar L'_d$,
  \[
   (\bar L'_1\cdots\bar L'_d\mid X')=[\mathbf Q(y):\mathbf Q(y\circ v)]\,(\bar L_1\cdots\bar L_d\mid X).
  \]
\item An arithmetic intersection system $(\cdot\mid\cdot)$ has the \emph{contraction property} if the following holds
  for every block map $(\kappa,v)$. Let $d_0<d$ and $r\ge0$, and let $X'$ be an arithmetic subvariety of dimension
  $d+r$ of $\prod_{k'}\mathbf P^{I'_{k'}}_{\mathbf Z}$, defined by the prime $P'$, such that $v_*^{-1}(P')$ is an
  arithmetic subvariety of dimension $d_0$. Then for smooth metrized line bundles $\bar L_1,\dots,\bar L_d$ on
  $\prod_k\mathbf P^{I_k}_{\mathbf Z}$ with pullbacks $\bar L'_1,\dots,\bar L'_d$, and smooth metrized line bundles
  $\bar N_1,\dots,\bar N_r$ on $\prod_{k'}\mathbf P^{I'_{k'}}_{\mathbf Z}$,
  $(\bar L'_1\cdots\bar L'_d\,\bar N_1\cdots\bar N_r\mid X')=0$.
\item An arithmetic intersection system is \emph{full} if it satisfies the projection formula and has the contraction
  property.
\end{enumerate}
\end{definition}
\leanref{ReciprocityTheory/ReductiveGroups/ArithmeticGeometry/BirationalGeometry/InterNamed.lean}{ReciprocityTheory.ChartRestriction.G.Geo.IsFullIntersectionSystem}
\leanref{ReciprocityTheory/ReductiveGroups/ArithmeticGeometry/BirationalGeometry/InterNamed.lean}{ReciprocityTheory.ChartRestriction.G.Geo.IsMultilinFam}
\leanref{ReciprocityTheory/ReductiveGroups/ArithmeticGeometry/BirationalGeometry/InterNamed.lean}{ReciprocityTheory.ChartRestriction.G.Geo.HasPFCoord}
\leanref{ReciprocityTheory/ReductiveGroups/ArithmeticGeometry/BirationalGeometry/InterNamed.lean}{ReciprocityTheory.ChartRestriction.G.Geo.HasContractMixed}
\leanref{ReciprocityTheory/ReductiveGroups/ArithmeticGeometry/ArithmeticIntersection/MultihomogeneousIdeals/IsFieldIntersection.lean}{ReciprocityTheory.ChartRestriction.G.AI.I1.IsBlockMap}
\leanref{ReciprocityTheory/ReductiveGroups/ArithmeticGeometry/ArithmeticIntersection/MixedFieldContraction/GenericFibreKerTarget.lean}{ReciprocityTheory.ChartRestriction.G.AI.I4.IsPullHerm}
\leanref{ReciprocityTheory/ReductiveGroups/ArithmeticGeometry/ArithmeticIntersection/MixedFieldContraction/GenericFibreKerTarget.lean}{ReciprocityTheory.ChartRestriction.G.AI.I4.imageZ}
\leanref{ReciprocityTheory/ReductiveGroups/ArithmeticGeometry/ArithmeticIntersection/MixedFieldContraction/GenericFibreKerTarget.lean}{ReciprocityTheory.ChartRestriction.G.AI.I4.blockFnField}

\begin{lemma}\label{cv-origin-intersection-projection-fibre}
Let $(\kappa,v)$ be a block map (Definition~\ref{cv-origin-intersection-full}(1)), let $X'$ be an arithmetic subvariety
of dimension $d+r$ of $\prod_{k'}\mathbf P^{I'_{k'}}_{\mathbf Z}$ defined by the prime $P'$, such that
$Q=v_*^{-1}(P')$ is an arithmetic subvariety of dimension $d$, and let $\Phi$ be the fraction field of $R/Q$. Then there
are a field $\Phi'$ containing $\Phi$ and a homomorphism $R'\to\Phi'$ with kernel $P'$ whose composite with $v_*$ is
$R\to R/Q\to\Phi$. For any such, the kernel $V$ of the induced homomorphism
$\Phi[x'_{k',i'}:k'\notin\kappa(K)]\to\Phi'$ is a relevant multihomogeneous prime that defines a projective variety of
dimension $r$ over $\Phi$, the \emph{generic fibre} of $X'$ over $Q$.
\end{lemma}
\leanref{ReciprocityTheory/ReductiveGroups/ArithmeticGeometry/ArithmeticIntersection/MixedFieldContraction/ContractMixedFieldStatement.lean}{ReciprocityTheory.ChartRestriction.G.AI.I4.GenFibreExistsStatement}
\leanref{ReciprocityTheory/ReductiveGroups/ArithmeticGeometry/ArithmeticIntersection/MixedFieldContraction/ContractMixedFieldStatement.lean}{ReciprocityTheory.ChartRestriction.G.AI.I4.FibreDimStatement}
\begin{proof}
Write $K'_0=K'\setminus\kappa(K)$ for the set of the blocks not hit by $\kappa$.

\emph{Existence.} The homomorphism $v_*$ induces a homomorphism $R/Q\to R'/P'$, which is injective as $Q=v_*^{-1}(P')$.
Both rings are domains. Let $\Phi'$ be the fraction field of $R'/P'$. The injection extends to the fraction fields,
$\Phi\to\Phi'$, and $R'\to R'/P'\to\Phi'$ has kernel $P'$ and composite with $v_*$ equal to $R\to R/Q\to\Phi\to\Phi'$.

Now let $\Phi'$ and $R'\to\Phi'$ be any such, and let $\xi_{k',i'}\in\Phi'$ be the image of $x'_{k',i'}$. Then
$\xi_{v(k,i)}$ is the image of $x_{k,i}$ in $\Phi\subset\Phi'$, and $V$ is the kernel of $f\mapsto f(\xi)$ on
$\Phi[x'_{k',i'}:k'\in K'_0]$.

\emph{A relevant multihomogeneous prime.} $V$ is prime, as the kernel of a homomorphism to a field. For $k'\in K'_0$, some
$x'_{k',i'}$ is not in $P'$, as $P'$ is relevant; then $\xi_{k',i'}\neq0$ and $x'_{k',i'}\notin V$. Let $f\in V$. Its
coefficients lie in $\Phi$, the fraction field of $R/Q$, so there are $b\in R\setminus Q$ and $a_m\in R$, for the
monomials $m$ of $f$, such that $\bar b\,f=\sum_m\bar a_mm$, where bars denote images in $\Phi$. Put
$F=\sum_mv_*(a_m)m\in R'$. Its image in $\Phi'$ is $\bar b\,f(\xi)=0$, so $F\in P'$. Let $k'\in K'_0$ and $n\ge0$. The
$v_*(a_m)$ involve only variables of the blocks $\kappa(K)$, so the component of $F$ of degree $n$ in the variables of the
block $k'$ is $F_n=\sum_mv_*(a_m)m$, the sum over the $m$ of degree $n$ in that block. As $P'$ is generated by
multihomogeneous elements, $F_n\in P'$. Its image is $\bar b\,f_n(\xi)$, where $f_n$ is the component of $f$ of degree $n$
in the block $k'$, so $f_n\in V$. Applying this to one block of $K'_0$ after another, every multihomogeneous component of
$f$ lies in $V$, and $V$ is generated by multihomogeneous elements.

\emph{The dimension.} Let $E\subset\Phi'$ be the subfield generated by $\Phi$ and the $\xi_{k',i'}$ with $k'\in K'_0$. The
ring $\Phi[x'_{k',i'}:k'\in K'_0]/V$ is the $\Phi$-algebra generated by these $\xi_{k',i'}$, a domain with fraction field
$E$, so its Krull dimension is the transcendence degree of $E$ over $\Phi$ \cite[Tag 00P0]{Sta}. For an arithmetic
subvariety of dimension $e$ defined by a prime $P_0$ of a ring $\mathbf Z[x_{k,i}:k\in K_0,\ i\in I_k]$, the fraction
field of the quotient has transcendence degree $\dim\mathbf Z[x]/P_0-1=e+|K_0|-1$ over $\mathbf Q$, as in the proof of
Lemma~\ref{cv-origin-weight-birational} \cite[Tag 02JU]{Sta}. So $\Phi$ has transcendence degree $d+|K|-1$. The field $E$
contains every $\xi_{k',i'}$: for $k'=\kappa(k)$, every $i'\in I'_{\kappa(k)}$ is a $v(k,i)$, so
$\xi_{k',i'}\in\Phi$. Conversely $\Phi$ is generated by the images of the $x_{k,i}$, which are the $\xi_{v(k,i)}$. So $E$ is
the fraction field of the image of $R'$, which is isomorphic to the fraction field of $R'/P'$, of transcendence degree
$d+r+|K'|-1$ over $\mathbf Q$. By the additivity of transcendence degrees, $E$ has transcendence degree
$(d+r+|K'|-1)-(d+|K|-1)=r+|K'_0|$ over $\Phi$, as $\kappa$ is injective. So
$\dim\Phi[x'_{k',i'}:k'\in K'_0]/V=r+|K'_0|$, which says that $V$ defines a projective variety of dimension $r$ over
$\Phi$.
\end{proof}
\leanref{GroundTheory/HeightsAndTheta/Heights/Intersections/QuotientMapImageZInjective.lean}{GroundTheory.Heights₂.Heights.Inter.genFibreOf}
\leanref{GroundTheory/HeightsAndTheta/Heights/SemipositiveHeights/FibreDim.lean}{GroundTheory.Heights₂.Heights.ForSecondHalf.FibreDim.proof}
\leanref{GroundTheory/HeightsAndTheta/Heights/OrbitLimits/PFibreMultiHomog.lean}{GroundTheory.Heights₂.Heights.Targets.FibreMultiHomog.proof}
% Lean: Targets.genFibreExists = GenericPointExists.proof, genFibreOf (FnField of imageZ v P' and of P', genPt,
% quotientMap_imageZ_injective and IsFractionRing.lift). Giv.fibreDim = ForSecondHalf.FibreDim.proof: RingHom.ker_isPrime,
% Targets.FibreMultiHomog.proof (exists_clear, liftZ, liftZ_component and cyc_isMultiHomogIdeal_of_isHomogeneous: the
% block-degree components), relevance from that of P', EvalKerKrullDim.proof (ringKrullDim_eq_trdeg_fractionRing),
% trdeg_of_isArithVar (IntDomainDim.proof) for imageZ v P' and P', the field E = adjoin of the fibre coordinates with
% hkerE (its coordinates have kernel P'), trdeg_add_eq and the count |K'_0| + |K| = |K'|. The statement takes any
% GenFibre G: fields Phi and Phi' with x of kernel Q generating Phi by quotients, and x' of kernel P' compatible with v.
% Size: QuotientMapImageZInjective.lean (part), FibreDim.lean (230 lines), PFibreMultiHomog.lean and LiftZComponent.lean,
% about 300 lines.

\begin{lemma}\label{cv-origin-intersection-projection-generic-rational}
Let $K$, $(I_k)_{k\in K}$ and $R$ be as in Definition~\ref{lit-multiprojective}, let $(\cdot\mid\cdot)$ be an arithmetic
intersection system, let $X$ be an arithmetic subvariety of dimension $n+1$ defined by $P$, and let
$\bar L_1,\dots,\bar L_n$ and $\bar A=(a,h_A)$ be smooth metrized line bundles; put
$\Phi=c_1(\bar L_1)\wedge\dots\wedge c_1(\bar L_n)$. For multihomogeneous $s\in R\setminus P$ let
$\mathcal C(s)=\mathcal C(\bar L_1\cdots\bar L_n;s\mid X)$ be the sum over the relevant minimal primes $W$ of $P+(s)$,
with their multiplicities in $R/(P+(s))$, of $(\bar L_1\cdots\bar L_n\mid W)$ for horizontal $W$ and of
$\deg_{L_1\cdots L_n}(W)\log p$ for $W\cap\mathbf Z=p\mathbf Z$. Let $f,g\in R\setminus P$ be multihomogeneous of
multidegrees $\alpha,\beta\in\mathbf N^K$ with $\alpha-\beta=a$. Then
\[
 (\bar L_1\cdots\bar L_n\bar A\mid X)=\mathcal C(f)-\mathcal C(g)-\int_{X(\mathbf C)}\log\frac{|f|}{|g|\,h_A}\,\Phi.
\]
\end{lemma}
\leanref{ReciprocityTheory/ReductiveGroups/ArithmeticGeometry/ArithmeticIntersection/RationalSections/ThmISt.lean}{ReciprocityTheory.ChartRestriction.G.AI.I2.RatRStatement}
\leanref{GroundTheory/HeightsAndTheta/Heights/Intersections/PConsequences.lean}{GroundTheory.Heights₂.Heights.Inter.RatR.proof}
\begin{proof}
Write $\bar L$ for $\bar L_1\cdots\bar L_n$, and let $\bar B=\bar{\mathcal O}(\beta)=(\beta,h_\beta)$ be the
Fubini--Study bundle (Definition~\ref{cv-origin-analytic-system-pullback}(1)), so that $\bar A\otimes\bar B=(\alpha,h_Ah_\beta)$.
By the recursion with the sections $f$ and $g$, and the additivity,
\[
 (\bar L\bar A\mid X)=(\bar L(\bar A\otimes\bar B)\mid X)-(\bar L\bar B\mid X)
 =\mathcal C(f)-\mathcal C(g)-\Big(\int_{X(\mathbf C)}\log\frac{|f|}{h_Ah_\beta}\,\Phi
 -\int_{X(\mathbf C)}\log\frac{|g|}{h_\beta}\,\Phi\Big).
\]
Let $\Psi$ be the good projection chosen for $X$. By
Lemma~\ref{cv-origin-intersection-poincare-lelong-all-double-independent}(1) the two integrands, the fibre sums of the
sheet densities through $\Psi$, are integrable. As in the proof of Lemma~\ref{cv-origin-intersection-family-product}, for
almost every $y$ neither $f$ nor $g$ vanishes at a point $z\in\hat X(\mathbf C)$ with $\Psi(\sigma(z))=(1,y)$; at the
chosen sheets over such $y$,
$\log(|f|/(h_Ah_\beta))-\log(|g|/h_\beta)=\log(|f|/(|g|h_A))$. So the difference of the two integrals is
$\int_{X(\mathbf C)}\log(|f|/(|g|h_A))\,\Phi$.
\end{proof}
% Lean: RatR.proof under PtsIntegrableStatement: hR (the recursion IsIntersectionFamily) for A.tensor (sqMetric beta)
% with f (tensor_sqMetric_a) and for sqMetric beta with g, hAdd, and hint (integral_sub with the two integrability
% statements, ae_eval_ne_zero for f and g, and Real.log_div, Real.log_mul at the chosen sheets); recRRat is
% cycR f - cycR g - ptsIntegral (logNormRat A f g). Size: about 60 lines in PConsequences.lean.

\begin{definition}\label{cv-origin-intersection-projection-generic-cycles-fibres}
Let $(\kappa,v)$ be a block map (Definition~\ref{cv-origin-intersection-full}(1)), and let $W'$ be a prime ideal of $R'$.
Put $Q_{W'}=v_*^{-1}(W')$, a prime of $R$, and let $\Phi_{W'}$ be the fraction field of $R/Q_{W'}$. The construction of
Lemma~\ref{cv-origin-intersection-projection-fibre}, with $W'$ in place of $P'$, gives the kernel $V_{W'}$ of the
homomorphism $\Phi_{W'}[x'_{k',i'}:k'\notin\kappa(K)]$ to the fraction field of $R'/W'$; it is the \emph{generic fibre}
of $W'$ over $Q_{W'}$. Let $\deg(V_{W'})$ be the degree of the zero-dimensional variety defined by $V_{W'}$ over
$\Phi_{W'}$ \cite[Tag 0BEP]{Sta} if $V_{W'}$ defines one, and $0$ otherwise.
\end{definition}
\leanref{GroundTheory/HeightsAndTheta/Heights/Intersections/IsPrimeSpanXb.lean}{GroundTheory.Heights₂.Heights.Inter.GenericPointExists.Statement}

\begin{lemma}\label{cv-origin-intersection-projection-generic-cycles-horizontal}
Let $(\cdot\mid\cdot)$ be an arithmetic intersection system, let $d\ge1$, and let $(\kappa,v)$, $X'$, $P'$ and $Q$ be as
in Lemma~\ref{cv-origin-intersection-projection-fibre} with $r=0$, so that $X'$ and $Q$ have dimension $d$. Suppose that
the conclusion of Lemma~\ref{cv-origin-intersection-projection-generic} holds in dimension $d-1$, for every block map.
Let $\bar L_1,\dots,\bar L_{d-1}$ be smooth metrized line bundles on $\prod_k\mathbf P^{I_k}_{\mathbf Z}$ with pullbacks
$\bar L'_1,\dots,\bar L'_{d-1}$, let $s\in R\setminus Q$ be multihomogeneous, and let $W'$ be a horizontal relevant
minimal prime of $P'+(v_*(s))$, with $Q_{W'}$ and $V_{W'}$ as in
Definition~\ref{cv-origin-intersection-projection-generic-cycles-fibres}. Then $Q_{W'}$ is horizontal, and either
$Q_{W'}$ is a relevant minimal prime of $Q+(s)$ and
$(\bar L'_1\cdots\bar L'_{d-1}\mid W')=\deg(V_{W'})\,(\bar L_1\cdots\bar L_{d-1}\mid Q_{W'})$, or it is not and
$(\bar L'_1\cdots\bar L'_{d-1}\mid W')=0$.
\end{lemma}
\leanref{GroundTheory/HeightsAndTheta/Heights/Intersections/IsHorizontalSpanXb.lean}{GroundTheory.Heights₂.Heights.Inter.ArithmeticIntersectionHorizontal.Statement}
\leanref{GroundTheory/HeightsAndTheta/Heights/Intersections/IsHorizontalImageZ.lean}{GroundTheory.Heights₂.Heights.Inter.ArithmeticIntersectionHorizontal.proof}
\begin{proof}
As $W'\cap\mathbf Z=0$ and $v_*$ is the identity on $\mathbf Z$, $Q_{W'}\cap\mathbf Z=0$. The prime $W'$ is an arithmetic
subvariety of dimension $d-1$, as for the components in Definition~\ref{lit-multiprojective}(3). By
Lemma~\ref{cv-origin-intersection-contraction-image}(1), $Q_{W'}$ is an arithmetic subvariety of some dimension
$e\le d-1$.

If $e<d-1$, then $(\bar L'_1\cdots\bar L'_{d-1}\mid W')=0$ by Lemma~\ref{cv-origin-intersection-contraction-pulled}, and
$Q_{W'}$ is not a relevant minimal prime of $Q+(s)$, as the horizontal ones are arithmetic subvarieties of dimension
$d-1$.

Let $e=d-1$. The prime $Q_{W'}$ contains $Q$ and $s$, as $W'$ contains $P'$ and $v_*(s)$. Let $M\subset Q_{W'}$ be a
minimal prime of $Q+(s)$. It is relevant and horizontal, as $Q_{W'}$ is, so it is an arithmetic subvariety of dimension
$d-1$. If some $x\in Q_{W'}$ were not in $M$, then $\dim R/Q_{W'}+1\le\dim R/M$, as the image of $x$ in the domain $R/M$ is
nonzero, against the equality of the dimensions. So $M=Q_{W'}$, and $Q_{W'}$ is a relevant minimal prime of $Q+(s)$. Now
$W'$ and $Q_{W'}=v_*^{-1}(W')$ are arithmetic subvarieties of dimension $d-1$, so the conclusion of
Lemma~\ref{cv-origin-intersection-projection-generic} in dimension $d-1$, for $(\kappa,v)$ and $W'$, gives
$(\bar L'_1\cdots\bar L'_{d-1}\mid W')=\deg(V_{W'})\,(\bar L_1\cdots\bar L_{d-1}\mid Q_{W'})$.
\end{proof}
% Lean: ArithmeticIntersectionHorizontal.proof with divSuppArith (hW'n), imageArith (d <= n), isHorizontal_imageZ; the
% case d < n through dim_eq and PushforwardProjectivePoint.full (contract_intersection_number, Lemma
% cv-origin-intersection-contraction-pulled); the case d = n through mem_minimalPrimes (dim_add_one_le,
% isRelevant_of_le) and the hypothesis hn = GenDegAt I n. Size: IsHorizontalImageZ.lean, about 160 lines, besides the
% lemmas cited.

\begin{lemma}\label{cv-origin-intersection-projection-generic-cycles-vertical-degree}
Let $(\kappa,v)$ be a block map (Definition~\ref{cv-origin-intersection-full}(1)), let $p$ be a prime number, let
$e\ge0$, and let $W'$ be a relevant multihomogeneous prime of $R'$ with $p\in W'$ and $\dim R'/W'=e+|K'|$, such that
$W=v_*^{-1}(W')$ has $\dim R/W=e+|K|$. Let $\xi_{k',i}$ be the image of $x'_{k',i}$ in the fraction field of $R'/W'$,
and let $\mathbf F_p(\xi)$ and $\mathbf F_p(\xi\circ v)$ be its subfields generated by the ratios $\xi_{k',i}/\xi_{k',j}$,
and $\xi_{v(k,i)}/\xi_{v(k,j)}$, with nonzero denominators. Then the generic fibre $V_{W'}$ of $W'$ over $W$
(Definition~\ref{cv-origin-intersection-projection-generic-cycles-fibres}) defines a zero-dimensional variety, and
$\deg(V_{W'})=[\mathbf F_p(\xi):\mathbf F_p(\xi\circ v)]$, which is finite.
\end{lemma}
\leanref{GroundTheory/HeightsAndTheta/Heights/Intersections/ImageVarIdBotTwo.lean}{GroundTheory.Heights₂.Heights.Inter.VerticalDegree.Statement}
\leanref{GroundTheory/HeightsAndTheta/Heights/Intersections/ImportsRest.lean}{GroundTheory.Heights₂.Heights.Inter.VerticalDegree.proof}
\leanref{GroundTheory/HeightsAndTheta/Heights/Intersections/ImportsTer.lean}{GroundTheory.Heights₂.Heights.Inter.FibreDimensionZero.proof}
\leanref{GroundTheory/HeightsAndTheta/Heights/Intersections/ChartCongr.lean}{GroundTheory.Heights₂.Heights.Inter.VerticalDegreeAfterAdjoining.proof}
\begin{proof}
Write $K'_0=K'\setminus\kappa(K)$, $E$ for the fraction field of $R'/W'$, and $\Phi_{W'}\subset E$ for that of $R/W$,
embedded through $v_*$. The relevant primes: as $W'$ is relevant and $v$ maps $I_k$ onto $I'_{\kappa(k)}$, $W$ is
relevant. The rings $R/W$ and $R'/W'$ are domains of finite type over $\mathbf F_p$, so their dimensions $e+|K|$ and
$e+|K'|$ are the transcendence degrees of $\Phi_{W'}$ and $E$ over $\mathbf F_p$ \cite[Lemma 10.116.1, Tag 00P0]{Sta}.

\emph{The dimension.} The parts ``A relevant multihomogeneous prime'' and ``The dimension'' of the proof of
Lemma~\ref{cv-origin-intersection-projection-fibre} use only that $W'$ is a relevant multihomogeneous prime and the
transcendence degrees of the two fields; with the degrees above they show that $V_{W'}$ is a relevant multihomogeneous
prime and that $\Phi_{W'}[x'_{k',i}:k'\in K'_0]/V_{W'}$ has dimension
$(e+|K'|)-(e+|K|)=|K'_0|$. So $V_{W'}$ defines a zero-dimensional projective variety over $\Phi_{W'}$.

\emph{The degree.} Its degree is the dimension over $\Phi_{W'}$ of its coordinate ring in a chart. The steps ``The
chart'' and ``The fibre and the ratios of $\xi$'' of the proof of Lemma~\ref{cv-origin-intersection-projection-count},
with $\mathbf F_p$ in place of $\mathbf Q$ and $W'$ in place of $P'$, use of the fields only that the kernel of
$\mathbf F_p[x'_{k',i}]\to E$, $f\mapsto f(\xi)$, is generated by multihomogeneous elements, which holds as it is the image
of $W'$. They show that this dimension is $[\mathbf F_p(\xi):\mathbf F_p(\xi\circ v)]$, the two being finite together; and
it is finite, as the coordinate ring of a zero-dimensional variety in a chart is a domain of dimension zero of finite type
over $\Phi_{W'}$.
\end{proof}
% Lean: VerticalDegree.proof: FibreDimensionZero.proof (ArithmeticIntersectionFibreMultiHomogeneous, EvalKerKrullDim and
% FibreTranscendenceDegreeAlt over ZMod p), hJ.1 for the value in the chart i', ChartRankAdjoin (Lemma
% cv-origin-intersection-projection-count-chart) and VerticalDegreeAfterAdjoining.proof (RankAdjoinFrac, FnDegAdjoin,
% ChartScalingIndep, AdjoinIndepRelfinrank, FinrankAdjoinRange: the argument of Lemma
% cv-origin-intersection-projection-count over ZMod p). The Lean's fnDeg is the rank of the chart ring of W' mod p over
% that of its image, which FnDegAdjoin and RankAdjoinFrac identify with the degree of the ratio fields. Size: about 1,000
% lines in ImportsRest.lean, ImportsTer.lean, ChartCongr.lean, PRankAdjoinFrac.lean, PFnDegAdjoin.lean, AevalMemAdjoin.lean
% and the files they use, besides the lemmas shared with Lemma cv-origin-intersection-projection-count.

\begin{lemma}\label{cv-origin-intersection-projection-generic-cycles-vertical}
Let $d\ge1$, and let $(\kappa,v)$, $X'$, $P'$ and $Q$ be as in Lemma~\ref{cv-origin-intersection-projection-fibre} with
$r=0$, so that $X'$ and $Q$ have dimension $d$. Let $a_1,\dots,a_{d-1}\in\mathbf Z^K$, and let $a'_j\in\mathbf Z^{K'}$
be $a_j$ extended by $0$, as in Definition~\ref{cv-origin-intersection-full}(1). Let $s\in R\setminus Q$ be
multihomogeneous, and let $W'$ be a relevant minimal prime of $P'+(v_*(s))$ with $W'\cap\mathbf Z=p\mathbf Z$ for a prime
number $p$, with $Q_{W'}$ and $V_{W'}$ as in Definition~\ref{cv-origin-intersection-projection-generic-cycles-fibres}.
Then either $Q_{W'}$ is a relevant minimal prime of $Q+(s)$ and
$\deg_{a'_1\cdots a'_{d-1}}(W')=\deg(V_{W'})\deg_{a_1\cdots a_{d-1}}(Q_{W'})$, or it is not and
$\deg_{a'_1\cdots a'_{d-1}}(W')=0$, with the degrees of Lemma~\ref{cv-origin-intersection-family-sections-degree}.
\end{lemma}
\leanref{GroundTheory/HeightsAndTheta/Heights/Intersections/NotIsHorizontalSpanTwo.lean}{GroundTheory.Heights₂.Heights.Inter.ArithmeticIntersectionVertical.Statement}
\leanref{GroundTheory/HeightsAndTheta/Heights/Intersections/ImportsAlt.lean}{GroundTheory.Heights₂.Heights.Inter.ArithmeticIntersectionVertical.proof}
\begin{proof}
Put $e=d-1$ and $W=Q_{W'}=v_*^{-1}(W')$. The prime $W'$ is multihomogeneous \cite[Tag 00JU]{Sta}, as in the proof of
Lemma~\ref{cv-origin-intersection-family-curves-vertical-closed}(4), and by (1) of that lemma, for $R'$,
$\dim R'/W'=\dim R'/P'-1=e+|K'|$. So $W'$ defines a projective variety $Z$ of dimension $e$ over $\mathbf F_p$, and
$\deg_{a'_1\cdots a'_e}(W')=(\mathcal O(a'_1)\cdots\mathcal O(a'_e)\cdot Z)$
(Lemma~\ref{cv-origin-intersection-family-sections-degree}). The prime $W$ contains $p$, and it is relevant and
multihomogeneous, as $v$ maps $I_k$ onto $I'_{\kappa(k)}$ and $v_*$ maps a multihomogeneous element of multidegree $a$ to
one of multidegree $a'$; it defines a projective variety $Y$ over $\mathbf F_p$ of dimension $\dim R/W-|K|$.

Let $f:\prod_{k'}\mathbf P^{I'_{k'}}_{\mathbf F_p}\to\prod_k\mathbf P^{I_k}_{\mathbf F_p}$ be the morphism
$(z'_{k',i})\mapsto(z'_{v(k,i)})_{k,i}$, the projection to the factors $\kappa(k)$ followed by the isomorphisms given by
the bijections $v:I_k\to I'_{\kappa(k)}$. Then $f^*\mathcal O(a)=\mathcal O(a')$ for $a\in\mathbf Z^K$, and $f(Z)=Y$: the
image of the integral scheme $Z$ is the closed subvariety defined by the kernel of
$\mathbf F_p[x_{k,i}]\to\mathbf F_p[x'_{k',i}]/W'$, $g\mapsto v_*(g)$, which is the reduction of $v_*^{-1}(W')=W$, as
$p\in W$. By the projection formula \cite[Lemma 33.45.7, Tag 0BET]{Sta},
\[
 \deg_{a'_1\cdots a'_e}(W')=[\kappa(Z):\kappa(Y)]\,(\mathcal O(a_1)\cdots\mathcal O(a_e)\cdot Y)
\]
if $\dim Y=e$, where $\kappa(Z)$ and $\kappa(Y)$ are the function fields, and $\deg_{a'_1\cdots a'_e}(W')=0$ if
$\dim Y<e$.

Suppose that $W$ is a relevant minimal prime of $Q+(s)$. By
Lemma~\ref{cv-origin-intersection-family-curves-vertical-closed}(1), $\dim R/W=\dim R/Q-1=e+|K|$, so $\dim Y=e$ and the
second factor is $\deg_{a_1\cdots a_e}(W)$. The function fields $\kappa(Z)$ and $\kappa(Y)$ are generated over
$\mathbf F_p$ by the ratios of the coordinates, so they are the fields $\mathbf F_p(\xi)$ and $\mathbf F_p(\xi\circ v)$
of Lemma~\ref{cv-origin-intersection-projection-generic-cycles-vertical-degree}, which gives
$[\kappa(Z):\kappa(Y)]=\deg(V_{W'})$.

Otherwise, as $v_*(s)\in W'$, the prime $W$ contains $Q+(s)$, and so it contains a minimal prime $M\neq W$ of $Q+(s)$. By
Lemma~\ref{cv-origin-intersection-family-curves-vertical-closed}(1), $\dim R/M=\dim R/Q-1=e+|K|$ and
$\dim R/W\le\dim R/M-1$. So $\dim Y<e$, and $\deg_{a'_1\cdots a'_e}(W')=0$.
\end{proof}
% Lean: ArithmeticIntersectionVertical.of_remaining = ArithmeticIntersectionVertical.proof with divSuppArith (hW'd),
% geoNumStatement (geoNum as the field intersection numbers over ZMod p, as in Lemma
% cv-origin-intersection-family-sections-degree), isMultiProjPrime_imageZ, imageVar_red (the image of W' mod p), and in
% the first case the projection formula ProjDegStatement over ZMod p (from ForSecondHalf.projDeg_all) and
% VerticalDegree.full (Lemma cv-origin-intersection-projection-generic-cycles-vertical-degree); in the second case
% Ideal.exists_minimalPrimes_le, dim_quot_minimal_add_one, krullDim_quot_add_one_le, exists_dim_F and
% Targets.ProjContract.full, the vanishing for an image of smaller dimension. The Lean proves the projection formula and
% the vanishing for its field intersection numbers (Targets.ExistsUnique.full, ProjDeg.of_found through
% DivPushForward and FnDegRelDeg, ProjContract.full); the text cites Stacks Tag 0BET for both. Size: ImportsAlt.lean,
% about 90 lines, besides the projection formula over ZMod p (most of the 11,000 lines that the Lean's proof uses, with
% Targets.ExistsUnique.full) and Lemma cv-origin-intersection-projection-generic-cycles-vertical-degree.

\begin{situation}\label{cv-origin-intersection-projection-generic-cycles-pushforward-situation}
Let $d\ge1$, and let $(\kappa,v)$, $X'$, $P'$, $Q$, $\Phi$ and $V$ be as in
Lemma~\ref{cv-origin-intersection-projection-fibre} with $r=0$, so that $X'$ and $Q$ are arithmetic subvarieties of
dimension $d$. Write $K'_0=K'\setminus\kappa(K)$. Let $s\in R\setminus Q$ be multihomogeneous of multidegree $\alpha$, and
let $W$ be a relevant minimal prime of $Q+(s)$. Let $F_W$ be the set of the relevant multihomogeneous primes $W'$ of $R'$
with $P'\subset W'$ and $v_*^{-1}(W')=W$. Let $A=\mathcal O_{Q,W}$
(Definition~\ref{cv-origin-intersection-family-curves-vertical-ring}), and let $B$ be the subring of the fraction field of
$R'/P'$ generated by the quotients $u/w$ of multihomogeneous $u,w\in R'$ of the same multidegree with $w\notin P'$ and
$w$ in no $W'\in F_W$. The injection $R/Q\to R'/P'$ induced by $v_*$ extends to the fraction fields and maps $A$ into $B$:
for $u,w\in R$ multihomogeneous of the same multidegree with $w\notin W$, $v_*(u)$ and $v_*(w)$ are multihomogeneous of
the same multidegree, $v_*(w)\notin P'$ as $w\notin Q$, and $v_*(w)$ lies in no $W'\in F_W$ as $v_*^{-1}(W')=W$. This
makes $B$ an $A$-algebra.
\end{situation}
\leanref{GroundTheory/HeightsAndTheta/Heights/Intersections/StatementsSemilocRing.lean}{GroundTheory.Heights₂.Heights.Inter.fibreOver}
\leanref{GroundTheory/HeightsAndTheta/Heights/Intersections/StatementsSemilocRing.lean}{GroundTheory.Heights₂.Heights.Inter.semiLocRing}
\leanref{GroundTheory/HeightsAndTheta/Heights/Intersections/StatementsSemilocRing.lean}{GroundTheory.Heights₂.Heights.Inter.locToSemiloc}

\begin{lemma}\label{cv-origin-intersection-projection-generic-cycles-pushforward-fibres}
In Situation~\ref{cv-origin-intersection-projection-generic-cycles-pushforward-situation}:
\begin{enumerate}
\item If $W_1,W_2\in F_W$ and $W_1\subset W_2$, then $W_1=W_2$.
\item $F_W$ is the set of the relevant minimal primes $W'$ of $P'+(v_*(s))$ with $v_*^{-1}(W')=W$. In particular $F_W$
  is finite.
\end{enumerate}
\end{lemma}
\leanref{GroundTheory/HeightsAndTheta/Heights/Intersections/StatementsFibreIncomp.lean}{GroundTheory.Heights₂.Heights.Inter.FibreIncomp.Statement}
\leanref{GroundTheory/HeightsAndTheta/Heights/SemipositiveHeights/FibreIncomp.lean}{GroundTheory.Heights₂.Heights.ForSecondHalf.FibreIncomp.full}
\leanref{GroundTheory/HeightsAndTheta/Heights/Intersections/PFibreDivSupp.lean}{GroundTheory.Heights₂.Heights.Inter.FibreDivSupp.proof}
\leanref{GroundTheory/HeightsAndTheta/Heights/Intersections/PFibreOverFinite.lean}{GroundTheory.Heights₂.Heights.Inter.FibreOverFinite.proof}
\begin{proof}
\emph{A dimension bound.} Let $W_2$ be a relevant multihomogeneous prime of $R'$ and $U=v_*^{-1}(W_2)$. We show
$\dim R/U+|K'_0|\le\dim R'/W_2$. Let $F_0$ be $\mathbf Q$ if $W_2\cap\mathbf Z=0$ and $\mathbf F_p$ if
$W_2\cap\mathbf Z=p\mathbf Z$; then $R/U$ and $R'/W_2$ are domains of finite type over $\mathbf Z$ that are torsion-free
or of characteristic $p$, and their dimensions are the transcendence degrees of their fraction fields over $F_0$, plus
one in the first case \cite[Tags 00P0 and 02JU]{Sta}, as in the proof of
Lemma~\ref{cv-origin-intersection-projection-fibre}. The homomorphism $v_*$ induces an injection $R/U\to R'/W_2$. For
$k'\in K'_0$ choose $i_{k'}$ with $x'_{k',i_{k'}}\notin W_2$, and let $y_{k'}\neq0$ be its image in $R'/W_2$. As $W_2$ is
multihomogeneous, the homomorphism $\theta:R'/W_2\to(R'/W_2)[T_{k'}:k'\in K'_0]$ that multiplies the image of
$x'_{k',i}$ by $T_{k'}$ for $k'\in K'_0$, and fixes the images of the other variables, is well defined; it multiplies the
image of a multihomogeneous element of multidegree $a$ by $\prod_{k'\in K'_0}T_{k'}^{a_{k'}}$, and it fixes the image of
$R/U$, whose variables lie in the blocks $\kappa(K)$. If $\sum_\beta c_\beta y^\beta=0$ with $c_\beta$ in the image of
$R/U$, applying $\theta$ gives $\sum_\beta c_\beta y^\beta T^\beta=0$, so every $c_\beta y^\beta$ and every $c_\beta$ is
$0$. So the $y_{k'}$ are algebraically independent over the fraction field of $R/U$, and the transcendence degrees give
the bound.

(1) Suppose that $W_1\subsetneq W_2$. As $s\in W=v_*^{-1}(W_1)$, $v_*(s)\in W_1$, while $v_*(s)\notin P'$, as
$s\notin Q$. By Lemma~\ref{cv-origin-intersection-family-curves-vertical-closed}(1), for $R'$ as well as for $R$,
$\dim R'/W_2\le\dim R'/W_1-1\le\dim R'/P'-2=d+|K'|-2$ and $\dim R/W=\dim R/Q-1=d+|K|-1$. The dimension bound for $W_2$,
with $v_*^{-1}(W_2)=W$, gives $d+|K|-1+|K'_0|\le d+|K'|-2$. As $\kappa$ is injective, $|K'|=|K|+|K'_0|$, a
contradiction.

(2) Let $W'\in F_W$. It contains $P'$ and $v_*(s)$, so it contains a minimal prime $U$ of $P'+(v_*(s))$. The prime $U$ is
multihomogeneous \cite[Tag 00JU]{Sta}, as in the proof of
Lemma~\ref{cv-origin-intersection-family-curves-vertical-closed}(4), and relevant, as $W'$ is. Its inverse image
$v_*^{-1}(U)$ is a prime that contains $Q+(s)$ and lies in $W$, so it is $W$, as $W$ is minimal over $Q+(s)$. So
$U\in F_W$, and $U=W'$ by (1). Conversely, a relevant minimal prime $W'$ of $P'+(v_*(s))$ is multihomogeneous by the
same reference and contains $P'$, so it lies in $F_W$ if $v_*^{-1}(W')=W$. The noetherian ring $R'/(P'+(v_*(s)))$ has
finitely many minimal primes.
\end{proof}
% Lean: FibreIncomp.full (ForSecondHalf): ringKrullDim_succ_le_of_surjective twice (hB, h21), divSuppArith for dim R/W,
% dim_add_card_le (the dimension bound: IntDomainDim.proof for horizontal W2; the reduction red p, isMultiProjPrime_red
% and exists_dim_trdeg_field over ZMod p for vertical W2; trdeg_add_card_le with algebraicIndependent_forgotten, the
% scaling homomorphism scaleWt) and card_add_card_forgottenBlocks. FibreDivSupp.proof (Ideal.exists_minimalPrimes_le,
% cyc_isMultiHomogIdeal_of_mem_minimalPrimes, hQi and FibreIncomp) and FibreOverFinite.proof (divSupp_finite). Size:
% FibreIncomp.lean (161 lines), FibreIncompTrdeg.lean (175), PFibreDivSupp.lean and PFibreOverFinite.lean, about 400 lines.

\begin{lemma}\label{cv-origin-intersection-projection-generic-cycles-pushforward-local}
Let $K$, $(I_k)_{k\in K}$ and $R$ be as in Definition~\ref{lit-multiprojective}, let $P$ be a relevant multihomogeneous
prime of $R$, let $s\in R\setminus P$ be multihomogeneous of multidegree $a$, let $W$ be a relevant minimal prime of
$P+(s)$, and let $t\in R\setminus W$ be multihomogeneous of multidegree $a$. Write $\mathcal O_W$ for $\mathcal O_{P,W}$
(Definition~\ref{cv-origin-intersection-family-curves-vertical-ring}). Then $\mathcal O_W$ is a noetherian local domain
of dimension $1$, $y=s/t$ is a nonzero element of $\mathcal O_W$, and the multiplicity $m_W$ of $W$ in $R/(P+(s))$ is
the length of $\mathcal O_W/y\mathcal O_W$.
\end{lemma}
\leanref{GroundTheory/HeightsAndTheta/Heights/Intersections/StatementsLocRingDimOne.lean}{GroundTheory.Heights₂.Heights.Inter.LocRingDimOne.Statement}
\leanref{GroundTheory/HeightsAndTheta/Heights/Intersections/PLocRingDimOne.lean}{GroundTheory.Heights₂.Heights.Inter.LocRingDimOne.proof}
\leanref{GroundTheory/HeightsAndTheta/Heights/Intersections/ImportsTerBis.lean}{GroundTheory.Heights₂.Heights.Inter.regularityDehomogenization_holds.Statement}
\begin{proof}
The prime $W$ is multihomogeneous \cite[Tag 00JU]{Sta}, as in the proof of
Lemma~\ref{cv-origin-intersection-family-curves-vertical-closed}(4). The proofs of (1), (2) and (5) of
Lemma~\ref{cv-origin-intersection-family-curves-vertical-component} use only that $P$ is a relevant multihomogeneous prime
and $W$ a relevant multihomogeneous minimal prime of $P+(s)$, through
Lemma~\ref{cv-origin-intersection-family-curves-vertical-local}; so they apply here. By (2), $\mathcal O_W$, a subring
of a field, is a noetherian local domain of dimension $1$; $y\neq0$ as $s\notin P$; and by (5), $m_W$ is the length of
$\mathcal O_W/y\mathcal O_W$.
\end{proof}
% Lean: LocRingDimOne.proof (LocRingFrac, LocRingNoeth, LocRingLocal, LocRingPow and
% Ideal.height_le_one_of_isPrincipal_of_mem_minimalPrimes, with t the coordinate product coordProd) and
% regularityDehomogenization_holds.full, both stated for any IsMultiProjPrime P. Size: PLocRingDimOne.lean, 55 lines,
% besides the lemmas of Lemma cv-origin-intersection-family-curves-vertical-component.

\begin{lemma}\label{cv-origin-intersection-projection-generic-cycles-pushforward-order}
Let $A$ be a noetherian local domain of dimension $1$ with fraction field $F$, let $M$ be a finite torsion-free
$A$-module, and let $y\in A$ be nonzero. Then $M/yM$ has finite length, and
\[
 \ell_A(M/yM)=\dim_F(M\otimes_AF)\,\ell_A(A/yA).
\]
\end{lemma}
\leanref{ReciprocityTheory/ReductiveGroups/ArithmeticGeometry/ArithmeticIntersection/LocalLengthAdditivityTorsionFree.lean}{ReciprocityTheory.ChartRestriction.G.AI.LocalLengthAdditivityTorsionFree}
\leanref{ReciprocityTheory/ReductiveGroups/ArithmeticGeometry/ArithmeticIntersection/KrullDimQuotientEqOneIffEqBot.lean}{ReciprocityTheory.ChartRestriction.G.AI.dimOneTorsionFreeFormula}
\begin{proof}
The module $(M/yM)\otimes_AF$ is $0$, as $y$ is a unit of $F$, so $M/yM$ is a finite module supported at the maximal
ideal and has finite length. Multiplication by $y$ on $M$ is injective, as $M$ is torsion-free. The primes $\mathfrak q$
of $A$ with $\dim A/\mathfrak q=1$ are the prime $0$ alone, as $A$ is a local domain of dimension $1$. So
\cite[Lemma 42.3.1, Tag 02QF]{Sta}, for $M$ and $x=y$, gives
$\ell_A(M/yM)-0=\ell_A(A/yA)\,\ell_F(M\otimes_AF)$, the multiplicity of the statement being
$\ell_A(M/yM)-\ell_A(\ker y)$, the order of $y$ on $A/0$ being $\ell_A(A/yA)$, and $M_{(0)}=M\otimes_AF$.
\end{proof}
% Lean: dimOneTorsionFreeFormula: dimOneFormula (the Lean's own proof of Stacks Tag 02QF, Lemma
% cv-origin-intersection-family-sections-cycles-local) with length_coker_ne_top, length_ker_eq_zero_of_torsionFree,
% krullDim_quotient_eq_one_iff_eq_bot for hset, length_localization_bot_eq_finrank and ordQuotient_bot. Size:
% KrullDimQuotientEqOneIffEqBot.lean, about 120 lines, besides LengthCokerEqKerAddDimOneOrderSum.lean.

\begin{lemma}\label{cv-origin-intersection-projection-generic-cycles-pushforward-integral-centre}
Let $K$, $(I_k)_{k\in K}$ and $R$ be as in Definition~\ref{lit-multiprojective}, let $P$ be a relevant multihomogeneous
prime of $R$, let $L$ be the fraction field of $R/P$, and let $\mathcal V\subset L$ be a valuation ring of $L$, with
maximal ideal $\mathfrak m_{\mathcal V}$. For $f\in R$ write $\bar f$ for its class in $L$. Then there is a relevant
multihomogeneous prime $Y\supset P$ of $R$, a \emph{centre} of $\mathcal V$, such that for all multihomogeneous $f,g\in R$
of the same multidegree with $g\notin Y$, $\bar f/\bar g$ lies in $\mathcal V$, and it lies in $\mathfrak m_{\mathcal V}$
if and only if $f\in Y$.
\end{lemma}
\leanref{GroundTheory/HeightsAndTheta/Heights/Intersections/StValCentreZ.lean}{GroundTheory.Heights₂.Heights.Inter.ValCentreZ.Statement}
\leanref{GroundTheory/HeightsAndTheta/Heights/Intersections/StValCentreZ.lean}{GroundTheory.Heights₂.Heights.Inter.IsValCentre}
\leanref{GroundTheory/HeightsAndTheta/Heights/Intersections/PValCentreZ.lean}{GroundTheory.Heights₂.Heights.Inter.ValCentreZ.proof}
\begin{proof}
\emph{A chart.} For $a,b\in L^\times$, $a/b$ or $b/a$ lies in $\mathcal V$ \cite[Tag 00IB]{Sta}, and if $a/b$ and $b/c$
lie in $\mathcal V$, so does $a/c$. So for each $k\in K$ the finitely many $\bar x_{k,i}\neq0$, which exist as $P$ is
relevant, have one, $\bar x_{k,j_k}$, with $r_{k,i}=\bar x_{k,i}/\bar x_{k,j_k}\in\mathcal V$ for every $i$ with
$\bar x_{k,i}\neq0$; put $r_{k,i}=0$ if $\bar x_{k,i}=0$. Let $\theta:R\to\mathcal V$ be the homomorphism with
$\theta(x_{k,i})=r_{k,i}$, $E=\mathcal V/\mathfrak m_{\mathcal V}$ the residue field, and $\psi:R\to E$ the composite of
$\theta$ with the reduction. Then $\psi(x_{k,j_k})=1$. For $g$ multihomogeneous of multidegree $a$,
$\bar x_{k,i}=r_{k,i}\bar x_{k,j_k}$ gives $\bar g=\theta(g)\,m_a$ with $m_a=\prod_k\bar x_{k,j_k}^{a_k}\neq0$.

\emph{The prime.} Let $Y$ be the kernel of $\varphi_\emptyset:R\to E[t_k:k\in K]$, $x_{k,i}\mapsto\psi(x_{k,i})t_k$, of
Lemma~\ref{cv-origin-intersection-family-curves-vertical-chain}, for $A_0=\mathbf Z$. By (1) of that lemma,
$\varphi_\emptyset(g)=\psi(g)\prod_kt_k^{a_k}$ for $g$ multihomogeneous of multidegree $a$; so such a $g$ lies in $Y$ if
and only if $\psi(g)=0$, that is, $\theta(g)\in\mathfrak m_{\mathcal V}$. The ideal $Y$ is prime, as $E[t]$ is a domain.
It is multihomogeneous: $\varphi_\emptyset$ maps the component of multidegree $a$ of an element into $E\prod_kt_k^{a_k}$, so
the components of an element of $Y$ lie in $Y$. It is relevant, as $\theta(x_{k,j_k})=1$ is not in
$\mathfrak m_{\mathcal V}$. It contains $P$: $P$ is generated by multihomogeneous elements $g$, and for these
$\theta(g)m_a=\bar g=0$, so $\theta(g)=0$.

\emph{The quotients.} Let $f,g$ be multihomogeneous of multidegree $b$ with $g\notin Y$. Then
$\bar f/\bar g=\theta(f)m_b/(\theta(g)m_b)=\theta(f)/\theta(g)$, and $\theta(g)\in\mathcal V\setminus\mathfrak
m_{\mathcal V}$ is a unit of $\mathcal V$. So $\bar f/\bar g\in\mathcal V$, and it lies in $\mathfrak m_{\mathcal V}$ if
and only if $\theta(f)$ does, that is, if and only if $f\in Y$.
\end{proof}
% Lean: ValCentreZ.proof: hj (j k maximizing R.valuation of the coordinates outside P, Finset.exists_max_image), r p from
% exists_funField_mul_eq with hrle (r p in R), chartHom r with hmem, psiE the residue map (IsLocalRing.residue), hform
% from RegularityFunctionFieldNumber.specMap_form (Lemma cv-origin-intersection-family-curves-vertical-chain (1)),
% RegularityValuationCentre.ker_isMultiHomogIdeal (specMap_component), hPW through chartHom_mul and toFrac_prod_ne_zero,
% and the quotients through e and hvg. The Lean's IsValCentre also lists IsMultiProjPrime y (prime, multihomogeneous,
% relevant).

\begin{lemma}\label{cv-origin-intersection-projection-generic-cycles-pushforward-integral-centre-fibre}
In Situation~\ref{cv-origin-intersection-projection-generic-cycles-pushforward-situation}, let $L$ be the fraction field
of $R'/P'$, let $\mathcal V\subset L$ be a valuation ring that contains the image of $A$, and let $Y$ be a centre of
$\mathcal V$ (Lemma~\ref{cv-origin-intersection-projection-generic-cycles-pushforward-integral-centre}, for $R'$ and
$P'$). Then $Y=P'$ or $Y\in F_W$. The proof does not use that $s$ is multihomogeneous or that $s\notin Q$.
\end{lemma}
\leanref{GroundTheory/HeightsAndTheta/Heights/Intersections/StCentreOverW.lean}{GroundTheory.Heights₂.Heights.Inter.CentreOverW.Statement}
\leanref{GroundTheory/HeightsAndTheta/Heights/Intersections/PCentreOverW.lean}{GroundTheory.Heights₂.Heights.Inter.CentreOverW.proof}
\begin{proof}
Write $\bar f$ for the class in $L$ of $f\in R'$, and put $Y_0=v_*^{-1}(Y)$, a prime of $R$ containing
$Q=v_*^{-1}(P')$. It is relevant: for $k\in K$ some $x'_{\kappa(k),i'}$ is not in $Y$, and $i'=v(k,i)$ for some
$i\in I_k$, so $x_{k,i}\notin Y_0$. For $f\in R$ and $D\in\mathbf N^K$ let $f_D$ be the component of $f$ of multidegree
$D$; as in the proof of Lemma~\ref{cv-origin-intersection-projection-generic-integral-section-count-elimination}, $v_*(f_D)$
is the component of $v_*(f)$ of multidegree $D'$, where $D'$ is $D$ on $\kappa(K)$ and $0$ elsewhere.

\emph{$Y_0\subset W$.} Let $a\in Y_0$, and suppose $a\notin W$. As $W$ is multihomogeneous, some $c=a_D$ is not in $W$.
It is multihomogeneous of multidegree $D$, and $v_*(c)\in Y$, as $Y$ is multihomogeneous and contains $v_*(a)$. As $Y_0$ is
relevant, there are $i_k$ with $x_{k,i_k}\notin Y_0$, and $u=\prod_kx_{k,i_k}^{D_k}$, of multidegree $D$, is not in the
prime $Y_0$. As $c\notin W$, $u/c$ lies in $A$, so $\overline{v_*(u)}/\overline{v_*(c)}$ lies in $\mathcal V$. As
$v_*(u)\notin Y$ and $v_*(c)\in Y$, both multihomogeneous of multidegree $D'$, $\overline{v_*(c)}/\overline{v_*(u)}$ lies in
$\mathfrak m_{\mathcal V}$ by the property of the centre; the two classes are nonzero, as $c\notin Q$ and $u\notin Q$.
So their product $1$ lies in $\mathfrak m_{\mathcal V}$, which is false.

\emph{$Y_0$ is $Q$ or $W$.} We have $Q\subset Y_0\subset W$. As $W$ is a minimal prime of $Q+(s)$ in the noetherian ring
$R$, $W/Q$ has height at most one in $R/Q$ \cite[Tag 00KV]{Sta}, so $Y_0=Q$ or $Y_0=W$. If $Y_0=W$, then $Y$ is a
relevant multihomogeneous prime containing $P'$ with $v_*^{-1}(Y)=W$, so $Y\in F_W$.

\emph{If $Y_0=Q$, then $Y=P'$.} Suppose $P'\subsetneq Y$, and let $f\in Y\setminus P'$. As $P'$ is multihomogeneous, a
multihomogeneous component $g$ of $f$, of some multidegree $d$, is not in $P'$, and $g\in Y$, as $Y$ is multihomogeneous.
As $Y$ is relevant, there are indices $c_{k'}$ with $x'_{k',c_{k'}}\notin Y$, hence $x'_{k',c_{k'}}\notin P'$.
Lemma~\ref{cv-origin-intersection-projection-generic-integral-section-count-relation}, with $n=d-1$, gives $b_0,\dots,b_r\in R$ with $b_0\notin Q$ and $\sum_iv_*(b_i)g^iT_c^{r-i}\in P'\subset Y$. The
terms with $i\ge1$ lie in $Y$, as $g$ does, so $v_*(b_0)T_c^r\in Y$; and $T_c\notin Y$, as $Y$ is prime. So
$v_*(b_0)\in Y$ and $b_0\in Y_0=Q$, which is false.
\end{proof}
% Lean: CentreOverW.proof: imageZ_le (the component c = weightedHomogeneousComponent (pullWeight v) D a, rename_component
% and LocRingPow.component_mem, exists_form_notMem for u, the membership hR of u/c and the centre property h2, with
% mul_lt_one_of_lt_of_le), eq_or_eq_of_mem_minimalPrimes (Ideal.map_height_le_one_of_mem_minimalPrimes in R/Q, Mathlib's
% form of the principal ideal theorem) and eq_of_imageZ_eq (LocRingPow.exists_component_not_mem, exists_elim, and the sum
% split by Finset.sum_range_succ'); imageZ_isRelevant. The hypotheses IsMultiHomog s alpha and s not in imageZ v P' are
% not used.

\begin{lemma}\label{cv-origin-intersection-projection-generic-cycles-pushforward-integral}
In Situation~\ref{cv-origin-intersection-projection-generic-cycles-pushforward-situation}, $B$ is integral over $A$.
\end{lemma}
\leanref{GroundTheory/HeightsAndTheta/Heights/Intersections/StatementsSemilocIntegral.lean}{GroundTheory.Heights₂.Heights.Inter.SemilocIntegral.Statement}
\leanref{GroundTheory/HeightsAndTheta/Heights/Intersections/PSemilocIntegral.lean}{GroundTheory.Heights₂.Heights.Inter.SemilocIntegral.proof}
\begin{proof}
Let $L$ be the fraction field of $R'/P'$, which contains $B$ and, by the injection of the situation, $A$. Let
$\mathcal V\subset L$ be a valuation ring of $L$ containing $A$, and let $Y$ be a centre of $\mathcal V$
(Lemma~\ref{cv-origin-intersection-projection-generic-cycles-pushforward-integral-centre}, for $R'$ and $P'$). By
Lemma~\ref{cv-origin-intersection-projection-generic-cycles-pushforward-integral-centre-fibre}, $Y=P'$ or $Y\in F_W$. The
generators $u/w$ of $B$ have $w\notin P'$ and $w$ in no prime of $F_W$, so $w\notin Y$, and $u/w\in\mathcal V$ by the
property of the centre. So $B\subset\mathcal V$. Hence $B$ lies in the intersection of the valuation rings of $L$ that
contain $A$, which is the integral closure of $A$ in $L$ \cite[Corollary~5.22]{AM69}.
\end{proof}
% Lean: SemilocIntegral.proof: iInf_valuationSubring_superset (Mathlib's form of the cited corollary), ValCentreZ (Lemma
% -integral-centre) and CentreOverW (Lemma -integral-centre-fibre) for each valuation subring V containing the range of
% locRing -> fraction field, hle (Subring.closure_le over the generators of semiLocRing), and the monic polynomial lifted
% through f.rangeRestrict (Polynomial.lifts_and_natDegree_eq_and_monic).

\begin{lemma}\label{cv-origin-intersection-projection-generic-cycles-pushforward-closure}
In Situation~\ref{cv-origin-intersection-projection-generic-cycles-pushforward-situation}, let $L$ be the fraction field of
$R'/P'$, and let $B_0\supset B$ be the subring of $L$ generated by the quotients $u/w$ of multihomogeneous $u,w\in R'$ of
the same multidegree with $w\notin P'$. Then there is a finite subset $T$ of $L$ whose $A$-span contains every element
of $B_0$ that is integral over $A$. The proof uses neither that $(\kappa,v)$ is a block map, nor the dimension of $Q$, nor
that $s$ is multihomogeneous and not in $Q$, nor that the elements lie in $B_0$: the $A$-span of $T$ contains every
element of $L$ integral over $A$.
\end{lemma}
\leanref{GroundTheory/HeightsAndTheta/Heights/Intersections/StatementsLocRingIntClosure.lean}{GroundTheory.Heights₂.Heights.Inter.LocRingIntClosureFinite.Statement}
\leanref{GroundTheory/HeightsAndTheta/Heights/Intersections/PLocRingIntClosure.lean}{GroundTheory.Heights₂.Heights.Inter.LocRingIntClosureFinite.proof}
\begin{proof}
\emph{A subring of finite type.} As $W$ is relevant, choose $c_k$ with $x_{k,c_k}\notin W$ for $k\in K$, and let
$D\subset A$ be the subring generated by the finitely many quotients $x_{k,i}/x_{k,c_k}$. By the proof of
Lemma~\ref{cv-origin-intersection-family-curves-vertical-local}(3), every $z\in A$ is $b/c$ with $b,c\in D$, $c\neq0$
and $c^{-1}\in A$. So $A$ is the localization of $D$ at the set $U$ of the elements of $D$ that are units of $A$. The ring
$D$ is a domain of finite type over $\mathbf Z$, and it maps injectively into $L$ through $A$; let $K_D\subset L$ be the
subfield generated by its image.

\emph{The integral closure of $D$.} The field $L$ is generated by the finitely many classes of the $x'_{k',i'}$, so it is
a finitely generated extension of $K_D$, and the algebraic closure $M$ of $K_D$ in $L$ is finite over $K_D$
\cite[Lemma 9.26.11, Tag 037J]{Sta}. As $\mathbf Z$ is universally Japanese \cite[Tag 0335]{Sta}, the integral closure
$D'$ of $D$ in $M$ is a finite $D$-module; let $T$ be a finite set of generators.

\emph{Elements integral over $A$.} Let $x\in L$ be integral over $A$, with $x^r+a_1x^{r-1}+\dots+a_r=0$ and $a_l\in A$.
Write $a_l=d_l/m_l$ with $d_l\in D$ and $m_l\in U$, and put $m=\prod_lm_l\in U$. Multiplying by $m^r$,
$(mx)^r+\sum_l(m^la_l)(mx)^{r-l}=0$ with $m^la_l\in D$; so $mx$ is integral over $D$. It is algebraic over $K_D$, so it
lies in $M$, and in $D'$, the $D$-span of $T$. As $m$ is a unit of $A$, $x=m^{-1}(mx)$ lies in the $A$-span of $T$.
\end{proof}
\leanref{GroundTheory/HeightsAndTheta/Heights/Intersections/PLocRingChartLoc.lean}{GroundTheory.Heights₂.Heights.Inter.LocRingChartLoc.proof}
\leanref{GroundTheory/HeightsAndTheta/Heights/Intersections/PIntClosureLocFinite.lean}{GroundTheory.Heights₂.Heights.Inter.IntClosureLocFinite.proof}
% Lean: LocRingIntClosureFinite.proof (the hypotheses _hv, _hQ, _hs and _hsQ are not used, and x in semiLocRing P' empty
% is dropped) from LocRingChartLoc.proof (the subring D of the ratios g p, hkey and LocRingFrac, as in Lemma
% cv-origin-intersection-family-curves-vertical-local (3)) and IntClosureLocFinite.proof with
% subfield_closure_coords_eq_top: RelAlgClosureFinite.proof (the Lean's own proof of Stacks Tag 037J, through a
% transcendence basis), Model.FiniteIntegralClosure.full (the Lean's own proof that domains of finite type over Z are
% Japanese, from FiniteIntegralClosureLocal, -CharZero and -CharP, about 800 lines), the localization U of D, and
% IsIntegral.exists_multiple_integral_of_isLocalization for m.

\begin{lemma}\label{cv-origin-intersection-projection-generic-cycles-pushforward-finite}
In Situation~\ref{cv-origin-intersection-projection-generic-cycles-pushforward-situation}, $B$ is a finite torsion-free
$A$-module.
\end{lemma}
\leanref{GroundTheory/HeightsAndTheta/Heights/Intersections/StatementsSemilocFinite.lean}{GroundTheory.Heights₂.Heights.Inter.SemilocFinite.Statement}
\leanref{GroundTheory/HeightsAndTheta/Heights/Intersections/PSemilocFinite.lean}{GroundTheory.Heights₂.Heights.Inter.SemilocFinite.proof}
\begin{proof}
The ring $B$ lies in a field into which $A$ maps injectively, so $B$ is torsion-free over $A$. By
Lemma~\ref{cv-origin-intersection-projection-generic-cycles-pushforward-integral}, every element of $B$ is integral over
$A$, so by Lemma~\ref{cv-origin-intersection-projection-generic-cycles-pushforward-closure}, $B$ lies in the $A$-span of a
finite set. This span is a finite module over $A$, which is noetherian by
Lemma~\ref{cv-origin-intersection-projection-generic-cycles-pushforward-local}; so its submodule $B$ is finite.
\end{proof}
% Lean: SemilocFinite.proof: hint (SemilocIntegral), LocRingIntClosureFinite for t, isNoetherian_span_of_finite and
% Module.Finite.of_injective; the torsion-freeness through the injectivity of fracMapZ. Size: PSemilocFinite.lean,
% 38 lines.

\begin{lemma}\label{cv-origin-intersection-projection-generic-cycles-pushforward-rank}
In Situation~\ref{cv-origin-intersection-projection-generic-cycles-pushforward-situation}, the rank
$\dim_{F_A}(B\otimes_AF_A)$ of $B$ over $A$, for the fraction field $F_A$ of $A$, is $\deg(V)$.
\end{lemma}
\leanref{GroundTheory/HeightsAndTheta/Heights/Intersections/StatementsSemilocRank.lean}{GroundTheory.Heights₂.Heights.Inter.SemilocRank.Statement}
\leanref{GroundTheory/HeightsAndTheta/Heights/Intersections/PSemilocRank.lean}{GroundTheory.Heights₂.Heights.Inter.SemilocRank.proof}
\leanref{GroundTheory/HeightsAndTheta/Heights/Intersections/PSemilocRankFrac.lean}{GroundTheory.Heights₂.Heights.Inter.SemilocRankFrac.proof}
\begin{proof}
Let $E$ be the fraction field of $R'/P'$, and $\xi_{k',i}\in E$ the image of $x'_{k',i}$, so that $\xi_{v(k,i)}$ is the
image of $x_{k,i}$ under the embedding of the fraction field of $R/Q$. The field $E$ has characteristic zero, as
$P'\cap\mathbf Z=0$, and the kernels of $\xi$ and of $\xi\circ v$ are $P'$ and $Q$. With the fields of ratios of
Definition~\ref{cv-origin-intersection-full}(2), Lemma~\ref{cv-origin-intersection-projection-count}, for the point
$\xi$, gives $\deg(V)=[\mathbf Q(\xi):\mathbf Q(\xi\circ v)]$, as $\deg(V)$ is the dimension of the coordinate ring of $V$
in a chart (as in the proof of Lemma~\ref{cv-origin-intersection-projection}). It remains to show that the rank of $B$
over $A$ is $[\mathbf Q(\xi):\mathbf Q(\xi\circ v)]$.

As $P'$ is relevant, choose $i_{k'}$ with $x'_{k',i_{k'}}\notin P'$ for $k'\in K'$, and put
$\rho_{k',i}=\xi_{k',i}/\xi_{k',i_{k'}}$. For $u$ multihomogeneous of multidegree $a$,
$u(\xi)=u(\rho)\prod_{k'}\xi_{k',i_{k'}}^{a_{k'}}$; so a quotient $u/w$ of the same multidegree is $u(\rho)/w(\rho)$, and
$B$ lies in $L=\mathbf Q(\xi)$, which is generated by the $\rho_{k',i}$. Likewise the image of $A$ lies in
$L_0=\mathbf Q(\xi\circ v)$: as $v$ maps $I_k$ onto $I'_{\kappa(k)}$, $\xi_{\kappa(k),i_{\kappa(k)}}=\xi_{v(k,i_0)}$ for
some $i_0$.

The field $L_0$ is generated by the image of $A$: as $W$ is relevant, for $k\in K$ there is $j$ with $x_{k,j}\notin W$,
and $\xi_{v(k,i)}/\xi_{v(k,i_0)}$ is the quotient of the images of $x_{k,i}/x_{k,j}$ and $x_{k,i_0}/x_{k,j}$, which lie
in $A$. The field $L$ is generated by $B$. Indeed, $F_W$ is finite by
Lemma~\ref{cv-origin-intersection-projection-generic-cycles-pushforward-fibres}(2), so
Lemma~\ref{cv-origin-intersection-family-sections-avoid}, for the primes of $F_W$ and $P'$, gives $d_0$ and, for a given
$k'$, a multihomogeneous $w$ of multidegree $e_{k'}+d_0$ that lies in none of them, where $e_{k'}$ is the unit vector of
$k'$. With $N=\prod_{l}x'^{d_{0,l}}_{l,i_l}\notin P'$, the elements $b=x'_{k',i}N/w$ and $b'=x'_{k',i_{k'}}N/w$ lie in
$B$, $b'\neq0$, and $\rho_{k',i}=b/b'$.

So the embeddings identify $L_0$ with the fraction field of $A$ and $L$ with that of $B$. As $B$ is a finite torsion-free
$A$-module (Lemma~\ref{cv-origin-intersection-projection-generic-cycles-pushforward-finite}) and a domain,
$B\otimes_AF_A$ is a domain finite over $F_A$, hence a field, the fraction field of $B$; so the rank of $B$ is
$[L:L_0]$.
\end{proof}
% Lean: SemilocRank.proof: hJ.1 (the value of the field intersection family on the zero-dimensional fibre in the chart
% i'), SemilocRankAdjoin.proof and ChartRankAdjoin.proof (Lemma cv-origin-intersection-projection-count-chart), with
% SemilocRankAdjoin.chain repeating the argument of Lemma cv-origin-intersection-projection-count (ChartScalingIndep,
% AdjoinIndepRelfinrank, FinrankAdjoinRange); the text cites that lemma for the point xi of the fraction field instead.
% SemilocRankFrac.proof: fracEval, aeval_scale_div (the quotients in rho), hL0 (exists j with X (k, j) not in W),
% SemilocChartRatio.proof (AvoidRing.exists_form_avoid with FibreOverFinite) for hL, and RankFracDegree.proof
% (isFractionRing_of_closure, IsFractionRing.finrank_eq) with SemilocFinite. Size: PSemilocRank.lean,
% PSemilocRankAdjoin.lean, PSemilocRankFrac.lean, PRankFracDegree.lean and PSemilocChartRatio.lean, about 380 lines.

\begin{lemma}\label{cv-origin-intersection-projection-generic-cycles-pushforward-max}
In Situation~\ref{cv-origin-intersection-projection-generic-cycles-pushforward-situation}, for $W'\in F_W$ let
$\mathfrak m_{W'}$ be the inverse image in $B$ of the maximal ideal of $\mathcal O_{P',W'}\supset B$
(Definition~\ref{cv-origin-intersection-family-curves-vertical-ring}). Then the $\mathfrak m_{W'}$ are maximal ideals of
$B$, distinct for distinct $W'$, every maximal ideal of $B$ is one of them, and $B_{\mathfrak m_{W'}}=\mathcal O_{P',W'}$.
\end{lemma}
\leanref{GroundTheory/HeightsAndTheta/Heights/Intersections/StatementsSemilocMaxIdeals.lean}{GroundTheory.Heights₂.Heights.Inter.SemilocMaxIdeals.Statement}
\leanref{GroundTheory/HeightsAndTheta/Heights/SemipositiveHeights/SemilocResidueDegOfFibreIncomp.lean}{GroundTheory.Heights₂.Heights.ForSecondHalf.semilocMaxIdeals_of_fibreIncomp}
\begin{proof}
Let $L$ be the fraction field of $R'/P'$, and write $\bar f$ for the class of $f\in R'$. Every $W'\in F_W$ is a relevant
multihomogeneous prime containing $P'$, the set $F_W$ is finite
(Lemma~\ref{cv-origin-intersection-projection-generic-cycles-pushforward-fibres}(2)), and $B\subset\mathcal O_{P',W'}$, as
the generators $u/w$ of $B$ have $w\notin W'$. By
Lemma~\ref{cv-origin-intersection-family-curves-vertical-local}(2), $\mathcal O_{P',W'}$ is local, so $\mathfrak m_{W'}$
is a prime ideal of $B$, and $1\notin\mathfrak m_{W'}$.

\emph{Quotients.} (a) Every $b\in B$ is $\bar u/\bar w$ with $u,w$ multihomogeneous of the same multidegree, $w\notin P'$
and $w$ in no $W'\in F_W$: such quotients include the generators, $0$ and $1$, and are closed under negation, sums and
products, as in the proof of Lemma~\ref{cv-origin-intersection-family-curves-vertical-local}(1), as $P'$ and the
$W'\in F_W$ are prime. (b) For such $b=\bar u/\bar w$ and $W'\in F_W$, $b\in\mathfrak m_{W'}$ if and only if $u\in W'$, by
Lemma~\ref{cv-origin-intersection-family-curves-vertical-local}(2). (c) If moreover $u\notin P'$ and $u$ lies in no
$W'\in F_W$, then $\bar w/\bar u\in B$, and $b$ is a unit of $B$.

\emph{Scaling.} (d) For $a\in\mathbf N^{K'}$ and a relevant prime $W''$ of $R'$ there are multihomogeneous $t$ and $h$,
of multidegrees $d$ and $a+d$ for some $d$, with $t\notin W''$ and $h$ in no $W'\in F_W$. Indeed,
Lemma~\ref{cv-origin-intersection-family-sections-avoid}, for the finite set $F_W$, gives $d$ and a multihomogeneous $h$
of multidegree $a+d\ge d$ in no $W'\in F_W$; and $t=\prod_{k'}x'^{d_{k'}}_{k',i_{k'}}$, with $x'_{k',i_{k'}}\notin W''$,
is not in the prime $W''$.

\emph{Comparison.} (e) If $W_1,W_2\in F_W$ and $\mathfrak m_{W_1}\subset\mathfrak m_{W_2}$, then $W_1\subset W_2$. As
$W_1$ is generated by multihomogeneous elements, let $f\in W_1$ be multihomogeneous of multidegree $a$, and take $t$ and
$h$ as in (d) for $a$ and $W''=W_2$. Then $h\notin P'$, as $P'\subset W_1$, so $b=\overline{ft}/\bar h\in B$. By (b),
$b\in\mathfrak m_{W_1}$, as $ft\in W_1$; so $b\in\mathfrak m_{W_2}$, and $ft\in W_2$ by (b). As $t\notin W_2$,
$f\in W_2$. In particular the $\mathfrak m_{W'}$ are distinct for distinct $W'$.

\emph{Localization.} (f) Let $W'\in F_W$. The elements of $B\setminus\mathfrak m_{W'}$ are units of
$\mathcal O_{P',W'}$, and $B\to\mathcal O_{P',W'}$ is injective, as both are subrings of $L$. Every
$z\in\mathcal O_{P',W'}$ is $\bar u/\bar w$ with $u,w$ multihomogeneous of some multidegree $a$ and $w\notin W'$
(Lemma~\ref{cv-origin-intersection-family-curves-vertical-local}(1)). Take $t$ and $h$ as in (d) for $a$ and $W''=W'$;
then $h\notin P'$, $b=\overline{ut}/\bar h$ and $c=\overline{wt}/\bar h$ lie in $B$, $c\notin\mathfrak m_{W'}$ by (b), as
$wt\notin W'$, and $z=b/c$. So $\mathcal O_{P',W'}$ is the localization $B_{\mathfrak m_{W'}}$.

\emph{Maximal ideals.} (g) Let $M$ be a maximal ideal of $B$. Suppose first that $F_W$ is empty. Then every nonzero
$b=\bar u/\bar w\in B$ has $u\notin P'$, so it is a unit by (c), and $B$ is a field. The ring $A$ maps injectively into
$B$, and $B$ is integral over $A$ by
Lemma~\ref{cv-origin-intersection-projection-generic-cycles-pushforward-integral}. So $A$ is a field: if $a\in A$ is
nonzero, $a^{-1}\in B$ satisfies $a^{-n}+c_1a^{-n+1}+\dots+c_n=0$ with $c_l\in A$, and multiplying by $a^{n-1}$ gives
$a^{-1}=-(c_1+c_2a+\dots+c_na^{n-1})\in A$. As $W$ is relevant, there is a multihomogeneous $t\notin W$ of the
multidegree of $s$, and by Lemma~\ref{cv-origin-intersection-projection-generic-cycles-pushforward-local} for $Q$, $s$, $W$
and $t$, $A$ has dimension $1$, which is false. So $F_W$ is not empty. Now $M$ lies in the union of the
$\mathfrak m_{W'}$, $W'\in F_W$: if $b=\bar u/\bar w\in M$, as in (a), lies in none of them, then $u$ lies in no
$W'\in F_W$ by (b), and $u\notin P'$, as $P'$ lies in some $W'\in F_W$; so $b$ is a unit by (c), which is false. By prime
avoidance \cite[Lemma 10.15.2, Tag 00DS]{Sta}, $M\subset\mathfrak m_{W'}$ for some $W'\in F_W$, and $M=\mathfrak m_{W'}$,
as $M$ is maximal and $\mathfrak m_{W'}$ is proper.

(h) Let $W'\in F_W$. The proper ideal $\mathfrak m_{W'}$ lies in a maximal ideal, which is $\mathfrak m_{W_2}$ for some
$W_2\in F_W$ by (g). By (e), $W'\subset W_2$, and by
Lemma~\ref{cv-origin-intersection-projection-generic-cycles-pushforward-fibres}(1), $W'=W_2$. So $\mathfrak m_{W'}$ is
maximal.
\end{proof}
\leanref{GroundTheory/HeightsAndTheta/Heights/Intersections/PSemilocMaxIdeals.lean}{GroundTheory.Heights₂.Heights.Inter.SemilocMaxIdeals.proof}
% Lean: semilocMaxIdeals_of_fibreIncomp = SemilocMaxIdeals.proof (FibreOverFinite.proof, FibreIncomp, SemilocIntegral.full,
% LocRingDimOne.full), with semiToLoc and semiMaxIdeal: (a) exists_frac (Subring.closure_induction), (b)
% mem_semiMaxIdeal_iff (mem_maximalIdeal_iff), (c) isUnit_of, (d) exists_scale (AvoidRing.exists_form_avoid and
% exists_form_notMem), (e) le_of_semiMaxIdeal_le, (f) isLocalization (isLocalization_iff), (g) hcov with
% exists_eq_of_isMaximal (Ideal.subset_union_prime, Mathlib's prime avoidance) and, for F_W empty, hfield,
% isField_of_isIntegral_of_isField and ringKrullDim_eq_zero_of_isField against LocRingDimOne, (h) Ideal.exists_le_maximal
% and FibreIncomp.

\begin{lemma}\label{cv-origin-intersection-projection-generic-cycles-pushforward-residue-fibre}
In Situation~\ref{cv-origin-intersection-projection-generic-cycles-pushforward-situation}, for every $W'\in F_W$ the
generic fibre $V_{W'}$ of $W'$ over $W$ (Definition~\ref{cv-origin-intersection-projection-generic-cycles-fibres})
defines a zero-dimensional variety over $\Phi_{W'}$.
\end{lemma}
\leanref{GroundTheory/HeightsAndTheta/Heights/Intersections/StSemilocResidueFibreDim.lean}{GroundTheory.Heights₂.Heights.Inter.SemilocResidueFibreDim.Statement}
\leanref{GroundTheory/HeightsAndTheta/Heights/Intersections/PSemilocResidueFibreDim.lean}{GroundTheory.Heights₂.Heights.Inter.SemilocResidueFibreDim.proof}
\begin{proof}
\openstep{The proof is not yet written out. As $s\in W$, $v_*(s)\in W'$, and $W'$ and $W$ are relevant minimal primes
of $P'+(v_*(s))$ and of $Q+(s)$ (Lemma~\ref{cv-origin-intersection-projection-generic-cycles-pushforward-fibres}); their
dimensions agree, and Lemma~\ref{cv-origin-intersection-projection-fibre}, for $W'$ in place of $P'$, if $W'$ is
horizontal, or Lemma~\ref{cv-origin-intersection-projection-generic-cycles-vertical-degree} if it is vertical, gives the
claim.}
% Lean: SemilocResidueFibreDim.proof: FibreDivSupp (hD', W' in divSupp P' (rename v s)), divSuppArith for W' and for W
% (the dimension of R/W), then FibreDimStatement (Lemma cv-origin-intersection-projection-fibre) with r = 0 in the
% horizontal case, and FibreDimensionZero after reduction modulo p = charOf W' (red, ringKrullDim_red, imageVar) in the
% vertical case. Size (tools/lean_size.py SemilocResidueFibreDim.proof, stopping at FibreDivSupp, FibreDimStatement and
% FibreDimensionZero, which the book cites): 237 lines.
\end{proof}

\begin{lemma}\label{cv-origin-intersection-projection-generic-cycles-pushforward-residue}
In Situation~\ref{cv-origin-intersection-projection-generic-cycles-pushforward-situation}, with $\mathfrak m_{W'}$ as in
Lemma~\ref{cv-origin-intersection-projection-generic-cycles-pushforward-max}, for every $W'\in F_W$ the $A$-module
$B/\mathfrak m_{W'}$ has finite length $\ell_A(B/\mathfrak m_{W'})=\deg(V_{W'})$, for the generic fibre $V_{W'}$ of $W'$
over $W$ (Definition~\ref{cv-origin-intersection-projection-generic-cycles-fibres}).
\end{lemma}
\leanref{GroundTheory/HeightsAndTheta/Heights/Intersections/StatementsSemilocResidueDeg.lean}{GroundTheory.Heights₂.Heights.Inter.SemilocResidueDeg.Statement}
\leanref{GroundTheory/HeightsAndTheta/Heights/SemipositiveHeights/SemilocResidueDegOfFibreIncomp.lean}{GroundTheory.Heights₂.Heights.ForSecondHalf.semilocResidueDeg_of_fibreIncomp}
\begin{proof}
By Lemma~\ref{cv-origin-intersection-projection-generic-cycles-pushforward-residue-fibre}, $V_{W'}$ defines a
zero-dimensional variety over $\Phi_{W'}$, and $\deg(V_{W'})$ is its degree, the dimension over $\Phi_{W'}$ of its
coordinate ring in a chart.
\openstep{The rest of the proof is not yet written out: the residue field of $B$ at $\mathfrak m_{W'}$ is finite over
that of $A$, and its degree, which is $\ell_A(B/\mathfrak m_{W'})$, is computed by the ratio fields as in
Lemma~\ref{cv-origin-intersection-projection-generic-cycles-pushforward-rank}, with
Lemma~\ref{cv-origin-intersection-projection-count-scaling}.}
% Lean: semilocResidueDeg_of_fibreIncomp = SemilocResidueDeg.proof from SemilocResidueFibreDim.proof (hfib, Lemma
% cv-origin-intersection-projection-generic-cycles-pushforward-residue-fibre) and SemilocResidueAdjoin.proof
% (SemilocResidueFrac.proof with SemilocFinite and SemilocMaxIdeals, and ChartScalingIndep.full), with hJ.1 for the value
% J' (fibreIdeal kappa G') 0 in the chart i' and ChartRankAdjoin.proof. Size (tools/lean_size.py
% semilocResidueDeg_of_fibreIncomp): 2,539 lines, of which 2,302 are not below SemilocResidueFibreDim.proof.
\end{proof}

\begin{lemma}\label{cv-origin-intersection-projection-generic-cycles-pushforward-length}
In Situation~\ref{cv-origin-intersection-projection-generic-cycles-pushforward-situation}, let $t\in R\setminus W$ be
multihomogeneous of multidegree $\alpha$, and put $y=s/t\in A$. Then $B/yB$ has finite length over $A$, and
\[
 \ell_A(B/yB)=\sum_{W'}m_{W'}\deg(V_{W'}),
\]
the sum over the relevant minimal primes $W'$ of $P'+(v_*(s))$ with $v_*^{-1}(W')=W$, where $m_{W'}$ is the multiplicity of
$W'$ in $R'/(P'+(v_*(s)))$.
\end{lemma}
\leanref{GroundTheory/HeightsAndTheta/Heights/Intersections/StatementsSemilocLength.lean}{GroundTheory.Heights₂.Heights.Inter.SemilocLength.Statement}
\leanref{GroundTheory/HeightsAndTheta/Heights/Intersections/PSemilocLength.lean}{GroundTheory.Heights₂.Heights.Inter.SemilocLength.proof}
\leanref{GroundTheory/HeightsAndTheta/Heights/Intersections/PSemilocLengthSum.lean}{GroundTheory.Heights₂.Heights.Inter.SemilocLengthSum.proof}
\begin{proof}
By Lemma~\ref{cv-origin-intersection-projection-generic-cycles-pushforward-fibres}(2), $F_W$ is finite, and its elements
are the $W'$ of the sum. For $W'\in F_W$ let $\mathcal O_{W'}=\mathcal O_{P',W'}$. The image of $y$ in $\mathcal O_{W'}$
is $v_*(s)/v_*(t)$, where $v_*(s)$ and $v_*(t)$ are multihomogeneous of the same multidegree, $v_*(s)\notin P'$ and
$v_*(t)\notin W'$, as $t\notin W=v_*^{-1}(W')$. So by
Lemma~\ref{cv-origin-intersection-projection-generic-cycles-pushforward-local}, for $P'$, $v_*(s)$, $W'$ and $v_*(t)$,
$\mathcal O_{W'}$ is a noetherian local domain of dimension $1$, $y\neq0$ there, and $m_{W'}=\ell(\mathcal O_{W'}/y)$.

By Lemma~\ref{cv-origin-intersection-projection-generic-cycles-pushforward-max}, the maximal ideals of $B$ are the
$\mathfrak m_{W'}$, finitely many, with $B_{\mathfrak m_{W'}}=\mathcal O_{W'}$; so the localizations of $B/yB$ at its
maximal ideals are the $\mathcal O_{W'}/y$, of length $m_{W'}$. By
Lemma~\ref{cv-origin-intersection-projection-generic-cycles-pushforward-finite}, $B/yB$ is a finite $A$-module, and
$(B/yB)\otimes_AF_A=0$ for the fraction field $F_A$ of $A$, as $y\neq0$; so it has finite length over $A$, and a fortiori
a finite length $\ell$ over $B$.
The ring $B/yB$ has finitely many maximal ideals, and its Jacobson radical is nilpotent: it kills every composition
factor, so its $\ell$-th power is $0$. So $B/yB$ is the product of its localizations at the $\mathfrak m_{W'}$
\cite[Lemma 10.53.5, Tag 00JA]{Sta}. A composition series of
$\mathcal O_{W'}/y$ over $B$ has $m_{W'}$ factors, each isomorphic to $B/\mathfrak m_{W'}$, of length $\deg(V_{W'})$ over
$A$ by Lemma~\ref{cv-origin-intersection-projection-generic-cycles-pushforward-residue}. Adding gives the statement; it is
finite as $F_W$ is.
\end{proof}
% Lean: SemilocLength.proof from SemilocLengthSum.proof (FieldNumbers.FiniteAlgebraLengthSum.proof for the decomposition
% of the length over the maximal ideals, with SemilocFinite, SemilocMaxIdeals, SemilocResidueDeg, FibreDivSupp,
% LocRingDimOne and FibreOverFinite; hterm, with the unit case for W' not containing rename v s, which does not occur),
% FibreDivSupp and regularityDehomogenization_holds for the W'. Size: PSemilocLength.lean and PSemilocLengthSum.lean,
% about 120 lines, besides FiniteAlgebraLengthSum (not measured) and the lemmas cited.

\begin{lemma}\label{cv-origin-intersection-projection-generic-cycles-pushforward}
Let $d\ge1$, and let $(\kappa,v)$, $X'$, $P'$, $Q$ and $V$ be as in Lemma~\ref{cv-origin-intersection-projection-fibre}
with $r=0$, so that $X'$ and $Q$ have dimension $d$. Let $s\in R\setminus Q$ be multihomogeneous, and let $W$ be a
relevant minimal prime of $Q+(s)$. Then
\[
 \sum_{W'}m_{W'}\deg(V_{W'})=\deg(V)\,m_W,
\]
the sum over the relevant minimal primes $W'$ of $P'+(v_*(s))$ with $Q_{W'}=W$, where $m_{W'}$ is the multiplicity of
$W'$ in $R'/(P'+(v_*(s)))$, $m_W$ that of $W$ in $R/(Q+(s))$, and $Q_{W'}$ and $V_{W'}$ are as in
Definition~\ref{cv-origin-intersection-projection-generic-cycles-fibres}.
\end{lemma}
\leanref{GroundTheory/HeightsAndTheta/Heights/Intersections/IdDeg.lean}{GroundTheory.Heights₂.Heights.Inter.ArithmeticIntersectionDivisorPushforward.Statement}
\leanref{GroundTheory/HeightsAndTheta/Heights/Intersections/ImportsBis.lean}{GroundTheory.Heights₂.Heights.Inter.ArithmeticIntersectionDivisorPushforward.proof}
\begin{proof}
Let $\alpha$ be the multidegree of $s$, and let $F_W$, $A=\mathcal O_{Q,W}$ and $B$ be as in
Situation~\ref{cv-origin-intersection-projection-generic-cycles-pushforward-situation}. As $W$ is relevant, there are
$i_k$ with $x_{k,i_k}\notin W$, and $t=\prod_kx_{k,i_k}^{\alpha_k}$ is multihomogeneous of multidegree $\alpha$ and not in
$W$. Put $y=s/t$. By Lemma~\ref{cv-origin-intersection-projection-generic-cycles-pushforward-local}, $A$ is a noetherian
local domain of dimension $1$, $y\in A$ is nonzero, and $m_W=\ell_A(A/yA)$. By
Lemmas~\ref{cv-origin-intersection-projection-generic-cycles-pushforward-finite}
and~\ref{cv-origin-intersection-projection-generic-cycles-pushforward-rank}, $B$ is a finite torsion-free $A$-module of
rank $\deg(V)$. So by Lemma~\ref{cv-origin-intersection-projection-generic-cycles-pushforward-order},
$\ell_A(B/yB)=\deg(V)\,\ell_A(A/yA)=\deg(V)\,m_W$. By
Lemma~\ref{cv-origin-intersection-projection-generic-cycles-pushforward-length}, $\ell_A(B/yB)$ is the left side of the
statement.
% Lean: ArithmeticIntersectionDivisorPushforward.of_remaining = ArithmeticIntersectionDivisorPushforward.proof with
% exists_form_notMem for t, regularityDehomogenization_holds and LocRingDimOne (Lemma
% cv-origin-intersection-projection-generic-cycles-pushforward-local), SemilocFinite (-finite), SemilocRank.of_remaining
% (-rank), SemilocLength.of_remaining (-length) and dimOneTorsionFreeFormula (-order). The statement allows any E with
% E W' = J' (fibreIdeal kappa G') 0, the text's deg(V_W'). Size: ImportsBis.lean, about 45 lines, besides these lemmas
% (about 5,100 lines in all, part of them shared with Lemmas cv-origin-intersection-family-curves-vertical-local and
% cv-origin-intersection-family-sections-cycles-local; the open step is in -residue).
\end{proof}

\begin{lemma}\label{cv-origin-intersection-projection-generic-cycles}
Let $(\cdot\mid\cdot)$ be an arithmetic intersection system, let $d\ge1$, and let $(\kappa,v)$, $X'$, $P'$, $Q$, $\Phi$
and $V$ be as in Lemma~\ref{cv-origin-intersection-projection-fibre} with $r=0$, so that $X'$ and $Q$ have dimension $d$.
Suppose that the conclusion of Lemma~\ref{cv-origin-intersection-projection-generic} holds in dimension $d-1$, for every
block map. Let $\bar L_1,\dots,\bar L_{d-1}$ be smooth metrized line bundles on $\prod_k\mathbf P^{I_k}_{\mathbf Z}$ with
pullbacks $\bar L'_1,\dots,\bar L'_{d-1}$, and let $s\in R\setminus Q$ be multihomogeneous. Then, with $\mathcal C$ as in
Lemma~\ref{cv-origin-intersection-projection-generic-rational},
\[
 \mathcal C(\bar L'_1\cdots\bar L'_{d-1};v_*(s)\mid X')=\deg(V)\,\mathcal C(\bar L_1\cdots\bar L_{d-1};s\mid Q).
\]
\end{lemma}
\leanref{GroundTheory/HeightsAndTheta/Heights/Intersections/IsFieldIntersectionNoBlockFam.lean}{GroundTheory.Heights₂.Heights.Inter.ArithmeticIntersectionCycle.Statement}
\leanref{GroundTheory/HeightsAndTheta/Heights/Intersections/Assembly.lean}{GroundTheory.Heights₂.Heights.Inter.ArithmeticIntersectionCycle.of_remaining}
\begin{proof}
Write $\bar L$ for $\bar L_1\cdots\bar L_{d-1}$ and $\bar L'$ for $\bar L'_1\cdots\bar L'_{d-1}$, and use $Q_{W'}$ and
$V_{W'}$ of Definition~\ref{cv-origin-intersection-projection-generic-cycles-fibres}. For a relevant multihomogeneous prime
$U$ of $R$ let $T(U)$ be $(\bar L\mid U)$ if $U$ is horizontal and $\deg_{L_1\cdots L_{d-1}}(U)\log p$ if
$U\cap\mathbf Z=p\mathbf Z$, and define $T'(U')$ for primes $U'$ of $R'$ in the same way with $\bar L'$. Then
$\mathcal C(\bar L';v_*(s)\mid X')=\sum_{W'}m_{W'}T'(W')$, the sum over the relevant minimal primes $W'$ of $P'+(v_*(s))$
with their multiplicities, and $\mathcal C(\bar L;s\mid Q)=\sum_Wm_WT(W)$, the sum over the relevant minimal primes $W$ of
$Q+(s)$.

We show that for each $W'$, $T'(W')=\deg(V_{W'})\,T(Q_{W'})$ if $Q_{W'}$ is a relevant minimal prime of $Q+(s)$, and
$T'(W')=0$ otherwise. If $W'$ is horizontal, this is
Lemma~\ref{cv-origin-intersection-projection-generic-cycles-horizontal}, by which $Q_{W'}$ is horizontal too. If
$W'\cap\mathbf Z=p\mathbf Z$, then $Q_{W'}\cap\mathbf Z=p\mathbf Z$, as $v_*$ is the identity on $\mathbf Z$; the
multidegrees of the $\bar L'_j$ are those of the $\bar L_j$ extended by $0$; and
Lemma~\ref{cv-origin-intersection-projection-generic-cycles-vertical}, multiplied by $\log p$, gives the claim.

Grouping the $W'$ according to $Q_{W'}$ and applying
Lemma~\ref{cv-origin-intersection-projection-generic-cycles-pushforward} to each relevant minimal prime $W$ of $Q+(s)$,
\[
 \mathcal C(\bar L';v_*(s)\mid X')=\sum_W\Big(\sum_{Q_{W'}=W}m_{W'}\deg(V_{W'})\Big)T(W)=\sum_W\deg(V)\,m_W\,T(W)
 =\deg(V)\,\mathcal C(\bar L;s\mid Q).
\]
\end{proof}
\leanref{GroundTheory/HeightsAndTheta/Heights/Intersections/IsHorizontalImageZIff.lean}{GroundTheory.Heights₂.Heights.Inter.ArithmeticIntersectionCycle.proof}
% Lean: ArithmeticIntersectionCycle.of_remaining = ArithmeticIntersectionCycle.proof with
% ArithmeticIntersectionHorizontal.full, ArithmeticIntersectionVertical.of_remaining (with ForSecondHalf.projDeg_all for
% ProjDegStatement), ArithmeticIntersectionDivisorPushforward.of_remaining (FibreIncomp.full,
% semilocMaxIdeals_of_fibreIncomp, semilocResidueDeg_of_fibreIncomp, Giv.fibreDim), GenericPointExists.full and
% Targets.ExistsUnique.full for E W' = J' (fibreIdeal kappa G') 0; hterm (isHorizontal_imageZ_iff, charOf_imageZ), key,
% Finset.sum_fiberwise_eq_sum_filter. The hypothesis is GenDegAt I n. Size: IsHorizontalImageZIff.lean, about 140 lines,
% besides Lemmas cv-origin-intersection-projection-generic-cycles-horizontal, -vertical and -pushforward.

\subsection{Integral sections and the projection formula}\label{subsection-cv-projection-formula}

\begin{definition}\label{cv-origin-intersection-projection-generic-integral-section-map}
Let $(\kappa,v)$, $X'$, $P'$ and $Q$ be as in Lemma~\ref{cv-origin-intersection-projection-fibre} with $r=0$, with $X'$
and $Q$ of dimension $n+1$. Let $J'$ be an ideal of $R'$, let $(\tilde X',(U_x),(\lambda_x))$ be an analytic resolution of
$X'$ off $J'$ (Definition~\ref{cv-origin-intersection-resolution}), and let $\Psi$ be a good projection for $Q$, with the
Segre map $\sigma$ of Definition~\ref{cv-origin-intersection-poincare-lelong-all-double-projection}.
\begin{enumerate}
\item For $x\in\tilde X'$ let $\pi(x)=(\lambda_x(x)_{v(k,i)})_{k,i}$, a point of $\hat Q(\mathbf C)$, and put
  $u(x)=\Psi(\sigma(\pi(x)))\in\mathbf C^{n+1}$, with coordinates $u_0(x),\dots,u_n(x)$.
\item Let $q:\tilde X'\to\mathbf C^n$ be the map $q(x)=(u_1(x)/u_0(x),\dots,u_n(x)/u_0(x))$, with $q(x)=0$ if
  $u_0(x)=0$, and for $u_0(x)\neq0$ put $W(x)=u_0(x)^{-1}\sigma(\pi(x))$.
\item The \emph{excluded set} $E$ is the set of the $x\in\tilde X'$ such that $\lambda_x(x)\in Z(J')$, or $u_0(x)=0$, or for
  some chart $\kappa$ of $\tilde X'$ at $x$ the map $q\circ\kappa^{-1}$ is not holomorphic near $\kappa(x)$ or has
  Jacobian determinant $0$ at $\kappa(x)$.
\end{enumerate}
\end{definition}
\leanref{GroundTheory/HeightsAndTheta/Heights/Intersections/PullBad.lean}{GroundTheory.Heights₂.Heights.Inter.pullPt}
\leanref{GroundTheory/HeightsAndTheta/Heights/Intersections/PullBad.lean}{GroundTheory.Heights₂.Heights.Inter.pullQ}
\leanref{GroundTheory/HeightsAndTheta/Heights/Intersections/PullBad.lean}{GroundTheory.Heights₂.Heights.Inter.pullW}
\leanref{GroundTheory/HeightsAndTheta/Heights/Intersections/PullBad.lean}{GroundTheory.Heights₂.Heights.Inter.pullBad}
% The Lean's pullQ divides by Psi_0 without a case distinction (division by 0 is 0 in Lean), which is the text's
% convention q(x) = 0 when u_0(x) = 0. pi(x) lies in the cone of Q, as Q = v_*^{-1}(P').

\begin{lemma}\label{cv-origin-intersection-projection-generic-integral-section-null-orbits}
Let $(\kappa,v)$, $X'$, $P'$ and $Q$ be as in Lemma~\ref{cv-origin-intersection-projection-fibre} with $r=0$, with $X'$
and $Q$ of dimension $n+1$, and write $z'\circ v=(z'_{v(k,i)})_{k,i}$. There is a multihomogeneous $f_0\in R'\setminus P'$
such that for every $u\in\mathbf C^{\sqcup_kI_k}$ there is a finite set $T$ of points of $\mathbf C^{\sqcup_{k'}I'_{k'}}$
with the following property: every $z'\in\hat X'(\mathbf C)$ with $f_0(z')\neq0$ and $z'\circ v$ in the orbit of $u$
under $(\mathbf C^\times)^K$ lies in the orbit of some $t\in T$ under $(\mathbf C^\times)^{K'}$.
\end{lemma}
\leanref{GroundTheory/HeightsAndTheta/Heights/Intersections/StatementsBlockGenFinite.lean}{GroundTheory.Heights₂.Heights.Inter.BlockGenFinite.Statement}
\leanref{GroundTheory/HeightsAndTheta/Heights/Intersections/PBlockGenFinite.lean}{GroundTheory.Heights₂.Heights.Inter.BlockGenFinite.proof}
\begin{proof}
\openstep{The proof is not yet written out: each ratio of coordinates of $X'$ in a block satisfies an equation over $R$
whose leading coefficient is not in $Q$ (Lemma~\ref{cv-origin-lift-finite}), $f_0$ is a product of these leading
coefficients, renamed by $v_*$, and of coordinates, and over a point $u$ where it does not vanish the ratios take finitely
many values.}
% Lean: BlockGenFinite.proof from BlockRatioAlg (PBlockGenFinite.lean). Size (tools/lean_size.py
% HilbertSamuel.Closed.Inter.BlockGenFinite): 249 lines.
\end{proof}

\begin{lemma}\label{cv-origin-intersection-projection-generic-integral-section-null-isolated}
Let $(\kappa,v)$ be a block map (Definition~\ref{cv-origin-intersection-full}(1)), let $X'$ be an arithmetic
subvariety of dimension $n+1$ of $\prod_{k'}\mathbf P^{I'_{k'}}_{\mathbf Z}$ defined by $P'$, and write
$z'\circ v=(z'_{v(k,i)})_{k,i}$. For every $u\in\mathbf C^{\sqcup_kI_k}$ there is a finite set $T$ of points of
$\mathbf C^{\sqcup_{k'}I'_{k'}}$ with the following property. Let $z'\in\hat X'(\mathbf C)$ with $z'\circ v$ in the
orbit of $u$, and suppose that there is $\varepsilon>0$ such that every $z''\in\hat X'(\mathbf C)$ with $z''\circ v$ in the
orbit of $u$ and $|z''-z'|<\varepsilon$ lies in the orbit of $z'$. Then $z'$ lies in the orbit of some $t\in T$.
\end{lemma}
\leanref{GroundTheory/HeightsAndTheta/Heights/Intersections/StatementsFibreIsolFinite.lean}{GroundTheory.Heights₂.Heights.Inter.FibreIsolFinite.Statement}
\leanref{GroundTheory/HeightsAndTheta/Heights/Intersections/PFibreIsolFinite.lean}{GroundTheory.Heights₂.Heights.Inter.FibreIsolFinite.proof}
\begin{proof}
\openstep{The proof is not yet written out: after normalizing the blocks of $\kappa(K)$ to those of $u$ and one
coordinate in each other block, the points over $u$ form a complex algebraic set in an affine space, and its isolated
points are finitely many.}
% Lean: FibreIsolFinite.proof (PFibreIsolFinite.lean: isolated_finite for the isolated points of the zero set of a set of
% polynomials in N variables, sameOrbit_comp_smul, minor_of_sameOrbit and sameOrbit_of_minor). Size (tools/lean_size.py
% HilbertSamuel.Closed.Inter.FibreIsolFinite): 445 lines.
\end{proof}

\begin{lemma}\label{cv-origin-intersection-projection-generic-integral-section-null-zero}
Let $X'$ be an arithmetic subvariety of dimension $n+1$ defined by the prime $P'$ of $R'$, let
$(\tilde X',(U_x),(\lambda_x))$ be an analytic resolution of $X'$ off an ideal $J'$
(Definition~\ref{cv-origin-intersection-resolution}), and let $f\in R'\setminus P'$ be multihomogeneous. Then for every
chart $\kappa\colon V\to\mathbf C^n$ of $\tilde X'$ the set of the $a\in\kappa(V)$ with $f(\lambda_x(x))=0$, for
$x=\kappa^{-1}(a)$, has measure zero.
\end{lemma}
\leanref{GroundTheory/HeightsAndTheta/Heights/Intersections/HasStrictFDerivAtEquiv.lean}{GroundTheory.Heights₂.Heights.Inter.pullbackPushforward_holds.fzero_null}
\begin{proof}
By the first paragraph of the proof of
Lemma~\ref{cv-origin-intersection-poincare-lelong-all-double-bridge-null}, it suffices that every $a_0\in\kappa(V)$ has a
ball $D_0\subset\kappa(V\cap U_{x_0})$ around it, $x_0=\kappa^{-1}(a_0)$, on which the set is the zero set of a
holomorphic function that does not vanish identically. Let $g(a)=f(\lambda_{x_0}(\kappa^{-1}(a)))$ on $D_0$. As
$\lambda_x(x)$ lies in the orbit of $\lambda_{x_0}(x)$ for $x\in U_{x_0}$ and $f$ is multihomogeneous, the set meets $D_0$
in the zero set of $g$. Suppose that $g$ vanishes on $D_0$. Let $\Psi'$ be a good projection for $X'$
(Lemma~\ref{cv-origin-intersection-poincare-lelong-all-double-good}(1)), with $v'$, $Q'$ and $B'$ as in
Definition~\ref{cv-origin-intersection-poincare-lelong-all-double-bridge-map} for $X'$. By
Lemma~\ref{cv-origin-intersection-poincare-lelong-all-double-bridge-null}(1), $\kappa(B'\cap V)$ has measure zero, so some
$a_1\in D_0$ has $x_1=\kappa^{-1}(a_1)\notin B'$. Then $Q'\circ\kappa^{-1}$ is holomorphic near $a_1$ with nonzero
Jacobian determinant, and by the holomorphic inverse function theorem it maps every neighbourhood of $a_1$ onto a
neighbourhood of $Q'(x_1)$. Take one, $O\subset D_0$, on which $v'_0\circ\kappa^{-1}$ does not vanish. For $a\in O$ and
$x=\kappa^{-1}(a)$, Lemma~\ref{cv-origin-intersection-poincare-lelong-all-double-bridge-fibres}(2) gives
$z\in\hat X'(\mathbf C)$, $\lambda_x(x)$ with one block rescaled, with $\Psi'(\sigma(z))=(1,Q'(x))$, and $f(z)=0$, as
$g(a)=0$. So a nonempty open set of values $y$ has a $z\in\hat X'(\mathbf C)$ with $f(z)=0$ and
$\Psi'(\sigma(z))=(1,y)$. The proof of Lemma~\ref{cv-origin-intersection-poincare-lelong-all-double-bridge-null}(2), for the
ideal $(f)$, which is generated by a multihomogeneous element and not contained in $P'$, shows that this set of $y$ has
measure zero, which is false.
\end{proof}
% Lean: fzero_null: goodProjExists for Psi', projB_null for span {f} (hnull), volume_eq_zero_of_locally_in_analytic_zero_set
% with g, psi0_null and ram_null (Lemma cv-origin-intersection-poincare-lelong-all-double-bridge-null (1)) for u1,
% analyticAt_brQ, hasStrictFDerivAt_equiv and himg (the image of O is a neighbourhood). Does not use J' beyond the
% resolution.

\begin{lemma}\label{cv-origin-intersection-projection-generic-integral-section-null-flat}
In Definition~\ref{cv-origin-intersection-projection-generic-integral-section-map}, let $\kappa\colon V\to\mathbf C^n$ be a
chart of $\tilde X'$, and let $O\subset\kappa(V)$ be a nonempty open set such that $u_0$ does not vanish on
$\kappa^{-1}(O)$ and $q\circ\kappa^{-1}$ is holomorphic on $O$. Then the Jacobian determinant of $q\circ\kappa^{-1}$ does
not vanish identically on $O$.
\end{lemma}
\leanref{GroundTheory/HeightsAndTheta/Heights/Intersections/ProjectionPullbackZeroSetNull.lean}{GroundTheory.Heights₂.Heights.Inter.pullbackPushforward_holds.not_flat_pull}
\begin{proof}
Let $f_0$ be given by Lemma~\ref{cv-origin-intersection-projection-generic-integral-section-null-orbits}, and let $O'$ be
the set of the $a\in O$ with $\lambda_x(x)\notin Z(J')$ and $f_0(\lambda_x(x))\neq0$ for $x=\kappa^{-1}(a)$. It is open,
as in the proof of Lemma~\ref{cv-origin-intersection-poincare-lelong-all-double-bridge-null}, as both conditions are
invariant under the orbits and open for $\lambda_{x_0}$ near $x_0$. It is nonempty, as $O$ has positive measure and
$O\setminus O'$ has measure zero, by (6) of Definition~\ref{cv-origin-intersection-resolution} and
Lemma~\ref{cv-origin-intersection-projection-generic-integral-section-null-zero}. Suppose that the Jacobian determinant of
$q\circ\kappa^{-1}$ vanishes on $O$. The constant rank argument of the proof of
Lemma~\ref{cv-origin-intersection-poincare-lelong-all-double-bridge-null} (\emph{No flat pieces}), applied on $O'$, gives
an infinite subset of $O'$ on which $q\circ\kappa^{-1}$ takes one value $y$; its image under $\kappa^{-1}$ is an infinite
set $S$ of points $x$ with $q(x)=y$. Let $x\in S$. By the proof of
Lemma~\ref{cv-origin-intersection-projection-generic-integral-section}, $W(x)\in F_\Psi(y)$, which is finite by
Lemma~\ref{cv-origin-intersection-poincare-lelong-all-double-good}(2); write each $w\in F_\Psi(y)$ as $\sigma(u_w)$ with
$u_w\in\hat Q(\mathbf C)$. As $\sigma(\pi(x))=u_0(x)\sigma(u_{W(x)})$, the argument of the proof of
Lemma~\ref{cv-origin-intersection-poincare-lelong-all-double-bridge-fibres}(3) puts $\pi(x)=\lambda_x(x)\circ v$ in the
orbit of $u_{W(x)}$. As $f_0(\lambda_x(x))\neq0$, $\lambda_x(x)$ lies in the orbit of a point of the finite set
$T(u_{W(x)})$ of Lemma~\ref{cv-origin-intersection-projection-generic-integral-section-null-orbits}. By (3) of
Definition~\ref{cv-origin-intersection-resolution}, as $\lambda_x(x)\notin Z(J')$, each orbit contains $\lambda_x(x)$ for at
most one $x$. So $S$ is finite, which is false.
\end{proof}
% Lean: not_flat_pull: BlockGenFinite for f0, isOpen_chart_notB and isOpen_chart_ne for O', R.null and fzero_null for its
% nonemptiness, exists_infinite_fibre (InterBridgeK, Lemma cv-origin-intersection-poincare-lelong-all-double-bridge-null),
% fibreFinite, pullW_mem_fibre, sameOrbit_of_segre and orbit_subsingleton (R.inj).

\begin{lemma}\label{cv-origin-intersection-projection-generic-integral-section-null}
In Definition~\ref{cv-origin-intersection-projection-generic-integral-section-map}, $\kappa(E\cap V)$ has measure zero
for every chart $\kappa:V\to\mathbf C^n$ of $\tilde X'$, and $q^{-1}(y)\setminus E$ is finite for every
$y\in\mathbf C^n$.
\end{lemma}
\leanref{GroundTheory/HeightsAndTheta/Heights/Intersections/PullBad.lean}{GroundTheory.Heights₂.Heights.Inter.pullbackPushforward_holds.Statement}
\leanref{GroundTheory/HeightsAndTheta/Heights/Intersections/Assembly.lean}{GroundTheory.Heights₂.Heights.Inter.pullbackPushforward_holds.full}
\begin{proof}
For $x_0\in\tilde X'$ and $x\in U_{x_0}$, $\lambda_x(x)$ lies in the orbit of $\lambda_{x_0}(x)$ under
$(\mathbf C^\times)^{K'}$, so $\pi(x)$
lies in the orbit of $\lambda_{x_0}(x)\circ v$ under $(\mathbf C^\times)^K$ and $u(x)$ is a nonzero multiple of
$\Psi(\sigma(\lambda_{x_0}(x)\circ v))$; by the first paragraph of the proof of
Lemma~\ref{cv-origin-intersection-projection-generic-integral-section-count-negligible}(2), $q\circ\kappa^{-1}$ is
holomorphic near $\kappa(x)$ when $u_0(x)\neq0$.

\emph{The excluded set.} As in the proof of Lemma~\ref{cv-origin-intersection-poincare-lelong-all-double-bridge-null}(1),
$\kappa(E\cap V)$ lies in the union of three sets: the $a$ with $\lambda_x(x)\in Z(J')$, for $x=\kappa^{-1}(a)$, of
measure zero by (6) of Definition~\ref{cv-origin-intersection-resolution}; the $a$ with $u_0(x)\neq0$ at which the
Jacobian determinant of $q\circ\kappa^{-1}$ vanishes, by the chain rule and the holomorphy above; and the $a$ with
$u_0(x)=0$. The second set has measure zero by the argument of that proof, with
Lemma~\ref{cv-origin-intersection-projection-generic-integral-section-null-flat} in place of its paragraph \emph{No flat
pieces}: near each of its points $q\circ\kappa^{-1}$ is holomorphic and its Jacobian determinant, a holomorphic function,
does not vanish identically. For the third, let $a_0$ be one of its points, $x_0=\kappa^{-1}(a_0)$ and $D_0$ a ball around
$a_0$ in $\kappa(V\cap U_{x_0})$. On $D_0$ the set is the zero set of the holomorphic function
$\phi(a)=\Psi(\sigma(\lambda_{x_0}(\kappa^{-1}(a))\circ v))_0$. Suppose that $\phi$ vanishes on $D_0$. As $\Psi$ is a good
projection for $Q$ and $\pi(x_0)\in\hat Q(\mathbf C)$, some coordinate $i_0$ of $\Psi(\sigma(\pi(x_0)))$ is nonzero, and
$i_0\neq0$. Let $\Psi'$ be $\Psi$ followed by the exchange of the coordinates $0$ and $i_0$, again a good projection for
$Q$, with $u'$ and $q'$ defined for it as $u$ and $q$ are for $\Psi$. On the open set $O$ of the $a\in D_0$ with
$\Psi'(\sigma(\lambda_{x_0}(\kappa^{-1}(a))\circ v))_0\neq0$, which contains $a_0$, $u'_0$ does not vanish on
$\kappa^{-1}(O)$, $q'\circ\kappa^{-1}$ is holomorphic, and its coordinate $i_0$ is $\phi$ divided by a nonzero function,
hence $0$. So the Jacobian determinant of $q'\circ\kappa^{-1}$ vanishes on $O$, which
Lemma~\ref{cv-origin-intersection-projection-generic-integral-section-null-flat}, for $\Psi'$, excludes.

\emph{Isolation.} Let $x\notin E$. We show that there is $\varepsilon>0$ such that every $z''\in\hat X'(\mathbf C)$ with
$z''\circ v$ in the orbit of $\pi(x)$ and $|z''-\lambda_x(x)|<\varepsilon$ lies in the orbit of $\lambda_x(x)$. Let $\kappa$
be a chart at $x$. As $x\notin E$, $q\circ\kappa^{-1}$ is a local biholomorphism at $\kappa(x)$, so $q$ is injective on an
open neighbourhood $N$ of $x$. For $a,b\in\mathbf C^{\prod_kI_k}$ with $a\neq0$ let
$\theta_b(a)=|(a_ib_l-a_lb_i)_{i,l}|/|a|$; it is continuous in $a$, unchanged when $a$ is multiplied by a nonzero scalar,
and $0$ exactly when $a$ is a multiple of $b$. Put $b=\sigma(\lambda_x(x))$, which is nonzero. The function
$x'\mapsto\theta_b(\sigma(\lambda_{x'}(x')))$ on $\tilde X'$ is well defined and continuous, by (1) of
Definition~\ref{cv-origin-intersection-resolution}, and it is positive off $N$: if it vanishes at $x'$, then
$\sigma(\lambda_{x'}(x'))$ is a multiple of $b$, so $\lambda_{x'}(x')$ lies in the orbit of $\lambda_x(x)$, as in the
proof of Lemma~\ref{cv-origin-intersection-poincare-lelong-all-double-bridge-fibres}(3), and $x'=x$ by (3) of
Definition~\ref{cv-origin-intersection-resolution}, as $\lambda_x(x)\notin Z(J')$. As $\tilde X'\setminus N$ is compact, the
function is at least some $m_0>0$ there. Some $g\in J'$ has $g(\lambda_x(x))\neq0$. Choose $\varepsilon>0$ such that
$|z''-\lambda_x(x)|<\varepsilon$ implies $\theta_b(\sigma(z''))<m_0$ and $g(z'')\neq0$. Let $z''$ be as above. Then
$z''\notin Z(J')$, so by (2) of Definition~\ref{cv-origin-intersection-resolution} it lies in the orbit of
$\lambda_{x''}(x'')$ for some $x''$; then $\theta_b(\sigma(\lambda_{x''}(x'')))=\theta_b(\sigma(z''))<m_0$, so $x''\in N$.
As $\pi(x'')$ lies in the orbit of $z''\circ v$, hence of $\pi(x)$, $u(x'')$ is a nonzero multiple of $u(x)$ and
$q(x'')=q(x)$. So $x''=x$, and $z''$ lies in the orbit of $\lambda_x(x)$.

\emph{Finite fibres.} Let $y\in\mathbf C^n$. The fibre $F_\Psi(y)$ is finite
(Lemma~\ref{cv-origin-intersection-poincare-lelong-all-double-good}(2)); write each $w\in F_\Psi(y)$ as $\sigma(u_w)$ with
$u_w\in\hat Q(\mathbf C)$, and let $T(u_w)$ be the finite set of
Lemma~\ref{cv-origin-intersection-projection-generic-integral-section-null-isolated}. Let $x\in q^{-1}(y)\setminus E$.
Then $u_0(x)\neq0$, $\lambda_x(x)\notin Z(J')$, and $W(x)\in F_\Psi(y)$, as in the proof of
Lemma~\ref{cv-origin-intersection-projection-generic-integral-section}. As in the proof of
Lemma~\ref{cv-origin-intersection-projection-generic-integral-section-null-flat}, $\pi(x)$ lies in the orbit of
$u_{W(x)}$. By the isolation, Lemma~\ref{cv-origin-intersection-projection-generic-integral-section-null-isolated} applies
to $z'=\lambda_x(x)$ and $u=u_{W(x)}$, and $\lambda_x(x)$ lies in the orbit of a point of $T(u_{W(x)})$. By (3) of
Definition~\ref{cv-origin-intersection-resolution}, each orbit contains $\lambda_x(x)$ for at most one such $x$. So
$q^{-1}(y)\setminus E$ is finite.
\end{proof}
\leanref{GroundTheory/HeightsAndTheta/Heights/Intersections/ImportsExtra.lean}{GroundTheory.Heights₂.Heights.Inter.pullbackPushforward_holds.proof}
\leanref{GroundTheory/HeightsAndTheta/Heights/Intersections/ProjectionPullbackZeroSetNull.lean}{GroundTheory.Heights₂.Heights.Inter.pullbackPushforward_holds.pullBad_null}
\leanref{GroundTheory/HeightsAndTheta/Heights/Intersections/ContinuousMinorRatioLift.lean}{GroundTheory.Heights₂.Heights.Inter.pullbackPushforward_holds.isolated_of_not_pullBad}
\leanref{GroundTheory/HeightsAndTheta/Heights/Intersections/ContinuousMinorRatioLift.lean}{GroundTheory.Heights₂.Heights.Inter.pullbackPushforward_holds.pullFib_finite}
% Lean: pullbackPushforward_holds.full = HilbertSamuel.Closed.Inter.pullbackPushforward_holds = pullbackPushforward_holds.proof
% with BlockGenFinite and FibreIsolFinite (Lemmas -null-orbits and -null-isolated, the open steps, 249 and 445 lines).
% pullBad_null: R.null, projection_pullback_zeroSet_null (the exchange Psi' by Equiv.swap, not_flat_pull for Psi') and
% ram_pull_null (det_pullQ_chart, analyticAt_det_fderiv, not_flat_pull), with
% volume_eq_zero_of_locally_in_analytic_zero_set. pullFib_finite: fibreFinite, pullW_mem_fibre, sameOrbit_of_segre,
% isolated_of_not_pullBad (hasStrictFDerivAt_equiv, minorRatio and continuous_minorRatio_lift, the compact complement of
% V, R.surj and R.inj) and orbit_subsingleton.

\begin{lemma}\label{cv-origin-intersection-projection-generic-integral-section-count-orbits}
Let $(\kappa,v)$, $X'$, $P'$ and $Q$ be as in Lemma~\ref{cv-origin-intersection-projection-fibre} with $r=0$, with $X'$
and $Q$ of dimension $n+1$, let $\Phi\subset\Phi'$ and $R'\to\Phi'$ be as there, with $\xi_{k',i'}\in\Phi'$ the image of
$x'_{k',i'}$, and let $(i_{k'})_{k'\notin\kappa(K)}$ be indices with $\xi_{k',i_{k'}}\neq0$. Let $S$ be the set of the
$\xi_{k',i'}/\xi_{k',i_{k'}}$ with $k'\notin\kappa(K)$, and let $\Psi$ be a good projection for $Q$, with the Segre map
$\sigma$ (Definition~\ref{cv-origin-intersection-poincare-lelong-all-double-projection}). For $z'\in\hat X'(\mathbf C)$
put $z'\circ v=(z'_{v(k,i)})_{k,i}$ and $u(z')=\Psi(\sigma(z'\circ v))$, and for $w\in\mathbf C^{\prod_kI_k}$ let $O(w)$
be the set of the orbits of $(\mathbf C^\times)^{K'}$ of the points $z'\in\hat X'(\mathbf C)$ with $u(z')_0\neq0$ and
$u(z')_0^{-1}\sigma(z'\circ v)=w$. Then for almost every $y\in\mathbf C^n$ and every $w\in F_\Psi(y)$, $O(w)$ has
$[\Phi(S):\Phi]$ elements.
\end{lemma}
\leanref{GroundTheory/HeightsAndTheta/Heights/Intersections/PtsOver.lean}{GroundTheory.Heights₂.Heights.Inter.GenericPointsCount.Statement}
\leanref{GroundTheory/HeightsAndTheta/Heights/Intersections/PtsOver.lean}{GroundTheory.Heights₂.Heights.Inter.ptsOver}
\leanref{GroundTheory/HeightsAndTheta/Heights/Intersections/ImportsMore.lean}{GroundTheory.Heights₂.Heights.Inter.GenericPointsCount.proof}
\begin{proof}
The block map is an instance of Situation~\ref{cv-origin-lift-situation}, with $K$, $(I_k)$, $K'$ and $(I'_{k'})$ in
the roles of $K_0$, $(I^0_k)$, $K$ and $(I_k)$, the injection $\kappa$, and the bijections $I_k\to I'_{\kappa(k)}$
induced by $v$. Then $f'=v_*(f)$, $P'_0=Q$, $z'|_0=z'\circ v$, and the fibre blocks are the $k'\notin\kappa(K)$. The
field $\Phi$ with the images of the $x_{k,i}$, and $\Phi'$ with the $\xi_{k',i'}$, form a generic fibre of $P'$
(Definition~\ref{cv-origin-lift-fibre}): $\Phi$ is the fraction field of $R/Q$, the kernel of $R'\to\Phi'$ is $P'$, and
the image of $x_{k,i}$ is $\xi_{v(k,i)}$. The indices $(i_{k'})$ form a chart of it, with chart field $\Phi(S)$. By
Lemma~\ref{cv-origin-lift-orbits} there is a multihomogeneous $D\in R\setminus Q$ such that for every
$z\in\hat Q(\mathbf C)$ with $D(z)\neq0$ the points of $\hat X'(\mathbf C)$ over $z$ form $[\Phi(S):\Phi]$ orbits of
$(\mathbf C^\times)^{K'}$.

The ideal $(D)$ is generated by a multihomogeneous element and is not contained in $Q$. The proof of
Lemma~\ref{cv-origin-intersection-poincare-lelong-all-double-bridge-null}(2) uses only these two properties of the ideal
$J$; for $Q$ and the ideal $(D)$ it shows that the set $N$ of the $y\in\mathbf C^n$ with $\Psi(\sigma(z))=(1,y)$ for some
$z\in\hat Q(\mathbf C)$ with $D(z)=0$ has measure zero. Let $y\notin N$ and $w\in F_\Psi(y)$, so that $w=\sigma(z)$ with
$z\in\hat Q(\mathbf C)$ and $\Psi(\sigma(z))=(1,y)$. Then $D(z)\neq0$. We show that $O(w)$ is the set of the orbits of the
points of $\hat X'(\mathbf C)$ over $z$, which has $[\Phi(S):\Phi]$ elements.

Let $z'\in\hat X'(\mathbf C)$ with $u(z')_0\neq0$ and $u(z')_0^{-1}\sigma(z'\circ v)=\sigma(z)$. Then
$\sigma(z'\circ v)=u(z')_0\,\sigma(z)$ with $z\in\hat C$, and the argument of the proof of
Lemma~\ref{cv-origin-intersection-poincare-lelong-all-double-bridge-fibres}(3) gives $c\in(\mathbf C^\times)^K$ with
$z'\circ v=(c_kz_{k,i})_{k,i}$; so $z'$ lies over $z$. Conversely, let $z'\in\hat X'(\mathbf C)$ lie over $z$, with
$z'\circ v=(c_kz_{k,i})_{k,i}$ and $c\in(\mathbf C^\times)^K$. Then $\sigma(z'\circ v)=\nu\sigma(z)$ with
$\nu=\prod_kc_k\neq0$, so $u(z')=\nu(1,y)$, $u(z')_0=\nu\neq0$, and $u(z')_0^{-1}\sigma(z'\circ v)=\sigma(z)=w$.
\end{proof}
% Lean: GenericPointsCount.proof: GenericPointsAvoidBadLocus (Lemma cv-origin-lift-orbits) for D, projB_null for the
% ideal span {D} (Lemma cv-origin-intersection-poincare-lelong-all-double-bridge-null (2)), sameOrbit_of_segre for the
% first inclusion and segre_smul for the second. The Lean's ptsOver P' v Psi w has no condition u_0 != 0; for w in the
% fibre the two sets agree, as Psi(w)_0 = 1.

\begin{lemma}\label{cv-origin-intersection-projection-generic-integral-section-count-relation}
Let $(\kappa,v)$, $X'$, $P'$ and $Q$ be as in Lemma~\ref{cv-origin-intersection-projection-fibre} with $r=0$, with $X'$
and $Q$ of dimension $n+1$. Let $g\in R'\setminus P'$ be multihomogeneous of multidegree $d\in\mathbf N^{K'}$, let
$c=(c_{k'})_{k'\in K'}$ be indices with $x'_{k',c_{k'}}\notin P'$ for every $k'$, and put
$T_c=\prod_{k'}x'^{d_{k'}}_{k',c_{k'}}$. Then there are $r\ge0$ and $b_0,\dots,b_r\in R$ with $b_0\notin Q$ and
\[
 \sum_{i=0}^rv_*(b_i)\,g^i\,T_c^{r-i}\in P'.
\]
\end{lemma}
\leanref{GroundTheory/HeightsAndTheta/Heights/Intersections/IsMultiHomogOfPullWeight.lean}{GroundTheory.Heights₂.Heights.Inter.PullbackElimination.exists_elim}
\begin{proof}
Let $F$ be the fraction field of $R'/P'$ and $\xi_{k',i'}\in F$ the class of $x'_{k',i'}$. By
Lemma~\ref{cv-origin-intersection-projection-fibre}, $v_*$ induces an embedding of the fraction field $\Phi$ of $R/Q$
into $F$, under which the class of $x_{k,i}$ is $\xi_{v(k,i)}$; as in the proof of
Lemma~\ref{cv-origin-intersection-projection-generic-integral-section-count-orbits}, $\Phi$ and $F$ with these elements
form a generic fibre of $P'$. The elements $\xi_{k',c_{k'}}$ are nonzero; put $\rho_{k',i}=\xi_{k',i}/\xi_{k',c_{k'}}$ and
$\gamma=g(\rho)$. As $g$ is multihomogeneous of multidegree $d$, $g(\xi)=t\gamma$ with
$t=\prod_{k'}\xi_{k',c_{k'}}^{d_{k'}}\neq0$, and $g(\xi)\neq0$ as $g\notin P'$; so $\gamma\neq0$. Each $\rho_{k',i}$ is
algebraic over $\Phi$. For $k'=\kappa(k)$ it lies in $\Phi$, as $v$ maps $I_k$ onto $I'_{\kappa(k)}$ and the
$\xi_{v(k,i)}$ lie in $\Phi$. For $k'\notin\kappa(K)$ it lies in the chart field of the chart
$(c_{k'})_{k'\notin\kappa(K)}$, which is finite over $\Phi$ by Lemma~\ref{cv-origin-lift-finite}. So $\gamma$ is
algebraic over $\Phi$. Its minimal polynomial has a nonzero constant term, as $\gamma\neq0$, and clearing denominators
gives $b_0,\dots,b_r\in R$ with $b_0\notin Q$ and $\sum_i\bar b_i\gamma^i=0$, where $\bar b_i$ is the class of $b_i$ in
$\Phi$. The image of $\sum_iv_*(b_i)g^iT_c^{r-i}$ in $F$ is $\sum_i\bar b_i(t\gamma)^it^{r-i}=t^r\sum_i\bar b_i\gamma^i=0$,
so it lies in $P'$.
\end{proof}
% Lean: exists_elim: the embedding phiA of R/Q into the fraction field of R'/P', aeval_eq_coordProd_mul for
% g(xi) = t gamma, the algebraicity hrho of each ratio, from BlockRatioAlg.proof for every block (the text uses Lemma
% cv-origin-lift-finite for the fibre blocks, whose proof is BlockRatioAlg, and the inclusion in Phi for the others), hgamma
% through Subalgebra.algebraicClosure, the polynomial with nonzero constant coefficient
% (exists_nonzero_coeff_and_aeval_eq_zero, lifted to R by Polynomial.map_surjective) and the sum hsum.

\begin{lemma}\label{cv-origin-intersection-projection-generic-integral-section-count-elimination}
Let $(\kappa,v)$, $X'$, $P'$ and $Q$ be as in Lemma~\ref{cv-origin-intersection-projection-fibre} with $r=0$, with $X'$
and $Q$ of dimension $n+1$, and let $J'$ be an ideal of $R'$ generated by multihomogeneous elements and not contained in
$P'$. Then there is a multihomogeneous $h\in R\setminus Q$ with $h\big((z'_{v(k,i)})_{k,i}\big)=0$ for every
$z'\in\hat X'(\mathbf C)$ at which every element of $J'$ vanishes.
\end{lemma}
\leanref{GroundTheory/HeightsAndTheta/Heights/Intersections/ImportsFurtherRest.lean}{GroundTheory.Heights₂.Heights.Inter.PullbackEliminationOfVariable.Statement}
\leanref{GroundTheory/HeightsAndTheta/Heights/Intersections/IsMultiHomogOfPullWeight.lean}{GroundTheory.Heights₂.Heights.Inter.PullbackEliminationOfVariable.proof}
\begin{proof}
Write $z'\circ v=(z'_{v(k,i)})_{k,i}$, so that $f(z'\circ v)=v_*(f)(z')$ for $f\in R$. As $J'$ is generated by
multihomogeneous elements and is not contained in $P'$, it has a multihomogeneous element $g\notin P'$, of some
multidegree $d\in\mathbf N^{K'}$. Let $C$ be the finite set of the families $c=(c_{k'})_{k'\in K'}$ of indices
$c_{k'}\in I'_{k'}$ with $x'_{k',c_{k'}}\notin P'$ for every $k'$.

\emph{One family.} Let $c\in C$. We find a multihomogeneous $h_c\in R\setminus Q$ with $h_c(z'\circ v)=0$ for every
$z'\in\hat X'(\mathbf C)$ with $g(z')=0$ and $z'_{k',c_{k'}}\neq0$ for every $k'$. By
Lemma~\ref{cv-origin-intersection-projection-generic-integral-section-count-relation} there are $b_0,\dots,b_r\in R$ with
$b_0\notin Q$ and $G=\sum_{i=0}^rv_*(b_i)g^iT_c^{r-i}\in P'$, where $T_c=\prod_{k'}x'^{d_{k'}}_{k',c_{k'}}$ is
multihomogeneous of multidegree $d$.

As $Q$ is multihomogeneous and $b_0\notin Q$, the component of $b_0$ of some multidegree $D\in\mathbf N^K$ is not in $Q$.
For $f\in R$ let $f_D$ be its component of multidegree $D$. The homomorphism $v_*$ maps a monomial of multidegree
$a\in\mathbf N^K$ to one of multidegree $a'$, where $a'$ is $a$ on $\kappa(K)$ and $0$ elsewhere, and $a\mapsto a'$ is
injective, as $\kappa$ is; so the component of $v_*(f)$ of multidegree $D'$ is $v_*(f_D)$. Each $g^iT_c^{r-i}$ is
multihomogeneous of multidegree $rd$. So the component of $G$ of multidegree $D'+rd$ is
$\sum_iv_*((b_i)_D)\,g^i\,T_c^{r-i}$, and it lies in $P'$, as $P'$ is multihomogeneous. Put $h_c=(b_0)_D$, a
multihomogeneous element of $R\setminus Q$. Let $z'\in\hat X'(\mathbf C)$ with $g(z')=0$ and $z'_{k',c_{k'}}\neq0$ for
every $k'$. Evaluating this component of $G$ at $z'$, the terms with $i\ge1$ vanish, so
$0=h_c(z'\circ v)\,T_c(z')^r$; as $T_c(z')\neq0$, $h_c(z'\circ v)=0$.

\emph{All families.} Put $h=\prod_{c\in C}h_c$. It is multihomogeneous, and it is not in the prime $Q$. Let
$z'\in\hat X'(\mathbf C)$ with every element of $J'$ vanishing at $z'$. As $z'\in\hat C'$, there are indices $c_{k'}$
with $z'_{k',c_{k'}}\neq0$ for every $k'$; then $x'_{k',c_{k'}}\notin P'$, as every element of $P'$ vanishes at $z'$, so
$c\in C$. As $g\in J'$, $g(z')=0$, so $h_c(z'\circ v)=0$ and $h(z'\circ v)=0$.
\end{proof}
\leanref{GroundTheory/HeightsAndTheta/Heights/Intersections/IsMultiHomogOfPullWeight.lean}{GroundTheory.Heights₂.Heights.Inter.PullbackElimination.exists_elim_homog}
% Lean: PullbackEliminationOfVariable.proof: key for each c with X (k', c k') not in P' (exists_elim_homog, eval_eq_zero,
% isMultiHomog_of_pullWeight), the product over the filter good of these c, and hgood for the chosen e.
% exists_elim_homog: exists_elim (Lemma -count-relation), the weighted homogeneous component for the weight pullWeight v
% (the text's D'), LocRingPow.component_mem, component_mul_right and rename_component.

\begin{lemma}\label{cv-origin-intersection-projection-generic-integral-section-count-negligible}
Let $(\kappa,v)$, $X'$, $P'$ and $Q$ be as in Lemma~\ref{cv-origin-intersection-projection-fibre} with $r=0$, with $X'$
and $Q$ of dimension $n+1$, let $\Psi$ be a good projection for $Q$, with the Segre map $\sigma$
(Definition~\ref{cv-origin-intersection-poincare-lelong-all-double-projection}), and let $J'$ be an ideal of $R'$
generated by multihomogeneous elements and not contained in $P'$. For $z'\in\hat X'(\mathbf C)$ put
$z'\circ v=(z'_{v(k,i)})_{k,i}$.
\begin{enumerate}
\item The set of the $y\in\mathbf C^n$ for which some $z'\in\hat X'(\mathbf C)$ at which every element of $J'$ vanishes
  has $u_0\neq0$ and $u_0^{-1}\sigma(z'\circ v)\in F_\Psi(y)$, where $u=\Psi(\sigma(z'\circ v))$, has measure zero.
\item Let $(\tilde X',(U_x),(\lambda_x))$ be an analytic resolution of $X'$ off $J'$
  (Definition~\ref{cv-origin-intersection-resolution}), with $u$, $q$ and $E$ as in
  Definition~\ref{cv-origin-intersection-projection-generic-integral-section-map}. Then the set of the $q(x)$ with $x\in E$
  and $u_0(x)\neq0$ has measure zero. This part uses only that $(\kappa,v)$ is a block map and the resolution: neither the
  dimensions of $X'$ and $Q$ nor that $\Psi$ is a good projection, which the definitions of $u$, $q$ and $E$ do not use.
\end{enumerate}
\end{lemma}
\leanref{GroundTheory/HeightsAndTheta/Heights/Intersections/ImportsFurtherTer.lean}{GroundTheory.Heights₂.Heights.Inter.PullbackVanishingOfNullCycle.Statement}
\leanref{GroundTheory/HeightsAndTheta/Heights/Intersections/ImportsFurtherMore.lean}{GroundTheory.Heights₂.Heights.Inter.pullbackBadImage_holds.Statement}
\begin{proof}
(1) Let $h$ be given by Lemma~\ref{cv-origin-intersection-projection-generic-integral-section-count-elimination}. The
ideal $(h)$ is generated by a multihomogeneous element and is not contained in $Q$, so, as in the proof of
Lemma~\ref{cv-origin-intersection-projection-generic-integral-section-count-orbits}, the proof of
Lemma~\ref{cv-origin-intersection-poincare-lelong-all-double-bridge-null}(2) shows that the set $N$ of the
$y\in\mathbf C^n$ with $\Psi(\sigma(z))=(1,y)$ for some $z\in\hat Q(\mathbf C)$ with $h(z)=0$ has measure zero. Let $y$
lie in the set of (1), with $z'$ and $u$ as there. As $\Psi(w)=(1,y)$ for $w\in F_\Psi(y)$,
$\Psi(u_0^{-1}\sigma(z'\circ v))=(1,y)$. The point $z'\circ v$ lies in $\hat Q(\mathbf C)$: it lies in $\hat C$, as $v$
maps $I_k$ onto $I'_{\kappa(k)}$, and $f(z'\circ v)=v_*(f)(z')=0$ for $f\in Q$. Fix $k_0\in K$, and let $z$ be
$z'\circ v$ with its block $k_0$ multiplied by $u_0^{-1}$. Then $z\in\hat Q(\mathbf C)$ and $h(z)=0$, as $Q$ and $h$ are
multihomogeneous and $h(z'\circ v)=0$, and $\sigma(z)=u_0^{-1}\sigma(z'\circ v)$. So $\Psi(\sigma(z))=(1,y)$, and
$y\in N$.

(2) \emph{Holomorphy.} Let $x\in\tilde X'$ with $u_0(x)\neq0$, and let $\kappa$ be a chart of $\tilde X'$ at $x$. For
$z\in U_x$, by (1) of Definition~\ref{cv-origin-intersection-resolution}, $\lambda_z(z)=(c_{k'}\lambda_x(z)_{k',i})$ for
some $c\in(\mathbf C^\times)^{K'}$, so $\pi(z)=(c_{\kappa(k)}\lambda_x(z)_{v(k,i)})_{k,i}$ and $u(z)$ is
$\prod_kc_{\kappa(k)}$ times $\Psi(\sigma(\lambda_x(z)\circ v))$. The latter is holomorphic in $z$, and its coordinate
$0$ does not vanish at $x$, hence near $x$. So, as in the proof of
Lemma~\ref{cv-origin-intersection-poincare-lelong-all-double-bridge-fibres}(1), $q\circ\kappa^{-1}$ is holomorphic near
$\kappa(x)$, the vector of the quotients of the coordinates $1,\dots,n$ of $\Psi(\sigma(\lambda_x(\kappa^{-1}(\cdot))\circ
v))$ by its coordinate $0$.

\emph{A cover.} Fix a finite atlas of charts $\kappa_j:V_j\to\mathbf C^n$ of $\tilde X'$. Let $S_{1,j}$ be the set of
the $a\in\kappa_j(V_j)$ with $\lambda_x(x)\in Z(J')$ and $u_0(x)\neq0$ for $x=\kappa_j^{-1}(a)$, and $S_{2,j}$ the set of
the $a\in\kappa_j(V_j)$ with $u_0(\kappa_j^{-1}(a))\neq0$ at which the Jacobian determinant of $q\circ\kappa_j^{-1}$ is
$0$. Let $x\in E$ with $u_0(x)\neq0$, and $j$ with $x\in V_j$. If $\lambda_x(x)\in Z(J')$, then $\kappa_j(x)\in S_{1,j}$.
Otherwise, by the definition of $E$, there is a chart $\kappa$ at $x$ such that $q\circ\kappa^{-1}$, which is holomorphic
near $\kappa(x)$ by the previous paragraph, has Jacobian determinant $0$ at $\kappa(x)$. By the chain rule, as the change
of charts $\kappa\circ\kappa_j^{-1}$ has an invertible derivative, the Jacobian determinant of $q\circ\kappa_j^{-1}$ at
$\kappa_j(x)$ is $0$, and $\kappa_j(x)\in S_{2,j}$. So the set of (2) lies in the union over $j$ of the images of
$S_{1,j}$ and $S_{2,j}$ under $q\circ\kappa_j^{-1}$.

\emph{The images.} By the first paragraph, $q\circ\kappa_j^{-1}$ is differentiable at every point of $S_{1,j}$ and of
$S_{2,j}$. The set $S_{1,j}$ has measure zero, by (6) of Definition~\ref{cv-origin-intersection-resolution}, so its image
has measure zero \cite[263D(ii)]{Fre01}. At the points of $S_{2,j}$ the real Jacobian determinant of
$q\circ\kappa_j^{-1}$, which is $|\det d(q\circ\kappa_j^{-1})|^2$, vanishes, so its image has measure zero by the same
reference, whose bound for the outer measure of the image is the integral of the real Jacobian determinant.
\end{proof}
\leanref{GroundTheory/HeightsAndTheta/Heights/Intersections/ImportsFurther.lean}{GroundTheory.Heights₂.Heights.Inter.PullbackVanishingOfNullCycle.proof}
\leanref{GroundTheory/HeightsAndTheta/Heights/Intersections/DetRestrictScalarsEqZero.lean}{GroundTheory.Heights₂.Heights.Inter.pullbackBadImage_holds.proof}
% Lean: PullbackVanishingOfNullCycle.proof: PullbackEliminationOfVariable for h, projB_null for span {h}, comp_mem_mcPts
% for z' o v, and mcPts_smul with the scaling c = update 1 k0 (u_0^-1) and segre_smul. pullbackBadImage_holds.proof:
% analyticAt_pullQ (through pullQ_eq_local, as in brQ_eq_local), the cover hsub by the sets S_1 and S_2 of the atlas
% charts, R.null j for S_1 with addHaar_image_eq_zero_of_differentiableOn_of_addHaar_eq_zero, and
% addHaar_image_eq_zero_of_det_fderivWithin_eq_zero with det_restrictScalars_eq_zero for S_2 (the Lean's form of the
% bound of Fremlin 263D(ii)). The Lean's statement (1) has no condition u_0 != 0; the two sets agree, as the elements of
% the fibre are nonzero. The Lean's pullBad uses the charts of the atlas R.A; the text's E uses any chart, and the chain
% rule moves between them.

\begin{lemma}\label{cv-origin-intersection-projection-generic-integral-section-count}
In Definition~\ref{cv-origin-intersection-projection-generic-integral-section-map}, let $V$ be the generic fibre of $X'$
over $Q$ (Lemma~\ref{cv-origin-intersection-projection-fibre}). For almost every $y\in\mathbf C^n$ and every
$w\in F_\Psi(y)$ (Definition~\ref{cv-origin-intersection-poincare-lelong-all-double-projection}(2), for $Q$), the number
of $x\in\tilde X'\setminus E$ with $q(x)=y$ and $W(x)=w$ is $\deg(V)$.
\end{lemma}
\leanref{GroundTheory/HeightsAndTheta/Heights/Intersections/ImportsFurtherFurther.lean}{GroundTheory.Heights₂.Heights.Inter.PullbackDegreeCount.Statement}
\leanref{GroundTheory/HeightsAndTheta/Heights/Intersections/NormSegreSameOrbit.lean}{GroundTheory.Heights₂.Heights.Inter.PullbackDegreeCount.proof}
\begin{proof}
\emph{The degree.} Let $\Phi\subset\Phi'$ and $R'\to\Phi'$ be as in
Lemma~\ref{cv-origin-intersection-projection-fibre}, with $\xi_{k',i'}$ the image of $x'_{k',i'}$. By that lemma, with
$r=0$, $V$ is a relevant prime generated by multihomogeneous elements and defines a projective variety of dimension $0$
over $\Phi$. Choose indices $i_{k'}$ with $x'_{k',i_{k'}}\notin V$, that is, $\xi_{k',i_{k'}}\neq0$, for
$k'\notin\kappa(K)$. As in the proof of Lemma~\ref{cv-origin-intersection-projection}, $\deg(V)$ is the dimension over
$\Phi$ of the coordinate ring of $V$ in this chart, which is $[\Phi(S):\Phi]$ by
Lemma~\ref{cv-origin-intersection-projection-count-chart}, for $F=\Phi$, $E=\Phi'$ and $z=(\xi_{k',i'})_{k'\notin\kappa(K)}$,
where $S$ is the set of the $\xi_{k',i'}/\xi_{k',i_{k'}}$.

\emph{The points over $w$.} Let $N_1$ and $N_2$ be the sets of measure zero of (1) and (2) of
Lemma~\ref{cv-origin-intersection-projection-generic-integral-section-count-negligible}, and $N_3$ the set of measure zero
off which the conclusion of Lemma~\ref{cv-origin-intersection-projection-generic-integral-section-count-orbits} holds,
for the chart $(i_{k'})$, with $u(z')$ and $O(w)$ as there. Let $y\notin N_1\cup N_2\cup N_3$ and $w\in F_\Psi(y)$, so that
$\Psi(w)=(1,y)$. Let $A_w$ be the set of the $x\in\tilde X'$ with $u_0(x)\neq0$ and $W(x)=w$. The points $x\in\tilde
X'\setminus E$ with $q(x)=y$ and $W(x)=w$ are those of $A_w$. Indeed, they lie in $A_w$, as $u_0$ does not vanish off $E$.
Conversely, let $x\in A_w$. Then $\Psi(W(x))=u_0(x)^{-1}u(x)=(1,q(x))$ and $\Psi(W(x))=\Psi(w)=(1,y)$, so $q(x)=y$; and
$x\notin E$, as $y\notin N_2$.

\emph{A bijection.} For $x\in\tilde X'$, $\lambda_x(x)\in\hat X'(\mathbf C)$ and $\pi(x)=\lambda_x(x)\circ v$, so
$u(x)=u(\lambda_x(x))$; for $x\in A_w$, $\lambda_x(x)$ is one of the points that define $O(w)$. So $x\mapsto$ the orbit
of $\lambda_x(x)$ maps $A_w$ to $O(w)$. As $y\notin N_1$, no point $z'\in\hat X'(\mathbf C)\cap Z(J')$ has $u(z')_0\neq0$
and $u(z')_0^{-1}\sigma(z'\circ v)=w$. The map is injective: if $x,x'\in A_w$ and $\lambda_x(x)$, $\lambda_{x'}(x')$ lie in
one orbit, then $\lambda_x(x)\notin Z(J')$ by the previous sentence, and $x=x'$ by (3) of
Definition~\ref{cv-origin-intersection-resolution}. It is onto: let $z'\in\hat X'(\mathbf C)$ with $u(z')_0\neq0$ and
$u(z')_0^{-1}\sigma(z'\circ v)=w$. Then $z'\notin Z(J')$, so by (2) of
Definition~\ref{cv-origin-intersection-resolution}, $z'=(c_{k'}\lambda_x(x)_{k',i})$ for some $x\in\tilde X'$ and
$c\in(\mathbf C^\times)^{K'}$. So $\sigma(z'\circ v)=\nu\,\sigma(\pi(x))$ with $\nu=\prod_kc_{\kappa(k)}\neq0$,
$u(x)=\nu^{-1}u(z')$, $u_0(x)\neq0$ and $W(x)=u(z')_0^{-1}\sigma(z'\circ v)=w$; and $\lambda_x(x)$ lies in the orbit of
$z'$.

So the number of the $x$ of the statement is the number of elements of $O(w)$, which is $[\Phi(S):\Phi]=\deg(V)$, as
$y\notin N_3$.
\end{proof}
\leanref{GroundTheory/HeightsAndTheta/Heights/Intersections/NormSegreSameOrbit.lean}{GroundTheory.Heights₂.Heights.Inter.PullbackDegreeCount.ncard_pullW_eq}
% Lean: PullbackDegreeCount.of_remaining = PullbackDegreeCount.proof with ChartRankAdjoin.full, GenericPointsCount
% (Lemma -count-orbits, under ChartNormCount.Statement through GenericPointsAvoidBadLocus.of_remaining),
% PullbackVanishingOfNullCycle.full and pullbackBadImage_holds.full (Lemma -count-negligible) and the hypothesis hDim, the
% FibreDimStatement (Lemma cv-origin-intersection-projection-fibre). hfib and hi for the chart, he for deg(V) (hJ.1 for the
% value of the field intersection family J in the chart i, and hRank), hN for N_1 and N_2, hc for N_3, hset for the
% points over w, and ncard_pullW_eq (with normSegre_sameOrbit, R.inj and R.surj) for the bijection.

\begin{lemma}\label{cv-origin-intersection-projection-generic-integral-section-density-sheet}
Let $(\kappa,v)$, $X'$, $P'$, $Q$, $J'$ and $(\tilde X',(U_x),(\lambda_x))$ be as in
Definition~\ref{cv-origin-intersection-projection-generic-integral-section-map}, and let
$\Psi:\mathbf C^{\prod_kI_k}\to\mathbf C^{n+1}$ be a $\mathbf C$-linear map. Define $\pi$, $u$, $q$, $W$ and $E$ as in that
definition, and sheets for $Q$ and $\Psi$ as in
Definition~\ref{cv-origin-intersection-poincare-lelong-all-double-projection}(2); neither uses that $\Psi$ is a good
projection. Let $\bar L_j=(a_j,h_j)$, $j=1,\dots,n$, be smooth metrized line bundles on
$\prod_k\mathbf P^{I_k}_{\mathbf Z}$ with pullbacks $\bar L'_j=(a'_j,h'_j)$, let $\bar A=(a,h_A)$ be a metrized line
bundle with pullback $\bar A'=(a',h_{A'})$, and let $s\in R\setminus Q$ be multihomogeneous of multidegree $a$. Let
$x\in\tilde X'\setminus E$, let $\kappa$ be a chart of $\tilde X'$ at $x$, let $M(x)$ be the Monge--Amp\`ere density at
$\kappa(x)$ of the functions $\log h'_j\circ\lambda_x\circ\kappa^{-1}$, and let $\Delta(x)$ be the Jacobian determinant of
$q\circ\kappa^{-1}$ at $\kappa(x)$. Then there is a sheet $\gamma$ through $W(x)$ at $q(x)$ with
\[
 \log\frac{|s(\gamma(q(x)))|}{h_A(\gamma(q(x)))}\operatorname{MA}(\log h_1\circ\gamma,\dots,\log h_n\circ\gamma)(q(x))
 =\log\frac{|v_*(s)(\lambda_x(x))|}{h_{A'}(\lambda_x(x))}\cdot\frac{M(x)}{|\Delta(x)|^2}.
\]
The proof does not use that $X'$ is an arithmetic subvariety of dimension $n+1$, beyond the existence of the resolution.
\end{lemma}
\leanref{GroundTheory/HeightsAndTheta/Heights/Intersections/ExistsSheetPullW.lean}{GroundTheory.Heights₂.Heights.Inter.pullbackDensity_holds.exists_sheet_pullW}
\begin{proof}
For $z'\in\hat X'(\mathbf C)$ write $z'\circ v=(z'_{v(k,i)})_{k,i}$, a point of $\hat Q(\mathbf C)$, as in the proof of
Lemma~\ref{cv-origin-intersection-projection-generic-integral-section-count-negligible}(1); so $\pi(z)=\lambda_z(z)\circ v$.
Put $b=\kappa(x)$, $\tilde q=q\circ\kappa^{-1}$ and $y_0=q(x)$. As $x\notin E$, $u_0(x)\neq0$, and $\tilde q$ is
holomorphic near $b$ with $\Delta(x)\neq0$; by the holomorphic inverse function theorem it has a holomorphic inverse $G$
near $y_0$ with $G(y_0)=b$ and $\det dG(y_0)=\Delta(x)^{-1}$. Near $y_0$ let
$\gamma_0(y')=\lambda_x(\kappa^{-1}(G(y')))\circ v$, holomorphic with $\gamma_0(y_0)=\pi(x)$, and
$\psi(y')=\Psi(\sigma(\gamma_0(y')))_0$, holomorphic with $\psi(y_0)=u_0(x)\neq0$. Fix $k_0\in K$, and for $y'$ near
$y_0$ let $\gamma(y')$ be $\gamma_0(y')$ with its block $k_0$ multiplied by $\psi(y')^{-1}$;
elsewhere let $\gamma$ be arbitrary.

\emph{It is a sheet.} Near $y_0$, $\gamma$ is holomorphic, and $\gamma(y')\in\hat Q(\mathbf C)$, as $\gamma_0(y')$ lies in
$\hat Q(\mathbf C)$ and $Q$ is generated by multihomogeneous elements. Its Segre image is
$\psi(y')^{-1}\sigma(\gamma_0(y'))$. By the first paragraph of the proof of
Lemma~\ref{cv-origin-intersection-projection-generic-integral-section-count-negligible}(2), $y'=\tilde q(G(y'))$ is the
vector of the quotients of the coordinates $1,\dots,n$ of $\Psi(\sigma(\gamma_0(y')))$ by its coordinate $0$, which is
$\psi(y')$. So $\Psi(\sigma(\gamma(y')))=(1,y')$. Finally $\sigma(\gamma(y_0))=u_0(x)^{-1}\sigma(\pi(x))=W(x)$.

\emph{The density.} As $s$ and $h_A$ are homogeneous of the same multidegree $a$, $|s|/h_A$ takes the same value at
$\gamma(y_0)$ and at $\gamma_0(y_0)=\pi(x)$, and this value is $|v_*(s)(\lambda_x(x))|/h_{A'}(\lambda_x(x))$, as
$v_*(s)(z')=s(z'\circ v)$ and $h_{A'}(z')=h_A(z'\circ v)$. Near $y_0$,
$\log h_j\circ\gamma=\log h_j\circ\gamma_0-a_{j,k_0}\log|\psi|$, and $\log|\psi/\psi(y_0)|$ is the real part of a
holomorphic function. So the complex Hessians of $\log h_j\circ\gamma$ and of $\log h_j\circ\gamma_0$ at $y_0$ agree, and
so do their Monge--Amp\`ere densities (Definition~\ref{cv-period-monge-ampere}(2)). As
$h'_j(z')=h_j(z'\circ v)$, $\log h_j\circ\gamma_0=\log h'_j\circ\lambda_x\circ\kappa^{-1}\circ G$, and the functions
$\log h'_j\circ\lambda_x\circ\kappa^{-1}$ are smooth near $b$, as $\pi(x)$ has a nonzero coordinate in every block. By
Lemma~\ref{cv-origin-intersection-poincare-lelong-all-double-bridge-ma}, applied to $G$, the density of the
$\log h_j\circ\gamma_0$ at $y_0$ is $M(x)\,|\det dG(y_0)|^2=M(x)/|\Delta(x)|^2$.
\end{proof}
% Lean: exists_sheet_pullW, in the chart chartIdx R.A x: hgood from x not in pullBad, the local inverse F of
% Q = pullQ o chart^-1 (toOpenPartialHomeomorph, analyticAt_symm, hasFDerivAt_symm), gamma0 = lift x o chart^-1 o F.symm
% composed with v, psi0, the block k0 scaled by psi0^-1 (bsmul with c = update 1 k0 psi0^-1), pullQ_eq_local and
% mcPts_smul with comp_mem_mcPts for the sheet, logNorm_bsmul and hAA' for the value of |s|/h_A, cxHessian_congr and
% cxHessian_add_re for the pluriharmonic term, maDens_comp with det_jacC and det_coe_symm, and hLL' (locDens of potOf R
% (L' j)) for the density. As for Lemma cv-origin-intersection-poincare-lelong-all-double-bridge-sheet, the text states it
% in any chart. The Lean's Psi is any linear map, and it does not assume IsArithVar P'.

\begin{lemma}\label{cv-origin-intersection-projection-generic-integral-section-density}
In Definition~\ref{cv-origin-intersection-projection-generic-integral-section-map}, let $\bar L_1,\dots,\bar L_n$ and
$\bar A=(a,h_A)$ be smooth metrized line bundles on $\prod_k\mathbf P^{I_k}_{\mathbf Z}$ with pullbacks
$\bar L'_j=(a'_j,h'_j)$ and $\bar A'=(a',h_{A'})$, and let $s\in R\setminus Q$ be multihomogeneous of multidegree $a$. For
$x\in\tilde X'\setminus E$, with a chart $\kappa$ at $x$, let $p(x)=M(x)/|\Delta(x)|^2$, where $M(x)$ is the
Monge--Amp\`ere density at $\kappa(x)$ of the potentials $\log h'_j\circ\lambda_x\circ\kappa^{-1}$ and $\Delta(x)$ is the
Jacobian determinant of $q\circ\kappa^{-1}$ at $\kappa(x)$. Then for almost every $y\in\mathbf C^n$ and every
$x\in q^{-1}(y)\setminus E$,
\[
 \log\frac{|v_*(s)(\lambda_x(x))|}{h_{A'}(\lambda_x(x))}\,p(x)
\]
is the sheet density at $(y,W(x))$ of $\log(|s|/h_A)$ against $\log h_1,\dots,\log h_n$, for $Q$ and $\Psi$
(Definition~\ref{cv-origin-intersection-poincare-lelong-all-double-projection}(3)). The proof does not use that $X'$ is
an arithmetic subvariety of dimension $n+1$, beyond the existence of the resolution.
\end{lemma}
\leanref{GroundTheory/HeightsAndTheta/Heights/Intersections/ImportsFurtherExtra.lean}{GroundTheory.Heights₂.Heights.Inter.pullbackDensity_holds.Statement}
\leanref{GroundTheory/HeightsAndTheta/Heights/Intersections/Assembly.lean}{GroundTheory.Heights₂.Heights.Inter.pullbackDensity_holds.full}
\begin{proof}
By Lemma~\ref{cv-origin-intersection-poincare-lelong-all-double-sheets}, for $Q$, $\Psi$, the $\bar L_j$, $\bar A$ and
$s$, almost every $y\in\mathbf C^n$ has the following property: any two sheets through the same point at $y$ give the same
value of $\log(|s|/h_A)$ at $y$ times the Monge--Amp\`ere density of the $\log h_j$ composed with the sheet. Let $y$ be
such a point and $x\in q^{-1}(y)\setminus E$. Lemma~\ref{cv-origin-intersection-projection-generic-integral-section-density-sheet}
gives a sheet $\gamma$ through $W(x)$ at $y$ for which this value is
$\log\big(|v_*(s)(\lambda_x(x))|/h_{A'}(\lambda_x(x))\big)\,p(x)$. As a sheet through $W(x)$ at $y$ exists, the sheet
density at $(y,W(x))$ is the value for the chosen sheet, which is the value for $\gamma$.
\end{proof}
\leanref{GroundTheory/HeightsAndTheta/Heights/Intersections/ExistsSheetPullW.lean}{GroundTheory.Heights₂.Heights.Inter.pullbackDensity_holds.proof}
% Lean: pullbackDensity_holds.full = pullbackDensity_holds.proof: sheetIndep (Lemma
% cv-origin-intersection-poincare-lelong-all-double-sheets) for imageZ v P', exists_sheet_pullW (Lemma -density-sheet), and
% sheetDensity with dif_pos for the chosen sheet. The hypothesis IsArithVar P' (n + 1) is not used (_hP').

\begin{lemma}\label{cv-origin-intersection-projection-generic-integral-section}
Let $(\kappa,v)$, $X'$, $P'$, $Q$, $\Phi$ and $V$ be as in Lemma~\ref{cv-origin-intersection-projection-fibre} with $r=0$,
with $X'$ and $Q$ of dimension $n+1$. Let $\bar L_1,\dots,\bar L_n$ and $\bar A=(a,h_A)$ be smooth metrized line bundles on
$\prod_k\mathbf P^{I_k}_{\mathbf Z}$ with pullbacks $\bar L'_1,\dots,\bar L'_n$ and $\bar A'=(a',h_{A'})$, and let
$s\in R\setminus Q$ be multihomogeneous of multidegree $a$. With $\Phi_Q$ and $\Phi_{X'}$ the products of the first Chern
forms of the $\bar L_j$ and of the $\bar L'_j$,
\[
 \int_{X'(\mathbf C)}\log\frac{|v_*(s)|}{h_{A'}}\,\Phi_{X'}=\deg(V)\int_{Q(\mathbf C)}\log\frac{|s|}{h_A}\,\Phi_Q.
\]
\end{lemma}
\leanref{GroundTheory/HeightsAndTheta/Heights/Intersections/ImportsFurtherBis.lean}{GroundTheory.Heights₂.Heights.Inter.IntersectionWithSection.Statement}
\leanref{GroundTheory/HeightsAndTheta/Heights/Intersections/PullWMemFibre.lean}{GroundTheory.Heights₂.Heights.Inter.IntersectionWithSection.proof}
\begin{proof}
The pullbacks $\bar A'$ and $\bar L'_j$ are smooth, and $v_*(s)$ is multihomogeneous of multidegree $a'$ and not in $P'$,
as $s\notin Q=v_*^{-1}(P')$. As $P'$ is relevant, some $x'_{k',i}$ is not in $P'$, and by
Lemma~\ref{cv-origin-intersection-resolution-exists} for the ideal $(x'_{k',i})$, $X'$ has an analytic resolution
$(\tilde X',(U_x),(\lambda_x))$ off some ideal $J'$. By
Lemma~\ref{cv-origin-intersection-poincare-lelong-all-double-good}(1), $Q$ has a good projection $\Psi$. Let $\pi$, $q$,
$W$ and $E$ be as in Definition~\ref{cv-origin-intersection-projection-generic-integral-section-map}, and for
$x\in\tilde X'$ let $g(x)=\log\big(|v_*(s)(\lambda_x(x))|/h_{A'}(\lambda_x(x))\big)$, with $\log0$ read as $0$.

\emph{The integral over $X'$.} By
Lemmas~\ref{cv-origin-intersection-poincare-lelong-all-double-independent}(2)
and~\ref{cv-origin-intersection-poincare-lelong-all-double-bridge}, the left side is the integral over $\tilde X'$ of $g$
against the Monge--Amp\`ere measure of the potentials $\log h'_j\circ\lambda_x$, and the first paragraph of the proof of
Lemma~\ref{cv-origin-intersection-poincare-lelong-all-double-bridge} shows that $g$ times the density is integrable on
compact subsets of the charts. Lemma~\ref{cv-origin-intersection-poincare-lelong-all-double-bridge-area} applies to
$\tilde X'$, $g$, $q$ and $E$: its hypotheses (1) and (3) are
Lemma~\ref{cv-origin-intersection-projection-generic-integral-section-null}, and (2) holds by the definition of $E$. So,
with $p$ as there, the left side is the integral over $y\in\mathbf C^n$ of $\sum_{x\in q^{-1}(y)\setminus E}g(x)p(x)$.

\emph{The fibre sums.} Let $x\notin E$, so that $u_0(x)\neq0$. The point $\pi(x)$ lies in $\hat Q(\mathbf C)$, as
$f(\pi(x))=v_*(f)(\lambda_x(x))=0$ for $f\in Q$ and its blocks are those of $\lambda_x(x)$ in $\kappa(K)$. Dividing one
block of $\pi(x)$ by $u_0(x)$ gives $z\in\hat Q(\mathbf C)$ with $\sigma(z)=W(x)$ and $\Psi(\sigma(z))=(1,q(x))$; so
$W(x)\in F_\Psi(q(x))$.
By Lemmas~\ref{cv-origin-intersection-projection-generic-integral-section-count}
and~\ref{cv-origin-intersection-projection-generic-integral-section-density}, for almost every $y$, every
$w\in F_\Psi(y)$ is $W(x)$ for $\deg(V)$ points $x\in q^{-1}(y)\setminus E$, and for each of them $g(x)p(x)$ is
the sheet density $S(y,w)$ of $\log(|s|/h_A)$ against $\log h_1,\dots,\log h_n$ at $(y,w)$. Grouping the finitely many
$x\in q^{-1}(y)\setminus E$ according to $W(x)$,
\[
 \sum_{x\in q^{-1}(y)\setminus E}g(x)p(x)=\deg(V)\sum_{w\in F_\Psi(y)}S(y,w)
\]
for almost every $y$. By Lemma~\ref{cv-origin-intersection-poincare-lelong-all-double-independent}(2), the integral of
the right side over $\mathbf C^n$ is $\deg(V)\int_{Q(\mathbf C)}\log(|s|/h_A)\,\Phi_Q$.
% Lean: IntersectionWithSection.proof (PullWMemFibre.lean) from ResExistsStatement (exists_anLogRes, Lemma
% cv-origin-intersection-resolution-exists), goodProjExists, pullbackPushforward_holds.full (hnull, hfin, Lemma
% -integral-section-null), ptsIntegral_eq_mInt', ptsIntegral_eq_at, pushForward_null (Lemma
% cv-origin-intersection-poincare-lelong-all-double-bridge-area) with brIntegrable, pullW_mem_fibre,
% PullbackDegreeCount.of_remaining (-count, with J (fibreIdeal kappa G) 0 for deg(V)), pullbackDensity_holds.full
% (-density) and Finset.sum_fiberwise_of_maps_to'. Size: PullWMemFibre.lean, about 120 lines, besides these lemmas (about
% 8,000 lines in all, besides Lemma cv-origin-intersection-resolution-exists; the open steps are in -null-orbits and
% -null-isolated).
\end{proof}

\begin{lemma}\label{cv-origin-intersection-projection-generic-integral}
Let $(\kappa,v)$, $X'$, $P'$, $Q$, $\Phi$ and $V$ be as in Lemma~\ref{cv-origin-intersection-projection-fibre} with $r=0$,
with $X'$ and $Q$ of dimension $n+1$. Let $\bar L_1,\dots,\bar L_n$ and $\bar A=(a,h_A)$ be smooth metrized line bundles on
$\prod_k\mathbf P^{I_k}_{\mathbf Z}$ with pullbacks $\bar L'_1,\dots,\bar L'_n$ and $\bar A'=(a',h_{A'})$, and let
$f,g\in R\setminus Q$ be multihomogeneous of multidegrees $\alpha,\beta\in\mathbf N^K$ with $\alpha-\beta=a$. Put
$\Phi_Q=c_1(\bar L_1)\wedge\dots\wedge c_1(\bar L_n)$ and $\Phi_{X'}=c_1(\bar L'_1)\wedge\dots\wedge c_1(\bar L'_n)$. Then
\[
 \int_{X'(\mathbf C)}\log\frac{|v_*(f)|}{|v_*(g)|\,h_{A'}}\,\Phi_{X'}
 =\deg(V)\int_{Q(\mathbf C)}\log\frac{|f|}{|g|\,h_A}\,\Phi_Q.
\]
\end{lemma}
\leanref{GroundTheory/HeightsAndTheta/Heights/Intersections/IsFieldIntersectionNoBlockFam.lean}{GroundTheory.Heights₂.Heights.Inter.ArithmeticIntersectionIntegral.Statement}
\leanref{GroundTheory/HeightsAndTheta/Heights/Intersections/PtsIntegralLogNormRat.lean}{GroundTheory.Heights₂.Heights.Inter.ArithmeticIntersectionIntegral.proof}
\begin{proof}
Let $\bar B=\bar{\mathcal O}(\beta)=(\beta,h_\beta)$ be the Fubini--Study bundle
(Definition~\ref{cv-origin-analytic-system-pullback}(1)), which is smooth, and let $\bar B'=(\beta',h_{\beta'})$ be its
pullback, with $h_{\beta'}(z')=h_\beta(z'\circ v)$. Then $f$ is a section of $\bar A\otimes\bar B=(\alpha,h_Ah_\beta)$ and
$g$ one of $\bar B$, neither in $Q$. The elements $v_*(f)$ and $v_*(g)$ are multihomogeneous of the multidegrees $\alpha$
and $\beta$ extended by $0$, which are those of the pullbacks $\bar A'\otimes\bar B'$ and $\bar B'$, and they are not in
$P'$, as $Q=v_*^{-1}(P')$. The pullbacks of smooth bundles are smooth.

As in the proof of Lemma~\ref{cv-origin-intersection-projection-generic-rational}, the two integrands below are integrable
(Lemma~\ref{cv-origin-intersection-poincare-lelong-all-double-independent}(1)), and for almost every point of the base
neither section vanishes at the chosen sheets; so
\[
 \int_{Q(\mathbf C)}\log\frac{|f|}{|g|\,h_A}\,\Phi_Q=\int_{Q(\mathbf C)}\log\frac{|f|}{h_Ah_\beta}\,\Phi_Q
 -\int_{Q(\mathbf C)}\log\frac{|g|}{h_\beta}\,\Phi_Q,
\]
and likewise on $X'$ with $v_*(f)$, $v_*(g)$, $h_{A'}$ and $h_{\beta'}$. By
Lemma~\ref{cv-origin-intersection-projection-generic-integral-section}, for $f$ and $\bar A\otimes\bar B$ and for $g$ and
$\bar B$, each of the two integrals on $X'$ is $\deg(V)$ times the corresponding integral on $Q$. Subtracting gives the
statement.
\end{proof}
\leanref{GroundTheory/HeightsAndTheta/Heights/Intersections/PtsIntegralLogNormRat.lean}{GroundTheory.Heights₂.Heights.Inter.ArithmeticIntersectionIntegral.ptsIntegral_logNormRat}
% Lean: ArithmeticIntersectionIntegral.proof (PtsIntegralLogNormRat.lean) under PtsIntegrableStatement
% (PtsIntegrable.of_resExists) and IntersectionWithSection.Statement: B = sqMetric beta, B' = pullHerm (isPullHerm_pullHerm,
% isSmoothHerm_pull), hfs and hgs on imageZ v P', hfs' and hgs' on P' (isMultiHomog_rename), ptsIntegral_logNormRat on
% both sides (integral_sub with the two integrability statements and ae_eval_ne_zero), hSec twice and mul_sub. Size:
% PtsIntegralLogNormRat.lean, about 200 lines, besides Lemma
% cv-origin-intersection-projection-generic-integral-section.

\begin{lemma}\label{cv-origin-intersection-projection-generic}
Let $(\cdot\mid\cdot)$ be an arithmetic intersection system, and let $(\kappa,v)$, $X'$, $P'$, $Q$, $\Phi$ and $V$ be as in
Lemma~\ref{cv-origin-intersection-projection-fibre} with $r=0$, so that $X'$ and $Q$ both have dimension $d$. Let
$\bar L_1,\dots,\bar L_d$ be smooth metrized line bundles on $\prod_k\mathbf P^{I_k}_{\mathbf Z}$ with pullbacks
$\bar L'_1,\dots,\bar L'_d$. Then
\[
 (\bar L'_1\cdots\bar L'_d\mid X')=\deg(V)\,(\bar L_1\cdots\bar L_d\mid Q),
\]
where $\deg(V)$ is the intersection number of the zero-dimensional variety defined by $V$ over $\Phi$
\cite[Tag 0BEP]{Sta}.
\end{lemma}
\leanref{ReciprocityTheory/ReductiveGroups/ArithmeticGeometry/ArithmeticIntersection/MixedFieldContraction/ContractMixedFieldStatement.lean}{ReciprocityTheory.ChartRestriction.G.AI.I4.GenDegStatement}
\leanref{Development/AutomorphicSpectrum/Proofs/MetricProofs/DegreeAndCoordinatesOfResolution.lean}{GlobalTheory.CompanionTheory.Ar.Pub.genDeg_of_resExists}
\begin{proof}
By induction on $d$, for all block maps at once. For $d=0$ there is nothing to prove, as there is no arithmetic
subvariety of dimension $0$ (Lemma~\ref{cv-origin-intersection-family-curves-points}(1)). Let $d\ge1$ and assume the
lemma in dimension $d-1$. Let $\bar A=(a,h_A)=\bar L_d$, with pullback $\bar A'=\bar L'_d=(a',h_{A'})$, where $a'$ is $a$
extended by $0$ and $h_{A'}(z')=h_A(z'\circ v)$; the pullbacks of smooth bundles are smooth. Put
$\alpha_k=\max(a_k,0)$ and $\beta_k=\max(-a_k,0)$. As $Q$ is relevant, there are $i_k$ with $x_{k,i_k}\notin Q$; let
$f=\prod_kx_{k,i_k}^{\alpha_k}$ and $g=\prod_kx_{k,i_k}^{\beta_k}$, multihomogeneous of multidegrees $\alpha$ and $\beta$,
with $\alpha-\beta=a$, and not in the prime $Q$. Then $v_*(f)$ and $v_*(g)$ are not in $P'$, as $Q=v_*^{-1}(P')$, and they
are multihomogeneous of the multidegrees $\alpha'$ and $\beta'$ obtained by extending $\alpha$ and $\beta$ by $0$, with
$\alpha'-\beta'=a'$. Write $\bar L$ for $\bar L_1\cdots\bar L_{d-1}$ and $\bar L'$ for $\bar L'_1\cdots\bar L'_{d-1}$, and
$\Phi_Q$, $\Phi_{X'}$ for the products of their first Chern forms. By
Lemma~\ref{cv-origin-intersection-projection-generic-rational} on $Q$ and on $X'$,
\[
 (\bar L\bar A\mid Q)=\mathcal C(\bar L;f\mid Q)-\mathcal C(\bar L;g\mid Q)
 -\int_{Q(\mathbf C)}\log\frac{|f|}{|g|\,h_A}\,\Phi_Q,
\]
\[
 (\bar L'\bar A'\mid X')=\mathcal C(\bar L';v_*(f)\mid X')-\mathcal C(\bar L';v_*(g)\mid X')
 -\int_{X'(\mathbf C)}\log\frac{|v_*(f)|}{|v_*(g)|\,h_{A'}}\,\Phi_{X'}.
\]
By Lemma~\ref{cv-origin-intersection-projection-generic-cycles}, with the induction hypothesis, each $\mathcal C$ term on
$X'$ is $\deg(V)$ times the corresponding one on $Q$, and by
Lemma~\ref{cv-origin-intersection-projection-generic-integral} so is the integral. So
$(\bar L'\bar A'\mid X')=\deg(V)\,(\bar L\bar A\mid Q)$.
\end{proof}
\leanref{GroundTheory/HeightsAndTheta/Heights/Intersections/IsMultiHomogRenameCoordPow.lean}{GroundTheory.Heights₂.Heights.Inter.PushforwardArithIntersection.proof}
\leanref{GroundTheory/HeightsAndTheta/Heights/Intersections/IsMultiHomogRenameCoordPow.lean}{GroundTheory.Heights₂.Heights.Inter.PushforwardArithIntersection.genDegAt_succ}
% Lean: genDeg_of_resExists under ResExistsStatement: GenDeg.of_pts (PtsIntegrable.of_resExists) = GenDeg.of_found
% (RatR.proof, ArithmeticIntersectionCycle.of_remaining, ArithmeticIntersectionIntegral.proof with
% IntersectionWithSection.proof) = PushforwardArithIntersection.proof arithVarZero: the induction on d of GenDegAt I d,
% genDegAt_zero (arithVarZero) and genDegAt_succ (alpha and beta by toNat, the coordinate powers coordPow with
% coordPow_not_mem, isMultiHomog_rename_coordPow, isSmoothHerm_pull, hRat on both sides, hCyc for f and g, hInt). The
% factor is J (fibreIdeal kappa G) 0 for any field intersection family J over G.Phi, the text's deg(V). Size: about 160
% lines in IsMultiHomogRenameCoordPow.lean, besides Lemmas cv-origin-intersection-projection-generic-rational, -cycles and
% -integral, which hold most of the 24,500 lines of Lemma cv-origin-intersection-projection.

\begin{lemma}\label{cv-origin-intersection-projection-count-chart}
Let $F\subset E$ be fields, let $K_0$ be a finite set, for $k\in K_0$ let $J_k$ be a finite nonempty set, let
$z\in E^{\sqcup_kJ_k}$, and let $j_k\in J_k$ with $z_{k,j_k}\neq0$ for every $k$. Suppose that the kernel $V$ of
$F[x_{k,i}]\to E$, $f\mapsto f(z)$, is generated by multihomogeneous elements, and put $u_{k,i}=z_{k,i}/z_{k,j_k}$. Then
$f\mapsto f(u)$ induces an isomorphism of $F$-algebras from $F[x_{k,i}]/(V+(x_{k,j_k}-1:k\in K_0))$ onto $F[u]$. Its
dimension over $F$ is finite if and only if $[F(u):F]$ is, and then the two are equal.
\end{lemma}
\leanref{GroundTheory/HeightsAndTheta/Heights/Intersections/StatementsChartRankAdjoin.lean}{GroundTheory.Heights₂.Heights.Inter.ChartRankAdjoin.Statement}
\leanref{GroundTheory/HeightsAndTheta/Heights/Intersections/PChartRankAdjoin.lean}{GroundTheory.Heights₂.Heights.Inter.ChartRankAdjoin.proof}
\begin{proof}
The homomorphism $f\mapsto f(u)$ is onto $F[u]$, and its kernel contains the $x_{k,j_k}-1$. It contains $V$: if
$f\in V$ is multihomogeneous of multidegree $a$, then $f(u)=f(z)/\prod_kz_{k,j_k}^{a_k}=0$. Conversely let $f(u)=0$.
Multiplying each monomial of $f$ by powers of the $x_{k,j_k}$ gives a multihomogeneous $h$ of some multidegree $N$ with
$h\equiv f$ modulo the $x_{k,j_k}-1$. Then $h(u)=f(u)=0$ and $h(z)=\prod_kz_{k,j_k}^{N_k}h(u)=0$, so $h\in V$ and
$f\in V+(x_{k,j_k}-1:k)$. If every $u_{k,i}$ is algebraic over $F$, the domain $F[u]$ is integral over $F$, hence a
field \cite[Tag 00GS]{Sta}, and $F[u]=F(u)$. Otherwise some $u_{k,i}$ is transcendental over $F$, and both $F[u]$ and
$F(u)$ contain the polynomial ring $F[u_{k,i}]$, of infinite dimension.
\end{proof}
% Lean: ChartRankAdjoin.proof: chartIdeal_eq_ker (aeval_mul_of_isMultiHomog for V, NumbersDVR.IntNumPoint.exists_homog_congr
% for h), quotientKerEquivRange, adjoin_toSubalgebra_of_isAlgebraic, and finrank_of_not_finite on both sides otherwise
% (the Lean's finrank is 0 for an infinite dimension). Size: PChartRankAdjoin.lean, 110 lines.

\begin{lemma}\label{cv-origin-intersection-projection-count-scaling}
Let $F\subset E$, $K_0$, $(J_k)$, $z$, $(j_k)$, $V$ and $u$ be as in
Lemma~\ref{cv-origin-intersection-projection-count-chart}. Then the elements $z_{k,j_k}$, $k\in K_0$, are algebraically
independent over the subfield $F(u)$ of $E$.
\end{lemma}
\leanref{GroundTheory/HeightsAndTheta/Heights/Intersections/StChartScalingIndep.lean}{GroundTheory.Heights₂.Heights.Inter.ChartScalingIndep.Statement}
\leanref{GroundTheory/HeightsAndTheta/Heights/Intersections/PChartScalingIndep.lean}{GroundTheory.Heights₂.Heights.Inter.ChartScalingIndep.proof}
\begin{proof}
Write $t_k=z_{k,j_k}$ and $t^\alpha=\prod_kt_k^{\alpha_k}$ for $\alpha\in\mathbf N^{K_0}$. First, let $S$ be a finite set of
exponents and $g_\alpha$, $\alpha\in S$, polynomials over $F$ with $\sum_{\alpha\in S}g_\alpha(u)t^\alpha=0$; then every
$g_\alpha(u)$ is $0$. Indeed, as in the proof of Lemma~\ref{cv-origin-intersection-projection-count-chart}, there are a
multidegree $M$ and multihomogeneous $H_\alpha$ of multidegree $M$ with $H_\alpha\equiv g_\alpha$ modulo the
$x_{k,j_k}-1$, so that $H_\alpha(z)=t^Mg_\alpha(u)$. The polynomial $H=\sum_\alpha H_\alpha\prod_kx_{k,j_k}^{\alpha_k}$ has
$H(z)=t^M\sum_\alpha g_\alpha(u)t^\alpha=0$, so $H\in V$. Its multihomogeneous component of multidegree $M+\beta$ is
$H_\beta\prod_kx_{k,j_k}^{\beta_k}$, and it lies in $V$, as $V$ is generated by multihomogeneous elements. So
$t^Mg_\beta(u)t^\beta=0$, and $g_\beta(u)=0$.

Now let $p=\sum_\alpha c_\alpha T^\alpha$ be a polynomial with coefficients $c_\alpha\in F(u)$ and $p(t)=0$. Write
$c_\alpha=a_\alpha(u)/b_\alpha(u)$ with polynomials $a_\alpha,b_\alpha$ over $F$ and $b_\alpha(u)\neq0$, and let $D$ be the
product of the $b_\alpha$. With $g_\alpha=a_\alpha\prod_{\beta\neq\alpha}b_\beta$, $g_\alpha(u)=D(u)c_\alpha$ and
$\sum_\alpha g_\alpha(u)t^\alpha=D(u)p(t)=0$. By the first paragraph $D(u)c_\alpha=0$, so $c_\alpha=0$ for every $\alpha$,
and $p=0$.
\end{proof}
% Lean: ChartScalingIndep.proof: algebraicIndependent_iff, IntermediateField.mem_adjoin_iff_div for the coefficients,
% the common denominator D, and ChartScalingIndep.coeff_eq_zero (exists_homog_chart for the H_alpha, aeval_blockScale,
% the weighted homogeneous component of multidegree M + beta, isHomogeneous_mw). Size: PChartScalingIndep.lean, 156 lines.

\begin{lemma}\label{cv-origin-intersection-projection-count-disjoint}
Let $F_0\subset F_1$ be subfields of a field $E$, and let $t_1,\dots,t_m\in E$ be algebraically independent over $F_1$.
Then $[F_1(t):F_0(t)]$ is finite if and only if $[F_1:F_0]$ is, and then the two degrees are equal.
\end{lemma}
\leanref{GroundTheory/HeightsAndTheta/Heights/Intersections/StAdjoinIndepRelfinrank.lean}{GroundTheory.Heights₂.Heights.Inter.AdjoinIndepRelfinrank.Statement}
\leanref{GroundTheory/HeightsAndTheta/Heights/Intersections/PAdjoinIndepRelfinrank.lean}{GroundTheory.Heights₂.Heights.Inter.AdjoinIndepRelfinrank.proof}
\begin{proof}
Let $(b_j)$ be a basis of $F_1$ over $F_0$. It is linearly independent over $F_0[t]$: if $\sum_jp_j(t)b_j=0$ with
polynomials $p_j$ over $F_0$, almost all $0$, then $\sum_jb_jp_j$ is a polynomial over $F_1$ that vanishes at $t$, hence
is $0$, and comparing coefficients every $p_j$ is $0$, as the $b_j$ are independent over $F_0$. Clearing denominators,
$(b_j)$ is linearly independent over $F_0(t)$. If $[F_1:F_0]=n$ is finite, the $F_0(t)$-span of $b_1,\dots,b_n$ is closed
under products, as each $b_ib_j$ is an $F_0$-combination of the $b_l$. It is a domain of dimension $n$ over the field
$F_0(t)$, hence a field, and it contains $F_1$ and $t$. So it is $F_1(t)$, and $[F_1(t):F_0(t)]=n$. If $[F_1:F_0]$ is
infinite, so is $[F_1(t):F_0(t)]$, by the linear independence. Conversely, if $[F_1(t):F_0(t)]$ is finite, the
independence gives $[F_1:F_0]\le[F_1(t):F_0(t)]$.
\end{proof}
% Lean: AdjoinIndepRelfinrank.proof: the basis b of L over L0 is linearly independent over Algebra.adjoin L0 (range t)
% (hR, through MvPolynomial coefficients) and over its fraction field (LinearIndependent.iff_fractionRing), so A and L'
% are linearly disjoint (LinearDisjoint.of_basis_right'); adjoin_rank_eq_rank_right_of_isAlgebraic_right in the finite
% case, finrank_of_not_finite in the other. Size: PAdjoinIndepRelfinrank.lean, 93 lines.

\begin{lemma}\label{cv-origin-intersection-projection-count}
Let $(\kappa,v)$ be a block map, let $\Omega$ be a field of characteristic zero, and let $y\in\Omega^{\sqcup_{k'}I'_{k'}}$
be a point whose kernel $P'$ and the kernel $Q=v_*^{-1}(P')$ of $y\circ v$ are arithmetic subvarieties of the same
dimension $d$. Let $V$ be the generic fibre of Lemma~\ref{cv-origin-intersection-projection-fibre}, and let
$(i_{k'})_{k'\notin\kappa(K)}$ be indices with $x'_{k',i_{k'}}\notin V$. Then the dimension over $\Phi$ of the
coordinate ring of $V$ in this chart, the quotient of $\Phi[x'_{k',i'}:k'\notin\kappa(K)]$ by $V$ and the
$x'_{k',i_{k'}}-1$, is the degree $[\mathbf Q(y):\mathbf Q(y\circ v)]$ of the fields of ratios of
Definition~\ref{cv-origin-intersection-full}(2).
\end{lemma}
\leanref{GroundTheory/HeightsAndTheta/Heights/OrbitLimits/CoordFibreDeg.lean}{GroundTheory.Heights₂.Heights.Targets.CoordFibreDeg.Statement}
\leanref{GroundTheory/HeightsAndTheta/Heights/Intersections/Assembly.lean}{GroundTheory.Heights₂.Heights.Inter.Closed.Targets.CoordFibreDeg}
\begin{proof}
Let $\Phi'$ and $R'\to\Phi'$ be as in Lemma~\ref{cv-origin-intersection-projection-fibre}, let $\xi_{k',i'}\in\Phi'$ be
the image of $x'_{k',i'}$, and write $K'_0=K'\setminus\kappa(K)$. Identify $\Phi$ with its image in $\Phi'$, so that
$\xi_{v(k,i)}$ is the image of $x_{k,i}$ in $\Phi$. The field $\Phi'$ has characteristic zero, as $P'\cap\mathbf Z=0$.
For a family $w$ of elements of a field of characteristic zero, indexed by $\sqcup_{k'}I'_{k'}$ or by
$\sqcup_kI_k$, write $\mathbf Q(w)$ for the subfield generated by the ratios $w_{k',a}/w_{k',b}$ of coordinates in one block
with $w_{k',b}\neq0$; so $\mathbf Q(y)$ and $\mathbf Q(y\circ v)$ are those of
Definition~\ref{cv-origin-intersection-full}(2).

\emph{The chart.} By Lemma~\ref{cv-origin-intersection-projection-fibre}, $V$ is generated by multihomogeneous elements,
and $\xi_{k',i_{k'}}\neq0$ for $k'\in K'_0$, as $x'_{k',i_{k'}}\notin V$. By
Lemma~\ref{cv-origin-intersection-projection-count-chart}, for $F=\Phi$, $E=\Phi'$ and $z=(\xi_{k',i'})_{k'\in K'_0}$, the
dimension of the statement is $[\Phi(S):\Phi]$, where $S$ is the set of the $\xi_{k',i'}/\xi_{k',i_{k'}}$ with
$k'\in K'_0$: the two are finite together, and then equal.

\emph{The fibre and the ratios of $\xi$.} We show $[\Phi(S):\Phi]=[\mathbf Q(\xi):\mathbf Q(\xi\circ v)]$. As $Q$ is
relevant, for each $k\in K$ there is $c_k$ with $x_{k,c_k}\notin Q$, so $\xi_{v(k,c_k)}\neq0$. Put $j_{k'}=v(k,c_k)$ for
$k'=\kappa(k)$ and $j_{k'}=i_{k'}$ for $k'\in K'_0$, and let $t_k=\xi_{\kappa(k),j_{\kappa(k)}}$ for $k\in K$. Let $E_0$
be the subfield generated by the $t_k$, $F_0=\mathbf Q(\xi\circ v)$ and $F_1=\mathbf Q(\xi)$, so that $F_0\subset F_1$.
\begin{enumerate}
\item[(a)] $\Phi=F_0E_0$. Each $t_k$ lies in $\Phi$, and so does $F_0$, generated by ratios of the $\xi_{v(k,i)}$.
  Conversely $\Phi$ is generated by the $\xi_{v(k,i)}=(\xi_{v(k,i)}/\xi_{v(k,c_k)})\,t_k$, which lie in $F_0E_0$.
\item[(b)] $\Phi(S)=F_1E_0$. By (a), $\Phi(S)=F_0E_0(S)\subset F_1E_0$. Conversely, a generator
  $\xi_{k',a}/\xi_{k',b}$ of $F_1$ lies in $F_0$ if $k'=\kappa(k)$, as $v$ maps $I_k$ onto $I'_{\kappa(k)}$, and is the
  quotient of two elements of $S$ if $k'\in K'_0$.
\item[(c)] The $t_k$ are algebraically independent over $F_1$. The kernel of $\mathbf Q[x'_{k',i'}]\to\Phi'$,
  $f\mapsto f(\xi)$, is generated by $P'$, as a multiple of an element of the kernel by a nonzero integer lies in $R'$;
  so it is generated by multihomogeneous elements. By
  Lemma~\ref{cv-origin-intersection-projection-count-scaling}, for $F=\mathbf Q$, $z=\xi$ and the indices $j_{k'}$, the
  $\xi_{k',j_{k'}}$, $k'\in K'$, are algebraically independent over the field generated by the
  $\xi_{k',a}/\xi_{k',j_{k'}}$, which is $F_1$, as $\xi_{k',a}/\xi_{k',b}$ is a quotient of two of these. The $t_k$ are
  among them, as $\kappa$ is injective.
\end{enumerate}
By (c) and Lemma~\ref{cv-origin-intersection-projection-count-disjoint}, $[F_1E_0:F_0E_0]=[F_1:F_0]$, and by (a) and (b)
this is $[\Phi(S):\Phi]$.

\emph{The ratios of $\xi$ and of $y$.} Let $A=R'/P'$, with fraction field $F_A$, and let $\alpha_{k',i'}$ be the class of
$x'_{k',i'}$. The homomorphism $R'\to\Phi'$ has kernel $P'$, so it induces a field embedding $F_A\to\Phi'$ with
$\alpha\mapsto\xi$. The homomorphism $R'\to\Omega[t_{k'}:k'\in K']$, $x'_{k',i'}\mapsto y_{k',i'}t_{k'}$, also has kernel
$P'$ (Definition~\ref{lit-multiprojective}(4)), so it induces a field embedding of $F_A$ into the fraction field $L$ of
$\Omega[t]$ with $\alpha_{k',i'}\mapsto y_{k',i'}t_{k'}$. In particular $\alpha_{k',b}$, $\xi_{k',b}$ and $y_{k',b}$
vanish together. The first embedding maps $\mathbf Q(\alpha)$ and $\mathbf Q(\alpha\circ v)$ onto $\mathbf Q(\xi)$ and
$\mathbf Q(\xi\circ v)$. The second maps them onto the subfields of $L$ generated by the
$y_{k',a}t_{k'}/y_{k',b}t_{k'}=y_{k',a}/y_{k',b}$, which are the images of $\mathbf Q(y)$ and $\mathbf Q(y\circ v)$ under
$\Omega\to L$. So $[\mathbf Q(\xi):\mathbf Q(\xi\circ v)]=[\mathbf Q(\alpha):\mathbf Q(\alpha\circ v)]
=[\mathbf Q(y):\mathbf Q(y\circ v)]$, which is finite by Definition~\ref{cv-origin-intersection-full}(2). With the first
two steps, this proves the lemma.
\end{proof}
\leanref{GroundTheory/HeightsAndTheta/Heights/OrbitLimits/PCoordFibreDeg.lean}{GroundTheory.Heights₂.Heights.Targets.CoordFibreDeg.proof}
\leanref{GroundTheory/HeightsAndTheta/Heights/OrbitLimits/PGenFibreChartDeg.lean}{GroundTheory.Heights₂.Heights.Targets.GenFibreChartDeg.proof}
\leanref{GroundTheory/HeightsAndTheta/Heights/OrbitLimits/PCoordRatioDeg.lean}{GroundTheory.Heights₂.Heights.Targets.CoordRatioDeg.proof}
% Lean: Closed.Targets.CoordFibreDeg = Targets.CoordFibreDeg.proof: CharZero of G.Phi', FibreMultiHomog.proof,
% ChartRankAdjoin.proof (the chart), GenFibreChartDeg.proof with ChartScalingIndep.full and AdjoinIndepRelfinrank.full
% (hPhi1 is (a), hPhiS is (b), hy and hindep are (c), isMultiHomogIdeal_ker_rat for the kernel over Q,
% FinrankAdjoinRange.proof for [Phi(S):Phi] as a relative degree inside Phi'), and CoordRatioDeg.proof (the embeddings
% f1 into Phi' and f2 into FractionRing (MvPolynomial K' Omega) of FractionRing (MRing ι' / P'), ratioField_map). The
% Lean's ratio fields blockFnField also contain the quotients by 0, which are 0. Size: PCoordFibreDeg.lean,
% PGenFibreChartDeg.lean, PCoordRatioDeg.lean, PChartRankAdjoin.lean, PChartScalingIndep.lean,
% PAdjoinIndepRelfinrank.lean and PFinrankAdjoinRange.lean, about 690 lines.

\begin{lemma}\label{cv-origin-intersection-projection}
Every arithmetic intersection system (Definition~\ref{lit-multiprojective}(3)) satisfies the projection formula of
Definition~\ref{cv-origin-intersection-full}(2).
\end{lemma}
\leanref{ReciprocityTheory/ReductiveGroups/ArithmeticGeometry/ArithmeticIntersection/MixedFieldContraction/ContractMixedFieldStatement.lean}{ReciprocityTheory.ChartRestriction.G.AI.I4.PFCoordStatement}
\leanref{Development/AutomorphicSpectrum/Proofs/MetricProofs/DegreeAndCoordinatesOfResolutionStatement.lean}{GlobalTheory.T.degree_and_coordinates_of_resolution}
\leanref{Development/AutomorphicSpectrum/Proofs/MetricProofs/DegreeAndCoordinatesOfResolution.lean}{GlobalTheory.CompanionTheory.Ar.Pub.genDeg_and_pfCoord}
\begin{proof}
Let $(\kappa,v)$, $\Omega$, $y$, $X'$, $X$ and the $\bar L_j$ be as in Definition~\ref{cv-origin-intersection-full}(2),
with $P'$ the kernel of $y$ and $Q=v_*^{-1}(P')$ that of $y\circ v$, both of dimension $d$. By
Lemma~\ref{cv-origin-intersection-projection-fibre} with $r=0$, $X'$ has a generic fibre $V$ over $Q$, a projective
variety of dimension $0$ over the fraction field $\Phi$ of $R/Q$. As $V$ is relevant, there are indices $i_{k'}$ with
$x'_{k',i_{k'}}\notin V$. The intersection number $\deg(V)$ of the zero-dimensional variety $V$ is the dimension over
$\Phi$ of its coordinate ring in this chart. By Lemma~\ref{cv-origin-intersection-projection-generic},
$(\bar L'_1\cdots\bar L'_d\mid X')=\deg(V)\,(\bar L_1\cdots\bar L_d\mid X)$, and by
Lemma~\ref{cv-origin-intersection-projection-count}, $\deg(V)=[\mathbf Q(y):\mathbf Q(y\circ v)]$.
\end{proof}
\leanref{Development/AutomorphicSpectrum/Proofs/MetricProofs/DegreeAndCoordinatesOfResolution.lean}{GlobalTheory.CompanionTheory.Ar.Pub.pfCoord_of_genDeg}
\leanref{GroundTheory/HeightsAndTheta/Heights/OrbitLimits/PFCoord.lean}{GroundTheory.Heights₂.Heights.Targets.PFCoord.of_found}
% Lean: Pub.genDeg_and_pfCoord (= GlobalTheory.T.degree_and_coordinates_of_resolution, its second clause), under
% ResExistsStatement (Lemma cv-origin-intersection-resolution-exists): genDeg_of_resExists gives GenDegStatement (the
% factor is J (fibreIdeal kappa G) 0 for a field intersection family J on the generic fibre G, from PtsIntegrable,
% IntersectionWithSection.proof and ArithmeticIntersectionCycle), and pfCoord_of_genDeg (Targets.PFCoord.of_found
% with Targets.ExistsUnique.full, genFibreExists, Giv.fibreDim and CoordFibreDeg) replaces it by the relfinrank of
% the block function fields. PFCoord.of_found: hGF gives G, hR4 the field intersection family J on G.Phi, hFD the
% dimension 0, hJ.1 the value ptDeg at the chart i, hGD and hCF. Size: about 24,500 lines, in Lemmas
% cv-origin-intersection-projection-fibre, -generic and -count.

\subsection{Contractions and full intersection systems}\label{subsection-cv-contractions}

\begin{lemma}\label{cv-origin-intersection-contraction-image}
Let $(\kappa,v)$ be a block map (Definition~\ref{cv-origin-intersection-full}(1)), let $X'$ be an arithmetic
subvariety of dimension $d'$ of $\prod_{k'\in K'}\mathbf P^{I'_{k'}}_{\mathbf Z}$, defined by the prime $P'$, and put
$Q=v_*^{-1}(P')$.
\begin{enumerate}
\item The prime $Q$ is an arithmetic subvariety of $\prod_{k\in K}\mathbf P^{I_k}_{\mathbf Z}$ of some dimension
  $d\le d'$.
\item Let $W\supset P'$ be a prime ideal of $R'$, and let $t\in R$ with $v_*(t)\in W$ and $t\notin Q$. Then
  $\dim R/v_*^{-1}(W)+1\le\dim R/Q$.
\end{enumerate}
\end{lemma}
\leanref{GroundTheory/HeightsAndTheta/Heights/OrbitLimits/ImageArithClosed.lean}{GroundTheory.Heights₂.Heights.imageArith}
\leanref{ReciprocityTheory/ReductiveGroups/ArithmeticGeometry/ArithmeticIntersection/MixedFieldContraction/ContractMixedFieldStatement.lean}{ReciprocityTheory.ChartRestriction.G.AI.I4.ImageArithStatement}
\leanref{GroundTheory/HeightsAndTheta/Heights/Intersections/TableDegreeThreeBase.lean}{GroundTheory.Heights₂.Heights.Inter.krullDim_quot_add_one_le}
\begin{proof}
(1) The ideal $Q$ is prime, and $Q\cap\mathbf Z=P'\cap\mathbf Z=0$. The map $v_*$ takes a monomial of multidegree
$a\in\mathbf Z^K$ to one of multidegree $a'$, where $a'_{\kappa(k)}=a_k$ and $a'_{k'}=0$ for $k'\notin\kappa(K)$. As
$\kappa$ is injective, $v_*$ takes the multihomogeneous components of an element of $Q$ to those of its image, which lie
in $P'$; so $Q$ is multihomogeneous. For each $k$ some $x'_{\kappa(k),i'}$ is not in $P'$, and $i'=v(k,i)$ for some
$i\in I_k$, so $x_{k,i}\notin Q$ and $Q$ is relevant.

For the dimension, let $L$ be the fraction field of $R'/P'$, and let $A\subset B\subset L$ be the $\mathbf Q$-subalgebras
generated by the images of the $x'_{v(k,i)}$ and of all the $x'_{k',i'}$. The rings $R/Q$ and $R'/P'$ are domains of
finite type over $\mathbf Z$ of characteristic zero, with the fraction fields of $A$ and $B$, so
$\dim R/Q=\operatorname{trdeg}_{\mathbf Q}A+1$ and $\dim R'/P'=\operatorname{trdeg}_{\mathbf Q}B+1$, by the dimension
formula as in the proof of Lemma~\ref{cv-origin-weight-birational} \cite[Tag 02JU]{Sta}. Let $K''$ be the set of
blocks $k'\notin\kappa(K)$, and for $k'\in K''$ choose $i'_{k'}$ with $x'_{k',i'_{k'}}\notin P'$. The images
$\xi_{k'}$ of these elements in $L$ are algebraically independent over $A$. Indeed, let $\sum_ec_e\xi^e=0$ with
$c_e\in A$, the sum over finitely many $e\in\mathbf N^{K''}$. After multiplication by a positive integer, $c_e$ is the
image of $v_*(g_e)$ for some $g_e\in R$, and then
$H=\sum_ev_*(g_e)\prod_{k'\in K''}x'^{\,e_{k'}}_{k',i'_{k'}}$ lies in $P'$. Grade $R'$ by the degrees in the blocks of
$K''$. The ideal $P'$ is homogeneous for this grading and $v_*(g_e)$ has degree $0$, so each term of $H$ lies in $P'$.
As $P'$ is prime and contains no $x'_{k',i'_{k'}}$, $v_*(g_e)\in P'$, that is $c_e=0$. Hence
$\operatorname{trdeg}A+|K'|-|K|\le\operatorname{trdeg}B$. The same argument with no blocks in the source, applied to $Q$,
gives $|K|\le\operatorname{trdeg}A$. So $d=\operatorname{trdeg}A+1-|K|$ is at least $1$, $\dim R/Q=d+|K|$, and
$d+|K|\le\operatorname{trdeg}B+1-|K'|+|K|=d'+|K|$.

(2) The ring $R/v_*^{-1}(W)$ is a quotient of $R/(Q+(t))$, since $Q\subset v_*^{-1}(W)$ and $t\in v_*^{-1}(W)$. The
image of $t$ in the domain $R/Q$ is nonzero. So a chain of primes of $R/Q$ that contain $t$ starts at a nonzero prime,
and it extends by the prime $0$. Hence $\dim R/(Q+(t))+1\le\dim R/Q$.
\end{proof}
\leanref{GroundTheory/HeightsAndTheta/Heights/OrbitLimits/ImageArith.lean}{GroundTheory.Heights₂.Heights.ImageArith.of_found}
\leanref{GroundTheory/HeightsAndTheta/Heights/OrbitLimits/PBlockTrdeg.lean}{GroundTheory.Heights₂.Heights.BlockTrdeg.proof}
% Lean: ImageArith.of_found: ContractRelevant.isMultiProjPrime_comap_rename (multihomogeneous and relevant),
% TrdegAdjoin.ringKrullDim_quotient_ker (dim = trdeg + 1) twice, and BlockTrdeg.proof (algebraicIndependent_forgotten
% with the grading forgetWeight), applied to (kappa, v) and to the empty block map into K. (2) is
% krullDim_quot_add_one_le (ringKrullDim_quotient_succ_le_of_nonZeroDivisor).

\begin{lemma}\label{cv-origin-intersection-contraction-ratios}
Let $K$, $(I_k)_{k\in K}$ and $R$ be as in Definition~\ref{lit-multiprojective}, let $Y$ be an arithmetic subvariety of
dimension $d$ defined by the prime $Q$, let $c=(c_k)\in\prod_kI_k$, and let $p_1,\dots,p_d\in\bigsqcup_kI_k$; for
$p\in I_k$ write $\bar p=(k,c_k)$. Then there is a nonzero $f\in\mathbf Z[t_1,\dots,t_d]$ with
\[
 f\big(z_{p_1}/z_{\bar p_1},\dots,z_{p_d}/z_{\bar p_d}\big)=0
\]
for every $z\in\hat Y(\mathbf C)$ with $z_{k,c_k}\neq0$ for all $k$.
\end{lemma}
\leanref{GroundTheory/HeightsAndTheta/Heights/Intersections/StRatioAlgDep.lean}{GroundTheory.Heights₂.Heights.Inter.RatioAlgDep.Statement}
\begin{proof}
If some $x_{k,c_k}$ lies in $Q$, then every $z\in\hat Y(\mathbf C)$ has $z_{k,c_k}=0$, and $f=1$ will do. Otherwise $Q$
is the kernel of a point $y\in\Omega^{\sqcup_kI_k}$ with function field $\Omega$
(Definition~\ref{lit-multiprojective}(4)): take for $y_{k,i}$ the class of $x_{k,i}/x_{k,c_k}$ in the fraction field of
$R/Q$, and for $\Omega$ the field these classes generate. Then $y_{k,c_k}=1$, and the fraction field of $R/Q$ is
$\Omega(y_{k,c_k}T_k:k\in K)$, of transcendence degree $\operatorname{trdeg}_{\mathbf Q}\Omega+|K|$, so
$d+|K|=\dim R/Q=\operatorname{trdeg}_{\mathbf Q}\Omega+|K|+1$ \cite[Tag 02JU]{Sta}. Hence $\Omega$ has transcendence
degree $d-1$, and the $d$ elements $y_{p_j}/y_{\bar p_j}$ of $\Omega$ are algebraically dependent: there is a nonzero
$f=\sum_mf_mt^m\in\mathbf Z[t_1,\dots,t_d]$ that vanishes at them. Let $\delta$ be the total degree of $f$, and put
\[
 G=\sum_mf_m\prod_{j=1}^dx_{p_j}^{m_j}x_{\bar p_j}^{\delta-m_j},
\]
a multihomogeneous element of $R$. Its image under $x_{k,i}\mapsto y_{k,i}T_k$ is
$\big(\prod_jy_{\bar p_j}T_{k(j)}\big)^\delta f\big(y_{p_1}/y_{\bar p_1},\dots,y_{p_d}/y_{\bar p_d}\big)=0$, where
$p_j\in I_{k(j)}$; so $G\in Q$. For $z\in\hat Y(\mathbf C)$ with all $z_{k,c_k}\neq0$ this gives
$0=G(z)=\big(\prod_jz_{\bar p_j}\big)^\delta f\big(z_{p_1}/z_{\bar p_1},\dots,z_{p_d}/z_{\bar p_d}\big)$, and the first
factor is nonzero.
\end{proof}
\leanref{GroundTheory/HeightsAndTheta/Heights/Intersections/PRatioAlgDep.lean}{GroundTheory.Heights₂.Heights.Inter.RatioAlgDep.proof}
% Lean: RatioAlgDep.proof: the point y with kernel Q and Birational y is HilbertSamuel.BirationalModelExists.proof;
% ConeTrdeg.proof and IntDomainDim.proof give dim R/Q = trdeg Omega + |K| + 1; algebraicIndependent_iff gives f;
% RatioAlgebraicDependence.homog and aeval_homog are G and its evaluation.

\begin{lemma}\label{cv-origin-intersection-contraction-rank}
Let $K$, $(I_k)_{k\in K}$ and $R$ be as in Definition~\ref{lit-multiprojective}, let $Y$ be an arithmetic subvariety of
dimension $d$, let $n\ge0$ and $y\in\mathbf C^n$, and let $F$ be a holomorphic map from a neighbourhood of $y$ to
$\prod_k\mathbf C^{I_k}$ with $F(y')\in\hat Y(\mathbf C)$ for all $y'$ near $y$. Let $E$ be the set of $x\in\mathbf C^n$
for which there is $c\in\mathbf C^K$ with $dF(y)(x)_{k,i}=c_kF(y)_{k,i}$ for all $(k,i)$; it is a linear subspace.
\begin{enumerate}
\item If $d\le n$, then $E\neq0$.
\item $\dim E+d>n$.
\end{enumerate}
\end{lemma}
\leanref{GroundTheory/HeightsAndTheta/Heights/Intersections/ExMapAlt.lean}{GroundTheory.Heights₂.Heights.Inter.pointFormRank_holds.Statement}
\leanref{GroundTheory/HeightsAndTheta/Heights/Intersections/PPFIntZero.lean}{GroundTheory.Heights₂.Heights.Inter.eulerDirs_finrank}
\begin{proof}
(1) As $F(y)\in\hat Y(\mathbf C)$, there are indices $c_k\in I_k$ with $F(y)_{k,c_k}\neq0$; for $p\in I_k$ write
$\bar p=(k,c_k)$. Near $y$ the ratios $r_p=F_p/F_{\bar p}$ are holomorphic; let $\ell_p$ be the derivative of $r_p$ at
$y$. Differentiating $F_p=r_pF_{\bar p}$ at $y$ gives
$dF(y)(x)_p=r_p(y)\,dF(y)(x)_{\bar p}+F(y)_{\bar p}\,\ell_p(x)$. So if $\ell_p(x)=0$ for every $p$, then $x\in E$, with
$c_k=dF(y)(x)_{k,c_k}/F(y)_{k,c_k}$. Suppose that $E=0$. Then the linear forms $\ell_p$ have no common zero other than
$0$, and as $d\le n$ there are $p_1,\dots,p_d$ for which $x\mapsto(\ell_{p_j}(x))_j$ maps $\mathbf C^n$ onto
$\mathbf C^d$. A holomorphic map whose derivative at $y$ is onto takes every neighbourhood of $y$ to a neighbourhood of
its value at $y$, by the implicit function theorem; this applies to $y'\mapsto(r_{p_j}(y'))_j$. The polynomial $f$ of
Lemma~\ref{cv-origin-intersection-contraction-ratios}, for $c$ and the $p_j$, vanishes at $(r_{p_j}(y'))_j$ for $y'$ near
$y$, so on a nonempty open subset of $\mathbf C^d$. Hence $f=0$, a contradiction.

(2) Suppose that $\dim E+d\le n$. Let $U$ be a complement of $E$ in $\mathbf C^n$, of dimension $u\ge d$, and let
$M:\mathbf C^u\to U$ be an isomorphism. By (1), applied to $y'\mapsto F(y+My')$ at $0$, there is a nonzero
$x\in\mathbf C^u$ with $Mx\in E$. Then $Mx\in E\cap U=0$, a contradiction.
\end{proof}
\leanref{GroundTheory/HeightsAndTheta/Heights/Intersections/EqZeroOfEventuallyEval.lean}{GroundTheory.Heights₂.Heights.Inter.pointFormRank_holds.proof}
% Lean: (1) is pointFormRank_holds.proof with exists_surj_coords (the choice of the p_j),
% HasStrictFDerivAt.map_nhds_eq_of_surj and eq_zero_of_eventually_eval; (2) is eulerDirs_finrank, with E = eulerDirs.

\begin{lemma}\label{cv-origin-intersection-contraction-euler}
Let $K$ and $(I_k)_{k\in K}$ be as in Definition~\ref{lit-multiprojective}, let $\bar A=(a,h)$ be a smooth metrized line
bundle, let $n\ge0$ and $y\in\mathbf C^n$, and let $F$ be a holomorphic map from a neighbourhood of $y$ to
$\prod_k\mathbf C^{I_k}$ with $F(y)\in\hat C$. Let $x\in\mathbf C^n$ and $c\in\mathbf C^K$ with
$dF(y)(x)_{k,i}=c_kF(y)_{k,i}$ for all $(k,i)$. Then the complex Hessian $H$ of $\log h\circ F$ at $y$
(Definition~\ref{cv-period-monge-ampere}(1)) satisfies $H\bar x=0$.
\end{lemma}
\leanref{GroundTheory/HeightsAndTheta/Heights/Intersections/ExMap.lean}{GroundTheory.Heights₂.Heights.Inter.pointFormEuler_holds.Statement}
\begin{proof}
This is a self-contained computation. Put $\varphi=\log h$, and for $e\in\mathbf C^K$ and $w\in\prod_k\mathbf C^{I_k}$
write $e\cdot w=(e_kw_{k,i})_{k,i}$. For $z\in\hat C$ and real $s$ near $0$,
$\varphi(z+s\,e\cdot z)=\sum_ka_k\log|1+se_k|+\varphi(z)$, so the derivative $d\varphi(z)(e\cdot z)$ does not depend on
$z$. Differentiating it in a direction $w$ gives $D^2\varphi(z)(w,e\cdot z)=-d\varphi(z)(e\cdot w)$. By the chain rule,
for $u,u'\in\mathbf C^n$,
\[
 D^2(\varphi\circ F)(y)(u,u')=D^2\varphi(F(y))\big(dF(y)u,dF(y)u'\big)+d\varphi(F(y))\big(D^2F(y)(u,u')\big).
\]
For the $l$-th unit vector $e_l$, the $l$-th entry of $H\bar x$ is
$\frac14\big(D^2(e_l,x)+D^2(ie_l,ix)+i\,D^2(e_l,ix)-i\,D^2(ie_l,x)\big)$, with $D^2=D^2(\varphi\circ F)(y)$.
Substitute $dF(y)(x)=c\cdot F(y)$, $dF(y)(ix)=(ic)\cdot F(y)$ and $D^2F(y)(u,ix)=i\,D^2F(y)(u,x)$, and the identity above
with $e=c$ and with $e=ic$. The eight terms cancel.
\end{proof}
\leanref{GroundTheory/HeightsAndTheta/Heights/Intersections/LeviOfRealSmulRight.lean}{GroundTheory.Heights₂.Heights.Inter.pointFormEuler_holds.proof}
% Lean: pointFormEuler_holds.proof: fderiv_logh_bsmul and hess_logh_bsmul (the identity), hess_comp_holo (the chain
% rule), cxHessian_mulVec_star (the entry of H xbar as leviOf), then ring.

\begin{lemma}\label{cv-origin-intersection-contraction-discriminant}
Let $D\ge0$, let $A_1,\dots,A_D$ be complex $D\times D$ matrices, let $V\subset\mathbf C^D$ be a subspace and
$S\subset\{1,\dots,D\}$ with $\dim V+|S|>D$, and suppose that $A_kx=0$ for all $k\in S$ and $x\in V$. Then
\[
 \sum_{T\subset\{1,\dots,D\}}(-1)^{D-|T|}\det\Big(\sum_{k\in T}A_k\Big)=0 .
\]
\end{lemma}
\leanref{GroundTheory/HeightsAndTheta/Heights/Intersections/PFIntZeroMixed.lean}{GroundTheory.Heights₂.Heights.Inter.mixedDisc_eq_zero_of_ker}
\begin{proof}
Let $P$ be an invertible matrix whose first $\dim V$ columns span $V$. Replacing each $A_k$ by $B_k=A_kP$ multiplies the
sum by $\det P$, and for $k\in S$ the columns of $B_k$ in the set $C$ of the first $\dim V$ positions vanish. Expand
\[
 \det\Big(\sum_{k\in T}B_k\Big)=\sum_\sigma\sum_g\operatorname{sgn}(\sigma)\prod_{l=1}^D(B_{g(l)})_{\sigma(l),l},
\]
the sum over the permutations $\sigma$ of $\{1,\dots,D\}$ and the maps $g:\{1,\dots,D\}\to T$, and exchange the sums. If
$g$ misses an index, the terms for $(\sigma,g)$ come with the factor $\sum_T(-1)^{D-|T|}$ over the sets $T$ containing
the image of $g$, which is $0$. If $g$ is a bijection, $g^{-1}(S)$ has $|S|$ elements and meets $C$, since
$|C|+|S|>D$; for $l$ in the intersection, $(B_{g(l)})_{\sigma(l),l}=0$.
\end{proof}
\leanref{GroundTheory/HeightsAndTheta/Heights/Intersections/PFIntZeroMixed.lean}{GroundTheory.Heights₂.Heights.Inter.mixedDisc_eq_zero_of_cols}
\leanref{GroundTheory/HeightsAndTheta/Heights/Intersections/PFIntZeroMixed.lean}{GroundTheory.Heights₂.Heights.Inter.mixedDisc_mul_right}

\begin{lemma}\label{cv-origin-intersection-contraction-forms}
Let $(\kappa,v)$ be a block map (Definition~\ref{cv-origin-intersection-full}(1)), let $X'$ be an arithmetic subvariety
of dimension $n+1$ of $\prod_{k'}\mathbf P^{I'_{k'}}_{\mathbf Z}$ defined by the prime $P'$, such that $v_*^{-1}(P')$ is
an arithmetic subvariety of dimension $d$, and let $g$ be a real function on $\hat X'(\mathbf C)$. For smooth metrized
line bundles $\bar M_1,\dots,\bar M_n$ on $\prod_{k'}\mathbf P^{I'_{k'}}_{\mathbf Z}$, let $J(\bar M_1,\dots,\bar M_n)$
be the integral of $g$ over $\hat X'(\mathbf C)$ against $c_1(\bar M_1)\wedge\dots\wedge c_1(\bar M_n)$.
\begin{enumerate}
\item If $d\le n$ and every $\bar M_j$ is the pullback of a smooth metrized line bundle, then
  $J(\bar M_1,\dots,\bar M_n)=0$.
\item If $n=d+r$ with $r\ge0$ and $\bar M_1,\dots,\bar M_d$ are the pullbacks of smooth metrized line bundles, then
  $J(\bar M_1,\dots,\bar M_n)=0$.
\end{enumerate}
\end{lemma}
\leanref{GroundTheory/HeightsAndTheta/Heights/Intersections/ImportsFurtherTail.lean}{GroundTheory.Heights₂.Heights.Inter.pointForm_holds.Statement}
\leanref{GroundTheory/HeightsAndTheta/Heights/Intersections/StatementsPFIntZero.lean}{GroundTheory.Heights₂.Heights.Inter.PFIntZero.Statement}
\begin{proof}
The integral over $\hat X'(\mathbf C)$ is computed through a linear projection of the Segre coordinates to
$\mathbf C^{n+1}$ that does not vanish on $\hat X'(\mathbf C)$: it is an integral over $y\in\mathbf C^n$ of finite sums of
terms $g(\gamma(y))$ times the Monge--Amp\`ere density at $y$ (Definition~\ref{cv-period-monge-ampere}(2)) of the
functions $\log h'_j\circ\gamma$, where $h'_j$ is the function of $\bar M_j$ and $\gamma$ is a holomorphic map from a
neighbourhood of $y$ to $\hat X'(\mathbf C)$. So it suffices to show that each such density vanishes.

Fix $y$ and $\gamma$, let $Y$ be the arithmetic subvariety defined by $v_*^{-1}(P')$, and let $F=(\gamma_{v(p)})_p$, a
holomorphic map to $\prod_k\mathbf C^{I_k}$. For $y'$ near $y$, $F(y')\in\hat Y(\mathbf C)$: the elements of
$v_*^{-1}(P')$ vanish at $F(y')$, and each block $k$ of $F(y')$ is the block $\kappa(k)$ of $\gamma(y')$, which is
nonzero. If $\bar M_j$ is the pullback of $(a_j,h_j)$, then $\log h'_j\circ\gamma=\log h_j\circ F$. Let $E$ be the
subspace of Lemma~\ref{cv-origin-intersection-contraction-rank} for $Y$, $F$ and $y$. By
Lemma~\ref{cv-origin-intersection-contraction-euler}, for every pulled-back $\bar M_j$ the complex Hessian of
$\log h'_j\circ\gamma$ at $y$ kills $\bar x$ for every $x\in E$.

(1) By Lemma~\ref{cv-origin-intersection-contraction-rank}(1) there is a nonzero $x\in E$. All the complex Hessians kill
$\bar x$, so every sum of them is singular, and every determinant in the Monge--Amp\`ere density vanishes.

(2) By Lemma~\ref{cv-origin-intersection-contraction-rank}(2), the conjugate space $\bar E$ has
$\dim\bar E+d>n$. Lemma~\ref{cv-origin-intersection-contraction-discriminant}, with $V=\bar E$ and
$S=\{1,\dots,d\}$, shows that the density vanishes.
\end{proof}
\leanref{GroundTheory/HeightsAndTheta/Heights/Intersections/MixedDiscEqZeroOfMulVec.lean}{GroundTheory.Heights₂.Heights.Inter.pointForm_holds.proof}
\leanref{GroundTheory/HeightsAndTheta/Heights/OrbitLimits/Contract.lean}{GroundTheory.Heights₂.Heights.Targets.Contract.pulledBackForms_topDensity_vanish}
\leanref{GroundTheory/HeightsAndTheta/Heights/Intersections/PPFIntZero.lean}{GroundTheory.Heights₂.Heights.Inter.PFIntZero.proof}
% Lean: ptsIntegral is ptsIntegralAt for a good projection, the integral over y of the sum over the fibre of
% sheetDensity, which is g(gamma y) times maDensity of the potentials along a sheet gamma (IsSheet), or 0. (1) is
% pointForm_holds.proof with maDensity_eq_zero (pointFormRank_holds, pointFormEuler_holds and
% mixedDisc_eq_zero_of_mulVec with u = star x); (2) is PFIntZero.proof with PFIntZero.maDensity_eq_zero, which applies
% mixedDisc_eq_zero_of_ker to the conjugate Hessians and eulerDirs (mixedDisc_map_star), the same statement.

\begin{lemma}\label{cv-origin-intersection-contraction-field}
Let $F$ be a field, let $(\kappa,v)$ be a block map (Definition~\ref{cv-origin-intersection-full}(1)), and write
$v_*:F[x]\to F[x']$ for the homomorphism of the polynomial rings $F[x]=F[x_{k,i}]$ and $F[x']=F[x'_{k',i'}]$ with
$v_*(x_{k,i})=x'_{v(k,i)}$. Let $V'\subset F[x']$ be a relevant multihomogeneous prime with
$\dim F[x']/V'=e+|K'|$, defining a projective variety $Z'$ of dimension $e$ over $F$, such that $V=v_*^{-1}(V')$ has
$\dim F[x]/V=d_0+|K|$. For $a\in\mathbf Z^K$ let $a'\in\mathbf Z^{K'}$ be its extension by zero, with
$a'_{\kappa(k)}=a_k$ and $a'_{k'}=0$ for $k'\notin\kappa(K)$. Let $c_1,\dots,c_e\in\mathbf Z^{K'}$, more than $d_0$ of
which are extensions by zero. Then the intersection number $(\mathcal O(c_1)\cdots\mathcal O(c_e)\cdot Z')$
\cite[Tag 0BEP]{Sta} is $0$.
\end{lemma}
\leanref{ReciprocityTheory/ReductiveGroups/ArithmeticGeometry/ArithmeticIntersection/MixedFieldContraction/ContractMixedFieldStatement.lean}{ReciprocityTheory.ChartRestriction.G.AI.I4.ContractMixedFieldStatement}
\leanref{ReciprocityTheory/ReductiveGroups/ArithmeticGeometry/ArithmeticIntersection/MultihomogeneousIdeals/ExistsUniqueStatement.lean}{ReciprocityTheory.ChartRestriction.G.AI.I1.ProjContractStatement}
\begin{proof}
By induction on $e$. For $e=0$ there are no $c_j$, while more than $d_0\ge0$ of them are extensions by zero; so there is
nothing to prove. Let $e\ge1$. The intersection number is symmetric in the $c_j$, so we may assume that $c_e=a'$ for
some $a\in\mathbf Z^K$. As $V$ is relevant, there are indices $i_k$ with $x_{k,i_k}\notin V$. Put
$f=\prod_kx_{k,i_k}^{\max(a_k,0)}$ and $g=\prod_kx_{k,i_k}^{\max(-a_k,0)}$. They are multihomogeneous, they do not lie in
the prime $V$, and their multidegrees differ by $a$; so $v_*(f)/v_*(g)$ is a nonzero rational section of
$\mathcal O(a')$ on $Z'$. Intersection numbers satisfy the recursion through the divisor of a rational section: by
additivity \cite[Tag 0BER]{Sta} and by \cite[Tags 0BEU and 0BES]{Sta} for the effective Cartier divisors of $v_*(f)$
and $v_*(g)$ on $Z'$,
\[
 (\mathcal O(c_1)\cdots\mathcal O(c_e)\cdot Z')=\sum_Wm_W\,(\mathcal O(c_1)\cdots\mathcal O(c_{e-1})\cdot W),
\]
where $W$ runs over the components of these divisors and $m_W$ is the multiplicity of $W$ in the divisor of $v_*(f)$
minus that in the divisor of $v_*(g)$. Each $W$ is a projective variety of dimension $e-1$, defined by a relevant
multihomogeneous prime that contains $v_*(f)$ or $v_*(g)$. Its image $v_*^{-1}(W)$ is a relevant multihomogeneous prime
that contains $f$ or $g$, which are not in $V$; as in Lemma~\ref{cv-origin-intersection-contraction-image}(2),
$\dim F[x]/v_*^{-1}(W)=d_1+|K|$ with $d_1<d_0$. Among $c_1,\dots,c_{e-1}$ more than $d_0-1\ge d_1$ are extensions by
zero. By induction every term of the sum vanishes.
\end{proof}
\leanref{GroundTheory/HeightsAndTheta/Heights/OrbitLimits/ContractMixedField.lean}{GroundTheory.Heights₂.Heights.Targets.contractMixedField}
\leanref{GroundTheory/HeightsAndTheta/Heights/OrbitLimits/ContractMixedField.lean}{GroundTheory.Heights₂.Heights.Targets.ContractMixedField.intNum_eq_zero}
\leanref{GroundTheory/HeightsAndTheta/Heights/OrbitLimits/ProjContract.lean}{GroundTheory.Heights₂.Heights.Targets.ProjContract.full}
% Lean: contractMixedField: NumbersDVR.Numbers.agree_of identifies the field intersection family with intNum, and
% ContractMixedField.intNum_eq_zero is the induction: IntNumSymm.full (symmetry), ExistsForm.full (f and g),
% IntNumIsFieldIntersection.full (the recursion with ratDiv), ContractMixedField.component (DivCycleSupport.full for
% the components, exists_dim and krullDim_quot_add_one_le for d_1 < d_0). The case where every c_j is an extension by
% zero is ProjContractStatement, which the Lean proves separately as ProjContract.full
% (FieldNumbers.ProjectionFormulaContracted.full); it is the case dim f(Z) < d of Stacks Tag 0BET.

\begin{lemma}\label{cv-origin-intersection-contraction-vertical}
Let $(\kappa,v)$ be a block map (Definition~\ref{cv-origin-intersection-full}(1)), let $X'$ be an arithmetic subvariety
of dimension $n+1$ of $\prod_{k'}\mathbf P^{I'_{k'}}_{\mathbf Z}$ defined by the prime $P'$, such that $v_*^{-1}(P')$ is
an arithmetic subvariety of dimension $d$, let $s\in R'\setminus P'$ be multihomogeneous, and let $W$ be a relevant
minimal prime of $P'+(s)$ that is not horizontal, with $W\cap\mathbf Z=p\mathbf Z$. Let $c_1,\dots,c_n\in\mathbf Z^{K'}$,
at least $d$ of which are extensions by zero of elements of $\mathbf Z^K$. Then the degree $\deg_{c_1\cdots c_n}(W)$ of
Definition~\ref{lit-multiprojective}(3), the intersection number $(\mathcal O(c_1)\cdots\mathcal O(c_n)\cdot W_{\mathbf F_p})$
on the projective variety $W_{\mathbf F_p}$ over $\mathbf F_p$ defined by $W$, is $0$.
\end{lemma}
\leanref{GroundTheory/HeightsAndTheta/Heights/Intersections/IsMultiHomogFirstCoordinate.lean}{GroundTheory.Heights₂.Heights.Inter.ProjectivePointVertical.Statement}
\leanref{GroundTheory/HeightsAndTheta/Heights/Intersections/StatementsPFVert.lean}{GroundTheory.Heights₂.Heights.Inter.PFVert.Statement}
\begin{proof}
The prime $W$ contains $p$ and $\dim R'/W=n+|K'|$, as for the vertical components in
Definition~\ref{lit-multiprojective}(3); so $W_{\mathbf F_p}$, defined by the image of $W$ in $\mathbf F_p[x']$, has
dimension $n$. Its image in $\prod_k\mathbf P^{I_k}_{\mathbf F_p}$ is defined by the image of $v_*^{-1}(W)$ in
$\mathbf F_p[x]$, a relevant multihomogeneous prime with $\dim\mathbf F_p[x]/v_*^{-1}(W)=D_0+|K|$ for some $D_0\ge0$. As
$p\in v_*^{-1}(W)$ and $p\notin v_*^{-1}(P')$, Lemma~\ref{cv-origin-intersection-contraction-image}(2) with $t=p$ gives
$D_0+|K|+1\le d+|K|$, so $D_0<d$. At least $d>D_0$ of the $c_j$ are extensions by zero, and
Lemma~\ref{cv-origin-intersection-contraction-field}, for $F=\mathbf F_p$, $e=n$ and $V'$ the image of $W$, gives the
claim. When every $c_j$ is an extension by zero, this is also \cite[Tag 0BET]{Sta}.
\end{proof}
\leanref{GroundTheory/HeightsAndTheta/Heights/Intersections/IsMultiProjPrimeImageZ.lean}{GroundTheory.Heights₂.Heights.Inter.ProjectivePointVertical.proof}
\leanref{GroundTheory/HeightsAndTheta/Heights/Intersections/PPFVert.lean}{GroundTheory.Heights₂.Heights.Inter.PFVert.proof}
% Lean: PFVert.proof with geoNum_mixed_eq_zero, and ProjectivePointVertical.proof with PointVertical.geoNum_eq_zero for
% the case where all c_j are extensions by zero: divSuppArith for W, isMultiProjPrime_red and ringKrullDim_red for the
% reduction, imageVar_red for its image, exists_dim_F for D_0, krullDim_quot_add_one_le with t = p;
% geoNumStatement and Targets.ExistsUnique.full identify geoNum with the field intersection number, and then
% contractMixedField (Lemma cv-origin-intersection-contraction-field), or ProjContract.full in the second case.

\begin{lemma}\label{cv-origin-intersection-contraction-pulled}
Let $(\cdot\mid\cdot)$ be an arithmetic intersection system (Definition~\ref{lit-multiprojective}(3)), let $(\kappa,v)$
be a block map (Definition~\ref{cv-origin-intersection-full}(1)), let $X'$ be an arithmetic subvariety of dimension $n$
of $\prod_{k'}\mathbf P^{I'_{k'}}_{\mathbf Z}$ defined by the prime $P'$, such that $v_*^{-1}(P')$ is an arithmetic
subvariety of dimension $d<n$, and let $\bar L_1,\dots,\bar L_n$ be smooth metrized line bundles on
$\prod_k\mathbf P^{I_k}_{\mathbf Z}$ with pullbacks $\bar L'_1,\dots,\bar L'_n$. Then
$(\bar L'_1\cdots\bar L'_n\mid X')=0$.
\end{lemma}
\leanref{ReciprocityTheory/ReductiveGroups/ArithmeticGeometry/ArithmeticIntersection/MixedFieldContraction/ContractMixedFieldStatement.lean}{ReciprocityTheory.ChartRestriction.G.AI.I4.ContractStatement}
\begin{proof}
By induction on $n$; for $n=0$ there is no $d<0$. Let $n\ge0$, let $X'$ have dimension $n+1$, so that $d\le n$, and
suppose the claim holds in dimension $n$. The pullback of a smooth metrized line bundle is smooth. Put
$Q=v_*^{-1}(P')$.

First let $\bar B=(b,h_B)$ be a smooth metrized line bundle on $\prod_k\mathbf P^{I_k}_{\mathbf Z}$ with
$b\in\mathbf N^K$ and pullback $\bar B'=(b',h'_B)$; we show $(\bar L'_1\cdots\bar L'_n\bar B'\mid X')=0$. As $Q$ is
relevant, there are indices $i_k$ with $x_{k,i_k}\notin Q$. Then $t=\prod_kx_{k,i_k}^{b_k}$ is multihomogeneous of
multidegree $b$ and not in the prime $Q$, so $s=v_*(t)$ is multihomogeneous of multidegree $b'$ and not in $P'$. The
recursion of Definition~\ref{lit-multiprojective}(3) with $s$ writes $(\bar L'_1\cdots\bar L'_n\bar B'\mid X')$ as a sum
over the relevant minimal primes $W$ of $P'+(s)$ minus the integral of $\log(|s|/h'_B)$ against
$c_1(\bar L'_1)\wedge\dots\wedge c_1(\bar L'_n)$. The integral vanishes by
Lemma~\ref{cv-origin-intersection-contraction-forms}(1), as $d\le n$. A horizontal $W$ is an arithmetic subvariety of
dimension $n$, and $v_*^{-1}(W)$ is one of some dimension $d_W$ by
Lemma~\ref{cv-origin-intersection-contraction-image}(1). As $s\in W$ and $t\notin Q$,
Lemma~\ref{cv-origin-intersection-contraction-image}(2) gives $d_W+1\le d\le n$, so
$(\bar L'_1\cdots\bar L'_n\mid W)=0$ by induction. For a vertical $W$ the degree vanishes by
Lemma~\ref{cv-origin-intersection-contraction-vertical}, as all $n\ge d$ bundles are pullbacks.

In general, let $a$ be the multidegree of $\bar L_{n+1}$, let $\bar M$ be the Fubini--Study bundle of multidegree
$(\max(-a_k,0))_k$ on $\prod_k\mathbf P^{I_k}_{\mathbf Z}$ (Definition~\ref{cv-origin-analytic-system-pullback}(1)), and
let $\bar M'$ be its pullback. The pullback of $\bar L_{n+1}\otimes\bar M$ is $\bar L'_{n+1}\otimes\bar M'$, and
$\bar L_{n+1}\otimes\bar M$ and $\bar M$ have multidegrees in $\mathbf N^K$. By additivity and the first step,
\[
 (\bar L'_1\cdots\bar L'_{n+1}\mid X')=(\bar L'_1\cdots\bar L'_n\,(\bar L'_{n+1}\otimes\bar M')\mid X')
 -(\bar L'_1\cdots\bar L'_n\,\bar M'\mid X')=0.
\]
\end{proof}
\leanref{GroundTheory/HeightsAndTheta/Heights/Intersections/TableDegreeThreeBase.lean}{GroundTheory.Heights₂.Heights.Inter.contract_intersection_number}
\leanref{GroundTheory/HeightsAndTheta/Heights/Intersections/TableDegreeThreeBase.lean}{GroundTheory.Heights₂.Heights.Inter.contract_of}
% Lean: contract_intersection_number (= PushforwardProjectivePoint.proof, contract_of): key with t = prod X^b,
% s = rename v t (hsm, hBa'); arInt = 0 by pointForm_holds (Lemma -forms(1)); cycR = vertR by imageArith and
% krullDim_quot_add_one_le with the induction hypothesis; vertR = 0 by ProjectivePointVertical (Lemma -vertical);
% then M = sqMetric of the negative part of the multidegree, M' = pullHerm, IsPullHerm.tensor and additivity hI.2.1.

\begin{lemma}\label{cv-origin-intersection-contraction}
Every arithmetic intersection system (Definition~\ref{lit-multiprojective}(3)) has the contraction property of
Definition~\ref{cv-origin-intersection-full}(3).
\end{lemma}
\leanref{ReciprocityTheory/ReductiveGroups/ArithmeticGeometry/ArithmeticIntersection/MixedFieldContraction/ContractMixedFieldStatement.lean}{ReciprocityTheory.ChartRestriction.G.AI.I4.ContractMixedStatement}
\leanref{Development/AutomorphicSpectrum/Proofs/MetricProofs/ArithIntersectionSystemOfPoincareLelong.lean}{GlobalTheory.CompanionTheory.Ar.Pub.contractMixed_closed}
\begin{proof}
Let $(\cdot\mid\cdot)$ be an arithmetic intersection system and $(\kappa,v)$ a block map. We show more generally the
following claim, for every $n\ge0$. Let $X'$ be an arithmetic subvariety of dimension $n$ of
$\prod_{k'}\mathbf P^{I'_{k'}}_{\mathbf Z}$ defined by the prime $P'$, such that $Q=v_*^{-1}(P')$ is an arithmetic
subvariety of dimension $d_0$, let $d_0<q\le n$, and let $\bar M_1,\dots,\bar M_n$ be smooth metrized line bundles on
$\prod_{k'}\mathbf P^{I'_{k'}}_{\mathbf Z}$ such that $\bar M_1,\dots,\bar M_q$ are the pullbacks of smooth metrized line
bundles. Then $(\bar M_1\cdots\bar M_n\mid X')=0$. The contraction property is the case $n=d+r$, $q=d$ and
$(\bar M_j)=(\bar L'_1,\dots,\bar L'_d,\bar N_1,\dots,\bar N_r)$, since the pullbacks of smooth metrized line bundles
are smooth.

The proof is by induction on $n$. For $n=0$ there is no $q$ with $d_0<q\le0$. Let $X'$ have dimension $n+1$ and suppose
the claim holds in dimension $n$. If $q=n+1$, every $\bar M_j$ is a pullback, and
Lemma~\ref{cv-origin-intersection-contraction-pulled} gives the claim. Let $q\le n$.

First let $\bar B=(b,h_B)$ be a smooth metrized line bundle on $\prod_{k'}\mathbf P^{I'_{k'}}_{\mathbf Z}$ with
$b\in\mathbf N^{K'}$; we show $(\bar M_1\cdots\bar M_n\bar B\mid X')=0$. As $P'$ is relevant, there are indices $i'_{k'}$
with $x'_{k',i'_{k'}}\notin P'$, and $s=\prod_{k'}x'^{\,b_{k'}}_{k',i'_{k'}}$ is multihomogeneous of multidegree $b$ and not
in $P'$. The recursion of Definition~\ref{lit-multiprojective}(3) with $s$ writes the number as a sum over the relevant
minimal primes $W$ of $P'+(s)$ minus the integral of $\log(|s|/h_B)$ against
$c_1(\bar M_1)\wedge\dots\wedge c_1(\bar M_n)$. As $d_0<q$, the bundles $\bar M_1,\dots,\bar M_{d_0}$ are pullbacks, so
the integral vanishes by Lemma~\ref{cv-origin-intersection-contraction-forms}(2), with $n=d_0+(n-d_0)$, and the degree of
each vertical $W$ vanishes by Lemma~\ref{cv-origin-intersection-contraction-vertical}. A horizontal $W$ is an arithmetic
subvariety of dimension $n$, and $v_*^{-1}(W)$ is one of some dimension $d_W$ by
Lemma~\ref{cv-origin-intersection-contraction-image}(1). As $R/v_*^{-1}(W)$ is a quotient of $R/Q$, $d_W\le d_0<q\le n$,
and $(\bar M_1\cdots\bar M_n\mid W)=0$ by induction.

In general, let $a$ be the multidegree of $\bar M_{n+1}$, and let $\bar M$ be the Fubini--Study bundle of multidegree
$(\max(-a_{k'},0))_{k'}$ on $\prod_{k'}\mathbf P^{I'_{k'}}_{\mathbf Z}$
(Definition~\ref{cv-origin-analytic-system-pullback}(1)). Both $\bar M_{n+1}\otimes\bar M$ and $\bar M$ have multidegrees
in $\mathbf N^{K'}$, so by additivity and the first step
$(\bar M_1\cdots\bar M_{n+1}\mid X')=(\bar M_1\cdots\bar M_n(\bar M_{n+1}\otimes\bar M)\mid X')
-(\bar M_1\cdots\bar M_n\bar M\mid X')=0$.
\end{proof}
\leanref{GroundTheory/HeightsAndTheta/Heights/OrbitLimits/ContractMixed.lean}{GroundTheory.Heights₂.Heights.Targets.ContractMixed.aux}
\leanref{GroundTheory/HeightsAndTheta/Heights/OrbitLimits/ContractMixed.lean}{GroundTheory.Heights₂.Heights.Targets.ContractMixed.of_contract}
\leanref{GroundTheory/HeightsAndTheta/Heights/OrbitLimits/ContractMixed.lean}{GroundTheory.Heights₂.Heights.Targets.ContractMixed.of_found}
% Lean: Pub.contractMixed_closed, proved outright: Targets.ContractMixed.of_found (imageArith,
% Contract.pulledBackForms_topDensity_vanish, ProjectivePointVertical.proof, PFIntZero.proof and PFVert.proof).
% Its proof, ContractMixed.aux with PrefixPull for the first q bundles: induction on the number of bundles; when all are
% pulled back, ContractStatement (contract_of, Lemma -pulled); otherwise the last bundle is A' (x) M minus M with M the
% sqMetric of the negative part, and the recursion with a monomial section has arInt = 0 (arInt_eq_zero, PFIntZero),
% vertical terms geoNum = 0 (geoNum_eq_zero, PFVert) and horizontal terms 0 by induction (imageArith and
% ringKrullDim_le_of_surjective). of_contract is the case n = d + r, q = d (IsPullHerm.isSmoothHerm).

\begin{lemma}\label{cv-origin-intersection-exists}
There is a full arithmetic intersection system (Definition~\ref{cv-origin-intersection-full}).
\end{lemma}
\leanref{ReciprocityTheory/ReductiveGroups/ArithmeticGeometry/BirationalGeometry/InterNamed.lean}{ReciprocityTheory.ChartRestriction.G.Geo.ArithIntersectionSystemStatement}
\leanref{Development/AutomorphicSpectrum/Proofs/MetricProofs/ArithIntersectionSystemOfPoincareLelongStatement.lean}{GlobalTheory.T.arith_intersection_system_of_poincare_lelong}
\leanref{Development/AutomorphicSpectrum/Proofs/MetricProofs/ArithIntersectionSystemOfPoincareLelong.lean}{GlobalTheory.CompanionTheory.Ar.Pub.arithIntersectionSystem_of_poincareLelong_resK}
\leanref{GroundTheory/HeightsAndTheta/Heights/OrbitLimits/ArithSystem.lean}{GroundTheory.Heights₂.Heights.Targets.ArithIntersectionSystem.of_found}
\leanref{Development/WeilBound/CoreStatements/HodgeTatePoints/AhSlim.lean}{PositivityTheory.HtPt.Ah.ais}
\begin{proof}
By Lemma~\ref{cv-origin-intersection-family} there is an arithmetic intersection system $(\cdot\mid\cdot)$. It is
additive in each of its arguments (Definition~\ref{lit-multiprojective}(3)). By
Lemma~\ref{cv-origin-intersection-projection} it satisfies the projection formula, and by
Lemma~\ref{cv-origin-intersection-contraction} it has the contraction property. So it is full.
\end{proof}
% Lean: PositivityTheory.HtPt.Ah.ais: resK from ResolutionCharZero and PrincipalisationCharZero (Lemma
% cv-resolution-coordinates), Tw.Main.poincareLelong_of_inv (Lemma cv-origin-intersection-poincare-lelong-all) from
% PoincareLelongInv (Theorem cv-origin-intersection-poincare-lelong), then
% Pub.arithIntersectionSystem_of_poincareLelong_resK: resolution_exists_of_resolution_statement (Lemma
% cv-origin-intersection-resolution-exists) and Targets.ArithIntersectionSystem.of_found with
% intersection_exists_of_poincare_lelong (Lemma cv-origin-intersection-family), degree_and_coordinates_of_resolution
% (Lemma cv-origin-intersection-projection) and contractMixed_closed (Lemma cv-origin-intersection-contraction).
% of_found takes IsMultilinFam from multilin_of_isFamily: the Lean's IsIntersectionFamily asks for additivity in the
% last argument only, and the symmetry moves argument j to the last place. Size of the whole construction, counted
% once: about 57,000 lines.

\subsection{Effective sections and ampleness}\label{subsection-cv-effective-sections-ampleness}

\begin{lemma}\label{cv-origin-weight-nef-generic}
Let $K$ and $(I_k)_{k\in K}$ be as in Definition~\ref{lit-multiprojective}, let $D\ge0$, let $X$ be an
arithmetic subvariety of dimension $D+1$ of $\prod_{k\in K}\mathbf P^{I_k}_{\mathbf Z}$, and let
$\bar L_j=(a_j,h_j)$, $j=1,\dots,D$, be smooth metrized line bundles. Then the integral over the complex points
of $X$ of $c_1(\bar L_1)\wedge\dots\wedge c_1(\bar L_D)$ is the intersection number
$(\mathcal O(a_1)\cdots\mathcal O(a_D)\cdot X_{\mathbf Q})$ on the generic fibre $X_{\mathbf Q}$ of $X$.
\end{lemma}
\leanref{Development/AutomorphicSpectrum/Proofs/MetricProofs/GenericIntersectionEqualsPointsIntegralStatement.lean}{GlobalTheory.T.generic_intersection_equals_points_integral}
\leanref{Development/WeilBound/CoreStatements/PluriharmonicCharts/X_deg_generic.lean}{PositivityTheory.Xgp.t_deg_generic}
\begin{proof}
Let $R$ be as in Definition~\ref{lit-multiprojective}, and let $P\subset R$ be the prime defining $X$. By
Lemma~\ref{cv-origin-intersection-exists} there is a full arithmetic intersection system $(\cdot\mid\cdot)$; only its
recursion (Definition~\ref{lit-multiprojective}(3)) is used below.
% Lean: DegreeStatement is a component of PB.Inputs; Giv.inputs_of_use supplies it as Giv.degree_of_system hII
% from ArithIntersectionSystemStatement (Lemma cv-origin-intersection-exists), of which only IsIntersectionSystem
% (hII.1) is used.
By Lemma~\ref{cv-origin-weight-nef-generic-fibre} with $p=2$, the integral of the statement is
$\sum_Vm_V\,(\mathcal O(a_1)\cdots\mathcal O(a_D)\cdot V)$, the sum over the relevant minimal primes $V$ of
$P+(2)$, where $m_V$ is the multiplicity of $V$ in $R/(P+(2))$. By
Lemma~\ref{cv-origin-weight-nef-generic-conservation} with $p=2$, this sum is
$(\mathcal O(a_1)\cdots\mathcal O(a_D)\cdot X_{\mathbf Q})$.
\end{proof}
% Lean: GlobalTheory.T.generic_intersection_equals_points_integral = PB.curvatureIntegral_eq_genericIntersection
% = PositivityTheory.Xgp.t_deg_generic, under PB.Inputs: Targets.ExistsUnique.full and
% IntNumIsFieldIntersection.full identify the field intersection family I with intNum over Q
% (isProjVarF_map_of_isArithVar); Ar.Giv.degreeGeneric_of_degree rewrites with DegreeStatement at p = 2 (Lemma
% cv-origin-weight-nef-generic-fibre) and Targets.FibreDegreeSum.proof at p = 2 (Lemma
% cv-origin-weight-nef-generic-conservation).

\begin{theorem}\label{cv-origin-weight-nef-kodaira}
Let $K$ and $(I_k)_{k\in K}$ be as in Definition~\ref{lit-multiprojective}, let $\Omega$ be a field of
characteristic zero, and let $y\in\Omega^{\sqcup_kI_k}$ be a point with function field $\Omega$ whose kernel is
an arithmetic subvariety $X$ of dimension $D+1$, with $D\ge1$. Let $a\in\mathbf Z^K$ be such that
$\mathcal O(a)$ has nonnegative degree on every curve in the generic fibre $X_{\mathbf Q}$ of $X$ and
$(\mathcal O(a)^D\cdot X_{\mathbf Q})>0$. Then there are an integer $k\ge1$ and $e\in\Omega$, $e\neq0$, such
that $e\prod_ly_{l,\kappa_l}^{-(ka_l-1)}\in\mathbf Q[y_{l,i}/y_{l,\kappa_l}:l\in K,\ i\in I_l]$ for every choice
of indices $\kappa_l\in I_l$ with all $y_{l,\kappa_l}\neq0$.
\end{theorem}
\leanref{GroundTheory/HeightsAndTheta/Heights/HilbertSamuel/StatementsKodairaGeneric.lean}{GroundTheory.Heights₂.Heights.HilbertSamuel.KodairaGeneric.Statement}
\leanref{Development/AutomorphicSpectrum/Proofs/AlgebraProofs/KodairaGenericOfAmpleSectionCount.lean}{GlobalTheory.CompanionTheory.Relay.kodaira_generic_of_ample_section_count}

The elements $e$ of the conclusion are the global sections of $\mathcal O(ka-\mathbf 1)$ on $X_{\mathbf Q}$, where
$\mathbf 1=(1,\dots,1)$, described as in Definition~\ref{lit-multiprojective}(4) with $\mathbf Q$ in place of
$\mathbf Z$, and $\mathcal O(\mathbf 1)$ is ample. A nef line bundle with positive top self-intersection is big
\cite[Theorem~2.2.16]{Laz04}, and for every ample line bundle $A$ some multiple of a big line bundle is
linearly equivalent to $A$ plus an effective divisor \cite[Corollary~2.2.7]{Laz04}. Lazarsfeld states these
for complex projective varieties; for irreducible projective varieties over any field they are
\cite[Theorem~4.6 and Lemma~4.3]{Gub08}.

\begin{lemma}\label{cv-origin-weight-nef-sections}
Let $K$ and $(I_k)_{k\in K}$ be as in Definition~\ref{lit-multiprojective}, let $\Omega$ be a field of
characteristic zero, and let $y\in\Omega^{\sqcup_kI_k}$ be a point with function field $\Omega$ whose kernel is
an arithmetic subvariety $X$ of dimension $D+1$, with $D\ge1$. Let $a\in\mathbf Z^K$ be such that
$\mathcal O(a)$ has nonnegative degree on every curve in the generic fibre $X_{\mathbf Q}$ of $X$ and
$(\mathcal O(a)^D\cdot X_{\mathbf Q})>0$. Then there are an integer $k\ge1$ and $e\in\Omega$, $e\neq0$, such
that $e\prod_ly_{l,\kappa_l}^{-(ka_l-1)}\in\mathbf Z[y_{l,i}/y_{l,\kappa_l}:l\in K,\ i\in I_l]$ for every choice
of indices $\kappa_l\in I_l$ with all $y_{l,\kappa_l}\neq0$.
\end{lemma}
\leanref{ReciprocityTheory/ReductiveGroups/ArithmeticGeometry/ArakelovGeometry/KleimanNefness/TargetSt.lean}{ReciprocityTheory.ChartRestriction.G.ArakelovGeometry.KH.SiuZStatement}
\leanref{Development/AutomorphicSpectrum/Proofs/AlgebraProofs/SiuInequalityIntegralHoldsStatement.lean}{GlobalTheory.T.siu_inequality_integral_holds}
\leanref{GroundTheory/HeightsAndTheta/Heights/SemipositiveHeights/OfRemaining.lean}{GroundTheory.Heights₂.Heights.ForSecondHalf.SiuZB.proof}
\leanref{GroundTheory/HeightsAndTheta/Heights/HilbertSamuel/PChartDenominator.lean}{GroundTheory.Heights₂.Heights.HilbertSamuel.ChartDenominator.proof}
\begin{proof}
Call a choice of indices $\kappa$ as in the statement a chart, and write $A_\kappa$ for the ring
$\mathbf Z[y_{l,i}/y_{l,\kappa_l}]$. By Theorem~\ref{cv-origin-weight-nef-kodaira} there are $k\ge1$ and
$e_0\neq0$ such that, for every chart $\kappa$, $s_\kappa=e_0\prod_ly_{l,\kappa_l}^{-(ka_l-1)}$ lies in
$\mathbf Q[y_{l,i}/y_{l,\kappa_l}]$. Then $s_\kappa$ is a polynomial with rational coefficients in the ratios
$y_{l,i}/y_{l,\kappa_l}$, so $n_\kappa s_\kappa\in A_\kappa$ for a common denominator $n_\kappa\ge1$ of these
coefficients. There are finitely many charts. Put $n=\prod_\kappa n_\kappa$ and $e=ne_0$. Then $e\neq0$, as
$\Omega$ has characteristic zero, and for every chart $\kappa$,
$e\prod_ly_{l,\kappa_l}^{-(ka_l-1)}=(n/n_\kappa)\,n_\kappa s_\kappa$ lies in $A_\kappa$.
\end{proof}
% Lean: SiuZB.proof: hagree (Targets.ExistsUnique.full) turns the field intersection family I into intNum
% over Q, so hPos is genDeg > 0; the case D >= 1 is KodairaGeneric (Theorem cv-origin-weight-nef-kodaira), with
% NefOnGenericDeg from NefOnF; ChartDenominator.proof clears the denominators (exists_nat_mul_mem_adjoin_int on
% each chart, then the product of the N c over all c). The Lean's case D = 0, through SectionsDimZero.proof, is
% not needed here.

\begin{lemma}\label{cv-origin-weight-nef-shift}
Let $K$ and $(I_k)_{k\in K}$ be as in Definition~\ref{lit-multiprojective}, let $\Omega$ be a field, and let
$y\in\Omega^{\sqcup_kI_k}$ be such that for every $l\in K$ some $y_{l,i}$ is nonzero. Call a choice of indices
$\kappa_l\in I_l$ with all $y_{l,\kappa_l}\neq0$ a chart. Let $a\in\mathbf Z^K$ and $d\in\mathbf Z_{\ge0}^K$, and
let $k\ge1$ and $e\in\Omega$, $e\neq0$, be such that
$e\prod_ly_{l,\kappa_l}^{-(ka_l-1)}\in\mathbf Z[y_{l,i}/y_{l,\kappa_l}]$ for every chart $\kappa$. Then there are
$k'\ge1$ and $e'\in\Omega$, $e'\neq0$, such that
$e'\prod_ly_{l,\kappa_l}^{-(k'a_l-1-d_l)}\in\mathbf Z[y_{l,i}/y_{l,\kappa_l}]$ for every chart $\kappa$.
\end{lemma}
\leanref{Development/AutomorphicSpectrum/Proofs/MetricProofs/ChartLatticeShiftNonzeroStatement.lean}{GlobalTheory.T.chart_lattice_shift_nonzero}
\leanref{Development/WeilBound/CoreStatements/PluriharmonicCharts/XPolSec.lean}{PositivityTheory.Xgp.t_pol_sec}
\begin{proof}
For each $l$ choose $i_l$ with $y_{l,i_l}\neq0$. Put $m=\sum_ld_l$, $k'=(m+1)k\ge1$ and
$e'=e^{m+1}\prod_ly_{l,i_l}^{m-d_l}$, which is nonzero. For every $l$,
$-(k'a_l-1-d_l)=-(m+1)(ka_l-1)-(m-d_l)$, so for every chart $\kappa$
\[
 e'\prod_ly_{l,\kappa_l}^{-(k'a_l-1-d_l)}
 =\Big(e\prod_ly_{l,\kappa_l}^{-(ka_l-1)}\Big)^{m+1}\prod_l\Big(\frac{y_{l,i_l}}{y_{l,\kappa_l}}\Big)^{m-d_l}.
\]
The first factor lies in $\mathbf Z[y_{l,i}/y_{l,\kappa_l}]$ by hypothesis, and so does the second, as
$m-d_l\ge0$ for every $l$.
\end{proof}

\begin{lemma}\label{cv-origin-weight-nef}
Let $\Psi$ be a finite nonempty set and $N\ge2$, put $D=|\Psi|N(N-1)/2$, let $S$ be an analytic system of
type $(\Psi,N)$ (Definition~\ref{cv-proportionality-analytic-system}), and let $\Sigma_0\subset\Sigma$ be a set
of indices $b$ for which, for every $t$, the regularized metric of $\bar L_b(t)$ has pluriharmonic local
potentials. Let $\Omega$ be a field of characteristic zero and $y\in\Omega^{\sqcup_kI_k}$ a point with
function field $\Omega$ whose kernel is an arithmetic subvariety $X$ of dimension $D+1$, and let
$\tilde X$, $\lambda_x$, $\pi$, $\bar y_b$ and $\bar y_b(t)$ be as in (3) to (5) of
Definition~\ref{cv-origin-analytic} for $y$ and $S$, except that $\pi$ is not assumed bijective and has some
degree $\deg\pi\ge1$ off a closed set of measure zero. Let $n\in\mathbf Z^\Sigma$ with $n_b\ge0$ for
$b\notin\Sigma_0$, let $a_n=\sum_bn_ba_b$, where $a_b$ is the multidegree of $\bar y_b$, and suppose that
$\deg_S(n)>0$. Then for every multihomogeneous $\beta\in R$ of a multidegree $d_\beta\in\mathbf Z_{\ge0}^K$
that is not in the kernel of $y$ there are $k\ge1$ and $e\in\Omega$, $e\neq0$, with
$e\prod_ly_{l,\kappa_l}^{-(ka_{n,l}-1-d_{\beta,l})}\in\mathbf Z[y_{l,i}/y_{l,\kappa_l}]$ for every chart
$\kappa$ of $y$.
\end{lemma}
\leanref{Development/AutomorphicSpectrum/Proofs/MetricProofs/NefPolarizationStatement.lean}{GlobalTheory.T.nef_polarization}
\begin{proof}
Let $\bar a_n=(a_n,h_n)$ with $h_n=\prod_bh_{\bar y_b}^{n_b}$, where $h_{\bar y_b}$ is the function of
$\bar y_b$. Of $\beta$, only its multidegree $d_\beta$ enters the proof.

\emph{Nefness.} By (3) of Definition~\ref{cv-origin-analytic}, the map from $\tilde X$ to the complex points
of $X$ given by the $\lambda_x$ is bijective and immersive. For $b\in\Sigma_0$ the first Chern form of
$\bar L_b(t)$ is $0$ for every $t$, since its local potentials are pluriharmonic. The proof of
Lemma~\ref{cv-nef-positivity} uses of the flag model and its realization only this map, the map $\pi$, and
the bundles $\bar y_b$ and $\bar y_b(t)$ with their relation to the $\bar L_b$ and $\bar L_b(t)$, which (5) of
Definition~\ref{cv-origin-analytic} gives here. So that proof shows that $\mathcal O(a_n)$ has nonnegative
degree on every curve in the generic fibre $X_{\mathbf Q}$ of $X$.

\emph{The degree.} By Lemma~\ref{cv-origin-weight-nef-degree}, $\bar a_n$ is smooth, and the integral over the
complex points of $X$ of the $D$-th power of its first Chern form is $\deg\pi\cdot\deg_S(n)$. This is positive,
since $\deg\pi\ge1$ and $\deg_S(n)>0$. By Lemma~\ref{cv-origin-weight-nef-generic} the integral is
$(\mathcal O(a_n)^D\cdot X_{\mathbf Q})$, so $(\mathcal O(a_n)^D\cdot X_{\mathbf Q})>0$.

\emph{The sections.} As $\Psi$ is nonempty and $N\ge2$, $D\ge1$. Lemma~\ref{cv-origin-weight-nef-sections},
applied to $a=a_n$, gives $k\ge1$ and $e\neq0$ with
$e\prod_ly_{l,\kappa_l}^{-(ka_{n,l}-1)}\in\mathbf Z[y_{l,i}/y_{l,\kappa_l}]$ for every chart $\kappa$ of $y$.
The kernel of $y$ is relevant: for every $l\in K$ some $x_{l,i}$ is not in it, and then $y_{l,i}\neq0$.
Lemma~\ref{cv-origin-weight-nef-shift}, applied with $d=d_\beta$, gives the claim.
\end{proof}
% Lean: PB.nefPol (GlobalTheory.T.nef_polarization), under PB.Inputs. The nefness is its first clause hnef
% (Lemma cv-nef-positivity): for a curve W of the generic fibre, arithmetic_variety_generic_fibre_projective
% gives the arithmetic curve over it, GlobalTheory.T.geometric_degree_nonnegative_psh its nonnegative degree and
% generic_intersection_equals_points_integral with D = 1 the comparison with I W 1. Then hRe =
% GlobalTheory.T.realisation_potential_degree_compatibility (clause 4: hsm, hdeg), hgp from Re.deg_pos and
% hpos, h1 and hI0 from generic_intersection_equals_points_integral with D, hsec = hIn.siuZ (KH.SiuZStatement,
% a component of PB.Inputs, with Re.arith and Birational y), and GlobalTheory.T.chart_lattice_shift_nonzero
% with exists_ne_zero_of_relevant (Re.arith.2.1.2.2). The hypotheses IsMultiHomog beta d_beta and
% beta not in mCoordIdeal y are not used.

\begin{lemma}\label{cv-origin-weight-gap-uniformized}
In the setting of Definition~\ref{cv-origin-reference}, suppose $p+q\ge1$, and let a reference geometry of
$\mathcal X$ with analytic data be given (Definition~\ref{cv-origin-analytic}), with analytic system $S$. Then $S$
is uniformized of type $p$ (Definition~\ref{cv-proportionality-analytic-system}(2)).
\end{lemma}
\leanref{Development/WeilBound/CoreStatements/CentralCharts/FlagImmersion.lean}{AnalyticReference.unif_of}
\begin{proof}
\openstep{The construction of period charts of type $p$ near each point of $S\setminus\partial S$, from the
period lattices of the fibres of $\mathcal A_Z$ and the frames of the flag of the reference geometry, and of a
point at which such a chart is a flag immersion, from the countability of the points near the generic point of
$Z$ whose periods have given flags, is not yet written out.}
% Lean: AnalyticReference.unif_of, under Top.Given and 1 <= p + q, with the inputs RootFieldNumerics.LocPer,
% Transport, FrameVals, PerMat, FlagImm and CountGen, which gap_polynomial_positive supplies as locPer,
% transport, frameVals, perMat_of perCoord matStep, flagImm_of rankLem and countGen_of (genRigid_of frameVals).
% Its proof: exists_inv (a holomorphic inverse psi of pi); ptsFam (injective complex points z of Z with
% coordinates the lift of psi x, holomorphic in the charts of S, and off the boundary in every open neighbourhood
% of the generic point); for x off the boundary, locChart (LocPer: a holomorphic family R of period lattices;
% FrameVals: the frames of the flag as lattice vectors with Gram matrix pi times R.gram; PerMat: matrices gm in
% the flag domain with stepPot the sums of grLog of Gram determinants) and periodChart give the period chart; the
% open set U of CountGen, a point of ptsFam in it, locChart with U, FlagImm and flagImmersion give the flag
% immersion. Size: most of the 56,000 lines of Lemma cv-origin-weight-gap; the directories FlagImmersions and
% CentralCharts of Development/WeilBound/CoreStatements alone have about 19,500 lines.
\end{proof}

\begin{lemma}\label{cv-origin-weight-gap}
In the setting of Definition~\ref{cv-origin-reference}, let the moduli schemes at all levels be given, let
$\mathcal X$ be a global model (Definition~\ref{cv-volume-global-model}), and suppose $p\ge1$ and $q\ge1$; the
hypotheses that $\mathcal X$ is a global model and that $q\ge1$ are not used in the proof. Let
a reference geometry of $\mathcal X$ with analytic data be given, with analytic system $S$, and let
$\nu\in\mathbf Z^{\Psi\times\{1,\dots,N\}}$ with $\nu_{\tau,i}>\nu_{\tau,k}$ for $i<k$. Let $c\in\mathbf R^\Sigma$
be given by $c_{\tau,i}=\nu_{\tau,i}-\nu_{\tau,i+1}$, with $\nu_{\tau,N+1}=0$, and $c_A=0$. Then $\deg_S(c)>0$.
\end{lemma}
\leanref{Development/AutomorphicSpectrum/Proofs/GeometryProofs/GlobalReferenceGapPositiveStatement.lean}{GlobalTheory.T.global_reference_gap_positive}
\leanref{Development/WeilBound/CoreStatements/CentralCharts/GapPolynomialPositive.lean}{AnalyticReference.gap_polynomial_positive}
\begin{proof}
Here $N=p+q$, so $0\le p\le N$; let $d=|\Psi|N(N-1)/2$. As $p+q\ge p\ge1$,
Lemma~\ref{cv-origin-weight-gap-uniformized} shows that $S$ is uniformized of type $p$. As $c_A=0$,
\[
 \deg_S(c)=\int_S\Big(\sum_{\tau,i}(\nu_{\tau,i}-\nu_{\tau,i+1})\,c_1(\bar L_{\tau,i})\Big)^d .
\]
Theorem~\ref{cv-proportionality-mumford}, with $j=p$ and $\lambda=\nu$, gives $v_S>0$ with
\[
 \deg_S(c)=v_S\prod_{\tau\in\Psi}\prod_{1\le i<k\le N}(\nu_{\tau,i}-\nu_{\tau,k}).
\]
Every factor $\nu_{\tau,i}-\nu_{\tau,k}$ with $i<k$ is positive, so $\deg_S(c)>0$.
\end{proof}
% Lean: AnalyticReference.gap_polynomial_positive: hU = AnalyticReference.unif_of (Lemma
% cv-origin-weight-gap-uniformized); GlobalTheory.T.uniformised_degree_polynomial_antisymmetric (=
% PB.anSys_proportional, Theorem cv-proportionality-mumford) for A.Sys.comap some with N = p + q, j = p and hdim
% (dimFl = antisymDegree), with PB.Inputs from Top.Given; hgap, PB.eval_toWt_gap and Deg.eval_formPoly_none
% (killNone_formPoly) for the evaluation at gapN nu, which is 0 at the index A; Deg.eval_antisymDelta_pos for the
% sign. The hypotheses IsGlobal and 1 <= q are not used: unif_of takes 1 <= p + q, which is le_add_right hp, and
% anSys_proportional takes only j <= N, which is Nat.le_add_right. Theorem cv-proportionality-mumford now states
% the hypotheses of anSys_proportional (Psi finite, possibly empty, any N, 0 <= j <= N), so nothing else is needed.

\begin{lemma}\label{cv-origin-weight-data}
In the setting of Definition~\ref{cv-origin-reference}, let a reference geometry of $\mathcal X$ be given,
with its notation, and let $n\in\mathbf Z^{\Psi\times\{1,\dots,N\}}$; put $\alpha=-\sum_{\tau,i}n_{\tau,i}a_{\tau,i}$.
Let $e,b\in\Omega$, $d_b\in\mathbf Z_{\ge0}^K$, and for each $k\in K$ an index $i_0(k)\in I_k$ with
$y_{k,i_0(k)}\neq0$. Suppose that for every chart $\kappa$ of $y$,
\[
 e\prod_ky_{k,\kappa_k}^{-(\alpha_k-1-d_{b,k})}\in\mathbf Z[y_{k,i}/y_{k,\kappa_k}]
 \quad\text{and}\quad
 b\prod_ky_{k,\kappa_k}^{-d_{b,k}}\in\mathbf Z[y_{k,i}/y_{k,\kappa_k}],
\]
and that at every point $P$ of $Z$ over no prime over which $\mathcal A_Z$ is not abelian, and every chart
$\kappa$ with $P\in Z_\kappa$, $b\prod_ky_{k,\kappa_k}^{-d_{b,k}}$ is not a unit at $P$. Put
\[
 s=e\,b\prod_{k\neq k_0}y_{k,i_0(k)}\prod_{\tau,i}e_{\tau,i}^{-n_{\tau,i}}.
\]
Then at every point $P$ of $Z$ over no prime, if $i\in I_{k_0}$ is a chart of the block $k_0$ at $P$ and
$t$ generates $\bigotimes_{\tau,l}\Lambda_{\tau,l}^{-n_{\tau,l}}$ at $P$, then $s\,y_{k_0,i}/t$ is regular at
$P$, and it is not a unit at $P$ if $\mathcal A_Z$ is not abelian at $P$.
\end{lemma}
\leanref{Development/AutomorphicSpectrum/Proofs/GeometryProofs/ExistsEffectiveWeightSection.lean}{GlobalTheory.CompanionTheory.Ar.ExRefEff.isEff_of_data}
\leanref{Development/AutomorphicSpectrum/Proofs/GeometryProofs/LatticeTensorPowerDegreeStatement.lean}{GlobalTheory.T.lattice_tensor_power_degree}
\begin{proof}
Let $P$, $i$ and $t$ be as in the statement. The point $P$ lies in some $Z_{c_0}$, and the chart $c$ equal
to $c_0$ with $c_{0,k_0}$ replaced by $i$ satisfies $P\in Z_c$: $y_{k_0,i}\neq0$, and the ratios of the
block $k_0$ are regular at $P$ by hypothesis and those of the other blocks because $P\in Z_{c_0}$. Write
$Y_k=y_{k,c_k}$, so $Y_{k_0}=y_{k_0,i}$. The root line $\Lambda_{\tau,l}$ is generated at $P$ by
$e_{\tau,l}\prod_kY_k^{a_{\tau,l,k}}$, so $\bigotimes\Lambda_{\tau,l}^{-n_{\tau,l}}$ is generated at $P$ by
$g=E_0\prod_kY_k^{\alpha_k}$ with $E_0=\prod e_{\tau,l}^{-n_{\tau,l}}\neq0$, and $g=vt$ for some $v$ regular
at $P$. Put $Q_1=\prod_kY_k^{-(\alpha_k-1-d_{b,k})}$, $Q_2=\prod_kY_k^{-d_{b,k}}$ and
$M=\prod_{k\neq k_0}y_{k,i_0(k)}/Y_k$. The exponents of $Y_k$ in $Q_1Q_2\prod_kY_k^{\alpha_k}$ are all $1$, so
\[
 \frac{s\,y_{k_0,i}}t=\frac{vs\,Y_{k_0}}{E_0\prod_kY_k^{\alpha_k}}=v\cdot(eQ_1)\cdot(bQ_2)\cdot M .
\]
The factors $eQ_1$ and $bQ_2$ lie in $\mathbf Z[y_{k,j}/Y_k]$ by hypothesis, and so does $M$; this ring
consists of functions regular at $P$, since $P\in Z_c$. Hence $s\,y_{k_0,i}/t$ is regular at $P$. If
$\mathcal A_Z$ is not abelian at $P$ and $s\,y_{k_0,i}/t$ were a unit at $P$, then $bQ_2$ would be nonzero
and $(bQ_2)^{-1}=(s\,y_{k_0,i}/t)^{-1}\,v\,(eQ_1)M$ would be regular at $P$, so $bQ_2$ would be a unit at
$P$, against the hypothesis on $b$ for the chart $c$.
\end{proof}

\begin{lemma}\label{cv-origin-weight}
In the setting of Definition~\ref{cv-origin-reference}, let the moduli schemes at all levels be given, let
$\mathcal X$ be a global model (Definition~\ref{cv-volume-global-model}), and suppose $p\ge1$ and $q\ge1$. Let
a reference geometry of $\mathcal X$ with analytic data be given, and let $\nu\in\mathbf Z^{\Psi\times\{1,\dots,N\}}$
with $\nu_{\tau,i}>\nu_{\tau,k}$ for $i<k$. Then there are integers $k_A,q\ge1$ and $s\in\Omega$, $s\neq0$, with
the property (2c)(iii) of Definition~\ref{cv-origin-tower} for $n_{\tau,i}=k_Aq\nu_{\tau,i}$, the root lines of
the reference geometry, the index $k_0$, and $s$ and $q$ in place of $g(s_c)$ and $q_c$: at every point $P$ of
$Z$ over no prime, if $y_{k_0,i}\neq0$, the ratios $y_{k_0,l}/y_{k_0,i}$ are regular at $P$ and $t$ generates
$\bigotimes_{\tau,l}\Lambda_{\tau,l}^{-n_{\tau,l}}$ at $P$, then $s\,y_{k_0,i}^q/t$ is regular at $P$, and it is
not a unit at $P$ if $\mathcal A_Z$ is not abelian at $P$.
\end{lemma}
\leanref{Development/AutomorphicSpectrum/Proofs/GeometryProofs/ExistsEffectiveWeightSectionStatement.lean}{GlobalTheory.T.exists_effective_weight_section}
\leanref{Development/AutomorphicSpectrum/Proofs/GeometryProofs/ExistsEffectiveWeightSection.lean}{GlobalTheory.CompanionTheory.Ar.ExRefEff.exists_eff}
\leanref{Development/AutomorphicSpectrum/Definitions/GlobalGeometry/EmbedsGeneric.lean}{GlobalTheory.CompanionTheory.Ar.RefGeo.IsEff}
\begin{proof}
The proof gives $q=1$. Let $S$, $\tilde X$, $\pi$ and $\bar y_b$ be those of the analytic data, and let $c$
be as in Lemma~\ref{cv-origin-weight-gap}. Let $\Sigma_0$ be the set of indices $(\tau,N)$. By
Definition~\ref{cv-origin-analytic}(v), the regularized metrics of the $\bar L_b(t)$, $b\in\Sigma_0$, have
pluriharmonic local potentials. For $b\notin\Sigma_0$, $c_b\ge0$: $c_A=0$, and for $i<N$,
$c_{\tau,i}=\nu_{\tau,i}-\nu_{\tau,i+1}>0$. By Lemma~\ref{cv-origin-weight-gap}, $\deg_S(c)>0$. By
Lemma~\ref{cv-origin-weight-birational}, $\Omega$ is the function field of the kernel $X$ of $y$; $X$ has
dimension $D+1$ by Definition~\ref{cv-origin-analytic}(3), and the maps $\lambda_x$, $\pi$ and the bundles of
the analytic data satisfy the hypotheses of Lemma~\ref{cv-origin-weight-nef}. Let $\beta$ and $d_\beta$ be
given by Lemma~\ref{cv-origin-weight-boundary}. Lemma~\ref{cv-origin-weight-nef} gives $k\ge1$ and $e\neq0$
with $e\prod_ly_{l,\kappa_l}^{-(ka_{c,l}-1-d_{\beta,l})}\in\mathbf Z[y_{l,i}/y_{l,\kappa_l}]$ for every chart
$\kappa$.

By Definition~\ref{cv-origin-analytic}(ii) and (iii), $a_c=-\sum_\tau\sum_ic_{\tau,i}\sum_{l\le i}a_{\tau,l}$, and
summation by parts gives $\sum_{i\ge l}c_{\tau,i}=\nu_{\tau,l}$, so $a_c=-\sum_{\tau,l}\nu_{\tau,l}a_{\tau,l}$. Put
$n=k\nu$; then $ka_c=-\sum_{\tau,l}n_{\tau,l}a_{\tau,l}$, which is the $\alpha$ of
Lemma~\ref{cv-origin-weight-data}. Choose for each $k'\in K$ an index $i_0(k')$ with $y_{k',i_0(k')}\neq0$;
the generic point of $Z$ lies in a chart. Lemma~\ref{cv-origin-weight-data}, applied to $n$, $e$,
$b=\beta(y)$ and $d_b=d_\beta$, shows that
$s=e\,\beta(y)\prod_{k'\neq k_0}y_{k',i_0(k')}\prod e_{\tau,i}^{-n_{\tau,i}}$ has the property with
$k_A=k$ and $q=1$. Each factor of $s$ is nonzero, so $s\neq0$.
\end{proof}
% Lean: ExRefEff.exists_eff: hodge_cpt_birational_arith_var (Birational y), reference_geometry_multihomog_form
% (beta, d_beta), Fl the indices (tau, N - 1) zero-based with A.flat, gapN nu >= 0 off Fl, nef_polarization (the
% second clause, with the field intersection family intNum) and global_reference_gap_positive; comb_a (Abel
% summation through A.fund_deg) for the multidegree, and isEff_of_data with q' = 1.

\begin{lemma}\label{cv-origin-reference-closure}
In the setting of Definition~\ref{cv-origin-reference}, let a reference geometry of $\mathcal X$ be given,
with its notation. Then $I_{k_0}$ is nonempty and the ratios $y_{k_0,i}/y_{k_0,j}$, $i,j\in I_{k_0}$,
generate the field $\Omega$.
\end{lemma}
\leanref{Development/AutomorphicSpectrum/Proofs/GeometryProofs/EffectiveScaledWeightOfRetraction.lean}{GlobalTheory.CompanionTheory.Ar.RefRoot.EmbedsGeneric.exists_closure}
\begin{proof}
The local ring of $Z$ at its generic point $\eta$ is $\Omega$, in which every prime is a unit, so $\eta$ lies
over no prime. By Definition~\ref{cv-origin-reference}(5), $\eta$ lies in a chart $i$ of the block $k_0$, so
$I_{k_0}$ is nonempty, and every $a\in\Omega$ has an element $s$ of $\mathbf Q[y_{k_0,j}/y_{k_0,i}]$, a unit at
$\eta$ and so nonzero, with $sa$ in the same ring. Then $a=(sa)/s$ lies in the field generated by the ratios.
\end{proof}

\begin{lemma}\label{cv-origin-orbit-forms}
Let $X$, $\Omega$, $\xi$ and $y\in\Omega^{\sqcup_{k\in K}I_k}$ be as in Lemma~\ref{cv-origin-snc-charts}, with $\Omega$ of
characteristic zero, and suppose that the block $k_0\in K$ of $y$ embeds the generic fibre in the sense of
Definition~\ref{cv-origin-reference}(5), with $X$ in place of $Z$. Let $L$ be a field of characteristic zero and
$z\in L^{\sqcup_kI_k}$ a point with $(z_{k,l})_{l\in I_k}\neq0$ for every $k$, at which every element of the kernel of $y$
vanishes. Let $i\in I_{k_0}$, $k\in K$ and $c\in I_k$ with $z_{k_0,i}\neq0$ and $z_{k,c}\neq0$. Then there are an integer
$n\ge0$ and, for each $j\in I_k$, a homogeneous polynomial $F_j$ of degree $n$ with integer coefficients in the variables
$x_{k_0,l}$, $l\in I_{k_0}$, such that $F_j(y_{k_0})\,y_{k,c}=F_c(y_{k_0})\,y_{k,j}$ for every $j$ and
$F_c(z_{k_0})\neq0$, where $y_{k_0}$ and $z_{k_0}$ are the blocks $k_0$ of $y$ and $z$.
\end{lemma}
\leanref{Development/AutomorphicSpectrum/Proofs/GeometryProofs/SmallPointsCoordRatioBoundStatement.lean}{GlobalTheory.T.small_points_coord_ratio_bound}
\begin{proof}
If $y_{k',l}=0$, then $x_{k',l}$ lies in the kernel of $y$, so $z_{k',l}=0$. Choose for each $k'$ an index $\kappa_{k'}$
with $z_{k',\kappa_{k'}}\neq0$; then $\kappa$ is a chart of $y$. Lemma~\ref{cv-origin-boundary-points}(2) gives a
morphism $\zeta:\operatorname{Spec}L\to X$ with coordinates $z$; let $P$ be its image. A positive integer has nonzero
image in $L$, and the homomorphism $\mathcal O_{X,P}\to L$ is local, so it is a unit at $P$: $P$ lies over no prime.

Let $B=\mathbf Q[y_{k_0,l}/y_{k_0,i}:l\in I_{k_0}]$. By the coordinates, $y_{k_0,i}\neq0$ and every $y_{k_0,l}/y_{k_0,i}$
is regular at $P$, with value $z_{k_0,l}/z_{k_0,i}$: $P$ lies in the chart $i$ of the block $k_0$. Likewise $y_{k,c}\neq0$
and every $y_{k,j}/y_{k,c}$ is regular at $P$. By Definition~\ref{cv-origin-reference}(5), for each $j$ there is
$s_j\in B$, a unit at $P$, with $s_jy_{k,j}/y_{k,c}\in B$. Then $s=\prod_js_j$ is a unit at $P$ and
$s\,y_{k,j}/y_{k,c}\in B$ for every $j$.

Every $b\in B$ has an integer $N\neq0$, an integer $m\ge0$ and a homogeneous polynomial $G$ of degree $m$ with integer
coefficients in the $x_{k_0,l}$ such that $G(y_{k_0})=Ny_{k_0,i}^m\,b$: for the generators take $N=m=1$ and
$G=x_{k_0,l}$, for a rational number take $m=0$ and $N$ its denominator, for a sum multiply the two polynomials by
an integer and a power of $x_{k_0,i}$ to bring them to a common $N$ and $m$, and for a product multiply them. Doing the
same for the finitely many elements $s\,y_{k,j}/y_{k,c}$ gives a common $N\neq0$ and $n$, and polynomials $F_j$ with
$F_j(y_{k_0})=Ny_{k_0,i}^n\,s\,y_{k,j}/y_{k,c}$. So $F_j(y_{k_0})\,y_{k,c}=Ny_{k_0,i}^n\,s\,y_{k,j}$ and
$F_c(y_{k_0})=Ny_{k_0,i}^n\,s$, which gives the relations.

By homogeneity $F_c(y_{k_0})=y_{k_0,i}^nF_c(y_{k_0,l}/y_{k_0,i})$, so $F_c(y_{k_0,l}/y_{k_0,i})=Ns$. Both sides are
regular at $P$, and their values at $\zeta$ give $F_c(z_{k_0,l}/z_{k_0,i})=N\sigma$, where $\sigma$ is the value of $s$,
nonzero since $s$ is a unit at $P$. As $N\neq0$ in $L$, $F_c(z_{k_0})=z_{k_0,i}^nF_c(z_{k_0,l}/z_{k_0,i})\neq0$.
\end{proof}
% Lean: BlockRatio.ratios_quotients_of_forms: local_ring_points_existsUnique for zz, unitAt_natCast, the EmbedsGeneric
% hypothesis hE, isFormQuot_of_mem and common_denominator for the forms, exists_unit_germ and aeval_ne_zero_of_unit for
% the last paragraph.

\begin{lemma}\label{cv-origin-reference-orbit}
In the setting of Definition~\ref{cv-origin-reference}, let a reference geometry of $\mathcal X$ be given,
with its notation, and let $X$ be the kernel of $y$. If $w,w'\in\hat X(\mathbf C)$ and $u\in\mathbf C$ satisfy
$w_{k_0,i}=uw'_{k_0,i}$ for all $i\in I_{k_0}$, then $w$ and $w'$ lie in one orbit of $(\mathbf C^\times)^K$.
\end{lemma}
\leanref{Development/AutomorphicSpectrum/Proofs/GeometryProofs/SameOrbitOfProportionalCoordsStatement.lean}{GlobalTheory.T.same_orbit_of_proportional_coords}
\begin{proof}
Let $i\in I_{k_0}$ with $w_{k_0,i}\neq0$; then $u\neq0$. Let $k\in K$ and $c\in I_k$ with $w_{k,c}\neq0$.
Lemma~\ref{cv-origin-orbit-forms}, applied to $Z$, $L=\mathbf C$ and $z=w$, gives $n$ and polynomials $F_j$. For each $j$
the polynomial $G_j=F_j(x_{k_0})x_{k,c}-F_c(x_{k_0})x_{k,j}$ satisfies
$G_j(y_{k',l}T_{k'})=T_{k_0}^nT_k\big(F_j(y_{k_0})y_{k,c}-F_c(y_{k_0})y_{k,j}\big)=0$, so it lies in the kernel of $y$ and
vanishes at $w$ and at $w'$:
\[
 F_j(w_{k_0})\,w_{k,c}=F_c(w_{k_0})\,w_{k,j},\qquad F_j(w'_{k_0})\,w'_{k,c}=F_c(w'_{k_0})\,w'_{k,j}.
\]
As $w_{k_0}=uw'_{k_0}$, $F_j(w_{k_0})=u^nF_j(w'_{k_0})$. So $F_c(w'_{k_0})\neq0$, and cancelling $u^n$,
$F_j(w'_{k_0})\,w_{k,c}=F_c(w'_{k_0})\,w_{k,j}$. Some $w'_{k,j}$ is nonzero, so the second relation gives
$w'_{k,c}\neq0$. Put $a_k=w_{k,c}/w'_{k,c}\neq0$. Then
$F_c(w'_{k_0})\,w_{k,j}=F_j(w'_{k_0})\,w_{k,c}=a_kF_j(w'_{k_0})\,w'_{k,c}=a_kF_c(w'_{k_0})\,w'_{k,j}$, so
$w_{k,j}=a_kw'_{k,j}$ for every $j$. So $w=a\cdot w'$ with $a=(a_k)\in(\mathbf C^\times)^K$.
\end{proof}
% Lean: OrbitBlock.sameOrbit_of_block: small_points_coord_ratio_bound for w (hpts), relation_at (with
% mConeMap_rename_block) for w and w', aeval_smul_of_homog for the scaling by u.

\begin{lemma}\label{cv-origin-ample-families}
Let $X$, $\Omega$, $\xi$, $y\in\Omega^{\sqcup_{k\in K}I_k}$ and $k_0$ be as in Lemma~\ref{cv-origin-orbit-forms}, and let
$k\in K$. Then there are finitely many families $(F_{t,j})_{j\in I_k}$, $1\le t\le l$, of polynomials with integer
coefficients in the variables $x_{k_0,i}$, $i\in I_{k_0}$, each family homogeneous of one degree $n_t$, such that
$F_{t,j}(y_{k_0})\,y_{k,j'}=F_{t,j'}(y_{k_0})\,y_{k,j}$ for all $t$, $j$ and $j'$, and such that for every field $L$ of
characteristic zero and every $z\in L^{\sqcup_{k'}I_{k'}}$ with $(z_{k',l})_l\neq0$ for every $k'$, at which every element of
the kernel of $y$ vanishes, some $F_{t,j}(z_{k_0})$ is nonzero.
\end{lemma}
\leanref{Development/AutomorphicSpectrum/Proofs/GeometryProofs/CoordRelationsHomogeneousStatement.lean}{GlobalTheory.T.coord_relations_homogeneous}
\begin{proof}
Every block of $y$ has a nonzero entry, as the image of $\xi$ lies in a chart, and $z_{k',l}=0$ whenever $y_{k',l}=0$,
since $x_{k',l}$ then lies in the kernel of $y$. Call a family $(F_j)_{j\in I_k}$ \emph{admissible} if it is homogeneous of
one degree and $F_j(y_{k_0})y_{k,j'}=F_{j'}(y_{k_0})y_{k,j}$ for all $j,j'$, and let $\mathcal S$ be the set of members
of admissible families. For $z$ as in the statement, choose $i$ and $c$ with $z_{k_0,i}\neq0$ and $z_{k,c}\neq0$. The
family of Lemma~\ref{cv-origin-orbit-forms} is admissible, since $y_{k,c}\neq0$ and
$F_j(y_{k_0})y_{k,j'}y_{k,c}=F_c(y_{k_0})y_{k,j}y_{k,j'}=F_{j'}(y_{k_0})y_{k,j}y_{k,c}$, and its member $F_c$ does not vanish at
$z_{k_0}$.

Fix $i\in I_{k_0}$ with $y_{k_0,i}\neq0$ and let $B=\mathbf Q[y_{k_0,l}/y_{k_0,i}]$. The ideal of $B$ generated by the
$F(y_{k_0,l}/y_{k_0,i})$, $F\in\mathcal S$, is $B$. Otherwise it lies in a maximal ideal $\mathfrak m$ of $B$, and a
valuation ring $W$ of $\Omega$ dominates $B_{\mathfrak m}$ \cite[Tag 00IA]{Sta}. Its residue field $\kappa_W$ has
characteristic zero, as $1/n\in B\subset W$ for integers $n\ge1$. For each $k'\neq k_0$ choose $c_{k'}$ with
$y_{k',c_{k'}}\neq0$ of least valuation in its block, and put $c_{k_0}=i$; then every $y_{k',l}/y_{k',c_{k'}}$ lies in $W$.
Every element $f$ of the kernel of $y$ vanishes at these ratios, as one sees by putting $T_{k'}=y_{k',c_{k'}}^{-1}$ in
$f(y_{k',l}T_{k'})=0$. So their residues form a point $z'$ over $\kappa_W$ as in the statement. By the first paragraph some
$F\in\mathcal S$ has $F(z'_{k_0})\neq0$; but $F(z'_{k_0})$ is the residue of $F(y_{k_0,l}/y_{k_0,i})\in\mathfrak m$, which is
$0$.

So $1=\sum_mb_mF_m(y_{k_0,l}/y_{k_0,i})$ for finitely many $F_m\in\mathcal S$ and $b_m\in B$. Let $z$ be as in the
statement with $z_{k_0,i}\neq0$, and suppose that every $F_m(z_{k_0})$ is $0$. Clearing the denominators of the $b_m$ gives
an integer $N\ge1$ and a polynomial $G$ with $G(y_{k_0,l}/y_{k_0,i})=N$ and, by the homogeneity of the $F_m$,
$G(z_{k_0,l}/z_{k_0,i})=0$. The homogenization $H=x_{k_0,i}^e(G-N)(x_{k_0,l}/x_{k_0,i})$, with $e$ the degree of $G$, vanishes
at $y_{k_0}$, so it lies in the kernel of $y$ and vanishes at $z$; dividing by $z_{k_0,i}^e$ gives $N=0$ in $L$, which is
false. If $y_{k_0,i}=0$, then $z_{k_0,i}=0$ for every such $z$. The admissible families of the $F_m$, for all $i$ with
$y_{k_0,i}\neq0$, are the required families.
\end{proof}
% Lean: FormFamilies.finite_families: small_points_coord_ratio_bound and mul_right_cancel for admissibility (Fam);
% exists_finset over the i (exists_finset_at, span_top with Ideal.image_subset_nonunits_valuationSubring,
% exists_ratio_mem, charZero_residue, ratio_pt_mem and aeval_ratio_eq_zero; exists_clear, homogAt and vanish).

\begin{lemma}\label{cv-origin-reference-ample}
In the setting of Definition~\ref{cv-origin-reference}, let a reference geometry of $\mathcal X$ be given,
with its notation. Then the multidegree $a_A$ equal to $1$ at $k_0$ and $0$ at the other indices is ample on
the generic fibre of the kernel of $y$, in the sense of Definition~\ref{cv-origin-tower}(3e).
\end{lemma}
\leanref{Development/AutomorphicSpectrum/Proofs/GeometryProofs/AmpleOfSingleCoordBlockStatement.lean}{GlobalTheory.T.ample_of_single_coord_block}
\begin{proof}
Let $K'=K\setminus\{k_0\}$. For $k\in K'$, Lemma~\ref{cv-origin-ample-families}, applied to $Z$, gives families
$(F^k_{t,j})_{j\in I_k}$, $t\le l_k$, of degrees $n^k_t$. Choose for each $k\in K'$ an index $j_k$ with $y_{k,j_k}\neq0$, and
let $N=\sum_{k\in K'}\sum_tn^k_t$. For $T=(t_k)_{k\in K'}$ and $i\in I_{k_0}$ put $e_T=N-\sum_{k\in K'}n^k_{t_k}\ge0$ and
\[
 s_{T,i}=y_{k_0,i}^{e_T}\prod_{k\in K'}F^k_{t_k,j_k}(y_{k_0})/y_{k,j_k},
\]
and for $J=(j'_k)\in\prod_{k\in K'}I_k$ put $g_{T,i,J}=x_{k_0,i}^{e_T}\prod_{k\in K'}F^k_{t_k,j'_k}(x_{k_0})$, multihomogeneous of
multidegree $N$ at $k_0$ and $0$ elsewhere, and $h_J=\prod_{k\in K'}x_{k,j'_k}$, of multidegree $0$ at $k_0$ and $1$
elsewhere. The relations $F^k_{t,j_k}(y_{k_0})y_{k,j'_k}=F^k_{t,j'_k}(y_{k_0})y_{k,j_k}$ give
$g_{T,i,J}(y)=s_{T,i}\,h_J(y)$. Let $m'=N+1$; the multidegree $m'a_A-(1,\dots,1)$ is $N$ at $k_0$ and $-1$ elsewhere, which
is the multidegree of $g_{T,i,J}$ minus that of $h_J$.

The $s_{T,i}$ lie in $H^0$ of this multidegree, and so in $\tilde H^0$ (Definition~\ref{lit-multiprojective}(4)). Indeed,
for a chart $\kappa$ of $y$ and $J=(\kappa_k)_{k\in K'}$,
\[
 s_{T,i}\prod_{k}y_{k,\kappa_k}^{-(m'a_{A,k}-1)}=s_{T,i}\,h_J(y)\,y_{k_0,\kappa_{k_0}}^{-N}=g_{T,i,J}(y)\,y_{k_0,\kappa_{k_0}}^{-N}
 =g_{T,i,J}\big(y_{k,l}/y_{k,\kappa_k}\big)\in\mathbf Z[y_{k,l}/y_{k,\kappa_k}],
\]
the last equality by the homogeneity of $g_{T,i,J}$.

They have no common zero on the generic fibre. Let $L$ be a field of characteristic zero and $z\in L^{\sqcup_kI_k}$ with
$(z_{k,l})_l\neq0$ for every $k$, at which every element of the kernel of $y$ vanishes. Choose $i$ with $z_{k_0,i}\neq0$,
and for each $k\in K'$, by Lemma~\ref{cv-origin-ample-families}, $t_k$ and $j'_k$ with
$F^k_{t_k,j'_k}(z_{k_0})\neq0$. The polynomial $F^k_{t_k,j'_k}(x_{k_0})x_{k,j}-F^k_{t_k,j}(x_{k_0})x_{k,j'_k}$ lies in the
kernel of $y$, as in the proof of Lemma~\ref{cv-origin-reference-orbit}, so it vanishes at $z$; taking $j$ with
$z_{k,j}\neq0$ shows $z_{k,j'_k}\neq0$. So $h_J(z)\neq0$ and $g_{T,i,J}(z)=z_{k_0,i}^{e_T}\prod_kF^k_{t_k,j'_k}(z_{k_0})\neq0$,
and $s_{T,i}=g_{T,i,J}(y)/h_J(y)$ has the nonzero value $g_{T,i,J}(z)/h_J(z)$ at $z$.
\end{proof}
% Lean: AmpleBlock.ampleQ_of_embeds: coord_relations_homogeneous (Lemma cv-origin-ample-families), then ampleQ_of_forms
% with m = N + 1, the sections sec (S), gForm, hForm, sec_mul, FibreNormDehomRoots.aeval_mem_chartLatNm for the
% membership in chartLatNm, and rel_at_point for z_{k, j'_k} != 0. Ht.AmpleQ asks for a value in the sense of
% Small.HasVal: g(y) = s h(y), deg g = deg h + (m a - 1), h(z) != 0, val = g(z)/h(z) != 0.

\begin{lemma}\label{cv-origin-reference-constants}
In the setting of Definition~\ref{cv-origin-hodge-component}, suppose that $F_c$ is algebraically closed in
the function field of every irreducible component of $\mathcal X$. Let $Z$ be a Hodge component of
$\mathcal X$ at $S$ with coordinates $y\in\Omega^{\sqcup_kI_k}$ and a flag of degree $1$. Then every element of
$\Omega$ algebraic over $\mathbf Q$ lies in the image of $F_c$ in $\Omega$.
\end{lemma}
\leanref{Development/AutomorphicSpectrum/Proofs/GeometryProofs/AlgebraicConstantsInRangeStatement.lean}{GlobalTheory.T.algebraic_constants_in_range}
\leanref{Development/AutomorphicSpectrum/Proofs/GeometryProofs/AlgebraicConstantsInRange.lean}{GlobalTheory.CompanionTheory.Ar.FlagConst.isAlgebraic_mem_cst_range}
\begin{proof}
Let $\Omega_0$ be the residue field of the generic point $\eta_0$ of the irreducible component under $Z$,
embedded in $\Omega$. The structure map of $\mathcal X$ at $\eta_0$ gives an embedding
$c_0:F_c\to\Omega_0$ whose composite with $\Omega_0\subset\Omega$ is the embedding of $F_c$ in $\Omega$; so
$\Omega$ has characteristic zero. By Definition~\ref{cv-origin-hodge-component}(2) with $\deg=1$,
$\Omega=\Omega_0(u)$ for a family $u$ of elements algebraically independent over $\Omega_0$. Let $a\in\Omega$ be
algebraic over $\mathbf Q$; then $a$ is algebraic over $\Omega_0$. The polynomial ring $\Omega_0[u]$ is
integrally closed with fraction field $\Omega$, and $a$ is integral over it, so $a\in\Omega_0[u]$. A polynomial
that is algebraic over $\Omega_0$ divides a nonzero element of $\Omega_0$, so it is a unit or zero, and hence
constant. So $a\in\Omega_0$, and $a$ is algebraic over $\mathbf Z$; as $F_c$, embedded by $c_0$, is
algebraically closed in $\Omega_0$, $a$ lies in the image of $F_c$.
\end{proof}

\begin{lemma}\label{cv-origin-reference-immersion}
In the setting of Definition~\ref{cv-origin-reference}, let a reference geometry of $\mathcal X$ be given,
with its notation, and let $\tilde X$ and the maps $\lambda_x$ be as in Definition~\ref{cv-origin-analytic}(3).
Let $x\in\tilde X$ and $p\in I_{k_0}$ with $\lambda_x(x)_{k_0,p}\neq0$, and let $\chi$ be a chart of $\tilde X$ at
$x$. Then the map $v\mapsto\big(\lambda_x(\chi^{-1}(v))_{k_0,i}/\lambda_x(\chi^{-1}(v))_{k_0,p}\big)_{i\in I_{k_0}}$ is
differentiable at $\chi(x)$, with injective derivative.
\end{lemma}
\leanref{Development/AutomorphicSpectrum/Proofs/GeometryProofs/HolomorphicAtlasCoordCompatStatement.lean}{GlobalTheory.T.holomorphic_atlas_coord_compat}
\begin{proof}
Let $G(v)=\lambda_x(\chi^{-1}(v))$ near $v_0=\chi(x)$; it is holomorphic, $G(v_0)=\lambda_x(x)=w$, and $G(v)\in\hat X(\mathbf C)$
for $v$ near $v_0$. Let $f(v)=\big(G(v)_{k_0,i}/G(v)_{k_0,p}\big)_{i\in I_{k_0}}$; it is differentiable at $v_0$, as
$w_{k_0,p}\neq0$. Choose a chart $p'$ of $w$, that is $w_{k,p'_k}\neq0$ for every $k$, with $p'_{k_0}=p$. For each $k$,
Lemma~\ref{cv-origin-orbit-forms}, applied to $Z$, $L=\mathbf C$, $z=w$, $i=p$ and $c=p'_k$, gives polynomials $F^k_d$,
$d\in I_k$, homogeneous of degree $n_k$, with $F^k_d(y_{k_0})y_{k,p'_k}=F^k_{p'_k}(y_{k_0})y_{k,d}$ and
$F^k_{p'_k}(w_{k_0})\neq0$. As in the proof of Lemma~\ref{cv-origin-reference-orbit}, the polynomials
$F^k_d(x_{k_0})x_{k,p'_k}-F^k_{p'_k}(x_{k_0})x_{k,d}$ lie in the kernel of $y$, so
$F^k_d(G(v)_{k_0})\,G(v)_{k,p'_k}=F^k_{p'_k}(G(v)_{k_0})\,G(v)_{k,d}$ for $v$ near $v_0$.

Let $\Phi(u)_{k,d}=F^k_d(u)/F^k_{p'_k}(u)$ for $u\in\mathbf C^{I_{k_0}}$; by homogeneity
$F^k_d(f(v))=F^k_d(G(v)_{k_0})\,G(v)_{k_0,p}^{-n_k}$, so $F^k_{p'_k}(f(v_0))\neq0$ and $\Phi$ is holomorphic near $f(v_0)$.
For $v$ near $v_0$, $G(v)_{k,p'_k}$ and $F^k_{p'_k}(f(v))$ are nonzero, and the relations give
\[
 G(v)_{k,d}/G(v)_{k,p'_k}=\Phi(f(v))_{k,d}.
\]
The left side, as a function of $v$, gives the affine coordinates of $\lambda_x$ in the chart $p'$, whose derivative at
$v_0$ is injective by Definition~\ref{cv-origin-analytic}(3). By the chain rule it is $d\Phi\circ df$ at $v_0$, so the
derivative of $f$ at $v_0$ is injective.
\end{proof}
% Lean: ImmBlock.immersion_block: hE.holo, hE.mem, small_points_coord_ratio_bound at lam x x for the chart p'
% (Function.update r k0 (ix p)), form_mem for the relations, aeval_mul_of_isHomogeneous (hscale), hfac
% (affCoord p' (G v) = Psi (f v) near v0), hE.immersion and ker_bot_of_factor.

\begin{lemma}\label{cv-origin-analytic-inverse}
Let $X_1$ and $X_2$ be complex manifolds of the same dimension, let $\pi:X_1\to X_2$ be a bijective
holomorphic map and $\psi:X_2\to X_1$ a holomorphic map with $\psi\circ\pi$ the identity, let $V$ be a normed
complex vector space and $\Phi:X_1\to V$ a map. Let $x\in X_2$, and let $\chi_1$ and $\chi_2$ be charts at
$\psi(x)$ and at $x$. If $\Phi\circ\chi_1^{-1}$ is differentiable at $\chi_1(\psi(x))$ with injective derivative,
then $\Phi\circ\psi\circ\chi_2^{-1}$ has injective derivative at $\chi_2(x)$.
\end{lemma}
\leanref{Development/AutomorphicSpectrum/Proofs/GeometryProofs/HolomorphicInverseChartValueStatement.lean}{GlobalTheory.T.holomorphic_inverse_chart_value}
\leanref{Development/AutomorphicSpectrum/Proofs/GeometryProofs/HolomorphicInverseChartValue.lean}{GlobalTheory.CompanionTheory.Ar.ImmersionBiholo.ker_fderiv_comp_biholo}
\begin{proof}
As $\pi$ is bijective and $\psi\circ\pi$ is the identity, $\pi\circ\psi$ is the identity. Put
$\sigma=\chi_2\circ\pi\circ\chi_1^{-1}$ and $\rho=\chi_1\circ\psi\circ\chi_2^{-1}$, holomorphic near
$\chi_1(\psi(x))$ and $\chi_2(x)$. Near $\chi_2(x)$, $\sigma\circ\rho$ is the identity, so by the chain rule
$d\sigma\circ d\rho$ is the identity and $d\rho$ is injective at $\chi_2(x)$. Near $\chi_2(x)$,
$\Phi\circ\psi\circ\chi_2^{-1}=(\Phi\circ\chi_1^{-1})\circ\rho$, whose derivative at $\chi_2(x)$ is the
composite of two injective maps.
\end{proof}

\subsection{Field extensions and coordinates}\label{subsection-cv-field-extensions-coordinates}

\begin{lemma}\label{cv-origin-extension-tate}
In Situation~\ref{cv-function-situation}, let $\Lambda$ be the lattice of the PEL datum, let $G$ be an
$\mathcal O_E$-semi-abelian scheme over a scheme $S$ (Definition~\ref{cv-model-toroidal-semiabelian}(2)), let $k$ be an
algebraically closed field of characteristic zero and $x:\operatorname{Spec}k\to S$ a point such that the fibre $G_x$ is
proper over $k$, and let $\sigma:k\to k'$ be an embedding into an algebraically closed field. For a point $g$ of $G_x$
with values in $k$ let $g_\sigma$ be its image with values in $k'$. If $((\zeta'_m),(\iota_m))$ is a Tate system of $G$
at $x$ (Definition~\ref{cv-toroidal-moduli-object}(2)), then $((\sigma(\zeta'_m)),(\iota'_m))$ with
$\iota'_m(a)=\iota_m(a)_\sigma$ is a Tate system of $G$ at $x\circ\operatorname{Spec}\sigma$.
\end{lemma}
\leanref{Development/AutomorphicSpectrum/Proofs/GeometryProofs/TateIsoBaseChangeStatement.lean}{GlobalTheory.T.tate_iso_base_change}
\leanref{Development/WeilBound/CoreStatements/HodgeTatePoints/MgTate.lean}{PositivityTheory.HtPt.Mg.torsion_desc}
\begin{proof}
The map $g\mapsto g_\sigma$ is a homomorphism compatible with the actions of $\mathcal O_E$, and it is injective: a point
with values in $k$ is determined by its image over $k'$, as $\sigma$ is injective. So $\sigma(\zeta'_m)$ is a primitive
$m$-th root of unity, $\sigma(\zeta'_{mm'})^{m'}=\sigma(\zeta'_m)$, the maps $\iota'_m$ are $\mathcal O_E$-linear with
kernel $m\Lambda$, $m'\iota'_{mm'}=\iota'_m$, and $\iota'_m$ takes values in the $m$-torsion.

\emph{The $m$-torsion.} Let $T\subset G_x$ be the fibre of multiplication by $m$ over the zero section, a closed subscheme
locally of finite type over $k$ whose points with values in an extension of $k$ are the points of $G_x$ of order dividing
$m$. Its points with values in $k$ are the $\iota_m(a)$; they are finitely many, as $\Lambda/m\Lambda$ is finite. The
scheme $T$ is Jacobson \cite[Tag 02J6]{Sta}, and the residue field of each of its closed points is finite over $k$
\cite[Tag 01TB]{Sta}, hence equal to $k$; so the closed points of $T$ are the images of its points with values in $k$, and
they are finitely many. Let $h$ be a point of $G_x$ with values in $k'$ and $mh=0$, a point of $T$ with image $t$. As $T$
is Jacobson, the closed points of the closure of $\{t\}$ are dense in it \cite[Tag 005U]{Sta}; they are finitely many, so
they form a closed set, which is the closure of $\{t\}$, and $t$ is a closed point. So the residue field of $t$ is $k$, and
$h$ factors through the point $\operatorname{Spec}k\to T$ at $t$ \cite[Tag 01J9]{Sta}, by an embedding $k\to k'$ which is
$\sigma$, as $h$ lies over $\operatorname{Spec}\sigma$. So $h=g_\sigma$ for a point $g$ of $T$ with values in $k$, that is,
$h=\iota_m(a)_\sigma=\iota'_m(a)$ for some $a\in\Lambda$.

\emph{The pairings.} As $G_x$ is proper over $k$, the commutator pairings of the theta groups of the powers of the line
bundle are compatible with the base change along $\sigma$: $e_m(g_\sigma,g'_\sigma)=\sigma(e_m(g,g'))$. So
$e_m(\iota'_m(a),\iota'_m(b))=\sigma((\zeta'_m)^{\langle a,b\rangle})=\sigma(\zeta'_m)^{\langle a,b\rangle}$.
\end{proof}
% Lean: Relay.tate_iso_base_change: Mg.extPt with extPt_inj (epi_spMap) for the injectivity, Mg.torsion_desc for the
% torsion (the kernel Kn of multiplication by n, locally of finite type, with finitely many sections by
% Module.finite_of_fg_torsion on La/nLa; exists_section_of_finite with LocallyOfFiniteType.jacobsonSpace,
% pointEquivClosedPoint, JacobsonSpace.closure_inter_closedPoints_eq_closure and pointOfClosedPoint), and
% Mg.weilIs_pull_of_proper (with proper_of_hom) for the pairings.

\begin{lemma}\label{cv-origin-extension-level}
In the setting of Definition~\ref{cv-origin-hodge-component}, let $\Lambda$ be the lattice of the PEL datum, let $K$ be a
field, $y_0$ a point of $X$ with values in $K$, and $N\ge1$. Then there are a finite extension $\Omega'$ of $K$, a unit
$z\in\Omega'^\times$ and points $\mathrm{ev}_0(a)$ of $\mathcal A$ at $\xi'_0=j\circ y_0\circ\operatorname{Spec}(\Omega'\supset K)$
for $a\in\Lambda$, such that for every algebraically closed field $k$ and embedding $\sigma:\Omega'\to k$, $\mathcal A$ has a
Tate system $((\zeta'_m),(\iota_m))$ at $\xi'_0\circ\operatorname{Spec}\sigma$ (Definition~\ref{cv-toroidal-moduli-object}(2))
with $\zeta'_N=\sigma(z)$ and $\iota_N(a)$ the image of $\mathrm{ev}_0(a)$ for every $a\in\Lambda$.
\end{lemma}
\leanref{Development/AutomorphicSpectrum/Proofs/GeometryProofs/ExistsFullLevelFiniteExtensionStatement.lean}{GlobalTheory.T.exists_full_level_finite_extension}
\leanref{Development/AutomorphicSpectrum/Proofs/GeometryProofs/ExistsFullLevelFiniteExtension.lean}{GlobalTheory.CompanionTheory.Ar.XR.LevelPt.exists_of_tateIso}
\begin{proof}
The field $K$ has characteristic zero, as $F_c$ embeds in it through $y_0$. Let $\bar K$ be an algebraic closure of $K$ and
$\bar y=y_0\circ\operatorname{Spec}(\bar K\supset K)$. The abelian scheme $A$ of $X$ is, with its level structures, an
object at level $l$ (Definition~\ref{cv-toroidal-moduli-object}(4)), so it has a Tate system $((\zeta'_m),(\iota_m))$ at
$\bar y$.

\emph{Fields of definition.} Let $A_{y_0}$ be the fibre of $A$ over $y_0$, a scheme of finite type over $K$, as $A$ is
proper over $X$. Let $g$ be a point of $A_{y_0}$ with values in $\bar K$, and $\operatorname{Spec}B$ an affine open
subscheme of $A_{y_0}$ that contains its image. Then $g$ is given by a homomorphism $B\to\bar K$ of $K$-algebras
\cite[Tag 01I1]{Sta}. The algebra $B$ is generated over $K$ by finitely many elements, whose images are algebraic over $K$
and generate a finite extension $F_g\subset\bar K$ of $K$ that contains the image of $B$; so $g$ is the image of a point of
$A_{y_0}$ with values in $F_g$. The group $\Lambda/N\Lambda$ is finite, as $\Lambda$ is finitely generated; let
$T\subset\Lambda$ be a finite set of representatives. Let $\Omega'\subset\bar K$ be the field generated over $K$ by
$\zeta'_N$ and the fields $F_{\iota_N(t)}$ for $t\in T$; it is finite over $K$. Put $z=\zeta'_N$.

\emph{The points.} For $t\in T$ let $e(t)$ be a point of $A_{y_0}$ with values in $\Omega'$ whose image over $\bar K$ is
$\iota_N(t)$. For $a\in\Lambda$ with representative $t\in T$, let $\mathrm{ev}_0(a)$ be the image of $e(t)$ under the
morphism $A\to\mathcal A$ over $j$, a point of $\mathcal A$ at $\xi'_0$. As $\iota_N$ is additive with kernel $N\Lambda$,
$\iota_N(a)=\iota_N(t)$, which is the image of $e(t)$ over $\bar K$.

\emph{The Tate systems.} Let $k$ be an algebraically closed field and $\sigma:\Omega'\to k$ an embedding. As $\bar K$ is
algebraic over $\Omega'$, $\sigma$ extends to an embedding $\bar\sigma:\bar K\to k$. The fibre of $A$ at $\bar y$ is
proper, so Lemma~\ref{cv-origin-extension-tate}, applied to $A$, $\bar y$ and $\bar\sigma$, gives a Tate system of $A$ at
$\bar y\circ\operatorname{Spec}\bar\sigma=y_0\circ\operatorname{Spec}(\Omega'\supset K)\circ\operatorname{Spec}\sigma$ with
$\zeta'_N=\bar\sigma(z)=\sigma(z)$ and $\iota_N(a)$ the image under $\bar\sigma$ of the old $\iota_N(a)$, that is, the image
of $e(t)$ under $\sigma$. The morphism $A\to\mathcal A$ of $\mathcal O_E$-semi-abelian schemes over $j$ identifies the
points of $A$ at a point $x$ of $X$ with values in a field with the points of $\mathcal A$ at $j\circ x$, by the universal
property of the fibre product, compatibly with the group laws, the actions of $\mathcal O_E$ and the commutator pairings
of the theta groups. So this Tate system of $A$ is one of $\mathcal A$ at $\xi'_0\circ\operatorname{Spec}\sigma$, with the
same roots of unity, and its $\iota_N(a)$ is the image of $\mathrm{ev}_0(a)$.
\end{proof}
% Lean: XR.exists_fullLevel_finite: CharZero K through the map from c.F, Y.oF.geom at the algebraic closure of K (the
% Tate system of the object oF at level l), then LevelPt.exists_of_tateIso: exists_reps (representatives of La/N La),
% exists_finite_subfield (an affine open containing the point, Algebra.FiniteType and the adjoined images of the
% generators), the compositum Omega of Fz = K(zeta_N) and the F tau in the algebraic closure, epi_spMap for hover,
% exists_extend (IsAlgClosed.lift), T.tate_iso_base_change (Lemma cv-origin-extension-tate, with Y.oF.proper) and
% T.tate_iso_pullback_equiv = XR.tateIsoAt_move along Y.m, with I.ker and I.add for a - u(r(a)) in N La.

\begin{lemma}\label{cv-origin-extension-charts}
Let $X$ be a scheme with a closed immersion $i:X\to\mathbf P^N_{\mathbf Z}$, where $\mathbf P^N_{\mathbf Z}$ has the
homogeneous coordinates $T_0,\dots,T_N$, let $\Omega$ be a field and $\xi:\operatorname{Spec}\Omega\to X$ a morphism with
image $\eta$, such that every point of $X$ specializes $\eta$ and $X$ satisfies condition (1) of
Definition~\ref{cv-origin-segre-coordinates}. Then there is a nonzero $t\in\Omega^{N+1}$ such that for every $d$, $i(\eta)$
lies in $D_+(T_d)$ if and only if $t_d\neq0$, and then the homomorphism $\mathbf Z[T_j/T_d]\to\Omega$ of $i\circ\xi$ maps
$T_j/T_d$ to $t_j/t_d$. The scheme $X$ has the coordinates $t$ (Definition~\ref{cv-origin-segre-coordinates}), and the
points in the chart $d$ of $t$ are those of $i^{-1}(D_+(T_d))$.
\end{lemma}
\leanref{Development/WeilBound/CoreStatements/HodgeTatePoints/MfCoords.lean}{PositivityTheory.HtPt.Mf.xr_coords}
\leanref{Development/WeilBound/CoreStatements/HodgeTatePoints/MfCoords.lean}{PositivityTheory.HtPt.Mf.coordTie_of_all}
\begin{proof}
\emph{The point $t$.} The open subschemes $D_+(T_c)=\operatorname{Spec}\mathbf Z[T_j/T_c]$ cover $\mathbf P^N_{\mathbf Z}$.
Let $i(\eta)\in D_+(T_c)$, and let $t_j$ be the image of $T_j/T_c$ under the homomorphism of $i\circ\xi$; so $t_c=1$ and
$t\neq0$. Inside $D_+(T_c)$, the open subset $D_+(T_c)\cap D_+(T_d)=D_+(T_cT_d)$ is the standard open of $T_d/T_c$
\cite[Tag 00JP]{Sta}, \cite[Tag 01M4]{Sta}. As the residue field of $i(\eta)$ embeds in $\Omega$, $i(\eta)$ lies in this
standard open, that is, in $D_+(T_d)$, if and only if $t_d\neq0$, and then $T_j/T_d=(T_j/T_c)(T_d/T_c)^{-1}$ maps to
$t_j/t_d$.

\emph{Charts.} For $t_d\neq0$ let $X_d=i^{-1}(D_+(T_d))$, an affine open subscheme of $X$ whose ring $B_d$ is a quotient of
$\mathbf Z[T_j/T_d]$ \cite[Tag 01IN]{Sta}. Let $P$ be a point of $X$. If $i(P)\in D_+(T_d)$, then the open set $X_d$
contains $P$, hence its generization $\eta$, so $t_d\neq0$; the germs at $P$ of the images of the $T_j/T_d$ map to
$t_j/t_d$, which are therefore regular at $P$, and $P$ lies in the chart $d$ of $t$. Conversely, let $P$ lie in the chart
$d$ of $t$, and let $P\in X_e$. The germ at $P$ of the image of $T_d/T_e$ maps to $t_d/t_e$, and $t_e/t_d$ is the image of
a germ at $P$; the product of the two germs maps to $1$, so it is $1$ by condition (1), and the image of $T_d/T_e$ is a unit
at $P$. So $P$ lies in the standard open of $T_d/T_e$ in $X_e$, which is $i^{-1}(D_+(T_eT_d))\subset X_d$. Hence the points
in the chart $d$ are those of the affine open subscheme $X_d$, and as the $X_e$ cover $X$, every point lies in a chart:
conditions (2) and (3) of Definition~\ref{cv-origin-segre-coordinates}.

\emph{Rings.} Let $\varphi:B_d\to\Omega$ be the homomorphism through the local ring at $\eta$. It maps the images of the
$T_j/T_d$ to the $t_j/t_d$, so $\varphi(B_d)=\mathbf Z[t_j/t_d]$, and the elements of $\varphi(B_d)$ are regular at every
point of $X_d$. Conversely, let $a\in\Omega$ be regular at every point of $X_d$, and let $I$ be the ideal of the $b\in B_d$
with $\varphi(b)a\in\varphi(B_d)$. For a maximal ideal $\mathfrak m$ of $B_d$, the local ring at its point is
$(B_d)_{\mathfrak m}$ \cite[Tag 01HV]{Sta}, so $a=\varphi(r)/\varphi(s)$ with $s\notin\mathfrak m$, and $s\in I$. So $I$
lies in no maximal ideal, $I=B_d$ and $a\in\varphi(B_d)$. This is condition (4).
\end{proof}
% Lean: HtPt.Mf.xr_coords: the open W of the points that specialize genPt xi (isClopen_setOf_genPt_specializes; here
% all of X), exists_point (a chart containing the image and the values of the ratios), and coordTie_of_all: inChart_iff
% (genMap_germ_eq, valR_pullRatio, the unit germ and basicOpen_pullRatio), the affine opens i^{-1}(chart c),
% exists_valR_eq_iff_mem_adjoin_ratio (app_surjective of the closed immersion, aeval_ratio_surjective) and
% exists_valR_eq_iff_forall_inChart_mem_regAt (integrality from xr_normal and PrimeSpectrum.iInf_localization_eq_bot,
% for which the book uses the ideal of denominators). The Lean states it for a noetherian coefficient ring o and without
% the hypothesis that every point specializes eta; the book states the case o = Z that is used.

\begin{lemma}\label{cv-origin-extension-normal}
Let $Y$ be a scheme with a closed immersion $i:Y\to\mathbf P^N_{\mathbf Z}$, let $\Omega$ be a field of characteristic zero
and $\xi:\operatorname{Spec}\Omega\to Y$ a morphism such that $\Omega$ is finite over the residue field of $Y$ at the image
of $\xi$. Then there are a finite family $x$ in $\Omega$, a scheme $M$ with a morphism $\xi_M:\operatorname{Spec}\Omega\to M$
of whose image every point of $M$ is a specialization, such that $M$ has the coordinates $x$
(Definition~\ref{cv-origin-segre-coordinates}) and its local rings are integrally closed in $\Omega$, and a finite morphism
$\pi:M\to Y$ with $\pi\circ\xi_M=\xi$.
\end{lemma}
\leanref{Development/AutomorphicSpectrum/Proofs/GeometryProofs/ExistsFiniteNormalModelOfResidueFiniteStatement.lean}{GlobalTheory.T.exists_finite_normal_model_of_residue_finite}
\leanref{Development/WeilBound/CoreStatements/HodgeTatePoints/MfCoords.lean}{PositivityTheory.HtPt.Mf.xr_norm}
\begin{proof}
\emph{The point.} As in the proof of Lemma~\ref{cv-origin-extension-charts}, there is a nonzero $t\in\Omega^{N+1}$ such
that $i\circ\xi$ factors through $D_+(T_d)=\operatorname{Spec}\mathbf Z[T_j/T_d]$ exactly when $t_d\neq0$, and then maps
$T_j/T_d$ to $t_j/t_d$. Let $F\subset\Omega$ be the field generated by the ratios $t_j/t_d$ with $t_d\neq0$, and fix $d$ with
$t_d\neq0$. The local ring of $\mathbf P^N_{\mathbf Z}$ at the image $Q$ of $i\circ\xi$ is a localization of
$\mathbf Z[T_j/T_d]$ \cite[Tag 01HV]{Sta}, so its image in $\Omega$ lies in $F$. As $i$ is a closed immersion, the local
ring of $Y$ at the image of $\xi$ is a quotient of that of $\mathbf P^N_{\mathbf Z}$ at $Q$ \cite[Tag 01IN]{Sta}; so the
residue field $\kappa$ of $Y$ there maps into $F$. As $\Omega$ is finite over $\kappa$, it is finite over $F$.

\emph{The closure.} Let $C\subset\mathbf P^N_{\mathbf Z}$ be the scheme theoretic image of $i\circ\xi$ \cite[Tag 01R7]{Sta}.
Its intersection with $D_+(T_e)$ is the scheme theoretic image of the restriction of $i\circ\xi$ \cite[Tag 01R8]{Sta}: it is
empty if $t_e=0$, and if $t_e\neq0$ it is $\operatorname{Spec}A_e$ with $A_e=\mathbf Z[t_j/t_e]\subset\Omega$, the image of
$\mathbf Z[T_j/T_e]$ \cite[Tag 056A]{Sta}. So $C$ is of finite type over $\mathbf Z$, the image of $\operatorname{Spec}\Omega$
is the zero ideal of each $A_e$, and the residue field of $C$ there is the fraction field $F$ of $A_d$. As $i\circ\xi$
factors through the closed subscheme $Y$, $C$ is a closed subscheme of $Y$; let $l:C\to Y$ be the closed immersion.

\emph{The normalization.} Let $\nu:M\to C$ be the normalization of $C$ in $\operatorname{Spec}\Omega$, with the
factorization $\xi_M:\operatorname{Spec}\Omega\to M$ of $\operatorname{Spec}\Omega\to C$ \cite[Tag 035H]{Sta}. Over
$\operatorname{Spec}A_e$ it is $\operatorname{Spec}A'_e$, where $A'_e$ is the integral closure of $A_e$ in $\Omega$
\cite[Tag 035F]{Sta}. As $\operatorname{Spec}\Omega$ is normal and irreducible, $M$ is integral and normal and $\xi_M$ is
dominant \cite[Tag 035L]{Sta}; so every point of $M$ specializes the image of $\xi_M$, its generic point. The scheme $C$ is
of finite type over $\mathbf Z$, hence Nagata \cite[Tag 035B]{Sta}, and $\Omega$ is finite over the residue field $F$ of $C$
at the image of $\operatorname{Spec}\Omega$; so $\nu$ is finite \cite[Tag 0AVK]{Sta}.

\emph{Local rings.} The local ring of $M$ at its generic point is the fraction field of $A'_d$, which is $\Omega$: every
$a\in\Omega$ is algebraic over $F$, the fraction field of $A_d$, so $ba$ is integral over $A_d$ for some nonzero $b\in A_d$.
The local rings of the integral scheme $M$ embed in its function field, so $M$ satisfies condition (1) of
Definition~\ref{cv-origin-segre-coordinates}; as $M$ is normal, they are integrally closed in their fraction field
$\Omega$.

\emph{Coordinates.} The morphism $\nu$ is finite, hence projective \cite[Tag 0B3I]{Sta}, and $C\to\operatorname{Spec}\mathbf Z$
is projective \cite[Tag 01W9]{Sta}; so $M$ is projective over $\mathbf Z$ \cite[Tag 0C4P]{Sta}, and, as the structure sheaf
of $\operatorname{Spec}\mathbf Z$ is ample \cite[Tag 01PS]{Sta}, $M$ has a closed immersion into some
$\mathbf P^{N'}_{\mathbf Z}$ \cite[Tag 087S]{Sta}. Lemma~\ref{cv-origin-extension-charts}, applied to $M$ and $\xi_M$, gives a
finite family $x$ in $\Omega$ such that $M$ has the coordinates $x$.

\emph{The morphism.} Put $\pi=l\circ\nu$, a finite morphism, as $\nu$ is finite and $l$ is a closed immersion. The
composites of $\pi\circ\xi_M$ and of $\xi$ with $i$ are both $\operatorname{Spec}\Omega\to C\to\mathbf P^N_{\mathbf Z}$, and
$i$ is a monomorphism \cite[Tag 01L7]{Sta}; so $\pi\circ\xi_M=\xi$.
\end{proof}
% Lean: HtPt.Mf.xr_norm: exists_point, hrange (genMap_point_mem: the stalk of PSpace at the image maps into
% TupleNorm.ratioField; i.stalkMap_surjective) and hFD; TupleNorm.Cl = (point t).image, Nm = its relative normalization
% (Scheme.Hom.normalization of clPoint), nmPoint and toCl; essFiniteType_of_finiteDimensional, isIntegral_Nm, IsDominant
% from Scheme.Hom.ker_toNormalization, all_of_isDominant, isIso_stalkMap_nmPoint and isCptGen_of_isDominant; xr_coords
% (Lemma cv-origin-extension-charts) with TupleNorm.isProjOver_str; xr_normal for the local rings; the closed immersion
% l = IsClosedImmersion.lift i clι; TupleNorm.isFinite_toCl; and cancel_mono i. The Lean proves the finiteness of the
% normalization (isFinite_fromNormalization_of_genericPoints) and the projectivity of Nm (exists_iso_Cl, through the
% integral closure of the cone of t) itself; the book cites Stacks for both.

\begin{lemma}\label{cv-origin-extension-cover}
In the setting of Definition~\ref{cv-origin-hodge-component}, let $\Lambda$ be the lattice of the PEL datum, let $e$ be
the generic point of an irreducible component of $\mathcal X$, and let $N\ge1$. Then there are:
a field $\Omega'$ of characteristic zero; a scheme $M$ with a morphism $\xi':\operatorname{Spec}\Omega'\to M$ of whose image
every point of $M$ is a specialization, such that $M$ has the coordinates $x$ for a finite family $x$ in $\Omega'$
(Definition~\ref{cv-origin-segre-coordinates}) and its local rings are integrally closed in $\Omega'$; a finite morphism
$f:M\to\mathcal X$ that maps the image of $\xi'$ to $e$, with $f\circ\xi'=j\circ y$ for a point $y$ of $X$ with values in
$\Omega'$; an $\mathcal O_E$-semi-abelian scheme $\mathcal A_M$ over $M$ with a morphism $\mathcal A_M\to\mathcal A$ of
$\mathcal O_E$-semi-abelian schemes over $f$ (Definition~\ref{cv-model-toroidal-semiabelian}(2)); and a unit
$z\in\Omega'^\times$ and points $\mathrm{ev}(a)$ of $\mathcal A_M$ at $\xi'$ for $a\in\Lambda$, such that:
\begin{enumerate}
\item for every algebraically closed field $k$ and embedding $\sigma:\Omega'\to k$, $\mathcal A_M$ has a Tate system
  $((\zeta'_m),(\iota_m))$ at $\xi'\circ\operatorname{Spec}\sigma$ (Definition~\ref{cv-toroidal-moduli-object}(2)) with
  $\zeta'_N=\sigma(z)$ and $\iota_N(a)$ the image of $\mathrm{ev}(a)$ for every $a\in\Lambda$;
\item every point of order dividing $N$ of $\mathcal A_M$ at $\xi'$ with values in an extension of $\Omega'$ is defined
  over $\Omega'$;
\item $z$ is a unit at every point of $M$.
\end{enumerate}
\end{lemma}
\leanref{Development/AutomorphicSpectrum/Statements/GeometryStatements/LevelCoverSt.lean}{GlobalTheory.CompanionTheory.Ar.XR.LevelCoverSt}
\leanref{Development/AutomorphicSpectrum/Proofs/GeometryProofs/LevelCoverStHoldsStatement.lean}{GlobalTheory.T.LevelCoverSt_holds}
\begin{proof}
Let $K$ be the residue field of $\mathcal X$ at $e$. Lemma~\ref{cv-origin-change-factor}, applied to the canonical
morphism $\operatorname{Spec}K\to\mathcal X$, gives a point $y_0$ of $X$ with values in $K$ whose composite with $j$ is that
morphism; $K$ has characteristic zero, as $F_c$ embeds in it.

\emph{The field.} Lemma~\ref{cv-origin-extension-level}, applied to $K$, $y_0$ and $N$, gives a finite extension
$\Omega'$ of $K$, a unit $z\in\Omega'^\times$ and points $\mathrm{ev}_0(a)$ of $\mathcal A$ at
$\xi'_0=j\circ y_0\circ\operatorname{Spec}(\Omega'\supset K)$ that satisfy (1) for $\mathcal A$ at $\xi'_0$. The field
$\Omega'$ has characteristic zero, as $K$ has.

\emph{The model.} The scheme $\mathcal X$ is a closed subscheme of a projective space over $\mathcal O_c$, hence projective
over $\mathcal O_c$ \cite[Tag 01W9]{Sta}, and
$\operatorname{Spec}\mathcal O_c$ is finite, hence projective, over $\mathbf Z$ \cite[Tag 0B3I]{Sta}; so $\mathcal X$ is
projective over $\mathbf Z$ \cite[Tag 0C4P]{Sta}, and, as the structure sheaf of $\operatorname{Spec}\mathbf Z$ is ample
\cite[Tag 01PS]{Sta}, it has a closed immersion into some $\mathbf P^{N'}_{\mathbf Z}$ \cite[Tag 087S]{Sta}. The morphism
$\xi'_0$ is the canonical morphism $\operatorname{Spec}K\to\mathcal X$ at $e$ followed by $\operatorname{Spec}(\Omega'\supset
K)$, and $\Omega'$ is finite over the residue field $K$. So Lemma~\ref{cv-origin-extension-normal}, applied to
$\mathcal X$ and $\xi'_0$, gives a finite family $x$ in $\Omega'$, a scheme $M$ with a morphism
$\xi':\operatorname{Spec}\Omega'\to M$ of whose image every point of $M$ is a specialization, such that $M$ has the
coordinates $x$ and its local rings are integrally closed in $\Omega'$, and a finite morphism $f:M\to\mathcal X$ with
$f\circ\xi'=\xi'_0$. Then $f$ maps the image of $\xi'$ to $e$, and $f\circ\xi'=j\circ y$ with
$y=y_0\circ\operatorname{Spec}(\Omega'\supset K)$.

\emph{The abelian scheme.} Let $\mathcal A_M=M\times_{\mathcal X}\mathcal A$, with the action of $\mathcal O_E$ and the line
bundle pulled back from $\mathcal A$; it is an $\mathcal O_E$-semi-abelian scheme over $M$, and the projection
$\mathcal A_M\to\mathcal A$ is a morphism of $\mathcal O_E$-semi-abelian schemes over $f$. For a point $x'$ of $M$ with
values in a field, the points of $\mathcal A_M$ at $x'$ are the points of $\mathcal A$ at $f\circ x'$, by the universal
property of the fibre product, compatibly with the group laws, the actions of $\mathcal O_E$ and the commutator pairings
of the theta groups of the line bundles. So a Tate system of $\mathcal A$ at $f\circ x'$ is a Tate system of
$\mathcal A_M$ at $x'$, with the same roots of unity. Let $\mathrm{ev}(a)$ be the point of $\mathcal A_M$ at $\xi'$ that
corresponds to $\mathrm{ev}_0(a)$. For $k$ algebraically closed and $\sigma:\Omega'\to k$, the Tate system of $\mathcal A$ at
$\xi'_0\circ\operatorname{Spec}\sigma$ given by (1) is one of $\mathcal A_M$ at $\xi'\circ\operatorname{Spec}\sigma$ with
$\zeta'_N=\sigma(z)$ and $\iota_N(a)$ the image of $\mathrm{ev}(a)$. This is (1).

\emph{Torsion points.} Let $\sigma:\Omega'\to k$ be an extension of fields and $g$ a point of $\mathcal A_M$ at
$\xi'\circ\operatorname{Spec}\sigma$ with $Ng=0$. Let $\tau:k\to\bar k$ be an algebraic closure. By (1) for $\tau\sigma$, the
image of $g$ over $\bar k$, of order dividing $N$, is $\iota_N(a)$ for some $a\in\Lambda$, the image of $\mathrm{ev}(a)$. A
point with values in $k$ is determined by its image over $\bar k$, as $k\to\bar k$ is injective, so $g$ is the image of
$\mathrm{ev}(a)$. This is (2).

\emph{The unit.} By (1) for an algebraic closure of $\Omega'$, $z$ maps to a primitive $N$-th root of unity, so $z^N=1$.
So $z$ and $z^{-1}=z^{N-1}$ are integral over every local ring of $M$, and these are integrally closed in $\Omega'$; so $z$
is a unit at every point of $M$. This is (3).
\end{proof}
% Lean: XR.levelCover: generic_point_normal_coords (the point y0; the book takes it from Lemma cv-origin-change-factor),
% exists_full_level_finite_extension (Lemma cv-origin-extension-level), T.projective_over_restrict_int (the book cites
% Stacks), exists_finite_normal_model_of_residue_finite (Lemma cv-origin-extension-normal) with
% finite_descResidueField_spMap_comp for its hypothesis, and
% full_level_generic_cover = XR.exists_full_level_over: ModFam.exists_sabPull for A_M,
% full_level_structure_pullback = XR.isFullLevelAt_move (through XR.tateIsoAt_move: LiePull.exists_pushEquiv and
% weilIs_grpPullback_of_comm'), full_level_structure_properties = HtPt.Mg.level_props (FullTorsion is full_desc, with
% epi_spMap), and unitAt_of_pow_eq_one = Flat.ZUnit.of_pow_of_normal.

\begin{lemma}\label{cv-origin-change-extension}
In the setting of Definition~\ref{cv-origin-hodge-component}, let $\Omega_1$ be a field of characteristic zero
and $\xi_1:\operatorname{Spec}\Omega_1\to\mathcal X$ a morphism whose image is the generic point of an
irreducible component, and let $S$ be a finite set of primes. Then there is a finite extension $\Omega$ of
$\Omega_1$ such that for every $p\in S$, $\Omega$ contains a primitive $p^4$-th root of unity, and every
point of order dividing $p^4$ of $\mathcal A$ at $\operatorname{Spec}\Omega\to\mathcal X$ with values in an
extension of $\Omega$ is defined over $\Omega$.
\end{lemma}
\leanref{Development/AutomorphicSpectrum/Proofs/GeometryProofs/ExistsFiniteExtensionPrimesStatement.lean}{GlobalTheory.T.exists_finite_extension_primes}
\begin{proof}
Put $N=\prod_{p\in S}p^4$, let $e$ be the image of $\xi_1$, the generic point of an irreducible component, and let $K$ be
the residue field of $\mathcal X$ at $e$.

\emph{The level cover.} Lemma~\ref{cv-origin-extension-cover}, applied to $e$ and $N$, gives a field $\Omega'$ of
characteristic zero, a scheme $M$ with a point $\xi':\operatorname{Spec}\Omega'\to M$ and coordinates, a finite morphism
$f:M\to\mathcal X$ that maps the image $\eta'$ of $\xi'$ to $e$, and a morphism $\mathcal A_M\to\mathcal A$ of
$\mathcal O_E$-semi-abelian schemes over $f$, such that every point of order dividing $N$ of $\mathcal A_M$ at $\xi'$ with
values in an extension of $\Omega'$ is defined over $\Omega'$.

\emph{The residue fields.} A morphism from the spectrum of a field to a scheme is a point with an embedding of its residue
field, as in the proof of Lemma~\ref{cv-origin-change-factor}. So $\xi_1$ and $f\circ\xi'$, both with image $e$, are given
by embeddings $\iota_1:K\to\Omega_1$ and $\iota_2:K\to\Omega'$. As $M$ has coordinates, the local ring of $M$ at $\eta'$
maps bijectively to $\Omega'$ (Definition~\ref{cv-origin-segre-coordinates}(1)), so $\Omega'$ is the residue field of $M$ at
$\eta'$. The morphism $f$ is finite, so its fibre over $e$ is the spectrum of a finite $K$-algebra, and $\Omega'$ is finite
over $K$ through $\iota_2$.

\emph{The field $\Omega$.} Let $\Omega_2$ be the quotient of $\Omega_1\otimes_K\Omega'$ by a maximal ideal, a field with
embeddings $a:\Omega_1\to\Omega_2$ and $b:\Omega'\to\Omega_2$ that agree on $K$; it is finite over $\Omega_1$, as $\Omega'$ is
finite over $K$. Let $\Omega$ be a splitting field over $\Omega_2$ of the $N$-th cyclotomic polynomial; as $\Omega_2$ has
characteristic zero, $\Omega$ contains a primitive $N$-th root of unity $\zeta$, and it is finite over $\Omega_2$, hence over
$\Omega_1$. For $p\in S$, $\zeta^{N/p^4}$ is a primitive $p^4$-th root of unity in $\Omega$.

\emph{Torsion points.} As $a$ and $b$ agree on $K$, the point $\operatorname{Spec}\Omega\to\operatorname{Spec}\Omega_1\to\mathcal X$
is $\operatorname{Spec}\Omega\to\operatorname{Spec}\Omega'\to M\to\mathcal X$. Let $k$ be an extension of $\Omega$ and $g$ a point of
$\mathcal A$ at $\operatorname{Spec}k\to\mathcal X$ of order dividing $p^4$. As $\mathcal A_M$ is the pullback of $\mathcal A$
along $f$, $g$ is the image of a point $g'$ of $\mathcal A_M$ at $\operatorname{Spec}k\to M$, and $g'$ has order dividing
$p^4$, hence $N$, as the map on points is an injective homomorphism. So $g'$ is the image of a point of $\mathcal A_M$ at
$\xi'$, and $g$ is the image of the corresponding point of $\mathcal A$ at $\operatorname{Spec}\Omega\to\mathcal X$: it is
defined over $\Omega$.
\end{proof}
% Lean: ChgField.exists_field: T.LevelCoverSt_holds = XR.levelCover (Lemma cv-origin-extension-cover), exists_factor
% for xi_1 and M.xi >> f, HodgeExists.finite_of_dominates (LocallyQuasiFinite and a quasi-finite map of fields) for the
% finiteness of i_2, exists_compositum (a maximal ideal of the tensor product), exists_root (CyclotomicField), and the
% transfer fullTorsion_ext (j.comp b) of HodgeExists.fullTorsion_push (exists_lift_torsion through m.cart).

\begin{lemma}\label{cv-origin-coords-join-map}
Let $\iota:\Omega\to\Omega'$ be an embedding of fields, and let $X$ and $X'$ be schemes with morphisms
$\xi:\operatorname{Spec}\Omega\to X$ and $\xi':\operatorname{Spec}\Omega'\to X'$, with images $\eta$ and $\eta'$.
\begin{enumerate}
\item Let $f:X'\to X$ be a morphism with $f\circ\xi'=\xi\circ\operatorname{Spec}\iota$, and let $P'$ be a point of $X'$
  that specializes $\eta'$. Then $f(P')$ specializes $\eta$, and the map $\mathcal O_{X,f(P')}\to\mathcal O_{X',P'}\to\Omega'$
  is the map $\mathcal O_{X,f(P')}\to\Omega$ of Definition~\ref{cv-origin-segre-coordinates} followed by $\iota$; so
  $\iota(a)$ is regular at $P'$ for every $a\in\Omega$ regular at $f(P')$. If $\mathcal O_{X',P'}\to\Omega'$ is injective,
  then every $a$ regular at $f(P')$ such that $\iota(a)$ is a unit at $P'$ is a unit at $f(P')$.
\item Let $x=(x_j)_{j\in J}$ be a finite family in $\Omega$ such that $X$ has the coordinates $x$
  (Definition~\ref{cv-origin-segre-coordinates}), and write $\iota(x)=(\iota(x_j))_{j\in J}$. Suppose that every point of
  $X'$ specializes $\eta'$, that $X'$ satisfies condition (1) of Definition~\ref{cv-origin-segre-coordinates}, and that
  every point of $X'$ lies in a chart of $\iota(x)$. Then there is a morphism $f:X'\to X$ with
  $f\circ\xi'=\xi\circ\operatorname{Spec}\iota$, and a point $P'$ of $X'$ lies in the chart $j$ of $\iota(x)$ if and only if
  $f(P')$ lies in the chart $j$ of $x$.
\end{enumerate}
\end{lemma}
\leanref{Development/AutomorphicSpectrum/Proofs/GeometryProofs/GenMapStalkMap.lean}{GlobalTheory.CompanionTheory.Ar.RegAtDom.regAt_dominates}
\leanref{Development/AutomorphicSpectrum/Proofs/GeometryProofs/DominatesOfCoordsRegular.lean}{GlobalTheory.CompanionTheory.Ar.CoordMap.exists_dominating_map}
\begin{proof}
(1) The closed point of $\operatorname{Spec}\Omega'$ maps to that of $\operatorname{Spec}\Omega$, so $f(\eta')=\eta$, as
$f\circ\xi'=\xi\circ\operatorname{Spec}\iota$. As $f$ is continuous, $f(P')$ specializes $f(\eta')=\eta$. Put $Q=f(P')$, let
$\varphi:\mathcal O_{X,Q}\to\Omega$ and $\varphi':\mathcal O_{X',P'}\to\Omega'$ be the maps of
Definition~\ref{cv-origin-segre-coordinates}, and let $f^\sharp:\mathcal O_{X,Q}\to\mathcal O_{X',P'}$ be the map of $f$ on
local rings. The morphism $\xi$ is $\operatorname{Spec}\varphi$ followed by the canonical morphism
$\operatorname{Spec}\mathcal O_{X,Q}\to X$, and likewise $\xi'$; and the canonical morphism
$\operatorname{Spec}\mathcal O_{X',P'}\to X'$ followed by $f$ is $\operatorname{Spec}f^\sharp$ followed by the canonical
morphism $\operatorname{Spec}\mathcal O_{X,Q}\to X$. So $\operatorname{Spec}(\varphi'\circ f^\sharp)$ and
$\operatorname{Spec}(\iota\circ\varphi)$, followed by $\operatorname{Spec}\mathcal O_{X,Q}\to X$, are $f\circ\xi'$ and
$\xi\circ\operatorname{Spec}\iota$, which are equal. Let $\operatorname{Spec}B$ be an affine open neighbourhood of $Q$. It
contains the image of $\operatorname{Spec}\mathcal O_{X,Q}$, which consists of generizations of $Q$ \cite[Tag 01J7]{Sta}, so
this morphism factors through $\operatorname{Spec}B$, and there it is given by the localization $B\to\mathcal O_{X,Q}$
\cite[Tag 01HV]{Sta}. The two composites $B\to\mathcal O_{X,Q}\to\Omega'$ give the same morphism
$\operatorname{Spec}\Omega'\to\operatorname{Spec}B$, so they are equal \cite[Tag 01I1]{Sta}, and as $\mathcal O_{X,Q}$ is a
localization of $B$, $\varphi'\circ f^\sharp=\iota\circ\varphi$. This is the compatibility. For $s\in\mathcal O_{X,Q}$ it
gives $\iota(\varphi(s))=\varphi'(f^\sharp(s))$, which is regular at $P'$.

Now let $\varphi'$ be injective, and let $a=\varphi(s)$ with $\iota(a)$ a unit at $P'$ and $\iota(a)^{-1}=\varphi'(t)$. Then
$\varphi'(f^\sharp(s)t)=\iota(a)\iota(a)^{-1}=1$, so $f^\sharp(s)t=1$, and $f^\sharp(s)$ is a unit. As $f^\sharp$ is a local
homomorphism, $s$ is a unit, so $a\neq0$ and $a^{-1}=\varphi(s^{-1})$ is regular at $Q$: $a$ is a unit at $Q$.

(2) In this part $j$ runs over the indices with $x_j\neq0$; then $\iota(x_j)\neq0$, and conversely. Let $U_j$ be the affine
open subscheme of $X$ of condition (3) of Definition~\ref{cv-origin-segre-coordinates}, and $V_j$ the set of points of $X'$
in the chart $j$ of $\iota(x)$, which is open by Lemma~\ref{cv-origin-segre-regular}(1). By condition (1), every element of
$\Omega$ is regular at $\eta$ and every element of $\Omega'$ at $\eta'$, so $\eta\in U_j$ and $\eta'\in V_j$. So $\xi$ and
$\xi'$ factor through morphisms $\theta_j:\operatorname{Spec}\Omega\to U_j$ and $\zeta_j:\operatorname{Spec}\Omega'\to V_j$.
The open subscheme $V_j$, with $\zeta_j$, has the local rings of $X'$ with the same maps to $\Omega'$.

\emph{A morphism on each chart.} By Lemma~\ref{cv-origin-segre-regular}(2), applied to $X'$ and $V_j$, the map
$\Gamma(V_j,\mathcal O_{X'})\to\Omega'$ is injective, and its image consists of the elements regular at every point of
$V_j$. Let $s\in\Gamma(U_j,\mathcal O_X)$ have image $a\in\Omega$. Then $a$ is regular at every point of $U_j$, so
$a\in\mathbf Z[x_i/x_j]$ by condition (4), and $\iota(a)$ lies in $\mathbf Z[\iota(x_i)/\iota(x_j)]$, whose elements are
regular at every point of $V_j$. So the map $\Gamma(U_j,\mathcal O_X)\to\Omega$ followed by $\iota$ factors through a ring
homomorphism $\Gamma(U_j,\mathcal O_X)\to\Gamma(V_j,\mathcal O_{X'})$, which, as $U_j$ is affine, is the map on global
sections of a morphism $g_j:V_j\to U_j$ \cite[Tag 01I1]{Sta}. The morphisms $g_j\circ\zeta_j$ and
$\theta_j\circ\operatorname{Spec}\iota$ from $\operatorname{Spec}\Omega'$ to $U_j$ have the same map on global sections, the
map $\Gamma(U_j,\mathcal O_X)\to\Omega$ followed by $\iota$, so they are equal, by the same reference. Let $G_j:V_j\to X$ be
$g_j$ followed by the inclusion of $U_j$; then $G_j\circ\zeta_j=\xi\circ\operatorname{Spec}\iota$.

\emph{Agreement on overlaps.} Let $P'\in V_j\cap V_k$. The point $G_k(P')$ lies in $U_k$, so in the chart $k$ of $x$, and
$x_j/x_k$ is regular at $G_k(P')$. By (1), applied to $G_k$ and $\zeta_k$, $\iota(x_j/x_k)$ is regular at $P'$. It is
nonzero, and its inverse $\iota(x_k)/\iota(x_j)$ is regular at $P'$, as $P'\in V_j$; so $\iota(x_j/x_k)$ is a unit at $P'$,
and by (1) $x_j/x_k$ is a unit at $G_k(P')$. So every $x_i/x_j=(x_i/x_k)(x_j/x_k)^{-1}$ is regular at $G_k(P')$, and
$G_k(P')\in U_j$. The restrictions of $G_j$ and $G_k$ to $W=V_j\cap V_k$ are therefore two morphisms $W\to U_j$, which agree
after composition with the factorization of $\xi'$ through $W$, as both composites are $\xi\circ\operatorname{Spec}\iota$.
The scheme $W$ is reduced, as its local rings embed in $\Omega'$, and $U_j$ is affine, hence separated. The locus where the
two morphisms agree is a closed subscheme of $W$ \cite[Tag 01KM]{Sta} through which $\xi'$ factors, so it contains $\eta'$
and hence every point of $W$; as $W$ is reduced, the identity of $W$ factors through it \cite[Tag 0356]{Sta}, and the two
morphisms are equal.

\emph{Gluing.} Every point of $X'$ lies in a chart of $\iota(x)$, so the $V_j$ cover $X'$, and the $G_j$ glue to a
morphism $f:X'\to X$ \cite[Tag 01JB]{Sta}. As $\xi'$ factors through $V_j$, $f\circ\xi'=G_j\circ\zeta_j=\xi\circ\operatorname{Spec}\iota$.
If $f(P')$ lies in the chart $j$ of $x$, then the ratios $\iota(x_i)/\iota(x_j)=\iota(x_i/x_j)$ are regular at $P'$ by (1),
so $P'$ lies in the chart $j$ of $\iota(x)$. Conversely, if $P'$ lies in the chart $j$ of $\iota(x)$, then $P'\in V_j$, and
$f(P')=G_j(P')\in U_j$ lies in the chart $j$ of $x$.
\end{proof}
% Lean: (1) is RegAtDom.regAt_dominates: map_genPt, genMap_stalkMap (spec_genMap and Scheme.SpecMap_stalkMap_fromSpecStalk;
% the Lean cancels the monomorphism X.fromSpecStalk and uses Spec.map_injective, which the book does through an affine
% neighbourhood), and isUnit_of_map_unit for the unit. (2) is CoordMap.exists_dominating_map, for a field embedding i and
% coefficient ring o (the book takes o = Z): tie.affine for U_j, isOpen_chart for V_j, exists_lift, exists_chart_map =
% exists_map_affine (val_injective, exists_val_eq), glue_compat (inChart_of_inChart, ext_of_land with Mathlib's
% ext_of_isDominant), Scheme.OpenCover.glueMorphisms, and inChart_pull for the charts. Lemmas cv-origin-segre-regular(3)
% and cv-origin-segre-map state the case i = id. Size: about 370 lines.

\begin{lemma}\label{cv-origin-coords-join}
Let $X$ be a scheme, $\Omega_1$ a field and $\xi:\operatorname{Spec}\Omega_1\to X$ a morphism of whose image every point
of $X$ is a specialization, and let $x=(x_j)_{j\in J}$ be a finite family in $\Omega_1$ such that $X$ has the coordinates
$x$ (Definition~\ref{cv-origin-segre-coordinates}). Let $\iota:\Omega_1\to\Omega$ be an embedding of fields and
$w=(w_l)_{l\in L}$ a finite family in $\Omega$ with some $w_l$ nonzero, such that $\Omega$ is generated as a field by
$\iota(\Omega_1)$ and the ratios $w_l/w_{l'}$ with $w_{l'}\neq0$. Then there are a scheme $X'$ with a morphism
$\xi':\operatorname{Spec}\Omega\to X'$ of whose image every point of $X'$ is a specialization, such that $X'$ has the
coordinates $(\iota(x_j)w_l)_{(j,l)\in J\times L}$, and a morphism $f:X'\to X$ with
$f\circ\xi'=\xi\circ\operatorname{Spec}\iota$.
\end{lemma}
\leanref{Development/AutomorphicSpectrum/Proofs/GeometryProofs/ExistsJoinModelDominatesStatement.lean}{GlobalTheory.T.exists_join_model_dominates}
\begin{proof}
Let $\eta$ be the image of $\xi$. It lies in a chart $c$ of $x$, by condition (2) of
Definition~\ref{cv-origin-segre-coordinates}, and $w_l\neq0$ for some $l$, so $\iota(x_c)w_l\neq0$. Let $A\subset\Omega$ be
the ring generated by the ratios $\iota(x_j)w_m/(\iota(x_c)w_l)$. It contains $\iota(x_j/x_c)=\iota(x_j)w_l/(\iota(x_c)w_l)$,
hence $\iota(\mathbf Z[x_j/x_c])$, and $w_m/w_l=\iota(x_c)w_m/(\iota(x_c)w_l)$. Every element of $\Omega_1$ is regular at
$\eta$, by condition (1), so it is a quotient of two elements of $\mathbf Z[x_j/x_c]$, by
Lemma~\ref{cv-origin-snc-charts}(1) applied to $x$ as a point with one block. So $\iota(\Omega_1)$ and the ratios
$w_m/w_{m'}=(w_m/w_l)/(w_{m'}/w_l)$ lie in the fraction field of $A$, which is therefore $\Omega$.
Lemma~\ref{cv-origin-segre-model}, applied to the family $(\iota(x_j)w_m)$ and the index $(c,l)$, gives a scheme $X'$ and
a morphism $\xi'$ of whose image every point of $X'$ is a specialization, such that $X'$ has the coordinates
$(\iota(x_j)w_m)$. Every point $P'$ of $X'$ lies in a chart $(j,m)$ of this family; then $x_j\neq0$, and every
$\iota(x_i/x_j)=\iota(x_i)w_m/(\iota(x_j)w_m)$ is regular at $P'$.

So every point of $X'$ lies in a chart of the family $(\iota(x_j))_{j\in J}$. As $X'$ has
coordinates, it satisfies condition (1) of Definition~\ref{cv-origin-segre-coordinates}, and every point of $X'$
specializes the image of $\xi'$. So Lemma~\ref{cv-origin-coords-join-map}(2), applied to $\iota$, $X$ with $\xi$ and $x$,
and $X'$ with $\xi'$, gives a morphism $f:X'\to X$ with $f\circ\xi'=\xi\circ\operatorname{Spec}\iota$.
\end{proof}
% Lean: join_exists: the chart c of x at the generic point (M.tie.cover), hmap and hgen for the ring A, hiF from
% regAt_chart_localization and regAt_local_order_fraction at the generic point, hle and hfrac (the fraction field),
% projective_model_nonempty (Lemma cv-origin-segre-model), join_coords_centre_chart = HtPt.Mb.join_charts (the charts of
% x at the points of X', ctrSet), and existsUnique_dominating_map with the embedding i, whose existence part is
% dominates_of_coords_regular = CoordMap.exists_dominating_map (Lemma cv-origin-coords-join-map(2)); its uniqueness
% part is not used here.

\begin{lemma}\label{cv-origin-coords-normal}
Let $\Omega$ be a field, let $X$ be a scheme with a morphism $\xi:\operatorname{Spec}\Omega\to X$ of whose image every
point of $X$ is a specialization, and let $x=(x_j)_{j\in J}$ be a finite family in $\Omega$ such that $X$ has the
coordinates $x$ (Definition~\ref{cv-origin-segre-coordinates}). For a point $P$ of $X$ let $\mathcal O_P\subset\Omega$ be
the image of its local ring and $B_P$ the integral closure of $\mathcal O_P$ in $\Omega$. Then there are a finite family $y$
in $\Omega$, a scheme $X'$ with a morphism $\xi':\operatorname{Spec}\Omega\to X'$ of whose image every point of $X'$ is a
specialization, such that $X'$ has the coordinates $y$ and its local rings are integrally closed in $\Omega$, and a finite
morphism $f:X'\to X$ with $f\circ\xi'=\xi$, such that for every point $P$ of $X$:
\begin{enumerate}
\item for every point $P'$ of $X'$ with $f(P')=P$, the local ring of $X'$ at $P'$ is $(B_P)_{\mathfrak m}$ for a maximal
  ideal $\mathfrak m$ of $B_P$, as subrings of $\Omega$;
\item for every maximal ideal $\mathfrak m$ of $B_P$ there is exactly one point $P'$ of $X'$ with $f(P')=P$ whose local
  ring is $(B_P)_{\mathfrak m}$.
\end{enumerate}
\end{lemma}
\leanref{Development/AutomorphicSpectrum/Proofs/GeometryProofs/ExistsNormalisationStatement.lean}{GlobalTheory.T.exists_normal_dominating_model}
\leanref{Development/AutomorphicSpectrum/Proofs/GeometryProofs/ExistsNormalisation.lean}{GlobalTheory.CompanionTheory.Ar.NormEx.exists_normalisation}
\leanref{Development/AutomorphicSpectrum/Proofs/GeometryProofs/IntegralClosureChartSpanStatement.lean}{GlobalTheory.T.integral_closure_chart_span}
\leanref{Development/AutomorphicSpectrum/Proofs/GeometryProofs/JoinCoordsChartIntegralStatement.lean}{GlobalTheory.T.join_coords_chart_integral}
\begin{proof}
Let $J_0$ be the set of $c\in J$ with $x_c\neq0$. For $c\in J_0$ put $A_c=\mathbf Z[x_j/x_c:j\in J]\subset\Omega$, let $A'_c$
be the integral closure of $A_c$ in $\Omega$, and let $U_c$ be the affine open subscheme of the chart $c$ of $x$.

\emph{Finiteness of $A'_c$.} Let $c\in J_0$. By condition (1) of Definition~\ref{cv-origin-segre-coordinates}, every element
of $\Omega$ is regular at the image of $\xi$, which lies in $U_c$; so it is a quotient of two elements of $A_c$, by
Lemma~\ref{cv-origin-snc-charts}(1) applied to $x$ as a point with one block, and $\Omega$ is the fraction field of $A_c$.
The domain $A_c$ is generated over $\mathbf Z$ by finitely many elements, and $\mathbf Z$ is universally Japanese
\cite[Tag 0335]{Sta}; so $A'_c$ is a finite $A_c$-module.

\emph{Common generators.} Let $c,c'\in J_0$. If $b\in A_c$, then $(x_c/x_{c'})^mb\in A_{c'}$ for all large $m$: this holds
for the generators $x_j/x_c$, as $(x_c/x_{c'})(x_j/x_c)=x_j/x_{c'}$, and for the integers, and it is preserved by sums and
products. If $a$ is integral over $A_c$, with $a^r+b_1a^{r-1}+\dots+b_r=0$ and $b_i\in A_c$, choose $m_0$ with
$(x_c/x_{c'})^{m_0}b_i\in A_{c'}$ for every $i$; multiplying the equation by $(x_c/x_{c'})^{m_0r}$ shows that
$(x_c/x_{c'})^{m_0}a$ is integral over $A_{c'}$, and so is $(x_c/x_{c'})^ma$ for $m\ge m_0$, as $x_c/x_{c'}\in A_{c'}$. For
each $c\in J_0$ choose generators $g_{c,1},\dots,g_{c,N_c}$ of the $A_c$-module $A'_c$, and choose $d\ge0$ such that
$(x_c/x_{c'})^dg_{c,k}$ is integral over $A_{c'}$ for all $c,c'\in J_0$ and all $k$. Let $L$ be the disjoint union of $J$
and the set of pairs $(c,k)$, and define $z\in\Omega^L$ by $z_j=x_j^d$ for $j\in J$ and $z_{(c,k)}=x_c^dg_{c,k}$. For
$c'\in J_0$, the elements $z_j/x_{c'}^d=(x_j/x_{c'})^d$ lie in $A_{c'}$, the elements
$z_{(c,k)}/x_{c'}^d=(x_c/x_{c'})^dg_{c,k}$ are integral over $A_{c'}$, and $g_{c',k}=z_{(c',k)}/x_{c'}^d$. So $A'_{c'}$ is the
$A_{c'}$-module generated by the $z_l/x_{c'}^d$, $l\in L$.

\emph{The scheme.} Let $y=(x_jz_l)_{(j,l)\in J\times L}$. For $c\in J_0$, $y_{(c,c)}=x_c^{d+1}\neq0$ and
$y_{(j,l)}/y_{(c,c)}=(x_j/x_c)(z_l/x_c^d)$, while $x_j/x_c=y_{(j,c)}/y_{(c,c)}$ and $z_l/x_c^d=y_{(c,l)}/y_{(c,c)}$. So
$\mathbf Z[y_q/y_{(c,c)}:q\in J\times L]$ is the $A_c$-algebra generated by the $z_l/x_c^d$, which is $A'_c$, by the
previous paragraph. Fix $c_0\in J_0$. Every element of $\Omega$ is a quotient of two elements of $A_{c_0}\subset A'_{c_0}$, so
Lemma~\ref{cv-origin-segre-model}, applied to $y$ and the index $(c_0,c_0)$, gives a scheme $X'$ with a morphism
$\xi':\operatorname{Spec}\Omega\to X'$ of whose image every point of $X'$ is a specialization, such that $X'$ has the
coordinates $y$.

\emph{Charts of $X'$.} Let $c\in J_0$, and let $P'$ be a point of $X'$ at which every $x_j/x_c$ is regular. The point $P'$
lies in a chart $(j,l)$ of $y$; then $z_l\neq0$, and every $z_k/z_l=y_{(j,k)}/y_{(j,l)}$ is regular at $P'$. The element
$v=z_l/x_c^d$ lies in $A'_c$, so it is integral over $A_c$, whose elements are regular at $P'$, and $v^{-1}=z_c/z_l$ is
regular at $P'$. A nonzero element $v$ integral over a ring $R$ with $v^{-1}\in R$ lies in $R$: dividing an equation
$v^r+b_1v^{r-1}+\dots+b_r=0$ by $v^{r-1}$ gives $v=-(b_1+b_2v^{-1}+\dots+b_rv^{1-r})$. So $v$ is regular at $P'$, and so are
the $z_k/x_c^d=(z_k/z_l)v$ and the $y_q/y_{(c,c)}$: $P'$ lies in the chart $(c,c)$ of $y$. Conversely, every $x_j/x_c=
y_{(j,c)}/y_{(c,c)}$ is regular at a point in the chart $(c,c)$. A point in a chart $(c,l)$ of $y$ has $x_c\neq0$, and every
$x_j/x_c=y_{(j,l)}/y_{(c,l)}$ is regular there; so every point of $X'$ lies in a chart $(c,c)$ with $c\in J_0$.

\emph{The morphism.} So every point of $X'$ lies in a chart of the family $x$, and $X'$, which has coordinates, satisfies
condition (1) of Definition~\ref{cv-origin-segre-coordinates}. Lemma~\ref{cv-origin-coords-join-map}(2), applied with the
identity of $\Omega$, gives a morphism $f:X'\to X$ with
$f\circ\xi'=\xi$ such that $P'$ lies in the chart $c$ of $x$ if and only if $f(P')$ does. By the previous paragraph,
$f(P')\in U_c$ if and only if $P'$ lies in the chart $(c,c)$ of $y$: $f^{-1}(U_c)$ is the affine open subscheme $X'_{(c,c)}$
of this chart.

\emph{Normality.} Let $P'$ be a point of $X'$, in the chart $(c,c)$ of $y$, and let $t\in\Omega$ be integral over the local
ring $\mathcal O_{X',P'}$. By Lemma~\ref{cv-origin-snc-charts}(1), applied to $y$ as a point with one block, every element
of $\mathcal O_{X',P'}$ is $u/v$ with $u,v\in A'_c$ and $v$ a unit at $P'$. Multiplying an equation of integral dependence of
$t$ over $\mathcal O_{X',P'}$ by a power of the product of the denominators of its coefficients gives $v\in A'_c$, a unit at
$P'$, with $vt$ integral over $A'_c$. As $A'_c$ is integral over $A_c$, $vt$ is integral over $A_c$, so $vt\in A'_c$ and
$t=(vt)v^{-1}$ is regular at $P'$. So the local rings of $X'$ are integrally closed in $\Omega$.

\emph{Finiteness of $f$.} Let $c\in J_0$. By Lemma~\ref{cv-origin-segre-regular}(2) and condition (4) of
Definition~\ref{cv-origin-segre-coordinates}, the rings of $U_c$ and of $f^{-1}(U_c)=X'_{(c,c)}$ map isomorphically onto
$A_c$ and $A'_c$, and as $f\circ\xi'=\xi$ the map of $f$ on these rings is the inclusion $A_c\subset A'_c$, which is finite.
As the $U_c$ cover $X$, $f$ is finite \cite[Tag 01WI]{Sta}.

\emph{The points over $P$.} Let $P$ be a point of $X$, in the chart $c$ of $x$, and put $R=\mathcal O_P$ and $B=B_P$. The
elements of $A_c$ are regular at $P$, so $A_c\subset R$ and $A'_c\subset B$. By Lemma~\ref{cv-origin-snc-charts}(1), every
element of $R$ is $u/v$ with $u,v\in A_c$ and $v$ a unit at $P$; so, as for the normality, for every $t\in B$ there is
$v\in A_c$, a unit of $R$ and hence of $B$, with $vt\in A'_c$. Hence two ideals of $B$ with the same intersection with
$A'_c$ are equal. Let $\mathfrak m$ be a maximal ideal of $B$. As $B$ is integral over $R$, the ring $R/(\mathfrak m\cap R)$
is a field \cite[Tag 00GR]{Sta}, so $\mathfrak m\cap R$ is the maximal ideal of $R$: an element of $R$ lies in
$\mathfrak m$ if and only if it is not a unit of $R$.

(1) Let $f(P')=P$ and $D=\mathcal O_{X',P'}$; then $P'$ lies in the chart $(c,c)$ of $y$. By
Lemma~\ref{cv-origin-coords-join-map}(1), with the identity of $\Omega$, $R\subset D$, and an element of $R$ that is a unit
at $P'$ is a unit at $P$. As $D$ is integrally closed in $\Omega$, $B\subset D$. Let $\mathfrak q$ be the set of elements of
$B$ that are not units of $D$, a prime ideal of $B$. Every element of $D$ is $u/v$ with $u,v\in A'_c\subset B$ and $v$ a
unit of $D$, so $v\notin\mathfrak q$; hence $D\subset B_{\mathfrak q}$, and $B_{\mathfrak q}\subset D$ as the elements of
$B\setminus\mathfrak q$ are units of $D$. So $D=B_{\mathfrak q}$. An element of $R$ lies in $\mathfrak q$ if and only if it
is not a unit of $R$; so $\mathfrak q\cap R$ is the maximal ideal of $R$, and $B/\mathfrak q$ is a domain integral over the
field $R/(\mathfrak q\cap R)$, hence a field \cite[Tag 00GS]{Sta}: $\mathfrak q$ is maximal.

(2) Let $\mathfrak m$ be a maximal ideal of $B$ and $\mathfrak p'=\mathfrak m\cap A'_c$. By
Lemma~\ref{cv-origin-snc-charts}(2), applied to $X'$ and the chart $(c,c)$ of $y$, whose ring of ratios is $A'_c$, there is
exactly one point $P'$ in this chart at which the elements of $A'_c$ that are units are those outside $\mathfrak p'$. The
point $f(P')$ lies in the chart $c$ of $x$. For $u\in A_c$: $u$ is a unit at $P$ if and only if $u\notin\mathfrak m$, if
and only if $u\notin\mathfrak p'$, if and only if $u$ is a unit at $P'$, if and only if $u$ is a unit at $f(P')$, by
Lemma~\ref{cv-origin-coords-join-map}(1). So $P$ and $f(P')$ are points of $U_c$ at which the same elements of $A_c$ are
units, and $f(P')=P$ by Lemma~\ref{cv-origin-snc-charts}(2). By (1), the local ring of $X'$ at $P'$ is $B_{\mathfrak q}$,
where $\mathfrak q$ is the set of elements of $B$ that are not units at $P'$; so $\mathfrak q\cap A'_c=\mathfrak p'=
\mathfrak m\cap A'_c$, and $\mathfrak q=\mathfrak m$. If $P''$ is a point with $f(P'')=P$ whose local ring is
$B_{\mathfrak m}$, then $P''$ lies in the chart $(c,c)$, and the elements of $A'_c$ that are units at $P''$ are those
outside $\mathfrak m$, that is, outside $\mathfrak p'$; so $P''=P'$.
\end{proof}
% Lean: NormEx.exists_normalisation, for o = Z and Omega' = Omega (as used in ChangeModel.exists_normal_model):
% chart_ring_integral_closure_finite (ChartFinite.finite_integralClosure_chRing: exists_eq_div_chRing and
% T.integral_closure_finite, which the book takes from Stacks 0335), integral_closure_chart_span =
% IntClGens.exists_twisted_gens (twist_mem, twist_integral with scaleRoots), join_chart and
% chart_ring_integral_fraction, projective_model_nonempty (Lemma cv-origin-segre-model), join_coords_chart_integral =
% ChNorm.join_chart_normalisation (chRing_join_iff, mem_of_isIntegral_of_inv_mem, join_coords_centre_chart), hcov,
% existsUnique_dominating_map (Lemma cv-origin-coords-join-map(2)), normal_of_charts (regAt_chart_localization,
% integral_over_localization_scaled = ClearDenominators.exists_mul_isIntegral, int_trans), dominating_model_chart_compat
% = ChartFin.isFinite_of_charts (chart_step, finite_of_values, HasAffineProperty.iff_of_iSup_eq_top), and
% normAt_of_charts: local_subring_integral_closure_max = Domination.eq_localisation_integralClosure (eq_ofPrime,
% isMaximal_of_comap), chart_prime_existsUnique_point (Lemma cv-origin-snc-charts(2)), maximal_ideal_membership_of_local,
% prime_mem_iff_not_unit_localization, ideal_ext, and fibre_iff with regular_functions_two_components. The Lean's chart_over
% and fibre_iff argue through centre_set_dominating_point and regular_functions_two_components = Centre.Sep.eq_of_dominated;
% here both points lie in one chart of x, where these reduce to Sep.eq_of_dominated_of_inChart, and the book takes the
% charts of f(P') from the chart statement of CoordMap.exists_dominating_map. Size: about 1,300 lines.

\begin{definition}\label{cv-origin-coords-crossings}
Let $\Omega$ be a field of characteristic zero, let $A$ be a subalgebra of $\Omega$ over $\mathbf Q$, let $h\in A$ and let
$e\ge0$. The element $h$ \emph{has normal crossings of dimension $e$ on} $A$ if there are $f_1,\dots,f_m$ and
$a_1,\dots,a_m$ in $A$ with $\sum_la_lf_l=1$ such that for each $l$ with $f_l\neq0$ the algebra $A[1/f_l]$ is smooth over
$\mathbf Q$ and has elements $z_1,\dots,z_e$, whose differentials form a basis of its module of differentials over
$\mathbf Q$, and an integer $n\le e$, such that a prime ideal of $A[1/f_l]$ contains $h$ if and only if it contains one of
$z_1,\dots,z_n$.
\end{definition}
\leanref{Development/AutomorphicSpectrum/Proofs/GeometryProofs/SncChartPoint.lean}{GlobalTheory.CompanionTheory.Ar.SncRes.SncChart}

\begin{lemma}\label{cv-origin-coords-snc-chart}
Let $\Omega$ be a field of characteristic zero, let $Z$ be a reduced scheme, smooth of relative dimension $e$ over
$\mathbf Q$, with a morphism $\zeta:\operatorname{Spec}\Omega\to Z$ whose image is dense, and let $T\subset Z$ be a subset.
Suppose that $Z$ is covered by finitely many affine open subschemes $V$, on each of which there are functions
$z_1,\dots,z_e$, whose differentials form a basis of $\Omega^1_{V/\mathbf Q}$, and an integer $k\le e$ such that $V\cap T$
is the set of points of $V$ at which one of $z_1,\dots,z_k$ vanishes. Let $W$ be an affine open subscheme of $Z$ that
contains the image of $\zeta$, let $A\subset\Omega$ be the image of the ring of $W$ under the homomorphism $\varphi$ of
$\zeta$, and let $\sigma$ be a function on $W$ whose zero set is $W\cap T$. Then $\varphi(\sigma)$ has normal crossings of
dimension $e$ on $A$ (Definition~\ref{cv-origin-coords-crossings}).
\end{lemma}
\leanref{Development/AutomorphicSpectrum/Proofs/GeometryProofs/SncChartOfSmoothCoordsStatement.lean}{GlobalTheory.T.snc_chart_of_smooth_coords}
\leanref{Development/AutomorphicSpectrum/Proofs/GeometryProofs/SncChartOfSmoothCoords.lean}{GlobalTheory.CompanionTheory.Ar.SncRes.sncChart_of_coords}
\begin{proof}
\emph{A cover.} Every point of $W$ lies in some $V$ of the cover, and it has an affine open neighbourhood that is a standard
open both of $W$ and of $V$ \cite[Tag 01IW]{Sta}. So the functions $u$ on $W$ whose standard open $D(u)$ is a standard open
of some $V$ generate the unit ideal of the ring of $W$: there are such $u_1,\dots,u_m$ and functions $c_1,\dots,c_m$ on $W$
with $\sum_lc_lu_l=1$. Put $f_l=\varphi(u_l)$ and $a_l=\varphi(c_l)$; then $\sum_la_lf_l=1$.

\emph{The localizations.} Let $f_l\neq0$, and let $D=D(u_l)$, a standard open $D(v)$ of some $V$ of the cover. As
$f_l\neq0$, $D$ contains the image of $\zeta$. The ring of $D$ is that of $W$ with $u_l$ inverted \cite[Tag 01HV]{Sta}, so
its image under $\zeta$ is $A[1/f_l]$. The map from it to $\Omega$ is injective: as the image of $\zeta$ is dense, $Z$ is
irreducible and the image of $\zeta$ is its generic point, whose prime in the ring of the reduced affine scheme $D$ is $0$.
So $A[1/f_l]$ is isomorphic to the ring of $D$, which is smooth over $\mathbf Q$, as $Z$ is. The ring of $D$ is also that of
$V$ with $v$ inverted, so the differentials of the restrictions of $z_1,\dots,z_e$ form a basis of its module of
differentials \cite[Tag 00RT]{Sta}; let $z'_1,\dots,z'_e$ be their images in $A[1/f_l]$, and $n=k$.

\emph{The primes.} A prime ideal $\mathfrak q$ of $A[1/f_l]$ is the prime of a point $t$ of $D$, and an element lies in
$\mathfrak q$ if and only if the corresponding function vanishes at $t$. So $\mathfrak q$ contains $\varphi(\sigma)$ if and
only if $t\in T$, if and only if one of $z_1,\dots,z_k$ vanishes at $t$, that is, if and only if $\mathfrak q$ contains one
of $z'_1,\dots,z'_k$.
\end{proof}
% Lean: SncRes.sncChart_of_coords: exists_cover_basicOpen (exists_basicOpen_le_affine_inter and
% IsAffineOpen.self_le_iSup_basicOpen_iff), then snc_on_basicOpen: range_valMap_basicOpen (the image A[1/f]),
% valMap_injective, smooth_secAlgebra, coords_of_formallyEtale for the localization and for the isomorphism with A[1/f]
% (KaehlerDifferential.tensorKaehlerEquivOfFormallyEtale), exists_point_of_prime, and the comparison of zero sets.

\begin{lemma}\label{cv-origin-coords-snc}
Let $\Omega$ be a field of characteristic zero, let $K$ and $I_k$ for $k\in K$ be finite sets, let
$y\in\Omega^{\sqcup_kI_k}$ be a point whose kernel $P$ in $R=\mathbf Z[x_{k,i}:k\in K,\ i\in I_k]$
(Definition~\ref{lit-multiprojective}(4)) is an arithmetic subvariety of dimension $e+1$, and let $\beta\in R$ be
multihomogeneous with $\beta\notin P$. Then there are $m\ge1$, a block $y_2\in\Omega^m$ with a nonzero entry, and a
multihomogeneous ideal $J$ of the polynomial ring $R_2$ in the variables of the blocks of the point $z=(y,y_2)$, not
contained in the kernel $P_2$ of $z$, such that:
\begin{enumerate}
\item the ratios of entries of one block of $z$ generate the same subfield of $\Omega$ as those of $y$;
\item $J$ satisfies condition (3) of Definition~\ref{cv-origin-reference} for $z$, with $e$ in place of $D$;
\item every point of the zero set of $P_2$ in $\hat C$ at which $\beta$ vanishes lies in the zero set of $J$.
\end{enumerate}
\end{lemma}
\leanref{Development/AutomorphicSpectrum/Proofs/GeometryProofs/AdjoinCoordsArithVar.lean}{GlobalTheory.CompanionTheory.Ar.SncRes.exists_snc_model}
\begin{proof}
\emph{A chart.} The prime $P$ is relevant, so for each $k$ there is $c_k\in I_k$ with $x_{k,c_k}\notin P$, that is,
$y_{k,c_k}\neq0$. Index the Segre products of $y$ by $j\in\prod_kI_k$, as $u_j=\prod_ky_{k,j_k}$, let $T_j$ be the
corresponding homogeneous coordinates of a projective space $\mathbf P^N$, and for a form $G$ in the $T_j$ let
$\tilde G\in R$ be $G$ evaluated at the monomials $\prod_kx_{k,j_k}$. Then $u_c\neq0$, and the ring
$A=\mathbf Q[u_j/u_c]$ is $\mathbf Q[y_{k,i}/y_{k,c_k}]$: $u_j/u_c=\prod_ky_{k,j_k}/y_{k,c_k}$, and $y_{k,i}/y_{k,c_k}=u_j/u_c$
for the $j$ that agrees with $c$ except that $j_k=i$. Its fraction field $F$ is the field generated by the ratios of
entries of one block of $y$.

\emph{Smoothness on an open set.} The ring $R/P$ is the subring of $\Omega[t_k:k\in K]$ generated by the $y_{k,i}t_k$
(Definition~\ref{lit-multiprojective}(4)), of dimension $e+1+|K|$. Its fraction field is $F(y_{k,c_k}t_k:k\in K)$, of
transcendence degree $\operatorname{trdeg}_{\mathbf Q}F+|K|$ over $\mathbf Q$, and, as in the proof of
Lemma~\ref{cv-origin-weight-birational}, the dimension formula \cite[Tag 02JU]{Sta} shows that $\dim R/P$ is one more than
this transcendence degree. So $F$ has transcendence degree $e$ over $\mathbf Q$, and $\operatorname{Spec}A$, an integral
scheme of finite type over $\mathbf Q$, has dimension $e$ at each of its points \cite[Tag 00P1]{Sta}. It is geometrically
reduced over $\mathbf Q$ \cite[Tag 020I]{Sta}, so it has a dense open subscheme smooth over $\mathbf Q$
\cite[Tag 056V]{Sta}. This contains a standard open $D(s)$, $s\neq0$, with $A[1/s]$ standard smooth over $\mathbf Q$
\cite[Tag 00TA]{Sta}, with module of differentials free of rank $e$ \cite[Tag 02G1]{Sta}. Call such an algebra smooth of
relative dimension $e$; a localization $A[1/(st)]$ with $t\neq0$ is again smooth of relative dimension $e$.

\emph{A form.} Every element of $A$ is $G(u/u_c)/m$ for a nonzero integer $m$ and a form $G$ in the $T_j$ with integer
coefficients: for the generator $u_j/u_c$ take $G=T_j$; for a rational number take a constant form and its denominator;
for sums and products, multiply by powers of $T_c$, which takes the value $u_c/u_c=1$, to make the degrees agree. Apply
this to $s$, with $m_1$ and $G_1$, and to $\beta(y_{k,i}/y_{k,c_k})\in A$, with $m_2$ and $G_2$ of degree $n_2$. As $\beta$
is multihomogeneous, of a multidegree $b$, $\beta(y)=\prod_ky_{k,c_k}^{b_k}\cdot\beta(y_{k,i}/y_{k,c_k})$, and an element of
multidegree $b$ lies in $P$ if and only if it vanishes at $y$, as it maps to its value at $y$ times $\prod_kt_k^{b_k}$. As
$\beta\notin P$, $\beta(y_{k,i}/y_{k,c_k})\neq0$. Put $G=T_cG_1G_2$, homogeneous of some degree $d+1$ and divisible by
$T_c$, and $g_0=G(u/u_c)=m_1m_2\,s\,\beta(y_{k,i}/y_{k,c_k})\neq0$. Then $A[1/g_0]$ is a localization of $A[1/s]$, so it is
smooth of relative dimension $e$. The element $\tilde G_2\prod_kx_{k,c_k}^{b_k}-m_2\beta(\prod_kx_{k,c_k})^{n_2}$ of $R$ is
multihomogeneous and vanishes at $y$, as $G_2(u/u_c)=m_2\beta(y_{k,i}/y_{k,c_k})$; so it lies in $P$. Let $v$ be a point of
the zero set of $P$ in $\hat C$ with $\beta(v)=0$, and $u(v)=(\prod_kv_{k,j_k})_j$. Then
$G_2(u(v))\prod_kv_{k,c_k}^{b_k}=0$, so $G_2(u(v))=0$ or $u(v)_c=\prod_kv_{k,c_k}=0$; in both cases $G(u(v))=0$.

\emph{Resolution.} Let $p:\operatorname{Spec}\Omega\to\mathbf P^N_{\mathbf Q}$ be the point given by $u$, and let
$C\subset\mathbf P^N_{\mathbf Q}$ be its scheme theoretic image, a reduced closed subscheme in which the image of $p$ is
dense \cite[Tag 056B]{Sta}, \cite[Tag 01R8]{Sta}. Let $W_0=C\cap D_+(T_c)$, an affine open subscheme whose ring maps onto
$A$ in $\Omega$, as it is a quotient of $\mathbf Q[T_j/T_c]$ \cite[Tag 01IN]{Sta}, and let $\sigma_0$ be the function
$G/T_c^{d+1}$ on $W_0$, with value $g_0$. As $T_c$ divides $G$, $D_+(G)\subset D_+(T_c)$, and $U=C\cap D_+(G)$ is the
standard open of $\sigma_0$ in $W_0$. Its ring is the localization of that of $W_0$ at $\sigma_0$, which maps isomorphically
onto $A[1/g_0]$, as $C$ is reduced and the image of $p$ is dense; so $U$ is smooth of relative dimension $e$ over
$\mathbf Q$. It contains the image of $p$, as $g_0\neq0$, so it is dense in $C$. Lemma~\ref{cv-resolution-coordinates},
applied to $\mathbf Q$, $C\subset\mathbf P^N_{\mathbf Q}$ and $U$, gives an integer $r$, a closed subscheme $X'$ of
$\mathbf P^N_{\mathbf Q}\times\mathbf P^r_{\mathbf Q}$ and the restriction $f:X'\to C$ of the first projection, with the
properties (1) to (3) of that lemma. As $f$ is an isomorphism over $U$, $p$ lifts to a morphism
$\zeta:\operatorname{Spec}\Omega\to f^{-1}(U)\subset X'$ with $f\circ\zeta=p$, whose image is dense in $X'$, as $f^{-1}(U)$
is. As in the first step of the proof of Lemma~\ref{cv-origin-extension-charts}, the second projection of $\zeta$ is given
by a nonzero $w\in\Omega^{r+1}$: with $T'_0,\dots,T'_r$ the coordinates of $\mathbf P^r$, it lies in $D_+(T'_l)$ exactly
when $w_l\neq0$, and then $T'_i/T'_l$ takes the value $w_i/w_l$.

\emph{The ratios of $w$.} Let $w_l\neq0$. The ratio $w_i/w_l$ is the value at $\zeta$ of a function on an open subset of
$X'$ that contains the image of $\zeta$. Through the isomorphism $f^{-1}(U)\to U$, it is the value at $p$ of a function on
an open subset $O$ of $C$ that contains the image of $p$. On a standard open $D(b)\subset O\cap W_0$ of $W_0$ that contains
this image, the function is $a/b^n$ with $a$ and $b$ functions on $W_0$; so its value is a quotient of two elements of $A$,
and $w_i/w_l\in F$.

\emph{The charts.} Let $a$ be an index with $u_a\neq0$ and $l$ one with $w_l\neq0$, and let $W$ be the open subscheme of
$X'$ where $T_a\neq0$ and $T'_l\neq0$. It is a closed subscheme of $D_+(T_a)\times D_+(T'_l)$, so it is affine and its ring
is a quotient of $\mathbf Q[T_j/T_a,T'_i/T'_l]$ \cite[Tag 01IN]{Sta}; it contains the image of $\zeta$, and the image of its
ring in $\Omega$ is $\mathbf Q[u_j/u_a,w_i/w_l]$. Let $\sigma$ be the function $G/T_a^{d+1}$ on $W$, with value $G(u/u_a)$
at $\zeta$. As $f^{-1}(U)$ is the set of points of $X'$ whose first projection lies in $D_+(G)$, the zero set of $\sigma$ in
$W$ is $W\setminus f^{-1}(U)$. The scheme $X'$ is smooth over $\mathbf Q$, hence normal \cite[Tag 033C]{Sta} and reduced.
By (1) and (3) of Lemma~\ref{cv-resolution-coordinates}, Lemma~\ref{cv-origin-coords-snc-chart} applies to $X'$, $\zeta$, the
set $X'\setminus f^{-1}(U)$, $W$ and $\sigma$: the element $G(u/u_a)$ has normal crossings of dimension $e$ on
$\mathbf Q[u_j/u_a,w_i/w_l]$ (Definition~\ref{cv-origin-coords-crossings}).

\emph{The block and the ideal.} Put $y_2=w$ and $z=(y,w)$, and let $J\subset R_2$ be the ideal generated by $\tilde G$,
a polynomial in the variables of the blocks of $y$. As the ratios $w_i/w_l$ lie in $F$, the ratios of entries of one block
of $z$ generate $F$: this is (1). Each monomial $\prod_kx_{k,j_k}$ has degree $1$ in each block of $y$, so $\tilde G$ is
multihomogeneous, of degree $d+1$ in each block of $y$ and $0$ in the new block, and $J$ is multihomogeneous. As
$\tilde G(z)=G(u)=u_c^{d+1}g_0\neq0$, $J$ is not contained in $P_2$.

(2) Let $\kappa$ be a chart of $z$, given by indices $\kappa_k$ for the blocks of $y$ and an index $l$ with $w_l\neq0$, and
let $a=(\kappa_k)_k$; then $u_a\neq0$. As in the first step, the ring of ratios of $z$ at $\kappa$ is
$A_\kappa=\mathbf Q[u_j/u_a,w_i/w_l]$, and, as $G$ is homogeneous, $\tilde G$ takes the value $G(u/u_a)$ at the ratios
$z_{k,i}/z_{k,\kappa_k}$. The values at these ratios of the elements $g\tilde G$ of $J$ are multiples of $G(u/u_a)$ in
$A_\kappa$, so a prime ideal of $A_\kappa[1/f_l]$ contains all of them if and only if it contains $G(u/u_a)$. So the normal
crossings of $G(u/u_a)$ on $A_\kappa$ give condition (3) of Definition~\ref{cv-origin-reference} at $\kappa$, with $e$ in
place of $D$.

(3) Let $v$ be a point of the zero set of $P_2$ in $\hat C$ with $\beta(v)=0$, and $v_y$ its blocks for $y$. The elements of
$P$, as polynomials in the variables of the blocks of $y$, lie in $P_2$; so $v_y$ lies in the zero set of $P$, and
$G(u(v_y))=0$ by the third step. So $\tilde G(v)=0$, and every element of $J$ vanishes at $v$.
\end{proof}
% Lean: SncRes.exists_snc_model, under ResKStatement (Lemma cv-resolution-coordinates): the chart c from the relevance in
% IsArithVar (hrel), the Segre indices enumerated by Fin (N + 1); exists_smooth = T.exists_standard_smooth_localization =
% SegreChart.exists_away_standardSmooth (adjoin_segre_eq, trdeg_blockFnField through CoordArithOfFnField.dim_eq, and
% AffinePoints.exists_ne_zero_isStandardSmoothOfRelativeDimension, for which the book cites Stacks); exists_form =
% T.standard_smooth_chart_of_nonzero = SncRes.exists_segre_form (SegreForm.exists_homog, aeval_blockRatio_mem,
% form_mem_mCoordIdeal_iff with aeval_eq_prod_mul, smooth_away_mul, vanish_of_value); snc_of_res: closure_facts =
% T.projective_point_image_reduced (isReduced_image_of_quasiCompact), chart_single, form_section, smooth_basicOpen =
% T.standard_smooth_basic_open, ResKStatement, lift_generic = T.lift_generic_point_over_iso_open (Coords.dense_point_lift),
% exists_coords = T.projective_point_coords_exist, isReduced_of_smooth_of_isReduced, val_mem_field =
% T.regular_values_in_adjoined_field (valR_mem_adjoin_of_modification and valR_mem_adjoin_of_affine), prod_chart =
% T.segre_product_affine_chart, and snc_transfer = T.snc_chart_of_smooth_coords (Lemma cv-origin-coords-snc-chart);
% and final_data = T.homogeneous_form_nonvanishing_point = SegreChart.final_data_main
% (J = span of the renamed segForm E Gz; blockFnField_joinPt, segForm_multiHomog, aeval_segForm, isSncModel_joinPt with
% adjoin_ratios_eq and aeval_ratios_segForm, mem_mcPts_span). The Lean's SncChart A e h is the book's "h has normal
% crossings of dimension e on A".

\begin{lemma}\label{cv-origin-change-coords}
Let $\Omega_1$ be a field and $\Omega$ a finite extension of $\Omega_1$ of characteristic zero, let
$y\in\Omega_1^{\sqcup_kI_k}$, and let $X$ be a scheme with a morphism $\xi:\operatorname{Spec}\Omega_1\to X$
of whose image every point of $X$ is a specialization, covered by the affine open subschemes of the charts
of $y$ as in Definition~\ref{cv-origin-hodge-component}(1). Then there are a point $y'\in\Omega^{\sqcup_kI'_k}$
that contains the image of $y$ as some of its blocks, a scheme $X'$ over $\mathbf Z$ with a morphism
$\xi':\operatorname{Spec}\Omega\to X'$ of whose image every point of $X'$ is a specialization, covered by the
affine open subschemes of the charts of $y'$, whose local rings are integrally closed in $\Omega$, and a
morphism $h:X'\to X$ with $h\circ\xi'=\xi\circ\operatorname{Spec}(\Omega\supset\Omega_1)$, such that
$\mathbf Q[y'_{k,i}/y'_{k,\kappa_k}]$ is smooth over $\mathbf Q$ for every chart $\kappa$ of $y'$.
\end{lemma}
\leanref{Development/AutomorphicSpectrum/Proofs/GeometryProofs/CoordsAfterFiniteExtensionStatement.lean}{GlobalTheory.T.coords_after_finite_extension}
\begin{proof}
Let $w=(w_0,\dots,w_n)$ with $w_0=1$ and $w_1,\dots,w_n$ a basis of $\Omega$ over $\Omega_1$; then $\Omega$ is generated as
a field by $\Omega_1$ and the ratios $w_a/w_b$. Let $x$ be the family of Segre products of $y$, so that $X$ has the
coordinates $x$ (Definition~\ref{cv-origin-segre-coordinates}).

\emph{The first new block.} Lemma~\ref{cv-origin-coords-join}, applied to $X$, $x$, the inclusion of $\Omega_1$ in
$\Omega$ and $w$, gives a scheme $X_1$ with a morphism $\xi_1:\operatorname{Spec}\Omega\to X_1$ of whose image every point
is a specialization, with the coordinates $(x_cw_a)_{c,a}$, and a morphism $f_1:X_1\to X$ with
$f_1\circ\xi_1=\xi\circ\operatorname{Spec}(\Omega\supset\Omega_1)$. The products $x_cw_a$ are the Segre products of the
point $y_1=(y,w)$, whose blocks are those of $y$ and the block $w$; so $X_1$ is covered by the affine open subschemes of
the charts of $y_1$. The image of $\xi_1$ lies in a chart of $y_1$, and by Lemma~\ref{cv-origin-snc-charts}(1) $\Omega$ is
the fraction field of the finitely generated ring of ratios of $y_1$ there; so $\Omega$ has a finite transcendence degree
$d$ over $\mathbf Q$. As in the proof of Lemma~\ref{cv-origin-weight-birational}, $\Omega$ is the function field of the
kernel of $y_1$, and this kernel is an arithmetic subvariety of dimension $d+1$.

\emph{The second new block.} Lemma~\ref{cv-origin-coords-snc}, applied to $y_1$ with $e=d$ and $\beta=1$, gives a block
$y_2$ with values in $\Omega$ and a nonzero entry such that the point $z=(y_1,y_2)$ satisfies condition (3) of
Definition~\ref{cv-origin-reference}, with $d$ in place of $D$, for a multihomogeneous ideal.
Then $A_\kappa=\mathbf Q[z_{k,i}/z_{k,\kappa_k}]$ is smooth over $\mathbf Q$ for every chart $\kappa$ of $z$: with the
$f_l$ and $a_l$ of that condition, every prime of $A_\kappa$ misses some $f_l$, as $\sum_la_lf_l=1$, and then $f_l\neq0$
and $A_\kappa[1/f_l]$ is smooth over $\mathbf Q$; so $A_\kappa$ is smooth at every prime \cite[Tag 00TB]{Sta}, hence smooth
\cite[Tag 00TC]{Sta}. Lemma~\ref{cv-origin-segre-join}, applied to $X_1$, the Segre products of $y_1$ and the block $y_2$, gives a scheme
$X_2$ with a morphism $\xi_2:\operatorname{Spec}\Omega\to X_2$ of whose image every point is a specialization, covered by
the affine open subschemes of the charts of $z$, and a morphism $f_2:X_2\to X_1$ with $f_2\circ\xi_2=\xi_1$.

\emph{The normalization.} For a point $P$ of $X_2$ let $\mathcal O_P\subset\Omega$ be its local ring.
Lemma~\ref{cv-origin-coords-normal}, applied to $X_2$, $\xi_2$ and the Segre products of $z$, gives a finite morphism
$f_3:X_3\to X_2$ and a morphism $\xi_3:\operatorname{Spec}\Omega\to X_3$ with $f_3\circ\xi_3=\xi_2$, such that every point
of $X_3$ specializes the image of $\xi_3$, $X_3$ has the coordinates $y_3$ for a finite family $y_3$ in $\Omega$, that is,
it is covered by the affine open subschemes of the charts of $y_3$ as a point with one block, and for every point $P$ of
$X_2$ the points of $X_3$ over $P$ correspond to the maximal ideals of the integral closure $B_P$ of $\mathcal O_P$ in
$\Omega$, with the localizations of $B_P$ at them as local rings. The local rings of $X_3$ are integrally closed in
$\Omega$. Every point $P_3$
of $X_3$ lies in a chart of $z$: the point $f_3(P_3)$ lies in one, and the elements regular at $f_3(P_3)$ are regular at
$P_3$, by Lemma~\ref{cv-origin-segre-regular}(3). So Lemma~\ref{cv-origin-reference-segre}, applied to $X_3$, the point
$y_3$ with one block and the blocks of $z$, shows that $X_3$ is covered by the affine open subschemes of the charts of
$y'=(y_3,z)$.

\emph{Smoothness.} Let $P$ be a point of $X_2$ over no prime, in a chart $\kappa$ of $z$, and put
$A=\mathbf Q[z_{k,i}/z_{k,\kappa_k}]$. The ring $A$ is smooth over $\mathbf Q$, hence normal \cite[Tag 033C]{Sta}; it is a
domain whose fraction field is $\Omega$, by Lemma~\ref{cv-origin-snc-charts}(1) at the image of $\xi_2$, so it is integrally
closed in $\Omega$. Positive integers are units at $P$, so $A\subset\mathcal O_P$, and by
Lemma~\ref{cv-origin-snc-charts}(1) every element of $\mathcal O_P$ is $u/v$ with $u,v\in A$ and $v$ a unit at $P$. Let
$a\in\Omega$ be integral over $\mathcal O_P$. Clearing the denominators of an equation of integral dependence gives $v\in A$,
a unit at $P$, with $va$ integral over $A$; so $va\in A$, and $a\in\mathcal O_P$. Hence $\mathcal O_P$ is integrally closed
in $\Omega$, $B_P=\mathcal O_P$ is local, and $f_3$ is an isomorphism over $P$ in the sense of
Lemma~\ref{cv-origin-change-projective}(2). Lemma~\ref{cv-origin-change-smooth}, applied to $f_3$, shows that
$\mathbf Q[y'_{k,i}/y'_{k,\kappa_k}]$ is smooth over $\mathbf Q$ for every chart $\kappa$ of $y'$. Put $X'=X_3$ and
$h=f_1\circ f_2\circ f_3$; the point $y'$ contains the image of $y$ as some of its blocks, through $y_1$ and $z$.
\end{proof}
% Lean: ChangeModel.exists_normal_model: closure_eq_top_of_basis, exists_join_model_dominates with Option (Fin n)
% (Lemma cv-origin-coords-join), tie_joinVec, exists_trdeg_eq and coordTie_birational_isArithVar (as in Lemma
% cv-origin-weight-birational), SncRes.exists_snc_model with beta = 1 and b = 0 (Lemma cv-origin-coords-snc) and
% Cmp.SncSmooth.smoothGenericFibreM_of_isSncModel
% (smooth_of_cover with exists_mul_pow_mem for A[1/f]; the book cites Stacks 00TB and 00TC), a second
% exists_join_model_dominates with the identity (Lemma cv-origin-segre-join), exists_normal_dominating_model with
% Omega' = Omega (Lemma cv-origin-coords-normal),
% RegAtDom.regAt_dominates for the charts (hchart), ExtraBlocks.exists_blocks_of_lines with no line bundles, whose
% CoordTie is join_coords_coordTie, the step of Lemma cv-origin-reference-segre, and T.smoothGenericFibre_of_dominates
% (Lemma cv-origin-change-smooth) with T.normal_iso_of_integrally_closed and mem_regAt_of_isIntegral (A integrally closed
% by isIntegrallyClosed_quotient_minimalPrime_of_smooth, ClearDenominators.exists_mul_isIntegral).

\subsection{Generic Weil pairings and Hodge--Tate periods}\label{subsection-cv-hodge-tate-periods}

\begin{lemma}\label{cv-origin-change-factor}
In the setting of Definition~\ref{cv-origin-hodge-component}, let $\Omega$ be a field and
$\xi:\operatorname{Spec}\Omega\to\mathcal X$ a morphism whose image is the generic point of an irreducible
component. Then $\xi=j\circ y_0$ for a point $y_0$ of $X$ with values in $\Omega$.
\end{lemma}
\leanref{Development/AutomorphicSpectrum/Proofs/GeometryProofs/GenericPointFactorsOpenStatement.lean}{GlobalTheory.T.generic_point_factors_open}
\begin{proof}
Let $e$ be the image of $\xi$, let $\pi:\mathcal X_{F_c}=\mathcal X\times_{\mathcal O_c}F_c\to\mathcal X$ be the
projection, and write $j$ also for the composite $\pi\circ j:X\to\mathcal X$.

\emph{The point $e$ lies in $\pi(j(X))$.} The model is flat over $\mathcal O_c$ (Definition~\ref{cv-model}), so
generizations lift along $\mathcal X\to\operatorname{Spec}\mathcal O_c$ \cite[Tag 03HV]{Sta}. The zero ideal of
$\mathcal O_c$ generizes the image of $e$, so some generization of $e$ lies over it; this is $e$, the generic point of
an irreducible component. Over an affine open subscheme $\operatorname{Spec}B$ of $\mathcal X$, $\pi$ is
$\operatorname{Spec}(T^{-1}B)\to\operatorname{Spec}B$ with $T=\mathcal O_c\setminus\{0\}$. So $\pi$ is injective, with
image the points over the zero ideal \cite[Tag 00E3]{Sta}, and its local rings are local rings of $\mathcal X$. Hence
$e=\pi(e')$ for a point $e'$ of $\mathcal X_{F_c}$, which has no generization other than itself since $\pi$ is
injective and continuous. The scheme $\mathcal X_{F_c}$ is of finite type over $F_c$, so its underlying space is
noetherian, and the irreducible component containing $e'$ has generic point $e'$ and contains a nonempty open subset.
This open subset meets the dense open subset $j(X)$ in a point $w$, and $e'$ specializes to $w$, so $e'\in j(X)$.

\emph{The residue field.} A morphism from the spectrum of a field to a scheme is a point with an embedding of its
residue field \cite[Tag 01J9]{Sta}. As $\pi$ identifies the residue fields at $e'$ and $e$, the canonical morphism
$\varepsilon:\operatorname{Spec}\kappa(e)\to\mathcal X$ is $\pi\circ\varepsilon'$ for a morphism $\varepsilon'$ with
image $e'$, and $\varepsilon'$ factors through the open subscheme $j(X)$: $\varepsilon=j\circ y$ for a morphism
$y:\operatorname{Spec}\kappa(e)\to X$.

\emph{The point $\xi$.} The local ring of $\mathcal X$ at $e$ is a domain, as $\mathcal X$ is normal, and its primes
correspond to the generizations of $e$ \cite[Tag 01J7]{Sta}, so it has only one; it is the field $\kappa(e)$. By
\cite[Tag 01J6]{Sta}, $\xi$ is given by a local homomorphism $\kappa(e)\to\Omega$, that is,
$\xi=\varepsilon\circ\operatorname{Spec}(i)$ for a field embedding $i:\kappa(e)\to\Omega$. So $\xi=j\circ y_0$ with
$y_0=y\circ\operatorname{Spec}(i)$.
\end{proof}
% Lean: ChangeGeo.exists_factor_XF_of_genericPoint: generic_point_normal_coords = XR.exists_cpt_of_genericPoint, whose
% y is field_point_lifts_generic_fibre (= HtPt.Ml.lift_generic_point: preimm_frac, lift_frac, IsOpenImmersion.lift) applied
% to fromSpecResidueField e, with generic_points_in_generic_fibre (= HtPt.Ml.generic_points_range: Flat.generalizingMap,
% generic_max, mem_of_dense_of_max) and isCptGen_of_maximal (injective_of_maximal_point); then TowerOfTest.factor.
% The CoordTie and NormalCpt parts of exists_cpt_of_genericPoint are not used here.

\begin{definition}\label{cv-origin-change-projective-periods}
Let $T$ be a scheme, $\Omega$ a field and $\xi:\operatorname{Spec}\Omega\to T$ a morphism with image $\eta$, let $G$ be a
commutative group scheme over $T$ with a line bundle $\mathcal L$, and let $n\ge0$ be an integer. For a field $k$ and an
embedding $i:\Omega\to k$, a point of $G$ at $\xi\circ\operatorname{Spec}i$ is a morphism $\operatorname{Spec}k\to G$ over
$\xi\circ\operatorname{Spec}i$, and a point $x$ of $G$ at $\xi$ gives the point $x_i=x\circ\operatorname{Spec}i$ at
$\xi\circ\operatorname{Spec}i$. Let $G_n$ be the group of the points $x$ of $G$ at $\xi$ with $x^n=1$ (all of them if
$n=0$), and $e_n$ the commutator pairing of the theta group of $\mathcal L^n$
(Definition~\ref{cv-toroidal-moduli-object}(2)).
\begin{enumerate}
\item $G$ has \emph{full $n$-torsion} at $\xi$ if for every field $k$ and every embedding $i:\Omega\to k$, every point $g$
  of $G$ at $\xi\circ\operatorname{Spec}i$ with $g^n=1$ is $x_i$ for a point $x$ of $G$ at $\xi$.
\item A \emph{generic Weil pairing} of order $n$ of $G$ at $\xi$ is a map $e:G_n\times G_n\to\Omega^\times$, where $G_n$
  is finite, that is multiplicative in each variable, such that every homomorphism $G_n\to\Omega^\times$ is $e(x,\cdot)$
  for some $x\in G_n$, such that $e(x,g)=1$ for every $x\in G_n$ only if $g=1$, and such that $e_n(x_i,g_i)=i(e(x,g))$
  for all $x,g\in G_n$, every algebraically closed field $k$ and every embedding $i:\Omega\to k$.
\item Let $\mathrm{Lie}$ be the Lie algebra of $G$ at $\xi$. For a point $P$ that specializes $\eta$ let $\Lambda(P)$ be
  the Lie algebra of $G$ at $P$ and $\Lambda^\vee(P)$ the $\mathcal O_P$-module of the linear forms $w$ on $\mathrm{Lie}$
  with $w(\Lambda(P))\subset\mathcal O_P$ (Definition~\ref{cv-origin-reference-lattices}). Let $x\in G_n$. A
  \emph{Hodge--Tate period} of $x$ at $P$ is a form $w\in\Lambda^\vee(P)$ for which there are an open subset $W$ of $G$
  that contains the image of $P$ under the unit section and a function $h$ on $W$ such that:
  \begin{enumerate}
  \item[(a)] $h(g)=e_n(x_i,g)$ for every algebraically closed field $k$, every embedding $i:\Omega\to k$ and every point
    $g$ of $G$ at $\xi\circ\operatorname{Spec}i$ with $g^n=1$ and image in $W$;
  \item[(b)] for every $e\in\mathrm{Lie}$ that extends to $P$ through a tangent vector $v$ of $G$ at
    $\operatorname{Spec}\mathcal O_{T,P}\to T$ with image in $W$, if $h\circ v=a+b\varepsilon$ with
    $a,b\in\mathcal O_{T,P}$, then $w(e)$ minus the image of $b$ in $\Omega$ lies in $n\mathcal O_P$.
  \end{enumerate}
  The module $H_n(P)$ \emph{spanned by the Hodge--Tate periods} of order $n$ at $P$ is the $\mathcal O_P$-submodule of
  the dual $\mathrm{Lie}^\vee$ generated by the Hodge--Tate periods at $P$ of the points $x\in G_n$.
\end{enumerate}
\end{definition}
\leanref{Development/AutomorphicSpectrum/Definitions/GlobalGeometry/MemEigSpOfMemHtDual.lean}{GlobalTheory.CompanionTheory.Ar.FullTorsion}
\leanref{Development/AutomorphicSpectrum/Definitions/GlobalGeometry/GenWeil.lean}{GlobalTheory.CompanionTheory.Ar.GenWeil}
\leanref{Development/AutomorphicSpectrum/Definitions/GlobalGeometry/EigLat.lean}{GlobalTheory.CompanionTheory.Ar.LieModel.omLat}
\leanref{Development/AutomorphicSpectrum/Definitions/GlobalGeometry/MemEigSpOfMemHtDual.lean}{GlobalTheory.CompanionTheory.Ar.LieModel.IsHTAt}
\leanref{Development/AutomorphicSpectrum/Definitions/GlobalGeometry/MemEigSpOfMemHtDual.lean}{GlobalTheory.CompanionTheory.Ar.LieModel.htOm}

\begin{lemma}\label{cv-origin-change-projective-weil-pullback-characters}
Let $C$ be a finite group and $\iota:K\to K'$ an embedding of fields, and suppose that the homomorphisms $C\to K^\times$
separate the points of $C$: if $\chi(c)=1$ for every such homomorphism $\chi$, then $c=1$. Then every homomorphism
$\chi':C\to K'^\times$ is $\iota\circ\chi$ for a homomorphism $\chi:C\to K^\times$.
\end{lemma}
\leanref{Development/WeilBound/CoreStatements/HodgeTatePoints/MkWeil.lean}{PositivityTheory.HtPt.Mk.char_descend}
\begin{proof}
By Dedekind's lemma, the homomorphisms $C\to K'^\times$ are linearly independent in the $K'$-vector space of functions
$C\to K'$, which has dimension $|C|$; so there are at most $|C|$ of them, and in the same way the group $\hat C$ of
homomorphisms $C\to K^\times$ has at most $|C|$ elements. As these separate points, $c\mapsto(\chi\mapsto\chi(c))$ is an
injective homomorphism from $C$ to the group of homomorphisms $\hat C\to K^\times$, which has at most $|\hat C|$ elements;
so $|C|\le|\hat C|$. The map $\chi\mapsto\iota\circ\chi$ from $\hat C$ to the homomorphisms $C\to K'^\times$ is injective,
as $\iota$ is, and its target has at most $|C|\le|\hat C|$ elements; so it is bijective.
\end{proof}
% Lean: char_descend with card_char_le (linearIndependent_monoidHom, the evaluation map ev and the injection m).

\begin{lemma}\label{cv-origin-change-projective-weil-pullback}
Let $T$ and $T'$ be schemes, $\Omega$ and $\Omega'$ fields, $\xi:\operatorname{Spec}\Omega\to T$ and
$\xi':\operatorname{Spec}\Omega'\to T'$ morphisms, $\iota:\Omega\to\Omega'$ a field embedding and $f:T'\to T$ a morphism
with $f\circ\xi'=\xi\circ\operatorname{Spec}\iota$. Let $G$ and $G'$ be commutative group schemes over $T$ and $T'$ and
$\varphi:G'\to G$ a homomorphism over $f$ that identifies $G'$ with $f^*G$; the commutativity of $G'$ follows from the
other hypotheses, as $f^*G$ is commutative. Let $n\ge0$ be an integer such that $G$
has full $n$-torsion at $\xi$ (Definition~\ref{cv-origin-change-projective-periods}). A point $x'$ of $G'$ at $\xi'$
\emph{lies over} a point $x$ of $G$ at $\xi$ if $\varphi\circ x'=x\circ\operatorname{Spec}\iota$.
\begin{enumerate}
\item $G'$ has full $n$-torsion at $\xi'$, each point of $G$ at $\xi$ has at most one point of $G'$ over it, each point of
  $G'$ at $\xi'$ lies over at most one point of $G$, $x'y'$ lies over $xy$ whenever $x'$ and $y'$ lie over $x$ and $y$,
  every $x\in G_n$ has a point of $G'_n$ over it, and every point of $G'_n$ lies over a point of $G_n$.
\item Let $\mathcal L$ and $\mathcal L'$ be line bundles on $G$ and $G'$ with $\mathcal L'$ isomorphic to
  $\varphi^*\mathcal L$. For every generic Weil pairing $e$ of order $n$ of $G$ at $\xi$, for $\mathcal L$, there is a
  generic Weil pairing $e'$ of order $n$ of $G'$ at $\xi'$, for $\mathcal L'$, with $e'(x',g')=\iota(e(x,g))$ whenever
  $x,g\in G_n$ and $x'$ and $g'$ lie over $x$ and $g$.
\end{enumerate}
\end{lemma}
\leanref{Development/AutomorphicSpectrum/Proofs/GeometryProofs/FullTorsionWeilPullbackStatement.lean}{GlobalTheory.T.full_torsion_weil_pullback}
\leanref{Development/AutomorphicSpectrum/Proofs/GeometryProofs/FullTorsionPullbackPointsStatement.lean}{GlobalTheory.T.full_torsion_pullback_points}
\begin{proof}
\emph{Points.} Let $k$ be a field and $a:\Omega'\to k$ an embedding. Then $f\circ\xi'\circ\operatorname{Spec}a=
\xi\circ\operatorname{Spec}(a\circ\iota)$, and as $\varphi$ identifies $G'$ with $f^*G$, composition with $\varphi$ is a
bijection $E_a$ from the points of $G'$ at $\xi'\circ\operatorname{Spec}a$ onto those of $G$ at
$\xi\circ\operatorname{Spec}(a\circ\iota)$; it is compatible with the group laws, as $\varphi$ is a homomorphism. In
particular a point of $G'$ is determined by its composite with $\varphi$, and $G'$ is commutative, as $G$ is. For a point
$x$ of $G$ at $\xi$ let $\tilde x=E_{\mathrm{id}}^{-1}(x\circ\operatorname{Spec}\iota)$, the point of $G'$ at $\xi'$ over
$x$. Then $x\mapsto\tilde x$ is a homomorphism, and it is injective, as a point with values in $\Omega$ is determined by its
composite with $\operatorname{Spec}\iota$. So each point of $G$ has exactly one point of $G'$ over it, each point of $G'$
lies over at most one point of $G$, and $x'y'$ lies over $xy$ if $x'$ and $y'$ lie over $x$ and $y$.

\emph{Torsion.} Let $x'\in G'_n$. Then $E_{\mathrm{id}}(x')$ is a point of $G$ at $\xi\circ\operatorname{Spec}\iota$ of
order dividing $n$, so by full $n$-torsion it is $x\circ\operatorname{Spec}\iota$ for a point $x$ of $G$ at $\xi$; then
$x'=\tilde x$, and $x\in G_n$, as $x\mapsto\tilde x$ is an injective homomorphism. Conversely $\tilde x\in G'_n$ for
$x\in G_n$. So $x\mapsto\tilde x$ is an isomorphism from $G_n$ onto $G'_n$. Now let $k$ be a field, $a:\Omega'\to k$ an
embedding and $g$ a point of $G'$ at $\xi'\circ\operatorname{Spec}a$ of order dividing $n$. Then $E_a(g)$ has order
dividing $n$, so it is $x\circ\operatorname{Spec}(a\circ\iota)$ for a point $x$ of $G$ at $\xi$, by full $n$-torsion. The
points $g$ and $\tilde x\circ\operatorname{Spec}a$ have the same composite $x\circ\operatorname{Spec}(a\circ\iota)$ with
$\varphi$, so they are equal. Hence $G'$ has full $n$-torsion at $\xi'$, which proves (1).

(2) Let $e$ be a generic Weil pairing of order $n$ of $G$ at $\xi$, and put $e'(\tilde x,\tilde g)=\iota(e(x,g))$ for
$x,g\in G_n$; by the previous step this defines $e'$ on $G'_n\times G'_n$. The group $G'_n$ is finite, as $G_n$ is, and
$e'$ is multiplicative in each variable, as $e$ is and $x\mapsto\tilde x$ is a homomorphism. If $e'(x',\tilde g)=1$ for
every $x'\in G'_n$, then $e(x,g)=1$ for every $x\in G_n$, as $\iota$ is injective, so $g=1$ and $\tilde g=1$.

\emph{Perfectness.} Let $\chi':G'_n\to\Omega'^\times$ be a homomorphism. Each $e(x,\cdot)$, $x\in G_n$, is a homomorphism
$G_n\to\Omega^\times$, and by the nondegeneracy of $e$ these separate the points of $G_n$. So, by
Lemma~\ref{cv-origin-change-projective-weil-pullback-characters}, $g\mapsto\chi'(\tilde g)$ is $\iota\circ\chi$ for a
homomorphism $\chi:G_n\to\Omega^\times$, which is $e(x,\cdot)$ for some $x\in G_n$. Then
$\chi'(\tilde g)=\iota(e(x,g))=e'(\tilde x,\tilde g)$ for every $g\in G_n$.

\emph{Geometric points.} Let $k$ be an algebraically closed field, $a:\Omega'\to k$ an embedding and $x,g\in G_n$. As
$\mathcal L'$ is isomorphic to $\varphi^*\mathcal L$ and $\varphi$ identifies $G'$ with $f^*G$, the commutator pairing
$e_n$ of $\mathcal L'^n$ at $(\tilde x\circ\operatorname{Spec}a,\tilde g\circ\operatorname{Spec}a)$ is that of
$\mathcal L^n$ at their composites $x\circ\operatorname{Spec}(a\circ\iota)$ and $g\circ\operatorname{Spec}(a\circ\iota)$
with $\varphi$. So it is $(a\circ\iota)(e(x,g))=a(e'(\tilde x,\tilde g))$. Hence $e'$ is a generic Weil pairing of order $n$
of $G'$ at $\xi'$. If $x'$ and $g'$ lie over $x$ and $g$, then $x'=\tilde x$ and $g'=\tilde g$, and
$e'(x',g')=\iota(e(x,g))$.
\end{proof}
% Lean: full_torsion_weil_pullback = HtPt.Mk.ht_weil_pull: exists_pushEquiv for E_a, inc for x -> x~ (hinc, injective by
% epi_fieldMap), hsurj, the full torsion of G' through Ek, and e' = i o e o (r x', r g') with the clauses finite, mul_left,
% mul_right, perfect (char_descend), nondeg and weil (weilIs_grpPullback_of_comm). full_torsion_pullback_points =
% TorsPull.torsion_of_pullback: pushHom and preHom (Pinj, Qinj by spMap_cancel), HcptPull.fullTorsion_of_dominates.

\begin{lemma}\label{cv-origin-change-projective-weil-pairing}
In Situation~\ref{cv-function-situation}, let $l$ be a level with full level $N$
(Definition~\ref{cv-model-toroidal-semiabelian}(1)), let $B$ be a ring in which every positive integer is a unit, with
elements $\rho$ and $\zeta_B$ such that $\rho^2=-d$ and $\zeta_B$ is a primitive $N$-th root of unity, and let
$(G_0,\eta,\eta')$ be an object at level $l$ over a scheme $S$ over $B$, defined as in
Definition~\ref{cv-toroidal-moduli-object}(4) with $B$, $\rho$ and $\zeta_B$ in place of $K_l$, $\sqrt{-d}$ and
$\zeta_N$. Let $G$ be an $\mathcal O_E$-semi-abelian scheme over a scheme $T$ with a morphism $G_0\to G$ over a morphism
$f:S\to T$. Let $\Omega$ be a field of characteristic zero, $\xi_0:\operatorname{Spec}\Omega\to S$ a morphism, and
$n\ge1$ an integer such that $G$ has full $n$-torsion at $f\circ\xi_0$ and $\Omega$ contains a primitive $n$-th root of
unity. Then $G$ has, for its line bundle, a generic Weil pairing of order $n$ at $f\circ\xi_0$
(Definition~\ref{cv-origin-change-projective-periods}).
\end{lemma}
\leanref{Development/AutomorphicSpectrum/Proofs/GeometryProofs/ExistsGenericWeilPairingStatement.lean}{GlobalTheory.T.exists_generic_weil_pairing}
\begin{proof}
Put $\xi=f\circ\xi_0$, let $\zeta\in\Omega$ be a primitive $n$-th root of unity, $G_n$ the group of the points of order
dividing $n$ of $G$ at $\xi$, $\Lambda$ the lattice of the PEL datum with its alternating form $\langle\,,\rangle$, and $K$
an algebraic closure of $\Omega$. For an embedding $i$ of $\Omega$ in a field, $x\mapsto x_i$ is an injective
homomorphism, as a point with values in $\Omega$ is determined by its image over an extension field.

\emph{Geometric points.} Let $k$ be an algebraically closed field and $i:\Omega\to k$ an embedding. The object $G_0$ has a
Tate system $((\zeta'_m),(\iota_m))$ at $\xi_0\circ\operatorname{Spec}i$ (Definition~\ref{cv-toroidal-moduli-object}). The
morphism $G_0\to G$ identifies $G_0$ with $f^*G$ and its line bundle with the pullback of that of $G$
(Definition~\ref{cv-model-toroidal-semiabelian}(2)); so composition with it is a bijection from the points of $G_0$ at
$\xi_0\circ\operatorname{Spec}i$ onto those of $G$ at $\xi\circ\operatorname{Spec}i$, compatible with the group laws and
the commutator pairings. So, with $\iota$ the map $\iota_n$ followed by this bijection and $\zeta'=\zeta'_n$, $\iota$ is a
homomorphism from $\Lambda$ onto the group of the points of $G$ at $\xi\circ\operatorname{Spec}i$ of order dividing $n$,
with kernel $n\Lambda$, and $e_n(\iota(a),\iota(b))=\zeta'^{\langle a,b\rangle}$, an $n$-th root of unity.

\emph{The pairing.} Let $x,g\in G_n$. By the previous step for $k=K$, $e_n(x_K,g_K)$ is an $n$-th root of unity of $K$, so
it is the image of $\zeta^m$ for some integer $m$, as the image of $\zeta$ is a primitive $n$-th root of unity; put
$e(x,g)=\zeta^m\in\Omega^\times$. Let $k$ be an algebraically closed field and $i:\Omega\to k$ an embedding. An algebraic
closure $M$ of the quotient of $k\otimes_\Omega K$ by a maximal ideal has embeddings $\alpha:k\to M$ and $\beta:K\to M$
that agree on $\Omega$. The commutator pairings are compatible with base change along embeddings of fields; so
$\alpha(e_n(x_i,g_i))$ and $\beta(e_n(x_K,g_K))=\alpha(i(e(x,g)))$ are both the commutator pairing of the images of $x$
and $g$ over $M$, and $e_n(x_i,g_i)=i(e(x,g))$, as $\alpha$ is injective.

\emph{Its properties.} Let $\iota$ and $\zeta'$ be as in the first step for $k=K$. For $x\in G_n$, $x_K=\iota(a)$ for some
$a\in\Lambda$; conversely, by full $n$-torsion, each $\iota(a)$ is $x_K$ for a point $x$ at $\xi$, which lies in $G_n$ as
$x\mapsto x_K$ is injective. Let $x_K=\iota(a)$, $x'_K=\iota(a')$ and $g_K=\iota(c)$. Then $(xx')_K=\iota(a+a')$, and the
images in $K$ of $e(xx',g)$ and of $e(x,g)e(x',g)$ are
$\zeta'^{\langle a+a',c\rangle}=\zeta'^{\langle a,c\rangle}\zeta'^{\langle a',c\rangle}$; so $e$ is multiplicative in
the first variable, and in the same way in the second. If $e(x,g)=1$ for every $x\in G_n$, then
$\zeta'^{\langle a,c\rangle}=1$ and $n$ divides $\langle a,c\rangle=-\langle c,a\rangle$ for every $a\in\Lambda$. As the
form is perfect, the linear form $a\mapsto\langle c,a\rangle/n$ on $\Lambda$ is $\langle c',\cdot\rangle$ for some
$c'\in\Lambda$, and then $c=nc'$; so $g_K=\iota(c)=1$ and $g=1$. In the same way $e(x,g)=1$ for every $g\in G_n$ only if
$x=1$. The group $G_n$ embeds into the image of $\Lambda/n\Lambda$, which is finite, as $\Lambda$ is finitely generated.
So $x\mapsto e(x,\cdot)$, followed by $\Omega\subset K$, is an injective homomorphism from $G_n$ to the group of
homomorphisms $G_n\to K^\times$, which has as many elements as $G_n$, as $K$ is algebraically closed of characteristic
zero. It is bijective, and every homomorphism $\chi:G_n\to\Omega^\times$ is $e(x,\cdot)$ for some $x\in G_n$, as
$\Omega\to K$ is injective. So $e$ is a generic Weil pairing of order $n$ of $G$ at $\xi$.
\end{proof}
% Lean: exists_generic_weil_pairing = WeilObj.genWeil_of_obj (obj_tate_at: o.geom at spMap i >> xi0 with the
% homomorphism B -> k, moved along the cartesian square m.cart by pushHom and weilIs_grpPullback_of_comm' with the
% line bundle isomorphism from m.lin; exk; the choice of e as a power of zeta in valK; char with exists_common_closed and
% T.weil_pairing_base_change_unit = weilIs_fieldExt; hl and hr; ef_dvd from D.unimod and D.alt for nondegR and nondegL;
% tors_finite; CommGroup.card_monoidHom_of_hasEnoughRootsOfUnity for perfect).

\begin{lemma}\label{cv-origin-change-projective-weil}
In the setting of Definition~\ref{cv-origin-hodge-component}, let $\Omega$ be a field of characteristic zero, $T$ a
scheme with a morphism $\xi:\operatorname{Spec}\Omega\to T$, and $f:T\to\mathcal X$ a morphism with
$f\circ\xi=j\circ y_0$ for a point $y_0$ of $X$ with values in $\Omega$. Let $G$ be an $\mathcal O_E$-semi-abelian scheme
over $T$ (Definition~\ref{cv-model-toroidal-semiabelian}(2)) with a morphism $\varphi:G\to\mathcal A$ over $f$, and let
$n\ge1$ be an integer such that $\mathcal A$ has full $n$-torsion at $j\circ y_0$ and $\Omega$ contains a primitive
$n$-th root of unity. Then $G$ has full $n$-torsion at $\xi$ and, for its line bundle, a generic Weil pairing of order $n$
at $\xi$ (Definition~\ref{cv-origin-change-projective-periods}).
\end{lemma}
\leanref{Development/AutomorphicSpectrum/Proofs/GeometryProofs/HodgeCptProjectiveModel.lean}{GlobalTheory.CompanionTheory.Ar.TopHodge.tors_weil}
\begin{proof}
Put $\xi_0=j\circ y_0$. The abelian scheme $A$ of $X$ is, with its level structures, an object at level $l$ over $X$
(Definition~\ref{cv-toroidal-moduli-object}(4)), with a morphism $A\to\mathcal A$ over $j$ (Definition~\ref{cv-model}). So
Lemma~\ref{cv-origin-change-projective-weil-pairing}, applied to this object, $\mathcal A$, $j$ and $y_0$, gives a generic
Weil pairing $e$ of order $n$ of $\mathcal A$ at $\xi_0$. The morphism $\varphi$ identifies $G$ with $f^*\mathcal A$, and
the line bundle of $G$ with the pullback of that of $\mathcal A$, and $f\circ\xi=\xi_0$. So
Lemma~\ref{cv-origin-change-projective-weil-pullback}, applied with $\iota$ the identity, shows that $G$ has full
$n$-torsion at $\xi$ and gives a generic Weil pairing of order $n$ of $G$ at $\xi$.
\end{proof}
% Lean: TopHodge.tors_weil: T.exists_generic_weil_pairing (Lemma cv-origin-change-projective-weil-pairing) for the object
% Y.oF and Y.m, then T.full_torsion_weil_pullback (Lemma cv-origin-change-projective-weil-pullback) along f with i the
% identity and m.cart, with the line bundle isomorphism from m.lin.

\begin{definition}\label{cv-origin-change-projective-local-torsion}
Let $R$ be a domain with fraction field $K$, and $C$ a commutative group.
\begin{enumerate}
\item A \emph{torsion scheme} of type $C$ over $R$ is a commutative group scheme $\Gamma$ over $\operatorname{Spec}R$ that
  is flat, locally quasi-finite, separated, quasi-compact and locally of finite presentation over $R$, with a homomorphism
  $t\mapsto\gamma_t$ from $C$ to the group of the points of $\Gamma$ with values in $K$, such that the images of $\gamma_s$
  and $\gamma_t$ are distinct points for $s\neq t$, every point of $\Gamma$ over the generic point of $\operatorname{Spec}R$
  is the image of some $\gamma_t$, and the generic fibre $\Gamma\times_RK$ is reduced.
\item Let $\Gamma$ be a torsion scheme of type $C$ over $R$, $S$ a scheme, $G$ a group scheme over $S$,
  $x:\operatorname{Spec}R\to S$ a morphism and $N\ge0$ an integer. A morphism $i:\Gamma\to G$ is a \emph{torsion embedding
  along $x$ of order $N$} if it lies over $x$ and, for every scheme $Y$ over $R$ and all points $y$ and $z$ of $\Gamma$ with
  values in $Y$: $i\circ(yz)=(i\circ y)(i\circ z)$, for the group law of $G$ over $x\circ(Y\to\operatorname{Spec}R)$;
  $i\circ y=i\circ z$ only if $y=z$; and a point $g$ of $G$ over $x\circ(Y\to\operatorname{Spec}R)$ has $g^N=1$ if and only
  if $g=i\circ y$ for some point $y$ of $\Gamma$ with values in $Y$.
\item Let $i:\Gamma\to G$ be as in (2), $L_0$ the Lie algebra of $G$ at $x$, an $R$-module
  (Definition~\ref{cv-origin-reference-lattices}(1)), and $\chi:C\to R^\times$ a homomorphism. Let $\pi:R\to R/NR$ be the
  quotient map and $\pi[\varepsilon]$ the map $\pi$ with $\varepsilon\mapsto\varepsilon$. A \emph{Hodge--Tate form} of
  $\chi$ is an $R$-linear form $w_0$ on $L_0$ for which there is a function $s$ on $\Gamma$ with $s\circ\gamma_t=\chi(t)$
  for every $t\in C$, such that every $e\in L_0$ has a morphism $v:\operatorname{Spec}(R/NR)[\varepsilon]\to\Gamma$ over
  $R$ with $i\circ v=e\circ\operatorname{Spec}\pi[\varepsilon]$ and $s\circ v=1+\pi(w_0(e))\varepsilon$.
\end{enumerate}
\end{definition}
\leanref{Development/AutomorphicSpectrum/Definitions/GlobalGeometry/ExistsUnitFormLocalizationEtaleNbhd.lean}{GlobalTheory.CompanionTheory.Ar.BT.QFGroup}
\leanref{Development/AutomorphicSpectrum/Definitions/GlobalGeometry/ExistsUnitFormLocalizationEtaleNbhd.lean}{GlobalTheory.CompanionTheory.Ar.BT.QFGroup.IsTorsionAlong}
\leanref{Development/AutomorphicSpectrum/Definitions/GlobalGeometry/HasValues.lean}{GlobalTheory.CompanionTheory.Ar.BT.QFGroup.HasValues}
\leanref{Development/AutomorphicSpectrum/Definitions/GlobalGeometry/HasValues.lean}{GlobalTheory.CompanionTheory.Ar.BT.QFGroup.IsHT}

\begin{lemma}\label{cv-origin-change-projective-local-chart}
Let $T$ be a locally noetherian scheme, $\Omega$ a field of characteristic zero and $\xi:\operatorname{Spec}\Omega\to T$ a
morphism with image $\eta$ such that the local ring at $\eta$ maps bijectively onto $\Omega$, and the local rings at the
specializations of $\eta$ map injectively into $\Omega$ and are integrally closed in $\Omega$. Then every point $P$ that
specializes $\eta$ lies in the image of an open immersion $x:\operatorname{Spec}R\to T$, where $R$ is a noetherian
integrally closed domain with an isomorphism $\kappa$ of its fraction field onto $\Omega$, such that every point of the
image of $x$ specializes $\eta$ and $\xi=x\circ\operatorname{Spec}\kappa_R$, where $\kappa_R:R\to\Omega$ is the
restriction of $\kappa$.
\end{lemma}
\leanref{Development/AutomorphicSpectrum/Proofs/GeometryProofs/NormalAffineChartExistsStatement.lean}{GlobalTheory.T.normal_affine_chart_exists}
\begin{proof}
As $T$ is locally noetherian, $P$ has an affine open neighbourhood $U$ for which the map $R=\Gamma(U,\mathcal O_T)\to
\mathcal O_{T,P}$ is injective: in a noetherian affine neighbourhood the kernel of this map is finitely generated, so it
is killed by a power of some $f$ that does not vanish at $P$, and the basic open set of $f$ will do. The ring $R$ is
noetherian. Let $k_0:R\to\Omega$ be the composite of $R\to\mathcal O_{T,P}$ with the map $\mathcal O_{T,P}\to\Omega$
through the local ring at $\eta$; it is injective, as both maps are, so $R$ is a domain. As $\xi$ is the composite of
$\operatorname{Spec}\Omega\to\operatorname{Spec}\mathcal O_{T,P}$ with the canonical morphism to $T$, $\xi=x\circ
\operatorname{Spec}k_0$, where $x:\operatorname{Spec}R\to T$ is the inclusion of $U$. So $\eta$ is the point of
$\operatorname{Spec}R$ given by the kernel of $k_0$, which is zero, and every point of $U$ specializes $\eta$. Every
element of $\Omega$ is $b/c$ with $b$ and $c$ regular at $P$, as $\mathcal O_{T,\eta}$, which maps onto $\Omega$, is a
localization of $\mathcal O_{T,P}$; and $\mathcal O_{T,P}$ is a localization of $R$. So $k_0$ identifies the fraction
field of $R$ with $\Omega$, which gives $\kappa$. The localization of $R$ at a maximal ideal is the local ring of $T$ at a
point of $U$, which specializes $\eta$; so it is integrally closed in $\Omega$, its fraction field, and $R$ is integrally
closed \cite[Proposition~5.13]{AM69}.
\end{proof}
% Lean: normal_affine_chart_exists = HtPt.Mk.ht_lat_chart: Scheme.exists_germ_injective (Mathlib, for a locally
% noetherian scheme), k0 = genMap o germ, hx through gen_stalk, hall, IsFractionRing through
% T.exists_stalk_fraction_of_generic and isLocalization_stalk, IsIntegrallyClosed.of_localization_maximal with stalk_normal0.

\begin{lemma}\label{cv-origin-change-projective-local-inputs-multiplication}
Let $S$ be a scheme and $G$ a semi-abelian scheme over $S$: a commutative group scheme, smooth, separated and
quasi-compact over $S$ with geometrically connected fibres, whose fibre at every point with values in an algebraically
closed field is an extension of an abelian variety by a torus. Let $N\ge1$ be an integer. Then multiplication by $N$,
$[N]:G\to G$, is flat and locally quasi-finite.
\end{lemma}
\leanref{Development/AbelianVarietyTheory/Proofs/IwasawaDecomposition/PowFlatLocallyQuasiFinite.lean}{GlobalTheory.SemiAbMulFlat.flat_and_locallyQuasiFinite_pow}
\begin{proof}
The homomorphism $[N]$ is locally of finite type and quasi-compact, as its composite with $G\to S$ is.

\emph{Quasi-finiteness.} Let $k$ be an algebraically closed field and $z:\operatorname{Spec}k\to S$ a point. The fibre
$G_z$ is an extension of an abelian variety $A$ by a torus of some dimension $r$. The torus has at most $N^r$ points with
values in $k$ of order dividing $N$, and $A$ finitely many, as multiplication by $N$ on $A$ is an isogeny
\cite{Mum70}; so $G_z$ has finitely many points with values in $k$ of order dividing $N$. Two points with values in
$k$ with the same image under $[N]$ differ by one of these, so every point of $G$ with values in $k$ has finitely many
points with values in $k$ over it. A scheme of finite type over $k$ with finitely many points with values in $k$ is
finite. The fibre of $[N]$ over a point of $G$ with residue field $\kappa$ is the image of its base change to an algebraic
closure of $\kappa$, which is such a scheme; so every fibre of $[N]$ is finite, and $[N]$ is locally quasi-finite
\cite[Tag 02NG]{Sta}.

\emph{Flatness.} As $G$ is flat and locally of finite presentation over $S$, by the fibrewise criterion of flatness
\cite[Th\'eor\`eme~11.3.10]{Gro66} it suffices that for every field $k$ and point $z:\operatorname{Spec}k\to S$ the base
change $[N]_z$ of $[N]$ to the fibre $G_z$ is flat. Flatness descends along $\operatorname{Spec}\bar k\to
\operatorname{Spec}k$ for an algebraic closure $\bar k$ of $k$ \cite[Proposition~2.5.1]{Gro65}, so we may assume that $k$
is algebraically closed. Then $G_z$ is of finite type over $k$, its local rings are domains, as it is smooth, and it is
connected; so it is integral. The morphism $[N]_z$ is locally quasi-finite and quasi-compact, being a base change of
$[N]$, and it fixes the unit point $o$. The generizations of $o$, ordered by specialization, form a set of finite length
$\dim\mathcal O_{G_z,o}$ whose greatest element is the generic point $\theta$. A locally quasi-finite morphism maps a
strict specialization to a strict specialization, as its fibres are discrete; so $[N]_z$ maps a chain of strict
specializations of length $\dim\mathcal O_{G_z,o}$ from $\theta$ to $o$ to such a chain ending at $o$, and
$[N]_z(\theta)=\theta$. By generic flatness \cite[Th\'eor\`eme~6.9.1]{Gro65} there is a nonempty open subset $U$ of $G_z$
over which $[N]_z$ is flat. Its inverse image contains $\theta$, so it is a nonempty open subset on which $[N]_z$ is flat,
and it contains a closed point $y$. For a closed point $x$ let $\alpha$ and $\beta$ be the translations by $x-y$ and by
$N(x-y)$, which are automorphisms of the scheme $G_z$. Then $[N]_z\circ\alpha=\beta\circ[N]_z$ and $\alpha(y)=x$, so
$[N]_z$ is flat at $x$. Every point of $G_z$ specializes to a closed point, and flatness at a point passes to its
generizations; so $[N]_z$ is flat.
\end{proof}
% Lean: flat_and_locallyQuasiFinite_pow: locallyQuasiFinite_of_finite_geometric_ker (finite_fiber_of_finite_ker,
% locallyQuasiFinite_of_finite_geometric_points) with HomExt.finite_torsion_points (finite_torsion_semiAbelian through
% semiAbelian_torsion_finite_and_closedPoint_finiteOrder: the torus extension, and AbTorsionFinite.finite_torsion for the
% abelian variety); flat_of_flatOver_fields, the fibrewise criterion for a locally quasi-finite morphism, proved its own
% way (the book cites EGA IV_3 11.3.10); IntDiv.flat_fibre (flat_of_flat_fieldBaseChange,
% grpObj_geometricallyIntegral_of_smooth) and IntDiv.flat_of_quasiFinite_grpEnd (FiniteEndFlat.apply_genericPoint,
% exists_open_flat, the translations alpha and beta, UnramifiedEndFlat.flat_of_closedPoints). About 6,400 lines, mostly the
% Lean's own proofs of the cited results.

\begin{lemma}\label{cv-origin-change-projective-local-inputs-torsion}
Let $R$ be a noetherian domain of characteristic zero with fraction field $K$, $S$ a scheme, $G$ a semi-abelian scheme
over $S$ and $x:\operatorname{Spec}R\to S$ a morphism. Let $C$ be a commutative group, $N\ge1$ an integer, and $\tau$ an
injective homomorphism from $C$ to the group of the points of $G$ at $x\circ\operatorname{Spec}(R\subset K)$, with
$\tau(t)^N=1$ for every $t$, such that for every algebraically closed field $k$ and embedding $a:K\to k$, every point $g$
of $G$ at $x\circ\operatorname{Spec}(R\subset K)\circ\operatorname{Spec}a$ with $g^N=1$ is
$\tau(t)\circ\operatorname{Spec}a$ for some $t\in C$. Then there are a torsion scheme $\Gamma$ of type $C$ over $R$ and a
torsion embedding $i:\Gamma\to G$ along $x$ of order $N$ with $i\circ\gamma_t=\tau(t)$ for every $t\in C$
(Definition~\ref{cv-origin-change-projective-local-torsion}).
\end{lemma}
\leanref{Development/AutomorphicSpectrum/Proofs/GeometryProofs/TorsionEmbeddingGeometricPointsStatement.lean}{GlobalTheory.T.torsion_embedding_geometric_points}
\begin{proof}
Let $G_R=G\times_S\operatorname{Spec}R$, with projection $q:G_R\to G$. For a scheme $Y$ over $R$, composition with $q$ is
a bijection from the points of $G_R$ with values in $Y$ over $R$ onto the points of $G$ over
$x\circ(Y\to\operatorname{Spec}R)$, compatible with the group laws. The group scheme $G_R$ is semi-abelian over $R$, as
the properties of $G$ pass to base changes, so by Lemma~\ref{cv-origin-change-projective-local-inputs-multiplication}
multiplication by $N$ on $G_R$ is flat and locally quasi-finite.

\emph{The kernel.} Let $\Gamma$ be the fibre product of $[N]:G_R\to G_R$ with the unit section
$\operatorname{Spec}R\to G_R$, and $j:\Gamma\to G_R$ the projection. For a scheme $Y$ over $R$, composition with $j$ is a
bijection from the points of $\Gamma$ with values in $Y$ over $R$ onto the points $g$ of $G_R$ with $g^N=1$. These form a
subgroup, as $G_R$ is commutative; so $\Gamma$ is a commutative group scheme over $R$, and $j$ is a homomorphism and a
monomorphism. The unit section is a closed immersion, as $G_R$ is separated over $R$ \cite[Tag 01KT]{Sta}, and so is its
base change $j$. The structure morphism of $\Gamma$ is the base change of $[N]$ along the unit section, so it is flat and
locally quasi-finite. It is separated and quasi-compact, as $j$ is a closed immersion and $G_R$ is separated and
quasi-compact over $R$, and it is locally of finite type, hence locally of finite presentation, as $R$ is noetherian.

\emph{The generic fibre.} As $R$ has characteristic zero, $N$ is a unit of $K$. So multiplication by $N$ on
$G_K=G_R\times_RK$ is formally unramified: if two points $g_1,g_2$ of $G_K$ with values in an affine scheme agree on a
closed subscheme defined by an ideal of square zero and $Ng_1=Ng_2$, then $g=g_1g_2^{-1}$ is trivial on that subscheme,
such points multiply by adding the derivations they define, and $g^N=1$ gives $g=1$. Its base change
$\Gamma_K=\Gamma\times_RK\to\operatorname{Spec}K$ along the unit section is formally unramified, flat and locally of
finite presentation, hence \'etale \cite[Th\'eor\`eme~17.6.1]{Gro67}, and so $\Gamma_K$ is reduced
\cite[Proposition~17.5.7]{Gro67}.

\emph{The points.} Put $i=q\circ j:\Gamma\to G$. For a scheme $Y$ over $R$, composition with $i$ is an injective
homomorphism from the points of $\Gamma$ with values in $Y$ over $R$ onto the points $g$ of $G$ over
$x\circ(Y\to\operatorname{Spec}R)$ with $g^N=1$. So each $\tau(t)$ is $i\circ\gamma_t$ for a unique point $\gamma_t$ of
$\Gamma$ with values in $K$ over $R$, and $t\mapsto\gamma_t$ is a homomorphism. A morphism from the spectrum of a field to
$\Gamma$ is a point of $\Gamma$ with an embedding of its residue field \cite[Tag 01J9]{Sta}. If $\gamma_s$ and $\gamma_t$
have the same image $g$, the two embeddings of the residue field $\kappa$ at $g$ into $K$ agree on the image of $R$; an
element $c$ of $\kappa$ that the first maps to $r/u$, with $r,u\in R$, satisfies $uc=r$ in $\kappa$, so the second also
maps it to $r/u$. Hence $\gamma_s=\gamma_t$, so $\tau(s)=\tau(t)$ and $s=t$. Let $g$ be a point of $\Gamma$ over the
generic point of $\operatorname{Spec}R$, $k$ an algebraic closure of its residue field and $c:\operatorname{Spec}k\to\Gamma$
the corresponding point. The map $R\to k$ given by $c$ is injective, as $g$ lies over the zero ideal, so it extends to an
embedding $a:K\to k$, and $c$ lies over $\operatorname{Spec}R$ through $\operatorname{Spec}(R\subset K)\circ
\operatorname{Spec}a$. The point $i\circ c$ of $G$ has $(i\circ c)^N=1$, so by hypothesis it is
$\tau(t)\circ\operatorname{Spec}a=i\circ\gamma_t\circ\operatorname{Spec}a$ for some $t$; so $c=\gamma_t\circ
\operatorname{Spec}a$, and $g$ is the image of $\gamma_t$. So $\Gamma$, with the points $\gamma_t$, is a torsion scheme of
type $C$ over $R$, and $i$, which lies over $x$, is a torsion embedding along $x$ of order $N$ with $i\circ\gamma_t=\tau(t)$.
\end{proof}
% Lean: torsion_embedding_geometric_points = BT.Tors.exists_qfGroup_torsion: IsSemiAbelian.pullback,
% SemiAbMulFlat.flat_and_locallyQuasiFinite_pow (Lemma cv-origin-change-projective-local-inputs-multiplication),
% exists_kernel_grp (PointQuot.exists_grpObj_of_mono), kernel_props (the unit a closed immersion, lofp_of_noeth),
% kernel_base_etale (formallyUnramified_mulA = formallyUnramified_mulN with eq_one_of_pow_eq_one,
% Etale.of_formallyUnramified_of_flat) and isReduced_of_smooth_of_isReduced, and exists_of_kernel (toPull and ofPull for the
% points of G_R, lift_of_torsion for gamma_t, eq_of_same_point for pt_inj, exists_geom_point and lift_to_frac for pt_onto).
% About 6,800 lines with Lemma cv-origin-change-projective-local-inputs-multiplication, of which about 460 are its own.

\begin{lemma}\label{cv-origin-change-projective-local-inputs-zero}
Let $R$ be a domain with fraction field $K$, $C$ a finite commutative group, $\Gamma$ a torsion scheme of type $C$ over $R$
(Definition~\ref{cv-origin-change-projective-local-torsion}), $U$ an open subset of $\Gamma$ and $s$ a function on $U$
with $s\circ\gamma_t=0$ for every $t\in C$ such that the image of $\gamma_t$ lies in $U$. Then $s=0$.
\end{lemma}
\leanref{Development/AutomorphicSpectrum/Proofs/GeometryProofs/SectionEqZeroOfValuesStatement.lean}{GlobalTheory.T.section_eq_zero_of_values}
\begin{proof}
Let $U_K=U\times_RK$, an open subset of the generic fibre $\Gamma_K$, with projection $p:U_K\to U$. Composition with $p$
is injective on functions: on an affine open subset $\operatorname{Spec}B$ of $U$ it is the map $B\to B\otimes_RK$, which
is injective as $B$ is flat over $R$ and $R\to K$ is injective, and the functions on $U$ are glued from those on an affine
cover. Let $z$ be a point of $U_K$. Then $p(z)$ lies over the generic point of $\operatorname{Spec}R$, so it is the image
of some $\gamma_t$, which therefore has its image in $U$. As $s\circ\gamma_t=0$, $s$ vanishes at $p(z)$, and $s\circ p$
vanishes at $z$. As $\Gamma_K$ is reduced, a function on $U_K$ that vanishes at every point is zero; so $s\circ p=0$, and
$s=0$.
\end{proof}
% Lean: section_eq_zero_of_values = BT.section_eq_zero_of_val_generic = HtPt.Unc.sec_inj: the generic fibre is
% scheme-theoretically dominant (IsSchemeTheoreticallyDominant.of_isDominant, preserved by the flat base change), so
% app_injective; eq_zero_of_basicOpen_eq_bot on the reduced generic fibre, pt_onto, and Scheme.basicOpen_appLE.

\begin{lemma}\label{cv-origin-change-projective-local-inputs-multiplicative}
Let $R$ be a domain with fraction field $K$, $C$ a finite commutative group, $\Gamma$ a torsion scheme of type $C$ over $R$,
$\chi:C\to R^\times$ a homomorphism and $s$ a function on $\Gamma$ with $s\circ\gamma_t=\chi(t)$ for every $t\in C$. Then
for every scheme $Y$ over $R$ and all points $y,z$ of $\Gamma$ with values in $Y$ over $R$, $s\circ(yz)=(s\circ y)(s\circ z)$,
and $s\circ1=1$ for the unit $1$ of $\Gamma$ with values in $Y$.
\end{lemma}
\leanref{Development/WeilBound/CoreStatements/HodgeTatePoints/HtatHopfMul.lean}{PositivityTheory.HtPt.Htat.ht_char_mul}
\begin{proof}
\emph{Products.} Let $\mu$ be the group law of $\Gamma$ and $p_1,p_2$ the projections of $\Gamma\times_R\Gamma$, and put
$u=s\circ\mu-(s\circ p_1)(s\circ p_2)$, a function on $\Gamma\times_R\Gamma$. It suffices to show $u=0$; composing with
the point $(y,z)$ then gives the claim. The value of $u$ at the point $(\gamma_a,\gamma_b)$ is
$\chi(ab)-\chi(a)\chi(b)=0$, as $t\mapsto\gamma_t$ is a homomorphism. Let $\operatorname{Spec}B_1$ and
$\operatorname{Spec}B_2$ be affine open subsets of $\Gamma$. For $k=1,2$ let $I_k$ be the set of the $t\in C$ for which
$\gamma_t$ has its image in $\operatorname{Spec}B_k$, and $e_t:B_k\to K$ the homomorphism of $R$-algebras it gives. By
Lemma~\ref{cv-origin-change-projective-local-inputs-zero}, an element of $B_k$ at which every $e_t$, $t\in I_k$, vanishes
is zero; so the map $B_k\to K^{I_k}$ is injective. As $B_1$ is flat over $R$, $B_1\otimes_RB_2\to B_1\otimes_RK^{I_2}$ is
injective, and as $K^{I_2}$ is flat over $R$, so is $B_1\otimes_RK^{I_2}\to K^{I_1}\otimes_RK^{I_2}$. The target is the
product, over $(a,b)\in I_1\times I_2$, of copies of $K\otimes_RK$, and multiplication $K\otimes_RK\to K$ is an
isomorphism, as $K$ is a localization of $R$. The resulting map from $B_1\otimes_RB_2$ to $K$ at $(a,b)$ is the value at
$(\gamma_a,\gamma_b)$. So a function on $\operatorname{Spec}B_1\times_R\operatorname{Spec}B_2$ that vanishes at every
$(\gamma_a,\gamma_b)$ with $a\in I_1$ and $b\in I_2$ is zero. These open subsets cover $\Gamma\times_R\Gamma$, so $u=0$.

\emph{The unit.} Let $\epsilon:\operatorname{Spec}R\to\Gamma$ be the unit section. Its composite with
$\operatorname{Spec}(R\subset K)$ is $\gamma_1$, so the image of $s\circ\epsilon\in R$ in $K$ is $\chi(1)=1$, and
$s\circ\epsilon=1$. The unit of $\Gamma$ with values in $Y$ is $\epsilon\circ(Y\to\operatorname{Spec}R)$, so $s\circ1=1$.
\end{proof}
% Lean: ht_char_mul = SeparatedTensorPullback.charValues_mul: mul_univ (sec_prod_zero over the cover
% Pullback.openCoverOfLeftRight, chart_pair, chart_sep with Twhopf.sec_inj, pullback_sep and tensor_sep with
% Algebra.IsEpi.injective_lift_mul) and one_univ (appTop_gen_injective). About 300 lines.

\begin{definition}\label{cv-origin-change-projective-local-inputs-hopf}
Let $R$ be a ring, $C$ a finite commutative group and $R^C$ the $R$-algebra of the functions $C\to R$.
\begin{enumerate}
\item A \emph{Hopf order} in $R^C$ is an $R$-subalgebra $A$ of $R^C$ that is a finitely generated projective $R$-module,
  such that there is a nonzero $d\in R$ with $df\in A$ for every $f\in R^C$, and such that for every $f\in A$ the function
  $(x,y)\mapsto f(xy)$ on $C\times C$ is a finite sum of functions $(x,y)\mapsto a(x)b(y)$ with $a,b\in A$, and
  $x\mapsto f(x^{-1})$ lies in $A$.
\item Let $R$ be a domain and $\Gamma$ a torsion scheme of type $C$ over $R$
  (Definition~\ref{cv-origin-change-projective-local-torsion}). Then $A(\Gamma)$ is the set of the $f\in R^C$ such that
  for every $t\in C$ there is a function $s$ on $\Gamma$ with $s\circ\gamma_h=f(th)$ for every $h\in C$. It is an
  $R$-subalgebra of $R^C$, as sums and products of functions on $\Gamma$ are functions on $\Gamma$.
\end{enumerate}
\end{definition}
\leanref{Development/AutomorphicSpectrum/Definitions/CommutativeAlgebra/TranslateMemIndOrder.lean}{GlobalTheory.CompanionTheory.Ar.BT.IsHopfOrder}
\leanref{Development/AutomorphicSpectrum/Definitions/GlobalGeometry/ExistsUnitFormLocalizationEtaleNbhd.lean}{GlobalTheory.CompanionTheory.Ar.BT.QFGroup.btOrder}

\begin{lemma}\label{cv-origin-change-projective-local-inputs-hopf-characters}
Let $R$ be an integrally closed domain, $C$ a finite commutative group and $A$ a Hopf order in $R^C$
(Definition~\ref{cv-origin-change-projective-local-inputs-hopf}). Then every homomorphism $\chi:C\to R^\times$ lies in $A$.
\end{lemma}
\leanref{Development/AutomorphicSpectrum/Proofs/AlgebraProofs/HopfOrderContainsCharacters.lean}{GlobalTheory.CompanionTheory.Ar.BT.char_mem_hopfOrder}
\begin{proof}
Let $K$ be the fraction field of $R$, $d$ as in the definition, and $\delta_t\in R^C$ the function that is $1$ at $t$ and
$0$ elsewhere. Let $D$ be the $R$-module of the maps $l:C\to K$ with $\sum_tl(t)f(t)\in R$ for every $f\in A$.

\emph{$D$ is finitely generated.} For a linear form $\lambda$ on $A$ put $\psi(\lambda)(t)=\lambda(d\delta_t)/d$. As
$df=\sum_tf(t)\,d\delta_t$, $\sum_t\psi(\lambda)(t)f(t)=\lambda(f)$ for $f\in A$; so $\psi(\lambda)\in D$, and every
$l\in D$ is $\psi(\lambda)$ for the linear form $\lambda(f)=\sum_tl(t)f(t)$. So $D$ is the image of the dual of $A$, which
is finitely generated, as $A$ is finitely generated and projective.

\emph{Convolution.} For $l,m\in D$ let $(l*m)(t)=\sum_xl(x)m(x^{-1}t)$. For $f\in A$ write $f(xy)=\sum_ka_k(x)b_k(y)$ with
$a_k,b_k\in A$. Then $\sum_t(l*m)(t)f(t)=\sum_{x,y}l(x)m(y)f(xy)=\sum_k\bigl(\sum_xl(x)a_k(x)\bigr)\bigl(\sum_ym(y)b_k(y)
\bigr)\in R$, so $l*m\in D$. Put $\langle l\rangle=\sum_tl(t)\chi(t)$. As $\chi$ is a homomorphism, $\langle l*m\rangle=
\langle l\rangle\langle m\rangle$, and $\delta_1\in D$ with $\langle\delta_1\rangle=1$. So the $\langle l\rangle$, $l\in
D$, form a finitely generated $R$-submodule of $K$ that contains $1$ and is closed under products, a ring finite over $R$;
so they are integral over $R$, and they lie in $R$, as $R$ is integrally closed.

\emph{The character.} As $A$ is finitely generated and projective, there are linear maps $g:A\to R^m$ and $p:R^m\to A$
with $p\circ g$ the identity. The function $d\chi$ lies in $A$, and the coordinates $\lambda_k$ of $g$ satisfy
$\lambda_k(d\chi)=\sum_t\psi(\lambda_k)(t)d\chi(t)=d\langle\psi(\lambda_k)\rangle=dr_k$ with $r_k\in R$. So
$d\chi=p(g(d\chi))=d\,p(r)$ with $r=(r_k)$, and $\chi=p(r)\in A$, as $d\neq0$ and $R$ is a domain.
\end{proof}
% Lean: char_mem_hopfOrder: CharExt.dualOrder, psi and range_psi, conv_mem with sum_conv_mul, pairWith_conv,
% isIntegral_pairWith (IsIntegral.of_mem_of_fg), and Module.Finite.exists_comp_eq_id_of_projective. The inverse clause of a
% Hopf order is not used. About 160 lines.

\begin{definition}\label{cv-origin-change-projective-local-inputs-hopf-order-subgroup}
Let $R$ be a domain with fraction field $K$, $C$ a finite commutative group and $\Gamma$ a torsion scheme of type $C$ over $R$
(Definition~\ref{cv-origin-change-projective-local-torsion}). For an open set $U\subset\Gamma$, let $C_U$ be the set of the
$h\in C$ such that the image of $\gamma_h$ lies in $U$.
\begin{enumerate}
\item An open set $U\subset\Gamma$ is a \emph{finite open subgroup} of $\Gamma$ if $U\to\operatorname{Spec}R$ is finite, the
  unit section of $\Gamma$ maps $\operatorname{Spec}R$ into $U$, and for every scheme $Y$ over $R$ and all points $y,z$ of
  $\Gamma$ with values in $Y$ that map $Y$ into $U$, the points $yz$ and $y^{-1}$ map $Y$ into $U$.
\item For an open set $U\subset\Gamma$, $A_U(\Gamma)$ is the set of the $f\in R^C$ such that for every $t\in C$ there is a
  function $s$ on $U$ with $s\circ\gamma_h=f(th)$ for every $h\in C_U$. It is an $R$-subalgebra of $R^C$, and
  $A_\Gamma(\Gamma)=A(\Gamma)$ (Definition~\ref{cv-origin-change-projective-local-inputs-hopf}(2)).
\end{enumerate}
\end{definition}
\leanref{Development/AutomorphicSpectrum/Definitions/GlobalGeometry/ExistsUnitFormLocalizationEtaleNbhd.lean}{GlobalTheory.CompanionTheory.Ar.BT.QFGroup.IsFinOpenSub}
\leanref{Development/AutomorphicSpectrum/Definitions/GlobalGeometry/ExistsUnitFormLocalizationEtaleNbhd.lean}{GlobalTheory.CompanionTheory.Ar.BT.QFGroup.btOrder}
% Lean: btOrder U is the algebra generated by btSet U, the f all of whose translates f(t .) are RegularAt every point of U
% (a function near the point with the values f(th) at the gamma_h that land there). Twhopf.mem_order (btSet is already an
% algebra, btSub) and Twhopf.reg_iff (Lemma cv-origin-change-projective-local-inputs-hopf-order-glue(1)) give the form of the
% text; ValuationInjectivePi.mem_btOrder_iff is the case U = Gamma.

\begin{lemma}\label{cv-origin-change-projective-local-inputs-hopf-order-glue}
Let $R$ be a domain with fraction field $K$, $C$ a finite commutative group, $\Gamma$ a torsion scheme of type $C$ over $R$,
$V\subset\Gamma$ an open set and $f:C\to K$ a map. For an open set $W\subset\Gamma$, say that a function $s$ on $W$ has
values $f$ if $s\circ\gamma_h=f(h)$ for every $h\in C_W$
(Definition~\ref{cv-origin-change-projective-local-inputs-hopf-order-subgroup}).
\begin{enumerate}
\item If every point of $V$ has an open neighbourhood $W\subset V$ with a function on $W$ that has values $f$, then there
  is a function on $V$ that has values $f$.
\item If $V$ is quasi-compact, there are a nonzero $d\in R$ and a function on $V$ that has values $df$.
\end{enumerate}
\end{lemma}
\leanref{Development/WeilBound/CoreStatements/HodgeTatePoints/TwhopfSec.lean}{PositivityTheory.HtPt.Twhopf.reg_iff}
\leanref{Development/AutomorphicSpectrum/Proofs/GeometryProofs/ExistsDenominatorRegular.lean}{GlobalTheory.CompanionTheory.Ar.BT.RegDenom.exists_denom_regularAt}
\leanref{Development/AutomorphicSpectrum/Proofs/GeometryProofs/FiniteOpenHopfOrder.lean}{GlobalTheory.CompanionTheory.Ar.BT.SubOrder.exists_den}
\begin{proof}
(1) Let $s_W$ be the functions on the sets $W$ of the hypothesis. On $W\cap W'$, the function $s_W-s_{W'}$ satisfies
$(s_W-s_{W'})\circ\gamma_h=0$ for every $h\in C_{W\cap W'}$, so it is zero by
Lemma~\ref{cv-origin-change-projective-local-inputs-zero}. So the $s_W$ glue to a function $s$ on $V$. For $h\in C_V$, the
image of $\gamma_h$ is one point, which lies in one of the sets $W$, and $s\circ\gamma_h=s_W\circ\gamma_h=f(h)$.

(2) First let $W\subset\Gamma$ be an affine open set and $B$ its ring. For $h\in C_W$ let $v_h:B\to K$, $b\mapsto b\circ\gamma_h$,
be the homomorphism of $\gamma_h$; on the image of $R$ it is the inclusion $R\subset K$, as $\gamma_h$ lies over
$\operatorname{Spec}K\to\operatorname{Spec}R$. The $v_h$ are pairwise distinct, as a morphism from $\operatorname{Spec}K$ to
the affine scheme $W$ is determined by its homomorphism and the $\gamma_h$ have distinct images. Distinct homomorphisms
from the multiplicative monoid of $B$ to $K$ are linearly independent over $K$ \cite[Chapter~VI, Theorem~4.1]{Lan02}. So a
linear form on $K^{C_W}$ that vanishes on the vectors $(v_h(b))_{h\in C_W}$, $b\in B$, is zero, and these vectors span
$K^{C_W}$. Write the restriction of $f$ to $C_W$ as $\sum_k\lambda_k(v_h(b_k))_h$ with $\lambda_k=a_k/c_k$, where
$a_k,c_k\in R$ and $c_k\neq0$. With $d=\prod_kc_k$ and $b=\sum_ka_k\bigl(\prod_{j\neq k}c_j\bigr)b_k$, the function $b$ on $W$
has values $df$.

In general, every point of $V$ lies in an affine open set of $\Gamma$, and finitely many of these, $W_1,\dots,W_n$, cover
$V$. Let $d_i$ and $b_i$ be as above for $W_i$, and $d=\prod_id_i$. The restriction of $\bigl(\prod_{j\neq i}d_j\bigr)b_i$ to
$W_i\cap V$ has values $df$, and (1) gives a function on $V$ that has values $df$.
\end{proof}
% Lean: (1) is Twhopf.reg_iff (existsUnique_gluing' with the compatibility from Twhopf.sec_inj, which is Lemma
% cv-origin-change-projective-local-inputs-zero). (2): RegDenom.exists_denom_regularAt on an affine open containing a point
% (exists_isAffineOpen_mem_and_subset), with exists_denom_of_injective (v injective through eq_of_valHom_eq and pt_inj;
% linearIndependent_monoidHom, Mathlib's form of Artin's theorem, then the induction on the span that clears denominators);
% SubOrder.exists_den, induction over the compact set with the product of the denominators (regularAt_mul). About 170 lines.

\begin{lemma}\label{cv-origin-change-projective-local-inputs-hopf-order-induced}
Let $R$ be a domain, $C$ a finite commutative group, $H\subset C$ a subgroup and $A$ a Hopf order in $R^H$
(Definition~\ref{cv-origin-change-projective-local-inputs-hopf}).
\begin{enumerate}
\item For $a\in A$ and $t\in H$, the function $h\mapsto a(th)$ lies in $A$.
\item The set $\operatorname{Ind}(A)$ of the $f\in R^C$ such that for every $t\in C$ the function $h\mapsto f(th)$ on $H$
  lies in $A$ is a Hopf order in $R^C$.
\end{enumerate}
\end{lemma}
\leanref{Development/WeilBound/CoreStatements/HodgeTatePoints/TwhopfInd.lean}{PositivityTheory.HtPt.Twhopf.translate_mem}
\leanref{Development/WeilBound/CoreStatements/HodgeTatePoints/TwhopfInd.lean}{PositivityTheory.HtPt.Twhopf.ind_hopf}
\begin{proof}
(1) Write $a(xy)=\sum_ka_k(x)b_k(y)$ with $a_k,b_k\in A$. Then $a(th)=\sum_ka_k(t)b_k(h)$, an $R$-linear combination of the
$b_k$.

(2) The set $\operatorname{Ind}(A)$ is an $R$-subalgebra of $R^C$. For each coset $c\in C/H$ choose a representative $\bar c$,
and for $x\in C$ let $\kappa(x)\in H$ be the element with $x=\bar c\,\kappa(x)$, where $c$ is the coset of $x$; then
$\kappa(xh)=\kappa(x)h$ for $h\in H$. For $c\in C/H$ and $a\in R^H$ let $e_c(a)\in R^C$ be the function that is
$a(\kappa(x))$ at $x\in c$ and $0$ elsewhere. If $a\in A$, then $e_c(a)\in\operatorname{Ind}(A)$: for $t\in C$, the
function $h\mapsto e_c(a)(th)$ is $h\mapsto a(\kappa(t)h)$ if $t\in c$, which lies in $A$ by (1), and $0$ otherwise. Every
$f\in R^C$ is $\sum_ce_c(f_c)$ with $f_c(h)=f(\bar ch)$, and $f_c\in A$ if $f\in\operatorname{Ind}(A)$.

So the $R$-linear maps $\Phi:\operatorname{Ind}(A)\to A^{C/H}$, $f\mapsto(f_c)_c$, and $\Psi:A^{C/H}\to\operatorname{Ind}(A)$,
$(a_c)_c\mapsto\sum_ce_c(a_c)$, satisfy $\Psi\circ\Phi=\mathrm{id}$. Hence $\operatorname{Ind}(A)$ is a direct summand of the
finitely generated projective module $A^{C/H}$, so it is finitely generated and projective. If $d\neq0$ satisfies $da\in A$
for every $a\in R^H$, then every translate $h\mapsto df(th)$ of $df$, for $f\in R^C$, lies in $A$, so $df\in\operatorname{Ind}(A)$.

Let $f\in\operatorname{Ind}(A)$. For $c,c'\in C/H$, the function $h\mapsto f(\bar c\bar c'h)$ lies in $A$, so
$f(\bar c\bar c'hh')=\sum_ka_{c,c',k}(h)\,b_{c,c',k}(h')$ with $a_{c,c',k},b_{c,c',k}\in A$. For $x\in c$ and $y\in c'$,
$xy=\bar c\bar c'\kappa(x)\kappa(y)$, so
\[
f(xy)=\sum_{c,c'\in C/H}\sum_ke_c(a_{c,c',k})(x)\,e_{c'}(b_{c,c',k})(y)
\]
for all $x,y\in C$, with all the $e_c(a_{c,c',k})$ and $e_{c'}(b_{c,c',k})$ in $\operatorname{Ind}(A)$. Finally, for $t\in C$
the function $h\mapsto f((th)^{-1})=f(t^{-1}h^{-1})$ is the composite of $h\mapsto f(t^{-1}h)$, which lies in $A$, with
$h\mapsto h^{-1}$, so it lies in $A$; hence $x\mapsto f(x^{-1})$ lies in $\operatorname{Ind}(A)$.
\end{proof}
% Lean: translate_mem through comul and the span induction; ind_hopf with exists_coset_coord (kappa), ext (e_c), hext_mem,
% hdecomp, Phi and Psi with Module.Projective.of_split, hlift and hsum for the comultiplication, and mul_inv for the
% inverse. About 190 lines.

\begin{lemma}\label{cv-origin-change-projective-local-inputs-hopf-order-finite}
Let $R$ be an integrally closed domain with fraction field $K$, $C$ a finite commutative group, $\Gamma$ a torsion scheme of
type $C$ over $R$ and $U$ a finite open subgroup of $\Gamma$
(Definition~\ref{cv-origin-change-projective-local-inputs-hopf-order-subgroup}). Then $C_U$ is a subgroup of $C$, and
$A_U(\Gamma)$ is a Hopf order in $R^C$.
\end{lemma}
\leanref{Development/WeilBound/CoreStatements/HodgeTatePoints/TwhopfSec.lean}{PositivityTheory.HtPt.Twhopf.sub_order}
\leanref{Development/WeilBound/CoreStatements/HodgeTatePoints/HtatHopfSub.lean}{PositivityTheory.HtPt.Htat.SubgroupHopfOrder.btOrder_isHopfOrder_of_finOpenSub}
\begin{proof}
As $t\mapsto\gamma_t$ is a homomorphism, and $U$ contains the image of the unit section and is stable under products and
inverses of points with values in $\operatorname{Spec}K$, $C_U$ is a subgroup of $C$.

The scheme $U$ is affine, as it is finite over $\operatorname{Spec}R$; let $B$ be its ring, a finite $R$-algebra. As $\Gamma$
is flat and locally of finite presentation over $R$, $B$ is flat over $R$ and of finite presentation as an $R$-algebra, hence a
finitely presented $R$-module \cite[Tag 0564]{Sta}; so $B$ is a finitely generated projective $R$-module
\cite[Tag 00NX]{Sta}. For $h\in C_U$ let $v_h:B\to K$, $b\mapsto b\circ\gamma_h$, a homomorphism of $R$-algebras. As $B$ is
finite over $R$, the elements $v_h(b)$ are integral over $R$, so they lie in $R$. Let $\Phi:B\to R^{C_U}$,
$b\mapsto(v_h(b))_h$, a homomorphism of $R$-algebras; it is injective by
Lemma~\ref{cv-origin-change-projective-local-inputs-zero}. Let $A=\Phi(B)$. Then $f\in R^{C_U}$ lies in $A$ if and only if
there is a function $s$ on $U$ with $s\circ\gamma_h=f(h)$ for every $h\in C_U$.

\emph{$A$ is a Hopf order in $R^{C_U}$.} It is isomorphic to $B$, so it is finitely generated and projective. For $h\in C_U$,
Lemma~\ref{cv-origin-change-projective-local-inputs-hopf-order-glue}(2), applied to the quasi-compact set $U$ and to the
function $\delta_h$ that is $1$ at $h$ and $0$ elsewhere, gives a nonzero $d_h\in R$ with $d_h\delta_h|_{C_U}\in A$. With
$d=\prod_hd_h$, $df=\sum_hf(h)\bigl(\prod_{h'\neq h}d_{h'}\bigr)d_h\delta_h|_{C_U}\in A$ for every $f\in R^{C_U}$.

For the comultiplication, let $x,y:\operatorname{Spec}(B\otimes_RB)\to\Gamma$ be the two projections to
$U=\operatorname{Spec}B$ followed by the inclusion of $U$. They map into $U$, so their product $xy$ does too, and for $s\in B$
the function $c=s\circ(xy)$ lies in $B\otimes_RB$. For $a,b\in C_U$, the homomorphism $v_a\otimes v_b:B\otimes_RB\to K$ gives
a point $w$ of $\operatorname{Spec}(B\otimes_RB)$ with values in $K$ such that $x\circ w=\gamma_a$ and $y\circ w=\gamma_b$; so
$(xy)\circ w=\gamma_a\gamma_b=\gamma_{ab}$ and $v_{ab}(s)=(v_a\otimes v_b)(c)$. Writing $c=\sum_kb_k\otimes b'_k$, this says
$\Phi(s)(ab)=\sum_k\Phi(b_k)(a)\,\Phi(b'_k)(b)$. For the inverse, let $x:\operatorname{Spec}B\to\Gamma$ be the inclusion of
$U$. Then $x^{-1}$ maps into $U$, and $c=s\circ x^{-1}\in B$ satisfies $v_a(c)=s\circ\gamma_a^{-1}=v_{a^{-1}}(s)$, that is,
$\Phi(c)(a)=\Phi(s)(a^{-1})$.

\emph{The induced order.} By the description of $A$, $f\in R^C$ lies in $A_U(\Gamma)$ if and only if for every $t\in C$ the
function $h\mapsto f(th)$ on $C_U$ lies in $A$, that is, $A_U(\Gamma)=\operatorname{Ind}(A)$ in the notation of
Lemma~\ref{cv-origin-change-projective-local-inputs-hopf-order-induced}. So $A_U(\Gamma)$ is a Hopf order in $R^C$ by that
lemma.
\end{proof}
% Lean: sub_order (= SubOrder.hopfOrder_of_finOpenSub): isAffineOpen_of_fin, secRho, sec_finite, sec_projective
% (Module.FinitePresentation.of_finite_of_finitePresentation and Module.Flat.projective_of_finitePresentation), hval by
% IsIntegrallyClosed.isIntegral_iff, Phi, hinj through Twhopf.sec_inj, SubOrder.exists_den with reg_iff for the full rank,
% SubOrder.exists_tensor (the points x, y of Spec(B (x) B), hF.mul, and w from productMap) and SubOrder.exists_inv; subH is C_U.
% Then btOrder_isHopfOrder_of_finOpenSub: btOrder U = indOrder H.subtype A (mem_order, reg_iff) and Twhopf.ind_hopf.
% About 320 lines.

\begin{lemma}\label{cv-origin-change-projective-local-inputs-hopf-order-closed-fibre}
Let $R$ be a local domain, $o$ the closed point of $\operatorname{Spec}R$, $C$ a finite commutative group and $\Gamma$ a torsion
scheme of type $C$ over $R$.
\begin{enumerate}
\item Let $U\subset\Gamma$ be an open set that contains every point of $\Gamma$ over $o$. Then every morphism $Y\to\Gamma$
  over $R$ from a scheme $Y$ finite over $R$ maps $Y$ into $U$. In particular, every open set $V\subset\Gamma$ that is finite
  over $R$ lies in $U$, and if $U$ is finite over $R$, then it is a finite open subgroup of $\Gamma$
  (Definition~\ref{cv-origin-change-projective-local-inputs-hopf-order-subgroup}).
\item If $R$ is Henselian, there is an open set $U\subset\Gamma$, finite over $R$, that contains every point of $\Gamma$ over
  $o$.
\end{enumerate}
\end{lemma}
\leanref{Development/WeilBound/CoreStatements/HodgeTatePoints/TwhopfFin.lean}{PositivityTheory.HtPt.Twhopf.finite_le}
\leanref{Development/WeilBound/CoreStatements/HodgeTatePoints/TwhopfFin.lean}{PositivityTheory.HtPt.Twhopf.finite_sub}
\leanref{Development/AutomorphicSpectrum/Proofs/GeometryProofs/IsFinOpenSubOfFinite.lean}{GlobalTheory.CompanionTheory.Ar.BT.QFGroup.isFinOpenSub_of_closedFibre}
\leanref{Development/WeilBound/CoreStatements/HodgeTatePoints/TwhopfFin.lean}{PositivityTheory.HtPt.Twhopf.finite_part}
\begin{proof}
(1) Let $y:Y\to\Gamma$ be a morphism over $R$ with $Y$ finite over $R$. The set $y^{-1}(U)$ is open in $Y$ and contains every
point of $Y$ over $o$. If it were not all of $Y$, the image in $\operatorname{Spec}R$ of its complement would be a nonempty
closed set, as $Y\to\operatorname{Spec}R$ is finite, hence closed; it would contain $o$, as every nonempty closed subset of the
spectrum of a local ring does; and a point of $Y$ outside $y^{-1}(U)$ would lie over $o$, a contradiction. Applied to the
inclusion of $V$, this gives $V\subset U$.

Let $U$ be finite over $R$, and $\iota:U\to\Gamma$ the inclusion. The unit section maps $\operatorname{Spec}R$ into $U$. Let
$Y$ be a scheme over $R$ and $y,z$ points of $\Gamma$ with values in $Y$ that map $Y$ into $U$, so that $y=\iota\circ y'$ and
$z=\iota\circ z'$ with $y',z':Y\to U$. Then $yz=m\circ(y',z')$, where $m:U\times_RU\to\Gamma$ is the product of the two
projections followed by $\iota$. As $U\times_RU$ is finite over $R$, $m$ maps it into $U$, so $yz$ maps $Y$ into $U$.
Similarly $y^{-1}=\iota^{-1}\circ y'$, and $\iota^{-1}:U\to\Gamma$ maps $U$ into $U$.

(2) Let $x$ be a point of $\Gamma$ over $o$ and $\operatorname{Spec}S\subset\Gamma$ an affine open set containing $x$. Then $S$
is of finite type over $R$ and quasi-finite at $x$. By \cite[Tag 04GJ]{Sta}, $S=A_1\times\cdots\times A_n\times B$, where each
$A_i$ is local and finite over $R$ and $R\to B$ is not quasi-finite at any prime of $B$ over the maximal ideal of $R$. So $x$
lies in some $\operatorname{Spec}A_i$, an open subset of $\Gamma$ that is finite over $R$. The points of $\Gamma$ over $o$ are
finitely many, as $\Gamma$ is quasi-compact and locally quasi-finite over $R$; let $W_1,\dots,W_m$ be open sets finite over
$R$ that contain them, and $U$ their union. The disjoint union of the $W_j$ is finite over $R$ and maps onto $U$, so $U$ is
universally closed over $R$ \cite[Tag 03GN]{Sta}. Its fibres over $\operatorname{Spec}R$ are finite, as images of those of the
disjoint union, and $U$ is separated and locally of finite type over $R$; so $U$ is finite over $R$ \cite[Tag 02LS]{Sta}.
\end{proof}
% Lean: (1) is preimage_eq_top_of_closedFibre_subset (opens_eq_top_of_closedFibre_subset, the closed map of a universally
% closed morphism and specializes_closedPoint), used as Twhopf.mem_of_finite, Twhopf.finite_le and, for the subgroup,
% Twhopf.finite_sub = QFGroup.isFinOpenSub_of_closedFibre (exists_factor_openIncl, the tensor product openOver U (x) openOver U
% and mem_of_finite_of_closedFibre). (2) is Twhopf.finite_part = finitePart (the largest open finite over R) with
% isFinite_finitePart and mem_finitePart_of_eq_closedPoint: exists_idempotent_finite_away_of_henselian, the Lean's own proof
% of the case of Stacks, Tag 04GJ (checked) used here (an idempotent e outside the prime with S_e finite over R, through
% Algebra.exists_etale_isIdempotentElem_forall_liesOver_eq and exists_algHom_lift_of_etale_of_isLocalRing);
% isFinite_sup_iota_comp through isFinite_of_comp_surjective (UniversallyClosed.of_comp_surjective, Stacks, Tag 03GN, checked;
% IsFinite.of_isProper_of_locallyQuasiFinite, Stacks, Tag 02LS, checked). About 300 lines.

\begin{lemma}\label{cv-origin-change-projective-local-inputs-hopf-order-base-change}
Let $R\to R'$ be an injective homomorphism of domains, $K$ and $K'$ their fraction fields and $\kappa:K\to K'$ the induced
homomorphism. Let $C$ be a finite commutative group, $\Gamma$ a torsion scheme of type $C$ over $R$,
$\Gamma'=\Gamma\times_{\operatorname{Spec}R}\operatorname{Spec}R'$ with projection $F:\Gamma'\to\Gamma$, and
$\gamma'_t:\operatorname{Spec}K'\to\Gamma'$, for $t\in C$, the morphism with components $\gamma_t\circ\operatorname{Spec}\kappa$
and $\operatorname{Spec}K'\to\operatorname{Spec}R'$.
\begin{enumerate}
\item The group scheme $\Gamma'$, with the points $\gamma'_t$, is a torsion scheme of type $C$ over $R'$
  (Definition~\ref{cv-origin-change-projective-local-torsion}); $F$ is compatible with the group laws, and
  $F\circ\gamma'_t=\gamma_t\circ\operatorname{Spec}\kappa$.
\item If $U$ is a finite open subgroup of $\Gamma$, then $F^{-1}(U)$ is a finite open subgroup of $\Gamma'$
  (Definition~\ref{cv-origin-change-projective-local-inputs-hopf-order-subgroup}).
\item Let $U\subset\Gamma$ be an open set, $f:C\to K$ a map and $s$ a function on $U$ with $s\circ\gamma_h=f(h)$ for every
  $h\in C_U$. Then $F^*s\circ\gamma'_h=\kappa(f(h))$ for every $h$ such that the image of $\gamma'_h$ lies in $F^{-1}(U)$.
\item Suppose that $R'$ is flat over $R$. Let $g'$ be a point of $\Gamma'$, $f:C\to K$ a map, $U'\subset\Gamma'$ an open
  neighbourhood of $g'$ and $s'$ a function on $U'$ with $s'\circ\gamma'_h=\kappa(f(h))$ for every $h$ such that the image of
  $\gamma'_h$ lies in $U'$. Then there are an open neighbourhood $U$ of $F(g')$ and a function $s$ on $U$ with
  $s\circ\gamma_h=f(h)$ for every $h\in C_U$.
\end{enumerate}
\end{lemma}
\leanref{Development/AutomorphicSpectrum/Proofs/GeometryProofs/ExistsBaseChangeGroupStatement.lean}{GlobalTheory.T.exists_base_change_group}
\leanref{Development/WeilBound/CoreStatements/HomogeneousNonvanishing/BtExistsBc.lean}{PositivityTheory.UnitarySelfGeneration.bt_exists_bc}
\leanref{Development/AutomorphicSpectrum/Proofs/GeometryProofs/BaseChangePointsCompatStatement.lean}{GlobalTheory.T.base_change_points_compat}
\leanref{Development/AutomorphicSpectrum/Proofs/GeometryProofs/IsFinOpenSubBaseChange.lean}{GlobalTheory.CompanionTheory.Ar.BT.isFinOpenSub_preimage_baseChange}
\leanref{Development/AutomorphicSpectrum/Proofs/GeometryProofs/RegularAtBaseChangeStatement.lean}{GlobalTheory.T.regular_at_base_change}
\leanref{Development/WeilBound/CoreStatements/HodgeTatePoints/HtatHopfReg.lean}{PositivityTheory.HtPt.Htat.StalkDivisibility.regularAt_of_baseChange}
\begin{proof}
(1) The scheme $\Gamma'$ is a commutative group scheme over $R'$, flat, locally quasi-finite, separated, quasi-compact and
locally of finite presentation, as these properties pass to base changes, and $F$ is compatible with the group laws. The map
$t\mapsto\gamma'_t$ is a homomorphism, and $F\circ\gamma'_t=\gamma_t\circ\operatorname{Spec}\kappa$ by construction. If
$\gamma'_s$ and $\gamma'_t$ have the same image, then so do $\gamma_s$ and $\gamma_t$, and $s=t$.

For the generic fibre, let $\Gamma_K=\Gamma\times_{\operatorname{Spec}R}\operatorname{Spec}K$, and let
$\delta_t:\operatorname{Spec}K\to\Gamma_K$ be the section given by $\gamma_t$. As $K$ is a localization of $R$, the projection
$\Gamma_K\to\Gamma$ is injective, and its image is the set of points over the generic point of $\operatorname{Spec}R$; so every
point of $\Gamma_K$ is the image of some $\delta_t$, and $\Gamma_K$ has finitely many points. Each $\delta_t$ is a closed
immersion, as a section of the separated morphism $\Gamma_K\to\operatorname{Spec}K$ \cite[Tag 01KT]{Sta}. So every point of
$\Gamma_K$ is closed, and the finite $T_1$ space $\Gamma_K$ is discrete. Hence the image of $\delta_t$ is an open point $W_t$,
and $\delta_t:\operatorname{Spec}K\to W_t$ is a surjective closed immersion into the reduced scheme $W_t$, hence an
isomorphism; so $\delta_t$ is an open immersion. The generic fibre of $\Gamma'$ is
$\Gamma'_{K'}=\Gamma_K\times_{\operatorname{Spec}K}\operatorname{Spec}K'$, and the section $\delta'_t$ of $\Gamma'_{K'}$ given by
$\gamma'_t$ is the base change of $\delta_t$, so it is an open immersion. A point of $\Gamma'_{K'}$ lies over a point
$\delta_t(\operatorname{Spec}K)$ of $\Gamma_K$, and then it is in the image of $\delta'_t$ \cite[Tag 01JT]{Sta}. So every point
of $\Gamma'$ over the generic point of $\operatorname{Spec}R'$ is the image of some $\gamma'_t$, and $\Gamma'_{K'}$ is reduced,
as its local rings are those of $\operatorname{Spec}K'$.

(2) The morphism $F^{-1}(U)\to\operatorname{Spec}R'$ is a base change of $U\to\operatorname{Spec}R$, so it is finite. For a scheme
$Y$ over $R'$, composition with $F$ takes points of $\Gamma'$ with values in $Y$ to points of $\Gamma$ with values in $Y$ over
$R$, compatibly with the unit, products and inverses. So if $y$ and $z$ map $Y$ into $F^{-1}(U)$, then $F\circ(yz)=(F\circ y)
(F\circ z)$ and $F\circ y^{-1}=(F\circ y)^{-1}$ map $Y$ into $U$, and $yz$ and $y^{-1}$ map $Y$ into $F^{-1}(U)$; and the unit
section of $\Gamma'$ maps into $F^{-1}(U)$.

(3) If the image of $\gamma'_h$ lies in $F^{-1}(U)$, then that of $\gamma_h$ lies in $U$, as $F\circ\gamma'_h=\gamma_h\circ
\operatorname{Spec}\kappa$, and $F^*s\circ\gamma'_h=s\circ\gamma_h\circ\operatorname{Spec}\kappa=\kappa(f(h))$.

(4) By Lemma~\ref{cv-origin-change-projective-local-inputs-hopf-order-glue}(2), applied to an affine open neighbourhood $U$ of
$F(g')$, there are a nonzero $d\in R$ and a function $s$ on $U$ with $s\circ\gamma_h=df(h)$ for $h\in C_U$. On
$W'=F^{-1}(U)\cap U'$, the functions $F^*s$ and $ds'$ have the same value $d\kappa(f(h))$ at every $\gamma'_h$ that lands in
$W'$, by (3); so $F^*s=ds'$ on $W'$, by Lemma~\ref{cv-origin-change-projective-local-inputs-zero} for $\Gamma'$. Let
$\mathcal O$ and $\mathcal O'$ be the local rings of $\Gamma$ at $F(g')$ and of $\Gamma'$ at $g'$. The morphism $F$ is flat, as
a base change of the flat morphism $\operatorname{Spec}R'\to\operatorname{Spec}R$, so the local homomorphism
$\mathcal O\to\mathcal O'$ is flat, hence faithfully flat \cite[Tag 00HR]{Sta}. The image in $\mathcal O'$ of the germ of $s$ is
$d$ times the germ of $s'$, so the germ of $s$ lies in $d\mathcal O'\cap\mathcal O=d\mathcal O$ \cite[Tag 05CK]{Sta}. So there
are an open neighbourhood $V\subset U$ of $F(g')$ and a function $s_0$ on $V$ with $s=ds_0$ on $V$. For $h\in C_V$,
$d\,(s_0\circ\gamma_h)=df(h)$, so $s_0\circ\gamma_h=f(h)$, as $d\neq0$ in $K$.
\end{proof}
% Lean: (1) is T.exists_base_change_group = UnitarySelfGeneration.bt_exists_bc: bcGroup (Over.pullback (spMap phi), the
% properties by base change), bcPt and bcPt_mul for gamma'_t, pt_inj through bcPt_fst, pt_onto and gen_reduced through
% bcSec_spec (the square of q, hs, genSec_open: Finite and T1Space of the generic fibre, the closed immersions genSec, and
% isIso_of_isClosedImmersion_of_surjective; genSec_cover with IsPreimmersion.of_isLocalization;
% exists_preimage_of_isPullback), and isGrpPullback_pullback for the group laws; F o gamma'_t is T.base_change_points_compat
% (genericPoint_baseChange). About 270 lines. (2) is isFinOpenSub_preimage_baseChange (pushPt, pushPt_one, pushPt_mul,
% pushPt_inv), about 75 lines. (3) is T.regular_at_base_change = Htat.reg_pull, about 40 lines. (4) is Htat.bt_reg_descent =
% StalkDivisibility.regularAt_of_baseChange: T.exists_denominator_regular, key through Twhopf.sec_inj, dvd_of_stalkMap_dvd
% (Module.FaithfullyFlat.of_flat_of_isLocalHom and Ideal.comap_map_eq_self_of_faithfullyFlat; Stacks, Tags 00HR and 05CK,
% checked), exists_germ_eq and germ_eq. About 130 lines.

\begin{lemma}\label{cv-origin-change-projective-local-inputs-hopf-order-translate}
Let $R$ be an integrally closed domain, $C$ a finite commutative group, $\Gamma$ a torsion scheme of type $C$ over $R$ and
$U'\subset U$ finite open subgroups of $\Gamma$
(Definition~\ref{cv-origin-change-projective-local-inputs-hopf-order-subgroup}). Then $A_U(\Gamma)=A_{U'}(\Gamma)$.
\end{lemma}
\leanref{Development/AutomorphicSpectrum/Proofs/GeometryProofs/BtOrderEqOfLe.lean}{GlobalTheory.CompanionTheory.Ar.BT.QFGroup.btOrder_eq_of_le}
\leanref{Development/AutomorphicSpectrum/Proofs/GeometryProofs/TranslateIntoSmallerOpen.lean}{GlobalTheory.CompanionTheory.Ar.BT.Transl.translates_cover}
\begin{proof}
Restricting functions from $U$ to $U'$ shows $A_U(\Gamma)\subset A_{U'}(\Gamma)$. Conversely, let $f\in A_{U'}(\Gamma)$ and
$t\in C$. By Lemma~\ref{cv-origin-change-projective-local-inputs-hopf-order-glue}(1), it suffices to find for each $g\in U$ an
open neighbourhood $W\subset U$ of $g$ and a function on $W$ whose value at each $\gamma_h$ that lands in $W$ is $f(th)$.

As $\Gamma$ is flat over $R$, $g$ has a generization $g_0$ over the generic point of $\operatorname{Spec}R$
\cite[Tag 03HV]{Sta}, and $g_0$ is the image of some $\gamma_a$. As $U$ is open and contains the specialization $g$ of $g_0$,
it contains $g_0$, so $a\in C_U$, and $b=a^{-1}\in C_U$ by
Lemma~\ref{cv-origin-change-projective-local-inputs-hopf-order-finite}. The scheme $U$ is affine with ring $B$ finite over
$R$, and the homomorphism $B\to K$ of $\gamma_b$ takes values integral over $R$, hence in $R$. So $\gamma_b$ extends to a
section $\sigma:\operatorname{Spec}R\to U\subset\Gamma$, with $\gamma_b=\sigma\circ\operatorname{Spec}(R\subset K)$. Let
$\tau:\Gamma\to\Gamma$ be the translation by $\sigma$, the product of $\sigma$ composed with the structure morphism and the
identity. Then $\tau\circ\gamma_h=\gamma_b\gamma_h=\gamma_{bh}$ for every $h\in C$. In particular $\tau(g_0)$, the image of
$\tau\circ\gamma_a=\gamma_1$, lies in the image of the unit section. As $g_0$ specializes to $g$, $\tau(g)$ lies in the closure
of the image of the unit section, which is closed, as the unit section is a closed immersion \cite[Tag 01KT]{Sta}. So $\tau(g)$
lies in the image of the unit section, which lies in $U'$.

As $f\in A_{U'}(\Gamma)$, there is a function $s'$ on $U'$ with $s'\circ\gamma_h=f(tb^{-1}h)$ for $h\in C_{U'}$. Let
$W=\tau^{-1}(U')\cap U$, an open neighbourhood of $g$. If $\gamma_h$ lands in $W$, then $\tau\circ\gamma_h=\gamma_{bh}$ lands
in $U'$, and $(s'\circ\tau)\circ\gamma_h=s'\circ\gamma_{bh}=f(tb^{-1}bh)=f(th)$. So the restriction of $s'\circ\tau$ to $W$ is the
required function.
\end{proof}
% Lean: QFGroup.btOrder_eq_of_le = btSet_eq_of_le: T.translate_into_smaller_open = Transl.translates_cover
% (Flat.generalizingMap, Stacks, Tag 03HV, checked; pt_onto; hU.inv; T.finite_point_extends_section =
% exists_section_of_generic_point with exists_section_of_isFinite and IsIntegrallyClosed.algebraMap_eq_of_integral; tau =
% (toUnit G >> sigma) * 1, hτpt, the closed range of the unit through UniversallyClosed.of_comp_of_isSeparated, and
% closure_minimal), then T.regular_at_translate (QFGroup.regularAt_translate) with the translate s * t^(-1). About 200 lines.

\begin{lemma}\label{cv-origin-change-projective-local-inputs-hopf-order-henselization}
Let $R$ be an integrally closed domain and $P$ a prime ideal of $R$. Then there are a Henselian local ring $R'$ that is an
integrally closed domain and a flat homomorphism $R\to R'$ under which the inverse image of the maximal ideal of $R'$ is $P$.
\end{lemma}
\leanref{Development/AutomorphicSpectrum/Proofs/AlgebraProofs/ExistsHenselianFlatOverPrimeStatement.lean}{GlobalTheory.T.exists_henselian_flat_over_prime}
\leanref{Development/WeilBound/CoreStatements/HodgeTatePoints/HtatHens.lean}{PositivityTheory.HtPt.Htat.hens_normal}
\begin{proof}
Let $R'$ be the henselization of the local ring $R_P$. It is Henselian, it is a filtered colimit of \'etale $R_P$-algebras, so
it is flat over $R_P$ and hence over $R$, and its maximal ideal is $PR'$ \cite[Tag 04GN]{Sta}; as $R_P\to R'$ is local, the
inverse image in $R$ of the maximal ideal of $R'$ is $P$. As $R_P$ is a normal domain, so is $R'$ \cite[Tag 06DI]{Sta}.
\end{proof}
% Lean: T.exists_henselian_flat_over_prime = CommAlg.exists_flat_henselian_normal_over_prime = Htat.hens_normal, with its own
% construction etaleColimit R kappa(P) (a colimit of etale R-algebras with a point in the residue field of P): isDomain,
% isIntegrallyClosed (stage_int, stage_idem, dvd_of_root), flat, henselianLocalRing, and the maximal ideal through toK. The text
% cites the henselization of R_P instead (Stacks, Tags 04GN and 06DI, checked). About 240 lines.

\begin{lemma}\label{cv-origin-change-projective-local-inputs-hopf-order-all}
Let $R$ be an integrally closed domain, $C$ a finite commutative group, $\Gamma$ a torsion scheme of type $C$ over $R$ and $U$
a finite open subgroup of $\Gamma$ (Definition~\ref{cv-origin-change-projective-local-inputs-hopf-order-subgroup}). Then
$A_U(\Gamma)=A(\Gamma)$.
\end{lemma}
\leanref{Development/WeilBound/CoreStatements/HodgeTatePoints/HtatHopfTop.lean}{PositivityTheory.HtPt.Htat.HopfOrderEqTop.btOrder_eq_top}
\begin{proof}
Restricting functions from $\Gamma$ to $U$ shows $A(\Gamma)\subset A_U(\Gamma)$. Conversely, let $f\in A_U(\Gamma)$ and $t\in C$.
By Lemma~\ref{cv-origin-change-projective-local-inputs-hopf-order-glue}(1), it suffices that every point $g$ of $\Gamma$ have
an open neighbourhood with a function on it whose value at each $\gamma_h$ that lands there is $f(th)$.

Let $P$ be the prime ideal of $R$ under $g$, and $R'$ as in
Lemma~\ref{cv-origin-change-projective-local-inputs-hopf-order-henselization}. The homomorphism $R\to R'$ is injective, as
multiplication by a nonzero element of $R$ is injective on $R$, hence on the flat $R$-module $R'$. Let $\kappa$, $\Gamma'$, $F$
and the $\gamma'_h$ be as in Lemma~\ref{cv-origin-change-projective-local-inputs-hopf-order-base-change}. As $g$ and the closed
point $o'$ of $\operatorname{Spec}R'$ both lie over $P$, there is a point $g'$ of $\Gamma'$ over $o'$ with $F(g')=g$
\cite[Tag 01JT]{Sta}. By Lemma~\ref{cv-origin-change-projective-local-inputs-hopf-order-closed-fibre}(2), there is an open set
$U'\subset\Gamma'$, finite over $R'$, that contains every point of $\Gamma'$ over $o'$, and it is a finite open subgroup of
$\Gamma'$ by (1) of that lemma. By Lemma~\ref{cv-origin-change-projective-local-inputs-hopf-order-base-change}(2),
$F^{-1}(U)$ is a finite open subgroup of $\Gamma'$, so $F^{-1}(U)\subset U'$ by
Lemma~\ref{cv-origin-change-projective-local-inputs-hopf-order-closed-fibre}(1). Let $f'\in R'^C$ be the image of $f$. By
Lemma~\ref{cv-origin-change-projective-local-inputs-hopf-order-base-change}(3), applied to the functions on $U$ that give the
translates of $f$, $f'\in A_{F^{-1}(U)}(\Gamma')$, and so $f'\in A_{U'}(\Gamma')$ by
Lemma~\ref{cv-origin-change-projective-local-inputs-hopf-order-translate}, as $R'$ is integrally closed. So there is a function
on the neighbourhood $U'$ of $g'$ whose value at each $\gamma'_h$ that lands in $U'$ is $f'(th)=\kappa(f(th))$. By
Lemma~\ref{cv-origin-change-projective-local-inputs-hopf-order-base-change}(4), there is a function on a neighbourhood of
$F(g')=g$ whose value at each $\gamma_h$ that lands there is $f(th)$.
\end{proof}
% Lean: HopfOrderEqTop.btOrder_eq_top, with the hypotheses hHenselian, hBaseChange and hRegular: hHenselian at the prime of
% g (Lemma cv-origin-change-projective-local-inputs-hopf-order-henselization), FaithfulSMul.algebraMap_injective, hBaseChange,
% Scheme.exists_preimage_of_isPullback for g', Twhopf.finite_part and finite_sub (U'), T.isFinOpenSub_base_change and
% Twhopf.finite_le (F^(-1) U <= U'), hRegular for h1, T.btOrder_eq_of_le_of_isFinOpenSub, and Htat.bt_reg_descent at g'.
% About 40 lines.

\begin{lemma}\label{cv-origin-change-projective-local-inputs-hopf-order-flat}
Let $R\to R'$ be a flat homomorphism of domains, $C$ a finite commutative group, $\Gamma$ a torsion scheme of type $C$ over $R$
and $\Gamma'$ its base change to $R'$, with the points $\gamma'_t$ of
Lemma~\ref{cv-origin-change-projective-local-inputs-hopf-order-base-change}. Then $A(\Gamma')$ is the $R'$-submodule of $R'^C$
spanned by the images of the elements of $A(\Gamma)$, which is the $R'$-subalgebra that they generate.
\end{lemma}
\leanref{Development/WeilBound/CoreStatements/HodgeTatePoints/HtatHopfFlat.lean}{PositivityTheory.HtPt.Htat.ValuationInjectivePi.btOrder_flat_baseChange}
\begin{proof}
The homomorphism $R\to R'$ is injective, as multiplication by a nonzero element of $R$ is injective on $R$, hence on the flat
$R$-module $R'$. Let $K$, $K'$, $\kappa$ and $F$ be as in
Lemma~\ref{cv-origin-change-projective-local-inputs-hopf-order-base-change}. The images of the elements of $A(\Gamma)$ contain
$1$ and are closed under products, so their $R'$-span is the $R'$-subalgebra that they generate.

For $f\in A(\Gamma)$ and $t\in C$, a function $s$ on $\Gamma$ with $s\circ\gamma_h=f(th)$ for all $h$ gives the function $F^*s$
on $\Gamma'$ with $F^*s\circ\gamma'_h=\kappa(f(th))$, by (3) of that lemma. So the image of $f$ lies in $A(\Gamma')$, and so does
the $R'$-span of these images.

Conversely, let $N\subset K^C$ be the $R$-module of the vectors $(s\circ\gamma_h)_h$ for functions $s$ on $\Gamma$, and for
$t\in C$ let $a_t:R^C\to K^C$, $f\mapsto(f(th))_h$. Then $A(\Gamma)$ is the kernel of
$b:R^C\to\prod_{t\in C}K^C/N$, $f\mapsto(a_t(f)\bmod N)_t$. As $R'$ is flat over $R$, $R'\otimes_RA(\Gamma)$ is the kernel of
$R'\otimes b$. Let $g\in A(\Gamma')$ and $z=\sum_hg(h)\otimes\delta_h\in R'\otimes_RR^C$, where $\delta_h$ is $1$ at $h$ and $0$
elsewhere. Fix $t$, and let $s'$ be a function on $\Gamma'$ with $s'\circ\gamma'_h=g(th)$ for all $h$. As $\Gamma$ is
quasi-compact and quasi-separated and $R'$ is flat over $R$, the functions on $\Gamma'$ are the $R'$-linear combinations of the
functions $F^*s$ with $s$ a function on $\Gamma$ \cite[Tag 02KH]{Sta}. Writing $s'=\sum_kr_kF^*s_k$ gives
$(g(th))_h=\sum_kr_k\kappa(n_k)$ with $n_k=(s_k\circ\gamma_h)_h\in N$. The map $\lambda:R'\otimes_RK^C\to K'^C$,
$r\otimes v\mapsto r\kappa(v)$, is injective, as every element of $R'\otimes_RK$ is $r\otimes e^{-1}$ with $r\in R'$ and $e\in R$
nonzero. Both $(R'\otimes a_t)(z)$ and the image of $\sum_kr_k\otimes n_k\in R'\otimes_RN$ map to $(g(th))_h$ under $\lambda$, so
they are equal, and the class of $(R'\otimes a_t)(z)$ in $R'\otimes_R(K^C/N)$ is zero. Hence $(R'\otimes b)(z)=0$, and $z$ lies in
$R'\otimes_RA(\Gamma)$. The image of $z$ in $R'^C$ is $g$, and the image of $R'\otimes_RA(\Gamma)$ is the $R'$-span of the images
of the elements of $A(\Gamma)$.
\end{proof}
% Lean: btOrder_flat_baseChange, for any Gamma' with IsBaseChange (the text's Gamma' is the one of hBaseChange): flat_inj,
% exists_kappa, btOrder_le_ext (alg_core with secFn = N, mem_btOrder_iff, the exact sequences hexS and hexN from
% Module.Flat.lTensor_exact, inj_pi through inj_one and tensor_frac_form, and val_mem_span: the flat base change of global
% sections, isIso_pushoutSection_of_isQuasiSeparated_of_flat_right, which is Stacks, Tag 02KH, checked, with
% pushout_induction) and ext_le_btOrder (base_change_points_compat). extOrder is the generated subalgebra, equal to the span by
% Htat.ext_span. About 400 lines.

\begin{lemma}\label{cv-origin-change-projective-local-inputs-hopf-order-local}
Let $R$ be a noetherian domain, $C$ a finite commutative group and $A\subset R^C$ an $R$-subalgebra for which there is a nonzero
$d\in R$ with $df\in A$ for every $f\in R^C$. Suppose that for every maximal ideal $P$ of $R$ there are a local domain $R'$ and
a flat homomorphism $\varphi:R\to R'$, under which the inverse image of the maximal ideal of $R'$ is $P$, such that the
$R'$-subalgebra $R'A$ of $R'^C$ generated by the functions $\varphi\circ f$, $f\in A$, is a Hopf order in $R'^C$
(Definition~\ref{cv-origin-change-projective-local-inputs-hopf}). Then $A$ is a Hopf order in $R^C$.
\end{lemma}
\leanref{Development/WeilBound/CoreStatements/HodgeTatePoints/HtatHopfDesc.lean}{PositivityTheory.HtPt.Htat.HopfOrderFlatLocal.isHopfOrder_of_flat_local}
\leanref{Development/WeilBound/CoreStatements/HodgeTatePoints/HtatDesc.lean}{PositivityTheory.HtPt.Htat.mem_descent}
\begin{proof}
\emph{Descent of membership.} Let $I$ be a finite set, $N\subset R^I$ a submodule and $f\in R^I$ such that for every maximal
ideal $P$, with $R'$ and $\varphi$ as in the statement, $\varphi\circ f$ is an $R'$-linear combination of the functions
$\varphi\circ n$, $n\in N$. We show $f\in N$. Let $J\subset R$ be the annihilator of the class $q$ of $f$ in $R^I/N$, and
suppose $J\neq R$. Let $P\supset J$ be a maximal ideal, with $R'$ and $\varphi$. The map $R/J\to R^I/N$, $1\mapsto q$, is
injective, and so is its tensor product with the flat $R$-module $R'$. Under $R'\otimes_RR^I=R'^I$, $1\otimes f$ is
$\varphi\circ f$, an $R'$-linear combination of the elements $1\otimes n$, $n\in N$; so $1\otimes q=0$ in
$R'\otimes_R(R^I/N)$, and $1\otimes1=0$ in $R'\otimes_RR/J=R'/JR'$. But $JR'\subset PR'$ lies in the maximal ideal of $R'$, a
contradiction. So $J=R$ and $f\in N$.

\emph{Finiteness and projectivity.} As $R$ is noetherian, the submodule $A$ of $R^C$ is finitely generated and finitely
presented. Let $P$ be a maximal ideal, with $R'$ and $\varphi$. The map $R'\otimes_RA\to R'^C$ is injective, as $R'$ is flat,
and its image is the $R'$-span of the functions $\varphi\circ f$, $f\in A$, which is $R'A$, as these functions contain $1$ and are
closed under products. As $R'A$ is projective over $R'$, it is flat, so $R'\otimes_RA$ is flat over $R'$. The elements of
$R\setminus P$ map to units of $R'$, so $\varphi$ factors through a homomorphism $R_P\to R'$, which is local and flat
\cite[Tag 00HT]{Sta}, hence faithfully flat \cite[Tag 00HR]{Sta}. As $R'\otimes_{R_P}A_P=R'\otimes_RA$ is flat over $R'$,
$A_P$ is flat over $R_P$ \cite[Tag 03C4]{Sta}. So $A$ is flat over $R$ \cite[Tag 00HT]{Sta}, and as it is finitely presented,
it is finitely generated and projective \cite[Tag 00NX]{Sta}. The condition on $d$ is part of the hypothesis.

\emph{Comultiplication and inverse.} Let $f\in A$, and let $N\subset R^{C\times C}$ be the $R$-span of the functions
$(x,y)\mapsto a(x)b(y)$ with $a,b\in A$. For a maximal ideal $P$, with $R'$ and $\varphi$: as $\varphi\circ f\in R'A$, the
function $(x,y)\mapsto\varphi(f(xy))$ is an $R'$-linear combination of functions $a(x)b(y)$ with $a,b\in R'A$, and each $a$ and
$b$ is an $R'$-linear combination of functions $\varphi\circ a_0$ with $a_0\in A$. So $(x,y)\mapsto\varphi(f(xy))$ is an
$R'$-linear combination of the $\varphi\circ n$, $n\in N$, and by the descent of membership $(x,y)\mapsto f(xy)$ lies in $N$.
Likewise $x\mapsto\varphi(f(x^{-1}))$ lies in $R'A$, the $R'$-span of the $\varphi\circ a_0$, $a_0\in A$, so
$x\mapsto f(x^{-1})$ lies in $A$.
\end{proof}
% Lean: isHopfOrder_of_flat_local: Module.Finite.of_injective and Module.finitePresentation_of_finite; flatness by
% Module.flat_of_localized_maximal with flat_baseChange_of_flat_extOrder (L injective, range = extOrder through ext_span)
% and flat_localized_of_flat_baseChange (Module.flat_iff_of_isLocalization, Module.FaithfullyFlat.of_flat_of_isLocalHom,
% Module.Flat.of_flat_tensorProduct; Stacks, Tags 00HT, 00HR and 03C4, checked); Module.Flat.projective_of_finitePresentation
% (Stacks, Tag 00NX, checked); comul and inv through Htat.mem_descent (the text's descent of membership) with
% prodSpan_extOrder_le and mul_mem_span_prodSpan_of_mem_span. The Lean's IsHopfOrder asks Module.Finite and
% Module.Projective, the text's "finitely generated projective". About 220 lines.

\begin{lemma}\label{cv-origin-change-projective-local-inputs-hopf-order}
Let $R$ be a noetherian integrally closed domain, $C$ a finite commutative group and $\Gamma$ a torsion scheme of type $C$
over $R$. Then $A(\Gamma)$ is a Hopf order in $R^C$ (Definition~\ref{cv-origin-change-projective-local-inputs-hopf}).
\end{lemma}
\leanref{Development/WeilBound/CoreStatements/HodgeTatePoints/HtatHopf.lean}{PositivityTheory.HtPt.Htat.HopfOrderTop.btOrder_top_isHopfOrder}
\begin{proof}
\emph{The Henselian case.} Suppose first that $R$ is a Henselian local ring and an integrally closed domain, not necessarily
noetherian. By Lemma~\ref{cv-origin-change-projective-local-inputs-hopf-order-closed-fibre}(2), there is an open set
$U\subset\Gamma$, finite over $R$, that contains every point of $\Gamma$ over the closed point of $\operatorname{Spec}R$, and $U$
is a finite open subgroup of $\Gamma$ by (1) of that lemma. By
Lemmas~\ref{cv-origin-change-projective-local-inputs-hopf-order-all}
and~\ref{cv-origin-change-projective-local-inputs-hopf-order-finite}, $A(\Gamma)=A_U(\Gamma)$ is a Hopf order in $R^C$.

\emph{Full rank.} Let $K$ be the fraction field of $R$, a Henselian local ring and an integrally closed domain, and $\Gamma_K$
the base change of $\Gamma$ to $K$, a torsion scheme of type $C$ over $K$
(Lemma~\ref{cv-origin-change-projective-local-inputs-hopf-order-base-change}(1)). By the Henselian case, $A(\Gamma_K)$ is a
Hopf order in $K^C$; so there is a nonzero $d'\in K$ with $d'g\in A(\Gamma_K)$ for every $g\in K^C$, and $A(\Gamma_K)=K^C$. As
$K$ is flat over $R$, Lemma~\ref{cv-origin-change-projective-local-inputs-hopf-order-flat} shows that every $g\in K^C$ is a
$K$-linear combination of elements of $A(\Gamma)$. For $g=\delta_t$, the function that is $1$ at $t$ and $0$ elsewhere, write
$\delta_t=\sum_k(a_k/c_k)f_k$ with $f_k\in A(\Gamma)$ and $a_k,c_k\in R$, $c_k\neq0$; then $s_t=\prod_kc_k$ is nonzero and
$s_t\delta_t=\sum_ka_k\bigl(\prod_{j\neq k}c_j\bigr)f_k$ lies in $A(\Gamma)$. With $d=\prod_ts_t$,
$df=\sum_tf(t)\bigl(\prod_{t'\neq t}s_{t'}\bigr)s_t\delta_t\in A(\Gamma)$ for every $f\in R^C$.

\emph{Locally.} Let $P$ be a maximal ideal of $R$, and $R'$ as in
Lemma~\ref{cv-origin-change-projective-local-inputs-hopf-order-henselization}; the homomorphism $R\to R'$ is flat, and it is
injective, as multiplication by a nonzero element of $R$ is injective on the flat $R$-module $R'$. Let $\Gamma'$ be the base
change of $\Gamma$ to $R'$, a torsion scheme of type $C$ over $R'$
(Lemma~\ref{cv-origin-change-projective-local-inputs-hopf-order-base-change}(1)). By the Henselian case, $A(\Gamma')$ is a Hopf
order in $R'^C$, and by Lemma~\ref{cv-origin-change-projective-local-inputs-hopf-order-flat} it is the $R'$-subalgebra generated
by the images of the elements of $A(\Gamma)$. As $R$ is noetherian,
Lemma~\ref{cv-origin-change-projective-local-inputs-hopf-order-local} shows that $A(\Gamma)$ is a Hopf order in $R^C$.
% Lean: btOrder_top_isHopfOrder, with the three statements T.exists_henselian_flat_over_prime (Lemma
% cv-origin-change-projective-local-inputs-hopf-order-henselization), T.exists_base_change_group (part (1) of Lemma
% cv-origin-change-projective-local-inputs-hopf-order-base-change) and T.regular_at_base_change (part (3) of it) as
% hypotheses, which ht_char_section supplies: HopfOrderFlatLocal.isHopfOrder_of_flat_local (Lemma
% cv-origin-change-projective-local-inputs-hopf-order-local); the full rank from the base change to FractionRing R
% (btOrder_flat_baseChange and exists_smul_mem_of_extOrder, the clearing of denominators of the text); at each maximal ideal
% the Henselian case btOrder_top_isHopfOrder' (isFinOpenSub_of_closedFibre with isFinite_finitePart and
% mem_finitePart_of_eq_closedPoint, HopfOrderEqTop.btOrder_eq_top and SubgroupHopfOrder.btOrder_isHopfOrder_of_finOpenSub) for
% the base change of Gamma. About 100 lines at the top; about 2,400 lines with the lemmas before this one.
\end{proof}

\begin{lemma}\label{cv-origin-change-projective-local-inputs-character}
Let $R$ be a noetherian integrally closed domain with fraction field $K$, $C$ a finite commutative group, $\Gamma$ a
torsion scheme of type $C$ over $R$ (Definition~\ref{cv-origin-change-projective-local-torsion}) and $\chi:C\to R^\times$ a
homomorphism. Then there is a function $s$ on $\Gamma$ with $s\circ\gamma_t=\chi(t)$ for every $t\in C$.
\end{lemma}
\leanref{Development/WeilBound/CoreStatements/HomogeneousNonvanishing/HtExists.lean}{PositivityTheory.UnitarySelfGeneration.ht_char_section}
\begin{proof}
By Lemma~\ref{cv-origin-change-projective-local-inputs-hopf-order}, $A(\Gamma)$ is a Hopf order in $R^C$, so by
Lemma~\ref{cv-origin-change-projective-local-inputs-hopf-characters} it contains $\chi$. By
Definition~\ref{cv-origin-change-projective-local-inputs-hopf}(2), with $t=1$, there is a function $s$ on $\Gamma$ with
$s\circ\gamma_h=\chi(h)$ for every $h\in C$.
\end{proof}
% Lean: ht_char_section = Htat.ht_char_section_of = SectionCharValues.exists_section_hasValues_char:
% HopfOrderTop.btOrder_top_isHopfOrder, char_mem_hopfOrder, and ValuationInjectivePi.mem_btOrder_iff (Twhopf.mem_order
% and Twhopf.reg_iff, which glue the local functions of the Lean's btOrder into the global ones of A(Gamma)).

\begin{lemma}\label{cv-origin-change-projective-local-inputs-tangent}
Let $R$ be a domain, $C$ a finite commutative group, $\Gamma$ a torsion scheme of type $C$ over $R$, $S$ a scheme, $G$ a
commutative group scheme over $S$, $x:\operatorname{Spec}R\to S$ a morphism, $N\ge0$ an integer and $i:\Gamma\to G$ a
torsion embedding along $x$ of order $N$ (Definition~\ref{cv-origin-change-projective-local-torsion}). Let $L_0$ be the
Lie algebra of $G$ at $x$, $\pi:R\to R/NR$ the quotient map and $\pi[\varepsilon]$ the map $\pi$ with
$\varepsilon\mapsto\varepsilon$. Then for every $e\in L_0$ there is exactly one morphism
$v:\operatorname{Spec}(R/NR)[\varepsilon]\to\Gamma$ over $R$ with $i\circ v=e\circ\operatorname{Spec}\pi[\varepsilon]$.
\end{lemma}
\leanref{Development/AutomorphicSpectrum/Proofs/GeometryProofs/ExistsUniqueTangentDualNumberStatement.lean}{GlobalTheory.T.existsUnique_tangent_dual_number}
\begin{proof}
Put $\bar e=e\circ\operatorname{Spec}\pi[\varepsilon]$, a point of $G$ over $x\circ(\operatorname{Spec}(R/NR)[\varepsilon]
\to\operatorname{Spec}R)$. As the sum in the Lie algebra is the product of points, $\bar e^N=(Ne)\circ\operatorname{Spec}
\pi[\varepsilon]$, and $Ne=e\circ\operatorname{Spec}\mu_N$ with $\mu_N(a+b\varepsilon)=a+Nb\varepsilon$. The composite
$\pi[\varepsilon]\circ\mu_N$ maps $a+b\varepsilon$ to $\pi(a)$, so it factors through $R[\varepsilon]\to R$, and $\bar e^N$
is the composite of the restriction of $e$ to $\operatorname{Spec}R$, the unit, with $\operatorname{Spec}(R/NR)[\varepsilon]
\to\operatorname{Spec}R$; so $\bar e^N=1$. By Definition~\ref{cv-origin-change-projective-local-torsion}(2),
$\bar e=i\circ v$ for a point $v$ of $\Gamma$ with values in $\operatorname{Spec}(R/NR)[\varepsilon]$ over $R$, and $v$ is
unique, as composition with $i$ is injective on points.
\end{proof}
% Lean: existsUnique_tangent_dual_number = BT.TanLift.tangent_mod_torsion: pt_nsmul, M.smul and dualMap_epsMul for
% e^N = 1, then htors (the kernel clause) and hinj.

\begin{lemma}\label{cv-origin-change-projective-local-inputs-forms}
Let $R$ be a noetherian integrally closed domain, $S$ a scheme, $G$ a commutative group scheme over $S$, smooth over $S$,
and $x:\operatorname{Spec}R\to S$ a morphism. Let $C$ be a finite commutative group, $\Gamma$ a torsion scheme of type $C$
over $R$ and $i:\Gamma\to G$ a torsion embedding along $x$ of order $N\ge0$
(Definition~\ref{cv-origin-change-projective-local-torsion}). Then every homomorphism $\chi:C\to R^\times$ has a
Hodge--Tate form.
\end{lemma}
\leanref{Development/AutomorphicSpectrum/Proofs/GeometryProofs/ExistsHodgeTateFunctionalStatement.lean}{GlobalTheory.T.exists_hodge_tate_functional}
\begin{proof}
Let $L_0$ be the Lie algebra of $G$ at $x$ and $\pi:R\to R/NR$ the quotient map. By
Lemma~\ref{cv-origin-change-projective-local-inputs-character} there is a function $s$ on $\Gamma$ with
$s\circ\gamma_t=\chi(t)$ for every $t\in C$. For $e\in L_0$ let $v_e$ be the morphism of
Lemma~\ref{cv-origin-change-projective-local-inputs-tangent}, and write $s\circ v_e=a_e+l(e)\varepsilon$ with
$a_e,l(e)\in R/NR$.

\emph{The constant term.} The composite of the restriction of $v_e$ to $\operatorname{Spec}(R/NR)$, along
$\varepsilon\mapsto0$, with $i$ is the restriction of $e\circ\operatorname{Spec}\pi[\varepsilon]$, which is the unit, as $e$
restricts to the unit section. As composition with $i$ is an injective homomorphism on points, this restriction of $v_e$
is the unit of $\Gamma$, so $a_e=1$ by Lemma~\ref{cv-origin-change-projective-local-inputs-multiplicative}.

\emph{Linearity.} For $e,e'\in L_0$, $i\circ(v_ev_{e'})=(i\circ v_e)(i\circ v_{e'})=(e+e')\circ\operatorname{Spec}
\pi[\varepsilon]$, as the sum in the Lie algebra is the product; so $v_{e+e'}=v_ev_{e'}$, and by
Lemma~\ref{cv-origin-change-projective-local-inputs-multiplicative}, $s\circ v_{e+e'}=(1+l(e)\varepsilon)(1+l(e')
\varepsilon)=1+(l(e)+l(e'))\varepsilon$. For $c\in R$ let $\mu$ be the endomorphism $a+b\varepsilon\mapsto
a+\pi(c)b\varepsilon$ of $(R/NR)[\varepsilon]$ and $\mu_c$ the endomorphism $a+b\varepsilon\mapsto a+cb\varepsilon$ of
$R[\varepsilon]$, so that $\pi[\varepsilon]\circ\mu_c=\mu\circ\pi[\varepsilon]$. Then $v_e\circ\operatorname{Spec}\mu$ lies
over $R$, and $i\circ v_e\circ\operatorname{Spec}\mu=e\circ\operatorname{Spec}\mu_c\circ\operatorname{Spec}\pi[\varepsilon]
=(ce)\circ\operatorname{Spec}\pi[\varepsilon]$; so $v_{ce}=v_e\circ\operatorname{Spec}\mu$, and $s\circ v_{ce}=\mu(1+l(e)
\varepsilon)=1+\pi(c)l(e)\varepsilon$. So $l:L_0\to R/NR$ is $R$-linear.

\emph{The form.} By Lemma~\ref{cv-origin-reference-lie-bundle-base}, as $G$ is smooth over $S$, $L_0$ is a projective
$R$-module, so $l=\pi\circ w_0$ for an $R$-linear form $w_0$ on $L_0$. Then $s\circ v_e=1+\pi(w_0(e))\varepsilon$ for every
$e\in L_0$, so $w_0$ is a Hodge--Tate form of $\chi$.
\end{proof}
% Lean: exists_hodge_tate_functional = Relay.exists_hodge_tate_functional = UnitarySelfGeneration.ht_exists =
% Htat.ExistsHodgeTate.exists_isHT: ht_char_section_of (Lemma cv-origin-change-projective-local-inputs-character), ht_linear
% = QuotientDualMap.char_on_tangent_linear (fst_valAt with tors_one, isTan_add, isTan_smul, valAt_mul and valAt_one from
% ht_char_mul), lie_proj, Module.projective_lifting_property and existsUnique_tangent_dual_number. About 600 lines besides
% Lemma cv-origin-change-projective-local-inputs-character.

\begin{lemma}\label{cv-origin-change-projective-local-inputs-chart}
Let $T$ be a locally noetherian scheme, $\Omega$ a field of characteristic zero and $\xi:\operatorname{Spec}\Omega\to T$ a
morphism with image $\eta$ such that the local ring at $\eta$ maps bijectively onto $\Omega$, and the local rings at the
specializations of $\eta$ map injectively into $\Omega$ and are integrally closed in $\Omega$. Let $G$ be a semi-abelian
scheme over $T$ with a line bundle $\mathcal L$, and $n\ge1$ an integer such that $G$ has full $n$-torsion at $\xi$ and a
generic Weil pairing $e$ of order $n$ at $\xi$ (Definition~\ref{cv-origin-change-projective-periods}). Let $R\subset\Omega$
be a domain with fraction field $\Omega$ and $x:\operatorname{Spec}R\to T$ an open immersion with
$\xi=x\circ\operatorname{Spec}(R\subset\Omega)$. Let $C$ be a finite commutative group, $\Gamma$ a torsion scheme of type
$C$ over $R$, $i:\Gamma\to G$ a torsion embedding along $x$ of order $n$
(Definition~\ref{cv-origin-change-projective-local-torsion}), $L_0$ the Lie algebra of $G$ at $x$ and
$j:L_0\to\mathrm{Lie}$ the map $v\mapsto v\circ\operatorname{Spec}(R[\varepsilon]\subset\Omega[\varepsilon])$. Let
$y\in G_n$, let $\chi:C\to R^\times$ be a homomorphism with $\chi(t)=e(y,i\circ\gamma_t)$ for every $t\in C$, $w_0$ a
Hodge--Tate form of $\chi$, and $w$ a linear form on $\mathrm{Lie}$ with $w(j(v))=w_0(v)$ for $v\in L_0$. Then $w$ is a
Hodge--Tate period of $y$ at every point of the image of $x$ that specializes $\eta$.
\end{lemma}
\leanref{Development/AutomorphicSpectrum/Proofs/GeometryProofs/HodgeTateWeilOpenChartStatement.lean}{GlobalTheory.T.hodge_tate_weil_open_chart}
\begin{proof}
Let $Q=x(\mathfrak p)$, for a prime $\mathfrak p$ of $R$, be a point of the image of $x$ that specializes $\eta$. As $x$ is an
open immersion, $\mathcal O_{T,Q}$ is the localization $R_{\mathfrak p}$, and its map to $\Omega$ through the local ring at
$\eta$ is the inclusion, as $\xi=x\circ\operatorname{Spec}(R\subset\Omega)$; so $\mathcal O_Q=R_{\mathfrak p}$. Let
$\pi:R\to R/nR$ be the quotient map. As $w_0$ is a Hodge--Tate form of $\chi$, there is a function $s$ on $\Gamma$ with
$s\circ\gamma_t=\chi(t)$ for every $t\in C$ such that every $c\in L_0$ has a morphism
$v_c:\operatorname{Spec}(R/nR)[\varepsilon]\to\Gamma$ over $R$ with $i\circ v_c=c\circ\operatorname{Spec}\pi[\varepsilon]$
and $s\circ v_c=1+\pi(w_0(c))\varepsilon$.

\emph{The function $h$.} Let $G_R=G\times_T\operatorname{Spec}R$, an open subscheme of $G$, as $x$ is an open immersion.
The morphism $i':\Gamma\to G_R$ given by $i$ and the structure morphism of $\Gamma$ is a closed immersion. Indeed, the
projection $G_R\to G$ is a point $u$ of $G$ with values in $G_R$, and the locus $Z$ where $u^n$ and the unit agree is a
closed subscheme of $G_R$, as $G$ is separated over $T$. The point $i$ of $G$ with values in $\Gamma$ has $i^n=1$, so $i'$
factors through a morphism $b:\Gamma\to Z$; conversely the inclusion of $Z$ gives a point of $G$ with $n$-th power $1$, which
is $i\circ y'$ for a point $y':Z\to\Gamma$ over $R$, and $y'$ is inverse to $b$, as composition with $i$ is injective on
points. The image of $Q$ under the unit section of $G$ lies in $G_R$. Let $W$ be an affine open neighbourhood of it in $G_R$.
As $i'$ is a closed immersion, the restriction of $s$ to $i'^{-1}(W)$ is the composite with $i'$ of a function $h$ on $W$.
So $W$ is an open subset of $G$ that contains the image of $Q$ under the unit section, and $h\circ(i\circ y')=s\circ y'$ for
every point $y'$ of $\Gamma$ such that $i\circ y'$ has its image in $W$.

\emph{Derivatives.} Let $e\in\mathrm{Lie}$ extend to $Q$ through a tangent vector $v$ of $G$ at
$\operatorname{Spec}\mathcal O_{T,Q}\to T$. We show that $w(e)\in\mathcal O_Q$ and that, if $v$ has its image in $W$ and
$h\circ v=a+b\varepsilon$, then $w(e)-b\in n\mathcal O_Q$. By Lemma~\ref{cv-origin-reference-lie-bundle-base}, applied to
$R\to R_{\mathfrak p}$, the Lie algebra of $G$ at $\operatorname{Spec}R_{\mathfrak p}\to T$ is $R_{\mathfrak p}\otimes_RL_0$;
so there are $u\in R\setminus\mathfrak p$ and $c\in L_0$ with $uv=c\circ\operatorname{Spec}(R[\varepsilon]\to
R_{\mathfrak p}[\varepsilon])$. Composing with $\operatorname{Spec}(R_{\mathfrak p}[\varepsilon]\subset\Omega[\varepsilon])$
gives $ue=j(c)$, so $uw(e)=w_0(c)$ and $w(e)=w_0(c)/u\in\mathcal O_Q$. Now suppose that $v$ has its image in $W$; then so
has $uv=v\circ\operatorname{Spec}(\varepsilon\mapsto u\varepsilon)$, and $h\circ(uv)=a+ub\varepsilon$. Let
$\pi_{\mathfrak p}:R_{\mathfrak p}\to R_{\mathfrak p}/nR_{\mathfrak p}$ be the quotient map, $\phi:R/nR\to
R_{\mathfrak p}/nR_{\mathfrak p}$ the map induced by the inclusion, and $v'=v_c\circ\operatorname{Spec}\phi[\varepsilon]$.
Then $i\circ v'=(uv)\circ\operatorname{Spec}\pi_{\mathfrak p}[\varepsilon]$, so $\pi_{\mathfrak p}[\varepsilon](h\circ(uv))
=s\circ v'=\phi[\varepsilon](s\circ v_c)=1+\pi_{\mathfrak p}(w_0(c))\varepsilon$. So $w_0(c)-ub\in nR_{\mathfrak p}$, and
$w(e)-b=(w_0(c)-ub)/u\in n\mathcal O_Q$.

\emph{The period.} The vectors that extend to $Q$ generate $\Lambda(Q)$ over $\mathcal O_Q$, so $w\in\Lambda^\vee(Q)$ by
the previous step, which also gives condition (b) of Definition~\ref{cv-origin-change-projective-periods}(3) for $w$, $W$
and $h$. For condition (a), let $k$ be an algebraically closed field, $a:\Omega\to k$ an embedding and $g$ a point of $G$
at $\xi\circ\operatorname{Spec}a$ with $g^n=1$ and image in $W$. By full $n$-torsion $g=g_0\circ\operatorname{Spec}a$ for a
point $g_0$ of $G$ at $\xi$, and $g_0\in G_n$, as $g_0\mapsto g_0\circ\operatorname{Spec}a$ is an injective homomorphism.
By Definition~\ref{cv-origin-change-projective-local-torsion}(2), $g_0=i\circ z$ for a point $z$ of $\Gamma$ with values in
$\Omega$ over $R$. The image of $z$ lies over the generic point of $\operatorname{Spec}R$, so it is the image of some
$\gamma_t$; its residue field embeds into $\Omega$ in only one way over $R$, as $\Omega$ is the fraction field of $R$, so
$z=\gamma_t$ and $g_0=i\circ\gamma_t$, with the same image as $g$. So $h(g)=a(h(g_0))=a(s(\gamma_t))=a(\chi(t))=
a(e(y,g_0))=e_n(y_a,g)$, by the definition of a generic Weil pairing. So $w$ is a Hodge--Tate period of $y$ at $Q$.
\end{proof}
% Lean: hodge_tate_weil_open_chart = HTChart.isHTAt_of_isHT: stalk_loc and hg for O_(T,Q) = R_p;
% BT.isClosedImmersion_torsion_lift (the equalizer of u^n and 1, h.kernel and h.inj) and lift_open for W and h, with hval;
% core_tan with exists_denom (lieModel_exists, T.lie_model_finite_base_change, isLocalizedModule_iff_isBaseChange),
% valAt_mod, valAt_spec and snd_of_mod for the derivatives; for clause (a) pow_descend, HTTie.tors_pt (the kernel clause and
% T.torus_point_over_fraction_field = QFGroup.exists_eq_comp_pt) and hw.weil. About 600 lines.

\begin{lemma}\label{cv-origin-change-projective-local-inputs-unique}
Let $T$ be a locally noetherian scheme, $\Omega$ a field of characteristic zero and $\xi:\operatorname{Spec}\Omega\to T$ a
morphism with image $\eta$ such that the local ring at $\eta$ maps bijectively onto $\Omega$, and the local rings at the
specializations of $\eta$ map injectively into $\Omega$ and are integrally closed in $\Omega$. Let $G$ be a semi-abelian
scheme over $T$ with a line bundle $\mathcal L$, and $n\ge1$ an integer such that $G$ has full $n$-torsion at $\xi$ and a
generic Weil pairing $e$ of order $n$ at $\xi$ (Definition~\ref{cv-origin-change-projective-periods}). Let $P$ be a point
that specializes $\eta$, $y\in G_n$ and $w$ a Hodge--Tate period of $y$ at $P$. Then the Hodge--Tate periods of $y$ at $P$
are the forms $w+nu$ with $u\in\Lambda^\vee(P)$.
\end{lemma}
\leanref{Development/AutomorphicSpectrum/Proofs/GeometryProofs/HodgeTateWeilUniqueStatement.lean}{GlobalTheory.T.hodge_tate_weil_unique}
\begin{proof}
Let $O=\mathcal O_{T,P}$ and $z:\operatorname{Spec}O\to T$ the canonical morphism. The map $O\to\Omega$ through the local
ring at $\eta$ is injective, and it identifies the fraction field of $O$ with $\Omega$, as the local ring at $\eta$, which
maps onto $\Omega$, is a localization of $O$. So $O=\mathcal O_P$ is a noetherian domain of characteristic zero with
fraction field $\Omega$, and $\xi=z\circ\operatorname{Spec}(O\subset\Omega)$.

\emph{A torsion scheme at $P$.} Let $k$ be an algebraically closed field, $a:\Omega\to k$ an embedding and $g$ a point of $G$
at $\xi\circ\operatorname{Spec}a$ with $g^n=1$. By full $n$-torsion $g=g_0\circ\operatorname{Spec}a$ for a point $g_0$ of
$G$ at $\xi$, which lies in $G_n$, as $g_0\mapsto g_0\circ\operatorname{Spec}a$ is an injective homomorphism. So
Lemma~\ref{cv-origin-change-projective-local-inputs-torsion}, applied to $O$, $G$, $z$, the group $C=G_n$ with its inclusion
into the points of $G$ at $\xi$, and $N=n$, gives a torsion scheme $\Gamma$ of type $G_n$ over $O$ and a torsion embedding
$i:\Gamma\to G$ along $z$ of order $n$ with $i\circ\gamma_g=g$ for every $g\in G_n$; the group $G_n$ is finite, as $e$ is a
generic Weil pairing.

\emph{Two periods.} Let $w'$ be a second Hodge--Tate period of $y$ at $P$, and $W$, $h$ and $W'$, $h'$ the open subsets and
functions of Definition~\ref{cv-origin-change-projective-periods}(3) for $w$ and $w'$. For $g\in G_n$ with image in $W$,
$h(g)=e(y,g)$: for an algebraic closure $a:\Omega\to\bar\Omega$, $a(h(g))=h(g\circ\operatorname{Spec}a)=e_n(y_a,g_a)=
a(e(y,g))$, by condition (a) and the definition of a generic Weil pairing. In the same way $h'(g)=e(y,g)$ if $g$ has its
image in $W'$. Put $V=W\cap W'$. The function $(h-h')\circ i$ on $i^{-1}(V)$ vanishes at every $\gamma_g$ with image in
$i^{-1}(V)$, as $i\circ\gamma_g=g$; so it is zero, by Lemma~\ref{cv-origin-change-projective-local-inputs-zero}. Let
$e_0\in\mathrm{Lie}$ extend to $P$ through a tangent vector $v$ of $G$ at $z$. Its image lies in $V$, as $v$ maps every
point of $\operatorname{Spec}O[\varepsilon]$ to a generization of the image of $P$ under the unit section. By
Lemma~\ref{cv-origin-change-projective-local-inputs-tangent}, with $\pi:O\to O/nO$, there is
$v':\operatorname{Spec}(O/nO)[\varepsilon]\to\Gamma$ over $O$ with $i\circ v'=v\circ\operatorname{Spec}\pi[\varepsilon]$.
Write $h\circ v=a+b\varepsilon$ and $h'\circ v=a'+b'\varepsilon$. Then $\pi[\varepsilon]((h-h')\circ v)=((h-h')\circ
i)\circ v'=0$, so $b-b'\in nO$.

\emph{The periods of $y$.} By condition (b), $w(e_0)-b$ and $w'(e_0)-b'$ lie in $n\mathcal O_P$, so $w'(e_0)-w(e_0)\in
n\mathcal O_P$. As $n\neq0$ in $\Omega$, $u=(w'-w)/n$ is a linear form on $\mathrm{Lie}$ with $u(e_0)\in\mathcal O_P$ for
every $e_0$ that extends to $P$. These generate $\Lambda(P)$, so $u\in\Lambda^\vee(P)$, and $w'=w+nu$. Conversely, let
$u\in\Lambda^\vee(P)$. Then $w+nu\in\Lambda^\vee(P)$, and $W$ and $h$ satisfy conditions (a) and (b) for $w+nu$: condition
(a) does not involve the form, and if $e_0$ extends to $P$ then $u(e_0)\in\mathcal O_P$, so $(w+nu)(e_0)-b=(w(e_0)-b)+
nu(e_0)\in n\mathcal O_P$. So $w+nu$ is a Hodge--Tate period of $y$ at $P$.
\end{proof}
% Lean: hodge_tate_weil_unique = HtUniq.isHTAt_iff: exists_frac (exists_frac_genMap); exists_gamma
% (T.torsion_embedding_geometric_points at the stalk, with tau the points of order dividing n moved to the fraction field
% by kappa); val_eq; tangent_congr (T.section_eq_zero_of_values on i^{-1}(W cap W'), stalkLie for the Lie algebra at the
% stalk, T.existsUnique_tangent_dual_number, secVal_comp, secVal_base); tan_le; and the two directions of isHTAt_iff.
% About 540 lines besides Lemmas cv-origin-change-projective-local-inputs-torsion, -zero and -tangent.

\begin{lemma}\label{cv-origin-change-projective-local}
Let $T$ be a locally noetherian scheme, $\Omega$ a field of characteristic zero and $\xi:\operatorname{Spec}\Omega\to T$ a
morphism with image $\eta$ such that the local ring at $\eta$ maps bijectively onto $\Omega$, and the local rings at the
specializations of $\eta$ map injectively into $\Omega$ and are integrally closed in $\Omega$. Let $G$ be a semi-abelian
scheme over $T$ with a line bundle $\mathcal L$, and $n\ge1$ an integer such that $G$ has full $n$-torsion at $\xi$ and a
generic Weil pairing $e$ of order $n$ at $\xi$ (Definition~\ref{cv-origin-change-projective-periods}), and let $P$ be a
point that specializes $\eta$.
\begin{enumerate}
\item $P$ has an open neighbourhood $U$ with, for each $x\in G_n$, a linear form $w_x$ on $\mathrm{Lie}$ that is a
  Hodge--Tate period of $x$ at every point of $U$ that specializes $\eta$.
\item If $w$ is a Hodge--Tate period of $x\in G_n$ at $P$, then the Hodge--Tate periods of $x$ at $P$ are the forms
  $w+nu$ with $u\in\Lambda^\vee(P)$.
\end{enumerate}
\end{lemma}
\leanref{Development/AutomorphicSpectrum/Proofs/GeometryProofs/HodgeTateWeilLocalStatement.lean}{GlobalTheory.T.hodge_tate_weil_local}
\begin{proof}
(2) is Lemma~\ref{cv-origin-change-projective-local-inputs-unique}.

(1) \emph{A chart.} Lemma~\ref{cv-origin-change-projective-local-chart} gives a noetherian integrally closed domain $R$
and an open immersion $x:\operatorname{Spec}R\to T$ whose image $U$ contains $P$ and consists of points that specialize
$\eta$, with an isomorphism of the fraction field of $R$ onto $\Omega$ through which $\xi=x\circ\operatorname{Spec}
(R\subset\Omega)$; we identify the fraction field of $R$ with $\Omega$. The group scheme $G$ is commutative and smooth over
$T$, being semi-abelian, and $G_n$ is finite, as $e$ is a generic Weil pairing.

\emph{The torsion scheme.} Let $k$ be an algebraically closed field, $a:\Omega\to k$ an embedding and $g$ a point of $G$
at $\xi\circ\operatorname{Spec}a$ with $g^n=1$. By full $n$-torsion, $g=t\circ\operatorname{Spec}a$ for a point $t$ of $G$
at $\xi$, and $t\in G_n$, as $t\mapsto t\circ\operatorname{Spec}a$ is an injective homomorphism. So
Lemma~\ref{cv-origin-change-projective-local-inputs-torsion}, applied to $R$, $G$, $x$, the group $C=G_n$ with its inclusion
$\tau$ into the points of $G$ at $\xi$, and $N=n$, gives a torsion scheme $\Gamma$ of type $G_n$ over $R$ and a torsion
embedding $i:\Gamma\to G$ along $x$ of order $n$ with $i\circ\gamma_t=t$ for every $t\in G_n$.

\emph{The forms.} Let $y\in G_n$. As $e$ is multiplicative in the second variable, $t\mapsto e(y,t)$ is a homomorphism
$G_n\to\Omega^\times$, and its values are $n$-th roots of unity, as $e(y,t)^n=e(y,t^n)=e(y,1)=1$. They are integral over
$R$, which is integrally closed with fraction field $\Omega$, so they lie in $R$, and they are units of $R$, as
$e(y,t)^{-1}=e(y,t)^{n-1}$. Let $\chi_y:G_n\to R^\times$ be this homomorphism. Let $L_0$ be the Lie algebra of $G$ at $x$.
By Lemma~\ref{cv-origin-change-projective-local-inputs-forms}, $\chi_y$ has a Hodge--Tate form $w_0$. By
Lemma~\ref{cv-origin-reference-lie-bundle-base}, applied to $R\subset\Omega$, the map $j:L_0\to\mathrm{Lie}$,
$v\mapsto v\circ\operatorname{Spec}(R[\varepsilon]\subset\Omega[\varepsilon])$, is $R$-linear and induces an isomorphism
$\Omega\otimes_RL_0\to\mathrm{Lie}$; so there is a linear form $w_y$ on $\mathrm{Lie}$ with $w_y(j(v))=w_0(v)$ for
$v\in L_0$. As $\chi_y(t)=e(y,t)=e(y,i\circ\gamma_t)$ for every $t\in G_n$,
Lemma~\ref{cv-origin-change-projective-local-inputs-chart} shows that $w_y$ is a Hodge--Tate period of $y$ at every point of
$U$. This gives (1), with the open neighbourhood $U$ of $P$ and the forms $w_y$, $y\in G_n$.
\end{proof}
% Lean: GlobalTheory.T.hodge_tate_weil_local = htLat_exists_lifts: (2) is T.hodge_tate_weil_unique; for (1),
% T.normal_affine_chart_exists (Lemma cv-origin-change-projective-local-chart) and, for each y, HTLift.exists_tors (the
% torsion Tors G xi n transported to the generic point of Spec R by kappa, hgeo from the full torsion, and
% T.torsion_embedding_geometric_points), HTLift.exists_char_units (the values of c = e(y, .) in R, as roots of X^n - 1),
% HTLift.exists_form (lieModel_exists for L0, T.exists_hodge_tate_functional for w0, T.lie_model_base_change_map and
% T.lie_model_finite_base_change for j and the extension w, IsBaseChange.lift) and T.hodge_tate_weil_open_chart. The
% inputs are Lemmas cv-origin-change-projective-local-inputs-torsion, -forms, -chart and -unique; nothing below them remains
% open, as Lemma cv-origin-change-projective-local-inputs-hopf-order (with the Lemmas -hopf-order-* before it) is written out.

\begin{lemma}\label{cv-origin-change-projective-endomorphism}
In the setting of Lemma~\ref{cv-origin-change-projective-local}, let $a$ be an endomorphism of the group scheme $G$ over
$T$, with differential $A:\mathrm{Lie}\to\mathrm{Lie}$, $v\mapsto a\circ v$, and let $x,x'\in G_n$ with
$e(x',g)=e(x,a\circ g)$ for every $g\in G_n$. If $w$ is a Hodge--Tate period of $x$ at $P$, then $w\circ A$ is a
Hodge--Tate period of $x'$ at $P$.
\end{lemma}
\leanref{Development/AutomorphicSpectrum/Proofs/GeometryProofs/HodgeTateWeilEndomorphismStatement.lean}{GlobalTheory.T.hodge_tate_weil_endomorphism}
\begin{proof}
Let $W$ and $h$ be as in Definition~\ref{cv-origin-change-projective-periods}(3) for $w$. If $e\in\mathrm{Lie}$ extends to
$P$ through a tangent vector $v$, then $Ae$ extends to $P$ through $a\circ v$, which is a tangent vector, as $a$ preserves
the unit section. So $A$ maps $\Lambda(P)$ into itself, and $w\circ A\in\Lambda^\vee(P)$. Put $W'=a^{-1}(W)$ and
$h'=h\circ a$; $W'$ contains the image of $P$ under the unit section, as $a$ preserves it.

(a) Let $k$ be an algebraically closed field, $i:\Omega\to k$ an embedding, and $g$ a point of $G$ at
$\xi\circ\operatorname{Spec}i$ of order dividing $n$ with image in $W'$. By full $n$-torsion $g=(g_0)_i$ for a point $g_0$,
which lies in $G_n$ as $g_0\mapsto(g_0)_i$ is injective. As $a\circ g=(a\circ g_0)_i$,
$e_n(x'_i,g)=i(e(x',g_0))=i(e(x,a\circ g_0))=e_n(x_i,a\circ g)$. The point $a\circ g$ has order dividing $n$ and image in
$W$, so $h'(g)=h(a\circ g)=e_n(x_i,a\circ g)=e_n(x'_i,g)$.

(b) Let $e\in\mathrm{Lie}$ extend to $P$ through a tangent vector $v$ with image in $W'$. Then $Ae$ extends to $P$ through
$a\circ v$, with image in $W$, and $h\circ(a\circ v)=h'\circ v$. So $(w\circ A)(e)=w(Ae)$ differs from the image of the
coefficient of $\varepsilon$ in $h'\circ v$ by an element of $n\mathcal O_P$.
\end{proof}
% Lean: LieModel.isHTAt_end: hExt and hlat, the open a^{-1}(W) with the section (a.app W) f, powDescend for g_0, hW.weil
% twice, and the clauses P1 and P2 of w.

\begin{lemma}\label{cv-origin-change-projective-coherent}
In the setting of Lemma~\ref{cv-origin-change-projective-local}, every point that specializes $\eta$ has an open
neighbourhood $U$ and finitely many linear forms that generate $H_n(Q)$ over $\mathcal O_Q$ for every point $Q$ of $U$
that specializes $\eta$. At every point $Q$ that specializes $\eta$, $n\Lambda^\vee(Q)\subset H_n(Q)\subset\Lambda^\vee(Q)$,
and $w\circ A\in H_n(Q)$ for every $w\in H_n(Q)$ and the differential $A$ of every endomorphism of the group scheme $G$.
\end{lemma}
\leanref{Development/AutomorphicSpectrum/Proofs/GeometryProofs/HodgeTateOmegaCoherentStatement.lean}{GlobalTheory.T.hodge_tate_omega_coherent}
\begin{proof}
The Hodge--Tate periods at $Q$ lie in the $\mathcal O_Q$-module $\Lambda^\vee(Q)$, so $H_n(Q)\subset\Lambda^\vee(Q)$.

\emph{Multiples.} Let $Q$ specialize $\eta$ and $u\in\Lambda^\vee(Q)$. By Lemma~\ref{cv-origin-change-projective-local}(1),
applied at $Q$, the unit $1\in G_n$ has a Hodge--Tate period $w_1$ at $Q$, and by (2) so has $w_1+nu$. So their difference
$nu$ lies in $H_n(Q)$.

\emph{Generators.} Let $P$ specialize $\eta$, and let $U_1$ and the forms $w_x$ be as in
Lemma~\ref{cv-origin-change-projective-local}(1). As $G$ is smooth over $T$, Lemma~\ref{cv-origin-reference-lie-bundle}
gives an open neighbourhood $U_2$ of $P$ and a basis $\beta_1,\dots,\beta_r$ of $\mathrm{Lie}^\vee$ with
$\Lambda^\vee(Q)=\sum_l\mathcal O_Q\beta_l$ for every $Q\in U_2$ that specializes $\eta$. Let $Q\in U_1\cap U_2$ specialize
$\eta$. Each $w_x$ is a Hodge--Tate period of $x$ at $Q$, so it lies in $H_n(Q)$, and so does each $n\beta_l$, by the
previous step. Conversely, by Lemma~\ref{cv-origin-change-projective-local}(2) at $Q$, a Hodge--Tate period of $x$ at $Q$
is $w_x+nu$ with $u\in\Lambda^\vee(Q)=\sum_l\mathcal O_Q\beta_l$. So the forms $w_x$, $x\in G_n$, and $n\beta_l$, which are
finitely many as $G_n$ is finite, generate $H_n(Q)$ for every $Q\in U_1\cap U_2$ that specializes $\eta$.

\emph{Endomorphisms.} Let $a$ be an endomorphism of $G$ with differential $A$, and $w$ a Hodge--Tate period of $x\in G_n$
at $Q$. As $a$ is a homomorphism and $e$ is multiplicative in the second variable, $g\mapsto e(x,a\circ g)$ is a
homomorphism $G_n\to\Omega^\times$, so it is $e(x',\cdot)$ for some $x'\in G_n$. By
Lemma~\ref{cv-origin-change-projective-endomorphism}, $w\circ A$ is a Hodge--Tate period of $x'$ at $Q$. As $w\mapsto
w\circ A$ is linear, it maps $H_n(Q)$, which these $w$ generate, into itself.
\end{proof}
% Lean: HTCoh.htOm_coherent: gen; mult (T.hodge_tate_weil_local at q for the unit); for IsCohOn, T.lie_omega_lattice_vector_bundle
% (the frame bv of omLat) and T.hodge_tate_weil_local (the forms w y), with the generators Fin.append b1 b2; htOm_le; and
% hW.perfect with T.hodge_tate_weil_endomorphism for the endomorphisms.

\subsection{Changing the reference geometry}\label{subsection-cv-change-reference-geometry}

\begin{lemma}\label{cv-origin-change-projective-pullback-transport}
Let $\iota:\Omega\to\Omega'$ be an embedding of fields, $C$ and $A$ commutative groups, $n\ge1$ an integer, and $C_n$ and
$A_n$ the sets of the elements $c$ with $c^n=1$. Let $\sim$ be a relation between $A$ and $C$ such that every $c\in C_n$
has some $a\in A_n$ with $a\sim c$, every $a\in A_n$ has some $c\in C_n$ with $a\sim c$, $a\sim c$ and $a'\sim c'$ imply
$aa'\sim cc'$, $a\sim c$ and $a'\sim c$ imply $a=a'$, and $a\sim c$ and $a\sim c'$ imply $c=c'$. Let $e:C\times C\to
\Omega^\times$ and $e':A\times A\to\Omega'^\times$ be maps such that:
\begin{enumerate}
\item[(a)] $e(x,gg')=e(x,g)e(x,g')$ for $x,g,g'\in C_n$, every map $\chi:C\to\Omega^\times$ with
  $\chi(gg')=\chi(g)\chi(g')$ for $g,g'\in C_n$ agrees on $C_n$ with $e(x,\cdot)$ for some $x\in C_n$, and $g\in C_n$ with
  $e(x,g)=1$ for every $x\in C_n$ is $1$;
\item[(b)] $e'(x,gg')=e'(x,g)e'(x,g')$ for $x,g,g'\in A_n$, and every map $\chi:A\to\Omega'^\times$ with
  $\chi(gg')=\chi(g)\chi(g')$ for $g,g'\in A_n$ agrees on $A_n$ with $e'(x,\cdot)$ for some $x\in A_n$.
\end{enumerate}
Then every $y'\in A_n$ has some $y\in C_n$, and every $y\in C_n$ has some $y'\in A_n$, such that
$e'(y',g')=\iota(e(y,g))$ whenever $g\in C_n$, $g'\in A_n$ and $g'\sim g$.
\end{lemma}
\leanref{Development/AutomorphicSpectrum/Proofs/GeometryProofs/WeilPairingTransportRelationStatement.lean}{GlobalTheory.T.weil_pairing_transport_relation}
\begin{proof}
Let $a_0\in A_n$ with $a_0\sim1$. Then $a_0^2\sim1$, so $a_0^2=a_0$ and $a_0=1$; so $1\sim1$, and by induction $a\sim c$
implies $a^k\sim c^k$ for $k\ge0$. Similarly $e(x,1)=e(x,1)^2$, so $e(x,1)=1$ and $e(x,g)^k=e(x,g^k)$ for $x,g\in C_n$;
the same holds for $e'$. Let $m$ be the number of $n$-th roots of unity in $\Omega$.

\emph{From $A$ to $C$.} Let $y'\in A_n$. For $x,g\in C_n$, $e(x,g)^n=e(x,g^n)=1$, so $e(x,g)$ is an $n$-th root of unity,
and $e(x,g^m)=e(x,g)^m=1$, as these roots form a group of order $m$; so $g^m=1$. If $a\in A_n$ and $a\sim c$ with $c\in
C_n$, then $a^m\sim c^m=1$, so $a^m=1$, as $1\sim1$. Hence $e'(y',a)^m=e'(y',a^m)=1$ for $a\in A_n$. The map
$u\mapsto\iota(u)$ from the $n$-th roots of unity of $\Omega$, a group of order $m$, to the $m$-th roots of unity of
$\Omega'$ is injective, and its target has at most $m$ elements, so it is onto. So for $c\in C_n$ there is
$\chi(c)\in\Omega^\times$ with $\iota(\chi(c))=e'(y',a)$, where $a$ is the element of $A_n$ with $a\sim c$; put
$\chi(c)=1$ for $c\notin C_n$. For $g,g'\in C_n$ with $a\sim g$ and $a'\sim g'$, $aa'\sim gg'$, so
$\iota(\chi(gg'))=e'(y',aa')=\iota(\chi(g)\chi(g'))$, and $\chi(gg')=\chi(g)\chi(g')$. By (a), $\chi$ agrees on $C_n$ with
$e(y,\cdot)$ for some $y\in C_n$, and then $e'(y',g')=\iota(\chi(g))=\iota(e(y,g))$ whenever $g'\sim g$.

\emph{From $C$ to $A$.} Let $y\in C_n$. For $a\in A_n$ put $\chi'(a)=\iota(e(y,c))$, where $c$ is the element of $C_n$
with $a\sim c$, and $\chi'(a)=1$ for $a\notin A_n$. As before $\chi'$ is multiplicative on $A_n$, so by (b) it agrees on
$A_n$ with $e'(y',\cdot)$ for some $y'\in A_n$, and $e'(y',g')=\chi'(g')=\iota(e(y,g))$ whenever $g'\sim g$.
\end{proof}
% Lean: weil_pairing_transport_relation = TorsionChars.chars_descend: hrel11, hrelpow, hCm (pow_card_eq_one' in
% rootsOfUnity n Omega and the nondegeneracy), hAm, hval with exists_units_map_eq (card_rootsOfUnity), and eperf, e'perf.

\begin{lemma}\label{cv-origin-change-projective-pullback-derivative}
Let $S$ be a scheme, $G$ a group scheme over $S$, $C$ a ring and $y:\operatorname{Spec}C\to S$ a morphism; tangent vectors
of $G$ at $y$ and its Lie algebra at $y$ are defined as in Definition~\ref{cv-origin-reference-lattices}(1). Let $W$ be an
open subset of $G$ and $h$ a function on $W$. For a morphism $v:\operatorname{Spec}C[\varepsilon]\to G$ with image in $W$
write $h\circ v=a+D(v)\varepsilon$ with $a,D(v)\in C$.
\begin{enumerate}
\item If $v$ and $v'$ are tangent vectors of $G$ at $y$, and $v$, $v'$, their product $v\cdot v'$ and the unit $1$ of $G$
  at $\operatorname{Spec}C[\varepsilon]\to\operatorname{Spec}C\to S$ have images in $W$, then
  $h\circ(v\cdot v')+h\circ1=h\circ v+h\circ v'$ in $C[\varepsilon]$.
\item If $W$ contains the images of all tangent vectors of $G$ at $y$, then $D$ is a $C$-linear map from the Lie algebra
  of $G$ at $y$ to $C$.
\end{enumerate}
\end{lemma}
\leanref{Development/AutomorphicSpectrum/Proofs/GeometryProofs/TangentSectionAdditiveStatement.lean}{GlobalTheory.T.tangent_section_additive}
\leanref{Development/AutomorphicSpectrum/Proofs/GeometryProofs/DualNumberSectionDerivationStatement.lean}{GlobalTheory.T.dual_number_section_derivation}
\begin{proof}
(1) This is a computation. Two elements of $C[\varepsilon]$ are equal if their images in $C_{\mathfrak m}[\varepsilon]$ are,
for every maximal ideal $\mathfrak m$ of $C$; composing with $\operatorname{Spec}C_{\mathfrak m}[\varepsilon]\to
\operatorname{Spec}C[\varepsilon]$, we may assume that $C$ is local. Let $\mu:G\times_SG\to G$ be the group law and $p_1$,
$p_2$ the projections. Then $v\cdot v'$, $v$, $v'$ and $1$ are the composites of $\mu$ with the four morphisms
$u_1=(v,v')$, $u_2=(v,1)$, $u_3=(1,v')$ and $u_4=(1,1)$ from $\operatorname{Spec}C[\varepsilon]$ to $G\times_SG$. These
agree modulo $\varepsilon$, as $v$ and $v'$ restrict to the unit; $u_1$ and $u_2$ have the same composite with $p_1$, as
do $u_3$ and $u_4$, and $u_1$ and $u_3$ have the same composite with $p_2$, as do $u_2$ and $u_4$. As $C[\varepsilon]$ is
local, the four morphisms factor through an affine open subset $Z$ of $\mu^{-1}(W)$ that lies in $p_1^{-1}(W_1)\cap
p_2^{-1}(W_2)$ for affine open subsets $W_1$ and $W_2$ of $G$, and they are given by ring maps
$\alpha_1,\dots,\alpha_4:\Gamma(Z,\mathcal O)\to C[\varepsilon]$ that agree modulo $\varepsilon$. As $(\varepsilon)^2=0$,
$\psi=\alpha_1-\alpha_2-\alpha_3+2\alpha_4$ is a ring map. On the functions pulled back from $W_1$ by $p_1$,
$\alpha_1=\alpha_2$ and $\alpha_3=\alpha_4$, so $\psi=\alpha_4$ there, and in the same way on those pulled back from $W_2$
by $p_2$. So the morphisms given by $\psi$ and $\alpha_4$ have the same composites with $p_1$ and $p_2$, hence are equal,
and $\psi=\alpha_4$. At the function $h\circ\mu$ this says $h\circ(v\cdot v')-h\circ v-h\circ v'+h\circ1=0$.

(2) The unit $1$ at $\operatorname{Spec}C[\varepsilon]$ is a tangent vector, the zero of the Lie algebra, and $h\circ1$
lies in $C$, as $1$ factors through $\operatorname{Spec}C[\varepsilon]\to\operatorname{Spec}C$; so $D(1)=0$. The sum of
two tangent vectors is their product, so (1) gives $D(v+v')=D(v)+D(v')$. For $c\in C$, $cv=v\circ\operatorname{Spec}
\mu_c$ with $\mu_c(a+b\varepsilon)=a+cb\varepsilon$, so $h\circ(cv)=\mu_c(h\circ v)$ and $D(cv)=cD(v)$.
\end{proof}
% Lean: tangent_section_additive = SqAdd.tan_mul_add: gamma_eq_of_localization (maximal ideals) and tan_local, which is
% T.first_order_points_commute = sqzero_add (the four points lift (x, y), lift (x, 1), lift (1, y), lift (1, 1) composed
% with mu; isPair_and_isPair_of_isTriple, localized at the maximal ideals, and the local computation in an affine chart).
% dual_number_section_derivation = HtPull.exists_linear_snd (pt_zero and one_snd for D(1) = 0, M.add with
% tangent_section_additive, M.smul with epsMul_snd).

\begin{lemma}\label{cv-origin-change-projective-pullback-extend}
Let $T$ be a scheme, $\Omega$ a field and $\xi:\operatorname{Spec}\Omega\to T$ a morphism with image $\eta$ such that the
local ring at $\eta$ maps bijectively onto $\Omega$ and those at its specializations injectively. Let $G$ be a
commutative group scheme over $T$ with Lie algebra $\mathrm{Lie}$ at $\xi$, $P$ a point that specializes $\eta$,
$g_P:\mathcal O_{T,P}\to\Omega$ the map through the local ring at $\eta$, $W$ an open subset of $G$ that contains the
image of $P$ under the unit section, and $h$ a function on $W$.
\begin{enumerate}
\item Every tangent vector of $G$ at $\operatorname{Spec}\mathcal O_{T,P}\to T$ has its image in $W$; write
  $h\circ v=a+D(v)\varepsilon$ with $a,D(v)\in\mathcal O_{T,P}$.
\item Let $b_1,\dots,b_r\in\mathrm{Lie}$ generate the Lie algebra $\Lambda(P)$ of $G$ at $P$ over $\mathcal O_P$
  (Definition~\ref{cv-origin-reference-lattices}(1)), with $b_k$ extending to $P$ through a tangent vector $v_k$. If
  $e\in\mathrm{Lie}$ extends to $P$ through a tangent vector $u$, then there are $a_1,\dots,a_r\in\mathcal O_{T,P}$ with
  $e=\sum_kg_P(a_k)b_k$ and $D(u)=\sum_ka_kD(v_k)$.
\end{enumerate}
\end{lemma}
\leanref{Development/AutomorphicSpectrum/Proofs/GeometryProofs/DualMapSnd.lean}{GlobalTheory.CompanionTheory.Ar.HtPull.tanin}
\leanref{Development/AutomorphicSpectrum/Proofs/GeometryProofs/HodgeTateWeilPullback.lean}{GlobalTheory.CompanionTheory.Ar.HtPull.deriv}
\begin{proof}
(1) A point of $\operatorname{Spec}\mathcal O_{T,P}[\varepsilon]$ lies over a point of $\operatorname{Spec}\mathcal
O_{T,P}$, whose image $Q$ in $T$ is a generization of $P$. A tangent vector restricts to the unit section, so it maps the
point to the image of $Q$ under the unit section, a generization of the image of $P$; this lies in the open set $W$.

(2) As $g_P$ is injective, a tangent vector at $\operatorname{Spec}\mathcal O_{T,P}\to T$ is determined by its composite
with $\operatorname{Spec}g_P[\varepsilon]$, and composition with $\operatorname{Spec}g_P[\varepsilon]$ carries sums to
sums and $av$ to $g_P(a)(v\circ\operatorname{Spec}g_P[\varepsilon])$. As $e$ extends to $P$, it lies in
$\Lambda(P)=\sum_k\mathcal O_Pb_k$, so $e=\sum_kg_P(a_k)b_k$ with $a_k\in\mathcal O_{T,P}$. Then $u$ and $\sum_ka_kv_k$ have
the same composite $e$ with $\operatorname{Spec}g_P[\varepsilon]$, so $u=\sum_ka_kv_k$. By (1) and
Lemma~\ref{cv-origin-change-projective-pullback-derivative}(2), applied to $\operatorname{Spec}\mathcal O_{T,P}\to T$, $D$
is $\mathcal O_{T,P}$-linear, and $D(u)=\sum_ka_kD(v_k)$.
\end{proof}
% Lean: tanin (LieCarrier.surjective_fst, range_fromSpecStalk, Specializes.mem_open); deriv = deriv_gen with the Lie model
% LiePull.extModel of the extendable vectors at Spec O_(T,P), T.dual_number_section_derivation for D, the coefficients
% from hl, and LiePull.val_inj and vec_spec.unique for u = sum a_k v_k.

\begin{lemma}\label{cv-origin-change-projective-pullback-lie}
Let $T$ and $T'$ be schemes, $\Omega$ and $\Omega'$ fields, and $\xi:\operatorname{Spec}\Omega\to T$ and
$\xi':\operatorname{Spec}\Omega'\to T'$ morphisms with images $\eta$ and $\eta'$, such that the local rings at $\eta$ and
$\eta'$ map bijectively onto $\Omega$ and $\Omega'$ and those at their specializations injectively. Let
$\iota:\Omega\to\Omega'$ be a field embedding and $f:T'\to T$ a morphism with $f\circ\xi'=\xi\circ\operatorname{Spec}
\iota$, let $G$ be a commutative group scheme over $T$, smooth over $T$, and $G'$ a commutative group scheme over $T'$ with
a homomorphism $\varphi:G'\to G$ over $f$ that identifies $G'$ with $f^*G$. Let $\mathrm{Lie}$ and $\mathrm{Lie}'$ be their
Lie algebras at $\xi$ and $\xi'$, $\Lambda$ and $\Lambda'$ the lattice families of their Lie algebras at the points that
specialize $\eta$ and $\eta'$, $\Lambda^\vee$ and $\Lambda'^\vee$ the dual lattice families
(Definition~\ref{cv-origin-change-projective-periods}(3)), and $j:\mathrm{Lie}\to\mathrm{Lie}'$ an additive map with
$\varphi\circ j(v)=v\circ\operatorname{Spec}\iota[\varepsilon]$ for $v\in\mathrm{Lie}$.
\begin{enumerate}
\item $j(ae)=\iota(a)j(e)$ for $a\in\Omega$, $j$ carries linearly independent families to linearly independent families,
  its image spans $\mathrm{Lie}'$, and $\Lambda'=f^*_j\Lambda$ (Definition~\ref{cv-origin-reference-lattices}(3)).
\item Let $P'$ be a point that specializes $\eta'$, $P=f(P')$, and $\lambda:\mathcal O_{T,P}\to\mathcal O_{T',P'}$ the local
  homomorphism of $f$. There are a basis $b_1,\dots,b_r$ of $\mathrm{Lie}$ with $\Lambda(P)=\sum_k\mathcal O_Pb_k$, such
  that $j(b_1),\dots,j(b_r)$ is a basis of $\mathrm{Lie}'$ with $\Lambda'(P')=\sum_k\mathcal O_{P'}j(b_k)$, and tangent
  vectors $v_k$ of $G$ at $\operatorname{Spec}\mathcal O_{T,P}\to T$ and $v'_k$ of $G'$ at $\operatorname{Spec}\mathcal
  O_{T',P'}\to T'$, through which $b_k$ extends to $P$ and $j(b_k)$ to $P'$, with $\varphi\circ v'_k=v_k\circ
  \operatorname{Spec}\lambda[\varepsilon]$.
\item If $j^\vee:\mathrm{Lie}^\vee\to\mathrm{Lie}'^\vee$ is an additive map with $j^\vee(w)(j(v))=\iota(w(v))$ for
  $v\in\mathrm{Lie}$ and $w\in\mathrm{Lie}^\vee$, then $\Lambda'^\vee=f^*_{j^\vee}\Lambda^\vee$.
\end{enumerate}
\end{lemma}
\leanref{Development/AutomorphicSpectrum/Proofs/GeometryProofs/LieLatticeDominantCompatStatement.lean}{GlobalTheory.T.lie_lattice_dominant_compat}
\leanref{Development/AutomorphicSpectrum/Proofs/GeometryProofs/LieModelPullbackSmoothStatement.lean}{GlobalTheory.T.lie_model_pullback_smooth}
\leanref{Development/AutomorphicSpectrum/Proofs/GeometryProofs/HodgeTateWeilPullback.lean}{GlobalTheory.CompanionTheory.Ar.HtPull.om_pull}
\begin{proof}
(1) and (2) are the argument of the proof of Lemma~\ref{cv-origin-reference-lie-pull}(1), which uses of the two group
schemes only that $G$ is smooth over $T$ and that the morphism identifies $G'$ with $f^*G$; the map $j$ there is the unique
additive map with $\varphi\circ j(e)=e\circ\operatorname{Spec}\iota[\varepsilon]$, as a tangent vector of $G'$ is
determined by its composite with $\varphi$. Its steps \emph{The map $j$} and \emph{The lattices} give (1). For (2), the
step \emph{The lattices} constructs a basis $b_1,\dots,b_r$ over $\mathcal O_{T,P}$ of the vectors that extend to $P$,
extending through tangent vectors $v_k$, so that $\Lambda(P)=\sum_k\mathcal O_Pb_k$; $b$ is a basis of $\mathrm{Lie}$ by
Lemma~\ref{cv-origin-reference-lie-bundle-base}, applied to $g_P$, and so $j(b)$ is a basis of $\mathrm{Lie}'$. It also
gives tangent vectors $v'_k$ of $G'$ with $\varphi\circ v'_k=v_k\circ\operatorname{Spec}\lambda[\varepsilon]$, through which
$j(b_k)$ extends to $P'$, that form a basis of the Lie algebra of $G'$ at $\operatorname{Spec}\mathcal O_{T',P'}\to T'$; so
$\Lambda'(P')=\sum_k\mathcal O_{P'}j(b_k)$.

(3) Let $P'$, $P$, $\lambda$ and $b$ be as in (2), and $(b^\vee_k)$ and $(b'^\vee_k)$ the bases dual to $(b_k)$ and
$(j(b_k))$. Then $j^\vee(b^\vee_k)=b'^\vee_k$, as both take the value $\iota(b^\vee_k(b_l))$ at $j(b_l)$. As
$\Lambda(P)=\sum_k\mathcal O_Pb_k$, $\Lambda^\vee(P)=\sum_k\mathcal O_Pb^\vee_k$, and in the same way
$\Lambda'^\vee(P')=\sum_k\mathcal O_{P'}b'^\vee_k$. With $g_P$ and $g_{P'}$ the maps from the local rings to $\Omega$ and
$\Omega'$, $g_{P'}\circ\lambda=\iota\circ g_P$, as $f\circ\xi'=\xi\circ\operatorname{Spec}\iota$; so $\iota$ maps
$\mathcal O_P$ into $\mathcal O_{P'}$. For $w\in\Lambda^\vee(P)$ the values $j^\vee(w)(j(b_l))=\iota(w(b_l))$ lie in
$\mathcal O_{P'}$, so $j^\vee(w)\in\Lambda'^\vee(P')$; hence $(f^*_{j^\vee}\Lambda^\vee)(P')\subset\Lambda'^\vee(P')$.
Conversely, $w'\in\Lambda'^\vee(P')$ is $\sum_kw'(j(b_k))j^\vee(b^\vee_k)$ with $w'(j(b_k))\in\mathcal O_{P'}$, so it lies
in the module generated by $j^\vee(\Lambda^\vee(P))$.
\end{proof}
% Lean: lie_lattice_dominant_compat = LiePull.main and lie_model_pullback_smooth = LiePull.frame (ext_basis, bc for B and
% B' = j(B), ext_up for the v'_k, bc for their spanning, lat_eq_span); om_pull (hcoord, hdual, and the two inclusions with
% RegAtDom.regAt_dominates).

\begin{situation}\label{cv-origin-change-projective-pullback-situation}
Let $T$ and $T'$ be schemes, $\Omega$ and $\Omega'$ fields, and $\xi:\operatorname{Spec}\Omega\to T$ and
$\xi':\operatorname{Spec}\Omega'\to T'$ morphisms with images $\eta$ and $\eta'$. Let $\iota:\Omega\to\Omega'$ be a field
embedding and $f:T'\to T$ a morphism with $f\circ\xi'=\xi\circ\operatorname{Spec}\iota$. Let $G$ and $G'$ be commutative
group schemes over $T$ and $T'$, with line bundles $\mathcal L$ and $\mathcal L'$, and $\varphi:G'\to G$ a homomorphism over
$f$ that identifies $G'$ with $f^*G$. Let $n\ge0$ be an integer such that $G$ has full $n$-torsion at $\xi$, and let $e$
and $e'$ be generic Weil pairings of order $n$ of $G$ at $\xi$ and of $G'$ at $\xi'$
(Definition~\ref{cv-origin-change-projective-periods}). A point of $G'$ at $\xi'$ \emph{lies over} a point of $G$ at $\xi$
as in Lemma~\ref{cv-origin-change-projective-weil-pullback}. Let $\mathrm{Lie}$ and $\mathrm{Lie}'$ be the Lie algebras of
$G$ at $\xi$ and of $G'$ at $\xi'$, and $j:\mathrm{Lie}\to\mathrm{Lie}'$ and $j^\vee:\mathrm{Lie}^\vee\to\mathrm{Lie}'^\vee$
additive maps with $\varphi\circ j(v)=v\circ\operatorname{Spec}\iota[\varepsilon]$ and $j^\vee(w)(j(v))=\iota(w(v))$ for
$v\in\mathrm{Lie}$ and $w\in\mathrm{Lie}^\vee$, where $\iota[\varepsilon]$ is $\iota$ with $\varepsilon\mapsto\varepsilon$.
\end{situation}

\begin{lemma}\label{cv-origin-change-projective-pullback-torsion}
In Situation~\ref{cv-origin-change-projective-pullback-situation}, let $y\in G_n$ and $y'\in G'_n$ be such that
$e'(y',g')=\iota(e(y,g))$ whenever $g\in G_n$, $g'\in G'_n$ and $g'$ lies over $g$. Let $W$ be an open subset of $G$ and $h$
a function on $W$ with $h(g)=e_n(y\circ\operatorname{Spec}\tau,g)$ for every algebraically closed field $k$, every
embedding $\tau:\Omega\to k$ and every point $g$ of $G$ at $\xi\circ\operatorname{Spec}\tau$ with $g^n=1$ and image in $W$,
where $e_n$ is the commutator pairing of $\mathcal L^n$. Then $h'=h\circ\varphi$, a function on $\varphi^{-1}(W)$, has
$h'(g')=e_n(y'\circ\operatorname{Spec}\tau,g')$ for every algebraically closed field $k$, every embedding
$\tau:\Omega'\to k$ and every point $g'$ of $G'$ at $\xi'\circ\operatorname{Spec}\tau$ with $g'^n=1$ and image in
$\varphi^{-1}(W)$, where $e_n$ is the commutator pairing of $\mathcal L'^n$.
\end{lemma}
\leanref{Development/AutomorphicSpectrum/Proofs/GeometryProofs/WeilPairingPullbackStatement.lean}{GlobalTheory.T.weil_pairing_pullback}
\begin{proof}
By Lemma~\ref{cv-origin-change-projective-weil-pullback}(1), $G'$ has full $n$-torsion at $\xi'$ and every point of
$G'_n$ lies over a point of $G_n$. Let $k$, $\tau$ and $g'$ be as in the statement. By full $n$-torsion,
$g'=g'_0\circ\operatorname{Spec}\tau$ for a point $g'_0$ of $G'$ at $\xi'$, which lies in $G'_n$, as $g'_0\mapsto
g'_0\circ\operatorname{Spec}\tau$ is an injective homomorphism; and $g'_0$ lies over some $g_0\in G_n$. By the last
condition on a generic Weil pairing, for $e'$ and then for $e$,
$e_n(y'\circ\operatorname{Spec}\tau,g')=\tau(e'(y',g'_0))=\tau(\iota(e(y,g_0)))=
e_n(y\circ\operatorname{Spec}(\tau\circ\iota),g_0\circ\operatorname{Spec}(\tau\circ\iota))$. Now
$\varphi\circ g'=\varphi\circ g'_0\circ\operatorname{Spec}\tau=g_0\circ\operatorname{Spec}(\tau\circ\iota)$, a point of $G$
at $\xi\circ\operatorname{Spec}(\tau\circ\iota)$ with $(\varphi\circ g')^n=1$ and image in $W$. So
$h'(g')=h(\varphi\circ g')=e_n(y'\circ\operatorname{Spec}\tau,g')$.
\end{proof}
% Lean: weil_pairing_pullback = HtTorPull.torsion_clause (hFT' for x', cancel_spMap for x'^n = 1, hdown for g0, hGW'.weil
% and hGW.weil, hl and HtPull.seccomp for the value of phi.app W s).

\begin{lemma}\label{cv-origin-change-projective-pullback-periods}
In Situation~\ref{cv-origin-change-projective-pullback-situation}, suppose that the local rings at $\eta$ and $\eta'$ map
bijectively onto $\Omega$ and $\Omega'$ and those at their specializations injectively, and that $G$ is smooth over $T$.
Let $P'$ be a point that specializes $\eta'$, $P=f(P')$, and let $y\in G_n$ and $y'\in G'_n$ be such that
$e'(y',g')=\iota(e(y,g))$ whenever $g\in G_n$, $g'\in G'_n$ and $g'$ lies over $g$. If $w$ is a Hodge--Tate period of $y$
at $P$, then $j^\vee(w)$ is a Hodge--Tate period of $y'$ at $P'$.
\end{lemma}
\leanref{Development/AutomorphicSpectrum/Proofs/GeometryProofs/HodgeTateWeilPullback.lean}{GlobalTheory.CompanionTheory.Ar.HtPull.ht_pull}
\begin{proof}
The point $P$ specializes $f(\eta')=\eta$, as $P'$ specializes $\eta'$. Let $W$ and $h$ be as in
Definition~\ref{cv-origin-change-projective-periods}(3) for $w$, and put $W'=\varphi^{-1}(W)$ and $h'=h\circ\varphi$. As
$\varphi$ is a homomorphism over $f$, it maps the image of $P'$ under the unit section of $G'$ to the image of $P$ under
the unit section of $G$, which lies in $W$; so $W'$ contains the image of $P'$ under the unit section. By
Lemma~\ref{cv-origin-change-projective-pullback-lie}(3), $j^\vee(w)\in\Lambda'^\vee(P')$, as $w\in\Lambda^\vee(P)$.

(a) By Lemma~\ref{cv-origin-change-projective-pullback-torsion}, $h'$ satisfies condition (a) of
Definition~\ref{cv-origin-change-projective-periods}(3) for $y'$ on $W'$.

(b) Let $b_k$, $v_k$ and $v'_k$ be as in Lemma~\ref{cv-origin-change-projective-pullback-lie}(2), and let $g_P$, $g_{P'}$
be the maps from the local rings to $\Omega$ and $\Omega'$, so that $g_{P'}\circ\lambda=\iota\circ g_P$. For tangent
vectors $v$ at $\operatorname{Spec}\mathcal O_{T,P}\to T$ and $v'$ at $\operatorname{Spec}\mathcal O_{T',P'}\to T'$ let
$D(v)$ and $D'(v')$ be the coefficients of $\varepsilon$ in $h\circ v$ and $h'\circ v'$; by
Lemma~\ref{cv-origin-change-projective-pullback-extend}(1) all these tangent vectors have their images in $W$ and $W'$. As
$\varphi\circ v'_k=v_k\circ\operatorname{Spec}\lambda[\varepsilon]$, $h'\circ v'_k=\lambda[\varepsilon](h\circ v_k)$, so
$D'(v'_k)=\lambda(D(v_k))$. By condition (b) for $w$ at $P$, $w(b_k)=g_P(D(v_k))+n\rho_k$ with $\rho_k\in\mathcal O_P$. Let
$e_0\in\mathrm{Lie}'$ extend to $P'$ through a tangent vector $u$. By
Lemma~\ref{cv-origin-change-projective-pullback-extend}(2), applied to $G'$, $P'$, the generators $j(b_k)$ of
$\Lambda'(P')$ and the $v'_k$, $e_0=\sum_kg_{P'}(a_k)j(b_k)$ and $D'(u)=\sum_ka_kD'(v'_k)$ with $a_k\in\mathcal O_{T',P'}$.
As $j^\vee(w)$ is $\Omega'$-linear and $j^\vee(w)(j(b_k))=\iota(w(b_k))$,
\[
j^\vee(w)(e_0)=\sum_kg_{P'}(a_k)\bigl(\iota(g_P(D(v_k)))+n\iota(\rho_k)\bigr)
=g_{P'}(D'(u))+n\sum_kg_{P'}(a_k)\iota(\rho_k),
\]
and $\sum_kg_{P'}(a_k)\iota(\rho_k)\in\mathcal O_{P'}$, as $\iota$ maps $\mathcal O_P$ into $\mathcal O_{P'}$. So
$j^\vee(w)(e_0)$ minus the image of $D'(u)$ lies in $n\mathcal O_{P'}$, which is condition (b) for $j^\vee(w)$ at $P'$.
\end{proof}
% Lean: ht_pull: T.lie_model_pullback_smooth for B, B', v, v' (hvv); RegAtDom.regAt_dominates; om_pull for the lattice
% clause; eta for the unit (heta'); T.weil_pairing_pullback for clause (a); tanin for hk and hk0, deriv for a and hsum,
% comp_ring and dualMap_snd for hcomp, the clause hD of w for rho and hrhoe, RegAtDom.genMap_stalkMap for hg.

\begin{lemma}\label{cv-origin-change-projective-pullback}
Let $T$, $\Omega$, $\xi$, $\eta$, $G$, $\mathcal L$, $n$ and $e$ be as in Lemma~\ref{cv-origin-change-projective-local}.
Let $T'$ be a locally noetherian scheme, $\Omega'$ a field and $\xi':\operatorname{Spec}\Omega'\to T'$ a morphism with
image $\eta'$ such that the local ring at $\eta'$ maps bijectively onto $\Omega'$, and the local rings at the
specializations of $\eta'$ map injectively into $\Omega'$ and are integrally closed in $\Omega'$. Let
$\iota:\Omega\to\Omega'$ be a field embedding and $f:T'\to T$ a morphism with $f\circ\xi'=\xi\circ\operatorname{Spec}\iota$,
let $G'$ be a commutative group scheme over $T'$ with a homomorphism $\varphi:G'\to G$ over $f$ that identifies $G'$ with
$f^*G$, and let $\mathcal L'$ be a line bundle on $G'$ for which $G'$ has a generic Weil pairing $e'$ of order $n$ at
$\xi'$. Let $\mathrm{Lie}'$, $\Lambda'^\vee$ and $H'_n$ be defined for $G'$ as $\mathrm{Lie}$, $\Lambda^\vee$ and $H_n$
are for $G$ in Definition~\ref{cv-origin-change-projective-periods}, and let $j:\mathrm{Lie}\to\mathrm{Lie}'$ and
$j^\vee:\mathrm{Lie}^\vee\to\mathrm{Lie}'^\vee$ be additive maps with $\varphi\circ j(v)=v\circ\operatorname{Spec}
\iota[\varepsilon]$ and $j^\vee(w)(j(v))=\iota(w(v))$ for $v\in\mathrm{Lie}$ and $w\in\mathrm{Lie}^\vee$, where
$\iota[\varepsilon]:\Omega[\varepsilon]\to\Omega'[\varepsilon]$ is $\iota$ with $\varepsilon\mapsto\varepsilon$. Then
$H'_n=f^*_{j^\vee}H_n$ and $\Lambda'^\vee=f^*_{j^\vee}\Lambda^\vee$
(Definition~\ref{cv-origin-reference-lattices}(3)).
\end{lemma}
\leanref{Development/AutomorphicSpectrum/Proofs/GeometryProofs/HodgeTateWeilPullbackStatement.lean}{GlobalTheory.T.hodge_tate_weil_pullback}
\begin{proof}
The data are those of Situation~\ref{cv-origin-change-projective-pullback-situation}. The field $\Omega'$ has
characteristic zero, as $\iota$ is injective. The group scheme $G$ is smooth over $T$, being semi-abelian, and $G'$ is
semi-abelian, as $\varphi$ identifies it with $f^*G$, whose fibres are base changes of those of $G$ to extensions of their
residue fields. By Lemma~\ref{cv-origin-change-projective-weil-pullback}(1), $G'$ has full $n$-torsion at $\xi'$, and the
relation \emph{lies over} between the points of $G'$ and of $G$ satisfies the hypotheses of
Lemma~\ref{cv-origin-change-projective-pullback-transport}, with $C$ and $A$ the groups of the points of $G$ at $\xi$ and
of $G'$ at $\xi'$. Conditions (a) and (b) there hold for the generic Weil pairings $e$ and $e'$. So that lemma shows that
every $y'\in G'_n$ has some $y\in G_n$, and every $y\in G_n$ has some $y'\in G'_n$, with $e'(y',g')=\iota(e(y,g))$
whenever $g\in G_n$, $g'\in G'_n$ and $g'$ lies over $g$; call such $y$ and $y'$ \emph{matched}.

By Lemma~\ref{cv-origin-change-projective-pullback-lie}(1) and (3), $j$ is $\iota$-semilinear, its image spans
$\mathrm{Lie}'$, and $\Lambda'^\vee=f^*_{j^\vee}\Lambda^\vee$, which is the second claim. The map $j^\vee$ is
$\iota$-semilinear: for $a\in\Omega$ and $w\in\mathrm{Lie}^\vee$, the linear forms $j^\vee(aw)$ and $\iota(a)j^\vee(w)$
take the same value $\iota(a)\iota(w(v))$ at every $j(v)$, and these span $\mathrm{Lie}'$.

Let $P'$ be a point that specializes $\eta'$, and $P=f(P')$, which specializes $\eta$. As $f\circ\xi'=\xi\circ
\operatorname{Spec}\iota$, the maps $g_P$ and $g_{P'}$ from the local rings to $\Omega$ and $\Omega'$ satisfy
$g_{P'}\circ\lambda=\iota\circ g_P$ for the local homomorphism $\lambda$ of $f$ at $P'$, so $\iota$ maps $\mathcal O_P$
into $\mathcal O_{P'}$.

\emph{The inclusion $H'_n(P')\subset(f^*_{j^\vee}H_n)(P')$.} Let $w'$ be a Hodge--Tate period of $y'\in G'_n$ at $P'$, and
$y\in G_n$ matched with $y'$. By Lemma~\ref{cv-origin-change-projective-local}(1) at $P$, $y$ has a Hodge--Tate period
$w_y$ at $P$, and by Lemma~\ref{cv-origin-change-projective-pullback-periods}, $j^\vee(w_y)$ is a Hodge--Tate period of
$y'$ at $P'$. Lemma~\ref{cv-origin-change-projective-local}(2), applied to $G'$, $\mathcal L'$ and $e'$, gives
$w'=j^\vee(w_y)+nu$ with $u\in\Lambda'^\vee(P')=(f^*_{j^\vee}\Lambda^\vee)(P')$. For $z\in\Lambda^\vee(P)$, $w_y+nz$ is a
Hodge--Tate period of $y$ at $P$, by Lemma~\ref{cv-origin-change-projective-local}(2) for $G$, so
$nz=(w_y+nz)-w_y\in H_n(P)$. As $nu$ is an $\mathcal O_{P'}$-combination of the forms $j^\vee(nz)$, $z\in\Lambda^\vee(P)$,
both $nu$ and $j^\vee(w_y)$ lie in the module generated by $j^\vee(H_n(P))$, and so does $w'$. These $w'$ generate
$H'_n(P')$.

\emph{The inclusion $(f^*_{j^\vee}H_n)(P')\subset H'_n(P')$.} Let $w$ be a Hodge--Tate period of $y\in G_n$ at $P$, and
$y'\in G'_n$ matched with $y$. By Lemma~\ref{cv-origin-change-projective-pullback-periods}, $j^\vee(w)$ is a Hodge--Tate
period of $y'$ at $P'$, so it lies in $H'_n(P')$. As these $w$ generate $H_n(P)$ over $\mathcal O_P$, $j^\vee$ is
$\iota$-semilinear and $\iota(\mathcal O_P)\subset\mathcal O_{P'}$, $j^\vee(H_n(P))\subset H'_n(P')$, and the module it
generates lies in $H'_n(P')$. So $H'_n=f^*_{j^\vee}H_n$.
\end{proof}
% Lean: GlobalTheory.T.hodge_tate_weil_pullback = HtPull.main: T.semi_abelian_of_pullback (isSemiAbelian_of_isGrpPullback:
% grpPullback_iso and IsSemiAbelian.pullback); T.full_torsion_pullback_points (Lemma
% cv-origin-change-projective-weil-pullback(1)); T.weil_pairing_transport_relation for c1 and c2; om_pull; ht_pull; LiePull.main
% for hspan and hjs; for the first inclusion T.hodge_tate_weil_local at f q', ht_pull and T.hodge_tate_weil_unique for G'
% and, in hnz, for G; for the second, c2, ht_pull and hjs with RegAtDom.regAt_dominates.

\begin{lemma}\label{cv-origin-change-projective-bundles-weil}
In the setting of Definition~\ref{cv-origin-hodge-component}, let $\Omega$ be a field of characteristic zero, $T$ a scheme
with a morphism $\xi:\operatorname{Spec}\Omega\to T$, and $f:T\to\mathcal X$ a morphism with $f\circ\xi=j\circ y_0$ for a
point $y_0$ of $X$ with values in $\Omega$. Let $G$ be an $\mathcal O_E$-semi-abelian scheme over $T$ with a morphism
$G\to\mathcal A$ over $f$, and $n\ge1$ an integer such that $G$ has full $n$-torsion at $\xi$ and $\Omega$ contains a
primitive $n$-th root of unity. Then $G$ has, for its line bundle, a generic Weil pairing of order $n$ at $\xi$
(Definition~\ref{cv-origin-change-projective-periods}).
\end{lemma}
\leanref{Development/AutomorphicSpectrum/Proofs/GeometryProofs/EigenLatticeDedekindTwist.lean}{GlobalTheory.CompanionTheory.Ar.HcVb.exists_genWeil}
\begin{proof}
Lemma~\ref{cv-origin-reference-pullback}(1) gives the $\mathcal O_E$-semi-abelian scheme $G_\Omega=\xi^*G$ over
$\operatorname{Spec}\Omega$ with a morphism $G_\Omega\to G$ over $\xi$, and by (2) its composite with $G\to\mathcal A$ is a
morphism over $f\circ\xi=j\circ y_0$. The morphism $A\to\mathcal A$ over $j$ identifies the abelian scheme $A$ of $X$ with
$j^*\mathcal A$ (Definition~\ref{cv-model}); as cartesian squares can be cancelled, $G_\Omega\to\mathcal A$ is the
composite of a morphism $G_\Omega\to A$ over $y_0$ with $A\to\mathcal A$. The abelian scheme $A$ with its level structures
is an object at level $l$ over $X$ (Definition~\ref{cv-toroidal-moduli-object}(4)), and $G_\Omega$ with the pullbacks of
these level structures is an object at level $l$ over $\operatorname{Spec}\Omega$: properness, the relative ampleness of
the line bundle, the principal polarization and the Tate systems at the geometric points pull back. So
Lemma~\ref{cv-origin-change-projective-weil-pairing}, applied to this object, to $G$, to $\xi$ in place of $f$ and to the
identity of $\operatorname{Spec}\Omega$ in place of $\xi_0$, gives the pairing.
\end{proof}
% Lean: HcVb.exists_genWeil: sabPull at xi (n1), sabComp with m (n2), rewritten by hj, sabCancel with Y.m (n3), and
% T.exists_generic_weil_pairing for the object ModelCompare.objOfSABMap Y.oF n3 over Spec Omega with f = xi and the identity.

\begin{lemma}\label{cv-origin-change-projective-bundles-predual}
Let $T$ be a scheme, $\Omega$ a field and $\xi:\operatorname{Spec}\Omega\to T$ a morphism with image $\eta$; let $V$ be an
$\Omega$-vector space, $W\subset V$ a finite-dimensional subspace, $H$ a lattice family in the dual $V^\vee$
(Definition~\ref{cv-origin-reference-lattices}(2)), $S$ a set of points that specialize $\eta$, and $r\ge0$. Let $H^\circ_W$
be the lattice family in $V$ with $H^\circ_W(Q)=\{v\in W:w(v)\in\mathcal O_Q$ for every $w\in H(Q)\}$. If the restrictions
to $W$ of the forms in $H$ form a vector bundle of rank $r$ on $S$ with generic fibre $W^\vee$, then $H^\circ_W$ is a vector
bundle of rank $r$ on $S$ with generic fibre $W$.
\end{lemma}
\leanref{Development/AutomorphicSpectrum/Proofs/GeometryProofs/PredualVectorBundleStatement.lean}{GlobalTheory.T.predual_vector_bundle}
\begin{proof}
Let $P\in S$, with an open neighbourhood $U$ and a basis $\beta_1,\dots,\beta_r$ of $W^\vee$ such that the restrictions to
$W$ of the forms in $H(Q)$ form $\sum_k\mathcal O_Q\beta_k$ for $Q\in S\cap U$. Let $c_1,\dots,c_r$ be the basis of $W$ dual
to $\beta$, and $Q\in S\cap U$. Each $\beta_k$ is the restriction of some $w_k\in H(Q)$. If $v\in H^\circ_W(Q)$, then
$v=\sum_k\beta_k(v)c_k$ with $\beta_k(v)=w_k(v)\in\mathcal O_Q$. Conversely, for $w\in H(Q)$ the restriction of $w$ to $W$ is
$\sum_ka_k\beta_k$ with $a_k\in\mathcal O_Q$, so $w(c_l)=a_l\in\mathcal O_Q$; hence each $c_l$ lies in $H^\circ_W(Q)$. So
$H^\circ_W(Q)=\sum_k\mathcal O_Qc_k$.
\end{proof}
% Lean: predual_vector_bundle = Lat.predualIn_isVBOn with preDualBasis and Lat.predualIn_isFrameAt.

\begin{lemma}\label{cv-origin-change-projective-bundles-saturate}
Let $T$ be a scheme, $\Omega$ a field and $\xi:\operatorname{Spec}\Omega\to T$ a morphism with image $\eta$; let $V$ be an
$\Omega$-vector space, $W\subset V$ a subspace, $L$ a lattice family in $V$, $S$ a set of points that specialize $\eta$, and
$r\ge0$. Let $O$ be a Dedekind domain with an injective homomorphism $\sigma:O\to\Omega$ whose values are regular at every
point of $S$, and $\mathfrak d$ a nonzero ideal of $O$. Let $L_{\mathfrak d}$ be the lattice family for which
$L_{\mathfrak d}(Q)$ is the $\mathcal O_Q$-module generated by the $e\in W$ with $\sigma(d)e\in L(Q)$ for every
$d\in\mathfrak d$. If $L$ is a vector bundle of rank $r$ on $S$ with generic fibre $W$, then so is $L_{\mathfrak d}$.
\end{lemma}
\leanref{Development/AutomorphicSpectrum/Proofs/GeometryProofs/IdealSaturatedVectorBundleStatement.lean}{GlobalTheory.T.ideal_saturated_vector_bundle}
\begin{proof}
Let $P\in S$, with an open neighbourhood $U$ and a basis $b$ of $W$ such that $L(Q)=\sum_k\mathcal O_Qb_k$ for $Q\in S\cap
U$. The ring $\mathcal O_P$ is local, as the image of the local ring of $T$ at $P$.

\emph{A local generator.} There are $d\in\mathfrak d$ and $u\in\Omega$ such that $\sigma(d')u\in\sigma(O)$ for every
$d'\in\mathfrak d$ and $\sigma(d)u$ is a unit at $P$. Indeed, let $x\in\mathfrak d$ be nonzero, and $\mathfrak a$ an ideal
of $O$ with $\mathfrak d\mathfrak a=xO$, which exists as $O$ is a Dedekind domain. For $a\in\mathfrak a$ and
$d'\in\mathfrak d$, $d'a\in xO$, so $\sigma(d')\sigma(a)/\sigma(x)\in\sigma(O)$, which lies in $\mathcal O_P$. If
$\sigma(d)\sigma(a)/\sigma(x)$ were a non-unit of $\mathcal O_P$ for all $d\in\mathfrak d$ and $a\in\mathfrak a$, then, as
$x$ is a sum of products $da$, $1=\sigma(x)/\sigma(x)$ would lie in the maximal ideal of $\mathcal O_P$. So some $d$ and
$u=\sigma(a)/\sigma(x)$ will do.

\emph{The basis.} Put $v=(\sigma(d)u)^{-1}$, which is regular at $P$, hence at the points of $S$ in an open neighbourhood
$U'\subset U$ of $P$ (Lemma~\ref{cv-origin-segre-regular}(1)). As $u\neq0$, $ub_1,\dots,ub_r$ is a basis of $W$. Let $Q\in
S\cap U'$. If $e\in W$ and $\sigma(d')e\in L(Q)$ for every $d'\in\mathfrak d$, then $\sigma(d)e=\sum_kc_kb_k$ with
$c_k\in\mathcal O_Q$, and $e=uv\sigma(d)e=\sum_k(vc_k)(ub_k)$ with $vc_k\in\mathcal O_Q$. Conversely, for $d'\in\mathfrak d$,
$\sigma(d')ub_k=\sigma(o)b_k$ for some $o\in O$, and $\sigma(o)\in\mathcal O_Q$, so $\sigma(d')ub_k\in L(Q)$ and $ub_k\in
L_{\mathfrak d}(Q)$. Hence $L_{\mathfrak d}(Q)=\sum_k\mathcal O_Qub_k$.
\end{proof}
% Lean: ideal_saturated_vector_bundle = Relay.ideal_saturated_vector_bundle: exists_mul_eq_one_of_dedekind_ideal_ne_bot
% (Ideal.dvd_iff_le for the ideal a with d a = (x), Submodule.mul_induction_on for the contradiction),
% exists_open_forall_mem_regAt' for v, and the frame u b.

\begin{lemma}\label{cv-origin-change-projective-bundles-generic}
Let $T$ be a scheme, $\Omega$ a field and $\xi:\operatorname{Spec}\Omega\to T$ a morphism with image $\eta$; let $G$ be a
commutative group scheme over $T$ with a line bundle $\mathcal L$, $n\ge0$ an integer, and $\mathrm{Lie}$, $\Lambda$,
$\Lambda^\vee$ and $H_n$ as in Definition~\ref{cv-origin-change-projective-periods}(3). Let $O$ be a commutative ring with homomorphisms
$\rho:O\to\operatorname{End}_\Omega\mathrm{Lie}$ and $\sigma:O\to\Omega$, $\mathrm{Lie}_\sigma$ the subspace on which each
$\rho(a)$ acts by $\sigma(a)$, and $Q$ a point that specializes $\eta$ such that $\Lambda$ is a vector bundle on $\{Q\}$
with generic fibre $\mathrm{Lie}$ and $H_n(Q)=\Lambda^\vee(Q)$. Let $H^\circ(Q)=\{e\in\mathrm{Lie}_\sigma:w(e)\in\mathcal
O_Q$ for every $w\in H_n(Q)\}$. Then $H^\circ(Q)=\Lambda(Q)\cap\mathrm{Lie}_\sigma$. If $\mathfrak d$ is an ideal of $O$
such that the values of $\sigma$ are regular at $Q$ and $\sigma(d)$ is a unit at $Q$ for some $d\in\mathfrak d$, then the
$\mathcal O_Q$-module generated by the $e\in\mathrm{Lie}_\sigma$ with $\sigma(d')e\in H^\circ(Q)$ for every
$d'\in\mathfrak d$ is also $\Lambda(Q)\cap\mathrm{Lie}_\sigma$.
\end{lemma}
\leanref{Development/AutomorphicSpectrum/Proofs/GeometryProofs/HodgeTateEigenLatticeEqStatement.lean}{GlobalTheory.T.hodge_tate_eigen_lattice_eq}
\begin{proof}
As $\Lambda(Q)=\sum_k\mathcal O_Qb_k$ for a basis $b$ of $\mathrm{Lie}$, $\Lambda^\vee(Q)=\sum_k\mathcal O_Qb^\vee_k$ for the
dual basis, and $e\in\mathrm{Lie}$ lies in $\Lambda(Q)$ if and only if $w(e)\in\mathcal O_Q$ for every $w\in\Lambda^\vee(Q)$,
as $b^\vee_k(e)$ are the coordinates of $e$. As $H_n(Q)=\Lambda^\vee(Q)$, $H^\circ(Q)=\Lambda(Q)\cap\mathrm{Lie}_\sigma$.
For the second claim, let $H'$ be the module generated by the $e\in\mathrm{Lie}_\sigma$ with $\sigma(d')e\in
\Lambda(Q)\cap\mathrm{Lie}_\sigma$ for every $d'\in\mathfrak d$. If $e\in\Lambda(Q)\cap\mathrm{Lie}_\sigma$, then so is
each $\sigma(d')e$, as $\sigma(d')\in\mathcal O_Q$; so $e\in H'$. Conversely, if $e$ is one of the generators of $H'$, then
$e=\sigma(d)^{-1}\cdot\sigma(d)e$ with $\sigma(d)^{-1}\in\mathcal O_Q$, so $e\in\Lambda(Q)\cap\mathrm{Lie}_\sigma$; hence
$H'=\Lambda(Q)\cap\mathrm{Lie}_\sigma$.
\end{proof}
% Lean: hodge_tate_eigen_lattice_eq = LieModel.htDual_htW_eq_eigLat_of_htOm_eq_omLat: Lat.isVBOn_dual_and_bidual (the
% bidual characterization at q), mem_htDual_iff, eigLat_le_htDual, and the unit sigma d for htW.

\begin{lemma}\label{cv-origin-change-projective-bundles-glue}
Let $T$ be a locally noetherian scheme, $\Omega$ a field of characteristic zero and $\xi:\operatorname{Spec}\Omega\to T$ a
morphism with image $\eta$ such that the local ring at $\eta$ maps bijectively onto $\Omega$ and those at its
specializations injectively. A point $Q$ that specializes $\eta$ \emph{lies over} a prime $p$ if $p$ is not a unit at $Q$.
Let $V$ be an $\Omega$-vector space, $W\subset V$ a subspace, $r\ge0$, $S$ a finite set of primes, and $L_0$ and $L_p$, for
$p\in S$, lattice families in $V$ such that $L_0$ is a vector bundle of rank $r$ with generic fibre $W$ on the set of the
points over no prime of $S$, each $L_p$ is one on the set of the points over no prime other than $p$, and
$L_p(Q)=L_0(Q)$ at every point $Q$ over no prime. Then there is a lattice family $L$ in $V$ that is a vector bundle of rank
$r$ with generic fibre $W$, with $L(Q)=L_0(Q)$ at the points over no prime of $S$ and $L(Q)=L_p(Q)$ at the points over $p$,
for $p\in S$.
\end{lemma}
\leanref{Development/AutomorphicSpectrum/Proofs/GeometryProofs/LatticeGlueOverPrimesStatement.lean}{GlobalTheory.T.lattice_glue_over_primes}
\begin{proof}
In this proof all points are points that specialize $\eta$.

\emph{Primes.} A point $Q$ lies over at most one prime: for primes $p\neq p'$ there are integers $a,a'$ with $ap+a'p'=1$,
and in the local ring $\mathcal O_Q$ one of $ap$ and $a'p'$ is a unit. So the lattice family $L$ equal to $L_p$ at the
points over $p\in S$ and to $L_0$ at the other points is well defined, and it has the stated values. If a prime is a unit at
$Q$, its inverse is regular at the points of an open neighbourhood of $Q$ (Lemma~\ref{cv-origin-segre-regular}(1)); so for
every finite set of primes none of which $Q$ lies over, $Q$ has an open neighbourhood whose points lie over none of them.

\emph{Points over no prime of $S$.} Let $Q$ be such a point. Then $L_0$ has a basis $b$ of $W$ on an open neighbourhood $U$
of $Q$: $L_0(Q')=\sum_k\mathcal O_{Q'}b_k$ at the points $Q'$ of $U$ over no prime of $S$. On the intersection of $U$ with
an open neighbourhood of $Q$ whose points lie over no prime of $S$, $L=L_0=\sum_k\mathcal O b_k$.

\emph{Integrality.} Let $W_0=\operatorname{Spec}A$ be an affine open subset of $T$ and $a\in\Omega$ regular at every point
of $W_0$ over no prime. Then $ma$ is regular at every point of $W_0$ for some integer $m\neq0$. We may assume that
$\eta\in W_0$, as otherwise $W_0$ has no points. Let $\varphi:A\to\Omega$ be the map through the local ring at $\eta$, and
$D$ the ideal of the $s\in A$ with $\varphi(s)a\in\varphi(A)$. If $D$ contains an integer $m\neq0$, then $ma\in\varphi(A)$
is regular on $W_0$. Otherwise $D$ does not meet the set of nonzero integers, so it lies in a prime ideal $\mathfrak q$ of
$A$ that does not meet it. The prime ideal $\mathfrak p_\eta$ of $\eta$ lies in $D$: for $s\in\mathfrak p_\eta$ the image of
$s$ in the local ring at $\eta$ is not a unit, so $\varphi(s)=0$, as that ring maps bijectively onto the field $\Omega$, and
$\varphi(s)a=\varphi(0)$. So the point $\mathfrak q$ specializes $\eta$, and it lies over no prime, as the nonzero integers
are units in $A_{\mathfrak q}$. Hence $a$ is regular at $\mathfrak q$: $a=\varphi(r')/\varphi(t)$ with $r',t\in A$ and
$t\notin\mathfrak q$. Then $\varphi(t)a=\varphi(r')$, so $t\in D\subset\mathfrak q$, a contradiction.

\emph{Comparing bases.} Let $W_0$ be an affine open subset of $T$ and $b$ and $c$ linearly independent families of $r$
vectors spanning the same subspace, with $\sum_k\mathcal O_Qb_k=\sum_k\mathcal O_Qc_k$ at every point $Q$ of $W_0$ over no
prime. Then there is an integer $m\neq0$ such that this holds at every point of $W_0$ over no prime dividing $m$. Indeed,
write $b_i=\sum_jA_{ij}c_j$ and $c_i=\sum_jB_{ij}b_j$ with $A_{ij},B_{ij}\in\Omega$. At the points over no prime these
coefficients are regular, as $b$ and $c$ are linearly independent; by the previous step there are nonzero integers
$m_{ij}$ and $m'_{ij}$ with $m_{ij}A_{ij}$ and $m'_{ij}B_{ij}$ regular on $W_0$. Let $m$ be the product of all of them. At a
point over no prime dividing $m$, $m$ is a unit, so the $A_{ij}$ and $B_{ij}$ are regular, and the two modules agree.

\emph{Points over a prime of $S$.} Let $Q$ lie over $p\in S$. It lies over no other prime, so $L_p$ has a basis $b$ of $W$
on an open neighbourhood $U_0$ of $Q$, at the points over no prime other than $p$. Let $U_2$ be an affine open
neighbourhood of $Q$ inside $U_0$ whose points lie over no prime of $S$ other than $p$. The set $C$ of the points of $U_2$
that do not lie over $p$ is quasi-compact, as $U_2$ is a noetherian topological space; its points lie over no prime of $S$.
Let $x\in C$. Then $L_0$ has a basis $c$ of $W$ on an open neighbourhood of $x$, at the points over no prime of $S$; let
$W_x$ be an affine open neighbourhood of $x$ inside it and inside $U_2$. At the points of $W_x$ over no prime,
$\sum_k\mathcal O b_k=L_p=L_0=\sum_k\mathcal O c_k$. By the previous step there is $m_x\neq0$ with
$L_0=\sum_k\mathcal Oc_k=\sum_k\mathcal Ob_k$ at the points of $W_x$ that do not lie over $p$ and lie over no prime dividing
$m_x$. Finitely many of the $W_x$ cover $C$; let $m$ be the product of their $m_x$, and $U_3$ an open neighbourhood of $Q$
whose points lie over no prime factor of $m$ other than $p$. Let $Q'$ be a point of $U_2\cap U_3$. If $Q'$ lies over $p$,
then $L(Q')=L_p(Q')=\sum_k\mathcal O_{Q'}b_k$. Otherwise $Q'\in C$ lies in one of the chosen $W_x$ and over no prime
dividing $m_x$, so $L(Q')=L_0(Q')=\sum_k\mathcal O_{Q'}b_k$. So $L$ is a vector bundle of rank $r$ with generic fibre $W$
near $Q$.
\end{proof}
% Lean: lattice_glue_over_primes = lat_glue: LatGlue.not_inFibre_of_inFibre, exists_open_avoid
% (T.regAt_unitAt_open), exists_nat_mul_mem (the ideal D, Ideal.exists_le_prime_disjoint, stalk_eq_germ_div_of_affine),
% exists_matrix and exists_nat_span_eq, and local_step (exists_isAffineOpen_mem_and_subset, noetherianSpace_of_isAffineOpen,
% elim_nhds_subcover, the product of the m_x and its primeFactors).

\begin{lemma}\label{cv-origin-change-projective-bundles}
In the setting of Definition~\ref{cv-origin-hodge-component}, let $T$ be a locally noetherian scheme, $\Omega$ a field of
characteristic zero and $\xi:\operatorname{Spec}\Omega\to T$ a morphism with image $\eta$ such that the local ring at
$\eta$ maps bijectively onto $\Omega$, and the local rings at the specializations of $\eta$ map injectively into $\Omega$
and are integrally closed in $\Omega$. Let $f:T\to\mathcal X$ be a morphism with $f\circ\xi=j\circ y_0$ for a point $y_0$
of $X$, and $F_c\to\Omega$ an embedding through which $f\circ\xi$ lies over $\mathcal O_c$. Let $G$ be an
$\mathcal O_E$-semi-abelian scheme over $T$ with a morphism $G\to\mathcal A$ over $f$, $\mathrm{Lie}$ its Lie algebra at
$\xi$ with the action of $\mathcal O_E$, and, for a label $s$, $\mathrm{Lie}_s$ the subspace on which each
$b\in\mathcal O_E$ acts by $s(b)$, through $F_c\to\Omega$. Let $S$ be a finite set of primes that contains every prime
dividing $d_E$, and $R_p\ge0$ integers, for $p\in S$, such that $G$ has full $p^{R_p}$-torsion at $\xi$ and $\Omega$
contains a primitive $p^{R_p}$-th root of unity. Suppose that for every $p\in S$ and every label $s$ the restrictions to
$\mathrm{Lie}_s$ of the forms in $H_{p^{R_p}}$ (Definition~\ref{cv-origin-change-projective-periods}) form a vector bundle
of rank $r_s$ with generic fibre $\mathrm{Lie}_s^\vee$ (Definition~\ref{cv-origin-reference-lattices}). Then for every
label $s$ the lattice family $\mathcal V_s$ of Definition~\ref{cv-origin-hodge-component}(1), for $G$, $S$ and the
$R_p$, is a vector bundle of rank $r_s$ with generic fibre $\mathrm{Lie}_s$.
\end{lemma}
\leanref{Development/AutomorphicSpectrum/Proofs/GeometryProofs/EigenLatticeDedekindTwistStatement.lean}{GlobalTheory.T.eigen_lattice_dedekind_twist}
\begin{proof}
Let $\sigma:F_c\to\Omega$ be the given embedding. For a label $s$ put $\sigma_s=\sigma\circ s:\mathcal O_E\to\Omega$, so
that $\mathrm{Lie}_s$ is the subspace on which each $b\in\mathcal O_E$ acts by $\sigma_s(b)$, and for $p\in S$ put
$n_p=p^{R_p}$. A point $Q$ that specializes $\eta$ \emph{lies over} a prime $p$ if $p$ is not a unit at $Q$; it lies over at
most one prime, as for primes $p\neq p'$ some integral combination of $p$ and $p'$ is $1$ and $\mathcal O_Q$ is local. For
$p\in S$ and a label $s$ let $H^p_s$ be the lattice family $Q\mapsto\{e\in\mathrm{Lie}_s:w(e)\in\mathcal O_Q$ for every
$w\in H_{n_p}(Q)\}$, and $\mathcal V^p_s$ that of Definition~\ref{cv-origin-hodge-component}(1): $\mathcal V^p_{\bar\tau}=
H^p_{\bar\tau}$, and $\mathcal V^p_\tau(Q)$ is the $\mathcal O_Q$-module generated by the $e\in\mathrm{Lie}_\tau$ with
$\sigma_\tau(\delta)e\in H^p_\tau(Q)$ for every $\delta$ in the different $\mathfrak d$ of $E$. Fix a label $s$, and let
$T_0$ be the set of the points that specialize $\eta$ and lie over no prime of $S$. Then $\mathcal V_s$ is
$\Lambda\cap\mathrm{Lie}_s$ at the points of $T_0$ and $\mathcal V^p_s$ at the points over $p\in S$.

\emph{Away from $S$.} As $G$ is smooth over $T$, $\Lambda$ is a vector bundle of rank $\dim\mathrm{Lie}$ with generic fibre
$\mathrm{Lie}$ by Lemma~\ref{cv-origin-reference-lie-bundle}, and $\dim\mathrm{Lie}_s=r_s$ by
Lemma~\ref{cv-origin-reference-rank}. The $2|\Psi|=[E:\mathbf Q]$ labels are distinct embeddings of $\mathcal O_E$ into
$F_c$, so the $\sigma_s$ are distinct. At every point $Q$ that specializes $\eta$, $\Lambda(Q)$ is stable under
$\mathcal O_E$, as in the proof of Lemma~\ref{cv-origin-reference-eigen}, and the values of the $\sigma_s$ are regular at
$Q$, as they are integral over $\mathbf Z$ and $\mathcal O_Q$ is integrally closed in $\Omega$. At the points of $T_0$ the
discriminant $d_E$ is a unit, as every prime dividing it lies in $S$. So Lemma~\ref{cv-origin-reference-eigen-split},
applied with $F=E$, $L=\Lambda$ and the set $T_0$, shows that $\Lambda\cap\mathrm{Lie}_s$ is a vector bundle of rank $r_s$
with generic fibre $\mathrm{Lie}_s$ on $T_0$.

\emph{At a prime $p\in S$.} By hypothesis the restrictions to $\mathrm{Lie}_{s'}$ of the forms in $H_{n_p}$ form a vector
bundle of rank $r_{s'}$ with generic fibre $\mathrm{Lie}_{s'}^\vee$, for every label $s'$. So, by
Lemma~\ref{cv-origin-change-projective-bundles-predual}, $H^p_{s'}$ is a vector bundle of rank $r_{s'}$ with generic fibre
$\mathrm{Lie}_{s'}$. If $s=\bar\tau$, then $\mathcal V^p_s=H^p_s$. If $s=\tau$, then $\sigma_\tau$ is injective on the
Dedekind domain $\mathcal O_E$, as $\tau$ and $\sigma$ are, its values are regular at every point that specializes $\eta$,
and $\mathfrak d\neq0$; so, by Lemma~\ref{cv-origin-change-projective-bundles-saturate}, $\mathcal V^p_\tau$ is a vector
bundle of rank $r_\tau$ with generic fibre $\mathrm{Lie}_\tau$.

\emph{At the points over no prime.} Let $p\in S$, and $Q$ a point that specializes $\eta$ and lies over no prime. By
Lemma~\ref{cv-origin-change-projective-bundles-weil}, $G$ has a generic Weil pairing of order $n_p$ at $\xi$, so by
Lemma~\ref{cv-origin-change-projective-coherent}, $n_p\Lambda^\vee(Q)\subset H_{n_p}(Q)\subset\Lambda^\vee(Q)$. As $n_p$
is a unit at $Q$, $H_{n_p}(Q)=\Lambda^\vee(Q)$. So Lemma~\ref{cv-origin-change-projective-bundles-generic}, applied with
$\rho$ the action of $\mathcal O_E$ and $\sigma_s$, gives $H^p_s(Q)=\Lambda(Q)\cap\mathrm{Lie}_s$. For $s=\tau$, the
integer $|d_E|$ lies in $\mathfrak d$, as it is the norm of $\mathfrak d$ \cite[Chapter~III, \S2]{Neu99} and an ideal of
$\mathcal O_E$ contains its norm, and $\sigma_\tau(|d_E|)=|d_E|$ is a unit at $Q$; so the same lemma gives
$\mathcal V^p_\tau(Q)=\Lambda(Q)\cap\mathrm{Lie}_\tau$. Hence $\mathcal V^p_s(Q)=(\Lambda\cap\mathrm{Lie}_s)(Q)$.

\emph{Gluing.} Lemma~\ref{cv-origin-change-projective-bundles-glue}, applied with $W=\mathrm{Lie}_s$,
$L_0=\Lambda\cap\mathrm{Lie}_s$ and $L_p=\mathcal V^p_s$, gives a lattice family $L$ that is a vector bundle of rank $r_s$
with generic fibre $\mathrm{Lie}_s$, equal to $\Lambda\cap\mathrm{Lie}_s$ on $T_0$ and to $\mathcal V^p_s$ at the points
over $p\in S$. So $L=\mathcal V_s$.
\end{proof}
% Lean: GlobalTheory.T.eigen_lattice_dedekind_twist = hodgeLattices_isVB: T.lie_omega_lattice_vector_bundle (hLat),
% T.eigen_lattice_constant_rank (hrk), T.label_embedding_injective, HcVb.integers_regular; hA from
% T.lattice_eigen_vector_bundle on the points over no prime of S (HcVb.lat_stable, HcVb.unitAt_discr); hDual from
% T.predual_vector_bundle with mem_htDual_iff; hB with T.ideal_saturated_vector_bundle for the different (inl tau); hC from
% HcVb.exists_genWeil, T.hodge_tate_omega_coherent and T.hodge_tate_eigen_lattice_eq (natAbs_discr_mem_different via
% Ideal.absNorm_mem and absNorm_differentIdeal); T.lattice_glue_over_primes; Lat.patch_of_unique with inFibre_unique.

\begin{lemma}\label{cv-origin-change-projective-frames}
In the setting of Definition~\ref{cv-origin-hodge-component}, let $T$, $\Omega$ and $\xi$ be as in
Lemma~\ref{cv-origin-reference-lie-bundle}, with $\Omega$ of characteristic zero and the local rings at the
specializations of the image of $\xi$ integrally closed in $\Omega$, let $F_c\to\Omega$ be an embedding, and let $G$ be an
$\mathcal O_E$-semi-abelian scheme over $T$ with $\mathrm{Lie}$, $\mathrm{Lie}_s$ and $\Lambda$ as in
Lemma~\ref{cv-origin-reference-eigen}. Let $P$ be a point over no prime that specializes the image of $\xi$. Then for every
label $s$ the restrictions to $\mathrm{Lie}_s$ of the forms in $\Lambda^\vee(P)$ form an $\mathcal O_P$-module generated by
linearly independent vectors of $\mathrm{Lie}_s^\vee$.
\end{lemma}
\leanref{Development/AutomorphicSpectrum/Proofs/GeometryProofs/HodgeCptProjectiveModel.lean}{GlobalTheory.CompanionTheory.Ar.TopHodge.frame_om}
\begin{proof}
\emph{The eigenspaces at $P$.} The $[E:\mathbf Q]$ labels are distinct embeddings of $\mathcal O_E$ into $F_c$, so their
composites $\sigma_s$ with $F_c\to\Omega$ are $[E:\mathbf Q]$ distinct homomorphisms $\mathcal O_E\to\Omega$, and the
subspaces $\mathrm{Lie}_s$ are independent. By Lemma~\ref{cv-origin-reference-lie-bundle}, $\Lambda$ is a vector bundle of
rank $\dim\mathrm{Lie}$ with generic fibre $\mathrm{Lie}$; in particular $\Lambda(P)$ spans $\mathrm{Lie}$. As in the proof
of Lemma~\ref{cv-origin-reference-eigen}, $\Lambda(P)$ is stable under $\mathcal O_E$, the values $\sigma_s(b)$ are regular
at $P$, and the discriminant of $E$ is a unit at $P$. So Lemma~\ref{cv-origin-reference-eigen-split}, applied with $F=E$,
$L=\Lambda$ and the set $\{P\}$, shows that $\Lambda(P)=\sum_t(\Lambda(P)\cap\mathrm{Lie}_t)$, and that
$\Lambda(P)\cap\mathrm{Lie}_s=\sum_i\mathcal O_Pb_i$ for a basis $b_1,\dots,b_r$ of $\mathrm{Lie}_s$.

\emph{The dual basis.} As $\Lambda(P)$ spans $\mathrm{Lie}$ and lies in $\sum_t\mathrm{Lie}_t$, $\mathrm{Lie}$ is the direct
sum of $\mathrm{Lie}_s$ and the $\mathrm{Lie}_t$ with $t\neq s$; let $\pi$ be the projection onto $\mathrm{Lie}_s$ along
the others. It maps $\Lambda(P)$ into $\Lambda(P)\cap\mathrm{Lie}_s$, as it fixes $\Lambda(P)\cap\mathrm{Lie}_s$ and kills
$\Lambda(P)\cap\mathrm{Lie}_t$ for $t\neq s$. Let $b^\vee_1,\dots,b^\vee_r$ be the dual basis of $\mathrm{Lie}_s^\vee$. For
$w\in\Lambda^\vee(P)$ the restriction of $w$ to $\mathrm{Lie}_s$ is $\sum_iw(b_i)b^\vee_i$, with $w(b_i)\in\mathcal O_P$ as
$b_i\in\Lambda(P)$. Conversely, $b^\vee_i\circ\pi$ lies in $\Lambda^\vee(P)$, since for $e\in\Lambda(P)$,
$\pi(e)\in\sum_l\mathcal O_Pb_l$ and $b^\vee_i(\pi(e))\in\mathcal O_P$; its restriction to $\mathrm{Lie}_s$ is
$b^\vee_i$. So the restrictions of the forms in $\Lambda^\vee(P)$ form $\sum_i\mathcal O_Pb^\vee_i$.
\end{proof}
% Lean: TopHodge.frame_om: T.label_embedding_injective, T.lie_omega_lattice_vector_bundle, T.lattice_eigen_vector_bundle on
% {q} with lat_stable, the integrality of the labels (hN) and the unit discriminant (hq), and frame_dualRestrict
% (eigSp_indep, the complement hc, the projection linearProjOfIsCompl with hpi, and the dual basis B.coord).

\begin{lemma}\label{cv-origin-change-projective}
In the setting of Definition~\ref{cv-origin-hodge-component}, let $S$ be a finite set of primes that contains
every prime dividing $d_E$, and let $R_p\ge4$ for $p\in S$. Let $\Omega$ be a field of characteristic zero,
$y\in\Omega^{\sqcup_kI_k}$, and $X_1$ a scheme over $\mathbf Z$ with a morphism $\xi_1:\operatorname{Spec}\Omega\to X_1$
of whose image every point is a specialization, covered by the affine open subschemes of the charts of $y$,
whose local rings are integrally closed in $\Omega$. Let $f_1:X_1\to\mathcal X$ be a morphism that maps the
image of $\xi_1$ to the generic point of an irreducible component, with $f_1\circ\xi_1=j\circ y_0$ for a point
$y_0$ of $X$. Suppose that for every $p\in S$, $\Omega$ contains a primitive $p^{R_p}$-th root of unity and
every point of order dividing $p^{R_p}$ of $\mathcal A$ at $f_1\circ\xi_1$ with values in an extension of
$\Omega$ is defined over $\Omega$. Then there are a point $y'$ that contains $y$ as some of its blocks, a Hodge
component $Z$ of $\mathcal X$ at $S$ with coordinates $y'$ and the integers $R_p$, and a morphism
$g:Z\to X_1$ that maps the generic point of $Z$ to the image of $\xi_1$, with $f=f_1\circ g$, such that:
\begin{enumerate}
\item for every $p\in S$, the module spanned by the Hodge--Tate periods of the points of order dividing
  $p^{R_p}$, and its images in the duals of the $\mathrm{Lie}_s$, are vector bundles of ranks
  $\dim\mathrm{Lie}$ and $r_s$ on the fibre of $Z$ over $p$;
\item $g$ is an isomorphism over each point $P$ of $X_1$ over no prime: $P$ has one preimage $P'$, and the
  local rings of $Z$ at $P'$ and of $X_1$ at $P$ are equal as subrings of $\Omega$.
\end{enumerate}
\end{lemma}
\leanref{Development/AutomorphicSpectrum/Proofs/GeometryProofs/HodgeCptProjectiveModelStatement.lean}{GlobalTheory.T.hodge_cpt_projective_model}
\leanref{Development/AutomorphicSpectrum/Definitions/GlobalGeometry/IntegralClosureHeightEqOneOfLiesOver.lean}{GlobalTheory.CompanionTheory.Ar.HT.HTFree}
\begin{proof}
Let $\eta_1$ be the image of $\xi_1$, and for $p\in S$ put $n_p=p^{R_p}$.

\emph{The Lie algebra.} The point $y_0$ of $X$, which lies over $F_c$, gives an embedding $F_c\to\Omega$, through which
$f_1\circ\xi_1=j\circ y_0$ lies over $\mathcal O_c$. Lemma~\ref{cv-origin-reference-pullback}(1) gives
$\mathcal A_1=f_1^*\mathcal A$ with a morphism $\mathcal A_1\to\mathcal A$ over $f_1$. Let $\mathrm{Lie}$ be the Lie
algebra of $\mathcal A_1$ at $\xi_1$, with its action of $\mathcal O_E$ (Lemma~\ref{cv-origin-reference-lie}), and for a
label $s$ let $\mathrm{Lie}_s$ be the subspace on which each $b\in\mathcal O_E$ acts by $s(b)$, through $F_c\to\Omega$. By
Lemma~\ref{cv-origin-reference-rank}, $\mathrm{Lie}$ is finite-dimensional and $\dim\mathrm{Lie}_s=r_s$. The scheme $X_1$
has the Segre products $x$ of $y$ as coordinates (Definition~\ref{cv-origin-segre-coordinates}); it is noetherian, being
covered by the finitely many affine open subschemes of the charts, whose rings $\mathbf Z[y_{k,i}/y_{k,\kappa_k}]$ are
finitely generated over $\mathbf Z$. Let $\Lambda$ and $\Lambda^\vee$ be the lattice families of the Lie algebras of
$\mathcal A_1$ at the points of $X_1$ and of their duals
(Definition~\ref{cv-origin-change-projective-periods}(3)). By Lemma~\ref{cv-origin-reference-lie-bundle}, $\Lambda^\vee$
is a vector bundle of rank $\dim\mathrm{Lie}$ with generic fibre $\mathrm{Lie}^\vee$.

\emph{Hodge--Tate periods on $X_1$.} Let $p\in S$. By hypothesis $\mathcal A$ has full $n_p$-torsion at
$j\circ y_0=f_1\circ\xi_1$ and $\Omega$ contains a primitive $n_p$-th root of unity, so by
Lemma~\ref{cv-origin-change-projective-weil} $\mathcal A_1$ has full $n_p$-torsion and a generic Weil pairing of order
$n_p$ at $\xi_1$. Let $H_p$ be the lattice family $P\mapsto H_{n_p}(P)$ of the modules spanned by the Hodge--Tate periods
of $\mathcal A_1$ of order $n_p$. As the local rings of $X_1$ are integrally closed in $\Omega$,
Lemma~\ref{cv-origin-change-projective-coherent} shows that $H_p$ is locally generated by finitely many forms and that
$n_p\Lambda^\vee(P)\subset H_p(P)\subset\Lambda^\vee(P)$ at every point $P$. As $n_p$ is invertible in $\Omega$ and
$\Lambda^\vee(P)$ spans $\mathrm{Lie}^\vee$, so does $H_p(P)$. If $P$ lies over no prime, then $n_p^{-1}$ is regular at
$P$, and $H_p(P)=\Lambda^\vee(P)$. For a label $s$ let $\pi_s:\mathrm{Lie}^\vee\to\mathrm{Lie}_s^\vee$ be the restriction,
which is surjective, and $H_{p,s}$ the lattice family $P\mapsto\pi_s(H_p(P))$. It is locally generated by the images of
local generators of $H_p$, and $H_{p,s}(P)$ spans $\mathrm{Lie}_s^\vee$.

\emph{A refinement.} Every point of $X_1$ lies in a chart of $x$. Lemma~\ref{cv-origin-reference-refinement}, applied to
$X_1$, $x$, the single family $w=x$, and the lattice families $H_p$ in $\mathrm{Lie}^\vee$ and $H_{p,s}$ in
$\mathrm{Lie}_s^\vee$, for $p\in S$ and the labels $s$, with these spaces as the subspaces $W_b$, gives a finite family
$x'$, a scheme $X'$ with a morphism $\xi':\operatorname{Spec}\Omega\to X'$ of whose image every point is a
specialization, with the coordinates $x'$ and with local rings integrally closed in $\Omega$, and a morphism
$g:X'\to X_1$ with $g\circ\xi'=\xi_1$, such that every point of $X'$ lies in a chart of $x$, each $g^*H_p$ is a vector
bundle of rank $\dim\mathrm{Lie}$ with generic fibre $\mathrm{Lie}^\vee$, each $g^*H_{p,s}$ is a vector bundle of rank
$r_s$ with generic fibre $\mathrm{Lie}_s^\vee$, and $g$ is a normalization over every point $P$ of $X_1$ at which every
$H_p(P)$ and every $H_{p,s}(P)$ is generated by linearly independent vectors.

\emph{The semi-abelian scheme on $X'$.} Lemma~\ref{cv-origin-reference-pullback} gives $\mathcal A'=g^*\mathcal A_1$, with
a morphism to $\mathcal A_1$ over $g$ whose composite with $\mathcal A_1\to\mathcal A$ is a morphism
$\mathcal A'\to\mathcal A$ over $f=f_1\circ g$. By Lemma~\ref{cv-origin-reference-lie-same}, composition with
$\mathcal A'\to\mathcal A_1$ identifies the Lie algebra of $\mathcal A'$ at $\xi'$ with $\mathrm{Lie}$, compatibly with the
actions of $\mathcal O_E$; we identify the two. As $f\circ\xi'=f_1\circ\xi_1=j\circ y_0$, the point $f\circ\xi'$ lies over
$\mathcal O_c$ through the same embedding $F_c\to\Omega$, and Lemma~\ref{cv-origin-change-projective-weil}, applied to
$X'$, $\xi'$, $f$ and $\mathcal A'\to\mathcal A$, shows that $\mathcal A'$ has full $n_p$-torsion and a generic Weil
pairing of order $n_p$ at $\xi'$, for $p\in S$. The scheme $X'$ is noetherian, as $X_1$ is. Let $H'_p$ be the lattice
family of the modules spanned by the Hodge--Tate periods of $\mathcal A'$ of order $n_p$.
Lemma~\ref{cv-origin-change-projective-pullback}, applied to $g$ with $\iota$, $j$ and $j^\vee$ the identities, gives
$H'_p=g^*H_p$. Pulling back commutes with $\pi_s$: $(g^*H_{p,s})(P')$ is generated over $\mathcal O_{P'}$ by
$\pi_s(H_p(g(P')))$, so it is the image under the linear map $\pi_s$ of the module $(g^*H_p)(P')$ that $H_p(g(P'))$
generates. So $H'_p$ is a vector bundle of rank $\dim\mathrm{Lie}$ with generic fibre $\mathrm{Lie}^\vee$, and
$P'\mapsto\pi_s(H'_p(P'))$ is the vector bundle $g^*H_{p,s}$, of rank $r_s$ with generic fibre $\mathrm{Lie}_s^\vee$.

\emph{The modules $\mathcal V_s$.} Lemma~\ref{cv-origin-change-projective-bundles}, applied to $X'$, $\xi'$, $f$,
$\mathcal A'$, $S$ and the $R_p$, shows that for every label $s$ the lattice family $\mathcal V_s$ of
Definition~\ref{cv-origin-hodge-component}(1) for $\mathcal A'$ is a vector bundle of rank $r_s$ with generic fibre
$\mathrm{Lie}_s$.

\emph{The Hodge component.} A point in the chart $c$ of $x$ lies in the chart $c_k$ of each block $k$ of $y$, as
$y_{k,i}/y_{k,c_k}$ is the Segre ratio for the chart $c$ with $c_k$ replaced by $i$. So
Lemma~\ref{cv-origin-reference-segre}, applied to $X'$, with $x'$ as a point with one block, and to the blocks of $y$,
shows that $X'$ is covered by the affine open subschemes of the charts of the point $y'$ with the blocks of $y$ and the
block $x'$. The scheme $X'$ is irreducible, as every point specializes the image of $\xi'$, and reduced, as its local
rings embed in $\Omega$; these local rings are integrally closed in their field of fractions $\Omega$, so $X'$ is a normal
integral scheme with function field $\Omega$. The morphism $f$ maps the image of $\xi'$ to $f_1(\eta_1)$, the generic
point of an irreducible component of $\mathcal X$; with $\Omega_0$ its residue field, $f\circ\xi'$ is given by an
embedding $\Omega_0\to\Omega$ \cite[Tag 01J9]{Sta}. So $Z=X'$, with the coordinates $y'$, the morphism $f$, the
semi-abelian scheme $\mathcal A_Z=\mathcal A'$ and the integers $R_p$, is a Hodge component of $\mathcal X$ at $S$: for
$p\in S$, $R_p\ge4$ and $\mathcal A'$ has full $n_p$-torsion at $\xi'$, and the $\mathcal V_s$ are vector bundles of rank
$r_s$. The morphism $g$ maps the generic point of $Z$ to $\eta_1$, as $g\circ\xi'=\xi_1$, and $f=f_1\circ g$.

(1) For $p\in S$, $H'_p$ and the lattice families $P'\mapsto\pi_s(H'_p(P'))$ are vector bundles of ranks
$\dim\mathrm{Lie}$ and $r_s$ on $Z$, so also on the fibre of $Z$ over $p$.

(2) Let $P$ be a point of $X_1$ over no prime, and $p\in S$. Then $H_p(P)=\Lambda^\vee(P)$, which is generated by a basis
of $\mathrm{Lie}^\vee$, as $\Lambda^\vee$ is a vector bundle; and $H_{p,s}(P)=\pi_s(\Lambda^\vee(P))$, which is generated
by linearly independent vectors by Lemma~\ref{cv-origin-change-projective-frames}. So $g$ is a normalization over $P$. As
the local ring $\mathcal O_{X_1,P}$ is integrally closed, Lemma~\ref{cv-origin-reference-normalization} shows that every
point $P'$ of $Z$ with $g(P')=P$ has $\mathcal O_{P'}=\mathcal O_{X_1,P}$, and that there is at most one such point. There
is one: $\mathcal O_{X_1,P}$ is its own integral closure in $\Omega$, and the normalization over $P$ has a point over $P$
for its maximal ideal (Definition~\ref{cv-origin-reference-lattices}(5)).
\end{proof}
% Lean: TopHodge.exists_hodgeCpt_htFree, under Out.WeilSt, step by step: exists_cst; sabPull; lieModel_exists and
% lieModel_exists_unique_act; T.eigen_lattice_constant_rank; T.regAt_noetherian_model_noetherian (X1 noetherian);
% T.lie_omega_lattice_vector_bundle (hom1); tors_weil (Lemma cv-origin-change-projective-weil); T.hodge_tate_omega_coherent
% (hcoh); hspan and heq; common2 = T.flag_model_common_refinement with Unit for the single family x and the families
% indexed by S + S x labels (T.coherent_lattice_closure for the images, span_map_top); sabPull, sabComp and
% LiePull.lieModel_pullback for N' with hrho' (Lemma cv-origin-reference-lie-same); tors_weil again (htw');
% T.hodge_tate_weil_pullback with i, j and jd the identities (hpull); pull_map and Subspace.dual_finrank_eq (hf1, hf2);
% T.eigen_lattice_dedekind_twist (hvb); T.segre_refine_line_bundles with no line bundles (Lemma cv-origin-reference-segre);
% T.generic_point_normal_coords and exists_factor for Omega_0; T.normal_cpt_iff_normal_scheme; isVBOn_mono for HTFree; and
% T.normal_iso_of_integrally_closed with hnorm, hom1, heq and frame_om (Lemma cv-origin-change-projective-frames) for IsoAt.
% Lemma cv-origin-change-projective-local-inputs-hopf-order, formerly the open part, is written out.

\begin{lemma}\label{cv-origin-change-comparison}
In the setting of Definition~\ref{cv-origin-hodge-component}, let $Z_1$ be a Hodge component of $\mathcal X$
at $S_1$ with coordinates $y_1$ with values in $\Omega_1$, and $Z$ a Hodge component at $S$ with coordinates
$y$ with values in $\Omega$, with Lie algebras $\mathrm{Lie}_1$ and $\mathrm{Lie}$ at their generic points.
Let $\iota:\Omega_1\to\Omega$ be a field embedding and $h:Z\to Z_1$ a morphism that maps the generic point
of $Z$ to that of $Z_1$ compatibly with $\iota$, with $f=f_1\circ h$.
\begin{enumerate}
\item There is an additive map $j_A:\mathrm{Lie}_1\to\mathrm{Lie}$ with $j_A(ae)=\iota(a)j_A(e)$, which carries
  linearly independent families to linearly independent families, commutes with the actions of
  $\mathcal O_E$ and is compatible with the maps to the Lie algebra of $\mathcal A$, such that at every point
  $P$ of $Z$ the Lie algebra of $\mathcal A_Z$ at $P$ is the $\mathcal O_{Z,P}$-module generated by $j_A$ of
  that of $\mathcal A_{Z_1}$ at $h(P)$.
\item Suppose moreover that $\Omega$ is finite over $\iota(\Omega_1)$, that the embedding of $F_c$ in
  $\Omega$ is $\iota$ composed with that in $\Omega_1$, and that $Z_1$ has a flag of degree $\deg$ with bases
  $b_{s,i}$. Then each family $j_A(b_s)$ spans $\mathrm{Lie}_s$, and the families $j_A(b_s)$ satisfy the
  genericity condition of Definition~\ref{cv-origin-hodge-component}(2) for $Z$ with degree
  $[\Omega:\iota(\Omega_1)]\deg$.
\item Under the hypotheses of (2), at every point $P$ of $Z$ over no prime, $\mathcal V_s$ has a basis at $P$
  adapted to the flag of $j_A(b_s)$.
\end{enumerate}
\end{lemma}
\leanref{Development/AutomorphicSpectrum/Proofs/GeometryProofs/HodgeCompComparisonLeftStatement.lean}{GlobalTheory.T.hodge_comp_comparison_left}
\leanref{Development/AutomorphicSpectrum/Proofs/GeometryProofs/HodgeCompComparisonMiddleStatement.lean}{GlobalTheory.T.hodge_comp_comparison_middle}
\leanref{Development/AutomorphicSpectrum/Proofs/GeometryProofs/HodgeCompComparisonRightStatement.lean}{GlobalTheory.T.hodge_comp_comparison_right}
\begin{proof}
For a group scheme $G$ over a scheme $T$ and a point $x:\operatorname{Spec}R\to T$, the Lie algebra of $G$ at $x$ is the
$R$-module of morphisms $\operatorname{Spec}R[\varepsilon]/(\varepsilon^2)\to G$ over $x$ that restrict to the unit section.
The Lie algebra of $\mathcal A_{Z_1}$ at a point $P$ is the one at $\operatorname{Spec}\mathcal O_{Z_1,P}$, an
$\mathcal O_{Z_1,P}$-submodule of $\mathrm{Lie}_1$ since $\mathcal O_{Z_1,P}\subset\Omega_1$, and in the same way for $Z$.

(1) Since $f=f_1\circ h$, and $\mathcal A_Z$ and $\mathcal A_{Z_1}$ are identified with $f^*\mathcal A$ and $f_1^*\mathcal A$
with their actions and cubical line bundles, $\mathcal A_Z$ is identified with $h^*\mathcal A_{Z_1}$ with these structures,
compatibly with the maps to $\mathcal A$. For $e\in\mathrm{Lie}_1$ let $j_A(e)$ be the composite of
$\operatorname{Spec}\Omega[\varepsilon]\to\operatorname{Spec}\Omega_1[\varepsilon]$, given by $\iota$, with $e$; it is a tangent
vector of $h^*\mathcal A_{Z_1}=\mathcal A_Z$ at the generic point of $Z$. Then $j_A$ is additive with
$j_A(ae)=\iota(a)j_A(e)$, and it commutes with the actions of $\mathcal O_E$ and with the maps to the Lie algebra of
$\mathcal A$, since these are given by composing with morphisms of $\mathcal A_{Z_1}$ and of $\mathcal A$. As
$\mathcal A_{Z_1}$ is smooth over $Z_1$, its sheaf of differentials is finite locally free \cite[Tag 02G1]{Sta}, and its
formation commutes with base change \cite[Tag 01V0]{Sta}. The Lie algebra at a point $x$ is the $R$-dual of the pullback
of the differentials along the unit section at $x$. So the Lie algebra of $\mathcal A_{Z_1}$ at a point $P$ is free over
$\mathcal O_{Z_1,P}$; $j_A$ carries a basis of $\mathrm{Lie}_1$ to a basis of $\mathrm{Lie}$, and so carries linearly
independent families to linearly independent families; and for a point $P$ of $Z$ the Lie algebra of $\mathcal A_Z$ at
$P$ is $\mathcal O_{Z,P}$ tensored with the Lie algebra of $\mathcal A_{Z_1}$ at $h(P)$, which in $\mathrm{Lie}$ is the
$\mathcal O_{Z,P}$-module generated by its image under $j_A$.

(2) Since $f=f_1\circ h$, the generic points of $Z_1$ and $Z$ map to the same generic point of an irreducible component of
$\mathcal X$, and the embedding of its residue field $\Omega_0$ in $\Omega$ is $\iota$ composed with that in $\Omega_1$. As
the embedding of $F_c$ in $\Omega$ is $\iota$ composed with that in $\Omega_1$ and $j_A$ commutes with the actions of
$\mathcal O_E$, a vector $e\in\mathrm{Lie}_1$ lies in $\mathrm{Lie}_{1,s}$ if and only if $j_A(e)$ lies in $\mathrm{Lie}_s$.
So the $r_s$ vectors $j_A(b_{s,i})$ are linearly independent in $\mathrm{Lie}_s$, which has dimension $r_s$, the rank of
$\mathcal V_s$, and they span it. For the genericity condition let $c_{s,l}$ and $u_{s,lk}$ be as in
Definition~\ref{cv-origin-hodge-component}(2) for the flag of $Z_1$. The same vectors $c_{s,l}$ serve for $Z$; the elements
$\iota(u_{s,lk})$ are algebraically independent over $\Omega_0$; $\langle j_A(b_{s,1}),\dots,j_A(b_{s,i})\rangle$ is spanned
by the images of the $c_{s,k}+\sum_{l>k}\iota(u_{s,lk})c_{s,l}$ with $k\le i$, since $j_A$ is $\iota$-semilinear and
compatible with the maps to the Lie algebra of $\mathcal A$; and
$[\Omega:\Omega_0(\iota(u))]=[\Omega:\iota(\Omega_1)]\,[\Omega_1:\Omega_0(u)]=[\Omega:\iota(\Omega_1)]\deg$.

(3) Let $P$ be a point of $Z$ over no prime. Every prime is a unit at $P$, so also at $h(P)$, since the homomorphism
$\mathcal O_{Z_1,h(P)}\to\mathcal O_{Z,P}$ is local; so $h(P)$ lies over no prime, and the discriminant of $E$ is a unit
at $h(P)$. At $P$ and at $h(P)$, $\mathcal V_s$ and $\mathcal V^1_s$ are the parts in $\mathrm{Lie}_s$ and
$\mathrm{Lie}_{1,s}$ of the Lie algebras $L$ and $L_1$ there. By the flag of $Z_1$, $\mathcal V^1_s$ has at $h(P)$ a basis
$b'$ adapted to the flag of $b_s$. The values $s(a)$ for $a\in\mathcal O_E$ are integral over $\mathbf Z$, so regular at
$h(P)$, as $Z_1$ is normal, and $L_1$ is stable under $\mathcal O_E$. As the discriminant of $E$ is invertible, the
$\mathcal O_{Z,P}$-span of $j_A(L_1)$ meets $\mathrm{Lie}_s$ in the $\mathcal O_{Z,P}$-span of
$j_A(L_1\cap\mathrm{Lie}_{1,s})$. By (1) the former is $L$, so $\mathcal V_s$ at $P$ is spanned by the family $j_A(b')$,
which is linearly independent, hence a basis, and adapted to the flag of $j_A(b_s)$ because $j_A$ carries spans to spans.
\end{proof}
% Lean: (1) is CoverLie.main: cancel (SABMap from f = f_1 o h), exists_hom (j from the cartesian square, through
% LieMove.pushModel and tan_up), map_diff for the actions, and T.lie_lattice_dominant_compat = LiePull.main (bc for the
% basis, from Smooth G.hom; lie_model_free_over_local_ring; lieLat_carrier, core and ext_up for the lattices). (2) is
% FlagRead.generic_read: exists_iso of the residue fields Omega_0 (hgen), mem_eigSp_iff, span_eq_of_indep with the rank of
% B.vb at the generic point, and genericFlag_pull with the degree finrank * Fl.deg. (3) is VbGen.isFlagVBOn_genFib:
% regAt_dominates for the units, Lat.patch_of_not (V = dR off the primes of S), discr_inv_mem, integral_element_regular,
% T.lattice_eigen_field_extension, flagStep_comp and span_image_span.

\begin{lemma}\label{cv-origin-change-bases}
In the setting of Definition~\ref{cv-origin-hodge-component}, let $S$ be a finite set of primes that contains
every prime dividing $d_E$, let $\Omega$ be a field of characteristic zero, and let $Z$ be a Hodge component of
$\mathcal X$ at $S$ with coordinates $y\in\Omega^{\sqcup_kI_k}$ whose generic point lies over a point $y_0$ of
$X$, $f\circ\xi=j\circ y_0$. Suppose that $Z$ satisfies condition (1) of Lemma~\ref{cv-origin-change-projective}
at every $p\in S$ and that $\Omega$ contains a primitive $p^{R_p}$-th root of unity for $p\in S$. Let $b_s$ be
bases of the $\mathrm{Lie}_s$ that satisfy the genericity condition of
Definition~\ref{cv-origin-hodge-component}(2) with some degree $\deg$. Then there are a point $y'$ that
contains $y$ as some of its blocks, a Hodge component $Z'$ of $\mathcal X$ at $S$ with coordinates $y'$, the
same integers $R_p$ and the same embedding of $F_c$, satisfying condition (1) of
Lemma~\ref{cv-origin-change-projective} at every $p\in S$, a morphism $g:Z'\to Z$ that maps the generic point
to the generic point by the identity of $\Omega$, with $f'=f\circ g$, and a flag of $Z'$ of degree $\deg$,
whose bases have the same images as the $b_{s,i}$ in the Lie algebra of $\mathcal A$, with root lines of some
multidegrees with generators, such that near every point of $Z'$ the vectors of $\mathrm{Lie}_s$ in the Lie
algebra of $\mathcal A_{Z'}$ have a basis adapted to the flag, and such that $g$ is an isomorphism over every
point $P$ of $Z$ over no prime at which every $\mathcal V_s$ has a basis adapted to the flag of $b_s$, in the
sense of Lemma~\ref{cv-origin-change-projective}(2).
\end{lemma}
\leanref{Development/AutomorphicSpectrum/Proofs/GeometryProofs/HodgeCptLabelBasesStatement.lean}{GlobalTheory.T.hodge_cpt_label_bases}
\begin{proof}
\openstep{The construction of a Hodge component over $Z$ on which the given generic bases extend to a flag
with root lines, isomorphic to $Z$ where the bases are adapted, is not yet written out.}
% Lean: GlobalTheory.T.hodge_cpt_label_bases, under Out.WeilSt. Size: about 37,000 lines, much of it shared
% with Lemma cv-origin-change-projective.
\end{proof}

\begin{lemma}\label{cv-origin-change-smooth}
Let $\Omega$ be a field of characteristic zero, let $X_1$ and $X_2$ be schemes with morphisms
$\xi_l:\operatorname{Spec}\Omega\to X_l$ of whose images every point is a specialization, covered by the affine
open subschemes of the charts of points $y_1$ and $y_2$ with values in $\Omega$, and let $g:X_2\to X_1$ be a
morphism with $g\circ\xi_2=\xi_1$. Suppose that $\mathbf Q[y_{1,k,i}/y_{1,k,\kappa_k}]$ is smooth over
$\mathbf Q$ for every chart $\kappa$ of $y_1$, and that $g$ is an isomorphism over every point of $X_1$ over no
prime, in the sense of Lemma~\ref{cv-origin-change-projective}(2). Then
$\mathbf Q[y_{2,k,i}/y_{2,k,\kappa_k}]$ is smooth over $\mathbf Q$ for every chart $\kappa$ of $y_2$.
\end{lemma}
\leanref{Development/AutomorphicSpectrum/Proofs/GeometryProofs/CoordsCompatibleAlongMapStatement.lean}{GlobalTheory.T.smoothGenericFibre_of_dominates}
\begin{proof}
Let $\kappa$ be a chart of $y_2$, put $A=\mathbf Q[y_{2,k,i}/y_{2,k,\kappa_k}]$, and let $A^{\mathbf Z}_\kappa$ be as in
Lemma~\ref{cv-origin-snc-charts} for $y_2$. The algebra $A$ is generated by finitely many elements, so it is of finite
type over the noetherian ring $\mathbf Q$. It is smooth if it is smooth at every prime \cite[Tag 00TC]{Sta}, and by
\cite[Tag 07VH]{Sta} it is smooth at a prime $\mathfrak P$ if $A_{\mathfrak P}$ is formally smooth over $\mathbf Q$ for the
$\mathfrak P$-adic topology. This holds if $A_{\mathfrak P}$ is formally smooth over $\mathbf Q$: a homomorphism from
$A_{\mathfrak P}$ to $C/J$, with $J^2=0$, that kills a power $\mathfrak P^m$ has lifts to $C$, and every lift kills
$\mathfrak P^{2m}$.

Let $\mathfrak P$ be a prime of $A$. Lemma~\ref{cv-origin-snc-charts}(2), applied to $X_2$ and the prime
$\mathfrak P\cap A^{\mathbf Z}_\kappa$, gives a point $P$ of $X_{2,\kappa}$ at which the elements of $A^{\mathbf Z}_\kappa$ outside
$\mathfrak P$ are the units. A positive integer is a unit of $A$, so it lies outside $\mathfrak P$ and is a unit at
$P$: $P$ lies over no prime. Every $a\in A$ has $na\in A^{\mathbf Z}_\kappa$ for some integer $n\ge1$. So $A$ lies in the local
ring of $X_2$ at $P$, the elements of $A$ that are units at $P$ are those outside $\mathfrak P$, and by
Lemma~\ref{cv-origin-snc-charts}(1) an element of $\Omega$ is regular at $P$ if and only if it is $r/s$ with $r,s\in A$
and $s\notin\mathfrak P$. That is, the local ring of $X_2$ at $P$ is $A_{\mathfrak P}$, as subrings of $\Omega$.

The point $Q=g(P)$ lies over no prime: a positive integer is regular at $Q$ and a unit at $P$, so it is a unit at $Q$ by
Lemma~\ref{cv-origin-model-pullback}(1), with the identity of $\Omega$. By hypothesis $g$ is an isomorphism over $Q$,
and $P$ lies over $Q$, so the local rings of $X_2$ at $P$ and of $X_1$ at $Q$ are equal. The point $Q$ lies in a chart
$\kappa'$ of $y_1$; let $A_1=\mathbf Q[y_{1,k,i}/y_{1,k,\kappa'_k}]$ and let $\mathfrak P_1$ be the set of its elements
that are not units at $Q$. As for $P$, Lemma~\ref{cv-origin-snc-charts}(1) shows that the local ring of $X_1$ at $Q$ is
$(A_1)_{\mathfrak P_1}$. The algebra $A_1$ is smooth, hence formally smooth, over $\mathbf Q$ \cite[Tag 00TN]{Sta}, the
localization $A_1\to(A_1)_{\mathfrak P_1}$ is formally smooth \cite[Tag 04EG]{Sta}, and so is the composite
\cite[Tag 031H]{Sta}. So $A_{\mathfrak P}=(A_1)_{\mathfrak P_1}$ is formally smooth over $\mathbf Q$.
\end{proof}
% Lean: SmoothChange.smoothGenericFibreM_of_isoAt. adjoin_blocks_eq identifies Q[y_{k,i}/y_{k,c_k}] with the ring of
% the Segre ratios; smooth_chRing (FiniteType, smoothLocus_eq_univ_iff) and formallySmooth_atPrime: exists_point_of_prime
% (prime_chart_existsUnique_point), exists_image_genFib (RegAtDom.regAt_dominates), IsoAt.regAt_eq, exists_prime_of_point,
% mem_regAt_iff and formallySmooth_of_mem_iff (isLocalization_of_mem_iff, FormallySmooth.of_isLocalization,
% FormallySmooth.comp). Mathlib's smoothLocus_eq_univ_iff is the local criterion that the book cites as Tags 00TC and 07VH.

\begin{lemma}\label{cv-origin-change}
In the setting of Definition~\ref{cv-origin-hodge-component}, let $S$ be an analytic system of type
$(\Psi,N)$, let data at one component as in Definition~\ref{cv-origin-tower}(2b) be given, with $S$ in place of
$S_c$, and let $R$ be a flag model of type $(\Psi,N)$ with data as in (2c) for $R$ in place of $M(B,c)$ and
the empty set in place of $S_B$, with these data of (2b) and $\deg=1$. Let $Z$, $y\in\Omega^{\sqcup_kI_k}$ and
$b_{s,i}$ be the Hodge component, the point of $R$ and the bases of the flag of these data. Let $S'$ be a
finite set of primes that contains every prime dividing $d_E$. Then there are a field $\Omega'$ of
characteristic zero with an embedding $\iota:\Omega\to\Omega'$, a point $y'\in\Omega'^{\sqcup_kI'_k}$ with an
injection $k\mapsto k'$ and bijections $i\mapsto i'$ from $I_k$ to $I'_{k'}$ with $y'_{k',i'}=\iota(y_{k,i})$, a
Hodge component $Z'$ of $\mathcal X$ at $S'$ with coordinates $y'$, a morphism $h:Z'\to Z$ that maps the
generic point of $Z'$ to that of $Z$ compatibly with $\iota$, and a flag of $Z'$ with bases $b'_{s,i}$ and
root lines of some multidegrees $a'_{\tau,i}$ with generators $e'_{\tau,i}$, such that:
\begin{enumerate}
\item $f'=f\circ h$, and the embedding of $F_c$ in $\Omega'$ is $\iota$ composed with that in $\Omega$;
\item each $b'_{s,i}$ has the same image in the Lie algebra of $\mathcal A$ as $b_{s,i}$, after the base change by
  $\iota$;
\item the flag of $Z'$ has degree $[\Omega':\iota(\Omega)]$;
\item $Z'$ satisfies condition (1) of Lemma~\ref{cv-origin-change-projective} at every $p\in S'$, and $R_p=4$ for
  $p\in S'$;
\item $\mathbf Q[y'_{k,i}/y'_{k,\kappa_k}]$ is smooth over $\mathbf Q$ for every chart $\kappa$ of $y'$.
\end{enumerate}
\end{lemma}
\leanref{Development/AutomorphicSpectrum/Proofs/GeometryProofs/RootOriginGoodPrimesFiniteStatement.lean}{GlobalTheory.T.root_origin_good_primes_finite}
\leanref{Development/AutomorphicSpectrum/Proofs/GeometryProofs/RootOriginGoodPrimesFinite.lean}{GlobalTheory.CompanionTheory.Ar.ChangeGeo.exists_cover}
\begin{proof}
The field $\Omega$ has characteristic zero, since $F_c$ embeds in it. The generic point of $Z$ maps to the
generic point $\eta_0$ of an irreducible component of $\mathcal X$, through $\Omega_0\subset\Omega$. Let
$\xi_1:\operatorname{Spec}\Omega\to\mathcal X$ be the composite. Lemma~\ref{cv-origin-change-extension}, applied to
$\xi_1$ and $S'$, gives a finite extension $\iota:\Omega\to\Omega'$ in which the points of order dividing $p^4$
are defined and which contains primitive $p^4$-th roots of unity, for $p\in S'$.
Lemma~\ref{cv-origin-change-coords}, applied to $Z$ with its charts and to $\Omega'$, gives a point $y_1$
containing $\iota(y)$ as blocks, a scheme $X_1$ with integrally closed local rings and smooth generic fibre in
the sense of (5), and $h_1:X_1\to Z$ over $\iota$. The generic point of $X_1$ maps to $\eta_0$, so by
Lemma~\ref{cv-origin-change-factor} its image in $\mathcal X$ is $j\circ y_0$ for a point $y_0$ of $X$.
Lemma~\ref{cv-origin-change-projective}, applied to $X_1$, $f\circ h_1$ and $R_p=4$ for $p\in S'$, gives a point
$y_2$ containing $y_1$ as blocks, a Hodge component $Z_2$ at $S'$ with coordinates $y_2$, $R_p=4$ and condition
(1) of that lemma, and $g_2:Z_2\to X_1$ that is an isomorphism over the points of $X_1$ over no prime. The
embedding of $F_c$ in $\Omega'$ given by $Z_2$ is $\iota$ composed with that of $Z$, since both describe the
morphism $\operatorname{Spec}\Omega'\to\operatorname{Spec}\mathcal O_c$ through $\mathcal X$, and $F_c$ is the
fraction field of $\mathcal O_c$.

Lemma~\ref{cv-origin-change-comparison}, applied to $Z_1=Z$, $Z_2$ and $h_1\circ g_2$, gives $j_A$, and the
families $j_A(b_s)$ satisfy the genericity condition with degree $[\Omega':\iota(\Omega)]\cdot1$, and are adapted
to $\mathcal V_s$ at every point of $Z_2$ over no prime. Lemma~\ref{cv-origin-change-bases}, applied to $Z_2$ and
the bases $j_A(b_s)$, gives a point $y'$ containing $y_2$ as blocks, a Hodge component $Z'$ at $S'$ with
coordinates $y'$, $R_p=4$ and condition (1), a morphism $g':Z'\to Z_2$ that is an isomorphism over every point
of $Z_2$ over no prime, and a flag of $Z'$ of degree $[\Omega':\iota(\Omega)]$ with root lines, whose bases have
the images of the $j_A(b_{s,i})$, which are those of the $b_{s,i}$ after base change by $\iota$. By
Lemma~\ref{cv-origin-change-smooth}, applied twice, to $g_2$ and to $g'$, the generic fibres of $y_2$ and of $y'$
are smooth. Put $h=h_1\circ g_2\circ g'$; the blocks of $y$ are carried into those of $y'$ through $y_1$ and
$y_2$, and (1) to (5) hold.
\end{proof}
% Lean: ChangeGeo.exists_cover: component_generic_is_generic_point, exists_finite_extension_primes,
% coords_after_finite_extension (ResKStatement), generic_point_factors_open, hodge_cpt_projective_model with
% R = 4 (Out.WeilSt), cst_ext for the constants, hodge_comp_comparison_left, _middle and _right,
% hodge_cpt_label_bases (Out.WeilSt), and smoothGenericFibre_of_dominates twice.

\subsection{Good origins and good towers}\label{subsection-cv-good-origins-towers}

\begin{lemma}[Grothendieck {\cite[Expos\'e~XII, Th\'eor\`eme~1.1 and Proposition~2.1]{Gro71}}]\label{cv-origin-model-resolution-coords}
Let $I$ be an ideal of $\mathbf C[x_1,\dots,x_n]$ and $w\in\mathbf C^n$ a common zero of $I$ such that the local
ring $\mathcal O$ of $\mathbf C[x_1,\dots,x_n]/I$ at $w$ is regular of dimension $e$, and let $r_1,\dots,r_e$ be
polynomials whose images generate the maximal ideal of $\mathcal O$. Then there are an open neighbourhood $U$ of $w$
in $\mathbf C^n$, an open set $V\subset\mathbf C^e$ and a holomorphic map $g:V\to\mathbf C^n$ that maps $V$
bijectively onto the set of common zeros of $I$ in $U$, with $(r_1(g(t)),\dots,r_e(g(t)))=t$ for every $t\in V$.
\end{lemma}
\leanref{ReciprocityTheory/AnalyticForms/RegularLocalRings/AnChart.lean}{ReciprocityTheory.AnForm.R.exists_holo_coords}

The completions of $\mathcal O$ and of the analytic local ring at $w$ of the zeros of $I$ are isomorphic
\cite[Expos\'e~XII, Th\'eor\`eme~1.1]{Gro71}, so the analytic local ring is regular of dimension $e$
\cite[Expos\'e~XII, Proposition~2.1]{Gro71}, with maximal ideal generated by the $r_l$, and the zeros of $I$ near $w$
form a complex manifold of dimension $e$ with the $r_l$ as holomorphic coordinates.

\begin{lemma}\label{cv-origin-model-resolution-regular}
Let $\Omega$ be a field of characteristic zero, $D\ge0$, and $y\in\Omega^{\sqcup_kI_k}$ a point whose kernel $P$ is an
arithmetic subvariety of dimension $D+1$. Let $\kappa$ be a chart of $y$ such that
$A_\kappa=\mathbf Q[y_{k,i}/y_{k,\kappa_k}]$ is smooth over $\mathbf Q$, and let $I_\kappa$ be the ideal of
$\mathbf C[x_{k,i}]$ generated by $P$ and the $x_{k,\kappa_k}-1$ for $k\in K$. Then at every common zero $w$ of
$I_\kappa$ the local ring $\mathcal O_w$ of $\mathbf C[x_{k,i}]/I_\kappa$ at $w$ is regular of dimension $D$.
\end{lemma}
\leanref{Development/AutomorphicSpectrum/Proofs/MetricProofs/ExistsAnalyticLogResolution.lean}{GlobalTheory.CompanionTheory.Ar.PB.SmoothMan.isRegularLocalRing_chart}
\leanref{Development/AutomorphicSpectrum/Proofs/MetricProofs/ExistsAnalyticLogResolution.lean}{GlobalTheory.CompanionTheory.Ar.PB.SmoothMan.card_le_of_smooth}
\begin{proof}
Since $P\cap\mathbf Z=0$, $P$ generates a prime of $\mathbf Q[x_{k,i}]$ with quotient of dimension $D+|K|$, and the
dimension of a scheme of finite type over a field at a point does not change under extension of the field
\cite[Tag 00P4]{Sta}; so every minimal prime of the ideal $P_{\mathbf C}$ of $\mathbf C[x_{k,i}]$ generated by $P$
has quotient of dimension $D+|K|$. So the maximal ideal of $w$ has height at least $D+|K|$ modulo $P_{\mathbf C}$, and
at least $D$ modulo $I_\kappa$, since the $|K|$ elements $x_{k,\kappa_k}-1$ lower the height by at most $|K|$. Conversely
let $\mathfrak q\supset I_\kappa$ be prime, $B=\mathbf C[x_{k,i}]/\mathfrak q$, and $\theta:\mathbf C[x_{k,i}]\to
B[t_k:k\in K]$, $x_{k,i}\mapsto t_k\bar x_{k,i}$. A multihomogeneous element of $P$ of multidegree $a$ maps to
$t^a$ times its image in $B$, which is $0$; so the image of $\theta$ has dimension at most $D+|K|$. It contains
$t_k=\theta(x_{k,\kappa_k})$ and the $t_k\bar x_{k,i}$, so $B[t_k]$ is algebraic over it and its transcendence degree is
$\dim B+|K|$. Hence $\dim B\le D$, and $\mathcal O_w$ has dimension $D$.

The cotangent space of $\mathcal O_w$ has dimension $|\sqcup_kI_k|-\dim J$, where $J\subset\mathbf C^{\sqcup_kI_k}$ is
the span of the gradients at $w$ of the elements of $I_\kappa$. Every polynomial $g$ over $\mathbf Q$ with
$g(y_{k,i}/y_{k,\kappa_k})=0$ lies in $I_\kappa$: homogenizing $g$ in each block with $x_{k,\kappa_k}$ and clearing
denominators gives an element of $P$ congruent to a nonzero integer multiple of $g$ modulo the $x_{k,\kappa_k}-1$. So
$g(y_{k,i}/y_{k,\kappa_k})\mapsto g(w)$ defines a homomorphism $\psi:A_\kappa\to\mathbf C$, and a vector $v$ with
$\partial_vh(w)=0$ for every $h\in I_\kappa$ defines a homomorphism $A_\kappa\to\mathbf C[\varepsilon]/(\varepsilon^2)$,
$g(y_{k,i}/y_{k,\kappa_k})\mapsto g(w)+\varepsilon\,\partial_vg(w)$, whose $\varepsilon$-part is a $\mathbf Q$-derivation
$A_\kappa\to\mathbf C$ along $\psi$; it determines $v$, whose entries are its values at the $y_{k,i}/y_{k,\kappa_k}$. These
derivations are the linear forms on $\mathbf C\otimes_{A_\kappa}\Omega_{A_\kappa/\mathbf Q}$. As $A_\kappa$ is smooth over
$\mathbf Q$, $\Omega_{A_\kappa/\mathbf Q}$ is projective of constant rank, equal to its rank at the generic point, the
transcendence degree $D$ of the fraction field of $A_\kappa$. So these vectors $v$, which form a space of dimension
$|\sqcup_kI_k|-\dim J$, lie in a space of dimension at most $D$, and the cotangent space of $\mathcal O_w$ has dimension at
most $D=\dim\mathcal O_w$. Hence $\mathcal O_w$ is regular.
\end{proof}
% Lean: dimension from pcEquidim (extQ_spec and ringKrullDim_of_mem_minimalPrimes_extC) and
% ringKrullDim_locAt_chart (ConeChart, ConeChartA); cotangent space from finrank_cotangentSpace_locAt and
% card_le_finrank_jacSpanAt: ker_le_chart, exists_unique_lift (the maps to C and to TrivSqZeroExt C C),
% exists_span_of_fibre, finrank_fibre_kaehler (rank of Omega at the stalks, isLocallyConstant_rankAtStalk,
% rank_kaehler_eq_trdeg), trdeg_blockFnField and card_le_finrank_add.

\begin{lemma}\label{cv-origin-model-resolution}
Let $\Omega$ be a field of characteristic zero, $D\ge0$, and $y\in\Omega^{\sqcup_kI_k}$ a point whose kernel is an
arithmetic subvariety $X$ of dimension $D+1$, such that $\mathbf Q[y_{k,i}/y_{k,\kappa_k}]$ is smooth over
$\mathbf Q$ for every chart $\kappa$ of $y$. Then there are a compact complex manifold $\tilde X$ of dimension $D$
and holomorphic maps $\lambda_x:U_x\to\hat X(\mathbf C)$ as in Definition~\ref{cv-origin-analytic}(3).
\end{lemma}
\leanref{Development/AutomorphicSpectrum/Proofs/MetricProofs/ExistsAnalyticLogResolutionStatement.lean}{GlobalTheory.T.exists_analytic_log_resolution}
\leanref{ReciprocityTheory/ReductiveGroups/ArithmeticGeometry/BirationalGeometry/InterRes.lean}{ReciprocityTheory.ChartRestriction.G.Geo.AnLogRes}
\begin{proof}
Let $P$ be the kernel of $y$, and let $\tilde X$ be the image of $\hat X(\mathbf C)$ in
$\prod_k\mathbf P^{I_k}(\mathbf C)$, with the induced topology. It is Hausdorff, and compact as the image of the
compact set of $z\in\hat X(\mathbf C)$ with $\sum_i|z_{k,i}|^2=1$ for every $k$. For a choice $\kappa=(\kappa_k)$ of
indices let $\tilde X_\kappa$ be the open set of classes with $z_{k,\kappa_k}\neq0$ for every $k$, and
$a_\kappa:\tilde X_\kappa\to\mathbf C^{\sqcup_kI_k}$ the map to the affine coordinates $z_{k,i}/z_{k,\kappa_k}$. It is
a homeomorphism onto the set of common zeros of the ideal $I_\kappa$ of
Lemma~\ref{cv-origin-model-resolution-regular}, and on $\tilde X_\kappa\cap\tilde X_{\kappa'}$ each coordinate of
$a_{\kappa'}$ is a polynomial in $a_\kappa$ divided by a polynomial in $a_\kappa$ that does not vanish there. If
$\tilde X_\kappa$ is not empty, then $\kappa$ is a chart of $y$: if $y_{k,\kappa_k}=0$, then $x_{k,\kappa_k}\in P$
vanishes on $\hat X(\mathbf C)$.

Let $x\in\tilde X_\kappa$. By Lemma~\ref{cv-origin-model-resolution-regular} the local ring of
$\mathbf C[x_{k,i}]/I_\kappa$ at $a_\kappa(x)$ is regular of dimension $D$, so its maximal ideal is generated by the
images of $D$ polynomials $r_1,\dots,r_D$. By Lemma~\ref{cv-origin-model-resolution-coords}, $x'\mapsto r(a_\kappa(x'))$
is a homeomorphism from an open neighbourhood of $x$ onto an open set $V\subset\mathbf C^D$, whose inverse composed
with $a_\kappa$ is holomorphic on $V$. Since the coordinates $a_{\kappa'}$ are holomorphic functions of $a_\kappa$
where both are defined, these charts are holomorphically compatible, and every polynomial in the coordinates
$a_\kappa$ is holomorphic on $\tilde X_\kappa$. As $\tilde X$ is compact, finitely many of the charts cover it. So
$\tilde X$ is a compact complex manifold of dimension $D$.

For $x\in\tilde X$ choose $\kappa$ with $x\in\tilde X_\kappa$, put $U_x=\tilde X_\kappa$, and let $\lambda_x(x')$ be the
point of $\hat C$ with coordinates $a_\kappa(x')$, whose coordinate at $(k,\kappa_k)$ is $1$. It lies in
$\hat X(\mathbf C)$, and it is holomorphic. At each point of $U_x\cap U_{x'}$ the lifts $\lambda_x$ and $\lambda_{x'}$
represent the same point of $\prod_k\mathbf P^{I_k}(\mathbf C)$, so they agree up to the action of
$(\mathbf C^\times)^K$. A point $z\in\hat X(\mathbf C)$ lies in the orbit of $\lambda_x(x)$ for its class $x$, and if
$\lambda_x(x)$ and $\lambda_{x'}(x')$ lie in one orbit, then $x=x'$. Finally let $\kappa'$ be a choice of indices with
$\lambda_x(x)_{k,\kappa'_k}\neq0$ for every $k$, so that $x\in\tilde X_{\kappa'}$, and let $\phi$ be a chart of
$\tilde X$ at $x$. Near $x$ the coordinates $a_\kappa$ are quotients of polynomials in $a_{\kappa'}$ by a polynomial
that does not vanish at $x$, and a chart $r\circ a_\kappa$ at $x$ as above is a holomorphic function of $a_\kappa$. So
there is a holomorphic map $G$ near $a_{\kappa'}(x)$ with $G\circ a_{\kappa'}\circ\phi^{-1}$ the identity near
$\phi(x)$, and the affine coordinates $a_{\kappa'}\circ\lambda_x\circ\phi^{-1}$ of $\lambda_x$ in the chart $\kappa'$
have injective derivative at $\phi(x)$.
\end{proof}
% Lean: nonempty_anLogRes_of_smoothGenericFibre = nonempty_anLogRes_of_algRes applied to SmoothMan.nonempty_algRes,
% with J the unit ideal, whose zero set is empty (not_mem_conePts_top'), and pcContract. nonempty_algRes: M = Pts
% (pC P), compact by Top.compactSpace_pts (the image under proj of the sphere bundle), src and coord =
% affCoord renumbered (compact_open_embedding), image and trans (chart_image_and_trans), lifts, compat, surj, inj
% (chart_lifts_package), inv (ChartInv.chart_coord_ratio), regular and dim (isRegularLocalRing_chart), local
% parameters (exists_params_of_regular), and the trivial boundary data (bd = top, rc = 0, sncAt with J empty).
% AlgRes.exists_atlas: the charts evalFam r o coord of AlgRes.exists_goodChart (exists_holo_coords) and a finite
% subcover; the immersion is AlgRes.fderiv_affCoord_ker_eq_bot (the left inverse G = chart o psi.symm o evalFam r o
% (F/q) from inv).

\begin{situation}\label{cv-origin-lift-situation}
Let $K_0$ and $K$ be finite sets, let $I^0_k$ ($k\in K_0$) and $I_k$ ($k\in K$) be finite nonempty sets, let
$k\mapsto k'$ be an injection from $K_0$ to $K$, and let $i\mapsto i'$ be bijections from $I^0_k$ to $I_{k'}$. Let
$R_0$ and $R$ be the rings $\mathbf Z[x_{k,i}]$ of Definition~\ref{lit-multiprojective} for the index sets
$(I^0_k)_{k\in K_0}$ and $(I_k)_{k\in K}$, with the sets $\hat C_0$ and $\hat C$ of that definition. For $f\in R_0$ let
$f'\in R$ be its renaming $x_{k,i}\mapsto x_{k',i'}$; for an ideal $P\subset R$ put $P_0=\{f\in R_0:f'\in P\}$; and for
$z\in\hat C$ put $z|_0=(z_{k',i'})_{k\in K_0,\,i\in I^0_k}$, a point of $\hat C_0$. The elements of $K$ not of the form
$k'$ are the \emph{fibre blocks}. For an ideal $P$ of $R_0$ or of $R$, let $\hat X_P(\mathbf C)$ be the set of points of
$\hat C_0$ or of $\hat C$ at which every element of $P$ vanishes. A point $z\in\hat C$ lies \emph{over} $w\in\hat C_0$ if
$z|_0$ lies in the orbit of $w$ under $(\mathbf C^\times)^{K_0}$, acting blockwise.
\end{situation}
\leanref{ReciprocityTheory/ReductiveGroups/ArithmeticGeometry/ArithmeticIntersection/MixedFieldContraction/GenericFibreKerTarget.lean}{ReciprocityTheory.ChartRestriction.G.AI.I4.imageZ}
\leanref{GroundTheory/HeightsAndTheta/Heights/Intersections/PtsOverPt.lean}{GroundTheory.Heights₂.Heights.Inter.ptsOverPt}

\begin{definition}\label{cv-origin-lift-fibre}
In Situation~\ref{cv-origin-lift-situation}, let $P\subset R$ be an ideal.
\begin{enumerate}
\item A \emph{generic fibre} of $P$ consists of fields $\Phi\subset\Phi'$ and families $x=(x_{k,i})_{k\in K_0}$ in $\Phi$
  and $x'=(x'_{k,i})_{k\in K}$ in $\Phi'$ such that the homomorphism $R_0\to\Phi$, $f\mapsto f(x)$, has kernel $P_0$;
  every element of $\Phi$ is $f(x)/g(x)$ with $f,g\in R_0$ and $g(x)\neq0$; the homomorphism $R\to\Phi'$,
  $f\mapsto f(x')$, has kernel $P$; and $x'_{k',i'}=x_{k,i}$ for $k\in K_0$ and $i\in I^0_k$.
\item A \emph{chart} of a generic fibre is a family $j=(j_k)$ of indices $j_k\in I_k$, one for each fibre block $k$,
  with $x'_{k,j_k}\neq0$. Its \emph{chart field} is the subfield $\Phi(r_j)$ of $\Phi'$ generated over $\Phi$ by the
  ratios $r_{j,k,i}=x'_{k,i}/x'_{k,j_k}$, for the fibre blocks $k$ and $i\in I_k$.
\end{enumerate}
\end{definition}
\leanref{ReciprocityTheory/ReductiveGroups/ArithmeticGeometry/ArithmeticIntersection/MixedFieldContraction/GenericFibreKerTarget.lean}{ReciprocityTheory.ChartRestriction.G.AI.I4.GenFibre}
\leanref{GroundTheory/HeightsAndTheta/Heights/Intersections/PtsOverChart.lean}{GroundTheory.Heights₂.Heights.Inter.ptsOverChart}

\begin{lemma}\label{cv-origin-lift-scaling}
Let $F\subset\Omega'$ be fields, let $K$ be a finite set with finite sets $I_k$, let $x\in\Omega'^{\sqcup_kI_k}$, and let
$c=(c_k)$ be indices with $x_{k,c_k}\neq0$ for every $k$. Suppose that the kernel of $F[x_{k,i}]\to\Omega'$,
$f\mapsto f(x)$, is generated by multihomogeneous elements. Then the elements $x_{k,c_k}$, $k\in K$, are algebraically
independent over the field $F(x_{k,i}/x_{k,c_k}:k\in K,\ i\in I_k)$.
\end{lemma}
\leanref{GroundTheory/HeightsAndTheta/Heights/Intersections/StChartScalingIndep.lean}{GroundTheory.Heights₂.Heights.Inter.ChartScalingIndep.Statement}
\begin{proof}
Put $t_k=x_{k,c_k}$ and $u=(x_{k,i}/x_{k,c_k})_{k,i}$, so that $x_{k,i}=t_ku_{k,i}$. Let $\sum_{\alpha\in S}a_\alpha t^\alpha=0$
for a finite set $S\subset\mathbf N^K$ and elements $a_\alpha$ of $F(u)$. Each $a_\alpha$ is a quotient of two elements
of $F[u]$; multiplying the relation by the product of the denominators, which is nonzero, we may assume that
$a_\alpha=g_\alpha(u)$ with $g_\alpha\in F[x_{k,i}]$. Let $M\in\mathbf N^K$ be at least the degree in each block of
every $g_\alpha$. Multiplying each monomial of $g_\alpha$ by powers of the $x_{k,c_k}$ gives a multihomogeneous
$G_\alpha$ of multidegree $M$, congruent to $g_\alpha$ modulo the ideal generated by the $x_{k,c_k}-1$, and
$G_\alpha(x)=\prod_kt_k^{M_k}g_\alpha(u)$ by homogeneity, as $u_{k,c_k}=1$. So
$H=\sum_\alpha G_\alpha\prod_kx_{k,c_k}^{\alpha_k}$ vanishes at $x$, and its multihomogeneous component of multidegree
$M+\beta$ is $G_\beta\prod_kx_{k,c_k}^{\beta_k}$. The kernel is generated by multihomogeneous elements, so it contains
the multihomogeneous components of its elements. Hence $\prod_kt_k^{M_k+\beta_k}g_\beta(u)=0$, and $a_\beta=0$ for every
$\beta$.
\end{proof}
\leanref{GroundTheory/HeightsAndTheta/Heights/Intersections/PChartScalingIndep.lean}{GroundTheory.Heights₂.Heights.Inter.ChartScalingIndep.proof}

\begin{lemma}\label{cv-origin-lift-degree}
In Situation~\ref{cv-origin-lift-situation}, let $\Omega_0$ and $\Omega$ be fields, $\Omega$ of characteristic zero, let
$\iota:\Omega_0\to\Omega$ be an embedding with $[\Omega:\iota(\Omega_0)]=d\ge1$, and let $y_0\in\Omega_0^{\sqcup_kI^0_k}$
and $y\in\Omega^{\sqcup_kI_k}$ be points with $y_{k',i'}=\iota(y_{0,k,i})$. Suppose that the kernels of $y_0$ and of $y$
(Definition~\ref{lit-multiprojective}(4)) are arithmetic subvarieties of dimension $D+1$ with function fields $\Omega_0$
and $\Omega$, and let $P$ be the kernel of $y$. Then $P_0$ is the kernel of $y_0$, and $P$ has a generic fibre with a
chart whose chart field has degree $d$ over $\Phi$ (Definition~\ref{cv-origin-lift-fibre}).
\end{lemma}
\leanref{Development/AutomorphicSpectrum/Proofs/GeometryProofs/ArithVarietyReindexedStatement.lean}{GlobalTheory.T.arith_variety_reindexed}
\leanref{Development/AutomorphicSpectrum/Proofs/GeometryProofs/LogResolutionReindexed.lean}{GlobalTheory.CompanionTheory.Ar.Chg.DegBlock.imageZ_eq}
\begin{proof}
\emph{The kernel of $y_0$.} Let $\Theta:\Omega_0[t_k:k\in K_0]\to\Omega[t_k:k\in K]$ be $\iota$ on the coefficients and
$t_k\mapsto t_{k'}$; it is injective, as $k\mapsto k'$ is. As $y_{k',i'}t_{k'}=\Theta(y_{0,k,i}t_k)$, the image of $f'$
under $x_{k,i}\mapsto y_{k,i}t_k$ is $\Theta$ of the image of $f$ under $x_{k,i}\mapsto y_{0,k,i}t_k$. So $f'\in P$ if and
only if $f$ lies in the kernel of $y_0$.

\emph{The generic fibre.} Let $\Phi'$ be the fraction field of $\Omega[t_k:k\in K]$, put $x'_{k,i}=y_{k,i}t_k$, put
$x_{k,i}=x'_{k',i'}$ for $k\in K_0$, and let $\Phi\subset\Phi'$ be the subfield generated by the $x_{k,i}$. By definition
the kernel of $f\mapsto f(x')$ is $P$, and as $f(x)=f'(x')$, the kernel of $f\mapsto f(x)$ is $P_0$. The field $\Phi$ has
characteristic zero and is generated by the $x_{k,i}$, so its elements are quotients $f(x)/g(x)$ with $f,g\in R_0$. The
ideal $P$ is relevant, so each fibre block $k$ has an index $j_k$ with $y_{k,j_k}\neq0$, and then $x'_{k,j_k}\neq0$: $j$ is
a chart.

\emph{The degree.} The kernel of $y_0$ is relevant, so for each $k\in K_0$ there is $c_k\in I^0_k$ with
$y_{0,k,c_k}\neq0$; put $t'_k=x'_{k',c'_k}=x_{k,c_k}\neq0$. Let $F_1\subset\Phi'$ be the field generated by the ratios
$x'_{k,i}/x'_{k,l}$ within the blocks of $K$, and $F_0\subset F_1$ the field generated by the ratios $x_{k,i}/x_{k,l}$
within the blocks of $K_0$.
\begin{enumerate}
\item[(i)] $\Phi=F_0(t')$: the ratios of $x$ and the $t'_k$ lie in $\Phi$, and $x_{k,i}=(x_{k,i}/x_{k,c_k})t'_k$.
\item[(ii)] $\Phi(r_j)=F_1(t')$: by (i), $\Phi\subset F_1(t')$, and $r_j$ lies in $F_1$; conversely, a ratio within a block
  $k'$ is a ratio of $x$, which lies in $\Phi$, and a ratio $x'_{k,i}/x'_{k,l}$ within a fibre block is
  $r_{j,k,i}/r_{j,k,l}$.
\item[(iii)] The $t'_k$ are algebraically independent over $F_1$. The kernel of $\mathbf Q[x_{k,i}]\to\Phi'$,
  $f\mapsto f(x')$, is generated by the multihomogeneous elements that generate $P$, and $F_1$ is generated by the ratios
  of $x'$ at its chart given by the $c'_k$ on the blocks $k'$ and by $j$ on the fibre blocks. So
  Lemma~\ref{cv-origin-lift-scaling} shows that the entries of $x'$ at this chart, among them the $t'_k$, are
  algebraically independent over $F_1$.
\item[(iv)] If $[F_1:F_0]$ is finite, then $[F_1(t'):F_0(t')]=[F_1:F_0]$. By (iii), a basis of $F_1$ over $F_0$ stays
  linearly independent over $F_0[t']$, hence over $F_0(t')$; and it spans $F_1(t')$ over $F_0(t')$, as $F_1$ is finite over
  $F_0$.
\end{enumerate}
Each ratio $x'_{k,i}/x'_{k,l}$ is $y_{k,i}/y_{k,l}$, an element of the image of $\Omega$ under the constant embedding
$\Omega\to\Phi'$. As $\Omega$ is the function field of $P$, every element of $\Omega$ is $g(y)/h(y)$ with $g,h$
multihomogeneous of one multidegree, and so a quotient of polynomials in the ratios $y_{k,i}/y_{k,l}$. Hence $F_1$ is the
image of $\Omega$, and in the same way $F_0$ is the image of $\iota(\Omega_0)$. With (ii), (i) and (iv),
\[
 [\Phi(r_j):\Phi]=[F_1(t'):F_0(t')]=[F_1:F_0]=[\Omega:\iota(\Omega_0)]=d.\qedhere
\]
\end{proof}
\leanref{Development/AutomorphicSpectrum/Proofs/GeometryProofs/ArithVarietyReindexed.lean}{GlobalTheory.CompanionTheory.Ar.Chg.GenFib.exists_genFibre_chart_degree}
\leanref{GroundTheory/HeightsAndTheta/Heights/OrbitLimits/PGenFibreChartDeg.lean}{GroundTheory.Heights₂.Heights.Targets.GenFibreChartDeg.proof}
\leanref{GroundTheory/HeightsAndTheta/Heights/Intersections/PAdjoinIndepRelfinrank.lean}{GroundTheory.Heights₂.Heights.Inter.AdjoinIndepRelfinrank.proof}

\begin{lemma}\label{cv-origin-lift-finite}
In Situation~\ref{cv-origin-lift-situation}, let $P\subset R$ be an arithmetic subvariety of dimension $n+1$ such that
$P_0$ is an arithmetic subvariety of dimension $n+1$, and let a generic fibre of $P$ with a chart $j$ be given
(Definition~\ref{cv-origin-lift-fibre}). Then the chart field $\Phi(r_j)$ is finite over $\Phi$.
\end{lemma}
\leanref{GroundTheory/HeightsAndTheta/Heights/Intersections/ImportsMoreRest.lean}{GroundTheory.Heights₂.Heights.Inter.ArithmeticIntersectionChartFinite.Statement}
\leanref{GroundTheory/HeightsAndTheta/Heights/Intersections/StatementsBlockRatioAlg.lean}{GroundTheory.Heights₂.Heights.Inter.BlockRatioAlg.Statement}
\begin{proof}
Let $k$ be a fibre block, $q\in I_k$ and $c=(k,j_k)$; then $x_c\notin P$, as $x'_c\neq0$. We find multihomogeneous
$A_0,\dots,A_m\in R_0$ of one multidegree with $A_m\notin P_0$ and $g=\sum_{l\le m}A'_lx_q^lx_c^{m-l}\in P$. Then
$g(x')=0$ gives $\sum_lA_l(x)\,r_{j,q}^l=0$ with $A_m(x)\neq0$, so $r_{j,q}$ is algebraic over $\Phi$, and $\Phi(r_j)$,
generated by finitely many algebraic elements, is finite over $\Phi$.

\emph{A point with kernel $P$.} Choose $c_{k_1}\in I_{k_1}$ with $x_{k_1,c_{k_1}}\notin P$ for each $k_1\in K$. Let $\Omega$
be the field of the quotients $g/h$, in the fraction field of $R/P$, of multihomogeneous $g,h$ of one multidegree with
$h\notin P$, and let $y'_{k_1,i}$ be the class of $x_{k_1,i}/x_{k_1,c_{k_1}}$. Then $P$ is the kernel of $y'$: substituting
$T_{k_1}=x_{k_1,c_{k_1}}$ in the image of $f$ under $x_{k_1,i}\mapsto y'_{k_1,i}T_{k_1}$ gives the class of $f$, and a
multihomogeneous generator $g$ of $P$ of multidegree $a$ has $g(y')=g(x)/\prod_{k_1}x_{k_1,c_{k_1}}^{a_{k_1}}=0$. The field
$\Omega$ is its function field. The ring $R/P$ is the subring of $\Omega[T_{k_1}]$ generated by the $y'_{k_1,i}T_{k_1}$, so its
fraction field is $\Omega(T_{k_1}:k_1\in K)$, of transcendence degree $\operatorname{trdeg}\Omega+|K|$ over $\mathbf Q$.
As in the proof of Lemma~\ref{cv-origin-weight-birational}, the dimension formula \cite[Tag 02JU]{Sta}, with
\cite[Tags 00G4, 00GB and 00NM]{Sta}, shows that the dimension of a domain of finite type over $\mathbf Z$ that contains
$\mathbf Z$ is the transcendence degree of its fraction field plus one. As $\dim R/P=n+1+|K|$, $\Omega$ has transcendence
degree $n$.

\emph{Algebraicity.} The point $y=(y'_{k',i'})_{k\in K_0}$ has kernel $P_0$, by the computation of the first step of
the proof of Lemma~\ref{cv-origin-lift-degree} with $\iota$ the identity. Choose $c^0_{k_1}\in I^0_{k_1}$ with
$y_{k_1,c^0_{k_1}}\neq0$ for $k_1\in K_0$, and let $A\subset\Omega$ be the ring generated by the ratios
$y_{k_1,i}/y_{k_1,c^0_{k_1}}$. The homomorphism $R_0\to A[T_{k_1}:k_1\in K_0]$,
$x_{k_1,i}\mapsto(y_{k_1,i}/y_{k_1,c^0_{k_1}})T_{k_1}$, also has kernel $P_0$, so $R_0/P_0$ embeds in $A[T_{k_1}]$. By the
dimension formula its fraction field has transcendence degree $n+|K_0|$, so the fraction field of $A$ has transcendence
degree at least $n$, and $\Omega$, of transcendence degree $n$, is algebraic over $A$. So $y'_q/y'_c$ is a root of a
polynomial $\sum_{l\le m}a_lT^l$ with $a_l\in A$ and $a_m\neq0$. Write $a_l=F_l(y_{k_1,i}/y_{k_1,c^0_{k_1}})$ with
$F_l\in R_0$, and multiply each monomial of $F_l$ by powers of the $x_{k_1,c^0_{k_1}}$ to obtain multihomogeneous $A_l$ of
one multidegree, with $A_l(y)=\delta a_l$ for one nonzero $\delta\in\Omega$. Then $g=\sum_lA'_lx_q^lx_c^{m-l}$ is
multihomogeneous with $g(y')=\delta y'^m_c\sum_la_l(y'_q/y'_c)^l=0$, so $g\in P$; and $A_m(y)=\delta a_m\neq0$, so
$A_m\notin P_0$.
\end{proof}
\leanref{GroundTheory/HeightsAndTheta/Heights/Intersections/AevalNeZeroOfNotMem.lean}{GroundTheory.Heights₂.Heights.Inter.ArithmeticIntersectionChartFinite.proof}
\leanref{GroundTheory/HeightsAndTheta/Heights/Intersections/PBlockRatioAlg.lean}{GroundTheory.Heights₂.Heights.Inter.BlockRatioAlg.proof}

\begin{lemma}\label{cv-origin-lift-spec}
Let $A$ be a commutative ring, $\Phi$ a field of characteristic zero and $\varphi:A\to\Phi$ a homomorphism such that every
element of $\Phi$ is $\varphi(f)/\varphi(g)$ with $f,g\in A$ and $\varphi(g)\neq0$. Let $\Phi'\supset\Phi$ be a field, $Q$ a
finite set and $r\in\Phi'^Q$ with $\Phi(r)$ finite over $\Phi$, of degree $e$. Then there is $E\in A$ with
$\varphi(E)\neq0$ such that for every homomorphism $\psi:A\to\mathbf C$ with $\ker\varphi\subset\ker\psi$ and
$\psi(E)\neq0$, the set $Z_\psi$ of the $v\in\mathbf C^Q$ with $\psi(F)(v)=0$ for every $F\in A[X_q:q\in Q]$ with
$\varphi(F)(r)=0$ has $e$ elements, and $v_q\neq0$ for every $v\in Z_\psi$ and every $q$ with $r_q\neq0$. Here $\varphi(F)$
and $\psi(F)$ are the polynomials obtained by applying $\varphi$ and $\psi$ to the coefficients.
\end{lemma}
\leanref{GroundTheory/HeightsAndTheta/Heights/Intersections/StChartSpecCount.lean}{GroundTheory.Heights₂.Heights.Inter.ChartSpecCount.Statement}
\begin{proof}
\emph{A primitive element.} The field $\Phi(r)$ is generated by finitely many elements algebraic over $\Phi$, so it is the
$\Phi$-algebra they generate, and by the primitive element theorem it is $\Phi(\theta)$ with $\theta=\tau(r)$ for a
polynomial $\tau$ over $\Phi$. Let $m$ be the minimal polynomial of $\theta$, of degree $e$, and write $r_q=g_q(\theta)$
with $g_q\in\Phi[T]$. As $m$ is irreducible and $\Phi$ has characteristic zero, $um'\equiv1$ modulo $m$ for some $u$;
$\tau(g(T))\equiv T$ modulo $m$, as both sides take the value $\theta$ at $\theta$; and if $r_q\neq0$, then $m$ does not
divide $g_q$, so $h_qg_q\equiv1$ modulo $m$ for some $h_q$.

\emph{Denominators.} Each of the finitely many coefficients of $m$, $u$, $\tau$, the $g_q$ and the $h_q$ is
$\varphi(a)/\varphi(b)$ with $\varphi(b)\neq0$; let $E$ be the product of these $b$, and $\varphi_E:A[1/E]\to\Phi$ the
extension of $\varphi$. Then these coefficients lie in $\varphi_E(A[1/E])$, and we lift $m$, $u$, $\tau$, $g_q$ and $h_q$ to
polynomials $m_A$, monic of degree $e$, $u_A$, $\tau_A$, $g_{A,q}$ and $h_{A,q}$ over $A[1/E]$. Let $\psi$ be as in the
statement, and $\psi_E:A[1/E]\to\mathbf C$ its extension. If $\varphi_E(a/E^n)=0$, then $\varphi(a)=0$, so $\psi(a)=0$:
$\ker\varphi_E\subset\ker\psi_E$.

\emph{Two transfers.} (a) If $p\in A[1/E][T]$ has $\varphi_E(p)(\theta)=0$, then $\psi_E(m_A)$ divides $\psi_E(p)$: $m$
divides $\varphi_E(p)$, so the remainder of $p$ after division by the monic $m_A$ maps to $0$ under $\varphi_E$, hence
under $\psi_E$. (b) Let $v\in Z_\psi$, and let $F$ be a polynomial over $A[1/E]$ with $\varphi_E(F)(r)=0$. For large $n$,
$E^nF$ is the image of a polynomial $F'$ over $A$, with $\varphi(F')(r)=0$, so $\psi(F')(v)=0$ and, as $\psi(E)\neq0$,
$\psi_E(F)(v)=0$.

\emph{The count.} Let $T_\psi$ be the set of roots of $\psi_E(m_A)$ in $\mathbf C$, and
$G(t)=(\psi_E(g_{A,q})(t))_{q\in Q}$. For $v\in Z_\psi$, (b) applied to $m_A(\tau_A(X))$ and to the
$X_q-g_{A,q}(\tau_A(X))$, which vanish at $r$, shows that $t=\psi_E(\tau_A)(v)$ lies in $T_\psi$ and that $v=G(t)$.
Conversely, let $t\in T_\psi$ and let $F$ be a polynomial over $A$ with $\varphi(F)(r)=0$. The polynomial
$F(g_{A,q}(T))_q$ in $T$ vanishes at $\theta$, as $g_q(\theta)=r_q$, so by (a) $\psi(F)(G(t))=0$, and $G(t)\in Z_\psi$. By
(a) applied to $\tau_A(g_A(T))-T$, $\psi_E(\tau_A)(G(t))=t$ for $t\in T_\psi$, so $G$ is injective on $T_\psi$, and
$Z_\psi=G(T_\psi)$ has $|T_\psi|$ elements. By (a) applied to $u_Am_A'-1$, $\psi_E(m_A)$ divides
$\psi_E(u_A)\psi_E(m_A)'-1$; so it is separable, and, being monic of degree $e$, it has $e$ roots. Finally, if $r_q\neq0$
and $v=G(t)$ with $t\in T_\psi$, (a) applied to $h_{A,q}g_{A,q}-1$ gives $\psi_E(h_{A,q})(t)\,v_q=1$, so $v_q\neq0$.
\end{proof}
\leanref{GroundTheory/HeightsAndTheta/Heights/Intersections/PChartSpecCount.lean}{GroundTheory.Heights₂.Heights.Inter.ChartSpecCount.proof}
\leanref{GroundTheory/HeightsAndTheta/Heights/Intersections/PChartPrimElt.lean}{GroundTheory.Heights₂.Heights.Inter.ChartPrimElt.proof}
\leanref{GroundTheory/HeightsAndTheta/Heights/Intersections/PLocAwaySpec.lean}{GroundTheory.Heights₂.Heights.Inter.LocAwaySpec.proof}

\begin{lemma}\label{cv-origin-lift-chart}
In Situation~\ref{cv-origin-lift-situation}, let $P\subset R$ be an arithmetic subvariety of dimension $n+1$ such that
$P_0$ is an arithmetic subvariety of dimension $n+1$, and let a generic fibre of $P$ with a chart $j$ be given whose chart
field has degree $e$ over $\Phi$ (Definition~\ref{cv-origin-lift-fibre}). Then there is a multihomogeneous
$D\in R_0\setminus P_0$ such that for every $w\in\hat X_{P_0}(\mathbf C)$ with $D(w)\neq0$:
\begin{enumerate}
\item the points $z\in\hat X_P(\mathbf C)$ over $w$ with $z_{k,j_k}\neq0$ for every fibre block $k$ form $e$ orbits of
  $(\mathbf C^\times)^K$;
\item each such $z$ has $z_{k,i}\neq0$ for every fibre block $k$ and every $i$ with $x'_{k,i}\neq0$.
\end{enumerate}
\end{lemma}
\leanref{GroundTheory/HeightsAndTheta/Heights/Intersections/PtsOverChart.lean}{GroundTheory.Heights₂.Heights.Inter.GenericPointsInAffineChart.Statement}
\begin{proof}
By Lemma~\ref{cv-origin-lift-finite}, $\Phi(r_j)$ is finite over $\Phi$, and $\Phi$ has characteristic zero, as the kernel
$P_0$ of $R_0\to\Phi$ meets $\mathbf Z$ in $0$. Let $Q$ be the set of pairs $(k,i)$ with $k$ a fibre block and $i\in I_k$.

\emph{The normalized fibre.} Lemma~\ref{cv-origin-lift-spec}, applied to $A=R_0$, $\varphi:f\mapsto f(x)$ and $r=r_j$, gives
$E\in R_0$ with $E(x)\neq0$, that is, $E\notin P_0$. Let $w\in\hat X_{P_0}(\mathbf C)$ with $E(w)\neq0$, and let $\psi$ be
$f\mapsto f(w)$, so that $\ker\varphi=P_0\subset\ker\psi$. Let $N_w$ be the set of $z\in\hat X_P(\mathbf C)$ with $z|_0=w$
and $z_{k,j_k}=1$ for every fibre block $k$, and for $z\in\hat C$ with $z|_0=w$ let $z_Q=(z_{k,i})_{(k,i)\in Q}$. We show
that for $z\in\hat C$ with $z|_0=w$ and $z_{k,j_k}=1$ for the fibre blocks, $z\in\hat X_P(\mathbf C)$ if and only if
$z_Q\in Z_\psi$. Let $z\in\hat X_P(\mathbf C)$, and let $F\in R_0[X_q]$ with $\varphi(F)(r_j)=0$. Let $d_q$ be the
degree of $F$ in $X_q$, and let $F^h\in R$ be obtained from $F$ by replacing each coefficient $a$ by $a'$ and each
monomial $\prod_qX_q^{\mu_q}$ by $\prod_{q=(k,i)}x_{k,i}^{\mu_q}x_{k,j_k}^{d_q-\mu_q}$. Then
$F^h(x')=\prod_{q=(k,i)}x'^{d_q}_{k,j_k}\cdot\varphi(F)(r_j)=0$, so $F^h\in P$, and $\psi(F)(z_Q)=F^h(z)=0$, as
$z_{k,j_k}=1$. Conversely, let $z_Q\in Z_\psi$. The point $z$ lies in $\hat C$, as $w\in\hat C_0$ and $z_{k,j_k}=1$. Let $g$ be one of the
multihomogeneous elements, of multidegree $a$, that generate $P$, and let $F_g\in R_0[X_q]$ be obtained from $g$ by
replacing each $x_{k',i'}$ by the coefficient $x_{k,i}$ and each $x_{k,i}$ with $k$ a fibre block by $X_{k,i}$. Then
$\varphi(F_g)(r_j)=g(\gamma x')=\prod_k\gamma_k^{a_k}g(x')=0$, where $\gamma_k$ is $1$ on the blocks $k'$ and
$x'^{-1}_{k,j_k}$ on the fibre blocks; so $g(z)=\psi(F_g)(z_Q)=0$, and $z\in\hat X_P(\mathbf C)$. So $z\mapsto z_Q$ is
injective on $N_w$, and its image is $Z_\psi$: a point $v\in Z_\psi$ has $v_{k,j_k}=1$, by the relation $X_{k,j_k}-1$, and
it is $z_Q$ for the $z$ with $z|_0=w$ and fibre coordinates $v$. By Lemma~\ref{cv-origin-lift-spec}, $N_w$ has $e$
elements, and its points have $z_{k,i}\neq0$ whenever $x'_{k,i}\neq0$.

\emph{The form $D$.} The ideal $P_0$ is generated by multihomogeneous elements, so it contains the multihomogeneous
components of its elements, and $E$ has a multihomogeneous component $D\notin P_0$. Let $w\in\hat C_0$ with $D(w)\neq0$.
For $\gamma\in\mathbf C^{K_0}$, $E(\gamma w)=\sum_bE_b(w)\gamma^b$, where the $E_b$ are the components of $E$; this
polynomial in $\gamma$ has the coefficient $D(w)\neq0$, so it does not vanish on all of $(\mathbf C^\times)^{K_0}$, and
$E(\gamma w)\neq0$ for some $\gamma\in(\mathbf C^\times)^{K_0}$.

\emph{The count.} Let $w\in\hat X_{P_0}(\mathbf C)$ with $D(w)\neq0$, $\gamma$ as above and $w_1=\gamma w$, a point of
$\hat X_{P_0}(\mathbf C)$ with $E(w_1)\neq0$. Let $z\in\hat X_P(\mathbf C)$ lie over $w$ with $z_{k,j_k}\neq0$ for every
fibre block $k$, and $z|_0=sw$. Multiplying the block $k'$ of $z$ by $\gamma_k/s_k$ and each fibre block $k$ by
$z^{-1}_{k,j_k}$ gives a point of $N_{w_1}$ in the orbit of $z$, as $P$ is generated by multihomogeneous elements. Two
points of $N_{w_1}$ in one orbit are equal: the factor on a block $k'$ is $1$, as $w_1$ has a nonzero coordinate in the
block $k$, and on a fibre block it is $1$, as both points have the coordinate $1$ at $j_k$. So these orbits correspond to
the $e$ points of $N_{w_1}$, which gives (1), and (2) follows, as the points of an orbit differ by nonzero factors.
\end{proof}
\leanref{GroundTheory/HeightsAndTheta/Heights/Intersections/ExistsNormalise.lean}{GroundTheory.Heights₂.Heights.Inter.GenericPointsInAffineChart.proof}
\leanref{GroundTheory/HeightsAndTheta/Heights/Intersections/Imports.lean}{GroundTheory.Heights₂.Heights.Inter.ChartNormCount.proof}
\leanref{GroundTheory/HeightsAndTheta/Heights/Intersections/PChartNormPts.lean}{GroundTheory.Heights₂.Heights.Inter.ChartNormPts.proof}
\leanref{GroundTheory/HeightsAndTheta/Heights/Intersections/POrbitCompNonvanish.lean}{GroundTheory.Heights₂.Heights.Inter.OrbitCompNonvanish.proof}

\begin{lemma}\label{cv-origin-lift-orbits}
In Situation~\ref{cv-origin-lift-situation}, let $P\subset R$ be an arithmetic subvariety of dimension $n+1$ such that
$P_0$ is an arithmetic subvariety of dimension $n+1$, and let a generic fibre of $P$ with a chart $i$ be given whose chart
field has degree $e$ over $\Phi$ (Definition~\ref{cv-origin-lift-fibre}). Then there is a multihomogeneous
$D\in R_0\setminus P_0$ such that for every $w\in\hat X_{P_0}(\mathbf C)$ with $D(w)\neq0$ the points of
$\hat X_P(\mathbf C)$ over $w$ form $e$ orbits of $(\mathbf C^\times)^K$.
\end{lemma}
\leanref{GroundTheory/HeightsAndTheta/Heights/Intersections/PtsOverPt.lean}{GroundTheory.Heights₂.Heights.Inter.GenericPointsAvoidBadLocus.Statement}
\begin{proof}
For each chart $j$ of the generic fibre, Lemma~\ref{cv-origin-lift-chart} gives a form $D_j$; let $D$ be their product,
which is multihomogeneous and lies outside the prime $P_0$. The chart fields of all charts are equal: a ratio
$x'_{k,l}/x'_{k,j_k}$ is $r_{i,k,l}/r_{i,k,j_k}$, and conversely. So each $D_j$ comes with the same count $e$. Let
$w\in\hat X_{P_0}(\mathbf C)$ with $D(w)\neq0$. For a chart $j$, let $O_j$ be the set of orbits of points
$z\in\hat X_P(\mathbf C)$ over $w$ with $z_{k,j_k}\neq0$ for every fibre block $k$. By
Lemma~\ref{cv-origin-lift-chart}(2) for the chart $i$, such a point with $z_{k,i_k}\neq0$ has $z_{k,j_k}\neq0$, as
$x'_{k,j_k}\neq0$; so $O_i\subset O_j$, and as both have $e$ elements, $O_i=O_j$. Every point $z\in\hat X_P(\mathbf C)$ over
$w$ lies in some $O_j$: for each fibre block $k$ choose $l_k$ with $z_{k,l_k}\neq0$; then $x'_{k,l_k}\neq0$, as otherwise
$x_{k,l_k}$ would lie in $P$, the kernel of $f\mapsto f(x')$, and vanish at $z$; so $l$ is a chart, and the orbit of $z$
lies in $O_l$. So the orbits over $w$ are those of $O_i$, and there are $e$ of them.
\end{proof}
\leanref{GroundTheory/HeightsAndTheta/Heights/Intersections/IsMultiHomogProdInt.lean}{GroundTheory.Heights₂.Heights.Inter.GenericPointsAvoidBadLocus.proof}

\begin{lemma}[Serre {\cite[Proposition~5]{Ser56}}]\label{cv-origin-lift-dense}
Let $n\ge0$, let $P$ be a prime ideal of $\mathbf Q[x_1,\dots,x_n]$ and $f\notin P$. Then the set of common zeros
$z\in\mathbf C^n$ of $P$ with $f(z)\neq0$ is dense, in the complex topology, in the set $V$ of common zeros of $P$.
\end{lemma}
\leanref{GroundTheory/HeightsAndTheta/Heights/HilbertSamuel/StatementsRatPrimeDense.lean}{GroundTheory.Heights₂.Heights.HilbertSamuel.RatPrimeDense.Statement}

Serre's proposition says that on a complex algebraic variety a Zariski open and Zariski dense subset is dense in the
complex topology. Here $V$, with its reduced structure, is such a variety, and the set where $f\neq0$ is Zariski open in
it. It is Zariski dense because $f$ vanishes on no irreducible component of $V$: these components correspond to the
minimal primes over $P\mathbf C[x]$, which are associated primes of $\mathbf C[x]/P\mathbf C[x]=(\mathbf Q[x]/P)\otimes_{\mathbf Q}\mathbf C$,
and multiplication by $f$ is injective on this ring, since it is injective on the domain $\mathbf Q[x]/P$ and
$\mathbf C$ is flat over $\mathbf Q$.

\begin{lemma}\label{cv-origin-lift-nonvanishing}
Let $P\subset R=\mathbf Z[x_{k,i}]$ be an arithmetic subvariety (Definition~\ref{lit-multiprojective}) and $f\in R\setminus P$.
Then the points of $\hat X(\mathbf C)$ at which $f$ does not vanish are dense in $\hat X(\mathbf C)$.
\end{lemma}
\leanref{GroundTheory/HeightsAndTheta/Heights/HilbertSamuel/StatementsNonvanishingDense.lean}{GroundTheory.Heights₂.Heights.HilbertSamuel.NonvanishingDense.Statement}
\begin{proof}
As $P\cap\mathbf Z=0$, the ideal $P_{\mathbf Q}$ that $P$ generates in $\mathbf Q[x_{k,i}]$ is a prime with
$P_{\mathbf Q}\cap R=P$, the localization of $P$ at the nonzero integers. So $f\notin P_{\mathbf Q}$, and $P$ and
$P_{\mathbf Q}$ have the same common zeros in $\prod_k\mathbf C^{I_k}$. Let $z\in\hat X(\mathbf C)$. By
Lemma~\ref{cv-origin-lift-dense}, $z$ is a limit of common zeros $z_n$ of $P$ with $f(z_n)\neq0$. The set $\hat C$ is open
in $\prod_k\mathbf C^{I_k}$ and contains $z$, so the $z_n$ eventually lie in $\hat C$, hence in $\hat X(\mathbf C)$.
\end{proof}
\leanref{GroundTheory/HeightsAndTheta/Heights/HilbertSamuel/PNonvanishingDense.lean}{GroundTheory.Heights₂.Heights.HilbertSamuel.NonvanishingDense.proof}

\begin{lemma}\label{cv-origin-lift-interior}
Let $P\subset R=\mathbf Z[x_{k,i}]$ be an arithmetic subvariety of dimension $e+1$ (Definition~\ref{lit-multiprojective}),
let $J\subset R$ be an ideal generated by multihomogeneous elements and not contained in $P$, and let $\hat Z_J$ be the set
of points of $\hat C$ at which every element of $J$ vanishes. Let $\tilde X$ be a compact complex manifold of dimension $e$
with, for each $x\in\tilde X$, an open neighbourhood $U_x$ and a holomorphic map $\lambda_x:U_x\to\hat X(\mathbf C)$, such
that any two of these maps agree up to the action of $(\mathbf C^\times)^K$ where both are defined, every point of
$\hat X(\mathbf C)\setminus\hat Z_J$ lies in the orbit of some $\lambda_x(x)$, $x=x'$ whenever $\lambda_x(x)\notin\hat Z_J$
lies in the orbit of $\lambda_{x'}(x')$, and the set of $x$ with $\lambda_x(x)\in\hat Z_J$ has measure zero in every chart
of $\tilde X$. Let $G\in R\setminus P$ be multihomogeneous. Then the set $S$ of $x\in\tilde X$ with $G(\lambda_x(x))=0$ has
empty interior.
\end{lemma}
\leanref{GroundTheory/HeightsAndTheta/Heights/HilbertSamuel/ArithRes.lean}{GroundTheory.Heights₂.Heights.HilbertSamuel.AnRes}
\leanref{GroundTheory/HeightsAndTheta/Heights/HilbertSamuel/PAnResNullForm.lean}{GroundTheory.Heights₂.Heights.HilbertSamuel.AnResNullForm.interior_eq_empty}
\begin{proof}
Suppose that the interior $W$ of $S$ is nonempty.

\emph{A point of $W$ off $\hat Z_J$.} Some $x\in W$ has $\lambda_x(x)\notin\hat Z_J$: otherwise, for a chart of $\tilde X$
whose domain meets $W$, the image of $W$ in the chart, a nonempty open set, would lie in the image of the set of $x$ with
$\lambda_x(x)\in\hat Z_J$, which has measure zero. Put $z=\lambda_x(x)$.

\emph{Approximation.} Let $T$ be the set of $z'\in\hat X(\mathbf C)\setminus\hat Z_J$ with $G(z')\neq0$. By
Lemma~\ref{cv-origin-lift-nonvanishing}, $z$ is a limit of points of $\hat X(\mathbf C)$ at which $G$ does not vanish; as
$\hat Z_J$ is closed in $\hat C$ and $z\notin\hat Z_J$, $z$ lies in the closure of $T$. For each $z'\in T$ choose
$x(z')\in\tilde X$ with $\lambda_{x(z')}(x(z'))$ in the orbit of $z'$. As $G$ is multihomogeneous,
$G(\lambda_{x(z')}(x(z')))\neq0$, so $x(z')\notin S$.

\emph{Contradiction.} As $\tilde X$ is compact, the points $x(z')$, for $z'\in T$ tending to $z$, have a cluster point
$x''$. For $x(z')$ in $U_{x''}$, the point $\lambda_{x''}(x(z'))$ lies in the orbit of $z'$, so the pair
$(\lambda_{x''}(x(z')),z')$ satisfies $a_{k,i}b_{k,i'}=a_{k,i'}b_{k,i}$ for all $k$, $i$ and $i'$. This condition on pairs
$(a,b)$ is closed, and $\lambda_{x''}$ is continuous, so it holds for $(\lambda_{x''}(x''),z)$; as both points lie in
$\hat C$, they lie in one orbit. As $z=\lambda_x(x)\notin\hat Z_J$, $x''=x$. So $x$ is a cluster point of the $x(z')$, and
as $W$ is an open neighbourhood of $x$, some $x(z')$ lies in $W\subset S$, a contradiction.
\end{proof}

\begin{lemma}\label{cv-origin-model-lift}
Let $y_0\in\Omega_0^{\sqcup_{k\in K_0}I^0_k}$ and $y\in\Omega^{\sqcup_{k\in K}I_k}$ be points with values in
fields, let $k\mapsto k'$ be an injection from $K_0$ to $K$ with bijections $i\mapsto i'$ from $I^0_k$ to
$I_{k'}$, and let $X_0$ and $X$ be the kernels of $y_0$ and $y$.
\begin{enumerate}
\item Suppose that every multihomogeneous polynomial in the kernel of $y_0$ is carried by the renaming
  $x_{k,i}\mapsto x_{k',i'}$ into the kernel of $y$. Let $\tilde X_0$ with maps $\lambda^0_x$ be as in
  Definition~\ref{cv-origin-analytic}(3) for $y_0$, of dimension $D$, and let $\tilde X$ with maps $\lambda_x$ be
  as there for $y$, of some dimension, except that the bijectivity and immersion are only asked off the zero
  set of a multihomogeneous ideal. Then there is a holomorphic map $h:\tilde X\to\tilde X_0$ such that for every
  $x$, $\lambda^0_{h(x)}(h(x))$ and $(\lambda_x(x)_{k',i'})_{k,i}$ lie in one orbit of $(\mathbf C^\times)^{K_0}$.
\item Suppose moreover that $\Omega$ has characteristic zero, that $y_{k',i'}=\iota(y_{0,k,i})$ for a field
  embedding $\iota:\Omega_0\to\Omega$ with $[\Omega:\iota(\Omega_0)]=d\ge1$, that $X_0$ and $X$ are arithmetic
  subvarieties of dimension $D+1$ with function fields $\Omega_0$ and $\Omega$, and that $\tilde X$ is as in
  Definition~\ref{cv-origin-analytic}(3) for $y$, of dimension $D$. Then $h$ has degree $d$ off a closed set of
  measure zero, as in Definition~\ref{cv-proportionality-analytic-system}(3).
\item If $\pi:\tilde X_0\to S$ is a bijective holomorphic map to a complex manifold of dimension $D$, and
  $h:\tilde X\to\tilde X_0$ is holomorphic of degree $d$ off a closed set $Q$ of measure zero in this sense, then
  $\pi\circ h$ is holomorphic of degree $d$ off $\pi(Q)$.
\end{enumerate}
\end{lemma}
\leanref{Development/AutomorphicSpectrum/Proofs/MetricProofs/LogResolutionFunctorialHolomorphicStatement.lean}{GlobalTheory.T.log_resolution_functorial_holomorphic}
\leanref{Development/AutomorphicSpectrum/Proofs/GeometryProofs/LogResolutionReindexedStatement.lean}{GlobalTheory.T.log_resolution_reindexed}
\leanref{Development/AutomorphicSpectrum/Proofs/GeometryProofs/HolomorphicDegreeCompositionStatement.lean}{GlobalTheory.T.holomorphic_degree_composition}
\begin{proof}
Here a holomorphic map $h$ between compact complex manifolds of dimension $D$ has \emph{degree $d$ off} a set $Q$ if
$Q$ and $h^{-1}(Q)$ are closed sets with empty interior, each locally the common zeros of finitely many holomorphic
functions, every point outside $Q$ has $d$ preimages, and the derivative of $h$ is invertible outside $h^{-1}(Q)$.
For $z\in\hat C$ write $z|_0=(z_{k',i'})_{k,i}$, and say that $z\in\hat X(\mathbf C)$ lies \emph{over}
$w\in\hat X_0(\mathbf C)$ if $z|_0$ lies in the orbit of $w$ under $(\mathbf C^\times)^{K_0}$.

(1) For $x\in\tilde X$ and $x'\in U_x$ put $\mu_x(x')=\lambda_x(x')|_0$. It lies in $\hat X_0(\mathbf C)$: for $k\in K_0$
some coordinate of $\lambda_x(x')$ in the block $k'$ is nonzero, and an element of the kernel of $y_0$ vanishes at
$\mu_x(x')$ because its renaming lies in the kernel of $y$. Two maps $\mu_x$ and $\mu_{x''}$ agree up to the action of
$(\mathbf C^\times)^{K_0}$ where both are defined. As the maps $\lambda^0$ are bijective everywhere, there is a unique
$h(x)\in\tilde X_0$ with $\lambda^0_{h(x)}(h(x))$ in the orbit of $\mu_x(x)$.

The map $h$ is continuous. Let $x\in\tilde X$, and let $x_0$ be a limit point of $h$ along a filter converging to
$x$; it exists because $\tilde X_0$ is compact. Along the filter, $h(x')\in U_{x_0}$ eventually, and
$\lambda^0_{x_0}(h(x'))$ lies in the orbit of $\mu_x(x')$, so the cross products
$\lambda^0_{x_0}(h(x'))_{k,i}\,\mu_x(x')_{k,j}-\lambda^0_{x_0}(h(x'))_{k,j}\,\mu_x(x')_{k,i}$ vanish. In the limit
they vanish for $\lambda^0_{x_0}(x_0)$ and $\mu_x(x)$, so these points lie in one orbit, and $x_0=h(x)$ by injectivity.
As $\tilde X_0$ is compact and Hausdorff, $h$ is continuous at $x$.

The map $h$ is holomorphic. Let $x\in\tilde X$, put $x_0=h(x)$, choose $\kappa$ with
$\lambda^0_{x_0}(x_0)_{k,\kappa_k}\neq0$ for every $k$, and let $\phi$ be a chart of $\tilde X_0$ at $x_0$. Write $a_\kappa$
for the affine coordinates $z_{k,i}/z_{k,\kappa_k}$, which are constant on orbits. The derivative of
$a_\kappa\circ\lambda^0_{x_0}\circ\phi^{-1}$ at $\phi(x_0)$ is injective; composing with a linear left inverse of it and
applying the holomorphic inverse function theorem gives a holomorphic $G$ near $a_\kappa(\lambda^0_{x_0}(x_0))$ with
$G\circ a_\kappa\circ\lambda^0_{x_0}\circ\phi^{-1}$ the identity near $\phi(x_0)$. Near $x$, $h(x')\in U_{x_0}$ by
continuity and $a_\kappa(\lambda^0_{x_0}(h(x')))=a_\kappa(\mu_x(x'))$, so $\phi\circ h=G\circ a_\kappa\circ\mu_x$, which is
holomorphic.

(2) By (1) and the bijectivity of the maps $\lambda$, $\lambda_x(x)$ lies over $\lambda^0_{x_0}(x_0)$ if and only if
$x_0=h(x)$, and the points $x$ with $h(x)=x_0$ correspond to the orbits of points of $\hat X(\mathbf C)$ over
$\lambda^0_{x_0}(x_0)$.

First, local sections. Let $q$ and $r$ be indices in one block of $K$ with $y_r\neq0$. The ratio $\alpha=y_q/y_r$ is
algebraic over $\iota(\Omega_0)$, as $d$ is finite, and its minimal polynomial $m$ is separable, so $um+vm'=1$ for
polynomials $u$ and $v$ over $\Omega_0$. As $\Omega_0$ is the function field of $X_0$, the coefficients of $m$, $u$ and
$v$ are quotients $g(y_0)/H(y_0)$ with $g$ multihomogeneous of one multidegree $b$ and one $H$ of multidegree $b$,
$H(y_0)\neq0$. This gives polynomials $M$, $U$ and $V$ in $\tau$ whose coefficients are multihomogeneous of
multidegree $b$, with $M(y_0)=H(y_0)\,m$, $U(y_0)=H(y_0)\,u$ and $V(y_0)=H(y_0)\,v$. The coefficients of
$UM+VM'-H^2$ lie in the kernel of $y_0$, so at every $w\in\hat X_0(\mathbf C)$ with $H(w)\neq0$ the polynomial
$M_w(\tau)=\sum_jM_j(w)\tau^j$ has only simple roots. Homogenizing $\sum_jM_j(y_0)\alpha^j=0$ in the block of $q$ and $r$
gives an element of the kernel of $y$; as the $M_j$ have one multidegree, $M_w(z_q/z_r)=0$ whenever
$z\in\hat X(\mathbf C)$ with $z_r\neq0$ lies over $w$. By the holomorphic implicit function theorem, a simple root of
$M_w$ is a holomorphic function of $w$ near a given point. Let $G_2$ be the product of the forms $H$ over all pairs
$(q,r)$. Let $z_0$ lie over $w_0$ with $G_2(w_0)\neq0$, and choose for each block $k\in K$ an index $r_k$ with
$z_{0,r_k}\neq0$. The roots near $z_{0,q}/z_{0,r_k}$ give neighbourhoods $V$ of $w_0$ and $V'$ of $z_0$ and a holomorphic
map $f$ on $V$ such that every $z\in V'\cap\hat X(\mathbf C)$ over a point $w\in V\cap\hat X_0(\mathbf C)$ lies in the
orbit of $f(w)$.

Next, the count. In the notation of Situation~\ref{cv-origin-lift-situation}, Lemma~\ref{cv-origin-lift-degree} shows
that $X_0$ is the ideal of the $f\in R_0$ with $f'\in X$, and that $X$ has a generic fibre with a chart whose chart field
has degree $d$. So Lemma~\ref{cv-origin-lift-orbits}, with $n=D$, gives a multihomogeneous form $G_1$ outside $X_0$ such
that the points of $\hat X(\mathbf C)$ over each $w\in\hat X_0(\mathbf C)$ with $G_1(w)\neq0$ form $d$ orbits of
$(\mathbf C^\times)^K$.
% Lean: exists_count = arith_variety_reindexed (Lemma cv-origin-lift-degree, with imageZ_eq) and
% generic_points_good_statement = Ggeo.blockProjection_genericPoints_good = BaseCount.generic_points_good_for_intersection
% (Lemma cv-origin-lift-orbits).

Put $G=G_1G_2$, which lies outside the kernel of $y_0$, a prime, and let $Q$ be the set of $x_0\in\tilde X_0$ with
$G(\lambda^0_{x_0}(x_0))=0$. Near $x_0$ it is the zero set of the holomorphic function $G\circ\lambda^0_{x_0}$, and it has
empty interior by Lemma~\ref{cv-origin-lift-interior}, applied to $X_0$ and $\tilde X_0$ with $J=R_0$, whose set of
common zeros in $\hat C_0$ is empty. Its preimage $h^{-1}(Q)$ is the set of $x$ at which the renaming of $G$, a
multihomogeneous form outside the kernel of $y$, vanishes at $\lambda_x(x)$; so it is closed, locally the zero set of
one holomorphic function, with empty interior by the same lemma, applied to $X$ and $\tilde X$ with $J=R$.
% Lean: both empty interiors are HilbertSamuel.AnResNullForm.interior_eq_empty with NonvanishingDense.full, for
% R0.toAnRes and R.toAnRes with J = top, inside log_resolution_compatible_pullback. For $x_0\notin Q$ the points $x$ with $h(x)=x_0$ correspond to the $d$ orbits over
$\lambda^0_{x_0}(x_0)$. Finally let $x\notin h^{-1}(Q)$, and let $f$ be a local section as above at
$w_0=\lambda^0_{h(x)}(h(x))$ and $z_0=\lambda_x(x)$. Near $x$, $\lambda_x$ lies in the orbit of
$f\circ\lambda^0_{h(x)}\circ h$, so the affine coordinates of $\lambda_x$ in a chart at $x$ are a holomorphic function of
$\phi_0\circ h$, for a chart $\phi_0$ of $\tilde X_0$ at $h(x)$. They have injective derivative at $x$, so $h$ has
injective, hence invertible, derivative at $x$. So $h$ has degree $d$ off $Q$.

(3) In charts, $\pi$ is an injective holomorphic map between open subsets of $\mathbf C^D$, so its derivative is
invertible everywhere \cite[Theorem~1.6.6]{Leb26}. As $\tilde X_0$ is compact and $S$ Hausdorff, the bijection $\pi$ is a
homeomorphism, and its inverse $g$ is holomorphic by the inverse function theorem. Then $\pi(Q)=g^{-1}(Q)$ is closed,
locally the common zeros of the holomorphic functions of $Q$ composed with $g$, and it has empty interior because $g$ is
a homeomorphism; $(\pi\circ h)^{-1}(\pi(Q))=h^{-1}(Q)$; the fibre of $\pi\circ h$ over $x\notin\pi(Q)$ is the fibre of $h$
over $g(x)\notin Q$, with $d$ points; and outside $h^{-1}(Q)$ the derivative of $\pi\circ h$ is the composite of two
invertible derivatives.
\end{proof}
% Lean: (1) is PB.Proj.holoMap_of_blockProj = exists_holoMap_of_local_lifts with mu = R'.lift o v (comp_mem_mcPts):
% continuous_of_orbit (ultrafilters, orbit_cross and orbit_of_cross) and analytic_of_orbit (exists_analytic_retraction
% from R.immersion). (2) is Chg.DegBlock.isDegreeOff_of_block: exists_local_section (algebraic_element_homog_relation =
% AlgPoly.core with the minimal polynomial, homog_relation_ratio_reindex, implicit_root_analytic), exists_count, and
% log_resolution_compatible_pullback = CovMap.Deg.isDegreeOff_of_correspondence with pi = id (count through the image of
% the fibre, det_ne_zero_of_local_inverse with det_fderiv_ne_zero_of_factor). (3) is
% DegComp.holoMap_comp_isDegreeOff: Cj.exists_holoMap_inverse_of_degree_one (InjHol.bijective_fderiv_of_injOn, the
% Lean's proof of Osgood's theorem, through det_ne_zero_of_injective) and the transport of each condition.

\begin{lemma}\label{cv-origin-model-pullback}
Let $X$ and $X_0$ be schemes, $\Omega$ and $\Omega_0$ fields, $\xi:\operatorname{Spec}\Omega\to X$ and
$\xi_0:\operatorname{Spec}\Omega_0\to X_0$ morphisms, $\iota:\Omega_0\to\Omega$ a field embedding and
$h:X\to X_0$ a morphism with $h\circ\xi=\xi_0\circ\operatorname{Spec}\iota$. Suppose that the local ring of $X$ at
the image of $\xi$ maps bijectively onto $\Omega$ and that the local rings at its specializations map
injectively into $\Omega$.
\begin{enumerate}
\item Let $P$ be a point of $X$ that specializes the image of $\xi$, and $a\in\Omega_0$ regular at $h(P)$.
  Then $\iota(a)$ is regular at $P$; if $\iota(a)$ is a unit at $P$, then $a$ is a unit at $h(P)$; and if
  $P$ is the image of a complex point $z$, the value of $\iota(a)$ at $z$ is the value of $a$ at $h\circ z$.
\item Let $y_0\in\Omega_0^{\sqcup_{k\in K_0}I^0_k}$ and $y\in\Omega^{\sqcup_{k\in K}I_k}$ with an injection
  $k\mapsto k'$ and bijections $i\mapsto i'$ such that $y_{k',i'}=\iota(y_{0,k,i})$. Suppose that $X_0$ is
  covered by the affine open subschemes of the charts of $y_0$ as in
  Definition~\ref{cv-origin-hodge-component}(1), and that every point of $X$ specializes the image of $\xi$.
  If $w$ are coordinates for $y$ of a complex point $z$ of $X$, then $(w_{k',i'})_{k,i}$ are coordinates for
  $y_0$ of $h\circ z$.
\end{enumerate}
\end{lemma}
\leanref{Development/AutomorphicSpectrum/Proofs/GeometryProofs/DominatingRegularPushStatement.lean}{GlobalTheory.T.dominating_regular_push}
\leanref{Development/AutomorphicSpectrum/Proofs/GeometryProofs/DominatingCoordsReindexedStatement.lean}{GlobalTheory.T.dominating_coords_reindexed}
\leanref{Development/AutomorphicSpectrum/Proofs/GeometryProofs/DominatingRegularPush.lean}{GlobalTheory.CompanionTheory.Ar.Chg.Stalk.regular_along_dominant}
\leanref{Development/AutomorphicSpectrum/Proofs/GeometryProofs/DominatingCoordsReindexed.lean}{GlobalTheory.CompanionTheory.Ar.Chg.coordsAt_blkV}
\begin{proof}
(1) The morphism $h$ gives a local homomorphism $\mathcal O_{X_0,h(P)}\to\mathcal O_{X,P}$, compatible with
$\iota$ on the images in $\Omega_0$ and $\Omega$, since $h\circ\xi=\xi_0\circ\operatorname{Spec}\iota$ and the local
rings of $X$ embed in $\Omega$. So $\iota(a)$ is regular at $P$. If $a$ were not a unit at $h(P)$, it would lie
in the maximal ideal, and its image $\iota(a)$ in the maximal ideal at $P$. The value at $z$ of the image of a
germ at $h(P)$ is the value at $h\circ z$ of the germ.

(2) For $k\in K_0$ some $w_{k',c'}$ is nonzero. Let $c\in I^0_k$ with $w_{k',c'}\neq0$; then
$\iota(y_{0,k,c})=y_{k',c'}\neq0$, so $y_{0,k,c}\neq0$. The point $h(z)$ lies in a chart $\kappa$ of $y_0$, so
$y_{0,k,j}/y_{0,k,\kappa_k}$ is regular at $h(z)$ for every $j$, the quotient of two Segre ratios for charts
that differ only at $k$. By (1), $r=\iota(y_{0,k,c}/y_{0,k,\kappa_k})$ is regular at $z$ with value $v$ equal to
the value of $y_{0,k,c}/y_{0,k,\kappa_k}$ at $h(z)$, and $\iota(y_{0,k,\kappa_k}/y_{0,k,c})=y_{k',\kappa'_k}/y_{k',c'}$ is
regular at $z$ by the coordinates $w$, with a value $v'$. As $rr^{-1}=1$, $vv'=1$, so $v\neq0$ and
$y_{0,k,c}/y_{0,k,\kappa_k}$ is a unit at $h(z)$. Hence $y_{0,k,j}/y_{0,k,c}$ is regular at $h(z)$, and by (1)
its value at $h\circ z$ is the value at $z$ of $y_{k',j'}/y_{k',c'}$, which is $w_{k',j'}/w_{k',c'}$.
\end{proof}

\begin{lemma}\label{cv-origin-model-hodge}
In the setting of Definition~\ref{cv-origin-hodge-component}, let $Z_0$ be a Hodge component of $\mathcal X$ at
$S_0$ with coordinates $y_0$ with values in $\Omega_0$, and $Z$ a Hodge component at $S$ with coordinates $y$
with values in $\Omega$, let $\iota:\Omega_0\to\Omega$ be a field embedding and $h:Z\to Z_0$ a morphism that maps
the generic point of $Z$ to that of $Z_0$ compatibly with $\iota$, with $f=f_0\circ h$ and the embedding of $F_c$
in $\Omega$ equal to $\iota$ composed with that in $\Omega_0$. Let flags of $Z_0$ and $Z$ have bases
$b^0_{s,a}$ and $b_{s,a}$ with the same images in the Lie algebra of $\mathcal A$, after the base change by
$\iota$. Then:
\begin{enumerate}
\item there is an additive map $j:\mathrm{Lie}_0\to\mathrm{Lie}$ with $j(ae)=\iota(a)j(e)$, which carries linearly
  independent families to linearly independent families, with $j(b^0_{s,a})=b_{s,a}$, such that at every point
  $P$ of $Z$ over no prime $\mathcal V_s(P)$ is the $\mathcal O_{Z,P}$-module generated by $j(\mathcal V^0_s(h(P)))$,
  and such that at every complex point $z$ of $Z$ a family in $\mathrm{Lie}_0$ with Hodge Gram matrix $\Gamma$ at
  $h\circ z$ is carried by $j$ to a family with Gram matrix $\Gamma$ at $z$;
\item at every point $P$ of $Z$ over no prime, if $t_0$ generates the $a$-th graded piece of the flag of
  $\mathcal V^0_s$ at $h(P)$, then $\iota(t_0)$ generates that of $\mathcal V_s$ at $P$;
\item if $\Omega$ has characteristic zero, then at every complex point $z$ of $Z$, if $t_0$ generates the $a$-th
  graded piece of $\mathcal V^0_s$ at $h\circ z$ with Hodge norm $\nu$, then $\iota(t_0)$ generates that of
  $\mathcal V_s$ at $z$ with Hodge norm $\nu$;
\item suppose that $\Omega$ has characteristic zero, that $y$ contains $\iota(y_0)$ as some of its blocks, with
  indices $k\mapsto k'$ and $i\mapsto i'$, and that the root lines $\Lambda^0_{\tau,i}$ and $\Lambda_{\tau,i}$ of the
  two flags have multidegrees $a^0$ and $a$ with generators $e^0$ and $e$, for one $(\tau,i)$. Let $z$ be a
  complex point of $Z$ with coordinates $w$, and $c$ and $c_0$ charts of $y$ and $y_0$ with $w_{k,c_k}\neq0$ and
  $w_{k',c'_{0,k}}\neq0$ for all $k$. Then
  $\iota\big(e^0\prod_ky_{0,k,c_{0,k}}^{a^0_k}\big)\big/\big(e\prod_ky_{k,c_k}^{a_k}\big)$ is regular at $z$ with a
  nonzero value.
\end{enumerate}
\end{lemma}
\leanref{Development/AutomorphicSpectrum/Proofs/GeometryProofs/HodgeCompDescentFirstStatement.lean}{GlobalTheory.T.hodge_comp_descent_first}
\leanref{Development/AutomorphicSpectrum/Proofs/GeometryProofs/HodgeCompDescentSecondStatement.lean}{GlobalTheory.T.hodge_comp_descent_second}
\leanref{Development/AutomorphicSpectrum/Proofs/GeometryProofs/HodgeCompDescentThirdStatement.lean}{GlobalTheory.T.hodge_comp_descent_third}
\leanref{Development/AutomorphicSpectrum/Proofs/GeometryProofs/HodgeCompDescentFourthStatement.lean}{GlobalTheory.T.hodge_comp_descent_fourth}
\begin{proof}
Write $L(P)$ and $L_0(P_0)$ for the Lie algebras of $\mathcal A_Z$ and $\mathcal A_{Z_0}$ at points $P$ of $Z$ and $P_0$ of
$Z_0$, as in the proof of Lemma~\ref{cv-origin-change-comparison}. If $P$ lies over no prime, so does $h(P)$: every prime
is a unit at $P$, hence at $h(P)$, as $\mathcal O_{Z_0,h(P)}\to\mathcal O_{Z,P}$ is local.

(1) Lemma~\ref{cv-origin-change-comparison}(1), with $Z_0$ in place of $Z_1$, gives an additive map
$j:\mathrm{Lie}_0\to\mathrm{Lie}$ with $j(ae)=\iota(a)j(e)$, which carries linearly independent families to linearly
independent families, commutes with the actions of $\mathcal O_E$ and is compatible with the maps to the Lie algebra of
$\mathcal A$, with $L(P)$ the $\mathcal O_{Z,P}$-module generated by $j(L_0(h(P)))$ for every $P$. The map from
$\mathrm{Lie}$ to the Lie algebra of $\mathcal A$ is injective, as $\mathcal A_Z=f^*\mathcal A$; so $j(b^0_{s,a})=b_{s,a}$, the
two having the same image. As the embedding of $F_c$ in $\Omega$ is $\iota$ composed with that in $\Omega_0$, $j$ carries
$\mathrm{Lie}_{0,s}$ into $\mathrm{Lie}_s$. Let $P$ lie over no prime. Then the discriminant of $E$ is a unit at $P$ and at
$h(P)$, and the Lie algebras there are the direct sums of their parts in the eigenspaces, which are the $\mathcal V_s$:
$L(P)=\bigoplus_s\mathcal V_s(P)$ and $L_0(h(P))=\bigoplus_s\mathcal V^0_s(h(P))$. So $L(P)$ is the sum over $s$ of the
$\mathcal O_{Z,P}$-spans of $j(\mathcal V^0_s(h(P)))\subset\mathrm{Lie}_s$, and as the eigenspaces $\mathrm{Lie}_s$ are
independent, $\mathcal V_s(P)=L(P)\cap\mathrm{Lie}_s$ is the span of $j(\mathcal V^0_s(h(P)))$. Finally let $z$ be a complex
point of $Z$. The fibre of $\mathcal A_Z$ at $z$, with its cubical line bundle, is that of $\mathcal A_{Z_0}$ at $h\circ z$,
compatibly with the maps to $\mathcal A$, and the value at $z$ of $j(e)$ is the tangent vector given by the value of $e$
at $h\circ z$. The Hodge form is defined from the fibre and its line bundle, so the two Gram matrices agree.

(2) Let $P$ lie over no prime and let $t_0$ generate the $a$-th graded piece of $\mathcal V^0_s$ at $h(P)$. Let $b'$ be a
basis of $\mathcal V^0_s$ near $h(P)$ adapted to the flag of $b^0_s$. The graded piece is generated at $h(P)$ by the
coefficient $t_1\neq0$ of $b^0_{s,a}$ in $b'_a$, so $t_0=ut_1$ with $u$ a unit at $h(P)$. By (1), $\mathcal V_s(P)$ is spanned
by $j(b')$, which is linearly independent and so a basis, adapted to the flag of $b_s=j(b^0_s)$, and the coefficient of
$b_{s,a}$ in $j(b'_a)$ is $\iota(t_1)$. So $\iota(t_1)$ generates the $a$-th graded piece of $\mathcal V_s$ at $P$. By
Lemma~\ref{cv-origin-model-pullback}(1), $\iota(u)$ and $\iota(u^{-1})$ are regular at $P$, so $\iota(t_0)=\iota(u)\iota(t_1)$
generates it too.

(3) Let $z$ be a complex point of $Z$; its centre $P$ lies over no prime, as $\mathbf C$ has characteristic zero, and
$h(P)$ is the centre of $h\circ z$. The Hodge norm $\nu$ of $t_0$ at $h\circ z$ is
$(\det\Gamma_a/\det\Gamma_{a-1})^{1/2}$, for the Gram matrix $\Gamma$ at $h\circ z$ of a basis $b'$ of $\mathcal V^0_s$ at
$h(P)$ adapted to the flag of $b^0_s$ with $b'_a-t_0b^0_{s,a}$ in the $(a-1)$-th step of the flag. As in (2), $j(b')$ is a
basis of $\mathcal V_s$ at $P$ adapted to the flag of $b_s$, with $j(b'_a)-\iota(t_0)b_{s,a}$ in the $(a-1)$-th step, and
by (1) its Gram matrix at $z$ is $\Gamma$. So $\iota(t_0)$ generates the $a$-th graded piece of $\mathcal V_s$ at $z$ with
Hodge norm $\nu$.

(4) Let $P$ be the centre of $z$, which lies over no prime. By Lemma~\ref{cv-origin-model-pullback}(2), the coordinates of
$h\circ z$ for $y_0$ are $(w_{k',i'})_{k,i}$, so $h(P)$ lies in the open subscheme of the chart $c_0$, and $P$ in that of
$c$. There the root lines are generated by $g_0=e^0\prod_ky_{0,k,c_{0,k}}^{a^0_k}$ and $g=e\prod_ky_{k,c_k}^{a_k}$
(Definition~\ref{cv-origin-hodge-component}(3)). The root line $\Lambda^0_{\tau,i}$ is a graded piece of the flag of some
$\mathcal V^0_s$ or its dual; so $g_0=u_0t_0^{\pm1}$ with $t_0$ a generator of that graded piece at $h(P)$ and $u_0$ a unit
at $h(P)$. By (2), $\iota(t_0)$ generates the corresponding graded piece of $\mathcal V_s$ at $P$, so
$\iota(t_0)^{\pm1}=vg$ with $v$ a unit at $P$. Hence $\iota(g_0)/g=\iota(u_0)v$, and $\iota(u_0)$ is a unit at $P$ by
Lemma~\ref{cv-origin-model-pullback}(1). A unit at the centre of $z$ is regular at $z$ with a nonzero value.
\end{proof}
% Lean: (1) is ChgLie.lie_of_cover: exists_sabMap (A_Z = h^*A_{Z_0}), T.lie_model_dominant_hodge (the map j with the
% conclusions of Lemma cv-origin-change-comparison(1)), pt_inj_of_pull for j(b0) = b and the actions,
% T.hodge_cpt_away_from_primes (V = dR off the primes, and the lattice is the sum of the dR s) and
% T.hodge_cpt_interior_generic (independence of the eigenspaces, mem_of_mem_iSup), exists_value_up and
% T.hodge_inner_product_map_iff for the Gram matrices. (2) is GP.genPull_of_frames = isGenAt_pull (exists_coeff,
% gr_isGen_of_frame, gen_ratio, isFrameAt_pull, gen_mul). (3) is Chg.normPull_of_stalks (IsGrNorm: frame, Gram matrix and
% determinant ratio). (4) is ChgUnit.comparison_unit (inChart_of_coordsAt, exists_gen_of_tens, exists_unit_of_span_eq,
% tens_unit_eq, exists_value_of_unit).

\begin{lemma}\label{cv-origin-model-metric-norm}
In the setting of Definition~\ref{cv-origin-hodge-component}, let $Z$ be a Hodge component of $\mathcal X$ at $S$ with
coordinates $y$ with values in $\Omega$, and let a flag of $Z$ have bases $b_s$. Let $s$ be a label, $1\le m\le r_s$, and
$z$ a complex point of $Z$ over which $\mathcal A_Z$ is abelian, with centre $P$. For $t'\in\Omega$ and $\nu'\in\mathbf R$,
say that $t'$ \emph{generates the $m$-th graded piece of the flag of $\mathcal V_s$ at $z$ with Hodge norm $\nu'$} if
there is a family $b'=(b'_1,\dots,b'_{r_s})$ in $\mathrm{Lie}$, linearly independent over $\Omega$, with
$\mathcal V_s(P)=\sum_l\mathcal O_{Z,P}b'_l$, $\langle b'_1,\dots,b'_l\rangle=\langle b_{s,1},\dots,b_{s,l}\rangle$ for every
$l$ and $b'_m-t'b_{s,m}\in\langle b_{s,1},\dots,b_{s,m-1}\rangle$, whose Hodge Gram matrix $\Gamma$ at $z$
(Definition~\ref{cv-origin-hodge-norm-gram}(3)) is positive definite with $\nu'=(\det\Gamma_m/\det\Gamma_{m-1})^{1/2}$, where
$\Gamma_l$ is the Gram matrix of $b'_1,\dots,b'_l$. Let $t\in\Omega$ generate the $m$-th graded piece of the flag of
$\mathcal V_s$ at $P$. Then there is a real $\nu>0$ such that, for all $t'\in\Omega$ and $\nu'\in\mathbf R$, $t'$
generates the $m$-th graded piece of the flag of $\mathcal V_s$ at $z$ with Hodge norm $\nu'$ if and only if $t'/t$ is
regular at $z$ with a nonzero value $u$ and $\nu'=|u|\nu$.
\end{lemma}
\leanref{Development/AutomorphicSpectrum/Proofs/GeometryProofs/HodgeCptInteriorFlagStatement.lean}{GlobalTheory.T.hodge_cpt_interior_flag}
\leanref{Development/AutomorphicSpectrum/Proofs/GeometryProofs/HodgeCptInteriorFlag.lean}{GlobalTheory.CompanionTheory.Ar.GrNorm.exists_norm}
\begin{proof}
Let $\mathcal O_P\subset\Omega$ be the ring of elements regular at $P$, and $\psi:\mathcal O_P\to\mathbf C$ the value at $z$.
A prime number $\ell$ has the nonzero value $\ell$ at $z$, so it is a unit at $P$, and $P$ lies over no prime. An element of
$\mathcal O_P$ with a nonzero value at $z$ is a unit of $\mathcal O_P$, as the homomorphism from the local ring at $P$ to
$\mathbf C$ is local.

\emph{Values and Gram matrices.} Let $L(P)$ be the Lie algebra of $\mathcal A_Z$ at $P$. By
Lemma~\ref{cv-origin-hodge-norm-eigen}(2), $\mathcal V_{s'}(P)=L(P)\cap\mathrm{Lie}_{s'}$ for every label $s'$ and
$L(P)=\sum_{s'}\mathcal V_{s'}(P)$. The group scheme $\mathcal A_Z$ is smooth over $Z$, so
Lemma~\ref{cv-origin-hodge-norm-values} gives the values $e(z)$, in the Lie algebra $T_z$ of $\mathcal A_Z$ at $z$, of the
$e\in L(P)$, and Lemma~\ref{cv-origin-hodge-norm-form} gives a positive definite Hermitian form $h$ on $T_z$ whose values are
the Hodge inner products. By Lemma~\ref{cv-origin-hodge-norm-values}(1), the values at $z$ of a family $b'$ in $L(P)$ are the
$b'_p(z)$, so $b'$ has Hodge Gram matrix $\Gamma$ at $z$ if and only if $\Gamma_{pq}=h(b'_p(z),b'_q(z))$ for all $p$ and $q$.

\emph{Independence of the values.} Let $b'$ be linearly independent over $\Omega$ with
$\mathcal V_s(P)=\sum_k\mathcal O_Pb'_k$. The $b'_k$ lie in $\mathrm{Lie}_s$, so by Lemma~\ref{cv-origin-hodge-norm-eigen}(1)
there are $\Omega$-linear forms $\pi_1,\dots,\pi_{r_s}$ on $\mathrm{Lie}$ that vanish on $\mathrm{Lie}_{s'}$ for $s'\neq s$, with
$\pi_j(b'_k)=\delta_{jk}$. They map $\mathcal V_{s'}(P)$ to $0$ for $s'\neq s$ and $\mathcal V_s(P)=\sum_k\mathcal O_Pb'_k$ into
$\mathcal O_P$, hence $L(P)$ into $\mathcal O_P$. By Lemma~\ref{cv-origin-hodge-norm-values}(4) there are $\mathbf C$-linear
forms $g_j$ on $T_z$ with $g_j(e(z))=\psi(\pi_j(e))$ for $e\in L(P)$. Then $g_j(b'_k(z))=\delta_{jk}$, and the $b'_k(z)$ are
linearly independent.

\emph{A basis through $t$.} By Definition~\ref{cv-origin-hodge-component}(2), $\mathcal V_s$ has near $P$ a basis $b$ with
$\langle b_1,\dots,b_l\rangle=\langle b_{s,1},\dots,b_{s,l}\rangle$ for every $l$ and $\mathcal V_s(P)=\sum_k\mathcal O_Pb_k$.
Let $a\in\Omega$ with $b_m-ab_{s,m}\in\langle b_{s,1},\dots,b_{s,m-1}\rangle$. By the first step of the proof of
Lemma~\ref{cv-origin-reference-descent-roots}, $a$ generates the $m$-th graded piece of the flag of $\mathcal V_s$ at $P$, as
$t$ does. So $t=da$ and $a=d't$ with $d,d'\in\mathcal O_P$, and $dd'=1$ as $t\neq0$. Replacing $b_m$ by $db_m$ keeps the
properties of $b$ and gives $b_m-tb_{s,m}\in\langle b_{s,1},\dots,b_{s,m-1}\rangle$.

\emph{The norm.} Apply Lemma~\ref{cv-origin-hodge-norm-adapted} to $A=\mathcal O_P$, $\varphi=\psi$, $W=\mathrm{Lie}$,
$L=L(P)$, $V=T_z$, $\mathrm{val}(e)=e(z)$, the form $h$, $Q=\mathcal V_s(P)$, $b^0=b_s$, the basis $b$, $m$ and $t$. Its
hypotheses hold: $\mathrm{val}$ is additive with $\mathrm{val}(ae)=\psi(a)\mathrm{val}(e)$ by
Lemma~\ref{cv-origin-hodge-norm-values}(2), the values of an adapted basis of $Q$ are linearly independent by the paragraph
on independence, and $b$ is an adapted basis by the previous paragraph. It gives $\nu>0$. Its condition (a) is, by the
paragraph on Gram matrices, the condition that $t'$ generates the $m$-th graded piece of the flag of $\mathcal V_s$ at $z$
with Hodge norm $\nu'$, and its condition (b) says that $t'/t$ is regular at $z$ with a nonzero value $u$ and $\nu'=|u|\nu$.
\end{proof}
% Lean: T.hodge_cpt_interior_flag = GrNorm.exists_norm, for a model over any O with O -> F_c injective: not_inFibre;
% T.hodge_cpt_interior_generic (hind) and T.hodge_cpt_away_from_primes (hV, hlat), Lemma cv-origin-hodge-norm-eigen;
% lieModel_exists and T.lie_lattice_eval_at_local_ring (ev, val, hsmul, hvalAt), Lemma cv-origin-hodge-norm-values;
% T.hodge_hermitian_form_exists with interior_pol (h, hHI, hh, hp) and hGram, Lemma cv-origin-hodge-norm-form; hfr from
% exists_dual and T.lattice_regular_at_local_point; Fl.sub, FlagGrGen.gr_isGenAt and the unit d (hdu) with scale_frame; and
% core, Lemma cv-origin-hodge-norm-adapted. Nothing below remains open. Under the Hermitian form of Lemma
% cv-origin-hodge-norm-form, Lemmas cv-origin-hodge-norm-form-descent (the field descent of relative ampleness by sections),
% cv-origin-hodge-norm-form-descent-ample (an ample line bundle on a scheme proper over a ring is relatively ample by
% sections, and all its nonvanishing loci are affine: relAmple_of_isAmpleLB and IsAmpleLB.isAffineOpen_nonzeroLocus, about
% 1,600 lines, with Lemmas -descent-ample-affine and -descent-ample-embedding) and cv-origin-hodge-norm-form-torus (the
% compact topological group A(C), the holomorphic chart, the doubling logarithm of the group law, and the divisibility of
% A(C)) are written out, as is the rest.

\begin{lemma}\label{cv-origin-model-metric-chart}
Let $X$ be a scheme, $\Omega$ a field, $\xi:\operatorname{Spec}\Omega\to X$ a morphism, $y\in\Omega^{\sqcup_{k\in K}I_k}$,
and $z$ a complex point of $X$ whose centre $P$ specializes the image of $\xi$, with coordinates $w$ for $y$ as in
Definition~\ref{cv-origin-reference}(4); here $g\in\Omega$ is \emph{regular at $z$ with value $r$} if it is the image in
$\Omega$ of an element of the local ring of $X$ at $P$ whose value at $z$ is $r$. Let $\bar A=(a,h)$ be a metrized line
bundle (Definition~\ref{lit-multiprojective}(2)), $e\in\Omega$ nonzero, $t\in\Omega$ and $\nu\in\mathbf R$, and let
$c=(c_k)$ and $c'=(c'_k)$ be choices of indices $c_k,c'_k\in I_k$ with $w_{k,c_k}\neq0$ and $w_{k,c'_k}\neq0$ for all $k$.
Suppose that $t/(e\prod_ky_{k,c_k}^{a_k})$ is regular at $z$ with a value $u$ such that
$\nu h(w)=|u|\prod_k|w_{k,c_k}|^{a_k}$. Then $t/(e\prod_ky_{k,c'_k}^{a_k})$ is regular at $z$ with a value $u'$ such that
$\nu h(w)=|u'|\prod_k|w_{k,c'_k}|^{a_k}$.
\end{lemma}
\leanref{Development/AutomorphicSpectrum/Definitions/GlobalGeometry/HodgeIPIffHodgeFormConj.lean}{GlobalTheory.CompanionTheory.Ar.secNorm_iff_chart}
\begin{proof}
Let $k\in K$. As $w_{k,c'_k}\neq0$, the coordinates give $y_{k,c'_k}\neq0$ and show that $y_{k,c_k}/y_{k,c'_k}$ is
regular at $z$ with value $r_k=w_{k,c_k}/w_{k,c'_k}$, which is nonzero. Values at $z$ are multiplicative. An element
regular at $z$ with a nonzero value is the image of a unit of the local ring at $P$, as the homomorphism from that ring to
$\mathbf C$ is local; so its inverse is regular at $z$ with the inverse value. Hence $\prod_k(y_{k,c_k}/y_{k,c'_k})^{a_k}$
is regular at $z$ with value $\prod_kr_k^{a_k}$, and
\[
\frac{t}{e\prod_ky_{k,c'_k}^{a_k}}=\frac{t}{e\prod_ky_{k,c_k}^{a_k}}\prod_k\Big(\frac{y_{k,c_k}}{y_{k,c'_k}}\Big)^{a_k}
\]
is regular at $z$ with value $u'=u\prod_kr_k^{a_k}$. Then
$|u'|\prod_k|w_{k,c'_k}|^{a_k}=|u|\prod_k|r_kw_{k,c'_k}|^{a_k}=|u|\prod_k|w_{k,c_k}|^{a_k}=\nu h(w)$.
\end{proof}
% Lean: secNorm_iff_chart: the r_k from CoordsAt at c', CxVal.zpow (through CxVal.inv and CxVal.isUnit) and CxVal.mul for
% the product, and norm_mul, norm_zpow for the identity. Size: about 45 lines.

\begin{lemma}\label{cv-origin-model-metric}
Let $Z_0$, $Z$, $\iota$, $h$ and the two flags be as in Lemma~\ref{cv-origin-model-hodge}, without the
conditions on $f$, on the embeddings of $F_c$ and on the bases, with $\Omega$ of characteristic zero and $y$
containing $\iota(y_0)$ as blocks. Suppose that the conclusions (2) and (3) of Lemma~\ref{cv-origin-model-hodge}
and those of Lemma~\ref{cv-origin-model-pullback} hold, that $h$ carries the complex points of $Z$ over which
$\mathcal A_Z$ is abelian to points over which $\mathcal A_{Z_0}$ is abelian, and let $(\tau,i)\in\Psi\times\{1,\dots,N\}$,
$a^0\in\mathbf Z^{K_0}$, $a\in\mathbf Z^K$, $e^0\in\Omega_0$, and $e\in\Omega$ nonzero. Let $\bar\sigma_0=(a^0,h_0)$ be a
Hodge metric on $\Lambda^0_{\tau,i}$ with generator $e^0$, and let $\bar T=(a-a^0,h_T)$ be a metrized line bundle, where
$a^0$ is extended by $0$ to the blocks not of the form $k'$, such that for $z$, $w$, $c$, $c_0$ as in
Lemma~\ref{cv-origin-model-hodge}(4), the element of (4) is regular at $z$ with a nonzero value $u$ and
$h_T(w)\prod_k|w_{k',c'_{0,k}}|^{a^0_k}=|u|\prod_k|w_{k,c_k}|^{a_k}$. Then the tensor
product of $\bar T$ with the pullback $(w\mapsto h_0((w_{k',i'})_{k,i}))$ of $\bar\sigma_0$ is a Hodge metric on
$\Lambda_{\tau,i}$ with generator $e$.
\end{lemma}
\leanref{Development/AutomorphicSpectrum/Proofs/GeometryProofs/HodgeCompDescentFifthStatement.lean}{GlobalTheory.T.hodge_comp_descent_fifth}
\begin{proof}
Let $\sigma=1$ and $(s,m)=(\tau,i)$ if $i\le p$, and $\sigma=-1$ and $(s,m)=(\bar\tau,N+1-i)$ if $i>p$. By
Definition~\ref{cv-origin-hodge-component}(3), $\Lambda_{\tau,i}$ and $\Lambda^0_{\tau,i}$ are the $m$-th graded pieces of
the flags of $\mathcal V_s$ and $\mathcal V^0_s$ if $\sigma=1$, and their duals if $\sigma=-1$. For $w\in\hat C$ put
$w_0=(w_{k',i'})_{k,i}$. The tensor product $\bar\sigma$ of $\bar T$ with the pullback of $\bar\sigma_0$ has multidegree
$a^0+(a-a^0)=a$ and function $H(w)=h_0(w_0)h_T(w)$. With the Hodge norms of
Lemma~\ref{cv-origin-model-metric-norm}, Definition~\ref{cv-origin-hodge-component}(4) asks the following of $\bar\sigma$:
for every complex point $z$ of $Z$ over which $\mathcal A_Z$ is abelian, with coordinates $w$, every $t\in\Omega$ that
generates the $m$-th graded piece of the flag of $\mathcal V_s$ at $z$ with Hodge norm $\nu$, and every chart $c$ of $y$
with $w_{k,c_k}\neq0$ for all $k$, the element $t^\sigma/(e\prod_ky_{k,c_k}^{a_k})$ is regular at $z$ with a value $U$
such that $\nu^\sigma H(w)=|U|\prod_k|w_{k,c_k}|^{a_k}$. As $t^\sigma$ has norm $\nu^\sigma$, this says that the generator
$e\prod_ky_{k,c_k}^{a_k}$ has norm $\prod_k|w_{k,c_k}|^{a_k}/H(w)$. The hypothesis on $\bar\sigma_0$ is the same statement
on $Z_0$. Fix $z$, $w$, $t$ and $\nu$, and let $P$ be the centre of $z$.

\emph{The point $h\circ z$.} By Lemma~\ref{cv-origin-model-pullback}(2), $w_0$ are coordinates of $h\circ z$ for $y_0$, and
by hypothesis $\mathcal A_{Z_0}$ is abelian over $h\circ z$, whose centre is $h(P)$. By
Definition~\ref{cv-origin-hodge-component}(2), $\mathcal V^0_s$ has near $h(P)$ a basis adapted to the flag of $b^0_s$, and
the coefficient $t_0$ of $b^0_{s,m}$ in its $m$-th vector generates the $m$-th graded piece of the flag of $\mathcal V^0_s$
at $h(P)$ (the first step of the proof of Lemma~\ref{cv-origin-reference-descent-roots}).
Lemma~\ref{cv-origin-model-metric-norm}, for $Z_0$, $h\circ z$ and $t_0$, gives $\nu_0>0$ such that $t'\in\Omega_0$ generates
the $m$-th graded piece of the flag of $\mathcal V^0_s$ at $h\circ z$ with Hodge norm $\nu'$ if and only if $t'/t_0$ is
regular at $h\circ z$ with a nonzero value $u$ and $\nu'=|u|\nu_0$. With $t'=t_0$ and $u=1$, $t_0$ generates it with Hodge
norm $\nu_0$.

\emph{The point $z$.} The centre $P$ lies over no prime: a prime number $\ell$ has the nonzero value $\ell$ at $z$, so it is
a unit at $P$. By conclusion (2) of Lemma~\ref{cv-origin-model-hodge}, $\iota(t_0)$ generates the $m$-th graded piece of the
flag of $\mathcal V_s$ at $P$; in particular $\iota(t_0)\neq0$. Lemma~\ref{cv-origin-model-metric-norm}, for $Z$, $z$ and
$\iota(t_0)$, gives $\nu_1>0$ with the same property on $Z$. By conclusion (3) of Lemma~\ref{cv-origin-model-hodge},
$\iota(t_0)$ generates the $m$-th graded piece at $z$ with Hodge norm $\nu_0$. So $\iota(t_0)/\iota(t_0)=1$ has a value
$u'$ at $z$ with $\nu_0=|u'|\nu_1$; the value of $1$ is $1$, so $\nu_1=\nu_0$. As $t$ generates the $m$-th graded piece at
$z$ with Hodge norm $\nu$, $t/\iota(t_0)$ is regular at $z$ with a nonzero value $u_1$, and $\nu=|u_1|\nu_0$.

\emph{Charts.} The coordinates $w$ have a nonzero entry in each block. Choose a chart $c$ of $y$ with $w_{k,c_k}\neq0$ for
all $k$, and a chart $c_0$ of $y_0$ with $w_{k',c'_{0,k}}\neq0$ for all $k$. As $\bar\sigma_0$ is a Hodge metric on
$\Lambda^0_{\tau,i}$ with generator $e^0$, and $t_0$ generates the $m$-th graded piece at $h\circ z$ with Hodge norm
$\nu_0$, the element $t_0^\sigma/(e^0\prod_ky_{0,k,c_{0,k}}^{a^0_k})$ is regular at $h\circ z$ with a value $u_0$ such that
$\nu_0^\sigma h_0(w_0)=|u_0|\prod_k|w_{k',c'_{0,k}}|^{a^0_k}$. By the hypothesis on $\bar T$, the element
$g=\iota(e^0\prod_ky_{0,k,c_{0,k}}^{a^0_k})/(e\prod_ky_{k,c_k}^{a_k})$ is regular at $z$ with a nonzero value $u$, and
$h_T(w)\prod_k|w_{k',c'_{0,k}}|^{a^0_k}=|u|\prod_k|w_{k,c_k}|^{a_k}$.

\emph{The norm.} Now
\[
\frac{t^\sigma}{e\prod_ky_{k,c_k}^{a_k}}=\Big(\frac{t}{\iota(t_0)}\Big)^\sigma\,
\iota\Big(\frac{t_0^\sigma}{e^0\prod_ky_{0,k,c_{0,k}}^{a^0_k}}\Big)\,g.
\]
The first factor is regular at $z$ with value $u_1^\sigma$, as $u_1\neq0$ makes $t/\iota(t_0)$ a unit at $P$; the second
has the value $u_0$ at $z$ by Lemma~\ref{cv-origin-model-pullback}(1); the third has the value $u$. So the left side is
regular at $z$ with value $U=u_1^\sigma u_0u$, and
\[
\nu^\sigma H(w)=|u_1|^\sigma\nu_0^\sigma h_0(w_0)h_T(w)=|u_1|^\sigma|u_0|\,h_T(w)\prod_k|w_{k',c'_{0,k}}|^{a^0_k}
=|u_1|^\sigma|u_0||u|\prod_k|w_{k,c_k}|^{a_k}=|U|\prod_k|w_{k,c_k}|^{a_k}.
\]
This is the required identity in the chart $c$. As $e\neq0$, Lemma~\ref{cv-origin-model-metric-chart} gives it in every
chart $c'$ of $y$ with $w_{k,c'_k}\neq0$ for all $k$.
\end{proof}
% Lean: GlobalTheory.T.hodge_comp_descent_fifth = Chg.isHodgeLine_pull_tensor: hCoord and hInt for z ≫ h,
% Flag.gr_gen_of_isFlagVBOn (gr_gen, about 110 lines) for t_0, T.hodge_cpt_interior_flag twice (Lemma
% cv-origin-model-metric-norm), Chg.not_inFibre and pushCpt for the centre, hGen and hNorm (conclusions (2) and (3) of
% Lemma cv-origin-model-hodge), CxVal.unique with CxVal.one for nu_1 = nu_0, hH at z ≫ h and the chart c0, hR (IsRatio for
% T) at c and c0, CxVal.zpow, hVal and CxVal.mul for the value u_1^sgn u_0 u, pullDeg a0 + T.a = a, and secNorm_iff_chart
% (Lemma cv-origin-model-metric-chart). Size: about 14,400 lines beyond the definitions in the statement, of which Lemma
% cv-origin-model-metric-norm is about 14,200.

\begin{lemma}\label{cv-origin-correction-extend}
Let $\Omega$ be a field and $y\in\Omega^{\sqcup_{k\in K}I_k}$, let $\tilde X$ with maps $\lambda_x$ be as in
Definition~\ref{cv-origin-analytic}(3) for $y$, and let $u:\tilde X\to\mathbf R$ be a smooth function. Then there is a
smooth metrized line bundle $\bar B=(0,h_B)$ (Definition~\ref{lit-multiprojective}(2)) with
$\log h_B(\lambda_x(x))=u(x)$ for every $x\in\tilde X$.
\end{lemma}
\leanref{Development/AutomorphicSpectrum/Proofs/MetricProofs/SmoothFunctionAsDegreeZeroMetricStatement.lean}{GlobalTheory.T.smooth_function_as_degree_zero_metric}
\begin{proof}
Let $F=\mathbf C^{\sqcup_kI_k\times I_k}$, a finite-dimensional real vector space. For $z\in\hat C$ put
$\|z_k\|^2=\sum_{i\in I_k}|z_{k,i}|^2$, which is positive for every $k$, and let $p(z)\in F$ have the entries
$p(z)_{k,(i,j)}=z_{k,i}\bar z_{k,j}/\|z_k\|^2$. Three facts about $p$ are used.
\begin{enumerate}
\item[(a)] The map $p$ is smooth on $\hat C$, and $p(cz)=p(z)$ for $c\in(\mathbf C^\times)^K$ acting blockwise, as
  $|c_k|^2=c_k\bar c_k$.
\item[(b)] If $z,w\in\hat C$ and $p(w)=p(z)$, then $z=cw$ for some $c\in(\mathbf C^\times)^K$. Indeed, put
  $c_k=\sum_j\bar w_{k,j}z_{k,j}/\|w_k\|^2$. Then $\sum_jp(w)_{k,(i,j)}z_{k,j}=c_kw_{k,i}$, while
  $\sum_jp(z)_{k,(i,j)}z_{k,j}=z_{k,i}$; so $z_{k,i}=c_kw_{k,i}$, and $c_k\neq0$ as $z_k\neq0$.
\item[(c)] Let $\kappa$ be a chart with $z_{k,\kappa_k}\neq0$ for every $k$. Then
  $p(z)_{k,(\kappa_k,\kappa_k)}=|z_{k,\kappa_k}|^2/\|z_k\|^2\neq0$ and
  $p(z)_{k,(i,\kappa_k)}/p(z)_{k,(\kappa_k,\kappa_k)}=z_{k,i}/z_{k,\kappa_k}$. So the map
  $q_\kappa(M)=(M_{k,(i,\kappa_k)}/M_{k,(\kappa_k,\kappa_k)})_{k,i}$, which is smooth on the open set of the $M\in F$
  with every $M_{k,(\kappa_k,\kappa_k)}\neq0$, sends $p(z)$ to the affine coordinates $(z_{k,i}/z_{k,\kappa_k})$ of $z$.
\end{enumerate}

\emph{The map $g$.} Let $g:\tilde X\to F$ be $g(x)=p(\lambda_x(x))$. For $x_0\in\tilde X$ and $y\in U_{x_0}$, the points
$\lambda_y(y)$ and $\lambda_{x_0}(y)$ lie in one orbit, so $g(y)=p(\lambda_{x_0}(y))$ by (a), and this is continuous in $y$.
So $g$ is continuous. If $g(x)=g(x')$, then $\lambda_x(x)$ and $\lambda_{x'}(x')$ lie in one orbit by (b), so $x=x'$:
$g$ is injective.

\emph{Local extensions.} Let $x_0\in\tilde X$. There is a function $f$ on $F$, smooth on a neighbourhood of $g(x_0)$,
with $f(g(y))=u(y)$ for $y$ near $x_0$. To see this, choose $\kappa$ with $\lambda_{x_0}(x_0)_{k,\kappa_k}\neq0$ for
every $k$, and a chart $\phi$ of $\tilde X$ at $x_0$. The affine coordinates $a(t)$ of $\lambda_{x_0}(\phi^{-1}(t))$ in
the chart $\kappa$ are holomorphic in $t$ near $\phi(x_0)$, and their derivative at $\phi(x_0)$ is injective by
Definition~\ref{cv-origin-analytic}(3). Composing with a linear left inverse of this derivative and applying the
holomorphic inverse function theorem gives a holomorphic $r$ near $a(\phi(x_0))$ with $r(a(t))=t$ for $t$ near
$\phi(x_0)$. Put $f=u\circ\phi^{-1}\circ r\circ q_\kappa$. By (c), $q_\kappa(g(x_0))=a(\phi(x_0))$ and $q_\kappa$ is
smooth near $g(x_0)$, so $f$ is smooth near $g(x_0)$. For $y$ near $x_0$, $\lambda_{x_0}(y)_{k,\kappa_k}\neq0$ and
$g(y)=p(\lambda_{x_0}(y))$, so $q_\kappa(g(y))=a(\phi(y))$ by (c), and $f(g(y))=u(y)$.

\emph{A global extension.} For each $w\in F$ choose an open neighbourhood $W_w$ and a smooth function $f_w$ on $W_w$
with $f_w(g(x))=u(x)$ whenever $g(x)\in W_w$. If $w=g(x_0)$, take $f$ as in the previous paragraph and let $N$ be the
interior of the set of $x$ with $f(g(x))=u(x)$, a neighbourhood of $x_0$. The set $g(\tilde X\setminus N)$ is compact,
as $\tilde X$ is, hence closed, and it does not contain $g(x_0)$, as $g$ is injective. Let $W_w$ be the open set of
points near $g(x_0)$ at which $f$ is smooth, with $g(\tilde X\setminus N)$ removed, and $f_w=f$: if $g(x)\in W_w$ then
$x\in N$ and $f(g(x))=u(x)$. If $w\notin g(\tilde X)$, take $W_w=F\setminus g(\tilde X)$, which is open as
$g(\tilde X)$ is compact, and $f_w=0$. Let $(\rho_m)$ be a smooth partition of unity on $F$ with the support of each
$\rho_m$ in some $W_{w_m}$, and put $v=\sum_m\rho_mf_{w_m}$. Then $v$ is smooth on $F$, and $v(g(x))=u(x)$ for every
$x$, as each $f_{w_m}$ with $\rho_m(g(x))\neq0$ takes the value $u(x)$ at $g(x)$.

\emph{The bundle.} Put $h_B(z)=\exp(v(p(z)))$ for $z\in\hat C$. It is positive and, by (a), smooth and invariant under
$(\mathbf C^\times)^K$, so $\bar B=(0,h_B)$ is a smooth metrized line bundle of multidegree $0$, and
$\log h_B(\lambda_x(x))=v(g(x))=u(x)$.
\end{proof}
% Lean: PB.exists_hermLine_log_eq_of_smooth, for an AnLogRes R with J = top. p is HermExtend.proj (contDiffOn_proj,
% proj_smul, eq_smul_of_proj_eq through secAt_proj and secAt_proj_self), q_kappa is SmoothExt.deproj (deproj_proj,
% contDiffAt_deproj); continuous_diag and injective_diag (R.compat, R.inj); local_ext (R.holo, R.immersion,
% exists_localLeftInverse with Mathlib's ContDiffAt.localInverse); exists_ext_of_local, through Mathlib's
% exists_contMDiffMap_forall_mem_convex_of_local (a smooth partition of unity) with the convex sets of the values u(x)
% over a point. Size: about 230 lines.

\begin{lemma}\label{cv-origin-model-correction}
Let $\iota:\Omega_0\to\Omega$ be an embedding of fields, let $y_0\in\Omega_0^{\sqcup_{k\in K_0}I^0_k}$ and
$y\in\Omega^{\sqcup_{k\in K}I_k}$ with an injection $k\mapsto k'$ and bijections $i\mapsto i'$ such that
$y_{k',i'}=\iota(y_{0,k,i})$, and let $X$ be the kernel of $y$.
\begin{enumerate}
\item Let $Z$ be a scheme with a morphism $\operatorname{Spec}\Omega\to Z$ of whose image every point is a
  specialization, covered by the affine open subschemes of the charts of $y$ as in
  Definition~\ref{cv-origin-hodge-component}(1), and let $\tilde X$ with maps $\lambda_x$ be as in
  Definition~\ref{cv-origin-analytic}(3) for $y$. Let $a\in\mathbf Z^K$, $a^0\in\mathbf Z^{K_0}$, $e\in\Omega$ and
  $e^0\in\Omega_0$ be such that for every complex point $z$ of $Z$ with coordinates $w$, and charts $c$ of $y$ and
  $c_0$ of $y_0$ with $w_{k,c_k}\neq0$ and $w_{k',c'_{0,k}}\neq0$, the element
  $\iota\big(e^0\prod_ky_{0,k,c_{0,k}}^{a^0_k}\big)\big/\big(e\prod_ky_{k,c_k}^{a_k}\big)$ is regular at $z$ with
  a nonzero value $u$. Then there is a smooth metrized line bundle $\bar T=(a-a^0,h_T)$, invariant under
  conjugation, whose local potentials on $\tilde X$ through the maps $\lambda_x$ are pluriharmonic, with
  $h_T(w)\prod_k|w_{k',c'_{0,k}}|^{a^0_k}=|u|\prod_k|w_{k,c_k}|^{a_k}$ for all such $z$, $w$, $c$ and $c_0$.
\item Let $\tilde X_0$ with maps $\lambda^0_x$ be as in Definition~\ref{cv-origin-analytic}(3) for $y_0$ and $\tilde X$
  as there for $y$, of the same dimension, let $\pi:\tilde X_0\to S$ be holomorphic to a complex manifold, and let
  $h:\tilde X\to\tilde X_0$ be holomorphic with $\lambda^0_{h(x)}(h(x))$ and $(\lambda_x(x)_{k',i'})$ in one orbit. Let
  $\bar A_0$ be a metrized line bundle for $y_0$ whose local potentials on $\tilde X_0$ differ by pluriharmonic
  functions from the pullback by $\pi$ of a family $\Theta$ of local potentials on $S$, and $\bar T$ a metrized line
  bundle for $y$ with pluriharmonic local potentials on $\tilde X$. Then the local potentials of the tensor
  product of $\bar T$ with the pullback of $\bar A_0$ differ by pluriharmonic functions from the pullback of
  $\Theta$ by $\pi\circ h$.
\item If metrized line bundles $\bar T_{\tau,i}$ for $y$ have pluriharmonic local potentials on $\tilde X$, so do the
  bundles $-\sum_{k\le i}\bar T_{\tau,k}$.
\end{enumerate}
\end{lemma}
\leanref{Development/AutomorphicSpectrum/Proofs/GeometryProofs/DegreeReindexedCoordsStatement.lean}{GlobalTheory.T.degree_reindexed_coords}
\leanref{Development/AutomorphicSpectrum/Proofs/GeometryProofs/PositiveHomogeneousPotentialStatement.lean}{GlobalTheory.T.positive_homogeneous_potential}
\leanref{Development/AutomorphicSpectrum/Proofs/GeometryProofs/LogResolutionPotentialPullbackStatement.lean}{GlobalTheory.T.log_resolution_potential_pullback}
\leanref{Development/AutomorphicSpectrum/Proofs/GeometryProofs/FlatOnLogResolutionCorrStatement.lean}{GlobalTheory.T.flat_on_log_resolution_corr}
\begin{proof}
(1) Write $r_{c,c_0}=\iota\big(e^0\prod_ky_{0,k,c_{0,k}}^{a^0_k}\big)\big/\big(e\prod_ky_{k,c_k}^{a_k}\big)$.

\emph{Points and values.} A point $w\in\hat C$ lies in $\hat X(\mathbf C)$ if and only if it is the coordinate vector of a
complex point of $Z$. If $w\in\hat X(\mathbf C)$, choose $\kappa$ with all $w_{k,\kappa_k}\neq0$; then $\kappa$ is a chart
of $y$, as $x_{k,i}$ lies in the kernel when $y_{k,i}=0$, and Lemma~\ref{cv-origin-boundary-points}(2) applies.
Conversely, let $\zeta$ be a complex point with image $P$ and coordinates $w$, and $\kappa$ a chart with all
$w_{k,\kappa_k}\neq0$; put $t_{k,i}=y_{k,i}/y_{k,\kappa_k}$, regular at $P$ with values $w_{k,i}/w_{k,\kappa_k}$. For $g$ in
the kernel, $g(t_{k,i}T_k)=0$ in $\Omega[T_k]$, hence in $\mathcal O_{Z,P}[T_k]$; mapping to $\mathbf C$ and putting
$T_k=w_{k,\kappa_k}$ gives $g(w)=0$. Next, if $\zeta$ has coordinates $w$ and $f\in\Omega$ is regular at $P$, then
$f=g(y)/h(y)$ for multihomogeneous $g$ and $h$ of one multidegree $d$ with $h(w)\neq0$, and for every such pair the value of
$f$ at $\zeta$ is $g(w)/h(w)$. Indeed, by Lemma~\ref{cv-origin-snc-charts}(1), $f=r/s$ with $r,s\in A^{\mathbf Z}_\kappa$ and $s$
a unit at $P$, and writing $r,s$ as polynomials in the $t_{k,i}$ and multiplying by $\prod_ky_{k,\kappa_k}^{d_k}$ gives such
$g,h$, with $h(w)\prod_kw_{k,\kappa_k}^{-d_k}$ the value of $s$, which is nonzero. For any such pair,
$h(t)=h(y)\prod_ky_{k,\kappa_k}^{-d_k}$ has the nonzero value $h(w)\prod_kw_{k,\kappa_k}^{-d_k}$, so it is a unit at $P$, and
$f=g(t)/h(t)$ has value $g(w)/h(w)$.

\emph{The function $\varphi$.} For $w\in\hat X(\mathbf C)$ let $\zeta$ be a complex point with coordinates $w$, and $c$ and
$c_0$ charts with all $w_{k,c_k}\neq0$ and $w_{k',c'_{0,k}}\neq0$. By hypothesis $r_{c,c_0}$ has a nonzero value $u$ at
$\zeta$; put
\[
 \psi(w)=u\prod_kw_{k,c_k}^{a_k}\Big/\prod_kw_{k',c'_{0,k}}^{a^0_k},\qquad\varphi(w)=|\psi(w)|.
\]
By the previous paragraph $u$ depends only on $w$, and $\psi(w)$ does not depend on $c$ and $c_0$: for other charts $\tilde c$
and $\tilde c_0$, $r_{\tilde c,\tilde c_0}=r_{c,c_0}\prod_k(y_{k',\tilde c'_{0,k}}/y_{k',c'_{0,k}})^{a^0_k}\prod_k(y_{k,c_k}/y_{k,\tilde c_k})^{a_k}$,
as $y_{k',i'}=\iota(y_{0,k,i})$, and these ratios have the corresponding ratios of $w$ as values. The function $\varphi$ is
positive. It is invariant under conjugation: writing $r_{c,c_0}=g(y)/h(y)$ as above,
$\psi(w)=g(w)h(w)^{-1}\prod_kw_{k,c_k}^{a_k}/\prod_kw_{k',c'_{0,k}}^{a^0_k}$ with integer polynomials, and $\bar w\in\hat X(\mathbf C)$.
For $s\in(\mathbf C^\times)^K$, $(s_kw_{k,i})$ are coordinates of the same complex point, so
$\varphi(s\cdot w)=\prod_k|s_k|^{(a-a^0)_k}\varphi(w)$, with $a^0$ extended by $0$. Near $w$ it is the absolute value of a
holomorphic function: on the open set $V$ of the $x$ with $h(x)\neq0$, all $x_{k,c_k}\neq0$ and all $x_{k',c'_{0,k}}\neq0$,
put $F(x)=g(x)h(x)^{-1}\prod_kx_{k,c_k}^{a_k}/\prod_kx_{k',c'_{0,k}}^{a^0_k}$; for $x\in V\cap\hat X(\mathbf C)$ the same $g$ and
$h$ compute the value of $r_{c,c_0}$, so $\psi(x)=F(x)\neq0$. Finally $\varphi(w)\prod_k|w_{k',c'_{0,k}}|^{a^0_k}=|u|\prod_k|w_{k,c_k}|^{a_k}$
by definition.

\emph{The bundle.} Let $\bar H$ be the smooth metrized line bundle of multidegree $a-a^0$ with
$h_H(z)=\prod_k\big(\sum_i|z_{k,i}|^2\big)^{(a-a^0)_k/2}$. The function $v(x)=\log\big(\varphi(\lambda_x(x))/h_H(\lambda_x(x))\big)$ on
$\tilde X$ does not change when $\lambda_x(x)$ is replaced by a point of its orbit, as $\varphi$ and $h_H$ have the same
homogeneity, so near $x$ it is $y\mapsto\log|F(\lambda_x(y))|-\log h_H(\lambda_x(y))$, which is smooth.
By Lemma~\ref{cv-origin-correction-extend} there is a smooth metrized line bundle $\bar B$ of multidegree $0$ with
$\log h_B(\lambda_x(x))=v(x)$ for every $x\in\tilde X$.
% Lean: GlobalTheory.T.smooth_function_as_degree_zero_metric, for an AnLogRes with J = top; the input is
% FlatExt.smooth_ratio.
Then $\bar T_0=\bar H\otimes\bar B$ has $h_{T_0}(\lambda_x(x))=\varphi(\lambda_x(x))$ for every $x$, and so $h_{T_0}=\varphi$ on
$\hat X(\mathbf C)$, as every point of $\hat X(\mathbf C)$ lies in the orbit of some $\lambda_x(x)$ and both functions have the
same homogeneity. Put $h_T(z)=\big(h_{T_0}(z)h_{T_0}(\bar z)\big)^{1/2}$. Then $\bar T=(a-a^0,h_T)$ is smooth, invariant under
conjugation and homogeneous of multidegree $a-a^0$, as $|\bar s_k|=|s_k|$, and $h_T=\varphi$ on $\hat X(\mathbf C)$, as
$\varphi$ is invariant under conjugation; this gives the identity of (1). Its local potentials are pluriharmonic: near
$x\in\tilde X$, $\log h_T(\lambda_x(y))=\log|F(\lambda_x(y))|$ is the real part of a branch of $\log F(\lambda_x(y))$, a
holomorphic function of $y$.

(2) Let $x\in\tilde X$. By hypothesis there are a neighbourhood $V_0$ of $h(x)$ and a holomorphic $G_0$ on it with
$\log h_{A_0}(\lambda^0_{h(x)}(x_0))-\Theta_{\pi(h(x))}(\pi(x_0))=\operatorname{Re}G_0(x_0)$ for $x_0\in V_0$, and a neighbourhood
$V_1$ of $x$ and a holomorphic $G_1$ with $\log h_T(\lambda_x(y))=\operatorname{Re}G_1(y)$ for $y\in V_1$. For $y$ near $x$ the
points $(\lambda_x(y)_{k',i'})_{k,i}$ and $\lambda^0_{h(x)}(h(y))$ lie in one orbit, by the compatibility of the lifts on
$\tilde X$ and on $\tilde X_0$ and the hypothesis at $y$: $\lambda_x(y)_{k',i'}=c_k(y)\lambda^0_{h(x)}(h(y))_{k,i}$ with
$c(y)\in(\mathbf C^\times)^{K_0}$. For each $k$ fix $i$ with $\lambda^0_{h(x)}(h(x))_{k,i}\neq0$; then near $x$, $c_k(y)$ is a
quotient of two nonvanishing holomorphic functions of $y$, so $\log|c_k|=\operatorname{Re}L_k$ for a holomorphic $L_k$. The
function of the tensor product at $\lambda_x(y)$ is
$h_{A_0}\big((\lambda_x(y)_{k',i'})\big)h_T(\lambda_x(y))=\prod_k|c_k(y)|^{a^0_k}h_{A_0}\big(\lambda^0_{h(x)}(h(y))\big)h_T(\lambda_x(y))$,
so its logarithm minus $\Theta_{\pi(h(x))}(\pi(h(y)))$ is $\operatorname{Re}\big(G_0(h(y))+G_1(y)+\sum_ka^0_kL_k(y)\big)$ near $x$.

(3) Near each point the logarithm of the function of each $\bar T_{\tau,k}$, composed with the lifts, is the real
part of a holomorphic function, and the logarithm of the function of $-\sum_{k\le i}\bar T_{\tau,k}$ is minus the sum
of these, which is the real part of a holomorphic function.
% Lean: (1) is degree_reindexed_coords = Chg.Ratio.exists_isRatio (complex_points_iff_centres, valAt_iff_value_at, cmpVal,
% cmpVal_eq, isCmp_unique, IsCmp.conj, IsCmp.smul, the local analytic F on V) and positive_homogeneous_potential =
% Chg.FlatExt.exists_flat_line (fsHerm, smooth_ratio, smooth_function_as_degree_zero_metric, the tensor product,
% conjugation_invariant_hermitian_metric = PB.exists_conjInvariant_line_averaged, isFlatOn_of_local); (2) is
% log_resolution_potential_pullback = Chg.Twist.samePot_twist (orbit_key, exists_coeff_log, log_tensor_pull, log_h_smul);
% (3) is flat_on_log_resolution_corr = Chg.isFlatOn_corr, with corr T none the trivial line.
\end{proof}

\begin{lemma}\label{cv-origin-model-section}
In the setting of Definition~\ref{cv-origin-hodge-component}, let a flag model $R$ with data at one component
as in Definition~\ref{cv-origin-tower}(2b), with point $\xi_c$ with values in $\Omega_c$, elements $y^c_l$, integers
$q_c$, $n^c_{\tau,i}$ and element $s_c$, and data as in (2c) for $R$ and a set $S_0$ in place of $S_B$ be given,
with Hodge component $Z_0$, field embedding $g:\Omega_c\to\Omega_0$ and flag, and let $y_0$, with values in
$\Omega_0$, be the point of $R$. Let $Z$ be a Hodge component of $\mathcal X$ at $S$ with coordinates $y$ with
values in $\Omega$, containing $\iota(y_0)$ as blocks for a field embedding $\iota:\Omega_0\to\Omega$, with a
morphism $h:Z\to Z_0$ that maps the generic point to the generic point compatibly with $\iota$, and a flag with
root lines of some multidegrees with generators. Suppose that the conclusion (2) of
Lemma~\ref{cv-origin-model-hodge} holds for the two flags, that the conclusion (1) of
Lemma~\ref{cv-origin-model-pullback} holds for $h$, and that $\mathcal A_Z$ is abelian at a point $P$ if and only
if $\mathcal A_{Z_0}$ is abelian at $h(P)$. Then property (2c)(iii) of Definition~\ref{cv-origin-tower} holds for
$Z$, its flag, the elements $\iota(g(y^c_l))$ in place of the $y_{k_0,l'}$, and $\iota\circ g$ in place of $g$.
\end{lemma}
\leanref{Development/AutomorphicSpectrum/Proofs/GeometryProofs/RootOriginCoordFieldStatement.lean}{GlobalTheory.T.root_origin_coord_field}
\leanref{Development/AutomorphicSpectrum/Proofs/GeometryProofs/RootOriginCoordField.lean}{GlobalTheory.CompanionTheory.Ar.Chg.CoverEff.section_on_cover}
\begin{proof}
Let $P$ be a point of $Z$ over no prime, let $i$ be an index with $\iota(g(y^c_i))\neq0$ and every
$\iota(g(y^c_l))/\iota(g(y^c_i))$ regular at $P$, and let $t$ generate
$\bigotimes_{\tau,l}\Lambda_{\tau,l}^{-n^c_{\tau,l}}$ at $P$. Put $P_0=h(P)$. A prime $r$ is a unit at $P$, so it is
a unit at $P_0$ by Lemma~\ref{cv-origin-model-pullback}(1), and $P_0$ lies over no prime. The point $P_0$ lies in
a chart $c_0$ of $y_0$, so the ratios $g(y^c_l)/g(y^c_{l_0})$ are regular at $P_0$, where $l_0'=c_{0,k_0}$, as
$g(y^c_l)=y_{0,k_0,l'}$. The ratio $g(y^c_i)/g(y^c_{l_0})$ is regular at $P_0$ and its image under $\iota$ is a unit at
$P$, so it is a unit at $P_0$, and every $g(y^c_l)/g(y^c_i)$ is regular at $P_0$.

The root line $\Lambda^0_{\tau,l}$ of the flag of $Z_0$ has a generator at $P_0$, and it is the graded piece of
$\mathcal V^0_s$ that defines it, or its dual, so that graded piece has a generator $\gamma_{\tau,l}$ at $P_0$, and
$t_0=\prod_{\tau,l}\gamma_{\tau,l}^{\mp n^c_{\tau,l}}$, with the sign of Definition~\ref{cv-origin-hodge-component}(3),
generates $\bigotimes\Lambda^{0,-n^c_{\tau,l}}_{\tau,l}$ at $P_0$. By the conclusion (2) of
Lemma~\ref{cv-origin-model-hodge}, each $\iota(\gamma_{\tau,l})$ generates the corresponding graded piece of
$\mathcal V_s$ at $P$, so $\iota(t_0)$ generates $\bigotimes\Lambda_{\tau,l}^{-n^c_{\tau,l}}$ at $P$, and $\iota(t_0)/t$ is a
unit at $P$. By (2c)(iii) for the data of $R$ at $P_0$, $\sigma_0=g(s_c)g(y^c_i)^{q_c}/t_0$ is regular at $P_0$, and it is
not a unit if $\mathcal A_{Z_0}$ is not abelian at $P_0$. Now
$\iota(g(s_c))\,\iota(g(y^c_i))^{q_c}/t=\iota(\sigma_0)\cdot\iota(t_0)/t$, which is regular at $P$. If $\mathcal A_Z$ is not
abelian at $P$, then $\mathcal A_{Z_0}$ is not abelian at $P_0$; if the product were a unit at $P$, $\iota(\sigma_0)$
would be a unit at $P$, and $\sigma_0$ a unit at $P_0$ by Lemma~\ref{cv-origin-model-pullback}(1), which it is not.
\end{proof}

\begin{lemma}\label{cv-origin-model}
Let the analytic system $S$, the data of (2b) and the flag model $R$ with data at the empty set of degree $1$
be as in Lemma~\ref{cv-origin-change}, and let $S'$ be a finite set of primes. Suppose that the map
$\tilde X\to X_R(\mathbf C)$ of these data is bijective and immersive everywhere and that $\pi:\tilde X\to S$ is
bijective. Let $\Omega'$, $\iota$, $y'$, $Z'$, $h$ and a flag of $Z'$ with root lines of multidegrees $a'_{\tau,i}$
with generators $e'_{\tau,i}$ satisfy (1), (2), (3) and (5) of Lemma~\ref{cv-origin-change}, and let $Z'$ satisfy
condition (1) of Lemma~\ref{cv-origin-change-projective} at every $p\in S'$. Then there are a flag model $M$ of
type $(\Psi,N)$ and data as in Definition~\ref{cv-origin-tower}(2c) for $M$ and $S'$, with the same data of
(2b), whose Hodge component satisfies condition (1) of Lemma~\ref{cv-origin-change-projective} at every
$p\in S'$ and has the integers $R_p$ of $Z'$, whose map $\tilde X'\to X_M(\mathbf C)$ is bijective and immersive
everywhere, and such that $\mathbf Q[y_{k,i}/y_{k,\kappa_k}]$ is smooth over $\mathbf Q$ for the point $y$ of $M$
and every chart $\kappa$.
\end{lemma}
\leanref{Development/AutomorphicSpectrum/Proofs/GeometryProofs/RootOriginCoordsOverFieldStatement.lean}{GlobalTheory.T.root_origin_coords_over_field}
\leanref{Development/AutomorphicSpectrum/Proofs/GeometryProofs/RootOriginCoordsOverField.lean}{GlobalTheory.CompanionTheory.Ar.root_model_on_cover}
\begin{proof}
Write $\Omega$, $y$, $Z$, $\tilde X$, $\lambda_x$, $\pi$, $\bar\sigma_{\tau,i}$ and $\bar y_b$ for the field, point, Hodge
component, manifold, lifts, map, Hodge metrics and bundles of the data of $R$, and $X$ and $X'$ for the kernels of
$y$ and $y'$; put $d'=[\Omega':\iota(\Omega)]$, the degree of the flag of $Z'$, and $D=|\Psi|N(N-1)/2$. By
Lemma~\ref{cv-origin-weight-birational}, $\Omega'$ is the function field of $X'$, an arithmetic subvariety of
dimension $D+1$. By (5) and Lemma~\ref{cv-origin-model-resolution} there is a manifold $\tilde X'$ with lifts
$\lambda'_x$ as in Definition~\ref{cv-origin-analytic}(3) for $y'$. The kernel of $y$ is carried into that of $y'$ by
the renaming of the blocks, since $y'$ contains $\iota(y)$. Lemma~\ref{cv-origin-model-lift}(1) gives a
holomorphic map $h':\tilde X'\to\tilde X$ compatible with the lifts up to orbits, (2) shows that $h'$ has degree
$d'$ off a closed set of measure zero, and (3) that $\pi'=\pi\circ h'$ is holomorphic of degree $d'$ off a closed
set of measure zero.

Lemma~\ref{cv-origin-model-pullback}, applied to $h:Z'\to Z$, transfers regularity and values of functions and
coordinates of complex points from $Z'$ to $Z$. Lemma~\ref{cv-origin-model-hodge}, applied to $Z$, $Z'$, $h$ and
the two flags, gives its conclusions (1) to (4), using (1) and (2) of Lemma~\ref{cv-origin-change}. For each
$(\tau,i)$, Lemma~\ref{cv-origin-model-correction}(1), with the element of
Lemma~\ref{cv-origin-model-hodge}(4) for $a'_{\tau,i}$, the multidegree $a_{\tau,i}$ of the root line of the data
of $R$ and the generators, gives a smooth metrized line bundle $\bar T_{\tau,i}$ of multidegree
$a'_{\tau,i}-a_{\tau,i}$, invariant under conjugation, with pluriharmonic local potentials on $\tilde X'$. Write
$\bar L^\ast$ for the pullback of a metrized line bundle $\bar L=(a,h_{\bar L})$ for $y$ to $y'$, with function
$w\mapsto h_{\bar L}((w_{k',i'})_{k,i})$ and the multidegree $a$ extended by $0$.

The flag model $M$ has the point $y'$, the arithmetic variety $X'$, the bundles
$\bar x'_{\tau,i}(t)=\bar x_{\tau,i}(t)^\ast+\bar T_{\tau,i}$, where $\bar x_{\tau,i}(t)$ are those of $R$, the
bundle $\bar A'=\bar A^\ast$, the integer $k_A$ of $R$, and $\delta'=[F_c:\mathbf Q]d'$. The data of (2c) for $M$ and
$S'$ are: the Hodge component $Z'$ with its flag of degree $d'$, root lines of multidegrees $a'_{\tau,i}$ with
generators $e'_{\tau,i}$, and Hodge metrics $\bar\sigma_{\tau,i}^\ast+\bar T_{\tau,i}$, which are Hodge metrics by
Lemma~\ref{cv-origin-model-metric} (the abelian loci correspond, since $\mathcal A_{Z'}$ and $\mathcal A_Z$ are the
pullbacks of $\mathcal A$ and $f'=f\circ h$); the embedding $\iota\circ g$, of degree $d'\cdot1$, which carries the
$b^c_{s,i}$ to the $b'_{s,i}$ by (2) of Lemma~\ref{cv-origin-change} and the elements algebraic over $\mathbf Q$
into $F_c$ by (1); the index $k_0'$ with the composite bijection; the manifold $\tilde X'$ with the lifts
$\lambda'_x$, the map $\pi'$ and the bundles $\bar y'_{\tau,i}=\bar y_{\tau,i}^\ast-\sum_{k\le i}\bar T_{\tau,k}$,
$\bar y'_A=\bar y_A^\ast$, and $\bar y'_b(t)$ defined in the same way from $\bar y_b(t)$. Condition (i) holds by
the choice of $\delta'$, and (ii) because pulling back keeps the multidegree $1$ at $k_0'$ and the metric. Condition
(iii) holds by Lemma~\ref{cv-origin-model-section}. In (iv), the coordinates of $\lambda'_x(x)$ in the block $k_0'$
are those of $\lambda_{h'(x)}(h'(x))$ in the block $k_0$ up to a scalar, which are proportional to
$\mathrm{pt}_c(\pi'(x))$, and a complex point of $Z'$ is abelian if and only if its image in $Z$ is. In (v), the
multidegrees add, the metric identities for $\bar y_b$ and $\bar y_b(t)$ pull back, and the logarithms of the functions
of the $\bar T_{\tau,k}$ cancel on the two sides; the first Chern forms of the $\bar y'_b$ on $\tilde X'$ are the
pullbacks by $\pi'$ of those of the $\bar L_b$ by Lemma~\ref{cv-origin-model-correction}(2) and (3). The Hodge
component $Z'$ satisfies condition (1) of Lemma~\ref{cv-origin-change-projective} and has the integers $R_p$ of
$Z'$, the map $\tilde X'\to X'(\mathbf C)$ is bijective and immersive everywhere by
Lemma~\ref{cv-origin-model-resolution}, and the generic fibre of $X'$ is smooth by (5).
\end{proof}
% Lean: root_model_on_cover: hodge_cpt_birational_arith_var, exists_analytic_log_resolution (R'), withTop o.Re.R,
% log_resolution_functorial_holomorphic (hh), log_resolution_reindexed (degree Fl.deg), holomorphic_degree_composition,
% dominating_regular_push, dominating_coords_reindexed, hodge_comp_descent_first to _fourth, degree_reindexed_coords
% and positive_homogeneous_potential (T r), RM with x t r = pullHerm (R.x t r) tensor T r, A = pullHerm R.A,
% delta = finrank Q F_c * Fl.deg, Re' with pi = o.Re.pi o hh, ref a = pullHerm (o.Re.ref a) tensor corr T a and
% ref_pot from log_resolution_potential_pullback and flat_on_log_resolution_corr, and o' with sing r = pullHerm
% (o.sing r) tensor T r (hodge_comp_descent_fifth), eff = root_origin_coord_field, interior from mem_abLocus_iff,
% fund_deg, fund_sing and fund_reg from log_tensor_pull and log_stepLine.

\begin{lemma}\label{cv-origin-data}
Let $E$ be a CM field with a host structure as in Theorem~\ref{cv-flat-degree}, let $\Phi$ be its set
of CM types, and let a PEL unitary datum of signature $(p,q)$ over $E$ with $p\ge1$ and $q\ge1$ be given,
together with its moduli schemes at all levels, a level, a coefficient field $F_c$, and a global model
$\mathcal Y$ of the moduli scheme over the ring of integers of $F_c$
(Definition~\ref{cv-volume-global-model}) such that $F_c$ is algebraically closed in the function field
of every irreducible component of $\mathcal Y$. Put $N=p+q$. Use
Definitions~\ref{cv-origin-hodge-component} and~\ref{cv-origin-tower} with $\mathcal X=\mathcal Y$ and
$\Psi=\Phi$.
\begin{enumerate}
\item Let $\eta$ be the generic point of an irreducible component of $\mathcal Y$. Then there are an
  analytic system $S$ of type $(\Phi,N)$ and data at one component as in
  Definition~\ref{cv-origin-tower}(2b), with $S$, $k_A$, $q$ and $n_{\tau,i}$ in place of $S_c$,
  $k_{A,c}$, $q_c$ and $n^c_{\tau,i}$, whose point $\xi$ lies over $\eta$ and whose integers satisfy
  $n_{\tau,i}=k_Aq(1-i)$ for all $\tau$ and $i$; and a flag model $R$ of type $(\Phi,N)$ with data as in
  (2c) for $R$ in place of $M(B,c)$ and the empty set in place of $S_B$, with these data of (2b). These
  satisfy: the degree $\deg$ is $1$; the map $\tilde X\to X_R(\mathbf C)$ is bijective and immersive
  everywhere, given near each point by holomorphic lifts to the coordinates; the map
  $\pi\colon\tilde X\to S$ is bijective; the multidegree of the bundle $\bar A$ of $R$ is ample on the
  generic fibre in the sense of (3e); and (3c) and (3d) hold for $S$ and these data.
\item Let $S$, the data of (2b) and $R$ with its data be as in (1), and let $S'$ be a finite set of
  primes that contains every prime dividing the discriminant $d_E$ of $E$. Then there are a flag model
  $M$ of type $(\Phi,N)$ and data as in (2c) for $M$ and $S'$, with the same data of (2b), whose Hodge
  component has $R_p=4$ and satisfies the condition of (3a) at every $p\in S'$, whose map
  $\tilde X\to X_M(\mathbf C)$ is bijective and immersive everywhere, and such that the generic fibre of
  the arithmetic variety of $M$ is smooth.
\end{enumerate}
\end{lemma}
\leanref{Development/AutomorphicSpectrum/Proofs/GeometryProofs/ExistsGoodRootTower.lean}{GlobalTheory.CompanionTheory.Ar.TopTower.comp_data}
\leanref{Development/AutomorphicSpectrum/Proofs/GeometryProofs/ExistsGoodRootTower.lean}{GlobalTheory.CompanionTheory.Ar.TopTower.level_data}
\begin{proof}
Put $D=|\Phi|N(N-1)/2$. (1) Let $\Omega_0$ be the residue field of $\mathcal Y$ at $\eta$. The local rings of
$\mathcal Y$ are domains, as $\mathcal Y$ is normal, and $\eta$ has no generization other than itself, so the local
ring at $\eta$ is the field $\Omega_0$ and the local rings at its specializations embed in it:
$\operatorname{Spec}\Omega_0\to\mathcal Y$ is as in Lemma~\ref{cv-origin-reference-char}. By
Lemma~\ref{cv-origin-reference-exists} there is a reference geometry of $\mathcal Y$ over the component of $\eta$,
with field $\Omega$, point $y$, Hodge component $Z$, flag of degree $1$ with bases $b_{s,i}$ and root lines of
multidegrees $a_{\tau,i}$ with generators $e_{\tau,i}$, and index $k_0$. By Lemma~\ref{cv-origin-analytic-exists} it
has analytic data, with Hodge metrics $\bar\sigma_{\tau,i}$, analytic system $S$, manifold $\tilde X$ with lifts
$\lambda_x$, map $\pi$ and bundles $\bar y_b$ and $\bar y_b(t)$, and a holomorphic $\psi:S\to\tilde X$ with
$\psi\circ\pi$ the identity; as $\pi$ is bijective, $\pi\circ\psi$ is the identity too. For the weight
$\nu_{\tau,i}=1-i$, which decreases strictly in $i$, Lemma~\ref{cv-origin-weight} gives integers $k_A,q\ge1$ and
$s\neq0$ in $\Omega$ with property (2c)(iii) for $n_{\tau,i}=k_Aq(1-i)$.

From these, the data of (1) are as follows. By Lemma~\ref{cv-origin-reference-closure}, $I_{k_0}$ is nonempty;
number it as $\{0',\dots,m'\}$. The data of Definition~\ref{cv-origin-tower}(2b) are: the system $S$, whose bundle
$\bar L_A$ has a positive metric, with $\varphi_A=\varphi_A(t)=0$, by Definition~\ref{cv-origin-analytic}(iv); the
point $\xi$ of $\mathcal Y$ given by the generic point of $Z$ followed by $f$, with values in $\Omega_c=\Omega$, which
lies over $\eta$; the vectors $b^c_{s,i}$, the images of the $b_{s,i}$ in the Lie algebra of $\mathcal A$ at $\xi$,
through the identification of $\mathcal A_Z$ with $f^*\mathcal A$; $y^c_i=y_{k_0,i'}$, whose ratios generate $\Omega$
by Lemma~\ref{cv-origin-reference-closure}; $h_c(u)=\big(\sum_i|u_i|^2\big)^{1/2}$;
$\mathrm{pt}_c(x)=(\lambda_{\psi(x)}(\psi(x))_{k_0,i'})_i$; $k_{A,c}=k_A$, $q_c=q$, $n^c_{\tau,i}=k_Aq(1-i)$; and
$s_c=s$. Here $\mathrm{pt}_c(x)\neq0$, since a point of $\hat X(\mathbf C)$ has a nonzero coordinate in every block. If
$\mathrm{pt}_c(x)=u\,\mathrm{pt}_c(x')$, then $\lambda_{\psi(x)}(\psi(x))$ and $\lambda_{\psi(x')}(\psi(x'))$ lie in one
orbit by Lemma~\ref{cv-origin-reference-orbit}, so $\psi(x)=\psi(x')$ by Definition~\ref{cv-origin-analytic}(3), and
$x=\pi(\psi(x))=\pi(\psi(x'))=x'$.

By Lemma~\ref{cv-origin-weight-birational} and Definition~\ref{cv-origin-analytic}(3), $\Omega$ is the function field of
the kernel $X$ of $y$, an arithmetic subvariety of dimension $D+1$. The flag model $R$ has the point $y$, the
arithmetic variety $X$, the bundles $\bar x_{\tau,i}(t)=\bar y_{\tau,i-1}(t)-\bar y_{\tau,i}(t)$, where
$\bar y_{\tau,0}(t)$ is the trivial bundle with function $1$, the bundle $\bar A=\bar y_A$, the integer $k_A$ and
$\delta=[F_c:\mathbf Q]$. The bundles $\bar x_{\tau,i}(t)$ are smooth and invariant under conjugation by
Definition~\ref{cv-origin-analytic}(5), and so is $\bar A$ by (iii) there; the multidegree of $\bar x_{\tau,i}(t)$ is
$-\sum_{k<i}a_{\tau,k}+\sum_{k\le i}a_{\tau,k}=a_{\tau,i}$ by (ii), and $-\sum_{k\le i}\bar x_{\tau,k}(t)=\bar y_{\tau,i}(t)$,
the sum telescoping. The data of (2c) for $R$ and the empty set are: the Hodge component $Z$ with its flag, of
degree $1$, the generators $e_{\tau,i}$ and the Hodge metrics $\bar\sigma_{\tau,i}$; the identity $g$ of $\Omega$, which
carries the elements of $\Omega$ algebraic over $\mathbf Q$ into $F_c$ by Lemma~\ref{cv-origin-reference-constants};
the index $k_0$ with $i\mapsto i'$; and $\tilde X$, $x\mapsto\lambda_x(x)$, $\pi$ and the bundles $\bar y_b$. Condition (i)
holds as $\delta=[F_c:\mathbf Q]\cdot1$; (ii) holds by Definition~\ref{cv-origin-analytic}(iii); (iii) is
Lemma~\ref{cv-origin-weight}; (iv) holds since $\psi(\pi(x))=x$, so the coordinates of $\lambda_x(x)$ in the block
$k_0$ are $\mathrm{pt}_c(\pi(x))$, and by Definition~\ref{cv-origin-analytic}(i); and (v) holds by
Definition~\ref{cv-origin-analytic}(ii) and (5) and the identity $-\sum_{k\le i}\bar x_{\tau,k}(t)=\bar y_{\tau,i}(t)$.

These data satisfy the conditions of (1). The degree is $1$. The map $\tilde X\to X(\mathbf C)$ is bijective and
immersive everywhere and given by the lifts $\lambda_x$, and $\pi$ is bijective, by
Definition~\ref{cv-origin-analytic}(3) and (4). The multidegree $a_A$ of $\bar A$ is ample on the generic fibre by
Lemma~\ref{cv-origin-reference-ample}, and (3c) holds by Definition~\ref{cv-origin-analytic}(v). For (3d), let
$x\in S$ and $u=\psi(x)$. On $\psi^{-1}(U_u)$ the coordinates of $\lambda_u(\psi(x'))$ in the block $k_0$ are
holomorphic in $x'$, and they are nonzero multiples of $\mathrm{pt}_c(x')$, since the lifts agree up to the action of
$(\mathbf C^\times)^K$. If the $p$-th of them is nonzero at $x$, the ratios to it have injective derivative at $u$ in a
chart of $\tilde X$ by Lemma~\ref{cv-origin-reference-immersion}, and so at $x$ in a chart of $S$ after composing with
$\psi$, by Lemma~\ref{cv-origin-analytic-inverse}.

(2) Let $S$, the data of (2b) and $R$ with its data be as in (1), and let $S'$ be given. By
Lemma~\ref{cv-origin-change}, with $\mathcal X=\mathcal Y$, there are $\Omega'$, $\iota$, $y'$, $Z'$, $h$ and a flag of
$Z'$ with root lines satisfying its conclusions (1) to (5). By (1), the map $\tilde X\to X(\mathbf C)$ of $R$ is
bijective and immersive everywhere and $\pi$ is bijective, so Lemma~\ref{cv-origin-model} gives a flag model $M$ of
type $(\Phi,N)$ and data as in (2c) for $M$ and $S'$, with the same data of (2b), whose Hodge component satisfies
condition (1) of Lemma~\ref{cv-origin-change-projective} at every $p\in S'$ with the integers $R_p=4$ of $Z'$, which
is the condition of Definition~\ref{cv-origin-tower}(3a), whose map $\tilde X'\to X_M(\mathbf C)$ is bijective and
immersive everywhere, and such that the generic fibre of the arithmetic variety of $M$ is smooth.
\end{proof}
% Lean, for the next writer. (1) is comp_data: isCptGen_residue_field_generic (the residue field at x),
% exists_reference_geometry_same_generic (Lemma cv-origin-reference-exists), exists_holomorphic_section_of_analytic
% (Lemma cv-origin-analytic-exists), exists_effective_weight_section with nu r = -r.2, zero-based (Lemma
% cv-origin-weight), and effective_scaled_weight_of_retraction (= RefRoot.refRoot, with refFixed, refModel,
% refOrigin, stepLine and log_line_eq), the assembly written above: EmbedsGeneric.exists_closure,
% same_orbit_of_proportional_coords, lie_model_push_along_map, hodge_cpt_birational_arith_var,
% algebraic_constants_in_range, ample_of_single_coord_block, holomorphic_atlas_coord_compat and
% holomorphic_inverse_chart_value. (2) is level_data: root_origin_good_primes_finite (Lemma cv-origin-change) and
% root_origin_coords_over_field (Lemma cv-origin-model).

\begin{theorem}\label{cv-volume-good-tower}
Let $E$ be a CM field with a host structure as in Theorem~\ref{cv-flat-degree}, let $\Phi$ be its set of
CM types, and let a PEL unitary datum of signature $(p,q)$ over $E$ with $p\ge1$ and $q\ge1$ be given,
together with its moduli schemes at all levels, a level, a coefficient field $F_c$, and a global model
$\mathcal Y$ of the moduli scheme over the ring of integers of $F_c$
(Definition~\ref{cv-volume-global-model}) such that $F_c$ is algebraically closed in the function field
of every irreducible component of $\mathcal Y$. Then there is a tower of type $(\Phi,p+q)$ with a good
origin in $\mathcal Y$ (Definition~\ref{cv-origin-tower}).
\end{theorem}
\leanref{Development/AutomorphicSpectrum/Statements/TowerExistence/GoodTowerExistsX.lean}{GlobalTheory.CompanionTheory.Ar.Top.GoodTowerExistsX}
\leanref{Development/AutomorphicSpectrum/Definitions/GlobalGeometry/SplitComps.lean}{GlobalTheory.CompanionTheory.Ar.SplitComps}
\leanref{Development/AutomorphicSpectrum/Proofs/GeometryProofs/ExistsGoodRootTower.lean}{GlobalTheory.CompanionTheory.Ar.TopTower.exists_good_tower}
\begin{proof}
Put $N=p+q$; as $p\ge1$ and $q\ge1$, $1\le p<N$. By Lemma~\ref{cv-origin-components} the points of
$\mathcal Y$ with no generization other than themselves form a finite nonempty set $C$ of generic points
of irreducible components; these are the components of the tower. For $c\in C$,
Lemma~\ref{cv-origin-data}(1) with $\eta=c$ gives an analytic system $S_c$, data of
Definition~\ref{cv-origin-tower}(2b) whose point $\xi_c$ lies over $c$, with integers $k_{A,c}$, $q_c$ and
$n^c_{\tau,i}=k_{A,c}q_c(1-i)$, and a flag model $R_c$ with data as there.

Let $M'$ be the least common multiple of the two integers of the level, which is nonzero, and let the
bad primes be the primes dividing $2d_EM'$; $2$ is one of them. For an integer $B$, the set $S_B$ of the
primes $p\le B$ together with the bad primes contains every prime dividing $d_E$. For each $B$ and each
$c\in C$, Lemma~\ref{cv-origin-data}(2), applied to the data of $c$ and to $S'=S_B$, gives a flag model
$M(B,c)$ and data of (2c) for it at $S_B$, with the data of (2b) at $c$.

Let $T$ be the tower of type $(\Phi,N)$ with index $j=p$, components $C$, these bad primes, the weight
$\nu_{\tau,i}=1-i$, which has rational entries, and the flag models $M(B,c)$. The data above form an
origin of $T$ in $\mathcal Y$. Condition (2a) holds by the choice of $j$ and of the bad primes. In (2b),
$n^c_{\tau,i}=k_{A,c}q_c\nu_{\tau,i}$, and the point $\eta_c$ under $\xi_c$ is $c$, so $c\mapsto\eta_c$ is
the identity of $C$, a bijection onto the generic points of the irreducible components. The data of (2c)
at $(B,c)$ are those of Lemma~\ref{cv-origin-data}(2). The origin is good: (3a) and (3b) hold by
Lemma~\ref{cv-origin-data}(2); (3c) and (3d) hold by Lemma~\ref{cv-origin-data}(1); and (3e) holds with
the flag model $R_c$ and its data, which have $\deg=1$, a map $\tilde X\to X_{R_c}(\mathbf C)$ that is
bijective and immersive everywhere, and an ample multidegree $a_A$, by Lemma~\ref{cv-origin-data}(1).
% Lean: GoodTowerExistsX is proved, together with VolumeTheoremX, by
% Link.good_tower_volume_theorem_bis through final_theorem_of_inputs_right, goodTower_volume_of_reduced and
% Top.goodTower_and_volume, as GlobalTheory.T.exists_good_root_tower = TopTower.exists_good_tower:
% generic_points_finite_nonempty, comp_data (Sys, Fx, the reference origin R, refOr with ref_deg,
% ref_top and ample, flat, and the immersion of GoodOrigin.Extra) and level_data (the flag models RM, the
% origins o' at S_B, omVB, R p = 4 from hR, and smooth). In the arithmetic volume step it is used through
% Plug.vs_tower, which takes the global model from Rec.exists_isGlobal and the splitting of components from
% T.split_components_of_alg_closed, and DeepStR (deepStR_of_volume).
\end{proof}

\section{The covolume bound}\label{section-cv-covolume-bound}

\subsection{Local shares and group determinants}\label{subsection-cv-local-shares}

\begin{situation}\label{cv-point}
Let $\varepsilon>0$, let $\mathcal F$ be a family of definite-side data with exponent
$\varepsilon$ (Chapter~\ref{chapter-assembly}) satisfying the family hypotheses, let $x$ be a scale of
$\mathcal F$ and let $P$ be a point of $\mathcal F$ at $x$ satisfying the point conditions
and the coefficient field bound of Chapter~\ref{chapter-assembly}. Write $\mathcal D$ for the
definite datum at $P$ (Definition~\ref{th-definite-datum}), with ranks $k$, $m$, $t$, $f$,
host fields $F\subset E$ with $[F:\mathbf Q]=f$, prime $d$ and Kronecker character $\chi$
of $\mathbf Q(\sqrt{-d})$, and Harish-Chandra parameters with weights $e_j(\rho)$ for
$1\le j\le m$ and $\rho$ running over the index set of the parameters, which has $f$
elements (Definition~\ref{th-hc-parameter}). Write $J\ge1$ for the remaining parameter of
the scale and $b=\exp(x^{\varepsilon})$.
Let $L$ be the coefficient field of $\mathcal D$, a totally complex Galois number field containing
$E$. For $1\le j\le m$ let $W_j$ be the Hermitian space of the $j$-th indefinite group
$U(W_j)$, of signature $(p_j,q_j)=(k+j-1,m-j)$ at every real place of $F$, and let
$M_j\subset L^{d_j}$ be the lattice of $j$-th lifted forms at $P$, with the metric $X_j$ on
$L^{d_j}$ given by the Petersson norms (Definition~\ref{cv-hermitian-norms}). Let $D$
be the rank of the glued lattice at $P$ (Chapter~\ref{chapter-integral-structure}). Put
\[
 \mathrm{main}_j=\sum_{\rho}(p_j-q_j)\,e_j(\rho)\log\Omega_E(\chi),\qquad
 u=f\,m\,b\,x^{\varepsilon}\log\log x,
\]
with $\log\Omega_E(\chi)$ as in Definition~\ref{cv-log-omega}.
\end{situation}
\leanref{Development/AbelianVarietyTheory/Definitions/Divisors/TupleReadBaseTupleEqFamily.lean}{GlobalTheory.BigTop.PtAt}
\leanref{Formalism/LocalToGlobal/IndefiniteForms/IndefSide.lean}{Formalism.L2.Indef.IndefData}
\leanref{Formalism/LocalToGlobal/Divisibility/ZscOfNum.lean}{Formalism.L2.DIV.ZscOfNum}
\leanref{Development/AbelianVarietyTheory/Definitions/Divisors/CovolXOfCovol.lean}{GlobalTheory.CovTop.mainT}
\leanref{Development/AbelianVarietyTheory/Definitions/Divisors/CovolXOfCovol.lean}{GlobalTheory.CovTop.unitE}

\begin{theorem}\label{cv-covolume-bound}
There is $\varepsilon_1>0$ with the following property. For every $\varepsilon$ with
$0<\varepsilon\le\varepsilon_1$ there is a largeness threshold such that, for every
family $\mathcal F$ of definite-side data with exponent $\varepsilon$ that satisfies the
family hypotheses and is large for that threshold, there are $C_0\ge0$ and $x_0$ with the
following property. Let $x\ge x_0$ be a scale of $\mathcal F$, let $P$ be a point of
$\mathcal F$ at $x$ satisfying the point conditions and the coefficient field bound, and
use the notation of Situation~\ref{cv-point}. Then for every $C\ge C_0$ and every
$1\le j\le m$,
\[
 \frac1D\operatorname{lcv}_{X_j}(M_j)\le\mathrm{main}_j+C\,u.
\]
\end{theorem}
\leanref{Development/AutomorphicArithmetic/ThetaCorrespondence/ThetaPairLifts/ThirdPlaceSeries/LiftLambdaBound.lean}{GlobalTheory.CovTop.CovolumeBoundOfX}
\leanref{Formalism/LocalToGlobal/IndefiniteForms/IndefSide.lean}{Formalism.L2.Indef.IndefData.L3}
\leanref{Development/AutomorphicArithmetic/ThetaCorrespondence/ThetaPairLifts/ThirdPlaceSeries/EveryCharacterOfLiteratureOfVWin.lean}{UpperTheory.LiteratureTheorems.every_character_of_literature_of_vWin}
\begin{proof}
This is Theorem~\ref{cv-covolume-bound-from-literature} together with the proofs of its
hypotheses. Each of them is a published theorem proved in full, collected in
Chapter~\ref{chapter-literature}; Theorems~\ref{cv-functorial-resolution}
and~\ref{cv-moriwaki-bigness} are two of them.
\end{proof}

The bound is one-sided and has the shape of the main term of
Chapter~\ref{chapter-inert-prime-balance}: summed over $j$, the terms $\mathrm{main}_j$ are
the multiples of $\log\Omega_E(\chi)$ against which the definite side is compared, and the
error is $m\,C\,u$.

\begin{definition}\label{cv-separated}
Let $g$ be a prime and $r\ge1$, fix a labelling of $\bar{\mathbf F}_g$
(Definition~\ref{lit-labelling}) with the cuspidal characters $\tau_P$ of
Corollary~\ref{lit-green-labelled}, and let $U\subset GL_r(\mathbf F_g)$ be the subgroup of upper
unitriangular matrices. For a finite group $\Gamma$, an element $a=\sum_ta_tt$ of $\mathbf C[\Gamma]$
and an irreducible character $\eta$ of $\Gamma$, write $a(1)=a_1$ and
$e_\eta=\frac{\eta(1)}{|\Gamma|}\sum_{t\in\Gamma}\eta(t^{-1})t$.
\begin{enumerate}
\item A \emph{label} is a monic polynomial $f\in\mathbf F_g[X]$ of degree $r$ with $f(0)\ne0$. An
  irreducible character $\eta$ of $GL_r(\mathbf F_g)$ \emph{has label} $f$ if $f=P_1\cdots P_s$ for
  monic irreducible polynomials $P_i$ with nonzero constant terms such that $\eta$ occurs in the
  Harish-Chandra induction of $\tau_{P_1}\otimes\cdots\otimes\tau_{P_s}$ from the standard Levi
  subgroup $\prod_iGL_{\deg P_i}(\mathbf F_g)$. The \emph{series idempotent} $e(f)$ is the sum of the
  $e_\eta$ over the irreducible characters $\eta$ with label $f$.
\item For a nontrivial additive character $\psi$ of $\mathbf F_g$, the \emph{Gelfand--Graev
  idempotent} is
  $e_U(\psi)=|U|^{-1}\sum_{u\in U}\psi\big(\sum_{i<r}u_{i,i+1}\big)^{-1}u$, and the \emph{weight} of
  a label $f$ is
  $\operatorname{wt}_\psi(f)=\operatorname{Re}\big((e_U(\psi)e(f))(1)\big)/\operatorname{Re}\big(e_U(\psi)(1)\big)$.
\item Let $D$ be a definite-side datum (Definition~\ref{th-definite-datum}) with ranks $k,m$ and
  $r=k+m-1$, level primes $g_0,g_1$, and level places and tame values $\chi_{k,v,j}$
  (Definition~\ref{asm-box-row-tame}). A \emph{tuple of labels at $g_k$} is a family $t=(t_v)$ of
  labels over $\mathbf F_{g_k}$, one for each level place $v$ over $g_k$. It is \emph{separated} if:
  \begin{enumerate}
  \item[(a)] every $t_v$ is squarefree;
  \item[(b)] for every prime $\ell\notin\mathcal B_r(g_k)$ (Definition~\ref{hf-good-host}) and every
    $v$, the roots of $t_v$ in $\bar{\mathbf F}_{g_k}^\times$ have pairwise distinct prime-to-$\ell$
    parts;
  \item[(c)] for every prime $\ell\ne g_k$ and all $v$ and $j$, no root of $t_v$ has the same
    prime-to-$\ell$ part as $\chi_{k,v,j}$ or as $\chi_{k,v,j}^{-1}$;
  \item[(d)] for every integer $a$ with $1\le2a\le r$ there is a level place $v$ such that $t_v$ has
    no factor of degree $a$;
  \item[(e)] some $t_v$ has an irreducible factor of prime degree $q$ with $r<2q$ and $q+3\le r$.
  \end{enumerate}
\item Let $G=\prod_{k=0,1}\prod_{v\mid g_k}GL_r(\mathbf F_{g_k})$, the product over the level places.
  An irreducible character $\bigotimes_{k,v}\eta_{k,v}$ of $G$ is \emph{separated} if, for $k=0$ and
  $k=1$, the labels of the $\eta_{k,v}$ form a separated tuple at $g_k$. The \emph{separated
  idempotent} $e_{\mathrm{sep}}\in\mathbf C[G]$ is the sum of the $e_\eta$ over the separated
  characters $\eta$. Given nontrivial additive characters $\psi_k$ of $\mathbf F_{g_k}$, the
  \emph{share} of a set $S$ of tuples of labels at $g_k$ is
  $s_k(S)=\sum_{t\in S}\prod_v\operatorname{wt}_{\psi_k}(t_v)$.
\end{enumerate}
\end{definition}
\leanref{Development/AbelianVarietyTheory/Definitions/NonGaloisDescent/LabelDatum.lean}{GlobalTheory.GLq.eSerC}
\leanref{Formalism/SharpRestriction/LevelLetterCocycles/ArithmeticCylinders.lean}{Formalism.MainNotions.ArithmeticCylinders.wtOf}
\leanref{Formalism/SharpRestriction/LevelLetterCocycles/ArithmeticCylinders.lean}{Formalism.MainNotions.ArithmeticCylinders.goodAt}
\leanref{Formalism/SharpRestriction/LevelLetterCocycles/ArithmeticCylinders.lean}{Formalism.MainNotions.ArithmeticCylinders.locShare}
\leanref{Formalism/CyclotomicNorms/StatementT.lean}{Formalism.CyclotomicNorms.IsGood}

Every irreducible character of $GL_r(\mathbf F_g)$ has exactly one label, so the series idempotents
of distinct labels are orthogonal central idempotents whose sum is $1$. Hence $e_{\mathrm{sep}}$ is
the product over $k$ of the sums, over the separated tuples $t$ at $g_k$, of
$\prod_{v\mid g_k}e(t_v)$.

\begin{lemma}\label{cv-bad-part}
Let $D$ be a definite-side datum with level primes $g_0,g_1$ and let $G$ and $e_{\mathrm{sep}}$ be as
in Definition~\ref{cv-separated}. Let $P\subset G$ be the subgroup of elements that are upper
unitriangular at every level place, let $\psi_{k,v}$ be nontrivial additive characters of
$\mathbf F_{g_k}$, one for each level place $v$, and let $\lambda:P\to\mathbf C^\times$ be
$\lambda(u)=\prod_{k,v}\psi_{k,v}\big(\sum_{i<r}u_{v,i,i+1}\big)$. For $k=0,1$ let $H_k\subset G$ be
the subgroup of elements that are upper unitriangular at the level places over $g_{1-k}$. Let
$\Gamma$ be a finite group with a surjection $\pi:\Gamma\to G$ and a homomorphism $\iota:G\to\Gamma$
with $\pi\iota=\mathrm{id}$ whose image commutes with $\ker\pi$. Let $\rho$ be a representation of
$\Gamma$ on a finite-dimensional complex vector space $V$, put $\Gamma_P=\pi^{-1}(P)$,
$\Gamma_k=\pi^{-1}(H_k)$ and $\tilde\lambda=\lambda\circ\pi$ on $\Gamma_P$, and let $V^{\tilde\lambda}$
be the subspace on which $\Gamma_P$ acts through $\tilde\lambda$. Let $c\ge0$, and suppose that for
$k=0,1$ there is $n_k\ge0$ such that the restriction of $V$ to $\Gamma_k$ embeds into the sum of $n_k$
copies of the regular representation of $\Gamma_k$, with
$n_k|\Gamma_k|\le c\,|\Gamma_P|\dim V^{\tilde\lambda}$. Let $\psi_k$ be nontrivial additive characters
of $\mathbf F_{g_k}$ and $S_k$ the set of tuples of labels at $g_k$ that are not separated. Then
\[
 \dim\big(V^{\tilde\lambda}\cap\ker\rho(\iota(e_{\mathrm{sep}}))\big)
 \le c\,\big(s_0(S_0)+s_1(S_1)\big)\dim V^{\tilde\lambda}.
\]
\end{lemma}
\leanref{Development/AbelianVarietyTheory/Proofs/DivisorProofs/BadPartDimLeSumLocalSharesWhittaker.lean}{GlobalTheory.BadCount.bad_part_le_locShare}
\leanref{Development/AbelianVarietyTheory/Proofs/DivisorProofs/FinrankCharacterEigenspaceInfKerLe.lean}{GlobalTheory.BadCount.finrank_eig_inf_ker_le}
\begin{proof}
Step 1: the idempotents. For $k=0,1$ let $b_k\in\mathbf C[G]$ be the sum over $t\in S_k$ of
$\prod_{v\mid g_k}e(t_v)$, placed in the factors over $g_k$. By the remark after
Definition~\ref{cv-separated}, $b_k$ is a central idempotent, $1-b_k$ is the corresponding sum over
the separated tuples, and $e_{\mathrm{sep}}=(1-b_0)(1-b_1)$. So
$1-e_{\mathrm{sep}}=b_0+b_1-b_0b_1$.

Step 2: the weights. Let $p=|P|^{-1}\sum_{u\in P}\lambda(u)^{-1}u$; it is the product over the level
places of the Gelfand--Graev idempotents $e_U(\psi_{k,v})$. Two nontrivial additive characters
satisfy $\psi'(z)=\psi(az)$ for some $a\in\mathbf F_g^\times$, and conjugation by
$\operatorname{diag}(a^{r-1},\dots,a,1)$ carries $e_U(\psi')$ to $e_U(\psi)$; as $e(f)$ is central,
$\operatorname{wt}_{\psi'}=\operatorname{wt}_\psi$. The element $b_kp$ is supported on $H_k$, because
$b_k$ is supported on the elements that are trivial at the places over $g_{1-k}$ and $p$ on $P$. Its
coefficient at $1$ is the product of $e_U(1)=|U|^{-1}$ over the places over $g_{1-k}$ and of
$\sum_{t\in S_k}\prod_{v\mid g_k}(e_U(\psi_{k,v})e(t_v))(1)$. As $e(f)$ is central,
$e_U(\psi)e(f)$ is an idempotent, so its coefficient at $1$ is real, being its rank on
$\mathbf C[GL_r(\mathbf F_g)]$ divided by $|GL_r(\mathbf F_g)|$; hence it equals
$\operatorname{wt}_\psi(f)/|U|$. Since $P$ is the product over the level places of the groups of
upper unitriangular matrices, this gives $(b_kp)(1)=s_k(S_k)/|P|$.

Step 3: lifting to $\Gamma$. Let $\tilde p=|\Gamma_P|^{-1}\sum_{\gamma\in\Gamma_P}\tilde\lambda(\gamma)^{-1}\gamma$.
It acts as the identity on $V^{\tilde\lambda}$, it equals $\iota(p)$ times the average over
$\ker\pi$, and it commutes with $\iota(b_k)$ because $\iota(G)$ commutes with $\ker\pi$. So
$a_k=\iota(b_k)\tilde p$ is an idempotent of $\mathbf C[\Gamma_k]$ with coefficient
$a_k(1)=(b_kp)(1)/|\ker\pi|=s_k(S_k)/|\Gamma_P|$.

Step 4: the rank of $a_k$. Let $j:V\to\mathbf C[\Gamma_k]^{n_k}$ be the embedding of the hypothesis.
It carries $\rho(a_k)V$ into $(a_k\mathbf C[\Gamma_k])^{n_k}$. Left multiplication by $a_k$ on
$\mathbf C[\Gamma_k]$ is an idempotent of trace $|\Gamma_k|a_k(1)$, since left multiplication by
$\gamma\ne1$ has trace $0$. Hence
$\dim\rho(a_k)V\le n_k|\Gamma_k|a_k(1)\le c\,s_k(S_k)\dim V^{\tilde\lambda}$.

Step 5: injectivity. Let $W=V^{\tilde\lambda}\cap\ker\rho(\iota(e_{\mathrm{sep}}))$ and consider
$x\mapsto(\rho(a_0)x,\rho(a_1)x)$ on $W$. For $x\in W$ we have $\rho(\tilde p)x=x$, so
$\rho(a_k)x=\rho(\iota(b_k))x$. If both vanish, Step 1 gives
$x=\rho(\iota(1-e_{\mathrm{sep}}))x=\rho(\iota(b_0+b_1-b_0b_1))x=0$. So the map is injective, and
$\dim W\le\dim\rho(a_0)V+\dim\rho(a_1)V$, which with Step 4 is the claim.
\end{proof}
\leanref{Development/AbelianVarietyTheory/Proofs/DivisorProofs/CoeffOneRealOfIdem.lean}{GlobalTheory.BadCount.badElt_mul_lvlIdem}

\begin{lemma}\label{cv-squarefree-mass}
Let $K$ be a finite field with $g=|K|$ elements and let $N\ge1$. Let $\mathcal S_N$ be the set of
monic squarefree polynomials $f\in K[X]$ of degree $N$ with $f(0)\ne0$, and for $f\in\mathcal S_N$ put
$w(f)=\prod_{P\mid f}(g^{\deg P}-1)^{-1}$, over the irreducible factors $P$ of $f$. Then
$\sum_{f\in\mathcal S_N}w(f)\ge1-2/g$.
\end{lemma}
\leanref{Formalism/PadicSubsetSums/StableGradedCounts/LocalEstimateIVOfTypeSetWeight.lean}{Formalism.IrreducibleCountingLemmas.GaloisIrreducibility.square_mass_lower_bound}
\begin{proof}
Each $f\in\mathcal S_N$ has $w(f)\ge\prod_{P\mid f}g^{-\deg P}=g^{-N}$. There are $g^N-g^{N-1}$ monic
squarefree polynomials of degree $N$ if $N\ge2$ \cite[Proposition~2.3]{Ros02}, and $g$ if $N=1$.
Those with $f(0)=0$ are $Xh$ with $h$ monic squarefree of degree $N-1$, so there are at most
$g^{N-1}$ of them. Hence $|\mathcal S_N|\ge g^N-2g^{N-1}$, and
$\sum_{f\in\mathcal S_N}w(f)\ge(g^N-2g^{N-1})g^{-N}=1-2/g$.
\end{proof}

\begin{lemma}\label{cv-local-estimates}
Let $K$ be a finite field with $g=|K|$ elements, let $\bar K$ be an algebraic closure of $K$, let
$N\ge1$, and let $\mathcal S_N$ and $w$ be as in Lemma~\ref{cv-squarefree-mass}. For a prime $\ell$
and $x\in\bar K^\times$, the \emph{prime-to-$\ell$ part} of $x$ is the component of $x$ of order prime
to $\ell$ in the cyclic group generated by $x$. Then:
\begin{enumerate}
\item if $g\ge2^{12}$, the sum of $w(f)$ over the $f\in\mathcal S_N$ for which some prime
  $\ell\notin\mathcal B_N(g)$ (Definition~\ref{hf-good-host}) and two distinct roots of $f$ in
  $\bar K$ have the same prime-to-$\ell$ part is at most $23\log g/\sqrt g$;
\item if $q_1\ne d$ are odd primes with $\max(q_1,d)\ge16$ such that $16$, $q_1$ and $d$ divide
  $g-1$, and $c\in K^\times$, then the sum of $w(f)$ over the $f\in\mathcal S_N$ for which some
  prime $\ell\ne g$ and some root of $f$ in $\bar K$ have the same prime-to-$\ell$ part as $c$ is at
  most $\frac1{4\min(q_1,d)}+\frac{1+\log N}{g-1}$;
\item if $0<a<N$, the sum of $w(f)$ over the $f\in\mathcal S_N$ that have a divisor of degree $a$
  in $K[X]$ is at most $7/8$;
\item if $N\ge100$, the sum of $w(f)$ over the $f\in\mathcal S_N$ that have no irreducible factor of
  prime degree $p$ with $N<2p$ and $p+3\le N$ is at most $1-\frac1{5\log N}$.
\end{enumerate}
\end{lemma}
\leanref{Formalism/PadicSubsetSums/StableGradedCounts/SquareSetMemExceptionalUnion.lean}{Formalism.IrreducibleCountingLemmas.GaloisIrreducibility.local_estimate_second}
\leanref{Formalism/PadicSubsetSums/StableGradedCounts/LocalEstimateThird.lean}{Formalism.IrreducibleCountingLemmas.GaloisIrreducibility.local_estimate_third}
\leanref{Formalism/PadicSubsetSums/StableGradedCounts/LocalEstimateIVOfTypeSetWeight.lean}{Formalism.IrreducibleCountingLemmas.GaloisIrreducibility.local_estimate_fourth}
\leanref{Formalism/PadicSubsetSums/StableGradedCounts/LocalEstimateIVOfTypeSetWeight.lean}{Formalism.IrreducibleCountingLemmas.GaloisIrreducibility.local_estimate_fifth}

Each part is a computation with the weights $w$, in which the polynomials are sorted by the
degrees of their irreducible factors and the number of monic irreducible polynomials of each
degree is bounded; parts (1) and (2) also use the factorizations of the orders $g^i-1$ of the
multiplicative groups of the finite extensions of $K$.

\begin{lemma}\label{cv-property-w}
Let $K$ be a finite field with $|K|\ge2^{12}$, let $N\ge100$, and let $q_1\ne d$ be odd primes with
$\max(q_1,d)\ge16$ such that $16$, $q_1$ and $d$ divide $|K|-1$. Let $I$ and $J$ be finite sets, let
$\xi_{v,j}\in K^\times$ for $v\in I$ and $j\in J$, and let $w$ be a probability measure on the set of
monic polynomials of degree $N$ over $K$ with $w(f)=\prod_{P\mid f}(|K|^{\deg P}-1)^{-1}$, over the
irreducible factors $P$, whenever $f$ is squarefree with $f(0)\ne0$. Call a family $(f_v)_{v\in I}$
of such polynomials separated if it satisfies conditions (a) to (e) of
Definition~\ref{cv-separated}(3) with $r=N$, $g_k=|K|$, $\bar{\mathbf F}_{g_k}$ an algebraic closure of
$K$, and the $\xi_{v,j}$ in place of the $\chi_{k,v,j}$, and with each $f_v(0)\ne0$. Then the
$w^{\otimes I}$-measure of the families that are not separated is at most
\[
 \frac{2|I|}{|K|}+\frac{23|I|\log|K|}{\sqrt{|K|}}
 +2|I||J|\Big(\frac1{4\min(q_1,d)}+\frac{1+\log N}{|K|-1}\Big)
 +\frac N2\Big(\frac78\Big)^{|I|}+e^{-|I|/(5\log N)}.
\]
\end{lemma}
\leanref{Formalism/PadicSubsetSums/StableGradedCounts/PropertyWOfEstimates.lean}{Formalism.IrreducibleCountingLemmas.GaloisIrreducibility.property_W_of_estimates}
\leanref{Formalism/PadicSubsetSums/CentralIdempotents/EigenbasisOfDistinctRoots.lean}{Formalism.IrreducibleCountingLemmas.GaloisIrreducibility.PropW}
\begin{proof}
Write $g=|K|$, let $\mathcal S_N$ be as in Lemma~\ref{cv-squarefree-mass}, and for a set $S$ of monic
polynomials of degree $N$ write $w(S)=\sum_{f\in S}w(f)$. Consider the following sets of such
polynomials:
\begin{itemize}
\item $S_1$, the complement of $\mathcal S_N$;
\item $S_2$, the $f\in\mathcal S_N$ that fail condition (b) of Definition~\ref{cv-separated}(3);
\item $S_3(c)$ for $c\in K^\times$, the $f\in\mathcal S_N$ for which some prime $\ell\ne g$ and some
  root of $f$ have the same prime-to-$\ell$ part as $c$;
\item $S_4(a)$ for $1\le a\le N/2$, the $f\in\mathcal S_N$ with a divisor of degree $a$;
\item $S_5$, the $f\in\mathcal S_N$ with no irreducible factor of prime degree $p$ with $N<2p$ and
  $p+3\le N$.
\end{itemize}
On $\mathcal S_N$ the measure $w$ is given by the formula of Lemma~\ref{cv-squarefree-mass}. Since
$w$ is a probability measure, that lemma and Lemma~\ref{cv-local-estimates} give
\[
 w(S_1)\le\frac2g,\quad w(S_2)\le\frac{23\log g}{\sqrt g},\quad
 w(S_3(c))\le\frac1{4\min(q_1,d)}+\frac{1+\log N}{g-1},\quad w(S_4(a))\le\frac78,\quad
 w(S_5)\le1-\frac1{5\log N},
\]
using $a\le N/2<N$ in the fourth bound and $N\ge100$ in the last.

If a family $(f_v)_{v\in I}$ is not separated, then one of the following holds: some $f_v$ lies in
$S_1$, which covers the failures of (a) and of $f_v(0)\ne0$; some $f_v$ lies in $S_2$; some $f_v$ lies
in $S_3(\xi_{v,j})$ or in $S_3(\xi_{v,j}^{-1})$ for some $j\in J$, which covers the failures of (c);
for some $a$ with $1\le a\le N/2$ every $f_v$ lies in $S_4(a)$, which is the failure of (d); or every
$f_v$ lies in $S_5$, which is the failure of (e).

Under the product measure $w^{\otimes I}$, the families with $f_v\in S$ for one given $v$ have measure
$w(S)$, because $w$ has total mass $1$ at the other places, and the families with $f_v\in S$ for every
$v$ have measure $w(S)^{|I|}$. Summing over the cases above, the measure of the families that are not
separated is at most
\[
 |I|\,\frac2g+|I|\,\frac{23\log g}{\sqrt g}
 +2|I||J|\Big(\frac1{4\min(q_1,d)}+\frac{1+\log N}{g-1}\Big)
 +\Big\lfloor\frac N2\Big\rfloor\Big(\frac78\Big)^{|I|}+\Big(1-\frac1{5\log N}\Big)^{|I|}.
\]
Finally $(1-t)^{|I|}\le e^{-t|I|}$ for $t=1/(5\log N)$, which gives the stated bound.
\end{proof}
\leanref{Formalism/PadicSubsetSums/StableGradedCounts/LocalEstimateIVOfTypeSetWeight.lean}{Formalism.IrreducibleCountingLemmas.GaloisIrreducibility.propW_of}

\begin{lemma}\label{cv-local-share}
For every $\varepsilon$ and every family as in Theorem~\ref{cv-covolume-bound} there is $x_0$ such
that at every point at a scale $x\ge x_0$, in the notation of Situation~\ref{cv-point} and
Definition~\ref{cv-separated}, with $\psi_k(a)=e^{2\pi ia/g_k}$ and $S_k$ the set of tuples of labels
at $g_k$ that are not separated,
\[
 s_0(S_0)+s_1(S_1)\le\exp\big(-x^{\varepsilon/(7J)}\big).
\]
\end{lemma}
\leanref{Development/AutomorphicArithmetic/ThetaCorrespondence/ThetaPairLifts/ThirdPlaceSeries/LocalShareCountLeExp.lean}{GlobalTheory.CovTop.LocShareNum.local_share_count_le_exp_alt}
\leanref{Development/AbelianVarietyTheory/Proofs/DivisorProofs/LocalShareSumBound.lean}{GlobalTheory.CovTop.LocShareNum.locShare_eUk_le_Wg}
\begin{proof}
Write $\varepsilon_z=\varepsilon/J$, let $q$ and $d$ be the auxiliary primes of the host, and fix $k$;
put $g=g_k$.

Step 1: the weights. Write $U$ for the upper unitriangular subgroup of $GL_r(\mathbf F_g)$ and
$e_U=e_U(\psi_k)$. The weights $\operatorname{wt}_{\psi_k}$ are nonnegative and sum to $1$ over the
labels. Indeed, each $e(f)$ is central, so $e_Ue(f)$ is an idempotent of $\mathbf C[GL_r(\mathbf F_g)]$,
and its coefficient at $1$ is its rank on the regular representation divided by
$|GL_r(\mathbf F_g)|$, which is real and nonnegative; and $e_U(1)=|U|^{-1}>0$. By
Corollary~\ref{lit-green-labelled} every irreducible character has exactly one label, so
$\sum_fe(f)=1$ and $\sum_f(e_Ue(f))(1)=e_U(1)$. Next, a squarefree label $f=P_1\cdots P_s$ has
weight $\prod_i(g^{\deg P_i}-1)^{-1}$. The Harish-Chandra induction of
$\tau_{P_1}\otimes\cdots\otimes\tau_{P_s}$ with distinct $P_i$ is the only character $\eta$ with
label $f$; it is irreducible and generic of degree
$\prod_{i=1}^{r}(g^i-1)/\prod_i(g^{\deg P_i}-1)$ (Theorem~\ref{lit-bundle-theta}(1)), so it occurs
once in the Gelfand--Graev representation $\mathbf C[GL_r(\mathbf F_g)]e_U$, which is multiplicity
free \cite{Gel70,Zel81}. Hence $(e_Ue(f))(1)=\eta(1)/|GL_r(\mathbf F_g)|$, and
$|GL_r(\mathbf F_g)|=|U|\prod_{i=1}^r(g^i-1)$.

Step 2: the share of the tuples that are not separated is the
$\operatorname{wt}^{\otimes f}$-measure of the families of labels that are not separated. Taking for
$w$ the weights of Step 1, extended by $0$ to the monic polynomials of degree $r$ that are not
labels, Lemma~\ref{cv-property-w} applies with $K=\mathbf F_g$, $N=r$, $q_1=q$, $I$ the set of the
$f$ level places over $g$, $J=\{1,\dots,m\}$ and $\xi_{v,j}=\chi_{k,v,j}$. Its hypotheses hold for the
datum. The primes $q$ and $d$ are odd and distinct, since $q\equiv3\pmod8$ and $d\equiv7\pmod{16}$
(Definition~\ref{hf-framed-host-field}). Both exceed $4tr+1\ge16$ (Definition~\ref{hf-host}). The
integer $Q_S$ of the host field is divisible by $2^{e+2}$ with $e\ge2$, so by $16$, and $g\equiv1$
modulo $Q_Sqd$, a condition on the level primes $g_0,g_1$ that is part of the definite-side datum
(compare Definition~\ref{hf-good-host}(2)); so $16qd$ divides $g-1$.
Finally $g\ge2^{12}$ and $r\ge100$ by the sizes in Step 3, since $8x^{\varepsilon_z}\ge12$ once
$x^{\varepsilon_z}\ge3/2$. So the share is at most the bound of Lemma~\ref{cv-property-w}, which is
$W_g-f/g$ for
\[
 W_g=\frac{3f}{g}+\frac{23f\log g}{\sqrt g}
 +2fm\Big(\frac1{4\min(q,d)}+\frac{1+\log r}{g-1}\Big)+\frac r2\Big(\frac78\Big)^f
 +e^{-f/(5\log r)}.
\]

Step 3: sizes. At every point at a large scale the size conditions of the family give
$g_k\ge\exp(8x^{\varepsilon_z})$, $x^2\le f\le x^6$, $m\le x^6$, $100\le r\le x^6$ and
$\min(q,d)\ge\exp(x^{\varepsilon_z/6})$. A computation with these bounds gives
$W_{g_k}\le x^{21}\exp(-x^{\varepsilon_z/6})$ for $x$ large, and
$x^{21}\exp(-x^{\varepsilon_z/6})\le\frac1{20}\exp(-x^{\varepsilon_z/7})$ for $x$ large, because
$x^{\varepsilon_z/6}\ge3x^{\varepsilon_z/7}$ and $21\log x\le x^{\varepsilon_z/6}/3$ eventually. Hence
$W_{g_0}+W_{g_1}\le\exp(-x^{\varepsilon_z/7})$.
\end{proof}
\leanref{Development/AbelianVarietyTheory/Proofs/InertiaEigenvalues/LemmaLSeriesWeightsOfGreenCuspidalParametrisatio.lean}{GlobalTheory.GLq.lemmaL_series_weights_of_greenCuspidalParametrisation}
\leanref{Formalism/LocalToGlobal/GlobalTheory/ConcNWeight.lean}{Formalism.L2.GlobalTheoryTheory.ConcNWeight.host_facts}
\leanref{Formalism/LocalToGlobal/GlobalTheory/ConcNWeight.lean}{Formalism.L2.GlobalTheoryTheory.ConcNWeight.rangeW_of_sizes}
\leanref{Development/AbelianVarietyTheory/Proofs/DivisorProofs/LocalShareSumBound.lean}{GlobalTheory.CovTop.LocShareNum.sizesRec_at_point}
\leanref{Formalism/LocalToGlobal/GlobalTheory/ConcNP.lean}{Formalism.L2.GlobalTheoryTheory.ConcN.weight_numbers_bound}
\leanref{Formalism/PadicSubsetSums/StableGradedCounts/DominantNumbersEighthOfSharp.lean}{Formalism.IrreducibleCountingLemmas.GaloisIrreducibility.dominant_numbers_sharp_holds}
\leanref{Development/AbelianVarietyTheory/Proofs/DivisorProofs/LocalShareSumBound.lean}{GlobalTheory.CovTop.LocShareNum.wg_sum_le}

\begin{lemma}\label{cv-lifting-map}
Let $K\subset L$ be number fields, let $V$ be a finite-dimensional $L$-vector space, and let
$W\subset V$ be a $K$-subspace with a finite $K$-basis that is linearly independent over $L$. For each
finite place $\lambda$ of $K$ let $P_\lambda$ be a set and, for $y\in P_\lambda$, let
$\nu_{\lambda,y}\colon W\to\mathbf R$ be a function; let $W^{\mathrm{int}}$ be the set of $w\in W$ with
$\nu_{\lambda,y}(w)\le1$ for all $\lambda$ and all $y\in P_\lambda$. Let $C$ be a $K$-vector space with,
for each $\lambda$, a set $P'_\lambda$ and functions $\nu'_{\lambda,z}\colon C\to\mathbf R$ for
$z\in P'_\lambda$; let $C^{\mathrm{int}}$ be the set of $f\in C$ with $\nu'_{\lambda,z}(f)\le1$ for all
$\lambda$ and all $z\in P'_\lambda$. Let $e\colon C\to W$ be a $K$-linear isomorphism and let
$\pi_\lambda\colon P_\lambda\to P'_\lambda$ be bijections with
$\nu_{\lambda,y}(e(f))=\nu'_{\lambda,\pi_\lambda(y)}(f)$ for all $\lambda$, all $y\in P_\lambda$ and all
$f\in C$. Let $\iota\colon C\to V$ be $e$ followed by the inclusion of $W$ in $V$. Then:
\begin{enumerate}
\item $\iota$ is additive and $\iota(af)=a\,\iota(f)$ for $a\in K$ and $f\in C$;
\item $\iota$ carries every $K$-linearly independent family in $C$ to an $L$-linearly independent
  family in $V$;
\item $\iota(C^{\mathrm{int}})=W^{\mathrm{int}}$.
\end{enumerate}
\end{lemma}
\leanref{Development/AbelianVarietyTheory/Proofs/DivisorProofs/ExistsSemilinearEmbSupLatticeOntoRealisation.lean}{GlobalTheory.CovTop.SupImage.exists_emb_image_latW}
\begin{proof}
(1) The map $e$ is $K$-linear, and the $K$-vector space structure of $V$ is the restriction of its
$L$-structure.

(2) Let $b_1,\dots,b_s$ be the given $K$-basis of $W$, and let $\phi\colon L\otimes_KW\to V$ be the
$L$-linear map with $\phi(a\otimes w)=aw$. The elements $1\otimes b_k$ form an $L$-basis of
$L\otimes_KW$, and $\phi$ sends them to the $b_k$, which are linearly independent over $L$; so $\phi$
is injective. Let $f_1,\dots,f_t\in C$ be linearly independent over $K$. Since $e$ is injective, the
$c_k=e(f_k)$ are linearly independent over $K$. Extending them to a $K$-basis of $W$ shows that the
$1\otimes c_k$ are linearly independent over $L$. Applying $\phi$, the $\iota(f_k)=c_k$ are linearly
independent over $L$.

(3) Let $f\in C$. Since each $\pi_\lambda$ is a bijection, $f\in C^{\mathrm{int}}$ if and only if
$\nu'_{\lambda,\pi_\lambda(y)}(f)\le1$ for all $\lambda$ and all $y\in P_\lambda$, that is, if and only
if $e(f)\in W^{\mathrm{int}}$. So $\iota(C^{\mathrm{int}})\subset W^{\mathrm{int}}$. Conversely, an
element of $W^{\mathrm{int}}$ is $e(f)$ for some $f\in C$, since $e$ is onto $W$, and this $f$ lies
in $C^{\mathrm{int}}$.
\end{proof}
\leanref{Development/AbelianVarietyTheory/Proofs/DivisorProofs/ExistsSemilinearEmbSupLatticeOntoRealisation.lean}{GlobalTheory.CovTop.SupImage.linearIndependent_of_form}

\begin{lemma}\label{cv-lifting-rational}
Let $K\subset L$ be number fields and $n\ge0$, and regard $K^n$ as a subset of $L^n$. For a subset
$X\subset L^n$ write $\mathcal O_LX$ for the $\mathcal O_L$-submodule and $LX$ for the $L$-subspace
of $L^n$ that $X$ generates.
\begin{enumerate}
\item Let $M\subset K^n$ be an $\mathcal O_K$-submodule and $U\subset K^n$ a $K$-subspace. Then
  $\mathcal O_LM\cap LU=\mathcal O_L(M\cap U)$.
\item Let $(A_p)_{p\in I}$ be a family of matrices in $M_n(K)$ and $(\theta_p)_{p\in I}$ a family of
  elements of $K$. Then the set of $y\in L^n$ with $A_py=\theta_py$ for every $p\in I$ is $LU$, where
  $U$ is the set of $x\in K^n$ with $A_px=\theta_px$ for every $p\in I$.
\end{enumerate}
\end{lemma}
\leanref{Development/AbelianVarietyTheory/Proofs/DivisorProofs/LatticeSpanInfRationalSubspaceIntegersExtStatement.lean}{GlobalTheory.T.lattice_span_inf_rational_subspace_integers_ext}
\leanref{Development/AbelianVarietyTheory/Proofs/DivisorProofs/RationalMatrixEigenspaceImageSpannedRationalStatement.lean}{GlobalTheory.T.rational_matrix_eigenspace_image_spanned_rational}
\begin{proof}
(1) The right side lies in the left. For the converse, $\mathcal O_L$ is a finitely generated
projective $\mathcal O_K$-module, so by the dual basis lemma there are $\omega_1,\dots,\omega_s\in\mathcal O_L$
and $\mathcal O_K$-linear maps $\phi_i\colon\mathcal O_L\to\mathcal O_K$ with $a=\sum_i\phi_i(a)\omega_i$
for $a\in\mathcal O_L$. Each $\phi_i$ extends to a $K$-linear map $L\to K$, which has the form
$a\mapsto\operatorname{Tr}_{L/K}(\eta_ia)$ for some $\eta_i\in L$ with
$\operatorname{Tr}_{L/K}(\eta_i\mathcal O_L)\subset\mathcal O_K$; and
$a=\sum_i\operatorname{Tr}_{L/K}(\eta_ia)\,\omega_i$ for all $a\in L$. Let $F_i\colon L^n\to K^n$ apply
$a\mapsto\operatorname{Tr}_{L/K}(\eta_ia)$ to each coordinate. For $c\in L$ and $x\in K^n$,
$F_i(cx)=\operatorname{Tr}_{L/K}(\eta_ic)\,x$. Hence $F_i$ maps $\mathcal O_LM$ into $M$ and $LU$ into
$U$. For $y$ in the intersection, each $F_i(y)$ lies in $M\cap U$, and
$y=\sum_i\omega_iF_i(y)$ lies in $\mathcal O_L(M\cap U)$.

(2) The elements of $LU$ satisfy the equations. Conversely, let $y\in L^n$ satisfy them, choose a basis
$(\beta_k)_k$ of $L$ over $K$, and write $y=\sum_k\beta_ky_k$ with $y_k\in K^n$. Since the $A_p$ and
$\theta_p$ have entries in $K$, comparing the coefficients of $\beta_k$ in $A_py=\theta_py$ gives
$A_py_k=\theta_py_k$ for every $k$. So each $y_k$ lies in $U$, and $y\in LU$.
\end{proof}
\leanref{Development/AbelianVarietyTheory/Proofs/DivisorProofs/LatticeSpanInfRationalSubspaceIntegersExt.lean}{GlobalTheory.span_inf_span_eq_span_inf_of_rational}
\leanref{Development/AbelianVarietyTheory/Proofs/DivisorProofs/DifferentIdealSmulMemSpanRational.lean}{GlobalTheory.exists_trace_dual_system}
\leanref{Development/AbelianVarietyTheory/Proofs/DivisorProofs/RationalMatrixEigenspaceImageSpannedRational.lean}{GlobalTheory.rational_eigen_image_span}

\begin{lemma}\label{cv-lifting-type}
Let $K\subset L$ be number fields, let $V$ be a finite-dimensional $L$-vector space, and let $W\subset V$
be a $K$-subspace with a finite $K$-basis that is an $L$-basis of $V$. Let $M\subset W$ be an
$\mathcal O_K$-submodule. Let $J$ be a group acting on $V$ by $L$-linear maps $\rho(u)$ with
$\rho(u)W\subset W$ for $u\in J$, let $\theta\colon J\to K$ be a function, and let $V^\theta$ be the set
of $v\in V$ with $\rho(u)v=\theta(u)v$ for every $u\in J$. Then
\[
 \mathcal O_LM\cap V^\theta=\mathcal O_L(M\cap V^\theta),
\]
where $\mathcal O_LX$ is the $\mathcal O_L$-submodule of $V$ generated by $X$.
\end{lemma}
\leanref{Development/AbelianVarietyTheory/Proofs/DivisorProofs/IntegralFormsEqSpanIntegralRationalFormsStatement.lean}{GlobalTheory.T.integralForms_eq_span_integralRationalForms}
\begin{proof}
The right side lies in the left, since $V^\theta$ is an $L$-subspace. Let $b_1,\dots,b_n$ be the given
basis, and let $\beta\colon L^n\to V$ be $y\mapsto\sum_ky_kb_k$, an $L$-linear isomorphism that maps
$K^n$ onto $W$. Since each $\rho(u)b_k$ lies in $W$, there are matrices $A_u\in M_n(K)$ with
$\rho(u)\beta(y)=\beta(A_uy)$ for $y\in L^n$. Hence $\beta^{-1}(V^\theta)$ is the set of $y\in L^n$ with
$A_uy=\theta(u)y$ for every $u\in J$, which is $LU$ by Lemma~\ref{cv-lifting-rational}(2), where $U$ is
the set of $x\in K^n$ with $A_ux=\theta(u)x$ for every $u$. Let $M_0=\beta^{-1}(M)\cap K^n$, an
$\mathcal O_K$-submodule of $K^n$ with $\beta(M_0)=M$, since $M\subset W$; so
$\beta^{-1}(\mathcal O_LM)=\mathcal O_LM_0$. Let $v\in\mathcal O_LM\cap V^\theta$. Then $\beta^{-1}(v)$ lies
in $\mathcal O_LM_0\cap LU$, which is $\mathcal O_L(M_0\cap U)$ by
Lemma~\ref{cv-lifting-rational}(1). For $x\in M_0\cap U$, $\beta(x)$ lies in $M$ and in $V^\theta$.
Hence $v\in\mathcal O_L(M\cap V^\theta)$.
\end{proof}
\leanref{Development/AbelianVarietyTheory/Proofs/DivisorProofs/IntegralFormsEqSpanIntegralRationalForms.lean}{GlobalTheory.mall_eq_span_rat}

\begin{lemma}\label{cv-realization-linear}
Let $\sigma\colon L\to M$ be a homomorphism of fields, let $W$ be a finite-dimensional $L$-vector
space, let $G$ be a finite group with a representation $\rho_0$ on $W$ and a character
$\chi\colon G\to L^\times$, let $P\subset G$ be a subgroup with a character
$\lambda_0\colon P\to L^\times$, and let $c\colon G\to L$ be a function. Put
\[
 W^{\mathrm{type}}=\{w\in W:\rho_0(p)w=\chi(p)\lambda_0(p)w\text{ for every }p\in P\},\qquad
 e=\sum_{g\in G}c(g)\chi(g)^{-1}\rho_0(g),
\]
and $e'=\sum_g\sigma(c(g))\,g\in M[G]$. For a representation $\rho$ of a group on an $M$-vector space
$V$, a subgroup $Q$ and a character $\mu$ of $Q$, let $V^\mu$ be the subspace of $V$ on which $Q$ acts
through $\mu$.
\begin{enumerate}
\item Let $V=M\otimes_LW$, with $M$ an $L$-algebra through $\sigma$, let $\rho$ be the
  representation of $G$ on $V$ with $\rho(g)=\sigma(\chi(g))^{-1}(1\otimes\rho_0(g))$, and let
  $\lambda=\sigma\circ\lambda_0$ on $Q=P$. Then
  \[
   \dim_LW^{\mathrm{type}}=\dim_MV^\lambda,\qquad
   \dim_L(W^{\mathrm{type}}\cap\ker e)=\dim_M\big(V^\lambda\cap\ker\rho(e')\big).
  \]
\item Let $K\subset L$ and $F\subset M$ be subfields and $j\colon K\to F$ a homomorphism of fields
  with $\sigma(a)=j(a)$ for $a\in K$. Let $C$ be a $K$-vector space with a finite basis $(b_k)_k$ and
  let $\varepsilon\colon C\to W$ be a $K$-linear map such that $(\varepsilon(b_k))_k$ is a basis of $W$
  over $L$. Let $V_0$ be an $F$-vector space and $\beta\colon C\to V_0$ an additive map with
  $\beta(af)=j(a)\beta(f)$ for $a\in K$, such that $(\beta(b_k))_k$ is linearly independent over $F$.
  Let $\Gamma$ be a finite group with a representation $\tau$ on $V_0$, and let $\pi\colon\Gamma\to G$
  and $\iota\colon G\to\Gamma$ be homomorphisms with $\pi\circ\iota=\mathrm{id}$. Suppose that for all
  $\gamma\in\Gamma$ and $f\in C$ there is $f'\in C$ with
  $\rho_0(\pi(\gamma))\varepsilon(f)=\varepsilon(f')$ and $\tau(\gamma)\beta(f)=\beta(f')$, and that
  every vector of $V_0$ fixed by $\ker\pi$ lies in the $F$-span of $\beta(C)$. Let $V=M\otimes_FV_0$
  with the representation $\rho(\gamma)=\sigma(\chi(\pi(\gamma)))^{-1}(1\otimes\tau(\gamma))$ of
  $\Gamma$, let $\Gamma_P=\pi^{-1}(P)$, and let $\tilde\lambda=\sigma\circ\lambda_0\circ\pi$ on
  $\Gamma_P$. Then
  \[
   \dim_LW^{\mathrm{type}}=\dim_MV^{\tilde\lambda},\qquad
   \dim_L(W^{\mathrm{type}}\cap\ker e)=\dim_M\big(V^{\tilde\lambda}\cap\ker\rho(\iota(e'))\big),
  \]
  where $\iota(e')=\sum_g\sigma(c(g))\,\iota(g)\in M[\Gamma]$.
\end{enumerate}
\end{lemma}
\leanref{Development/AbelianVarietyTheory/Proofs/DivisorProofs/FinrankCharacterEigenspaceTwistedComplexificatio.lean}{GlobalTheory.DefRead.finrank_eig_twBC}
\leanref{Development/AbelianVarietyTheory/Proofs/DivisorProofs/FinrankIInfKerSemilinear.lean}{GlobalTheory.BadCount.RankRead.finrank_eq}
\begin{proof}
(1) For $p\in P$ let $a_p=\rho_0(p)-\chi(p)\lambda_0(p)$, so that $W^{\mathrm{type}}=\bigcap_p\ker a_p$.
Then $\rho(p)-\lambda(p)=\sigma(\chi(p))^{-1}(1\otimes a_p)$, so $V^\lambda=\bigcap_p\ker(1\otimes a_p)$.
Also $\rho(e')=\sum_g\sigma(c(g))\sigma(\chi(g))^{-1}(1\otimes\rho_0(g))=1\otimes e$. Since $M$ is flat
over $L$, extending scalars commutes with kernels: the kernel of the base change of
$a=(a_p)_p\colon W\to\bigoplus_{p\in P}W$ is $M\otimes_L\ker a$, and the same holds for the map
$(e,a)$. Taking dimensions gives (1).

(2) Let $\Phi\colon W\to V$ be the additive map with $\Phi(aw)=\sigma(a)\Phi(w)$ for $a\in L$ and
$\Phi(\varepsilon(b_k))=1\otimes\beta(b_k)$ for every $k$. As $\varepsilon$ is $K$-linear,
$\beta(af)=j(a)\beta(f)$ and $\sigma$ agrees with $j$ on $K$, $\Phi(\varepsilon(f))=1\otimes\beta(f)$
for every $f\in C$. The vectors $1\otimes\beta(b_k)$ are linearly independent over $M$, because the
$\beta(b_k)$ are linearly independent over $F$ and $M$ is flat over $F$. So the map $M\otimes_LW\to V$,
$z\otimes w\mapsto z\Phi(w)$, is injective, and its image $U$ is the $M$-span of $\Phi(W)$.

For $\gamma\in\Gamma$ and $w\in W$,
\[
 \rho(\gamma)\Phi(w)=\Phi\big(\chi(\pi(\gamma))^{-1}\rho_0(\pi(\gamma))w\big).
\]
Both sides are additive in $w$ and change by $\sigma(a)$ when $w$ is multiplied by $a\in L$, so it
suffices to take $w=\varepsilon(b_k)$. With $f'$ as in the hypothesis for $\gamma$ and $b_k$, both sides
then equal $\sigma(\chi(\pi(\gamma)))^{-1}(1\otimes\beta(f'))$. For $\gamma\in\Gamma_P$ let
$a_\gamma=\chi(\pi(\gamma))^{-1}\rho_0(\pi(\gamma))-\lambda_0(\pi(\gamma))$ on $W$ and
$a'_\gamma=\rho(\gamma)-\tilde\lambda(\gamma)$ on $V$, so that $a'_\gamma\circ\Phi=\Phi\circ a_\gamma$.
Every $p\in P$ is $\pi(\iota(p))$ with $\iota(p)\in\Gamma_P$, so $W^{\mathrm{type}}=\bigcap_\gamma\ker a_\gamma$;
and $V^{\tilde\lambda}=\bigcap_\gamma\ker a'_\gamma$.

Next, $V^{\tilde\lambda}\subset U$. Let $x\in V^{\tilde\lambda}$ and $\gamma\in\ker\pi$. Then
$\gamma\in\Gamma_P$, $\tilde\lambda(\gamma)=1$ and $\chi(\pi(\gamma))=1$, so $(1\otimes\tau(\gamma))x=x$. As
in (1), the vectors of $M\otimes_FV_0$ fixed by $1\otimes\tau(\gamma)$ for every $\gamma\in\ker\pi$ form the
$M$-span of the $1\otimes v$ with $v\in V_0$ fixed by $\ker\pi$. By hypothesis these $v$ lie in the
$F$-span of $\beta(C)$, so $x$ lies in the $M$-span of the vectors $1\otimes\beta(f)=\Phi(\varepsilon(f))$,
which is contained in $U$.

The $a'_\gamma$ preserve $U$ and act on it as the base changes of the $a_\gamma$. By the argument of
(1), $\dim_L\bigcap_\gamma\ker a_\gamma=\dim_M\big(U\cap\bigcap_\gamma\ker a'_\gamma\big)=\dim_MV^{\tilde\lambda}$,
which is the first equality. For the second, $\pi(\iota(g))=g$ gives
$\rho(\iota(e'))\Phi(w)=\sum_g\sigma(c(g))\rho(\iota(g))\Phi(w)=\Phi(e(w))$, and the same argument applies
with $e$ and $\rho(\iota(e'))$ added to the families $(a_\gamma)$ and $(a'_\gamma)$.
\end{proof}
\leanref{Development/AbelianVarietyTheory/Proofs/DivisorProofs/FinrankIInfKerSemilinear.lean}{GlobalTheory.BadCount.RankRead.finrank_iInf_ker_semilinear}

\begin{lemma}\label{cv-realization-split}
Let $\sigma\colon K\to\mathbf C$ be a homomorphism of fields, and regard $\mathbf C$ as a $K$-algebra
through $\sigma$. Let $B$ be a $K$-algebra that is a finite product of finite separable field
extensions of $K$, and let $W$ be a $B$-module. Let $H$ be a group acting on $B$ by $K$-algebra
automorphisms $\tau(h)$, and for $h\in H$ let $T(h)\colon W\to W$ be a $K$-linear map with
$T(h)(bw)=\tau(h)(b)\,T(h)(w)$ for $b\in B$ and $w\in W$. Let $Y$ be the finite set of $K$-algebra
homomorphisms $B\to\mathbf C$, with $H$ acting by $h\cdot y=y\circ\tau(h)^{-1}$, and suppose that this
action is free. For $y\in Y$ let $A_y$ be the set of $v\in\mathbf C\otimes_KW$ with $(1\otimes b)v=y(b)v$
for every $b\in B$. Then
\[
 \mathbf C\otimes_KW=\bigoplus_{y\in Y}A_y,\qquad (1\otimes T(h))A_y\subset A_{h\cdot y}\quad
 \text{for }h\in H\text{ and }y\in Y.
\]
\end{lemma}
\leanref{Development/AbelianVarietyTheory/Proofs/DivisorProofs/HalfLevelSubmodulePermutation.lean}{GlobalTheory.DefSplit.split_of_free}
\begin{proof}
The operators $1\otimes b$ on $\mathbf C\otimes_KW$ commute, and the $A_y$ are their joint eigenspaces
for distinct characters $y$ of $B$, so the sum of the $A_y$ is direct. The map
$\mathbf C\otimes_KB\to\mathbf C^Y$, $z\otimes b\mapsto(z\,y(b))_y$, is an isomorphism of
$\mathbf C$-algebras. It suffices to check this for a finite separable extension $B=K(\alpha)$. Then
$\mathbf C\otimes_KB=\mathbf C[X]/(m)$, where $m$ is the image under $\sigma$ of the minimal polynomial of
$\alpha$; its roots are distinct and are the $y(\alpha)$, and the Chinese remainder theorem gives the
isomorphism. Let $e_y\in\mathbf C\otimes_KB$ correspond to the unit vector at $y$. Then $\sum_ye_y=1$
and $(1\otimes b)e_y=y(b)e_y$, so every $v\in\mathbf C\otimes_KW$ is the sum of the $e_yv$, and $e_yv$
lies in $A_y$. Finally, let $v\in A_y$, $h\in H$ and $b\in B$. Then
\[
 (1\otimes b)(1\otimes T(h))v=(1\otimes T(h))(1\otimes\tau(h)^{-1}(b))v
 =y\big(\tau(h)^{-1}(b)\big)(1\otimes T(h))v,
\]
so $(1\otimes T(h))v$ lies in $A_{h\cdot y}$.
\end{proof}
\leanref{Development/AbelianVarietyTheory/Proofs/DivisorProofs/HalfLevelSubmodulePermutation.lean}{GlobalTheory.DefSplit.iSup_part}
\leanref{Development/AbelianVarietyTheory/Proofs/DivisorProofs/HalfLevelSubmodulePermutation.lean}{GlobalTheory.DefSplit.indep_part}

\begin{lemma}\label{cv-regular-group-determinant}
Let $K$ be a field with a homomorphism of fields $K\to\mathbf C$, let $R$ be a commutative
$K$-algebra, and let $\Gamma$ be a finite group acting on $R$ by $K$-algebra automorphisms, written
$x\mapsto\gamma\cdot x$. For a subgroup $S\subset\Gamma$ and $u\in R$, the \emph{group determinant} of $u$
on $S$ is $\det(\gamma\gamma'\cdot u)_{\gamma,\gamma'\in S}\in R$, which is defined up to sign. Let
$S_0,S\subset\Gamma$ be subgroups and let $u_0,u\in R$ be elements whose group determinants on $S_0$
and on $S$ are not zero divisors in $R$. Let $V_-,V,V_1,V_2\subset R$ be $K$-subspaces with $V$ and
$V_2$ finite-dimensional, such that
\begin{enumerate}
\item[(a)] $V$ is stable under $\Gamma$, and $V_1$ is stable under $S$;
\item[(b)] $x\,(\gamma\cdot u_0)\in V$ for $x\in V_-$ and $\gamma\in S_0$;
\item[(c)] $x\,(\gamma\cdot u)\in V_1$ for $x\in V$ and $\gamma\in S$, and $x\,(\gamma\cdot u)\in V_2$ for
  $x\in V_1$ and $\gamma\in S$.
\end{enumerate}
Let $\nu\colon\Gamma\to\mathbf C^\times$ be a character, let $\rho$ be the representation of $\Gamma$ on
$\mathbf C\otimes_KV$ with $\rho(\gamma)(z\otimes x)=\nu(\gamma)\,z\otimes(\gamma\cdot x)$, let
$\lambda\colon S_0\to\mathbf C^\times$ be a character, and let $(\mathbf C\otimes_KV)^\lambda$ be the
subspace on which $S_0$ acts through $\lambda$. Then:
\begin{enumerate}
\item there is $n\ge0$ with $n|S|\le\dim_KV_2$ such that the restriction of $\rho$ to $S$ embeds into
  the sum of $n$ copies of the regular representation of $S$;
\item $\dim_KV_-\le|S_0|\dim_{\mathbf C}(\mathbf C\otimes_KV)^\lambda$.
\end{enumerate}
In particular, if $\dim_KV_2\le t\dim_KV_-$ with $t\ge0$, then
$n|S|\le t\,|S_0|\dim_{\mathbf C}(\mathbf C\otimes_KV)^\lambda$.
\end{lemma}
\leanref{Development/AbelianVarietyTheory/Proofs/DivisorProofs/EmbRegCountOfRanks.lean}{GlobalTheory.BadCount.embReg_count_of_ranks}
\leanref{Development/AbelianVarietyTheory/Proofs/DivisorProofs/EmbedsInRegularRepresentationCopiesOfGroupDeterm.lean}{GlobalTheory.BadCount.embReg_of_groupDet}
\leanref{Development/AbelianVarietyTheory/Proofs/DivisorProofs/FinrankLeCardMulFinrankCharacterSpaceBaseChange.lean}{GlobalTheory.BadCount.finrank_le_card_mul_finrank_eig_baseChange}
\begin{proof}
Step 1: group determinants. Let $H\subset\Gamma$ be a subgroup and $v\in R$ an element whose group
determinant $D$ on $H$ is not a zero divisor, and let $N=(\gamma\gamma'\cdot v)_{\gamma,\gamma'\in H}$. If
$(a_\gamma)_{\gamma\in H}\in R^H$ satisfies $\sum_{\gamma'}(\gamma\gamma'\cdot v)\,a_{\gamma'}=0$ for every
$\gamma$, or $\sum_\gamma(\gamma\gamma'\cdot v)\,a_\gamma=0$ for every $\gamma'$, then $a=0$: multiplying by
the adjugate of $N$, or of its transpose, gives $Da_\gamma=0$ for every $\gamma$. The same holds in
$R'=\mathbf C\otimes_KR$, with $\Gamma$ acting on the second factor, for $1\otimes v$: its group
determinant is $1\otimes D$, and multiplication by $1\otimes D$ on $R'$ is the base change of the
injective multiplication by $D$ on $R$, so it is injective.

Step 2: the embedding. Let $I\subset V_1$ be the subspace of elements fixed by $S$. For $h\in S$ and
$x\in V$ put $T_h(x)=\sum_{g\in S}g\cdot(x\,(h\cdot u))$, which lies in $I$ by (a) and (c). The map
$I^S\to V_2$, $(F_h)_h\mapsto\sum_hF_h\,(h\cdot u)$, is injective: if the sum vanishes, applying
$g\in S$ gives $\sum_h(gh\cdot u)F_h=0$ for every $g$, since the $F_h$ are fixed by $S$, and Step 1
gives $F=0$. Hence $|S|\dim_KI\le\dim_KV_2$. Put $n=\dim_KI$. The map $x\mapsto(T_h(x))_h$ is injective
on $V$: if $T_h(x)=\sum_g(gh\cdot u)(g\cdot x)$ vanishes for every $h$, Step 1 gives $g\cdot x=0$ for
every $g$, so $x=0$. For $\gamma\in S$ the substitution $g\mapsto g\gamma^{-1}$ gives
$T_h(\gamma\cdot x)=T_{\gamma^{-1}h}(x)$. Choose coordinate functions $\ell_1,\dots,\ell_n$ for a basis of $I$
and put $J(x)=\big(\sum_{h\in S}\ell_i(T_h(x))\,h\big)_{i\le n}\in K[S]^n$. Then $J$ is injective and
$J(\gamma\cdot x)=\gamma J(x)$ for $\gamma\in S$. Extending scalars, $1\otimes J$ embeds $\mathbf C\otimes_KV$,
with $S$ acting through the second factor, into $\mathbf C[S]^n$ with $S$ acting by left
multiplication. The automorphism $\psi$ of $\mathbf C[S]$ with $\psi(h)=\nu(h)^{-1}h$ satisfies
$\psi(\gamma y)=\nu(\gamma)^{-1}\gamma\psi(y)$, so applying $\psi$ in each coordinate after $1\otimes J$
gives an embedding of the restriction of $\rho$ to $S$ into $\mathbf C[S]^n$. This is (1).

Step 3: the eigenspace. The algebra $R'$ of Step 1 contains $\mathbf C\otimes_KV_-$ and
$\mathbf C\otimes_KV$, as $\mathbf C$ is flat over $K$. Let $\lambda'=\lambda\nu^{-1}$ on $S_0$; then
$(\mathbf C\otimes_KV)^\lambda$ is the set of $x$ with $\gamma\cdot x=\lambda'(\gamma)x$ for $\gamma\in S_0$,
for the action on the second factor. The operator $E=\sum_{g\in S_0}\lambda'(g)^{-1}\,g$ maps
$\mathbf C\otimes_KV$ into $(\mathbf C\otimes_KV)^\lambda$. The map
\[
 \mathbf C\otimes_KV_-\to\big((\mathbf C\otimes_KV)^\lambda\big)^{S_0},\qquad
 y\mapsto\big(E(y\,(h\cdot u_0))\big)_{h\in S_0},
\]
is defined by (b) and is injective: if every component vanishes, then
$\sum_g(gh\cdot u_0)\,\lambda'(g)^{-1}(g\cdot y)=0$ in $R'$ for every $h\in S_0$, and Step 1 gives
$g\cdot y=0$ for every $g$, so $y=0$. Hence
$\dim_KV_-=\dim_{\mathbf C}(\mathbf C\otimes_KV_-)\le|S_0|\dim_{\mathbf C}(\mathbf C\otimes_KV)^\lambda$,
which is (2). The last assertion follows from $n|S|\le\dim_KV_2\le t\dim_KV_-$ and (2).
\end{proof}
\leanref{Development/AbelianVarietyTheory/Proofs/DivisorProofs/TwistedTraceEquivariantAndInjectiveOfGroupDet.lean}{GlobalTheory.twistedTrace_equivariant_and_injective_of_groupDet}

\subsection{Decay at the cusp}\label{subsection-cv-cusp-decay}

\begin{situation}\label{cv-cusp-setting}
Let $\mathcal D$ be a definite datum (Definition~\ref{th-definite-datum}) with ranks $k$ and $m$, host
fields $F\subset E$ with $f=[F:\mathbf Q]$, prime $d$ and Kronecker character $\chi$ of
$\mathbf Q(\sqrt{-d})$. Let $r=k+m-1$, let $1\le j\le m$, and let $W_j$ be the Hermitian space of the
$j$-th indefinite group of $\mathcal D$, of signature $(p,q)=(k+j-1,m-j)$. Let $X$ be the PEL moduli
scheme of $U(W_j)$ at the principal level of an integer $M\ge3$, of full level $1$
(Definition~\ref{cv-model-toroidal-semiabelian}(1)), with coefficient field $F_c$ and with
uniformizations at the complex embeddings of $F_c$, normalized by a real number $\xi>0$, that satisfy the
compatibility conditions of the moduli problem. Let $\Phi$ be the CM type of the embeddings
$\tau:E\to\mathbf C$ that send $\sqrt{-d}$ to the square root of $-d$ fixed by $\mathcal D$. An embedding
$\sigma:F_c\to\mathbf C$ is \emph{good} if it sends the square root of $-d$ in $F_c$ given by the datum
to that root; put $\sigma^g=\sigma$ if $\sigma$ is good and $\sigma^g=\bar\sigma$ otherwise.

\emph{Frames and forms.} At a good $\sigma$ the uniformization gives a finite set of components $Z$, and
for each $Z$ an invertible matrix $G_Z\in M_r(E)$ with $\bar G_{Z,ab}=G_{Z,ba}$, the bar denoting complex
conjugation on $E$ and on $\mathbf C$, and a subgroup $\Gamma_Z$ of the unitary group of $G_Z$ over $E$.
Here the \emph{unitary group} of a Hermitian matrix $H$, over $E$ or over $\mathbf C$, is the group of
invertible $\gamma$ with $\gamma^{\mathsf T}H\bar\gamma=H$: those that preserve the form
$(u,v)\mapsto u^{\mathsf T}H\bar v=\sum_{a,b}u_aH_{ab}\bar v_b$. A \emph{frame} is a family
$x=(x_\tau)_{\tau\in\Phi}$ of matrices in $GL_r(\mathbf C)$. Its \emph{Gram matrix} at $\tau$ is
$x_\tau^{\mathsf T}\,\tau(G_Z)\,\bar x_\tau$, whose $(a,b)$ entry is the value of the form of $\tau(G_Z)$ at
the $a$-th and the $b$-th columns of $x_\tau$; it is Hermitian, so its minors are real. The family $x$ is a
frame of $Z$ if for every $\tau$ and every $0\le i\le r$ the leading principal $i\times i$ minor
$\delta_{\tau,i}(x)$ of $x_\tau^{\mathsf T}\,\tau(G_Z)\,\bar x_\tau$ has the sign $(-1)^{\max(0,i-q)}$, and it
is \emph{orthonormal} if moreover each $x_\tau^{\mathsf T}\tau(G_Z)\bar x_\tau$ is diagonal with entries $1$,
$q$ times, followed by $-1$. The group $\Gamma_Z$ acts on the frames of $Z$ by $(\gamma x)_\tau=\tau(\gamma)x_\tau$, and
the group $B$ of families of invertible upper triangular matrices acts by right multiplication;
$\gamma x\in xB$ only for $\gamma=1$. A \emph{weight} on $Z$ is a family $\lambda=(\lambda_{\tau,i})$ of
integers, $\tau\in\Phi$, $1\le i\le r$; its \emph{spread} is
$\sum_\tau(\max_i\lambda_{\tau,i}-\min_i\lambda_{\tau,i})$, and the uniformization attaches to it a
character $\chi_\lambda$ of $\Gamma_Z$ of finite order. A \emph{form} of weight $\lambda$ on $Z$ is a
holomorphic function $P$ on the frames of $Z$ with $P(xb)=\prod_{\tau,i}b_{\tau,ii}^{\lambda_{\tau,i}}P(x)$
for $b\in B$ and $P(\gamma x)=\chi_\lambda(\gamma)P(x)$ for $\gamma\in\Gamma_Z$. Its \emph{pointwise norm}
\[
 |P|(x)=|P(x)|\prod_{\tau\in\Phi}\prod_{i=1}^r\Big|\frac{\delta_{\tau,i}(x)}{\delta_{\tau,i-1}(x)}\Big|^{-\lambda_{\tau,i}/2}
\]
is invariant under $B$, and under $\Gamma_Z$ when $\chi_\lambda$ is trivial. A \emph{cusp form} is a form
$P$ with $|P|$ bounded whose constant Fourier--Jacobi coefficient along every isotropic vector vanishes.

\emph{Charts and fundamental sets.} For a frame $x$ of $Z$ and a family $z=(z_{\tau,ab})$ of complex
numbers, $\tau\in\Phi$ and $1\le b<a\le r$, let $n(z)$ be the family of lower unitriangular matrices with
entries $z_{\tau,ab}$ below the diagonal. Let $\Omega_x$ be the open set of those $z$ for which $xn(z)$ is
a frame of $Z$, and let $\mu_x$ be the measure on $\Omega_x$ with density
$\prod_\tau\prod_{i=0}^{r-1}|\delta_{\tau,i}(xn(z))|^{-2}$ with respect to the Lebesgue measure. A
\emph{fundamental set} for $\Gamma_Z$ at $x$ is a measurable set $\mathcal F\subset\Omega_x$ such that
$xn(z')\in\Gamma_Z\,xn(z)B$ with $z,z'\in\mathcal F$ only if $z=z'$, and such that for almost every
$z\in\Omega_x$ there is $z'\in\mathcal F$ with $xn(z)\in\Gamma_Z\,xn(z')B$.

\emph{The adelic quotient.} The uniformization at $\sigma$ also gives a map $g\mapsto(Z(g),x(g))$ from
$U(W_j)(\mathbf A_F)$ to frames of components, with $|\delta_{\tau,i}(x(g))|=1$ for all $\tau$ and $i$,
through which every function $h$ on frames that is invariant under the groups $\Gamma_Z$ and $B$ gives a
function $g\mapsto h(x(g))$ on $U(W_j)(F)\backslash U(W_j)(\mathbf A_F)$. Let $\mu$ be the probability
measure on this quotient induced by a Haar measure. For a family $P=(P_Z)_Z$ of forms of weights
$\lambda_Z$ with trivial characters put
\[
 \sup|P|=\sup_Z\sup_x|P_Z|(x),\qquad\|P\|_{\mathrm{Pet}}^2=\int|P_{Z(g)}|(x(g))^2\,d\mu(g).
\]

\emph{Forms on $X$.} A \emph{dominant weight} is a pair $\kappa=(\alpha,n)$ of a family of nonnegative
integers $\alpha_{\tau,i}$, for the embeddings $\tau:E\to\mathbf C$ and $0\le i<m_\tau$, where $m_\tau=p$
for $\tau\in\Phi$ and $m_\tau=q$ otherwise, nonincreasing in $i$, and an even integer $n\ge0$. Put
$|\alpha|=\sum_{\tau,i}\alpha_{\tau,i}$ and $w_\kappa=|\alpha|+frn+1$. At a good $\sigma$, $\kappa$ gives on
every component the weight $\lambda(\kappa)$ with $\lambda(\kappa)_{\tau,i}=\alpha_{\tau',i-q-1}+n/2$ for
$i>q$ and $\lambda(\kappa)_{\tau,i}=-\alpha_{\bar\tau',q-i}-n/2$ for $i\le q$, where $\tau'$ and $\bar\tau'$
are the embeddings of $E$ that the uniformization at $\sigma$ matches with $\tau$ and with $\bar\tau$; the
compatibility conditions include that $\chi_{\lambda(\kappa)}$ is trivial on every $\Gamma_Z$. Let
$S_\kappa$ be the $F_c$-vector space of rational sections $f$ of the automorphic line bundle of weight
$\kappa$ on $X$ whose reading $f_\sigma=(f_{\sigma,Z})_Z$ in the uniformization at every good $\sigma$ is a
family of cusp forms of weight $\lambda(\kappa)$. For $f\in S_\kappa$ and every embedding
$\sigma:F_c\to\mathbf C$, the \emph{supremum norm} and the \emph{Petersson norm} of $f$ at $\sigma$ are
\[
 \|f\|_{\sigma,\infty}=(2\pi\xi)^{-|\alpha|/2}\sup|f_{\sigma^g}|,\qquad
 \|f\|_\sigma=(2\pi\xi)^{-|\alpha|/2}\|f_{\sigma^g}\|_{\mathrm{Pet}}.
\]
The form $f$ is \emph{integral} if its valuation is nonnegative at every integral point of $X$, algebraic
over $F_c$, with values in a complete discrete valuation ring above a finite place of $F_c$; the integral
forms make up an $\mathcal O_{F_c}$-submodule $\Lambda_\kappa$ of $S_\kappa$.

\emph{Constants.} For $w>0$ put $N=fr(r-1)/2$ and
\[
 \theta_w=\min\Big(1,\frac{\pi}{2M\sqrt d\,(fw+2)\log2}\Big),\qquad
 \rho_w=\frac1{200r^3(w+f)},\qquad
 \ell_w=\Big(\frac{1600}3f^{7/2}r^4w^{1/2}|d_E|\,\theta_w^{-1/2}\Big)^{2q+1},
\]
\[
 C^{\mathrm{loc}}_w=\frac{8\,N!}{(\pi\rho_w^2)^N}\Big(1+\frac{8f^2r^2|d_E|}{\theta_w}\Big)^{2fr^2},\qquad
 C^{\mathrm{mass}}=(15M)^{2fr^2}\big(12f^2r^4\operatorname{rd}_E^2\big)^{6fr^2},\qquad
 C_w=C^{\mathrm{loc}}_wC^{\mathrm{mass}},
\]
and for a dominant weight $\kappa$, with $\ell=\ell_{w_\kappa}$ and $\log\Omega_E(\chi)$ as in
Definition~\ref{cv-log-omega},
\[
 \eta^{\sup}_\kappa=\Big(\frac{1+|\log\Omega_E(\chi)|}2+\frac12\log\ell\Big)(|\alpha|+frn)
 +|\alpha|\log(r\ell),\qquad
 \eta_\kappa=\eta^{\sup}_\kappa+\frac12\log C_{w_\kappa}.
\]
\end{situation}
\leanref{Development/AbelianVarietyTheory/Definitions/Divisors/UnifOfSpaceEq.lean}{GlobalTheory.LevGeo.Laws}
\leanref{ReciprocityTheory/ReductiveGroups/ArithmeticGeometry/FundamentalDomains/Signature.lean}{ReciprocityTheory.ChartRestriction.G.FD.Data}
\leanref{ReciprocityTheory/ReductiveGroups/ArithmeticGeometry/FundamentalDomains/Signature.lean}{ReciprocityTheory.ChartRestriction.G.FD.Data.pnorm}
\leanref{ReciprocityTheory/ReductiveGroups/ArithmeticGeometry/ArakelovVolumes/UpperBounds/FundSetMeasStatement.lean}{ReciprocityTheory.ChartRestriction.G.AR.Up.IsFundSet}
\leanref{Development/AbelianVarietyTheory/Definitions/Foundations/IsProbabilityMeasureMuPart.lean}{GlobalTheory.AllWt.cuspW}
\leanref{Development/AbelianVarietyTheory/Definitions/Foundations/IsProbabilityMeasureMuPart.lean}{GlobalTheory.AllWt.latW}
\leanref{Development/AbelianVarietyTheory/Definitions/Foundations/IsProbabilityMeasureMuPart.lean}{GlobalTheory.AllWt.normW}
\leanref{Formalism/LocalToGlobal/ProjectiveLattices/L2DefiniteCuspForms.lean}{Formalism.L2.Pj.supC}
\leanref{Development/WeilBound/RepresentationTheoryGL/TubeMeasures/DisjointTubes/ConstantSupremum/SupremumConstantLogBound.lean}{MainStatements.Tam.Plug.supConst2}
\leanref{Development/WeilBound/RepresentationTheoryGL/TubeMeasures/DisjointTubes/MassLev2.lean}{MainStatements.Tam.massLev2}
\leanref{Development/AbelianVarietyTheory/Definitions/Divisors/UnifOfSpaceEq.lean}{GlobalTheory.LevGeo.locConst}
\leanref{Development/AbelianVarietyTheory/Definitions/Divisors/UnifOfSpaceEq.lean}{GlobalTheory.LevGeo.cmLen}
\leanref{Development/AbelianVarietyTheory/Definitions/Divisors/UnifOfSpaceEq.lean}{GlobalTheory.LevGeo.lamSup}
\leanref{Development/AutomorphicArithmetic/ThetaCorrespondence/FockModels/TamagawaPeterssonLambdaSqNonneg.lean}{UpperTheory.Lit.Tam.lamPet2}

\begin{definition}\label{cv-decay-majorant}
In Situation~\ref{cv-cusp-setting}, let $\sigma$ be a good embedding and $Z$ a component at $\sigma$, and
let $c$ be the complex conjugation of $E$. For $v,v'\in E^r$ put
$G_Z(v,v')=\sum_{a,b}v_a\,G_{Z,ab}\,c(v'_b)$, and let $\tau(v)\in\mathbf C^r$ be the vector of the
$\tau(v_a)$. For $\tau\in\Phi$ and $u,u'\in\mathbf C^r$ put
$\tau(G_Z)(u,u')=\sum_{a,b}u_a\,\tau(G_{Z,ab})\,\overline{u'_b}$, so that the Gram matrix
$x_\tau^{\mathsf T}\tau(G_Z)\bar x_\tau$ of a frame $x$ has the entries $\tau(G_Z)(x_{\tau,a},x_{\tau,b})$, where
$x_{\tau,a}$ is the $a$-th column of $x_\tau$.
\begin{enumerate}
\item The \emph{majorant} of a frame $x$ of $Z$ at $\tau\in\Phi$ is the positive definite Hermitian form
  \[
   h_{x,\tau}(u)=\tau(G_Z)(u^+,u^+)-\tau(G_Z)(u^-,u^-),
  \]
  where $u=u^++u^-$ with $u^+$ in the span $U^+_{x,\tau}$ of the first $q$ columns of $x_\tau$ and $u^-$
  orthogonal to $U^+_{x,\tau}$ for $\tau(G_Z)$.
\item Let $v\in E^r$ with $G_Z(v,v)=0$. For $s=(s_\tau)_{\tau\in\Phi}\in\mathbf C^\Phi$, let $n_v(s)$ be the
  family of the matrices $n_v(s)_\tau$ with
  \[
   n_v(s)_\tau u=u-2i\xi s_\tau\,\tau(G_Z)(u,\tau(v))\,\tau(v)\qquad(u\in\mathbf C^r),
  \]
  and for a frame $x$ let $n_v(s)x$ be the family $(n_v(s)_\tau x_\tau)_\tau$. For $\zeta\in\mathbf C$,
  $n_v(\zeta)$ denotes $n_v(s)$ with $s_\tau=\zeta$ for every $\tau$.
\item Let $\lambda$ be a weight on $Z$. Its \emph{spread at} $\tau$ is
  $\operatorname{sp}_\tau(\lambda)=\max_i\lambda_{\tau,i}-\min_i\lambda_{\tau,i}$, so that its spread is
  $\sum_\tau\operatorname{sp}_\tau(\lambda)$. The \emph{normalizing factor} of a frame $x$ of $Z$ is
  \[
   \beta_\lambda(x)=\prod_{\tau\in\Phi}\prod_{i=1}^r
   \Big|\frac{\delta_{\tau,i}(x)}{\delta_{\tau,i-1}(x)}\Big|^{-\lambda_{\tau,i}/2},
  \]
  so that $|P|(x)=|P(x)|\,\beta_\lambda(x)$ for every form $P$ of weight $\lambda$ on $Z$.
\end{enumerate}
\end{definition}
\leanref{ReciprocityTheory/ReductiveGroups/ArithmeticGeometry/FundamentalDomains/Signature.lean}{ReciprocityTheory.ChartRestriction.G.FD.Data.maj}
\leanref{ReciprocityTheory/ReductiveGroups/ArithmeticGeometry/FundamentalDomains/Signature.lean}{ReciprocityTheory.ChartRestriction.G.FD.Data.flowMat}
\leanref{ReciprocityTheory/ReductiveGroups/ArithmeticGeometry/FundamentalDomains/ModerateCoefficientPositivity.lean}{ReciprocityTheory.ChartRestriction.G.FD.Data.spread}
\leanref{ReciprocityTheory/ReductiveGroups/ArithmeticGeometry/FundamentalDomains/ModerateCoefficientPositivity.lean}{ReciprocityTheory.ChartRestriction.G.FD.Data.nfac}

\begin{lemma}\label{cv-decay-flow}
In Situation~\ref{cv-cusp-setting}, let $\sigma$ be a good embedding, $Z$ a component at $\sigma$ and $x$ a
frame of $Z$, with the notation of Definition~\ref{cv-decay-majorant}. Let $v\in E^r$ be nonzero with
$G_Z(v,v)=0$, and put $m_\tau=h_{x,\tau}(\tau(v))$ for $\tau\in\Phi$. Then $m_\tau>0$ for every $\tau$, and:
\begin{enumerate}
\item $n_v(0)_\tau=1$ and $n_v(s)_\tau n_v(s')_\tau=n_v(s+s')_\tau$ for all $s,s'\in\mathbf C^\Phi$ and
  $\tau\in\Phi$;
\item for $s\in\mathbf C^\Phi$, the family $n_v(s)x$ is a frame of $Z$ if and only if
  $q_\tau=1+2\xi m_\tau\operatorname{Im}s_\tau$ is positive for every $\tau$;
\item if $n_v(s)x$ is a frame of $Z$, then for every weight $\lambda$ on $Z$
  \[
   \prod_\tau\min(q_\tau,q_\tau^{-1})^{\operatorname{sp}_\tau(\lambda)/2}
   \le\frac{\beta_\lambda(x)}{\beta_\lambda(n_v(s)x)}
   \le\prod_\tau\max(q_\tau,q_\tau^{-1})^{\operatorname{sp}_\tau(\lambda)/2}.
  \]
\end{enumerate}
\end{lemma}
\leanref{ReciprocityTheory/ReductiveGroups/ArithmeticGeometry/FundamentalDomains/EmbeddingHermitianPairs.lean}{ReciprocityTheory.ChartRestriction.G.FD.EmbeddingHermitianPairs.flowMat_mul}
\leanref{ReciprocityTheory/ReductiveGroups/ArithmeticGeometry/FundamentalDomains/FlowDomainCriterionHolds.lean}{ReciprocityTheory.ChartRestriction.G.FD.flowDomainCriterion_holds}
\leanref{ReciprocityTheory/ReductiveGroups/ArithmeticGeometry/FundamentalDomains/M3bNormHolds.lean}{ReciprocityTheory.ChartRestriction.G.FD.m3bNorm_holds}
\begin{proof}
Since $h_{x,\tau}$ is positive definite and $\tau(v)\ne0$, $m_\tau>0$. Part (1) uses only
$\tau(G_Z)(\tau(v),\tau(v))=\tau(G_Z(v,v))=0$: the product $n_v(s)_\tau n_v(s')_\tau$ sends $u$ to
$u-2i\xi(s_\tau+s'_\tau)\,\tau(G_Z)(u,\tau(v))\,\tau(v)$.

Fix $\tau$ and $s$. For $0\le i\le r$ let $e_i$ be the orthogonal projection of $\tau(v)$, for $\tau(G_Z)$,
onto the span of the first $i$ columns of $x_\tau$, and put $m_{\tau,i}=\tau(G_Z)(e_i,e_i)$. A computation
with the leading minors gives
\[
 \delta_{\tau,i}(n_v(s)x)=\delta_{\tau,i}(x)\,a_i,\qquad a_i=1+4\xi m_{\tau,i}\operatorname{Im}s_\tau .
\]
The numbers $m_{\tau,i}$ increase for $i\le q$ and decrease for $i\ge q$. Moreover $m_{\tau,0}=0$,
$m_{\tau,r}=\tau(G_Z)(\tau(v),\tau(v))=0$, and $m_{\tau,q}=\tau(G_Z)(\tau(v)^+,\tau(v)^+)=m_\tau/2$, because
$\tau(v)$ is isotropic. So $a_0=a_r=1$, $a_q=q_\tau$, and every $a_i$ lies between $1$ and $q_\tau$. Hence the
minors of $n_v(s)x$ have the signs required of a frame if and only if every $q_\tau$ is positive; as
$n_v(s)_\tau$ is invertible by (1), this is (2).

For (3), the factor at $\tau$ of $\beta_\lambda(x)/\beta_\lambda(n_v(s)x)$ is
$\prod_i(a_i/a_{i-1})^{\lambda_{\tau,i}/2}$, with logarithm $\frac12\sum_i\lambda_{\tau,i}(\log a_i-\log a_{i-1})$.
The increments $\log a_i-\log a_{i-1}$ have one sign for $i\le q$ and the opposite sign for $i>q$, and they
add up to $\log q_\tau$ over $i\le q$ and to $-\log q_\tau$ over $i>q$. So the logarithm lies between
$-\frac12\operatorname{sp}_\tau(\lambda)|\log q_\tau|$ and $\frac12\operatorname{sp}_\tau(\lambda)|\log q_\tau|$, and
taking the product over $\tau$ gives (3).
\end{proof}
\leanref{ReciprocityTheory/ReductiveGroups/ArithmeticGeometry/FundamentalDomains/MajorantFamilyNorms.lean}{ReciprocityTheory.ChartRestriction.G.FD.MajorantFamilyNorms.dmin_flow}

\begin{lemma}\label{cv-decay-fourier}
In Situation~\ref{cv-cusp-setting}, let $\sigma$ be a good embedding, $Z$ a component at $\sigma$ and $x$ a
frame of $Z$, with the notation of Definition~\ref{cv-decay-majorant}. Let $\lambda$ be a weight on $Z$ and
$P$ a form of weight $\lambda$ on $Z$ with $|P|$ bounded. Let $v\in E^r$ be nonzero with $G_Z(v,v)=0$, and
let $\Pi\subset\mathbf R^\Phi$ be a lattice of rank $f$ such that for every $p\in\Pi$ there is
$\gamma\in\Gamma_Z$ with $\chi_\lambda(\gamma)=1$ and $\tau(\gamma)=n_v(p)_\tau$ for every $\tau\in\Phi$. Write
$\langle\nu,s\rangle=\sum_\tau\nu_\tau s_\tau$, let $\Pi^*$ be the set of $\nu\in\mathbf R^\Phi$ with
$2\langle\nu,p\rangle\in\mathbf Z$ for every $p\in\Pi$, let $D$ be a fundamental domain for $\Pi$ in
$\mathbf R^\Phi$, and for $\nu\in\Pi^*$ put
\[
 a_\nu=\frac1{\operatorname{vol}(D)}\int_DP(n_v(t)x)\,e^{-4\pi i\langle\nu,t\rangle}\,dt .
\]
\begin{enumerate}
\item If $s\in\mathbf C^\Phi$ and $n_v(s)x$ is a frame of $Z$, then
  $P(n_v(s)x)=\sum_{\nu\in\Pi^*}a_\nu e^{4\pi i\langle\nu,s\rangle}$, and the series converges absolutely.
\item If $a_\nu\ne0$, then $\nu_\tau\ge0$ for every $\tau$.
\item If $a_0=0$, then $P(n_v(\zeta)x)\to0$ as $\operatorname{Im}\zeta\to\infty$.
\end{enumerate}
\end{lemma}
\leanref{ReciprocityTheory/ReductiveGroups/ArithmeticGeometry/FundamentalDomains/ModerateCoefficientPositivity.lean}{ReciprocityTheory.ChartRestriction.G.FD.Data.fj}
\leanref{ReciprocityTheory/ReductiveGroups/ArithmeticGeometry/FundamentalDomains/M3bExpHolds.lean}{ReciprocityTheory.ChartRestriction.G.FD.m3bExp_holds}
\leanref{ReciprocityTheory/ReductiveGroups/ArithmeticGeometry/FundamentalDomains/ModerateCoefficientPositivityHolds.lean}{ReciprocityTheory.ChartRestriction.G.FD.moderateCoefficientPositivity_holds}
\leanref{ReciprocityTheory/ReductiveGroups/ArithmeticGeometry/FundamentalDomains/MaximumBoundField.lean}{ReciprocityTheory.ChartRestriction.G.FD.flowTermNormDecay.tendsto_zero}
\begin{proof}
Let $m_\tau=h_{x,\tau}(\tau(v))$, and let $\Omega$ be the set of $y\in\mathbf R^\Phi$ with
$1+2\xi m_\tau y_\tau>0$ for every $\tau$, an open convex set containing $[0,\infty)^\Phi$. By
Lemma~\ref{cv-decay-flow}(2), $n_v(s)x$ is a frame of $Z$ if and only if $\operatorname{Im}s\in\Omega$, so
$F(s)=P(n_v(s)x)$ is holomorphic on the tube $\mathbf R^\Phi+i\Omega$. It is periodic under $\Pi$: for
$p\in\Pi$ and $\gamma$ as in the statement, $n_v(s+p)x=\gamma\,n_v(s)x$ by Lemma~\ref{cv-decay-flow}(1), so
$F(s+p)=\chi_\lambda(\gamma)F(s)=F(s)$.

(1) For $y\in\Omega$ let
$a_\nu(y)=\operatorname{vol}(D)^{-1}\int_DF(t+iy)e^{-4\pi i\langle\nu,t\rangle}dt$ be the Fourier coefficients of
the $\Pi$-periodic smooth function $t\mapsto F(t+iy)$, so that $a_\nu(0)=a_\nu$. Since $F$ is holomorphic,
$\partial F/\partial y_\tau=i\,\partial F/\partial t_\tau$, and integration by parts on $\mathbf R^\Phi/\Pi$ gives
$\partial a_\nu(y)/\partial y_\tau=-4\pi\nu_\tau a_\nu(y)$; as $\Omega$ is convex and contains $0$,
$a_\nu(y)=a_\nu e^{-4\pi\langle\nu,y\rangle}$. Let $y\in\Omega$, and choose $\delta>0$ with
$y+\delta\epsilon\in\Omega$ for every $\epsilon\in\{-1,0,1\}^\Phi$. Let $C$ be the largest of the suprema of $|F|$
on the sets $\mathbf R^\Phi+i(y+\delta\epsilon)$, which are finite by periodicity. For $\nu\in\Pi^*$ let
$\epsilon_\tau=-1$ if $\nu_\tau>0$ and $\epsilon_\tau=1$ otherwise; then
\[
 \big|a_\nu e^{-4\pi\langle\nu,y\rangle}\big|=|a_\nu(y+\delta\epsilon)|\,e^{4\pi\delta\langle\nu,\epsilon\rangle}
 \le Ce^{-4\pi\delta\sum_\tau|\nu_\tau|}.
\]
The right side is summable over the lattice $\Pi^*$. So the Fourier series of $t\mapsto F(t+iy)$ converges
absolutely, and its sum is $F(t+iy)$; at $t+iy=s$ this is (1).

(2) Let $S=\sup|P|$. For $y\in[0,\infty)^\Phi$ the numbers $q_\tau=1+2\xi m_\tau y_\tau$ are at least $1$, and
Lemma~\ref{cv-decay-flow}(3) gives
$\beta_\lambda(n_v(t+iy)x)^{-1}\le\beta_\lambda(x)^{-1}\prod_\tau q_\tau^{\operatorname{sp}_\tau(\lambda)/2}$. Since
$|F(t+iy)|=|P|(n_v(t+iy)x)/\beta_\lambda(n_v(t+iy)x)$, this gives
\[
 |a_\nu|=|a_\nu(y)|\,e^{4\pi\langle\nu,y\rangle}
 \le S\beta_\lambda(x)^{-1}\prod_\tau(1+2\xi m_\tau y_\tau)^{\operatorname{sp}_\tau(\lambda)/2}\,e^{4\pi\langle\nu,y\rangle}.
\]
If $\nu_{\tau_0}<0$ for some $\tau_0$, take $y_{\tau_0}=T$ and $y_\tau=0$ for $\tau\ne\tau_0$: the right side is
a power of $1+2\xi m_{\tau_0}T$ times $e^{4\pi\nu_{\tau_0}T}$, which tends to $0$ as $T\to\infty$. So $a_\nu=0$.

(3) For $\operatorname{Im}\zeta\ge0$ the family $n_v(\zeta)x$ is a frame, and by (1)
$P(n_v(\zeta)x)=\sum_\nu a_\nu e^{4\pi i\zeta\sum_\tau\nu_\tau}$. The term of index $\nu$ has absolute value
$|a_\nu|e^{-4\pi(\sum_\tau\nu_\tau)\operatorname{Im}\zeta}$. By (2) and $a_0=0$, every $\nu$ with $a_\nu\ne0$ is
nonzero with $\nu_\tau\ge0$ for all $\tau$, so $\sum_\tau\nu_\tau>0$ and the term tends to $0$ as
$\operatorname{Im}\zeta\to\infty$. The terms are bounded by $|a_\nu|$, which is summable by (1) at $s=0$, so the
sum tends to $0$ by dominated convergence.
\end{proof}

\begin{lemma}\label{cv-decay-schwarz}
Let $M'>0$, $0<t<T$ and $B\ge0$, and let $g:\mathbf C\to\mathbf C$ be holomorphic on
$\{z:\operatorname{Im}z>-T\}$, with $g(z+M')=g(z)$ for every $z$ and $g(z)\to0$ as $\operatorname{Im}z\to\infty$.
If $|g(z)|\le B$ whenever $\operatorname{Im}z=-t$, then $|g(0)|\le e^{-2\pi t/M'}B$.
\end{lemma}
\leanref{ReciprocityTheory/ReductiveGroups/ArithmeticGeometry/FundamentalDomains/PeriodicSchwarzLemma.lean}{ReciprocityTheory.ChartRestriction.G.FD.PeriodicSchwarzLemma.periodic_schwarz}
\begin{proof}
Put $R=e^{2\pi t/M'}>1$. By periodicity there is a holomorphic function $h$ on $0<|q|<e^{2\pi T/M'}$ with
$g(z)=h(e^{2\pi iz/M'})$, and $h(q)\to0$ as $q\to0$, so $h$ extends holomorphically to $q=0$ with $h(0)=0$.
The circle $|q|=R$ is the image of the line $\operatorname{Im}z=-t$, so $|h|\le B$ on it. The maximum
principle for $h(q)/q$ on $|q|\le R$ gives $|h(q)|\le B|q|/R$ there. At $q=1$, the image of $z=0$, this is
$|g(0)|\le B/R$.
\end{proof}

\begin{lemma}\label{cv-decay-maximum}
In Situation~\ref{cv-cusp-setting}, let $\sigma$ be a good embedding and $Z$ a component at $\sigma$, with the
notation of Definition~\ref{cv-decay-majorant}. Let $w>0$, let $\lambda$ be a weight on $Z$ of spread at most
$w$, and let $P$ be a cusp form of weight $\lambda$ on $Z$. Let $L'\subset E^r$ be a finitely generated
$\mathcal O_E$-submodule with $G_Z(u,u')\in\mathcal O_E$ for all $u,u'\in L'$, and let $M'>0$ be such that for
every nonzero $v\in L'$ with $G_Z(v,v)=0$ and every $b\in\mathcal O_E$ with $c(b)=b$ there is
$\gamma\in\Gamma_Z$ with $\chi_\lambda(\gamma)=1$ and $\tau(\gamma)=n_v(M'\tau(b))_\tau$ for every $\tau\in\Phi$.
Let $z_0$ be a frame of $Z$ with $|P|(z_0)>0$ and $|P|(y)\le2|P|(z_0)$ for every frame $y$ of $Z$, and put
\[
 \theta=\min\Big(1,\frac{\pi}{2\xi M'(fw+2)\log2}\Big).
\]
Then every nonzero $v\in L'$ has $h_{z_0,\tau}(\tau(v))\ge\theta$ for some $\tau\in\Phi$.
\end{lemma}
\leanref{ReciprocityTheory/ReductiveGroups/ArithmeticGeometry/FundamentalDomains/CuspidalExponentialDecay.lean}{ReciprocityTheory.ChartRestriction.G.FD.MaxBoundStatement}
\leanref{ReciprocityTheory/ReductiveGroups/ArithmeticGeometry/FundamentalDomains/MaxBoundPf.lean}{ReciprocityTheory.ChartRestriction.G.FD.maxBound_holds}
\begin{proof}
Suppose that $v\in L'$ is nonzero and that $m_\tau=h_{z_0,\tau}(\tau(v))<\theta$ for every $\tau$.

Step 1: $v$ is isotropic. The number $a=G_Z(v,v)$ lies in $\mathcal O_E$ and satisfies $c(a)=a$, so its
conjugates under $\tau$ and $\bar\tau$ have the same absolute value. For $\tau\in\Phi$,
$|\tau(a)|=|\tau(G_Z)(\tau(v),\tau(v))|\le h_{z_0,\tau}(\tau(v))<\theta\le1$, since
$\tau(G_Z)(u,u)=\tau(G_Z)(u^+,u^+)+\tau(G_Z)(u^-,u^-)$ with the two terms of opposite signs. So every conjugate
of the algebraic integer $a$ has absolute value less than $1$, its norm is an integer of absolute value less
than $1$, and $a=0$.

Step 2: a periodic function. By Step 1 and the hypothesis with $b=1$, there is $\gamma\in\Gamma_Z$ with
$\chi_\lambda(\gamma)=1$ and $\tau(\gamma)=n_v(M')_\tau$ for every $\tau$. Put $T=1/(2\xi\theta)$, $t=T/2$ and
$y(\zeta)=n_v(\zeta)z_0$. If $\operatorname{Im}\zeta>-T$, then $1+2\xi m_\tau\operatorname{Im}\zeta>1-2\xi\theta T=0$
for every $\tau$, so $y(\zeta)$ is a frame by Lemma~\ref{cv-decay-flow}(2). Let $g(\zeta)=P(y(\zeta))$ for
$\operatorname{Im}\zeta>-T$ and $g(\zeta)=0$ otherwise. Then $g$ is holomorphic on $\operatorname{Im}\zeta>-T$, and
$y(\zeta+M')=\gamma\,y(\zeta)$ by Lemma~\ref{cv-decay-flow}(1), so $g(\zeta+M')=\chi_\lambda(\gamma)g(\zeta)=g(\zeta)$.

Step 3: decay. Let $\Pi$ be the set of the vectors $(M'\tau(b))_{\tau\in\Phi}$ with $b\in\mathcal O_E$ and
$c(b)=b$. The fixed field of $c$ has degree $f$, and its real embeddings are the restrictions of the
$\tau\in\Phi$, so $\Pi$ is a lattice of rank $f$ in $\mathbf R^\Phi$, and it satisfies the hypothesis of
Lemma~\ref{cv-decay-fourier} for $x=z_0$. The constant coefficient $a_0$ of $P$ along $v$ vanishes, since $P$ is
a cusp form. By Lemma~\ref{cv-decay-fourier}(3), $g(\zeta)\to0$ as $\operatorname{Im}\zeta\to\infty$.

Step 4: a bound on the line $\operatorname{Im}\zeta=-t$. There
$q_\tau=1-2\xi m_\tau t$ lies in $(1/2,1]$, since $2\xi m_\tau t<2\xi\theta t=1/2$. By
Lemma~\ref{cv-decay-flow}(3), $\beta_\lambda(z_0)/\beta_\lambda(y(\zeta))\le\prod_\tau q_\tau^{-\operatorname{sp}_\tau(\lambda)/2}
\le2^{w/2}$. Hence
\[
 |g(\zeta)|=\frac{|P|(y(\zeta))}{\beta_\lambda(y(\zeta))}\le\frac{2|P|(z_0)\,2^{w/2}}{\beta_\lambda(z_0)} .
\]

Step 5: the contradiction. By Lemma~\ref{cv-decay-schwarz} and $y(0)=z_0$,
$|P(z_0)|=|g(0)|\le e^{-2\pi t/M'}\,2^{1+w/2}|P|(z_0)/\beta_\lambda(z_0)$. Multiplying by $\beta_\lambda(z_0)$ and
dividing by $|P|(z_0)>0$ gives $1\le e^{-2\pi t/M'}2^{1+w/2}$. But
$2\pi t/M'=\pi/(2\xi M'\theta)\ge(fw+2)\log2$ by the choice of $\theta$, so
$e^{-2\pi t/M'}2^{1+w/2}\le2^{1+w/2-fw-2}<1$, as $f\ge1$ and $w>0$.
\end{proof}
\leanref{ReciprocityTheory/ReductiveGroups/ArithmeticGeometry/FundamentalDomains/MaximumBoundField.lean}{ReciprocityTheory.ChartRestriction.G.FD.flowTermNormDecay.herm_eq_zero}

\begin{lemma}\label{cv-submean-chart}
In Situation~\ref{cv-cusp-setting}, let $\sigma$ be a good embedding, $Z$ a component at $\sigma$ and $x$ an
orthonormal frame of $Z$, with the notation of Definition~\ref{cv-decay-majorant}. Let $w>0$ and let $\lambda$
be a weight on $Z$ of spread at most $w$. Let $z=(z_{\tau,ab})$, $\tau\in\Phi$, $1\le b<a\le r$, be a family of
complex numbers with $|z|^2=\sum|z_{\tau,ab}|^2\le\rho_w^2$. Then $xn(z)$ is a frame of $Z$,
$\beta_\lambda(xn(z))^2\ge e^{-1/20}$, and the density of $\mu_x$ at $z$ is at least $e^{-1/10}$.
\end{lemma}
\leanref{Development/AbelianVarietyTheory/Proofs/DivisorProofs/NormFactorAndDensityNearOneInFlagChartStatement.lean}{GlobalTheory.T.normFactor_and_density_near_one_in_flagChart}

This is a self-contained computation. Since $x$ is orthonormal, $|\delta_{\tau,i}(x)|=1$ for all $\tau$ and $i$.
Since $n(z)_\tau$ is unipotent, $|\delta_{\tau,r}(xn(z))|=1$ as well, so $\beta_\lambda(xn(z))$ depends on
$\lambda$ only through the differences of its entries at each $\tau$, and the bound needs only the spread. The
minors $\delta_{\tau,i}(xn(z))$ are polynomials in $z$ with value $\delta_{\tau,i}(x)$ at $z=0$, and the radius
$\rho_w=1/(200r^3(w+f))$ keeps them close enough to $\pm1$ for both bounds.

\begin{lemma}\label{cv-returns-majorant}
In Situation~\ref{cv-cusp-setting}, let $\sigma$ be a good embedding, $Z$ a component at $\sigma$, $x$ a frame
of $Z$, $\tau\in\Phi$ and $u\in\mathbf C^r$, with the notation of Definition~\ref{cv-decay-majorant}.
\begin{enumerate}
\item $h_{xb,\tau}(u)=h_{x,\tau}(u)$ for every $b\in B$.
\item $h_{\gamma x,\tau}(\tau(\gamma)u)=h_{x,\tau}(u)$ for every $\gamma\in\Gamma_Z$.
\item If $x$ is orthonormal and $z=(z_{\tau,ab})$ is a family of complex numbers indexed as in
  Lemma~\ref{cv-submean-chart} with $r^2|z|\le1/100$, then $h_{xn(z),\tau}(u)\le2h_{x,\tau}(u)$ and $h_{x,\tau}(u)\le2h_{xn(z),\tau}(u)$.
\end{enumerate}
\end{lemma}
\leanref{Development/AbelianVarietyTheory/Proofs/DivisorProofs/CardGroupElementsChartReturnLe.lean}{GlobalTheory.ChartMult.maj_rmul_any}
\leanref{Development/AbelianVarietyTheory/Proofs/TorusHaarMeasure/MajLactRmul.lean}{GlobalTheory.MajInv.maj_lact}
\leanref{Development/AbelianVarietyTheory/Proofs/DivisorProofs/MajorantChartPointCentreFactorTwoStatement.lean}{GlobalTheory.T.majorant_chart_point_centre_factor_two}
\begin{proof}
(1) Since $b_\tau$ is upper triangular, the first $q$ columns of $x_\tau b_\tau$ span $U^+_{x,\tau}$, so the
decomposition $u=u^++u^-$ is the same for $xb$ and for $x$. (2) The matrix $\tau(\gamma)$ preserves
$\tau(G_Z)$ and maps $U^+_{x,\tau}$ onto $U^+_{\gamma x,\tau}$, so it maps the decomposition of $u$ for $x$ to
that of $\tau(\gamma)u$ for $\gamma x$. (3) is a self-contained computation comparing the projections onto
$U^+_{xn(z),\tau}$ and onto $U^+_{x,\tau}$; for an orthonormal $x$, $h_{x,\tau}$ is the standard Hermitian norm in
the coordinates given by the columns of $x_\tau$.
\end{proof}

\begin{lemma}\label{cv-returns-short-basis}
In Situation~\ref{cv-cusp-setting}, let $\sigma$ be a good embedding, $Z$ a component at $\sigma$ and $x$ a
frame of $Z$, with the notation of Definition~\ref{cv-decay-majorant}. Let $L'\subset E^r$ be a finitely
generated $\mathcal O_E$-submodule that is self-dual for $G_Z$: a vector $u\in E^r$ lies in $L'$ if and only
if $G_Z(u,u')\in\mathcal O_E$ for every $u'\in L'$. Let $\theta>0$, and suppose that
$\sum_{\tau\in\Phi}h_{x,\tau}(\tau(v))\ge\theta$ for every nonzero $v\in L'$. Then $L'$ contains a basis
$w_1,\dots,w_r$ of $E^r$ with
\[
 \sum_{\tau\in\Phi}h_{x,\tau}(\tau(w_i))\le\frac{(2fr)^4|d_E|^2}{4\theta}\qquad(1\le i\le r).
\]
\end{lemma}
\leanref{Development/AbelianVarietyTheory/Proofs/CuspFormEvaluations/ExistsThickPointHyperbolic.lean}{GlobalTheory.exists_short_basis_at}
\begin{proof}
Let $V=E^r$, a $\mathbf Q$-vector space of dimension $2fr$, and $T(u,u')=\operatorname{Tr}_{E/\mathbf Q}G_Z(u,u')$,
a nondegenerate symmetric $\mathbf Q$-bilinear form on $V$. Since $\Phi$ contains one embedding of each
conjugate pair, $u\mapsto(\tau(u))_{\tau\in\Phi}$ identifies $V\otimes\mathbf R$ with
$\bigoplus_{\tau\in\Phi}\mathbf C^r$. On this space let $\langle\,,\rangle_x$ be the real part of the sum over $\tau$
of the Hermitian forms of the majorants $h_{x,\tau}$, so that $|u|_x^2=\sum_\tau h_{x,\tau}(\tau(u))$ for
$u\in V$, and let $J$ act on the $\tau$-component by $1$ on $U^+_{x,\tau}$ and by $-1$ on its orthogonal
complement for $\tau(G_Z)$. Then $J$ is an isometry for $\langle\,,\rangle_x$, and
$\tau(G_Z)(u,u')$ is the value at $(Ju,u')$ of the Hermitian form of $h_{x,\tau}$; summing $\tau$ and $\bar\tau$,
\[
 \langle 2Ju,u'\rangle_x=\sum_{\tau\in\Phi}2\operatorname{Re}\tau(G_Z(u,u'))=T(u,u')\qquad(u,u'\in V).
\]

Let $L'^T=\{u\in V:T(u,L')\subset\mathbf Z\}$; its dual for $T$ is $L'$. Let $u\in L'^T$ be nonzero. For
$u'\in L'$ and $a\in\mathcal O_E$, $\operatorname{Tr}_{E/\mathbf Q}(a\,G_Z(u,u'))=T(u,c(a)u')\in\mathbf Z$, so
$G_Z(u,u')$ lies in the inverse different $\mathfrak d^{-1}$ of $E$, and $|d_E|G_Z(u,u')\in\mathcal O_E$
because $|d_E|=N(\mathfrak d)$ lies in $\mathfrak d$. By self-duality $|d_E|u\in L'$, so
$\theta\le|d_E|^2|u|_x^2$, and $|2Ju|_x^2=4|u|_x^2\ge4\theta/|d_E|^2$. So every nonzero vector of the lattice
$2J(L'^T)$ of the Euclidean space $(V\otimes\mathbf R,\langle\,,\rangle_x)$ has length at least
$2\sqrt\theta/|d_E|$, and its dual lattice is $L'$, by the displayed identity. By
Lemma~\ref{cv-transference} with $N=2fr$, $L'$ contains $2fr$ vectors, linearly independent over
$\mathbf Q$, with $|w|_x\le(2fr)^2|d_E|/(2\sqrt\theta)$, that is,
$\sum_\tau h_{x,\tau}(\tau(w))\le(2fr)^4|d_E|^2/(4\theta)$. They span $V$ over $\mathbf Q$, so they span $E^r$
over $E$, and $r$ of them form a basis of $E^r$.
\end{proof}
\leanref{Development/AbelianVarietyTheory/Proofs/CuspFormEvaluations/ExistsThickPointHyperbolic.lean}{GlobalTheory.ThickPoint.transference_pair}

\begin{lemma}\label{cv-returns-count}
In Situation~\ref{cv-cusp-setting}, let $\sigma$ be a good embedding, $Z$ a component at $\sigma$ and $x$ a
frame of $Z$, with the notation of Definition~\ref{cv-decay-majorant}. Let $L'\subset E^r$ be a finitely
generated $\mathcal O_E$-submodule, and let $\theta>0$ be such that $\sum_\tau h_{x,\tau}(\tau(v))\ge\theta$ for
every nonzero $v\in L'$. Then for every $R\ge0$ the number of $v\in L'$ with
$\sum_\tau h_{x,\tau}(\tau(v))\le R$ is at most $(1+2\sqrt{R/\theta})^{2fr}$.
\end{lemma}
\leanref{Development/AbelianVarietyTheory/Proofs/DivisorProofs/EncardSelfDualLatticeBoundedMajorantLeStatement.lean}{GlobalTheory.T.encard_selfDualLattice_bounded_majorant_le}
\begin{proof}
Consider the Euclidean norm $|u|_x^2=\sum_\tau h_{x,\tau}(u_\tau)$ on $\bigoplus_{\tau\in\Phi}\mathbf C^r$, a real
space of dimension $2fr$, in which $v\in L'$ is represented by $(\tau(v))_\tau$. Two distinct vectors $v,v'$ in
the set to be counted have $|v-v'|_x\ge\sqrt\theta$, as $v-v'\in L'$ is nonzero, and each has
$|v|_x\le\sqrt R$. So the balls of radius $\sqrt\theta/2$ around them are disjoint and lie in the ball of
radius $\sqrt R+\sqrt\theta/2$ around $0$. Comparing volumes, their number is at most
$\big((\sqrt R+\sqrt\theta/2)/(\sqrt\theta/2)\big)^{2fr}=(1+2\sqrt{R/\theta})^{2fr}$.
\end{proof}

\begin{lemma}\label{cv-returns-unfold}
In Situation~\ref{cv-cusp-setting}, let $\sigma$ be a good embedding and $Z$ a component at $\sigma$. Let $x_0$
and $x$ be orthonormal frames of $Z$, let $\mathcal F$ be a fundamental set for $\Gamma_Z$ at $x_0$, and let
$A\subset\Omega_x$ be measurable. Let $k\ge0$ be a measurable function on the frames of $Z$ with
$k(\gamma y)=k(y)=k(yb)$ for all frames $y$, $\gamma\in\Gamma_Z$ and $b\in B$, and let $m$ be at least the number
of $\gamma\in\Gamma_Z$ with $\gamma xn(z)\in xn(z')B$ for some $z,z'\in A$. Then
\[
 \int_Ak(xn(z))\,d\mu_x(z)\le m\int_{\mathcal F}k(x_0n(z))\,d\mu_{x_0}(z).
\]
\end{lemma}
\leanref{Development/AbelianVarietyTheory/Proofs/DivisorProofs/LintegralInvariantChartSetLeMulFundamentalSetStatement.lean}{GlobalTheory.T.lintegral_invariant_chartSet_le_mul_fundamentalSet}
\begin{proof}
Two facts about the charts are used. First, if $k'\ge0$ is measurable on the frames of $Z$ with
$k'(yb)=k'(y)$ for $b\in B$, and $x_1,x_2$ are orthonormal frames of $Z$, then
$\int_{\Omega_{x_1}}k'(x_1n(z))\,d\mu_{x_1}(z)=\int_{\Omega_{x_2}}k'(x_2n(z))\,d\mu_{x_2}(z)$. This is a
self-contained computation, which says that the measures $\mu_{x_1}$ are the readings in the charts of one
measure on the frames modulo $B$. Second, for such $k'$,
\[
 \int_{\Omega_{x_0}}k'(x_0n(z))\,d\mu_{x_0}(z)=\int_{\mathcal F}\sum_{\gamma\in\Gamma_Z}k'(\gamma x_0n(z))\,d\mu_{x_0}(z),
\]
by the definition of a fundamental set and the first fact, applied to the orthonormal frames $\gamma x_0$.

Let $S$ be the set of frames $xn(z)b$ with $z\in A$ and $b\in B$, and let $k'=\mathbf 1_Sk$, which is invariant
under $B$. As a lower unitriangular matrix in $B$ is $1$, $xn(z)\in S$ only for $z\in A$. So the two facts give
\[
 \int_Ak(xn(z))\,d\mu_x(z)=\int_{\Omega_x}k'(xn(z))\,d\mu_x(z)
 =\int_{\mathcal F}\sum_{\gamma\in\Gamma_Z}k'(\gamma x_0n(z))\,d\mu_{x_0}(z).
\]
Fix $z\in\mathcal F$, put $y=x_0n(z)$, and let $T_y$ be the set of $\gamma\in\Gamma_Z$ with $\gamma y\in S$. If
$\gamma_0\in T_y$, then $\gamma\mapsto\gamma\gamma_0^{-1}$ maps $T_y$ injectively into the set counted by $m$: from
$\gamma_0y=xn(z_0)b_0$ and $\gamma y=xn(z_1)b_1$ with $z_0,z_1\in A$ follows
$\gamma\gamma_0^{-1}xn(z_0)=xn(z_1)b_1b_0^{-1}$. So $\sum_\gamma k'(\gamma y)\le|T_y|\,k(y)\le m\,k(y)$, using
$k(\gamma y)=k(y)$, and integrating over $\mathcal F$ gives the claim.
\end{proof}
\leanref{ReciprocityTheory/ReductiveGroups/ArithmeticGeometry/ArakelovVolumes/UpperBounds/MeasInvStatementPf.lean}{ReciprocityTheory.ChartRestriction.G.AR.Up.measInv_holds}
\leanref{Development/AbelianVarietyTheory/Proofs/TorusHaarMeasure/MeasurableSetSatFrames.lean}{GlobalTheory.Unfold.lintegral_unfold}

\begin{lemma}\label{cv-tube-chart}
In Situation~\ref{cv-cusp-setting}, let $G\in M_r(E)$ be a Hermitian matrix and $\Gamma$ a subgroup of its
unitary group over $E$ with the properties of the pairs $(G_Z,\Gamma_Z)$ in the compatibility conditions, and
suppose that $\det\gamma=1$ for every $\gamma\in\Gamma$. Define frames for $G$, the sets $\Omega_x$, the measures
$\mu_x$ and fundamental sets for $\Gamma$ as in Situation~\ref{cv-cusp-setting}, with $G$ and $\Gamma$ in place of
$G_Z$ and $\Gamma_Z$. Let $SU_\Phi(G)=\prod_{\tau\in\Phi}SU(\tau(G))$, into which $\Gamma$ embeds by
$\gamma\mapsto(\tau(\gamma))_\tau$. For a family $x=(x_\tau)_{\tau\in\Phi}$ in $GL_r(\mathbf C)$ let $\nu_x$ be the
Haar measure on $SU_\Phi(G)$ normalized at the identity by the Lebesgue measure on its Lie algebra for the
Euclidean norm $|Y|^2=\sum_\tau\|x_\tau^{-1}Y_\tau x_\tau\|^2$, where $\|\cdot\|$ is the Frobenius norm. Let $x$
be an orthonormal frame for $G$ and $\mathcal F$ a fundamental set for $\Gamma$ at $x$. Then
\[
 \big(2^{r(r-1)/2}(2\pi)^{r-1}\sqrt r\big)^f\mu_x(\mathcal F)=\operatorname{vol}(\Gamma\backslash SU_\Phi(G),\nu_x),
\]
where the right side is the $\nu_x$-measure of a measurable fundamental domain for $\Gamma$ in
$SU_\Phi(G)$, and $0$ if there is none.
\end{lemma}
\leanref{Development/AbelianVarietyTheory/PublishedTheorems/Divisors₂/EmbGLMemRealSU.lean}{GlobalTheory.SUVol.riemHaar}
\leanref{Development/AbelianVarietyTheory/Proofs/DivisorProofs/CovolumeRealSpecialUnitaryEqFlagDomainMassStatement.lean}{GlobalTheory.T.covolume_realSpecialUnitary_eq_flagDomain_mass}

This is a self-contained computation. For an orthonormal $x$, the stabilizer of $x_\tau$ modulo upper
triangular matrices in $SU(\tau(G))$ is $x_\tau$ times the group of diagonal unitary matrices of determinant
$1$ times $x_\tau^{-1}$, whose volume for the norm $|Y|$ is $(2\pi)^{r-1}\sqrt r$. The factor $2^{r(r-1)/2}$ at
each $\tau$ compares the Lebesgue measure in the chart coordinates $z_{\tau,ab}$ with the norm $|Y|$ on the
remaining part of the Lie algebra, and the density of $\mu_x$ makes it the invariant measure on the frames
modulo $B$.

\begin{lemma}\label{cv-tube-det}
Let $E$ be a CM field with complex conjugation $c$, let $G\in M_r(E)$ be a Hermitian matrix with
$\det G\ne0$, let $\Lambda\subset E^r$ be a finitely generated $\mathcal O_E$-submodule spanning $E^r$, and let
$M\ge3$. Let $\gamma\in GL_r(E)$ preserve the Hermitian form of $G$, with $\gamma\Lambda=\Lambda$ and
$(\gamma-1)\Lambda\subset M\Lambda$. Then $\det\gamma=1$.
\end{lemma}
\leanref{Development/AbelianVarietyTheory/Proofs/DivisorProofs/DetEqOneUnitaryPrincipalCongruenceLevelThreeStatement.lean}{GlobalTheory.T.det_eq_one_unitary_principal_congruence_level_three}
\begin{proof}
Let $u=\det\gamma$. Since $\gamma$ preserves the form of $G$, $c(u)u\det G=\det G$, so $c(u)u=1$ and every
conjugate of $u$ has absolute value $1$. The module $\Lambda$ is projective of rank $r$, and $\gamma$ acts on the
invertible $\mathcal O_E$-module $I=\wedge^r\Lambda$ by $u$. From $\gamma\Lambda=\Lambda$ follows $uI=I$, so
$u\in\mathcal O_E$; from $(\gamma-1)\Lambda\subset M\Lambda$ follows that $\gamma$ acts trivially on
$\Lambda/M\Lambda$, hence on $I/MI$, so $u-1\in M\mathcal O_E$. Then $(u-1)/M$ lies in $\mathcal O_E$ and all its
conjugates have absolute value at most $2/M<1$, so its norm is an integer of absolute value less than $1$, and
$u=1$.
\end{proof}

\begin{lemma}\label{cv-tube-index}
Let $E$ be a CM field with $[E:\mathbf Q]=2f$, let $\Lambda\subset E^r$ be a finitely generated
$\mathcal O_E$-submodule spanning $E^r$, put $\hat\Lambda=\Lambda\otimes\hat{\mathbf Z}$, and let $M\ge1$. Let
$K$ be a group of $\hat{\mathcal O}_E$-linear automorphisms of $\hat\Lambda$, and let $K(M)$ be the subgroup of
the $k\in K$ with $(k-1)\hat\Lambda\subset M\hat\Lambda$. Then $K$ is the union of at most $M^{2fr^2}$ cosets
$sK(M)$ with $s\in K$.
\end{lemma}
\leanref{Development/WeilBound/RepresentationTheoryGL/TubeMeasures/DisjointTubes/TubeSt.lean}{MainStatements.Tam.IdxSt}
\leanref{Development/WeilBound/RepresentationTheoryGL/TubeMeasures/DisjointTubes/InverseMultiplication/IsKMInvMul.lean}{MainStatements.Tam.idx}
\begin{proof}
The module $\Lambda$ is projective of rank $r$, so $\Lambda/M\Lambda\cong(\mathcal O_E/M)^r$, which has $M^{2fr}$
elements and is generated by $r$ elements $e_1,\dots,e_r$; and $\hat\Lambda/M\hat\Lambda=\Lambda/M\Lambda$. For
$k\in K$ let $t(k)$ be the family of the classes of $k(e_i)$ in $\Lambda/M\Lambda$; there are at most
$M^{2fr^2}$ such families. If $t(k)=t(k')$, then $k'^{-1}k$ fixes the generators $e_i$ of
$\hat\Lambda/M\hat\Lambda$, hence acts trivially on it, so $k\in k'K(M)$. Choosing one $s$ for each family that
occurs gives the cosets.
\end{proof}

\begin{lemma}\label{cv-tube-archimedean}
In Situation~\ref{cv-cusp-setting}, let $c$ be the complex conjugation of $E$, let $r\ge1$, and let
$H\in M_r(E)$ be an invertible matrix with $c(H_{ab})=H_{ba}$ for all $a,b$ such that $E^r$, with the Hermitian
form $h(v,v')=\sum_{a,b}c(v_a)H_{ab}v'_b$, has signature $(P,Q)$ at every real place of $F$. Put $G=-H^{\mathsf T}$.
Its form $G(v,v')=\sum_{a,b}v_aG_{ab}c(v'_b)$ is $-h(v',v)$, of signature $(Q,P)$, and an \emph{orthonormal frame}
for $G$ is a family $x=(x_\tau)_{\tau\in\Phi}$ in $GL_r(\mathbf C)$ such that each
$x_\tau^{\mathsf T}\tau(G)\bar x_\tau$ is diagonal with entries $1$, $Q$ times, followed by $-1$. Let $m$ be the Haar
measure on $U(P,Q)$ of Theorem~\ref{cv-volume-mass}, and $m^\Phi$ the product measure on $U(P,Q)^\Phi$. There are an
orthonormal frame $x_0$ for $G$, isomorphisms $U(\tau(G))\cong U(P,Q)$ for $\tau\in\Phi$, given by $x_0$, and a continuous map
$\pi:U(P,Q)^\Phi\to SU_\Phi(G)$ with $\pi(sb)=s\pi(b)$ for $s\in SU_\Phi(G)$, viewed in $U(P,Q)^\Phi$ through
these isomorphisms, and $b\in U(P,Q)^\Phi$, such that for every $\tau_0\in\Phi$ and every measurable
$S\subset SU_\Phi(G)$, with $\nu_{x_0}$ as in Lemma~\ref{cv-tube-chart},
\[
 16^{-fr^2}\nu_{x_0}(S)\le m^\Phi\big(\pi^{-1}(S)\cap\{b:\operatorname{Im}\det b_{\tau_0}>0\}\big).
\]
\end{lemma}
\leanref{Development/AbelianVarietyTheory/Proofs/DivisorProofs/ExistsFlagDomainOrthonormalFrameIdentificationStatement.lean}{GlobalTheory.T.exists_flagDomain_orthonormalFrame_identification}
\leanref{Development/AbelianVarietyTheory/Proofs/DivisorProofs/ExistsContinuousEquivariantDetOneFactorProjectioStatement.lean}{GlobalTheory.T.exists_continuous_equivariant_detOneFactorProjection}
\leanref{Development/AbelianVarietyTheory/Proofs/DivisorProofs/ArchimedeanMassProductLowerBoundUnitaryPqStatement.lean}{GlobalTheory.T.archimedean_mass_product_lower_bound_unitary_pq}

This is a self-contained computation, comparing the normalizations of $m$ and of $\nu_{x_0}$ at the
identity through $\pi$.

\begin{lemma}\label{cv-tube-bound}
In Situation~\ref{cv-cusp-setting}, let $c$ be the complex conjugation of $E$, and let $H\in M_r(E)$ be an
invertible matrix with $c(H_{ab})=H_{ba}$ for all $a,b$ such that $V=E^r$, with the Hermitian form
$h(v,v')=\sum_{a,b}c(v_a)H_{ab}v'_b$, has signature $(P,Q)$ at every real place of $F$. Put $G=-H^{\mathsf T}$, whose
form $G(v,v')=\sum_{a,b}v_aG_{ab}c(v'_b)=-h(v',v)$ has signature $(Q,P)$. Suppose that $V$ contains an
$\mathcal O_E$-lattice $L$ that is self-dual for $h$, equivalently for $G$. The unitary groups of $h$ and of $G$ are
the same group $U(V)$. Let $K\subset U(V)(\mathbf A_{F,f})$ be the stabilizer of $\hat L$, and $K(M)$ the
group of the $k\in K$ with $(k-1)\hat L\subset M\hat L$. For $g\in U(V)(\mathbf A_{F,f})$ let
$L_g=V\cap g\hat L$, and let $\Gamma_g$ be the group of the $\gamma\in U(V)(F)$ with $\gamma L_g=L_g$ and
$(\gamma-1)L_g\subset ML_g$; by Lemma~\ref{cv-tube-det} it embeds in $SU_\Phi(G)$. Let
$g_1,\dots,g_n\in U(V)(\mathbf A_{F,f})$ lie in distinct double cosets $U(V)(F)g_kK(M)$, and let
$x_1,\dots,x_n$ be orthonormal frames for $G$: families $x_k=(x_{k,\tau})_{\tau\in\Phi}$ in $GL_r(\mathbf C)$ with
each $x_{k,\tau}^{\mathsf T}\tau(G)\bar x_{k,\tau}$ diagonal with entries $1$, $Q$ times, followed by $-1$. Then, with
$\nu_x$ as in Lemma~\ref{cv-tube-chart},
\[
 \sum_{k=1}^n\operatorname{vol}(\Gamma_{g_k}\backslash SU_\Phi(G),\nu_{x_k})
 \le16^{fr^2}M^{2fr^2}\operatorname{vol}U(r)^f\operatorname{mass}(I_r),
\]
with $\operatorname{vol}U(r)$ as in Theorem~\ref{cv-volume-mass} and $I_r$ as in Definition~\ref{cv-mass}.
\end{lemma}
\leanref{Development/WeilBound/RepresentationTheoryGL/TubeMeasures/DisjointTubes/TubeSt.lean}{MainStatements.Tam.TubeSt}
\leanref{Development/WeilBound/RepresentationTheoryGL/TubeMeasures/DisjointTubes/FrameCovolume/CovolFrame.lean}{MainStatements.Tam.tube}
\begin{proof}
Let $x_0$, the isomorphisms and $\pi$ be as in Lemma~\ref{cv-tube-archimedean}, and fix $\tau_0\in\Phi$. Write
$G_\infty=U(P,Q)^\Phi$, identified with $\prod_\rho U(V)(F_\rho)$ through these isomorphisms, and
$h=(h_\infty,h_f)$ for the components of $h\in U(V)(\mathbf A_F)$. Let $\mu$ be a Haar measure on
$U(V)(\mathbf A_F)$, $\mathcal F_V$ a measurable fundamental domain for $U(V)(F)$, and
$X=\operatorname{vol}U(r)^f\operatorname{mass}(I_r)$.

Step 1: packing sets. The measure $\nu_x$ is the same for all orthonormal frames $x$: for two of them,
$x'^{-1}_\tau x_\tau$ preserves the form of the diagonal matrix $x_\tau^{\mathsf T}\tau(G)\bar x_\tau$, and conjugation by an element of the
unimodular group of that form preserves the Lebesgue measure on its Lie algebra. So $\nu_{x_k}=\nu_{x_0}$. For
each $k$ choose a measurable $F_k\subset SU_\Phi(G)$ with
$\nu_{x_0}(F_k)\ge\operatorname{vol}(\Gamma_{g_k}\backslash SU_\Phi(G),\nu_{x_0})$ such that $\gamma y\in F_k$ with
$y\in F_k$ and $\gamma\in\Gamma_{g_k}$ only if $\gamma=1$: remove from a measurable fundamental domain the null set
of its points that lie in another of its countably many translates, or take $F_k=\emptyset$ if there is no
fundamental domain.

Step 2: disjoint tubes. Put $B_k=\pi^{-1}(F_k)\cap\{b\in G_\infty:\operatorname{Im}\det b_{\tau_0}>0\}$ and
$T_k=\{h:g_k^{-1}h_f\in K(M),\ h_\infty\in B_k\}$. Let $\delta\in U(V)(F)$, $h\in T_k$ and $\delta h\in T_{k'}$.
Then $g_{k'}^{-1}\delta g_k=(g_{k'}^{-1}\delta h_f)(g_k^{-1}h_f)^{-1}$ lies in $K(M)$, so $g_k$ and $g_{k'}$ lie
in the same double coset and $k=k'$. Then $\delta\in g_kK(M)g_k^{-1}$, so $\delta$ preserves $g_k\hat L$ and acts
trivially on $g_k\hat L/Mg_k\hat L$, and $\delta\in\Gamma_{g_k}$. Its image $\delta_\infty$ lies in $SU_\Phi(G)$, so
$\pi((\delta h)_\infty)=\delta_\infty\pi(h_\infty)$, and both $\pi(h_\infty)$ and $\delta_\infty\pi(h_\infty)$ lie in
$F_k$; by Step 1, $\delta=1$. Hence the sets $T_k$ are disjoint, and the translates of their union by the
elements of $U(V)(F)$ are disjoint, so $\sum_k\mu(T_k)\le\mu(\mathcal F_V)$.

Step 3: the index. Left translation by $g_k$, whose archimedean component is $1$, gives
$\mu(T_k)=\mu\{h:h_f\in K(M),\ h_\infty\in B_k\}$. By Lemma~\ref{cv-tube-index}, $K$ is the union of at most
$M^{2fr^2}$ cosets $sK(M)$ with $s\in K$, and translating by these $s$ gives
$\mu\{h:h_f\in K,\ h_\infty\in B_k\}\le M^{2fr^2}\mu(T_k)$.

Step 4: the volume formula. By Theorem~\ref{cv-volume-mass},
$\mu(\mathcal F_V)\,m^\Phi(B_k)=\mu\{h:h_f\in K,\ h_\infty\in B_k\}\,X$ for every $k$. Summing over $k$ and using
Steps 2 and 3,
\[
 \mu(\mathcal F_V)\sum_km^\Phi(B_k)\le M^{2fr^2}X\sum_k\mu(T_k)\le M^{2fr^2}X\,\mu(\mathcal F_V).
\]
Here $\mu(\mathcal F_V)$ is finite, since $U(V)(F)\backslash U(V)(\mathbf A_F)$ has finite volume
\cite{BHC62}, and positive, since otherwise Theorem~\ref{cv-volume-mass} with $B=G_\infty$ would give
$\mu\{h:h_f\in K\}=0$ for this nonempty open set. So $\sum_km^\Phi(B_k)\le M^{2fr^2}X$.

Step 5: the comparison. By Lemma~\ref{cv-tube-archimedean}, $\nu_{x_0}(F_k)\le16^{fr^2}m^\Phi(B_k)$. With
Steps 1 and 4,
\[
 \sum_k\operatorname{vol}(\Gamma_{g_k}\backslash SU_\Phi(G),\nu_{x_k})\le\sum_k\nu_{x_0}(F_k)
 \le16^{fr^2}\sum_km^\Phi(B_k)\le16^{fr^2}M^{2fr^2}X.
\]
\end{proof}
\leanref{Development/AbelianVarietyTheory/Proofs/DivisorProofs/ExistsPackingSetMeasureGeCovolumeStatement.lean}{GlobalTheory.T.exists_packingSet_measure_ge_covolume}
\leanref{Development/AbelianVarietyTheory/Proofs/DivisorProofs/RiemannianHaarIndependentOfOrthonormalFrameStatement.lean}{GlobalTheory.T.riemannianHaar_independent_of_orthonormal_frame}
\leanref{Development/WeilBound/RepresentationTheoryGL/TubeMeasures/DisjointTubes/FrameCovolume/TubeDisj.lean}{MainStatements.Tam.T2.tube_disj}
\leanref{Development/WeilBound/RepresentationTheoryGL/TubeMeasures/DisjointTubes/FrameCovolume/MeasTransl.lean}{MainStatements.Tam.T2.meas_low}

\begin{lemma}\label{cv-cusp-frame-graph}
Let $0\le p\le r$ with $r\ge1$. For $u\in\mathbf C^r$ write $u^+=(u_1,\dots,u_p)$ and
$u^-=(u_{p+1},\dots,u_r)$, let $|u|$ be the standard Hermitian norm, and put
$j(u,u')=\sum_{a\le p}\bar u_au'_a-\sum_{a>p}\bar u_au'_a$. Let $\mu>0$ and $0<\rho\le1$, put
$R=4\sqrt r\,\mu/\rho$, and let $X_1,\dots,X_p\in\mathbf C^r$ satisfy $|X_a-Re_a|\le\mu$, where $e_a$ is the
$a$-th standard basis vector. Let $U$ be the span of $X_1,\dots,X_p$.
\begin{enumerate}
\item If $u=\sum_aw_aX_a$ with $w\in\mathbf C^p$, then $|u^-|\le\frac\rho3|u^+|$, and $u^+\ne0$ if $w\ne0$.
\item If $y\in\mathbf C^r$ is nonzero and $j(X_a,y)=0$ for every $a\le p$, then $j(y,y)<0$.
\item If $y,v\in\mathbf C^r$ satisfy $j(X_a,y)=0$ for every $a\le p$ and $v-y\in U$, then
  $|j(y,y)|\le2|v|^2$.
\end{enumerate}
\end{lemma}
\leanref{ReciprocityTheory/SemistableReduction/ApproxGraph.lean}{ReciprocityTheory.MultiplicativeHeights.approxGraph_holds}
\leanref{ReciprocityTheory/SemistableReduction/TorsionNonvanishStatement.lean}{ReciprocityTheory.MultiplicativeHeights.ApproxGraphStatement}
\begin{proof}
(1) Put $s=\sqrt r\,\mu|w|$, and view $w$ as the vector of $\mathbf C^r$ with $w_a=0$ for $a>p$. Then
$u=Rw+\varepsilon$ with $\varepsilon=\sum_aw_a(X_a-Re_a)$, and $|\varepsilon|\le\mu\sum_a|w_a|\le s$ by the
Cauchy--Schwarz inequality. So $|u^-|=|\varepsilon^-|\le s$ and $|u^+|\ge R|w|-s=4s/\rho-s\ge3s$. As $\rho\le1$,
$s\le\frac\rho3(4s/\rho-s)$, which gives $|u^-|\le\frac\rho3|u^+|$. If $w\ne0$, then $s>0$ and $|u^+|\ge3s>0$.

A reduction for (2) and (3). Let $y\in\mathbf C^r$ with $j(X_a,y)=0$ for every $a$. By (1), the linear map
$w\mapsto(\sum_aw_aX_a)^+$ from $\mathbf C^p$ to itself is injective, hence onto, so there is $u\in U$ with
$u^+=y^+$. Write $\langle\,,\rangle$ for the standard Hermitian products, conjugate-linear in the first
variable. Then $0=j(u,y)=\langle u^+,y^+\rangle-\langle u^-,y^-\rangle$, so
$|y^+|^2=\operatorname{Re}\langle u^-,y^-\rangle\le|u^-||y^-|\le\frac\rho3|y^+||y^-|$ by (1). Hence
$|y^+|\le\frac\rho3|y^-|\le\frac13|y^-|$.

(2) If $y\ne0$, then $y^-\ne0$ by the reduction, and $j(y,y)=|y^+|^2-|y^-|^2<0$.

(3) Write $v=y+u$ with $u\in U$. By (1), $|u^-|\le\frac13|u^+|$, and $j(u,y)=0$ gives
$\operatorname{Re}\langle y^+,u^+\rangle=\operatorname{Re}\langle y^-,u^-\rangle$. Hence
\[
 |v|^2=|y^+|^2+|u^+|^2+|y^-|^2+|u^-|^2+4\operatorname{Re}\langle y^-,u^-\rangle
 \ge|y^-|^2+|u^+|^2-\tfrac43|y^-||u^+|\ge\tfrac59|y^-|^2,
\]
the last step because $(|u^+|-\frac23|y^-|)^2\ge0$. With the reduction,
$|j(y,y)|=|y^-|^2-|y^+|^2\le|y^-|^2\le2|v|^2$.
\end{proof}
\leanref{ReciprocityTheory/SemistableReduction/ApproxGraphB.lean}{ReciprocityTheory.MultiplicativeHeights.ApproximationGraph.graph_bound}
\leanref{ReciprocityTheory/SemistableReduction/ApproxGraphB.lean}{ReciprocityTheory.MultiplicativeHeights.ApproximationGraph.exists_graph_pos}
\leanref{ReciprocityTheory/SemistableReduction/ApproxGraphC.lean}{ReciprocityTheory.MultiplicativeHeights.ApproximationGraph.pos_small_of_orth}

\begin{lemma}\label{cv-cusp-frame-covering}
In Situation~\ref{cv-cusp-setting}, let $\sigma$ be a good embedding, $Z$ a component at $\sigma$ and $x$ an
orthonormal frame of $Z$, with the notation of Definition~\ref{cv-decay-majorant}, and write $|u|$ for the
standard Hermitian norm of $u\in\mathbf C^r$. Let $L'\subset E^r$ be a finitely generated
$\mathcal O_E$-submodule that spans $E^r$ and is self-dual for $G_Z$: a vector $u\in E^r$ lies in $L'$ if and
only if $G_Z(u,u')\in\mathcal O_E$ for every $u'\in L'$. Let $\lambda_*>0$, and suppose that every nonzero
$v\in L'$ has $|x_\tau^{-1}\tau(v)|\ge\lambda_*$ for some $\tau\in\Phi$. Put $\mu=4(fr)^3|d_E|/\lambda_*$. Then:
\begin{enumerate}
\item for every family $(w_\tau)_{\tau\in\Phi}$ in $\mathbf C^r$ there is $l\in L'$ with
  $\sum_\tau|x_\tau^{-1}\tau(l)-w_\tau|^2\le\mu^2$;
\item $L'$ contains a basis $v_1,\dots,v_r$ of $E^r$ with $\sum_\tau|x_\tau^{-1}\tau(v_i)|^2\le\mu^2$ for every
  $i$.
\end{enumerate}
\end{lemma}
\leanref{Formalism/LocalToGlobal/ProjectiveLattices/ArchimedeanQuadraticData/FundamentalDomainD5.lean}{Formalism.L2.Pj.FundamentalDomainD5.d5_translation}
\leanref{ReciprocityTheory/ReductiveGroups/ArithmeticGeometry/FundamentalDomains/CoveringRadiusStatement.lean}{ReciprocityTheory.ChartRestriction.G.FD.CoveringRadiusStatement}
\begin{proof}
Let $V=E^r$, a $\mathbf Q$-vector space of dimension $N=2fr$. Let $W=(\mathbf C^r)^\Phi$, a real vector
space of dimension $N$ with the norm $|w|^2=\sum_\tau|w_\tau|^2$ and the inner product
$\langle w,w'\rangle=\operatorname{Re}\sum_\tau\sum_a\bar w_{\tau,a}w'_{\tau,a}$, and let $\iota:V\to W$ be the
map $v\mapsto(x_\tau^{-1}\tau(v))_\tau$. Let $\Theta$ be the isometry of $W$ that changes the signs of the
coordinates of index $a>q$ in each $w_\tau$. Since $x$ is orthonormal,
$\tau(G_Z)(x_\tau c,x_\tau c')=\sum_{a\le q}\bar c'_ac_a-\sum_{a>q}\bar c'_ac_a$ for $c,c'\in\mathbf C^r$, so
$\operatorname{Re}\tau(G_Z(v,y))$ is the $\tau$-term of $\langle\Theta\iota(y),\iota(v)\rangle$. As $\Phi$
contains one embedding of each conjugate pair, $\operatorname{Tr}_{E/\mathbf Q}a=2\sum_\tau\operatorname{Re}\tau(a)$
for $a\in E$, and
\[
 \langle\Theta\iota(y),\iota(v)\rangle=\tfrac12T(v,y),\qquad T(v,y)=\operatorname{Tr}_{E/\mathbf Q}G_Z(v,y)
 \qquad(v,y\in V).
\]
The form $T$ is nondegenerate: if $T(v,y)=0$ for every $v$, then, replacing $v$ by $av$ with $a\in E$,
$G_Z(v,y)=0$ for every $v$, and $y=0$ as $x_\tau^{\mathsf T}\tau(G_Z)\bar x_\tau$ is invertible.

Step 1: the lattice. The module $L'$ has a $\mathbf Z$-basis $e_1,\dots,e_N$, which is a $\mathbf Q$-basis
of $V$. The $\mathbf R$-linear extension of $\iota$ to $V\otimes\mathbf R$ is injective by the identity above,
as $T$ is nondegenerate, so $\iota(e_1),\dots,\iota(e_N)$ is an $\mathbf R$-basis of $W$ and $\Gamma=\iota(L')$ is a
lattice of rank $N$ in $W$.

Step 2: the dual lattice. Let $u\in W$ with $\langle u,\iota(l)\rangle\in\mathbf Z$ for every $l\in L'$. As $T$ is
nondegenerate, there is $y\in V$ with $T(e_k,y)=2\langle u,\iota(e_k)\rangle$ for every $k$; then
$\langle\Theta\iota(y),\iota(e_k)\rangle=\langle u,\iota(e_k)\rangle$ for every $k$, so $u=\Theta\iota(y)$, and
$T(l,y)\in2\mathbf Z$ for $l\in L'$. For $a\in\mathcal O_E$ and $l\in L'$,
$\operatorname{Tr}_{E/\mathbf Q}(a\,G_Z(l,y))=T(al,y)\in\mathbf Z$, so $G_Z(l,y)$ lies in the inverse different
of $E$, and $|d_E|\,G_Z(l,y)\in\mathcal O_E$. Hence $G_Z(|d_E|y,l)=|d_E|\,\overline{G_Z(l,y)}\in\mathcal O_E$ for
every $l\in L'$, and $|d_E|y\in L'$ by self-duality. If $u\ne0$, then $y\ne0$, and the hypothesis gives some
$\tau$ with $\lambda_*^2\le|x_\tau^{-1}\tau(|d_E|y)|^2\le|d_E|^2|\iota(y)|^2=|d_E|^2|u|^2$. So every nonzero
vector of the dual lattice $\Gamma^\vee$ has length at least $\mu'=\lambda_*/|d_E|$.

Step 3: short vectors. By Banaszczyk's transference theorem \cite[Theorem 2.1]{Ban93}, the product of the
first successive minimum of $\Gamma^\vee$ and the $N$-th successive minimum of $\Gamma$ is at most $N$. So
$\Gamma$ contains $N$ linearly independent vectors of length at most $\ell=2N/\mu'$: there are
$v_1,\dots,v_N\in L'$, linearly independent over $\mathbf Q$, with $|\iota(v_i)|\le\ell$.

Step 4: covering. Let $w\in W$, and write $w=\sum_ia_i\iota(v_i)$ with $a_i\in\mathbf R$, as the $\iota(v_i)$ form
an $\mathbf R$-basis of $W$. Choose integers $n_i$ with $|a_i-n_i|\le\frac12$, and put $l=\sum_in_iv_i\in L'$. Then
$|\iota(l)-w|\le\sum_i|a_i-n_i||\iota(v_i)|\le N\ell/2=4(fr)^2|d_E|/\lambda_*\le\mu$. This is (1).

Step 5: a basis. The $v_i$ span $V$ over $\mathbf Q$, hence over $E$, so $r$ of them form a basis of $E^r$.
Each has $\sum_\tau|x_\tau^{-1}\tau(v_i)|^2=|\iota(v_i)|^2\le\ell^2$, and $\ell=4fr|d_E|/\lambda_*\le\mu$. This
is (2).
\end{proof}
\leanref{Formalism/SharpRestriction/Divisibility/GeneratorPowers/PhiCMinF.lean}{Formalism.MainNotions.ModuliIntegralCharts.H1.PhiCMin.shortDual_holds}
\leanref{Formalism/LocalToGlobal/ProjectiveLattices/ArchimedeanQuadraticData/DotSelfEqSqNorm.lean}{Formalism.L2.Pj.FundamentalDomainD5.cover_round}

\begin{lemma}\label{cv-cusp-frame-basis}
In Situation~\ref{cv-cusp-setting}, let $\sigma$ be a good embedding, $Z$ a component at $\sigma$ and $x$ an
orthonormal frame of $Z$, with the notation of Definition~\ref{cv-decay-majorant}. For $u\in\mathbf C^r$ let
$u^+$ and $u^-$ be its first $q$ and its last $r-q$ coordinates, let $|u|$ be its standard Hermitian norm, and
put $j(u,u')=\sum_{a\le q}\bar u_au'_a-\sum_{a>q}\bar u_au'_a$. Let $L'\subset E^r$ be a finitely generated
$\mathcal O_E$-submodule with $G_Z(u,u')\in\mathcal O_E$ for all $u,u'\in L'$. Let $\mu>0$ and $0<\rho\le1$, put
$R=4\sqrt r\,\mu/\rho$, and suppose that:
\begin{enumerate}
\item[(i)] for every family $(w_\tau)_{\tau\in\Phi}$ in $\mathbf C^r$ there is $l\in L'$ with
  $\sum_\tau|x_\tau^{-1}\tau(l)-w_\tau|^2\le\mu^2$;
\item[(ii)] $L'$ contains a basis $v_1,\dots,v_r$ of $E^r$ with $\sum_\tau|x_\tau^{-1}\tau(v_i)|^2\le\mu^2$ for
  every $i$.
\end{enumerate}
Then $L'$ contains a basis $\lambda_1,\dots,\lambda_r$ of $E^r$ with the following properties.
\begin{enumerate}
\item $G_Z(\lambda_a,\lambda_b)=0$ for $a\le q<b$.
\item Let $\tau\in\Phi$ and $w\in\mathbf C^r$ with $w_a=0$ for $a>q$, and put
  $u=x_\tau^{-1}\sum_aw_a\tau(\lambda_a)$. Then $|u^-|\le\frac\rho3|u^+|$, and $j(u,u)>0$ if $w\ne0$.
\item Let $\tau\in\Phi$ and $w\in\mathbf C^r$ nonzero with $w_a=0$ for $a\le q$, and put
  $u=x_\tau^{-1}\sum_aw_a\tau(\lambda_a)$. Then $j(u,u)<0$.
\item For every $\tau\in\Phi$, $|\tau(G_Z(\lambda_a,\lambda_a))|\le(2R)^2$ if $a\le q$, and
  $|\tau(G_Z(\lambda_a,\lambda_a))|\le2(2R)^{4q}\mu^2$ if $a>q$.
\end{enumerate}
\end{lemma}
\leanref{ReciprocityTheory/SemistableReduction/LatAll.lean}{ReciprocityTheory.MultiplicativeHeights.Closed.approx}
\leanref{ReciprocityTheory/SemistableReduction/TorsionNonvanishStatement.lean}{ReciprocityTheory.MultiplicativeHeights.ApproxStatement}
\begin{proof}
Since $x$ is orthonormal, $\tau(G_Z)(u,u')=j(x_\tau^{-1}u',x_\tau^{-1}u)$ for $u,u'\in\mathbf C^r$, so
$\tau(G_Z(v,v'))=j(x_\tau^{-1}\tau(v'),x_\tau^{-1}\tau(v))$ for $v,v'\in E^r$.

Step 1: vectors near the axes. For $a\le q$, (i) applied to the family with $w_\tau=Re_a$ for every $\tau$ gives
$l_a\in L'$ with $\sum_\tau|X_{\tau,a}-Re_a|^2\le\mu^2$, where $X_{\tau,a}=x_\tau^{-1}\tau(l_a)$. So for each
$\tau$ the vectors $X_{\tau,1},\dots,X_{\tau,q}$ satisfy the hypothesis of Lemma~\ref{cv-cusp-frame-graph} with
$p=q$. Also $|X_{\tau,a}|\le R+\mu\le2R$, as $\mu\le R$.

Step 2: the Gram matrix. Let $M$ be the $q\times q$ matrix with entries $M_{ab}=G_Z(l_b,l_a)$, which lie in
$\mathcal O_E$, and let $\Delta=\det M$. Let $g\in E^q$ with $Mg=0$, and put $u=\sum_bg_bl_b$. Then
$G_Z(u,l_a)=0$ for every $a\le q$, so $G_Z(u,u)=0$. At any $\tau$, $x_\tau^{-1}\tau(u)=\sum_b\tau(g_b)X_{\tau,b}$,
and by Lemma~\ref{cv-cusp-frame-graph}(1) its value under $j$ is at least $(1-\rho^2/9)$ times the square of
the norm of its first $q$ coordinates, which is positive unless $g=0$. So $g=0$, $M$ is invertible and
$\Delta\ne0$.

Step 3: the basis. For $y\in E^r$ let $\beta(y)\in E^q$ have entries $\beta(y)_a=G_Z(y,l_a)$, let
$c(y)=\operatorname{adj}(M)\beta(y)$, and put $T(y)=\Delta y-\sum_bc(y)_bl_b$. As $Mc(y)=\Delta\beta(y)$,
$G_Z(T(y),l_a)=\Delta\beta(y)_a-(Mc(y))_a=0$ for $a\le q$. If $y\in L'$, then $\beta(y)$, $M$ and $\Delta$ are
integral, so $T(y)\in L'$. If $T(y)=0$, then $y=\Delta^{-1}\sum_bc(y)_bl_b$, so the kernel of $T$ lies in the span
of $l_1,\dots,l_q$, and $T$ has rank at least $r-q$. As $v_1,\dots,v_r$ span $E^r$, there are indices
$i_{q+1},\dots,i_r$ such that $T(v_{i_{q+1}}),\dots,T(v_{i_r})$ are linearly independent. Put $\lambda_a=l_a$ for
$a\le q$ and $\lambda_a=T(v_{i_a})$ for $a>q$. These lie in $L'$, and (1) holds. They are linearly independent:
if $\sum_ag_a\lambda_a=0$, then $u=\sum_{a\le q}g_al_a=-\sum_{a>q}g_a\lambda_a$ has $G_Z(u,l_b)=0$ for $b\le q$, so
$g_a=0$ for $a\le q$ by Step~2, and then $g_a=0$ for $a>q$.

Step 4: signs. In (2), $u=\sum_{a\le q}w_aX_{\tau,a}$, and Lemma~\ref{cv-cusp-frame-graph}(1) gives
$|u^-|\le\frac\rho3|u^+|$, and $u^+\ne0$ if $w\ne0$; then $j(u,u)=|u^+|^2-|u^-|^2\ge(1-\rho^2/9)|u^+|^2>0$. In (3),
$u$ is nonzero, as the $\lambda_a$ are linearly independent, and
$j(X_{\tau,b},u)=\sum_{a>q}w_a\tau(G_Z(\lambda_a,l_b))=0$ for $b\le q$ by (1); so $j(u,u)<0$ by
Lemma~\ref{cv-cusp-frame-graph}(2).

Step 5: sizes. For $a\le q$, $|\tau(G_Z(\lambda_a,\lambda_a))|=|j(X_{\tau,a},X_{\tau,a})|\le|X_{\tau,a}|^2\le(2R)^2$. For
$a>q$, let $v=v_{i_a}$, $y=x_\tau^{-1}\tau(T(v))$ and $v'=x_\tau^{-1}\tau(\Delta v)$. Then $j(X_{\tau,b},y)=0$ for
$b\le q$, and $v'-y=\sum_b\tau(c(v)_b)X_{\tau,b}$ lies in the span of the $X_{\tau,b}$. By
Lemma~\ref{cv-cusp-frame-graph}(3) and (ii),
\[
 |\tau(G_Z(\lambda_a,\lambda_a))|=|j(y,y)|\le2|v'|^2=2|\tau(\Delta)|^2|x_\tau^{-1}\tau(v)|^2\le2|\tau(\Delta)|^2\mu^2.
\]
The matrix $\tau(M)$ has entries $j(X_{\tau,a},X_{\tau,b})$. It is positive semidefinite, since
$j(u,u)=|u^+|^2-|u^-|^2\ge0$ on the span of the $X_{\tau,a}$ by Lemma~\ref{cv-cusp-frame-graph}(1), and its
diagonal entries are at most $(2R)^2$. By the inequality of arithmetic and geometric means for its eigenvalues,
$|\tau(\Delta)|\le(2R)^{2q}$, and (4) follows.
\end{proof}
\leanref{ReciprocityTheory/SemistableReduction/ApproxBasis.lean}{ReciprocityTheory.MultiplicativeHeights.approxBasis_holds}
\leanref{ReciprocityTheory/SemistableReduction/ApproxBasisB.lean}{ReciprocityTheory.MultiplicativeHeights.approxB_split}
\leanref{ReciprocityTheory/SemistableReduction/ApproxBasisA.lean}{ReciprocityTheory.MultiplicativeHeights.approxB_gram_posSemidef}
\leanref{ReciprocityTheory/SemistableReduction/ApproxBasisA.lean}{ReciprocityTheory.MultiplicativeHeights.approxB_det_le_of_posSemidef}

\begin{definition}\label{cv-cusp-frame-block}
In Situation~\ref{cv-cusp-setting}, let $\sigma$ be a good embedding and $Z$ a component at $\sigma$. For a
family $Y=(Y_\tau)_{\tau\in\Phi}$ of complex $(r-q)\times q$ matrices put $|Y|=\max_\tau\|Y_\tau\|$, where
$\|Y_\tau\|$ is the operator norm for the standard Hermitian norms, and for a frame $x$ of $Z$ let $x[Y]$ be the
family of the matrices
\[
 x[Y]_\tau=x_\tau\begin{pmatrix}1_q&0\\Y_\tau&1_{r-q}\end{pmatrix}\qquad(\tau\in\Phi).
\]
\end{definition}
\leanref{ReciprocityTheory/ReductiveGroups/ArithmeticGeometry/FundamentalDomains/IsotropicCoefficientSplitting.lean}{ReciprocityTheory.ChartRestriction.G.FD.Data.xZ}
\leanref{ReciprocityTheory/ReductiveGroups/ArithmeticGeometry/FundamentalDomains/IsotropicCoefficientSplitting.lean}{ReciprocityTheory.ChartRestriction.G.FD.Data.zN}

\begin{lemma}\label{cv-cusp-frame-lattice}
In Situation~\ref{cv-cusp-setting}, let $\sigma$ be a good embedding and $Z$ a component at $\sigma$, with the
notation of Definitions~\ref{cv-decay-majorant} and~\ref{cv-cusp-frame-block}. Let $z_0$ be a frame of $Z$ and
$x_0$ an orthonormal frame of $Z$ such that, for every $\tau\in\Phi$, the first $q$ columns of $z_{0,\tau}$ and
those of $x_{0,\tau}$ span the same subspace of $\mathbf C^r$. Let $L'\subset E^r$ be a finitely generated
$\mathcal O_E$-submodule that spans $E^r$ and is self-dual for $G_Z$: a vector $u\in E^r$ lies in $L'$ if and
only if $G_Z(u,u')\in\mathcal O_E$ for every $u'\in L'$. Let $0<\lambda_*\le1$ and $0<\rho\le1$, and suppose that
$\sum_\tau h_{z_0,\tau}(\tau(v))\ge\lambda_*^2$ for every nonzero $v\in L'$. Put
$\mu=4f^{7/2}r^3|d_E|/\lambda_*$ and $T=8\sqrt r\,\mu/\rho$. Then there are a family $Y$ with $|Y|\le\rho/3$ and a
basis $v_1,\dots,v_r$ of $E^r$ in $L'$ such that:
\begin{enumerate}
\item $G_Z(v_a,v_b)=0$ for $a\le q<b$;
\item for every $\tau\in\Phi$, $\tau(G_Z(v_a,v_a))>0$ for $a\le q$ and $\tau(G_Z(v_a,v_a))<0$ for $a>q$;
\item $|\tau(G_Z(v_a,v_a))|\le T^{2(2q+1)}$ for every $a$ and every $\tau\in\Phi$;
\item for every $\tau\in\Phi$, the first $q$ columns of $x_0[Y]_\tau$ span the same subspace of $\mathbf C^r$ as
  $\tau(v_1),\dots,\tau(v_q)$.
\end{enumerate}
\end{lemma}
\leanref{ReciprocityTheory/ReductiveGroups/ArithmeticGeometry/FundamentalDomains/LatticeBoundAtMaximumOfTransference.lean}{ReciprocityTheory.ChartRestriction.G.FD.latticeBoundAtMaximum_of_transference}
\leanref{ReciprocityTheory/ReductiveGroups/ArithmeticGeometry/FundamentalDomains/NonvanishingNearMaximum.lean}{ReciprocityTheory.ChartRestriction.G.FD.LatticeBoundAtMaximum}
\leanref{ReciprocityTheory/ReductiveGroups/ArithmeticGeometry/FundamentalDomains/NonvanishingNearMaximum.lean}{ReciprocityTheory.ChartRestriction.G.FD.LamBasis}
\begin{proof}
Write $|u|$ for the standard Hermitian norm and $u^\pm$, $j$ as in Lemma~\ref{cv-cusp-frame-basis}. The
majorant of $z_0$ at $\tau$ depends only on the span of the first $q$ columns of $z_{0,\tau}$, so
$h_{z_0,\tau}=h_{x_0,\tau}$, and $h_{x_0,\tau}(u)=|x_{0,\tau}^{-1}u|^2$ as $x_0$ is orthonormal. So every nonzero
$v\in L'$ has $\sum_\tau|x_{0,\tau}^{-1}\tau(v)|^2\ge\lambda_*^2$, and, as $\Phi$ has $f$ elements,
$|x_{0,\tau}^{-1}\tau(v)|\ge\lambda_*/\sqrt f$ for some $\tau$.

Step 1: a basis. The values $G_Z(u,u')$ with $u,u'\in L'$ are integral, by self-duality. By
Lemma~\ref{cv-cusp-frame-covering} for $x_0$, with $\lambda_*/\sqrt f$ in place of $\lambda_*$, the module $L'$
satisfies (i) and (ii) of Lemma~\ref{cv-cusp-frame-basis} with
$4(fr)^3|d_E|\sqrt f/\lambda_*=\mu$. That lemma, for $x_0$, gives a basis $v_1,\dots,v_r$ of $E^r$ in $L'$ with
its properties (1) to (4), where $2R=T$. Property (1) here is (1) there. For $a\le q$, (2) there with $w=e_a$
gives $\tau(G_Z(v_a,v_a))=j(u,u)>0$ for $u=x_{0,\tau}^{-1}\tau(v_a)$; for $a>q$, (3) there gives
$\tau(G_Z(v_a,v_a))<0$ in the same way. This is (2).

Step 2: sizes. As $\lambda_*/\sqrt f\le1$, $\mu\ge4$, and as $\rho\le1$, $T\ge8\mu\ge1$ and $2\mu^2\le T^2$. So
$T^2\le T^{4q+2}$ and $2T^{4q}\mu^2\le T^{4q+2}$, and (4) of Lemma~\ref{cv-cusp-frame-basis} gives (3).

Step 3: the chart. Fix $\tau$, and let $c_a=x_{0,\tau}^{-1}\tau(v_a)$ for $a\le q$. Let $A$ be the $q\times q$
matrix with columns $c_1^+,\dots,c_q^+$ and $B$ the $(r-q)\times q$ matrix with columns $c_1^-,\dots,c_q^-$. For
$w\in\mathbf C^q$, (2) of Lemma~\ref{cv-cusp-frame-basis} gives $|Bw|\le\frac\rho3|Aw|$, and $|Aw|^2>|Bw|^2$ if
$w\ne0$. So $A$ is invertible, and $Y_\tau=BA^{-1}$ has $\|Y_\tau\|\le\rho/3$. The first $q$ columns of
$x_0[Y]_\tau$ form the matrix $x_{0,\tau}\binom{1_q}{Y_\tau}$, and
$x_{0,\tau}\binom{1_q}{Y_\tau}A=x_{0,\tau}\binom AB$ has columns $\tau(v_1),\dots,\tau(v_q)$. As $A$ is invertible,
the two families span the same subspace. This is (4).
\end{proof}
\leanref{ReciprocityTheory/ReductiveGroups/ArithmeticGeometry/FundamentalDomains/SmallOperatorExistence.lean}{ReciprocityTheory.ChartRestriction.G.FD.SmallOperatorExistence.exists_Z}
\leanref{ReciprocityTheory/ReductiveGroups/ArithmeticGeometry/FundamentalDomains/IsotropicCombinations.lean}{ReciprocityTheory.ChartRestriction.G.FD.IsotropicCombinations.la_bound}

\begin{lemma}\label{cv-cusp-frame-chart}
In Situation~\ref{cv-cusp-setting}, let $\sigma$ be a good embedding, $Z$ a component at $\sigma$ and $x$ an
orthonormal frame of $Z$, with the notation of Definitions~\ref{cv-decay-majorant}
and~\ref{cv-cusp-frame-block}. Let $Y$ be a family with $|Y|<1$. Then $x[Y]$ is a frame of $Z$. If moreover
$|Y|\le\rho$ with $0<\rho\le1/10$, then:
\begin{enumerate}
\item $|\log\beta_\lambda(x[Y])-\log\beta_\lambda(x)|\le4wr|Y|^2$ for every $w>0$ and every weight $\lambda$ on $Z$ of
  spread at most $w$;
\item $\frac{1-\rho}{1+\rho}h_{x,\tau}(u)\le h_{x[Y],\tau}(u)\le\frac{1+\rho}{1-\rho}h_{x,\tau}(u)$ for every
  $\tau\in\Phi$ and $u\in\mathbf C^r$.
\end{enumerate}
\end{lemma}
\leanref{ReciprocityTheory/ReductiveGroups/ArithmeticGeometry/FundamentalDomains/SignTypeCancelIdentity.lean}{ReciprocityTheory.ChartRestriction.G.FD.OrthonormalChartBounds.xZ_mem_dom}
\leanref{ReciprocityTheory/ReductiveGroups/ArithmeticGeometry/FundamentalDomains/LogDOfPosBound.lean}{ReciprocityTheory.ChartRestriction.G.FD.OrthonormalChartBounds.nfac_xZ}
\leanref{ReciprocityTheory/ReductiveGroups/ArithmeticGeometry/FundamentalDomains/MajXZ.lean}{ReciprocityTheory.ChartRestriction.G.FD.OrthonormalChartBounds.maj_xZ}

This is a self-contained computation. As $x$ is orthonormal, $x[Y]_\tau^{\mathsf T}\tau(G_Z)\overline{x[Y]_\tau}$ is the matrix
$\begin{pmatrix}1-Y_\tau^{\mathsf T}\bar Y_\tau&-Y_\tau^{\mathsf T}\\-\bar Y_\tau&-1\end{pmatrix}$, whose leading minors have the signs required
of a frame when $\|Y_\tau\|<1$, and whose minors differ from those of the diagonal matrix of $x$ by factors
$1+O(\|Y_\tau\|^2)$. For (2), the first $q$ columns of $x[Y]_\tau$ span the graph of $Y_\tau$, in the
coordinates given by the columns of $x_\tau$, over the span of the first $q$ of them.

\begin{lemma}\label{cv-cusp-frame-nonvanishing}
In Situation~\ref{cv-cusp-setting}, let $\sigma$ be a good embedding and $Z$ a component at $\sigma$, with the
notation of Definitions~\ref{cv-decay-majorant} and~\ref{cv-cusp-frame-block}. Let $w\ge1$, let $\lambda$ be a
weight on $Z$ of spread at most $w$, and let $P$ be a form of weight $\lambda$ on $Z$. Let $z_0$ be a frame of $Z$
with $|P|(z_0)>0$ and $|P|(y)\le2|P|(z_0)$ for every frame $y$ of $Z$, and let $b\in B$ be such that
$x_0=z_0b$ is orthonormal. Then $P(x_0[Y])\ne0$ for every family $Y$ with $|Y|\le1/(50\sqrt{wr})$.
\end{lemma}
\leanref{ReciprocityTheory/ReductiveGroups/ArithmeticGeometry/FundamentalDomains/NonvanishingNearHalfMaxHolds.lean}{ReciprocityTheory.ChartRestriction.G.FD.nonvanishingNearHalfMax_holds}
\leanref{ReciprocityTheory/ReductiveGroups/ArithmeticGeometry/FundamentalDomains/NonvanishingNearHalfMax.lean}{ReciprocityTheory.ChartRestriction.G.FD.NonvanishingNearHalfMax}
\begin{proof}
Put $\rho_0=1/(10\sqrt{wr})$. As $wr\ge1$, $\rho_0\le1/10$, and $4wr\rho_0^2=1/25$. Since $|P|$ is invariant
under $B$, $|P|(x_0)=|P|(z_0)>0$, so $P(x_0)\ne0$, and $|P|\le2|P|(x_0)$ on all frames of $Z$. Suppose that
$P(x_0[Y])=0$ with $|Y|\le\rho_0/5$. For $t\in\mathbf C$ with $|t|<5$, $|tY|<\rho_0$, so by
Lemma~\ref{cv-cusp-frame-chart}, $x_0[tY]$ is a frame of $Z$ and
$\beta_\lambda(x_0[tY])\ge e^{-1/25}\beta_\lambda(x_0)$. Hence
\[
 |P(x_0[tY])|=\frac{|P|(x_0[tY])}{\beta_\lambda(x_0[tY])}\le\frac{2e^{1/25}|P|(x_0)}{\beta_\lambda(x_0)}
 =2e^{1/25}|P(x_0)|<3|P(x_0)|,
\]
as $e^{1/25}<3/2$. The function $g(t)=P(x_0[tY])$ is holomorphic on $|t|<5$, as the entries of $x_0[tY]$ are
affine in $t$ and $P$ is holomorphic on the frames. On the disc $|t-1|<4$ it satisfies $|g|<3|P(x_0)|$, and
$g(1)=0$. By the Schwarz lemma, $|g(t)|\le\frac34|P(x_0)|\,|t-1|$ on that disc. At $t=0$, where
$g(0)=P(x_0)$, this gives $|P(x_0)|\le\frac34|P(x_0)|$, a contradiction.
\end{proof}

\begin{lemma}\label{cv-cusp-frame}
In Situation~\ref{cv-cusp-setting}, let $\sigma$ be a good embedding and $Z$ a component at $\sigma$, with the
notation of Definition~\ref{cv-decay-majorant}. Let $w\ge1$, let $\lambda$ be a weight on $Z$ of spread at most
$w$, let $P$ be a form of weight $\lambda$ on $Z$, and let $z_0$ be a frame of $Z$ with $|P|(z_0)>0$ and
$|P|(y)\le2|P|(z_0)$ for every frame $y$ of $Z$. Let $L'\subset E^r$ be a finitely generated
$\mathcal O_E$-submodule that spans $E^r$ and is self-dual for $G_Z$: a vector $u\in E^r$ lies in $L'$ if and
only if $G_Z(u,u')\in\mathcal O_E$ for every $u'\in L'$. Let $0<\lambda_*\le1$, and suppose that
$\sum_\tau h_{z_0,\tau}(\tau(v))\ge\lambda_*^2$ for every nonzero $v\in L'$. Put
\[
 \ell=\Big(\frac{1600}3f^{7/2}r^4w^{1/2}|d_E|\,\lambda_*^{-1}\Big)^{2q+1}.
\]
Then there are a frame $x$ of $Z$ with $P(x)\ne0$ and a basis $v_1,\dots,v_r$ of $E^r$ in $L'$ such that:
\begin{enumerate}
\item $G_Z(v_a,v_b)=0$ for $a\le q<b$;
\item for every $\tau\in\Phi$, $\tau(G_Z(v_a,v_a))>0$ for $a\le q$ and $\tau(G_Z(v_a,v_a))<0$ for $a>q$;
\item $|\tau(G_Z(v_a,v_a))|\le\ell^2$ for every $a$ and every $\tau\in\Phi$;
\item for every $\tau\in\Phi$, the first $q$ columns of $x_\tau$ span the same subspace of $\mathbf C^r$ as
  $\tau(v_1),\dots,\tau(v_q)$;
\item $h_{x,\tau}(u)\le2h_{z_0,\tau}(u)$ for every $\tau\in\Phi$ and $u\in\mathbf C^r$.
\end{enumerate}
\end{lemma}
\leanref{ReciprocityTheory/ReductiveGroups/ArithmeticGeometry/FundamentalDomains/NonvanishingNearHalfMax.lean}{ReciprocityTheory.ChartRestriction.G.FD.FrameAtHalfMaximum}
\leanref{Formalism/LocalToGlobal/ProjectiveLattices/ArchimedeanQuadraticData/D6P4Statement.lean}{Formalism.L2.Pj.FundamentalDomainD5.t41H_statement}
\begin{proof}
Put $\rho_0=1/(10\sqrt{wr})\le1/10$ and $\rho=3\rho_0/5$. The leading minors of
$z_{0,\tau}^{\mathsf T}\tau(G_Z)\bar z_{0,\tau}$ are nonzero, so the Gram--Schmidt process along the first $i$ columns,
$i=1,\dots,r$, gives $b\in B$ for which $x_0=z_0b$ is orthonormal. As $b_\tau$ and $b_\tau^{-1}$ are upper
triangular, the first $q$ columns of $z_{0,\tau}$ and of $x_{0,\tau}$ span the same subspace, so
$U^+_{x_0,\tau}=U^+_{z_0,\tau}$. The majorant at $\tau$ of a frame is defined through this span alone
(Definition~\ref{cv-decay-majorant}(1)), so $h_{x_0,\tau}=h_{z_0,\tau}$.

By Lemma~\ref{cv-cusp-frame-lattice} with this $\rho$ there are a family $Y$ with $|Y|\le\rho/3=\rho_0/5$ and a
basis $v_1,\dots,v_r$ of $E^r$ in $L'$ with (1) and (2), with (3) for the bound $T^{2(2q+1)}$, and with (4) for
$x=x_0[Y]$. Here $T=8\sqrt r\,\mu/\rho$ with $\mu=4f^{7/2}r^3|d_E|/\lambda_*$ and $\rho=3/(50\sqrt{wr})$, so
$T=\frac{1600}3f^{7/2}r^4w^{1/2}|d_E|/\lambda_*$ and $T^{2(2q+1)}=\ell^2$. By Lemma~\ref{cv-cusp-frame-chart}
with $\rho_0/5\le1/10$ in place of $\rho$, $x$ is a frame of $Z$ and
$h_{x,\tau}\le\frac{1+\rho_0/5}{1-\rho_0/5}h_{x_0,\tau}\le2h_{z_0,\tau}$, which is (5). By
Lemma~\ref{cv-cusp-frame-nonvanishing}, $P(x)\ne0$.
\end{proof}
\leanref{ReciprocityTheory/ReductiveGroups/ArithmeticGeometry/FundamentalDomains/MaxPointFrameOfTransference.lean}{ReciprocityTheory.ChartRestriction.G.FD.maxPointFrame_of_lower_bound}
\leanref{ReciprocityTheory/ReductiveGroups/ArithmeticGeometry/FundamentalDomains/MaxPointFrameOfTransference.lean}{ReciprocityTheory.ChartRestriction.G.FD.maxPointFrame_of_transference}
\leanref{ReciprocityTheory/ReductiveGroups/ArithmeticGeometry/FundamentalDomains/OrthonormalFlagTranslateHolds.lean}{ReciprocityTheory.ChartRestriction.G.FD.orthonormalFlagTranslate_holds}
\leanref{ReciprocityTheory/ReductiveGroups/ArithmeticGeometry/FundamentalDomains/FlagSpanBounds.lean}{ReciprocityTheory.ChartRestriction.G.FD.FlagSpanBounds.sameFlag_rmul}

\subsection{Metrized lines and good CM points}\label{subsection-cv-good-cm-points}

\begin{definition}\label{cv-cusp-cm-lines}
Let $K$ be a number field, and write $\mathrm N(I)$ for the absolute norm of a fractional ideal $I$ of $K$. A
\emph{metrized line} over $K$ is a pair $L=(I_L,n(L))$ of a nonzero fractional ideal $I_L$ of $K$ and a family
$n(L)=(n_s(L))_s$ of positive real numbers, one for each embedding $s\colon K\to\mathbf C$. Its \emph{degree}
and its \emph{normalized degree} are
\[
 \deg L=-\log\mathrm N(I_L)-\sum_s\log n_s(L),\qquad\widehat\deg L=\frac{\deg L}{[K:\mathbf Q]}.
\]
The \emph{dual} of $L$ is $L^\vee=(I_L^{-1},(n_s(L)^{-1})_s)$, the \emph{tensor product} of $L$ with a
metrized line $L'$ is $L\otimes L'=(I_LI_{L'},(n_s(L)n_s(L'))_s)$, and $L^k=(I_L^k,(n_s(L)^k)_s)$ for integers
$k\ge0$. As $\mathrm N$ is multiplicative, $\deg(L\otimes L')=\deg L+\deg L'$, $\deg L^\vee=-\deg L$ and
$\deg L^k=k\deg L$. For a nonzero prime $\mathfrak p$ of $\mathcal O_K$, a \emph{local generator} of $L$ at
$\mathfrak p$ is an element $c\in K^\times$ with $\operatorname{ord}_{\mathfrak p}c=\operatorname{ord}_{\mathfrak p}I_L$.
\end{definition}
\leanref{Formalism/LocalToGlobal/CMDeterminantIdeals/DetIdeal.lean}{Formalism.L2.LamCM.MLine}
\leanref{Formalism/LocalToGlobal/CMDeterminantIdeals/DetIdeal.lean}{Formalism.L2.LamCM.MLine.deg}
\leanref{Formalism/LocalToGlobal/CMDeterminantIdeals/DetIdeal.lean}{Formalism.L2.LamCM.MLine.ndeg}
\leanref{Formalism/LocalToGlobal/CMDeterminantIdeals/DetIdeal.lean}{Formalism.L2.LamCM.IsLocGen}
\leanref{Formalism/LocalToGlobal/CMDeterminantIdeals/CMIdealValuations/OneLeAbsNormOfLeOne.lean}{Formalism.L2.LamCM.degree_algebraic}

\begin{lemma}\label{cv-cusp-cm-cauchy}
Let $V$ be a finite set, let $G$ be a polynomial over $\mathbf C$ in variables $z_v$, $v\in V$, and let
$\mu\in\mathbf N^V$, with $c_\mu$ the coefficient of $z^\mu=\prod_vz_v^{\mu_v}$ in $G$. Let $\rho_v>0$ for
$v\in V$ and let $S$ be a real number such that $|G(z)|\le S$ whenever $|z_v|\le\rho_v$ for every $v$. Then
$|c_\mu|\prod_v\rho_v^{\mu_v}\le S$.
\end{lemma}
\leanref{Formalism/LocalToGlobal/CMDeterminantIdeals/StatementC.lean}{Formalism.L2.LamCM.CoeffCauchy}
\begin{proof}
This is Cauchy's estimate for the coefficients of a polynomial, obtained by averaging over roots of unity. Let
$N$ be larger than $\mu_v$ for every $v$ and than every exponent of a variable in a monomial of $G$, and put
$\zeta=e^{2\pi i/N}$. For $k\in\{0,\dots,N-1\}^V$ let $z_k$ be the point with coordinates $\rho_v\zeta^{k_v}$.
For integers $0\le a,b<N$ the sum $\sum_{j=0}^{N-1}\zeta^{j(a-b)}$ is $N$ if $a=b$ and $0$ otherwise. Expanding
$G$ in monomials therefore gives
\[
 \sum_kG(z_k)\prod_v\zeta^{-k_v\mu_v}=N^{|V|}c_\mu\prod_v\rho_v^{\mu_v}.
\]
The left side is a sum of $N^{|V|}$ terms of absolute value at most $S$.
\end{proof}
\leanref{Formalism/LocalToGlobal/CMDeterminantIdeals/CMIdealValuations/AbsNormSpanGalois.lean}{Formalism.L2.LamCM.coeff_cauchy_bound}

\begin{lemma}\label{cv-cusp-cm-numbers}
Let $K$ be a number field, let $g\ge0$ be an integer, and let $\Delta$, $D$ and $\mathcal L_1,\dots,\mathcal L_g$
be metrized lines over $K$ (Definition~\ref{cv-cusp-cm-lines}). Let $\ell\ge1$ and $t$ be real numbers, put
$p_t=(1+|t|)/2$, and suppose that:
\begin{enumerate}
\item[(i)] $\sum_j\deg\mathcal L_j-g[K:\mathbf Q]\log\ell\le\deg\Delta$;
\item[(ii)] $\widehat\deg\mathcal L_j=t/2$ for every $j$;
\item[(iii)] $\widehat\deg D\le\frac12\widehat\deg\Delta^\vee+\frac g4\log2\pi$.
\end{enumerate}
Let $V$ be a finite set, $c\colon V\to\{1,\dots,g\}$ a map and $n\ge0$ an integer, and let $G$ be a polynomial
over $K$ in variables $z_v$, $v\in V$. Let $\mu\in\mathbf N^V$ and $|\mu|=\sum_v\mu_v$, and suppose that the
coefficient $c_\mu$ of $z^\mu$ in $G$ is nonzero and lies in $I_{\mathcal N}$, where
$\mathcal N=D^n\otimes\big(\bigotimes_v\mathcal L_{c(v)}^{\mu_v}\big)^\vee$. Let $R>0$, and let $S_s>0$ for every
embedding $s\colon K\to\mathbf C$ be such that
\[
 |s(G)(z)|\,n_s(D)^n\le S_sR^{|\mu|}\qquad\text{whenever }|z_v|\,n_s(\mathcal L_{c(v)})\le1\text{ for every }v,
\]
where $s(G)$ is $G$ with $s$ applied to its coefficients. Then
\[
 -\big(p_t+\tfrac12\log\ell\big)(|\mu|+gn)-|\mu|\log R\le\frac1{[K:\mathbf Q]}\sum_s\log S_s.
\]
\end{lemma}
\leanref{Formalism/LocalToGlobal/CMDeterminantIdeals/CoeffMemOfLoc.lean}{Formalism.L2.LamCM.LamNumber0}
\begin{proof}
Step 1: the degree of $D$. By (ii), $\deg\mathcal L_j=\frac t2[K:\mathbf Q]$, so (i) gives
$\widehat\deg\Delta\ge g(\frac t2-\log\ell)$. As $\widehat\deg\Delta^\vee=-\widehat\deg\Delta$, (iii) gives
$\widehat\deg D\le g(-\frac t4+\frac12\log\ell+\frac14\log2\pi)$. Since $\log2\pi\le2$ and $-t\le|t|$, this is
at most $g(p_t+\frac12\log\ell)$.

Step 2: the degree of $\mathcal N$. By the rules for degrees and (ii),
$\widehat\deg\mathcal N=n\widehat\deg D-|\mu|\frac t2$. As $\log\ell\ge0$, $-\frac t2\le p_t+\frac12\log\ell$, and
Step 1 gives $\widehat\deg\mathcal N\le(p_t+\frac12\log\ell)(|\mu|+gn)$.

Step 3: the sizes of $c_\mu$. Fix $s$. Lemma~\ref{cv-cusp-cm-cauchy}, applied to $s(G)$ with
$\rho_v=n_s(\mathcal L_{c(v)})^{-1}$ and the bound $S_sR^{|\mu|}n_s(D)^{-n}$, gives
\[
 |s(c_\mu)|\,n_s(\mathcal N)=|s(c_\mu)|\,n_s(D)^n\prod_vn_s(\mathcal L_{c(v)})^{-\mu_v}\le S_sR^{|\mu|}.
\]

Step 4: the product formula. As $c_\mu$ is a nonzero element of $I_{\mathcal N}$, $c_\mu\mathcal O_K\subset I_{\mathcal N}$,
so $\mathrm N(I_{\mathcal N})\le\mathrm N(c_\mu\mathcal O_K)=\prod_s|s(c_\mu)|$. With Step 3,
\[
 -\deg\mathcal N=\log\mathrm N(I_{\mathcal N})+\sum_s\log n_s(\mathcal N)\le\sum_s\log\big(|s(c_\mu)|\,n_s(\mathcal N)\big)
 \le\sum_s\log S_s+[K:\mathbf Q]\,|\mu|\log R.
\]
Dividing by $[K:\mathbf Q]$ and using Step 2 gives the claim.
\end{proof}
\leanref{Formalism/LocalToGlobal/CMDeterminantIdeals/CMIdealValuations/AbsNormSpanGalois.lean}{Formalism.L2.LamCM.lambda_number_base}
\leanref{Formalism/LocalToGlobal/CMDeterminantIdeals/CMIdealValuations/AbsNormSpanGalois.lean}{Formalism.L2.LamCM.degree_delta_base}
\leanref{Formalism/LocalToGlobal/CMDeterminantIdeals/CMIdealValuations/AbsNormSpanGalois.lean}{Formalism.L2.LamCM.degree_norm_base}
\leanref{Formalism/LocalToGlobal/CMDeterminantIdeals/CMIdealValuations/AbsNormSpanGalois.lean}{Formalism.L2.LamCM.lambda_endomorphism}

\begin{lemma}\label{cv-cusp-cm-hadamard}
Let $K$ be a number field, let $g\ge0$ be an integer, let $\Delta$ and $\mathcal L_1,\dots,\mathcal L_g$ be
metrized lines over $K$ (Definition~\ref{cv-cusp-cm-lines}), let $\ell>0$, and let $U\in GL_g(K)$ have rows
$U_1,\dots,U_g$. For an embedding $s\colon K\to\mathbf C$ write $s(U_j)\in\mathbf C^g$ for the column vector of
the $s(U_{jk})$, and let $H_s$ be a positive definite Hermitian $g\times g$ matrix with
$n_s(\Delta)=(\det H_s)^{1/2}$ and $s(U_j)^*H_s\,s(U_j)\le(\ell\,n_s(\mathcal L_j))^2$ for every $j$. Suppose that
at every nonzero prime $\mathfrak p$ of $\mathcal O_K$ there are local generators $c_\Delta$ of $\Delta$ and $c_j$ of
$\mathcal L_j$ with $\operatorname{ord}_{\mathfrak p}(c_1\cdots c_g\det U)\ge\operatorname{ord}_{\mathfrak p}c_\Delta$.
Then
\[
 \sum_j\deg\mathcal L_j-g[K:\mathbf Q]\log\ell\le\deg\Delta.
\]
\end{lemma}
\leanref{Formalism/LocalToGlobal/CMDeterminantIdeals/HadamardLoc0.lean}{Formalism.L2.LamCM.HadamardLoc0}
\begin{proof}
Step 1: the ideals. Let $J=(\det U)\prod_jI_{\mathcal L_j}$. At every nonzero prime $\mathfrak p$,
$\operatorname{ord}_{\mathfrak p}J=\operatorname{ord}_{\mathfrak p}(c_1\cdots c_g\det U)\ge
\operatorname{ord}_{\mathfrak p}c_\Delta=\operatorname{ord}_{\mathfrak p}I_\Delta$. So $J\subset I_\Delta$, and
$\mathrm N(I_\Delta)\le\mathrm N(J)=\prod_j\mathrm N(I_{\mathcal L_j})\prod_s|s(\det U)|$.

Step 2: Hadamard's inequality. Fix $s$, and let $\Gamma$ be the $g\times g$ matrix with entries
$s(U_i)^*H_s\,s(U_j)$. It is $A^*H_sA$ for the matrix $A$ with columns $s(U_1),\dots,s(U_g)$, which is
invertible, so $\Gamma$ is positive definite and $\det\Gamma=|s(\det U)|^2\det H_s$. Hadamard's inequality
$\det\Gamma\le\prod_j\Gamma_{jj}$ and the bounds on the $\Gamma_{jj}$ give
\[
 \log|s(\det U)|+\log n_s(\Delta)\le g\log\ell+\sum_j\log n_s(\mathcal L_j).
\]

Step 3: summing the inequalities of Step 2 over $s$ and adding the logarithm of the inequality of Step 1,
\[
 \log\mathrm N(I_\Delta)+\sum_s\log n_s(\Delta)\le\sum_j\Big(\log\mathrm N(I_{\mathcal L_j})+\sum_s\log n_s(\mathcal L_j)\Big)
 +g[K:\mathbf Q]\log\ell,
\]
which is the claim.
\end{proof}
\leanref{Formalism/LocalToGlobal/CMDeterminantIdeals/CMIdealValuations/DetValuationLeOne.lean}{Formalism.L2.LamCM.hadamard_local_base}
\leanref{Formalism/LocalToGlobal/CMDeterminantIdeals/CMIdealValuations/DetValuationLeOne.lean}{Formalism.L2.LamCM.prod_span_le_of_loc}
\leanref{Formalism/LocalToGlobal/CMDeterminantIdeals/CMIdealValuations/DetValuationLeOne.lean}{Formalism.L2.LamCM.hadamard_of_le}

\begin{definition}\label{cv-cusp-cm-key}
Let $k$ be a field of characteristic $0$ and $K$ a number field.
\begin{enumerate}
\item A \emph{key datum} over $k$ is a triple $(A,\theta,v)$ of a principally polarized abelian variety $A$ over
  $k$ of dimension $g\ge1$, with a symmetric line bundle that defines the polarization and is rigidified along
  the zero section, a generator $\theta$ of its space of sections, and a basis $v=(v_1,\dots,v_g)$ of the Lie
  algebra of $A$, as in Theorem~\ref{cv-definite-key-formula}. For $M\in GL_g(k)$ let $Mv$ be the basis whose
  $i$-th vector is $\sum_jM_{ij}v_j$.
\item Let $(A,\theta,v)$ be a key datum over $K$ and $\Theta$, $\Delta$ metrized lines over $K$
  (Definition~\ref{cv-cusp-cm-lines}). They are \emph{adapted} to it at a nonzero prime $\mathfrak p$ of
  $\mathcal O_K$ if there are a field $K'$ of characteristic $0$, a discrete valuation ring $R\subset K'$ with
  fraction field $K'$, an embedding $\iota\colon K\to K'$ with
  $\iota^{-1}(R)=\{y\in K:\operatorname{ord}_{\mathfrak p}y\ge0\}$, local generators $c_\Theta$ of $\Theta$ and
  $c_\Delta$ of $\Delta$ at $\mathfrak p$, units $u_\Theta$ and $u_\Delta$ of $R$, and $M\in GL_g(K')$ with
  $\det M=\iota(c_\Delta)u_\Delta$, such that the key datum $(A_{K'},\iota(c_\Theta)u_\Theta\theta,Mv)$, where
  $A_{K'}$ and its data are the base change along $\iota$, extends over $R$: to a principally polarized abelian
  scheme over $R$ with a rigidified symmetric line bundle, a generator over $R$ of its sections and a basis over
  $R$ of its Lie algebra.
\item They are \emph{adapted} to it at an embedding $s\colon K\to\mathbf C$ if the following holds. Write the
  base change of $A$ along $s$ as $\mathbf C^g/\Lambda_s$, let $H$ be the Hermitian form of its polarization, let
  $\theta$ also denote the function on $\mathbf C^g$ that represents the generator, let $W$ be the matrix whose
  $i$-th row is the vector $v_i\in\mathbf C^g$, and let $F$ be a fundamental domain for $\Lambda_s$. Then
  \[
   n_s(\Theta)^2=\frac1{\operatorname{vol}F}\int_F|\theta(z)|^2e^{-\pi H(z,z)}\,dz,\qquad
   n_s(\Delta)^2=\frac{\pi^g|\det W|^2}{\operatorname{vol}F}.
  \]
\end{enumerate}
\end{definition}
\leanref{ReciprocityTheory/ReductiveGroups/ArithmeticGeometry/BirationalGeometry/ToroidalCompactification/IsLieBasis.lean}{ReciprocityTheory.ChartRestriction.G.Geo.Tor.KeyIn}
\leanref{Formalism/LocalToGlobal/CMDeterminantIdeals/StatementB.lean}{Formalism.L2.LamCM.KeyIntegral}
\leanref{Formalism/LocalToGlobal/CMDeterminantIdeals/StatementB.lean}{Formalism.L2.LamCM.KeyLocAt}
\leanref{Formalism/LocalToGlobal/CMDeterminantIdeals/StatementB.lean}{Formalism.L2.LamCM.KeyArchAt}

\begin{lemma}\label{cv-cusp-cm-theta-degree}
Let $K$ be a number field, let $(A,\theta,v)$ be a key datum of dimension $g$ over $K$, and let $\Theta$ and
$\Delta$ be metrized lines over $K$ adapted to it at every nonzero prime of $\mathcal O_K$ and at every embedding
$K\to\mathbf C$ (Definition~\ref{cv-cusp-cm-key}). Then
\[
 \deg\Theta^\vee=\tfrac12\deg\Delta^\vee+\tfrac g4[K:\mathbf Q]\log2\pi.
\]
\end{lemma}
\leanref{Development/AbelianVarietyTheory/Proofs/CubicalLineBundles/ThetaLineDualDegreeOfKeyFormulaPowerFormStatement.lean}{GlobalTheory.T.thetaLine_dual_degree_of_keyFormulaPowerForm}
\begin{proof}
Let $\kappa$ be the rule of Theorem~\ref{cv-definite-key-formula}, and put $k=\kappa(A,\theta,v)\in K^\times$.

Step 1: the ideal of $k$. Let $\mathfrak p$ be a nonzero prime of $\mathcal O_K$, with the data $K'$, $R$,
$\iota$, $c_\Theta$, $c_\Delta$, $u_\Theta$, $u_\Delta$ and $M$ of Definition~\ref{cv-cusp-cm-key}(2), and put
$a=\iota(c_\Theta)u_\Theta$. By parts (1) and (2) of the theorem,
$a^8\kappa(A_{K'},a\theta,Mv)=\det(M)^4\iota(k)$, and by part (3), applied to the extension over $R$,
$\kappa(A_{K'},a\theta,Mv)$ is a unit $u$ of $R$. Hence $\iota(kc_\Delta^4c_\Theta^{-8})=u_\Theta^8uu_\Delta^{-4}$
is a unit of $R$. As $\iota^{-1}(R)$ is the valuation ring of $\mathfrak p$,
$\operatorname{ord}_{\mathfrak p}k=8\operatorname{ord}_{\mathfrak p}c_\Theta-4\operatorname{ord}_{\mathfrak p}c_\Delta
=8\operatorname{ord}_{\mathfrak p}I_\Theta-4\operatorname{ord}_{\mathfrak p}I_\Delta$. This holds at every
$\mathfrak p$, so $k\mathcal O_K=I_\Theta^8I_\Delta^{-4}$ and
$\sum_s\log|s(k)|=\log\mathrm N(k\mathcal O_K)=8\log\mathrm N(I_\Theta)-4\log\mathrm N(I_\Delta)$.

Step 2: the infinite places. Let $s\colon K\to\mathbf C$ be an embedding. By part (2) of the theorem, $\kappa$ of
the base change of $(A,\theta,v)$ along $s$ is $s(k)$, and part (5), with Definition~\ref{cv-cusp-cm-key}(3),
gives $|s(k)|\,n_s(\Theta)^8n_s(\Delta)^{-4}=(2\pi)^{2g}$.

Step 3: taking logarithms in Step 2, summing over $s$ and using Step 1,
\[
 8\log\mathrm N(I_\Theta)-4\log\mathrm N(I_\Delta)+8\sum_s\log n_s(\Theta)-4\sum_s\log n_s(\Delta)
 =2g[K:\mathbf Q]\log2\pi,
\]
that is, $-8\deg\Theta+4\deg\Delta=2g[K:\mathbf Q]\log2\pi$. As $\deg\Theta^\vee=-\deg\Theta$ and
$\deg\Delta^\vee=-\deg\Delta$, this is the claim.
\end{proof}
\leanref{Development/AbelianVarietyTheory/Proofs/CubicalLineBundles/ThetaLineDualDegreeOfKeyFormulaPowerForm.lean}{GlobalTheory.thetaLine_dual_deg_of_keyFormulaPow}
\leanref{Formalism/LocalToGlobal/CMDeterminantIdeals/CMIdealValuations/DetValuationLeOne.lean}{Formalism.L2.LamCM.keyDegCore_core}

\begin{lemma}\label{cv-cusp-cm-qr}
Let $B$ be a finite set, let $m_b\ge0$ be integers for $b\in B$, and let $a_{b,i}\ge0$ be integers for $b\in B$
and $1\le i\le m_b$; put $|a|=\sum_{b,i}a_{b,i}$. A complex function $P$ on the families $X=(X_b)_{b\in B}$ of
complex $m_b\times m_b$ matrices has the \emph{transformation law with exponents $a$} if
\[
 P(X')=\prod_it_{ii}^{a_{b,i}}\,P(X)
\]
whenever every $X_{b'}$ is invertible, $b\in B$, $t\in GL_{m_b}(\mathbf C)$ is upper triangular and $X'$ is $X$
with $X_b$ replaced by $X_bt$. Let $P$ have this law, and let $C$ and $R\ge0$ be real numbers such that
$|P(Q)|\le C$ for every family $Q$ of unitary matrices. Then $|P(X)|\le CR^{|a|}$ for every family $X$ of
invertible matrices whose columns all have Hermitian norm at most $R$.
\end{lemma}
\leanref{ReciprocityTheory/SemistableReduction/TorsionNonvanishStatement.lean}{ReciprocityTheory.MultiplicativeHeights.LemmaHStatement}
\begin{proof}
The Gram--Schmidt process applied to the columns of $X_b$ gives $X_b=Q_bT_b$ with $Q_b$ unitary and $T_b$ upper
triangular, and $|(T_b)_{jj}|$ is at most the norm of the $j$-th column of $X_b$, hence at most $R$. As $X_b$ is
invertible, so is $T_b$. Starting from the family $Q$ and replacing $Q_b$ by $Q_bT_b$ for one $b$ after the other,
the transformation law gives $P(X)=\prod_{b,i}(T_b)_{ii}^{a_{b,i}}P(Q)$. Hence $|P(X)|\le R^{|a|}C$.
\end{proof}
\leanref{ReciprocityTheory/SemistableReduction/LemmaH.lean}{ReciprocityTheory.MultiplicativeHeights.lemmaH_holds}

\begin{lemma}\label{cv-cusp-cm-archimedean}
Let $B$ be a finite set, let $m_b\ge0$ be integers for $b\in B$, let $a_{b,i}\ge0$ be integers for $b\in B$ and
$1\le i\le m_b$, with $|a|=\sum_{b,i}a_{b,i}$, and let $r$ be an integer with $m_b\le r$ for every $b$. Let $P$ be
a complex polynomial in the entries of families $X=(X_b)_{b\in B}$ of $m_b\times m_b$ matrices that has the
transformation law with exponents $a$ (Lemma~\ref{cv-cusp-cm-qr}). For each $b$ let $H_b$ be a positive definite
Hermitian $m_b\times m_b$ matrix, let $C$ and $\ell\ge0$ be real numbers, and let $\nu_{b,k}>0$ for $b\in B$ and
$1\le k\le m_b$, such that $(H_b)_{kk}\le(\ell\nu_{b,k})^2$ for all $b$ and $k$ and $|P(X)|\le C$ for every family
$X$ with $X_b^*H_bX_b=1$ for every $b$. Then $|P(Z)|\le C(r\ell)^{|a|}$ for every family $Z$ with
$|Z_{b,ki}|\,\nu_{b,k}\le1$ for all $b$, $k$ and $i$.
\end{lemma}
\leanref{Formalism/LocalToGlobal/CMDeterminantIdeals/StatementG.lean}{Formalism.L2.LamCM.ArchTie}
\begin{proof}
By the spectral theorem, $H_b=S_b^*S_b$ for an invertible matrix $S_b$. Put $P'(Y)=P((S_b^{-1}Y_b)_b)$. Then $P'$
has the transformation law with exponents $a$, as $(S_b^{-1}Y_b)t=S_b^{-1}(Y_bt)$. If every $Q_b$ is unitary,
then $X_b=S_b^{-1}Q_b$ satisfies $X_b^*H_bX_b=Q_b^*Q_b=1$, so $|P'(Q)|\le C$.

Let $Z$ be as in the statement, and put $Y_b=S_bZ_b$, so that $P'(Y)=P(Z)$. The $i$-th column of $Y_b$ is
$\sum_kZ_{b,ki}S_be_k$, and $|S_be_k|^2=e_k^*H_be_k=(H_b)_{kk}\le(\ell\nu_{b,k})^2$; so its norm is at most
$\sum_k|Z_{b,ki}|\,\ell\nu_{b,k}\le m_b\ell\le r\ell$. For all but finitely many $u\in\mathbf C$ every $Y_b+u\cdot1$
is invertible, and its columns have norm at most $r\ell+|u|$, so Lemma~\ref{cv-cusp-cm-qr} gives
$|P'((Y_b+u\cdot1)_b)|\le C(r\ell+|u|)^{|a|}$. As $P'$ is a polynomial, letting $u$ tend to $0$ gives
$|P(Z)|=|P'(Y)|\le C(r\ell)^{|a|}$.
\end{proof}
\leanref{Formalism/LocalToGlobal/CMDeterminantIdeals/CMIdealValuations/FiniteDetAddSmulOneEqZero.lean}{Formalism.L2.LamCM.ArchTie_of}

\begin{lemma}\label{cv-cusp-cm-integral}
Let $V$ be a complete discrete valuation ring with finite residue field and fraction field $K'$. An
\emph{extension} of $V$ is a complete discrete valuation ring $W$ with finite residue field, together with a
local homomorphism $V\to W$ and an embedding of $K'$ in the fraction field $K_W$ of $W$ that extends it. Let $B$
be a finite set and $m_b\ge0$ integers for $b\in B$, and consider polynomials over $K'$ in the entries $y_{b,ki}$,
$1\le k,i\le m_b$, of families $Y=(Y_b)_{b\in B}$ of $m_b\times m_b$ matrices.
\begin{enumerate}
\item Let $P$ be such a polynomial with $P(X)\in W$ for every extension $W$ of $V$ and every family $X$ of
  matrices over $W$ with $\det X_b\in W^\times$ for every $b$. Then the coefficients of $P$ lie in $V$.
\item Let $G$ be such a polynomial, let $\tau\in K'^\times$, and let $U_b\in GL_{m_b}(K')$ and $c_{b,k}\in K'$ be
  such that $c_{b,j}(U_b)_{jk}\in V$ for all $b$, $j$ and $k$. Suppose that $G(Y)\in\tau W$ for every extension
  $W$ of $V$ and every family $Y$ over $K_W$ for which every $U_b^{\mathsf T}Y_b$ has entries in $W$ and
  determinant in $W^\times$. Then the coefficient $c_\mu$ of every monomial
  $y^\mu=\prod_{b,k,i}y_{b,ki}^{\mu_{b,ki}}$ in $G$ satisfies $c_\mu\prod_{b,k,i}c_{b,k}^{\mu_{b,ki}}\in\tau V$.
\end{enumerate}
\end{lemma}
\leanref{Formalism/LocalToGlobal/CMDeterminantIdeals/StatementG.lean}{Formalism.L2.LamCM.LemmaG}
\leanref{Formalism/LocalToGlobal/CMDeterminantIdeals/StatementS.lean}{Formalism.L2.LamCM.CoeffIntegralLoc}
\begin{proof}
(1) Let $d$ be the total degree of $P$. There is an extension $W$ of $V$ whose residue field $k_W$ has more than
$d+\sum_bm_b$ elements, an unramified extension of $V$ of large degree \cite[Chapter~III, \S5]{Ser79}. As $V\to W$
is local, $K'\cap W=V$ in $K_W$: if $y\in K'$ does not lie in $V$, then $y^{-1}$ lies in the maximal ideal of $V$,
so its image is not a unit of $W$ and $y\notin W$. So it suffices to show that the coefficients of $P$ lie in
$W$. Suppose that they do not, let $\varpi$ be a uniformizer of $W$, and let $e\ge1$ be the least integer for
which $Q=\varpi^eP$ has coefficients in $W$. Then some coefficient of $Q$ is a unit, so the reduction $\bar Q$ of
$Q$ over $k_W$ is nonzero. The polynomial $\bar Q\prod_b\det Y_b$ over $k_W$ is nonzero, of total degree at most
$d+\sum_bm_b<|k_W|$, so it does not vanish at some family of matrices over $k_W$. A lift $X$ of this family to $W$
has $\det X_b\in W^\times$ for every $b$ and $Q(X)\in W^\times$. But $Q(X)=\varpi^eP(X)\in\varpi W$ by hypothesis,
a contradiction.

(2) Let $P$ be the polynomial with $P(X)=\tau^{-1}G\big(((U_b^{-1})^{\mathsf T}X_b)_b\big)$. Let $W$ be an extension
of $V$ and $X$ a family over $W$ with $\det X_b\in W^\times$ for every $b$. The family $Y_b=(U_b^{-1})^{\mathsf T}X_b$
has $U_b^{\mathsf T}Y_b=X_b$, so $G(Y)\in\tau W$ and $P(X)\in W$. By (1), $P$ has coefficients in $V$. Now
$G(Y)=\tau P((U_b^{\mathsf T}Y_b)_b)$, so, with $D_b$ the diagonal matrix with entries $c_{b,1},\dots,c_{b,m_b}$,
\[
 G\big((D_bY_b)_b\big)=\tau P\big((U_b^{\mathsf T}D_bY_b)_b\big).
\]
The matrix $U_b^{\mathsf T}D_b$ has entries $c_{b,j}(U_b)_{jk}$ in $V$, so the right side is $\tau$ times a
polynomial in the $y_{b,ki}$ with coefficients in $V$. The coefficient of $y^\mu$ on the left is
$c_\mu\prod_{b,k,i}c_{b,k}^{\mu_{b,ki}}$.
\end{proof}
\leanref{Formalism/LocalToGlobal/CMDeterminantIdeals/CMIdealValuations/ExistsEvalNeZeroOfTotalDegreeLt.lean}{Formalism.L2.LamCM.LemmaG_of}
\leanref{Formalism/LocalToGlobal/CMDeterminantIdeals/CMIdealValuations/ExistsEvalNeZeroOfTotalDegreeLt.lean}{Formalism.L2.LamCM.coeffIntegralOfUnitEval}
\leanref{Formalism/LocalToGlobal/CMDeterminantIdeals/CMIdealValuations/LargeResidueExtensionLemma.lean}{Formalism.L2.LamCM.coeff_integral_local}

\begin{proposition}\label{cv-cusp-cm-product}
Let $K$ be a number field, let $(A,\theta,v)$ be a key datum of dimension $g$ over $K$, and let $\Theta$ and
$\Delta$ be metrized lines over $K$ adapted to it at every nonzero prime of $\mathcal O_K$ and at every embedding
$K\to\mathbf C$ (Definition~\ref{cv-cusp-cm-key}). Let $B$ be a finite set with integers $m_b\ge0$, and let $e$
be a bijection from the set of pairs $(b,j)$ with $b\in B$ and $1\le j\le m_b$ onto $\{1,\dots,g\}$. Let $t$ be a
real number and $\mathcal L_1,\dots,\mathcal L_g$ metrized lines over $K$ with $\widehat\deg\mathcal L_j=t/2$ for
every $j$, and let $U\in GL_g(K)$ have rows $U_1,\dots,U_g$; for an embedding $s$ write $s(U_j)$ for the column
vector of the $s(U_{jk})$. Let $a_{b,i}\ge0$, for $b\in B$ and $1\le i\le m_b$, $n\ge0$ and $r\ge1$ be integers
with $m_b\le r$ for every $b$, put $|a|=\sum_{b,i}a_{b,i}$, let $\ell\ge1$, and let $S_s>0$ for every embedding
$s\colon K\to\mathbf C$. Let $G$ be a nonzero polynomial over $K$ in the entries $y_{b,ki}$ of families
$Y=(Y_b)_{b\in B}$ of $m_b\times m_b$ matrices, and write $\iota(G)$ and $s(G)$ for $G$ with $\iota$ or $s$
applied to its coefficients. Suppose:
\begin{enumerate}
\item[(L)] for every nonzero prime $\mathfrak p$ of $\mathcal O_K$ there are local generators $d$ of
  $\Theta^\vee$, $c_\Delta$ of $\Delta$ and $c_j$ of $\mathcal L_j$ at $\mathfrak p$
  (Definition~\ref{cv-cusp-cm-lines}) with
  $\operatorname{ord}_{\mathfrak p}(c_1\cdots c_g\det U)\ge\operatorname{ord}_{\mathfrak p}c_\Delta$, a complete
  discrete valuation ring $V$ with finite residue field and fraction field $K'$, an embedding
  $\iota\colon K\to K'$ with $\iota^{-1}(V)=\{y\in K:\operatorname{ord}_{\mathfrak p}y\ge0\}$, and matrices
  $U'_b\in GL_{m_b}(K')$ with $\iota(c_{e(b,j)})(U'_b)_{jk}\in V$ for all $b$, $j$ and $k$, such that
  $\iota(G)(Y)\in\iota(d)^nW$ for every extension $W$ of $V$ (Lemma~\ref{cv-cusp-cm-integral}) and every family
  $Y$ over $K_W$ for which every $U'^{\mathsf T}_bY_b$ has entries in $W$ and determinant in $W^\times$;
\item[(A)] for every embedding $s\colon K\to\mathbf C$ there is a positive definite Hermitian $g\times g$ matrix
  $H_s$ with $n_s(\Delta)=(\det H_s)^{1/2}$ and $s(U_j)^*H_s\,s(U_j)\le(\ell\,n_s(\mathcal L_j))^2$ for every $j$,
  such that $s(G)$ has the transformation law with exponents $a$ (Lemma~\ref{cv-cusp-cm-qr}) and
  $|s(G)(Y)|\,n_s(\Theta^\vee)^n\le S_s$ for every family $Y$ with $Y_b^*\Gamma_{s,b}Y_b=1$ for every $b$; here
  $\Gamma_{s,b}$ is the $m_b\times m_b$ matrix with entries $s(U_{e(b,j)})^*H_s\,s(U_{e(b,j')})$.
\end{enumerate}
Then
\[
 -\Big(\frac{1+|t|}2+\frac12\log\ell\Big)(|a|+gn)-|a|\log(r\ell)\le\frac1{[K:\mathbf Q]}\sum_s\log S_s.
\]
\end{proposition}
\leanref{Development/AbelianVarietyTheory/Proofs/CubicalLineBundles/CmPointLowerBoundFree.lean}{GlobalTheory.cmPoint_lowerBound_free}
\leanref{Development/AbelianVarietyTheory/Proofs/CubicalLineBundles/CmPointLowerBoundOfCoefficientBoundStatement.lean}{GlobalTheory.T.cmPoint_lower_bound_of_coefficient_bound}
\leanref{Formalism/LocalToGlobal/CMDeterminantIdeals/StatementS.lean}{Formalism.L2.LamCM.Step5Loc}
\leanref{Formalism/LocalToGlobal/CMDeterminantIdeals/StatementS.lean}{Formalism.L2.LamCM.Step5Arch}
\begin{proof}
As $G\ne0$, some monomial $y^\mu=\prod_{b,k,i}y_{b,ki}^{\mu_{b,ki}}$ has a nonzero coefficient $c_\mu$ in $G$; put
$|\mu|=\sum_{b,k,i}\mu_{b,ki}$.

Step 1: homogeneity. Fix an embedding $s_0$. By (A), $s_0(G)$ has the transformation law with exponents $a$.
Applying it with the upper triangular matrix $2\cdot1$ at each $b$ in turn gives $s_0(G)(2Y)=2^{|a|}s_0(G)(Y)$ for every family $Y$ of
invertible matrices. So the polynomial $\big(s_0(G)(2Y)-2^{|a|}s_0(G)(Y)\big)\prod_b\det Y_b$ vanishes at every
family, hence is zero, and as $\prod_b\det Y_b\ne0$, $s_0(G)(2Y)=2^{|a|}s_0(G)(Y)$ as polynomials. The coefficients
of $y^\mu$ give $2^{|\mu|}s_0(c_\mu)=2^{|a|}s_0(c_\mu)$, so $|\mu|=|a|$.

Step 2: $c_\mu$ lies in $I_{\mathcal N}$, where
$\mathcal N=(\Theta^\vee)^n\otimes\big(\bigotimes_{b,k,i}\mathcal L_{e(b,k)}^{\mu_{b,ki}}\big)^\vee$. Let $\mathfrak p$ be a
nonzero prime, with the data of (L). Lemma~\ref{cv-cusp-cm-integral}(2), for $\iota(G)$ with
$\tau=\iota(d)^n$, the matrices $U'_b$ and $c_{b,k}=\iota(c_{e(b,k)})$, gives
$\iota\big(c_\mu\prod_{b,k,i}c_{e(b,k)}^{\mu_{b,ki}}\big)\in\iota(d)^nV$. As $\iota^{-1}(V)$ is the valuation ring
of $\mathfrak p$,
\[
 \operatorname{ord}_{\mathfrak p}c_\mu\ge n\operatorname{ord}_{\mathfrak p}d-\sum_{b,k,i}\mu_{b,ki}\operatorname{ord}_{\mathfrak p}c_{e(b,k)}
 =\operatorname{ord}_{\mathfrak p}I_{\mathcal N}.
\]
This holds at every $\mathfrak p$, so $c_\mu\in I_{\mathcal N}$.

Step 3: the archimedean bound. Fix $s$. As $H_s$ is positive definite and the vectors $s(U_{e(b,j)})$ are
linearly independent, $\Gamma_{s,b}$ is positive definite, and its diagonal entries are
$s(U_{e(b,k)})^*H_s\,s(U_{e(b,k)})\le(\ell\,n_s(\mathcal L_{e(b,k)}))^2$. Lemma~\ref{cv-cusp-cm-archimedean}, with
$P=s(G)$, $H_b=\Gamma_{s,b}$, $\nu_{b,k}=n_s(\mathcal L_{e(b,k)})$ and $C=S_s\,n_s(\Theta^\vee)^{-n}$, gives
\[
 |s(G)(Z)|\,n_s(\Theta^\vee)^n\le S_s(r\ell)^{|a|}\qquad\text{whenever }|Z_{b,ki}|\,n_s(\mathcal L_{e(b,k)})\le1
 \text{ for all }b,k,i.
\]

Step 4: the degrees. By Lemma~\ref{cv-cusp-cm-theta-degree}, dividing by $[K:\mathbf Q]$,
$\widehat\deg\Theta^\vee=\frac12\widehat\deg\Delta^\vee+\frac g4\log2\pi$. By Lemma~\ref{cv-cusp-cm-hadamard},
with the matrices $H_s$ of (A) and the local generators $c_\Delta$ and $c_j$ of (L),
$\sum_j\deg\mathcal L_j-g[K:\mathbf Q]\log\ell\le\deg\Delta$.

Step 5: Lemma~\ref{cv-cusp-cm-numbers}, with $D=\Theta^\vee$, the variables $y_{b,ki}$ with $c(y_{b,ki})=e(b,k)$,
the coefficient $c_\mu$, $R=r\ell$, and Steps 2 to 4, gives the claim with $|\mu|$ in place of $|a|$. By Step 1,
$|\mu|=|a|$.
\end{proof}
\leanref{Development/AbelianVarietyTheory/Proofs/CubicalLineBundles/CmPointLowerBoundOfCoefficientBound.lean}{GlobalTheory.cmLowerBound_of_coeffBound}
\leanref{Development/AbelianVarietyTheory/Proofs/CubicalLineBundles/ThetaPolynomialLowerBoundOfThetaLineDegree.lean}{GlobalTheory.thetaLine_lower_bound_of_degree}
\leanref{Formalism/LocalToGlobal/CMDeterminantIdeals/CMIdealValuations/FiniteDetAddSmulOneEqZero.lean}{Formalism.L2.LamCM.subst_homogeneous}
\leanref{Formalism/LocalToGlobal/CMDeterminantIdeals/CMIdealValuations/AbsNormSpanGalois.lean}{Formalism.L2.LamCM.coeff_mem_of_local}

\begin{lemma}\label{cv-cusp-cm-point-polynomial}
In Situation~\ref{cv-cusp-setting}, recall that $\tau(E)\subset F_c$ for every embedding $\tau\colon E\to\mathbf C$,
and that $m_\tau=p$ for $\tau\in\Phi$ and $m_\tau=q$ otherwise. A \emph{framed point} of $X$ over an $F_c$-algebra
$S$ is a point $y$ of $X$ with values in $S$ together with, for every $\tau$, a basis
$u_\tau=(u_{\tau,1},\dots,u_{\tau,m_\tau})$ of the eigenspace of $\mathcal O_E$ on the Lie algebra of $A_y$ on which
$\mathcal O_E$ acts through $\tau$, and a generator $\vartheta$ of the module of sections of the line bundle of the
polarization of $A_y$. For a family $g=(g_\tau)$ of matrices $g_\tau\in GL_{m_\tau}(S)$ and $c\in S^\times$, the
framed point $y\cdot(g,c)$ has the point of $y$, the bases whose $i$-th vectors are $\sum_j(g_\tau)_{ji}u_{\tau,j}$,
and the generator $c\vartheta$; put $y\cdot g=y\cdot(g,1)$. A homomorphism of $F_c$-algebras $\varphi\colon S\to S'$
carries $y$ to its base change, a framed point over $S'$. Let $\kappa=(\alpha,n)$ consist of integers
$\alpha_{\tau,i}\ge0$, for all $\tau$ and $0\le i<m_\tau$, and $n\ge0$, and put $a_{\tau,i}=\alpha_{\tau,i-1}$ for
$1\le i\le m_\tau$. Let $f$ be a rational section of $X$ of weight $\kappa$, such as an element of $S_\kappa$ when
$\kappa$ is dominant. It assigns to every framed point $y$ over an $F_c$-algebra $S$ a value $f(y)\in S$,
compatibly with base change, and it satisfies the transformation law of its weight: over a field,
$f(y\cdot(T,c))=c^n\prod_{\tau,i}T_{\tau,ii}^{a_{\tau,i}}f(y)$ for every family $T$ of invertible upper triangular
matrices and every $c\ne0$. Let $K\supset F_c$ be a number field and $y$ a framed point of $X$ over $K$, and consider
polynomials in the entries of families $Y=(Y_\tau)$ of $m_\tau\times m_\tau$ matrices.
\begin{enumerate}
\item There is a polynomial $G$ over $K$ such that $f(y_S\cdot Y)=\varphi(G)(Y)$ for every $F_c$-algebra $S$, every
  homomorphism of $F_c$-algebras $\varphi\colon K\to S$ and every family $Y$ of matrices over $S$ with every
  $\det Y_\tau$ a unit; here $y_S$ is the base change of $y$ along $\varphi$, and $\varphi(G)$ is $G$ with $\varphi$
  applied to its coefficients.
\item Let $G$ be as in (1). For every embedding $s\colon K\to\mathbf C$, $s(G)$ has the transformation law with
  exponents $a$ (Lemma~\ref{cv-cusp-cm-qr}).
\item Let $G$ be as in (1), let $\sigma\colon F_c\to\mathbf C$ be an embedding, $s\colon K\to\mathbf C$ an
  embedding that extends $\sigma$, and $w$ a framed point of $X$ over $\mathbf C$, an $F_c$-algebra through
  $\sigma$, whose point is the image of the point of $y$ under $s$. If $f(w)\ne0$, then $G\ne0$.
\end{enumerate}
\end{lemma}
\leanref{Development/AbelianVarietyTheory/Proofs/CubicalLineBundles/ExistsIsFramePolynomialOfRuleNumberFieldStatement.lean}{GlobalTheory.T.exists_isFramePolynomialOfRule_numberField}
\leanref{Development/AbelianVarietyTheory/Proofs/CubicalLineBundles/FramePolynomialEquivariantAndNeZeroStatement.lean}{GlobalTheory.T.framePolynomial_equivariant_and_ne_zero}
\leanref{Development/AbelianVarietyTheory/Definitions/CMLattices/IntegralOf.lean}{GlobalTheory.LamT.IsFrPoly}
\begin{proof}
(1) Step 1: a polynomial with a denominator. Let $A$ be the $K$-algebra of polynomials in the entries of a family
$\mathcal X=(\mathcal X_\tau)$ of $m_\tau\times m_\tau$ matrices, with $\prod_\tau\det\mathcal X_\tau$ inverted. Let
$u_0$ be the base change of $y$ to $A$, and $u=u_0\cdot\mathcal X$. The value $f(u)\in A$ is
$P(\mathcal X)/(\prod_\tau\det\mathcal X_\tau)^N$ for a polynomial $P$ over $K$ and an integer $N\ge0$. Let $S$,
$\varphi$ and $Y$ be as in (1), and let $k\colon A\to S$ be the homomorphism of $F_c$-algebras that extends
$\varphi$ and sends $\mathcal X$ to $Y$. Base change of framed points is unique and transitive, so $k$ carries
$u_0$ to $y_S$; it commutes with changes of frame, so $k$ carries $u$ to $y_S\cdot Y$; and the values of $f$ commute
with base change. Hence
\[
 \Big(\prod_\tau\det Y_\tau\Big)^Nf(y_S\cdot Y)=k\Big(\Big(\prod_\tau\det\mathcal X_\tau\Big)^Nf(u)\Big)
 =\varphi(P)(Y).\tag{$*$}
\]

Step 2: the law of $P$. Let $X$ be a family over $K$ with every $\det X_\tau\ne0$, let $\tau_0$ be an embedding,
$t\in GL_{m_{\tau_0}}(K)$ upper triangular, $X'$ the family $X$ with $X_{\tau_0}$ replaced by $X_{\tau_0}t$, and $T$
the family with $T_{\tau_0}=t$ and $T_\tau=1$ for $\tau\ne\tau_0$. Changes of frame compose, so
$y\cdot X'=(y\cdot X)\cdot T$, and the transformation law gives $f(y\cdot X')=\prod_it_{ii}^{a_{\tau_0,i}}f(y\cdot X)$.
By $(*)$ for $S=K$ and the identity of $K$, and as $\det(X_{\tau_0}t)=\det X_{\tau_0}\prod_it_{ii}$,
\[
 P(X')=\prod_it_{ii}^{a_{\tau_0,i}+N}\,P(X).
\]

Step 3: the polynomial $G$. As $K$ is infinite, Lemma~\ref{cv-volume-forms-triangular}, with the exponents $a$ and
$n_\tau=N$ for every $\tau$, gives a polynomial $G$ with $P=(\prod_\tau D_\tau)^NG$, where $D_\tau$ is the polynomial
$\det Y_\tau$. Then $\varphi(P)(Y)=(\prod_\tau\det Y_\tau)^N\varphi(G)(Y)$ for all $S$, $\varphi$ and $Y$ as in (1),
and as $\prod_\tau\det Y_\tau$ is a unit of $S$, $(*)$ gives $f(y_S\cdot Y)=\varphi(G)(Y)$.

(2) Let $z$ be the base change of $y$ to $\mathbf C$ along $s$, and let $X$, $\tau_0$, $t$, $X'$ and $T$ be as in
Step 2, over $\mathbf C$. By (1), $f(z\cdot X)=s(G)(X)$ and $f(z\cdot X')=s(G)(X')$. As $z\cdot X'=(z\cdot X)\cdot T$,
the transformation law gives $s(G)(X')=\prod_it_{ii}^{a_{\tau_0,i}}s(G)(X)$.

(3) Let $z$ be the base change of $y$ to $\mathbf C$ along $s$. The framed points $z$ and $w$ have the same point,
so there are a family $h$ of invertible complex matrices and $c\in\mathbf C^\times$ with
$w=z\cdot(h,c)=(z\cdot h)\cdot(1,c)$. By the transformation law and (1), $f(w)=c^nf(z\cdot h)=c^ns(G)(h)$. So
$G=0$ would give $f(w)=0$.
\end{proof}
\leanref{Development/AbelianVarietyTheory/Proofs/CubicalLineBundles/ExistsIsFramePolynomialOfRuleNumberField.lean}{GlobalTheory.LamT.exists_isFrPoly}
\leanref{Development/AbelianVarietyTheory/Proofs/CubicalLineBundles/ExistsPolyDivDet.lean}{GlobalTheory.LamT.RatFn.exists_poly_div_det}
\leanref{Development/AbelianVarietyTheory/Proofs/CubicalLineBundles/EvalFrEqProdDetPowMul.lean}{GlobalTheory.LamT.evalFr_eq_prod_det_pow_mul}
\leanref{Development/AbelianVarietyTheory/Proofs/CubicalLineBundles/FramePolynomialEquivariantAndNeZero.lean}{GlobalTheory.LamT.isFrPoly_equiv_and_ne_zero}
% Lean: Step 1 is RatFn.exists_poly_div_det (the generic ring RatFn.exists_genericRing,
% Frames.FramedBaseChange.exists_frMap, existsUnique_frChg, frMap_unique_refl_iff, ChgBC.frChg_baseChange and
% evR_frMap); Step 2 is the hypothesis passed to evalFr_eq_prod_det_pow_mul in exists_isFrPoly
% (frChg_refl_and_comp and evR_frChg_mul_triangular); Step 3 is evalFr_eq_prod_det_pow_mul
% (FieldPoly.exists_eq_prod_blockDetPoly_pow_mul). (2) is evalFr_mul_tri; (3) uses frChg_exists_unique,
% existsUnique_frChg and evR_frChg_tri.

\begin{definition}\label{cv-cusp-cm-point-good}
In Situation~\ref{cv-cusp-setting}, where $r=p+q$, use the framed points of
Lemma~\ref{cv-cusp-cm-point-polynomial}, let $K\supset F_c$ be a number field, and for an embedding
$\tau\colon E\to\mathbf C$ write $\tau_K\colon E\to K$ for $\tau$ followed by the inclusion $F_c\subset K$.
\begin{enumerate}
\item Let $A$ be an abelian scheme over $\mathcal O_K$ and $v=(v_1,\dots,v_h)$ a basis of the Lie algebra of its
  generic fibre. A \emph{model matrix} of $(A,v)$ at a nonzero prime $\mathfrak p$ of $\mathcal O_K$ is a matrix
  $M\in GL_h(K)$ for which there are a principally polarized abelian scheme over the localization
  $\mathcal O_{K,\mathfrak p}$, with a rigidified symmetric line bundle, a generator of its sections and a basis
  $v^{\mathfrak p}$ of its Lie algebra, whose abelian scheme with its line bundle is the base change of $A$, such that
  $v^{\mathfrak p}_k=\sum_lM_{kl}v_l$ over $K$ for every $k$.
\item Let $\ell>0$. A \emph{good CM point} of $X$ over $K$ with bound $\ell$ consists of the following data, with
  the following properties.
  \begin{enumerate}
  \item[(a)] (The point.) A point $o$ of $X$ with values in $K$, with abelian variety $A_o$, and a key datum
    $(A_o,\theta,v)$ over $K$ (Definition~\ref{cv-cusp-cm-key}) of dimension $g=fr$, whose abelian variety, with its
    action of $\mathcal O_E$, its polarization and its rigidified line bundle, is that of $o$.
  \item[(b)] (Integral models.) Principally polarized abelian schemes $B$, $A^\Phi$ and $A'$ over $\mathcal O_K$
    with actions of $\mathcal O_E$, for each of which the Rosati involution of the polarization induces complex
    conjugation on $\mathcal O_E$; the generic fibre of $B$, with its action, polarization and line bundle, is
    $A_o$. Key data $(A^\Phi_K,\theta^\Phi,v^\Phi)$ and $(A'_K,\theta',v')$ of dimension $f$ over $K$, whose abelian
    varieties with their line bundles are the generic fibres of $A^\Phi$ and $A'$, and bijections $\tau^\Phi$ from
    $\{1,\dots,f\}$ onto $\Phi$ and $\tau'$ from $\{1,\dots,f\}$ onto the set of the other embeddings, such that every
    $b\in\mathcal O_E$ multiplies $v^\Phi_i$ by $\tau^\Phi(i)_K(b)$ and $v'_i$ by $\tau'(i)_K(b)$.
  \item[(c)] (The isogeny.) Homomorphisms $\varphi_j\colon A^\Phi\to B$ for $1\le j\le p$ and
    $\varphi'_j\colon A'\to B$ for $1\le j\le q$ that commute with $\mathcal O_E$, morphisms $\psi_j\colon B\to A^\Phi$
    and $\psi'_j\colon B\to A'$, and an integer $N\ge1$, such that on points with values in any $\mathcal O_K$-scheme
    \[
     \sum_j\varphi_j\psi_j(b)+\sum_j\varphi'_j\psi'_j(b)=Nb,\qquad
     \psi_i(\beta)=Na_i,\qquad\psi'_i(\beta)=Na'_i,\qquad
     \beta=\sum_j\varphi_j(a_j)+\sum_j\varphi'_j(a'_j),
    \]
    for all points $b$ of $B$, $a_j$ of $A^\Phi$ and $a'_j$ of $A'$, and all $i$.
  \item[(d)] (Polarizations.) Elements $h_j$, $1\le j\le p$, and $h'_j$, $1\le j\le q$, of $\mathcal O_E$ fixed by
    complex conjugation, with $|\tau(h_j)|\le\ell^2$ and $|\tau(h'_j)|\le\ell^2$ for every $\tau$, such that the
    polarization of $B$ pulls back along $\varphi_j$ to that of $A^\Phi$ composed with the action of $h_j$, and along
    $\varphi'_j$ to that of $A'$ composed with the action of $h'_j$.
  \item[(e)] (Lines.) Metrized lines $\mathcal M^\Phi_i$ and $\mathcal M'_i$, $1\le i\le f$, $\Theta$ and $\Delta$
    over $K$ (Definition~\ref{cv-cusp-cm-lines}). For every $i$ and every nonzero prime $\mathfrak p$ there are a
    model matrix $M$ of $(A^\Phi,v^\Phi)$ at $\mathfrak p$ and a local generator $c$ of $\mathcal M^\Phi_i$ at
    $\mathfrak p$ such that, for $c'\in K$, the entries of $c'e_iM^{-1}$, with $e_i$ the $i$-th unit row, are integral
    at $\mathfrak p$ exactly when $\operatorname{ord}_{\mathfrak p}c'\ge\operatorname{ord}_{\mathfrak p}c$; and for
    every embedding $s\colon K\to\mathbf C$, $n_s(\mathcal M^\Phi_i)^2$ is the squared Hodge norm
    (Definition~\ref{cv-cm-point-periods}(4)) of $v^\Phi_i$ in the base change of $A^\Phi_K$ along $s$. The same
    holds for $\mathcal M'_i$, $A'$ and $v'$. The lines $\Theta$ and $\Delta$ are adapted to $(A_o,\theta,v)$ at every
    nonzero prime and every embedding (Definition~\ref{cv-cusp-cm-key}), and at every nonzero prime $\mathfrak p$
    some model matrix of $(B,v)$ at $\mathfrak p$ has determinant a local generator of $\Delta$.
  \item[(f)] (Letters.) A bijection $e$ from the set of pairs $(\tau,i)$, with $\tau\colon E\to\mathbf C$ and
    $1\le i\le m_\tau$, onto $\{1,\dots,g\}$, a matrix $U\in M_g(K)$ and metrized lines
    $\mathcal L_1,\dots,\mathcal L_g$ over $K$ such that, for $k=e(\tau,i)$: if $\tau=\tau^\Phi(j)$, then
    $\sum_lU_{kl}v_l=d\varphi_i(v^\Phi_j)$ and $\mathcal L_k=\mathcal M^\Phi_j$; if $\tau=\tau'(j)$, then
    $\sum_lU_{kl}v_l=d\varphi'_i(v'_j)$ and $\mathcal L_k=\mathcal M'_j$.
  \item[(g)] (The framed point.) A framed point $y$ of $X$ over $K$ with point $o$, whose generator is $\theta$ and
    whose basis at every $\tau$ has $i$-th vector $\sum_lU_{e(\tau,i),l}v_l$.
  \end{enumerate}
\end{enumerate}
\end{definition}
\leanref{Development/AbelianVarietyTheory/Definitions/CMLattices/IntegralOf.lean}{GlobalTheory.LamT.CMPt}
\leanref{Development/AbelianVarietyTheory/Definitions/CMLattices/IntegralOf.lean}{GlobalTheory.LamT.CMPt.Good}
\leanref{Development/AbelianVarietyTheory/Definitions/CMLattices/LamLetters.lean}{GlobalTheory.LamF.LamObjGood}
\leanref{Development/AbelianVarietyTheory/Definitions/CMLattices/LamLetters.lean}{GlobalTheory.LamF.LamLetters.IsGood}
\leanref{Development/AbelianVarietyTheory/Definitions/CMLattices/LamLetters.lean}{GlobalTheory.LamF.IsModelAt}
\leanref{Development/AbelianVarietyTheory/Definitions/CMLattices/LamLetters.lean}{GlobalTheory.LamF.IsVecLine}
\leanref{Development/AbelianVarietyTheory/Definitions/CMLattices/LamLetters.lean}{GlobalTheory.LamF.IsIsogSum}
\leanref{Development/AbelianVarietyTheory/Definitions/CMLattices/LamLetters.lean}{GlobalTheory.LamF.PolClause}

\begin{lemma}\label{cv-cusp-cm-point-letters}
In Situation~\ref{cv-cusp-setting}, let $K\supset F_c$ be a number field, let $\ell>0$, and let a good CM point of
$X$ over $K$ with bound $\ell$ be given, with the notation of Definition~\ref{cv-cusp-cm-point-good}. Then
$\det U\ne0$.
\end{lemma}
\leanref{Development/AbelianVarietyTheory/Proofs/CubicalLineBundles/EqZeroOfPow.lean}{GlobalTheory.LamF.Ind.lamInd}
\begin{proof}
Let $K[\varepsilon]$ be the ring of dual numbers over $K$. The Lie algebra of an abelian variety over $K$ is the
set of its points with values in $K[\varepsilon]$ that reduce to the origin; on these points the group law is the
addition of the Lie algebra, and a homomorphism acts on them through its derivative.

Suppose that $\det U=0$. Then $\sum_kc_kU_{kl}=0$ for all $l$, for some nonzero $c\in K^g$. Put
$w_k=\sum_lU_{kl}v_l$, so that $\sum_kc_kw_k=0$. For $1\le j\le p$ and $1\le j'\le q$ put
\[
 a_j=\sum_{i=1}^fc_{e(\tau^\Phi(i),j)}v^\Phi_i,\qquad a'_{j'}=\sum_{i=1}^fc_{e(\tau'(i),j')}v'_i.
\]
Since $e$ is a bijection, $m_\tau=p$ for $\tau\in\Phi$ and $m_\tau=q$ otherwise, every $k$ is $e(\tau^\Phi(i),j)$ for
one pair $(i,j)$ or $e(\tau'(i),j')$ for one pair $(i,j')$, and Definition~\ref{cv-cusp-cm-point-good}(2)(f) gives
\[
 0=\sum_kc_kw_k=\sum_jd\varphi_j(a_j)+\sum_{j'}d\varphi'_{j'}(a'_{j'}).
\]
View the $a_j$ and $a'_{j'}$ as points over $K[\varepsilon]$, an $\mathcal O_K$-algebra. By (c) of the definition,
applied to these points and to the zero points, $Na_j=\psi_j(0)=N\cdot0=0$, and likewise $Na'_{j'}=0$, in the Lie
algebras. As $K$ has characteristic $0$ and $v^\Phi$ is a basis, $a_j=0$ gives $c_{e(\tau^\Phi(i),j)}=0$ for every $i$,
and in the same way $a'_{j'}=0$ gives $c_{e(\tau'(i),j')}=0$. Hence $c=0$, a contradiction.
\end{proof}
\leanref{Development/AbelianVarietyTheory/Proofs/CubicalLineBundles/EqZeroOfPow.lean}{GlobalTheory.LamF.Ind.eq_zero_of_pow}
% Lean: lamInd: Matrix.exists_vecMul_eq_zero_iff gives c; hone is sum_k c_k w_k = 0 (tanExp_tanExp,
% lieExp_eq_one); exists_mulEquiv_points_of_isGrpPullback passes between points of the generic fibres and points of
% B, A^Phi, A' over K[eps]; hs and hs0 are the sums for a, a' and for the zero points; p1, q1 (and p2, q2) are the
% identities of IsIsogSum for psi_j (psi'_j); eq_zero_of_pow concludes.

\begin{lemma}\label{cv-cusp-cm-point-model}
Let $K$ be a number field, let $A$ be an abelian scheme over $\mathcal O_K$ with a basis $v=(v_1,\dots,v_h)$ of the
Lie algebra of its generic fibre, let $\mathfrak p$ be a nonzero prime of $\mathcal O_K$, and let $M$ be a model
matrix of $(A,v)$ at $\mathfrak p$ (Definition~\ref{cv-cusp-cm-point-good}(1)). Write $\operatorname{Lie}A$ for the
Lie algebra of $A$, an $\mathcal O_K$-lattice in that of the generic fibre. For $x\in K^h$, the entries of $xM^{-1}$
are integral at $\mathfrak p$ if and only if $s\sum_lx_lv_l\in\operatorname{Lie}A$ for some
$s\in\mathcal O_K\setminus\mathfrak p$.
\end{lemma}
\leanref{Development/AbelianVarietyTheory/Proofs/CubicalLineBundles/IsModelAtIntegralIffMemLieLat.lean}{GlobalTheory.LamF.isModelAt_integral_iff_mem_lieLat}
\begin{proof}
The abelian scheme over the localization $\mathcal O_{K,\mathfrak p}$ that comes with $M$ is the base change of $A$,
so its Lie algebra is the localization of $\operatorname{Lie}A$: the set of vectors $w$ with
$sw\in\operatorname{Lie}A$ for some $s\in\mathcal O_K\setminus\mathfrak p$. The basis $v^{\mathfrak p}$ of
Definition~\ref{cv-cusp-cm-point-good}(1) is a basis of it over $\mathcal O_{K,\mathfrak p}$, and
$v^{\mathfrak p}_k=\sum_lM_{kl}v_l$. Hence $\sum_lx_lv_l=\sum_k(xM^{-1})_kv^{\mathfrak p}_k$ lies in it exactly when
the coordinates $(xM^{-1})_k$ lie in $\mathcal O_{K,\mathfrak p}$, that is, are integral at $\mathfrak p$.
\end{proof}
% Lean: isModelAt_integral_iff_mem_lieLat, through GlobalTheory.T.lieLattice_localization_modelLieAlgebra.

\begin{lemma}\label{cv-cusp-cm-point-local-ideal}
Let $K$ be a number field, let $A$ be an abelian scheme over $\mathcal O_K$ with a basis $v=(v_1,\dots,v_h)$ of the
Lie algebra of its generic fibre, let $1\le m\le h$, and let $I$ be a nonzero fractional ideal of $K$. Suppose that
for every nonzero prime $\mathfrak p$ of $\mathcal O_K$ there are a model matrix $M_{\mathfrak p}$ of $(A,v)$ at
$\mathfrak p$ (Definition~\ref{cv-cusp-cm-point-good}(1)) and $c_{\mathfrak p}\in K^\times$ with
$\operatorname{ord}_{\mathfrak p}c_{\mathfrak p}=\operatorname{ord}_{\mathfrak p}I$ such that, for $c\in K$, the
entries of $ce_mM_{\mathfrak p}^{-1}$, with $e_m$ the $m$-th unit row, are integral at $\mathfrak p$ exactly when
$\operatorname{ord}_{\mathfrak p}c\ge\operatorname{ord}_{\mathfrak p}c_{\mathfrak p}$. Then, with
$\operatorname{Lie}A$ as in Lemma~\ref{cv-cusp-cm-point-model}, $I$ is the set of the $c\in K$ with
$cv_m\in\operatorname{Lie}A$.
\end{lemma}
\leanref{Development/AbelianVarietyTheory/Proofs/CubicalLineBundles/MemLieVectorLineIdealIffMemLieLattice.lean}{GlobalTheory.LamF.isVecLine_mem_iff_mem_lieLat}
\begin{proof}
Let $c\in K$. As $I$ is a nonzero fractional ideal, $c\in I$ if and only if
$\operatorname{ord}_{\mathfrak p}c\ge\operatorname{ord}_{\mathfrak p}I=\operatorname{ord}_{\mathfrak p}c_{\mathfrak p}$
for every $\mathfrak p$. By the hypothesis, this holds if and only if the entries of $ce_mM_{\mathfrak p}^{-1}$ are
integral at $\mathfrak p$ for every $\mathfrak p$, and by Lemma~\ref{cv-cusp-cm-point-model} for $x=ce_m$, if and
only if for every $\mathfrak p$ some $s\in\mathcal O_K\setminus\mathfrak p$ has $scv_m\in\operatorname{Lie}A$. This
last condition holds if and only if $cv_m\in\operatorname{Lie}A$: if it holds, the ideal of the $a\in\mathcal O_K$
with $acv_m\in\operatorname{Lie}A$ lies in no maximal ideal of $\mathcal O_K$, so it contains $1$.
\end{proof}
\leanref{Development/AbelianVarietyTheory/Proofs/CubicalLineBundles/ValuationLeIffExistsSmulMem.lean}{GlobalTheory.LamF.local_global_ideal_lattice}
% Lean: isVecLine_mem_iff_mem_lieLat, from the second clause of IsVecLine (the hypothesis), the first and third parts
% of local_global_ideal_lattice (Idele.mem_iff_valuation_le and mem_iff_forall_exists_smul_mem) and
% isModelAt_integral_iff_mem_lieLat at every prime.

\begin{lemma}\label{cv-cusp-cm-point-local-row}
In Situation~\ref{cv-cusp-setting}, let $K\supset F_c$ be a number field, let $\ell>0$, and let a good CM point of
$X$ over $K$ with bound $\ell$ be given, with the notation of Definition~\ref{cv-cusp-cm-point-good}. Let
$\mathfrak p$ be a nonzero prime of $\mathcal O_K$, let $M$ be a model matrix of $(B,v)$ at $\mathfrak p$, let
$1\le j\le g$, and let $c_j$ be a local generator of $\mathcal L_j$ at $\mathfrak p$
(Definition~\ref{cv-cusp-cm-lines}). Then the entries of $c_jU_jM^{-1}$, where $U_j$ is the $j$-th row of $U$, are
integral at $\mathfrak p$.
\end{lemma}
\leanref{Development/AbelianVarietyTheory/Proofs/CubicalLineBundles/LieLetterIntegralOfLocalGeneratorStatement.lean}{GlobalTheory.T.lieLetter_integral_of_local_generator}
\leanref{Development/AbelianVarietyTheory/Proofs/CubicalLineBundles/LieLetterIntegralOfLocalGenerator.lean}{GlobalTheory.LamF.letter_integral_of_isLocGen}
\begin{proof}
Write $j=e(\tau,i)$, and suppose that $\tau=\tau^\Phi(m)$; if $\tau=\tau'(m)$, replace $A^\Phi$, $v^\Phi$, $\varphi_i$
and $\mathcal M^\Phi_m$ by $A'$, $v'$, $\varphi'_i$ and $\mathcal M'_m$ below. By (f) of
Definition~\ref{cv-cusp-cm-point-good}, $\sum_lU_{jl}v_l=d\varphi_i(v^\Phi_m)$ and $\mathcal L_j=\mathcal M^\Phi_m$.
By (e) and
Lemma~\ref{cv-cusp-cm-point-local-ideal}, the ideal $I_{\mathcal L_j}$ of $\mathcal M^\Phi_m$ is the set of the
$c'\in K$ with $c'v^\Phi_m\in\operatorname{Lie}A^\Phi$. The fractional ideal $c_j^{-1}I_{\mathcal L_j}$ has order $0$
at $\mathfrak p$, so its intersection with $\mathcal O_K$ is an ideal of $\mathcal O_K$ that is not contained in
$\mathfrak p$. For $s$ in it with $s\notin\mathfrak p$, $sc_j$ lies in $I_{\mathcal L_j}$, so
$sc_jv^\Phi_m\in\operatorname{Lie}A^\Phi$. The homomorphism $\varphi_i$ of abelian schemes over $\mathcal O_K$ maps
$\operatorname{Lie}A^\Phi$ into $\operatorname{Lie}B$, so $sc_j\sum_lU_{jl}v_l=d\varphi_i(sc_jv^\Phi_m)$ lies in
$\operatorname{Lie}B$, and Lemma~\ref{cv-cusp-cm-point-model} for $M$ shows that the entries of $c_jU_jM^{-1}$ are
integral at $\mathfrak p$.
\end{proof}
% Lean: letter_integral_of_isLocGen = image_integral_of_isLocGen: s from the second part of
% local_global_ideal_lattice, isVecLine_mem_iff_mem_lieLat (Lemma cv-cusp-cm-point-local-ideal), smul_image_mem_lieLat
% through GlobalTheory.T.existsUnique_differential_monHom_lieBasis, and isModelAt_integral_iff_mem_lieLat.

\begin{definition}\label{cv-cusp-cm-point-local-model}
In Situation~\ref{cv-cusp-setting}, use the framed points of Lemma~\ref{cv-cusp-cm-point-polynomial}. Let $V$ be a
complete discrete valuation ring of characteristic $0$ whose fraction field $K'$ is an $F_c$-algebra, and let $y$ be a
framed point of $X$ over $K'$, with bases $u_\tau$ and generator $\vartheta$; its abelian variety $A_y$, with its
action of $\mathcal O_E$ and its line bundle, is an $\mathcal O_E$-semi-abelian scheme over $K'$
(Definition~\ref{cv-model-toroidal-semiabelian}(2)).
\begin{enumerate}
\item A \emph{model} of $y$ over $V$ is an $\mathcal O_E$-semi-abelian scheme $G$ over $V$ together with a morphism
  $A_y\to G$ of $\mathcal O_E$-semi-abelian schemes over $\operatorname{Spec}K'\to\operatorname{Spec}V$, which identifies
  $A_y$ with the generic fibre of $G$. For an embedding $\tau\colon E\to\mathbf C$, the \emph{$\tau$-eigenspace}
  $\operatorname{Lie}_\tau G$ of the Lie algebra $\operatorname{Lie}G$ of $G$ over $V$ is the $V$-module of the
  $w\in\operatorname{Lie}G$ with $bw=\tau(b)w$ for every $b\in\mathcal O_E$; here $\tau(b)\in F_c$ is viewed in
  $K'$, and it lies in $V$, being integral over $\mathbf Z$.
\item The bases of $y$ are \emph{integral for} a model $G$ if for every $\tau$ there is a basis
  $b_{\tau,1},\dots,b_{\tau,m_\tau}$ of $\operatorname{Lie}_\tau G$ over $V$ whose images in the Lie algebra of the
  generic fibre are those of $u_{\tau,1},\dots,u_{\tau,m_\tau}$ under the morphism. The generator of $y$ is
  \emph{integral for} $G$ if some section of the line bundle of $G$ generates the $V$-module of its sections and has
  the pullback $\vartheta$ to $A_y$. The framed point $y$ is \emph{integral for} $G$ if its bases and its generator
  are, and \emph{integral} if it is integral for some model over $V$.
\end{enumerate}
\end{definition}
\leanref{Formalism/LocalToGlobal/ProjectiveLattices/VAtVNeBot.lean}{Formalism.L2.Pj.Integral}
\leanref{Development/AbelianVarietyTheory/Definitions/CMLattices/IntegralOf.lean}{GlobalTheory.LamT.LieInt}
\leanref{Development/AbelianVarietyTheory/Definitions/CMLattices/IntegralOf.lean}{GlobalTheory.LamT.SecInt}

\begin{lemma}\label{cv-cusp-cm-point-local-eigenbasis}
In Situation~\ref{cv-cusp-setting}, use the framed points of Lemma~\ref{cv-cusp-cm-point-polynomial}. Let $V$ be a
complete discrete valuation ring of characteristic $0$ whose fraction field $K'$ is an $F_c$-algebra, let $y$ be a
framed point of $X$ over $K'$, and let $G$ be a model of $y$ over $V$
(Definition~\ref{cv-cusp-cm-point-local-model}).
\begin{enumerate}
\item There is a family $h=(h_\tau)$ of invertible matrices over $K'$ such that the bases of $y\cdot(h,c)$ are
  integral for $G$ for every $c\in K'^\times$.
\item If the bases of $y$ are integral for $G$, then so are those of $y\cdot(Q,c)$ for every family $Q=(Q_\tau)$ of
  matrices over $V$ with unit determinants and every $c\in K'^\times$.
\end{enumerate}
\end{lemma}
\leanref{Development/AbelianVarietyTheory/Proofs/CubicalLineBundles/ExistsIntegralEigenbasisOfFrameAtLocalModelStatement.lean}{GlobalTheory.T.exists_integralEigenbasis_of_frame_at_localModel}
\leanref{Development/AbelianVarietyTheory/Proofs/CubicalLineBundles/ExistsIntegralEigenbasisOfFrameAtLocalModel.lean}{GlobalTheory.LamT.lieInt_exists_and_change}
\begin{proof}
(1) By Lemma~\ref{cv-origin-reference-lie-bundle-base}, $\operatorname{Lie}G$ is a finite projective, hence free,
$V$-module, and $K'\otimes_V\operatorname{Lie}G$ is the Lie algebra of the generic fibre, which the model identifies
with $\operatorname{Lie}A_y$. So $\operatorname{Lie}G$ maps injectively to $\operatorname{Lie}A_y$, compatibly with the
actions of $\mathcal O_E$ (Lemma~\ref{cv-origin-reference-lie}(2) and (3)). Fix $\tau$ and put $m=m_\tau$. The image of
$\operatorname{Lie}_\tau G$ lies in the $\tau$-eigenspace of $\operatorname{Lie}A_y$, of which
$u_{\tau,1},\dots,u_{\tau,m}$ is a basis; for $w\in\operatorname{Lie}_\tau G$ let $[w]\in K'^m$ be the coordinates
of its image in this basis. Then $w\mapsto[w]$ is $V$-linear and injective, and its image $\Lambda$ is a finitely
generated $V$-module, as $V$ is noetherian. For each $i$ some nonzero $d\in V$ has $du_{\tau,i}$ in the image of
$\operatorname{Lie}G$, by clearing the denominators of its coordinates in a basis of $\operatorname{Lie}G$; the
preimage lies in $\operatorname{Lie}_\tau G$, as $\operatorname{Lie}G$ maps injectively. So $\Lambda$ contains a
vector whose coordinates after the $i$-th vanish and whose $i$-th coordinate is not $0$.

For each $i$, the $i$-th coordinates of the vectors of $\Lambda$ whose coordinates after the $i$-th vanish form a
nonzero finitely generated $V$-submodule of $K'$, which is $Vt_i$ for some $t_i\ne0$, as $V$ is a principal ideal
domain; choose $\beta_i\in\Lambda$ with these coordinates $0$ and with $i$-th coordinate $t_i$. The $\beta_i$ form a
basis of $\Lambda$ over $V$. They are linearly independent, as the matrix of their coordinates is triangular with
nonzero diagonal. They span $\Lambda$: if $\alpha\in\Lambda$ has its coordinates after the $i$-th equal to $0$, then
$\alpha_i=at_i$ with $a\in V$, and $\alpha-a\beta_i$ has its coordinates from the $i$-th on equal to $0$; induction on
$i$ finishes. Let $C_\tau$ be the matrix with the rows $\beta_1,\dots,\beta_m$, lower triangular with nonzero
diagonal, and put $h_\tau=C_\tau^{\mathsf T}$. The $i$-th basis vector of $y\cdot(h,c)$ at $\tau$ is
$\sum_j(h_\tau)_{ji}u_{\tau,j}=\sum_j(C_\tau)_{ij}u_{\tau,j}$, the image of the $w_i\in\operatorname{Lie}_\tau G$ with
$[w_i]=\beta_i$; and $w_1,\dots,w_m$ form a basis of $\operatorname{Lie}_\tau G$ over $V$.

(2) Let $b_{\tau,j}$ be bases of the $\operatorname{Lie}_\tau G$ over $V$ with the images $u_{\tau,j}$. As $\det Q_\tau$
is a unit of $V$, the vectors $\sum_j(Q_\tau)_{ji}b_{\tau,j}$ form a basis of $\operatorname{Lie}_\tau G$ over $V$, and
their images $\sum_j(Q_\tau)_{ji}u_{\tau,j}$ are the bases of $y\cdot(Q,c)$.
\end{proof}
\leanref{Formalism/LocalToGlobal/ProjectiveLattices/LieFreePf.lean}{Formalism.L2.Pj.lieFree}
\leanref{Formalism/LocalToGlobal/ProjectiveLattices/LieAlg.lean}{Formalism.L2.Pj.LieAlg.eig_basis}
\leanref{Formalism/LocalToGlobal/ProjectiveLattices/LieAlg.lean}{Formalism.L2.Pj.LieAlg.herm}
\leanref{Development/AbelianVarietyTheory/Proofs/CubicalLineBundles/ExistsIntegralEigenbasisOfFrameAtLocalModel.lean}{GlobalTheory.LamT.LocLie.lieInt_change}
% Lean: (1) is lieFree: a chart of G at the origin (univChartLoc, tanBC_alg) writes the tangent spaces over V and K'
% as V^n and K'^n, the actions of O_E as the matrices tmap, and tau(b) as chi tau b in V (int_val); eig_basis is the
% submodule E, its coordinate map f, the vectors d u_i (IsLocalization.exist_integer_multiples) and herm (the
% triangular basis); lieInt_exists_and_change takes g = C (lower triangular, FrChg g a y y' with the Lean's row
% convention, so h = C^T in the text's). (2) is lieInt_change with k = X^T, through comb_eq and reframe_basis.

\begin{lemma}\label{cv-cusp-cm-point-local-generator}
In Situation~\ref{cv-cusp-setting}, use the framed points of Lemma~\ref{cv-cusp-cm-point-polynomial}. Let $V$ be a
complete discrete valuation ring of characteristic $0$ whose fraction field $K'$ is an $F_c$-algebra, let $y$ be a
framed point of $X$ over $K'$, and let $G$ be a model of $y$ over $V$
(Definition~\ref{cv-cusp-cm-point-local-model}) that is an abelian scheme whose line bundle defines a principal
polarization. Then there is $s\in K'^\times$ such that the generator of $y\cdot(g,s)$ is integral for $G$ for every
family $g$ of invertible matrices over $K'$.
\end{lemma}
\leanref{Development/AbelianVarietyTheory/Proofs/CubicalLineBundles/ExistsUnitIntegralGeneratorPrincipalLocalModelStatement.lean}{GlobalTheory.T.exists_unit_integral_generator_principal_local_model}
\leanref{Development/AbelianVarietyTheory/Proofs/CubicalLineBundles/ExistsUnitIntegralGeneratorPrincipalLocalModel.lean}{GlobalTheory.LamT.LocSec.exists_unit_secInt}
\begin{proof}
By Theorem~\ref{cv-volume-key-sections} for $n=1$ over $\operatorname{Spec}V$, whose only open set containing the
closed point is $\operatorname{Spec}V$, the sections of the line bundle of $G$ form a free $V$-module of rank one,
with a generator $\theta_G$ whose pullback along every base change generates the sections there. So the pullback
$\sigma$ of $\theta_G$ to $A_y$ generates the sections of the line bundle of $A_y$, and $\sigma=s\vartheta$ for the
generator $\vartheta$ of $y$ and some $s\in K'$, with $s\ne0$ as $\sigma\ne0$. For every $g$, the generator of
$y\cdot(g,s)$ is $s\vartheta=\sigma$.
\end{proof}
% Lean: exists_unit_secInt: KeyOnPol.exists_generator for theta_G (from the law SecBaseChangeSt,
% Theorem cv-volume-key-sections), SectionBaseChange.generates_and_baseChange for its pullback sigma, exists_secTie
% for the section of the point, y.gen for s, and s /= 0 by TieCalc.secTie_scale.

\begin{lemma}\label{cv-cusp-cm-point-local-base}
In Situation~\ref{cv-cusp-setting}, use the framed points of Lemma~\ref{cv-cusp-cm-point-polynomial}. Let $V$ and
$W$ be complete discrete valuation rings of characteristic $0$ with fraction fields $K'$ and $K_W$, and let
$\psi\colon V\to W$ and $i\colon K'\to K_W$ be ring homomorphisms that agree on $V$; let $K'$ be an $F_c$-algebra,
and $K_W$ one through $i$. Let $y$ be a framed point of $X$ over $K'$ and $G$ a model of $y$ over $V$
(Definition~\ref{cv-cusp-cm-point-local-model}) that is an abelian scheme whose line bundle defines a principal
polarization, and let $y_W$ be the image of $y$ under $i$ (Lemma~\ref{cv-volume-key-base-change}). Then the base
change $G_W$ of $G$ along $\psi$ is a model of $y_W$; and if $y$ is integral for $G$, then $y_W\cdot Q$ is integral
for $G_W$ for every family $Q$ of matrices over $W$ with unit determinants.
\end{lemma}
\leanref{Development/AbelianVarietyTheory/Proofs/CubicalLineBundles/LocalLowerBoundCmPointFramePolynomialEveryPrime.lean}{GlobalTheory.LamT.LocClause.integral_baseChange_chg}
\leanref{Development/AbelianVarietyTheory/Proofs/CubicalLineBundles/IntegralEigenbasisBaseChangeDvrStatement.lean}{GlobalTheory.T.integral_eigenbasis_baseChange_dvr}
\leanref{Development/AbelianVarietyTheory/Proofs/CubicalLineBundles/IntegralGeneratorBaseChangeValuationRingStatement.lean}{GlobalTheory.T.integral_generator_baseChange_valuationRing}
\begin{proof}
The abelian variety of $y_W$ is the base change of $A_y$ along $i$, and the base change of the morphism $A_y\to G$
is a morphism from it to $G_W$ over $\operatorname{Spec}K_W\to\operatorname{Spec}W$; so $G_W$ is a model of $y_W$. Let
$y$ be integral for $G$, with bases $b_{\tau,j}$ of the $\operatorname{Lie}_\tau G$ and a generator $\theta_G$ of the
sections of the line bundle of $G$ as in Definition~\ref{cv-cusp-cm-point-local-model}(2).

Bases. By Lemma~\ref{cv-origin-reference-lie-bundle-base}, $\operatorname{Lie}G_W=W\otimes_V\operatorname{Lie}G$,
compatibly with the actions of $\mathcal O_E$. The homomorphism $\psi$ is injective, as $i$ is, so $W$ is torsion
free over the discrete valuation ring $V$, hence flat. So $\operatorname{Lie}_\tau G_W$, the intersection of the
kernels of the maps $b-\tau(b)$, is $W\otimes_V\operatorname{Lie}_\tau G$, and the images of the $b_{\tau,j}$ form a
basis of it over $W$. Their images in the generic fibre are those of the images under $i$ of the $u_{\tau,j}$, the
bases of $y_W$. So the bases of $y_W$ are integral for $G_W$, and by
Lemma~\ref{cv-cusp-cm-point-local-eigenbasis}(2) over $W$ so are those of $y_W\cdot Q$.

Generator. By Theorem~\ref{cv-volume-key-sections}, the pullback of $\theta_G$ to $G_W$ generates the sections of
its line bundle. Its pullback to the abelian variety of $y_W$ is the pullback of the generator $\vartheta$ of $y$,
which is the generator of $y_W$ and of $y_W\cdot Q$; here the identifications of line bundles are those that respect
the rigidifications, and two such identifications agree, as the global functions on an abelian variety are constant.
So the generator of $y_W\cdot Q$ is integral for $G_W$.
\end{proof}
% Lean: integral_baseChange_chg: the base-changed model M' of the second part of locMod_exists (pat_pullback,
% sab_comp_exists, sab_cancel_exists); integral_eigenbasis_baseChange_dvr = lieInt_baseChange (exists_charts,
% flat_of_inj, eig_flat); the second part of lieInt_exists_and_change for Q with a = 1; secBC_one_of_frMap and
% SecCalc.sec_calc for the section of w; integral_generator_baseChange_valuationRing = secInt_baseChange
% (SectionBaseChange.generates, LiftBaseChange.rigIso_comp and iso_unique_of_rig with GlobalFnsSt); integral_of.

\begin{situation}\label{cv-cusp-cm-point-local-situation}
In Situation~\ref{cv-cusp-setting}, let $K\supset F_c$ be a number field, let $\ell>0$, and let a good CM point of
$X$ over $K$ with bound $\ell$ be given, with the notation of Definition~\ref{cv-cusp-cm-point-good}. Let
$\mathfrak p$ be a nonzero prime of $\mathcal O_K$, let $V$ be a complete discrete valuation ring of characteristic
$0$ with fraction field $K'$, and let $\iota\colon K\to K'$ be an embedding with
$\iota^{-1}(V)=\{y\in K:\operatorname{ord}_{\mathfrak p}y\ge0\}$; make $K'$ an $F_c$-algebra through $\iota$. As
$\iota$ maps $\mathcal O_K$ into $V$, it restricts to $\rho\colon\mathcal O_K\to V$; let $B_V$ be the base change of
$B$ along $\rho$, and $y_{K'}$ the image under $\iota$ of the framed point $y$ of (g) of that definition
(Lemma~\ref{cv-volume-key-base-change}). The point of $y$ is $o$, so the abelian variety of $y_{K'}$ is the base
change of $A_o$ along $\iota$, and the base change along $\iota$ of the identification in (b) of that definition of
$A_o$ with the generic fibre of $B$ makes $B_V$ a model of $y_{K'}$ (Definition~\ref{cv-cusp-cm-point-local-model}), an abelian
scheme whose line bundle defines a principal polarization.
\end{situation}
\leanref{Development/AbelianVarietyTheory/Definitions/CMLattices/IntegralOf.lean}{GlobalTheory.LamT.LocMod}
\leanref{Development/AbelianVarietyTheory/Proofs/CubicalLineBundles/ExistsCmPointLocalModelOverDvrStatement.lean}{GlobalTheory.T.exists_cmPoint_localModel_over_dvr}
% Lean: the first part of locMod_exists: rho from iK (h a for the integrality of O_K), G = pat_pullback of z.B along
% rho, t' = pat_pullback of z.Phi.T along iK, m1 from z.Phi.m and m2 from z.m' (sab_comp_exists, sab_cancel_exists).

\begin{lemma}\label{cv-cusp-cm-point-local-unit}
In Situation~\ref{cv-cusp-cm-point-local-situation}, let $c_\Theta$ be a local generator of $\Theta$ at $\mathfrak p$,
let $g$ be a family of invertible matrices over $K'$, and let $a\in K'^\times$ be such that the generator of
$y_{K'}\cdot(g,a)$ is integral for $B_V$ (Definition~\ref{cv-cusp-cm-point-local-model}). Then
$a=\iota(c_\Theta)u$ for a unit $u$ of $V$.
\end{lemma}
\leanref{Development/AbelianVarietyTheory/Proofs/CubicalLineBundles/IntegralGeneratorAssociatedLocalGeneratorCMLineStatement.lean}{GlobalTheory.T.integralGenerator_associated_localGenerator_CM_line}
\leanref{Development/AbelianVarietyTheory/Proofs/CubicalLineBundles/IntegralGeneratorAssociatedLocalGeneratorCMLine.lean}{GlobalTheory.LamT.LocGen.integral_generator_unit}
\begin{proof}
Let $\theta_V$ be a generator of the sections of the line bundle of $B_V$ whose pullback is the generator of
$y_{K'}\cdot(g,a)$. By (g) of Definition~\ref{cv-cusp-cm-point-good}, the generator of $y$ is $\theta$, so that of
$y_{K'}\cdot(g,a)$ is $a\theta_{K'}$, with $\theta_{K'}$ the image of $\theta$.

The basis. By (e) there is a model matrix $M$ of $(B,v)$ at $\mathfrak p$ whose determinant is a local generator of
$\Delta$: an abelian scheme $B_{\mathfrak p}$ over the localization $R=\mathcal O_{K,\mathfrak p}$, the base change of
$B$, with a basis $v^{\mathfrak p}$ of its Lie algebra such that $v^{\mathfrak p}_k=\sum_lM_{kl}v_l$ over $K$
(Definition~\ref{cv-cusp-cm-point-good}(1)). As $\iota(R)\subset V$, $B_V$ is the base change of $B_{\mathfrak p}$
along $\iota|_R$, so by
Lemma~\ref{cv-origin-reference-lie-bundle-base} the images $v^V_k$ of the $v^{\mathfrak p}_k$ form a basis of
$\operatorname{Lie}B_V$ over $V$, and their images in the generic fibre are $\sum_l\iota(M_{kl})v_l$, where $v_l$ also
denotes the image of $v_l$ under $\iota$.

The key formula. So the key datum $(A_{K'},a\theta_{K'},\iota(M)v)$ over $K'$, where $(A_{K'},\theta_{K'},v)$ is the
base change of $(A_o,\theta,v)$ along $\iota$, extends over $V$, to $B_V$ with $\theta_V$ and $v^V$. Let $\kappa$ be
the rule of Theorem~\ref{cv-definite-key-formula} and $k=\kappa(A_o,\theta,v)\in K^\times$. By parts (1) and (2) of
the theorem, $a^8\kappa(A_{K'},a\theta_{K'},\iota(M)v)=\iota(\det M)^4\iota(k)$, and by part (3)
$\kappa(A_{K'},a\theta_{K'},\iota(M)v)$ is a unit of $V$.

The valuation. By (e), $\Theta$ and $\Delta$ are adapted to $(A_o,\theta,v)$ at $\mathfrak p$
(Definition~\ref{cv-cusp-cm-key}(2)). As in Step 1 of the proof of Lemma~\ref{cv-cusp-cm-theta-degree}, the theorem
then gives $\operatorname{ord}_{\mathfrak p}(kc'^4_\Delta c'^{-8}_\Theta)=0$ for the local generators $c'_\Theta$ of
$\Theta$ and $c'_\Delta$ of $\Delta$ of that definition. As $c_\Theta$ and $c'_\Theta$ are local generators of
$\Theta$, and $\det M$ and $c'_\Delta$ of $\Delta$, $\operatorname{ord}_{\mathfrak p}(k(\det M)^4c_\Theta^{-8})=0$, and
$\iota(k(\det M)^4c_\Theta^{-8})$ is a unit of $V$, as $\iota^{-1}(V)$ is the valuation ring of $\mathfrak p$. Hence
\[
 \Big(\frac a{\iota(c_\Theta)}\Big)^8=\iota\big(k(\det M)^4c_\Theta^{-8}\big)\,\kappa(A_{K'},a\theta_{K'},\iota(M)v)^{-1}
\]
is a unit of $V$. As $V$ is integrally closed, $u=a/\iota(c_\Theta)$ lies in $V$, and it is a unit, as $u^8$ is.
\end{proof}
\leanref{Development/AbelianVarietyTheory/Proofs/CubicalLineBundles/ValuationEqOfIsLocGen.lean}{GlobalTheory.KeyVal.integral_section_unit}
\leanref{Development/AbelianVarietyTheory/Proofs/CubicalLineBundles/ExistsLieBasisLocalModelStatement.lean}{GlobalTheory.T.exists_lieBasis_local_model}
% Lean: integral_generator_unit: hmod p (the model matrix A with IsLocGen p A.det), exists_isLieBasis_locMod for v^V
% (exists_restrict_ringHom, IsGrpPullback.cancel, exists_isLieBasis_baseChange), eq_mul_of_secTie for
% sigma = a theta' (hsigma_a), the key data x' and xR with KeyPB and KeyIntegral, and KeyVal.integral_section_unit:
% exists_pow_rel (the key formula in the power form keyFormulaPow_and_keyDegree, exponents 8 N and 4 N; the text uses
% Theorem cv-definite-key-formula, N = 1), unit_iff for h1 and h2 (the adaptedness KeyLocAt, and
% valuation_eq_of_isLocGen), and exists_unit_of_pow.

\begin{lemma}\label{cv-cusp-cm-point-local-letters}
In Situation~\ref{cv-cusp-cm-point-local-situation}, let $h$ be a family of invertible matrices over $K'$ and
$a\in K'^\times$ such that the bases of $y_{K'}\cdot(h,a)$ are integral for $B_V$
(Definition~\ref{cv-cusp-cm-point-local-model}). Let $\tau\colon E\to\mathbf C$ be an embedding, let
$1\le j\le m_\tau$, and let $c$ be a local generator of $\mathcal L_{e(\tau,j)}$ at $\mathfrak p$. Then
$\iota(c)(h_\tau^{-1})_{kj}\in V$ for $1\le k\le m_\tau$.
\end{lemma}
\leanref{Development/AbelianVarietyTheory/Proofs/CubicalLineBundles/LieLetterIntegralInIntegralEigenbasisStatement.lean}{GlobalTheory.T.lieLetter_integral_in_integral_eigenbasis}
\leanref{Development/AbelianVarietyTheory/Proofs/CubicalLineBundles/LieLetterIntegralInIntegralEigenbasis.lean}{GlobalTheory.LamT.LetterInt.letter_integral_in_eigenbasis}
\begin{proof}
Put $q=e(\tau,j)$, let $u_{\tau,1},\dots,u_{\tau,m_\tau}$ be the basis of $y_{K'}$ at $\tau$, and write $v_l$ also for
the images of the $v_l$ under $\iota$; by (g) of Definition~\ref{cv-cusp-cm-point-good},
$u_{\tau,j}=\sum_l\iota(U_{ql})v_l$. By (e) there is a model matrix $M$ of $(B,v)$ at
$\mathfrak p$, and by Lemma~\ref{cv-cusp-cm-point-local-row} the entries $r_k$ of $cU_qM^{-1}$ are integral at
$\mathfrak p$, so that the $\iota(r_k)$ lie in $V$. As in the proof of Lemma~\ref{cv-cusp-cm-point-local-unit}, the
images $v^V_k$ in $\operatorname{Lie}B_V$ of the basis of the model matrix $M$ form a basis over $V$, with images
$\sum_l\iota(M_{kl})v_l$ in the generic fibre. Put $w=\sum_k\iota(r_k)v^V_k\in\operatorname{Lie}B_V$. Its image in the
generic fibre is
\[
 \sum_k\iota(r_k)\sum_l\iota(M_{kl})v_l=\iota(c)\sum_l\iota(U_{ql})v_l=\iota(c)u_{\tau,j},
\]
an eigenvector of $\mathcal O_E$ for $\tau$; as $\operatorname{Lie}B_V$ maps injectively to the Lie algebra of the
generic fibre, $w\in\operatorname{Lie}_\tau B_V$. The bases of $y_{K'}\cdot(h,a)$ are integral, so
$w=\sum_kx_kb_k$ with $x_k\in V$, for the basis $b_1,\dots,b_{m_\tau}$ of $\operatorname{Lie}_\tau B_V$ whose images
are the vectors $u'_k=\sum_i(h_\tau)_{ik}u_{\tau,i}$ of $y_{K'}\cdot(h,a)$. Then
$\iota(c)u_{\tau,j}=\sum_kx_ku'_k=\sum_i\big(\sum_k(h_\tau)_{ik}x_k\big)u_{\tau,i}$, so $h_\tau x=\iota(c)e_j$ for the
column $x=(x_k)_k$ and the $j$-th unit vector $e_j$, and $x_k=\iota(c)(h_\tau^{-1})_{kj}$ lies in $V$.
\end{proof}
% Lean: letter_integral_in_eigenbasis: the model matrix A of hobj, letter_integral_of_isLocGen (Lemma
% cv-cusp-cm-point-local-row), exists_lieBasis_local_model for vV, the letter w of hlet and hu, LieComb.lieComb_algebra
% and comb_match, comb_eq_scale for iota(c) u, isEigV_of_generic for the eigenvector, the integral eigenbasis of
% hint, and coord_eq for the coordinates; the Lean's FrChg g with g = h^T, so its (g b)^-1 j k is (h^-1)_kj.

\begin{lemma}\label{cv-cusp-cm-point-local}
In Situation~\ref{cv-cusp-setting}, let $\kappa=(\alpha,n)$ be a weight as in
Lemma~\ref{cv-cusp-cm-point-polynomial}, and let $f$ be a rational section of $X$ of weight $\kappa$ that is
integral: its valuation is nonnegative at every integral point of $X$, algebraic over $F_c$, with values in a
complete discrete valuation ring above a finite place of $F_c$. Let $K\supset F_c$ be a number field, let $\ell>0$,
let a good CM point of $X$ over $K$ with bound $\ell$ be given, with the notation of
Definition~\ref{cv-cusp-cm-point-good}, and let $G$ be a polynomial as in
Lemma~\ref{cv-cusp-cm-point-polynomial}(1) for $f$ and its framed point $y$. Let $\mathfrak p$ be a nonzero prime of
$\mathcal O_K$. Then condition (L) of Proposition~\ref{cv-cusp-cm-product} holds at $\mathfrak p$ for the metrized
lines $\Theta$, $\Delta$ and $\mathcal L_1,\dots,\mathcal L_g$, the matrix $U$, the bijection $e$, the polynomial
$G$, the set $B$ of embeddings $\tau\colon E\to\mathbf C$ with the integers $m_\tau$, and the integer $n$.
\end{lemma}
\leanref{Development/AbelianVarietyTheory/Proofs/CubicalLineBundles/LocalLowerBoundCmPointFramePolynomialEveryPrimeStatement.lean}{GlobalTheory.T.localLowerBound_cmPoint_framePolynomial_every_prime}
\leanref{Development/AbelianVarietyTheory/Proofs/CubicalLineBundles/LocalLowerBoundCmPointFramePolynomialEveryPrime.lean}{GlobalTheory.LamT.LocClause.step5Loc_cmPt}
\begin{proof}
Step 1: the local generators and the determinant. Choose local generators $c_\Theta$ of $\Theta$, $c_\Delta$ of
$\Delta$ and $c_j$ of $\mathcal L_j$ at $\mathfrak p$ (Definition~\ref{cv-cusp-cm-lines}); then $d=c_\Theta^{-1}$ is a
local generator of $\Theta^\vee$. By (e) of the definition there is a model matrix $M$ of $(B,v)$ at $\mathfrak p$
whose determinant is a local generator of $\Delta$, so
$\operatorname{ord}_{\mathfrak p}\det M=\operatorname{ord}_{\mathfrak p}c_\Delta$. By
Lemma~\ref{cv-cusp-cm-point-local-row}, the matrix with rows $c_jU_jM^{-1}$ has entries integral at $\mathfrak p$,
so its determinant $c_1\cdots c_g\det U/\det M$ has nonnegative order at $\mathfrak p$. Hence
$\operatorname{ord}_{\mathfrak p}(c_1\cdots c_g\det U)\ge\operatorname{ord}_{\mathfrak p}\det M
=\operatorname{ord}_{\mathfrak p}c_\Delta$.

Step 2: the integral framed point. Let $K'$ be the completion of $K$ at $\mathfrak p$, $V$ its valuation ring and
$\iota\colon K\to K'$ the inclusion. Then $V$ is a complete discrete valuation ring of characteristic $0$ with the
finite residue field $\mathcal O_K/\mathfrak p$, and $\iota^{-1}(V)=\{y\in K:\operatorname{ord}_{\mathfrak p}y\ge0\}$
\cite[Chapter~II]{Ser79}. Make $K'$ an $F_c$-algebra through $\iota$, and let $B_V$ and $y_{K'}$ be as in
Situation~\ref{cv-cusp-cm-point-local-situation}, so that $B_V$ is a model of $y_{K'}$. By
Lemma~\ref{cv-cusp-cm-point-local-eigenbasis}(1) there is a family $h=(h_\tau)$ of invertible matrices over $K'$
such that the bases of $y_{K'}\cdot(h,c)$ are integral for $B_V$ for every $c\in K'^\times$, and by
Lemma~\ref{cv-cusp-cm-point-local-generator} there is $s\in K'^\times$ such that the generator of $y_{K'}\cdot(h,s)$
is integral for $B_V$. So $y'=y_{K'}\cdot(h,s)$ is integral for $B_V$. By
Lemma~\ref{cv-cusp-cm-point-local-unit}, $s=\iota(c_\Theta)u$ for a unit $u$ of $V$. For every $\tau$ put
$U'_\tau=(h_\tau^{-1})^{\mathsf T}\in GL_{m_\tau}(K')$. By Lemma~\ref{cv-cusp-cm-point-local-letters},
$\iota(c_{e(\tau,j)})(U'_\tau)_{jk}=\iota(c_{e(\tau,j)})(h_\tau^{-1})_{kj}$ lies in $V$ for all $\tau$, $j$ and $k$.

Step 3: the values of $G$. Let $W$ be an extension of $V$ (Lemma~\ref{cv-cusp-cm-integral}), with fraction field
$K_W$, and let $Y=(Y_\tau)$ be a family of matrices over $K_W$ such that every $Q_\tau=U'^{\mathsf T}_\tau
Y_\tau=h_\tau^{-1}Y_\tau$ has entries in $W$ and determinant in $W^\times$; here $h_\tau$, $s$ and $u$ also denote
their images under $K'\to K_W$, and $\iota$ the composite $K\to K_W$. Make $K_W$ an $F_c$-algebra through $K'$. Then
$Y_\tau=h_\tau Q_\tau$ is invertible. Let $y'_W$ be the image of $y'$ under $K'\to K_W$
(Lemma~\ref{cv-volume-key-base-change}), and put $w=y'_W\cdot Q$. By Lemma~\ref{cv-cusp-cm-point-local-base} for
$V\to W$, $y'$ and $B_V$, the framed point $w$ is integral (Definition~\ref{cv-cusp-cm-point-local-model}). As its
residue field is finite, $W$ lies above a finite place of $F_c$, and as $f$ is integral, $f(w)\in W$.

Let $y_W$ be the image of $y$ under $K\to K_W$. Base change of framed points is transitive and commutes with changes
of frame, and changes of frame compose; so $y'_W=y_W\cdot(h,s)$ and
\[
 w=y_W\cdot(hQ,s)=(y_W\cdot Y)\cdot(1,s).
\]
By the transformation law of $f$ over the field $K_W$ and Lemma~\ref{cv-cusp-cm-point-polynomial}(1) for $S=K_W$,
$f(w)=s^nf(y_W\cdot Y)=s^n\iota(G)(Y)$. As $s=\iota(c_\Theta)u=\iota(d)^{-1}u$,
\[
 \iota(G)(Y)=s^{-n}f(w)=\iota(d)^n\,u^{-n}f(w)\in\iota(d)^nW,
\]
since $u$ is a unit of $V$ and so of $W$.

Steps 1 to 3 give condition (L) at $\mathfrak p$, with the local generators $d$, $c_\Delta$ and $c_j$, the ring $V$,
the embedding $\iota$ and the matrices $U'_\tau$.
% Lean: Step 1 is GlobalTheory.LamF.letters_det_valuation_le (valuation_prod_mul_det_le, with
% T.lieLetter_integral_of_local_generator, Lemma cv-cusp-cm-point-local-row); d is ca^{-1} (isLocGen_inv). Steps 2
% and 3 are the rest of step5Loc_cmPt. Step 2: T, iK from GlobalTheory.T.exists_completeDVR_completion_integral_iff
% (= exists_cdvr_completion_integral_iff, from locFld); the local model M of
% GlobalTheory.T.exists_cmPoint_localModel_over_dvr (Situation cv-cusp-cm-point-local-situation); yT by
% FramedBaseChange.exists_frMap; g from GlobalTheory.T.exists_integralEigenbasis_of_frame_at_localModel (h = g^T);
% s from GlobalTheory.T.exists_unit_integral_generator_principal_local_model; y' = existsUnique_frChg yT g s; u from
% GlobalTheory.T.integralGenerator_associated_localGenerator_CM_line; U'_b = (g b)^-1 with
% GlobalTheory.T.lieLetter_integral_in_integral_eigenbasis. Step 3: T' = W, kappa and i; yW and y'' by exists_frMap;
% hc = ChgBC.frChg_baseChange (y'' = yW changed by (g, s)); w = existsUnique_frChg y'' X^T 1; hint =
% integral_baseChange_chg (Lemma cv-cusp-cm-point-local-base); w0 from
% GlobalTheory.T.ruleValue_integral_of_supLattice_values_nonneg (= evR_mem_range_of_valSet_nonneg: exists_overLa
% for the finite place, allTest_of_secBaseChange, and mem_range_of_nval_nonneg); hYX (Y^T = X^T g); hval = evR_chg
% (existsUnique_frChg, frChg_refl_and_comp, evR_scale for the generator, and IsFrPoly); the witness u^-n w0, with hu'
% for iota(ca^-1) u^-1 s = 1, and map_evalFr.
\end{proof}
\leanref{Development/AbelianVarietyTheory/Proofs/CubicalLineBundles/ValuationDetLeOneOfEntries.lean}{GlobalTheory.LamF.letters_det_valuation_le}
\leanref{Development/AbelianVarietyTheory/Proofs/CubicalLineBundles/ExistsCompleteDVRCompletionIntegralIffStatement.lean}{GlobalTheory.T.exists_completeDVR_completion_integral_iff}
\leanref{Development/AbelianVarietyTheory/Proofs/CubicalLineBundles/RuleValueIntegralOfSupLatticeValuesNonnegStatement.lean}{GlobalTheory.T.ruleValue_integral_of_supLattice_values_nonneg}
\leanref{Development/AbelianVarietyTheory/Proofs/CubicalLineBundles/LocalLowerBoundCmPointFramePolynomialEveryPrime.lean}{GlobalTheory.LamT.LocClause.evR_chg}

\subsection{Archimedean bounds at CM points}\label{subsection-cv-cm-point-archimedean-bounds}

\begin{definition}\label{cv-cusp-cm-point-metric}
In Situation~\ref{cv-cusp-setting}, use the framed points of Lemma~\ref{cv-cusp-cm-point-polynomial}. Let
$\sigma\colon F_c\to\mathbf C$ be an embedding and $y$ a framed point of $X$ over $\mathbf C$, an $F_c$-algebra
through $\sigma$, with bases $u_\tau$ and generator $\vartheta$. Write its abelian variety as $\mathbf C^h/\Lambda$,
let $H$ be the Hermitian form of its polarization, linear in the first variable, let $\vartheta$ also denote the
function on $\mathbf C^h$ that represents the generator, and view each $u_{\tau,i}$ as a vector of $\mathbf C^h$. Let
$N$ be a real number and $\Gamma=(\Gamma_\tau)$ a family of complex $m_\tau\times m_\tau$ matrices. The framed point
$y$ is \emph{metric with constant $N$ and Gram blocks $\Gamma$} if $(\Gamma_\tau)_{ik}=\pi H(u_{\tau,k},u_{\tau,i})$
for all $\tau$, $i$ and $k$, and
\[
 N=\frac1{\operatorname{vol}F}\int_F|\vartheta(z)|^2e^{-\pi H(z,z)}\,dz
\]
for a fundamental domain $F$ of $\Lambda$.
\end{definition}
\leanref{Development/AbelianVarietyTheory/Definitions/CMLattices/IntegralOf.lean}{GlobalTheory.LamT.IsMetric}
\leanref{Development/AbelianVarietyTheory/Definitions/CMLattices/IntegralOf.lean}{GlobalTheory.LamT.IsMetricAt}

\begin{lemma}\label{cv-cusp-cm-point-metric-change}
In Situation~\ref{cv-cusp-setting}, let $\sigma\colon F_c\to\mathbf C$ be an embedding and let $y$ be a framed point
of $X$ over $\mathbf C$ through $\sigma$ that is metric with constant $N$ and Gram blocks $\Gamma$
(Definition~\ref{cv-cusp-cm-point-metric}). Let $Y=(Y_\tau)$ be a family of invertible complex $m_\tau\times m_\tau$
matrices and $a\in\mathbf C^\times$. Then $y\cdot(Y,a)$ is metric with constant $|a|^2N$ and Gram blocks
$(Y_\tau^*\Gamma_\tau Y_\tau)_\tau$.
\end{lemma}
\leanref{Development/AbelianVarietyTheory/Proofs/CubicalLineBundles/MetricChangeOfFrameStatement.lean}{GlobalTheory.T.metric_change_of_frame}
\leanref{Development/AbelianVarietyTheory/Proofs/CubicalLineBundles/MetricChangeOfFrame.lean}{GlobalTheory.MetricChange.isMetric_frChg}
\begin{proof}
The framed point $y\cdot(Y,a)$ has the point of $y$, so the uniformization, the form $H$ and the fundamental domain
$F$ of $y$ serve for it. Its generator $a\vartheta$ is represented by the function $a\vartheta$, whose mean is
$|a|^2N$. Its $i$-th basis vector at $\tau$ is $u'_{\tau,i}=\sum_j(Y_\tau)_{ji}u_{\tau,j}$, and the vector of
$\mathbf C^h$ of a linear combination of tangent vectors is the same combination of their vectors. As $H$ is linear
in the first variable and conjugate linear in the second,
\[
 \pi H(u'_{\tau,k},u'_{\tau,i})=\sum_{j,l}\overline{(Y_\tau)_{ji}}\,(\Gamma_\tau)_{jl}\,(Y_\tau)_{lk}
 =(Y_\tau^*\Gamma_\tau Y_\tau)_{ik}.
\]
\end{proof}
\leanref{Development/AbelianVarietyTheory/Proofs/CubicalLineBundles/MetricChangeOfFrame.lean}{GlobalTheory.MetricChange.metricAt_change}
% Lean: metricAt_change keeps u, Hm, chi and Fd, takes theta' = a theta (scaleSec) and W' = sum_j Y_ji W_j
% (lift_comb and GlobalTheory.T.tangentVector_existsUnique_uniformisingVector), and computes with hermForm_sum. The
% Lean's FrChg (Y^T) a is the text's y.(Y,a).

\begin{lemma}\label{cv-cusp-cm-point-hodge}
In Situation~\ref{cv-cusp-setting}, let $K\supset F_c$ be a number field, let $\ell>0$, and let a good CM point of
$X$ over $K$ with bound $\ell$ be given, with the notation of Definition~\ref{cv-cusp-cm-point-good}. Let
$s\colon K\to\mathbf C$ be an embedding, write $s(U_j)\in\mathbf C^g$ for the column vector of the $s(U_{jk})$, and
let $y_s$ be the base change of $y$ along $s$, a framed point of $X$ over $\mathbf C$ through the restriction of $s$
to $F_c$. Then there is a positive definite Hermitian $g\times g$ matrix $H_s$ with $n_s(\Delta)=(\det H_s)^{1/2}$
and $s(U_j)^*H_s\,s(U_j)\le(\ell\,n_s(\mathcal L_j))^2$ for every $j$, such that $y_s$ is metric with constant
$n_s(\Theta)^2$ and Gram blocks $(\Gamma_{s,\tau})_\tau$ (Definition~\ref{cv-cusp-cm-point-metric}), where
$\Gamma_{s,\tau}$ is the $m_\tau\times m_\tau$ matrix with entries $s(U_{e(\tau,i)})^*H_s\,s(U_{e(\tau,k)})$.
\end{lemma}
\leanref{Development/AbelianVarietyTheory/Proofs/CubicalLineBundles/ExistsMetricDataCmPointComplexEmbeddingStatement.lean}{GlobalTheory.T.exists_metricData_cmPoint_complexEmbedding}
\leanref{Development/AbelianVarietyTheory/Proofs/CubicalLineBundles/ExistsMetricDataCmPointComplexEmbedding.lean}{GlobalTheory.LamT.CMMetric.exists_isMetric}
\begin{proof}
Write the base change of $A_o$ along $s$ as $\mathbf C^g/\Lambda$, with the Hermitian form $H$ of its polarization,
linear in the first variable, a fundamental domain $F$ of $\Lambda$ and the function $\vartheta$ that represents
$\theta$. Let $W_k\in\mathbf C^g$ be the vector of the image of $v_k$, and $W$ the matrix with rows
$W_1,\dots,W_g$. For tangent vectors $x$ and $x'$ put $h(x,x')=\pi H(x',x)$, so that $h(x,x)$ is the squared Hodge
norm of $x$ (Definition~\ref{cv-cm-point-periods}(4)), and let $H_s$ be the matrix with entries
$(H_s)_{kl}=h(W_k,W_l)$. For $c,c'\in\mathbf C^g$, $c^*H_sc'=h\big(\sum_kc_kW_k,\sum_lc'_lW_l\big)$.

Step 1: $H_s$ is positive definite, as $H$ is and $W_1,\dots,W_g$ are linearly independent.

Step 2: the determinant. As $\Delta$ is adapted to $(A_o,\theta,v)$ at $s$ (Definition~\ref{cv-cusp-cm-key}(3)),
$n_s(\Delta)^2=\pi^g|\det W|^2/\operatorname{vol}F$. The polarization is principal, so the imaginary part of $H$ is
unimodular on $\Lambda$, and $\operatorname{vol}F\cdot\det H=1$, where $\det H$ is the determinant of the matrix
$(H(e_a,e_b))_{a,b}$ of $H$ in the standard basis. As $(H_s)_{kl}=\pi H(W_l,W_k)$, $\det H_s=\pi^g|\det W|^2\det H$.
Hence $n_s(\Delta)^2=\det H_s$.

Step 3: the bounds. Fix $j=e(\tau,i)$, and suppose that $\tau=\tau^\Phi(m)$; if $\tau=\tau'(m)$, replace $A^\Phi$,
$v^\Phi$, $\varphi_i$, $h_i$ and $\mathcal M^\Phi_m$ by $A'$, $v'$, $\varphi'_i$, $h'_i$ and $\mathcal M'_m$. Let
$h_{A^\Phi}$ be the form defined as $h$ for the base change of $A^\Phi_K$ along $s$, so that $h_{A^\Phi}(x,x)$ is
the squared Hodge norm of a tangent vector $x$ of it. By (f) of the definition,
$\sum_lU_{jl}v_l=d\varphi_i(v^\Phi_m)$ and $\mathcal L_j=\mathcal M^\Phi_m$, so
$s(U_j)^*H_s\,s(U_j)=h\big(d\varphi_i(v^\Phi_m),d\varphi_i(v^\Phi_m)\big)$, with $v^\Phi_m$ and $\varphi_i$ base
changed along $s$. By (d), the polarization of $B$ pulls back along $\varphi_i$ to that of $A^\Phi$ composed with
$h_i$. Over $\mathbf C$ the Hermitian form of a pulled back polarization is the pulled back form, so
$h(d\varphi_i(x),d\varphi_i(x'))=h_{A^\Phi}(x,h_ix')$ for all tangent vectors $x$ and $x'$ of the base change of
$A^\Phi_K$ along $s$. By (b), $h_i$ multiplies $v^\Phi_m$ by $t=s(\tau_K(h_i))$, and $h_{A^\Phi}$ is linear in its
second variable, so $s(U_j)^*H_s\,s(U_j)=t\,h_{A^\Phi}(v^\Phi_m,v^\Phi_m)$. Here the left side is real, and
$h_{A^\Phi}(v^\Phi_m,v^\Phi_m)$ is positive and equal to $n_s(\mathcal M^\Phi_m)^2$ by (e); taking real parts,
\[
 s(U_j)^*H_s\,s(U_j)=\operatorname{Re}\big(t\,h_{A^\Phi}(v^\Phi_m,v^\Phi_m)\big)
 \le|t|\,h_{A^\Phi}(v^\Phi_m,v^\Phi_m)=|t|\,n_s(\mathcal M^\Phi_m)^2.
\]
As $s\circ\tau_K$ is an embedding of $E$ in $\mathbf C$, (d) gives $|t|\le\ell^2$, and
$s(U_j)^*H_s\,s(U_j)\le(\ell\,n_s(\mathcal L_j))^2$. The inequality in the display is in fact an equality: $t$ is
real, as $h_i$ is fixed by complex conjugation and every embedding of the CM field $E$ in $\mathbf C$ commutes
with it, and $t>0$, as the left side is positive.

Step 4: the metric. By (g), $y_s$ has the base change of $o$ as its point, the generator $\theta$, and at $\tau$ the
basis vectors $u_{\tau,i}$ with vectors $\sum_ls(U_{e(\tau,i),l})W_l$. As $\Theta$ is adapted to $(A_o,\theta,v)$ at
$s$, the mean of $|\vartheta(z)|^2e^{-\pi H(z,z)}$ over $F$ is $n_s(\Theta)^2$. And
$\pi H(u_{\tau,k},u_{\tau,i})=h(u_{\tau,i},u_{\tau,k})=s(U_{e(\tau,i)})^*H_s\,s(U_{e(\tau,k)})$, which is the
$(i,k)$ entry of $\Gamma_{s,\tau}$.
\end{proof}
\leanref{Development/AbelianVarietyTheory/Proofs/CubicalLineBundles/LamLenOfIsHerm.lean}{GlobalTheory.LamF.lamLen_of_isHerm}
\leanref{Development/AbelianVarietyTheory/Proofs/CubicalLineBundles/TopLieVectorNormSqEqGramDet.lean}{GlobalTheory.CM.TopLieVector.sq_eq_det_gram}
\leanref{Development/AbelianVarietyTheory/Proofs/CubicalLineBundles/FormLetterLe.lean}{GlobalTheory.LamF.form_letter_le}
% Lean: exists_isMetric takes h from GlobalTheory.CM.LieMetric.exists_unique_isHerm (the Hodge metric of B) and
% H_s = h.gram s of the standard basis; Step 1 is gram_stdBasis_posDef_and_hq. Steps 2 and 3 are lamLen_of_isHerm:
% Step 2 is TopLieVector.sq_eq_det_gram (KeyArchAt and volume_mul_det_eq_one_of_unimodular), Step 3 is form_letter_le
% (lamEq_of_weilIs_all_points from PolClause, GlobalTheory.T.hermitianLength_image_eigenBasisVector_le =
% LenOne.form_image_le, which states the comparison as the inequality of the display: the pulled back form is
% HermHom.form_diff_eq, and Re(t h) <= |t| h is re_form_smul_le; vecLine_norm_sq_eq_form for n_s; hb for |t| <= l^2).
% The closing equality is not in the Lean. Step 4 uses exists_fibre and exists_unif (from
% KeyArchAt) and LieMetric.form_eq_pi_mul_hermForm.

\begin{lemma}\label{cv-cusp-cm-point-cusp-bound-ortho-hermitian}
In Situation~\ref{cv-cusp-setting}, let $\sigma$ be a good embedding and $Z$ a component at $\sigma$, with the lattice
$L_Z$, its alternating form $\langle\ ,\ \rangle$ and the Hermitian model $(G,\eta,\varphi_Z)$ of $L_Z$
(Definition~\ref{cv-uniformization-model}) that the compatibility conditions of the uniformization attach to $Z$, as
listed at the beginning of the proof of Lemma~\ref{cv-cover-uniformization}. Suppose that $|\tau(\eta)|=\xi$ for every
embedding $\tau\colon E\to\mathbf C$. Let $x$ be an orthonormal frame of $Z$, and let the following data at $x$ be given,
which the compatibility conditions provide.
\begin{enumerate}
\item[(W)] An indexing of the coordinates of $\mathbf C^g$ by the pairs $(\tau,a)$ of
  Situation~\ref{cv-uniformization-situation}; a family uniformization of $X$ at $\sigma$ restricted to
  $\mathbf Q(\sqrt{-d})$, for the type that gives the coordinate $(\tau,a)$ the embedding $\sigma\circ\tau$
  (Theorem~\ref{cv-uniformization-family}), with exponentials $\exp_{i,P}$; for every component $Z'$ at $\sigma$ an index
  $i_{Z'}$ of the family, with $L_{Z'}=L_{i_{Z'}}$, and the Hermitian model $(G,\eta,\varphi_{Z'})$ of $L_{Z'}$; the
  period $P=P^Z_x$ of $x$ on $L_Z$, given by the formula for periods of Situation~\ref{cv-uniformization-situation}
  through $\varphi_Z$, which is a period of $L_Z$ of this type; the equality of the point of the framed point $w_x$ at
  $x$ with $y_{i_Z}(P)$; the derivative of $\exp_{i_Z,P}$ at $0$ in the direction of the unit vector of the coordinate
  $(\tau,a)$, which is the basis vector $u_{\tau,a+1}$ of $w_x$; and, if $|\delta_{\tau,i}(x)|=1$ for all $\tau$ and $i$,
  the existence of data as in (U).
\item[(U)] A Hermitian form $H$ on $\mathbf C^g$, linear in the first variable and positive definite, with
  $\operatorname{Im}H(Pa,Pa')=\langle a,a'\rangle$ for all $a,a'\in L_Z$; numbers $\chi(u)$ of absolute value $1$ for
  $u\in P(L_Z)$; a function $\vartheta$ on $\mathbf C^g$; for the trivializations of the line bundle of the polarization
  over the sets $U_j$ of Theorem~\ref{cv-uniformization-family}(4), nowhere vanishing holomorphic functions $F_j$ on the
  sets of the $z$ with $\exp_{i_Z,P}(z)\in U_j$, with $F_j(z)=g_{jj'}(\exp_{i_Z,P}(z))F_{j'}(z)$ where both sides are
  defined, $F_j(0)=\varrho_j(y_{i_Z}(P))$, and
  $F_j(z)=\chi(u)\exp\big(\pi H(z,u)+\tfrac\pi2H(u,u)\big)F_j(z+u)$ for $u\in P(L_Z)$, such that the section of $w_x$
  has the value $\vartheta(z)F_j(z)$ at $\exp_{i_Z,P}(z)$ in the trivialization over $U_j$; and a fundamental domain of
  $P(L_Z)$ over which the mean of $|\vartheta(z)|^2e^{-\pi H(z,z)}$ is $1$.
\end{enumerate}
Of these data the proof uses only the indexing, the model $(G,\eta,\varphi_Z)$, the period $P$ and its formula, that
$P$ is a period, and the form $H$ with its three properties; the other data in (W) and (U) are not used. Then
$H(z,w)=2\xi\sum_kz_k\bar w_k$ for all $z,w\in\mathbf C^g$.
\end{lemma}
\leanref{Development/AbelianVarietyTheory/Proofs/CubicalLineBundles/ImDiagEqTrace.lean}{GlobalTheory.OrthoHerm.hm_eq_smul_one}
\begin{proof}
Step 1: $H$ is determined by the values of $\operatorname{Im}H$ on $P(L_Z)\times P(L_Z)$. Indeed
$\operatorname{Re}H(z,w)=\operatorname{Im}H(iz,w)$, so $H$ is determined by the $\mathbf R$-bilinear form
$\operatorname{Im}H$, and an $\mathbf R$-bilinear form is determined by its values on a set that spans $\mathbf C^g$ over
$\mathbf R$. The set $P(L_Z)$ does, as $P$ is a period of $L_Z$ (Definition~\ref{cv-cm-point-periods}(3)), whose image
is a lattice in $\mathbf C^g$.

Step 2: a diagonal form. Index the coordinates of $\mathbf C^g$ as in (W), and attach to the coordinates $(\tau,a)$ and $(\bar\tau,a)$ the embedding $\tau\in\Phi$. Let $D$ be the Hermitian form with
$D(z,w)=\sum_k2\operatorname{Im}\tau_k(\eta)\,z_k\bar w_k$, where $\tau_k$ is the embedding attached to the coordinate
$k$. We show that $\operatorname{Im}D(Pa,Pa')=\langle a,a'\rangle$ for $a,a'\in L_Z$. As $\Phi$ is a CM type,
$\operatorname{Tr}_{E/\mathbf Q}(t)=\sum_{\tau\in\Phi}2\operatorname{Re}\tau(t)$ for $t\in E$. Fix $\tau\in\Phi$ and put
$c=x_\tau^{-1}\tau(\varphi_Za)$ and $c'=x_\tau^{-1}\tau(\varphi_Za')$. As $x$ is orthonormal,
$x_\tau^{\mathsf T}\tau(G)\bar x_\tau$ is the diagonal matrix $J$ with entries $1$, $q$ times, followed by $-1$; so
$\tau(h_G(\varphi_Za,\varphi_Za'))=c^{\mathsf T}J\bar c'$, and
$h_{\tau(G)}(x_{\tau,b},\tau(\varphi_Za))=(J\bar c)_b=\bar c_b$ for $b\le q$. By the formula for periods, the
coordinates of $Pa$ attached to $\tau$ are the $c_b$ with $b>q$ and the $\bar c_b$ with $b\le q$. Hence
\[
 \sum_{\tau_k=\tau}\operatorname{Im}\big((Pa)_k\overline{(Pa')_k}\big)=\sum_{b>q}\operatorname{Im}(c_b\bar c'_b)+
 \sum_{b\le q}\operatorname{Im}(\bar c_bc'_b)=-\operatorname{Im}\tau\big(h_G(\varphi_Za,\varphi_Za')\big).
\]
As $\bar\eta=-\eta$, $\tau(\eta)=i\operatorname{Im}\tau(\eta)$, and
$2\operatorname{Re}\big(\tau(\eta)\tau(t)\big)=-2\operatorname{Im}\tau(\eta)\operatorname{Im}\tau(t)$ for $t\in E$.
Summing over $\tau\in\Phi$,
$\operatorname{Im}D(Pa,Pa')=\operatorname{Tr}_{E/\mathbf Q}\big(\eta\,h_G(\varphi_Za,\varphi_Za')\big)=\langle a,a'\rangle$,
the last equality as $(G,\eta,\varphi_Z)$ is a Hermitian model.

Step 3: by Steps 1 and 2, $H=D$. As $H$ is positive definite, $2\operatorname{Im}\tau_k(\eta)=H(e_k,e_k)>0$ for every $k$,
and as $\tau_k(\eta)$ is purely imaginary, $\operatorname{Im}\tau_k(\eta)=|\tau_k(\eta)|=\xi$.
\end{proof}
\leanref{ReciprocityTheory/ReductiveGroups/ArithmeticGeometry/PullbackSections/HolomorphicFamilies/UnitHmA.lean}{ReciprocityTheory.ChartRestriction.G.PullbackSections.UH.herm_eq_of_im}
\leanref{Development/AbelianVarietyTheory/Proofs/CubicalLineBundles/ImDiagEqTrace.lean}{GlobalTheory.OrthoHerm.im_diag_eq_trace}
\leanref{Development/AbelianVarietyTheory/Proofs/UnitaryExistence/MemIffConjugateNotMem.lean}{GlobalTheory.UnitFrame.im_sum_periods}
% Lean: hm_eq_smul_one, for any data W : UnfAt U k x (the text's (W): e, F, ik, phi, isom, Pm, pm_eq, per, pt, der,
% unit) and Wu : UnitAt W (the text's (U): Hm, chi, fr, theta, Fd with herm, pos, c1 to c9); it uses of Wu only herm,
% pos and c1 (with W.isom for Im H(Pa, Pa') = Tr(-xi_0 h(phi a, phi a')), so eta = -xi_0 of the Lean), and of W the
% period W.per (W.per.2.2.2.1, the spanning over R, in Step 1) and W.pm_eq with relab_good and
% FrameUnif.exists_letter_equiv for the indexing of the coordinates (hper). Step 1 is UH.herm_eq_of_im; Step 2 is
% im_diag_eq_trace (UnitFrame.trace_eq_sum_re and im_sum_periods, which is the display); Step 3 is hd with Wu.pos and
% hxi. Size: about 290 lines.

\begin{lemma}\label{cv-cusp-cm-point-cusp-bound-ortho}
In Situation~\ref{cv-cusp-setting}, let $\sigma$ be a good embedding, $Z$ a component at $\sigma$ and $x$ an
orthonormal frame of $Z$, and let $w_x$ be the complex framed point of $X$ over $\sigma$ that the uniformization
attaches to $x$. Then $w_x$ is metric with constant $1$ and with the Gram blocks $2\pi\xi$ times the identity
matrices (Definition~\ref{cv-cusp-cm-point-metric}).
\end{lemma}
\leanref{Development/AbelianVarietyTheory/Proofs/CubicalLineBundles/UnifMetricScalarStatement.lean}{GlobalTheory.T.unif_metric_scalar}
\leanref{Development/AbelianVarietyTheory/Proofs/CubicalLineBundles/UnifMetricScalar.lean}{GlobalTheory.OrthoMetric.isMetric_of_isOrtho}
\begin{proof}
Let $y=w_x$, with abelian variety $A_y$, bases $u_{\tau,i}$ and section $\vartheta$. As $x$ is orthonormal,
$|\delta_{\tau,i}(x)|=1$ for all $\tau$ and $i$. So the compatibility conditions of the uniformization at $\sigma$, as
listed at the beginning of the proof of Lemma~\ref{cv-cover-uniformization}, give at $x$: the period $P=P^Z_x$ on the
lattice $L_Z$ of $Z$, with its Hermitian model $(G,\eta,\varphi_Z)$; the map $p\colon\mathbf C^g\to A_y(\mathbf C)$ that
sends $z$ to the image of $z$ under the isomorphism $\mathbf C^g/P(L_Z)\to A_y$, which carries the unit vector $e_k$ of
the coordinate $k=(\tau,i)$, as a tangent vector at $0$, to $u_{\tau,i}$; a positive definite Hermitian form $H_x$ with
$\operatorname{Im}H_x(Pa,Pa')=\langle a,a'\rangle$ on $L_Z$; numbers $\chi_x(u)$; a function $\vartheta_x$; and nowhere
vanishing holomorphic functions $F^x_j$ that satisfy the transition rule of the line bundle, take the values of the
rigidification at $0$, satisfy the functional equation of Definition~\ref{cv-origin-hodge-norm-gram}(1) for
$(H_x,\chi_x)$, and represent $\vartheta$ through $\vartheta_x$; and the mean of $|\vartheta_x(z)|^2e^{-\pi H_x(z,z)}$
over a fundamental domain $F$ of $P(L_Z)$ is $1$. The normalization of the uniformization by $\xi$ says that
$|\tau(\eta)|=\xi$ for every embedding $\tau\colon E\to\mathbf C$.

Step 1: a key datum. The universal abelian scheme over $X$ is a principally polarized abelian scheme of relative
dimension $g=fr$, with the symmetric line bundle of its polarization rigidified along the zero section. Its fibre $A_y$
at the point of $y$, the pullback along that point, is therefore a principally polarized abelian variety over
$\mathbf C$ of dimension $g$, with the pulled back line bundle, symmetric and rigidified along the zero section: being an
abelian scheme of relative dimension $g$, a relatively ample line bundle defining a principal polarization, the symmetry
and the rigidification are all preserved by base change. As $p+q\ge1$, $g\ge1$. The Lie algebra of $A_y$ is a complex
vector space of dimension $g$, so it has a basis $v$. The section $\vartheta$ of $y$ generates the space of sections of
this line bundle, so it is nonzero. So $(A_y,\vartheta,v)$ is a key datum (Definition~\ref{cv-cusp-cm-key}(1)).

Step 2: the uniformization. The map $p$ is a surjective homomorphism with kernel $P(L_Z)$, holomorphic in the sense of
Definition~\ref{cv-origin-hodge-norm-gram}(1) (Theorem~\ref{cv-uniformization-family}). With the trivializing functions
$F^x_j$ it is a uniformization of $A_y$ for its line bundle with Appell--Humbert data $(H_x,\chi_x)$, whose trivializing
functions take the values of the rigidification at $0$, and in which $\vartheta_x$ represents $\vartheta$. The kernel
$P(L_Z)$ spans $\mathbf C^g$ over $\mathbf R$, as $A_y(\mathbf C)$ is compact. The tangent vector $u_{\tau,i}$ is the
derivative of $p$ at $0$ in the direction $e_{(\tau,i)}$, so $e_{(\tau,i)}$ is its differential, which is unique by
Lemma~\ref{cv-cusp-cm-point-eigen-degree-rank-one-norm-riemann}(2).

Step 3: the form. The data at $x$ listed at the beginning of the proof are data (W) and (U) of
Lemma~\ref{cv-cusp-cm-point-cusp-bound-ortho-hermitian}, with $H=H_x$, and $|\tau(\eta)|=\xi$; by that lemma,
$H_x(z,w)=2\xi\sum_kz_k\bar w_k$.

Step 4: the metric. The data of Steps 1 and 2, with the fundamental domain $F$, compute the constant and the Gram blocks
of Definition~\ref{cv-cusp-cm-point-metric} for $y$: the constant is the mean over $F$, which is $1$, and the $(i,k)$
entry of the Gram block at $\tau$ is $\pi H_x(e_{(\tau,k)},e_{(\tau,i)})$, which by Step 3 is $2\pi\xi$ if $i=k$ and $0$
otherwise.
\end{proof}
\leanref{Development/AbelianVarietyTheory/Proofs/CubicalLineBundles/UniversalTripleFibreKeyDataNonemptyField.lean}{GlobalTheory.LamT.nonempty_fibKey}
\leanref{Development/AbelianVarietyTheory/Proofs/CubicalLineBundles/RigidifiedSectionCorrespondenceExistsUniqueNonva.lean}{GlobalTheory.SecTieEx.secTie_exists_unique}
\leanref{Development/AbelianVarietyTheory/Proofs/CubicalLineBundles/UnifMetricScalar.lean}{GlobalTheory.OrthoMetric.exists_unif}
% Lean: isMetric_of_isOrtho: the data at x are PullbackSections.nonempty_unfAt (W: the family uniformization W.F, the
% period W.Pm, the point and the derivatives W.der) and nonempty_unitAt with PfSupAtUnit.dmin_ortho (Wu: Hm, chi, the
% trivializing functions fr, theta and Fd, with c1 to c9; c9 is the mean 1), fields of IsUnf. Step 1 is
% LamT.nonempty_fibKey (sabPull and PAT.exists_of_sabMap for the fibre, principalPolarisation_exists_unique_section,
% KeyPow.exists_isLieBasis; g >= 1 from 1 <= p + q), with SecTieEx.secTie_exists_unique for the section th tied to the
% section of y and Frames.secNZ_of_secGen for th nonzero. Step 2 is OrthoMetric.exists_unif (fibreUnif, the exponential
% F.ex lifted to the fibre through the cartesian square, carry with ConjUnif.exists_inv for the trivializing functions,
% val_tie and val_rig for IsRig and IsFn, AHWeil.tangent_and_weil_of_isAH and eq_of_isDerivAt0 for the differentials).
% Step 3 is OrthoHerm.hm_eq_smul_one. Step 4 is isMetricAt_of, with hG. Size: about 5,600 lines, of which 4,772 (by
% tools/lean_size.py) are below nonempty_fibKey, 510 below secTie_exists_unique, 450 below exists_unif and 290 below
% hm_eq_smul_one. Step 1 follows the Lean's proof of nonempty_fibKey, which has been read: exists_pat_of_obj (the
% universal object is a PAT of dimension [E:Q]/2 (p + q)), sabPull (the pullback along the point),
% PAT.exists_of_sabMap (it is a PAT of the same dimension: PolPull.pol_pull_of_relAmple_local for the ampleness and the
% principal polarization, smoothness by base change, the symmetry), g >= 1 from 1 <= p + q, KeyPow.exists_isLieBasis
% (a Lie basis); the Lean's section th comes from principalPolarisation_exists_unique_section and is tied to the section
% of y by secTie_exists_unique, nonzero by secNZ_of_secGen, while the text takes the section of y on the fibre, which is
% a generator, directly. The 4,772 lines are mostly base-change facts (the pullback of the semi-abelian scheme with its
% cubical line bundle, relative ampleness locally on the base); the text's short argument is a correct proof of Step 1.

\begin{lemma}\label{cv-cusp-cm-point-cusp-bound-unique-key}
Let $(A,\theta,v)$ be a key datum of dimension $g$ over $\mathbf C$ (Definition~\ref{cv-cusp-cm-key}(1)), with the line
bundle $\mathcal L$ of its polarization trivialized on an open cover $(U_\alpha)$ and its rigidification, a generator
$r$ of $0^*\mathcal L$ with values $r_\alpha$ in the trivializations. For $j=1,2$ let $u_j\colon\mathbf C^g\to
A(\mathbf C)$ be a uniformization of $A$ for $\mathcal L$ (Definition~\ref{cv-origin-hodge-norm-gram}(1), with
$T=\operatorname{Spec}\mathbf C$, $G=A$ and the identity) with kernel $\Lambda_j$, trivializing functions
$\theta^{(j)}_\alpha$ and Appell--Humbert data $(H_j,\chi_j)$, such that $\theta^{(j)}_\alpha(0)=r_\alpha$ whenever
$0\in U_\alpha$. Let $\vartheta_j$ be a function on $\mathbf C^g$ with $\theta_\alpha(u_j(z))=\vartheta_j(z)\,
\theta^{(j)}_\alpha(z)$ for every $\alpha$ and $z\in u_j^{-1}(U_\alpha)$, where $\theta_\alpha$ is $\theta$ in the
trivialization over $U_\alpha$, and let $F_j$ be a fundamental domain of $\Lambda_j$. Then:
\begin{enumerate}
\item $\displaystyle\frac1{\operatorname{vol}F_1}\int_{F_1}|\vartheta_1(z)|^2e^{-\pi z^{\mathsf T}H_1\bar z}\,dz=
  \frac1{\operatorname{vol}F_2}\int_{F_2}|\vartheta_2(z)|^2e^{-\pi z^{\mathsf T}H_2\bar z}\,dz$;
\item if $x$ and $x'$ are tangent vectors of $A$ at the origin, with differentials $a$ and $b$ for $u_1$ and $a'$ and
  $b'$ for $u_2$, then $a^{\mathsf T}H_1\bar b=a'^{\mathsf T}H_2\bar b'$.
\end{enumerate}
\end{lemma}
\leanref{Development/AbelianVarietyTheory/Proofs/CubicalLineBundles/ThetaSectionNormAndHermitianFormIndepUniformisatStatement.lean}{GlobalTheory.T.thetaSection_norm_and_hermitianForm_indep_uniformisation}
\leanref{Development/AbelianVarietyTheory/Proofs/CubicalLineBundles/ThetaSectionNormAndHermitianFormIndepUniformisat.lean}{GlobalTheory.KeyUniq.norm_eq_and_hermForm_eq}
\begin{proof}
Step 1: the comparison. By Lemma~\ref{cv-cusp-cm-point-eigen-degree-rank-one-norm-riemann}(2), for $u_1$ and for $u_2$
every tangent vector of $A$ at the origin has exactly one differential, and every vector of $\mathbf C^g$ is the
differential of exactly one tangent vector. Lemma~\ref{cv-origin-hodge-norm-form-lift}, applied to $u_2$ and the
homomorphism $u_1$, gives a $\mathbf C$-linear map $F\colon\mathbf C^g\to\mathbf C^g$ with $u_2\circ F=u_1$. As in Step 2
of the proof of Lemma~\ref{cv-cusp-cm-point-eigen-degree-rank-one-norm-riemann}, $F^{-1}(\Lambda_2)=\Lambda_1$ and the
kernel of $F$, a subspace contained in the discrete group $\Lambda_1$, is $0$; so $F$ is bijective and
$F(\Lambda_1)=\Lambda_2$. By the chain rule, as in the proof of Lemma~\ref{cv-origin-hodge-norm-form-compare}, a tangent
vector has differential $a$ for $u_1$ if and only if it has differential $F(a)$ for $u_2$.

Step 2: the Hermitian forms. Lemma~\ref{cv-origin-hodge-norm-form-pullback}, for the identity of $A$, $u_1$, $u_2$, $F$
and $n=1$, gives $a^{\mathsf T}H_1\bar b=F(a)^{\mathsf T}H_2\overline{F(b)}$ for all $a,b\in\mathbf C^g$. In (2),
$a'=F(a)$ and $b'=F(b)$ by Step 1, as differentials are unique; this is (2).

Step 3: the means. Let $N_j$ be the mean in (1) for $j=1,2$. Let $W$ be the matrix whose $i$-th row is the
differential for $u_1$ of $v_i$; by Step 1 the matrix $W'$ of the differentials for $u_2$ has the rows $F(W_i)$, so
$\det W'=\det F\cdot\det W$. The image $F(F_1)$ is a fundamental domain of $\Lambda_2$, of volume
$|\det F|^2\operatorname{vol}F_1$ for the Lebesgue measure of $\mathbf C^g=\mathbf R^{2g}$, and all fundamental domains of
$\Lambda_2$ have the same volume; so $\operatorname{vol}F_2=|\det F|^2\operatorname{vol}F_1$. Hence
\[
 \eta=\frac{\operatorname{vol}F_1}{\pi^g|\det W|^2}=\frac{\operatorname{vol}F_2}{\pi^g|\det W'|^2}.
\]
Theorem~\ref{cv-definite-key-formula}(5), for $(A,\theta,v)$ written through $u_1$ and through $u_2$, gives
$|\kappa(A,\theta,v)|N_1^4\eta^2=(2\pi)^{2g}=|\kappa(A,\theta,v)|N_2^4\eta^2$. As $\kappa(A,\theta,v)\ne0$ and $\eta>0$,
$N_1^4=N_2^4$, and $N_1=N_2$ as both are at least $0$.
\end{proof}
\leanref{Development/AbelianVarietyTheory/Proofs/CubicalLineBundles/ThetaSectionNormAndHermitianFormIndepUniformisat.lean}{GlobalTheory.KeyUniq.volume_fund}
\leanref{Development/AbelianVarietyTheory/Proofs/ThetaLawProofs/UnimodularAndTangentOfIsAH.lean}{GlobalTheory.UnifCompare.hermForm_eq}
% Lean: norm_eq_and_hermForm_eq: unimodular_and_tangent_of_isAH (with thetaUnifAV_all) for the tangent bijections t1,
% t1', t2' (Lemma cv-cusp-cm-point-eigen-degree-rank-one-norm-riemann(2)); existsUnique_linear_lift_of_isHolTo (Lemma
% cv-origin-hodge-norm-form-lift) for F; UnifCompare.injective, mem_iff and isDiff_iff (Step 1);
% UnifCompare.hermForm_eq (Step 2, the argument of Lemma cv-origin-hodge-norm-form-pullback). Step 3 is the last clause of
% keyFormulaPow_and_keyDegree.1, the key formula in its power form with exponents 4 N(g) and 2 N(g) (Theorem
% cv-definite-key-formula(5) is the case N(g) = 1 of the text), with volume_fund (volume_image, through
% addHaar_image_linearMap, and IsAddFundamentalDomain.measure_eq for two fundamental domains) and hdet for det W'.
% Size: about 190 lines below norm_eq_and_hermForm_eq beyond the statements the book cites.

\begin{lemma}\label{cv-cusp-cm-point-cusp-bound-unique}
In Situation~\ref{cv-cusp-setting}, let $\sigma\colon F_c\to\mathbf C$ be an embedding and $y$ a framed point of $X$
over $\mathbf C$ through $\sigma$. If $y$ is metric with constant $N$ and Gram blocks $\Gamma$ and also with constant
$N'$ and Gram blocks $\Gamma'$ (Definition~\ref{cv-cusp-cm-point-metric}), then $N=N'$ and $\Gamma=\Gamma'$.
\end{lemma}
\leanref{Development/AbelianVarietyTheory/Proofs/CubicalLineBundles/MetricDataUniqueStatement.lean}{GlobalTheory.T.metric_data_unique}
\leanref{Development/AbelianVarietyTheory/Proofs/CubicalLineBundles/MetricDataUnique.lean}{GlobalTheory.LamT.MetUniq.isMetric_unique}
\begin{proof}
The constant and the Gram blocks of Definition~\ref{cv-cusp-cm-point-metric} are computed from the following data: a key
datum $(A_1,\theta_1,v_1)$ over $\mathbf C$ (Definition~\ref{cv-cusp-cm-key}(1)) with an isomorphism $m_1$ of $A_1$ onto
the abelian variety of $y$ that respects the polarizations and the rigidifications and carries $\theta_1$ to the generator
$\vartheta$ of $y$; a uniformization $u_1$ of $A_1$ for its line bundle with Appell--Humbert data $(H_1,\chi_1)$, whose
trivializing functions take the values of the rigidification at $0$; a function $\vartheta_1$ representing $\theta_1$ and a
fundamental domain $F_1$ of the kernel of $u_1$, from which $N$ is the mean of the definition; and for each $\tau$ and
$i$ a tangent vector $w^{(1)}_{\tau,i}$ of $A_1$ with $m_1\circ w^{(1)}_{\tau,i}=u_{\tau,i}$, with differential
$W^{(1)}_{\tau,i}$ for $u_1$, so that $(\Gamma_\tau)_{ik}=\pi\,W^{(1)\mathsf T}_{\tau,k}H_1\overline{W^{(1)}_{\tau,i}}$. Let
$(A_2,\theta_2,v_2)$, $m_2$, $u_2$, $(H_2,\chi_2)$, $\vartheta_2$, $F_2$, $w^{(2)}_{\tau,i}$ and $W^{(2)}_{\tau,i}$ be the data
for $N'$ and $\Gamma'$.

Step 1: the two models. The isomorphism $\psi=m_1^{-1}\circ m_2\colon A_2\to A_1$ respects the polarizations and the
rigidifications and carries $\theta_2$ to $\theta_1$. Let $v'_2$ be the basis of the Lie algebra of $A_2$ that $\psi$
carries to $v_1$; then $(A_2,\theta_2,v'_2)$ is a key datum, and $\psi$ is an isomorphism of key data onto
$(A_1,\theta_1,v_1)$. A tangent vector of $A_2$ is determined by its image under $m_2$, so $\psi$ carries
$w^{(2)}_{\tau,i}$ to $w^{(1)}_{\tau,i}$.

Step 2: transport. The map $u_3=\psi^{-1}\circ u_1$ is a uniformization of $A_2$ for its line bundle, with the kernel
of $u_1$, Appell--Humbert data $(H_1,\chi_3)$ for some $\chi_3$, and trivializing functions those of $u_1$ multiplied by
the isomorphism of the line bundles in the trivializations; they take the values of the rigidification of $A_2$ at $0$,
and $\vartheta_1$ represents $\theta_2$ through $u_3$, as $\psi$ carries $\theta_2$ to $\theta_1$. So the mean computed
from $u_3$, $\vartheta_1$ and $F_1$ is $N$. A tangent vector $w$ of $A_2$ has the differential for $u_3$ that $\psi\circ w$
has for $u_1$; so $W^{(1)}_{\tau,i}$ is the differential of $w^{(2)}_{\tau,i}$ for $u_3$.

Step 3: one key datum. Lemma~\ref{cv-cusp-cm-point-cusp-bound-unique-key}, for $(A_2,\theta_2,v'_2)$, the uniformizations
$u_3$ and $u_2$, the functions $\vartheta_1$ and $\vartheta_2$ and the domains $F_1$ and $F_2$, gives $N=N'$ by (1), and
by (2), for the tangent vectors $w^{(2)}_{\tau,k}$ and $w^{(2)}_{\tau,i}$,
$W^{(1)\mathsf T}_{\tau,k}H_1\overline{W^{(1)}_{\tau,i}}=W^{(2)\mathsf T}_{\tau,k}H_2\overline{W^{(2)}_{\tau,i}}$, that is,
$\Gamma=\Gamma'$.
\end{proof}
\leanref{Development/AbelianVarietyTheory/Proofs/CubicalLineBundles/UniversalTripleFibreKeyDataBaseChangeKeyPullbackStatement.lean}{GlobalTheory.T.universalTripleFibreKeyData_baseChange_keyPullback}
\leanref{Development/AbelianVarietyTheory/Proofs/CubicalLineBundles/ExistsAppellHumbertUniformisationOfKeyDatumIsoStatement.lean}{GlobalTheory.T.exists_appellHumbert_uniformisation_of_keyDatumIso}
% Lean: isMetricAt_unique, with the fibre key data Phi and Phi' (FibKey: a key datum Phi.key with a cartesian map
% Phi.m onto the fibre of the universal abelian scheme at y.x, and SecTie for the section). Step 1 is
% T.universalTripleFibreKeyData_baseChange_keyPullback = LamT.fibKey_keyPB for phi the identity of C (with
% SecCalc.sec_calc for SecBC (id) 1 y.x y.x s s): sabMap_cancel_and_comp_over (the map n of the two models), pat_pullback
% (the line bundles, with the rigidification), ArcQBC.exists_pullSec and TieCalc.secTie_calc (the sections) and
% exists_isLieBasis_baseChange (v'_2), which gives KeyPB; the tangent vectors through Chart.comb_push and Chart.tan_push.
% Step 2 is T.exists_appellHumbert_uniformisation_of_keyDatumIso = KeyTransport.uniformisation_transport (keyPB_cancel to
% invert the key pullback, exists_push_key for u_3, with the same theta and fundamental domain). Step 3 is uniq_cast,
% that is, Lemma cv-cusp-cm-point-cusp-bound-unique-key. Size: about 5,400 lines below isMetric_unique, of which 5,606
% (tools/lean_size.py) are below fibKey_keyPB, which works over any ring homomorphism S -> S' and verifies the cartesian
% squares, and 870 below uniformisation_transport. The proof of fibKey_keyPB has been read: for phi the identity of C its
% steps are those of Step 1 (n = psi from the cancellation of the two cartesian maps m_1 and m_2 onto the same fibre; the
% line bundles and rigidifications from pat_pullback; theta_2 to theta_1 as both are tied to the section of y; v'_2 the
% basis carried to v_1), and its lines are the general base-change facts, which the identity does not need. The text's
% Step 1 is a correct proof of what is used.

\begin{lemma}\label{cv-cusp-cm-point-cusp-bound-good}
In Situation~\ref{cv-cusp-setting}, let $\kappa=(\alpha,n)$ be a dominant weight and $f\in S_\kappa$, let $\sigma$ be a
good embedding, and let $y$ be a framed point of $X$ over $\mathbf C$ through $\sigma$ that is metric with constant $N$
and with the identity matrices as Gram blocks (Definition~\ref{cv-cusp-cm-point-metric}). Then
$|f(y)|\le N^{n/2}\|f\|_{\sigma,\infty}$.
\end{lemma}
\leanref{Development/AbelianVarietyTheory/Proofs/CubicalLineBundles/CuspFormValueLeSupNormAtOrthonormalFrameStatement.lean}{GlobalTheory.T.cuspForm_value_le_supNorm_at_orthonormal_frame}
\leanref{Development/AbelianVarietyTheory/Proofs/CubicalLineBundles/CuspFormValueLeSupNormAtOrthonormalFrame.lean}{GlobalTheory.LamT.OrthoBound.cabs_evR_le_supC}
\begin{proof}
For a frame $x$ of a component $Z$ at $\sigma$, let $w_x$ be the complex framed point that the uniformization attaches
to $x$; the reading of $f$ at $x$ is $f_{\sigma,Z}(x)=f(w_x)$. Write $a_{\tau,i}=\alpha_{\tau,i-1}$.

Step 1: an orthonormal frame. By the compatibility conditions of the uniformization at $\sigma$, $y=w_x\cdot(T,c)$ for
a frame $x$ of a component $Z$, a family $T$ of invertible upper triangular matrices and $c\in\mathbf C^\times$
(Lemma~\ref{cv-cusp-cm-point-polynomial}). The leading principal minors of the Gram matrices of $x$ are nonzero, the
$i$-th with the sign $(-1)^{\max(0,i-q)}$, so the Gram--Schmidt process gives $b\in B$ with $xb$ orthonormal. The
framed points $w_x$ and $w_{xb}$ are triangularly equivalent, by the compatibility conditions, as $xb\in\Gamma_ZxB$,
and changes of frame compose. Replacing $x$ by $xb$, we may assume that $x$ is orthonormal.

Step 2: the change of frame. By Lemma~\ref{cv-cusp-cm-point-cusp-bound-ortho}, $w_x$ is metric with constant $1$ and
Gram blocks $2\pi\xi$ times the identity. By Lemma~\ref{cv-cusp-cm-point-metric-change}, $y=w_x\cdot(T,c)$ is metric with
constant $|c|^2$ and Gram blocks $2\pi\xi\,T_\tau^*T_\tau$. By hypothesis and
Lemma~\ref{cv-cusp-cm-point-cusp-bound-unique}, $|c|^2=N$ and $2\pi\xi\,T_\tau^*T_\tau=1$ for every $\tau$. Then
$T_\tau^*=(2\pi\xi)^{-1}T_\tau^{-1}$ is both lower and upper triangular, so $T_\tau$ is diagonal and
$|T_{\tau,ii}|^2=(2\pi\xi)^{-1}$.

Step 3: the value. By the transformation law of $f$ (Lemma~\ref{cv-cusp-cm-point-polynomial}),
$f(y)=c^n\prod_{\tau,i}T_{\tau,ii}^{a_{\tau,i}}f(w_x)$. By Step 2, as $\sum_{\tau,i}a_{\tau,i}=|\alpha|$,
\[
 |f(y)|=N^{n/2}(2\pi\xi)^{-|\alpha|/2}\,|f_{\sigma,Z}(x)|.
\]

Step 4: the supremum. As $x$ is orthonormal, its Gram matrices are diagonal with entries $\pm1$, so
$|\delta_{\tau,i}(x)|=1$ for all $\tau$ and $i$, and the pointwise norm of $f_{\sigma,Z}$ at $x$ is $|f_{\sigma,Z}(x)|$. It
is at most $\sup|f_\sigma|$. As $\sigma$ is good, $\sigma^g=\sigma$, and
$|f(y)|\le N^{n/2}(2\pi\xi)^{-|\alpha|/2}\sup|f_\sigma|=N^{n/2}\|f\|_{\sigma,\infty}$.
\end{proof}
\leanref{Development/AbelianVarietyTheory/Proofs/CubicalLineBundles/ExistsOrthonormalFrameBorelChangeComplexFramedPoStatement.lean}{GlobalTheory.T.exists_orthonormalFrame_borelChange_complexFramedPoint}
\leanref{Development/AbelianVarietyTheory/Proofs/CubicalLineBundles/ExistsOrthonormalFrameBorelChangeComplexFramedPo.lean}{GlobalTheory.LamT.exists_ortho_frBor}
% Lean: cabs_evR_le_supC. Step 1 is T.exists_orthonormalFrame_borelChange_complexFramedPoint = LamT.exists_ortho_frBor
% (the field of IsUnf that makes every complex framed point triangularly equivalent to some un k x, the orthonormal
% translate FD.orthonormalFlagTranslate_holds, the fibre field of IsUnf for x and xb, and frBor_calculus). Step 2 is
% T.unif_metric_scalar (Lemma cv-cusp-cm-point-cusp-bound-ortho), T.metric_change_of_frame with the transposes of g,
% T.metric_data_unique (Lemma cv-cusp-cm-point-cusp-bound-unique) and norm_diag_of_triangular. Step 3 is the weight
% law hwt of Frames.algebraic_laws, with Comparison₂.cabs_wchar and hlie for (2 pi xi)^(-|alpha|/2) = lieFac. Step 4
% is pnorm_unit with dmin_ortho, and Pj.abelian₂.pnorm_le_supN; supC = lieFac * supN. Size: about 13,800 lines,
% most of them in Lemmas cv-cusp-cm-point-cusp-bound-ortho and -unique.

\begin{lemma}\label{cv-cusp-cm-point-cusp-bound-conjugate}
In Situation~\ref{cv-cusp-setting}, let $\sigma\colon F_c\to\mathbf C$ be an embedding and $y$ a framed point of $X$
over $\mathbf C$ through $\sigma$ that is metric with constant $N$ and Gram blocks $(\Gamma_\tau)_\tau$
(Definition~\ref{cv-cusp-cm-point-metric}). Complex conjugation is a homomorphism of $F_c$-algebras from $\mathbf C$
through $\sigma$ to $\mathbf C$ through $\bar\sigma$; let $\bar y$ be the base change of $y$ along it. Then $\bar y$ is
metric with constant $N$ and Gram blocks $(\bar\Gamma_\tau)_\tau$, the complex conjugate matrices.
\end{lemma}
\leanref{Development/AbelianVarietyTheory/Proofs/CubicalLineBundles/IsMetricConjugateFramedPointStatement.lean}{GlobalTheory.T.isMetric_conjugate_framedPoint}
\leanref{Development/AbelianVarietyTheory/Proofs/CubicalLineBundles/IsMetricConjugateFramedPoint.lean}{GlobalTheory.LamT.isMetric_conj}
\begin{proof}
Write $c$ for complex conjugation, on $\mathbf C$ and on $\mathbf C^g$. As in the proof of
Lemma~\ref{cv-cusp-cm-point-cusp-bound-unique}, $N$ and $\Gamma$ are computed from a key datum $(A,\theta,v)$ over
$\mathbf C$ with an isomorphism $m$ of $A$ onto the abelian variety of $y$ respecting the polarizations and the
rigidifications and carrying $\theta$ to the section of $y$; a uniformization $u$ of $A$ with Appell--Humbert data
$(H,\chi)$, whose trivializing functions $\theta_\alpha$ take the values of the rigidification at $0$; a function
$\vartheta$ representing $\theta$; a fundamental domain $F$ of the kernel $\Lambda$ of $u$; and tangent vectors
$w_{\tau,i}$ of $A$ with $m\circ w_{\tau,i}=u_{\tau,i}$ and differentials $W_{\tau,i}$, with
$(\Gamma_\tau)_{ik}=\pi W_{\tau,k}^{\mathsf T}H\overline{W_{\tau,i}}$.

Step 1: a key datum at $\bar y$. As in Step 1 of the proof of Lemma~\ref{cv-cusp-cm-point-cusp-bound-ortho}, the
abelian variety $A'$ of $\bar y$, with the symmetric line bundle of its principal polarization rigidified along the zero
section, the section $\theta'$ of $\bar y$, which generates its space of sections and so is nonzero, and a basis of its
Lie algebra, is a key datum; this uses $p+q\ge1$. The framed point $\bar y$ is the base change of $y$ along $c$, so $m$
and the identification of $A'$ with the fibre at the point of $\bar y$ give a morphism $\varphi\colon A'\to A$ over
$\operatorname{Spec}c$ that identifies $A'$ with the base change of $A$ along $c$, compatibly with the line bundles,
the rigidifications and the sections $\theta'$ and $\theta$. Replace the basis of the Lie algebra of $A'$ by the base
change $v'$ of $v$. A tangent vector of $A'$ is determined by its image in the fibre at the point of $\bar y$, so the
tangent vector $w'_{\tau,i}$ of $A'$ over $u_{\tau,i}$ of $\bar y$, the base change of $u_{\tau,i}$, is the base change of
$w_{\tau,i}$.

Step 2: the conjugate uniformization. A point of $A'$ with values in $\mathbf C$ is the base change along $c$ of a point
of $A$, through $\varphi$; let $u'(z)$ be the base change of $u(\bar z)$. Then $u'$ is a uniformization of $A'$ for its
line bundle, with kernel $\bar\Lambda=c(\Lambda)$ and trivializing functions $z\mapsto\overline{\theta_\alpha(\bar z)}$,
carried to the trivializations of the line bundle of $A'$ by the isomorphism of Step 1: the regular functions on $A'$
are the base changes of those on $A$, whose values at $u'(z)$ are the conjugates of their values at $u(\bar z)$. It
has the Appell--Humbert data $(\bar H,\chi')$, $\chi'(\bar l)=\overline{\chi(l)}$, as conjugating the functional equation
of the $\theta_\alpha$ at $\bar z$ and $l$ gives it for $z$ and $\bar l$ with $\bar H$. Its trivializing functions take
the values of the rigidification of $A'$ at $0$, and $\vartheta'(z)=\overline{\vartheta(\bar z)}$ represents $\theta'$,
as $\varphi$ carries $\theta'$ to $\theta$. The set $\bar F=c(F)$ is a fundamental domain of $\bar\Lambda$. As $c$
preserves the Lebesgue measure and $\bar z^{\mathsf T}\bar H z=\overline{z^{\mathsf T}H\bar z}=z^{\mathsf T}H\bar z$,
the mean of $|\vartheta'|^2e^{-\pi\bar H}$ over $\bar F$ is $N$. Finally, if a tangent vector $w$ of $A$ has differential
$W$ for $u$, its base change has differential $\bar W$ for $u'$.

Step 3: the metric. By Steps 1 and 2, the data $(A',\theta',v')$, $u'$, $(\bar H,\chi')$, $\vartheta'$, $\bar F$ and the
tangent vectors $w'_{\tau,i}$ with differentials $\overline{W_{\tau,i}}$ compute for $\bar y$ the constant $N$ and the
Gram blocks with entries
$\pi\,\overline{W_{\tau,k}}^{\mathsf T}\bar H\,W_{\tau,i}=\overline{\pi W_{\tau,k}^{\mathsf T}H\overline{W_{\tau,i}}}
=\overline{(\Gamma_\tau)_{ik}}$.
\end{proof}
\leanref{Development/AbelianVarietyTheory/Proofs/CubicalLineBundles/ExistsUniformisationOfComplexConjugateKeyDatumStatement.lean}{GlobalTheory.T.exists_uniformisation_of_complexConjugate_keyDatum}
\leanref{Development/AbelianVarietyTheory/Proofs/CubicalLineBundles/ExistsUniformisationOfComplexConjugateKeyDatum.lean}{GlobalTheory.ConjUnif.exists_conj_unif}
% Lean: isMetric_conj unpacks FrMap (conjAlg sigma) y y' (hx, the tangent vectors hu, the sections) and applies
% isMetricAt_conj. Step 1 is T.universalTripleFibreKeyData_nonempty_field (= LamT.nonempty_fibKey) for Phi0 at x',
% Frames.secNZ_of_secGen and T.rigidifiedSectionCorrespondence_existsUnique_nonvanishing (= SecTieEx.secTie_exists_unique)
% for th, then T.universalTripleFibreKeyData_baseChange_keyPullback (= fibKey_keyPB) for starRingEnd C, which gives
% KeyPB of the two key data along conjugation; the tangent vectors w' from Frames.lieBasis_iff_tanBasis, with
% Chart.comb_push for the combinations cf and isLieComb_of_two_pullbacks. Step 2 is
% T.exists_uniformisation_of_complexConjugate_keyDatum = ConjUnif.exists_conj_unif (conjPull, carry with exists_inv,
% conjPull_isAH, val_rig and val_via for IsRig and IsFn, cj_domain, measurePreserving_star with hermSq_star, and the
% differentials). Step 3 is isMetricAt_of_dim with hermForm_map_conj. Size: about 8,000 lines below isMetric_conj, of
% which 5,606 are below fibKey_keyPB and 4,772 below nonempty_fibKey (tools/lean_size.py; they overlap; both are general
% facts on pulling back the universal triple, here to a point and along conjugation), 810 below exists_conj_unif and 510
% below secTie_exists_unique. The proofs of nonempty_fibKey and fibKey_keyPB have been read, and Step 1 follows them:
% the key datum on the fibre as in Step 1 of Lemma cv-cusp-cm-point-cusp-bound-ortho; then, for phi = starRingEnd C,
% sabMap_cancel_and_comp_over gives the morphism phi : A' -> A over Spec c between the two pullbacks of the universal
% triple (x' = Spec c >> x), pat_pullback the identification of the line bundles with their rigidifications,
% ArcQBC.exists_pullSec and TieCalc.secTie_calc the sections (the section of y-bar is the base change of that of y,
% SecBC), and exists_isLieBasis_baseChange the basis v'. The text's Step 1 is a correct proof of what is used.

\begin{lemma}\label{cv-cusp-cm-point-cusp-bound}
In Situation~\ref{cv-cusp-setting}, let $\kappa=(\alpha,n)$ be a dominant weight and $f\in S_\kappa$, let
$\sigma\colon F_c\to\mathbf C$ be an embedding, and let $y$ be a framed point of $X$ over $\mathbf C$ through
$\sigma$ that is metric with constant $N$ and with the identity matrices as Gram blocks
(Definition~\ref{cv-cusp-cm-point-metric}). Then $|f(y)|\le N^{n/2}\|f\|_{\sigma,\infty}$.
\end{lemma}
\leanref{Development/AbelianVarietyTheory/Proofs/CubicalLineBundles/CuspFormBoundedBySupNormAtGoodEmbeddingAllEmbeddStatement.lean}{GlobalTheory.T.cuspForm_bounded_by_supNorm_at_good_embedding_all_embeddings}
\leanref{Development/AbelianVarietyTheory/Proofs/CubicalLineBundles/CuspFormBoundedBySupNormAtGoodEmbeddingAllEmbedd.lean}{GlobalTheory.LamT.cusp_bound_allEmb}
\begin{proof}
If $\sigma$ is good, this is Lemma~\ref{cv-cusp-cm-point-cusp-bound-good}. Otherwise $\sigma$ sends the square root
$\rho$ of $-d$ in $F_c$ given by the datum to $-\rho_0$, where $\rho_0$ is the root fixed by $\mathcal D$. As $\rho_0$ is
purely imaginary, $\bar\sigma(\rho)=-\bar\rho_0=\rho_0$, so $\bar\sigma$ is good, and $\sigma^g=\bar\sigma=\bar\sigma^g$,
which gives
$\|f\|_{\sigma,\infty}=\|f\|_{\bar\sigma,\infty}$. Let $\bar y$ be the base change of $y$ along complex conjugation, as in
Lemma~\ref{cv-cusp-cm-point-cusp-bound-conjugate}. By that lemma, $\bar y$ is metric with constant $N$ and Gram blocks
the conjugates of the identity matrices, which are the identity matrices. Lemma~\ref{cv-cusp-cm-point-cusp-bound-good},
for $\bar\sigma$ and $\bar y$, gives $|f(\bar y)|\le N^{n/2}\|f\|_{\bar\sigma,\infty}$. The values of $f$ commute with base
change (Lemma~\ref{cv-cusp-cm-point-polynomial}), so $f(\bar y)=\overline{f(y)}$, and
$|f(y)|=|f(\bar y)|\le N^{n/2}\|f\|_{\sigma,\infty}$.
\end{proof}
% Lean: cusp_bound_allEmb: by_cases on GoodEmb c sigma. At a good sigma, goodOf c sigma = sigma and
% T.cuspForm_value_le_supNorm_at_orthonormal_frame (Lemma cv-cusp-cm-point-cusp-bound-good). Otherwise goodOf c sigma
% is the conjugate (goodEmb_goodOf), y' is FramedBaseChange.exists_frMap along FrameNorm.conjAlg sigma,
% T.isMetric_conjugate_framedPoint (Lemma cv-cusp-cm-point-cusp-bound-conjugate) with Matrix.map_one, and evR_frMap
% (with Frames.algebraic_laws, cusp_le_sec and allTest_of_secBaseChange) and Complex.norm_conj for the value. Size:
% about 14,700 lines below cusp_bound_allEmb beyond the laws of the moduli problem and the uniformization of abelian
% varieties, by following its imports and stopping at what the book cites; the earlier figure of 32,000 lines counted
% more. Most of it is in Lemmas cv-cusp-cm-point-cusp-bound-ortho (about 5,600 lines), -unique (about 5,400) and
% -conjugate (about 8,000), which overlap and are now written out.

\begin{lemma}\label{cv-cusp-cm-point-arch}
In Situation~\ref{cv-cusp-setting}, let $\kappa=(\alpha,n)$ be a dominant weight and $f\in S_\kappa$, let
$K\supset F_c$ be a number field, let $\ell>0$, let a good CM point of $X$ over $K$ with bound $\ell$ be given, with
the notation of Definition~\ref{cv-cusp-cm-point-good}, and let $G$ be a polynomial as in
Lemma~\ref{cv-cusp-cm-point-polynomial}(1) for $f$ and its framed point $y$. Let $s\colon K\to\mathbf C$ be an
embedding. Then condition (A) of Proposition~\ref{cv-cusp-cm-product} holds at $s$ for the metrized lines $\Theta$,
$\Delta$ and $\mathcal L_1,\dots,\mathcal L_g$, the matrix $U$, the bijection $e$, the polynomial $G$, the set $B$ of
embeddings $\tau\colon E\to\mathbf C$ with the integers $m_\tau$, the exponents $a_{\tau,i}=\alpha_{\tau,i-1}$, the
integer $n$, the number $\ell$ and $S_s=\|f\|_{s|_{F_c},\infty}$.
\end{lemma}
\leanref{Development/AbelianVarietyTheory/Proofs/CubicalLineBundles/ArchimedeanHypothesisLowerBoundCmPointStatement.lean}{GlobalTheory.T.archimedeanHypothesis_lowerBound_cmPoint}
\leanref{Development/AbelianVarietyTheory/Proofs/CubicalLineBundles/ArchimedeanHypothesisLowerBoundCmPoint.lean}{GlobalTheory.LamT.step5Arch_cm}
\begin{proof}
Let $y_s$ be the base change of $y$ along $s$, and let $H_s$ be the matrix of Lemma~\ref{cv-cusp-cm-point-hodge}. It
has the properties required in (A) apart from those of $s(G)$, and $y_s$ is metric with constant $n_s(\Theta)^2$
and Gram blocks $\Gamma_{s,\tau}$, the matrices of (A). By Lemma~\ref{cv-cusp-cm-point-polynomial}(2), $s(G)$ has
the transformation law with exponents $a$.

Let $Y=(Y_\tau)$ be a family with $Y_\tau^*\Gamma_{s,\tau}Y_\tau=1$ for every $\tau$. Taking determinants, every
$Y_\tau$ is invertible. By Lemma~\ref{cv-cusp-cm-point-metric-change} with $a=1$, $y_s\cdot Y$ is metric with
constant $n_s(\Theta)^2$ and with the identity matrices as Gram blocks, so Lemma~\ref{cv-cusp-cm-point-cusp-bound}
gives $|f(y_s\cdot Y)|\le n_s(\Theta)^nS_s$. By Lemma~\ref{cv-cusp-cm-point-polynomial}(1), with $S=\mathbf C$ and
$\varphi=s$, $f(y_s\cdot Y)=s(G)(Y)$. As $n_s(\Theta^\vee)=n_s(\Theta)^{-1}$, this gives
$|s(G)(Y)|\,n_s(\Theta^\vee)^n\le S_s$.
\end{proof}
% Lean: step5Arch_cm: exists_metricData_cmPoint_complexEmbedding (Lemma cv-cusp-cm-point-hodge) for the base change
% y_sigma (FramedBaseChange.exists_frMap), framePolynomial_equivariant_and_ne_zero.1 for the law, existsUnique_frChg
% for y' = y_sigma changed by (Y^T, 1), metric_change_of_frame, cuspForm_bounded_by_supNorm_at_good_embedding_all_embeddings
% with N = n_s(Theta)^2, and hG (IsFrPoly) for f(y') = s(G)(Y).

\subsection{Serre tensor products over the complex numbers}\label{subsection-cv-serre-tensor-complex}

\begin{definition}\label{cv-cusp-cm-point-eigen-degree-line}
Let $K$ be a number field and $A$ a polarized abelian scheme over $\mathcal O_K$, and write $\operatorname{Lie}A$ for
its Lie algebra, an $\mathcal O_K$-lattice in the Lie algebra of the generic fibre $A_K$. For a nonzero vector $v$ of
the Lie algebra of $A_K$, the \emph{metrized line of $v$} is the metrized line $\mathcal M(A,v)$ over $K$
(Definition~\ref{cv-cusp-cm-lines}) whose ideal is the set of $c\in K$ with $cv\in\operatorname{Lie}A$, and for which
$n_s(\mathcal M(A,v))^2$ is the squared Hodge norm (Definition~\ref{cv-cm-point-periods}(4)) of the image of $v$ in the
base change of $A_K$ along $s$, for every embedding $s\colon K\to\mathbf C$. If a ring $O$ acts on $A$ and
$\tau\colon O\to K$ is a ring homomorphism, the \emph{$\tau$-eigenspace} of the Lie algebra of $A_K$ is the subspace
on which every $b\in O$ acts by $\tau(b)$.
\end{definition}
\leanref{Development/AbelianVarietyTheory/Definitions/CMLattices/IsDiffOf.lean}{GlobalTheory.CM.GenFib.lieLat}
\leanref{Development/AbelianVarietyTheory/Definitions/CMLattices/IsDiffOf.lean}{GlobalTheory.CM.GenFib.eigLat}
\leanref{Development/AbelianVarietyTheory/Definitions/CMLattices/IsDiffOf.lean}{GlobalTheory.CM.GenFib.IsHerm}

\begin{lemma}\label{cv-cusp-cm-point-eigen-degree-curve}
Let $d\equiv3\pmod4$ be a prime, let $\theta\in\mathbf C$ with $\theta^2=-d$, let $M\subset k\subset\mathbf C$ be
number fields, and let $C$ be a CM curve with root $\theta$ over $M$, with basis $v$ of its Lie algebra and squared
Hodge norms $n(s)$ (Definition~\ref{cv-cm-point-curve}). Let $C_k$ be the base change of $C$ to $\mathcal O_k$, with
the base changes of its polarization and of its action of $\mathcal O=\mathbf Z[(1+\sqrt{-d})/2]$, and let $v_k$ be
the image of $v$. Then $C_k$ with $v_k$ is a CM curve with root $\theta$ over $k$, the squared Hodge norm of $v_k$ at
an embedding $s\colon k\to\mathbf C$ is $n(s|_M)$, and
\[
 \frac1{[k:\mathbf Q]}\sum_{s\colon k\to\mathbf C}\log n(s|_M)=\frac1{[M:\mathbf Q]}\sum_{s\colon M\to\mathbf C}\log n(s).
\]
\end{lemma}
\leanref{Development/AbelianVarietyTheory/Proofs/CubicalLineBundles/EllPullExistsOAct.lean}{GlobalTheory.CM.exists_ellMap}
\leanref{Development/AbelianVarietyTheory/Proofs/CubicalLineBundles/CmCurveBaseChangeComp.lean}{GlobalTheory.CM.EllComp.ellMap_comp}
\leanref{Development/AbelianVarietyTheory/Proofs/CubicalLineBundles/HasPeriodOfEllMap.lean}{GlobalTheory.CM.hasPeriod_of_ellMap}
\begin{proof}
The base change of the action of $\mathcal O$ is an action on $C_k$, and the Rosati involution of the base changed
polarization is complex conjugation, since these are identities between homomorphisms and line bundles, which base
change preserves. The base change of a basis of the Lie algebra of an abelian scheme is a basis of the Lie algebra of
its base change, and $\sqrt{-d}$ acts on $v_k$ by $\theta$, as it does on $v$. For an embedding $s\colon k\to\mathbf C$
the complex curve of $C_k$ at $s$, with the image of $v_k$, is the complex curve of $C$ at $s|_M$ with the image of
$v$. So the squared Hodge norm of $v_k$ at $s$ is $n(s|_M)$, and at the inclusion $k\subset\mathbf C$ the
uniformization and the number $z$ of $C$ serve for $C_k$. Finally every embedding of $M$ in $\mathbf C$ has $[k:M]$
extensions to $k$, and $[k:\mathbf Q]=[k:M]\,[M:\mathbf Q]$.
\end{proof}
% Lean: exists_ellMap (exists_isPolFibre for the base change, ellPull_exists_oAct for the action,
% exists_isLieBasis_baseChange and lieComb_baseChange for v_k, lamEq_baseChange for the Rosati condition, ellPull_prim
% for the normalization prim of CM.Ell); ellMap_comp (composition of base changes, IsPrincipalAt, and isLieNormSq_iff
% for the Hodge norms); hasPeriod_of_ellMap, with n'(s) = n(s.comp j) and meanLog_comp for the means.

\begin{lemma}\label{cv-cusp-cm-point-eigen-degree-base}
In Situation~\ref{cv-cusp-setting}, let $j\colon K\to K'$ be an embedding of number fields, let $A$ be a principally
polarized abelian scheme over $\mathcal O_K$ with an action of $\mathcal O_E$ for which the Rosati involution of the
polarization induces complex conjugation on $\mathcal O_E$, and let $A'$ be its base change to $\mathcal O_{K'}$
along $j$, with the base changes of the action and of the polarization. Let $\tau\colon E\to K$ be an embedding such
that the $\tau$-eigenspace of the Lie algebra of $A_K$ (Definition~\ref{cv-cusp-cm-point-eigen-degree-line}) is a
line, spanned by $v$, and let $v'$ be the image of $v$ in the Lie algebra of $A'_{K'}$. Then the $j\tau$-eigenspace
of the Lie algebra of $A'_{K'}$ is the line spanned by $v'$, and
$\widehat\deg\mathcal M(A',v')=\widehat\deg\mathcal M(A,v)$.
\end{lemma}
\leanref{Development/AbelianVarietyTheory/Proofs/CubicalLineBundles/DegreeLieLatticesCmTripleBaseChangeStatement.lean}{GlobalTheory.T.degree_lieLattices_cmTriple_baseChange}
\leanref{Development/AbelianVarietyTheory/Proofs/CubicalLineBundles/DegreeLieLatticesCmTripleBaseChange.lean}{GlobalTheory.CM.LieBaseChange.deg_lieLat_eigLat_baseChange}
\begin{proof}
Identify the Lie algebra of $A'_{K'}$ with $K'\otimes_K$ the Lie algebra of $A_K$, through $j$.

Step 1: the Lie algebras. The image of $\operatorname{Lie}A$ lies in $\operatorname{Lie}A'$. Conversely, there are
$r_1,\dots,r_n\in\mathcal O_K\setminus\{0\}$ that generate the unit ideal such that, for every $i$, the Lie algebra
of the base change of $A$ to $\mathcal O_K[1/r_i]$ has a basis, whose vectors lie in $r_i^{-l}\operatorname{Lie}A$
for some $l$. The base change of such a basis to $\mathcal O_{K'}[1/r_i]$ is a basis of the Lie algebra of the base
change of $A'$. So for $c'\in\operatorname{Lie}A'$ and every $i$, some power of $r_i$ times $c'$ lies in the
$\mathcal O_{K'}$-span of the image of $\operatorname{Lie}A$. As the $r_i$ generate the unit ideal of
$\mathcal O_{K'}$, so does $c'$. Hence $\operatorname{Lie}A'$ is the $\mathcal O_{K'}$-span of the image of
$\operatorname{Lie}A$.

Step 2: the eigenvectors. The eigenspace condition is a family of $K$-linear equations, so the $j\tau$-eigenspace of
the Lie algebra of $A'_{K'}$ is $K'\otimes_K$ the $\tau$-eigenspace, the line spanned by $v'$. Let $I$ be the ideal of
$\mathcal M(A,v)$, so that the $\tau$-eigenvectors in $\operatorname{Lie}A$ are the $cv$ with $c\in I$. The
$\mathcal O_{K'}$-span of their images is $I\mathcal O_{K'}v'$, and it consists of $j\tau$-eigenvectors in
$\operatorname{Lie}A'$ by Step~1. Conversely, the ring $\mathcal O_{K'}$ is a finitely generated projective
$\mathcal O_K$-module, so there are $b_1,\dots,b_m\in\mathcal O_{K'}$ and $\mathcal O_K$-linear maps
$\varphi_k\colon\mathcal O_{K'}\to\mathcal O_K$ with $y=\sum_k\varphi_k(y)b_k$ for every $y$; each $\varphi_k$
extends to a $K$-linear map $K'\to K$. Let $c'\in\operatorname{Lie}A'$ be a $j\tau$-eigenvector, and by Step~1 write
$c'=\sum_is_ij(l_i)$ with $s_i\in\mathcal O_{K'}$ and $l_i\in\operatorname{Lie}A$. Then $c'=\sum_kb_kj(l'_k)$ with
$l'_k=\sum_i\varphi_k(s_i)l_i\in\operatorname{Lie}A$. For $b\in\mathcal O_E$ the vector
$\sum_is_ij\big(bl_i-\tau(b)l_i\big)=bc'-j\tau(b)c'$ is $0$, and applying $\varphi_k$ to its coordinates gives
$bl'_k=\tau(b)l'_k$. So every $l'_k$ is a $\tau$-eigenvector in $\operatorname{Lie}A$, and $c'\in I\mathcal O_{K'}v'$.
Hence the ideal of $\mathcal M(A',v')$ is $I\mathcal O_{K'}$.

Step 3: the norms. For an embedding $s'\colon K'\to\mathbf C$, the base change of $A'_{K'}$ along $s'$, with the
image of $v'$, is the base change of $A_K$ along $s'j$ with the image of $v$. So
$n_{s'}(\mathcal M(A',v'))=n_{s'j}(\mathcal M(A,v))$.

Step 4: the degrees. The norm of $I\mathcal O_{K'}$ is $\mathrm N(I)^{[K':K]}$, and every embedding of $K$ in
$\mathbf C$ has $[K':K]$ extensions to $K'$. So $\deg\mathcal M(A',v')=[K':K]\deg\mathcal M(A,v)$, and the normalized
degrees agree, as $[K':\mathbf Q]=[K':K]\,[K:\mathbf Q]$.
\end{proof}
\leanref{Development/AbelianVarietyTheory/Proofs/CubicalLineBundles/DegreeLieLatticesCmTripleBaseChange.lean}{GlobalTheory.CM.LieBaseChange.deg_eq_of_isPolFibre}
\leanref{Development/AbelianVarietyTheory/Proofs/CubicalLineBundles/LieLatticeLocalBasisDescriptionStatement.lean}{GlobalTheory.T.lie_lattice_local_basis_description}
\leanref{Development/AbelianVarietyTheory/Proofs/CubicalLineBundles/EigenLatticeNumberFieldBaseChangeStatement.lean}{GlobalTheory.T.eigenLattice_numberField_baseChange}
\leanref{Development/AbelianVarietyTheory/Proofs/CubicalLineBundles/DegLatticeNumberFieldExtensionEqStatement.lean}{GlobalTheory.T.deg_lattice_numberField_extension_eq}
% Lean: deg_lieLat_eigLat_baseChange = deg_eq_of_isPolFibre, which works with the Lie lattice and every eigen-lattice
% (the text states the case of an eigenline). Step 1 is LieBaseChange.lieLat_eq_span: image_subset_lieLat,
% lie_lattice_local_basis_description (= LieGlue.exists_cover_lieLat: the r_i and the matrices M_i of local bases),
% exists_isLieBasis_baseChange, lieLattice_localization_modelLieAlgebra on both sides, exists_pow_smul_mem and
% Submodule.mem_of_span_eq_top_of_smul_pow_mem. Step 2 is LieBaseChange.eigLat_eq_span with
% eigenLattice_numberField_baseChange (= CM.eigenLattice_baseChange: the frame b_k, phi_k is eigenBC_frame, the
% relation is eigenBC_relation) and exists_lieRepresentation_eigenLattices for the matrices of the action. Step 3 is
% form_single_eq (LieMetric.hermForm_eq). Step 4 is deg_lattice_numberField_extension_eq (= deg_span_image_eq);
% degree_lie_lattice_hermitian_invariant (= FibIndep.deg_lieLat_eq) removes the choice of a basis of the generic
% fibre, and LieMetric.exists_unique_isHerm gives the Hodge metric.

\begin{definition}\label{cv-cusp-cm-point-eigen-degree-serre}
In Situation~\ref{cv-cusp-setting}, let $\sqrt{-d}\in E$ be the square root of $-d$ of the datum, and put
$w_0=(1+\sqrt{-d})/2$ and $\mathcal O=\mathbf Z[w_0]\subset\mathcal O_E$. Let $R$ be a ring, let $C$ be a
principally polarized elliptic curve over $R$ with an action of $\mathcal O$, let $L$ be an $\mathcal O_E$-module and
$h\colon L\times L\to\mathcal O$ a map. A principally polarized abelian scheme $A$ over $R$ with an action of
$\mathcal O_E$ is a \emph{Serre tensor product} of $C$ and $L$ with form $h$ if there are isomorphisms of
$\mathcal O$-modules $L\otimes_{\mathcal O}C(T)\to A(T)$, natural in the $R$-scheme $T$, where $\mathcal O$ acts on
$A(T)$ through $\mathcal O\subset\mathcal O_E$, such that the homomorphisms $t_a\colon C\to A$, $x\mapsto a\otimes x$,
for $a\in L$, satisfy:
\begin{enumerate}
\item $t_{ba}=b\circ t_a$ for every $b\in\mathcal O_E$;
\item the polarization of $A$ pulls back along the pair $(t_a,t_{a'})$ to that of $C$ composed with the action of
  $h(a,a')$, for all $a,a'\in L$.
\end{enumerate}
\end{definition}
\leanref{Development/AbelianVarietyTheory/Definitions/CMLattices/IsDiffOf.lean}{GlobalTheory.CM.IsTen}
\leanref{Development/AbelianVarietyTheory/Definitions/ShimuraVarieties/PowObjGeometricallyConnected.lean}{GlobalTheory.Serre.IsTensor}
\leanref{Development/AbelianVarietyTheory/Definitions/CMLattices/IsDiffOf.lean}{GlobalTheory.CM.IsKForm}

\begin{definition}\label{cv-cusp-cm-point-eigen-degree-tensor-family}
In Situation~\ref{cv-cusp-setting}, let $\sqrt{-d}\in E$ be the square root of $-d$ of the datum, and put
$w_0=(1+\sqrt{-d})/2$ and $\mathcal O=\mathbf Z[w_0]\subset\mathcal O_E$. Let $L$ be an $\mathcal O_E$-module, $M\subset L$
an $\mathcal O_E$-submodule and $h\colon L\times L\to\mathcal O$ a map. Let $R$ be a ring, $C$ a principally polarized
elliptic curve over $R$ with an action of $\mathcal O$, and $A$ a principally polarized abelian scheme over $R$ with an
action of $\mathcal O_E$, and let $\mathcal L_C$ and $\mathcal L_A$ be the line bundles of their polarizations. A
\emph{family of homomorphisms} from $C$ to $A$ for $h$ on $M$ is a family $(t_a)_{a\in M}$ of homomorphisms
$t_a\colon C\to A$ such that:
\begin{enumerate}
\item each $t_a$ commutes with the actions of $\mathcal O$, which acts on $A$ through $\mathcal O\subset\mathcal O_E$;
\item $t_{a+a'}=t_a+t_{a'}$ and $t_{ba}=b\circ t_a$ for all $a,a'\in M$ and $b\in\mathcal O_E$;
\item $\Lambda_{\mathcal L_A}(t_ap_1,t_{a'}p_2)\cong\Lambda_{\mathcal L_C}(h(a,a')p_1,p_2)$ on $C\times_RC$ for all
  $a,a'\in M$, with $p_1$ and $p_2$ the projections (Definition~\ref{cv-toroidal-abelian-dvr-ample-torsion-weil}).
\end{enumerate}
For a ring homomorphism $\sigma\colon R\to\mathbf C$, the family is \emph{onto at $\sigma$} if every point of $A$ over
$\mathbf C$ through $\sigma$ is a finite sum $\sum_mt_{a_m}(x_m)$, with $a_m\in M$ and points $x_m$ of $C$ over
$\mathbf C$ through $\sigma$.
\end{definition}
\leanref{Development/AbelianVarietyTheory/Definitions/CMLattices/ExistsIsThetaDegOfIsLocalModel.lean}{GlobalTheory.CM.IsFam}
\leanref{Development/AbelianVarietyTheory/Definitions/CMLattices/ExistsIsThetaDegOfIsLocalModel.lean}{GlobalTheory.CM.OntoAt}

\begin{lemma}\label{cv-cusp-cm-point-eigen-degree-tensor-complex-positive}
In Situation~\ref{cv-cusp-setting}, let $\sqrt{-d}\in E$ be the square root of $-d$ of the datum, and put
$w_0=(1+\sqrt{-d})/2$. Let $L$ be the lattice of a PEL datum over $E$ with its perfect alternating form
$\langle\ ,\ \rangle$, let $\sigma_1,\dots,\sigma_g$ be embeddings $E\to\mathbf C$, and let $P\colon L\to\mathbf C^g$ be
a period of $L$ for them (Definition~\ref{cv-cm-point-periods}(3)). Let $\omega\in\mathbf C$, and let $x\in L$ be
nonzero with $P(x)_k=0$ for every $k$ with $\sigma_k(\sqrt{-d})\neq2\omega-1$. Then
$\langle w_0x,x\rangle=\operatorname{Im}(\omega)\,r$ for a real number $r>0$.
\end{lemma}
\leanref{Development/AbelianVarietyTheory/Proofs/CubicalLineBundles/ExistsCmCurveHomFamilyToFibreAtPeriod.lean}{GlobalTheory.CM.exists_hermitianPairing_eq_im_mul_pureType}
\begin{proof}
Let $H$ be a positive definite Hermitian form on $\mathbf C^g$ with $\operatorname{Im}H(P(y),P(y'))=\langle y,y'\rangle$
for all $y,y'\in L$. As $2w_0=1+\sqrt{-d}$, $P(2w_0x)_k=(1+\sigma_k(\sqrt{-d}))P(x)_k$, which is $2\omega P(x)_k$ for every
$k$, since $\sigma_k(\sqrt{-d})=2\omega-1$ wherever $P(x)_k\neq0$. So $P(w_0x)=\omega P(x)$. As $P$ is injective and
$x\neq0$, $P(x)\neq0$, so $r=H(P(x),P(x))$ is a positive real number. Then
$\langle w_0x,x\rangle=\operatorname{Im}H(\omega P(x),P(x))=\operatorname{Im}(\omega r)=\operatorname{Im}(\omega)\,r$.
\end{proof}
% Lean: exists_hermitianPairing_eq_im_mul_pureType: hP.2.2.2.2 for Hm, two_weight_eq_one_add_root and
% PureType.period_rt_smul for hw0, hP.2.1 for P(x) != 0, PureTypeVector.hermForm_smul_smul and hermForm_self_im.

\begin{definition}\label{cv-cusp-cm-point-eigen-degree-tensor-laws}
Let $S$ be a scheme and $A$ a principally polarized abelian scheme over $S$ of relative dimension $g$, with the line
bundle $\mathcal L$ of its polarization. For a scheme $T$ over $S$ let $A_T=T\times_SA$, with projections $p_T$ and $p_A$;
for $x\in A(T)$ let $\tau_x=x\,p_T+p_A\colon A_T\to A$ and let $T_x$ be the translation of $A_T$ by $x$, with
$p_A\circ T_x=\tau_x$; let $[n]$ be the multiplication by $n$; and let $e_k$ be the commutator pairing of the theta
group of $\mathcal L^k$ (Definition~\ref{cv-cusp-cm-point-eigen-degree-rank-one-norm-weil-translate-group}). A line
bundle on $A_T$ is \emph{trivial locally on $T$} if it is trivial over the inverse images of the members of an open
cover of $T$. The \emph{laws of abelian schemes} are the following statements, for all such $S$, $A$ and $\mathcal L$;
in (2) to (8), (10), (12) and (15) the scheme $S$ is locally noetherian.
\begin{enumerate}
\item If $S=\operatorname{Spec}k$ for an algebraically closed field $k$, then $\dim_k\Gamma(A,\mathcal L^n)=n^g$ for
  every $n\ge1$.
\item Let $n\ge1$. Every point of $S$ has an open neighbourhood $V$ with sections $b_1,\dots,b_{n^g}$ of
  $p_A^*\mathcal L^n$ on $A_V$ whose pullbacks, for every ring $R$ and every morphism $\operatorname{Spec}R\to V$, form
  a basis of the $R$-module of sections of $p_A^*\mathcal L^n$ on $A_R$.
\item Let $n\ge4$ and let $U\subset S$ be an affine open subscheme. Then $\mathcal L^n$ has sections $s_0,\dots,s_M$
  over the inverse image $A_U$ of $U$ without common zeros, such that each nonvanishing locus $D(s_i)$ is affine and its
  ring of functions is generated over $\Gamma(U,\mathcal O_S)$ by the ratios $s_j/s_i$.
\item Let $N\ge1$, $T$ a scheme over $S$ and $x\in A(T)$. Then $p_A^*\mathcal L^N\otimes\tau_x^*\mathcal L^{-N}$ is
  trivial locally on $T$ if and only if $Nx=0$.
\item $[n]^*\mathcal L\cong\mathcal L^{n^2}$ for every $n\ge0$.
\item Let $n\ge1$. Then $[n]$ is finite, flat and surjective, and \'etale if $n$ is a unit on $S$; and for every
  algebraically closed field $k$ and every morphism $\operatorname{Spec}k\to S$ with $n\neq0$ in $k$, the points of $A$
  over $k$ through it that are killed by $n$ are $n^{2g}$ in number.
\item For every generator of the restriction of $\mathcal L$ to the zero section, $\mathcal L$ with this
  trivialization has a cubical structure, and any two cubical structures are equal.
\item Let $T$ be a scheme over $S$ and $x,y\in A(T)$. Then
  $\tau_{x+y}^*\mathcal L\otimes p_A^*\mathcal L\otimes\tau_x^*\mathcal L^{-1}\otimes\tau_y^*\mathcal L^{-1}$ is
  trivial locally on $T$.
\item For abelian schemes $A'$ and $B'$ over any scheme $S'$, every morphism $A'\to B'$ over $S'$ that carries the zero
  section to the zero section is a homomorphism.
\item Let $k\ge0$ and let $T$ be a scheme over $S$. (a) For $x,y\in A(T)$ with $kx=ky=0$ there is exactly one unit
  $c$ of $\Gamma(T,\mathcal O_T)$ with $e_k(x,y)=c$. (b) For $x,x',y\in A(T)$, if $e_k(x,y)=c$ and $e_k(x',y)=c'$, then
  $e_k(x+x',y)=cc'$. (c) If $kx=0$, then $e_k(x,x)=1$. (d) If $l\ge0$, $kl$ is a unit of $\Gamma(T,\mathcal O_T)$,
  $klx=0$ and $ky=0$, then $e_{kl}(x,y)=c$ if and only if $e_k(lx,y)=c$.
\item If $S=\operatorname{Spec}k$ for an algebraically closed field $k$ and $n\neq0$ in $k$, then every $x\in A(k)$ with
  $nx=0$ and $e_n(x,y)=1$ for all $y\in A(k)$ with $ny=0$ is $0$.
\item Let $k$ be odd and a unit on $S$, $T$ a scheme over $S$ and $x\in A(T)$ with $kx=0$. Call an element
  $(x,\varphi)$ of the theta group of $\mathcal L^k$ over an open subscheme $V$ of $T$ \emph{symmetric of order $k$}
  if $\prod_{j=0}^{k-1}T_{jx}^*\varphi=1$ and if $T_{-x}^*\varphi\otimes\nu\otimes\iota^*\varphi\otimes
  T_{-x}^*\nu^{-1}=1$ for some generator $\nu$ of $p_A^*\mathcal L^k\otimes(-p_A)^*\mathcal L^{-k}$, where
  $\iota=(p_V,-p_A)$ (here $x$ stands for its restriction to $V$). Then every point of $T$ has an open neighbourhood $V$
  over which $x$ lies under an element of the theta group of $\mathcal L^k$ that is symmetric of order $k$, and any two
  elements of the theta group of $\mathcal L^k$ over $T$ that lie over $x$ and are symmetric of order $k$ are equal.
\item If $S=\operatorname{Spec}K$ for an algebraically closed field $K$, and $k=1$ or $k\ge4$, with $k\neq0$ in $K$,
  then the only subspaces of $\Gamma(A,\mathcal L^k)$ that are stable under the maps $s\mapsto\varphi\cdot T_x^*s$, for
  the elements $(x,\varphi)$ of the theta group of $\mathcal L^k$ over $K$, are $0$ and $\Gamma(A,\mathcal L^k)$.
\item If $S=\operatorname{Spec}K$ for an algebraically closed field $K$, then for every nonzero section $\vartheta$ of
  $\mathcal L$, every $n\ge2$ with $n\neq0$ in $K$ and every $a\in A(K)$ there is $x\in A(K)$ with $nx=0$ such that
  $\vartheta$ does not vanish at $a+x$.
\item Let $M\ge3$, and let $\sigma$ be an automorphism of $A$ that is a homomorphism, such that $\sigma^*\mathcal
  L\otimes\mathcal L^{-1}$ is trivial over the inverse images in $A$ of the members of an open cover of $S$ and
  $\sigma\circ x=x$ for every scheme $T$ over $S$ and every $x\in A(T)$ with $Mx=0$. Then $\sigma$ is the identity.
\item Let $V$ be a discrete valuation ring with fraction field $K$, and $G$ a commutative group scheme over $V$ that is
  smooth, separated and quasi-compact over $V$ with geometrically connected fibres, with generic fibre $G_K$. Then
  every line bundle on $G_K$ that is trivial along the unit section is isomorphic to the restriction of a line bundle on
  $G$ that is trivial along the unit section.
\item Let $V$, $K$, $G$ and $G_K$ be as in (16), and let $\mathcal N$ be a line bundle on $G$ with a generator along the
  unit section. Then every generator of $\Theta(\mathcal N)$ over $G_K\times_KG_K\times_KG_K$ whose value at the origin
  is $1$ is the restriction of exactly one generator of $\Theta(\mathcal N)$ over $G\times_VG\times_VG$.
\item Let $V$ and $G$ be as in (16), let $\pi$ be a uniformizer of $V$, and let $s$ be a section of a line bundle
  $\mathcal N$ on $G$. If, on each member of an open cover of $G$ on which $\mathcal N$ is trivialized, the function
  that represents $s$ is divisible by every power of $\pi$, then $s=0$.
\item Let $N\ge0$, $T$ a scheme over $S$, $z\in A(T)$ and $\epsilon$ a generator of $z^*\mathcal L$. Then for every
  section $s$ of $p_A^*\mathcal L^N$ on $A_T$ there is exactly one $v\in\Gamma(T,\mathcal O_T)$ with
  $(1,z)^*s=v\,\epsilon^N$.
\item Let $m$ be odd and a unit on $S$, and $\zeta$ a unit of $\Gamma(S,\mathcal O_S)$. Suppose that there are points
  $P_1,\dots,P_{2g}\in A(S)$ killed by $m^2$, with $\zeta$ a primitive $m^2$-th root of unity, such that on every
  geometric fibre their images form a basis of the points killed by $m^2$ over $\mathbf Z/m^2$, and
  $e_{m^2}(P_i,P_j)=\zeta^{J_{ij}}$, with $J_{i,i+g}=1$, $J_{i+g,i}=-1$ and $J_{ij}=0$ otherwise. Then every point of every
  scheme $T$ over $S$ has an affine open neighbourhood $V$ such that the following $\Gamma(V,\mathcal O_V)$-module is
  free of rank $m^{2g}$: the module of sections of $p_A^*\mathcal L^{m^2}$ on $A_V$ spanned by the sections that are,
  near every point of $V$, linear combinations of sections $\varphi\cdot T_x^*(\alpha\cdot[m]^*\vartheta)$, with
  $\vartheta$ a section of $p_A^*\mathcal L$, $\alpha$ a generator of $p_A^*\mathcal L^{m^2}\otimes(m\,p_A)^*\mathcal
  L^{-1}$ and $(x,\varphi)$ an element of the theta group of $\mathcal L^{m^2}$.
\item For every abelian scheme $A'$ over any scheme $S'$, the map $\Gamma(S',\mathcal O_{S'})\to\Gamma(A',\mathcal
  O_{A'})$ is bijective.
\item For every morphism $S'\to S$ there are a principally polarized abelian scheme $A'$ over $S'$ of relative
  dimension $g$, with the line bundle $\mathcal L'$ of its polarization, and a morphism $A'\to A$ that makes $A'$ the
  base change of $A$ along $S'\to S$, respects the group laws, and pulls $\mathcal L$ back to a line bundle isomorphic
  to $\mathcal L'$.
\item If $S=\operatorname{Spec}R$ for a noetherian ring $R$, then for every group scheme $H$ over $S$ that is finite and
  flat over $S$ and every homomorphism $i\colon H\to A$ that is a monomorphism, there are an abelian scheme $A'$ over
  $S$ and a finite, flat and surjective homomorphism $q\colon A\to A'$ such that $i$ identifies $H$ with the fibre
  product of $q$ and the zero section of $A'$.
\item For every complete discrete valuation ring $V$ and every finite separable extension $k'$ of its residue field
  there are a complete discrete valuation ring $W$ that is a finite $V$-algebra, with $V\to W$ injective, local and
  formally unramified, and an isomorphism of the residue field of $W$ with $k'$ over the residue field of $V$.
\end{enumerate}
\end{definition}
\leanref{ReciprocityTheory/ReductiveGroups/ArithmeticGeometry/ExtensionDescent/GlobalSectionsRankAlgClosed.lean}{ReciprocityTheory.ChartRestriction.G.Out.Laws}
\leanref{ReciprocityTheory/ReductiveGroups/ArithmeticGeometry/ExtensionDescent/GlobalSectionsRankAlgClosed.lean}{ReciprocityTheory.ChartRestriction.G.Out.PolBaseChangeSt}

\begin{lemma}\label{cv-cusp-cm-point-eigen-degree-tensor-complex-pairing-hermitian}
Assume the conclusion of Theorem~\ref{cv-cusp-cm-point-eigen-degree-rank-one-norm-riemann-theta} and the laws of
abelian schemes (Definition~\ref{cv-cusp-cm-point-eigen-degree-tensor-laws}); the laws are not used in the proof.
Let $X$, $Y$, $A$ and $B$ be principally polarized abelian varieties over $\mathbf C$, with the line bundles of their
polarizations trivialized on open covers, and for $n\ge1$ let $e^A_n$ and $e^B_n$ be the commutator pairings of the theta
groups of the $n$-th powers of the line bundles of $A$ and $B$
(Definition~\ref{cv-cusp-cm-point-eigen-degree-rank-one-norm-weil-translate-group}). Let $u_X\colon\mathbf
C^{g_X}\to X(\mathbf C)$, $u_Y$, $u_A\colon\mathbf C^{g_A}\to A(\mathbf C)$ and $u_B\colon\mathbf C^{g_B}\to B(\mathbf C)$ be
uniformizations for these line bundles (Definition~\ref{cv-origin-hodge-norm-gram}(1)), with kernels $\Lambda_X$,
$\Lambda_Y$, $\Lambda_A$ and $\Lambda_B$, and with Appell--Humbert data $(H_A,\chi_A)$ for $u_A$ and $(H_B,\chi_B)$ for
$u_B$. Let $f_1\colon X\to A$, $f_2\colon Y\to A$, $g_1\colon X\to B$ and $g_2\colon Y\to B$ be morphisms and $F_1$, $F_2$,
$G_1$, $G_2$ $\mathbf C$-linear maps with $u_A\circ F_1=f_1\circ u_X$, $u_A\circ F_2=f_2\circ u_Y$, $u_B\circ G_1=g_1\circ
u_X$ and $u_B\circ G_2=g_2\circ u_Y$. Then $e^A_n(f_1x,f_2y)=e^B_n(g_1x,g_2y)$ for all $n\ge1$ and all $x\in X(\mathbf
C)$, $y\in Y(\mathbf C)$ with $nx=ny=0$ if and only if
$F_1(z)^{\mathsf T}H_A\overline{F_2(w)}=G_1(z)^{\mathsf T}H_B\overline{G_2(w)}$ for all $z\in\mathbf C^{g_X}$ and
$w\in\mathbf C^{g_Y}$.
\end{lemma}
\leanref{Development/AbelianVarietyTheory/Proofs/CubicalLineBundles/WeilEqIffHermFormEqComplexStatement.lean}{GlobalTheory.T.weilEq_iff_hermForm_eq_complex}
\leanref{Development/AbelianVarietyTheory/Proofs/CubicalLineBundles/WeilEqIffHermFormEqComplex.lean}{GlobalTheory.CM.weilEq_iff_hermForm_eq}
\begin{proof}
Step 1: the lattices. For $l\in\Lambda_X$, $u_A(F_1(l))=f_1(u_X(l))=f_1(u_X(0))=u_A(F_1(0))$ is the origin, so
$F_1(\Lambda_X)\subset\Lambda_A$; in the same way $F_2(\Lambda_Y)\subset\Lambda_A$, $G_1(\Lambda_X)\subset\Lambda_B$ and
$G_2(\Lambda_Y)\subset\Lambda_B$.

Step 2: the values. Let $a\in\Lambda_X$ and $b\in\Lambda_Y$. By
Lemma~\ref{cv-cusp-cm-point-eigen-degree-rank-one-norm-riemann}(1) for $u_A$ and $u_B$, the numbers
$m=\operatorname{Im}\big(F_1(a)^{\mathsf T}H_A\overline{F_2(b)}\big)$ and
$m'=\operatorname{Im}\big(G_1(a)^{\mathsf T}H_B\overline{G_2(b)}\big)$ are integers, and by (3) of that lemma, for $n\ge1$,
\[
 e^A_n\big(f_1u_X(a/n),f_2u_Y(b/n)\big)=e^A_n\big(u_A(F_1(a)/n),u_A(F_2(b)/n)\big)=\exp(-2\pi im/n),
\]
and in the same way $e^B_n(g_1u_X(a/n),g_2u_Y(b/n))=\exp(-2\pi im'/n)$. The points $u_X(a/n)$ and $u_Y(b/n)$ are killed
by $n$.

Step 3: from the pairings to the forms. Suppose that the pairings agree. By Step~2, $\exp(-2\pi im/n)=\exp(-2\pi im'/n)$
for every $n\ge1$, so $n$ divides $m-m'$ for every $n$, and $m=m'$. So the $\mathbf R$-bilinear forms
$(z,w)\mapsto\operatorname{Im}\big(F_1(z)^{\mathsf T}H_A\overline{F_2(w)}\big)$ and
$(z,w)\mapsto\operatorname{Im}\big(G_1(z)^{\mathsf T}H_B\overline{G_2(w)}\big)$ agree on $\Lambda_X\times\Lambda_Y$, which
spans $\mathbf C^{g_X}\times\mathbf C^{g_Y}$ over $\mathbf R$, so they agree everywhere. For a form $\beta$ that is linear in
the first variable, $\operatorname{Re}\beta(z,w)=\operatorname{Im}\beta(iz,w)$; so the two forms of the statement agree.

Step 4: from the forms to the pairings. Suppose that the forms agree, and let $x\in X(\mathbf C)$ and $y\in Y(\mathbf C)$
be killed by $n\ge1$. As $u_X$ is onto, $x=u_X(z)$ for some $z$, and $a=nz$ lies in $\Lambda_X$, as $u_X(nz)=nx=0$; so
$x=u_X(a/n)$. In the same way $y=u_Y(b/n)$ with $b\in\Lambda_Y$. By hypothesis $m=m'$, so the pairings agree by Step~2.
\end{proof}
% Lean: weilEq_iff_hermForm_eq, under ThetaUnifAV and Out.Laws: HermWeil.map_mem (Step 1); key with HermWeil.weilIs_iff
% (Lemma cv-cusp-cm-point-eigen-degree-rank-one-norm-riemann(1) and (3), through AHWeil.tangent_and_weil_of_isAH and
% the uniqueness weilIs_unique_fieldPoint) and map_smul (Step 2); hermForm_eq_of_im with eq_of_rootVal and ex_div_pow
% (Step 3); exists_eq_ex_div (Step 4). The Lean phrases the agreement of the pairings as WeilEq, an equivalence of the
% relations WeilIs for every value.

\begin{lemma}\label{cv-cusp-cm-point-eigen-degree-tensor-complex-datum-structure}
In Situation~\ref{cv-cusp-setting}, let $\sqrt{-d}\in E$ be the square root of $-d$ of the datum, $\theta_0\in\mathbf C$ the
square root of $-d$ fixed by the datum, and $s\mapsto\bar s$ complex conjugation on $E$. Let $L$ be an
$\mathcal O_E$-module with a $\mathbf Z$-bilinear form $\langle\ ,\ \rangle\colon L\times L\to\mathbf Z$ such that
$x\mapsto\langle x,\cdot\rangle$ is a bijection of $L$ onto $\operatorname{Hom}(L,\mathbf Z)$ and $\langle
sx,y\rangle=\langle x,\bar sy\rangle$ for all $s\in\mathcal O_E$ and $x,y\in L$. Let $(\sigma_k)_{1\le k\le g}$ be a type of
signature $(p,q)$ for $\theta_0$ (Definition~\ref{cv-cm-point-periods}(2)), and let $P\colon L\to\mathbf C^g$ be an
injective additive map with $P(sx)_k=\sigma_k(s)P(x)_k$ for all $s$, $x$ and $k$, whose image is discrete and spans
$\mathbf C^g$ over $\mathbf R$. Let $H\in M_g(\mathbf C)$ be a positive definite Hermitian matrix with
$\operatorname{Im}\big(P(x)^{\mathsf T}H\,\overline{P(y)}\big)=\langle x,y\rangle$ for all $x,y\in L$. Then $L$, with
$\langle\ ,\ \rangle$, is the lattice of a PEL datum over $E$ of signature $(p,q)$ (Definition~\ref{cv-cm-point-periods}).
\end{lemma}
\leanref{Development/AbelianVarietyTheory/Proofs/ModuliProofs/TwoMulEqOfIsTypeFn.lean}{GlobalTheory.PeriodDatum.exists_datum}
\begin{proof}
Write $B(z,w)=z^{\mathsf T}H\bar w$.

Step 1: the lattice. The image $P(L)$ is discrete and spans $\mathbf C^g$ over $\mathbf R$, so it is a free
$\mathbf Z$-module of rank $2g$, and so is $L\cong P(L)$. The $\mathbf R$-linear extension
$e\colon L\otimes\mathbf R\to\mathbf C^g$ of $P$ is onto, as $P(L)$ spans, so it is an isomorphism, as both sides have
dimension $2g$ over $\mathbf R$.

Step 2: the module. The $\mathcal O_E$-module $L$ is finitely generated, as it is over $\mathbf Z$, and torsion free: if
$sx=0$ with $s\neq0$, then $\sigma_k(s)P(x)_k=0$ with $\sigma_k(s)\neq0$ for every $k$, so $P(x)=0$ and $x=0$. A finitely
generated torsion-free module over the Dedekind domain $\mathcal O_E$ is projective. Its rank $n$ satisfies
$[E:\mathbf Q]\,n=2g$. Counting the values of the type, $g=(p+q)[E:\mathbf Q]/2$, as half of the embeddings send
$\sqrt{-d}$ to $\theta_0$. So $n=p+q$.

Step 3: the form. For $x\in L$, $\langle x,x\rangle=\operatorname{Im}B(P(x),P(x))=0$, as $H$ is Hermitian; so the form is
alternating, and it is perfect and adjoint for the conjugation by hypothesis. By bilinearity, its extension to
$L\otimes\mathbf R$ is $(x,y)\mapsto\operatorname{Im}B(e(x),e(y))$.

Step 4: the complex structure. Let $J=e^{-1}\circ i\circ e$ on $L\otimes\mathbf R$, with $i$ the multiplication by $i$ on
$\mathbf C^g$. Then $J^2=-1$; $\langle Jx,Jy\rangle=\operatorname{Im}B(ie(x),ie(y))=\langle x,y\rangle$; and for $x\neq0$,
$\langle Jx,x\rangle=\operatorname{Im}\big(iB(e(x),e(x))\big)=B(e(x),e(x))>0$. For $s\in\mathcal O_E$, the action of $s$ on
$L\otimes\mathbf R$ is $e^{-1}\circ D_s\circ e$, with $D_s$ the diagonal map $z\mapsto(\sigma_k(s)z_k)_k$; it commutes with
$J$, and its trace on the complex vector space $L\otimes\mathbf R$, with $i$ acting as $J$, is
$\sum_k\sigma_k(s)=\sum_\sigma\big(p\,\sigma(s)+q\,\overline{\sigma(s)}\big)$, the sum over the embeddings $\sigma$ with
$\sigma(\sqrt{-d})=\theta_0$, as the type takes these $p$ times and their conjugates $q$ times. So $J$ is a complex
structure of signature $(p,q)$ for $\theta_0$ (Definition~\ref{cv-cm-point-periods}(5)).
\end{proof}
% Lean: exists_datum: lattice_equiv (Step 1: ZLattice.module_finite, module_free and ZLattice.rank, with the
% liftBaseChange e); htor, IsDedekindDomain.flat_iff_torsion_eq_bot and Module.Flat.projective_of_finitePresentation,
% and hrank with finrank_int_eq_mul_of_finite and two_mul_eq_of_isTypeFn (= PeriodNonempty.sum_type, Step 2); alt from
% him and hermForm_self_im, and hB (Step 3); J with hJ, and the trace conditions through hdt
% (PeriodNonempty.trace_real) and sum_type (Step 4). The Lean's complex structure of signature (p, q) asks for the real
% traces of A and of J A, 2 Re t and -2 Im t, which is the text's complex trace t. The datum is La itself, with the
% identity as the isomorphism.

\begin{lemma}\label{cv-cusp-cm-point-eigen-degree-tensor-complex-datum}
Assume the conclusion of Theorem~\ref{cv-cusp-cm-point-eigen-degree-rank-one-norm-riemann-theta} and the laws of
abelian schemes (Definition~\ref{cv-cusp-cm-point-eigen-degree-tensor-laws}).
In Situation~\ref{cv-cusp-setting}, let $\sqrt{-d}\in E$ be the square root of $-d$ of the datum and $\theta_0\in\mathbf C$
the square root of $-d$ fixed by the datum. Let $A$ be a principally polarized abelian variety over $\mathbf C$ of
dimension $h$ with an action of $\mathcal O_E$ for which the Rosati involution of the polarization induces complex
conjugation on $\mathcal O_E$. Let $p,q\ge0$, let $(\sigma_k)_{1\le k\le g}$ be a type of signature $(p,q)$ for
$\theta_0$ (Definition~\ref{cv-cm-point-periods}(2)), and suppose that for every embedding $\tau\colon E\to\mathbf C$
the $\tau$-eigenspace of the Lie algebra of $A$ has dimension $p$ if $\tau(\sqrt{-d})=\theta_0$ and $q$ otherwise. Let
$u\colon\mathbf C^h\to A(\mathbf C)$ be a uniformization of $A$ for the line bundle of its polarization
(Definition~\ref{cv-origin-hodge-norm-gram}(1)) with kernel $\Lambda$ and Appell--Humbert data $(H,\chi)$. Then there
are the lattice $L$ of a PEL datum over $E$ of signature $(p,q)$, with its perfect alternating form
$\langle\ ,\ \rangle$, a period $P\colon L\to\mathbf C^g$ of type $(\sigma_k)$ (Definition~\ref{cv-cm-point-periods}(3))
and a $\mathbf C$-linear bijection $F\colon\mathbf C^g\to\mathbf C^h$ such that:
\begin{enumerate}
\item $F(z)\in\Lambda$ if and only if $z\in P(L)$;
\item $\operatorname{Im}\big(F(P(a))^{\mathsf T}H\,\overline{F(P(b))}\big)=\langle a,b\rangle$ for all $a,b\in L$;
\item $u(F(s\cdot z))=s\big(u(F(z))\big)$ for all $s\in\mathcal O_E$ and $z\in\mathbf C^g$, where
  $s\cdot z=(\sigma_k(s)z_k)_k$.
\end{enumerate}
\end{lemma}
\leanref{Development/AbelianVarietyTheory/Proofs/CubicalLineBundles/ExistsHomFamilyOfPureTypeComplexTriple.lean}{GlobalTheory.CM.PureTriple.exists_datum_period}
\begin{proof}
Write $\bar s$ for the complex conjugate of $s\in\mathcal O_E$, and $s\cdot z=(\sigma_k(s)z_k)_k$ for $z\in\mathbf C^g$.

Step 1: the lattice. By Lemma~\ref{cv-cusp-cm-point-eigen-degree-rank-one-norm-riemann}, $E(x,y)=\operatorname{Im}(x^{\mathsf
T}H\bar y)$ takes integer values on $\Lambda\times\Lambda$, $x\mapsto E(x,\cdot)$ is a bijection of $\Lambda$ onto
$\operatorname{Hom}(\Lambda,\mathbf Z)$, and every $z\in\mathbf C^h$ is the differential of exactly one tangent vector
$v_z$ of $A$ at the origin, and every tangent vector has exactly one differential. The map $z\mapsto v_z$ is additive
and $\mathbf C$-linear for the operations of the Lie algebra.

Step 2: the action on $\mathbf C^h$. For $s\in\mathcal O_E$, Lemma~\ref{cv-origin-hodge-norm-form-lift}, applied to $u$
and $s\circ u$, gives a $\mathbf C$-linear map $\rho(s)$ with $u\circ\rho(s)=s\circ u$. A $\mathbf C$-linear map $D$ with
$u\circ D=0$ maps $\mathbf C^h$ into the discrete group $\Lambda$, so it is $0$; hence $\rho(s)$ is unique, and
$s\mapsto\rho(s)$ is a ring homomorphism $\mathcal O_E\to\operatorname{End}_{\mathbf C}(\mathbf C^h)$, as $s\mapsto s\circ u$
respects sums and products. As in Step~1 of the proof of
Lemma~\ref{cv-cusp-cm-point-eigen-degree-tensor-complex-pairing-hermitian}, $\rho(s)(\Lambda)\subset\Lambda$.

Step 3: diagonalization. By the chain rule, as in the proof of Lemma~\ref{cv-origin-hodge-norm-form-compare},
$s\circ v_z=v_{\rho(s)z}$; so $z\mapsto v_z$ carries the $\tau$-eigenspace of $\rho$, for an embedding
$\tau\colon E\to\mathbf C$, onto the $\tau$-eigenspace of the Lie algebra of $A$, which has dimension $p$ if
$\tau(\sqrt{-d})=\theta_0$ and $q$ otherwise. The diagonal action $s\mapsto(z\mapsto s\cdot z)$ on $\mathbf C^g$ has the
same eigenspace dimensions, as the type takes each embedding $p$ or $q$ times. As $\mathcal O_E\otimes_{\mathbf Z}\mathbf C$
is the product of copies of $\mathbf C$ indexed by the embeddings $E\to\mathbf C$, both are the direct sums of their
eigenspaces; so there is a $\mathbf C$-linear bijection $F\colon\mathbf C^g\to\mathbf C^h$ with $F(s\cdot z)=\rho(s)F(z)$
for all $s$ and $z$.

Step 4: adjointness. For $s\in\mathcal O_E$, the Rosati condition is
$\Lambda_{\mathcal L}(s\,p_1,p_2)\cong\Lambda_{\mathcal L}(p_1,\bar s\,p_2)$ on $A\times A$, for the line bundle $\mathcal L$ of
the polarization. By Lemma~\ref{cv-cusp-cm-point-eigen-degree-rank-one-norm-weil}, the commutator pairings
$e_n(sx,y)$ and $e_n(x,\bar sy)$ agree; by
Lemma~\ref{cv-cusp-cm-point-eigen-degree-tensor-complex-pairing-hermitian}, for the morphisms $s$, the identity, the
identity and $\bar s$ with the lifts $\rho(s)$, the identity, the identity and $\rho(\bar s)$,
$(\rho(s)z)^{\mathsf T}H\bar w=z^{\mathsf T}H\,\overline{\rho(\bar s)w}$ for all $z,w\in\mathbf C^h$. Hence
$E(\rho(s)x,y)=E(x,\rho(\bar s)y)$.

Step 5: the datum. Make $\Lambda$ an $\mathcal O_E$-module through $\rho$, and let $P_0\colon\Lambda\to\mathbf C^g$ be the
restriction of $F^{-1}$. It is injective and additive, $P_0(sx)=s\cdot P_0(x)$ by Step~3, and its image is discrete and
spans $\mathbf C^g$ over $\mathbf R$, as $\Lambda$ does in $\mathbf C^h$ and $F$ is a linear bijection. The matrix $H'$ with
$z^{\mathsf T}H'\bar w=F(z)^{\mathsf T}H\,\overline{F(w)}$ is positive definite and Hermitian, with
$\operatorname{Im}(P_0(x)^{\mathsf T}H'\,\overline{P_0(y)})=E(x,y)$. By Steps 1 and 4,
Lemma~\ref{cv-cusp-cm-point-eigen-degree-tensor-complex-datum-structure} shows that $\Lambda$, with $E$, is the lattice
$L$ of a PEL datum over $E$ of signature $(p,q)$, with $\langle\ ,\ \rangle=E$; and $P=P_0$ is a period of $L$ of type
$(\sigma_k)$, for the form $H'$.

Step 6: the properties. (1) $F(z)\in\Lambda$ if and only if $z=P(F(z))\in P(L)$. (2)
$\operatorname{Im}\big(F(P(a))^{\mathsf T}H\,\overline{F(P(b))}\big)=E(a,b)=\langle a,b\rangle$. (3)
$u(F(s\cdot z))=u(\rho(s)F(z))=s\big(u(F(z))\big)$.
\end{proof}
% Lean: exists_datum_period: AHWeil.tangent_and_weil_of_isAH (Lemma cv-cusp-cm-point-eigen-degree-rank-one-norm-riemann:
% E_f = Im H on Lambda, unimodular, and the tangent bijections); isExp_push and IsExp.exists_rep (the action of O_E on
% C^h by linear lifts rho(a), from HomLinear.exists_linear_of_isMonHom); exists_equiv_of_finrank_eq with
% isSplitting_embeddings, finrank_eigenSub_diagRep and IsExp.finrank_eigenSub (a C-linear h : C^g = C^h carrying the
% diagonal action of the type to rho, from the eigenspace dimensions); HermWeil.map_mem (rho(a) preserves Lambda);
% weilEq_iff_hermForm_eq with weilEq_of_lamEq and T.t.adj (H(rho(a)z, w) = H(z, rho(bar a)w), from the Rosati
% condition); exists_hermForm_comp (the form H pulled back along h); and PeriodDatum.exists_datum (the datum D with
% lattice Lambda as an O_E-module, form E_f and signature (p, q), with an O_E-linear isomorphism eps). F is h, and P is
% h^{-1} restricted to Lambda, composed with eps (the identity in PeriodDatum.exists_datum). Step 1 is
% tangent_and_weil_of_isAH; Step 2 is isExp_push and IsExp.exists_rep (uniqueness by linear_eq_zero_of_discrete) with
% HermWeil.map_mem; Step 3 is IsExp.finrank_eigenSub, finrank_eigenSub_diagRep, isSplitting_embeddings and
% exists_equiv_of_finrank_eq; Step 4 is hadj; Step 5 is the module structure, hP, hrange, hPs, hPi, hPd, hPf,
% exists_hermForm_comp, hpos, him and PeriodDatum.exists_datum; Step 6 is the final refine.

\begin{lemma}\label{cv-cusp-cm-point-eigen-degree-tensor-complex-homs-line}
In Situation~\ref{cv-cusp-setting}, let $\sqrt{-d}\in E$ be the square root of $-d$ of the datum, and put
$w_0=(1+\sqrt{-d})/2$. Let $L$ be the lattice of a PEL datum over $E$ with its perfect alternating form
$\langle\ ,\ \rangle$, let $\sigma_1,\dots,\sigma_g$ be embeddings $E\to\mathbf C$, let $P\colon L\to\mathbf C^g$ be a
period of $L$ for them (Definition~\ref{cv-cm-point-periods}(3)), and let $H\in M_g(\mathbf C)$ be a positive definite
Hermitian matrix with $\operatorname{Im}\big(P(x)^{\mathsf T}H\,\overline{P(y)}\big)=\langle x,y\rangle$ for all $x,y\in L$.
Let $\rho\in\mathbf C$ with $\rho^2=-d$, and let $a\in L$ with $P(a)_k=0$ for every $k$ with $\sigma_k(\sqrt{-d})\neq\rho$.
Then:
\begin{enumerate}
\item $\big(m+n\frac{1+i\sqrt d}2\big)P(a)\in P(L)$ for all $m,n\in\mathbf Z$;
\item $(zP(a))^{\mathsf T}H\,\overline{z'P(a)}=|\langle w_0a,a\rangle|\cdot2z\bar z'/\sqrt d$ for all $z,z'\in\mathbf C$.
\end{enumerate}
\end{lemma}
\leanref{Development/AbelianVarietyTheory/Proofs/CubicalLineBundles/PureTypeVectorSpansQuadraticLattice.lean}{GlobalTheory.pureType_vector_spans_quadratic_lattice}
\leanref{Development/AbelianVarietyTheory/Proofs/CubicalLineBundles/ExistsIsHomFamilyComplex.lean}{GlobalTheory.PureType.hermForm_line}
\begin{proof}
As $2w_0=1+\sqrt{-d}$, $P(w_0a)_k=\sigma_k(w_0)P(a)_k=\frac{1+\rho}2P(a)_k$ for every $k$, since $P(a)_k=0$ unless
$\sigma_k(\sqrt{-d})=\rho$. So $P(w_0a)=\frac{1+\rho}2P(a)$. As $\rho^2=-d$, $\rho=\pm i\sqrt d$.

(1) If $\rho=i\sqrt d$, then $\big(m+n\frac{1+i\sqrt d}2\big)P(a)=P(ma+nw_0a)$. If $\rho=-i\sqrt d$, then
$\frac{1+i\sqrt d}2=1-\frac{1+\rho}2$, and $\big(m+n\frac{1+i\sqrt d}2\big)P(a)=P((m+n)a-nw_0a)$.

(2) Put $r=P(a)^{\mathsf T}H\,\overline{P(a)}$, a real number with $r\ge0$, as $H$ is Hermitian and positive definite. Then
$(zP(a))^{\mathsf T}H\,\overline{z'P(a)}=z\bar z'r$, and
$\langle w_0a,a\rangle=\operatorname{Im}\big(\frac{1+\rho}2r\big)=\pm\frac{\sqrt d}2r$. So $r=2|\langle
w_0a,a\rangle|/\sqrt d$.
\end{proof}
% Lean: pureType_vector_spans_quadratic_lattice: hw0 and hw0' (with two_weight_eq_one_add_root), hsq for rho = +-i sqrt d,
% the two cases of (1), and hermForm_smul_smul, hermForm_self_im and habs for (2). PureType.hermForm_line is the same
% statement for the form H pulled back along a linear bijection F (hermForm_mulVec), which is how Lemma
% cv-cusp-cm-point-eigen-degree-tensor-complex-homs applies it.

\begin{lemma}\label{cv-cusp-cm-point-eigen-degree-tensor-complex-homs-curve}
Let $d\ge1$, $c\in\mathbf C$ nonzero, and $\Lambda=c(\mathbf Z+\mathbf Z\frac{1+i\sqrt d}2)\subset\mathbf C$. Let $H_0>0$ be a
real number such that $E(x,y)=\operatorname{Im}(xH_0\bar y)$ takes integer values on $\Lambda\times\Lambda$ and $x\mapsto
E(x,\cdot)$ maps $\Lambda$ onto $\operatorname{Hom}(\Lambda,\mathbf Z)$. Then $(cz)H_0\overline{(cw)}=2z\bar w/\sqrt d$ for all
$z,w\in\mathbf C$.
\end{lemma}
\leanref{Development/AbelianVarietyTheory/Proofs/CubicalLineBundles/ExistsIsHomFamilyComplex.lean}{GlobalTheory.PureType.hermForm_of_isQuadLattice}
\begin{proof}
Put $r=|c|^2H_0>0$ and $w=\frac{1+i\sqrt d}2$, so that $(cz)H_0\overline{(cw')}=z\bar w'r$ for all $z,w'$. Let $e_1=c$ and
$e_2=cw$ in $\Lambda$. Then $E(e_2,e_1)=\operatorname{Im}(w)r=\frac{\sqrt d}2r>0$. For $x=c(m+nw)$ and $y=c(m'+n'w)$ in
$\Lambda$, $E(x,y)=\operatorname{Im}\big((m+nw)(m'+n'\bar w)\big)r=(nm'-mn')E(e_2,e_1)$, so $E(e_2,e_1)$ divides every value of
$E$ on $\Lambda$. The map $y\mapsto E(e_2,y)/E(e_2,e_1)$ is a homomorphism $\Lambda\to\mathbf Z$; by hypothesis it is
$E(x,\cdot)$ for some $x\in\Lambda$, so $E(x,e_1)=1$. As $E(e_2,e_1)$ divides $E(x,e_1)=1$ and is positive, it is $1$, and
$r=2/\sqrt d$.
\end{proof}
% Lean: hermForm_of_isQuadLattice, for g = 1, a lattice with IsQuadLattice Lambda d c, a positive Hermitian 1 x 1 matrix
% and Ef with Ef = Im hermForm on Lambda and Ef surjective: key, e_1 and e_2, hN, hN0, hdvd, f and hb, h1, h2, h3.

\begin{lemma}\label{cv-cusp-cm-point-eigen-degree-tensor-complex-homs-points}
Let $X$ be a reduced scheme, locally of finite type over $\mathbf C$, whose complex points are its points with values in
$\mathbf C$ over $\mathbf C$, and let $U\subset X$ be an open subscheme.
\begin{enumerate}
\item An open subscheme of $X$ that contains every complex point of $U$ contains $U$.
\item A regular function on $U$ that vanishes at every complex point of $U$ is $0$, and a regular function on $U$ that
  vanishes at no complex point of $U$ is a unit.
\item Two morphisms from $U$ to a separated scheme that agree at every complex point of $U$ are equal.
\end{enumerate}
\end{lemma}
\leanref{Development/AbelianVarietyTheory/Proofs/CubicalLineBundles/SurjectiveOfComplexPointsLiftOfUniversallyClosed.lean}{GlobalTheory.ComplexPoints.opens_le_of_complexPoints}
\leanref{Development/AbelianVarietyTheory/Proofs/CubicalLineBundles/SurjectiveOfComplexPointsLiftOfUniversallyClosed.lean}{GlobalTheory.ComplexPoints.eq_of_valAt'_eq}
\leanref{Development/AbelianVarietyTheory/Proofs/CubicalLineBundles/SurjectiveOfComplexPointsLiftOfUniversallyClosed.lean}{GlobalTheory.ComplexPoints.isUnit_of_valAt'_ne_zero}
\leanref{Development/AbelianVarietyTheory/Proofs/CubicalLineBundles/SurjectiveOfComplexPointsLiftOfUniversallyClosed.lean}{GlobalTheory.ComplexPoints.hom_ext_of_points}
\begin{proof}
The scheme $U$ is locally of finite type over $\mathbf C$, so it is a Jacobson scheme and its closed points have residue
fields algebraic over $\mathbf C$, hence equal to $\mathbf C$ \cite[Tag 0478]{Sta}; so each closed point of $U$ is the image
of a complex point. As $U$ is Jacobson, every closed subset of $U$ is the closure of the closed points of $U$ that it
contains \cite[Tag 01P2]{Sta}, \cite[Tag 005U]{Sta}; so a nonempty closed subset of $U$ contains a complex point, and a
closed subset that contains every complex point is $U$.

(1) If $V$ is such an open subscheme, $U\setminus V$ is a closed subset of $U$ without complex points, so it is empty.

(2) Let $f$ vanish at every complex point. Its zero set is closed in $U$ and contains every complex point, so it is $U$, and
$f$ lies in every prime ideal of the ring of each affine open subscheme of $U$; so it is nilpotent there, hence $0$, as $U$
is reduced. If $f$ vanishes at no complex point, the open set where $f$ does not vanish contains every complex point of $U$,
so it is $U$ by (1), and $f$ is a unit.

(3) Let $a,b\colon U\to Y$ be the morphisms. As $Y$ is separated, the inverse image $E$ of the diagonal of $Y$ under
$(a,b)\colon U\to Y\times Y$ is a closed subscheme of $U$, and a complex point at which $a$ and $b$ agree factors through
$E$. So the closed set $E$ contains every complex point, hence is $U$. The ideal of $E$ then consists of functions that
vanish at every point, which are nilpotent, hence $0$, as $U$ is reduced. So $E=U$ and $a=b$.
\end{proof}
% Lean: opens_le_of_complexPoints (= opens_le_of_forall_complex_section, through the Jacobson property of schemes locally
% of finite type over C), eq_of_valAt'_eq (eq_zero_of_basicOpen_eq_bot), isUnit_of_valAt'_ne_zero (isUnit_res_basicOpen)
% and hom_ext_of_points (with hom_ext_of_complexPoints for morphisms over C).

\begin{lemma}\label{cv-cusp-cm-point-eigen-degree-tensor-complex-homs-sign}
Let $A$ be a principally polarized abelian variety over $\mathbf C$ of dimension $g$, with the line bundle $\mathcal L$ of
its polarization trivialized on an open cover, and let $u\colon\mathbf C^g\to A(\mathbf C)$ be a uniformization of $A$
for $\mathcal L$ (Definition~\ref{cv-origin-hodge-norm-gram}(1), with $T=\operatorname{Spec}\mathbf C$, $G=A$ and
$\varphi$ the identity) with kernel $\Lambda$ and Appell--Humbert data $(H,\chi)$. Then $\chi(l)^2=1$ for every
$l\in\Lambda$.
\end{lemma}
\leanref{Development/AbelianVarietyTheory/Proofs/CubicalLineBundles/SemicharacterSqEqOneOfSymmetric.lean}{GlobalTheory.semicharacter_sq_eq_one_of_symmetric}
\begin{proof}
Let $(U_\alpha)$ be the open cover, $g_{\alpha\beta}$ the transition functions and $\theta_\alpha$ the trivializing
functions of $u$, and put $e_l(x)=\exp\big(\pi x^{\mathsf T}H\bar l+\frac\pi2l^{\mathsf T}H\bar l\big)$, so that
$\theta_\alpha(x)=\chi(l)e_l(x)\theta_\alpha(x+l)$ where defined. Let $l\in\Lambda$.

Step 1: $\chi(l)\chi(-l)=1$. Let $\alpha$ be an index with $u(0)\in U_\alpha$. As $u(l)=u(0)$, the functional equation at
$x=0$ for $l$ and at $x=l$ for $-l$ gives $\theta_\alpha(0)=\chi(l)e_l(0)\theta_\alpha(l)$ and
$\theta_\alpha(l)=\chi(-l)e_{-l}(l)\theta_\alpha(0)$. As $e_l(0)e_{-l}(l)=1$ and $\theta_\alpha(0)\neq0$,
$\chi(l)\chi(-l)=1$.

Step 2: $\chi(-l)=\chi(l)$. The line bundle of the polarization is symmetric: locally on the base, which here has one point,
$[-1]^*\mathcal L\otimes\mathcal L^{-1}$ has a generator $\tau$. In the trivializations, $\tau$ is given by regular functions
$\tau_{\alpha\beta}$ without zeros on $[-1]^{-1}(U_\alpha)\cap U_\beta$, with
$\tau_{\alpha\beta}=(g_{\alpha\alpha'}\circ[-1])\,g_{\beta\beta'}^{-1}\tau_{\alpha'\beta'}$ where defined. For $x\in\mathbf C^g$
with $u(-x)\in U_\alpha$ and $u(x)\in U_\beta$ put $r(x)=\tau_{\alpha\beta}(u(x))\,\theta_\beta(x)/\theta_\alpha(-x)$. By the
transition rules of $\tau$ and of the trivializing functions, $r(x)$ does not depend on $\alpha$ and $\beta$, so $r$ is a
holomorphic function on $\mathbf C^g$ without zeros. As $u(x+l)=u(x)$ and $e_{-l}(-x)=e_l(x)$, the functional
equations of $\theta_\beta$ and $\theta_\alpha$ give $r(x+l)=\chi(-l)\chi(l)^{-1}r(x)$.
Lemma~\ref{cv-origin-hodge-norm-form-pullback-trivial}, for $\Lambda$, the zero matrix, $c(l)=\chi(-l)\chi(l)^{-1}$, which
has absolute value $1$, and $r$, gives $\chi(-l)=\chi(l)$.

By Steps 1 and 2, $\chi(l)^2=\chi(l)\chi(-l)=1$.
\end{proof}
% Lean: semicharacter_sq_eq_one_of_symmetric: mul_neg_eq_one (Step 1, from h.2.2.2 at (l, 0) and (-l, l)) and neg_eq
% (Step 2: exists_symmetry_values for T i j = tau_{ij} o u from the field sym of the polarization, R i j and r, hd, hne, and
% appellHumbert_trivial_of_nowhere_vanishing, Lemma cv-origin-hodge-norm-form-pullback-trivial).

\begin{lemma}\label{cv-cusp-cm-point-eigen-degree-tensor-complex-homs-glue}
Let $X$ be a reduced scheme, locally of finite type over $\mathbf C$, and $Y$ a scheme with a separated morphism to an
affine scheme. Let $Z$ be a set, $e$ a map from $Z$ onto the complex points of $X$, and $e'$ a map from $Z$ to the points
of $Y$ with values in $\mathbf C$. Suppose that for every $z\in Z$ there are an open subscheme $W\subset X$ containing
$e(z)$ and an affine open subscheme $D\subset Y$ such that $e'(w)$ lies in $D$ for every $w$ with $e(w)$ in $W$, and such
that for every regular function $f$ on $D$ there is a regular function $h$ on $W$ with $h(e(w))=f(e'(w))$ for every $w$
with $e(w)$ in $W$. Then there is a morphism $\phi\colon X\to Y$ with $\phi\circ e(z)=e'(z)$ for every $z\in Z$.
\end{lemma}
\leanref{Development/AbelianVarietyTheory/Proofs/CubicalLineBundles/DenseClosedPoints.lean}{GlobalTheory.SchemeHomsFromPoints.exists_hom}
\leanref{Development/AbelianVarietyTheory/Proofs/CubicalLineBundles/DenseClosedPoints.lean}{GlobalTheory.SchemeHomsFromPoints.exists_local}
\begin{proof}
The scheme $Y$ is separated, as it is separated over an affine scheme.

Step 1: local morphisms. Let $z\in Z$, with $W$ and $D$ as in the hypothesis. For a regular function $f$ on $D$ the function
$h$ is unique, by Lemma~\ref{cv-cusp-cm-point-eigen-degree-tensor-complex-homs-points}(2) applied to the difference of two
choices, as $e$ is onto the complex points of $W$. So $f\mapsto h$ is a map $\rho\colon\Gamma(D,\mathcal O_Y)\to\Gamma(W,\mathcal
O_X)$, and it is a ring homomorphism, by the same uniqueness, as values at points respect sums and products. Let
$\phi_z\colon W\to D$ be the canonical morphism $W\to\operatorname{Spec}\Gamma(W,\mathcal O_X)$ followed by
$\operatorname{Spec}\rho$ and the isomorphism $\operatorname{Spec}\Gamma(D,\mathcal O_Y)\cong D$. For $w$ with $e(w)$ in $W$,
the point $\phi_z\circ e(w)$ of $D$ corresponds to the homomorphism $f\mapsto h(e(w))=f(e'(w))$ to $\mathbf C$, so
$\phi_z\circ e(w)=e'(w)$.

Step 2: the cover. The open subschemes $W$ of Step~1, for $z\in Z$, contain every complex point of $X$, so they cover $X$,
by Lemma~\ref{cv-cusp-cm-point-eigen-degree-tensor-complex-homs-points}(1).

Step 3: gluing. For $z,z'\in Z$, the morphisms $\phi_z$ and $\phi_{z'}$ agree at every complex point $e(w)$ of $W\cap W'$,
as both carry it to $e'(w)$; so they agree on $W\cap W'$, by
Lemma~\ref{cv-cusp-cm-point-eigen-degree-tensor-complex-homs-points}(3). They glue to a morphism $\phi\colon X\to Y$, and
$\phi\circ e(z)=\phi_z\circ e(z)=e'(z)$.
\end{proof}
% Lean: exists_hom: exists_local (Step 1: the ring homomorphism phi from hF' and the injectivity inj, through
% eq_zero_of_points, and W.toSpecGamma, Spec.map phi and hD.fromSpec), hcov (Step 2: the Jacobson space X, through
% nonempty_inter_closedPoints and pointOfClosedPoint), hglue with hom_ext_of_points (Step 3) and glueMorphisms.

\begin{lemma}\label{cv-cusp-cm-point-eigen-degree-tensor-complex-homs-hom}
Let $P$ and $P'$ be principally polarized abelian varieties over $\mathbf C$ of dimensions $g$ and $g'$, with the line
bundles $\mathcal L$ and $\mathcal L'$ of their polarizations trivialized on open covers, and let $u\colon\mathbf C^g\to
P(\mathbf C)$ and
$u'\colon\mathbf C^{g'}\to P'(\mathbf C)$ be uniformizations for them (Definition~\ref{cv-origin-hodge-norm-gram}(1)) with
kernels $\Lambda$ and $\Lambda'$ and Appell--Humbert data $(H,\chi)$ and $(H',\chi')$, such that, for every $k\ge0$, every
theta function of weight $k$ for $u$ represents exactly one section of $\mathcal L^{\otimes k}$, and every section of
$\mathcal L'^{\otimes k}$ is represented by exactly one theta function of weight $k$ for $u'$
(Theorem~\ref{cv-cusp-cm-point-eigen-degree-rank-one-norm-riemann-theta}(4)). Let $F\colon\mathbf C^g\to\mathbf C^{g'}$ be
a $\mathbf C$-linear map and $n\ge0$ with $F(\Lambda)\subset\Lambda'$ and $F(z)^{\mathsf T}H'\,\overline{F(w)}=n\,z^{\mathsf
T}H\bar w$ for all $z,w$. Then there is a homomorphism $f\colon P\to P'$ with $u'\circ F=f\circ u$.
\end{lemma}
\leanref{Development/AbelianVarietyTheory/Proofs/CubicalLineBundles/ExistsHomOfLinearScalingHermForm.lean}{GlobalTheory.exists_hom_of_linear_scaling_hermForm}
\begin{proof}
Let $(U_\alpha)$ and $(U'_\beta)$ be the open covers on which $\mathcal L$ and $\mathcal L'$ are trivialized, and
$\theta_\alpha$ and $\theta'_\beta$ the trivializing functions of $u$ and $u'$. For a section $s$ of a power of $\mathcal L$
or of $\mathcal L'$, write $s_\alpha$ or $s_\beta$ for the function that represents it over $U_\alpha$ or $U'_\beta$; so a
theta function $\vartheta$ of weight $k$ for $u$ represents a section $s$ of $\mathcal L^{\otimes k}$ when
$s_\alpha(u(z))=\vartheta(z)\theta_\alpha(z)^k$ for all $\alpha$ and $z$ with $u(z)\in U_\alpha$. Put
$e_l(z)=\exp\big(\pi z^{\mathsf T}H\bar l+\frac\pi2l^{\mathsf T}H\bar l\big)$ for $l\in\Lambda$ and $z\in\mathbf C^g$.

Step 1: products of theta functions. By
Lemma~\ref{cv-cusp-cm-point-eigen-degree-tensor-complex-homs-sign}, $\chi(l)^2=1$ for $l\in\Lambda$ and $\chi'(l')^2=1$ for
$l'\in\Lambda'$. Let $N\ge0$ and let $\vartheta_1$ and $\vartheta_2$ be theta functions of weight $N$ for $u'$. Then
$\vartheta(z)=\vartheta_1(F(z))\vartheta_2(F(z))$ is a theta function of weight $2nN$ for $u$. It is holomorphic, as $F$ is
linear. For $l\in\Lambda$, $F(l)\in\Lambda'$ and $F(z)^{\mathsf T}H'\,\overline{F(l)}=n\,z^{\mathsf T}H\bar l$, so the
functional equations of $\vartheta_1$ and $\vartheta_2$ at $F(z)$ and $F(l)$ give
\[
 \vartheta(z+l)=\chi'(F(l))^{2N}e_l(z)^{2nN}\vartheta(z)=\big(\chi(l)e_l(z)\big)^{2nN}\vartheta(z),
\]
as $\chi'(F(l))^{2N}=1=\chi(l)^{2nN}$.

Step 2: sections. The line bundle $\mathcal L'$ is relatively ample by sections
(Definition~\ref{cv-origin-hodge-norm-form-ample}). For the affine open subscheme $\operatorname{Spec}\mathbf C$ of the
base, this gives $N\ge1$ and sections $s_0,\dots,s_M$ of $\mathcal L'^{\otimes N}$ without common zeros, such that each
open set $D(s_i)$ on which $s_i$ does not vanish is affine and its ring is generated over $\mathbf C$ by regular functions
$r_{i0},\dots,r_{iM}$ with $s_j=r_{ij}s_i$ on $D(s_i)$. By the hypothesis on $u'$, each $s_i$ is represented by a theta
function $\vartheta_i$ of weight $N$ for $u'$; by Step~1 and the hypothesis on $u$, the theta function
$z\mapsto\vartheta_i(F(z))\vartheta_j(F(z))$ represents a section $t_{ij}$ of $\mathcal L^{\otimes2nN}$. Let $w\in\mathbf
C^{g'}$ and $\beta$ with $u'(w)\in U'_\beta$. As $(s_i)_\beta(u'(w))=\vartheta_i(w)\theta'_\beta(w)^N$ and
$\theta'_\beta(w)\neq0$, the section $s_i$ does not vanish at $u'(w)$ if and only if $\vartheta_i(w)\neq0$, and then
$r_{ij}(u'(w))=\vartheta_j(w)/\vartheta_i(w)$.

Step 3: the local condition. The scheme $P$ is smooth over $\mathbf C$, hence reduced, and $u$ is onto the complex points of
$P$. Let $z\in\mathbf C^g$. As the $s_i$ have no common zeros, there is $i$ with $\vartheta_i(F(z))\neq0$, by Step~2. Let
$\alpha$ be an index with $u(z)\in U_\alpha$, and let $W\subset U_\alpha$ be the open set on which $(t_{ii})_\alpha$ does not
vanish. As $(t_{ii})_\alpha(u(w))=\vartheta_i(F(w))^2\theta_\alpha(w)^{2nN}$, a point $u(w)$ of $U_\alpha$ lies in $W$ if and
only if $\vartheta_i(F(w))\neq0$. So $u(z)\in W$, and $u'(F(w))\in D(s_i)$ whenever $u(w)\in W$, by Step~2. For each $j$
the function $b_j=(t_{ij})_\alpha/(t_{ii})_\alpha$ is regular on $W$, with
$b_j(u(w))=\vartheta_j(F(w))/\vartheta_i(F(w))=r_{ij}(u'(F(w)))$ for $u(w)\in W$. A regular function $f$ on $D(s_i)$ is a
polynomial $Q$ over $\mathbf C$ in the $r_{ij}$, and $h=Q(b_{0},\dots,b_{M})$ is a regular function on $W$ with
$h(u(w))=f(u'(F(w)))$ whenever $u(w)\in W$.

Step 4: the morphism. Lemma~\ref{cv-cusp-cm-point-eigen-degree-tensor-complex-homs-glue}, for $X=P$, $Y=P'$, which is
separated over $\operatorname{Spec}\mathbf C$, $Z=\mathbf C^g$, $e=u$ and $e'=u'\circ F$, applies by Step~3 and gives a
morphism of schemes $f\colon P\to P'$ with $f\circ u=u'\circ F$. The two morphisms from $P$ to $\operatorname{Spec}\mathbf
C$, $f$ followed by the structure morphism of $P'$ and the structure morphism of $P$, agree at every complex point $u(z)$,
as $u'(F(z))$ is a complex point of $P'$; so they are equal, by
Lemma~\ref{cv-cusp-cm-point-eigen-degree-tensor-complex-homs-points}(3), and $f$ is a morphism over $\mathbf C$.

Step 5: $f$ is a homomorphism. First $f(0)=f(u(0))=u'(F(0))=0$. The morphisms $(x,y)\mapsto f(x+y)$ and $(x,y)\mapsto
f(x)+f(y)$ from $P\times P$ to $P'$ agree at the points $(u(z),u(w))$, which are all the complex points of $P\times P$, as
$f(u(z)+u(w))=f(u(z+w))=u'(F(z)+F(w))=u'(F(z))+u'(F(w))$. The scheme $P\times P$ is smooth over $\mathbf C$, hence
reduced, and $P'$ is separated, so they are equal by
Lemma~\ref{cv-cusp-cm-point-eigen-degree-tensor-complex-homs-points}(3).
\end{proof}
\leanref{Development/AbelianVarietyTheory/Proofs/CubicalLineBundles/ExistsHomOfLinearScalingHermForm.lean}{GlobalTheory.AnalyticHom.exists_hom}
\leanref{Development/AbelianVarietyTheory/Proofs/CubicalLineBundles/ExistsHomOfLinearScalingHermForm.lean}{GlobalTheory.AnalyticHom.isThetaFn_comp}
% Lean: exists_hom_of_linear_scaling_hermForm: semicharacter_sq_eq_one_of_symmetric for u and u' (Lemma
% cv-cusp-cm-point-eigen-degree-tensor-complex-homs-sign) and AnalyticHom.exists_hom; then hw (Step 4, through
% ComplexPoints.hom_ext_of_points) and isMonHom_of_additive (Step 5, over_hom_ext_of_complexPoints with
% isReduced_of_smooth_complex for A x A). exists_hom: P'.ample at the affine open top (Step 2, with exists_sec_pow to move
% the sections from the inverse image of top to A'), the theta functions theta of hθ and t of hs with isThetaFn_comp
% (Step 1), hval and hmem (Step 2), and exists_hom_of_values (= SchemeHomsFromPoints.exists_hom, Lemma
% cv-cusp-cm-point-eigen-degree-tensor-complex-homs-glue) with W the basic open of t i i in the chart ch z and D = D(s_i);
% the functions h are built by MvPolynomial.induction_on, with exists_quot for b_j (Step 3). Only the existence halves of
% the hypotheses on theta functions are used.

\begin{lemma}\label{cv-cusp-cm-point-eigen-degree-tensor-complex-homs-bundle}
Let $P$, $P'$, $\mathcal L$, $\mathcal L'$, $u$, $u'$, $(H,\chi)$, $(H',\chi')$, $F$ and $n$ be as in
Lemma~\ref{cv-cusp-cm-point-eigen-degree-tensor-complex-homs-hom}, and let $f\colon P\to P'$ be a morphism with
$u'\circ F=f\circ u$. Then there is $N\ge1$ with $(f^*\mathcal L')^{\otimes N}\cong\mathcal L^{\otimes Nn}$.
\end{lemma}
\leanref{Development/AbelianVarietyTheory/Proofs/CubicalLineBundles/LbIsoPowOfHermFormPullEq.lean}{GlobalTheory.lbIso_pow_of_hermForm_pull_eq}
\begin{proof}
Let $(U_\alpha)$ and $(U'_\beta)$, with transition functions $g_{\alpha\alpha'}$ and $g'_{\beta\beta'}$, be the open covers on
which $\mathcal L$ and $\mathcal L'$ are trivialized, $\theta_\alpha$ and $\theta'_\beta$ the trivializing functions of $u$
and $u'$, and $s_\alpha$ and $s_\beta$ the functions that represent sections, as in the proof of
Lemma~\ref{cv-cusp-cm-point-eigen-degree-tensor-complex-homs-hom}.

Step 1: sections. As $\mathcal L'$ is relatively ample by sections (Definition~\ref{cv-origin-hodge-norm-form-ample}), for
the affine open subscheme $\operatorname{Spec}\mathbf C$ of the base there are $N_0\ge1$ and sections $s_0,\dots,s_m$ of
$\mathcal L'^{\otimes N_0}$ without common zeros. By the hypothesis on $u'$, each $s_i$ is represented by a theta function
$\vartheta_i$ of weight $N_0$ for $u'$. For every $w\in\mathbf C^{g'}$ some $s_i$ does not vanish at $u'(w)$, and then
$\vartheta_i(w)\neq0$, as $(s_i)_\beta(u'(w))=\vartheta_i(w)\theta'_\beta(w)^{N_0}$.

Step 2: squares. By Step~1 of the proof of Lemma~\ref{cv-cusp-cm-point-eigen-degree-tensor-complex-homs-hom}, which uses
Lemma~\ref{cv-cusp-cm-point-eigen-degree-tensor-complex-homs-sign}, the function $z\mapsto\vartheta_i(F(z))^2$ is a theta
function of weight $2N_0n$ for $u$. By the hypothesis on $u$, it represents a section $t_i$ of $\mathcal L^{\otimes2N_0n}$:
$(t_i)_\alpha(u(z))=\vartheta_i(F(z))^2\theta_\alpha(z)^{2N_0n}$ whenever $u(z)\in U_\alpha$.

Step 3: local ratios. Put $N=2N_0$. The scheme $P$ is smooth over $\mathbf C$, hence reduced, and $u$ is onto its complex
points. Fix $\alpha$ and $\beta$, and let $W=U_\alpha\cap f^{-1}(U'_\beta)$. The regular function $\sigma_i=(s_i)_\beta\circ f$
on $W$ has $\sigma_i(u(z))=\vartheta_i(F(z))\theta'_\beta(F(z))^{N_0}$ for $u(z)\in W$, as $f(u(z))=u'(F(z))$. On the open set
$W_i\subset W$ where $\sigma_i$ does not vanish, the regular function $(t_i)_\alpha/\sigma_i^2$ has the value
$\theta_\alpha(z)^{Nn}/\theta'_\beta(F(z))^N$ at $u(z)$, which does not depend on $i$. By
Lemma~\ref{cv-cusp-cm-point-eigen-degree-tensor-complex-homs-points}(2) these functions agree on $W_i\cap W_j$, so they glue
to a regular function $c_{\alpha\beta}$ on the union of the $W_i$. This union is $W$, by
Lemma~\ref{cv-cusp-cm-point-eigen-degree-tensor-complex-homs-points}(1), as every point $u(z)$ of $W$ lies in some $W_i$, by
Step~1. The values of $c_{\alpha\beta}$ at the complex points are nonzero, so $c_{\alpha\beta}$ is a unit, by
Lemma~\ref{cv-cusp-cm-point-eigen-degree-tensor-complex-homs-points}(2).

Step 4: the isomorphism. As $\theta_\alpha=(g_{\alpha\alpha'}\circ u)\,\theta_{\alpha'}$ and
$\theta'_\beta(F(z))=g'_{\beta\beta'}(f(u(z)))\,\theta'_{\beta'}(F(z))$, the two sides of
\[
 c_{\alpha\beta}\,(g'_{\beta\beta'}\circ f)^N=g_{\alpha\alpha'}^{Nn}\,c_{\alpha'\beta'}
\]
have the same value at every complex point where both are defined, so the identity holds, by
Lemma~\ref{cv-cusp-cm-point-eigen-degree-tensor-complex-homs-points}(2). The line bundle $(f^*\mathcal L')^{\otimes N}$ is
trivialized on the $f^{-1}(U'_\beta)$ with transition functions $(g'_{\beta\beta'}\circ f)^N$, and $\mathcal
L^{\otimes Nn}$ on the $U_\alpha$ with transition functions $g_{\alpha\alpha'}^{Nn}$. So multiplying the function that
represents a section of $(f^*\mathcal L')^{\otimes N}$ over $f^{-1}(U'_\beta)$ by $c_{\alpha\beta}$, on $U_\alpha\cap
f^{-1}(U'_\beta)$, gives a well defined isomorphism $(f^*\mathcal L')^{\otimes N}\to\mathcal L^{\otimes Nn}$, with inverse
given by the $c_{\alpha\beta}^{-1}$.
\end{proof}
\leanref{Development/AbelianVarietyTheory/Proofs/CubicalLineBundles/LbIsoPowOfHermFormPullEq.lean}{GlobalTheory.PullIso.lbIso_of_frames}
\leanref{Development/AbelianVarietyTheory/Proofs/CubicalLineBundles/LbIsoPowOfHermFormPullEq.lean}{GlobalTheory.PullIso.exists_of_local_quotient}
% Lean: lbIso_pow_of_hermForm_pull_eq: exists_generating_sections (from P'.ample at the affine open top, Step 1), theta
% from hfn and hθ0 (Step 1), hT with semicharacter_sq_eq_one_of_symmetric (Step 2), t from hsec; lbIso_of_frames for the
% frames fr_J^(2Nn) of L^(2Nn) and fr'_K(F z)^(2N) of the pullback (Step 4, the units c and their compatibility checked by
% eq_of_valAt'_eq), with the functions c_{JK} from exists_of_local_quotient (Step 3: a_i = t_i in the chart J and
% b_i = sigma_i^2 with sigma_i the pullback of s_i, glued over the basic opens of the b_i). Only the existence halves of
% the hypotheses on theta functions are used.

\begin{lemma}\label{cv-cusp-cm-point-eigen-degree-tensor-complex-homs}
In Situation~\ref{cv-cusp-setting}, let $\sqrt{-d}\in E$ be the square root of $-d$ of the datum, and put
$w_0=(1+\sqrt{-d})/2$. Let $A$ be a principally polarized abelian variety over $\mathbf C$ of dimension $h$ with an
action of $\mathcal O_E$ for which the Rosati involution of the polarization induces complex conjugation on
$\mathcal O_E$, with the line bundle $\mathcal L_A$ of its polarization, and let $u\colon\mathbf C^h\to
A(\mathbf C)$ be a uniformization of $A$ for $\mathcal L_A$ (Definition~\ref{cv-origin-hodge-norm-gram}(1)) with kernel
$\Lambda$ and Appell--Humbert data $(H,\chi)$, such that every $z\in\mathbf C^h$ is the differential of exactly one
tangent vector of $A$ at the origin and every section of $\mathcal L_A^{\otimes n}$, for $n\ge0$, is represented by
exactly one theta function of weight $n$ (Theorem~\ref{cv-cusp-cm-point-eigen-degree-rank-one-norm-riemann-theta}(4)).
Let $L$ be the lattice of a PEL datum over $E$ with its perfect alternating form $\langle\ ,\ \rangle$, let
$\sigma_1,\dots,\sigma_g$ be embeddings $E\to\mathbf C$, let $P\colon L\to\mathbf C^g$ be a period of $L$ for them
(Definition~\ref{cv-cm-point-periods}(3)), and let $F\colon\mathbf C^g\to\mathbf C^h$ be a $\mathbf C$-linear bijection
with (1), (2) and (3) of Lemma~\ref{cv-cusp-cm-point-eigen-degree-tensor-complex-datum}. Let $C$ be a principally
polarized elliptic curve over $\mathbf C$ with the line bundle $\mathcal L_C$ of its polarization, and
$u_0\colon\mathbf C\to C(\mathbf C)$ a uniformization of $C$ for $\mathcal L_C$ with Appell--Humbert data for which (2),
(4) and (5) of Theorem~\ref{cv-cusp-cm-point-eigen-degree-rank-one-norm-riemann-theta} hold, whose kernel is
$c(\mathbf Z+\mathbf Z\frac{1+i\sqrt d}2)$ with $c\neq0$. Let $\rho\in\mathbf C$ with $\rho^2=-d$, and let $M\subset L$
be an $\mathcal O_E$-submodule such that $P(a)_k=0$ for every $a\in M$ and every $k$ with
$\sigma_k(\sqrt{-d})\neq\rho$. Then there are homomorphisms $t_a\colon C\to A$, for $a\in M$, with
$t_a(u_0(cz))=u(F(zP(a)))$ for all $a\in M$ and $z\in\mathbf C$, such that for every tangent vector $v$ of $C$ at the
origin:
\begin{enumerate}
\item $t_{a+b}=t_a+t_b$ and $t_{sa}=s\circ t_a$ for all $a,b\in M$ and $s\in\mathcal O_E$;
\item $\sqrt{-d}\circ t_a\circ v=\rho\,(t_a\circ v)$ in the Lie algebra of $A$, for every $a\in M$;
\item for every $a\in M$ there is $N\ge1$ with $(t_a^*\mathcal L_A)^{\otimes N}\cong\mathcal L_C^{\otimes
  N|\langle w_0a,a\rangle|}$.
\end{enumerate}
\end{lemma}
\leanref{Development/AbelianVarietyTheory/Proofs/CubicalLineBundles/ExistsHomFamilyOfPureTypeComplexTriple.lean}{GlobalTheory.CM.PureTriple.exists_isHomFamily}
\begin{proof}
Let $H_0>0$ be the Appell--Humbert form of $u_0$, a positive real number as $C$ has dimension $1$, and let $\Lambda_0$ be
the kernel of $u_0$. For $a\in M$ let $G_a\colon\mathbf C\to\mathbf C^h$ be the $\mathbf C$-linear map
$x\mapsto F(c^{-1}xP(a))$, so that $G_a(cz)=F(zP(a))$.

Step 1: lattices and forms. Let $a\in M$. The form $(z,w)\mapsto F(z)^{\mathsf T}H\,\overline{F(w)}$ on $\mathbf C^g$ is
positive definite and Hermitian, as $F$ is bijective, and its imaginary part on $P(L)$ is $\langle\ ,\ \rangle$ by (2) of
Lemma~\ref{cv-cusp-cm-point-eigen-degree-tensor-complex-datum}. By
Lemma~\ref{cv-cusp-cm-point-eigen-degree-tensor-complex-homs-line}(1) for this form, $(m+n\frac{1+i\sqrt d}2)P(a)\in P(L)$
for $m,n\in\mathbf Z$, so by (1) of Lemma~\ref{cv-cusp-cm-point-eigen-degree-tensor-complex-datum}, $G_a(\Lambda_0)\subset
\Lambda$. By Lemma~\ref{cv-cusp-cm-point-eigen-degree-tensor-complex-homs-line}(2), $F(zP(a))^{\mathsf
T}H\,\overline{F(z'P(a))}=|\langle w_0a,a\rangle|\cdot2z\bar z'/\sqrt d$, and by
Lemma~\ref{cv-cusp-cm-point-eigen-degree-tensor-complex-homs-curve}, which applies by (5) of
Theorem~\ref{cv-cusp-cm-point-eigen-degree-rank-one-norm-riemann-theta} for $u_0$, $(cz)H_0\overline{(cz')}=2z\bar
z'/\sqrt d$. So $G_a(x)^{\mathsf T}H\,\overline{G_a(y)}=|\langle w_0a,a\rangle|\,xH_0\bar y$ for all $x,y\in\mathbf C$.

Step 2: the homomorphisms. For $a\in M$,
Lemma~\ref{cv-cusp-cm-point-eigen-degree-tensor-complex-homs-hom}, for $C$, $A$, $u_0$, $u$, $G_a$ and
$n=|\langle w_0a,a\rangle|$, by Step~1 and (4) of
Theorem~\ref{cv-cusp-cm-point-eigen-degree-rank-one-norm-riemann-theta} for $u_0$, gives a homomorphism $t_a\colon C\to A$
with $u\circ G_a=t_a\circ u_0$, that is, $t_a(u_0(cz))=u(F(zP(a)))$ for every $z\in\mathbf C$.

Step 3: points. Two morphisms $C\to A$ that agree at the points $u_0(cz)$, $z\in\mathbf C$, are equal: these are all the
complex points of $C$, as $u_0$ is onto and $c\neq0$, $C$ is smooth over $\mathbf C$, hence reduced, and $A$ is separated.

Step 4: condition (1). For $a,b\in M$ and $z\in\mathbf C$, as $u$ is a homomorphism and $P$ and $F$ are additive,
$t_{a+b}(u_0(cz))=u(F(zP(a)))+u(F(zP(b)))=(t_a+t_b)(u_0(cz))$, so $t_{a+b}=t_a+t_b$ by Step~3. For $s\in\mathcal O_E$,
$zP(sa)=(\sigma_k(s)zP(a)_k)_k$, so by (3) of Lemma~\ref{cv-cusp-cm-point-eigen-degree-tensor-complex-datum},
$t_{sa}(u_0(cz))=s\big(u(F(zP(a)))\big)=s\big(t_a(u_0(cz))\big)$, and $t_{sa}=s\circ t_a$ by Step~3.

Step 5: condition (2). Let $v$ be a tangent vector of $C$ at the origin, with differential $cz$ for $u_0$, by (2) of
Theorem~\ref{cv-cusp-cm-point-eigen-degree-rank-one-norm-riemann-theta}. By the chain rule, as in the proof of
Lemma~\ref{cv-origin-hodge-norm-form-compare}, $t_b\circ v$ has differential $G_b(cz)=F(zP(b))$ for $u$, for every
$b\in M$. As $P(\sqrt{-d}\,a)_k=\sigma_k(\sqrt{-d})P(a)_k=\rho P(a)_k$ for every $k$, $t_{\sqrt{-d}a}\circ v$ has
differential $\rho F(zP(a))$, which is also the differential of $\rho\,(t_a\circ v)$. By the uniqueness in the
hypothesis on $u$, $t_{\sqrt{-d}a}\circ v=\rho\,(t_a\circ v)$, and $t_{\sqrt{-d}a}=\sqrt{-d}\circ t_a$ by Step~4.

Step 6: condition (3). For $a\in M$,
Lemma~\ref{cv-cusp-cm-point-eigen-degree-tensor-complex-homs-bundle}, for the data of Step~2 and $f=t_a$, gives $N\ge1$
with $(t_a^*\mathcal L_A)^{\otimes N}\cong\mathcal L_C^{\otimes N|\langle w_0a,a\rangle|}$.
\end{proof}
% Lean: exists_isHomFamily: hu (IsThetaUnif u_0) for H_0, chi_0, the tangent bijection hD_0, the theta functions
% hth_0, E_0 and its bijectivity hb; G a = F o toSpanSingleton (Pe a) o (c^{-1} proj i_0) and hG (Step 0); hall (every
% vector of C^1 is c z); hline = PureType.hermForm_line, hGLambda and hGH with hermForm_of_isQuadLattice (Step 1);
% exists_hom_of_linear_scaling_hermForm (Step 2); ext with ComplexPoints.over_hom_ext_of_complexPoints (Step 3); the
% additivity and hlin with hact (Step 4); HomLinear.isDiff_comp, period_rt_smul, isTangent_scaleVec,
% CMCurveTest.isDiff_scaleVec and the uniqueness hD (Step 5); and lbIso_pow_of_hermForm_pull_eq (Step 6). What remains
% open is in Lemmas cv-cusp-cm-point-eigen-degree-tensor-complex-homs-hom (about 400 lines) and -homs-bundle (about
% 250).

\begin{lemma}\label{cv-cusp-cm-point-eigen-degree-tensor-complex-pairing-polarize}
In Situation~\ref{cv-cusp-setting}, let $\sqrt{-d}\in E$ be the square root of $-d$ of the datum, and put
$w_0=(1+\sqrt{-d})/2$. Let $L$ be the lattice of a PEL datum over $E$ with its perfect alternating form
$\langle\ ,\ \rangle$, $M\subset L$ an $\mathcal O_E$-submodule, $H\in M_g(\mathbf C)$ a matrix, and $\omega,c\in\mathbf C$
with $(2\omega-1)^2=-d$. Let $x\colon M\to\mathbf C^g$ be an additive map with $x(\sqrt{-d}\,a)=(2\omega-1)x(a)$ and
$x(a)^{\mathsf T}H\,\overline{x(a)}=\langle w_0a,a\rangle c$ for every $a\in M$. Then
$x(a)^{\mathsf T}H\,\overline{x(b)}=\big(\langle w_0a,b\rangle-\langle a,b\rangle+\langle a,b\rangle\omega\big)c$ for all
$a,b\in M$.
\end{lemma}
\leanref{Development/AbelianVarietyTheory/Proofs/CubicalLineBundles/IsCMPairingFamilyOfPolarisationDiagonalComplex.lean}{GlobalTheory.CM.FamPairing.polarise}
\begin{proof}
Write $B(z,w)=z^{\mathsf T}H\bar w$, which is linear in $z$ and conjugate linear in $w$, and $\theta=2\omega-1$. As
$\theta^2=-d<0$, $\theta$ is purely imaginary and nonzero, so $\bar\theta=-\theta$.

The form. For $u,v\in L$ and $s\in\mathcal O_E$, $\langle u,v\rangle=-\langle v,u\rangle$ and $\langle su,v\rangle=\langle
u,\bar sv\rangle$ (Definition~\ref{cv-cm-point-periods}), with $\bar w_0=1-w_0$, $\sqrt{-d}=2w_0-1$ and
$4\bar w_0w_0=1+d$. Fix $a,b\in M$ and put $\alpha=\langle w_0a,b\rangle$, $e=\langle a,b\rangle$ and
$\mu=\langle w_0a,w_0b\rangle$. Then $\langle a,w_0b\rangle=\langle(1-w_0)a,b\rangle=e-\alpha$, and:
\begin{enumerate}
\item[(i)] $\langle w_0b,a\rangle=-\langle a,w_0b\rangle=\alpha-e$;
\item[(ii)] $\langle w_0a,\sqrt{-d}\,b\rangle=2\mu-\alpha$;
\item[(iii)] $\langle w_0\sqrt{-d}\,b,a\rangle=\langle w_0a,\sqrt{-d}\,b\rangle-\langle a,\sqrt{-d}\,b\rangle=(2\mu-\alpha)-(2(e-\alpha)-e)
  =2\mu+\alpha-e$, by (i) for $\sqrt{-d}\,b$ in place of $b$;
\item[(iv)] $4\mu=\langle a,4\bar w_0w_0b\rangle=(1+d)e$.
\end{enumerate}

The diagonal. For $p,q\in M$, expanding $B(x(p+q),x(p+q))=\langle w_0(p+q),p+q\rangle c$ by the additivity of $x$ and
subtracting the values at $p$ and at $q$ gives
\[
 B(x(p),x(q))+B(x(q),x(p))=\big(\langle w_0p,q\rangle+\langle w_0q,p\rangle\big)c.\tag{S}
\]
For $(p,q)=(a,b)$, by (i), $B(x(a),x(b))+B(x(b),x(a))=(2\alpha-e)c$. For $(p,q)=(a,\sqrt{-d}\,b)$, as $x(\sqrt{-d}\,b)=\theta
x(b)$ and $\bar\theta=-\theta$, (ii), (iii) and (iv) give $-\theta B(x(a),x(b))+\theta B(x(b),x(a))=(4\mu-e)c=de\,c$.
Multiplying the first identity by $\theta$ and subtracting the second gives
$2\theta B(x(a),x(b))=\theta(2\alpha-e)c-de\,c=\theta(2\alpha-e+\theta e)c=2\theta(\alpha-e+e\omega)c$, as $-d=\theta^2$.
Dividing by $2\theta\neq0$ gives the claim.
\end{proof}
% Lean: polarise: EigenNorm.conj_eq_neg_of_sq (conj theta = -theta); S1 (the identity (S), from hq at p + q); Sab and
% Sarb; g1 to g4 (ef_w0_smul_left_swap, ef_rt_smul_right and four_mul_ef_w0_smul_w0_smul, through D.adj and
% cmHost_bar_halfIntegerGenerator); key and hr0.

\begin{lemma}\label{cv-cusp-cm-point-eigen-degree-tensor-complex-pairing-lambda-fibre}
Assume laws (7) and (21) of Definition~\ref{cv-cusp-cm-point-eigen-degree-tensor-laws}. Let $k$ be a field, $P$ a
principally polarized abelian variety over $k$ with the line bundle $\mathcal L$ of its polarization, $X$ and $Y$ group
schemes over $k$, with their group laws written multiplicatively, and $f_1\colon X\to P$ and $f_2\colon Y\to P$
homomorphisms. Let $n\ge0$ and $x\in X(k)$ with $x^n=1$, let $j\colon Y\to X\times Y$ be the morphism $w\mapsto(x,w^n)$, and
for $y\in Y(k)$ let $T_y$ be the translation $w\mapsto yw$ of $Y$. Then $j^*\Lambda_{\mathcal L}(f_1p_1,f_2p_2)$, with $p_1,p_2$ the projections of $X\times Y$
(Definition~\ref{cv-toroidal-abelian-dvr-ample-torsion-weil}), has a generator $u$ such that $T_y^*u=c\,u$ for every
$y\in Y(k)$ with $y^n=1$ and every $c\in k^\times$ with $e_n(f_1x,f_2y)=c$, for the commutator pairing $e_n$ of the theta
group of $\mathcal L^n$ (Definition~\ref{cv-cusp-cm-point-eigen-degree-rank-one-norm-weil-translate-group}).
\end{lemma}
\leanref{Development/AbelianVarietyTheory/Proofs/CubicalLineBundles/LambdaBundleIsoOfWeilPairingCompatible.lean}{GlobalTheory.CM.LamOfWeil.exists_gen_fibre}
\begin{proof}
As $f_2$ is a homomorphism to the commutative group scheme $P$ and $y^n=1$,
$f_2((yw)^n)=n\,f_2(y)+n\,f_2(w)=f_2(w^n)$; so $T_y^*j^*\Lambda_{\mathcal L}(f_1p_1,f_2p_2)$ is $j^*\Lambda_{\mathcal
L}(f_1p_1,f_2p_2)$, and the identity of the statement compares two generators of one line bundle. Let
$T=\operatorname{Spec}k$, which has one point.
The line bundle $\mathcal L$ is trivialized along the zero section, and by law (7) it has a cubical structure. As $n\,f_1x=0$,
Lemma~\ref{cv-cusp-cm-point-eigen-degree-rank-one-norm-weil-translate-pairing}, for $P$, $\mathcal L$, $T$, $n$ and the
point $f_1x$, gives a generator $u_0$ of $(f_1x\,p_T+n\,p_P)^*\mathcal L\otimes(n\,p_P)^*\mathcal L^{-1}$ on $P$ with the
translation property of that lemma. As $T$ has one point, $(f_1x)^*\mathcal L^{-1}\otimes0^*\mathcal L$ has a generator
$\sigma$. Put $u=f_2^*u_0\otimes\sigma$. As $f_2(w^n)=n\,f_2(w)$, $u$ is a generator of
\[
 (f_1x+n f_2)^*\mathcal L\otimes(nf_2)^*\mathcal L^{-1}\otimes(f_1x)^*\mathcal L^{-1}\otimes0^*\mathcal L
 =j^*\Lambda_{\mathcal L}(f_1p_1,f_2p_2).
\]
Let $y\in Y(k)$ with $y^n=1$ and $c$ with $e_n(f_1x,f_2y)=c$. Then $n\,f_2y=0$, and that lemma, for the point $f_2y$, for
$Z=Y$ with its structure morphism and for $w=f_2$, gives $u_0(f_2y+f_2)=c\,u_0(f_2)$, that is,
$T_y^*f_2^*u_0=c\,f_2^*u_0$, as $f_2\circ T_y=f_2y+f_2$. As $\sigma$ is pulled back from $T$, $T_y^*u=c\,u$.
\end{proof}
% Lean: exists_gen_fibre (with GlobalFnsSt and CubeSt): the cubical structure cb from hC; PairingDetermines.exists_gen_transl
% for the point x f_1, a copy of AdjointOfPairing.exists_gen_pairing (Lemma
% cv-cusp-cm-point-eigen-degree-rank-one-norm-weil-translate-pairing) with the same statement and proof, which uses
% GlobalFnsSt; exists_frame for sigma; exists_gen for u (the line bundles compared through prod_sq); and w, w_0 for the
% translation identity, compared in charts at every point of Y.

\begin{lemma}\label{cv-cusp-cm-point-eigen-degree-tensor-complex-pairing-lambda-quotient}
Assume law (6) of Definition~\ref{cv-cusp-cm-point-eigen-degree-tensor-laws}. Let $K$ be an algebraically closed field,
$A$ a principally polarized abelian variety over $K$, $n\ge0$ with $n\neq0$ in $K$, and $U\subset A$ an open subscheme.
Let $f$ be a regular function on $[n]^{-1}(U)$ such that $f\circ T_y=f$ on $[n]^{-1}(U)$ for every $y\in A(K)$ with $ny=0$,
where $T_y$ is the translation by $y$. Then there is exactly one regular function $g$ on $U$ with $g\circ[n]=f$.
\end{lemma}
\leanref{Development/AbelianVarietyTheory/Proofs/AbelianSchemeLawProofs/ExistsUniqueDescendOfTorsionInvariant.lean}{GlobalTheory.torsion_invariant_descends}
\leanref{Development/AbelianVarietyTheory/Proofs/AbelianSchemeLawProofs/ExistsUniqueDescendOfTorsionInvariant.lean}{GlobalTheory.existsUnique_descend_of_torsion_invariant}
\begin{proof}
As $n\neq0$ in $K$, $n\ge1$ and $n$ is a unit on $\operatorname{Spec}K$, so by law (6) the morphism $[n]$ is finite, flat,
surjective and \'etale. Let $V=[n]^{-1}(U)$ and $m\colon V\to U$ the restriction of $[n]$.

Step 1: equal images. Let $Z$ be a scheme and $a,b\colon Z\to V$ morphisms with $m\circ a=m\circ b$. Then $f\circ a=f\circ b$.
Indeed, give $Z$ the structure of a scheme over $K$ through $a$ and the structure morphism of $A$; the point $d=a-b$ of $A$
with values in $Z$ has $nd=0$, so it factors through the kernel of $[n]$, which is finite and \'etale over the
algebraically closed field $K$, hence a disjoint union of copies of $\operatorname{Spec}K$, one for each $y\in A(K)$ with
$ny=0$. So $Z$ is covered by open subsets on which $d$ is constant with a value $y$ of this kind, and there $a=T_y\circ b$,
so $f\circ a=f\circ T_y\circ b=f\circ b$.

Step 2: descent. Let $U'\subset U$ be an affine open subscheme, with ring $R'$. As $m$ is finite, $m^{-1}(U')$ is affine, with
ring $B'$, and $R'\to B'$ is flat, and faithfully flat as $m$ is surjective. By Step~1 for the two projections of
$m^{-1}(U')\times_{U'}m^{-1}(U')$, the restriction of $f$ to $m^{-1}(U')$ has the same images under the two maps $B'\to
B'\otimes_{R'}B'$, so it is the image of exactly one element of $R'$ \cite[Tag 023M]{Sta}. By this uniqueness, these
elements, for the affine open subschemes $U'$ of $U$, agree on overlaps and glue to a regular function $g$ on $U$ with
$g\circ m=f$; and $g$ is unique, as $R'\to B'$ is injective for every $U'$.
\end{proof}
% Lean: torsion_invariant_descends: H (the law MulNSt) for IsFinite, Surjective and Etale, and
% existsUnique_descend_of_torsion_invariant: key (Step 1, with exists_local_rational_torsion for the cover of Z on which
% a - b is a rational n-torsion point y_p, and TLinPf.comp_mOf with hf) and existsUnique_appTop_eq_of_flat for the flat,
% surjective, quasi-compact restriction of [n] (Step 2; the text proves it on affine opens by the exactness of the Amitsur
% complex).

\begin{lemma}\label{cv-cusp-cm-point-eigen-degree-tensor-complex-pairing-lambda-descent}
Let $Y$ and $Y'$ be schemes, $\mathcal M$ and $\mathcal M'$ line bundles on $Y$, $m\colon Y'\to Y$ a morphism, and
$(t_i)_{i\in I}$ morphisms $Y'\to Y'$ with $m\circ t_i=m$. Suppose that for every open subscheme $U\subset Y$ and every
regular function $f$ on $m^{-1}(U)$ with $f\circ t_i=f$ on $m^{-1}(U)$ for every $i$, there is exactly one regular function
$g$ on $U$ with $g\circ m=f$. Let $u$ and $u'$ be generators of $m^*\mathcal M$ and $m^*\mathcal M'$ such that
$t_i^*(u'\otimes u^{-1})=u'\otimes u^{-1}$ for every $i$, where $t_i^*m^*=m^*$. Then $\mathcal M\cong\mathcal M'$.
\end{lemma}
\leanref{Development/AbelianVarietyTheory/Proofs/ModuliProofs/ExistsGenOfInvariant.lean}{GlobalTheory.IsoDescent.lbIso_of_invariant}
\leanref{Development/AbelianVarietyTheory/Proofs/ModuliProofs/ExistsGenOfInvariant.lean}{GlobalTheory.IsoDescent.exists_gen_of_invariant}
\begin{proof}
Put $\mathcal N=\mathcal M'\otimes\mathcal M^{-1}$ and $\rho=u'\otimes u^{-1}$, a generator of $m^*\mathcal N$ with
$t_i^*\rho=\rho$. Let $(U_j)$ be an open cover of $Y$ on which $\mathcal N$ is trivialized, with transition functions
$h_{jj'}$. On $m^{-1}(U_j)$ the generator $\rho$ is represented by a unit $a_j$, with $a_j=(h_{jj'}\circ m)\,a_{j'}$ on
overlaps. As the $h_{jj'}\circ m$ are invariant under the $t_i$, so is each $a_j$, and so are $a_j^{-1}$. By the hypothesis,
$a_j=b_j\circ m$ and $a_j^{-1}=b'_j\circ m$ for regular functions $b_j$, $b'_j$ on $U_j$; the hypothesis also shows that
$g\mapsto g\circ m$ is injective, so $b_jb'_j=1$ and $b_j=h_{jj'}b_{j'}$ on $U_j\cap U_{j'}$. So the units $b_j$ represent a
generator of $\mathcal N$, and $\mathcal M\cong\mathcal M'$.
\end{proof}
% Lean: lbIso_of_invariant, through exists_gen_of_invariant for LB.div M' M with a_j = u'/u in the charts (hc for the
% transition rule, hinv for the invariance, through gen_trans); hinj, hunit and hex for the b_j.

\begin{lemma}\label{cv-cusp-cm-point-eigen-degree-tensor-complex-pairing-lambda-dense}
Assume laws (4) and (13) of Definition~\ref{cv-cusp-cm-point-eigen-degree-tensor-laws}. Let $K$ be an algebraically closed
field of characteristic $0$ and $A$ a principally polarized abelian variety over $K$. Then the points of $A(K)$ that are
killed by some $n\ge1$ are dense in $A$.
\end{lemma}
\leanref{Development/AbelianVarietyTheory/Proofs/ModuliProofs/ExistsMemNonzeroLocus.lean}{GlobalTheory.TorsionDense.dense_torsion}
\begin{proof}
\openstep{The density of the torsion points, from the irreducibility of the action of the theta group on sections of a
power of the polarization and from the theta group of torsion points, is not yet written out.}
% Lean: TorsionDense.dense_torsion (with IrredSt and KerPowSt): isAmpleLB_of_relAmple_of_isAffine, exists_affine_le for
% an affine nonvanishing locus D(s) inside a given open set, exists_affine at the origin for t, nonzeroLocus_secPow, and
% exists_mem_nonzeroLocus (a torsion point b with s^(4m) b, through the irreducibility law and the theta group of
% L^(n m 4) at torsion points). Size (tools/lean_size.py): 592 lines.
\end{proof}

\begin{lemma}\label{cv-cusp-cm-point-eigen-degree-tensor-complex-pairing-lambda}
Assume the laws of abelian schemes (Definition~\ref{cv-cusp-cm-point-eigen-degree-tensor-laws}).
Let $k$ be an algebraically closed field of characteristic $0$, let $X$, $Y$, $A$ and $B$ be principally polarized
abelian varieties over $k$, let $\mathcal L_A$ and $\mathcal L_B$ be the line bundles of the polarizations of $A$ and
$B$, and for $n\ge1$ let $e^A_n$ and $e^B_n$ be the commutator pairings of the theta groups of $\mathcal L_A^n$ and
$\mathcal L_B^n$. Let $f_1\colon X\to A$, $f_2\colon Y\to A$, $g_1\colon X\to B$ and $g_2\colon Y\to B$ be
homomorphisms with $e^A_n(f_1x,f_2y)=e^B_n(g_1x,g_2y)$ for all $n\ge1$ and all $x\in X(k)$, $y\in Y(k)$ with $nx=ny=0$.
Then $\Lambda_{\mathcal L_A}(f_1p_1,f_2p_2)\cong\Lambda_{\mathcal L_B}(g_1p_1,g_2p_2)$ on $X\times Y$, with $p_1$ and
$p_2$ the projections (Definition~\ref{cv-toroidal-abelian-dvr-ample-torsion-weil}).
\end{lemma}
\leanref{Development/AbelianVarietyTheory/Proofs/CubicalLineBundles/LambdaBundleIsoOfWeilPairingCompatibleStatement.lean}{GlobalTheory.T.lambdaBundleIso_of_weilPairing_compatible}
\leanref{Development/AbelianVarietyTheory/Proofs/CubicalLineBundles/LambdaBundleIsoOfWeilPairingCompatible.lean}{GlobalTheory.CM.lamEq_of_weilEq}
\begin{proof}
Let $\mathcal D=\Lambda_{\mathcal L_B}(g_1p_1,g_2p_2)\otimes\Lambda_{\mathcal L_A}(f_1p_1,f_2p_2)^{-1}$ on $X\times Y$; the
claim is that $\mathcal D$ is trivial. For $x\in X(k)$ let $j_x\colon Y\to X\times Y$ be the morphism $y\mapsto(x,y)$, and
let $\sigma=(1,0)\colon X\to X\times Y$.

Step 1: along $\sigma$. The pullback $\sigma^*\Lambda_{\mathcal L_A}(f_1p_1,f_2p_2)=\Lambda_{\mathcal L_A}(f_1,0)$ is
$f_1^*\mathcal L_A\otimes f_1^*\mathcal L_A^{-1}\otimes0^*\mathcal L_A^{-1}\otimes0^*\mathcal L_A$, which is trivial; in the
same way $\sigma^*\Lambda_{\mathcal L_B}(g_1p_1,g_2p_2)$ is trivial. So $\sigma^*\mathcal D$ is trivial.

Step 2: the fibres over torsion points. Let $x\in X(k)$ and $n\ge1$ with $nx=0$. As $j_x\circ[n]$ is the morphism
$y\mapsto(x,ny)$, Lemma~\ref{cv-cusp-cm-point-eigen-degree-tensor-complex-pairing-lambda-fibre}, for $A$, $f_1$, $f_2$ and for
$B$, $g_1$, $g_2$, which applies by laws (7) and (21), gives generators $u$ of $[n]^*j_x^*\Lambda_{\mathcal L_A}(f_1p_1,f_2p_2)$
and $u'$ of $[n]^*j_x^*\Lambda_{\mathcal L_B}(g_1p_1,g_2p_2)$ with $T_y^*u=c\,u$ and $T_y^*u'=c'\,u'$ whenever $y\in Y(k)$, $ny=0$,
$e^A_n(f_1x,f_2y)=c$ and $e^B_n(g_1x,g_2y)=c'$. For such $y$, as $n\,f_1x=n\,f_2y=0$, law (10)(a) gives a unit $c$ with
$e^A_n(f_1x,f_2y)=c$, and by hypothesis $e^B_n(g_1x,g_2y)=c$. So $T_y^*(u'\otimes u^{-1})=u'\otimes u^{-1}$. As $k$ has
characteristic $0$, $n\neq0$ in $k$, so by Lemma~\ref{cv-cusp-cm-point-eigen-degree-tensor-complex-pairing-lambda-quotient},
which applies by law (6), the morphism $m=[n]$ of $Y$ and the translations $T_y$, for $y\in Y(k)$ with $ny=0$, satisfy the
hypothesis of Lemma~\ref{cv-cusp-cm-point-eigen-degree-tensor-complex-pairing-lambda-descent}. That lemma gives
$j_x^*\Lambda_{\mathcal L_A}(f_1p_1,f_2p_2)\cong j_x^*\Lambda_{\mathcal L_B}(g_1p_1,g_2p_2)$, so $j_x^*\mathcal D$ is trivial.

Step 3: density. By Lemma~\ref{cv-cusp-cm-point-eigen-degree-tensor-complex-pairing-lambda-dense}, which applies by laws (4)
and (13), the points $x\in X(k)$ killed by some $n\ge1$ are dense in $X$.

Step 4: seesaw. The scheme $X$ is smooth over $k$, hence reduced. Apply
Lemma~\ref{cv-toroidal-theta-family-pel-bundles-seesaw} with $S_0=\operatorname{Spec}k$, the abelian scheme $Y$ over $k$ with
the relatively ample line bundle $\mathcal L_Y$ of its polarization, the reduced scheme $S=X$, so that the product there is
$X\times Y$ with its projection to $X$, the line bundle $\mathcal D$ and the section $\sigma$. For a point $x\in X(k)$, the
fibre of $X\times Y$ over $x$ is $Y$, embedded by $j_x$; so by Steps 2 and 3 the points of $X$ over which the pullback of
$\mathcal D$ to the fibre is trivial are dense. By Step~1, $\sigma^*\mathcal D$ is trivial. So $\mathcal D$ is trivial, that is,
$\Lambda_{\mathcal L_A}(f_1p_1,f_2p_2)\cong\Lambda_{\mathcal L_B}(g_1p_1,g_2p_2)$.
\end{proof}
\leanref{Development/AbelianVarietyTheory/Proofs/ModuliProofs/SeesawTrivialityOnDenseFibresProductAbelian.lean}{GlobalTheory.triv_of_triv_on_dense_fibres_prod_abelian}
% Lean: lamEq_of_weilEq: axis (exists_gen_axis, through AdjointOfPairing.exists_gen_of_formEq_zero) for Step 1, at the
% section lift (id, 1); the triviality on 0 x Y that the Lean also proves is an unused hypothesis of
% triv_of_triv_on_dense_fibres_prod_abelian. Step 2 is exists_gen_fibre with hL.globalFns and hL.cube (Lemma
% cv-cusp-cm-point-eigen-degree-tensor-complex-pairing-lambda-fibre), (hL.weil _ A n _).1 for c, hW, and lbIso_of_invariant
% with torsion_invariant_descends and hL.mulN; Step 3 is TorsionDense.dense_torsion with hL.irred and hL.kerPow; Step 4 is
% T.seesaw_triviality_on_dense_fibres_product_abelian = triv_of_triv_on_dense_fibres_prod_abelian, which applies
% T.seesaw_reduced_base_abelianScheme_of_dense_fibres (Lemma cv-toroidal-theta-family-pel-bundles-seesaw) with S := X,
% the pullback square isPullback_lift_toUnit_fst_left for the fibres and hL.globalFns. What remains open is in Lemma
% cv-cusp-cm-point-eigen-degree-tensor-complex-pairing-lambda-dense (592 lines) and in the local triviality step of Lemma
% cv-toroidal-theta-family-pel-bundles-seesaw. This is the converse of Lemma cv-cusp-cm-point-eigen-degree-rank-one-norm-weil.

\begin{lemma}\label{cv-cusp-cm-point-eigen-degree-tensor-complex-pairing}
Assume the conclusion of Theorem~\ref{cv-cusp-cm-point-eigen-degree-rank-one-norm-riemann-theta} and the laws of
abelian schemes (Definition~\ref{cv-cusp-cm-point-eigen-degree-tensor-laws}).
In Situation~\ref{cv-cusp-setting}, let $\sqrt{-d}\in E$ be the square root of $-d$ of the datum, and put
$w_0=(1+\sqrt{-d})/2$ and $\mathcal O=\mathbf Z[w_0]$. Let $\omega\in\mathbf C$, and let $C$ be a principally polarized
elliptic curve over $\mathbf C$ with an action of $\mathcal O$ whose Rosati involution is complex conjugation and a basis
$v$ of its Lie algebra on which $w_0$ acts by $\omega$, with the line bundle $\mathcal L_C$ of its polarization. Let $L$
be the lattice of a PEL datum over $E$ with its perfect alternating form $\langle\ ,\ \rangle$, let
$\varepsilon\in\{\pm1\}$, and put $h(x,y)=\varepsilon(\langle w_0x,y\rangle-\langle x,y\rangle)+\varepsilon\langle
x,y\rangle w_0$. Let $A$ be a principally polarized abelian variety over $\mathbf C$ with an action of $\mathcal O_E$ for
which the Rosati involution of the polarization induces complex conjugation on $\mathcal O_E$, with the line bundle
$\mathcal L_A$ of its polarization. Let $M\subset L$ be an $\mathcal O_E$-submodule with $\varepsilon\langle
w_0x,x\rangle>0$ for every nonzero $x\in M$, and let $t_a\colon C\to A$, for $a\in M$, be homomorphisms such that
$t_{a+b}=t_a+t_b$ and $t_{sa}=s\circ t_a$ for all $a,b\in M$ and $s\in\mathcal O_E$, such that $\sqrt{-d}\circ t_a\circ
v=(2\omega-1)(t_a\circ v)$ for every $a\in M$, and such that for every $a\in M$ there is $N\ge1$ with
$(t_a^*\mathcal L_A)^{\otimes N}\cong\mathcal L_C^{\otimes N|\langle w_0a,a\rangle|}$. Then $(t_a)_{a\in M}$ is a family
of homomorphisms from $C$ to $A$ for $h$ on $M$ (Definition~\ref{cv-cusp-cm-point-eigen-degree-tensor-family}).
\end{lemma}
\leanref{Development/AbelianVarietyTheory/Proofs/CubicalLineBundles/IsCMPairingFamilyOfPolarisationDiagonalComplexStatement.lean}{GlobalTheory.T.isCMPairingFamily_of_polarisationDiagonal_complex}
\leanref{Development/AbelianVarietyTheory/Proofs/CubicalLineBundles/IsCMPairingFamilyOfPolarisationDiagonalComplex.lean}{GlobalTheory.CM.FamPairing.isFam_of_isHomFamily}
\begin{proof}
Step 1: uniformizations and lifts. The line bundles $\mathcal L_C$ and $\mathcal L_A$ are trivialized along the zero
sections. By Theorem~\ref{cv-cusp-cm-point-eigen-degree-rank-one-norm-riemann-theta}, $C$ and $A$ have uniformizations
$u_C\colon\mathbf C^{g_C}\to C(\mathbf C)$ and $u_A\colon\mathbf C^h\to A(\mathbf C)$ for $\mathcal L_C$ and
$\mathcal L_A$, with kernels $\Lambda_C$ and $\Lambda_A$ and Appell--Humbert data $(H_C,\chi_C)$ and $(H_A,\chi_A)$, for
which (2) of that theorem holds. For $a\in M$, Lemma~\ref{cv-origin-hodge-norm-form-lift}, applied to $u_A$ and to
$t_a\circ u_C$, gives a $\mathbf C$-linear map $F_a$ with $u_A\circ F_a=t_a\circ u_C$. Then $F_a(\Lambda_C)\subset\Lambda_A$,
as $u_A(F_a(l))=t_a(u_C(l))$ is the origin for $l\in\Lambda_C$; if a tangent vector $v'$ of $C$ at the origin has
differential $w$ for $u_C$, then $t_a\circ v'$ has differential $F_a(w)$ for $u_A$, by the chain rule as in the proof of
Lemma~\ref{cv-origin-hodge-norm-form-compare}; and a morphism $t'\colon C\to A$ with $u_A\circ F_a=t'\circ u_C$ is $t_a$, as
the two agree on complex points, $C$ is reduced and $A$ is separated. Differentials are additive and $\mathbf C$-linear in
the tangent vector, for the sum and the multiples of the Lie algebra (Definition~\ref{cv-origin-reference-lattices}(1)).

Step 2: the character of $\mathcal O$. As $v$ is a basis of the Lie algebra of $C$, for $s\in\mathcal O$ there is a unique
$\tau_0(s)\in\mathbf C$ with $s\circ v=\tau_0(s)v$. As $\mathcal O$ acts through a ring homomorphism to the endomorphisms
of $C$, which act additively on tangent vectors, $\tau_0\colon\mathcal O\to\mathbf C$ is a ring homomorphism, and
$\tau_0(w_0)=\omega$. So $(2\omega-1)^2=\tau_0((2w_0-1)^2)=-d$. Let $w_v$ be the differential of $v$. Every
$z\in\mathbf C^{g_C}$ is a multiple of $w_v$: it is the differential of a tangent vector, which is $sv$ for some $s$, with
differential $sw_v$. For $s\in\mathcal O$, Lemma~\ref{cv-origin-hodge-norm-form-lift}, applied to $u_C$ and to $s\circ u_C$,
gives a $\mathbf C$-linear map $G_s$ with $u_C\circ G_s=s\circ u_C$; the tangent vector $s\circ v=\tau_0(s)v$ has the
differentials $G_s(w_v)$ and $\tau_0(s)w_v$, so $G_s(w_v)=\tau_0(s)w_v$.

Step 3: the vectors $x(a)$. For $a\in M$ put $x(a)=F_a(w_v)$, the differential of $t_a\circ v$. As
$t_{a+b}\circ v=t_a\circ v+t_b\circ v$ in the Lie algebra of $A$, $x(a+b)=x(a)+x(b)$. As $t_{\sqrt{-d}a}\circ
v=\sqrt{-d}\circ t_a\circ v=(2\omega-1)(t_a\circ v)$, $x(\sqrt{-d}\,a)=(2\omega-1)x(a)$. Hence
$2x(w_0a)=x(a)+x(\sqrt{-d}\,a)=2\omega x(a)$, and, as $\mathcal O=\mathbf Z[w_0]$, $x(sa)=\tau_0(s)x(a)$ for every
$s\in\mathcal O$.

Step 4: the diagonal values. Let $a\in M$, and put $m=|\langle w_0a,a\rangle|$, which is $\varepsilon\langle w_0a,a\rangle$
by the positivity. Let $N\ge1$ with $(t_a^*\mathcal L_A)^{\otimes N}\cong\mathcal L_C^{\otimes Nm}$. If $a=0$, then $t_a=0$,
$F_a=0$ and $x(a)=0$. Otherwise $m\ge1$, and $u_C$, with the $m$-th powers of its trivializing functions, is a
uniformization of $C$ for $\mathcal L_C^{\otimes m}$ with Appell--Humbert data $(mH_C,\chi_C^m)$.
Lemma~\ref{cv-origin-hodge-norm-form-pullback}, for $\mathcal L_C^{\otimes m}$, $\mathcal L_A$, $t_a$, $F_a$ and $n=N$,
gives $m\,z^{\mathsf T}H_C\bar z'=F_a(z)^{\mathsf T}H_A\overline{F_a(z')}$ for all $z,z'$. In both cases
\[
 x(a)^{\mathsf T}H_A\,\overline{x(a)}=\langle w_0a,a\rangle\,\varepsilon\,w_v^{\mathsf T}H_C\bar w_v.
\]

Step 5: all values. By Steps 3 and 4, Lemma~\ref{cv-cusp-cm-point-eigen-degree-tensor-complex-pairing-polarize}, for
$H=H_A$ and $c=\varepsilon\,w_v^{\mathsf T}H_C\bar w_v$, gives
$x(a)^{\mathsf T}H_A\overline{x(b)}=(\langle w_0a,b\rangle-\langle a,b\rangle+\langle a,b\rangle\omega)c$ for all
$a,b\in M$.

Step 6: condition (1). Let $a\in M$ and $s\in\mathcal O$. For $z=rw_v$, $F_{sa}(z)=r\,x(sa)=r\tau_0(s)x(a)$ and
$F_a(G_s(z))=r\tau_0(s)F_a(w_v)=r\tau_0(s)x(a)$; so $F_{sa}=F_a\circ G_s$, by Step~2. Then
$u_A\circ F_{sa}=t_a\circ u_C\circ G_s=t_a\circ s\circ u_C$, so $t_{sa}=t_a\circ s$ by Step~1. As $t_{sa}=s\circ t_a$, $t_a$
commutes with the actions of $\mathcal O$.

Step 7: condition (3). Let $a,b\in M$ and put $s=h(a,b)\in\mathcal O$, so that
$\tau_0(s)=\varepsilon(\langle w_0a,b\rangle-\langle a,b\rangle)+\varepsilon\langle a,b\rangle\omega$. For $z=rw_v$ and
$z'=r'w_v$, by Step~5,
\[
 F_a(z)^{\mathsf T}H_A\overline{F_b(z')}=r\bar r'x(a)^{\mathsf T}H_A\overline{x(b)}=r\bar r'\tau_0(s)\,w_v^{\mathsf
 T}H_C\bar w_v=G_s(z)^{\mathsf T}H_C\bar z'.
\]
By Lemma~\ref{cv-cusp-cm-point-eigen-degree-tensor-complex-pairing-hermitian}, for $X=Y=B=C$, $A$, the uniformizations
$u_C$, $u_C$, $u_A$ and $u_C$, the morphisms $t_a$, $t_b$, $s$ and the identity, and the maps $F_a$, $F_b$, $G_s$ and the
identity, the commutator pairings $e^A_n(t_ax,t_by)$ and $e^C_n(sx,y)$ agree. By
Lemma~\ref{cv-cusp-cm-point-eigen-degree-tensor-complex-pairing-lambda},
$\Lambda_{\mathcal L_A}(t_ap_1,t_bp_2)\cong\Lambda_{\mathcal L_C}(sp_1,p_2)$ on $C\times C$.

Condition (2) of Definition~\ref{cv-cusp-cm-point-eigen-degree-tensor-family} is a hypothesis, so $(t_a)_{a\in M}$ is a
family of homomorphisms from $C$ to $A$ for $h$ on $M$.
\end{proof}
% Lean: isFam_of_isHomFamily, under ThetaUnifAV and Out.Laws. Step 1 is ThetaUnifAV for C.P and T.pol with their
% rigidifications, and HomLinear.exists_linear_of_isMonHom for F a (with hFex, hFLambda, hFdiff, hFuniq). Step 2 is
% LieChar.ell_char (= exists_char: the unique coefficient d(c), with LieComb.lieComb_algebra for its additivity and
% multiplicativity), span (UnifDifferential.isDiff_of_isLieComb and C.basis) and hG. Step 3 is hxadd
% (LieChar.isLieComb_mul), hxrt and smul_of_rt, with char_sq. Step 4 is hq, through hermForm_pull_eq_of_lbIso_pow with
% the powers N and N m, which the text reduces to Lemma cv-origin-hodge-norm-form-pullback for the m-th power of L_C,
% and habs for the positivity. Step 5 is FamPairing.polarise. Step 6 is hFuniq with hG and hxc. Step 7 is
% T.weilEq_iff_hermForm_eq_complex (with hh for tau_0(h(a, b)), from IsKForm) and
% T.lambdaBundleIso_of_weilPairing_compatible. What remains open is in Lemma
% cv-cusp-cm-point-eigen-degree-tensor-complex-pairing-lambda.

\begin{lemma}\label{cv-cusp-cm-point-eigen-degree-tensor-complex}
Assume the conclusion of Theorem~\ref{cv-cusp-cm-point-eigen-degree-rank-one-norm-riemann-theta} and the laws of
abelian schemes (Definition~\ref{cv-cusp-cm-point-eigen-degree-tensor-laws}).
In Situation~\ref{cv-cusp-setting}, let $\sqrt{-d}\in E$ be the square root of $-d$ of the datum, and put
$w_0=(1+\sqrt{-d})/2$ and $\mathcal O=\mathbf Z[w_0]$. Let $\theta\in\mathbf C$ with $\theta^2=-d$, and let $C$ be a
principally polarized elliptic curve over $\mathbf C$ with an action of $\mathcal O$ whose Rosati involution is complex
conjugation and a basis of its Lie algebra on which $\sqrt{-d}$ acts by $\theta$, such that the lattice of every
uniformization of $C$ is $c(\mathbf Z+\mathbf Z\frac{1+i\sqrt d}2)$ for some $c\in\mathbf C^\times$. Let $r\ge0$, and
let $A$ be a principally polarized abelian variety over $\mathbf C$ with an action of $\mathcal O_E$ for which the
Rosati involution of the polarization induces complex conjugation on $\mathcal O_E$. Suppose that for every
algebraically closed field $\Omega$ and every embedding $\iota\colon\mathbf C\to\Omega$, the eigenspace of an
embedding $\sigma\colon E\to\Omega$ in the Lie algebra of the base change of $A$ along $\iota$ has dimension $r$ if
$\sigma(\sqrt{-d})=\iota(\theta)$ and dimension $0$ otherwise. Then there are the lattice $L$ of a PEL datum over $E$ of
signature $(p,q)$ with $p+q=r$, with its perfect alternating form $\langle\ ,\ \rangle$
(Definition~\ref{cv-cm-point-periods}), a sign $\varepsilon\in\{\pm1\}$ with $\varepsilon\langle w_0x,x\rangle>0$ for
every nonzero $x\in L$, and a family of homomorphisms from $C$ to $A$ for the form
$h(x,y)=\varepsilon(\langle w_0x,y\rangle-\langle x,y\rangle)+\varepsilon\langle x,y\rangle w_0$ on $L$
(Definition~\ref{cv-cusp-cm-point-eigen-degree-tensor-family}) that is onto at the identity of $\mathbf C$.
\end{lemma}
\leanref{Development/AbelianVarietyTheory/Proofs/CubicalLineBundles/ExistsHomFamilyOfPureTypeComplexTripleStatement.lean}{GlobalTheory.T.exists_homFamily_of_pureType_complex_triple}
\leanref{Development/AbelianVarietyTheory/Proofs/CubicalLineBundles/ExistsHomFamilyOfPureTypeComplexTriple.lean}{GlobalTheory.CM.exists_isFam_of_pureType}
\begin{proof}
Let $\theta_0$ be the square root of $-d$ fixed by the datum, and put $\omega=(1+\theta)/2$, so that $\theta=2\omega-1$
and $w_0$ acts on the basis $v$ of the Lie algebra of $C$ by $\omega$.

Step 1: the signature. Put $(p,q)=(r,0)$ if $\theta=\theta_0$ and $(p,q)=(0,r)$ otherwise, so that $p+q=r$; in the
second case $\theta=-\theta_0$. By the hypothesis on eigenspaces, for $\Omega=\mathbf C$ and $\iota$ the identity, the
$\tau$-eigenspace of the Lie algebra of $A$, for an embedding $\tau\colon E\to\mathbf C$, has dimension $r$ if
$\tau(\sqrt{-d})=\theta$ and $0$ otherwise; in both cases this is $p$ if $\tau(\sqrt{-d})=\theta_0$ and $q$ otherwise.
Moreover an embedding $\tau$ with $\tau(\sqrt{-d})\neq\theta$ has $p=0$ if $\tau(\sqrt{-d})=\theta_0$ and $q=0$ if not.

Step 2: uniformizations. Let $(\sigma_k)_{1\le k\le g}$ be a type of signature $(p,q)$ for $\theta_0$
(Definition~\ref{cv-cm-point-periods}(2)), which lists each embedding as often as its multiplicity. The line bundles of
the polarizations of $A$ and $C$ are trivialized along the zero sections. By
Theorem~\ref{cv-cusp-cm-point-eigen-degree-rank-one-norm-riemann-theta}, $A$ has a uniformization $u\colon\mathbf
C^h\to A(\mathbf C)$ with Appell--Humbert data $(H,\chi)$ for which (2) and (4) of that theorem hold, and $C$ has a
uniformization $u_0\colon\mathbf C\to C(\mathbf C)$ with Appell--Humbert data for which (2), (4) and (5) hold. By the
hypothesis on $C$, the kernel of $u_0$ is $c(\mathbf Z+\mathbf Z\frac{1+i\sqrt d}2)$ for some $c\neq0$.

Step 3: the datum. Lemma~\ref{cv-cusp-cm-point-eigen-degree-tensor-complex-datum}, by Step~1, gives the lattice $L$ of a
PEL datum over $E$ of signature $(p,q)$, with its form $\langle\ ,\ \rangle$, a period $P\colon L\to\mathbf C^g$ of type
$(\sigma_k)$ and a $\mathbf C$-linear bijection $F\colon\mathbf C^g\to\mathbf C^h$ with (1), (2) and (3) of that lemma.

Step 4: purity. Every $a\in L$ has $P(a)_k=0$ for every $k$ with $\sigma_k(\sqrt{-d})\neq\theta$, since there is no
such $k$: by Step~1 the embedding $\sigma_k$ would occur in the type with multiplicity $0$.

Step 5: the homomorphisms. As $\theta^2=-d$, Lemma~\ref{cv-cusp-cm-point-eigen-degree-tensor-complex-homs}, for $u$, $L$,
$P$, $F$, $C$, $u_0$, $c$, $\rho=\theta$ and $M=L$, which is allowed by Step~4, gives homomorphisms $t_a\colon C\to A$,
for $a\in L$, with $t_a(u_0(cz))=u(F(zP(a)))$ and the properties (1), (2) and (3) of that lemma for the tangent vector
$v$.

Step 6: the sign. As $\theta^2=-d<0$, $\theta$ is not real, so $\operatorname{Im}\omega=\operatorname{Im}\theta/2$ is
nonzero. Let $\varepsilon\in\{\pm1\}$ with $\varepsilon\operatorname{Im}\omega>0$. For a nonzero $x\in L$,
Lemma~\ref{cv-cusp-cm-point-eigen-degree-tensor-complex-positive}, with Step~4, gives
$\langle w_0x,x\rangle=\operatorname{Im}(\omega)\,r$ with $r>0$, so $\varepsilon\langle w_0x,x\rangle>0$.

Step 7: the family. With $h$ as in the statement,
Lemma~\ref{cv-cusp-cm-point-eigen-degree-tensor-complex-pairing}, for $C$, $v$, $L$, $\varepsilon$, $A$, $M=L$ and the
$t_a$, by Steps 5 and 6, shows that $(t_a)_{a\in L}$ is a family of homomorphisms from $C$ to $A$ for $h$ on $L$.

Step 8: onto. The image $P(L)$ spans $\mathbf C^g$ over $\mathbf R$, hence over $\mathbf C$, and $F$ is a bijection, so
the vectors $F(P(a))$, $a\in L$, span $\mathbf C^h$ over $\mathbf C$. Every point of $A$ over $\mathbf C$ is $u(w)$ for
some $w\in\mathbf C^h$, as $u$ is onto. Write $w=\sum_mz_mF(P(a_m))$, a finite sum with $z_m\in\mathbf C$ and $a_m\in L$.
As $u$ is a homomorphism, $u(w)=\sum_mu(F(z_mP(a_m)))=\sum_mt_{a_m}(u_0(cz_m))$. So the family is onto at the identity
of $\mathbf C$.
\end{proof}
\leanref{Development/AbelianVarietyTheory/Proofs/CubicalLineBundles/ExistsHomFamilyOfPureTypeComplexTriple.lean}{GlobalTheory.CM.PureTriple.ontoAt_of_span}
% Lean: exists_isFam_of_pureType, under ThetaUnifAV and Out.Laws. The character tau of O on Lie C with tau(rt) =
% 2 omega - 1 is LieChar.ell_char (the text's basis on which sqrt(-d) acts by theta). Step 1 is the case split on
% 2 omega - 1 = rt0 (hdet', hpure). Step 2 is AllClasses.exists_isTypeFn, PureType.exists_isThetaUnif (from
% ThetaUnifAV) for T.pol and for C.P, and hC (IsPrincipalAt at the identity, through polFibre_id_comp_cancel). Step 3 is
% exists_datum_period. Step 4 is hlat. Step 5 is exists_isHomFamily with L = top. Step 6 is the choice of eps (the sign
% of Im omega, via Re((2 omega - 1)^2) = -d) and exists_hermitianPairing_eq_im_mul_pureType. Step 7 is
% T.isCMPairingFamily_of_polarisationDiagonal_complex with C.basis.1 0 for the tangent vector. Step 8 is ontoAt_of_span
% with Submodule.mem_span_set' and u.ex_add. What remains open is in Lemmas
% cv-cusp-cm-point-eigen-degree-tensor-complex-homs-hom (about 400 lines), -complex-homs-bundle (about 250) and
% -complex-pairing-lambda (about 900).

\subsection{Descent of Serre tensor products}\label{subsection-cv-serre-tensor-descent}

\begin{lemma}\label{cv-cusp-cm-point-eigen-degree-tensor-identities-section}
Let $j\colon R\to k$ be a flat homomorphism from a ring to a field, let $X$ be a quasi-compact and quasi-separated
scheme over $R$, and let $\pi\colon X_k\to X$ be the projection from $X_k=X\times_R\operatorname{Spec}k$. Suppose that
$X_k$ is nonempty and that every global function on $X_k$ comes from $k$. Let $\mathcal N$ be a line bundle on $X$ with
$\pi^*\mathcal N$ trivial. Then $\mathcal N$ has a global section $\sigma$ such that $\pi^*\sigma$ generates
$\pi^*\mathcal N$ at every point of $X_k$.
\end{lemma}
\leanref{Development/AbelianVarietyTheory/Proofs/CubicalLineBundles/ExistsGenOfGenericFibre.lean}{GlobalTheory.GenericTriv.exists_sec_isUnit}
\begin{proof}
Let $t$ be a generator of $\pi^*\mathcal N$. By flat base change \cite[Tag 02KH]{Sta},
$\Gamma(X_k,\pi^*\mathcal N)=\Gamma(X,\mathcal N)\otimes_Rk$, so $t=\sum_ac_a\pi^*u_a$ for finitely many $c_a\in k$ and
$u_a\in\Gamma(X,\mathcal N)$. For each $a$, the ratio of $\pi^*u_a$ to $t$ is a global function on $X_k$, so
$\pi^*u_a=\kappa_at$ with $\kappa_a\in k$. If every $\kappa_a$ were $0$, then $t=\sum_ac_a\kappa_at$ would be $0$, which
is impossible, as $X_k$ is nonempty and $t$ is a generator. So $\kappa_a\neq0$ for some $a$, and $\sigma=u_a$ has
$\pi^*\sigma=\kappa_at$, a generator at every point.
\end{proof}
% Lean: exists_sec_isUnit: FlatSec.exists_comb_of_flat (the flat base change of sections, as in Lemma
% cv-toroidal-abelian-dvr-ample-descent), exists_ratio (the ratio d_a, a global function by Out.TLinPf.glue), the
% surjectivity hs for kappa_a, and the nonvanishing of t at a point of the nonempty X_k.

\begin{lemma}\label{cv-cusp-cm-point-eigen-degree-tensor-identities-field}
Let $\beta\colon k\to k'$ be a homomorphism of fields, $X$ a proper and geometrically integral scheme over $k$, and
$\pi\colon X'\to X$ the projection from $X'=X\times_k\operatorname{Spec}k'$. Let $\mathcal N$ be a line bundle on $X$ with
$\pi^*\mathcal N$ trivial. Then $\mathcal N$ is trivial.
\end{lemma}
\leanref{Development/AbelianVarietyTheory/Proofs/CubicalLineBundles/ExistsGenOfFieldExtension.lean}{GlobalTheory.FieldTriv.exists_gen_of_isPullback}
\begin{proof}
The scheme $X$ is quasi-compact and quasi-separated, and $X'$ is integral, in particular nonempty. The global functions
on $X'$ are those of $k'$. Indeed, let $\bar k$ be an algebraic closure of $k$; the scheme $X_{\bar k}$ is integral and
proper over $\bar k$, so $\Gamma(X_{\bar k},\mathcal O)$ is a domain finite over $\bar k$
\cite[Th\'eor\`eme~3.2.1]{Gro61}, hence $\bar k$. By flat base change \cite[Tag 02KH]{Sta},
$\Gamma(X,\mathcal O)\otimes_k\bar k=\bar k$, so $\Gamma(X,\mathcal O)=k$, and then
$\Gamma(X',\mathcal O)=\Gamma(X,\mathcal O)\otimes_kk'=k'$. The homomorphism $\beta$ is flat. By
Lemma~\ref{cv-cusp-cm-point-eigen-degree-tensor-identities-section}, $\mathcal N$ has a global section $\sigma$ such
that $\pi^*\sigma$ generates $\pi^*\mathcal N$ at every point. Let $x\in X$. The fibre of $\pi$ over $x$ is the spectrum
of the nonzero ring $\kappa(x)\otimes_kk'$, so $x=\pi(y)$ for some $y\in X'$. In a trivialization of $\mathcal N$ near
$x$, the local homomorphism $\mathcal O_{X,x}\to\mathcal O_{X',y}$ maps the germ of $\sigma$ to a unit, so this germ is a
unit. So $\sigma$ generates $\mathcal N$ at every point, and $\mathcal N$ is trivial.
\end{proof}
\leanref{Development/AbelianVarietyTheory/Proofs/HomogeneousIdeals/IsIsoAppTopOfIsPullbackOfSurjective.lean}{GlobalTheory.ProperFunctions.isIso_appTop_baseChange_of_field}
% Lean: exists_gen_of_isPullback: GeometricallyIntegral.geometrically_isIntegral (X' integral),
% ProperFunctions.isIso_appTop_baseChange_of_field (isIso_appTop_of_field through the algebraic closure, with
% isIso_appTop_of_isAlgClosed and isIso_appTop_of_isPullback_of_surjective, then isIso_appTop_of_isPullback), the
% flatness of beta, GenericTriv.exists_sec_isUnit, Scheme.exists_preimage_of_isPullback for y, and
% isUnit_of_isUnit_germ with gs_stalkMap and isUnit_map_iff.

\begin{lemma}\label{cv-cusp-cm-point-eigen-degree-tensor-identities-dedekind}
Let $R$ be a Dedekind domain with fraction field $K$, $A$ an abelian scheme over $R$ with unit section $e$, and
$\pi\colon A_K\to A$ the projection from its generic fibre $A_K=A\times_R\operatorname{Spec}K$. Let $\mathcal N$ be a line
bundle on $A$ such that $e^*\mathcal N$ and $\pi^*\mathcal N$ are trivial. Then $\mathcal N$ is trivial.
\end{lemma}
\leanref{Development/AbelianVarietyTheory/Proofs/CubicalLineBundles/SurjectiveAppTopOfSection.lean}{GlobalTheory.TrivDedekind.exists_gen_of_generic}
\leanref{Development/AbelianVarietyTheory/Proofs/CubicalLineBundles/ExistsGenOfGenericFibre.lean}{GlobalTheory.GenericTriv.exists_gen_of_fibre}
\begin{proof}
Step 1: the scheme $A$. The scheme $A$ is proper and smooth over $R$, with geometrically connected fibres; its fibres
over algebraically closed fields are regular \cite[Tag 056S]{Sta}, hence integral \cite[Tag 033N]{Sta}, so $A$ has
geometrically integral fibres. It is quasi-compact and quasi-separated. The generic fibre $A_K$ is integral and proper
over $K$, so $\Gamma(A_K,\mathcal O)$ is a domain finite over $K$ \cite[Th\'eor\`eme~3.2.1]{Gro61}, hence a field. The
unit section of $A_K$ gives a homomorphism $\Gamma(A_K,\mathcal O)\to K$ that is a left inverse of
$K\to\Gamma(A_K,\mathcal O)$; it is injective, as $\Gamma(A_K,\mathcal O)$ is a field. So the global functions on $A_K$
are those of $K$.

Step 2: a section. The homomorphism $R\to K$ is flat, and $A_K$ is nonempty. By
Lemma~\ref{cv-cusp-cm-point-eigen-degree-tensor-identities-section}, $\mathcal N$ has a global section $\sigma$ such
that $\pi^*\sigma$ generates $\pi^*\mathcal N$ at every point. Let $\xi$ be the generic point of $\operatorname{Spec}R$
and let $x\in A$ lie over $\xi$. Then $x=\pi(y)$ for a point $y$ of $A_K$, and in a trivialization of $\mathcal N$ near
$x$ the local homomorphism $\mathcal O_{A,x}\to\mathcal O_{A_K,y}$ maps the germ of $\sigma$ to a unit; so $\sigma$
generates $\mathcal N$ at $x$.

Step 3: the constant $a$. Let $\rho$ be a generator of $e^*\mathcal N$. Then $e^*\sigma=a\rho$ for some
$a\in\Gamma(\operatorname{Spec}R,\mathcal O)=R$. The point $e(\xi)$ lies over $\xi$, so $\sigma$ generates $\mathcal N$
there by Step~2, and $a\neq0$.

Step 4: local divisibility. Let $x\in A$, let $\epsilon$ be a generator of $\mathcal N$ near $x$, and write
$\sigma=s\epsilon$ near $x$. Then $s$ and $a$ divide each other in $\mathcal O_{A,x}$. If $x$ lies over $\xi$, then $s$ is
a unit at $x$ by Step~2, and so is $a$, whose image in $\mathcal O_{A_K,y}$, for $y$ as in Step~2, is a nonzero element
of $K$. Otherwise $x$ lies over a nonzero prime $\mathfrak p$ of $R$, and $V=R_{\mathfrak p}$ is a discrete valuation
ring with fraction field $K$ \cite[Theorem~9.3]{AM69}. Let $A_V=A\times_R\operatorname{Spec}V$, with projection
$g\colon A_V\to A$ and unit section $e_V$; it is smooth over $V$ with integral fibres, and its generic fibre is $A_K$. Let
$y\in A_V$ be a point over $x$ and over the closed point of $\operatorname{Spec}V$. On $A_K$, $a^{-1}\pi^*\sigma$ is a
generator of $\pi^*\mathcal N$ whose pullback along the unit section of $A_K$ is the pullback of $\rho$. By
Lemma~\ref{cv-toroidal-abelian-dvr-cubical-dvr}(3), for $H=A_V$, the line bundle $g^*\mathcal N$, the section $e_V$ and
the pullback of $\rho$, it extends to a generator $c$ of $g^*\mathcal N$ over $A_V$. The sections $g^*\sigma$ and $ac$ of
$g^*\mathcal N$ agree on $A_K$, so they are equal, as $A_V$ is integral
(Lemma~\ref{cv-toroidal-abelian-dvr-cubical-dvr}(1)). So near $y$ the pullback of $s$ is $a$ times a unit, and $s$ and $a$
divide each other in $\mathcal O_{A_V,y}$. The local homomorphism $\mathcal O_{A,x}\to\mathcal O_{A_V,y}$ is flat, as
$g$ is a base change of the flat morphism $\operatorname{Spec}V\to\operatorname{Spec}R$, hence faithfully flat
\cite[Tag 00HR]{Sta}, so every ideal of $\mathcal O_{A,x}$ is the contraction of its extension
\cite[Tag 05CK]{Sta}. So $s$ and $a$ divide each other in $\mathcal O_{A,x}$.

Step 5: division by $a$. As $A$ is flat over the domain $R$ and $a\neq0$, $a$ is a nonzerodivisor in every local ring
of $A$. By Step~4, on each open set where $\mathcal N$ has a generator $\epsilon$ and $\sigma=s\epsilon$, the local
quotients of $s$ by $a$ glue to a function $s'$ with $s=as'$, and the sections $s'\epsilon$ glue to a global section
$\sigma'$ of $\mathcal N$ with $\sigma=a\sigma'$. Let $x\in A$. By Step~4, $a=sw$ for some $w\in\mathcal O_{A,x}$, so
$a=as'w$, and $s'w=1$ as $a$ is a nonzerodivisor. So $\sigma'$ generates $\mathcal N$ at every point, and $\mathcal N$ is
trivial.
\end{proof}
\leanref{Development/AbelianVarietyTheory/Proofs/CubicalLineBundles/ExistsGenOfGenericFibre.lean}{GlobalTheory.GenericTriv.exists_unit_dvr}
% Lean: TrivDedekind.exists_gen_of_generic: forall_geometricallyIntegral_of_smooth_of_geometricallyConnected,
% DvrModel.isIntegral_of_geometricallyIntegral and surjective_appTop_of_section (Mathlib's isField_of_universallyClosed
% and the section; Step 1), then GenericTriv.exists_gen_of_fibre for j the map to the fraction field: exists_sec_isUnit
% (Step 2), exists_ratio for e with Gamma(Spec R) = R (Step 3), dvd (Step 4: lift for the generic points;
% IsLocalization.AtPrime.isDiscreteValuationRing_of_dedekind_domain, Scheme.exists_preimage_of_isPullback and
% exists_unit_dvr, whose extension is DvrModel.existsUnique_isUnit_appLE_eq, the step of Lemma
% cv-toroidal-abelian-dvr-cubical-dvr(3), with SecSpread.ext_generic for g^*sigma = a c; dvd_descend for the faithful
% flatness), and Geo.Tor.exists_div_sec with cst_regular (Step 5).

\begin{lemma}\label{cv-cusp-cm-point-eigen-degree-tensor-identities-line}
Let $R$ be a Dedekind domain, $\beta\colon R\to k$ an injective homomorphism to a field, $A$ an abelian scheme over $R$
with unit section $e$, and $\pi\colon A_k\to A$ the projection from $A_k=A\times_R\operatorname{Spec}k$. Let
$\mathcal N$ and $\mathcal N'$ be line bundles on $A$ with $e^*\mathcal N\cong e^*\mathcal N'$ and
$\pi^*\mathcal N\cong\pi^*\mathcal N'$. Then $\mathcal N\cong\mathcal N'$.
\end{lemma}
\leanref{Development/AbelianVarietyTheory/Proofs/CubicalLineBundles/LineBundleIsoDescendsFieldToDedekindDomainStatement.lean}{GlobalTheory.T.lineBundleIso_descends_field_to_dedekindDomain}
\leanref{Development/AbelianVarietyTheory/Proofs/CubicalLineBundles/LineBundleIsoDescendsFieldToDedekindDomain.lean}{GlobalTheory.IsoDescent.lbIso_of_field}
\begin{proof}
Let $K$ be the fraction field of $R$ and $\delta\colon K\to k$ the homomorphism extending $\beta$. Let
$\pi_0\colon A_K\to A$ be the projection from the generic fibre. Then $A_k=A_K\times_K\operatorname{Spec}k$, with a
projection $q\colon A_k\to A_K$ such that $\pi=\pi_0q$. Put $\mathcal D=\mathcal N\otimes\mathcal N'^{-1}$. Then
$q^*\pi_0^*\mathcal D=\pi^*\mathcal D$ is trivial. The scheme $A_K$ is proper and geometrically integral over $K$, as in
Step~1 of the proof of Lemma~\ref{cv-cusp-cm-point-eigen-degree-tensor-identities-dedekind}, so
$\pi_0^*\mathcal D$ is trivial by Lemma~\ref{cv-cusp-cm-point-eigen-degree-tensor-identities-field}. Also
$e^*\mathcal D$ is trivial. By Lemma~\ref{cv-cusp-cm-point-eigen-degree-tensor-identities-dedekind}, $\mathcal D$ is
trivial, that is, $\mathcal N\cong\mathcal N'$.
\end{proof}
% Lean: lbIso_of_field: forall_geometricallyIntegral_of_smooth_of_geometricallyConnected, IsFractionRing.lift for delta,
% the pullback square H_1 (IsPullback.of_right), exists_gen_pull_div_iff and the classes CechPic.cls for D = N/N',
% FieldTriv.exists_gen_of_isPullback (h_2) and TrivDedekind.exists_gen_of_generic (h_3).

\begin{lemma}\label{cv-cusp-cm-point-eigen-degree-tensor-identities-lambda}
Let $R$ be a Dedekind domain and $f\colon R\to k$ an injective homomorphism to a field. Let $X$ and $Y$ be abelian schemes
over $R$, $A$ and $B$ commutative group schemes over $R$, $\mathcal L$ and $\mathcal M$ line bundles on $A$ and $B$, and
$f_1\colon X\to A$, $f_2\colon Y\to A$, $g_1\colon X\to B$ and $g_2\colon Y\to B$ homomorphisms. Write $X_k$, $Y_k$,
$A_k$, $B_k$, $\mathcal L_k$, $\mathcal M_k$, $f'_i$ and $g'_i$ for the base changes along $f$, and $p_1,p_2$ and
$p'_1,p'_2$ for the projections of $X\times_RY$ and of $X_k\times_kY_k$. For a line bundle $\mathcal P$ on a commutative
group scheme $G$ over a scheme $S$ and points $x,y\in G(T)$, put $\Lambda_{\mathcal P}(x,y)=(x+y)^*\mathcal P\otimes
x^*\mathcal P^{-1}\otimes y^*\mathcal P^{-1}\otimes0^*\mathcal P$, as in
Definition~\ref{cv-toroidal-abelian-dvr-ample-torsion-weil}. If
$\Lambda_{\mathcal L_k}(f'_1p'_1,f'_2p'_2)\cong\Lambda_{\mathcal M_k}(g'_1p'_1,g'_2p'_2)$ on $X_k\times_kY_k$, then
$\Lambda_{\mathcal L}(f_1p_1,f_2p_2)\cong\Lambda_{\mathcal M}(g_1p_1,g_2p_2)$ on $X\times_RY$.
\end{lemma}
\leanref{Development/AbelianVarietyTheory/Proofs/CubicalLineBundles/FamilyIdentitiesDescendFieldToDedekind.lean}{GlobalTheory.FamDescent.lamEq_descend}
\begin{proof}
Put $P=X\times_RY$, an abelian scheme over $R$ with unit section $e$, and
$\mathcal N=\Lambda_{\mathcal L}(f_1p_1,f_2p_2)$ and $\mathcal N'=\Lambda_{\mathcal M}(g_1p_1,g_2p_2)$. We compare the
classes of the line bundles in the Picard groups, on which pullback acts by homomorphisms.

Along the unit section: as $f_1$ and $f_2$ are homomorphisms, $f_1p_1e$, $f_2p_2e$ and their sum are the unit section
$0$ of $A$, so $e^*\mathcal N=0^*\mathcal L\otimes0^*\mathcal L^{-1}\otimes0^*\mathcal L^{-1}\otimes0^*\mathcal L$ is
trivial. In the same way $e^*\mathcal N'$ is trivial, so $e^*\mathcal N\cong e^*\mathcal N'$.

After base change: the base change of $P$ along $f$ is $P_k=X_k\times_kY_k$, with projection $\pi\colon P_k\to P$. Let
$\alpha\colon A_k\to A$ be the projection, so that $\mathcal L_k=\alpha^*\mathcal L$. Then $\alpha f'_1p'_1=f_1p_1\pi$
and $\alpha f'_2p'_2=f_2p_2\pi$, and as $\alpha$ is a homomorphism over $\operatorname{Spec}k\to\operatorname{Spec}R$, it
carries the sum $f'_1p'_1+f'_2p'_2$ to $(f_1p_1+f_2p_2)\pi$ and the unit section of $A_k$ to $0\pi$. As pullback
commutes with tensor products and duals, $\pi^*\mathcal N\cong\Lambda_{\mathcal L_k}(f'_1p'_1,f'_2p'_2)$. In the same way
$\pi^*\mathcal N'\cong\Lambda_{\mathcal M_k}(g'_1p'_1,g'_2p'_2)$, so $\pi^*\mathcal N\cong\pi^*\mathcal N'$ by hypothesis.

By Lemma~\ref{cv-cusp-cm-point-eigen-degree-tensor-identities-line}, for the abelian scheme $P$ and $\beta=f$,
$\mathcal N\cong\mathcal N'$.
\end{proof}
% Lean: lamEq_descend: CechPicMul.exists_picardGroup (the Picard groups Pic, the classes cl and the pullbacks pl, with
% the rules hiso, hcomp and hlin); zero (the unit section, through LamCalc.class_cubeLB and IsMonHom.one_hom); gen (the
% base change, through lamBC_psum_push and the isomorphisms iL, iM) with LamCalc.lbIso_of_lamEq for the hypothesis H
% pulled back along (u, v); AbelianSchemeGerms.isAbelianScheme_tensor for P; and
% T.lineBundleIso_descends_field_to_dedekindDomain. The hypotheses that A and B are commutative and that L, M are line
% bundles are used; the morphisms over k need not be homomorphisms, as they are the base changes. Size of lamEq_descend
% with the four lemmas before it: about 2,200 lines.

\begin{lemma}\label{cv-cusp-cm-point-eigen-degree-tensor-identities}
Assume the laws of abelian schemes (Definition~\ref{cv-cusp-cm-point-eigen-degree-tensor-laws}); this hypothesis is
not used in the proof.
In Situation~\ref{cv-cusp-setting}, let $\mathcal O$, $L$, $M$ and $h$ be as in
Definition~\ref{cv-cusp-cm-point-eigen-degree-tensor-family}. Let $R$ be a Dedekind domain, $f\colon R\to k$ an
injective homomorphism to a field, $C$ a principally polarized elliptic curve over $R$ with an action of $\mathcal O$,
and $A$ a principally polarized abelian scheme over $R$ with an action of $\mathcal O_E$. Let $(t_a)_{a\in M}$ be
morphisms $C\to A$ over $R$ whose base changes along $f$ form a family of homomorphisms from the base change of $C$ to
that of $A$ for $h$ on $M$. Then $(t_a)_{a\in M}$ is a family of homomorphisms from $C$ to $A$ for $h$ on $M$.
\end{lemma}
\leanref{Development/AbelianVarietyTheory/Proofs/CubicalLineBundles/FamilyIdentitiesDescendFieldToDedekindStatement.lean}{GlobalTheory.T.family_identities_descend_field_to_dedekind}
\leanref{Development/AbelianVarietyTheory/Proofs/CubicalLineBundles/FamilyIdentitiesDescendFieldToDedekind.lean}{GlobalTheory.FamDescent.isFam_descend}
\begin{proof}
As $f$ is injective, $\operatorname{Spec}k\to\operatorname{Spec}R$ is schematically dominant, and so is its base change
$C_k\to C$, since $C$ is flat over $R$. As $A$ is separated over $R$, two morphisms $C\to A$ over $R$ that agree on
$C_k$ are equal. Each $t_a$ is therefore a homomorphism, as the two morphisms $C\times_RC\to A$ that express this agree
after base change, and in the same way $t_a$ commutes with the actions of $\mathcal O$, $t_{a+a'}=t_a+t_{a'}$ and
$t_{ba}=b\circ t_a$, as these identities hold after base change. It remains to see that the isomorphisms (3) of
Definition~\ref{cv-cusp-cm-point-eigen-degree-tensor-family} over $k$ descend to isomorphisms over $R$. Let
$a,a'\in M$. The curve $C$ is an abelian scheme over $R$, the action of $h(a,a')$ on $C$ is a homomorphism, and the
line bundles of the polarizations of the base changes $C_k$ and $A_k$ are the pullbacks of the line bundles
$\mathcal L_C$ and $\mathcal L_A$ of the polarizations of $C$ and $A$.
Lemma~\ref{cv-cusp-cm-point-eigen-degree-tensor-identities-lambda}, for $X=Y=C$, the group schemes $A$
and $C$ with $\mathcal L_A$ and $\mathcal L_C$, $f_1=t_a$, $f_2=t_{a'}$, $g_1$ the action of $h(a,a')$ and $g_2$ the
identity, whose base changes are the base changes of $t_a$ and $t_{a'}$, the action of $h(a,a')$ on $C_k$ and the
identity, turns (3) over $k$ into (3) over $R$.
% Lean: isFam_descend: ExtUnique.isSchemeTheoreticallyDominant_spec_of_injective, ExtUnique.isMonHom_of_isGrpPullback
% and ExtUnique.over_hom_ext_of_isPullback for the identities; FamDescent.lamEq_descend (Lemma
% cv-cusp-cm-point-eigen-degree-tensor-identities-lambda) for (3), with C.act.hom (h a b), the isomorphism m.lin of
% the SABMap for L_A, iC of the EllMap for L_C, and hact for the action on C_k.
\end{proof}

\begin{lemma}\label{cv-cusp-cm-point-eigen-degree-tensor-descent-specialize}
Let $k\subset L$ be fields, let $R$ be a finitely generated $k$-subalgebra of $L$, and let $r\in R$ be nonzero.
\begin{enumerate}
\item If $L$ is algebraically closed, there are a field $K_1$ with $k\subset K_1\subset L$ and $[K_1:k]$ finite and a
  homomorphism of $k$-algebras $\rho\colon R\to K_1$ with $\rho(r)\neq0$ and $\rho(c)=c$ for every $c\in R$ algebraic
  over $k$.
\item If $L=\mathbf C$ and $k$ is a number field, and $c\in R$ is such that the values $\sigma(c)$, for the
  homomorphisms of $k$-algebras $\sigma\colon R\to\mathbf C$, lie in a countable set, then $c$ is algebraic over $k$.
\end{enumerate}
\end{lemma}
\leanref{Development/AbelianVarietyTheory/Proofs/CubicalLineBundles/SpecialisationNumberFieldAndCountableValues.lean}{GlobalTheory.exists_specialisation_finite}
\leanref{Development/AbelianVarietyTheory/Proofs/CubicalLineBundles/SpecialisationNumberFieldAndCountableValues.lean}{GlobalTheory.specialisation_numberField_and_countable_values}
\begin{proof}
Step 1: extending homomorphisms. Let $F$ be a field with $k\subset F\subset L$, $M$ an algebraically closed field,
$\tau\colon F\to M$ a homomorphism of $k$-algebras and $r'\in R$ nonzero. Then there is a homomorphism of $k$-algebras
$\rho\colon R\to M$ with $\rho(x)=\tau(x)$ for $x\in R\cap F$ and $\rho(r')\neq0$. Indeed, let $A\subset L$ be the
$F$-algebra generated by $R$ and $r'^{-1}$; it is finitely generated, as $R$ is finitely generated over $k$. Let
$\mathfrak m$ be a maximal ideal of $A$. By Zariski's lemma \cite[Corollary~5.24]{AM69}, $A/\mathfrak m$ is a finite
extension of $F$, generated by finitely many elements $\alpha_1,\dots,\alpha_n$. So $\tau$ extends to a homomorphism
$A/\mathfrak m\to M$: if $\tau_i$ extends $\tau$ to $F_i=F(\alpha_1,\dots,\alpha_i)$, the image under $\tau_i$ of the
minimal polynomial $f$ of $\alpha_{i+1}$ over $F_i$ has a root $\beta$ in $M$, as $M$ is algebraically closed, and
$F_{i+1}\cong F_i[X]/(f)$ maps to $M$ by $\tau_i$ and $X\mapsto\beta$. Composed with $A\to A/\mathfrak m$, it gives a homomorphism $\varphi\colon A\to M$ that is $\tau$ on $F$. Let $\rho$ be
its restriction to $R$, a homomorphism of $k$-algebras, as $\tau$ is one. As $\varphi(r')\varphi(r'^{-1})=1$,
$\rho(r')\neq0$.

Step 2: proof of (1). Let $\bar k\subset L$ be the field of elements algebraic over $k$; it is algebraically closed, as
$L$ is. Step~1, for $F=M=\bar k$, $\tau$ the identity and $r'=r$, gives $\rho_0\colon R\to\bar k$ with $\rho_0(r)\neq0$
and $\rho_0(c)=c$ for every $c\in R$ algebraic over $k$. Let $t_1,\dots,t_n$ generate $R$ as a $k$-algebra, and let $K_1$
be the field generated over $k$ by $\rho_0(t_1),\dots,\rho_0(t_n)$. These are algebraic over $k$, so $[K_1:k]$ is finite,
and $\rho_0(R)\subset K_1$. Take for $\rho$ the map $\rho_0$ with values in $K_1$.

Step 3: proof of (2). Let $S\subset\mathbf C$ be a countable set containing every $\sigma(c)$, and suppose that $c$ is
transcendental over $k$. Let $z\in\mathbf C$ be transcendental over $\mathbf Q$; it is transcendental over $k$, which is
algebraic over $\mathbf Q$. There is an isomorphism of $k$-algebras $k(c)\to k(z)$ with $c\mapsto z$, as both are fields
of rational functions in one variable over $k$. Step~1, for $F=k(c)$, $M=\mathbf C$, this isomorphism followed by the
inclusion $k(z)\subset\mathbf C$, and $r'=1$, gives a homomorphism of $k$-algebras $\sigma\colon R\to\mathbf C$ with
$\sigma(c)=z$; so $z\in S$. Hence $\mathbf C$ is the union of $S$ and of the set of algebraic numbers, both countable,
which is impossible, as $\mathbf C$ is uncountable.
\end{proof}
% Lean: exists_specialisation_finite: exists_algHom_extend_of_fg_ne_zero (Step 1, through
% exists_ringHom_extend_of_finiteType: Ideal.exists_maximal, finite_of_finite_type_of_isJacobsonRing and
% IsAlgClosed.lift) with F = algebraicClosure k L, and K_1 = IntermediateField.adjoin k (rho_0 '' t) (Step 2).
% specialisation_numberField_and_countable_values: its first half is the case of (1) with one fixed c; (2) is its second
% half (RatFunc.algEquivOfTranscendental, exists_algHom_extend_of_fg_ne_zero with r = 1, Algebraic.countable and
% not_countable_complex). Size: SpecialisationNumberFieldAndCountableValues.lean, 184 lines.

\begin{lemma}\label{cv-cusp-cm-point-eigen-degree-tensor-descent-countable}
Assume the conclusion of Theorem~\ref{cv-cusp-cm-point-eigen-degree-rank-one-norm-riemann-theta} and law (22) of
Definition~\ref{cv-cusp-cm-point-eigen-degree-tensor-laws}, on base change.
Let $S$ be a scheme, $t\colon\operatorname{Spec}\mathbf C\to S$ a morphism, and $P$ and $P'$ principally polarized
abelian schemes over $S$, with base changes $P_{\mathbf C}$ and $P'_{\mathbf C}$ along $t$. Then the homomorphisms
$P_{\mathbf C}\to P'_{\mathbf C}$ form a countable set, and two homomorphisms $f,f'\colon P_{\mathbf C}\to P'_{\mathbf C}$
with $f\circ v=f'\circ v$ for every tangent vector $v$ of $P_{\mathbf C}$ at the origin
(Definition~\ref{cv-origin-reference-lattices}(1), with $T=\operatorname{Spec}\mathbf C$ and $x$ the identity) are
equal.
\end{lemma}
\leanref{Development/AbelianVarietyTheory/Proofs/CubicalLineBundles/PolCountableAndTangentExt.lean}{GlobalTheory.HomDescent.countable_and_tangent_ext}
\leanref{Development/AbelianVarietyTheory/Proofs/CubicalLineBundles/PolCountableAndTangentExt.lean}{GlobalTheory.HomDescent.pol_countable_and_tangent_ext}
\begin{proof}
Step 1: the base changes. The base change of a principally polarized abelian scheme is one, by law (22) of
Definition~\ref{cv-cusp-cm-point-eigen-degree-tensor-laws}, so there are principally
polarized abelian varieties $Q$ and $Q'$ over $\mathbf C$ with isomorphisms of group schemes $e\colon Q\to P_{\mathbf C}$
and $e'\colon Q'\to P'_{\mathbf C}$. A homomorphism $f\colon P_{\mathbf C}\to P'_{\mathbf C}$ gives the homomorphism
$e'^{-1}fe$, and a tangent vector at the origin followed by a homomorphism is a tangent vector at the origin, as
homomorphisms preserve the unit section. So it suffices to prove the lemma for $Q$ and $Q'$ in place of $P_{\mathbf C}$
and $P'_{\mathbf C}$.

Step 2: uniformizations. The line bundles of the polarizations of $Q$ and $Q'$ are trivialized along the zero
sections. By Theorem~\ref{cv-cusp-cm-point-eigen-degree-rank-one-norm-riemann-theta}(2), $Q$ and $Q'$ have
uniformizations $u\colon\mathbf C^g\to Q(\mathbf C)$ and $u'\colon\mathbf C^{g'}\to Q'(\mathbf C)$ for these line bundles
(Definition~\ref{cv-origin-hodge-norm-gram}(1)), with kernels $\Lambda$ and $\Lambda'$, such that every tangent vector at
the origin has exactly one differential and every vector is the differential of exactly one tangent vector at the
origin.

Step 3: linear lifts. Let $f\colon Q\to Q'$ be a homomorphism. Lemma~\ref{cv-origin-hodge-norm-form-lift}, applied to $u'$
and to the homomorphism $f\circ u$, which satisfies its hypothesis as $f$ is a morphism, gives a $\mathbf C$-linear map
$F_f\colon\mathbf C^g\to\mathbf C^{g'}$ with $u'\circ F_f=f\circ u$. For $l\in\Lambda$, $u'(F_f(l))=f(u(l))$ is the
origin, so $F_f(\Lambda)\subset\Lambda'$. By the chain rule, as in the proof of
Lemma~\ref{cv-origin-hodge-norm-form-compare}, if a tangent vector $v$ of $Q$ at the origin has differential $w$ for $u$,
then $f\circ v$ has differential $F_f(w)$ for $u'$. A morphism $f'\colon Q\to Q'$ with $u'\circ F_f=f'\circ u$ is $f$:
the two agree on the complex points, as $u$ is onto, $Q$ is smooth over $\mathbf C$, hence reduced, and $Q'$ is
separated.

Step 4: countability. The group $\Lambda$ is discrete in $\mathbf C^g$, so it is a finitely generated $\mathbf Z$-module;
let $s_1,\dots,s_n$ generate it. In the same way $\Lambda'$ is finitely generated, hence countable. The map
$f\mapsto(F_f(s_1),\dots,F_f(s_n))$ from the homomorphisms $Q\to Q'$ to $\Lambda'^n$ is injective: if $f$ and $f'$ have
the same image, then $F_f$ and $F_{f'}$ agree on $\Lambda$, which spans $\mathbf C^g$ over $\mathbf R$, so $F_f=F_{f'}$,
and $f'=f$ by Step~3. So the homomorphisms form a countable set.

Step 5: tangent vectors. Let $f,f'\colon Q\to Q'$ be homomorphisms with $f\circ v=f'\circ v$ for every tangent vector $v$
at the origin. Let $z\in\mathbf C^g$, and let $v$ be the tangent vector at the origin with differential $z$. By Step~3,
the tangent vector $f\circ v=f'\circ v$ has the differentials $F_f(z)$ and $F_{f'}(z)$, which are equal by the uniqueness
of differentials. So $F_f=F_{f'}$, and $f'=f$ by Step~3.
\end{proof}
% Lean: countable_and_tangent_ext: exists_pol_iso (the law PolBaseChangeSt, with grpPullback_iso) for Q and Q', and
% isTangent_comp (Step 1). pol_countable_and_tangent_ext: the hypothesis ThetaUnifAV for Q and Q' with their
% rigidifications (Step 2); HomLinear.exists_linear_of_isMonHom (Step 3: Lemma cv-origin-hodge-norm-form-lift, the
% kernels, isDiff_comp, and ComplexPoints.over_hom_ext_of_complexPoints with isReduced_of_smooth_complex for the
% uniqueness); u.discrete, Module.Finite.exists_fin, Finsupp.Countable.of_moduleFinite and the span u.full (Step 4); h2,
% h1' and hu' (Step 5).

\begin{lemma}\label{cv-cusp-cm-point-eigen-degree-tensor-descent-points}
Let $k\subset\mathbf C$ be a number field, $R$ a finitely generated $k$-subalgebra of $\mathbf C$ with inclusion
$\iota\colon R\to\mathbf C$, $B$ a $k$-algebra, $Z$ a scheme and $\Phi\colon\operatorname{Spec}(R\otimes_kB)\to Z$ a
morphism that maps every point into an affine open subscheme $U$ of $Z$. For a homomorphism of $k$-algebras
$\sigma\colon R\to\mathbf C$, put $\Phi_\sigma=\Phi\circ\operatorname{Spec}(\sigma\otimes\mathrm{id}_B)$, a morphism
$\operatorname{Spec}(\mathbf C\otimes_kB)\to Z$, and suppose that the $\Phi_\sigma$ form a countable set. Let
$\rho_1\colon R\to\mathbf C$ be a homomorphism of $k$-algebras with $\rho_1(c)=c$ for every $c\in R$ algebraic over $k$.
Then $\Phi_{\rho_1}=\Phi_\iota$.
\end{lemma}
\leanref{Development/AbelianVarietyTheory/Proofs/CubicalLineBundles/PolCountableAndTangentExt.lean}{GlobalTheory.HomDescent.specialise_eq_of_countable}
\begin{proof}
As $\Phi$ maps every point into $U$, it factors through the open immersion $j\colon U\to Z$, and as $U$ is affine,
$\Phi=j\circ\operatorname{Spec}\varphi$ for a ring homomorphism $\varphi\colon\Gamma(U,\mathcal O_Z)\to R\otimes_kB$. Then
$\Phi_\sigma=j\circ\operatorname{Spec}\big((\sigma\otimes\mathrm{id}_B)\circ\varphi\big)$ for every $\sigma$. As $j$ is a
monomorphism and $\operatorname{Spec}$ is faithful, $\Phi_\sigma$ determines $(\sigma\otimes\mathrm{id}_B)\circ\varphi$, so
these homomorphisms form a countable set. Let $(b_i)$ be a basis of the $k$-vector space $B$. The elements $1\otimes b_i$
form bases of $R\otimes_kB$ over $R$ and of $\mathbf C\otimes_kB$ over $\mathbf C$, and the coordinates of
$(\sigma\otimes\mathrm{id}_B)(x)$ are the images under $\sigma$ of those of $x$. Let $y\in\Gamma(U,\mathcal O_Z)$, and let
$c\in R$ be the $i$-th coordinate of $\varphi(y)$. The values $\sigma(c)$ are the $i$-th coordinates of
$(\sigma\otimes\mathrm{id}_B)(\varphi(y))$, so they lie in a countable set, and $c$ is algebraic over $k$ by
Lemma~\ref{cv-cusp-cm-point-eigen-degree-tensor-descent-specialize}(2). So $\rho_1(c)=c=\iota(c)$. Hence
$(\rho_1\otimes\mathrm{id}_B)(\varphi(y))=(\iota\otimes\mathrm{id}_B)(\varphi(y))$ for every $y$, and
$\Phi_{\rho_1}=\Phi_\iota$.
\end{proof}
% Lean: specialise_eq_of_countable: IsOpenImmersion.lift for Phi' through hU.fromSpec, Spec.map_surjective for phi, the
% injective map G (Spec.map_injective and cancel_mono), the basis Module.Basis.ofVectorSpace k B with
% Algebra.TensorProduct.basis and basis_repr_map, and specialisation_numberField_and_countable_values (its second half,
% with r = 1).

\begin{lemma}\label{cv-cusp-cm-point-eigen-degree-tensor-descent-hom}
Assume the conclusion of Theorem~\ref{cv-cusp-cm-point-eigen-degree-rank-one-norm-riemann-theta} and law (22) of
Definition~\ref{cv-cusp-cm-point-eigen-degree-tensor-laws}, on base change.
Let $k\subset\mathbf C$ be a number field and $P$ and $P'$ principally polarized abelian varieties over $k$; for a
$k$-algebra $A$, write $P_A$ and $P'_A$ for their base changes to $A$. Let $f\colon P_{\mathbf C}\to P'_{\mathbf C}$ be a
homomorphism. Then there are a field $K_1$ with $k\subset K_1\subset\mathbf C$ and $[K_1:k]$ finite and a homomorphism
$f_1\colon P_{K_1}\to P'_{K_1}$ whose base change to $\mathbf C$ is $f$.
\end{lemma}
\leanref{Development/AbelianVarietyTheory/Proofs/CubicalLineBundles/PolCountableAndTangentExt.lean}{GlobalTheory.hom_descends_to_numberField}
\begin{proof}
Step 1: spreading out. As $P$ is proper and smooth over $k$, it is quasi-compact, quasi-separated and flat over $k$, and
$P'$ is of finite presentation and separated over $k$. The field $\mathbf C$ is the filtered union of its finitely
generated $k$-subalgebras, so by \cite[Th\'eor\`eme~8.8.2]{Gro66} there are a finitely generated $k$-subalgebra
$R\subset\mathbf C$ and a morphism $f_R\colon P_R\to P'_R$ over $R$ whose base change along the inclusion
$\iota\colon R\to\mathbf C$ is $f$. As $\iota$ is injective, $\operatorname{Spec}\mathbf C\to\operatorname{Spec}R$ is
schematically dominant, and so are its base changes to the flat $R$-schemes $P_R$ and $P_R\times_RP_R$. As $P'_R$ is
separated over $R$, two morphisms from either of these to $P'_R$ over $R$ that agree after base change to $\mathbf C$
are equal. So $f_R$ is a homomorphism, as $f$ is.

Step 2: specialization. By Lemma~\ref{cv-cusp-cm-point-eigen-degree-tensor-descent-specialize}(1), for $L=\mathbf C$ and
$r=1$, there are a field $K_1$ with $k\subset K_1\subset\mathbf C$ and $[K_1:k]$ finite and a homomorphism of
$k$-algebras $\rho\colon R\to K_1$ with $\rho(c)=c$ for every $c\in R$ algebraic over $k$. Let $f_1\colon P_{K_1}\to
P'_{K_1}$ be the base change of $f_R$ along $\rho$, a homomorphism, and let $\rho_1\colon R\to\mathbf C$ be $\rho$
followed by the inclusion $K_1\subset\mathbf C$. For a homomorphism of $k$-algebras $\sigma\colon R\to\mathbf C$, let
$g_\sigma\colon P_{\mathbf C}\to P'_{\mathbf C}$ be the base change of $f_R$ along $\sigma$, a homomorphism. Then
$g_\iota=f$, and the base change of $f_1$ to $\mathbf C$ is $g_{\rho_1}$; it remains to show that $g_{\rho_1}=f$. By
Lemma~\ref{cv-cusp-cm-point-eigen-degree-tensor-descent-countable}, for $S=\operatorname{Spec}k$, the $g_\sigma$ form a
countable set, and it suffices to show that $g_{\rho_1}\circ v=f\circ v$ for every tangent vector $v$ of $P_{\mathbf C}$
at the origin.

Step 3: a family of points. Let $v$ be such a tangent vector, put $B=\mathbf C[\varepsilon]$, a $k$-algebra, and let
$v_0\colon\operatorname{Spec}B\to P$ be $v$ followed by the projection $P_{\mathbf C}\to P$, a morphism over $k$. For a
$k$-algebra $A$, let $x_A\colon\operatorname{Spec}(A\otimes_kB)\to P_A$ be the morphism with components
$v_0\circ\operatorname{Spec}(b\mapsto1\otimes b)$ and $\operatorname{Spec}(a\mapsto a\otimes1)$. For a homomorphism of
$k$-algebras $\sigma\colon A\to A'$, $x_A\circ\operatorname{Spec}(\sigma\otimes\mathrm{id}_B)$ is $x_{A'}$ followed by the
projection $P_{A'}\to P_A$. Let $q\colon P'_R\to P'$ and $q_{\mathbf C}\colon P'_{\mathbf C}\to P'$ be the projections,
and put $\Phi=q\circ f_R\circ x_R\colon\operatorname{Spec}(R\otimes_kB)\to P'$. For every $\sigma\colon R\to\mathbf C$,
$\Phi\circ\operatorname{Spec}(\sigma\otimes\mathrm{id}_B)=q_{\mathbf C}\circ g_\sigma\circ x_{\mathbf C}$, as $g_\sigma$ is
the base change of $f_R$ along $\sigma$. So these morphisms form a countable set.

Step 4: the image of $\Phi$. Let $U$ be an affine open subscheme of $P'$ containing the origin. Every point of
$\operatorname{Spec}(R\otimes_kB)$ is the image of a point of $\operatorname{Spec}(R\otimes_k\mathbf C)$, as the kernel of
the map $R\otimes_kB\to R\otimes_k\mathbf C$ given by $\varepsilon\mapsto0$ consists of elements of square zero. On
$\operatorname{Spec}(R\otimes_k\mathbf C)$, $x_R$ is the unit section of $P_R$ composed with the structure morphism, as
$v$ restricts to the origin, and $f_R$ preserves the unit section. So $\Phi$ maps every point to the origin of $P'$, which
lies in $U$.

Step 5: the comparison. By Lemma~\ref{cv-cusp-cm-point-eigen-degree-tensor-descent-points}, for $Z=P'$, $U$, $\Phi$ and
$\rho_1$, $q_{\mathbf C}\circ g_{\rho_1}\circ x_{\mathbf C}=q_{\mathbf C}\circ f\circ x_{\mathbf C}$. Let
$m\colon\operatorname{Spec}B\to\operatorname{Spec}(\mathbf C\otimes_kB)$ be given by $a\otimes b\mapsto ab$. Then
$x_{\mathbf C}\circ m=v$, as both have the components $v_0$ and $\operatorname{Spec}(\mathbf C\to B)$. So
$q_{\mathbf C}\circ g_{\rho_1}\circ v=q_{\mathbf C}\circ f\circ v$. As $g_{\rho_1}\circ v$ and $f\circ v$ are morphisms over
$\mathbf C$ to $P'_{\mathbf C}=P'\times_k\operatorname{Spec}\mathbf C$, they are equal.
\end{proof}
\leanref{Development/AbelianVarietyTheory/Proofs/CubicalLineBundles/PolCountableAndTangentExt.lean}{GlobalTheory.HomDescent.exists_fg_isMonHom}
\leanref{Development/AbelianVarietyTheory/Proofs/CubicalLineBundles/PolCountableAndTangentExt.lean}{GlobalTheory.HomDescent.tangent_range}
\leanref{Development/AbelianVarietyTheory/Proofs/CubicalLineBundles/PolCountableAndTangentExt.lean}{GlobalTheory.HomDescent.tangent_factor}
% Lean: hom_descends_to_numberField, under the hypotheses ThetaUnifAV and PolBaseChangeSt. Step 1 is exists_fg_isMonHom
% (Spread.exists_fg_comp_eq, from Mathlib's Scheme.exists_pi_app_comp_eq_of_locallyOfFinitePresentation, the case of EGA
% IV 8.8.2 used here, and ExtUnique.isMonHom_of_isGrpPullback with isSchemeTheoreticallyDominant_spec_of_injective).
% Step 2 is exists_specialisation_finite with r = 1, exists_specialise for f_1 and for g, and countable_and_tangent_ext.
% Step 3 is ptBC, spMap_map_ptBC and the equations e. Step 4 is tangent_range (surjective_spMap_map_fst, eSec_comm and
% IsMonHom.one_hom) with exists_isAffineOpen_mem_and_subset. Step 5 is specialise_eq_of_countable, tangent_factor and
% pullback.hom_ext. Size: PolCountableAndTangentExt.lean, 409 lines, with SpecialisationNumberFieldAndCountableValues.lean
% (184) and MorphismsSpreadOutToFgSubalgebra.lean (336), beyond Lemma cv-origin-hodge-norm-form-lift.

\begin{lemma}\label{cv-cusp-cm-point-eigen-degree-tensor-descent}
Assume the conclusion of Theorem~\ref{cv-cusp-cm-point-eigen-degree-rank-one-norm-riemann-theta} and the laws of
abelian schemes (Definition~\ref{cv-cusp-cm-point-eigen-degree-tensor-laws}); of the laws, the proof uses only (22).
In Situation~\ref{cv-cusp-setting}, let $\mathcal O$, $L$, $M$ and $h$ be as in
Definition~\ref{cv-cusp-cm-point-eigen-degree-tensor-family}, with $L$ finitely generated. Let $k\subset\mathbf C$ be a
number field, $C$ a principally polarized elliptic curve over $k$ with an action of $\mathcal O$, and $A$ a principally
polarized abelian variety over $k$ with an action of $\mathcal O_E$, and let $(t_a)_{a\in M}$ be a family of
homomorphisms from the base change of $C$ to $\mathbf C$ to that of $A$ for $h$ on $M$. Then there is a number field
$k_1$ with $k\subset k_1\subset\mathbf C$ such that for every number field $k_2$ with $k_1\subset k_2\subset\mathbf C$
there is a family of homomorphisms $(t'_a)_{a\in M}$ from the base change of $C$ to $k_2$ to that of $A$ for $h$ on $M$
whose base change to $\mathbf C$ is $(t_a)_{a\in M}$.
\end{lemma}
\leanref{Development/AbelianVarietyTheory/Proofs/CubicalLineBundles/CmFamilyDescendsToNumberFieldStatement.lean}{GlobalTheory.T.cmFamily_descends_to_numberField}
\leanref{Development/AbelianVarietyTheory/Proofs/CubicalLineBundles/CmFamilyDescendsToNumberField.lean}{GlobalTheory.CM.isFam_descends_to_numberField}
\begin{proof}
Step 1: one homomorphism. For each $a\in M$ there is a number field $K_a$ with $k\subset K_a\subset\mathbf C$ such that
$t_a$ is the base change of a homomorphism from $C_{K'}$ to $A_{K'}$, for every number field $K'$ with $K_a\subset
K'\subset\mathbf C$; here $C_{K'}$ and $A_{K'}$ are the base changes to $K'$. Indeed,
Lemma~\ref{cv-cusp-cm-point-eigen-degree-tensor-descent-hom}, which applies by law (22) of
Definition~\ref{cv-cusp-cm-point-eigen-degree-tensor-laws}, for $C$, $A$ and the homomorphism $t_a$, gives a number
field $K_a$ with $k\subset K_a\subset\mathbf C$ and a homomorphism $f_a\colon C_{K_a}\to A_{K_a}$ whose base change to
$\mathbf C$ is $t_a$. For $K'$ as above, the base change of $f_a$ to $K'$ is a homomorphism whose base change to
$\mathbf C$ is $t_a$, as base change along $K_a\subset K'$ followed by base change along $K'\subset\mathbf C$ is base
change along $K_a\subset\mathbf C$.
% Lean: FamDesc.exists_field: hom_descends_to_numberField (Lemma cv-cusp-cm-point-eigen-degree-tensor-descent-hom) for
% e^{-1} t_a d, with e and d the comparisons (grpPullback_iso) of the given base changes to C with the pullbacks; the
% homomorphism over k_2 is the lift c_2.lift through the given base change T_2, and the equations key, a_1 and r compare
% it with t_a.

Step 2: the field $k_1$. The $\mathcal O_E$-module $L$ is finitely generated, hence so is the group $M$, and $M$ is
generated as a monoid by a finite set $S$. Let $k_1$ be the number field generated by $k$, $K_0$ and the $K_a$ for
$a\in S$.

Step 3: all homomorphisms. Let $k_2$ be a number field with $k_1\subset k_2\subset\mathbf C$. For $a\in S\cup\{0\}$,
$t_a$ is the base change of a homomorphism $t'_a$ over $k_2$ by Step~1. If $t_a$ and $t_{a'}$ are the base changes of
$t'_a$ and $t'_{a'}$, then $t_{a+a'}=t_a+t_{a'}$ is the base change of $t'_a+t'_{a'}$. As $S$ generates $M$ as a
monoid, every $t_a$ is the base change of a homomorphism $t'_a$ over $k_2$.

Step 4: the family. A field is a Dedekind domain, and $k_2\to\mathbf C$ is injective. So
Lemma~\ref{cv-cusp-cm-point-eigen-degree-tensor-identities}, for $R=k_2$ and the base changes of $C$ and $A$ to $k_2$,
shows that $(t'_a)_{a\in M}$ is a family of homomorphisms for $h$ on $M$.
\end{proof}
\leanref{Development/AbelianVarietyTheory/Proofs/CubicalLineBundles/CmFamilyDescendsToNumberField.lean}{GlobalTheory.CM.FamDesc.exists_field}
\leanref{Development/AbelianVarietyTheory/Proofs/CubicalLineBundles/CmFamilyDescendsToNumberField.lean}{GlobalTheory.CM.FamDesc.exists_mul}
% Lean: isFam_descends_to_numberField: exists_field for each a (Step 1, with choose), the finite generating set S of
% hg.fg_top (Step 2), k_1 = k + K 0 + sup of K a over S, closure_induction with exists_mul (Step 3) and
% T.family_identities_descend_field_to_dedekind (Step 4). The Lean's statement also fixes the comparisons of the base
% changes (EllMap, SABMap), which the text leaves implicit.

\begin{lemma}\label{cv-cusp-cm-point-eigen-degree-tensor-extend}
Assume the laws of abelian schemes (Definition~\ref{cv-cusp-cm-point-eigen-degree-tensor-laws}) and the conclusion of
Theorem~\ref{cv-model-toroidal-extension}; the laws are not used in the proof, except to apply
Lemma~\ref{cv-cusp-cm-point-eigen-degree-tensor-identities}, which does not use them.
In Situation~\ref{cv-cusp-setting}, let $\mathcal O$, $L$, $M$ and $h$ be as in
Definition~\ref{cv-cusp-cm-point-eigen-degree-tensor-family}. Let $R$ be a Dedekind domain with fraction field $K$,
$C$ a principally polarized elliptic curve over $R$ with an action of $\mathcal O$, and $A$ a principally polarized
abelian scheme over $R$ with an action of $\mathcal O_E$, and let $(t'_a)_{a\in M}$ be a family of homomorphisms from
$C_K$ to $A_K$ for $h$ on $M$. Then there is a family of homomorphisms $(t_a)_{a\in M}$ from $C$ to $A$ for $h$ on $M$
whose restriction to $C_K$ is $(t'_a)_{a\in M}$.
\end{lemma}
\leanref{Development/AbelianVarietyTheory/Proofs/CubicalLineBundles/CmFamilyExtendsOverDedekindDomainStatement.lean}{GlobalTheory.T.cmFamily_extends_over_dedekindDomain}
\leanref{Development/AbelianVarietyTheory/Proofs/CubicalLineBundles/CmFamilyExtendsOverDedekindDomain.lean}{GlobalTheory.isFam_extends_dedekind}
\begin{proof}
Each homomorphism $t'_a$ of abelian varieties over $K$ extends to a homomorphism $t_a\colon C\to A$ of abelian schemes
over $R$ \cite[Chapter~I, Proposition~2.7]{FC90}; over the localization of $R$ at a nonzero prime, a discrete valuation
ring, this is the assumed conclusion of Theorem~\ref{cv-model-toroidal-extension}. The base change of $t_a$ along
$R\to K$ is $t'_a$, and $R\to K$ is
injective, so Lemma~\ref{cv-cusp-cm-point-eigen-degree-tensor-identities} shows that $(t_a)_{a\in M}$ is a family of
homomorphisms for $h$ on $M$.
\end{proof}
% Lean: isFam_extends_dedekind: abelianScheme_hom_extends_dedekind_of_isGrpPullback for each a, which the Lean derives
% from its case of a discrete valuation ring, the hypothesis HomExtDvr (Theorem cv-model-toroidal-extension), and
% T.family_identities_descend_field_to_dedekind.

\begin{lemma}\label{cv-cusp-cm-point-eigen-degree-tensor-recognition-tensor}
Assume the laws of abelian schemes (Definition~\ref{cv-cusp-cm-point-eigen-degree-tensor-laws}).
In Situation~\ref{cv-cusp-setting}, let $\sqrt{-d}\in E$ be the square root of $-d$ of the datum, and put
$w_0=(1+\sqrt{-d})/2$ and $\mathcal O=\mathbf Z[w_0]$. Let $R$ be a noetherian ring, $\omega\in R$, and $C$ a
principally polarized elliptic curve over $R$ with an action of $\mathcal O$ whose Rosati involution is complex
conjugation and a basis of its Lie algebra on which $w_0$ acts by $\omega$. Let $L$ be the lattice of a PEL datum over $E$
of signature $(p,q)$ with its perfect alternating form $\langle\ ,\ \rangle$, let $\varepsilon\in\{\pm1\}$ with
$\varepsilon\langle w_0x,x\rangle>0$ for every nonzero $x\in L$, and put $h(x,y)=\varepsilon(\langle
w_0x,y\rangle-\langle x,y\rangle)+\varepsilon\langle x,y\rangle w_0$. Then there is a principally polarized abelian
scheme $A_0$ over $R$ of relative dimension $\frac{[E:\mathbf Q]}2(p+q)$, with an action of $\mathcal O_E$ for which the
Rosati involution of the polarization induces complex conjugation on $\mathcal O_E$, that is a Serre tensor product of
$C$ and $L$ with form $h$ (Definition~\ref{cv-cusp-cm-point-eigen-degree-serre}).
\end{lemma}
\leanref{Development/AbelianVarietyTheory/Proofs/CubicalLineBundles/ExistsCmTensorLatticeCurvePrincipalTripleStatement.lean}{GlobalTheory.T.exists_cm_tensor_lattice_curve_principal_triple}
\leanref{Development/AbelianVarietyTheory/Proofs/CubicalLineBundles/ExistsCmTensorLatticeCurvePrincipalTriple.lean}{GlobalTheory.CM.exists_isTen}
\begin{proof}
\openstep{Serre's construction of the tensor product of the curve with the lattice, its line bundle with the form $h$,
and the proofs that this line bundle is relatively ample and defines a principal polarization, are not yet written
out.}
% Lean: exists_isTen, under Out.Laws: T.exists_unique_cmHermitianLatticeForm (= CM.exists_kForm, the form h_0 of D with
% values in O, equal to h by its uniqueness); OK_isDedekindDomain_and_projective (O is a Dedekind domain, L is finite
% projective over it); Serre.isTensor (Serre's tensor construction L tensor_O C, in PowObjGeometricallyConnected.lean,
% 1,445 lines) with the action of O_E (algAct) and IsTensor.isAbelianScheme; the cubical structure of the line bundle of
% C (laws.cube); TensorBundle.exists_tensor_bundle (a line bundle N on the tensor product with Theta(N) trivial and
% LamEq for the pure tensors, 163 lines); PrimNorm.exists_primitive_normalisation (the rigidified, cubical, symmetric
% and primitive line bundle with the Rosati condition, 397 lines, with LamCalc.lamEq_of_tmul); and the fields of the PAT:
% T.cmTensor_smoothOfRelativeDimension (= CM.tensor_smoothOfRelativeDimension, 96 lines),
% T.relAmple_lineBundle_of_positive_hermitian_cm_form (= CM.relAmple_of_kForm_pos, 112 lines, which uses the
% positivity) and PerfectTensorKernel.kTriv_of_perfect (64 lines, the principal polarization from the perfectness of
% the form). Size: about 4,300 lines beyond the laws of the moduli problem, most of the step.
\end{proof}

\begin{lemma}\label{cv-cusp-cm-point-eigen-degree-tensor-recognition-cube}
Let $S$ be a scheme, $\mathcal O$ a ring, $G$ a commutative group scheme over $S$ with an action of $\mathcal O$, and $A$
and $A'$ commutative group schemes over $S$ with actions of $\mathcal O$. Let $M$ and $M'$ be $\mathcal O$-modules with
isomorphisms of $\mathcal O$-modules $M\otimes_{\mathcal O}G(T)\to A(T)$ and $M'\otimes_{\mathcal O}G(T)\to A'(T)$, natural
in the $S$-scheme $T$, and let $t_m\colon G\to A$ and $t'_{m'}\colon G\to A'$ be the homomorphisms $x\mapsto m\otimes x$
and $x\mapsto m'\otimes x$. Let $B$ and $B'$ be commutative group schemes over $S$ with line bundles $\mathcal L$ and
$\mathcal L'$ such that $\Theta(\mathcal L)$ and $\Theta(\mathcal L')$ are trivial, and let $f_1\colon A\to B$,
$f_2\colon A'\to B$, $g_1\colon A\to B'$ and $g_2\colon A'\to B'$ be homomorphisms. Suppose that
$\Lambda_{\mathcal L}(f_1t_mp_1,f_2t'_{m'}p_2)\cong\Lambda_{\mathcal L'}(g_1t_mp_1,g_2t'_{m'}p_2)$ on $G\times_SG$ for
all $m\in M$ and $m'\in M'$, with $p_1,p_2$ the projections. Then
$\Lambda_{\mathcal L}(f_1p_1,f_2p_2)\cong\Lambda_{\mathcal L'}(g_1p_1,g_2p_2)$ on $A\times_SA'$.
\end{lemma}
\leanref{Development/AbelianVarietyTheory/Proofs/CubicalLineBundles/TensorPairingIdentityOfPureTensorsStatement.lean}{GlobalTheory.T.tensor_pairing_identity_of_pure_tensors}
\leanref{Development/AbelianVarietyTheory/Proofs/CubicalLineBundles/TensorPairingIdentityOfPureTensors.lean}{GlobalTheory.LamCalc.lamEq_of_tmul}
\begin{proof}
Here $\Theta(\mathcal L)$ is the line bundle on $B\times_SB\times_SB$ whose fibre at $(x_1,x_2,x_3)$ is
$\bigotimes_I\mathcal L(x_I)^{\otimes(-1)^{|I|+1}}$, over the subsets $I\subset\{1,2,3\}$, with $x_I=\sum_{i\in I}x_i$,
and $\Lambda$ is as in Definition~\ref{cv-toroidal-abelian-dvr-ample-torsion-weil}. Expanding the definitions, for points
$x,x',y$ of $B$ with values in an $S$-scheme $T$,
\[
 \Theta(\mathcal L)(x,x',y)\cong\Lambda_{\mathcal L}(x+x',y)\otimes\Lambda_{\mathcal L}(x,y)^{-1}\otimes\Lambda_{\mathcal L}(x',y)^{-1},
\]
so $\Lambda_{\mathcal L}(x+x',y)\cong\Lambda_{\mathcal L}(x,y)\otimes\Lambda_{\mathcal L}(x',y)$, and in the same way in the
second variable, as $\Lambda_{\mathcal L}(x,y)=\Lambda_{\mathcal L}(y,x)$; the same holds for $\mathcal L'$.

For an $S$-scheme $T$ and points $x\in A(T)$ and $y\in A'(T)$, say that $(x,y)$ is good if
$\Lambda_{\mathcal L}(f_1x,f_2y)\cong\Lambda_{\mathcal L'}(g_1x,g_2y)$ on $T$. If $(x,y)$ and $(x',y)$ are good, so is
$(x+x',y)$: as $f_1$ and $g_1$ are homomorphisms, $f_1(x+x')=f_1x+f_1x'$ and $g_1(x+x')=g_1x+g_1x'$, and the
isomorphisms above give $\Lambda_{\mathcal L}(f_1(x+x'),f_2y)\cong\Lambda_{\mathcal L}(f_1x,f_2y)\otimes
\Lambda_{\mathcal L}(f_1x',f_2y)$ and the same for $\mathcal L'$. In the same way, if $(x,y)$ and $(x,y')$ are good, so is
$(x,y+y')$. For $x=t_m(u)$ and $y=t'_{m'}(v)$ with $u,v\in G(T)$, $(x,y)$ is good: pull back the isomorphism of the
hypothesis along $(u,v)\colon T\to G\times_SG$. By the natural isomorphisms, every point of $A(T)$ is a finite sum of
points $t_m(u)$, and $0=t_0(0)$; the same holds for $A'(T)$. So every pair $(x,y)$ is good. For $T=A\times_SA'$, $x=p_1$
and $y=p_2$, this is the claim.
\end{proof}
\leanref{Development/AbelianVarietyTheory/Proofs/CubicalLineBundles/LbIsoCubeLBMulRightOf.lean}{GlobalTheory.LamCalc.lbIso_cubeLB_mul_right_of}
% Lean: lamEq_of_tmul: isTensor_induction (Serre.IsTensor.induction_on, with t.tmul_zero for the base case) twice;
% lbIso_of_lamEq (the pullback along (u, v)); lbIso_cubeLB_mul_left_of and lbIso_cubeLB_mul_right_of (the additivity
% of Lambda in each variable from the triviality of Theta, LbIsoCubeLBMulRightOf.lean, 178 lines, a computation with
% the classes of the line bundles); and key at (fst A A', snd A A').

\begin{lemma}\label{cv-cusp-cm-point-eigen-degree-tensor-recognition-lift}
In Situation~\ref{cv-cusp-setting}, let $\sqrt{-d}\in E$ be the square root of $-d$ of the datum, and put
$w_0=(1+\sqrt{-d})/2$ and $\mathcal O=\mathbf Z[w_0]$. Let $R$ be a ring, $\omega\in R$, and $C$ a principally polarized
elliptic curve over $R$ with an action of $\mathcal O$ whose Rosati involution is complex conjugation and a basis of its
Lie algebra on which $w_0$ acts by $\omega$, with the line bundle $\mathcal L_C$ of its polarization. Let $L_1$ and $L$
be the lattices of PEL data over $E$ with their perfect alternating forms $\langle\ ,\ \rangle_1$ and
$\langle\ ,\ \rangle$, let $\varepsilon\in\{\pm1\}$, and let $h_1$ and $h$ be the forms
$h_1(x,y)=\varepsilon(\langle w_0x,y\rangle_1-\langle x,y\rangle_1)+\varepsilon\langle x,y\rangle_1w_0$ and
$h(x,y)=\varepsilon(\langle w_0x,y\rangle-\langle x,y\rangle)+\varepsilon\langle x,y\rangle w_0$. Let $A_1$ and $A$ be
principally polarized abelian schemes over $R$ with actions of $\mathcal O_E$ for which the Rosati involutions induce
complex conjugation on $\mathcal O_E$, with the line bundles $\mathcal L_{A_1}$ and $\mathcal L_A$ of their
polarizations, such that $A_1$ is a Serre tensor product of $C$ and $L_1$ with form $h_1$, with homomorphisms
$t^1_a\colon C\to A_1$ (Definition~\ref{cv-cusp-cm-point-eigen-degree-serre}). Let $M\subset L$ be an
$\mathcal O_E$-submodule, $(t_a)_{a\in M}$ a family of homomorphisms from $C$ to $A$ for $h$ on $M$
(Definition~\ref{cv-cusp-cm-point-eigen-degree-tensor-family}), $f\colon L_1\to M$ an $\mathcal O_E$-linear map and
$c\in\mathcal O_E$ with $\langle f(a),f(b)\rangle=\langle ca,b\rangle_1$ for all $a,b\in L_1$. Then there is a
homomorphism $\Phi\colon A_1\to A$ that commutes with the actions of $\mathcal O_E$, with $\Phi\circ t^1_a=t_{f(a)}$ for
every $a\in L_1$, and with $\Lambda_{\mathcal L_A}(\Phi p_1,\Phi p_2)\cong\Lambda_{\mathcal L_{A_1}}(cp_1,p_2)$ on
$A_1\times_RA_1$.
\end{lemma}
\leanref{Development/AbelianVarietyTheory/Proofs/CubicalLineBundles/ExistsTensorProductHomOfCmCurveHomFamilyStatement.lean}{GlobalTheory.T.exists_tensorProduct_hom_of_cmCurveHomFamily}
\leanref{Development/AbelianVarietyTheory/Proofs/CubicalLineBundles/ExistsTensorProductHomOfCmCurveHomFamily.lean}{GlobalTheory.CM.exists_famLift}
\begin{proof}
Step 1: the forms. For $a,b\in L_1$, $h(f(a),f(b))=h_1(ca,b)$. Indeed, $f$ is $\mathcal O_E$-linear, so
$\langle w_0f(a),f(b)\rangle=\langle f(w_0a),f(b)\rangle=\langle cw_0a,b\rangle_1=\langle w_0(ca),b\rangle_1$, and
$\langle f(a),f(b)\rangle=\langle ca,b\rangle_1$.

Step 2: the homomorphism. For an $R$-scheme $T$, the map $L_1\times C(T)\to A(T)$, $(a,x)\mapsto t_{f(a)}(x)$, is
additive in each variable, by (2) of Definition~\ref{cv-cusp-cm-point-eigen-degree-tensor-family} and as each
$t_{f(a)}$ is a homomorphism. It is $\mathcal O$-balanced: for $b\in\mathcal O$, $t_{f(ba)}(x)=b\,t_{f(a)}(x)=t_{f(a)}(bx)$,
by (2) and (1) of that definition. So it induces a homomorphism of $\mathcal O$-modules
$L_1\otimes_{\mathcal O}C(T)\to A(T)$, natural in $T$. Composed with the inverse of the natural isomorphism
$L_1\otimes_{\mathcal O}C(T)\to A_1(T)$, it gives homomorphisms $A_1(T)\to A(T)$, natural in $T$, so by the Yoneda lemma a
homomorphism $\Phi\colon A_1\to A$ with $\Phi\circ t^1_a=t_{f(a)}$ for every $a$.

Step 3: the actions. Two homomorphisms from $A_1$ that agree after composition with every $t^1_a$ are equal, as every
point of $A_1$ is a finite sum of points $t^1_a(x)$. For $s\in\mathcal O_E$ and $a\in L_1$,
$\Phi\circ s\circ t^1_a=\Phi\circ t^1_{sa}=t_{f(sa)}=s\circ t_{f(a)}=s\circ\Phi\circ t^1_a$, by (1) of
Definition~\ref{cv-cusp-cm-point-eigen-degree-serre} and (2) of
Definition~\ref{cv-cusp-cm-point-eigen-degree-tensor-family}. So $\Phi$ commutes with the actions of $\mathcal O_E$.

Step 4: the line bundles. The line bundles $\Theta(\mathcal L_A)$ and $\Theta(\mathcal L_{A_1})$ are trivial, by the
cubical structures of the polarizations. For $a,b\in L_1$, by (3) of
Definition~\ref{cv-cusp-cm-point-eigen-degree-tensor-family}, Step~1, and (1) and (2) of
Definition~\ref{cv-cusp-cm-point-eigen-degree-serre} for $A_1$,
\[
 \Lambda_{\mathcal L_A}(\Phi t^1_ap_1,\Phi t^1_bp_2)=\Lambda_{\mathcal L_A}(t_{f(a)}p_1,t_{f(b)}p_2)
 \cong\Lambda_{\mathcal L_C}(h_1(ca,b)p_1,p_2)\cong\Lambda_{\mathcal L_{A_1}}(t^1_{ca}p_1,t^1_bp_2)
 =\Lambda_{\mathcal L_{A_1}}(c\,t^1_ap_1,t^1_bp_2)
\]
on $C\times_RC$. Lemma~\ref{cv-cusp-cm-point-eigen-degree-tensor-recognition-cube}, for $G=C$, $A=A'=A_1$, $M=M'=L_1$,
$B=A$, $B'=A_1$, $f_1=f_2=\Phi$, $g_1=c$ and $g_2$ the identity, gives
$\Lambda_{\mathcal L_A}(\Phi p_1,\Phi p_2)\cong\Lambda_{\mathcal L_{A_1}}(cp_1,p_2)$.
\end{proof}
% Lean: exists_famLift: IsKForm.map_eq (Step 1); F : D_1.La -> Pts with t.ten.lift F hF, t.ten.isLinearHom_lift and
% t.ten.tmul_lift (Step 2: the Lean's Serre tensor product has the universal property as part of its structure);
% t.ten.hom_ext with t.lin and hsmul (Step 3); T.tensor_pairing_identity_of_pure_tensors (Lemma
% cv-cusp-cm-point-eigen-degree-tensor-recognition-cube) with the cubical structures T.t.cub and T_1.t.cub, hlam and
% t.lam (Step 4).

\begin{lemma}\label{cv-cusp-cm-point-eigen-degree-tensor-recognition-onto}
Let $R$ be a domain, $\sigma\colon R\to\mathbf C$ an injective homomorphism, $A$ and $B$ abelian schemes over $R$, and
$\Phi\colon B\to A$ a morphism over $R$ such that every point of $A$ over $\mathbf C$ through $\sigma$ is $\Phi\circ x$ for
a point $x$ of $B$ over $\mathbf C$ through $\sigma$. Then $\Phi$ is surjective, and for every algebraically closed field
$k$ and every morphism $t\colon\operatorname{Spec}k\to\operatorname{Spec}R$, every point of $A$ over $k$ through $t$ is
$\Phi\circ x$ for a point $x$ of $B$ over $k$ through $t$.
\end{lemma}
\leanref{Development/AbelianVarietyTheory/Proofs/CubicalLineBundles/AbelianSchemeSurjectiveOfOntoComplexPointsStatement.lean}{GlobalTheory.T.abelianScheme_surjective_of_onto_complex_points}
\leanref{Development/AbelianVarietyTheory/Proofs/CubicalLineBundles/AbelianSchemeSurjectiveOfOntoComplexPoints.lean}{GlobalTheory.OntoSurj.surjective_of_onto_complex_points}
\begin{proof}
Let $a\colon A\to\operatorname{Spec}R$ be the structure morphism. As $B$ is proper over $R$ and $A$ is separated, $\Phi$ is
proper, so its image $Z$ is closed. Let $A_\sigma=A\times_R\operatorname{Spec}\mathbf C$, with projection
$\pi\colon A_\sigma\to A$. It is locally of finite type over $\mathbf C$, so it is a Jacobson scheme, and its closed
points have residue fields algebraic over $\mathbf C$, hence equal to $\mathbf C$ \cite[Tag 0478]{Sta}; in a Jacobson
scheme every closed subset is the closure of the closed points it contains \cite[Tag 01P2]{Sta}, \cite[Tag 005U]{Sta}.
For a closed point $z$ of $A_\sigma$, the
point $\operatorname{Spec}\mathbf C\to A_\sigma\to A$ lies over $\sigma$, so it is $\Phi\circ x$ for a point $x$ of $B$,
and $\pi(z)\in Z$. So the closed set $\pi^{-1}(Z)$ contains every closed point of $A_\sigma$, hence all of $A_\sigma$. As the
image of $\pi$ is the inverse image under $a$ of the image of $\operatorname{Spec}\sigma$, $Z$ contains
$a^{-1}(\operatorname{Spec}\sigma(\operatorname{Spec}\mathbf C))$. The image of $\operatorname{Spec}\sigma$ contains the
generic point of $\operatorname{Spec}R$, as $\sigma$ is injective, so it is dense; and $a$ is open, as $A$ is smooth over
$R$ \cite[Tag 01UA]{Sta}, so the inverse image of a dense set is dense. So the closed set $Z$ is all of $A$, and $\Phi$ is
surjective.

Let $k$ be algebraically closed and $y$ a point of $A$ over $k$. The fibre of $\Phi$ over $y$ is a nonempty scheme
locally of finite type over $k$, as $\Phi$ is surjective and of finite type; it has a closed point, whose residue field is
finite over $k$ \cite[Corollary~5.24]{AM69}, hence $k$. This point gives a point $x$ of $B$ over $k$ with $y=\Phi\circ x$.
\end{proof}
\leanref{Development/AbelianVarietyTheory/Proofs/CubicalLineBundles/ProdSurjectiveGeometricPoints.lean}{GlobalTheory.GeomSurj.exists_lift_of_surjective}
% Lean: surjective_of_onto_complex_points: IsProper.of_comp for Phi, surjective_of_points (the density of the image of
% Spec sigma through PrimeSpectrum.denseRange_comap_iff_ker_le_nilRadical, LocallyOfFiniteType.jacobsonSpace for the
% fibre A_sigma, pointOfClosedPoint, closure_closedPoints, Scheme.Pullback.range_fst and the openness of A.hom), and
% GeomSurj.exists_lift_of_surjective (the fibre over y, a closed point with residue field finite over k by
% finite_of_finite_type_of_isJacobsonRing, and IsAlgClosed.lift).

\begin{lemma}\label{cv-cusp-cm-point-eigen-degree-tensor-recognition-iso}
Let $S$ be a reduced scheme, $A$ and $B$ abelian schemes over $S$ with line bundles $\mathcal L$ on $A$ and $\mathcal M$ on
$B$, and $\Phi\colon A\to B$ a surjective homomorphism. Suppose that the kernel of $\mathcal L$ is trivial: a point $x$ of
$A$ with values in an affine $S$-scheme $T$ is $0$ as soon as $t_x^*\mathcal L\otimes\mathcal L^{-1}$, the line bundle
$(xp_1+p_2)^*\mathcal L\otimes p_2^*\mathcal L^{-1}$ on $T\times_SA$, is trivial over the inverse images of the members of
an open cover of $T$. If $\Lambda_{\mathcal M}(\Phi p_1,\Phi p_2)\cong\Lambda_{\mathcal L}(p_1,p_2)$ on $A\times_SA$,
then $\Phi$ is an isomorphism.
\end{lemma}
\leanref{Development/AbelianVarietyTheory/Proofs/CubicalLineBundles/IsIsoOfSurjectivePrincipalPairingStatement.lean}{GlobalTheory.T.isIso_of_surjective_principalPairing}
\leanref{Development/AbelianVarietyTheory/Proofs/CubicalLineBundles/IsIsoOfSurjectivePrincipalPairing.lean}{GlobalTheory.CM.PairingIso.isIso_of_lamEq}
\begin{proof}
Step 1: the kernel. Let $T$ be an $S$-scheme and $x\in A(T)$ with $\Phi x=0$. Pulling the isomorphism of the
hypothesis back along $(xp_1,p_2)\colon T\times_SA\to A\times_SA$ gives $\Lambda_{\mathcal M}(0,\Phi
p_2)\cong\Lambda_{\mathcal L}(xp_1,p_2)$ on $T\times_SA$. The left side is $(\Phi p_2)^*\mathcal M\otimes0^*\mathcal
M^{-1}\otimes(\Phi p_2)^*\mathcal M^{-1}\otimes0^*\mathcal M$, which is trivial. So
$\Lambda_{\mathcal L}(xp_1,p_2)=(xp_1+p_2)^*\mathcal L\otimes(xp_1)^*\mathcal L^{-1}\otimes p_2^*\mathcal L^{-1}\otimes
0^*\mathcal L$ is trivial. The line bundles $(xp_1)^*\mathcal L$ and $0^*\mathcal L$ on $T\times_SA$ are pulled back from
$T$, so they are trivial over the inverse images of the members of an open cover of $T$, and hence so is
$(xp_1+p_2)^*\mathcal L\otimes p_2^*\mathcal L^{-1}$. By the hypothesis on $\mathcal L$, applied to the restrictions of
$x$ to the affine open subschemes of $T$, $x=0$.

Step 2: $\Phi$ is a monomorphism. If $x,y\in A(T)$ have $\Phi x=\Phi y$, then $\Phi(x-y)=0$, so $x=y$ by Step~1.

Step 3: $\Phi$ is an isomorphism. As $A$ is proper over $S$ and $B$ is separated, $\Phi$ is proper. A proper monomorphism
is a closed immersion \cite[Tag 04XV]{Sta}. The scheme $B$ is smooth over the reduced scheme $S$, so it is
reduced, and a surjective closed immersion into a reduced scheme is an isomorphism.
\end{proof}
% Lean: isIso_of_lamEq: eq_one_of_comp_eq_one (Step 1: h' is the hypothesis pulled back along (x p_1, p_2) through
% PrimNorm.chartSquare, g_1 the triviality of Lambda_M(0, -), exists_gen_of_lbIso, locGen_base for the bundles pulled
% back from T, and KerPow.eq_one_of_locGen, which applies the hypothesis KTriv on the affine opens of T); mul_inv_eq_one
% (Step 2); IsProper.of_comp, IsClosedImmersion.iff_isProper_and_mono, isReduced_of_smooth_of_isReduced and
% isIso_of_isClosedImmersion_of_surjective (Step 3).

\begin{lemma}\label{cv-cusp-cm-point-eigen-degree-tensor-recognition}
Assume the laws of abelian schemes (Definition~\ref{cv-cusp-cm-point-eigen-degree-tensor-laws}).
In Situation~\ref{cv-cusp-setting}, let $\sqrt{-d}\in E$ be the square root of $-d$ of the datum, and put
$w_0=(1+\sqrt{-d})/2$ and $\mathcal O=\mathbf Z[w_0]$. Let $L$ be the lattice of a PEL datum over $E$ with its perfect
alternating form $\langle\ ,\ \rangle$, let $\varepsilon\in\{\pm1\}$ with $\varepsilon\langle w_0x,x\rangle>0$ for every
nonzero $x\in L$, and put $h(x,y)=\varepsilon(\langle w_0x,y\rangle-\langle x,y\rangle)+\varepsilon\langle
x,y\rangle w_0$. Let $R$ be a noetherian domain and $\sigma\colon R\to\mathbf C$ an injective homomorphism, let $C$ be a
principally polarized elliptic curve over $R$ with an action of $\mathcal O$ whose Rosati involution is complex
conjugation, and let $A$ be a principally polarized abelian scheme over $R$ with an action of $\mathcal O_E$. Let
$(t_a)_{a\in L}$ be a family of homomorphisms from $C$ to $A$ for $h$ on $L$
(Definition~\ref{cv-cusp-cm-point-eigen-degree-tensor-family}) that is onto at $\sigma$. Then $A$ is a Serre tensor
product of $C$ and $L$ with form $h$ (Definition~\ref{cv-cusp-cm-point-eigen-degree-serre}).
\end{lemma}
\leanref{Development/AbelianVarietyTheory/Proofs/CubicalLineBundles/NonemptyTensorOfOntoFullFamilyStatement.lean}{GlobalTheory.T.nonempty_tensor_of_onto_full_family}
\leanref{Development/AbelianVarietyTheory/Proofs/CubicalLineBundles/NonemptyTensorOfOntoFullFamily.lean}{GlobalTheory.CM.nonempty_isTen_of_isFam_top}
\begin{proof}
Step 1: a tensor product. By Lemma~\ref{cv-cusp-cm-point-eigen-degree-tensor-recognition-tensor}, there is a
principally polarized abelian scheme $A_0$ over $R$ with an action of $\mathcal O_E$ for which the Rosati involution
induces complex conjugation, that is a Serre tensor product of $C$ and $L$ with form $h$; let $t^0_a$ be its
homomorphisms.

Step 2: the comparison. Lemma~\ref{cv-cusp-cm-point-eigen-degree-tensor-recognition-lift}, for $L_1=L$, $A_1=A_0$,
$M=L$, $f$ the identity and $c=1$, gives a homomorphism $\Phi\colon A_0\to A$ that commutes with the actions of
$\mathcal O_E$, with $\Phi\circ t^0_a=t_a$ for every $a\in L$ and
$\Lambda_{\mathcal L_A}(\Phi p_1,\Phi p_2)\cong\Lambda_{\mathcal L_{A_0}}(p_1,p_2)$ on $A_0\times_RA_0$, for the line
bundles $\mathcal L_A$ and $\mathcal L_{A_0}$ of the polarizations.

Step 3: $\Phi$ is surjective. Every point of $A$ over $\mathbf C$ through $\sigma$ is a sum $\sum_mt_{a_m}(x_m)$, as the
family is onto at $\sigma$, which is $\Phi\big(\sum_mt^0_{a_m}(x_m)\big)$, as $\Phi$ is a homomorphism. As $\sigma$ is
injective, Lemma~\ref{cv-cusp-cm-point-eigen-degree-tensor-recognition-onto} shows that $\Phi$ is surjective.

Step 4: $\Phi$ is an isomorphism. The base $\operatorname{Spec}R$ is reduced, the polarization of $A_0$ is principal, so
the kernel of $\mathcal L_{A_0}$ is trivial, and by Step~2 that of $A$ pulls back along $\Phi$ to it. With Step~3,
Lemma~\ref{cv-cusp-cm-point-eigen-degree-tensor-recognition-iso} shows that $\Phi$ is an isomorphism.
% Lean: Step 1 is T.exists_cm_tensor_lattice_curve_principal_triple (= CM.exists_isTen, which uses the positivity for
% the polarization); Step 2 is T.exists_tensorProduct_hom_of_cmCurveHomFamily (= CM.exists_famLift) with f the
% identity of L (Submodule.topEquiv) and c = 1, with T_1.t.act.one; Step 3 is
% T.abelianScheme_surjective_of_onto_complex_points with OntoAt.exists_eq_comp; Step 4 is
% T.isIso_of_surjective_principalPairing (= CM.PairingIso.isIso_of_lamEq), with T_1.ktriv for the principal
% polarization. What remains open is in Lemma cv-cusp-cm-point-eigen-degree-tensor-recognition-tensor, about 4,300 of
% the 5,000 lines counted for the four inputs.

Step 5: the tensor structure. Composing the natural isomorphisms $L\otimes_{\mathcal O}C(T)\to A_0(T)$ with $\Phi$ gives
natural isomorphisms $L\otimes_{\mathcal O}C(T)\to A(T)$ of $\mathcal O$-modules, whose homomorphisms
$x\mapsto a\otimes x$ are $\Phi\circ t^0_a=t_a$. Conditions (1) and (2) of
Definition~\ref{cv-cusp-cm-point-eigen-degree-serre} for them are (2) and (3) of
Definition~\ref{cv-cusp-cm-point-eigen-degree-tensor-family}.
\end{proof}
% Lean: nonempty_isTen_of_isFam_top: Step 5 is Serre.IsTensor.ofIso along asIso Phi, with key for t_a and the fields
% hfam.2.2.1 and hfam.2.2.2 of the family for lin and lam.

\begin{lemma}\label{cv-cusp-cm-point-eigen-degree-tensor}
Assume the conclusion of Theorem~\ref{cv-cusp-cm-point-eigen-degree-rank-one-norm-riemann-theta}, the laws of abelian
schemes (Definition~\ref{cv-cusp-cm-point-eigen-degree-tensor-laws}) and the conclusion of
Theorem~\ref{cv-model-toroidal-extension}.
In Situation~\ref{cv-cusp-setting}, let $\sqrt{-d}\in E$ be the square root of $-d$ of the datum, and put
$w_0=(1+\sqrt{-d})/2$ and $\mathcal O=\mathbf Z[w_0]$. Let $k\subset\mathbf C$ be a number field, let
$\theta\in\mathbf C$ with $\theta^2=-d$, and let $C$ be a CM curve with root $\theta$ over $k$
(Definition~\ref{cv-cm-point-curve}). Let $r\ge0$, and let $A$ be a principally polarized abelian scheme over
$\mathcal O_k$ with an action of $\mathcal O_E$ for which the Rosati involution of the polarization induces complex
conjugation on $\mathcal O_E$. Suppose that for every algebraically closed field $\Omega$ and every embedding
$\iota\colon k\to\Omega$, the eigenspace of an embedding $\sigma\colon E\to\Omega$ in the Lie algebra of the base
change of $A$ to $\Omega$ along $\iota$ has dimension $r$ if $\sigma(\sqrt{-d})=\iota(\theta)$ and dimension $0$
otherwise. Then
there are a number field $k_1$ with $k\subset k_1\subset\mathbf C$, the lattice $L$ of a PEL datum over $E$ of
signature $(p,q)$ with $p+q=r$, with its perfect alternating form $\langle\ ,\ \rangle$
(Definition~\ref{cv-cm-point-periods}), and a sign $\varepsilon\in\{\pm1\}$ with $\varepsilon\langle w_0x,x\rangle>0$
for every nonzero $x\in L$, such that the base change of $A$ to $\mathcal O_{k_1}$ is a Serre tensor product
(Definition~\ref{cv-cusp-cm-point-eigen-degree-serre}) of the base change of $C$ to $\mathcal O_{k_1}$ and $L$ with
the form
\[
 h(x,y)=\varepsilon\big(\langle w_0x,y\rangle-\langle x,y\rangle\big)+\varepsilon\langle x,y\rangle w_0.
\]
\end{lemma}
\leanref{Development/AbelianVarietyTheory/Proofs/CubicalLineBundles/ExistsTensorDecompositionPureTypeCmTripleStatement.lean}{GlobalTheory.T.exists_tensor_decomposition_pure_type_cm_triple}
\leanref{Development/AbelianVarietyTheory/Proofs/CubicalLineBundles/ExistsTensorDecompositionPureTypeCmTriple.lean}{GlobalTheory.CM.Recog.exists_isTen_of_hasType}
\begin{proof}
For a ring homomorphism from $\mathcal O_k$ to a ring $R'$, write $C_{R'}$ and $A_{R'}$ for the base changes of $C$
and $A$ to $R'$, with the base changes of the actions, of the polarizations and of the basis of the Lie algebra of
$C$; base changes along composites are the composites of the base changes.

Step 1: the complex fibres. Base change along $\mathcal O_k\subset k\subset\mathbf C$. Then $C_{\mathbf C}$ is a
principally polarized elliptic curve over $\mathbf C$ with an action of $\mathcal O$ whose Rosati involution is complex
conjugation, $\sqrt{-d}$ acts on the basis of its Lie algebra by $\theta$, and the lattice of its uniformizations is
$c(\mathbf Z+\mathbf Z\frac{1+i\sqrt d}2)$ by Definition~\ref{cv-cm-point-curve}(3). The Rosati involution of $A_{\mathbf
C}$ induces complex conjugation on $\mathcal O_E$. For an algebraically closed field $\Omega$ and an embedding
$\iota\colon\mathbf C\to\Omega$, the base change of $A_{\mathbf C}$ along $\iota$ is that of $A$ along the restriction
of $\iota$ to $k$, so $A_{\mathbf C}$ satisfies the hypothesis of
Lemma~\ref{cv-cusp-cm-point-eigen-degree-tensor-complex} on eigenspaces.

Step 2: the complex family. Lemma~\ref{cv-cusp-cm-point-eigen-degree-tensor-complex}, for $\theta$, $C_{\mathbf C}$, $r$
and $A_{\mathbf C}$, gives the lattice $L$ of a PEL datum of signature $(p,q)$ with $p+q=r$, the sign $\varepsilon$
with the positivity of the statement, and a family of homomorphisms $(t_a)_{a\in L}$ from $C_{\mathbf C}$ to $A_{\mathbf
C}$ for $h$ on $L$ that is onto at the identity of $\mathbf C$.

Step 3: descent. Lemma~\ref{cv-cusp-cm-point-eigen-degree-tensor-descent}, for $k$, $C_k$, $A_k$ and the family of
Step~2, gives a number field $k_1$ with $k\subset k_1\subset\mathbf C$; applied with $k_2=k_1$, it gives a family of
homomorphisms $(t'_a)_{a\in L}$ from $C_{k_1}$ to $A_{k_1}$ for $h$ on $L$ whose base change to $\mathbf C$ is
$(t_a)_{a\in L}$.

Step 4: extension. Lemma~\ref{cv-cusp-cm-point-eigen-degree-tensor-extend}, for the Dedekind domain $\mathcal O_{k_1}$
with fraction field $k_1$, $C_{\mathcal O_{k_1}}$, $A_{\mathcal O_{k_1}}$ and the family of Step~3, gives a family of
homomorphisms $(t''_a)_{a\in L}$ from $C_{\mathcal O_{k_1}}$ to $A_{\mathcal O_{k_1}}$ for $h$ on $L$ whose restriction
to $k_1$ is $(t'_a)_{a\in L}$.

Step 5: onto. Let $\sigma\colon\mathcal O_{k_1}\to\mathbf C$ be the inclusion. The points of $A_{\mathcal O_{k_1}}$
and of $C_{\mathcal O_{k_1}}$ over $\mathbf C$ through $\sigma$ are those of $A_{\mathbf C}$ and of $C_{\mathbf C}$, and
on them $t''_a$ is $t_a$, by Steps 3 and 4. As $(t_a)_{a\in L}$ is onto at the identity of $\mathbf C$, $(t''_a)_{a\in
L}$ is onto at $\sigma$.

Step 6: recognition. The ring $\mathcal O_{k_1}$ is a noetherian domain, $\sigma$ is injective, and the Rosati
involution of $C_{\mathcal O_{k_1}}$ is complex conjugation. Lemma~\ref{cv-cusp-cm-point-eigen-degree-tensor-recognition}
shows that $A_{\mathcal O_{k_1}}$ is a Serre tensor product of $C_{\mathcal O_{k_1}}$ and $L$ with form $h$.
\end{proof}
\leanref{Development/AbelianVarietyTheory/Proofs/CubicalLineBundles/ExistsTensorDecompositionPureTypeCmTriple.lean}{GlobalTheory.CM.Recog.ontoAt_of_isGrpPullback}
\leanref{Development/AbelianVarietyTheory/Proofs/CubicalLineBundles/HasTypeBaseChange.lean}{GlobalTheory.CM.hasType_baseChange}
% Lean: Recog.exists_isTen_of_hasType. The base changes of C and A to k and to C are exists_ellMap and
% exists_pat_pullback_g_eq; ellMap_comp carries IsPrincipalAt (the lattice condition for every uniformization, which
% the text reads into Definition cv-cm-point-curve(3)) to C_C, and hasType_baseChange gives the type of A_C (Step 1).
% Step 2 is T.exists_homFamily_of_pureType_complex_triple. Step 3 is T.cmFamily_descends_to_numberField with k_2 = k_1,
% for the base changes C_2 and T_2 of C_1 = C_{O_{k_1}} and T_1 to k_1; Factor.ellMap_cancel,
% Factor.ellMap_cancel_and_sabMap_cancel and ObjMaps.exists_sab_comp identify them with the base changes of C_k and
% A_k, and C_C and A_C with their base changes to C. Step 4 is T.cmFamily_extends_over_dedekindDomain. Step 5 is
% ontoAt_of_isGrpPullback. Step 6 is T.nonempty_tensor_of_onto_full_family. Size: about 9,600 lines beyond the laws of
% the moduli problem and the uniformization of abelian varieties, as counted before these lemmas were stated. The
% descent and the identities are written to the floor. What remains is in the open steps of Lemmas
% cv-cusp-cm-point-eigen-degree-tensor-complex-homs-hom (about 400 lines), -complex-homs-bundle (about 250),
% -complex-pairing-lambda (about 900) and -recognition-tensor (about 4,300).

\subsection{Degrees of eigenvectors and existence of CM points}\label{subsection-cv-eigenvector-degrees}

\begin{lemma}\label{cv-cusp-cm-point-eigen-degree-ideal}
In Situation~\ref{cv-cusp-setting}, let $\sqrt{-d}\in E$ be the square root of $-d$ of the datum, put
$w_0=(1+\sqrt{-d})/2$, and write $b\mapsto\bar b$ for complex conjugation on $E$. Let $L$ be the lattice of a PEL
datum over $E$ of signature $(p,q)$ with $p+q=1$, with its perfect alternating form $\langle\ ,\ \rangle$
(Definition~\ref{cv-cm-point-periods}), let $\varepsilon\in\{\pm1\}$ with $\varepsilon\langle w_0x,x\rangle>0$ for
every nonzero $x\in L$, and put $h(x,y)=\varepsilon(\langle w_0x,y\rangle-\langle x,y\rangle)+\varepsilon\langle
x,y\rangle w_0$. Then there are a nonzero fractional ideal $\mathfrak a$ of $E$, an $\mathcal O_E$-linear bijection
$e\colon L\to\mathfrak a$ and $\beta\in E$ with $\bar\beta=\beta$ and $s(\beta)>0$ for every embedding
$s\colon E\to\mathbf C$, such that:
\begin{enumerate}
\item for every field $\Omega$ of characteristic $0$ with $[E:\mathbf Q]$ embeddings $E\to\Omega$, every embedding
  $\tau_0\colon E\to\Omega$ and all $x,y\in L$,
  \[
   \tau_0(h(x,y))=\sum_\sigma\sigma\big(\beta\,e(x)\,\overline{e(y)}\big),
  \]
  the sum over the embeddings $\sigma\colon E\to\Omega$ with $\sigma(\sqrt{-d})=\tau_0(\sqrt{-d})$;
\item $2\log\mathrm N(\mathfrak a)+\log\mathrm N_{E/\mathbf Q}(\beta)=[E:\mathbf Q]\big(\frac12\log
  d-\log\operatorname{rd}_E\big)$, with $\operatorname{rd}_E=|d_E|^{1/[E:\mathbf Q]}$ the root discriminant.
\end{enumerate}
\end{lemma}
\leanref{Development/AbelianVarietyTheory/Proofs/CubicalLineBundles/ExistsFractionalIdealTraceFormRankOneStatement.lean}{GlobalTheory.T.exists_fractionalIdeal_traceForm_rank_one}
\leanref{Development/AbelianVarietyTheory/Proofs/CubicalLineBundles/ExistsFractionalIdealTraceFormRankOne.lean}{GlobalTheory.CM.exists_ideal_traceForm_of_rank_one}
\begin{proof}
Put $\xi=\varepsilon/\sqrt{-d}$. As $\overline{\sqrt{-d}}=-\sqrt{-d}$ (proof of Lemma~\ref{cv-volume-key-xi}(1)),
$\bar\xi=-\xi\ne0$. Lemma~\ref{cv-volume-key-xi}(2), for $L$ of rank $1$ and $\xi$, gives $\beta\in E^\times$ with
$\bar\beta=\beta$ and an injective $\mathcal O_E$-linear map $e\colon L\to E$ whose image spans $E$, such that
$\langle x,y\rangle=\operatorname{Tr}_{E/\mathbf Q}\big(\xi\beta\,e(x)\overline{e(y)}\big)$ for all $x,y$. The image
$\mathfrak a=e(L)$ is a nonzero finitely generated $\mathcal O_E$-submodule of $E$, a nonzero fractional ideal. As
$\varepsilon^2=1$,
\[
 \varepsilon\langle x,y\rangle=\operatorname{Tr}_{E/\mathbf Q}\Big(\frac{\beta\,e(x)\overline{e(y)}}{\sqrt{-d}}\Big),\qquad
 \varepsilon\langle w_0x,y\rangle=\operatorname{Tr}_{E/\mathbf Q}\Big(\frac{w_0\beta\,e(x)\overline{e(y)}}{\sqrt{-d}}\Big).
 \tag{$*$}
\]

Positivity. Let $v\in E^\times$, and fix a nonzero $x_0\in L$. A nonzero integer multiple $nv/e(x_0)$ of
$v/e(x_0)$ lies in $\mathcal O_E$; then $x=(nv/e(x_0))x_0$ is nonzero and $e(x)=nv$. Put
$z=\beta\,e(x)\overline{e(x)}=n^2\beta v\bar v$. As $w_0=\frac12+\frac12\sqrt{-d}$, $(*)$ gives
$\varepsilon\langle w_0x,x\rangle=\frac12\varepsilon\langle x,x\rangle+\frac12\operatorname{Tr}_{E/\mathbf Q}(z)
=\frac12\operatorname{Tr}_{E/\mathbf Q}(z)$, since the form is alternating. So
$\operatorname{Tr}_{E/\mathbf Q}(\beta v\bar v)>0$ for every $v\ne0$. Every $s(\beta)$ is real, as
$s(\bar b)=\overline{s(b)}$ for the CM field $E$. Suppose that $s(\beta)<0$ for some $s$. By the approximation
theorem for the inequivalent absolute values of $E$ at its infinite places there is $a\in E$ with $|s(a)|>1$ and
$|s'(a)|<1$ for every embedding $s'\ne s,\bar s$. Then
$\operatorname{Tr}_{E/\mathbf Q}(\beta a^n\bar a^n)=\sum_{s'}s'(\beta)|s'(a)|^{2n}$, in which the terms of $s$ and
$\bar s$ are $s(\beta)|s(a)|^{2n}\le s(\beta)<0$ and the others tend to $0$ as $n\to\infty$. So this trace is
negative for large $n$, a contradiction. Hence $s(\beta)>0$ for every $s$.

(1) Let $z=\beta\,e(x)\overline{e(y)}$. For an embedding $\sigma\colon E\to\Omega$,
$\sigma(w_0)=(1+\sigma(\sqrt{-d}))/2$, so
\[
 \sigma\Big(\frac{w_0z}{\sqrt{-d}}\Big)-\sigma\Big(\frac z{\sqrt{-d}}\Big)+\sigma\Big(\frac z{\sqrt{-d}}\Big)\tau_0(w_0)
 =\sigma(z)\,\frac{\sigma(\sqrt{-d})+\tau_0(\sqrt{-d})}{2\sigma(\sqrt{-d})},
\]
which is $\sigma(z)$ if $\sigma(\sqrt{-d})=\tau_0(\sqrt{-d})$ and $0$ if $\sigma(\sqrt{-d})=-\tau_0(\sqrt{-d})$; as
$\sigma(\sqrt{-d})^2=-d$, these are the two cases. The trace is the sum over the $[E:\mathbf Q]$ embeddings into
$\Omega$, so by $(*)$
\[
 \tau_0(h(x,y))=\varepsilon\langle w_0x,y\rangle-\varepsilon\langle x,y\rangle+\varepsilon\langle x,y\rangle\tau_0(w_0)
 =\sum_{\sigma(\sqrt{-d})=\tau_0(\sqrt{-d})}\sigma(z).
\]

(2) With $G=(\beta)$, the triple $(G,\xi,e)$ is a Hermitian model of $L$ (Definition~\ref{cv-uniformization-model}),
so by Lemma~\ref{cv-volume-key-models-dual} an element $v\in E$ lies in $\mathfrak a$ if and only if
$\operatorname{Tr}_{E/\mathbf Q}(\xi\beta v\bar y)\in\mathbf Z$ for every $y\in\mathfrak a$. So $\xi\beta\mathfrak a$
is the dual of $\bar{\mathfrak a}$ for the trace form, $\xi\beta\mathfrak a=\bar{\mathfrak a}^{-1}\mathfrak d_E^{-1}$
with $\mathfrak d_E$ the different of $E$. Taking norms, with $\mathrm N(\bar{\mathfrak a})=\mathrm N(\mathfrak a)$
and $\mathrm N(\mathfrak d_E)=|d_E|$,
$2\log\mathrm N(\mathfrak a)+\log|\mathrm N_{E/\mathbf Q}(\xi\beta)|=-\log|d_E|=-[E:\mathbf Q]\log\operatorname{rd}_E$.
Finally $\mathrm N_{E/\mathbf Q}(\sqrt{-d})^2=\mathrm N_{E/\mathbf Q}(-d)=d^{[E:\mathbf Q]}$, so
$\log|\mathrm N_{E/\mathbf Q}(\xi)|=-\frac12[E:\mathbf Q]\log d$, and $\mathrm N_{E/\mathbf Q}(\beta)>0$ as $\beta$ is
totally positive.
\end{proof}
\leanref{Development/AbelianVarietyTheory/Proofs/CubicalLineBundles/ExistsFractionalIdealTraceFormRankOne.lean}{GlobalTheory.CM.RankOne.re_pos_of_trace_pos}
\leanref{Development/AbelianVarietyTheory/Proofs/CubicalLineBundles/ExistsFractionalIdealTraceFormRankOne.lean}{GlobalTheory.CM.RankOne.relTrace_eq}
\leanref{Development/AbelianVarietyTheory/Proofs/CubicalLineBundles/ExistsFractionalIdealTraceFormRankOne.lean}{GlobalTheory.CM.RankOne.log_absNorm_of_selfDual}
% Lean: exists_ideal_traceForm_of_rank_one: exists_herm_coord_of_datum with xi = eps/rt and r = 1 (Lemma
% cv-volume-key-xi(2)), HermModel.rank_integral_selfDual for the self-duality (Lemma cv-volume-key-models-dual); hE and
% hE' are (*); the positivity is re_pos_of_trace_pos (AbsoluteValue.exists_one_lt_lt_one_pi_of_not_isEquiv for a, and
% the limit of the sums g); (1) is relTrace_eq; (2) is log_absNorm_of_selfDual (the dual of conjFrac a is beta xi a,
% dual_eq_mul_inv, absNorm_differentIdeal) with log_norm_rt.

\begin{lemma}\label{cv-cusp-cm-point-eigen-degree-dual}
In Situation~\ref{cv-cusp-setting}, let $\sqrt{-d}\in E$ be the square root of $-d$ of the datum, and put
$w_0=(1+\sqrt{-d})/2$ and $\mathcal O=\mathbf Z[w_0]$. Let $L$ be the lattice of a PEL datum over $E$, with its
perfect alternating form $\langle\ ,\ \rangle$ (Definition~\ref{cv-cm-point-periods}), let $\varepsilon\in\{\pm1\}$,
and put $h(x,y)=\varepsilon(\langle w_0x,y\rangle-\langle x,y\rangle)+\varepsilon\langle x,y\rangle w_0$ for
$x,y\in L$. Then $h$ is $\mathcal O$-linear in its first variable, every $\mathcal O$-linear map
$f\colon L\to\mathcal O$ is $h(\cdot,y)$ for a unique $y\in L$, and there are $x_1,\dots,x_N$ and $y_1,\dots,y_N$ in
$L$ with $a=\sum_mh(a,y_m)x_m$ for every $a\in L$.
\end{lemma}
\leanref{Development/AbelianVarietyTheory/Proofs/CubicalLineBundles/ExistsUniqueCmHermitianLatticeFormStatement.lean}{GlobalTheory.T.exists_unique_cmHermitianLatticeForm}
\leanref{Development/AbelianVarietyTheory/Proofs/CubicalLineBundles/ExistsUniqueCmHermitianLatticeForm.lean}{GlobalTheory.CM.exists_kForm}
\begin{proof}
Linearity. The form $h$ is additive in each variable. Put $u=\langle w_0x,y\rangle-\langle x,y\rangle$ and
$t=\langle x,y\rangle$. As $w_0^2=w_0-\frac{1+d}4$, $\langle w_0^2x,y\rangle=\langle w_0x,y\rangle-\frac{1+d}4t$, so
$h(w_0x,y)=\varepsilon\big(-\frac{1+d}4t+(u+t)w_0\big)=\varepsilon(uw_0+tw_0^2)=w_0h(x,y)$.

Perfectness. Every element of $\mathcal O$ is $m+nw_0$ for unique integers $m$ and $n$. Let $f\colon L\to\mathcal O$
be $\mathcal O$-linear, and write $f(x)=f_1(x)+f_2(x)w_0$. As the form is perfect and $f_2$ is additive, there is
$y\in L$ with $\varepsilon\langle x,y\rangle=f_2(x)$ for every $x$. Then $f(x)$ and $h(x,y)$ have the same coefficient
at $w_0$, so $f(x)-h(x,y)$ is an integer $m(x)$. Both $f$ and $h(\cdot,y)$ are $\mathcal O$-linear, so
$m(w_0x)=w_0m(x)$, and an integer multiple of $w_0$ is an integer only if it is $0$; hence $f=h(\cdot,y)$. If also
$f=h(\cdot,y')$, the coefficients at $w_0$ give $\langle x,y-y'\rangle=0$ for every $x$, and $y=y'$ as the form is
perfect.

Dual families. The module $L$ is finitely generated and projective over $\mathcal O_E$, hence over $\mathcal O$.
Choose an $\mathcal O$-linear surjection $\pi\colon\mathcal O^N\to L$ and an $\mathcal O$-linear section $s$ of it,
with coordinates $s_1,\dots,s_N$. Put $x_m=\pi(e_m)$, and let $y_m$ be the element with $h(\cdot,y_m)=s_m$. Then
$a=\pi(s(a))=\sum_ms_m(a)x_m=\sum_mh(a,y_m)x_m$.
\end{proof}
\leanref{Development/AbelianVarietyTheory/Proofs/CubicalLineBundles/ExistsUniqueCmHermitianLatticeForm.lean}{GlobalTheory.CM.kLin_perfect}
% Lean: exists_kForm: kLin (kFun_smul for the linearity), kLin_perfect (the coordinates c of
% integers_of_cm_field_conj_structure, D.unimod.2 for y, kform_coeff_unique for the integer m), and for the dual
% families integers_of_cm_field_conj_structure (L projective over OK H), Module.Finite.exists_fin',
% Module.projective_lifting_property and kLin_perfect for the y_i.

\begin{definition}\label{cv-cusp-cm-point-eigen-degree-rank-one-norm-weil-translate-group}
Let $S$ be a scheme, $A$ a commutative group scheme over $S$, $\mathcal L$ a line bundle on $A$ and $k$ an integer.
Let $T$ be a scheme over $S$, let $A_T=T\times_SA$ with projections $p_T$ and $p_A$, and for $x\in A(T)$ let
$t_x\colon A_T\to A_T$ be the translation by $x$, the morphism over $T$ with $p_A\circ t_x=x\,p_T+p_A$.
\begin{enumerate}
\item An \emph{element of the theta group of $\mathcal L^k$ over $T$} is a pair $(x,\varphi)$ of a point
  $x\in A(T)$ and a generator $\varphi$ of the line bundle $p_A^*\mathcal L^k\otimes(x\,p_T+p_A)^*\mathcal L^{-k}$ on
  $A_T$; it \emph{lies over} $x$.
\item Two elements $(x,\varphi)$ and $(y,\psi)$ of the theta group of $\mathcal L^k$ over $T$ have \emph{commutator}
  $c$, a unit of $\Gamma(T,\mathcal O_T)$, if $\varphi\otimes t_x^*\psi=c\cdot\psi\otimes t_y^*\varphi$; both sides are
  generators of $p_A^*\mathcal L^k\otimes(x\,p_T+y\,p_T+p_A)^*\mathcal L^{-k}$.
\item For $x,y\in A(T)$ and a unit $c$ of $\Gamma(T,\mathcal O_T)$, the \emph{commutator pairing} of the theta group
  of $\mathcal L^k$ has the value $e_k(x,y)=c$ if every point of $T$ has an open neighbourhood $V$ such that the
  restrictions of $x$ and $y$ to $V$ lie under elements of the theta group of $\mathcal L^k$ over $V$, and any two such
  elements, over $x|_V$ and over $y|_V$, have commutator $c|_V$.
\end{enumerate}
\end{definition}
\leanref{ReciprocityTheory/ReductiveGroups/ArithmeticGeometry/ExtensionDescent/ConstGen.lean}{ReciprocityTheory.ChartRestriction.G.Out.ThetaGrp}
\leanref{ReciprocityTheory/ReductiveGroups/ArithmeticGeometry/ExtensionDescent/ConstGen.lean}{ReciprocityTheory.ChartRestriction.G.Out.CommIsOf}
\leanref{ReciprocityTheory/ReductiveGroups/ArithmeticGeometry/ExtensionDescent/ConstGen.lean}{ReciprocityTheory.ChartRestriction.G.Out.WeilIs}

\begin{lemma}\label{cv-cusp-cm-point-eigen-degree-rank-one-norm-weil-translate-pairing}
Let $S$ be a scheme, let $A$ be a commutative group scheme over $S$ that is an abelian scheme, and let $\mathcal L$ be
a line bundle on $A$, trivialized along the zero section, with a cubical structure. Let $T$ be a scheme over $S$ with
exactly one point, let $A_T=T\times_SA$ with projections $p_T$ and $p_A$, let $n\ge0$ be an integer and let $x\in
A(T)$ with $nx=0$. Then the line bundle $(x\,p_T+n\,p_A)^*\mathcal L\otimes(n\,p_A)^*\mathcal L^{-1}$ on $A_T$ has a
generator $u$ with the following property. Let $y\in A(T)$ with $ny=0$ and let $c$ be a unit of $\Gamma(T,\mathcal O_T)$
with $e_n(x,y)=c$ for the commutator pairing of the theta group of $\mathcal L^n$
(Definition~\ref{cv-cusp-cm-point-eigen-degree-rank-one-norm-weil-translate-group}). Let $Z$ be a scheme over $T$, with
structure morphism $t$, and $w\in A(Z)$. Then the pullbacks $u(t,w)$ and $u(t,yt+w)$ of $u$ along the points $(t,w)$
and $(t,yt+w)$ of $A_T$ are generators of the same line bundle $(xt+nw)^*\mathcal L\otimes(nw)^*\mathcal L^{-1}$ on $Z$,
and $u(t,yt+w)=t^*c\cdot u(t,w)$.
\end{lemma}
\leanref{Development/AbelianVarietyTheory/Proofs/ModuliProofs/LbIsoCubeLBOfWeilIsAdjoint.lean}{GlobalTheory.AdjointOfPairing.exists_gen_pairing}
\begin{proof}
Let $\theta$ be the cubical structure, and for points $a$, $m$, $x'$ of $A$ with values in a scheme let
$c_k(a;m,x')=\prod_{i=0}^{k-1}\theta(a+im,m,x')$ be the generators of
Definition~\ref{cv-toroidal-abelian-dvr-ample-torsion-weil}. By the computation there,
$c_n(0;m,x')$ is a generator of the tensor product of $\Lambda_{\mathcal L}(m,x')^{-n}$ and
$(nm+x')^*\mathcal L\otimes(nm)^*\mathcal L^{-1}\otimes x'^*\mathcal L^{-1}\otimes0^*\mathcal L$. As $T$ has one point,
every line bundle on $T$ has a generator: it has one near the point of $T$, and the only open neighbourhood of that point
is $T$.

Step 1: lifts to the theta group. Let $\sigma_x$ be a generator of $0^*\mathcal L^n\otimes x^*\mathcal L^{-n}$ on $T$.
For $m=x\,p_T$ and $x'=p_A$, as $nx=0$, the second factor above is trivial, so $c_n(0;x\,p_T,p_A)$ is a generator of
$\Lambda_{\mathcal L}(x\,p_T,p_A)^{-n}$, which is $p_A^*\mathcal L^n\otimes(x\,p_T+p_A)^*\mathcal L^{-n}$ tensored with the
pullback of $x^*\mathcal L^n\otimes0^*\mathcal L^{-n}$. Hence $\varphi_x=c_n(0;x\,p_T,p_A)\otimes p_T^*\sigma_x$ is a
generator of $p_A^*\mathcal L^n\otimes(x\,p_T+p_A)^*\mathcal L^{-n}$, and $(x,\varphi_x)$ is an element of the theta
group of $\mathcal L^n$ over $T$. In the same way every $y\in A(T)$ with $ny=0$ lies under the element $(y,\varphi_y)$,
with $\varphi_y=c_n(0;y\,p_T,p_A)\otimes p_T^*\sigma_y$.

Step 2: the generator. Let $\sigma'$ be a generator of $0^*\mathcal L^{n-1}\otimes x^*\mathcal L^{1-n}$ on $T$, and put
\[
 u=c_n(0;p_A,x\,p_T)\otimes\varphi_x^{-1}\otimes p_T^*\sigma'.
\]
Counting the exponents of the pullbacks of $\mathcal L$ along $p_A$, $x\,p_T$, $x\,p_T+p_A$, $0$, $x\,p_T+n\,p_A$ and
$n\,p_A$ in the three factors, $u$ is a generator of $(x\,p_T+n\,p_A)^*\mathcal L\otimes(n\,p_A)^*\mathcal L^{-1}$. As
$n(yt)=0$, the pullbacks of this line bundle along $(t,w)$ and $(t,yt+w)$ are both
$(xt+nw)^*\mathcal L\otimes(nw)^*\mathcal L^{-1}$.

Step 3: the commutator. The generators $\varphi_x\otimes t_x^*\varphi_y$ and $\varphi_y\otimes t_y^*\varphi_x$ of one line
bundle on $A_T$ differ by a unit of $\Gamma(A_T,\mathcal O)$, which is the pullback of a unit $c_0$ of
$\Gamma(T,\mathcal O_T)$, as $A_T$ is an abelian scheme over $T$ \cite[Tag 0E0L]{Sta}; so $(x,\varphi_x)$ and
$(y,\varphi_y)$ have commutator $c_0$. Elements of the theta group over an open set $V\subset T$ that lie over $x|_V$ and
$y|_V$ differ from the restrictions of $(x,\varphi_x)$ and $(y,\varphi_y)$ by units of $\Gamma(V\times_SA,\mathcal O)$,
which again come from $V$ and are invariant under translations; so they have the commutator $c_0|_V$, and
$e_n(x,y)=c_0$. Two values of $e_n(x,y)$ agree near each point of $T$, as both are commutators of the same elements over
the intersection of the neighbourhoods. Hence $c=c_0$:
\[
 \varphi_x\otimes t_x^*\varphi_y=c\cdot\varphi_y\otimes t_y^*\varphi_x.
\]

Step 4: the translation. Write $x_Z=xt$ and $y_Z=yt$, and $\varphi(t,w)$ for the pullback of a generator $\varphi$ on
$A_T$ along $(t,w)$; then $t_x\circ(t,w)=(t,x_Z+w)$. Identities (i) and (ii) of the proof of
Lemma~\ref{cv-toroidal-abelian-dvr-ample-torsion-basic}, with $ny_Z=0$ and the symmetry $\theta(y_Z,x_Z,w)=
\theta(y_Z,w,x_Z)$, give
\[
 c_n(0;y_Z+w,x_Z)\,c_n(0;y_Z,w)=c_n(0;w,x_Z)\,c_n(0;y_Z,x_Z+w).
\]
By Steps 1 and 2, $c_n(0;y_Z,w)=\varphi_y(t,w)\otimes t^*\sigma_y^{-1}$,
$c_n(0;y_Z,x_Z+w)=\varphi_y(t,x_Z+w)\otimes t^*\sigma_y^{-1}$, and
$c_n(0;w',x_Z)=u(t,w')\otimes\varphi_x(t,w')\otimes t^*\sigma'^{-1}$ for $w'=w$ and $w'=y_Z+w$. Substituting,
\[
 u(t,y_Z+w)\otimes\varphi_x(t,y_Z+w)\otimes\varphi_y(t,w)=u(t,w)\otimes\varphi_x(t,w)\otimes\varphi_y(t,x_Z+w).
\]
By Step 3, pulled back along $(t,w)$, $\varphi_x(t,w)\otimes\varphi_y(t,x_Z+w)=t^*c\cdot\varphi_y(t,w)\otimes
\varphi_x(t,y_Z+w)$. Cancelling this generator gives $u(t,y_Z+w)=t^*c\cdot u(t,w)$.
\end{proof}
\leanref{Development/AbelianVarietyTheory/Proofs/AbelianSchemeLawProofs/ExistsGenTranslEqPairing.lean}{GlobalTheory.DescentPairing.exists_theta}
\leanref{Development/AbelianVarietyTheory/Proofs/AbelianSchemeLawProofs/ExistsGenTranslEqPairing.lean}{GlobalTheory.DescentPairing.exists_triv}
\leanref{Development/AbelianVarietyTheory/Proofs/AbelianSchemeLawProofs/WeilIsIffCommIsOfAndLiftSubsingleton.lean}{GlobalTheory.GlobalLift.weilIs_iff_commIsOf}
\leanref{Development/AbelianVarietyTheory/Proofs/AbelianSchemeLawProofs/ExistsGenTranslEqPairing.lean}{GlobalTheory.DescentPairing.cubN_transl}
% Lean: exists_gen_pairing (with GlobalFnsSt). Step 1 is DescentPairing.exists_theta (sigma from exists_frame on the
% one-point T, the cube generator from exists_cubN, and exists_gen for the line bundle); Step 2 is exists_triv (the cube
% generator for (w, x), the inverse of theta_x and the frame sigma'). Step 3 is GlobalLift.weilIs_iff_commIsOf:
% CommutatorPairing.exists_commIsOf (exists_ratio with surj_fst, from GlobalFnsSt), weilIs_of_commIsOf
% (commIsOf_change) and weilIs_unique (commVal_indep); then commIsOf_iff for the identity in charts. Step 4 is
% cubN_transl (cubN_mul, which is (i), and IsCub.N_add, which is (ii), with the symmetry s12), with the values of
% theta_x, theta_y and u from exists_theta and exists_triv. The Lean compares values in charts of L at each point of Z;
% the text compares the generators.

\begin{lemma}\label{cv-cusp-cm-point-eigen-degree-rank-one-norm-weil-translate}
Let $S$ be a scheme, let $A$ be a commutative group scheme over $S$ that is an abelian scheme, and let $\mathcal L$ be
a line bundle on $A$, trivialized along the zero section, with a cubical structure. Let $T$ be a scheme over $S$ with
exactly one point, let $X$ be a scheme over $S$ and $Y$ a commutative group scheme over $S$, let $a\colon X\to A$ be a
morphism and $g\colon Y\to A$ a homomorphism over $S$, let $n\ge0$ be an integer and let $x\in X(T)$ with $n\,a(x)=0$.
Let $q_T$ and $q_Y$ be the projections of $T\times_SY$, and let $\mathcal M$ be the line bundle on $T\times_SY$ obtained
by pulling back $\Lambda_{\mathcal L}(a\,p_1,g\,p_2)$ (Definition~\ref{cv-toroidal-abelian-dvr-ample-torsion-weil}) on
$X\times_SY$, with $p_1$ and $p_2$ the projections, along $(x\,q_T,n\,q_Y)$. For $z\in Y(T)$ with $nz=0$, the pullbacks
of $\mathcal M$ along the points $(1,z)$ and $(1,0)$ of $T\times_SY$ are both the line bundle
$\Lambda_{\mathcal L}(a(x),0)$ on $T$. Then $\mathcal M$ has a generator $t$ such that $t(1,z)=e\,t(1,0)$ for every
$z\in Y(T)$ with $nz=0$ and every unit $e$ of $\Gamma(T,\mathcal O_T)$ with $e_n(a(x),g(z))=e$, for the commutator pairing
of the theta group of $\mathcal L^n$
(Definition~\ref{cv-cusp-cm-point-eigen-degree-rank-one-norm-weil-translate-group}).
\end{lemma}
\leanref{Development/AbelianVarietyTheory/Proofs/CubicalLineBundles/WeilPairingEqOfIsomorphicCubePairings.lean}{GlobalTheory.CM.exists_gen_lam_two}
\begin{proof}
Let $u$ be the generator of Lemma~\ref{cv-cusp-cm-point-eigen-degree-rank-one-norm-weil-translate-pairing} for the point
$a(x)$ of $A$, a generator of $(a(x)\,p_T+n\,p_A)^*\mathcal L\otimes(n\,p_A)^*\mathcal L^{-1}$ on $A_T$. As $T$ has one
point, the line bundle $a(x)^*\mathcal L^{-1}\otimes0^*\mathcal L$ on $T$ has a generator $\sigma$. Put
$t=(q_T,g\,q_Y)^*u\otimes q_T^*\sigma$. As $g$ is a homomorphism, $n\,g\,q_Y=g(n\,q_Y)$, so $t$ is a generator of
\[
 (a(x)\,q_T+g(n\,q_Y))^*\mathcal L\otimes(a(x)\,q_T)^*\mathcal L^{-1}\otimes g(n\,q_Y)^*\mathcal L^{-1}\otimes
 0^*\mathcal L=\mathcal M.
\]
Let $z\in Y(T)$ with $nz=0$ and $e$ with $e_n(a(x),g(z))=e$. Then $n\,g(z)=0$, and that lemma, for $y=g(z)$, for $Z=T$
with the identity as structure morphism, and for $w=0$, gives $u(1,g(z))=e\,u(1,0)$. Hence $t(1,z)=u(1,g(z))\otimes\sigma=e\,u(1,0)\otimes\sigma=
e\,t(1,0)$.
\end{proof}
% Lean: exists_gen_lam_two: exists_gen_pairing for x = y a (u_0), exists_frame for sigma on the one-point T, exists_gen
% for t = (pullback of u_0 along lift (fst T Y) (snd T Y g)) times (pullback of sigma along fst), with prod_mulT and
% NormRos.prod_square for the line bundles; then the values w (val_mul, val_pull) and w_0 (exists_gen_pairing at Z = T,
% t = 1 and w = 1).

\begin{lemma}\label{cv-cusp-cm-point-eigen-degree-rank-one-norm-weil}
Let $k$ be an algebraically closed field of characteristic $0$, let $X$, $Y$, $A$ and $B$ be principally polarized
abelian varieties over $k$, let $\mathcal L_A$ and $\mathcal L_B$ be the line bundles of the polarizations of $A$ and
$B$, and for $n\ge1$ let $e^A_n$ and $e^B_n$ be the commutator pairings of the theta groups of $\mathcal L_A^n$ and
$\mathcal L_B^n$. Let $f_1\colon X\to A$, $f_2\colon Y\to A$, $g_1\colon X\to B$ and $g_2\colon Y\to B$ be
homomorphisms such that, with $p_1$ and $p_2$ the projections of $X\times Y$,
\[
 \Lambda_{\mathcal L_A}(f_1p_1,f_2p_2)\cong\Lambda_{\mathcal L_B}(g_1p_1,g_2p_2)
\]
on $X\times Y$ (Definition~\ref{cv-toroidal-abelian-dvr-ample-torsion-weil}). Then $e^A_n(f_1x,f_2y)=e^B_n(g_1x,g_2y)$
for all $n\ge1$ and $x\in X(k)$, $y\in Y(k)$ with $nx=ny=0$.
\end{lemma}
\leanref{Development/AbelianVarietyTheory/Proofs/CubicalLineBundles/WeilPairingEqOfIsomorphicCubePairingsStatement.lean}{GlobalTheory.T.weil_pairing_eq_of_isomorphic_cube_pairings}
\leanref{Development/AbelianVarietyTheory/Proofs/CubicalLineBundles/WeilPairingEqOfIsomorphicCubePairings.lean}{GlobalTheory.CM.weilEq_of_lamEq}
\begin{proof}
The line bundles $\mathcal L_A$ and $\mathcal L_B$ are trivialized along the zero sections, and by the theorem of the
cube they have cubical structures. Fix $n$, $x$ and $y$, and let $j\colon Y\to X\times Y$ be the morphism
$w\mapsto(x,nw)$. As $nx=ny=0$, the points $f_1x$, $f_2y$, $g_1x$ and $g_2y$ are killed by $n$.

Lemma~\ref{cv-cusp-cm-point-eigen-degree-rank-one-norm-weil-translate}, for $S=T=\operatorname{Spec}k$, $A$,
$\mathcal L_A$, $a=f_1$ and $g=f_2$,
gives a generator $t_A$ of the pullback $\mathcal M_A$ of $\Lambda_{\mathcal L_A}(f_1p_1,f_2p_2)$ along $j$ with
$t_A(y)=e^A_n(f_1x,f_2y)\,t_A(0)$; for $B$, $\mathcal L_B$, $a=g_1$ and $g=g_2$, it gives a generator $t_B$ of the
pullback $\mathcal M_B$ of $\Lambda_{\mathcal L_B}(g_1p_1,g_2p_2)$ along $j$ with $t_B(y)=e^B_n(g_1x,g_2y)\,t_B(0)$.

Let $\varphi$ be the isomorphism of the hypothesis. Then $j^*\varphi$ carries $t_A$ to a generator of $\mathcal M_B$,
so $t_B=c\cdot j^*\varphi(t_A)$ for an invertible function $c$ on $Y$. The global functions on the abelian variety
$Y$ are the constants, so $c(y)=c(0)$. As $ny=0$, $j(y)=(x,0)=j(0)$, so at $y$ and at $0$ the map $j^*\varphi$ is the
map $\varphi_{(x,0)}$ of the fibres of the two line bundles at $(x,0)$. Hence
\[
 c(0)=c(y)=\frac{t_B(y)}{\varphi_{(x,0)}(t_A(y))}
 =\frac{e^B_n(g_1x,g_2y)}{e^A_n(f_1x,f_2y)}\cdot\frac{t_B(0)}{\varphi_{(x,0)}(t_A(0))}
 =\frac{e^B_n(g_1x,g_2y)}{e^A_n(f_1x,f_2y)}\,c(0),
\]
and $e^A_n(f_1x,f_2y)=e^B_n(g_1x,g_2y)$.
\end{proof}
\leanref{Development/AbelianVarietyTheory/Proofs/CubicalLineBundles/WeilPairingEqOfIsomorphicCubePairings.lean}{GlobalTheory.CM.weil_value_eq_of_lamEq}
\leanref{Development/AbelianVarietyTheory/Proofs/CubicalLineBundles/WeilIsOfLbIsoCubeLB.lean}{GlobalTheory.PairingOfLam.val_ratio}
% Lean: weilEq_of_lamEq: the rigidifications A.rigid and B.rigid, the cubical structures from the law Out.Laws.cube,
% and weil_value_eq_of_lamEq for T = Spec k with its single point; the uniqueness of the commutator pairing is the law
% Out.Laws.weil. weil_value_eq_of_lamEq: exists_gen_lam_two twice (Lemma
% cv-cusp-cm-point-eigen-degree-rank-one-norm-weil-translate) and PairingOfLam.val_ratio, which glues the ratio c of
% the two generators from the charts and uses WeilCube.surj_fst (the law GlobalFnsSt: the global functions on
% T x Y come from T) for c(y) = c(0), with e_1 the equality j(y) = j(0).

\begin{theorem}[Mumford \cite{Mum70}, Birkenhake--Lange \cite{BL04}]\label{cv-cusp-cm-point-eigen-degree-rank-one-norm-riemann-theta}
Let $A$ be a principally polarized abelian variety over $\mathbf C$ of dimension $g$, let $\mathcal L$ be the line
bundle of its polarization, trivialized on an open cover $(U_\alpha)$, and let $r$ be a generator of $0^*\mathcal L$,
with value $r_\alpha\in\mathbf C^\times$ in the trivialization over $U_\alpha$ when $0\in U_\alpha$. For $k\ge1$ let
$e_k$ be the commutator pairing of the theta group of $\mathcal L^k$
(Definition~\ref{cv-cusp-cm-point-eigen-degree-rank-one-norm-weil-translate-group}). Then $A$ has a uniformization
$u\colon\mathbf C^g\to A(\mathbf C)$ for $\mathcal L$ (Definition~\ref{cv-origin-hodge-norm-gram}(1), with
$T=\operatorname{Spec}\mathbf C$, $G=A$ and $\varphi$ the identity), with kernel $\Lambda$, trivializing functions
$\theta_\alpha$ and Appell--Humbert data $(H,\chi)$, such that, with $E(a,b)=\operatorname{Im}(a^{\mathsf T}H\bar b)$:
\begin{enumerate}
\item $\theta_\alpha(0)=r_\alpha$ for every $\alpha$ with $0\in U_\alpha$;
\item every tangent vector of $A$ at the origin has exactly one differential, and every $a\in\mathbf C^g$ is the
  differential of exactly one tangent vector of $A$ at the origin;
\item for every homomorphism $f\colon A\to A$ there is a $\mathbf C$-linear map $F\colon\mathbf C^g\to\mathbf C^g$ with
  $u\circ F=f\circ u$;
\item for every $n\ge0$, call a holomorphic function $\vartheta$ on $\mathbf C^g$ with
  $\vartheta(z+l)=\big(\chi(l)\exp(\pi z^{\mathsf T}H\bar l+\frac\pi2l^{\mathsf T}H\bar l)\big)^n\vartheta(z)$ for all
  $l\in\Lambda$ and $z$ a \emph{theta function of weight $n$}, and say that it \emph{represents} a section $s$ of
  $\mathcal L^n$ if $s_\alpha(u(z))=\vartheta(z)\,\theta_\alpha(z)^n$ for every $\alpha$ and $z\in u^{-1}(U_\alpha)$, where
  $s_\alpha$ is $s$ in the trivialization over $U_\alpha$; then every section of $\mathcal L^n$ is represented by
  exactly one theta function of weight $n$, and every theta function of weight $n$ represents exactly one section;
\item $E$ takes integer values on $\Lambda\times\Lambda$, and $a\mapsto E(a,\cdot)$ is a bijection of $\Lambda$ onto
  $\operatorname{Hom}(\Lambda,\mathbf Z)$;
\item $e_k\big(u(a/k),u(b/k)\big)=\exp(-2\pi iE(a,b)/k)$ for all $k\ge1$ and $a,b\in\Lambda$.
\end{enumerate}
\end{theorem}
\leanref{ReciprocityTheory/ReductiveGroups/ArithmeticGeometry/BirationalGeometry/ToroidalCompactification/ThetaUnifAV.lean}{ReciprocityTheory.ChartRestriction.G.Geo.Tor.ThetaUnifAV}
\leanref{Development/AbelianVarietyTheory/Proofs/ModuliProofs/ThetaUniformisationComplexAbelianVarietiesStatement.lean}{GlobalTheory.T.theta_uniformisation_complex_abelian_varieties}
\leanref{Development/AbelianVarietyTheory/Proofs/ToricGeometry/ThetaUnifAVOfLaws.lean}{GlobalTheory.thetaUnifAV_of_laws}

The uniformization of a complex abelian variety and the lifts of its homomorphisms are \cite[\S1]{Mum70} and
\cite[Chapter~1]{BL04}; the Appell--Humbert data and the theta functions representing sections are \cite[\S\S2--3]{Mum70}
and \cite[Chapters~2 and~3]{BL04}; for a principal polarization $E$ is unimodular on $\Lambda$
\cite[Chapters~3 and~4]{BL04}; and the commutator pairings of the theta groups are computed from $E$ in
\cite[Chapter~6]{BL04}.
% The sign in (6) is the Lean's: ThetaUnifAV has c = exp(-(2 pi i) Ef a b / k) with Ef a b = Im (hermForm Hm a b) and
% hermForm Hm z w = z^T Hm conj(w), and Out.WeilIs through CommIsOf, the text's commutator of
% Definition cv-cusp-cm-point-eigen-degree-rank-one-norm-weil-translate-group. Remark
% cv-cusp-cm-point-eigen-degree-rank-one-norm-riemann-theta-sign checks it from the conventions of the text. The
% locators in the cited books for the sign were not available to check (open1-cv-fidelity).

\begin{remark}\label{cv-cusp-cm-point-eigen-degree-rank-one-norm-riemann-theta-sign}
The sign in (6) agrees with the conventions of the theorem. By (4) for $n=1$, a section of $\mathcal L$ is represented
by a function $\vartheta$ with $\vartheta(z+l)=f_l(z)\vartheta(z)$, where
$f_l(z)=\chi(l)\exp(\pi z^{\mathsf T}H\bar l+\frac\pi2l^{\mathsf T}H\bar l)$; so sections of a line bundle built from
$\mathcal L$ by tensor powers and translations are represented by functions with the corresponding factors, and the
pullback along the translation by $u(v)$ of the section represented by $\vartheta$ is represented by
$z\mapsto\vartheta(z+v)$. Write $H(z,w)=z^{\mathsf T}H\bar w$. Let $v=a/k$ and $w=b/k$ with $a,b\in\Lambda$, and put
$x=u(v)$ and $y=u(w)$. As $f_l(z)/f_l(z+v)=\exp(-\pi H(v,l))$, a generator of
$\mathcal L^k\otimes t_x^*\mathcal L^{-k}$ is represented by a nowhere vanishing function $\Phi$ with
$\Phi(z+l)=\exp(-\pi kH(v,l))\Phi(z)$. The function $\Phi(z)=\exp(-\pi kH(z,v))$ is one: the quotient of the two factors
is $\exp(\pi k(H(v,l)-H(l,v)))=\exp(2\pi iE(a,l))=1$, by (5). In the same way $\Psi(z)=\exp(-\pi kH(z,w))$ represents a
generator for $y$. The equation $\varphi\otimes t_x^*\psi=c\cdot\psi\otimes t_y^*\varphi$ of
Definition~\ref{cv-cusp-cm-point-eigen-degree-rank-one-norm-weil-translate-group}(2) reads
$\Phi(z)\Psi(z+v)=c\,\Psi(z)\Phi(z+w)$, so
\[
 c=\exp\big(-\pi k(H(v,w)-H(w,v))\big)=\exp\big(-2\pi ik\operatorname{Im}H(v,w)\big)=\exp(-2\pi iE(a,b)/k).
\]
Any other generators are constant multiples of these, as the global functions on $A$ are the constants, and give the
same $c$.
\end{remark}

\begin{lemma}\label{cv-cusp-cm-point-eigen-degree-rank-one-norm-riemann}
Let $A$ be a principally polarized abelian variety over $\mathbf C$ of dimension $h$, with the line bundle $\mathcal L$
of its polarization trivialized on an open cover, and for $n\ge1$ let $e_n$ be the commutator pairing of the theta
group of $\mathcal L^n$ (Definition~\ref{cv-cusp-cm-point-eigen-degree-rank-one-norm-weil-translate-group}). Let
$u\colon\mathbf C^g\to A(\mathbf C)$, with its trivializing functions, be a
uniformization of $A$ for $\mathcal L$ (Definition~\ref{cv-origin-hodge-norm-gram}(1), with
$T=\operatorname{Spec}\mathbf C$, $G=A$ and $\varphi$ the identity) with kernel $\Lambda$ and Appell--Humbert data
$(H,\chi)$, and put $E(a,b)=\operatorname{Im}(a^{\mathsf T}H\bar b)$ for $a,b\in\mathbf C^g$. Then:
\begin{enumerate}
\item $E$ takes integer values on $\Lambda\times\Lambda$, and $a\mapsto E(a,\cdot)$ is a bijection of $\Lambda$ onto
  $\operatorname{Hom}(\Lambda,\mathbf Z)$;
\item $g=h$, every tangent vector of $A$ at the origin has exactly one differential, and every $a\in\mathbf C^g$ is the
  differential of exactly one tangent vector of $A$ at the origin;
\item $e_n\big(u(a/n),u(b/n)\big)=\exp(-2\pi iE(a,b)/n)$ for all $n\ge1$ and $a,b\in\Lambda$.
\end{enumerate}
\end{lemma}
\leanref{Development/AbelianVarietyTheory/Proofs/CubicalLineBundles/AppellHumbertUnifTangentWeilPairingStatement.lean}{GlobalTheory.T.appellHumbert_unif_tangent_weilPairing}
\leanref{Development/AbelianVarietyTheory/Proofs/CubicalLineBundles/AppellHumbertUnifTangentWeilPairing.lean}{GlobalTheory.AHWeil.tangent_and_weil_of_isAH}
\begin{proof}
Step 1: a theta uniformization. The line bundle $\mathcal L$ is trivialized along the zero section; let $r$ be a
generator of $0^*\mathcal L$. Theorem~\ref{cv-cusp-cm-point-eigen-degree-rank-one-norm-riemann-theta}, for $A$,
$\mathcal L$ and $r$, gives a uniformization $u_0\colon\mathbf C^h\to A(\mathbf C)$ for $\mathcal L$ with kernel
$\Lambda_0$ and Appell--Humbert data $(H_0,\chi_0)$ for which (1), (2) and (3) hold, with
$E_0(a,b)=\operatorname{Im}(a^{\mathsf T}H_0\bar b)$ in place of $E$: they are (5), (2) and (6) of that theorem.
% Lean: the hypothesis ThetaUnifAV of the Lean statement (proved as T.theta_uniformisation_complex_abelian_varieties),
% applied to A with its rigidification P.rigid; tangent_and_weil_of_isAH keeps IsAH, the two tangent bijections, the
% formula for E_0 and the Weil formula, and unimodular_and_tangent_of_isAH (for (1)) uses ThetaUnifAV again.

Step 2: the comparison. By (2) for $u_0$, Lemma~\ref{cv-origin-hodge-norm-form-lift} applies to $u_0$ and to the
homomorphism $u$, which is holomorphic by Definition~\ref{cv-origin-hodge-norm-gram}(1), and gives a $\mathbf C$-linear
map $F\colon\mathbf C^g\to\mathbf C^h$ with $u_0\circ F=u$. For $z\in\mathbf C^g$, $u(z)$ is the origin if and only if
$u_0(F(z))$ is, so $F^{-1}(\Lambda_0)=\Lambda$. The kernel of $F$ is a subspace contained in the discrete group
$\Lambda$, hence $0$. Every $w\in\mathbf C^h$ has $u_0(w)=u(z)=u_0(F(z))$ for some $z$, so $\mathbf C^h$ is the union of
the closed sets $F(\mathbf C^g)+l$, $l\in\Lambda_0$; as $\Lambda_0$ is countable, one of them has an interior point by
Baire's theorem, and the subspace $F(\mathbf C^g)$ is $\mathbf C^h$. So $F$ is an isomorphism, $g=h$, and
$F(\Lambda)=\Lambda_0$.

Step 3: the forms. Lemma~\ref{cv-origin-hodge-norm-form-pullback}, applied to the identity of $A$, to $u$ and $u_0$, to
$F$ and to $n=1$, gives $z^{\mathsf T}H\bar w=F(z)^{\mathsf T}H_0\overline{F(w)}$ for all $z,w$. So
$E(a,b)=E_0(F(a),F(b))$, and (1) for $u$ follows from (1) for $u_0$, as $F$ restricts to a bijection of $\Lambda$ onto
$\Lambda_0$.

Step 4: differentials and pairings. By the chain rule, as in the proof of Lemma~\ref{cv-origin-hodge-norm-form-compare},
a morphism $\operatorname{Spec}\mathbf C[\varepsilon]\to A$ over $\mathbf C$ has differential $a$ for $u$ if and only if it
has differential $F(a)$ for $u_0$; as $F$ is bijective, (2) for $u$ follows from (2) for $u_0$. Finally, for $a,b\in\Lambda$,
$u(a/n)=u_0(F(a)/n)$ and $u(b/n)=u_0(F(b)/n)$, and (3) for $u_0$ at $F(a),F(b)\in\Lambda_0$ gives (3) for $u$.
\end{proof}
\leanref{Development/AbelianVarietyTheory/Proofs/ThetaLawProofs/UnimodularAndTangentOfIsAH.lean}{GlobalTheory.unimodular_and_tangent_of_isAH}
\leanref{Development/AbelianVarietyTheory/Proofs/CubicalLineBundles/AppellHumbertUnifTangentWeilPairing.lean}{GlobalTheory.AHWeil.dim_eq}
% Lean: tangent_and_weil_of_isAH: hT (ThetaUnifAV) gives u0, H0, chi0, E0 and the Weil formula hW;
% existsUnique_linear_lift_of_isHolTo (Lemma cv-origin-hodge-norm-form-lift) gives F; dim_eq (Step 2, through
% AnyUnif.injective_of_ker_le_discrete and eq_top_of_add_discrete, the Baire argument); unimodular_and_tangent_of_isAH
% for (1) and (2), with UnifCompare.hermForm_eq (Step 3; it repeats the argument of Lemma
% cv-origin-hodge-norm-form-pullback, gluing the ratios of the trivializing functions and applying
% appellHumbert_trivial_of_nowhere_vanishing), UnifCompare.mem_iff for F^{-1}(Lambda_0) = Lambda, and isDiff_iff
% (Step 4); and hW transported along F for (3).

\begin{lemma}\label{cv-cusp-cm-point-eigen-degree-rank-one-norm}
Let $R$ be a ring and $s\colon R\to\mathbf C$ a ring homomorphism. Let $A$ and $A'$ be principally polarized abelian
schemes over $R$, with the line bundles $\mathcal L$ and $\mathcal L'$ of their polarizations, and let $f\colon A\to
A'$ and $g\colon A\to A$ be homomorphisms such that, with $p_1$ and $p_2$ the projections of $A\times_RA$,
\[
 \Lambda_{\mathcal L'}(fp_1,fp_2)\cong\Lambda_{\mathcal L}(gp_1,p_2)
\]
on $A\times_RA$ (Definition~\ref{cv-toroidal-abelian-dvr-ample-torsion-weil}). Let $v$ be a tangent vector of $A$ over
$R$ and $a\in R$ with $g\circ v=av$. Let $r$ and $r'$ be the squared Hodge norms
(Definition~\ref{cv-cm-point-periods}(4)) of the images of $v$ and of $f\circ v$ in the base changes of $A$ and $A'$
along $s$, computed in uniformizations with Appell--Humbert data (Definition~\ref{cv-origin-hodge-norm-gram}(1)). Then
$r'=\operatorname{Re}(s(a))\,r$.
\end{lemma}
\leanref{Development/AbelianVarietyTheory/Proofs/CubicalLineBundles/LieAlgebraSquaredLengthCompPairingHomStatement.lean}{GlobalTheory.T.lieAlgebraSquaredLength_comp_pairing_hom}
\leanref{Development/AbelianVarietyTheory/Proofs/CubicalLineBundles/LieAlgebraSquaredLengthCompPairingHom.lean}{GlobalTheory.NormHom.lieNormSq_comp}
\begin{proof}
Step 1: the base changes. Let $Q$ and $Q'$ be the base changes of $A$ and $A'$ along $s$, principally polarized abelian
varieties over $\mathbf C$ with the base changes $\mathcal L_Q$ and $\mathcal L_{Q'}$ of the line bundles. Let $v_1$,
$f_1\colon Q\to Q'$ and $g_1\colon Q\to Q$ be the base changes of $v$, $f$ and $g$; $f_1$ and $g_1$ are homomorphisms,
$f_1\circ v_1$ is the image of $f\circ v$, and $g_1\circ v_1=s(a)v_1$. Pulling back commutes with the formation of
$\Lambda$, so the isomorphism of the hypothesis gives an isomorphism
$\Lambda_{\mathcal L_{Q'}}(f_1p_1,f_1p_2)\cong\Lambda_{\mathcal L_Q}(g_1p_1,p_2)$ on $Q\times Q$. Let
$u\colon\mathbf C^g\to Q(\mathbf C)$ and $u'\colon\mathbf C^{g'}\to Q'(\mathbf C)$ be the uniformizations in which the
squared Hodge norms are computed, with kernels $\Lambda$ and $\Lambda'$ and Appell--Humbert data $(H,\chi)$ and
$(H',\chi')$, and let $w$ and $w'$ be the differentials of $v_1$ and $f_1\circ v_1$, so that
$r=\pi\operatorname{Re}(w^{\mathsf T}H\bar w)$ and $r'=\pi\operatorname{Re}(w'^{\mathsf T}H'\bar w')$.

Step 2: linear lifts. By Lemma~\ref{cv-cusp-cm-point-eigen-degree-rank-one-norm-riemann}(2), $u$ and $u'$ satisfy the
hypotheses on differentials of Lemma~\ref{cv-origin-hodge-norm-form-lift}. That lemma, applied to $u'$ and the
homomorphism $f_1\circ u$, and to $u$ and $g_1\circ u$, which are holomorphic as $f_1$ and $g_1$ are morphisms, gives
$\mathbf C$-linear maps $F\colon\mathbf C^g\to\mathbf C^{g'}$ and $G\colon\mathbf C^g\to\mathbf C^g$ with
$u'\circ F=f_1\circ u$ and $u\circ G=g_1\circ u$. They carry $\Lambda$ into $\Lambda'$ and into $\Lambda$, as $f_1$ and
$g_1$ fix the origin. By the chain rule, as in the proof of Lemma~\ref{cv-origin-hodge-norm-form-compare}, $f_1\circ v_1$
has differential $F(w)$ and $g_1\circ v_1$ has differential $G(w)$. The differential of $s(a)v_1$ is $s(a)w$. As
differentials are unique, $w'=F(w)$ and $G(w)=s(a)w$.

Step 3: the Hermitian forms. Put $E(z,z')=\operatorname{Im}(z^{\mathsf T}H\bar z')$ and
$E'(z,z')=\operatorname{Im}(z^{\mathsf T}H'\bar z')$. Let $a_1,a_2\in\Lambda$ and $n\ge1$, and put $x=u(a_1/n)$ and
$y=u(a_2/n)$, points killed by $n$. Then $f_1x=u'(F(a_1)/n)$, $f_1y=u'(F(a_2)/n)$ and $g_1x=u(G(a_1)/n)$, with
$F(a_1),F(a_2)\in\Lambda'$ and $G(a_1)\in\Lambda$. Lemma~\ref{cv-cusp-cm-point-eigen-degree-rank-one-norm-weil}, for
$X=Y=Q$, $A=Q'$ and $B=Q$, with $f_1$ in place of both of its homomorphisms to $A$, $g_1$ and the identity in place of
its homomorphisms to $B$, and the isomorphism of Step 1, gives equal commutator pairings at $(f_1x,f_1y)$ and at
$(g_1x,y)$. By
Lemma~\ref{cv-cusp-cm-point-eigen-degree-rank-one-norm-riemann}(1) and (3) for $u'$ and for $u$, the integers
$m'=E'(F(a_1),F(a_2))$ and $m=E(G(a_1),a_2)$ satisfy $\exp(-2\pi im'/n)=\exp(-2\pi im/n)$. So $n$ divides $m'-m$ for
every $n\ge1$, and $m'=m$. The $\mathbf R$-bilinear forms $(z,z')\mapsto E'(F(z),F(z'))$ and $(z,z')\mapsto
E(G(z),z')$ therefore agree on $\Lambda\times\Lambda$, which spans $\mathbf C^g\times\mathbf C^g$ over $\mathbf R$, so
they agree everywhere. For a form $B$ that is linear in the first variable and conjugate linear in the second,
$\operatorname{Re}B(z,z')=\operatorname{Im}B(iz,z')$. Hence
\[
 F(z)^{\mathsf T}H'\,\overline{F(z')}=G(z)^{\mathsf T}H\bar z'\qquad(z,z'\in\mathbf C^g).
\]

Step 4: the norms. As $H$ is Hermitian, $w^{\mathsf T}H\bar w$ is real. By Steps 2 and 3,
$w'^{\mathsf T}H'\bar w'=F(w)^{\mathsf T}H'\,\overline{F(w)}=G(w)^{\mathsf T}H\bar w=s(a)\,w^{\mathsf T}H\bar w$. Taking
real parts, $r'=\pi\operatorname{Re}(s(a))\,w^{\mathsf T}H\bar w=\operatorname{Re}(s(a))\,r$.
\end{proof}
\leanref{Development/AbelianVarietyTheory/Proofs/CubicalLineBundles/LieAlgebraSquaredLengthCompPairingHom.lean}{GlobalTheory.NormHom.exists_monHom_lift}
\leanref{Development/AbelianVarietyTheory/Proofs/CubicalLineBundles/ExistsLinearOfIsMonHom.lean}{GlobalTheory.HomLinear.exists_linear_of_isMonHom}
\leanref{Development/AbelianVarietyTheory/Proofs/CubicalLineBundles/IsCubeLBPairingIsoBaseChangeStatement.lean}{GlobalTheory.T.isCubeLBPairingIso_baseChange}
\leanref{Development/AbelianVarietyTheory/Proofs/CubicalLineBundles/WeilEqIffHermFormEqComplexStatement.lean}{GlobalTheory.T.weilEq_iff_hermForm_eq_complex}
\leanref{Development/AbelianVarietyTheory/Proofs/CubicalLineBundles/WeilEqIffHermFormEqComplex.lean}{GlobalTheory.CM.HermWeil.hermForm_eq_of_im}
% Lean: lieNormSq_comp, with IsLieNormSq unfolded for v and f o v (Step 1: the fibres Q and Q', u, Hm, chi, w and the
% equation r = pi Re(hermForm Hm w w)). Step 1 also uses exists_monHom_lift for f_1 and g_1 and
% isCubeLBPairingIso_baseChange (= CM.lamEq_baseChange) for the isomorphism. Step 2 is
% HomLinear.exists_linear_of_isMonHom (Lemma cv-origin-hodge-norm-form-lift with isDiff_comp), with the tangent
% bijections t1, t2 from appellHumbert_unif_tangent_weilPairing, existsUnique_isVecOver for v'_1 = v_1 f_1, and
% lieComb_baseChange with UnifDifferential.isDiff_of_isLieComb for G(w) = s(a) w. Step 3 is
% weil_pairing_eq_of_isomorphic_cube_pairings (Lemma cv-cusp-cm-point-eigen-degree-rank-one-norm-weil) and
% weilEq_iff_hermForm_eq_complex (= CM.weilEq_iff_hermForm_eq: HermWeil.weilIs_iff, map_mem, ex_div_pow,
% eq_of_rootVal and hermForm_eq_of_im). Step 4 is him, e_4 and Complex.mul_re. Size: about 2,600 lines below
% lieNormSq_comp (1,956 found by following its imports, and the 712 below AdjointOfPairing.exists_gen_pairing), of
% which about 1,100 are in Lemma cv-cusp-cm-point-eigen-degree-rank-one-norm-weil-translate.

\begin{lemma}\label{cv-cusp-cm-point-eigen-degree-rank-one}
In Situation~\ref{cv-cusp-setting}, let $\sqrt{-d}\in E$ be the square root of $-d$ of the datum, and put
$w_0=(1+\sqrt{-d})/2$ and $\mathcal O=\mathbf Z[w_0]$. Let $K$ be a number field with $[E:\mathbf Q]$ embeddings
$E\to K$, let $\omega\in\mathcal O_K$, and let $C$ be a principally polarized elliptic curve over $\mathcal O_K$ with
an action of $\mathcal O$ whose Rosati involution is complex conjugation, and a basis $v_C$ of its Lie algebra on which
$w_0$ acts by $\omega$. For an embedding $s\colon K\to\mathbf C$ let $n(s)$ be the squared Hodge norm
(Definition~\ref{cv-cm-point-periods}(4)) of the image of $v_C$ in the complex curve of $C$ at $s$, and put
$x=-\frac1{[K:\mathbf Q]}\sum_s\log n(s)$. Let $L$ be the lattice of a PEL datum over $E$ of signature $(p,q)$ with
$p+q=1$, with its perfect alternating form $\langle\ ,\ \rangle$, let $\varepsilon\in\{\pm1\}$ with
$\varepsilon\langle w_0x',x'\rangle>0$ for every nonzero $x'\in L$, and put
$h(x',y')=\varepsilon(\langle w_0x',y'\rangle-\langle x',y'\rangle)+\varepsilon\langle x',y'\rangle w_0$. Let $A$ be a
principally polarized abelian scheme over $\mathcal O_K$ with an action of $\mathcal O_E$ for which the Rosati
involution induces complex conjugation on $\mathcal O_E$, and which is a Serre tensor product of $C$ and $L$ with form
$h$ (Definition~\ref{cv-cusp-cm-point-eigen-degree-serre}). Let $\tau\colon E\to K$ be an embedding with
$\tau(w_0)=\omega$. Then the $\tau$-eigenspace of the Lie algebra of $A_K$
(Definition~\ref{cv-cusp-cm-point-eigen-degree-line}) is a line, and for every nonzero $v$ in it
\[
 2\,\widehat\deg\mathcal M(A,v)=x+\tfrac12\log d-\log\operatorname{rd}_E.
\]
\end{lemma}
\leanref{Development/AbelianVarietyTheory/Proofs/CubicalLineBundles/TwoMulDegreeEigenLatticeTensorRankOneStatement.lean}{GlobalTheory.T.two_mul_degree_eigenLattice_tensor_rankOne}
\leanref{Development/AbelianVarietyTheory/Proofs/CubicalLineBundles/TwoMulDegreeEigenLatticeTensorRankOne.lean}{GlobalTheory.CM.TensorDeg.two_mul_deg_eigLat}
\begin{proof}
Write $\operatorname{Lie}A$ for the Lie algebra of $A$ over $\mathcal O_K$, and $t_a$ for the homomorphisms of
Definition~\ref{cv-cusp-cm-point-eigen-degree-serre}.

Step 1: the character of $\mathcal O$. The Lie algebra of $C$ is free of rank one with basis $v_C$, so every
$k\in\mathcal O$ acts on it by an element $\tau_0(k)\in\mathcal O_K$, and $\tau_0\colon\mathcal O\to\mathcal O_K$ is a
ring homomorphism with $\tau_0(w_0)=\omega$. Let $I$ be the set of embeddings $\sigma\colon E\to K$ with
$\sigma(\sqrt{-d})=\tau(\sqrt{-d})$; it contains $\tau$. Every $\sigma\in I$ agrees with $\tau_0$ on $\mathcal O$, as
$\sigma(w_0)=(1+\sigma(\sqrt{-d}))/2=\tau(w_0)=\omega$.

Step 2: the lattice. Take $\mathfrak a$, $e$ and $\beta$ as in Lemma~\ref{cv-cusp-cm-point-eigen-degree-ideal}. By
its part (1) for $\Omega=K$ and the embedding $\tau$, and by Step~1, as $h(a,b)\in\mathcal O$,
\[
 \tau_0(h(a,b))=\sum_{\sigma\in I}\sigma\big(\beta\,e(a)\overline{e(b)}\big)\qquad(a,b\in L).\tag{$\dagger$}
\]
Take $x_1,\dots,x_N$ and $y_1,\dots,y_N$ as in Lemma~\ref{cv-cusp-cm-point-eigen-degree-dual}.

Step 3: the Lie algebra. For $a\in L$ let $c(a)\in\operatorname{Lie}A$ be the image of $v_C$ under the derivative of
$t_a$. The Lie algebra of an abelian scheme over $\mathcal O_K$ is the kernel of the map from its group of points
over the dual numbers $\mathcal O_K[\epsilon]$ to its group of points over $\mathcal O_K$, a map that the inclusion
$\mathcal O_K\subset\mathcal O_K[\epsilon]$ splits. So the natural isomorphisms of the Serre tensor product identify
$\operatorname{Lie}A$ with
$L\otimes_{\mathcal O}\mathcal O_Kv_C$, where $\mathcal O$ acts on $\mathcal O_K$ through $\tau_0$, with $a\otimes
v_C\mapsto c(a)$. So $\operatorname{Lie}A$ is the $\mathcal O_K$-span of the $c(a)$, and $c$ is additive. By (1) of
the definition, $c(ba)=b\,c(a)$ for $b\in\mathcal O_E$. As $t_a$ commutes with $\mathcal O$, $t_{ka}=t_a\circ k$ for
$k\in\mathcal O$, and $c(ka)=\tau_0(k)c(a)$.

Step 4: the Hodge forms. For an embedding $s\colon K\to\mathbf C$ and tangent vectors $u,u'$ of the base change of
$A_K$ along $s$, put $\langle u,u'\rangle_s=\pi H(u',u)$, where $H$ is the Hermitian form of the polarization, linear
in the first variable; so $\langle u,u\rangle_s$ is the squared Hodge norm of $u$, and $\langle\ ,\ \rangle_s$ is
linear in its second variable. Let $a\in L$ and put $k=h(a,a)$. For the line bundles $\mathcal L_A$ and $\mathcal L_C$
of the polarizations, (2) of the definition is an isomorphism
$\Lambda_{\mathcal L_A}(t_ap_1,t_ap_2)\cong\Lambda_{\mathcal L_C}(kp_1,p_2)$ on $C\times C$, with $p_1$ and $p_2$ the
projections (Definition~\ref{cv-toroidal-abelian-dvr-ample-torsion-weil}), and by Step~1 the action of $k$ multiplies
$v_C$ by $\tau_0(k)$. Lemma~\ref{cv-cusp-cm-point-eigen-degree-rank-one-norm}, for $R=\mathcal O_K$, the restriction
of $s$, $C$, $A$, $f=t_a$, the action of $k$ as $g$, $v=v_C$ and $\tau_0(k)$, gives
$\langle c(a),c(a)\rangle_s=\operatorname{Re}(s(\tau_0(k)))\,n(s)$. As the form is alternating,
$k=\varepsilon\langle w_0a,a\rangle$ is an integer, and $s(\tau_0(k))=k$ is real. So
\[
 \langle c(a),c(a)\rangle_s=s(\tau_0(h(a,a)))\,n(s)\qquad(a\in L).\tag{$\ddagger$}
\]

Step 5: coordinates. For $a,b\in L$ put $\operatorname{ev}(a)=(\sigma(e(a)))_{\sigma\in I}$ and
$\operatorname{dv}(b)=(\sigma(\beta\overline{e(b)}))_{\sigma\in I}$, vectors in $K^I$; by $(\dagger)$,
$\tau_0(h(a,b))=\sum_\sigma\operatorname{ev}(a)_\sigma\operatorname{dv}(b)_\sigma$. Let $\Phi$ be the $K$-linear map
from $K^I$ to the Lie algebra of $A_K$ given by
\[
 \Phi(z)=\sum_m\Big(\sum_{\sigma\in I}z_\sigma\operatorname{dv}(y_m)_\sigma\Big)c(x_m).
\]
By Step~3 and the dual families,
$\Phi(\operatorname{ev}(a))=\sum_m\tau_0(h(a,y_m))c(x_m)=c\big(\sum_mh(a,y_m)x_m\big)=c(a)$.
\begin{enumerate}
\item[(a)] The $\operatorname{ev}(a)$ span $K^I$. If $g\in K^I$ has $\sum_\sigma g_\sigma\sigma(e(a))=0$ for every
  $a$, fix $a_0$ with $e(a_0)\ne0$; replacing $a$ by $ba_0$ gives
  $\sum_\sigma g_\sigma\sigma(e(a_0))\sigma(b)=0$ for every $b\in\mathcal O_E$. The restrictions of the distinct
  $\sigma$ to $\mathcal O_E$ are distinct characters, hence linearly independent (Artin), so
  $g_\sigma\sigma(e(a_0))=0$ and $g=0$.
\item[(b)] For every $s$ and all $z,z'\in K^I$,
  $\langle\Phi z,\Phi z'\rangle_s=\sum_\sigma b_{s,\sigma}\overline{s(z_\sigma)}\,s(z'_\sigma)$ with
  $b_{s,\sigma}=s(\sigma(\beta))\,n(s)>0$. By (a) and sesquilinearity it suffices to take $z=\operatorname{ev}(a)$ and
  $z'=\operatorname{ev}(a')$; let $B(a,a')$ be the difference of the two sides. Then $B$ is additive in $a'$,
  $B(a',a)=\overline{B(a,a')}$, and $B(a,w_0a')=s(\omega)B(a,a')$, since $c(w_0a')=\omega c(a')$ and
  $\sigma(e(w_0a'))=\omega\,\sigma(e(a'))$ for $\sigma\in I$. Moreover $B(a,a)=0$: by $(\dagger)$ applied with $s$, and as
  $s\sigma(\bar y)=\overline{s\sigma(y)}$, $\sum_\sigma b_{s,\sigma}|s\sigma(e(a))|^2=s(\tau_0(h(a,a)))\,n(s)$, which
  is $\langle c(a),c(a)\rangle_s$ by $(\ddagger)$. Expanding $B(a+a',a+a')=0$ gives $\operatorname{Re}B(a,a')=0$ for all
  $a,a'$. Then
  $0=\operatorname{Re}B(a,w_0a')=-\operatorname{Im}s(\omega)\operatorname{Im}B(a,a')$, and
  $\operatorname{Im}s(\omega)\ne0$, as $2s(\omega)-1=s\tau(\sqrt{-d})$ is a square root of $-d$. So $B=0$.
\item[(c)] By (b), $\Phi$ is injective; its image contains the $c(a)$, which span the Lie algebra of $A_K$ over $K$ by
  Step~3. So $\Phi$ is an isomorphism, and it carries the $\mathcal O_K$-span of the $\operatorname{ev}(a)$ onto
  $\operatorname{Lie}A$.
\item[(d)] Every $b\in\mathcal O_E$ acts on the Lie algebra of $A_K$ by
  $\Phi(z)\mapsto\Phi\big((\sigma(b)z_\sigma)_\sigma\big)$: for $z=\operatorname{ev}(a)$ both are $c(ba)$, and both
  are $K$-linear in $z$.
\item[(e)] For $\sigma_0\in I$ and $t\in K$, the vector $te_{\sigma_0}$, with $e_{\sigma_0}$ the unit vector at
  $\sigma_0$, lies in the $\mathcal O_K$-span of the $\operatorname{ev}(a)$ if and only if
  $t\operatorname{dv}(b)_{\sigma_0}\in\mathcal O_K$ for every $b\in L$. The linear form
  $z\mapsto\sum_\sigma z_\sigma\operatorname{dv}(b)_\sigma$ takes the value $\tau_0(h(a,b))\in\mathcal O_K$ at
  $\operatorname{ev}(a)$, so it takes values in $\mathcal O_K$ on the span, and its value at $te_{\sigma_0}$ is
  $t\operatorname{dv}(b)_{\sigma_0}$. Conversely,
  $z=\sum_m\big(\sum_\sigma z_\sigma\operatorname{dv}(y_m)_\sigma\big)\operatorname{ev}(x_m)$ for every $z\in K^I$: for
  $z=\operatorname{ev}(a)$ the right side is $\sum_m\tau_0(h(a,y_m))\operatorname{ev}(x_m)=\operatorname{ev}(a)$, as
  $\sigma(e(ka))=\tau_0(k)\sigma(e(a))$ for $k\in\mathcal O$, and both sides are linear in $z$. For $z=te_{\sigma_0}$
  the coefficients are the $t\operatorname{dv}(y_m)_{\sigma_0}$.
\end{enumerate}

Step 6: the eigenline. Put $v_0=\Phi(e_\tau)$. For $\sigma\in I$ with $\sigma\ne\tau$ some $b\in\mathcal O_E$ has
$\sigma(b)\ne\tau(b)$, as $\mathcal O_E$ spans $E$. So by (c) and (d) the $\tau$-eigenspace of the Lie algebra of $A_K$
is $\Phi(Ke_\tau)=Kv_0$, a line, and by (b) $\langle v_0,v_0\rangle_s=s(\tau(\beta))\,n(s)$. By (c) and (e),
$tv_0\in\operatorname{Lie}A$ if and only if $t\tau(\beta\overline{e(b)})\in\mathcal O_K$ for every $b\in L$, that is,
$t\tau(\beta\bar{\mathfrak a})\subset\mathcal O_K$. So the ideal of $\mathcal M(A,v_0)$ is $\mathfrak b^{-1}$, where
$\mathfrak b$ is the fractional ideal of $K$ generated by $\tau(\beta\bar{\mathfrak a})$, and
$n_s(\mathcal M(A,v_0))^2=s(\tau(\beta))\,n(s)$.

Step 7: the degree. The norm of $\mathfrak b$ is
$\mathrm N(\beta\bar{\mathfrak a})^{[K:E]}=\big(\mathrm N_{E/\mathbf Q}(\beta)\mathrm N(\mathfrak a)\big)^{[K:E]}$, and
every embedding of $E$ in $\mathbf C$ has $[K:E]$ extensions to $K$, so
$\sum_s\log s(\tau(\beta))=[K:E]\log\mathrm N_{E/\mathbf Q}(\beta)$. Hence
\[
 \deg\mathcal M(A,v_0)=\log\mathrm N(\mathfrak b)-\tfrac12\sum_s\log s(\tau(\beta))-\tfrac12\sum_s\log n(s)
 =[K:E]\big(\log\mathrm N(\mathfrak a)+\tfrac12\log\mathrm N_{E/\mathbf Q}(\beta)\big)-\tfrac12\sum_s\log n(s),
\]
and, as $[K:\mathbf Q]=[K:E]\,[E:\mathbf Q]$, Lemma~\ref{cv-cusp-cm-point-eigen-degree-ideal}(2) gives
\[
 2\,\widehat\deg\mathcal M(A,v_0)=\frac{2\log\mathrm N(\mathfrak a)+\log\mathrm N_{E/\mathbf Q}(\beta)}{[E:\mathbf Q]}+x
 =x+\tfrac12\log d-\log\operatorname{rd}_E.
\]
A nonzero $v$ in the eigenspace is $tv_0$ with $t\in K^\times$. The ideal of $\mathcal M(A,v)$ is $t^{-1}\mathfrak
b^{-1}$ and $n_s(\mathcal M(A,v))=|s(t)|\,n_s(\mathcal M(A,v_0))$, so the degree changes by
$\log|\mathrm N_{K/\mathbf Q}(t)|-\sum_s\log|s(t)|=0$.
\end{proof}
\leanref{Development/AbelianVarietyTheory/Proofs/CubicalLineBundles/TwoMulDegreeEigenLatticeTensorRankOne.lean}{GlobalTheory.CM.TensorDeg.exists_eigenline}
\leanref{Development/AbelianVarietyTheory/Proofs/CubicalLineBundles/SingleMemSpanIffOfSelfDual.lean}{GlobalTheory.CM.single_mem_span_iff_of_selfDual}
\leanref{Development/AbelianVarietyTheory/Proofs/CubicalLineBundles/LieTensorLine.lean}{GlobalTheory.CM.lie_tensor_line}
\leanref{Development/AbelianVarietyTheory/Proofs/CubicalLineBundles/LieAlgebraSquaredLengthCompPairingHomStatement.lean}{GlobalTheory.T.lieAlgebraSquaredLength_comp_pairing_hom}
\leanref{Development/AbelianVarietyTheory/Proofs/CubicalLineBundles/TwoMulDegDualLineCmIdealStatement.lean}{GlobalTheory.T.two_mul_deg_dual_line_cmIdeal}
% Lean: two_mul_deg_eigLat. Step 1 is LieChar.ell_char (exists_char), with existsUnique_extend_ringOfIntegers_hom for
% tau' and hext. Step 2 is exists_fractionalIdeal_traceForm_rank_one with hform, and exists_unique_cmHermitianLatticeForm
% (its uniqueness hku identifies the form with h). Step 3 is lie_tensor_line (Serre.IsTensor.lieEquivLine, through the
% splitting of the points over the dual numbers), with integral_tangent_coordinates for the coordinates co and the map
% c, t.lin and cosmul; hlat. Step 4 is lieAlgebraSquaredLength_comp_pairing_hom = NormHom.lieNormSq_comp (Lemma
% cv-cusp-cm-point-eigen-degree-rank-one-norm), applied with R = O_K, f = t.ten.tmul a, g = C.act.act (h a a), the
% isomorphism t.lam a a and C.basis, hτ₀ for the action on v_C; HermNorm.isLieNormSq_form identifies the Hodge metric
% h.form with squared Hodge norms, and the Lean's exists_eigenline takes the real part. Lemma cv-cm-point-trace(3) is
% not used: its Lean statement (lie_norm_square_of_family) is about Hodge Gram matrices of families with IsFam, a
% different declaration from lieNormSq_comp. Steps 5 and 6 are exists_eigenline: (a) span_embeddings_eq_top (linearIndependent_monoidHom),
% (b) form_eq_of_span and eq_zero_of_re_diag, (c) hinj and hmap, (d) hact, (e) single_mem_span_iff_of_selfDual; the
% eigen-lattice is described through exists_lieRepresentation_eigenLattices. Step 7 is two_mul_deg_dual_line_cmIdeal
% (= CM.two_mul_deg_dual_line, through exists_fractionalIdeal_eq_span_image, meanLog_comp and deg_ideal_smul), with
% hnorm. The Lean states the formula for deg of the eigen-lattice; that the eigenspace is a line, and the last
% paragraph, are in the proof here (Phi is an isomorphism; lieLattice_fg_and_span_eq_top_ringOfIntegers), and the
% independence of v is deg_eq_ndeg_of_line with deg_ideal_smul, as in Lemma cv-cusp-cm-point-degree.

\begin{lemma}\label{cv-cusp-cm-point-eigen-degree}
In Situation~\ref{cv-cusp-setting}, let $\sqrt{-d}\in E$ be the square root of $-d$ of the datum. Let $K$ be a
number field with $[E:\mathbf Q]$ embeddings $E\to K$, let $\rho\in\mathcal O_K$ with $\rho^2=-d$, and let $A$ be a
principally polarized abelian scheme over $\mathcal O_K$ with an action of $\mathcal O_E$ for which the Rosati
involution of the polarization induces complex conjugation on $\mathcal O_E$. Suppose that the Lie algebra of the
generic fibre $A_K$ has a basis of eigenvectors of $\mathcal O_E$ whose eigencharacters are the embeddings
$\tau\colon E\to K$ with $\tau(\sqrt{-d})=\rho$, each once. Let $\tau$ be one of these embeddings and $v$ a nonzero
vector of the Lie algebra of $A_K$ on which $\mathcal O_E$ acts through $\tau$. Let $\mathcal M$ be the metrized line
over $K$ (Definition~\ref{cv-cusp-cm-lines}) whose ideal is the set of $c\in K$ with $cv$ in the Lie algebra of $A$,
and for which $n_s(\mathcal M)^2$ is the squared Hodge norm (Definition~\ref{cv-cm-point-periods}(4)) of the image of
$v$ in the base change of $A_K$ along $s$, for every embedding $s\colon K\to\mathbf C$. Then
$\widehat\deg\mathcal M=\frac12\log\Omega_E(\chi)$.
\end{lemma}
\leanref{Development/AbelianVarietyTheory/Proofs/CubicalLineBundles/TwoMulDegEigenLatticePolarizedAbelianSchemeStatement.lean}{GlobalTheory.T.two_mul_deg_eigenLattice_polarizedAbelianScheme}
\leanref{Development/AbelianVarietyTheory/Statements/PrincipalLattices/ConclusionsAtIdentity.lean}{GlobalTheory.CM.DegSt}
\begin{proof}
Put $w_0=(1+\sqrt{-d})/2$ and $x=\gamma+\log(2\pi/\sqrt d)-\frac{L'}L(1,\chi)$, a real number. By hypothesis the
$\tau$-eigenspace of the Lie algebra of $A_K$ (Definition~\ref{cv-cusp-cm-point-eigen-degree-line}) is the line
spanned by $v$, and $\mathcal M$ is the metrized line $\mathcal M(A,v)$ of $v$.

Step 1: a CM curve. Fix an embedding $\sigma\colon K\to\mathbf C$ and put $\theta=\sigma(\rho)$, so that
$\theta^2=-d$. By Lemma~\ref{cv-cm-point}(1) there is a CM curve $C$ with root $\theta$ over a number field
$M\subset\mathbf C$ (Definition~\ref{cv-cm-point-curve}) with $-\frac1{[M:\mathbf Q]}\sum_s\log n(s)=x$. Let
$k\subset\mathbf C$ be the number field generated by $\sigma(K)$ and $M$, and let $j\colon K\to k$ be the map
$\sigma$. By Lemma~\ref{cv-cusp-cm-point-eigen-degree-curve}, the base change $C_k$ of $C$ to $\mathcal O_k$ is a CM
curve with root $\theta$ over $k$, and the mean of the logarithms of its squared Hodge norms is $-x$. Its root
$\theta=j(\rho)$ lies in $\mathcal O_k$, and $w_0$ acts on the basis of its Lie algebra by $\omega=(1+\theta)/2$.

Step 2: the base change of $A$. Let $A_k$ be the base change of $A$ to $\mathcal O_k$ along $j$, with the base changes
of the action and of the polarization, whose Rosati involution induces complex conjugation on $\mathcal O_E$. The Lie
algebra of its generic fibre is $k\otimes_K$ that of $A_K$, so the images of the basis of eigenvectors of the
hypothesis form a basis of eigenvectors whose eigencharacters are the $j\tau'$, for the embeddings
$\tau'\colon E\to K$ with $\tau'(\sqrt{-d})=\rho$, each once. As $K$ has $[E:\mathbf Q]$ embeddings $E\to K$, for an
algebraically closed field $\Omega$ and an embedding $\iota\colon k\to\Omega$ the $\iota j\tau'$ are the embeddings
$E\to\Omega$ with value $\iota(\theta)$ at $\sqrt{-d}$. So $A_k$ satisfies the hypothesis of
Lemma~\ref{cv-cusp-cm-point-eigen-degree-tensor} for $C_k$ and $r=1$.

Step 3: the tensor decomposition. Lemma~\ref{cv-cusp-cm-point-eigen-degree-tensor} gives a number field $k_1$ with
$k\subset k_1\subset\mathbf C$, the lattice $L$ of a PEL datum over $E$ of signature $(p,q)$ with $p+q=1$, a sign
$\varepsilon$ and the form $h$ of that lemma, such that the base change $A_1$ of $A_k$ to $\mathcal O_{k_1}$ is a
Serre tensor product of the base change $C_1$ of $C_k$ to $\mathcal O_{k_1}$ and $L$ with form $h$. By
Lemma~\ref{cv-cusp-cm-point-eigen-degree-curve}, $C_1$ is a CM curve with root $\theta$ over $k_1$, $w_0$ acts on the
basis of its Lie algebra by $\omega$, and the mean of the logarithms of its squared Hodge norms is $-x$.

Step 4: the degrees. Let $\tau_1\colon E\to k_1$ be $\tau$ followed by $j$ and the inclusion $k\subset k_1$. As
$2\tau(w_0)=1+\tau(\sqrt{-d})=1+\rho$ and $j(\rho)=\theta=2\omega-1$, $\tau_1(w_0)=\omega$, and $k_1$ has
$[E:\mathbf Q]$ embeddings of $E$, as $K$ has. Let $v_1$ be the image of $v$ in the Lie algebra of the generic fibre of
$A_1$. Lemma~\ref{cv-cusp-cm-point-eigen-degree-base}, for $j$ and then for $k\subset k_1$, shows that the
$\tau_1$-eigenspace of that Lie algebra is the line spanned by $v_1$, and that
$\widehat\deg\mathcal M(A_1,v_1)=\widehat\deg\mathcal M$. Lemma~\ref{cv-cusp-cm-point-eigen-degree-rank-one}, for
$k_1$, $\omega$, $C_1$, $L$, $\varepsilon$, $h$, $A_1$ and $\tau_1$, gives
$2\,\widehat\deg\mathcal M(A_1,v_1)=x+\frac12\log d-\log\operatorname{rd}_E$.

Step 5: the constant. By Definition~\ref{cv-log-omega},
\[
 x+\tfrac12\log d-\log\operatorname{rd}_E
 =\gamma+\log2\pi-\tfrac12\log d-\frac{L'}L(1,\chi)+\tfrac12\log d-\log\operatorname{rd}_E=\log\Omega_E(\chi),
\]
so $\widehat\deg\mathcal M=\frac12\log\Omega_E(\chi)$.
\end{proof}
\leanref{Development/AbelianVarietyTheory/Proofs/CubicalLineBundles/TwoMulDegEigenLatticePolarizedAbelianScheme.lean}{GlobalTheory.LamF.two_mul_deg_eigLat_of_polAV}
\leanref{Development/AbelianVarietyTheory/Proofs/CubicalLineBundles/CmEigenLineDegreeFormula.lean}{GlobalTheory.CM.PureDeg.degSt}
\leanref{Development/AbelianVarietyTheory/Proofs/CubicalLineBundles/ExistsCmEllipticCurveIntegersWithPeriodStatement.lean}{GlobalTheory.T.exists_cmEllipticCurve_integers_with_period}
\leanref{Development/AbelianVarietyTheory/Proofs/CubicalLineBundles/ExistsPatOfPolAV.lean}{GlobalTheory.LamF.exists_pat_of_polAV}
\leanref{Development/AbelianVarietyTheory/Proofs/CubicalLineBundles/DegEigLatEqOfIso.lean}{GlobalTheory.LamF.deg_eigLat_eq_of_iso}
\leanref{Development/AbelianVarietyTheory/Proofs/CubicalLineBundles/HasTypeBaseChange.lean}{GlobalTheory.CM.hasType_baseChange}
% Lean: two_mul_deg_eigLat_of_polAV first replaces the line bundle of the polarization of A by a normalized one with
% the same polarization (exists_pat_of_polAV, through exists_normalised_polarisation), which keeps the type
% (hasType_of_iso) and the degrees (deg_eigLat_eq_of_iso); in the text a polarization is not tied to a line bundle. It
% then applies DegSt, proved as PureDeg.degSt. Step 1: sigma from IsAlgClosed.lift, hrhoc, the curve from
% exists_cmEllipticCurve_integers_with_period (= CM.EllExists.exists_ell, the statement under Lemma cv-cm-point(1); its
% IsPrincipalAt is the lattice condition of Definition cv-cm-point-curve(3) for every uniformization),
% exists_compositum for k, j and i, then exists_ellMap, ellMap_comp (hprin') and hasPeriod_of_ellMap (hper'). Step 2:
% PAT.exists_pullback and hasType_baseChange (htype'), with hrhoomega for j(rho) = 2 omega - 1. Step 3:
% exists_tensor_decomposition_pure_type_cm_triple with r = 1, hasPeriod_of_ellMap (hper1), and hsp1 for the embeddings
% of E in k_1. Step 4: degree_lieLattices_cmTriple_baseChange twice (d1, d2), h2, h3, h4 and hw for tau_1(w0) =
% omega_1, and two_mul_degree_eigenLattice_tensor_rankOne (key); the Hodge metrics of the base changes come from
% GenFib.nonempty_of_baseChange and LieMetric.exists_unique_isHerm. Step 5 is hr with logOm and hx.

\begin{lemma}\label{cv-cusp-cm-point-degree}
In Situation~\ref{cv-cusp-setting}, let $K\supset F_c$ be a number field, let $\ell>0$, and let a good CM point of
$X$ over $K$ with bound $\ell$ be given, with the notation of Definition~\ref{cv-cusp-cm-point-good}. Then
$\widehat\deg\mathcal M^\Phi_i=\widehat\deg\mathcal M'_i=\frac12\log\Omega_E(\chi)$ for $1\le i\le f$.
\end{lemma}
\leanref{Development/AbelianVarietyTheory/Proofs/CubicalLineBundles/LamDegStLogOm.lean}{GlobalTheory.LamF.lamDegSt_logOm}
\leanref{Development/AbelianVarietyTheory/Definitions/CMLattices/LamLetters.lean}{GlobalTheory.LamF.LamDegSt}
\begin{proof}
Let $\rho\in\mathcal O_{F_c}\subset\mathcal O_K$ be the square root of $-d$ given by the datum. For an embedding
$\tau\colon E\to\mathbf C$, $\tau_K(\sqrt{-d})=\tau(\sqrt{-d})$ is $\rho$ if $\tau\in\Phi$ and $-\rho$ otherwise, and
the $\tau_K$ are $[E:\mathbf Q]$ distinct embeddings $E\to K$, so they are all of them.

Fix $i$. By (b) of the definition, the Rosati involution of the polarization of $A^\Phi$ induces complex conjugation
on $\mathcal O_E$, and every $b\in\mathcal O_E$ multiplies $v^\Phi_j$ by $\tau^\Phi(j)_K(b)$. As $\tau^\Phi$ is a
bijection onto $\Phi$, $v^\Phi$ is a basis of eigenvectors whose eigencharacters are the embeddings $E\to K$ with
value $\rho$ at $\sqrt{-d}$, each once. Lemma~\ref{cv-cusp-cm-point-eigen-degree}, for $A^\Phi$, $\rho$,
$\tau^\Phi(i)_K$ and $v=v^\Phi_i$, gives a metrized line $\mathcal M$ with
$\widehat\deg\mathcal M=\frac12\log\Omega_E(\chi)$, and it remains to see that $\mathcal M=\mathcal M^\Phi_i$. For
every $s$, $n_s(\mathcal M)^2$ and, by (e), $n_s(\mathcal M^\Phi_i)^2$ are the squared Hodge norm of the image of
$v^\Phi_i$ in the base change of $A^\Phi_K$ along $s$. Let $\mathfrak p$ be a nonzero prime, and let
$M$ and $c$ be the model matrix of $(A^\Phi,v^\Phi)$ and the local generator of $\mathcal M^\Phi_i$ at $\mathfrak p$
given by (e). For $c'\in K$, Lemma~\ref{cv-cusp-cm-point-model} and (e) show that $sc'v^\Phi_i$ lies in
$\operatorname{Lie}A^\Phi$ for some $s\in\mathcal O_K\setminus\mathfrak p$ if and only if
$\operatorname{ord}_{\mathfrak p}c'\ge\operatorname{ord}_{\mathfrak p}c$. The left side says that $c'$ lies in the
localization of $I_{\mathcal M}$ at $\mathfrak p$, so
$\operatorname{ord}_{\mathfrak p}I_{\mathcal M}=\operatorname{ord}_{\mathfrak p}c=\operatorname{ord}_{\mathfrak p}I_{\mathcal M^\Phi_i}$.
This holds at every $\mathfrak p$, so $I_{\mathcal M}=I_{\mathcal M^\Phi_i}$ and $\mathcal M=\mathcal M^\Phi_i$.

The same argument for $A'$, $v'$, $-\rho$ and $\tau'(i)_K$ gives $\widehat\deg\mathcal M'_i=\frac12\log\Omega_E(\chi)$,
as the $\tau_K$ with $\tau\notin\Phi$ are the embeddings $E\to K$ with value $-\rho$ at $\sqrt{-d}$.
\end{proof}
\leanref{Development/AbelianVarietyTheory/Proofs/CubicalLineBundles/CmLieLineDegreeOfEigenLineDegreeFormula.lean}{GlobalTheory.LamF.lamDegSt_of_degSt}
\leanref{Development/AbelianVarietyTheory/Proofs/CubicalLineBundles/NormalisedDegreeCmEigenLineEq.lean}{GlobalTheory.LamF.ndeg_vecLine_eq}
% Lean: lamDegSt_logOm = cmLieLineDegree_of_eigenLineDegreeFormula = lamDegSt_of_degSt, with the labels and rho of
% GlobalTheory.LamF.labK_spec, applies ndeg_vecLine_eq (= GlobalTheory.T.normalisedDegree_cmEigenLine_eq) to A^Phi
% with rho and to A' with -rho. ndeg_vecLine_eq: LieType.hasType_of_isLieType (the type), rosatiAdj_iff_lamEq (the
% Rosati condition), CM.LieMetric.exists_unique_isHerm (the Hodge metric), then deg_eq_ndeg_of_line with
% EigLine.mem_eigLat_iff, mem_lieVector_line_ideal_iff_mem_lie_lattice (= isVecLine_mem_iff_mem_lieLat: the ideals,
% through isModelAt_integral_iff_mem_lieLat) and LieMetric.vecLine_norm_sq_eq_form (the norms), and
% two_mul_deg_eigenLattice_polarizedAbelianScheme (Lemma cv-cusp-cm-point-eigen-degree).

\begin{lemma}\label{cv-cusp-cm-point-exists-lattice-pair}
In Situation~\ref{cv-cusp-setting}, let $\sigma_0$ be the inclusion of $F_c$ in $\mathbf C$ and $Z$ a component at
$\sigma_0$, with the lattice $L_Z$, its alternating form $\langle\ ,\ \rangle$ and the Hermitian model
$(G,\eta,\varphi_Z)$ (Definition~\ref{cv-uniformization-model}) that the compatibility conditions of the uniformization
attach to $Z$, as listed at the beginning of the proof of Lemma~\ref{cv-cover-uniformization}. Let $L'\subset E^r$ be
the lattice, self-dual for $G_Z$, that the uniformization provides for $Z$, and suppose that $\varphi_Z(L_Z)$ is the
$\mathcal O_E$-span of $\mathfrak AL'$ for a nonzero fractional ideal $\mathfrak A$ of $E$. Then there are lattices $L_+$
and $L_-$ of PEL data over $E$ of signatures $(1,0)$ and $(0,1)$, with forms $\langle\ ,\ \rangle_\pm$,
$\mathcal O_E$-linear isomorphisms $e_\pm\colon L_\pm\to\mathfrak A$ and signs $s_\pm\in\{1,-1\}$ with $s_-=-s_+$, such
that $\langle a,b\rangle_\pm=\operatorname{Tr}_{E/\mathbf Q}\big(-s_\pm\eta\,e_\pm(a)\overline{e_\pm(b)}\big)$ for
$a,b\in L_\pm$.
\end{lemma}
\leanref{Development/AbelianVarietyTheory/Proofs/CubicalLineBundles/ExistsCmPointLatticeDataOfAdaptedBasis.lean}{GlobalTheory.LamT.CMBasis.exists_pair}
\begin{proof}
\openstep{The construction is not yet written out: the form $\operatorname{Tr}_{E/\mathbf Q}(\pm\eta\,a\bar b)$ is
perfect on $\mathfrak A$, because $\langle\ ,\ \rangle$ is perfect on $L_Z$ and $L'$ is self-dual, and $\mathfrak A$ with
it is a lattice of PEL data of signature $(1,0)$ for one sign and $(0,1)$ for the other.}
% Lean: CMBasis.exists_pair: the sign t from X.kappa0_xi (kappa0 = 2 t xi0), trace_xi_mul_conj_int_iff_mem (the
% duality of A for the trace form, from hLambda and the perfectness EF_perf and L'_dual), SelfDualDatum.exists_datum
% (with exists_form, finrank_int and trace_eq) for the datum on A with xi = t xi0, the signs from X.xi0_im and
% X.kappa0_sign, and CM.oppDatum for L_-. The text's eta is the Lean's -xi0. The Lean states it for any FD.DataX X with
% hPsi and hLambda; the text states it for the data of the component Z. Size (tools/lean_size.py CMBasis.exists_pair):
% 552 lines.
\end{proof}

\begin{lemma}\label{cv-cusp-cm-point-exists-lattice-pure}
In Situation~\ref{cv-cusp-setting}, let $\sigma_0$, $Z$, $L_Z$, $(G,\eta,\varphi_Z)$ and $L'$ be as in
Lemma~\ref{cv-cusp-cm-point-exists-lattice-pair}, and let $\rho_0$ be the square root of $-d$ fixed by $\mathcal D$. Let
$x$ be a frame of $Z$, with period $P=P^Z_x$ on $L_Z$ (Situation~\ref{cv-uniformization-situation}), let $\ell>0$, and
let $v_1,\dots,v_r$ be a basis of $E^r$ in $L'$ with properties (1) to (4) of Lemma~\ref{cv-cusp-frame} for this $\ell$.
Let $c\in E$ and $1\le i\le r$ with $cv_i\in\varphi_Z(L_Z)$, and put $u=\varphi_Z^{-1}(cv_i)$. If $i>q$, then $P(u)$
vanishes at every coordinate whose embedding sends $\sqrt{-d}$ to $-\rho_0$; if $i\le q$, then $P(u)$ vanishes at every
coordinate whose embedding sends $\sqrt{-d}$ to $\rho_0$.
\end{lemma}
\leanref{Development/AbelianVarietyTheory/Proofs/CubicalLineBundles/InvMulVecEqZero.lean}{GlobalTheory.PurePer.perFr_pure}
\begin{proof}
\openstep{The proof is not yet written out: by property (4), $\tau(v_1),\dots,\tau(v_q)$ span the first $q$ columns of
$x_\tau$, so the coordinates of $x_\tau^{-1}\tau(cv_i)$ that the period formula reads vanish in the two cases.}
% Lean: PurePer.perFr_pure (for any a in E and j, in terms of perFr X.toData x (a lam j) tau i), with
% PurePer.inv_mulVec_eq_zero and emb_smul', through the span clause hsp of FD.LamBasis; the text reads it through W.pm_eq
% and the coordinates ix of the uniformization (hix: ix k < q off Phi), as in CMBasis.exists_cmLat. Size
% (tools/lean_size.py PurePer.perFr_pure): 113 lines.
\end{proof}

\begin{lemma}\label{cv-cusp-cm-point-exists-lattice}
In Situation~\ref{cv-cusp-setting}, let $\sigma_0$ be the inclusion of $F_c$ in $\mathbf C$, $\rho_0$ the square root
of $-d$ fixed by $\mathcal D$, and $Z$ a component at $\sigma_0$, with the lattice $L_Z$, its alternating form
$\langle\ ,\ \rangle$ and the Hermitian model $(G,\eta,\varphi_Z)$ that the compatibility conditions of the
uniformization attach to $Z$, as listed at the beginning of the proof of Lemma~\ref{cv-cover-uniformization}. Let
$L'\subset E^r$ be the lattice, self-dual for $G_Z$, that the uniformization provides for $Z$, and suppose that
$\varphi_Z(L_Z)$ is the $\mathcal O_E$-span of $\mathfrak AL'$ for a nonzero fractional ideal $\mathfrak A$ of $E$. Let
$x$ be a frame of $Z$, with period $P=P^Z_x$ on $L_Z$ (Situation~\ref{cv-uniformization-situation}), let $\ell>0$,
and let $v_1,\dots,v_r$ be a basis of $E^r$ in $L'$ with properties (1) to (4) of Lemma~\ref{cv-cusp-frame} for this
$\ell$. Then there are lattices $L_+$ and $L_-$ of PEL data over $E$ of signatures $(1,0)$ and $(0,1)$, with forms
$\langle\ ,\ \rangle_\pm$, $\mathcal O_E$-linear maps $f_j\colon L_+\to L_Z$ for $1\le j\le p$ and
$f'_j\colon L_-\to L_Z$ for $1\le j\le q$, elements $h_j$ and $h'_j$ of $\mathcal O_E$, and an integer $n\ge1$, such
that:
\begin{enumerate}
\item for $a\in L_+$ the period $P(f_ja)$ vanishes at every coordinate whose embedding sends $\sqrt{-d}$ to $-\rho_0$,
  and for $a\in L_-$ the period $P(f'_ja)$ vanishes at every coordinate whose embedding sends $\sqrt{-d}$ to $\rho_0$;
\item $\langle f_ja,f_jb\rangle=\langle h_ja,b\rangle_+$ for $a,b\in L_+$, and
  $\langle f'_ja,f'_jb\rangle=\langle h'_ja,b\rangle_-$ for $a,b\in L_-$;
\item for every $a\in L_Z$, $na$ is $\sum_jf_j(b_j)+\sum_jf'_j(b'_j)$ for some $b_j\in L_+$ and $b'_j\in L_-$;
\item $|\tau(h_j)|\le\ell^2$ and $|\tau(h'_j)|\le\ell^2$ for all $j$ and every embedding $\tau\colon E\to\mathbf C$.
\end{enumerate}
\end{lemma}
\leanref{Development/AbelianVarietyTheory/Proofs/CubicalLineBundles/ExistsCmPointLatticeDataOfAdaptedBasisStatement.lean}{GlobalTheory.T.exists_cmPoint_lattice_data_of_adapted_basis}
\leanref{Development/AbelianVarietyTheory/Proofs/CubicalLineBundles/ExistsCmPointLatticeDataOfAdaptedBasis.lean}{GlobalTheory.LamT.CMBasis.exists_cmLat_unfAt}
\begin{proof}
Write $G_Z(u,u')=\sum_{a,b}u_aG_{ab}\bar u'_b$ for the Hermitian form of the model, so that
$\langle a,b\rangle=\operatorname{Tr}_{E/\mathbf Q}(\eta\,G_Z(\varphi_Za,\varphi_Zb))$, and $\Lambda=\varphi_Z(L_Z)$. As
$L'$ is self-dual for $G_Z$, the values $G_Z(u,u')$ with $u,u'\in L'$ lie in $\mathcal O_E$, and $G_Z(u,u)$ is fixed by
complex conjugation.

\emph{Rank one lattices.} Lemma~\ref{cv-cusp-cm-point-exists-lattice-pair} gives lattices $L_+$ and $L_-$ of PEL data
over $E$, of signatures $(1,0)$ and $(0,1)$, with forms $\langle\ ,\ \rangle_\pm$, $\mathcal O_E$-linear isomorphisms
$e_\pm\colon L_\pm\to\mathfrak A$ and signs $s_\pm\in\{1,-1\}$ with
$\langle a,b\rangle_\pm=\operatorname{Tr}_{E/\mathbf Q}(-s_\pm\eta\,e_\pm(a)\overline{e_\pm(b)})$.

\emph{The maps.} For $c\in\mathfrak A$ and $1\le i\le r$, $cv_i$ lies in $\mathfrak AL'\subset\Lambda$. Number the
$p$ indices $i>q$ as $\sigma(1),\dots,\sigma(p)$ and the $q$ indices $i\le q$ as $\sigma'(1),\dots,\sigma'(q)$, and put
\[
 f_j(a)=\varphi_Z^{-1}\big(e_+(a)\,v_{\sigma(j)}\big)\quad(a\in L_+),\qquad
 f'_j(a)=\varphi_Z^{-1}\big(e_-(a)\,v_{\sigma'(j)}\big)\quad(a\in L_-),
\]
which are $\mathcal O_E$-linear. Property (1) of the statement is
Lemma~\ref{cv-cusp-cm-point-exists-lattice-pure}, for the indices $\sigma(j)>q$ and $\sigma'(j)\le q$.

\emph{The multipliers.} Put $h_j=-s_+G_Z(v_{\sigma(j)},v_{\sigma(j)})$ and $h'_j=-s_-G_Z(v_{\sigma'(j)},v_{\sigma'(j)})$,
elements of $\mathcal O_E$. For $a,b\in L_+$ and $v=v_{\sigma(j)}$, $G_Z(e_+(a)v,e_+(b)v)=e_+(a)\overline{e_+(b)}G_Z(v,v)$,
so, as $s_+^2=1$,
\[
 \langle f_ja,f_jb\rangle=\operatorname{Tr}_{E/\mathbf Q}\big(\eta\,G_Z(v,v)\,e_+(a)\overline{e_+(b)}\big)
 =\operatorname{Tr}_{E/\mathbf Q}\big(-s_+\eta\,e_+(h_ja)\overline{e_+(b)}\big)=\langle h_ja,b\rangle_+,
\]
and in the same way for $L_-$; this is (2). For $\tau\in\Phi$, $|\tau(h_j)|=|\tau(G_Z(v,v))|\le\ell^2$ by property (3)
of the basis; for $\tau\notin\Phi$, $\bar\tau\in\Phi$, and $|\tau(G_Z(v,v))|=|\bar\tau(\overline{G_Z(v,v)})|
=|\bar\tau(G_Z(v,v))|\le\ell^2$. This is (4).

\emph{The integer $n$.} Let $S$ be a finite set of generators of the $\mathcal O_E$-module $L_Z$. The coordinates of the
$\varphi_Z(a)$, $a\in S$, in the basis $v_1,\dots,v_r$ of $E^r$ are finitely many elements of $E$. As $\mathfrak A\neq0$,
there is an integer $n\ge1$ with $n$ times each of them in $\mathfrak A$: choose $0\neq c\in\mathfrak A$; the quotients
by $c$ have a common denominator $n$ in $\mathbf Z$, and $nc'=(nc'/c)c\in\mathcal O_E\mathfrak A=\mathfrak A$. So for
$a\in S$, $n\varphi_Z(a)=\sum_ic_iv_i$ with $c_i\in\mathfrak A$, and
$na=\sum_jf_j(e_+^{-1}(c_{\sigma(j)}))+\sum_jf'_j(e_-^{-1}(c_{\sigma'(j)}))$. The elements $a$ for which $na$ has such a
form make up an $\mathcal O_E$-submodule, so it is all of $L_Z$. This is (3).
\end{proof}
% Lean: T.exists_cmPoint_lattice_data_of_adapted_basis = LamT.CMBasis.exists_cmLat_unfAt, which applies
% CMBasis.exists_cmLat to the data of the uniformization at x (W.per, W.phi k, W.isom k, W.pm_eq); the conclusion is a
% CMLat W.per p q (fields D1, sigP, sigM, f, pure, hh, form, span) with |tau(hh j)| <= La^2. exists_cmLat: exists_pair
% (Lemma -lattice-pair), perFr_pure (Lemma -lattice-pure) through the coordinates ix of pm, the maps m j a = a lam j and
% f j = phi^-1 o m (sigma j) o e with sigma = sumComm then finSumFinEquiv (the L_+ indices at positions >= q), hh j =
% -s * herm Gm (lam (sigma j)) (lam (sigma j)) with h3 for the form, CMBasis.exists_nsmul_mem and
% basisOfLinearIndependentOfCardEqFinrank for the span, and the conjugate embedding for the bound off Psi. The text's eta
% is the Lean's -Y.xi0. Size of the remaining open steps: Lemmas -lattice-pair (552 lines) and -lattice-pure (113 lines).

\begin{lemma}\label{cv-cusp-cm-point-exists-object}
In Situation~\ref{cv-function-situation}, let $L$ be the lattice of the PEL datum, of signature $(p,q)$, and $\theta$ the
square root of $-d$ of the datum. Let $l$ be a level of the datum and $X'$ its moduli scheme at level $l$, over the field
$K_l\subset\mathbf C$ (Definition~\ref{cv-model-toroidal-semiabelian}(1)), and let $\nu\mapsto\sigma_\nu$, $1\le\nu\le g$,
be a map to the embeddings $E\to\mathbf C$. Let a family uniformization of $X'$ at the inclusion of $K_l$ in $\mathbf C$ for
this map be given: lattices $L_i$, maps $t_i$, points $y_i(P)$ and homomorphisms $\exp_{i,P}$ with the properties of
Theorem~\ref{cv-uniformization-family} for this map in place of a type, with periods as in
Definition~\ref{cv-cm-point-periods}(3) for it. Let $P$ be a period of one of the lattices $L_i$, and let $L_\pm$, $f_j$,
$f'_j$, $h_j$, $h'_j$ and $n$ satisfy (1) to (3) of Lemma~\ref{cv-cusp-cm-point-exists-lattice} for $L_i$ with its form,
$P$ and $\theta$ in place of $L_Z$ with its form, $P$ and $\rho_0$. Let $k\subset\mathbf C$ be a number field. Then there
are a number field $K_1\subset\mathbf C$ containing $K_l$ and $k$, an element $\rho\in\mathcal O_{K_1}$ equal to $\theta$,
and the following data.
\begin{enumerate}
\item[(a)] Over $\mathcal O_{K_1}$, principally polarized abelian schemes $A_+$, $A_-$ and $B$ with actions of
  $\mathcal O_E$, of types $(1,0)$, $(0,1)$ and $(p,q)$ at $\rho$, and $\mathcal O_E$-linear homomorphisms
  $\varphi_j\colon A_+\to B$ for $1\le j\le p$ and $\varphi'_j\colon A_-\to B$ for $1\le j\le q$, such that $E$ has
  $[E:\mathbf Q]$ embeddings into $K_1$, $\rho^2=-d$, every point of $B$ over $\mathbf C$ is a sum of images of points of
  $A_+$ and $A_-$ over $\mathbf C$ under the $\varphi_j$ and $\varphi'_j$, and over $\mathbf C$ the commutator pairings of
  $B$ pulled back along $\varphi_j$ and $\varphi'_j$ are those of $A_+$ composed with the action of $h_j$ and of $A_-$
  composed with the action of $h'_j$ (Lemma~\ref{cv-cusp-cm-point-eigen-degree-rank-one-norm-weil}).
\item[(b)] A point $o$ of $X'$ with values in $K_1$, over $K_l$, whose image in $X'(\mathbf C)$ is $y_i(P)$, and an
  identification of the abelian variety of $o$ with the generic fibre of $B$ that respects the actions of
  $\mathcal O_E$, the polarizations and the rigidified line bundles.
\item[(c)] Types of signatures $(1,0)$ and $(0,1)$ for $\theta$, of lengths $g_+$ and $g_-$, periods $P_\pm$ of $L_\pm$ of
  these types, and surjective holomorphic homomorphisms $e_\pm$ from $\mathbf C^{g_\pm}$ onto the points of $A_\pm$ over
  $\mathbf C$, with kernels $P_\pm(L_\pm)$, compatible with the actions of $\mathcal O_E$ through these types as in
  Theorem~\ref{cv-uniformization-family}(a), and with the Weil pairings given by the forms of $L_\pm$ as in (b) there;
  and $\mathbf C$-linear maps $F_j\colon\mathbf C^{g_+}\to\mathbf C^g$ and $F'_j\colon\mathbf C^{g_-}\to\mathbf C^g$ with
  $F_j(P_+(a))=P(f_ja)$ and $F'_j(P_-(a))=P(f'_ja)$, such that for every $z$ the point $\varphi_j(e_+(z))$ of $B$
  corresponds under (b) to the point $\exp_{i,P}(F_j(z))$ of $A_{y_i(P)}$, and $\varphi'_j(e_-(z))$ to
  $\exp_{i,P}(F'_j(z))$.
\end{enumerate}
\end{lemma}
\leanref{Development/AbelianVarietyTheory/Statements/PrincipalLattices/ExistsPrincipalLattice.lean}{GlobalTheory.CM.CMExistsSt}
\leanref{Development/AbelianVarietyTheory/Proofs/CubicalLineBundles/CmExistsHolds.lean}{GlobalTheory.CM.cmExists_holds}
\begin{proof}
\openstep{The construction of the integral CM object from the rank one lattices, over the integers of a number field,
and of the point of $X'$ under it is not yet written out.}
% Lean: CM.cmExists_holds = ExObj.cmExists, from cmCurve_closed (CM elliptic curves), cmPoint_closed (CM points are
% algebraic), thetaUnifAV_all, the laws firstFiveLaws, homExtension_semiAbelian_dvr_unconditional (homomorphisms extend
% over discrete valuation rings) and kernelFiniteFlat_holds. CMExistsSt is stated, as in the text, for any E, H, D, l,
% Xo, g, any ty : Fin g -> (E -> C), a family uniformization FU : FamilyUnif Xo (algebraMap (KA l) C) ty, any period Pe
% of FU.Di i with IsPeriod and any CMLat of it (the text's (1) to (3) of Lemma cv-cusp-cm-point-exists-lattice); the
% text's (a) is CMObj (split, rho_sq, A, B, typeP, typeM, typeB, hh = Z.hh, phi, lin, onto, weil), (b) is y with y >>
% Xo.s = spMap (inclusion hk) over KA l, spMap (K_1 -> C) >> y = FU.u i Pe, and t_0 with the SAB maps m to the universal
% triple and m' to X.B, and (c) is gA, tyA, PA, e, F with IsTypeFn, IsExp and the last clause. Only (a), (b), k <= K_1
% and rho = rt0 are used below. Size: about 22,700 lines below cmExists_holds beyond the statements the book cites.
\end{proof}

\begin{lemma}\label{cv-cusp-cm-point-exists-good}
In Situation~\ref{cv-cusp-setting}, let $K\supset F_c$ be a number field with an embedding $s\colon K\to\mathbf C$, let
$\rho\in\mathcal O_K$, let $\ell>0$, and let $A_+$, $A_-$, $B$, $\varphi_j$, $\varphi'_j$, $h_j$ and $h'_j$ over
$\mathcal O_K$ have the properties (a) of Lemma~\ref{cv-cusp-cm-point-exists-object} for $K$, $\rho$ and $s$, with
$|\tau(h_j)|\le\ell^2$ and $|\tau(h'_j)|\le\ell^2$ for all $j$ and $\tau$. Suppose that an embedding
$\tau\colon E\to\mathbf C$ lies in $\Phi$ if and only if $\tau_K(\sqrt{-d})=\rho$. Let $o$ be a point of $X$ with values in
$K$, over $F_c$, with a key datum on its abelian variety as in Definition~\ref{cv-cusp-cm-point-good}(2)(a), and with an
identification of its abelian variety with the generic fibre of $B$ that respects the actions of $\mathcal O_E$, the
polarizations and the rigidified line bundles. Then for every $F_c$-algebra $S$ and every homomorphism of $F_c$-algebras
$\sigma_1\colon K\to S$ there are a number field $K'\supset F_c$, a good CM point of $X$ over $K'$ with bound $\ell$
(Definition~\ref{cv-cusp-cm-point-good}), with point $o'$, and a homomorphism of $F_c$-algebras $\sigma'\colon K'\to S$
with $\sigma'(o')=\sigma_1(o)$.
\end{lemma}
\leanref{Development/AbelianVarietyTheory/Proofs/CubicalLineBundles/ExistsGoodCmPointOfIntegralCmObjectStatement.lean}{GlobalTheory.T.exists_good_cm_point_of_integral_cm_object}
\leanref{Development/AbelianVarietyTheory/Proofs/CubicalLineBundles/ExistsGoodCmPointOfIntegralCmObject.lean}{GlobalTheory.LamT.exists_cmPt_good_of_cmObj}
\begin{proof}
\openstep{The construction of the key data of the CM factors, of the lines, the letters and the framed point of a good CM
point from the integral CM object is not yet written out.}
% Lean: T.exists_good_cm_point_of_integral_cm_object = LamT.exists_cmPt_good_of_cmObj, for X : CMObj H K rho sigma p q,
% a FibKey Phi at x and a SABMap Phi.T.t -> X.B.t over Spec K -> Spec O_K; the hypotheses psiSet (the text's condition on
% Phi), |tau(hh j)| <= La^2 and x >> Xo.s = sOf c K. Size: about 28,000 lines below exists_cmPt_good_of_cmObj beyond the
% statements the book cites; it overlaps with Lemma cv-cusp-cm-point-exists-object.
\end{proof}

\begin{lemma}\label{cv-cusp-cm-point-exists}
In Situation~\ref{cv-cusp-setting}, let $\sigma_0$ be the inclusion of $F_c$ in $\mathbf C$, let $Z$ be a component
at $\sigma_0$, and let $L'\subset E^r$ be the lattice, self-dual for $G_Z$, that the uniformization at $\sigma_0$
provides for $Z$. Let $x$ be a frame of $Z$, let $w_x$ be the complex framed point of $X$ over $\sigma_0$
(Lemma~\ref{cv-cusp-cm-point-polynomial}) that the uniformization attaches to $x$, let $\ell>0$, and let
$v_1,\dots,v_r$ be a basis of $E^r$ in $L'$ with properties (1) to (4) of Lemma~\ref{cv-cusp-frame} for this $\ell$.
Then there are a number field $K\supset F_c$, a good CM point of $X$ over $K$ with bound $\ell$
(Definition~\ref{cv-cusp-cm-point-good}), and an embedding $s_1\colon K\to\mathbf C$ that extends $\sigma_0$ and
carries its point $o$ to the point of $w_x$.
\end{lemma}
\leanref{Development/AbelianVarietyTheory/Proofs/CubicalLineBundles/ExistsGoodCmPointAtCmBasisFrameStatement.lean}{GlobalTheory.T.exists_good_cmPoint_at_cmBasis_frame}
\leanref{Development/AbelianVarietyTheory/Proofs/CubicalLineBundles/ExistsGoodCmPointAtCmBasisFrame.lean}{GlobalTheory.LamT.exists_good_cmPt_at_frame}
\begin{proof}
Let $L_Z$, with its Hermitian model $(G,\eta,\varphi_Z)$, be the lattice that the compatibility conditions attach to
$Z$, and $P=P^Z_x$ the period of $x$ on it; the point of $w_x$ is the point that the uniformization attaches to $P$, as
listed at the beginning of the proof of Lemma~\ref{cv-cover-uniformization}. The lattice $L'$ that the uniformization
provides for $Z$ comes with a nonzero fractional ideal $\mathfrak A$ of $E$ for which $\varphi_Z(L_Z)$ is the
$\mathcal O_E$-span of $\mathfrak AL'$.

Step 1: lattice data. Lemma~\ref{cv-cusp-cm-point-exists-lattice} gives lattices $L_\pm$ of signatures $(1,0)$ and
$(0,1)$, maps $f_j$ and $f'_j$ into $L_Z$, multipliers $h_j$ and $h'_j$ with $|\tau(h_j)|,|\tau(h'_j)|\le\ell^2$, and
an integer $n$, with properties (1) to (4) of that lemma.

Step 2: the integral CM object. The compatibility conditions attach to $Z$ the lattice $L_Z=L_{i_Z}$ of a family
uniformization of $X$ at $\sigma_0$ restricted to $\mathbf Q(\sqrt{-d})$, which is the field $K_l$ of the level of $X$,
of full level $1$ (Theorem~\ref{cv-uniformization-family}); $P$ is a period of $L_Z$, and the point of $w_x$ is
$y_{i_Z}(P)$. Lemma~\ref{cv-cusp-cm-point-exists-object}, for the PEL datum and the level of $X$, this family
uniformization, the period $P$ of $L_Z$, the data of Step 1, which satisfy (1) to (3) of
Lemma~\ref{cv-cusp-cm-point-exists-lattice}, and $k=F_c$, gives a number field $K_1\subset\mathbf C$ containing $F_c$, an
element $\rho\in\mathcal O_{K_1}$ equal to $\rho_0$, abelian schemes $A_+$, $A_-$ and $B$ over $\mathcal O_{K_1}$ with the
homomorphisms $\varphi_j$ and $\varphi'_j$, which have the properties (a) of that lemma with the multipliers $h_j$ and
$h'_j$ of Step 1, and a point $o$ of $X$ with values in $K_1$, over $\mathbf Q(\sqrt{-d})$, whose image in $X(\mathbf C)$
is the point of $w_x$, with an identification of its abelian variety with the generic fibre of $B$. As
$F_c\subset K_1\subset\mathbf C$, $o$ is also a point of $X$ with values in the $F_c$-algebra $K_1$, over $F_c$.

Step 3: a key datum. As in Step 1 of the proof of Lemma~\ref{cv-cusp-cm-point-cusp-bound-ortho}, the abelian variety of
$o$, the fibre at $o$ of the universal abelian scheme, is a principally polarized abelian variety over $K_1$ of dimension
$g=fr\ge1$, with the symmetric line bundle of its polarization rigidified along the zero section. By
Theorem~\ref{cv-volume-key-sections}, for $S=\operatorname{Spec}K_1$ and $n=1$, the sections of this line bundle form a
$K_1$-vector space with a basis of one element, which is a nonzero section. With it and a basis of its Lie algebra, the
abelian variety of $o$ is a key datum over $K_1$. The identification of Step 2 is with the same fibre, so it is an identification of this key
datum's abelian variety with the generic fibre of $B$ that respects the actions of $\mathcal O_E$, the polarizations
and the rigidified line bundles.

Step 4: the good CM point. For an embedding $\tau\colon E\to\mathbf C$, $\tau_{K_1}(\sqrt{-d})$ is the complex number
$\tau(\sqrt{-d})$, as $K_1\subset\mathbf C$; it equals $\rho=\rho_0$ if and only if $\tau\in\Phi$. So
Lemma~\ref{cv-cusp-cm-point-exists-good} applies to $K=K_1$, $s$ and $\sigma_1$ the inclusion of $K_1$ in
$S=\mathbf C$, $\rho$, the data of Step 2, the bound $\ell$ and the point $o$ with the key datum of Step 3. It gives a
number field $K\supset F_c$, a good CM point of $X$ over $K$ with bound $\ell$, with point $o'$, and a homomorphism of
$F_c$-algebras $s_1\colon K\to\mathbf C$, which extends $\sigma_0$, with $s_1(o')$ the image of $o$ in
$X(\mathbf C)$, the point of $w_x$.
\end{proof}
\leanref{Development/AbelianVarietyTheory/Proofs/CubicalLineBundles/SabMapCancelAndCompOver.lean}{GlobalTheory.sabMap_cancel_and_comp_over}
% Lean: exists_good_cmPt_at_frame: PullbackSections.nonempty_unfAt (the uniformization data W at x0, with W.pt for the
% point of w_x); Step 1 is T.exists_cmPoint_lattice_data_of_adapted_basis (Lemma cv-cusp-cm-point-exists-lattice), for
% data Xx of the component with Xx.Lambda = AA L'; Step 2 is CM.cmExists_holds (Lemma cv-cusp-cm-point-exists-object)
% with k = F_c (hfd), keeping K_1, rho, X, y, t_0, m, m', hle, hrho, hhh, hy and hpt; Step 3 is LamT.nonempty_fibKey
% (with 1 <= p + q from hU.1 and 0 < r; its proof has been read, and the text follows it, with Theorem
% cv-volume-key-sections for n = 1 in place of principalPolarisation_exists_unique_section, from the law SecTwistSt,
% for the nonzero section; its 4,772 lines are base-change facts), sabMap_cancel_and_comp_over (the map n between the two pullbacks of the
% universal triple to y) and sabComp (n composed with m'); Step 4 is T.exists_good_cm_point_of_integral_cm_object
% (Lemma cv-cusp-cm-point-exists-good) with the psiSet condition from hrho, the bound from hhh and the lattice data,
% then he, hpt and W.pt. Size: about 35,000 lines below exists_good_cmPt_at_frame beyond the statements the book cites,
% nearly all in Lemmas cv-cusp-cm-point-exists-object (about 22,700) and -good (about 28,000), which overlap, and
% -lattice (about 910).

\begin{lemma}\label{cv-cusp-cm-point}
In Situation~\ref{cv-cusp-setting}, let $\kappa=(\alpha,n)$ be a dominant weight, let $\ell\ge1$ be a real
number, and let $f\in\Lambda_\kappa$ be nonzero with $\|f\|_{\sigma,\infty}>0$ for every embedding
$\sigma\colon F_c\to\mathbf C$. Let $\sigma_0$ be the inclusion of $F_c$ in $\mathbf C$, let $Z$ be a component at
$\sigma_0$, and let $L'\subset E^r$ be the lattice, self-dual for $G_Z$, that the uniformization at $\sigma_0$
provides for $Z$. Let $x$ be a frame of $Z$ with $f_{\sigma_0,Z}(x)\ne0$, and let $v_1,\dots,v_r$ be a basis of
$E^r$ in $L'$ with properties (1) to (4) of Lemma~\ref{cv-cusp-frame} for this $\ell$. Then there are a number
field $K\supset F_c$, a key datum of dimension $g=fr$ over $K$, metrized lines $\Theta$, $\Delta$ and
$\mathcal L_1,\dots,\mathcal L_g$ over $K$, a bijection $e$, a matrix $U\in GL_g(K)$ and a nonzero polynomial $G$
over $K$ that satisfy the hypotheses of Proposition~\ref{cv-cusp-cm-product} with the following choices: $B$ is
the set of embeddings $\tau\colon E\to\mathbf C$, with $m_\tau$ as in Situation~\ref{cv-cusp-setting};
$t=\log\Omega_E(\chi)$ (Definition~\ref{cv-log-omega}); $a_{\tau,i}=\alpha_{\tau,i-1}$; $n$ is the integer of
$\kappa$ and $r$ that of the setting; $\ell$ is as above; and $S_s=\|f\|_{s|_{F_c},\infty}$ for every embedding
$s\colon K\to\mathbf C$.
\end{lemma}
\leanref{Development/AbelianVarietyTheory/Proofs/CubicalLineBundles/ExistsGoodCmPointAtCmBasisFrameStatement.lean}{GlobalTheory.T.exists_good_cmPoint_at_cmBasis_frame}
\leanref{Development/AbelianVarietyTheory/Proofs/CubicalLineBundles/LocalLowerBoundCmPointFramePolynomialEveryPrimeStatement.lean}{GlobalTheory.T.localLowerBound_cmPoint_framePolynomial_every_prime}
\leanref{Development/AbelianVarietyTheory/Proofs/CubicalLineBundles/ArchimedeanHypothesisLowerBoundCmPointStatement.lean}{GlobalTheory.T.archimedeanHypothesis_lowerBound_cmPoint}
\leanref{Development/AbelianVarietyTheory/Proofs/CubicalLineBundles/LamDegStLogOm.lean}{GlobalTheory.LamF.lamDegSt_logOm}
\begin{proof}
Put $g=fr$ and $t=\log\Omega_E(\chi)$, which is real as $\chi$ is real, and $a_{\tau,i}=\alpha_{\tau,i-1}$.
The uniformization at $\sigma_0$ attaches to the frame $x$ a complex framed point $w_x$ of $X$ over
$\sigma_0$ (Lemma~\ref{cv-cusp-cm-point-polynomial}), whose point is the point of the frame $x$, and the reading of
$f$ at $x$ is $f_{\sigma_0,Z}(x)=f(w_x)$.

Step 1: the CM point. By Lemma~\ref{cv-cusp-cm-point-exists}, there are a number field $K\supset F_c$, a good CM
point of $X$ over $K$ with bound $\ell$ (Definition~\ref{cv-cusp-cm-point-good}), and an embedding
$s_1\colon K\to\mathbf C$ that extends $\sigma_0$ and carries its point $o$ to the point of $w_x$. Use the notation
of the definition.

Step 2: the polynomial. Let $G$ be a polynomial as in Lemma~\ref{cv-cusp-cm-point-polynomial}(1), for $f$ and the
framed point $y$ of the CM point. The point of $w_x$ is the image of $o$ under $s_1$, and $f(w_x)\ne0$, so
Lemma~\ref{cv-cusp-cm-point-polynomial}(3) gives $G\ne0$.

Step 3: condition (L). As $f\in\Lambda_\kappa$ is integral, Lemma~\ref{cv-cusp-cm-point-local} shows that for every
nonzero prime $\mathfrak p$ of $\mathcal O_K$, condition (L) of Proposition~\ref{cv-cusp-cm-product} holds for
$\Theta$, $\Delta$, the lines $\mathcal L_k$, $U$, $e$, the integer $n$ and $G$.

Step 4: condition (A). As $f\in S_\kappa$, Lemma~\ref{cv-cusp-cm-point-arch} shows that for every embedding
$s\colon K\to\mathbf C$, condition (A) of Proposition~\ref{cv-cusp-cm-product} holds at $s$, with the exponents $a$
and $S_s=\|f\|_{s|_{F_c},\infty}$.

Step 5: the degrees. By Lemma~\ref{cv-cusp-cm-point-degree}, $\widehat\deg\mathcal M^\Phi_i=\widehat\deg\mathcal M'_i=t/2$
for every $i$. By (f) of the definition, every $\mathcal L_k$ is one of these lines, so $\widehat\deg\mathcal L_k=t/2$.

Step 6: $\det U\ne0$, by Lemma~\ref{cv-cusp-cm-point-letters}, so $U\in GL_g(K)$.

Step 7: conclusion. The key datum $(A_o,\theta,v)$ has dimension $g=fr$, and by (e) of the definition $\Theta$ and
$\Delta$ are adapted to it at every nonzero prime and every embedding. For $B$ the set of embeddings
$\tau\colon E\to\mathbf C$, $e$ is a bijection from the pairs $(\tau,i)$ with $1\le i\le m_\tau$ onto
$\{1,\dots,g\}$, and $m_\tau\le r$. The integers $a_{\tau,i}$ and $n$ are nonnegative, $r\ge1$ and $\ell\ge1$, and
$S_s>0$ for every $s$ by hypothesis. With Steps 2 to 6, these are the hypotheses of
Proposition~\ref{cv-cusp-cm-product}.
\end{proof}
\leanref{Development/AbelianVarietyTheory/Proofs/CubicalLineBundles/MeanLogSupNormLowerBoundCmFrameNonvanishing.lean}{GlobalTheory.LamT.letters_lines}
% Lean: GlobalTheory.LamT.meanLogSup_lowerBound, before its call to GlobalTheory.cmPoint_lowerBound_free. Steps 1,
% 3, 4 and 5 are GlobalTheory.T.exists_good_cmPoint_at_cmBasis_frame,
% GlobalTheory.T.localLowerBound_cmPoint_framePolynomial_every_prime,
% GlobalTheory.T.archimedeanHypothesis_lowerBound_cmPoint and GlobalTheory.LamF.lamDegSt_logOm (Lemmas
% cv-cusp-cm-point-exists, cv-cusp-cm-point-local, cv-cusp-cm-point-arch and cv-cusp-cm-point-degree), with the data
% Xx of the component, Xx.Lambda = AA L' for a nonzero fractional ideal AA
% (GlobalTheory.T.exists_selfDual_data_transvection_principal_level), and LamBasis Xx La x0 lam for the basis. Step 2 is
% GlobalTheory.T.exists_isFramePolynomialOfRule_numberField and the second half of
% GlobalTheory.T.framePolynomial_equivariant_and_ne_zero, with w_x = (U c.F).un k x0 (hev: rdU is the value of F at
% it). Step 5 after the degrees is letters_lines; Step 6 is GlobalTheory.LamF.Ind.lamInd; Step 7 is the call to
% cmPoint_lowerBound_free with hg (g = [E:Q]/2 r from the field hgB and r = p + q) and hmr (m_tau <= r).

\subsection{Norm bounds at the cusp}\label{subsection-cv-cusp-norm-bounds}

\begin{lemma}\label{cv-cusp-inputs}
In Situation~\ref{cv-cusp-setting}, let $\sigma$ be a good embedding, and for every component $Z$ at
$\sigma$ let $x_Z$ be an orthonormal frame of $Z$ and $\mathcal F_Z$ a fundamental set for $\Gamma_Z$ at
$x_Z$.
\begin{enumerate}
\item (Local bound.) Let $Z$ be a component, let $w>0$, let $\lambda$ be a weight on $Z$ of spread at
  most $w$ with $\chi_\lambda$ trivial, and let $P$ be a cusp form of weight $\lambda$ on $Z$. Then every
  frame $y$ of $Z$ satisfies
  \[
   |P|(y)^2\le C^{\mathrm{loc}}_w\int_{\mathcal F_Z}|P|(x_Zn(z))^2\,d\mu_{x_Z}(z).
  \]
\item (Covolumes.) $\sum_Z\mu_{x_Z}(\mathcal F_Z)\le C^{\mathrm{mass}}$.
\item (CM points.) Let $\kappa$ be a dominant weight and let $f\in\Lambda_\kappa$ be nonzero. Then
  \[
   \frac1{[F_c:\mathbf Q]}\sum_{\sigma':F_c\to\mathbf C}\log\|f\|_{\sigma',\infty}\ge-\eta^{\sup}_\kappa.
  \]
\end{enumerate}
\end{lemma}
\leanref{Development/AbelianVarietyTheory/Proofs/DivisorProofs/WeightedFrameSumPositive.lean}{GlobalTheory.SupModel.comp_bound}
\leanref{Development/WeilBound/RepresentationTheoryGL/TubeMeasures/DisjointTubes/ConstantSupremum/Imports.lean}{MainStatements.Tam.Plug.comp}
\leanref{Development/AbelianVarietyTheory/Proofs/DivisorProofs/MeanLogSupNormLowerBoundPrincipalLevelStatement.lean}{GlobalTheory.T.meanLogSupNorm_lowerBound_principalLevel}
\begin{proof}
(1) If $|P|(y)=0$ there is nothing to prove. Otherwise $s=\sup_x|P|(x)$ is finite and positive. Choose a
frame $z_0$ of $Z$ with $|P|(z_0)>s/2$, so that $|P|\le2|P|(z_0)$ on all frames of $Z$; put
$\theta=\theta_w$ and $\rho=\rho_w$, and use the notation of Definition~\ref{cv-decay-majorant}.
\begin{enumerate}
\item[(a)] No short vectors at $z_0$. The uniformization provides a lattice $L'\subset E^r$, self-dual for
  $G_Z$ and stable under $\Gamma_Z$, and a real number $M'>0$ with $M'\xi=M\sqrt d$, such that for every
  nonzero $v\in L'$ with $G_Z(v,v)=0$ and every $b\in\mathcal O_E$ with $c(b)=b$ some $\gamma\in\Gamma_Z$ has
  $\tau(\gamma)=n_v(M'\tau(b))_\tau$ for every $\tau\in\Phi$; and $\chi_\lambda(\gamma)=1$, as $\chi_\lambda$ is
  trivial. With this $M'$, the number $\theta$ of Lemma~\ref{cv-decay-maximum} is $\theta_w$. By that lemma,
  every nonzero $v\in L'$ has $h_{z_0,\tau}(\tau(v))\ge\theta$ for some $\tau$.
\item[(b)] An orthonormal frame. For each $\tau$, the Gram--Schmidt process for the positive definite form
  $h_{z_0,\tau}$, applied to the columns of $z_{0,\tau}$ in order, gives an invertible upper triangular
  $b_\tau$ for which the columns of $z_{0,\tau}b_\tau$ are orthonormal for $h_{z_0,\tau}$; put $b=(b_\tau)\in B$
  and $x=z_0b$. The first $q$ columns of $x_\tau$ span $U^+_{z_0,\tau}$, on which $h_{z_0,\tau}$ is
  $\tau(G_Z)$, and the others are orthogonal to it for $h_{z_0,\tau}$, hence lie in $U^-_{z_0,\tau}$, the
  orthogonal of $U^+_{z_0,\tau}$ for $\tau(G_Z)$, on which $h_{z_0,\tau}$ is $-\tau(G_Z)$. So
  $x_\tau^{\mathsf T}\tau(G_Z)\bar x_\tau$ is diagonal with entries $1$ and then $-1$, and $x$ is orthonormal. Writing
  $c_a(u)$ for the coordinates of $u\in\mathbf C^r$ in the columns of $x_\tau$, orthonormality for
  $h_{z_0,\tau}$ gives $h_{z_0,\tau}(u)=\sum_a|c_a(u)|^2$, and, as $x$ is orthonormal, $h_{x,\tau}(u)$ is the
  same sum (proof of Lemma~\ref{cv-returns-majorant}(3)); so $h_{x,\tau}=h_{z_0,\tau}$. Also $|P|(x)=|P|(z_0)$,
  since $|P|$ is invariant under $B$. By (a), $\sum_\tau h_{x,\tau}(\tau(v))\ge\theta$ for every nonzero
  $v\in L'$.
\item[(c)] Sub-mean value. Let $B_\rho$ be the set of $z\in\mathbf C^N$ with $|z|<\rho$. By
  Lemma~\ref{cv-submean-chart}, $B_\rho\subset\Omega_x$, so $\varphi(z)=P(xn(z))$ is holomorphic on $B_\rho$.
  For a unit vector $e\in\mathbf C^N$ and $0<t<\rho$, the function $\zeta\mapsto\varphi(\zeta e)$ is holomorphic
  on $|\zeta|<\rho$, so $|\varphi(0)|^2\le\frac1{2\pi}\int_0^{2\pi}|\varphi(te^{i\vartheta}e)|^2\,d\vartheta$.
  Integrating over $e$ in the unit sphere and over $t$ in polar coordinates gives
  $|\varphi(0)|^2\le\operatorname{vol}(B_\rho)^{-1}\int_{B_\rho}|\varphi(z)|^2\,dz$ for the Lebesgue measure $dz$,
  and $\operatorname{vol}(B_\rho)=\pi^N\rho^{2N}/N!$. As $x$ is orthonormal, $\beta_\lambda(x)=1$ and
  $|\varphi(0)|=|P|(x)$. By Lemma~\ref{cv-submean-chart}, on $B_\rho$ we have
  $|\varphi(z)|^2=|P|(xn(z))^2\beta_\lambda(xn(z))^{-2}\le e^{1/20}|P|(xn(z))^2$ and $dz\le e^{1/10}d\mu_x(z)$.
  Since $e^{1/20+1/10}\le2$,
  \[
   |P|(x)^2\le\frac{2\,N!}{(\pi\rho^2)^N}\int_{B_\rho}|P|(xn(z))^2\,d\mu_x(z).
  \]
\item[(d)] Returns. Let $m$ be the number of $\gamma\in\Gamma_Z$ with $\gamma xn(z)\in xn(z')B$ for some
  $z,z'$ with $|z|,|z'|\le\rho$. Then $m\le(1+8f^2r^2|d_E|/\theta)^{2fr^2}$. Indeed, by (b) and
  Lemma~\ref{cv-returns-short-basis}, $L'$ contains a basis $w_1,\dots,w_r$ of $E^r$ with
  $\sum_\tau h_{x,\tau}(\tau(w_i))\le B_0=(2fr)^4|d_E|^2/(4\theta)$. Let $\gamma$ be counted by $m$, with
  $\gamma xn(z)=xn(z')b$. Since $r^2\rho=1/(200r(w+f))\le1/100$, Lemma~\ref{cv-returns-majorant} gives, for
  $u\in\mathbf C^r$,
  \[
   h_{x,\tau}(\tau(\gamma)u)\le2h_{xn(z'),\tau}(\tau(\gamma)u)=2h_{\gamma xn(z),\tau}(\tau(\gamma)u)
   =2h_{xn(z),\tau}(u)\le4h_{x,\tau}(u).
  \]
  So the vectors $\gamma w_i$ lie in $L'$ and satisfy $\sum_\tau h_{x,\tau}(\tau(\gamma w_i))\le4B_0$, and they
  determine $\gamma$, as the $w_i$ are a basis. By Lemma~\ref{cv-returns-count} with $R=4B_0$,
  $m\le(1+2\sqrt{4B_0/\theta})^{2fr\cdot r}$, and $2\sqrt{4B_0/\theta}=2(2fr)^2|d_E|/\theta=8f^2r^2|d_E|/\theta$.
\item[(e)] Unfolding. The function $|P|^2$ is measurable and invariant under $\Gamma_Z$ and $B$, as
  $\chi_\lambda$ is trivial. By Lemma~\ref{cv-returns-unfold} with $x_0=x_Z$, $\mathcal F=\mathcal F_Z$ and
  $A=B_\rho$, the integral in (c) is at most $m\int_{\mathcal F_Z}|P|(x_Zn(z))^2\,d\mu_{x_Z}(z)$.
\end{enumerate}
As $|P|(y)\le2|P|(x)$, (c), (d) and (e) give (1), since $4\cdot2N!(\pi\rho^2)^{-N}$ times the bound in (d)
is $C^{\mathrm{loc}}_w$.

(2) The uniformization at $\sigma$ attaches the components to double cosets. It provides a basis of $W_j$, with
Gram matrix $H\in M_r(E)$ for the form of $W_j$, such that $G_Z=G=-H^{\mathsf T}$ for every component $Z$. So
$V=E^r$ with the form $h(v,v')=\sum_{a,b}c(v_a)H_{ab}v'_b$ is $W_j$ written in this basis, and it satisfies the
hypotheses of Lemma~\ref{cv-tube-bound} with $(P,Q)=(p,q)$, the form of $G$ having signature $(q,p)$, and
$U(V)=U(W_j)$. The uniformization also provides elements $g_Z\in U(V)(\mathbf A_{F,f})$, one for
each component, in distinct double cosets $U(V)(F)g_ZK(M)$; and for each $Z$ a family
$u_Z=(u_{Z,\tau})_{\tau\in\Phi}$ of matrices preserving the forms $\tau(G)$, such that for every
$\gamma\in\Gamma_Z$ there is $\gamma'\in\Gamma_{g_Z}$ with $u_{Z,\tau}\tau(\gamma)=\tau(\gamma')u_{Z,\tau}$ for every
$\tau$, and conversely. Here $K(M)$ and $\Gamma_{g_Z}$ are as in Lemma~\ref{cv-tube-bound}, and the pairs
$(G,\Gamma_{g_Z})$ have the properties of the pairs $(G_Z,\Gamma_Z)$.
\begin{enumerate}
\item[(a)] Conjugation. For each $Z$ choose an orthonormal frame $x'_Z$ for $G$ and a fundamental set
  $\mathcal F'_Z$ for $\Gamma_{g_Z}$ at $x'_Z$, in the sense of Lemma~\ref{cv-tube-chart}. Then
  $\mu_{x_Z}(\mathcal F_Z)=\mu_{x'_Z}(\mathcal F'_Z)$. Indeed, $y\mapsto u_Zy=(u_{Z,\tau}y_\tau)_\tau$ maps the frames of
  $Z$ onto the frames for $G$, keeps the minors $\delta_{\tau,i}$, commutes with the right action of $B$ and
  with the charts, and carries $\Gamma_Z$-orbits to $\Gamma_{g_Z}$-orbits. So $\mathcal F_Z$ is a fundamental set
  for $\Gamma_{g_Z}$ at the orthonormal frame $u_Zx_Z$, and $\mu_{u_Zx_Z}=\mu_{x_Z}$. Two fundamental sets for
  $\Gamma_{g_Z}$ at two orthonormal frames have the same measure, by the two facts in the proof of
  Lemma~\ref{cv-returns-unfold}.
\item[(b)] Covolumes. As $M\ge3$, the elements of $\Gamma_{g_Z}$ have determinant $1$ by
  Lemma~\ref{cv-tube-det}. By Lemma~\ref{cv-tube-chart},
  $c_r^f\mu_{x'_Z}(\mathcal F'_Z)=\operatorname{vol}(\Gamma_{g_Z}\backslash SU_\Phi(G),\nu_{x'_Z})$ with
  $c_r=2^{r(r-1)/2}(2\pi)^{r-1}\sqrt r\ge1$.
\item[(c)] The sum. By (a), (b) and Lemma~\ref{cv-tube-bound},
  \[
   \sum_Z\mu_{x_Z}(\mathcal F_Z)\le\sum_Z\operatorname{vol}(\Gamma_{g_Z}\backslash SU_\Phi(G),\nu_{x'_Z})
   \le16^{fr^2}M^{2fr^2}\operatorname{vol}U(r)^f\operatorname{mass}(I_r),
  \]
  with $\operatorname{vol}U(r)$ as in Theorem~\ref{cv-volume-mass} and $I_r$ the standard lattice
  (Definition~\ref{cv-mass}).
\end{enumerate}
As in the proof of Proposition~\ref{cv-level-volume},
$\operatorname{vol}U(r)\le13^{r^2}$, and Lemma~\ref{cv-mass-bounds} with $[E:\mathbf Q]=2f$ gives
$\operatorname{mass}(I_r)\le(12f^2r^4\operatorname{rd}_E^2)^{6fr^2}$. Since $16\cdot13\le15^2$, the right
side is at most $(15M)^{2fr^2}(12f^2r^4\operatorname{rd}_E^2)^{6fr^2}=C^{\mathrm{mass}}$.

(3) Let $\sigma_0$ be the inclusion of $F_c$ in $\mathbf C$, which is good, and put $w=w_\kappa$ and
$\ell=\ell_w$. As $f\ge1$, $r\ge1$, $w\ge1$, $|d_E|\ge1$ and $\theta_w\le1$, every factor in the definition of
$\ell_w$ is at least $1$, so $\ell\ge1$.
\begin{enumerate}
\item[(a)] Near a maximum. Since $f\ne0$, its reading at $\sigma_0$ is nonzero, so $\sup|f_{\sigma_0}|$ is
  positive, and it is finite. Choose a component $Z$ and a frame $z_0$ of $Z$ with
  $|f_{\sigma_0,Z}|\le2|f_{\sigma_0,Z}|(z_0)$ on all frames of $Z$.
\item[(b)] A CM frame. With $L'$ as in (1)(a), there are a frame $x$ of $Z$ with $f_{\sigma_0,Z}(x)\ne0$
  and a basis $v_1,\dots,v_r$ of $E^r$ in $L'$ such that $G_Z(v_a,v_b)=0$ for $a\le q<b$,
  $G_Z(v_a,v_a)$ is totally positive for $a\le q$ and totally negative for $a>q$,
  $|\tau(G_Z(v_a,v_a))|\le\ell^2$ for all $a$ and $\tau$, and for every $\tau$ the first $q$ columns of
  $x_\tau$ span the same space as $\tau(v_1),\dots,\tau(v_q)$. Indeed, the weight $\lambda(\kappa)$ has spread
  less than $w$ and trivial characters, by Step~1 of the proof of Proposition~\ref{cv-cusp-rank-bound}, and
  $w\ge1$. By Lemma~\ref{cv-decay-maximum} at $z_0$, with this $w$ and the data $L'$, $M'$ of (1)(a), every
  nonzero $v\in L'$ has $h_{z_0,\tau}(\tau(v))\ge\theta_w$ for some $\tau$ (notation of
  Definition~\ref{cv-decay-majorant}), so $\sum_\tau h_{z_0,\tau}(\tau(v))\ge\theta_w$. Lemma~\ref{cv-cusp-frame},
  applied to $P=f_{\sigma_0,Z}$ with $\lambda_*=\theta_w^{1/2}\le1$, for which its constant is $\ell_w$, gives
  $x$ and $v_1,\dots,v_r$; the values $G_Z(v_a,v_a)$ lie in $F$, whose real embeddings are the restrictions of
  the $\tau\in\Phi$.
\item[(c)] The value at the CM point. Since $f\ne0$, its reading at every good embedding is nonzero, so
  $\|f\|_{\sigma,\infty}>0$ for every embedding $\sigma\colon F_c\to\mathbf C$. By
  Lemma~\ref{cv-cusp-cm-point}, applied to $\ell$ and to $x$ and $v_1,\dots,v_r$ of (b), there are a number field
  $K\supset F_c$ and data over $K$ that satisfy the hypotheses of Proposition~\ref{cv-cusp-cm-product} with
  $g=fr$, $t=\log\Omega_E(\chi)$, $|a|=|\alpha|$, the integers $n$ and $r\ge1$, the number $\ell\ge1$ and
  $S_s=\|f\|_{s|_{F_c},\infty}$ for the embeddings $s\colon K\to\mathbf C$. The proposition gives
  \[
   -\Big(\frac{1+|\log\Omega_E(\chi)|}2+\frac12\log\ell\Big)(|\alpha|+frn)-|\alpha|\log(r\ell)
   \le\frac1{[K:\mathbf Q]}\sum_{s\colon K\to\mathbf C}\log\|f\|_{s|_{F_c},\infty}.
  \]
  The left side is $-\eta^{\sup}_\kappa$. Every embedding of $F_c$ has $[K:F_c]$ extensions to $K$, so the
  right side is the mean of $\log\|f\|_{\sigma,\infty}$ over the embeddings $\sigma$ of $F_c$.
\end{enumerate}
% Lean: (3) GlobalTheory.lamSup_of_laws: half_max, then (b) is Formalism.L2.Pj.abelian₂.cm_frame_of_eq
% (Lemma cv-cusp-frame), and (c) is GlobalTheory.T.meanLogSupNorm_lowerBound_cmFrame_nonvanishing =
% GlobalTheory.LamT.meanLogSup_lowerBound: the data of Lemma cv-cusp-cm-point, then
% GlobalTheory.cmPoint_lowerBound_free (Proposition cv-cusp-cm-product) and
% Formalism.L2.Pj.abelian₂.mean_restrict. The positivity of the norms is hSpos in lamSup_of_laws, from
% GlobalTheory.supNorm_definite_on_cusp; the hypothesis on Xx.Lambda and the ideal AA come from
% GlobalTheory.T.exists_selfDual_data_transvection_principal_level.
\end{proof}
\leanref{Development/AbelianVarietyTheory/Proofs/DivisorProofs/CuspFormSupNormLePeterssonIntegralComponentStatement.lean}{GlobalTheory.T.cuspForm_sup_norm_le_petersson_integral_component}
\leanref{ReciprocityTheory/ReductiveGroups/ArithmeticGeometry/ArakelovVolumes/UpperBounds/ModeHalfStatementPfC.lean}{ReciprocityTheory.ChartRestriction.G.AR.Up.ModeHalfC.orthonormalize}
\leanref{Development/AbelianVarietyTheory/Proofs/DivisorProofs/SubMeanValueFlagChartInvariantMeasureStatement.lean}{GlobalTheory.T.subMeanValue_flagChart_invariantMeasure}
\leanref{Development/AbelianVarietyTheory/Proofs/CovolumeBounds/EnormPowLeInvMeasureMulLintegralBall.lean}{GlobalTheory.enorm_pow_le_inv_measure_mul_lintegral_ball}
\leanref{Development/AbelianVarietyTheory/Proofs/DivisorProofs/CardGroupElementsChartReturnLeStatement.lean}{GlobalTheory.T.card_group_elements_chart_return_le}
\leanref{Development/AbelianVarietyTheory/Proofs/DivisorProofs/CardGroupElementsChartReturnLe.lean}{GlobalTheory.ChartMult.maj_le_four}
\leanref{Development/WeilBound/RepresentationTheoryGL/TubeMeasures/DisjointTubes/MeasureSumBounds/SumMuFLe.lean}{MainStatements.Tam.C.sum_muF_le}
\leanref{Development/AbelianVarietyTheory/Proofs/DivisorProofs/CovolumeEqOfRealUnitaryConjStatement.lean}{GlobalTheory.T.covolume_eq_of_realUnitary_conj}
\leanref{Formalism/LocalToGlobal/ProjectiveLattices/RosatiLatticeNorms/CmFrameOfEq.lean}{Formalism.L2.Pj.abelian₂.cm_frame_of_eq}
\leanref{Development/AbelianVarietyTheory/Proofs/CubicalLineBundles/MeanLogSupNormLowerBoundCmFrameNonvanishingStatement.lean}{GlobalTheory.T.meanLogSupNorm_lowerBound_cmFrame_nonvanishing}
\leanref{Development/AbelianVarietyTheory/Proofs/CubicalLineBundles/MeanLogSupNormLowerBoundCmFrameNonvanishing.lean}{GlobalTheory.LamT.meanLogSup_lowerBound}
\leanref{Formalism/LocalToGlobal/ProjectiveLattices/RosatiLatticeNorms/LinkNum.lean}{Formalism.L2.Pj.abelian₂.mean_restrict}

\begin{definition}\label{cv-cusp-torus-bound}
The \emph{torus class bound} is the following statement. Let $E'$ be a CM field with totally real
subfield $F'$, unramified over $F'$ at every finite prime, with class number $h_{E'}$, and let $T$ be the
group of $t\in E'\otimes_{F'}\mathbf A_{F',f}$ with $t\bar t=1$. Let $r,M\ge1$, let $G\in M_r(E')$ be an
invertible Hermitian matrix and $\Lambda\subset E'^r$ a finitely generated $\mathcal O_{E'}$-module that
spans $E'^r$ and is self-dual for $h_G$, and let $K_\Lambda(M)$ be the group of unitary
$k\in GL_r(E'\otimes_{F'}\mathbf A_{F',f})$ with $k\widehat\Lambda=\widehat\Lambda$ and
$(k-1)\widehat\Lambda\subset M\widehat\Lambda$. Then the number of classes of $T$ modulo multiplication by
the elements $u\in E'$ with $u\bar u=1$ and by the determinants of the elements of $K_\Lambda(M)$ is at most
$h_{E'}M^{[E':\mathbf Q]}$.
\end{definition}
\leanref{Development/AbelianVarietyTheory/Statements/PlaceShiftStatements/DefiniteBound.lean}{GlobalTheory.UCls.TorusBound}
\leanref{Development/AbelianVarietyTheory/Definitions/Divisors/IsKMOne.lean}{GlobalTheory.UCls.torusNo}

\begin{lemma}\label{cv-cusp-torus-hilbert90}
Let $E'$ be a CM field with totally real subfield $F'$, let $\mathbf A_{E',f}=E'\otimes_{F'}\mathbf A_{F',f}$,
and write $z\mapsto\bar z$ for the complex conjugation of $E'$ and for its action on $\mathbf A_{E',f}$
through the first factor. Let $t\in\mathbf A_{E',f}$ with $t\bar t=1$. Then there is a unit $x$ of
$\mathbf A_{E',f}$ with $x=t\bar x$.
\end{lemma}
\leanref{Development/AbelianVarietyTheory/Proofs/DivisorProofs/FiniteAdeleQuadraticExistsUnitEqMulSg.lean}{GlobalTheory.UCls.NormOne.finiteAdele_quadratic_exists_unit_eq_mul_sg}
\begin{proof}
For $z\in\mathbf A_{E',f}$ the element $z+\bar z$ lies in $\mathbf A_{F',f}$, as it does for $z=y\otimes a$
with $y\in E'$ and $a\in\mathbf A_{F',f}$. So $s=(t+\bar t)/2$ lies in $\mathbf A_{F',f}$. At a finite place
$v$ of $F'$, $(1+s_v)+(1-s_v)=2$, so by the ultrametric inequality we may choose $e_v\in\{0,1\}$ with
$|e_v(1+s_v)+(1-e_v)(1-s_v)|_v\ge|2|_v$. Then $e=(e_v)_v$ is an idempotent of $\mathbf A_{F',f}$, and
$a=e(1+s)+(1-e)(1-s)$ is a unit of $\mathbf A_{F',f}$: it is nonzero at every place, and it is a unit at the
almost all places where $|2|_v=1$ and $a_v$ is integral. Choose $y\in E'$ with $\bar y\ne y$, and put
$d=y-\bar y$, a unit with $\bar d=-d$. Let
\[
 x=e(1+t)+(1-e)d(1-t).
\]
Since $\bar e=e$, $\bar x=e(1+\bar t)-(1-e)d(1-\bar t)$. Using $e^2=e$, $e(1-e)=0$, $t\bar t=1$ and
$t+\bar t=2s$,
\[
 x\bar x=2\big(e(1+s)-(1-e)d^2(1-s)\big)=2a\big(e-(1-e)d^2\big),
\]
and $e-(1-e)d^2$ is a unit, with inverse $e-(1-e)d^{-2}$. So $x$ is a unit. Finally
$t\bar x=e(t+t\bar t)-(1-e)d(t-t\bar t)=e(1+t)+(1-e)d(1-t)=x$.
\end{proof}
\leanref{Development/AbelianVarietyTheory/Proofs/UnitaryExistence/ExistsIdempotentIsUnitOneAddOrSubAdele.lean}{GlobalTheory.AdIdem.finite_part}
\leanref{Development/AbelianVarietyTheory/Proofs/UnitaryExistence/Hilbert90AdeleRingTensorQuadratic.lean}{GlobalTheory.AdLat.hilbert90_involution_of_idempotent_unit}

\begin{lemma}\label{cv-cusp-torus-ideal-classes}
Let $E'$ be a CM field with totally real subfield $F'$ and class number $h_{E'}$. Let
$\mathbf A_{E',f}=E'\otimes_{F'}\mathbf A_{F',f}$, with the complex conjugation $z\mapsto\bar z$ acting through
the first factor, let $\widehat{\mathcal O}_{E'}\subset\mathbf A_{E',f}$ be the subring of the elements integral
at every finite place, and let $T$ be the group of $t\in\mathbf A_{E',f}$ with $t\bar t=1$. Then there is a set
$S\subset T$ with at most $h_{E'}$ elements such that every $t\in T$ is $usk$ with $u\in E'$, $u\bar u=1$,
$s\in S$ and $k\in T\cap\widehat{\mathcal O}_{E'}$.
\end{lemma}
\leanref{Development/AbelianVarietyTheory/Proofs/DivisorProofs/FiniteAdeleQuadraticExistsUnitEqMulSg.lean}{GlobalTheory.UCls.normOne_classes_le_classNumber}
\leanref{Development/AbelianVarietyTheory/Proofs/DivisorProofs/FiniteAdeleQuadraticExistsUnitEqMulSg.lean}{GlobalTheory.UCls.NormOne.exists_class}
\begin{proof}
Step 1: a factorization. Let $t\in T$. By Lemma~\ref{cv-cusp-torus-hilbert90} there is a unit $x$ of
$\mathbf A_{E',f}$ with $x=t\bar x$. Choose an integer $d\ge1$ with $z=dx\in\widehat{\mathcal O}_{E'}$, and
put $w=z^{-1}$. As $\bar d=d$, $z=t\bar z$, so $t=z\bar w$.

Step 2: an ideal. Let $I$ be the set of $a\in\mathcal O_{E'}$ with $wa\in\widehat{\mathcal O}_{E'}$, an ideal of
$\mathcal O_{E'}$. It contains an integer $e\ge1$, chosen with $ew\in\widehat{\mathcal O}_{E'}$, so it is
nonzero. Write $e^{-1}z=y+h$ with $y\in E'$ and $h\in\widehat{\mathcal O}_{E'}$, as every element of
$\mathbf A_{E',f}$ is such a sum. Then $ey=z-eh$ lies in
$E'\cap\widehat{\mathcal O}_{E'}=\mathcal O_{E'}$, and in $I$, as $w\cdot ey=1-weh$. So $z=ey+eh$ with $ey$ and $e$
in $I$ and $h$ integral.

Make the choices of Steps 1 and 2 for every $t\in T$, giving $z_t$, $w_t$ and $I_t$, and let $\phi(t)$ be the
class of $I_t$ in the class group of $E'$.

Step 3: equal classes. Let $t,t'\in T$ with $\phi(t)=\phi(t')$. There are nonzero $a,b\in\mathcal O_{E'}$ with
$aI_t=bI_{t'}$. Write $z_t=x_1+eh$ with $x_1,e\in I_t$ and $h$ integral, as in Step 2, and write $ax_1=bx_1'$
and $ae=be'$ with $x_1',e'\in I_{t'}$. Then $az_tw_{t'}=b(x_1'w_{t'}+e'w_{t'}h)=bk$ with $k$ integral, by the
definition of $I_{t'}$. In the same way $bz_{t'}w_t=ak'$ with $k'$ integral. As $z_tw_t=z_{t'}w_{t'}=1$, this
gives $z_t=(b/a)z_{t'}k$, $w_t=(a/b)w_{t'}k'$ and $kk'=1$. Hence, by Step 1,
\[
 t=z_t\bar w_t=u\,t'\,k\bar k',\qquad u=\frac ba\,\overline{\Big(\frac ab\Big)},
\]
where $u\bar u=1$, $k\bar k'$ is integral, and $(k\bar k')\overline{(k\bar k')}=kk'\,\overline{kk'}=1$.

Step 4: the set $S$. For every class in the image of $\phi$ choose one $t'\in T$ with that class, and let $S$
be the set of these. It has at most $h_{E'}$ elements. Every $t\in T$ is $usk$ as claimed by Step 3, with
$s$ the element of $S$ with $\phi(s)=\phi(t)$.
\end{proof}
\leanref{Development/AbelianVarietyTheory/Proofs/DivisorProofs/FiniteAdeleQuadraticExistsUnitEqMulSg.lean}{GlobalTheory.UCls.NormOne.exists_ideal}
\leanref{Development/AbelianVarietyTheory/Proofs/DivisorProofs/FiniteAdeleQuadraticExistsUnitEqMulSg.lean}{GlobalTheory.UCls.NormOne.exists_quot}

\begin{lemma}\label{cv-cusp-torus-congruence}
Let $E'$ be a CM field with totally real subfield $F'$. Let $\mathbf A_{E',f}=E'\otimes_{F'}\mathbf A_{F',f}$,
with the complex conjugation $z\mapsto\bar z$ acting through the first factor, let
$\widehat{\mathcal O}_{E'}\subset\mathbf A_{E',f}$ be the subring of the elements integral at every finite
place, and let $T$ be the group of $t\in\mathbf A_{E',f}$ with $t\bar t=1$. Let $M\ge1$. Then there is a set
$S'\subset T$ with at most $M^{[E':\mathbf Q]}$ elements such that every $k\in T\cap\widehat{\mathcal O}_{E'}$ is
$sk'$ with $s\in S'$, $k'\in T$ and $k'-1\in M\widehat{\mathcal O}_{E'}$.
\end{lemma}
\leanref{Development/AbelianVarietyTheory/Proofs/DivisorProofs/ExistsIntegerAddNsmul.lean}{GlobalTheory.UCls.NormOneMod.normOne_mod}
\leanref{Development/AbelianVarietyTheory/Proofs/DivisorProofs/ExistsIntegerAddNsmul.lean}{GlobalTheory.UCls.NormOneMod.exists_reps_normOne}
\begin{proof}
Every element of $\widehat{\mathcal O}_{F'}$ is congruent modulo $M\widehat{\mathcal O}_{F'}$ to an element of
$\mathcal O_{F'}$, and $\widehat{\mathcal O}_{E'}$ is generated as an additive group by the products $ay$ with
$a\in\widehat{\mathcal O}_{F'}$ and $y\in\mathcal O_{E'}$. So every element of $\widehat{\mathcal O}_{E'}$ is
congruent modulo $M\widehat{\mathcal O}_{E'}$ to an element of $\mathcal O_{E'}$. In a $\mathbf Z$-basis of
$\mathcal O_{E'}$, the elements with coordinates in $\{0,\dots,M-1\}$ form a set $R$ of at most
$M^{[E':\mathbf Q]}$ representatives of $\mathcal O_{E'}$ modulo $M\mathcal O_{E'}$. For each $\rho\in R$ such
that some element of $T\cap\widehat{\mathcal O}_{E'}$ is congruent to $\rho$ modulo $M\widehat{\mathcal O}_{E'}$,
choose one such element $s_\rho$, and let $S'$ be the set of the $s_\rho$.

Let $k\in T\cap\widehat{\mathcal O}_{E'}$. It is congruent to some $\rho\in R$, so $s=s_\rho$ is defined and
$k-s\in M\widehat{\mathcal O}_{E'}$. Put $k'=k\bar s$. Then $k'\in T$, $k=sk'$ as $s\bar s=1$, and
$k'-1=(k-s)\bar s$ lies in $M\widehat{\mathcal O}_{E'}$, as $\bar s$ is integral.
\end{proof}
\leanref{Development/AbelianVarietyTheory/Proofs/DivisorProofs/ExistsIntegerAddNsmul.lean}{GlobalTheory.UCls.NormOneMod.exists_integer_add_nsmul}
\leanref{Development/AbelianVarietyTheory/Proofs/DivisorProofs/ExistsIntegerAddNsmul.lean}{GlobalTheory.UCls.NormOneMod.exists_reps}

\begin{lemma}\label{cv-cusp-torus-norms}
Let $E'$ be a CM field with totally real subfield $F'$ and complex conjugation $z\mapsto\bar z$, unramified
over $F'$ at every finite prime. Let $r\ge1$, let $G\in M_r(E')$ be an invertible matrix with
$\bar G_{ab}=G_{ba}$ for all $a,b$, and put $h_G(u,v)=\sum_{a,b}u_aG_{ab}\bar v_b$ for $u,v\in E'^r$. Let
$\Lambda\subset E'^r$ be a finitely generated $\mathcal O_{E'}$-submodule that spans $E'^r$ and is self-dual
for $h_G$: a vector $u\in E'^r$ lies in $\Lambda$ if and only if $h_G(u,v)\in\mathcal O_{E'}$ for every
$v\in\Lambda$. Then the numbers $h_G(x,x)$ with $x\in\Lambda$ lie in $\mathcal O_{F'}$ and generate the unit
ideal of $\mathcal O_{F'}$.
\end{lemma}
\leanref{Development/AbelianVarietyTheory/Proofs/DivisorProofs/GaloisAutMemAdelicCompletionIntegral.lean}{GlobalTheory.UCls.Det.span_norm_eq_top}
\begin{proof}
For $u,v\in E'^r$, $\overline{h_G(u,v)}=h_G(v,u)$, as $\bar G_{ab}=G_{ba}$. For $x,y\in\Lambda$ the number
$h_G(x,y)$ lies in $\mathcal O_{E'}$ by self-duality; so $h_G(x,x)$, which is fixed by conjugation, lies in
$\mathcal O_{F'}$. Suppose that these numbers lie in a maximal ideal $\mathfrak m$ of $\mathcal O_{F'}$, and let
$P$ be a prime of $\mathcal O_{E'}$ above $\mathfrak m$.

Step 1: there is $t\in\mathcal O_{E'}$ with $t-\bar t\notin P$. As no finite prime of $F'$ ramifies in $E'$, the
ring map $\mathcal O_{F'}\to\mathcal O_{E'}$ is unramified. Let $\mu$ be the multiplication map
$x\otimes y\mapsto xy$ of $\mathcal O_{E'}\otimes_{\mathcal O_{F'}}\mathcal O_{E'}$. There is an idempotent $e$ of
this tensor product such that $\mu$ is the localization at $e$ \cite[Tag 02FL]{Sta}. So $\mu(e)=1$, and $e$ kills
the kernel of $\mu$, which contains $1\otimes b-b\otimes1$ for every $b\in\mathcal O_{E'}$. Let $\psi$ be the ring
map $x\otimes y\mapsto x\bar y$. Some $b\in\mathcal O_{E'}$ has $\bar b\ne b$, as conjugation is not the identity
on $E'$, the field of fractions of $\mathcal O_{E'}$. Applying $\psi$ to $(1\otimes b-b\otimes1)\,e=0$ gives
$(\bar b-b)\,\psi(e)=0$, so $\psi(e)=0$, as $\mathcal O_{E'}$ is a domain. Suppose that $t-\bar t\in P$ for every
$t\in\mathcal O_{E'}$. Then $\mu(x\otimes y)-\psi(x\otimes y)=x(y-\bar y)$ lies in $P$ for all
$x,y\in\mathcal O_{E'}$, so $\mu(z)-\psi(z)\in P$ for every $z$ in the tensor product, as both maps are additive.
For $z=e$ this puts $1=\mu(e)-\psi(e)$ in $P$, which is absurd.

Step 2: $h_G(x,y)\in P$ for all $x,y\in\Lambda$. By polarization,
$h_G(x,y)+h_G(y,x)=h_G(x+y,x+y)-h_G(x,x)-h_G(y,y)$ lies in $P$. Hence so does
\[
 (t-\bar t)\,h_G(x,y)=\big(h_G(tx,y)+h_G(y,tx)\big)-\bar t\big(h_G(x,y)+h_G(y,x)\big),
\]
and $h_G(x,y)\in P$, since $P$ is prime and $t-\bar t\notin P$.

Step 3: the contradiction. As $P$ is a nonzero proper ideal, there is $a\in P^{-1}$ with
$a\notin\mathcal O_{E'}$. For $x,y\in\Lambda$, $h_G(ax,y)=a\,h_G(x,y)$ lies in $aP\subset\mathcal O_{E'}$ by
Step 2, so $ax\in\Lambda$ by self-duality. Thus $a\Lambda\subset\Lambda$. The $\mathcal O_{E'}$-module
$\Lambda$ is finitely generated, and it is faithful for $\mathcal O_{E'}[a]$, being a nonzero submodule of
$E'^r$, as it spans $E'^r$ with $r\ge1$. So $a$ is integral over $\mathcal O_{E'}$
\cite[Proposition~5.1]{AM69}, and $a\in\mathcal O_{E'}$, a contradiction.
\end{proof}
\leanref{Development/AbelianVarietyTheory/Proofs/DivisorProofs/GaloisAutMemAdelicCompletionIntegral.lean}{GlobalTheory.UCls.Det.exists_sub_notMem}
\leanref{Development/AbelianVarietyTheory/Proofs/DivisorProofs/GaloisAutMemAdelicCompletionIntegral.lean}{GlobalTheory.UCls.Det.exists_herm_notMem}
% Lean: Step 1 is exists_sub_notMem for A = O_{F'}, B = O_{E'} and sigma = IsCMField.ringOfIntegersComplexConj (Unr E
% is Algebra.Unramified): the idempotent is Algebra.FormallyUnramified.elem (one_tmul_mul_elem, lmul_elem), psi is
% Algebra.TensorProduct.productMap id sigma, and the element moved by sigma comes from IsCMField.complexConj_ne_one in
% span_norm_eq_top.

\begin{lemma}\label{cv-cusp-torus-determinant}
Let $E'$ be a CM field with totally real subfield $F'$ and complex conjugation $z\mapsto\bar z$, unramified
over $F'$ at every finite prime. Let $\mathbf A_{E',f}=E'\otimes_{F'}\mathbf A_{F',f}$, with conjugation acting
through the first factor, and let $\widehat{\mathcal O}_{E'}\subset\mathbf A_{E',f}$ be the subring of the
elements integral at every finite place. Let $r,M\ge1$, let $G\in M_r(E')$ be an invertible matrix with
$\bar G_{ab}=G_{ba}$ for all $a,b$, and put $h_G(u,v)=\sum_{a,b}u_aG_{ab}\bar v_b$ for
$u,v\in\mathbf A_{E',f}^r$. Let $\Lambda\subset E'^r$ be a finitely generated $\mathcal O_{E'}$-submodule that
spans $E'^r$ and is self-dual for $h_G$: a vector $u\in E'^r$ lies in $\Lambda$ if and only if
$h_G(u,v)\in\mathcal O_{E'}$ for every $v\in\Lambda$. Let $\widehat\Lambda\subset\mathbf A_{E',f}^r$ be the
$\widehat{\mathcal O}_{E'}$-submodule generated by $\Lambda$, and let $K_\Lambda(M)$ be the group of the
$g\in GL_r(\mathbf A_{E',f})$ with $h_G(gu,gv)=h_G(u,v)$ for all $u,v$, $g\widehat\Lambda=\widehat\Lambda$ and
$(g-1)\widehat\Lambda\subset M\widehat\Lambda$. Let $k\in\mathbf A_{E',f}$ with $k\bar k=1$ and
$k-1\in M\widehat{\mathcal O}_{E'}$. Then $k=\det g$ for some $g\in K_\Lambda(M)$.
\end{lemma}
\leanref{Development/AbelianVarietyTheory/Proofs/DivisorProofs/GaloisAutMemAdelicCompletionIntegral.lean}{GlobalTheory.UCls.Det.exists_isKM_det_eq}
\begin{proof}
Step 1: a vector whose length is a unit. By Lemma~\ref{cv-cusp-torus-norms} there are
$x_1,\dots,x_n\in\Lambda$ such that the numbers $N_i=h_G(x_i,x_i)\in\mathcal O_{F'}$ generate the unit ideal.
So at each finite place $v$ of $F'$ some $N_i$ is a unit; choose such an index $i(v)$, and let
$\varepsilon_i\in\widehat{\mathcal O}_{F'}$ be the idempotent with $(\varepsilon_i)_v=1$ if $i(v)=i$ and
$(\varepsilon_i)_v=0$ otherwise. Put $e=\sum_i\varepsilon_ix_i\in\widehat\Lambda$. As
$\varepsilon_i\varepsilon_j=0$ for $i\ne j$, $\varepsilon_i^2=\varepsilon_i$ and
$\bar\varepsilon_i=\varepsilon_i$, $h_G(e,e)=\sum_i\varepsilon_iN_i$, whose component at $v$ is the unit
$N_{i(v)}$. So $h_G(e,e)$ has an inverse $b\in\widehat{\mathcal O}_{F'}$, and $\bar b=b$.

Step 2: the element $g$. Let $g$ be the matrix with $gx=x+(k-1)\,b\,h_G(x,e)\,e$ for
$x\in\mathbf A_{E',f}^r$. By the formula $\det(1+u\ell)=1+\ell(u)$ for a vector $u$ and a linear form $\ell$,
$\det g=1+(k-1)\,b\,h_G(e,e)=k$, which is a unit since $k\bar k=1$. For $x,y\in\mathbf A_{E',f}^r$, using
$\overline{h_G(y,e)}=h_G(e,y)$, $\bar b=b$ and $b\,h_G(e,e)=1$,
\[
 h_G(gx,gy)=h_G(x,y)+\big((k-1)+(\bar k-1)+(k-1)(\bar k-1)\big)\,b\,h_G(x,e)\,h_G(e,y)=h_G(x,y),
\]
as the coefficient is $k\bar k-1=0$.

Step 3: the lattice. Since $\Lambda$ is self-dual, $h_G(x,e)\in\widehat{\mathcal O}_{E'}$ for
$x\in\widehat\Lambda$. Write $k-1=My$ with $y\in\widehat{\mathcal O}_{E'}$. For $x\in\widehat\Lambda$,
$gx-x=M\,y\,b\,h_G(x,e)\,e$ lies in $M\widehat\Lambda$; in particular $g\widehat\Lambda\subset\widehat\Lambda$.
Conversely, $x'=x+(\bar k-1)\,b\,h_G(x,e)\,e$ lies in $\widehat\Lambda$, as $\bar k-1=M\bar y$ is integral, and
$h_G(x',e)=\bar k\,h_G(x,e)$ gives $gx'=x+\big((\bar k-1)+(k-1)\bar k\big)\,b\,h_G(x,e)\,e=x$. So
$g\widehat\Lambda=\widehat\Lambda$, and $g\in K_\Lambda(M)$.
\end{proof}
\leanref{Development/AbelianVarietyTheory/Proofs/DivisorProofs/GaloisAutMemAdelicCompletionIntegral.lean}{GlobalTheory.UCls.Det.exists_unit_norm}
\leanref{Development/AbelianVarietyTheory/Proofs/DivisorProofs/GaloisAutMemAdelicCompletionIntegral.lean}{GlobalTheory.UCls.Det.exists_isKM_of_unit_norm}

\begin{lemma}\label{cv-cusp-torus-bound-holds}
The torus class bound of Definition~\ref{cv-cusp-torus-bound} holds.
\end{lemma}
\leanref{Development/AbelianVarietyTheory/Proofs/DivisorProofs/NormOneTorusClassNumberBoundProofStatement.lean}{GlobalTheory.T.normOneTorusClassNumberBound_proof}
\leanref{Development/AbelianVarietyTheory/Proofs/DivisorProofs/NormOneTorusClassNumberBoundProof.lean}{GlobalTheory.UCls.torusBound}
\begin{proof}
Let $E'$, $T$, $G$, $\Lambda$ and $M$ be as in Definition~\ref{cv-cusp-torus-bound}. There are three facts.
First, by Lemma~\ref{cv-cusp-torus-ideal-classes}, there is a set $S\subset T$ with at most $h_{E'}$ elements
such that every $t\in T$ is $sus'$ with $s\in S$, $u\in E'$ of norm $1$, and $s'\in T$ integral at every
finite place. Second, by Lemma~\ref{cv-cusp-torus-congruence}, there is a set $S'\subset T$ with at most
$M^{[E':\mathbf Q]}$ elements such that every integral $s'\in T$ is $s''k$ with $s''\in S'$ and $k\in T$
integral and congruent to $1$ modulo $M$. Third, by Lemma~\ref{cv-cusp-torus-determinant}, every such $k$ is
the determinant of an element of $K_\Lambda(M)$. So the products $ss''$, $s\in S$, $s''\in S'$, meet every
class, and there are at most $h_{E'}M^{[E':\mathbf Q]}$ of them.
\end{proof}

\begin{lemma}\label{cv-cusp-sup-bound}
In Situation~\ref{cv-cusp-setting}, suppose that the torus class bound
(Definition~\ref{cv-cusp-torus-bound}) holds; this hypothesis is not used in the proof. Let $\sigma$ be a
good embedding and $w>0$, and let $\lambda_Z$ be a
weight on each component $Z$ at $\sigma$, of spread at most $w$ and with $\chi_{\lambda_Z}$ trivial on
$\Gamma_Z$. Then every family $P=(P_Z)_Z$ of cusp forms of weights $\lambda_Z$ satisfies
\[
 (\sup|P|)^2\le C_w\,\|P\|_{\mathrm{Pet}}^2.
\]
\end{lemma}
\leanref{Development/AutomorphicArithmetic/ThetaCorrespondence/FockModels/TamagawaPeterssonLambdaSqNonneg.lean}{UpperTheory.Lit.Tam.tamagawa_local_volume_factor}
\begin{proof}
For every component $Z$ choose an orthonormal frame $x_Z$ of $Z$ and a fundamental set $\mathcal F_Z$
for $\Gamma_Z$ at $x_Z$; both exist. Put $S=\sum_Z\mu_{x_Z}(\mathcal F_Z)$.

Step 1: unfolding. For every family $h=(h_Z)_Z$ of nonnegative measurable functions on the frames of the
components, each invariant under $\Gamma_Z$ and $B$,
\[
 S\int h_{Z(g)}(x(g))\,d\mu(g)=\sum_Z\int_{\mathcal F_Z}h_Z(x_Zn(z))\,d\mu_{x_Z}(z),
\]
and $0<S<\infty$. Indeed, through the map $g\mapsto(Z(g),x(g))$, the integral of such a function over
the quotient against a Haar measure of $U(W_j)(\mathbf A_F)$ is a constant $c_1>0$ times the sum over $Z$
of its integrals over the quotients by $\Gamma_Z$ of the real unitary groups of the $G_Z$. Each of these
is a constant $c_2>0$, depending only on $r$, $q$ and $f$, times its integral over $\mathcal F_Z$ against
$\mu_{x_Z}$, because $\mathcal F_Z$ is a fundamental set and $\mu_{x_Z}$ is the invariant measure of the
flag domain read in the chart at $x_Z$. Normalizing to the probability measure $\mu$, the left side
without the factor $S$ is a constant $c>0$, independent of $h$, times the right side; taking every
$h_Z=1$ gives $cS=1$.

Step 2: the bound. Let $Z$ be a component and $y$ a frame of $Z$. The functions $|P_{Z'}|^2$ are
measurable and invariant under $\Gamma_{Z'}$ and $B$, because the characters are trivial. By
Lemma~\ref{cv-cusp-inputs}(1), Step 1 and Lemma~\ref{cv-cusp-inputs}(2),
\[
 |P_Z|(y)^2\le C^{\mathrm{loc}}_w\int_{\mathcal F_Z}|P_Z|(x_Zn(z))^2\,d\mu_{x_Z}(z)
 \le C^{\mathrm{loc}}_w\,S\,\|P\|_{\mathrm{Pet}}^2\le C^{\mathrm{loc}}_wC^{\mathrm{mass}}\|P\|_{\mathrm{Pet}}^2.
\]
Taking the supremum over $Z$ and $y$ gives the claim, since $C_w=C^{\mathrm{loc}}_wC^{\mathrm{mass}}$.
\end{proof}
\leanref{Development/AbelianVarietyTheory/Proofs/DivisorProofs/AdelicProbabilityEqNormalisedFundamentalSetsStatement.lean}{GlobalTheory.T.adelicProbability_eq_normalised_fundamentalSets}
\leanref{Development/AbelianVarietyTheory/Proofs/DivisorProofs/AdelicProbabilityEqNormalisedFundamentalSets.lean}{GlobalTheory.LevGeo.adelicProb_mul_volFund}

\begin{lemma}\label{cv-cusp-bergman}
Let $(Q,\mu)$ be a probability space and $V$ a complex vector space, and for every $t\in Q$ let
$\ell_t:V\to\mathbf C$ be a linear map. Suppose that $t\mapsto|\ell_t(v)|$ is measurable for every
$v\in V$, and that $v=0$ whenever $\ell_t(v)=0$ for all $t$. Let $C\ge0$ be such that
\[
 |\ell_t(v)|^2\le C\int_Q|\ell_s(v)|^2\,d\mu(s)<\infty\qquad(v\in V,\ t\in Q).
\]
Then $V$ is finite-dimensional and $\dim_{\mathbf C}V\le C$.
\end{lemma}
\leanref{Development/AbelianVarietyTheory/Proofs/CovolumeBounds/FiniteAndFinrankLeOfPointwise.lean}{GlobalTheory.Bergman.finite_and_finrank_le_of_pointwise}
\begin{proof}
For $v,v'\in V$ put $\langle v,v'\rangle=\int_Q\ell_t(v)\overline{\ell_t(v')}\,d\mu(t)$, which is finite
by the Cauchy--Schwarz inequality. This is a Hermitian inner product on $V$: if $\langle v,v\rangle=0$,
the hypothesis gives $\ell_t(v)=0$ for every $t$, so $v=0$. Let $W\subset V$ be a finite-dimensional
subspace with an orthonormal basis $b_1,\dots,b_n$. For $t\in Q$, the restriction of $\ell_t$ to $W$
satisfies $|\ell_t(v)|^2\le C\langle v,v\rangle$, so the square of its norm as a linear form on $W$,
which is $\sum_i|\ell_t(b_i)|^2$, is at most $C$. Integrating over the probability space $Q$,
\[
 n=\sum_i\langle b_i,b_i\rangle=\int_Q\sum_i|\ell_t(b_i)|^2\,d\mu(t)\le C.
\]
So every finite-dimensional subspace of $V$ has dimension at most $C$, which gives both claims.
\end{proof}
\leanref{Development/AbelianVarietyTheory/Proofs/CovolumeBounds/FiniteAndFinrankLeOfPointwise.lean}{GlobalTheory.Bergman.finite_and_finrank_le}

\begin{proposition}\label{cv-cusp-rank-bound}
In Situation~\ref{cv-cusp-setting}, suppose that the torus class bound
(Definition~\ref{cv-cusp-torus-bound}) holds; this hypothesis is not used in the proof. Let
$\kappa=(\alpha,n)$ be a dominant weight. Then $S_\kappa$ is finite-dimensional over $F_c$, and
$\dim_{F_c}S_\kappa\le C_{w_\kappa}$.
\end{proposition}
\leanref{Development/AutomorphicArithmetic/ThetaCorrespondence/FockModels/TamagawaPeterssonLambdaSqNonneg.lean}{UpperTheory.Lit.Tam.finrank_cuspW_le_supConst2}
\begin{proof}
Let $\sigma$ be a good embedding, for example the inclusion of $F_c$ in $\mathbf C$, put $w=w_\kappa$, and
write $\lambda=\lambda(\kappa)$ for the weight on each component at $\sigma$.

Step 1: the weight. Since $\alpha$ is nonnegative, the entries of $\lambda$ at
$\tau$ lie between $-\sum_i\alpha_{\bar\tau',i}-n/2$ and $\sum_i\alpha_{\tau',i}+n/2$. As $\tau$ runs over
$\Phi$, the embeddings $\tau'$ and $\bar\tau'$ run over all embeddings of $E$, each once, so the spread
of $\lambda$ is at most $|\alpha|+fn\le|\alpha|+frn<w$. The characters $\chi_\lambda$ are trivial by the
compatibility conditions.

Step 2: the readings. Let $S_\sigma$ be the complex span of the readings $f_\sigma$, $f\in S_\kappa$.
Cusp forms of weight $\lambda$ on a component form a complex vector space, so every $P\in S_\sigma$ is a
family of cusp forms of weight $\lambda$, and Lemma~\ref{cv-cusp-sup-bound} and Step 1 give
$(\sup|P|)^2\le C_w\|P\|_{\mathrm{Pet}}^2$.

Step 3: the dimension of $S_\sigma$. Apply Lemma~\ref{cv-cusp-bergman} to
$Q=U(W_j)(F)\backslash U(W_j)(\mathbf A_F)$ with the measure $\mu$, $V=S_\sigma$, and
$\ell_t(P)=P_{Z(g)}(x(g))$ for a chosen representative $g$ of each $t\in Q$. Since
$|\delta_{\tau,i}(x(g))|=1$, $|\ell_t(P)|=|P_{Z(g)}|(x(g))$, which is measurable in $t$ because $P$ is
continuous, and which is at most $\sup|P|$. Hence, by Step 2,
$|\ell_t(P)|^2\le C_w\|P\|_{\mathrm{Pet}}^2=C_w\int_Q|\ell_s(P)|^2\,d\mu(s)$, and this is finite. If
$\ell_t(P)=0$ for every $t$, then $\|P\|_{\mathrm{Pet}}=0$, so $\sup|P|=0$ by Step 2, and $P=0$ because
$|P_Z|(x)$ is $|P_Z(x)|$ times a positive number. So $\dim_{\mathbf C}S_\sigma\le C_w$.

Step 4: from $F_c$ to $\mathbf C$. Let $f_1,\dots,f_s\in S_\kappa$ be linearly independent over $F_c$,
and let $w_1,\dots,w_s\in\mathbf C$ with $\sum_iw_if_{i,\sigma}=0$. Let $y$ be a point of $X$ with values in a
reduced finitely generated $F_c$-algebra $R$. For every $F_c$-algebra map $\psi:R\to\mathbf C$ over
$\sigma$ there are a component $Z$, a frame $x$ of $Z$ and $c\in\mathbf C$ with
$\psi(f(y))=c\,f_{\sigma,Z}(x)$ for every $f$ in $S_\kappa$. So the element $\sum_if_i(y)\otimes w_i$ of
$R\otimes_{F_c,\sigma}\mathbf C$, a reduced finitely generated $\mathbf C$-algebra, vanishes at all its
complex points, and is zero by the Nullstellensatz. Applying $1\otimes\theta$ for an $F_c$-linear map
$\theta:\mathbf C\to F_c$ gives $\sum_i\theta(w_i)f_i(y)=0$. This holds at every such point $y$, so
$\sum_i\theta(w_i)f_i=0$, hence $\theta(w_i)=0$ for every $i$ and every $\theta$, and every $w_i=0$. The
same argument with $s=1$ shows that a nonzero $f\in S_\kappa$ has a nonzero reading at $\sigma$. So
$\dim_{F_c}S_\kappa\le\dim_{\mathbf C}S_\sigma\le C_w$, and $S_\kappa$ is finite-dimensional.
\end{proof}
\leanref{Development/AbelianVarietyTheory/Proofs/DivisorProofs/SumWeightSpreadLeExponentNormAddStatement.lean}{GlobalTheory.T.sum_weightSpread_le_exponentNorm_add}
\leanref{Development/AbelianVarietyTheory/Proofs/DivisorProofs/CuspFormsFiniteDimensional.lean}{GlobalTheory.LevGeo.RankBd.finrank_le_of_sup}
\leanref{Development/AbelianVarietyTheory/Proofs/UnitaryExistence/ReadingLinearIndependentAndInjective.lean}{GlobalTheory.reading_linearIndependent_and_injective}

\begin{proposition}\label{cv-cusp-norm-bound}
In Situation~\ref{cv-cusp-setting}, suppose that the torus class bound
(Definition~\ref{cv-cusp-torus-bound}) holds; this hypothesis is not used in the proof. Let $\kappa$ be a
dominant weight and let $f\in\Lambda_\kappa$ be nonzero. Then $\|f\|_\sigma>0$ for every embedding $\sigma:F_c\to\mathbf C$, and
\[
 \frac1{[F_c:\mathbf Q]}\sum_{\sigma:F_c\to\mathbf C}\log\|f\|_\sigma\ge-\eta_\kappa.
\]
\end{proposition}
\leanref{Development/AutomorphicArithmetic/ThetaCorrespondence/FockModels/TamagawaPeterssonLambdaSqNonneg.lean}{UpperTheory.Lit.Tam.petersson_norm_tamagawa_formula}
\begin{proof}
Put $w=w_\kappa$, and let $\sigma$ be an embedding of $F_c$; then $\sigma^g$ is good. The reading
$P=f_{\sigma^g}$ is a family of cusp forms of weight $\lambda(\kappa)$, whose spread is less than $w$ and
whose characters are trivial, by Step 1 of the proof of Proposition~\ref{cv-cusp-rank-bound}. It is
nonzero since $f\ne0$, by Step 4 of that proof, so $\sup|P|>0$; and $\sup|P|$ is finite since cusp forms
are bounded. By Lemma~\ref{cv-cusp-sup-bound}, $(\sup|P|)^2\le C_w\|P\|_{\mathrm{Pet}}^2$. Multiplying by
$(2\pi\xi)^{-|\alpha|}$ gives $\|f\|_{\sigma,\infty}^2\le C_w\|f\|_\sigma^2$, so $\|f\|_\sigma>0$ and
\[
 \log\|f\|_{\sigma,\infty}\le\tfrac12\log C_w+\log\|f\|_\sigma.
\]
Averaging over the $[F_c:\mathbf Q]$ embeddings and using Lemma~\ref{cv-cusp-inputs}(3),
\[
 -\eta^{\sup}_\kappa\le\frac1{[F_c:\mathbf Q]}\sum_\sigma\log\|f\|_{\sigma,\infty}
 \le\frac12\log C_w+\frac1{[F_c:\mathbf Q]}\sum_\sigma\log\|f\|_\sigma,
\]
which is the claim, because $\eta_\kappa=\eta^{\sup}_\kappa+\frac12\log C_w$.
\end{proof}

\subsection{Realizations and the point conditions}\label{subsection-cv-realizations-point-conditions}

\begin{definition}\label{cv-realization-level-move}
Let $E$ be a CM field and let $X$ be the moduli scheme of a PEL datum over $E$ at a level
$(M,\mathfrak M)$ (Definition~\ref{cv-deep-cover}). Let $L$ be the lattice of the datum, a finitely
generated projective $\mathcal O_E$-module, and let $A$ be the universal abelian scheme over $X$, with its
action of $\mathcal O_E$, its polarization and its level structure: for every point $x$ of $X$, an
$\mathcal O_E$-linear map $\eta_x$ from $L$ to the points of the fibre $A_x$ killed by $\mathfrak M$, with
kernel $\mathfrak ML$, which is onto when $x$ has values in an algebraically closed field. Let $\gamma$ be
an $\mathcal O_E$-linear endomorphism of $L$.
\begin{enumerate}
\item A \emph{level automorphism} of $X$ with matrix $\gamma$ is a pair $(\rho,\nu)$ of an automorphism
  $\rho$ of $X$ over its base and a morphism $\nu\colon A\to A$ over $\rho$ that restricts to isomorphisms
  $A_x\to A_{\rho(x)}$ compatible with the actions of $\mathcal O_E$ and the polarizations, such that
  $\nu(\eta_x(\gamma a))=\eta_{\rho(x)}(a)$ for every point $x$ and every $a\in L$.
\item Let $y$ be a framed point of $X$ over a ring $S$, as in Situation~\ref{cv-definite-situation} and
  Definition~\ref{cv-deep-cover}, with point $x$. Its \emph{move} by a level automorphism $(\rho,\nu)$ is the framed point over $S$ with point
  $\rho\circ x$ whose bases and generator are the images under $\nu$ of those of $y$. Through $\nu$, the
  move is $y$ with its level structure $\eta_x$ replaced by $\eta_x\circ\gamma$.
\end{enumerate}
\end{definition}
\leanref{Formalism/LocalToGlobal/ProjectiveLattices/IsLevMove.lean}{Formalism.L2.Pj.IsLevMove}
\leanref{Formalism/LocalToGlobal/ProjectiveLattices/IsLevMove.lean}{Formalism.L2.Pj.FrMove}
\leanref{Formalism/LocalToGlobal/ProjectiveLattices/IsLevMove.lean}{Formalism.L2.Pj.levR}

\begin{lemma}\label{cv-realization-rigidity}
Let $R$ be a Noetherian integral domain, $N$ a finitely generated torsion-free $R$-module and $I\ne R$ an
ideal. Let $g$ be a prime number with $g\in I$ and $g\ne0$ in $R$, let $r\ge0$ be an integer with
$r+1<g$, and let $s\in R$ satisfy $1-s\in I$ and $g\mid sa$ for every $a\in I$. Let $\alpha$ be an
$R$-linear endomorphism of $N$ with $\alpha^n=1$ for some $n\ge1$ and $(\alpha-1)^rN\subset IN$. Then
$\alpha=1$.
\end{lemma}
\leanref{Development/AbelianVarietyTheory/Proofs/DivisorProofs/SubMemSmulTopOfLevel.lean}{GlobalTheory.NeatOne.eq_one_of_finite_order}
\begin{proof}
Step 1. If a submodule $P\subset N$ satisfies $P\subset IP$, then $P=0$. Indeed $P$ is finitely generated,
as $R$ is Noetherian, so there is $t\in R$ with $t-1\in I$ and $tP=0$ \cite[Corollary~2.5]{AM69}. As
$I\ne R$, $t\ne0$, and $N$ is torsion-free.

Step 2: order prime to $g$. Let $\beta$ be an endomorphism of $N$ with $\beta^b=1$ for some $b\ge1$ and
$(\beta-1)^rN\subset IN$, and suppose that $1-cb^{r+1}\in I$ for some $c\in R$. Then $\beta=1$. Put
$D=\beta-1$, $P=D(N)$, $S=\sum_{i<b}\beta^i$ and $e=b-S$. Then $SD=\beta^b-1=0$, so $e$ acts on $P$ as
multiplication by $b$. Also $e=DQ$ with $Q=-\sum_{i<b}\sum_{j<i}\beta^j$, which commutes with $D$, so
$e(N)\subset P$. For $z\in N$, $e^r(Dz)=D^r(Q^rDz)\in IN$, hence
$b^{r+1}Dz=e(e^r(Dz))\in e(IN)\subset IP$, and $Dz=(1-cb^{r+1})Dz+c\,b^{r+1}Dz\in IP$. So $P\subset IP$, and
$P=0$ by Step~1.

Step 3: powers of $g$. If $(\alpha-1)^rN\subset IN$, then $(\alpha^{g^a}-1)^rN\subset IN$ for every $a\ge0$.
Indeed the binomial coefficients $\binom{g^a}i$ with $0<i<g^a$ are divisible by $g$, so
$\alpha^{g^a}-1=(\alpha-1)^{g^a}+gD'$ for an endomorphism $D'$ that commutes with $\alpha$. In the expansion of
the $r$-th power, every term other than $(\alpha-1)^{g^ar}$ has the factor $g\in I$, and
$(\alpha-1)^{g^ar}N\subset(\alpha-1)^rN\subset IN$.

Step 4: order $g$. Let $\beta$ be an endomorphism of $N$ with $\beta^g=1$ and $(\beta-1)^rN\subset IN$. Then
$\beta=1$. As $\binom gi$ is divisible by $g$ for $0<i<g$ and $\binom g1=g$, there is a polynomial
$Q\in\mathbf Z[X]$ with $(X+1)^g=X^g+1+gX(1+XQ(X))$. Put $D=\beta-1$ and $q=Q(D)$. As $\beta^g=1$,
$g(D+D^2q)=-D^g$. Let $z\in N$, and write $g=2+d+r$ with $d\ge0$, which is possible as $r+1<g$. Since
$D^rz\in IN$ and $g\mid sa$ for $a\in I$, $sD^rz=gw$ for some $w\in N$. Then
\[
 g\,s\,(Dz+D^2qz)=-sD^gz=-D^{2+d}(sD^rz)=-g\,D^{2+d}w,
\]
and as $N$ is torsion-free and $g\ne0$, $sDz=D^2z'$ with $z'=-D^dw-sqz$. By induction on $k$,
$s^kDz\in D^{k+1}(N)$ for every $z$ and $k$: if $s^kDz=D^{k+1}z_1$, then
$s^{k+1}Dz=D^k(sDz_1)\in D^{k+2}(N)$. For $k=r$ this gives $s^rDz\in D(D^rN)\subset D(IN)=I\,D(N)$. As
$1-s\in I$, also $1-s^r\in I$, and $Dz=(1-s^r)Dz+s^rDz\in IP$ for $P=D(N)$. So $P\subset IP$, and $P=0$ by
Step~1.

Step 5. Write $n=g^ab$ with $g\nmid b$. By Step~3, $\beta=\alpha^{g^a}$ satisfies $(\beta-1)^rN\subset IN$,
and $\beta^b=1$. As $g$ is prime to $b^{r+1}$, there are integers $u,v$ with $ug+vb^{r+1}=1$, so
$1-vb^{r+1}=ug\in I$, and Step~2 gives $\alpha^{g^a}=1$. If $a\ge1$ and $\alpha^{g^a}=1$, then
$\beta=\alpha^{g^{a-1}}$ satisfies $\beta^g=1$ and, by Step~3, $(\beta-1)^rN\subset IN$, so
$\alpha^{g^{a-1}}=1$ by Step~4. Descending to $a=0$ gives $\alpha=1$.
\end{proof}
\leanref{Development/AbelianVarietyTheory/Proofs/DivisorProofs/SubMemSmulTopOfLevel.lean}{GlobalTheory.NeatOne.eq_one_of_prime_order}
\leanref{Development/AbelianVarietyTheory/Proofs/DivisorProofs/SubMemSmulTopOfLevel.lean}{GlobalTheory.NeatOne.exists_poly}
\leanref{Development/AbelianVarietyTheory/Proofs/DivisorProofs/EqOneOfFiniteOrderUnipotentModTwoPrimes.lean}{GlobalTheory.Neat.eq_one_of_unipotent_mod_of_unit}
\leanref{Development/AbelianVarietyTheory/Proofs/DivisorProofs/EqOneOfFiniteOrderUnipotentModTwoPrimes.lean}{GlobalTheory.Neat.pow_prime_pow_unipotent}

\begin{lemma}\label{cv-realization-fixed}
Let $X$, $(M,\mathfrak M)$ and $L$ be as in Definition~\ref{cv-realization-level-move}, and let $r\ge0$ be
an integer. Let $\mathfrak P\supset\mathfrak M$ be a prime ideal of $\mathcal O_E$ over a prime number $g$
with $r+1<g$, and suppose that $M\mathcal O_E=\mathfrak PJ$ for an ideal $J$ with
$\mathfrak P+J=\mathcal O_E$. Let $(\rho,\nu)$ be a level automorphism of $X$ with matrix $\gamma$ such that
$(\gamma-1)^rL\subset\mathfrak PL$, and suppose that $\rho$ fixes a point of $X$ with values in a field.
Then $(\gamma-1)L\subset\mathfrak ML$.
\end{lemma}
\leanref{Development/AbelianVarietyTheory/Proofs/DivisorProofs/SubMemSmulTopOfLevel.lean}{GlobalTheory.NeatOne.free_of_unip_one}
\begin{proof}
Let $x$ be a point of $X$ with values in a field $k$ such that $\rho\circ x=x$.

Step 1: a complex point. The point $x$ factors through the residue field $\kappa$ of a point $\xi$ of $X$,
and $\rho\circ x=x$ says that $\rho$ fixes $\xi$ and induces the identity on $\kappa$. As $X$ is of finite
type over a number field, $\kappa$ is a finitely generated extension of $\mathbf Q$, so it embeds in
$\mathbf C$, and the resulting complex point is fixed by $\rho$. So we may take $k=\mathbf C$.

Step 2: the fibre. Put $A=A_x$ and $\eta=\eta_x$. As $\rho(x)=x$, $\nu$ restricts to an automorphism
$\theta$ of $A$ compatible with the action of $\mathcal O_E$ and the polarization, with
$\theta(\eta(\gamma a))=\eta(a)$ for $a\in L$. The automorphisms of a polarized abelian variety form a
finite group \cite[\S21]{Mum70}, so $\theta^N=1$ for some $N\ge1$. Write $A(\mathbf C)=V/\Lambda$, with $V$
the Lie algebra of $A$ and $\Lambda$ the period lattice \cite[\S1]{Mum70}. Then $\Lambda$ is a finitely
generated torsion-free $\mathcal O_E$-module; the differential of $\theta$ restricts to an
$\mathcal O_E$-linear automorphism $F$ of $\Lambda$ with $F^N=1$; and the map
$\psi\colon\Lambda\to A(\mathbf C)$, $\lambda\mapsto\lambda/M\bmod\Lambda$, is $\mathcal O_E$-linear, with kernel
$M\Lambda$ and image the points killed by $M$, and $\psi\circ F=\theta\circ\psi$.

Step 3: transfer to $\Lambda$. As $\mathfrak M\bar{\mathfrak M}=M\mathcal O_E$, $M\in\mathfrak M$ kills every
$\eta(a)$, so $\eta(L)\subset\psi(\Lambda)$. Put $F'=F^{-1}$.
\begin{enumerate}
\item[(a)] If $\psi(l)=\eta(a)$, then $\psi(F'l)=\eta(\gamma a)$. Indeed, choose $\mu\in\Lambda$ with
  $\psi(\mu)=\eta(\gamma a)$. Then $\psi(F(\mu-F'l))=\theta(\eta(\gamma a))-\psi(l)=0$, so
  $F(\mu-F'l)\in M\Lambda$, hence $\mu-F'l\in M\Lambda$ and $\psi(F'l)=\psi(\mu)$. In particular
  $\psi((F'-1)l)=\eta((\gamma-1)a)$.
\item[(b)] Let $W=\psi^{-1}(\eta(L))$, an $\mathcal O_E$-submodule of $\Lambda$, finitely generated and
  torsion-free, with $\psi(W)=\eta(L)$. By (a), $F'W\subset W$, and $F'^N=1$.
\item[(c)] $M\Lambda\subset\mathfrak PW$. Let $c'\in\bar{\mathfrak M}$ and $l\in\Lambda$. For $c\in\mathfrak M$,
  $cc'\in M\mathcal O_E$, so $c\,\psi(c'l)=0$. Thus $\psi(c'l)$ is killed by $\mathfrak M$, so it lies in
  $\eta(L)$, and $c'l\in W$. Hence $M\Lambda=\mathfrak M\bar{\mathfrak M}\Lambda\subset\mathfrak MW\subset\mathfrak PW$.
\item[(d)] $(F'-1)^rW\subset\mathfrak PW$. Let $l\in W$ with $\psi(l)=\eta(a)$. Applying (a) $r$ times gives
  $\psi((F'-1)^rl)=\eta((\gamma-1)^ra)$. Write $(\gamma-1)^ra=\sum_ic_ia_i$ with $c_i\in\mathfrak P$ and
  $a_i\in L$, and choose $l_i\in W$ with $\psi(l_i)=\eta(a_i)$. Then $(F'-1)^rl-\sum_ic_il_i$ lies in
  $\ker\psi=M\Lambda$, and $(F'-1)^rl\in\mathfrak PW$ by (c).
\end{enumerate}

Step 4: conclusion. As $\mathfrak P+J=\mathcal O_E$, there is $s\in J$ with $1-s\in\mathfrak P$. For
$a\in\mathfrak P$, $sa\in\mathfrak PJ=M\mathcal O_E$, and $g\mid M$ because $M\in\mathfrak M\subset\mathfrak P$;
so $g\mid sa$. Lemma~\ref{cv-realization-rigidity}, with $R=\mathcal O_E$, $N=W$, $I=\mathfrak P$ and $\alpha$
the restriction of $F'$ to $W$, gives $F'l=l$ for $l\in W$. For $a\in L$ choose $l\in W$ with
$\psi(l)=\eta(a)$; by (a), $\eta(\gamma a)=\psi(F'l)=\psi(l)=\eta(a)$. So $\gamma a-a\in\ker\eta=\mathfrak ML$.
\end{proof}
\leanref{Development/AbelianVarietyTheory/Proofs/DivisorProofs/SubMemSmulTopOfLevel.lean}{GlobalTheory.NeatOne.sub_mem_smul_top_of_level}
\leanref{Development/AbelianVarietyTheory/Proofs/DivisorProofs/ExistsFixedComplexPointOfFixedFieldPointStatement.lean}{GlobalTheory.T.exists_fixed_complexPoint_of_fixed_fieldPoint}
\leanref{Development/AbelianVarietyTheory/Proofs/DivisorProofs/ExistsFiniteOrderFibreAutAtLevelMoveFixedPointStatement.lean}{GlobalTheory.T.exists_finiteOrder_fibreAut_at_levelMove_fixedPoint}
\leanref{Development/AbelianVarietyTheory/Proofs/DivisorProofs/PeriodLatticeActionOfFiniteOrderAutomorphismStatement.lean}{GlobalTheory.T.periodLattice_action_of_finiteOrder_automorphism}

\begin{lemma}\label{cv-realization-definite}
In Situation~\ref{cv-definite-situation}, let $L$ be the lattice of the PEL datum and let $V$ be the
$F_c$-space of forms of a weight $(\nu,n)$. Let $G$ be a group and $(\gamma_g)_{g\in G}$ a family of
$\mathcal O_E$-linear endomorphisms of $L$ such that $X$ has a level automorphism with matrix $\gamma_g$ for
every $g$ (Definition~\ref{cv-realization-level-move}). For a framed point $y$ over an $F_c$-algebra, write
$y_g$ for its move by such a level automorphism. Suppose that the moves of $y$ by any two level
automorphisms with matrix $\gamma_g$ are isomorphic, that $y_1$ is isomorphic to $y$ and $(y_g)_h$ to
$y_{gh}$, and that the function $gf\colon y\mapsto f(y_g)$ lies in $V$ for every $f\in V$. For a function
$b$ on $X$ and a point $x$ of $X$ with values in $S$, write $b(x)\in S$ for the value of $b$ at $x$.
\begin{enumerate}
\item $X$ is the spectrum of an $F_c$-algebra $B$ that is a finite product of finite separable field
  extensions of $F_c$.
\item For every $g\in G$ there is a unique $F_c$-algebra automorphism $\tau(g)$ of $B$ such that
  $b(x_g)=\tau(g)(b)(x)$ for every $b\in B$ and every framed point $y$ with point $x$, where $x_g$ is the point
  of $y_g$. The map $g\mapsto\tau(g)$ is a homomorphism.
\item For $b\in B$ and $f\in V$, the function $bf\colon y\mapsto b(x)f(y)$, where $x$ is the point of $y$,
  lies in $V$, and $g(bf)=\tau(g)(b)\cdot gf$.
\item Let $\sigma\colon F_c\to\mathbf C$ be an embedding and $H\subset G$ a subgroup such that $h=1$ whenever
  $h\in H$ and some framed point $y$ over $\mathbf C$ via $\sigma$ has the same point as $y_h$. Then $H$ acts
  freely on the set of $F_c$-algebra homomorphisms $\chi\colon B\to\mathbf C$, with $\mathbf C$ an
  $F_c$-algebra through $\sigma$, by $h\cdot\chi=\chi\circ\tau(h)^{-1}$.
\end{enumerate}
\end{lemma}
\leanref{Development/AbelianVarietyTheory/Proofs/DivisorProofs/HalfLevelSubmodulePermutation.lean}{GlobalTheory.DefSplit.split_rules}
\begin{proof}
(1) The moduli scheme of a PEL datum of signature $(p,q)$ is smooth and quasi-projective over its
coefficient field, of relative dimension $fpq$ \cite{Kot92b,Lan13}. In signature $(p,0)$ this is $0$, so
$X$ is \'etale and quasi-compact over $F_c$. An \'etale scheme over a field is a disjoint union of open
points, each the spectrum of a finite separable extension; as $X$ is quasi-compact there are finitely many,
and $X$ is the spectrum of the product of these extensions.

(2) Let $(\rho,\nu)$ be a level automorphism with matrix $\gamma_g$. Its move of a framed point $y$ has point
$\rho\circ x$, and $y_g$ is isomorphic to it, so $x_g=\rho\circ x$ and $b(x_g)=(\rho^*b)(x)$; take
$\tau(g)=\rho^*$. An element of $B$ is determined by its values at the points of framed points, because
$B$ is reduced and every point of $X$ with values in an algebraically closed field is the point of a framed
point: the eigenspaces of the Lie algebra and the module of sections are vector spaces over that field.
This gives uniqueness. For framed points $y$, the point of $(y_g)_h$ is $x_{gh}$, so
$\tau(gh)(b)(x)=b(x_{gh})=\tau(h)(b)(x_g)=\tau(g)(\tau(h)(b))(x)$, and $\tau(gh)=\tau(g)\tau(h)$ by
uniqueness. In the same way $\tau(1)=1$, so every $\tau(g)$ is invertible.

(3) The function $y\mapsto b(x)$ is a form of weight $(0,0)$: it is compatible with base change and does
not depend on the bases and the generator of $y$. The product of forms of weights $(\nu,n)$ and $(0,0)$ is
a form of weight $(\nu,n)$, so $bf\in V$. Finally
$(g(bf))(y)=b(x_g)f(y_g)=\tau(g)(b)(x)\,(gf)(y)$.

(4) Let $h\in H$ and let $\chi$ satisfy $h\cdot\chi=\chi$, that is, $\chi\circ\tau(h)=\chi$. Let $x$ be the
point of $X=\operatorname{Spec}B$ with values in $\mathbf C$ given by $\chi$, and let $y$ be a framed point
with point $x$, which exists as in (2). For every $b\in B$, $b(x_h)=\chi(\tau(h)(b))=\chi(b)=b(x)$, so
$x_h=x$, as $X$ is affine. So $y$ and $y_h$ have the same point, and $h=1$.
\end{proof}
\leanref{Development/AbelianVarietyTheory/Proofs/DivisorProofs/HalfLevelSubmodulePermutation.lean}{GlobalTheory.DefSplit.exists_map}
\leanref{Development/AbelianVarietyTheory/Proofs/DivisorProofs/HalfLevelSubmodulePermutation.lean}{GlobalTheory.DefSplit.isAffine_XF}
\leanref{Development/AbelianVarietyTheory/Proofs/DivisorProofs/HalfLevelSubmodulePermutation.lean}{GlobalTheory.DefSplit.etale_cst}
\leanref{Development/AbelianVarietyTheory/Proofs/FoundationProofs/IsOpenImmersionOfSection.lean}{GlobalTheory.FlatLine.DefinitePoint.etale_strK}

\begin{lemma}\label{cv-realization-levels}
Let $E$ be a number field, $L$ a finitely generated $\mathcal O_E$-module and $r\ge1$ an integer. Let
$\mathfrak P_v$, for $v$ in a finite set $T$, be distinct maximal ideals of $\mathcal O_E$ with residue fields
$\mathbf F_{g_v}$ of prime order $g_v$, and let $I\subset\mathcal O_E$ be an ideal with
$I+\mathfrak P_v=\mathcal O_E$ for every $v$. Put $\mathfrak M=\prod_v\mathfrak P_v$ and
$\mathfrak M_Y=\mathfrak MI$, let $\Gamma$ be the group of $\mathcal O_E$-linear automorphisms of
$L/\mathfrak M_YL$, and put $G=\prod_{v\in T}GL_r(\mathbf F_{g_v})$. For each $v$ let
$\varphi_v\colon L/\mathfrak P_vL\to\mathbf F_{g_v}^r$ be an isomorphism of $\mathbf F_{g_v}$-vector spaces,
and for $\gamma\in\Gamma$ let $\pi(\gamma)_v$ be the matrix, through $\varphi_v$, of the automorphism of
$L/\mathfrak P_vL$ induced by $\gamma$.
\begin{enumerate}
\item The map $\pi\colon\Gamma\to G$ is a homomorphism, and its kernel is the group of the $\gamma$ that induce
  the identity on $L/\mathfrak ML$.
\item There is a homomorphism $\iota\colon G\to\Gamma$ with $\pi\circ\iota=\mathrm{id}$ whose image commutes
  with every element of $\ker\pi$.
\item Let $\gamma_g$, for $g\in G$, be endomorphisms of $L$ such that $\gamma_g$ induces on $L/\mathfrak P_vL$
  the automorphism with matrix $g_v$ through $\varphi_v$, for every $v$. If an endomorphism $\tilde\gamma$ of
  $L$ induces $\gamma\in\Gamma$ on $L/\mathfrak M_YL$, then $(\tilde\gamma-\gamma_{\pi(\gamma)})L\subset\mathfrak ML$.
\end{enumerate}
\end{lemma}
\leanref{Development/AbelianVarietyTheory/Definitions/Divisors/IsCoprimeOfNatCast.lean}{GlobalTheory.LevDesc.typeRed}
\leanref{Development/AbelianVarietyTheory/Definitions/Divisors/IsCoprimeOfNatCast.lean}{GlobalTheory.LevDesc.typeEmb}
\leanref{Development/AbelianVarietyTheory/Definitions/Divisors/IsCoprimeOfNatCast.lean}{GlobalTheory.LevDesc.IsTypeLevel.commute}
\leanref{Development/AbelianVarietyTheory/Proofs/DivisorProofs/ExistsTypeFrameCoverLevelMapsReductionsStatement.lean}{GlobalTheory.T.exists_typeFrame_cover_levelMaps_reductions}
\begin{proof}
The ideals $\mathfrak P_v$ and $I$ are pairwise coprime, so their product $\mathfrak M_Y$ is their
intersection, and by the Chinese remainder theorem the map
\[
 L/\mathfrak M_YL\to\bigoplus_{v\in T}L/\mathfrak P_vL\oplus L/IL
\]
is an isomorphism of $\mathcal O_E$-modules. Its summands correspond to the submodules of $L/\mathfrak M_YL$
killed by the $\mathfrak P_v$ and by $I$, so every $\gamma\in\Gamma$ preserves them, and $\Gamma$ is the product
of the groups of $\mathcal O_E$-linear automorphisms of the summands. As $\mathcal O_E$ acts on
$L/\mathfrak P_vL$ through $\mathbf F_{g_v}$, these are its $\mathbf F_{g_v}$-linear automorphisms, and
$\varphi_v$ identifies their group with $GL_r(\mathbf F_{g_v})$. Under these identifications $\pi$ is the
projection of $\Gamma$ onto the factors of the $v\in T$, which is a homomorphism. Its kernel consists of
the $\gamma$ that induce the identity on every $L/\mathfrak P_vL$, that is, on
$L/\mathfrak ML=\bigoplus_vL/\mathfrak P_vL$. This is (1). For (2), let $\iota(g)$ act as $g_v$ on
$L/\mathfrak P_vL$, through $\varphi_v$, and as the identity on $L/IL$; it commutes with the elements of
$\ker\pi$, which act only on $L/IL$. For (3), $\tilde\gamma$ and $\gamma_{\pi(\gamma)}$ induce the same
automorphism of every $L/\mathfrak P_vL$, as $\mathfrak M_Y\subset\mathfrak P_v$. So
$\tilde\gamma-\gamma_{\pi(\gamma)}$ maps $L$ into $\bigcap_v\mathfrak P_vL=\mathfrak ML$.
\end{proof}
\leanref{Development/AbelianVarietyTheory/Definitions/Divisors/IsCoprimeOfNatCast.lean}{GlobalTheory.LevDesc.IsTypeLevel.parts}
\leanref{Development/AbelianVarietyTheory/Definitions/Divisors/IsCoprimeOfNatCast.lean}{GlobalTheory.LevDesc.IsTypeLevel.typeRed_bijective}
\leanref{Development/AbelianVarietyTheory/Definitions/Divisors/IsCoprimeOfNatCast.lean}{GlobalTheory.LevDesc.typeRed_eq_one}

\begin{lemma}\label{cv-realization-pullback}
In Situation~\ref{cv-cusp-setting}, suppose that the level of $X$ is a level $(M,\mathfrak M)$ with $M$
prime to $15$, let $L$ be the lattice of the PEL datum of $X$, and let $Y$ be a deep cover of $X$ of
exponent $e$, with deep level $(15^eM,\mathfrak MI)$ and level map $p\colon Y\to X$
(Definition~\ref{cv-deep-cover}). For a framed point $y$ of $Y$ over a field, let $p(y)$ be the framed point
of $X$ whose point is the image of that of $y$, with the bases and the generator of $y$ carried over by the
identification of the abelian varieties. Let $\Gamma$ be the group of $\mathcal O_E$-linear automorphisms of
$L/\mathfrak MIL$, and $\Gamma_1\subset\Gamma$ the subgroup of those that induce the identity on
$L/\mathfrak ML$.
\begin{enumerate}
\item Let $\gamma\in\Gamma$. For every endomorphism $\tilde\gamma$ of $L$ that induces $\gamma$, $Y$ has a
  level automorphism with matrix $\tilde\gamma$ (Definition~\ref{cv-realization-level-move}), and the moves of
  a framed point $y$ of $Y$ by all these level automorphisms are isomorphic. Write $y_\gamma$ for any of them.
\item Let $\gamma\in\Gamma$, let $\tilde\gamma$ induce $\gamma$, and let $\gamma'$ be an endomorphism of $L$
  with $(\tilde\gamma-\gamma')L\subset\mathfrak ML$ such that $X$ has a level automorphism with matrix
  $\gamma'$. Then, for every framed point $y$ of $Y$ over a field, $p(y_\gamma)$ is isomorphic to the move of
  $p(y)$ by that level automorphism.
\item Let $y$ and $y'$ be framed points of $Y$ over an algebraically closed field such that $p(y)$ and
  $p(y')$ are isomorphic. Then $y'$ is isomorphic to $y_\gamma$ for some $\gamma\in\Gamma_1$.
\end{enumerate}
\end{lemma}
\leanref{Development/AbelianVarietyTheory/Proofs/JetHarmonics/ActLevFullLevelGroupFramedPoints.lean}{GlobalTheory.FullLevelAction.actLaws_levFull}
\leanref{Development/AbelianVarietyTheory/Proofs/DivisorProofs/PullActLev.lean}{GlobalTheory.LevPullEquiv.pull_actLev}
\leanref{Development/AbelianVarietyTheory/Proofs/DivisorProofs/LevelCoverFibreTorsorTrivialLevelChanges.lean}{GlobalTheory.LevOrbit.fibre_torsor}
\begin{proof}
Let $\eta$ be the level structure of $Y$.

(1) As $\gamma$ is an automorphism of $L/\mathfrak MIL$, $\eta\circ\tilde\gamma$ is again a level structure:
its kernel is $\mathfrak MIL$, and it is onto the points killed by $\mathfrak MI$ over algebraically closed
fields. It depends only on $\gamma$. As $Y$ represents its moduli problem
(Theorem~\ref{cv-deep-pel-facts}(1)), the universal abelian scheme with $\eta$ replaced by
$\eta\circ\tilde\gamma$ is identified, by an isomorphism $\nu$, with the pullback of the universal object by a
unique automorphism $\rho$ of $Y$, and $(\rho,\nu)$ is a level automorphism with matrix $\tilde\gamma$. Through
its $\nu$, the move of $y$ by any level automorphism with matrix $\tilde\gamma$ is $y$ with $\eta$ replaced by
$\eta\circ\tilde\gamma$. These data depend only on $\gamma$, and framed points with the same data are
isomorphic, as $Y$ represents its moduli problem.

(2) The level structure of $p(y)$ is $15^e\eta$ (Definition~\ref{cv-deep-cover}(4)). Through $\nu$,
$y_\gamma$ is $y$ with $\eta$ replaced by $\eta\circ\tilde\gamma$, so $p(y_\gamma)$ is $p(y)$ with $15^e\eta$
replaced by $15^e\eta\circ\tilde\gamma$. As $15^e\eta$ is the level structure of $X$ at a point, it kills
$\mathfrak ML$, so $15^e\eta\circ\tilde\gamma=15^e\eta\circ\gamma'$. The move of $p(y)$ by a level automorphism
with matrix $\gamma'$ is $p(y)$ with $15^e\eta$ replaced by $15^e\eta\circ\gamma'$, and framed points of $X$
with the same data are isomorphic.

(3) Identify the abelian varieties of $y$ and $y'$ through an isomorphism of $p(y)$ with $p(y')$. Then $y'$
is $y$ with $\eta$ replaced by a level structure $\eta'$ with $15^e\eta'=15^e\eta$. Both induce isomorphisms
of $L/\mathfrak MIL$ onto the points killed by $\mathfrak MI$, so $\eta'=\eta\circ\gamma$ for some
$\gamma\in\Gamma$. For $a\in L$, $\eta(\gamma a-a)$ is killed by $15^e$ and by $\mathfrak MI$, hence by
$\mathfrak MI+15^e\mathcal O_E$. This ideal is $I$: $15^e\mathcal O_E=I\bar I$, and $\mathfrak M+\bar I=\mathcal O_E$
because $M$ is prime to $15$, so $\mathfrak MI+I\bar I=I$. The elements of $L$ that $\eta$ maps to points
killed by $I$ are those $b$ with $Ib\subset\mathfrak MIL$, which form $\mathfrak ML$ as $I$ is invertible. So
$\gamma a-a\in\mathfrak ML$, $\gamma\in\Gamma_1$, and $y'$ is isomorphic to $y_\gamma$ by (1).
\end{proof}

\begin{lemma}\label{cv-point-group-determinant}
Let $K$ be a field of characteristic $0$ with a homomorphism of fields $K\to\mathbf C$, let $V$ be a
$K$-vector space, let $S$ be a finite group and let $I$ be a finite set. For each $i\in I$ let
$\mathrm{ev}_i\colon V\to\mathbf C^S$ be a $K$-linear map, let $J_i$ be a finite set with a map
$r_i\colon J_i\to S$ whose image contains $1$, and let $W_i\subset V$ be a $K$-subspace, such that:
\begin{enumerate}
\item[(a)] $\mathrm{ev}_i(v)(h)=0$ for $v\in W_i$ and $h\in S$ outside the image of $r_i$;
\item[(b)] if complex numbers $a_j$, $j\in J_i$, satisfy $\sum_ja_j\,\mathrm{ev}_i(v)(r_i(j))=0$ for every
  $v\in W_i$, then $a_j=0$ for every $j$.
\end{enumerate}
Then there is $v\in V$ with $\det\big(\mathrm{ev}_i(v)(hg)\big)_{h,g\in S}\ne0$ for every $i\in I$.
\end{lemma}
\leanref{Development/AbelianVarietyTheory/Proofs/DivisorProofs/ExistsVectorGroupDeterminantsNeZeroStatement.lean}{GlobalTheory.T.exists_vector_groupDeterminants_ne_zero}
\leanref{Development/AbelianVarietyTheory/Proofs/DivisorProofs/ExistsVectorGroupDeterminantsNeZero.lean}{GlobalTheory.GenericGroupDet.exists_vec_det_ne_zero}
\begin{proof}
Step 1: a delta function. Fix $i$. The vectors $(\mathrm{ev}_i(v)(r_i(j)))_{j\in J_i}$, $v\in W_i$, span
$\mathbf C^{J_i}$: otherwise a nonzero linear form $x\mapsto\sum_ja_jx_j$ on $\mathbf C^{J_i}$ vanishes on
all of them, and (b) gives $a=0$. So there are $w_1,\dots,w_n\in W_i$ and $c_1,\dots,c_n\in\mathbf C$ such
that $\sum_kc_k\,\mathrm{ev}_i(w_k)(r_i(j))$ is $1$ if $r_i(j)=1$ and $0$ otherwise, for every $j$. Then
$\sum_kc_k\,\mathrm{ev}_i(w_k)$ is the function $\delta$ on $S$ with $\delta(1)=1$ and $\delta(h)=0$ for
$h\ne1$: on the image of $r_i$ by the choice of the $c_k$, and outside it by (a), since $1$ lies in the
image.

Step 2: the group determinant of $\delta$. The matrix $(\delta(hg))_{h,g\in S}$ has the entry $1$ where
$g=h^{-1}$ and $0$ elsewhere. It is a permutation matrix, so its determinant is $\pm1$.

Step 3: a common vector. Let $v_1,\dots,v_N$ be the vectors $w_k$ of Step 1 for all $i$ together, and for
$x\in\mathbf C^N$ put $P_i(x)=\det\big(\sum_px_p\,\mathrm{ev}_i(v_p)(hg)\big)_{h,g\in S}$, a polynomial in
$x$. It is not zero: at the $x$ with $x_p=c_k$ when $v_p$ is the vector $w_k$ chosen for $i$, and $x_p=0$
for the other $p$, the matrix is $(\delta(hg))$ by Step 1, and Step 2 applies. So $P=\prod_iP_i$ is a
nonzero polynomial. The image $T$ of $K$ in $\mathbf C$ is infinite, as $K$ has characteristic $0$, and a
nonzero polynomial in $N$ variables does not vanish on all of $T^N$. So there is $t\in K^N$ with
$P(x)\ne0$ for the image $x$ of $t$. Put $v=\sum_pt_pv_p$. By $K$-linearity,
$\mathrm{ev}_i(v)=\sum_px_p\,\mathrm{ev}_i(v_p)$, so $\det\big(\mathrm{ev}_i(v)(hg)\big)=P_i(x)\ne0$ for every
$i$.
\end{proof}
% Lean: Step 1 is exists_combination_delta (the span by Submodule.exists_le_ker_of_lt_top), Step 2
% det_delta_ne_zero (the permutation matrix of inversion), Step 3 the polynomials P i in
% MvPolynomial over C and MvPolynomial.funext_set on the infinite range of algebraMap F C.

\begin{lemma}\label{cv-point-regular-section}
In Situation~\ref{cv-cusp-setting}, suppose that the level of $X$ is a level $(M,\mathfrak M)$ with $M$
prime to $15$, let $Y$ be a deep cover of $X$ (Definition~\ref{cv-deep-cover}) with coefficient field
$F_Y$, and let $R$ be the $F_Y$-algebra of rational sections of all weights on $Y$
(Lemma~\ref{cv-cover-level-automorphisms}). Let $I$ be the set of pairs $(\tau,i)$ that index the basis
vectors of the frames, together with one further index $0$. For a framed point $y$ of $Y$ over a field
$k\supset F_Y$ and $t\in(k^\times)^I$, let $y\cdot t$ be the framed point with the same point whose basis
vectors are multiplied by the $t_{\tau,i}$ and whose generator of sections is multiplied by $t_0$. An
element $g\in R$ has \emph{exponent} $e\in\mathbf Z^I$ if $g(y\cdot t)=t^e\,g(y)$ for all such $y$ and
$t$, where $t^e=\prod_it_i^{e_i}$. Let $F\in R$ have an exponent, and suppose that every irreducible
component $Z$ of $Y\otimes F_Y$ contains the point of a complex framed point $y$ of $Y$ with $F(y)\ne0$.
Then $F$ is not a zero divisor in $R$.
\end{lemma}
\leanref{Development/AbelianVarietyTheory/Proofs/DivisorProofs/IntegerExponentSectionMulEqZeroImpEqZeroStatement.lean}{GlobalTheory.T.integerExponentSection_mul_eq_zero_imp_eq_zero}
\leanref{Development/AbelianVarietyTheory/Proofs/DivisorProofs/IntegerExponentSectionMulEqZeroImpEqZero.lean}{GlobalTheory.SecRingNonZeroDiv.eq_zero_of_mul_eq_zero}
\begin{proof}
A section of a weight has an exponent, determined by the weight, so every $f\in R$ is a finite sum
$\sum_{d\in D}f_d$ with $D\subset\mathbf Z^I$ finite and $f_d$ of exponent $d$. Each component $Z$ has a
\emph{generic framed point} $y_Z$: a framed point over the generic point of $Z$, with values in the
field $\Omega_Z$ of rational functions over the residue field of that point in the entries of a generic
upper unitriangular change of frame. Two facts about it are used, for $g\in R$ with an exponent.
\begin{enumerate}
\item[(i)] If $g(y_Z)=0$, then $g$ vanishes at every framed point whose point lies on $Z$.
\item[(ii)] If $g(y_Z)=0$ for every component $Z$, then $g=0$.
\end{enumerate}
Let $f\in R$ with $Ff=0$, and write $f=\sum_{d\in D}f_d$ as above. Let $e$ be the exponent of $F$. By (i)
and the hypothesis, $F(y_Z)\ne0$ for every $Z$. Fix $Z$ and let $t\in(\Omega_Z^\times)^I$. The framed point
$y_Z\cdot t$ has the point of $y_Z$, and evaluating $Ff=0$ there gives
\[
 0=F(y_Z\cdot t)\sum_{d\in D}f_d(y_Z\cdot t)=t^eF(y_Z)\sum_{d\in D}t^df_d(y_Z).
\]
As $t^eF(y_Z)\ne0$, $\sum_df_d(y_Z)\,t^d=0$ for every $t$. The maps $t\mapsto t^d$, $d\in D$, are
homomorphisms $(\Omega_Z^\times)^I\to\Omega_Z^\times$, and they are distinct: if $d_i\ne d'_i$, then at the
$t$ with $t_i=2$ and all other entries $1$ they take the values $2^{d_i}\ne2^{d'_i}$, as $\Omega_Z$ has
characteristic $0$. By Artin's theorem on the linear independence of distinct characters
\cite[Theorem VI.4.1]{Lan02}, $f_d(y_Z)=0$ for every $d$. This holds for every $Z$, so $f_d=0$ by (ii) for
every $d$, and $f=0$.
\end{proof}
% Lean: the decomposition is toSubmodule_secRing (secRing is the sum of the secZ e); the generic framed
% point is Formalism.L2.Pj.genFr (IsGenFr, over Om Z = FractionRing (MvPolynomial (FlagVar Dw) (kap Z))),
% from GlobalTheory.Reg.genFrEx_and_genFromTest; (i) is GlobalTheory.CompHt.genVal_eq_zero_iff (through
% secZ_eq_zero_of_genVal applied to idem Z * f) and (ii) is GlobalTheory.CompHt.secZ_eq_zero_of_genVal
% (Jacobson rings and the atlas of generic frames GlobalTheory.Reg.evLaws_genFr_atlas). The change of frame
% is GlobalTheory.SecSeparation.exists_frBor with GlobalTheory.FlagTan.evR_frBor, and Artin's theorem is
% Mathlib's linearIndependent_monoidHom.

\begin{lemma}[Alp\"oge]\label{cv-point-conditions-separation}
In Situation~\ref{cv-cusp-setting}, suppose that $f>2$, $p\ge1$ and $q\ge1$, and that the level of $X$ is a level
$(M,\mathfrak M)$ with $M$ prime to $15$. Let $Y$ be a deep cover of $X$ (Definition~\ref{cv-deep-cover}), with deep
level $(M_Y,\mathfrak M_Y)$ and coefficient field $F_Y\subset\mathbf C$. Let $L$ be the lattice of the PEL datum of
$X$, let $\Gamma$ be the group of $\mathcal O_E$-linear automorphisms of $L/\mathfrak M_YL$, and for $\gamma\in\Gamma$
and a framed point $y$ of $Y$ let $y_\gamma$ be its move (Lemma~\ref{cv-realization-pullback}(1)). Let
$\mathfrak P$ be a prime ideal of $\mathcal O_E$ that contains $\mathfrak M_Y$ and whose absolute norm is a prime
number $g$, with $g\notin\mathfrak P^2$ and $p+q+1<g$, and let $S\subset\Gamma$ be a subgroup whose elements induce
unipotent automorphisms of $L/\mathfrak PL$. Let $h$ be the Hodge weight (Definition~\ref{cv-volume-forms-weights}),
let $s\ge p+q+\binom{fpq+1}2+1$ be an integer, and let $Z$ be an irreducible component of $Y\otimes F_Y$. Call a
cusp form of weight $sh$ on $Y$ (Situation~\ref{cv-cusp-setting}, for $Y$ and its uniformizations) \emph{supported
on $Z$} if it equals its product with the function that is $1$ on $Z$ and $0$ on the other components. Then there
is a framed point $y$ of $Y$ over $\mathbf C$, through the inclusion $F_Y\subset\mathbf C$, whose point lies on $Z$,
with the following property: if complex numbers $a_\gamma$, for the $\gamma\in S$ such that the point of $y_\gamma$
lies on $Z$, satisfy $\sum_\gamma a_\gamma v(y_\gamma)=0$ for every cusp form $v$ of weight $sh$ on $Y$ supported on
$Z$, then every $a_\gamma$ is $0$.
\end{lemma}
\leanref{Development/AutomorphicArithmetic/ThetaCorrespondence/ThetaPairLifts/HodgeBundleWeights/ExistsCuspRegularLevelDeterminant.lean}{GlobalTheory.OrbitSep.separate_orbits_by_cusp}
\leanref{Development/AbelianVarietyTheory/Proofs/DivisorProofs/GapGeSubKapBase.lean}{GlobalTheory.OrbitSep.NeatPrime}
\leanref{Development/AbelianVarietyTheory/Proofs/DivisorProofs/GapGeSubKapBase.lean}{GlobalTheory.OrbitSep.UnipAt}
\begin{proof}
\openstep{The proof of this lemma is not yet written out. It compactifies the flag space over $Z$ toroidally,
passes to the quotient by the elements of $S$ that preserve $Z$, to which the Hodge line bundle descends as a nef
and big line bundle, separates the points of that quotient over its generic fibre by sections of the canonical
bundle twisted by a multiple of the Hodge line bundle, turns these sections into cusp forms of weight $sh$ through
a section that does not vanish on the interior, and shows with the unipotence of $S$ modulo $\mathfrak P$ that
distinct elements of $S$ move $y$ to distinct points.}
% Lean: separate_orbits_by_cusp, with the published inputs PublishedHodgeBundleTheorems, GrothendieckDeformation,
% HSCore.Quoted, SiegelToroidalIntegral, ToroidalForFanKS and SemiAbelianLieSeminegative as hypotheses, Y : Deep T j
% with Y.Good e o, hf : 2 < finrank Q E / 2, NeatPrime I g (p + q) (g prime, absNorm I = g, g not in I^2,
% p + q + 1 < g), Y.d.l2.Mb <= I, UnipAt I gamma for gamma in S, and the index set {h in S : LiesOn Z (move h^-1 y)}
% (the text writes gamma = h^-1). Steps: T.covering_topology_refinement; T.torFlag_nonempty_of_one_le_dimFlag (the
% toroidal flag space TF); T.flagSpace_component_smooth_projective_connected and
% T.compactificationComponent_integral_smooth_projective; T.exists_complexFramedPoint_over_componentGenericPoint (y);
% T.componentStabilizer_acts_by_levelChange (the stabilizer Gamma of Z and its action rho);
% T.equivariant_normal_projective_compactification (X', the quotient W, q); T.hodgeLineBundle_rational_linear_equiv_class
% (delta) and T.exists_descent_linearised_lineBundle_finiteGroup_quotient (N, using UnipAt and NeatPrime);
% T.exists_toroidalCanonicalBundle_weight_line (kappa and theta) with K the canonical bundle; m = binom(dim + 1, 2) + 1,
% s = m + t, N' = t delta - kappa, B = m delta + N'; nefness from T.isNefLineBundle_of_pullback_surjective_projective
% and T.base_weight_line_bundle_is_nef; T.exists_cuspRule_of_boundaryVanishing_componentSection (the cusp forms Phi);
% T.sections_separate_quotient_generic_fibre_descends with T.hodgeLineBundle_bigness_on_shimuraVariety (the key
% separation); T.levelAutomorphism_unipotent_fixes_complexPoint_eq_one (distinct h give distinct points). Size: about
% 168,000 lines of Lean below separate_orbits_by_cusp that the book does not yet cite; with Lemma
% cv-point-conditions-spanning about 343,000 lines together.
\end{proof}

\begin{lemma}[Alp\"oge]\label{cv-point-conditions-spanning}
In the setting of Proposition~\ref{cv-point-conditions}, let $P$ be a point of the family at a scale $x$, in the
notation of Situation~\ref{cv-point}, and let $1\le j\le m$. Then the integral cusp forms of weight $\kappa_j$ on the
PEL moduli scheme of $U(W_j)$ span, over its coefficient field $F_c$, the $F_c$-space $V_j^{\mathrm{aut}}$ of all its
cusp forms of weight $\kappa_j$ (Situation~\ref{cv-cusp-setting}).
\end{lemma}
\leanref{Development/AutomorphicArithmetic/ThetaCorrespondence/ThetaPairLifts/ThirdPlaceSeries/GoodWAt.lean}{GlobalTheory.CovTop.goodW_at_alt}
\leanref{Development/AbelianVarietyTheory/Proofs/DivisorProofs/ExistsSmulNonnegOfPow.lean}{GlobalTheory.SupFull.fullAt_of_pins}
\begin{proof}
\openstep{The proof of this lemma is not yet written out. For a cusp form $f$ of weight $\kappa_j$, it writes a power
of $f$, a bounded section of large weight, as a combination of short theta series of a level datum, shows that
each of these has a nonzero integral multiple, and concludes that a nonzero multiple of $f$ is integral.}
% Lean: the third condition of GoodW (span of latW = top) in goodW_at_alt, which is GoodWCore.goodW_core applied to
% GlobalTheory.GoodLat.Obj.latFull = fullLat, the second part of SupFull.fullAt_of_pins: pow_mem_cusp (the power
% f^(4 f r + 4)), exists_datum_span (a DeltaDat d at level l with short generators shortGen spanning the bounded
% sections, via exists_unf, isBddSec_of_isUnf and holSub_SS), exists_smul_short_nonneg (a nonzero integral multiple
% of each short generator) and exists_smul_nonneg_of_pow (from the power to f). Size: about 241,000 lines of Lean below
% these that the book does not yet cite, about 175,000 of them not used by Lemma cv-point-conditions-separation.
\end{proof}

\begin{proposition}\label{cv-point-conditions}
For every $\varepsilon$ and every family as in Theorem~\ref{cv-covolume-bound}, there are
constants $A,B,C_1,C_2\ge0$ and $x_0$ such that at every point $P$ at a scale $x\ge x_0$,
in the notation of Situation~\ref{cv-point}, the following hold for every $j$. Let $V_j$
be the $L$-vector space of all $j$-th lifted forms, with its Hermitian norm system $N_j$
given by the Petersson norms, let $\Lambda_j^{\mathrm{all}}\subset V_j$ be the lattice
spanned by the integral lifted forms, let $V_j^{\mathrm{type}}\subset V_j$ be the subspace of
lifted forms of the prescribed type, and let
$\Lambda_j=\Lambda_j^{\mathrm{all}}\cap V_j^{\mathrm{type}}$. Let $G$ be the finite level group,
identified with the group $G$ of Definition~\ref{cv-separated}, let $\rho_j$ be its action on
$V_j$ and $\chi\colon G\to L^\times$ the character of the rational structure, let $e_j$ be the
action on $V_j$ of the separated idempotent $e_{\mathrm{sep}}$ through the representation
$g\mapsto\chi(g)^{-1}\rho_j(g)$, and let $\Lambda_j^{\mathrm{sep}}=\Lambda_j\cap e_jV_j$.
Let $\Lambda_j^{\mathrm{aut}}$ be the lattice of integral cusp forms of weight $\kappa_j$
and the prescribed type on the PEL unitary Shimura variety of $U(W_j)$, with its Petersson
norm, and put $\mu_j^{\mathrm{aut}}=\mu(\Lambda_j^{\mathrm{aut}})$. Let
$K_j=PP\cdot\sum_{\tau,i}|\kappa_{j,\tau,i}|$, where $PP$ is the size parameter of the
family.
\begin{enumerate}
\item (Recursion.) $\frac1D\operatorname{lcv}_{X_j}(M_j)\le-\mu(\Lambda_j^{\mathrm{sep}})$.
\item (Cut.) $e_j$ is an orthogonal idempotent, $|G|\,e_j(\Lambda_j)\subset\Lambda_j$ and
  $2\log|G|\le b$.
\item (Bad set.) $\operatorname{rk}(\Lambda_j\cap\ker e_j)\le3e^{-x^{\varepsilon/(7J)}}\operatorname{rk}\Lambda_j$.
\item (Size.) $\mu_{\max}(\Lambda_j^{\mathrm{all}})\le x^Ab$, $1\le b\le u\le x^Bb$,
  $|\mathrm{main}_j|\le x^Bb$, $|\mathrm{main}_j|\le K_j$ and $b\le K_j$.
\item (Comparison with automorphic forms.) $\mu(\Lambda_j)\ge\mu_j^{\mathrm{aut}}-C_1u$.
\item (Arithmetic volume.) $\mu_j^{\mathrm{aut}}\ge-\mathrm{main}_j-C_2u$.
\end{enumerate}
\end{proposition}
\leanref{Development/AbelianVarietyTheory/Definitions/Divisors/CovolXOfCovol.lean}{GlobalTheory.CovTop.RecAt}
\leanref{Development/AbelianVarietyTheory/Definitions/Divisors/CovolXOfCovol.lean}{GlobalTheory.CovTop.CutAt}
\leanref{Development/AbelianVarietyTheory/Definitions/Divisors/CovolXOfCovol.lean}{GlobalTheory.CovTop.BadAt}
\leanref{Development/AbelianVarietyTheory/Definitions/Divisors/CovolXOfCovol.lean}{GlobalTheory.CovTop.SizeAt}
\leanref{Development/AbelianVarietyTheory/Definitions/Divisors/CovolXOfCovol.lean}{GlobalTheory.CovTop.CoreAt}
\leanref{Development/AbelianVarietyTheory/Definitions/Divisors/CovolXOfCovol.lean}{GlobalTheory.CovTop.LowAt}
\leanref{Development/AutomorphicArithmetic/ThetaCorrespondence/ThetaPairLifts/ThirdPlaceSeries/TopCovolG.lean}{SelfIntersectionVolumeEstimates.topCovolG_alt_alt}
\begin{proof}
Each condition holds at all points at scales $x\ge x_0$ for a threshold $x_0$ and constants
of its own, and one takes the largest. Write $r=k+m-1$ for the dimension of $W_j$, $f$ for the
degree of the totally real host field, and $x^\varepsilon$ for the logarithm of $b$. The
family's size conditions give a constant $A'$ with $fr^2x^{\varepsilon}\le x^{A'}$ and
$x^{3A'}\le b$, and $\log g_k\le9x^\varepsilon+1$ for the level primes $g_0,g_1$.

(1) Recursion. The lattice $M_j$ is $\Lambda_j^{\mathrm{sep}}$ written in a basis of the
$L$-space $V_j^{\mathrm{type}}\cap e_jV_j$, which contains it, the metric $X_j$ is $s=|P_0|^{-1/2}\le1$ times the norm system $N_j$, where $P_0$ is the
level subgroup of $G$, and the point conditions give $d_j=D$. Lemma~\ref{cv-rescaling} gives
(1).

(2) Cut. The finite level group is $G=\prod_{k=0,1}\prod_{v\mid g_k}GL_r(\mathbf F_{g_k})$,
over the $f$ places $v$ of the totally real host field above each $g_k$. Let $\mathcal X$ be the
set of separated characters, which are absolutely irreducible characters of $G$ with values in $L$,
put $e_\eta=\frac{\eta(1)}{|G|}\sum_{t\in G}\eta(t^{-1})t$ for $\eta\in\mathcal X$, and let
$e=\sum_{\eta\in\mathcal X}e_\eta=\sum_gc(g)g$ in $L[G]$, which is $e_{\mathrm{sep}}$; then
$e_j=\sum_gc(g)\chi(g)^{-1}\rho_j(g)$.
\begin{enumerate}
\item[(a)] The orthogonality relations for the distinct irreducible characters in $\mathcal X$
  give $e^2=e$, that is $\sum_gc(g)c(g^{-1}k)=c(k)$. The coefficients are class functions,
  $|G|c(g)=\sum_\eta\eta(1)\eta(g^{-1})$ is an algebraic integer, and
  $\sigma(c(g^{-1}))=\overline{\sigma(c(g))}$ for every embedding $\sigma$ of $L$, because
  $\eta(g^{-1})=\overline{\eta(g)}$ for characters of finite groups.
\item[(b)] Since $g\mapsto\chi(g)^{-1}\rho_j(g)$ is a homomorphism, $e_j\circ e_j=e_j$.
\item[(c)] The action $\rho_j$ preserves the Hermitian form of $N_j$ at every embedding
  $\sigma$, and $|\sigma(\chi(g))|=1$ because $\chi(g)$ is a root of unity. By the last property
  in (a), the adjoint of $e_j$ at $\sigma$ is
  $\sum_g\overline{\sigma(c(g)\chi(g)^{-1})}\rho_j(g^{-1})
  =\sum_g\sigma(c(g^{-1})\chi(g^{-1})^{-1})\rho_j(g^{-1})=e_j$. A self-adjoint
  idempotent has orthogonal image and kernel, which is the identity required of an orthogonal
  idempotent in Lemma~\ref{cv-cut}.
\item[(d)] The lattice $\Lambda_j$ consists of the elements of $\Lambda_j^{\mathrm{all}}$ on which
  the level subgroup acts through a given character. The action $\rho_j$ preserves integral forms,
  and $e_j$ commutes with $\rho_j$ because $c$ is a class function, so it preserves the type. As
  $|G|c(g)\in\mathcal O_L$ and $\chi(g)^{-1}$ is a unit of $\mathcal O_L$,
  $|G|e_j(x)=\sum_g|G|c(g)\chi(g)^{-1}\rho_j(g)x$ lies in $\Lambda_j$ for $x\in\Lambda_j$.
\item[(e)] Since $|GL_r(\mathbf F_g)|\le g^{r^2}$, $\log|G|\le2fr^2\max_k\log g_k\le2fr^2(9x^\varepsilon+1)$,
  and $4fr^2(9x^\varepsilon+1)\le b$ by the size conditions.
\end{enumerate}

(3) Bad set. Let $P$, $\lambda$ and $H_0,H_1$ be as in Lemma~\ref{cv-bad-part}. The ranks are
realized as dimensions: there are a finite group $\Gamma$ with a surjection $\pi:\Gamma\to G$ and a
section $\iota$ as in Lemma~\ref{cv-bad-part}, a representation $\rho$ of $\Gamma$ on a
finite-dimensional complex vector space $V$, and nontrivial additive characters $\psi_{k,v}$ defining
$\lambda$, such that
\[
 \operatorname{rk}\Lambda_j=\dim V^{\tilde\lambda},\qquad
 \operatorname{rk}(\Lambda_j\cap\ker e_j)=\dim\big(V^{\tilde\lambda}\cap\ker\rho(\iota(e_{\mathrm{sep}}))\big).
\]
For $j=m$, moreover, for each $k$ the group $\Gamma_k=\pi^{-1}(H_k)$ acts freely on a finite set
$Y_k$ and permutes the summands of a decomposition $V=\bigoplus_{y\in Y_k}A_y$, with
$\rho(\gamma)A_y=A_{\gamma y}$. For $j<m$, for each $k$ the restriction of $V$ to $\Gamma_k$ embeds into
the sum of $n_k$ copies of the regular representation with
$n_k|\Gamma_k|\le\frac32|\Gamma_P|\dim V^{\tilde\lambda}$.

Ranks and types. Let $F_c$ and $V_j^{\mathrm{aut}}$ be as in (5), and take the rational lifted
forms $V_j^{\mathrm{rat}}\subset V_j$, the lifting isomorphism
$e\colon V_j^{\mathrm{aut}}\to V_j^{\mathrm{rat}}$, the level subgroup $P_0=P$ with its character
$\psi_0$, and the embedding $\sigma_0\colon L\to\mathbf C$ of (5a). The level automorphisms give an
action of $G$ on $V_j^{\mathrm{aut}}$, written $f\mapsto gf$, with $\rho_j(g)e(f)=e(gf)$. A lifted form
$w$ has the prescribed type if $\rho_j(u)w=\chi(u)\psi_0(u)^{-1}w$ for $u\in P$; these forms make up
$V_j^{\mathrm{type}}$. The character
$\sigma_0\circ\psi_0^{-1}$ of $P$ is $\lambda$ for the nontrivial additive characters
$\psi_{k,v}(a)=\psi_v(-a)$. With $c$ as in (2), $e_j=\sum_gc(g)\chi(g)^{-1}\rho_j(g)$ and
$e_{\mathrm{sep}}=\sum_g\sigma_0(c(g))\,g$ in $\mathbf C[G]$. The integral lifted forms span $V_j$ over
$L$: by (5a) they are the lifts of the integral cusp forms, the integral cusp forms span
$V_j^{\mathrm{aut}}$ over $F_c$ by Lemma~\ref{cv-point-conditions-spanning}, and the lifts of a basis of $V_j^{\mathrm{aut}}$ form a basis of $V_j$.
So for $w\in V_j^{\mathrm{type}}$ there is a nonzero $d\in\mathcal O_L$ with $dw\in\Lambda_j^{\mathrm{all}}$,
and then $dw\in\Lambda_j$. Hence $\Lambda_j$ spans $V_j^{\mathrm{type}}$ and $\Lambda_j\cap\ker e_j$ spans
$V_j^{\mathrm{type}}\cap\ker e_j$ over $L$, and
\[
 \operatorname{rk}\Lambda_j=\dim_LV_j^{\mathrm{type}},\qquad
 \operatorname{rk}(\Lambda_j\cap\ker e_j)=\dim_L(V_j^{\mathrm{type}}\cap\ker e_j).
\]
% Lean: GlobalTheory.IntFormsRank.rkK_Mall with the spanning GlobalTheory.CovTop.SpanMK.span_MK_eq_top_alt;
% e_j is GlobalTheory.CovTop.eGood; the Whittaker form of the type character is
% GlobalTheory.ShareOn.isWhit_and_eq_embLL_inv_and_halfUnip; e_sep is read from the point field grpRecPin.

The case $j=m$. Take $\Gamma=G$ and $\pi=\iota=\mathrm{id}$, so that $\Gamma_P=P$ and
$\tilde\lambda=\lambda$. Let $V=\mathbf C\otimes_LV_m$, with $\mathbf C$ an $L$-algebra through
$\sigma_0$, and $\rho(g)=\sigma_0(\chi(g))^{-1}(1\otimes\rho_m(g))$.
Lemma~\ref{cv-realization-linear}(1), with $W=V_m$, $\rho_0=\rho_m$, $\lambda_0=\psi_0^{-1}$ and $c$,
so that the operator $e$ of the lemma is $e_m$ and $e'=e_{\mathrm{sep}}$, gives
$\dim_LV_m^{\mathrm{type}}=\dim V^\lambda$ and
$\dim_L(V_m^{\mathrm{type}}\cap\ker e_m)=\dim(V^\lambda\cap\ker\rho(e_{\mathrm{sep}}))$, which with the
ranks above is the realization. For the decompositions, let $X$ be the Shimura variety of the definite
group $U(W_m)$ over $F_c$, which is as in Situation~\ref{cv-definite-situation} with $p=r$, let $L$ be the
lattice of its PEL datum, and let $\gamma_g$, for $g\in G$, be the matrices of the level automorphisms
(Definition~\ref{cv-realization-level-move}) through which $G$ acts on $V_m^{\mathrm{aut}}$; they satisfy
the hypotheses of Lemma~\ref{cv-realization-definite}, and every form of weight $\kappa_m$ is a cusp form.
That lemma gives an $F_c$-algebra $B$, a finite product of finite separable extensions of $F_c$, with
$X=\operatorname{Spec}B$; the structure of a $B$-module on $V_m^{\mathrm{aut}}$, by multiplication with
functions; and $F_c$-algebra automorphisms $\tau(g)$ of $B$ with $g(bf)=\tau(g)(b)\,gf$. The group $H_k$
acts freely on the set $Y_k$ of $F_c$-algebra homomorphisms $B\to\mathbf C$, with $\mathbf C$ an
$F_c$-algebra through $\sigma_0$, by Lemma~\ref{cv-realization-definite}(4). Indeed, the level of $X$ is
$(g_0g_1,\mathfrak M)$ with $\mathfrak M=\prod_v\mathfrak P_v$, the product over the level places $v$ of primes
$\mathfrak P_v$ of $E$ above $v$ with residue fields $\mathbf F_{g_k}$ for $v\mid g_k$, and $\gamma_g$ induces
on $L/\mathfrak P_vL$ the automorphism with matrix $g_v$ in a basis over $\mathbf F_{g_k}$. Let $h\in H_k$,
and let $y$ be a framed point over $\mathbf C$ that has the same point as its move $y_h$; then a level
automorphism with matrix $\gamma_h$ fixes that point. Let $v$ be a level place over $g_{1-k}$. As $h_v$ is
upper unitriangular, $(\gamma_h-1)^rL\subset\mathfrak P_vL$. Moreover $g_0g_1\mathcal O_E=\mathfrak P_vJ$ with
$J=\bar{\mathfrak M}\prod_{v'\ne v}\mathfrak P_{v'}$, which is prime to $\mathfrak P_v$, and $r+1<g_{1-k}$ by
the size conditions. Lemma~\ref{cv-realization-fixed} gives $(\gamma_h-1)L\subset\mathfrak ML$. So $\gamma_h$
induces the identity on every $L/\mathfrak P_{v'}L$, $h_{v'}=1$ at every level place $v'$, and $h=1$.
% Lean: Lemma cv-realization-definite is applied in GlobalTheory.DefSplit.exists_free_action_permuting_summands_thirdPlace
% through split_rules, with the action (p.pins i).2.act.1 and mem_cusp_of_sec for the cusp forms; the
% freeness hypothesis is GlobalTheory.NeatOne.free_half (via neat_divisor_free_half_lifts_alt), which
% applies free_of_unip_one with the cofactor of GlobalTheory.NeatOne.exists_cofactor, the unipotence
% GlobalTheory.NeatOne.unip_mod_place and the faithfulness GlobalTheory.NeatOne.eq_one_of_lev_cong; r + 1 < g_k
% is GlobalTheory.NeatOne.rank_add_one_lt.
Lemma~\ref{cv-realization-split}, with $K=F_c$, $H=H_k$ and
$T(h)f=hf$, decomposes $\mathbf C\otimes_{F_c}V_m^{\mathrm{aut}}$ into summands $A'_y$, $y\in Y_k$, with
$(1\otimes T(h))A'_y\subset A'_{h\cdot y}$. Since $e$ carries a basis of $V_m^{\mathrm{aut}}$ over $F_c$ to
a basis of $V_m$ over $L$, the map $z\otimes f\mapsto z\otimes e(f)$ is an isomorphism
$\mathbf C\otimes_{F_c}V_m^{\mathrm{aut}}\to V$ that carries $1\otimes T(g)$ to $1\otimes\rho_m(g)$. Let
$A_y$ be the image of $A'_y$. As $\rho(g)$ is a multiple of $1\otimes\rho_m(g)$,
$V=\bigoplus_{y\in Y_k}A_y$ and $\rho(\gamma)A_y\subset A_{\gamma y}$ for $\gamma\in\Gamma_k=H_k$; applying
this to $\gamma^{-1}$ gives equality.
% Lean: GlobalTheory.DefRead.point_read_alt (V is VAt, rho is repAt), with
% GlobalTheory.T.finrank_characterEigenspace_twistedComplexification; the decomposition is
% GlobalTheory.DefHalf.halfH_acts_freely_on_internal_sum_decomposition, from
% GlobalTheory.DefSplit.split_rules (through split_complex) and
% GlobalTheory.T.transport_permuted_directSum_complexification.

The case $j<m$. Let $Y$ be the deep cover of the Shimura variety of $U(W_j)$ in the proof of
Proposition~\ref{cv-arithmetic-volume-step}, with coefficient field $F_Y\supset F_c$ and an embedding
$F_Y\to\mathbf C$ that extends $\sigma_0$ on $F_c$. Let $L$ be the lattice of the PEL datum of the
Shimura variety $X$ of $U(W_j)$. As in the case $j=m$, the level of $X$ is $(g_0g_1,\mathfrak M)$ with
$\mathfrak M=\prod_v\mathfrak P_v$, and $\gamma_g$ induces $g_v$ on $L/\mathfrak P_vL$ in a basis
$\varphi_v$ over $\mathbf F_{g_k}$, for the matrices $\gamma_g$ of the level automorphisms through which $G$
acts on $V_j^{\mathrm{aut}}$. Let $(15^eg_0g_1,\mathfrak MI)$ be the deep level of $Y$, and $\Gamma$ the
group of $\mathcal O_E$-linear automorphisms of $L/\mathfrak MIL$. By
Lemma~\ref{cv-realization-pullback}(1), every $\gamma\in\Gamma$ moves the framed points of $Y$, and it acts on
the $F_Y$-space $V_Y$ of cusp forms of weight $\kappa_j$ on $Y$ by the translation
$(\tau(\gamma)f)(y)=f(y_\gamma)$. Lemma~\ref{cv-realization-levels}, with the level places $v$, the ideal $I$
and the bases $\varphi_v$, gives a homomorphism $\pi\colon\Gamma\to G$ whose kernel is the group $\Gamma_1$
of Lemma~\ref{cv-realization-pullback}, a homomorphism $\iota\colon G\to\Gamma$ with
$\pi\circ\iota=\mathrm{id}$ whose image commutes with $\ker\pi$, and
$(\tilde\gamma-\gamma_{\pi(\gamma)})L\subset\mathfrak ML$ whenever $\tilde\gamma$ induces $\gamma$. For
$f\in V_j^{\mathrm{aut}}$ let $\beta(f)$ be its pullback to $Y$, the form $y\mapsto f(p(y))$ for the level map
$p$. By (Y7) in the proof of Proposition~\ref{cv-arithmetic-volume-step}, $\beta$ is an additive map
$V_j^{\mathrm{aut}}\to V_Y$ with $\beta(af)=a\beta(f)$ for $a\in F_c$, which carries $F_c$-independent
families to $F_Y$-independent ones. By Lemma~\ref{cv-realization-pullback}(2), with $\gamma'=\gamma_{\pi(\gamma)}$,
\[
 (\tau(\gamma)\beta(f))(y)=f(p(y_\gamma))=(\pi(\gamma)f)(p(y))=\beta(\pi(\gamma)f)(y)
\]
at every framed point $y$ of $Y$ over an algebraically closed field, so
$\tau(\gamma)\beta(f)=\beta(\pi(\gamma)f)$. A form $f$ on $Y$ fixed by $\ker\pi$ takes equal values at framed
points $y$ and $y'$ of $Y$ with $p(y)$ and $p(y')$ isomorphic: by Lemma~\ref{cv-realization-pullback}(3),
$f(y')=f(y_\gamma)=(\tau(\gamma)f)(y)=f(y)$ for some $\gamma\in\Gamma_1$. In particular such a form is
constant on fibres in the sense of Lemma~\ref{cv-cover-pullback}, so by part (4) of that lemma, with
$\kappa=\kappa_j$ and $S_\kappa=V_j^{\mathrm{aut}}$, whose map of pullbacks is $\beta$, it lies in the
$F_Y$-span of $\beta(V_j^{\mathrm{aut}})$. Let $V=\mathbf C\otimes_{F_Y}V_Y$ with
$\rho(\gamma)=\sigma_0(\chi(\pi(\gamma)))^{-1}(1\otimes\tau(\gamma))$. Lemma~\ref{cv-realization-linear}(2),
with $K=F_c$, $F=F_Y$, $M=\mathbf C$, $W=V_j$, $\rho_0=\rho_j$, $\lambda_0=\psi_0^{-1}$,
$C=V_j^{\mathrm{aut}}$ with a basis that $e$ carries to a basis of $V_j$, $\varepsilon$ the lifting
isomorphism followed by the inclusion of $V_j^{\mathrm{rat}}$ in $V_j$, $f'=\pi(\gamma)f$ and the map
$\beta$, gives the realization; the operator $e$ of the lemma is $e_j$, and $e'=e_{\mathrm{sep}}$. The
map $\pi$ is onto because $\pi\circ\iota=\mathrm{id}$.
% Lean: GlobalTheory.CovTop.BadRep.badRep_at, through rank_cover, rep_data and rank_gen; pi and iota are
% c o t o typeRed fr and typeEmb fr o t^-1 o c^-1 from GlobalTheory.CovTop.BadRep.frame_data, whose congruence
% CongMod l.Mb (liftLa gamma) (lev (t (typeRed fr gamma))) is Lemma cv-realization-levels(3); beta and its
% properties are GlobalTheory.CovTop.BadRep.pull_data (levelCover_cuspForm_pullback_injective_descent,
% GlobalTheory.LevPullEquiv.pull_actLev, GlobalTheory.T.fibreConstant_iff_coveringGroup_invariant).

For the regular representations, let $R$ be the ring of sections of all weights on $Y$, an
$F_Y$-algebra on which $\Gamma$ acts by translation, and for a weight $\kappa$ let
$S_Y(\kappa)\subset R$ be the space of cusp forms of weight $\kappa$, so that $V_Y=S_Y(\kappa_j)$. Let
$h$ be the Hodge weight and $c_0$ the integer of the proof of
Proposition~\ref{cv-arithmetic-volume-step}. Products of cusp forms of weight $c_0h$ with cusp forms
of weight $\kappa_j-c_0h$, $\kappa_j$ or $\kappa_j+c_0h$ are cusp forms of the sum of the weights;
translation preserves $S_Y(c_0h)$, $S_Y(\kappa_j)$ and $S_Y(\kappa_j+c_0h)$; and for $s=-1,0,2$ the space
$S_Y(\kappa_j+sc_0h)$ is finite-dimensional, of dimension the rank $\operatorname{rk}(\kappa_j+sc_0h)$
of its integral cusp forms, by (Y1) in the proof of
Proposition~\ref{cv-arithmetic-volume-step}. There are $u_0,u\in S_Y(c_0h)$ whose group determinants on
$\Gamma_P$ and on $\Gamma_k$, respectively, are not zero divisors in $R$. Such a form $u_S$ is found as
follows for $S=\Gamma_P$ and for $S=\Gamma_k$.
\begin{enumerate}
\item[(i)] (Separation.) For every irreducible component $Z$ of $Y\otimes F_Y$ there is a complex framed
  point $y_Z$ of $Y$, over the inclusion $F_Y\subset\mathbf C$, whose point lies on $Z$, with this property:
  if complex numbers $a_\gamma$, for the $\gamma\in S$ such that the point of $(y_Z)_\gamma$ lies on $Z$,
  satisfy $\sum_\gamma a_\gamma\,v((y_Z)_\gamma)=0$ for every $v\in S_Y(c_0h)$ supported on $Z$, then every
  $a_\gamma$ is $0$. This is Lemma~\ref{cv-point-conditions-separation} with $s=c_0$, which is at least
  $p+q+\binom{fpq+1}2+1$ by the definition of $c_0$, and with the following prime. If $S=\Gamma_P$, let $v$ be a
  level place over $g_0$; if $S=\Gamma_k$, let $v$ be a level place over $g_{1-k}$; and let $\mathfrak P=\mathfrak P_v$.
  Its absolute norm is the prime $g$ over which $v$ lies, as its residue field is $\mathbf F_g$; $g\notin\mathfrak P^2$,
  as $g_0g_1\mathcal O_E=\mathfrak PJ$ with $J$ prime to $\mathfrak P$, as in the case $j=m$; $r+1<g$ by the size
  conditions; and $\mathfrak MI\subset\mathfrak P$. For $\gamma\in S$, $\pi(\gamma)$ is upper unitriangular at $v$,
  by the definitions of $P$ and $H_k$ in Lemma~\ref{cv-bad-part}; so the automorphism that $\gamma$ induces on
  $L/\mathfrak PL$, whose matrix in the basis $\varphi_v$ is $\pi(\gamma)_v$, is unipotent. The other hypotheses of
  the lemma, $f>2$, $p\ge1$ and $q\ge1$, hold at all large scales by the size conditions of the family.
\item[(ii)] (A common form.) Apply Lemma~\ref{cv-point-group-determinant} with $K=F_Y$, $V=S_Y(c_0h)$, the
  group $S$, the finite set $I$ of components, $\mathrm{ev}_Z(v)(\gamma)=v((y_Z)_\gamma)$, $J_Z$ the set of
  $\gamma\in S$ such that the point of $(y_Z)_\gamma$ lies on $Z$, with its inclusion $r_Z$ in $S$, and
  $W_Z$ the subspace of forms supported on $Z$. The set $J_Z$ contains $1$. A form supported on $Z$ vanishes
  at the framed points whose point does not lie on $Z$, which is (a), and (b) is (i). So there is
  $u_S\in S_Y(c_0h)$ with $\det\big(u_S((y_Z)_{\gamma\gamma'})\big)_{\gamma,\gamma'\in S}\ne0$ for every $Z$.
\item[(iii)] (Regularity.) The group determinant $\Delta=\det\big(\tau(\gamma\gamma')u_S\big)_{\gamma,\gamma'\in S}$
  is a sum of products of $|S|$ translates of $u_S$. As a section of the weight $c_0h$, $u_S$ has an
  exponent in the sense of Lemma~\ref{cv-point-regular-section}; translations do not change the frame of a
  framed point; and values at a framed point are multiplicative. So $\Delta$ has an exponent. As $(\tau(\gamma)f)(y)=f(y_\gamma)$,
  the value of $\Delta$ at $y_Z$ is $\det\big(u_S((y_Z)_{\gamma\gamma'})\big)\ne0$, for every $Z$. By
  Lemma~\ref{cv-point-regular-section}, $\Delta$ is not a zero divisor in $R$.
\end{enumerate}
Put $u_0=u_{\Gamma_P}$ and $u=u_{\Gamma_k}$.
% Lean: GlobalTheory.OrbitSep.exists_cusp_regular_level_determinant (= GlobalTheory.T.orbits_separated_by_regular_cusp)
% for S = grpHR fr and S = grpGHR fr k: (i) is separate_orbits_by_cusp with s = cZeroN D j (which is at least
% p + q + binom(dim + 1, 2) + 1 by cZeroN_isLeast), indexed by move h^-1 for h in S; (ii) is
% GlobalTheory.T.exists_vector_groupDeterminants_ne_zero with W Z = suppOn Z and AllWt.liesOff_iff for (a);
% (iii) is LevDesc.levDet_mem_secZ, GlobalTheory.LevelKeeps.evR_transl and RingHom.map_det at y Z, then
% GlobalTheory.T.integerExponentSection_mul_eq_zero_imp_eq_zero.
Apply Lemma~\ref{cv-regular-group-determinant} with
$K=F_Y$, these $u_0$ and $u$, $S_0=\Gamma_P$, $S=\Gamma_k$, $V_-=S_Y(\kappa_j-c_0h)$, $V=S_Y(\kappa_j)$,
$V_1=S_Y(\kappa_j+c_0h)$, $V_2=S_Y(\kappa_j+2c_0h)$, $\nu(\gamma)=\sigma_0(\chi(\pi(\gamma)))^{-1}$, the
character $\tilde\lambda$ and $t=\frac32$. The inequality $\dim V_2\le\frac32\dim V_-$ is the rank
comparison in the proof of Proposition~\ref{cv-arithmetic-volume-step}(2) for $s=2$:
$\operatorname{rk}(\kappa_j+2c_0h)\le(1+\min(\frac12,b/K_j))\operatorname{rk}(\kappa_j-c_0h)
\le\frac32\operatorname{rk}(\kappa_j-c_0h)$. The lemma gives the embeddings, with
$n_k|\Gamma_k|\le\frac32|\Gamma_P|\dim V^{\tilde\lambda}$.
% Lean: GlobalTheory.BadEmb.regular_embeddings_cover_bad_embedding_locus_alt, through embReg_deep,
% embReg_of_cusp and embReg_of_submodules; the products are GlobalTheory.SlopeRec.mul_mem_cuspW, the
% translations GlobalTheory.LevelIsometry.keeps_cusp_lattice_norm, the dimensions
% GlobalTheory.BadEmb.finrank_of_goodAt;
% the rank comparison is GlobalTheory.T.rank_growth_bound_small_scale_alt_alt with s = 2, and
% kappa_j + 2 c_0 h is good by GlobalTheory.T.covering_shift_preserves_goodAt.
% Lean: the descent for j < m is the part "span iff FibConst" of levRing_descent, which is
% Lemma cv-cover-pullback(4). Separation (i) is Lemma cv-point-conditions-separation
% (OrbitSep.separate_orbits_by_cusp), whose prime and unipotence come from
% GlobalTheory.T.exists_neatPrime_levelAutSubgroup_unipotent = OrbitSep.exists_neatPrime_unipAt: exists_index (j = 0
% for grpHR fr, j = k + 1 for grpGHR fr k), exists_plOver (the place v), neat_typePrime (the prime typePrime l v,
% through absNorm_eq_of_split, GSplit.g_splits and hg) and unipAt_of_unitri (isNilpotent_sub_one_of_unitri); hf, hp
% and hq come from T.weight_conditions_at_index' in regular_embeddings_cover_bad_embedding_locus_alt. Spanning is Lemma
% cv-point-conditions-spanning, the last condition of GoodW in GlobalTheory.CovTop.GramAt. On Y,
% the finite dimension, finite generation and spanning for kappa_j + s c_0 h (s = -1, 0, 2) are (Y1) in the
% proof of Proposition cv-arithmetic-volume-step: Deep.GoodAt from DeepGoodness for kappa_j and
% kappa_j - c_0 h, and GlobalTheory.T.covering_shift_preserves_goodAt for kappa_j + 2 c_0 h.

For $j=m$ the hypothesis of Lemma~\ref{cv-bad-part} holds with $c=1$. Choose representatives
$y_1,\dots,y_s$ of the $\Gamma_k$-orbits on $Y_k$. Since the action is free,
$V=\bigoplus_i\bigoplus_{\gamma\in\Gamma_k}\rho(\gamma)A_{y_i}$, so $V$ is the sum of
$n_k=\sum_i\dim A_{y_i}$ copies of the regular representation of $\Gamma_k$ and
$n_k|\Gamma_k|=\dim V$. The subgroup $\Gamma_P\subset\Gamma_k$ also acts freely, so $V$ restricted to
$\Gamma_P$ is a sum of $\dim V/|\Gamma_P|$ copies of its regular representation, in which every
character occurs once; hence $|\Gamma_P|\dim V^{\tilde\lambda}=\dim V=n_k|\Gamma_k|$.

In both cases Lemma~\ref{cv-bad-part}, with $\psi_k(a)=e^{2\pi ia/g_k}$, and
Lemma~\ref{cv-local-share} give, at all large scales,
\[
 \operatorname{rk}(\Lambda_j\cap\ker e_j)\le\tfrac32\big(s_0(S_0)+s_1(S_1)\big)\operatorname{rk}\Lambda_j
 \le\tfrac32e^{-x^{\varepsilon/(7J)}}\operatorname{rk}\Lambda_j.
\]

(4) Size. The bounds $1\le b\le u\le x^Bb$, $|\mathrm{main}_j|\le x^Bb$,
$|\mathrm{main}_j|\le K_j$ and $b\le K_j$ follow from the size conditions of the family and the
shape of the weights $\kappa_j$, since $u$, $\mathrm{main}_j$ and $K_j$ are explicit in the ranks,
the weights and $x$. For the maximal slope, let $\Lambda_j^{\mathrm{aut,all}}$ be the lattice of
all integral cusp forms of weight $\kappa_j$ on the Shimura variety of $U(W_j)$. By (5c) below
and Lemma~\ref{cv-slope-base-change}(2),
$\mu_{\max}(\Lambda_j^{\mathrm{all}})\le\mu_{\max}(\Lambda_j^{\mathrm{aut,all}})+2C_Mu$. Three facts
bound $\mu_{\max}(\Lambda_j^{\mathrm{aut,all}})$. In the notation of Situation~\ref{cv-cusp-setting},
$\Lambda_j^{\mathrm{aut,all}}=\Lambda_{\kappa_j}$. First, by Proposition~\ref{cv-cusp-norm-bound}, every
nonzero $f\in\Lambda_j^{\mathrm{aut,all}}$ satisfies
$\frac1{[F_c:\mathbf Q]}\sum_\sigma\log\|f\|_\sigma\ge-\lambda_j$, where $F_c$ is the coefficient field
of the Shimura variety and $\lambda_j=\eta_{\kappa_j}$ is explicit in $\log\Omega_E(\chi)$, the size of
$\kappa_j$, the rank, the prime $d$ and the level. Second, the rank of $\Lambda_j^{\mathrm{aut,all}}$ is
at most $\dim_{F_c}S_{\kappa_j}$, which is at most $C_j=C_{w_{\kappa_j}}$, explicit in the same
quantities, by Proposition~\ref{cv-cusp-rank-bound}.
Third, Lemma~\ref{cv-minkowski-slope} applies to $\Lambda_j^{\mathrm{aut,all}}$ with $K=F_c$ and
$\lambda=\lambda_j$; Petersson norms satisfy $N_{\bar\sigma}=N_\sigma$. From these,
$\mu_{\max}(\Lambda_j^{\mathrm{aut,all}})\le2\lambda'_j+2\log C_j+\log|d_{F_c}|/[F_c:\mathbf Q]$ for a
number $\lambda'_j$ explicit in the same quantities, and the
size conditions of the family bound the right side by $x^{A_1}b$ for a constant $A_1$ of the family.
As $u\le x^{B}b$, this gives $\mu_{\max}(\Lambda_j^{\mathrm{all}})\le x^Ab$ for a suitable $A$.

(5) Comparison. Let $F_c\subset L$ be the coefficient field of the Shimura variety of $U(W_j)$
and $V_j^{\mathrm{aut}}$ the $F_c$-space of its cusp forms of weight $\kappa_j$.
\begin{enumerate}
\item[(a)] Lifting forms defines an additive map $\iota:V_j^{\mathrm{aut}}\to V_j$ with
  $\iota(af)=a\,\iota(f)$ for $a\in F_c$, which carries $F_c$-independent families to
  $L$-independent ones and maps $\Lambda_j^{\mathrm{aut,all}}$ onto the integral lifted forms.
  The realization at the point and its weight tie (Definitions~\ref{is-realisation}
  and~\ref{is-one-map-data}) provide the following data. First, an $F_c$-subspace
  $V_j^{\mathrm{rat}}\subset V_j$ of \emph{rational lifted forms}, with a finite $F_c$-basis that is an
  $L$-basis of $V_j$. Second, for each finite place $\lambda$ of $F_c$, a set $P_\lambda$ of points and,
  for $y\in P_\lambda$, a function $\nu_{\lambda,y}$ on $V_j^{\mathrm{rat}}$; the \emph{integral lifted
  forms} are the $w\in V_j^{\mathrm{rat}}$ with $\nu_{\lambda,y}(w)\le1$ for all $\lambda$ and $y$, and
  $\Lambda_j^{\mathrm{all}}$ is their $\mathcal O_L$-span. Third, an $F_c$-linear isomorphism
  $e:V_j^{\mathrm{aut}}\to V_j^{\mathrm{rat}}$, the lifting of forms. Fourth, for each $\lambda$, a
  bijection $\pi_\lambda$ from $P_\lambda$ onto the set $P'_\lambda$ of valuations of rational sections
  at the algebraic integral points of the Shimura variety over complete discrete valuation rings
  above $\lambda$, such that $\nu_{\lambda,y}(e(f))=N\lambda^{-z(f)}$ for $z=\pi_\lambda(y)$ and
  $f\in V_j^{\mathrm{aut}}$, read as $0$ when $z(f)=\infty$. A cusp form $f$ is integral when
  $z(f)\ge0$ for every $\lambda$ and every $z\in P'_\lambda$; since $N\lambda>1$, this holds if and only
  if $N\lambda^{-z(f)}\le1$ for all of them. Let $\iota$ be $e$ followed by the inclusion of
  $V_j^{\mathrm{rat}}$ in $V_j$. Lemma~\ref{cv-lifting-map}, with $K=F_c$, $W=V_j^{\mathrm{rat}}$,
  $C=V_j^{\mathrm{aut}}$ and $\nu'_{\lambda,z}(f)=N\lambda^{-z(f)}$, gives the three properties of
  $\iota$. So $\Lambda_j^{\mathrm{all}}$ is the $\mathcal O_L$-span of
  $\iota(\Lambda_j^{\mathrm{aut,all}})$.

  It remains to identify $\Lambda_j$. Write $\rho=\rho_j$ for the action of $G$ on $V_j$, and
  $\chi$ for the character of the rational structure, as in the statement. The level subgroup $P_0$ is the
  group of elements of $G$ that are upper unitriangular at every level place, and its character
  $\psi_0\colon P_0\to L^\times$ satisfies
  $\sigma_0(\psi_0(u))=\prod_v\psi_v\big(\sum_{i<r}u_{v,i,i+1}\big)$, for the embedding $\sigma_0$ of
  $L$ at the base place and additive characters $\psi_v$ of the residue fields at the level places
  $v$. A lifted form $w$ has the prescribed type if $\rho(u)w=\chi(u)\psi_0(u)^{-1}w$ for every
  $u\in P_0$; these forms make up the subspace $V_j^{\mathrm{type}}$ of the statement, and
  $\Lambda_j=\Lambda_j^{\mathrm{all}}\cap V_j^{\mathrm{type}}$. A cusp form $f$ has the prescribed type
  when $\iota(f)$ does; these cusp forms make up an $F_c$-subspace, and $\Lambda_j^{\mathrm{aut}}$ is the
  set of elements of $\Lambda_j^{\mathrm{aut,all}}$ of the prescribed type. Apply
  Lemma~\ref{cv-lifting-type} with $K=F_c$, $V=V_j$, $W=V_j^{\mathrm{rat}}$, $M$ the set of integral
  lifted forms, $J=P_0$ and $\theta=\chi\psi_0^{-1}$, so that $V^\theta=V_j^{\mathrm{type}}$ and
  $\mathcal O_LM=\Lambda_j^{\mathrm{all}}$. Its hypotheses hold. By the point data, either the action of
  $G$ on cusp forms through the level preserves $V_j^{\mathrm{aut}}$ and $e$ carries it to $\rho$, so
  that $\rho(g)e(f)=e(gf)$ lies in $V_j^{\mathrm{rat}}$, or $G$ acts trivially on $V_j$; in both cases
  $\rho(g)V_j^{\mathrm{rat}}\subset V_j^{\mathrm{rat}}$. By the weak characterization of the realization,
  $\chi(g)$ is a product, over the level places $v$, of values of characters at $\det g_v$; so
  $\chi(u)=1$ for $u\in P_0$. By the weight tie, the level of the Shimura variety is $g_0g_1$, and
  $F_c$ contains a primitive root of unity $\zeta$ of an order $n_0$ divisible by $g_0g_1$. As $n_0$ is
  $0$ in the residue fields, of characteristic $g_0$ or $g_1$, each $\psi_v(s)^{n_0}=\psi_v(n_0s)$ is
  $1$, so $\psi_0(u)^{n_0}=1$ and $\psi_0(u)^{-1}$ is a power of $\zeta$. Hence $\theta$ takes values in $F_c$.
  The lemma shows that $\Lambda_j$ is the $\mathcal O_L$-span of the integral lifted forms of the
  prescribed type. These are the $\iota(f)$ with $f\in\Lambda_j^{\mathrm{aut,all}}$ of the prescribed
  type, that is, the elements of $\iota(\Lambda_j^{\mathrm{aut}})$. So $\Lambda_j$ is the
  $\mathcal O_L$-span of $\iota(\Lambda_j^{\mathrm{aut}})$.
  % Lean: iota is GlobalTheory.CovTop.SupImage.exists_emb_image_latW, from the fields e, ks, pt and
  % nrm_eq of the weight tie Formalism.IndefR.FWT.FWTie and Formalism.L2.Pj.Core.nrmM_le_one_iff.
  % The identification of Lambda_j is the step hS of GlobalTheory.CoreSlope.coreAt_of, with RatAt
  % from GlobalTheory.CovTop.ratAt_of_rows_alt; its hypotheses are GlobalTheory.CovTop.rho_mem_WK
  % (from the point field rhoGeo) and GlobalTheory.CovTop.chi_mul_psi_inv_mem_core_alt.
\item[(b)] Let $\mu$ be the Haar measure on $U(W_j)(\mathbf A_F)$ normalized by the level subgroup
  as in Proposition~\ref{cv-level-volume}, $\mathcal F$ a fundamental domain, $d_\tau\ge1$ the
  dimension of the representation of the maximal compact subgroup of $U(p_j,q_j)$ of weight
  $\kappa_j$ at the real place $\tau$, and $R_j=\mu(\mathcal F)\prod_\tau d_\tau$. The norm
  $N_{j,\sigma}(\iota f)$ is the $L^2$-norm over $\mathcal F$ of the adelic function of $f$
  paired with a unit vector of that representation. By Schur orthogonality for the maximal
  compact subgroup this is $R_j^{1/2}$ times the Petersson norm of $f$, computed with the
  invariant probability measure: $N_{j,\sigma}(\iota f)=R_j^{1/2}\|f\|_{\sigma|F_c}$.
\item[(c)] The level subgroup $K$ satisfies the hypotheses of Proposition~\ref{cv-level-volume}, with
  $r=k+m-1$, $s=2x^\varepsilon$ and $K_0$ the full level group of $U(W_j)$, which contains the stabilizer $K_V$
  of the self-dual lattice. The level subgroup lies in $K_V$, and it is the inverse image in $K_0$
  of a subgroup under a homomorphism $K_0\to\prod_{\mathfrak P}GL_r(\mathcal O_E/\mathfrak P)$, the
  product over the primes $\mathfrak P$ of $E$ above $g_0g_1$. So
  $[K_0:K]\le\prod_{\mathfrak P}|GL_r(\mathcal O_E/\mathfrak P)|\le\prod_{\mathfrak P}N\mathfrak P^{r^2}$,
  and $\prod_{\mathfrak P}N\mathfrak P\le(g_0g_1)^{[E:\mathbf Q]}=(g_0g_1)^{2f}$, so
  $[K_0:K]\le(g_0g_1)^{2fr^2}$. It remains to show that $s=2x^\varepsilon$ is a size bound for the
  datum (Definition~\ref{cv-level-sizes}); recall $b=\exp(x^\varepsilon)$, so that $\exp(s)=b^2$, and the
  size conditions give $x^\varepsilon\ge10$.
  \begin{enumerate}
  \item[(c1)] The size conditions give $\log(fr)\le x^\varepsilon$ and $\log(r+t)\le x^\varepsilon$, so
    $fr(r+t)\le b^2=\exp(s)$.
  \item[(c2)] The size conditions bound by $4b$ the absolute values of the entries of the weight attached
    to the dominant weight $\kappa_j$ at each $\tau\in\mathcal T$. These entries are the numbers inside
    the integer parts in Definition~\ref{cv-level-sizes}, plus $e_j(\tau)$. An integer part differs from
    its argument by less than $1$, so $|\kappa^\lambda_j(\tau)_a+e_j(\tau)|\le4b+1$ for each of the $fr$
    pairs $(\tau,a)$. As $fr\le b$ by (c1) and $b\ge x^\varepsilon+1\ge3$, the sum in
    Definition~\ref{cv-level-sizes}(2) is at most $fr(4b+1)\le5b^2\le b^4=\exp(2s)$.
  \item[(c3)] The size conditions give $\log g_k\le9x^\varepsilon+1$ for $k=0,1$, so
    $\log(g_0g_1)\le18x^\varepsilon+2\le60x^\varepsilon=30s$.
  \item[(c4)] The parameter choices of the point give $\log\operatorname{rd}_E\le x^\varepsilon$ and
    $|\operatorname{Re}\frac{L'}{L}(1,\chi_d)|\le x^\varepsilon$, so their sum is at most
    $2x^\varepsilon\le2s$.
  \end{enumerate}
  % Lean: GlobalTheory.CovTop.MeasUp.sizes_two (Sizes D (2 L)) from HostNum (logfr, logg, big) and
  % WtNum.up; (c2) is wtSize_le with abs_ent_le, through GlobalTheory.CovTop.WtShape.laOf_read and
  % kapQ_inl, kapQ_inr, abs_wt_le; (c4) is Formalism.BasicNotions.log_rdE_le_scL (from condE_le of
  % ParameterChoice') and the field logDerivative_le_power. The bound is applied in
  % GlobalTheory.CovTop.MeasUp.ratAt_pos_le through OneMeasure (log_covolume_ratio_le).
  That proposition gives
  $\log\mu(\mathcal F)\le fr^2\log(r+3)+238fr^2x^\varepsilon$ and
  $\log\mu(\mathcal F)\ge-3(fr^2(\log(fr+3)+1)+\log|d_E|)$. The Weyl dimension formula and the
  size of the weights give $\log d_\tau\le3r^2x^\varepsilon$. As $fr^2x^\varepsilon\le u$ and
  $\log|d_E|\le u$ by the size conditions, $|\log R_j|\le2C_Mu$ for a constant $C_M$, so
  $|\log R_j^{1/2}|\le C_Mu$.
\item[(d)] Lemma~\ref{cv-slope-base-change}(1), with $K=F_c$, $K'=L$ and $B=C_Mu$, gives
  $\mu(\Lambda_j)\ge\mu(\Lambda_j^{\mathrm{aut}})-2C_Mu$, which is (5) with $C_1=2C_M$. The
  hypothesis $N_{\bar\sigma}=N_\sigma$ of the lemma holds for Petersson norms.
\end{enumerate}

(6) This is Proposition~\ref{cv-arithmetic-volume-step}.
\end{proof}
\leanref{Development/AutomorphicArithmetic/ThetaCorrespondence/ThetaPairLifts/ThirdPlaceSeries/RecAtOfGramAt.lean}{GlobalTheory.CovTop.recAt_of_gramAt_alt}
\leanref{Development/AutomorphicArithmetic/ThetaCorrespondence/ThetaPairLifts/ThirdPlaceSeries/CutAtOfRows.lean}{GlobalTheory.CovTop.cutAt_of_rows_alt}
\leanref{Development/AbelianVarietyTheory/Proofs/DivisorProofs/GroupIdempotentCoefficients.lean}{GlobalTheory.TieIdem.idemOf_props}
\leanref{Development/AutomorphicArithmetic/ThetaCorrespondence/ThetaPairLifts/ThirdPlaceSeries/TwoLogCardLevelsLe.lean}{GlobalTheory.CovTop.two_log_card_levels_le_alt}
\leanref{Development/AutomorphicArithmetic/ThetaCorrespondence/ThetaPairLifts/ThirdPlaceSeries/CoveringBadLocusOfDescentSteps.lean}{GlobalTheory.CovTop.covering_bad_locus_of_descent_steps_alt}
\leanref{Development/AutomorphicArithmetic/ThetaCorrespondence/ThetaPairLifts/ThirdPlaceSeries/PointRead.lean}{GlobalTheory.DefRead.point_read_alt}
\leanref{Development/AutomorphicArithmetic/ThetaCorrespondence/ThetaPairLifts/ThirdPlaceSeries/FreeHalfOfDeformationSpace.lean}{GlobalTheory.DefHalf.halfH_acts_freely_on_internal_sum_decomposition}
\leanref{Development/AutomorphicArithmetic/LocalEmbeddings/BadRepAt.lean}{GlobalTheory.CovTop.BadRep.badRep_at}
\leanref{Development/AutomorphicArithmetic/LocalEmbeddings/KeepsOfGood.lean}{GlobalTheory.CovTop.BadRep.rank_gen}
\leanref{Development/AutomorphicArithmetic/ThetaCorrespondence/ThetaPairLifts/ThirdPlaceSeries/RegularEmbeddingsCoverBadEmbeddingLocus.lean}{GlobalTheory.BadEmb.regular_embeddings_cover_bad_embedding_locus_alt}
\leanref{Development/AbelianVarietyTheory/Proofs/DivisorProofs/CoveringRepresentationTypeFrame.lean}{GlobalTheory.BadEmb.embReg_deep}
\leanref{Development/AutomorphicArithmetic/ThetaCorrespondence/ThetaPairLifts/ThirdPlaceSeries/CoveringLambdaInvariantBound.lean}{GlobalTheory.CovTop.LamAll.covering_lambda_all_places_alt}
\leanref{Development/AutomorphicArithmetic/ThetaCorrespondence/ThetaPairLifts/ThirdPlaceSeries/CoveringMeasureReductionBound.lean}{GlobalTheory.CovTop.MeasRd.covering_measure_reduction_alt}
\leanref{Development/AbelianVarietyTheory/Proofs/DivisorProofs/CoveringLeFromMeasure.lean}{GlobalTheory.CoreSlope.coreAt_of}
\leanref{Development/AutomorphicArithmetic/ThetaCorrespondence/ThetaPairLifts/HodgeBundleWeights/LevelSubgroupIndexBounded.lean}{GlobalTheory.CovTop.LevelIndex.level_index_le_bound}
\leanref{Development/WeilBound/RepresentationTheoryGL/TubeMeasures/DisjointTubes/LevelStabilisers/MemStabOfLev.lean}{MainStatements.Tam.stab}
\leanref{Development/AutomorphicArithmetic/ThetaCorrespondence/ThetaPairLifts/ThirdPlaceSeries/SizeAltBis.lean}{GlobalTheory.CovTop.SizeRows.size_of_alt}
\leanref{Development/AbelianVarietyTheory/Proofs/DivisorProofs/SizeBoundStatement.lean}{GlobalTheory.CovTop.SizeRows.sizeAt_of}

\subsection{Definite fibres and the key formula}\label{subsection-cv-definite-fibres-key-formula}

\begin{lemma}\label{cv-definite-orbits}
Let $K$ be a totally complex number field, let $V$ be a finite-dimensional $K$-vector space with a
Hermitian norm system $N$ (Definition~\ref{cv-hermitian-norms}), and let $C$ be a nonempty finite set.
Let $(e_Z)_{Z\in C}$ be $K$-linear endomorphisms of $V$ with $e_Ze_Z=e_Z$, $e_Ze_{Z'}=0$ for $Z\ne Z'$
and $\sum_Ze_Z=1$. Let $G$ be a finite group acting freely on $C$, let $\rho$ be a representation of
$G$ on $V$ with $\rho(g)e_Z=e_{gZ}\rho(g)$ for all $g$ and $Z$, and let $\chi\colon G\to K^\times$ be a
homomorphism. Let $M_0\subset V$ be a finitely generated $\mathcal O_K$-submodule stable under
$\rho(G)$, and let $M$ be the set of $v\in M_0$ with $\rho(g)v=\chi(g)v$ for all $g\in G$. For each
$Z\in C$ let $I_Z$ be a nonempty finite set, let $L_Z\subset K^{I_Z}$ be a finitely generated
$\mathcal O_K$-submodule that spans $K^{I_Z}$, let $Y_{Z,\sigma}$ be invertible complex matrices indexed
by $I_Z$, one for each embedding $\sigma\colon K\to\mathbf C$, with the Hermitian norm system
$N^Z_\sigma(x)=|Y_{Z,\sigma}\sigma(x)|$ on $K^{I_Z}$, and let $\varphi_Z\colon K^{I_Z}\to V$ be an injective
$K$-linear map with image $e_Z(V)$ and $\varphi_Z(L_Z)\subset M_0$. Let $s$ be a real number with
$\mu(L_Z)\ge s$ for every $Z$. Suppose that, for every embedding $\sigma$, every $v\in V$ and every real
number $B$, $N_\sigma(v)\le B$ holds as soon as, for every $Z\in C$, there are $g\in G$, $o\in C$,
$x\in K^{I_o}$ and $u\in K$ with $Z=go$, $|\sigma(u)|=1$, $e_Zv=u\,\rho(g)\varphi_o(x)$ and
$N^o_\sigma(x)\le B$. Then $M\ne0$ and $\mu(M)\ge s$.
\end{lemma}
\leanref{Development/AbelianVarietyTheory/Proofs/DivisorProofs/LevelMoveEigenlatticeSlopeDefinite.lean}{GlobalTheory.DefPt.OrbitGeo.slope_eig}
\leanref{Development/AbelianVarietyTheory/Proofs/DivisorProofs/FreeOrbitProjectorEigen.lean}{GlobalTheory.freeOrbit_projector_eigen}
\leanref{Development/AbelianVarietyTheory/Proofs/DivisorProofs/ISupIndepRangeInc.lean}{GlobalTheory.SumSlope.le_slope_sq_of_sum}
\begin{proof}
Choose a set $O\subset C$ of representatives of the orbits of $G$. Since the action is free, every
$Z\in C$ is $go$ for exactly one pair $(g,o)\in G\times O$. Each $\chi(g)$ is a root of unity, since
$\chi(g)^{|G|}=1$; so it lies in $\mathcal O_K$, and $|\sigma(\chi(g))|=1$ for every $\sigma$. Define
\[
 \Phi\colon\bigoplus_{o\in O}K^{I_o}\to V,\qquad
 \Phi(x)=\sum_{o\in O}\sum_{g\in G}\chi(g)^{-1}\rho(g)\varphi_o(x_o).
\]
As $\varphi_o(x_o)=e_o\varphi_o(x_o)$, the term of $(g,o)$ is $e_{go}\chi(g)^{-1}\rho(g)\varphi_o(x_o)$,
and the $e_Z$ are orthogonal idempotents. Hence $e_{go}\Phi(x)=\chi(g)^{-1}\rho(g)\varphi_o(x_o)$ for all
$g\in G$ and $o\in O$. In particular $e_o\Phi(x)=\varphi_o(x_o)$, so $\Phi$ is injective. Replacing $g$ by
$hg$ in the sum gives $\rho(h)\Phi(x)=\chi(h)\Phi(x)$. If $x_o\in L_o$ for every $o$, then
$\Phi(x)\in M_0$, since $M_0$ is stable under $\rho(G)$ and the $\chi(g)^{-1}$ are integral; so
$\Phi(x)\in M$. Conversely, let $v\in M$. Then
$e_{go}v=e_{go}\chi(g)^{-1}\rho(g)v=\chi(g)^{-1}\rho(g)e_ov$, and $e_ov=\varphi_o(x_o)$ for some $x_o$
because $e_ov$ lies in $e_o(V)$. Summing over $Z=go$ gives $v=\sum_Ze_Zv=\Phi(x)$. So $KM$ lies in the
image of $\Phi$, and the image of $\bigoplus_oL_o$ is a submodule of $M$ of the same rank.

Fix $\sigma$ and $x$, and put $B=\big(\sum_{o\in O}N^o_\sigma(x_o)^2\big)^{1/2}$. Each $e_{go}\Phi(x)$ is
$\chi(g)^{-1}\rho(g)\varphi_o(x_o)$ with $N^o_\sigma(x_o)\le B$, so the hypothesis gives
$N_\sigma(\Phi(x))^2\le\sum_oN^o_\sigma(x_o)^2$. Thus $\Phi$, from the orthogonal sum of the norm
systems $N^o$ to $N$, does not increase norms at any embedding, and so does not increase Gram
determinants. Hence the covolume of $\Phi(\bigoplus_oL_o)$ for $N$ is at most that of
$\bigoplus_oL_o$ for the orthogonal sum. The covolume of $M$ is at most that of its submodule
$\Phi(\bigoplus_oL_o)$ of the same rank. By Definition~\ref{cv-hermitian-norms},
$\mu(M)\ge\mu\big(\bigoplus_oL_o\big)$, and the slope of an orthogonal sum is the average of the slopes
$\mu(L_o)$ weighted by the ranks, which is at least $s$. Finally $M\ne0$, since it contains
$\Phi(\bigoplus_oL_o)$, $\Phi$ is injective and $L_o\ne0$.
\end{proof}

\begin{lemma}\label{cv-definite-fibre-lattice}
Let $K$ be a totally complex number field and $T$ a nonempty finite set. For each $t\in T$ let
$N_t\ge1$, let $\lambda_t\in\mathbf N^{N_t}$ be non-increasing with $|\lambda_t|=\sum_i\lambda_{t,i}$,
let $(f_I)_{I\in\mathcal I_t}$ be the basis of bideterminants of $\nabla(\lambda_t)$ of
Lemma~\ref{is-standard-basis}, whose size is $\dim\lambda_t=\dim\nabla(\lambda_t)$, and for an
$N_t\times N_t$ matrix $g$ let $\rho_t(g)$ be the matrix, in this basis, of the map $f(x)\mapsto f(gx)$
of $\nabla(\lambda_t)$ (Definition~\ref{is-polynomial-model}). Let $Q_t$ be the matrix of the basis
$(f_I)$ of $\nabla(\lambda_t)\otimes\mathbf C$ in an orthonormal basis for the inner product
$\langle f,f'\rangle=\dim\lambda_t\int_{U(N_t)}f(u)\overline{f'(u)}\,du$, for the Haar probability
measure on the unitary group, and let $I=\prod_t\mathcal I_t$. For each finite place $w$ of $K$ let
$a_w\in K^\times$ and invertible matrices $B_{w,t}$ over $K$ be given, with $\operatorname{ord}_w(a_w)=0$
and with $B_{w,t}$ and $B_{w,t}^{-1}$ integral at $w$ for all but finitely many $w$. For each embedding
$\sigma\colon K\to\mathbf C$ let $s_\sigma>0$ and invertible complex matrices $X_{\sigma,t}$ be given, with
$s_{\bar\sigma}=s_\sigma$ and $|\det X_{\bar\sigma,t}|=|\det X_{\sigma,t}|$. Let $L\subset K^I$ be the set
of $x$ such that the vector $a_w\big(\bigotimes_t\rho_t(B_{w,t})\big)x$ is integral at $w$ for every $w$,
and let $Y_\sigma=s_\sigma\bigotimes_tQ_t\rho_t(X_{\sigma,t})$. Then $L$ is a finitely generated
$\mathcal O_K$-submodule spanning $K^I$, every $Y_\sigma$ is invertible, and, for the Hermitian norm
system $x\mapsto|Y_\sigma\sigma(x)|$ (Definition~\ref{cv-hermitian-norms}),
\[
 \frac{\mu(L)}2\ge\frac1{[K:\mathbf Q]}\Big(\sum_w\operatorname{ord}_w(a_w)\log Nw-\sum_\sigma\log s_\sigma\Big)
 +\sum_{t\in T}\Big(\frac{|\lambda_t|}{N_t}\deg_t-\frac12\log\dim\lambda_t\Big),
\]
where
$\deg_t=\frac1{[K:\mathbf Q]}\big(\sum_w\operatorname{ord}_w(\det B_{w,t})\log Nw-\sum_\sigma\log|\det X_{\sigma,t}|\big)$.
\end{lemma}
\leanref{Development/AbelianVarietyTheory/Proofs/DivisorProofs/SlopeDualWeylFibreLatticeGeStatement.lean}{GlobalTheory.T.slope_dualWeylFibreLattice_ge}
\leanref{Development/AbelianVarietyTheory/Proofs/DivisorProofs/IsLatticeLatticeOfLocalMatrixFamilyStatement.lean}{GlobalTheory.T.isLattice_latticeOfLocalMatrixFamily}
\leanref{Development/AbelianVarietyTheory/Proofs/DivisorProofs/LatticeSlopeSqLocalMatricesEqDeterminantsStatement.lean}{GlobalTheory.T.latticeSlopeSq_localMatrices_eq_determinants}
\leanref{Development/AbelianVarietyTheory/Proofs/DivisorProofs/DualWeylGramDeterminantLeDimPowDimStatement.lean}{GlobalTheory.T.dualWeyl_gramDeterminant_le_dim_pow_dim}
\begin{proof}
Write $A_w=a_w\bigotimes_t\rho_t(B_{w,t})$.
For local matrices $A_w$ that are integral with integral inverse at almost all $w$, and complex matrices
$Y_\sigma$ with $|\det Y_{\bar\sigma}|=|\det Y_\sigma|$, the module of $x$ with $A_wx$ integral at every $w$
is obtained from $\mathcal O_K^I$ by changing it at finitely many places, and
Definition~\ref{cv-hermitian-norms} gives its slope as
\[
 \frac{2}{|I|\,[K:\mathbf Q]}\Big(\sum_w\operatorname{ord}_w(\det A_w)\log Nw-\sum_\sigma\log|\det Y_\sigma|\Big),
\]
since $\mathcal O_K^I$ with the standard norms has slope $0$. The determinant of $\rho_t(g)$ is a
polynomial character of $GL_{N_t}$, hence a power of $\det g$; a scalar $c$ acts on $\nabla(\lambda_t)$ by
$c^{|\lambda_t|}$, so $\det\rho_t(g)=(\det g)^{|\lambda_t|\dim\lambda_t/N_t}$. The determinant of a
Kronecker product gives
$\operatorname{ord}_w\det A_w=|I|\big(\operatorname{ord}_w(a_w)+\sum_t\frac{|\lambda_t|}{N_t}\operatorname{ord}_w\det B_{w,t}\big)$ and
$\log|\det Y_\sigma|=|I|\big(\log s_\sigma+\sum_t\big(\frac{\log|\det Q_t|}{\dim\lambda_t}+\frac{|\lambda_t|}{N_t}\log|\det X_{\sigma,t}|\big)\big)$.
Finally, every $f_I$ is a product of minors, so $|f_I(u)|\le1$ for unitary $u$ and
$\langle f_I,f_I\rangle\le\dim\lambda_t$; Hadamard's inequality for the columns of $Q_t$ gives
$|\det Q_t|^2\le(\dim\lambda_t)^{\dim\lambda_t}$, that is,
$\log|\det Q_t|/\dim\lambda_t\le\frac12\log\dim\lambda_t$. Substituting gives the claim.
\end{proof}

\begin{theorem}\label{cv-definite-key-formula}
There is a rule $\kappa$ with the following properties. For every field $K$ of characteristic $0$ and
every principally polarized abelian variety $A$ of dimension $g$ over $K$, with a symmetric line bundle
defining the polarization and rigidified along the zero section, a generator $\theta$ of its space of
sections and a basis $v=(v_1,\dots,v_g)$ of the Lie algebra of $A$, it gives a number
$\kappa(A,\theta,v)\in K^\times$.
\begin{enumerate}
\item If $a\in K^\times$, $M\in GL_g(K)$ and $v'_i=\sum_jM_{ij}v_j$, then
  $a^8\kappa(A,a\theta,v')=\det(M)^4\kappa(A,\theta,v)$.
\item $\kappa$ is compatible with homomorphisms of fields of characteristic $0$.
\item If $K$ is the fraction field of a discrete valuation ring $R$ and the data extend over $R$, to a
  principally polarized abelian scheme with a rigidified symmetric line bundle, a generator of its
  sections over $R$ and a basis over $R$ of its Lie algebra, then $\kappa(A,\theta,v)$ is a unit $u$ of
  $R$, and for every surjection $r$ of $R$ onto a field of characteristic $0$, $\kappa$ of the data
  over that field is $r(u)$.
\item Over $\mathbf C$, let $A=\mathbf C^g/(\mathbf Z^g+\Omega\mathbf Z^g)$ for a symmetric matrix
  $\Omega$ with positive definite imaginary part, let $\theta$ be the theta function with a
  characteristic in $\{0,\frac12\}^{2g}$, and let $v_i$ be the vector $(2\pi i)^{-1}e_i$. Then
  $\kappa(A,\theta,v)=1$.
\item Over $\mathbf C$, let $A=\mathbf C^g/\Lambda$, let $H$ be the Hermitian form of the
  polarization, let $\theta$ also denote the function on $\mathbf C^g$ that represents the generator, let
  $W$ be the matrix whose $i$-th row is the vector $v_i\in\mathbf C^g$, and let $F$ be a fundamental
  domain for $\Lambda$. Then
  \[
   |\kappa(A,\theta,v)|\Big(\frac1{\operatorname{vol}F}\int_F|\theta(z)|^2e^{-\pi H(z,z)}\,dz\Big)^4
   \Big(\frac{\operatorname{vol}F}{\pi^g|\det W|^2}\Big)^2=(2\pi)^{2g}.
  \]
\end{enumerate}
\end{theorem}
\leanref{ReciprocityTheory/ReductiveGroups/ArithmeticGeometry/BirationalGeometry/ToroidalCompactification/IsLieBasis.lean}{ReciprocityTheory.ChartRestriction.G.Geo.Tor.KeyFormula}
\leanref{Development/AbelianVarietyTheory/Proofs/ModuliProofs/ExactKeyFormulaProofStatement.lean}{GlobalTheory.T.exactKeyFormula_proof}

Parts (1) to (3) are the key formula of Moret-Bailly \cite{MB85}: the eighth power of the line of
theta sections is the fourth power of the dual of the determinant of the Lie algebra, over any base.
Parts (4) and (5) are its archimedean form, with the normalization given by the classical theta
functions.

\begin{situation}\label{cv-definite-situation}
Let $E$ be a CM field with a CM type $\Phi$, and let a PEL datum of signature $(p,0)$ over $E$ with
$p\ge1$, a level of full level $1$ (Definition~\ref{cv-model-toroidal-semiabelian}(1)), its moduli
scheme $X$ over its coefficient field $F_c$, and uniformizations of $X$ at the complex embeddings of
$F_c$, normalized by a real number, that satisfy the compatibility conditions of the moduli problem, be
given (Definition~\ref{cv-constant-sections}). Suppose that the
prime-to-$\ell$ part of the level is at least $3$ for every prime $\ell$, and that $F_c$ is
algebraically closed in the function field of every connected component of $X$. For a point $y$ of $X$
over an $F_c$-algebra $S$, let $A_y$ be the abelian scheme over $S$ that it carries, for $\tau\in\Phi$
let $\operatorname{Lie}_\tau(y)$ be the $\tau$-eigenspace of $\mathcal O_E$ on the Lie algebra of $A_y$,
free of rank $p$, and let $\Theta(y)$ be the module of sections of the symmetric line bundle of the
principal polarization of $A_y$, free of rank one.
\begin{enumerate}
\item A \emph{framed point} over $S$ is such a point $y$ together with a basis
  $u_\tau=(u_{\tau,1},\dots,u_{\tau,p})$ of $\operatorname{Lie}_\tau(y)$ for each $\tau\in\Phi$ and a
  generator $\vartheta$ of $\Theta(y)$. For invertible $p\times p$ matrices $Y=(Y_\tau)_{\tau\in\Phi}$
  over $S$ and $t\in S^\times$, the framed point $y\cdot(Y,t)$ has the same point, the bases whose
  $i$-th vectors are $\sum_j(Y_\tau)_{ji}u_{\tau,j}$, and the generator $t\vartheta$.
\item Let $\nu=(\nu_\tau)_{\tau\in\Phi}$ be non-increasing sequences in $\mathbf N^p$ and $n\ge0$. A
  \emph{form of weight $(\nu,n)$} is a function $f$ that assigns to every framed point $y$ over an
  $F_c$-algebra $S$ an element $f(y)\in S$, compatibly with base change, such that
  $f(y\cdot(Y,t))=t^n\prod_{\tau,i}(Y_\tau)_{ii}^{\nu_{\tau,i}}f(y)$ whenever every $Y_\tau$ is upper
  triangular. It is \emph{supported} on a connected component $Z$ of $X$ if $f(y)=0$ for every framed
  point $y$ over a field whose point does not lie on $Z$.
\item Let $T$ be a complete discrete valuation ring above a finite place of $F_c$. A framed point over the
  fraction field of $T$ is \emph{integral} if it extends over $T$: to an abelian scheme over $T$ with
  bases of the eigenspaces of its Lie algebra and a generator of the sections of its line bundle. A form
  $f$ is \emph{integral} if $f(y)\in T$ for every such $T$ and every integral framed point $y$.
\item Let $\sigma$ be a complex embedding of $F_c$ and $N>0$. A framed point over $\mathbf C$ via
  $\sigma$ is \emph{metric with constant $N$} if, in the uniformization at $\sigma$, each basis $u_\tau$
  is orthonormal for $\pi H$, where $H$ is the Hermitian form of the polarization, and the mean of
  $|\vartheta(z)|^2e^{-\pi H(z,z)}$ over a fundamental domain of the period lattice is $N$.
\end{enumerate}
\end{situation}
\leanref{Formalism/LocalToGlobal/ProjectiveLattices/IsCombAt.lean}{Formalism.L2.Pj.FrPt}
\leanref{Development/AbelianVarietyTheory/Definitions/ExistenceOfLattices/UnivFrame.lean}{GlobalTheory.Frames.FrChg}
\leanref{Formalism/LocalToGlobal/ProjectiveLattices/IsCombAt.lean}{Formalism.L2.Pj.sec}
\leanref{Formalism/LocalToGlobal/ProjectiveLattices/VAtVNeBot.lean}{Formalism.L2.Pj.Integral}
\leanref{Development/AbelianVarietyTheory/Definitions/CMLattices/IntegralOf.lean}{GlobalTheory.LamT.IsMetricAt}
\leanref{Development/AbelianVarietyTheory/Definitions/Divisors/EigLatLe.lean}{GlobalTheory.DefPt.DefCase}

\begin{lemma}\label{cv-definite-point}
In Situation~\ref{cv-definite-situation}, let $Z$ be a connected component of $X$.
\begin{enumerate}
\item There is a framed point $y_0$ over $F_c$ on $Z$ such that every point of $X$ over an
  $F_c$-algebra that lies on $Z$ is the base change of the point of $y_0$.
\item For every finite place $w$ of $F_c$ there are $a_w\in F_c^\times$ and invertible matrices
  $B_w=(B_{w,\tau})_{\tau\in\Phi}$ over $F_c$ with the following property: for every complete discrete
  valuation ring $T$ above $w$, every integral framed point over the fraction field of $T$ that lies on
  $Z$ is $y_0\cdot(B_wY,a_wt)$, after base change, for some matrices $Y=(Y_\tau)$ over $T$ with unit
  determinants and some $t\in T^\times$. One can take $a_w=1$ and $B_w=1$ for all but finitely many
  $w$.
\item For every complex embedding $\sigma$ of $F_c$ there are $N_\sigma>0$ and invertible complex
  matrices $X_\sigma=(X_{\sigma,\tau})_{\tau\in\Phi}$ such that the base change of $y_0$ to $\mathbf C$
  via $\sigma$, changed by $(X_\sigma,1)$, is metric with constant $N_\sigma$.
\end{enumerate}
Choose such data, and put, with sums over the finite places $w$ and the complex embeddings $\sigma$ of
$F_c$,
\[
 \deg_\Theta=\frac{1}{[F_c:\mathbf Q]}\Big(\sum_w\operatorname{ord}_w(a_w)\log Nw+\frac12\sum_\sigma\log N_\sigma\Big),\quad
 \deg_\tau=\frac{1}{[F_c:\mathbf Q]}\Big(\sum_w\operatorname{ord}_w(\det B_{w,\tau})\log Nw-\sum_\sigma\log|\det X_{\sigma,\tau}|\Big).
\]
\begin{enumerate}
\item[(4)] If $s$ is a flat section of type $\tau$ on $Z$ (Definition~\ref{cv-constant-sections}), then
  $\deg_\tau=-\deg s$.
\item[(5)] $\deg_\Theta\ge\frac12\sum_{\tau\in\Phi}\deg_\tau+p\,|\Phi|\big(\frac14\log2\pi-\frac12\log\operatorname{rd}_E\big)$.
\end{enumerate}
\end{lemma}
\leanref{Development/AbelianVarietyTheory/Definitions/Divisors/EigLatLe.lean}{GlobalTheory.DefPt.PtData}
\leanref{Development/AbelianVarietyTheory/Proofs/DivisorProofs/NonemptyDefiniteRationalPointDataComponentStatement.lean}{GlobalTheory.T.nonempty_definiteRationalPointData_component}
\leanref{Development/AbelianVarietyTheory/Proofs/FoundationProofs/IsOpenImmersionOfSection.lean}{GlobalTheory.FlatLine.DefinitePoint.main}
\leanref{Development/AbelianVarietyTheory/Proofs/DivisorProofs/DegreeDualDetLieEqNegFlatSectionDegreeStatement.lean}{GlobalTheory.T.degree_dualDetLie_eq_neg_flatSection_degree}
\leanref{Development/AbelianVarietyTheory/Proofs/DivisorProofs/ThetaLineDegreeGeHalfDetLieDefinitePointStatement.lean}{GlobalTheory.T.thetaLine_degree_ge_half_detLie_definitePoint}
\begin{proof}
(1) In signature $(p,0)$ the moduli scheme $X$ is finite and \'etale over $F_c$. So $Z$ is the spectrum
of a finite extension of $F_c$ inside its function field, which is $F_c$ because $F_c$ is algebraically
closed there. Thus $Z$ is an $F_c$-rational point and an open subscheme of $X$, through which every point
lying on $Z$ factors. Its eigenspaces and its module $\Theta$ are free; choose bases and a generator.

(2) Two framed points with the same point over the same ring differ by a unique change $(Y,t)$. Let
$w$ lie over the prime $\ell$. As the prime-to-$\ell$ part of the level is at least $3$, the abelian
variety of $y_0$ has good reduction at $w$, and the base change of $y_0$ to the completion
$F_{c,w}$ has the same point as an integral framed point $y_w$; so $y_w$ is $y_0$ changed by some
$(M,b)$ over $F_{c,w}$. Since $F_c$ is dense in $F_{c,w}$, there are $B_w$ over $F_c$ and $a_w\in
F_c^\times$ with $B_w^{-1}M$ integral with unit determinant and $b/a_w$ a unit. An integral framed point
over the fraction field of $T$ on $Z$ and the base change of $y_w$ have the same point and both extend
over $T$, so they differ by a change $(Y',t')$ with $Y'$ invertible over $T$ and $t'\in T^\times$. This
gives the form $y_0\cdot(B_wY,a_wt)$. The framed point $y_0$ extends over the ring of integers of $F_c$
away from finitely many places, and at the other places $a_w=1$ and $B_w=1$ will do.

(3) Orthonormalize the bases of $y_0$ for $\pi H$ in each eigenspace; this gives $X_\sigma$, and
$N_\sigma$ is the mean of $|\vartheta|^2e^{-\pi H}$ for the generator of $y_0$.

(4) The section $s$ has weight $-1$ at $\tau$, so, read on framed points, $s(y\cdot(Y,t))=\det(Y_\tau)^{-1}s(y)$,
and there is $b\in F_c^\times$ with $s(y_0\cdot(Y,t))=b\det(Y_\tau)^{-1}$ for all $(Y,t)$. In signature $(p,0)$
the defect $\delta_\tau$ of flat sections is $0$. At an integral framed point $y_0\cdot(B_wY,a_wt)$ above
$w$, the valuation of $s$ is $\operatorname{ord}_w(b)-\operatorname{ord}_w(\det B_{w,\tau})$, which is therefore
$\rho_s(w)$. At the metric framed point of (3), the norm of $s$ is $c_\sigma(s)=|\sigma(b)|/|\det X_{\sigma,\tau}|$.
By the definition of the degree and the product formula for $b$, $\deg s=-\deg_\tau$.

(5) Let $g=p|\Phi|$, let $v$ be the basis of the Lie algebra of $A_{y_0}$ formed by the bases
$u_\tau$ of $y_0$, and let $\kappa_0=\kappa(A_{y_0},\vartheta,v)\in F_c^\times$ be given by
Theorem~\ref{cv-definite-key-formula}.

At a complex embedding $\sigma$, the framed point $y$ of (3) has the basis $v$ changed by the
matrices $X_{\sigma,\tau}$ and the same generator. By parts (1) and (2) of the theorem, $\kappa$ at $y$ is
$\prod_\tau\det(X_{\sigma,\tau})^4\sigma(\kappa_0)$. As $y$ is metric, the rows of $W$ are orthonormal
for $\pi H$, so $\pi^g|\det W|^2\det H=1$; as the polarization is principal,
$\operatorname{vol}F\cdot\det H=1$. So the last factor in part (5) is $1$, and part (5) gives
\[
 4\sum_\tau\log|\det X_{\sigma,\tau}|+\log|\sigma(\kappa_0)|+4\log N_\sigma=2g\log2\pi.
\]

At a finite place $w$, let $T$ be the ring of integers of the completion $F_{c,w}$ and $y$ an integral
framed point over $F_{c,w}$ on $Z$, which exists by the proof of (2); then $y=y_0\cdot(B_wY,a_wt)$, and parts (1) and (2) give
$a_w^8t^8\kappa(y)=\prod_\tau\det(B_{w,\tau}Y_\tau)^4\kappa_0$, where $\kappa(y)$ is $\kappa$ at $y$. Let
$v'$ be a basis over $T$ of the Lie algebra of the abelian scheme over $T$; by part (3),
$\kappa$ is a unit at $v'$ and the generator of $y$. The bases of $y$ together form $Mv'$ for a matrix $M$ over
$T$. For each $\tau$ there is an element $\delta_\tau$ of the different of $E$, embedded by $\tau$, with
$\operatorname{ord}_w(\delta_\tau)=\operatorname{ord}_w(\mathfrak d_\tau)$, and multiplying by
$\delta_\tau$ carries the Lie algebra over $T$ into the sum of its eigenspaces; so $\det M$ divides
$\prod_\tau\delta_\tau^p$. By part (1), $\operatorname{ord}_w\kappa(y)=4\operatorname{ord}_w\det M$, which
lies between $0$ and $4p\sum_\tau\operatorname{ord}_w(\mathfrak d_\tau)$. Hence
\[
 4\sum_\tau\operatorname{ord}_w(\det B_{w,\tau})+\operatorname{ord}_w(\kappa_0)-8\operatorname{ord}_w(a_w)
 \le4p\sum_\tau\operatorname{ord}_w(\mathfrak d_\tau).
\]

Multiply the finite inequalities by $\log Nw/[F_c:\mathbf Q]$ and sum over $w$, and subtract the mean over
$\sigma$ of the archimedean identities. The terms in $\kappa_0$ cancel by the product formula, and
$\frac1{[F_c:\mathbf Q]}\sum_w\operatorname{ord}_w(\mathfrak d_\tau)\log Nw=\log\operatorname{rd}_E$, since
the norm of the different of $E$ is $|d_E|$. The result is
$4\sum_\tau\deg_\tau-8\deg_\Theta\le4g\log\operatorname{rd}_E-2g\log2\pi$, which is (5).
\end{proof}

\begin{lemma}\label{cv-definite-fibre}
In Situation~\ref{cv-definite-situation}, let $\nu=(\nu_\tau)_{\tau\in\Phi}$ be non-increasing sequences
in $\mathbf N^p$, let $n\ge0$ be even, and let $c_1$ be a real number such that on every connected
component of $X$, for every $\tau\in\Phi$, flat sections of type $\tau$ exist and have degree $-c_1$
(Definition~\ref{cv-constant-sections}). Let $Z$ be a connected component of $X$, and put
\[
 B=\sum_{\tau\in\Phi}\Big(\Big(\frac{|\nu_\tau|}p+\frac n2\Big)c_1-\frac12\log\dim\nu_\tau\Big)
 +n\,p\,|\Phi|\Big(\frac14\log2\pi-\frac12\log\operatorname{rd}_E\Big),
\]
where $|\nu_\tau|=\sum_i\nu_{\tau,i}$ and $\dim\nu_\tau=\dim\nabla(\nu_\tau)$
(Definition~\ref{is-polynomial-model}). Then there are a nonempty finite set $I$, a finitely generated
$\mathcal O_{F_c}$-submodule $L$ spanning $F_c^I$, invertible complex matrices $Y_\sigma$ indexed by $I$,
one for each complex embedding $\sigma$ of $F_c$, and an injective $F_c$-linear map $\varphi$ from $F_c^I$
onto the space of forms of weight $(\nu,n)$ supported on $Z$, such that:
\begin{enumerate}
\item $\mu(L)\ge2B$ for the Hermitian norm system $x\mapsto|Y_\sigma\sigma(x)|$;
\item $\varphi(x)$ is integral for every $x\in L$;
\item $|\varphi(x)(y)|\le|Y_\sigma\sigma(x)|$ for every $x\in F_c^I$, every $\sigma$ and every framed point
  $y$ over $\mathbf C$ via $\sigma$ on $Z$ that is metric with constant $1$.
\end{enumerate}
\end{lemma}
\leanref{Development/AbelianVarietyTheory/Proofs/DivisorProofs/DefiniteCaseCuspFormFibreSlopeGeBoundStatement.lean}{GlobalTheory.T.definiteCase_cuspForm_fibre_slope_ge_bound}
\leanref{Development/AbelianVarietyTheory/Definitions/Divisors/EigLatLe.lean}{GlobalTheory.DefPt.FibreAt}
\leanref{Development/AbelianVarietyTheory/Proofs/DivisorProofs/ExistsInjectiveDualWeylTensorCuspFormsRationalPoStatement.lean}{GlobalTheory.T.exists_injective_dualWeylTensor_cuspForms_rationalPoint}
\leanref{Development/AbelianVarietyTheory/Proofs/DivisorProofs/IntegralCoordinatesFormsAndFibreNormBoundStatement.lean}{GlobalTheory.T.integral_coordinates_forms_and_fibre_norm_bound}
\leanref{Development/AbelianVarietyTheory/Proofs/DivisorProofs/MulNormDualWeylTensorEvaluationLeFibreNormStatement.lean}{GlobalTheory.T.mul_norm_dualWeylTensorEvaluation_le_fibreNorm}
\begin{proof}
Let $y_0$, $a_w$, $B_w$, $N_\sigma$ and $X_\sigma$ be as in Lemma~\ref{cv-definite-point}. For
$\tau\in\Phi$ let $(f_I)_{I\in\mathcal I_\tau}$ be the basis of $\nabla(\nu_\tau)$ of
Lemma~\ref{is-standard-basis}, let $\rho_\tau$ and $Q_\tau$ be as in
Lemma~\ref{cv-definite-fibre-lattice}, let $I=\prod_\tau\mathcal I_\tau$, and for $x\in F_c^I$ let
$P_x(Y)=\sum_{I}x_I\prod_\tau f_{I_\tau}(Y_\tau)$, a polynomial in matrices $Y=(Y_\tau)$. Then
$P_x(BY)=P_{\rho(B)x}(Y)$ with $\rho(B)=\bigotimes_\tau\rho_\tau(B_\tau)$.

Forms on $Z$. By Lemma~\ref{cv-definite-point}(1), every framed point on $Z$ over an $F_c$-algebra $S$
is, after base change, $y_0\cdot(Y,t)$ for exactly one invertible $Y$ over $S$ and $t\in S^\times$. A form
$f$ supported on $Z$ is therefore determined by $F(Y,t)=f(y_0\cdot(Y,t))$. Compatibility with base
change, applied to the universal change of frames, makes $F(Y,t)=t^nG(Y)$ for a polynomial $G$, and the
transformation law under upper triangular matrices says that $G$ is homogeneous of degree $\nu_{\tau,i}$
in the $i$-th column of $Y_\tau$ and invariant under $Y_\tau\mapsto Y_\tau u$ for upper unitriangular
$u$. By Definition~\ref{is-polynomial-model} and Lemma~\ref{is-standard-basis}, $G=P_x$ for exactly one
$x\in F_c^I$. Conversely each $P_x$ defines in this way a form $\varphi(x)$ of weight $(\nu,n)$ supported
on $Z$. So $\varphi$ is an injective linear map onto the forms supported on $Z$.

Integrality. Let $L$ be the module of Lemma~\ref{cv-definite-fibre-lattice} for $T=\Phi$,
$N_\tau=p$, $\lambda_\tau=\nu_\tau$, the elements $a_w^n$ and the matrices $B_{w,\tau}$. Let $x\in L$, and
let $y$ be an integral framed point on $Z$ over the fraction field of a complete discrete valuation ring
above $w$. By Lemma~\ref{cv-definite-point}(2), $y=y_0\cdot(B_wY,a_wt)$ with $Y$ integral and $t$ a unit,
so $\varphi(x)(y)=t^nP_{x'}(Y)$ with $x'=a_w^n\rho(B_w)x$ integral at $w$. The $f_I$ have integer
coefficients, so $\varphi(x)(y)$ is integral. This is (2).

Norms. Let $y$ be a framed point over $\mathbf C$ via $\sigma$ on $Z$, metric with constant $1$, and let
$z$ be the framed point of Lemma~\ref{cv-definite-point}(3), metric with constant $N_\sigma$. Both
carry orthonormal bases of the same eigenspaces, so $y=z\cdot(U,t)$ with every $U_\tau$ unitary and
$|t|^2N_\sigma=1$. Hence $|\varphi(x)(y)|=N_\sigma^{-n/2}|P_{\sigma(x)}(X_\sigma U)|
=N_\sigma^{-n/2}|P_{x''}(U)|$ with $x''=\rho(X_\sigma)\sigma(x)$. For a polynomial $h$ in the tensor
product of the $\nabla(\nu_\tau)\otimes\mathbf C$ and unitary $U$, $|h(U)|$ is at most the norm of $h$ for
the tensor product of the inner products of Lemma~\ref{cv-definite-fibre-lattice}: the space is an
irreducible representation of $\prod_\tau U(p)$, evaluation at $U$ is evaluation at $1$ after
translating by $U$, and Schur orthogonality with the Cauchy--Schwarz inequality bounds evaluation at $1$
by this norm. With $Y_\sigma=N_\sigma^{-n/2}\bigotimes_\tau Q_\tau\rho_\tau(X_{\sigma,\tau})$ this gives
$|\varphi(x)(y)|\le|Y_\sigma\sigma(x)|$, which is (3).

Slope. The complex conjugate of a metric framed point is metric with the same constant, so
$N_{\bar\sigma}=N_\sigma$ and $|\det X_{\bar\sigma,\tau}|=|\det X_{\sigma,\tau}|$.
Lemma~\ref{cv-definite-fibre-lattice}, with $s_\sigma=N_\sigma^{-n/2}$, shows that $L$ is a lattice and
\[
 \frac{\mu(L)}2\ge n\deg_\Theta+\sum_\tau\Big(\frac{|\nu_\tau|}p\deg_\tau-\frac12\log\dim\nu_\tau\Big),
\]
with $\deg_\Theta$ and $\deg_\tau$ as in Lemma~\ref{cv-definite-point}. By part (4) of that lemma and the
hypothesis on flat sections, $\deg_\tau=c_1$ for every $\tau$. By part (5),
$\deg_\Theta\ge\frac12|\Phi|c_1+p|\Phi|\big(\frac14\log2\pi-\frac12\log\operatorname{rd}_E\big)$. As
$n\ge0$, $\mu(L)/2\ge B$, which is (1).
\end{proof}

\begin{lemma}\label{cv-isotypic-twisted-trace}
Let $K$ be a field, let $R$ be a commutative $K$-algebra, and let $H$ be a finite group acting on
$R$ on the right by $K$-algebra automorphisms, written $x\mapsto\gamma^*x$, so that
$(\gamma\gamma')^*x=\gamma'^*(\gamma^*x)$. Let $\chi\colon H\to K^\times$ be a character, and for $F\in R$
put
\[
 T(F)=\sum_{g\in H}\chi(g)^{-1}g^*F.
\]
Let $v\in R$ be an element whose group determinant $\det\big((\gamma\gamma')^*v\big)_{\gamma,\gamma'\in H}$
is not a zero divisor in $R$.
\begin{enumerate}
\item For every $F\in R$ and $g\in H$, $g^*T(F)=\chi(g)T(F)$.
\item If $f\in R$ satisfies $T(f\cdot h^*v)=0$ for every $h\in H$, then $f=0$.
\item Let $V,V_+\subset R$ be finite-dimensional $K$-subspaces with $h^*x\in V$ and $x\cdot h^*v\in V_+$
  for every $h\in H$ and $x\in V$, and let $V^\chi\subset V$ be the subspace of the $x$ with
  $g^*x=\chi(g)x$ for every $g\in H$. Then $|H|\dim_KV^\chi\le\dim_KV_+$.
\end{enumerate}
\end{lemma}
\leanref{Development/AbelianVarietyTheory/Proofs/DivisorProofs/TwistedTraceEquivariantAndInjectiveOfGroupDet.lean}{GlobalTheory.twistedTrace_equivariant_and_injective_of_groupDet}
\leanref{Development/AbelianVarietyTheory/Proofs/DivisorProofs/FinrankLeCardMulFinrankEigenspaceLe.lean}{GlobalTheory.finrank_le_card_mul_finrank_eigenspace_le}
\begin{proof}
Let $\Delta$ be the group determinant of $v$. If $(a_\gamma)_{\gamma\in H}\in R^H$ satisfies
$\sum_{\gamma'}\big((\gamma\gamma')^*v\big)a_{\gamma'}=0$ for every $\gamma$, or
$\sum_\gamma\big((\gamma\gamma')^*v\big)a_\gamma=0$ for every $\gamma'$, then multiplying by the adjugate
of the matrix $\big((\gamma\gamma')^*v\big)_{\gamma,\gamma'}$, or of its transpose, gives $\Delta a_\gamma=0$
for every $\gamma$, so $a=0$.

(1) Since $g^*h^*F=(hg)^*F$, $g^*T(F)=\sum_h\chi(h)^{-1}(hg)^*F$. Substituting $h'=hg$ gives
$\sum_{h'}\chi(h'g^{-1})^{-1}h'^*F=\chi(g)T(F)$.

(2) As $g^*$ is a ring automorphism, $T(f\cdot h^*v)=\sum_g\chi(g)^{-1}(g^*f)\,(hg)^*v$. If this
vanishes for every $h$, the first paragraph, with $a_g=\chi(g)^{-1}g^*f$, gives $g^*f=0$ for every $g$;
for $g=1$ this is $f=0$.

(3) The $K$-linear map $(V^\chi)^H\to V_+$, $(F_h)_{h\in H}\mapsto\sum_hF_h\cdot h^*v$, is injective. Indeed,
if $\sum_hF_h\cdot h^*v=0$, applying $g^*$ and using $g^*F_h=\chi(g)F_h$ gives
$\chi(g)\sum_h\big((hg)^*v\big)F_h=0$ for every $g$, and the first paragraph gives $F=0$. Comparing
dimensions gives $|H|\dim_KV^\chi\le\dim_KV_+$.
\end{proof}

\begin{lemma}\label{cv-lattices-galois-different}
Let $F'/F$ be a Galois extension of number fields, let $\mathfrak D_{F'/F}$ be its different and
$\mathfrak D_{F'}$ the different of $F'$ over $\mathbf Q$, and let $m\ge0$ be an integer. Let
$M\subset F'^m$ be an $\mathcal O_{F'}$-submodule such that $(s(x_1),\dots,s(x_m))\in M$ for every
$x\in M$ and every $s\in\operatorname{Gal}(F'/F)$. Then for every $\delta\in\mathfrak D_{F'/F}$ and every
$x\in M$, the vector $\delta x$ lies in the $\mathcal O_{F'}$-span of $M\cap F^m$. The same holds for
every $\delta\in\mathfrak D_{F'}$, since $\mathfrak D_{F'}=\mathfrak D_{F'/F}\,\mathfrak D_F\mathcal O_{F'}$ is
contained in $\mathfrak D_{F'/F}$.
\end{lemma}
\leanref{Development/AbelianVarietyTheory/Proofs/DivisorProofs/GaloisStableLatticeRationalSpanInclusionsStatement.lean}{GlobalTheory.T.galois_stable_lattice_rational_span_inclusions}
\leanref{Development/AbelianVarietyTheory/Proofs/DivisorProofs/DifferentIdealSmulMemSpanRational.lean}{GlobalTheory.differentIdeal_smul_mem_span_rational}
\begin{proof}
Write $\operatorname{Tr}$ for the trace of $F'/F$. As $\mathcal O_{F'}$ is a finitely generated projective
$\mathcal O_F$-module, there are $e_1,\dots,e_N\in\mathcal O_{F'}$ and $\mathcal O_F$-linear maps
$h_k\colon\mathcal O_{F'}\to\mathcal O_F$ with $a=\sum_kh_k(a)e_k$ for every $a\in\mathcal O_{F'}$. Extend
each $h_k$ to an $F$-linear form on $F'$; then $w=\sum_kh_k(w)e_k$ for every $w\in F'$. The trace form
is nondegenerate, so $h_k(w)=\operatorname{Tr}(c_kw)$ for some $c_k\in F'$, and $c_k$ lies in the inverse
different $\mathfrak D_{F'/F}^{-1}$, the set of $c$ with $\operatorname{Tr}(c\,\mathcal O_{F'})\subset\mathcal O_F$.
For $y,w\in F'$, multiplying $w=\sum_k\operatorname{Tr}(c_kw)e_k$ by $y$ and taking traces gives
$\operatorname{Tr}(yw)=\operatorname{Tr}\big(\sum_k\operatorname{Tr}(e_ky)c_k\,w\big)$, so
$y=\sum_k\operatorname{Tr}(e_ky)c_k$ by nondegeneracy. Now let $x\in M$ and $\delta\in\mathfrak D_{F'/F}$,
and put $t_k=(\operatorname{Tr}(e_kx_i))_{i\le m}\in F^m$. As $t_k=\sum_s(s(e_kx_i))_i$, the sum over
$\operatorname{Gal}(F'/F)$, and $M$ is a Galois-stable $\mathcal O_{F'}$-module, $t_k\in M\cap F^m$. Finally
$\delta c_k\in\mathcal O_{F'}$ and $\delta x=\sum_k(\delta c_k)t_k$.
\end{proof}

\begin{lemma}\label{cv-lattices-different-slope}
Let $K\subset K'$ be totally complex number fields, and let $\mathfrak D_{K'}$ and $d_{K'}$ be the
different and the discriminant of $K'$. Let $V$ be a $K$-vector space and $V'$ a $K'$-vector space
with Hermitian norm systems $N$ and $N'$ (Definition~\ref{cv-hermitian-norms}), such that
$N_{\bar\sigma}=N_\sigma$ for every embedding $\sigma:K\to\mathbf C$. Let $\iota:V\to V'$ be an additive
map with $\iota(ax)=a\,\iota(x)$ for $a\in K$ and $x\in V$, which carries every $K$-linearly
independent family to a $K'$-linearly independent family and satisfies
$N'_{\sigma'}(\iota(x))=N_{\sigma'|K}(x)$ for every embedding $\sigma':K'\to\mathbf C$ and every $x\in V$.
Let $L\subset V$ be an $\mathcal O_K$-submodule and $M\subset V'$ a nonzero finitely generated
$\mathcal O_{K'}$-submodule with $\iota(L)\subset M$, such that $\delta M$ is contained in the
$\mathcal O_{K'}$-span of $\iota(L)$ for every $\delta\in\mathfrak D_{K'}$. Then
\[
 \mu(M)-\frac{2\log|d_{K'}|}{[K':\mathbf Q]}\le\mu(L).
\]
\end{lemma}
\leanref{Development/AbelianVarietyTheory/Proofs/DivisorProofs/SlopeSqSubLeOfDifferentSmul.lean}{GlobalTheory.SlopeDifferent.slope_sq_sub_le_of_different_smul}
\begin{proof}
The map $\iota$ is injective: a nonzero $x\in V$ is a linearly independent family, so $\iota(x)\ne0$.
So $L$ embeds in $M$, which is finitely generated over $\mathcal O_K$, and $L$ is finitely generated. Let $L'$ be the
$\mathcal O_{K'}$-span of $\iota(L)$, so that $L'\subset M$. For a nonzero $\delta\in\mathfrak D_{K'}$,
$\delta M\subset L'$, so $M$ lies in the $K'$-span of $\iota(L)$; as $M\ne0$, $L\ne0$. Let
$j=\operatorname{rk}L$, let $e_1,\dots,e_j\in L$ be a basis of $KL$, and put $L_0=\sum_i\mathcal O_Ke_i$ and
$M_0=\sum_i\mathcal O_{K'}\iota(e_i)$. The vectors $\iota(e_i)$ are linearly independent over $K'$ and
span the $K'$-span of $\iota(L)$, so $\operatorname{rk}M=j$ and $M_0\subset L'\subset M$. As in the proof
of Lemma~\ref{cv-slope-base-change}(1), $[L':M_0]=[L:L_0]^{[K':K]}$. The ideal $\mathfrak D_{K'}$ has norm
$|d_{K'}|$, so $[M:\mathfrak D_{K'}M]=|d_{K'}|^j$, and $\mathfrak D_{K'}M\subset L'$ gives
$[M:L']\le|d_{K'}|^j$. By the norm hypothesis the Gram matrices satisfy
$G_{\sigma'}(\iota e)=G_{\sigma'|K}(e)$, so they agree at $\sigma'$ and at $\bar\sigma'$, and each
embedding $\sigma$ of $K$ has $[K':K]$ extensions $\sigma'$. The slope formula in the proof of
Lemma~\ref{cv-slope-base-change}, for $M$ with the vectors $\iota(e_i)$ and for $L$ with the $e_i$,
gives
\[
 j\,\mu(M)=\frac{2\log[M:L']+2[K':K]\log[L:L_0]-[K':K]\sum_\sigma\log\det G_\sigma(e)}{[K':\mathbf Q]}
 \le\frac{2j\log|d_{K'}|}{[K':\mathbf Q]}+j\,\mu(L).
\]
Dividing by $j$ gives the claim.
\end{proof}

\begin{lemma}\label{cv-height-components-numerics}
Let $I$ be a finite nonempty set. Let $\ell$, $D$, $A$, $c_s$, $X$, $c_c$, $\zeta$, $\rho$, $R$, $\beta$,
$\delta'$, $E$, $C_1$, $C_2$ and $\mu$ be real numbers with $D\ge0$, $\ell\ge3D+5$, $A>0$,
$c_s,X,c_c,R\ge0$, $0\le\zeta<1$, $0\le\rho<1$, $\beta>0$, $\delta'>0$ and $E\ge1$. Let $h$, $s$,
$\delta$, $r$ and $g$ be real functions on $I$ such that for every $Z\in I$:
\begin{enumerate}
\item $\delta'\le\delta_Z\le\delta'/(1-\rho)$;
\item $\beta\delta_Z\le r_Z$ and $(1-\zeta)r_Z\le\beta\delta_Z$;
\item $h_Z+c_s\ge0$ and $(1-\rho)\max(0,s_Z+c_s)\le h_Z+c_s+c_c$;
\item $\beta\delta_Z\big(\ell h_Z-3(D+1)|h_Z|-A\log(2+|h_Z|)-R\big)\le g_Z$.
\end{enumerate}
Suppose moreover that $\frac1{|I|}\sum_Zs_Z\ge-X$, that
\[
 \ell X\le C_1E,\quad \ell c_c\le C_1E,\quad c_s\big(6D+10+\ell(\zeta+2\rho)\big)\le C_1E,\quad
 A\log\max(A,1)+R\le C_2E,
\]
and that $2\sum_Zg_Z/\sum_Zr_Z\le\mu$. Then $\mu\ge-2(3C_1+C_2+4)E$.
\end{lemma}
\leanref{Development/AbelianVarietyTheory/Proofs/DivisorProofs/UniformLowerBoundOfSums.lean}{GlobalTheory.CovTop.FromComp.core}
\leanref{Formalism/MeanHeightDescent/DivisionInequalityWithCovolume.lean}{Formalism.ExtensionFiltration.mean_height_sum_bound}

\subsection{Uniformization and markings}\label{subsection-cv-uniformization-markings}

\begin{definition}\label{cv-deep-cover}
In Situation~\ref{cv-cusp-setting}, let $w=(1+\sqrt{-d})/2\in\mathcal O_E$. A \emph{level} is a pair
$(M',\mathfrak M')$ of an odd integer $M'\ge3$ and an ideal $\mathfrak M'\subset\mathcal O_E$ generated
by elements of $\mathbf Z[w]$, with $\mathfrak M'\bar{\mathfrak M}'=M'\mathcal O_E$ and
$\mathfrak M'+\bar{\mathfrak M}'=\mathcal O_E$; it stands for the level $(M',\mathfrak M',1)$ of full
level $1$ of Definition~\ref{cv-model-toroidal-semiabelian}(1), and all levels written as pairs have
full level $1$. Suppose that the level of $X$ is a level
$(M,\mathfrak M)$ with $M$ prime to $15$. A \emph{framed point} of a PEL moduli scheme over a field $k$
is a $k$-point $y$ together with a basis of every eigenspace of $\mathcal O_E$ on the Lie algebra of the
abelian variety $A_y$ of $y$ and a generator of the one-dimensional space of sections of the line
bundle of the polarization on $A_y$. Let $e\ge0$ be an integer and $o$ a real number. A \emph{deep
cover} of $X$ of exponent $e$ and size $o$ consists of the following data, with the following
properties.
\begin{enumerate}
\item (Deep level.) An ideal $I\subset\mathcal O_E$ with $I\bar I=15^e\mathcal O_E$ and
  $I+\bar I=\mathcal O_E$, such that $(M_Y,\mathfrak M_Y)=(15^eM,\mathfrak MI)$ is a level, the
  \emph{deep level}.
\item (Coefficient field.) A number field $F_Y\subset\mathbf C$ that contains $F_c$ and the $M_Y$-th
  roots of unity, that every automorphism of $\mathbf C$ fixing $F_c$ maps into itself, and that
  satisfies $\log[F_Y:\mathbf Q]+\log\operatorname{rd}_{F_Y}\le o$.
\item (Moduli scheme.) The PEL moduli scheme $Y$ of the PEL datum of $X$ at the deep level, with
  coefficient field $F_Y$. The field $F_Y$ is algebraically closed in the function field of every
  connected component of $Y\otimes F_Y$.
\item (Level map.) A morphism $\pi\colon Y\to X$ and a homomorphism, over $\pi$, from the universal
  abelian scheme over $Y$ to the one over $X$ that identifies the first with the pullback of the
  second, compatibly with the actions of $\mathcal O_E$ and the polarizations, and under which $15^e$
  times the level structure of $Y$ is the level structure of $X$. Every framed point of $X$ over an
  algebraically closed field containing $F_Y$ lifts, through this identification, to a framed point
  of $Y$ over the same field.
\item (Uniformization.) For every good embedding $\sigma\colon F_Y\to\mathbf C$, one that sends the
  square root of $-d$ given by the datum to the square root fixed by $\mathcal D$: a uniformization of
  $Y$ at $\sigma$, normalized by $\xi$, that satisfies the compatibility conditions of the moduli
  problem and in which every component has the Hermitian matrix of $W_j$; and a map
  $g\mapsto(Z(g),x(g))$ from $U(W_j)(\mathbf A_F)$ to frames of components with
  $|\delta_{\tau,i}(x(g))|=1$ for all $\tau$ and $i$. Let $K_Y\subset U(W_j)(\mathbf A_{F,f})$ be the
  principal congruence subgroup of level $M_Y$ of the stabilizer of the lattice of $W_j$. For
  $\gamma\in U(W_j)(F)$ and $k\in K_Y$, the elements $\gamma g$ and $gk$ have the component $Z(g)$ and
  frames in $\Gamma_{Z(g)}x(g)$. Right translation by $h\in U(W_j)(F\otimes\mathbf R)$ corresponds,
  up to $\Gamma_{Z(g)}$, to right multiplication of $x(g)$ by $h$ written in a base frame of $Z(g)$. If
  $Z(g')=Z(g)$ and $x(g')\in\Gamma_{Z(g)}x(g)$, then $g'\in U(W_j)(F)\,g\,K_Y$. Every frame of a
  component $Z$ lies in $\Gamma_Z\,x(g)B$ for some $g$ with $Z(g)=Z$. The data at all good embeddings
  are compatible with one identification of the PEL datum of $X$ with the lattice of $W_j$, with
  orthonormal base frames.
\end{enumerate}
\end{definition}
\leanref{Development/AbelianVarietyTheory/Definitions/Divisors/NormNegCmConstantEqAtAllEmbeddings.lean}{GlobalTheory.CovTop.Deep}
\leanref{Development/AbelianVarietyTheory/Definitions/Divisors/NormNegCmConstantEqAtAllEmbeddings.lean}{GlobalTheory.CovTop.Deep.Good}
\leanref{Formalism/LocalToGlobal/ProjectiveLattices/ThShort.lean}{Formalism.L2.Pj.DeltaDat}
\leanref{Formalism/LocalToGlobal/ProjectiveLattices/ThShort.lean}{Formalism.L2.Pj.DeltaDat.IsLev}
\leanref{ReciprocityTheory/ReductiveGroups/ArithmeticGeometry/BirationalGeometry/ToroidalCompactification/TwoWeightEqOneAddRoot.lean}{ReciprocityTheory.ChartRestriction.G.Geo.Tor.Level}
\leanref{ReciprocityTheory/ReductiveGroups/ArithmeticGeometry/BirationalGeometry/ToroidalCompactification/PushInvertible.lean}{ReciprocityTheory.ChartRestriction.G.Geo.Tor.Coeff}
\leanref{Formalism/LocalToGlobal/ProjectiveLattices/CompHt.lean}{Formalism.L2.Pj.FnAlgClosed}
\leanref{Formalism/LocalToGlobal/ProjectiveLattices/CuspidalL2Subspace.lean}{Formalism.L2.Pj.IsLevelK}

\begin{definition}\label{cv-uniformization-model}
In Situation~\ref{cv-cusp-setting}, write $b\mapsto\bar b$ for complex conjugation on $E$, and for
$G\in M_r(E)$ and $u,v\in E^r$ put $h_G(u,v)=\sum_{a,b}u_aG_{ab}\bar v_b$. Let $L'$ be the lattice of a
PEL datum over $E$: a finitely generated projective $\mathcal O_E$-module with a perfect alternating form
$\langle\ ,\ \rangle'$ as in Definition~\ref{cv-cm-point-periods}.
\begin{enumerate}
\item A \emph{Hermitian model} of $L'$ is a triple $(G,\eta,\varphi)$ of an invertible matrix
  $G\in M_r(E)$ with $\bar G_{ab}=G_{ba}$ for all $a,b$, an element $\eta\in E^\times$ with
  $\bar\eta=-\eta$, and an injective $\mathcal O_E$-linear map $\varphi\colon L'\to E^r$ whose image
  spans $E^r$ over $E$, such that
  \[
   \langle x,y\rangle'=\operatorname{Tr}_{E/\mathbf Q}\big(\eta\,h_G(\varphi x,\varphi y)\big)\qquad(x,y\in L').
  \]
\item Let $(M',\mathfrak M')$ be a level (Definition~\ref{cv-deep-cover}) and put $\Lambda'=\varphi(L')$.
  The \emph{principal congruence group} $\Gamma(\varphi)$ of the model is the group of
  $\gamma\in GL_r(E)$ with $h_G(\gamma u,\gamma v)=h_G(u,v)$ for all $u,v\in E^r$, with
  $\gamma\Lambda'=\Lambda'$, and with $\gamma\varphi(a)\in\varphi(a)+M'\Lambda'$ for every
  $a\in\bar{\mathfrak M}'L'$.
\end{enumerate}
\end{definition}
\leanref{Development/AbelianVarietyTheory/Definitions/ExistenceOfLattices/NormConst.lean}{GlobalTheory.FrameUnif.LatModel}
\leanref{Development/AbelianVarietyTheory/Definitions/ExistenceOfLattices/NormConst.lean}{GlobalTheory.FrameUnif.princ}

\begin{lemma}\label{cv-uniformization-signature}
In Situation~\ref{cv-cusp-setting}, let $L'$ be the lattice of a PEL datum of signature $(p',q')$ over
$E$, let $(G,\eta,\varphi)$ be a Hermitian model of $L'$ (Definition~\ref{cv-uniformization-model}), and
let $\tau\in\Phi$. If $\operatorname{Im}\tau(\eta)>0$, the Hermitian matrix $\tau(G)$ has $q'$ positive
and $p'$ negative eigenvalues. If $\operatorname{Im}\tau(\eta)<0$, it has $p'$ positive and $q'$
negative eigenvalues.
\end{lemma}
\leanref{Development/AbelianVarietyTheory/Proofs/UnitaryExistence/SignatureHermitianModelOfDatumStatement.lean}{GlobalTheory.T.signature_hermitian_model_of_datum}
\begin{proof}
Let $\theta$ be the square root of $-d$ fixed by $\mathcal D$. The lattice $L'$ has a period $P$ of a
type $(\sigma_k)$ of signature $(p',q')$ for $\theta$, with a positive definite Hermitian form $H$
(Definition~\ref{cv-cm-point-periods}). Multiplication by $i$ on $\mathbf C^g$, carried to
$L'\otimes\mathbf R$ by $P$, is a complex structure that commutes with $\mathcal O_E$. Through $\varphi$
and $\tau$ it gives a $\mathbf C$-linear map $J$ of $\mathbf C^r$ with $J^2=-1$. Its eigenspaces for $i$
and for $-i$ have dimensions $p'$ and $q'$, since $\mathcal O_E$ acts on the coordinate $k$ of
$\mathbf C^g$ through $\sigma_k$, and the type takes the value $\tau$ $p'$ times and the value $\bar\tau$
$q'$ times. As $\tau(\eta)=i\operatorname{Im}\tau(\eta)$, the identity
$\operatorname{Im}H(Px,Py)=\langle x,y\rangle'=\operatorname{Tr}_{E/\mathbf Q}(\eta\,h_G(\varphi x,\varphi y))$
shows that the two eigenspaces are orthogonal for $\tau(G)$, and that $\operatorname{Im}\tau(\eta)\,\tau(G)$
is negative definite on the first and positive definite on the second. This gives both signatures.
\end{proof}
% Lean: HermSign.signature_of_model, from exists_complexStructure (the map J with J^2 = -1, the
% positivity, and trace i(p'-q')) and exists_diag_of_complexStructure; the period is
% PeriodExists.typeFn_period_exist.

\begin{lemma}\label{cv-uniformization-lattice}
In Situation~\ref{cv-cusp-setting}, let $L$ with the form $\langle\ ,\ \rangle$ be the lattice of the PEL
datum of $X$, of signature $(p,q)$. Let $G_W$ be the Gram matrix of the basis of $W_j$, let
$h_{W_j}$ be the Hermitian form of $W_j$ in the coordinates of this basis, put $G=-G_W^{\mathsf T}$, and
let $L_W\subset E^r$ be the $\mathcal O_E$-span of $\mathfrak AL_0$, for the lattice $L_0$ of $W_j$ in
these coordinates and the ideal $\mathfrak A$ fixed by the host. Let $\xi_0\in E$ be the element of
the host. Suppose that $\iota\colon L\to E^r$ is an injective $\mathcal O_E$-linear map with image $L_W$
and that $\langle x,y\rangle=\epsilon\operatorname{Tr}_{E/\mathbf Q}(\xi_0\,h_{W_j}(\iota x,\iota y))$ for
all $x,y\in L$, for a sign $\epsilon$. Then $\epsilon=1$, and $(G,\xi_0,\iota)$ is a Hermitian model of
$L$ (Definition~\ref{cv-uniformization-model}).
\end{lemma}
\leanref{Development/AbelianVarietyTheory/Proofs/UnitaryExistence/ExistsLatticeEmbeddingLiftDatumHermitianSpace.lean}{GlobalTheory.Member.exists_iotaOK}
\leanref{Development/AbelianVarietyTheory/Proofs/UnitaryExistence/LiftDatumAlternatingFormSignEqOneStatement.lean}{GlobalTheory.T.liftDatum_alternatingForm_sign_eq_one}
\begin{proof}
The lattice $L_W$ spans $E^r$, so $\iota$ is a coordinate system for $L$. For $u,v\in E^r$,
$h_G(u,v)=-h_{W_j}(v,u)$; since $h_{W_j}$ is Hermitian and $\bar\xi_0=-\xi_0$, this gives
$\operatorname{Tr}_{E/\mathbf Q}(\xi_0h_G(u,v))=\operatorname{Tr}_{E/\mathbf Q}(\xi_0h_{W_j}(u,v))$. So
$(G,\epsilon\xi_0,\iota)$ is a Hermitian model of $L$. Suppose $\epsilon=-1$. For the embedding
$\varphi_0\in\Phi$ of the host, $\operatorname{Im}\varphi_0(-\xi_0)<0$, so by
Lemma~\ref{cv-uniformization-signature} the matrix $\varphi_0(G)$ has $p$ positive eigenvalues. The
signature of $W_j$ at $\varphi_0$ gives $q$ positive eigenvalues for $\varphi_0(G)$. Hence $p=q$, which
contradicts $2q<r=p+q$. So $\epsilon=1$.
\end{proof}
% Lean: the hypothesis is the pinning of the datum to W_j in T.Q.Pins (fields phi, eps, hinj, hrange,
% hiso); exists_iotaOK composes phi with the coordinates; GlobalTheory.Member.sign_eq_one uses
% HermSign.signature_of_model and RecDict.posIndex_gmW; 2q < r is part of exists_iotaOK.

\begin{lemma}\label{cv-uniformization-models}
In Situation~\ref{cv-cusp-setting}, let $(M',\mathfrak M')$ be a level of full level $1$
(Definition~\ref{cv-deep-cover}), let $L$ be the lattice of a PEL
datum of signature $(p,q)$ over $E$, and let $(G,\eta,\varphi)$ be a Hermitian model of $L$ with
$\operatorname{Im}\tau(\eta)>0$ for every $\tau\in\Phi$ (Definition~\ref{cv-uniformization-model}).
Then every lattice $L'$ of a PEL datum of signature $(p,q)$ over $E$ has a Hermitian model
$(G,\eta,\varphi')$ with the same $G$ and $\eta$. For every such model, with $\Lambda'=\varphi'(L')$:
\begin{enumerate}
\item there is $c\in\mathcal O_E\setminus\{0\}$ such that $c\,h_G(u,v)$ is integral for all
  $u,v\in\Lambda'$;
\item the group $\Gamma(\varphi')$ has finite index in the group of $\gamma\in GL_r(E)$ that preserve
  $h_G$ and $\Lambda'$;
\item if $\gamma\in\Gamma(\varphi')$, $x$ is a frame for $G$ (a frame of $Z$ in
  Situation~\ref{cv-cusp-setting}, with $G$ in place of $G_Z$), $b\in B$ and $\gamma x=xb$, then
  $\gamma=1$;
\item there is an integer $m\ge1$ such that, for every $u\in\Lambda'$ with $h_G(u,u)=0$, every
  totally real integer $s$ of $E$ and $\kappa=\pm2\eta$, the transvection $v\mapsto v+ms\kappa\,h_G(v,u)\,u$
  lies in $\Gamma(\varphi')$.
\end{enumerate}
\end{lemma}
\leanref{Development/AbelianVarietyTheory/Proofs/UnitaryExistence/ExistsFlagDomainModelsOfFamilyUniformisation.lean}{GlobalTheory.ModelsExist.exists_modelsOf}
\leanref{Development/AbelianVarietyTheory/Proofs/UnitaryExistence/LatticeModelOfSameSignatureHasseNormStatement.lean}{GlobalTheory.T.latticeModel_of_same_signature_hasseNorm}
\begin{proof}
The datum of $L'$ has a Hermitian model $(G',\eta,\psi)$ for some matrix $G'$, obtained by choosing
coordinates on $L'\otimes\mathbf Q$ and solving for the Hermitian form whose trace against $\eta$ is
$\langle\ ,\ \rangle'$. By Lemma~\ref{cv-uniformization-signature} the matrices $\tau(G)$ and $\tau(G')$
have the same signature at every embedding $\tau$, since $L$ and $L'$ have the same signature and the
models share $\eta$. Both forms are self-dual for the trace pairing against $\eta$ on lattices of
the same rank, so the quotient $\det G'/\det G$, which lies in the totally real subfield $E^+$, is a
local norm from $E$ at every place of $E^+$: at the real places by the equality of signatures, at the
finite places by the self-duality. By the Hasse norm theorem for the cyclic extension
$E/E^+$ \cite[Chapter~VI, (4.5)]{Neu99} it is a norm, $\det G'=z\bar z\det G$. Landherr's
theorem \cite{Lan36} then gives $\gamma\in GL_r(E)$ with $h_G(\gamma u,\gamma v)=h_{G'}(u,v)$, and
$\varphi'=\gamma\psi$ is the required model.

(1) As $\eta\ne0$, some nonzero integer $n$ makes $n\eta$ integral. Since
$\operatorname{Tr}_{E/\mathbf Q}(\eta h_G(u,v))\in\mathbf Z$ on $\Lambda'$ and $\Lambda'$ is an
$\mathcal O_E$-module, $\eta h_G(u,v)$ lies in the inverse different, so $c=d_En\eta$ works.

(2) The group of isometries of $h_G$ preserving $\Lambda'$ acts on the finite set
$\Lambda'/M'\Lambda'$, and $\Gamma(\varphi')$ contains the kernel of this action.

(3) Since $\gamma$ preserves $\langle\ ,\ \rangle'$ and $\mathfrak M'+\bar{\mathfrak M}'=\mathcal O_E$
with $\mathfrak M'\bar{\mathfrak M}'=M'\mathcal O_E$, the congruence $\gamma\equiv1$ modulo $M'$ on the
part $\bar{\mathfrak M}'\Lambda'$ gives $\gamma\equiv1$ modulo $M'\Lambda'$ on all of $\Lambda'$. The
relation $\gamma x=xb$ says that $\gamma$ fixes the flag of $x$ at every $\tau\in\Phi$, so $\gamma$
lies in a compact subgroup of $\prod_{\tau\in\Phi}U(\tau(G))$; preserving the lattice $\Lambda'$, it has
finite order. An automorphism of finite order of a lattice that is congruent to $1$ modulo an integer
$M'\ge3$ is the identity \cite{Min87}.

(4) Put $\alpha=ms\kappa$. Expanding $h_G$ at the images of $v$ and $w$ leaves
$h_G(v,w)+(\alpha+\bar\alpha)h_G(v,u)h_G(u,w)$, since $h_G(u,u)=0$; and $\bar\alpha=-\alpha$ because $s$
is totally real and $\bar\eta=-\eta$. So the transvection preserves $h_G$. As $\eta h_G(v,u)$ lies in
the inverse different for $u,v\in\Lambda'$, a multiple $m$ of $M'd_E$ makes $\alpha h_G(v,u)\in
M'\mathcal O_E$, so the transvection and its inverse $v\mapsto v-\alpha h_G(v,u)u$ preserve $\Lambda'$
and are congruent to $1$ modulo $M'\Lambda'$.
\end{proof}
% Lean: latModel_of_same_signature: the model is GlobalTheory.model_in_same_space (with
% exists_herm_coord_of_datum, signature_hermitian_model_of_datum, DetIdeal.det_ideal for the local
% norms, LocalQuad.landherrSt_of_hasseNormSt and the published HasseNormSt); (1) exists_den_herm_model;
% (2) principalLevelGroup_finiteIndex_unitaryStabiliser; (3) principalGroup_acts_freely_on_flags
% (LevelHalf.trivial_mod_of_trivial_on_half, GamFree.gam_free with latMinkowski);
% (4) transvection_mem_principalCongruenceGroup. With l.A = 1 the Lean's second congruence is empty.
% exists_modelsOf assembles the models of all lattices of a uniformization family, with trivial
% characters chi = 1 and kappa0 = 2 xi0.

\begin{theorem}\label{cv-uniformization-family}
In Situation~\ref{cv-function-situation}, let $L$ be the lattice of the PEL datum, of signature $(p,q)$ with $p+q\ge1$,
and let $\theta$ be the square root of $-d$ of the datum. Let $l=(M',\mathfrak M',N)$ be a level of the datum, with
full level $N$, let $X'$ be the moduli scheme at level $l$ over $K_l=\mathbf Q(\sqrt{-d},\zeta_N)$, with level structure
$\eta$ and full level structure $\eta'$ (Definition~\ref{cv-model-toroidal-semiabelian}(1)), and put
$M''=\operatorname{lcm}(M',N)$. Let $\sigma\colon K_l\to\mathbf C$ be an embedding with $\sigma(\sqrt{-d})=\theta$, and
let $k\mapsto\sigma_k$, $1\le k\le g$, be a type of signature $(p,q)$ for $\theta$
(Definition~\ref{cv-cm-point-periods}). Then there are lattices $L_1,\dots,L_n$ of PEL data of signature $(p,q)$ over
$E$, with $n\ge1$, and $\mathcal O_E$-linear maps $t_i\colon L\to L_i$ that induce isomorphisms
$L/M''L\to L_i/M''L_i$, and for every $i$ and every period $P$ of $L_i$ of this type a complex point $y_i(P)$ of $X'$
over $\sigma$ and a surjective homomorphism $\exp_{i,P}\colon\mathbf C^g\to A_{y_i(P)}(\mathbf C)$ with kernel $P(L_i)$,
such that:
\begin{enumerate}
\item[(a)] $\exp_{i,P}\big((\sigma_k(b)z_k)_k\big)=b\cdot\exp_{i,P}(z)$ for $b\in\mathcal O_E$ and $z\in\mathbf C^g$;
\item[(b)] for every integer $m\ge1$, all $a,b\in L_i$, and the commutator pairing $e_m$ of the theta group of
  $\mathcal L^m$ (Definition~\ref{cv-cusp-cm-point-eigen-degree-rank-one-norm-weil-translate-group}), where $\mathcal L$ is
  the line bundle of the polarization,
  $e_m\big(\exp_{i,P}(P(a)/m),\exp_{i,P}(P(b)/m)\big)=\exp(-2\pi i\langle a,b\rangle/m)$;
\item[(c)] the level structure of $y_i(P)$ sends $a\in\bar{\mathfrak M}'L$ to $\exp_{i,P}(P(t_ia)/M')$, and its full
  level structure sends $a\in L$ to $\exp_{i,P}(P(t_ia)/N)$.
\end{enumerate}
By (3) and (5) below, $\exp_{i,P}$ induces an isomorphism of complex Lie groups $\mathbf C^g/P(L_i)\to A_{y_i(P)}$,
compatible with the actions of $\mathcal O_E$; it is \emph{the isomorphism} of the theorem. Moreover:
\begin{enumerate}
\item every complex point of $X'$ over $\sigma$ is $y_i(P)$ for some $i$ and $P$, and the point
  determines $i$;
\item $y_i(P)=y_i(P')$ if and only if $P'\circ\gamma=h\circ P$ for a $\mathbf C$-linear automorphism $h$
  of $\mathbf C^g$ and an $\mathcal O_E$-linear automorphism $\gamma$ of $L_i$ that preserves its form and satisfies
  $\gamma(a)\in a+M'L_i$ for $a\in\bar{\mathfrak M}'L_i$ and $\gamma(a)\in a+NL_i$ for $a\in L_i$;
\item for an open set $V\subset\mathbf C^m$ and a holomorphic family of periods $P_b$, $b\in V$, of $L_i$ of this type
  (each $P_b$ a period, and $b\mapsto P_b(a)$ holomorphic for every $a\in L_i$), the point $y_i(P_b)$ depends
  holomorphically on $b$ and $\exp_{i,P_b}(z)$ on $(b,z)\in V\times\mathbf C^g$; and every holomorphic family of complex
  points of $X'$ over $\sigma$, on an open set $V\subset\mathbf C^m$, is, near each of its points, $y_i(P_b)$ for some
  $i$ and a holomorphic family of periods $P_b$ of $L_i$ of this type;
\item let $\mathcal L$ be the line bundle of the polarization on the universal abelian scheme
  $\mathcal A$ over $X'$, with its rigidification along the zero section, and fix trivializations of
  $\mathcal L$ over the sets $U_j$ of an open cover of $\mathcal A$: the values of a section in the
  trivializations over $U_j$ and $U_{j'}$ satisfy $s_j=g_{jj'}s_{j'}$ for transition functions $g_{jj'}$,
  and the rigidification has the value $\varrho_j(y)\ne0$ at a point $y$ of $X'$ whose zero lies in
  $U_j$. Let $V\subset\mathbf C^m$ be an open set that is connected or empty, $P_b$, $b\in V$, a holomorphic family of
  periods of $L_i$ of this type, and $p_b(z)=\exp_{i,P_b}(z)\in A_{y_i(P_b)}$, the image of $z\in\mathbf C^g$ under the
  isomorphism above. Let $(e_1,\dots,e_g,f_1,\dots,f_g)$ be a basis of $L_i$ with
  $\langle e_k,e_l\rangle=\langle f_k,f_l\rangle=0$ and $\langle e_k,f_l\rangle=\epsilon\,\delta_{kl}$ for a
  sign $\epsilon$, and let $h_b,Z_b\in M_g(\mathbf C)$, $b\in V$, be matrices with $\operatorname{Im}Z_b$
  positive definite, $h_bP_b(e_k)$ the $k$-th unit vector and $h_bP_b(f_k)$ the $k$-th column of $Z_b$.
  Then there are a characteristic $c\in\{0,\frac12\}^{2g}$ and, for every $j$, a nowhere vanishing
  holomorphic function $F_j$ on the set of $(b,z)\in V\times\mathbf C^g$ with $p_b(z)\in U_j$, such that
  \[
   F_j(b,z)=g_{jj'}(p_b(z))\,F_{j'}(b,z)\quad\text{if }p_b(z)\in U_j\cap U_{j'},\qquad
   F_j(b,0)=\varrho_j(y_i(P_b)),
  \]
  and such that for every $b\in V$ the line bundle $\mathcal L$ on $A_{y_i(P_b)}$ has a section whose
  value at $p_b(z)$ in the trivialization over $U_j$ is $\vartheta[c](Z_b,h_bz)\,F_j(b,z)$. The family
  $F=(F_j)$ is the \emph{factor} of this section.
\item every vector $w\in\mathbf C^g$ is the derivative at $0$ of $\exp_{i,P}$ in the direction $w$ for exactly one
  tangent vector of $A_{y_i(P)}$ at the origin, and every tangent vector of $A_{y_i(P)}$ at the origin is this
  derivative for exactly one $w\in\mathbf C^g$.
\end{enumerate}
Data $L_i$, $t_i$, $y_i(P)$ and $\exp_{i,P}$ with these properties are a \emph{family uniformization} of $X'$ at
$\sigma$ for the type $(\sigma_k)$.
\end{theorem}
\leanref{ReciprocityTheory/ReductiveGroups/ArithmeticGeometry/BirationalGeometry/ToroidalCompactification/IsHolPeriods.lean}{ReciprocityTheory.ChartRestriction.G.Geo.Tor.N1c}
\leanref{ReciprocityTheory/ReductiveGroups/ArithmeticGeometry/BirationalGeometry/ToroidalCompactification/IsHolPeriods.lean}{ReciprocityTheory.ChartRestriction.G.Geo.Tor.FamilyUnif}
\leanref{Development/AbelianVarietyTheory/Proofs/ModuliProofs/CmModuliConsequenceOfAbelianSchemeLaws.lean}{GlobalTheory.closing_of_abelianSchemeLaws}

This is the complex uniformization of PEL moduli schemes by Shimura's families of abelian varieties
\cite{Shi63}, in the adelic form of \cite[\S8]{Kot92b}, with the theta functions of the polarization
\cite[Chapter~8]{BL04}.

\begin{lemma}\label{cv-uniformization-adapted-basis}
Let $\Lambda$ be a free $\mathbf Z$-module of rank $2g$ with a perfect alternating form
$B\colon\Lambda\times\Lambda\to\mathbf Z$, and let $N\ge1$ be an integer. Let $P,Q\subset\Lambda$ be
subgroups that contain $N\Lambda$, with $P+Q=\Lambda$, such that $B(x,y)$ is divisible by $N$ whenever
$x,y\in P$ and whenever $x,y\in Q$. Then $\Lambda$ has a basis $(e_1,\dots,e_g,f_1,\dots,f_g)$ with
$B(e_a,e_b)=B(f_a,f_b)=0$ and $B(e_a,f_b)=\delta_{ab}$ for all $a,b$, such that an element of $\Lambda$
lies in $P$ if and only if its coordinates at $f_1,\dots,f_g$ are divisible by $N$, and lies in $Q$ if
and only if its coordinates at $e_1,\dots,e_g$ are divisible by $N$.
\end{lemma}
\leanref{Development/AbelianVarietyTheory/Proofs/FoundationProofs/ExistsIsSymplecticBasisToMatrixOffDiagDvd.lean}{GlobalTheory.exists_isSymplecticBasis_adapted_dvd}
\leanref{Development/AbelianVarietyTheory/Proofs/FoundationProofs/ExistsIsSymplecticBasisToMatrixOffDiagDvd.lean}{GlobalTheory.IsSymplecticBasis.exists_adapted_mod}
\leanref{Development/AbelianVarietyTheory/Proofs/FoundationProofs/ExistsIsSymplecticBasisToMatrixOffDiagDvd.lean}{GlobalTheory.IsSymplecticBasis.exists_adapted}
\leanref{Development/AbelianVarietyTheory/Proofs/FoundationProofs/ExistsIsSymplecticBasisToMatrixOffDiagDvd.lean}{GlobalTheory.exists_mul_mul_eq_one_of_isUnit_sub_transpose}
\begin{proof}
A basis with these values of $B$ is a \emph{symplectic basis}. One exists: as $B$ is perfect, a primitive
vector $e_1$ has a vector $f_1$ with $B(e_1,f_1)=1$; the orthogonal of $\mathbf Ze_1+\mathbf Zf_1$ is a
complement of it on which $B$ is again perfect, and induction on $g$ finishes. Fix a symplectic basis $b$
and put $R=\mathbf Z/N\mathbf Z$. Reducing the coordinates in $b$ gives a surjection
$\pi\colon\Lambda\to R^{2g}$ with kernel $N\Lambda$, and $\bar B(\pi x,\pi y)=B(x,y)$ modulo $N$ defines an
alternating form $\bar B$ on $R^{2g}$ for which $\bar b=\pi(b)$ is a symplectic basis. The images
$\bar P=\pi(P)$ and $\bar Q=\pi(Q)$ are isotropic for $\bar B$, and $\bar P+\bar Q=R^{2g}$. Since $P$ and
$Q$ contain the kernel of $\pi$, an element $x$ lies in $P$ if and only if $\pi x\in\bar P$, and likewise
for $Q$.

Step 1: an adapted symplectic basis of $R^{2g}$. Write $\bar b_k=p_k+q_k$ with $p_k\in\bar P$ and
$q_k\in\bar Q$ for $1\le k\le2g$, and let $G$ be the $2g\times2g$ matrix with entries $\bar B(p_k,q_l)$.
As $\bar P$ and $\bar Q$ are isotropic, $\bar B(\bar b_k,\bar b_l)=G_{kl}-G_{lk}$, so $G-G^{\mathsf T}$ is
the Gram matrix of a symplectic basis and is invertible. For each prime $\ell$ dividing $N$ the reduction
of $G$ modulo $\ell$ therefore has rank at least $g$, since $2g=\operatorname{rank}(G-G^{\mathsf T})\le
2\operatorname{rank}G$ over $\mathbf F_\ell$; so there are matrices $U_\ell$ of size $g\times2g$ and
$V_\ell$ of size $2g\times g$ over $\mathbf F_\ell$ with $U_\ell GV_\ell=1$. By the Chinese remainder
theorem there are matrices $U$ and $V'$ over $R$ that reduce to $U_\ell$ and $V_\ell$ for every $\ell$.
Then $\det(UGV')$ is a unit modulo every $\ell$, hence a unit of $R$, and $V=V'(UGV')^{-1}$ gives
$UGV=1$. Put $x_a=\sum_kU_{ak}p_k\in\bar P$ and $y_a=\sum_kV_{ka}q_k\in\bar Q$ for $1\le a\le g$. Then
$\bar B(x_a,y_b)=(UGV)_{ab}=\delta_{ab}$, and $\bar B(x_a,x_b)=\bar B(y_a,y_b)=0$ because $\bar P$ and
$\bar Q$ are isotropic. So the matrix of the family $(x,y)$ in the basis $\bar b$ is symplectic; it has
determinant $1$, and $(x,y)$ is a symplectic basis of $R^{2g}$.

Step 2: lifting. The matrix of $(x,y)$ in $\bar b$ lies in $Sp_{2g}(R)$, and the reduction map
$Sp_{2g}(\mathbf Z)\to Sp_{2g}(\mathbf Z/N\mathbf Z)$ is surjective \cite[Theorem~1]{NS64}. A lift
$T\in Sp_{2g}(\mathbf Z)$ is the matrix in $b$ of a symplectic basis $(e,f)$ of $\Lambda$ with
$\pi(e_a)=x_a$ and $\pi(f_a)=y_a$.

Step 3: the two conditions. The coordinate at $y_a$ of a vector $\bar v\in R^{2g}$ in the basis $(x,y)$ is
$\bar B(x_a,\bar v)$. For $\bar v\in\bar P$ it vanishes, as $\bar P$ is isotropic; conversely, if these
coordinates vanish for all $a$, then $\bar v$ is a combination of the $x_a$ and lies in $\bar P$. So
$\bar v\in\bar P$ if and only if its coordinates at the $y_a$ vanish, and in the same way $\bar v\in\bar Q$
if and only if its coordinates at the $x_a$ vanish. The coordinates of $\pi(v)$ in $(x,y)$ are the
reductions modulo $N$ of those of $v$ in $(e,f)$, and the last sentence of the first paragraph gives the
two conditions.
\end{proof}
% Lean: exists_isSymplecticBasis_adapted_dvd takes the symplectic basis from
% exists_isSymplecticBasis_of_finrank and applies IsSymplecticBasis.exists_adapted_mod (reduction phi,
% exists_lift_adapted = exists_adapted over ZMod N, then exists_lift with symplecticGroupMap_surjective);
% exists_adapted is Step 1, with exists_mul_mul_eq_one_of_isUnit_sub_transpose (rank at each maximal
% ideal, lifted by exists_matrix_lift_maximalSpectrum) and exists_basis_of_toMatrix_mem; Step 3 is
% mem_iff_repr_inr and mem_iff_repr_inl with repr_semilinear_of_map_basis.

\begin{lemma}\label{cv-uniformization-mod-eight}
Let $E$ be a CM field with complex conjugation $b\mapsto\bar b$, let $d\equiv7\pmod{16}$ be a prime and
$\sqrt{-d}\in\mathcal O_E$ an element with square $-d$, and put $w=(1+\sqrt{-d})/2$, which lies in
$\mathcal O_E$ since $d\equiv3\pmod4$. Let $L$ be a finitely generated projective $\mathcal O_E$-module of
rank $n$ with a perfect alternating form $\langle\ ,\ \rangle\colon L\times L\to\mathbf Z$ such that
$\langle bx,y\rangle=\langle x,\bar by\rangle$ for $b\in\mathcal O_E$, and put $g=n[E:\mathbf Q]/2$. Then
$L$ has a $\mathbf Z$-basis $(e_1,\dots,e_g,f_1,\dots,f_g)$ with
\[
 \langle e_a,e_b\rangle=\langle f_a,f_b\rangle=0,\qquad\langle e_a,f_b\rangle=-\delta_{ab},\qquad
 w^4e_k\in e_k+8L,\qquad w^4f_k\in8L
\]
for all $a$, $b$ and $k$.
\end{lemma}
\leanref{Development/AbelianVarietyTheory/Proofs/UnitaryExistence/HostGeneratorPowFourTraceNormModEight.lean}{GlobalTheory.exists_symplectic_basis_adapted_mod_eight}
\begin{proof}
The $\mathbf Z$-module $L$ is free of rank $n[E:\mathbf Q]=2g$. Put $m=(d+1)/8$, an integer. As
$w^2=w-(d+1)/4$, we have $w^2=w-2m$ and $w(1-w)=2m$, and from these the identities
\[
 w^4+(1-w)^4=1+8(m^2-m),\qquad w^4(1-w)^4=16m^4,\qquad w^8=w^4+8\big((m^2-m)w^4-2m^4\big).
\]
Since $\overline{\sqrt{-d}}=-\sqrt{-d}$, $\bar w=1-w$, so $\langle w^4x,y\rangle=\langle x,(1-w)^4y\rangle$
and $\langle(1-w)^4x,y\rangle=\langle x,w^4y\rangle$. Let $P$ be the set of $x\in L$ with
$w^4x\in x+8L$, and $Q$ the set of $x\in L$ with $w^4x\in8L$. Both are subgroups of $L$ containing $8L$.

The form is divisible by $8$ on $P$. Let $x,y\in P$, with $w^4x=x+8a$ and $w^4y=y+8b$. Then
$(1-w)^4y=(1-w)^4(w^4y-8b)=8\big(2m^4y-(1-w)^4b\big)$, so
$\langle x,y\rangle=\langle w^4x,y\rangle-8\langle a,y\rangle=\langle x,(1-w)^4y\rangle-8\langle a,y\rangle$
is divisible by $8$.

The form is divisible by $8$ on $Q$. Let $x,y\in Q$, with $w^4x=8a$ and $w^4y=8b$. By the first identity
$x=8\big(a-(m^2-m)x\big)+(1-w)^4x$, so $\langle x,y\rangle\equiv\langle(1-w)^4x,y\rangle=
\langle x,w^4y\rangle=8\langle x,b\rangle\pmod8$.

$P+Q=L$. Let $x\in L$ and $c=\big((m^2-m)w^4-2m^4\big)x$. By the third identity $w^4(w^4x)=w^4x+8c$ and
$w^4(x-w^4x)=-8c$, so $w^4x\in P$, $x-w^4x\in Q$, and $x$ is their sum.

Lemma~\ref{cv-uniformization-adapted-basis}, for the form $-\langle\ ,\ \rangle$ and $N=8$, gives a basis
$(e,f)$ with $\langle e_a,f_b\rangle=-\delta_{ab}$ and the other pairings $0$, in which an element lies in
$P$ when its coordinates at the $f_k$ are divisible by $8$, and in $Q$ when its coordinates at the $e_k$
are. So $e_k\in P$ and $f_k\in Q$, which are the last two conditions.
\end{proof}
% Lean: hsq = w0_sq_eq_w0_sub_two_mul (with d mod 16 = 7, m = 2 (d/16) + 1), hA, hB, hC, host_bar_w0;
% hPiso, hQiso, hsup; then exists_isSymplecticBasis_adapted_dvd for -EF (unimodular by
% bijective_neg_bilinForm) with N = 8, and datum_free_int, datum_finrank_int for the rank.

\begin{situation}\label{cv-uniformization-situation}
In Situation~\ref{cv-cusp-setting}, let $L$ be the lattice of the PEL datum of $X$, of signature $(p,q)$
with $p+q=r$, let $G$, $L_W$ and $\xi_0$ be as in Lemma~\ref{cv-uniformization-lattice}, and let
$(M',\mathfrak M')$ be a level of full level $1$ with $M'\ge3$ (Definition~\ref{cv-deep-cover}). Let
$X'$ be the PEL moduli scheme of the datum at this
level, $F'\subset\mathbf C$ a number field containing $F_c$ and the $M'$-th roots of unity, and
$\sigma\colon F'\to\mathbf C$ a good embedding. Put $g=fr$ and index the coordinates of $\mathbf C^g$ by
the pairs $(\tau,a)$ of an embedding $\tau\colon E\to\mathbf C$ and an integer $0\le a<m_\tau$, with
$m_\tau=p$ for $\tau\in\Phi$ and $m_\tau=q$ otherwise; giving the coordinate $(\tau,a)$ the embedding
$\sigma\circ\tau$ is a type of signature $(p,q)$ for $\theta$. Let $L_i$, $t_i$ and $y_i(P)$ be as in
Theorem~\ref{cv-uniformization-family} for this type, and $(G,\xi_0,\varphi_i)$ Hermitian models of the
$L_i$ (Lemma~\ref{cv-uniformization-models}); put $\Lambda_i=\varphi_i(L_i)$ and
$\Gamma_i=\Gamma(\varphi_i)$.

\emph{Frames and periods.} A \emph{frame} is a frame of $Z$ in Situation~\ref{cv-cusp-setting} with $G$
in place of $G_Z$; write $x_{\tau,b}$ for the $b$-th column of $x_\tau$ and
$h_{\tau(G)}(u,v)=\sum_{a,b}u_a\tau(G_{ab})\bar v_b$ on $\mathbf C^r$. The \emph{period} of a frame $x$ on
$L_i$ is the additive map $P^i_x\colon L_i\to\mathbf C^g$ with coordinates, for $\tau\in\Phi$,
\[
 P^i_x(v)_{(\tau,a)}=\big(x_\tau^{-1}\tau(\varphi_iv)\big)_{q+a+1},\qquad
 P^i_x(v)_{(\bar\tau,a)}=h_{\tau(G)}\big(x_{\tau,q-a},\tau(\varphi_iv)\big).
\]

\emph{Normalized framed points.} Fix for each $i$ a symplectic basis of $L_i$ as in
Lemma~\ref{cv-uniformization-mod-eight}, for the prime $d$, which is $7$ modulo $16$, and
$w=(1+\sqrt{-d})/2$, and, holomorphically in
the frame $x$, matrices $h_x$ and $Z_x$ and a characteristic $c_i$ as in
Theorem~\ref{cv-uniformization-family}(4) for the periods $P^i_x$. A \emph{normalized framed point} at
$x$ is a complex framed point $(y,v,w)$ of $X'$ over $\sigma$ with $y=y_i(P^i_x)$, with $v$ the images
in $\operatorname{Lie}A_y$ of the unit vectors of $\mathbf C^g$, and with $w$ the section of the line
bundle of the polarization whose pullback is $\kappa\,\vartheta[c_i](Z_x,h_xz)$ times the factor of
Theorem~\ref{cv-uniformization-family}(4), for a number $\kappa$ with
$\kappa^2=N_\xi\det h_x$, where $N_\xi=\big(2^g/\prod_{k}2|\sigma_k(\xi_0)|\big)^{1/2}$.
Two complex framed points are \emph{triangularly equivalent} if they have the same point, each vector
of the second basis of an eigenspace of $\operatorname{Lie}$ is a combination of the vectors of the first
basis of that eigenspace with index at most its own and a nonzero coefficient at its own index, and
their sections differ by a nonzero factor.
\end{situation}
\leanref{Formalism/LocalToGlobal/ProjectiveLattices/UnFree.lean}{Formalism.L2.PjFree.perFr}
\leanref{Development/AbelianVarietyTheory/Definitions/ExistenceOfLattices/NormConst.lean}{GlobalTheory.FrameUnif.ThetaData}
\leanref{Development/AbelianVarietyTheory/Definitions/ExistenceOfLattices/NormConst.lean}{GlobalTheory.FrameUnif.IsNormUn}

\begin{lemma}\label{cv-uniformization-adapted-blocks}
In Situation~\ref{cv-cusp-setting}, let $L$ be the lattice of a PEL datum over $E$ with a Hermitian model
$(G',\eta,\varphi)$ (Definition~\ref{cv-uniformization-model}), and let
$(e_1,\dots,e_g,f_1,\dots,f_g)$ be a basis of $L$ as in Lemma~\ref{cv-uniformization-mod-eight}, for
the prime $d$, which is $7$ modulo $16$, and $w=(1+\sqrt{-d})/2$. Let $M\ge3$ be an integer and
$\gamma'$ an $\mathcal O_E$-linear automorphism of $L$ with
$\langle\gamma'x,\gamma'y\rangle=\langle x,y\rangle$ and $\gamma'x\in x+ML$ for all $x,y\in L$. Write
the matrix of $\gamma'$ in this basis in blocks of size $g$ as
$N=\begin{pmatrix}N_{11}&N_{12}\\N_{21}&N_{22}\end{pmatrix}$, the first rows and columns belonging to
$e_1,\dots,e_g$. Then the entries of $N_{12}$ and $N_{21}$ are divisible by $8$, and
$\det N_{11}\equiv1\pmod8$.
\end{lemma}
\leanref{Development/AbelianVarietyTheory/Proofs/UnitaryExistence/CmUnitPowFourIdempotentModEight.lean}{GlobalTheory.AdaptedDet.blocks_det_mod_eight}
\leanref{Development/AbelianVarietyTheory/Proofs/UnitaryExistence/CmUnitPowFourIdempotentModEight.lean}{GlobalTheory.AdaptedDet.det_interp_datum}
\leanref{Development/AbelianVarietyTheory/Proofs/UnitaryExistence/IsIntegralDetSubOne.lean}{GlobalTheory.CongruenceDet.det_eq_one}
\begin{proof}
For $x\in L$ let $\pi(x)\in(\mathbf Z/8\mathbf Z)^{2g}$ be its vector of coordinates modulo $8$. As
$w^4e_k\in e_k+8L$ and $w^4f_k\in8L$, $\pi(w^4x)$ is $\pi(x)$ with its coordinates at $f_1,\dots,f_g$
replaced by $0$.

The blocks $N_{12}$ and $N_{21}$. As $\gamma'$ is $\mathcal O_E$-linear,
$w^4\gamma'f_j=\gamma'(w^4f_j)\in8L$, so the coordinates of $\gamma'f_j$ at $e_1,\dots,e_g$ are divisible
by $8$; they form the $j$-th column of $N_{12}$. Likewise $w^4\gamma'e_j\in\gamma'e_j+8L$, so the
coordinates of $\gamma'e_j$ at $f_1,\dots,f_g$ are divisible by $8$; they form the $j$-th column of
$N_{21}$.

Reduction of the determinant. Put $u=w^4$ and $T=u\gamma'+(1-u)$, an $\mathcal O_E$-linear endomorphism of
$L$. By the first paragraph, $\pi(Te_j)$ has the entries of the $j$-th column of $N_{11}$ at the $e_k$
and $0$ at the $f_k$, and $\pi(Tf_j)$ is $\pi(f_j)$ plus a vector that vanishes at the $f_k$. So the
matrix of $T$ modulo $8$ is block upper triangular with diagonal blocks $N_{11}$ and $1$, and
$\det T\equiv\det N_{11}\pmod8$, where $\det T$ is the determinant of $T$ on the $\mathbf Z$-module $L$.

The determinant of $\gamma'$ over $E$. In the coordinates $\varphi$, $\gamma'$ is a matrix
$\tilde\gamma\in GL_r(E)$ with $\tilde\gamma\varphi(x)=\varphi(\gamma'x)$. Since $\gamma'$ preserves the
form and commutes with $\mathcal O_E$,
$\operatorname{Tr}_{E/\mathbf Q}(b\eta\,h_{G'}(\tilde\gamma v,\tilde\gamma v'))=
\operatorname{Tr}_{E/\mathbf Q}(b\eta\,h_{G'}(v,v'))$ for all $b\in\mathcal O_E$ and $v,v'\in\varphi(L)$,
so $h_{G'}(\tilde\gamma v,\tilde\gamma v')=h_{G'}(v,v')$ on $E^r$, that is
$\tilde\gamma^{\mathsf T}G'\bar{\tilde\gamma}=G'$. As $\det G'\ne0$, $\delta=\det\tilde\gamma$ satisfies
$\delta\bar\delta=1$. The map $(\tilde\gamma-1)/M$ preserves the lattice $\varphi(L)$, so for each
eigenvalue $\mu$ of $\tilde\gamma$ the number $(\mu-1)/M$ is an algebraic integer; writing
$\delta=\prod_\mu\big(1+M\cdot(\mu-1)/M\big)$ shows that $(\delta-1)/M$ is an algebraic integer. Every
conjugate of $(\delta-1)/M$ has absolute value at most $2/M\le2/3$, because every conjugate of $\delta$
has absolute value $1$. So the norm of $(\delta-1)/M$ is a rational integer of absolute value less than
$1$, hence $0$, and $\delta=1$.

The determinant of $T$. The eigenvalues $\mu_1,\dots,\mu_r$ of $\tilde\gamma$, with multiplicities, are
algebraic integers, as $\tilde\gamma$ preserves $\varphi(L)$, and $\det\tilde T=\prod_k(u\mu_k+1-u)$ for
the matrix $\tilde T=u\tilde\gamma+(1-u)$ of $T$. By the third identity in the proof of
Lemma~\ref{cv-uniformization-mod-eight}, $u^2-u=8t$ with $t\in\mathcal O_E$. For algebraic integers $\mu$,
$\Pi$ and $y$,
\[
 (u\mu+1-u)(u\Pi+1-u+8y)=u\mu\Pi+1-u+8\big(t(\mu\Pi-\mu-\Pi+1)+y(u\mu+1-u)\big),
\]
so by induction on the number of factors $\det\tilde T=u\delta+1-u+8y=1+8y$ for an algebraic integer $y$,
which lies in $E$ and so in $\mathcal O_E$. The determinant of $T$ on $L$ is the determinant of $\tilde T$
as a $\mathbf Q$-linear map of $E^r$, which is $N_{E/\mathbf Q}(\det\tilde T)=N_{E/\mathbf Q}(1+8y)$, the
determinant of multiplication by $1+8y$ on $\mathcal O_E$; it is congruent to $1$ modulo $8$. Hence
$\det N_{11}\equiv1\pmod8$.
\end{proof}
% Lean: blocks_det_mod_eight: he (w0^4 acts mod 8 as the projection), the off-diagonal blocks, and the
% determinant through T = w0^4 gam + (1 - w0^4) and det_interp_datum (coordinates from
% exists_herm_coord_of_datum with eta = rt; CongruenceDet.det_eq_one with isIntegral_det_sub_one and
% det_mul_conj; cmUnit_pow_four_idempotent_mod_eight; det_interp with prod_interp and det_smul_add;
% det_eq_norm; norm_one_add). The Lean takes its own Hermitian coordinates for L; any model serves.

\begin{lemma}\label{cv-uniformization-period-change}
Let $L$ be a free $\mathbf Z$-module with a basis $(e_1,\dots,e_g,f_1,\dots,f_g)$ and an alternating form
with $\langle e_a,e_b\rangle=\langle f_a,f_b\rangle=0$ and $\langle e_a,f_b\rangle=-\delta_{ab}$. Let
$\gamma'$ be an automorphism of $L$ that preserves the form, with matrix
$N=\begin{pmatrix}N_{11}&N_{12}\\N_{21}&N_{22}\end{pmatrix}$ in this basis, written in blocks as in
Lemma~\ref{cv-uniformization-adapted-blocks}, and put
\[
 A=N_{22}^{\mathsf T},\quad B=N_{12}^{\mathsf T},\quad C=N_{21}^{\mathsf T},\quad D=N_{11}^{\mathsf T},
 \qquad\gamma=\begin{pmatrix}A&B\\C&D\end{pmatrix}.
\]
Let $P\colon L\to\mathbf C^g$ be an additive map, $h,h'\in M_g(\mathbf C)$ with $\det h\ne0$, and $Z,Z'$
symmetric complex $g\times g$ matrices, such that for every $k$ the vectors $hP(e_k)$ and
$h'P(\gamma'e_k)$ are the $k$-th unit vector and $hP(f_k)$ and $h'P(\gamma'f_k)$ are the $k$-th columns
of $Z$ and of $Z'$. Then $\gamma\in Sp_{2g}(\mathbf Z)$, the matrix $CZ+D$ is invertible, and
\[
 Z'=(AZ+B)(CZ+D)^{-1},\qquad h'=\big((CZ+D)^{-1}\big)^{\mathsf T}h,\qquad
 \det h=\det(CZ+D)\det h'.
\]
\end{lemma}
\leanref{Development/AbelianVarietyTheory/Proofs/UnitaryExistence/ChangeOfNormalisation.lean}{GlobalTheory.PeriodChange.period_change}
\begin{proof}
In coordinates the form is $\langle x,y\rangle=-x^{\mathsf T}Jy$ with
$J=\begin{pmatrix}0&1\\-1&0\end{pmatrix}$, and $\gamma'$ preserves it, so $N^{\mathsf T}JN=J$. With
$S=\begin{pmatrix}0&1\\1&0\end{pmatrix}$ we have $\gamma=SN^{\mathsf T}S$ and $SJS=-J$, so
$\gamma J\gamma^{\mathsf T}=SN^{\mathsf T}(SJS)NS=-S(N^{\mathsf T}JN)S=-SJS=J$, and
$\gamma\in Sp_{2g}(\mathbf Z)$.

Since $\gamma'e_l=\sum_kN_{kl}e_k+\sum_kN_{g+k,l}f_k$, the matrix with columns $hP(\gamma'e_l)$ is
$A_0=N_{11}+ZN_{21}$, and the matrix with columns $hP(\gamma'f_l)$ is $B_0=N_{12}+ZN_{22}$. Let $X$ be the
matrix with columns $P(\gamma'e_l)$. Then $hX=A_0$ and $h'X=1$, so $X$ is invertible, $A_0=hX$ is
invertible and $h'=A_0^{-1}h$; and $Z'=h'h^{-1}B_0=A_0^{-1}B_0$. As
$A_0^{\mathsf T}=D+CZ$, using $Z^{\mathsf T}=Z$, the matrix $CZ+D$ is invertible and
$h'=((CZ+D)^{-1})^{\mathsf T}h$. As $Z'$ is symmetric, $Z'=B_0^{\mathsf T}(A_0^{\mathsf T})^{-1}=
(AZ+B)(CZ+D)^{-1}$, since $B_0^{\mathsf T}=B+AZ$. Finally $h=A_0h'$ gives
$\det h=\det(CZ+D)\det h'$.
\end{proof}
% Lean: period_change: hNv, swap_mem_symplectic (gamma = S N^T S), hcol, change_of_normalisation.

\begin{theorem}\label{cv-uniformization-theta-transformation}
Let $g\ge1$ and let $\mathfrak H_g$ be the set of symmetric complex $g\times g$ matrices with positive
definite imaginary part. For $c=(a,b)$ with $a,b\in\mathbf R^g$, $Z\in\mathfrak H_g$ and
$z\in\mathbf C^g$ put
\[
 \vartheta[c](Z,z)=\sum_{n\in\mathbf Z^g}\exp\big(\pi i(n+a)^{\mathsf T}Z(n+a)
 +2\pi i(n+a)^{\mathsf T}(z+b)\big),
\]
and $\vartheta=\vartheta[0]$. Let $\gamma=\begin{pmatrix}A&B\\C&D\end{pmatrix}\in Sp_{2g}(\mathbf Z)$,
with blocks of size $g$, such that the entries of $B$ and $C$ are divisible by $8$. Then for every
$Z\in\mathfrak H_g$ the matrix $CZ+D$ is invertible, $\gamma Z=(AZ+B)(CZ+D)^{-1}$ lies in $\mathfrak H_g$,
and for every $z\in\mathbf C^g$
\[
 \vartheta\big(\gamma Z,((CZ+D)^{-1})^{\mathsf T}z\big)^2=\chi_{-4}(\det D)\,\det(CZ+D)\,
 \exp\big(2\pi i\,z^{\mathsf T}(CZ+D)^{-1}Cz\big)\,\vartheta(Z,z)^2,
\]
where $\chi_{-4}(k)$ is $(-1)^{(k-1)/2}$ for odd $k$ and $0$ for even $k$.
\end{theorem}
\leanref{Development/AbelianVarietyTheory/Proofs/UnitaryExistence/ThetaSqTransformLevelEight.lean}{GlobalTheory.theta_sq_transform_level_eight}
\leanref{Development/AbelianVarietyTheory/Proofs/HomogeneousIdeals/FromBlocksMapIntCastMemSymplecticGroup.lean}{GlobalTheory.theta_halfVec_sq_transform}
\leanref{ReciprocityTheory/ReductiveGroups/ArithmeticGeometry/SymplecticGenerators/SymplecticCharacters/SubConjTransposeOfSymm.lean}{ReciprocityTheory.ChartRestriction.G.Key.Chi.jAut_isUnit}
\leanref{ReciprocityTheory/ReductiveGroups/ArithmeticGeometry/SymplecticGenerators/SymplecticCharacters/SubConjTransposeOfSymm.lean}{ReciprocityTheory.ChartRestriction.G.Key.Chi.spAct_mem}

These matrices lie in the theta group, and the formula is the square of the theta transformation formula
of Igusa \cite{Igu72} and Mumford \cite[Chapter~II, (5.1)]{Mum83} (see also
\cite[Proposition~8.6.1]{BL04}), with the square of its eighth root of unity written as
$\chi_{-4}(\det D)$, as Alp\"oge evaluates it.

\begin{lemma}\label{cv-uniformization-theta-ratio}
Let $\gamma=\begin{pmatrix}A&B\\C&D\end{pmatrix}\in Sp_{2g}(\mathbf Z)$ with the entries of $B$ and $C$
divisible by $8$ and with $\det D\equiv1\pmod4$. Let $Z\in\mathfrak H_g$, let $c=(a,b)$ with
$a,b\in\{0,\frac12\}^g$, let $\rho$ be a nowhere vanishing holomorphic function on $\mathbf C^g$ with
$\rho(0)=1$, and let $\kappa\in\mathbf C^\times$, such that, in the notation of
Theorem~\ref{cv-uniformization-theta-transformation},
\[
 \vartheta[c]\big(\gamma Z,((CZ+D)^{-1})^{\mathsf T}z\big)\,\rho(z)=\kappa\,\vartheta[c](Z,z)
 \qquad(z\in\mathbf C^g).
\]
Then $\kappa^2=\det(CZ+D)$.
\end{lemma}
\leanref{Development/AbelianVarietyTheory/Proofs/UnitaryExistence/RatioSqEqDetJAut.lean}{GlobalTheory.ThetaRatio.ratio_sq_eq_det_jAut}
\leanref{Development/AbelianVarietyTheory/Proofs/UnitaryExistence/TwoDvdMultiplierOfThetaZActEqUnitMul.lean}{GlobalTheory.halfChar_fixed_of_theta_transform_unit}
\leanref{Development/AbelianVarietyTheory/Proofs/UnitaryExistence/ThetaSqLawOfZero.lean}{GlobalTheory.ThetaCharMult.theta_sq_law_of_zero}
\begin{proof}
Put $n=2a$ and $n'=2b$, in $\{0,1\}^g$, $Z'=\gamma Z$, $J=(CZ+D)^{-1}$, $z'=J^{\mathsf T}z$ and
$q(z)=z^{\mathsf T}JCz$. As $\det D\equiv1\pmod4$, $\chi_{-4}(\det D)=1$. The function
$\vartheta[c](Z,\cdot)$ is entire and not identically $0$, so by the identity theorem a holomorphic
function on $\mathbf C^g$ whose product with it, or with its square, vanishes is itself $0$.

Step 1: the characteristic is fixed, that is, $A^{\mathsf T}n\equiv n$ and $D^{\mathsf T}n'\equiv n'$
modulo $2$. For $m,m'\in\mathbf Z^g$ the theta function satisfies
\[
 \vartheta[c](Z,z+Zm+m')=\exp\big(2\pi i(a^{\mathsf T}m'-b^{\mathsf T}m)\big)
 \exp\big(-\pi i\,m^{\mathsf T}Zm-2\pi i\,m^{\mathsf T}z\big)\,\vartheta[c](Z,z),
\]
and $z\mapsto J^{\mathsf T}z$ maps $z+Zm+m'$ to $z'+Z'p+p'$ with $p=Dm-Cm'$ and $p'=Am'-Bm$. Evaluate the
relation at $z+Zm+m'$, expand both theta functions by this rule, and divide by
$\vartheta[c](Z,z)$. With the identity for $q(z+Zm+m')-q(z)$ that follows from $\gamma$ being symplectic,
this shows that $\sigma(z)=\rho(z)\exp(-\pi i\,q(z))$ satisfies
$\sigma(z+Zm+m')=(-1)^{N(m,m')}\sigma(z)$, with the integer
\[
 N(m,m')=n^{\mathsf T}p'-n'^{\mathsf T}p-n^{\mathsf T}m'+n'^{\mathsf T}m-m^{\mathsf T}m'+p^{\mathsf T}p'.
\]
So $\sigma^2$ is an entire function, periodic for the lattice $Z\mathbf Z^g+\mathbf Z^g$, which spans
$\mathbf C^g$ over $\mathbf R$. It is bounded, hence constant, and $(\sigma-\sigma(0))(\sigma+\sigma(0))=0$
with $\sigma(0)=\rho(0)\ne0$ forces $\sigma=\sigma(0)$. So $N(m,m')$ is even for all $m,m'$. For $m=0$ and
$m'=\varepsilon_i$, the $i$-th unit vector,
$(A^{\mathsf T}n)_i-n_i=N(0,\varepsilon_i)-n'^{\mathsf T}C\varepsilon_i+(A\varepsilon_i)^{\mathsf T}
C\varepsilon_i$, and for $m=\varepsilon_i$ and $m'=0$,
$(D^{\mathsf T}n')_i-n'_i=-N(\varepsilon_i,0)-n^{\mathsf T}B\varepsilon_i-(D\varepsilon_i)^{\mathsf T}
B\varepsilon_i$. As $B$ and $C$ are even, both are even.

Step 2: the formula of Theorem~\ref{cv-uniformization-theta-transformation} holds for $\vartheta[c]$:
\[
 \vartheta[c](Z',z')^2=\det(CZ+D)\exp\big(2\pi i\,q(z)\big)\,\vartheta[c](Z,z)^2.
\]
This is a computation from the theorem at the point $z+Z\frac m2+\frac{m'}2$, with
$m=A^{\mathsf T}n+C^{\mathsf T}n'$ and $m'=B^{\mathsf T}n+D^{\mathsf T}n'$, which $z\mapsto J^{\mathsf T}z$
maps to $z'+Z'a+b$. It uses that $\vartheta(Z,z+Z\alpha+\beta)^2$ is $\vartheta[\alpha,\beta](Z,z)^2$ times
an explicit exponential, that $\vartheta[\frac k2,\frac{k'}2](Z,z)^2$ depends only on $k$ and $k'$ modulo
$2$, that $m\equiv n$ and $m'\equiv n'$ modulo $2$ by Step 1 and the evenness of $B$ and $C$, and that
$n^{\mathsf T}n'-m^{\mathsf T}m'$ is divisible by $8$, from $AD^{\mathsf T}=1+BC^{\mathsf T}$ and the
divisibility of $B$ and $C$ by $8$; then the exponential factors cancel.

Step 3. Square the relation and multiply Step 2 by $\rho(z)^2$:
$\big(\det(CZ+D)\exp(2\pi i\,q(z))\rho(z)^2-\kappa^2\big)\vartheta[c](Z,z)^2=0$ for all $z$. So the
first factor vanishes identically, and at $z=0$ it gives $\det(CZ+D)=\kappa^2$.
\end{proof}
% Lean: ratio_sq_eq_det_jAut, with exists_eq_halfVec. Step 1 is halfChar_fixed_of_theta_transform_unit
% (transl from thetaLevelChar_one, mulVec_lattice_translate_eq, quadratic_lattice_translate_eq,
% two_dvd_of_entire_sign_periodic with entire_apply_eq_of_periodic_of_span_real and
% span_real_period_lattice_full, then the two unit vectors); Step 2 is theta_sq_law_of_zero
% (shift_algebra, theta_zero_shift_sq, theta_sq_red, phase_dvd with sp_rel_int); Step 3 uses
% entire_eq_zero_of_mul_continuous_eq_zero and theta_ne_zero.

\begin{lemma}\label{cv-uniformization-theta-change}
In Situation~\ref{cv-uniformization-situation}, let $x$ be a frame, $\gamma\in\Gamma_i$ and $x'=\gamma x$.
Then there are a nowhere vanishing holomorphic function $\rho$ on $\mathbf C^g$ with $\rho(0)=1$ and a
number $\kappa_0\ne0$ with the following properties.
\begin{enumerate}
\item The factor of Theorem~\ref{cv-uniformization-family}(4) at $x'$ is $\rho(z)$ times the factor at
  $x$, both read in the same local trivialization of the line bundle of the polarization over the
  universal abelian scheme.
\item $\vartheta[c_i](Z_{x'},h_{x'}z)\,\rho(z)=\kappa_0\,\vartheta[c_i](Z_x,h_xz)$ for all
  $z\in\mathbf C^g$.
\item $\kappa_0^2\det h_{x'}=\det h_x$.
\end{enumerate}
\end{lemma}
\leanref{Development/AbelianVarietyTheory/Proofs/UnitaryExistence/ExistsSecGeneratorFibre.lean}{GlobalTheory.thetaFrame_ratio}
\leanref{Development/AbelianVarietyTheory/Proofs/UnitaryExistence/ThetaFrameDetMulSqInvariant.lean}{GlobalTheory.thetaFrame_det_mul_sq_invariant}
\begin{proof}
Step 1: the periods. Let $\gamma'=\varphi_i^{-1}\gamma^{-1}\varphi_i$. It is an $\mathcal O_E$-linear
automorphism of $L_i$ that preserves the form, and it is congruent to $1$ modulo $M'$ on
$\bar{\mathfrak M}'L_i$ by the definition of $\Gamma_i$, hence on $L_i$, as in the proof of
Lemma~\ref{cv-uniformization-models}(3). As $\tau(\gamma)$ preserves $h_{\tau(G)}$, the formula for
periods gives $P^i_{x'}=P^i_x\circ\gamma'$.

Step 2: the factors and the sections, (1) and (2). By Theorem~\ref{cv-uniformization-family}(2), with
$h=1$, the periods $P^i_x$ and $P^i_{x'}$ give the same point $y$. The two isomorphisms
$\mathbf C^g/P^i_x(L_i)\to A_y$ and $\mathbf C^g/P^i_{x'}(L_i)\to A_y$ have the same source, since
$P^i_{x'}(L_i)=P^i_x(L_i)$, and they differ by an automorphism of $A_y$ that preserves the polarization
and the level structure; it has finite order and is congruent to $1$ modulo $M'\ge3$ on the lattice, so
it is the identity \cite{Min87}. So the maps from $\mathbf C^g$ to the universal abelian scheme defined by
$x$ and by $x'$ agree. The factor of Theorem~\ref{cv-uniformization-family}(4) is the value of the section
in a local trivialization of the line bundle over the universal abelian scheme; the values in two
trivializations differ by the transition function, and at $z=0$ the value is that of the rigidification
of the line bundle along the zero section. Hence the quotient of the factors at $x'$ and at $x$ does not
depend on the trivialization; it is a nowhere vanishing holomorphic function $\rho$ of $z$, and
$\rho(0)=1$. This is (1). The two sections of Theorem~\ref{cv-uniformization-family}(4) at $x$ and at
$x'$ are sections of the line bundle of the polarization on $A_y$, whose sections form a space of
dimension $1$. So the second is $\kappa_0$ times the first for some $\kappa_0\ne0$, and in a common
trivialization this is (2).

Step 3: the change of period matrix. Let $N$ be the matrix of $\gamma'$ in the basis of $L_i$ fixed in
Situation~\ref{cv-uniformization-situation}. By Lemma~\ref{cv-uniformization-adapted-blocks}, for the
model $(G,\xi_0,\varphi_i)$ and $M=M'$, the entries of $N_{12}$ and $N_{21}$ are divisible by $8$ and
$\det N_{11}\equiv1\pmod8$. By Step 1, Lemma~\ref{cv-uniformization-period-change} applies to
$P=P^i_x$, $h=h_x$, $Z=Z_x$, $h'=h_{x'}$ and $Z'=Z_{x'}$: the matrix
$\gamma_0=\begin{pmatrix}A&B\\C&D\end{pmatrix}$ built from $N$ lies in $Sp_{2g}(\mathbf Z)$,
\[
 Z_{x'}=\gamma_0Z_x,\qquad h_{x'}=\big((CZ_x+D)^{-1}\big)^{\mathsf T}h_x,\qquad
 \det h_x=\det(CZ_x+D)\det h_{x'}.
\]
Its blocks $B=N_{12}^{\mathsf T}$ and $C=N_{21}^{\mathsf T}$ are divisible by $8$, and
$\det D=\det N_{11}\equiv1\pmod4$.

Step 4: (3). Put $\tilde\rho(v)=\rho(h_x^{-1}v)$ for $v\in\mathbf C^g$. For $z=h_x^{-1}v$, Step 3 gives
$h_{x'}z=((CZ_x+D)^{-1})^{\mathsf T}v$, so (2) reads
\[
 \vartheta[c_i]\big(\gamma_0Z_x,((CZ_x+D)^{-1})^{\mathsf T}v\big)\,\tilde\rho(v)
 =\kappa_0\,\vartheta[c_i](Z_x,v).
\]
The characteristic $c_i$ lies in $\{0,\frac12\}^{2g}$, $\tilde\rho$ is holomorphic and nowhere vanishing
with $\tilde\rho(0)=1$, and $Z_x\in\mathfrak H_g$, so Lemma~\ref{cv-uniformization-theta-ratio} gives
$\kappa_0^2=\det(CZ_x+D)$. With Step 3, $\kappa_0^2\det h_{x'}=\det h_x$.
\end{proof}
% Lean: thetaPt_change (GlobalTheory.T.thetaPoints_change_levelIsometry_linearMap, here with h0 = 1)
% through thetaFrame_det_mul_sq_invariant(_uncond). Step 1 is per_gam with
% LevelHalf.trivial_mod_of_trivial_on_half; Step 2 is uniformization_period_lattice_change
% (FamRigid.ex_eq_of_level_group: familyUnif_exp_unique_alt, finite_order and latMinkowski) and
% thetaFrame_ratio (the factors fr j of IsThetaFrame, with the cocycle fr_coc and the rigidification
% fr_rig, and exists_sec_generator_fibre for the one-dimensional space of sections); Step 3 is
% blocks_det_mod_eight for gam.symm and PeriodChange.period_change; Step 4 is ratio_sq_eq_det_jAut with
% the published law theta_sq_transform_level_eight. The factor of Theorem cv-uniformization-family(4) is
% the family fr j of the field theta of FamilyUnif (IsThetaFrame), with the cocycle (fr_coc) and the
% rigidification (fr_rig) stated there.

\begin{lemma}\label{cv-uniformization-frames}
In Situation~\ref{cv-uniformization-situation}, the following hold for every $i$.
\begin{enumerate}
\item For every frame $x$, $P^i_x$ is a period of $L_i$ of the type of the situation, holomorphic in $x$,
  and every period of $L_i$ of this type is $P^i_x$ for some frame $x$.
\item There is a normalized framed point $u_i(x)$ at every frame $x$.
\item If $|\delta_{\tau,a}(x)|=1$ for all $\tau$ and $a$, the mean of $|w|^2e^{-\pi H(z,z)}$ over a
  fundamental domain of $P^i_x(L_i)$ is $1$, for the section $w$ of $u_i(x)$ read on $\mathbf C^g$ and the
  Hermitian form $H$ of the polarization.
\item If $b\in B$ and $xb$ is a frame, then $u_i(xb)$ is triangularly equivalent to $u_i(x)$, by a change
  that multiplies a section of weight $\lambda$ by $\prod_{\tau,a}b_{\tau,aa}^{\lambda_{\tau,a}}$.
\item If $\gamma\in\Gamma_i$, then $u_i(\gamma x)$ and $u_i(x)$ have the same point and the same tangent
  vectors, and their sections differ by a sign; so the change multiplies sections of every weight with
  even $n$ by $1$.
\item For every rational section $f$ of an automorphic line bundle on $X'$, the function
  $x\mapsto f(u_i(x))$ is holomorphic on the set of frames.
\item Every complex framed point of $X'$ over $\sigma$ is triangularly equivalent to $u_i(x)$ for some $i$
  and some frame $x$; and $u_i(x)$ is triangularly equivalent to $u_{i'}(x')$ if and only if $i=i'$ and
  $x'\in\Gamma_ixB$.
\end{enumerate}
\end{lemma}
\leanref{Development/AbelianVarietyTheory/Proofs/UnitaryExistence/ExistsFramedPointsAllNormalisationsFamilyUnif.lean}{GlobalTheory.FrameUnif.exists_framesOK}
\leanref{Development/AbelianVarietyTheory/Proofs/UnitaryExistence/FamilyUniformisationIsComplexUniformisation.lean}{GlobalTheory.FrameUnif.isUnf_built}
\begin{proof}
(1) The vector $\tau(\varphi_iv)$ has coordinates $x_\tau^{-1}\tau(\varphi_iv)$ in the basis of columns of
$x_\tau$, and $P^i_x$ keeps the last $p$ of them at $\tau$ and the pairings with the first $q$ columns at
$\bar\tau$. The relation $\langle\ ,\ \rangle_i=\operatorname{Tr}_{E/\mathbf Q}(\xi_0h_G)$ and the signs of
the minors of a frame give a positive definite Hermitian form $H$ with $\operatorname{Im}H(P^i_xv,P^i_xv')=
\langle v,v'\rangle_i$, using Lemma~\ref{cv-uniformization-signature} with $\operatorname{Im}\tau(\xi_0)>0$. Conversely
a period of this type determines at each $\tau\in\Phi$ the subspace of $\mathbf C^r$ on which the first
$q$ coordinates vanish, and a frame adapted to it.

(2) The symplectic basis adapted modulo $8$ exists by Lemma~\ref{cv-uniformization-mod-eight}, applied to
$L_i$, and the matrices $h_x$, $Z_x$ and the characteristic exist by
Theorem~\ref{cv-uniformization-family}(4) applied to the holomorphic family $P^i_x$; the number
$N_\xi\det h_x$ is nonzero, so $\kappa$ exists, and the section with this pullback exists by the same
theorem.

(3) This is a computation with the Gaussian integral of $|\vartheta[c_i](Z_x,h_xz)|^2e^{-\pi H(z,z)}$
over a fundamental domain: with the factor $\kappa^2=N_\xi\det h_x$ and minors of absolute value $1$
the mean is $1$.

(4) For $b\in B$ the period $P^i_{xb}$ is $h\circ P^i_x$ for a matrix $h$ that is block triangular for
the eigenspaces, with diagonal entries those of $b^{-1}$ at the coordinates $(\tau,a)$ and those of
$\bar b$ at $(\bar\tau,a)$; so the two points agree by Theorem~\ref{cv-uniformization-family}(2), the
tangent vectors change by triangular matrices, and $\kappa$ changes by $(\det h)^{-1/2}$ up to sign. The
action on sections of weight $\lambda$ is the stated character.

(5) For $\gamma\in\Gamma_i$, $\gamma'=\varphi_i^{-1}\gamma^{-1}\varphi_i$ is an automorphism of $L_i$ as in
Theorem~\ref{cv-uniformization-family}(2), and $P^i_{\gamma x}=P^i_x\circ\gamma'$. So the points agree,
with $h=1$, and the tangent vectors agree. Let $\rho$ and $\kappa_0$ be as in
Lemma~\ref{cv-uniformization-theta-change}, and $\kappa$ and $\kappa'$ the numbers of $u_i(x)$ and
$u_i(\gamma x)$. Read in a common trivialization, with $F$ the factor of
Theorem~\ref{cv-uniformization-family}(4) at $x$, the section of $u_i(x)$ is
$\kappa\vartheta[c_i](Z_x,h_xz)F(z)$, and by parts (1) and (2) of that lemma the section of $u_i(\gamma x)$
is $\kappa'\vartheta[c_i](Z_{\gamma x},h_{\gamma x}z)\rho(z)F(z)=\kappa'\kappa_0\vartheta[c_i](Z_x,h_xz)F(z)$.
So the two sections differ by the factor $t=\kappa'\kappa_0/\kappa$, and by part (3) of that lemma
$t^2\kappa^2\det h_{\gamma x}=\kappa'^2\det h_x$. As $\kappa^2=N_\xi\det h_x$ and
$\kappa'^2=N_\xi\det h_{\gamma x}$, $t^2=1$.
% Lean: GlobalTheory.T.thetaPoints_change_levelIsometry_linearMap (ThetaPtChg.thetaPt_change with
% h0 = 1, val_of_secBC for the factor t0), with the adapted basis ThetaData.adapt from
% GlobalTheory.exists_symplectic_basis_adapted_mod_eight.

(6) By Theorem~\ref{cv-uniformization-family}(3) the points $u_i(x)$ and their framings depend
holomorphically on $x$.

(7) Let $(y,v,w)$ be a complex framed point over $\sigma$. By
Theorem~\ref{cv-uniformization-family}(1), $y=y_i(P)$. The tangent vectors $v$ are images of vectors of
$\mathbf C^g$ in the eigenspaces, so a matrix $V\in GL_g(\mathbf C)$, block diagonal for the eigenspaces,
takes them to unit vectors; $VP$ is a period presenting $y$ (Theorem~\ref{cv-uniformization-family}(2)),
and by (1) it is $P^i_x$ for a frame $x$. Then $(y,v,w)$ and $u_i(x)$ have the same point and tangent
vectors, so they are triangularly equivalent. For the second statement, if $u_i(x)$ and $u_{i'}(x')$
are triangularly equivalent, then $i=i'$ by Theorem~\ref{cv-uniformization-family}(1), and
Theorem~\ref{cv-uniformization-family}(2) gives $P^i_{x'}=h\circ P^i_x\circ\gamma'$ with $h$ triangular
for the tangent vectors. Then $\gamma=\varphi_i\gamma'^{-1}\varphi_i^{-1}\in\Gamma_i$, and
$P^i_{x'}=h\circ P^i_{\gamma x}$ with $h$ triangular forces $x'=\gamma xb$ for some $b\in B$. The converse
is (4) and (5).
\end{proof}
% Lean: (1) GlobalTheory.T.exists_framePeriods_alt (onto: hon) with cmType_letter_frame_column_dictionary;
% (2) nonempty_thetaData (exists_symplectic_basis_adapted_mod_eight, FrameTheta.exists_thetaFrame_frames)
% and exists_isNormUn (exists_isThetaPt'); (3) NormFrame.frameOK_of_isNormUn; (4) NormUnBor.normUn_bor;
% (5) normUn_gamma (per_gam, thetaPoints_change_levelIsometry_linearMap, t0^2 = 1, wchar_ex_sign);
% (6) NormUnHol.differentiableOn_rdU; (7) cover_unfOf with CoverFrame.coverTri, and fibre_unfOf with
% FibreFlag.fibTri. isUnf_built assembles these as IsUnf (unfOf e F M un).

\begin{definition}\label{cv-uniformization-marking}
In Situation~\ref{cv-uniformization-situation}, let $\iota\colon L\to L_W$ be an isomorphism as in
Lemma~\ref{cv-uniformization-lattice}, so that $(G,\xi_0,\iota)$ is a Hermitian model of $L$. Let
$\hat L_0$ be the closure of $L_0$ in $(E\otimes_F\mathbf A_{F,f})^r$, and write $h=h_\infty h_f$ for
$h\in U(W_j)(\mathbf A_F)$, acting on coordinate vectors.
\begin{enumerate}
\item The \emph{lattice} of $h$ is $L_h=\mathfrak A\,(E^r\cap h_f\hat L_0)$, and its completion is
  $\hat L_h=\mathfrak A\,h_f\hat L_0$; so $L_1=L_W$.
\item For $\tau\in\Phi$ let $\theta_\tau\colon(F\otimes\mathbf R)\otimes_FE\to\mathbf C$ be the
  homomorphism that extends $\tau$, and fix a frame $x^0$ that is orthonormal for $G$. The
  \emph{frame} of $h$ is $\omega(h)=(\theta_\tau(h_\infty)x^0_\tau)_{\tau\in\Phi}$.
\item A \emph{marking} of $h$ is a pair $(i,\gamma)$ of a component $i$ and $\gamma\in U(W_j)(F)$ with
  $\Lambda_i=L_{\gamma h}$ and
  $\varphi_i(t_ia)-(\gamma h)_f\,\iota(a)\in M'\hat L_{\gamma h}$ for every $a\in\bar{\mathfrak M}'L$.
\item $K'\subset U(W_j)(\mathbf A_{F,f})$ is the principal congruence subgroup of level $M'$ of the
  stabilizer of $L_0$: the group of $k$ with $(k-1)\hat L_0\subset M'\hat L_0$.
\end{enumerate}
\end{definition}
\leanref{Development/AbelianVarietyTheory/Definitions/ExistenceOfLattices/CoeSelfDualLatticeInFrameCoordinatesEq.lean}{GlobalTheory.Member.Ctx.Mark}
\leanref{Development/AbelianVarietyTheory/Definitions/ExistenceOfLattices/CoeSelfDualLatticeInFrameCoordinatesEq.lean}{GlobalTheory.Member.Ctx.fr}
\leanref{Formalism/LocalToGlobal/ProjectiveLattices/FamilyUniformisationBasePointOrthogonalBis.lean}{Formalism.L2.Pj.latG}
\leanref{Development/AbelianVarietyTheory/Proofs/UnitaryExistence/ExistsPrincipalLevelSubgroupOpen.lean}{GlobalTheory.Member.PrincLevel.exists_levelK_and_nhds}

The frame $x^0$ exists because $\tau(G)$ has $q$ positive and $p$ negative eigenvalues for $\tau\in\Phi$,
by Lemma~\ref{cv-uniformization-signature} with $\operatorname{Im}\tau(\xi_0)>0$. The set $K'$ is a group,
since $(kk'-1)=k(k'-1)+(k-1)$ and $k^{-1}-1=-k^{-1}(k-1)$ for $k,k'$ preserving $\hat L_0$.

\begin{lemma}\label{cv-uniformization-markings}
In Definition~\ref{cv-uniformization-marking}:
\begin{enumerate}
\item every $h\in U(W_j)(\mathbf A_F)$ has a marking;
\item for every component $i$ there is $h\in U(W_j)(\mathbf A_{F,f})$ with the marking $(i,1)$;
\item if $(i,\gamma)$ and $(i',\gamma')$ are markings of $h$, then $i=i'$, and
  $\omega(\gamma'h)=\delta\,\omega(\gamma h)$ for some $\delta\in\Gamma_i$;
\item if $(i,\gamma)$ is a marking of $h$, then $(i,\gamma\gamma_0^{-1})$ is a marking of $\gamma_0h$ for
  $\gamma_0\in U(W_j)(F)$, $(i,\gamma)$ is a marking of $hu$ for $u\in U(W_j)(F\otimes\mathbf R)$ and for
  $u\in K'$, and for $\delta\in\Gamma_i$ there is $\gamma_1\in U(W_j)(F)$ such that $(i,\gamma_1\gamma)$ is
  a marking of $h$ and $\omega(\gamma_1\gamma h)=\delta\,\omega(\gamma h)$;
\item if $(i,\gamma)$ is a marking of $h$, $(i,\gamma')$ one of $h'$, and $\omega(\gamma h)=\omega(\gamma'h')$,
  then $(\gamma h)^{-1}\gamma'h'\in K'$.
\end{enumerate}
\end{lemma}
\leanref{Development/AbelianVarietyTheory/Proofs/UnitaryExistence/ExistsComponentLatticeMarkingAdelicStatement.lean}{GlobalTheory.T.exists_component_lattice_marking_adelic}
\leanref{Development/AbelianVarietyTheory/Proofs/UnitaryExistence/ExistsFiniteAdelicElementMarkingComponentStatement.lean}{GlobalTheory.T.exists_finiteAdelic_element_marking_component}
\leanref{Development/AbelianVarietyTheory/Proofs/UnitaryExistence/MarkingComponentUniqueFrameUniqueArithmeticGroupStatement.lean}{GlobalTheory.T.marking_component_unique_frame_unique_arithmeticGroup}
\leanref{Development/AbelianVarietyTheory/Proofs/UnitaryExistence/MarkingRuleTranslateStableStatement.lean}{GlobalTheory.T.markingRule_translate_stable}
\leanref{Development/AbelianVarietyTheory/Proofs/UnitaryExistence/MemLevelKOfSameMarkAndFrameStatement.lean}{GlobalTheory.T.mem_levelK_of_same_mark_and_frame}
\begin{proof}
A \emph{level structure} on a lattice $\Lambda\subset E^r$ is a family of $\mathcal O_E$-linear maps
$f_N\colon L\to\Lambda$, $N\ge1$, with $f_{NN'}\equiv f_N$ modulo $N\Lambda$, inducing isomorphisms
$L/NL\to\Lambda/N\Lambda$, and with
$\operatorname{Tr}_{E/\mathbf Q}(\xi_0h_G(f_Na,f_Na'))\equiv\langle a,a'\rangle$ modulo $N$.

(1) The lattice $L_h$ is self-dual for the pairing $\operatorname{Tr}_{E/\mathbf Q}(\xi_0h_G)$, because
$L_W$ is (the form of $L$ is perfect) and $h$ preserves $h_G$. Approximating $h_f$ by rational matrices
gives a level structure $(f_N)$ on $L_h$ from $\iota$. A self-dual lattice for this pairing is the
lattice of a PEL datum of signature $(p,q)$, by Lemma~\ref{cv-uniformization-signature}; with its level
structure it defines a complex point of $X'$ over $\sigma$, its rational Hermitian space being that of
$L$ by Landherr's theorem and the Hasse norm theorem as in the proof of
Lemma~\ref{cv-uniformization-models}. By Theorem~\ref{cv-uniformization-family}(1) this point is some
$y_i(P)$, which gives an isomorphism of $L_h$ onto $L_i$ preserving the forms and carrying $f_{M'}$ to
$t_i$ modulo $M'$ on $\bar{\mathfrak M}'L$. As both lattices have Hermitian models with $(G,\xi_0)$, this
isomorphism is a matrix $\gamma'\in GL_r(E)$ preserving $h_G$, with $\gamma'L_h=\Lambda_i$. Such a matrix is
the image of some $\gamma\in U(W_j)(F)$, and $L_{\gamma h}=\gamma L_h=\Lambda_i$; the congruence for $f_{M'}$
gives the congruence of the marking.

(2) By Theorem~\ref{cv-uniformization-family}, applied to a period of $L_i$, the maps $t_i$ extend to a
level structure on $\Lambda_i$ from $\iota$. Its limit over $N$ is an isometry of the completions, which
is $g\hat L_W=\hat\Lambda_i$ for some $g\in U(W_j)(\mathbf A_F)$. The finite part $h$ of $g$ has
$L_h=\Lambda_i$ and satisfies the congruence of the marking with $\gamma=1$.

(3) The matrix $\gamma'\gamma^{-1}$ carries $\Lambda_i=L_{\gamma h}$ onto $\Lambda_{i'}=L_{\gamma'h}$, preserves
$h_G$ and is compatible with the two congruences; through $\varphi_i$ and $\varphi_{i'}$ it is an
isomorphism of $L_i$ onto $L_{i'}$ carrying $t_i$ to $t_{i'}$ modulo $M'$ on $\bar{\mathfrak M}'L$.
Transporting a period of $L_i$ along it presents one complex point both as $y_i(P)$ and as $y_{i'}(P')$,
so $i=i'$ by Theorem~\ref{cv-uniformization-family}(1). Then $\delta=\gamma'\gamma^{-1}\in\Gamma_i$, and
$\theta_\tau(\delta)=\tau(\delta)$ gives $\omega(\gamma'h)=\delta\,\omega(\gamma h)$.

(4) The lattice and the congruence of a marking depend on $\gamma h$ only through its finite part, and
are unchanged under right multiplication by $u\in K'$, since $u\hat L_0=\hat L_0$ and
$u\equiv1$ modulo $M'\hat L_0$. For $\delta\in\Gamma_i$, $\delta$ is the image of some
$\gamma_1\in U(W_j)(F)$; as $\delta\Lambda_i=\Lambda_i$ and $\delta\equiv1$ modulo $M'$ on
$\varphi_i(\bar{\mathfrak M}'L_i)$, $(i,\gamma_1\gamma)$ is a marking, and
$\omega(\gamma_1\gamma h)=\delta\,\omega(\gamma h)$ as in (3).

(5) Put $u=(\gamma h)^{-1}\gamma'h'$. The equality of frames gives
$\theta_\tau(u_\infty)=1$ for every $\tau\in\Phi$, so $u\in U(W_j)(\mathbf A_{F,f})$. The two markings give
$u\hat L_0=\hat L_0$ and $u\equiv1$ modulo $M'$ on $\bar{\mathfrak M}'\hat L_0$; as $u$ preserves the
form, $u\equiv1$ modulo $M'\hat L_0$, as in the proof of Lemma~\ref{cv-uniformization-models}(3).
\end{proof}
% Lean: (1) Member.exists_mark: adelicUnitaryLattice_selfDual_traceForm, AdLat.exists_linear_approx_tower
% with finitePart_preserves_form_det_norm_one, selfDual_lattice_occurs
% (exists_latticeDatum_of_selfDual_lattice, allClassesSt_of_hasseNorm,
% exists_unitary_matrix_of_lattice_isometry), rational_points_adelicUnitary_eq_unitary and
% adelicUnitaryLattice_rationalTranslate_invariance; (2) Member.exists_finite_mark:
% familyUnif_lattice_genus, onto_adelic_of_tower, adelicUnitary_eq_real_mul_finite_split,
% onto_marking_sub; (3) MarkUniq.mark_unique with ClassUniq.eq_of_matched; (4) and (5) are the two
% statements cited.

\begin{lemma}\label{cv-uniformization-adelic-frames}
In Definition~\ref{cv-uniformization-marking}, let $h\in U(W_j)(\mathbf A_F)$. Then $\omega(h)$ is a frame
that is orthonormal for $G$, so $|\delta_{\tau,a}(\omega(h))|=1$ for all $\tau$ and $a$; the map
$h\mapsto\omega(h)$ is continuous; $\omega(hu)=\omega(h)$ for $u\in U(W_j)(\mathbf A_{F,f})$;
$\omega(hu)=\omega(h)\,\big((x^0_\tau)^{-1}\theta_\tau(u_\infty)x^0_\tau\big)_\tau$ for every $u$; and every frame
that is orthonormal for $G$ is $\omega(u)$ for some $u\in U(W_j)(F\otimes\mathbf R)$.
\end{lemma}
\leanref{Development/AbelianVarietyTheory/Proofs/UnitaryExistence/MemberDictionaryFramesOrthonormalUnitStatement.lean}{GlobalTheory.T.memberDictionary_frames_orthonormal_unit}
\begin{proof}
Each $\theta_\tau$ is a continuous ring homomorphism, trivial on the finite adeles, and the matrix
$\theta_\tau(h_\infty)$ preserves the form of $\tau(G)$. So $\omega(h)$ is orthonormal, and the leading
minors of a diagonal matrix with entries $\pm1$ are $\pm1$. The two product rules follow from the
multiplicativity of $\theta_\tau$. Two frames orthonormal for $G$ differ at each $\tau\in\Phi$ by a matrix
that preserves the form of $\tau(G)$, and every family of such matrices, $\tau\in\Phi$, is
$(\theta_\tau(u))_\tau$ for some $u\in U(W_j)(F\otimes\mathbf R)$, because the $\theta_\tau$ identify
$(F\otimes\mathbf R)\otimes_FE$ with $\prod_{\tau\in\Phi}\mathbf C$.
\end{proof}
% Lean: Member.fr_orthonormal_unit, with archimedeanMatrix_unitary_continuous_multiplicative (from
% infiniteAdele_bijective), UnitFrames.unit_frame_eq_orthonormal_mul_triangular and
% UnitFrames.frame_mul_triangular_orthonormal.

\begin{lemma}\label{cv-uniformization-dictionary}
In Definition~\ref{cv-uniformization-marking}, there is a map $g\mapsto(Z(g),x(g))$ from
$U(W_j)(\mathbf A_F)$ to pairs of a component and a frame such that $x(g)=\omega(\gamma g)$ for a marking
$(Z(g),\gamma)$ of $g$, and such that for every component $i$ some $g\in U(W_j)(\mathbf A_{F,f})$ has
$(Z(g),x(g))=(i,x^0)$. Every such map has the following properties.
\begin{enumerate}
\item $|\delta_{\tau,a}(x(g))|=1$ for all $g$, $\tau$ and $a$.
\item For $\gamma\in U(W_j)(F)$ and $k\in K'$, $Z(\gamma g)=Z(gk)=Z(g)$, and $x(\gamma g)$ and $x(gk)$ lie in
  $\Gamma_{Z(g)}x(g)$.
\item For $h\in U(W_j)(F\otimes\mathbf R)$, $Z(gh)=Z(g)$ and
  $x(gh)\in\Gamma_{Z(g)}\,x(g)\big((x^0_\tau)^{-1}\theta_\tau(h)x^0_\tau\big)_\tau$.
\item If $Z(g')=Z(g)$ and $x(g')\in\Gamma_{Z(g)}x(g)$, then $g'\in U(W_j)(F)\,g\,K'$.
\item For every component $i$, every frame lies in $\Gamma_i\,x(g)B$ for some $g$ with $Z(g)=i$.
\item If $(i,\gamma)$ is a marking of $g$, then $Z(g)=i$ and $x(g)\in\Gamma_i\,\omega(\gamma g)$.
\end{enumerate}
\end{lemma}
\leanref{Development/AbelianVarietyTheory/Proofs/UnitaryExistence/ExistsDictionaryByMarkingRuleStatement.lean}{GlobalTheory.T.exists_dictionary_by_marking_rule}
\leanref{Development/AbelianVarietyTheory/Proofs/UnitaryExistence/MemberDictionaryInvariantRationalLevelStatement.lean}{GlobalTheory.T.memberDictionary_invariant_rational_level}
\leanref{Development/AbelianVarietyTheory/Proofs/UnitaryExistence/MemberDictionaryRealEquivariantOntoTriangularStatement.lean}{GlobalTheory.T.memberDictionary_real_equivariant_onto_triangular}
\leanref{Development/AbelianVarietyTheory/Proofs/UnitaryExistence/DictIsUnifOfMarkingRuleDictionary.lean}{GlobalTheory.Member.dictIsUnif_of_isDict}
\begin{proof}
For the existence, choose by Lemma~\ref{cv-uniformization-markings}(1) a marking $(i,\gamma)$ of every $g$,
with $\gamma=1$ whenever some $(i,1)$ is a marking of $g$, and put $(Z(g),x(g))=(i,\omega(\gamma g))$. For a
component $i$, let $h$ be as in Lemma~\ref{cv-uniformization-markings}(2); by part (3) of that lemma $Z(h)=i$,
and $x(h)=\omega(h)=\omega(1)=x^0$ by Lemma~\ref{cv-uniformization-adelic-frames}.

(6) is Lemma~\ref{cv-uniformization-markings}(3), and (1) is
Lemma~\ref{cv-uniformization-adelic-frames}. For (2), if $(i,\gamma_0)$ is the marking of $g$ used for
$x(g)$, then $(i,\gamma_0\gamma^{-1})$ marks $\gamma g$ and $(i,\gamma_0)$ marks $gk$
(Lemma~\ref{cv-uniformization-markings}(4)); by (6), $x(\gamma g)\in\Gamma_i\omega(\gamma_0g)$, and
$x(gk)\in\Gamma_i\omega(\gamma_0gk)=\Gamma_i\omega(\gamma_0g)$ by
Lemma~\ref{cv-uniformization-adelic-frames}. Part (3) follows in the same way from the marking
$(i,\gamma_0)$ of $gh$ and the product rule of Lemma~\ref{cv-uniformization-adelic-frames}. For (4), let
$x(g')=\delta x(g)$ with $\delta\in\Gamma_i$, and let $(i,\gamma)$ and $(i,\gamma')$ be the markings used. By
Lemma~\ref{cv-uniformization-markings}(4) there is $\gamma_1$ with $(i,\gamma_1\gamma)$ a marking of $g$ and
$\omega(\gamma_1\gamma g)=\delta\omega(\gamma g)=\omega(\gamma'g')$, and part (5) of that lemma gives
$(\gamma_1\gamma g)^{-1}\gamma'g'\in K'$. For (5), write a frame as $x'b$ with $x'$ orthonormal and
$b\in B$, by the Gram--Schmidt process for $\tau(G)$; by Lemma~\ref{cv-uniformization-adelic-frames},
$x'=\omega(u)$ with $u\in U(W_j)(F\otimes\mathbf R)$. If $h$ is finite with $(Z(h),x(h))=(i,x^0)$, then by (3)
$x(hu)\in\Gamma_i\,x^0\big((x^0_\tau)^{-1}\theta_\tau(u)x^0_\tau\big)_\tau=\Gamma_i\,x'$, so $x'b\in\Gamma_ix(hu)B$.
\end{proof}
% Lean: Member.Ctx.exists_isDict (the choice, with dif on Mark h k 1); Member.dict_invariant gives (1),
% (2) and (6); memberDictionary_real_equivariant_onto_triangular gives (3) and (5); (4) is the last
% field of dictIsUnif_of_isDict, from markingRule_translate_stable and mem_levelK_of_same_mark_and_frame.

\subsection{Existence of deep covers}\label{subsection-cv-deep-covers-exist}

\begin{theorem}\label{cv-deep-pel-facts}
In Situation~\ref{cv-cusp-setting}, let $(M',\mathfrak M')$ be a level
(Definition~\ref{cv-deep-cover}).
\begin{enumerate}
\item The PEL moduli problem of the PEL datum of $X$ at level $(M',\mathfrak M')$ is represented, on
  reduced schemes over $\mathbf Q(\sqrt{-d})$, by a smooth quasi-projective scheme $X'$ over
  $\mathbf Q(\sqrt{-d})$ of relative dimension $fpq$, with a universal abelian scheme.
\item Suppose that the level of $X$ is a level $(M,\mathfrak M)$ with $M\mid M'$ and
  $\mathfrak M'\subset\mathfrak M$. Then there are a morphism $\pi\colon X'\to X$ and a homomorphism,
  over $\pi$, from the universal abelian scheme over $X'$ to the one over $X$ that identifies the
  first with the pullback of the second, compatibly with the actions of $\mathcal O_E$ and the
  polarizations, and under which $M'/M$ times the level structure of $X'$ is that of $X$. For every
  number field $F'\subset\mathbf C$ that contains $F_c$ and the $M'$-th roots of unity, the base change
  $X'\otimes F'\to X\otimes F'$ of $\pi$ is finite, \'etale and surjective.
\item Let $F'\subset\mathbf C$ be a number field that contains $F_c$ and the $M'$-th roots of unity, and
  let $K'\subset U(W_j)(\mathbf A_{F,f})$ be the principal congruence subgroup of level $M'$ of the
  stabilizer of the lattice of $W_j$. Then $X'$, with coefficient field $F'$, has at every good
  embedding $\sigma\colon F'\to\mathbf C$ a uniformization and a map from $U(W_j)(\mathbf A_F)$ to
  frames of components with the properties listed in Definition~\ref{cv-deep-cover}(5), for $K'$ in
  place of $K_Y$, compatible with one identification of the PEL datum of $X$ with the lattice of $W_j$
  that does not depend on $\sigma$.
\end{enumerate}
\end{theorem}
\leanref{Development/AbelianVarietyTheory/Proofs/ModuliProofs/ExistsModuliSchemePolarizedAbelianLevel.lean}{GlobalTheory.Moduli.moduli_scheme_exists}
\leanref{ReciprocityTheory/ReductiveGroups/ArithmeticGeometry/BirationalGeometry/ToroidalCompactification/PushInvertible.lean}{ReciprocityTheory.ChartRestriction.G.Geo.Tor.Moduli}
\leanref{Development/AbelianVarietyTheory/Proofs/ModuliProofs/ModuliSchemeLevelMapsHeckePointsStatement.lean}{GlobalTheory.T.moduliScheme_levelMaps_heckePoints'}
\leanref{ReciprocityTheory/ReductiveGroups/ArithmeticGeometry/BirationalGeometry/ToroidalCompactification/HeckeKer.lean}{ReciprocityTheory.ChartRestriction.G.Geo.Tor.ToroidalLevelCompatibility}
\leanref{Development/AbelianVarietyTheory/Proofs/UnitaryExistence/ExistsMemberDictionaryContextGoodEmbedding.lean}{GlobalTheory.Member.exists_ctx}
\leanref{Development/AbelianVarietyTheory/Proofs/UnitaryExistence/ExistsUniformisationDictionaryCharacterLaws.lean}{GlobalTheory.Member.exists_pkgLaws}
\leanref{Development/AbelianVarietyTheory/Proofs/UnitaryExistence/ExistsPrincipalLevelSubgroupOpen.lean}{GlobalTheory.Member.PrincLevel.exists_levelK_and_nhds}

\begin{proof}
Parts (1) and (2) are the representability of PEL moduli problems at a level of at least $3$ and the
maps between moduli schemes at a level and at a deeper one \cite{Kot92b,Lan13}.

For (3), the datum of $X$ comes with an isomorphism $\iota$ of its lattice $L$ onto $L_W$ as in
Lemma~\ref{cv-uniformization-lattice}, so $(G,\xi_0,\iota)$ is a Hermitian model of $L$; it does not
depend on $\sigma$. Fix a frame $x^0$ orthonormal for $G$ (Definition~\ref{cv-uniformization-marking}),
and let $\sigma$ be a good embedding of $F'$. Theorem~\ref{cv-uniformization-family} gives the lattices
$L_i$, the maps $t_i$ and the points $y_i(P)$ at $\sigma$, and Lemma~\ref{cv-uniformization-models},
applied to $(G,\xi_0,\iota)$, gives Hermitian models $(G,\xi_0,\varphi_i)$ of the $L_i$ and the groups
$\Gamma_i$; this is Situation~\ref{cv-uniformization-situation}.

The uniformization at $\sigma$ has the components $i$, each with the Hermitian matrix $G$ of $W_j$, the
lattice $\Lambda_i$, the group $\Gamma_i$, trivial characters and the element $-\xi_0$, with
$|\tau(-\xi_0)|=\xi$ for every $\tau$; it sends a frame $x$ of the component $i$ to the normalized
framed point $u_i(x)$ of Lemma~\ref{cv-uniformization-frames}. By that lemma, these data describe the
complex framed points over $\sigma$ up to triangular equivalence (part (7)), the readings of sections are
holomorphic and transform under $B$ by the weight characters and under $\Gamma_i$ trivially (parts (4)
to (6)), and the sections are normalized at frames with unit minors (part (3)). By
Lemma~\ref{cv-uniformization-models}, $\Gamma_i$ acts freely on frames modulo $B$ and contains the
transvections. These are the compatibility conditions of the moduli problem.

The map $g\mapsto(Z(g),x(g))$ is that of Lemma~\ref{cv-uniformization-dictionary}, with the principal
congruence subgroup $K'$ of Definition~\ref{cv-uniformization-marking}. Its parts (1) to (5) are the
properties listed in Definition~\ref{cv-deep-cover}(5), and it is built from markings for the one
identification $\iota$, with the orthonormal base frame $x^0$ at every good $\sigma$.
\end{proof}
% Lean: GlobalTheory.CovTop.Inhab.exists_good_pin: GlobalTheory.Member.exists_ctx (exists_family_models:
% closing_of_abelianSchemeLaws and ModelsExist.exists_modelsOf; the base frame x0 from the signature of
% the model) and GlobalTheory.Member.exists_pkgLaws (FrameUnif.exists_framesOK, isUnf_built,
% exists_dictionary_by_marking_rule, memberDictionary_invariant_rational_level for the unit minors,
% dictIsUnif_of_markingRule_dictionary), with GlobalTheory.Member.PrincLevel.exists_levelK_and_nhds and
% GlobalTheory.Member.exists_iotaOK.

\begin{lemma}\label{cv-deep-ideal}
Let $d\equiv7\pmod{16}$ be a prime such that $-d$ is a nonzero square modulo $3$ and modulo $5$. Let $E$
be a CM field containing a square root $\sqrt{-d}$ of $-d$, with complex conjugation $a\mapsto\bar a$, and
let $w=(1+\sqrt{-d})/2$. Then there is an integer $s$ with $15\mid s^2+d$. For such an $s$ and every
integer $e\ge0$, the ideal $I=J^e$, where $J=(15,s+\sqrt{-d})\subset\mathcal O_E$, is generated by
elements of $\mathbf Z[w]$, and $I\bar I=15^e\mathcal O_E$ and $I+\bar I=\mathcal O_E$.
\end{lemma}
\leanref{Development/AbelianVarietyTheory/Definitions/Divisors/NormNegCmConstantEqAtAllEmbeddings.lean}{GlobalTheory.CovTop.Inhab.exists_root}
\leanref{Development/AbelianVarietyTheory/Definitions/Divisors/NormNegCmConstantEqAtAllEmbeddings.lean}{GlobalTheory.CovTop.Inhab.exists_half}
\begin{proof}
The only nonzero square modulo $3$ is $1$, so $-d\equiv1\pmod3$, and $-d\equiv1$ or $4\pmod5$. Take
$s=1$ in the first case and $s=2$ in the second. Then $s^2\equiv1\equiv-d\pmod3$ and
$s^2\equiv-d\pmod5$, so $15\mid s^2+d$.

Since $d\equiv7\pmod{16}$, $d$ is neither $3$ nor $5$, so $15$ is prime to $d$, and then to $s$, because
$15\mid s^2+d$. Hence $15$ and $2s$ are coprime, and $15a+2sb=1$ for some integers $a,b$. Complex
conjugation sends $\sqrt{-d}$ to $-\sqrt{-d}$, so $\bar J=(15,s-\sqrt{-d})$. As
$2s=(s+\sqrt{-d})+(s-\sqrt{-d})$ lies in $J+\bar J$, so does $1=15a+2sb$. The product $J\bar J$ is
generated by $15^2$, $15(s-\sqrt{-d})$, $15(s+\sqrt{-d})$ and $(s+\sqrt{-d})(s-\sqrt{-d})=s^2+d$, which
all lie in $15\mathcal O_E$. Conversely $15\in J\cap\bar J$, and $J\cap\bar J=J\bar J$ because $J$ and
$\bar J$ are coprime. So $J\bar J=15\mathcal O_E$. Then $I\bar I=(J\bar J)^e=15^e\mathcal O_E$, and
$I+\bar I=\mathcal O_E$ because powers of coprime ideals are coprime. Finally $15$ and
$s+\sqrt{-d}=s-1+2w$ lie in $\mathbf Z[w]$, and $J^e$ is generated by the products of $e$ of them.
\end{proof}

\begin{lemma}\label{cv-deep-ray-field}
Let $d\equiv7\pmod{16}$ be a prime, let $\sqrt{-d}\in\mathbf C$ be a square root of $-d$, let
$K_0=\mathbf Q(\sqrt{-d})\subset\mathbf C$, and let $n\ge3$ be an integer. Let $\wp$ be the Weierstrass
function of the lattice $\mathcal O_{K_0}\subset\mathbf C$, with invariants $g_2,g_3$, and put
$\Delta=g_2^3-27g_3^2$. The \emph{Weber values of level $n$} are $g_2^3/\Delta$ and the numbers
$g_2g_3\wp(t)/\Delta$ for $t\in n^{-1}\mathcal O_{K_0}\setminus\mathcal O_{K_0}$. Let $R\subset\mathbf C$
be the ray class field of $K_0$ of modulus $n\mathcal O_{K_0}$. Then $R$ is Galois over $\mathbf Q$, it
contains $\sqrt{-d}$, the $n$-th roots of unity and the Weber values of level $n$, and
\[
 \log[R:\mathbf Q]+\log\operatorname{rd}_R\le20(\log n+\log d+1).
\]
\end{lemma}
\leanref{Development/AbelianVarietyTheory/Proofs/DivisorProofs/ExistsGoodCoverOverSizedField.lean}{GlobalTheory.CovTop.SizedCover.exists_ray}
\leanref{Development/AbelianVarietyTheory/Proofs/IwasawaDecomposition/RayClassFieldDegreeLeAndRootDiscrBound.lean}{GlobalTheory.Mem.Fld.RaySize.degOver_le_and_sz}
\leanref{Development/AbelianVarietyTheory/Proofs/HeightProofs/WeberInvariantTorsionValuesSubsetRayClassFieldStatement.lean}{GlobalTheory.T.weberInvariantTorsionValues_subset_rayClassField}
\leanref{Development/AbelianVarietyTheory/Definitions/HeightFunctions/XiValIffXiAt.lean}{GlobalTheory.FinalJoin.WeberIn}
\begin{proof}
The ray class field exists by class field theory \cite[Chapter~VI]{Neu99}: $R$ is abelian over $K_0$,
unramified at the primes not dividing $n$, and a principal prime $\alpha\mathcal O_{K_0}$ with
$\alpha\equiv1\pmod n$ splits completely in $R$. The modulus is stable under complex conjugation, so
the conjugate field $\bar R$ is a ray class field of $K_0$ of the same modulus, and $\bar R=R$. Hence
$R$ is Galois over $\mathbf Q$. It contains $K_0$, and so $\sqrt{-d}$. A prime $\alpha\mathcal O_{K_0}$
with $\alpha\equiv1\pmod n$ has norm $\alpha\bar\alpha\equiv1\pmod n$, so its Frobenius fixes the $n$-th
roots of unity, and $K_0$ adjoined these roots lies in $R$. Since $d>4$, the units of
$\mathcal O_{K_0}$ are $\pm1$, and the Weber values of level $n$ lie in $R$ by Weber's theorem
\cite[Theorem~II.5.6]{Sil94}.

Degree. Let $m=[R:K_0]$, the order of the ray class group of modulus $n\mathcal O_{K_0}$. This group is
an extension of the class group by a quotient of $(\mathcal O_{K_0}/n)^\times$, which has at most
$n^2$ elements, so $m\le h(K_0)n^2$. As $d\equiv3\pmod4$, $K_0$ has discriminant $-d$, and the class
number of a quadratic field is at most the absolute value of its discriminant, so $h(K_0)\le d$. Hence
$m\le dn^2$ and $\log[R:\mathbf Q]=\log2+\log m\le\log2+\log d+2\log n$.

Root discriminant. Let $\mathfrak P$ be a prime of $R$ over a prime $p$, with ramification index
$e_{\mathfrak P}$ over $K_0$. If $p\nmid n$, $\mathfrak P$ is unramified and does not divide the
different $\mathfrak D_{R/K_0}$. If $p\mid n$, the exponent of $\mathfrak P$ in $\mathfrak D_{R/K_0}$
is at most $e_{\mathfrak P}-1+v_{\mathfrak P}(e_{\mathfrak P})$ \cite[Chapter~III, \S6]{Ser79}; here
$v_{\mathfrak P}(n)\ge e_{\mathfrak P}$, and $v_{\mathfrak P}(m)\ge v_{\mathfrak P}(e_{\mathfrak P})$
because $e_{\mathfrak P}$ divides $m$. So $\mathfrak D_{R/K_0}$ divides $nm\mathcal O_R$. Then
$|d_R|=|d_{K_0}|^mN(\mathfrak D_{R/K_0})\le d^m(nm)^{2m}$, and taking $2m$-th roots,
$\log\operatorname{rd}_R\le\frac12\log d+\log n+\log m$.

Adding the two bounds, $\log[R:\mathbf Q]+\log\operatorname{rd}_R\le\log2+\frac52\log d+5\log n$, which
is at most $20(\log n+\log d+1)$.
\end{proof}

\begin{lemma}\label{cv-deep-compositum}
Let $F,F'\subset\mathbf C$ be number fields, and let $L=FF'$ be their compositum. Then
$[L:\mathbf Q]\le[F:\mathbf Q][F':\mathbf Q]$ and $\operatorname{rd}_L\le\operatorname{rd}_F\operatorname{rd}_{F'}$.
\end{lemma}
\leanref{Development/AbelianVarietyTheory/Proofs/HeightProofs/FieldSizeSupAdd.lean}{GlobalTheory.FieldSize.sz_sup}
\leanref{Development/AbelianVarietyTheory/Proofs/HeightProofs/MapDifferentLeDifferentSupIntermediateField.lean}{GlobalTheory.DiffBaseChange.map_different_le_sup}
\begin{proof}
A basis of $F'$ over $\mathbf Q$ spans $L$ over $F$, which gives the first inequality.

For the second, let $(b_i)$ be a $\mathbf Z$-basis of $\mathcal O_{F'}$ and $(b_i^*)$ its dual basis for
the trace form of $F'$ over $\mathbf Q$. For embeddings $\sigma,\tau$ of $F'$ in $\mathbf C$,
$\sum_i\sigma(b_i)\tau(b_i^*)$ is $1$ if $\sigma=\tau$ and $0$ otherwise. Let $y\in L$, let $u$ be the
inclusion of $L$ in $\mathbf C$, and write $\operatorname{Tr}_{L/F}(x)=\sum_\rho\rho(x)$ over the
embeddings $\rho$ of $L$ that agree with $u$ on $F$. Then
\[
 \sum_iu\big(\operatorname{Tr}_{L/F}(yb_i)\,b_i^*\big)=\sum_\rho\rho(y)\sum_i\rho(b_i)u(b_i^*)=u(y),
\]
because the inner sum vanishes unless $\rho$ agrees with $u$ on $F'$, and an embedding of $L=FF'$ that
agrees with $u$ on $F$ and on $F'$ is $u$. So $y=\sum_i\operatorname{Tr}_{L/F}(yb_i)\,b_i^*$.

Now let $\delta$ lie in the different $\mathfrak D_{F'/\mathbf Q}$, and let $y$ lie in the inverse
different of $L$ over $F$, so that $\operatorname{Tr}_{L/F}(y\mathcal O_L)\subset\mathcal O_F$. The
$b_i^*$ lie in the inverse different of $F'$, so $\delta b_i^*\in\mathcal O_{F'}$, and
$\delta y=\sum_i\operatorname{Tr}_{L/F}(yb_i)\,\delta b_i^*\in\mathcal O_L$. Hence $\delta$ lies in
$\mathfrak D_{L/F}$, that is, $\mathfrak D_{L/F}$ divides $\mathfrak D_{F'/\mathbf Q}\mathcal O_L$. As
$\mathfrak D_{L/\mathbf Q}=\mathfrak D_{L/F}\mathfrak D_{F/\mathbf Q}\mathcal O_L$
\cite[Chapter~III, \S2]{Neu99}, taking norms from $L$ to $\mathbf Q$ gives
$|d_L|\le|d_F|^{[L:F]}|d_{F'}|^{[L:F']}$, and taking $[L:\mathbf Q]$-th roots gives
$\operatorname{rd}_L\le\operatorname{rd}_F\operatorname{rd}_{F'}$.
\end{proof}

\begin{definition}\label{cv-components-conjugation}
Let $d\equiv7\pmod{16}$ be a prime, let $\theta\in\mathbf C$ with $\theta^2=-d$, and let
$\mathcal O=\mathbf Z+\mathbf Z\frac{1+\theta}2\subset\mathbf C$, the ring of integers of $\mathbf Q(\theta)$.
For $w\in\mathcal O$ write $\bar w$ for its complex conjugate and $N(w)=w\bar w$. Let $\tau$ be an
automorphism of $\mathbf C$. The \emph{conjugate} of a complex point $y\colon\operatorname{Spec}\mathbf C\to Y$
of a scheme $Y$ is $y^\tau=y\circ\operatorname{Spec}\tau$; if $Y$ is a scheme over a field $k$ and $y$ lies over
an embedding $j\colon k\to\mathbf C$, then $y^\tau$ lies over $\tau\circ j$, and its coordinates are the images
under $\tau$ of those of $y$. The \emph{conjugate} of a scheme $Y$ over $\mathbf C$ is its base change $Y^\tau$
along $\operatorname{Spec}\tau$, of which $y^\tau$ is a complex point; if $Y$ comes by base change from a
scheme over a subfield of $\mathbf C$ that $\tau$ fixes, then $Y^\tau=Y$.
\begin{enumerate}
\item A \emph{CM curve of level one} for $\theta$ is an elliptic curve $C$ over $\mathbf C$ with a
  holomorphic isomorphism of groups $\mathbf C/\mathcal O\to C(\mathbf C)$, written $z\mapsto[z]$.
\item A \emph{$\tau$-isomorphism} of $C$ is an isomorphism of elliptic curves $\iota\colon C\to C^\tau$ over
  $\mathbf C$. For an integer $n\ge1$ and $w\in\mathcal O$, $\iota$ \emph{acts by $w$ on the points of order
  $n$} if $\iota([wz])=[z]^\tau$ for every $z\in n^{-1}\mathcal O$.
\item The automorphism $\tau$ \emph{fixes a CM curve with its points of order $n$} if some CM curve of level
  one for $\theta$ has a $\tau$-isomorphism that acts by $1$ on its points of order $n$.
\end{enumerate}
\end{definition}
\leanref{Development/AbelianVarietyTheory/Definitions/ExistenceOfLattices/PermutesRefl.lean}{GlobalTheory.CompRat.Curve}
\leanref{Development/AbelianVarietyTheory/Definitions/ExistenceOfLattices/PermutesRefl.lean}{GlobalTheory.CompRat.Curve.ActsBy}
\leanref{Development/AbelianVarietyTheory/Definitions/ExistenceOfLattices/PermutesRefl.lean}{GlobalTheory.CompRat.FixesCurve}
\leanref{Development/AbelianVarietyTheory/Definitions/ExistenceOfLattices/PermutesRefl.lean}{GlobalTheory.CompRat.CurveOver}

\begin{lemma}\label{cv-components-weber-curve}
Let $d$, $\theta$ and $\mathcal O$ be as in Definition~\ref{cv-components-conjugation}, and let $n\ge1$ be an
integer. Let $\wp$ be the Weierstrass function of the lattice $\mathcal O\subset\mathbf C$, with invariants
$g_2,g_3$, and put $\Delta=g_2^3-27g_3^2$. Let $\tau$ be an automorphism of $\mathbf C$ with
$\tau(\theta)=\theta$ that fixes $g_2^3/\Delta$ and the numbers $g_2g_3\wp(t)/\Delta$ for
$t\in n^{-1}\mathcal O\setminus\mathcal O$. (The hypothesis $\tau(\theta)=\theta$ is not used in the proof.)
Then $\tau$ fixes a CM curve with its points of order $n$. If an automorphism of $\mathbf C$ fixes a CM curve with its
points of order $n$, it fixes it with its points of order $m$ for every $m$ dividing $n$.
\end{lemma}
\leanref{Development/AbelianVarietyTheory/Proofs/UnitaryExistence/CmHostRootComplexEqPmSqrt.lean}{GlobalTheory.CompRat.Weber.fixesCurve_of_weber}
\leanref{Development/AbelianVarietyTheory/Proofs/UnitaryExistence/CmHostRootComplexEqPmSqrt.lean}{GlobalTheory.CompRat.Weber.exists_conj}
\leanref{Development/AbelianVarietyTheory/Proofs/UnitaryExistence/CmHostRootComplexEqPmSqrt.lean}{GlobalTheory.CompRat.Weber.exists_model}
\leanref{Development/AbelianVarietyTheory/Proofs/HeightProofs/CmCurveTorsionDefinedOverOfWeberValues.lean}{GlobalTheory.FinalJoin.CurveField.actsBy_of_dvd}
\begin{proof}
First, $g_2\ne0$ and $g_3\ne0$. For $c\in\mathbf C^\times$ the lattice $c\mathcal O$ has the invariants
$c^{-4}g_2$ and $c^{-6}g_3$, and a lattice is determined by its invariants
\cite[Theorem~VI.5.1]{Sil09}. If $g_2=0$, then $\zeta=(1+\sqrt{-3})/2$ satisfies $\zeta^6=1$, so $\zeta\mathcal O$
has the invariants of $\mathcal O$ and $\zeta\mathcal O=\mathcal O$; if $g_3=0$, the same holds for $\zeta=i$. In
both cases $\zeta=\zeta\cdot1\in\mathcal O$. But $m+k\frac{1+\theta}2\in\mathcal O$ with $k\ne0$ has absolute
value at least $|k|\sqrt d/2>1$, so the elements of $\mathcal O$ of absolute value $1$ are real, and $\zeta$ is
not.

Put $j=g_2^3/\Delta$ and $\lambda=g_2g_3/\Delta\ne0$, and let $R\subset\mathbf C$ be the subfield fixed by
$\tau$; it contains $j$. Since $j-1=27g_3^2/\Delta$, the elements $a_4=-j(j-1)/108$ and $a_6=-j(j-1)^2/2916$
of $R$ are $a_4=-\lambda^2g_2/4$ and $a_6=-\lambda^3g_3/4$, and $-16(4a_4^3+27a_6^2)=\lambda^6\Delta\ne0$. So
$y^2=x^3+a_4x+a_6$ is an elliptic curve $W$ over $R$. Choose $\mu\in\mathbf C$ with $\mu^2=\lambda$. As
$\wp'^2=4\wp^3-g_2\wp-g_3$, the map that sends $z\notin\mathcal O$ to the point
$(\lambda\wp(z),\mu^3\wp'(z)/2)$ and $\mathcal O$ to the origin is a holomorphic isomorphism of groups
$\mathbf C/\mathcal O\to W(\mathbf C)$ \cite[Proposition~VI.3.6]{Sil09}. So $C=W\otimes_R\mathbf C$ with this
map is a CM curve of level one, and $x([z])=\lambda\wp(z)$ for $z\notin\mathcal O$.

Since $C$ comes from $W$ over $R$, $C^\tau=C$, and $y\mapsto y^\tau$ is the action of $\tau$ on the
coordinates of the points of $W$. It is an automorphism of the group $C(\mathbf C)$, as the group law of $W$
is given by rational functions with coefficients in $R$. Let $t\in n^{-1}\mathcal O$. If $t\in\mathcal O$, then
$[t]^\tau=[t]$ is the origin. Otherwise $[t]^\tau=[z]$ for some $z\notin\mathcal O$, and
$\lambda\wp(z)=x([t]^\tau)=\tau(\lambda\wp(t))=\lambda\wp(t)$, because $\lambda\wp(t)$ is one of the numbers
that $\tau$ fixes. As $\lambda\ne0$, $\wp(z)=\wp(t)$, so $z\equiv\pm t\pmod{\mathcal O}$
\cite[\S VI.3]{Sil09}. Hence $[t]^\tau=\pm[t]$ for every $t\in n^{-1}\mathcal O$.

If $[t]^\tau=[t]$ for every $t\in n^{-1}\mathcal O$, the identity of $C=C^\tau$ is a $\tau$-isomorphism that
acts by $1$ on the points of order $n$. Otherwise some $a\in n^{-1}\mathcal O$ has $[a]^\tau=-[a]\ne[a]$, and
then $[b]^\tau=-[b]$ for every $b\in n^{-1}\mathcal O$. Indeed, if $[b]^\tau=[b]$, then $[a+b]^\tau=[b]-[a]$ is
$[a+b]$ or $-[a+b]$; the first gives $[a]=-[a]$, against the choice of $a$, and the second gives
$[b]=-[b]$. So the inverse map $\iota=-1$ of $C=C^\tau$ is a $\tau$-isomorphism with
$\iota([t])=-[t]=[t]^\tau$ for every $t\in n^{-1}\mathcal O$, and it acts by $1$ on the points of order $n$.

Finally, if $\iota$ acts by $1$ on the points of order $n$ and $m$ divides $n$, then
$m^{-1}\mathcal O\subset n^{-1}\mathcal O$, so $\iota$ acts by $1$ on the points of order $m$.
\end{proof}

\begin{lemma}\label{cv-components-multiplier}
Let $d$, $\theta$ and $\mathcal O$ be as in Definition~\ref{cv-components-conjugation}, let $C$ be a CM curve of
level one for $\theta$, let $\tau$ be an automorphism of $\mathbf C$ with $\tau(\theta)=\theta$, let $\iota$ be
a $\tau$-isomorphism of $C$, and let $n\ge1$ be an integer.
\begin{enumerate}
\item There is $w\in\mathcal O$ by which $\iota$ acts on the points of order $n$. For every such $w$ there are
  $w',a\in\mathcal O$ with $w'w=1+na$, and $\iota$ acts by $w$ on the points of order $m$ for every $m$
  dividing $n$.
\item If $\iota$ acts by $w$ on the points of order $n$, then $\tau(\zeta_n)=\zeta_n^{N(w)}$, where
  $\zeta_n=e^{2\pi i/n}$.
\end{enumerate}
\end{lemma}
\leanref{Development/AbelianVarietyTheory/Proofs/UnitaryExistence/ConjugationActsOnCmTorsionByMultiplicationStatement.lean}{GlobalTheory.T.conjugation_acts_on_cm_torsion_by_multiplication}
\leanref{Development/AbelianVarietyTheory/Proofs/UnitaryExistence/ConjugationActsOnCmTorsionByMultiplication.lean}{GlobalTheory.CompRat.Acts.corr}
\leanref{Development/AbelianVarietyTheory/Proofs/UnitaryExistence/ExistsCmEndomorphismCommConjugationStatement.lean}{GlobalTheory.T.exists_cmEndomorphism_comm_conjugation}
\leanref{Development/AbelianVarietyTheory/Proofs/UnitaryExistence/CyclotomicCharacterEqCmMultiplierNormStatement.lean}{GlobalTheory.T.cyclotomicCharacter_eq_cmMultiplier_norm}
\begin{proof}
(1) Let $w_+=(1+\theta)/2$. Since $w_+\mathcal O\subset\mathcal O$, the map $[z]\mapsto[w_+z]$ is an
endomorphism $[w_+]$ of $C$; it acts on the tangent space at the origin by $w_+$. Its conjugate $[w_+]^\tau$
acts on the tangent space of $C^\tau$ by $\tau(w_+)=w_+$, so $\iota^{-1}\circ[w_+]^\tau\circ\iota$ is an
endomorphism of $C$ that acts on the tangent space by $w_+$. An endomorphism of a complex elliptic curve is
determined by its action on the tangent space, so $\iota\circ[w_+]=[w_+]^\tau\circ\iota$.

For $z,z'\in\mathbf C$ write $z\to z'$ if $\iota([z'])=[z]^\tau$. Every $z$ has some $z\to z'$ and every $z'$
some $z\to z'$, as $\iota$ and $y\mapsto y^\tau$ are bijections. If $z\to z'$ and $u\to u'$, then
$z+u\to z'+u'$, because both maps are additive, and $z-u\in\mathcal O$ if and only if $z'-u'\in\mathcal O$. By
the commutation above, $z\to z'$ gives $w_+z\to w_+z'$, so $cz\to cz'$ for every $c\in\mathcal O=\mathbf
Z+\mathbf Zw_+$. Also $0\to0$. Choose $z'$ with $1/n\to z'$. Then $1\to nz'$ and $0\to0$, and
$1-0\in\mathcal O$, so $w=nz'$ lies in $\mathcal O$. For $c\in\mathcal O$, $c/n\to wc/n$, that is,
$\iota([wc/n])=[c/n]^\tau$, and $\iota$ acts by $w$ on the points of order $n$.

Now let $\iota$ act by $w$ on the points of order $n$, and choose $u$ with $u\to1/n$. Then $nu\to1$ and
$0\to0$, so $w'=nu\in\mathcal O$. By assumption $w'/n\to ww'/n$, and $w'/n\to1/n$, so
$ww'/n-1/n\in\mathcal O$, that is, $ww'=1+na$ with $a\in\mathcal O$. Finally, for $m$ dividing $n$,
$m^{-1}\mathcal O\subset n^{-1}\mathcal O$, so $\iota$ acts by $w$ on the points of order $m$.

(2) Let $e_n$ be the Weil pairing on the points of order $n$ of a complex elliptic curve. It is a
nondegenerate alternating pairing with values in the $n$-th roots of unity; it is defined algebraically, so
$e_n(y^\tau,y'^\tau)=\tau(e_n(y,y'))$ on the conjugate curve, and isomorphisms preserve it; and
$e_n(\alpha y,\alpha y')=e_n(y,y')^{\deg\alpha}$ for an endomorphism $\alpha$ \cite[\S III.8]{Sil09}. The
endomorphism $[w]$ of $C$ has degree $N(w)$, the order of $w^{-1}\mathcal O/\mathcal O$. The points of order
$n$ of $C$ form a group isomorphic to $(\mathbf Z/n\mathbf Z)^2$, on which $e_n$ is nondegenerate, so
$e_n([b/n],[b'/n])=\zeta_n$ for some $b,b'\in\mathcal O$. Then
\[
 \tau(\zeta_n)=e_n([b/n]^\tau,[b'/n]^\tau)=e_n(\iota[wb/n],\iota[wb'/n])=e_n([w][b/n],[w][b'/n])
 =\zeta_n^{N(w)}.
\]
\end{proof}

\begin{situation}\label{cv-components-situation}
In Situation~\ref{cv-cusp-setting}, let $\theta$ be the square root of $-d$ fixed by $\mathcal D$, let
$\sigma_0$ be the embedding of $\mathbf Q(\sqrt{-d})$ in $\mathbf C$ with $\sigma_0(\sqrt{-d})=\theta$, and let
$\mathcal O$ be as in Definition~\ref{cv-components-conjugation}. Put $w_0=(1+\sqrt{-d})/2\in\mathcal O_E$ and
$w_\pm=(1\pm\theta)/2\in\mathcal O$, and for $w\in\mathcal O$ let $\hat w\in\mathcal O_E$ be its image under the
isomorphism $\mathbf Q(\theta)\to\mathbf Q(\sqrt{-d})\subset E$ that sends $\theta$ to $\sqrt{-d}$. Let $L$ be
the lattice of the PEL datum of $X$, let $(M',\mathfrak M')$ be a level of full level $1$
(Definition~\ref{cv-deep-cover}), and let $X'$ be the PEL moduli scheme of the PEL datum of $X$ at this level (Theorem~\ref{cv-deep-pel-facts}(1)),
with universal abelian scheme $A$ and the sections $\eta(a)$, $a\in\bar{\mathfrak M}'L$, of its level
structure. For a number field $F''\subset\mathbf C$ containing $\theta$, write $X'_{F''}$ for the base change of
$X'$ along $\mathbf Q(\sqrt{-d})\to F''$, $\sqrt{-d}\mapsto\theta$. Fix a type $k\mapsto\sigma_k$, $1\le k\le g$,
of signature $(p,q)$ for $\theta$ (Definition~\ref{cv-cm-point-periods}), through which $b\in\mathcal O_E$ acts
on $\mathbf C^g$ by $(bz)_k=\sigma_k(b)z_k$. Let the lattices $L_1,\dots,L_s$, with forms
$\langle\ ,\ \rangle_i$, the maps $t_i\colon L\to L_i$ and the complex points $y_i(P)$ of $X'$ over $\sigma_0$
be as in Theorem~\ref{cv-uniformization-family} for this type. For a period $P$ of $L_i$ of this type, let
$\exp_{i,P}\colon\mathbf C^g\to A_{y_i(P)}(\mathbf C)$ be the projection to $\mathbf C^g/P(L_i)$ followed by the
isomorphism of that theorem, and let $L_\pm(P)$ be the pure sublattices of $P$. The \emph{component} $Y_i$ is
the set of the points $y_i(P)$.
\end{situation}

\begin{lemma}\label{cv-components-points}
In Situation~\ref{cv-components-situation}, let $C$ be a CM curve of level one for $\theta$
(Definition~\ref{cv-components-conjugation}).
\begin{enumerate}
\item Let $P$ and $P'$ be periods of $L_i$ and $L_{i'}$. Then $y_i(P)=y_{i'}(P')$ if and only if there are an
  $\mathcal O_E$-linear isomorphism $\gamma\colon L_i\to L_{i'}$ with
  $\langle\gamma a,\gamma b\rangle_{i'}=\langle a,b\rangle_i$ for $a,b\in L_i$ and
  $\gamma(t_i(a))-t_{i'}(a)\in M'L_{i'}$ for $a\in\bar{\mathfrak M}'L$, and a $\mathbf C$-linear automorphism $h$
  of $\mathbf C^g$ with $P'\circ\gamma=h\circ P$.
\item Let $P$ be a period of $L_i$, let $\varepsilon\in\{1,-1\}$ and let $a\in L_\varepsilon(P)$. There is a
  homomorphism $\phi\colon C\to A_{y_i(P)}$ with $\phi([z])=\exp_{i,P}(zP(a))$ for $z\in\mathbf C$.
  Conversely, if $\phi\colon C\to A_{y_i(P)}$ is a homomorphism with $\phi([w_\varepsilon z])=w_0\phi([z])$
  for $z\in\mathbf C$, then $\phi([z])=\exp_{i,P}(zP(a))$ for a unique $a\in L_\varepsilon(P)$.
\item Let $F''\subset\mathbf C$ be a number field containing $\theta$, let $j\colon F''\to\mathbf C$ be an
  embedding with $j(\theta)=\theta$, and let $f$ be a regular function on $X'_{F''}$ that takes finitely many
  values at the complex points of $X'_{F''}$ over $j$. Then for each $i$, $f$ takes the same value at all the
  complex points of $X'_{F''}$ over $j$ that lie over points of $Y_i$.
\item Let $F''\subset\mathbf C$ be a number field containing $F_c$ and the $M'$-th roots of unity. There are a
  number field $F'''\supset F''$ and, for each $i$, a regular function $e_i$ on $X'_{F'''}$ with $e_i^2=e_i$
  whose value at a complex point of $X'_{F'''}$ over the inclusion $F'''\subset\mathbf C$ is $1$ if the point
  lies over a point of $Y_i$ and $0$ if not.
\end{enumerate}
\end{lemma}
\leanref{Development/AbelianVarietyTheory/Proofs/ModuliProofs/PointEqIffIsometryClosed.lean}{GlobalTheory.UnifPoint.point_eq_iff_isometry_closed}
\leanref{Development/AbelianVarietyTheory/Proofs/UnitaryExistence/ExistsCmCurveHomOfPureTypeVectorAnyLevelStatement.lean}{GlobalTheory.T.exists_cmCurve_hom_of_pureType_vector_anyLevel}
\leanref{Development/AbelianVarietyTheory/Proofs/UnitaryExistence/CmCurveHomExistsPureTypeVectorAnyLevelStatement.lean}{GlobalTheory.T.cmCurveHom_exists_pureTypeVector_anyLevel}
\leanref{Development/AbelianVarietyTheory/Proofs/UnitaryExistence/RegularFunctionConstOfFiniteValuesAnyLevelStatement.lean}{GlobalTheory.T.regularFunction_const_of_finite_values_any_level}
\leanref{Development/AbelianVarietyTheory/Proofs/UnitaryExistence/ExistsCoeffFieldIdempotentsPeriodSpacesAnyLevelStatement.lean}{GlobalTheory.T.exists_coeffField_idempotents_periodSpaces_anyLevel}
\begin{proof}
(1) Let $\gamma$ and $h$ be given; $h$ induces an isomorphism $\mathbf C^g/P(L_i)\to\mathbf C^g/P'(L_{i'})$. It
commutes with the actions of $\mathcal O_E$: for $b\in\mathcal O_E$ and $x\in L_i$,
$h(bP(x))=h(P(bx))=P'(b\gamma x)=bh(P(x))$, and $P(L_i)$ spans $\mathbf C^g$ over $\mathbf R$. It respects the
Weil pairings, which the forms give, because $\gamma$ preserves the forms, and so it respects the
polarizations. It carries the image of $P(t_i(a))/M'$ to that of $P'(\gamma t_i(a))/M'$, which is the image of
$P'(t_{i'}(a))/M'$ by the congruence. By Theorem~\ref{cv-uniformization-family}, the abelian varieties of
$y_i(P)$ and $y_{i'}(P')$ with their structures are therefore isomorphic, and the two points are equal because
$X'$ represents the moduli problem (Theorem~\ref{cv-deep-pel-facts}(1)). Conversely, suppose that
$y_i(P)=y_{i'}(P')$. Then $\exp_{i,P}$ and $\exp_{i',P'}$ are two uniformizations of the same complex
abelian variety, so there is a $\mathbf C$-linear automorphism $h$ of $\mathbf C^g$ with
$\exp_{i',P'}\circ h=\exp_{i,P}$ \cite[Chapter~1]{BL04}. It carries the kernel $P(L_i)$ of $\exp_{i,P}$ onto
the kernel $P'(L_{i'})$ of $\exp_{i',P'}$, so $\gamma=P'^{-1}\circ h\circ P$ is a bijection $L_i\to L_{i'}$ with
$P'\circ\gamma=h\circ P$. It is $\mathcal O_E$-linear, since both uniformizations are compatible with the
actions of $\mathcal O_E$, and it preserves the forms, since both give the Weil pairings of the same
polarization. For $a\in\bar{\mathfrak M}'L$ the level structure gives
$\exp_{i',P'}(P'(t_{i'}(a))/M')=\eta(a)=\exp_{i,P}(P(t_i(a))/M')=\exp_{i',P'}(P'(\gamma t_i(a))/M')$, so
$P'(\gamma t_i(a)-t_{i'}(a))/M'\in P'(L_{i'})$, that is, $\gamma(t_i(a))-t_{i'}(a)\in M'L_{i'}$.

(2) The coordinates of $P(a)$ vanish outside the $k$ with $\sigma_k(\sqrt{-d})=\varepsilon\theta$, where
$\sigma_k(w_0)=w_\varepsilon$, so $P(w_0a)=w_\varepsilon P(a)$. Hence $(m+nw_\varepsilon)P(a)=P((m+nw_0)a)$ lies in
$P(L_i)$ for $m,n\in\mathbf Z$, and the holomorphic homomorphism $z\mapsto\exp_{i,P}(zP(a))$ vanishes on
$\mathcal O=\mathbf Z+\mathbf Zw_\varepsilon$. It factors through $C(\mathbf C)$, and a holomorphic homomorphism
of complex abelian varieties is algebraic \cite[Chapter~I]{Mum70}. Conversely, $z\mapsto\phi([z])$ is a
holomorphic homomorphism $\mathbf C\to A_{y_i(P)}(\mathbf C)$, so it lifts through $\exp_{i,P}$ to a $\mathbf
C$-linear map $z\mapsto zv$ \cite[Chapter~1]{BL04}, and $v=P(a)$ for some $a\in L_i$ because $\phi([1])=0$.
The relation $\phi([w_\varepsilon z])=w_0\phi([z])=\exp_{i,P}(zP(w_0a))$ gives
$z\big(w_\varepsilon P(a)-P(w_0a)\big)\in P(L_i)$ for every $z\in\mathbf C$. A discrete subgroup contains no
complex line, so $w_\varepsilon P(a)=P(w_0a)$, that is, $(w_\varepsilon-\sigma_k(w_0))P(a)_k=0$ for every $k$.
So $P(a)_k=0$ unless $\sigma_k(\sqrt{-d})=\varepsilon\theta$, and $a\in L_\varepsilon(P)$. The element $a$ is
unique: if $\exp_{i,P}(zP(a))=\exp_{i,P}(zP(a'))$ for every $z$, then $\mathbf C(P(a)-P(a'))\subset P(L_i)$, so
$P(a)=P(a')$ and $a=a'$.

(3) The periods of $L_i$ of the fixed type form, up to $\mathbf C$-linear automorphisms of $\mathbf C^g$, the
symmetric domain of the unitary group of $L_i\otimes\mathbf R$ \cite{Shi63}: there are a connected open set
$B$ in some $\mathbf C^m$ and periods $P_b$ of $L_i$, depending holomorphically on $b\in B$, such that every
period of $L_i$ of this type is $h\circ P_b$ for some $b$ and some $\mathbf C$-linear automorphism $h$. By
Theorem~\ref{cv-uniformization-family}(2), $y_i(h\circ P)=y_i(P)$. By part (3) of that theorem the complex point
of $X'_{F''}$ over $j$ that lies over $y_i(P_b)$ depends holomorphically on $b$, so the value of $f$ there is a
continuous function of $b$ with values in a finite set. As $B$ is connected, it is constant.

(4) Since $X'$ is of finite type over $\mathbf Q(\sqrt{-d})$, there is a number field $F'''\supset F''$ such
that every connected component of $X'_{F'''}$ has an $F'''$-rational point. Such a component $Z$ is
connected with a rational point, hence geometrically connected \cite[\S4.5]{Gro65}, so its complex points
over the inclusion $F'''\subset\mathbf C$ form a connected space in the complex topology
\cite[Expos\'e~XII, \S2]{Gro71}. On this space, the index $i$ with $y\in Y_i$ of the image $y$ of a point in $X'$
is locally constant, by Theorem~\ref{cv-uniformization-family}(1) and~(3), so it is constant. The union $U_i$ of
the components on which it equals $i$ is open and closed in $X'_{F'''}$, and $e_i$ is the regular function
equal to $1$ on $U_i$ and to $0$ on its complement.
\end{proof}

\begin{lemma}\label{cv-components-split}
In Situation~\ref{cv-components-situation}, suppose that there is a PEL datum of signature $(1,0)$ over $E$.
Then every $L_i$ has a period $P$ of the fixed type with $L_i=L_+(P)+L_-(P)$.
\end{lemma}
\leanref{Development/AbelianVarietyTheory/Proofs/UnitaryExistence/ExistsSplitPeriodAnyLevelStatement.lean}{GlobalTheory.T.exists_split_period_any_level}
\leanref{Development/AbelianVarietyTheory/Definitions/ExistenceOfLattices/PermutesRefl.lean}{GlobalTheory.CompRat.SplitLat}
\leanref{Development/AbelianVarietyTheory/Proofs/UnitaryExistence/ExistsSplitPureTypePeriodOfSignatureSplit.lean}{GlobalTheory.CompRat.exists_isSplit_of_splitLat}
\leanref{Development/AbelianVarietyTheory/Proofs/UnitaryExistence/HermitianLatticeSignatureSplitOfDefinite.lean}{GlobalTheory.CompRat.splitLat_of_definite}
\leanref{Development/AbelianVarietyTheory/Proofs/UnitaryExistence/HermitianLatticeSignatureSplitOfIndefinite.lean}{GlobalTheory.SplitGenus.splitLat_of_indefinite}
\begin{proof}
Call a lattice $L'$ with an alternating form as in Definition~\ref{cv-cm-point-periods}, of rank $p+q$,
\emph{split} if $L'=S_+\oplus S_-$ for orthogonal submodules $S_\pm$ of ranks $p$ and $q$ such that
$\operatorname{Im}(\theta)\langle\sqrt{-d}\,a,a\rangle$ is positive for nonzero $a\in S_+$ and negative for nonzero
$a\in S_-$. If $P$ is a period of $L'$ with a positive definite Hermitian form $H$,
$\operatorname{Im}H(P(x),P(y))=\langle x,y\rangle$, then for $a\in L'_\pm(P)$,
$\langle\sqrt{-d}\,a,a\rangle=\operatorname{Im}H(\pm\theta P(a),P(a))=\pm\operatorname{Im}(\theta)H(P(a),P(a))$.
Conversely, a split $L'$ has a period $P$ of the fixed type with $L'_\pm(P)=S_\pm$: $S_+\otimes\mathbf R$ is the
product of the spaces $S_+\otimes_{\mathcal O_E,\sigma}\mathbf C$ over the embeddings $\sigma$ with
$\sigma(\sqrt{-d})=\theta$, each of dimension $p$, and $S_-\otimes\mathbf R$ that of the spaces
$S_-\otimes_{\mathcal O_E,\sigma}\mathbf C$ over the other embeddings, each of dimension $q$; bases of these
spaces give an injective $P$ of the fixed type, and the signs make the Hermitian form $H$ with
$\operatorname{Im}H(P(x),P(y))=\langle x,y\rangle$ positive definite. So it suffices to show that every $L_i$ is
split.

If $q=0$, the computation above for a period of $L_i$ shows that $L_i$ is split with $S_-=0$. Let now $q\ge1$;
here $p\ge1$ and $p+q=r\ge6$, since $m\ge2$ and $k\ge m+3$. Alp\"oge's argument is as follows. The lattice of the
PEL datum of signature $(1,0)$ is split with $S_-=0$. The sum $L_0$ of $p$ copies of it and $q$ copies of it with
$\mathcal O_E$ acting through complex conjugation is split, of rank $p+q$; write $L_0=S^0_+\oplus S^0_-$. The
Hermitian spaces over $E$ of $L_i$ and $L_0$ have the same dimension and the same signatures and carry
unimodular lattices, so they are isometric, by Landherr's theorem \cite{Lan36} and the Hasse norm theorem.
Identify them as one space $V$ containing $L_i$ and $L_0$, and let $\pi$ be the orthogonal projection of $V$
onto $V_+=S^0_+\otimes\mathbf Q$, so that $\pi$ preserves $L_0$. Let $b\ge1$ be an integer with $bL_i\subset L_0$
and $bL_0\subset L_i$; at the primes $\ell\nmid b$ the completions $L_{i,\ell}$ and $L_{0,\ell}$ agree. As $E$ is
unramified over $F$ at the finite primes, unimodular lattices in $V$ are locally isometric \cite{Jac62}, and by
weak approximation there is a unitary automorphism $\gamma_0$ of $V$ with $\gamma_0(L_{i,\ell})=L_{0,\ell}$ for
every prime $\ell\mid b$. So $\gamma_0^{-1}\pi\gamma_0$ preserves $L_{i,\ell}$ for $\ell\mid b$, and $\pi$ preserves
$L_{i,\ell}=L_{0,\ell}$ for $\ell\nmid b$. The space $V$ is indefinite of dimension at least $3$, so it contains a
nonzero isotropic vector, and $\gamma_0=s\gamma_1$ with $s$ a unitary automorphism that commutes with $\pi$ and
$\gamma_1$ in the group generated by the unitary transvections. By strong approximation in this group
\cite{Kne65}, it contains an element $u$ such that $u^{-1}\gamma_1$ preserves $L_{i,\ell}$ for $\ell\mid b$ and
$u$ preserves $L_{i,\ell}$ for $\ell\nmid b$. Then
$u^{-1}\pi u=(u^{-1}\gamma_1)(\gamma_0^{-1}\pi\gamma_0)(\gamma_1^{-1}u)$ preserves $L_{i,\ell}$ for every $\ell$,
hence $L_i$, and $L_i$ is the orthogonal sum of $L_i\cap u^{-1}V_+$ and $L_i\cap u^{-1}V_-$, where $V_-$ is the
kernel of $\pi$. These have ranks $p$ and $q$ and the signs of $S^0_+$ and $S^0_-$, so $L_i$ is split.
\end{proof}

\begin{lemma}\label{cv-components-cm-point}
In Situation~\ref{cv-components-situation}, let $C$ be a CM curve of level one for $\theta$, let $\tau$ be an
automorphism of $\mathbf C$ with $\tau(\theta)=\theta$, and let $\iota$ be a $\tau$-isomorphism of $C$
(Definition~\ref{cv-components-conjugation}). Let $P$ be a period of $L_i$ with $L_i=L_+(P)+L_-(P)$, and let
$P'$ be a period of $L_{i'}$ with $y_i(P)^\tau=y_{i'}(P')$. Then there is an $\mathcal O_E$-linear isomorphism
$\Phi\colon L_i\to L_{i'}$ with the following properties.
\begin{enumerate}
\item $\langle\Phi a,\Phi b\rangle_{i'}=\langle a,b\rangle_i$ for $a,b\in L_i$.
\item For $\varepsilon=\pm1$ and $a\in L_i$, $\Phi(a)\in L_\varepsilon(P')$ if and only if $a\in L_\varepsilon(P)$.
\item For every $n\ge1$, every $w\in\mathcal O$ by which $\iota$ acts on the points of order $n$, and all
  $a\in L_+(P)$ and $b\in L_-(P)$,
  \[
   \exp_{i,P}\big(P(a+b)/n\big)^\tau=\exp_{i',P'}\big(P'(\Phi(\hat wa+\bar{\hat w}b))/n\big).
  \]
\item If $\iota$ acts by $w$ on the points of order $M'$, then $t_{i'}(v)-\Phi(\hat wa+\bar{\hat w}b)\in
  M'L_{i'}$ for every $v\in\bar{\mathfrak M}'L$ and all $a\in L_+(P)$ and $b\in L_-(P)$ with $t_i(v)=a+b$.
\end{enumerate}
\end{lemma}
\leanref{Development/AbelianVarietyTheory/Proofs/UnitaryExistence/MarkingAtConjugateSplitPointTransportedTwistedStatement.lean}{GlobalTheory.T.marking_at_conjugate_splitPoint_transported_twisted}
\leanref{Development/AbelianVarietyTheory/Proofs/UnitaryExistence/ExistsLatticeIsometryConjSplitPointAnyLevelStatement.lean}{GlobalTheory.T.exists_lattice_isometry_conj_split_point_any_level}
\leanref{Development/AbelianVarietyTheory/Proofs/UnitaryExistence/ExistsLatticeLinearMapConjugateSplitPointStatement.lean}{GlobalTheory.T.exists_lattice_linearMap_conjugate_splitPoint}
\leanref{Development/AbelianVarietyTheory/Proofs/UnitaryExistence/ExistsConjSplitPointLatticeLinearMap.lean}{GlobalTheory.ConjSplit.exists_param}
\leanref{Development/AbelianVarietyTheory/Proofs/UnitaryExistence/ExistsLatticeLinearMapConjugateSplitPoint.lean}{GlobalTheory.ConjSplitAnyLevel.exists_lin}
\leanref{Development/AbelianVarietyTheory/Proofs/UnitaryExistence/ExistsIsometryOfConjugateSplitPoint.lean}{GlobalTheory.CompRat.ConjIsometry.conj_pairing}
\leanref{Development/AbelianVarietyTheory/Proofs/UnitaryExistence/ExistsIsometryOfConjugateSplitPoint.lean}{GlobalTheory.CompRat.ConjIsometry.surjective_of_isometry}
\begin{proof}
Put $y=y_i(P)$ and $y'=y_{i'}(P')$. The abelian scheme $A$, the action of $\mathcal O_E$, the polarization and
the sections $\eta(v)$ come from $X'$ over $\mathbf Q(\sqrt{-d})$, and $\tau$ fixes $\theta$, so conjugation by
$\tau$ is an isomorphism of groups $A_y(\mathbf C)\to A_{y'}(\mathbf C)$ that commutes with $\mathcal O_E$,
carries $\eta(v)(y)$ to $\eta(v)(y')$, and transforms Weil pairings by $\tau$.

Step 1: a parameter for $\iota$. For $z\in\mathbf C$ choose $f(z)\in\mathbf C$ with $\iota([z])=[f(z)]^\tau$. As in
the proof of Lemma~\ref{cv-components-multiplier}(1), $\iota$ and conjugation are isomorphisms of groups that
commute with $[w_+]$, so $f(z+z')-f(z)-f(z')$ and $f(cz)-cf(z)$ lie in $\mathcal O$ for $z,z'\in\mathbf C$ and
$c\in\mathcal O$, and every complex number is congruent to some $f(z)$ modulo $\mathcal O$. If $\iota$ acts by $w$
on the points of order $n$, then $f(w/n)-1/n\in\mathcal O$.

Step 2: the maps $\Phi_\pm$. Let $\varepsilon=\pm1$ and $a\in L_\varepsilon(P)$. By the proof of
Lemma~\ref{cv-components-points}(2), $\exp_{i,P}(zP(a))$ depends only on $z$ modulo $\mathcal O$, and that lemma
gives a homomorphism $\phi\colon C\to A_y$ with $\phi([z])=\exp_{i,P}(zP(a))$. Then $\psi=\phi^\tau\circ\iota\colon
C\to A_{y'}$ is a homomorphism with $\psi([z])=\phi([f(z)])^\tau=\exp_{i,P}(f(z)P(a))^\tau$, and by Step 1
\[
 \psi([w_\varepsilon z])=\exp_{i,P}\big(w_\varepsilon f(z)P(a)\big)^\tau=\exp_{i,P}\big(f(z)P(w_0a)\big)^\tau
 =w_0\psi([z]).
\]
By Lemma~\ref{cv-components-points}(2) there is a unique $\Phi_\varepsilon(a)\in L_\varepsilon(P')$ with
\[
 \exp_{i,P}\big(f(z)P(a)\big)^\tau=\exp_{i',P'}\big(zP'(\Phi_\varepsilon(a))\big)\qquad(z\in\mathbf C).
\]
By the uniqueness, $\Phi_\varepsilon\colon L_\varepsilon(P)\to L_\varepsilon(P')$ is $\mathcal O_E$-linear. It is
injective: if $\Phi_\varepsilon(a)=0$, then $\exp_{i,P}(f(z)P(a))=0$ for every $z$, so $\exp_{i,P}(zP(a))=0$ for
every $z$ by Step 1, and $\mathbf CP(a)\subset P(L_i)$ forces $a=0$.

Step 3: the map $\Phi$. Since $L_+(P)\cap L_-(P)=0$ and $L_i=L_+(P)+L_-(P)$, putting
$\Phi(a+b)=\Phi_+(a)+\Phi_-(b)$ for $a\in L_+(P)$ and $b\in L_-(P)$ defines an $\mathcal O_E$-linear map
$\Phi\colon L_i\to L_{i'}$. As $L_+(P')\cap L_-(P')=0$ and $\Phi_\pm$ are injective, $\Phi$ is injective,
$\Phi(a+b)\in L_+(P')$ only if $b=0$, and $\Phi(a+b)\in L_-(P')$ only if $a=0$. This gives (2).

Step 4: torsion points. Let $\iota$ act by $w$ on the points of order $n$, and let $a\in L_+(P)$. By Step 1,
$f(w/n)-1/n\in\mathcal O$, so by Step 2
\[
 \exp_{i,P}\big(P(a)/n\big)^\tau=\exp_{i,P}\big(f(w/n)P(a)\big)^\tau=\exp_{i',P'}\big((w/n)P'(\Phi_+(a))\big)
 =\exp_{i',P'}\big(P'(\hat w\Phi_+(a))/n\big),
\]
because $wP'(a')=P'(\hat wa')$ for $a'\in L_+(P')$. In the same way
$\exp_{i,P}(P(b)/n)^\tau=\exp_{i',P'}(P'(\bar{\hat w}\Phi_-(b))/n)$ for $b\in L_-(P)$, because
$wP'(b')=P'(\bar{\hat w}b')$ for $b'\in L_-(P')$. Adding the two gives (3).

Step 5: the forms. Let $x=a+b$ and $x'=a'+b'$ with $a,a'\in L_+(P)$ and $b,b'\in L_-(P)$, let $n\ge1$, let $w$ be
as in Lemma~\ref{cv-components-multiplier}(1), and put $\zeta_n=e^{2\pi i/n}$. The Weil pairing of
$\exp_{i,P}(P(x)/n)$ and $\exp_{i,P}(P(x')/n)$ is $\zeta_n^{-\langle x,x'\rangle_i}$, the form giving the Weil
pairings (Theorem~\ref{cv-uniformization-family}). Conjugating by $\tau$ and using (3) and the Weil pairing on
$A_{y'}$ gives
\[
 \tau(\zeta_n)^{-\langle x,x'\rangle_i}
 =\zeta_n^{-\langle\Phi(\hat wa+\bar{\hat w}b),\,\Phi(\hat wa'+\bar{\hat w}b')\rangle_{i'}}.
\]
The modules $L_+(P')$ and $L_-(P')$ are orthogonal (Lemma~\ref{cv-cm-point-gram}(5)), and
$\langle\hat wu,\hat wu'\rangle_{i'}=\langle u,\bar{\hat w}\hat wu'\rangle_{i'}=N(w)\langle u,u'\rangle_{i'}$, and the
same for $\bar{\hat w}$; so the exponent on the right is $-N(w)\langle\Phi x,\Phi x'\rangle_{i'}$. By
Lemma~\ref{cv-components-multiplier}(2), $\zeta_n^{N(w)}=\tau(\zeta_n)$, a primitive $n$-th root of unity.
Hence $n$ divides $\langle\Phi x,\Phi x'\rangle_{i'}-\langle x,x'\rangle_i$. As this holds for every $n$, (1)
follows.

Step 6: $\Phi$ is surjective. Let $u\in L_{i'}$. The form of $L_i$ is perfect, so some $x\in L_i$ has
$\langle x,z\rangle_i=\langle u,\Phi z\rangle_{i'}$ for every $z\in L_i$, and by (1) $u-\Phi x$ is orthogonal to
$\Phi(L_i)$. Both lattices have rank $p+q$, so $c(u-\Phi x)=\Phi v$ for some nonzero $c\in\mathcal O_E$ and some
$v\in L_i$. Then $\langle v,z\rangle_i=\langle\Phi v,\Phi z\rangle_{i'}=\langle u-\Phi x,\bar c\,\Phi z\rangle_{i'}=0$
for every $z$, so $v=0$, $c(u-\Phi x)=0$ and $u=\Phi x$.

Step 7: the level structure. Let $\iota$ act by $w$ on the points of order $M'$, let $v\in\bar{\mathfrak M}'L$, and
let $t_i(v)=a+b$ with $a\in L_+(P)$ and $b\in L_-(P)$. By Theorem~\ref{cv-uniformization-family} and (3),
\[
 \exp_{i',P'}\big(P'(t_{i'}(v))/M'\big)=\eta(v)(y')=\eta(v)(y)^\tau=\exp_{i,P}\big(P(t_i(v))/M'\big)^\tau
 =\exp_{i',P'}\big(P'(\Phi(\hat wa+\bar{\hat w}b))/M'\big).
\]
The kernel of $\exp_{i',P'}$ is $P'(L_{i'})$ and $P'$ is injective, so
$t_{i'}(v)-\Phi(\hat wa+\bar{\hat w}b)\in M'L_{i'}$, which is (4).
\end{proof}

\begin{lemma}\label{cv-components-cm-fixed}
In Situation~\ref{cv-components-situation}, let $\tau$ be an automorphism of $\mathbf C$ with
$\tau(\theta)=\theta$ that fixes a CM curve with its points of order $M'$
(Definition~\ref{cv-components-conjugation}), and let $P$ be a period of $L_i$ with $L_i=L_+(P)+L_-(P)$. Then
$y_i(P)^\tau=y_i(P)$.
\end{lemma}
\leanref{Development/AbelianVarietyTheory/Proofs/UnitaryExistence/PeriodSpacesFixedByCmCurveTorsionFixingAutomorph.lean}{GlobalTheory.CompRat.AnyLevel.conj_point_eq_of_isSplit}
\leanref{Development/AbelianVarietyTheory/Proofs/UnitaryExistence/ExistsComplexLinearOfPureTypes.lean}{GlobalTheory.CompRat.exists_complex_linear_of_pure_types}
\begin{proof}
Since $\tau(\theta)=\theta$, $y_i(P)^\tau$ is a complex point of $X'$ over $\sigma_0$, so
$y_i(P)^\tau=y_{i'}(P')$ for some $i'$ and some period $P'$ of $L_{i'}$ (Theorem~\ref{cv-uniformization-family}(1)).
Let $\iota$ be a $\tau$-isomorphism of a CM curve of level one that acts by $1$ on the points of order $M'$,
and let $\Phi$ be as in Lemma~\ref{cv-components-cm-point}. By its part (4) with $w=1$,
$\Phi(t_i(v))-t_{i'}(v)\in M'L_{i'}$ for $v\in\bar{\mathfrak M}'L$, and $\Phi$ preserves the forms by part (1).

Let $h$ be the $\mathbf R$-linear automorphism of $\mathbf C^g$ with $h(P(x))=P'(\Phi x)$ for $x\in L_i$; it exists
because $P(L_i)$ and $P'(L_{i'})$ are lattices in $\mathbf C^g$ and $\Phi$ is bijective. Let $\varepsilon=\pm1$
and $a\in L_\varepsilon(P)$. Then $\Phi(a)\in L_\varepsilon(P')$ by part (2), and $\theta P(a)=\varepsilon
P(\sqrt{-d}\,a)$, so
\[
 h(\theta P(a))=\varepsilon P'(\Phi(\sqrt{-d}\,a))=\varepsilon P'(\sqrt{-d}\,\Phi a)=\theta P'(\Phi a)=\theta
 h(P(a)).
\]
As $L_i=L_+(P)+L_-(P)$ and $P(L_i)$ spans $\mathbf C^g$ over $\mathbf R$, $h$ commutes with multiplication by
$\theta\notin\mathbf R$, so $h$ is $\mathbf C$-linear. Then $P'\circ\Phi=h\circ P$, and
Lemma~\ref{cv-components-points}(1) gives $y_i(P)=y_{i'}(P')=y_i(P)^\tau$.
\end{proof}

\begin{lemma}\label{cv-components-permutation}
In Situation~\ref{cv-components-situation}, let $\tau$ be an automorphism of $\mathbf C$ with
$\tau(\theta)=\theta$. Then there is a permutation $\pi$ of $\{1,\dots,s\}$ with $y^\tau\in Y_{\pi(i)}$ for every
$i$ and every $y\in Y_i$.
\end{lemma}
\leanref{Development/AbelianVarietyTheory/Proofs/UnitaryExistence/ExistsPeriodSpacePermOfComplexAutStatement.lean}{GlobalTheory.T.exists_periodSpace_perm_of_complexAut}
\leanref{Development/AbelianVarietyTheory/Proofs/UnitaryExistence/ExistsPeriodSpacePermOfComplexAut.lean}{GlobalTheory.ConjPermutesAnyLevel.exists_map}
\begin{proof}
Since $\tau(\theta)=\theta$, the conjugate of a complex point of $X'$ over $\sigma_0$ is again one, and it lies in
exactly one component (Theorem~\ref{cv-uniformization-family}(1)). For each $i$ fix $y^0_i\in Y_i$, and let
$\pi(i)$ be the index with $(y^0_i)^\tau\in Y_{\pi(i)}$. Let $y\in Y_i$, with $y^\tau\in Y_{i'}$.
Lemma~\ref{cv-components-points}(4), for $F''=F_c(e^{2\pi i/M'})$, gives a number field $F'''$ and an idempotent
regular function $e=e_{\pi(i)}$ on $X'_{F'''}$. Let $j=\tau^{-1}|_{F'''}$, an embedding with $j(\theta)=\theta$,
and let $x$ and $x_i$ be the complex points of $X'_{F'''}$ over $j$ that lie over $y$ and $y^0_i$. Their
conjugates $x^\tau$ and $x_i^\tau$ lie over the inclusion $F'''\subset\mathbf C$ and over $y^\tau$ and $(y^0_i)^\tau$, so
$e(x^\tau)$ is $1$ if $i'=\pi(i)$ and $0$ if not, and $e(x_i^\tau)=1$. The values of $e$ lie in $\{0,1\}$, so
Lemma~\ref{cv-components-points}(3) gives $e(x)=e(x_i)$. The value of a regular function on $X'_{F'''}$ at a
conjugate point is the conjugate of its value, so $e(x^\tau)=\tau(e(x))=\tau(e(x_i))=e(x_i^\tau)=1$, and
$i'=\pi(i)$. The same argument for $\tau^{-1}$ gives a map $\pi'$ with $\pi'(\pi(i))=i$, since $y^0_i$ is the
conjugate of $(y^0_i)^\tau$ by $\tau^{-1}$. So $\pi$ is injective, hence a permutation.
\end{proof}

\begin{lemma}\label{cv-deep-components}
In Situation~\ref{cv-cusp-setting}, let $(M',\mathfrak M')$ be a level
(Definition~\ref{cv-deep-cover}) and let $X'$ be the PEL moduli scheme of the PEL datum of $X$ at this
level (Theorem~\ref{cv-deep-pel-facts}(1)). Let $F'\subset\mathbf C$ be a number field that contains
$F_c$, the $M'$-th roots of unity and the Weber values of level $M'$ of
Lemma~\ref{cv-deep-ray-field}, for the square root of $-d$ fixed by $\mathcal D$. Then $F'$ is
algebraically closed in the function field of every connected component of $X'\otimes F'$.
\end{lemma}
\leanref{Development/AbelianVarietyTheory/Proofs/UnitaryExistence/GeometricallyIrreducibleComponentsWeberValues.lean}{GlobalTheory.FinalJoin.FieldAlg.fnAlgClosed_of_weberIn}
\leanref{Development/AbelianVarietyTheory/Proofs/HeightProofs/CmCurveTorsionDefinedOverOfWeberValues.lean}{GlobalTheory.FinalJoin.CurveField.curveOver_of_weberIn}
\leanref{Development/AbelianVarietyTheory/Proofs/UnitaryExistence/FnAlgClosedOfCurveOverLcm.lean}{GlobalTheory.CompRat.fnAlgClosed_of_curveOver_lcm}
\leanref{Development/AbelianVarietyTheory/Proofs/UnitaryExistence/PeriodSpacesFixedByCmCurveTorsionFixingAutomorphStatement.lean}{GlobalTheory.T.periodSpaces_fixed_by_cmCurveTorsionFixingAutomorphism}
\leanref{Development/AbelianVarietyTheory/Proofs/UnitaryExistence/CoeffFieldAlgClosedInFunctionFieldOfAutFixPeriodStatement.lean}{GlobalTheory.T.coeffField_algClosed_in_functionField_of_aut_fix_periods}

\begin{proof}
Let $\theta$ be the square root of $-d$ fixed by $\mathcal D$, and let $\tau$ be an automorphism of
$\mathbf C$ that fixes $F'$. Fix a type of signature $(p,q)$ for $\theta$, and let the lattices $L_i$, the
points $y_i(P)$ and the components $Y_1,\dots,Y_s$ be as in Situation~\ref{cv-components-situation}.

Step 1: the elliptic curve. The ring of integers of $\mathbf Q(\sqrt{-d})$ is the lattice $\mathcal O$ of
Definition~\ref{cv-components-conjugation}, so the Weber values of level $M'$ of
Lemma~\ref{cv-deep-ray-field} are the numbers of Lemma~\ref{cv-components-weber-curve} for $n=M'$. The
field $F_c$ contains a square root of $-d$, so $F'$ contains $\theta$ and $\tau(\theta)=\theta$. The
automorphism $\tau$ fixes the Weber values, so it fixes a CM curve with its points of order $M'$.

Step 2: the components. As $\tau(\theta)=\theta$, Lemma~\ref{cv-components-permutation} gives a permutation $\pi$ with
$y^\tau\in Y_{\pi(i)}$ for every $i$ and every $y\in Y_i$. Alp\"oge constructs from $\mathcal D$ a PEL datum
of signature $(1,0)$ over $E$, so by Lemma~\ref{cv-components-split} each $L_i$ has a period $P$ with
$L_i=L_+(P)+L_-(P)$. By Step 1 and Lemma~\ref{cv-components-cm-fixed}, $y_i(P)^\tau=y_i(P)\in Y_i$, so
$\pi(i)=i$. Hence $y^\tau\in Y_i$ for every $i$ and every $y\in Y_i$.

Step 3: the function fields. Let $Z$ be a connected component of $X'\otimes F'=X'_{F'}$, with function
field $F'(Z)$, and let $r\in F'(Z)$ be algebraic over $F'$, with minimal polynomial $\mu$. The field
$F'(Z)$ is countable, so it has an embedding $\varphi\colon F'(Z)\to\mathbf C$ over $F'$. It gives a complex
point $x$ of $Z$ over the inclusion $F'\subset\mathbf C$, at which a regular function takes the image
under $\varphi$ of its value at the generic point of $Z$. Since $X'$ is smooth, $X'_{F'}$ is normal, and
$r$ is integral over $F'$, so there is a regular function $u$ on $X'_{F'}$ with value $r$ at the generic
point of $Z$ and with $u\,\mu(u)=0$. At the complex points of $X'_{F'}$ over the inclusion, $u$ takes
values in the finite set of $0$ and the roots of $\mu$. The point $x$ lies over a point of some $Y_i$,
and by Step 2 so does $x^\tau$, which is again a point over the inclusion because $\tau$ fixes $F'$.
Lemma~\ref{cv-components-points}(3) gives $u(x^\tau)=u(x)$. The value of $u$ at $x^\tau$ is the
conjugate of its value at $x$, so $\tau(\varphi(r))=\varphi(r)$. This holds for every automorphism of
$\mathbf C$ that fixes $F'$, and every embedding of $F'(\varphi(r))$ in $\mathbf C$ over $F'$ extends to such
an automorphism; so $\varphi(r)$ has no conjugate over $F'$ other than itself, and $\varphi(r)\in F'$.
Hence $r\in F'$.
\end{proof}
% Lean: fnAlgClosed_of_weberIn (6 <= r from two_le_m and m_add_three_le_k, the datum of signature (1,0)
% from RankOneDatum.exists_datum_rank_one, Step 1 from curveOver_of_weberIn), then
% CompRat.fnAlgClosed_of_curveOver_lcm with the family of closing_of_abelianSchemeLaws at the inclusion
% of KA l; Step 2 is T.periodSpaces_fixed_by_cmCurveTorsionFixingAutomorphism
% (CompRat.AnyLevel.permutes_id_of_fixesCurve) and Step 3 is
% T.coeffField_algClosed_in_functionField_of_aut_fix_periods
% (CompRat.FixedConstLevel.fnAlgClosed_of_permutes_id, with exists_section, ev_conj, mem_range_of_fixed).

\begin{remark}\label{cv-deep-components-literature}
Deligne's reciprocity law for the connected components of Shimura varieties \cite{Del79} describes the
action of the automorphisms of $\mathbf C$ that fix the reflex field on the set of connected components
through class field theory. Lemma~\ref{cv-deep-components} says that this action is trivial on the
automorphisms that fix $F_c$, the $M'$-th roots of unity and the Weber values of level $M'$; these
values generate the ray class field of $\mathbf Q(\sqrt{-d})$ of modulus $M'$ over
$\mathbf Q(\sqrt{-d})$ \cite[Theorem~II.5.6]{Sil94}.
\end{remark}

\begin{proposition}\label{cv-deep-cover-exists}
In Situation~\ref{cv-cusp-setting}, suppose that the level of $X$ is a level $(M,\mathfrak M)$ with $M$
prime to $15$ (Definition~\ref{cv-deep-cover}), and that $-d$ is a nonzero square modulo $3$ and modulo
$5$. Then for every integer $e\ge0$, $X$ has a deep cover of exponent $e$ and size
\[
 \log[F_c:\mathbf Q]+\log\operatorname{rd}_{F_c}+20\big(\log(15^eM)+\log d+1\big).
\]
\end{proposition}
\leanref{Development/AbelianVarietyTheory/Proofs/DivisorProofs/ExistsGoodCoverOverSizedFieldStatement.lean}{GlobalTheory.T.exists_good_cover_over_sized_field}
\leanref{Development/AbelianVarietyTheory/Proofs/DivisorProofs/ExistsGoodCoverOverSizedField.lean}{GlobalTheory.CovTop.SizedCover.exists_good_sized}
\leanref{Development/AbelianVarietyTheory/Proofs/FoundationProofs/ExistsHalfOfBasePins.lean}{GlobalTheory.CoverHalf.exists_coverLevel}
\leanref{Development/AbelianVarietyTheory/Proofs/UnitaryExistence/ExistsDeltaDatumHigherLevel.lean}{GlobalTheory.LevelCover.exists_deltaDat_isLev}
\leanref{Development/AbelianVarietyTheory/Proofs/CuspFormEvaluations/FramedPointsLiftToLevelCoverStatement.lean}{GlobalTheory.T.framedPoints_lift_to_levelCover}
\leanref{Development/AbelianVarietyTheory/Definitions/Divisors/NormNegCmConstantEqAtAllEmbeddings.lean}{GlobalTheory.CovTop.Inhab.exists_good_of_cover}
\leanref{Development/AbelianVarietyTheory/Definitions/Divisors/NormNegCmConstantEqAtAllEmbeddings.lean}{GlobalTheory.CovTop.Inhab.exists_good_pin}
\begin{proof}
Let $n=15^eM$.

Step 1: the deep level. Lemma~\ref{cv-deep-ideal} gives an ideal $I\subset\mathcal O_E$, generated by
elements of $\mathbf Z[w]$, with $I\bar I=15^e\mathcal O_E$ and $I+\bar I=\mathcal O_E$. Put
$(M_Y,\mathfrak M_Y)=(n,\mathfrak MI)$. This is a level. Indeed $n\ge M\ge3$ is odd; $\mathfrak MI$ is
generated by the products of generators of $\mathfrak M$ and of $I$, which lie in $\mathbf Z[w]$;
$\mathfrak M_Y\bar{\mathfrak M}_Y=(\mathfrak M\bar{\mathfrak M})(I\bar I)=n\mathcal O_E$; and
$\mathfrak M_Y+\bar{\mathfrak M}_Y=\mathcal O_E$, because each of $\mathfrak M$ and $I$ is coprime to
each of $\bar{\mathfrak M}$ and $\bar I$: $\mathfrak M+\bar{\mathfrak M}=I+\bar I=\mathcal O_E$, and
$\mathfrak M$ and $\bar{\mathfrak M}$ contain $M$ while $I$ and $\bar I$ contain $15^e$, which is prime
to $M$. Moreover $M\mid M_Y$ and $\mathfrak M_Y\subset\mathfrak M$.

Step 2: the moduli scheme. Theorem~\ref{cv-deep-pel-facts}(1) gives the PEL moduli scheme $Y$ of the
PEL datum of $X$ at the deep level; here $p+q=r\ge1$.

Step 3: the coefficient field. Let $R$ be the field of Lemma~\ref{cv-deep-ray-field} for $n$ and the
square root of $-d$ fixed by $\mathcal D$, and put $F_Y=F_cR$. It contains $F_c$ and the $n$-th roots
of unity. An automorphism of $\mathbf C$ that fixes $F_c$ maps $R$ to itself, since $R$ is Galois over
$\mathbf Q$, and so maps $F_cR$ to itself. Lemmas~\ref{cv-deep-compositum}
and~\ref{cv-deep-ray-field} give
\[
 \log[F_Y:\mathbf Q]+\log\operatorname{rd}_{F_Y}\le\log[F_c:\mathbf Q]+\log\operatorname{rd}_{F_c}
 +20(\log n+\log d+1).
\]

Step 4: the components. The field $F_Y$ contains $R$, hence the Weber values of level $n=M_Y$, and
Lemma~\ref{cv-deep-components} shows that $F_Y$ is algebraically closed in the function field of every
connected component of $Y\otimes F_Y$.

Step 5: the level map. Theorem~\ref{cv-deep-pel-facts}(2), for the levels $(M,\mathfrak M)$ and
$(M_Y,\mathfrak M_Y)$, gives $\pi$ and the homomorphism of universal abelian schemes, under which
$M_Y/M=15^e$ times the level structure of $Y$ is that of $X$. Let $y$ be a framed point of $X$ over an
algebraically closed field $k$ containing $F_Y$. Since $Y\otimes F_Y\to X\otimes F_Y$ is finite and
surjective, the fibre over the point of $y$ is a nonempty finite $k$-scheme, and it has a $k$-point
$y_2$. The universal abelian scheme over $Y$ is the pullback of the one over $X$, so the abelian
variety of $y_2$ is identified with that of $y$, compatibly with the action of $\mathcal O_E$ and with
the line bundle of the polarization. The bases of the eigenspaces and the generator of sections of $y$
carry over to $y_2$, which gives a framed point of $Y$ over $y$.

Step 6: the uniformizations. Theorem~\ref{cv-deep-pel-facts}(3), for the level $(M_Y,\mathfrak M_Y)$
and the field $F_Y$, gives at every good embedding the uniformization and the map with the properties
of Definition~\ref{cv-deep-cover}(5), for the principal congruence subgroup $K_Y$ of level $M_Y$.

These are the data of a deep cover of exponent $e$, and Step 3 bounds its size.
\end{proof}

\begin{corollary}\label{cv-deep-cover-sized}
Let $\varepsilon$ and a family be as in Theorem~\ref{cv-covolume-bound}. There is $x_0$ with the
following property. Let $P$ be a point at a scale $x\ge x_0$, use the notation of
Situation~\ref{cv-point}, and put $r=k+m-1$, $z=frx^\varepsilon$ and $e=\lfloor3x^\varepsilon/\log3\rfloor$.
Then for every $1\le j<m$, the moduli scheme $X$ of $U(W_j)$ at $P$, with its coefficient field $F_c$ and
its uniformizations, has a deep cover of exponent $e$ and size $391z^2$
(Definition~\ref{cv-deep-cover}).
\end{corollary}
\leanref{Development/AutomorphicArithmetic/ThetaCorrespondence/ThetaPairLifts/ThirdPlaceSeries/ExistsDeepCoveringLevel.lean}{GlobalTheory.T.covering_deepEx_of_prints}
\leanref{Development/AutomorphicArithmetic/ThetaCorrespondence/ThetaPairLifts/ThirdPlaceSeries/ExistsDeepCoveringLevel.lean}{GlobalTheory.CovTop.DeepExist.exists_deep_covering_level_alt}
\leanref{Development/AutomorphicArithmetic/ThetaCorrespondence/ThetaPairLifts/ThirdPlaceSeries/LiftLambdaBound.lean}{GlobalTheory.CovTop.DeepExistence}
\leanref{Development/AbelianVarietyTheory/Definitions/Divisors/CovolXOfCovol.lean}{GlobalTheory.CovTop.GoodCov}
\leanref{Development/AbelianVarietyTheory/Proofs/DivisorProofs/DeepExistenceStatement.lean}{GlobalTheory.CovTop.DeepExist.size_le}
\leanref{Development/AbelianVarietyTheory/Proofs/DivisorProofs/DeepExistenceStatement.lean}{GlobalTheory.CovTop.DeepExist.floorThreeLOverLogThree_mul_log_fifteen_le}
\leanref{Development/AbelianVarietyTheory/Definitions/Divisors/CovolXOfCovol.lean}{GlobalTheory.CovTop.RecTied.core_size_le}
\leanref{Development/AbelianVarietyTheory/Definitions/Divisors/CovolXOfCovol.lean}{GlobalTheory.CovTop.two_PO_le}
\begin{proof}
By the size conditions of the family there is $x_0$ such that at every point at a scale $x\ge x_0$:
$x^\varepsilon\ge10$, $f\ge8$, $\log g_k\le9x^\varepsilon+1$ for the level primes $g_0,g_1$,
$\log d\le x^\varepsilon$, $-d$ is a nonzero square modulo $3$ and modulo $5$, and $15$ is prime to
$g_0g_1$. Fix such a point and $j<m$. The level of $X$ is a level $(g_0g_1,\mathfrak M)$ in the sense of
Definition~\ref{cv-deep-cover}. The coefficient field bound gives $\log[F_c:\mathbf Q]\le P_0$ and
$\log\operatorname{rd}_{F_c}\le P_0$ with $P_0=19fr^2x^\varepsilon+198f^2x^\varepsilon$. Since
$fx^\varepsilon\ge2$ and $r^2x^\varepsilon\ge2$,
\[
 2P_0\le19fr^2x^\varepsilon\cdot fx^\varepsilon+198f^2x^\varepsilon\cdot r^2x^\varepsilon=217z^2.
\]
Proposition~\ref{cv-deep-cover-exists} gives a deep cover of exponent $e$ and size
$217z^2+20(\log(15^eg_0g_1)+\log d+1)$. As $e\le3x^\varepsilon/\log3$ and $\log15\le3\log3$,
$e\log15\le9x^\varepsilon$, so
\[
 \log(15^eg_0g_1)+\log d+1\le9x^\varepsilon+2(9x^\varepsilon+1)+x^\varepsilon+1=28x^\varepsilon+3.
\]
Since $fr\ge8$, $z\ge8x^\varepsilon\ge80$, and $20(28x^\varepsilon+3)\le70z+60\le174z^2$. A deep cover of
size $o$ is a deep cover of every size $o'\ge o$, so $X$ has a deep cover of exponent $e$ and size
$391z^2$.
\end{proof}

\subsection{Frames and theta functions}\label{subsection-cv-frames-theta-functions}

\begin{lemma}\label{cv-uniformization-torus-isometry}
Let $E$, $d$ and $\theta$ be as in Definition~\ref{cv-cm-point-periods}, let $L$ and $L'$ be lattices of
PEL data of signature $(p,q)$ over $E$ with forms $\langle\ ,\ \rangle$ and $\langle\ ,\ \rangle'$, and let
$P$ and $P'$ be periods of $L$ and $L'$ of the same type $k\mapsto\sigma_k$. Let
$\psi\colon\mathbf C^g/P(L)\to\mathbf C^g/P'(L')$ be an isomorphism of complex Lie groups that commutes
with the actions $z\mapsto(\sigma_k(b)z_k)_k$ of $b\in\mathcal O_E$ and preserves the Weil pairings on
the points of order $n$ for every $n\ge1$, where the Weil pairing of the classes of $P(x)/n$ and $P(y)/n$
is $\exp(-2\pi i\langle x,y\rangle/n)$, and likewise for $P'$ and $\langle\ ,\ \rangle'$. Then there are an $\mathcal O_E$-linear
isomorphism $g_1\colon L\to L'$ with $\langle g_1x,g_1y\rangle'=\langle x,y\rangle$ for all $x,y\in L$ and a
$\mathbf C$-linear automorphism $h$ of $\mathbf C^g$ with $P'\circ g_1=h\circ P$ and
$\psi(z+P(L))=h(z)+P'(L')$ for $z\in\mathbf C^g$.
\end{lemma}
\leanref{Development/AbelianVarietyTheory/Proofs/UnitaryExistence/LevelAutomorphismPeriodFamily.lean}{GlobalTheory.LevStab.exists_isometry_of_hol}
\leanref{Development/AbelianVarietyTheory/Proofs/JetHarmonics/ExistsIsometryOfSabMap.lean}{GlobalTheory.LevelCusp.exists_isometry_of_sabMap}
\begin{proof}
The isomorphism $\psi$ lifts to a $\mathbf C$-linear automorphism $h$ of $\mathbf C^g$ with
$h(P(L))=P'(L')$ \cite[Chapter~1]{BL04}. For $b\in\mathcal O_E$, the $\mathbf R$-linear maps
$h\circ m_b$ and $m_b\circ h$, with $m_b(z)=(\sigma_k(b)z_k)_k$, agree modulo the discrete group
$P'(L')$, so they are equal. Hence $g_1=P'^{-1}\circ h\circ P$ is an $\mathcal O_E$-linear bijection of $L$
onto $L'$, and $P'\circ g_1=h\circ P$. For $x,y\in L$ and $n\ge1$ the Weil pairings give
$\exp(-2\pi i\langle g_1x,g_1y\rangle'/n)=\exp(-2\pi i\langle x,y\rangle/n)$, so the integers
$\langle g_1x,g_1y\rangle'$ and $\langle x,y\rangle$ are congruent modulo every $n$, hence equal.
\end{proof}
% Lean: exists_isometry_of_hol (exists_linearEquiv_of_period for h, realLinear_eq_of_sub_mem_discrete for
% the O_E-linearity, int_eq_of_exp_eq with weilIs_unique_fieldPoint for the forms); exists_isometry_of_sabMap
% applies it to a map of universal abelian schemes over the moduli scheme (exists_pointEquiv,
% weilIs_grpPullback_of_comm).

\begin{lemma}\label{cv-uniformization-model-match}
In Situation~\ref{cv-cusp-setting}, let $L$ and $L'$ be lattices of PEL data over $E$ with Hermitian models
$(G,\eta,\varphi)$ and $(G,\eta,\varphi')$ with the same $G$ and $\eta$
(Definition~\ref{cv-uniformization-model}), and let $g_1\colon L\to L'$ be an $\mathcal O_E$-linear
isomorphism that preserves the forms. Then there is $\delta\in GL_r(E)$ with
$\varphi'\circ g_1=\delta\circ\varphi$. It satisfies $h_G(\delta u,\delta v)=h_G(u,v)$ for $u,v\in E^r$, and
it carries $\varphi(L)$ onto $\varphi'(L')$.
\end{lemma}
\leanref{Development/AbelianVarietyTheory/Proofs/DivisorProofs/MovedSectionReadingEqMovedFrameReading.lean}{GlobalTheory.GeoLevel.exists_isMatch_of_models}
\leanref{Development/AbelianVarietyTheory/Proofs/UnitaryExistence/ExistsUnitaryMatrixOfLatticeIsometry.lean}{GlobalTheory.IsoExt.exists_unitary_of_isometry}
\begin{proof}
The images $\varphi(L)$ and $\varphi'(L')$ span $E^r$ and $\varphi$ and $\varphi'$ are injective, so
$\varphi'\circ g_1\circ\varphi^{-1}$ on $\varphi(L)$ extends to an $E$-linear automorphism $\delta$ of $E^r$.
For $x,y\in L$ and $b\in\mathcal O_E$,
\[
 \operatorname{Tr}_{E/\mathbf Q}\big(b\eta\,h_G(\delta\varphi x,\delta\varphi y)\big)
 =\langle g_1(bx),g_1y\rangle'=\langle bx,y\rangle
 =\operatorname{Tr}_{E/\mathbf Q}\big(b\eta\,h_G(\varphi x,\varphi y)\big).
\]
As the trace form is nondegenerate and $\eta\ne0$, $h_G(\delta\varphi x,\delta\varphi y)=h_G(\varphi x,\varphi y)$,
and by linearity $h_G(\delta u,\delta v)=h_G(u,v)$ on $E^r$. Finally
$\delta\varphi(L)=\varphi'(g_1L)=\varphi'(L')$.
\end{proof}
% Lean: exists_unitary_of_isometry (exists_gl_of_equiv, FrameFibre.herm_of_trace); exists_isMatch_of_models
% adds the lattice condition and IsMatch.

\begin{lemma}\label{cv-uniformization-frames-connected}
In Situation~\ref{cv-cusp-setting}, let $G\in M_r(E)$ be a matrix with $\bar G_{ab}=G_{ba}$ for all $a,b$
such that $\tau(G)$ has $q$ positive and $p$ negative eigenvalues for every $\tau\in\Phi$. Then the
frames for $G$ (the frames of $Z$ in Situation~\ref{cv-cusp-setting}, with $G$ in place of $G_Z$) form a
nonempty connected open subset of $\prod_{\tau\in\Phi}GL_r(\mathbf C)$.
\end{lemma}
\leanref{Development/AbelianVarietyTheory/Proofs/CuspFormEvaluations/ContinuousRmulPair.lean}{GlobalTheory.frames_isConnected}
\leanref{ReciprocityTheory/ReductiveGroups/ArithmeticGeometry/FundamentalDomains/FourierCoefficientDomain.lean}{ReciprocityTheory.ChartRestriction.G.FD.FourierCoefficientDomain.isOpen_dom}
\begin{proof}
For a family $x=(x_\tau)_{\tau\in\Phi}$ of invertible matrices write $\Gamma_\tau(x)=x_\tau^{\mathsf T}\tau(G)\bar x_\tau$, the
Gram matrix of the columns of $x_\tau$ for the Hermitian form $(u,v)\mapsto u^{\mathsf T}\tau(G)\bar v$, whose leading minors
are the $\delta_{\tau,i}(x)$, and let $\Gamma^+_\tau(x)$ be its leading $q\times q$ block.

The positive block. A family $x$ of invertible matrices is a frame for $G$ if and only if $\Gamma^+_\tau(x)$ is
positive definite for every $\tau$. If $x$ is a frame, the leading minors of $\Gamma^+_\tau(x)$ are the
$\delta_{\tau,i}(x)>0$ with $i\le q$, so it is positive definite by Sylvester's criterion. Conversely, suppose
that $\Gamma^+_\tau(x)$ is positive definite. Then $\delta_{\tau,i}(x)>0$ for $i\le q$. Let $i=q+k$ with $k\ge1$,
and write the leading $i\times i$ block of $\Gamma_\tau(x)$ as $\begin{pmatrix}A&B\\C&D\end{pmatrix}$ with
$A=\Gamma^+_\tau(x)$. Then $\delta_{\tau,i}(x)=\det A\det(D-CA^{-1}B)$, and $D-CA^{-1}B$ is the Gram matrix of the
projections of the columns $q+1,\dots,q+k$ of $x_\tau$ to the orthogonal complement of the span of the first $q$
columns. These projections are linearly independent, as $x_\tau$ is invertible. The form is positive definite
on the span of the first $q$ columns and $\tau(G)$ has only $q$ positive eigenvalues, so the form is negative
definite on that complement, and $\det(D-CA^{-1}B)$ has the sign $(-1)^k$.

Openness and the group $B$. The frames are the families with $\det x_\tau\ne0$ that satisfy finitely many strict
inequalities between the continuous real functions $\delta_{\tau,i}$ and $0$, so they form an open set. For
$b\in B$, $\Gamma_\tau(xb)=b_\tau^{\mathsf T}\Gamma_\tau(x)\bar b_\tau$, and as $b_\tau$ is upper triangular the leading $i\times i$
block of this matrix involves only the leading $i\times i$ blocks of $b_\tau$ and of $\Gamma_\tau(x)$. So
$\delta_{\tau,i}(xb)=\delta_{\tau,i}(x)\prod_{a\le i}|b_{\tau,aa}|^2$, and $xb$ is a frame if $x$ is.

Nonempty. As $\tau(G)$ has $q$ positive and $p$ negative eigenvalues, so has
$\overline{\tau(G)}=\tau(G)^{\mathsf T}$; a unitary diagonalization $U^*\,\overline{\tau(G)}\,U$ of it, with the
positive eigenvalues first, rescaled column by column by positive real numbers, gives $x^0_\tau$ with
$\overline{\Gamma_\tau(x^0)}=(x^0_\tau)^*\,\overline{\tau(G)}\,x^0_\tau=J$, the
diagonal matrix with entries $1$, $q$ times, followed by $-1$. As $J$ is real, $\Gamma_\tau(x^0)=J$, and $x^0$
is an orthonormal frame.

The chart. Let $n(z)$ and $\Omega=\Omega_{x^0}$ be as in Situation~\ref{cv-cusp-setting}, for $x^0$. Write
$n(z)_\tau=\begin{pmatrix}n_1&0\\w&n_2\end{pmatrix}$ in blocks of sizes $q$ and $p$, and put $Z_\tau=wn_1^{-1}$. Then
$\Gamma_\tau(x^0n(z))=n(z)_\tau^{\mathsf T}J\,\overline{n(z)_\tau}$, whose leading $q\times q$ block is
$n_1^{\mathsf T}\bar n_1-w^{\mathsf T}\bar w=n_1^{\mathsf T}(1-Z_\tau^{\mathsf T}\bar Z_\tau)\bar n_1$. As $x^0n(z)$ is
invertible, $z\in\Omega$ if and only if $1-Z_\tau^{\mathsf T}\bar Z_\tau$ is positive definite for every $\tau$. For $0\le t\le1$ let $z_t$ be $z$ with the entries of the blocks $w$
multiplied by $t$, and $z'_t$ be $z$ with these entries replaced by $0$ and the others multiplied by $t$. For
$z_t$ the blocks $n_1$ are those of $z$ and $Z_\tau$ becomes $tZ_\tau$, and
$1-t^2Z_\tau^{\mathsf T}\bar Z_\tau=(1-t^2)+t^2(1-Z_\tau^{\mathsf T}\bar Z_\tau)$ is positive definite if $z\in\Omega$; for $z'_t$, $Z_\tau=0$. So
for $z\in\Omega$ the paths $t\mapsto z_t$ and $t\mapsto z'_t$ in $\Omega$ join $z=z_1$ to $z_0=z'_1$ and $z'_1$ to
$z'_0=0$, and $\Omega$ is connected.

Dense decompositions. Let $\Delta(x)=\prod_\tau\prod_{i=1}^r\Delta_{\tau,i}(x)$, where $\Delta_{\tau,i}(x)$ is the
leading $i\times i$ minor of $(x^0_\tau)^{-1}x_\tau$. If $\Delta(x)\ne0$, Gaussian elimination gives a lower
unitriangular $n_\tau$ and an invertible upper triangular $b_\tau$ with $(x^0_\tau)^{-1}x_\tau=n_\tau b_\tau$, so
$x=x^0n(z)b$ with $b\in B$; if $x$ is moreover a frame, so is $x^0n(z)=xb^{-1}$, and $z\in\Omega$. Let $x$ be a frame
and $\varepsilon>0$. The function $t\mapsto\Delta(x+t(x^0-x))$ is a polynomial in $t$ with the value
$\Delta(x^0)=1$ at $t=1$, so it has finitely many zeros. As the frames are open, there is a real $t>0$, as small
as needed, that is not a zero and for which $x+t(x^0-x)$ is a frame within $\varepsilon$ of $x$.

Connectedness. The sets $\Omega$ and $B$ are connected, $B$ being a product of copies of $\mathbf C^\times$, of
$\mathbf C$ and of $\{0\}$. The map $(z,b)\mapsto x^0n(z)b$ on $\Omega\times B$ is continuous, so its image $S$ is
connected. By the paragraph on $B$, $S$ consists of frames, and by the previous paragraph every frame lies in the
closure of $S$. A set that lies between a connected set and its closure is connected, so the frames form a
connected set, which contains $x^0=x^0n(0)$.
\end{proof}
% Lean: frames_isConnected (GGConnStatement): the positive block is Data.gramPos_posDef and
% ExF.mem_dom_of_gramPos (psgn_dmin_le, and psgn_dmin_gt by the Schur complement fromBlocks); the group B is
% OntoS.dom_rmul; the orthonormal frame is x0F from the field signature of FD.Data; the chart is
% ExF.mem_OmF_iff with gramPos_eq (Am = 1 - Zt Zt^H, Zt the transpose of W n1^{-1}) and OFC.omFConn_of_cm (the
% paths S1 and S2); the decomposition is LUPf.lu_holds (DeltaF, split_eq, bF_mem, zF_mem, dense_delta through
% linePoly); B is connected by isPreconnected_bor; the end is IsPreconnected.subset_closure. The Lean's Gram
% matrix gram tau G x has entries hermC tau G (col x a) (col x b): it is the text's Gamma_tau(x).

\begin{lemma}\label{cv-uniformization-unit-minors}
In the setting of Lemma~\ref{cv-uniformization-frames-connected}, let $x$ be a frame for $G$. Then there is
a family $b\in B$ of diagonal matrices with positive real diagonal entries such that $xb$ is a frame for
$G$ with $|\delta_{\tau,i}(xb)|=1$ for all $\tau\in\Phi$ and $0\le i\le r$.
\end{lemma}
\leanref{Development/AbelianVarietyTheory/Proofs/GenericRegularity/ExistsUnitGramFrameBorelChange.lean}{GlobalTheory.FrameNorm.exists_unitFr_frBor}
\leanref{Development/AbelianVarietyTheory/Proofs/GenericRegularity/ExistsUnitGramFrameBorelChange.lean}{GlobalTheory.FrameNorm.UnitFr.norm_dmin_unitF}
\begin{proof}
For a family $b$ of diagonal matrices with positive real entries $b_{\tau,a}$, the leading $i\times i$
block of $b_\tau^{\mathsf T}x_\tau^{\mathsf T}\tau(G)\bar x_\tau\bar b_\tau$ is that of $x_\tau^{\mathsf T}\tau(G)\bar x_\tau$ multiplied on both sides
by the diagonal matrix of $b_{\tau,1},\dots,b_{\tau,i}$, so
$\delta_{\tau,i}(xb)=\delta_{\tau,i}(x)\prod_{a\le i}b_{\tau,a}^2$. The minors of $xb$ have the signs of
those of $x$, so $xb$ is a frame. Put $b_{\tau,a}=|\delta_{\tau,a-1}(x)/\delta_{\tau,a}(x)|^{1/2}$; the
product telescopes, and as $\delta_{\tau,0}=1$, $|\delta_{\tau,i}(xb)|=|\delta_{\tau,i}(x)|\cdot
|\delta_{\tau,i}(x)|^{-1}=1$.
\end{proof}
% Lean: unitF = rmul x (unitB x), with unitB the diagonal of the real numbers bt, bt_sq and prod_gN.

\begin{lemma}\label{cv-uniformization-modulus}
Let $O$ be a connected open subset of a finite-dimensional complex vector space, and let $P$ and $Q$ be
holomorphic functions on $O$ with $|P(x)|=|Q(x)|$ for every $x\in O$. Then there is $c\in\mathbf C$
with $P=cQ$ on $O$.
\end{lemma}
\leanref{Development/AbelianVarietyTheory/Proofs/JetHarmonics/ExistsEqConstMulOfForallExistsNormEq.lean}{GlobalTheory.exists_eq_const_mul_of_forall_exists_norm_eq}
\leanref{Development/AbelianVarietyTheory/Proofs/JetHarmonics/ExistsEqConstMulOfForallExistsNormEq.lean}{GlobalTheory.exists_eventually_eq_mul_of_norm_eq}
\begin{proof}
If $O$ is empty there is nothing to prove; let $x\in O$. If $Q$ vanishes on a neighbourhood of $x$, so
does $P$, and $c=0$ works near $x$. Otherwise there is $z\in O$ near $x$ with $Q(z)\ne0$, and a ball
$D\subset O$ around $z$ on which $Q$ does not vanish. On $D$ the function $P/Q$ is holomorphic and has
absolute value $1$, so it attains its maximum modulus at $z$ and is constant, equal to some $c$, by the
maximum modulus principle; so $P=cQ$ on $D$. In both cases $P-cQ$ is holomorphic on the connected set
$O$ and vanishes on a nonempty open subset, so it vanishes on $O$ by the identity theorem.
\end{proof}

\begin{lemma}\label{cv-uniformization-factor-unique}
Let $g\ge1$, let $\Pi\subset\mathbf C^g$ be a subgroup that spans $\mathbf C^g$ over $\mathbf R$, and
let $(O_j)_{j\in J}$ be a cover of $\mathbf C^g$ by open sets with $O_j+a=O_j$ for $a\in\Pi$. Let
$c_{jj'}$ be functions on $O_j\cap O_{j'}$ and $m$ a function on $\Pi\times\mathbf C^g$. Let $F_j$ and
$F'_j$, $j\in J$, be holomorphic functions on $O_j$, with $F_j$ nowhere zero, such that for all $j,j'$, all
$a\in\Pi$ and all $z\in O_j$:
\begin{enumerate}
\item $F_j(z)=c_{jj'}(z)F_{j'}(z)$ and $F'_j(z)=c_{jj'}(z)F'_{j'}(z)$ if $z\in O_{j'}$;
\item $F'_j(0)=F_j(0)$ if $0\in O_j$;
\item $|F_j(z)|=m(a,z)|F_j(z+a)|$ and $|F'_j(z)|=m(a,z)|F'_j(z+a)|$.
\end{enumerate}
Then $F'_j=F_j$ on $O_j$ for every $j$.
\end{lemma}
\leanref{Development/AbelianVarietyTheory/Proofs/GenericRegularity/UnitFramesGeneratorRatioModulusOne.lean}{GlobalTheory.UnitFrameGen.frames_eq}
\leanref{Development/AbelianVarietyTheory/Proofs/GenericRegularity/UnitFramesGeneratorRatioModulusOne.lean}{GlobalTheory.UnitFrameGen.eq_of_norm_periodic}
\begin{proof}
By (1), the quotient $\rho(z)=F'_j(z)/F_j(z)$ for $z\in O_j$ does not depend on $j$, so it is a
holomorphic function on $\mathbf C^g$, and $\rho(0)=1$ by (2). By (3), $|\rho(z+a)|=|\rho(z)|$ for
$a\in\Pi$. As $\Pi$ spans $\mathbf C^g$ over $\mathbf R$, it contains an $\mathbf R$-basis
$v_1,\dots,v_{2g}$ of $\mathbf C^g$, and every $z$ is a point of the compact set
$\{\sum_kt_kv_k:0\le t_k\le1\}$ plus an element of $\Pi$. So $|\rho|$ is bounded, $\rho$ is constant by
Liouville's theorem, and $\rho=\rho(0)=1$.
\end{proof}

\begin{lemma}\label{cv-uniformization-level-lift}
In Situation~\ref{cv-cusp-setting}, let $\Lambda$ be the lattice of the PEL datum of $X$, let $L_0$, $L_W$ and
the ideal $\mathfrak A$ be as in Lemma~\ref{cv-uniformization-lattice}, let $\iota\colon\Lambda\to L_W$ be an
$\mathcal O_E$-linear isomorphism, and let $\hat L_0$ and $\hat L_W=\mathfrak A\hat L_0$ be the closures of $L_0$
and $L_W$ in $(E\otimes_F\mathbf A_{F,f})^r$. Let $(M',\mathfrak M')$ be a level
(Definition~\ref{cv-deep-cover}), let $\gamma$ be an $\mathcal O_E$-linear automorphism of
$\Lambda/\mathfrak M'\Lambda$, and let $s\colon\Lambda\to\Lambda$ be an $\mathcal O_E$-linear map that induces
$\gamma$. Then there is $k\in U(W_j)(\mathbf A_{F,f})$ with $k\hat L_0=\hat L_0$ and
$k\iota(a)-\iota(s(a))\in M'\hat L_W$ for every $a\in\bar{\mathfrak M}'\Lambda$.
\end{lemma}
\leanref{Development/AbelianVarietyTheory/Proofs/DivisorProofs/ExistsUnitaryPart.lean}{GlobalTheory.LevelReduction.exists_isLevLift}
\leanref{Development/AbelianVarietyTheory/Proofs/DivisorProofs/ExistsUnitaryPart.lean}{GlobalTheory.LevelReduction.exists_unitary_part}
\begin{proof}
Write $\hat{\mathcal O}_E=\mathcal O_E\otimes\hat{\mathbf Z}$, write $z\mapsto\bar z$ for the conjugation of
$E\otimes_F\mathbf A_{F,f}$, and extend $h_G$ to $(E\otimes_F\mathbf A_{F,f})^r$. The group
$U(W_j)(\mathbf A_{F,f})$ consists of the invertible matrices $k$ with $k^\dagger k=1$, where $z^\dagger$ is the
adjoint of a matrix $z$ for $h_G$, so that $h_G(zv,w)=h_G(v,z^\dagger w)$. As
$\mathfrak M'+\bar{\mathfrak M}'=\mathcal O_E$ and $\mathfrak M'\bar{\mathfrak M}'=M'\mathcal O_E$, an element of
$\bar{\mathfrak M}'\hat{\mathcal O}_E$ lies in $M'\hat{\mathcal O}_E$ if and only if it lies in
$\mathfrak M'\hat{\mathcal O}_E$, and $\hat{\mathcal O}_E/\mathfrak M'\hat{\mathcal O}_E=\mathcal O_E/\mathfrak M'$.

Coordinates. The ring $\hat{\mathcal O}_E$ is the product of the completions of $\mathcal O_E$ at its primes,
which are discrete valuation rings. So the $\hat{\mathcal O}_E$-module $\hat L_0$ is free of rank $r$.
Let $\beta\in GL_r(E\otimes_F\mathbf A_{F,f})$ be a matrix whose columns form a basis of $\hat L_0$. A matrix $k$
satisfies $k\hat L_0=\hat L_0$ if and only if $\beta^{-1}k\beta\in GL_r(\hat{\mathcal O}_E)$. As $\mathfrak A$ is a
nonzero ideal and every ideal class of $E$ contains an integral ideal prime to $M'$, there is $t\in E$ with
$t\mathfrak A\subset\mathcal O_E$ and $t\mathfrak A+M'\mathcal O_E=\mathcal O_E$; write $1=M'm+ta_0$ with
$m\in\mathcal O_E$ and $a_0\in\mathfrak A$. For $a\in\Lambda$ put $\rho(a)=t\beta^{-1}\iota(a)$. As
$\iota(a)\in L_W\subset\hat L_W=\mathfrak A\hat L_0$ and $t\mathfrak A\subset\mathcal O_E$, $\rho(a)$ lies in
$\hat{\mathcal O}_E^r$, and $\rho$ is $\mathcal O_E$-linear.

(i) If $z\in\hat L_W$ and $t\beta^{-1}z=M'u$ with $u\in\hat{\mathcal O}_E^r$, then $z\in M'\hat L_W$: indeed
$z=M'mz+a_0tz=M'(mz+a_0\beta u)$, and $mz$ and $a_0\beta u$ lie in $\hat L_W$.

(ii) $\rho(\Lambda)+M'\hat{\mathcal O}_E^r=\hat{\mathcal O}_E^r$. Let $s\in\hat{\mathcal O}_E^r$. Then
$a_0\beta s\in\hat L_W$, and as $L_W$ is dense in $\hat L_W$, $a_0\beta s=\iota(a)+M'w$ with $a\in\Lambda$ and
$w\in\hat L_W$. So $s=M'ms+ta_0s=\rho(a)+M'(ms+t\beta^{-1}w)$, and $t\beta^{-1}w\in\hat{\mathcal O}_E^r$.

(iii) If $\rho(a)\in\mathfrak M'\hat{\mathcal O}_E^r$, then $a\in\mathfrak M'\Lambda$. For $c\in\bar{\mathfrak M}'$,
$\rho(ca)=c\rho(a)$ has entries in $\mathfrak M'\bar{\mathfrak M}'\hat{\mathcal O}_E=M'\hat{\mathcal O}_E$, so
$\iota(ca)\in M'\hat L_W$ by (i). Then $\iota(ca)/M'$ lies in $E^r\cap\hat L_W=L_W$, so $ca\in M'\Lambda$. For
$b\in\mathfrak M'$ it follows that $bca\in M'\mathfrak M'\Lambda$, so every $x\in\mathfrak M'\bar{\mathfrak M}'$ has
$xa=M'a'$ with $a'\in\mathfrak M'\Lambda$. For $x=M'$ this gives $a=a'$, as $\Lambda$ is torsion free.

By (ii), as $M'\in\mathfrak M'$, and by (iii), reducing $\rho$ modulo $\mathfrak M'$ gives an
$\mathcal O_E$-linear isomorphism $\phi\colon\Lambda/\mathfrak M'\Lambda\to(\mathcal O_E/\mathfrak M')^r$. Let
$y\in GL_r(\hat{\mathcal O}_E)$, $k=\beta y\beta^{-1}$, and $a,a'\in\bar{\mathfrak M}'\Lambda$; then
$k\hat L_W=\hat L_W$, and
\[
 k\iota(a)-\iota(a')\in M'\hat L_W\iff\phi(a')=\bar y\,\phi(a),\tag{$*$}
\]
where $\bar y$ is the reduction of $y$ modulo $\mathfrak M'$. Indeed,
$t\beta^{-1}\big(k\iota(a)-\iota(a')\big)=y\rho(a)-\rho(a')$, whose entries lie in
$\bar{\mathfrak M}'\hat{\mathcal O}_E$. If $k\iota(a)-\iota(a')=M'w$ with $w\in\hat L_W$, then
$y\rho(a)-\rho(a')=M't\beta^{-1}w\in M'\hat{\mathcal O}_E^r$, and $\phi(a')=\bar y\,\phi(a)$. Conversely, if
$\phi(a')=\bar y\,\phi(a)$, the entries of $y\rho(a)-\rho(a')$ lie in $\mathfrak M'\hat{\mathcal O}_E$, hence in
$M'\hat{\mathcal O}_E$ by the first paragraph, and (i) applies to $z=k\iota(a)-\iota(a')\in\hat L_W$.

The matrix. The automorphism $\phi\gamma\phi^{-1}$ of $(\mathcal O_E/\mathfrak M')^r$ is given by the
reduction of a matrix $Y\in M_r(\mathcal O_E)$, and $\det Y$ is a unit modulo $\mathfrak M'$.

The idempotent. Let $T$ be the finite set of primes of $\mathcal O_E$ that contain $\mathfrak M'$. As
$\mathfrak M'+\bar{\mathfrak M}'=\mathcal O_E$, no prime of $T$ is the conjugate of a prime of $T$. Let
$e\in E\otimes_F\mathbf A_{F,f}$ have component $1$ at the primes of $T$ and $0$ at the others. Then $e^2=e$,
$e\bar e=0$, $e\in\hat{\mathcal O}_E$ and $e\equiv1$ modulo $\mathfrak M'\hat{\mathcal O}_E$. As $\det Y$ is a unit
at the primes of $T$, there is $u\in\hat{\mathcal O}_E$ with $u\det Y=e$. Put $m=\beta Y\beta^{-1}$ and
$m'=\beta(u\operatorname{adj}Y)\beta^{-1}$, with $\operatorname{adj}Y$ the adjugate. As
$Y\cdot u\operatorname{adj}Y=e$, $e\,mm'=e$, and then also $e\,m'm=e$. Both $m$ and $m'$ map $\hat L_0$ into
itself, as $Y$ and $u\operatorname{adj}Y$ have entries in $\hat{\mathcal O}_E$.

The element $k$. Put $k=1+e(m-1)+\bar e(m'^\dagger-1)$. As $(ez)^\dagger=\bar ez^\dagger$ and $e\bar e=0$,
$k^\dagger=1+\bar e(m^\dagger-1)+e(m'-1)$ and
\[
 k^\dagger k=1+e(m'm-1)+\bar e\big((m'm)^\dagger-1\big)=1,
\]
so $k\in U(W_j)(\mathbf A_{F,f})$, with $k^{-1}=k^\dagger$, $ek=em$ and $ek^{-1}=em'$. Let $w\in\hat L_0$. Then
$ekw=emw\in\hat L_0$ and $(1-e-\bar e)kw=(1-e-\bar e)w\in\hat L_0$. For $z\in\hat L_0$,
$h_G(z,\bar ekw)=e\,h_G(k^{-1}z,w)=h_G(ek^{-1}z,w)=h_G(em'z,w)\in\hat{\mathcal O}_E$. As $L_0$ is unimodular
(Definition~\ref{th-spaces}) and $h_G(v,w)=-h_{W_j}(w,v)$, $\hat L_0$ is the set of $w'$ with $h_G(z,w')\in\hat{\mathcal O}_E$ for all
$z\in\hat L_0$, so $\bar ekw\in\hat L_0$. Hence $kw\in\hat L_0$. The same argument with $m$ and $m'$ exchanged
gives $k^{-1}\hat L_0\subset\hat L_0$, so $k\hat L_0=\hat L_0$.

The congruence. The matrix $y=\beta^{-1}k\beta$ lies in $GL_r(\hat{\mathcal O}_E)$ and $ey=eY$; as
$e\equiv1$ modulo $\mathfrak M'\hat{\mathcal O}_E$, $\bar y$ is the reduction of $Y$. Let
$a\in\bar{\mathfrak M}'\Lambda$. Then $s(a)\in\bar{\mathfrak M}'\Lambda$, and $s(a)$ has the class $\gamma(a)$
modulo $\mathfrak M'\Lambda$, so $\phi(s(a))=\bar y\,\phi(a)$. By $(*)$, $k\iota(a)-\iota(s(a))\in M'\hat L_W$.
\end{proof}
% Lean: exists_isLevLift: the basis beta is T.exists_lattice_basis_hermitian_space; the coordinates and (*) are
% GeoActCompare.exists_coord, from exists_lattice_coord = coord_adelic, coord_frame and coord_core (rho = t phi
% with phi = beta^{-1} iota; t, a_0, m from exists_scale_coprime; (i) is key, (ii) the surjectivity through
% hat_rational_points_and_approx, (iii) the kernel through hrat) and exists_coord_of (the isomorphism and (*),
% with Ideal.mul_eq_inf_of_coprime); the matrix is exists_matrix; the idempotent and k are exists_stab_lift
% (exists_half_idem, conj_inv_part, exists_unitary_part, mem_St_of_part); the congruence is map_eq_of_smul_eq and
% graph_iff.

\begin{lemma}\label{cv-uniformization-scalar-orbits}
In the setting of Lemma~\ref{cv-uniformization-frames-connected}, let $\Gamma\subset GL_r(E)$ be a subgroup and
$\gamma\in GL_r(E)$, and suppose that for every frame $x$ for $G$ there are $\gamma'\in\Gamma$ and $b\in B$ with
$\gamma x=\gamma'xb$. Then there are $\gamma'\in\Gamma$ and $u\in E^\times$ with $\gamma'=u\gamma$.
\end{lemma}
\leanref{Development/AbelianVarietyTheory/Proofs/CubicalLineBundles/ExistsSmulOfPreservesFrameOrbits.lean}{GlobalTheory.OrbitScalar.exists_smul_of_orbits}
\begin{proof}
Let $V=\prod_{\tau\in\Phi}M_r(\mathbf C)$, and let $O\subset V$ be the set of frames for $G$, which is open and
nonempty (Lemma~\ref{cv-uniformization-frames-connected}). For $x\in V$ let $x_{\tau,1}$ be the first column of
$x_\tau$. For $\gamma'\in\Gamma$ let $C_{\gamma'}$ be the closed set of $x\in V$ such that for every $\tau$ the
vectors $\tau(\gamma)x_{\tau,1}$ and $\tau(\gamma')x_{\tau,1}$ are linearly dependent. If $x\in O$ and
$\gamma x=\gamma'xb$ with $b\in B$, then the first column of $x_\tau b_\tau$ is $b_{\tau,11}x_{\tau,1}$, as $b_\tau$
is upper triangular, so $\tau(\gamma)x_{\tau,1}=b_{\tau,11}\tau(\gamma')x_{\tau,1}$ and $x\in C_{\gamma'}$. So $V$ is
the union of the closed set $V\setminus O$ and the sets $C_{\gamma'}$, $\gamma'\in\Gamma$, which are countably many
since $E$ is countable. By the Baire category theorem the union of their interiors is dense in $V$. It meets
the nonempty open set $O$, which does not meet $V\setminus O$, so some $x\in O$ lies in the interior of some
$C_{\gamma'}$.

Fix $\tau_0\in\Phi$. Changing only the first column of $x$ at $\tau_0$ gives a nonempty open set
$U\subset\mathbf C^r$ such that $\tau_0(\gamma)v$ and $\tau_0(\gamma')v$ are linearly dependent for every $v\in U$.
As $\tau_0(\gamma')$ is invertible, every nonzero $v\in U$ is an eigenvector of
$A=\tau_0(\gamma')^{-1}\tau_0(\gamma)$. Choose $v_1\ne0$ in $U$, with $Av_1=cv_1$. For $w$ small, $v_1+w$ and
$v_1-w$ lie in $U$ and are nonzero, with eigenvalues $a$ and $a'$ say. Adding gives
$(a-c)(v_1+w)=(c-a')(v_1-w)$. If $w$ is not a multiple of $v_1$, then $a=c$ and $Aw=cw$; if it is, $Aw=cw$
directly. So $A=c$ on a neighbourhood of $0$, hence everywhere, and $\tau_0(\gamma)=c\,\tau_0(\gamma')$. Choose an
entry with $\gamma'_{kl}\ne0$ and put $u_0=\gamma_{kl}/\gamma'_{kl}\in E$. Then $c=\tau_0(u_0)$, and $\gamma=u_0\gamma'$
because $\tau_0$ is injective. As $\gamma$ is invertible, $u_0\ne0$, and $u=u_0^{-1}$ works.
\end{proof}

\begin{lemma}\label{cv-uniformization-root-congruence}
Let $E$ be a number field, let $r\ge1$ and $M\ge3$ be integers, and let $\Lambda,\Lambda'\subset E^r$ be finitely
generated $\mathcal O_E$-submodules that span $E^r$. Let $\gamma,\delta\in M_r(E)$ with $\gamma v-v\in M\Lambda$ for
$v\in\Lambda$ and $\delta v-v\in M\Lambda'$ for $v\in\Lambda'$. Let $u\in E$ be a root of unity,
$\tau\colon E\to\mathbf C$ an embedding and $Q\in GL_r(\mathbf C)$ with $\tau(u\gamma)Q=Q\tau(\delta)$. Then
$u=1$.
\end{lemma}
\leanref{Development/AbelianVarietyTheory/Proofs/CubicalLineBundles/RootOfUnityEqOneOfConjCongruence.lean}{GlobalTheory.root_of_unity_eq_one_of_conj_congruence}
\begin{proof}
Step 1: if $w$ is a nonzero row vector in $\mathbf C^r$ with $w\,\tau(\gamma)=lw$, then $(l-1)/M$ is an algebraic
integer. The additive map $\psi(v)=w\,\tau(v)$ on $E^r$ satisfies $\psi(\gamma v)=l\psi(v)$ and
$\psi(bv)=\tau(b)\psi(v)$ for $b\in E$. Its image $\psi(\Lambda)$ is a finitely generated subgroup of
$\mathbf C$, nonzero because $\Lambda$ spans $E^r$ and $w\ne0$. For $v\in\Lambda$,
$\frac{l-1}M\psi(v)=\psi\big((\gamma v-v)/M\big)\in\psi(\Lambda)$. So multiplication by $(l-1)/M$ preserves a
nonzero finitely generated subgroup of $\mathbf C$, and $(l-1)/M$ is an algebraic integer. The same holds for
$\delta$ and $\Lambda'$.

Step 2. Let $l$ be an eigenvalue of $\tau(\gamma)$, with a nonzero row eigenvector $w$. Then $wQ\ne0$ and
$(wQ)\tau(\delta)=w\,\tau(u\gamma)Q=\tau(u)l\,(wQ)$. By Step 1, $(l-1)/M$ and $(\tau(u)l-1)/M$ are algebraic
integers, and so is $\tau(u)$. Hence
\[
 \tau\Big(\frac{u-1}M\Big)=\frac{\tau(u)l-1}M-\tau(u)\,\frac{l-1}M
\]
is an algebraic integer, and so is $(u-1)/M$. Every conjugate of $u$ has absolute value $1$, so every conjugate
of $(u-1)/M$ has absolute value at most $2/M<1$. If $u\ne1$, then $(u-1)/M$ is a nonzero algebraic integer
whose conjugates have absolute value at most $1$; by Kronecker's theorem \cite{Kro57} it is a root of unity,
and its conjugates have absolute value $1$, a contradiction. So $u=1$.
\end{proof}

\begin{lemma}\label{cv-uniformization-block-free}
In the setting of Lemma~\ref{cv-uniformization-frames-connected}, let $\Gamma\subset GL_r(E)$ be a subgroup such
that, for $\gamma\in\Gamma$ and a frame $x$ for $G$, $\gamma x\in xB$ only if $\gamma=1$. Let $\gamma\in\Gamma$, let
$x$ be a frame for $G$, and let $m=(m_\tau)_{\tau\in\Phi}$ be a family of matrices in $GL_r(\mathbf C)$ with
$(m_\tau)_{ac}=0$ whenever $c\le q<a$, such that $\gamma x=xm$. Then $\gamma=1$.
\end{lemma}
\leanref{Development/AbelianVarietyTheory/Proofs/DivisorProofs/ExistsLevelBlock.lean}{GlobalTheory.GammaLevel.eq_one_of_block}
\leanref{Development/AbelianVarietyTheory/Proofs/IwasawaDecomposition/DomainPointStabilizerTrivial.lean}{GlobalTheory.FreeOnDomain.exists_conj_block}
\begin{proof}
Fix $\tau$ and write $m_\tau=\begin{pmatrix}A&C\\0&D\end{pmatrix}$ in blocks of sizes $q$ and $p$. By Schur's
triangularization there are invertible matrices $P$ and $R$ with $P^{-1}AP$ and $R^{-1}DR$ upper triangular. Put
$b_\tau=\begin{pmatrix}P&0\\0&R\end{pmatrix}$. Then $t_\tau=b_\tau^{-1}m_\tau b_\tau$ is upper triangular and
invertible, so $t=(t_\tau)_\tau\in B$. The leading $q\times q$ block of $x_\tau^{\mathsf T}\tau(G)\bar x_\tau$ is positive definite,
as the first $q$ leading minors of a frame are positive, and that of $b_\tau^{\mathsf T}x_\tau^{\mathsf T}\tau(G)\bar x_\tau\bar b_\tau$ is
$P^{\mathsf T}$ times it times $\bar P$, so it is positive definite too. As $\tau(G)$ has $q$ positive eigenvalues, the
orthogonal complement of the span of the first $q$ columns of $x_\tau b_\tau$ is negative definite, so the
leading $i\times i$ minors for $i>q$ have the signs $(-1)^{i-q}$. So $xb$ is a frame for $G$, and
$\gamma(xb)=xmb=(xb)t$ with $t\in B$. Hence $\gamma=1$.
\end{proof}
% Lean: eq_one_of_block applies Data.Gamma_free to the frame rmul x b (PointOfFrames.rmul_mem_dom, through
% AR.Up.ExF.mem_dom_of_gramPos) with b and t from exists_conj_block (exists_conj_triangular on the two diagonal
% blocks).

\begin{lemma}\label{cv-uniformization-unitary-orbit}
Let $r\ge1$, let $A$ and $B$ be Hermitian forms on $\mathbf C^r$ with $B$ nondegenerate, and let
$v\in\mathbf C^r$. Suppose that $|A(hv,hv)|=1$ for every linear automorphism $h$ of $\mathbf C^r$ that preserves
$B$. Then there is a real number $t\ne0$ with $A=tB$.
\end{lemma}
\leanref{Development/AbelianVarietyTheory/Proofs/CubicalLineBundles/EqZeroOfEventuallyDotProductStarEqZero.lean}{GlobalTheory.UnitaryOrbitForm.eq_smul_of_norm_eq_one_on_orbit}
\begin{proof}
Taking $h=1$ gives $|A(v,v)|=1$, so $v\ne0$. Let $w\ne0$ with $B(w,w)=B(v,v)$. By Witt's theorem for Hermitian
forms \cite[Chapter~7]{Sch85}, applied to the isometry $v\mapsto w$ of the lines spanned by $v$ and $w$, there
is a linear automorphism $h$ preserving $B$ with $hv=w$; so $|A(w,w)|=1$. Hence $c=B(v,v)$ is not $0$:
otherwise $w=2v$ would give $|A(2v,2v)|=4$. Put $s=A(v,v)$, a real number with $s=\pm1$. For $w$ near $v$,
$w\ne0$ and $cB(w,w)>0$ by continuity, and $w'=(c/B(w,w))^{1/2}w$ satisfies $B(w',w')=c$, so
$|A(w,w)c/B(w,w)|=|A(w',w')|=1$. The continuous function $w\mapsto A(w,w)c/B(w,w)$ takes only the values $\pm1$
near $v$ and the value $s$ at $v$, so it equals $s$ near $v$. Hence the Hermitian form $C=A-(s/c)B$ satisfies
$C(w,w)=0$ for all $w$ in a neighbourhood of $v$. For $u\in\mathbf C^r$ and small real $t\ne0$, adding
$C(v+tu,v+tu)=0$ and $C(v-tu,v-tu)=0$ and using $C(v,v)=0$ gives $2t^2C(u,u)=0$. So $C(u,u)=0$ for every
$u$, and $C=0$ by polarization. This is the claim with $t=s/c$.
\end{proof}
% Lean: eq_smul_of_norm_eq_one_on_orbit, with A and B as matrices and the form (u, v) -> u^T A conj(v); the
% transitivity is UnitaryOrbitForm.exists_unitary_mulVec_eq through HermWitt.exists_gl_isometry_mul_eq, and the
% last step eq_zero_of_eventually_dotProduct_star_eq_zero with eq_zero_of_forall_dotProduct_star_eq_zero.

\begin{lemma}\label{cv-uniformization-uniform-constant}
Let $O$ be a nonempty connected open subset of a finite-dimensional complex vector space, let $I$ be a set, and
for $i\in I$ let $n_i\ge0$ be an integer and $P_i$ and $Q_i$ holomorphic functions on $O$. Let $c$ be a function
on $O$ with $|c(x)|=1$ and $P_i(x)=c(x)^{n_i}Q_i(x)$ for all $i\in I$ and $x\in O$. Then there is $u\in\mathbf C$
with $|u|=1$ and $P_i=u^{n_i}Q_i$ on $O$ for every $i\in I$.
\end{lemma}
\leanref{Development/AbelianVarietyTheory/Proofs/DivisorProofs/LevelCoverPullbackReadingEqUnitConstPowMul.lean}{GlobalTheory.LevRead.eq_pow_mul_of_norm_one}
\begin{proof}
For each $i$, $|P_i|=|Q_i|$ on $O$, so $P_i=u_iQ_i$ for some $u_i\in\mathbf C$ by
Lemma~\ref{cv-uniformization-modulus}. For $i\in I$ let $C_i$ be the set of $t\in\mathbf C$ with $|t|=1$ and
$P_i=t^{n_i}Q_i$ on $O$; it is closed in the unit circle. Let $s\subset I$ be finite. There is a nonempty open
set $W\subset O$ on which no $Q_i$ with $i\in s$ that is not identically $0$ on $O$ vanishes: by induction on
$s$, if $W$ serves for $s$ and $Q_i$ is not identically $0$ on $O$, then $Q_i$ does not vanish identically on $W$,
by the identity theorem on the connected set $O$, and $W\cap\{Q_i\ne0\}$ serves for $s\cup\{i\}$. Let $x_0\in W$.
Then $c(x_0)\in C_i$ for every $i\in s$: if $Q_i$ is not identically $0$, then $u_iQ_i(x_0)=P_i(x_0)=
c(x_0)^{n_i}Q_i(x_0)$ with $Q_i(x_0)\ne0$, so $u_i=c(x_0)^{n_i}$; otherwise $P_i=c^{n_i}Q_i$ vanishes as well. So
the closed subsets $C_i$ of the unit circle have the finite intersection property, and as the circle is compact
they have a common point $u$.
\end{proof}
% Lean: eq_pow_mul_of_norm_one, with eq_const_mul_of_norm_one for u_i, exists_open_ne_zero for W, and
% IsCompact.inter_iInter_nonempty for the sphere.

\begin{lemma}\label{cv-uniformization-theta-marking}
In Situation~\ref{cv-cusp-setting}, let $L$ be the lattice of the PEL datum of $X$, and let $X'$ be the PEL
moduli scheme of this datum at a level $(M',\mathfrak M')$ (Definition~\ref{cv-deep-cover}), with a coefficient
field $F'$ and, at every good embedding of $F'$, a uniformization normalized by $\xi$ that satisfies the
compatibility conditions of the moduli problem, as $X$ has when its level is a level and as a deep cover of $X$
has (Definition~\ref{cv-deep-cover}(5)). Let $\sigma$ be a good embedding of $F'$ and $Z$ a component of the
uniformization at $\sigma$, with lattice $\Lambda_Z=\varphi_Z(L_Z)\subset E^r$ for the Hermitian model
$(G,\eta,\varphi_Z)$ of the lattice $L_Z$ of Theorem~\ref{cv-uniformization-family} that the compatibility
conditions attach to $Z$. For a frame $x$ of $Z$ and $w\in E^r$ let $P_x(w)\in\mathbf C^g$ be given by the
formula for periods of Situation~\ref{cv-uniformization-situation} with $\tau(w)$ in place of $\tau(\varphi_iv)$,
so that $P_x\circ\varphi_Z$ is the period $P^Z_x$ of $x$ on $L_Z$. Let $x$ be a frame of $Z$ with
$|\delta_{\tau,i}(x)|=1$ for all $\tau$ and $i$, and let $y$ be the framed point at $x$, with abelian variety
$A_y$.
\begin{enumerate}
\item There are a surjective homomorphism $p_x\colon\mathbf C^g\to A_y(\mathbf C)$ of complex Lie groups that
  commutes with $\mathcal O_E$, has the kernel $P_x(\Lambda_Z)$ and carries the unit vectors of $\mathbf C^g$,
  as tangent vectors at $0$, to the bases of $y$; a positive definite Hermitian form $H_x$ on $\mathbf C^g$,
  linear in the first variable; numbers $\chi_x(u)$ of absolute value $1$ for $u\in P_x(\Lambda_Z)$; a function
  $\vartheta_x$ on $\mathbf C^g$; and, for the trivializations of the line bundle of the polarization over the
  sets $U_j$ of Theorem~\ref{cv-uniformization-family}(4), nowhere vanishing holomorphic functions $F^x_j$ on
  the sets of the $z$ with $p_x(z)\in U_j$, such that the section of $y$ has the value
  $\vartheta_x(z)F^x_j(z)$ at $p_x(z)$ in the trivialization over $U_j$,
  \[
   F^x_j(z)=\chi_x(u)\exp\big(\pi H_x(z,u)+\tfrac\pi2H_x(u,u)\big)F^x_j(z+u)\qquad(u\in P_x(\Lambda_Z)),
  \]
  and the mean of $|\vartheta_x(z)|^2e^{-\pi H_x(z,z)}$ over a fundamental domain of $P_x(\Lambda_Z)$ is $1$.
\item There is an $\mathcal O_E$-linear map $m_Z\colon L\to\Lambda_Z$, the \emph{marking} of $Z$, which does
  not depend on $x$, such that the level structure of $y$ sends $a\in\bar{\mathfrak M}'L$ to
  $p_x\big(P_x(M'^{-1}m_Z(a))\big)$, and such that every $v\in\bar{\mathfrak M}'\Lambda_Z$ is $m_Z(a)+M'u$ with
  $a\in\bar{\mathfrak M}'L$ and $u\in\Lambda_Z$.
\item For every class $q\in L/\mathfrak M'L$ there are $a\in\bar{\mathfrak M}'L$ in the class $q$, a vector
  $Y\in E^r$ and an integer $j\ge0$ with $Y-M'^{-1}m_Z(a)\in\Lambda_Z$ and $sY\in2M'\Lambda_Z$ for every
  $s\in\mathfrak M'^j$.
\end{enumerate}
\end{lemma}
\leanref{Development/AbelianVarietyTheory/Proofs/SiegelThetaTransformation/AnalyticFibreUnitFrameUniformisation.lean}{GlobalTheory.Dict.uw_anFib}
\leanref{Development/AbelianVarietyTheory/Proofs/SiegelThetaTransformation/UniformisationComponentLatticesMarkings.lean}{GlobalTheory.Dict.uw_lattice}
\leanref{Development/AbelianVarietyTheory/Proofs/SiegelThetaTransformation/DualHalfLevelMarkingExistsAndMatch.lean}{GlobalTheory.Dict.DualMark.dual_marking}
\begin{proof}
(1) By the compatibility conditions, listed at the beginning of the proof of
Lemma~\ref{cv-cover-uniformization}, the point of $y$ is $y_Z(P^Z_x)$, and its bases are the images of the unit
vectors under the isomorphism $\mathbf C^g/P^Z_x(L_Z)\to A_y$ of Theorem~\ref{cv-uniformization-family}, which
commutes with $\mathcal O_E$. As $P^Z_x=P_x\circ\varphi_Z$ and $\varphi_Z(L_Z)=\Lambda_Z$, the projection of
$\mathbf C^g$ onto $\mathbf C^g/P^Z_x(L_Z)$ followed by this isomorphism is a homomorphism $p_x$ as stated. The
other data are those of the compatibility conditions at frames whose minors have absolute value $1$
(Lemma~\ref{cv-uniformization-frames}(3)).

(2) Let $t_Z\colon L\to L_Z$ be the map of Theorem~\ref{cv-uniformization-family} that the compatibility
conditions attach to $Z$, and put $m_Z=\varphi_Z\circ t_Z$, an $\mathcal O_E$-linear map $L\to\Lambda_Z$ that
depends only on $Z$. By that theorem, the level structure of the point $y_Z(P^Z_x)$ of $y$ sends
$a\in\bar{\mathfrak M}'L$ to the image of $P^Z_x(t_Za)/M'$ under the isomorphism
$\mathbf C^g/P^Z_x(L_Z)\to A_y$, that is, to $p_x\big(P^Z_x(t_Za)/M'\big)$. The coordinates of $P_x(w)$ are
$\mathbf Q$-linear in $w$, and $P^Z_x=P_x\circ\varphi_Z$, so
$P^Z_x(t_Za)/M'=P_x\big(M'^{-1}\varphi_Z(t_Za)\big)=P_x\big(M'^{-1}m_Z(a)\big)$, which is the first assertion.

For the second, $\bar{\mathfrak M}'\Lambda_Z$ is generated as an additive group by the products $sw$ with
$s\in\bar{\mathfrak M}'$ and $w\in\Lambda_Z$, and the vectors $m_Z(a)+M'u$ with $a\in\bar{\mathfrak M}'L$ and
$u\in\Lambda_Z$ are closed under addition; so it suffices to treat $v=sw$. Write $w=\varphi_Z(b)$ with $b\in L_Z$.
As $t_Z$ induces a surjection $L/M'L\to L_Z/M'L_Z$, $b=t_Z(a)+M'e$ with $a\in L$ and $e\in L_Z$. Then
$sw=m_Z(sa)+M'\varphi_Z(se)$, with $sa\in\bar{\mathfrak M}'L$ and $\varphi_Z(se)\in\Lambda_Z$.

(3) As $\mathfrak M'+\bar{\mathfrak M}'=\mathcal O_E$, there is $c\in\bar{\mathfrak M}'$ with
$c-1\in\mathfrak M'$. There are also $t\in\mathcal O_E$ and $j\ge0$ with $t-1\in\mathfrak M'$ and
$t\,\mathfrak M'^j\subset2M'\mathcal O_E$. Indeed, in the finite ring $R=\mathcal O_E/2M'\mathcal O_E$ the
decreasing chain of the ideals $\mathfrak M'^iR$ is constant from some $i=j$ on, so $N=\mathfrak M'^jR$ satisfies
$N=\mathfrak M'N$, and there is $t$ with $t-1\in\mathfrak M'$ and $tN=0$ \cite[Corollary~2.5]{AM69}.

First, $s\,m_Z(a)\in M'\Lambda_Z$ for $s\in\mathfrak M'$ and $a\in\bar{\mathfrak M}'L$: as $m_Z$ is
$\mathcal O_E$-linear, $m_Z(a)$ lies in $\bar{\mathfrak M}'\Lambda_Z$, and $\mathfrak M'\bar{\mathfrak M}'=M'\mathcal O_E$.

Let $q$ be a class in $L/\mathfrak M'L$ and $b\in L$ a vector in it, and put $a=cb$ and $Y=tM'^{-1}m_Z(a)$. Then
$a\in\bar{\mathfrak M}'L$, and $a-b=(c-1)b\in\mathfrak M'L$, so $a$ lies in the class $q$. Next,
$Y-M'^{-1}m_Z(a)=M'^{-1}(t-1)m_Z(a)$ lies in $\Lambda_Z$ by the first remark, as $t-1\in\mathfrak M'$. Finally,
let $s\in\mathfrak M'^{j+1}$, a sum of products $s_1s_2$ with $s_1\in\mathfrak M'^j$ and $s_2\in\mathfrak M'$.
For such a product, $ts_1=2M'w$ with $w\in\mathcal O_E$, and $s_2M'^{-1}m_Z(a)=u$ lies in $\Lambda_Z$ by the
first remark, so $s_1s_2Y=(ts_1)\,(s_2M'^{-1}m_Z(a))=2M'wu\in2M'\Lambda_Z$. Summing, $sY\in2M'\Lambda_Z$, and
$j+1$ is the required integer.
% Lean: the setting is Formalism.L2.PjFree.IsUnf U with the uniformization witness UW U of Dict.nonempty_uw (its
% fields F, ik, phi, isom and pt are the lattices L_Z, phi_Z and the data at frames). (1) is the first part of
% uw_anFib: W.pt k x gives the period Pm = P_x o phi_Z, the point, the derivatives and, at unit frames, Hm = H_x,
% chil, fr = F^x, theta and Fd with mean 1; Dict.exists_ex_transport gives ex = p_x (ex_add, ex_onto, ex_ker,
% ex_lin, der), packaged as an AnFib. (2) is A.Marked (W.mark k) from uw_anFib: m_Z = W.mark k = emb k o F.t (ik k)
% (UW.mark), the level structure is the field ex_eta of the FamilyUnif (through exists_ex_transport), and hM
% rewrites M^-1 Pm(t a) as pv(M^-1 mark a) (hsm, hrow, hPv); the second assertion is IsMark.surj of uw_lattice,
% proved by UWLat.isMark_of from W.F.t_surj by induction on the products s w. (3) is the last clause of
% DualMark.dual_marking with N = M', which is DualMark.exists_roots and uses no other hypothesis: c = cSep l
% (LevelPairing.cSep_spec), t and j from DualMark.exists_pow_le (IsArtinian.monotone_stabilizes and
% Submodule.exists_sub_one_mem_and_smul_eq_zero_of_fg_of_le_smul), the first remark is key (smul_eq_nsmul with
% prod_le), the representative b is levRep q, and the integer is j + 1.
\end{proof}
\leanref{Development/AbelianVarietyTheory/Proofs/SiegelThetaTransformation/UniformisationComponentLatticesMarkings.lean}{GlobalTheory.Dict.UWLat.isMark_of}
\leanref{Development/AbelianVarietyTheory/Proofs/SiegelThetaTransformation/DualHalfLevelMarkingExistsAndMatch.lean}{GlobalTheory.Dict.DualMark.exists_roots}

\begin{lemma}\label{cv-uniformization-theta-half-basis}
Let $g,N\ge1$, let $Z\in\mathfrak H_g$, and let $\vartheta[c]$ be as in
Theorem~\ref{cv-uniformization-theta-transformation}. Write $e_1,\dots,e_g$ for the unit vectors of $\mathbf C^g$
and $\mathbf e(w)=\exp(2\pi iw)$, and for $c\in\{0,\dots,N-1\}^g$ put $\vartheta_c(z)=\vartheta[(c/N,0)](NZ,Nz)$.
\begin{enumerate}
\item Each $\vartheta_c$ is holomorphic on $\mathbf C^g$, and
  $\vartheta_c(z+Zn+n')=\mathbf e\big(-N(\tfrac12n^{\mathsf T}Zn+n^{\mathsf T}z)\big)\vartheta_c(z)$ for all
  $z\in\mathbf C^g$ and $n,n'\in\mathbf Z^g$.
\item Let $F$ be a holomorphic function on $\mathbf C^g$ with $F(z+e_i)=F(z)$ and
  $F(z+Ze_j)=\mathbf e\big(-\tfrac N2Z_{jj}-Nz_j\big)F(z)$ for all $z$, $i$ and $j$. Then $F$ is a linear
  combination of the functions $\vartheta_c$.
\end{enumerate}
\end{lemma}
\leanref{Development/AbelianVarietyTheory/Proofs/ToricGeometry/ClassicalThetaLevelAndLinearIndependent.lean}{GlobalTheory.classicalTheta_level_and_linearIndependent}
\leanref{Development/AbelianVarietyTheory/Proofs/DivisorProofs/ExistsThetaDivisionPointNeZero.lean}{GlobalTheory.ThetaDivisionPoints.mem_span_of_laws}

The $\vartheta_c$ are the classical theta functions of level $N$, and the $N^g$ of them form a basis of the
space of the functions $F$ of (2) \cite[Chapter~II, \S1]{Mum83}.

\begin{lemma}\label{cv-uniformization-theta-half-point}
Let $g\ge1$ and $N\ge2$, let $Z\in\mathfrak H_g$, and let $\vartheta_c$, for $c\in\{0,\dots,N-1\}^g$, be as in
Lemma~\ref{cv-uniformization-theta-half-basis}. Then for every $z_0\in\mathbf C^g$ there is $c$ with
$\vartheta_c(z_0)\ne0$.
\end{lemma}
\leanref{Development/AbelianVarietyTheory/Proofs/DivisorProofs/ExistsThetaDivisionPointNeZero.lean}{GlobalTheory.ThetaDivisionPoints.exists_theta_ne_zero}
\begin{proof}
Write $\vartheta(z)=\vartheta[0](Z,z)$ and $\mathbf e(w)=\exp(2\pi iw)$. This is the function $\vartheta_0$ of
Lemma~\ref{cv-uniformization-theta-half-basis} for the level $1$, so by (1) of that lemma $\vartheta$ is holomorphic
and
\[
 \vartheta(z+Zn+n')=\mathbf e\big(-\tfrac12n^{\mathsf T}Zn-n^{\mathsf T}z\big)\vartheta(z)\qquad(n,n'\in\mathbf Z^g).
\]
It is not identically $0$: by Lemma~\ref{cv-volume-key-normalized-theta} for the characteristic $0$ and $z_0=0$,
$\vartheta(t)\ne0$ for some $t\in[0,1]^g$. The functions $u\mapsto\vartheta(z_0+u)$ and $u\mapsto\vartheta(z_0-(N-1)u)$ are holomorphic and not identically
$0$, so the sets where they do not vanish are open and dense, and some $u$ lies in both. Put
$F(z)=\vartheta(z+u)^{N-1}\vartheta(z-(N-1)u)$. It is holomorphic, $F(z_0)\ne0$, and by the display
\[
 F(z+Zn+n')=\mathbf e\big(-(N-1)(\tfrac12n^{\mathsf T}Zn+n^{\mathsf T}(z+u))-\tfrac12n^{\mathsf T}Zn
 -n^{\mathsf T}(z-(N-1)u)\big)F(z)=\mathbf e\big(-N(\tfrac12n^{\mathsf T}Zn+n^{\mathsf T}z)\big)F(z).
\]
With $n=0$ and $n'=e_i$, and with $n=e_j$ and $n'=0$, this gives the laws of
Lemma~\ref{cv-uniformization-theta-half-basis}(2), so $F$ is a linear combination of the $\vartheta_c$. As
$F(z_0)\ne0$, some $\vartheta_c(z_0)$ is not $0$.
\end{proof}
\leanref{Development/AbelianVarietyTheory/Proofs/ToricGeometry/ExistsApplySubSmulNeZeroAndApplyAddNeZero.lean}{GlobalTheory.exists_level_apply_ne_zero_of_two_le}
\leanref{ReciprocityTheory/ThetaQuasiPeriods/Transl.lean}{ReciprocityTheory.PTheta.thetaQuasiPeriod_holds}
\leanref{Development/AbelianVarietyTheory/Proofs/ToricGeometry/RiemannThetaLevelOneNeZero.lean}{GlobalTheory.riemannTheta_level_one_ne_zero}
\leanref{Development/AbelianVarietyTheory/Proofs/CovolumeBounds/EqZeroOfSemicharacterMultiplier.lean}{GlobalTheory.theta_ne_zero}
% Lean: exists_theta_ne_zero: F is exists_level_apply_ne_zero_of_two_le (u from
% dense_setOf_ne_zero_of_differentiable_ne_zero for the two functions, the law from hq' and eC_nat_mul), whose
% hypotheses on the Riemann theta function are riemannTheta_level_one_ne_zero: the holomorphy and the law are
% classicalTheta_level_and_linearIndependent' at N = 1 (Lemma cv-uniformization-theta-half-basis(1)), and the
% nonvanishing is GlobalTheory.theta_ne_zero = IsSiegel.riemannTheta_ne_zero, which is
% IsSiegel.exists_riemannTheta_add_ofReal_ne_zero at v = 0 (Lemma cv-volume-key-normalized-theta for z_0 = 0); then
% mem_span_of_laws with laws_of_level.

\begin{lemma}\label{cv-uniformization-theta-half-classical}
Let $g\ge1$ and $N\ge2$, let $Z\in\mathfrak H_g$ (Theorem~\ref{cv-uniformization-theta-transformation}), let
$\beta_1,\dots,\beta_g$ be nonzero complex numbers, and write $e_1,\dots,e_g$ for the unit vectors of
$\mathbf C^g$. Let $f$ be a holomorphic function on $\mathbf C^g$, not identically $0$, with $f(z+e_i)=f(z)$ and
$f(z+Ze_j)=\beta_j\exp(-2\pi iNz_j)f(z)$ for all $z$, $i$ and $j$. Then $f\big(N^{-1}(m+Zn)\big)\ne0$ for some
$m,n\in\mathbf Z^g$.
\end{lemma}
\leanref{Development/AbelianVarietyTheory/Proofs/DivisorProofs/ExistsThetaDivisionPointNeZeroStatement.lean}{GlobalTheory.T.exists_theta_division_point_ne_zero}
\begin{proof}
Write $\mathbf e(w)=\exp(2\pi iw)$, $C=\{0,\dots,N-1\}^g$, and let $\vartheta_c$, $c\in C$, be as in
Lemma~\ref{cv-uniformization-theta-half-basis}.

Step 1: normalization. Choose $w\in\mathbf C^g$ with $\beta_j\mathbf e(-Nw_j)=\mathbf e(-\frac N2Z_{jj})$ for every
$j$, and put $F(z)=f(z+w)$. Then $F$ is holomorphic and not identically $0$, $F(z+e_i)=F(z)$ and
$F(z+Ze_j)=\mathbf e(-\frac N2Z_{jj}-Nz_j)F(z)$. By Lemma~\ref{cv-uniformization-theta-half-basis}(2),
$F=\sum_{c\in C}a_c\vartheta_c$ for some numbers $a_c$, and $a_{c_0}\ne0$ for some $c_0\in C$.

Step 2: two translations. For $a,b\in\mathbf R^g$ and $Z'\in\mathfrak H_g$, a computation with the series gives
\[
 \vartheta[(a,0)](Z',z+m)=\mathbf e(a^{\mathsf T}m)\,\vartheta[(a,0)](Z',z)\ \ (m\in\mathbf Z^g),\qquad
 \vartheta[(a,0)](Z',z+Z'b)=\mathbf e\big(-\tfrac12b^{\mathsf T}Z'b-b^{\mathsf T}z\big)\,\vartheta[(a+b,0)](Z',z).
\]
With $Z'=NZ$ they give, for $c\in C$ and $m,n\in\mathbf Z^g$ with $c_0+n\in C$,
\[
 \vartheta_c(z+m/N)=\mathbf e(c^{\mathsf T}m/N)\,\vartheta_c(z),\qquad
 \vartheta_{c_0}(z+Zn/N)=\mathbf e\big(-\tfrac1{2N}n^{\mathsf T}Zn-n^{\mathsf T}z\big)\,\vartheta_{c_0+n}(z).
\]

Step 3: suppose that $f$ vanishes at all the points $N^{-1}(m+Zn)$. Fix $n\in\mathbf Z^g$ and put
$z_n=-w+Zn/N$. For every $m\in\mathbf Z^g$, $z_n+m/N+w=N^{-1}(m+Zn)$, so by Steps 1 and 2
\[
 0=F(z_n+m/N)=\sum_{c\in C}a_c\vartheta_c(z_n)\,\mathbf e(c^{\mathsf T}m/N).
\]
The maps $m\mapsto\mathbf e(c^{\mathsf T}m/N)$, $c\in C$, are distinct homomorphisms $\mathbf Z^g\to\mathbf C^\times$, as
their values $\mathbf e(c_i/N)$ at the $e_i$ determine $c$. By Artin's theorem on the linear independence of
distinct characters \cite[Theorem VI.4.1]{Lan02}, $a_c\vartheta_c(z_n)=0$ for every $c$. In particular
$\vartheta_{c_0}(z_n)=0$ for every $n\in\mathbf Z^g$.

Step 4: the contradiction. By Lemma~\ref{cv-uniformization-theta-half-point} there is $c\in C$ with
$\vartheta_c(-w)\ne0$. For $n=c-c_0$, Step 2 gives
$\vartheta_{c_0}(z_n)=\vartheta_{c_0}(-w+Zn/N)=\mathbf e\big(-\tfrac1{2N}n^{\mathsf T}Zn+n^{\mathsf T}w\big)\vartheta_c(-w)\ne0$,
which contradicts Step 3.
\end{proof}
\leanref{Development/AbelianVarietyTheory/Proofs/DivisorProofs/ExistsThetaDivisionPointNeZero.lean}{GlobalTheory.ThetaDivisionPoints.exists_division_point_ne_zero}
\leanref{Development/AbelianVarietyTheory/Proofs/DivisorProofs/ExistsThetaDivisionPointNeZero.lean}{GlobalTheory.ThetaDivisionPoints.theta_add_int_div}
\leanref{Development/AbelianVarietyTheory/Proofs/DivisorProofs/ExistsThetaDivisionPointNeZero.lean}{GlobalTheory.ThetaDivisionPoints.theta_add_tau_div}
\leanref{ReciprocityTheory/ThetaQuasiPeriods/Transl.lean}{ReciprocityTheory.PTheta.thetaRealTransl_holds}
% Lean: exists_division_point_ne_zero: w and hw for Step 1, mem_span_of_laws and c0 with a c0 /= 0; Step 2 is
% theta_add_int_div (thetaQuasiPeriod_holds) and theta_add_tau_div (thetaRealTransl_holds); Step 3 is hvan, with
% intChar, intChar_injective and Mathlib's linearIndependent_monoidHom; Step 4 is exists_theta_ne_zero at -w and
% theta_add_tau_div with n = c - c0.

\begin{lemma}\label{cv-uniformization-theta-half-normal-form}
Let $g\ge1$, let $\Lambda\subset\mathbf C^g$ be a lattice, a discrete subgroup that spans $\mathbf C^g$ over
$\mathbf R$, and let $H$ be a positive definite Hermitian form on $\mathbf C^g$, linear in the first variable, such
that $E=\operatorname{Im}H$ takes integer values on $\Lambda\times\Lambda$ and $l\mapsto E(l,\cdot)$ is a bijection
of $\Lambda$ onto $\operatorname{Hom}(\Lambda,\mathbf Z)$. Let $\chi\colon\Lambda\to\mathbf C$ be a function with
$|\chi(l)|=1$ for every $l$, put $a(l,z)=\chi(l)\exp\big(\pi H(z,l)+\frac\pi2H(l,l)\big)$, and write
$e_1,\dots,e_g$ for the unit vectors of $\mathbf C^g$. Then there are a $\mathbf C$-linear automorphism $\varphi$ of
$\mathbf C^g$, a matrix $Z\in\mathfrak H_g$ (Theorem~\ref{cv-uniformization-theta-transformation}), a holomorphic
function $Q$ on $\mathbf C^g$ with $Q(0)=0$, and nonzero numbers $b_1,\dots,b_g$, such that
$\Lambda=\varphi(\mathbf Z^g+Z\mathbf Z^g)$ and, for all $z\in\mathbf C^g$ and all $i$ and $j$,
\[
 a\big(\varphi(e_i),\varphi(z)\big)=\exp\big(Q(z+e_i)-Q(z)\big),\qquad
 a\big(\varphi(Ze_j),\varphi(z)\big)=b_j\exp(-2\pi iz_j)\exp\big(Q(z+Ze_j)-Q(z)\big).
\]
\end{lemma}
\leanref{Development/AbelianVarietyTheory/Proofs/ToricGeometry/IsUnitOfLagrangian.lean}{GlobalTheory.SiegelNormalForm.exists_normal_form}

This is the normal form of a principally polarized complex torus in a symplectic basis \cite[Chapter~8]{BL04}:
$\varphi$ sends the unit vectors to half of a symplectic basis of $\Lambda$ for $E$, $Z$ is the normalized period
matrix, and $Q$ is a quadratic form plus a linear form, which makes the factor $a$ equivalent to the classical factor
of automorphy on these generators of $\Lambda$.

\begin{lemma}\label{cv-uniformization-theta-half-principal}
Let $g\ge1$, let $\Lambda\subset\mathbf C^g$ be a lattice, a discrete subgroup that spans $\mathbf C^g$ over
$\mathbf R$, and let $H$ be a positive definite Hermitian form on $\mathbf C^g$, linear in the first variable, such
that $E=\operatorname{Im}H$ takes integer values on
$\Lambda\times\Lambda$ and $l\mapsto E(l,\cdot)$ is a bijection of $\Lambda$ onto $\operatorname{Hom}(\Lambda,\mathbf Z)$.
Let $N\ge2$, let $c\colon\Lambda\to\mathbf C$ be a function with $|c(l)|=1$ for every $l$, and let $f$ be a
holomorphic function on $\mathbf C^g$, not identically $0$, with
\[
 f(z+l)=c(l)\exp\big(N\big(\pi H(z,l)+\tfrac\pi2H(l,l)\big)\big)f(z)\qquad(z\in\mathbf C^g,\ l\in\Lambda).
\]
Then $f(l/N)\ne0$ for some $l\in\Lambda$.
\end{lemma}
\leanref{Development/AbelianVarietyTheory/Proofs/DivisorProofs/ThetaFunctionNeZeroAtDivisionPointSelfDualStatement.lean}{GlobalTheory.T.thetaFunction_ne_zero_at_divisionPoint_selfDual}
\begin{proof}
For each $l\in\Lambda$ choose $\chi(l)\in\mathbf C$ with $\chi(l)^N=c(l)$; then $|\chi(l)|=1$. Let $\varphi$, $Z$,
$Q$ and $b_j$ be as in Lemma~\ref{cv-uniformization-theta-half-normal-form} for $\chi$, and put
$F(z)=f(\varphi(z))\exp(-NQ(z))$, a holomorphic function, not identically $0$ as $\varphi$ is onto. For
$l=\varphi(e_i)$ the hypothesis and that lemma give
$f(\varphi(z)+l)=a(l,\varphi(z))^Nf(\varphi(z))=\exp\big(N(Q(z+e_i)-Q(z))\big)f(\varphi(z))$, so $F(z+e_i)=F(z)$.
In the same way $F(z+Ze_j)=b_j^N\exp(-2\pi iNz_j)F(z)$. By Lemma~\ref{cv-uniformization-theta-half-classical} with
$\beta_j=b_j^N$, $F\big(N^{-1}(m+Zn)\big)\ne0$ for some $m,n\in\mathbf Z^g$. Then $l=\varphi(m+Zn)$ lies in
$\Lambda$, and $f(l/N)=F\big(N^{-1}(m+Zn)\big)\exp\big(NQ(N^{-1}(m+Zn))\big)\ne0$.
\end{proof}
\leanref{Development/AbelianVarietyTheory/Proofs/DivisorProofs/ThetaFunctionNeZeroAtDivisionPointSelfDual.lean}{GlobalTheory.ThetaDivisionPoints.exists_ne_zero_of_principal}
\leanref{Development/AbelianVarietyTheory/Proofs/DivisorProofs/ThetaFunctionNeZeroAtDivisionPointSelfDual.lean}{GlobalTheory.ThetaDivisionPoints.exists_ne_zero_of_selfDual}
% Lean: exists_ne_zero_of_selfDual takes the lattice as the image p(L) of a self-dual Z-lattice L in a Q-vector
% space V and builds Ef (integral, bijective by Module.Baer and the nondegeneracy of B) and the N-th roots chi
% (IsAlgClosed.exists_pow_nat_eq); exists_ne_zero_of_principal is the rest: exists_normal_form, F = f o phi times
% exp(-n Q), exists_theta_division_point_ne_zero, and map_smul for f(phi(z/N)) = f(phi(z)/N).

\begin{lemma}\label{cv-uniformization-theta-half-sign-liouville}
Let $V$ be a finite-dimensional normed complex vector space, let $S\subset V$ be a set and $R$ a real number such
that every $z\in V$ has some $l\in S$ with $\|z+l\|\le R$. Let $\psi$ be a holomorphic function on $V$ without
zeros, and let $D\colon V\times V\to\mathbf C$ be a function such that $z\mapsto D(z,z)$ is continuous and
\[
 D(z+z',w)=D(z,w)+D(z',w),\qquad D(tz,w)=t\,D(z,w),\qquad D(w,z)=\overline{D(z,w)}
\]
for all $z,z',w\in V$ and $t\in\mathbf C$. Suppose that for every $l\in S$ there is $\rho\in\mathbf C$ with
$|\rho|=1$ and
\[
 \psi(z+l)=\rho\exp\big(\pi D(z,l)+\tfrac\pi2D(l,l)\big)\psi(z)\qquad(z\in V).
\]
Then $D(z,w)=0$ for all $z,w\in V$, and $\psi$ is constant.
\end{lemma}
\leanref{Development/AbelianVarietyTheory/Proofs/CuspFormEvaluations/TripleProductTorusPeriodicBounded.lean}{GlobalTheory.FrameUnits.form_eq_zero_and_const}
\begin{proof}
Step 1: periodic functions are bounded. Let $\phi$ be a continuous real function on $V$ with $\phi(z+l)=\phi(z)$
for all $z\in V$ and $l\in S$. For $z\in V$ choose $l\in S$ with $\|z+l\|\le R$. Then $\phi(z)=\phi(z+l)$ is at most
the maximum of $\phi$ on the closed ball of radius $R$ about $0$, which is compact. So $\phi$ is bounded above.

Step 2: the weight. As $D(w,z)=\overline{D(z,w)}$, $\operatorname{Re}D(w,z)=\operatorname{Re}D(z,w)$, and additivity
in the first variable gives
\[
 \operatorname{Re}D(z+l,z+l)=\operatorname{Re}D(z,z)+2\operatorname{Re}D(z,l)+\operatorname{Re}D(l,l).
\]
Put $q(z)=\frac\pi2\operatorname{Re}D(z,z)$, a continuous function. For $t\in\mathbf C$,
$D(tw,tw)=t\,D(w,tw)=t\,\overline{D(tw,w)}=t\bar t\,\overline{D(w,w)}=|t|^2D(w,w)$, as $D(w,w)$ is real; so
$q(tw)=|t|^2q(w)$. Taking absolute values in the law of $\psi$ gives
$|\psi(z+l)|=\exp\big(q(z+l)-q(z)\big)|\psi(z)|$ for $l\in S$.

Step 3: $q\ge0$. Suppose that $q(v)<0$ for some $v\in V$. By Step 2 the continuous function $|\psi(z)|e^{-q(z)}$ is
invariant under $S$, so by Step 1 it is at most some $C$. The entire function $f(t)=\psi(tv)$ of $t\in\mathbf C$
then satisfies $|f(t)|\le Ce^{q(tv)}=Ce^{|t|^2q(v)}$. So $f$ is bounded, hence constant by Liouville's theorem,
and $|f(0)|=|f(n)|\le Ce^{n^2q(v)}$ for every integer $n\ge1$. The right side tends to $0$, so
$\psi(0)=f(0)=0$, a contradiction.

Step 4: $q\le0$. The function $1/\psi$ is holomorphic without zeros, $-q(tw)=-|t|^2q(w)$, and
$|1/\psi(z+l)|e^{q(z+l)}=|1/\psi(z)|e^{q(z)}$ for $l\in S$ by Step 2. The argument of Step 3 for $1/\psi$ and $-q$
gives $-q\ge0$.

Step 5: the conclusion. By Steps 3 and 4, $\operatorname{Re}D(z,z)=0$ for every $z$, and the identity of Step 2
gives $\operatorname{Re}D(z,w)=0$ for all $z,w$. As $D(iz,w)=iD(z,w)$,
$\operatorname{Im}D(z,w)=-\operatorname{Re}D(iz,w)=0$. So $D=0$. Then $q=0$ and $|\psi(z+l)|=|\psi(z)|$ for
$l\in S$, so $|\psi|$ is bounded by Step 1, and $\psi$ is constant by Liouville's theorem.
\end{proof}
\leanref{Development/AbelianVarietyTheory/Proofs/CuspFormEvaluations/TripleProductTorusPeriodicBounded.lean}{GlobalTheory.FrameUnits.exists_le_of_periodic}
\leanref{Development/AbelianVarietyTheory/Proofs/CuspFormEvaluations/TripleProductTorusPeriodicBounded.lean}{GlobalTheory.FrameUnits.weight_nonneg_of_periodic}
\leanref{Development/AbelianVarietyTheory/Proofs/CubicalLineBundles/AppellHumbertTrivialOfNowhereVanishing.lean}{GlobalTheory.AppellHumbertTrivial.eq_zero_of_gaussian_decay}
% Lean: form_eq_zero_and_const (59 lines): hre, hsq, hs and hn are Step 2; h1 and h2 are weight_nonneg_of_periodic
% (17 lines) for psi with q and for 1/psi with -q (Steps 3 and 4), which takes C from exists_le_of_periodic (9 lines,
% Step 1) and applies AppellHumbertTrivial.eq_zero_of_gaussian_decay (25 lines: f bounded by max(C, 0),
% Differentiable.apply_eq_apply_of_bounded, and the limit of C exp(n^2 a)) to t -> psi(t v); h3, h4 and the value at
% I z give D = 0; exists_le_of_periodic for |psi| and apply_eq_apply_of_bounded give the constancy.

\begin{lemma}\label{cv-uniformization-theta-half-sign-glue}
Let $Y$ be a scheme and $\mathcal L$ a line bundle on $Y$ with trivializations over the sets $U_j$ of an open cover
of $Y$: the values of a section in the trivializations over $U_j$ and $U_{j'}$ satisfy $s_j=g_{jj'}s_{j'}$ for
transition functions $g_{jj'}$. Let $g\ge1$, and let $p$ be a map from $\mathbf C^g$ to the complex points of $Y$
that is holomorphic: for every open set $U\subset Y$ and every regular function $f$ on $U$, the set $O_U$ of the $z$
with $p(z)\in U$ is open, and the value $f(p(z))$ is a holomorphic function of $z\in O_U$. Write $O_j=O_{U_j}$. Let
$K$ be a finite set, and for $k\in K$ let $\phi_k\colon Y\to Y$ be a morphism, $\varepsilon_k$ an integer and
$\lambda_k\colon\mathbf C^g\to\mathbf C^g$ a holomorphic map with $\phi_k(p(z))=p(\lambda_k(z))$ for every $z$. For
each $j$ let $F_j$ be a holomorphic function on $O_j$ without zeros, with $F_j(z)=g_{jj'}(p(z))F_{j'}(z)$ for
$z\in O_j\cap O_{j'}$. Suppose that $\mathcal L$ is isomorphic to
$\bigotimes_{k\in K}\phi_k^*\mathcal L^{\varepsilon_k}$. Then there is a holomorphic function $\psi$ on
$\mathbf C^g$ without zeros with the following property. Let $z,z'\in\mathbf C^g$ with $p(z')=p(z)$, and let $a$
and $a_k$, $k\in K$, be complex numbers such that $F_j(z)=aF_j(z')$ for every $j$ with $z\in O_j$, and
$F_j(\lambda_k(z))=a_kF_j(\lambda_k(z'))$ for every $k$ and every $j$ with $\lambda_k(z)\in O_j$. Then
\[
 \psi(z')\,a=\Big(\prod_{k\in K}a_k^{\varepsilon_k}\Big)\psi(z).
\]
\end{lemma}
\leanref{Development/AbelianVarietyTheory/Proofs/SiegelThetaTransformation/AnalyticFibrePrincipalBundleSignEq.lean}{GlobalTheory.Dict.FrameSign.exists_unit}
\begin{proof}
For a map $\kappa$ from $K$ to the indices of the cover put $W_\kappa=\bigcap_k\phi_k^{-1}(U_{\kappa(k)})$. The line
bundle $\mathcal M=\bigotimes_k\phi_k^*\mathcal L^{\varepsilon_k}$ has trivializations over the sets $W_\kappa$,
with the transition functions $\prod_k(g_{\kappa(k)\kappa'(k)}\circ\phi_k)^{\varepsilon_k}$. So an isomorphism of
$\mathcal L$ with $\mathcal M$ is a family of nowhere vanishing regular functions $t_{\kappa j}$ on
$W_\kappa\cap U_j$, by which it multiplies the values of sections of $\mathcal L$ over $U_j$, with
\[
 t_{\kappa j}=\prod_k(g_{\kappa(k)\kappa'(k)}\circ\phi_k)^{\varepsilon_k}\,g_{jj'}^{-1}\,t_{\kappa'j'}
 \quad\text{on }W_\kappa\cap U_j\cap W_{\kappa'}\cap U_{j'}.
\]

For $\kappa$ and $j$ let $D_{\kappa j}$ be the set of the $z$ with $p(z)\in W_\kappa\cap U_j$, an open set. As
$\phi_k(p(z))=p(\lambda_k(z))$, $z\in D_{\kappa j}$ if and only if $z\in O_j$ and $\lambda_k(z)\in O_{\kappa(k)}$ for
every $k$; and the sets $D_{\kappa j}$ cover $\mathbf C^g$, as the $U_j$ cover $Y$. On $D_{\kappa j}$ put
\[
 G_{\kappa j}(z)=t_{\kappa j}(p(z))\,F_j(z)\prod_{k\in K}F_{\kappa(k)}(\lambda_k(z))^{-\varepsilon_k},
\]
a holomorphic function without zeros. Substituting the transition law of $t$,
$F_j(z)=g_{jj'}(p(z))F_{j'}(z)$ and $F_{\kappa(k)}(\lambda_kz)=g_{\kappa(k)\kappa'(k)}(p(\lambda_kz))
F_{\kappa'(k)}(\lambda_kz)$ in $G_{\kappa j}$ gives $G_{\kappa'j'}$ on $D_{\kappa j}\cap D_{\kappa'j'}$. So the
$G_{\kappa j}$ glue to a holomorphic function $\psi$ on $\mathbf C^g$ without zeros. It is the ratio of two nowhere
vanishing sections of the pullback of $\mathcal M$ to $\mathbf C^g$: the image under the isomorphism of the section
of the pullback of $\mathcal L$ with the values $F_j$, and the product of the pullbacks of that section under the
$\lambda_k$ with the exponents $\varepsilon_k$.

Let $z$, $z'$, $a$ and the $a_k$ be as in the statement, and choose $\kappa$ and $j$ with $z\in D_{\kappa j}$. As
$p(z')=p(z)$ and $p(\lambda_k(z'))=\phi_k(p(z'))=p(\lambda_k(z))$, also $z'\in D_{\kappa j}$, and the factor
$t_{\kappa j}$ has the same value at $p(z)$ and at $p(z')$. The numbers $a$ and $a_k$ are not $0$, as the $F_j$
have no zeros, so $F_j(z')=a^{-1}F_j(z)$ and $F_{\kappa(k)}(\lambda_k(z'))^{-\varepsilon_k}=a_k^{\varepsilon_k}
F_{\kappa(k)}(\lambda_k(z))^{-\varepsilon_k}$. Hence $\psi(z')=a^{-1}\prod_ka_k^{\varepsilon_k}\psi(z)$.
\end{proof}
\leanref{Development/AbelianVarietyTheory/Proofs/ToricGeometry/AppellHumbertDataOfIsoCube.lean}{GlobalTheory.SelfCube.exists_glue}
% Lean: exists_unit (92 lines): the isomorphism is a generator t of LB.div (LinComb L phi eps) L (LBIso), whose
% index set is pairs (kappa, j) and whose transition functions are those of LinComb divided by those of L; D and hD
% are the sets D_{kappa j} (le_preimage_inf, le_preimage_wchart); T and G are the functions G_{kappa j}; they glue by
% SelfCube.exists_glue (holomorphy from IsHolTo and DifferentiableOn.zpow), agreeing on overlaps by val_gen, val_div
% and val_linComb (values at a complex point of the transition functions of the generator, of LB.div and of LinComb,
% from ValueRules) and fr_g (c0, c1); nonvanishing by val_unit_ne_zero; the law by eT and c1.

\begin{lemma}\label{cv-uniformization-theta-half-sign-fibre}
Let a PEL datum over a CM field $E$ with prime $d$ be given, as in Situation~\ref{cv-function-situation}, and
put $w_0=(1+\sqrt{-d})/2\in\mathcal O_E$. Let $X$ be data (1) to (6) of flag data of rank $r$ over $E$
(Definition~\ref{is-onto-share-flag-data}), with set $\Psi$, a symplectic basis of its module and an indexing
(Definition~\ref{is-eigenvalue-conjugate-frame-fibres}(1) and~(2)), and write $g=|\Psi|r$ for the integer written
$d$ there. Let $y$ be a complex framed point of the moduli scheme of the datum at a level, $j\mapsto\ell_j$ a
bijection from $\{1,\dots,g\}$ onto the letters of the datum, $x\in\mathcal D_X$, $L\subset E^r$ an
$\mathcal O_E$-submodule that is self-dual for $X$ (Definition~\ref{is-eigenvalue-conjugate-frame-fibres}(7)), and
$(\operatorname{ex},\operatorname{fr},\chi)$ an analytic fibre at $y$ for these data, $x$ and $L$
(Definition~\ref{is-eigenvalue-conjugate-frame-fibres}(3)). Then $\chi(v)=\exp(\pi i\langle v,w_0v\rangle_X)$ for
every $v\in L$.
\end{lemma}
\leanref{Development/AbelianVarietyTheory/Proofs/SiegelThetaTransformation/AnalyticFibrePrincipalBundleSignEq.lean}{GlobalTheory.Dict.analyticFibre_principalBundle_sign_eq}
\begin{proof}
Write $P=p_x$, $H=H_x$, $p=\operatorname{ex}$, $F_\beta=\operatorname{fr}_\beta$,
$\langle\ ,\ \rangle=\langle\ ,\ \rangle_X$ and $Pv$ for $P(v)$, and let $b\in E$ act on $\mathbf C^g$ by
$b\cdot z=(\rho_j(b)z_j)_j$, through the row types $\rho_j$ of the indexing; then $P$ is additive and
$P(bv)=b\cdot Pv$ for $b\in E$ (Definition~\ref{is-eigenvalue-conjugate-frame-fibres}(1)). Put
$\varepsilon(v)=\exp(\pi i\langle v,w_0v\rangle)$ for $v\in E^r$. Let $\mathcal A$ be the semi-abelian scheme of
the universal object over the moduli scheme $S$ of $y$ (Definition~\ref{cv-toroidal-moduli-object}(4) and~(5)), an
$\mathcal O_E$-semi-abelian scheme with the line bundle $\mathcal L$
(Definition~\ref{cv-model-toroidal-semiabelian}(2)), whose fibre over the point under $y$ is $A_y$. The map $p$ is
holomorphic in the sense of Lemma~\ref{cv-uniformization-theta-half-sign-glue}: for an open set
$U\subset\mathcal A$, the points of $A_y(\mathbf C)$ in $U$ form an open subset, and a regular function on $U$
restricts to a holomorphic function on it. Write $O_\beta$ for the open set of the $z$ with $p(z)\in U_\beta$.

Step 0: the form $H$ and the span of $P(L)$. By Lemma~\ref{is-eigenvalue-conjugate-frame-hermitian}(1) and~(2),
whose proof uses only data (1) to (6) of $X$, $H$ is a Hermitian form, linear in the first variable, with
$\operatorname{Im}H(Pu,Pv)=\langle u,v\rangle$ for all $u,v\in E^r$. The module $L$ spans $E^r$ over $E$. Otherwise
some nonzero $E$-linear form $f$ on $E^r$ vanishes on $L$. Write $f(u)=\sum_au_a\bar z_a$ with $z\in E^r$, and let
$w\in E^r$ satisfy $\sum_aw_aA_{ab}=z_b$ for every $b$, for the matrix $A$ and the form $h$ of $X$; $w$ exists as $A$
is invertible. For $c\in E$ and $u\in L$, $h(cw,u)=c\sum_bz_b\bar u_b=c\,\overline{f(u)}=0$, so $\langle cw,u\rangle=0$,
and $cw\in L$ by self-duality. If $w_k\ne0$ for some $k$, the $k$-th coordinates of the vectors $cw$, $c\in E$, are all
of $E$, so $E$ is a quotient of the finitely generated $\mathcal O_E$-module $L$; then $E$ is integral over
$\mathcal O_E$, and $\mathcal O_E$ is a field, which it is not. So $w=0$, hence $z=0$ and $f=0$, a contradiction. By Lemma~\ref{is-eigenvalue-conjugate-frame-hermitian}(4),
the vectors $Pv$, $v\in L$, therefore span $\mathbf C^g$ over $\mathbf R$.

Step 1: for $u,v\in L$, $\chi(u+v)=\chi(u)\chi(v)\exp(\pi i\langle u,v\rangle)$, and $\chi(0)=1$. As the
$U_\beta$ cover $\mathcal A$, there is $\beta$ with $p(0)\in U_\beta$. For $w\in L$, $Pw$ lies in the kernel of $p$,
so $p(Pw)=p(0)$ and $Pw\in O_\beta$. The law (c) of the analytic fibre for $u$ at $0$, for $v$ at $Pu$ and for $u+v$
at $0$ gives, as $P(u+v)=Pu+Pv$ and $H(0,\cdot)=0$,
\begin{align*}
 F_\beta(0)&=\chi(u)e^{\frac\pi2H(Pu,Pu)}F_\beta(Pu),\qquad
 F_\beta(Pu)=\chi(v)e^{\pi H(Pu,Pv)+\frac\pi2H(Pv,Pv)}F_\beta(Pu+Pv),\\
 F_\beta(0)&=\chi(u+v)e^{\frac\pi2H(Pu+Pv,Pu+Pv)}F_\beta(Pu+Pv).
\end{align*}
As $F_\beta(Pu+Pv)\ne0$ and $H(Pu+Pv,Pu+Pv)=H(Pu,Pu)+H(Pu,Pv)+H(Pv,Pu)+H(Pv,Pv)$, this gives
$\chi(u+v)=\chi(u)\chi(v)\exp\big(\frac\pi2(H(Pu,Pv)-H(Pv,Pu))\big)$. As $H$ is Hermitian, by Step 0,
$H(Pu,Pv)-H(Pv,Pu)=2i\operatorname{Im}H(Pu,Pv)=2i\langle u,v\rangle$, which is the formula. For $u=v=0$ it gives
$\chi(0)=\chi(0)^2$, as $\langle0,0\rangle=0$; and $\chi(0)\ne0$, as $|\chi(0)|=1$. So $\chi(0)=1$.

Step 2: the function $\psi$. Let $F$ be the $\mathbf C$-linear map $z\mapsto w_0\cdot z$ of $\mathbf C^g$. Then
$P(w_0v)=F(Pv)$ for $v\in E^r$, $p(F(z))=w_0\,p(z)$, and $w_0L\subset L$ as $w_0\in\mathcal O_E$. For
$J\subset\{1,2\}$ let $\lambda_J$ be the sum over $k\in J$ of $\mu_1=1$ and $\mu_2=F$, so that $\lambda_\emptyset=0$
and $\lambda_{\{1,2\}}=1+F$; let $m_J\colon\mathcal A\to\mathcal A$ be the sum over $k\in J$ of the identity ($k=1$)
and the action of $w_0$ ($k=2$) in the group scheme $\mathcal A$ over $S$, so that $m_\emptyset$ is the zero section
composed with $\mathcal A\to S$; and put $n_J=(-1)^{|J|}$. For the identity $o$ of $\mathcal A$, as a point of
$\mathcal A$ with values in $\mathcal A$,
\[
 \Lambda_{\mathcal L}(o,w_0o)=m_{\{1,2\}}^*\mathcal L\otimes m_{\{1\}}^*\mathcal L^{-1}\otimes
 m_{\{2\}}^*\mathcal L^{-1}\otimes m_\emptyset^*\mathcal L=\bigotimes_Jm_J^*\mathcal L^{n_J},
\]
and $\mathcal L$ is isomorphic to it, as $\mathcal A$ is normalized. As $p$ is a homomorphism with
$p(F(z))=w_0\,p(z)$, $m_J(p(z))=p(\lambda_J(z))$. Lemma~\ref{cv-uniformization-theta-half-sign-glue}, for
$Y=\mathcal A$, the four sets $J$, the morphisms $m_J$, the exponents $n_J$, the maps $\lambda_J$ and the
$F_\beta$, gives a holomorphic function $\psi$ on $\mathbf C^g$ without zeros with the property stated there.

Step 3: the law of $\psi$. Put $D(z,w)=H(z,Fw)+H(Fz,w)-H(z,w)$ for $z,w\in\mathbf C^g$. Then for $v\in L$ and
$z\in\mathbf C^g$,
\[
 \psi(z+Pv)=\frac{\varepsilon(v)}{\chi(v)}\exp\big(\pi D(z,Pv)+\tfrac\pi2D(Pv,Pv)\big)\psi(z).
\]
For $J\subset\{1,2\}$ let $u_J$ be the sum over $k\in J$ of $v_1=v$ and $v_2=w_0v$; then $u_J\in L$ and
$P(u_J)=\lambda_J(Pv)$. Put $z'=z+Pv$. Then $p(z')=p(z)$, as $p(Pv)=0$, and $\lambda_J(z')=\lambda_J(z)+P(u_J)$. So
the laws of the $F_\beta$ for $v$ at $z$ and for $u_J$ at $\lambda_J(z)$ give the hypotheses of the property of $\psi$
with
\[
 a=\chi(v)e^{\pi H(z,Pv)+\frac\pi2H(Pv,Pv)},\qquad
 a_J=\chi(u_J)e^{\pi H(\lambda_Jz,\lambda_JPv)+\frac\pi2H(\lambda_JPv,\lambda_JPv)},
\]
and that property gives $\psi(z+Pv)\,a=a_{\{1,2\}}a_{\{1\}}^{-1}a_{\{2\}}^{-1}a_\emptyset\,\psi(z)$. Here
$a_\emptyset=\chi(0)=1$ by Step 1, $a_{\{1\}}=a$, $a_{\{2\}}=\chi(w_0v)e^{\pi H(Fz,FPv)+\frac\pi2H(FPv,FPv)}$, and
$\chi(v+w_0v)=\chi(v)\chi(w_0v)\varepsilon(v)$ by Step 1. Expanding
$H(z+Fz,Pv+FPv)$ and $H(Pv+FPv,Pv+FPv)$ by additivity in each variable gives
\[
 \frac{a_{\{1,2\}}}{a_{\{1\}}a_{\{2\}}}=\varepsilon(v)\exp\Big(\pi\big(H(z,FPv)+H(Fz,Pv)\big)
 +\frac\pi2\big(H(Pv,FPv)+H(FPv,Pv)\big)\Big),
\]
and dividing by $a$ gives the law.

Step 4: $D=0$ and $\psi$ is constant. Apply Lemma~\ref{cv-uniformization-theta-half-sign-liouville} to
$V=\mathbf C^g$, the set $P(L)$, $\psi$ and $D$. As $P(L)$ spans $\mathbf C^g$ over $\mathbf R$ by Step 0, it contains a basis
$f_1,\dots,f_{2g}$ of $\mathbf C^g$ over $\mathbf R$. For $z=\sum_kc_kf_k$ with real $c_k$, the vector
$l=-\sum_k\lfloor c_k\rfloor f_k$ lies in the group $P(L)$, and $\|z+l\|\le\sum_k\|f_k\|$; so the bound $R$ exists.
The function $z\mapsto D(z,z)$ is continuous; $D$ is additive and $\mathbf C$-linear in the first variable, as $H$ is
and $F$ is $\mathbf C$-linear; and $D(w,z)=\overline{D(z,w)}$, as $H$ is Hermitian. For $l=Pv$ with $v\in L$, Step 3
is the law of that lemma with $\rho=\varepsilon(v)/\chi(v)$, and $|\rho|=1$, as $|\chi(v)|=1$ and $\langle v,w_0v\rangle$
is real. So $D=0$ and $\psi$ is constant.

Step 5: the conclusion. Let $v\in L$. Step 3 at $z=0$, with $D=0$, gives $\psi(Pv)=\varepsilon(v)\chi(v)^{-1}\psi(0)$.
As $\psi(Pv)=\psi(0)\ne0$, $\chi(v)=\varepsilon(v)$.
\end{proof}
\leanref{Development/AbelianVarietyTheory/Proofs/SiegelThetaTransformation/AnalyticFibrePrincipalBundleSignEq.lean}{GlobalTheory.Dict.FrameSign.chi_add}
\leanref{Development/AbelianVarietyTheory/Proofs/SiegelThetaTransformation/AnalyticFibrePrincipalBundleSignEq.lean}{GlobalTheory.Dict.FrameSign.comp_psum}
\leanref{Development/AbelianVarietyTheory/Proofs/SiegelThetaTransformation/AnalyticFibrePrincipalBundleSignEq.lean}{GlobalTheory.Dict.FrameSign.exists_norm_add_le}
% Lean: analyticFibre_principalBundle_sign_eq (99 lines; AnalyticFibrePrincipalBundleSignEq.lean, 294 lines; about
% 2,200 lines with the lemmas it uses, counted through the theorems it reaches in the files it imports), for X :
% Data E r, bas, idx, e, x, L, y : FrPt Xo c C and A : AnFib X bas idx e x L y with IsSympl X bas, x in X.dom and
% SelfD X L (p is A.ex into ptO Xo y.x -> Xo.univ.t.G, F_beta is A.fr, chi is A.chi, the law is A.fr_per with
% hx = hermForm Hm by hx_eq_hermForm). Step 0: Hm_isHermitian and im_hermForm_Hm (Lemma
% is-eigenvalue-conjugate-frame-hermitian(1), (2)), TensLat.span_of_selfD (by_contra: exists_le_ker_of_lt_top,
% DualFamily.exists_herm_eq for z, w = z Gm^-1, DualFamily.eq_zero_of_line) and span_image_pv (that lemma's (4)).
% Step 1 is chi_add (27 lines: fr_per three times, Hm_isHermitian, im_hermForm_Hm) and hz. Step 2: F is
% z -> rowType * z (pv_smul, ex_lin), lam J the sum over J of ![id, F], the morphisms psum G ![1, act w0] J with the
% exponents (-1)^(2 - |J|) of cubeLB in Xo.univ.t.prim (Definition cv-model-toroidal-semiabelian(2)), hlam from
% comp_psum (20 lines), and exists_unit. Step 3 is key (u J, hup, prod_finset_fin_two, a2 = chi_add for v and w0 v).
% Step 4: exists_norm_add_le (16 lines, ZSpan.floor and ZSpan.norm_fract_le) for the image of L and
% FrameUnits.form_eq_zero_and_const. Step 5: key at 0 with hD and hc. The holomorphy of p in the sense of exists_unit
% is the field A.hol (IsHolTo). The fields ex_onto, der and fr_0 of A are not used.

\begin{lemma}\label{cv-uniformization-theta-half-sign}
In the setting of Lemma~\ref{cv-uniformization-theta-marking}, with $P_x$ and $\chi_x$ as there, let
$(G,\eta,\varphi_Z)$ be the Hermitian model of $L_Z$ with $\Lambda_Z=\varphi_Z(L_Z)$, put
$\langle u,v\rangle=\operatorname{Tr}_{E/\mathbf Q}(\eta\,h_G(u,v))$ for $u,v\in E^r$, with $h_G$ as in
Definition~\ref{cv-uniformization-model}, and let $w_0=(1+\sqrt{-d})/2\in\mathcal O_E$ for the prime $d$ of
$\mathcal D$. Then $\chi_x(P_x(v))=\exp(\pi i\langle v,w_0v\rangle)$ for every $v\in\Lambda_Z$.
\end{lemma}
\leanref{Development/AbelianVarietyTheory/Proofs/SiegelThetaTransformation/AnalyticFibrePrincipalBundleSignEq.lean}{GlobalTheory.Dict.analyticFibre_principalBundle_sign_eq}
\begin{proof}
The setting at $\sigma$ is that of Situation~\ref{is-eigenvalue-conjugate-frame-uniformization}, for the datum of $X$
at the level $(M',\mathfrak M',1)$, the coefficient field $F'$, the good embedding $\sigma$ and the uniformization
at $\sigma$, with the lattices $L_Z$ and the maps $t_Z$ and $\varphi_Z$ of the compatibility conditions as a
uniformization witness. Let $X_Z$ be the flag data of $Z$ there, with module $\Lambda_Z$, matrix $G$ and $p_0=q$,
and let $j\mapsto\ell_j$ be the bijection of the situation from $\{1,\dots,g\}$ onto the letters of the datum, which
index the coordinates of $\mathbf C^g$ (Situation~\ref{cv-uniformization-situation}); number the coordinates by it.

The form and the module. By Lemma~\ref{is-eigenvalue-conjugate-frame-component-lattice}(1) and~(3), $\Lambda_Z$ is
self-dual for $X_Z$, and the form $\langle\ ,\ \rangle_Z$ of $L_Z$ is
$(a,a')\mapsto\langle\varphi_Za,\varphi_Za'\rangle_{X_Z}$. By Definition~\ref{cv-uniformization-model}(1) it is also
$(a,a')\mapsto\langle\varphi_Za,\varphi_Za'\rangle$. Both forms on $E^r$ are $\mathbf Q$-bilinear, and $\Lambda_Z$,
an $\mathcal O_E$-module that spans $E^r$ over $E$, spans $E^r$ over $\mathbf Q$; so
$\langle\ ,\ \rangle_{X_Z}=\langle\ ,\ \rangle$.

The periods. By Lemma~\ref{is-eigenvalue-conjugate-frame-letters}, for $n=0$ and $a=0$, there is a letter dictionary
for $X_Z$ with the bijection $j\mapsto\ell_j$, as $r=p+q$, $p_0=q$, $\Psi=\Phi$ and $\Lambda_Z$ is self-dual; fix
its symplectic basis of $\Lambda_Z$ and its indexing, which are matched with $j\mapsto\ell_j$ for $\sigma$. For the
letter $\ell_j=(\tau',i)$, Definition~\ref{is-eigenvalue-conjugate-frame-fibres}(5) gives
$p_x(w)_j=(x_\tau^{-1}\tau(w))_{q+i}$ if the embedding $\tau$ of $\ell_j$ lies in $\Phi$, and
$p_x(w)_j=h_{\bar\tau}(x_{\bar\tau}e_{q+1-i},\bar\tau(w))$ otherwise, which is the formula for the coordinate
$\ell_j$ of $P_x(w)$ in Situation~\ref{cv-uniformization-situation}, with $a=i-1$. So $P_x=p_x$, and the action of
$E$ on $\mathbf C^g$ through the type is the action $b\cdot z=(\rho_j(b)z_j)_j$ through the row types. The frame $x$
of $Z$ lies in the flag domain $\mathcal D_{X_Z}$: both are defined by the signs $(-1)^{\max(0,i-q)}$ of the leading
minors of the Gram matrices of the columns of the $x_\tau$ for $G$.

The form $H_x$. By the compatibility conditions, listed at the beginning of the proof of
Lemma~\ref{cv-cover-uniformization}, $H_x$ is a Hermitian form, linear in the first variable, with
$\operatorname{Im}H_x(P^Z_xa,P^Z_xa')=\langle a,a'\rangle_Z$ for $a,a'\in L_Z$. As $P^Z_x=P_x\circ\varphi_Z$,
$\operatorname{Im}H_x(p_xu,p_xv)=\langle u,v\rangle_{X_Z}$ for $u,v\in\Lambda_Z$. Let $H'_x$ be the form $H_x$ of
Definition~\ref{is-eigenvalue-conjugate-frame-fibres}(2) for $X_Z$, the symplectic basis and the indexing. As
$\Lambda_Z$ spans $E^r$ over $E$, Lemma~\ref{is-eigenvalue-conjugate-frame-hermitian}(4) applies: the periods
$p_x(v)$, $v\in\Lambda_Z$, span $\mathbf C^g$ over $\mathbf R$, and so $H_x=H'_x$.

The analytic fibre. Put $\chi(v)=\chi_x(P_x(v))$ for $v\in\Lambda_Z$ and $\chi(v)=1$ for the other $v\in E^r$. Then
$(p_x,F^x,\chi)$ is an analytic fibre at $y$ for $X_Z$, the symplectic basis, the indexing, the bijection
$j\mapsto\ell_j$, $x$ and $\Lambda_Z$ (Definition~\ref{is-eigenvalue-conjugate-frame-fibres}(3)), with the line
bundle $\mathcal L$ of the polarization and its trivializations over the sets $U_j$ of
Theorem~\ref{cv-uniformization-family}(4). Indeed, by Lemma~\ref{cv-uniformization-theta-marking}(1), $p_x$ is a
surjective holomorphic homomorphism onto $A_y(\mathbf C)$ with the kernel $P_x(\Lambda_Z)=p_x(\Lambda_Z)$, it commutes
with $\mathcal O_E$, and its derivative at $0$ in the direction of the $j$-th unit vector is the vector of the bases
of $y$ at the letter $\ell_j$, which is (a). The $F^x_j$ are holomorphic and without zeros there, and by the
compatibility conditions $F^x_j(z)=g_{jj'}(p_x(z))F^x_{j'}(z)$ where both sides are defined and $F^x_j(0)$ is the
value of the rigidification at $y$, which is (b). And $|\chi(v)|=1$ for $v\in\Lambda_Z$, and the law of the $F^x_j$
in Lemma~\ref{cv-uniformization-theta-marking}(1), with $H_x=H'_x$, is (c).

So Lemma~\ref{cv-uniformization-theta-half-sign-fibre} gives
$\chi_x(P_x(v))=\chi(v)=\exp(\pi i\langle v,w_0v\rangle_{X_Z})=\exp(\pi i\langle v,w_0v\rangle)$ for every
$v\in\Lambda_Z$.
\end{proof}
% Lean: analyticFibre_principalBundle_sign_eq (Lemma cv-uniformization-theta-half-sign-fibre), as
% exists_sec_rdU_ne_zero applies it through isTheta1_of_reads_generator: X = U.X k, L = (U.X k).Lambda, SelfD from
% uw_lattice (hL), bas, idx and e = W.e from the letter dictionary of exists_dic (Dc.sympl, he, Dc.per_eq,
% Dc.ty_eq), x.2 for x in X.dom, and A the AnFib of uw_anFib (Lemma cv-uniformization-theta-marking(1)). The form:
% uw_lattice gives EF X (phi a) (phi b) = <a, b>_Z (hEF of uniformisation_component_lattices_markings). The
% identification H_x = hx is hH in uw_anFib: the Hermitian Hm of the compatibility conditions with hIm is hx by the
% uniqueness clause of hx_laws, whose real span is span_image_pv (span_range_pv and LatFin.exists_int_smul_mem) for
% L = Lambda_Z, which spans E^r over E (hspan of uniformisation_component_lattices_markings). chi is
% if |chil (pv v)| = 1 then chil (pv v) else 1, as in uw_anFib. The Lean's eps1 X w0 v is exp(pi i EF(v, w0 v)).

\begin{lemma}\label{cv-uniformization-theta-half}
In the setting of Lemma~\ref{cv-uniformization-theta-marking}, with $p_x$, $H_x$, $\chi_x$ and $\vartheta_x$ as
there:
\begin{enumerate}
\item $\vartheta_x(z+u)=\chi_x(u)\exp\big(\pi H_x(z,u)+\frac\pi2H_x(u,u)\big)\vartheta_x(z)$ for every
  $z\in\mathbf C^g$ and every $u\in P_x(\Lambda_Z)$;
\item there is $u\in\bar{\mathfrak M}'\Lambda_Z$ with $\vartheta_x\big(P_x(M'^{-1}u)\big)\ne0$.
\end{enumerate}
\end{lemma}
\leanref{Development/AbelianVarietyTheory/Proofs/SiegelThetaTransformation/ReadsExistsUniqueTheta.lean}{GlobalTheory.Dict.reads_existsUnique_theta}
\leanref{Development/AbelianVarietyTheory/Proofs/SiegelThetaTransformation/FramedPointGeneratorIsThetaLevelOne.lean}{GlobalTheory.Dict.isTheta1_of_reads_generator}
\leanref{Development/AbelianVarietyTheory/Proofs/DivisorProofs/ExistsLevelOneThetaNeZeroOnIdealHalfLatticeStatement.lean}{GlobalTheory.T.exists_levelOne_theta_ne_zero_on_ideal_half_lattice}
\begin{proof}
(1) Let $z\in\mathbf C^g$ and $u\in P_x(\Lambda_Z)$, and choose $j$ with $p_x(z)\in U_j$. As $p_x(u)=0$,
$p_x(z+u)=p_x(z)$, and the section of $y$ has one value there, so
$\vartheta_x(z+u)F^x_j(z+u)=\vartheta_x(z)F^x_j(z)$. Replacing $F^x_j(z)$ by
$\chi_x(u)\exp\big(\pi H_x(z,u)+\frac\pi2H_x(u,u)\big)F^x_j(z+u)$ and dividing by $F^x_j(z+u)\ne0$ gives (1).

(2) The function $\vartheta_x$ is not $0$, since the mean of $|\vartheta_x(z)|^2e^{-\pi H_x(z,z)}$ is $1$. Write
$P=P_x$, $H=H_x$ and $\Lambda'=\bar{\mathfrak M}'\Lambda_Z$.

Step 1: $\vartheta_x$ is holomorphic. On the open set of the $z$ with $p_x(z)\in U_j$ it is the value of the section
of $y$ at $p_x(z)$ in the trivialization over $U_j$, a holomorphic function of $z$, divided by the nowhere vanishing
holomorphic function $F^x_j$; and these open sets cover $\mathbf C^g$.

Step 2: the form on $E^r$. For $u,v\in E^r$ put $\langle u,v\rangle=\operatorname{Tr}_{E/\mathbf Q}(\eta\,h_G(u,v))$,
with $h_G$ as in Definition~\ref{cv-uniformization-model}. As $h_G$ is linear in the first variable and conjugate
linear in the second, $\langle au,bv\rangle=\langle a\bar bu,v\rangle$ for $a,b\in E$. By that definition,
$\varphi_Z$ carries the perfect alternating form $\langle\ ,\ \rangle_Z$ of $L_Z$ to $\langle\ ,\ \rangle$ on
$\Lambda_Z$; so $\langle\ ,\ \rangle$ takes integer values on $\Lambda_Z$, and $u\mapsto\langle u,\cdot\rangle$ is a
bijection of $\Lambda_Z$ onto $\operatorname{Hom}(\Lambda_Z,\mathbf Z)$. As $\Lambda_Z$ spans $E^r$ over
$\mathbf Q$, the form is nondegenerate on $E^r$, and a vector $u\in E^r$ with $\langle u,\Lambda_Z\rangle\subset\mathbf Z$
lies in $\Lambda_Z$: some $u'\in\Lambda_Z$ has $\langle u',\cdot\rangle=\langle u,\cdot\rangle$ on $\Lambda_Z$, hence
on $E^r$, and $u=u'$.

Step 3: for $u\in E^r$, $u\in\Lambda'$ if and only if $M'^{-1}\langle u,v\rangle\in\mathbf Z$ for every $v\in\Lambda'$.
Let $u=au_0$ and $v=bv_0$ with $a,b\in\bar{\mathfrak M}'$ and $u_0,v_0\in\Lambda_Z$. Then
$a\bar b\in\bar{\mathfrak M}'\mathfrak M'=M'\mathcal O_E$, say $a\bar b=M't$, and
$M'^{-1}\langle u,v\rangle=\langle tu_0,v_0\rangle\in\mathbf Z$; by additivity the condition holds for every
$u\in\Lambda'$. Conversely, suppose that it holds for $u$. For $a\in\bar{\mathfrak M}'$ and $v_0\in\Lambda_Z$,
$\langle\bar aM'^{-1}u,v_0\rangle=M'^{-1}\langle u,av_0\rangle\in\mathbf Z$, so $\bar aM'^{-1}u\in\Lambda_Z$ by
Step 2. The conjugates of the elements of $\bar{\mathfrak M}'$ make up $\mathfrak M'$, so
$\mathfrak M'M'^{-1}u\subset\Lambda_Z$; and as $M'\in\bar{\mathfrak M}'\mathfrak M'$, $u=M'\cdot M'^{-1}u$ lies in
$\bar{\mathfrak M}'\Lambda_Z=\Lambda'$.

Step 4: the lattice $P(\Lambda')$. By the compatibility conditions, listed at the beginning of the proof of
Lemma~\ref{cv-cover-uniformization}, $\operatorname{Im}H(P^Z_xa,P^Z_xa')=\langle a,a'\rangle_Z$ for $a,a'\in L_Z$; so
$\operatorname{Im}H(Pu,Pv)=\langle u,v\rangle$ for $u,v\in\Lambda_Z$, and then for $u,v\in E^r$, as $P$ is
$\mathbf Q$-linear. So $P$ is injective on $E^r$: if $Pu=0$, then $\langle u,\cdot\rangle=0$, and $u=0$ by Step 2.
Put $\Lambda''=P(\Lambda')$, $H'=H/M'$ and $E'=\operatorname{Im}H'$.

The set $\Lambda''$ spans $\mathbf C^g$ over $\mathbf R$. First, $P(E^r)$ spans $\mathbf C^g$ over $\mathbf R$. Let
$(e_1,\dots,e_g,f_1,\dots,f_g)$ be a basis of $\Lambda_Z$ with $\langle e_k,e_l\rangle=\langle f_k,f_l\rangle=0$ and
$\langle e_k,f_l\rangle=-\delta_{kl}$, which exists by Lemma~\ref{cv-uniformization-adapted-basis} for $N=1$, as
$\langle\ ,\ \rangle$ is perfect on $\Lambda_Z$, after replacing each $f_k$ by $-f_k$; and let $A$ and $C$ be the matrices with the columns $P(e_k)$ and $P(f_k)$. As $H$ is
positive definite with $\operatorname{Im}H(Pu,Pv)=\langle u,v\rangle$, Riemann's relations, as in the proof of
Lemma~\ref{cv-volume-key-normalized-riemann}(1), show that $A$ is invertible and that $\Omega=A^{-1}C$ has positive
definite imaginary part. For $z\in\mathbf C^g$ put $\zeta=A^{-1}z$ and choose $t\in\mathbf R^g$ with
$(\operatorname{Im}\Omega)t=\operatorname{Im}\zeta$; then $\zeta-\Omega t$ is real, and $z=A(\zeta-\Omega t)+Ct$ is a real
combination of the $P(e_k)$ and the $P(f_k)$. Second, every
$w\in E^r$ has a nonzero integer $k$ with $kw\in\Lambda'$: as $\Lambda_Z$ spans $E^r$ over $E$, some nonzero integer
$k_0$ has $k_0w\in\Lambda_Z$, and $k=M'k_0$ will do, as $M'\in\bar{\mathfrak M}'$. Then $P(w)=k^{-1}P(kw)$ lies in the
real span of $\Lambda''$.

For $u,v\in\Lambda'$, $E'(Pu,Pv)=M'^{-1}\langle u,v\rangle$, an integer by Step 3. If $l\in\mathbf C^g$ has
$E'(l,m)=0$ for every $m\in\Lambda''$, then $E'(l,\cdot)=0$ on $\mathbf C^g$, as $\Lambda''$ spans it over $\mathbf R$;
in particular $H'(l,l)=-E'(l,il)=0$, and $l=0$. The group $\Lambda'$ is finitely generated, as a subgroup of the
finitely generated group $\Lambda_Z$; let $m_1,\dots,m_n$ generate $\Lambda''$. An $l\in\Lambda''$ with
$|E'(l,m_j)|<1$ for every $j$ has $E'(l,m_j)=0$, as these are integers, so $E'(l,\cdot)=0$ on $\Lambda''$ and $l=0$.
So $0$ is isolated in $\Lambda''$, and $\Lambda''$ is a discrete subgroup of $\mathbf C^g$ that spans it over
$\mathbf R$: a lattice.

The map $l\mapsto E'(l,\cdot)$ from $\Lambda''$ to $\operatorname{Hom}(\Lambda'',\mathbf Z)$ is injective by the
previous paragraph. It is onto. For $u\in E^r$ with $M'^{-1}\langle u,\cdot\rangle=0$, $E'(Pu,\cdot)$ vanishes on
$\Lambda''$, so $Pu=0$ and $u=0$; as $E^r$ is finite-dimensional over $\mathbf Q$, $u\mapsto M'^{-1}\langle u,\cdot\rangle$
is a bijection of $E^r$ onto its dual. A homomorphism $\psi\colon\Lambda''\to\mathbf Z$ gives $\psi\circ P$ on
$\Lambda'$, which extends to a $\mathbf Q$-linear form on $E^r$, as $\mathbf Q$ is divisible; this form is
$M'^{-1}\langle u,\cdot\rangle$ for some $u\in E^r$. Then $u\in\Lambda'$ by Step 3, and $\psi=E'(Pu,\cdot)$.

Step 5: the conclusion. Let $l\in\Lambda''$, and write $l=Pv$ with $v\in\Lambda'$. With $w_0=(1+\sqrt{-d})/2$, put
$\varepsilon(v)=\exp(\pi i\langle v,w_0v\rangle)$; Lemma~\ref{cv-uniformization-theta-half-sign} gives
$\chi_x(l)=\varepsilon(v)$. As $w_0v\in\Lambda_Z$, $\langle v,w_0v\rangle$ is an integer by Step 2, so
$\varepsilon(v)^{-1}=\varepsilon(v)$, and (1) gives
\[
 \vartheta_x(z+l)=\varepsilon(v)^{-1}\exp\big(M'\big(\pi H'(z,l)+\tfrac\pi2H'(l,l)\big)\big)\vartheta_x(z).
\]
As $\langle v,w_0v\rangle$ is real, $|\varepsilon(v)^{-1}|=1$. By Step 1, $\vartheta_x$ is holomorphic, and it is not
$0$. As $M'\ge3$, Lemma~\ref{cv-uniformization-theta-half-principal} for $\Lambda''$, $H'$, the function
$c(Pv)=\varepsilon(v)^{-1}$, $N=M'$ and $f=\vartheta_x$ gives $l=Pu$ with $u\in\Lambda'$ and $\vartheta_x(l/M')\ne0$.
As $P$ is $\mathbf Q$-linear, $l/M'=P(M'^{-1}u)$.
% Lean: (1) is the last clause of reads_existsUnique_theta with k = 1 (the values of the section at ex z and
% ex (z + pv v), fr_per, and cancelling fr), in the form (chi v)^{-1} fock x 1 (pv v) theta = theta. (2) is
% GlobalTheory.Dict.ThetaHalf.exists_ne_zero_on_half for the ideal Mfr l = conj(Mb) and M' (with
% conj(Mfr l) * Mfr l = (M'), from dual_marking), applied to the level-one theta function given by
% isTheta1_of_reads_generator and to theta /= 0 (AdTh.GenConst.ne_zero_of_meanOne); this is the first step of
% AdTh.GenConst.exists_root_ne_zero. Step 1 is the differentiability clause of reads_existsUnique_theta
% (isHolFam_fibrePoint, fr_hol, fr_ne). Step 2 is SelfD of uw_lattice (UWLat.selfD_of_unimod) for the trace form EF.
% Step 3 is ThetaHalf.mem_smul_iff. Step 4: him and hinj of hx_laws (hinj is proved from him, as in the text);
% the real span is hspan of hx_laws (span_image_pv: span_range_pv, from isUnit_det_perMat and tauMat_im_posDef,
% and GlobalTheory.LatFin.exists_int_smul_mem) with ThetaHalf.span_smul; then, in
% ThetaDivisionPoints.exists_ne_zero_of_selfDual, hint and hk (Ef integral), hker (eq_zero_of_im_eq_zero), hdisc
% (the finitely many generators), hBinj and hBsurj, and hbij (Module.Baer.of_divisible). Step 5 is key and
% ThetaHalf.exists_ne_zero_of_scaled (B and Hm divided by M), then
% T.thetaFunction_ne_zero_at_divisionPoint_selfDual (Lemma cv-uniformization-theta-half-principal), and hsm for
% P(M'^-1 u) = P(u)/M'. The multiplier is that of IsTheta1 from isTheta1_of_reads_generator: (1), then
% analyticFibre_principalBundle_sign_eq (Lemma cv-uniformization-theta-half-sign, chi = eps1) and
% thetaSignExpEF_inv_of_selfDual_lattice (eps1^-1 = eps1); exists_ne_zero_on_half takes c = eps1^-1 and uses only
% |c| = 1 (eps1 is exp of i times a real number).
\end{proof}
\leanref{Development/AbelianVarietyTheory/Proofs/DivisorProofs/ExistsLevelOneThetaNeZeroOnIdealHalfLattice.lean}{GlobalTheory.Dict.ThetaHalf.exists_ne_zero_on_half}
\leanref{Development/AbelianVarietyTheory/Proofs/DivisorProofs/ExistsLevelOneThetaNeZeroOnIdealHalfLattice.lean}{GlobalTheory.Dict.ThetaHalf.mem_smul_iff}
\leanref{Development/AbelianVarietyTheory/Proofs/DivisorProofs/ExistsLevelOneThetaNeZeroOnIdealHalfLattice.lean}{GlobalTheory.Dict.ThetaHalf.exists_ne_zero_of_scaled}
\leanref{Development/AbelianVarietyTheory/Proofs/SiegelThetaTransformation/HermitianFormAtFrameLaws.lean}{GlobalTheory.Dict.span_image_pv}
\leanref{Development/AbelianVarietyTheory/Proofs/SiegelThetaTransformation/FramedPointGeneratorIsThetaLevelOne.lean}{GlobalTheory.Dict.thetaSignExpEF_inv_of_selfDual_lattice}

\subsection{Taylor families and theta values}\label{subsection-cv-taylor-families}

\begin{definition}\label{cv-uniformization-theta-rule-values}
In the setting of Lemma~\ref{cv-uniformization-theta-marking}, let $\mathcal A$ be the semi-abelian scheme of the
universal object over $X'$, with the line bundle $\mathcal L$ of the polarization, its rigidification along the zero
section and the level structure $\eta$ (Definition~\ref{cv-toroidal-moduli-object}(4) and~(5)), and put
$w_0=(1+\sqrt{-d})/2$. Framed points over an $F'$-algebra $S$ are as in
Definition~\ref{cv-volume-twisting-section}(3), with their tangent vectors indexed by the letters of the datum
(Definition~\ref{is-eigenvalue-conjugate-frame-fibres}(3)). For such a framed point $y$ let $A_y$ be the fibre of
$\mathcal A$ over the point under $y$, $\mathcal L_y$ the restriction of $\mathcal L$ to it, $\vartheta_y$ the
generator of $y$, and $\eta_y(c)\in A_y(S)$ the point given by $\eta(c)$, for $c\in L$. A framed point $z$ over an
$F'$-algebra $T$ is an \emph{image} of $y$ under a homomorphism of $F'$-algebras $\varphi\colon S\to T$ if the point
under $z$ is the point under $y$ composed with $\operatorname{Spec}\varphi$, the tangent vectors of $z$ are those of
$y$ base changed along $\varphi$, and $\vartheta_z$ is the pullback of $\vartheta_y$ to $A_z=A_y\times_ST$.
\begin{enumerate}
\item (Translates.) Let $y$ be a framed point over $S$, $N\ge1$ an integer and $t,t'\in A_y(S)$ with $Nt'=t$. An
  \emph{$N$-translate} of $\vartheta_y$ over $t$ with \emph{root} $t'$ is a section $F=Us_1$ of $\mathcal L_y^{N^2}$,
  where $s_1$ is the section of $\mathcal L_y^{N^2}$ that corresponds to $[N]^*\vartheta_y$ under the isomorphism
  $[N]^*\mathcal L_y\cong\mathcal L_y^{N^2}$ compatible with the rigidifications, and $U$ is a symmetric element of
  the theta group of $\mathcal L_y^{N^2}$ that lies over $t'$ and satisfies $U^{N^2}=1$
  (Definition~\ref{cv-uniformization-theta-rule-unique-translates-symmetric}, for $A_y$ over $S$; this asks
  $N^2t'=0$). A \emph{level translate} of
  $\vartheta_y$ over $t$ is an $M'$-translate of $\vartheta_y$ over $t$ with a root $t'$ such that $st'=0$ for every
  $s\in\mathfrak M'^2$.
\item (Taylor families.) A \emph{Taylor family} $\nabla$ assigns to every framed point $y$ over an $F'$-algebra $S$,
  every integer $k$, every section $F$ of $\mathcal L_y^k$ and every family $\alpha=(\alpha_\ell)$ of nonnegative
  integers indexed by the letters an element $\nabla_y(F;\alpha)\in S$.
\item (Theta values.) Let $n\ge0$, $c_1,\dots,c_n\in L$, $P$ a polynomial in the letter variables of~(6) with
  coefficients in $F'$, and $\nabla$ a Taylor family. For a framed point $y'$ over $S'$ and sections
  $F_1,\dots,F_n$ of $\mathcal L_{y'}^{M'^2}$ put
  \[
   \nabla_{y'}(P;F_1,\dots,F_n)=M'^{-\deg P}\sum_mP_m\prod_{b=1}^n\nabla_{y'}(F_b;m_b),
  \]
  where $m=(m_{\ell,b})$ runs over the exponents of the monomials $\prod_{\ell,b}X_{\ell,b}^{m_{\ell,b}}$ of $P$, $P_m$
  is the coefficient of that monomial, $m_b=(m_{\ell,b})_\ell$, and $\deg P$ is the total degree of $P$; for $P=1$
  this is $\prod_b\nabla_{y'}(F_b;0)$. A \emph{theta value} of $c_1,\dots,c_n$ and $P$ at a framed point $y$ over $S$
  for $\nabla$ is an element $v\in S$ for which there are a faithfully flat homomorphism of $F'$-algebras
  $\varphi\colon S\to S'$, an image $y'$ of $y$ under $\varphi$, and level translates $F_b$ of $\vartheta_{y'}$ over
  $\eta_{y'}(c_b)$, $1\le b\le n$, such that $\varphi(v)=\nabla_{y'}(P;F_1,\dots,F_n)$.
\item (Theta rules.) A \emph{theta rule} of $c_1,\dots,c_n$ and $P$ for $\nabla$ assigns to every framed point $y$
  over a test ring, a quotient of a polynomial ring over $F'$ by a radical ideal, a theta value $f(y)$ of
  $c_1,\dots,c_n$ and $P$ at $y$ for $\nabla$.
\item (Canonical families.) A \emph{chart datum} at a framed point $y$ over an $F'$-algebra $R$ consists of the
  following data. First, an integer $d'\ge0$, a formal germ $B=(F,[\,\cdot\,],\Phi_1,\Phi_2,S,P)$ of dimension $d'$
  over $R$ (Definition~\ref{cv-hodge-lattice-germ}(1)), points $x_m$ of $\mathcal A$ with values in the rings $J^1_m$
  over the point under $y$ and generators $e_m$ of $\Lambda(x_m(t),x_m(t'))$ on $\operatorname{Spec}J^2_m$, for
  $m\ge0$, in the notation of Definition~\ref{cv-germs-chart}(1) and~(2) with $d=d'$, that satisfy conditions (a),
  (d) and (e) of Definition~\ref{cv-germs-chart}(3). Let $v_j$, for $1\le j\le d'$, be the image of $x_1$ under
  $t_i\mapsto\delta_{ij}\varepsilon$, a tangent vector at the point under $y$. Second, a matrix $(u_{\ell j})$ over
  $R$, indexed by the letters $\ell$ and $1\le j\le d'$, such that the tangent vector of $y$ at each letter $\ell$ is
  $\sum_ju_{\ell j}v_j$, and that is invertible: for every $c\in R^{d'}$ there is exactly one family $(a_\ell)$ of
  elements of $R$ indexed by the letters with $\sum_\ell a_\ell u_{\ell j}=c_j$ for every $j$. Third, a
  generator $w$ of $\mathcal L\otimes\Lambda(o,w_0o)^{-1}$ on $\mathcal A$, with $o$ the identity of $\mathcal A$,
  whose pullback along the zero section is the rigidification. Fourth, a power series $h(t,t')$ over $R$ with
  $h(0,t')=h(t,0)=0$ whose exponentials give the series of $B$:
  \begin{align*}
   \Phi_1(t,t',t'')&=e^{h(t,t'')+h(t',t'')-h(t+_Ft',t'')},&\Phi_2(t,t',t'')&=e^{h(t,t')+h(t,t'')-h(t,t'+_Ft'')},\\
   S(t,t')&=e^{h(t',t)-h(t,t')},&P_a(t,t')&=e^{h([a]t,t')-h(t,[\bar a]t')}.
  \end{align*}
  Along each $x_m$ the line bundle $\mathcal L$ is trivialized by $w(x_m(t))\,e_m(t,[w_0]t)$, and for $m=0$ this
  trivialization is the rigidification. A \emph{germ} of a section $s$ of $\mathcal L_y^k$, $k\ge0$, for the chart
  datum is a power series $f_s$ over $R$ whose class modulo $I^{m+1}$ is the value of $s$ at $x_m$ in the $k$-th power
  of this trivialization, for every $m$. A Taylor family $\nabla$ is \emph{canonical} if every framed point $y$ over
  an $F'$-algebra $R$ has a chart datum for which every section $s$ of $\mathcal L_y^k$, $k\ge0$, has a germ, and such
  that for every germ $f_s$ of such a section and every $\alpha$, $\nabla_y(s;\alpha)$ is the constant term of
  $\big(\prod_\ell D_\ell^{\alpha_\ell}\big)\big(f_s(t)\,e^{k\,h(t,[w_0]t)}\big)$, where
  $D_\ell=\sum_j\big(\sum_iu_{\ell i}\,\partial F_j(t,t')/\partial t'_i|_{t'=0}\big)\,\partial/\partial t_j$.
\item (Harmonic polynomials.) Let $n\ge0$. The \emph{letter variables} are the $X_{\ell,b}$, for the letters $\ell$
  and $1\le b\le n$. Let $a=(a_\ell)$ be a family of nonnegative integers indexed by the letters. A polynomial $P$ in
  the letter variables with coefficients in $F'$ is \emph{harmonic of row degrees $a$} if:
  \begin{enumerate}
  \item[(a)] $\sum_{b=1}^n\partial^2P/\partial X_{\ell,b}\partial X_{\ell',b}=0$ for all letters $\ell=(\tau,i)$ with
    $\tau\in\Phi$ and $\ell'=(\bar\tau,i')$;
  \item[(b)] for every letter $(\tau,i)$, every $i'<i$ and every $\lambda\in F'$, $P$ is unchanged when
    $X_{(\tau,i),b}$ is replaced by $X_{(\tau,i),b}+\lambda X_{(\tau,i'),b}$ for every $b$;
  \item[(c)] for every letter $\ell$, every monomial of $P$ has degree $a_\ell$ in the variables
    $X_{\ell,1},\dots,X_{\ell,n}$;
  \item[(d)] the coefficients of $P$ are algebraic integers.
  \end{enumerate}
\end{enumerate}
\end{definition}
\leanref{Formalism/LocalToGlobal/ProjectiveLattices/IsThRule.lean}{Formalism.L2.Pj.IsVtOver}
\leanref{Formalism/LocalToGlobal/ProjectiveLattices/IsThRule.lean}{Formalism.L2.Pj.NabV}
\leanref{Formalism/LocalToGlobal/ProjectiveLattices/IsThRule.lean}{Formalism.L2.Pj.thExpr}
\leanref{Formalism/LocalToGlobal/ProjectiveLattices/IsThRule.lean}{Formalism.L2.Pj.ThValAt}
\leanref{Formalism/LocalToGlobal/ProjectiveLattices/IsThRule.lean}{Formalism.L2.Pj.IsThRule}
\leanref{Formalism/LocalToGlobal/ProjectiveLattices/ThGerm.lean}{Formalism.L2.Pj.ChartAt}
\leanref{Formalism/LocalToGlobal/ProjectiveLattices/ThGerm.lean}{Formalism.L2.Pj.IsGerm}
\leanref{Formalism/LocalToGlobal/ProjectiveLattices/ThGerm.lean}{Formalism.L2.Pj.IsNab}
\leanref{Formalism/LocalToGlobal/ProjectiveLattices/IsCombAt.lean}{Formalism.L2.Pj.FrMap}
\leanref{Formalism/LocalToGlobal/ProjectiveLattices/IsThRule.lean}{Formalism.L2.Pj.harmA}
\leanref{Formalism/LocalToGlobal/ProjectiveLattices/RowDeg.lean}{Formalism.L2.Pj.HHa}
\leanref{Formalism/LocalToGlobal/ProjectiveLattices/RowDeg.lean}{Formalism.L2.Pj.HaO}
% IsVtOver G PP N ok s t F: t' with t'^N = t and ok t', the normalized alpha0 (IsNormMul), its pullback alpha, the
% twisted pullback s_1 of s along [N] (IsTwistPull), U with SymLift (over t', U^(N^2) = 1, symmetric) and
% IsActOf U s_1 F; a level translate is IsVtOver with N = M' and ok = KilledBy (Mb^2). NabV is the Taylor family (k an
% integer, alpha in Lt1 Dw ->0 N). ThValAt: thExpr is M'^(-deg P) sum_m P_m prod_b nab y' (M'^2) F_b (colOf m b),
% with deg P = P.totalDegree, which for P = 1 is prod_b nab y' (M'^2) F_b 0; FrMap phi y y' is the image (the point
% spMap phi >> y.x, the tangent vectors through dualMap phi, and the section pulled back along the induced map of the
% fibres); the level point is levPt (toUnit then Xo.univ.eta). (6) is harmA c a n = HaO: HHa (harmK with cc the
% complex conjugation, isX = InPsi H, tau(sqrt(-d)) = rt0, and hinv = 1, so the Laplacian sum_b d^2/dX_(tau,i,b)
% dX_(conj tau,i',b); RowInv with rowShift; RowDeg with a (tau, i)) and integral coefficients; the letter variables
% are LtN Dw n, triples (tau, i, b).
% IsThRule is on TRing (MvPolynomial (Fin n) F' modulo a radical ideal). IsNab: star = bar on O_E, and for every
% framed point over a Q-algebra a ChartAt C (IsChart: zero, trunc, law, act, condition (a) of Definition
% cv-germs-chart(3); IsFrame: gen, ph1, ph2, sy, ad, conditions (d) and (e); neither the basis property of the v_j in
% (b) nor FrTrunc and FrAxes of (c); uvec with tan and basis, the invertibility of the matrix; wP with
% IsPrim; h with Germ.FB.IsCanonGauge: ax0, ax1, ph1, ph2, sy, ad; rig0 for eL at m = 0), every F has a germ
% (IsGerm through eL = wP(x_m) e_m(t, [w0] t)), and nab y k F alpha = Germ.FB.nabVal: the constant coefficient of
% DpolyOf FK uvec (monomial alpha) (f * exp(k hL)), with hL(t) = h(t, [w0] t) and Dof the invariant derivation.

\begin{lemma}\label{cv-uniformization-theta-rule-connection-charts-frames}
In the setting of Lemma~\ref{cv-uniformization-theta-marking}, let $x_0$ be a point of $X'$ over a ring $R$, and let
$x_m$, $m\ge0$, be points of $\mathcal A$ with values in the rings $J^1_m$ over $x_0$, with a commutative formal group
law $F$ of dimension $d$ and series $[a]$, $a\in\mathcal O_E$, satisfying conditions (a) and (b) of
Definition~\ref{cv-germs-chart}(3). Then there are generators $e_m$ of $\Lambda(x_m(t),x_m(t'))$ on
$\operatorname{Spec}J^2_m$, power series $\Phi_1$, $\Phi_2$, $S$ and $P_a$, $a\in\mathcal O_E$, over $R$ with $P_a(0,0)=1$,
and a generator $w$ of $\mathcal L\otimes\Lambda(o,w_0o)^{-1}$ on $\mathcal A$ whose pullback along the zero section is
the rigidification, such that conditions (c) to (e) of Definition~\ref{cv-germs-chart}(3) hold and the trivialization
$w(x_0(t))\,e_0(t,[w_0]t)$ of $\mathcal L$ along $x_0$ is the rigidification.
\end{lemma}
\leanref{Development/AbelianVarietyTheory/Proofs/CuspFormEvaluations/ExistsFramesOfFormalChart.lean}{GlobalTheory.exists_frames_of_chart}
\begin{proof}
\openstep{The construction of the generators $e_m$ and of the series from the cubical structure along the points
$x_m$ is not yet written out.}
% Lean: exists_frames_of_chart: T.exists_frames_cubeDefect_along_chart_over_ring (e with IsGen, FrTrunc and FrAxes),
% T.exists_firstPartialLaw_series_of_frames_alt and T.exists_second_partial_law_alt (Phi_1, Phi_2),
% T.exists_symmetry_series_of_frames_alt (S), T.exists_adjointness_series_of_frames (P_a with constant term 1),
% GermB.exists_wP (w with IsPrimG) and T.frame_rigidification_of_canonical_axes (Rig0G). The Lean states it for any
% semi-abelian scheme with the structures of the moduli problem over a scheme. Compare the construction in the proof
% of Lemma cv-germs-local, over an affine open set with a trivialization of the line bundle. Size beyond the laws of
% abelian schemes and chart independence: about 1,700 lines.
\end{proof}

\begin{lemma}\label{cv-uniformization-theta-rule-connection-charts-identities}
In the setting of Lemma~\ref{cv-uniformization-theta-marking}, let $x_0$, the points $x_m$ and the series $F$ and
$[a]$ be as in Lemma~\ref{cv-uniformization-theta-rule-connection-charts-frames}. Let $e_m$ be generators of
$\Lambda(x_m(t),x_m(t'))$ on $\operatorname{Spec}J^2_m$ with $e_m(0,t)=1=e_m(t,0)$, and $\Phi_1$, $\Phi_2$, $S$ and
$P_a$, $a\in\mathcal O_E$, power series over $R$ with $P_a(0,0)=1$, that satisfy conditions (d) and (e) of
Definition~\ref{cv-germs-chart}(3). Then $(F,[\,\cdot\,],\Phi_1,\Phi_2,S,P)$ is a formal germ
(Definition~\ref{cv-hodge-lattice-germ}(1)).
\end{lemma}
\leanref{Development/AbelianVarietyTheory/Proofs/CuspFormEvaluations/FrameBiextDictionarySecondPartialLaw.lean}{GlobalTheory.FrameBiextId.fbId}
\leanref{Development/AbelianVarietyTheory/Proofs/CuspFormEvaluations/ExistsBiextensionChartWithGermsFramedPoint.lean}{GlobalTheory.ChartRing.bijective_appTop_powObj}
\begin{proof}
The global functions. For every $n\ge1$ every function on the $n$-fold fibre product $\mathcal A^n$ of $\mathcal A$ over
$X'$ is the pullback of a function on $X'$, by induction on $n$: $\mathcal A^n$ is an abelian scheme over
$\mathcal A^{n-1}$, the base change of $\mathcal A$, so its global functions are those of $\mathcal A^{n-1}$.

Reduction to the rings $J^k_m$. A power series over $R$ is determined by its reductions modulo $I^{m+1}$ for all $m$,
so each identity of Definition~\ref{cv-hodge-lattice-germ}(1) may be checked in $J^k_m$, with $k\le4$, for every $m$.
Substituting into the relations (d) and (e) writes each side as a product of values of the cubical structure
$\theta$, of the $e_m$ and of the $w_a$ at points $x_m(u)$, which are added and multiplied by elements of
$\mathcal O_E$ through $x_m(u+_Fu')=x_m(u)+x_m(u')$ and $x_m([a]u)=a\,x_m(u)$. The factors $e_m$ cancel, using
$e_m=1$ on the axes, and each identity becomes one between values of $\theta$ and of the $w_a$.

The identities in $\Phi_1$, $\Phi_2$ and $S$. These follow from the conditions on $\theta$ in
Definition~\ref{cv-germs-chart}(1): it is $1$ when one of its arguments is $0$, symmetric, and satisfies the cocycle
identity.

The identities in $P$. First, $w_a(0,0)=1$: condition (e) at $t=t'=0$ reads
$w_a(0,0)e_m(0,0)=P_a(0,0)e_m(0,0)$, with $P_a(0,0)=1$. The identities with $P$ reduce to the five relations displayed
in the proof of Lemma~\ref{cv-germs-local}, between values of the $w_a$ and of $\theta$ at points $y,y',y''$ of
$\mathcal A$. For the universal points, the projections of $\mathcal A^3$, the two sides of each relation are
generators of the same line bundle, so their ratio is an invertible function on $\mathcal A^3$; by the first
paragraph it is the pullback of a unit on $X'$, and it is $1$ on the unit section, as $w_a(0,0)=1$ and $\theta$ is
$1$ when one of its arguments is $0$. So each relation holds, and then so does each identity with $P$.
\end{proof}
% Lean: GlobalTheory.FrameBiextId.fbId (FBIdSt: AxesOK and FBRest for mkFB G A ph1 ph2 sy ad), for any semi-abelian
% scheme t with bijective appTop on its fibre powers (ChartRing.bijective_appTop_powObj, from the law GlobalFnsSt),
% star = bar, IsChartG, IsFrameG, FrAxes (e_m = 1 on the axes) and constantCoeff (ad a) = 1. Each of the identities is
% VV_ext (a series is determined by its values at all points of the jet schemes) with the dictionaries
% frame_biext_dictionary_first_partial_law, frame_biext_dictionary_second_partial_law, dict_sy and dict_ad (the
% relations (d) and (e) at the points), the face, symmetry and cocycle laws of the cubical structure (isCub_cubV),
% dict_ax for e on the axes, and for the identities with ad the relations cubical_law_left, cubical_law_right,
% adv_one, adv_mul, adv_add and adv_sy of the w_a, from hp = (hPow 3).2 and h1 (w_a(1, 1) = 1, from dict_ad at 0,
% had and FrAxes). It does not use FrTrunc. The same argument is the Identities paragraph of the proof of Lemma
% cv-germs-local, which cites FBIdAd.fbId.

\begin{lemma}\label{cv-uniformization-theta-rule-connection-charts-gauge}
Let $E$ be a CM field, let $B$ be a formal germ of dimension $d$ over a commutative ring $R$
(Definition~\ref{cv-hodge-lattice-germ}), and let $k$ be an $R$-algebra that contains $\mathbf Q$. Then there is a power
series $h(t,t')$ over $k$ with $h(0,t')=h(t,0)=0$ whose exponentials give the series of $B$ over $k$, as in
Definition~\ref{cv-uniformization-theta-rule-values}(5).
\end{lemma}
\leanref{Development/AbelianVarietyTheory/Proofs/FormalGroupCharts/CanonGaugeExists.lean}{GlobalTheory.canonGauge_exists}
\begin{proof}
\openstep{The construction of the canonical gauge, through the logarithm of the formal group law over $k$, is not yet
written out.}
% Lean: canonGauge_exists : CanonGaugeSt = canonGauge_of_log FormalLog.logSt_and_logDerivSt.1, for any commutative
% ring O with a star structure, free and finite over Z and reduced, in place of O_E. Size beyond the laws of abelian
% schemes, chart independence and the lemmas before this one: about 2,960 lines.
\end{proof}

\begin{lemma}\label{cv-uniformization-theta-rule-connection-charts-germs}
In the setting of Lemma~\ref{cv-uniformization-theta-marking}, let $y$ be a framed point of $X'$ over an
$F'$-algebra $R$, and let the points $x_m$, the generators $e_m$ and $w$ and the series be as in
Lemmas~\ref{cv-uniformization-theta-rule-connection-charts-frames} and~\ref{cv-uniformization-theta-rule-connection-charts-identities}
for the point under $y$. Then for every $k\ge0$, every section of $\mathcal L_y^k$ has a germ for these data, in the
sense of Definition~\ref{cv-uniformization-theta-rule-values}(5).
\end{lemma}
\leanref{Formalism/SharpRestriction/Divisibility/GeneratorPowers/GermDPf.lean}{Formalism.MainNotions.ModuliIntegralCharts.H1.germSt_holds}
\begin{proof}
\openstep{The proof that the values of a section at the points $x_m$, read in the trivializations
$w(x_m(t))\,e_m(t,[w_0]t)$, are the truncations of one power series is not yet written out.}
% Lean: germSt_holds (GermSt), through the generators eLG of the trivializations, their powers, and the
% compatibility of the e_m under truncation (FrTrunc). Size: about 2,000 lines, nearly all shared with Lemma
% cv-uniformization-theta-rule-connection-charts-frames; the translation between the Lean's two forms of charts,
% frames and germs (isChart_iff, isFrame_iff, isPrim_iff, isGerm_iff) adds about 1,230 lines to Lemma
% cv-uniformization-theta-rule-connection-charts.
\end{proof}

\begin{lemma}\label{cv-uniformization-theta-rule-connection-charts}
In the setting of Lemma~\ref{cv-uniformization-theta-marking}, every framed point $y$ of $X'$ over an $F'$-algebra
$R$ has a chart datum (Definition~\ref{cv-uniformization-theta-rule-values}(5)) whose generators $e_m$ satisfy: $e_{m+1}$
reduces to $e_m$ modulo $I^{m+1}$ for every $m\ge0$; and for which every section of $\mathcal L_y^k$, $k\ge0$, has a
germ.
\end{lemma}
\leanref{Development/AbelianVarietyTheory/Proofs/CuspFormEvaluations/ExistsBiextensionChartWithGermsFramedPointStatement.lean}{GlobalTheory.T.exists_biextension_chart_with_germs_framed_point}
\leanref{Development/AbelianVarietyTheory/Proofs/CuspFormEvaluations/ExistsBiextensionChartWithGermsFramedPoint.lean}{GlobalTheory.ChartRing.exists_chartAt_with_germs}
\begin{proof}
The points. By Lemma~\ref{cv-volume-key-lie}(2), the tangent vectors of $y$ together form a basis of the
tangent vectors at the point under $y$; number them $v_1,\dots,v_d$ through an ordering of the letters.
Theorem~\ref{cv-germs-formal}(1), for the pullback of $\mathcal A$ to $\operatorname{Spec}R$, gives points $x_m$ with
values in $J^1_m$ such that $x_m(0)$ is the unit, $x_{m+1}$ reduces to $x_m$, and $v_l$ is the image of $x_1$ under
$t_j\mapsto\delta_{jl}\varepsilon$. As in the proof of Lemma~\ref{cv-germs-local}, the uniqueness in
Theorem~\ref{cv-germs-formal}(2) gives a commutative formal group law $F$ with $x_m(t+_Ft')=x_m(t)+x_m(t')$ and series
$[a]$ with $a\,x_m(t)=x_m([a]t)$. So condition (a) of Definition~\ref{cv-germs-chart}(3) holds. Take for
$(u_{\ell j})$ the identity matrix, which is invertible, so that the tangent vector of $y$ at the $j$-th letter is
$v_j$.

The generators and the formal germ. Lemma~\ref{cv-uniformization-theta-rule-connection-charts-frames} gives generators
$e_m$, series $\Phi_1$, $\Phi_2$, $S$ and $P_a$ and a generator $w$ that satisfy conditions (c) to (e) of
Definition~\ref{cv-germs-chart}(3), with the trivialization along $x_0$ the rigidification. By
Lemma~\ref{cv-uniformization-theta-rule-connection-charts-identities}, $B=(F,[\,\cdot\,],\Phi_1,\Phi_2,S,P)$ is a formal
germ. So $B$, the points $x_m$ and the generators $e_m$ satisfy conditions (a), (d) and (e) of
Definition~\ref{cv-germs-chart}(3), and by condition (c) $e_{m+1}$ reduces to $e_m$.

The gauge and the germs. As $R$ contains $\mathbf Q$,
Lemma~\ref{cv-uniformization-theta-rule-connection-charts-gauge} gives a power series $h$ over $R$ as in
Definition~\ref{cv-uniformization-theta-rule-values}(5). With $B$, the $x_m$, the $e_m$, the identity matrix, $w$ and
$h$ this is a chart datum at $y$, and by Lemma~\ref{cv-uniformization-theta-rule-connection-charts-germs} every section
of $\mathcal L_y^k$, $k\ge0$, has a germ for it.
\end{proof}
% Lean: T.exists_biextension_chart_with_germs_framed_point = ChartRing.exists_chartAt_with_germs (Laws, star = bar):
% ChartData.exists_isChartG_of_basis (d the number of letters, uvec the identity; Theorem cv-germs-formal and the
% series as in the proof of Lemma cv-germs-local) with FrameTangent.frame_tangent_basis (Lemma cv-volume-key-lie(2)),
% exists_frames_of_chart (Lemma cv-uniformization-theta-rule-connection-charts-frames), FrameBiextId.fbId with
% bijective_appTop_powObj (Lemma cv-uniformization-theta-rule-connection-charts-identities), canonGauge_exists (Lemma
% cv-uniformization-theta-rule-connection-charts-gauge) and germSt_holds (Lemma
% cv-uniformization-theta-rule-connection-charts-germs), with isChart_iff, isFrame_iff, isPrim_iff and isGerm_iff
% between the Lean's two forms of these notions. The reduction of e_(m+1) to e_m in the statement is FrTruncC y C (htr
% of exists_frames_of_chart). Sizes are counted through the theorems a
% declaration reaches in the files it imports, resolved through the namespaces and open lines of its file: about
% 10,100 lines beyond the laws of abelian schemes (Out.Laws, through firstFiveLaws and allSixLaws, about 57,100 lines,
% which the chapter uses throughout) and chart independence (Lemma cv-uniformization-theta-rule-unique-chart), of
% which about 3,220 in the formal chart (exists_isChartG_of_basis and frame_tangent_basis, the statements of Theorem
% cv-germs-formal and Lemma cv-volume-key-lie), 1,700, 1,030 and 2,960 in Lemmas
% cv-uniformization-theta-rule-connection-charts-frames, -identities and -gauge, and 1,230 in the rest.

\begin{lemma}\label{cv-uniformization-theta-rule-connection}
In the setting of Lemma~\ref{cv-uniformization-theta-marking}, there is a canonical Taylor family on $X'$
(Definition~\ref{cv-uniformization-theta-rule-values}(5)).
\end{lemma}
\leanref{Development/AbelianVarietyTheory/Proofs/FormalGroupCharts/ExistsIsNabUncond.lean}{GlobalTheory.CanonConn.exists_isNab_uncond}
\leanref{Development/AbelianVarietyTheory/Proofs/CuspFormEvaluations/ExistsIsNab.lean}{GlobalTheory.CanonConn.exists_isNab}
\begin{proof}
Let $y$ be a framed point of $X'$ over an $F'$-algebra $R$. By
Lemma~\ref{cv-uniformization-theta-rule-connection-charts}, choose a chart datum at $y$ for which every section of
$\mathcal L_y^k$, $k\ge0$, has a germ, and choose a germ $f_s$ of every such section $s$. For $k\ge0$ let
$\nabla_y(s;\alpha)$ be the constant term of $\big(\prod_\ell D_\ell^{\alpha_\ell}\big)\big(f_s(t)\,e^{k\,h(t,[w_0]t)}\big)$,
with $D_\ell$ as in Definition~\ref{cv-uniformization-theta-rule-values}(5), and for $k<0$ put $\nabla_y(s;\alpha)=0$.
A section $s$ has only one germ for the chart datum: the class of a germ modulo $I^{m+1}$ is the value of $s$ at $x_m$
in the $k$-th power of the trivialization, for every $m$, and a power series is determined by these classes. So the
formula holds for every germ of $s$, and $\nabla$ is canonical.
\end{proof}
% Lean: CanonConn.exists_isNab_uncond = CanonConn.exists_isNab with firstFiveLaws.1, for every moduli scheme of a
% datum at a level: the star structure H.starRing (star = bar), C and g from
% T.exists_biextension_chart_with_germs_framed_point (Lemma cv-uniformization-theta-rule-connection-charts), nab the
% nabVal of g for k = Int.ofNat k and 0 for negative k (and 0 over rings that are not Q-algebras, which do not occur
% over F'), and isGermFun_holds (pull_unique, gen_eq_of_corr, value_unique, psJet_inj) for the independence of the
% germ. Size beyond Lemma cv-uniformization-theta-rule-connection-charts: about 90 lines.

\begin{lemma}\label{cv-uniformization-theta-rule-unique-base-change}
In the setting of Lemma~\ref{cv-uniformization-theta-marking}, let $\varphi\colon S\to T$ and $\psi\colon T\to U$ be
homomorphisms of $F'$-algebras and $y$ a framed point of $X'$ over $S$, with images as in
Definition~\ref{cv-uniformization-theta-rule-values}.
\begin{enumerate}
\item The framed point $y$ has an image under $\varphi$.
\item The framed point $y$ is an image of itself under the identity of $S$.
\item If $z$ is an image of $y$ under $\varphi$ and $w$ an image of $z$ under $\psi$, then $w$ is an image of $y$
  under $\psi\circ\varphi$.
\item If $z$ is an image of $y$ under $\varphi$ and $w$ an image of $y$ under $\psi\circ\varphi$, then $w$ is an
  image of $z$ under $\psi$.
\end{enumerate}
\end{lemma}
\leanref{Development/AbelianVarietyTheory/Proofs/CuspFormEvaluations/ExistsFramedPointBaseChangeAlgHomStatement.lean}{GlobalTheory.T.exists_framedPoint_baseChange_algHom}
\leanref{Development/AbelianVarietyTheory/Proofs/CuspFormEvaluations/ExistsFramedPointBaseChangeAlgHom.lean}{GlobalTheory.FramedBaseChange.exists_frMap}
\leanref{Development/AbelianVarietyTheory/Proofs/CuspFormEvaluations/FramedPointBaseChangeLaws.lean}{GlobalTheory.FramedBaseChange.refl}
\leanref{Development/AbelianVarietyTheory/Proofs/CuspFormEvaluations/FramedPointBaseChangeLaws.lean}{GlobalTheory.FramedBaseChange.trans}
\leanref{Development/AbelianVarietyTheory/Proofs/CuspFormEvaluations/FramedPointBaseChangeLaws.lean}{GlobalTheory.FramedBaseChange.cancel}
\begin{proof}
(2) to (4) hold because base change along the identity is the identity, and base change along $\psi\circ\varphi$ is
base change along $\varphi$ followed by base change along $\psi$, for the point, for the tangent vectors and for the
pullback of the generator; in (4) the maps of fibres $A_w\to A_z\to A_y$ compose to the map $A_w\to A_y$, as both
are determined by their two projections.

(1) Let $z$ have the point under $y$ composed with $\operatorname{Spec}\varphi$, the tangent vectors of $y$ base
changed along $\varphi$, and the pullback $\vartheta'$ of $\vartheta_y$ to $A_z=A_y\times_ST$.

The tangent vectors. The base changed vectors are tangent vectors at the point of $z$, and each lies in the
eigenspace of the same embedding as the vector of $y$ it comes from, as the action of $\mathcal O_E$ commutes with
base change. By Lemma~\ref{cv-volume-key-lie}(2) the vectors of $y$ together form a basis of the Lie algebra of $A_y$
over $S$. By part (1) of that lemma their images form a basis of the Lie algebra of $A_z$ over $T$, and by part (2)
again, as they lie in the eigenspaces, the images at each embedding $\tau$ form a basis of the $\tau$-eigenspace.

The generator. The scheme $X'$ is locally noetherian, and its universal abelian scheme with the line bundle of the
polarization satisfies the hypotheses of Theorem~\ref{cv-volume-key-sections}. Let $q$ be a point of
$\operatorname{Spec}T$. By that theorem for $n=1$ there are an open neighbourhood $V$ of the image of $q$ in $X'$ and a
section $b$ of the line bundle over the abelian scheme over $V$ whose pullbacks along all morphisms from affine
schemes to $V$ are generators. For $g$ in a ring $R$ with a point of $X'$, write $D(g)$ for the open set of
$\operatorname{Spec}R$ where $g$ is invertible and $b_g$ for the pullback of $b$ to the abelian scheme over $D(g)$, when
$D(g)$ maps into $V$. The abelian schemes $A_y$ and $A_z$ are proper over affine schemes, hence quasi-compact and
quasi-separated; so a section of the line bundle over the part of $A_y$ over $D(g)$, multiplied by a power of $g$, is
the restriction of a section over $A_y$, and a section over $A_y$ that vanishes over $D(g)$ is killed by a power of
$g$ \cite[Tag 01PW]{Sta}, and the same over $A_z$.

Every section $s$ of $\mathcal L_z$ is a multiple of $\vartheta'$. Choose $f\in S$ such that $D(f)$ contains the image
of $q$ and maps into $V$, and put $f'=\varphi(f)$. Some $f^nb_f$ is the restriction of a section of $\mathcal L_y$,
which is $r\vartheta_y$ for some $r\in S$, as $\vartheta_y$ generates; pulled back, $f'^nb_{f'}=\varphi(r)\vartheta'$
over $D(f')$. Over $D(f')$, $s=d\,b_{f'}$ for a function $d$ on $D(f')$, and $f'^md$ is the restriction of some
$d_0\in T$. So $f'^{m+n}s$ and $d_0\varphi(r)\vartheta'$ agree over $D(f')$, and for some $N$,
$f'^{N+m+n}s=f'^Nd_0\varphi(r)\vartheta'$ over $A_z$. Thus the ideal of the $a\in T$ with $as\in T\vartheta'$ contains
an element outside $q$. As $q$ is arbitrary, this ideal is $T$, and $s\in T\vartheta'$.

The multiple is unique. Let $a\in T$ with $a\vartheta'=0$. Choose $g\in T$ such that $D(g)$ contains $q$ and maps into
$V$. Some $g^nb_g$ is the restriction of a section of $\mathcal L_z$, which is $r'\vartheta'$ with $r'\in T$ by the
previous paragraph. So $ag^nb_g=r'a\vartheta'=0$ over $D(g)$, and as $b_g$ is a generator of a free module of rank one
there, $ag^n=0$ on $D(g)$. So $a$ is zero in the local ring at $q$. As $q$ is arbitrary, $a=0$.
% Lean: FramedBaseChange.exists_frMap (Laws): the tangent vectors spMap (dualMap psi) >> y.u with isTanAt_pull and
% isEigAt_pull; their basis property T.existsUnique_isCombAt_baseChange_of_isEigAt = frame_baseChange_eigen_basis,
% through FrameTangent.frame_tangent_basis (the vectors form a basis of all tangent vectors),
% FrameBaseChange.basis_baseChange (through exists_isLieBasis_baseChange, as Lemma cv-volume-key-lie(1)) and
% FrameBaseChange.eigen_basis (the eigenvector argument of Lemma cv-volume-key-lie(2), with exists_embK_ne); the
% section from GermD.exists_pullSec; and its generation T.framedPoint_generator_ringBaseChange =
% framedPoint_generator_baseChange, which is GenBaseChange.existsUnique_of_isPullSec with the base change law
% hL.secBaseChange (Theorem cv-volume-key-sections): isGenSec_pullS for existence and eq_zero_of_smul for uniqueness,
% with VVRank.exists_basis, EmbPull.exists_pow_extend_secM and EmbPull.exists_pow_kill_secM (Stacks, Tag 01PW).
% (2) to (4) are FramedBaseChange.refl, trans and cancel (pulledGroup_hom_ext and appLE_comp_appLE). Sizes beyond the
% laws of abelian schemes and chart independence: (1) about 3,700 lines, of which about 1,770 for the bases (the
% statements of Lemma cv-volume-key-lie) and 1,850 for the generator; (2) to (4) about 130 more.
\end{proof}

\begin{definition}\label{cv-uniformization-theta-rule-unique-translates-symmetric}
Let $S$ be a scheme, $A$ a commutative group scheme over $S$, $\mathcal L$ a line bundle on $A$, $k$ an integer and
$T$ a scheme over $S$, with $A_T=T\times_SA$, its projections $p_T$ and $p_A$ and the translations $t_x$ of
Definition~\ref{cv-cusp-cm-point-eigen-degree-rank-one-norm-weil-translate-group}, and let $\iota$ be the automorphism
of $A_T$ over $T$ with $p_A\circ\iota=-p_A$. Let $(x,\varphi)$ be an element of the theta group of $\mathcal L^k$ over
$T$.
\begin{enumerate}
\item Let $n\ge1$ with $nx=0$. Then $(x,\varphi)^n=1$ if the generator $\prod_{j=0}^{n-1}t_{jx}^*\varphi$ of
  $p_A^*\mathcal L^k\otimes(nx\,p_T+p_A)^*\mathcal L^{-k}=\mathcal O_{A_T}$ is $1$.
\item The element $(x,\varphi)$ is \emph{symmetric} if for some generator $\nu$ of
  $p_A^*\mathcal L^k\otimes(-p_A)^*\mathcal L^{-k}$ on $A_T$ the generator
  $t_{-x}^*\varphi\otimes\nu\otimes\iota^*\varphi\otimes t_{-x}^*\nu^{-1}$ is $1$. Its line bundle is $\mathcal O_{A_T}$:
  each of $(-x\,p_T+p_A)^*\mathcal L^k$, $p_A^*\mathcal L^k$, $(-p_A)^*\mathcal L^k$ and $(x\,p_T-p_A)^*\mathcal L^k$
  occurs in it once and its inverse once.
\end{enumerate}
\end{definition}
\leanref{ReciprocityTheory/ReductiveGroups/ArithmeticGeometry/ExtensionDescent/ThetaPt.lean}{ReciprocityTheory.ChartRestriction.G.Out.SymLift}
\leanref{ReciprocityTheory/ReductiveGroups/ArithmeticGeometry/ExtensionDescent/ThetaPt.lean}{ReciprocityTheory.ChartRestriction.G.Out.powGen}
\leanref{ReciprocityTheory/ReductiveGroups/ArithmeticGeometry/ExtensionDescent/ThetaPt.lean}{ReciprocityTheory.ChartRestriction.G.Out.symGen}
\leanref{ReciprocityTheory/ReductiveGroups/ArithmeticGeometry/ExtensionDescent/ThetaPt.lean}{ReciprocityTheory.ChartRestriction.G.Out.negLB}
% Lean: SymLift A L k n x theta: pt_eq (theta lies over x), pow_eq (powGen theta n, the product of the pullbacks of
% theta.phi along translT (x^j), j < n, corresponds to oneGen) and sym (some generator nu of negLB A L k T with symGen
% theta nu = theta.phi pulled back along translT x^-1, times nu, times theta.phi pulled back along negT, times the
% inverse of nu pulled back along translT x^-1, corresponding to oneGen). These are Mumford's symmetric elements of
% order dividing n.

\begin{lemma}\label{cv-uniformization-theta-rule-unique-translates-lift}
Let $S$ be a scheme, $A$ an abelian scheme over $S$, $\mathcal L$ a line bundle on $A$, $k$ an integer, $n\ge1$ an odd
integer and $T$ a scheme over $S$, with the notation of
Definition~\ref{cv-uniformization-theta-rule-unique-translates-symmetric}. Let $x\in A(T)$ with $nx=0$, let
$(x,\varphi)$ be an element of the theta group of $\mathcal L^k$ over $T$, and let $\nu$ be a generator of
$p_A^*\mathcal L^k\otimes(-p_A)^*\mathcal L^{-k}$ on $A_T$. Then there is a symmetric element $U$ of the theta group of
$\mathcal L^k$ over $T$ that lies over $x$ and satisfies $U^n=1$.
\end{lemma}
\leanref{Development/AbelianVarietyTheory/Proofs/AbelianSchemeLawProofs/SymmetricLiftOddOrderExistsUnique.lean}{GlobalTheory.SymLiftOdd.exists_symLift}
\leanref{Development/AbelianVarietyTheory/Proofs/AbelianSchemeLawProofs/SymmetricLiftOddOrderExistsUnique.lean}{GlobalTheory.SymLiftOdd.defect_pow}
\begin{proof}
As $A_T$ is an abelian scheme over $T$, its global functions are those of $T$. So a generator of the trivial line
bundle $\mathcal O_{A_T}$ is a unit of $\Gamma(T,\mathcal O_T)$, and it is unchanged by $t_y^*$, for $y\in A(T)$, and
by $\iota^*$.

The constants. Let $a$ be the unit $t_{-x}^*\varphi\otimes\nu\otimes\iota^*\varphi\otimes t_{-x}^*\nu^{-1}$, and $b$
the unit $\prod_{j=0}^{n-1}t_{jx}^*\varphi$.

The relation $a^n=b^2$. Let $j\ge0$. As $t_{-x}\circ t_{(j+1)x}=t_{jx}$ and $\iota\circ t_{(j+1)x}=t_{-(j+1)x}\circ\iota$,
\[
 a=t_{(j+1)x}^*a=t_{jx}^*\varphi\otimes t_{(j+1)x}^*\nu\otimes\iota^*t_{-(j+1)x}^*\varphi\otimes t_{jx}^*\nu^{-1}.
\]
Take the product over $0\le j<n$. The factors $t_{jx}^*\varphi$ give $b$. As $nx=0$, $t_{-(j+1)x}=t_{(n-1-j)x}$, so
the factors $\iota^*t_{-(j+1)x}^*\varphi$ give $\iota^*b=b$. The factors in $\nu$ give $t_{nx}^*\nu\otimes\nu^{-1}=1$.
So $a^n=b^2$.

The scaling. Write $n=2l+1$, put $c=b^{-1}a^l$, a unit of $\Gamma(T,\mathcal O_T)$, and $U=(x,c\varphi)$. As $c$
comes from $T$, $t_y^*(c\varphi)=c\,t_y^*\varphi$ and $\iota^*(c\varphi)=c\,\iota^*\varphi$. So
$\prod_{j<n}t_{jx}^*(c\varphi)=c^nb=b^{1-n}(a^n)^l=b^{1-n+2l}=1$, and
$t_{-x}^*(c\varphi)\otimes\nu\otimes\iota^*(c\varphi)\otimes t_{-x}^*\nu^{-1}=c^2a=b^{-2}a^n=1$. So $U$ is symmetric,
lies over $x$ and satisfies $U^n=1$.
\end{proof}
% Lean: SymLiftOdd.exists_symLift (GlobalFnsSt, IsAbelianScheme A, L.IsLB, n odd, x^n = 1, theta over x, nu a
% generator of negLB): exists_symConst for a (symGen_isGen, exists_constT with surj_fst from the law of global
% functions), exists_constT with form_pow and powGen_isGen for b, defect_pow for a^n = b^2 (a computation in a
% commutative group with the values val_symGen and val_powGen of the generators in charts of L, which is the product
% over j of the text), exists_scale for theta' = (x, c phi) with c = b^-1 a^l, and the two identities for pow_eq and
% sym.

\begin{lemma}\label{cv-uniformization-theta-rule-unique-translates-theta}
Let $X$ be the PEL moduli scheme of a PEL datum at a level (Definition~\ref{cv-model-toroidal-semiabelian}(1)), such
as $X'$ of Lemma~\ref{cv-uniformization-theta-marking} or the scheme $X$ of
Definition~\ref{cv-volume-twisting-section}. Let $\mathcal A$ be the abelian scheme of its universal object, with the
line bundle $\mathcal L$ of the polarization and its rigidification along the zero section; for $X'$ these are those of
Definition~\ref{cv-uniformization-theta-rule-values}. Let $N\ge1$ be odd.
\begin{enumerate}
\item There is a unique isomorphism $[N]^*\mathcal L\cong\mathcal L^{N^2}$ on $\mathcal A$ compatible with the
  rigidifications along the zero section.
\item Let $T$ be a scheme over $X$, $\mathcal A_T=\mathcal A\times_XT$, and $t'$ a point of $\mathcal A$ over $T$
  with $N^2t'=0$. Then there is a symmetric element $U$ of the theta group of $\mathcal L^{N^2}$ over $T$ that lies
  over $t'$ and satisfies $U^{N^2}=1$ (Definition~\ref{cv-uniformization-theta-rule-unique-translates-symmetric}), and
  any two such elements are equal.
\end{enumerate}
\end{lemma}
\leanref{Development/AbelianVarietyTheory/Proofs/CuspFormEvaluations/UniversalNormalisedMulByNatPullbackIsoStatement.lean}{GlobalTheory.T.universal_normalised_mulByNat_pullback_iso}
\leanref{Development/AbelianVarietyTheory/Proofs/CuspFormEvaluations/UniversalSymLiftExistsUniqueGlobalStatement.lean}{GlobalTheory.T.universal_symLift_exists_unique_global}
\leanref{Development/AbelianVarietyTheory/Proofs/AbelianSchemeLawProofs/SymmetricLiftOddOrderExistsUnique.lean}{GlobalTheory.SymLiftOdd.symLift_unique}
\begin{proof}
The scheme $X$ lies over a field of characteristic $0$, so $N$ is invertible on it. For every scheme $T$ over $X$ the
global functions on the abelian scheme $\mathcal A_T$ are those of $T$.

(1) As $\mathcal L$ is symmetric and trivialized along the zero section, the theorem of the cube gives an isomorphism
$[N]^*\mathcal L\cong\mathcal L^{N^2}$ on $\mathcal A$ \cite{Mum70}. Dividing it by its value along the zero section,
read in the rigidifications, makes it compatible with them. Two compatible isomorphisms differ by an invertible
function on $\mathcal A$, which comes from $X$ and is $1$ along the zero section; so they are equal.

(2) Put $k=N^2$, and use the notation of
Definition~\ref{cv-uniformization-theta-rule-unique-translates-symmetric} for $\mathcal A$ over $X$.

Uniqueness. Two elements of the theta group over $t'$ differ by an invertible function on $\mathcal A_T$, which is a
unit $c$ of $T$. The generator of Definition~\ref{cv-uniformization-theta-rule-unique-translates-symmetric}(2) is
multiplied by $c^2$ when $\varphi$ is multiplied by $c$, and it is unchanged when $\nu$ is multiplied by a unit of
$T$, which is invariant under $t_{-x}^*$; and any two generators $\nu$ differ by such a unit. So if both elements are
symmetric, then $c^2=1$, and if both have $k$-th power $1$, then $c^k=1$. As $N$ is odd, $c=c^k(c^2)^{-(k-1)/2}=1$.

Existence. The generator $\nu$. As $\mathcal L$ is symmetric, every point of $X$ has an open neighbourhood $V$ over
which $[-1]^*\mathcal L\otimes\mathcal L^{-1}$ has a generator on $\mathcal A_V$. Its pullback along the zero section is
a generator of $0^*\mathcal L\otimes0^*\mathcal L^{-1}=\mathcal O_V$; dividing by it, it becomes $1$ along the zero
section. Two such generators over $V$ and $V'$ differ over $V\cap V'$ by an invertible function on
$\mathcal A_{V\cap V'}$, which comes from $V\cap V'$ and is $1$ along the zero section; so they agree, and they glue
to a generator $u$ of $[-1]^*\mathcal L\otimes\mathcal L^{-1}$ on $\mathcal A$. Its pullback along $p_A$, to the power
$-k$, is a generator $\nu$ of $p_A^*\mathcal L^k\otimes(-p_A)^*\mathcal L^{-k}$ on $\mathcal A_T$.

Elements of the theta group locally on $T$. Let $\theta$ be the cubical structure of $\mathcal L$
(Definition~\ref{cv-germs-chart}(1)), and let $x,z\in\mathcal A(T)$. The pullback of $\theta$ along
$(x\,p_T,z\,p_T,p_A)$ is a generator, on $\mathcal A_T$, of
\[
 \big((x+z)\,p_T+p_A\big)^*\mathcal L\otimes p_A^*\mathcal L\otimes(x\,p_T+p_A)^*\mathcal L^{-1}\otimes
 (z\,p_T+p_A)^*\mathcal L^{-1}
\]
tensored with the pullback along $p_T$ of $(x+z)^*\mathcal L^{-1}\otimes x^*\mathcal L\otimes z^*\mathcal L\otimes
0^*\mathcal L^{-1}$ on $T$, which has a generator near each point of $T$. So every point of $T$ has an open
neighbourhood over which the first line bundle has a generator. Applying this to $(jt',t')$ for $j=0,\dots,k-1$ and
multiplying, every point of $T$ has an open neighbourhood $V$ over which
$(kt'\,p_T+p_A)^*\mathcal L\otimes p_A^*\mathcal L^{k-1}\otimes(t'p_T+p_A)^*\mathcal L^{-k}$ has a generator; as
$kt'=0$, this is $p_A^*\mathcal L^k\otimes(t'p_T+p_A)^*\mathcal L^{-k}$, and the generator is an element of the theta
group of $\mathcal L^k$ over $V$ that lies over $t'$.

Gluing. By Lemma~\ref{cv-uniformization-theta-rule-unique-translates-lift}, applied over each such $V$ with the
restriction of $\nu$, there is a symmetric element $U_V$ of the theta group of $\mathcal L^k$ over $V$ that lies over
$t'$ and satisfies $U_V^k=1$. The restrictions of $U_V$ and $U_{V'}$ to $V\cap V'$ are two such elements over
$V\cap V'$, so they are equal by the uniqueness; hence the $U_V$ glue to an element $U_0$ of the theta group of
$\mathcal L^k$ over $T$ that lies over $t'$. Lemma~\ref{cv-uniformization-theta-rule-unique-translates-lift}, applied
over $T$ to $U_0$ and $\nu$, gives a symmetric element $U$ over $t'$ with $U^k=1$.
\end{proof}
% Lean: (1) is T.universal_normalised_mulByNat_pullback_iso = moduli_univ_normMul: the law mulPull for a generator,
% NormMulUniv.exists_normalised (the rigidification), NormMulUniv.normalised_unique (the law of global functions) and
% NormMulUniv.exists_pull (the pullbacks to every T, which the text uses through the pullback). (2) is
% T.universal_symLift_exists_unique_global = moduli_univ_symLift_global: symLift_exists_unique_of_cast with N^2 a unit
% on X (moduli_univ_polAV); existence exists_symLift_global (the law symLift for local lifts,
% SymLiftGlobal.exists_lift_of_local for the gluing, through tval_agree and symLift_unique, exists_gen_negLB and
% SymLiftOdd.exists_symLift over T), uniqueness SymLiftOdd.symLift_unique (exists_ratio for phi and for nu, then c^2 = 1
% from sym and c^n = 1 from pow_eq). The law symLift is proved in the archive as AbLaw.FromCube.symLift =
% T.symmetric_lift_odd_order_exists_unique (symLift_odd), from the local elements of the theta group of
% kerPow_of_square (AbSquare.square_of_cube: square_local, the cube generator pulled back along (x, z, p_A) times a
% local generator on T; KerPow.locGen_pow by induction, locGen_theta_iff and exists_theta), exists_gen_negLB
% (exists_gen_neg_of_sym from the symmetry P.sym, glued by exists_gen_of_local_on_base, and exists_gen_zsmul) and
% SymLiftOdd.exists_symLift (Lemma cv-uniformization-theta-rule-unique-translates-lift); the text follows this proof.
% It does not use that N is invertible, which the law asks. The Lean states (1) and (2) for every moduli scheme of a
% datum at a level, as the text does. Size of (1) beyond the laws of abelian schemes: about 80 lines.

\begin{lemma}\label{cv-uniformization-theta-rule-unique-translates-rosati}
In the setting of Lemma~\ref{cv-uniformization-theta-marking}, let $N\ge1$ be odd, and let $f$ and $g$ be
endomorphisms of the group scheme $\mathcal A$ over $X'$ that are adjoint for $\mathcal L$: locally on $X'$ the line
bundles $\Lambda(f(y_1),y_2)$ and $\Lambda(y_1,g(y_2))$ on $\mathcal A\times_{X'}\mathcal A$ are isomorphic, where
$y_1,y_2$ are the two projections and $\Lambda$ is as in Definition~\ref{cv-germs-chart}(1). Let $T$ be a scheme over
$X'$, $s$ a section of the pullback of $\mathcal L$ to $\mathcal A_T$ and $t\in\mathcal A(T)$, and let $F$ and $F'$ be
$N$-translates of $s$ over $t$, defined as in Definition~\ref{cv-uniformization-theta-rule-values}(1) with $s$ in place
of $\vartheta_y$, whose roots $x$ satisfy $f(x)=0$, $g(x)=x$ and $N^2x=0$. Then $F=F'$.
\end{lemma}
\leanref{ReciprocityTheory/ReductiveGroups/ArithmeticGeometry/BirationalGeometry/ToroidalCompactification/VtUniq.lean}{ReciprocityTheory.ChartRestriction.G.Geo.Tor.vtOver_unique}
\begin{proof}
\openstep{The proof that two such translates are equal, by comparing the two symmetric elements of the theta group
through the adjunction of $f$ and $g$, is not yet written out.}
% Lean: vtOver_unique (GlobalFnsSt, IsAbelianScheme, L symmetric, N odd, RosatiAdj A L f g): the two normalized
% isomorphisms agree (gen_unique), D = U'.pt^-1 * U.pt with D^N = 1 and D g = D, ratio_const for the ratios of
% generators (constant by the law of global functions), val_theta, exists_ros and ros_pow' for the adjunction, sym_val
% for the symmetry, and sec_ext_val to compare F and F' through their values. Size beyond the laws of abelian
% schemes, chart independence and the lemmas before this one from Lemma cv-uniformization-theta-rule-connection-charts
% on: about 510 lines.
\end{proof}

\begin{lemma}\label{cv-uniformization-theta-rule-unique-translates}
In the setting of Lemma~\ref{cv-uniformization-theta-marking}, let $y$ be a framed point of $X'$ over an
$F'$-algebra $S$, with translates as in Definition~\ref{cv-uniformization-theta-rule-values}(1).
\begin{enumerate}
\item Let $N\ge1$ be odd and $t,t'\in A_y(S)$ with $Nt'=t$ and $N^2t'=0$. Then there is an $N$-translate of
  $\vartheta_y$ over $t$ with root $t'$.
\item Two level translates of $\vartheta_y$ over the same point $t\in A_y(S)$ are equal.
\item Let $z$ be an image of $y$ under a homomorphism of $F'$-algebras $S\to T$, and let $F$ be a level translate of
  $\vartheta_y$ over $t\in A_y(S)$. Then the pullback of $F$ to $A_z$ is a level translate of $\vartheta_z$ over the
  pullback of $t$.
\end{enumerate}
\end{lemma}
\leanref{Development/AbelianVarietyTheory/Proofs/CuspFormEvaluations/ExistsTranslationSectionsUniversalObject.lean}{GlobalTheory.exists_isVtOver_and_existsUnique_isUUOver}
\leanref{ReciprocityTheory/ReductiveGroups/ArithmeticGeometry/BirationalGeometry/ToroidalCompactification/VtUniq.lean}{ReciprocityTheory.ChartRestriction.G.Geo.Tor.vtUniq_obj}
\leanref{Development/AbelianVarietyTheory/Proofs/CuspFormEvaluations/IsVtOverIsUUOverBaseChange.lean}{GlobalTheory.VtBaseChange.isVtOver_isUUOver_baseChange}
\begin{proof}
(1) Let $\alpha$ be the pullback to $A_y$ of the isomorphism of
Lemma~\ref{cv-uniformization-theta-rule-unique-translates-theta}(1), which is compatible with the rigidifications, and
$s_1$ the section of $\mathcal L_y^{N^2}$ that corresponds to $[N]^*\vartheta_y$ under it. By part (2) of that lemma,
for $T=\operatorname{Spec}S$, there is a symmetric element $U$ of the theta group of $\mathcal L_y^{N^2}$ over $t'$ with
$U^{N^2}=1$. Then $Us_1$ is an $N$-translate of $\vartheta_y$ over $t$ with root $t'$.

(2) As $\mathfrak M'+\bar{\mathfrak M}'=\mathcal O_E$, also $\mathfrak M'^2+\bar{\mathfrak M}'^2=\mathcal O_E$; write
$1=p+\bar a$ with $p,a\in\mathfrak M'^2$. Complex conjugation on $\mathcal O_E$ is the Rosati involution of $\mathcal L$
(Definition~\ref{cv-model-toroidal-semiabelian}(2)), so $[a]$ and $[\bar a]$ are adjoint for $\mathcal L$ in the sense
of Lemma~\ref{cv-uniformization-theta-rule-unique-translates-rosati}. A root $x$ of a level translate is killed by
$\mathfrak M'^2$; so $ax=0$, $\bar ax=x-px=x$, and $M'^2x=0$, as $M'^2\in\mathfrak M'^2$. That lemma, for the odd
integer $N=M'$, with $f=[a]$, $g=[\bar a]$, $T=\operatorname{Spec}S$ and $s=\vartheta_y$, gives the claim.

(3) Let $F=Us_1$ with root $t'$, as in Definition~\ref{cv-uniformization-theta-rule-values}(1), and let $\vartheta_z$,
$t'_z$ and $U_z$ be the pullbacks to $A_z=A_y\times_ST$ of $\vartheta_y$, $t'$ and $U$. Pulling back commutes with
$[M']^*$ and carries the pullback of the isomorphism of
Lemma~\ref{cv-uniformization-theta-rule-unique-translates-theta}(1), for $N=M'$, to $A_y$ to its pullback to $A_z$; so
it carries $s_1$ to the section that corresponds to $[M']^*\vartheta_z$. It carries $U$ to an element $U_z$ of the
theta group over $t'_z$, symmetric and with $U_z^{M'^2}=1$, and $Us_1$ to $U_z$ applied to the pullback of $s_1$. Finally $M't'_z$ is
the pullback of $t$, and $t'_z$ is killed by $\mathfrak M'^2$, as $t'$ is.
\end{proof}
% Lean: (1) is the first part of exists_isVtOver_and_existsUnique_isUUOver (Laws): alpha0 from
% T.universal_normalised_mulByNat_pullback_iso (Lemma cv-uniformization-theta-rule-unique-translates-theta(1)) with its
% pullback alpha, s_1 from existsUnique_isTwistPull, U from T.universal_symLift_exists_unique_global (part (2) of that
% lemma) and F from existsUnique_isTwistPull for U.phi (IsActOf). (2) is vtUniq_obj (GlobalFnsSt): cop
% (Ideal.sup_pow_eq_top' and the coprimality of Mb and its conjugate), p and a with p + conj(a) = 1, o.t.rosatiAdj a,
% and vtOver_unique (Lemma cv-uniformization-theta-rule-unique-translates-rosati) with f = [a], g = [conj a]. (3) is
% the first part of isVtOver_isUUOver_baseChange: exists_mulTwist for alpha and s_1,
% T.exists_symLift_pullback_isActOf (ThetaPullback.exists_thetaGrp_pullback) for U and the action, and KilledBy for
% the pulled back root. The Lean states (1) for any condition on the root and any section of L in place of the
% generator, (2) for any section of L, and (3) for every N, every ideal in place of Mb^2 and every morphism over X'.
% Sizes beyond the laws of abelian schemes, chart independence and the lemmas before this one from Lemma
% cv-uniformization-theta-rule-connection-charts on: (1) about 290 lines beyond Lemma
% cv-uniformization-theta-rule-unique-translates-theta, (2) about 60 beyond Lemma
% cv-uniformization-theta-rule-unique-translates-rosati, and (3) about 290.

\begin{lemma}\label{cv-uniformization-theta-rule-unique-chart}
In the setting of Lemma~\ref{cv-uniformization-theta-marking}, let $y$ be a framed point of $X'$ over an $F'$-algebra
$S$, and let $C$ and $C'$ be chart data at $y$ (Definition~\ref{cv-uniformization-theta-rule-values}(5)). For a chart
datum $C$ with power series $h$, an integer $k\ge0$, a power series $f$ over $S$ and a family $\alpha$ of nonnegative
integers indexed by the letters, let $V_C(f;\alpha)$ be the constant term of
$\big(\prod_\ell D_\ell^{\alpha_\ell}\big)\big(f(t)\,e^{k\,h(t,[w_0]t)}\big)$, with $D_\ell$ as in
Definition~\ref{cv-uniformization-theta-rule-values}(5) for $C$. Let $k,n\ge0$, let $F_1,\dots,F_n$ be sections of
$\mathcal L_y^k$, and for each $b$ let $f_b$ be a germ of $F_b$ for $C$ and $f'_b$ one for $C'$. Let $a$ be a family of
nonnegative integers indexed by the letters and $P$ a polynomial harmonic of row degrees $a$
(Definition~\ref{cv-uniformization-theta-rule-values}(6)). Then, with the notation of
Definition~\ref{cv-uniformization-theta-rule-values}(3),
\[
 \sum_mP_m\prod_{b=1}^nV_C(f_b;m_b)=\sum_mP_m\prod_{b=1}^nV_{C'}(f'_b;m_b).
\]
\end{lemma}
\leanref{Development/AbelianVarietyTheory/Proofs/FormalGroupCharts/ChartIndependenceUnconditional.lean}{GlobalTheory.chartIndependence_unconditional}
\leanref{Formalism/SharpRestriction/Divisibility/GeneratorPowers/ThExprUniqOfChartIndep2.lean}{Formalism.MainNotions.ModuliIntegralCharts.H1.ChartIndepSt2}
\leanref{Formalism/SharpRestriction/Divisibility/GeneratorPowers/CIndTop.lean}{Formalism.MainNotions.ModuliIntegralCharts.H1.chartIndep2_of_glob}
\begin{proof}
\openstep{The proof that these harmonic combinations of Taylor values do not depend on the chart datum is not yet
written out.}
% Lean: chartIndependence_unconditional = chartIndep2_of_glob with the law of global functions
% (lawsFieldsSeparately, GlobalFnsSt), through chartIndep2_of_gauge and chartIndep2_of_fix: ChartIndepTruncSt
% (chartIndepTrunc_of), FrTruncFixSt (frTruncFix_holds), GermFixSt (germFix_holds) and HarmInvSt (harmInv_of). The
% truncated statement chartIndepTrunc_of passes to exponential charts and gauges: for each chart datum an exponential
% chart (ExpChartOfSt, expChartOf_of with expActOf_of and linEigSt_holds) and an exponential gauge (ExpGaugeOfSt,
% ExpGauge.expGaugeOf_of), the Taylor values as values in the exponential chart (TieSt, tie_holds), the uniqueness of
% the exponential chart given the tangent vectors (ExpUniqChart.expUniqTan), the germs for two exponential gauges
% differing by the factor exp(-k q) (gaugeUniq_of_global, with the law of global functions), and the harmonic
% polynomial killing that factor (KillAddValSt, killAddVal_of with psiCover_holds). Sizes, counted as in Lemma
% cv-uniformization-theta-rule-connection-charts, each beyond the laws of abelian schemes and the ones before it:
% linEigSt_holds 1,610, expActOf_of 2,190, expChartOf_of 1,180, expGaugeOf_of 3,200, tie_holds 1,150, expUniqTan
% 1,800, gaugeUniq_of_global 1,340, killAddVal_of 1,440, chartIndepTrunc_of 60, frTruncFix_holds 3,390,
% germFix_holds 780, harmInv_of 640, chartIndep2_of_fix 140 and the rest 160: about 19,100 lines in all.
\end{proof}

\begin{lemma}\label{cv-uniformization-theta-rule-unique-taylor-chart}
In the setting of Lemma~\ref{cv-uniformization-theta-marking}, let $z$ be an image of a framed point $y$ over an
$F'$-algebra $S$ under a homomorphism of $F'$-algebras $\varphi\colon S\to T$, and let $C$ be a chart datum at $y$
(Definition~\ref{cv-uniformization-theta-rule-values}(5)), with formal germ $B$, points $x_m$, generators $e_m$, matrix
$(u_{\ell j})$, generator $w$ and power series $h$. For a power series $f$ over $S$ write $\varphi(f)$ for $f$ with
$\varphi$ applied to the coefficients.
\begin{enumerate}
\item There is a chart datum $C'$ at $z$ with the formal germ $B\otimes T$ (Definition~\ref{cv-hodge-lattice-germ}(3)),
  the points $x_m\circ\operatorname{Spec}(\varphi_m)$, where $\varphi_m\colon J^1_m\to T\otimes_SJ^1_m$ is the base
  change, the pullbacks of the $e_m$ as generators, the matrix $(\varphi(u_{\ell j}))$, the generator $w$ and the power
  series $\varphi(h)$, such that for every $k\ge0$, every section $F$ of $\mathcal L_y^k$ and every germ $f$ of $F$ for
  $C$, $\varphi(f)$ is a germ for $C'$ of the pullback of $F$ to $A_z$.
\item With $V_C$ and $V_{C'}$ as in Lemma~\ref{cv-uniformization-theta-rule-unique-chart}, for every $k\ge0$, every power
  series $f$ over $S$ and every family $\alpha$ of nonnegative integers indexed by the letters,
  $V_{C'}(\varphi(f);\alpha)=\varphi\big(V_C(f;\alpha)\big)$.
\end{enumerate}
\end{lemma}
\leanref{Development/AbelianVarietyTheory/Proofs/CuspFormEvaluations/RawChartFormalBiextensionBaseChangeStatement.lean}{GlobalTheory.T.rawChart_formalBiextension_baseChange}
\leanref{Development/AbelianVarietyTheory/Proofs/CuspFormEvaluations/RawChartFormalBiextensionBaseChange.lean}{GlobalTheory.rawChart_baseChange}
\leanref{Formalism/SharpRestriction/Divisibility/GeneratorPowers/CIndBCNab.lean}{Formalism.MainNotions.ModuliIntegralCharts.H1.nabMap_holds}
\begin{proof}
(1) Each datum of $C$ is carried to $T$ by base change along $\operatorname{Spec}\varphi$, and each condition on it
is preserved. The conditions of Definition~\ref{cv-germs-chart}(3)(a) on the $x_m$ are equalities of points with
values in the rings $J^k_m$, and substituting power series commutes with applying $\varphi$ to their coefficients; so
the points $x'_m=x_m\circ\operatorname{Spec}(\varphi_m)$ satisfy them for the series of $B\otimes T$. The image of $x'_1$
under $t_j\mapsto\delta_{jl}\varepsilon$ is the base change $v'_l$ of $v_l$. Base change carries a combination of
tangent vectors to the combination of their base changes with the coefficients mapped by $\varphi$, so the tangent
vector of $z$ at a letter $\ell$, the base change of that of $y$, is $\sum_j\varphi(u_{\ell j})v'_j$. The matrix
$(\varphi(u_{\ell j}))$ is invertible: the inverse of the bijection $a\mapsto(\sum_\ell a_\ell u_{\ell j})_j$ is
$S$-linear, given by a matrix $(n_{j\ell})$ with $\sum_jn_{j\ell}u_{\ell'j}=\delta_{\ell\ell'}$ and
$\sum_\ell u_{\ell j}n_{j'\ell}=\delta_{jj'}$, and these identities hold after applying $\varphi$. The pullbacks $e'_m$
of the $e_m$ are generators, and pulling back preserves products of generators and the identities between them, so
the $e'_m$ satisfy conditions (d) and (e) of Definition~\ref{cv-germs-chart}(3) with the series of $B\otimes T$; the
generator $w$ lives on $\mathcal A$ and is
unchanged. The conditions on $h$ are identities of power series and hold after applying $\varphi$. Finally, the
trivialization $w(x'_m(t))\,e'_m(t,[w_0]t)$ is the pullback of $w(x_m(t))\,e_m(t,[w_0]t)$; so for $m=0$ it is the
rigidification, and the value of the pullback $F'$ of $F$ at $x'_m$ in its $k$-th power is the image under $\varphi$ of
the value of $F$ at $x_m$, which is the class of $\varphi(f)$ modulo $I^{m+1}$.

(2) The derivations $D_\ell$ of $C'$ are those of $C$ with $\varphi$ applied to the coefficients, as the group law of
$B\otimes T$ and the matrix of $C'$ are those of $C$ mapped by $\varphi$. As $h(t,[w_0]t)$ has no constant term,
$\varphi$ commutes with the exponential series of $k\,h(t,[w_0]t)$, and it commutes with products and with taking the
constant term.
\end{proof}
% Lean: T.rawChart_formalBiextension_baseChange = rawChart_baseChange (RawBCSt, for the raw chart ChartAt.toRaw C):
% B' = FBMap.mapFB psi C.B, x' m = rjBC psi x0 d m 1 >> C.x m, e' from T.exists_rawChart_frames_baseChange
% (RawFrameBC.exists_frames: the pulled back generators), T.isRawChart_baseChange_jetPoints
% (RawChartBC.rIsChart_baseChange: sq_aug, sq_trunc, sq_subst and sq_copy for the conditions on x', isCombAt_map for
% the tangent vectors, basis_map for the matrix), T.isFrame_rawChart_baseChange, RawChartBaseChange.isCanonGauge_mapFB
% and T.rawChart_baseChange_rig_germ_alt (the rigidification and the germs, through isGerm_iff_raw). The invertibility
% of the mapped matrix is basis_map (the inverse matrix P with P * M = 1 and M * P = 1, mapped). (2) is nabMap_holds
% (NabMapSt, with
% map_DpolyOf, map_expS, map_hL); its hypothesis that h is a canonical gauge enters only through const_hL, h(t,[w0]t)
% without constant term. Size beyond the laws and chart independence: about 1,480 lines.

\begin{lemma}\label{cv-uniformization-theta-rule-unique-taylor}
In the setting of Lemma~\ref{cv-uniformization-theta-marking}, let $\nabla$ be a canonical Taylor family, let
$n,N\ge0$, let $a$ be a family of nonnegative integers indexed by the letters, and let $P$ be a polynomial in the
letter variables with coefficients in $F'$ that satisfies conditions (a) to (c) of
Definition~\ref{cv-uniformization-theta-rule-values}(6) for $a$. Let $z$ be an image of a framed point $y$ over an
$F'$-algebra $S$ under a homomorphism of $F'$-algebras $\varphi\colon S\to T$, let $F_1,\dots,F_n$ be sections of
$\mathcal L_y^{N^2}$, and let $F'_b$ be the pullback of $F_b$ to $A_z$. Then, with the notation of
Definition~\ref{cv-uniformization-theta-rule-values}(3),
\[
 \varphi\Big(\sum_mP_m\prod_{b=1}^n\nabla_y(F_b;m_b)\Big)=\sum_mP_m\prod_{b=1}^n\nabla_z(F'_b;m_b).
\]
\end{lemma}
\leanref{Development/AbelianVarietyTheory/Proofs/CuspFormEvaluations/HarmonicValueBaseChangeStatement.lean}{GlobalTheory.T.harmonicValue_baseChange_alt}
\leanref{Development/AbelianVarietyTheory/Proofs/CuspFormEvaluations/HarmonicValueBaseChange.lean}{GlobalTheory.harmonicValue_baseChange}
\leanref{Development/AbelianVarietyTheory/Proofs/CuspFormEvaluations/HarmonicValueBaseChange.lean}{GlobalTheory.HarmonicValueBaseChange.sum_map_of_harmA}
\begin{proof}
Put $k=N^2$, and let $V_C$ be as in Lemma~\ref{cv-uniformization-theta-rule-unique-chart} for a chart datum $C$. As
$F'$ is a number field, some integer $D\ge1$ makes the coefficients of $DP$ algebraic integers; then $DP$ is harmonic
of row degrees $a$, with the monomials of $P$. As $\nabla$ is canonical, there are a chart datum $C$ at $y$ and a chart
datum $C_2$ at $z$ for which every section has a germ, with $\nabla_y(F;\alpha)=V_C(f;\alpha)$ and
$\nabla_z(F';\alpha)=V_{C_2}(g;\alpha)$ for germs $f$ of $F$ for $C$ and $g$ of $F'$ for $C_2$. Choose germs $f_b$ of $F_b$
for $C$ and $g_b$ of $F'_b$ for $C_2$.

Let $C'$ be the chart datum at $z$ of Lemma~\ref{cv-uniformization-theta-rule-unique-taylor-chart}(1); then
$\varphi(f_b)$ is a germ of $F'_b$ for $C'$. Lemma~\ref{cv-uniformization-theta-rule-unique-chart} for $z$, the chart
data $C'$ and $C_2$, these germs and the polynomial $DP$ gives, after cancelling $D$, which is a unit in $T$,
\[
 \sum_mP_m\prod_bV_{C'}(\varphi(f_b);m_b)=\sum_mP_m\prod_bV_{C_2}(g_b;m_b)=\sum_mP_m\prod_b\nabla_z(F'_b;m_b).
\]
By Lemma~\ref{cv-uniformization-theta-rule-unique-taylor-chart}(2), and as $\varphi$ is a homomorphism of
$F'$-algebras, the left side is
$\varphi\big(\sum_mP_m\prod_bV_C(f_b;m_b)\big)=\varphi\big(\sum_mP_m\prod_b\nabla_y(F_b;m_b)\big)$.
\end{proof}
% Lean: T.harmonicValue_baseChange_alt = harmonicValue_baseChange, with the hypothesis ChartIndepSt2 (Lemma
% cv-uniformization-theta-rule-unique-chart): D from exists_integral_multiples with D . P in harmA, then
% HarmonicValueBaseChange.sum_map_of_harmA: C from IsNab (hnabC S y), C' from
% T.rawChart_formalBiextension_baseChange (Lemma cv-uniformization-theta-rule-unique-taylor-chart(1)), C2 from IsNab at
% the image, ChartIndepSt2 for C' and C2, and nabMap_holds (Lemma cv-uniformization-theta-rule-unique-taylor-chart(2));
% then the cancellation of D. Size beyond the laws, chart independence and the lemmas before this one from Lemma
% cv-uniformization-theta-rule-connection-charts on: about 1,600 lines, of which about 1,480 in Lemma
% cv-uniformization-theta-rule-unique-taylor-chart.

\begin{lemma}\label{cv-uniformization-theta-rule-unique-compare}
In the setting of Lemma~\ref{cv-uniformization-theta-marking}, let $\nabla$ be a canonical Taylor family, $n\ge0$,
$c_1,\dots,c_n\in L$, $a$ a family of nonnegative integers indexed by the letters and $P$ a polynomial harmonic of
row degrees $a$ (Definition~\ref{cv-uniformization-theta-rule-values}(6)). Let $y$ be a framed point of $X'$ over an
$F'$-algebra $S$. For $i=1,2$ let $\varphi_i\colon S\to S_i$ be a homomorphism of $F'$-algebras, $y_i$ an image of $y$
under $\varphi_i$, and $F_{i,b}$, $1\le b\le n$, level translates of $\vartheta_{y_i}$ over $\eta_{y_i}(c_b)$. Let
$\psi_i\colon S_i\to T$ be homomorphisms of $F'$-algebras with $\psi_1\circ\varphi_1=\psi_2\circ\varphi_2$. Then
\[
 \psi_1\big(\nabla_{y_1}(P;F_{1,1},\dots,F_{1,n})\big)=\psi_2\big(\nabla_{y_2}(P;F_{2,1},\dots,F_{2,n})\big).
\]
\end{lemma}
\leanref{Development/AbelianVarietyTheory/Proofs/CuspFormEvaluations/ThetaExpressionEqOfWitnessesStatement.lean}{GlobalTheory.T.thetaExpression_eq_of_witnesses}
\leanref{Development/AbelianVarietyTheory/Proofs/CuspFormEvaluations/ThetaExpressionEqOfWitnesses.lean}{GlobalTheory.ThetaCompare.thExpr_eq_of_witnesses}
\begin{proof}
By Lemma~\ref{cv-uniformization-theta-rule-unique-base-change}(1), $y$ has an image $w$ under
$\psi_1\circ\varphi_1=\psi_2\circ\varphi_2$, and by part (4) of that lemma $w$ is an image of $y_i$ under $\psi_i$, for
$i=1,2$. Let $G_{i,b}$ be the pullback of $F_{i,b}$ to $A_w$. The pullback of $\eta_{y_i}(c_b)$ to $A_w$ is
$\eta_w(c_b)$, as both are given by $\eta(c_b)$, so $G_{i,b}$ is a level translate of $\vartheta_w$ over
$\eta_w(c_b)$ by Lemma~\ref{cv-uniformization-theta-rule-unique-translates}(3). By
Lemma~\ref{cv-uniformization-theta-rule-unique-taylor}, with $N=M'$, and as $\psi_i$ is $F'$-linear,
$\psi_i\big(\nabla_{y_i}(P;F_{i,1},\dots,F_{i,n})\big)=\nabla_w(P;G_{i,1},\dots,G_{i,n})$. By
Lemma~\ref{cv-uniformization-theta-rule-unique-translates}(2), $G_{1,b}=G_{2,b}$ for every $b$, which gives the
claim.
\end{proof}
% Lean: thExpr_eq_of_witnesses: w from T.exists_framedPoint_baseChange_algHom, FramedBaseChange.cancel for hw1 and
% hw2, exists_pull for G1 and G2 (GermD.exists_pullSec, VtBaseChange.isVtOver_isUUOver_baseChange with the level
% point h >> levPt = levPt, and T.harmonicValue_baseChange_alt with map_mul and AlgHom.commutes for the factor
% M'^(-deg P); this step asks only conditions (a) to (c), hP.1), and vtUniq_obj for G1 = G2. Size beyond the laws,
% chart independence and the lemmas before this one from Lemma cv-uniformization-theta-rule-connection-charts on:
% about 80 lines.

\begin{lemma}\label{cv-uniformization-theta-rule-unique-cover}
In the setting of Lemma~\ref{cv-uniformization-theta-marking}, let $y$ be a framed point of $X'$ over an
$F'$-algebra $S$, $n\ge0$ and $c_1,\dots,c_n\in L$. Then there are a faithfully flat homomorphism of $F'$-algebras
$\varphi\colon S\to S'$, an image $y'$ of $y$ under $\varphi$ and level translates $F_b$ of $\vartheta_{y'}$ over
$\eta_{y'}(c_b)$, $1\le b\le n$ (Definition~\ref{cv-uniformization-theta-rule-values}).
\end{lemma}
\leanref{Development/AbelianVarietyTheory/Proofs/CuspFormEvaluations/ExistsFaithfullyFlatCoverLevelSectionsStatement.lean}{GlobalTheory.T.exists_faithfullyFlat_cover_level_sections}
\leanref{Development/AbelianVarietyTheory/Proofs/CuspFormEvaluations/ExistsFaithfullyFlatCoverLevelSections.lean}{GlobalTheory.exists_faithfullyFlat_isVtOver}
\leanref{Development/AbelianVarietyTheory/Proofs/CuspFormEvaluations/ExistsFaithfullyFlatRoots.lean}{GlobalTheory.exists_faithfullyFlat_roots}
\leanref{Development/AbelianVarietyTheory/Proofs/CuspFormEvaluations/ExistsPowMulCoprimeLe.lean}{GlobalTheory.exists_root_killedBy_pow}
\begin{proof}
The level points. The level structure $\eta$ is additive and vanishes on $\mathfrak M'L$
(Definition~\ref{cv-toroidal-moduli-object}(4)), and it is $\mathcal O_E$-linear: writing $1=u+v$ with
$u\in\mathfrak M'$ and $v\in\bar{\mathfrak M}'$, $\eta(c)=\eta(vc)$ for every $c\in L$; at every geometric point,
$\eta(vc)$ is the value of the $\mathcal O_E$-linear map $\iota_{M'}$ of a Tate system at $vc\in\bar{\mathfrak M}'L$;
and two points of $\mathcal A$ over the reduced scheme $X'$ that agree at every geometric point are equal. So
$s\eta_y(c)=\eta_y(sc)=0$ for $c\in L$ and $s\in\mathfrak M'$, and $M'\eta_y(c)=0$, as $M'\in\mathfrak M'$.

Roots. Multiplication by $M'$ on $\mathcal A$ is finite, flat and surjective: $\mathcal A$ is proper over $X'$, and on
each fibre multiplication by $M'$ is an isogeny \cite{Mum70}, so it is flat by the fibrewise criterion of flatness,
and it is proper and quasi-finite, hence finite. Let $B_0=S$. For $b=1,\dots,n$ in turn, the fibre product of $[M']$
with the image of the point $\eta_y(c_b)$ over $\operatorname{Spec}B_{b-1}$ is finite, hence affine, over
$\operatorname{Spec}B_{b-1}$, say $\operatorname{Spec}B_b$, and $B_{b-1}\to B_b$ is flat and surjective on spectra,
hence faithfully flat; its second projection is a point $w_b$ of $\mathcal A$ over $B_b$ with $M'w_b$ the image of
$\eta_y(c_b)$. Put $B=B_n$, which is faithfully flat over $S$, and let $\varphi\colon S\to B$ be the structure map. By
Lemma~\ref{cv-uniformization-theta-rule-unique-base-change}(1), $y$ has an image $y'$ under $\varphi$. Then
$A_{y'}=A_y\times_SB$, the pullback of $\eta_y(c_b)$ is $\eta_{y'}(c_b)$, and the images of the $w_b$ over $B$ give
$v_b\in A_{y'}(B)$ with $M'v_b=\eta_{y'}(c_b)$.

Roots killed by $\mathfrak M'^2$. As $\mathfrak M'+\bar{\mathfrak M}'=\mathcal O_E$, also
$\mathfrak M'^2+\bar{\mathfrak M}'=\mathcal O_E$; write $1=i+j$ with $i\in\mathfrak M'^2$ and
$j\in\bar{\mathfrak M}'$, and put $t'_b=jv_b$. Then $M't'_b=j\eta_{y'}(c_b)=\eta_{y'}(c_b)-i\eta_{y'}(c_b)=
\eta_{y'}(c_b)$, as $i\in\mathfrak M'$. For $s\in\mathfrak M'^2$, $sj\in\mathfrak M'^2\bar{\mathfrak M}'=M'\mathfrak M'$,
say $sj=M'u$ with $u\in\mathfrak M'$, so $st'_b=uM'v_b=u\eta_{y'}(c_b)=0$. Also $M'^2t'_b=M'\eta_{y'}(c_b)=0$.

Translates. As $M'$ is odd, Lemma~\ref{cv-uniformization-theta-rule-unique-translates}(1) gives an $M'$-translate
$F_b$ of $\vartheta_{y'}$ over $\eta_{y'}(c_b)$ with root $t'_b$. As $t'_b$ is killed by $\mathfrak M'^2$, $F_b$ is a
level translate.
\end{proof}
% Lean: exists_faithfullyFlat_isVtOver (Laws): hpow and hkill from moduli_univ_eta_killed (UnivPol.isReduced,
% LevelPoints.eta_smul with section_ext and the Tate systems of o.geom, eta_ker, eta_pow);
% T.exists_faithfullyFlat_nthRoots_of_torsion = exists_faithfullyFlat_roots (the law mulN: [M'] finite, flat and
% surjective; FlatRoots.exists_lift, the affine fibre product with flat_and_surjective_SpecMap_iff, and exists_lifts
% by induction on the points, composing faithfully flat maps); T.exists_framedPoint_baseChange_algHom for y';
% exists_root_killedBy_pow with N = M' (exists_root_killedBy with I = Mb^2 and J = conj(Mb), t' = j w and
% comp_act_eq_one_of_pow_eq_killed); and the first part of T.exists_translation_sections_universal_object (Lemma
% cv-uniformization-theta-rule-unique-translates(1)) with ok = KilledBy (Mb^2). Size beyond the laws and Lemmas
% cv-uniformization-theta-rule-unique-base-change(1) and -translates(1): about 310 lines; beyond the laws, chart
% independence and the lemmas before this one from Lemma cv-uniformization-theta-rule-connection-charts on, about 290.

\begin{lemma}\label{cv-uniformization-theta-rule-unique}
In the setting of Lemma~\ref{cv-uniformization-theta-marking}, let $\nabla$ be a canonical Taylor family, $n\ge0$,
$c_1,\dots,c_n\in L$, $a$ a family of nonnegative integers indexed by the letters and $P$ a polynomial harmonic of
row degrees $a$ (Definition~\ref{cv-uniformization-theta-rule-values}(6)). Then every framed point $y$ of $X'$ over an
$F'$-algebra has exactly one theta value of $c_1,\dots,c_n$ and $P$ at $y$ for $\nabla$
(Definition~\ref{cv-uniformization-theta-rule-values}(3)).
\end{lemma}
\leanref{Development/AbelianVarietyTheory/Proofs/CuspFormEvaluations/ThetaValueExistsUniqueOfChartIndependence.lean}{GlobalTheory.thValTot_of_chartIndep}
\leanref{Development/AbelianVarietyTheory/Proofs/FormalGroupCharts/ChartIndependenceUnconditional.lean}{GlobalTheory.chartIndependence_unconditional}
\begin{proof}
Let $y$ be a framed point over $S$.

Existence. By Lemma~\ref{cv-uniformization-theta-rule-unique-cover} there are a faithfully flat homomorphism
$\varphi\colon S\to S'$, an image $y'$ of $y$ under $\varphi$ and level translates $F_b$ of $\vartheta_{y'}$ over
$\eta_{y'}(c_b)$. Write $V=\nabla_{y'}(P;F_1,\dots,F_n)$. Let $\psi_1,\psi_2\colon S'\to S'\otimes_SS'$ be the maps
$s\mapsto s\otimes1$ and $s\mapsto1\otimes s$, so that $\psi_1\circ\varphi=\psi_2\circ\varphi$.
Lemma~\ref{cv-uniformization-theta-rule-unique-compare}, applied to $(\varphi,y',F)$ twice, gives $V\otimes1=1\otimes V$.
As $\varphi$ is faithfully flat, the sequence $S\to S'\to S'\otimes_SS'$, with the second map $s\mapsto1\otimes s-s\otimes1$,
is exact \cite[Tag 023M]{Sta}. So $V=\varphi(v)$ for some $v\in S$, and $v$ is a theta value.

Uniqueness. Let $w$ be another theta value, with a faithfully flat $\varphi_2\colon S\to S_2$, an image $y_2$ and level
translates $F_{2,b}$. Lemma~\ref{cv-uniformization-theta-rule-unique-compare}, with $T=S'\otimes_SS_2$ and the maps
$s\mapsto s\otimes1$ and $s\mapsto1\otimes s$, gives $\varphi(v)\otimes1=1\otimes\varphi_2(w)$, that is
$1\otimes\varphi_2(v)=1\otimes\varphi_2(w)$. As $S'$ is faithfully flat over $S$, the map $S_2\to S'\otimes_SS_2$,
$s\mapsto1\otimes s$, is injective \cite[Tag 023M]{Sta}, and so is $\varphi_2$. Hence $v=w$.
\end{proof}
% Lean: thValTot_of_chartIndep (with allSixLaws.1 and ChartIndepSt2, chartIndependence_unconditional): the cover
% T.exists_faithfullyFlat_cover_level_sections (Lemma cv-uniformization-theta-rule-unique-cover); thExpr_tmul_lemma,
% which is T.thetaExpression_eq_of_witnesses (Lemma cv-uniformization-theta-rule-unique-compare) for T = S1 (x)_S S2
% and the two inclusions; Algebra.IsEffective.of_faithfullyFlat for v; Algebra.TensorProduct.tmul_one_eq_one_tmul,
% Module.FaithfullyFlat.tensorProduct_mk_injective and FaithfulSMul.algebraMap_injective for uniqueness. Size beyond
% the laws, chart independence, the connection and Lemmas cv-uniformization-theta-rule-unique-base-change to -cover:
% about 90 lines.

\begin{lemma}\label{cv-uniformization-theta-rule-weight-natural}
In the setting of Lemma~\ref{cv-uniformization-theta-marking}, let $\nabla$ be a canonical Taylor family, $n\ge0$,
$c_1,\dots,c_n\in L$, $a$ a family of nonnegative integers indexed by the letters and $P$ a polynomial harmonic of
row degrees $a$ (Definition~\ref{cv-uniformization-theta-rule-values}(6)). Let $\varphi\colon S\to T$ be a
homomorphism of $F'$-algebras, $y$ a framed point of $X'$ over $S$ and $z$ an image of $y$ under $\varphi$. If $v$ is a
theta value of $c_1,\dots,c_n$ and $P$ at $y$ for $\nabla$ and $w$ is one at $z$, then $\varphi(v)=w$.
\end{lemma}
\leanref{Development/AbelianVarietyTheory/Proofs/CuspFormEvaluations/ThetaValueNaturalAlgebraMapStatement.lean}{GlobalTheory.T.thetaValue_natural_algebraMap}
\leanref{Development/AbelianVarietyTheory/Proofs/CuspFormEvaluations/ThetaValueNaturalAlgebraMap.lean}{GlobalTheory.thetaValue_natural}
\begin{proof}
Let $v$ come with a faithfully flat $\varphi_0\colon S\to S'$, an image $y_0$ of $y$ and level translates $F_{0,b}$,
and $w$ with a faithfully flat $\varphi_1\colon T\to T'$, an image $z_1$ of $z$ and level translates $F_{1,b}$. By
Lemma~\ref{cv-uniformization-theta-rule-unique-base-change}(3), $z_1$ is an image of $y$ under $\varphi_1\circ\varphi$.
Let $U=T'\otimes_SS'$, with $S$ acting on $T'$ through $\varphi_1\circ\varphi$, and let $\psi_1\colon S'\to U$ and
$\psi_2\colon T'\to U$ be the maps $s\mapsto1\otimes s$ and $s\mapsto s\otimes1$; then
$\psi_1\circ\varphi_0=\psi_2\circ\varphi_1\circ\varphi$. By Lemma~\ref{cv-uniformization-theta-rule-unique-compare},
\[
 \psi_2\big(\varphi_1(\varphi(v))\big)=\psi_1\big(\varphi_0(v)\big)=\psi_1\big(\nabla_{y_0}(P;F_{0,1},\dots,F_{0,n})\big)
 =\psi_2\big(\nabla_{z_1}(P;F_{1,1},\dots,F_{1,n})\big)=\psi_2\big(\varphi_1(w)\big).
\]
As $\varphi_0$ is faithfully flat, $\psi_2$ is injective \cite[Tag 023M]{Sta}, and as $\varphi_1$ is faithfully flat, it
is injective. So $\varphi(v)=w$.
\end{proof}
% Lean: thetaValue_natural (Laws, ChartIndepSt2): FramedBaseChange.trans for hm01, the algebra U = T' (x)_S S' with
% psi1 = includeRight and psi2 = algebraMap, hcomp, T.thetaExpression_eq_of_witnesses (Lemma
% cv-uniformization-theta-rule-unique-compare), FaithfulSMul.algebraMap_injective for psi2 and hff1.injective. Size
% beyond the laws, chart independence and the lemmas before this one from Lemma
% cv-uniformization-theta-rule-connection-charts on: about 50 lines.

\begin{lemma}\label{cv-uniformization-theta-rule-weight-commute}
Let $S$ be a commutative ring and $F$ a commutative formal group law of dimension $d$ over $S$, in the variables
$t=(t_1,\dots,t_d)$ and $t'=(t'_1,\dots,t'_d)$. For $v\in S^d$ let $D_v$ be the $S$-linear map of $S[[t]]$ to itself
\[
 D_v=\sum_{j=1}^d\Big(\sum_{i=1}^dv_iJ_{ji}(t)\Big)\frac{\partial}{\partial t_j},\qquad
 J_{ji}(t)=\frac{\partial F_j(t,t')}{\partial t'_i}\Big|_{t'=0}.
\]
Then $D_{\sum_kc_kw_k}=\sum_kc_kD_{w_k}$ for finitely many $c_k\in S$ and $w_k\in S^d$, and $D_vD_w=D_wD_v$ for all
$v,w\in S^d$.
\end{lemma}
\leanref{Development/AbelianVarietyTheory/Proofs/FormalGroupCharts/InvariantDerivationsCommuteFormalConnection.lean}{GlobalTheory.FormalConnection.dof_comm_law}
\leanref{Development/AbelianVarietyTheory/Proofs/CuspFormEvaluations/InvariantOperatorValueTriangularChange.lean}{GlobalTheory.RowTriangular.dEnd_comb}
\begin{proof}
The first claim can be read off the formula. For the second, write $\partial_j=\partial/\partial t_j$ and
$D_v=\sum_j\xi^v_j\partial_j$ with $\xi^v_j=\sum_iv_iJ_{ji}$. In $D_vD_w-D_wD_v$ the terms with second derivatives
cancel, as $\partial_j\partial_{j'}=\partial_{j'}\partial_j$, so
\[
 D_vD_w-D_wD_v=\sum_j\big(D_v(\xi^w_j)-D_w(\xi^v_j)\big)\partial_j,\qquad D_v(\xi^w_j)=\sum_{b,i}v_bw_iT_{bij},
\]
where $T_{bij}=\sum_cJ_{cb}\,\partial_cJ_{ji}$. It is enough that $T_{bij}=T_{ibj}$. Differentiating the associativity
$F(F(t,s),s')=F(t,F(s,s'))$ in $s'_i$ at $s'=0$ gives
$J_{ji}(F(t,s))=\sum_c\big(\partial F_j(t,s)/\partial s_c\big)J_{ci}(s)$. Differentiating this in $s_b$ at $s=0$, with
$F(t,0)=t$ and $J_{ci}(0)=\delta_{ci}$, gives
\[
 T_{bij}=\frac{\partial^2F_j(t,s)}{\partial s_b\,\partial s_i}\Big|_{s=0}+\sum_cJ_{jc}(t)\,\kappa_{cbi},
\]
where $\kappa_{cbi}$, the coefficient of $t_b$ in $J_{ci}$, is the coefficient of $t_bs_i$ in $F_c(t,s)$. The first term
is symmetric in $b$ and $i$, and so is $\kappa_{cbi}$, as $F_c(t,s)=F_c(s,t)$.
\end{proof}
% Lean: FormalConnection.dof_comm_law: dof_vf (D_v is the vector field vf (co G v)), vf_comm (the bracket of two
% vector fields), co_bracket with vf_co_jac and T_symm, which is T_eq (from FormalLog.jac_F, the derivative of the
% associativity) with cc_d_jac and kappa_symm (G.comm). The linearity is RowTriangular.dEnd_comb. Size beyond the laws
% of abelian schemes, chart independence and the lemmas before this one from Lemma
% cv-uniformization-theta-rule-connection-charts on: about 470 lines.

\begin{lemma}\label{cv-uniformization-theta-rule-weight-substitution}
In the setting of Lemma~\ref{cv-uniformization-theta-marking}, let $n\ge0$, let $a$ be a family of nonnegative integers
indexed by the letters, and let $P$ be a polynomial in the letter variables with coefficients in $F'$ that satisfies
conditions (b) and (c) of Definition~\ref{cv-uniformization-theta-rule-values}(6) for $a$. Let $T$ be a commutative
$F'$-algebra, and let $c_{\tau,i,j}\in T$, for letters $(\tau,i)$ and $(\tau,j)$, with $c_{\tau,i,j}=0$ for $j>i$ and
$c_{\tau,i,i}$ a unit for every letter $(\tau,i)$. Then for every family $y=(y_{\ell,b})$ of elements of $T$ indexed by
the letters $\ell$ and $1\le b\le n$,
\[
 P\big((z_{\ell,b})_{\ell,b}\big)=\Big(\prod_{(\tau,i)}c_{\tau,i,i}^{a_{(\tau,i)}}\Big)P\big((y_{\ell,b})_{\ell,b}\big),
 \qquad z_{(\tau,i),b}=\sum_jc_{\tau,i,j}\,y_{(\tau,j),b}.
\]
\end{lemma}
\leanref{Development/AbelianVarietyTheory/Proofs/CuspFormEvaluations/InvariantOperatorValueTriangularChange.lean}{GlobalTheory.RowTriangular.eval₂_triangular}
\begin{proof}
For $j<i$ put $N_{\tau,i,j}=c_{\tau,i,i}^{-1}c_{\tau,i,j}$, and put $N_{\tau,i,j}=0$ for $j\ge i$. Let
$y'_{(\tau,i),b}=y_{(\tau,i),b}+\sum_jN_{\tau,i,j}\,y_{(\tau,j),b}$, so that $z_{(\tau,i),b}=c_{\tau,i,i}\,y'_{(\tau,i),b}$.

The diagonal. By condition (c), each monomial of $P$ has degree $a_{(\tau,i)}$ in the variables $X_{(\tau,i),b}$,
$1\le b\le n$, for every letter $(\tau,i)$. So
$P\big((z_{\ell,b})\big)=\big(\prod_{(\tau,i)}c_{\tau,i,i}^{a_{(\tau,i)}}\big)P\big((y'_{\ell,b})\big)$.

The unipotent part. Let $(\tau,i)$ be a letter and $i'<i$. Let $Q$ be $P$ with $X_{(\tau,i),b}$ replaced by
$X_{(\tau,i),b}+ZX_{(\tau,i'),b}$ for every $b$, a polynomial in the letter variables and one more variable $Z$. By
condition (b), $Q-P$ vanishes at $Z=\lambda$ for every $\lambda\in F'$; as $F'$ is infinite, $Q=P$. Substituting a
family in $T$ for the letter variables and an element of $T$ for $Z$ shows that $P$ is unchanged, as a function of
families in $T$, when the row $(y_{(\tau,i),b})_b$ is replaced by $(y_{(\tau,i),b}+\lambda y_{(\tau,i'),b})_b$, for
any $\lambda\in T$. These replacements for the $i'<i$, with $\lambda=N_{\tau,i,i'}$, change only the row $(\tau,i)$ and
read only other rows; done one after the other, they replace the row $(\tau,i)$ by $\big(y_{(\tau,i),b}+\sum_{j<i}N_{\tau,i,j}y_{(\tau,j),b}\big)_b$ and
keep the others, and $P$ is unchanged. Now do this for the rows one at a time, in order of decreasing $i$. Each step
changes one row $(\tau,i)$ using the rows $(\tau,j)$ with $j<i$, which no earlier step has changed. So the steps
together carry $y$ to $y'$, and $P\big((y'_{\ell,b})\big)=P\big((y_{\ell,b})\big)$.
\end{proof}
% Lean: RowTriangular.eval₂_triangular (for any infinite field k, ring homomorphism k -> T, RowInv and RowDeg):
% eval₂_diag for the diagonal (RowDeg) and eval₂_unip for the unipotent part, which is
% eval₂Hom_unipotent_block_substitution_eq (induction on the rows by Finset.induction_on_max_value, the row of
% largest index first) through eval₂Hom_unipotent_substitution_eq (one row, by induction on the i') and eval₂_shift
% (RowInv extended to coefficients in T by MvPolynomial.funext over the infinite field). Size: about 170 lines.

\begin{lemma}\label{cv-uniformization-theta-rule-weight-operators}
In the setting of Lemma~\ref{cv-uniformization-theta-marking}, let $S$ be an $F'$-algebra, let $B$ be a formal germ of
dimension $d$ over $S$ (Definition~\ref{cv-hodge-lattice-germ}) with group law $F$, and let $h(t,t')$ be a power series
over $S$ with $h(0,t')=h(t,0)=0$ whose exponentials give the series of $B$ as in
Definition~\ref{cv-uniformization-theta-rule-values}(5); this hypothesis on $h$ is not used in the proof. Let $k\ge0$.
For a power series $f(t)$ over $S$, a family $u=(u_\ell)$ of vectors of $S^d$ indexed by the letters and a family
$\alpha$ of nonnegative integers indexed by the letters, let $V_u(f;\alpha)$ be the constant term of
$\big(\prod_\ell D_{u_\ell}^{\alpha_\ell}\big)\big(f(t)\,e^{k\,h(t,[w_0]t)}\big)$, with $D_v$ as in
Lemma~\ref{cv-uniformization-theta-rule-weight-commute} for $F$; by that lemma the order of the factors does not
matter. Let $n\ge0$, let $a$ be a family of nonnegative integers indexed by the letters, and let $P$ be a polynomial
in the letter variables with coefficients in $F'$ that satisfies conditions (a) to (c) of
Definition~\ref{cv-uniformization-theta-rule-values}(6) for $a$; condition (a) is not used in the proof. Let $u$ be a
family of vectors of $S^d$ indexed by the letters, and let $c_{\tau,i,j}\in S$, for letters $(\tau,i)$ and $(\tau,j)$,
with $c_{\tau,i,j}=0$ for $j>i$ and $c_{\tau,i,i}$ a unit for every letter $(\tau,i)$; put
$u'_{(\tau,i)}=\sum_jc_{\tau,i,j}u_{(\tau,j)}$. Then for all power series $f_1,\dots,f_n$ over $S$, with the notation of
Definition~\ref{cv-uniformization-theta-rule-values}(3),
\[
 \sum_mP_m\prod_{b=1}^nV_{u'}(f_b;m_b)=\Big(\prod_{(\tau,i)}c_{\tau,i,i}^{a_{(\tau,i)}}\Big)
 \sum_mP_m\prod_{b=1}^nV_u(f_b;m_b).
\]
\end{lemma}
\leanref{Development/AbelianVarietyTheory/Proofs/CuspFormEvaluations/InvariantOperatorValueTriangularChangeStatement.lean}{GlobalTheory.T.invariantOperator_value_triangular_change}
\leanref{Development/AbelianVarietyTheory/Proofs/CuspFormEvaluations/InvariantOperatorValueTriangularChange.lean}{GlobalTheory.nabVal_triangular}
\begin{proof}
Put $G_b=f_b(t)\,e^{k\,h(t,[w_0]t)}$. By Lemma~\ref{cv-uniformization-theta-rule-weight-commute} the maps $D_v$ commute
with one another, so the $S$-subalgebra $A$ of $\operatorname{End}_S(S[[t]])$ that they generate is commutative. Let
$A^{\otimes n}$ be the tensor product over $S$ of $n$ copies of $A$, and let $\lambda\colon A^{\otimes n}\to S$ be the
$S$-linear map with $\lambda(z_1\otimes\dots\otimes z_n)=\prod_b\lambda_b(z_b)$, where $\lambda_b(z)$ is the constant
term of $z(G_b)$. For a family $x=(x_\ell)$ of elements of $A$ indexed by the letters, let $x_{\ell,b}\in A^{\otimes n}$
be $x_\ell$ in the $b$-th factor and $1$ in the others. If $x_\ell=D_{u_\ell}$ for a family $u$ of vectors, then
$V_u(f_b;\alpha)=\lambda_b\big(\prod_\ell x_\ell^{\alpha_\ell}\big)$, and the monomial with exponents $m$ takes at the
$x_{\ell,b}$ the value $\bigotimes_b\prod_\ell x_\ell^{m_{\ell,b}}$. As $\lambda$ is $S$-linear,
\[
 \sum_mP_m\prod_bV_u(f_b;m_b)=\lambda\big(P\big((x_{\ell,b})_{\ell,b}\big)\big).
\]
Let $x_\ell=D_{u_\ell}$ and $x'_\ell=D_{u'_\ell}$. By Lemma~\ref{cv-uniformization-theta-rule-weight-commute},
$x'_{(\tau,i)}=\sum_jc_{\tau,i,j}x_{(\tau,j)}$, so $x'_{(\tau,i),b}=\sum_jc_{\tau,i,j}x_{(\tau,j),b}$ in $A^{\otimes n}$.
Lemma~\ref{cv-uniformization-theta-rule-weight-substitution}, in the commutative $F'$-algebra $A^{\otimes n}$ with the
images of the $c_{\tau,i,j}$, gives
$P\big((x'_{\ell,b})\big)=\big(\prod_{(\tau,i)}c_{\tau,i,i}^{a_{(\tau,i)}}\big)P\big((x_{\ell,b})\big)$. Applying
$\lambda$, which is $S$-linear, gives the claim.
\end{proof}
% Lean: nabVal_triangular (the gauge hypothesis and the harmonicity hP.1 are not used): hcomm from
% FormalConnection.dof_comm_law (Lemma cv-uniformization-theta-rule-weight-commute) and t4 (FK = G.F over S), the
% commutative Algebra.adjoin S s with Algebra.adjoinCommRingOfComm, lam and prodForm on the PiTensorProduct, hval
% (dpoly_monomial, dmon_eq, prodForm_prod), hx' from dEnd_comb, and RowTriangular.eval₂_triangular (Lemma
% cv-uniformization-theta-rule-weight-substitution). Size beyond Lemmas cv-uniformization-theta-rule-weight-commute
% and -weight-substitution: about 140 lines.

\begin{lemma}\label{cv-uniformization-theta-rule-weight-taylor}
In the setting of Lemma~\ref{cv-uniformization-theta-marking}, let $\nabla$ be a canonical Taylor family, let
$n,N\ge0$, let $a$ be a family of nonnegative integers indexed by the letters, and let $P$ be a polynomial in the
letter variables with coefficients in $F'$ that satisfies conditions (a) to (c) of
Definition~\ref{cv-uniformization-theta-rule-values}(6) for $a$. Let $y$ be a framed point of $X'$ over an
$F'$-algebra $S$, let $t=(t_0,(t_\ell)_\ell)$ be a family of units of $S$ indexed by $0$ and the letters, and let
$y'=y\cdot t$ be a framed point with the point of $y$ whose tangent vector at each letter $(\tau,i)$ is $t_{(\tau,i)}$
times that of $y$ at $(\tau,i)$ plus a combination of those of $y$ at the letters $(\tau,i')$ with $i'<i$, and whose
generator is $t_0\vartheta_y$; the condition on the generator is not used in the proof. Then for all sections
$F_1,\dots,F_n$ of $\mathcal L_y^{N^2}=\mathcal L_{y'}^{N^2}$, with the notation of
Definition~\ref{cv-uniformization-theta-rule-values}(3),
\[
 \sum_mP_m\prod_{b=1}^n\nabla_{y'}(F_b;m_b)=\Big(\prod_\ell t_\ell^{a_\ell}\Big)\sum_mP_m\prod_{b=1}^n\nabla_y(F_b;m_b).
\]
\end{lemma}
\leanref{Development/AbelianVarietyTheory/Proofs/CuspFormEvaluations/HarmonicSumTriangularFrameChangeStatement.lean}{GlobalTheory.T.harmonic_sum_triangular_frameChange}
\leanref{Development/AbelianVarietyTheory/Proofs/CuspFormEvaluations/HarmonicSumTriangularFrameChange.lean}{GlobalTheory.TriangularFrame.harmonic_sum_frBor}
\leanref{Development/AbelianVarietyTheory/Proofs/CuspFormEvaluations/HarmonicSumTriangularFrameChange.lean}{GlobalTheory.TriangularFrame.sum_eq}
\begin{proof}
Write the tangent vector of $y'$ at each letter $(\tau,i)$ as $\sum_jc_{\tau,i,j}$ times that of $y$ at $(\tau,j)$,
with $c_{\tau,i,j}\in S$, $c_{\tau,i,j}=0$ for $j>i$ and $c_{\tau,i,i}=t_{(\tau,i)}$. Write $V_C(f;\alpha)$ for the
constant term of $\big(\prod_\ell D_\ell^{\alpha_\ell}\big)\big(f(t)e^{k\,h(t,[w_0]t)}\big)$ for a chart datum $C$ with
power series $h$ and a power series $f$, as in Definition~\ref{cv-uniformization-theta-rule-values}(5), with
$k=N^2$. As $\nabla$ is canonical, there are a chart datum $C$ at $y$ and a chart datum $C'$ at $y'$ for which every
section has a germ, with $\nabla_y(F;\alpha)=V_C(f;\alpha)$ and $\nabla_{y'}(F;\alpha)=V_{C'}(f';\alpha)$ for germs $f$
of $F$ for $C$ and $f'$ for $C'$. Choose germs $f_b$ of $F_b$ for $C$ and $f'_b$ of $F_b$ for $C'$.

A third chart datum. Let $(u_{\ell e})$ be the matrix of $C$ and $v_1,\dots,v_{d'}$ the tangent vectors given by its
points $x_m$, and let $D$ be the chart datum with the formal germ, the points $x_m$, the generators $e_m$ and $w$ and the
power series $h$ of $C$ and the matrix $u'_{(\tau,i),e}=\sum_jc_{\tau,i,j}u_{(\tau,j),e}$. It is a chart datum at $y'$,
which has the point of $y$: the
tangent vector of $y'$ at $(\tau,i)$ is $\sum_jc_{\tau,i,j}\sum_eu_{(\tau,j),e}v_e=\sum_eu'_{(\tau,i),e}v_e$, by the
linearity of combinations of tangent vectors (Definition~\ref{cv-germs-chart}(2)); and the matrix $u'$ is the product
of $(c_{\tau,i,j})$, which is block lower triangular with units on the diagonal, and the invertible matrix $u$, so it
is invertible. The germs of a section for a chart datum depend only on its points $x_m$ and its generators $e_m$ and
$w$, so $f_b$ is a germ of $F_b$ for $D$.

Chart independence. As $F'$ is a number field, some integer $N_0\ge1$ makes the coefficients of $N_0P$ algebraic
integers; then $N_0P$ is harmonic of row degrees $a$, with the monomials of $P$. Lemma~\ref{cv-uniformization-theta-rule-unique-chart}
for $y'$, the chart data $C'$ and $D$, the germs $f'_b$ and $f_b$ and the polynomial $N_0P$ gives, after cancelling
$N_0$, which is a unit in $S$,
\[
 \sum_mP_m\prod_b\nabla_{y'}(F_b;m_b)=\sum_mP_m\prod_bV_{C'}(f'_b;m_b)=\sum_mP_m\prod_bV_D(f_b;m_b).
\]

The triangular change. The chart data $C$ and $D$ have the same formal germ and power series $h$, and the matrix of
$D$ comes from that of $C$ by the $c_{\tau,i,j}$. Lemma~\ref{cv-uniformization-theta-rule-weight-operators}, with
$c_{\tau,i,i}=t_{(\tau,i)}$, gives
\[
 \sum_mP_m\prod_bV_D(f_b;m_b)=\Big(\prod_\ell t_\ell^{a_\ell}\Big)\sum_mP_m\prod_bV_C(f_b;m_b)
 =\Big(\prod_\ell t_\ell^{a_\ell}\Big)\sum_mP_m\prod_b\nabla_y(F_b;m_b).
\]
\end{proof}
% Lean: T.harmonic_sum_triangular_frameChange = TriangularFrame.harmonic_sum_frBor (ChartIndepSt2), for FrBor t y y'
% (its third clause, on the generator, is discarded) and the identity h, with the F' the pullbacks of the F along it
% (eq_of_isPullSec_id); then TriangularFrame.sum_eq: C and C' from IsNab (hnab.2), the chart datum D with
% uv = cf . uvec (tan from TangentLinear.isCombAt_mulVec with isTanAt_jetTan, basis from basis_of_triangular),
% exists_natCast_smul_mem_harmA for N0, ChartIndepSt2 (Lemma cv-uniformization-theta-rule-unique-chart) for C' and D,
% the cancellation of N0, and T.invariantOperator_value_triangular_change (Lemma
% cv-uniformization-theta-rule-weight-operators). Size beyond the laws, chart independence and the lemmas before
% this one from Lemma cv-uniformization-theta-rule-connection-charts on: about 1,200 lines, of which about 780 in
% Lemmas cv-uniformization-theta-rule-weight-commute to -weight-operators and 180 in TangentLinear.isCombAt_mulVec.

\begin{lemma}\label{cv-uniformization-theta-rule-weight-triangular}
In the setting of Lemma~\ref{cv-uniformization-theta-marking}, let $\nabla$ be a canonical Taylor family, $n\ge0$,
$c_1,\dots,c_n\in L$, $a$ a family of nonnegative integers indexed by the letters and $P$ a polynomial harmonic of
row degrees $a$ (Definition~\ref{cv-uniformization-theta-rule-values}(6)). Let $y$ be a framed point of $X'$ over $S$,
$t$ a family of units of $S$ indexed by $0$ and the letters, and $y'=y\cdot t$
(Lemma~\ref{cv-uniformization-theta-rule-weight-taylor}). If $v$ is a theta value of $c_1,\dots,c_n$ and $P$ at $y$
for $\nabla$, then $t_0^n\prod_\ell t_\ell^{a_\ell}\,v$ is one at $y'$.
\end{lemma}
\leanref{Development/AbelianVarietyTheory/Proofs/CuspFormEvaluations/ThetaValueWeightTriangularFrameStatement.lean}{GlobalTheory.T.thetaValue_weight_triangular_frame}
\leanref{Development/AbelianVarietyTheory/Proofs/CuspFormEvaluations/ThetaValueWeightTriangularFrame.lean}{GlobalTheory.thetaValue_weight}
\leanref{Development/AbelianVarietyTheory/Proofs/CuspFormEvaluations/ConnectionLinearInSection.lean}{GlobalTheory.nab_smul_section}
\begin{proof}
Let $v$ come with a faithfully flat $\varphi\colon S\to S'$, an image $z$ of $y$ under $\varphi$ and level translates
$F_b$ of $\vartheta_z$ over $\eta_z(c_b)$, with $\varphi(v)=\nabla_z(P;F_1,\dots,F_n)$. By
Lemma~\ref{cv-uniformization-theta-rule-unique-base-change}(1), $y'$ has an image $z'$ under $\varphi$. The point
under $z'$ is the point under $y'$, which is that under $y$, composed with $\operatorname{Spec}\varphi$, so $z'$ has
the point of $z$. Base change carries the triangular combinations of the tangent vectors of $y$ that give those of $y'$
to the same combinations of those of $z$, with coefficients mapped by $\varphi$, and the pullback of $t_0\vartheta_y$
is $\varphi(t_0)\vartheta_z$; so $z'=z\cdot\varphi(t)$.

Write $r=\varphi(t_0)$. The sections $rF_b$ are level translates of $\vartheta_{z'}=r\vartheta_z$ over
$\eta_{z'}(c_b)=\eta_z(c_b)$, with the roots of the $F_b$: if $F_b=Us_1$ with $s_1$ corresponding to
$[M']^*\vartheta_z$, then $rs_1$ corresponds to $[M']^*(r\vartheta_z)$, and $rF_b=U(rs_1)$. Taylor values are linear
in the section: if $f$ is a germ of a section $F$ of $\mathcal L_{z'}^{M'^2}$ for a chart datum at $z'$, then $rf$ is a
germ of $rF$, as values in a trivialization are linear, and the constant term in
Definition~\ref{cv-uniformization-theta-rule-values}(5) is linear in the germ; so
$\nabla_{z'}(rF_b;\alpha)=r\nabla_{z'}(F_b;\alpha)$. With Lemma~\ref{cv-uniformization-theta-rule-weight-taylor} for
$z$, $z'$, $N=M'$ and the sections $F_b$, and as each product in $\nabla_{z'}(P;rF_1,\dots,rF_n)$ has the $n$ factors
$\nabla_{z'}(rF_b;m_b)$,
\[
 \nabla_{z'}(P;rF_1,\dots,rF_n)=r^n\prod_\ell\varphi(t_\ell)^{a_\ell}\,\nabla_z(P;F_1,\dots,F_n)
 =\varphi\Big(t_0^n\prod_\ell t_\ell^{a_\ell}\,v\Big).
\]
So $t_0^n\prod_\ell t_\ell^{a_\ell}v$ is a theta value at $y'$, with $\varphi$, $z'$ and the sections $rF_b$.
\end{proof}
% Lean: the first clause of thetaValue_weight (Laws, ChartIndepSt2): exists_reframe,
% T.exists_framedPoint_baseChange_algHom for z', frBor_units_map_of_frMap (z' = z . phi(t)), gen_eq and
% isVtOver_scale (scaleSec, isTwistPull_scale) for the sections r F_b, map_wchar, and sum_eq, which combines
% T.harmonic_sum_triangular_frameChange (Lemma cv-uniformization-theta-rule-weight-taylor, with h the identity and
% isPullSec_id) with T.connection_linear_in_section = nab_smul_section (isValueAt_scale and nabVal_C_mul, about 150
% lines), and wchar (ex Dw a n) t = t_0^n prod_l t_l^(a_l). The second clause of thetaValue_weight, on ShortValAt, is
% not used here. Size beyond the laws, chart independence and the lemmas before this one from Lemma
% cv-uniformization-theta-rule-connection-charts on: about 450 lines.

\begin{lemma}\label{cv-uniformization-theta-rule-weight}
In the setting of Lemma~\ref{cv-uniformization-theta-marking}, let $\nabla$ be a canonical Taylor family, $n\ge0$,
$c_1,\dots,c_n\in L$, $a$ a family of nonnegative integers indexed by the letters, $P$ a polynomial harmonic of row
degrees $a$, and $f$ a theta rule of $c_1,\dots,c_n$ and $P$ for $\nabla$
(Definition~\ref{cv-uniformization-theta-rule-values}(4) and~(6)). Then $f$ is a rational section of $X'$ of exponent
$e(a,n)$, the family with the entry $a_\ell$ at each letter $\ell$ and $n$ at $0$: $f(z)=\varphi(f(y))$ whenever $z$
is an image of a framed point $y$ over a test ring $R$ under a homomorphism of $F'$-algebras $\varphi$ from $R$ to a
test ring, and $f(y\cdot t)=t_0^n\prod_\ell t_\ell^{a_\ell}f(y)$ for every framed point $y$ over a test ring $R$, every
family $t$ of units of $R$ indexed by $0$ and the letters, and every $y\cdot t$ as in
Lemma~\ref{cv-uniformization-theta-rule-weight-taylor}.
\end{lemma}
\leanref{Development/AbelianVarietyTheory/Proofs/CuspFormEvaluations/ThetaRuleMemSec.lean}{GlobalTheory.thetaRule_mem_sec}
\begin{proof}
Naturality. Let $z$ be an image of $y$ under $\varphi\colon R\to R'$. Then $f(y)$ is a theta value of
$c_1,\dots,c_n$ and $P$ at $y$, and $f(z)$ is one at $z$, so $f(z)=\varphi(f(y))$ by
Lemma~\ref{cv-uniformization-theta-rule-weight-natural}.

Weight. Let $y'=y\cdot t$ over $R$. By Lemma~\ref{cv-uniformization-theta-rule-weight-triangular},
$t_0^n\prod_\ell t_\ell^{a_\ell}f(y)$ is a theta value at $y'$, and so is $f(y')$. As $y'$ is an image of itself under
the identity of $R$ (Lemma~\ref{cv-uniformization-theta-rule-unique-base-change}(2)),
Lemma~\ref{cv-uniformization-theta-rule-weight-natural} gives $f(y')=t_0^n\prod_\ell t_\ell^{a_\ell}f(y)$.
\end{proof}
% Lean: thetaRule_mem_sec (30 lines): f in sec Xo c a n (secW: IsNatural for FrMap, IsWeight for FrBor with
% wchar (ex Dw a n) t); naturality from T.thetaValue_natural_algebraMap (= thetaValue_natural, Lemma
% cv-uniformization-theta-rule-weight-natural), the weight from the first clause of
% T.thetaValue_weight_triangular_frame (Lemma cv-uniformization-theta-rule-weight-triangular) and naturality for the
% identity (hid). Size beyond the laws, chart independence, the connection and the theta values (Lemmas
% cv-uniformization-theta-rule-unique-base-change to -unique): about 1,800 lines, of which about 50 in Lemma
% cv-uniformization-theta-rule-weight-natural, 1,200 in -weight-taylor, 450 in -weight-triangular and 80 here.

\subsection{Theta rules}\label{subsection-cv-theta-rules}

\begin{lemma}\label{cv-uniformization-theta-rule-complex-charts}
In the setting of Lemma~\ref{cv-uniformization-theta-marking}, let $S_0$ be an $F'$-algebra and $x_0$ a point of $X'$
over $S_0$. Then there are $g_1,\dots,g_n\in S_0$ that generate the unit ideal and, for each $i$, a test ring $R_i$
(Definition~\ref{cv-uniformization-theta-rule-values}(4)), an open immersion $j_i\colon\operatorname{Spec}R_i\to X'$ and
a framed point $y_{0,i}$ over $R_i$ with the point $j_i$, such that for every homomorphism of $F'$-algebras
$\lambda\colon S_0\to T$ with $\lambda(g_i)$ a unit, $x_0\circ\operatorname{Spec}\lambda=j_i\circ\operatorname{Spec}\psi$
for some homomorphism of $F'$-algebras $\psi\colon R_i\to T$.
\end{lemma}
\leanref{Development/AbelianVarietyTheory/Proofs/UnitaryExistence/FramedPointReducedFromTestRing.lean}{GlobalTheory.Frames.TestDescent.exists_chart_cover}
\begin{proof}
\openstep{The proof is not yet written out: every point of $X'$ lies in an open chart, the spectrum of a test ring with
bases of the eigenspaces of the Lie algebra and a generator of the sections of the line bundle over it, and finitely
many such charts cover the image of $\operatorname{Spec}S_0$.}
% Lean: TestDescent.exists_chart_cover (SecBaseChangeSt): T.exists_openChart_framedPoint_testRing =
% Frames.exists_openChart_testRing (ModuliChart.exists_chart_eigenframe, as in Lemma cv-volume-key-lie(4);
% T.eigenBasis_iff_tangentBasis; T.exists_fibreSection_nowhereZero_of_pullbackTriple and
% T.fibreSectionGenerates_of_nowhereZero, as in Lemma cv-volume-key-generator(3); reduced_finiteType_iff_testRing),
% Frames.openChart_factor_near_prime at each prime, and a finite subcover (Ideal.span_eq_top_iff_finite). Size beyond
% the laws of abelian schemes, chart independence and the lemmas before this one from Lemma
% cv-uniformization-theta-rule-connection-charts on: about 4,220 lines.
\end{proof}

\begin{lemma}\label{cv-uniformization-theta-rule-complex-universal}
In the setting of Lemma~\ref{cv-uniformization-theta-marking}, let $y_0$ be a framed point of $X'$ over a test ring
$R$. For an $F'$-algebra $T$, a family $Y=(Y_\tau)$ of invertible matrices over $T$, indexed by the embeddings $\tau$
of $E$, with $Y_\tau$ of size the number $m_\tau$ of letters $(\tau,i)$, a unit $t$ of $T$ and a framed point $w$ over
$T$, let $w\cdot(Y,t)$ be the framed point with the point of $w$, the tangent vector $\sum_j(Y_\tau)_{ij}w_{(\tau,j)}$ at
each letter $(\tau,i)$, where the $w_\ell$ are the tangent vectors of $w$, and the generator $t\vartheta_w$. Then there
are a test ring $A$, a homomorphism of $F'$-algebras $\pi\colon R\to A$, such a family $X=(X_\tau)$ over $A$, a unit
$s$ of $A$ and a framed point $u$ over $A$ such that:
\begin{enumerate}
\item for every $F'$-algebra $T$, every homomorphism of $F'$-algebras $\psi\colon R\to T$, every family $g=(g_\tau)$ of
  invertible matrices over $T$ of these sizes and every unit $a$ of $T$, there is exactly one homomorphism of
  $F'$-algebras $k\colon A\to T$ with $k\circ\pi=\psi$, $k(X_\tau)=g_\tau$ for every $\tau$ and $k(s)=a$;
\item $u=u_0\cdot(X,s)$ for an image $u_0$ of $y_0$ under $\pi$
  (Definition~\ref{cv-uniformization-theta-rule-values});
\item for every $F'$-algebra $T$, every framed point $z$ over $T$ and every homomorphism of $F'$-algebras
  $\psi\colon R\to T$ such that the point under $z$ is the point under $y_0$ composed with $\operatorname{Spec}\psi$,
  there is exactly one homomorphism of $F'$-algebras $a\colon A\to T$ with $a\circ\pi=\psi$ such that $z$ is an image
  of $u$ under $a$.
\end{enumerate}
\end{lemma}
\leanref{Development/AbelianVarietyTheory/Proofs/UnitaryExistence/ExistsUniversalFrameTestRingStatement.lean}{GlobalTheory.T.exists_universalFrame_testRing}
\leanref{Development/AbelianVarietyTheory/Proofs/UnitaryExistence/UniversalFrameClassifiesFramedPointsStatement.lean}{GlobalTheory.T.universalFrame_classifies_framedPoints}
\begin{proof}
The ring. Let $A_0$ be the polynomial ring over $R$ in variables $x_{\tau,ij}$, for the embeddings $\tau$ and
$1\le i,j\le m_\tau$, and one more variable $s$, localized at $s\prod_\tau\det(x_{\tau,ij})_{i,j}$. It is of finite type
over $F'$, as $R$ is, and reduced, as a localization of a polynomial ring over the reduced ring $R$. As in the proof
of Lemma~\ref{cv-uniformization-theta-rule-complex-tested}, there are a test ring $A$ and an isomorphism of
$F'$-algebras $A_0\to A$; let $\pi$ be $R\to A_0\to A$, and $X_\tau$ and $s$ the images of the matrices
$(x_{\tau,ij})$ and of $s$. They are invertible, as their determinants divide the inverted element. Part (1) is the
universal property of polynomial rings and of localization: $k$ must send $x_{\tau,ij}$ to the entries of $g_\tau$ and
$s$ to $a$, and this extends $\psi$ to $A_0$, as the determinants of the $g_\tau$ and $a$ are units of $T$.

The framed point. Let $u_0$ be an image of $y_0$ under $\pi$
(Lemma~\ref{cv-uniformization-theta-rule-unique-base-change}(1)), and $u=u_0\cdot(X,s)$; its vectors at the letters of
each $\tau$ form a basis of the eigenspace, as those of $u_0$ do and $X_\tau$ is invertible. This gives (2).

Frame changes. Two framed points $w$ and $w'$ over an $F'$-algebra $T$ with the same point satisfy
$w'=w\cdot(Y,t)$ for exactly one pair $(Y,t)$, as in the proof of
Lemma~\ref{cv-uniformization-theta-rule-complex-glue-injective}. If $k\colon A\to T$ is a homomorphism and $w$ is an
image of $u_0$ under $k$, then the image of $u$ under $k$ is $w\cdot(k(X),k(s))$, as base change carries combinations of
tangent vectors to combinations with the mapped coefficients. An image is determined by the framed point and the
homomorphism.

(3) Let $T$, $z$ and $\psi$ be as in (3), and let $z_0$ be an image of $y_0$ under $\psi$. Then $z$ and $z_0$ have the
same point, so $z=z_0\cdot(g,a_0)$ for exactly one pair $(g,a_0)$. Let $k\colon A\to T$ with $k\circ\pi=\psi$. An image
of $u_0$ under $k$ is an image of $y_0$ under $k\circ\pi=\psi$
(Lemma~\ref{cv-uniformization-theta-rule-unique-base-change}(3)), so it is $z_0$. Hence $z$ is an image of $u$ under
$k$ if and only if $z=z_0\cdot(k(X),k(s))$, that is, if and only if $k(X_\tau)=g_\tau$ for every $\tau$ and
$k(s)=a_0$. By (1) there is exactly one such $k$.
\end{proof}
% Lean: T.exists_universalFrame_testRing = Frames.exists_univFrame_testRing (UnivFrame y0 pi X s u: FrameRing, the
% image u0 of y0 under pi and FrChg X s u0 u): exists_frameRing_testRing (exists_frameRing_reduced, the ring A0,
% and reduced_finiteType_iff_testRing with FrameRing_of_algEquiv), FramedBaseChange.exists_frMap for u0 and
% existsUnique_frChg for u. T.universalFrame_classifies_framedPoints = Frames.univFrame_classifies: z0 from
% T.exists_framedPoint_baseChange_algebraMap, frameChange_exists_unique_of_same_point for (g, a0), key (frMap_trans,
% framedPoint_baseChange_unique_refl_iff for w = z0, and T.frameChange_comm_baseChange), and the universal property
% of FrameRing for existence and uniqueness. Both with SecBaseChangeSt. The Lean's FrChg g a y y' has y'.u tau i the
% combination of the y.u tau j with the coefficients g tau i j, as in the text.

\begin{lemma}\label{cv-uniformization-theta-rule-complex-glue-injective}
In the setting of Lemma~\ref{cv-uniformization-theta-marking}, let $\varphi\colon S\to T$ be an injective homomorphism
of $F'$-algebras, let $z$ and $z'$ be framed points of $X'$ over $S$ with the same point, and let $w$ be a framed point
over $T$ that is an image of $z$ and of $z'$ under $\varphi$ (Definition~\ref{cv-uniformization-theta-rule-values}).
Then $z=z'$.
\end{lemma}
\leanref{Development/AbelianVarietyTheory/Proofs/UnitaryExistence/FrameEqOfBaseChangeInjectiveStatement.lean}{GlobalTheory.T.frame_eq_of_baseChange_injective}
\begin{proof}
The frame change. Let $u_\ell$ and $u'_\ell$ be the tangent vectors of $z$ and $z'$ and $\vartheta$, $\vartheta'$ their
generators. By Lemma~\ref{cv-volume-key-lie}(2), for every embedding $\tau$ the vectors $u_{(\tau,i)}$ form a basis
of $\operatorname{Lie}_\tau$ at the point, and so do the $u'_{(\tau,i)}$; so $u'_{(\tau,i)}=\sum_jY_{\tau,ji}u_{(\tau,j)}$
for a unique invertible matrix $Y_\tau$ over $S$. The generator $\vartheta$ generates the module of sections freely,
as in the proof of Lemma~\ref{cv-uniformization-theta-rule-unique-base-change}(1); so $\vartheta'=a\vartheta$ and
$\vartheta=b\vartheta'$ with $a,b\in S$, $ab=1$, and $a$ is unique.

Base change. Base change along $\varphi$ carries a combination of tangent vectors to the combination of their base
changes with the coefficients mapped by $\varphi$, and the pullback of $a\vartheta$ is $\varphi(a)$ times that of
$\vartheta$. So the tangent vectors $w_\ell$ and the generator $\vartheta_w$ of $w$ satisfy
$w_{(\tau,i)}=\sum_j\varphi(Y_{\tau,ji})w_{(\tau,j)}$ and $\vartheta_w=\varphi(a)\vartheta_w$. By the uniqueness above,
for $w$ over $T$, $\varphi(Y_\tau)=1$ and $\varphi(a)=1$. As $\varphi$ is injective, $Y_\tau=1$ and $a=1$, so $z=z'$.
\end{proof}
% Lean: T.frame_eq_of_baseChange_injective = Frames.frPt_eq_of_frMap_injective (SecBaseChangeSt):
% T.frameChange_exists_unique_of_same_point (the frame change (g, a) from z to z'), T.frameChange_comm_baseChange
% (it is carried to the frame change from w to w), the uniqueness for w with frameChangeByMatrices_refl_and_comp, the
% injectivity of phi for a and the entries of g, and existsUnique_frameChangeByMatrices.

\begin{lemma}\label{cv-uniformization-theta-rule-complex-glue-local}
In the setting of Lemma~\ref{cv-uniformization-theta-marking}, let $S$ be an $F'$-algebra, $h_1,\dots,h_n\in S$
elements that generate the unit ideal, and $y$ and $y'$ framed points of $X'$ over $S$ with the same point. For each
$i$ let $z_i$ be a framed point over $S_{h_i}$ that is an image of $y$ and of $y'$ under $S\to S_{h_i}$
(Definition~\ref{cv-uniformization-theta-rule-values}). Then $y=y'$.
\end{lemma}
\leanref{Development/AbelianVarietyTheory/Proofs/UnitaryExistence/FramedPointReducedFromTestRing.lean}{GlobalTheory.Frames.TestDescent.eq_of_cover}
\begin{proof}
As in the proof of Lemma~\ref{cv-uniformization-theta-rule-complex-glue-injective}, $y'$ comes from $y$ by invertible
matrices $Y_\tau$ and a unit $a$, and the images of $Y_\tau$ and $a$ in each $S_{h_i}$ are $1$, as $z_i$ is an image of
both. An element of $S$ whose image in every $S_{h_i}$ is $0$ is $0$: it is killed by a power of each $h_i$, and these
powers generate the unit ideal. So $Y_\tau=1$ and $a=1$, and $y=y'$.
\end{proof}
% Lean: TestDescent.eq_of_cover: frameChange_exists_unique_of_same_point for (g, a), key (the frame change carried to
% each z i is the identity, by frameChange_comm_baseChange and the uniqueness), eq_of_away for the entries of g and for
% a, and existsUnique_frameChangeByMatrices.

\begin{lemma}\label{cv-uniformization-theta-rule-complex-glue-cover}
In the setting of Lemma~\ref{cv-uniformization-theta-marking}, let $S$ be an $F'$-algebra, $x$ a point of $X'$ over
$S$, and $h_1,\dots,h_n\in S$ elements that generate the unit ideal. For each $i$ let $y_i$ be a framed point over
$S_{h_i}$ whose point is the restriction of $x$, and suppose that for all $i$ and $j$ every image of $y_i$ over
$S_{h_ih_j}$ is an image of $y_j$, under the maps induced by the identity of $S$
(Definition~\ref{cv-uniformization-theta-rule-values}). Then there is a framed point $y_0$ over $S$ with the point
$x$ such that each $y_i$ is an image of $y_0$ under $S\to S_{h_i}$.
\end{lemma}
\leanref{Development/AbelianVarietyTheory/Proofs/UnitaryExistence/FramedPointsGluePrincipalCoverStatement.lean}{GlobalTheory.T.framedPoints_glue_principal_cover}
\begin{proof}
Let $A_x$ be the abelian scheme of $x$, with the line bundle $L_x$; its part over the open set $D(h_i)$ of
$\operatorname{Spec}S$ is the abelian scheme of the restriction of $x$. For $i$ and $j$ let $w_{ij}$ be an image of $y_i$
over $S_{h_ih_j}$ (Lemma~\ref{cv-uniformization-theta-rule-unique-base-change}(1)); it is an image of $y_j$ as well. So
the tangent vectors of $y_i$ and of $y_j$ at each letter have the same base change to $S_{h_ih_j}$, and the generators
of $y_i$ and $y_j$ have the same restriction to the part of $A_x$ over $D(h_ih_j)$.

The tangent vectors. Fix a letter $\ell$. The tangent vectors of the $y_i$ at $\ell$ are morphisms
$\operatorname{Spec}S_{h_i}[\varepsilon]\to A_x$ that agree over $\operatorname{Spec}S_{h_ih_j}[\varepsilon]$. These open
sets cover $\operatorname{Spec}S[\varepsilon]$, so the morphisms glue to a morphism $u_\ell\colon\operatorname{Spec}S[\varepsilon]\to
A_x$ over $\operatorname{Spec}S$; it is a tangent vector, and it lies in the eigenspace of the letter, as both are
equalities of morphisms that hold on the cover. The Lie algebra $\operatorname{Lie}(x)$ is finite locally free over $S$
and its formation commutes with base change (Lemma~\ref{cv-volume-key-lie}(1) and its proof), and the $u_\ell$ form a
basis of it after base change to every $S_{h_i}$; so the map $S^{\{\text{letters}\}}\to\operatorname{Lie}(x)$ that they
define is an isomorphism, as it is after localizing at every $h_i$. By Lemma~\ref{cv-volume-key-lie}(2), the $u_\ell$
at the letters of each embedding form a basis of its eigenspace.

The generator. The generators $\vartheta_i$ of the $y_i$ are sections of $L_x$ over the parts of $A_x$ over the
$D(h_i)$, which cover $A_x$, and they agree on the overlaps; so they glue to a section $\vartheta$ of $L_x$ over $A_x$.
Each $\vartheta_i$ generates, so it vanishes on no fibre (Lemma~\ref{cv-volume-key-generator}(3)); every point of
$\operatorname{Spec}S$ lies in some $D(h_i)$, so by part (1) of that lemma $\vartheta$ vanishes on no fibre, and by
part (3) it generates the sections of $L_x$.

So $y_0=(x,(u_\ell),\vartheta)$ is a framed point over $S$, and by construction each $y_i$ is an image of $y_0$ under
$S\to S_{h_i}$.
\end{proof}
% Lean: T.framedPoints_glue_principal_cover = Frames.frPt_glue (SecBaseChangeSt): ovl (an image z over S_(hi hj) and
% framedPoint_baseChange_unique_refl_iff for the vectors and sections of y i and y j), nf, hu with
% T.existsUnique_glue_dualNumber for the tangent vectors, eig0 with dual_ext, tb with T.tanBasis_of_local_tanBasis and
% T.eigenBasis_iff_tangentBasis, T.existsUnique_glue_sections_principal_cover for the section, nz with
% secBC_nonvanishing_iff and T.fibreSectionNowhereZero_of_generates, and T.fibreSectionGenerates_of_nowhereZero
% (Lemma cv-volume-key-generator(1) and (3)).

\begin{lemma}\label{cv-uniformization-theta-rule-complex-glue}
In the setting of Lemma~\ref{cv-uniformization-theta-marking}, let $y$ be a framed point of $X'$ over an
$F'$-algebra $S$, let $S_1\subset S$ be an $F'$-subalgebra and $x_1$ a point of $X'$ over $S_1$ whose composition with
$\operatorname{Spec}S\to\operatorname{Spec}S_1$ is the point under $y$, and let $h_1,\dots,h_n\in S_1$ generate the unit
ideal of $S_1$ and of $S$. For each $i$ let $z_i$ be an image of $y$ over $S_{h_i}$ and $v_i$ a framed point over
$(S_1)_{h_i}$ whose point is the restriction of $x_1$, such that $z_i$ is an image of $v_i$ under the map
$(S_1)_{h_i}\to S_{h_i}$ induced by the inclusion. Then there is a framed point $y_1$ over $S_1$ with the point $x_1$ of
which $y$ is an image under the inclusion.
\end{lemma}
\leanref{Development/AbelianVarietyTheory/Proofs/UnitaryExistence/FramedPointReducedFromTestRing.lean}{GlobalTheory.Frames.TestDescent.exists_descent_of_local}
\begin{proof}
Write $\iota$ for the inclusion $S_1\to S$, and $\kappa_i\colon(S_1)_{h_i}\to S_{h_i}$ and
$\kappa_{ij}\colon(S_1)_{h_ih_j}\to S_{h_ih_j}$ for the maps it induces; they are injective, as localization is exact.
Images are taken under the maps induced by the identities of $S_1$ and $S$ unless said otherwise, and the parts of
Lemma~\ref{cv-uniformization-theta-rule-unique-base-change} are used throughout: images exist (1), and they compose
(3) and cancel (4).

The overlaps. Let $i$ and $j$ be indices, $u$ an image of $v_i$ and $u'$ an image of $v_j$ over $(S_1)_{h_ih_j}$, and
$w$ an image of $y$ over $S_{h_ih_j}$. As $z_i$ is an image of $y$, $w$ is an image of $z_i$ under
$S_{h_i}\to S_{h_ih_j}$ by (4). As $z_i$ is an image of $v_i$ under $\kappa_i$, $w$ is an image of $v_i$ under the
composite $(S_1)_{h_i}\to S_{h_ih_j}$ by (3), which is also the composite of $(S_1)_{h_i}\to(S_1)_{h_ih_j}$ and
$\kappa_{ij}$; so $w$ is an image of $u$ under $\kappa_{ij}$ by (4). In the same way $w$ is an image of $u'$ under
$\kappa_{ij}$. The points of $u$ and $u'$ are both the restriction of $x_1$, so $u=u'$ by
Lemma~\ref{cv-uniformization-theta-rule-complex-glue-injective}. Hence every image of $v_i$ over $(S_1)_{h_ih_j}$ is an
image of $v_j$.

The gluing. By Lemma~\ref{cv-uniformization-theta-rule-complex-glue-cover} there is a framed point $y_1$ over $S_1$ with
the point $x_1$ such that each $v_i$ is an image of $y_1$ over $(S_1)_{h_i}$.

The comparison. Let $w$ be an image of $y_1$ under $\iota$, by (1). Its point is $x_1$ composed with
$\operatorname{Spec}\iota$, the point under $y$. For each $i$, $z_i$ is an image of $v_i$ under $\kappa_i$, so by (3) it
is an image of $y_1$ under the composite $S_1\to S_{h_i}$, which is also the composite of $\iota$ and $S\to S_{h_i}$;
so $z_i$ is an image of $w$ over $S_{h_i}$ by (4). It is an image of $y$ as well. As the $h_i$ generate the unit ideal
of $S$, Lemma~\ref{cv-uniformization-theta-rule-complex-glue-local} gives $w=y$. So $y$ is an image of $y_1$ under
$\iota$.
\end{proof}
% Lean: TestDescent.exists_descent_of_local (SecBaseChangeSt): T.framedPoints_glue_principal_cover (Lemma
% cv-uniformization-theta-rule-complex-glue-cover), with the compatibility on the overlaps from exists_away_val (the
% maps (S1)_(hi hj) -> S_(hi hj) are injective), frMap_restrict (the compositions and cancellations of the text),
% x_restrict and T.frame_eq_of_baseChange_injective (Lemma cv-uniformization-theta-rule-complex-glue-injective); then
% w from T.exists_framedPoint_baseChange_algebraMap and eq_of_cover (Lemma
% cv-uniformization-theta-rule-complex-glue-local) with cancel and trans. The Lean also asks the maps kappa_i, with
% kappa_i(s) = s on S1, and the existence of the z_i, which the text takes from the images.

\begin{lemma}\label{cv-uniformization-theta-rule-complex-descent}
In the setting of Lemma~\ref{cv-uniformization-theta-marking}, let $y$ be a framed point of $X'$ over an
$F'$-algebra $S$ (Definition~\ref{cv-uniformization-theta-rule-values}). Then there are a finitely generated
$F'$-subalgebra $S_1\subset S$ and a framed point $y_1$ over $S_1$ of which $y$ is an image under the inclusion.
\end{lemma}
\leanref{Development/AbelianVarietyTheory/Proofs/UnitaryExistence/FramedPointReducedFromTestRing.lean}{GlobalTheory.Frames.TestDescent.exists_fg_descent}
\begin{proof}
The point. The scheme $X'$ is of finite type over $F'$, and $S$ is the union of its finitely generated
$F'$-subalgebras; so the point under $y$ is the composition of $\operatorname{Spec}S\to\operatorname{Spec}S_0$ with a
point $x_0$ of $X'$ over a finitely generated $F'$-subalgebra $S_0$ of $S$ \cite[\S8]{Gro66}.

Local universal framed points. Let $g_i$, $R_i$, $j_i$ and $y_{0,i}$ be given by
Lemma~\ref{cv-uniformization-theta-rule-complex-charts} for $x_0$, and let $h_i$ be the image of $g_i$ in $S$; the
$h_i$ generate the unit ideal of $S$. For each $i$ let $A_i$, $\pi_i$ and $u_i$ be given by
Lemma~\ref{cv-uniformization-theta-rule-complex-universal} for $y_{0,i}$, and let $z_i$ be an image of $y$ over $S_{h_i}$
(Lemma~\ref{cv-uniformization-theta-rule-unique-base-change}(1)). As $h_i$ is a unit in $S_{h_i}$, the point under $z_i$
is $j_i\circ\operatorname{Spec}\psi_i$ for a homomorphism $\psi_i\colon R_i\to S_{h_i}$, which is the point under
$y_{0,i}$ composed with $\operatorname{Spec}\psi_i$. So, by part (3) of that lemma, $z_i$ is an image of $u_i$ under a
homomorphism $a_i\colon A_i\to S_{h_i}$ with $a_i\circ\pi_i=\psi_i$.

A finitely generated subalgebra. Each $A_i$ is finitely generated over $F'$, and each element of $S_{h_i}$ is a quotient
of an element of $S$ by a power of $h_i$. So there is a finitely generated $F'$-subalgebra $S_1$ of $S$ containing
$S_0$ such that each $a_i$ is the composition of a homomorphism $a'_i\colon A_i\to(S_1)_{h_i}$ with the map
$(S_1)_{h_i}\to S_{h_i}$ induced by the inclusion, which is injective, and such that the $h_i$ generate the unit ideal
of $S_1$: add to $S_0$ the numerators of the images of generators of the $A_i$ and the coefficients of an expression
of $1$ in the $h_i$. Let $x_1$ be $x_0$ composed with $\operatorname{Spec}S_1\to\operatorname{Spec}S_0$, and let
$v_i$ be an image of $u_i$ under $a'_i$ (Lemma~\ref{cv-uniformization-theta-rule-unique-base-change}(1)). Its point is
the restriction of $x_1$: by Lemma~\ref{cv-uniformization-theta-rule-complex-charts} both are $j_i$ composed with the
spectrum of a homomorphism $R_i\to(S_1)_{h_i}$, the two homomorphisms agree after the injective map to $S_{h_i}$, and
$j_i$ is an open immersion. By
Lemma~\ref{cv-uniformization-theta-rule-unique-base-change}(4), $z_i$ is an image of $v_i$.

Gluing. Lemma~\ref{cv-uniformization-theta-rule-complex-glue} gives a framed point $y_1$ over $S_1$ of which $y$ is an
image under the inclusion.
\end{proof}
% Lean: TestDescent.exists_fg_descent (SecBaseChangeSt): moduli_point_from_fg_subalgebra (the point, about 50 lines),
% exists_chart_cover (Lemma cv-uniformization-theta-rule-complex-charts), T.exists_universalFrame_testRing (Lemma
% cv-uniformization-theta-rule-complex-universal), T.exists_framedPoint_baseChange_algebraMap for z_i,
% T.universalFrame_classifies_framedPoints for a_i, exists_fg_subalgebra_desc_away (the subalgebra S1, the
% factorizations a'_i through kappa_i and the unit ideal, about 100 lines), v_i, x_of_descended and cancel, and
% exists_descent_of_local (Lemma cv-uniformization-theta-rule-complex-glue). Size beyond the laws of abelian schemes,
% chart independence and the lemmas before this one from Lemma cv-uniformization-theta-rule-connection-charts on:
% about 240 lines beyond Lemmas cv-uniformization-theta-rule-complex-charts to -complex-glue.

\begin{lemma}\label{cv-uniformization-theta-rule-complex-tested}
In the setting of Lemma~\ref{cv-uniformization-theta-marking}, every framed point $y$ of $X'$ over a reduced
$F'$-algebra $S$ is an image, under a homomorphism of $F'$-algebras, of a framed point over a test ring, a quotient of
a polynomial ring over $F'$ by a radical ideal.
\end{lemma}
\leanref{Development/AbelianVarietyTheory/Proofs/UnitaryExistence/FramedPointReducedFromTestRing.lean}{GlobalTheory.Frames.allTest_of_secBaseChange}
\begin{proof}
By Lemma~\ref{cv-uniformization-theta-rule-complex-descent} there are a finitely generated $F'$-subalgebra $S_1$ of $S$
and a framed point $y_1$ over $S_1$ of which $y$ is an image under the inclusion $j\colon S_1\to S$. As $S$ is
reduced, so is $S_1$. Choose generators $s_1,\dots,s_k$ of $S_1$; the kernel of $F'[x_1,\dots,x_k]\to S_1$,
$x_i\mapsto s_i$, is a radical ideal, as $S_1$ is reduced, so there are a test ring $R$ and an isomorphism of
$F'$-algebras $e\colon R\to S_1$. By Lemma~\ref{cv-uniformization-theta-rule-unique-base-change}(1), $y_1$ has an image
$y'$ under $e^{-1}$. As $y$ is an image of $y_1$ under $j=(j\circ e)\circ e^{-1}$, part (4) of that lemma shows that
$y$ is an image of $y'$ under $j\circ e$.
\end{proof}
% Lean: Frames.allTest_of_secBaseChange (SecBaseChangeSt, from lawsFieldsSeparately.2.1): TestDescent.exists_fg_descent
% (Lemma cv-uniformization-theta-rule-complex-descent), isReduced_of_injective, reduced_finiteType_iff_testRing,
% T.exists_framedPoint_baseChange_algebraMap (the existence of images, here from the base change law) and
% FramedBaseChange.cancel. Size beyond Lemma cv-uniformization-theta-rule-complex-descent: about 20 lines.

\begin{lemma}\label{cv-uniformization-theta-rule-complex}
In the setting of Lemma~\ref{cv-uniformization-theta-marking}, let $\nabla$ be a canonical Taylor family, $n\ge0$,
$c_1,\dots,c_n\in L$, $a$ a family of nonnegative integers indexed by the letters, $P$ a polynomial harmonic of row
degrees $a$, $f$ a theta rule of $c_1,\dots,c_n$ and $P$ for $\nabla$
(Definition~\ref{cv-uniformization-theta-rule-values}(4) and~(6)), and $y$ a framed point of $X'$ over a reduced
$F'$-algebra $S$. Then the value $f(y)$ of $f$ at $y$, read through a framed point over a test ring of which $y$ is an
image, is a theta value of $c_1,\dots,c_n$ and $P$ at $y$ for $\nabla$. In particular, if $f$ is a rational section,
its reading $f_Z(x)$ at a frame $x$ of a component $Z$ is a theta value at the framed point at $x$.
\end{lemma}
\leanref{Development/AbelianVarietyTheory/Proofs/SiegelThetaTransformation/ReadingEqThP.lean}{GlobalTheory.Dict.ThRead.thValAt_evR}
\begin{proof}
By Lemma~\ref{cv-uniformization-theta-rule-complex-tested}, $y$ is an image of a framed point $y_0$ over a test ring
$R$ under a homomorphism $\varphi\colon R\to S$, and the value of $f$ at $y$ read through $y_0$ is $\varphi(f(y_0))$.
By Lemma~\ref{cv-uniformization-theta-rule-unique}, $y$ has a theta value $w$ of $c_1,\dots,c_n$ and $P$. As $f(y_0)$
is a theta value at $y_0$, Lemma~\ref{cv-uniformization-theta-rule-weight-natural} gives $\varphi(f(y_0))=w$. The
reading $f_Z(x)$ is the value of $f$ at the framed point at $x$, a framed point over $\mathbf C$, which is a reduced
$F'$-algebra.
\end{proof}
% Lean: ThRead.thValAt_evR (15 lines): Frames.allTest_of_secBaseChange (Lemma
% cv-uniformization-theta-rule-complex-tested) with lawsFieldsSeparately.2.1, thValTot_of_chartIndep for w (Lemma
% cv-uniformization-theta-rule-unique), and thetaValue_natural (Lemma cv-uniformization-theta-rule-weight-natural) for
% evR f y = w, where evR (evalAt) reads f through the presentation chosen by dif_pos. The reading rdU U f k x is evR f
% at the framed point unC U k x. Size beyond the theta values and their naturality: about 5,900 lines, nearly all in
% Lemma cv-uniformization-theta-rule-complex-descent.

\begin{lemma}\label{cv-uniformization-theta-rule-translate}
Let $X$ be data (1) to (6) of flag data of rank $r$ over a CM field $E$ (Definition~\ref{is-onto-share-flag-data}), with
set $\Psi$, a symplectic basis of its module and an indexing (Definition~\ref{is-eigenvalue-conjugate-frame-fibres}(1)
and~(2)), and put $g=|\Psi|r$. Let $n\ge0$ be even, $x\in\mathcal D_X$, $\varphi$ a Schwartz--Bruhat function on
$M_{n\times g}(\mathbf Q)$ (Definition~\ref{is-eigenvalue-conjugate-frame-fibres}(6)) and $y\in M_{n\times r}(E)$ with
rows $y_1,\dots,y_n$. Then $\rho_X(y)\varphi$ is Schwartz--Bruhat
(Definition~\ref{is-eigenvalue-conjugate-frame-heisenberg}(1)), and for every $z=(z_1,\dots,z_n)\in(\mathbf C^g)^n$
\[
 \theta_{\rho_X(y)\varphi}(x;z)=\exp\Big(-\pi\sum_{k=1}^nH_x\big(z_k+\tfrac12p_x(y_k),p_x(y_k)\big)\Big)\,
 \theta_\varphi\big(x;z_1+p_x(y_1),\dots,z_n+p_x(y_n)\big).
\]
\end{lemma}
\leanref{Development/AbelianVarietyTheory/Proofs/SiegelThetaTransformation/HMatMulVecPerCrd.lean}{GlobalTheory.SW.heisG_laws}
% Lean: the first and last clauses of SW.heisG_laws, a computation with the series: heisG is heisS at the coordinates
% of y in the symplectic basis, T.heisenbergOperatorRationalPoint_laws (h1, h2) for thS, hMat_mulVec_per_crd for the
% periods of the coordinates and heis_exponent for the exponents (about 1,750 lines with what it uses, all shared
% with the vacuum lemmas of Chapter 8).

\begin{lemma}\label{cv-uniformization-theta-rule-product}
Let $X$, $n$, $x$, $\varphi$ and $y$ be as in Lemma~\ref{cv-uniformization-theta-rule-translate}, let
$\vartheta_1,\dots,\vartheta_n$ be functions on $\mathbf C^g$, and let $\kappa\in\mathbf C$ satisfy
$\prod_b\vartheta_b(z_b)=\kappa\,\theta_\varphi(x;z)$ for every $z\in(\mathbf C^g)^n$. Then for every $z$
\[
 \prod_{b=1}^ne^{-\pi H_x(z_b+\frac12p_x(y_b),\,p_x(y_b))}\,\vartheta_b\big(z_b+p_x(y_b)\big)
 =\kappa\,\theta_{\rho_X(y)\varphi}(x;z).
\]
\end{lemma}
\leanref{Development/AbelianVarietyTheory/Proofs/SiegelThetaTransformation/FockHxDifferentiable.lean}{GlobalTheory.Dict.fock_laws}
\begin{proof}
By Lemma~\ref{cv-uniformization-theta-rule-translate},
$\kappa\,\theta_{\rho_X(y)\varphi}(x;z)=\exp\big(-\pi\sum_kH_x(z_k+\frac12p_x(y_k),p_x(y_k))\big)\,
\kappa\,\theta_\varphi(x;z')$, where $z'$ has the components $z_k+p_x(y_k)$. By the hypothesis at $z'$,
$\kappa\,\theta_\varphi(x;z')=\prod_b\vartheta_b(z_b+p_x(y_b))$. The exponential of the sum is the product of the
exponentials.
\end{proof}
% Lean: the last clause of Dict.fock_laws, with fock x 1 a th z = exp(-pi hx(z + a/2, a)) th(z + a): heisG_laws at z,
% the hypothesis h at z + pv y, Finset.prod_mul_distrib and Complex.exp_sum.

\begin{lemma}\label{cv-uniformization-theta-rule-fibre-reading-generator}
In the setting of Lemma~\ref{cv-uniformization-theta-marking}, let $X_0$ be data (1) to (6) of flag data of rank $r$
over $E$ (Definition~\ref{is-onto-share-flag-data}), with set $\Psi$, a symplectic basis of its module and an
indexing (Definition~\ref{is-eigenvalue-conjugate-frame-fibres}(1) and~(2)), and put $g=|\Psi|r$. Let $y$ be a framed
point of $X'$ over $\mathbf C$, an $F'$-algebra through an embedding of $F'$, let $j\mapsto\ell_j$ be a bijection from
$\{1,\dots,g\}$ onto the letters, $x\in\mathcal D_{X_0}$, $\Lambda\subset E^r$ an $\mathcal O_E$-submodule self-dual
for $X_0$, and $(\operatorname{ex},\operatorname{fr},\chi)$ an analytic fibre at $y$ for these data, $x$ and $\Lambda$
(Definition~\ref{is-eigenvalue-conjugate-frame-fibres}(3)), with the trivializing opens $U_\beta$ of $\mathcal L_y$.
Write $P=p_x$ and $H=H_x$. Let $K$ be a finite set, and for $k\in K$ let $\phi_k\colon A_y\to A_y$ be a morphism,
$\varepsilon_k$ an integer and $\lambda_k\colon\mathbf C^g\to\mathbf C^g$ a holomorphic map with
$\phi_k(\operatorname{ex}(z))=\operatorname{ex}(\lambda_k(z))$ for every $z$. Let $t$ be a generator of
$\bigotimes_k\phi_k^*\mathcal L_y^{\varepsilon_k}$ on $A_y$.
\begin{enumerate}
\item There is a holomorphic function $v$ on $\mathbf C^g$ without zeros such that, for every $z$ and all indices
  $\beta_k$ with $\operatorname{ex}(\lambda_k(z))\in U_{\beta_k}$, the value of $t$ at $\operatorname{ex}(z)$ in the
  trivializations over the $U_{\beta_k}$ is $v(z)\prod_k\operatorname{fr}_{\beta_k}(\lambda_k(z))^{\varepsilon_k}$. For
  $z,z'\in\mathbf C^g$ with $\operatorname{ex}(z')=\operatorname{ex}(z)$ and $u_k\in\Lambda$ with
  $\lambda_k(z')=\lambda_k(z)+Pu_k$ for every $k$,
  \[
   v(z')=v(z)\prod_k\Big(\chi(u_k)\exp\big(\pi H(\lambda_kz,Pu_k)+\tfrac\pi2H(Pu_k,Pu_k)\big)\Big)^{\varepsilon_k}.
  \]
\item Suppose moreover that each $\lambda_k$ is $\mathbf C$-linear with $\lambda_k(P(\Lambda))\subset P(\Lambda)$, and
  that $\sum_k\varepsilon_kH(\lambda_kz,\lambda_kw)=0$ for all $z,w\in\mathbf C^g$. Then $v$ is constant, and
  $\prod_k\chi(u_k)^{\varepsilon_k}=1$ for all $u,u_k\in\Lambda$ with $\lambda_k(Pu)=Pu_k$ for every $k$.
\end{enumerate}
\end{lemma}
\leanref{Development/AbelianVarietyTheory/Proofs/CuspFormEvaluations/ExistsHolomorphicGeneratorReadingAutomorphyStatement.lean}{GlobalTheory.T.exists_holomorphic_generator_reading_automorphy}
\leanref{Development/AbelianVarietyTheory/Proofs/CuspFormEvaluations/AbelianHolomorphicFrameReadingConstOfFormsCancelStatement.lean}{GlobalTheory.T.abelianHolomorphicFrame_reading_const_of_forms_cancel}
\begin{proof}
Write $p=\operatorname{ex}$, $F_\beta=\operatorname{fr}_\beta$, $g_{\beta\beta'}$ for the transition functions of
$\mathcal L_y$ and $O_\beta$ for the open set of the $z$ with $p(z)\in U_\beta$.

(1) For $z$ and indices $\beta=(\beta_k)$ with $\lambda_k(z)\in O_{\beta_k}$ for every $k$, let $v_\beta(z)$ be the value
of $t$ at $p(z)$ in the trivializations over the $U_{\beta_k}$, divided by $\prod_kF_{\beta_k}(\lambda_k(z))^{\varepsilon_k}$,
which is not $0$. For other indices $\beta'$ the value of $t$ is multiplied by
$\prod_kg_{\beta'_k\beta_k}(\phi_k(p(z)))^{\varepsilon_k}$, and so is the product, as
$F_{\beta'_k}=g_{\beta'_k\beta_k}(p(\cdot))F_{\beta_k}$ and $\phi_k(p(z))=p(\lambda_k(z))$; so $v_\beta(z)=v_{\beta'}(z)$,
and these numbers define a function $v$ on $\mathbf C^g$. On the open set where $\lambda_k(z)\in O_{\beta_k}$ for
every $k$ the value of $t$ at $p(z)$ is holomorphic in $z$, as $p$ is holomorphic (as in the proof of
Lemma~\ref{cv-uniformization-theta-half-sign-fibre}), and so is each $F_{\beta_k}\circ\lambda_k$; so $v$ is holomorphic,
and it has no zeros, as $t$ is a generator. Let $z$, $z'$ and the $u_k$ be as in the statement. The law (c) of the
analytic fibre for $u_k$ at $\lambda_k(z)$ gives $F_{\beta_k}(\lambda_k(z'))=c_k^{-1}F_{\beta_k}(\lambda_k(z))$, with
$c_k=\chi(u_k)\exp\big(\pi H(\lambda_kz,Pu_k)+\frac\pi2H(Pu_k,Pu_k)\big)$, for every $\beta_k$ with
$\lambda_k(z)\in O_{\beta_k}$, which then also contains $\lambda_k(z')$. As $t$ has the same value at $p(z')$ and
$p(z)$, $v(z')\prod_kc_k^{-\varepsilon_k}=v(z)$, which is the claim.

(2) Let $u,u_k$ be as in the statement and $z'=z+Pu$. As $Pu$ lies in the kernel of $p$, $p(z')=p(z)$, and
$\lambda_k(z')=\lambda_k(z)+Pu_k$ by linearity. The hypothesis for $w=Pu$, at $z$ and at $Pu$, gives
$\sum_k\varepsilon_kH(\lambda_kz,Pu_k)=0$ and $\sum_k\varepsilon_kH(Pu_k,Pu_k)=0$. So (1) gives $v(z+Pu)=\rho\,v(z)$
with $\rho=\prod_k\chi(u_k)^{\varepsilon_k}$, and $|\rho|=1$ by (c). As $\Lambda$ is self-dual, the
vectors $Pu$, $u\in\Lambda$, span $\mathbf C^g$ over $\mathbf R$, and a bound $R$ as in
Lemma~\ref{cv-uniformization-theta-half-sign-liouville} exists for the group $P(\Lambda)$ (Steps 0 and 4 of the proof
of Lemma~\ref{cv-uniformization-theta-half-sign-fibre}). That lemma, for $V=\mathbf C^g$, $S=P(\Lambda)$, $\psi=v$ and
$D=0$, shows that $v$ is constant. As $v$ has no zeros, $\rho=1$.
\end{proof}
% Lean: T.exists_holomorphic_generator_reading_automorphy = AnValGenRd.exists_reading, for an analytic frame Phi
% (AnFr) of any abelian scheme, any scheme Z with a holomorphic family q of points and differentiable maps p_k with
% q z >> phi_k = e(p_k z): v with GenRd (the reading), its uniqueness, and the law v z' = v z prod fac(a_k, p_k z)^eps_k
% when q z' = q z and p_k z' = p_k z + a_k with a_k in the lattice. T.abelianHolomorphicFrame_reading_const_of_forms_cancel
% = AnVal.genRd_const: hprod (the exponentials cancel), hstep, hnorm (|chi| = 1) and norm_periodic_const_range (a
% holomorphic function with periodic absolute value under a group spanning over R is constant), then the product of
% chi at the period. The text states both for the analytic frame of an analytic fibre with Z = A_y, q = ex and linear
% maps, the case used, and gives the Liouville step through Lemma cv-uniformization-theta-half-sign-liouville with
% D = 0 (FrameUnits.form_eq_zero_and_const) in place of norm_periodic_const_range.

\begin{lemma}\label{cv-uniformization-theta-rule-fibre-reading-sign}
In the setting of Lemma~\ref{cv-uniformization-theta-marking}, let $X_0$ be data (1) to (6) of flag data of rank $r$
over $E$ (Definition~\ref{is-onto-share-flag-data}), with set $\Psi$, a symplectic basis of its module and an
indexing (Definition~\ref{is-eigenvalue-conjugate-frame-fibres}(1) and~(2)), and put $g=|\Psi|r$. Let $y$ be a framed
point of $X'$ over $\mathbf C$, an $F'$-algebra through an embedding of $F'$, let $j\mapsto\ell_j$ be a bijection from
$\{1,\dots,g\}$ onto the letters, $x\in\mathcal D_{X_0}$, $\Lambda\subset E^r$ an $\mathcal O_E$-submodule self-dual
for $X_0$, and $(\operatorname{ex},\operatorname{fr},\chi)$ an analytic fibre at $y$ for these data, $x$ and $\Lambda$
(Definition~\ref{is-eigenvalue-conjugate-frame-fibres}(3)). For $b\in\mathcal O_E$ and $z\in\mathbf C^g$ put
$b\cdot z=(\rho_j(b)z_j)_j$, with the row types $\rho_j$ of the indexing. Then for every $v\in\Lambda$,
\[
 \chi(v)=\exp\big(\pi i\operatorname{Im}H_x\big(p_x(v),w_0\cdot p_x(v)\big)\big)\qquad\text{and}\qquad\chi(v)^2=1.
\]
\end{lemma}
\leanref{Development/AbelianVarietyTheory/Proofs/CuspFormEvaluations/HalfIntegralGeneratorTraceEqOne.lean}{GlobalTheory.AnVal.chi_eq_sign}
\leanref{Development/AbelianVarietyTheory/Proofs/SiegelThetaTransformation/ExistsAnalyticFrameOfAnalyticFibre.lean}{GlobalTheory.Dict.AnFib.exists_anPt}
\begin{proof}
Write $P=p_x$, $H=H_x$, $\langle\ ,\ \rangle=\langle\ ,\ \rangle_{X_0}$ and $Pv$ for $P(v)$. By
Lemma~\ref{is-eigenvalue-conjugate-frame-hermitian}(1) to~(3), whose proof uses only data (1) to (6) of $X_0$, $H$ is
Hermitian and linear in the first variable, $\operatorname{Im}H(Pu,Pv)=\langle u,v\rangle$ for $u,v\in E^r$,
$H(b\cdot z,w)=H(z,\bar b\cdot w)$ for $b\in\mathcal O_E$, and $P(bv)=b\cdot Pv$. By
Definition~\ref{is-eigenvalue-conjugate-frame-fibres}(3)(a), $\operatorname{ex}(b\cdot z)=b\,\operatorname{ex}(z)$. By
Step 1 of the proof of Lemma~\ref{cv-uniformization-theta-half-sign-fibre}, for $u,v\in\Lambda$,
$\chi(u+v)=\chi(u)\chi(v)\exp(\pi i\langle u,v\rangle)$, and $\chi(0)=1$.

The generator. Let $\mathcal A$ be the abelian scheme over $X'$, with the line bundle $\mathcal L$, and $o$ its identity
as a point with values in $\mathcal A$. For $J\subset\{1,2\}$ let $m_J\colon\mathcal A\to\mathcal A$ be the sum over
$k\in J$ of the identity ($k=1$) and the action of $w_0$ ($k=2$), and $\lambda_J$ the sum over $k\in J$ of the identity
and $z\mapsto w_0\cdot z$, so that $m_J(\operatorname{ex}(z))=\operatorname{ex}(\lambda_J(z))$. Then
$\mathcal L\otimes\Lambda(o,w_0o)^{-1}=\mathcal L\otimes\bigotimes_Jm_J^*\mathcal L^{-(-1)^{|J|}}$, and it has a generator
$t$, as $\mathcal A$ is normalized (Definition~\ref{cv-model-toroidal-semiabelian}(2)); restrict it to $A_y$. The maps
$\lambda_J$ and the identity carry $P(\Lambda)$ into itself, as $w_0\Lambda\subset\Lambda$.

The forms cancel. For $z,w\in\mathbf C^g$, expanding $H(z+w_0\cdot z,w+w_0\cdot w)$ by additivity,
\[
 H(z,w)-H(0,0)+H(z,w)+H(w_0\cdot z,w_0\cdot w)-H(z+w_0\cdot z,w+w_0\cdot w)=H(z,w)-H(z,w_0\cdot w)-H(w_0\cdot z,w).
\]
As $H(w_0\cdot z,w)=H(z,\bar w_0\cdot w)$, $H$ is additive in the second variable and $w_0+\bar w_0=1$, this is
$H(z,w)-H(z,(w_0+\bar w_0)\cdot w)=0$.

The sign. Let $v\in\Lambda$. Lemma~\ref{cv-uniformization-theta-rule-fibre-reading-generator}(2), for $t$, the
identity and the $m_J$ as morphisms, with the exponents $1$ and $-(-1)^{|J|}$ and the maps $1$ and $\lambda_J$, and
for $u=v$, with $v$ for the identity and the sum $u_J$ over $k\in J$ of $v$ ($k=1$) and $w_0v$ ($k=2$) for $m_J$,
gives
\[
 \chi(v)\,\chi(0)^{-1}\chi(v)\,\chi(w_0v)\,\chi(v+w_0v)^{-1}=1.
\]
With $\chi(0)=1$ and $\chi(v+w_0v)=\chi(v)\chi(w_0v)\exp(\pi i\langle v,w_0v\rangle)$, this is
$\chi(v)=\exp(\pi i\langle v,w_0v\rangle)=\exp\big(\pi i\operatorname{Im}H(Pv,w_0\cdot Pv)\big)$. As $\Lambda$ is
self-dual and $v,w_0v\in\Lambda$, $\langle v,w_0v\rangle$ is an integer, so $\chi(v)^2=\exp(2\pi i\langle v,w_0v\rangle)=1$.
\end{proof}
% Lean: AnFib.exists_anPt makes the analytic fibre an analytic point Phi (AnVal.AnPt: the analytic frame with lin,
% adj and der, the lattice {pv v : v in L}, Hm = hx, chi (pv v) = A.chi v); the text uses these facts directly, from
% the fibre and Lemma is-eigenvalue-conjugate-frame-hermitian. Then AnVal.chi_eq_sign:
% T.exists_normalisedIso_reading_one_analyticFrame (the generator wP of L (x) Lambda(o, w0 o)^-1 and its reading),
% hlat (w0 L in L), htr (halfIntegralGenerator_trace_eq_one: w0 + conj w0 = 1), the maps Lambda (Sum.elim id and
% psumL over the subsets of {id, tyAct w0}), the cancellation of the forms (Phi.adj and semichar_herm_add_left/right,
% the computation of the text), T.abelianHolomorphicFrame_reading_const_of_forms_cancel (Lemma
% cv-uniformization-theta-rule-fibre-reading-generator(2)) for prod chi^eps = 1, chi_add and chi_zero (the Lean's
% AnFr.chi_add; the text cites Step 1 of the proof of Lemma cv-uniformization-theta-half-sign-fibre, the same
% computation from fr_per), and herm_im_int for chi^2 = 1. The text takes the generator t without the
% normalization along the zero section that wP has; the proof does not use it.

\begin{lemma}\label{cv-uniformization-theta-rule-fibre-reading-pullback}
Let $y$, $x$, $\Lambda$ and the analytic fibre $(\operatorname{ex},\operatorname{fr},\chi)$ be as in
Lemma~\ref{cv-uniformization-theta-rule-fibre-reading-sign}, and let $N\ge1$. Let $s$ be a section of $\mathcal L_y$
with a theta reading $\vartheta$ (Definition~\ref{is-eigenvalue-conjugate-frame-fibres}(4)), and let $s_1$ be the section
of $\mathcal L_y^{N^2}$ that corresponds to $[N]^*s$ under the isomorphism $[N]^*\mathcal L_y\cong\mathcal L_y^{N^2}$
compatible with the rigidifications. Then $z\mapsto\vartheta(Nz)$ is a theta reading of $s_1$.
\end{lemma}
\leanref{Development/AbelianVarietyTheory/Proofs/CuspFormEvaluations/NormalisedPullbackMulNReadsThetaStatement.lean}{GlobalTheory.T.normalisedPullbackMulN_reads_theta}
\begin{proof}
Write $p=\operatorname{ex}$ and $F_\beta=\operatorname{fr}_\beta$. Let $\alpha$ be the isomorphism
$[N]^*\mathcal L_y\cong\mathcal L_y^{N^2}$ compatible with the rigidifications, a generator of
$\mathcal L_y^{N^2}\otimes[N]^*\mathcal L_y^{-1}$.

The reading of $\alpha$. Apply Lemma~\ref{cv-uniformization-theta-rule-fibre-reading-generator} to $\alpha$, with the
morphisms $1$ and $[N]$, the exponents $N^2$ and $-1$, and the maps $z\mapsto z$ and $z\mapsto Nz$: as $p$ is a
homomorphism, $N\,p(z)=p(Nz)$, and $NP(\Lambda)\subset P(\Lambda)$. The forms cancel:
$N^2H_x(z,w)-H_x(Nz,Nw)=0$, as $N$ is real. So by that lemma there is a constant $c$ such that the value of $\alpha$
at $p(z)$, in the trivializations over $U_\beta$ at $p(z)$ and over $U_{\beta'}$ at $p(Nz)$, is
$c\,F_\beta(z)^{N^2}F_{\beta'}(Nz)^{-1}$. At $z=0$, $p(0)$ is the origin, and $\alpha$ carries the rigidifying section
at the origin to its $N^2$-th power; its values in $U_\beta$ are $F_\beta(0)$ and $F_\beta(0)^{N^2}$
(Definition~\ref{is-eigenvalue-conjugate-frame-fibres}(3)(b)). So $c\,F_\beta(0)^{N^2}F_\beta(0)^{-1}F_\beta(0)=F_\beta(0)^{N^2}$,
and $c=1$.

The reading of $s_1$. The section $s_1$ is $\alpha([N]^*s)$. Its value at $p(z)$ in the trivialization over $U_\beta$
is the value of $\alpha$ there times the value of $s$ at $N\,p(z)=p(Nz)$ over $U_{\beta'}$, that is
\[
 F_\beta(z)^{N^2}F_{\beta'}(Nz)^{-1}\cdot\vartheta(Nz)F_{\beta'}(Nz)=\vartheta(Nz)F_\beta(z)^{N^2}.
\]
So $z\mapsto\vartheta(Nz)$ is a theta reading of $s_1$.
\end{proof}
% Lean: T.normalisedPullbackMulN_reads_theta = AnVal.rd_mulN, for an analytic frame (AnFr) of any abelian scheme with
% a rigidified line bundle at a complex point, a normalized alpha0 (IsNormMul) with its pullback alpha, and
% IsTwistPull for s_1; through AnFib.exists_anPt for the analytic fibre. The reading of alpha is mulN_genRd:
% T.abelianHolomorphicFrame_reading_const_of_forms_cancel (Lemma cv-uniformization-theta-rule-fibre-reading-generator)
% with Lambda = (id, N), hper and hherm (N^2 - N^2 = 0), then the constant 1 from the normalization at the zero
% section (hrig with fr_0, genRd_corr_mul_inv and genRd_rigPow). The reading of s_1 is AnFr.rd_twistPull with hpt.

\begin{lemma}\label{cv-uniformization-theta-rule-fibre-reading-symmetric}
Let $y$, $x$, $\Lambda$ and the analytic fibre $(\operatorname{ex},\operatorname{fr},\chi)$ be as in
Lemma~\ref{cv-uniformization-theta-rule-fibre-reading-sign}, and suppose that $\chi(v)^2=1$ for every $v\in\Lambda$. Let
$k$ be an integer, $n\ge1$ odd, and $w\in\mathbf C^g$ with $nw=2p_x(v)$ for some $v\in\Lambda$. Let $U$ be a symmetric
element of the theta group of $\mathcal L_y^k$ over $\mathbf C$ that lies over $\operatorname{ex}(w)$ and satisfies
$U^n=1$ (Definition~\ref{cv-uniformization-theta-rule-unique-translates-symmetric}), and let $s$ be a section of
$\mathcal L_y^k$ with a theta reading $\psi$. Then the function
\[
 z\mapsto\exp\big(-k\pi H_x\big(z+\tfrac12w,w\big)\big)\,\psi(z+w)
\]
is a theta reading of $Us$, the section that corresponds to the pullback of $s$ under the translation by
$\operatorname{ex}(w)$ through $U$.
\end{lemma}
\leanref{Development/AbelianVarietyTheory/Proofs/CuspFormEvaluations/SymmetricLiftOddOrderActsAsTranslationStatement.lean}{GlobalTheory.T.symmetric_lift_odd_order_acts_as_translation}
\begin{proof}
Write $P=p_x$, $H=H_x$, $p=\operatorname{ex}$ and $F_\beta=\operatorname{fr}_\beta$, and let $U=(p(w),\varphi)$, with
$\varphi$ a generator of $\mathcal L_y^k\otimes t_{p(w)}^*\mathcal L_y^{-k}$ on $A_y$. Values at a point $p(z)$ are
read in trivializations over sets $U_\beta$ containing the points involved; by
Lemma~\ref{is-eigenvalue-conjugate-frame-hermitian}(1), $H$ is Hermitian.

Step 1: the theta constant. Lemma~\ref{cv-uniformization-theta-rule-fibre-reading-generator}(1) for $\varphi$, with
the identity and the translation by $p(w)$, the exponents $k$ and $-k$ and the maps $z\mapsto z$ and $z\mapsto w+z$,
gives a holomorphic function $v$ without zeros with: the value of $\varphi$ at $p(z)$ is
$v(z)F_\beta(z)^kF_{\beta'}(w+z)^{-k}$. For $u\in\Lambda$ and $z'=z+Pu$, with $u_1=u_2=u$, it gives
$v(z+Pu)=v(z)\exp(-k\pi H(w,Pu))$. Put $\psi_0(z)=v(z)\exp(k\pi H(z,w))$, holomorphic without zeros. Then
$\psi_0(z+Pu)=\psi_0(z)\exp\big(k\pi(H(Pu,w)-H(w,Pu))\big)$, and the factor has absolute value $1$, as $H$ is
Hermitian. Lemma~\ref{cv-uniformization-theta-half-sign-liouville}, for $S=P(\Lambda)$ and $D=0$, as in the proof of
Lemma~\ref{cv-uniformization-theta-rule-fibre-reading-generator}(2), shows that $\psi_0$ is constant. So there is
$c\ne0$ with
\[
 v(z)=c\,\exp\big(-k\pi H(z+\tfrac12w,w)\big)\qquad(z\in\mathbf C^g).
\]

Step 2: the reading of $Us$. The value of $Us$ at $p(z)$ over $U_\beta$ is the value of $\varphi$ there times the value
of $s$ at $p(w+z)$ over $U_{\beta'}$, that is
$c\,e^{-k\pi H(z+\frac12w,w)}F_\beta(z)^kF_{\beta'}(w+z)^{-k}\cdot\psi(z+w)F_{\beta'}(w+z)^k$. So $c$ times the function
of the statement is a theta reading of $Us$.

Step 3: $c^2=1$. Let $\nu$ be a generator of $\mathcal L_y^k\otimes[-1]^*\mathcal L_y^{-k}$ with
$t_{-p(w)}^*\varphi\otimes\nu\otimes[-1]^*\varphi\otimes t_{-p(w)}^*\nu^{-1}=1$. By
Lemma~\ref{cv-uniformization-theta-rule-fibre-reading-generator}(2), with the maps $z\mapsto z$ and $z\mapsto-z$ and
the exponents $k$ and $-k$, whose forms cancel, the value of $\nu$ at $p(z)$ is $bF_\beta(z)^kF_{\beta'}(-z)^{-k}$ for
a constant $b$. The values at the origin $p(0)$ of the four factors are those of $\varphi$ at $p(-w)$, of $\nu$ at
$p(0)$, of $\varphi$ at $p(0)$ and of $\nu^{-1}$ at $p(-w)$:
\[
 c\,e^{\frac{k\pi}2H(w,w)}\frac{F(-w)^k}{F(0)^k},\qquad b,\qquad c\,e^{-\frac{k\pi}2H(w,w)}\frac{F(0)^k}{F(w)^k},\qquad
 b^{-1}\frac{F(w)^k}{F(-w)^k},
\]
with the indices of the trivializations left out. Their product is $c^2$, so $c^2=1$.

Step 4: $c^n=1$. As $U^n=1$, $\prod_{i=0}^{n-1}t_{ip(w)}^*\varphi=1$. Its value at $p(0)$ is the product of the values
of $\varphi$ at the points $p(iw)$,
\[
 \prod_{i=0}^{n-1}c\,e^{-k\pi H((i+\frac12)w,w)}\frac{F(iw)^k}{F((i+1)w)^k}=c^n\,e^{-\frac{k\pi}2n^2H(w,w)}
 \frac{F(0)^k}{F(nw)^k},
\]
as $\sum_{i<n}(i+\frac12)=\frac{n^2}2$. As $nw=P(2v)$, the law (c) of the analytic fibre at $0$ gives
$F(0)=\chi(2v)e^{\frac\pi2n^2H(w,w)}F(nw)$. So $c^n\chi(2v)^k=1$. By Step 1 of the proof of
Lemma~\ref{cv-uniformization-theta-half-sign-fibre}, $\chi(2v)=\chi(v)^2\exp(\pi i\langle v,v\rangle_{X_0})$, with
$\langle v,v\rangle_{X_0}=\operatorname{Im}H(Pv,Pv)=0$ as $H$ is Hermitian; so $\chi(2v)=\chi(v)^2=1$, and $c^n=1$.

Step 5: the conclusion. As $n$ is odd, $c=c^n(c^2)^{-(n-1)/2}=1$, and Step 2 gives the claim.
\end{proof}
% Lean: T.symmetric_lift_odd_order_acts_as_translation = AnVal.rd_fock_of_symLift, for an analytic frame (AnFr) of any
% abelian scheme with a rigidified line bundle at a complex point, with the hypothesis chi^2 = 1 on the lattice,
% SymLift k n (e w) U and IsActOf U s s'; the reading is AnVal.fock Hm k w psi; through AnFib.exists_anPt for the
% analytic fibre. Step 1 is T.exists_thetaConstant_of_thetaGroup = AnVal.exists_thConst (exists_reading, Lemma
% cv-uniformization-theta-rule-fibre-reading-generator(1), hper, and thConst_liouville for the constant), Step 2
% T.reading_of_thetaGroup_action = AnFr.rd_of_isActOf, Step 3 T.thetaConst_symmetric_square_eq_one =
% AnValSym.sq_const_eq_one (neg_const for nu, the values R1 to R4 at the origin and genRd_corr_mul_inv), Step 4
% T.thetaConstant_pow_eq_one_of_power_trivial = AnValThPow.thConst_pow (the values at the points U.pt^m, telescope and
% fr_per at 0) with chi(2 delta) = 1 from T.appellHumbert_frame_semicharacter (anFr_semichar), and Step 5 hc1. The
% text gives Steps 3 and 4 as computations of values at the origin; the Lean's proofs are these computations in
% charts of L.

\begin{lemma}\label{cv-uniformization-theta-rule-fibre-reading}
In the setting of Lemma~\ref{cv-uniformization-theta-marking}, let $X_0$ be data (1) to (6) of flag data of rank $r$
over $E$ (Definition~\ref{is-onto-share-flag-data}), with set $\Psi$, a symplectic basis of its module and an
indexing (Definition~\ref{is-eigenvalue-conjugate-frame-fibres}(1) and~(2)), and put $g=|\Psi|r$. Let $y$ be a framed
point of $X'$ over $\mathbf C$, an $F'$-algebra through an embedding of $F'$, let $j\mapsto\ell_j$ be a bijection from
$\{1,\dots,g\}$ onto the letters, $x\in\mathcal D_{X_0}$, $\Lambda\subset E^r$ an $\mathcal O_E$-submodule self-dual
for $X_0$, $(\operatorname{ex},\operatorname{fr},\chi)$ an analytic fibre at $y$ for these data, $x$ and $\Lambda$, and
$\vartheta$ a theta reading of the section of $y$ for it (Definition~\ref{is-eigenvalue-conjugate-frame-fibres}(3)
and~(4)). Let $N\ge1$ be odd, $a\in E^r$ with $N^2a\in2\Lambda$, $t\in A_y(\mathbf C)$, and $F$ an $N$-translate of
$\vartheta_y$ over $t$ with root $\operatorname{ex}(p_x(a))$
(Definition~\ref{cv-uniformization-theta-rule-values}(1)). Then the function
\[
 z\mapsto\exp\big(-\pi H_x\big(Nz+\tfrac12p_x(Na),p_x(Na)\big)\big)\,\vartheta\big(Nz+p_x(Na)\big)
\]
on $\mathbf C^g$ is a theta reading of $F$ for the analytic fibre.
\end{lemma}
\leanref{Development/AbelianVarietyTheory/Proofs/SiegelThetaTransformation/SectionsReadAsTranslatesAtScaledPointStatement.lean}{GlobalTheory.T.sections_read_as_translates_at_scaled_point}
\begin{proof}
Write $F=Us_1$ as in Definition~\ref{cv-uniformization-theta-rule-values}(1): $s_1$ is the section of
$\mathcal L_y^{N^2}$ that corresponds to $[N]^*\vartheta_y$ under the isomorphism $[N]^*\mathcal L_y\cong\mathcal L_y^{N^2}$
compatible with the rigidifications, and $U$ is a symmetric element of the theta group of $\mathcal L_y^{N^2}$ over
$\operatorname{ex}(p_x(a))$ with $U^{N^2}=1$. By Lemma~\ref{cv-uniformization-theta-rule-fibre-reading-pullback}, the
function $\psi(z)=\vartheta(Nz)$ is a theta reading of $s_1$. By Lemma~\ref{cv-uniformization-theta-rule-fibre-reading-sign},
$\chi(v)^2=1$ for every $v\in\Lambda$. Write $N^2a=2u$ with $u\in\Lambda$; as the periods are additive,
$N^2p_x(a)=2p_x(u)$. Lemma~\ref{cv-uniformization-theta-rule-fibre-reading-symmetric}, with $k=n=N^2$, which is odd,
$w=p_x(a)$ and $v=u$, shows that
\[
 z\mapsto\exp\big(-N^2\pi H_x\big(z+\tfrac12p_x(a),p_x(a)\big)\big)\,\vartheta\big(Nz+Np_x(a)\big)
\]
is a theta reading of $F$. The form $H_x$ is linear in its first variable and conjugate linear in its second, and $N$
is real, so $N^2H_x(z+\frac12p_x(a),p_x(a))=H_x(Nz+\frac12p_x(Na),p_x(Na))$, with $Np_x(a)=p_x(Na)$. This is the
claim.
\end{proof}
% Lean: the second clause of T.sections_read_as_translates_at_scaled_point = Dict.reads_isUUOver_isVtOver, with
% IsSympl, x in X.dom, SelfD X L, N odd, A.Reads 1 y.s theta, a with N^2 a = 2u for some u in L, and IsVtOver N with
% the condition t' = A.ex (pv a) on the root; the reading is fock x 1 (pv (N a)) theta (N z), with
% fock x 1 a th z = exp(-pi hx(z + a/2, a)) th(z + a). The first clause, on IsUUOver, is not used here. The proof:
% AnFib.exists_anPt for Phi, hchi from AnVal.chi_eq_sign (Lemma cv-uniformization-theta-rule-fibre-reading-sign), htw
% from T.normalisedPullbackMulN_reads_theta (Lemma cv-uniformization-theta-rule-fibre-reading-pullback), hsy from
% T.symmetric_lift_odd_order_acts_as_translation (Lemma cv-uniformization-theta-rule-fibre-reading-symmetric), hpv
% (hx_laws: pv (n a) = n pv a) and hfk, the scaling clause of Dict.fock_laws (fock x (N^2) p (phi(N .)) =
% (fock x 1 (N p) phi)(N .)), which is the last computation of the text. Size beyond the laws of abelian schemes, chart
% independence, the connection and Lemmas cv-uniformization-theta-rule-fibre-reading-sign to -symmetric: about 130
% lines, besides Dict.fock_laws, shared with Lemma cv-uniformization-theta-rule-product.

\begin{lemma}\label{cv-uniformization-theta-rule-fibre-taylor-chart}
In the setting of Lemma~\ref{cv-uniformization-theta-marking}, let $X_0$, $g$, $j\mapsto\ell_j$, $x$, $\Lambda$, the
framed point $y$ over $\mathbf C$ and the analytic fibre be as in
Lemma~\ref{cv-uniformization-theta-rule-fibre-reading-sign}. Then there is a chart datum at $y$
(Definition~\ref{cv-uniformization-theta-rule-values}(5)) of dimension $g$, with the group law $F(t,t')=t+t'$, the power
series $h=0$ and the matrix with $u_{\ell_j,j'}=\delta_{jj'}$, such that for every $k\ge0$, every section of
$\mathcal L_y^k$ and every theta reading $\vartheta$ of it (Definition~\ref{is-eigenvalue-conjugate-frame-fibres}(4)),
the Taylor series of $\vartheta$ at $0$ is a germ of the section.
\end{lemma}
\leanref{Development/AbelianVarietyTheory/Proofs/CuspFormEvaluations/ExistsAnalyticBiextensionChartComplexPointStatement.lean}{GlobalTheory.T.exists_analytic_biextension_chart_complex_point}
\begin{proof}
\openstep{The construction of the chart datum from the analytic fibre, with the coordinates of $\mathbf C^g$ as its
formal chart, generators from the functions $\operatorname{fr}_\beta$ and the gauge $0$, is not yet written out.}
% Lean: T.exists_analytic_biextension_chart_complex_point = AnVal.exists_analytic_chart, for an analytic point
% (AnVal.AnPt, from the analytic fibre by T.exists_analyticFrame_of_analyticFibre) and star = bar: a ChartAt C with
% C.d = g, C.h = 0, the additive group law, C.uvec the identity through el, and IsGerm for the Taylor series tay theta 0
% of every reading. Size beyond the laws of abelian schemes, chart independence, the connection and Lemmas
% cv-uniformization-theta-rule-fibre-reading-sign to -fibre-reading: about 11,600 lines.
\end{proof}

\begin{lemma}\label{cv-uniformization-theta-rule-fibre-taylor}
In the setting of Lemma~\ref{cv-uniformization-theta-marking}, let $X_0$, $g$, $j\mapsto\ell_j$, $x$, $\Lambda$, the
framed point $y$ over $\mathbf C$ and the analytic fibre be as in Lemma~\ref{cv-uniformization-theta-rule-fibre-reading},
with $\mathbf C$ an $F'$-algebra through an embedding $\iota$ of $F'$. Let $\nabla$ be a canonical Taylor family, let
$n,N\ge0$, let $a$ be a family of nonnegative integers indexed by the letters, and let $P$ be a polynomial in the
letter variables with coefficients in $F'$ that satisfies conditions (a) to (c) of
Definition~\ref{cv-uniformization-theta-rule-values}(6) for $a$. Let $F_1,\dots,F_n$ be sections of
$\mathcal L_y^{N^2}$ with theta readings $\vartheta_1,\dots,\vartheta_n$. Then
\[
 \sum_m\iota(P_m)\prod_{b=1}^n\nabla_y(F_b;m_b)=\Big(P^\iota(\partial)\prod_{b=1}^n\vartheta_b(z_b)\Big)\Big|_{z=0},
\]
with the notation of Definition~\ref{cv-uniformization-theta-rule-values}(3), where $z=(z_1,\dots,z_n)\in(\mathbf C^g)^n$,
$P^\iota$ is $P$ with $\iota$ applied to the coefficients, and $P^\iota(\partial)$ replaces each letter variable
$X_{\ell_j,b}$ by $\partial/\partial z_{b,j}$.
\end{lemma}
\leanref{Development/AbelianVarietyTheory/Proofs/SiegelThetaTransformation/HarmonicConnectionValueEqPolynomialDerivativeStatement.lean}{GlobalTheory.T.harmonicConnectionValue_eq_polynomial_derivative}
\begin{proof}
Put $k=N^2$, and let $V_C$ be as in Lemma~\ref{cv-uniformization-theta-rule-unique-chart} for a chart datum $C$. As
$F'$ is a number field, some integer $N_0\ge1$ makes the coefficients of $N_0P$ algebraic integers; then $N_0P$ is
harmonic of row degrees $a$, with the monomials of $P$. As $\nabla$ is canonical, there is a chart datum $C_1$ at $y$
for which every section has a germ, with $\nabla_y(F;\alpha)=V_{C_1}(f;\alpha)$ for every germ $f$ of $F$; choose germs
$f_b$ of the $F_b$. Let $C$ be the chart datum of Lemma~\ref{cv-uniformization-theta-rule-fibre-taylor-chart}, for which
the Taylor series $\hat\vartheta_b$ of $\vartheta_b$ at $0$ is a germ of $F_b$. Lemma~\ref{cv-uniformization-theta-rule-unique-chart}
for $C_1$, $C$ and $N_0P$ gives, after cancelling $N_0$,
\[
 \sum_m\iota(P_m)\prod_b\nabla_y(F_b;m_b)=\sum_m\iota(P_m)\prod_bV_C(\hat\vartheta_b;m_b).
\]

The Taylor values of $C$. The functions $\vartheta_b$ are holomorphic: near each point $z$ with
$\operatorname{ex}(z)\in U_\beta$, $\vartheta_b$ is the value of $F_b$ at $\operatorname{ex}(z)$ in the trivialization over
$U_\beta$, a holomorphic function of $z$, divided by the nowhere vanishing holomorphic function
$\operatorname{fr}_\beta^k$. For the group law $t+t'$, $\partial F_j(t,t')/\partial t'_i=\delta_{ij}$, so the derivation
$D_{\ell_j}$ of $C$ is $\partial/\partial t_j$; and $h=0$. So $V_C(\hat\vartheta_b;\alpha)$ is the constant term of
$\prod_j(\partial/\partial t_j)^{\alpha_{\ell_j}}\hat\vartheta_b$, which is the derivative
$\big(\prod_j(\partial/\partial z_j)^{\alpha_{\ell_j}}\vartheta_b\big)(0)$.

The product. For a monomial $\prod_{j,b}X_{\ell_j,b}^{m_{\ell_j,b}}$, the operator
$\prod_{j,b}(\partial/\partial z_{b,j})^{m_{\ell_j,b}}$ applied to $\prod_b\vartheta_b(z_b)$ gives
$\prod_b\big(\prod_j(\partial/\partial z_j)^{m_{\ell_j,b}}\vartheta_b\big)(z_b)$, as the $b$-th factor depends only on
$z_b$. At $z=0$ this is $\prod_bV_C(\hat\vartheta_b;m_b)$. Summing with the coefficients $\iota(P_m)$ gives the claim.
\end{proof}
% Lean: T.harmonicConnectionValue_eq_polynomial_derivative = Dict.nab_harm_val_eq_polyDeriv, with ChartIndepSt2,
% IsSympl, x in X.dom, SelfD X L, IsNab and P in HHa (conditions (a) to (c)); the derivative is I2.polyDeriv with the
% directions Pi.single b (Pi.single (e.symm l) 1). The proof: T.exists_analyticFrame_of_analyticFibre, then
% T.harmonicValue_canonicalConnection_eq_polyDeriv = AnVal.harm_nab_eq_polyDeriv: exists_natCast_smul_mem_harmA for
% N0, C1 from IsNab, C from T.exists_analytic_biextension_chart_complex_point (Lemma
% cv-uniformization-theta-rule-fibre-taylor-chart), the holomorphy from T.section_reading_existsUnique_isTheta,
% ChartIndepSt2 (Lemma cv-uniformization-theta-rule-unique-chart), PullbackSections.nabVal_additive for the Taylor
% values of C, and T.polyDeriv_prod_column_functions = AnVal.polyDeriv_prod_cols for the product. Size beyond the laws,
% chart independence, the connection and Lemmas cv-uniformization-theta-rule-fibre-reading-sign to -fibre-reading:
% about 12,200 lines, of which about 11,600 in Lemma cv-uniformization-theta-rule-fibre-taylor-chart.

\begin{lemma}\label{cv-uniformization-theta-rule-fibre}
In the setting of Lemma~\ref{cv-uniformization-theta-marking}, with the signature $(p,q)$ of the datum, let $\nabla$
be a canonical Taylor family, $n\ge0$ even, $c_1,\dots,c_n\in L$, $a$ a family of nonnegative integers indexed by the
letters and $P$ a polynomial harmonic of row degrees $a$ (Definition~\ref{cv-uniformization-theta-rule-values}(6)),
and put $w_0=(1+\sqrt{-d})/2$. Let $X_0$ be data (1) to (6) of flag data of rank $r=p+q$ over $E$
(Definition~\ref{is-onto-share-flag-data}), with set $\Psi$, a symplectic basis of its module and an indexing
$j\mapsto(\tau_j,i_j)$ (Definition~\ref{is-eigenvalue-conjugate-frame-fibres}(1) and~(2)), and put $g=|\Psi|r$. Let $y$
be a framed point of $X'$ over $\mathbf C$, an $F'$-algebra through an embedding $\iota$ of $F'$, let $j\mapsto\ell_j$
be a bijection from $\{1,\dots,g\}$ onto the letters, $x\in\mathcal D_{X_0}$, $\Lambda\subset E^r$ an
$\mathcal O_E$-submodule self-dual for $X_0$, $(\operatorname{ex},\operatorname{fr},\chi)$ an analytic fibre at $y$
for these data, $x$ and $\Lambda$, and $\vartheta$ a theta reading of the section of $y$ for it
(Definition~\ref{is-eigenvalue-conjugate-frame-fibres}(3) and~(4)). Let $\varepsilon$ be a bijection from the letter
variables onto the variables $Y_{\sigma,i,k}$ of $X_0$, $\sigma\in\Psi$, $1\le i\le r$, $1\le k\le n$, that sends
$X_{\ell_j,k}$ to $Y_{\tau_j,i_j,k}$ for all $j$ and $k$; for a polynomial $Q$ in these variables with complex
coefficients let $Q(\partial)$ replace each $Y_{\tau_j,i_j,k}$ by $\partial/\partial z_{k,j}$, for
$z=(z_1,\dots,z_n)\in(\mathbf C^g)^n$. Let $\varphi$ be a vacuum of $\mathcal O_E^n\otimes\Lambda$ for $w_0$
(Definition~\ref{is-eigenvalue-conjugate-frame-heisenberg}(4)) and $\kappa\in\mathbf C$ with
$\prod_{b=1}^n\vartheta(z_b)=\kappa\,\theta_\varphi(x;z)$ for every $z\in(\mathbf C^g)^n$. Let $m\colon L\to E^r$ be an
$\mathcal O_E$-linear map with $m(L)\subset\Lambda$, such that $a'\in M'L$ whenever $a'\in\bar{\mathfrak M}'L$ and
$m(a')\in M'\Lambda$, such that every $v\in\bar{\mathfrak M}'\Lambda$ is $m(a')+M'u$ with $a'\in\bar{\mathfrak M}'L$
and $u\in\Lambda$, and such that the level structure of $y$ sends $a'\in\bar{\mathfrak M}'L$ to
$\operatorname{ex}\big(p_x(M'^{-1}m(a'))\big)$. Let $Y\in M_{n\times r}(E)$ have rows $Y_b$ such that for every $b$
there are $a'\in\bar{\mathfrak M}'L$ with $a'-c_b\in\mathfrak M'L$ and $Y_b-M'^{-1}m(a')\in\Lambda$, and an integer
$e\ge0$ with $sY_b\in2M'\Lambda$ for every $s\in\mathfrak M'^e$. Then every theta value of $c_1,\dots,c_n$ and $P$ at
$y$ for $\nabla$ is
\[
 \kappa\,\Big(\big(\varepsilon(P^\iota)\big)(\partial)\,\theta_{\rho_{X_0}(Y)\varphi}(x;z)\Big)\Big|_{z=0},
\]
where $\varepsilon(P^\iota)$ is $P$ with $\iota$ applied to the coefficients and $\varepsilon$ to the variables. For
$P=1$ this is $\kappa\,\theta_{\rho_{X_0}(Y)\varphi}(x;0)$.
\end{lemma}
\leanref{Development/AbelianVarietyTheory/Proofs/SiegelThetaTransformation/ThetaValueStandardLatticeEqClassicalTheta.lean}{GlobalTheory.Dict.thVal_eq_thP}
\begin{proof}
For $1\le b\le n$ fix $a'_b\in\bar{\mathfrak M}'L$ as in the statement, and put
$t'_b=\operatorname{ex}(p_x(M'^{-1}Y_b))$. The map $v\mapsto\operatorname{ex}(p_x(v))$ on $E^r$ is additive, vanishes
on $\Lambda$ and satisfies $\operatorname{ex}(p_x(sv))=s\operatorname{ex}(p_x(v))$ for $s\in\mathcal O_E$, by the
properties of the periods and of $\operatorname{ex}$ (Definition~\ref{is-eigenvalue-conjugate-frame-fibres}(1)
and~(3)).

Step 1: the vectors $M'^{-1}Y_b$. Fix $b$. For $s\in\mathfrak M'$, $sY_b\in\Lambda$: as
$\mathfrak M'\bar{\mathfrak M}'=M'\mathcal O_E$, $sa'_b=M'a''$ with $a''\in L$, so $sM'^{-1}m(a'_b)=m(a'')\in\Lambda$, and
$sY_b=s\big(Y_b-M'^{-1}m(a'_b)\big)+m(a'')$. For $N\in\{M',2\}$ let $N'=2M'/N$, and let $J_N$ be the ideal of the
$s\in\mathcal O_E$ with $sN^{-1}Y_b\in\Lambda$. It contains $\mathfrak M'^e$: if $sY_b=2M'u$ with $u\in\Lambda$, then
$sN^{-1}Y_b=N'u$. It contains $N\mathfrak M'$, by the first remark. For $N=M'$, $\mathfrak M'^2\bar{\mathfrak M}'=
M'\mathfrak M'\subset J_{M'}$ and $\mathfrak M'^e+\bar{\mathfrak M}'=\mathcal O_E$, so
$\mathfrak M'^2=\mathfrak M'^2\mathfrak M'^e+\mathfrak M'^2\bar{\mathfrak M}'\subset J_{M'}$: $sM'^{-1}Y_b\in\Lambda$ for
$s\in\mathfrak M'^2$. For $N=2$, $2\mathfrak M'\subset J_2$ and $\mathfrak M'^e+2\mathcal O_E=\mathcal O_E$, as $M'$
is odd and lies in $\mathfrak M'$; so $\mathfrak M'=\mathfrak M'\mathfrak M'^e+2\mathfrak M'\subset J_2$, and in
particular $u_b=\frac12M'Y_b\in\Lambda$. Then $M'^2\cdot M'^{-1}Y_b=2u_b$.

Step 2: the roots. By the hypothesis on the level structure, and as $\eta$ vanishes on $\mathfrak M'L$ and
$a'_b-c_b\in\mathfrak M'L$,
\[
 M't'_b=\operatorname{ex}(p_x(Y_b))=\operatorname{ex}\big(p_x(Y_b-M'^{-1}m(a'_b))\big)
 +\operatorname{ex}\big(p_x(M'^{-1}m(a'_b))\big)=\eta_y(a'_b)=\eta_y(c_b).
\]
By Step 1, $M'^2t'_b=\operatorname{ex}(p_x(2u_b))=0$, and $st'_b=\operatorname{ex}(p_x(sM'^{-1}Y_b))=0$ for
$s\in\mathfrak M'^2$.

Step 3: the theta value. As $M'$ is odd, Lemma~\ref{cv-uniformization-theta-rule-unique-translates}(1) gives an
$M'$-translate $F_b$ of $\vartheta_y$ over $\eta_y(c_b)$ with root $t'_b$, and by Step 2 it is a level translate.
The identity of $\mathbf C$ is faithfully flat, and $y$ is an image of itself under it
(Lemma~\ref{cv-uniformization-theta-rule-unique-base-change}(2)), so $\nabla_y(P;F_1,\dots,F_n)$ is a theta value of
$c_1,\dots,c_n$ and $P$ at $y$. By Lemma~\ref{cv-uniformization-theta-rule-unique} it is the only one.

Step 4: the readings. By Lemma~\ref{cv-uniformization-theta-rule-fibre-reading}, with $N=M'$, the vector $M'^{-1}Y_b$
in place of $a$ (Step 1) and $t=\eta_y(c_b)$, the section $F_b$ has the theta reading
\[
 G_b(z)=\exp\big(-\pi H_x\big(M'z+\tfrac12p_x(Y_b),p_x(Y_b)\big)\big)\,\vartheta\big(M'z+p_x(Y_b)\big).
\]
The vacuum $\varphi$ is Schwartz--Bruhat, so Lemma~\ref{cv-uniformization-theta-rule-product}, for
$\vartheta_b=\vartheta$ and at the point $M'z$, gives $\prod_bG_b(z_b)=\kappa\,\theta_{\rho_{X_0}(Y)\varphi}(x;M'z)$ for
every $z\in(\mathbf C^g)^n$.

Step 5: the Taylor values. By Lemma~\ref{cv-uniformization-theta-rule-fibre-taylor}, with $N=M'$,
\[
 \sum_m\iota(P_m)\prod_b\nabla_y(F_b;m_b)=\Big(P^\iota(\partial)\,\kappa\,\theta_{\rho_{X_0}(Y)\varphi}(x;M'z)\Big)\Big|_{z=0},
\]
where $P^\iota(\partial)$ replaces $X_{\ell_j,b}$ by $\partial/\partial z_{b,j}$. The function
$\theta_{\rho_{X_0}(Y)\varphi}(x;\cdot)$ is holomorphic on $(\mathbf C^g)^n$: $\rho_{X_0}(Y)\varphi$ is Schwartz--Bruhat
(Lemma~\ref{cv-uniformization-theta-rule-translate}), and its theta series converges absolutely and uniformly on
compact sets \cite[Chapter~II, \S1]{Mum83}. By (c) of Definition~\ref{cv-uniformization-theta-rule-values}(6), every
monomial of $P$ has degree $\deg P$, so by the chain rule the right side is
$M'^{\deg P}\kappa\,\big(P^\iota(\partial)\theta_{\rho_{X_0}(Y)\varphi}(x;z)\big)|_{z=0}$. Multiplying by
$M'^{-\deg P}$, the theta value of Step 3 is $\kappa\,\big(P^\iota(\partial)\theta_{\rho_{X_0}(Y)\varphi}(x;z)\big)|_{z=0}$.
As $\varepsilon$ sends $X_{\ell_j,k}$ to $Y_{\tau_j,i_j,k}$, which $Q(\partial)$ replaces by $\partial/\partial z_{k,j}$,
$P^\iota(\partial)$ is $\big(\varepsilon(P^\iota)\big)(\partial)$, which gives the claim. For $P=1$ the operator is
the identity.
\end{proof}
% Lean: thVal_eq_thP (40 lines), with ChartIndepSt2 and firstFiveLaws.1: hroot from root_laws and smul_mem_of_rep
% (Step 1: the key lemma for N N' = 2M', an ideal I with Mb + I = 1 and K I <= Mb (N), used for (M', 2, conj Mb,
% Mb^2) and (2, M', (2), Mb), with level_mem); hpow from ex_pv_pow, pv_add, ex_pv_mem, A.Marked (hA) and levPt_congr
% (Step 2); hF from exists_isVtOver_and_existsUnique_isUUOver with the condition t' = ex(pv(M'^-1 Y_b)) on the root;
% hval with AlgHom.id, FaithfullyFlat.of_bijective, FramedBaseChange.refl and ex_pv_act; huniq from
% thValTot_of_chartIndep (Step 3); hread from T.sections_read_as_translates_at_scaled_point (Step 4); hnv from
% T.harmonicConnectionValue_eq_polynomial_derivative and hfock from the last clause of fock_laws (Lemma
% cv-uniformization-theta-rule-product), with heisS_laws (IsSB) and SW.theta_holomorphic_weight (analyticity of thA,
% which the Lean proves; the text cites Mumford); then polyDeriv_scale with card_eq_totalDegree,
% I2.polyDeriv_smul_fun, the definition of thP and polyDeriv_rename with hdir (Step 5). The Lean's variables of X_0
% are Harm.Var X.Psi q0 X.p0 (Fin n), with hr : p0 + q0 = r and the row rowIdx (inl i) = p0 + i, rowIdx (inr i) =
% p0 - 1 - i; the letter map eL with letterDir (eL l) the unit vector at (b, e^-1 l) is the hypothesis on epsilon,
% and the conclusion is v0 = kappa thP X hr bas idx (heisG X bas Y v) (rename eL (map P)) x. IsMark l L mk is the
% three conditions on m (mem, inj, surj), and A.Marked mk the last one. Size beyond the laws, chart independence, the
% connection, the theta values and Lemmas cv-uniformization-theta-rule-translate, -product, -fibre-reading and
% -fibre-taylor, counted as in Lemma cv-uniformization-theta-rule-connection-charts: about 6,000 lines, mostly reached
% through SW.theta_holomorphic_weight, which the Lean proves together with the transformation laws of the theta
% series (Siegel theta transformation, theta characteristics); the text uses only the holomorphy, cited to Mumford.

\begin{lemma}\label{cv-uniformization-theta-rule-one}
In the setting of Lemma~\ref{cv-uniformization-theta-marking}, let $n\ge0$. The polynomial $1$ in the letter variables
is harmonic of row degrees $0$ (Definition~\ref{cv-uniformization-theta-rule-values}(6)).
\end{lemma}
\leanref{Development/AbelianVarietyTheory/Proofs/DivisorProofs/PowGeneratorEqConstMulVacuumTheta.lean}{GlobalTheory.AdTh.GenConst.one_mem_harmA}
\begin{proof}
Its derivatives vanish, the substitutions of (b) fix it, its only monomial has degree $0$ in every variable, and $1$
is an algebraic integer.
\end{proof}
% Lean: AdTh.GenConst.one_mem_harmA (14 lines): harmK and lap of a constant, map_one for RowInv, the support {0} for
% RowDeg, and one_mem for the coefficients.

\begin{lemma}\label{cv-uniformization-theta-rule}
In the setting of Lemma~\ref{cv-uniformization-theta-marking}, let $n\ge0$ be even and $c_1,\dots,c_n\in L$.
There is a rational section $f$ of $X'$ of the dominant weight $(0,n)$ with the following property. Let
$\sigma$, $Z$, $x$, $H_x$, $\vartheta_x$ and $m_Z$ be as in Lemma~\ref{cv-uniformization-theta-marking}, and for
$1\le b\le n$ let $a_b\in\bar{\mathfrak M}'L$ with $a_b-c_b\in\mathfrak M'L$, and let $Y_b\in E^r$ with
$Y_b-M'^{-1}m_Z(a_b)\in\Lambda_Z$ and $sY_b\in2M'\Lambda_Z$ for every $s$ in some power of $\mathfrak M'$. Then
the reading of $f$ at $x$ is
\[
 f_Z(x)=\prod_{b=1}^n\exp\big(-\pi H_x\big(\tfrac12P_x(Y_b),P_x(Y_b)\big)\big)\,\vartheta_x\big(P_x(Y_b)\big).
\]
\end{lemma}
\leanref{Formalism/LocalToGlobal/ProjectiveLattices/ThProps.lean}{Formalism.L2.Pj.existsUnique_thRule}
\leanref{Development/AbelianVarietyTheory/Proofs/CuspFormEvaluations/ThetaRuleMemSec.lean}{GlobalTheory.thetaRule_mem_sec}
\leanref{Development/AbelianVarietyTheory/Proofs/SiegelThetaTransformation/ThetaValueStandardLatticeEqClassicalTheta.lean}{GlobalTheory.Dict.thVal_eq_thP}
\leanref{Development/AbelianVarietyTheory/Proofs/SiegelThetaTransformation/FockHxDifferentiable.lean}{GlobalTheory.Dict.fock_laws}
\begin{proof}
The section. Let $\nabla$ be a canonical Taylor family on $X'$ (Lemma~\ref{cv-uniformization-theta-rule-connection}).
The polynomial $1$ is harmonic of row degrees $0$ (Lemma~\ref{cv-uniformization-theta-rule-one}). By
Lemma~\ref{cv-uniformization-theta-rule-unique}, every framed point $y$ of $X'$ over a test ring has exactly one
theta value of $c_1,\dots,c_n$ and $1$ for $\nabla$; let $f(y)$ be this value. Then $f$ is a theta rule of
$c_1,\dots,c_n$ and $1$ for $\nabla$, and the only one (Definition~\ref{cv-uniformization-theta-rule-values}(4)). By
Lemma~\ref{cv-uniformization-theta-rule-weight}, for the row degrees $0$, $f$ is a rational section of $X'$ of exponent
$e(0,n)$, which is the exponent of the dominant weight $(0,n)$ (Definition~\ref{cv-volume-key-theta-point}), as $n$ is
even.

The reading. Let $\sigma$, $Z$, $x$, the $a_b$ and the $Y_b$ be as in the statement, let $y$ be the framed point at
$x$, over $\mathbf C$ through $\sigma$, and let $Y\in M_{n\times r}(E)$ have the rows $Y_b$. By
Lemma~\ref{cv-uniformization-theta-rule-complex}, $f_Z(x)$ is a theta value of $c_1,\dots,c_n$ and $1$ at $y$ for
$\nabla$. Let $X_Z$, its module $\Lambda_Z$ and the bijection $j\mapsto\ell_j$ be as in the proof of
Lemma~\ref{cv-uniformization-theta-half-sign}, and take there the letter dictionary of
Lemma~\ref{is-eigenvalue-conjugate-frame-letters} for $n$ and $a=0$; its letter map sends each letter variable
$X_{\ell_j,k}$ to the variable $Y_{\tau_j,i_j,k}$ of $X_Z$, for the indexing $j\mapsto(\tau_j,i_j)$ of the dictionary
(Definition~\ref{is-eigenvalue-conjugate-frame-letter-dictionary}). The argument of that proof shows that
$\Lambda_Z$ is self-dual for $X_Z$, that $x\in\mathcal D_{X_Z}$, that $P_x=p_x$, that $H_x$ is the form $H_x$ of
Definition~\ref{is-eigenvalue-conjugate-frame-fibres}(2), and that $(p_x,F^x,\chi)$ is an analytic fibre at $y$ for
these data, $x$ and $\Lambda_Z$. By Lemma~\ref{cv-uniformization-theta-marking}(1), $\vartheta_x$ is a theta reading of
the section of $y$ for it. By Lemma~\ref{is-eigenvalue-conjugate-frame-level-one}, $\vartheta_x$ is a level-one theta
function for $w_0$, $x$ and $\Lambda_Z$. As $w_0$ is an algebraic integer with $w_0+\bar w_0=1$,
Lemma~\ref{is-eigenvalue-conjugate-frame-vacua}(1), for $\Lambda_S=\mathcal O_E^n$, gives a vacuum $\varphi$ of
$\mathcal O_E^n\otimes\Lambda_Z$ for $w_0$, and Lemma~\ref{is-eigenvalue-conjugate-frame-theta-product} gives
$\kappa\in\mathbf C$ with $\prod_b\vartheta_x(z_b)=\kappa\,\theta_\varphi(x;z)$ for every $z\in(\mathbf C^g)^n$.

The marking $m_Z$ maps $L$ into $\Lambda_Z$, every $v\in\bar{\mathfrak M}'\Lambda_Z$ is $m_Z(a)+M'u$ with
$a\in\bar{\mathfrak M}'L$ and $u\in\Lambda_Z$, and the level structure of $y$ sends $a\in\bar{\mathfrak M}'L$ to
$p_x\big(P_x(M'^{-1}m_Z(a))\big)$ (Lemma~\ref{cv-uniformization-theta-marking}(2)); and an $a\in\bar{\mathfrak M}'L$
with $m_Z(a)\in M'\Lambda_Z$ lies in $M'L$ (Lemma~\ref{is-eigenvalue-conjugate-frame-component-lattice}(2), as
$m_Z=\varphi_Z\circ t_Z$). So Lemma~\ref{cv-uniformization-theta-rule-fibre} applies, with $X_0=X_Z$, $\Lambda=\Lambda_Z$,
$m=m_Z$, $\operatorname{ex}=p_x$, $P=1$ and the letter map of the dictionary as $\varepsilon$, and gives
$f_Z(x)=\kappa\,\theta_{\rho_{X_Z}(Y)\varphi}(x;0)$. By
Lemma~\ref{cv-uniformization-theta-rule-product}, for $\vartheta_b=\vartheta_x$ and $z=0$, and as $p_x(Y_b)=P_x(Y_b)$,
\[
 \kappa\,\theta_{\rho_{X_Z}(Y)\varphi}(x;0)=\prod_{b=1}^n\exp\big(-\pi H_x\big(\tfrac12P_x(Y_b),P_x(Y_b)\big)\big)\,
 \vartheta_x\big(P_x(Y_b)\big),
\]
which is the formula of the statement.
\end{proof}
% Lean: the theta rule part of DictLevel.exists_sec_rdU_ne_zero. The section: CanonConn.exists_isNab_uncond (nab),
% allSixLaws.1, chartIndependence_unconditional, one_mem_harmA (Lemma cv-uniformization-theta-rule-one),
% existsUnique_thRule (11 lines: the values of thValTot_of_chartIndep at the test rings) and thetaRule_mem_sec for
% a = 0, whose sec Xo c 0 n is the exponent ex Dw 0 n. The reading: ThRead.thValAt_evR at unC U k x
% (rdU = evR by dif_pos); nonempty_uw and uw_lattice (Lemma is-eigenvalue-conjugate-frame-component-lattice: SelfD
% and IsMark of W.mark k), exists_dic with this n and a = 0 (Lemma is-eigenvalue-conjugate-frame-letters), uw_anFib
% (Lemma cv-uniformization-theta-marking(1), with the analytic fibre as in the proof of Lemma
% cv-uniformization-theta-half-sign), isTheta1_of_reads_generator (Lemma is-eigenvalue-conjugate-frame-level-one),
% SW.TensVac.vacua with stdL_mem_Vset (Lemma is-eigenvalue-conjugate-frame-vacua(1)), VacProd.prod_eq_mul_thA
% (Lemma is-eigenvalue-conjugate-frame-theta-product), thVal_eq_thP with q0 = Dw.p, Dc.hr, the letter map Dc.eL with
% Dc.dir_eq, P = 1 and polyDeriv_one (Lemma cv-uniformization-theta-rule-fibre) and the last clause of fock_laws (Lemma
% cv-uniformization-theta-rule-product). Sizes, each beyond the ones before it, counted through the theorems a
% declaration reaches in the files it imports, resolved through the namespaces and open lines of its file: the laws
% of abelian schemes about 57,100 lines (shared with the chapter); chart independence 19,100 (Lemma
% cv-uniformization-theta-rule-unique-chart); the connection 10,200 (Lemma -connection-charts 10,100); the theta
% values 5,700 (the base change of framed points 1,300, the translates 2,300, the base change of Taylor sums 1,600,
% the rest 450); the weight 1,800 (the triangular change of Taylor sums 1,200); the reading as a value 5,900 (the
% descent to a finitely generated subalgebra); the reading of translates 5,200 and the Taylor sums at a complex point
% 12,100 (Lemmas -fibre-reading and -fibre-taylor); the analytic fibre and the lattice 3,700 (Lemma
% cv-uniformization-theta-marking), the level-one law 700, the dictionary 1,500, the vacua 2,000, the theta product
% 5,900, and the rest of the theta value at an analytic fibre 2,200 (Lemma cv-uniformization-theta-rule-fibre); Lemmas
% cv-uniformization-theta-rule-translate and -product add none, and Lemma -one about 100. Without the laws, about
% 76,000 lines in all. The open parts are Lemmas cv-uniformization-theta-rule-unique-chart (chart independence,
% 19,100), -connection-charts-frames (1,700), -connection-charts-identities (1,030), -connection-charts-gauge (2,960),
% -connection-charts-germs (shared with -frames), -unique-translates-rosati (510), -complex-charts (4,220),
% -complex-universal (790), -complex-glue (640), -fibre-reading-sign (1,060, and 660 for the analytic point),
% -fibre-reading-pullback (290), -fibre-reading-symmetric (1,170) and -fibre-taylor-chart (11,600), and the inputs
% from Chapter 8.

\begin{lemma}\label{cv-uniformization-theta-section}
In the setting of Lemma~\ref{cv-uniformization-theta-marking}, let $x$ be a frame of $Z$ with
$|\delta_{\tau,i}(x)|=1$ for all $\tau$ and $i$, and let $n\ge0$ be even. Then there is a rational section $f$ of
$X'$ of the dominant weight $(0,n)$ whose reading $f_Z(x)$ at $x$ is not $0$.
\end{lemma}
\leanref{Development/AbelianVarietyTheory/Proofs/TwistedOperators/ExistsSecRdUNeZero.lean}{GlobalTheory.DictLevel.exists_sec_rdU_ne_zero}
\leanref{Development/AbelianVarietyTheory/Proofs/DivisorProofs/PowGeneratorEqConstMulVacuumTheta.lean}{GlobalTheory.AdTh.GenConst.exists_root_ne_zero}
\begin{proof}
Let $H_x$, $\chi_x$, $\vartheta_x$ and $m_Z$ be as in Lemma~\ref{cv-uniformization-theta-marking}.

The point. By Lemma~\ref{cv-uniformization-theta-half}(2) there is $u\in\bar{\mathfrak M}'\Lambda_Z$ with
$\vartheta_x(P_x(M'^{-1}u))\ne0$, and by Lemma~\ref{cv-uniformization-theta-marking}(2), $u=m_Z(a)+M't$ with
$a\in\bar{\mathfrak M}'L$ and $t\in\Lambda_Z$. Let $q$ be the class of $a$ in $L/\mathfrak M'L$ and $c\in L$ a
representative of $q$.

The section. For $1\le b\le n$, Lemma~\ref{cv-uniformization-theta-marking}(3) gives $a_b\in\bar{\mathfrak M}'L$ in
the class $q$ and $Y_b\in E^r$ with $Y_b-M'^{-1}m_Z(a_b)\in\Lambda_Z$ and $sY_b\in2M'\Lambda_Z$ for every $s$ in some
power of $\mathfrak M'$. Let $f$ be the section of Lemma~\ref{cv-uniformization-theta-rule} for the points
$c_1=\dots=c_n=c$, of the dominant weight $(0,n)$. As $a_b$ and $c$ lie in the class $q$, $a_b-c\in\mathfrak M'L$,
and that lemma gives
\[
 f_Z(x)=\prod_{b=1}^n\exp\big(-\pi H_x\big(\tfrac12P_x(Y_b),P_x(Y_b)\big)\big)\,\vartheta_x\big(P_x(Y_b)\big).
\]
It remains to show that $\vartheta_x(P_x(Y_b))\ne0$ for every $b$.

The factors. Fix $b$. The element $a_b-a$ lies in $\bar{\mathfrak M}'L$, as $a_b$ and $a$ do, and in
$\mathfrak M'L$, as both lie in the class $q$. As $\mathfrak M'+\bar{\mathfrak M}'=\mathcal O_E$, write $1=s+s'$
with $s\in\mathfrak M'$ and $s'\in\bar{\mathfrak M}'$. Then $s(a_b-a)$ and $s'(a_b-a)$ lie in
$\mathfrak M'\bar{\mathfrak M}'L=M'L$, so $a_b-a=M'c'$ with $c'\in L$, and $m_Z(a_b)=m_Z(a)+M'm_Z(c')$. Hence
\[
 \ell=Y_b-M'^{-1}u=\big(Y_b-M'^{-1}m_Z(a_b)\big)+m_Z(c')-t
\]
lies in $\Lambda_Z$. As $P_x$ is additive, $P_x(Y_b)=P_x(M'^{-1}u)+P_x(\ell)$, and
Lemma~\ref{cv-uniformization-theta-half}(1) gives
\[
 \vartheta_x\big(P_x(Y_b)\big)=\chi_x(P_x\ell)\exp\big(\pi H_x(P_x(M'^{-1}u),P_x\ell)+\tfrac\pi2H_x(P_x\ell,P_x\ell)\big)
 \,\vartheta_x\big(P_x(M'^{-1}u)\big),
\]
which is not $0$.
\end{proof}
% Lean: exists_sec_rdU_ne_zero: the data at x are nonempty_uw, uw_lattice (IsMark of W.mark k, SelfD of the
% lattice), exists_dic and uw_anFib (Lemma cv-uniformization-theta-marking); isTheta1_of_reads_generator gives hT;
% the vacuum and kappa (SW.TensVac.vacua, VacProd.prod_eq_mul_thA) and the reading enter Lemma
% cv-uniformization-theta-rule. The point and the factors are AdTh.GenConst.exists_root_ne_zero: u from
% T.exists_levelOne_theta_ne_zero_on_ideal_half_lattice, a and t from IsMark.surj, q the constant family of the
% class of a, the Y_b from DualMark.dual_marking (IsRoots), 1 = s + s' from LevelFibre.level_split, the two
% multiples of M' from MarkLift.modN_of_mem_half and MarkLift.modN_of_mem_bar, and the factor relation from the
% clause of IsTheta1 at the vector l (fock and pv additive). The section is the theta rule of the points
% levRep (q b) = c and P = 1 (existsUnique_thRule, thetaRule_mem_sec); its reading is thVal_eq_thP with
% polyDeriv_one, and fock_laws turns it into the product, whose factors are nonzero (h3).

\subsection{Uniformization of deep covers}\label{subsection-cv-deep-cover-uniformization}

\begin{lemma}\label{cv-uniformization-fibre-count-tangent}
Let $k$ be a field of characteristic $0$, let $A$ and $B$ be abelian varieties over $k$, and let
$f,f'\colon A\to B$ be homomorphisms. Let $k[\varepsilon]$ be the ring of dual numbers, and suppose that
$f\circ v=f'\circ v$ for every tangent vector $v$ of $A$ at the origin, that is, every point $v$ of $A$ with values
in $k[\varepsilon]$ that reduces to the origin. Then $f=f'$.
\end{lemma}
\leanref{Development/AbelianVarietyTheory/Proofs/DivisorProofs/AbelianSchemeHomEqOfTangentEqStatement.lean}{GlobalTheory.T.abelianScheme_hom_eq_of_tangent_eq}
\begin{proof}
Put $g=f-f'$, a homomorphism with $g\circ v=0$ for every tangent vector $v$ of $A$ at the origin. Let $R_A$ and
$R_B$ be the local rings of $A$ and $B$ at the origin, with maximal ideals $\mathfrak m_A$ and $\mathfrak m_B$, and
write $x(0)\in k$ for the residue of $x$. As $g$ maps the origin to the origin, it induces a local homomorphism of
$k$-algebras $\psi\colon R_B\to R_A$. Multiplication by $2$ induces local endomorphisms $\phi_A$ of $R_A$ and
$\phi_B$ of $R_B$, and $\psi\circ\phi_B=\phi_A\circ\psi$, as $g\circ[2]=[2]\circ g$. For a linear form $\lambda$ on
$\mathfrak m_A/\mathfrak m_A^2$, the map $v_\lambda\colon R_A\to k[\varepsilon]$,
$x\mapsto x(0)+\lambda(x-x(0))\varepsilon$, is a homomorphism of $k$-algebras, as $\lambda$ vanishes on
$\mathfrak m_A^2$; so it is a tangent vector of $A$ at the origin. An element of $\mathfrak m_A$ killed by every
$\lambda$ lies in $\mathfrak m_A^2$. The same holds for $B$.

Step 1: $\phi_A(x)-2x\in\mathfrak m_A^2$ for $x\in\mathfrak m_A$, and likewise for $B$. The group law of $A$ induces
addition on tangent vectors at the origin, so $[2]\circ v_\lambda=v_{2\lambda}$, that is,
$v_\lambda\circ\phi_A=v_{2\lambda}$. For $x\in\mathfrak m_A$ this says $\lambda(\phi_A(x))=2\lambda(x)$, for every
$\lambda$.

Step 2: $\psi(\mathfrak m_B)\subset\mathfrak m_A^2$. For every $\lambda$, $g\circ v_\lambda$ is the origin of $B$,
that is, $v_\lambda(\psi(x))=x(0)$ for $x\in R_B$. For $x\in\mathfrak m_B$ this gives $\lambda(\psi(x))=0$.

Step 3: for $n\ge1$ and $y\in\mathfrak m_A^n$, $\phi_A(y)-2^ny\in\mathfrak m_A^{n+1}$. This is Step 1 for $n=1$, and
it passes from $n$ to $n+1$ by the identity
$\phi_A(ab)-2^{n+1}ab=(\phi_A(a)-2^na)\phi_A(b)+2^na(\phi_A(b)-2b)$ for $a\in\mathfrak m_A^n$ and
$b\in\mathfrak m_A$, as $\phi_A(\mathfrak m_A)\subset\mathfrak m_A$.

Step 4: $\psi(\mathfrak m_B)\subset\mathfrak m_A^n$ for every $n\ge2$. The case $n=2$ is Step 2. Suppose that
$\psi(\mathfrak m_B)\subset\mathfrak m_A^n$, and let $x\in\mathfrak m_B$. By Step 1 for $B$,
$\phi_B(x)-2x\in\mathfrak m_B^2$, so $\psi(\phi_B(x)-2x)\in\mathfrak m_A^{2n}\subset\mathfrak m_A^{n+1}$, and by Step 3,
$\phi_A(\psi(x))-2^n\psi(x)\in\mathfrak m_A^{n+1}$. As $\psi\circ\phi_B=\phi_A\circ\psi$, the difference is
$(2^n-2)\psi(x)\in\mathfrak m_A^{n+1}$. As $n\ge2$ and $k$ has characteristic $0$, $2^n-2$ is invertible in $k$, so
$\psi(x)\in\mathfrak m_A^{n+1}$.

Step 5: the conclusion. The local ring $R_A$ is Noetherian, so the intersection of the powers of $\mathfrak m_A$ is
$0$ \cite[Corollary~10.20]{AM69}, and $\psi(\mathfrak m_B)=0$ by Step 4. Hence $\psi(x)=x(0)$ for every $x\in R_B$:
the composite of $\operatorname{Spec}R_A\to A$ with $g$ is the constant morphism with value the origin. The
morphism $\operatorname{Spec}R_A\to A$ is dominant, as $A$ is integral. Two morphisms from the reduced scheme $A$ to
the separated scheme $B$ that agree after composing with a dominant morphism are equal, so $g$ is the constant
morphism with value the origin. So $g=0$ and $f=f'$.
\end{proof}
\leanref{Development/AbelianVarietyTheory/Proofs/DivisorProofs/AbelianSchemeHomEqOfTangentEq.lean}{GlobalTheory.LieFaithful.hom_eq_of_tangent_eq}
\leanref{Development/AbelianVarietyTheory/Proofs/DivisorProofs/MapEqZeroOfCommuteDoubling.lean}{GlobalTheory.LieFaithful.map_eq_zero_of_commute_doubling}
% Lean: hom_eq_of_tangent_eq: g = f * f'^{-1}; psi, phi_A and phi_B are exists_germHom; v_lambda is tanPt (tanHom of
% germEv and lambda, with smT_zero_tanPt); Step 1 is germ_doubling (smT_add, smT_unit, mem_sq_of_tanHom_eq_zero),
% Step 2 is germ_sq_of_tan, Steps 3 and 4 are map_sub_two_pow_mul_mem and map_eq_zero_of_commute_doubling (with
% Ideal.iInf_pow_eq_bot_of_isLocalRing for Krull), and Step 5 is germMap_comp_eq_one with
% ext_of_isDominant_of_isSeparated (A integral by isIntegral_left_of_isAbelianScheme_field).

\begin{lemma}\label{cv-uniformization-fibre-count}
In Situation~\ref{cv-cusp-setting}, suppose that the level of $X$ is a level $(M,\mathfrak M)$ with $M$ prime to
$15$, and let $Y$ be a deep cover of $X$ of exponent $e$, with deep level $(15^eM,\mathfrak MI)$ and coefficient
field $F_Y$ (Definition~\ref{cv-deep-cover}). Let $L$, $\Gamma$, $\Gamma_1$, $p(y_2)$ and $(y_2)_\gamma$ be as in
Lemma~\ref{cv-realization-pullback}.
\begin{enumerate}
\item Let $k$ be an algebraically closed field that is an algebra over $F_Y$, hence of characteristic $0$. If $y_2$
  is a framed point of $Y$ over $k$ and $\gamma,\gamma'\in\Gamma$ satisfy $(y_2)_\gamma=(y_2)_{\gamma'}$, then
  $\gamma=\gamma'$.
\item Consider framed points over $\mathbf C$, an algebra over $F_Y$ through an embedding $\sigma'$ and over $F_c$
  through the restriction of $\sigma'$. For every framed point $y$ of $X$ over $\mathbf C$, the number of framed
  points $y_2$ of $Y$ over $\mathbf C$ with $p(y_2)=y$ is $|\Gamma_1|$. In particular it is the same for all framed
  points of $X$ over $\mathbf C$.
\end{enumerate}
\end{lemma}
\leanref{Development/AbelianVarietyTheory/Proofs/DivisorProofs/CardFramedLiftsLevelCoverConstStatement.lean}{GlobalTheory.T.card_framed_lifts_levelCover_const}
\leanref{Development/AbelianVarietyTheory/Proofs/DivisorProofs/CardFramedLiftsLevelCoverConst.lean}{GlobalTheory.LevDegree.card_lifts_const}
\leanref{Development/AbelianVarietyTheory/Proofs/DivisorProofs/LevelStructureChangeFramedGeometricPointUnique.lean}{GlobalTheory.LevelFree.move_unique}
\begin{proof}
Framed points over a field that are isomorphic are equal: they have the same point, as $Y$ and $X$ represent their
moduli problems (Theorem~\ref{cv-deep-pel-facts}(1)), and their bases and generators correspond under an
automorphism of the abelian variety of that point that preserves the action of $\mathcal O_E$, the polarization and
the level structure; such an automorphism has finite order and is congruent to $1$ modulo the level, which is at
least $3$, so it is trivial \cite{Min87}. So $(y_2)_\gamma$ is a well defined framed point, and
Lemma~\ref{cv-realization-pullback} holds with equalities in place of isomorphisms.

(1) Let $(\rho,\nu)$ and $(\rho',\nu')$ be level automorphisms of $Y$ with matrices $\tilde\gamma$ and $\tilde\gamma'$
that induce $\gamma$ and $\gamma'$ (Definition~\ref{cv-realization-level-move}), and let $y'_2$ be the common move of
$y_2$, with point $\rho\circ x=\rho'\circ x$, where $x$ is the point of $y_2$. The isomorphisms
$\nu_x,\nu'_x\colon A_x\to A_{\rho(x)}$ both carry the bases of $y_2$, which form a basis of the Lie algebra of
$A_x$, to the bases of $y'_2$; so they have the same differential at the origin, and they agree on the tangent
vectors of $A_x$ at the origin. As $k$ has characteristic $0$, Lemma~\ref{cv-uniformization-fibre-count-tangent}
gives $\nu_x=\nu'_x$. For $a\in L$,
\[
 \nu_x\big(\eta_x(\tilde\gamma a)\big)=\eta_{\rho(x)}(a)=\nu'_x\big(\eta_x(\tilde\gamma'a)\big),
\]
so $\eta_x(\tilde\gamma a-\tilde\gamma'a)=0$, and $\tilde\gamma a-\tilde\gamma'a\in\mathfrak MIL$, as $\eta_x$ has the
kernel $\mathfrak MIL$. So $\gamma$ and $\gamma'$ agree on $L/\mathfrak MIL$.

(2) Let $y$ be a framed point of $X$ over $\mathbf C$. By Definition~\ref{cv-deep-cover}(4) there is a framed point
$y_0$ of $Y$ over $\mathbf C$ with $p(y_0)=y$. Let $\gamma\in\Gamma_1$ and let $\tilde\gamma$ be an endomorphism of $L$
that induces $\gamma$. As $\mathfrak MIL\subset\mathfrak ML$ and $\gamma$ induces the identity on $L/\mathfrak ML$,
$(\tilde\gamma-1)L\subset\mathfrak ML$, and the identity of $X$, with the identity of its universal abelian scheme,
is a level automorphism with matrix $1$, whose move of $p(y_0)$ is $p(y_0)$. So
Lemma~\ref{cv-realization-pullback}(2) gives $p((y_0)_\gamma)=p(y_0)=y$. Hence $\gamma\mapsto(y_0)_\gamma$ maps
$\Gamma_1$ to the set of the framed points $y_2$ of $Y$ with $p(y_2)=y$. It is injective by (1). It is onto: if
$p(y_2)=y=p(y_0)$, then $y_2=(y_0)_\gamma$ for some $\gamma\in\Gamma_1$ by Lemma~\ref{cv-realization-pullback}(3). So
there are $|\Gamma_1|$ such framed points.
\end{proof}
% Lean: card_lifts_const (T.card_framed_lifts_levelCover_const) counts, for every framed point y of X over CAt
% sigma_2, the framed points over y as the gamma in Gam Dw l_2 (automorphisms of La / Mb_2 La) with
% CongMod l.Mb (lift gamma) id, that is Gamma_1, with lift from LevPullModel.exists_isLift. The moves are total
% by LevelActionOfLifts.actLev_of_lifts_of_eq_one and unique by LevelSpanCompare.levR_unique (the text: Lemma
% cv-realization-pullback(1) and the first paragraph); y_0 is the lifting field of DeltaDat.IsLev (Definition
% cv-deep-cover(4)); the bijection uses T.levelCover_fibre_torsor_trivial_level_changes = LevOrbit.fibre_torsor:
% its backward direction (T.levMove_comp_r_eq_of_trivial_mod and frOver_of_square_of_frMove with the identity move
% of T.framedPointMove_calculus) is Lemma cv-realization-pullback(2) with gamma' = 1, its forward direction
% (exists_levMove_of_frOver) is Lemma cv-realization-pullback(3), and its uniqueness is
% T.levelStructureChange_framedGeometricPoint_unique = LevelFree.move_unique, which is (1), over any algebraically
% closed field k with an algebra structure over c.F: rho x = rho' x, the two nu agree on the tangent vectors through
% the frame (FrTan.frTanBasis_holds), hence on the fibre (LevelFree.sections_eq through
% T.abelianScheme_hom_eq_of_tangent_eq, over any field of characteristic 0: Lemma
% cv-uniformization-fibre-count-tangent), and eta_injective.

\begin{lemma}\label{cv-cover-uniformization}
In Situation~\ref{cv-cusp-setting}, suppose that the level of $X$ is a level $(M,\mathfrak M)$ with $M$
prime to $15$, and let $Y$ be a deep cover of $X$ (Definition~\ref{cv-deep-cover}), with deep level
$(M_Y,\mathfrak M_Y)$, coefficient field $F_Y$ and level map $\pi_Y\colon Y\to X$. Let $\Lambda$ be the
lattice of the PEL datum of $X$, let $\gamma$ be an $\mathcal O_E$-linear automorphism of
$\Lambda/\mathfrak M_Y\Lambda$, which acts on $Y$ by changing the level structure of each point by
$\gamma$, and let $\gamma^*$ be the pullback of rational sections of $Y$ along this action. Let $\sigma'$ be
a good embedding of $F_Y$ and $\sigma=\sigma'|F_c$; components, frames, arithmetic groups $\Gamma_Z$ and
readings are those of the uniformizations of $Y$ at $\sigma'$ and of $X$ at $\sigma$.
\begin{enumerate}
\item For every component $Z$ of $Y$ there are a component $Z'$ of $Y$ and a unitary transformation $\delta$
  of the Hermitian matrix of $W_j$ that carries the lattice of $Z$ to that of $Z'$, such that for every
  rational section $f$ of $Y$ of a weight with even last entry there is $c\in\mathbf C$ with
  $(\gamma^*f)_Z(x)=c\,f_{Z'}(\delta x)$ for every frame $x$ of $Z$. The group $\Gamma_Z$ consists of unitary
  transformations congruent to $1$ modulo $M_Y$ on the lattice of $Z$, and it contains every unitary
  transformation that preserves this lattice and is congruent to $1$ modulo $M_Y$ on it.
\item There is an element $k_\gamma$ of the stabilizer of the lattice of $W_j$ in $U(W_j)(\mathbf A_{F,f})$
  such that for every rational section $f$ of $Y$ of a weight with even last entry and every
  $g\in U(W_j)(\mathbf A_F)$, the reading of $\gamma^*f$ at $x(g)$ and the reading of $f$ at $x(gk_\gamma)$
  have the same absolute value.
\item There are a map $Z_2\mapsto Z$ from the components of $Y$ onto those of $X$, a number $N>0$, and
  for each $Z_2$ a unitary transformation $\delta$ of the Hermitian matrix of $W_j$ and $t\in\mathbf C$ with
  $|t|=1$, such that $\delta\Gamma_{Z_2}\delta^{-1}$ is a subgroup of $\Gamma_Z$ of some finite index
  $w(Z_2)$, the indices $w(Z_2)$ of the components above each $Z$ add up to $N$, and $g_{Z_2}(x)=t^nf_Z(\delta x)$
  for every frame $x$ of $Z_2$, whenever $f$ is a rational section of a dominant weight $(\alpha,n)$ on $X$ and
  $g$ is a rational section of $Y$ with $g(y_2)=f(y)$ for every framed point $y_2$ of $Y$ over an
  algebraically closed field containing $F_Y$ and the framed point $y$ of $X$ under it.
\end{enumerate}
\end{lemma}
\leanref{Development/AbelianVarietyTheory/Proofs/FoundationProofs/KeepsCuspOfIsUnifAt.lean}{GlobalTheory.LevelKeeps.exists_carries}
\leanref{Development/AbelianVarietyTheory/Proofs/FoundationProofs/KeepsCuspOfIsUnifAt.lean}{GlobalTheory.LevelKeeps.exists_const}
\leanref{Development/AbelianVarietyTheory/Proofs/FoundationProofs/KeepsCuspOfIsUnifAt.lean}{GlobalTheory.LevelKeeps.mem_group}
\leanref{Development/AbelianVarietyTheory/Proofs/DivisorProofs/KeepsCuspLatticeNorm.lean}{GlobalTheory.LevelIsometry.exists_induces}
\leanref{Development/AbelianVarietyTheory/Proofs/DivisorProofs/LevelCoverComplexModelsComponentsIsometriesDegreStatement.lean}{GlobalTheory.T.levelCover_complexModels_components_isometries_degree}
\begin{proof}
Write $\Lambda_Z$ for the lattice of a component $Z$ of $Y$ at $\sigma'$. The compatibility conditions of
the moduli problem for the uniformization of $Y$ at $\sigma'$ include the following, which in the
construction of Theorem~\ref{cv-deep-pel-facts}(3) come from Situation~\ref{cv-uniformization-situation}
and Lemma~\ref{cv-uniformization-frames}. Every component $Z$ has a lattice $L_Z$ of
Theorem~\ref{cv-uniformization-family}, for the deep level and the type of
Situation~\ref{cv-uniformization-situation} at $\sigma'$, with maps $t_Z$, and a Hermitian model
$(G,\eta,\varphi_Z)$ of $L_Z$ with $\varphi_Z(L_Z)=\Lambda_Z$, where $G$ is the Hermitian matrix of $W_j$ and
$\eta$ does not depend on $Z$. At a frame $x$ of $Z$ the framed point has the point $y_Z(P^Z_x)$ and, as
tangent vectors, the images of the unit vectors of $\mathbf C^g$ under the isomorphism
$\mathbf C^g/P^Z_x(L_Z)\to A_{y_Z(P^Z_x)}$ of that theorem, where $P^Z_x$ is the period of $x$ on $L_Z$
defined through $\varphi_Z$ as in Situation~\ref{cv-uniformization-situation}; the reading $f_Z(x)$ is the
value of $f$ at this framed point. Every complex framed point of $Y$ over $\sigma'$ is triangularly
equivalent to the framed point at a frame of some component, and the framed points at a frame $x$ of $Z$ and
a frame $x'$ of $Z'$ are triangularly equivalent if and only if $Z'=Z$ and $x'\in\Gamma_ZxB$. The group
$\Gamma_Z$ preserves $h_G$ and $\Lambda_Z$, it has finite index in the group of all elements of $GL_r(E)$
that do, and $\gamma x\in xB$ for $\gamma\in\Gamma_Z$ and a frame $x$ of $Z$ only if $\gamma=1$
(Lemma~\ref{cv-uniformization-models}(2) and (3) in the construction). For a rational section $f$ of a dominant
weight $\kappa$, the reading $f_Z$ is holomorphic on the frames of $Z$, and
$f_Z(xb)=\prod_{\tau,i}b_{\tau,ii}^{\lambda_{\tau,i}}f_Z(x)$ for $b\in B$, with the weight
$\lambda=\lambda(\kappa)$ of Situation~\ref{cv-cusp-setting}, which is the same on all components.
The map $g\mapsto(Z(g),x(g))$ of Definition~\ref{cv-deep-cover}(5) comes from markings for an isomorphism
$\iota\colon\Lambda\to L_W$ as in Lemma~\ref{cv-uniformization-lattice}: for every
$g\in U(W_j)(\mathbf A_F)$ there is a marking $(Z(g),\gamma_g)$ of $g$ in the sense of
Definition~\ref{cv-uniformization-marking}(3), for the deep level and the data $\Lambda_Z$, $\varphi_Z$,
$t_Z$ of the components, with $x(g)=\omega(\gamma_gg)$. Each component $Z$ has a base frame $x^0_Z$,
orthonormal for $G$, with $Z(g_Z)=Z$ and $x(g_Z)=x^0_Z$ for some $g_Z$, and for $h\in U(W_j)(F\otimes\mathbf R)$,
$Z(gh)=Z(g)$ and $x(gh)$ lies in $\Gamma_{Z(g)}\,x(g)\big((x^0_\tau)^{-1}\theta_\tau(h)x^0_\tau\big)_\tau$ with
$x^0=x^0_{Z(g)}$ (Lemma~\ref{cv-uniformization-dictionary} in the construction).
Finally, let $x$ be a frame of $Z$ with $|\delta_{\tau,i}(x)|=1$ for all $\tau$ and $i$, put $P=P^Z_x$, let
$p(z)$ be the image of $z\in\mathbf C^g$ in $A_{y_Z(P)}$, and take the line bundle $\mathcal L$, its
trivializations over the sets $U_j$, the transition functions $g_{jj'}$ and the rigidification
$\varrho_j$ as in Theorem~\ref{cv-uniformization-family}(4). Then there are a positive definite Hermitian
form $H_x$ on $\mathbf C^g$, linear in the first variable, with
$\operatorname{Im}H_x(Pa,Pa')=\langle a,a'\rangle$ for $a,a'\in L_Z$, a function $\vartheta_x$ on
$\mathbf C^g$, numbers $\chi_x(u)$ of absolute value $1$ for $u\in P(L_Z)$, and nowhere vanishing
holomorphic functions $F^x_j$ on the sets $\{z:p(z)\in U_j\}$, with $F^x_j(z)=g_{jj'}(p(z))F^x_{j'}(z)$
where both sides are defined, $F^x_j(0)=\varrho_j(y_Z(P))$, and
\[
 F^x_j(z)=\chi_x(u)\exp\big(\pi H_x(z,u)+\tfrac\pi2H_x(u,u)\big)F^x_j(z+u)\qquad(u\in P(L_Z)),
\]
such that the section of the framed point at $x$ has the value $\vartheta_x(z)F^x_j(z)$ at $p(z)$ in the
trivialization over $U_j$, and the mean of $|\vartheta_x(z)|^2e^{-\pi H_x(z,z)}$ over a fundamental domain
of $P(L_Z)$ is $1$ (Lemma~\ref{cv-uniformization-frames}(3)).

The action. The value of $\gamma^*f$ at a framed point $y$ is the value of $f$ at the framed point $\gamma y$,
which has the abelian variety of $y$ through an isomorphism $\nu$ compatible with $\mathcal O_E$ and the
polarizations, the tangent vectors and section carried by $\nu$, and the level structure
$a\mapsto\nu(\eta_y(s(a)))$, where $\eta_y$ is the level structure of $y$ and $s$ is an $\mathcal O_E$-linear
endomorphism of $\Lambda$ that induces the change of level structure.

(1), the component $Z'$ and $\delta$. Fix a frame $x_0$ of $Z$, let $y$ be its framed point and
$P=P^Z_{x_0}$. The framed point $\gamma y$ is triangularly equivalent to the framed point at a frame $x'$ of a
component $Z'$, so with $P'=P^{Z'}_{x'}$ its point is $y_{Z'}(P')$. The isomorphisms of
Theorem~\ref{cv-uniformization-family} and $\nu$ give an isomorphism
$\mathbf C^g/P(L_Z)\to\mathbf C^g/P'(L_{Z'})$ of complex Lie groups that commutes with $\mathcal O_E$ and
preserves the Weil pairings. Lemma~\ref{cv-uniformization-torus-isometry} gives an $\mathcal O_E$-linear
isomorphism $g_1\colon L_Z\to L_{Z'}$ that preserves the forms and a linear automorphism $h$ of
$\mathbf C^g$ with $P'\circ g_1=h\circ P$, through which $\nu$ is read. For
$a\in\bar{\mathfrak M}_Y\Lambda$, the level structure of $\gamma y$ sends $a$ to the image of
$P'(t_{Z'}(a))/M_Y$, and also to the image under $\nu$ of the image of $P(t_Z(s(a)))/M_Y$, which is the
image of $P'(g_1t_Z(s(a)))/M_Y$. As $P'$ is injective and $P'(L_{Z'})$ is the kernel of
$\mathbf C^g\to A_{\gamma y}$, $g_1(t_Z(s(a)))\in t_{Z'}(a)+M_YL_{Z'}$. By
Lemma~\ref{cv-uniformization-model-match} there is $\delta\in GL_r(E)$ with
$\varphi_{Z'}\circ g_1=\delta\circ\varphi_Z$; it is unitary for $G$ and carries $\Lambda_Z$ onto
$\Lambda_{Z'}$.

Say that $\delta\in GL_r(E)$ \emph{carries} $Z$ to $Z'$ if it preserves $h_G$ and
$\varphi_{Z'}\circ g_1=\delta\circ\varphi_Z$ for an $\mathcal O_E$-linear isomorphism $g_1\colon L_Z\to L_{Z'}$
that preserves the forms and satisfies $g_1(t_Z(s(a)))\in t_{Z'}(a)+M_YL_{Z'}$ for
$a\in\bar{\mathfrak M}_Y\Lambda$. The $\delta$ just found carries $Z$ to $Z'$, and the readings and the
constant below are proved for any $\delta$ that carries $Z$ to $Z'$, with its $g_1$.

(1), the readings. Let $x$ be a frame of $Z$. As $\delta$ preserves $h_G$, $\delta x$, with
$(\delta x)_\tau=\tau(\delta)x_\tau$, is a frame of $Z'$ with the same leading minors, and as $\tau(\delta)$
preserves $h_{\tau(G)}$, the formula for periods gives $P^{Z'}_{\delta x}\circ g_1=P^Z_x$. The complex points
of $Y$ classify the abelian varieties with their structures up to isomorphism. The framed point at $x$,
carried by $\nu$ with its level structure changed through $s$, and the framed point at $\delta x$ are
identified by $g_1$, which preserves the forms, carries one period to the other, and by the congruence
above matches the level structures. So $\gamma$ applied to the framed point at $x$ has the point
$y_{Z'}(P^{Z'}_{\delta x})$; and since an automorphism of these data that fixes the level structure has
finite order and is congruent to $1$ modulo $M_Y\ge3$, hence is trivial \cite{Min87}, $\nu$ is read by the
identity of $\mathbf C^g$. So the two framed points have the same tangent vectors, and the section of the
first is $t(x)$ times that of the second for some $t(x)\in\mathbf C^\times$. Hence
$(\gamma^*f)_Z(x)=t(x)^nf_{Z'}(\delta x)$ for every rational section $f$ of a weight with last entry $n$.

(1), the constant. Let $\kappa_2$ be the dominant weight with $\alpha=0$ and $n=2$. The action changes
only the level structures, so $\gamma^*f$ has weight $\kappa_2$ when $f$ has.

Step 1: if $|\delta_{\tau,i}(x)|=1$ for all $\tau$ and $i$, then $|t(x)|=1$. The frame $\delta x$ of $Z'$ has
the same minors, so the data listed at the beginning of the proof exist at $x$ for $Z$, with $P=P^Z_x$,
$H_x$, $\chi_x$, $\vartheta_x$ and $F^x$, and at $\delta x$ for $Z'$, with $P'=P^{Z'}_{\delta x}$,
$H'=H_{\delta x}$, $\chi'=\chi_{\delta x}$, $\vartheta'=\vartheta_{\delta x}$ and $F'=F^{\delta x}$. As
$P'\circ g_1=P$, both periods have the image $\Pi=P(L_Z)$, and
$\operatorname{Im}H'(Pa,Pa')=\langle g_1a,g_1a'\rangle=\langle a,a'\rangle=\operatorname{Im}H_x(Pa,Pa')$ for
$a,a'\in L_Z$. A Hermitian form $H$ is determined by its imaginary part, since
$\operatorname{Re}H(z,w)=\operatorname{Im}H(iz,w)$, and $\Pi$ spans $\mathbf C^g$ over $\mathbf R$; so
$H'=H_x$. Let $y_1=y_Z(P)$ and $y_2=y_{Z'}(P')$ be the two points, $p\colon\mathbf C^g\to A_{y_1}$ and
$p'\colon\mathbf C^g\to A_{y_2}$ the maps of the uniformization; as $\nu$ is read by the identity,
$p'=\nu\circ p$. As $\nu$ preserves the polarizations, it identifies the line bundles $\mathcal L$ on
$A_{y_1}$ and on $A_{y_2}$, compatibly with their rigidifications. On the open set $O_{j'j}$ of the
$z\in\mathbf C^g$ with $p(z)\in U_j$ and $p'(z)\in U_{j'}$, this identification multiplies values in the
trivialization over $U_j$ by a nowhere vanishing holomorphic function $\omega_{j'j}$, to give values in the
trivialization over $U_{j'}$; it depends on $z$ only through $p(z)$, satisfies
$\omega_{j'j}(z)g_{jk}(p(z))=g_{j'k'}(p'(z))\omega_{k'k}(z)$ on $O_{j'j}\cap O_{k'k}$, and satisfies
$\omega_{j'j}(0)\varrho_j(y_1)=\varrho_{j'}(y_2)$. The sets $O_{j'j}$ cover $\mathbf C^g$ and are stable
under translation by $\Pi$. The two families $(\omega_{j'j}F^x_j)$ and $(F'_{j'})$, indexed by the pairs
$(j',j)$, satisfy the hypotheses of Lemma~\ref{cv-uniformization-factor-unique} with the factors
$g_{j'k'}(p'(z))$ for the change from $(j',j)$ to $(k',k)$ and with
$m(u,z)=\exp(\pi\operatorname{Re}H_x(z,u)+\frac\pi2H_x(u,u))$: hypothesis (1) by the transition laws of
$F^x$ and $F'$ and the relation for $\omega$, (2) because both take the value $\varrho_{j'}(y_2)$ at
$z=0$, and (3) by the formulas for $F^x$ and $F'$, as $H'=H_x$, $|\chi_x|=|\chi'|=1$ and $\omega_{j'j}$ is
periodic for $\Pi$. So $F'_{j'}=\omega_{j'j}F^x_j$ on $O_{j'j}$. The section of $\gamma$ applied to the
framed point at $x$ is carried by $\nu$, so its value at $p'(z)$ over $U_{j'}$ is
$\omega_{j'j}(z)\vartheta_x(z)F^x_j(z)$. It is also $t(x)$ times the section of the framed point at
$\delta x$, whose value there is $\vartheta'(z)F'_{j'}(z)$. Hence $\vartheta_x=t(x)\vartheta'$. Since the
values of a section at $p(z)$ and $p(z+u)$ agree, the formula for $F^x$ shows that
$|\vartheta_x(z)|^2e^{-\pi H_x(z,z)}$ is periodic for $\Pi$, and so is $|\vartheta'(z)|^2e^{-\pi H_x(z,z)}$;
so their means over a fundamental domain of $\Pi$ do not depend on the domain. Both means are $1$, so
$|t(x)|^2=1$.

Step 2: for every frame $x$ of $Z$ there is $c(x)$ with $|c(x)|=1$ and $(\gamma^*f)_Z(x)=c(x)f_{Z'}(\delta x)$
for every rational section $f$ of weight $\kappa_2$. By Lemma~\ref{cv-uniformization-unit-minors} there
is $b\in B$ such that $x_0=xb$ has minors of absolute value $1$. The readings and Step 1 give
$(\gamma^*f)_Z(x_0)=t(x_0)^2f_{Z'}(\delta x_0)$ with $|t(x_0)|=1$. Both $\gamma^*f$ and $f$ have weight
$\kappa_2$, the weight $\lambda(\kappa_2)$ is the same on $Z$ and $Z'$, and $\delta x_0=(\delta x)b$; so both
sides change by the same nonzero factor $\prod_{\tau,i}b_{\tau,ii}^{\lambda(\kappa_2)_{\tau,i}}$ between $x$
and $x_0$, and $c(x)=t(x_0)^2$ works.

Step 3: there is $\zeta$ with $|\zeta|=1$ and $(\gamma^*f)_Z(x)=\zeta f_{Z'}(\delta x)$ for every frame $x$
of $Z$ and every rational section $f$ of weight $\kappa_2$. The frames of $Z$ form a connected open set $O$
(Lemma~\ref{cv-uniformization-frames-connected}), on which the functions $P_f=(\gamma^*f)_Z$ and
$Q_f(x)=f_{Z'}(\delta x)$ are holomorphic. If every $Q_f$ vanishes on $O$, then so does every $P_f$ by
Step 2, and $\zeta=1$ works. Otherwise choose $f_0$ and $x_1\in O$ with $Q_{f_0}(x_1)\ne0$. By Step 2,
$|P_{f_0}|=|Q_{f_0}|$ on $O$, so Lemma~\ref{cv-uniformization-modulus} gives $\zeta$ with
$P_{f_0}=\zeta Q_{f_0}$. On a neighbourhood of $x_1$ where $Q_{f_0}$ does not vanish, Step 2 forces
$c(x)=\zeta$; so $|\zeta|=1$, and for every $f$ the holomorphic function $P_f-\zeta Q_f$ vanishes near
$x_1$, hence on $O$ by the identity theorem.

Step 4: $t(x)^2=\zeta$ for every frame $x$ of $Z$. By Lemma~\ref{cv-uniformization-unit-minors} there is
$b\in B$ such that $(\delta x)b$ has minors of absolute value $1$. By Lemma~\ref{cv-uniformization-theta-section},
for $Y$ at $\sigma'$, the component $Z'$, the frame $(\delta x)b$ and $n=2$, some rational section $f_2$ of $Y$ of
weight $\kappa_2$ has a nonzero reading at $(\delta x)b$. Then $f_{2,Z'}(\delta x)\ne0$, since the two
readings differ by the factor $\prod_{\tau,i}b_{\tau,ii}^{\lambda(\kappa_2)_{\tau,i}}$, and the readings and
Step 3 give $t(x)^2f_{2,Z'}(\delta x)=(\gamma^*f_2)_Z(x)=\zeta f_{2,Z'}(\delta x)$, so $t(x)^2=\zeta$. Hence
for a weight with even last entry $n$, $(\gamma^*f)_Z(x)=\zeta^{n/2}f_{Z'}(\delta x)$ for every frame $x$ of
$Z$, which is the first assertion of (1) with $c=\zeta^{n/2}$, and $|c|=1$.

(1), the group $\Gamma_Z$. Let $\Gamma^1_Z$ be the group of $\beta\in GL_r(E)$ that preserve $h_G$ and
$\Lambda_Z$ and satisfy $\beta v-v\in M_Y\Lambda_Z$ for every $v\in\Lambda_Z$. The second assertion of (1) is
$\Gamma_Z=\Gamma^1_Z$.

Step A: for $\beta\in\Gamma^1_Z$ and every frame $x$ of $Z$, $\beta x\in\Gamma_ZxB$. As each $\tau(\beta)$
preserves the form of $\tau(G)$, $\beta x$ is a frame of $Z$ with the minors of $x$. The map
$\beta'=\varphi_Z^{-1}\circ\beta\circ\varphi_Z$ is an $\mathcal O_E$-linear automorphism of $L_Z$. It preserves
the form of $L_Z$, which is $\operatorname{Tr}_{E/\mathbf Q}(\eta\,h_G(\varphi_Z\,\cdot,\varphi_Z\,\cdot))$, and
$\beta'(a)\in a+M_YL_Z$ for $a\in L_Z$. The formula for periods gives $P^Z_{\beta x}\circ\beta'=P^Z_x$. By
Theorem~\ref{cv-uniformization-family}(2), with $h=1$, the two periods give the same point $y$. The
isomorphisms $\mathbf C^g/P^Z_x(L_Z)\to A_y$ and $\mathbf C^g/P^Z_{\beta x}(L_Z)\to A_y$ have the same source
and differ by an automorphism of $A_y$ that preserves the polarization and the level structure; it has
finite order and is congruent to $1$ modulo $M_Y\ge3$ on the lattice, so it is the identity \cite{Min87}. So
the framed points at $x$ and at $\beta x$ have the same point and the same tangent vectors, and their
sections, which span the same line, differ by a nonzero factor: the two framed points are triangularly
equivalent. By the compatibility conditions, $\beta x\in\Gamma_ZxB$.

Step B: for $\beta\in\Gamma^1_Z$ there are $\beta''\in\Gamma_Z$ and $u\in E^\times$ with $\beta''=u\beta$. This
is Lemma~\ref{cv-uniformization-scalar-orbits} for the group $\Gamma_Z$, whose hypothesis is Step A.

Step C: $u$ is a root of unity. By the compatibility conditions there is a finite set $S$ of elements of
$GL_r(E)$ that preserve $h_G$ and $\Lambda_Z$ such that every such element lies in $s\Gamma_Z$ for some
$s\in S$. For $i\ge0$ choose $s_i\in S$ with $s_i^{-1}\beta^i\in\Gamma_Z$. As the scalar $u$ commutes with
$\beta$, the element $\beta''^i(s_i^{-1}\beta^i)^{-1}=u^is_i$ lies in $\Gamma_Z$. Two indices $i<i'$ have
$s_i=s_{i'}$, and then $u^{i'-i}=(u^{i'}s_{i'})(u^is_i)^{-1}\in\Gamma_Z$. For a frame $x$ of $Z$,
$u^{i'-i}x=xb$ with $b$ the family of scalar matrices $\tau(u^{i'-i})$, which lies in $B$. So $u^{i'-i}=1$.

Step D: $u=1$, so $\Gamma^1_Z\subset\Gamma_Z$. Let $x^0=x^0_Z$ be the base frame of $Z$ and $g_0=g_Z$, so that
$Z(g_0)=Z$ and $x(g_0)=x^0$. For $\tau\in\Phi$ let $A_\tau(u,v)=u^{\mathsf T}\tau(G)\bar v=h_{\tau(G)}(u,v)$ on
$\mathbf C^r$: the form of $\tau(G)$ in Step A, which $\tau(\beta)$ preserves for $\beta\in\Gamma_Z$, and whose
Gram matrix $x_\tau^{\mathsf T}\tau(G)\bar x_\tau$ on the columns of $x_\tau$ has the leading minors
$\delta_{\tau,i}(x)$, for every frame $x$ (Situation~\ref{cv-cusp-setting}). Let $G_W$ be the Gram matrix of the
basis of $W_j$, so that the Hermitian form of $W_j$ at $\tau$ is $(u,v)\mapsto u^*\tau(G_W)v$, and let
$B_\tau(u,v)=u^{\mathsf T}\tau(G_W)^{\mathsf T}\bar v$, which is this form with its two arguments exchanged; a
linear automorphism of $\mathbf C^r$ preserves $B_\tau$ if and only if it preserves the form of $W_j$ at $\tau$.
As the $\theta_\tau$ identify $(F\otimes\mathbf R)\otimes_FE$ with $\prod_{\tau\in\Phi}\mathbf C$, they identify
$U(W_j)(F\otimes\mathbf R)$ with the product over $\tau\in\Phi$ of the groups of linear automorphisms of
$\mathbf C^r$ that preserve $B_\tau$.

The unitarity of $\beta''$: $\tau(\beta'')$ preserves $B_\tau$ for every $\tau\in\Phi$. Fix $\tau$, let
$v$ be the first column of $x^0_\tau$, and let $w$ be a linear automorphism of $\mathbf C^r$ that preserves
$B_\tau$. Take $h\in U(W_j)(F\otimes\mathbf R)$ with $\theta_\tau(h)=w$ and $\theta_{\tau'}(h)=1$ for $\tau'\ne\tau$.
By the compatibility conditions, $x(g_0h)\in\Gamma_Zx'$ for the family $x'$ with $x'_\tau=wx^0_\tau$ and
$x'_{\tau'}=x^0_{\tau'}$ for $\tau'\ne\tau$. The minors of $x(g_0h)$ have absolute value $1$
(Definition~\ref{cv-deep-cover}(5)), and elements of $\Gamma_Z$ do not change minors, so
$|A_\tau(wv,wv)|=|\delta_{\tau,1}(x')|=1$. By Lemma~\ref{cv-uniformization-unitary-orbit}, $A_\tau=tB_\tau$ for a
real $t\ne0$. As $\beta''\in\Gamma_Z$, $\tau(\beta'')$ preserves $A_\tau$, hence $B_\tau$.

The congruence. By the unitarity of $\beta''$ there is $h\in U(W_j)(F\otimes\mathbf R)$ with
$\theta_\tau(h)=\tau(\beta'')$ for every $\tau\in\Phi$. By the compatibility conditions,
$Z(g_0h)=Z$ and $x(g_0h)\in\Gamma_Z\,x^0\big((x^0_\tau)^{-1}\tau(\beta'')x^0_\tau\big)_\tau=\Gamma_Z\beta''x^0=
\Gamma_Zx(g_0)$. By Definition~\ref{cv-deep-cover}(5), $g_0h=\beta_0g_0k$ with $\beta_0\in U(W_j)(F)$ and
$k\in K_Y$. Let $\delta\in GL_r(E)$ be the matrix of $\beta_0$. As $h$ is archimedean and $k$ is finite, the
finite parts give $\delta g_{0,f}=g_{0,f}k^{-1}$, and the archimedean parts give
$\tau(\delta)\theta_\tau(g_{0,\infty})=\theta_\tau(g_{0,\infty})\theta_\tau(h)$ for $\tau\in\Phi$. Let
$\Lambda'=E^r\cap g_{0,f}\hat L_0$, a lattice in $E^r$. For $v=g_{0,f}z\in\Lambda'$ with $z\in\hat L_0$,
$\delta v-v=g_{0,f}(k^{-1}z-z)\in M_Yg_{0,f}\hat L_0$, since $k^{-1}\in K_Y$; as $\delta v-v\in E^r$, it lies in
$M_Y\Lambda'$. Fix $\tau\in\Phi$ and put $Q=\theta_\tau(g_{0,\infty})^{-1}$. Then
$\tau(u\beta)Q=\theta_\tau(h)Q=Q\tau(\delta)$. Lemma~\ref{cv-uniformization-root-congruence}, for
$\Lambda_Z$, $\Lambda'$, $\beta$, $\delta$ and $M_Y\ge3$, gives $u=1$. So $\beta=\beta''\in\Gamma_Z$.

Step E: $\Gamma_Z\subset\Gamma^1_Z$. Let $\beta\in\Gamma_Z$; it preserves $h_G$ and $\Lambda_Z$. Let $x$ be a
frame of $Z$. By the compatibility conditions the framed points at $x$ and at $\beta x$ are triangularly
equivalent, so they have the same point. By Theorem~\ref{cv-uniformization-family}(2),
$P^Z_{\beta x}=h\circ P^Z_x\circ\beta_1$ for a $\mathbf C$-linear automorphism $h$ of $\mathbf C^g$ and an
automorphism $\beta_1$ of $L_Z$ that preserves the form and is congruent to $1$ modulo $M_Y$ on
$\bar{\mathfrak M}_YL_Z$, hence on $L_Z$, as in the proof of Lemma~\ref{cv-uniformization-models}(3). The
$E$-linear extension $\delta$ of $\varphi_Z\circ\beta_1^{-1}\circ\varphi_Z^{-1}$ preserves $\Lambda_Z$ and is
congruent to $1$ modulo $M_Y$ on it, and it preserves $h_G$ as in the proof of
Lemma~\ref{cv-uniformization-model-match}; so $\delta\in\Gamma_Z$ by Step D. The formula for periods gives
$P^Z_{\delta x}=P^Z_x\circ\beta_1$, so $P^Z_{\beta x}=h\circ P^Z_{\delta x}$. The map
$v\mapsto(\tau(\varphi_Zv))_{\tau\in\Phi}$ identifies $L_Z\otimes\mathbf R$ with $\prod_{\tau\in\Phi}\mathbf C^r$.
At each $\tau\in\Phi$ the coordinates $(\tau,a)$ of the period of a frame $x'$ are $\mathbf C$-linear functions
of $\tau(\varphi_Zv)$, and the coordinates $(\bar\tau,a)$ are the conjugate-linear functions
$h_{\tau(G)}(x'_{\tau,q-a},\cdot)$ of it. As $h$ is $\mathbf C$-linear, comparing the parts of both sides that
are conjugate-linear in the component at $\tau$ shows that each function
$h_{\tau(G)}(\tau(\beta)x_{\tau,c},\cdot)$ with $c\le q$ is a combination of the functions
$h_{\tau(G)}(\tau(\delta)x_{\tau,c'},\cdot)$ with $c'\le q$. As $h_{\tau(G)}$ is nondegenerate, the first $q$
columns of $\tau(\beta)x_\tau$ lie in the span of the first $q$ columns of $\tau(\delta)x_\tau$. That is,
$\beta x=\delta x\,m$ for a family $m$ of invertible matrices with $(m_\tau)_{ac}=0$ for $c\le q<a$. Then
$\delta^{-1}\beta\in\Gamma_Z$ and $(\delta^{-1}\beta)x=xm$, so $\delta^{-1}\beta=1$ by
Lemma~\ref{cv-uniformization-block-free}. Hence $\beta=\delta\in\Gamma^1_Z$.

(2). Let $\hat L_0$, $\hat L_h$ and $\omega$ be as in Definition~\ref{cv-uniformization-marking}, write
$\hat L_W=\hat L_1$ for the completion of $L_W=L_1$, and let $k_\gamma$ be an element of the stabilizer of
$\hat L_0$ in
$U(W_j)(\mathbf A_{F,f})$ with $k_\gamma\iota(a)-\iota(s(a))\in M_Y\hat L_W$ for every
$a\in\bar{\mathfrak M}_Y\Lambda$, which exists by Lemma~\ref{cv-uniformization-level-lift} for the level
$(M_Y,\mathfrak M_Y)$. Let $g\in U(W_j)(\mathbf A_F)$, put $Z=Z(g)$ and $Z'=Z(gk_\gamma)$, and let
$\gamma_1=\gamma_g$ and $\gamma_2=\gamma_{gk_\gamma}$ be as in the markings listed at the beginning of the
proof, so that $x(g)=\omega(u)$ with $u=\gamma_1g$, and $x(gk_\gamma)=\omega(\gamma_2gk_\gamma)$. Let
$\delta\in GL_r(E)$ be the matrix of $\gamma_2\gamma_1^{-1}\in U(W_j)(F)$, so that
$\gamma_2gk_\gamma=\delta uk_\gamma$ and $\delta$ preserves $h_G$. As $\delta$ is rational and $k_\gamma$
stabilizes $\hat L_0$, $\Lambda_{Z'}=L_{\delta uk_\gamma}=\delta L_u=\delta\Lambda_Z$ and
$\hat L_{\delta uk_\gamma}=\delta\hat L_u$. So $g_1=\varphi_{Z'}^{-1}\circ\delta\circ\varphi_Z$ is an
$\mathcal O_E$-linear isomorphism of $L_Z$ onto $L_{Z'}$, and it preserves the forms, since the models of
$L_Z$ and $L_{Z'}$ have the same $G$ and $\eta$ and $\delta$ preserves $h_G$. For
$a\in\bar{\mathfrak M}_Y\Lambda$, also $s(a)\in\bar{\mathfrak M}_Y\Lambda$, and
\[
 \delta\varphi_Z(t_Z(s(a)))-\varphi_{Z'}(t_{Z'}(a))
 =\delta\big(\varphi_Z(t_Z(s(a)))-u_f\iota(s(a))\big)
 -\big(\varphi_{Z'}(t_{Z'}(a))-(\delta uk_\gamma)_f\iota(a)\big)
 -\delta u_f\big(k_\gamma\iota(a)-\iota(s(a))\big).
\]
The three terms on the right lie in $M_Y\hat L_{\delta uk_\gamma}$: the first by the marking of $g$, the
second by the marking of $gk_\gamma$, the third by the choice of $k_\gamma$, as
$\delta u_f\hat L_W=\hat L_{\delta u}=\hat L_{\delta uk_\gamma}$. The left side lies in $E^r$, hence in
$E^r\cap M_Y\hat L_{\delta uk_\gamma}=M_Y\Lambda_{Z'}$, as a lattice in $E^r$ is the intersection of $E^r$
with its completion; so $g_1(t_Z(s(a)))\in t_{Z'}(a)+M_YL_{Z'}$, and $\delta$ carries $Z$ to $Z'$. As
$k_\gamma$ is finite and $\delta$ is rational, $\theta_\tau((\delta uk_\gamma)_\infty)=\tau(\delta)\theta_\tau(u_\infty)$
for $\tau\in\Phi$, so $x(gk_\gamma)=\delta\,\omega(u)=\delta x(g)$. By (1) for this $\delta$, for a rational
section $f$ of a weight with even last entry there is $c$ with $|c|=1$ and
$(\gamma^*f)_Z(x(g))=c\,f_{Z'}(\delta x(g))=c\,f_{Z(gk_\gamma)}(x(gk_\gamma))$. This is (2).

(3), the setting. The compatibility conditions listed at the beginning of the proof hold also for the
uniformization of $X$ at $\sigma$, with the level $(M,\mathfrak M)$ in place of the deep level
(Theorem~\ref{cv-deep-pel-facts}(3) for this level); write $L_Z$, $t_Z$, $\varphi_Z$, $\Lambda_Z$ and $P^Z_x$ for its
data, with the same $G$ and $\eta$. For a framed point $y_2$ of $Y$ over $\mathbf C$, the \emph{framed point of
$X$ under} $y_2$ has the point $\pi_Y(y_2)$ and the tangent vectors and section of $y_2$ carried by the
identification of Definition~\ref{cv-deep-cover}(4). As this identification is an isomorphism of abelian
varieties compatible with $\mathcal O_E$ and the polarizations, two framed points of $Y$ over the same framed
point of $X$ with the same point are equal; for every point of $Y$ over the point of a framed point $y$ of
$X$ there is a framed point of $Y$ with that point over $y$; and if two framed points of $Y$ are
triangularly equivalent, so are the framed points of $X$ under them, while conversely a triangular change
of a framed point of $X$ lifts, with the same factors, to framed points over it with the same point.

(3), the components and $\delta$. Let $Z_2$ be a component of $Y$. Fix a frame $x_0$ of $Z_2$, with framed point
$y_2$ and period $P_0=P^{Z_2}_{x_0}$. The framed point of $X$ under $y_2$ is triangularly equivalent to the framed
point at a frame $x_1$ of a component $Z$ of $X$, so with $P_1=P^Z_{x_1}$, $\pi_Y(y_{Z_2}(P_0))=y_Z(P_1)$. The
identification of Definition~\ref{cv-deep-cover}(4) and the isomorphisms of
Theorem~\ref{cv-uniformization-family} give an isomorphism $\mathbf C^g/P_0(L_{Z_2})\to\mathbf C^g/P_1(L_Z)$ of
complex Lie groups that commutes with $\mathcal O_E$ and preserves the Weil pairings, and
Lemma~\ref{cv-uniformization-torus-isometry} gives an $\mathcal O_E$-linear isomorphism $g_1\colon L_{Z_2}\to L_Z$
that preserves the forms and a linear automorphism $h$ of $\mathbf C^g$ with $P_1\circ g_1=h\circ P_0$, through
which the identification is read. Let $a\in\bar{\mathfrak M}_Y\Lambda$. As $M_Y=15^eM$, $15^e$ times the level
structure of $y_2$ sends $a$ to the image of $P_0(t_{Z_2}(a))/M$; the identification carries it to the image of
$P_1(g_1t_{Z_2}(a))/M$, and it is the level structure of $X$ at $a$, the image of $P_1(t_Z(a))/M$. So
$g_1(t_{Z_2}(a))\in t_Z(a)+ML_Z$. By Lemma~\ref{cv-uniformization-model-match} there is $\delta\in GL_r(E)$ with
$\varphi_Z\circ g_1=\delta\circ\varphi_{Z_2}$; it preserves $h_G$ and carries $\Lambda_{Z_2}$ onto $\Lambda_Z$. Send $Z_2$
to $Z$.

Let $x$ be a frame of $Z_2$. As $\delta$ preserves $h_G$, $\delta x$ is a frame of $Z$ with the minors of $x$, and the
formula for periods gives $P^Z_{\delta x}\circ g_1=P^{Z_2}_x$. The isomorphism $g_1$ preserves the forms, carries
one period to the other, and by the congruence above matches $15^e$ times the level structure of the point
$y_{Z_2}(P^{Z_2}_x)$ with that of $y_Z(P^Z_{\delta x})$. As the complex points of $X$ classify the abelian varieties
with their structures up to isomorphism, $\pi_Y(y_{Z_2}(P^{Z_2}_x))=y_Z(P^Z_{\delta x})$; and as an automorphism of
these data has finite order and is congruent to $1$ modulo $M\ge3$, hence is trivial \cite{Min87}, the
identification of Definition~\ref{cv-deep-cover}(4) is read by the identity of $\mathbf C^g$. So the framed point of
$X$ under the framed point at $x$ has the point and the tangent vectors of the framed point at $\delta x$, and its
section is $t(x)$ times the section of the framed point at $\delta x$, for some $t(x)\in\mathbf C^\times$. If
$|\delta_{\tau,i}(x)|=1$ for all $\tau$ and $i$, then $|t(x)|=1$: the argument of Step~1 of (1) applies, with the
data at unit frames of $Y$ at $x$ and of $X$ at $\delta x$, and the identification of
Definition~\ref{cv-deep-cover}(4) in place of $\nu$.

(3), the readings. Let $f$ be a rational section of a dominant weight $\kappa=(\alpha,n)$ on $X$ and $g$ a rational
section of $Y$ with $g(y_2)=f(y)$ as in the statement. The identification of Definition~\ref{cv-deep-cover}(4)
carries changes of frame to the same changes of frame, so $g$ has the transformation law of the weight
$\kappa$. For a frame $x$ of $Z_2$, the framed point of $X$ under the framed point at $x$ has the section $t(x)$
times that at $\delta x$, so $g_{Z_2}(x)=t(x)^nf_Z(\delta x)$. By the Gram--Schmidt process for the form of $\tau(G)$
there is $b\in B$ such that $xb$ is orthonormal for $G$; its minors are $\pm1$, so $c(x)=t(xb)$ has $|c(x)|=1$. The
readings $g_{Z_2}$ and $f_Z$ have the transformation law under $B$ of the weight $\lambda(\kappa)$, which is the same
on $Z_2$ and $Z$, and $\delta(xb)=(\delta x)b$; so $g_{Z_2}(x)=c(x)^nf_Z(\delta x)$, with $c(x)$ independent of $f$ and
$g$. The frames of $Z_2$ form a nonempty connected open set (Lemma~\ref{cv-uniformization-frames-connected}), on
which $g_{Z_2}$ and $x\mapsto f_Z(\delta x)$ are holomorphic. Lemma~\ref{cv-uniformization-uniform-constant}, for the
family of all such pairs $(f,g)$ and the function $c$, gives $t$ with $|t|=1$ and $g_{Z_2}(x)=t^nf_Z(\delta x)$ for
every frame $x$ of $Z_2$ and every pair.

(3), the groups. Let $\gamma_2\in\Gamma_{Z_2}$ and $h=\delta\gamma_2\delta^{-1}$. Let $x'$ be a frame of $Z$. Then
$x=\delta^{-1}x'$ is a frame of $Z_2$, as $\delta^{-1}$ preserves $h_G$, and $hx'=\delta(\gamma_2x)$. The framed points of
$Y$ at $x$ and at $\gamma_2x$ are triangularly equivalent, by the compatibility conditions, and so are the framed
points of $X$ under them. By the paragraph on the components these are triangularly equivalent to the framed
points at $x'$ and at $hx'$, so the framed points at $x'$ and $hx'$ are triangularly equivalent, and
$hx'\in\Gamma_Zx'B$ by the compatibility conditions for $X$. As this holds for every frame $x'$ of $Z$,
Lemma~\ref{cv-uniformization-scalar-orbits} gives $\gamma'\in\Gamma_Z$ and $u\in E^\times$ with $\gamma'=uh$.

The congruence (C): for a component $Z$ of $X$, $\gamma\in\Gamma_Z$, $\tau\in\Phi$ and an eigenvalue $l$ of
$\tau(\gamma)$, the number $(l-1)/M$ is an algebraic integer; and likewise for the components of $Y$, with $M_Y$ in
place of $M$. Consider $X$; for $Y$ the argument is the same at the deep level. Let $K$ be the principal congruence
subgroup of level $M$ of the stabilizer of the lattice of $W_j$, let $x^0=x^0_Z$ be the base frame of $Z$, and let
$g_0$ have $Z(g_0)=Z$ and $x(g_0)=x^0$. As $\gamma$ preserves $h_G$ and $h_G(u,v)=-h_{W_j}(v,u)$
(Lemma~\ref{cv-uniformization-lattice}), $\gamma$ preserves $h_{W_j}$, and there is $h_\gamma\in U(W_j)(F\otimes\mathbf R)$
with $\theta_\tau(h_\gamma)=\tau(\gamma)$ for every $\tau\in\Phi$. The paragraph on the congruence in Step D of (1), with
$\gamma$ in place of $\beta''$, $h_\gamma$ in place of $h$ and $K$ in place of $K_Y$, gives $\delta'\in GL_r(E)$, a lattice
$\Lambda'\subset E^r$ with $\delta'v-v\in M\Lambda'$ for $v\in\Lambda'$, and an invertible matrix $Q$ with
$\tau(\gamma)Q=Q\tau(\delta')$. So $l$ is an eigenvalue of $\tau(\delta')$, and $(l-1)/M$ is an algebraic integer by
Step~1 of the proof of Lemma~\ref{cv-uniformization-root-congruence}.

Then $u=1$. As $h$ and $\gamma'$ preserve $h_G$, $h_G(v,w)=h_G(\gamma'v,\gamma'w)=u\bar u\,h_G(hv,hw)=u\bar u\,h_G(v,w)$,
and as $G\ne0$, $u\bar u=1$; as $E$ is a CM field, $|\varphi(u)|=1$ for every embedding $\varphi\colon E\to\mathbf C$.
Fix $\tau\in\Phi$, an eigenvalue $l$ of $\tau(\gamma')$ and a nonzero row vector $w$ with $w\tau(\gamma')=lw$. Then
$w\tau(\gamma'^{-1})=l^{-1}w$, and the nonzero row vector $w\tau(\delta)$ satisfies
$w\tau(\delta)\tau(\gamma_2)=w\tau(h)\tau(\delta)=\tau(u)^{-1}l\,w\tau(\delta)$. By (C), $a'=(l^{-1}-1)/M$ and
$b'=(\tau(u)^{-1}l-1)/M_Y$ are algebraic integers. As $\tau(u)^{-1}=(1+M_Yb')(1+Ma')$ and $M_Y=15^eM$,
$\tau\big((u^{-1}-1)/M\big)=a'+15^eb'+M_Ya'b'$ is an algebraic integer, and so is $(u^{-1}-1)/M$. Its conjugates have
absolute value at most $2/M<1$, and a nonzero algebraic integer has a conjugate of absolute value at least $1$, as
its norm is a nonzero integer. So $u=1$, and $h=\gamma'\in\Gamma_Z$. Hence $\delta\Gamma_{Z_2}\delta^{-1}\subset\Gamma_Z$.

The index. The lattices $\delta\Lambda_{Z_2}$ and $\Lambda_Z$ are finitely generated and span $E^r$, so
$m\delta\Lambda_{Z_2}\subset\Lambda_Z$ and $m\Lambda_Z\subset\delta\Lambda_{Z_2}$ for some integer $m\ge1$. The group
$\Gamma_Z$ preserves $\Lambda_Z$ and acts on the finite group $\Lambda_Z/m^2\Lambda_Z$; let $K'$ be the kernel, a subgroup of
finite index in $\Gamma_Z$. For $\gamma\in K'$ and $v\in\delta\Lambda_{Z_2}$, $mv\in\Lambda_Z$, so
$\gamma(mv)-mv\in m^2\Lambda_Z$ and $\gamma v-v\in m\Lambda_Z\subset\delta\Lambda_{Z_2}$; as $K'$ is a group, it preserves
$\delta\Lambda_{Z_2}$. By the
compatibility conditions for $Z_2$, conjugated by $\delta$, $\delta\Gamma_{Z_2}\delta^{-1}$ has finite index in the
group of elements of $GL_r(E)$ that preserve $h_G$ and $\delta\Lambda_{Z_2}$, which contains $K'$. So
$\delta\Gamma_{Z_2}\delta^{-1}$ has finite index $w(Z_2)$ in $\Gamma_Z$.

(3), the indices. Let $Z$ be a component of $X$, $x'$ a frame of $Z$ and $y'$ its framed point. For a framed point $y_2$
of $Y$ over $y'$, choose a frame $x_2$ of a component $Z_2$ of $Y$ whose framed point is triangularly equivalent to
$y_2$. The framed point $y'$ is then triangularly equivalent to the framed point of $X$ under the framed point at
$x_2$, hence to the framed point at $\delta x_2$, a frame of the image of $Z_2$; so by the compatibility conditions
for $X$, $Z_2$ maps to $Z$ and $x'=\gamma\delta x_2b$ with $\gamma\in\Gamma_Z$ and $b\in B$. The component $Z_2$ is
determined by $y_2$, by the compatibility conditions for $Y$. If $x'_2$ is another such frame, with $\gamma'$ and $b'$,
then $x'_2=\gamma_2x_2b_2$ with $\gamma_2\in\Gamma_{Z_2}$ and $b_2\in B$, so $\gamma\delta x_2b=\gamma'(\delta\gamma_2\delta^{-1})\delta x_2b_2b'$,
and as $\Gamma_Z$ acts freely modulo $B$, $\gamma=\gamma'\delta\gamma_2\delta^{-1}$. So $y_2\mapsto(Z_2,\gamma\,\delta\Gamma_{Z_2}\delta^{-1})$
is a well defined map to the pairs of a component $Z_2$ over $Z$ and a coset of $\delta\Gamma_{Z_2}\delta^{-1}$ in
$\Gamma_Z$. It is injective: if $y_2$ and $y'_2$ give the same pair, with frames $x_2$ and $x'_2$ of $Z_2$ and
$\gamma^{-1}\gamma'=\delta\gamma_2\delta^{-1}$, then $\gamma_2^{-1}x_2$ and $x'_2$ differ by an element of $B$, so the framed
points at $x_2$ and $x'_2$ have the same point, and $y_2$ and $y'_2$, which lie over $y'$ with that point, are equal.
It is surjective: given $Z_2$ over $Z$ and $\gamma\in\Gamma_Z$, let $x_2=\delta^{-1}\gamma^{-1}x'$, a frame of $Z_2$. The
framed point of $X$ under the framed point at $x_2$ is triangularly equivalent to the framed point at
$\gamma^{-1}x'$, hence to $y'$; so the point of the framed point at $x_2$ lies over the point of $y'$, and a framed
point $y_2$ of $Y$ with this point over $y'$ is triangularly equivalent to the framed point at $x_2$, and is sent to
the pair of $Z_2$ and $\gamma$. Hence the number of framed points of $Y$ over $y'$ is $\sum_{Z_2\mapsto Z}w(Z_2)$. Fix
a frame of one component of $X$, and let $N$ be the number of framed points of $Y$ over its framed point (if $X$
has no components, put $N=1$). By Lemma~\ref{cv-uniformization-fibre-count}(2), the number of framed points of $Y$
over $\mathbf C$ that lie over a framed point of $X$ over $\mathbf C$ is the same for all framed points of $X$. So
$\sum_{Z_2\mapsto Z}w(Z_2)=N$ for every component $Z$ of $X$, and $N>0$ as every framed point of $X$ over
$\mathbf C$ lifts to $Y$ (Definition~\ref{cv-deep-cover}(4)). In particular every component of $X$ is the image of
a component of $Y$. This proves (3).
% Lean: the first two paragraphs of (1) are exists_carries (Carries, via exists_isometry_of_sabMap,
% UnifPoint.div_eq_iff with ex_eta, and exists_isMatch_of_models; the data of the components are the field
% of IsUnf with HasPeriods and hef, and hG, hxi from the deep cover) and DictMove.move_eq (exp_of_levMove
% with ArcCM.point_exp_eq_of_isometry, TwistTie.perFr_of_unitary), with GeoLevel.exists_frBor_gen and
% FlagTan.evR_frBor for t(x); gamma^* is pull gamma = transl gamma^{-1}, and the action is levR_move.
% The listed compatibility conditions are fields of Formalism.L2.PjFree.IsUnf: the form rule is IsSec of
% laU (with LevelCusp.laU_match), the data at unit frames are its last clause (Hm, chil, fr, theta, Fd).
% The constant is LevelKeeps.exists_const through DictMove.exists_const_of_carries for the weight (0, 2):
% Step 1 is the clause cabs (t none) = 1 of LevelCusp.exists_match, which is
% LevelCusp.cabs_gen_eq_one_of_move (exists_isometry_of_sabMap with h = 1 here, hermForm_eq_of_im,
% MoveValues.frMove_values for omega, UnitFrameGen.frames_eq, ThetaPtChg.val_of_secBC, normSq_periodic,
% mean_comp_eq); Step 2 is the pointwise clause of exists_const_of_carries (exists_unitFr_frBor,
% Frames.algebraic_laws, the IsSec rule with Dict.Pull.cast_bchar); Step 3 is exists_const_of_pointwise;
% Step 4 is exists_sec_two and the end of exists_const.
% "Carries" is DictMove.Carries (IsMatch, equal Gm, and lat: the isometry gam with the congruence through s
% and phi' o gam = delta o phi). (2) is LevelIsometry.petE_transl: the markings are the last clause of
% Formalism.L2.Pj.DictLat (adPin of the deep cover); the paragraph on (2) is StabMove.exists_carries
% (MarkCoset.rational_matrix for delta, TwistTie.marked_isometry_half with the lemma transport of
% MarkedLatticeIsometryDoubleCosetMove for the three terms and MarkCoset.finVec_div,
% T.archimedeanMatrix_unitary_continuous_multiplicative for the frames) and
% LevelIsometry.norm_rdU_transl (exists_const_of_carries with |z| = 1, exists_frBor_gen, exists_sec_two).
% The group paragraph is LevelKeeps.mem_group (Steps A to D) and GammaLevel.trivial_mod_of_mem (Step E):
% Step A is GlobalTheory.congr_lact_mem_orbit (FamRigid.ex_eq_of_level_group with h0 = 1, perFr_lact,
% frBor_of_eq with VerticalDeriv.isDerivAt0_unique, and the fibre field of IsUnf); Step B is
% OrbitScalar.exists_smul_of_orbits; Step C is the rest of unf_group_contains_congruence_upto_rootsOfUnity
% (Data.Gamma_index, isUStab_pow', Data.Gamma_free); Step D is PrincCore.root_eq_one (the unitarity hun is
% unitary_of_first_minor with the last field of MuUnit, dmin_lact and hsur of
% T.adelicUnitary_exists_with_prescribed_parts; then the fields of
% DictIsUnif for x0, archMat and the fibre, T.adelicUnitary_exists_with_prescribed_parts for h, IsLevelK,
% exists_lattice_congr with T.exists_lattice_rationalVectorsInAdelicTranslate for Lambda', and
% root_of_unity_eq_one_of_conj_congruence); Step E is exists_level_block (FrameFibre.exists_level_block_of_eq:
% the fibre of FamilyUnif, LevelHalf.trivial_mod_of_trivial_on_half, exists_gl, herm_of_trace, per_move and
% FramePeriod.frame_of_period_linear) and eq_one_of_block. k_gamma is Lemma cv-uniformization-level-lift.
% (3) is GlobalTheory.LevModel.levCover_model. The setting: the IsUnf and ModelW fields of Laws for X
% (hL.unf, hL.modelW); "the framed point of X under y_2" is FrOver d y_2 y; the four facts after it are
% LevModel.frOver_inj, FrameLift.exists_frOver, CoverSec.frBor_down (or LevModel.exists_frBor_of_frOver)
% and FrameLift.frBor_lift. The components and delta are LevModel.over_comp (the predicate Over):
% the frame x_1 from the cover field of IsUnf for X and NonIntegralFrameExample.exists_framed,
% CoverFrameMod.exists_isometry_of_cover for g_1 and h, LevModel.mark_of_exp for the congruence,
% TraceIsometry.exists_gl_of_trace_isometry for delta (epsilon = 1 as both models have xi0, by ModelW),
% FrameDom.mem_dom_of_isometry, CoverFrameRead.point_over (CoverExp.ex_comp_eq) and perFr_move,
% exists_frOver_down, GeoLevel.exists_frBor_gen for t(x), and CoverFrameMod.norm_gen_eq_one_of_frOver
% for |t(x)| = 1 (the argument of Step 1). The readings are LevRead.reading_pull: read_pointwise
% (evR_pull_genScal) and exists_pow_of_pointwise (PullbackSections.exists_borel_orthonormal_move, laU_eq,
% TB.WickVanishing.flag_domain_frames_connected and eq_pow_mul_of_norm_one). The groups are
% LevModel.conj_arithGroup_le: LevModel.lact_conj_eq_rmul (fibre field of IsUnf for Y and X,
% exists_frBor_of_frOver, the groupoid frBor_groupoid), LevModel.eq_mul_scalar_of_moved (the Lean's own
% Baire argument orbit_open_scalar on the open set of moved frames; the text applies Lemma
% cv-uniformization-scalar-orbits, as the moved frames are the frames of Z), LevScalar.scalar_eq_one_of_conj
% (norm_eq_one_of_herm, the eigenvector identity hint, and eq_one_of_congr through
% NumberField.exists_conjugate_one_le_norm), with (C) LevScalar.isIntegral_eig: through
% T.principalLevel_components_doubleCosets_congruence = LevGeo.Reps.core, the conjugation is
% T.arithmeticGroup_conj_rational_principalCongruence = MsRepsFwd.conj_rat_congr (unitarity from ModelW,
% hsur, the dictionary fields x0, archMat and the fibre, IsLevelK, latOf_congr for Lambda' = latOf g latC),
% then row_eig_conj and isIntegral_eig_of_congr (Step 1 of root_of_unity_eq_one_of_conj_congruence); and the
% index ArithIndex.relIndex_map_conj_ne_zero
% (LatticeCongr.exists_relIndex_ne_zero_stab, the kernel on Lambda/n^2 Lambda, and Data.Gamma_index).
% The indices are LevModel.cover_degree_const: LevModel.lifts_equiv through nonempty_equiv_sigma_quotient
% (with eq_rmul_lact, eq_of_rmul_lact_eq, moved_iff and mem_dom_of_moved) and LevDegree.card_lifts_const
% (Lemma cv-uniformization-fibre-count). The section of Step 4 is LevelKeeps.exists_sec_two, that is
% DictLevel.exists_sec_rdU_ne_zero with n = 2 at the unit frame (Lemma cv-uniformization-theta-section, whose open
% inputs are in Lemmas cv-uniformization-theta-marking, -theta-half and -theta-rule); Lemma
% cv-volume-twisting-nonvanishing gives weight (0, 2m + 2) only.
\end{proof}

\begin{lemma}\label{cv-cover-level-automorphisms}
In Situation~\ref{cv-cusp-setting}, suppose that $q\ge1$ and $f\ge3$, and let $Y$ be a deep cover of $X$
(Definition~\ref{cv-deep-cover}), with deep level $(M_Y,\mathfrak M_Y)$ and coefficient field $F_Y$.
Let $\Lambda$ be the lattice of the PEL datum of $X$, a projective $\mathcal O_E$-module of rank $r$, and
let $\Gamma_Y$ be the group of $\mathcal O_E$-linear automorphisms of $\Lambda/\mathfrak M_Y\Lambda$. An
element $\gamma\in\Gamma_Y$ acts on $Y$ over $F_Y$ by changing the level structure of each point by
$\gamma$, and $\gamma^*$ denotes the pullback along this action on the $F_Y$-algebra $R$ of rational
sections of all weights on $Y$. A \emph{weight} is a pair $\kappa=(\alpha,n)$ of a family of integers
$\alpha_{\tau,i}$, indexed as for the dominant weights of Situation~\ref{cv-cusp-setting}, and an even
integer $n$. For a weight $\kappa$ let $S_Y(\kappa)\subset R$ be the $F_Y$-space of rational sections of
weight $\kappa$ on $Y$ whose readings at the good embeddings are cusp forms, with its integral forms and
its Petersson norms $N_\sigma$, defined as in Situation~\ref{cv-cusp-setting} for $Y$.
\begin{enumerate}
\item $\Gamma_Y$ is finite, and $|\Gamma_Y|\le M_Y^{4f^2r^2}$.
\item Each $\gamma^*$ is an $F_Y$-algebra automorphism of $R$ that preserves weights, and
  $(\gamma\gamma')^*=\gamma'^*\gamma^*$ for $\gamma,\gamma'\in\Gamma_Y$.
\item Let $\kappa$ be a weight and $\gamma\in\Gamma_Y$. Then $\gamma^*$ maps $S_Y(\kappa)$ to itself and
  integral forms to integral forms, and $N_\sigma(\gamma^*f)=N_\sigma(f)$ for every $f\in S_Y(\kappa)$ and
  every embedding $\sigma\colon F_Y\to\mathbf C$.
\item Let $\kappa$ and $\kappa'$ be weights. The product of an element of $S_Y(\kappa)$ and an element of
  $S_Y(\kappa')$ lies in $S_Y(\kappa+\kappa')$, and it is integral if both factors are.
\end{enumerate}
\end{lemma}
\leanref{Development/AbelianVarietyTheory/Proofs/DivisorProofs/CardLeOfDetEqZero.lean}{GlobalTheory.ShortDet.card_levelAut_le}
\leanref{Development/AbelianVarietyTheory/Definitions/Divisors/MemLevEigenIffTransl.lean}{GlobalTheory.LevDesc.isRightAct_pull}
\leanref{Development/AbelianVarietyTheory/Proofs/DivisorProofs/KeepsCuspLatticeNorm.lean}{GlobalTheory.LevelIsometry.keeps_cusp_lattice_norm}
\leanref{Development/AbelianVarietyTheory/Definitions/Foundations/ExistsIsDeepOfField.lean}{GlobalTheory.SlopeRec.mul_mem_cuspW}
\begin{proof}
(1) The lattice $\Lambda$ is free of rank $2fr$ over $\mathbf Z$. Since
$M_Y\mathcal O_E=\mathfrak M_Y\bar{\mathfrak M}_Y\subset\mathfrak M_Y$, the group $\Lambda/\mathfrak M_Y\Lambda$
is a quotient of $\Lambda/M_Y\Lambda$, so it has at most $M_Y^{2fr}$ elements. An $\mathcal O_E$-linear
automorphism of $\Lambda/\mathfrak M_Y\Lambda$ is determined by the images of the classes of a
$\mathbf Z$-basis of $\Lambda$, which gives $|\Gamma_Y|\le(M_Y^{2fr})^{2fr}$.

(2) The action of $\gamma$ changes only the level structures of the points of $Y$, and keeps their
abelian varieties, the bases of the eigenspaces on their Lie algebras and the generators of sections.
So the value of $\gamma^*f$ at a framed point is the value of $f$ at a framed point with the same frame,
and $\gamma^*f$ has the weight of $f$. Pullback along an action is an algebra map, and
$(\gamma\gamma')^*=\gamma'^*\gamma^*$ because $\gamma\mapsto\gamma^*$ is the pullback along a left action.

(3) Integral forms. An integral point of $Y$ over a complete discrete valuation ring, with its level
structure changed by $\gamma$, is again an integral point, and the value of $\gamma^*f$ at the first is the
value of $f$ at the second. So the valuations of $\gamma^*f$ at integral points are valuations of $f$ at
integral points, and $\gamma^*f$ is integral if $f$ is.

Cusp forms. Let $\sigma$ be a good embedding and $Z$ a component of the uniformization at $\sigma$, and let
$Z'$, $\delta$ and $c$ be as in Lemma~\ref{cv-cover-uniformization}(1), so that
$(\gamma^*f)_Z(x)=c\,f_{Z'}(\delta x)$ for all frames $x$ of $Z$. For $\gamma'\in\Gamma_Z$, the transformation
$\delta\gamma'\delta^{-1}$ is unitary, preserves the lattice of $Z'$ and is congruent to $1$ modulo $M_Y$ on
it, since $\gamma'$ has these properties for $Z$; so $\delta\Gamma_Z\delta^{-1}\subset\Gamma_{Z'}$ by
Lemma~\ref{cv-cover-uniformization}(1). Hence the reading of $\gamma^*f$ on $Z$ is holomorphic, has the
transformation law of the weight under $B$, is invariant under $\Gamma_Z$ because the characters of weights
with even $n$ are trivial on the arithmetic groups (Definition~\ref{cv-deep-cover}(5)), has bounded
pointwise norm, and has vanishing constant Fourier--Jacobi coefficients, since $f_{Z'}$ has these
properties. So $\gamma^*f\in S_Y(\kappa)$.

Norms. Let $k_\gamma$ be as in Lemma~\ref{cv-cover-uniformization}(2). Since $|\delta_{\tau,i}(x(g))|=1$ for
all $g$, the pointwise norm of the reading of $\gamma^*f$ at $x(g)$ is that of the reading of $f$ at
$x(gk_\gamma)$. The Haar measure is invariant under right translation by $k_\gamma$, so the integrals that
define $N_\sigma(\gamma^*f)$ and $N_\sigma(f)$ are equal.

(4) Weights add under products, since the value of a product at a framed point is the product of the
values. The readings of a product are the products of the readings; a product of holomorphic functions
with the transformation laws of $\lambda$ and $\lambda'$ is holomorphic with the law of $\lambda+\lambda'$,
and its pointwise norm is the product of the pointwise norms, so it is bounded. Along an isotropic vector
the Fourier--Jacobi coefficients of a product are the convolutions of those of the factors, and those of a
form with bounded pointwise norm vanish outside the positive semidefinite indices; so the constant
coefficient of the product is the product of the constant coefficients, which vanishes. So the product is a
cusp form. The
valuation of a product at an integral point is the sum of the valuations of the factors, which gives
the last claim.
\end{proof}
% Lean: (1) card_levelAut_le; (2) LevDesc.transl and isRightAct_pull, with pull h = transl h^{-1} and
% move (the level structure changed through levActR (levFull Dw l)); (3) keeps_cusp_lattice_norm:
% integrality GlobalTheory.LevelIsometry.val_transl_nonneg (integral_of_levR), cusp forms
% GlobalTheory.LevelKeeps.keepsCusp_of_isUnf (rdU_transl_mem with exists_carries, exists_const, conj_mem,
% and the arithmetic groups from GlobalTheory.LevelKeeps.mem_group, levelK and DictIsUnif), norms
% GlobalTheory.LevelIsometry.petE_transl (exists_induces, eLpNorm_right); (4) mul_mem_cuspW (mul_mem_secZ,
% holSub_mul with siegelSpaceAlgebra_holds = BoundedSectionAlgebra.siegelSpaceAlgebra_of from
% moderateCoefficientPositivity_holds and coefficientConvolution_holds) and the additivity of valuations
% (AllWt.mul_mem_supLat, IsVal.mul).

\begin{lemma}\label{cv-cover-pullback}
In Situation~\ref{cv-cusp-setting}, suppose that the level of $X$ is a level $(M,\mathfrak M)$ with $M$
prime to $15$, and let $Y$ be a deep cover of $X$ (Definition~\ref{cv-deep-cover}), with coefficient field
$F_Y$, level map $\pi_Y\colon Y\to X$, and $S_Y(\kappa)$ and $N_\sigma$ as in
Lemma~\ref{cv-cover-level-automorphisms}. Let $\kappa$ be a dominant weight, with $S_\kappa$, its norms
$\|\cdot\|_\sigma$ and its integral forms as in Situation~\ref{cv-cusp-setting}. A rational section $g$ of
weight $\kappa$ on $Y$ is a \emph{pullback} of a rational section $f$ of weight $\kappa$ on $X$ if
$g(y_2)=f(y)$ whenever $y_2$ is a framed point of $Y$ over an algebraically closed field containing $F_Y$
and $y$ is the framed point of $X$ under it, with the frame carried over by the identification of
abelian varieties of Definition~\ref{cv-deep-cover}(4). A rational section of $Y$ is \emph{constant on
fibres} if it takes the same value at any two such framed points of $Y$ over the same framed point of $X$.
\begin{enumerate}
\item $F_Y$ is a Galois extension of $F_c$.
\item Every rational section $f$ of weight $\kappa$ on $X$ has a unique pullback $\pi(f)$, and $\pi(f)=0$
  only if $f=0$. The section $\pi(f)$ is a cusp form if and only if $f$ is, and it is integral if and only
  if $f$ is.
\item The map $\pi\colon S_\kappa\to S_Y(\kappa)$ is additive, $\pi(af)=a\,\pi(f)$ for $a\in F_c$, it
  carries $F_c$-linearly independent families to $F_Y$-linearly independent families, and
  $N_{\sigma'}(\pi(f))=\|f\|_{\sigma'|F_c}$ for every $f\in S_\kappa$ and every embedding
  $\sigma'\colon F_Y\to\mathbf C$.
\item A form $g\in S_Y(\kappa)$ lies in the $F_Y$-span of $\pi(S_\kappa)$ if and only if it is constant on
  fibres.
\item Let $b_1,\dots,b_m\in S_\kappa$, $x\in F_Y^m$ and $s\in\operatorname{Gal}(F_Y/F_c)$. If
  $\sum_ix_i\pi(b_i)$ is integral, then so is $\sum_is(x_i)\pi(b_i)$.
\end{enumerate}
\end{lemma}
\leanref{Development/AbelianVarietyTheory/Proofs/DivisorProofs/LevelCoverCuspFormPullbackInjectiveDescentStatement.lean}{GlobalTheory.T.levelCover_cuspForm_pullback_injective_descent}
\leanref{Development/AbelianVarietyTheory/Proofs/DivisorProofs/LevelCoverCuspFormPullbackInjectiveDescent.lean}{GlobalTheory.LevRingDesc.levRing_descent}
\leanref{Development/AbelianVarietyTheory/Proofs/DivisorProofs/IntValuationPrimeBack.lean}{GlobalTheory.GalSup.intAt_galois}
\begin{proof}
(1) By Definition~\ref{cv-deep-cover}(2), every automorphism of $\mathbf C$ that fixes $F_c$ maps $F_Y$ into
itself. So $F_Y$ is a normal extension of $F_c$, and it is separable since the characteristic is $0$.

(2) Existence. For a framed point $y_2$ of $Y$, let $\pi(f)(y_2)$ be the value of $f$ at the framed point of
$X$ under $y_2$. The identification of Definition~\ref{cv-deep-cover}(4) is compatible with the actions of
$\mathcal O_E$ and the polarizations, so it carries the frames of $y_2$ to frames of the point below and
changes of frame to the same changes of frame; hence $\pi(f)$ has the transformation laws of weight
$\kappa$, and it is a pullback of $f$.

Uniqueness. Every framed point of $Y$ over an algebraically closed field containing $F_Y$ lies over a
framed point of $X$, and a rational section of $Y$ is determined by its values at these points. So any
two pullbacks of $f$ agree. If $\pi(f)=0$, then $f$ vanishes at every framed point of $X$ over an
algebraically closed field containing $F_Y$, since by Definition~\ref{cv-deep-cover}(4) each of them lifts
to a framed point of $Y$; so $f=0$.

Cusp forms. Let $\sigma'$ be a good embedding of $F_Y$; its restriction $\sigma$ to $F_c$ is good. By
Lemma~\ref{cv-cover-uniformization}(3), every component $Z$ of $X$ at $\sigma$ is the image of a component
$Z_2$ of $Y$ at $\sigma'$, and the reading of $\pi(f)$ on $Z_2$ is $x\mapsto t^nf_Z(\delta x)$ with $|t|=1$,
for a unitary transformation $\delta$ of the common Hermitian matrix with $\delta\Gamma_{Z_2}\delta^{-1}$ of
finite index in $\Gamma_Z$. Holomorphy, the transformation laws, boundedness of the pointwise norm and the
vanishing of the constant Fourier--Jacobi coefficients pass along such a transformation in both
directions, so the readings of $\pi(f)$ are cusp forms if and only if those of $f$ are.

Integral forms. Let $y$ be an integral point of $X$ over a complete discrete valuation ring $T$ above a
finite place of $F_c$. After a finite extension $T'$ of $T$, $y$ lifts to an integral point $y_2$ of $Y$,
and $\pi(f)(y_2)=f(y)$ in $T'$. The valuation of $T'$ restricts to a positive multiple of that of $T$, so
if $\pi(f)$ is integral the valuation of $f(y)$ is nonnegative. Conversely, every integral point of $Y$
over a complete discrete valuation ring lies over an integral point of $X$ over the same ring, at which
$f$ has the value of $\pi(f)$. Finally, whether $f$ is integral does not change when its coefficient field
is extended from $F_c$ to $F_Y$, since the integral points over complete discrete valuation rings above
places of $F_Y$ restrict to integral points above places of $F_c$ and conversely after a finite extension.

(3) Additivity and $\pi(af)=a\,\pi(f)$ follow from the uniqueness in (2), since sums and multiples of
pullbacks are pullbacks. Let $b_1,\dots,b_s\in S_\kappa$ be linearly independent over $F_c$, and let
$t\in F_Y^s$ with $\sum_kt_k\pi(b_k)=0$. Then $\sum_kt_kb_k$, a section of $X$ with coefficients in $F_Y$,
has pullback $0$, so it vanishes by (2), applied over $F_Y$. For every $F_c$-linear map
$\lambda\colon F_Y\to F_c$, evaluating $\sum_kt_kb_k=0$ at the framed points of $X$, with values in tensor
products with $F_Y$ over $F_c$, and applying $\lambda$ to the coefficients gives $\sum_k\lambda(t_k)b_k=0$.
So $\lambda(t_k)=0$ for every $\lambda$ and every $k$, and $t=0$.
For the norms, it suffices to treat a good $\sigma'$, since both sides are computed at the good one of
$\sigma'$ and $\bar\sigma'$; let $\sigma=\sigma'|F_c$. By Step~1 of the proof of
Lemma~\ref{cv-cusp-sup-bound}, the square of a Petersson norm is the sum over the components of the
integrals of the squared pointwise norm over fundamental sets, divided by the sum of the volumes of these
sets. Let $Z_2\mapsto Z$, $N$, $\delta$, $t$ and $w(Z_2)$ be as in Lemma~\ref{cv-cover-uniformization}(3). The
pointwise norm of the reading of $\pi(f)$ on $Z_2$ is that of $f$ on $Z$, transported by $\delta$, so the
integral and the volume for $Z_2$ are $w(Z_2)$ times those for $Z$. The indices $w(Z_2)$ of the components
above each $Z$ add up to $N$, so the numerator and the denominator for $\pi(f)$ are $N$ times those for $f$,
and $N_{\sigma'}(\pi(f))=\|f\|_\sigma$; the factors $(2\pi\xi)^{-|\alpha|/2}$ agree.

(4) Pullbacks are constant on fibres, and so are their $F_Y$-linear combinations. Conversely, let
$g\in S_Y(\kappa)$ be constant on fibres. Let $f'$ assign to a framed point of $X$ the common value of $g$
at the framed points of $Y$ above it, after extending to an algebraically closed field; such points exist
by Definition~\ref{cv-deep-cover}(4), and the value does not depend on the choice since $g$ is constant on
fibres. Lifting changes of frame to $Y$ shows that $f'$ has the transformation laws of weight $\kappa$, so
$f'$ is a rational section of weight $\kappa$ on $X$ with coefficients in $F_Y$, and $g$ is its pullback. By
(2) over $F_Y$, $f'$ is a cusp form. The cusp forms on $X$ with coefficients in $F_Y$ are the $F_Y$-linear
combinations of cusp forms over $F_c$, by flat base change from $F_c$ to $F_Y$ for the sections of the
automorphic line bundle. So $f'=\sum_ka_kf_k$ with $a_k\in F_Y$ and $f_k\in S_\kappa$, and
$g=\sum_ka_k\pi(f_k)$ by uniqueness.

(5) By (2) it suffices to show that $\sum_is(x_i)b_i$, a section of $X$ with coefficients in $F_Y$, is
integral when $\sum_ix_ib_i$ is. Let $y$ be an integral point of $X$ over a complete discrete valuation ring
$T$, with $F_Y$ embedded in the fraction field of $T$ by $\iota$. The same point with $F_Y$ embedded by
$\iota\circ s$ is again an integral point, above another place of $F_Y$. The sections $b_i$ have
coefficients in $F_c$, which $s$ fixes, so the value of $\sum_is(x_i)b_i$ at $y$ read through $\iota$ is
$\sum_i\iota(s(x_i))b_i(y)$, the value of $\sum_ix_ib_i$ at $y$ read through $\iota\circ s$. Its valuation is
nonnegative.
\end{proof}
% Lean: (1) GlobalTheory.T.deepLevelCover_principalLevel_laws_galois (from the field normal of Deep.Good),
% with the laws of X from GlobalTheory.T.principalLevelLiftModuliLaws_of_core. (2)
% GlobalTheory.CoverSec.section_pullback_existsUnique (ruleUp, exists_frOver_below, eq_of_geometric, and the
% lifting field of IsLev), GlobalTheory.T.cuspForm_iff_levelCover_pullback (LevCusp.cusp_iff_pull, through
% levelCover_complexModels_components_isometries_degree and ss_transfer),
% GlobalTheory.CoverInt.pull_intAt_iff (CoverPoint.exists_ext_lift_to_cover, CoverDown.exists_intFr_below)
% and GlobalTheory.IntCoefExt.intAt_rulBC_iff. (3) levRing_descent: exists_pullMap,
% GlobalTheory.CuspFieldChange.linearIndependent_rulBC (evaluation at points and coordinate functionals of
% the tensor product), GlobalTheory.T.peterssonNorm_eq_of_levelCover_pullback
% (LevPullNorm.levPull_normW). (4) GlobalTheory.T.section_constant_on_fibres_descends
% (CoverDesc.section_descends: ruleDown, val_ruleDown, ruleDown_mem) and
% GlobalTheory.CuspFieldChange.cusp_unfBC_eq_span (SecFieldChange.cuspBC_holds from the base change law
% secBaseChange, and rulBC_cusp). The norms follow
% GlobalTheory.LevPullNorm.petE_pull_good: GlobalTheory.T.levelCover_complexModels_components_isometries_degree
% (the map of components, delta, the indices and the degree N), petE_sq_mul_vol, and
% GlobalTheory.T.lintegral_fundamentalSet_conj_finiteIndex_subgroup with sum_mul_fibre. (5) intAt_galois
% (changeEmbedding, primeBack). The uniqueness of pullbacks of 0 is also
% GlobalTheory.T.existsUnique_cuspForm_pullback_goodCover, and the last field of levRing_descent,
% span_integral_descent, is the descent of Lemma cv-lattices-galois-different used in the proof of (6) of
% Proposition cv-arithmetic-volume-step.

\begin{lemma}\label{cv-cover-type}
In the situation of Lemma~\ref{cv-cover-pullback}, let $\Lambda$ and $\Gamma_Y$ be as in
Lemma~\ref{cv-cover-level-automorphisms}; an element of $\Gamma_Y$ is \emph{congruent} to an
$\mathcal O_E$-linear endomorphism $\ell$ of $\Lambda$ modulo $\mathfrak M$ if its lifts to
$\mathcal O_E$-linear endomorphisms of $\Lambda$ agree with $\ell$ modulo $\mathfrak M\Lambda$, which
makes sense since $\mathfrak M_Y\subset\mathfrak M$. Let $G$ be a group and $u\mapsto\ell_u$ a map from $G$ to
the $\mathcal O_E$-linear endomorphisms of $\Lambda$. Write $\gamma\cdot g=(\gamma^{-1})^*g$ for the action
of $\Gamma_Y$ on the rational sections of $Y$, and suppose that $G$ acts on the rational sections of $X$,
written $u\cdot f$, in the same way through the change of level structures by $\ell_u$, and that every
$u\in G$ maps $S_\kappa$ to itself. Let $P\subset G$ be a finite subgroup, let $\lambda\colon P\to F_c^\times$ be a homomorphism, and let
$\theta\colon\Gamma_Y\to G$ be a homomorphism such that:
\begin{enumerate}
\item[(a)] every $\gamma\in\Gamma_Y$ is congruent to $\ell_{\theta(\gamma)}$ modulo $\mathfrak M$;
\item[(b)] $\theta(\gamma)=1$ for every $\gamma\in\Gamma_Y$ that is congruent to the identity modulo
  $\mathfrak M$;
\item[(c)] every element of $P$ is $\theta(\gamma)$ for some $\gamma\in\Gamma_Y$.
\end{enumerate}
Put $H=\theta^{-1}(P)$. Then:
\begin{enumerate}
\item $\gamma\cdot\pi(f)=\pi(\theta(\gamma)\cdot f)$ for every $\gamma\in\Gamma_Y$ and every $f\in S_\kappa$.
\item A form $g\in S_Y(\kappa)$ is constant on fibres if and only if $\gamma\cdot g=g$ for every
  $\gamma\in\Gamma_Y$ that is congruent to the identity modulo $\mathfrak M$.
\item A form $g\in S_Y(\kappa)$ lies in the $F_Y$-span of the set of $\pi(f)$, for $f\in S_\kappa$ with
  $u\cdot f=\lambda(u)f$ for all $u\in P$, if and only if $\gamma\cdot g=\lambda(\theta(\gamma))g$ for every
  $\gamma\in H$.
\end{enumerate}
\end{lemma}
\leanref{Development/AbelianVarietyTheory/Proofs/DivisorProofs/LatticeLevelCharacterBound.lean}{GlobalTheory.CovTop.TypeLattice.mem_span_pull_iff}
\leanref{Development/AbelianVarietyTheory/Proofs/DivisorProofs/LatticeLevelCharacterBound.lean}{GlobalTheory.CovTop.TypeLattice.mem_span_image_eigen}
\leanref{Development/AbelianVarietyTheory/Proofs/DivisorProofs/FibreConstantIffCoveringGroupInvariantStatement.lean}{GlobalTheory.T.fibreConstant_iff_coveringGroup_invariant}
\begin{proof}
(1) Let $y_2$ be a framed point of $Y$ over a framed point $y$ of $X$, and let $y_2'$ and $y'$ be these
points with their level structures changed by $\gamma$ and by $\ell_{\theta(\gamma)}$. Conjugating the level
map of Definition~\ref{cv-deep-cover}(4) by these changes gives a morphism with the same properties, since
$\gamma$ is congruent to $\ell_{\theta(\gamma)}$ modulo $\mathfrak M$ by (a), the level structure of $X$
only sees $\Lambda$ modulo $\mathfrak M\Lambda$, and $15^e$ times the changed level structure of $Y$ is the
changed level structure of $X$; under it $y_2'$ lies over $y'$. It differs from the level map by a change of
level structure on $Y$ by an element that is the identity modulo $\mathfrak M$, which does not change the
point of $X$ below, so $y_2'$ lies over $y'$. Hence $\gamma\cdot\pi(f)$ is a pullback of
$\theta(\gamma)\cdot f$, and the uniqueness in Lemma~\ref{cv-cover-pullback}(2) gives the claim.

(2) Two framed points of $Y$ over the same framed point $y$ of $X$ have the same abelian variety, frame and
identification with that of $y$, and level structures that induce the level structure of $y$; so they
differ by a change of level structure by an element of $\Gamma_Y$ congruent to the identity modulo
$\mathfrak M$. Conversely, changing the level structure of a point of $Y$ by such an element does not
change the point of $X$ below it. So the fibres are the orbits of the group $K$ of these elements, and $g$
is constant on fibres if and only if it is invariant under $K$.

(3) Suppose first that $f\in S_\kappa$ satisfies $u\cdot f=\lambda(u)f$ for $u\in P$. For $\gamma\in H$,
(1) gives $\gamma\cdot\pi(f)=\pi(\theta(\gamma)\cdot f)=\pi(\lambda(\theta(\gamma))f)=\lambda(\theta(\gamma))\pi(f)$,
by Lemma~\ref{cv-cover-pullback}(3). The condition is linear in $g$, so it holds on the $F_Y$-span.

Conversely, let $g$ satisfy the condition. For $\gamma\in\Gamma_Y$ congruent to the identity modulo
$\mathfrak M$, $\theta(\gamma)=1$ by (b), so $\gamma\in H$ and $\lambda(\theta(\gamma))=1$, and
$\gamma\cdot g=g$; by (2) and Lemma~\ref{cv-cover-pullback}(4),
$g=\sum_ka_k\pi(f_k)$ with $a_k\in F_Y$ and $f_k\in S_\kappa$. For $u\in P$ choose $\gamma_u\in\Gamma_Y$
with $\theta(\gamma_u)=u$, by (c), and let $T=\sum_{u\in P}\lambda(u)^{-1}\gamma_u\cdot$, an $F_Y$-linear map on
the rational sections of $Y$. Then $Tg=|P|g$. For $f\in S_\kappa$, (1) gives $T\pi(f)=\pi(ef)$ with
$ef=\sum_{u\in P}\lambda(u)^{-1}u\cdot f\in S_\kappa$, and for $u'\in P$,
\[
 u'\cdot ef=\sum_{u\in P}\lambda(u)^{-1}(u'u)\cdot f=\lambda(u')\sum_{v\in P}\lambda(v)^{-1}v\cdot f=\lambda(u')\,ef.
\]
So $g=|P|^{-1}Tg=|P|^{-1}\sum_ka_k\pi(ef_k)$ lies in the span.
\end{proof}
% Lean: mem_span_pull_iff, for a group G with lev : G -> End(Lambda) and ActLev, hkeep (every element
% keeps cuspA), P finite, lam, and theta with the hypotheses hθ (a), hker (b) and hs (c). (1) is the
% step hc, through GlobalTheory.LevPullEquiv.pull_actLev and the base change
% GlobalTheory.FieldBC.actLev_and_actBC; (2) is fibreConstant_iff_coveringGroup_invariant
% (LevInv.fibConst_iff_invariant), with K the elements whose lift is congruent to the identity
% modulo Mb; (3) uses levelCover_cuspForm_pullback_injective_descent (the field p5, Lemma cv-cover-pullback(4))
% and mem_span_image_eigen, whose operator T is the text's T.

\subsection{Generators and twisting sections}\label{subsection-cv-twisting-sections}

\begin{definition}\label{cv-volume-forms-weights}
In Situation~\ref{cv-cusp-setting}, suppose that $p\ge1$ and $q\ge1$. Let $\lambda=(\lambda_{\tau,i})$ be a
family of integers, for $\tau\in\Phi$ and $0\le i<r$, and let $m$ be an integer. The weight
$\kappa(\lambda,m)$ is the pair of the family of integers whose entry at $\tau\in\Phi$ in position $i<p$ is
$\lambda_{\tau,i}-m$ and whose entry at $\bar\tau$ in position $i<q$ is $-\lambda_{\tau,r-1-i}-m$, and of the
even integer $2m$; thus $\kappa(\lambda,m)+\kappa(\lambda',m')=\kappa(\lambda+\lambda',m+m')$. The family
$\lambda$ is \emph{regular} if $\lambda_{\tau,0}\ge\lambda_{\tau,1}\ge\dots\ge\lambda_{\tau,r-1}$ and
$\lambda_{\tau,p-1}\ge\lambda_{\tau,p}+1$ for every $\tau\in\Phi$. Let $\mathbf h$ be the family with
$\mathbf h_{\tau,i}=0$ for $i<p$ and $\mathbf h_{\tau,i}=-1$ for $i\ge p$; the \emph{Hodge weight} is
$h=\kappa(\mathbf h,0)$.
\end{definition}
\leanref{Development/AbelianVarietyTheory/Definitions/Foundations/IsProbabilityMeasureMuPart.lean}{GlobalTheory.AllWt.wtOf}
\leanref{Development/AbelianVarietyTheory/Definitions/Divisors/BarCompSubsetClosure.lean}{GlobalTheory.CompHS.gap}
\leanref{Development/AbelianVarietyTheory/Definitions/Divisors/IsHomogeneousCompIdeal.lean}{GlobalTheory.CompHt.hodgeWt}

\begin{lemma}[Harish-Chandra \cite{HC68}]\label{cv-volume-forms-finite}
In Situation~\ref{cv-cusp-setting}, let $\sigma$ be a good embedding of $F_c$, and let $\lambda_Z$ be a weight
on each component $Z$ at $\sigma$. Then the families $(P_Z)_Z$ of cusp forms $P_Z$ of weight $\lambda_Z$ on
$Z$ form a finite-dimensional complex vector space.
\end{lemma}
\leanref{Development/AbelianVarietyTheory/Proofs/UnitaryExistence/FiniteDimensionalCuspSectionsVanishingOffDomainStatement.lean}{GlobalTheory.T.finiteDimensional_cuspSections_vanishing_offDomain}

\begin{lemma}\label{cv-volume-generators-functionals}
Let $K$ be a number field, let $V$ be a finite-dimensional $K$-vector space, and let $L\subset V$ be an
$\mathcal O_K$-submodule. Suppose that for every nonzero $x\in V$ there are a number field $K'$ containing
$K$ and a $K$-linear map $\varphi\colon V\to K'$ with $\varphi(x)\ne0$ and $\varphi(L)\subset\mathcal O_{K'}$.
Then $L$ is a finitely generated $\mathcal O_K$-module.
\end{lemma}
\leanref{Development/AbelianVarietyTheory/Proofs/DivisorProofs/FgIntOfForallIsIntegral.lean}{GlobalTheory.fg_of_integral_functionals}
\begin{proof}
It suffices to show that $L\cap W$ is a finitely generated abelian group for every $K$-subspace $W\subset V$;
the case $W=V$ is the claim. The proof is by induction on $\dim_KW$. If $W=0$ there is nothing to show.
Otherwise let $x\in W$ be nonzero, and let $K'$ and $\varphi$ be as in the hypothesis. The subspace
$W'=W\cap\ker\varphi$ does not contain $x$, so $\dim_KW'<\dim_KW$, and $L\cap W'$ is finitely generated by
induction. The image $\varphi(L\cap W)$ is a subgroup of $\mathcal O_{K'}$, a free abelian group of finite
rank, so it is finitely generated. The kernel of $\varphi$ on $L\cap W$ is $L\cap W'$. An abelian group with
a finitely generated subgroup and a finitely generated quotient by it is finitely generated, so $L\cap W$
is finitely generated.
\end{proof}

\begin{lemma}\label{cv-volume-generators-power}
Let $K$ be a number field and $R$ a commutative $K$-algebra. For every nonzero prime $\mathfrak p$ of
$\mathcal O_K$ let $v_{\mathfrak p}$ be its valuation on $K$, and let $\mathcal V_{\mathfrak p}$ be a set of
maps $v\colon R\to\mathbf R\cup\{+\infty\}$ with $v(0)=+\infty$, $v(f+g)\ge\min(v(f),v(g))$,
$v(fg)=v(f)+v(g)$ and $v(af)=v_{\mathfrak p}(a)+v(f)$ for $f,g\in R$ and $a\in K^\times$. Call $f\in R$
\emph{integral} if $v(f)\ge0$ for every $\mathfrak p$ and every $v\in\mathcal V_{\mathfrak p}$. Let $S\subset R$
be a set each of whose elements $g$ has an integral multiple $ag$ with $a\in K^\times$, and let $f\in R$
and $k\ge0$ be such that $f^{k+1}$ lies in the $K$-span of $S$. Then $f$ has an integral multiple $bf$ with
$b\in K^\times$.
\end{lemma}
\leanref{Development/AbelianVarietyTheory/Proofs/DivisorProofs/ExistsSmulNonnegOfPow.lean}{GlobalTheory.SupFull.exists_smul_nonneg_of_pow}
\begin{proof}
The integral elements form an $\mathcal O_K$-submodule $L$ of $R$, since $v(0)=+\infty$, every $v$ is
ultrametric, and $v(af)=v_{\mathfrak p}(a)+v(f)\ge v(f)$ for nonzero $a\in\mathcal O_K$. Every element of $S$
lies in the $K$-span of $L$, so $f^{k+1}=\sum_ic_i\ell_i$ with $c_i\in K$ and $\ell_i\in L$, and a common
denominator $t\in\mathcal O_K\smallsetminus\{0\}$ of the $c_i$ gives $tf^{k+1}\in L$. Let $\mathfrak p$ be a
prime and $v\in\mathcal V_{\mathfrak p}$. If $v(f)\ge0$ or $v(f)=+\infty$, then
$v(tf)=v_{\mathfrak p}(t)+v(f)\ge0$, as $v_{\mathfrak p}(t)\ge0$. Otherwise $v(f)<0$, and
$v_{\mathfrak p}(t)+(k+1)v(f)=v(tf^{k+1})\ge0$ gives $v(tf)=v_{\mathfrak p}(t)+v(f)\ge-k\,v(f)\ge0$. So $tf$ is
integral, and $b=t$ will do.
\end{proof}

\begin{lemma}\label{cv-volume-forms-cm-frame}
In Situation~\ref{cv-cusp-setting}, let $\sigma$ be a good embedding and $Z$ a component at $\sigma$, and let
$L'\subset E^r$ be a finitely generated $\mathcal O_E$-submodule that spans $E^r$ and is self-dual for $G_Z$, as
in Lemma~\ref{cv-cusp-frame}. Let $\lambda$ be a weight on $Z$, and let $P$ be a form of weight $\lambda$ on $Z$
with $|P|$ bounded and $P(x_1)\ne0$ for some frame $x_1$ of $Z$. Then there are $\ell>0$, a frame $x$ of $Z$
with $P(x)\ne0$, and a basis $v_1,\dots,v_r$ of $E^r$ in $L'$ with properties (1) to (4) of
Lemma~\ref{cv-cusp-frame}.
\end{lemma}
\leanref{Development/AbelianVarietyTheory/Proofs/DivisorProofs/ExistsGoodCmPtFrPolyNeZero.lean}{GlobalTheory.CMDetect.exists_cm_frame}
\begin{proof}
Near a maximum. The pointwise norm $|P|(x_1)$ is $|P(x_1)|$ times a positive number, so it is positive. Let
$s_0$ be the supremum of $|P|$ over the frames of $Z$, which is finite and positive, and choose a frame $z_0$
of $Z$ with $|P|(z_0)>s_0/2$. Then $|P|(y)\le2|P|(z_0)$ for every frame $y$ of $Z$, and $|P|(z_0)>0$.

No short vectors. In the notation of Definition~\ref{cv-decay-majorant},
$v\mapsto\sum_\tau h_{z_0,\tau}(\tau(v))$ extends to a positive definite quadratic form on
$E^r\otimes_{\mathbf Q}\mathbf R$, in which $L'$ is a lattice. So its minimum $m$ over the nonzero vectors of
$L'$ is positive. Put $s=\min(1,m)^{1/2}$. Then $0<s\le1$, and $\sum_\tau h_{z_0,\tau}(\tau(v))\ge s^2$ for every
nonzero $v\in L'$.

The frame. Lemma~\ref{cv-cusp-frame}, with $w=\max(1,w')$ for the spread $w'$ of $\lambda$, the frame $z_0$
and $\lambda_*=s$, gives its number $\ell$, a frame $x$ of $Z$ with $P(x)\ne0$, and the basis.
\end{proof}
\leanref{Development/AbelianVarietyTheory/Proofs/DivisorProofs/ExistsGoodCmPtFrPolyNeZero.lean}{GlobalTheory.CMDetect.exists_maj_lower}
% Lean: exists_cm_frame: exists_half_max, then exists_maj_lower, which takes the minimum over the lattice Lambda
% of the data X (with N_C L' inside Lambda and the weight k_0 of the majorant, through
% AR.Up.ModeHalfC.first_minimum_pos), and Formalism.L2.Pj.FundamentalDomainD5.t41H_statement (Lemma
% cv-cusp-frame) with w = max 1 (sum of the spreads); the basis is FD.LamBasis X La x lam.

\begin{lemma}\label{cv-volume-forms-triangular}
Let $K$ be an infinite field, $B$ a finite set and $m_b\ge0$ integers for $b\in B$. For a polynomial $P$ over
$K$ in the entries of a family $X=(X_b)_{b\in B}$ of square matrices of sizes $m_b$, write $P(X)$ for its
value at such a family over $K$, and let $D_b$ be the polynomial $\det X_b$. Let $a_{b,i}\ge0$, for $b\in B$ and
$1\le i\le m_b$, and $n_b\ge0$ be integers, and let $P$ be a polynomial such that
\[
 P(X')=\prod_iT_{ii}^{a_{b,i}+n_b}\,P(X)
\]
whenever $X$ is a family over $K$ with $\det X_{b'}\ne0$ for every $b'$, $b\in B$, $T\in GL_{m_b}(K)$ is upper
triangular and $X'$ is $X$ with $X_b$ replaced by $X_bT$. Then $P=\big(\prod_bD_b^{n_b}\big)G$ for a polynomial $G$
over $K$.
\end{lemma}
\leanref{Development/AbelianVarietyTheory/Proofs/CubicalLineBundles/ExistsEqProdBlockDetPolyPowMul.lean}{GlobalTheory.LamT.FieldPoly.exists_eq_prod_blockDetPoly_pow_mul}

This is a computation with polynomials, by induction on $\sum_bn_b$. When $n_b\ge1$, all the exponents at $b$
are positive, $D_b$ divides $P$, and $P/D_b$ satisfies the hypothesis with $n_b-1$ in place of $n_b$.

\begin{lemma}\label{cv-volume-forms-frame-polynomial}
In Situation~\ref{cv-cusp-setting}, suppose that the level of $X$ is a level $(M,\mathfrak M)$ with $M$ prime
to $15$, and let $Y$ be a deep cover of $X$ (Definition~\ref{cv-deep-cover}) with coefficient field $F_Y$.
Framed points of $Y$ over $F_Y$-algebras are defined as over fields, and the eigenspace of an embedding
$\tau$ of $E$ on the Lie algebra has rank $m_\tau$. For a framed point $y$ over an $F_Y$-algebra $S$, whose
basis of the eigenspace of $\tau$ is $u_{\tau,1},\dots,u_{\tau,m_\tau}$, and a family $g=(g_\tau)$ of matrices
$g_\tau\in GL_{m_\tau}(S)$, let $y\cdot g$ be the framed point with the point and the section of $y$ whose basis
of the eigenspace of $\tau$ has $i$-th vector $\sum_j(g_\tau)_{ji}u_{\tau,j}$. A rational section $f$ of weight
$\kappa=(\alpha,n)$ satisfies the transformation law of its weight: over a field, $f(y\cdot T)$ is
$\prod_{\tau,i}T_{\tau,ii}^{\alpha_{\tau,i}}f(y)$ for every family $T$ of invertible upper triangular matrices,
and multiplying the section of $y$ by $c$ multiplies $f(y)$ by $c^n$. Polynomials $P$ in the entries of
such families $g$, and their values $P(g)$, are as in Lemma~\ref{cv-volume-forms-triangular}. Let $K\supset F_Y$
be a field, $y$ a framed point of $Y$ over $K$ and $f$ a rational section of $Y$.
\begin{enumerate}
\item There are a polynomial $P$ over $K$ and an integer $N\ge0$ such that
  $\big(\prod_\tau\det g_\tau\big)^Nf(y_S\cdot g)=P(g)$ for every $K$-algebra $S$ and every family $g$ of
  invertible matrices over $S$, where $y_S$ is the base change of $y$ to $S$ and $P(g)$ is computed in $S$.
\item Suppose that $f$ has weight $(\alpha,n)$, and let $k_\tau\ge0$ be integers with
  $a_{\tau,i}=\alpha_{\tau,i}+k_\tau\ge0$ for all $\tau$ and $i$. Then there is a polynomial $G$ over $K$,
  homogeneous of degree $\sum_{\tau,i}a_{\tau,i}$, with $\big(\prod_\tau\det g_\tau^{k_\tau}\big)f(y\cdot g)=G(g)$
  for every family $g$ of invertible matrices over $K$.
\end{enumerate}
\end{lemma}
\leanref{Development/AbelianVarietyTheory/Proofs/DivisorProofs/ExistsNumberFieldFrameValueNeZeroOfComplex.lean}{GlobalTheory.FrameDenseZ.exists_poly_div_det_secZ}
\leanref{Development/AbelianVarietyTheory/Proofs/DivisorProofs/ExistsIsHomogeneousSecZGenFix.lean}{GlobalTheory.FieldPolyZ.exists_isHomogeneous_secZ_genFix}
\begin{proof}
(1) Let $A$ be the $K$-algebra of polynomials in the entries of a family $\mathcal X$ of matrices of sizes
$m_\tau$, with the determinants $\det\mathcal X_\tau$ inverted, and let $u=y_A\cdot\mathcal X$. The value
$f(u)\in A$ is $P(\mathcal X)/(\prod_\tau\det\mathcal X_\tau)^N$ for a polynomial $P$ over $K$ and some $N\ge0$.
Given $S$ and $g$, the homomorphism of $K$-algebras $A\to S$ that sends $\mathcal X$ to $g$ carries $u$ to
$y_S\cdot g$, and values of rational sections commute with base change. So
$f(y_S\cdot g)=P(g)/(\prod_\tau\det g_\tau)^N$.

(2) The field $K$ contains $F_Y$, so it is infinite. Let $P$ and $N$ be as in (1) for $S=K$, put
$s=\sum_\tau k_\tau$, and let $P'=(\prod_\tau D_\tau)^sP$, where $D_\tau$ is the polynomial $\det g_\tau$. Then
\[
 \Big(\prod_\tau\det g_\tau\Big)^{s+N}f(y\cdot g)=P'(g)\tag{$*$}
\]
for every family $g$ of invertible matrices over $K$. Let $\tau$ be an embedding, $T\in GL_{m_\tau}(K)$ upper
triangular, $g$ a family of invertible matrices over $K$, and $g'$ the family $g$ with $g_\tau$ replaced by
$g_\tau T$. Then $y\cdot g'=(y\cdot g)\cdot T'$, where $T'$ is $T$ at $\tau$ and the identity elsewhere, so the
transformation law gives $f(y\cdot g')=\prod_iT_{ii}^{\alpha_{\tau,i}}f(y\cdot g)$. With $(*)$ and
$\det(g_\tau T)=\det g_\tau\prod_iT_{ii}$,
\[
 P'(g')=\prod_iT_{ii}^{a_{\tau,i}+s+N-k_\tau}\,P'(g).
\]
Lemma~\ref{cv-volume-forms-triangular}, with the exponents $a_{\tau,i}$ and $n_\tau=s+N-k_\tau\ge0$, gives
$P'=(\prod_\tau D_\tau^{s+N-k_\tau})G$ for a polynomial $G$. Dividing $(*)$ by
$\prod_\tau\det g_\tau^{s+N-k_\tau}$ gives $(\prod_\tau\det g_\tau^{k_\tau})f(y\cdot g)=G(g)$.

Homogeneity. Put $d=\sum_{\tau,i}a_{\tau,i}$. Let $g$ be a family of invertible matrices over $K$ and $2g$ the
family of the $2g_\tau$. Then $y\cdot2g=(y\cdot g)\cdot T$ for the family $T$ of the matrices $2\cdot1$, so
$f(y\cdot2g)=2^{\sum\alpha_{\tau,i}}f(y\cdot g)$, and
\[
 G(2g)=\prod_\tau\det(2g_\tau)^{k_\tau}f(y\cdot2g)=2^{\sum_\tau m_\tau k_\tau+\sum_{\tau,i}\alpha_{\tau,i}}G(g)=2^dG(g).
\]
So the polynomial $\big(G(2g)-2^dG(g)\big)\prod_\tau D_\tau$ vanishes at every family over the infinite field
$K$, hence is zero. As $\prod_\tau D_\tau\ne0$, $G(2g)=2^dG(g)$ as polynomials, and comparing coefficients,
every monomial of $G$ has degree $d$.
\end{proof}
% Lean: (1) is exists_poly_div_det_secZ = FieldPolyZ.exists_poly_div_det_of_natural, with the generic ring
% GlobalTheory.LamT.RatFn.exists_genericRing, Frames.FramedBaseChange.exists_frMap and existsUnique_frChg for
% u, and evR_frMap for the base change. (2) is exists_isHomogeneous_secZ_genFix: hA is (*), hB is the
% transformation law (GlobalTheory.evR_frBor_of_mem_secZ), hP'tr the equivariance, then
% exists_eq_prod_blockDetPoly_pow_mul and isHomogeneous_of_scaling with the diagonal matrices 2. The Lean's
% change of frame by FrChg (Y^T) 1 is the text's y.g, with g = Y.

\begin{lemma}\label{cv-volume-forms-frame-nonzero}
In the setting of Lemma~\ref{cv-volume-forms-frame-polynomial}, let $K\supset F_Y$ be a number field, let
$\sigma\colon F_Y\to\mathbf C$ be an embedding and $\sigma_1\colon K\to\mathbf C$ an embedding that extends
$\sigma$. Let $y$ be a framed point of $Y$ over $K$, $f$ a rational section of $Y$ of some weight, and $w$ a
framed point of $Y$ over $\mathbf C$, an $F_Y$-algebra through $\sigma$, whose point is the image of the point
of $y$ under $\sigma_1$. If $f(w)\ne0$, there is a framed point $y'$ of $Y$ over $K$ with the point of $y$ and
$f(y')\ne0$.
\end{lemma}
\leanref{Development/AbelianVarietyTheory/Proofs/DivisorProofs/ExistsNumberFieldFrameValueNeZeroOfComplexStatement.lean}{GlobalTheory.T.exists_numberFieldFrame_value_ne_zero_of_complex}
\begin{proof}
Let $P$ and $N$ be as in Lemma~\ref{cv-volume-forms-frame-polynomial}(1) for $y$ and $f$, and let $n$ be the
last entry of the weight of $f$. Let $z$ be the base change of $y$ to $\mathbf C$ along $\sigma_1$. The framed
points $z$ and $w$ have the same point, so there are a family $h$ of invertible complex matrices and
$c\in\mathbf C^\times$ such that $w$ is $z\cdot h$ with its section multiplied by $c$. By the transformation law,
$f(w)=c^nf(z\cdot h)$, so $f(z\cdot h)\ne0$. By Lemma~\ref{cv-volume-forms-frame-polynomial}(1), with $S=\mathbf C$
through $\sigma_1$, $(\prod_\tau\det h_\tau)^Nf(z\cdot h)$ is the value at $h$ of the polynomial obtained from $P$
by applying $\sigma_1$ to its coefficients. Hence $P\ne0$. As $K$ is infinite and $P\prod_\tau D_\tau$, with $D_\tau$
the polynomial $\det g_\tau$, is a nonzero polynomial, it does not vanish at some family $g$ of matrices over $K$. Then every $\det g_\tau$ is
nonzero and $P(g)\ne0$, and $y'=y\cdot g$ has the point of $y$ and $(\prod_\tau\det g_\tau)^Nf(y')=P(g)\ne0$.
\end{proof}
\leanref{Development/AbelianVarietyTheory/Proofs/DivisorProofs/ExistsNumberFieldFrameValueNeZeroOfComplex.lean}{GlobalTheory.FrameDenseZ.exists_frame_evR_ne_zero_of_complex}
% Lean: exists_frame_evR_ne_zero_of_complex: Frames.FramedBaseChange.exists_frMap gives z,
% Frames.frChg_exists_unique gives h and c (the Lean's g and s), evR_frBor_of_mem_secZ the factor c^n, and
% exists_frame_evalFr_ne_zero the family g (a point where P times the block determinants does not vanish).

\begin{lemma}\label{cv-volume-forms-denominator}
In the setting of Lemma~\ref{cv-volume-forms-frame-polynomial}, let $K\supset F_Y$ be a number field and $y_0$ a
framed point of $Y$ over $K$. Let $\kappa=(\alpha,n)$ be a weight and $k_\tau=\sum_i\max(0,-\alpha_{\tau,i})$.
Suppose that there is a nonzero $t\in\mathcal O_K$ such that $tG$ has coefficients integral over $\mathbf Z$
whenever $f$ is an integral rational section of weight $\kappa$ and $G$ is a polynomial over $K$ with
$(\prod_\tau\det g_\tau^{k_\tau})f(y_0\cdot g)=G(g)$ for every family $g$ of invertible matrices over $K$. Then for
every framed point $y$ of $Y$ over $K$ with the point of $y_0$ there is $N\in K^\times$ such that $Nf(y)$ is
integral over $\mathbf Z$ for every integral rational section $f$ of weight $\kappa$.
\end{lemma}
\leanref{Development/AbelianVarietyTheory/Proofs/DivisorProofs/ExistsFrameDenominatorDominantWeightPowerCm.lean}{GlobalTheory.LatFull.CMInt.exists_frame_denom_dominant_pow}
\begin{proof}
Put $a_{\tau,i}=\alpha_{\tau,i}+k_\tau\ge0$ and $d=\sum_{\tau,i}a_{\tau,i}$. The framed points $y_0$ and $y$ have the
same point, so there are a family $g$ of invertible matrices over $K$ and $c\in K^\times$ such that $y$ is
$w_0=y_0\cdot g$ with its section multiplied by $c$. By the transformation law, $f(y)=c^nf(w_0)$ for every
rational section $f$ of weight $\kappa$. For each of the finitely many monomials $g^\mu$ of degree $d$ in the
entries of $g$ choose a nonzero $N_\mu\in\mathcal O_K$ with $N_\mu t^{-1}g^\mu\in\mathcal O_K$, and put
$N_0=\prod_\mu N_\mu$ and $N=N_0\big(\prod_\tau\det g_\tau^{k_\tau}\big)c^{-n}$.

Let $f$ be an integral rational section of weight $\kappa$, and let $G$ be the polynomial of degree $d$ of
Lemma~\ref{cv-volume-forms-frame-polynomial}(2) for $y_0$, $f$ and the $k_\tau$. By hypothesis $tG$ has
coefficients in $\mathcal O_K$, so $N_0G(g)=\sum_\mu N_0c_\mu g^\mu$, with $c_\mu$ the coefficients of $G$, lies
in $\mathcal O_K$. Finally
$Nf(y)=N_0\big(\prod_\tau\det g_\tau^{k_\tau}\big)f(w_0)=N_0G(g)$.
\end{proof}
% Lean: exists_frame_denom_dominant_pow, with t from GlobalTheory.T.exists_commonDenom_framePolynomials_cmPoint_generator
% at the good CM point z (the hypothesis here; its y_0 is z.y), Frames.frChg_exists_unique and
% FieldPolyZ.exists_frChg_one_evR_eq for g (the Lean's Y), c and w_0, LatFull.CMInt.exists_denom_eval for N_0 (with the
% fractional ideals t^{-1} O_K), and exists_isHomogeneous_secZ_genFix for G.

\begin{lemma}\label{cv-volume-generators-exponents}
In the setting of Definition~\ref{cv-volume-forms-weights}, let $\lambda$ be a regular family, and let $N\ge1$,
$M\ge0$ and $k\ge0$ be integers with $2M+N\le k+1$. Then there is $w\colon\Phi\to\mathbf Z$ such that
$\kappa((k+1)\lambda+\lambda_w,M)$, where $\lambda_{w,\tau,i}=Nw_\tau$ for all $\tau$ and $i$, is a dominant
weight $(\alpha,2M)$ in the sense of Situation~\ref{cv-cusp-setting}.
\end{lemma}
\leanref{Development/AbelianVarietyTheory/Proofs/DivisorProofs/ExistsDominantExponentsMatchingWeightBarGap.lean}{GlobalTheory.WtMatchArithBar.exists_dom_exponents_bar}
\begin{proof}
For $\tau\in\Phi$ let $X_\tau=(k+1)\lambda_{\tau,p-1}-M$ and $w_\tau=-\lfloor X_\tau/N\rfloor$, so that
$-X_\tau\le Nw_\tau\le-X_\tau+N-1$. The weight has the entries $(k+1)\lambda_{\tau,i}+Nw_\tau-M$ at $\tau$ in position
$i<p$, the entries $-(k+1)\lambda_{\tau,r-1-i}-Nw_\tau-M$ at $\bar\tau$ in position $i<q$, and the last entry $2M$.
For $i<p$, $\lambda_{\tau,i}\ge\lambda_{\tau,p-1}$, so the entry at $\tau$ is at least $X_\tau+Nw_\tau\ge0$. For
$i<q$, $r-1-i\ge p$, so $\lambda_{\tau,r-1-i}\le\lambda_{\tau,p}\le\lambda_{\tau,p-1}-1$ as $\lambda$ is regular, and
the entry at $\bar\tau$ is at least $-X_\tau-Nw_\tau+k+1-2M\ge k+2-2M-N\ge1$. The entries do not increase with the
position, since $\lambda_{\tau,i}$ does not increase with $i$, and $2M$ is even and nonnegative.
\end{proof}
% Lean: exists_dom_exponents_bar, with WtMatchArithBar.exists_shift for w_tau (w = -(X / N)), vbig and vsmall
% for the entries, and mono and nonneg; the Lean writes the weight as (k + 1) wtOf la 0 + wtOf (parWt (N w)) M.

\begin{lemma}\label{cv-volume-generators-integers}
Let $K$ be a number field of degree $n$ with root discriminant $\operatorname{rd}_K=|d_K|^{1/n}$. There are
$\omega_1,\dots,\omega_n\in\mathcal O_K$, linearly independent over $\mathbf Q$, with
$|\sigma(\omega_k)|\le n^3\operatorname{rd}_K$ for every embedding $\sigma\colon K\to\mathbf C$ and every $k$.
\end{lemma}
\leanref{Development/AbelianVarietyTheory/Proofs/CuspFormEvaluations/ExistsShortIndependentIntegers.lean}{GlobalTheory.exists_short_independent_integers}
\begin{proof}
Let $\varphi,\psi\colon K\to\mathbf R^n$ be the $\mathbf Q$-linear maps with coordinates
$\operatorname{Re}\sigma(x)+\operatorname{Im}\sigma(x)$ and $\operatorname{Re}\sigma(x)-\operatorname{Im}\sigma(x)$, one
for each embedding $\sigma$. As the embeddings come in complex conjugate pairs and real ones,
$\langle\varphi(x),\psi(y)\rangle=\operatorname{Tr}_{K/\mathbf Q}(xy)$ and
$|\varphi(x)|^2=|\psi(x)|^2=\sum_\sigma|\sigma(x)|^2$. Let $\mathfrak D^{-1}$ be the inverse different, the set of
$y\in K$ with $\operatorname{Tr}_{K/\mathbf Q}(xy)\in\mathbf Z$ for all $x\in\mathcal O_K$. The images
$\Gamma=\psi(\mathfrak D^{-1})$ and $\varphi(\mathcal O_K)$ are lattices of rank $n$ in $\mathbf R^n$, and the pairing
shows that $\varphi(\mathcal O_K)$ is the dual lattice of $\Gamma$.

Let $y\in\mathfrak D^{-1}$ be nonzero. The ideal $y\mathfrak D$ is integral and $\mathfrak D$ has norm $|d_K|$, so
$|N_{K/\mathbf Q}(y)|\ge|d_K|^{-1}$, and by the inequality of arithmetic and geometric means
$\sum_\sigma|\sigma(y)|^2\ge n|N_{K/\mathbf Q}(y)|^{2/n}\ge n|d_K|^{-2/n}$. So every nonzero vector of $\Gamma$ has
length at least $|d_K|^{-1/n}$. By Lemma~\ref{cv-transference}, $\varphi(\mathcal O_K)$ contains $n$ linearly
independent vectors $\varphi(\omega_k)$ of length at most $n^2|d_K|^{1/n}$. The $\omega_k$ are linearly independent
over $\mathbf Q$, and $|\sigma(\omega_k)|\le|\varphi(\omega_k)|\le n^2\operatorname{rd}_K\le n^3\operatorname{rd}_K$.
\end{proof}
% Lean: exists_short_independent_integers, with ShortIntegers.realCoords e 1 and e (-1) for phi and psi
% (inner_realCoords_one_neg_one_eq_trace, norm_sq_realCoords_eq_sum_norm_sq), the dual basis beta of an integral
% basis alpha, GlobalTheory.finrank_mul_discr_rpow_le_sum_sq_of_mem_traceDual for the lower bound, and
% ReciprocityTheory.ChartRestriction.G.T0.transference_dual (Lemma cv-transference).

\begin{lemma}\label{cv-volume-generators-short}
Let $K$ be a number field of degree $d$, let $V$ be a finite-dimensional $K$-vector space, and for every
embedding $\sigma\colon K\to\mathbf C$ let $N_\sigma\colon V\to\mathbf R_{\ge0}$ satisfy
$N_\sigma(ax)=|\sigma(a)|N_\sigma(x)$ and $N_\sigma(x+x')\le N_\sigma(x)+N_\sigma(x')$. Let $L\subset V$ be a finitely
generated $\mathcal O_K$-submodule, let $A$ be a real number, and let $S\subset L$ be a subset that spans $V$
over $K$ with $N_\sigma(y)\le A$ for all $y\in S$ and all $\sigma$. Let $m=\dim_KV$. Then $L$ is generated as an
abelian group by $n\le2dm$ elements $w_1,\dots,w_n$ with
$N_\sigma(w_i)\le\max(1,d^4m\operatorname{rd}_KA)$ for all $i$ and $\sigma$.
\end{lemma}
\leanref{Development/AbelianVarietyTheory/Proofs/DivisorProofs/ExistsShortGeneratorsOfShortSpanningFamilyStatement.lean}{GlobalTheory.T.exists_short_generators_of_short_spanning_family}
\begin{proof}
Choose a basis $t_1,\dots,t_m$ of $V$ over $K$ in $S$, and let $\omega_1,\dots,\omega_d$ be as in
Lemma~\ref{cv-volume-generators-integers}. The $dm$ elements $e_{k,j}=\omega_kt_j$ form a basis of $V$ over
$\mathbf Q$, lie in $L$, and satisfy $N_\sigma(e_{k,j})=|\sigma(\omega_k)|N_\sigma(t_j)\le C$ with
$C=d^3\operatorname{rd}_KA$. For $x\in V$ write $x=\sum_{k,j}c_{k,j}e_{k,j}$ with $c_{k,j}\in\mathbf Q$, and let
$\rho(x)=\sum_{k,j}\{c_{k,j}\}e_{k,j}$, with the fractional parts $\{c_{k,j}\}\in[0,1)$. Then
$x-\rho(x)$ lies in the subgroup generated by the $e_{k,j}$, and $N_\sigma(\rho(x))\le dmC$.

The abelian group $L$ is finitely generated and torsion-free, so it has a basis $b_1,\dots,b_{m'}$ over
$\mathbf Z$, which is linearly independent over $\mathbf Q$; so $m'\le dm$. Take the elements $e_{k,j}$ and
$\rho(b_1),\dots,\rho(b_{m'})$, at most $2dm$ in all. They lie in $L$, since $\rho(b_l)=b_l-(b_l-\rho(b_l))$, and
they generate $L$, since $b_l=(b_l-\rho(b_l))+\rho(b_l)$. Their norms are at most $dmC=d^4m\operatorname{rd}_KA$, as
$dm\ge1$.
\end{proof}
\leanref{Development/AbelianVarietyTheory/Proofs/DivisorProofs/ExistsShortGeneratorsOfShortSpanningFamily.lean}{GlobalTheory.exists_isShortGen_of_short_span}
% Lean: exists_isShortGen_of_short_span: the basis t from exists_linearIndependent, the Q-basis E = ob.smulTower vb,
% red for rho, the Z-basis ZB of L (Module.free_of_finite_type_torsion_free'), and the family Sum.elim E (red o bZ);
% the bound is IsShortGen with abelian_2 = max 1 (d^4 m rd A), where m = rkK L = dim V.

\begin{definition}\label{cv-volume-twisting-section}
Let $E$ be a CM field with CM type $\Phi$, and let $X$ be the moduli scheme of a PEL datum of
signature $(p,q)$ over $E$ at a level $(M,\mathfrak M,N)$ (Definition~\ref{cv-model-toroidal-semiabelian}),
over its coefficient field $F_c$, with uniformizations at the complex embeddings of $F_c$ normalized by a
real number $\xi>0$. Every complex embedding $\sigma$ of $F_c$ has a good embedding
$\sigma^g\in\{\sigma,\bar\sigma\}$, and forms, cusp forms, frames of a connected component and pointwise
norms are read in the uniformization at a good embedding (Situation~\ref{cv-cusp-setting}). For an
embedding $\tau\colon E\to\mathbf C$ let $m_\tau=p$ if $\tau\in\Phi$ and $m_\tau=q$ otherwise, and put
$g=\sum_\tau m_\tau$.
\begin{enumerate}
\item A \emph{weight} $G$ is a family of integers $G_{\tau,i}$, for the embeddings $\tau$ of $E$ and
  $0\le i<m_\tau$, with a further integer, its last entry. Put $|G|=\sum_{\tau,i}G_{\tau,i}$. For
  $w\colon\Phi\to\mathbf Z$ and $m\in\mathbf Z$ write $\kappa(w,m)$ for the weight $\kappa(\lambda,m)$ of
  Definition~\ref{cv-volume-forms-weights} with $\lambda_{\tau,i}=w_\tau$ for all $i$: its entries are
  $w_\tau-m$ at $\tau\in\Phi$ and $-w_\tau-m$ at $\bar\tau$, and its last entry is $2m$.
\item For a rational section $s$ of weight $G$ on $X$, a good embedding $\sigma$ and a frame $x$ of a
  connected component at $\sigma$, put $|s|_\xi(x)=(2\pi\xi)^{-|G|/2}|s_\sigma|(x)$, where $|s_\sigma|$ is
  the pointwise norm of the reading of $s$ at $\sigma$.
\item Let $y$ be a point of $X$ with values in a field $K\supset F_c$ of characteristic $0$, with abelian
  variety $A_y$. A \emph{framed point} over $K$ is $y$ together with a basis of each eigenspace
  $\operatorname{Lie}_\tau(y)$ of the Lie algebra of $A_y$, of dimension $m_\tau$, and a generator
  $\vartheta$ of the sections of the symmetric line bundle of the principal polarization of $A_y$. These
  bases, in a fixed order, form a basis $v$ of the Lie algebra; write the framed point as $(y,v,\vartheta)$.
  Framed points over an $F_c$-algebra $S$ are defined in the same way. A framed point over $S$ is
  \emph{tested} if it is the image, under a homomorphism of $F_c$-algebras $R\to S$, of a framed point over
  an $F_c$-algebra $R$ that is the quotient of a polynomial ring $F_c[x_1,\dots,x_n]$ by a radical ideal.
\item A \emph{key section} of $X$ is an invertible rational section of weight $\kappa(0,-1)$.
\item Let $Z$ be a connected component of $X$, with idempotent $1_Z$, and let $G$ be a weight with even
  last entry. A \emph{twisting section} $t$ of weight $G$ on $Z$ consists of rational sections $s$ of
  weight $G$ and $s'$ of weight $-G$ with $1_Zs=s$, $1_Zs'=s'$ and $ss'=1_Z$, numbers $c_\sigma>0$ for the
  complex embeddings $\sigma$ of $F_c$, and real functions $l$ and $h$ on the finite places of $F_c$, zero
  at almost all places, such that:
  \begin{enumerate}
  \item[(a)] $|s|_\xi(x)=c_\sigma$ at every frame $x$ of $Z$ at $\sigma^g$, for every $\sigma$;
  \item[(b)] $l(w)\le v(s)\le h(w)$ for the valuation $v$ at every integral point of $Z$ above every
    finite place $w$;
  \item[(c)] for every weight $G'$ with even last entry, $fs$ is a cusp form on $Z$ of weight $G'+G$ for
    every cusp form $f$ on $Z$ of weight $G'$, and $fs'$ is a cusp form on $Z$ of weight $G'$ for every
    cusp form $f$ on $Z$ of weight $G'+G$;
  \item[(d)] $Z$ has frames at every good embedding, the framed point over $\mathbf C$ at every frame of $Z$
    at a good embedding is tested, and $Z$ has integral points above every finite place.
  \end{enumerate}
  Its \emph{degree} is
  \[
   \deg t=\frac1{[F_c:\mathbf Q]}\Big(\sum_wl(w)\log Nw-\sum_\sigma\log c_\sigma\Big).
  \]
\end{enumerate}
\end{definition}
\leanref{Development/AbelianVarietyTheory/Definitions/Foundations/IsProbabilityMeasureMuPart.lean}{GlobalTheory.AllWt.homN}
\leanref{Development/AbelianVarietyTheory/Definitions/AutomorphicMembership/FlagComponentFramedPointFrame.lean}{GlobalTheory.ShFam.KeySec}
\leanref{Development/AbelianVarietyTheory/Definitions/Divisors/MulSmulIdemMem.lean}{GlobalTheory.FlatTw.TwSec}
\leanref{Development/AbelianVarietyTheory/Definitions/Divisors/MulSmulIdemMem.lean}{GlobalTheory.FlatTw.TwSec.degLo}

\begin{lemma}\label{cv-volume-twisting-products}
In the setting of Definition~\ref{cv-volume-twisting-section}, let $Z$ be a connected component of $X$,
and let $t_0$ and $t_1$ be twisting sections on $Z$ of weights $G_0$ and $G_1$, with sections $s_0$ and
$s_1$, numbers $c_{0,\sigma}$ and $c_{1,\sigma}$ and functions $l_0$ and $l_1$.
\begin{enumerate}
\item There is a twisting section of weight $G_0+G_1$ on $Z$ with section $s_0s_1$, numbers
  $c_{0,\sigma}c_{1,\sigma}$ and function $l_0+l_1$. Its degree is $\deg t_0+\deg t_1$.
\item For every integer $n\ge0$ there is a twisting section of weight $nG_0$ on $Z$ with function $nl_0$
  and degree $n\deg t_0$.
\end{enumerate}
\end{lemma}
\leanref{Development/AbelianVarietyTheory/Proofs/DivisorProofs/TwSecMulPow.lean}{GlobalTheory.FlatTw.twSec_mul_pow}
\begin{proof}
(1) Let $s_0'$, $s_1'$, $h_0$ and $h_1$ be the other data of $t_0$ and $t_1$. Take the sections $s_0s_1$
and $s_0's_1'$, the numbers $c_{0,\sigma}c_{1,\sigma}$ and the functions $l_0+l_1$ and $h_0+h_1$. Then
$1_Zs_0s_1=s_0s_1$, $1_Zs_0's_1'=s_0's_1'$ and $(s_0s_1)(s_0's_1')=1_Z1_Z=1_Z$. The pointwise norm is
multiplicative and the factor $(2\pi\xi)^{-|G|/2}$ is multiplicative in $G$, so
$|s_0s_1|_\xi=c_{0,\sigma}c_{1,\sigma}$ at the frames of $Z$. Valuations are additive, which gives (b).
Multiplication by $s_0s_1$ is multiplication by $s_0$ followed by multiplication by $s_1$, and
multiplication by $s_0's_1'$ is multiplication by $s_1'$ followed by multiplication by $s_0'$, which gives
(c); (d) holds for $Z$. As $\log(c_{0,\sigma}c_{1,\sigma})=\log c_{0,\sigma}+\log c_{1,\sigma}$, the degree
is $\deg t_0+\deg t_1$.

(2) For $n=0$ take the sections $1_Z$ and $1_Z$, the numbers $1$ and the functions $0$: the pointwise norm of
$1_Z$ is $1$ at the frames of $Z$, its valuation at the integral points of $Z$ is $0$, multiplication by
$1_Z$ fixes the cusp forms on $Z$, and the degree is $0$. For $n+1$ apply (1) to the twisting section for
$n$ and $t_0$.
\end{proof}

\begin{lemma}\label{cv-volume-twisting-points}
In the setting of Definition~\ref{cv-volume-twisting-section}, suppose that $p+q\ge1$, that the prime-to-$\ell$
part of $M$ is at least $3$ for every prime $\ell$, and that the uniformizations at the good embeddings
satisfy the compatibility conditions of the moduli problem. Then every connected component $Z$ of $X$ has
frames at every good embedding, the framed point over $\mathbf C$ at every frame of $Z$ at a good embedding is
tested (Definition~\ref{cv-volume-twisting-section}(3)), and $Z$ has integral points above every finite place
of $F_c$.
\end{lemma}
\leanref{Development/AbelianVarietyTheory/Proofs/FoundationProofs/ExistsNumberFieldPoint.lean}{GlobalTheory.FlatLine.hasPts}
\leanref{Development/AbelianVarietyTheory/Definitions/Foundations/ExistsFlatSecRankZero.lean}{GlobalTheory.FlatLine.HasPts}
\begin{proof}
Let $Z$ be a connected component of $X$.

A point over a number field. The scheme $X$ is smooth, hence of finite type, over $F_c$, and $Z$ is open in
$X$ and not empty. Let $\operatorname{Spec}A$ be a nonempty affine open subscheme of $Z$ and $\mathfrak m$ a
maximal ideal of $A$. The $F_c$-algebra $A$ is finitely generated, so the field $L=A/\mathfrak m$ is a finite
extension of $F_c$ by Zariski's lemma \cite[Corollary~5.24]{AM69}. This gives a point $y_L$ of $X$ over the
number field $L$ that lies on $Z$.

Frames. Let $\sigma$ be a good embedding of $F_c$. Extend $\sigma$ to an embedding of $L$ in $\mathbf C$; it
carries $y_L$ to a point $y$ of $X$ over $\mathbf C$ via $\sigma$ that lies on $Z$. Over a field the
eigenspaces $\operatorname{Lie}_\tau(y)$ are vector spaces of dimension $m_\tau$, and the sections of the
symmetric line bundle of the principal polarization form a space of dimension one \cite[\S16]{Mum70};
bases and a generator make $y$ a framed point over $\mathbf C$ via $\sigma$. One of the compatibility
conditions of the moduli problem is that every framed point over $\mathbf C$ via $\sigma$ is triangularly
equivalent to the framed point at a frame of a component of the uniformization at $\sigma$, as in
Lemma~\ref{cv-uniformization-frames}(7). So the framed point $y$ is triangularly equivalent to the framed
point at a frame $x$. The two have the same point, which lies on $Z$; so $x$ is a frame of $Z$ at $\sigma$.

Tested framed points. Every framed point $y$ over a reduced $F_c$-algebra $S$ is tested, and $\mathbf C$ is
reduced. Indeed, $X$ is of finite type over $F_c$, its universal abelian scheme, the Lie algebra and the line
bundle are of finite presentation, and $S$ is the filtered union of its finitely generated $F_c$-subalgebras;
so the point of $y$, its bases and its generator descend to a framed point $y_1$ over one of them, $S_1$, whose
image under the inclusion $S_1\subset S$ is $y$ \cite[\S8]{Gro66}. The algebra $S_1$ is reduced, as a subring
of $S$, and finitely generated, so it is isomorphic to the quotient of a polynomial ring over $F_c$ by a
radical ideal; composing the inclusion with this isomorphism shows that $y$ is tested.

Integral points. Let $w$ be a finite place of $F_c$, over the prime $\ell$, let $\mathfrak P$ be a prime of
$\mathcal O_L$ above $w$, and let $T$ be the completion of $\mathcal O_L$ at $\mathfrak P$, a complete
discrete valuation ring above $w$ with finite residue field, whose fraction field we write $L_{\mathfrak P}$.
The point $y_L$ gives a point $y$ of $X$ over $L_{\mathfrak P}$ that lies on $Z$. Let $O=F_c\cap T$ be the
valuation ring of $F_c$ at $w$. As the prime-to-$\ell$ part of the level is at least $3$, the proof of
Theorem~\ref{cv-model-toroidal}, for an embedding of $F_c$ in $\mathbf C_\ell$ through $L_{\mathfrak P}$, gives a
model $(\mathcal X,\mathcal A,j)$ of $X$ over $O$ whose semi-abelian scheme $\mathcal A$ is theta-split, so that
Lemma~\ref{cv-model-toroidal-pushforward} applies to it. As $\mathcal X$ is projective over $O$, the
valuative criterion of properness extends $j\circ y$ to a point $\hat y$ of $\mathcal X$ over $T$. The
pullback $G=\hat y^*\mathcal A$ is a semi-abelian scheme over $T$ whose generic fibre is $A_y$, with its
action of $\mathcal O_E$ and its line bundle. Since $G$ is smooth over $T$, its Lie algebra is a free
$T$-module of finite rank. For an embedding $\tau$ of $E$ let $\tau'\colon\mathcal O_E\to T$ be $\tau$
followed by $F_c\to T$. The elements of the Lie algebra on which every $a\in\mathcal O_E$ acts by
$\tau'(a)$ form a $T$-submodule; it is free, as $T$ is a discrete valuation ring, and its rank is the
dimension $m_\tau$ of the corresponding eigenspace $\operatorname{Lie}_\tau(y)$ of the generic fibre. The
line bundle of $\mathcal A$ is generated near the image of the closed point of $\operatorname{Spec}T$ by a
section $\theta$ as in Lemma~\ref{cv-model-toroidal-pushforward}, and so the sections over $T$ of the line
bundle of $G$ are the $T$-multiples of the pullback of $\theta$. Bases of these modules and this generator
form an integral framed point over $L_{\mathfrak P}$ whose point is $y$. It lies on $Z$, and it is algebraic
over $F_c$ since $L$ is a number field. So $Z$ has an integral point above $w$.
\end{proof}
% Lean: hasPts takes the frames from framesOn_nonempty (exists_numberField_point through exists_finite_point and
% finite_of_finite_type_of_isJacobsonRing; IsAlgClosed.lift for the embedding; exists_frPt_local with the law
% SecTwistSt for the framed point; the field hU.2.2.1 of IsUnf for the frame) and the integral points from
% valOn_nonempty (locFld for T, then IntFrExt, from GlobalTheory.Exist.frame_laws_of_moduli, which is
% GlobalTheory.Exist.intFrExt_of_good: exists_model_pushInvertible_over_completeDVR = exists_model_pushInvertible_under_cdvr
% takes the model Y p W of the local system of PartI over the valuation ring c.V W of F_c at w, with PushInvertible;
% intFr_of_model extends the point by extPt and builds the integral frame by intFrames_of_laws, with lieFree for
% the bases of the eigenvectors and thetaFrame_of for the generator). HasPts also asks that the framed points at
% the frames be tested (FlatLine.Tested = Formalism.L2.Pj.FromTest over a TRing, a polynomial ring modulo a
% radical ideal); this is AllTest of frame_laws_of_moduli (GlobalTheory.Frames.allTest_of_secBaseChange:
% TestDescent.exists_fg_descent, reduced_finiteType_iff_testRing, exists_framedPoint_baseChange_algebraMap and
% FramedBaseChange.cancel). The Lean proves the descent to a finitely generated subalgebra its own way, through
% moduli_point_from_fg_subalgebra and the universal frames over the test rings of a chart cover
% (exists_chart_cover, exists_universalFrame_testRing); the text cites EGA IV_3 Section 8 for it.

\begin{lemma}\label{cv-volume-twisting-coset}
Let $K$ be an algebraically closed field, let $A$ be an abelian variety over $K$ with a symmetric line bundle
$L$ that defines a principal polarization, and let $\vartheta$ be a nonzero section of $L$. Let $n\ge2$ be an
integer that is invertible in $K$, and let $a\in A(K)$. Then there is $x\in A(K)$ with $nx=0$ such that
$\vartheta$ does not vanish at $a+x$.
\end{lemma}
\leanref{ReciprocityTheory/ReductiveGroups/ArithmeticGeometry/ExtensionDescent/GlobalSectionsRankAlgClosed.lean}{ReciprocityTheory.ChartRestriction.G.Out.GoodPtSt}
\leanref{Development/AbelianVarietyTheory/Proofs/AbelianSchemeLawProofs/GoodPtOfIrredOfMulPull.lean}{GlobalTheory.GoodPoint.goodPt_of_irred_of_mulPull}
\leanref{Development/AbelianVarietyTheory/Proofs/UniformizationProofs/ExistsGeneratorNeZero.lean}{GlobalTheory.UT.ExactTerm.exists_good_point}
\begin{proof}
Multiplication by $n$ on $A$ is finite, \'etale and surjective, as $n$ is invertible in $K$, so there is
$b\in A(K)$ with $nb=a$. As $L$ is symmetric, $[n]^*L\cong L^{n^2}$ \cite[\S6, Corollary~3]{Mum70}; let $s_1$
be the section of $L^{n^2}$ that corresponds to $[n]^*\vartheta$. Suppose that $\vartheta$ vanishes at $a+x$
for every $x$ with $nx=0$. Let $B=b+A[n^2](K)$, and let $V$ be the space of sections of $L^{n^2}$ that vanish
at every point of $B$. For $c\in A[n^2](K)$, $n(b+c)=a+nc$ with $n(nc)=0$, so $s_1$ vanishes at $b+c$; thus
$s_1\in V$.

The theta group $\mathcal G(L^{n^2})$ acts on the sections of $L^{n^2}$, and an element lying over a point $c$
carries a section that vanishes at $z+c$ to a section that vanishes at $z$. Every element of
$\mathcal G(L^{n^2})$ lies over a point of $A[n^2](K)$, the kernel of the polarization of $L^{n^2}$, and $B$ is
stable under translation by such points; so $V$ is stable under $\mathcal G(L^{n^2})$. As $n^2$ is invertible in
$K$, the sections of $L^{n^2}$ form an irreducible representation of $\mathcal G(L^{n^2})$ \cite[\S23]{Mum70}.
So $V=0$ or $V$ is the whole space. If $V=0$, then $s_1=0$, so $[n]^*\vartheta=0$ and $\vartheta=0$, as $[n]$ is
faithfully flat; this is excluded. Otherwise every section of $L^{n^2}$ vanishes at $b$; but $L^{n^2}$ is a
power of an ample line bundle with exponent at least $2$, so it is generated by its global sections
\cite[\S17]{Mum70}. This contradiction proves the lemma.
\end{proof}
% Lean: GoodPtSt is a field of the Lean's laws of abelian schemes, proved by goodPt_of_irred_of_mulPull from the
% irreducibility of the theta group representation (IrredSt), the normalized isomorphism [n]^*L = L^{n^2}, the
% finite flat surjective multiplication by n (etale when n is invertible) and the global generation of L^n for
% n >= 2 (the last hypothesis, used for L^{n n}); exists_pow_eq gives b, thetaGrp_pt_pow the points of the theta
% group, and vanishing P.L _ B the space V. exists_good_point transports it to the universal family.

\begin{lemma}\label{cv-volume-twisting-combination}
Let $K$ be a field of characteristic $0$, let $I$ and $J$ be finite sets, let $V$ and $W_i$, for $i\in I$, be
vector spaces over $K$, let $\ell_i\colon V\to W_i$ be linear maps, and let $s_j\in V$ for $j\in J$. Suppose
that for every $i$ there is $j$ with $\ell_i(s_j)\ne0$. Then there are integers $0\le c_j\le|I|$ such that
$\ell_i\big(\sum_jc_js_j\big)\ne0$ for every $i\in I$.
\end{lemma}
\leanref{Development/AbelianVarietyTheory/Proofs/CovolumeBounds/EqZeroOfEvalGridEqZeroOfTotalDegreeLe.lean}{GlobalTheory.exists_sum_natCast_smul_apply_ne_zero}
\begin{proof}
First, for $v,w\in V$ there is an integer $0\le c\le|I|$ with $\ell_i(cw+v)\ne0$ for every $i$ such that
$\ell_i(v)\ne0$ or $\ell_i(w)\ne0$. Otherwise each of the $|I|+1$ integers $0\le c\le|I|$ has an index $i_c$
with $\ell_{i_c}(v)$ or $\ell_{i_c}(w)$ nonzero and $\ell_{i_c}(cw+v)=0$. If $i_c=i_{c'}=i$, then
$\ell_i(w)\ne0$, since otherwise $\ell_i(v)=\ell_i(cw+v)=0$; and $(c-c')\ell_i(w)=0$, so $c=c'$, as $K$ has
characteristic $0$. So $c\mapsto i_c$ is an injective map from a set of $|I|+1$ elements to $I$, which is
impossible.

Next, for every finite set $T\subset J$ there are integers $0\le c_j\le|I|$ with
$\ell_i\big(\sum_{j\in T}c_js_j\big)\ne0$ for every $i$ such that $\ell_i(s_j)\ne0$ for some $j\in T$. This is
proved by induction on $T$. For $T$ empty there is nothing to prove. If it holds for $T$, with integers $c_j$,
and $a\notin T$, apply the first step to $v=\sum_{j\in T}c_js_j$ and $w=s_a$, and put $c_a=c$. For $T=J$ this is
the claim.
\end{proof}
% Lean: exists_sum_natCast_smul_apply_ne_zero is exists_sum_nsmul_apply_ne_zero (vector spaces over a field of
% characteristic 0 are torsion free), the induction exists_sum_nsmul_apply_ne_zero_finset, with the first step
% exists_nsmul_add_apply_ne_zero.

\begin{definition}\label{cv-volume-twisting-level-structures}
In the setting of Definition~\ref{cv-volume-twisting-section}, let $\Lambda$ be the lattice of the PEL datum
(Definition~\ref{cv-model-toroidal-semiabelian}), and for an integer $k\ge1$ let $G_k$ be the finite group of
$\mathcal O_E$-linear automorphisms of $\Lambda/k\Lambda$. Let $y$ be a point of $X$ over an $F_c$-algebra $S$,
with abelian scheme $A_y$. A \emph{full level-$k$ structure} on $y$ is an $\mathcal O_E$-linear homomorphism
$e\colon\Lambda/k\Lambda\to A_y(S)$ such that for distinct $q,q'\in\Lambda/k\Lambda$ the points $e(q)$ and
$e(q')$ have distinct images under every homomorphism from $S$ to a field.
\end{definition}
\leanref{Development/AbelianVarietyTheory/Definitions/ThetaValues/IsNullVal.lean}{GlobalTheory.EtaForm.IsFullLev}
\leanref{Development/AbelianVarietyTheory/Definitions/ThetaValues/IsNullVal.lean}{GlobalTheory.EtaForm.LevGrp}

\begin{lemma}\label{cv-volume-twisting-levels}
In the setting of Definition~\ref{cv-volume-twisting-section}, let $k\ge1$, let $\Lambda$ and $G_k$ be as in
Definition~\ref{cv-volume-twisting-level-structures}, and let $y$ be a point of $X$ over an $F_c$-algebra $S$.
\begin{enumerate}
\item There is a faithfully flat homomorphism of $F_c$-algebras $S\to S'$ such that the point of $X$ over $S'$
  given by $y$ has a full level-$k$ structure.
\item If $e$ is a full level-$k$ structure on $y$ and $S\to T$ is a homomorphism of $F_c$-algebras, then
  $e$ followed by $A_y(S)\to A_y(T)$ is a full level-$k$ structure on the point over $T$, and $ke(q)=0$ for every
  $q$.
\item If $S$ is a local ring and $e$ and $e'$ are full level-$k$ structures on $y$, then every $p\in A_y(S)$ with
  $kp=0$ is $e(q)$ for some $q$, and there is $g\in G_k$ with $e'=e\circ g$.
\item If $S$ is an algebraically closed field, $y$ has a full level-$k$ structure.
\end{enumerate}
\end{lemma}
\leanref{Development/AbelianVarietyTheory/Proofs/CuspFormEvaluations/FullLinearLevelCoverBaseChangeTorsor.lean}{GlobalTheory.fullLevel_cover_baseChange_torsor}
These are properties of PEL moduli problems \cite{Kot92b,Lan13}. Over an algebraically closed field of
characteristic $0$ the $k$-torsion of $A_y$ is isomorphic to $\Lambda/k\Lambda$ as an $\mathcal O_E$-module,
which is (4). The $k$-torsion of $A_y$ is finite and \'etale over $S$, as $k$ is invertible in $F_c$, and the
full level-$k$ structures form a finite \'etale cover of $\operatorname{Spec}S$ that is a torsor under $G_k$;
this gives (1) and (3), and (2) is immediate.

\begin{definition}\label{cv-volume-twisting-null-values}
In the setting of Definition~\ref{cv-volume-twisting-section}, let $y=(y,v,\vartheta)$ be a framed point over
an $F_c$-algebra $S$, let $L$ be the symmetric line bundle of the principal polarization of $A_y$, rigidified
along the zero section, and let $p\in A_y(S)$ with $81p=0$. There are a unique isomorphism
$[9]^*L\cong L^{81}$ compatible with the rigidifications, and a unique symmetric element $U_p$ of the theta
group of $L^{81}$ over $S$ that lies over $p$ and satisfies $U_p^{81}=1$
(Definition~\ref{cv-uniformization-theta-rule-unique-translates-symmetric}): by
Lemma~\ref{cv-uniformization-theta-rule-unique-translates-theta} for $X$ and $N=9$, the pullback of the isomorphism of
its part (1) is one, two such differ by a unit of $S$ that is $1$ along the zero section, and part (2), for
$T=\operatorname{Spec}S$, gives $U_p$. Let $s_1$ be the section of
$L^{81}$ that corresponds to $[9]^*\vartheta$. The \emph{null value} $w_y(p)\in S$ is the value of the section
$U_ps_1$ of $L^{81}$ at the zero section, read in the rigidification.
\end{definition}
\leanref{Development/AbelianVarietyTheory/Definitions/ThetaValues/IsNullVal.lean}{GlobalTheory.EtaForm.IsNullVal}

\begin{lemma}\label{cv-volume-twisting-null-facts}
In the setting of Definition~\ref{cv-volume-twisting-section}, let $y=(y,v,\vartheta)$ be a framed point over an
$F_c$-algebra $S$ and $p\in A_y(S)$ with $81p=0$, and let $w_y(p)$ be the null value
(Definition~\ref{cv-volume-twisting-null-values}).
\begin{enumerate}
\item If $\psi\colon S\to T$ is a homomorphism of $F_c$-algebras, $z$ is the image of $y$ over $T$ and $p'$ the
  image of $p$, then $w_z(p')=\psi(w_y(p))$.
\item If $y'$ is a framed point over $S$ with the point of $y$, whose bases are obtained from those of $y$ by
  upper triangular matrices and whose generator is $t\vartheta$ with $t\in S^\times$, then
  $w_{y'}(p)=t\,w_y(p)$.
\item If $S$ is a field and $\vartheta$ does not vanish at $9p$, then $w_y(p)\ne0$.
\end{enumerate}
\end{lemma}
\leanref{Development/AbelianVarietyTheory/Proofs/CuspFormEvaluations/ThetaNullValueExistsUniqueNaturalLinear.lean}{GlobalTheory.nullValue_existsUnique_natural_linear}
\leanref{Development/AbelianVarietyTheory/Proofs/CuspFormEvaluations/ExistsShiftNullValNeZero.lean}{GlobalTheory.EtaBase.exists_nullVal_ne_zero}
\begin{proof}
(1) The normalized isomorphism and the symmetric element $U_p$ are unique, so they are carried to those over
$T$; so are $s_1$, the action of $U_p$ and the value at the zero section.

(2) The construction does not involve the bases, and $s_1$, $U_ps_1$ and its value are linear in $\vartheta$.

(3) The element $U_p$ is an isomorphism of $L^{81}$ with its translate by $p$, so the value of $U_ps_1$ at the zero
section is a unit times the value of $s_1$ at $p$. The value of $s_1$ at $p$ is a unit times the value of
$\vartheta$ at $9p$, which is not zero.
\end{proof}
% Lean: nullValue_existsUnique_natural_linear for N = 9 = 3^2 (odd): existence Model.exists_isNullVal (the section F of
% exists_isVtOver_and_existsUnique_isUUOver and its value GermD.exists_value at the rigidification rigT), uniqueness
% Model.isNullVal_unique (the normalized pullback isomorphism universal_normalised_mulByNat_pullback_iso and the unique
% symmetric lift universal_symLift_exists_unique_global), (1) NullValue.baseChange, (2) NullValue.triangular. In
% IsNullVal, IsVtOver takes t' = p and t = p^9, the twisted pullback s_1, and U with SymLift (over p, U^81 = 1,
% symmetric). (3) is the computation in exists_nullVal_ne_zero (sv_act, sv_mulPull, sv_of_isValueAt).

\begin{lemma}\label{cv-volume-twisting-eta}
In the setting of Definition~\ref{cv-volume-twisting-section}, let $\Lambda$ and $G=G_{81}$ be as in
Definition~\ref{cv-volume-twisting-level-structures}. Let $J$ be a finite set, $n\ge0$, $c\colon J\to\mathbf Z$,
and $x_{j,b}\in\Lambda/81\Lambda$ for $j\in J$ and $1\le b\le n$. For a framed point $y$ over an $F_c$-algebra $S$
and a full level-$81$ structure $e$ on $y$ put
\[
 E(y,e)=\prod_{g\in G}\sum_{j\in J}c_j\prod_{b=1}^nw_y\big(e(g(x_{j,b}))\big),
\]
with the null values of Definition~\ref{cv-volume-twisting-null-values}. Then there is a unique rational
section $f$ of $X$, of the weight whose entries $G_{\tau,i}$ are all $0$ and whose last entry is $n|G|$, such
that for every framed point $y$ over an $F_c$-algebra $S$, every faithfully flat homomorphism of
$F_c$-algebras $\varphi\colon S\to S'$ and every full level-$81$ structure $e$ on the image $y'$ of $y$ over $S'$,
$\varphi(f(y))=E(y',e)$.
\end{lemma}
\leanref{Development/AbelianVarietyTheory/Definitions/ThetaValues/IsNullVal.lean}{GlobalTheory.EtaForm.IsEtaRule}
\leanref{Development/AbelianVarietyTheory/Proofs/CuspFormEvaluations/EtaValAtExistsUniqueNatural.lean}{GlobalTheory.EtaForm.etaValAt_existsUnique_natural}
\leanref{Development/AbelianVarietyTheory/Proofs/CuspFormEvaluations/IsEtaRuleExistsUniqueMemSec.lean}{GlobalTheory.EtaForm.isEtaRule_existsUnique_mem_sec}
\begin{proof}
Independence of the level structure. Let $z$ be a framed point over an $F_c$-algebra $T$ with full level-$81$
structures $e$ and $e'$. Then $E(z,e)=E(z,e')$. It suffices to check this in the localization $T_P$ at every
maximal ideal $P$. By Lemma~\ref{cv-volume-twisting-levels}(2) and~(3), the images of $e$ and $e'$ over $T_P$
are full level-$81$ structures with $e'=e\circ g_0$ for some $g_0\in G$. By
Lemma~\ref{cv-volume-twisting-null-facts}(1), the images of $E(z,e)$ and $E(z,e')$ in $T_P$ are the expressions
$E$ for the images; and $E(z_P,e\circ g_0)=E(z_P,e)$, because $g\mapsto g_0g$ is a bijection of $G$.

Existence. Let $y$ be a framed point over $S$. By Lemma~\ref{cv-volume-twisting-levels}(1) there are a faithfully
flat homomorphism $\varphi\colon S\to B$ and a full level-$81$ structure $e$ on the image $y_B$ of $y$. The two
homomorphisms $\iota_1,\iota_2\colon B\to B\otimes_SB$ agree on $S$, so they carry $y_B$ to one framed point $z$,
and $e$ to two full level-$81$ structures on it (Lemma~\ref{cv-volume-twisting-levels}(2)). By
Lemma~\ref{cv-volume-twisting-null-facts}(1) and the first step,
$\iota_1(E(y_B,e))=E(z,\iota_1e)=E(z,\iota_2e)=\iota_2(E(y_B,e))$. By faithfully flat descent
\cite[Expos\'e~VIII, \S1]{Gro71} there is a unique $v\in S$ with $\varphi(v)=E(y_B,e)$; put $f(y)=v$.

Compatibility. Let $\psi\colon S\to T$ be a homomorphism, $z$ the image of $y$ over $T$, and let $\varphi'\colon
T\to T'$ and $e'$ be a faithfully flat homomorphism and a full level-$81$ structure on the image $z'$ of $z$. The
homomorphism $T'\to B\otimes_ST'$ is faithfully flat, hence injective, and over $B\otimes_ST'$ the images of
$e$ and $e'$ are two full level-$81$ structures on one framed point. By the first step and
Lemma~\ref{cv-volume-twisting-null-facts}(1), the images of $\varphi'(\psi(f(y)))$ and $E(z',e')$ in
$B\otimes_ST'$ agree, so $\varphi'(\psi(f(y)))=E(z',e')$. For $T=S$ and $\psi$ the identity this is the property
in the statement, which determines $f(y)$, as $\varphi'$ is injective; and with $\varphi'$ the homomorphism of
the existence step for $z$ it gives $f(z)=\psi(f(y))$.

Weight. Let $y'$ be a framed point over $S$ with the point of $y$, whose bases are obtained from those of $y$ by
upper triangular matrices and whose generator is $t\vartheta$. A full level-$81$ structure $e$ over $B$ serves for
the images of both. By Lemma~\ref{cv-volume-twisting-null-facts}(2) each null value is multiplied by
$\varphi(t)$, so $E(y'_B,e)=\varphi(t)^{n|G|}E(y_B,e)$, and $f(y')=t^{n|G|}f(y)$. This is the transformation
rule of the weight whose entries are $0$ and whose last entry is $n|G|$.
\end{proof}
% Lean: the rule is IsEtaRule (EtaValAt at each framed point over a test ring, with etaExpr); existence is
% exists_etaValAt (fullLinearLevel_cover_baseChange_torsor for the cover, existsUnique_preimage_of_faithfullyFlat for
% the descent), uniqueness and naturality map_eq_of_etaValAt (exists_faithfullyFlat_pushout), the independence
% etaExpr_eq_of_levels (Module.eq_of_localization_maximal, exists_levGrp and etaExpr_comp), and the weight
% ScalarWeight.etaValAt_frBor (etaExpr_scale), which gives membership in sec Xo c 0 (n |LevGrp Dw 81|).

\begin{lemma}\label{cv-volume-twisting-nonvanishing}
In the setting of Definition~\ref{cv-volume-twisting-section}, let $\sigma$ be a good embedding of $F_c$ and
$x_0$ a frame of a component of the uniformization at $\sigma$. Then there are an integer $m\ge0$ and a
rational section $f_0$ on $X$ of the weight whose entries $G_{\tau,i}$ are all $0$ and whose last entry is
$2m+2$, such that the reading of $f_0$ at $\sigma$ does not vanish at $x_0$.
\end{lemma}
\leanref{Development/AbelianVarietyTheory/Proofs/IwasawaDecomposition/UnfoldedIsSecZeroKeySection.lean}{GlobalTheory.CuspFrom.KeyHol.exists_sec_rdU_ne_zero}
\begin{proof}
Let $y=(y,v,\vartheta)$ be the framed point over $\mathbf C$ at $x_0$, and let $\Lambda$ and $G=G_{81}$ be as in
Definition~\ref{cv-volume-twisting-level-structures}. By Lemma~\ref{cv-volume-twisting-levels}(4), $y$ has a full
level-$81$ structure $e$ over $\mathbf C$. For $u\in\Lambda/3\Lambda$ fix a lift $\tilde u\in\Lambda$, and write
$[3\tilde u]$ for the class of $3\tilde u$ in $\Lambda/81\Lambda$.

For every $g\in G$ there is $u\in\Lambda/3\Lambda$ with $w_y\big(e(g[3\tilde u])\big)\ne0$. Indeed, by
Lemma~\ref{cv-volume-twisting-coset}, with $n=3$ and $a=0$, there is $z\in A_y(\mathbf C)$ with $3z=0$ at which
$\vartheta$ does not vanish. As $81z=0$, $z=e(q)$ for some $q$ by Lemma~\ref{cv-volume-twisting-levels}(3), and
$e(3q)=3z=0=e(0)$, so $3q=0$. Let $m\in\Lambda$ represent $g^{-1}(q)$. Then $3m\in81\Lambda$, and as $\Lambda$ is
torsion free, $m=27m'$ with $m'\in\Lambda$; let $u$ be the class of $m'$ in $\Lambda/3\Lambda$. As
$\tilde u-m'\in3\Lambda$, $27\tilde u-m\in81\Lambda$, so $g^{-1}(q)=[27\tilde u]$. The point
$p=e(g[3\tilde u])$ satisfies $9p=e(g[27\tilde u])=e(q)=z$, so $w_y(p)\ne0$ by
Lemma~\ref{cv-volume-twisting-null-facts}(3).

Let $J$ be the set of pairs $j=(j_1,j_2)$ of elements of $\Lambda/3\Lambda$, and put $x_{j,b}=[3\tilde j_b]$ for
$b=1,2$. Let $V$ be the space of complex functions on $G$, let $\ell_g\colon V\to\mathbf C$ be the evaluation at
$g$, and let $s_j\in V$ be the function
$g\mapsto w_y\big(e(g(x_{j,1}))\big)\,w_y\big(e(g(x_{j,2}))\big)$. By the preceding paragraph, for every $g$ there is
$u$ with $\ell_g(s_{(u,u)})\ne0$. By Lemma~\ref{cv-volume-twisting-combination} there are integers
$0\le c_j\le|G|$ with $\sum_jc_js_j(g)\ne0$ for every $g\in G$. Let $f_0$ be the rational section of
Lemma~\ref{cv-volume-twisting-eta} for $J$, $n=2$, these $c_j$ and the $x_{j,b}$. Its weight has entries $0$ and
last entry $2|G|=2m+2$, with $m=|G|-1\ge0$. With $\varphi$ the identity of $\mathbf C$,
\[
 f_0(y)=E(y,e)=\prod_{g\in G}\sum_{j\in J}c_js_j(g)\ne0,
\]
and the reading of $f_0$ at $x_0$ is its value $f_0(y)$ at the framed point at $x_0$.
\end{proof}
% Lean: exists_sec_rdU_ne_zero with R = 3: GlobalTheory.EtaForm.exists_small_coeff_etaValAt_ne_zero, with one point
% (Nc = 1), n = 2 and x_0 = 0, gives the coefficients cf <= |LevGrp Dw 3^4| (the level structure from
% fullLevel_cover_baseChange_torsor at an algebraically closed field, the null values from
% nullValue_existsUnique_natural_linear, the nonvanishing exists_shift_nullVal_ne_zero, whose u comes from
% EtaBase.exists_nullVal_ne_zero and EtaBase.exists_shift, and the coefficients exists_sum_natCast_smul_apply_ne_zero);
% GlobalTheory.EtaForm.isEtaRule_existsUnique_mem_sec gives the section, in sec Xo c 0 (2 |LevGrp Dw 3^4|), with
% x_{j,b} = shiftFam 3 0 j b, the class of 3 qRep(j b). The reading at x_0 is the value at the framed point there
% (etaValAt_evR_of_fromTest, with AllTest).

\begin{lemma}\label{cv-volume-twisting-holomorphic}
In the setting of Definition~\ref{cv-volume-twisting-section}, suppose that $p+q\ge1$ and that $X$ has a key
section, and let $\sigma$ be a good embedding at which the uniformization satisfies the compatibility
conditions of the moduli problem. Let $G$ be a weight with even last entry whose entries $G_{\tau,i}$ do not
increase with $i$, for every embedding $\tau$ of $E$, and write $\lambda(G)$ for the weight that $G$ gives on
the components of the uniformization at $\sigma$, by the formula of Situation~\ref{cv-cusp-setting} with the
entries of $G$ in place of $\alpha$. Then for every rational section $f$ of weight $G$ on $X$ and every
component of the uniformization at $\sigma$, the reading $Q$ of $f$ is a form of weight $\lambda(G)$: it is
holomorphic on the frames of the component, $Q(xb)=\prod_{\tau,i}b_{\tau,ii}^{\lambda(G)_{\tau,i}}Q(x)$ for
every frame $x$ and every family $b$ of invertible upper triangular matrices, and
$Q(\gamma x)=\chi_{\lambda(G)}(\gamma)Q(x)$ for every $\gamma$ in the arithmetic group of the component.
\end{lemma}
\leanref{Development/AbelianVarietyTheory/Proofs/DivisorProofs/IsSecReadingDecreasingIntegerExponentsStatement.lean}{GlobalTheory.T.isSec_reading_decreasing_integer_exponents}
\leanref{Development/AbelianVarietyTheory/Proofs/DivisorProofs/IsSecReadingDecreasingIntegerExponents.lean}{GlobalTheory.FlatTw.Hol.isSec_modelWt_rdU}
\begin{proof}
For integers $a$ and $n$ write $(a;n)$ for the weight whose entries $G_{\tau,i}$ all equal $a$ and whose last
entry is $n$. Let $\rho$ be a key section of $X$, of weight $\kappa(0,-1)=(1;-2)$. One of the compatibility
conditions of the moduli problem is that the reading at $\sigma$ of a rational section of a dominant weight
$\kappa$ is a form of weight $\lambda(\kappa)$ on every component, as in
Lemma~\ref{cv-uniformization-frames}(4) to~(6).
The map $\kappa\mapsto\lambda(\kappa)$ is additive on weights with even last entries, and the reading of a
product of rational sections is the product of their readings.

Fix a component and a frame $x_0$ of it. By Lemma~\ref{cv-volume-twisting-nonvanishing} there are $m\ge0$ and
a rational section $f_0$ of weight $(0;2m+2)$ whose reading does not vanish at $x_0$. Write the last entry of
$G$ as $2g_2$, put $N=|2g_2|$, let $B\ge0$ be an integer with $G_{\tau,i}\ge-B$ for all $\tau$ and $i$, and put
$t=(B+1)m+(N+1)(m+1)$ and
\[
 g=\rho^{B+1}f_0^{B+N+2}.
\]
The weight of $g$ is $(B+1)(1;-2)+(B+N+2)(0;2m+2)=(B+1;2t)$, as
$2(m+1)(B+N+2)-2(B+1)=2t$. The weight of $fg$ is $G+(B+1;2t)$: its entries $G_{\tau,i}+B+1\ge1$ do not
increase with $i$, and its last entry is $2(g_2+t)$, with $g_2+t\ge0$ since $t\ge N+1$. So $g$ and $fg$ are
rational sections of dominant weights, and their readings $P'$ and $P$ are forms of weights $\mu=\lambda(B+1;2t)$
and $\lambda(G)+\mu$. The reading $Q$ of $f$ satisfies $QP'=P$ on the frames of the component. The product of
the readings of $\rho$ and $\rho^{-1}$ at $x_0$ is the reading of $1$, which does not vanish at $x_0$, since its
product with the reading of $f_0$ is the reading of $f_0$; so the reading of $\rho$ does not vanish at $x_0$,
and $P'(x_0)\ne0$.

Hence $Q=P/P'$ is holomorphic near $x_0$. For a family $b$ of invertible upper triangular matrices and a
weight $\lambda$ on the components write $b^\lambda=\prod_{\tau,i}b_{\tau,ii}^{\lambda_{\tau,i}}$. Then
$Q(x_0b)\,b^\mu P'(x_0)=P(x_0b)=b^{\lambda(G)}b^\mu Q(x_0)P'(x_0)$, and cancelling the nonzero factor
$b^\mu P'(x_0)$ gives $Q(x_0b)=b^{\lambda(G)}Q(x_0)$. The same computation with an
element $\gamma$ of the arithmetic group gives $Q(\gamma x_0)=\chi_{\lambda(G)}(\gamma)Q(x_0)$, as
$\chi_{\lambda(G)+\mu}=\chi_{\lambda(G)}\chi_\mu$ on that group and $\chi_\mu(\gamma)$, of finite order, is not
zero. As $x_0$ was arbitrary, $Q$ is a form of weight $\lambda(G)$.
\end{proof}
% Lean: isSec_modelWt_rdU applies isSec_of_local_quot at each x_0 with la = laU U ag (2 tg), P = rdU (f g) and
% P' = rdU g, where g = rho^(B+1) f_0^(B+1+N+1), ag = B + 1 and tg = (B+1) m + (N+1)(m+1); the readings of
% sections of dominant weight are sections by the field hU.2.2.2.2.1 of IsUnf, and modelWt_add, modelWt_wtZ and
% modelWt_congr give the weights. The nonvanishing of the reading of rho is hrhone, through h1ne.

\begin{lemma}\label{cv-volume-twisting-cusp}
In the setting of Definition~\ref{cv-volume-twisting-section}, suppose that $p+q\ge1$, that the
uniformizations at the good embeddings satisfy the compatibility conditions of the moduli problem, and
that $X$ has a key section. Let $Z$ be a connected component of $X$, let $G$ be a weight with even last
entry whose entries $G_{\tau,i}$ do not increase with $i$, for every embedding $\tau$ of $E$, and let $s$ be
a rational section of weight $G$ supported on $Z$, that is, vanishing at every framed point whose point
does not lie on $Z$, such that $|s|_\xi$ is bounded on the frames of $Z$ at
every good embedding. Then $fs$ is a cusp form on $Z$ of weight $G'+G$ for every weight $G'$ with even last
entry and every cusp form $f$ on $Z$ of weight $G'$.
\end{lemma}
\leanref{Development/AbelianVarietyTheory/Proofs/DivisorProofs/MulMemCuspFormsOfBoundedSectionComponentStatement.lean}{GlobalTheory.T.mul_mem_cuspForms_of_boundedSection_component}
\begin{proof}
Let $G'$ be a weight with even last entry and $f$ a cusp form on $Z$ of weight $G'$. The product $fs$ is a
rational section of weight $G'+G$, and it is supported on $Z$, as its value at a framed point is the product
of the values of $f$ and $s$. Let $\sigma$ be a good embedding, consider a component of the uniformization at
$\sigma$, and let $P$ and $Q$ be the readings of $f$ and $s$ on it, so that $PQ$ is the reading of $fs$. As the
last entries of $G'$ and $G$ are even, $\lambda(G'+G)=\lambda(G')+\lambda(G)$, in the notation of
Lemma~\ref{cv-volume-twisting-holomorphic}. The reading $P$ is a cusp form of weight $\lambda(G')$. By
Lemma~\ref{cv-volume-twisting-holomorphic}, $Q$ is a form of weight $\lambda(G)$. Its pointwise norm is
bounded: at a frame of $Z$ it is $(2\pi\xi)^{|G|/2}|s|_\xi$, which is bounded by hypothesis, and at a frame
whose framed point does not lie on $Z$ the reading $Q$ vanishes, as $s$ is supported on $Z$. The product of a
cusp form and a form with bounded pointwise norm is a cusp form of the sum of the weights, by the argument for
products in the proof of Lemma~\ref{cv-cover-level-automorphisms}(4) (Lemma~\ref{is-onto-share-sections}(1)).
So $PQ$ is a cusp form of weight $\lambda(G'+G)$ at every component and every good embedding, and $fs$ is a
cusp form on $Z$ of weight $G'+G$.
\end{proof}
% Lean: mul_mem_cuspOn_of_bdd is mul_mem_cuspOn_of_isSec (ExistsUnramPtsDefinite.lean) with the readings of s
% given by GlobalTheory.T.isSec_reading_decreasing_integer_exponents; the product is mul_mem_secZ and
% mul_mem_suppOn, the weights modelWt_add, the cusp property the field 2.1 of siegelSpaceAlgebra_holds, and the
% bound on the frames off Z rdU_eq_zero.

\begin{lemma}\label{cv-volume-twisting-flat}
In the setting of Definition~\ref{cv-volume-twisting-section}, suppose that $p+q\ge1$, that the
uniformizations at the good embeddings satisfy the compatibility conditions of the moduli problem, that
$X$ has a key section, and that there is a real number $c$ such that on every connected component of $X$,
for every $\tau\in\Phi$, flat sections of type $\tau$ exist and all have degree $-c$
(Definition~\ref{cv-constant-sections}, with its prime $d$). Let $B$ be a finite set of primes such that
for every connected component $Z'$ of $X$ and every finite place $w$ of $F_c$ over no prime of $B$ there
is an integral point of $Z'$ above $w$ at which the valuation of every rational section is an integer or
$+\infty$. Let $Z$ be a connected component with frames at every good embedding and integral points above
every finite place, and let $w\colon\Phi\to\mathbf Z$. Then there is a twisting section $t$ of weight
$\kappa(w,0)$ on $Z$ whose function $l$ takes integer values at the places over no prime of $B$, with
\[
 \deg t=c\sum_\tau w_\tau-q\big(\log\operatorname{rd}_E-\log d\big)\sum_\tau w_\tau
 \ge c\sum_\tau w_\tau-q\log\operatorname{rd}_E\sum_\tau|w_\tau|.
\]
\end{lemma}
\leanref{Development/AbelianVarietyTheory/Proofs/DivisorProofs/ExistsTwSecOfConst.lean}{GlobalTheory.FlatTwisting.twSec_of_flat}
\begin{proof}
For $\tau\in\Phi$ choose a flat section $s_\tau$ of type $\tau$ on $Z$, of weight $\lambda_\tau$ and defect
$\delta_\tau$ (Definition~\ref{cv-constant-sections}). Products and integer powers of constant sections on
$Z$ are constant sections: weights, defects and the functions $\rho$ add, the numbers $c_\sigma$
multiply, and so degrees add. Hence $s=1_Z\prod_\tau s_\tau^{-w_\tau}$, with inverse
$s'=1_Z\prod_\tau s_\tau'^{\,-w_\tau}$, is a constant section on $Z$ of weight $-\sum_\tau w_\tau\lambda_\tau$,
which is the weight with value $w_\tau$ at $\tau$, and of defect $\delta=-\sum_\tau w_\tau\delta_\tau$, with
$c_\sigma(s)=\prod_\tau c_\sigma(s_\tau)^{-w_\tau}$, $\rho_s=-\sum_\tau w_\tau\rho_{s_\tau}$ and
$\deg s=-\sum_\tau w_\tau\deg s_\tau=c\sum_\tau w_\tau$. As a rational section, $s$ has weight
$\kappa(w,0)$, and $s'$ has weight $\kappa(-w,0)=-\kappa(w,0)$.

Let $t$ consist of $s$, $s'$, the numbers $c_\sigma(s)$ and the functions $l=h=\delta+\rho_s$, which vanish
at almost all places. By Definition~\ref{cv-constant-sections}, $|s|_\xi=c_\sigma(s)$ at the frames of $Z$
at $\sigma^g$, and $v(s)=\delta(w)+\rho_s(w)$ at every integral point of $Z$ above $w$. The weights
$\kappa(\pm w,0)$ have last entry $0$ and entries $\pm w_\tau$ at $\tau$ and $\mp w_\tau$ at $\bar\tau$,
constant in the index; $|s|_\xi$ and $|s'|_\xi=c_\sigma(s)^{-1}$ are constant on the frames of $Z$; and $s$
and $s'$ are supported on $Z$, as $1_Zs=s$, $1_Zs'=s'$ and $1_Z$ vanishes off $Z$. So
Lemma~\ref{cv-volume-twisting-cusp}, applied to $s$ and to $s'$, shows that $fs$ is a cusp form of weight
$G'+\kappa(w,0)$ for every cusp form $f$ on $Z$ of weight $G'$, and that $fs'$ is a cusp form of weight
$G'+\kappa(w,0)-\kappa(w,0)=G'$ for every cusp form $f$ on $Z$ of weight $G'+\kappa(w,0)$. So $t$ is a
twisting section of weight $\kappa(w,0)$ on $Z$.

Let $w$ be a finite place over no prime of $B$, and choose an integral point of $Z$ above $w$ at which all
valuations are integers or $+\infty$. There $v(s)=\delta(w)+\rho_s(w)$ is a real number, hence an
integer; so $l(w)$ is an integer.

By the definition of the degree,
$\deg t=\deg s+\frac1{[F_c:\mathbf Q]}\sum_w\delta(w)\log Nw$. The defect sums to
$\frac1{[F_c:\mathbf Q]}\sum_w\delta_\tau(w)\log Nw=q(\log\operatorname{rd}_E-\log d)$ for every $\tau$
(the identity after Definition~\ref{cv-constant-sections}), so the second term is
$-q(\log\operatorname{rd}_E-\log d)\sum_\tau w_\tau$. This is the equality. For the inequality, $E$
contains $\mathbf Q(\sqrt{-d})$, whose root discriminant is $\sqrt d$, and the root discriminant does not
decrease in finite extensions, so $\frac12\log d\le\log\operatorname{rd}_E$. As $\log d\ge0$, this gives
$|\log\operatorname{rd}_E-\log d|\le\log\operatorname{rd}_E$, hence
$(\log\operatorname{rd}_E-\log d)\sum_\tau w_\tau\le\log\operatorname{rd}_E\sum_\tau|w_\tau|$.
\end{proof}
% Lean: twSec_of_flat, second part: the constant section is FlatSec.shift (ConstSec.prod and ConstSec.zpow,
% with parWt_eq_sum), of degree sum_tau w_tau c1 by FlatExists.deg_shift; the twisting section is
% exists_twSec_of_const = TwSec.ofConst, with lo = hi = delta + rho_vw, the multiplication property
% mul_sig_mem_cuspOn (BddCusp.mul_mem_cuspOn_of_bdd with wtOf_parWt_anti, and ConstSec.symm for the inverse),
% and IntOff from intOff_of_unramPts. The degree is lineDeg_mul with finsum_sum_flDefect (sum_flDefect_vw),
% and the bound uses GlobalTheory.T.abs_logRootDisc_sub_log_hostPrime_le = abs_lrd_sub_log_d_le_lrd
% (rootDiscr_and_different_quadField and rootDiscr_le_of_intermediateField).

\begin{lemma}\label{cv-volume-twisting-colength}
Let $E$ be a number field with different $\mathfrak D_E$, let $T$ be an integrally closed domain whose
fraction field $F$ has characteristic $0$, and let $\tau_i\colon\mathcal O_E\to T$, for $i$ in a finite set
$I$ with $|I|=[E:\mathbf Q]$, be distinct ring homomorphisms. Let $m_i\ge0$ be integers, $g=\sum_im_i$, let
$A\colon\mathcal O_E\to M_g(T)$ be a ring homomorphism, and for each $i$ let $u_{i,1},\dots,u_{i,m_i}\in T^g$
satisfy $A(a)u_{i,k}=\tau_i(a)u_{i,k}$ for all $a\in\mathcal O_E$, such that every $x\in T^g$ with
$A(a)x=\tau_i(a)x$ for all $a$ is a combination of $u_{i,1},\dots,u_{i,m_i}$ with coefficients in $T$. Let
$U$ be the $g\times g$ matrix whose rows are the vectors $u_{i,k}$, in some order. Then $\det U$ divides
$\prod_i\tau_i(\delta_i)^{m_i}$ in $T$ for all $\delta_i\in\mathfrak D_E$.
\end{lemma}
\leanref{Development/AbelianVarietyTheory/Proofs/DivisorProofs/ColengthDvdOfNeZero.lean}{GlobalTheory.ColengthDifferent.colength_dvd}
\begin{proof}
If $\delta_i=0$ for some $i$ with $m_i>0$, the product is $0$. Otherwise, replacing each $\delta_i=0$, which
has $m_i=0$, by a nonzero element of $\mathfrak D_E$ does not change the product; so suppose that every
$\delta_i$ is nonzero. Let $\sigma_i\colon E\to F$ extend $\tau_i$. These are $[E:\mathbf Q]$ distinct
embeddings, so $\sum_i\sigma_i(x)=\operatorname{Tr}_{E/\mathbf Q}(x)$ for $x\in E$.

Let $\omega_1,\dots,\omega_N$ be a $\mathbf Z$-basis of $\mathcal O_E$ and $\omega_1^\vee,\dots,\omega_N^\vee$
the dual basis of $E$ for the trace form, so that $y=\sum_k\operatorname{Tr}(y\omega_k^\vee)\omega_k
=\sum_k\operatorname{Tr}(y\omega_k)\omega_k^\vee$ for $y\in E$. The different is the inverse of the module
dual to $\mathcal O_E$, so $c_{i,k}=\delta_i\omega_k^\vee$ lies in $\mathcal O_E$. Put
$P_i=\sum_k\tau_i(c_{i,k})A(\omega_k)$, a matrix over $T$.

First, $A(a)P_i=\tau_i(a)P_i$ for $a\in\mathcal O_E$. Indeed
$A(a)P_i=\sum_k\tau_i(c_{i,k})A(a\omega_k)=\sum_l\big(\sum_k\operatorname{Tr}(a\omega_k\omega_l^\vee)\tau_i(c_{i,k})\big)A(\omega_l)$,
and $\sum_k\operatorname{Tr}(a\omega_l^\vee\omega_k)c_{i,k}=\delta_ia\omega_l^\vee=ac_{i,l}$; so
$A(a)P_i=\sum_l\tau_i(ac_{i,l})A(\omega_l)=\tau_i(a)P_i$. Hence the $j$-th column of $P_i$ is
$\sum_kG_{j,(i,k)}u_{i,k}$ for some $G_{j,(i,k)}\in T$; let $G$ be the $g\times g$ matrix of these
coefficients, with columns indexed like the rows of $U$.

Second, over $F$,
\[
 \sum_i\sigma_i(\delta_i)^{-1}P_i=\sum_k\Big(\sum_i\sigma_i(\omega_k^\vee)\Big)A(\omega_k)
 =A\Big(\sum_k\operatorname{Tr}(\omega_k^\vee)\omega_k\Big)=A(1)=1.
\]
Let $D$ be the diagonal matrix whose entry at the index $(i,k)$ is $\sigma_i(\delta_i)^{-1}$. The
$(j,j')$ entry of $GDU$ is $\sum_i\sigma_i(\delta_i)^{-1}\sum_kG_{j,(i,k)}(u_{i,k})_{j'}$, the $(j',j)$
entry of $\sum_i\sigma_i(\delta_i)^{-1}P_i$; so $GDU=1$. Taking determinants,
$\det G\det U=\prod_i\tau_i(\delta_i)^{m_i}$, and $\det G\in T$.
\end{proof}
% Lean: colength_dvd reduces to colength_dvd_of_ne_zero; c_{i,k} is exists_mul_traceDual, the eigen
% property of P_i is mul_proj (mul_dual_expand), the identity sum_proj_eq_one uses sum_eq_trace, and hmat is
% G D U = 1.

\begin{lemma}\label{cv-volume-twisting-key-unit}
Let $\kappa$ be the rule of Theorem~\ref{cv-definite-key-formula}. Let $S$ be a noetherian integrally closed
domain of characteristic $0$, and let $A$ be a principally polarized abelian scheme over $S$ with a symmetric
line bundle defining the polarization and rigidified along the zero section, a generator $\theta$ of its
sections over $S$ and a basis $v$ over $S$ of its Lie algebra. Then there is a unit $u\in S^\times$ such that
for every homomorphism $f$ from $S$ to a field $k$ of characteristic $0$, $\kappa$ of the data over $k$ obtained
by $f$ is $f(u)$.
\end{lemma}
\leanref{Development/AbelianVarietyTheory/Definitions/KeyFormula/KeyRootOfKeyFormula.lean}{GlobalTheory.KeyCl.FamUnit}
\leanref{Development/AbelianVarietyTheory/Proofs/ThetaLawProofs/FamUnit.lean}{GlobalTheory.KeyFamUnit.famUnit}
\begin{proof}
Let $K$ be the fraction field of $S$ and $c=\kappa(A_K,\theta,v)\in K^\times$. For every prime $P$ of height one,
$S_P$ is a discrete valuation ring with fraction field $K$, and the data extend over it; by
Theorem~\ref{cv-definite-key-formula}(3), $c$ is a unit of $S_P$. As $S$ is the intersection of the $S_P$
over the primes $P$ of height one \cite[Theorem~11.5]{Mat86}, $c$ and $c^{-1}$ lie in $S$; let $u=c$.

Let $f\colon S\to k$ be a homomorphism to a field of characteristic $0$, with kernel $P$. If $P=0$, $f$ extends to
$K$, and the claim is Theorem~\ref{cv-definite-key-formula}(2). Otherwise there are a complete discrete
valuation ring $O$ and an injective homomorphism $\varphi\colon S\to O$ whose composite with the residue map
has kernel $P$ \cite[7.1.7]{Gro61b}. The residue field of $O$ has characteristic $0$, since the integers lie
outside $P$. By Theorem~\ref{cv-definite-key-formula}(3), $\kappa$ of the data over $O$ is a unit $u'$ of $O$
whose residue is $\kappa$ of the data over the residue field; by the case $P=0$ applied to $\varphi$, followed by
the inclusion of $O$ in its fraction field, $u'=\varphi(u)$. So for the composite $f'$ of $\varphi$ with the residue
map, $\kappa$ of the data over the residue field is $f'(u)$. Both $f$ and $f'$ have kernel $P$, so they factor
through injective homomorphisms from the fraction field $k_P$ of $S/P$. By
Theorem~\ref{cv-definite-key-formula}(2), applied to $k_P$ and the residue field, $\kappa$ of the data over $k_P$
is the image of $u$; applied to $k_P$ and $k$, $\kappa$ of the data over $k$ is $f(u)$.
\end{proof}
% Lean: famUnit from Funct and UnitSp (parts (2) and (3) of the key formula): KeySmoothBase.exists_unit_normal
% (normal_domain_height_one for the height-one localizations and the intersection, exists_unit_dvr for each) and
% KeySpec.isValue_of_generic (exists_complete_discreteValuationRing_of_prime_lt, valueAt_of_point).

\begin{lemma}\label{cv-volume-twisting-key-theta}
In Situation~\ref{cv-uniformization-situation}, let $\kappa$ be the rule of Theorem~\ref{cv-definite-key-formula},
and let $(y,v,w)$ be a normalized framed point at a frame $x$, with number $\nu$, $\nu^2=N_\xi\det h_x$. Then
\[
 \nu^8\kappa(A_y,w,v)=\big((2\pi i)^g\det h_x\big)^4,\qquad\text{so}\qquad
 \kappa(A_y,w,v)=\Big(\frac{(2\pi i)^g}{N_\xi}\Big)^4.
\]
\end{lemma}
\leanref{Development/AbelianVarietyTheory/Proofs/IwasawaDecomposition/ThetaKeyRuleValueAtThetaPointStatement.lean}{GlobalTheory.T.theta_key_rule_value_at_theta_point}
\leanref{Development/AbelianVarietyTheory/Proofs/IwasawaDecomposition/ThetaKeyRuleValueAtThetaPoint.lean}{GlobalTheory.KeyFourth.value_at_thetaPt}
This is a computation with parts (1) and (4) of Theorem~\ref{cv-definite-key-formula}. In the coordinates
$z'=h_xz$ the abelian variety is $\mathbf C^g/(\mathbf Z^g+Z_x\mathbf Z^g)$, the section $w$ is $\nu$ times the
theta function $\vartheta[c_i](Z_x,z')$ with a characteristic in $\{0,\frac12\}^{2g}$, and the basis $v$ of unit
vectors is obtained from the basis $(2\pi i)^{-1}e_i$ of part (4) by the matrix $(2\pi i)h_x^{\mathsf T}$, of
determinant $(2\pi i)^g\det h_x$. Part (4) gives the value $1$ for the theta function and that basis, and part
(1), with $a=\nu$, gives the first identity; the second follows from $\nu^2=N_\xi\det h_x$.

\subsection{Roots of theta values and gluing}\label{subsection-cv-theta-roots-gluing}

\begin{lemma}\label{cv-volume-key-fixed-field}
Let $K$ be a perfect field and $L$ an algebraically closed field containing $K$. An element of $L$ fixed by
every automorphism of $L$ over $K$ lies in $K$. In particular, for a subfield $K$ of $\mathbf C$, an element
of $\mathbf C$ fixed by every automorphism of $\mathbf C$ over $K$ lies in $K$.
\end{lemma}
\leanref{Development/AbelianVarietyTheory/Proofs/HomogeneousIdeals/ExistsRingEquivOfIsTranscendenceBasis.lean}{GlobalTheory.GaloisDescent.mem_range_algebraMap_of_forall_algEquiv}
\leanref{Development/AbelianVarietyTheory/Proofs/HomogeneousIdeals/ExistsRingEquivOfIsTranscendenceBasis.lean}{GlobalTheory.GaloisDescent.complex_mem_range_of_forall_algEquiv}
\begin{proof}
First, let $F\subset L$ and $F'\subset L'$ be fields with $L$ and $L'$ algebraically closed, let $s$ be a
transcendence basis of $L$ over $F$ and $s'$ one of $L'$ over $F'$, and let $e\colon F[s]\to F'[s']$ be an
isomorphism of the polynomial rings. Then $e$ extends to an isomorphism $L\to L'$: it extends to the fields
of fractions $F(s)\to F'(s')$, of which $L$ and $L'$ are algebraic closures
\cite[Chapter~V, Theorem~2.8]{Lan02}.

Let $x\in L$ be fixed by every automorphism of $L$ over $K$. Suppose first that $x$ is algebraic over $K$, and
let $\bar K$ be the algebraic closure of $K$ in $L$. Let $f$ be an automorphism of $\bar K$ over $K$ and $s$ a
transcendence basis of $L$ over $\bar K$. The automorphism of $\bar K[s]$ that is $f$ on the coefficients and
fixes $s$ extends to an automorphism of $L$, which is the identity on $K$, so it fixes $x$; hence $f(x)=x$.
As $K$ is perfect, the minimal polynomial of $x$ over $K$ is separable. If its degree were at least $2$, it
would have a root $x'\ne x$ in $\bar K$, and the isomorphism $K(x)\to K(x')$ over $K$ with $x\mapsto x'$
would extend to an automorphism $f$ of $\bar K$ with $f(x)\ne x$ \cite[Chapter~V, Theorem~2.8]{Lan02}. So
$x\in K$.

Suppose now that $x$ is transcendental over $K$, and extend $\{x\}$ to a transcendence basis $t$ of $L$ over
$K$ \cite[Chapter~VIII, \S1]{Lan02}. The automorphism of $K[t]$ over $K$ with $x\mapsto x+1$ that fixes the
other elements of $t$ extends to an automorphism $\sigma$ of $L$, which is the identity on $K$, with
$\sigma(x)=x+1\ne x$. This is impossible.

A subfield of $\mathbf C$ has characteristic $0$, so it is perfect, and $\mathbf C$ is algebraically closed.
\end{proof}
% Lean: exists_ringEquiv_of_isTranscendenceBasis (the extension through IsAlgClosure.equivOfEquiv); the
% algebraic case uses InfiniteGalois.mem_range_algebraMap_iff_fixed on algebraicClosure K L, which the text
% replaces by the minimal-polynomial argument; the transcendental case uses exists_isTranscendenceBasis_superset.

\begin{lemma}\label{cv-volume-key-descent-spaces}
Let $L$ be a field, $K\subset L$ a subfield and $G$ a set of endomorphisms of $L$ over $K$ such that every
$x\in L$ with $g(x)=x$ for all $g\in G$ lies in $K$. Let $V$ be a $K$-vector space and $W$ an $L$-subspace of
$L\otimes_KV$. Then $(g\otimes1)(W)\subset W$ for every $g\in G$ if and only if $W=L\otimes_KN$ for a
$K$-subspace $N$ of $V$; and $N$ is then unique, the set of $v\in V$ with $1\otimes v\in W$. In particular
this holds for a subfield $K$ of $\mathbf C$, $L=\mathbf C$ and the group $G$ of automorphisms of
$\mathbf C$ over $K$.
\end{lemma}
\leanref{Development/AbelianVarietyTheory/Proofs/HomogeneousIdeals/ExistsOneTmulOfForallFixed.lean}{GlobalTheory.GaloisDescent.stable_iff_exists_baseChange}
\leanref{Development/AbelianVarietyTheory/Proofs/HomogeneousIdeals/FinrankFixedSubmoduleComplexification.lean}{GlobalTheory.GaloisDescent.complex_stable_iff_exists_baseChange}
\begin{proof}
If $W=L\otimes_KN$, then $(g\otimes1)(a\otimes v)=g(a)\otimes v$ lies in $W$ for $a\in L$ and $v\in N$, so $W$
is stable. Conversely, suppose that $W$ is stable, and let $N$ be the set of $v$ with $1\otimes v\in W$. Fix a
basis $(b_i)$ of $V$ over $K$. The $1\otimes b_i$ form a basis of $L\otimes_KV$ over $L$; for $w$ in
$L\otimes_KV$ let $w_i$ be its coordinates and $\operatorname{supp}w$ the set of $i$ with $w_i\ne0$. The
coordinates of $(g\otimes1)(w)$ are the $g(w_i)$.

(a) Every nonzero $w\in W$ has a $u\in W$ with coordinates in $K$, with $\operatorname{supp}u\subset
\operatorname{supp}w$ and with $u_i=1$ for some $i$. We argue by induction on the size of
$\operatorname{supp}w$. Choose $i$ with $w_i\ne0$ and put $w'=w_i^{-1}w$, so that $w'_i=1$. If
$(g\otimes1)(w')=w'$ for every $g\in G$, the coordinates of $w'$ are fixed by $G$, hence lie in $K$, and
$u=w'$ will do. Otherwise $w''=(g\otimes1)(w')-w'$ is a nonzero element of $W$ for some $g$, with coordinates
$g(w'_k)-w'_k$. These vanish for $k=i$, as $g(1)=1$, and for $k\notin\operatorname{supp}w$. So
$\operatorname{supp}w''\subset\operatorname{supp}w\smallsetminus\{i\}$, and the induction hypothesis for
$w''$ gives $u$.

(b) $W$ is spanned over $L$ by its elements with coordinates in $K$. We argue by induction on the size of
$\operatorname{supp}w$ for $w\in W$. For $w\ne0$ take $u$ and $i$ as in (a); then $i\in\operatorname{supp}w$,
and $w-w_iu\in W$ has support in $\operatorname{supp}w\smallsetminus\{i\}$. By induction it lies in the span,
and so does $w$.

An element of $L\otimes_KV$ with coordinates in $K$ is $1\otimes v$ with $v\in V$, and it lies in $W$ if and
only if $v\in N$. By (b), $W\subset L\otimes_KN$; and $L\otimes_KN\subset W$, as $W$ is an $L$-subspace
containing $1\otimes N$. For the uniqueness let $N'$ be a $K$-subspace with $L\otimes_KN'=W$, and let
$\lambda\colon L\to K$ be a $K$-linear map with $\lambda(1)=1$. The map $\lambda\otimes1\colon L\otimes_KV\to V$
maps $a\otimes v$ to $\lambda(a)v$, hence $L\otimes_KN'$ into $N'$. So $1\otimes v\in W$ gives
$v=(\lambda\otimes1)(1\otimes v)\in N'$, and $N\subset N'$; the inclusion $N'\subset N$ is clear. The last
assertion follows from Lemma~\ref{cv-volume-key-fixed-field}.
\end{proof}
% Lean: baseChange_rationalPart with span_image_comap_eq; (a) is exists_rational_support_subset and (b)
% span_rational_eq, with repr_apply_of_semilinear for the coordinates; the uniqueness is
% rationalPart_baseChange (comap_span_image, with Module.Projective.exists_dual_eq_one for lambda).

\begin{lemma}\label{cv-volume-key-roots-descent}
Let $K\subset\mathbf C$ be a subfield and $A$ a $K$-algebra such that $\mathbf C\otimes_KA$ is a domain. Let
$n\ge1$, $w\in\mathbf C\otimes_KA$, $c\in\mathbf C^\times$ and $u\in A$ nonzero with $w^n=c\otimes u$. Then
$u=ab^n$ with $a\in K^\times$ and $b\in A$.
\end{lemma}
\leanref{Development/AbelianVarietyTheory/Proofs/IwasawaDecomposition/ExistsBaseFieldMultipleOfNthRootComplexificationStatement.lean}{GlobalTheory.T.exists_base_field_multiple_of_nth_root_complexification}
\leanref{Development/AbelianVarietyTheory/Proofs/IwasawaDecomposition/ExistsBaseFieldMultipleOfNthRootComplexification.lean}{GlobalTheory.KeyFourth.exists_eq_smul_pow_of_pow_eq_tmul}
\begin{proof}
Let $\lambda\colon\mathbf C\to K$ be a $K$-linear map with $\lambda(1)=1$, and let
$T=\lambda\otimes1\colon\mathbf C\otimes_KA\to A$, so that $T(z\otimes a)=\lambda(z)a$. If $z\ne0$ and $a\ne0$
then $z\otimes a\ne0$, since $T\big(z^{-1}(z\otimes a)\big)=T(1\otimes a)=a$. In particular $c\otimes u\ne0$, so
$w\ne0$.

Two elements $x$ and $y$ of a domain containing $\mathbf C$ with $x^n=y^n$ satisfy $x=\zeta y$ for an $n$-th
root of unity $\zeta$, as $x^n-y^n=\prod_\zeta(x-\zeta y)$. The line $\mathbf Cw$ is stable under
$\sigma\otimes1$ for every automorphism $\sigma$ of $\mathbf C$ over $K$. Indeed, $S=\sigma\otimes1$ is a ring
endomorphism of $\mathbf C\otimes_KA$ with $S(zx)=\sigma(z)S(x)$ for $z\in\mathbf C$. Choose $\mu\in\mathbf C$ with
$\mu^n=\sigma(c)/c$. Then $S(w)^n=S(c\otimes u)=\sigma(c)\otimes u=(\mu w)^n$, so $S(w)=\zeta\mu w$ for some
$\zeta$, and $S(zw)=\sigma(z)\zeta\mu w\in\mathbf Cw$.

By Lemma~\ref{cv-volume-key-descent-spaces}, $\mathbf Cw=\mathbf C\otimes_KN$ for a $K$-subspace $N$ of $A$,
which is nonzero as $w\ne0$. Let $b\in N$ be nonzero. Then $1\otimes b=lw$ for some $l\in\mathbf C$, and
$1\otimes b^n=(1\otimes b)^n=l^nw^n=l^nc\otimes u$. Applying $T$ gives $b^n=\lambda(l^nc)\,u$. As
$\mathbf C\otimes_KA$ is a domain and $1\otimes b\ne0$, $1\otimes b^n\ne0$, so $b^n\ne0$ and
$\lambda(l^nc)\ne0$. Hence $u=\lambda(l^nc)^{-1}b^n$.
\end{proof}
% Lean: exists_eq_smul_of_pow_eq_pow for x = zeta y; the stability is hst, and the descent is
% GaloisDescent.complex_stable_iff_exists_baseChange.

\begin{theorem}[Grothendieck \cite{Gro71}]\label{cv-volume-key-connected}
Let $S$ be a finitely generated $\mathbf C$-algebra that is a domain. The set of homomorphisms of
$\mathbf C$-algebras $S\to\mathbf C$, with the topology of pointwise convergence, is connected.
\end{theorem}
\leanref{Development/AbelianVarietyTheory/Proofs/ThetaLawProofs/ThetaFrameValueEqUnitSmoothConnectedBase.lean}{GlobalTheory.RootConst.preconnectedSpace_algHom}
Generators $a_1,\dots,a_m$ of $S$ identify these homomorphisms with the complex points of the irreducible
variety $\operatorname{Spec}S$ in $\mathbf C^m$, and the topology of pointwise convergence with the complex
topology, as every element of $S$ is a polynomial in the $a_i$. These points form a connected space
\cite[Expos\'e~XII, \S2]{Gro71}.

\begin{lemma}\label{cv-volume-key-continuity}
Let $S$ be a smooth $\mathbf C$-algebra, let $P$ be the set of homomorphisms of $\mathbf C$-algebras
$S\to\mathbf C$ with the topology of pointwise convergence, and let $t\colon P\to\mathbf C$. Suppose that
$b\mapsto t(\psi_b)$ is continuous on $B$ whenever $B$ is an open subset of some $\mathbf C^m$ and
$(\psi_b)_{b\in B}$ is a family in $P$ with $b\mapsto\psi_b(a)$ holomorphic on $B$ for every $a\in S$. Then
$t$ is continuous.
\end{lemma}
\leanref{Development/AbelianVarietyTheory/Proofs/ThetaLawProofs/ThetaFrameValueEqUnitSmoothConnectedBase.lean}{GlobalTheory.RootConst.continuous_of_holomorphic_families}
\begin{proof}
Let $p\in P$. The points of $P$ are the complex points of the smooth variety $\operatorname{Spec}S$, which form
a complex manifold on which the elements of $S$ are holomorphic functions \cite[Expos\'e~XII, \S2]{Gro71}.
So there are $d\ge0$, an open neighbourhood $U$ of $p$, an open set $V\subset\mathbf C^d$ and a homeomorphism
$e\colon U\to V$ such that $b\mapsto e^{-1}(b)(a)$ is holomorphic on $V$ for every $a\in S$. By the hypothesis,
applied to the family $e^{-1}(b)$ for $b\in V$, the function $t\circ e^{-1}$ is continuous on $V$. So $t$ is
continuous on $U$, and in particular at $p$.
\end{proof}
% Lean: the chart is exists_chart, from Algebra.IsSmoothAt.exists_notMem_isStandardSmooth (a standard smooth
% localization at p, [Sta, Tag 00TA]) and AffinePoints.exists_isAlgebraicChart_of_away (the holomorphic
% implicit function theorem; the coordinates of e are quotients p(x_i)/p(s) of elements of S).

\begin{lemma}\label{cv-volume-key-roots-complex}
Let $B$ be a smooth $\mathbf C$-algebra that is a domain, $u\in B^\times$ and $n\ge1$. Suppose that a function
$\rho$ on the homomorphisms of $\mathbf C$-algebras $B\to\mathbf C$ satisfies $\rho(\psi)^n=\psi(u)$ for every
$\psi$, and that it is continuous along holomorphic families, as in Lemma~\ref{cv-volume-key-continuity}. Then
$u=b^n$ for some $b\in B$.
\end{lemma}
\leanref{Development/AbelianVarietyTheory/Proofs/IwasawaDecomposition/ExistsPowEqUnitOfContinuousComplexRootStatement.lean}{GlobalTheory.T.exists_pow_eq_unit_of_continuous_complex_root}
\leanref{Development/AbelianVarietyTheory/Proofs/IwasawaDecomposition/ExistsPowEqUnitOfContinuousComplexRoot.lean}{GlobalTheory.ContRoot.exists_pow_eq_of_continuous_root}
\begin{proof}
By Lemma~\ref{cv-volume-key-continuity}, $\rho$ is continuous. The ring $B$ is normal \cite[Tag 033C]{Sta},
so it is integrally closed in its field of fractions $K$. Let $\beta$ be an $n$-th root of $u$ in an algebraic
closure of $K$. It is integral over $B$, so its minimal polynomial $P$ over $K$ is monic with coefficients in
$B$ \cite[Proposition~5.15]{AM69}, of degree at least $1$. It divides $X^n-u$ in $K[X]$, hence in $B[X]$, as
$P$ is monic. The ring $B'=B[X]/(P)$ is free over $B$ with basis $1,X,\dots,X^{\deg P-1}$, so it embeds in the
field $K[X]/(P)$ and is a domain; it is finitely generated over $\mathbf C$. Let $\theta$ be the class of $X$,
so that $\theta^n=u$ in $B'$.

As $u$ is a unit, $\psi(u)\ne0$ and so $\rho(\psi)\ne0$ for every $\psi$. For a homomorphism
$\chi\colon B'\to\mathbf C$ let $\chi|_B$ be its restriction to $B$. The function
$\chi\mapsto\chi(\theta)/\rho(\chi|_B)$ is continuous for the topologies of pointwise convergence, and its
$n$-th power is $\chi(u)/\chi(u)=1$. By Theorem~\ref{cv-volume-key-connected} the homomorphisms
$B'\to\mathbf C$ form a connected space, and the $n$-th roots of unity a discrete one, so this function is a
constant $\zeta$.

Choose a homomorphism $\psi\colon B\to\mathbf C$, which exists by the Nullstellensatz, and let $P^\psi$ be the
image of $P$. For every root $z$ of $P^\psi$ there is a homomorphism $\chi$ extending $\psi$ with
$\chi(\theta)=z$, so $z=\zeta\rho(\psi)$, and $P^\psi$ has only one root. It divides $X^n-\psi(u)$, which has
distinct roots as $\psi(u)\ne0$. So $P^\psi$, which is monic of the degree of $P$, has $\deg P$ distinct roots,
and $\deg P=1$. Thus $P=X+c$ with $c\in B$, and $(-c)^n=u$, since $P$ divides $X^n-u$.
\end{proof}
% Lean: isIntegrallyClosed_of_smooth; exists_monic_prime_dvd (minpoly.isIntegrallyClosed_dvd and
% minpoly.prime_of_isIntegrallyClosed, so AdjoinRoot P is a domain); the constant is RootConst.eq_of_pow_eq_one;
% the extensions chi are AdjoinRoot.liftAlgHom; separable_X_pow_sub_C for the distinct roots.

\begin{lemma}\label{cv-volume-key-roots}
Let $K\subset\mathbf C$ be a subfield, $A$ a smooth $K$-algebra such that $\mathbf C\otimes_KA$ is a domain,
$u\in A^\times$, $n\ge1$ and $c\in\mathbf C^\times$. Suppose that a function $\rho$ on the homomorphisms of
$K$-algebras $A\to\mathbf C$ satisfies $\rho(\psi)^n=c\,\psi(u)$ for every $\psi$, and that $b\mapsto\rho(\psi_b)$
is continuous on $B$ whenever $B$ is an open subset of some $\mathbf C^m$ and $(\psi_b)_{b\in B}$ is a family of
such homomorphisms with $b\mapsto\psi_b(a)$ holomorphic on $B$ for every $a\in A$. Then $u=ab^n$ with
$a\in K^\times$ and $b\in A$.
\end{lemma}
\leanref{Development/AbelianVarietyTheory/Proofs/IwasawaDecomposition/ExistsKeySectionUnif.lean}{GlobalTheory.KeyFourth.exists_eq_smul_pow}
\begin{proof}
The $\mathbf C$-algebra $B=\mathbf C\otimes_KA$ is smooth, as a base change of $A$, and a domain. Restricting
along $a\mapsto1\otimes a$ maps the homomorphisms of $\mathbf C$-algebras $B\to\mathbf C$ to the homomorphisms
of $K$-algebras $A\to\mathbf C$, and holomorphic families to holomorphic families. Put $u'=c\otimes u$, a unit
of $B$, and $\rho'(\chi)=\rho(\chi|_A)$. Then $\rho'(\chi)^n=c\,\chi(1\otimes u)=\chi(u')$, and $\rho'$ is
continuous along holomorphic families. By Lemma~\ref{cv-volume-key-roots-complex} there is $w\in B$ with
$w^n=c\otimes u$. As $B\ne0$, $A\ne0$, so the unit $u$ is nonzero, and
Lemma~\ref{cv-volume-key-roots-descent} gives $u=ab^n$ with $a\in K^\times$ and $b\in A$.
\end{proof}

\begin{definition}\label{cv-volume-key-theta-point}
In Situation~\ref{cv-uniformization-situation}, let $I$ be the set of the coordinates $(\tau,a)$ of $\mathbf C^g$
together with one further index $0$. For $t\in(\mathbf C^\times)^I$ and $e\in\mathbf Z^I$ put
$t^e=\prod_{\iota\in I}t_\iota^{e_\iota}$ and $\chi(t)=t_0^{-2}\prod_{\tau,a}t_{\tau,a}$. The \emph{exponent} of a dominant
weight $\kappa=(\alpha,n)$ (Situation~\ref{cv-cusp-setting}) is the family $e(\kappa)\in\mathbf Z^I$ with
$e(\kappa)_{\tau,a}=\alpha_{\tau,a}$ and $e(\kappa)_0=n$.
\begin{enumerate}
\item For complex framed points $y$ and $y'$ of $X'$ over $\sigma$ and $t\in(\mathbf C^\times)^I$, write
  $y'=y\cdot t$ if $y'$ has the point of $y$, if for every $\tau$ and $a$ the vector of index $a$ of the basis of
  $y'$ in the eigenspace of $\tau$ is a combination of the vectors of $y$ in that eigenspace with index at most
  $a$, the vector of index $a$ having the coefficient $t_{\tau,a}$, and if the section of $y'$ is $t_0$ times that
  of $y$. So $y$ and $y'$ are triangularly equivalent if and only if $y'=y\cdot t$ for some $t$.
\item Let $i$ be an index, $x$ a frame and $\mu\in\mathbf C$. A \emph{theta point} at $x$ for $i$ with
  \emph{scalar} $\mu$ is a complex framed point $(y,v,w)$ of $X'$ over $\sigma$ with $y=y_i(P^i_x)$, with $v$
  the images in $\operatorname{Lie}A_y$ of the unit vectors of $\mathbf C^g$, the unit vector of the coordinate
  $(\tau,a)$ giving the vector of index $a$ in the eigenspace of $\tau$, and with $w$ the section whose
  pullback is $\mu\,\vartheta[c_i](Z_x,h_xz)$ times the factor of Theorem~\ref{cv-uniformization-family}(4). A
  normalized framed point at $x$ is a theta point at $x$ for $i$ with a scalar $\nu$ such that
  $\nu^2=N_\xi\det h_x$.
\item Let $u_i(x)$ be a normalized framed point at $x$, for every index $i$ and every frame $x$
  (Lemma~\ref{cv-uniformization-frames}(2)). A \emph{root value} of a complex framed point $y$ of $X'$ over
  $\sigma$ is a number $\chi(t)$ such that $y=u_i(x)\cdot t$ for some index $i$, some frame $x$ and some
  $t\in(\mathbf C^\times)^I$.
\item Let $V\subset\mathbf C^m$ be an open set. A \emph{holomorphic family of framed points} over $V$ is a
  family $(y_b)_{b\in V}$ of complex framed points of $X'$ over $\sigma$ for which there are an $F'$-algebra
  $A$, a framed point $y_A$ of $X'$ over $A$ and homomorphisms of $F'$-algebras $\psi_b\colon A\to\mathbf C$,
  with $\mathbf C$ an $F'$-algebra through $\sigma$, such that $b\mapsto\psi_b(f)$ is holomorphic on $V$ for every
  $f\in A$ and $y_b$ is the image of $y_A$ under $\psi_b$ for every $b\in V$.
\end{enumerate}
\end{definition}
\leanref{Development/AbelianVarietyTheory/Proofs/IwasawaDecomposition/KchMul.lean}{GlobalTheory.KeyFourth.kch}
\leanref{Formalism/LocalToGlobal/ProjectiveLattices/IsCombAt.lean}{Formalism.L2.Pj.FrBor}
\leanref{Development/AbelianVarietyTheory/Definitions/ExistenceOfLattices/NormConst.lean}{GlobalTheory.FrameUnif.IsThetaPt}
\leanref{Development/AbelianVarietyTheory/Proofs/IwasawaDecomposition/KchMul.lean}{GlobalTheory.KeyFourth.RootVal}
\leanref{Formalism/LocalToGlobal/ProjectiveLattices/IsCombAt.lean}{Formalism.L2.Pj.FrMap}

\begin{lemma}\label{cv-volume-key-root-values}
In Situation~\ref{cv-uniformization-situation}, let $u_i(x)$ be normalized framed points as in
Definition~\ref{cv-volume-key-theta-point}(3), and let $y$ be a complex framed point of $X'$ over $\sigma$.
\begin{enumerate}
\item The point $y$ has at most one root value.
\item If $y$ is a theta point at a frame $x$ for an index $i$ with a scalar $\mu\ne0$, then $N_\xi\det h_x/\mu^2$
  is a root value of $y$.
\end{enumerate}
\end{lemma}
\leanref{Development/AbelianVarietyTheory/Proofs/IwasawaDecomposition/KchEqOneOfWchar.lean}{GlobalTheory.KeyFourth.rootVal_unique}
\leanref{Development/AbelianVarietyTheory/Proofs/IwasawaDecomposition/ContinuousOnKeySectionComplexValue.lean}{GlobalTheory.KeyFourth.rootVal_of_isThetaPt}
\begin{proof}
The relation $y'=y\cdot t$ is reflexive, with $t=1$; it composes, $y'=y\cdot t$ and $y''=y'\cdot t'$ giving
$y''=y\cdot tt'$, and it is symmetric, $y'=y\cdot t$ giving $y=y'\cdot t^{-1}$, as products and inverses of
invertible upper triangular matrices are upper triangular, with the products and inverses of the diagonal
entries. The family $t$ is determined by $y$ and $y\cdot t$, as the vectors of $y$ form bases and its section is
not zero. Also $\chi(tt')=\chi(t)\chi(t')$.

(1) Let $\chi(t)$ and $\chi(t')$ be root values of $y$, with $u_i(x)\cdot t=y=u_{i'}(x')\cdot t'$. Then
$u_{i'}(x')=u_i(x)\cdot s$ with $s=tt'^{-1}$. By Lemma~\ref{cv-uniformization-frames}(7), $i'=i$ and
$x'=\gamma x\beta$ with $\gamma\in\Gamma_i$ and $\beta\in B$. By part (5) of that lemma, $u_i(\gamma x)=u_i(x)\cdot t_1$
for a family $t_1$ that multiplies sections of every dominant weight $\kappa$ by $1$, that is
$t_1^{e(\kappa)}=1$. By part (4), $u_i(\gamma x\beta)=u_i(\gamma x)\cdot t_2$ for a family $t_2$ that multiplies a
section of every dominant weight $\kappa$ by $\prod_{\tau,a}\beta_{\tau,aa}^{\lambda(\kappa)_{\tau,a}}$, that is
$t_2^{e(\kappa)}=\prod_{\tau,a}\beta_{\tau,aa}^{\lambda(\kappa)_{\tau,a}}$. So $u_i(x)\cdot t_1t_2=u_i(x')=u_i(x)\cdot s$,
and $s=t_1t_2$. Let $\kappa_1$ be the dominant weight with all $\alpha_{\tau,a}=1$ and $n=0$, and $\kappa_2$ the one
with all $\alpha_{\tau,a}=0$ and $n=2$. Their exponents give $\chi(t)=t^{e(\kappa_1)}/t^{e(\kappa_2)}$ for every $t$.
The weights $\lambda(\kappa_1)$ and $\lambda(\kappa_2)$ of Situation~\ref{cv-cusp-setting} agree: both have the
entry $-1$ at the positions $l\le q$ and $1$ at the positions $l>q$. Hence $\chi(t_1)=1$ and $\chi(t_2)=1$, and
$\chi(t)/\chi(t')=\chi(s)=\chi(t_1)\chi(t_2)=1$.

(2) Let $\nu$ be the scalar of $u_i(x)$, with $\nu^2=N_\xi\det h_x$; it is not zero, as $N_\xi\ne0$ and $h_x$ is
invertible. The framed points $u_i(x)$ and $y$ are theta points at $x$ with the scalars $\nu$ and $\mu$. So they
have the point $y_i(P^i_x)$ and the same bases, the images of the unit vectors. Their sections $w_0$ and $w$
have the pullbacks $\nu\vartheta[c_i](Z_x,h_xz)$ and $\mu\vartheta[c_i](Z_x,h_xz)$ times the factor. As $w_0$
generates the sections of the line bundle, $w=cw_0$ for some $c\in\mathbf C$, and comparing the pullbacks at a
point $z$ where $w_0$ does not vanish gives $c=\mu/\nu$. So $y=u_i(x)\cdot t$ with $t_{\tau,a}=1$ for all
$(\tau,a)$ and $t_0=\mu/\nu$, and $\chi(t)=(\mu/\nu)^{-2}=N_\xi\det h_x/\mu^2$.
\end{proof}
% Lean: the groupoid rules are GlobalTheory.Reg.BorelGroupoid.frBor_groupoid, and chi(t t') = chi(t) chi(t') is
% kch_mul. (1) is rootVal_unique: fibre_unfOf with FibreFlag.fibTri for (7), the Gamma rule hun.2.2.1 for (5),
% the Borel rule hun.2.1 for (4), and kch_eq_one_of_wchar for the weights ex Dw 1 0 and ex Dw 0 2, whose laU
% agree. (2) is rootVal_of_isThetaPt through GlobalTheory.frBor_of_isThetaPt (the tangent vectors agree by
% isDerivAt0_ex_isTanAt, and the sections by sec_eq_scal_of_values) and kch_eq.

\begin{lemma}\label{cv-volume-key-adapted-periods}
In Situation~\ref{cv-uniformization-situation}, let $(y_b)_{b\in V}$ be a holomorphic family of framed points over
an open set $V\subset\mathbf C^m$ (Definition~\ref{cv-volume-key-theta-point}(4)). Then every $b_0\in V$ has an
open neighbourhood $V'\subset V$ with an index $i$ and a holomorphic family of periods $P_b$, $b\in V'$, of $L_i$
of the type of the situation, such that for every $b\in V'$ the point of $y_b$ is $y_i(P_b)$ and the basis of
$y_b$ consists of the images in the Lie algebra of $A_{y_i(P_b)}$ of the unit vectors of $\mathbf C^g$ under the
map $\mathbf C^g\to A_{y_i(P_b)}$ of Theorem~\ref{cv-uniformization-family} for $P_b$, the unit vector of the
coordinate $(\tau,a)$ giving the vector of index $a$ in the eigenspace of $\tau$.
\end{lemma}
\leanref{Development/AbelianVarietyTheory/Proofs/GenericRegularity/ExistsHolomorphicPeriodsAdapted.lean}{GlobalTheory.AdaptedPeriods.exists_holPeriods_adapted}
\begin{proof}
Let $A$, $y_A$ and $\psi_b$ be as in Definition~\ref{cv-volume-key-theta-point}(4). The point of $y_b$ is the
point of $y_A$ composed with $\operatorname{Spec}\psi_b$, so near it its affine coordinates are values of
$\psi_b$, and the points of the $y_b$ form a holomorphic family of complex points of $X'$ over $\sigma$. By
Theorem~\ref{cv-uniformization-family}(3) there are an open neighbourhood $V'\subset V$ of $b_0$, an index $i$
and a holomorphic family of periods $P^0_b$, $b\in V'$, of $L_i$ of this type, such that $y_i(P^0_b)$ is the
point of $y_b$ for every $b\in V'$. Let $p^0_b\colon\mathbf C^g\to A_{y_b}$ be the map given by $P^0_b$.

The adapted periods. Let $b\in V'$. The derivative of $p^0_b$ at $0$ is an isomorphism of $\mathbf C^g$ onto the
Lie algebra of $A_{y_b}$, and it maps a vector into the eigenspace of $\tau$ if and only if the vector vanishes at
the coordinates of the other embeddings. For a coordinate $\iota=(\tau,a)$ let $w_{b,\iota}\in\mathbf C^g$ be the
vector that it maps to the vector $v_{b,\iota}$ of index $a$ in the eigenspace of $\tau$ of $y_b$. For a fixed $\tau$ the
$w_{b,(\tau,a)}$ form a basis of the vectors that vanish at the coordinates of the other embeddings, as the
$v_{b,(\tau,a)}$ form a basis of the eigenspace. So the matrix $W_b$ with columns $w_{b,\iota}$ is invertible and
block diagonal for the embeddings, and so is $Q_b=W_b^{-1}$, which maps $w_{b,\iota}$ to the unit vector
$\varepsilon_\iota$. As $Q_b$ commutes with the action of $\mathcal O_E$ on $\mathbf C^g$, $P_b=Q_b\circ P^0_b$ is a
period of $L_i$ of the same type, with the Hermitian form $H(Q_b^{-1}z,Q_b^{-1}z')$ for the form $H$ of $P^0_b$.
By Theorem~\ref{cv-uniformization-family}(2), with $h=Q_b$ and $\gamma=1$, $y_i(P_b)=y_i(P^0_b)$. The isomorphism
$\mathbf C^g/P_b(L_i)\to A_{y_b}$ composed with $z\mapsto Q_bz$ and the isomorphism of $P^0_b$ differ by an
automorphism of $A_{y_b}$ that preserves the polarization and the level structure, which is the identity, as in
Step~2 of the proof of Lemma~\ref{cv-uniformization-theta-change}. So the map $p_b$ given by $P_b$ satisfies
$p_b(Q_bz)=p^0_b(z)$, and the derivative of $p_b$ at $0$ maps $\varepsilon_\iota=Q_bw_{b,\iota}$ to $v_{b,\iota}$.

Holomorphy. It remains to show that $b\mapsto P_b(v)=Q_bP^0_b(v)$ is holomorphic for every $v\in L_i$. As
$P^0_b(v)$ is holomorphic in $b$ and $Q_b=W_b^{-1}$, it is enough that each $w_{b,\iota}$ is holomorphic in $b$ near
every point $b_1\in V'$. Let $f$ be a regular function on an open set of the universal abelian scheme that
contains the origin of $A_{y_{b_1}}$, and for $b$ near $b_1$ let $L_f(b)$ be the derivative at $0$ of
$z\mapsto f(p^0_b(z))$, a linear form on $\mathbf C^g$. It is holomorphic in $b$, as $p^0_b(z)$ is holomorphic in
$(b,z)$ by Theorem~\ref{cv-uniformization-family}(3). The number $L_f(b)(w_{b,\iota})$ is the derivative of $f$
along $v_{b,\iota}$. This tangent vector is the image under $\psi_b$ of the corresponding tangent vector of $y_A$,
so the derivative is a quotient of values of $\psi_b$ at elements of $A$, with a denominator that does not
vanish at $b_1$; it is holomorphic in $b$ near $b_1$. The forms $L_f(b_1)$ span the dual of $\mathbf C^g$, as the
derivative of $p^0_{b_1}$ at $0$ is injective. Choose $f_1,\dots,f_g$ with $L_{f_1}(b_1),\dots,L_{f_g}(b_1)$ a basis
of the dual. The matrix with rows $L_{f_k}(b)$ is holomorphic in $b$ and invertible near $b_1$, and it maps
$w_{b,\iota}$ to the vector of the $L_{f_k}(b)(w_{b,\iota})$, which is holomorphic. So $w_{b,\iota}$ is holomorphic near
$b_1$.
\end{proof}
% Lean: the points are holomorphic by ThetaValueCont.isHolTo_spMap and LocLift.isHolTo_comp_hom; the local
% periods P0 are the field lift of FamilyUnif. The adapted periods are exists_adapted, the same argument
% as CoverFrame.coverTri: isDerivAt0_ex_isTanAt, isEigAt_iff_of_isDerivAt0, exists_inverse_block for Q_b,
% PeriodLin.exists_period_mulVec, FamRigid.ex_eq_of_level_group (with gam = 1 and h0 = Q_b) and
% ExpChain.isDerivAt0_ex_iff. The holomorphy of w is HolFrameMatrix.differentiableAt_coord
% (HolLieFrame.differentiableAt_of_forms with the forms of differentiableAt_fderiv and the values of
% exists_hol_snd), and that of Q_b is differentiableAt_of_mul_eq_one.

\begin{lemma}\label{cv-volume-key-frame-lift}
In Situation~\ref{cv-uniformization-situation}, let $i$ be an index, $V\subset\mathbf C^m$ an open set, $P_b$,
$b\in V$, a holomorphic family of periods of $L_i$ of the type of the situation, and $b_0\in V$. Then there are
frames $x_b$, $b\in V$, with $P^i_{x_b}=P_b$ for every $b\in V$, such that $b\mapsto x_b$ is continuous at $b_0$.
\end{lemma}
\leanref{Development/AbelianVarietyTheory/Proofs/IwasawaDecomposition/ExistsFrameLift.lean}{GlobalTheory.FrameLiftAt.exists_frame_lift}
\begin{proof}
By Lemma~\ref{cv-uniformization-frames}(1) there is for every $b\in V$ a frame $x^1_b$ with $P^i_{x^1_b}=P_b$; put
$x^0=x^1_{b_0}$. By the formula for periods in Situation~\ref{cv-uniformization-situation}, $P^i_x$ depends only
on the columns $x_{\tau,k}$ with $k\le q$ and on the rows of $x_\tau^{-1}$ of index $k>q$, for $\tau\in\Phi$.

Step 1: convergence. Let $\tau\in\Phi$. For $k\le q$ the map $u\mapsto h_{\tau(G)}(x_{\tau,k},\tau(u))$ on $E^r$ is
additive and is multiplied by $\overline{\tau(c)}$ when $u$ is multiplied by $c\in E$; for $k>q$ the map
$u\mapsto(x_\tau^{-1}\tau(u))_k$ is additive and is multiplied by $\tau(c)$. As $b\to b_0$, $P_b(v)\to P_{b_0}(v)$
for every $v\in L_i$. So these maps for $x=x^1_b$ converge to those for $x^0$ at every $u\in\Lambda_i$, and, as
$\Lambda_i$ spans $E^r$ over $E$, at every $u\in E^r$. For $u$ a unit vector this says that the rows of
$(x^1_{b,\tau})^{-1}$ of index $k>q$ converge to those of $(x^0_\tau)^{-1}$. For $u$ with entries
$u_l=\overline{(G^{-1})_{lk'}}$,
\[
 h_{\tau(G)}(z,\tau(u))=\sum_{a,l}z_a\tau(G_{al})\tau\big((G^{-1})_{lk'}\big)=z_{k'}\qquad(z\in\mathbf C^r),
\]
as $\overline{\tau(\bar c)}=\tau(c)$ on the CM field $E$; so the columns $x^1_{b,\tau,k}$ with $k\le q$ converge to
those of $x^0_\tau$.

Step 2: nearby frames. Let $\tau\in\Phi$, and write $A_b=x^1_{b,\tau}$ and $A_0=x^0_\tau$. Let $R_b$ be the matrix
whose columns of index $k\le q$ are those of $A_b$ and whose other columns are those of $A_0$, and let $S_b$ be the
matrix whose rows of index $k\le q$ are those of the identity and whose other rows are those of $A_b^{-1}R_b$. By
Step 1, $R_b\to A_0$, and the rows of index $k>q$ of $A_b^{-1}R_b$ converge to those of $A_0^{-1}A_0=1$, so
$S_b\to1$. For $b$ near $b_0$ the matrices $R_b$ and $S_b$ are invertible; put $C_b=R_bS_b^{-1}$, so that
$C_b\to A_0$. For $l\le q$ the column $l$ of $S_b$ is the unit vector $\varepsilon_l$: in the rows $k\le q$ by
definition, and in the rows $k>q$ because the column $l$ of $R_b$ is that of $A_b$, so that the column $l$ of
$A_b^{-1}R_b$ is $\varepsilon_l$. Hence $C_b\varepsilon_l=R_b\varepsilon_l=A_b\varepsilon_l$, and the columns of
$C_b$ of index $l\le q$ are those of $A_b$. The row $k>q$ of $C_b^{-1}=S_bR_b^{-1}$ is the row $k$ of
$A_b^{-1}R_bR_b^{-1}=A_b^{-1}$. Let $x^2_b$ be the family of these $C_b$ for $\tau\in\Phi$. Then $x^2_b\to x^0$ as
$b\to b_0$, and $x^2_b$ has the columns and rows of $x^1_b$ on which the period depends. The frames form an open
set, as the conditions on the signs of the minors $\delta_{\tau,l}$ are open; so for $b$ near $b_0$, $x^2_b$ is a
frame with $P^i_{x^2_b}=P_b$.

Step 3: a selection. Fix a norm on families of matrices. For $b\in V$ let $d_b$ be the infimum of $\|x-x^0\|$ over
the frames $x$ with $P^i_x=P_b$, of which $x^1_b$ is one. Put $x_{b_0}=x^0$, and for $b\ne b_0$ choose a frame
$x_b$ with $P^i_{x_b}=P_b$ and $\|x_b-x^0\|\le d_b+|b-b_0|$. By Step 2, $d_b\le\|x^2_b-x^0\|$ for $b$ near $b_0$,
which tends to $0$; so $x_b\to x^0$ as $b\to b_0$.
\end{proof}
% Lean: the frames x1 are from GlobalTheory.GenericFrames.dictionary (honto), and the period formula is its hper
% (pc of the frame, with T.per). Step 1 is in PeriodLift.period_lift, with tendsto_of_span for the spanning and
% hermC_emb_inv_col for the vector u; Step 2 is PeriodLift.exists_matrix_lift (Q, D and Q D^{-1} in the Lean,
% R_b, S_b and C_b here) with FourierCoefficientDomain.isOpen_dom for the openness of the frames; Step 3 is
% FrameLiftAt.exists_selection.

\begin{lemma}\label{cv-volume-key-theta-scalar}
In Situation~\ref{cv-uniformization-situation}, use the terms of Definition~\ref{cv-volume-key-theta-point}.
\begin{enumerate}
\item Let $i$ be an index, $x$ a frame and $y$ a complex framed point of $X'$ over $\sigma$ with the point
  $y_i(P^i_x)$ whose basis consists of the images in its Lie algebra of the unit vectors of $\mathbf C^g$, the unit
  vector of the coordinate $(\tau,a)$ giving the vector of index $a$ in the eigenspace of $\tau$. Then $y$ is a theta point at
  $x$ for $i$ with a scalar $\mu\ne0$.
\item Let $(y_b)_{b\in V}$ be a holomorphic family of framed points over an open set $V\subset\mathbf C^m$, let $i$
  be an index and $P_b$, $b\in V$, a holomorphic family of periods of $L_i$ of the type of the situation, and for
  every $b\in V$ let $x_b$ be a frame with $P^i_{x_b}=P_b$ and $\mu_b\in\mathbf C$ a number such that $y_b$ is a
  theta point at $x_b$ for $i$ with scalar $\mu_b$. If $b\mapsto x_b$ is continuous at a point $b_0\in V$, then so
  is $b\mapsto\mu_b$.
\end{enumerate}
\end{lemma}
\leanref{Development/AbelianVarietyTheory/Proofs/IwasawaDecomposition/ContinuousOnKeySectionComplexValue.lean}{GlobalTheory.KeyFourth.exists_isThetaPt_of_adapted}
\leanref{Development/AbelianVarietyTheory/Proofs/IwasawaDecomposition/ContinuousAtScalar.lean}{GlobalTheory.ThetaScalar.continuousAt_scalar}
\begin{proof}
(1) By Theorem~\ref{cv-uniformization-family}(4), for the matrices $h_x$ and $Z_x$ and the characteristic $c_i$
fixed in Situation~\ref{cv-uniformization-situation}, the line bundle of the polarization on $A_y$ has a section
$s$ whose pullback is $\vartheta[c_i](Z_x,h_xz)$ times the factor. The theta function $\vartheta[c_i](Z_x,\cdot)$ is
not identically zero, as its Fourier coefficients in $z$ are not zero (Theorem~\ref{cv-uniformization-theta-transformation}
gives the series). As $h_x$ is invertible and the factor does not vanish, $s\ne0$. The section of $y$ generates
the sections of this line bundle, so $s$ is $c$ times it for some $c$, and $c\ne0$. Hence the pullback of the
section of $y$ is $c^{-1}\vartheta[c_i](Z_x,h_xz)$ times the factor, and $y$ is a theta point at $x$ with scalar
$\mu=c^{-1}$.

(2) Let $A$, $y_A$ and $\psi_b$ be as in Definition~\ref{cv-volume-key-theta-point}(4). For $b\in V$ let
$p_b\colon\mathbf C^g\to A_{y_b}$ be the map given by $P_b$, and $w_b$ the section of $y_b$. The section $w_{b_0}$
generates the sections of the line bundle on $A_{y_{b_0}}$, so it is not zero, and as $p_{b_0}$ is onto there are
$z_0\in\mathbf C^g$ and an index $j$ of the open cover with $p_{b_0}(z_0)\in U_j$ at which the value of $w_{b_0}$ in
the trivialization over $U_j$ is not zero. Let $\mathcal A_A$ be the abelian scheme of $y_A$ over
$\operatorname{Spec}A$, with the section $w_A$ of $y_A$. For $b\in V$ the point $p_b(z_0)$ lies over the point of
$y_b$, which is the image of the point of $y_A$ under $\psi_b$; with $\psi_b$ it gives a complex point $q_b$ of
$\mathcal A_A$, which depends holomorphically on $b$, as $\psi_b$ and $p_b(z_0)$ do
(Theorem~\ref{cv-uniformization-family}(3)). As $y_b$ is the image of $y_A$, the value $f(b)$ of $w_A$ at $q_b$ in
the trivialization over $U_j$ is the value of $w_b$ at $p_b(z_0)$. Let $W$ be the set of $b\in V$ at which this
value is defined and not zero. It is open, it contains $b_0$, and $f$ is holomorphic on $W$. For $b\in W$, as
$y_b$ is a theta point at $x_b$ with scalar $\mu_b$ and $P^i_{x_b}=P_b$,
\[
 f(b)=\mu_b\,\vartheta[c_i](Z_{x_b},h_{x_b}z_0)\,F_j(x_b,z_0),
\]
where $F_j(x,z)$ is the factor at $x$; so the last two factors do not vanish. As $Z_x$, $h_x$ and $F_j(x,z_0)$
are holomorphic in $x$, $\vartheta[c_i](Z,z)$ is holomorphic in $(Z,z)$, and $b\mapsto x_b$ is continuous at
$b_0$, these factors are continuous at $b_0$. So
$\mu_b=f(b)\big/\big(\vartheta[c_i](Z_{x_b},h_{x_b}z_0)F_j(x_b,z_0)\big)$ for $b\in W$, and as $W$ is a
neighbourhood of $b_0$, $b\mapsto\mu_b$ is continuous at $b_0$.
\end{proof}
% Lean: (1) is exists_isThetaPt_of_adapted: T.norm, ThetaHalf.thetaChar_holomorphic_ne_zero (the nonvanishing is
% IsSiegel.exists_riemannTheta_add_ofReal_ne_zero, an integral of theta times an exponential over the unit cube),
% the section of the factor T.fr_sec, and GenFactor.exists_factor (FibreSection.exists_sec_same_values and the
% generator property FrPt.gen). (2) is continuousAt_scalar: z0 and j from HolGen.exists_val_ne, the points q_b
% from exists_pt, holomorphic by HolPullback.isHolTo_of_components, the values by val_pull, W a basic open set,
% and the continuity of theta and of the factor from thetaChar_holomorphic_ne_zero and T.fr_hol.

\begin{lemma}\label{cv-volume-key-root-continuity}
In Situation~\ref{cv-uniformization-situation}, let $u_i(x)$ be normalized framed points as in
Definition~\ref{cv-volume-key-theta-point}(3), let $(y_b)_{b\in V}$ be a holomorphic family of framed points over
an open set $V\subset\mathbf C^m$, and let $\rho\colon V\to\mathbf C$ be a function such that $\rho(b)$ is a root value
of $y_b$ for every $b\in V$. Then $\rho$ is continuous on $V$.
\end{lemma}
\leanref{Development/AbelianVarietyTheory/Proofs/IwasawaDecomposition/ContinuousOnKeySectionComplexValueStatement.lean}{GlobalTheory.T.continuousOn_keySection_complexValue}
\leanref{Development/AbelianVarietyTheory/Proofs/IwasawaDecomposition/ContinuousOnKeySectionComplexValue.lean}{GlobalTheory.KeyFourth.continuousOn_rootVal}
\begin{proof}
Let $b_0\in V$. By Lemma~\ref{cv-volume-key-adapted-periods} there are an open neighbourhood $V'\subset V$ of $b_0$,
an index $i$ and a holomorphic family of periods $P_b$, $b\in V'$, of $L_i$ such that every $y_b$ with $b\in V'$
has the point $y_i(P_b)$ and the images of the unit vectors as its basis. By
Lemma~\ref{cv-volume-key-frame-lift} there are frames $x_b$, $b\in V'$, with $P^i_{x_b}=P_b$, such that
$b\mapsto x_b$ is continuous at $b_0$. By Lemma~\ref{cv-volume-key-theta-scalar}(1), each $y_b$ is a theta point at
$x_b$ for $i$ with a scalar $\mu_b\ne0$, and by part (2) of that lemma $b\mapsto\mu_b$ is continuous at $b_0$. By
Lemma~\ref{cv-volume-key-root-values}(2), $N_\xi\det h_{x_b}/\mu_b^2$ is a root value of $y_b$, so by part (1) of
that lemma
\[
 \rho(b)=\frac{N_\xi\det h_{x_b}}{\mu_b^2}\qquad(b\in V').
\]
As $h_x$ is holomorphic in $x$ and $b\mapsto x_b$ is continuous at $b_0$, $\det h_{x_b}$ is continuous at $b_0$;
and $\mu_{b_0}\ne0$. So $\rho$ is continuous at $b_0$, as $V'$ is a neighbourhood of $b_0$.
\end{proof}
% Lean: continuousOn_rootVal: AdaptedPeriods.exists_holPeriods_adapted, FrameLiftAt.exists_frame_lift,
% exists_isThetaPt_of_adapted, ThetaScalar.continuousAt_scalar, the continuity of (T k).h from (T k).hol and
% FourierCoefficientDomain.isOpen_dom, then rootVal_of_isThetaPt and rootVal_unique. The Lean states this for a
% uniformization family F of X at any level (the field of Formalism.L2.PjFree.IsUnf), with models M (ModelBase,
% NegXi), theta data T and normalized points un with IsNormUn and FramesOK; the text states it in
% Situation cv-uniformization-situation, where the book has these objects. Level 1 suffices for the book: it
% applies Lemma cv-volume-twisting-key-section(2) only to deep covers Y (proof of Proposition
% cv-arithmetic-volume-step), whose levels have full level 1.

\begin{theorem}[Mumford \cite{Mum70}]\label{cv-volume-key-sections}
Let $S$ be a locally noetherian scheme, $\pi\colon A\to S$ an abelian scheme of relative dimension $g$, and $L$ a
line bundle on $A$ that is relatively ample over $S$, trivial along the zero section and, locally on $S$,
isomorphic to its pullback under $-1$. Suppose that $L$ defines a principal polarization: a point $x$ of $A$ over
a ring $R$ is zero whenever the line bundle $t_x^*L\otimes L^{-1}$ on $A\times_S\operatorname{Spec}R$, where $t_x$
is the translation by $x$, is trivial locally on $\operatorname{Spec}R$. Let $n\ge1$. Then every point of $S$ has
an open neighbourhood $V$ with sections $b_1,\dots,b_{n^g}$ of $L^n$ over $\pi^{-1}(V)$ such that, for every ring
$R$ and every morphism $\operatorname{Spec}R\to V$, the pullbacks of $b_1,\dots,b_{n^g}$ form a basis of the
$R$-module of sections of the pullback of $L^n$ to $A\times_S\operatorname{Spec}R$.
\end{theorem}
\leanref{ReciprocityTheory/ReductiveGroups/ArithmeticGeometry/ExtensionDescent/GlobalSectionsRankAlgClosed.lean}{ReciprocityTheory.ChartRestriction.G.Out.SecBaseChangeSt}
\leanref{Development/AbelianVarietyTheory/Proofs/GenericRegularity/ExistsIntegralSmoothChartFrameModuli.lean}{GlobalTheory.LocalFrame.secBaseChange_of_secTwist}

On every geometric fibre of $A$ the line bundle $L$ is ample and its group $K(L)$ is trivial, so
$\chi(L)^2=\deg\phi_L=1$; by the vanishing theorem and the Riemann--Roch theorem the higher cohomology of $L^n$
vanishes and its space of sections has dimension $\chi(L^n)=n^g$ \cite[\S16]{Mum70}. Over an affine noetherian
open set $\operatorname{Spec}B\subset S$, the cohomology of $L^n$ over $A\times_S\operatorname{Spec}R$, for every
$B$-algebra $R$, is that of $K^\bullet\otimes_BR$ for a finite complex $K^\bullet$ of finitely generated
projective $B$-modules \cite[\S5, Theorem]{Mum70}. By Corollary~3 of that section, applied in each degree, and
Nakayama's lemma, the cohomology of $K^\bullet$ vanishes in positive degrees. So $K^\bullet$ is split exact in
positive degrees, and its cohomology $H$ in degree $0$ is a projective $B$-module of rank $n^g$ such that
$R\otimes_BH$ is the module of sections over $A\times_S\operatorname{Spec}R$ for every $B$-algebra $R$, as in the
proof of Corollary~2 there. Over an affine neighbourhood on which $H$ is free, a basis gives the sections $b_i$.

\begin{lemma}\label{cv-volume-key-generator}
In the setting of Definition~\ref{cv-volume-twisting-section}, let $x$ be a point of $X$ over a ring $S$, with
abelian scheme $A_x$ over $S$ and line bundle $L_x$ of the polarization. A section of $L_x$ over $A_x$
\emph{vanishes on no fibre} if, for every point $\mathfrak p$ of $\operatorname{Spec}S$, it is nonzero at some
point of the fibre of $A_x$ over $\mathfrak p$.
\begin{enumerate}
\item Let $\varphi\colon S\to T$ be a ring homomorphism, $s$ a section of $L_x$ over $A_x$ and $s'$ its pullback to
  $A_x\times_S\operatorname{Spec}T$. Then $s'$ is nonzero at some point of the fibre over a point $\mathfrak q$ of
  $\operatorname{Spec}T$ if and only if $s$ is nonzero at some point of the fibre over $\varphi^{-1}(\mathfrak q)$.
\item Every point $\mathfrak p$ of $\operatorname{Spec}S$ has an element $g\in S\setminus\mathfrak p$ and a section
  $\theta$ of $L_x$ over $A_x\times_S\operatorname{Spec}S_g$ whose pullback along every ring homomorphism
  $S_g\to T$ generates the $T$-module of sections of the pullback of $L_x$ and vanishes on no fibre.
\item A section of $L_x$ over $A_x$ generates the $S$-module of sections of $L_x$ over $A_x$ if and only if it
  vanishes on no fibre.
\end{enumerate}
\end{lemma}
\leanref{Development/AbelianVarietyTheory/Proofs/UnitaryExistence/ScaledBaseChangeFibreSectionNonvanishingIff.lean}{GlobalTheory.Frames.secBC_nonvanishing_iff}
\leanref{Development/AbelianVarietyTheory/Proofs/UnitaryExistence/ExistsLocalGeneratorStableBaseChangeStatement.lean}{GlobalTheory.T.exists_local_generator_stable_baseChange}
\leanref{Development/AbelianVarietyTheory/Proofs/UnitaryExistence/FibreSectionGeneratesOfNowhereZeroStatement.lean}{GlobalTheory.T.fibreSectionGenerates_of_nowhereZero}
\leanref{Development/AbelianVarietyTheory/Proofs/UnitaryExistence/FibreSectionNowhereZeroOfGeneratesStatement.lean}{GlobalTheory.T.fibreSectionNowhereZero_of_generates}
\leanref{Development/AbelianVarietyTheory/Proofs/UnitaryExistence/MemNonzeroLocusOfPull.lean}{GlobalTheory.SectionBaseChange.generates}
\begin{proof}
The scheme $X$ is of finite type over $F_c$, hence locally noetherian, and its universal abelian scheme with the
line bundle of the polarization satisfies the hypotheses of Theorem~\ref{cv-volume-key-sections}.

(1) Put $\mathfrak p=\varphi^{-1}(\mathfrak q)$. The fibre of $A_x\times_S\operatorname{Spec}T$ over $\mathfrak q$ is
the base change of the fibre of $A_x$ over $\mathfrak p$ along the field extension
$k(\mathfrak p)\to k(\mathfrak q)$, so every point of the second fibre is the image of a point of the first. At a
point $z$ of the first fibre with image $z_0$, the value of $s'$ at $z$ is the image of the value of $s$ at $z_0$
under the injective map of residue fields, so $s'$ is nonzero at $z$ if and only if $s$ is nonzero at $z_0$.

(3), the condition is sufficient. Let $\theta$ vanish on no fibre and let $s$ be a section of $L_x$ over $A_x$. By
Theorem~\ref{cv-volume-key-sections} for $n=1$, every point of $X$ has an open neighbourhood $V$ with a section
$b$ of the line bundle over the abelian scheme over $V$ whose pullbacks along all morphisms from affine schemes
to $V$ are generators. Cover $\operatorname{Spec}S$ by affine open sets $U=\operatorname{Spec}S_U$ that map into
such sets $V$, and let $b_U$ be the pullback of $b$ to $U$. Then $\theta=c\,b_U$ and $s=d\,b_U$ over $U$ with
$c,d\in S_U$. If $c$ lay in a prime $\mathfrak p$ of $S_U$, then $\theta$ would vanish at every point of the fibre
over $\mathfrak p$; so $c$ is a unit, and $s=(d/c)\theta$ over $U$. The coefficient $a_U=d/c$ is the only $a\in
S_U$ with $s=a\theta$ over $U$, as $b_U$ is a generator and $c$ a unit; so the $a_U$ agree on the overlaps and glue
to the unique $a\in S$ with $s=a\theta$.

(2) Let $p$ be the image of $\mathfrak p$ in $X$, and $V$ and $b$ as in the proof of (3). Shrink $V$ to an affine
open neighbourhood $\operatorname{Spec}B$ of $p$, with $B$ noetherian. The section $b$ vanishes on no fibre over
$\operatorname{Spec}B$: at a point $\mathfrak r$ of $\operatorname{Spec}B$, the pullback of $b$ to the fibre over
$k(\mathfrak r)$ is a generator of a free module of rank one, hence nonzero; that fibre is smooth over
$k(\mathfrak r)$, hence reduced, so a section that vanishes at all of its points is zero. Choose $g\notin\mathfrak
p$ such that $\operatorname{Spec}S_g$ maps into $\operatorname{Spec}B$, and let $\theta$ be the pullback of $b$.
For $S_g\to T$, the pullback of $\theta$ is the pullback of $b$ along $\operatorname{Spec}T\to\operatorname{Spec}B$;
it vanishes on no fibre by (1), and so it generates by the first part of (3), applied over $T$.

(3), the condition is necessary. Let $s$ generate and let $\mathfrak p$ be a point of $\operatorname{Spec}S$. Let
$g$ and $\theta$ be as in (2). As $A_x$ is quasi-compact and separated over $\operatorname{Spec}S$, some
$g^N\theta$ is the pullback of a section $s_0$ of $L_x$ over $A_x$ \cite[Tag 01PW]{Sta}. Since $g$ is a unit in
$S_g$ and $\theta$ vanishes on no fibre, $g^N\theta$ is nonzero at some point of the fibre over the prime of
$S_g$ above $\mathfrak p$, and by (1) $s_0$ is nonzero at some point of the fibre over $\mathfrak p$. As $s$
generates, $s_0=as$ for some $a\in S$, so $s$ is nonzero at that point too.
\end{proof}
% Lean: (1) is secBC_nonvanishing_iff (exists_mem_basicOpen_iff_of_isPullback). The sufficiency in (3) is
% FrameGenerator.secGen_of_secNZ = SectionBaseChange.generates (generates_local, existsUnique_of_local); (2) is
% Frames.exists_local_generator (PushLoc.exists_secNZ through PushInv.pushInvertible and
% PushInv.exists_nonvanishing, exists_away_lift, and fibreSectionGenerates_of_nowhereZero_noetherian_baseChange =
% FrameGenerator.secGen_of_secNZ_baseChange for the generation); the necessity in (3) is Frames.secNZ_of_secGen
% (SecLoc.exists_secBC_pow_mul for s0). The Lean's Frames.secGen_of_secNZ, used by exists_frMap, proves the
% sufficiency a second time by gluing the local generators of (2); the text gives only the first proof. All of
% them assume the base change law SecBaseChangeSt (Theorem cv-volume-key-sections).

\begin{lemma}\label{cv-volume-key-lie}
In the setting of Definition~\ref{cv-volume-twisting-section}, let $x$ be a point of $X$ over an $F_c$-algebra $S$.
Let $\operatorname{Lie}(x)$ be the Lie algebra of $A_x$ over $S$, the $S$-module of morphisms
$\operatorname{Spec}S[\varepsilon]/(\varepsilon^2)\to A_x$ over $\operatorname{Spec}S$ that restrict to the zero
section, and for an embedding $\tau\colon E\to\mathbf C$ let $\operatorname{Lie}_\tau(x)$ be its submodule on which
each $b\in\mathcal O_E$ acts by $\tau(b)$.
\begin{enumerate}
\item For every homomorphism $S\to T$ of $F_c$-algebras, the images of a basis of $\operatorname{Lie}(x)$ form a
  basis of the Lie algebra of $A_x\times_S\operatorname{Spec}T$ over $T$.
\item $\operatorname{Lie}(x)$ is the direct sum of the $\operatorname{Lie}_\tau(x)$. Vectors
  $u_{\tau,i}\in\operatorname{Lie}_\tau(x)$, for the embeddings $\tau$ and $0\le i<m_\tau$, form a basis of
  $\operatorname{Lie}_\tau(x)$ for every $\tau$ if and only if together they form a basis of $\operatorname{Lie}(x)$.
\item If $S$ is a field, $\operatorname{Lie}_\tau(x)$ has dimension $m_\tau$ for every $\tau$.
\item Every point of $X$ has an affine open neighbourhood $\operatorname{Spec}R$, with $R$ of finite type over
  $F_c$, such that for the inclusion $\iota$ of $\operatorname{Spec}R$ every $\operatorname{Lie}_\tau(\iota)$ has a
  basis over $R$.
\end{enumerate}
\end{lemma}
\leanref{Development/AbelianVarietyTheory/Proofs/UnitaryExistence/IsTangentVectorBasisAtPointBaseChange.lean}{GlobalTheory.Frames.tanBasis_baseChange}
\leanref{Development/AbelianVarietyTheory/Proofs/UnitaryExistence/EigenBasisIffTangentBasis.lean}{GlobalTheory.Frames.eigBasis_iff_tanBasis}
\leanref{Formalism/LocalToGlobal/ProjectiveLattices/GenFrPfB.lean}{Formalism.L2.Pj.gfp_finrank}
\leanref{Development/AbelianVarietyTheory/Proofs/UnitaryExistence/ExistsChartEigenframe.lean}{GlobalTheory.ModuliChart.exists_eigen_chart}
\begin{proof}
(1) The abelian scheme $A_x$ is smooth over $S$, so its sheaf of differentials is finite locally free
\cite[Tag 02G1]{Sta}, and its formation commutes with base change \cite[Tag 01V0]{Sta}. The Lie algebra is the
$S$-dual of the pullback $\omega$ of the differentials along the zero section, a finite locally free $S$-module, and
the Lie algebra over $T$ is the $T$-dual of $T\otimes_S\omega$, which is $T\otimes_S\operatorname{Lie}(x)$. So a basis
goes to a basis.

(2) Every embedding of $E$ takes values in $F_c$, and the discriminant of $E$ is invertible in $F_c$, so
$\mathcal O_E\otimes_{\mathbf Z}S$ is the product of copies of $S$, one for each embedding. Its idempotents
decompose every module over it uniquely into the submodules on which $\mathcal O_E$ acts through the embeddings;
so $\operatorname{Lie}(x)$ is the direct sum of the $\operatorname{Lie}_\tau(x)$. Hence the family $u$ spans
$\operatorname{Lie}(x)$ if and only if, for every $\tau$, the $u_{\tau,i}$ span $\operatorname{Lie}_\tau(x)$; and a
linear relation among all the $u_{\tau,i}$ is a sum of relations among the $u_{\tau,i}$ of one $\tau$ at a time.

(3) By (1) and (2), the $\tau$-eigenspace of the Lie algebra after an extension of the field $S$ is the extension of
$\operatorname{Lie}_\tau(x)$, so its dimension does not change. Over an algebraic closure of $S$ it is $m_\tau$, by
the condition of the moduli problem on the Lie algebra (Definition~\ref{cv-volume-twisting-section}(3)).

(4) Let $p$ be the point. As in the proof of (1), the Lie algebra of the universal abelian scheme over an affine
open set is the dual of a finite locally free module of rank $g$, the relative dimension. So there are an affine
open neighbourhood $\operatorname{Spec}R$ of $p$, of finite type over $F_c$, and a basis $v_1,\dots,v_g$ of
$\operatorname{Lie}(\iota)$ over $R$. In this
basis $\mathcal O_E$ acts through matrices over $R$, and by (2) each $\operatorname{Lie}_\tau(\iota)$ is the image
of an idempotent. Let $k$ be an algebraic closure of the residue field of $p$. By (1) and (2) the
$\tau$-eigenspace of the Lie algebra over $k$ is spanned by the images of $\operatorname{Lie}_\tau(\iota)$, and by
(3) it has dimension $m_\tau$. Choose for every $\tau$ vectors $u_{\tau,i}\in\operatorname{Lie}_\tau(\iota)$,
$0\le i<m_\tau$, whose images form a basis of this eigenspace; by (2) all their images form a basis over $k$. Let
$r\in R$ be the determinant of their coordinates in the basis $v$. Its image in $k$ is nonzero, so $r$ does not
vanish at $p$. By (1) the Lie algebra over $R_r$ has the basis of the images of the $v_j$, in which the images of
the $u_{\tau,i}$ have an invertible matrix of coordinates; so they form a basis, and by (2) bases of the
eigenspaces. Replace $R$ by $R_r$.
\end{proof}
% Lean: (1) is tanBasis_baseChange, through lieBasis_iff_tangentBasis_pullback_universal and
% exists_isLieBasis_baseChange (LieBaseChange.isLieBasis_bcVec, from the bijection expBij between J^n and the
% points of the group over B that reduce to the unit modulo a square-zero ideal J); the text cites the
% differentials instead. (2) is eigBasis_iff_tanBasis (existsUnique_cmEmbedding_eigenvector_sum, basis_iff_of_split).
% (3) is gfp_finrank (finrank_eig_bc for the algebraic closure and the field geom of the universal object for the
% dimension). (4) is exists_eigen_chart: LieLoc.lie_locally_free for the basis v (from a standard smooth
% presentation, exists_univChart_away; the text uses (1) and Nakayama's lemma), then exists_isLieBasis_eigen
% (LieEigen.exists_eigenvectors and LieEigen.exists_basis_of_isUnit, with EigenSplit.isSplitting_of_isDomain and
% the field geom of the universal object for the dimensions).

\begin{lemma}\label{cv-volume-key-base-change}
In the setting of Definition~\ref{cv-volume-twisting-section}, let $y$ be a framed point of $X$ over an
$F_c$-algebra $S$ and $\varphi\colon S\to T$ a homomorphism of $F_c$-algebras. Then $y$ has an \emph{image} under
$\varphi$: a framed point over $T$ whose point is the point of $y$ composed with $\operatorname{Spec}\varphi$,
whose basis vectors are the images of those of $y$, and whose section is the pullback of the section of $y$.
\end{lemma}
\leanref{Development/AbelianVarietyTheory/Proofs/UnitaryExistence/ExistsFramedPointBaseChangeAlgebraMap.lean}{GlobalTheory.Frames.FramedBaseChange.exists_frMap}
\begin{proof}
The images of the basis vectors are tangent vectors in the same eigenspaces, as the action of $\mathcal O_E$
commutes with base change. By Lemma~\ref{cv-volume-key-lie}(2) the vectors of $y$ together form a basis of the
Lie algebra over $S$; by part (1) of that lemma their images form a basis over $T$, and by part (2) again the
images are bases of the eigenspaces. The section of $y$ generates the sections of the line bundle of the
polarization over $S$, so it vanishes on no fibre (Lemma~\ref{cv-volume-key-generator}(3)). By part (1) of that
lemma its pullback vanishes on no fibre over $T$, and by part (3) the pullback generates the sections over $T$.
\end{proof}
% Lean: exists_frMap: existsUnique_scaledBaseChange_fibreSection for the pullback of the section, tan_bc and eig_bc,
% framedPoint_eigenbasis_baseChange (eigBasis_iff_tanBasis, tanBasis_baseChange) for the bases,
% fibreSectionNowhereZero_of_generates, secBC_nonvanishing_iff and fibreSectionGenerates_of_nowhereZero. It assumes
% the base change law SecBaseChangeSt (Theorem cv-volume-key-sections).

\begin{lemma}\label{cv-volume-key-generic-points}
In the setting of Definition~\ref{cv-volume-twisting-section}, let $Z$ be a connected component of $X$, let
$\kappa_Z$ be the residue field of its generic point $\eta_Z$, and let $\Omega_Z$ be the field of fractions of the
polynomial ring over $\kappa_Z$ in variables $x_{\tau,ij}$, for the embeddings $\tau\colon E\to\mathbf C$ and
$0\le i<j<m_\tau$. Then there are a framed point $y^0_Z$ over $\kappa_Z$ with the point $\eta_Z$ and a framed point
$y_Z$ over $\Omega_Z$ with the point $\eta_Z$ whose vector of index $i$ in the eigenspace of $\tau$ is
$v_{\tau,i}+\sum_{j>i}x_{\tau,ij}v_{\tau,j}$, where $v_{\tau,j}$ are the images of the vectors of $y^0_Z$. Such a
$y_Z$ is a \emph{generic framed point} of $Z$.
\end{lemma}
\leanref{Formalism/LocalToGlobal/ProjectiveLattices/GenFrPf.lean}{Formalism.L2.Pj.genFrEx_of}
\leanref{Formalism/LocalToGlobal/ProjectiveLattices/IsGenFrGenFr.lean}{Formalism.L2.Pj.IsGenFr}
\begin{proof}
Let $K$ be $\kappa_Z$ or $\Omega_Z$, and $\eta_K$ the point $\eta_Z$ over $K$. The sections of the line bundle of the
polarization over the abelian variety of $\eta_K$ form a free $K$-module of rank one, by
Theorem~\ref{cv-volume-key-sections} for $n=1$ and $S=\operatorname{Spec}K$; choose a generator $\vartheta_K$.

By Lemma~\ref{cv-volume-key-lie}(3), each eigenspace $\operatorname{Lie}_\tau(\eta_Z)$ over $\kappa_Z$ has
dimension $m_\tau$. Choose bases $v^0_{\tau,j}$, $0\le j<m_\tau$; with $\vartheta_{\kappa_Z}$ they form the framed
point $y^0_Z$. Let $v_{\tau,j}$ be their images over $\Omega_Z$. By parts (1) and (2) of
Lemma~\ref{cv-volume-key-lie} they are linearly independent, and they lie in the $\tau$-eigenspace over
$\Omega_Z$, which has dimension $m_\tau$ by part (3) of that lemma. The vectors $w_{\tau,i}=v_{\tau,i}+\sum_{j>i}x_{\tau,ij}v_{\tau,j}$ lie in the same eigenspace and
are obtained from the $v_{\tau,j}$ by a unitriangular matrix, so they are $m_\tau$ linearly independent vectors of
it and form a basis. With $\vartheta_{\Omega_Z}$ they form the framed point $y_Z$.
\end{proof}
% Lean: genFrEx_of: gfp_frPt_gen over the tower kap Z -> Om Z (gfp_tower) at the generic point xGen Z, with the
% coefficients genCoef (1 at i, the variable X (i, j) at j > i, 0 below). gfp_frPt_gen works in a chart of the
% tangent space (univChartLoc, gfp_bc): bases bas of the eigenspaces gfpEig of dimension Dw.mult tau (gfp_finrank,
% Lemma cv-volume-key-lie(3)), their images (li_bc_lemma for the independence), the unitriangular combinations
% (gfp_li_tri) and gfp_basis; the generators are gfp_gen over each field, from globalSections_rank_one_fibre and
% the hypothesis SecTwistSt, which the text replaces by Theorem cv-volume-key-sections. The components are the
% irreducible components of Xo.XF c.

\begin{definition}\label{cv-volume-key-chart}
In the setting of Definition~\ref{cv-volume-twisting-section}, let $I$ be the set of pairs $(\tau,i)$ of an
embedding $\tau$ of $E$ and an integer $0\le i<m_\tau$, together with one further index $0$. For framed points $y$
and $y'$ of $X$ over an $F_c$-algebra $S$ and $t\in(S^\times)^I$, write $y'=y\cdot t$ if $y'$ has the point of $y$,
if its vector of index $i$ in the eigenspace of $\tau$ is a combination of the vectors of $y$ in that eigenspace
with index at most $i$, with the coefficient $t_{\tau,i}$ at index $i$, and if its section is $t_0$ times that of
$y$. Choose a generic framed point $y_Z$ of every connected component $Z$ of $X$
(Lemma~\ref{cv-volume-key-generic-points}). A \emph{chart} $C$ consists of a connected component $Z_C$, a smooth
$F_c$-subalgebra $R_C$ of $\Omega_{Z_C}$ whose field of fractions is $\Omega_{Z_C}$, a framed point $y_C$ over $R_C$
and $t_C\in(\Omega_{Z_C}^\times)^I$ such that an image of $y_C$ over $\Omega_{Z_C}$ is $y_{Z_C}\cdot t_C$, with the
following properties.
\begin{enumerate}
\item If $\psi,\psi'\colon R_C\to L$ are homomorphisms of $F_c$-algebras to a field, and $z$ and $z'$ are images of
  $y_C$ under them with $z'=z\cdot t$ for some $t$, then $\psi=\psi'$.
\item Let $w$ be a framed point over an $F_c$-algebra $S$, and $\chi\colon S\to L$ a homomorphism of
  $F_c$-algebras to a field such that an image of $w$ under $\chi$ is $z\cdot t$ for an image $z$ of $y_C$ under
  some homomorphism $R_C\to L$ and some $t$. Then there is $g\in S$ with $\chi(g)\ne0$ such that, for every
  homomorphism $\lambda\colon S\to S'$ of $F_c$-algebras with $\lambda(g)$ a unit, every image of $w$ under $\lambda$
  is $z'\cdot t'$ for an image $z'$ of $y_C$ under some homomorphism $R_C\to S'$ and some $t'$.
\end{enumerate}
\end{definition}
\leanref{Development/AbelianVarietyTheory/Definitions/RegularCharts/Reframe.lean}{GlobalTheory.Reg.FChart}

\begin{lemma}\label{cv-volume-key-local-frames}
In the setting of Definition~\ref{cv-volume-twisting-section}, every point of $X$ has an affine open
neighbourhood $\operatorname{Spec}A\subset X$, with $A$ a smooth $F_c$-algebra that is a domain, such that the
inclusion $\operatorname{Spec}A\to X$ is the point of a framed point $y_A$ of $X$ over $A$.
\end{lemma}
\leanref{Development/AbelianVarietyTheory/Proofs/GenericRegularity/ExistsIntegralSmoothChartFrameModuliStatement.lean}{GlobalTheory.T.exists_integral_smooth_chart_frame_moduli}
\begin{proof}
Let $p$ be the point. By Lemma~\ref{cv-volume-key-lie}(4) there are an affine open neighbourhood
$\operatorname{Spec}R_1$ of $p$, with $R_1$ of finite type over $F_c$, and bases of the eigenspaces
$\operatorname{Lie}_\tau$ for its inclusion. By Lemma~\ref{cv-volume-key-generator}(2), for this inclusion and the
prime $p$, there are $h\in R_1$ with $h(p)\ne0$ and a section of the line bundle of the polarization over
$R=(R_1)_h$ that vanishes on no fibre. By parts (1) and (2) of Lemma~\ref{cv-volume-key-lie} the bases restrict
to bases of the eigenspaces over $R$. As $X$ is smooth over $F_c$, so is $R$, and its local ring at $p$ is
regular, hence a domain. As $R$ is noetherian, the kernel of $R\to R_p$ is generated by finitely many elements,
each killed by some element of $R$ outside $p$; for their product $g$, $R_g$ embeds in $R_p$ and is a domain, and
$g(p)\ne0$. Put $A=R_g$. The bases restrict to bases of the eigenspaces over $A$, as before, and the section
restricts to a section that vanishes on no fibre (Lemma~\ref{cv-volume-key-generator}(1)), so it generates the
sections over $A$ (part (3) of that lemma). These data form the framed point $y_A$.
\end{proof}
% Lean: LocalFrame.exists_integral_chart_frame: ModuliChart.exists_chart_eigenframe for R, the bases and the
% section (PushInv.pushInvertible and exists_sec_of_secOn for the section, which vanishes on no fibre, then
% exists_eigen_chart, Lemma cv-volume-key-lie(4), localizing further; the text takes the eigenbases first and then
% localizes for the section, with the same two inputs), SecTransport.exists_sec_nonzero, smooth_of_chart,
% isRegularRing_of_smooth, isDomain_of_isRegularLocalRing, exists_away_isDomain for g, tanBasis_baseChange with
% eigBasis_iff_tanBasis for the bases, and FrameGenerator.secGen_of_secNZ_baseChange for the generator. Its
% hypothesis SecTwistSt is used only through LocalFrame.secBaseChange_of_secTwist, the Lean's derivation of the
% base change law SecBaseChangeSt (about 22,900 of the lemma's 27,800 lines), which the text cites as Theorem
% cv-volume-key-sections.

\begin{lemma}\label{cv-volume-key-reframe}
In the setting of Definition~\ref{cv-volume-twisting-section}, let $y$ be a framed point of $X$ over an
$F_c$-algebra $S$, with basis vectors $v_{\tau,i}$.
\begin{enumerate}
\item For every family $g=(g_\tau)$ of matrices $g_\tau\in GL_{m_\tau}(S)$ there is a framed point with the point
  and the section of $y$ whose vector of index $i$ in the eigenspace of $\tau$ is $\sum_j(g_\tau)_{ji}v_{\tau,j}$.
\item If $y'$ is a framed point over $S$ with the point of $y$, there is a unique family $g=(g_\tau)$ of
  invertible matrices such that the vector of index $i$ of $y'$ in the eigenspace of $\tau$ is
  $\sum_j(g_\tau)_{ji}v_{\tau,j}$ for all $\tau$ and $i$.
\item Every $g\in GL_n(L)$, for a field $L$, is $g=ub$ with $b$ upper triangular with nonzero diagonal entries and
  $u$ a matrix such that, for some permutation $\sigma$ of $\{0,\dots,n-1\}$, $u_{\sigma(i),i}=1$ and
  $u_{\sigma(i),j}=0$ for $j>i$.
\end{enumerate}
\end{lemma}
\leanref{Development/AbelianVarietyTheory/Proofs/GenericRegularity/IsTriangularReframeExistsAndUniqueStatement.lean}{GlobalTheory.T.isTriangularReframe_exists_and_unique}
\leanref{Development/AbelianVarietyTheory/Proofs/GenericRegularity/FlagChartsCoverFramedFieldPoints.lean}{GlobalTheory.Reg.exists_perm_lower_upper}
\begin{proof}
(1) The combinations $\sum_j(g_\tau)_{ji}v_{\tau,j}$ are tangent vectors in the eigenspace of $\tau$. They form a
basis of it: a vector of the eigenspace has coordinates $a$ in the $v_{\tau,j}$, and its coordinates in the new
family are $g_\tau^{-1}a$.

(2) The coordinates of the vectors of $y'$ in the bases of $y$ give $g$, and they are unique. The coordinates of
the vectors of $y$ in the bases of $y'$ give matrices $h_\tau$ with $g_\tau h_\tau=1$, so $g_\tau$ is invertible.

(3) By induction on $k$, the first $k$ columns of $g$ are $ub$ for a matrix $u$
with $k$ columns, rows $\sigma(0),\dots,\sigma(k-1)$ as required, and $b$ upper triangular of size $k$ with
nonzero diagonal entries. Given this for $k$, the matrix of the entries of $u$ in the rows $\sigma(l)$ is lower
triangular with ones on the diagonal, so there is a combination $uc$ of the columns of $u$ that agrees with the
column $k$ of $g$ in the rows $\sigma(0),\dots,\sigma(k-1)$. The difference $v$ of the column $k$ of $g$ and $uc$
vanishes in these rows, and it is not zero, as the columns of $g$ are linearly independent and the columns of
$u$ are combinations of the first $k$ of them. Let $\sigma(k)$ be a row with $v_{\sigma(k)}\ne0$. The new column of
$u$ is $v/v_{\sigma(k)}$, which vanishes in the rows $\sigma(l)$, $l<k$, and the new column of $b$ is $c$ above the
entry $v_{\sigma(k)}$.
\end{proof}
% Lean: (1) and (2) are ReframeEx.exists_reframe and existsUnique_reframe (comb_exists, comb_unique,
% TangentLinear.isCombAt_mulVec); (3) is exists_perm_lower_upper through exists_integral_pivot_decomposition
% (exists_pivot_unipotent_mul_triangular, here with the valuation ring W = L).

\begin{lemma}\label{cv-volume-key-cells}
Let $S$ be a ring, $m\ge0$ an integer and $\sigma$ a permutation of $\{0,\dots,m-1\}$. For a family
$x=(x_{ij})_{0\le i<j<m}$ of elements of $S$ let $n(x)\in M_m(S)$ be the lower unitriangular matrix with the entry
$x_{ij}$ in row $j$ and column $i$, and let $u_\sigma(x)$ be the matrix whose row $\sigma(j)$ is the row $j$ of
$n(x)$ for every $j$; these are the matrices $u$ with $u_{\sigma(i),i}=1$ and $u_{\sigma(i),j}=0$ for $j>i$. Let
$g\in M_m(S)$.
\begin{enumerate}
\item $g=u_\sigma(x)b$ for some $x$ and some upper triangular $b$ with units on the diagonal if and only if, for
  $1\le k\le m$, the minor of $g$ on the rows $\sigma(0),\dots,\sigma(k-1)$ and the columns $0,\dots,k-1$ is a
  unit.
\item If $u_\sigma(x)b=u_\sigma(x')b'$ with $b$ and $b'$ upper triangular and units on the diagonal of $b$, then
  $x=x'$ and $b=b'$.
\end{enumerate}
\end{lemma}
\leanref{Development/AbelianVarietyTheory/Proofs/GenericRegularity/ExistsLowerMulUpperOfLeadingMinors.lean}{GlobalTheory.cell_coordinates_exists_unique}
\begin{proof}
Let $h$ be the matrix whose row $i$ is the row $\sigma(i)$ of $g$. Then $g=u_\sigma(x)b$ if and only if
$h=n(x)b$, and the minors in (1) are the leading principal minors of $h$; so we may take $\sigma$ to be the
identity.

(1) If $h=lb$ with $l$ lower unitriangular and $b$ upper triangular, the leading $k\times k$ block of $h$ is the
product of those of $l$ and of $b$, so its determinant is the product of the first $k$ diagonal entries of $b$, a
unit. Conversely, suppose that the leading principal minors of $h$ are units; we argue by induction on $m$, the
case $m=0$ being empty. The leading block $h'$ of size $m-1$ is $l'b'$ with $l'$ lower unitriangular and $b'$ upper
triangular, by induction, and $\det b'=\det h'$ is a unit. Let $c$ be the last row of $h$ and $d$ its last column,
without their last entries, and $h_{mm}$ the last entry. Then $h=lb$ with
\[
 l=\begin{pmatrix}l'&0\\cb'^{-1}&1\end{pmatrix},\qquad
 b=\begin{pmatrix}b'&l'^{-1}d\\0&h_{mm}-cb'^{-1}l'^{-1}d\end{pmatrix}.
\]
As $\det b=\det h$ is a unit and $b$ is upper triangular, its diagonal entries are units.

(2) The diagonal entries of $b$ are units, so $b$ is invertible, and $n(x')^{-1}n(x)=b'b^{-1}$ is both lower
unitriangular and upper triangular, hence $1$. So $n(x)=n(x')$ and $b=b'$.
\end{proof}
% Lean: cell_coordinates_exists_unique, for an embedding iota of Fin r in Fin n and an n x r matrix; the text states
% the square case, which is the one used (through minor_iff and cellM_unique, with cellM sigma x = u_sigma(x)).
% Its proof is the text's: exists_lower_mul_upper_of_leading_minors (the induction above),
% isUnit_det_of_upper_of_isUnit_diag and eq_one_of_lower_unitriangular_eq_mul.

\begin{lemma}\label{cv-volume-key-cell-change}
Let $\kappa$ be an infinite field and $T$ a finite set, and for $\tau\in T$ let $m_\tau\ge0$ be an integer,
$G_\tau\in GL_{m_\tau}(\kappa)$ and $\sigma_\tau$ a permutation of $\{0,\dots,m_\tau-1\}$. Let $\Omega$ be the field
of fractions of the polynomial ring over $\kappa$ in variables $x_{\tau,ij}$, for $\tau\in T$ and
$0\le i<j<m_\tau$, and write $x_\tau=(x_{\tau,ij})_{i<j}$. Then there are an automorphism $e$ of $\Omega$ over
$\kappa$ and upper triangular matrices $b_\tau\in M_{m_\tau}(\Omega)$ with nonzero diagonal entries such that
\[
 G_\tau\,n(x_\tau)=u_{\sigma_\tau}\big(e(x_\tau)\big)\,b_\tau\qquad(\tau\in T),
\]
in the notation of Lemma~\ref{cv-volume-key-cells}.
\end{lemma}
\leanref{Development/AbelianVarietyTheory/Proofs/GenericRegularity/ExistsFlagCoordinatesCellChangeBirationalStatement.lean}{GlobalTheory.T.exists_flagCoordinates_cellChange_birational}
\begin{proof}
Let $A_\tau$ be the matrix whose row $i$ is the row $\sigma_\tau(i)$ of $G_\tau$; it is invertible.

The leading principal minors of $A_\tau n(x_\tau)$ are nonzero. Fix $\tau$, write $A=A_\tau$ and $m=m_\tau$, and let
$s$ be a variable. The matrix $P=s+A^{-1}$ over $\kappa[s]$ has as leading principal minors the characteristic
polynomials of the leading blocks of $-A^{-1}$, which are monic, hence nonzero; and $AP=1+sA$ has leading
principal minors whose value at $s=0$ is $1$, hence nonzero. By Lemma~\ref{cv-volume-key-cells}(1) over the field
$\kappa(s)$, $P=lc$ and $1+sA=l'c'$ with $l,l'$ lower unitriangular and $c,c'$ upper triangular with nonzero
diagonal entries. Then $Al=APc^{-1}=l'c'c^{-1}$, so by the same lemma the leading principal minors of $Al$ are
nonzero. The homomorphism of $\kappa$-algebras from $\kappa[x_\tau]$ to $\kappa(s)$ that sends $x_\tau$ to the
entries of $l$ below the diagonal, so that $n(x_\tau)$ goes to $l$, takes the leading principal minors of
$An(x_\tau)$ to those of $Al$; so they are nonzero.

The new coordinates. By Lemma~\ref{cv-volume-key-cells}(1) over $\Omega$, $A_\tau n(x_\tau)=l_\tau b_\tau$ with $l_\tau$
lower unitriangular and $b_\tau$ upper triangular with nonzero diagonal entries. Let $y_{\tau,ij}$ be the entry of
$l_\tau$ in row $j$ and column $i$, and $F\subset\Omega$ the subfield generated over $\kappa$ by all the
$y_{\tau,ij}$. The matrix $A_\tau^{-1}l_\tau$ has entries in $F$ and equals $n(x_\tau)b_\tau^{-1}$, so by
Lemma~\ref{cv-volume-key-cells}(1) its leading principal minors are nonzero; by the same lemma over $F$ it is
$l''b''$ with $l''$ lower unitriangular over $F$, and by part (2) over $\Omega$, $n(x_\tau)=l''$. So every
$x_{\tau,ij}$ lies in $F$, and $F=\Omega$. The $y_{\tau,ij}$ are as many as the variables and generate $\Omega$, whose
transcendence degree over $\kappa$ is the number of variables, so they are algebraically independent. Hence
$x_{\tau,ij}\mapsto y_{\tau,ij}$ extends to an embedding $e$ of $\Omega$ into itself over $\kappa$, which is onto, as
its image contains the $y_{\tau,ij}$. Finally the row $\sigma_\tau(i)$ of $G_\tau n(x_\tau)$ is the row $i$ of
$l_\tau b_\tau$, which is the row $\sigma_\tau(i)$ of $u_{\sigma_\tau}(e(x_\tau))b_\tau$.
\end{proof}
% Lean: exists_flagCoordinates_cellChange_birational = Reg.cell_change_birational: A = G.submatrix sigma id,
% CellChange.generic_minors (through exists_lower_minors, with charmatrix (-A^{-1}) and 1 + X A), minors_iff for
% l and b, lower_mem_range with upper_inv for F = Omega, and exists_algEquiv_of_generates (a transcendence basis
% count) for e. The hypothesis that kappa is infinite is not used in the proof.

\begin{lemma}\label{cv-volume-key-cell-charts}
In the setting of Definition~\ref{cv-volume-twisting-section}, let $\operatorname{Spec}A\subset X$ be an affine open
set with $A$ a smooth $F_c$-algebra that is a domain, $y_A$ a framed point of $X$ over $A$ whose point is the
inclusion, and $\sigma=(\sigma_\tau)$ a family of permutations of the sets $\{0,\dots,m_\tau-1\}$. Then there is a
chart $C$ (Definition~\ref{cv-volume-key-chart}) with the following property. Let $L\supset F_c$ be a field,
$f\colon A\to L$ a homomorphism of $F_c$-algebras, $z_0$ an image of $y_A$ under $f$, and $y$ a framed point over $L$
with the point of $z_0$ whose vectors are obtained from those of $z_0$ by matrices $g_\tau=u_\tau b_\tau$ as in
Lemma~\ref{cv-volume-key-reframe}(2), where $b_\tau$ is upper triangular with nonzero diagonal entries and
$(u_\tau)_{\sigma_\tau(i),i}=1$ and $(u_\tau)_{\sigma_\tau(i),j}=0$ for $j>i$. Then $y=z\cdot t$ for an image $z$ of
$y_C$ under some homomorphism $R_C\to L$ and some $t\in(L^\times)^I$.
\end{lemma}
\leanref{Development/AbelianVarietyTheory/Proofs/GenericRegularity/ExistsFlagCellChartOverFramedAffineOpenStatement.lean}{GlobalTheory.T.exists_flagCellChart_over_framedAffineOpen}
\begin{proof}
Reframing. Say that a framed point $y'$ is obtained from $y$ by $g=(g_\tau)$ if both have the same point and
the vectors are related as in Lemma~\ref{cv-volume-key-reframe}(2). If $y'$ is obtained from $y$ by $g$ and $y''$
from $y'$ by $g'$, then $y''$ is obtained from $y$ by $(g_\tau g'_\tau)$; and if $y'$ is obtained from $y$ by $g$,
their images under a homomorphism $\varphi$ are obtained from one another by $(\varphi(g_\tau))$. An image of an
image is an image, and two images of a framed point under one homomorphism agree. Two framed points $y,y'$ over a
ring with the same point satisfy $y'=y\cdot t$ for some $t$ if and only if $y'$ is obtained from $y$ by upper
triangular matrices with units on the diagonal, since their sections generate the same free module of rank one
and so differ by a unit. In the notation of Lemma~\ref{cv-volume-key-cells}, the matrices $u_\tau$ of the statement
are $u_{\sigma_\tau}(x'_\tau)$ with $x'_{\tau,ij}=(u_\tau)_{\sigma_\tau(j),i}$; write $u_\sigma(x')$ for the family.

The component. Let $Z$ be the connected component of $X$ that contains $\operatorname{Spec}A$; as $A$ is a domain,
the generic point $\eta_Z$ lies in $\operatorname{Spec}A$. Write $\kappa=\kappa_Z$ and $\Omega=\Omega_Z$, and let
$f_0\colon A\to\kappa$ be the homomorphism given by $\eta_Z$; it is injective, and $\kappa$ is the field of fractions
of its image. Let $y_Z$ be the chosen generic framed point and $y^0_Z$ the framed point over $\kappa$ of
Lemma~\ref{cv-volume-key-generic-points} for it, so that $y_Z$ is obtained from the image of $y^0_Z$ over $\Omega$
by the matrices $n(x_\tau)$.

The matrix $G$. Let $z_\kappa$ be the image of $y_A$ under $f_0$ (Lemma~\ref{cv-volume-key-base-change}). It has the
point $\eta_Z$, as $y^0_Z$ does, so $y^0_Z$ is obtained from $z_\kappa$ by invertible matrices $G_\tau$ over
$\kappa$ (Lemma~\ref{cv-volume-key-reframe}(2)). The field $\kappa$ contains $F_c$, so it is infinite, and
Lemma~\ref{cv-volume-key-cell-change} for $G$ and $\sigma$ gives an automorphism $e$ of $\Omega$ over $\kappa$ and
upper triangular $b_\tau$ with nonzero diagonal entries such that $G_\tau n(x_\tau)=u_{\sigma_\tau}(e(x_\tau))b_\tau$.

The ring $R_C$. Let $A[x]$ be the polynomial ring over $A$ in the variables $x_{\tau,ij}$, and
$\theta\colon A[x]\to\Omega$ the homomorphism of $F_c$-algebras with $\theta(a)=f_0(a)$ for $a\in A$ and
$\theta(x_{\tau,ij})=e(x_{\tau,ij})$. As $e$ fixes $\kappa$, $\theta(A[x])$ is the image under $e$ of $f_0(A)[x]$, so
$\theta$ is injective and the field of fractions of its image $R_C$ is $e(\Omega)=\Omega$. The ring $R_C$ is
isomorphic to $A[x]$, so it is smooth over $F_c$. Let $c\colon A\to R_C$ be the restriction of $\theta$.

The framed point $y_C$. Let $u_0$ be the image of $y_A$ under $c$, and $y_C$ the framed point over $R_C$ obtained
from $u_0$ by the matrices $u_{\sigma_\tau}(e(x_\tau))$, whose determinants are $\pm1$
(Lemma~\ref{cv-volume-key-reframe}(1)). Let $\psi\colon R_C\to S$ be a homomorphism of $F_c$-algebras, $z$ an image
of $y_C$ under $\psi$ and $z_0$ the image of $y_A$ under $\psi\circ c$. The image of $u_0$ under $\psi$ is $z_0$, so
\[
 z\text{ is obtained from }z_0\text{ by the matrices }u_{\sigma_\tau}\big(\psi(e(x_\tau))\big).\tag{$*$}
\]

The family $t_C$. Let $y_g$ be the image of $y_C$ over $\Omega$. By $(*)$, $y_g$ is obtained by the matrices
$u_{\sigma_\tau}(e(x_\tau))$ from the image of $y_A$ under $\theta\circ c$, which is the image $z_\Omega$ of $z_\kappa$
over $\Omega$. The image of $y^0_Z$ over $\Omega$ is obtained from $z_\Omega$ by $G$, so $y_Z$ is obtained from
$z_\Omega$ by $G_\tau n(x_\tau)=u_{\sigma_\tau}(e(x_\tau))b_\tau$. By Lemma~\ref{cv-volume-key-reframe}(2), $y_Z$ is
obtained from $y_g$ by $b$, so $y_Z=y_g\cdot t$ for some $t$, and $y_g=y_Z\cdot t_C$ with $t_C=t^{-1}$. So $Z$,
$R_C$, $y_C$ and $t_C$ have the form required of a chart.

The stated property, over any $F_c$-algebra $S$ in place of $L$, with units on the diagonals of the upper
triangular matrices. Let $f\colon A\to S$, $z_0$ the image of $y_A$ under $f$, and $y$ obtained from $z_0$ by the
matrices $u_{\sigma_\tau}(x'_\tau)b'_\tau$, with $b'_\tau$ upper triangular with units on the diagonal. Let
$\psi\colon R_C\to S$ be the homomorphism with $\psi\circ c=f$ and $\psi(e(x_{\tau,ij}))=x'_{\tau,ij}$, and $z$ an
image of $y_C$ under $\psi$ (Lemma~\ref{cv-volume-key-base-change}). By $(*)$, $z$ is obtained from $z_0$ by
$u_\sigma(x')$, so $y$ is obtained from $z$ by $b'$, and $y=z\cdot t$ for some $t$. Over a field the nonzero
diagonal entries are units, which is the statement.

Property (1). Let $\psi,\psi'\colon R_C\to L$ and images $z$, $z'$ of $y_C$ with $z'=z\cdot t$. The points of $z$ and
$z'$ agree, and they are the points $\operatorname{Spec}L\to\operatorname{Spec}A\to X$ given by $\psi\circ c$ and
$\psi'\circ c$; as $\operatorname{Spec}A\to X$ is an open immersion, $\psi\circ c=\psi'\circ c$. Let $z_0$ be the image
of $y_A$ under this homomorphism. By $(*)$, $z$ and $z'$ are obtained from $z_0$ by
$u_\sigma(\psi(e(x)))$ and $u_\sigma(\psi'(e(x)))$, and $z'$ is obtained from $z$ by upper triangular matrices
$b'_\tau$ with units on the diagonal. By Lemma~\ref{cv-volume-key-reframe}(2),
$u_{\sigma_\tau}(\psi(e(x_\tau)))b'_\tau=u_{\sigma_\tau}(\psi'(e(x_\tau)))$, and by Lemma~\ref{cv-volume-key-cells}(2),
$\psi(e(x_{\tau,ij}))=\psi'(e(x_{\tau,ij}))$. As $R_C$ is generated by $c(A)$ and the $e(x_{\tau,ij})$, $\psi=\psi'$.

Property (2). Let $w$ be a framed point over $S$ and $\chi\colon S\to L$ such that an image $w_L$ of $w$ under $\chi$
is $z\cdot t$ for an image $z$ of $y_C$ under some $\psi\colon R_C\to L$. The point of $w_L$ is that of $z$, which
factors through $\operatorname{Spec}A$. So the preimage of $\operatorname{Spec}A$ under the point of $w$ is an open
subset of $\operatorname{Spec}S$ that contains the prime $\ker\chi$, and it contains a
principal open set $D(g_1)$ with $\chi(g_1)\ne0$. Over $S_1=S_{g_1}$ the point of $w$ factors through
$\operatorname{Spec}A$, by a homomorphism $f_1\colon A\to S_1$. Let $w_1$ be the image of $w$ over $S_1$ and $z_1$
the image of $y_A$ under $f_1$; they have the same point, so $w_1$ is obtained from $z_1$ by invertible matrices
$G'_\tau$ over $S_1$. Let $\chi_1\colon S_1\to L$ extend $\chi$. Then $w_L$ is obtained by $\chi_1(G')$ from the image
of $y_A$ under $\chi_1\circ f_1$, and $\chi_1\circ f_1=\psi\circ c$, as both give the point of $w_L$ and
$\operatorname{Spec}A\to X$ is an open immersion. By $(*)$, $w_L$ is also obtained from that image by matrices
$u_{\sigma_\tau}(\psi(e(x_\tau)))b'_\tau$ with $b'_\tau$ upper triangular with units on the diagonal. By
Lemma~\ref{cv-volume-key-reframe}(2) these are the $\chi_1(G'_\tau)$, so by Lemma~\ref{cv-volume-key-cells}(1) the
minors of $\chi_1(G'_\tau)$ named there, for $\sigma_\tau$, are nonzero. Write the product of these minors of the
$G'_\tau$, over all $\tau$ and $k$, as $g_2/g_1^N$ with $g_2\in S$, and put $g=g_1g_2$. Then $\chi(g)\ne0$. Let
$\lambda\colon S\to S'$ be a homomorphism of $F_c$-algebras with $\lambda(g)$ a unit, and $w'$ an image of $w$ under
$\lambda$. Then $\lambda$ factors through $S_1$, by $\nu\colon S_1\to S'$, and $\nu$ takes the minors of the $G'_\tau$
to units; by Lemma~\ref{cv-volume-key-cells}(1), $\nu(G'_\tau)=u_{\sigma_\tau}(x''_\tau)b''_\tau$ with $b''_\tau$ upper
triangular with units on the diagonal. Now $w'$ is the image of $w_1$ under $\nu$, which is obtained by $\nu(G')$
from the image of $y_A$ under $\nu\circ f_1$. By the stated property over $S'$, $w'=z'\cdot t'$ for an image $z'$ of
$y_C$ under some homomorphism $R_C\to S'$ and some $t'$.
\end{proof}
% Lean: GlobalTheory.Reg.exists_fChart_cell. The component is closure_range_mem with exists_residue_embed and
% exists_generic for f0; y0 is from isGenFr_genFr; FrBC (Lemma cv-volume-key-base-change) gives the images; G is
% isTriangularReframe_exists_and_unique; e and b are exists_flagCoordinates_cellChange_birational; theta is
% exists_poly_embed (polyChart_smooth_fraction_lift for the fraction field); the reframing facts are
% triangularReframe_comp_baseChange_borel (reframe_calculus) and frameBorel_groupoid_properties (the inverse t^-1);
% (*) is CellChart.reframe_pt, with reframe_cancel and isUnit_det_cellM; the stated property is CellChart.catches,
% which the Lean proves over any c.F-algebra; property (1) is CellChart.mono_chart (algHom_eq_of_point, cellM_unique)
% and property (2) is CellChart.open_chart (exists_away_fac, minor_iff, IsLocalization.Away.surj for g2). The
% matrices u are cellM sigma x = (lowM x tau).submatrix (sigma tau).symm id.

\begin{lemma}\label{cv-volume-key-atlas}
In the setting of Definition~\ref{cv-volume-twisting-section}, with generic framed points $y_Z$ chosen as in
Definition~\ref{cv-volume-key-chart}, the following hold.
\begin{enumerate}
\item Every framed point $y$ of $X$ over a field $L\supset F_c$ is $z\cdot t$ for an image $z$ of $y_C$ under a
  homomorphism $R_C\to L$, for some chart $C$ and some $t\in(L^\times)^I$.
\item Let $C$ be a chart and $Z$ a connected component of $X$. If an image $z$ of $y_C$ under a homomorphism
  $\psi\colon R_C\to\Omega_Z$ satisfies $y_Z=z\cdot t$ for some $t$, then $Z_C=Z$.
\end{enumerate}
\end{lemma}
\leanref{Development/AbelianVarietyTheory/Proofs/GenericRegularity/FlagChartsCoverFramedFieldPointsStatement.lean}{GlobalTheory.T.flagCharts_cover_framed_fieldPoints}
\leanref{Development/AbelianVarietyTheory/Proofs/GenericRegularity/ExistsSectionOfGenericValuesRegularOnCharts.lean}{GlobalTheory.Reg.SecGlue.chart_comp}
\begin{proof}
(1) Let $q\colon\operatorname{Spec}L\to X$ be the point of $y$. By Lemma~\ref{cv-volume-key-local-frames} its image
has an affine open neighbourhood $\operatorname{Spec}A$ with a framed point $y_A$ as there. As $\operatorname{Spec}L$
is a point, $q$ factors through $\operatorname{Spec}A$, by a homomorphism $f\colon A\to L$ of $F_c$-algebras. Let
$z_0$ be an image of $y_A$ under $f$ (Lemma~\ref{cv-volume-key-base-change}); its point is $q$. By
Lemma~\ref{cv-volume-key-reframe}(2) the vectors of $y$ are obtained from those of $z_0$ by a family $g$ of
invertible matrices, and by part (3) of that lemma $g_\tau=u_\tau b_\tau$ with $u_\tau$, $b_\tau$ and a permutation
$\sigma_\tau$ as there. Lemma~\ref{cv-volume-key-cell-charts}, for $A$, $y_A$ and $\sigma$, gives a chart $C$, and its
property, for $f$, $z_0$ and $y$, gives $y=z\cdot t$ as stated.

(2) Let $q\colon\operatorname{Spec}R_C\to X$ be the point of $y_C$. Let $\eta$ be the point of
$\operatorname{Spec}R_C$ given by the zero ideal, which is the image of the point of $\operatorname{Spec}\Omega_{Z_C}$,
as $R_C\subset\Omega_{Z_C}$, and let $\eta'$ be the image of the point of $\operatorname{Spec}\Omega_Z$ under
$\psi$. As every prime ideal contains the zero ideal, $\eta'$ is a specialization of $\eta$. An image of $y_C$ over
$\Omega_{Z_C}$ is $y_{Z_C}\cdot t_C$, so $q(\eta)$ is the point of $y_{Z_C}$, the generic point of $Z_C$. Likewise
$q(\eta')$ is the point of $z$, which is that of $y_Z$, the generic point of $Z$. So the generic point of $Z$ is a
specialization of that of $Z_C$ and lies in $Z_C$. Hence $Z$, its closure, is contained in $Z_C$, and as $Z$ is an
irreducible component of $X$, $Z=Z_C$.
\end{proof}
% Lean: (1) is GlobalTheory.Reg.atlas: exists_integral_smooth_chart_frame_moduli at the point of y,
% exists_algHom_of_mem_range for f, the base change hBC (FrBC, from exists_frMap in evLaws_genFr_atlas) for z0,
% isTriangularReframe_exists_and_unique, exists_perm_lower_upper and exists_flagCellChart_over_framedAffineOpen
% with x the entries of the u_tau below the pivots. (2) is chart_comp (liftF, liftF_gen_apply, closure_gen and the
% maximality of irreducible components). The connected components of the text are the Lean's irreducible
% components of Xo.XF c, which agree as X is smooth.

\begin{lemma}\label{cv-volume-key-chart-compat}
In the setting of Definition~\ref{cv-volume-twisting-section}, with generic framed points $y_Z$ and charts as in
Definition~\ref{cv-volume-key-chart}, put $s^e=\prod_is_i^{e_i}$ for $s\in(S^\times)^I$, $S$ a ring, and
$e\in\mathbf Z^I$. Let $C$ and $C'$ be charts, let $L\supset F_c$ be a field, let $\psi\colon R_C\to L$ and
$\psi'\colon R_{C'}\to L$ be homomorphisms of $F_c$-algebras, let $z$ be an image of $y_C$ under $\psi$ and $z'$ an
image of $y_{C'}$ under $\psi'$ (Lemma~\ref{cv-volume-key-base-change}), and let $t\in(L^\times)^I$ with $z'=z\cdot t$.
Then $Z_C=Z_{C'}$, and
\[
 \psi'(t_{C'}^eb)=t^e\,\psi(t_C^eb)
\]
for every $e\in\mathbf Z^I$ and every $b\in\Omega_{Z_C}$ with $t_C^eb\in R_C$ and $t_{C'}^eb\in R_{C'}$.
\end{lemma}
\leanref{Development/AbelianVarietyTheory/Proofs/IwasawaDecomposition/ExistsIntegerExponentSectionOfChartRegular.lean}{GlobalTheory.SecGlueInt.chart_compat}
\begin{proof}
Four rules are used throughout. The relation $y'=y\cdot t$ is reflexive, composes and is symmetric, and it
determines $t$ (proof of Lemma~\ref{cv-volume-key-root-values}). An image of a framed point $y$ under a homomorphism
$\varphi$ is unique, as its point, its vectors and its section are those of $y$ carried along $\varphi$. Images
compose: an image under $\mu$ of an image of $y$ under $\lambda$ is an image of $y$ under $\mu\lambda$, and an image
of $y$ under $\mu\lambda$ is an image under $\mu$ of the image of $y$ under $\lambda$. If $y'=y\cdot t$, the images of
$y$ and $y'$ under $\varphi$ satisfy the same relation with $\varphi(t)$.

The component. As $z=z'\cdot t^{-1}$, property~(2) of the chart $C'$ (Definition~\ref{cv-volume-key-chart}), for
$w=y_C$ and $\chi=\psi$, gives $g\in R_C$ with $\psi(g)\ne0$ such that, for every homomorphism
$\lambda\colon R_C\to S'$ of $F_c$-algebras with $\lambda(g)$ a unit, every image of $y_C$ under $\lambda$ is
$z''\cdot t''$ for an image $z''$ of $y_{C'}$ under some homomorphism $R_{C'}\to S'$ and some $t''$. Then $g\ne0$, so
$g$ is a unit of the field $\Omega_{Z_C}$. For the inclusion $R_C\subset\Omega_{Z_C}$ this gives
$y^\Omega_C=z_0\cdot t_0$, where $y^\Omega_C$ is the image of $y_C$ over $\Omega_{Z_C}$ and $z_0$ is an image of
$y_{C'}$ under some $\psi_0\colon R_{C'}\to\Omega_{Z_C}$. As $y^\Omega_C=y_{Z_C}\cdot t_C$, the point of $z_0$ is the
generic point $\eta$ of $Z_C$, and it is the image under $\psi_0$ of a point of $\operatorname{Spec}R_{C'}$. Every
point of $\operatorname{Spec}R_{C'}$ maps into $Z_{C'}$: it is a specialization of the zero ideal, which maps to the
point of the image of $y_{C'}$ over $\Omega_{Z_{C'}}$, the generic point of $Z_{C'}$, and $Z_{C'}$ is closed. So
$\eta\in Z_{C'}$, hence $Z_C\subset Z_{C'}$, and $Z_C=Z_{C'}$ as both are connected components of $X$.

The values. Put $Z=Z_C$, let $S=(R_C)_g$ and let $\lambda\colon R_C\to S$ be the localization map. As $g\ne0$ in
$\Omega_Z$ and $\psi(g)\ne0$, there are homomorphisms $\mu\colon S\to\Omega_Z$ and $\nu\colon S\to L$ with $\mu\lambda$
the inclusion of $R_C$ and $\nu\lambda=\psi$; and $\mu$ is injective, since every element of $S$ is
$\lambda(a)/\lambda(g)^n$ and $R_C\to\Omega_Z$ is injective. Let $w$ be an image of $y_C$ under $\lambda$. By the choice
of $g$, $w=v\cdot s$ for an image $v$ of $y_{C'}$ under some $\varphi\colon R_{C'}\to S$ and some $s\in(S^\times)^I$.

Apply $\mu$. The framed point $y^\Omega_C$ is an image of $w$ under $\mu$, so $y^\Omega_C=v_\mu\cdot\mu(s)$ for the
image $v_\mu$ of $v$ under $\mu$, which is an image of $y_{C'}$ under $\mu\varphi$. With the image
$y^\Omega_{C'}=y_Z\cdot t_{C'}$ of $y_{C'}$ over $\Omega_Z$ this gives $y^\Omega_{C'}=v_\mu\cdot\mu(s)t_C^{-1}t_{C'}$.
Property~(1) of the chart $C'$ shows that $\mu\varphi$ is the inclusion of $R_{C'}$, so $v_\mu=y^\Omega_{C'}$. Then
$y_Z\cdot t_C=y^\Omega_C=y_Z\cdot t_{C'}\mu(s)$, and $t_C=t_{C'}\mu(s)$.

Apply $\nu$. The point $z$ is an image of $w$ under $\nu$, so $z=v_\nu\cdot\nu(s)$ for the image $v_\nu$ of $v$ under
$\nu$, which is an image of $y_{C'}$ under $\nu\varphi$; and $z'=z\cdot t=v_\nu\cdot\nu(s)t$. Property~(1) of $C'$
shows that $\nu\varphi=\psi'$, so $v_\nu=z'$, and $z'=z'\cdot1$ gives $\nu(s)t=1$.

Now $\mu(\lambda(t_C^eb))=t_C^eb=\mu(s)^et_{C'}^eb=\mu\big(s^e\varphi(t_{C'}^eb)\big)$, so
$\lambda(t_C^eb)=s^e\varphi(t_{C'}^eb)$, as $\mu$ is injective. Applying $\nu$ gives
$\psi(t_C^eb)=\nu(s)^e\psi'(t_{C'}^eb)=t^{-e}\psi'(t_{C'}^eb)$.
\end{proof}
% Lean: chart_compat: C'.opn for g, with the inverse rule of frBor_groupoid; comp_eq for the component (chart_mem,
% liftF_gen_apply, closure_genPt' and the maximality of irreducible components); away_lift for mu and nu, away_inj;
% then the base change hev.bc (Lemma cv-volume-key-base-change), frMap_cancel and frMap_trans for composing images,
% frBorBC, the groupoid rules frBor_groupoid, the property mono' of C' (e1, e4), the uniqueness of images
% frMap_unique_refl_iff (e2, e5), and the identities e3, e6, e7, a3 and a4.

\begin{lemma}\label{cv-volume-key-local-values}
In the setting of Definition~\ref{cv-volume-twisting-section}, with generic framed points $y_Z$ and charts as in
Definition~\ref{cv-volume-key-chart}, and with $s^e$ as in Lemma~\ref{cv-volume-key-chart-compat}, let
$e\in\mathbf Z^I$, and let $b_Z\in\Omega_Z$ for every connected component $Z$ of $X$, such that
$t_C^eb_{Z_C}\in R_C$ for every chart $C$. Let $y$ be a framed point over an $F_c$-algebra $S$. An element $v\in S$ is
a \emph{local value} of $y$ if there are a chart $C$, a homomorphism $\psi\colon R_C\to S$ of $F_c$-algebras, an image
$z$ of $y_C$ under $\psi$ and $t\in(S^\times)^I$ with $y=z\cdot t$ and $v=t^e\psi(t_C^eb_{Z_C})$.
\begin{enumerate}
\item If $S$ is a field, $y$ has at most one local value.
\item If $v$ is a local value of $y$ and $y'$ is an image of $y$ under a homomorphism $\varphi\colon S\to T$ of
  $F_c$-algebras, then $\varphi(v)$ is a local value of $y'$.
\item If $v$ is a local value of $y$, $t\in(S^\times)^I$ and $y'=y\cdot t$, then $t^ev$ is a local value of $y'$.
\item If $S$ is a field, $y$ has a local value.
\item For every component $Z$, $b_Z$ is a local value of $y_Z$.
\end{enumerate}
\end{lemma}
\leanref{Development/AbelianVarietyTheory/Proofs/IwasawaDecomposition/ExistsIntegerExponentSectionOfChartRegular.lean}{GlobalTheory.SecGlueInt.Val}
\leanref{Development/AbelianVarietyTheory/Proofs/IwasawaDecomposition/ExistsIntegerExponentSectionOfChartRegular.lean}{GlobalTheory.SecGlueInt.val_unique}
\leanref{Development/AbelianVarietyTheory/Proofs/IwasawaDecomposition/ExistsIntegerExponentSectionOfChartRegular.lean}{GlobalTheory.SecGlueInt.val_map}
\leanref{Development/AbelianVarietyTheory/Proofs/IwasawaDecomposition/ExistsIntegerExponentSectionOfChartRegular.lean}{GlobalTheory.SecGlueInt.val_bor}
\leanref{Development/AbelianVarietyTheory/Proofs/IwasawaDecomposition/ExistsIntegerExponentSectionOfChartRegular.lean}{GlobalTheory.SecGlueInt.val_ex}
\leanref{Development/AbelianVarietyTheory/Proofs/IwasawaDecomposition/ExistsIntegerExponentSectionOfChartRegular.lean}{GlobalTheory.SecGlueInt.val_gen}
\begin{proof}
The rules at the start of the proof of Lemma~\ref{cv-volume-key-chart-compat} are used.

(1) Let $v$ and $v'$ be local values, given by $C,\psi,z,t$ and by $C',\psi',z',t'$. As $z\cdot t=y=z'\cdot t'$,
$z'=z\cdot tt'^{-1}$. By Lemma~\ref{cv-volume-key-chart-compat}, $Z_C=Z_{C'}$ and
$\psi'(t_{C'}^eb_{Z_C})=(tt'^{-1})^e\psi(t_C^eb_{Z_C})$, so $v'=t'^e\psi'(t_{C'}^eb_{Z_{C'}})=t^e\psi(t_C^eb_{Z_C})=v$.

(2) Let $v$ be given by $C,\psi,z,t$, and let $z_1$ be an image of $z$ under $\varphi$. It is an image of $y_C$ under
$\varphi\psi$, and $y'=z_1\cdot\varphi(t)$. So $\varphi(v)=\varphi(t)^e(\varphi\psi)(t_C^eb_{Z_C})$ is a local value of
$y'$.

(3) If $v$ is given by $C,\psi,z,t_1$, then $y'=z\cdot t_1t$ and $t^ev=(t_1t)^e\psi(t_C^eb_{Z_C})$.

(4) By Lemma~\ref{cv-volume-key-atlas}(1), $y=z\cdot t$ for an image $z$ of $y_C$ under some homomorphism
$\psi\colon R_C\to S$, for some chart $C$ and some $t$; then $t^e\psi(t_C^eb_{Z_C})$ is a local value of $y$.

(5) By Lemma~\ref{cv-volume-key-atlas}(1) for $y_Z$ and part~(2) of that lemma, there is a chart $C$ with $Z_C=Z$.
The image $y^\Omega_C$ of $y_C$ under the inclusion $R_C\subset\Omega_Z$ is $y_Z\cdot t_C$, so
$y_Z=y^\Omega_C\cdot t_C^{-1}$, and the local value of $y_Z$ given by $C$, the inclusion, $y^\Omega_C$ and $t_C^{-1}$ is
$t_C^{-e}t_C^eb_Z=b_Z$.
\end{proof}
% Lean: Val is the text's local value (with hb the hypothesis on the charts); val_unique (chart_compat), val_map
% (hl.bc, frMap_trans, frBorBC, zchar_map), val_bor (the product rule of frBor_groupoid, zchar_mul), val_ex (the
% atlas, Lemma cv-volume-key-atlas(1)) and val_gen (the atlas with chart_comp, Lemma cv-volume-key-atlas(2), then
% C.gen and the inverse of C.bor).

\begin{lemma}\label{cv-volume-key-representatives}
In the setting of Definition~\ref{cv-volume-twisting-section}, with generic framed points $y_Z$ and charts as in
Definition~\ref{cv-volume-key-chart}, let $e\in\mathbf Z^I$, and let $b_Z\in\Omega_Z$ for every connected component
$Z$ of $X$, such that $t_C^eb_{Z_C}\in R_C$ for every chart $C$; local values are those of
Lemma~\ref{cv-volume-key-local-values} for $e$ and $(b_Z)$. Let $y$ be a framed point over an $F_c$-algebra $S$. Then
there is $s\in S$ with $\chi(s)=v$ whenever $\chi\colon S\to L$ is a homomorphism of $F_c$-algebras to a field, $w$ is
an image of $y$ under $\chi$ and $v$ is a local value of $w$.
\end{lemma}
\leanref{Development/AbelianVarietyTheory/Proofs/IwasawaDecomposition/ExistsIntegerExponentSectionOfChartRegular.lean}{GlobalTheory.SecGlueInt.exists_rep}
\begin{proof}
Let $J$ be the set of $g\in S$ for which there is $s\in S$ with $\chi(g)v=\chi(s)$ for all $\chi$, $w$ and $v$ as in
the statement. It is an ideal of $S$: $0$ lies in it with $s=0$, $g+g'$ with $s+s'$, and $rg$ with $rs$. The claim
is that $1\in J$. Suppose not, and let $P\supset J$ be a maximal ideal, $k=S/P$ and $\chi_0\colon S\to k$ the quotient
map. Let $w_0$ be an image of $y$ under $\chi_0$ (Lemma~\ref{cv-volume-key-base-change}). By
Lemma~\ref{cv-volume-key-atlas}(1), $w_0=z_0\cdot t_0$ for an image $z_0$ of $y_C$ under a homomorphism $R_C\to k$,
for some chart $C$ and some $t_0$. Property~(2) of $C$ (Definition~\ref{cv-volume-key-chart}) gives $g\in S$ with
$\chi_0(g)\ne0$ such that, for every homomorphism $\lambda\colon S\to S'$ of $F_c$-algebras with $\lambda(g)$ a unit,
every image of $y$ under $\lambda$ is $z\cdot t$ for an image $z$ of $y_C$ under some homomorphism $R_C\to S'$ and some
$t$.

Let $\lambda\colon S\to S_g$ be the localization map and $w'$ an image of $y$ under $\lambda$. Then $w'=z\cdot t$ for an
image $z$ of $y_C$ under some $\psi\colon R_C\to S_g$, so $u=t^e\psi(t_C^eb_{Z_C})$ is a local value of $w'$. Write
$u\lambda(g)^n=\lambda(a)$ with $a\in S$ and $n\ge0$. Then $g^{n+1}\in J$, with $s=ga$. Indeed, let $\chi$, $w$ and $v$
be as in the statement. If $\chi(g)=0$, then $\chi(g^{n+1})v=0=\chi(ga)$. Otherwise $\chi=\mu\lambda$ for a homomorphism
$\mu\colon S_g\to L$, and $w$ is an image of $w'$ under $\mu$. By Lemma~\ref{cv-volume-key-local-values}(2), $\mu(u)$
is a local value of $w$, so $v=\mu(u)$ by part~(1) of that lemma, and
$\chi(g)^{n+1}v=\chi(g)\mu(u\lambda(g)^n)=\chi(g)\chi(a)=\chi(ga)$.

So $g^{n+1}\in J\subset P$, and $g\in P$ as $P$ is prime. This contradicts $\chi_0(g)\ne0$.
\end{proof}
% Lean: exists_rep: the ideal I of the g with some s and Rep hb y g s, Ideal.exists_le_maximal, chi0 to
% FractionRing (S/P), the atlas hA, C.opn, IsLocalization.Away.surj for a, then frMap_cancel, val_map and
% val_unique; hgP is the primality of P.

\begin{lemma}\label{cv-volume-key-gluing}
In the setting of Definition~\ref{cv-volume-twisting-section}, with generic framed points $y_Z$ and charts as in
Definition~\ref{cv-volume-key-chart}, and with $s^e$ as in Lemma~\ref{cv-volume-key-chart-compat}, let
$e\in\mathbf Z^I$, and let $b_Z\in\Omega_Z$ for every connected component $Z$ of $X$, such that
$t_C^eb_{Z_C}\in R_C$ for every chart $C$. Then there is a rational section $f$ of exponent $e$ with $f(y_Z)=b_Z$
for every $Z$. Here a rational section is a rule $f$ that assigns to every framed point $y$ over a \emph{test ring},
an $F_c$-algebra $F_c[x_1,\dots,x_n]/\mathfrak a$ with $\mathfrak a$ a radical ideal, an element $f(y)$ of that ring,
such that $f(y')=\varphi(f(y))$ whenever $y'$ is an image of $y$ under a homomorphism $\varphi$ of $F_c$-algebras
between test rings; it has exponent $e$ if $f(y\cdot t)=t^ef(y)$ for every framed point $y$ over a test ring $R$ and
every $t\in(R^\times)^I$; and its value at a framed point $w$ over an $F_c$-algebra $S$ that is the image of a framed
point $y$ over a test ring under a homomorphism $\varphi$ is $f(w)=\varphi(f(y))$, for a choice of $y$ and $\varphi$.
\end{lemma}
\leanref{Development/AbelianVarietyTheory/Proofs/IwasawaDecomposition/ExistsIntegerExponentSectionOfChartRegularStatement.lean}{GlobalTheory.T.exists_integer_exponent_section_of_chart_regular}
\leanref{Development/AbelianVarietyTheory/Proofs/IwasawaDecomposition/ExistsIntegerExponentSectionOfChartRegular.lean}{GlobalTheory.SecGlueInt.exists_sec_of_regular}
\begin{proof}
Let local values be those of Lemma~\ref{cv-volume-key-local-values}. For every test ring $R$ and every framed point
$y$ over $R$, Lemma~\ref{cv-volume-key-representatives} gives $f(y)\in R$ with $\chi(f(y))=v$ whenever
$\chi\colon R\to L$ is a homomorphism of $F_c$-algebras to a field, $w$ is an image of $y$ under $\chi$ and $v$ is a
local value of $w$. Test rings are reduced, and an element $a$ of a reduced $F_c$-algebra $S$ with $\chi(a)=0$ for
every homomorphism $\chi$ of $F_c$-algebras from $S$ to a field is zero: otherwise $a$ is not nilpotent, so it lies
outside some prime ideal $\mathfrak p$, and its image in the fraction field of $S/\mathfrak p$ is not zero. So the
identities below may be checked after every such $\chi$; every framed point over a field has a local value
(Lemma~\ref{cv-volume-key-local-values}(4)).

Compatibility. Let $y'$ be an image of $y$ under $\varphi\colon R\to R'$, let $\chi\colon R'\to L$, let $w$ be an image
of $y'$ under $\chi$ and $v$ a local value of $w$. Then $w$ is an image of $y$ under $\chi\varphi$, so
$\chi(f(y'))=v=\chi(\varphi(f(y)))$.

Exponent. Let $y'=y\cdot t$ over a test ring $R$, let $\chi\colon R\to L$, and let $w$ and $w'$ be images of $y$ and
$y'$ under $\chi$, so that $w'=w\cdot\chi(t)$. If $v$ is a local value of $w$, then $\chi(t)^ev$ is a local value of
$w'$ (Lemma~\ref{cv-volume-key-local-values}(3)), so $\chi(f(y'))=\chi(t)^ev=\chi(t^ef(y))$.

Values. The field $\Omega_Z$ is reduced, so $y_Z$ is tested (proof of Lemma~\ref{cv-volume-twisting-points}): it is
the image of a framed point $y'$ over a test ring under a homomorphism $\varphi$, and $f(y_Z)=\varphi(f(y'))$ for such
$y'$ and $\varphi$. As $b_Z$ is a local value of $y_Z$ (Lemma~\ref{cv-volume-key-local-values}(5)),
$\varphi(f(y'))=b_Z$ for every such choice.
\end{proof}
% Lean: exists_sec_of_regular, through exists_secZ_of_regular, which takes the laws EvLaws and the atlas from
% evLaws_genFr_atlas: the choice g of exists_rep on every TRing (hg'), eq_of_fields with the reduced test rings
% (Ideal.isRadical_iff_quotient_reduced), val_ex, frMap_trans and val_bor for membership in secZ (mem_secZ), and
% hl.test (AllTest, framedPoint_reduced_from_testRing) with val_gen for evC. A rational section of the text is an
% element of Rul (Rules on TRing), naturality is IsNatural, the exponent law is that of secZ, and the value at a
% tested framed point is evR = evalAt (FromTest, with the chosen presentation). Of the laws EvLaws only bc (Lemma
% cv-volume-key-base-change) and test are used. The earlier count of 45,000 lines for (K2) included the laws wd and
% wt (framedModuli_algebraicLaws_chart_weight_integrality, about 11,800 lines), which the gluing does not use, and
% the atlas of (K1).

\subsection{Uniformization families and normalized framed points}\label{subsection-cv-normalized-framed-points}

\begin{lemma}\label{cv-volume-key-xi}
Let $E$ be a CM field with complex conjugation $b\mapsto\bar b$, let $d$ be a prime, let $\sqrt{-d}\in E$ be an
element with square $-d$, and let $\theta\in\mathbf C$ with $\theta^2=-d$.
\begin{enumerate}
\item There is $\xi_0\in E$ with $\bar\xi_0=-\xi_0$, $\xi_0\ne0$ and $\operatorname{Im}\tau(\xi_0)<0$ for every
  embedding $\tau\colon E\to\mathbf C$ with $\tau(\sqrt{-d})=\theta$.
\item Let $L$ be a finitely generated projective $\mathcal O_E$-module of rank $r$ with a perfect alternating form
  $\langle\ ,\ \rangle\colon L\times L\to\mathbf Z$ such that $\langle bx,y\rangle=\langle x,\bar by\rangle$ for
  $b\in\mathcal O_E$, and let $\xi\in E$ with $\bar\xi=-\xi\ne0$. For $G\in M_r(E)$ and $u,v\in E^r$ put
  $h_G(u,v)=\sum_{a,b}u_aG_{ab}\bar v_b$. Then there are an invertible matrix $G\in M_r(E)$ with $\bar G_{ab}=G_{ba}$
  for all $a,b$ and an injective $\mathcal O_E$-linear map $\varphi\colon L\to E^r$ whose image spans $E^r$ over $E$,
  such that $\langle x,y\rangle=\operatorname{Tr}_{E/\mathbf Q}\big(\xi\,h_G(\varphi x,\varphi y)\big)$ for all
  $x,y\in L$.
\end{enumerate}
\end{lemma}
\leanref{Development/AbelianVarietyTheory/Proofs/UnitaryExistence/ExistsIsUnf.lean}{GlobalTheory.IsUnfEx.exists_xi}
\leanref{Development/AbelianVarietyTheory/Proofs/UnitaryExistence/ExistsHermCoordOfDatum.lean}{GlobalTheory.exists_herm_coord_of_datum}
\begin{proof}
(1) For every embedding $\tau$, $\tau(\sqrt{-d})$ is a square root of $-d$, so it is purely imaginary and not $0$. As
$\tau(\bar b)=\overline{\tau(b)}$ for the CM field $E$, $\tau(\overline{\sqrt{-d}})=-\tau(\sqrt{-d})$, so
$\overline{\sqrt{-d}}=-\sqrt{-d}$. Likewise $\theta$ is purely imaginary and not $0$. Take $\xi_0=\sqrt{-d}$ if
$\operatorname{Im}\theta<0$ and $\xi_0=-\sqrt{-d}$ otherwise. For $\tau$ with $\tau(\sqrt{-d})=\theta$,
$\operatorname{Im}\tau(\xi_0)=\pm\operatorname{Im}\theta<0$.

(2) For $x,y\in L$ the map $b\mapsto\langle bx,y\rangle$ is additive from $\mathcal O_E$ to $\mathbf Z$, so, as the
trace form of $E$ is nondegenerate, there is a unique $h_0(x,y)\in E$ with
$\operatorname{Tr}_{E/\mathbf Q}(b\,h_0(x,y))=\langle bx,y\rangle$ for all $b\in\mathcal O_E$. By the uniqueness,
$h_0$ is additive in each variable, $h_0(bx,y)=b\,h_0(x,y)$ and $h_0(x,by)=\bar b\,h_0(x,y)$; and, as the form is
alternating and the trace is invariant under conjugation, $h_0(y,x)=-\overline{h_0(x,y)}$. If $h_0(x,y)=0$ for all
$y$, then $\langle x,y\rangle=\operatorname{Tr}_{E/\mathbf Q}h_0(x,y)=0$ for all $y$, and $x=0$, as the form is perfect.

Choose $x_1,\dots,x_r\in L$ linearly independent over $\mathcal O_E$. Every $x\in L$ has a multiple $tx$, with
$t\in\mathcal O_E\setminus\{0\}$, in their span. Put $\varphi(x)=(h_0(x,x_j))_j$ and let $A$ be the matrix
$(h_0(x_i,x_j))_{i,j}$. The map $\varphi$ is injective: if $\varphi(x)=0$, then for every $y$, with $sy$ in the span,
$\bar s\,h_0(x,y)=h_0(x,sy)=0$, so $x=0$. The matrix $A$ is invertible: a row vector $c$ with $cA=0$, scaled to have
entries in $\mathcal O_E$, gives $\varphi(\sum_ic_ix_i)=cA=0$, so $\sum_ic_ix_i=0$ and $c=0$. The image of $\varphi$
contains the rows of $A$, so it spans $E^r$. Put $G=-\xi^{-1}A^{-1}$. As $\bar A_{ij}=-A_{ji}$ and $\bar\xi=-\xi$,
$\bar G_{ab}=G_{ba}$, and $G$ is invertible. Let $x,y\in L$, with $tx=\sum_ic_ix_i$ and $sy=\sum_jd_jx_j$. Then
$t\varphi(x)=cA$ and $\overline{s\varphi(y)}=-A\bar d$, as a row and a column, so
\[
 t\bar s\,\xi\,h_G(\varphi x,\varphi y)=\xi\,(cA)(-\xi^{-1}A^{-1})(-A\bar d)=cA\bar d=h_0(tx,sy)=t\bar s\,h_0(x,y),
\]
and $\operatorname{Tr}_{E/\mathbf Q}\big(\xi\,h_G(\varphi x,\varphi y)\big)=\operatorname{Tr}_{E/\mathbf Q}h_0(x,y)
=\langle x,y\rangle$.
\end{proof}
% Lean: (1) is exists_xi (rt or -rt by the sign of the imaginary part of rt0, with cc rt = -rt). (2) is
% exists_herm_coord_of_datum: h0 from exists_trace_mul_eq_of_addMonoidHom with uniqueness by
% DatumCoord.eq_of_trace_int, the family l from exists_linearIndependent_of_le_finrank with hmax, psi and gram (hdet),
% and Gm = -xi^-1 gram^-1. The Lean states (2) for the lattice D.La of any datum D, of rank r = D.p + D.q.

\begin{lemma}\label{cv-volume-key-models-half}
In the setting of Definition~\ref{cv-volume-twisting-section}, with the level $(M,\mathfrak M,N)$ of $X$, let $L'$ be the
lattice of a PEL datum over $E$, with its form $\langle\ ,\ \rangle'$, and let $\gamma$ be an $\mathcal O_E$-linear
automorphism of $L'$ that preserves the form, such that $\gamma(a)\in a+ML'$ for every $a\in\bar{\mathfrak M}L'$. Then
$\gamma(a)\in a+ML'$ for every $a\in L'$.
\end{lemma}
\leanref{Development/AbelianVarietyTheory/Proofs/UnitaryExistence/ExistsEqNsmulOfDvdPairing.lean}{GlobalTheory.LevelHalf.trivial_mod_of_trivial_on_half}
\begin{proof}
As $\mathfrak M+\bar{\mathfrak M}=\mathcal O_E$, write $1=i_0+j_0$ with $i_0\in\mathfrak M$ and $j_0\in\bar{\mathfrak M}$.
Let $a\in L'$ and $y=\gamma(a)-a$. As $j_0a\in\bar{\mathfrak M}L'$, $j_0y=\gamma(j_0a)-j_0a\in ML'$. For $z\in L'$,
$\bar i_0z\in\bar{\mathfrak M}L'$, so $\gamma(\bar i_0z)=\bar i_0z+Me$ for some $e\in L'$, and as $\gamma$ preserves the
form,
\[
 \langle i_0y,z\rangle'=\langle y,\bar i_0z\rangle'=\langle\gamma(a),\bar i_0z\rangle'-\langle\gamma(a),\gamma(\bar i_0z)\rangle'
 =-M\langle\gamma(a),e\rangle'.
\]
So $z\mapsto\langle i_0y,z\rangle'/M$ is an integer-valued linear form on $L'$. As the form is perfect, it is
$\langle w,\cdot\rangle'$ for some $w\in L'$, and $i_0y=Mw$. Hence $y=i_0y+j_0y\in ML'$.
\end{proof}
% Lean: trivial_mod_of_trivial_on_half, for gam : D.La equiv D.La over O_E: i and j from the field of the level
% (Mb + bar Mb = top), hjy, hdiv (with the field hbar of the datum, <b x, y> = <x, bar b y>), and the functional psi
% with the perfectness hbij of the form.

\begin{lemma}\label{cv-volume-key-models-volume}
Let $R$ be a Dedekind domain with fraction field $K$, and let $r\ge0$. For a set $S\subset R^r$ of row vectors let
$\operatorname{vol}(S)$ be the ideal of $R$ generated by the determinants of the $r\times r$ matrices whose rows lie in
$S$. Let $\mathfrak c$ be a nonzero fractional ideal of $R$, let $A\in M_r(K)$ with $\det A\ne0$, and let
$S,S'\subset R^r$ with $\operatorname{vol}(S')\ne0$, such that a vector $v\in K^r$ lies in $S$ if and only if
$vAy^{\mathsf T}\in\mathfrak c$ for every $y\in S'$. Then
$(\det A)\operatorname{vol}(S)\operatorname{vol}(S')=\mathfrak c^r$.
\end{lemma}
\leanref{Development/AbelianVarietyTheory/Proofs/UnitaryExistence/DetMulVolMulConj.lean}{GlobalTheory.DetIdeal.det_mul_vol_eq}
\begin{proof}
The inclusion. Let $X$ and $Y$ be $r\times r$ matrices with rows in $S$ and in $S'$. The entries of $XAY^{\mathsf T}$ are
the numbers $xAy^{\mathsf T}$ for the rows $x$ of $X$ and $y$ of $Y$, which lie in $\mathfrak c$. Every term of the
expansion of a determinant is a product of $r$ entries up to sign, so
$\det X\det A\det Y=\det(XAY^{\mathsf T})\in\mathfrak c^r$. Hence
$J=(\det A)\operatorname{vol}(S)\operatorname{vol}(S')\mathfrak c^{-r}$ is an ideal of $R$. Suppose that $J\ne R$, and
let $P$ be a maximal ideal of $R$ that contains $J$.

Local generators. Let $V$ be a nonzero fractional ideal generated by a set $\Sigma$. The products $au$ with
$a\in\Sigma$ and $u\in V^{-1}$ lie in $R$ and generate $VV^{-1}=R$, so one of them lies outside $P$. For
$V=\mathfrak c$ this gives $t\in\mathfrak c$ and $w\in\mathfrak c^{-1}$ with $s'=tw\notin P$. For
$V=\operatorname{vol}(S')$, generated by the determinants of the matrices with rows in $S'$, it gives such a matrix $Y$
and $u\in\operatorname{vol}(S')^{-1}$ with $s=u\det Y\notin P$.

A matrix with rows in $S$. Put $X=ut\operatorname{adj}(Y^{\mathsf T})A^{-1}$, so that
$XA=ut\operatorname{adj}(Y^{\mathsf T})$. Let $y\in S'$ and $1\le i\le r$, and let $Y_i(y)$ be $Y$ with its $i$-th row
replaced by $y$. By Cramer's rule the $i$-th entry of $\operatorname{adj}(Y^{\mathsf T})y^{\mathsf T}$ is
$\det Y_i(y)$, which lies in $\operatorname{vol}(S')$. So the $i$-th row $X_i$ of $X$ satisfies
$X_iAy^{\mathsf T}=t\cdot u\det Y_i(y)\in tR\subset\mathfrak c$. By the hypothesis, every row of $X$ lies in $S$.

The contradiction. As $\operatorname{adj}(Y^{\mathsf T})Y^{\mathsf T}=(\det Y)\cdot1$, $XAY^{\mathsf T}=ut\det Y\cdot1$,
so $(ut\det Y)^r=\det X\det A\det Y$ lies in $(\det A)\operatorname{vol}(S)\operatorname{vol}(S')$. Multiplying by
$w^r\in\mathfrak c^{-r}$ gives $(ss')^r=(u\det Y)^r(tw)^r\in J\subset P$. As $P$ is prime, $s\in P$ or $s'\in P$, a
contradiction. So $J=R$.
\end{proof}
% Lean: det_mul_vol_eq: det_mul_vol_le (with det_mem_pow) for the inclusion, FractionalIdeal.le_one_iff_exists_coeIdeal
% for J, Ideal.exists_le_maximal for P, DetIdeal.exists_local_generator twice (for c, and for vol L' spanned by the
% determinants), Xk = (u t) adj(Yk^T) A^-1 with Matrix.cramer_transpose_apply, and det(Xk A Yk^T) = (u t det Y)^r.

\begin{lemma}\label{cv-volume-key-models-dual}
Let $E$ be a CM field with complex conjugation $b\mapsto\bar b$, let $L'$ be a finitely generated projective
$\mathcal O_E$-module with a perfect alternating form $\langle\ ,\ \rangle'\colon L'\times L'\to\mathbf Z$ such that
$\langle bx,y\rangle'=\langle x,\bar by\rangle'$ for $b\in\mathcal O_E$, and let $(G,\eta,\varphi)$ be a Hermitian model of
$L'$ (Definition~\ref{cv-uniformization-model}), with $G\in M_r(E)$. Then a vector $v\in E^r$ lies in $\varphi(L')$ if and
only if $\operatorname{Tr}_{E/\mathbf Q}\big(\eta\,h_G(v,u)\big)\in\mathbf Z$ for every $u\in\varphi(L')$.
\end{lemma}
\leanref{Development/AbelianVarietyTheory/Proofs/UnitaryExistence/FinrankEqOfInjectiveOfSpan.lean}{GlobalTheory.HermModel.rank_integral_selfDual}
\begin{proof}
If $v=\varphi(x)$, then $\operatorname{Tr}_{E/\mathbf Q}(\eta\,h_G(v,\varphi y))=\langle x,y\rangle'\in\mathbf Z$ for
$y\in L'$. Conversely, let $v$ satisfy the condition, and put
$n(y)=\operatorname{Tr}_{E/\mathbf Q}(\eta\,h_G(v,\varphi y))$ for $y\in L'$. The map $n\colon L'\to\mathbf Z$ is additive,
so, as the form is perfect, there is $x\in L'$ with $n(y)=\langle x,y\rangle'$ for all $y$. Put $v'=v-\varphi(x)$; then
$\operatorname{Tr}_{E/\mathbf Q}(\eta\,h_G(v',\varphi y))=0$ for all $y\in L'$. For $b\in\mathcal O_E$,
$h_G(v',\varphi(\bar by))=b\,h_G(v',\varphi y)$, so $\operatorname{Tr}_{E/\mathbf Q}(b\,\eta\,h_G(v',\varphi y))=0$ for
every $b\in\mathcal O_E$. As the trace form of $E$ is nondegenerate and $\eta\ne0$, $h_G(v',\varphi y)=0$ for all $y$. As
$\varphi(L')$ spans $E^r$, $h_G(v',\cdot)$ vanishes on $E^r$, and $v'=0$, as $G$ is invertible. So
$v=\varphi(x)\in\varphi(L')$.
\end{proof}
% Lean: the last part of HermModel.rank_integral_selfDual: D.unimod.2 for x, eq_zero_of_trace_mul_integer with
% herm_smul_right for the multiples by H.bar.symm a, the span induction, and eq_zero_of_herm_eq_zero (hdet).

\begin{lemma}\label{cv-volume-key-models-ideal}
Let $E$ be a CM field with complex conjugation $b\mapsto\bar b$, let $r\ge0$, and let $\eta\in E^\times$ with
$\bar\eta=-\eta$. Let $L'$ and $L''$ be finitely generated projective $\mathcal O_E$-modules with perfect alternating forms
as in Lemma~\ref{cv-volume-key-models-dual}, with Hermitian models $(G,\eta,\varphi)$ of $L'$ and $(G',\eta,\psi)$ of
$L''$ (Definition~\ref{cv-uniformization-model}), where $G,G'\in M_r(E)$. Then there are nonzero ideals $I$ and $J$ of
$\mathcal O_E$ with
\[
 \frac{\det G'}{\det G}\,I\bar I=J\bar J
\]
as fractional ideals, where $\bar I$ is the ideal of the conjugates of the elements of $I$.
\end{lemma}
\leanref{Development/AbelianVarietyTheory/Proofs/UnitaryExistence/DetMulVolMulConj.lean}{GlobalTheory.DetIdeal.det_ideal}
\leanref{Development/AbelianVarietyTheory/Proofs/UnitaryExistence/DetMulVolMulConj.lean}{GlobalTheory.DetIdeal.det_mul_vol_mul_conj}
\begin{proof}
Let $\mathfrak d^{-1}$ be the inverse different of $E$, the fractional ideal of the $b\in E$ with
$\operatorname{Tr}_{E/\mathbf Q}(b\mathcal O_E)\subset\mathbf Z$, and let $\operatorname{vol}$ be as in
Lemma~\ref{cv-volume-key-models-volume} for $R=\mathcal O_E$. The modules $\varphi(L')$ and $\psi(L'')$ are finitely
generated, so there is $n\in\mathcal O_E\setminus\{0\}$ with $n\varphi(L')\subset\mathcal O_E^r$ and
$n\psi(L'')\subset\mathcal O_E^r$.

One model. Let $S=n\varphi(L')$, and let $\bar S$ be the set of the vectors of $S$ with conjugated entries. We show that
$\operatorname{vol}(S)\ne0$ and
\[
 \Big(\frac{\eta}{n\bar n}\Big)^r\det G\cdot\operatorname{vol}(S)\,\overline{\operatorname{vol}(S)}=\mathfrak d^{-r}.
\]
As $\varphi(L')$ spans $E^r$, it contains the rows of an invertible matrix $Y_0$; the matrix $nY_0$ has rows in $S$ and
determinant $n^r\det Y_0\ne0$, so $\operatorname{vol}(S)\ne0$. Conjugating the entries of the matrices with rows in $S$
gives the matrices with rows in $\bar S$, so $\operatorname{vol}(\bar S)=\overline{\operatorname{vol}(S)}\ne0$. Put
$A=\frac{\eta}{n\bar n}G$, with $\det A=(\eta/n\bar n)^r\det G\ne0$. For $v\in E^r$ and $y=n\varphi(z)\in S$,
\[
 vA\bar y^{\mathsf T}=\frac{\eta}{n\bar n}\,vG\,\bar n\,\overline{\varphi(z)}^{\mathsf T}=\eta\,h_G(n^{-1}v,\varphi z).
\]
As $b\,h_G(n^{-1}v,\varphi z)=h_G(n^{-1}v,\varphi(\bar bz))$ for $b\in\mathcal O_E$, the numbers $vA\bar y^{\mathsf T}$
with $y\in S$ all lie in $\mathfrak d^{-1}$ if and only if
$\operatorname{Tr}_{E/\mathbf Q}(\eta\,h_G(n^{-1}v,u))\in\mathbf Z$ for every $u\in\varphi(L')$, that is, by
Lemma~\ref{cv-volume-key-models-dual}, if and only if $n^{-1}v\in\varphi(L')$, or $v\in S$.
Lemma~\ref{cv-volume-key-models-volume}, for $\mathfrak c=\mathfrak d^{-1}$, $A$, $S$ and $S'=\bar S$, gives the identity.

The quotient. Let $J=\operatorname{vol}(n\varphi(L'))$ and $I=\operatorname{vol}(n\psi(L''))$. By the identity for the
two models,
\[
 \Big(\frac{\eta}{n\bar n}\Big)^r\det G\cdot J\bar J=\mathfrak d^{-r}=\Big(\frac{\eta}{n\bar n}\Big)^r\det G'\cdot I\bar I,
\]
and dividing by $(\eta/n\bar n)^r\det G\ne0$ gives $J\bar J=(\det G'/\det G)\,I\bar I$.
\end{proof}
% Lean: det_ideal: exists_denominator for the sum of the ranges (finitely generated by rank_integral_selfDual), and
% det_mul_vol_mul_conj for each model: L = {n phi z} in O_E^r, hvol from exists_det_ne_zero, map_vol for the conjugate,
% det_mul_vol_eq with c = FractionalIdeal.dual Z Q 1 and A = (xi/(n bar n)) Gm, hpair for the pairing and hdual (Lemma
% cv-volume-key-models-dual). The case r = 0 and the derivation of xi != 0 from the perfect form are separate in the Lean.

\begin{lemma}\label{cv-volume-key-models-generator}
Let $B$ be a Dedekind domain, and let $I$ and $Q$ be nonzero ideals of $B$. Then there are $i\in I$ and $r\in B$ with
$r-1\in Q$ and $rI\subset iB$.
\end{lemma}
\leanref{Development/AbelianVarietyTheory/Proofs/UnitaryExistence/ExistsBinaryEqAdicCompletion.lean}{GlobalTheory.LocalNorm.exists_local_generator}
\begin{proof}
The ideal $QI$ is nonzero and contained in $I$, so $I=QI+iB$ for some $i\in I$: by the Chinese remainder theorem there
is $i\in I$ outside $I\mathfrak p$ for each of the finitely many primes $\mathfrak p$ that contain $QI$, and then $QI+iB$
and $I$ have the same valuation at every prime. So the finitely generated module $I/iB$ satisfies $Q(I/iB)=I/iB$, and by
Nakayama's lemma \cite[Corollary~2.5]{AM69} there is $r\in B$ with $r-1\in Q$ and $r(I/iB)=0$, that is, $rI\subset iB$.
\end{proof}
% Lean: IsDedekindDomain.exists_sup_span_eq for QI + (i) = I, and
% Submodule.exists_sub_one_mem_and_smul_le_of_fg_of_le_sup for r.

\begin{lemma}\label{cv-volume-key-models-unit}
Let $E$ be a CM field with totally real subfield $E^+$ and complex conjugation $b\mapsto\bar b$, let $x\in E^+$ be
nonzero, and let $I$ and $J$ be nonzero ideals of $\mathcal O_E$ with $xI\bar I=J\bar J$ as fractional ideals. Let
$\mathfrak p$ be a nonzero prime ideal of $\mathcal O_{E^+}$. Then there are $z_1\in E^\times$, $z_2\in E$ and
$u\in\mathcal O_{E^+}\setminus\mathfrak p$ with $x\,z_1\bar z_1=u\,z_2\bar z_2$.
\end{lemma}
\leanref{Development/AbelianVarietyTheory/Proofs/UnitaryExistence/ExistsBinaryEqAdicCompletion.lean}{GlobalTheory.LocalNorm.exists_mul_norm_eq}
\begin{proof}
Let $P$ be a prime ideal of $\mathcal O_E$ with $P\cap\mathcal O_{E^+}=\mathfrak p$, and let $Q=\mathfrak p\mathcal O_E$, a
nonzero ideal with $Q\subset P$ and $\bar Q=Q$. If $r-1\in Q$, then $r\notin P$; and $\bar r-1\in\bar Q=Q$, so
$\bar r\notin P$ and $r\bar r\notin P$.

By Lemma~\ref{cv-volume-key-models-generator} there are $i\in I$ and $r\in\mathcal O_E$ with $r-1\in Q$ and
$rI\subset i\mathcal O_E$, and $j\in J$ and $s\in\mathcal O_E$ with $s-1\in Q$ and $sJ\subset j\mathcal O_E$. As $r\ne0$
and $I\ne0$, $i\ne0$; likewise $j\ne0$, and $s\ne0$. Conjugating gives $\bar r\bar I\subset\bar i\mathcal O_E$, so
$r\bar rI\bar I\subset i\bar i\mathcal O_E$, and likewise $s\bar sJ\bar J\subset j\bar j\mathcal O_E$. As
$i\bar i\in I\bar I$, $m=x\,i\bar i$ lies in $J\bar J$; as $j\bar j\in J\bar J$, $j\bar j=xm'$ with $m'\in I\bar I$. So
there are $c,c'\in\mathcal O_E$ with
\[
 c\,j\bar j=s\bar s\,m=s\bar s\,x\,i\bar i,\qquad c'\,i\bar i=r\bar r\,m'.
\]
Multiplying, $cc'\,i\bar i\,j\bar j=r\bar r\,s\bar s\,i\bar i\,(xm')=r\bar r\,s\bar s\,i\bar i\,j\bar j$, so
$cc'=r\bar r\,s\bar s\notin P$, and $c\notin P$. Conjugating the first identity and using $\bar x=x$ gives
$\bar c\,j\bar j=c\,j\bar j$, so $\bar c=c$ and $c=u\in\mathcal O_{E^+}$; and $u\notin\mathfrak p$, since otherwise
$c\in Q\subset P$. With $z_1=is$ and $z_2=j$, $x\,z_1\bar z_1=s\bar s\,x\,i\bar i=u\,j\bar j=u\,z_2\bar z_2$.
\end{proof}
% Lean: exists_mul_norm_eq: P from Ideal.exists_ideal_over_maximal_of_isIntegral, Q = v O_E, hone and hnorm, the local
% generators (exists_local_generator), hgen for i, j != 0, m and m' from h, exists_mul_eq_of_mem_mul_map for c and c',
% hcc, hcP, hfix and IsCMField.ringOfIntegersComplexConj_eq_self_iff for u.

\begin{lemma}\label{cv-volume-key-models-binary}
Let $F$ be a number field, let $\mathfrak p$ be a nonzero prime ideal of $\mathcal O_F$, with completion $F_{\mathfrak p}$,
and let $a,n,u\in\mathcal O_F$ with $a^2-4n\notin\mathfrak p$ and $u\notin\mathfrak p$. Then there are
$r,s\in F_{\mathfrak p}$ with $r^2+ars+ns^2=u$.
\end{lemma}
\leanref{Development/AbelianVarietyTheory/Proofs/UnitaryExistence/ExistsBinaryEqAdicCompletion.lean}{GlobalTheory.LocalNorm.exists_binary_eq_adicCompletion}
\leanref{Development/AbelianVarietyTheory/Proofs/AdelicLatticeProofs/BinaryFormUniversal.lean}{GlobalTheory.BinaryFormUniversal.binary_form_universal}
\leanref{Development/AbelianVarietyTheory/Proofs/UnitaryExistence/ExistsBinaryEqAdicCompletion.lean}{GlobalTheory.LocalNorm.exists_add_mul_sq_eq_of_norm}
\begin{proof}
The residue field. Let $k=\mathcal O_F/\mathfrak p$, a finite field, and let $a_0$, $n_0$ and $u_0$ be the images of
$a$, $n$ and $u$ in $k$, so that $D_0=a_0^2-4n_0\ne0$. There are $r_0,s_0\in k$ with
$r_0^2+a_0r_0s_0+n_0s_0^2=u_0$. If $k$ has characteristic $2$, every element of $k$ is a square, and $s_0=0$ will do.
Otherwise $|k|$ is odd, and the sets of the values of $x^2$ and of $D_0s^2+4u_0$, for $x,s\in k$, each have
$(|k|+1)/2$ elements, so they meet: $x^2=D_0s_0^2+4u_0$ for some $x$ and $s_0$. Then $r_0=(x-a_0s_0)/2$ satisfies
$r_0^2+a_0r_0s_0+n_0s_0^2=(x^2-D_0s_0^2)/4=u_0$. Lift $r_0$ and $s_0$ to $r,s\in\mathcal O_F$, and put
$N=r^2+ars+ns^2$. Then $u-N\in\mathfrak p$, so $N\notin\mathfrak p$.

A quadratic equation. Let $K$ be a complete field with an ultrametric absolute value, and let $c,e\in K$ with $|c|\le1$
and $|e|<1$. Then $t+ct^2=e$ for some $t\in K$: the map $t\mapsto e-ct^2$ sends the closed ball $B$ of radius $|e|$
about $0$ into itself, as $|ct^2|\le|e|^2\le|e|$, and it is a contraction of $B$, as
$|c(t^2-t'^2)|\le|t+t'|\,|t-t'|\le|e|\,|t-t'|$; its fixed point will do. In $F_{\mathfrak p}$ put $D=a^2-4n$,
$c=-n/D$ and $e=(u-N)/N$. As $D$ and $N$ are units at $\mathfrak p$, $|c|\le1$ and $|e|<1$. Let $t\in F_{\mathfrak p}$
with $t+ct^2=e$.

The solution. Put $\lambda=1-2nt/D$, $\mu=at/D$, $r'=r\lambda-ns\mu$ and $s'=r\mu+s\lambda+as\mu$. In the algebra
$F_{\mathfrak p}[X]/(X^2-aX+n)$, $(r+sX)(\lambda+\mu X)=r'+s'X$, and the norm of $\alpha+\beta X$ is
$\alpha^2+a\alpha\beta+n\beta^2$. As the norm is multiplicative,
\[
 r'^2+ar's'+ns'^2=N\,(\lambda^2+a\lambda\mu+n\mu^2)=N\Big(1+t-\frac{nt^2}D\Big)=N(1+e)=u.
\]
\end{proof}
% Lean: exists_binary_eq_adicCompletion: the residue field step is UnramifiedHermitian.exists_binary_eq_of_finite =
% BinaryFormUniversal.binary_form_universal (FiniteField.isSquare_of_char_two, and FiniteField.exists_root_sum_quadratic
% for X^2 and -D X^2 - 4u); the quadratic equation is exists_add_mul_sq_eq_of_norm (ContractingWith.exists_fixedPoint'
% on the closed ball); the solution is the final linear_combination, with y = D^-1.

\begin{lemma}\label{cv-volume-key-models-units}
Let $E$ be a CM field with totally real subfield $E^+$ and complex conjugation $b\mapsto\bar b$, let $d$ be a prime, and
let $\sqrt{-d}\in E$ be an element with square $-d$.
\begin{enumerate}
\item For $z\in E$, the elements $\alpha=(z+\bar z)/2$ and $\beta=(z-\bar z)/(2\sqrt{-d})$ lie in $E^+$, and
  $\alpha^2+d\beta^2=z\bar z$.
\item Suppose that no finite prime of $E^+$ ramifies in $E$. Let $\mathfrak p$ be a nonzero prime ideal of
  $\mathcal O_{E^+}$, with completion $E^+_{\mathfrak p}$, and let $u\in\mathcal O_{E^+}\setminus\mathfrak p$. Then
  $u=a^2+db^2$ for some $a,b\in E^+_{\mathfrak p}$.
\end{enumerate}
\end{lemma}
\leanref{Development/AbelianVarietyTheory/Proofs/UnitaryExistence/ExistsBinaryEqAdicCompletion.lean}{GlobalTheory.LocalNorm.exists_coords}
\leanref{Development/AbelianVarietyTheory/Proofs/UnitaryExistence/ExistsBinaryEqAdicCompletion.lean}{GlobalTheory.LocalNorm.exists_eq_sq_add_of_notMem}
\begin{proof}
(1) By the proof of Lemma~\ref{cv-volume-key-xi}(1), $\overline{\sqrt{-d}}=-\sqrt{-d}$, so $\bar\alpha=\alpha$ and
$\bar\beta=\beta$. As $d/(2\sqrt{-d})^2=-1/4$, $\alpha^2+d\beta^2=\big((z+\bar z)^2-(z-\bar z)^2\big)/4=z\bar z$.

(2) Let $P$ be a prime ideal of $\mathcal O_E$ with $P\cap\mathcal O_{E^+}=\mathfrak p$. First, there is
$w\in\mathcal O_E$ with $w-\bar w\notin P$. As no prime of $E^+$ ramifies in $E$, the ring map
$\mathcal O_{E^+}\to\mathcal O_E$ is unramified. Let $\mu$ be the multiplication map $x\otimes y\mapsto xy$ of
$\mathcal O_E\otimes_{\mathcal O_{E^+}}\mathcal O_E$. There is an idempotent $t$ of this tensor product such that $\mu$
is the localization at $t$ \cite[Tag 02FL]{Sta}. So $\mu(t)=1$, and $t$ kills the kernel of $\mu$, which contains
$1\otimes b-b\otimes1$ for every $b\in\mathcal O_E$. Let $\psi$ be the ring map $x\otimes y\mapsto x\bar y$. The element
$\sqrt{-d}$ lies in $\mathcal O_E$, as a root of $X^2+d$, and $\overline{\sqrt{-d}}=-\sqrt{-d}\ne\sqrt{-d}$ as $d\ne0$.
Applying $\psi$ to $(1\otimes b-b\otimes1)\,t=0$ for $b=\sqrt{-d}$ gives $(\bar b-b)\,\psi(t)=0$, so $\psi(t)=0$, as
$\mathcal O_E$ is a domain. Suppose that $b-\bar b\in P$ for every $b\in\mathcal O_E$. Then
$\mu(x\otimes y)-\psi(x\otimes y)=x(y-\bar y)$ lies in $P$ for all $x,y\in\mathcal O_E$, so $\mu(z)-\psi(z)\in P$ for
every $z$ in the tensor product, as both maps are additive. For $z=t$ this puts $1=\mu(t)-\psi(t)$ in $P$, which is
absurd. This gives $w$.

Put $a_1=w+\bar w$ and $n=w\bar w$, which lie in $\mathcal O_{E^+}$. Then
$a_1^2-4n=(w-\bar w)^2\notin P$, so $a_1^2-4n\notin\mathfrak p$. By Lemma~\ref{cv-volume-key-models-binary} there are
$r,s\in E^+_{\mathfrak p}$ with $r^2+a_1rs+ns^2=u$. With $\alpha$ and $\beta$ as in (1) for $z=w$, $a_1=2\alpha$ and
$n=\alpha^2+d\beta^2$, so $u=(r+\alpha s)^2+d(\beta s)^2$.
\end{proof}
% Lean: (1) exists_coords with complexConj_rt; (2) exists_eq_sq_add_of_notMem: P from
% Ideal.exists_ideal_over_maximal_of_isIntegral, w from LocalNorm.exists_sub_notMem (A = O_{E^+}, B = O_E, sigma the
% conjugation IsCMField.ringOfIntegersComplexConj, moved element H.rt, unramified by the field H.unram), a and n from
% IsCMField.ringOfIntegersComplexConj_eq_self_iff, hD, exists_binary_eq_adicCompletion and exists_coords.
% exists_sub_notMem (about 30 lines): t from Algebra.FormallyUnramified.iff_exists_tensorProduct (Mathlib: t with
% (1 (x) s - s (x) 1) t = 0 and lmul' t = 1; the text cites Stacks 02FL, whose idempotent has both properties);
% psi = productMap (id) sigma, h0: psi(t) = 0 from ht b and sigma b != b; h1: lmul' z - psi z in P by induction on z
% (tmul, add); then 1 = lmul' t - psi t in P.

\begin{lemma}\label{cv-volume-key-models-local}
Let $E$, $E^+$, $d$ and $\sqrt{-d}$ be as in Lemma~\ref{cv-volume-key-models-units}, and suppose that no finite prime of
$E^+$ ramifies in $E$. Let $x\in E^+$ be nonzero, and let $I$ and $J$ be nonzero ideals of $\mathcal O_E$ with
$xI\bar I=J\bar J$ as fractional ideals. Then for every finite place $v$ of $E^+$ there are $a,b\in E^+_v$ with
$x=a^2+db^2$; that is, $x$ is a norm from $E\otimes_{E^+}E^+_v=E^+_v[X]/(X^2+d)$.
\end{lemma}
\leanref{Development/AbelianVarietyTheory/Proofs/UnitaryExistence/ExistsBinaryEqAdicCompletion.lean}{GlobalTheory.local_norm_of_ideal_norm}
\begin{proof}
Let $\mathfrak p$ be the prime ideal of $\mathcal O_{E^+}$ of $v$. By Lemma~\ref{cv-volume-key-models-unit} there are
$z_1\in E^\times$, $z_2\in E$ and $u\in\mathcal O_{E^+}\setminus\mathfrak p$ with $x\,z_1\bar z_1=u\,z_2\bar z_2$. By
Lemma~\ref{cv-volume-key-models-units}, $u=a_0^2+db_0^2$ with $a_0,b_0\in E^+_v$, and $z_k\bar z_k=a_k^2+db_k^2$ with
$a_k,b_k\in E^+$ for $k=1,2$; here $a_1^2+db_1^2=z_1\bar z_1\ne0$. The identity
$(a^2+db^2)(a'^2+db'^2)=(aa'-dbb')^2+d(ab'+a'b)^2$ shows that the values of $a^2+db^2$ on $E^+_v$ are closed under
products, and they are closed under division by nonzero squares. So
\[
 x=\frac{(a_0^2+db_0^2)(a_2^2+db_2^2)(a_1^2+db_1^2)}{(a_1^2+db_1^2)^2}
\]
is such a value.
\end{proof}
% Lean: local_norm_of_ideal_norm: exists_mul_norm_eq (Lemma cv-volume-key-models-unit), exists_eq_sq_add_of_notMem and
% exists_coords (Lemma cv-volume-key-models-units), and exists_eq_sq_add_of_mul_eq, which writes x as a^2 + d b^2 by an
% explicit formula for the quotient.

\begin{lemma}\label{cv-volume-key-models}
In the setting of Definition~\ref{cv-volume-twisting-section}, let $L$ be the lattice of the PEL datum of $X$, of rank
$r=p+q$, let $\sigma$ be a good embedding at which the uniformization satisfies the compatibility conditions of the
moduli problem, and let $L_1,\dots,L_n$, with forms $\langle\ ,\ \rangle_i$, be the lattices of the uniformization
family at $\sigma$ that these conditions include: lattices of PEL data of signature $(p,q)$ over $E$, with maps $t_i$
and points $y_i(P)$ that have the properties of Theorem~\ref{cv-uniformization-family} at the level $(M,\mathfrak M,N)$
of $X$. Let $\xi_0$ be as in Lemma~\ref{cv-volume-key-xi}(1), and $G$ and $\varphi$ as in part~(2) of that lemma for $L$
and $\xi=-\xi_0$. A \emph{frame} is a family $x=(x_\tau)_{\tau\in\Phi}$ of matrices in $GL_r(\mathbf C)$ such that, for
$0\le j\le r$, the leading principal $j\times j$ minor of the Gram matrix $\big(\tau h_G(x_{\tau,a},x_{\tau,b})\big)_{a,b}$
of the columns $x_{\tau,a}$ of $x_\tau$ has the sign $(-1)^{\max(0,j-q)}$. Then for every $i$ there is an injective
$\mathcal O_E$-linear map $\varphi_i\colon L_i\to E^r$, whose image $\Lambda_i$ spans $E^r$, with
$\langle x,y\rangle_i=\operatorname{Tr}_{E/\mathbf Q}\big(-\xi_0h_G(\varphi_ix,\varphi_iy)\big)$ for $x,y\in L_i$, with
the following properties. Let $\Gamma_i$ be the group of $\gamma\in GL_r(E)$ with $h_G(\gamma u,\gamma v)=h_G(u,v)$ for
$u,v\in E^r$ and $\gamma\Lambda_i=\Lambda_i$, such that $\gamma\varphi_i(a)\in\varphi_i(a)+M\Lambda_i$ for
$a\in\bar{\mathfrak M}L_i$ and $\gamma\varphi_i(a)\in\varphi_i(a)+N\Lambda_i$ for $a\in L_i$.
\begin{enumerate}
\item There is $c\in\mathcal O_E\setminus\{0\}$ such that $c\,h_G(u,v)$ is integral for all $u,v\in\Lambda_i$.
\item $\Gamma_i$ has finite index in the group of $\gamma\in GL_r(E)$ that preserve $h_G$ and $\Lambda_i$.
\item If $\gamma\in\Gamma_i$, $x$ is a frame, $b$ is a family of invertible upper triangular matrices and
  $\tau(\gamma)x_\tau=x_\tau b_\tau$ for every $\tau\in\Phi$, then $\gamma=1$.
\item There is an integer $m\ge1$ such that, for every $u\in\Lambda_i$ with $h_G(u,u)=0$, every integral $s\in E$
  with $\bar s=s$ and $\kappa_0=\pm2\xi_0$, the transvection $v\mapsto v+ms\kappa_0\,h_G(v,u)\,u$ lies in $\Gamma_i$.
\item If $x$ is a frame, so are $(\tau(\gamma)x_\tau)_\tau$ for $\gamma\in\Gamma_i$ and $(x_\tau b_\tau)_\tau$ for every
  family $b$ of invertible upper triangular matrices.
\end{enumerate}
\end{lemma}
\leanref{Development/AbelianVarietyTheory/Proofs/UnitaryExistence/ExistsFlagDomainModelsOfFamilyUniformisation.lean}{GlobalTheory.ModelsExist.exists_modelsOf}
\leanref{Development/AbelianVarietyTheory/Definitions/ExistenceOfLattices/NormConst.lean}{GlobalTheory.FrameUnif.ModelBase}
\leanref{Development/AbelianVarietyTheory/Definitions/ExistenceOfLattices/NormConst.lean}{GlobalTheory.FrameUnif.NegXi}
\begin{proof}
Notation. Fix $i$ and put $\eta=-\xi_0$, so that $\bar\eta=-\eta\ne0$ and $\operatorname{Im}\tau(\eta)>0$ for
$\tau\in\Phi$. Let $E^+$ be the totally real subfield of $E$; as $\overline{\sqrt{-d}}=-\sqrt{-d}$,
$E=E^+(\sqrt{-d})$.

The models. Step 1: coordinates. The lattice $L_i$ has signature $(p,q)$, so it has rank $r=p+q$. By
Lemma~\ref{cv-volume-key-xi}(2), for $L_i$ and $\xi=\eta$, there are an invertible $G'\in M_r(E)$ with
$\bar G'_{ab}=G'_{ba}$ and an injective $\mathcal O_E$-linear map $\psi\colon L_i\to E^r$ whose image spans $E^r$, with
$\langle x,y\rangle_i=\operatorname{Tr}_{E/\mathbf Q}(\eta\,h_{G'}(\psi x,\psi y))$.

Step 2: signatures. The triples $(G,\eta,\varphi)$ and $(G',\eta,\psi)$ are Hermitian models of $L$ and of $L_i$
(Definition~\ref{cv-uniformization-model}), both of signature $(p,q)$. By Lemma~\ref{cv-uniformization-signature}, for
$\tau\in\Phi$ both $\tau(G)$ and $\tau(G')$ have $q$ positive and $p$ negative eigenvalues. The matrices $\bar\tau(G)$ and
$\bar\tau(G')$ are the complex conjugates of these, with the same signatures.

Step 3: the determinant at the real places. As $\bar G=G^{\mathsf T}$, $\det G$ is fixed by complex conjugation, and so is
$\det G'$. So $\Delta=\det G'/\det G$ is a nonzero element of $E^+$. A real place of $E^+$ extends to an embedding
$\tau$ of $E$, and $\tau(\det G)$ and $\tau(\det G')$ are real of the sign $(-1)^p$, by Step 2. So $\Delta$ is positive at
every real place, a norm from $\mathbf C$.

Step 4: the determinant at the finite places. By Lemma~\ref{cv-volume-key-models-ideal}, applied to the Hermitian
models $(G,\eta,\varphi)$ of $L$ and $(G',\eta,\psi)$ of $L_i$, there are nonzero ideals $I$ and $J$ of $\mathcal O_E$
with $\Delta I\bar I=J\bar J$. By Step 3, $\Delta$ is a nonzero element of $E^+$, and no finite prime of $E^+$ ramifies in
$E$ (Situation~\ref{cv-function-situation}). So, by Lemma~\ref{cv-volume-key-models-local}, for every finite place $v$
of $E^+$ there are $a,b\in E^+_v$ with $\Delta=a^2+db^2$, the norm of $a+b\sqrt{-d}$ from $E\otimes_{E^+}E^+_v$.
% Lean: GlobalTheory.DetIdeal.det_ideal (Lemma cv-volume-key-models-ideal: the volume ideals vol of the self-dual
% lattices phi(L) and psi(L_i)) and GlobalTheory.local_norm_of_ideal_norm (Lemma cv-volume-key-models-local, with the
% field H.unram of the host), used in GlobalTheory.model_in_same_space as hI and hloc. Size: about 740 lines, the files
% DetMulVolMulConj.lean and ExistsBinaryEqAdicCompletion.lean.

Step 5: the norm. By Steps 3 and 4 and the Hasse norm theorem for the cyclic extension $E/E^+$
\cite[Chapter~VI, (4.5)]{Neu99}, $\Delta$ is a norm: $\Delta=a^2+db^2$ with $a,b\in E^+$. So $\det G'=z\bar z\det G$ with
$z=a+b\sqrt{-d}$.

Step 6: the isometry. By Landherr's theorem \cite{Lan36}, two nondegenerate Hermitian forms over $E/E^+$ of the same
dimension, with the same signatures at the real places and with determinants that differ by a norm, are isometric. By
Steps 2 and 5 there is $\gamma\in GL_r(E)$ with $h_G(\gamma u,\gamma v)=h_{G'}(u,v)$ for $u,v\in E^r$. Put
$\varphi_i=\gamma\circ\psi$. It is injective and $\mathcal O_E$-linear, its image $\Lambda_i$ spans $E^r$, and
$\langle x,y\rangle_i=\operatorname{Tr}_{E/\mathbf Q}(\eta\,h_{G'}(\psi x,\psi y))
=\operatorname{Tr}_{E/\mathbf Q}(-\xi_0h_G(\varphi_ix,\varphi_iy))$.
% Lean: exists_modelsOf applies GlobalTheory.T.latticeModel_of_same_signature_hasseNorm =
% FrameUnif.latModel_of_same_signature to each lattice F.Di i; the model is GlobalTheory.model_in_same_space, with
% exists_herm_coord_of_datum (Step 1), signature_hermitian_model_of_datum and gram_conj (Step 2), cc_det and det_gram
% (Step 3), the published GlobalTheory.Published.HasseNormSt for the quadratic extension E^+(sqrt(-d)) of E^+
% (Step 5), and LocalQuad.landherrSt_of_hasseNormSt, the Lean's proof of Landherr's theorem in this case (Step 6).

(1) Let $u,v\in\Lambda_i$. For $b\in\mathcal O_E$,
$\operatorname{Tr}_{E/\mathbf Q}\big(b\,\eta h_G(u,v)\big)=\langle b\varphi_i^{-1}u,\varphi_i^{-1}v\rangle_i\in\mathbf Z$, so
$\eta h_G(u,v)$ lies in the inverse different of $E$, and $D_E\eta h_G(u,v)$ is integral for the discriminant $D_E$ of
$E$. Choose an integer $n\ge1$ with $n\xi_0$ integral. Then $c=D_En\xi_0$ is a nonzero integer of $E$, and
$c\,h_G(u,v)=-nD_E\eta h_G(u,v)$ is integral.
% Lean: FrameUnif.exists_den_herm_model, den = discr E * (n xi): HermModel.rank_integral_selfDual (the trace dual and
% isIntegral_discr_mul_of_mem_traceDual) and IsAlgebraic.exists_integral_multiple.

(2) The group $U$ of the $\gamma\in GL_r(E)$ that preserve $h_G$ and $\Lambda_i$ acts on the finite group
$\Lambda_i/MN\Lambda_i$. Let $K$ be the kernel of this action. An element $\gamma\in K$ satisfies
$\gamma\varphi_i(a)\in\varphi_i(a)+MN\Lambda_i$ for every $a\in L_i$, so it satisfies both congruences in the definition of
$\Gamma_i$, and $K\subset\Gamma_i$. As $K$ has finite index in $U$, so has $\Gamma_i$.
% Lean: FrameUnif.princ_finiteIndex with CongruenceIndex.exists_finset_cosets for the modulus l.M * l.A.

(3) Let $\gamma\in\Gamma_i$, $x$ a frame and $b$ a family of invertible upper triangular matrices with
$\tau(\gamma)x_\tau=x_\tau b_\tau$ for every $\tau\in\Phi$.

The congruence. The map $\gamma'=\varphi_i^{-1}\gamma\varphi_i$ is an $\mathcal O_E$-linear automorphism of $L_i$ that
preserves the form, as $\gamma$ preserves $h_G$, and $\gamma'(a)\in a+ML_i$ for $a\in\bar{\mathfrak M}L_i$. By
Lemma~\ref{cv-volume-key-models-half}, $\gamma'(a)\in a+ML_i$ for every $a\in L_i$, that is $\gamma v\in v+M\Lambda_i$
for every $v\in\Lambda_i$.

Bounded powers. Let $\tau\in\Phi$ and $h=h_{\tau(G)}$. The leading principal minors of the Gram matrix of $x_\tau$ are
not zero, so by the Gram--Schmidt process there is an upper triangular $U$ with diagonal entries $1$ such that the
columns $x'_1,\dots,x'_r$ of $x'=x_\tau U$ satisfy $h(x'_a,x'_c)=0$ for $a\ne c$ and $d_c=h(x'_c,x'_c)\ne0$. Put
$b'=U^{-1}b_\tau U$, an upper triangular matrix with $\tau(\gamma)x'=x'b'$. As $\tau(\gamma)$ preserves $h$, the columns
$\sum_{a\le c}b'_{ac}x'_a$ of $x'b'$ satisfy the same relations as the $x'_c$. By induction on $c$ this gives
$b'_{ac}=0$ for $a<c$ and $|b'_{cc}|=1$: if $b'$ is diagonal with entries of absolute value $1$ in its first $c-1$
columns, the orthogonality of the column $c$ to the columns $a<c$ gives $b'_{ac}\overline{b'_{aa}}d_a=0$, and then
$h$ of the column $c$ with itself gives $|b'_{cc}|^2d_c=d_c$. So $\tau(\gamma)^k=x'b'^kx'^{-1}$ has entries bounded
independently of $k\ge0$.

Finite order. As $\Lambda_i$ spans $E^r$ and is finitely generated, there are an integer $l\ge1$ with
$l\varepsilon_c\in\Lambda_i$ for the unit vectors $\varepsilon_c$ of $E^r$, and an integer $l'\ge1$ with
$l'\Lambda_i\subset\mathcal O_E^r$. As $\gamma^k\Lambda_i=\Lambda_i$, the entries of $ll'\gamma^k$ are integers of $E$. Their
images under the embeddings in $\Phi$ are bounded independently of $k$, and so are those under the conjugate
embeddings. The integers of $E$ whose images under all embeddings are bounded form a finite set. So $\gamma^k$ takes
finitely many values, and $\gamma$ has finite order.

Minkowski. The restriction of $\gamma$ to $\Lambda_i$ has finite order and is congruent to $1$ modulo $M\ge3$, so it is
the identity \cite{Min87}. As $\Lambda_i$ spans $E^r$, $\gamma=1$.
% Lean: PrincFree.princ_free: exists_equiv_of_stab for gamma', trivial_mod_of_trivial_on_half (Lemma
% cv-volume-key-models-half), and GamFree.gam_free: pow_bounded (TriangularIsometryBoundedPowers.
% bounded_pow_of_triangular_isometry: exists_diag for U and d, diag_of_isometry for b', bounded_pow) with
% gram_unitary, isOfFinOrder_of_bounded (LatFin.exists_common_den, exists_int_smul_mem and finite_bounded), and the
% published Minkowski statement latMinkowski for the restriction to Lambda_i.

(4) Put $m=MN|D_E|$. Let $u\in\Lambda_i$ with $h_G(u,u)=0$, $s$ an integer of $E$ with $\bar s=s$, $\kappa_0=\pm2\xi_0$, and
$\alpha=ms\kappa_0$, so that $\bar\alpha=-\alpha$. Let $T(v)=v+\alpha h_G(v,u)u$. As $h_G(u,u)=0$ and $h_G$ is Hermitian,
\[
 h_G(Tv,Tw)=h_G(v,w)+(\alpha+\bar\alpha)\,h_G(v,u)h_G(u,w)=h_G(v,w).
\]
For $v\in\Lambda_i$, $D_E\kappa_0h_G(v,u)=\mp2D_E\eta h_G(v,u)$ is integral by the proof of (1), so
$s\,D_E\kappa_0h_G(v,u)\,u\in\Lambda_i$, as $\Lambda_i$ is an $\mathcal O_E$-module, and
$Tv-v=\pm MN\cdot s\,D_E\kappa_0h_G(v,u)\,u\in MN\Lambda_i$. As $h_G(Tv,u)=h_G(v,u)$, the inverse of $T$ is
$v\mapsto v-\alpha h_G(v,u)u$, which also maps $\Lambda_i$ into $\Lambda_i$. So $T\Lambda_i=\Lambda_i$, and $T$ satisfies both
congruences in the definition of $\Gamma_i$; hence $T\in\Gamma_i$.
% Lean: PrincTransv.transv_mem_princ with m = l.M * l.A * |discr E|: rank_integral_selfDual for hint,
% FlagData.smul_mem_of_isIntegral, FlagData.herm_transvMat, transvMat_mulVec and herm_add_smul_left.

(5) Right multiplication by an upper triangular $b_\tau$ replaces the first $j$ columns of $x_\tau$ by their
combinations with the leading $j\times j$ block of $b_\tau$, so it multiplies the $j$-th leading principal minor of
the Gram matrix by $|b_{\tau,11}\cdots b_{\tau,jj}|^2>0$. An element $\gamma\in\Gamma_i$ preserves $h_G$, so
$\tau(\gamma)x_\tau$ has the same Gram matrix as $x_\tau$; and $\tau(\gamma)x_\tau$ is invertible.
\end{proof}
% Lean: ModelBase is the rank r = p + q, the type (p0 = q, Psi = Phi), the trivial characters chi, ef (the form of
% L_i), gam (Gamma_i is princ l Gm), and domStab, which is (5) (Formalism.L2.Pj.dom_rmul and dom_lact with ustab_mat);
% NegXi is Im tau(xi_0) < 0 for tau in Phi. A frame is an element of GG Psi Gm q (dmin and psgn).

\begin{situation}\label{cv-volume-key-rebuilt}
In the setting of Definition~\ref{cv-volume-twisting-section}, let $\sigma$ be a good embedding of $F_c$ at which the
uniformization satisfies the compatibility conditions of the moduli problem, and let $L_i$, $t_i$, $y_i(P)$,
$\xi_0$, $G$, $\varphi_i$, $\Lambda_i$, $\Gamma_i$ and the frames be those of Lemma~\ref{cv-volume-key-models} for
$\sigma$. Index the coordinates of $\mathbf C^g$ by the pairs $(\tau,a)$ of an embedding $\tau\colon E\to\mathbf C$ and
an integer $0\le a<m_\tau$, the coordinate $(\tau,a)$ having the embedding $\sigma\circ\tau$, and write $\sigma_k$ for
the embedding of the coordinate $k$.

\emph{Periods.} Write $x_{\tau,b}$ for the $b$-th column of $x_\tau$ and
$h_{\tau(G)}(u,v)=\sum_{a,b}u_a\tau(G_{ab})\bar v_b$ on $\mathbf C^r$. The \emph{period} of a frame $x$ on $L_i$ is the
additive map $P^i_x\colon L_i\to\mathbf C^g$ with coordinates, for $\tau\in\Phi$,
\[
 P^i_x(v)_{(\tau,a)}=\big(x_\tau^{-1}\tau(\varphi_iv)\big)_{q+a+1},\qquad
 P^i_x(v)_{(\bar\tau,a)}=h_{\tau(G)}\big(x_{\tau,q-a},\tau(\varphi_iv)\big).
\]

\emph{Theta data.} \emph{Theta data} for $i$ are a basis $(e_1,\dots,e_g,f_1,\dots,f_g)$ of $L_i$ as in
Lemma~\ref{cv-uniformization-mod-eight}, for the prime $d$ and $w=(1+\sqrt{-d})/2$; functions $x\mapsto h_x$ and
$x\mapsto Z_x$ from the frames to $M_g(\mathbf C)$, holomorphic in $x$, with $\det h_x\ne0$, $Z_x$ symmetric with
positive definite imaginary part, $h_xP^i_x(e_k)$ the $k$-th unit vector and $h_xP^i_x(f_k)$ the $k$-th column of
$Z_x$; a characteristic $c_i\in\{0,\frac12\}^{2g}$; and a factor $F^i$ for the periods $P^i_x$, as in
Theorem~\ref{cv-uniformization-family}(4), such that for every frame $x$ the line bundle of the polarization on
$A_{y_i(P^i_x)}$ has a section whose value at the image $p_x(z)$ of $z\in\mathbf C^g$, in the trivialization over $U_j$,
is $\vartheta[c_i](Z_x,h_xz)F^i_j(x,z)$. Fix theta data for every $i$.

\emph{Normalized framed points.} Put $N_\xi=\big(2^g/\prod_k2|\sigma_k(\xi_0)|\big)^{1/2}$. A \emph{normalized framed
point} at a frame $x$, for $i$, is a complex framed point $(y,v,w)$ of $X$ over $\sigma$ with $y=y_i(P^i_x)$, with $v$
the images in $\operatorname{Lie}A_y$ of the unit vectors of $\mathbf C^g$, and with $w$ the section of the line bundle
of the polarization whose value at $p_x(z)$ is $\nu\,\vartheta[c_i](Z_x,h_xz)F^i_j(x,z)$, for a number $\nu$ with
$\nu^2=N_\xi\det h_x$.
\end{situation}
\leanref{Development/AbelianVarietyTheory/Definitions/ExistenceOfLattices/NormConst.lean}{GlobalTheory.FrameUnif.normConst}
\leanref{Development/AbelianVarietyTheory/Definitions/ExistenceOfLattices/NormConst.lean}{GlobalTheory.FrameUnif.IsThetaPt}

\begin{lemma}\label{cv-volume-key-normalized-tangent}
In Situation~\ref{cv-volume-key-rebuilt}, let $i$ be an index, let $P$ be a period of $L_i$ of the type $(\sigma_k)$ of
the coordinates (Definition~\ref{cv-cm-point-periods}), let $y=y_i(P)$, and let $p\colon\mathbf C^g\to A_y$ be the
quotient map $\mathbf C^g\to\mathbf C^g/P(L_i)$ followed by the isomorphism of Theorem~\ref{cv-uniformization-family}.
For $w\in\mathbf C^g$ let $dp(w)\in\operatorname{Lie}A_y$ be the tangent vector at the origin along which the derivative
of every regular function $f$ near the origin is the derivative of $f\circ p$ at $0$ along $w$, and for a coordinate
$(\tau,a)$ let $\varepsilon_{\tau,a}$ be its unit vector.
\begin{enumerate}
\item The map $w\mapsto dp(w)$ is an isomorphism of complex vector spaces from $\mathbf C^g$ onto
  $\operatorname{Lie}A_y$.
\item Let $\tau$ be an embedding of $E$. Then $dp(w)$ lies in the eigenspace $\operatorname{Lie}_\tau(y)$, on which
  $\mathcal O_E$ acts through $\sigma\circ\tau$, if and only if $w$ vanishes at every coordinate that is not of the form
  $(\tau,a)$.
\item For every $\tau$ the vectors $v_{\tau,a}=dp(\varepsilon_{\tau,a})$, $0\le a<m_\tau$, form a basis of
  $\operatorname{Lie}_\tau(y)$.
\end{enumerate}
\end{lemma}
\leanref{Development/AbelianVarietyTheory/Proofs/UnitaryExistence/ExistsTangentEigenbasisExponentialStatement.lean}{GlobalTheory.T.exists_tangent_eigenbasis_exponential}
\leanref{Development/AbelianVarietyTheory/Proofs/UnitaryExistence/ExponentialDerivativeAtZeroIsTangent.lean}{GlobalTheory.isDerivAt0_ex_isTanAt}
\leanref{Development/AbelianVarietyTheory/Proofs/UnitaryExistence/IsEigAtIffOfIsDerivAt0.lean}{GlobalTheory.isEigAt_iff_of_isDerivAt0}
\leanref{Development/AbelianVarietyTheory/Proofs/UnitaryExistence/IsCombAtIffIsDerivAt0.lean}{GlobalTheory.DerLin.isCombAt_iff_isDerivAt0}
\begin{proof}
(1) By Theorem~\ref{cv-uniformization-family}, $p$ is a holomorphic homomorphism of complex Lie groups onto $A_y$
whose kernel $P(L_i)$ is discrete. So $p$ maps a neighbourhood of $0$ biholomorphically onto a neighbourhood of the
origin of $A_y$. Hence $dp(w)$ exists and is unique for every $w$, it depends linearly on $w$, and every tangent vector
of $A_y$ at the origin is $dp(w)$ for one and only one $w$.

(2) For $b\in\mathcal O_E$ and $z\in\mathbf C^g$ let $b\cdot z$ be the vector with the coordinates $\sigma_k(b)z_k$, and
let $[b]$ be the action of $b$ on $A_y$. The isomorphism of Theorem~\ref{cv-uniformization-family} is compatible with
the actions of $\mathcal O_E$, so $p(b\cdot z)=[b]p(z)$, and the derivative at $0$ along $w$ gives
$dp(b\cdot w)=[b]_*dp(w)$. If $w$ vanishes at the coordinates not of the form $(\tau,a)$, then
$b\cdot w=(\sigma\circ\tau)(b)w$ for every $b$, so $dp(w)\in\operatorname{Lie}_\tau(y)$. Conversely, if
$dp(w)\in\operatorname{Lie}_\tau(y)$, then $dp(b\cdot w)=dp((\sigma\circ\tau)(b)w)$, and by (1)
$\sigma_k(b)w_k=(\sigma\circ\tau)(b)w_k$ for every $b$ and every coordinate $k$. At a coordinate $k=(\tau',a)$ with
$w_k\ne0$ the embeddings $\sigma\circ\tau'$ and $\sigma\circ\tau$ agree on $\mathcal O_E$, hence on $E$, and $\tau'=\tau$
as $\sigma$ is injective.

(3) By (1) and (2), $dp$ maps the subspace of the vectors that vanish outside the coordinates $(\tau,a)$, of which the
$\varepsilon_{\tau,a}$ form a basis, isomorphically onto $\operatorname{Lie}_\tau(y)$.
\end{proof}
% Lean: GlobalTheory.LieFrame.exists_eigenbasis: (1) is isDerivAt0_ex_isTanAt, from the field ex_diff of FamilyUnif,
% with the linearity from DerLin.isCombAt_iff_isDerivAt0 (exists_tangent_module_universal_group and
% derivative_along_tangent_combination_moduli); (2) is isEigAt_iff_of_isDerivAt0 (ex_lin, isDerivAt_of_comp_linear,
% eq_of_forall_fderiv_eq and DerivEig.eq_of_relab_eq); (3) is eq_sum_single and sum_single_apply. The map p is
% z => (F.ex k P z).left, and dp(w) is the v with IsDerivAt0 for w.

\begin{lemma}\label{cv-volume-key-normalized-theta}
Let $g\ge1$, let $Z\in\mathfrak H_g$, let $c=(a,b)$ with $a,b\in\mathbf R^g$, and let $\vartheta[c]$ be as in
Theorem~\ref{cv-uniformization-theta-transformation}. Then $z\mapsto\vartheta[c](Z,z)$ is holomorphic on
$\mathbf C^g$, and for every $z_0\in\mathbf C^g$ there is $t\in[0,1]^g$ with $\vartheta[c](Z,z_0+t)\ne0$.
\end{lemma}
\leanref{Development/AbelianVarietyTheory/Proofs/UnitaryExistence/ThetaCharHalfDifferentiableNeZeroFamiliesStatement.lean}{GlobalTheory.T.thetaChar_half_differentiable_ne_zero_families}
\leanref{Development/AbelianVarietyTheory/Definitions/ThetaMetrics/ExistsEventuallyNormRiemannThetaTermLe.lean}{GlobalTheory.IsSiegel.exists_riemannTheta_add_ofReal_ne_zero}
\begin{proof}
Let $T_n(z)$ be the term of index $n\in\mathbf Z^g$ of the series. As $\operatorname{Im}Z$ is positive definite, the
series converges absolutely and uniformly on compact subsets of $\mathbf C^g$ \cite[Chapter~II, \S1]{Mum83}, so its sum
is holomorphic. For $t\in\mathbf R^g$, $T_n(z_0+t)=T_n(z_0)e^{2\pi i(n+a)^{\mathsf T}t}$, and
$\int_{[0,1]^g}e^{2\pi in^{\mathsf T}t}\,dt$ is $1$ for $n=0$ and $0$ otherwise. Integrating term by term, which the
uniform convergence allows,
\[
 \int_{[0,1]^g}\vartheta[c](Z,z_0+t)\,e^{-2\pi ia^{\mathsf T}t}\,dt=T_0(z_0)
 =\exp\big(\pi ia^{\mathsf T}Za+2\pi ia^{\mathsf T}(z_0+b)\big)\ne0.
\]
So $\vartheta[c](Z,z_0+t)\ne0$ for some $t\in[0,1]^g$.
\end{proof}
% Lean: ThetaHalf.thetaChar_holomorphic_ne_zero, for characteristics in {0, 1/2}^{2g}: CanTh.differentiable_thetaChar,
% and CanTh.thetaChar_ne_zero through theta_ne_zero, IsSiegel.riemannTheta_ne_zero and
% exists_riemannTheta_add_ofReal_ne_zero (integral_riemannTheta_mul_cexp at n = 0). Its second part, the holomorphy
% in families, is not used here.

\begin{lemma}\label{cv-volume-key-normalized-gram}
In Situation~\ref{cv-volume-key-rebuilt}, let $x$ be a frame and $\tau\in\Phi$, and write $h=h_{\tau(G)}$. Then there is
an upper triangular matrix $U\in GL_r(\mathbf C)$ with all diagonal entries $1$ such that the columns
$x'_1,\dots,x'_r$ of $x_\tau U$ satisfy $h(x'_a,x'_b)=0$ for $a\ne b$, $h(x'_b,x'_b)>0$ for $b\le q$ and
$h(x'_b,x'_b)<0$ for $b>q$.
\end{lemma}
\leanref{Development/AbelianVarietyTheory/Proofs/UnitaryExistence/HermFormConjDiag.lean}{GlobalTheory.FramePeriod.exists_upper_diag}
\begin{proof}
Let $\delta_j$ be the leading principal $j\times j$ minor of the Gram matrix $(h(x_{\tau,a},x_{\tau,b}))_{a,b}$; then
$\delta_0=1$, and by the definition of a frame $\delta_j$ is real and not zero, with the sign $(-1)^{\max(0,j-q)}$.
Construct the $x'_b$ in turn. Suppose that $x'_1,\dots,x'_{b-1}$ are pairwise orthogonal and that each $x'_c$ is
$x_{\tau,c}$ plus a combination of $x_{\tau,1},\dots,x_{\tau,c-1}$. The Gram matrix of $x'_1,\dots,x'_{b-1}$ is that of
$x_{\tau,1},\dots,x_{\tau,b-1}$ changed by a unitriangular matrix, so its determinant
$\prod_{c<b}h(x'_c,x'_c)$ is $\delta_{b-1}\ne0$. Put
$x'_b=x_{\tau,b}-\sum_{c<b}h(x_{\tau,b},x'_c)\,h(x'_c,x'_c)^{-1}x'_c$; it is orthogonal to $x'_1,\dots,x'_{b-1}$. In the end
$\delta_b=\prod_{c\le b}h(x'_c,x'_c)$ for every $b$, so $h(x'_b,x'_b)=\delta_b/\delta_{b-1}$ has the sign
$(-1)^{\max(0,b-q)-\max(0,b-1-q)}$, which is $1$ for $b\le q$ and $-1$ for $b>q$.
\end{proof}
% Lean: exists_upper_diag, with U = FrameDegreeData.uMat (uGu_lower, uMat_upper, uMat_diag, lminor_conj and
% lminor_diagonal); the Lean writes the Gram matrix of the columns of x U as U^T A conj(U).

\begin{lemma}\label{cv-volume-key-normalized-periods}
In Situation~\ref{cv-volume-key-rebuilt}, let $i$ be an index.
\begin{enumerate}
\item For every frame $x$, $P^i_x$ is a period of $L_i$ of the type $(\sigma_k)$ of the coordinates
  (Definition~\ref{cv-cm-point-periods}).
\item For every $v\in L_i$, $x\mapsto P^i_x(v)$ is a holomorphic function on the set of frames.
\item Every period of $L_i$ of the type $(\sigma_k)$ is $P^i_x$ for some frame $x$.
\end{enumerate}
\end{lemma}
\leanref{Development/AbelianVarietyTheory/Proofs/UnitaryExistence/ExistsFramePeriodsStatement.lean}{GlobalTheory.T.exists_framePeriods_alt}
\leanref{Development/AbelianVarietyTheory/Proofs/UnitaryExistence/HermFormConjDiag.lean}{GlobalTheory.FramePeriod.isPeriod_frame}
\leanref{Development/AbelianVarietyTheory/Proofs/UnitaryExistence/HermFormConjDiag.lean}{GlobalTheory.FramePeriod.exists_loc}
\leanref{Development/AbelianVarietyTheory/Proofs/UnitaryExistence/DetOfTriangular.lean}{GlobalTheory.per_move}
\leanref{Development/AbelianVarietyTheory/Proofs/UnitaryExistence/ExistsFrameOfIsPeriod.lean}{GlobalTheory.exists_frame_of_isPeriod}
\leanref{Development/AbelianVarietyTheory/Proofs/UnitaryExistence/ExistsFrameOfIsPeriod.lean}{GlobalTheory.PeriodOnto.exists_frame}
\begin{proof}
Notation. Put $\eta=-\xi_0$, so that $\bar\eta=-\eta$, $\operatorname{Im}\tau(\eta)>0$ for $\tau\in\Phi$
(Lemma~\ref{cv-volume-key-xi}(1)) and $\langle u,u'\rangle_i=\operatorname{Tr}_{E/\mathbf Q}(\eta\,h_G(\varphi_iu,\varphi_iu'))$.
For $\tau\in\Phi$ write $h_\tau=h_{\tau(G)}$. The matrix $\tau(G)$ is Hermitian, and
$\tau(h_G(u,u'))=h_\tau(\tau(u),\tau(u'))$ for $u,u'\in E^r$, as $\tau(\bar b)=\overline{\tau(b)}$ on the CM field $E$.
The embeddings of $E$ are those of $\Phi$ and their conjugates, and $\tau(\eta)=i\operatorname{Im}\tau(\eta)$, so
\[
 \langle u,u'\rangle_i=\sum_{\tau\in\Phi}2\operatorname{Re}\big(\tau(\eta)\,h_\tau(\tau\varphi_iu,\tau\varphi_iu')\big)
 =-\sum_{\tau\in\Phi}2\operatorname{Im}\tau(\eta)\operatorname{Im}h_\tau(\tau\varphi_iu,\tau\varphi_iu').\tag{$*$}
\]
For $z\in\mathbf C^g$ and $\tau\in\Phi$ let $\zeta_\tau(z)\in\mathbf C^r$ have the entries $z_{(\bar\tau,q-b)}$ for $b\le q$
and $z_{(\tau,b-q-1)}$ for $b>q$; the map $z\mapsto(\zeta_\tau(z))_{\tau\in\Phi}$ is a $\mathbf C$-linear bijection of
$\mathbf C^g$ onto $\prod_{\tau\in\Phi}\mathbf C^r$. For a frame $x$ and $\tau\in\Phi$ let $\lambda_{x,\tau}$ be the
$\mathbf R$-linear map of $\mathbf C^r$ with $\lambda_{x,\tau}(u)_b=h_\tau(x_{\tau,b},u)$ for $b\le q$ and
$\lambda_{x,\tau}(u)_b=(x_\tau^{-1}u)_b$ for $b>q$. By the definition of $P^i_x$,
$\zeta_\tau(P^i_x(v))=\lambda_{x,\tau}(\tau(\varphi_iv))$. Let $D$ be the diagonal matrix with the entries $1$ for $b\le q$
and $-1$ for $b>q$. Hermitian forms are taken linear in the first variable.

(1) Let $x$ be a frame and $\tau\in\Phi$, let $U$ and $x'_b$ be as in Lemma~\ref{cv-volume-key-normalized-gram}, and put
$d_b=h_\tau(x'_b,x'_b)$. For $u\in\mathbf C^r$ let $c(u)=(x_\tau U)^{-1}u$, so that $u=\sum_bc_b(u)x'_b$. Then
$h_\tau(u,u')=\sum_bd_bc_b(u)\overline{c_b(u')}$ and $h_\tau(x'_b,u)=d_b\overline{c_b(u)}$. As
$x'_b=\sum_{c\le b}U_{cb}x_{\tau,c}$ and $U^{-1}$ is upper triangular, the matrix $S$ with
$(Sw)_b=d_b^{-1}\sum_{c\le b}U_{cb}w_c$ for $b\le q$ and $(Sw)_b=\sum_{c\ge b}(U^{-1})_{bc}w_c$ for $b>q$ satisfies
\[
 (S\lambda_{x,\tau}(u))_b=\overline{c_b(u)}\quad(b\le q),\qquad(S\lambda_{x,\tau}(u))_b=c_b(u)\quad(b>q).
\]
If $\lambda_{x,\tau}(u)=0$, then $c(u)=0$ and $u=0$; so the $\mathbf R$-linear map $\lambda_{x,\tau}$ is bijective, and $S$
is invertible. Put $H_\tau(w,w')=\sum_b|d_b|(Sw)_b\overline{(Sw')_b}$, a positive definite Hermitian form on
$\mathbf C^r$. As $|d_b|=d_b$ for $b\le q$ and $|d_b|=-d_b$ for $b>q$,
\[
 \operatorname{Im}H_\tau(\lambda_{x,\tau}u,\lambda_{x,\tau}u')=\sum_{b\le q}d_b\operatorname{Im}\big(\overline{c_b(u)}\,c_b(u')\big)
 -\sum_{b>q}d_b\operatorname{Im}\big(c_b(u)\overline{c_b(u')}\big)=-\operatorname{Im}h_\tau(u,u').
\]
Let $H(z,z')=\sum_{\tau\in\Phi}2\operatorname{Im}\tau(\eta)\,H_\tau(\zeta_\tau(z),\zeta_\tau(z'))$, a positive definite
Hermitian form on $\mathbf C^g$. By $(*)$, $\operatorname{Im}H(P^i_xv,P^i_xv')=\langle v,v'\rangle_i$ for $v,v'\in L_i$. For
$b\in\mathcal O_E$, $\tau(\varphi_i(bv))=\tau(b)\tau(\varphi_iv)$, and the entries of $\lambda_{x,\tau}$ of index $b>q$ are
linear and those of index $b\le q$ conjugate-linear; so $P^i_x(bv)_k=\sigma_k(b)P^i_x(v)_k$ at every coordinate $k$.
Finally $P^i_x$ is the composite of $\varphi_i$, the map $u\mapsto(\tau(u))_{\tau\in\Phi}$ from $E^r$ to
$\prod_{\tau\in\Phi}\mathbf C^r$, the bijection $(\lambda_{x,\tau})_\tau$ and the inverse of $(\zeta_\tau)_\tau$. As $\Phi$ is
a CM type, the second map identifies $E^r\otimes_{\mathbf Q}\mathbf R$ with $\prod_{\tau\in\Phi}\mathbf C^r$ and maps the
lattice $\Lambda_i$ of $E^r$ onto a lattice. Hence $P^i_x$ is injective and its image is a lattice, and $P^i_x$ is a
period of this type.

(2) The entries of $x_\tau^{-1}$ are holomorphic functions of $x_\tau\in GL_r(\mathbf C)$, and $h_\tau(x_{\tau,b},u)$ is
linear in $x_\tau$. So every coordinate of $P^i_x(v)$ is holomorphic in $x$.

(3) Let $P$ be a period of $L_i$ of this type, with a positive definite Hermitian form $H$ such that
$\operatorname{Im}H(Pv,Pv')=\langle v,v'\rangle_i$. For $w\in\mathbf C^r$ let $\kappa(w)$ have the entries $\overline{w_b}$
for $b\le q$ and $w_b$ for $b>q$.

Step 1: the coordinates. Let $k$ be a coordinate. The map $v\mapsto P(v)_k$ is additive with
$P(bv)_k=\sigma_k(b)P(v)_k$ for $b\in\mathcal O_E$. As $\varphi_i$ is injective and $\Lambda_i$ spans $E^r$ over $E$, it is
the composite of $\varphi_i$ with a map on $E^r$ that is additive and is multiplied by $\sigma_k(b)$ when its argument is
multiplied by $b\in E$; so there is $c_k\in\mathbf C^r$ with $P(v)_k=\sum_jc_{k,j}\sigma_k(\varphi_i(v)_j)$ for all $v$.
Let $\tau\in\Phi$. The entries of $\zeta_\tau(z)$ of index $b>q$ are the coordinates $(\tau,a)$, of the embedding $\tau$,
and those of index $b\le q$ are the coordinates $(\bar\tau,a)$, of the embedding $\bar\tau$, as in (1). So there is a
matrix $M_\tau\in M_r(\mathbf C)$ with $\kappa(\zeta_\tau(Pv))=M_\tau\,\tau(\varphi_iv)$ for all $v$.

Step 2: $M_\tau$ is invertible. Let $w$ be a row vector with $wM_\tau=0$. The $\mathbf R$-linear form
$z\mapsto w\,\kappa(\zeta_\tau(z))$ vanishes on $P(L_i)$, which spans $\mathbf C^g$ over $\mathbf R$; so it vanishes on
$\mathbf C^g$, and as $\kappa\circ\zeta_\tau$ maps onto $\mathbf C^r$, $w=0$.

Step 3: the form. Put $A_\tau=(M_\tau^{-1})^{\mathsf T}\tau(G)\,\overline{M_\tau^{-1}}$, a Hermitian matrix with
$h_\tau(u,u')=(M_\tau u)^{\mathsf T}A_\tau\overline{M_\tau u'}$. By $(*)$ the identity
\[
 \operatorname{Im}H(z,z')=-\sum_{\tau\in\Phi}2\operatorname{Im}\tau(\eta)\,
 \operatorname{Im}\big(\kappa\zeta_\tau(z)^{\mathsf T}A_\tau\overline{\kappa\zeta_\tau(z')}\big)\tag{$**$}
\]
holds for $z=Pv$ and $z'=Pv'$. Both sides are $\mathbf R$-bilinear in $(z,z')$, and $P(L_i)$ spans $\mathbf C^g$ over
$\mathbf R$, so $(**)$ holds for all $z,z'\in\mathbf C^g$.

Step 4: the signs. Fix $\tau\in\Phi$, put $s=\operatorname{Im}\tau(\eta)$ and $A=A_\tau$, and for $w\in\mathbf C^r$ let
$\omega(w)\in\mathbf C^g$ be the vector with $\kappa\zeta_\tau(\omega(w))=w$ and $\zeta_{\tau'}(\omega(w))=0$ for
$\tau'\ne\tau$. Then $i\,\omega(w)=\omega(-iDw)$, and by $(**)$,
$\operatorname{Im}H(\omega(w),\omega(w'))=-2s\operatorname{Im}(w^{\mathsf T}A\overline{w'})$. As $H(iz,iz')=H(z,z')$, the
imaginary parts of $w^{\mathsf T}A\overline{w'}$ and $(Dw)^{\mathsf T}A\overline{Dw'}$ agree for all $w,w'$, and with $iw$ in
place of $w$ so do the real parts. So $DAD=A$, and $A$ commutes with $D$. As $\operatorname{Re}H(z,z)=\operatorname{Im}H(iz,z)$,
for $u\ne0$
\[
 0<\operatorname{Re}H(\omega(\bar u),\omega(\bar u))=\operatorname{Im}H(\omega(-iD\bar u),\omega(\bar u))
 =-2s\operatorname{Im}\big(-i(D\bar u)^{\mathsf T}Au\big)=2s\,\bar u^{\mathsf T}DAu.
\]
So the Hermitian matrix $DA$ is positive definite.

Step 5: the frame. Let $N$ be the matrix whose rows of index $b\le q$ are those of $A$ and whose other rows are those of
the identity. If $Nu=0$, then $(Au)_b=0$ for $b\le q$ and $u_b=0$ for $b>q$, so $\bar u^{\mathsf T}DAu=0$ and $u=0$; so $N$
is invertible. As $A$ commutes with $D$, so do $N$ and $Y=N^{-1}$. Put $x_\tau=M_\tau^{-1}\bar Y$. Then
$x_\tau^{\mathsf T}\tau(G)\bar x_\tau=Y^*AY$, and $D\,x_\tau^{\mathsf T}\tau(G)\bar x_\tau=Y^*(DA)Y$ is positive definite. Its
leading principal submatrices are positive definite, so the leading principal $j\times j$ minor of the Gram matrix
$x_\tau^{\mathsf T}\tau(G)\bar x_\tau=(h_\tau(x_{\tau,a},x_{\tau,b}))_{a,b}$ has the sign of the product of the first $j$
diagonal entries of $D$, which is $(-1)^{\max(0,j-q)}$. So $x=(x_\tau)_{\tau\in\Phi}$ is a frame.

Step 6: the period. Let $u\in\mathbf C^r$. For $b\le q$,
$\overline{h_\tau(x_{\tau,b},u)}=(u^{\mathsf T}\tau(G)\bar x_\tau)_b=\big((M_\tau u)^{\mathsf T}AY\big)_b$, and the entry
$(AY)_{cb}$ is $1$ for $c=b$ and $0$ for the other $c\le q$, as these rows of $A$ are those of $N$, and $0$ for $c>q$, as
$AY$ commutes with $D$. So $\overline{h_\tau(x_{\tau,b},u)}=(M_\tau u)_b$. For $b>q$ the row $b$ of
$x_\tau^{-1}=\bar NM_\tau$ is the row $b$ of $M_\tau$, so $(x_\tau^{-1}u)_b=(M_\tau u)_b$. With $u=\tau(\varphi_iv)$ and Step 1
this says that $\kappa\zeta_\tau(P^i_xv)=\kappa\zeta_\tau(Pv)$ for every $\tau\in\Phi$ and every $v$, so $P^i_x=P$.
\end{proof}
% Lean: exists_framePeriods (the period Pf = perFr from cmType_letter_frame_column_dictionary, whose index map idx
% and column formula hper are the text's zeta_tau and lambda_{x,tau}) through periods_eq_frame_periods. (1) is
% FramePeriod.isPeriod_frame: exists_loc for H_tau (with S, the local period locPer = lambda_{x,tau} and the
% coefficients c = U^-1 M^-1 u), allPer_bijective, CMLattice.exists_basis_of_cmType for the lattice, trace_re for (*),
% and hermForm_block for H. (2) is per_move. (3) is exists_frame_of_isPeriod: Step 1 is exists_row
% (ExistsSmulMemRange.exists_coeffs) with blockOf = kappa o zeta_tau, Step 2 is isUnit_block, Step 3 is im_form_eq,
% Step 4 is block_sign (with ofBlock = omega and I_smul_ofBlock), Step 5 is exists_frame (N, Y, x = Mi conj(Y)) with
% lminor_sign, and Step 6 is the end of exists_frames with hermC_col.

\begin{lemma}\label{cv-volume-key-normalized-riemann}
In Situation~\ref{cv-volume-key-rebuilt}, let $i$ be an index and $(e_1,\dots,e_g,f_1,\dots,f_g)$ a basis of $L_i$ with
$\langle e_k,e_l\rangle_i=\langle f_k,f_l\rangle_i=0$ and $\langle e_k,f_l\rangle_i=-\delta_{kl}$ for all $k,l$. For a
period $P$ of $L_i$ of the type $(\sigma_k)$ of the coordinates, let $A_P$ and $C_P$ be the matrices with the columns
$P(e_1),\dots,P(e_g)$ and $P(f_1),\dots,P(f_g)$, and let $\varepsilon_k$ be the $k$-th unit vector of $\mathbf C^g$.
\begin{enumerate}
\item $A_P$ is invertible, $h=A_P^{-1}$ and $Z=A_P^{-1}C_P$ are the only matrices with $hP(e_k)=\varepsilon_k$ and
  $hP(f_k)=Z\varepsilon_k$ for all $k$, and $Z$ is symmetric with positive definite imaginary part.
\item If $P_b$, $b\in V$, is a family of such periods on an open set $V\subset\mathbf C^m$ with $b\mapsto P_b(v)$
  holomorphic for every $v\in L_i$, then the matrices $h_b$ and $Z_b$ of (1) for $P_b$ are holomorphic in $b$.
\end{enumerate}
\end{lemma}
\leanref{Development/AbelianVarietyTheory/Proofs/UnitaryExistence/DifferentiableOnLeftInverse.lean}{GlobalTheory.RiemannRel.riemann_relations}
\leanref{Development/AbelianVarietyTheory/Proofs/UnitaryExistence/DifferentiableOnLeftInverse.lean}{GlobalTheory.RiemannRel.det_ne_zero}
\leanref{Development/AbelianVarietyTheory/Proofs/UnitaryExistence/DifferentiableOnLeftInverse.lean}{GlobalTheory.RiemannRel.symm_posDef}
\begin{proof}
(1) Write $A=A_P$ and $C=C_P$, and let $H$ be a positive definite Hermitian form of $P$, linear in the first variable,
with $\operatorname{Im}H(Pv,Pv')=\langle v,v'\rangle_i$. The $(j,k)$ entries of $A^{\mathsf T}H\bar A$, $A^{\mathsf T}H\bar C$ and
$C^{\mathsf T}H\bar C$ are $H(Pe_j,Pe_k)$, $H(Pe_j,Pf_k)$ and $H(Pf_j,Pf_k)$, with the imaginary parts $0$, $-\delta_{jk}$
and $0$.

$A$ is invertible. Let $Av=0$, and write $v=x+iy$ with $x,y\in\mathbf R^g$. For $r,s\in\mathbf R^g$,
$\operatorname{Im}H(Ar,As)=r^{\mathsf T}\operatorname{Im}(A^{\mathsf T}H\bar A)s=0$. As $Ax=-iAy$, $H(Ax,Ax)=iH(Ax,Ay)$, whose real
part is $-\operatorname{Im}H(Ax,Ay)=0$; so $Ax=0$, and then $Ay=0$. For $r\in\mathbf R^g$,
$\operatorname{Im}H(Ar,C\varepsilon_k)=-r_k$; with $r=x$ and $r=y$ this gives $x=y=0$.

Uniqueness. If $hP(e_k)=\varepsilon_k$ for all $k$, then $hA=1$ and $h=A^{-1}$, and $Z=hC$.

$Z$. Let $S=A^{\mathsf T}H\bar A$. It is Hermitian with zero imaginary part, so it is real and symmetric, and it is positive
definite, as $r^{\mathsf T}Sr=H(Ar,Ar)>0$ for $0\ne r\in\mathbf R^g$. Write $Z=X+iY$ with $X$ and $Y$ real. As $C=AZ$,
$A^{\mathsf T}H\bar C=S\bar Z$, whose imaginary part $-SY$ is $-1$; so $Y=S^{-1}$ is symmetric and positive definite. And
$C^{\mathsf T}H\bar C=Z^{\mathsf T}S\bar Z$ has the imaginary part $Y^{\mathsf T}SX-X^{\mathsf T}SY=X-X^{\mathsf T}$, as
$SY=1=Y^{\mathsf T}S$; so $X$ is symmetric, and so is $Z$.

(2) The entries of $h_b=A_{P_b}^{-1}$ are the cofactors of $A_{P_b}$ divided by its determinant, which does not vanish;
both are polynomials in the entries $P_b(e_k)$, which are holomorphic in $b$. So $h_b$ is holomorphic, and so is
$Z_b=h_bC_{P_b}$.
\end{proof}
% Lean: riemann_relations for a period P of D.La with a basis bas such that EF (bas a) (bas b) = -omegaStd a b: hM,
% hN and hK are the three imaginary parts, det_ne_zero (hreal, hcross, x and y), symm_posDef (S, hSY, hYinv, hKim,
% hXs) and hnorm with cols_iff for the uniqueness; (2) is differentiableOn_left_inverse (the adjugate,
% differentiableOn_det).

\begin{lemma}\label{cv-volume-key-normalized}
In Situation~\ref{cv-volume-key-rebuilt}, the following hold for every $i$.
\begin{enumerate}
\item There are theta data for $i$.
\item For the theta data fixed for $i$, there is a normalized framed point $u_i(x)$ at every frame $x$.
\end{enumerate}
\end{lemma}
\leanref{Development/AbelianVarietyTheory/Proofs/UnitaryExistence/NonemptyThetaDataFlagDomainModelsStatement.lean}{GlobalTheory.T.nonempty_thetaData_flagDomainModels}
\leanref{Development/AbelianVarietyTheory/Proofs/UnitaryExistence/ExistsNormalisedFramedPointOnEachComponentStatement.lean}{GlobalTheory.T.exists_normalisedFramedPoint_on_each_component}
\begin{proof}
(1) Fix $i$. Lemma~\ref{cv-uniformization-mod-eight}, applied to $L_i$, gives a basis
$(e_1,\dots,e_g,f_1,\dots,f_g)$ of $L_i$ as in the theta data; its $g=[E:\mathbf Q](p+q)/2$ is the number of
coordinates. The triple $(G,-\xi_0,\varphi)$ is a Hermitian model of $L$ (Definition~\ref{cv-uniformization-model},
Lemma~\ref{cv-volume-key-xi}(2)) with $\operatorname{Im}\tau(-\xi_0)>0$ for $\tau\in\Phi$. So by
Lemma~\ref{cv-uniformization-signature} every $\tau(G)$ with $\tau\in\Phi$ has $q$ positive and $p$ negative
eigenvalues, and by Lemma~\ref{cv-uniformization-frames-connected}, for this $G$, the frames form a nonempty connected
open subset $\mathcal X$ of the complex vector space $\prod_{\tau\in\Phi}M_r(\mathbf C)$. By
Lemma~\ref{cv-volume-key-normalized-periods}(1) and (2), $P^i_x$, $x\in\mathcal X$, is a holomorphic family of periods
of $L_i$ of the type of the coordinates. By Lemma~\ref{cv-volume-key-normalized-riemann}, for the basis $(e,f)$, every
frame $x$ has matrices $h_x$ and $Z_x$ with $h_xP^i_x(e_k)$ the $k$-th unit vector and $h_xP^i_x(f_k)$ the $k$-th column
of $Z_x$; moreover $\det h_x\ne0$, $Z_x$ is symmetric with positive definite imaginary part, and $h_x$ and $Z_x$ are
holomorphic in $x$. Theorem~\ref{cv-uniformization-family}(4), for $V=\mathcal X$, the family $P^i_x$, the basis $(e,f)$
with the sign $-1$ and these $h_x$ and $Z_x$, gives a characteristic $c_i\in\{0,\frac12\}^{2g}$ and a factor $F^i$
such that for every frame $x$ the line bundle of the polarization on $A_{y_i(P^i_x)}$ has a section whose value at
$p_x(z)$ in the trivialization over $U_j$ is $\vartheta[c_i](Z_x,h_xz)F^i_j(x,z)$. These are theta data for $i$.
% Lean: GlobalTheory.FrameUnif.nonempty_thetaData: the periods Pf, their holomorphy and IsPeriod from
% GlobalTheory.T.exists_framePeriods_alt (Lemma cv-volume-key-normalized-periods, with
% cmType_letter_frame_column_dictionary), the basis from exists_symplectic_basis_adapted_mod_eight (with
% PeriodExists.typeFn_period_exist and Relab.relab_good for g), frames_isConnected for the data M.X k, whose signature
% field the models take from signature_hermitian_model_of_datum, the linear coordinates co of the frames, and
% FrameTheta.exists_thetaFrame_frames: RiemannRel.riemann_relations for h and Z (Lemma
% cv-volume-key-normalized-riemann), isOpen_frames, and the field theta of FamilyUnif with sgn = -1.

(2) Let $x$ be a frame, put $P=P^i_x$ and $y=y_i(P)$, and let $p=p_x$ and $v_{\tau,a}$ be as in
Lemma~\ref{cv-volume-key-normalized-tangent} for $P$, which is a period of $L_i$ of the type of the coordinates by
Lemma~\ref{cv-volume-key-normalized-periods}(1). By Theorem~\ref{cv-uniformization-family}, $y$ is a complex point of $X$
over $\sigma$, and by Lemma~\ref{cv-volume-key-normalized-tangent}(3) the $v_{\tau,a}$, $0\le a<m_\tau$, form a basis of
$\operatorname{Lie}_\tau(y)$ for every $\tau$. The number $N_\xi\det h_x$ is not zero, as $N_\xi>0$ for $\xi_0\ne0$ and
$\det h_x\ne0$; let $\nu$ be a square root of it, so that $\nu\ne0$. By the theta data the line bundle of the
polarization on $A_y$ has a section $s$ whose value at $p(z)$ in the trivialization over $U_j$ is
$\vartheta[c_i](Z_x,h_xz)F^i_j(x,z)$; put $w=\nu s$. By Lemma~\ref{cv-volume-key-normalized-theta} there is $z_1$ with
$\vartheta[c_i](Z_x,z_1)\ne0$. Put $z=h_x^{-1}z_1$, and let $j$ be an index with $p(z)\in U_j$. As $F^i_j$ does not
vanish, $w$ is not zero at $p(z)$. The scheme $\operatorname{Spec}\mathbf C$ has one point, so by
Lemma~\ref{cv-volume-key-generator}(3), for the point $y$ over $\mathbf C$, $w$ generates the module of sections of the
line bundle of the polarization on $A_y$. So $(y,v,w)$, with $v$ the basis of the $v_{\tau,a}$, is a complex framed
point of $X$ over $\sigma$, and it is a normalized framed point at $x$.
\end{proof}
% Lean: (2) is GlobalTheory.FrameUnif.exists_isNormUn: the square root nu (IsAlgClosed.exists_pow_nat_eq), not zero by
% normConst_pos_of_ne_zero and T.norm, and GlobalTheory.T.exists_isThetaPt' = ThetaPtEx.exists_thetaPt: the point
% F.u k P over sigma (F.over_sigma with T.isPer), the vectors from exists_tangent_eigenbasis_exponential (Lemma
% cv-volume-key-normalized-tangent), the section from T.fr_sec, the nonvanishing of theta from
% thetaChar_half_differentiable_ne_zero_families (Lemma cv-volume-key-normalized-theta), and exists_gen:
% exists_fibreSection_other_model_same_values (the section moved to the other fibre product Gx y), the multiple
% scalSec y nu, and sec_values_complex_fibre, whose SecGen criterion (a section nonzero at a complex point generates)
% comes from gfp_gen under SecTwistSt and the reducedness of the smooth fibre; the text takes the criterion from Lemma
% cv-volume-key-generator(3), which rests on Theorem cv-volume-key-sections.

\begin{lemma}\label{cv-volume-key-cover}
In Situation~\ref{cv-volume-key-rebuilt}, let $y'=(y,v',w')$ be a complex framed point of $X$ over $\sigma$.
\begin{enumerate}
\item There are $i$ and a frame $x$ with $y=y_i(P^i_x)$ such that, if $v_{\tau,a}$ denotes the image in
  $\operatorname{Lie}A_y$ of the unit vector of $\mathbf C^g$ at the coordinate $(\tau,a)$, every vector $v'_{\tau,a}$
  of $y'$ is a combination $\sum_{b\le a}c_bv_{\tau,b}$ with $c_a\ne0$.
\item For such $i$ and $x$, if $u$ is a normalized framed point at $x$ for $i$, then $y'=u\cdot t$ for some
  $t\in(\mathbf C^\times)^I$.
\end{enumerate}
\end{lemma}
\leanref{Development/AbelianVarietyTheory/Proofs/UnitaryExistence/ExistsInverseBlock.lean}{GlobalTheory.CoverFrame.coverTri}
\leanref{Development/AbelianVarietyTheory/Proofs/UnitaryExistence/FamilyUniformisationIsComplexUniformisation.lean}{GlobalTheory.FrameUnif.cover_unfOf}
\begin{proof}
(1) By Theorem~\ref{cv-uniformization-family}(1) there are an index $i$ and a period $P$ of $L_i$ of the type of the
coordinates with $y=y_i(P)$. Let $p$, $dp$ and $\varepsilon_{\tau,a}$ be as in Lemma~\ref{cv-volume-key-normalized-tangent}
for $P$.

The adapted period. By Lemma~\ref{cv-volume-key-normalized-tangent}(1) each vector $v'_{\tau,a}$ of $y'$ is
$dp(w_{\tau,a})$ for one $w_{\tau,a}\in\mathbf C^g$, and by part (2) of that lemma $w_{\tau,a}$ vanishes outside the
coordinates $(\tau,b)$. As the $v'_{\tau,a}$, $0\le a<m_\tau$, form a basis of $\operatorname{Lie}_\tau(y)$, the
$w_{\tau,a}$ form a basis of the vectors that vanish outside these coordinates. So the matrix $W$ with the columns
$w_{\tau,a}$ is invertible and block diagonal for the embeddings, and so is $V=W^{-1}$, which maps $w_{\tau,a}$ to
$\varepsilon_{\tau,a}$. As $V$ commutes with the action of $\mathcal O_E$ on $\mathbf C^g$, through $\sigma_k$ at the
coordinate $k$, $P'=V\circ P$ is a period of $L_i$ of the same type, with the Hermitian form $H(V^{-1}z,V^{-1}z')$ for a
form $H$ of $P$.

The same point and the same map. By Theorem~\ref{cv-uniformization-family}(2), with $h=V$ and the identity of $L_i$,
$y_i(P')=y$. Let $p'\colon\mathbf C^g\to A_y$ be the map of $P'$. The isomorphisms $\mathbf C^g/P(L_i)\to A_y$ and
$\mathbf C^g/P'(L_i)\to A_y$ of Theorem~\ref{cv-uniformization-family} are compatible with $\mathcal O_E$, the forms and
the level structures, so the inverse of the second composed with the first lifts to a $\mathbf C$-linear isomorphism
$h_1$ of $\mathbf C^g$ with $p'(h_1z)=p(z)$ and $h_1\circ P=P'\circ\gamma_1$, for an $\mathcal O_E$-linear automorphism
$\gamma_1$ of $L_i$ that preserves the form and is congruent to $1$ modulo $M$ on $\bar{\mathfrak M}L_i$. Put
$\delta=\gamma_1^{-1}$ and $\beta=Vh_1^{-1}$; then $P'(\delta a)=\beta(P'(a))$ for $a\in L_i$. As $\delta$ preserves the
form, which is $\operatorname{Im}H'$ on $P'(L_i)$ for the form $H'$ of $P'$, and $P'(L_i)$ spans $\mathbf C^g$ over
$\mathbf R$, $\beta$ preserves $\operatorname{Im}H'$, and so also $H'(z,z)=\operatorname{Im}H'(iz,z)$, as $\beta$ is
$\mathbf C$-linear. So the orbit of each $P'(a)$ under $\beta$ lies in the finite set of points of the lattice $P'(L_i)$ of bounded norm, and as $L_i$ is
finitely generated, $\delta$ has finite order. As $\delta$
preserves the form and is congruent to $1$ modulo $M$ on $\bar{\mathfrak M}L_i$, it is congruent to $1$ modulo $M$ on
$L_i$ by Lemma~\ref{cv-volume-key-models-half}; as $M\ge3$, $\delta=1$ \cite{Min87}. So
$h_1\circ P=P'=V\circ P$, and $h_1=V$, as $P(L_i)$ spans $\mathbf C^g$ over $\mathbf R$. Hence $p'(Vz)=p(z)$ for all $z$, and
for the derivatives $dp'$ of $p'$ at $0$, $dp'(\varepsilon_{\tau,a})=dp'(Vw_{\tau,a})=dp(w_{\tau,a})=v'_{\tau,a}$.

The frame. By Lemma~\ref{cv-volume-key-normalized-periods}(3), $P'=P^i_x$ for a frame $x$. So $y=y_i(P^i_x)$, and every
vector $v'_{\tau,a}$ of $y'$ is the image $dp'(\varepsilon_{\tau,a})$ of the unit vector of the coordinate $(\tau,a)$:
the combination of the statement with $c_a=1$ and $c_b=0$ for $b<a$.
% Lean: coverTri: F.onto for i and P; isDerivAt0_ex_isTanAt for the w as derivatives; isEigAt_iff_of_isDerivAt0, the
% basis field of y and DerLin.isCombAt_iff_isDerivAt0 for the spanning (hsp); exists_inverse_block for V;
% PeriodLin.exists_period_mulVec (PeriodConn.isPeriod_mulVec) for P'; FamRigid.ex_eq_of_level_group with gam = 1 for
% the point (the field fibre of FamilyUnif) and the map: familyUnif_exp_unique_alt (ExpUniqFamily.fibrewise, from
% ex_add, ex_ker, ex_lin, ex_weil, ex_eta and ex_etaA) for h1 and gamma1, finite_order for delta,
% LevelHalf.trivial_mod_of_trivial_on_half for the congruence on L_i, and the published Minkowski statement
% latMinkowski; ExpChain.isDerivAt0_ex_iff for dp'; exists_framePeriods_alt (Lemma cv-volume-key-normalized-periods(3))
% for the frame; and isCombAt_iff_isDerivAt0 with the coefficients Pi.single i 1.

(2) Let $u=(y_u,v,w)$. By the definition of a normalized framed point, $y_u=y_i(P^i_x)=y$ and the vectors of $u$ are
the $v_{\tau,a}$. By (1), each vector of $y'$ is a combination of the vectors of $u$ in its eigenspace with index at
most its own and a nonzero coefficient at its own index. The sections $w'$ and $w$ both generate the module of
sections of the line bundle of the polarization over $\mathbf C$, which is free of rank one, so $w'=t_0w$ with
$t_0\in\mathbf C^\times$. Hence $y'=u\cdot t$, where $t_{\tau,a}$ is the coefficient of $v'_{\tau,a}$ at $v_{\tau,a}$.
\end{proof}
% Lean: cover_unfOf, with FrameOK only through its first clauses (the period perFr, the point and the tangent
% vectors as derivatives, which are part of IsNormUn through IsThetaPt and of the theta data through per), and
% Frames.frBor_calculus (the second clause: frameChange_exists_unique_of_same_point for the generators of the
% sections).

\subsection{The arithmetic volume step}\label{subsection-cv-arithmetic-volume-step}

\begin{lemma}\label{cv-volume-twisting-key-section}
In the setting of Definition~\ref{cv-volume-twisting-section}, suppose that $p+q\ge1$, and let $\kappa$ be
the rule of Theorem~\ref{cv-definite-key-formula}.
\begin{enumerate}
\item There is an invertible rational section of weight $\kappa(0,-4)$ on $X$ whose value at every framed
  point $(y,v,\vartheta)$ over a field of characteristic $0$ is $\kappa(A_y,\vartheta,v)$. Let $K_X$ be such
  a section.
\item If the uniformization at some good embedding satisfies the compatibility conditions of the moduli
  problem and $F_c$ is algebraically closed in the function field of every connected component of $X$,
  then $X$ has a key section $\rho$ such that for every connected component $Z$ there is
  $a_Z\in F_c^\times$ with $1_Z\rho^4=a_Z1_ZK_X$.
\item If the uniformizations at the good embeddings satisfy the compatibility conditions of the moduli
  problem and $X$ has a key section, then for every good embedding $\sigma$ there is $C_\sigma>0$ with
  $|\log C_\sigma|\le2g\log2\pi$ and $|K_X|_\xi(x)=C_\sigma$ at every frame $x$ at $\sigma$.
\end{enumerate}
\end{lemma}
\leanref{Development/AbelianVarietyTheory/Proofs/DivisorProofs/ExistsUnitSectionThetaKeyRuleWeightFourStatement.lean}{GlobalTheory.T.exists_unit_section_thetaKeyRule_weightFour}
\leanref{Development/AbelianVarietyTheory/Proofs/DivisorProofs/ExistsKeySectionFourthRootPerComponentStatement.lean}{GlobalTheory.T.exists_keySection_fourthRoot_perComponent}
\leanref{Development/AbelianVarietyTheory/Proofs/DivisorProofs/ThetaKeySectionNormComplexPointsEqTwoPiPowStatement.lean}{GlobalTheory.T.thetaKeySection_norm_complexPoints_eq_two_pi_pow}
\begin{proof}
Notation. Let $I$ be the set of pairs $(\tau,i)$ with $0\le i<m_\tau$, together with one further index $0$. For
framed points $y$ and $y'$ over an $F_c$-algebra $S$ and $t\in(S^\times)^I$, write $y'=y\cdot t$ if $y'$ has the
point of $y$, its bases are obtained from those of $y$ by upper triangular matrices whose diagonal entries are the
$t_{\tau,i}$, and its generator is $t_0$ times that of $y$. For $e\in\mathbf Z^I$ put $t^e=\prod_it_i^{e_i}$; a
rational section $f$ of a weight $G$ satisfies $f(y\cdot t)=t^ef(y)$, where $e_{\tau,i}=G_{\tau,i}$ and $e_0$ is the
last entry of $G$, and $e$ is its \emph{exponent}. Let $\chi(t)=t^{e_1}$ for the exponent $e_1$ of $\kappa(0,-1)$,
that is $\chi(t)=t_0^{-2}\prod_{\tau,i}t_{\tau,i}$. For a framed point $y=(y,v,\vartheta)$ over a field of
characteristic $0$ write $\kappa(y)=\kappa(A_y,\vartheta,v)$. By Theorem~\ref{cv-definite-key-formula}(1), with
$a=t_0$ and a triangular matrix of determinant $\prod_{\tau,i}t_{\tau,i}$,
\[
 \kappa(y\cdot t)=\chi(t)^4\kappa(y).\tag{$*$}
\]
For each connected component $Z$ let $y_Z$ be its generic framed point, over the field $\Omega_Z$ (proof of
Lemma~\ref{cv-point-regular-section}), and put $v_Z=\kappa(y_Z)\in\Omega_Z^\times$. The proofs of (1) and (2) use
the following facts, which are proved after the list.
\begin{enumerate}
\item[(K1)] (Charts.) There is a family of \emph{charts} $C$, each consisting of a connected component $Z_C$, a
  smooth $F_c$-subalgebra $R_C\subset\Omega_{Z_C}$ with fraction field $\Omega_{Z_C}$, a framed point $y_C$ over $R_C$
  and $t_C\in(\Omega_{Z_C}^\times)^I$ such that the image of $y_C$ over $\Omega_{Z_C}$ is $y_{Z_C}\cdot t_C$. Every
  framed point over a field $L\supset F_c$ is $y'\cdot t$ for the image $y'$ of some $y_C$ under some homomorphism
  of $F_c$-algebras $R_C\to L$ and some $t\in(L^\times)^I$; for $L=\Omega_Z$ and $y_Z$, the chart has $Z_C=Z$.
\item[(K2)] (Gluing.) Let $e\in\mathbf Z^I$ and $b_Z\in\Omega_Z$ for every component $Z$, with
  $t_C^e\,b_{Z_C}\in R_C$ for every chart $C$. Then there is a rational section $f$ of exponent $e$ with
  $f(y_Z)=b_Z$ for every $Z$.
\item[(K3)] (Normalized points.) Let $\sigma$ be a good embedding at which the uniformization satisfies the
  compatibility conditions, and put $g=\sum_\tau m_\tau$. There are $N\in\mathbf C^\times$ and a set $\mathcal U$ of
  complex framed points over $\sigma$ with $\kappa(u)=\big((2\pi i)^g/N\big)^4$ for every $u\in\mathcal U$, such that
  every complex framed point over $\sigma$ is $u\cdot t$ for some $u\in\mathcal U$ and $t\in(\mathbf C^\times)^I$.
\item[(K3$'$)] (Continuous roots.) In (K3), choose for every complex framed point $y$ over $\sigma$ some
  $u_y\in\mathcal U$ and $t_y$ with $y=u_y\cdot t_y$, and put $\rho(y)=\chi(t_y)$. Then $b\mapsto\rho(y_b)$ is continuous
  on $B$ whenever $B$ is an open subset of some $\mathbf C^m$ and $y_b$ is the image of a framed point over an
  $F_c$-algebra $A$ under homomorphisms $\psi_b\colon A\to\mathbf C$ over $\sigma$ with $b\mapsto\psi_b(a)$ holomorphic
  on $B$ for every $a\in A$.
\end{enumerate}
Moreover, a rational section of a weight whose values at all the $y_Z$ vanish is zero; this is fact (ii) of the
proof of Lemma~\ref{cv-point-regular-section}.

The roots. In (K3$'$), put $c_\sigma=\big(N/(2\pi i)^g\big)^4$. For a complex framed point $y$ over $\sigma$, by
$(*)$ and (K3), $\kappa(y)=\kappa(u_y\cdot t_y)=\chi(t_y)^4\kappa(u_y)=\rho(y)^4\big((2\pi i)^g/N\big)^4$, so
$\rho(y)^4=c_\sigma\kappa(y)$.

The charts. Fact (K1) holds for the generic framed points $y_Z$ of Lemma~\ref{cv-volume-key-generic-points} and
the charts of Definition~\ref{cv-volume-key-chart}: the cover is Lemma~\ref{cv-volume-key-atlas}(1), and the last
clause is part (2) of that lemma, applied to a chart and a homomorphism $R_C\to\Omega_Z$ given by part (1) for
$y_Z$.
% Lean: (K1) is the structure FChart and GlobalTheory.Reg.Atlas (evLaws_genFr_atlas: genFrEx_of,
% flagCharts_cover_framed_fieldPoints with the base change FrBC from exists_frMap), with chart_comp for the last
% clause, as used in KeyRoot.exists_keySec_pow_four (hgen). The text proves it down to the base change law for
% sections of the polarization, which it cites (Theorem cv-volume-key-sections); the Lean derives that law from its
% hypothesis SecTwistSt (LocalFrame.secBaseChange_of_secTwist, about 22,900 lines).

The gluing. Fact (K2), for the generic framed points and the charts of (K1), is Lemma~\ref{cv-volume-key-gluing}.
% Lean: (K2) is GlobalTheory.T.exists_integer_exponent_section_of_chart_regular = SecGlueInt.exists_secZ_of_regular,
% proved in Lemmas cv-volume-key-chart-compat to cv-volume-key-gluing. The vanishing is
% GlobalTheory.Reg.sec_eq_zero_of_generic_values.

The normalized points. Let $\sigma$ be as in (K3), and use Situation~\ref{cv-volume-key-rebuilt} for $\sigma$, with
theta data for every $i$ (Lemma~\ref{cv-volume-key-normalized}(1)). Let $\mathcal U$ be the set of the normalized
framed points $u_i(x)$ at the frames $x$, for all $i$, given by part~(2) of that lemma, and put $N=N_\xi$, which is
positive as $\xi_0\ne0$. By Lemma~\ref{cv-volume-key-cover}, every complex framed point over $\sigma$ is $u\cdot t$ for
some $u\in\mathcal U$ and some $t$. Let $u=(y,v,w)$ be a normalized framed point at a frame $x$, with number $\nu$. The
computation of Lemma~\ref{cv-volume-twisting-key-theta}, from this form of $u$ and parts (1) and (4) of
Theorem~\ref{cv-definite-key-formula}, gives $\nu^8\kappa(u)=\big((2\pi i)^g\det h_x\big)^4$; as $\nu^2=N\det h_x$,
$\kappa(u)=\big((2\pi i)^g/N\big)^4$.
% Lean: (K3) is in KeyRoot.exists_keySec_pow_four: the family F of hU (the compatibility conditions IsUnf), xi_0
% from IsUnfEx.exists_xi and Gm, phi from exists_herm_coord_of_datum (Lemma cv-volume-key-xi), the models M from
% ModelsExist.exists_modelsOf (Lemma cv-volume-key-models), the theta data T k and the normalized points un k (Lemma
% cv-volume-key-normalized), and the rebuilt uniformization isUnf_built, of which (K3) uses the field hU0.2.2.1,
% cover_unfOf with CoverFrame.coverTri (Lemma cv-volume-key-cover). The other parts of FramesOK (the normalization at
% frames with unit minors, the Borel and Gamma rules, the holomorphy of readings) enter isUnf_built and (K3'), but
% not (K3). N = normConst xi_0, the same for all k and nonzero (hC); kappa(u) is
% theta_key_rule_value_at_theta_point, as in rule_eq_of_rootVal. The one open step left for (K3) is
% Step 4 of the proof of Lemma cv-volume-key-models, the local norm condition at the finite places (DetIdeal.det_ideal
% and local_norm_of_ideal_norm, about 740 lines); Lemmas cv-volume-key-normalized and cv-volume-key-cover are written
% out. The earlier count of 34,000 lines also counted isUnf_built and the parts of FramesOK that (K3) does not use.

The continuous roots. In (K3) the set $\mathcal U$ is the set of normalized framed points $u_i(x)$ of the
uniformization at $\sigma$ rebuilt from theta data, as in Situation~\ref{cv-uniformization-situation} and
Lemma~\ref{cv-uniformization-frames}(2) and (7); at every level it is constructed in
Situation~\ref{cv-volume-key-rebuilt} and Lemma~\ref{cv-volume-key-normalized}. For this set,
$\rho(y)=\chi(t_y)$ is a root value of $y$ (Definition~\ref{cv-volume-key-theta-point}(3)), which by
Lemma~\ref{cv-volume-key-root-values}(1) does not depend on the choice of $u_y$ and $t_y$, and a family
$(y_b)_{b\in B}$ as in (K3$'$) is a holomorphic family of framed points
(Definition~\ref{cv-volume-key-theta-point}(4)). So (K3$'$) is Lemma~\ref{cv-volume-key-root-continuity}, applied to
the function $b\mapsto\rho(y_b)$.
% Lean: rho is any function with KeyFourth.RootVal (unfOf e0 F M un) y (rho y) at every y (hex and the choice of
% rho in KeyRoot.exists_keySec_pow_four, through hU0.2.2.1); the continuity is
% GlobalTheory.T.continuousOn_keySection_complexValue, applied to the families fun p => z (s p), proved in Lemmas
% cv-volume-key-root-values to cv-volume-key-root-continuity. The Lean proves it for the uniformization family F of
% X at its own level, which the text has in Situation cv-uniformization-situation; the uniformization of X rebuilt
% at its own level is Situation cv-volume-key-rebuilt, used for (K3). The roots paragraph is rule_eq_of_rootVal (e1 is
% (*), e2 and e3 the value at the normalized point), with c = (N / (2 pi i)^g)^4.

The values on the charts. Let $C$ be a chart. The ring $R_C$ is smooth over $F_c$ and a subring of a field, so it is
a noetherian integrally closed domain of characteristic $0$. By Lemma~\ref{cv-volume-twisting-key-unit}, applied to
the data of $y_C$, there is $u_C\in R_C^\times$ such that $\kappa$ of the image of $y_C$ under every homomorphism
$\psi$ from $R_C$ to a field of characteristic $0$ is $\psi(u_C)$. For the inclusion in $\Omega_{Z_C}$ this is
$\kappa(y_{Z_C}\cdot t_C)=\chi(t_C)^4v_{Z_C}$ by $(*)$; so $u_C=\chi(t_C)^4v_{Z_C}$ in $\Omega_{Z_C}$.

(1) Let $e=4e_1$, the exponent of $\kappa(0,-4)$, so that $t^e=\chi(t)^4$. For every chart,
$t_C^e\,v_{Z_C}=u_C$ and $t_C^{-e}v_{Z_C}^{-1}=u_C^{-1}$ lie in $R_C$. By (K2) there are rational sections $K_X$ of
weight $\kappa(0,-4)$ and $K'$ of weight $-\kappa(0,-4)$ with $K_X(y_Z)=v_Z$ and $K'(y_Z)=v_Z^{-1}$ for every $Z$.
The section $K_XK'-1$, of weight $0$, vanishes at every $y_Z$, so $K_XK'=1$ and $K_X$ is invertible.

Let $y$ be a framed point over a field $F$ of characteristic $0$. By (K1), $y=y'\cdot t$ with $y'$ the image of
some $y_C$ under some $\psi\colon R_C\to F$. The image of $K_X(y_C)$ in $\Omega_{Z_C}$ is
$K_X(y_{Z_C}\cdot t_C)=\chi(t_C)^4v_{Z_C}=u_C$, and $R_C\subset\Omega_{Z_C}$, so $K_X(y_C)=u_C$. Hence
$K_X(y')=\psi(u_C)=\kappa(y')$, and by $(*)$, $K_X(y)=\chi(t)^4K_X(y')=\chi(t)^4\kappa(y')=\kappa(y)$.

(2) Let $\sigma$ be the good embedding of the hypothesis, and let $\rho$ be as in (K3$'$) and $c_\sigma$ as in the
roots paragraph. Let $C$ be a chart. Its complex points over $\sigma$ are the homomorphisms $\psi\colon R_C\to\mathbf C$
over $\sigma$; for each, let $y_\psi$ be the image of $y_C$, so that $\rho(y_\psi)^4=c_\sigma\kappa(y_\psi)=c_\sigma\psi(u_C)$.
As $F_c$ is algebraically closed in the function field $\Omega$ of $Z_C$ and $R_C$ is smooth,
$\mathbf C\otimes_{F_c}R_C$ is a domain. By (K3$'$) and Lemma~\ref{cv-volume-key-roots}, with $K=\sigma(F_c)$, $A=R_C$,
$u=u_C$, $n=4$, $c=c_\sigma$
and the function $\psi\mapsto\rho(y_\psi)$, $u_C=a_Cb_C^4$ with $a_C\in F_c^\times$ and $b_C\in R_C$.

For a component $Z$, take a chart $C$ with $Z_C=Z$ (K1), and put $a_Z=a_C^{-1}$ and $b_Z=b_C/\chi(t_C)$. Then
$v_Z=\chi(t_C)^{-4}u_C=a_Z^{-1}b_Z^4$, and $b_Z\ne0$. For every chart $C'$ with $Z_{C'}=Z$,
\[
 \big(\chi(t_{C'})b_Z\big)^4=a_Z\chi(t_{C'})^4v_Z=a_Zu_{C'},
\]
a unit of $R_{C'}$. As $R_{C'}$ is integrally closed, $\chi(t_{C'})b_Z$ and its inverse, whose fourth powers lie in
$R_{C'}$, lie in $R_{C'}$. By (K2), for the exponents $e_1$ and $-e_1$, there are rational sections $\rho$ of weight
$\kappa(0,-1)$ and $\rho'$ of weight $\kappa(0,1)$ with $\rho(y_Z)=b_Z$ and $\rho'(y_Z)=b_Z^{-1}$; as in (1),
$\rho\rho'=1$, so $\rho$ is a key section.

Finally let $Z$ be a component, with idempotent $1_Z$. The section $1_Z\rho^4-a_Z1_ZK_X$ has weight
$\kappa(0,-4)$. Its value at $y_Z$ is $b_Z^4-a_Zv_Z=0$, as $1_Z(y_Z)=1$; and at $y_{Z'}$ for $Z'\ne Z$ it is $0$, as
$1_Z$ vanishes off $Z$. So it is zero, and $1_Z\rho^4=a_Z1_ZK_X$ with $a_Z\in F_c^\times$.
% Lean: (1) is KeyFourth.exists_unit_secZ_keyRule: exists_chart_units gives u_C (FamUnit, with
% exists_universalTripleFibre_at_framedPoint and universal_triple_fibres_base_change_key_data for the key data of a
% framed point, and theta_key_rule_triangular_frame_change_fourth_power for (*)), the sections f and f' come from
% exists_integer_exponent_section_of_chart_regular, f f' = 1 from sec_eq_zero_of_generic_values, and the values from
% the atlas (k_1 to k_6). (2) is KeyRoot.exists_keySec_pow_four: hroot (exists_eq_smul_pow with
% GeomConn.isDomain_tensor), hgen, hbC (IsIntegrallyClosed.exists_algebraMap_eq_of_isIntegral_pow), the key section
% exists_keySec_of_generic, and the comparison by sec_eq_zero_of_generic_values with evC_eq_zero off Z.

(3) Let $\sigma$ be a good embedding and $x$ a frame of a component of the uniformization at $\sigma$, and put
$g=\sum_\tau m_\tau$. The weight $\kappa(0,-4)$ has all its entries equal to $4$ and last entry $-8$; so, by
Lemma~\ref{cv-volume-twisting-holomorphic}, the reading $Q$ of $K_X$ is a form of weight
$\lambda(\kappa(0,-4))$, and its pointwise norm is invariant under right multiplication by families of
invertible upper triangular matrices (Situation~\ref{cv-cusp-setting}). By the Gram--Schmidt process there is
such a family $b$ with $z=xb$ an orthonormal frame. Every minor $\delta_{\tau,i}(z)$ has absolute value $1$, so
the pointwise norm of $Q$ at $x$ is $|Q(z)|=|\kappa(A_y,\vartheta,v)|$, where $(y,v,\vartheta)$ is the framed
point at $z$.

The compatibility conditions of the moduli problem present $A_y$ as $\mathbf C^g/\Lambda$ with the vectors of
$v$ the unit vectors of $\mathbf C^g$, in some order, so that the matrix $W$ whose rows are the vectors of $v$
has $|\det W|=1$. The Hermitian form $H$ of the polarization is $2\xi$ times the identity. Indeed, the
imaginary part of $H$ on $\Lambda$ is the alternating form $\operatorname{Tr}_{E/\mathbf Q}(-\xi_0h)$ of the
lattice of the component, and at an orthonormal frame the period formula of the uniformization, as in
Lemma~\ref{cv-uniformization-frames}(1), shows that the diagonal matrix with entries
$2\operatorname{Im}\tau(-\xi_0)$, where $\tau\in\Phi$ is the embedding to which the coordinate belongs, has the
same imaginary part on $\Lambda$; so $H$ is this diagonal matrix. Its entries are positive, as $H$ is positive
definite, and $\tau(-\xi_0)$ is purely imaginary, as $\bar\xi_0=-\xi_0$, of absolute value $\xi$ by the
normalization; so every entry is $2\xi$. The mean of $|\vartheta(u)|^2e^{-\pi H(u,u)}$ over a fundamental domain
$F$ of $\Lambda$ is $1$, as in Lemma~\ref{cv-uniformization-frames}(3), since the minors of $z$ have absolute
value $1$. As the polarization is principal, $\operatorname{vol}F\cdot\det H=1$, so
$\operatorname{vol}F=(2\xi)^{-g}$. Theorem~\ref{cv-definite-key-formula}(5) gives
\[
 |\kappa(A_y,\vartheta,v)|\Big(\frac{(2\xi)^{-g}}{\pi^g}\Big)^2=(2\pi)^{2g},\qquad\text{that is,}\qquad
 |\kappa(A_y,\vartheta,v)|=(2\pi)^{2g}(2\pi\xi)^{2g}.
\]
The entries of $\kappa(0,-4)$ add up to $4g$, so
$|K_X|_\xi(x)=(2\pi\xi)^{-2g}(2\pi)^{2g}(2\pi\xi)^{2g}=(2\pi)^{2g}$. Take $C_\sigma=(2\pi)^{2g}$; then
$|\log C_\sigma|=2g\log2\pi$.
\end{proof}
% Lean: (3) is KeyArchNorm.homN_key_const with C = (2 pi)^(2g): isSec_modelWt_rdU for the reading of K (the
% weight wtOf 0 (-4), exZ_wtOf), pnorm_rmul with the orthonormalizing bOf (ortho_bOf, rmul_mem_dom), dmin_ortho,
% norm_rule_of_isOrtho (the uniformization unfAt and unitAt at an orthonormal frame, OrthoHerm.hm_eq_smul_one
% through im_diag_eq_trace and herm_eq_of_im, and norm_rule_of_coord with the ArchNorm field, part (5) of the key
% formula, and volume_mul_det_eq_one_of_unimodular) and lieFacZ_mul.

\begin{lemma}\label{cv-volume-twisting-key-valuation}
In the setting of Definition~\ref{cv-volume-twisting-section}, suppose that $p+q\ge1$ and that the
prime-to-$\ell$ part of $M$ is at least $3$ for every prime $\ell$. Let $K_X$ be an invertible rational
section of weight $\kappa(0,-4)$ on $X$ whose value at every framed point $(y,v,\vartheta)$ over a field of
characteristic $0$ is $\kappa(A_y,\vartheta,v)$, for the rule $\kappa$ of
Theorem~\ref{cv-definite-key-formula}. For an embedding $\tau$ of $E$ let $\mathfrak d_\tau$ be the
ideal of $\mathcal O_{F_c}$ generated by the image of the different of $E$ under the embedding of $E$ in
$F_c$ given by $\tau$. Then, for every finite place $w$ of $F_c$ and the valuation $v$ at every integral
point of $X$ above $w$, normalized to extend $\operatorname{ord}_w$,
\[
 0\le v(K_X)\le4\sum_\tau m_\tau\operatorname{ord}_w(\mathfrak d_\tau).
\]
\end{lemma}
\leanref{Development/AbelianVarietyTheory/Proofs/DivisorProofs/ThetaKeySectionValuationBoundsDifferentStatement.lean}{GlobalTheory.T.theta_key_section_valuation_bounds_different}
\begin{proof}
Let $y$ be an integral framed point over the fraction field of a complete discrete valuation ring $T$
above $w$, with ramification index $e$ over $w$, so that $v(K_X)=\operatorname{ord}_T(K_X(y))/e$. Its
abelian variety extends to an abelian scheme $\mathcal A$ over $T$, its generator $\vartheta$ to a
generator of the sections over $T$, and its frame vectors to elements of the Lie algebra of $\mathcal A$
over $T$. This Lie algebra is free over $T$; let $v'$ be a basis. By
Theorem~\ref{cv-definite-key-formula}(3), $\kappa(A_y,\vartheta,v')$ is a unit $u$ of $T$. Write the
frame vectors in the basis $v'$: $v=Uv'$ for a $g\times g$ matrix $U$ over $T$. By part (1) of the
theorem, with $a=1$, $K_X(y)=\kappa(A_y,\vartheta,v)=\det(U)^4u$.

In the basis $v'$, $\mathcal O_E$ acts on $T^g$ through a ring homomorphism $A$. For $\tau\colon E\to\mathbf C$
let $\tau'\colon\mathcal O_E\to T$ be $\tau$ followed by $F_c\to T$; these are $[E:\mathbf Q]$ distinct
homomorphisms. The rows of $U$ of type $\tau$ are eigenvectors of $A$ for $\tau'$, and, as $y$ is integral,
every eigenvector in $T^g$ for $\tau'$ is a combination of them with coefficients in $T$. Choose $\delta_\tau$
in the different of $E$ with $\operatorname{ord}_w(\tau(\delta_\tau))=\operatorname{ord}_w(\mathfrak d_\tau)$,
which exists because $\mathfrak d_\tau$ is generated by the $\tau(\delta)$. By
Lemma~\ref{cv-volume-twisting-colength}, $\det U$ divides $\prod_\tau\tau'(\delta_\tau)^{m_\tau}$ in $T$.
Hence $0\le\operatorname{ord}_T\det U\le e\sum_\tau m_\tau\operatorname{ord}_w(\mathfrak d_\tau)$, and
$v(K_X)=4\operatorname{ord}_T(\det U)/e$ gives the claim.
\end{proof}
% Lean: GlobalTheory.KeyIntegralVal.key_valuation_bounds: exists_tied_det (the Lie basis w' of the model from
% Chart.exists_isLieBasis, the matrices A and U from exists_lie_matrices, the unit u from key_unit, and the
% field Homog of the key rule), exists_mem_different_vLa, ColengthDifferent.colength_dvd and nval_bounds.

\begin{lemma}\label{cv-volume-twisting-key}
In the setting of Definition~\ref{cv-volume-twisting-section}, suppose that $p+q\ge1$, that the
uniformizations at the good embeddings satisfy the compatibility conditions of the moduli problem, that
$F_c$ is algebraically closed in the function field of every connected component of $X$, and that the
prime-to-$\ell$ part of $M$ is at least $3$ for every prime $\ell$. Let $d$ be the prime of
Definition~\ref{cv-constant-sections}, let $B$ be the set of primes dividing $2d|d_E|M'$, where $M'$ is the
least common multiple of $M$ and $N$, and suppose that for every connected component $Z$ and every finite
place $w$ of $F_c$ over no prime of $B$ there is an integral point of $Z$ above $w$ at which the valuation of
every rational section is an integer or $+\infty$. Then $X$ has a key section, and every connected
component $Z$ of $X$ has a twisting section $t$ of weight $\kappa(0,1)$ whose function $l$ takes integer
values at the places over no prime of $B$, with
\[
 \deg t\ge-g\Big(\frac12\log2\pi+\log\operatorname{rd}_E\Big),\qquad g=\sum_\tau m_\tau=\frac12[E:\mathbf Q](p+q).
\]
\end{lemma}
\leanref{Development/AbelianVarietyTheory/Proofs/DivisorProofs/OrdConjDifferentEqZero.lean}{GlobalTheory.KeyRootTw.exists_twSec_keyRoot}
\begin{proof}
Let $\kappa$ be the rule of Theorem~\ref{cv-definite-key-formula}, and let $K_X$ be a section given by
Lemma~\ref{cv-volume-twisting-key-section}(1). By part (2) of that lemma, $X$ has a key section $\rho$, and
$1_Z\rho^4=a_Z1_ZK_X$ with $a_Z\in F_c^\times$ for every component $Z$. By part (3), for every complex
embedding $\sigma$ of $F_c$ there is $C_\sigma>0$ with $|\log C_\sigma|\le2g\log2\pi$ and
$|K_X|_\xi=C_\sigma$ at every frame at $\sigma^g$. By Lemma~\ref{cv-volume-twisting-key-valuation},
$0\le v(K_X)\le4\beta(w)$ at every integral point above $w$, where
$\beta(w)=\sum_\tau m_\tau\operatorname{ord}_w(\mathfrak d_\tau)$.

Fix $Z$, put $a=a_Z$, and let $s=1_Z\rho^{-1}$ and $s'=1_Z\rho$, of weights $\kappa(0,1)$ and $\kappa(0,-1)$.
Then $1_Zs=s$, $1_Zs'=s'$ and $ss'=1_Z$.

Norms. Let $x$ be a frame of $Z$ at $\sigma^g$. The pointwise norm of $1_Z$ there is $1$, and
$|\sigma^g(a)|=|\sigma(a)|$, so $|a1_ZK_X|_\xi(x)=|\sigma(a)|C_\sigma$. As $1_Z\rho^4=(1_Z\rho)^4$ and the
pointwise norm is multiplicative, $|s'|_\xi(x)^4=|\sigma(a)|C_\sigma$, so
$|s'|_\xi(x)=(|\sigma(a)|C_\sigma)^{1/4}$; and $|s|_\xi(x)|s'|_\xi(x)=|1_Z|_\xi(x)=1$ gives
$|s|_\xi(x)=c_\sigma$ with $c_\sigma=(|\sigma(a)|C_\sigma)^{-1/4}$.

Valuations. At an integral point of $Z$ above $w$, $v(1_Z)=0$, $v(\rho)+v(\rho^{-1})=0$ and
$4v(\rho)=v_w(a)+v(K_X)$. Hence $v(s)=-v(\rho)$ lies between $l(w)=-\frac14v_w(a)-\beta(w)$ and
$h(w)=-\frac14v_w(a)$, which vanish at almost all $w$.

Multiplication. The weights $\kappa(0,1)$ and $\kappa(0,-1)$ have entries $-1$ and $1$ at every embedding,
constant in the index, and even last entries $2$ and $-2$; $|s|_\xi$ and $|s'|_\xi$ are constant on the
frames of $Z$; and $s$ and $s'$ are supported on $Z$, as multiples of $1_Z$. So
Lemma~\ref{cv-volume-twisting-cusp}, with the key section $\rho$, applies to $s$ and to
$s'$: $fs$ is a cusp form of weight $G'+\kappa(0,1)$ for every cusp form $f$ on $Z$ of weight $G'$, and
$fs'$ is a cusp form of weight $G'+\kappa(0,1)+\kappa(0,-1)=G'$ for every cusp form $f$ on $Z$ of weight
$G'+\kappa(0,1)$. By Lemma~\ref{cv-volume-twisting-points}, $Z$ has frames at every good embedding and
integral points above every finite place. So these data form a twisting section $t$ of weight
$\kappa(0,1)$ on $Z$.

Integrality. Let $w$ lie over no prime of $B$. Then $w$ lies over no prime dividing $d_E$; as the norm of
$\mathfrak d_\tau$ is a power of $|d_E|$, $\operatorname{ord}_w(\mathfrak d_\tau)=0$ for every $\tau$, so
$\beta(w)=0$ and $l(w)=h(w)$. At an integral point of $Z$ above $w$ at which all valuations are integers or
$+\infty$, $v(s)$ lies between $l(w)$ and $h(w)$, so it is the real number $l(w)$, which is therefore an
integer.

Degree. By the product formula, $\sum_wv_w(a)\log Nw=\sum_\sigma\log|\sigma(a)|$, so the terms in $a$
cancel and
\[
 \deg t=-\frac1{[F_c:\mathbf Q]}\sum_w\beta(w)\log Nw+\frac1{4[F_c:\mathbf Q]}\sum_\sigma\log C_\sigma.
\]
The norm of $\mathfrak d_\tau$ is $|d_E|^{[F_c:E]}$, so
$\frac1{[F_c:\mathbf Q]}\sum_w\operatorname{ord}_w(\mathfrak d_\tau)\log Nw=\log\operatorname{rd}_E$, and
the first term is $-g\log\operatorname{rd}_E$. The second is at least $-\frac14\cdot2g\log2\pi$. This gives
the bound.
\end{proof}
% Lean: exists_twSec_keyRoot. kappa is KeyFamClosing.keyFam_one; K is
% GlobalTheory.T.exists_unit_section_thetaKeyRule_weightFour, rho and a_Z are
% GlobalTheory.T.exists_keySection_fourthRoot_perComponent, the valuation bound
% GlobalTheory.T.theta_key_section_valuation_bounds_different and the norms
% GlobalTheory.T.thetaKeySection_norm_complexPoints_eq_two_pi_pow. The section is t with sig = idem Z * rho^-1
% and inv = idem Z * rho; its norms are homN_root, its valuations val_root, its multiplication properties
% GlobalTheory.T.mul_mem_cuspForms_of_boundedSection_component, its points FlatLine.hasPts, the integrality
% ord_conjDifferent_eq_zero, and the degree le_lineDeg_root (lineDeg_num, sum_ord), with g from
% card_frame_index.

\begin{lemma}\label{cv-volume-twisting-unramified}
In Situation~\ref{cv-cusp-setting}, let $(M',\mathfrak M')$ be a level (Definition~\ref{cv-deep-cover}) such
that the prime-to-$\ell$ part of $M'$ is at least $3$ for every prime $\ell$, let $X'$ be the PEL moduli
scheme of the PEL datum of $X$ at this level (Theorem~\ref{cv-deep-pel-facts}(1)), and let
$F'\subset\mathbf C$ be a number field that contains $F_c$ and the $M'$-th roots of unity. Let $Z$ be a
connected component of $X'\otimes F'$, and let $w$ be a finite place of $F'$ over a prime $\ell$ that does
not divide $2d|d_E|M'$. Then $Z$ has an integral point above $w$, over a complete discrete valuation ring
whose valuation restricts to that of $w$ on $F'$, at which the valuation of every rational section of
$X'\otimes F'$ is an integer or $+\infty$.
\end{lemma}
\leanref{Development/AutomorphicArithmetic/ThetaCorrespondence/ThetaPairLifts/HodgeBundleWeights/UnramifiedPointsAtDeepLevel.lean}{GlobalTheory.T.unramified_points_of_deep_cover}
\leanref{Development/AutomorphicArithmetic/ThetaCorrespondence/ThetaPairLifts/HodgeBundleWeights/UnramifiedPointsOfCMDatum.lean}{GlobalTheory.T.cm_points_unramified}
\leanref{Development/AutomorphicArithmetic/ThetaCorrespondence/ThetaPairLifts/HodgeBundleWeights/UnramifiedPointsOfCMDatum.lean}{GlobalTheory.UnramCM.unramified_points_of_CMDatum}
\begin{proof}
Let $\theta$ be the square root of $-d$ fixed by $\mathcal D$, and $K_0=\mathbf Q(\theta)$, of discriminant
$-d$.

The field. Let $R\subset\mathbf C$ be the ray class field of $K_0$ of modulus $M'\mathcal O_{K_0}$
(Lemma~\ref{cv-deep-ray-field}). It is unramified over $K_0$ at the primes not dividing $M'$
\cite[Chapter~VI]{Neu99}, and $K_0$ is unramified over $\mathbf Q$ at $\ell$, as $\ell\nmid d$; so $\ell$ is
unramified in $R$. By Lemma~\ref{cv-deep-ray-field}, $R$ contains the Weber values of level $M'$, so by
Lemma~\ref{cv-components-weber-curve} every automorphism of $\mathbf C$ that fixes $R$ fixes a CM curve with
its points of order $M'$ (Definition~\ref{cv-components-conjugation}). Let $K=F'R\subset\mathbf C$ be the
compositum. As $\ell$ is unramified in $R$, there is a prime $\mathfrak p$ of $K$ above $w$ at which $K$ is
unramified over $F'$, so that the valuation of $\mathfrak p$ restricts to that of $w$ on $F'$.

A point over $K$. Fix a type of signature $(p,q)$ for $\theta$, and let $L_i$, $y_i(P)$ and $Y_i$ be as in
Situation~\ref{cv-components-situation} for the level $(M',\mathfrak M')$. As in the first step of the proof
of Lemma~\ref{cv-volume-twisting-points}, $Z$ has a point over a number field containing $F'$, and an
embedding of this field in $\mathbf C$ over $F'$ gives a complex point $x$ of $Z$ over the inclusion
$F'\subset\mathbf C$. Its image in $X'$ is $y_i(P)$ for some $i$ and some period $P$ of $L_i$
(Theorem~\ref{cv-uniformization-family}(1)). Alp\"oge constructs from $\mathcal D$ a PEL datum of signature
$(1,0)$ over $E$, so by Lemma~\ref{cv-components-split} $L_i$ has a period $P'$ with $L_i=L_+(P')+L_-(P')$. Let
$x'$ be the complex point of $X'\otimes F'$ over the inclusion with image $y_i(P')$. As $Z$ is open and closed,
the function $e_Z$ that is $1$ on $Z$ and $0$ on its complement is regular on $X'\otimes F'$, and by
Lemma~\ref{cv-components-points}(3) it takes the same value at $x$ and at $x'$; so $x'$ lies on $Z$. Let
$\tau$ be an automorphism of $\mathbf C$ that fixes $K$. It fixes $F'$, hence $\theta$, and it fixes a CM
curve with its points of order $M'$; so $y_i(P')^\tau=y_i(P')$ by Lemma~\ref{cv-components-cm-fixed}, and
$x'^\tau=x'$, as $\tau$ fixes $F'$. So the residue field of the image of $x'$ embeds in $\mathbf C$ with an
image fixed by every automorphism of $\mathbf C$ over $K$, hence contained in $K$, and $x'$ comes from a
point of $Z$ over $K$. Bases of the eigenspaces of its Lie algebra and a generator of the sections of its
line bundle, which exist over a field as in the proof of Lemma~\ref{cv-volume-twisting-points}, make it a
framed point $y_K$ over $K$ on $Z$.

The integral point. Let $T$ be the completion of $\mathcal O_K$ at $\mathfrak p$; its valuation restricts to
that of $w$ on $F'$. The prime-to-$\ell$ part of $M'$ is at least $3$, so, as in the proof of
Lemma~\ref{cv-volume-twisting-points}, a model of $X'\otimes F'$ over the valuation ring of $F'$ at $w$ and
the valuative criterion of properness give an integral framed point $y$ over the fraction field of $T$ with
the point of $y_K$. It lies on $Z$, and it is algebraic over $F'$.

Valuations. Let $f$ be a rational section. The valuation of $f$ at $y$ is the valuation of $f(y)$ in the
fraction field of $T$, divided by the ramification index of $T$ over $w$, which is $1$. So it is an integer,
or $+\infty$ when $f(y)=0$.
\end{proof}
% Lean: unramified_points_of_deep_cover = UnramCover.unramified_points_at_deep_level applies cm_points_unramified
% with the datum of signature (1,0) of complexUniformizationLaws_of_goodDeepCover (hLaws.line) and 6 <= r.
% unramified_points_of_CMDatum: UnramPts.of_index_one (valuations nval T.V 1); the field
% GlobalTheory.T.unramified_curve_field_existence = exists_field_of_rayField (the ray class field of RayField,
% not_dvd_discr_of_unramIn, curveOver_of_weberIn); the prime compositum_intermediateField_exists_place_valuation_eq;
% the point GlobalTheory.T.exists_framedPoint_liesOn_component_cmCurve_field = FixedPoint.exists_frPt_liesOn
% (exists_fixed_point: exists_numberField_point, F.onto, exists_isSplit_period_anyLevel, exists_indicator with
% valAt_eq_of_finite_values_anyLevel, conj_point_eq_of_isSplit; then exists_descent and exists_framed); the
% integral point GlobalTheory.T.exists_integralFramedPoint_completion_index_one = CompletionPoint.exists_intFr_index_one
% (modelAtSt_of_good, exists_initial_completion, intFr_of_model).

\begin{lemma}\label{cv-volume-generators-key-valuation-level}
In Situation~\ref{cv-function-situation}, let $(M,\mathfrak M,N)$ be a level of the datum
(Definition~\ref{cv-model-toroidal-semiabelian}(1)). Let $c$ be the largest divisor of $M$ prime to $d$, and put
$m=4c^2+d$ and $\alpha=2c+\sqrt{-d}\in\mathcal O_E$. Then $(Mm,\mathfrak M\alpha,N)$ is a level, with $M\mid Mm$ and
$\mathfrak M\alpha\subset\mathfrak M$, and for every prime $\ell$ the prime-to-$\ell$ part of $Mm$ is at least $3$.
\end{lemma}
\leanref{Development/AbelianVarietyTheory/Proofs/DivisorProofs/ExistsGoodLevelMultipleSameAdelicPartStatement.lean}{GlobalTheory.T.exists_goodLevel_multiple_same_adelicPart}
\leanref{Development/AbelianVarietyTheory/Proofs/DivisorProofs/ExistsGoodLevelMultipleSameAdelicPart.lean}{GlobalTheory.GoodMultipleLevel.exists_good_multiple}
\begin{proof}
Write $w_0=(1+\sqrt{-d})/2$. By the proof of Lemma~\ref{cv-volume-key-xi}(1), $\overline{\sqrt{-d}}=-\sqrt{-d}$; and
$d\ge7$ is odd, as $d\equiv7\pmod{16}$.

The integers. $M$ and $m$ are coprime. Indeed, let $\ell$ be a prime that divides both. If $\ell\mid c$, then
$\ell\mid m-4c^2=d$, so $\ell=d$, which does not divide $c$. Otherwise $\ell$ divides $M/c$, a power of $d$, so $\ell=d$,
and then $d\mid4c^2$, which is impossible as $d$ is odd and prime to $c$. The integer $m$ is odd, and $m\ge d\ge7$. Also
$4c$ and $m$ are coprime: $m$ is odd, and a prime that divides $c$ and $m$ divides $d$, which is prime to $c$.

The element $\alpha$. Its conjugate is $\bar\alpha=2c-\sqrt{-d}$, so $\alpha\bar\alpha=4c^2+d=m$ and
$\alpha+\bar\alpha=4c$. Choose integers $u$ and $v$ with $4cu+mv=1$; then $\alpha(u+v\bar\alpha)+\bar\alpha u=1$, so
$\alpha\mathcal O_E+\bar\alpha\mathcal O_E=\mathcal O_E$. Also $\alpha=(2c-1)+2w_0\in\mathbf Z[w_0]$.

The level. The ideal $\mathfrak M\alpha$ is generated by the products of $\alpha$ with generators of $\mathfrak M$ in
$\mathbf Z[w_0]$, which lie in $\mathbf Z[w_0]$. Next,
$(\mathfrak M\alpha)\overline{(\mathfrak M\alpha)}=\mathfrak M\bar{\mathfrak M}\,\alpha\bar\alpha=Mm\mathcal O_E$. Each
of $\mathfrak M$ and $\alpha\mathcal O_E$ is coprime to each of $\bar{\mathfrak M}$ and $\bar\alpha\mathcal O_E$:
$\mathfrak M+\bar{\mathfrak M}=\mathcal O_E$ by hypothesis and $\alpha\mathcal O_E+\bar\alpha\mathcal O_E=\mathcal O_E$ by
the above, while $\mathfrak M$ and $\bar{\mathfrak M}$ contain $\mathfrak M\bar{\mathfrak M}=M\mathcal O_E$,
$\alpha\mathcal O_E$ and $\bar\alpha\mathcal O_E$ contain $m$, and $M$ and $m$ are coprime. So
$\mathfrak M\alpha+\overline{\mathfrak M\alpha}=\mathcal O_E$. As $Mm$ is odd and at least $3$,
$(Mm,\mathfrak M\alpha,N)$ is a level; and $M\mid Mm$ and $\mathfrak M\alpha\subset\mathfrak M$.

Every prime. Let $\ell$ be a prime. As $M$ and $m$ are coprime, $\ell$ does not divide one of them, $k$ say. Then $k$
divides the prime-to-$\ell$ part of $Mm$, which is therefore at least $k\ge3$.
\end{proof}
% Lean: exists_good_multiple, with c = l.M / d^(v_d l.M), a = 2c, N = a^2 + d and alpha = a + rt: coprime_M_N,
% LevelAllPrimes.add_rt_mul_sub_rt (alpha bar(alpha) = N), LevelAllPrimes.isCoprime_add_rt_sub_rt (from
% gcd(2a, a^2 + d) = 1), HostModel.rt_mem_adjoin, the coprimalities c1 to c4, and GlobalModelGuards.goodAt_of_coprime.
% The level l_2 has M = l.M * N, Mb = l.Mb * (alpha) and A = l.A. Size: about 120 lines.

\begin{lemma}\label{cv-volume-generators-key-valuation-lift}
In the setting of Definition~\ref{cv-volume-twisting-section}, suppose that $p+q\ge1$, and let $(M_2,\mathfrak M_2,N)$
be a level with $M\mid M_2$ and $\mathfrak M_2\subset\mathfrak M$ (Definition~\ref{cv-model-toroidal-semiabelian}(1)).
\begin{enumerate}
\item The PEL datum of $X$ has a moduli scheme $X_2$ at this level, and there are a morphism $\pi\colon X_2\to X$ and
  a homomorphism, over $\pi$, from the universal abelian scheme over $X_2$ to the one over $X$ that identifies the
  first with the pullback of the second, compatibly with the actions of $\mathcal O_E$ and the polarizations. For
  every number field $F'\subset\mathbf C$ that contains $F_c$ and the roots of unity of order
  $\operatorname{lcm}(M_2,N)$, the base change $X_2\otimes F'\to X\otimes F'$ of $\pi$ is finite, \'etale and
  surjective.
\item Let $F_2\subset\mathbf C$ be the field generated over $F_c$ by the roots of unity of order
  $\operatorname{lcm}(M_2,N)$. Let $T$ be a complete discrete valuation ring whose fraction field $K$ has
  characteristic $0$, with an embedding $F_c\to K$ under which the valuation of $T$ restricts to a positive multiple
  of that of a finite place $w$ of $F_c$, and let $y$ be a point of $X$ over $K$ over this embedding. Then there are
  a complete discrete valuation ring $T'$ with fraction field $K'$, an injective local homomorphism $T\to T'$ with
  $\operatorname{ord}_{T'}(a)=e'\operatorname{ord}_T(a)$ for an integer $e'\ge1$ and every $a\in K$, an embedding
  $F_2\to K'$ that extends $F_c\to K\to K'$ and under which the valuation of $T'$ restricts to a positive multiple of
  that of a finite place of $F_2$, and a point $y_2$ of $X_2$ over $K'$ over this embedding such that $\pi(y_2)$ is
  the image of $y$ in $X(K')$.
\end{enumerate}
\end{lemma}
\leanref{Development/AbelianVarietyTheory/Proofs/ModuliProofs/ModuliSchemeLevelMapsHeckePoints.lean}{GlobalTheory.moduliScheme_levelMaps_heckePoints}
\leanref{Development/AbelianVarietyTheory/Proofs/DivisorProofs/ModulipointLiftToFinerLevelCompleteDVRStatement.lean}{GlobalTheory.T.modulipoint_lift_to_finerLevel_completeDVR}
\leanref{Development/AbelianVarietyTheory/Proofs/DivisorProofs/ModulipointLiftToFinerLevelCompleteDVR.lean}{GlobalTheory.LevelLift.lift_to_finer_level}
\begin{proof}
(1) This is the representability of PEL moduli problems at a level of at least $3$, and the maps between the moduli
schemes at a level and at a finer one \cite{Kot92b,Lan13}. For levels of full level $1$ it is
Theorem~\ref{cv-deep-pel-facts}(1) and (2).

(2) The field $F_2$ is a finite extension of $F_c$, so the $K$-algebra $K\otimes_{F_c}F_2$ is finite and nonzero; its
quotient by a maximal ideal is a finite extension $L_0$ of $K$ with an embedding $F_2\to L_0$ that extends
$F_c\to K\to L_0$. With this embedding, the image of $y$ in $X(L_0)$ is a point $z$ of $X\otimes F_2$ over $L_0$. By
(1), the base change $\pi_2\colon X_2\otimes F_2\to X\otimes F_2$ of $\pi$ is finite and surjective, so the fibre of
$\pi_2$ over $z$ is the spectrum of a nonzero finite $L_0$-algebra. The residue field $L$ at one of its points is a
finite extension of $L_0$, and the point gives a point $y'$ of $X_2\otimes F_2$ over $L$ whose image under $\pi_2$ is
the image of $z$.

The field $L$ is a finite extension of $K$, so the integral closure $T'$ of $T$ in $L$ is a complete discrete
valuation ring with fraction field $K'=L$, and $\operatorname{ord}_{T'}(a)=e'\operatorname{ord}_T(a)$ for $a\in K$,
where $e'\ge1$ is the ramification index \cite[Chapter~II, \S2, Proposition~3]{Ser79}; the inclusion $T\to T'$ is an
injective local homomorphism. The residue field of $T$ has positive characteristic, as $T$ lies above $w$, and so has
that of $T'$. The ring of integers of $F_2$ is integral over $\mathbf Z$, so the embedding $F_2\to K'$ maps it into
$T'$, and the preimage of the maximal ideal of $T'$ is a nonzero prime ideal, as it contains the residue
characteristic. So the valuation of $T'$ restricts to a positive multiple of that of the finite place of $F_2$ given
by this prime. Let $y_2$ be the image of $y'$ in $X_2$, a point over $K'$ over the embedding $F_2\to K'$. Then
$\pi(y_2)$ is the image of $y$ in $X(K')$.
\end{proof}
% Lean: (1) is ((moduliScheme_levelMaps_heckePoints).2 E H Dw hpq).1, ToroidalLevelCompatibility for Level.Dvd l l_2:
% the level map r, the SABMap m and, for coefficient fields c and c_2 with the same field, rF finite, etale and surjective;
% the scheme X_2 is GlobalTheory.Moduli.moduli_scheme_exists (N1ab). (2) is lift_to_finer_level, with l_2.A = l.A:
% F_2 = c.F with zeta (lcm(M_2, A)) adjoined, exists_field_over for L_0, exists_point_finite_surjective for L and y',
% PlaceExtField.exists_cdvr_finite_extension for T' (GlobalTheory.SS.exists_complete_dvr_finite_extension), and
% LocalDepthCompare.exists_overLa_of_not_charZero (with not_charZero_of_overLa) for the place of F_2. Size: about 130
% lines, besides the cited level maps.

\begin{lemma}\label{cv-volume-generators-key-valuation}
In the setting of Definition~\ref{cv-volume-twisting-section}, suppose that $p+q\ge1$. Let $K_X$ be an
invertible rational section of weight $\kappa(0,-4)$ on $X$ whose value at every framed point $(y,v,\vartheta)$
over a field of characteristic $0$ is $\kappa(A_y,\vartheta,v)$, for the rule $\kappa$ of
Theorem~\ref{cv-definite-key-formula}. Then $v(K_X)\ge0$ for the valuation $v$ at every integral point of $X$
above every finite place of $F_c$.
\end{lemma}
\leanref{Development/AbelianVarietyTheory/Proofs/DivisorProofs/ThetaKeyRuleValuationNonnegIntegralFramedPointStatement.lean}{GlobalTheory.T.thetaKeyRule_valuation_nonneg_integral_framed_point}
\begin{proof}
Let $y$ be an integral framed point over the fraction field of a complete discrete valuation ring $T$ above a
finite place $w$, with ramification index $e$ over $w$, and let $p$ be the prime below $w$. Let $v'$, $U$ and
$u=\kappa(A_y,\vartheta,v')$ be as in the proof of Lemma~\ref{cv-volume-twisting-key-valuation}. That proof
gives $\det U\in T$ and $K_X(y)=\det(U)^4u$, and it uses the hypothesis on $M$ only to show that $u$ is a unit
of $T$. Once $u$ is a unit, $v(K_X)=4\operatorname{ord}_T(\det U)/e\ge0$.

If the prime-to-$p$ part of $M$ is at least $3$, $u$ is a unit by Theorem~\ref{cv-definite-key-formula}(3),
as in that proof: Alp\"oge takes the data over $T$ from a model of $X$ over the valuation ring of $F_c$ at
$w$, and this model exists under the condition at $p$ alone.

Otherwise, Alp\"oge passes to a finer level. There is a level $(M_2,\mathfrak M_2,N_2)$ with $M\mid M_2$,
$\mathfrak M_2\subset\mathfrak M$ and $N_2=N$ whose prime-to-$\ell$ part is at least $3$ for every prime $\ell$
(Lemma~\ref{cv-volume-generators-key-valuation-level}); let $X_2$ be the moduli scheme of the PEL datum of $X$ at
this level, with its level map to $X$ (Lemma~\ref{cv-volume-generators-key-valuation-lift}(1)). By part (2) of that
lemma there are a complete discrete valuation ring $T'$ and an injective local homomorphism $T\to T'$, with
$\operatorname{ord}_{T'}(a)=e'\operatorname{ord}_T(a)$ for some integer $e'\ge1$ and every $a$ in the fraction
field of $T$, and a point of $X_2$ over the fraction field of $T'$, above a finite place of the coefficient
field of $X_2$, that maps to the point of $y$. The abelian scheme over $T$ with its line bundle, the generator
and the basis $v'$ become, by base change to $T'$, data over $T'$ for this point, and by
Theorem~\ref{cv-definite-key-formula}(2) the value of $\kappa$ at them is the image $u_2$ of $u$. At the level
of $X_2$ the first case applies to every prime, so $u_2$ is a unit of $T'$. Hence
$e'\operatorname{ord}_T(u)=\operatorname{ord}_{T'}(u_2)=0$, so $u$ and $u^{-1}$ lie in $T$, and $u$ is a unit.
% Lean: thetaKeyRule_valuation_nonneg_integral_framed_point = GlobalTheory.KeyValAny.key_valuation_nonneg, the
% proof of KeyIntegralVal.key_valuation_bounds with exists_tied_det_any in place of exists_tied_det. The two
% differ only in the unit: key_unit_any splits on Level.GoodAt p for the prime p in w; key_unit_at uses
% exists_model_at (GoodAt p only) where key_unit uses exists_model (GoodAt at every prime). The bad case is
% GlobalTheory.T.thetaKeyRule_isUnit_at_bad_level_place = KeyBadPrime.key_unit_bad: the level is
% GlobalTheory.T.exists_goodLevel_multiple_same_adelicPart (Level.Dvd l l_2, l_2.A = l.A, GoodAt l_2 at every
% prime), the moduli scheme GlobalTheory.Moduli.moduli_scheme_exists, the lift
% GlobalTheory.T.modulipoint_lift_to_finerLevel_completeDVR (T', iT, phi, e' > 0, vV (iT a) = e' vV a), the data
% over T' GlobalTheory.T.exists_integralFramedData_completeDVR_extension (KeyPB), the unit KeyIntegralVal.key_unit
% at l_2, and the base change hkappa.2.1 (part (2) of the key formula). The level and the lift are Lemmas
% cv-volume-generators-key-valuation-level and cv-volume-generators-key-valuation-lift. The data over T' are
% FramedExt.exists_framed_ext: pullbacks along T -> T' (sabPull), exists_isLieBasis_baseChange for the basis, and
% generates_of_flat (flat base change of the sections, flatBC_iso) for the generator; with sabComp and sabCancel they
% are data for the point x_2.
\end{proof}

\begin{lemma}\label{cv-volume-generators-key-inverse}
In Situation~\ref{cv-cusp-setting}, suppose that the level of $X$ is a level $(M,\mathfrak M)$ with $M$ prime
to $15$ and that $f\ge3$; the hypothesis $f\ge3$ is not used in the proof. Let $Y$ be a deep cover of $X$
(Definition~\ref{cv-deep-cover}) with coefficient field $F_Y$, in the setting of
Definition~\ref{cv-volume-twisting-section} for $Y$ over $F_Y$. Then there are a rational section $u$ of $Y$ of
weight $\kappa(0,1)$, a rational section $u'$ of weight $\kappa(0,-1)$ whose valuation at every integral point of
$Y$ above every finite place of $F_Y$ is nonnegative, and $t\in F_Y^\times$, with $uu'=t$.
\end{lemma}
\leanref{Development/AutomorphicArithmetic/ThetaCorrespondence/ThetaPairLifts/HodgeBundleWeights/ExistsInverseIntegralInvariant.lean}{GlobalTheory.KeyInvInt.exists_inverse_integral_invariant}
\begin{proof}
Write $X$ for $Y$ and $F_c$ for $F_Y$ in Definition~\ref{cv-volume-twisting-section}. Then $p+q=r\ge1$; the
uniformization at every good embedding satisfies the compatibility conditions of the moduli problem
(Definition~\ref{cv-deep-cover}(5)); and $F_Y$ is algebraically closed in the
function field of every connected component of $Y$ (Definition~\ref{cv-deep-cover}(3)).

Let $\kappa$ be the rule of Theorem~\ref{cv-definite-key-formula} and $K_X$ a section given by
Lemma~\ref{cv-volume-twisting-key-section}(1). By part (2) of that lemma, $X$ has a key section $\rho$, an
invertible rational section of weight $\kappa(0,-1)$, and $1_Z\rho^4=a_Z1_ZK_X$ with $a_Z\in F_c^\times$ for every
connected component $Z$. If the prime-to-$\ell$ part of the deep level $M_Y$ is at least $3$ for every prime
$\ell$, Lemma~\ref{cv-volume-twisting-key-valuation} gives $v(K_X)\ge0$ at every integral point; otherwise
Lemma~\ref{cv-volume-generators-key-valuation} does. The components are finitely many, so there is a nonzero
$s\in\mathcal O_{F_c}$ with $sa_Z\in\mathcal O_{F_c}$ for every $Z$.

Let $v$ be the valuation at an integral point above a finite place $w$, normalized to extend
$\operatorname{ord}_w$. The point lies on some component $Z$, and $v(1_Z)=0$. As
$1_Z(s\rho)^4=s^3(sa_Z)\,1_ZK_X$,
\[
 4v(s\rho)=v\big(1_Z(s\rho)^4\big)=\operatorname{ord}_w\big(s^3(sa_Z)\big)+v(K_X)\ge0.
\]
So $u=\rho^{-1}$, $u'=s\rho$ and $t=s$ will do: $u$ has weight $-\kappa(0,-1)=\kappa(0,1)$, and $uu'=s$.
\end{proof}
\leanref{Development/AbelianVarietyTheory/Proofs/DivisorProofs/ExistsConstantSmulIntegralOfFourthPower.lean}{GlobalTheory.RuleIntegral.exists_smul_nonneg}
% Lean: exists_inverse_integral_invariant: kappa is KeyFamClosing.keyFam_one, K_X is
% KeyFourth.exists_unit_secZ_keyRule, rho and a_Z are KeyRoot.exists_keySec_pow_four at the good inclusion. The
% valuation bound is KeyIntegralVal.key_valuation_bounds when every prime is good, and otherwise
% GlobalTheory.T.thetaKeyRule_valuation_nonneg_integral_framed_point (Lemma cv-volume-generators-key-valuation).
% p + q >= 1 is from D.R.two_le_m, the uniformization hY.unf at goodOf of the inclusion, and FnAlgClosed hY.geo;
% the hypotheses PublishedHodgeBundleTheorems and 2 < [E:Q]/2 (the text's f >= 3) are not used. The
% scalar is GlobalTheory.T.exists_constant_smul_integral_of_fourth_power = RuleIntegral.exists_smul_nonneg
% (exists_common_denom, exists_idem_zero, nonneg_of_four), and u = rho^{-1} by AllWt.inv_mem_secZ with
% exOf 0 1 = -keyEx.

\begin{lemma}\label{cv-volume-theta-index-ideal}
Let $E$ be a number field and $n\ge0$. For an $\mathcal O_E$-submodule $M$ of $E^n$ let $\det M$ be the
$\mathcal O_E$-submodule of $E$ generated by the determinants of the $n\times n$ matrices whose rows lie in $M$.
Let $M'\subset M$ be $\mathcal O_E$-submodules of $E^n$ such that $kM\subset\mathcal O_E^n$ for some nonzero
$k\in\mathcal O_E$, and such that every $m\in M$ has a multiple $am\in M'$ with $a\in\mathcal O_E$ nonzero. Then
there is a nonzero ideal $\mathfrak a$ of $\mathcal O_E$ with $\det M'=\mathfrak a\det M$ and
$N\mathfrak a=[M:M']$.
\end{lemma}
\leanref{Development/AbelianVarietyTheory/Proofs/UnitaryExistence/ExistsIndexIdeal.lean}{GlobalTheory.BetaEx.exists_index_ideal}
\begin{proof}
Suppose first that $M=M'+\mathcal O_Em$ for some $m$, and let $\mathfrak c$ be the ideal of the $a\in\mathcal O_E$
with $am\in M'$. The map $a\mapsto am$ induces an isomorphism of $\mathcal O_E/\mathfrak c$ with $M/M'$, so
$N\mathfrak c=[M:M']$. We show that $\det M'=\mathfrak c\det M$.

Let $a\in\mathfrak c$, and let $w$ be a matrix with rows $w_i=v_i+t_im$ in $M$, where $v_i\in M'$ and
$t_i\in\mathcal O_E$. Expanding the determinant in each row, and using that a matrix with two rows proportional
to $m$ has determinant $0$, gives $\det w=\det v+\sum_jt_j\det v^{(j)}$, where $v^{(j)}$ is $v$ with its $j$-th row
replaced by $m$. So $a\det w=a\det v+\sum_j\det v'^{(j)}$, where $v'^{(j)}$ is $v$ with its $j$-th row replaced by
$t_jam\in M'$, and $a\det w\in\det M'$. Hence $\mathfrak c\det M\subset\det M'$.

Conversely, $\det M'\subset\det M$, so we may suppose that $\det M\ne0$. It is a fractional ideal, as
$k^n\det M\subset\mathcal O_E$, so it has an inverse $J$, with $\det M\cdot J=\mathcal O_E$. Let $v$ be a matrix
with rows $v_i\in M'$, and let $y\in J$. By Cramer's rule, $(\det v)\,m=\sum_i\det v^{(i)}\,v_i$. Here
$y\det v^{(i)}\in\det M\cdot J=\mathcal O_E$, as $v^{(i)}$ has rows in $M$; so $y\det v\cdot m\in M'$, and also
$y\det v\in\mathcal O_E$. So $y\det v\in\mathfrak c$. Writing $1=\sum_kx_ky_k$ with $x_k\in\det M$ and $y_k\in J$
gives $\det v=\sum_kx_k\,(y_k\det v)\in\mathfrak c\det M$.

In general, $M$ is finitely generated, since $kM$ is a submodule of the noetherian module $\mathcal O_E^n$. Let
$s_1,\dots,s_r$ generate it, and put $M_0=M'$ and $M_l=M_{l-1}+\mathcal O_Es_l$, so that $M_r=M$. Each $M_l$
satisfies $kM_l\subset\mathcal O_E^n$, and the case above gives ideals $\mathfrak c_l$, nonzero since $s_l$ has a
nonzero multiple in $M'$, with $\det M_{l-1}=\mathfrak c_l\det M_l$ and $N\mathfrak c_l=[M_l:M_{l-1}]$. Then
$\mathfrak a=\mathfrak c_1\cdots\mathfrak c_r$ has $\det M'=\mathfrak a\det M$ and
$N\mathfrak a=\prod_l[M_l:M_{l-1}]=[M:M']$.
\end{proof}
% Lean: the cyclic case is detSpan_cyclic (smul_det_mem for the first inclusion, cramer_rows and the inverse of the
% fractional ideal detSpan M for the second) and index_cyclic; the induction over generators is key in
% exists_index_ideal. The Lean's detSet takes the columns of the matrix; the determinants are the same.

\begin{lemma}\label{cv-volume-theta-small-scalar}
Let $E$ be a CM field with complex conjugation $x\mapsto\bar x$, let $d\equiv7\pmod{16}$ be a prime and
$\sqrt{-d}\in E$ an element with square $-d$, such that no finite prime of the totally real subfield of $E$
ramifies in $E$, let $\theta$ be a square root of $-d$ in $\mathbf C$, and let $\Phi$ be the CM type of the
embeddings $\tau\colon E\to\mathbf C$ with $\tau(\sqrt{-d})=\theta$. Let a PEL datum of signature $(p,q)$ over $E$
and a level (Definition~\ref{cv-model-toroidal-semiabelian}(1)) be given, with coefficient field
$F\subset\mathbf C$, which contains $\tau(E)$ for every $\tau$. Put $m_\tau=p$ for $\tau\in\Phi$ and $m_\tau=q$ for
the other embeddings $\tau$. Let $n\ge0$ and let $\Lambda\subset E^n$
be an $\mathcal O_E$-submodule that is its own dual for the Hermitian form $\sum_b\bar v_bw_b$: a vector
$v\in E^n$ lies in $\Lambda$ if and only if $\sum_b\bar v_bw_b\in\mathcal O_E$ for every $w\in\Lambda$. Put
$\Lambda'=\Lambda\cap\mathcal O_E^n$, and for an $\mathcal O_E$-submodule $L$ of $E^n$ let $I_\tau(L)$ be the
$\mathcal O_F$-submodule of $F$ generated by the $\tau(\delta)$ for $\delta\in\det L$
(Lemma~\ref{cv-volume-theta-index-ideal}).
Then there is $s\in F^\times$ with
\[
 s\prod_\tau I_\tau(\Lambda')^{m_\tau}\subset\prod_\tau I_\tau(\Lambda)^{m_\tau}\quad\text{and}\quad
 |\sigma(s)|^2[\Lambda:\Lambda']^{p+q}\le\exp\big(2\log[F:\mathbf Q]+2\log\operatorname{rd}_F\big)
\]
for every embedding $\sigma\colon F\to\mathbf C$.
\end{lemma}
\leanref{Development/AbelianVarietyTheory/Proofs/DivisorProofs/ExistsSmallScalarReflexIdealSelfDualLatticeStatement.lean}{GlobalTheory.T.exists_small_scalar_reflexIdeal_selfDualLattice}
\leanref{Development/AbelianVarietyTheory/Proofs/DivisorProofs/ExistsSmallScalarReflexIdealSelfDualLattice.lean}{GlobalTheory.SmallScal.exists_small_scalar}
\begin{proof}
First, $M=\Lambda$ and $M'=\Lambda'$ satisfy the hypotheses of Lemma~\ref{cv-volume-theta-index-ideal}. The
lattice $\Lambda$ spans $E^n$ over $E$. Otherwise some nonzero $E$-linear form $f$ vanishes on $\Lambda$; write
$f(w)=\sum_b\bar u_bw_b$. Then $u\in\Lambda$ by self-duality, as $f(w)=0$ for $w\in\Lambda$, so
$\sum_b\bar u_bu_b=f(u)=0$. At an embedding $\varphi$ of $E$ this reads $\sum_b|\varphi(u_b)|^2=0$, as
$\varphi(\bar x)=\overline{\varphi(x)}$ on a CM field; so $u=0$ and $f=0$, a contradiction. Hence for each $b$ there
is a nonzero $a_b\in\mathcal O_E$ with $a_be_b\in\Lambda$, for the unit vectors $e_b$. Let $a=\prod_ba_b$, so that
$ae_b\in\Lambda$ for every $b$. For $v\in\Lambda$, self-duality applied to $ae_b\in\Lambda$ gives
$\bar a\,v_b\in\mathcal O_E$; so $\bar a\Lambda\subset\mathcal O_E^n$ with $\bar a\ne0$. Every $v\in\Lambda$ has a
nonzero multiple $cv\in\mathcal O_E^n$ with $c\in\mathcal O_E$, and $cv\in\Lambda'$.

Lemma~\ref{cv-volume-theta-index-ideal} gives a nonzero ideal $\mathfrak a$ of $\mathcal O_E$ with
$\det\Lambda'=\mathfrak a\det\Lambda$ and $N\mathfrak a=[\Lambda:\Lambda']$. Applying $\tau$,
$I_\tau(\Lambda')=\tau(\mathfrak a)\mathcal O_F\cdot I_\tau(\Lambda)$. Let
$B=\prod_\tau\big(\tau(\mathfrak a)\mathcal O_F\big)^{m_\tau}$, so that
$\prod_\tau I_\tau(\Lambda')^{m_\tau}=B\prod_\tau I_\tau(\Lambda)^{m_\tau}$. Let $2f=[E:\mathbf Q]$ and let $k=[F:\mathbf Q]/2$, the
number of complex places of $F$, which is totally complex as it contains the CM field $\tau(E)$. Then
$N\big(\tau(\mathfrak a)\mathcal O_F\big)=N\mathfrak a^{[F:\mathbf Q]/2f}$ and $\sum_\tau m_\tau=f(p+q)$, so
$NB=[\Lambda:\Lambda']^{(p+q)k}$.

Put $R=\exp\big(2\log[F:\mathbf Q]+2\log\operatorname{rd}_F\big)$ and $X=[\Lambda:\Lambda']^{p+q}$. By Minkowski's
theorem for the fractional ideal $B^{-1}$ \cite[Chapter~I, Theorem~5.3]{Neu99}, there is a nonzero $s\in B^{-1}$ with
$|\sigma(s)|^2<R/X$ at every embedding $\sigma$, provided that
\[
 \Big(\frac2\pi\Big)^k\sqrt{|d_F|}\,N(B^{-1})<\Big(\frac RX\Big)^k,\qquad\text{that is,}\qquad
 2^k\sqrt{|d_F|}<\pi^kR^k.
\]
This holds, as $2<\pi$ and $\sqrt{|d_F|}\le|d_F|=\operatorname{rd}_F^{[F:\mathbf Q]}\le R^k$. Then $sB\subset\mathcal O_F$, so
$s\prod_\tau I_\tau(\Lambda')^{m_\tau}=sB\prod_\tau I_\tau(\Lambda)^{m_\tau}\subset\prod_\tau I_\tau(\Lambda)^{m_\tau}$, and
$|\sigma(s)|^2X<R$.
\end{proof}
% Lean: bdd_of_unimod (the spanning and the bound) and exists_smul_std; the ideal identity e uses span_image_span
% and span_image_smul; the norm is absNorm_map_ringOfIntegers_hom with RefSum.sum_mult; Minkowski's theorem is
% Formalism.ExtensionFiltration.mk_lt, with g = sqrt(R / X) at every place and R = exp(fieldSz F)^2.

\begin{proposition}\label{cv-arithmetic-volume-step}
In the setting of Proposition~\ref{cv-point-conditions}, after enlarging $C_2$ and $x_0$,
$\mu_j^{\mathrm{aut}}\ge-\mathrm{main}_j-C_2u$ at every point and for every $j$.
\end{proposition}
\leanref{Development/AutomorphicArithmetic/ThetaCorrespondence/ThetaPairLifts/ThirdPlaceSeries/LowOfCoverHypotheses.lean}{GlobalTheory.CovTop.LowCover.low_of_cover_hypotheses}
\leanref{Development/AbelianVarietyTheory/Proofs/DivisorProofs/LowerBoundOfAllEstimates.lean}{GlobalTheory.CovTop.LowCover.lowAt_of_cover}
\leanref{Development/AutomorphicArithmetic/ThetaCorrespondence/ThetaPairLifts/ThirdPlaceSeries/CoveringLowDefectBound.lean}{GlobalTheory.CovTop.DefLow.covering_low_defect_alt}
\begin{proof}
Write the weight $\kappa_j$ as a pair $(\alpha,n)$ of a vector $\alpha$ and an even integer $n$, let $h$
be the Hodge weight and $c_0$ the integer of the datum by which weights are shifted, and let
$a_\tau$, for $\tau$ in the CM type $\Phi$, be the anchor weights of $\kappa_j$, its entries in
position $k-1$. Put
$x^\varepsilon=\log b$ and $z=fr\,x^\varepsilon$ with $r=k+m-1$.

The case $j=m$. The signature is $(k+m-1,0)$, the group $U(W_m)$ is definite, and its Shimura
variety is zero-dimensional. Let $c_1=-\frac p2\log\Omega_E(\chi)$ be minus the degree of the flat
sections of Theorem~\ref{cv-flat-degree}, for $q=0$. For each $\tau\in\Phi$ let $\nu_\tau$
be the highest weight of the representation of $U(p)$ attached to $\kappa_m$ at $\tau$, with size
$|\nu_\tau|$ and dimension $\dim\nu_\tau$, and put
\[
 B_m=\sum_\tau\Big(\Big(\frac{|\nu_\tau|}p+\frac n2\Big)c_1-\frac12\log\dim\nu_\tau\Big)
 +n\,f\,p\Big(\frac14\log2\pi-\frac12\log\operatorname{rd}_E\Big).
\]
The moduli scheme $X$ of this Shimura variety over its coefficient field $F_c$, with its
uniformizations, is as in Situation~\ref{cv-definite-situation} with $p=k+m-1$: its level is
divisible by the two distinct level primes $g_0$ and $g_1$, so its prime-to-$\ell$ part is at least $3$
for every prime $\ell$, and $F_c$ is algebraically closed in the function field of every connected
component. By Theorem~\ref{cv-flat-degree} with $q=0$, flat sections of every type $\tau\in\Phi$ exist
on every component and have degree $-c_1$. The forms of weight $\kappa_m$ are the forms of weight
$(\nu,n)$ of the situation, and all of them are cusp forms.

Let $C$ be the set of connected components of $X$, a finite nonempty set of points
(Lemma~\ref{cv-definite-point}(1)). Let $V$ be the $F_c$-space of forms of weight $\kappa_m$ with its
Petersson norms, let $M_0\subset V$ be the finitely generated module of integral forms, and for $Z\in C$
let $e_Z$ be multiplication by the function that is $1$ on $Z$ and $0$ on the other components. These
are orthogonal idempotents with sum $1$, and $e_Z(V)$ is the space of forms supported on $Z$. Let $P_0$
be the level subgroup, of the elements of the finite level group of
Proposition~\ref{cv-point-conditions} that are upper unitriangular at every level place. It acts
on $X$, hence on $C$, and on $V$ through the level by a representation $\rho$, with
$\rho(g)e_Z=e_{gZ}\rho(g)$. Each $\rho(g)$ preserves $M_0$, and for every framed point $y$ over
$\mathbf C$ on $gZ$ that is metric with constant $1$ there is a framed point $y'$ on $Z$ with the same
property and $(\rho(g)f)(y)=f(y')$ for all $f$. The action of $P_0$ on $C$ is free: each $g\in P_0$ acts
on the lattice of the datum by an endomorphism $\gamma$ with $(\gamma-1)^r$ divisible by the level
$g_0g_1$, and a level automorphism that fixes a point of $X$ and is unipotent modulo two distinct primes
is the identity. By (5a) in the proof of Proposition~\ref{cv-point-conditions},
$\Lambda_m^{\mathrm{aut}}$ is the set of $f\in M_0$ with $\rho(u)f=\theta(u)f$ for all $u\in P_0$, where
$\theta$ is the character of $P_0$ with values in $F_c^\times$ defined there.

For each $Z\in C$, Lemma~\ref{cv-definite-fibre} with $\nu$, $n$ and $c_1$ gives $I_Z$, $L_Z$,
$Y_{Z,\sigma}$ and $\varphi_Z$ with $\mu(L_Z)\ge2B_m$; its number $B$ is $B_m$, since $|\Phi|=f$. The
Petersson norm of $v\in V$ at $\sigma$ is the square root of a mean, for a probability measure, of
$|v(y)|^2$ over framed points $y$ over $\mathbf C$ via $\sigma$ that are metric with constant $1$.
Suppose that every $e_Zv$ is $u\,\rho(g)\varphi_o(x)$ with $Z=go$, $|\sigma(u)|=1$ and
$|Y_{o,\sigma}\sigma(x)|\le B$. At such a framed point $y$ on $Z$,
$|v(y)|=|(e_Zv)(y)|=|\varphi_o(x)(y')|\le B$ for a framed point $y'$ on $o$, metric with constant $1$,
by Lemma~\ref{cv-definite-fibre}(3); so the Petersson norm of $v$ at $\sigma$ is at most $B$.
Lemma~\ref{cv-definite-orbits}, with $K=F_c$, the group $P_0$, the character $\theta$, the module
$M=\Lambda_m^{\mathrm{aut}}$ and $s=2B_m$, gives $\Lambda_m^{\mathrm{aut}}\ne0$ and
$\mu^{\mathrm{aut}}_m\ge2B_m$.
% Lean: GlobalTheory.T.two_mul_defBound_sub_le_muCore_of_gramAt
% = GlobalTheory.CovTop.DefCore.lifts_two_mul_defBound_sub_unitE_le_muCore, with constant 0. Its case
% FnAlgClosed is GlobalTheory.CovTop.DefCore.defect_bound_of_connected_alt; the other case,
% GlobalTheory.T.two_defBound_sub_unitError_le_muCore_alt
% (= GlobalTheory.CoreDesc.isDefined_of_not_functionField_algClosed_alt), is vacuous by
% GlobalTheory.CoreDesc.fnAlgClosed_core_alt. The situation is DefCase, from
% GlobalTheory.T.finite_nonempty_components_and_dict_mem_dom (DefCore.defCase_core_alt, with
% goodAt_of_two_primes); the flat sections are GlobalTheory.T.flatExistsAt_last_index_of_fnAlgClosed;
% P_0, the character and eigLat = latTypeK are GlobalTheory.T.exists_finite_group_actLev_char
% (DefEig.exists_finite_group_lev_char_slope); freeness is GlobalTheory.T.freePt_last_index_of_fnAlgClosed
% (DefFree.free_base_alt); the orbit step is GlobalTheory.T.levelMove_eigenlattice_slope_definite
% = DefPt.OrbitGeo.eigLat_slope (actLev_idem, actLev_int, homN_actLev, petAt_le), applied in
% GlobalTheory.T.eigenlattice_slope_lower_bound_definite_case. B_m is the definition defBound.
To bound $2B_m$ below, note that the entries of $\kappa_m$ at
$\tau$ satisfy $\sum_i\alpha_{\tau,i}=|\nu_\tau|+pn/2$. As $c_1=-\frac p2\log\Omega_E(\chi)$, the
term of $2B_m$ at $\tau$ is
$-\big(\sum_i\alpha_{\tau,i}\big)\log\Omega_E(\chi)-\log\dim\nu_\tau$. Write
$\sum_i\alpha_{\tau,i}=p\,a_\tau+\sum_i(\alpha_{\tau,i}-a_\tau)$ and bound
$\big|\sum_i(\alpha_{\tau,i}-a_\tau)\log\Omega_E(\chi)\big|$ by
$|\log\Omega_E(\chi)|\sum_i|\alpha_{\tau,i}-a_\tau|$. The last term of $2B_m$ is
$nfp\big(\frac12\log2\pi-\log\operatorname{rd}_E\big)\ge-nfp\,|\log\operatorname{rd}_E|$. Summing over
$\tau$,
\[
 2B_m\ge-\sum_\tau p\,a_\tau\log\Omega_E(\chi)-|\log\Omega_E(\chi)|\sum_{\tau,i}|\alpha_{\tau,i}-a_\tau|
 -\sum_\tau\log\dim\nu_\tau-n\,f\,p\,|\log\operatorname{rd}_E|.
\]
The anchors are chosen so that $\big|\mathrm{main}_m-\sum_\tau p\,a_\tau\log\Omega_E(\chi)\big|\le C_Au$,
and the size conditions bound the other three terms by multiples of $u$. Hence
$\mu^{\mathrm{aut}}_m\ge-\mathrm{main}_m-C_2u$.

The case $j<m$. By Corollary~\ref{cv-deep-cover-sized}, after enlarging $x_0$, at every point the
moduli scheme of $U(W_j)$ has a deep cover $Y$ of exponent $\lfloor3x^\varepsilon/\log3\rfloor$ and size
$391z^2$ (Definition~\ref{cv-deep-cover}), with coefficient field $F_Y$. A cusp form on $Y$ is integral if
its valuation is nonnegative at every integral point of $Y$, algebraic over $F_Y$, with values in a
complete discrete valuation ring above a finite place of $F_Y$. The size of the cover and the size
conditions give $\log|d_{F_Y}|/[F_Y:\mathbf Q]\le391z^2\le b$.
% Lean: the cover is GlobalTheory.T.covering_deepEx_of_prints (Corollary cv-deep-cover-sized).
Let $\kappa^\pm=\kappa_j\pm c_0h$, let $A$ be the lattice of
integral cusp forms on $Y$ of weight $\kappa^-$, and let $B$ be the lattice of integral cusp forms on
$Y$ of weight $\kappa_j$ in the span of the pullbacks of the forms of the prescribed type. On $Y$
the following hold.
\begin{enumerate}
\item (Maps.) There are $n'$ linear maps $\varphi_h$ from forms of weight $\kappa^-$ to forms of
  weight $\kappa_j$ with $\varphi_h(A)\subset B$, jointly injective, with
  $\sum_hN_\sigma(\varphi_h(f))^2\le e^{2b}N_\sigma(f)^2$ at every embedding $\sigma$, and with
  $n'\operatorname{rk}B\le\operatorname{rk}(\kappa^+)$, the rank of the integral cusp forms of weight
  $\kappa^+$.
\item (Ranks.) $\operatorname{rk}(\kappa^+)\le(1+b/K_j)\operatorname{rk}A$.
\item (Generators.) $B$ is spanned over $F_Y$ by elements of $B$ of norm at most $e^{K_j}$ at
  every embedding.
\item (Flat sections.) $\mu(A)\ge\mu_0-\mathrm{main}_j-C_Fu$, where $\mu_0$ is the slope of the
  integral sections of the weight $u_-$ obtained from $\kappa^-$ by removing the anchor.
\item (Uniform lower bound.) $\mu_0\ge-C_Uu$.
\item (Ring.) $\mu(B)-2b\le\mu_j^{\mathrm{aut}}$.
\end{enumerate}

Notation for (1) to (6). Let $R$ be the ring of sections of all weights on $Y$, an $F_Y$-algebra, and
let $\Gamma_Y$ be the finite group of level automorphisms of $Y$. An element $\gamma\in\Gamma_Y$ acts on
$R$ by the pullback $f\mapsto\gamma^*f$, an $F_Y$-algebra automorphism, and
$(\gamma\gamma')^*=\gamma'^*\gamma^*$. For a weight $\kappa$ let $S_Y(\kappa)\subset R$ be the $F_Y$-space
of cusp forms of weight $\kappa$ on $Y$, with the Petersson norms $N_\sigma$, and let
$\operatorname{rk}(\kappa)$, $\mu(\kappa)$ and $\deg(\kappa)=\frac12\operatorname{rk}(\kappa)\mu(\kappa)$ be
the rank, the slope and the degree of its lattice of integral forms; for a connected component $Z$
of $Y$ let $\operatorname{rk}_Z(\kappa)$ and $\deg_Z(\kappa)$ be the rank and the degree of the integral
forms of weight $\kappa$ over $Z$. Thus $\mu(A)=\mu(\kappa^-)$ and $\mu_0=\mu(u_-)$. Let $a_\tau$ be the
anchor weights of $\kappa_j$ and $u_W$ the weight obtained from $\kappa_j$ by removing the anchor, so
that $u_-=u_W-c_0h$. The proofs use the following properties of $Y$, in which $C_Y$, $C_1$ and $C_2$
are nonnegative constants of the family; the size conditions give $2C_Y+200\le z$, $z^{12}\le b$ and
$z^{12}\le PP$ at all large scales.
\begin{enumerate}
\item[(Y1)] (Translations and products.) For every even weight $\kappa$ and every $\gamma\in\Gamma_Y$,
  $\gamma^*$ maps $S_Y(\kappa)$ to itself and preserves integral forms and the norms $N_\sigma$. The
  product of cusp forms of weights $\kappa$ and $\kappa'$ is a cusp form of weight $\kappa+\kappa'$,
  integral if both factors are. For $\kappa=\kappa_j+sc_0h$ ($s=-1,0,1,2$), for $\kappa=u_-$ and for
  $\kappa=u_W+sc_0h$ ($s=1,2$) the space $S_Y(\kappa)$ is finite-dimensional, $N$ is a Hermitian norm
  system on it (Definition~\ref{cv-hermitian-norms}) with $N_{\bar\sigma}=N_\sigma$, and its integral forms
  are a finitely generated $\mathcal O_{F_Y}$-module that spans it; the integral forms supported on a
  connected component $Z$ are a finitely generated $\mathcal O_{F_Y}$-module that spans the forms supported
  on $Z$. Moreover $\operatorname{rk}(\kappa^-)>0$, and $Y$ has finitely many connected components, at
  least one.
\item[(Y2)] (Type.) There are a subgroup $H\subset\Gamma_Y$ with $\log|H|\le45z^2$ and a character
  $\chi_H\colon H\to F_Y^\times$ such that $B$ is the set of integral $f\in S_Y(\kappa_j)$ with
  $\gamma^*f=\chi_H(\gamma)f$ for every $\gamma\in H$.
\item[(Y3)] (Auxiliary form.) There is an integral $v\in S_Y(c_0h)$ whose group determinant
  $\det\big((\gamma\gamma')^*v\big)_{\gamma,\gamma'\in H}$ is not a zero divisor in $R$, such that
  $N_\sigma(f\cdot\gamma^*v)\le e^{C_Yz^9}N_\sigma(f)$ for every $\gamma\in\Gamma_Y$, every cusp form $f$
  of even weight and every $\sigma$.
\item[(Y4)] (Short forms.) The space $S_Y(\kappa_j)$ is spanned over $F_Y$ by integral forms $y$ with
  $N_\sigma(y)\le\exp\big(C_Yz^2\operatorname{size}(\kappa_j)\big)$ for every $\sigma$, where
  $\operatorname{size}(\kappa_j)=\sum_{\tau,i}|\kappa_{j,\tau,i}|\ge1$, so that
  $K_j=PP\cdot\operatorname{size}(\kappa_j)$.
\item[(Y5)] (Twisting.) Let $c_\tau$ be minus the mean over the components of $Y$ of the degrees of
  the flat sections of type $\tau$ (Definition~\ref{cv-constant-sections}), let $\lambda_Y$ be the sum
  of $\log\ell$ over the bad primes $\ell$ of $Y$, and put
  \[
   \Theta=\sum_\tau a_\tau c_\tau-q\log\operatorname{rd}_E\sum_\tau|a_\tau|
   -\frac n2fr\Big(\frac12\log2\pi+\log\operatorname{rd}_E\Big)-\lambda_Y.
  \]
  Then $\operatorname{rk}(\kappa^-)=\operatorname{rk}(u_-)$, $\deg(u_-)+\operatorname{rk}(u_-)\Theta\le\deg(\kappa^-)$,
  and $\operatorname{rk}(\kappa_j+sc_0h)=\sum_Z\operatorname{rk}_Z(u_W+sc_0h)$ for $s=-1,1,2$. Moreover
  $\lambda_Y\le40x^\varepsilon$.
\item[(Y6)] (Components.) Let $D$ be the dimension of the flag variety of $U(W_j)$. There are a
  weight $u_1$, an integer $\ell\ge3D+5$, numbers $0\le\rho<1$ and $\zeta=7(D^2/\ell+D\rho)$ with
  $\zeta\le\frac12$ and $K_j\zeta\le b$, and for each component $Z$ a geometric degree $\delta_Z>0$
  and a height $h_Z$ of $u_1$ on $Z$, such that, with $\beta=\binom{\ell-1}D$:
  \begin{enumerate}
  \item[(a)] there is $\delta'>0$ with $\delta'\le\delta_Z\le\delta'/(1-\rho)$ for every $Z$;
  \item[(b)] $\beta\delta_Z\le\operatorname{rk}_Z(u_-)$, $(1-\zeta)\operatorname{rk}_Z(u_-)\le\beta\delta_Z$, and
    $\operatorname{rk}_Z(u_W+sc_0h)\le(1+\zeta)\beta\delta_Z$ for $s=1,2$;
  \item[(c)] there are $A'>0$ and $R'\ge0$ with $A'\log\max(A',1)+R'\le C_2u$ and
    $\beta\delta_Z\big(\ell h_Z-3(D+1)|h_Z|-A'\log(2+|h_Z|)-R'\big)\le\deg_Z(u_-)$ for every $Z$;
  \item[(d)] there are $c_s,c_c,X\ge0$ and numbers $s_Z$ with $h_Z+c_s\ge0$ and
    $(1-\rho)\max(0,s_Z+c_s)\le h_Z+c_s+c_c$ for every $Z$, with mean of the $s_Z$ at least $-X$,
    and with $\ell X\le C_1u$, $\ell c_c\le C_1u$ and $c_s\big(6D+10+\ell(\zeta+2\rho)\big)\le C_1u$;
  \item[(e)] $\sum_Z\operatorname{rk}_Z(u_-)=\operatorname{rk}(u_-)$ and $\sum_Z\deg_Z(u_-)\le\deg(u_-)$.
  \end{enumerate}
\item[(Y7)] (Pullback.) The field $F_Y$ is a Galois extension of $F_c$. Pulling back to $Y$ is an
  additive map $\pi$ from the cusp forms of weight $\kappa_j$ on the Shimura variety of $U(W_j)$ to
  $S_Y(\kappa_j)$, with $\pi(af)=a\,\pi(f)$ for $a\in F_c$, which carries $F_c$-linearly independent
  families to $F_Y$-linearly independent families, satisfies $N_{\sigma'}(\pi(f))=N_{\sigma'|F_c}(f)$ for
  every embedding $\sigma'$ of $F_Y$, and maps a form to an integral form if and only if the form is
  integral. For cusp forms $b_1,\dots,b_m$ of weight $\kappa_j$ on the Shimura variety,
  $x\in F_Y^m$ and $s\in\operatorname{Gal}(F_Y/F_c)$, if $\sum_ix_i\pi(b_i)$ is integral then so is
  $\sum_is(x_i)\pi(b_i)$.
\end{enumerate}
Derivation of (Y1) to (Y7). At every point at a scale $x\ge x_0$, the signature of $W_j$ has $q=m-j\ge1$,
$f\ge8$, the level of the Shimura variety of $U(W_j)$ is a level $(g_0g_1,\mathfrak M)$ with $g_0g_1$ prime
to $15$, and $Y$ is a deep cover of it (proof of Corollary~\ref{cv-deep-cover-sized}). So
Lemmas~\ref{cv-cover-level-automorphisms}, \ref{cv-cover-pullback} and~\ref{cv-cover-type} apply, with $X$
the moduli scheme of $U(W_j)$ and with the group $\Gamma_Y$ and the action $\gamma^*$ of
Lemma~\ref{cv-cover-level-automorphisms}. The weights $\kappa_j+sc_0h$, $u_-$ and $c_0h$ have even last
entry. As $M_Y=15^eg_0g_1$ with $e\log15\le9x^\varepsilon$, $\log g_k\le9x^\varepsilon+1$ and
$x^\varepsilon\ge10$,
\[
 \log M_Y\le27x^\varepsilon+2\le30x^\varepsilon.
\]

(Y7) This is Lemma~\ref{cv-cover-pullback}(1), (3) and~(5) for $\kappa=\kappa_j$, with the integrality
statement of Lemma~\ref{cv-cover-pullback}(2).

(Y2) Let $G$ be the finite level group of Proposition~\ref{cv-point-conditions}, acting on the forms of
$X$ through the matrices $\ell_u$ of level automorphisms as in the proof of that proposition, so that
$u\cdot f$ is the action $f\mapsto uf$ of $G$ on cusp forms there, which maps $S_{\kappa_j}$ to itself. Let
$G_X$ be the group of $\mathcal O_E$-linear automorphisms of $\Lambda/\mathfrak M\Lambda$. Bases of the
spaces $\Lambda/\mathfrak P\Lambda$ for the primes $\mathfrak P\supset\mathfrak M$ of $E$ above the level
places identify $G_X$ with $G$, compatibly with the matrices $\ell_u$; let $\theta\colon\Gamma_Y\to G$ be
reduction modulo $\mathfrak M$, defined since $\mathfrak M_Y\subset\mathfrak M$, followed by this
identification. Conditions (a) and (b) of Lemma~\ref{cv-cover-type} hold by construction. For (c), write
$\mathfrak M_Y=\mathfrak MI$ as in Definition~\ref{cv-deep-cover}(1); the ideals $\mathfrak M$ and $I$ are
coprime, since $\mathfrak M$ contains $M$ and $I$ contains $15^e$, so
$\Lambda/\mathfrak M_Y\Lambda=\Lambda/\mathfrak M\Lambda\times\Lambda/I\Lambda$. Hence reduction modulo
$\mathfrak M$ maps the subgroup of elements of $\Gamma_Y$ that are the identity modulo $I$ bijectively onto
$G_X$, and the inverse of this bijection, followed by the inclusion, is a section of $\theta$. By (5a) of
the proof of Proposition~\ref{cv-point-conditions}, the forms of the prescribed type are the $f\in S_{\kappa_j}$ with
$u\cdot f=\psi_0(u)^{-1}f$ for every $u$ in the level subgroup $P_0$, and $\psi_0$ takes values in
$F_c^\times$. Put $H=\theta^{-1}(P_0)$, the group of $\gamma\in\Gamma_Y$ whose reduction modulo $\mathfrak M$
is upper unitriangular at every level place, and $\chi_H(\gamma)=\psi_0(\theta(\gamma))$.
Lemma~\ref{cv-cover-type}(3), with $P=P_0$ and $\lambda=\psi_0^{-1}$, shows that a form
$g\in S_Y(\kappa_j)$ lies in the span of the pullbacks of the forms of the prescribed type if and only if
$\gamma\cdot g=\psi_0(\theta(\gamma))^{-1}g$ for $\gamma\in H$, that is, $\gamma^*g=(\gamma^{-1})\cdot g=\chi_H(\gamma)g$
for $\gamma\in H$. Taking integral forms gives the description of $B$. By
Lemma~\ref{cv-cover-level-automorphisms}(1),
$\log|H|\le4f^2r^2\log M_Y\le120f^2r^2x^\varepsilon\le45z^2$, as $x^\varepsilon\ge3$.
% Lean: (Y2) is GlobalTheory.CovTop.TypeLattice.latDeep_eq_latN (the bases through
% GlobalTheory.T.exists_typeFrame_cover_levelMaps_reductions and GlobalTheory.TypeBase.base_level_unip_char,
% with typeK_eq for the prescribed type; chiHR = genChar o typeRed, and lam = (genChar o s)^{-1} is the text's
% lambda) and GlobalTheory.CovTop.TypeLattice.log_card_le. In latDeep_eq_latN, theta is t o typeRed fr;
% (a) and (b) are the facts h1 and h2 of exists_typeFrame_cover_levelMaps_reductions, and (c) is the
% section GlobalTheory.LevDesc.typeEmb with typeRed_typeEmb, the inverse of the bijection
% IsTypeLevel.typeRed_bijective on grpG (the elements trivial modulo the cofactor), which rests on the
% decomposition IsTypeLevel.parts of the level automorphisms.

(Y1) The translations and the products are Lemma~\ref{cv-cover-level-automorphisms}(3) and~(4). For every
weight $\kappa$ the space $S_Y(\kappa)$ is finite-dimensional: the readings at a good embedding are
families of cusp forms of a given weight on the finitely many components, which form a finite-dimensional
complex space by Lemma~\ref{cv-volume-forms-finite} for the uniformization of $Y$, and readings carry
$F_Y$-linearly independent families to $\mathbf C$-linearly
independent families, as in Step~4 of the proof of Proposition~\ref{cv-cusp-rank-bound}. The norm
$N_\sigma$ comes from a Hermitian inner product, the integral of the product of the readings at the good one
of $\sigma$ and $\bar\sigma$, which is definite since a nonzero form has a nonzero reading. So $N$ is a
Hermitian norm system with $N_{\bar\sigma}=N_\sigma$. The components of $Y$, the irreducible components of
$Y\otimes F_Y$, are finitely many because $Y\otimes F_Y$ is a noetherian scheme, and there is at least one
because $Y$ is not empty.

The integral forms split over the components. At a good $\sigma$, the frames of a component $Z$ are those of
the components of the uniformization that lie over $Z$, and a form is supported on $Z$ if it is the product
of itself with the idempotent of $Z$. This idempotent is an integral section of weight $0$, so the integral
forms of a weight $\kappa$ are the direct sum of the integral forms supported on the components, and their
rank is the sum of the ranks $\operatorname{rk}_Z(\kappa)$. Forms supported on distinct components are
orthogonal for the Petersson inner products, and $N_\sigma(f)^2=m_{Z,\sigma}N_{Z,\sigma}(f)^2$ for $f$ supported
on $Z$, where $N_{Z,\sigma}$ is the Petersson norm on $Z$, computed with the probability measure on the part
of the adelic quotient over $Z$, and $m_{Z,\sigma}\le1$ is the measure of that part, which is positive when
$\operatorname{rk}_Z(\kappa)>0$. Hence, when the integral forms of weight $\kappa$ are finitely generated,
\[
 \deg(\kappa)=\sum_Z\Big(\deg_Z(\kappa)-\frac{\operatorname{rk}_Z(\kappa)}{2[F_Y:\mathbf Q]}
 \sum_\sigma\log m_{Z,\sigma}\Big)\ge\sum_Z\deg_Z(\kappa).\tag{$*$}
\]
% Lean: GlobalTheory.Deep.finiteDimensional_V_of_good (GlobalTheory.cuspW_finiteDimensional_of_isUnf, which proves
% the finiteness through a finite determining set of frames, exists_finite_determining_set_SS in
% GlobalTheory.T.finiteDimensional_cuspSections_vanishing_offDomain), GlobalTheory.HSApp.hasGram_nrm_of_good
% (exists_semilinear_L2_independent, goodOf_conj), GlobalTheory.IndefiniteRowSideFullHecke.ComponentMassSum.finite_comp
% and GlobalTheory.CovTop.good_comp_nonempty; DeepGoodAt is GlobalTheory.CovTop.DeepGoodAll.deepGood_alt_alt.
% The splitting is GlobalTheory.T.componentDegrees_sum_le_total_degree (cuspForms_supLattice_split_over_components,
% components_peterssonOrthogonal_and_masses, component_ranks_degrees_add_with_masses).

The integral forms of (Y1). By Definition~\ref{cv-volume-forms-weights}, $\kappa_j=\kappa(\lambda_j,n/2)$ for
the family $\lambda_j$ with $\lambda_{j,\tau,i}=\alpha_{\tau,i}+n/2$ for $i<p$ and
$\lambda_{j,\tau,i}=-\alpha_{\bar\tau,r-1-i}-n/2$ for $i\ge p$. Let $a$ be the family whose entries at $\tau$
are all $a_\tau$. Then $\kappa_j+sc_0h=\kappa(\lambda_j+sc_0\mathbf h,n/2)$ and
$u_W+sc_0h=\kappa(\lambda_j-a+sc_0\mathbf h,0)$; in particular $u_-=\kappa(\lambda_j-a-c_0\mathbf h,0)$, and
$\kappa_a=\kappa(a,n/2)$. The size conditions give that the families $\lambda_j+sc_0\mathbf h$ for
$s=-1,0,1$ and $\lambda_j-a+sc_0\mathbf h$ for $s=-1,1,2$ are regular. Adding $c_0\mathbf h$ to a regular
family gives a regular family, since $c_0\ge1$ and $\mathbf h$ is nonincreasing with
$\mathbf h_{\tau,p-1}-\mathbf h_{\tau,p}=1$; so $\lambda_j+2c_0\mathbf h$ is regular, and so is $c_0\mathbf h$.
Four inputs are used; their proofs follow the rank below.
\begin{enumerate}
\item[(i)] (CM points.) Let $\lambda$ be a regular family, $m\ge0$ and $\kappa=\kappa(\lambda,m)$. For every
  nonzero $F\in S_Y(\kappa)$ there are a number field $K\supset F_Y$, a framed point $y$ of $Y$ over $K$ and
  $N\in K^\times$ such that $F(y)\ne0$ and $N\,G(y)$ is integral over $\mathbf Z$ for every integral
  $G\in S_Y(\kappa)$.
\item[(ii)] (Powers.) Let $\lambda$ be a regular family and $f\in S_Y(\kappa(\lambda,0))$. There are $k\ge0$,
  a dominant weight $\kappa'=(\alpha',n')$ with $n'\ge8fr+4$, elements $g,h\in R$ and $t\in F_Y^\times$ such
  that $f^{k+1}g\in S_Y(\kappa')$, $gh=t$, and $h$ is integral: its valuation is nonnegative at every
  integral point of $Y$.
\item[(iii)] (Integral multiples.) Let $\kappa'=(\alpha',n')$ be a dominant weight with $n'\ge8fr+4$. Every
  $F\in S_Y(\kappa')$ has an integral multiple $cF$ with $c\in F_Y^\times$.
\item[(iv)] (Nonzero forms.) Let $D=f\binom r2$ be the dimension of the flag variety of $U(W_j)$, and let
  $\varkappa$ be the family with $\varkappa_{\tau,i}=q-p+2i+1$ for $i<p$ and $\varkappa_{\tau,i}=2(i-p)+1-p-q$
  for $i\ge p$. Suppose that $D\ge1$, and let $\lambda$ be a family with
  $(\lambda-\varkappa)_{\tau,i}-(\lambda-\varkappa)_{\tau,i+1}\ge\binom{D+1}2+1$ for every $\tau$ and every
  $0\le i<r-1$. Then every connected component $Z$ of $Y$ carries a nonzero cusp form of weight
  $\kappa(\lambda,0)$ supported on $Z$.
\end{enumerate}

Finite generation. Let $\lambda$ be a regular family, $m\ge0$ and $\kappa=\kappa(\lambda,m)$. For a nonzero
$F\in S_Y(\kappa)$ let $K$, $y$ and $N$ be as in (i). The map $G\mapsto N\,G(y)$ from $S_Y(\kappa)$ to $K$ is
$F_Y$-linear, since values at framed points are additive and $(cG)(y)=c\,G(y)$ for $c\in F_Y$; it does not
vanish at $F$, and it maps the integral forms into $\mathcal O_K$. As $S_Y(\kappa)$ is finite-dimensional,
Lemma~\ref{cv-volume-generators-functionals} shows that the integral forms of weight $\kappa$ are a finitely
generated $\mathcal O_{F_Y}$-module.

Spanning. Let $\lambda$ be a regular family and $f\in S_Y(\kappa(\lambda,0))$, and let $k$, $\kappa'$, $g$, $h$
and $t$ be as in (ii). By (iii) there is $c\in F_Y^\times$ with $cf^{k+1}g$ integral. The valuation of a
product at an integral point is the sum of the valuations, so $ct\,f^{k+1}=(cf^{k+1}g)h$ is integral. The
valuations at the integral points of $Y$ above the primes of $\mathcal O_{F_Y}$ satisfy the hypotheses of
Lemma~\ref{cv-volume-generators-power}, which, with $K=F_Y$ and $S=\{f^{k+1}\}$, gives $b\in F_Y^\times$ with
$bf$ integral. So $f$ lies in the $F_Y$-span of the integral forms of weight $\kappa(\lambda,0)$.

Now let $m\ge1$ and $F\in S_Y(\kappa(\lambda,m))$. Every connected component $Z$ carries a twisting section
$\theta_Z$ of weight $\kappa(0,1)$, the section of weight $(0,1)$ of (Y5); its power $\theta_Z^m$ is a
twisting section of weight $\kappa(0,m)$, and $\kappa(\lambda,m)=\kappa(\lambda,0)+\kappa(0,m)$. As in (Y5),
multiplication by $\theta_Z^m$ is a bijection from the cusp forms of weight $\kappa(\lambda,0)$ supported on
$Z$ to those of weight $\kappa(\lambda,m)$; and since the valuations of $\theta_Z^m$ at integral points are
bounded below, by $0$ above all but finitely many places, there is $c_Z\in F_Y^\times$ such that
multiplication by $c_Z\theta_Z^m$ maps integral forms to integral forms. Write $F=\sum_Ze_ZF$ with the
idempotents $e_Z$ of the components, and let $G_Z$ be the form of weight $\kappa(\lambda,0)$ supported on $Z$
with $\theta_Z^mG_Z=e_ZF$. By the case $m=0$ there is a nonzero $d_Z\in\mathcal O_{F_Y}$ with $d_ZG_Z$
integral, and then $c_Zd_Z\,e_ZF=c_Z\theta_Z^m(d_ZG_Z)$ is integral. So every $e_ZF$, and hence $F$, lies in
the span of the integral forms.

The weights of (Y1). For $\kappa_j+sc_0h=\kappa(\lambda_j+sc_0\mathbf h,n/2)$ with $s=-1,0,1,2$, and for
$u_W+sc_0h$ with $s=-1,1,2$, finite generation and spanning are the cases above, as $n$ is even. For a
connected component $Z$, the integral forms of one of these weights supported on $Z$ are a submodule of its
integral forms, hence finitely generated over the noetherian ring $\mathcal O_{F_Y}$; and a form supported
on $Z$ has a nonzero integral multiple, which is supported on $Z$, so they span the forms supported on $Z$.

The rank. The size conditions give $D\ge1$ and that the family $\lambda_j-a-c_0\mathbf h$ of $u_-$ satisfies
the hypothesis of (iv). Let $Z$ be a connected component, and let $g$ be a nonzero cusp form of weight $u_-$
supported on $Z$, given by (iv). Let $t_Z$ be a twisting section of weight $\kappa_a$ on $Z$, as in (Y5),
and $t'_Z$ its inverse on $Z$, so that $t_Zt'_Z=e_Z$. Then $gt_Z$ is a cusp form of weight
$u_-+\kappa_a=\kappa^-$, and it is nonzero, since $(gt_Z)t'_Z=ge_Z=g$. As the integral forms of weight
$\kappa^-$ span $S_Y(\kappa^-)$, $\operatorname{rk}(\kappa^-)=\dim_{F_Y}S_Y(\kappa^-)>0$.
% Lean: (Y1) integral forms. For kappa_j + s c_0 h (s = -1, 0, 1): GlobalTheory.CovTop.ShiftGood.covering_goodAt_shift
% (= GlobalTheory.T.covering_shift_preserves_goodAt), through goodAt_of_exZ and GlobalTheory.CovTop.Model.exZ_shiftG
% (kappa_j + s c_0 h = kappa(lambda_j + s c_0 h, n/2), lambda_j = laOf a n); finite generation
% GlobalTheory.CovTop.ShiftFG.weight_lattice_pow_fg (fg_of_integral_functionals with the map N * evR F y,
% ShiftFG.exists_good_cmPt_frame_ne_zero and GlobalTheory.T.exists_frame_denominator_dominant_weight_power_cm);
% spanning GlobalTheory.GenShift.span_weight_lattice_of_weights (span_latW_add with the twisting sections
% GlobalTheory.KeyTwDeep.twSec_keyRoot of weight wtOf 0 1, their powers exists_twPow, twSec_mulEquiv_integral and
% exists_twist_scalar). For m = 0: GlobalTheory.CovTop.GoodDom.covering_regular_at_domain
% (= GlobalTheory.T.covering_domain_lattice_regularity), with weight_lattice_fg_on_domain
% (GlobalTheory.T.cmPoint_detects_dominant_covering, GlobalTheory.T.exists_frame_denominator_dominant_weight_cm),
% weight_lattice_span_on_domain (GlobalTheory.T.weight_match_pow_mul_cuspidal,
% GlobalTheory.T.divisor_multiple_nonneg_eventually, GlobalTheory.SupFull.exists_smul_nonneg_of_pow) and, for the
% components, GlobalTheory.HSApp.goodOn_of_goodAt (goodOn_of_goodW). kappa_j + 2 c_0 h is
% covering_shift_preserves_goodAt with antitone_add_hodge and gap_add_hodge in
% GlobalTheory.BadEmb.regular_embeddings_cover_bad_embedding_locus_alt; rk_le_alt_alt reaches it instead through
% GlobalTheory.CovTop.HodgeSteps.covering_goodAt_shift, which twists u_W + 2 c_0 h by
% GlobalTheory.T.flatLine_thetaRoot_twist_twistingSection. The regular families come from
% GlobalTheory.T.weight_conditions_at_index' and the field hanti of weight_parameter_strict_antitone_alt (used in
% GlobalTheory.CovTop.HodgeSteps.rk_le_alt_alt); c_0 >= 1 is cZeroN_isLeast. rk(kappa^-) > 0 is
% GlobalTheory.CovTop.LowPos.lowered_rank_pos_alt_alt: GlobalTheory.T.exists_ne_zero_cuspOn_of_covers (input (iv),
% with D >= 1 and the gap condition from wtBig_run_alt and gapGe_wt_add_parWt), exists_twist_ne_zero_of_lowered
% (the Lean multiplies by the inverse of the twisting section s t^(-m) of exists_twSec_comb, of weight u_- - kappa^-)
% and rk_pos.

Proofs of (i) to (iv). Let $\sigma_0$ be the inclusion of $F_Y$ in $\mathbf C$, a good embedding. For a
component $Z$ at $\sigma_0$, the uniformization provides, as in (1)(a) of the proof of
Lemma~\ref{cv-cusp-inputs}, a finitely generated $\mathcal O_E$-submodule $L'_Z\subset E^r$ that spans $E^r$ and
is self-dual for $G_Z$. Framed points over $F_Y$-algebras, the framed points $y\cdot g$ and the transformation
law are as in Lemma~\ref{cv-volume-forms-frame-polynomial}. The proofs use the following facts about $Y$, which
are part of the open step below.
\begin{enumerate}
\item[(S1)] (CM points.) Let $Z$ be a component at $\sigma_0$, $x$ a frame of $Z$, and $v_1,\dots,v_r$ a basis of
  $E^r$ in $L'_Z$ with properties (1) to (4) of Lemma~\ref{cv-cusp-frame} for some $\ell>0$. There are a number
  field $K\supset F_Y$, an embedding $\sigma_1\colon K\to\mathbf C$ that is the identity on $F_Y$, and a framed
  point $y_0$ of $Y$ over $K$ such that $\sigma_1$ carries the point of $y_0$ to the point of the framed point of
  $Y$ at $x$, with the following property: for every weight $\kappa$ and all integers $k_\tau\ge0$, one for each
  embedding $\tau$ of $E$, there is a nonzero $t\in\mathcal O_K$ such that $tG$ has coefficients integral over
  $\mathbf Z$ whenever $f$ is an integral rational section of weight $\kappa$ and $G$ is a polynomial over $K$
  with $(\prod_\tau\det g_\tau^{k_\tau})f(y_0\cdot g)=G(g)$ for every family $g$ of invertible matrices over $K$.
\item[(S2)] (Flat units.) There is an integer $N\ge1$ such that for every $w\colon\Phi\to\mathbf Z$ there are
  rational sections $\varphi$ of weight $\kappa(\lambda_w,0)$, where $\lambda_{w,\tau,i}=Nw_\tau$, and $\varphi'$, and
  $t\in F_Y^\times$, with $\varphi\varphi'=t$, with $\varphi'$ integral, and with the pointwise norm of the reading
  of $\varphi$ bounded on the frames of every component at every good embedding.
\item[(S3)] (Products.) Let $f$ be a cusp form of weight $\kappa(\lambda,0)$ on $Y$, for a family $\lambda$, and
  let $f'\in R$ be a rational section of some weight such that $ff'$ is a rational section of a dominant weight
  $\kappa'$. Then $ff'\in S_Y(\kappa')$.
\item[(S4)] (Integral multiples.) Let $w\ge1$ be an integer with $u^w=1$ for every root of unity $u\in E$, and
  let $\kappa'=(\alpha',n')$ be a dominant weight with $n'\ge8fr+4$ such that $w$ divides every entry of the
  weights $\lambda(\kappa')$ of the readings at the good embeddings (Situation~\ref{cv-cusp-setting}). Then every
  $F\in S_Y(\kappa')$ has an integral multiple $sF$ with $s\in\mathcal O_{F_Y}$ nonzero.
\item[(S5)] (Compactified flag spaces.) Let $D\ge1$ be as in (iv). For every connected component $Z$ of $Y$
  there are a geometrically connected projective variety $\bar Z$ over $F_Y$, smooth of dimension $D$, with a
  dense open subset $Z^\circ$, and for every family $\lambda$ a line bundle $\mathcal L_\lambda$ on $\bar Z$ and an
  injective $F_Y$-linear map $c\mapsto c^\circ$ from the cusp forms of weight $\kappa(\lambda,0)$ supported on $Z$
  to the sections of $\mathcal L_\lambda$ over $Z^\circ$. Every $c^\circ$ extends to a section of $\mathcal L_\lambda$
  over $\bar Z$ that vanishes on $\bar Z\smallsetminus Z^\circ$, and every such section restricts to some
  $c^\circ$.
\item[(S6)] (No base points.) In (S5), let $\lambda$ be a family as in (iv). Then the sections of
  $\mathcal L_\lambda$ over $\bar Z$ whose restrictions to $Z^\circ$ are of the form $c^\circ$ have no common zero on
  $Z^\circ$.
\end{enumerate}

Proof of (i). Let $F\in S_Y(\kappa)$ be nonzero.

A nonzero reading. Some component $Z$ at $\sigma_0$ and some frame $x_1$ of $Z$ have $F_{\sigma_0,Z}(x_1)\ne0$.
Otherwise, let $y$ be a framed point of $Y$ over a reduced finitely generated $F_Y$-algebra $S$, and
$\psi\colon S\to\mathbf C$ a homomorphism of $F_Y$-algebras. The framed point $\psi(y)$ over $\mathbf C$ is
triangularly equivalent to the framed point at a frame $x$ of some component $Z$, by the compatibility
conditions recalled in the proof of Lemma~\ref{cv-cover-uniformization}. So by the transformation law
$\psi(F(y))=c\,F_{\sigma_0,Z}(x)=0$ for some $c\in\mathbf C$. By the Nullstellensatz, an element of $S$ killed by
every such $\psi$ is $0$. So $F$ vanishes at every framed point, and $F=0$.

A CM frame. The reading $F_{\sigma_0,Z}$ is a cusp form, so $|F_{\sigma_0,Z}|$ is bounded.
Lemma~\ref{cv-volume-forms-cm-frame}, with $L'=L'_Z$, gives $\ell>0$, a frame $x$ of $Z$ with
$F_{\sigma_0,Z}(x)\ne0$, and a basis $v_1,\dots,v_r$ of $E^r$ in $L'_Z$ with (1) to (4) of
Lemma~\ref{cv-cusp-frame}.

A point over a number field. Let $K$, $\sigma_1$ and $y_0$ be as in (S1). The framed point $w$ of $Y$ at $x$ has
$F(w)=F_{\sigma_0,Z}(x)\ne0$, and its point is the image of the point of $y_0$ under $\sigma_1$. By
Lemma~\ref{cv-volume-forms-frame-nonzero}, there is a framed point $y$ of $Y$ over $K$ with the point of $y_0$ and
$F(y)\ne0$.

The denominator. Write $\kappa=(\alpha,n)$, let $k_\tau=\sum_i\max(0,-\alpha_{\tau,i})$, and let $t$ be given by
(S1) for $\kappa$ and these $k_\tau$. Lemma~\ref{cv-volume-forms-denominator} gives $N\in K^\times$ with $NG(y)$
integral over $\mathbf Z$ for every integral $G\in S_Y(\kappa)$.

Proof of (ii). Let $f\in S_Y(\kappa(\lambda,0))$ with $\lambda$ regular. The cover $Y$ is a deep cover of a moduli
scheme whose level is prime to $15$, and $[F:\mathbf Q]\ge8$ (derivation of (Y1) to (Y7)), so
Lemma~\ref{cv-volume-generators-key-inverse} gives $u$ of weight $\kappa(0,1)$, an integral $u'$ and
$t_1\in F_Y^\times$ with $uu'=t_1$. Let $N$ be as in (S2), put $M_0=8fr+4$ and $k=2M_0+N$, and let $w$ be given by
Lemma~\ref{cv-volume-generators-exponents} for $\lambda$, $N$, $M=M_0$ and $k$, so that
$\kappa'=\kappa((k+1)\lambda+\lambda_w,M_0)$ is a dominant weight $(\alpha',2M_0)$. Let $\varphi$, $\varphi'$ and $t_2$ be
given by (S2) for $w$, and put $g=u^{M_0}\varphi$, $h=u'^{M_0}\varphi'$ and $t=t_1^{M_0}t_2$. Then $f^{k+1}g$ is a
rational section of weight $(k+1)\kappa(\lambda,0)+M_0\kappa(0,1)+\kappa(\lambda_w,0)=\kappa'$, so
$f^{k+1}g\in S_Y(\kappa')$ by (S3), applied to $f$ and $f'=f^kg$. Moreover $gh=(uu')^{M_0}\varphi\varphi'=t$, the
valuation of $h$ at an integral point is $M_0v(u')+v(\varphi')\ge0$, and $n'=2M_0\ge8fr+4$.

Proof of (iii). Let $w$ be the number of roots of unity in $E$, so that $u^w=1$ for every root of unity $u\in E$,
and let $F\in S_Y(\kappa')$. By Lemma~\ref{cv-cover-level-automorphisms}(4), $F^w\in S_Y(w\kappa')$, and
$w\kappa'=(w\alpha',wn')$ is a dominant weight with $wn'\ge8fr+4$. The entries of the weights $\lambda(w\kappa')$
of the readings are $\pm(w\alpha'_{\tau',i}+wn'/2)$, which $w$ divides since $n'$ is even. By (S4) there is a
nonzero $s\in\mathcal O_{F_Y}$ with $sF^w$ integral. Then $sF$ is integral: at an integral point, if $v(sF)<0$
then $v\big((sF)^w\big)=wv(sF)<0$, while $(sF)^w=s^{w-1}(sF^w)$ has nonnegative valuation, as $s$ is integral.

Proof of (iv). Let $Z$ be a connected component and $\lambda$ as in (iv), and let $\bar Z$, $Z^\circ$,
$\mathcal L_\lambda$ and $c\mapsto c^\circ$ be as in (S5). As $\bar Z$ is geometrically connected, it is not empty,
so its dense open subset $Z^\circ$ contains a point $y$. By (S6) there is a section of $\mathcal L_\lambda$ over
$\bar Z$ that does not vanish at $y$ and restricts to $c^\circ$ on $Z^\circ$ for a cusp form $c$ of weight
$\kappa(\lambda,0)$ supported on $Z$. Then $c^\circ(y)\ne0$, so $c\ne0$.
% Lean: (i) is ShiftFG.exists_good_cmPt_frame_ne_zero (for m = 0 the same steps in
% CovDetectDom.exists_cmPoint_frame_ne_zero_of_dominant) with GlobalTheory.T.exists_frame_denominator_dominant_weight_power_cm
% (and _cm for m = 0). The nonzero reading is exists_rdU_ne_zero_of_mem_secZ (test_value_eq_const_mul_reading_secZ,
% Formalism.L2.Pj.eq_zero_of_cpts); L'_Z is CovData.exists_dataX_good, which also gives a nonzero fractional ideal
% AA with Lambda = AA L'_Z, used by (S1); the bound is abelian2.holSub_SS; the CM frame is CMDetect.exists_cm_frame;
% the point over K is FrameDenseZ.exists_frame_evR_ne_zero_of_complex; the denominator is
% LatFull.CMInt.exists_frame_denom_dominant_pow. (ii) is WtMatchBar.exists_pow_mul_mem_cuspForms_of_weight_match
% with M = 4 [E:Q] r + 4 and k = max K (2M + N), where K = 0 is the bound of exists_pow_mul_unit_mem_cusp; u, u'
% are GlobalTheory.T.inverse_invariant_integral (Lemma cv-volume-generators-key-inverse), w is
% exists_dominant_exponents_matching_weight_bar_gap. (iii) is LargeCover.exists_smul_nonneg_all: cusp_mul_of_M5
% for the powers, exists_pow_root_eq_one (the torsion order of the units of E), exists_smul_nonneg_of_dvd for (S4),
% and IntCuspLat.nonneg_of_pow. (iv) is LowNonzero.exists_cusp_value_ne_zero_low_level: dense_U0, ne_zero_of_noBase
% for the section, and ofSpace, ofSpace_injective for the cusp form.

(Y3) Let $H$ be as in (Y2). For a form $\varphi$ of weight $(\alpha,n)$ on $Y$, a good $\sigma$ and a frame $x$,
write $|\varphi|_\xi(x)=(2\pi\xi)^{-|\alpha|/2}|\varphi_\sigma|(x)$ for the pointwise norm of its reading with
the factor of the Petersson norms; it is multiplicative. Three inputs are used.
\begin{enumerate}
\item[(a)] (Sup norms and ranks.) There is a constant $C'$ of the family with
  $|\varphi|_\xi(x)^2\le e^{C'z^3}N_{Z,\sigma}(\varphi)^2$, hence $|\varphi|_\xi(x)\le e^{C'z^3}N_{Z,\sigma}(\varphi)$,
  for every cusp form $\varphi$ of weight $c_0h$ supported on a component $Z$, every good $\sigma$ and every
  frame $x$ of $Z$, and with $\log\operatorname{rk}_Z(c_0h)\le C'z^3$ for every $Z$. Indeed, let $w_0$ be a number
  larger than the spread of the weight. Lemma~\ref{cv-cusp-sup-bound} for $Y$, applied to the family that is
  $\varphi_\sigma$ on $Z$ and $0$ elsewhere, gives
  $|\varphi|_\xi(x)^2\le C_{w_0}m_{Z,\sigma}N_{Z,\sigma}(\varphi)^2$ with $m_{Z,\sigma}\le1$.
  Lemma~\ref{cv-cusp-bergman}, applied as in Steps~3 and~4 of the proof of
  Proposition~\ref{cv-cusp-rank-bound} to the readings at a good $\sigma$ of the forms on $Z$, with this bound,
  gives $\operatorname{rk}_Z(c_0h)\le C_{w_0}$. Finally $1\le C_{w_0}$ and $\log C_{w_0}\le C'z^3$ with $C'=4000$,
  a computation from the definition of $C_{w_0}$, $\log M_Y\le30x^\varepsilon$ and the size conditions. The
  second form of the sup bound follows from the first, as $e^{C'z^3}\ge1$.
% Lean: (a) is UpperTheory.Lit.Tam.tamagawa_number_formula, whose constant 4000 serves all three parts: the
% Bergman bound BergLe (tamagawaPeterssonLambdaSq_nonneg, with w = 2[E:Q] size + 1), the sup bound SupBdWt
% (SupBdWt.of_bergOn) and log rkC <= C z^3 (GlobalTheory.T.cuspFormRank_le_bergmanBound_goodCover
% = RankBergman.rkC_le_of_bergLe); the computation is UpperTheory.Lit.Tam.supConst2Bounds. The bound
% |phi| <= beta N from |phi|^2 <= beta N^2 and beta >= 1 is the step in GlobalTheory.ShortDet.exists_short.
\item[(b)] (Short generators.) There is a constant $C''$ of the family such that for every $Z$ the integral
  forms of weight $c_0h$ on $Z$, which span the forms on $Z$ by the arguments for the integral forms of
  (Y1) with $\lambda=c_0\mathbf h$ and $m=0$, are generated over $\mathbf Z$ by $n_Z\le e^{C''z^3}$ of them, each with
  $N_{Z,\sigma}\le e^{2C''z^9}$ at every $\sigma$.
\item[(c)] (A separating form.) There is $u_0\in S_Y(c_0h)$ whose group determinant over $H$ is not a zero
  divisor in $R$. This is the form $u_0$ of the proof of Proposition~\ref{cv-point-conditions}, for the
  group $S=\Gamma_P$ there, which is $H$: both are the level automorphisms of $Y$ whose reduction modulo
  $\mathfrak M$ is upper unitriangular at every level place. Its group determinant there is the one here:
  the translation there is $\tau(\gamma)=(\gamma^{-1})^*$, so the matrix $\big(\tau(\gamma\gamma')u_0\big)$
  becomes $\big((\gamma\gamma')^*u_0\big)$ after both indices are replaced by their inverses and the matrix
  is transposed, which does not change the determinant. It is obtained there from
  Lemmas~\ref{cv-point-group-determinant} and~\ref{cv-point-regular-section} and the separation (i) of
  that proof.
\end{enumerate}
Proof of (b). It uses one more fact about $Y$, which is part of the open step below.
\begin{enumerate}
\item[(S7)] (Short spanning families.) There is a constant $C_2\ge0$ of the family such that for every
  component $Z$ the integral forms of weight $c_0h$ on $Z$ contain a family that spans the forms of weight
  $c_0h$ on $Z$ over $F_Y$ and whose members satisfy $N_{Z,\sigma}\le\exp\big(C_2(z^3S+z^5)\big)$ at every
  $\sigma$, where $S=\sum_{\tau,i}|c_0\mathbf h_{\tau,i}|$.
\end{enumerate}
Fix $Z$, let $L$ be the module of integral forms of weight $c_0h$ on $Z$, of rank $\varrho=\operatorname{rk}_Z(c_0h)$,
and put $d=[F_Y:\mathbf Q]$. By (a), $\log\varrho\le C'z^3$. The norms $N_{Z,\sigma}$ satisfy $N_{Z,\sigma}(a\varphi)=|\sigma(a)|N_{Z,\sigma}(\varphi)$ and the
triangle inequality, and $L$ is finitely generated and spans the forms on $Z$.
Lemma~\ref{cv-volume-generators-short}, with the family of (S7) and $A=\exp(C_2(z^3S+z^5))$, gives generators
$w_1,\dots,w_{n_Z}$ of $L$ over $\mathbf Z$ with $n_Z\le2d\varrho$ and
$N_{Z,\sigma}(w_i)\le\max(1,d^4\varrho\operatorname{rd}_{F_Y}A)$. By Definition~\ref{cv-deep-cover}(2),
$\log d\le391z^2$ and $\log\operatorname{rd}_{F_Y}\le391z^2$; and $S=c_0fq\le22f^3r^5\le z^5$, as $c_0\le22f^2r^4$
and $f\ge8$. With $C''=2000+C'+C_2$ and $z\ge1$,
\[
 \log n_Z\le\log2+391z^2+C'z^3\le C''z^3,\qquad
 \log\big(d^4\varrho\operatorname{rd}_{F_Y}A\big)\le1955z^2+C'z^3+C_2(z^3S+z^5)\le C''(z^3S+z^5)\le2C''z^9.
\]
So $N_{Z,\sigma}(w_i)\le e^{2C''z^9}$.
% Lean: (b) is ShortGenComp.short_generator_comparison, with C_3 = 2000 + C_b + C_2 (the text's C'' = 2000 + C' +
% C_2, as C_b is the constant C' of tamagawa_number_formula): C_2 and the family s are
% GlobalTheory.T.short_span_composition_closed (S7), the rank bound log rkC <= C_b z^3 is the third part of
% UpperTheory.Lit.Tam.tamagawa_number_formula (RankBergman.rkC_le_of_bergLe from BergLe, the bound behind (a)),
% the span of the integral forms is covering_domain_lattice_regularity, and the generators are
% GlobalTheory.T.exists_short_generators_of_short_spanning_family (Lemma cv-volume-generators-short); the
% hypotheses on c_0 h (gap at least 8fr + 10, c_0 <= 22 f^2 r^4) are cZeroN_isLeast and cZero_spec in
% GlobalTheory.ShortDet.exists_short_integral_vector.
An element of $R$ is not a zero divisor if and only if its restriction to the function field of every
component is nonzero. Let $w_1,\dots,w_N$ be the generators of (b), for all components together; they span
$S_Y(c_0h)$, so $u_0=\sum_ia_iw_i$ with $a_i\in F_Y$. Replace the coefficients one at a time by integers
$0\le m_i\le|H|^2$, keeping the group determinant a non-zero-divisor. Let $v'$ be the sum of the other
terms. Since $w_i$ is supported on one component $W$, the matrix $((\gamma\gamma')^*w_i)_{\gamma,\gamma'\in H}$
vanishes on every component other than the at most $|H|$ components $\gamma W$, $\gamma\in H$. On each of
these, the group determinant of $v'+tw_i$ is a polynomial in $t$ of degree at most $|H|$ over the function
field, nonzero at the current coefficient, so it has at most $|H|$ roots; on the other components it does
not depend on $t$. So some integer $0\le m\le|H|^2$ is not a root on any component. The resulting
$u=\sum_im_iw_i$ is integral, its group determinant is not a zero divisor, and at every frame $x$ of a
component $Z$, by (a) and (b),
\[
 |u|_\xi(x)\le|H|^2\,n_Z\,e^{C'z^3}e^{2C''z^9}\le e^{(240+C'+3C'')z^9},
\]
using $\log|H|\le120f^2r^2x^\varepsilon\le120z^2$ from (Y2) and $z\ge1$. For a cusp form $f$ of even weight and
$\gamma\in\Gamma_Y$, Lemma~\ref{cv-cover-level-automorphisms}(2) and~(3) give
$N_\sigma(f\cdot\gamma^*u)=N_\sigma\big((\gamma^{-1})^*f\cdot u\big)$, which is at most $\sup_x|u|_\xi(x)$ times
$N_\sigma((\gamma^{-1})^*f)=N_\sigma(f)$, as $|\cdot|_\xi$ is multiplicative. This is (Y3) with $v=u$, for any
$C_Y\ge240+C'+3C''$.
% Lean: (Y3) is GlobalTheory.ShortDet.exists_short_integral_vector, with C = C_3 + 240 + C_1 + 2 C_3: (a) is
% UpperTheory.Lit.Tam.tamagawa_number_formula (SupBdWt from BergLe through tamagawaPeterssonLambdaSq_nonneg,
% with the numerical bound UpperTheory.Lit.Tam.supConst2Bounds), (b) is
% GlobalTheory.T.short_generator_norm_comparison (IsShortGen; log of the norm bound at most
% C(z^3 size + z^5), with size(c_0 h) <= (fr)^5 from c_0 <= 22 f^2 r^4) with the span from
% covering_domain_lattice_regularity, (c) is GlobalTheory.T.orbits_separated_by_regular_cusp; the integer
% coefficients are GlobalTheory.ShortDet.exists_short (exists_nat_coeff, exists_nat_step with
% card_le_of_det_eq_zero, nonZeroDiv_iff), and the product bound is GlobalTheory.ShortDet.normW_mul_pull_le
% (petAt_mul_le). The size condition 2 C_Y + 200 <= z is from hostN_alt.

(Y4) Alp\"oge obtains the forms from theta series attached to lattices. Write $S=\operatorname{size}(\kappa_j)$,
so $S\ge1$, and recall that $n$ is the last entry of $\kappa_j$. A \emph{theta datum} $(x_0,\Lambda,N,P)$ consists
of a family $x_0=(x_{0,1},\dots,x_{0,n})$ of elements of the level module of $Y$ that generates it over
$\mathcal O_E$; an $\mathcal O_E$-lattice $\Lambda\subset E^n$ that is its own dual for the Hermitian form
$\sum_b\bar v_bw_b$ and satisfies $k\Lambda\subset\mathcal O_E^n\subset k^{-1}\Lambda$ for an integer $k$ prime to
$M_Y$; an odd integer $N$ divisible by $M_Y$ with $N\mathcal O_E^n\subset\Lambda'=\Lambda\cap\mathcal O_E^n$; and a
polynomial $P$ over $F_Y$, harmonic of weight $\kappa_j$, in variables $X_{\tau,i,b}$, that lies in the ring
generated over the integers of $F_Y$ by the linear forms $\sum_b\tau(v_b)X_{\tau,i,b}$ with $v\in\Lambda$. Its
\emph{theta series} is the rational section of weight $\kappa_j$ whose value at a framed point $y$ is a finite sum,
over $n$-tuples of points of order dividing $N^2$ of $A_y$ attached to $\Lambda$ and over the level module, of
products of derivatives, given by the monomials of $P$, of translates of the theta function of $y$, with
coefficients given by $x_0$. Put $\|P\|=\big(\sum_\sigma\sum_\mu|\sigma(P_\mu)|^2\big)^{1/2}$, over the embeddings
$\sigma$ of $F_Y$ and the coefficients $P_\mu$ of $P$, and for an $\mathcal O_E$-lattice $L\subset E^n$ and an
embedding $\tau$ of $E$ let $I_\tau(L)$ be the fractional ideal of $F_Y$ generated by the images under $\tau$ of
the determinants of $n\times n$ matrices with columns in $L$. The derivation of (Y4) uses the following facts,
for constants $C_A,C_B,C_C\ge0$ of the family; (T1), (T2) and (T3) are the open steps below.
\begin{enumerate}
\item[(T1)] (Spanning.) The cusp forms of weight $\kappa_j$ that are theta series of theta data with
  $\|P\|\le e^{C_Az^2S}$, and with $k$ also prime to a fixed integer of the datum, span $S_Y(\kappa_j)$ over $F_Y$.
\item[(T2)] (Integrality.) There is an integer $N_a\ge1$ with $\log N_a\le C_Cz^2S$ such that, for every cusp form
  $\theta$ of weight $\kappa_j$ that is the theta series of a theta datum as in (T1), without the bound on $\|P\|$,
  and every $s\in F_Y$ with
  $s\prod_\tau I_\tau(\Lambda')^{m_\tau}\subset\prod_\tau I_\tau(\Lambda)^{m_\tau}$, the form $N_as\theta$ is integral.
\item[(T3)] (Size.) For every cusp form $\theta$ of weight $\kappa_j$ that is the theta series of a theta datum
  with $\|P\|\le\rho$, the pointwise norm satisfies $|\theta|_\xi\le[\Lambda:\Lambda']^{(p+q)/2}\rho\,e^{C_Bz^2S}$ at
  every frame at every good embedding.
\item[(T4)] (Small scalar.) For every lattice $\Lambda$ as in a theta datum there is $s\in F_Y^\times$ with
  $s\prod_\tau I_\tau(\Lambda')^{m_\tau}\subset\prod_\tau I_\tau(\Lambda)^{m_\tau}$ and
  $|\sigma(s)|^2[\Lambda:\Lambda']^{p+q}\le\exp\big(2\log[F_Y:\mathbf Q]+2\log\operatorname{rd}_{F_Y}\big)$ for every
  $\sigma$. This is Lemma~\ref{cv-volume-theta-small-scalar} for the PEL datum of $W_j$ at the deep level, with
  coefficient field $F_Y$, the CM type $\Phi$ of the embeddings that send $\sqrt{-d}$ to the square root of $-d$
  fixed by $\mathcal D$, and the signature $(p,q)$ of $W_j$.
\end{enumerate}
The scaled bound. Let $\theta$ be as in (T3), and let $s\in F_Y$ and $B\ge0$ satisfy
$|\sigma(s)|^2[\Lambda:\Lambda']^{p+q}\le B^2$ for every $\sigma$. At a frame at a good embedding $\sigma$ the reading of
$s\theta$ is $\sigma(s)$ times that of $\theta$, so $|s\theta|_\xi=|\sigma(s)|\,|\theta|_\xi\le
|\sigma(s)|[\Lambda:\Lambda']^{(p+q)/2}\rho\,e^{C_Bz^2S}\le B\rho\,e^{C_Bz^2S}$.

Let $\Theta$ be the set of forms of (T1), and let $\theta\in\Theta$, with theta datum $(x_0,\Lambda,N,P)$. Choose $s$ by
(T4). By (T2), $N_as\theta$ is integral. As $\log[F_Y:\mathbf Q]+\log\operatorname{rd}_{F_Y}\le391z^2\le391z^2S$ by the
size of the cover, the scaled bound with $B=e^{391z^2S}$ and $\rho=e^{C_Az^2S}$ gives
$|s\theta|_\xi\le e^{(C_A+C_B+391)z^2S}$ at every frame. The Petersson norm $N_\sigma$ is the square root of a mean,
for a probability measure, of the squares of the pointwise norms, so
\[
 N_\sigma(N_as\theta)\le N_a\,e^{(C_A+C_B+391)z^2S}\le e^{(C_A+C_B+C_C+391)z^2S}
\]
for every $\sigma$. Each $\theta\in\Theta$ is $(N_as)^{-1}$ times the integral form $N_as\theta$, so by (T1) these forms
span $S_Y(\kappa_j)$. This is (Y4) once $C_Y\ge C_A+C_B+C_C+391$. The hypotheses of (T1) to (T3) on the weight,
$n\ge8fr+4$, that the entries of $\kappa_j$ in the first $p$ positions exceed $n/2$, and
$\log S\le20x^\varepsilon$, follow from the size conditions.
% Lean: GlobalTheory.T.exists_good_cover_at_minimal_level = GlobalTheory.CovGmin.minimal_level_cover_is_good,
% with C = CA + CB + CC + 391 (exp_consts). A theta datum is (x0, La, N, P) of thFam: x0 : Fin n -> LevQc with span
% top, La : GenLat E n with Unimod, TrivAt (badC Y) and TrivAt M (badC Y = QS q d M_Y; the text's "fixed integer of
% the datum"), LevOK N M, P in harmL La a (HHa: harmK, RowInv, RowDeg a; symO La); its theta series is the cusp
% form y with thLaValAt nab (phiHat (indPhi x0)) La N P at every framed point over a test ring (ShortValAt with
% beta = 1: the sum over q of w q times sum_i sum_m P_m prod_b nab^(m_b) F q i b, with F of IsUUOver at the
% N^2-torsion points z i of InZh), nab the canonical connection of CanonConn.exists_isNab_uncond. ||P|| is
% ThetaWeightData.polyNorm, Lambda' is La.Ln and I_tau is idealAt. (T1) is GlobalTheory.CovGmin.span_thFam (CA = 20000;
% PrismSpan.cusp_sub_span_short with harm_short), (T2) GlobalTheory.T.exists_bounded_integer_goodCover_thetaSeries_integral
% (CC), (T3) GlobalTheory.RawThetaSup.homN_raw_le (CB = 9000), and the scaled bound is
% GlobalTheory.RawThetaSup.homN_smul_raw_le (homN_smul for the factor |sigma(s)| at goodOf sigma), (T4)
% GlobalTheory.T.exists_small_scalar_reflexIdeal_selfDualLattice (exp (fieldSz F_Y)^2). The span is
% span_eq_top_of_multiples and the norms nrm_le_of_eq_smul (petW_le_of_homN_le); the size of F_Y is hG.size with
% o <= 391 z^2. The constant C satisfies C + 45 <= z by family_sizes_eventually_large.
\openstep{The spanning (T1) of $S_Y(\kappa_j)$ by the theta series of theta data with $\|P\|\le e^{C_Az^2S}$ is
not yet written out.}
% Lean: (T1) is GlobalTheory.CovGmin.span_thFam (C_A = 20000): GlobalTheory.CovGmin.PrismSpan.cusp_sub_span_short at
% the good inclusion sigma_0, with the principal level structures of GlobalTheory.LevGeo.hasPrinc_laws, ChiOne for the
% dominant weights (hG.chiTriv), badC Y > 0, and GlobalTheory.CovGmin.PrismSpan.harm_short for the bound on ||P||.
% Size: about 150,000 lines counted with everything the proof uses and the book does not yet cite; about 58,000 of
% them are not used by (T2) and (T3), and the three together use about 265,000.
\openstep{The integrality (T2) of the multiples $N_as\theta$ of theta series is not yet written out.}
% Lean: (T2) is GlobalTheory.CovInt.integral_smul_thLa: N_a = intZ^2 with GlobalTheory.T.log_shortThetaSeriesDenominatorBound_le,
% the short rule GlobalTheory.RawSpan.isShortRule_of_raw, and at each prime of F_Y the case of a prime dividing M_Y
% (GlobalTheory.T.intersectedLattice_detIdeal_integral_trivAt_prime, with GlobalTheory.CovInt.SqLev.smul_intZ_sq_nonneg_of_sq_dvd
% when its square divides M_Y) and the other primes. Size: about 170,000 lines, counted as for (T1); about 107,000
% of them are not used by (T1) and (T3).
\openstep{The bound (T3) for the pointwise norms of theta series is not yet written out.}
% Lean: (T3) is GlobalTheory.RawThetaSup.homN_raw_le (C_B = 9000). Size: about 89,000 lines, counted as for (T1);
% about 7,000 of them are not used by (T1) and (T2).

(Y5) Write $\kappa_a=\kappa_j-u_W$, the weight $\kappa(a,n/2)$ of
Definition~\ref{cv-volume-twisting-section}, given by the anchor weights $a_\tau$ and the entry $n$. For
every component $Z$ there is a twisting section $t_Z$ of weight $\kappa_a$ on $Z$
(Definition~\ref{cv-volume-twisting-section}), with section $s_Z$, numbers $c_{Z,\sigma}$ and function
$l_Z$, such that $l_Z(w)$ is an integer when $w$ lies over no bad prime of $Y$, that is, no prime dividing
$2d|d_E|M_Y$, and
\[
 \deg t_Z=\sum_wl_Z(w)\frac{\log Nw}{[F_Y:\mathbf Q]}-\frac1{[F_Y:\mathbf Q]}\sum_\sigma\log c_{Z,\sigma}
 \ge\Theta+\lambda_Y.
\]
Indeed, $Y$ satisfies the hypotheses of Lemmas~\ref{cv-volume-twisting-flat}
and~\ref{cv-volume-twisting-key}: $p+q=r\ge1$; its uniformizations at the good embeddings satisfy the
compatibility conditions, with the normalization $\xi$; $F_Y$ is algebraically closed in the function field
of every component (Definition~\ref{cv-deep-cover}(3)); its level has full level $1$ and prime-to-$\ell$
part at least $3$ for every prime $\ell$; by Theorem~\ref{cv-flat-degree}, flat sections of every type
exist on every component, all of one degree $-c$, so that $c_\tau=c$ for every $\tau$; and every component
has, above each finite place over no bad prime, an integral point at which the valuation of every rational
section is an integer or $+\infty$, by Lemma~\ref{cv-volume-twisting-unramified} for the deep level and
$F'=F_Y$. By Lemma~\ref{cv-volume-twisting-key}, $Y$ has a key section and $Z$
has a twisting section $t_2$ of weight $\kappa(0,1)$, with function $l_2$ integral at the places over no bad
prime and $\deg t_2\ge-fr(\frac12\log2\pi+\log\operatorname{rd}_E)$, since $\frac12[E:\mathbf Q](p+q)=fr$.
By Lemma~\ref{cv-volume-twisting-points}, $Z$ has frames at every good embedding and integral points above
every finite place, so Lemma~\ref{cv-volume-twisting-flat}, with $w=a$ and the set of bad primes, gives a
twisting section $t_1$ of weight $\kappa(a,0)$, with function $l_1$ integral at the places over no bad prime
and $\deg t_1\ge\sum_\tau a_\tau c_\tau-q\log\operatorname{rd}_E\sum_\tau|a_\tau|$. By
Lemma~\ref{cv-volume-twisting-products}, $t_Z=t_1t_2^{n/2}$ is a twisting section of weight
$\kappa(a,0)+\frac n2\kappa(0,1)=\kappa_a$, with function $l_Z=l_1+\frac n2l_2$, integral at the places over
no bad prime, and with $\deg t_Z=\deg t_1+\frac n2\deg t_2\ge\Theta+\lambda_Y$.

Let $\kappa$ be a weight for which the conclusions of (Y1) hold. Multiplication by $t_Z$ is a bijection from
the forms of weight $\kappa$ on $Z$ to those of weight $\kappa+\kappa_a$, with inverse the multiplication by the
inverse of $t_Z$. The valuations of $t_Z$ at integral points are bounded above and below, so a nonzero
multiple of the image of the integral forms consists of integral forms and conversely, and
$\operatorname{rk}_Z(\kappa+\kappa_a)=\operatorname{rk}_Z(\kappa)$. Rounding the numbers $l_Z(w)$, which are integers
outside the bad primes, gives $a\in F_Y^\times$ and a nonzero ideal $\mathfrak a\subset\mathcal O_{F_Y}$ with
$v_w(ba)+l_Z(w)\ge0$ for every nonzero $b\in\mathfrak a$ and every $w$, and
\[
 \sum_wl_Z(w)\frac{\log Nw}{[F_Y:\mathbf Q]}-\lambda_Y\le-\frac1{[F_Y:\mathbf Q]}\sum_\sigma\log|\sigma(a)|
 -\frac{\log N\mathfrak a}{[F_Y:\mathbf Q]}.
\]
The map $\psi(f)=a\,t_Zf$ multiplies $N_{Z,\sigma}$ by $|\sigma(a)|c_{Z,\sigma}$, and $\mathfrak a\,\psi$ maps the
integral forms of weight $\kappa$ on $Z$ into the integral forms of weight $\kappa+\kappa_a$, a lattice of the
same rank. Comparing covolumes,
\[
 \deg_Z(\kappa)+\operatorname{rk}_Z(\kappa)\Big(-\frac1{[F_Y:\mathbf Q]}\sum_\sigma\log\big(|\sigma(a)|c_{Z,\sigma}\big)
 -\frac{\log N\mathfrak a}{[F_Y:\mathbf Q]}\Big)\le\deg_Z(\kappa+\kappa_a),
\]
and with the rounding, $\deg_Z(\kappa+\kappa_a)\ge\deg_Z(\kappa)+\operatorname{rk}_Z(\kappa)(\deg t_Z-\lambda_Y)$. Summing over the
components with $(*)$, whose masses $m_{Z,\sigma}$ do not depend on the weight, gives
$\operatorname{rk}(\kappa+\kappa_a)=\operatorname{rk}(\kappa)$ and
$\deg(\kappa)+\operatorname{rk}(\kappa)\Theta\le\deg(\kappa+\kappa_a)$. For $\kappa=u_-$ these are the first two
claims, as $u_-+\kappa_a=\kappa^-$; for $\kappa=u_W+sc_0h$ the equality of ranks on each component, with the
splitting of ranks, gives $\operatorname{rk}(\kappa_j+sc_0h)=\sum_Z\operatorname{rk}_Z(u_W+sc_0h)$. The bad primes
of $Y$ are the primes dividing $2d|d_E|M_Y$, so $\lambda_Y\le\log(2d|d_E|M_Y)\le40x^\varepsilon$ by the size
conditions.
% Lean: (Y5) is GlobalTheory.T.twisted_flat_bundle_ranks (TwistRank.flat_twist_rank_eq: exists_twSec_par,
% rankOn_twist, GlobalTheory.T.exists_twisted_section_with_deep_root, supLattice_rank_eq_sum_component_ranks)
% and GlobalTheory.T.flatLine_thetaRoot_twist_twistingSection (FlatRootTwist.flat_root_twist: the twisting
% section TwSec, with sig, inv, cInf, lo and degLo = lineDeg cInf lo, is t_1 t_2^m from
% GlobalTheory.FlatTwisting.twSec_of_flat and GlobalTheory.KeyRootTw.exists_twSec_keyRoot through twSec_mul_pow,
% with IntOff badPrimes lo; then twist_goodOn (Mul.twSec_mulEquiv_integral, exists_twist_scalar,
% lattice_transfer), degOn_add_rank_mul_degLo_le (the rounding GlobalTheory.T.exists_ideal_rounding_realExponents_badPrimes
% = Rounding.exists_rounding, a self-contained computation, and the covolume comparison
% GlobalTheory.T.deg_add_le_deg_of_ideal_smul_le_alt) and
% GlobalTheory.Proof.TwistedWeightDegree.degW_twist, which sums with Adds), with w the anchor, m = n/2, c_tau = c1, badPrimes the prime
% factors of 2 d |disc E| M', and lambda_Y <= 40 L from GlobalTheory.CovTop.FlatVal.flat_val_alt
% (BadLog.logSum_le_log, log_level_le).

(Y6) The weight $u_1$, the integer $\ell$ and the numbers $\rho$, $\zeta$ and $\beta$, with their bounds, come
from the size conditions. Part (e) is $(*)$ for $\kappa=u_-$ with the splitting of ranks. Parts (a) to (d)
bound the geometric degrees $\delta_Z$ of the components, compare the ranks on the components with
$\beta\delta_Z$, bound the degrees $\deg_Z(u_-)$ below by an arithmetic Hilbert--Samuel estimate in terms of
the heights $h_Z$, and bound these heights below through Theorem~\ref{cv-volume-theorem}. That theorem
is applied to a global model of the moduli scheme $Y$ of the deep cover, whose level has full level $1$
and prime-to-$\ell$ part at least $3$ at every prime $\ell$, with the tower and good origin of
Theorem~\ref{cv-volume-good-tower}; $F_Y$ is algebraically closed in the function field of every
connected component of $Y\otimes F_Y$ by Definition~\ref{cv-deep-cover}(3).
% Lean: (Y6): u_1, ell = lenT, zeta = zed, rho = rhoB, beta = binT from SizeAt and
% GlobalTheory.T.size_bounds_for_family_alt; (a) DegAt from GlobalTheory.T.coverDegree_of_grothendieck_deformation,
% (b) ShareAt from GlobalTheory.T.share_of_deep_good_size_flat_deg_cover_degree and the jet bound inside
% GlobalTheory.CovTop.HodgeSteps.rk_le_alt_alt (GlobalTheory.T.jet_rank_complex_le,
% componentImageDegree_between_of_proportionality, exists_jet_order), (c) HSAt from
% GlobalTheory.T.coverHeightSum_of_flatDeg_coverDegree_share, (d) HeightAt from
% SelfIntersectionVolumeEstimates.height_lower_bound_general_alt (with DeepStR from the volume theorem), (e)
% SumAt from GlobalTheory.T.sum_component_degrees_div_le_latticeSlope = GlobalTheory.CovSum.sum_deg_div_le_mu.
\openstep{The facts (S1) to (S7) about $Y$, used in the proofs of inputs (i) to (iv) of (Y1) and of the
short generators (b) of (Y3), are not yet written out: the CM points with bounded denominators (S1), the
flat units (S2), the cusp forms obtained as products (S3), the integral multiples at weights divisible by
the number of roots of unity of $E$ (S4), the compactified flag spaces (S5) and their base points (S6), and
the short spanning families of weight $c_0h$ (S7).}
\openstep{Part (a) of (Y6), the bounds $\delta'\le\delta_Z\le\delta'/(1-\rho)$ for the geometric degrees of the
components, is not yet written out.}
% Lean: DegAt from GlobalTheory.CovTop.DegProof.covering_degree_eq_alt: the proportionality principle on the
% compactification (GlobalTheory.T.proportionality_principle_compactification), the top Chern number as a volume
% (GlobalTheory.T.compactification_top_chern_eq_volume), the birationality of the map given by the forms of weight u_1
% (GlobalTheory.T.birOn_wtOf_of_consecutive_gap_bound) and GlobalTheory.T.componentImageDegree_between_of_proportionality,
% with delta' = CompHt.delPrime Y.U Z' u_1 = (1 - rho) topDeg Z' u_1 (delPrime_eq). Here delta_Z is
% CompHt.compDeg Y.U Z (wtOf u_1 0): the degree, through Chow forms (ChowHt.imageDeg, IsChowFormInt), of the
% projective variety of dimension D cut out by the relations among the cusp forms of weight u_1 on Z at the generic
% point of Z (compIdeal). Size: about 630,000 lines counted with everything the proof uses and the book does not yet
% cite, all but a few dozen of them also used by (b), (c) or (d); (a) to (d) together use about 810,000.
\openstep{Part (b) of (Y6), the comparison of the ranks $\operatorname{rk}_Z(u_-)$ and
$\operatorname{rk}_Z(u_W+sc_0h)$ with $\beta\delta_Z$, is not yet written out.}
% Lean: ShareAt from GlobalTheory.CovTop.ShareRowsNef.share_alt_alt: nef_share_lower_bound (with
% GlobalTheory.T.twisted_flat_bundle_ranks) for beta delta_Z <= rk_Z(u_-), nef_share_upper_bound for
% (1 - zeta) rk_Z(u_-) <= beta delta_Z, with GlobalTheory.T.birOn_weightPair_both_components_alt and
% GlobalTheory.T.valuation_multiplicity_run_estimate_alt; the bound for u_W + s c_0 h is
% GlobalTheory.JetRank.jet_bundle_rank_le (GlobalTheory.T.jet_rank_complex_le) with
% componentImageDegree_between_of_proportionality and exists_jet_order inside GlobalTheory.CovTop.HodgeSteps.rk_le_alt_alt.
% Size: about 660,000 lines for share_alt_alt and jet_bundle_rank_le together, counted as in (a); about 1,800 of them
% are not used by (a), (c) and (d).
\openstep{Part (c) of (Y6), the arithmetic Hilbert--Samuel lower bound for $\deg_Z(u_-)$ in terms of the heights
$h_Z$, is not yet written out.}
% Lean: HSAt from GlobalTheory.CovTop.ShortRoad.hilbertSamuel_of_component_estimates_alt
% (GlobalTheory.T.twisted_bundle_degree_lower_bound, twisted_flat_bundle_ranks, covering_domain_lattice_regularity,
% GlobalTheory.T.exists_degree_parameters_bound_alt_alt), with A' = 26 (26 D^2 + 2 zeta ell) and
% R' = 4000 b + 2 C_p z^2. Here h_Z is CompHt.compHeight Y.U Z (wtOf u_1 0): the height (ChowHt.imageHeight) of the
% variety of (a) for the lattice of integral cusp forms of weight u_1 on Z (latOn) with their Petersson norms
% (normOn). Size: about 650,000 lines, counted as in (a); about 18,000 of them are not used by (a), (b) and (d).
\openstep{Part (d) of (Y6), the lower bound for the heights $h_Z$ through Theorem~\ref{cv-volume-theorem}, is not
yet written out.}
% Lean: HeightAt from SelfIntersectionVolumeEstimates.height_lower_bound_general_alt =
% GlobalTheory.CovHt.HtTopG.height_of_meanAll_bound, with GlobalTheory.T.covering_short_of_deep_size_flatDeg,
% GlobalTheory.T.covHt_geom_of_siegel_toroidal_semiAbelian, GlobalTheory.T.carrying_section_integral_frame_alt (from
% GlobalTheory.T.refW_of_prints_deepGood_size_flatDeg and GlobalTheory.T.exists_section_bounded_by_bergman_kernel_alt),
% GlobalTheory.T.budget_of_hostN_of_liftWeight_of_size, and DeepStR from the volume theorem (Plug.deepStR_of_volume).
% Size: about 750,000 lines, counted as in (a); about 102,000 of them are not used by (a), (b) and (c).
% Lean: (Y5) has its inputs in Lemmas cv-volume-twisting-points, cv-volume-twisting-cusp,
% cv-volume-twisting-key-section and cv-volume-twisting-unramified, whose remaining open steps are in their
% own proofs. In flat_at_of_split_covering the other hypotheses of flat_root_twist are the complex data,
% GoodAt at every prime and FlatExists of its hypothesis hc, and hY.1.geo.
% Lean: (S1) to (S7). (S1) is GlobalTheory.LamT.exists_good_cmPt_at_frame (= GlobalTheory.T.exists_good_cmPoint_at_cmBasis_frame,
% the CM point of the open step in Lemma cv-cusp-inputs(3); through exists_cmPoint_lattice_data_of_adapted_basis,
% CM.cmExists_holds and exists_good_cm_point_of_integral_cm_object, with the ideal AA of CovData.exists_dataX_good),
% with y_0 = z.y of the good CM point z, together with
% GlobalTheory.T.exists_commonDenom_framePolynomials_cmPoint_generator (LatFull.CMIntZ.exists_coeff_denom_intWt_any).
% (S2) is GlobalTheory.T.detLine_flat_unit_exists = FlatUnit.exists_flat_unit_section: N = 1 from
% FlatAnyLevel.exists_flat_unit_bounded_any_level (FlatAnyLevel.archEx on each component and for each tau), the
% integer multiples of GlobalTheory.T.inverse_of_flat_unit_isIntegral = FlatDenom.unit_pow_clears_denominator_flat_model
% (with CovTop.koecher_principle_unit_constant), the products GlobalTheory.T.exists_boundedUnit_sum_component, and
% the sum over the components. (S3) is CuspMulUnit.pow_mul_mem_cusp_of_unit (= GlobalTheory.T.exists_pow_mul_unit_mem_cusp),
% whose proof uses only that f is a cusp form and that the product is a section of the dominant weight:
% NonDomVanish.eq_zero_of_lt (the weight of the second factor does not increase), FlatTw.Hol.isSec_modelWt_rdU with
% the key section HSApp.VanishingBound.keySec (its readings are holomorphic), CuspFrom.isModerate, fj_mul_eq_zero
% and FJZero.mem_SS_of_isModerate. (S4) is LargeCover.exists_smul_nonneg_of_dvd: a pinned DeltaDat
% (LevelCover.exists_deltaDat_isLev, exists_phiPin_of_character), the short generators shortGen with
% holSub0_le_span_shortGen_upto and shortValTot_of_pinned, and exists_smul_nonneg_of_shortRule_any_level. (S5) is
% GlobalTheory.T.torFlag_nonempty_of_one_le_dimFlag (after covering_topology_refinement),
% GlobalTheory.T.flagSpace_component_smooth_projective_connected (IsVariety0: smooth of relative dimension dimFlag,
% projective, geometrically connected), CompHS.cuspSec_injective, GlobalTheory.T.cuspForms_extend_vanish_compactifiedFlagSpace
% (Extends) and vanBd_le_space; (S6) is GlobalTheory.T.hodge_bundle_base_point_free_on_open
% (HSApp.globally_generated_on_open_of_weight_gap_ge, FreeOnOpen for the space of sections from cusp forms). Their
% published inputs are PublishedHodgeBundleTheorems, GrothendieckDeformation, HSCore.Quoted, SiegelToroidalIntegral,
% ToroidalChartsForFan and SemiAbelianLieSeminegative, and for (S6) also NadelVanishingAndNonKltDivisor and
% Principalisation. (S7) is GlobalTheory.T.short_span_composition_closed = ShortSpanComp.short_span_comp_mem: the theta
% series of exists_good_cover_at_minimal_level, which is (Y4), at the dominant weight of
% exists_thetaWeight_exponents_of_large_gap, multiplied by the twisting sections of
% exists_keyUnit_pow_twistingSection_on_component and exists_short_integral_flatLine_twistingSection through
% exists_spanning_family_mul_twistingSection, with log_sum_bad_primes_bound_of_covering and
% covolume_numerical_bounds_collect. (Y3) (c) is GlobalTheory.T.orbits_separated_by_regular_cusp for S = grpHR fr,
% the same declaration and group as u_0 in the proof of Proposition cv-point-conditions, whose remaining input
% OrbitSep.separate_orbits_by_cusp is Lemma cv-point-conditions-separation. (Y4) and (Y6) (a) to (d) are as in the comments above.

Proof of (1). Let $H$, $\chi_H$ and $v$ be as in (Y2) and (Y3), and for $F\in R$ put
$T(F)=\sum_{g\in H}\chi_H(g)^{-1}g^*F$. Take $n'=|H|$ and, for $h\in H$ and $f\in S_Y(\kappa^-)$, put
$\varphi_h(f)=T(f\cdot h^*v)$. Since $\kappa^-+c_0h=\kappa_j$, (Y1) shows that $f\cdot h^*v$ and its
translates lie in $S_Y(\kappa_j)$, so $\varphi_h$ is an $F_Y$-linear map $S_Y(\kappa^-)\to S_Y(\kappa_j)$.
The values of $\chi_H$ are roots of unity, as $H$ is finite; so each $\chi_H(g)^{-1}$ is integral and
$|\sigma(\chi_H(g))|=1$ for every $\sigma$. Lemma~\ref{cv-isotypic-twisted-trace} applies with $K=F_Y$,
the action of $H$ on $R$ by pullbacks, $\chi=\chi_H$ and $v$.
\begin{enumerate}
\item[(i)] Let $f\in A$. By (Y1), $f\cdot h^*v$ and its translates are integral, so $\varphi_h(f)$ is
  integral. By Lemma~\ref{cv-isotypic-twisted-trace}(1), $g^*\varphi_h(f)=\chi_H(g)\varphi_h(f)$ for
  $g\in H$. By (Y2), $\varphi_h(f)\in B$.
\item[(ii)] The maps $\varphi_h$ are jointly injective by Lemma~\ref{cv-isotypic-twisted-trace}(2).
\item[(iii)] For every $\sigma$, the triangle inequality for $N_\sigma$, the invariance of $N_\sigma$
  under translations (Y1) and (Y3) give
  $N_\sigma(\varphi_h(f))\le\sum_gN_\sigma\big(g^*(f\cdot h^*v)\big)=|H|N_\sigma(f\cdot h^*v)\le
  |H|e^{C_Yz^9}N_\sigma(f)$. Summing the squares over $h$ gives
  $\sum_hN_\sigma(\varphi_h(f))^2\le|H|^3e^{2C_Yz^9}N_\sigma(f)^2$. By (Y2) and the size conditions
  $2C_Y+200\le z$ and $z^{12}\le b$, $3\log|H|+2C_Yz^9\le(135+2C_Y)z^9\le z^{10}\le b\le2b$. Hence
  $\sum_hN_\sigma(\varphi_h(f))^2\le e^{2b}N_\sigma(f)^2$.
\item[(iv)] Apply Lemma~\ref{cv-isotypic-twisted-trace}(3) with $V=S_Y(\kappa_j)$ and
  $V_+=S_Y(\kappa^+)$; its hypotheses hold by (Y1), as $\kappa_j+c_0h=\kappa^+$. By (Y2), $B\subset V^\chi$,
  so $\operatorname{rk}B\le\dim V^\chi$, and $\dim S_Y(\kappa^+)=\operatorname{rk}(\kappa^+)$ by (Y1). Hence
  $|H|\operatorname{rk}B\le\operatorname{rk}(\kappa^+)$.
\end{enumerate}

Proof of (2). Let $s\in\{1,2\}$. For every component $Z$, (Y6)(b) gives
$\operatorname{rk}_Z(u_W+sc_0h)\le(1+\zeta)\beta\delta_Z\le(1+\zeta)\operatorname{rk}_Z(u_-)$. Summing over
$Z$ and using (Y5) for $s$ and for $-1$,
$\operatorname{rk}(\kappa_j+sc_0h)\le(1+\zeta)\operatorname{rk}(\kappa^-)=(1+\zeta)\operatorname{rk}A$. As
$\zeta\le\frac12$ and $\zeta\le b/K_j$ by (Y6),
$\operatorname{rk}(\kappa_j+sc_0h)\le\big(1+\min(\frac12,b/K_j)\big)\operatorname{rk}A$. The case $s=1$
is (2).

Proof of (3). Let $s_0$ be the set of forms of (Y4), let $T$ be as in the proof of (1), and put
$s=\{T(y):y\in s_0\}$. For $y\in s_0$, $T(y)$ is integral by (Y1), as the $\chi_H(g)^{-1}$ are integral,
and $g^*T(y)=\chi_H(g)T(y)$ for $g\in H$ by Lemma~\ref{cv-isotypic-twisted-trace}(1); so $T(y)\in B$ by
(Y2). As in (iii) of the proof of (1), $N_\sigma(T(y))\le|H|\exp\big(C_Yz^2\operatorname{size}(\kappa_j)\big)$.
Every $x\in B$ lies in the $F_Y$-span of $s$: write $x=\sum_ka_ky_k$ with $a_k\in F_Y$ and $y_k\in s_0$;
since $g^*x=\chi_H(g)x$ for $g\in H$, $T(x)=|H|x$, so $x=|H|^{-1}\sum_ka_kT(y_k)$. Finally, by (Y2), (Y4) and the size conditions
$C_Y+45\le z$ and $z^{12}\le PP$,
\[
 \log|H|+C_Yz^2\operatorname{size}(\kappa_j)\le(45+C_Y)z^2\operatorname{size}(\kappa_j)
 \le z^3\operatorname{size}(\kappa_j)\le PP\cdot\operatorname{size}(\kappa_j)=K_j,
\]
so the elements of $s$ have norm at most $e^{K_j}$ at every embedding.

Proof of (4). Let $\Theta$ be as in (Y5). By (Y5) and (Y1),
$\operatorname{rk}A=\operatorname{rk}(\kappa^-)=\operatorname{rk}(u_-)>0$, and dividing
$\deg(u_-)+\operatorname{rk}(u_-)\Theta\le\deg(\kappa^-)$ by $\frac12\operatorname{rk}(u_-)$ gives
$\mu_0+2\Theta\le\mu(A)$. It remains to show $2\Theta\ge-\mathrm{main}_j-(11C_A+83)u$. By
Theorem~\ref{cv-flat-degree} every flat section of type $\tau$ has degree
$\frac{p-q}2\log\Omega_E(\chi)+q(\log d-\log\operatorname{rd}_E)$, so
$c_\tau+\frac{p-q}2\log\Omega_E(\chi)=-q(\log d-\log\operatorname{rd}_E)$, which has absolute value at
most $3qx^\varepsilon$ by the size conditions. Hence
$a_\tau c_\tau\ge-\frac12(p-q)a_\tau\log\Omega_E(\chi)-3qx^\varepsilon|a_\tau|$ for each $\tau$, and the
anchor conditions $\big|\mathrm{main}_j-\sum_\tau(p-q)a_\tau\log\Omega_E(\chi)\big|\le C_Au$ and
$q\,x^\varepsilon\sum_\tau|a_\tau|\le C_Au$ give
\[
 2\sum_\tau a_\tau c_\tau\ge-\mathrm{main}_j-C_Au-6C_Au.
\]
The size conditions give $10\le x^\varepsilon\le u$, $\log\operatorname{rd}_E\le2x^\varepsilon$ and
$fr\,n\,x^\varepsilon\le b\le u$. So $2q\log\operatorname{rd}_E\sum_\tau|a_\tau|\le4C_Au$; as
$\frac12\log2\pi\le x^\varepsilon$,
$2\cdot\frac n2fr\big(\frac12\log2\pi+\log\operatorname{rd}_E\big)\le3fr\,n\,x^\varepsilon\le3u$; and
$2\lambda_Y\le80x^\varepsilon\le80u$ by (Y5). Adding up, $2\Theta\ge-\mathrm{main}_j-(11C_A+83)u$, which
is (4) with $C_F=11C_A+83$.

Proof of (5). By (Y5) and (Y1), $\operatorname{rk}(u_-)=\operatorname{rk}(\kappa^-)>0$. By (Y6)(e) and
$\deg(u_-)=\frac12\operatorname{rk}(u_-)\mu_0$,
\[
 2\,\frac{\sum_Z\deg_Z(u_-)}{\sum_Z\operatorname{rk}_Z(u_-)}\le\mu_0.
\]
Apply Lemma~\ref{cv-height-components-numerics} with $I$ the set of components of $Y$, finite and
nonempty by (Y1), $E=u$, which is at least $1$ because $1\le b\le u$, the numbers and functions
$\ell$, $D$, $\zeta$, $\rho$, $\beta$, $\delta'$, $\delta$, $h$, $s$, $c_s$, $c_c$ and $X$ of (Y6), $A=A'$,
$R=R'$, $r_Z=\operatorname{rk}_Z(u_-)$, $g_Z=\deg_Z(u_-)$ and $\mu=\mu_0$. Its hypotheses are (Y6)(a) to
(d) and the display. It gives $\mu_0\ge-2(3C_1+C_2+4)u$, which is (5) with $C_U=2(3C_1+C_2+4)$.

Proof of (6). Let $C$ be the $F_c$-space of cusp forms of weight $\kappa_j$ and the prescribed type on
the Shimura variety of $U(W_j)$. Then $\Lambda_j^{\mathrm{aut}}$ is the set of integral forms in $C$,
and $B$ is the set of integral forms in the $F_Y$-span of $\pi(C)$, for the pullback $\pi$ of (Y7).
Apply Lemma~\ref{cv-lattices-different-slope} with $K=F_c$, $K'=F_Y$, $\iota=\pi$,
$L=\Lambda_j^{\mathrm{aut}}$ and $M=B$. The fields are totally complex, as they contain a CM field;
the Petersson norms satisfy $N_{\bar\sigma}=N_\sigma$; $\pi$ has the required properties by (Y7); $B$ is
finitely generated by (Y1); and $B\ne0$, because $A\ne0$ by (Y1) and the maps of (1) are jointly
injective with $\varphi_h(A)\subset B$. By (Y7), $\pi$ maps $\Lambda_j^{\mathrm{aut}}$ into $B$. For the
last hypothesis let $g\in B$, and write $g=\sum_{i\le m}x_i\pi(b_i)$ with $b_i\in C$ and $x\in F_Y^m$.
Let $M'\subset F_Y^m$ be the set of $y$ for which $\sum_iy_i\pi(b_i)$ is integral. It is an
$\mathcal O_{F_Y}$-submodule containing $x$, and it is stable under $\operatorname{Gal}(F_Y/F_c)$ by
(Y7). By Lemma~\ref{cv-lattices-galois-different}, for $\delta$ in the different $\mathfrak D_{F_Y}$ of
$F_Y$, the vector $\delta x$ is an $\mathcal O_{F_Y}$-linear combination of elements $y\in M'\cap F_c^m$.
For such $y$, $\sum_iy_i\pi(b_i)=\pi\big(\sum_iy_ib_i\big)$, and $\sum_iy_ib_i\in C$ is integral by
(Y7), so it lies in $\Lambda_j^{\mathrm{aut}}$. Hence $\delta g$ lies in the $\mathcal O_{F_Y}$-span of
$\pi(\Lambda_j^{\mathrm{aut}})$. The lemma gives
$\mu(B)-2\log|d_{F_Y}|/[F_Y:\mathbf Q]\le\mu(\Lambda_j^{\mathrm{aut}})=\mu_j^{\mathrm{aut}}$, and
$\log|d_{F_Y}|/[F_Y:\mathbf Q]\le b$ gives (6).
% Lean: (1) GlobalTheory.CovTop.DescentMaps.mapsCore, with GlobalTheory.LevDesc.cuspPh_mem_chiLat (i),
% GlobalTheory.levPh_injective_of_levDet_nonZeroDivisor (ii), GlobalTheory.sum_sq_normW_cuspPh_le and the bound hb
% of covering_descent_maps_compatible_alt (iii), GlobalTheory.CovolRank.card_mul_rkK_chiLat_le_rankW (iv);
% (2) GlobalTheory.CovTop.HodgeSteps.rk_le_alt_alt with s = 1 (covering_rank_of_descent_steps_alt);
% (3) GlobalTheory.CovGminTop.gminAt_of_gens through exists_short_eigen and num_step; (4)
% GlobalTheory.FlatSplit.flat_at_of_split_covering and the numerical lemma GlobalTheory.FlatSplit.twice_twist_ge;
% (5) GlobalTheory.CovTop.FromComp.atPoint; (6) GlobalTheory.CoreRing.slope_sq_latOver_sub_le, whose step hDiff
% is span_integral_descent.

Conclusion. Apply Lemma~\ref{cv-slope-descent} with $P=b$, $\Lambda=K_j$, $\varepsilon=b/K_j$ and
$t=-\mathrm{main}_j-(C_F+C_U)u-2b$. Its hypotheses are (1), (2) and (3), and $t\le\mu(A)-2b$ by (4)
and (5). Statement (4) of Proposition~\ref{cv-point-conditions} gives $b\le u$, $b\le K_j$ and
$|\mathrm{main}_j|\le K_j$. Hence
$\varepsilon(|t|+2\Lambda)\le|\mathrm{main}_j|\cdot b/K_j+(C_F+C_U)u+2b+2b\le(C_F+C_U)u+5u$,
and so $\mu(B)\ge-\mathrm{main}_j-(2(C_F+C_U)+7)u$. With (6) this gives
$\mu_j^{\mathrm{aut}}\ge-\mathrm{main}_j-(2(C_F+C_U)+9)u$. Take $C_2$ the larger of this constant and
the constant of the case $j=m$.
\end{proof}
\leanref{Development/AbelianVarietyTheory/Definitions/Divisors/CovolXOfCovol.lean}{GlobalTheory.CovTop.MapsAt}
\leanref{Development/AbelianVarietyTheory/Definitions/Divisors/CovolXOfCovol.lean}{GlobalTheory.CovTop.RkAt}
\leanref{Development/AbelianVarietyTheory/Definitions/Divisors/CovolXOfCovol.lean}{GlobalTheory.CovTop.GminAt}
\leanref{Development/AbelianVarietyTheory/Definitions/Divisors/CovolXOfCovol.lean}{GlobalTheory.CovTop.FlatAt}
\leanref{Development/AbelianVarietyTheory/Definitions/Divisors/CovolXOfCovol.lean}{GlobalTheory.CovTop.UAt}
\leanref{Development/AbelianVarietyTheory/Definitions/Divisors/CovolXOfCovol.lean}{GlobalTheory.CovTop.RingAt}
\leanref{Development/AbelianVarietyTheory/Proofs/DivisorProofs/MapsBoundOfDeepGood.lean}{GlobalTheory.CovTop.DescentMaps.mapsCore}
\leanref{Development/AutomorphicArithmetic/ThetaCorrespondence/ThetaPairLifts/HodgeBundleWeights/FlatAtOfSplitCovering.lean}{GlobalTheory.FlatSplit.flat_at_of_split_covering}
\leanref{Development/AbelianVarietyTheory/Proofs/DivisorProofs/UniformLowerBoundOfSums.lean}{GlobalTheory.CovTop.FromComp.atPoint}
\leanref{Development/AutomorphicArithmetic/ThetaCorrespondence/ThetaPairLifts/ThirdPlaceSeries/RingOfCoveringDescent.lean}{GlobalTheory.CovTop.RingStep.ring_of_covering_descent_alt}
\leanref{Development/AbelianVarietyTheory/Proofs/DivisorProofs/TwoDefiniteBoundGeAnchoredMainTermSubCrudeErrorsStatement.lean}{GlobalTheory.T.two_definiteBound_ge_anchoredMainTerm_sub_crude_errors}
\leanref{Development/AutomorphicSpectrum/Definitions/GlobalGeometry/ModelSmall.lean}{GlobalTheory.CompanionTheory.Ar.Plug.DeepStR}
\leanref{Development/AbelianVarietyTheory/Proofs/DivisorProofs/TwoDefiniteBoundGeAnchoredMainTermSubCrudeErrors.lean}{GlobalTheory.DefArith.two_defBound_ge}

\begin{theorem}\label{cv-covolume-bound-from-literature}
Assume the following results from the literature (Chapter~\ref{chapter-literature}):
resolution of singularities and principalization in
characteristic zero (Theorem~\ref{cv-functorial-resolution}); the Poincar\'e--Lelong
formula; the upper bound in the arithmetic Hilbert--Samuel formula; the bigness inequality
of Theorem~\ref{cv-moriwaki-bigness}; exactness of the Hodge--Tate sequence for
$p$-divisible groups; Zarhin's trick; asymptotic Riemann--Roch for nef divisors; the
integral toroidal compactifications of Siegel and of unitary PEL moduli spaces;
seminegativity of the Lie algebra of semi-abelian schemes; separation of pairs of points by
sections; the Arakelov degree formula; the published results of arithmetic geometry
collected for the Hodge bundle; and Grothendieck--Messing deformation theory. Then the
conclusion of Theorem~\ref{cv-covolume-bound} holds.
\end{theorem}
\leanref{Development/AutomorphicArithmetic/ThetaCorrespondence/ThetaPairLifts/ThirdPlaceSeries/TopCovolSlim2Use.lean}{SelfIntersectionVolumeEstimates.topCovolSlim2_use_alt_alt}
\begin{proof}
Let $\varepsilon$, $\mathcal F$ and $C_1,C_2,A,B,x_0$ be as in
Proposition~\ref{cv-point-conditions}, whose proof uses the listed results through
Theorems~\ref{cv-flat-degree} and~\ref{cv-volume-theorem}. Put
$\delta=e^{-x^{\varepsilon/(7J)}}$, where $\varepsilon/J$ is the same at every point of
$\mathcal F$. Enlarge $x_0$ so that
$12\delta\,(x^A+x^B+C_1+C_2+1)\le1$ for $x\ge x_0$, and put $C_0=C_1+C_2+1$.

Fix a point $P$ at a scale $x\ge x_0$ and an index $j$, and put
$Q=(x^A+x^B+C_1+C_2+1)u$. Apply Lemma~\ref{cv-cut} with $L=\Lambda_j$,
$M=\Lambda_j^{\mathrm{all}}$, $e=e_j$, $g=|G|$ and this $\delta$ and $Q$. Its hypotheses
hold: $\delta\le1/12$; the rank condition is
Proposition~\ref{cv-point-conditions}(3); $\mu_{\max}(M)\le x^Ab\le x^Au\le Q$ and
$2\log g\le b\le Q$ by (2) and (4); and by (5), (6) and (4),
$\mu(\Lambda_j)\ge-\mathrm{main}_j-(C_1+C_2)u\ge-x^Bb-(C_1+C_2)u\ge-Q$. So
$\mu(\Lambda_j^{\mathrm{sep}})\ge\mu(\Lambda_j)-12\delta Q\ge\mu(\Lambda_j)-u$. With (1),
(5) and (6),
\[
 \frac1D\operatorname{lcv}_{X_j}(M_j)\le-\mu(\Lambda_j^{\mathrm{sep}})
 \le-\mu_j^{\mathrm{aut}}+C_1u+u\le\mathrm{main}_j+(C_1+C_2+1)u.
\]
Since $u\ge0$, the right side is at most $\mathrm{main}_j+Cu$ for every $C\ge C_0$.
\end{proof}
\leanref{Development/AbelianVarietyTheory/Proofs/DivisorProofs/CovolumeFromGramAndBounds.lean}{GlobalTheory.CovTop.Assembly.lcv_le}
\leanref{Development/AutomorphicArithmetic/ThetaCorrespondence/ThetaPairLifts/ThirdPlaceSeries/CovolXOfSteps.lean}{GlobalTheory.CovTop.Assembly.covolX_of_lift_bounds_and_coreLe}

\begin{remark}\label{cv-bound-as-used}
The hypotheses of Theorem~\ref{cv-covolume-bound} are met at the points where the bound is
used. On the definite side the bound is applied at the points of the selected summands of
the families of Chapter~\ref{chapter-assembly}, and at those points the point conditions
and the coefficient field bound are proved before the bound is invoked. So the bound is
used at points satisfying its hypotheses, and the remaining question, that such points
exist at all large scales, is the existence of the families
(Chapter~\ref{chapter-host-fields} and Chapter~\ref{chapter-assembly}).

The inequality itself is not empty at those points. The point conditions give $d_j=D$, so
$M_j$ is a lattice of full rank $D$ in $L^D$; its realization is discrete, so
$\frac1D\operatorname{lcv}_{X_j}(M_j)=-\mu(M_j)$
is a finite number that can be arbitrarily large for a sparse lattice; and the right side
is finite, since $|\mathrm{main}_j|\le x^Bb$ and $u\le x^Bb$. The main term
$\mathrm{main}_j$ is the term that matters, and the downstream use keeps it: the definite
side uses only the sum over $j$, in the form
$\sum_j\frac1D\operatorname{lcv}_{X_j}(M_j)\le\sum_j\mathrm{main}_j+mCu$, and absorbs
$mCu$ into its error budget, of size a constant times $m\,x^{\varepsilon}(fmbx^{\varepsilon}+N_c)$
for a constant $N_c$ of the family.
\end{remark}
\leanref{Development/AutomorphicArithmetic/ThetaCorrespondence/ThetaPairLifts/ThirdPlaceSeries/LinkedBundlesL3OfOpen.lean}{GlobalTheory.Glob.CovolRow.linked_bundles_l3_of_open_alt}
\leanref{Development/AutomorphicArithmetic/ThetaCorrespondence/ThetaPairLifts/ThirdPlaceSeries/EventualL3OfCovolume.lean}{GlobalTheory.Fc.Ind.eventual_L3_of_covolume_alt}
\leanref{Formalism/LocalToGlobal/Divisibility/DivisionHeightBoundStatement.lean}{Formalism.L2.DIV.l3main_of_scale}
